%% parts/p01/symbol-registry.tex --- scoped notation and lookup tables.
\chapter{Symbol Registry, Binding, and Precedence}
\chaptermark{Symbol Registry}
\label{sec:p01-symbol-registry}

The registry is normative for this handbook, not a claim that the literature
uses one notation.  Each row fixes a glyph, its formal or grammatical
standing, its scope, and (when it is an operator) its precedence.  Normally a
glyph receives another role only in a disjoint scope declared before use.  A
named primary calculus, logic, model, or method instead keeps its published
syntax: every overlapping local role must be typed and registered honestly,
not renamed to fit handbook notation.
\index{symbol registry}\index{notation!scope}

The editorial defaults are vertical bars for set cardinality,
\(\dom\mathcal M\) for a structure domain,
\(\mathcal M\equiv\mathcal N\) for elementary equivalence,
\(\varphi\Leftrightarrow\psi\) for formula equivalence,
\(\defeq\) for a definitional sign, and \(\forces\), with negation
\(\nforces\), for the declared forcing relation.  A source presentation's
explicitly typed local notation governs that presentation.  For example,
the inherited propositional semantics uses \(\varphi\equiv\psi\) for
equality of truth values under every valuation, whereas the displayed
structure operands in \(\mathcal M\equiv\mathcal N\) give the
elementary-equivalence role.  The register records both roles; neither
declaration renames the other presentation.
The inherited first-order semantics of Part~2 (\cref{sec:fol-semantics-01})
and the inherited cardinality results that follow it keep their own bar
notation for domains.  There the operand fixes the role: bars around a
structure denote its domain, so that \(\lvert\mathcal M\rvert=\dom\mathcal M\),
while bars around a set denote the cardinality of that set.  In the
term \(\abs{\lvert\mathcal N\rvert}\) of the downward L\"owenheim--Skolem bound
(\cref{thm:downward-ls}) the inner bars therefore give the domain of
\(\mathcal N\) and the outer bars its cardinality; in the editorial
notation the same cardinal is written \(\abs{\dom\mathcal N}\).  New
editorial text writes structure domains only as \(\dom\mathcal M\).

In the inherited CPL semantics the truth-value set is
\(\mathbb B=\{\mathbf F,\mathbf T\}\).
For numerical Boolean tables, the coding sends \(\mathbf F\) to \(0\)
and \(\mathbf T\) to \(1\); a valuation written with numerical values
uses this coding.  The coding preserves the declared truth functions and
does not identify the two sets literally.  The Boolean-table section uses
the numerical presentation; the source CPL valuation and substitution
sections retain the labelled values.

\section{Generated master register}
\label{sec:p01-generated-symbols}

The following scope tables are generated from
\texttt{registry/symbols.csv}.  Their role and scope columns are part of each
notation's declaration; the prose notes record local restrictions and
typed collision boundaries across model theory, computability and arithmetic,
set theory, and nonclassical logic.

% GENERATED by tools/symbols2tex.py from handbook/registry/symbols.csv -- do not hand-edit.
\begin{landscape}
\Needspace{.4\textheight}
\subsection*{Global scope}
\begingroup\scriptsize\setlength{\tabcolsep}{3pt}
\begin{longtable}{@{}p{0.9in}p{1.3in}p{1.0in}p{0.5in}p{0.3in}p{0.4in}p{0.35in}p{1.65in}@{}}
\toprule
\textbf{Glyph} & \textbf{Macro} & \textbf{Role} & \textbf{Scope} & \textbf{Prec./form} & \textbf{First} & \textbf{Reg. page} & \textbf{Principal uses / notes} \\
\midrule
\endfirsthead
\toprule
\textbf{Glyph} & \textbf{Macro} & \textbf{Role} & \textbf{Scope} & \textbf{Prec./form} & \textbf{First} & \textbf{Reg. page} & \textbf{Principal uses / notes} \\
\midrule
\endhead
\midrule\multicolumn{8}{r}{\emph{continued}}\\\endfoot
\bottomrule\endlastfoot
\phantomsection\label{entry:symbol-registry-00002}\(\mathcal L\) & \texttt{\detokenize{\Lang}} & signature & global & n/a & \hyperref[sec:prop-deduction-01]{\ref*{sec:prop-deduction-01}} & \thepage & Nonlogical symbols with declared arities or sorts; governs text after a language declaration. \\
\phantomsection\label{entry:symbol-registry-00003}\(x,y,z\) & \texttt{\detokenize{x,}\allowbreak\detokenize{y,}\allowbreak\detokenize{z}} & object-language variables & global & n/a & \hyperref[sec:prop-deduction-01]{\ref*{sec:prop-deduction-01}} & \thepage & Variables occurring in terms and formulas. \\
\phantomsection\label{entry:symbol-registry-00004}\(r,t\) & \texttt{\detokenize{r,}\allowbreak\detokenize{t}} & term metavariables & global & n/a & \hyperref[sec:prop-deduction-01]{\ref*{sec:prop-deduction-01}} & \thepage & Members of \(\mathrm{Term}(\mathcal L)\). \\
\phantomsection\label{entry:symbol-registry-00005}\(\varphi,\psi,\chi\) & \texttt{\detokenize{\varphi,}\allowbreak\detokenize{\psi,}\allowbreak\detokenize{\chi}} & formula metavariables & global & n/a & \hyperref[sec:prop-syntax-semantics-01]{\ref*{sec:prop-syntax-semantics-01}} & \thepage & Members of the active formula set.  A source-bound presentation may instead declare its own formula metavariables locally; the source declaration governs that presentation. \\
\phantomsection\label{entry:symbol-registry-00006}\(\neg\varphi\) & \texttt{\detokenize{\neg}} & object-language negation & global & 3 & \hyperref[sec:prop-syntax-semantics-01]{\ref*{sec:prop-syntax-semantics-01}} & \thepage & Primitive unary connective; prefix and governs the immediately following formula. \\
\phantomsection\label{entry:symbol-registry-00007}\(\varphi\land\psi\) & \texttt{\detokenize{\land}} & object-language conjunction & global & 4 & \hyperref[sec:prop-syntax-semantics-01]{\ref*{sec:prop-syntax-semantics-01}} & \thepage & Primitive or explicitly defined binary connective in the active language; left associative in handbook prose. \\
\phantomsection\label{entry:symbol-registry-00008}\(\varphi\lor\psi\) & \texttt{\detokenize{\lor}} & object-language disjunction & global & 5 & \hyperref[sec:prop-syntax-semantics-01]{\ref*{sec:prop-syntax-semantics-01}} & \thepage & Primitive or explicitly defined binary connective in the active language; left associative in handbook prose. \\
\phantomsection\label{entry:symbol-registry-00009}\(\varphi\to\psi\) & \texttt{\detokenize{\to}} & object-language implication & global & 6 & \hyperref[sec:prop-syntax-semantics-01]{\ref*{sec:prop-syntax-semantics-01}} & \thepage & Primitive or explicitly defined binary connective in the active language; right associative in handbook prose. \\
\phantomsection\label{entry:symbol-registry-00010}\(\sigma,\tau\) & \texttt{\detokenize{\sigma,}\allowbreak\detokenize{\tau}} & sentence metavariables & global & n/a & \hyperref[sec:prop-deduction-01]{\ref*{sec:prop-deduction-01}} & \thepage & Members of \(\mathrm{Sent}(\mathcal L)\). \\
\phantomsection\label{entry:symbol-registry-00011}\(\forall x\,\varphi\) & \texttt{\detokenize{\forall}} & universal quantifier & global & 3 & \hyperref[sec:prop-deduction-01]{\ref*{sec:prop-deduction-01}} & \thepage & Binds occurrences of x in its declared scope; prefix at the unary rank. \\
\phantomsection\label{entry:symbol-registry-00012}\(\exists x\,\varphi\) & \texttt{\detokenize{\exists}} & existential quantifier & global & 3 & \hyperref[sec:prop-deduction-01]{\ref*{sec:prop-deduction-01}} & \thepage & Binds occurrences of x in its declared scope; prefix at the unary rank. \\
\phantomsection\label{entry:symbol-registry-00014}\(\Gamma,\Delta,\Sigma,\Pi\) & \texttt{\detokenize{\Gamma,}\allowbreak\detokenize{\Delta,}\allowbreak\detokenize{\Sigma,}\allowbreak\detokenize{\Pi}} & formula contexts & global & n/a & \hyperref[sec:prop-syntax-semantics-01]{\ref*{sec:prop-syntax-semantics-01}} & \thepage & Finite sets, multisets, or sequences exactly as declared by the active calculus.  Source-bound locally typed uses such as the Buchi alphabet Sigma are recorded separately and are not contexts. \\
\phantomsection\label{entry:symbol-registry-00021}\(T,U\) & \texttt{\detokenize{T,}\allowbreak\detokenize{U}} & theories & global & n/a & \hyperref[sec:prop-deduction-01]{\ref*{sec:prop-deduction-01}} & \thepage & Sets of sentences in the fixed signature. \\
\phantomsection\label{entry:symbol-registry-00022}\(\mathcal M,\mathcal N\) & \texttt{\detokenize{\mathcal M,}\allowbreak\detokenize{\mathcal N}} & structures & global & n/a & \hyperref[sec:prop-deduction-01]{\ref*{sec:prop-deduction-01}} & \thepage & Structures for the active signature with nonempty domains unless locally declared otherwise; in exact Part V, \(\mathcal N\) is the standard structure for the fixed arithmetic signature. \\
\phantomsection\label{entry:symbol-registry-00023}\(\operatorname{dom}X\) & \texttt{\detokenize{\dom}} & domain operator & global & n/a & \hyperref[sec:prop-deduction-01]{\ref*{sec:prop-deduction-01}} & \thepage & Editorial domain notation; operand type fixes structure, function or name domains. A source presentation may retain its own explicitly typed domain notation. \\
\phantomsection\label{entry:symbol-registry-00024}\(v\) & \texttt{\detokenize{v}} & valuation & global & n/a & \hyperref[sec:prop-syntax-semantics-01]{\ref*{sec:prop-syntax-semantics-01}} & \thepage & Map from propositional variables to the active truth-value set in propositional and finite-matrix semantics; its homomorphic extension \(\widehat v\) is fixed by the active semantics. \\
\phantomsection\label{entry:symbol-registry-00025}\(s\) & \texttt{\detokenize{s}} & variable assignment & global & n/a & \hyperref[sec:prop-deduction-01]{\ref*{sec:prop-deduction-01}} & \thepage & Map from first-order variables to \(\operatorname{dom}\mathcal M\). \\
\phantomsection\label{entry:symbol-registry-00026}\(\mathcal M,s\models\varphi\) & \texttt{\detokenize{\entails}} & satisfaction & global & meta & \hyperref[sec:prop-deduction-01]{\ref*{sec:prop-deduction-01}} & \thepage & \(\varphi\) is satisfied by \(\mathcal M\) under \(s\). \\
\phantomsection\label{entry:symbol-registry-00027}\(\mathcal M\models\sigma\) & \texttt{\detokenize{\entails}} & truth in a structure & global & meta & \hyperref[sec:prop-deduction-01]{\ref*{sec:prop-deduction-01}} & \thepage & Sentence \(\sigma\) is true in \(\mathcal M\). \\
\phantomsection\label{entry:symbol-registry-00028}\(\Gamma\models\varphi\) & \texttt{\detokenize{\entails}} & semantic consequence & global & meta & \hyperref[sec:prop-syntax-semantics-01]{\ref*{sec:prop-syntax-semantics-01}} & \thepage & Every interpretation and assignment in the declared class satisfying \(\Gamma\) satisfies \(\varphi\). \\
\phantomsection\label{entry:symbol-registry-00029}\(\Gamma\vdash_{S}\varphi\) & \texttt{\detokenize{\derives_{S}}} & derivability (per-system turnstile) & global & meta & \hyperref[sec:prop-syntax-semantics-01]{\ref*{sec:prop-syntax-semantics-01}} & \thepage & A derivation exists in the named calculus \(S\); the subscript is mandatory outside a declared local abbreviation. \\
\phantomsection\label{entry:symbol-registry-00030}\(\varphi\Leftrightarrow\psi\) & \texttt{\detokenize{\Leftrightarrow}} & editorial formula equivalence & global & meta & \hyperref[sec:prop-syntax-semantics-01]{\ref*{sec:prop-syntax-semantics-01}} & \thepage & Metalinguistic equivalence in presentations declaring this notation. Source-local formula equivalence written with \(\equiv\) is retained and registered separately. \\
\phantomsection\label{entry:symbol-registry-00031}\(\varphi\leftrightarrow\psi\) & \texttt{\detokenize{\leftrightarrow}} & object-language biconditional & global & 7 & \hyperref[sec:prop-syntax-semantics-01]{\ref*{sec:prop-syntax-semantics-01}} & \thepage & Object-language constructor; rank 7 of the extended precedence ladder. \\
\phantomsection\label{entry:symbol-registry-00032}\(a\mathrel{:\equiv}b\) & \texttt{\detokenize{\defeq}} & definitional abbreviation & global & meta & \hyperref[sec:levels-01]{\ref*{sec:levels-01}} & \thepage & The left expression is defined by the right expression; this sign is metalinguistic. \\
\phantomsection\label{entry:symbol-registry-00033}\(\varphi[t/x]\) & \texttt{\detokenize{[t/}\allowbreak\detokenize{x]}} & capture-avoiding substitution & global & n/a & \hyperref[sec:prop-deduction-01]{\ref*{sec:prop-deduction-01}} & \thepage & Substitution of \(t\) for free occurrences of \(x\) in \(\varphi\). \\
\phantomsection\label{entry:symbol-registry-00034}\(\operatorname{FV}(\varphi),\operatorname{FV}(\Gamma),\operatorname{FV}(M),\operatorname{FV}(\alpha)\) & \texttt{\detokenize{\FV}} & free-variable operation & global & n/a & \hyperref[sec:prop-deduction-01]{\ref*{sec:prop-deduction-01}} & \thepage & Variables with a free occurrence in the locally declared formula, context, term, or type syntax; context action takes the union over members. \\
\phantomsection\label{entry:symbol-registry-00035}\(\operatorname{Mod}(T)\) & \texttt{\detokenize{\Mod}} & model class & global & n/a & \hyperref[sec:prop-deduction-01]{\ref*{sec:prop-deduction-01}} & \thepage & Structures satisfying every sentence of \(T\). \\
\phantomsection\label{entry:symbol-registry-00036}\(\operatorname{Th}(\mathcal M)\) & \texttt{\detokenize{\Th}} & theory of a structure & global & n/a & \hyperref[sec:prop-deduction-01]{\ref*{sec:prop-deduction-01}} & \thepage & Sentences true in \(\mathcal M\). \\
\phantomsection\label{entry:symbol-registry-00037}\(X^*\) & \texttt{\detokenize{{}^*}} & finite iteration or closure & global & n/a & \hyperref[sec:strings-01]{\ref*{sec:strings-01}} & \thepage & For an alphabet this is finite-word Kleene star; for a relation it is reflexive-transitive closure. Operand type fixes the reading. \\
\phantomsection\label{entry:symbol-registry-00038}\(\bot,\top\) & \texttt{\detokenize{\bottom,}\allowbreak\detokenize{\topv}} & falsum and verum or declared CLL nullary units & global & n/a & \hyperref[sec:prop-syntax-semantics-01]{\ref*{sec:prop-syntax-semantics-01}} & \thepage & Outside the P8.6 CLL grammar these are bottom and the defined top abbreviation; in that grammar \(\bot\) is the par unit and \(\top\) is the with unit, with the ambient syntax fixing the role. \\
\phantomsection\label{entry:symbol-registry-00039}\(\mathbb N\) & \texttt{\detokenize{\N}} & set of natural numbers & global & n/a & \hyperref[sec:strings-01]{\ref*{sec:strings-01}} & \thepage & Contains 0 under the Part-1 natural-number convention. \\
\phantomsection\label{entry:symbol-registry-00040}\(\mathcal P(A)\) & \texttt{\detokenize{\mathcal P}} & power set & global & n/a & \hyperref[sec:strings-01]{\ref*{sec:strings-01}} & \thepage & Set of all subsets of the set A; set existence comes from the ambient foundation. \\
\phantomsection\label{entry:symbol-registry-00041}\(|X|\) & \texttt{\detokenize{\abs}} & cardinality & global & n/a & \hyperref[sec:strings-01]{\ref*{sec:strings-01}} & \thepage & Cardinality for set operands. Source-local uses for domains or symbolic length require their own explicit type and scope; the editorial domain form is \(\operatorname{dom}\). \\
\phantomsection\label{entry:symbol-registry-00042}\(\mathcal M\equiv\mathcal N\) & \texttt{\detokenize{\equiv}} & elementary equivalence of structures & global & meta & \hyperref[sec:prop-deduction-01]{\ref*{sec:prop-deduction-01}} & \thepage & The structures satisfy the same sentences. Structure operands distinguish this role from declared source-local formula equivalence. \\
\phantomsection\label{entry:symbol-registry-00043}\(\mathcal M\preccurlyeq\mathcal N\) & \texttt{\detokenize{\preccurlyeq}} & elementary substructure & global & meta & \hyperref[sec:prop-deduction-01]{\ref*{sec:prop-deduction-01}} & \thepage & Every formula with parameters from \(\mathcal M\) has the same truth value in \(\mathcal M\) and \(\mathcal N\). \\
\phantomsection\label{entry:symbol-registry-00048}\(\mathcal M\cong\mathcal N\) & \texttt{\detokenize{\cong}} & isomorphism & global & meta & \hyperref[sec:prop-deduction-01]{\ref*{sec:prop-deduction-01}} & \thepage & Structure-preserving bijection; \(\simeq\) remains separately reserved for Kleene equality. \\
\phantomsection\label{entry:symbol-registry-00075}\(\Gamma\Rightarrow\Delta\) & \texttt{\detokenize{\Rightarrow}} & sequent & global & meta & \hyperref[sec:prop-syntax-semantics-01]{\ref*{sec:prop-syntax-semantics-01}} & \thepage & Metalevel punctuation between antecedent and succedent contexts. \\
\phantomsection\label{entry:symbol-registry-00089}\(>\) & \texttt{\detokenize{>}} & strict ordering & global & 2 & \hyperref[sec:levels-01]{\ref*{sec:levels-01}} & \thepage & Contract list; strict dual of \(\leq\); ``positive'' means \(>0\) by the natural-numbers convention. \\
\phantomsection\label{entry:symbol-registry-00114}\(\Gamma\models_{S}\varphi\) & \texttt{\detokenize{\entails_{S}}} & semantic consequence (per-system turnstile) & global & meta & \hyperref[sec:strings-01]{\ref*{sec:strings-01}} & \thepage & The subscript names the semantics (\(\models_{\mathcal M}\) matrix, \(\models_{\mathrm{K3}}\), \(\models_{\mathrm{LP}}\), \(\models_{\mathsf{S4}}\), ...). \\
\phantomsection\label{entry:symbol-registry-00122}\(\mathcal M\models T\) & \texttt{\detokenize{\entails}} & model of a theory & global & meta & \hyperref[sec:prop-deduction-01]{\ref*{sec:prop-deduction-01}} & \thepage & \(\mathcal M\) satisfies every member of \(T\). \\
\phantomsection\label{entry:symbol-registry-00155}\(\mathsf{Var}\) & \texttt{\detokenize{\mathsf{Var}}} & first-order variable alphabet & global & n/a & \hyperref[sec:prop-deduction-01]{\ref*{sec:prop-deduction-01}} & \thepage & The fixed countably infinite object-language variable set; distinct from the occurrence-set operator \(\operatorname{Var}\). \\
\phantomsection\label{entry:symbol-registry-00156}\(\mathcal F,\mathcal R\) & \texttt{\detokenize{\mathcal F,}\allowbreak\detokenize{\mathcal R}} & function- and relation-symbol sets & global & n/a & \hyperref[sec:prop-deduction-01]{\ref*{sec:prop-deduction-01}} & \thepage & The disjoint nonlogical-symbol families of a first-order signature; constants are zero-ary function symbols. \\
\phantomsection\label{entry:symbol-registry-00157}\(\operatorname{ar}_{\mathcal F},\operatorname{ar}_{\mathcal R}\) & \texttt{\detokenize{\operatorname{ar}_{\mathcal F},}\allowbreak\detokenize{\operatorname{ar}_{\mathcal R}}} & signature arity maps & global & n/a & \hyperref[sec:prop-deduction-01]{\ref*{sec:prop-deduction-01}} & \thepage & Maps the function- and relation-symbol families to their finite arities. \\
\phantomsection\label{entry:symbol-registry-00158}\(\operatorname{Term}(\Lang),\operatorname{Fm}(\Lang),\operatorname{Sent}(\Lang)\) & \texttt{\detokenize{\operatorname{Term},}\allowbreak\detokenize{\operatorname{Fm},}\allowbreak\detokenize{\operatorname{Sent}}} & first-order term formula and sentence sets & global & n/a & \hyperref[sec:prop-deduction-01]{\ref*{sec:prop-deduction-01}} & \thepage & Least syntactic sets generated by the active first-order grammar; sentences are formulas with no free variables. \\
\phantomsection\label{entry:symbol-registry-00159}\(\operatorname{Var}(e)\) & \texttt{\detokenize{\operatorname{Var}}} & occurring-variable set & global & n/a & \hyperref[sec:prop-deduction-01]{\ref*{sec:prop-deduction-01}} & \thepage & All variables occurring in the declared term, formula, literal, clause, or finite clause set \(e\), whether free or bound; for a finite clause set the operator is the union of the occurrence sets of its clauses. \\
\phantomsection\label{entry:symbol-registry-00160}\(\varphi=_{\alpha}\psi\) & \texttt{\detokenize{=_{\alpha}}} & alpha-equivalence & global & meta & \hyperref[sec:prop-deduction-01]{\ref*{sec:prop-deduction-01}} & \thepage & Least constructor congruence generated by capture-avoiding bound-variable renaming; distinct from identity formula equivalence and elementary equivalence. \\
\phantomsection\label{entry:symbol-registry-00161}\(s[x\mapsto a]\) & \texttt{\detokenize{[x\mapsto a]}} & assignment update & global & n/a & \hyperref[sec:prop-deduction-01]{\ref*{sec:prop-deduction-01}} & \thepage & Agrees with assignment \(s\) except that variable \(x\) is sent to \(a\). \\
\phantomsection\label{entry:symbol-registry-00162}\(\operatorname{val}^{\mathcal M}_s(t)\) & \texttt{\detokenize{\operatorname{val}}} & term value & global & n/a & \hyperref[sec:prop-deduction-01]{\ref*{sec:prop-deduction-01}} & \thepage & Recursive denotation of the first-order term \(t\) in structure \(\mathcal M\) under assignment \(s\). \\
\phantomsection\label{entry:symbol-registry-00163}\(X\mathbin{\upharpoonright}Y\) & \texttt{\detokenize{\upharpoonright}} & restriction or reduct operator & global & n/a & \hyperref[sec:prop-deduction-01]{\ref*{sec:prop-deduction-01}} & \thepage & Operand types distinguish function or set restriction, a structure reduct, and the P8.3 public-announcement submodel induced by the surviving worlds. \\
\phantomsection\label{entry:symbol-registry-00164}\(\operatorname{Cn}(T)\) & \texttt{\detokenize{\operatorname{Cn}}} & semantic consequence closure & global & n/a & \hyperref[sec:prop-deduction-01]{\ref*{sec:prop-deduction-01}} & \thepage & Sentences of the active signature semantically entailed by theory \(T\). \\
\phantomsection\label{entry:symbol-registry-00188}\(\mathbb R\) & \texttt{\detokenize{\mathbb R}} & real ordered field & global & n/a & \hyperref[sec:prop-metatheory-01]{\ref*{sec:prop-metatheory-01}} & \thepage & The ordinary field of real numbers; the P2.3 categoricity claim names its ordered-field structure. \\
\phantomsection\label{entry:symbol-registry-00189}\(\mathbb Q\) & \texttt{\detokenize{\mathbb Q}} & rational ordered field & global & n/a & \hyperref[sec:arithmetic-coding-01]{\ref*{sec:arithmetic-coding-01}} & \thepage & The ordinary field of rational numbers; used as an order type in P5 and as the scalar domain of linear rational arithmetic in P12.2. \\
\phantomsection\label{entry:symbol-registry-00231}\(\mathbb Z\) & \texttt{\detokenize{\mathbb Z}} & set of integers & global & n/a & \hyperref[sec:maps-elementarity-01]{\ref*{sec:maps-elementarity-01}} & \thepage & Integers; used for tape positions and integer-coefficient polynomials. \\
\phantomsection\label{entry:symbol-registry-00232}\(f:A\rightharpoonup B\) & \texttt{\detokenize{\rightharpoonup}} & partial-function type arrow & global & n/a & \hyperref[sec:strings-01]{\ref*{sec:strings-01}} & \thepage & The function may be undefined outside a declared subset of \(A\). \\
\phantomsection\label{entry:symbol-registry-06555}\(a\mathrel{:=}b\) & \texttt{\detokenize{\defeq}} & source definitional sign & global & meta & \hyperref[sec:set-prerequisites-01]{\ref*{sec:set-prerequisites-01}} & \thepage & Printed only inside the restoration provenance scopes (environment inheritedrestorationscope), which occur throughout the book; there the macro \texttt{\detokenize{\defeq}} is redefined to print := and the editorial \(\mathrel{:\equiv}\) is restored on exit. The left expression abbreviates the right one; not object-language identity. Same role as the global editorial \(a\mathrel{:\equiv}b\), different glyph. \\
\end{longtable}
\endgroup
\Needspace{.4\textheight}
\subsection*{Stable part P1 scopes}
\begingroup\scriptsize\setlength{\tabcolsep}{3pt}
\begin{longtable}{@{}p{0.9in}p{1.3in}p{1.0in}p{0.5in}p{0.3in}p{0.4in}p{0.35in}p{1.65in}@{}}
\toprule
\textbf{Glyph} & \textbf{Macro} & \textbf{Role} & \textbf{Scope} & \textbf{Prec./form} & \textbf{First} & \textbf{Reg. page} & \textbf{Principal uses / notes} \\
\midrule
\endfirsthead
\toprule
\textbf{Glyph} & \textbf{Macro} & \textbf{Role} & \textbf{Scope} & \textbf{Prec./form} & \textbf{First} & \textbf{Reg. page} & \textbf{Principal uses / notes} \\
\midrule
\endhead
\midrule\multicolumn{8}{r}{\emph{continued}}\\\endfoot
\bottomrule\endlastfoot
\phantomsection\label{entry:symbol-registry-00079}\(\leq\) & \texttt{\detokenize{\leq}} & canonical ordering & Parts 1–4 from set-prerequisites to universal-reductions & 2 & \hyperref[sec:levels-01]{\ref*{sec:levels-01}} & \thepage & Ambient order on declared numerical and order-theoretic objects; P4.3 is excluded because bare \(\leq\) there is object-language arithmetic syntax. \\
\phantomsection\label{entry:symbol-registry-00085}\(<\) & \texttt{\detokenize{<}} & canonical strict ordering & Parts 1–4 from set-prerequisites to universal-reductions & 2 & \hyperref[sec:levels-01]{\ref*{sec:levels-01}} & \thepage & Ambient strict order on declared numerical and order-theoretic objects; P4.3 is excluded because bare \(<\) there is object-language arithmetic syntax. \\
\phantomsection\label{entry:symbol-registry-04969}\(X\) & \texttt{\detokenize{X}} & set carrying the generating binary relation & P1 Equivalence closure and quotients & object & \hyperref[sec:p01-equivalence-closure-and-quotients]{\ref*{sec:p01-equivalence-closure-and-quotients}} & \thepage & May be empty or infinite; no well-order or choice of paths is assumed. \\
\phantomsection\label{entry:symbol-registry-04970}\(A\) & \texttt{\detokenize{A}} & carrier of a general equivalence relation and quotient & P1 Equivalence closure and quotients & object & \hyperref[def:p01-ec-quotient]{\ref*{def:p01-ec-quotient}} & \thepage & The generated case specializes A to X without changing the quotient notation. \\
\phantomsection\label{entry:symbol-registry-04971}\(B\) & \texttt{\detokenize{B}} & codomain of a given function to be factored & P1 Equivalence closure and quotients & object & \hyperref[def:p01-ec-composition]{\ref*{def:p01-ec-composition}} & \thepage & May be empty when the stated input function exists; not assumed equal to the range. \\
\phantomsection\label{entry:symbol-registry-04972}\(R\) & \texttt{\detokenize{R}} & generating binary relation as a set of ordered pairs & P1 Equivalence closure and quotients & relation & \hyperref[sec:p01-equivalence-closure-and-quotients]{\ref*{sec:p01-equivalence-closure-and-quotients}} & \thepage & Infix xRy means membership of the ordered pair; preserves the inherited R notation. \\
\phantomsection\label{entry:symbol-registry-04973}\(R^{-1}\) & \texttt{\detokenize{R^{-1}}} & converse of the generating relation & P1 Equivalence closure and quotients & object & \hyperref[def:p01-ec-zigzag]{\ref*{def:p01-ec-zigzag}} & \thepage & Swaps coordinates of pairs; not a presumed inverse function. \\
\phantomsection\label{entry:symbol-registry-04974}\(R^*\) & \texttt{\detokenize{R^*}} & reflexive-transitive relation closure & P1 Equivalence closure and quotients & meta & \hyperref[thm:p01-ec-construction]{\ref*{thm:p01-ec-construction}} & \thepage & Relation scope only; not the earlier star on an alphabet's finite strings. \\
\phantomsection\label{entry:symbol-registry-04975}\(E\) & \texttt{\detokenize{E}} & equivalence relation constructed by paths or used for a quotient & P1 Equivalence closure and quotients & relation & \hyperref[sec:p01-equivalence-closure-and-quotients]{\ref*{sec:p01-equivalence-closure-and-quotients}} & \thepage & The path-defined E is proved equivalent before being used in quotient constructions; general quotient E retains the inherited type. \\
\phantomsection\label{entry:symbol-registry-04976}\(\mathcal F\) & \texttt{\detokenize{\mathcal F}} & set of all equivalence relations on X containing R & P1 Equivalence closure and quotients & object & \hyperref[def:p01-ec-constructions]{\ref*{def:p01-ec-constructions}} & \thepage & Bounded by the power set of X times X and proved nonempty using the full relation. \\
\phantomsection\label{entry:symbol-registry-04977}\(F\) & \texttt{\detokenize{F}} & an equivalence relation in the intersection family & P1 Equivalence closure and quotients & relation & \hyperref[def:p01-ec-constructions]{\ref*{def:p01-ec-constructions}} & \thepage & Not a finitary constructor or formula; inclusion in this family fixes its type. \\
\phantomsection\label{entry:symbol-registry-04978}\(I\) & \texttt{\detokenize{I}} & bounded intersection relation on X & P1 Equivalence closure and quotients & relation & \hyperref[def:p01-ec-constructions]{\ref*{def:p01-ec-constructions}} & \thepage & Its equality with the path construction is proved in both directions. \\
\phantomsection\label{entry:symbol-registry-04979}\(p\) & \texttt{\detokenize{p}} & ordered pair in the ambient product & P1 Equivalence closure and quotients & object & \hyperref[def:p01-ec-constructions]{\ref*{def:p01-ec-constructions}} & \thepage & Not a propositional variable in this local mathematical construction. \\
\phantomsection\label{entry:symbol-registry-04980}\(\N\) & \texttt{\detokenize{\N}} & natural-number set including zero & P1 Equivalence closure and quotients & object & \hyperref[sec:p01-equivalence-closure-and-quotients]{\ref*{sec:p01-equivalence-closure-and-quotients}} & \thepage & The inherited finite ordinal convention supplies path indices and lengths. \\
\phantomsection\label{entry:symbol-registry-04981}\(n\) & \texttt{\detokenize{n}} & number of steps in a finite path & P1 Equivalence closure and quotients & object & \hyperref[def:p01-ec-zigzag]{\ref*{def:p01-ec-zigzag}} & \thepage & A path of length zero has one vertex; its vertex sequence has length n plus one. \\
\phantomsection\label{entry:symbol-registry-04982}\(i\) & \texttt{\detokenize{i}} & finite path vertex or step index & P1 Equivalence closure and quotients & object & \hyperref[def:p01-ec-zigzag]{\ref*{def:p01-ec-zigzag}} & \thepage & Step indices satisfy i less than n; endpoint index n is not an additional edge. \\
\phantomsection\label{entry:symbol-registry-04983}\(x_i\) & \texttt{\detokenize{x_i}} & vertex of a specified path & P1 Equivalence closure and quotients & object & \hyperref[def:p01-ec-zigzag]{\ref*{def:p01-ec-zigzag}} & \thepage & The sequence is a function from the finite ordinal n plus one to X. \\
\phantomsection\label{entry:symbol-registry-04984}\((x,y,z,a,a')\) & \texttt{\detokenize{x;y;z;a;a'}} & elements of the declared relation carrier & P1 Equivalence closure and quotients & object & \hyperref[def:p01-ec-zigzag]{\ref*{def:p01-ec-zigzag}} & \thepage & Mathematical element variables rather than object-language variable tokens. \\
\phantomsection\label{entry:symbol-registry-04985}\([a]_E\) & \texttt{\detokenize{[a]_E}} & equivalence class with inherited orientation xEa & P1 Equivalence closure and quotients & object & \hyperref[def:p01-ec-quotient]{\ref*{def:p01-ec-quotient}} & \thepage & Class equality and nonempty intersection are proved equivalent to aEa'. \\
\phantomsection\label{entry:symbol-registry-04986}\(A/E\) & \texttt{\detokenize{A/}\allowbreak\detokenize{E}} & quotient set of equivalence classes & P1 Equivalence closure and quotients & object & \hyperref[def:p01-ec-quotient]{\ref*{def:p01-ec-quotient}} & \thepage & Bounded within the power set of A; empty A gives the empty quotient rather than a singleton containing the empty set. \\
\phantomsection\label{entry:symbol-registry-04987}\(q\) & \texttt{\detokenize{q}} & quotient projection from the carrier to its classes & P1 Equivalence closure and quotients & object & \hyperref[def:p01-ec-quotient]{\ref*{def:p01-ec-quotient}} & \thepage & Totality and surjectivity are proved by the graph and the class definition; no section is assumed. \\
\phantomsection\label{entry:symbol-registry-04988}\(C\) & \texttt{\detokenize{C}} & equivalence class used as a quotient input & P1 Equivalence closure and quotients & object & \hyperref[def:p01-ec-quotient]{\ref*{def:p01-ec-quotient}} & \thepage & Its local representative witness does not constitute a global choice function. \\
\phantomsection\label{entry:symbol-registry-04989}\(h\) & \texttt{\detokenize{h}} & given or candidate function from the quotient to B & P1 Equivalence closure and quotients & object & \hyperref[def:p01-ec-composition]{\ref*{def:p01-ec-composition}} & \thepage & Used in the uniqueness comparison with the constructed factor. \\
\phantomsection\label{entry:symbol-registry-04990}\(\circ\) & \texttt{\detokenize{\circ}} & composition of functions of matching types & P1 Equivalence closure and quotients & infix & \hyperref[def:p01-ec-composition]{\ref*{def:p01-ec-composition}} & \thepage & Declared pointwise before use; not the inherited metavariable for a binary connective. \\
\phantomsection\label{entry:symbol-registry-04991}\(f\) & \texttt{\detokenize{f}} & given function on the carrier with codomain B & P1 Equivalence closure and quotients & object & \hyperref[thm:p01-ec-factorization]{\ref*{thm:p01-ec-factorization}} & \thepage & Descent requires equality on equivalent inputs; equality on R-generators suffices only after the path argument. \\
\phantomsection\label{entry:symbol-registry-04992}\(\bar f\) & \texttt{\detokenize{\bar f}} & unique quotient factor when the constancy condition holds & P1 Equivalence closure and quotients & object & \hyperref[thm:p01-ec-factorization]{\ref*{thm:p01-ec-factorization}} & \thepage & Preserves the inherited bar-f notation and defines its graph without choosing representatives. \\
\phantomsection\label{entry:symbol-registry-04993}\((b,b')\) & \texttt{\detokenize{b;b'}} & candidate outputs in the codomain B & P1 Equivalence closure and quotients & object & \hyperref[thm:p01-ec-factorization]{\ref*{thm:p01-ec-factorization}} & \thepage & Equality of outputs is proved from equality on equivalent input elements. \\
\phantomsection\label{entry:symbol-registry-04994}\((C_0,C_1)\) & \texttt{\detokenize{C_0;C_1}} & explicit finite subsets later proved to be the two classes & P1 Equivalence closure and quotients & object & \hyperref[def:p01-ec-finite-data]{\ref*{def:p01-ec-finite-data}} & \thepage & Class standing is not inferred just from the appearance of the generator list. \\
\phantomsection\label{entry:symbol-registry-04995}\((u,v)\) & \texttt{\detokenize{u;v}} & two fixed distinct target elements in the finite example & P1 Equivalence closure and quotients & object & \hyperref[def:p01-ec-finite-data]{\ref*{def:p01-ec-finite-data}} & \thepage & Not the path sequence u\_i used in the concatenation proof. \\
\phantomsection\label{entry:symbol-registry-04996}\(g\) & \texttt{\detokenize{g}} & identity-valued function on the finite carrier & P1 Equivalence closure and quotients & object & \hyperref[thm:p01-ec-finite-example]{\ref*{thm:p01-ec-finite-example}} & \thepage & It violates representative independence and cannot factor through this quotient. \\
\phantomsection\label{entry:symbol-registry-04997}\(\bar g\) & \texttt{\detokenize{\bar g}} & rejected proposed quotient function formula & P1 Equivalence closure and quotients & meta & \hyperref[thm:p01-ec-finite-example]{\ref*{thm:p01-ec-finite-example}} & \thepage & No such function is established; one class would receive distinct outputs. \\
\phantomsection\label{entry:symbol-registry-04998}\(t\) & \texttt{\detokenize{t}} & explicit minimum-representative map for the finite example & P1 Equivalence closure and quotients & object & \hyperref[thm:p01-ec-finite-example]{\ref*{thm:p01-ec-finite-example}} & \thepage & A justified finite section is a different map and does not repair the failed factorization of g. \\
\phantomsection\label{entry:symbol-registry-04999}\(y_i\) & \texttt{\detokenize{y_i}} & vertex of a reversed or second specified path & P1 Equivalence closure and quotients & object & \hyperref[annex:p01.ec.construction]{\ref*{annex:p01.ec.construction}} & \thepage & Each proof paragraph supplies the corresponding sequence and its endpoints before use. \\
\phantomsection\label{entry:symbol-registry-05000}\((m,j)\) & \texttt{\detokenize{m;j}} & first-path length and second-path index in concatenation & P1 Equivalence closure and quotients & object & \hyperref[annex:p01.ec.construction]{\ref*{annex:p01.ec.construction}} & \thepage & The two vertex clauses agree at their shared index m including zero-length cases. \\
\phantomsection\label{entry:symbol-registry-05001}\(u_i\) & \texttt{\detokenize{u_i}} & vertex of the concatenated path & P1 Equivalence closure and quotients & object & \hyperref[annex:p01.ec.construction]{\ref*{annex:p01.ec.construction}} & \thepage & Defined on indices zero through m plus n; not a target label u in the finite example. \\
\phantomsection\label{entry:symbol-registry-05002}\(T\) & \texttt{\detokenize{T}} & auxiliary binary relation used for closure comparisons & P1 Equivalence closure and quotients & relation & \hyperref[annex:p01.ec.construction]{\ref*{annex:p01.ec.construction}} & \thepage & Declared reflexive-transitive in the star comparison and equivalence in the finite-block comparison. \\
\phantomsection\label{entry:symbol-registry-05003}\(c\) & \texttt{\detokenize{c}} & existential representative for one specified class in the factor proof & P1 Equivalence closure and quotients & object & \hyperref[annex:p01.ec.factorization]{\ref*{annex:p01.ec.factorization}} & \thepage & Not a function selecting representatives for all classes. \\
\phantomsection\label{entry:symbol-registry-05004}\(\mathsf{BPI}\) & \texttt{\detokenize{\mathsf{BPI}}} & Boolean Prime Ideal principle & p01-bpi-principles & meta & \hyperref[def:p01-bpi-principles]{\ref*{def:p01-bpi-principles}} & \thepage & Every nontrivial set-sized Boolean algebra has a prime ideal \\
\phantomsection\label{entry:symbol-registry-05005}\(0_B\) & \texttt{\detokenize{0_B}} & bottom element of a Boolean algebra & p01-bpi-algebra & object & \hyperref[def:p01-bpi-boolean-algebra]{\ref*{def:p01-bpi-boolean-algebra}} & \thepage & Typed algebraic constant; need not be the empty set \\
\phantomsection\label{entry:symbol-registry-05006}\(1_B\) & \texttt{\detokenize{1_B}} & top element of a Boolean algebra & p01-bpi-algebra & object & \hyperref[def:p01-bpi-boolean-algebra]{\ref*{def:p01-bpi-boolean-algebra}} & \thepage & Empty finite meet; distinct from bottom only in nontrivial algebras \\
\phantomsection\label{entry:symbol-registry-05007}\(\mathbf2\) & \texttt{\detokenize{\mathbf2}} & two-element Boolean algebra & p01-bpi-algebra & object & \hyperref[def:p01-bpi-boolean-algebra]{\ref*{def:p01-bpi-boolean-algebra}} & \thepage & Set of the two Boolean values with Boolean operations \\
\phantomsection\label{entry:symbol-registry-05008}\(\wedge\) & \texttt{\detokenize{\wedge}} & Boolean meet & p01-bpi-algebra & infix & \hyperref[def:p01-bpi-boolean-algebra]{\ref*{def:p01-bpi-boolean-algebra}} & \thepage & Algebraic operation; logical conjunction has a separately typed propositional scope \\
\phantomsection\label{entry:symbol-registry-05009}\(\vee\) & \texttt{\detokenize{\vee}} & Boolean join & p01-bpi-algebra & infix & \hyperref[def:p01-bpi-boolean-algebra]{\ref*{def:p01-bpi-boolean-algebra}} & \thepage & Algebraic operation; logical disjunction has a separately typed propositional scope \\
\phantomsection\label{entry:symbol-registry-05010}\(\neg\) & \texttt{\detokenize{\neg}} & Boolean complement & p01-bpi-algebra & object & \hyperref[def:p01-bpi-boolean-algebra]{\ref*{def:p01-bpi-boolean-algebra}} & \thepage & Unary algebraic operation rather than negation in an object theory \\
\phantomsection\label{entry:symbol-registry-05011}\(\leq_B\) & \texttt{\detokenize{\leq_B}} & Boolean lattice order & p01-bpi-algebra & relation & \hyperref[def:p01-bpi-boolean-algebra]{\ref*{def:p01-bpi-boolean-algebra}} & \thepage & Defined by meet absorption \\
\phantomsection\label{entry:symbol-registry-05012}\(\cdot\) & \texttt{\detokenize{\cdot}} & Koppelberg Boolean product & p01-bpi-source-dictionary & infix & \hyperref[def:p01-bpi-boolean-algebra]{\ref*{def:p01-bpi-boolean-algebra}} & \thepage & Published meet notation; not numerical multiplication \\
\phantomsection\label{entry:symbol-registry-05013}\(+\) & \texttt{\detokenize{+}} & Koppelberg Boolean sum & p01-bpi-source-dictionary & infix & \hyperref[def:p01-bpi-boolean-algebra]{\ref*{def:p01-bpi-boolean-algebra}} & \thepage & Published join notation; not numerical addition \\
\phantomsection\label{entry:symbol-registry-05014}\(-\) & \texttt{\detokenize{-}} & Koppelberg Boolean complement & p01-bpi-source-dictionary & object & \hyperref[def:p01-bpi-boolean-algebra]{\ref*{def:p01-bpi-boolean-algebra}} & \thepage & Published unary Boolean operation; not numerical subtraction \\
\phantomsection\label{entry:symbol-registry-05015}\(d(E)\) & \texttt{\detokenize{d(E)}} & Boolean-complement image of a subset & p01-bpi-algebra & meta & \hyperref[def:p01-bpi-filter-ideal]{\ref*{def:p01-bpi-filter-ideal}} & \thepage & Filter-ideal duality; not the set complement of the subset \\
\phantomsection\label{entry:symbol-registry-05016}\(\chi_U\) & \texttt{\detokenize{\chi_U}} & characteristic homomorphism of an ultrafilter & p01-bpi-algebra & meta & \hyperref[thm:p01-bpi-duality]{\ref*{thm:p01-bpi-duality}} & \thepage & Maps a Boolean algebra to the two-element algebra \\
\phantomsection\label{entry:symbol-registry-05017}\(\triangle\) & \texttt{\detokenize{\triangle}} & Boolean symmetric difference & p01-bpi-algebra & infix & \hyperref[def:p01-bpi-quotient]{\ref*{def:p01-bpi-quotient}} & \thepage & Meet-join-complement expression defining the ideal congruence \\
\phantomsection\label{entry:symbol-registry-05018}\(\sim_I\) & \texttt{\detokenize{\sim_I}} & congruence modulo a Boolean ideal & p01-bpi-algebra & relation & \hyperref[def:p01-bpi-quotient]{\ref*{def:p01-bpi-quotient}} & \thepage & Symmetric difference belongs to the ideal; not elementary equivalence \\
\phantomsection\label{entry:symbol-registry-05019}\(\equiv\) & \texttt{\detokenize{\equiv}} & Koppelberg local Boolean congruence & p01-bpi-source-dictionary & relation & \hyperref[def:p01-bpi-quotient]{\ref*{def:p01-bpi-quotient}} & \thepage & Explicitly quoted source variant corresponding to the ideal-parameterized congruence; not elementary equivalence \\
\phantomsection\label{entry:symbol-registry-05020}\([a]_I\) & \texttt{\detokenize{[a]_I}} & Boolean ideal equivalence class & p01-bpi-algebra & object & \hyperref[def:p01-bpi-quotient]{\ref*{def:p01-bpi-quotient}} & \thepage & Quotient class; no chosen representative selector \\
\phantomsection\label{entry:symbol-registry-05021}\(B/I\) & \texttt{\detokenize{B/}\allowbreak\detokenize{I}} & Boolean algebra quotient by an ideal & p01-bpi-algebra & object & \hyperref[def:p01-bpi-quotient]{\ref*{def:p01-bpi-quotient}} & \thepage & Underlying set of congruence classes with induced operations \\
\phantomsection\label{entry:symbol-registry-05022}\(e_\varepsilon\) & \texttt{\detokenize{e_\varepsilon}} & finite Boolean partition cell & p01-bpi-algebra & object & \hyperref[thm:p01-bpi-finite-local]{\ref*{thm:p01-bpi-finite-local}} & \thepage & Meet of literals indexed by one finite binary assignment \\
\phantomsection\label{entry:symbol-registry-05023}\(b_k^{\varepsilon_k}\) & \texttt{\detokenize{b_k^{\varepsilon_k}}} & Boolean literal selected by a bit & p01-bpi-algebra & object & \hyperref[thm:p01-bpi-finite-local]{\ref*{thm:p01-bpi-finite-local}} & \thepage & Superscript selects an element or its complement; not exponentiation \\
\phantomsection\label{entry:symbol-registry-05024}\(\bot\) & \texttt{\detokenize{\bot}} & propositional false constant & p01-bpi-propositional & object & \hyperref[def:p01-bpi-propositional-compactness]{\ref*{def:p01-bpi-propositional-compactness}} & \thepage & Truth-table constant; distinct from algebraic bottom notation \\
\phantomsection\label{entry:symbol-registry-05025}\(\top\) & \texttt{\detokenize{\top}} & propositional true constant & p01-bpi-propositional & object & \hyperref[def:p01-bpi-propositional-compactness]{\ref*{def:p01-bpi-propositional-compactness}} & \thepage & Explicit constant also for an empty variable set \\
\phantomsection\label{entry:symbol-registry-05026}\(\neg\) & \texttt{\detokenize{\neg}} & propositional negation & p01-bpi-propositional & object & \hyperref[def:p01-bpi-propositional-compactness]{\ref*{def:p01-bpi-propositional-compactness}} & \thepage & Formula constructor interpreted by Boolean complement \\
\phantomsection\label{entry:symbol-registry-05027}\(\land\) & \texttt{\detokenize{\land}} & propositional conjunction & p01-bpi-propositional & infix & \hyperref[def:p01-bpi-propositional-compactness]{\ref*{def:p01-bpi-propositional-compactness}} & \thepage & Formula constructor interpreted by Boolean meet \\
\phantomsection\label{entry:symbol-registry-05028}\(\lor\) & \texttt{\detokenize{\lor}} & propositional disjunction & p01-bpi-propositional & infix & \hyperref[def:p01-bpi-propositional-compactness]{\ref*{def:p01-bpi-propositional-compactness}} & \thepage & Formula constructor interpreted by Boolean join \\
\phantomsection\label{entry:symbol-registry-05029}\(\Leftrightarrow\) & \texttt{\detokenize{\Leftrightarrow}} & defined formula equivalence & p01-bpi-propositional & infix & \hyperref[def:p01-bpi-propositional-compactness]{\ref*{def:p01-bpi-propositional-compactness}} & \thepage & Abbreviation for equality of two truth values; not elementary equivalence \\
\phantomsection\label{entry:symbol-registry-05030}\(\nu\) & \texttt{\detokenize{\nu}} & total propositional valuation & p01-bpi-propositional & meta & \hyperref[def:p01-bpi-propositional-compactness]{\ref*{def:p01-bpi-propositional-compactness}} & \thepage & Function from the declared variable set to the two Boolean values \\
\phantomsection\label{entry:symbol-registry-05031}\(\operatorname{Con}\) & \texttt{\detokenize{\operatorname{Con}}} & syntactic consistency assertion & p01-bpi-choice-strength & meta & \hyperref[thm:p01-bpi-choice-separation]{\ref*{thm:p01-bpi-choice-separation}} & \thepage & Metatheoretic relative consistency; not an unconditional model-existence assertion \\
\phantomsection\label{entry:symbol-registry-06398}\(\operatorname{seg}(w)\) & \texttt{\detokenize{\operatorname{seg}}} & initial segment of a strictly well-ordered set below a member & \hyperref[sec:p01-w02-p01-011-finite-coding]{\ref*{sec:p01-w02-p01-011-finite-coding}} & n/a & \hyperref[def:p01-w02-p01-011-readings]{\ref*{def:p01-w02-p01-011-readings}} & \thepage & The set of predecessors of w under the given strict well-order, ordered by its restriction. \\
\phantomsection\label{entry:symbol-registry-06399}\(\max(\alpha,\beta)\) & \texttt{\detokenize{\max}} & larger of two ordinals & \hyperref[sec:p01-w02-p01-011-finite-coding]{\ref*{sec:p01-w02-p01-011-finite-coding}} & n/a & \hyperref[def:p01-w02-p01-011-readings]{\ref*{def:p01-w02-p01-011-readings}} & \thepage & Ordinal instance of the inherited max of a finite list of cardinals (sec:set-prerequisites-06); defined by ordinal trichotomy. \\
\phantomsection\label{entry:symbol-registry-06400}\(p\lhd q\) & \texttt{\detokenize{\lhd}} & canonical well-order on a square of ordinals & \hyperref[sec:p01-w02-p01-011-finite-coding]{\ref*{sec:p01-w02-p01-011-finite-coding}} & 2 & \hyperref[def:p01-w02-p01-011-canonical-order]{\ref*{def:p01-w02-p01-011-canonical-order}} & \thepage & Compares maximum, then first coordinate, then second coordinate. \\
\phantomsection\label{entry:symbol-registry-06401}\(k(p)\) & \texttt{\detokenize{k}} & lexicographic key of a pair of ordinals & \hyperref[def:p01-w02-p01-011-canonical-order]{\ref*{def:p01-w02-p01-011-canonical-order}} & n/a & \hyperref[def:p01-w02-p01-011-canonical-order]{\ref*{def:p01-w02-p01-011-canonical-order}} & \thepage & k(alpha,beta) is the triple (max(alpha,beta),alpha,beta); local name. \\
\phantomsection\label{entry:symbol-registry-06402}\(a<_{\mathrm{lex}}b\) & \texttt{\detokenize{<_{\mathrm{lex}}}} & lexicographic order on triples of ordinals & \hyperref[annex:p01-w02-p01-011-canonical-order]{\ref*{annex:p01-w02-p01-011-canonical-order}} & 2 & \hyperref[annex:p01-w02-p01-011-canonical-order]{\ref*{annex:p01-w02-p01-011-canonical-order}} & \thepage & Decided at the first index where the triples differ. \\
\phantomsection\label{entry:symbol-registry-06403}\(F(w)\) & \texttt{\detokenize{F}} & ordinal order isomorphic to an initial segment & \hyperref[thm:p01-w02-p01-011-representation]{\ref*{thm:p01-w02-p01-011-representation}} & n/a & \hyperref[thm:p01-w02-p01-011-representation]{\ref*{thm:p01-w02-p01-011-representation}} & \thepage & Unique by item 1 of that lemma; local name. \\
\phantomsection\label{entry:symbol-registry-06404}\(\operatorname{card}(\beta)\) & \texttt{\detokenize{\operatorname{card}}} & least ordinal bijective with an ordinal & \hyperref[sec:p01-w02-p01-011-finite-coding]{\ref*{sec:p01-w02-p01-011-finite-coding}} & n/a & \hyperref[thm:p01-w02-p01-011-infinite-ordinals]{\ref*{thm:p01-w02-p01-011-infinite-ordinals}} & \thepage & Defined without Choice for ordinals; it is a cardinal. Under Choice it agrees with the inherited |beta| of cor:cardinal-assignment. \\
\phantomsection\label{entry:symbol-registry-06405}\(A^{<\omega}\) & \texttt{\detokenize{A^{<\omega}}} & set of finite sequences from a set & \hyperref[sec:p01-w02-p01-011-finite-coding]{\ref*{sec:p01-w02-p01-011-finite-coding}} & n/a & \hyperref[def:p01-w02-p01-011-readings]{\ref*{def:p01-w02-p01-011-readings}} & \thepage & Inherited notation of cor:syntax-counting; equal to the string set A\^{}* of sec:strings-01. \\
\phantomsection\label{entry:symbol-registry-06406}\((B,\mathcal F)\) & \texttt{\detokenize{(B,}\allowbreak\detokenize{\mathcal F)}} & finitary grammar given by base strings and formation rules & \hyperref[sec:p01-w02-p01-011-syntax-counting]{\ref*{sec:p01-w02-p01-011-syntax-counting}} & n/a & \hyperref[def:p01-w02-p01-011-grammar]{\ref*{def:p01-w02-p01-011-grammar}} & \thepage & The inductive definition of sec:strings-02 with universe Gamma\^{}*. \\
\phantomsection\label{entry:symbol-registry-06407}\(e:A\to\N\) & \texttt{\detokenize{e}} & fixed injection coding the letters of an alphabet by natural numbers & \hyperref[sec:p01-w02-p01-012-string-coding]{\ref*{sec:p01-w02-p01-012-string-coding}} & n/a & \hyperref[def:p01-w02-p01-012-coding-maps]{\ref*{def:p01-w02-p01-012-coding-maps}} & \thepage & Source notation of the countability lemma in sec:strings-01; a metalanguage function, not an object-language symbol. \\
\phantomsection\label{entry:symbol-registry-06408}\(e_n\) & \texttt{\detokenize{e_n}} & stage code of strings of length n & \hyperref[sec:p01-w02-p01-012-string-coding]{\ref*{sec:p01-w02-p01-012-string-coding}} & n/a & \hyperref[def:p01-w02-p01-012-coding-maps]{\ref*{def:p01-w02-p01-012-coding-maps}} & \thepage & Source notation; defined by recursion on n from e and the numerical pairing map; domain A\^{}n. \\
\phantomsection\label{entry:symbol-registry-06409}\(\operatorname{pair}_{\N}(a,b)\) & \texttt{\detokenize{\operatorname{pair}_{\N}}} & numerical pairing function & \hyperref[sec:p01-w02-p01-012-string-coding]{\ref*{sec:p01-w02-p01-012-string-coding}} & n/a & \hyperref[sec:p01-w02-p01-012-string-coding]{\ref*{sec:p01-w02-p01-012-string-coding}} & \thepage & Same glyph and role as the registered P4 row (first use sec:turing-machines-01); the source fixes it first in sec:strings-01. Distinct from the set-theoretic ordered pair. \\
\phantomsection\label{entry:symbol-registry-06410}\(\operatorname{unpair}(z)\) & \texttt{\detokenize{\operatorname{unpair}}} & inverse of the numerical pairing function returning the pair of components & \hyperref[sec:p01-w02-p01-012-decoding-procedure]{\ref*{sec:p01-w02-p01-012-decoding-procedure}} & n/a & \hyperref[sec:p01-w02-p01-012-decoding-procedure]{\ref*{sec:p01-w02-p01-012-decoding-procedure}} & \thepage & Local name for the inverse computation of sec:strings-01 (t\_s interval index, b=z-t\_s, a=s-b). \\
\phantomsection\label{entry:symbol-registry-06411}\(t_s\) & \texttt{\detokenize{t_s}} & triangular number s(s+1)/2 & \hyperref[sec:p01-w02-p01-012-string-coding]{\ref*{sec:p01-w02-p01-012-string-coding}} & n/a & \hyperref[sec:p01-w02-p01-012-assumptions]{\ref*{sec:p01-w02-p01-012-assumptions}} & \thepage & Source notation of the pairing bijectivity argument in sec:strings-01. \\
\phantomsection\label{entry:symbol-registry-06412}\(\operatorname{len}(w)\) & \texttt{\detokenize{\operatorname{len}}} & length of a finite sequence & \hyperref[sec:p01-w02-p01-012-string-coding]{\ref*{sec:p01-w02-p01-012-string-coding}} & n/a & \hyperref[sec:p01-w02-p01-012-assumptions]{\ref*{sec:p01-w02-p01-012-assumptions}} & \thepage & Editorial spelling of the source command \texttt{\detokenize{\length}} of sec:strings-01, which is restricted to inherited scopes; same glyph as the registered finite word length row of P8. \\
\phantomsection\label{entry:symbol-registry-06413}\(w\mathbin{\upharpoonright}n\) & \texttt{\detokenize{\upharpoonright}} & restriction of a finite sequence to its first n arguments & \hyperref[sec:p01-w02-p01-012-string-coding]{\ref*{sec:p01-w02-p01-012-string-coding}} & n/a & \hyperref[sec:p01-w02-p01-012-assumptions]{\ref*{sec:p01-w02-p01-012-assumptions}} & \thepage & Instance of the registered global restriction operator; operand types fix the reading. \\
\phantomsection\label{entry:symbol-registry-06414}\(\langle x\rangle\) & \texttt{\detokenize{\langle x\rangle}} & string of length 1 with value x at 0 & \hyperref[sec:p01-w02-p01-012-decoding-procedure]{\ref*{sec:p01-w02-p01-012-decoding-procedure}} & n/a & \hyperref[sec:p01-w02-p01-012-decoding-procedure]{\ref*{sec:p01-w02-p01-012-decoding-procedure}} & \thepage & Local; the registered \(\langle C\rangle\) (clone generated, P2) has a non-overlapping scope. \\
\phantomsection\label{entry:symbol-registry-06415}\(d_k\) & \texttt{\detokenize{d_k}} & partial decoding map from numbers to strings of length k & \hyperref[def:p01-w02-p01-012-decoding]{\ref*{def:p01-w02-p01-012-decoding}} & n/a & \hyperref[def:p01-w02-p01-012-decoding]{\ref*{def:p01-w02-p01-012-decoding}} & \thepage & Local name in the decoding procedure; left inverse of e\_k on its image. \\
\phantomsection\label{entry:symbol-registry-06416}\(\mathtt a\) & \texttt{\detokenize{\mathtt a}} & letters of the two-letter example alphabet symbol family member \(\mathtt a\) & \hyperref[def:p01-w02-p01-012-two-letters]{\ref*{def:p01-w02-p01-012-two-letters}} & n/a & \hyperref[def:p01-w02-p01-012-two-letters]{\ref*{def:p01-w02-p01-012-two-letters}} & \thepage & Example-local letters. Explicit displayed family; member types and common actual first-use exposed, not one atomic glyph. Family listing, not one atomic glyph. Author macro \texttt{\detokenize{\mathtt a,}\allowbreak\detokenize{\mathtt b}} not carried to the members. \\
\phantomsection\label{entry:symbol-registry-06417}\(\mathtt b\) & \texttt{\detokenize{\mathtt b}} & letters of the two-letter example alphabet symbol family member \(\mathtt b\) & \hyperref[def:p01-w02-p01-012-two-letters]{\ref*{def:p01-w02-p01-012-two-letters}} & n/a & \hyperref[def:p01-w02-p01-012-two-letters]{\ref*{def:p01-w02-p01-012-two-letters}} & \thepage & Example-local letters. Explicit displayed family; member types and common actual first-use exposed, not one atomic glyph. Family listing, not one atomic glyph. Author macro \texttt{\detokenize{\mathtt a,}\allowbreak\detokenize{\mathtt b}} not carried to the members. \\
\phantomsection\label{entry:symbol-registry-06418}\(\Pi\) & \texttt{\detokenize{\Pi}} & background set of sentences fixed before translation & \hyperref[sec:p01-w02-p01-022-translation]{\ref*{sec:p01-w02-p01-022-translation}} & n/a & \hyperref[def:p01-w02-p01-022-setting]{\ref*{def:p01-w02-p01-022-setting}} & \thepage & Set of L-sentences recording connections between declared symbols; validity relative to Pi is consequence from Gamma together with Pi. Other registered Pi roles have disjoint scopes. \\
\phantomsection\label{entry:symbol-registry-06419}\(\varphi^{\mathcal M,s}_x\) & \texttt{\detokenize{\varphi^{\mathcal M,s}_x}} & solution set of a formula in a variable & \hyperref[sec:p01-w02-p01-022-translation]{\ref*{sec:p01-w02-p01-022-translation}} & n/a & \hyperref[def:p01-w02-p01-022-setting]{\ref*{def:p01-w02-p01-022-setting}} & \thepage & The set of a in dom M with M,s[x->a] satisfying phi; depends on s only on FV(phi) minus x. \\
\phantomsection\label{entry:symbol-registry-06420}\(\exists^{\leq1}x\,\varphi\) & \texttt{\detokenize{\exists^{\leq1}}} & at-most-one abbreviation & \hyperref[sec:p01-w02-p01-022-translation]{\ref*{sec:p01-w02-p01-022-translation}} & n/a & \hyperref[def:p01-w02-p01-022-uniqueness]{\ref*{def:p01-w02-p01-022-uniqueness}} & \thepage & Metalinguistic abbreviation with a comparison variable that does not occur in phi; not a new quantifier of the language. \\
\phantomsection\label{entry:symbol-registry-06421}\(\exists!x\,\varphi\) & \texttt{\detokenize{\exists!}} & unique-existence abbreviation & \hyperref[sec:p01-w02-p01-022-translation]{\ref*{sec:p01-w02-p01-022-translation}} & n/a & \hyperref[def:p01-w02-p01-022-uniqueness]{\ref*{def:p01-w02-p01-022-uniqueness}} & \thepage & Same typed role as the registered unique-existence row scoped to thm:p09-fagin1976-boundaries and as the abbreviation stated in sec:zf-zfc-01; this row records its use in this section, with the freshness condition of thm:p01-w02-p01-022-freshness. \\
\phantomsection\label{entry:symbol-registry-06422}\(\mathsf{the}_x(\varphi,\chi)\) & \texttt{\detokenize{\mathsf{the}_x}} & Russellian contextual elimination of a definite description & \hyperref[sec:p01-w02-p01-022-translation]{\ref*{sec:p01-w02-p01-022-translation}} & n/a & \hyperref[def:p01-w02-p01-022-description]{\ref*{def:p01-w02-p01-022-description}} & \thepage & Abbreviates exists x (phi and forall y (phi[y/x] -> y=x) and chi) with a comparison variable y not occurring in phi; not a term former. Distinct from the registered iota description term former. \\
\phantomsection\label{entry:symbol-registry-06423}\(s\leadsto t\) & \texttt{\detokenize{\leadsto}} & one-step relation of a step system & \hyperref[sec:p01-w02-p01-023-invariants]{\ref*{sec:p01-w02-p01-023-invariants}} & meta & \hyperref[def:p01-w02-p01-023-step-system]{\ref*{def:p01-w02-p01-023-step-system}} & \thepage & A binary relation on the state set \(S\) of a step system \(\langle S,I,\leadsto\rangle\); not an object-language connective. \\
\phantomsection\label{entry:symbol-registry-06424}\(\langle S,I,\leadsto\rangle\) & \texttt{\detokenize{S,}\allowbreak\detokenize{I}} & state set and initial states of a step system & \hyperref[sec:p01-w02-p01-023-invariants]{\ref*{sec:p01-w02-p01-023-invariants}} & n/a & \hyperref[def:p01-w02-p01-023-step-system]{\ref*{def:p01-w02-p01-023-step-system}} & \thepage & \(I\subseteq S\); local letters of the invariant subsection only. \\
\phantomsection\label{entry:symbol-registry-06425}\(\operatorname{Reach}\) & \texttt{\detokenize{\operatorname{Reach}}} & set of reachable states of a step system & \hyperref[sec:p01-w02-p01-023-invariants]{\ref*{sec:p01-w02-p01-023-invariants}} & n/a & \hyperref[def:p01-w02-p01-023-step-system]{\ref*{def:p01-w02-p01-023-step-system}} & \thepage & Ends of runs; distinct from the P2 second-order reachability formula \(\operatorname{Reach}(a,b)\), which is a formula with two arguments in another scope. \\
\phantomsection\label{entry:symbol-registry-06426}\(P\) & \texttt{\detokenize{P}} & properties (subsets of the state set) and the set of terminal states symbol family member \(P\) & \hyperref[sec:p01-w02-p01-023-invariants]{\ref*{sec:p01-w02-p01-023-invariants}} & n/a & \hyperref[def:p01-w02-p01-023-step-system]{\ref*{def:p01-w02-p01-023-step-system}} & \thepage & \(P\) inductive property, \(Q\) target property, \(E\) even integers in the non-inductive example, \(T\) terminal states. Explicit displayed family; member types and common actual first-use exposed, not one atomic glyph. Family listing, not one atomic glyph. Author macro \texttt{\detokenize{P,}\allowbreak\detokenize{Q,}\allowbreak\detokenize{E,}\allowbreak\detokenize{T}} not carried to the members. \\
\phantomsection\label{entry:symbol-registry-06427}\(Q\) & \texttt{\detokenize{Q}} & properties (subsets of the state set) and the set of terminal states symbol family member \(Q\) & \hyperref[sec:p01-w02-p01-023-invariants]{\ref*{sec:p01-w02-p01-023-invariants}} & n/a & \hyperref[def:p01-w02-p01-023-step-system]{\ref*{def:p01-w02-p01-023-step-system}} & \thepage & \(P\) inductive property, \(Q\) target property, \(E\) even integers in the non-inductive example, \(T\) terminal states. Explicit displayed family; member types and common actual first-use exposed, not one atomic glyph. Family listing, not one atomic glyph. Author macro \texttt{\detokenize{P,}\allowbreak\detokenize{Q,}\allowbreak\detokenize{E,}\allowbreak\detokenize{T}} not carried to the members. \\
\phantomsection\label{entry:symbol-registry-06428}\(E\) & \texttt{\detokenize{E}} & properties (subsets of the state set) and the set of terminal states symbol family member \(E\) & \hyperref[sec:p01-w02-p01-023-invariants]{\ref*{sec:p01-w02-p01-023-invariants}} & n/a & \hyperref[def:p01-w02-p01-023-step-system]{\ref*{def:p01-w02-p01-023-step-system}} & \thepage & \(P\) inductive property, \(Q\) target property, \(E\) even integers in the non-inductive example, \(T\) terminal states. Explicit displayed family; member types and common actual first-use exposed, not one atomic glyph. Family listing, not one atomic glyph. Author macro \texttt{\detokenize{P,}\allowbreak\detokenize{Q,}\allowbreak\detokenize{E,}\allowbreak\detokenize{T}} not carried to the members. \\
\phantomsection\label{entry:symbol-registry-06429}\(T\) & \texttt{\detokenize{T}} & properties (subsets of the state set) and the set of terminal states symbol family member \(T\) & \hyperref[sec:p01-w02-p01-023-invariants]{\ref*{sec:p01-w02-p01-023-invariants}} & n/a & \hyperref[def:p01-w02-p01-023-step-system]{\ref*{def:p01-w02-p01-023-step-system}} & \thepage & \(P\) inductive property, \(Q\) target property, \(E\) even integers in the non-inductive example, \(T\) terminal states. Explicit displayed family; member types and common actual first-use exposed, not one atomic glyph. Family listing, not one atomic glyph. Author macro \texttt{\detokenize{P,}\allowbreak\detokenize{Q,}\allowbreak\detokenize{E,}\allowbreak\detokenize{T}} not carried to the members. \\
\phantomsection\label{entry:symbol-registry-06430}\(\mu\) & \texttt{\detokenize{\mu}} & variant (termination measure) into the natural numbers & \hyperref[sec:p01-w02-p01-023-invariants]{\ref*{sec:p01-w02-p01-023-invariants}} & n/a & \hyperref[def:p01-w02-p01-023-step-system]{\ref*{def:p01-w02-p01-023-step-system}} & \thepage & Decreases strictly on every step from a state of the inductive property; \(\mu_\Phi\) is the forward-chaining instance. \\
\phantomsection\label{entry:symbol-registry-06431}\(\mathcal R\) & \texttt{\detokenize{\mathcal R}} & rule system (set of rule instances) & \hyperref[thm:p01-w02-p01-023-derivation-invariant]{\ref*{thm:p01-w02-p01-023-derivation-invariant}} & n/a & \hyperref[thm:p01-w02-p01-023-derivation-invariant]{\ref*{thm:p01-w02-p01-023-derivation-invariant}} & \thepage & Pairs \((\langle u_1,\ldots,u_k\rangle,u)\) of premise tuple and conclusion over an item set \(U\); introduced in the unlabelled rule-system paragraph immediately before the theorem; local to that paragraph and theorem. \\
\phantomsection\label{entry:symbol-registry-06432}\((B,h)\) & \texttt{\detokenize{B,}\allowbreak\detokenize{h}} & Horn clause in pair form (negative atoms and positive literal or falsum) & \hyperref[def:p01-w02-p01-023-forward-chaining]{\ref*{def:p01-w02-p01-023-forward-chaining}} & n/a & \hyperref[def:p01-w02-p01-023-forward-chaining]{\ref*{def:p01-w02-p01-023-forward-chaining}} & \thepage & Represents \(\neg b_1\lor\cdots\lor\neg b_k\lor h\); \(h=\bottom\) when the clause has no positive literal. The pair form is introduced in the unlabelled paragraph immediately before the definition. \\
\phantomsection\label{entry:symbol-registry-06433}\(\Phi\) & \texttt{\detokenize{\Phi}} & Horn clause set & \hyperref[def:p01-w02-p01-023-forward-chaining]{\ref*{def:p01-w02-p01-023-forward-chaining}} & n/a & \hyperref[def:p01-w02-p01-023-forward-chaining]{\ref*{def:p01-w02-p01-023-forward-chaining}} & \thepage & Local to the forward-chaining definition, theorem and example. Blank macro filled with the glyph's TeX spelling. \\
\phantomsection\label{entry:symbol-registry-06434}\(A\) & \texttt{\detokenize{A}} & atoms of a Horn clause set & \hyperref[def:p01-w02-p01-023-forward-chaining]{\ref*{def:p01-w02-p01-023-forward-chaining}} & n/a & \hyperref[def:p01-w02-p01-023-forward-chaining]{\ref*{def:p01-w02-p01-023-forward-chaining}} & \thepage & Local to the forward-chaining definition, theorem and example. Blank macro filled with the glyph's TeX spelling. \\
\phantomsection\label{entry:symbol-registry-06435}\(M\) & \texttt{\detokenize{M}} & forward-chaining state & \hyperref[def:p01-w02-p01-023-forward-chaining]{\ref*{def:p01-w02-p01-023-forward-chaining}} & n/a & \hyperref[def:p01-w02-p01-023-forward-chaining]{\ref*{def:p01-w02-p01-023-forward-chaining}} & \thepage & Local to the forward-chaining definition, theorem and example. Blank macro filled with the glyph's TeX spelling. \\
\phantomsection\label{entry:symbol-registry-06436}\(P_\Phi\) & \texttt{\detokenize{P_\Phi}} & invariant of forward chaining & \hyperref[def:p01-w02-p01-023-forward-chaining]{\ref*{def:p01-w02-p01-023-forward-chaining}} & n/a & \hyperref[def:p01-w02-p01-023-forward-chaining]{\ref*{def:p01-w02-p01-023-forward-chaining}} & \thepage & Local to the forward-chaining definition, theorem and example. Blank macro filled with the glyph's TeX spelling. \\
\phantomsection\label{entry:symbol-registry-06437}\(\mu_\Phi\) & \texttt{\detokenize{\mu_\Phi}} & variant of forward chaining & \hyperref[def:p01-w02-p01-023-forward-chaining]{\ref*{def:p01-w02-p01-023-forward-chaining}} & n/a & \hyperref[def:p01-w02-p01-023-forward-chaining]{\ref*{def:p01-w02-p01-023-forward-chaining}} & \thepage & Local to the forward-chaining definition, theorem and example. Blank macro filled with the glyph's TeX spelling. \\
\phantomsection\label{entry:symbol-registry-06438}\(K_2\) & \texttt{\detokenize{K_2}} & two-world and three-world intuitionistic Kripke models symbol family member \(K_2\) & \hyperref[annex:p01-w02-p01-023-classical-principles]{\ref*{annex:p01-w02-p01-023-classical-principles}} & n/a & \hyperref[annex:p01-w02-p01-023-classical-principles]{\ref*{annex:p01-w02-p01-023-classical-principles}} & \thepage & Countermodels for the classical principles; worlds \(w_0,w_1,w_2\). Explicit displayed family; member types and common actual first-use exposed, not one atomic glyph. Family listing, not one atomic glyph. Author macro \texttt{\detokenize{K\_2,}\allowbreak\detokenize{K\_3}} not carried to the members. \\
\phantomsection\label{entry:symbol-registry-06439}\(K_3\) & \texttt{\detokenize{K_3}} & two-world and three-world intuitionistic Kripke models symbol family member \(K_3\) & \hyperref[annex:p01-w02-p01-023-classical-principles]{\ref*{annex:p01-w02-p01-023-classical-principles}} & n/a & \hyperref[annex:p01-w02-p01-023-classical-principles]{\ref*{annex:p01-w02-p01-023-classical-principles}} & \thepage & Countermodels for the classical principles; worlds \(w_0,w_1,w_2\). Explicit displayed family; member types and common actual first-use exposed, not one atomic glyph. Family listing, not one atomic glyph. Author macro \texttt{\detokenize{K\_2,}\allowbreak\detokenize{K\_3}} not carried to the members. \\
\phantomsection\label{entry:symbol-registry-06440}\(\Gamma\models_{\mathsf{CPL}}\varphi\) & \texttt{\detokenize{\entails_{\mathsf{CPL}}}} & classical two-valued propositional consequence & \hyperref[sec:p01-w02-p01-023-governing-logic]{\ref*{sec:p01-w02-p01-023-governing-logic}} & meta & \hyperref[sec:p01-w02-p01-023-governing-logic]{\ref*{sec:p01-w02-p01-023-governing-logic}} & \thepage & Instance of the registered per-system pattern \(\models_{S}\); negation read as \(\varphi\to\bottom\), which has the classical truth value of primitive negation. \\
\phantomsection\label{entry:symbol-registry-06441}\(\Gamma\vdash_{\mathsf{IPL}}\varphi\) & \texttt{\detokenize{\Gamma\vdash_{\mathsf{IPL}}\varphi}} & IPL derivability and Kripke consequence used in P1 symbol family member \(\Gamma\vdash_{\mathsf{IPL}}\varphi\) & \hyperref[sec:p01-w02-p01-023-governing-logic]{\ref*{sec:p01-w02-p01-023-governing-logic}} & meta & \hyperref[sec:p01-w02-p01-023-governing-logic]{\ref*{sec:p01-w02-p01-023-governing-logic}} & \thepage & Same relations as the registered P8 rows; this row extends their scope to the P1 governing-logic subsection without changing the role. Explicit displayed family; member types and common actual first-use exposed, not one atomic glyph. Family listing, not one atomic glyph. Author macro \texttt{\detokenize{\derives\_{\mathsf{IPL}},}\allowbreak\detokenize{\entails\_{\mathsf{IPL}}}} not carried to the members. \\
\phantomsection\label{entry:symbol-registry-06442}\(\ \Gamma\models_{\mathsf{IPL}}\varphi\) & \texttt{\detokenize{\ \Gamma\models_{\mathsf{IPL}}\varphi}} & IPL derivability and Kripke consequence used in P1 symbol family member \(\ \Gamma\models_{\mathsf{IPL}}\varphi\) & \hyperref[sec:p01-w02-p01-023-governing-logic]{\ref*{sec:p01-w02-p01-023-governing-logic}} & meta & \hyperref[sec:p01-w02-p01-023-governing-logic]{\ref*{sec:p01-w02-p01-023-governing-logic}} & \thepage & Same relations as the registered P8 rows; this row extends their scope to the P1 governing-logic subsection without changing the role. Explicit displayed family; member types and common actual first-use exposed, not one atomic glyph. Family listing, not one atomic glyph. Author macro \texttt{\detokenize{\derives\_{\mathsf{IPL}},}\allowbreak\detokenize{\entails\_{\mathsf{IPL}}}} not carried to the members. \\
\phantomsection\label{entry:symbol-registry-06443}\(X\) & \texttt{\detokenize{X}} & the formula \(\neg(p\land q)\to(\neg p\lor\neg q)\) & \hyperref[annex:p01-w02-p01-023-repairs]{\ref*{annex:p01-w02-p01-023-repairs}} & n/a & \hyperref[annex:p01-w02-p01-023-repairs]{\ref*{annex:p01-w02-p01-023-repairs}} & \thepage & Local abbreviation in one annex proof. \\
\phantomsection\label{entry:symbol-registry-06444}\(G(g)\) & \texttt{\detokenize{G}} & functional rule assigning a set to each function whose domain is an ordinal below a fixed ordinal & \hyperref[sec:p01-w02-p01-013-transfinite-recursion]{\ref*{sec:p01-w02-p01-013-transfinite-recursion}} & n/a & \hyperref[def:p01-w02-p01-013-reading]{\ref*{def:p01-w02-p01-013-reading}} & \thepage & Source notation of thm:transfinite-recursion; given by a formula theta(g,y,p) with set parameters. Local; overlaps other local G rows (for example the P6 Gale-Stewart game G(A)). \\
\phantomsection\label{entry:symbol-registry-06445}\(F\) & \texttt{\detokenize{F}} & function defined by transfinite recursion on an ordinal & \hyperref[sec:p01-w02-p01-013-transfinite-recursion]{\ref*{sec:p01-w02-p01-013-transfinite-recursion}} & n/a & \hyperref[thm:p01-w02-p01-013-recursion]{\ref*{thm:p01-w02-p01-013-recursion}} & \thepage & Source notation of thm:transfinite-recursion; local. Distinct from the local F of thm:p01-w02-p01-011-representation. \\
\phantomsection\label{entry:symbol-registry-06446}\(\delta\) & \texttt{\detokenize{\delta}} & length of a recursion (an ordinal) & \hyperref[sec:p01-w02-p01-013-transfinite-recursion]{\ref*{sec:p01-w02-p01-013-transfinite-recursion}} & n/a & \hyperref[def:p01-w02-p01-013-reading]{\ref*{def:p01-w02-p01-013-reading}} & \thepage & Source notation of thm:transfinite-recursion; local. \\
\phantomsection\label{entry:symbol-registry-06447}\(f_\alpha\) & \texttt{\detokenize{f_\alpha}} & the unique attempt of length alpha & \hyperref[thm:p01-w02-p01-013-assembly]{\ref*{thm:p01-w02-p01-013-assembly}} & n/a & \hyperref[thm:p01-w02-p01-013-assembly]{\ref*{thm:p01-w02-p01-013-assembly}} & \thepage & Defined only when an attempt of length alpha exists; unique by thm:p01-w02-p01-013-agreement. \\
\phantomsection\label{entry:symbol-registry-06448}\(h_\gamma\) & \texttt{\detokenize{h_\gamma}} & assembled attempt of length gamma & \hyperref[thm:p01-w02-p01-013-assembly]{\ref*{thm:p01-w02-p01-013-assembly}} & n/a & \hyperref[thm:p01-w02-p01-013-assembly]{\ref*{thm:p01-w02-p01-013-assembly}} & \thepage & Formed by one Replacement instance; local. \\
\phantomsection\label{entry:symbol-registry-06449}\(F_\gamma\) & \texttt{\detokenize{F_\gamma}} & recursion function for the shorter length gamma & \hyperref[thm:p01-w02-p01-013-coherence]{\ref*{thm:p01-w02-p01-013-coherence}} & n/a & \hyperref[thm:p01-w02-p01-013-coherence]{\ref*{thm:p01-w02-p01-013-coherence}} & \thepage & Equals the restriction of F to gamma; local. \\
\phantomsection\label{entry:symbol-registry-06450}\(S(\eta,x)\) & \texttt{\detokenize{S}} & successor clause of a three-clause recursion & \hyperref[thm:p01-w02-p01-013-cases]{\ref*{thm:p01-w02-p01-013-cases}} & n/a & \hyperref[thm:p01-w02-p01-013-cases]{\ref*{thm:p01-w02-p01-013-cases}} & \thepage & Functional rule on pairs of an ordinal below delta and a set; local. \\
\phantomsection\label{entry:symbol-registry-06451}\(L(g)\) & \texttt{\detokenize{L}} & limit clause of a three-clause recursion & \hyperref[thm:p01-w02-p01-013-cases]{\ref*{thm:p01-w02-p01-013-cases}} & n/a & \hyperref[thm:p01-w02-p01-013-cases]{\ref*{thm:p01-w02-p01-013-cases}} & \thepage & Functional rule on functions whose domain is a limit ordinal below delta; local. Not a language or the constructible hierarchy. \\
\phantomsection\label{entry:symbol-registry-06452}\(\Phi(X)\) & \texttt{\detokenize{\Phi}} & closure operator X maps to {0} union successors of X on subsets of N & \hyperref[def:p01-w02-p01-013-closure-example]{\ref*{def:p01-w02-p01-013-closure-example}} & n/a & \hyperref[def:p01-w02-p01-013-closure-example]{\ref*{def:p01-w02-p01-013-closure-example}} & \thepage & Example-local. \\
\phantomsection\label{entry:symbol-registry-06453}\(U_\alpha\) & \texttt{\detokenize{U_\alpha}} & union of the values of F below alpha & \hyperref[annex:p01-w02-p01-013-closure-stages]{\ref*{annex:p01-w02-p01-013-closure-stages}} & n/a & \hyperref[annex:p01-w02-p01-013-closure-stages]{\ref*{annex:p01-w02-p01-013-closure-stages}} & \thepage & Proof-local. \\
\phantomsection\label{entry:symbol-registry-06454}\(T_n\) & \texttt{\detokenize{T_n}} & set of natural numbers above n & \hyperref[sec:p01-w02-p01-013-limits]{\ref*{sec:p01-w02-p01-013-limits}} & n/a & \hyperref[sec:p01-w02-p01-013-limits]{\ref*{sec:p01-w02-p01-013-limits}} & \thepage & Positions earlier than n in the converse order on N; local. \\
\phantomsection\label{entry:symbol-registry-06455}\(V(g)\) & \texttt{\detokenize{V}} & set of values of a function & \hyperref[thm:p01-w02-p01-013-ill-founded]{\ref*{thm:p01-w02-p01-013-ill-founded}} & n/a & \hyperref[thm:p01-w02-p01-013-ill-founded]{\ref*{thm:p01-w02-p01-013-ill-founded}} & \thepage & Lemma-local; not the cumulative hierarchy V\_alpha. \\
\phantomsection\label{entry:symbol-registry-06456}\(G_1\) & \texttt{\detokenize{G_1}} & rules of the two counterexamples symbol family member \(G_1\) & \hyperref[thm:p01-w02-p01-013-ill-founded]{\ref*{thm:p01-w02-p01-013-ill-founded}} & n/a & \hyperref[thm:p01-w02-p01-013-ill-founded]{\ref*{thm:p01-w02-p01-013-ill-founded}} & \thepage & Lemma-local. Explicit displayed family; member types and common actual first-use exposed, not one atomic glyph. Family listing, not one atomic glyph. Author macro \texttt{\detokenize{G\_1,}\allowbreak\detokenize{G\_2}} not carried to the members. \\
\phantomsection\label{entry:symbol-registry-06457}\(G_2\) & \texttt{\detokenize{G_2}} & rules of the two counterexamples symbol family member \(G_2\) & \hyperref[thm:p01-w02-p01-013-ill-founded]{\ref*{thm:p01-w02-p01-013-ill-founded}} & n/a & \hyperref[thm:p01-w02-p01-013-ill-founded]{\ref*{thm:p01-w02-p01-013-ill-founded}} & \thepage & Lemma-local. Explicit displayed family; member types and common actual first-use exposed, not one atomic glyph. Family listing, not one atomic glyph. Author macro \texttt{\detokenize{G\_1,}\allowbreak\detokenize{G\_2}} not carried to the members. \\
\phantomsection\label{entry:symbol-registry-06458}\(\Gamma\models_{\Lang}\varphi\) & \texttt{\detokenize{\entails_{\Lang}}} & semantic consequence over all structures of a named first-order signature & \hyperref[sec:p01-w02-p01-003-language-change]{\ref*{sec:p01-w02-p01-003-language-change}} & meta & \hyperref[def:p01-w02-p01-003-signature-subscripts]{\ref*{def:p01-w02-p01-003-signature-subscripts}} & \thepage & Instance of the registered per-system pattern \(\models_{S}\) with the subscript naming the signature whose structures and assignments form the permitted data; the unsubscripted inherited \(\Gamma\models\varphi\) is the case of one fixed signature. \\
\phantomsection\label{entry:symbol-registry-06459}\(\mathsf{ND}_{\mathsf{FOL}}[\Lang]\) & \texttt{\detokenize{\mathsf{ND}_{\mathsf{FOL}}[\Lang]}} & first-order natural deduction with rule instances over a named signature & \hyperref[sec:p01-w02-p01-003-language-change]{\ref*{sec:p01-w02-p01-003-language-change}} & n/a & \hyperref[def:p01-w02-p01-003-signature-subscripts]{\ref*{def:p01-w02-p01-003-signature-subscripts}} & \thepage & Same rule schemata as the registered \(\mathsf{ND}_{\mathsf{FOL}}\); the bracket records that formulas and instantiating terms are drawn from \(\operatorname{Fm}(\Lang)\) and \(\operatorname{Term}(\Lang)\). \\
\phantomsection\label{entry:symbol-registry-06460}\(\Gamma\vdash_{\mathsf{ND}_{\mathsf{FOL}}[\Lang]}\varphi\) & \texttt{\detokenize{\derives_{\mathsf{ND}_{\mathsf{FOL}}[\Lang]}}} & derivability in first-order natural deduction over a named signature & \hyperref[sec:p01-w02-p01-003-language-change]{\ref*{sec:p01-w02-p01-003-language-change}} & meta & \hyperref[def:p01-w02-p01-003-signature-subscripts]{\ref*{def:p01-w02-p01-003-signature-subscripts}} & \thepage & Instance of the registered per-system turnstile \(\vdash_{S}\); the subscript names the calculus \(\mathsf{ND}_{\mathsf{FOL}}[\Lang]\). \\
\phantomsection\label{entry:symbol-registry-06461}\(\mathcal S_{\Lang}\) & \texttt{\detokenize{\mathcal S_{\Lang}}} & first-order semantics over a named signature & \hyperref[sec:p01-w02-p01-003-language-change]{\ref*{sec:p01-w02-p01-003-language-change}} & n/a & \hyperref[def:p01-w02-p01-003-signature-subscripts]{\ref*{def:p01-w02-p01-003-signature-subscripts}} & \thepage & Permits exactly the pairs of an \(\Lang\)-structure and an assignment in it; the third component of the logical-system triple of sec:levels-01. \\
\phantomsection\label{entry:symbol-registry-06462}\(\Lang'\) & \texttt{\detokenize{\Lang'}} & extending first-order signature & \hyperref[sec:p01-w02-p01-003-language-change]{\ref*{sec:p01-w02-p01-003-language-change}} & n/a & \hyperref[sec:p01-w02-p01-003-language-change]{\ref*{sec:p01-w02-p01-003-language-change}} & \thepage & A signature extending \(\Lang\) in the sense of the inherited signature-extension definition; not a derivative or a dual. \\
\phantomsection\label{entry:symbol-registry-06463}\(\Lang_0\) & \texttt{\detokenize{\Lang_0}} & common sub-signature and two compared signatures symbol family member \(\Lang_0\) & \hyperref[thm:p01-w02-p01-003-common-subsignature]{\ref*{thm:p01-w02-p01-003-common-subsignature}} & n/a & \hyperref[thm:p01-w02-p01-003-common-subsignature]{\ref*{thm:p01-w02-p01-003-common-subsignature}} & \thepage & \(\Lang_0\) keeps the symbols shared by \(\Lang_1\) and \(\Lang_2\) with equal arity. Explicit displayed family; member types and common actual first-use exposed, not one atomic glyph. Family listing, not one atomic glyph. Author macro \texttt{\detokenize{\Lang\_0,}\allowbreak\detokenize{\Lang\_1,}\allowbreak\detokenize{\Lang\_2}} not carried to the members. \\
\phantomsection\label{entry:symbol-registry-06464}\(\Lang_1\) & \texttt{\detokenize{\Lang_1}} & common sub-signature and two compared signatures symbol family member \(\Lang_1\) & \hyperref[thm:p01-w02-p01-003-common-subsignature]{\ref*{thm:p01-w02-p01-003-common-subsignature}} & n/a & \hyperref[thm:p01-w02-p01-003-common-subsignature]{\ref*{thm:p01-w02-p01-003-common-subsignature}} & \thepage & \(\Lang_0\) keeps the symbols shared by \(\Lang_1\) and \(\Lang_2\) with equal arity. Explicit displayed family; member types and common actual first-use exposed, not one atomic glyph. Family listing, not one atomic glyph. Author macro \texttt{\detokenize{\Lang\_0,}\allowbreak\detokenize{\Lang\_1,}\allowbreak\detokenize{\Lang\_2}} not carried to the members. \\
\phantomsection\label{entry:symbol-registry-06465}\(\Lang_2\) & \texttt{\detokenize{\Lang_2}} & common sub-signature and two compared signatures symbol family member \(\Lang_2\) & \hyperref[thm:p01-w02-p01-003-common-subsignature]{\ref*{thm:p01-w02-p01-003-common-subsignature}} & n/a & \hyperref[thm:p01-w02-p01-003-common-subsignature]{\ref*{thm:p01-w02-p01-003-common-subsignature}} & \thepage & \(\Lang_0\) keeps the symbols shared by \(\Lang_1\) and \(\Lang_2\) with equal arity. Explicit displayed family; member types and common actual first-use exposed, not one atomic glyph. Family listing, not one atomic glyph. Author macro \texttt{\detokenize{\Lang\_0,}\allowbreak\detokenize{\Lang\_1,}\allowbreak\detokenize{\Lang\_2}} not carried to the members. \\
\phantomsection\label{entry:symbol-registry-06466}\(\mathcal K_A\) & \texttt{\detokenize{\mathcal K_A}} & class of structures satisfying the axiom A & \hyperref[thm:p01-w02-p01-003-restricted-class]{\ref*{thm:p01-w02-p01-003-restricted-class}} & n/a & \hyperref[thm:p01-w02-p01-003-restricted-class]{\ref*{thm:p01-w02-p01-003-restricted-class}} & \thepage & Restricted semantic class; \(\Delta\models_{\mathcal K_A}\chi\) quantifies over its members and their assignments. \\
\phantomsection\label{entry:symbol-registry-06467}\(A\) & \texttt{\detokenize{A}} & axiom with the new constant and old-signature sentence in the failure case symbol family member \(A\) & \hyperref[thm:p01-w02-p01-003-restricted-class]{\ref*{thm:p01-w02-p01-003-restricted-class}} & n/a & \hyperref[thm:p01-w02-p01-003-restricted-class]{\ref*{thm:p01-w02-p01-003-restricted-class}} & \thepage & \(A=\forall x\,(x=c\lor R(c,x))\) and \(\theta=\exists y\,\forall x\,(x=y\lor R(y,x))\); local names only. Explicit displayed family; member types and common actual first-use exposed, not one atomic glyph. Family listing, not one atomic glyph. Author macro \texttt{\detokenize{A,}\allowbreak\detokenize{\theta}} not carried to the members. \\
\phantomsection\label{entry:symbol-registry-06468}\(\theta\) & \texttt{\detokenize{\theta}} & axiom with the new constant and old-signature sentence in the failure case symbol family member \(\theta\) & \hyperref[thm:p01-w02-p01-003-restricted-class]{\ref*{thm:p01-w02-p01-003-restricted-class}} & n/a & \hyperref[thm:p01-w02-p01-003-restricted-class]{\ref*{thm:p01-w02-p01-003-restricted-class}} & \thepage & \(A=\forall x\,(x=c\lor R(c,x))\) and \(\theta=\exists y\,\forall x\,(x=y\lor R(y,x))\); local names only. Explicit displayed family; member types and common actual first-use exposed, not one atomic glyph. Family listing, not one atomic glyph. Author macro \texttt{\detokenize{A,}\allowbreak\detokenize{\theta}} not carried to the members. \\
\phantomsection\label{entry:symbol-registry-06469}\(c\) & \texttt{\detokenize{c}} & new constant symbol and binary relation symbol of the worked example symbol family member \(c\) & \hyperref[def:p01-w02-p01-003-add-constant]{\ref*{def:p01-w02-p01-003-add-constant}} & n/a & \hyperref[def:p01-w02-p01-003-add-constant]{\ref*{def:p01-w02-p01-003-add-constant}} & \thepage & Object-language signature symbols local to the example and the failure case. Explicit displayed family; member types and common actual first-use exposed, not one atomic glyph. Family listing, not one atomic glyph. Author macro \texttt{\detokenize{c,}\allowbreak\detokenize{R}} not carried to the members. \\
\phantomsection\label{entry:symbol-registry-06470}\(R\) & \texttt{\detokenize{R}} & new constant symbol and binary relation symbol of the worked example symbol family member \(R\) & \hyperref[def:p01-w02-p01-003-add-constant]{\ref*{def:p01-w02-p01-003-add-constant}} & n/a & \hyperref[def:p01-w02-p01-003-add-constant]{\ref*{def:p01-w02-p01-003-add-constant}} & \thepage & Object-language signature symbols local to the example and the failure case. Explicit displayed family; member types and common actual first-use exposed, not one atomic glyph. Family listing, not one atomic glyph. Author macro \texttt{\detokenize{c,}\allowbreak\detokenize{R}} not carried to the members. \\
\phantomsection\label{entry:symbol-registry-06471}\(\mathcal N\) & \texttt{\detokenize{\mathcal N}} & two-element structure with empty relation & \hyperref[def:p01-w02-p01-003-add-constant]{\ref*{def:p01-w02-p01-003-add-constant}} & n/a & \hyperref[def:p01-w02-p01-003-add-constant]{\ref*{def:p01-w02-p01-003-add-constant}} & \thepage & Domain \(\{0,1\}\) and \(R^{\mathcal N}=\varnothing\); reused in the failure case. \\
\phantomsection\label{entry:symbol-registry-06472}\(\xi_k;\ \Xi\) & \texttt{\detokenize{\xi_k,}\allowbreak\detokenize{\Xi}} & schematic letters and their set & \hyperref[sec:p01-w02-p01-004-schemas-codes]{\ref*{sec:p01-w02-p01-004-schemas-codes}} & n/a & \hyperref[def:p01-w02-p01-004-schema]{\ref*{def:p01-w02-p01-004-schema}} & \thepage & Metalinguistic symbols of the schema syntax, distinct from every propositional-alphabet symbol; displays write \(\varphi,\psi,\chi\) for \(\xi_0,\xi_1,\xi_2\) in the schematic reading. \\
\phantomsection\label{entry:symbol-registry-06473}\(\mathfrak s\) & \texttt{\detokenize{\mathfrak s}} & formula-schema metavariables symbol family member \(\mathfrak s\) & \hyperref[sec:p01-w02-p01-004-schemas-codes]{\ref*{sec:p01-w02-p01-004-schemas-codes}} & n/a & \hyperref[def:p01-w02-p01-004-schema]{\ref*{def:p01-w02-p01-004-schema}} & \thepage & Range over \(\mathrm{Sch}\); chosen to avoid the registered sentence metavariables \(\sigma,\tau\). Explicit displayed family; member types and common actual first-use exposed, not one atomic glyph. Family listing, not one atomic glyph. Author macro \texttt{\detokenize{\mathfrak s,}\allowbreak\detokenize{\mathfrak t}} not carried to the members. \\
\phantomsection\label{entry:symbol-registry-06474}\(\mathfrak t\) & \texttt{\detokenize{\mathfrak t}} & formula-schema metavariables symbol family member \(\mathfrak t\) & \hyperref[sec:p01-w02-p01-004-schemas-codes]{\ref*{sec:p01-w02-p01-004-schemas-codes}} & n/a & \hyperref[def:p01-w02-p01-004-schema]{\ref*{def:p01-w02-p01-004-schema}} & \thepage & Range over \(\mathrm{Sch}\); chosen to avoid the registered sentence metavariables \(\sigma,\tau\). Explicit displayed family; member types and common actual first-use exposed, not one atomic glyph. Family listing, not one atomic glyph. Author macro \texttt{\detokenize{\mathfrak s,}\allowbreak\detokenize{\mathfrak t}} not carried to the members. \\
\phantomsection\label{entry:symbol-registry-06475}\(\mathrm{Sch}\) & \texttt{\detokenize{\mathrm{Sch}}} & set of propositional formula schemas & \hyperref[sec:p01-w02-p01-004-schemas-codes]{\ref*{sec:p01-w02-p01-004-schemas-codes}} & n/a & \hyperref[def:p01-w02-p01-004-schema]{\ref*{def:p01-w02-p01-004-schema}} & \thepage & Least set of strings over the propositional alphabet plus \(\Xi\) closed under the formation clauses. \\
\phantomsection\label{entry:symbol-registry-06476}\(\operatorname{MV}(\mathfrak s)\) & \texttt{\detokenize{\operatorname{MV}}} & schematic letters occurring in a schema & \hyperref[sec:p01-w02-p01-004-schemas-codes]{\ref*{sec:p01-w02-p01-004-schemas-codes}} & n/a & \hyperref[def:p01-w02-p01-004-schema]{\ref*{def:p01-w02-p01-004-schema}} & \thepage & Finite set. \\
\phantomsection\label{entry:symbol-registry-06477}\(\iota;\ \mathfrak s[\iota]\) & \texttt{\detokenize{\iota,}\allowbreak\detokenize{[\iota]}} & instantiation and instance of a schema & \hyperref[sec:p01-w02-p01-004-schemas-codes]{\ref*{sec:p01-w02-p01-004-schemas-codes}} & n/a & \hyperref[def:p01-w02-p01-004-instance]{\ref*{def:p01-w02-p01-004-instance}} & \thepage & \(\iota:\Xi\to\mathsf{Fm}_{\mathsf{CPL}}\); bracket postfix action by structural recursion, parallel to the registered uniform-substitution notation \(\alpha[\upsilon]\) of def:p02-uniform-substitution. \\
\phantomsection\label{entry:symbol-registry-06478}\(\operatorname{Inst}(\mathfrak s)\) & \texttt{\detokenize{\operatorname{Inst}}} & instance set of a schema & \hyperref[sec:p01-w02-p01-004-schemas-codes]{\ref*{sec:p01-w02-p01-004-schemas-codes}} & n/a & \hyperref[def:p01-w02-p01-004-instance]{\ref*{def:p01-w02-p01-004-instance}} & \thepage &  \\
\phantomsection\label{entry:symbol-registry-06479}\(\operatorname{match}(\mathfrak s,\varphi);\ \mathsf{fail}\) & \texttt{\detokenize{\operatorname{match},}\allowbreak\detokenize{\mathsf{fail}}} & matching procedure and its failure answer & \hyperref[sec:p01-w02-p01-004-schemas-codes]{\ref*{sec:p01-w02-p01-004-schemas-codes}} & n/a & \hyperref[def:p01-w02-p01-004-matching]{\ref*{def:p01-w02-p01-004-matching}} & \thepage &  \\
\phantomsection\label{entry:symbol-registry-06480}\(\operatorname{tr}\) & \texttt{\detokenize{\operatorname{tr}}} & translation of schemas into formulas & \hyperref[sec:p01-w02-p01-004-schemas-codes]{\ref*{sec:p01-w02-p01-004-schemas-codes}} & n/a & \hyperref[sec:p01-w02-p01-004-schemas]{\ref*{sec:p01-w02-p01-004-schemas}} & \thepage & Tokenwise \(p_n\mapsto p_{2n}\), \(\xi_k\mapsto p_{2k+1}\). \\
\phantomsection\label{entry:symbol-registry-06481}\(\upsilon_\iota\) & \texttt{\detokenize{\upsilon_\iota}} & uniform substitution induced by an instantiation & \hyperref[thm:p01-w02-p01-004-substitution-link]{\ref*{thm:p01-w02-p01-004-substitution-link}} & n/a & \hyperref[thm:p01-w02-p01-004-substitution-link]{\ref*{thm:p01-w02-p01-004-substitution-link}} & \thepage & Instance of the uniform-substitution letter \(\upsilon\) of def:p02-uniform-substitution. \\
\phantomsection\label{entry:symbol-registry-06482}\(\#\varphi\) & \texttt{\detokenize{\#}} & numerical code of a propositional formula & \hyperref[sec:p01-w02-p01-004-schemas-codes]{\ref*{sec:p01-w02-p01-004-schemas-codes}} & n/a & \hyperref[def:p01-w02-p01-004-code]{\ref*{def:p01-w02-p01-004-code}} & \thepage & Distinct from the arithmetic code \(\ulcorner\cdot\urcorner\) of sec:arithmetic-coding-04; same pairing function, different tags and sorts. \\
\phantomsection\label{entry:symbol-registry-06483}\(\pi;\ \pi^{-1}\) & \texttt{\detokenize{\pi,}\allowbreak\detokenize{\pi^{-1}}} & local abbreviation for the numerical pairing and its inverse & \hyperref[sec:p01-w02-p01-004-schemas-codes]{\ref*{sec:p01-w02-p01-004-schemas-codes}} & n/a & \hyperref[def:p01-w02-p01-004-code]{\ref*{def:p01-w02-p01-004-code}} & \thepage & \(\pi=\operatorname{pair}_{\N}\) of sec:strings-01; local only. \\
\phantomsection\label{entry:symbol-registry-06484}\(t_s\) & \texttt{\detokenize{t_s}} & triangular number & \hyperref[sec:p01-w02-p01-004-schemas-codes]{\ref*{sec:p01-w02-p01-004-schemas-codes}} & n/a & \hyperref[def:p01-w02-p01-004-decoding]{\ref*{def:p01-w02-p01-004-decoding}} & \thepage & \(t_s=s(s+1)/2\), as in the bijectivity proof of sec:strings-01. \\
\phantomsection\label{entry:symbol-registry-06485}\(\operatorname{dec};\ \mathsf{none}\) & \texttt{\detokenize{\operatorname{dec},}\allowbreak\detokenize{\mathsf{none}}} & decoding procedure and its non-code answer & \hyperref[sec:p01-w02-p01-004-schemas-codes]{\ref*{sec:p01-w02-p01-004-schemas-codes}} & n/a & \hyperref[def:p01-w02-p01-004-decoding]{\ref*{def:p01-w02-p01-004-decoding}} & \thepage &  \\
\phantomsection\label{entry:symbol-registry-06486}\(\circ_k\) & \texttt{\detokenize{\circ_k}} & binary connective with code tag k & \hyperref[def:p01-w02-p01-004-decoding]{\ref*{def:p01-w02-p01-004-decoding}} & n/a & \hyperref[def:p01-w02-p01-004-decoding]{\ref*{def:p01-w02-p01-004-decoding}} & \thepage & \(\circ_3=\land,\circ_4=\lor,\circ_5=\to\). \\
\phantomsection\label{entry:symbol-registry-06487}\(\lambda;\ \mathfrak s^{\lambda}\) & \texttt{\detokenize{\lambda,}\allowbreak\detokenize{^{\lambda}}} & letter renaming and renamed schema & \hyperref[annex:p01-w02-p01-004-distinction]{\ref*{annex:p01-w02-p01-004-distinction}} & n/a & \hyperref[annex:p01-w02-p01-004-distinction]{\ref*{annex:p01-w02-p01-004-distinction}} & \thepage & \(\lambda:\Xi\to\Xi\); annex-local. \\
\phantomsection\label{entry:symbol-registry-06488}\(\operatorname{tk}\) & \texttt{\detokenize{\operatorname{tk}}} & token map underlying tr & \hyperref[sec:annex-p01-w02-p01-004-schemas-codes]{\ref*{sec:annex-p01-w02-p01-004-schemas-codes}} & n/a & \hyperref[sec:annex-p01-w02-p01-004-schemas-codes]{\ref*{sec:annex-p01-w02-p01-004-schemas-codes}} & \thepage & Annex-local; defined in the unlabelled preamble of the annex section before annex:p01-w02-p01-004-schema-readability; sends \(p_n\) to \(p_{2n}\) and \(\xi_k\) to \(p_{2k+1}\), fixes every other token. \\
\phantomsection\label{entry:symbol-registry-06489}\(\epsilon;\ i\cdot a\) & \texttt{\detokenize{\epsilon,}\allowbreak\detokenize{\cdot}} & empty address and address with a prefixed index & \hyperref[sec:p01-w02-p01-024-binding]{\ref*{sec:p01-w02-p01-024-binding}} & n/a & \hyperref[def:p01-w02-p01-024-addresses]{\ref*{def:p01-w02-p01-024-addresses}} & \thepage & Addresses are finite sequences of natural numbers; the same empty-address role as the P2 occurrence-address row, in a disjoint scope. \\
\phantomsection\label{entry:symbol-registry-06490}\(b\sqsubseteq a\) & \texttt{\detokenize{\sqsubseteq}} & prefix order on addresses & \hyperref[sec:p01-w02-p01-024-binding]{\ref*{sec:p01-w02-p01-024-binding}} & meta & \hyperref[def:p01-w02-p01-024-addresses]{\ref*{def:p01-w02-p01-024-addresses}} & \thepage & Metalinguistic; unrelated to the description-logic concept inclusion and the scoped information orders that use the same glyph elsewhere. \\
\phantomsection\label{entry:symbol-registry-06491}\(\operatorname{VO}(e)\) & \texttt{\detokenize{\operatorname{VO}(e)}} & variable-occurrence addresses of a term or formula and the variable at an address symbol family member \(\operatorname{VO}(e)\) & \hyperref[sec:p01-w02-p01-024-binding]{\ref*{sec:p01-w02-p01-024-binding}} & n/a & \hyperref[def:p01-w02-p01-024-addresses]{\ref*{def:p01-w02-p01-024-addresses}} & \thepage & The address 0 of a quantified formula is the binding occurrence after the quantifier. Explicit displayed family; member types and common actual first-use exposed, not one atomic glyph. Family listing, not one atomic glyph. Author macro \texttt{\detokenize{\operatorname{VO},}\allowbreak\detokenize{e[a]}} not carried to the members. \\
\phantomsection\label{entry:symbol-registry-06492}\(\ e[a]\) & \texttt{\detokenize{\ e[a]}} & variable-occurrence addresses of a term or formula and the variable at an address symbol family member \(\ e[a]\) & \hyperref[sec:p01-w02-p01-024-binding]{\ref*{sec:p01-w02-p01-024-binding}} & n/a & \hyperref[def:p01-w02-p01-024-addresses]{\ref*{def:p01-w02-p01-024-addresses}} & \thepage & The address 0 of a quantified formula is the binding occurrence after the quantifier. Explicit displayed family; member types and common actual first-use exposed, not one atomic glyph. Family listing, not one atomic glyph. Author macro \texttt{\detokenize{\operatorname{VO},}\allowbreak\detokenize{e[a]}} not carried to the members. \\
\phantomsection\label{entry:symbol-registry-06493}\(\operatorname{QO}(\varphi)\) & \texttt{\detokenize{\operatorname{QO}(\varphi)}} & quantifier-occurrence addresses and the subformula occurrence at an address symbol family member \(\operatorname{QO}(\varphi)\) & \hyperref[sec:p01-w02-p01-024-binding]{\ref*{sec:p01-w02-p01-024-binding}} & n/a & \hyperref[def:p01-w02-p01-024-addresses]{\ref*{def:p01-w02-p01-024-addresses}} & \thepage & Braces select a subformula occurrence; square brackets select a variable occurrence. Explicit displayed family; member types and common actual first-use exposed, not one atomic glyph. Family listing, not one atomic glyph. Author macro \texttt{\detokenize{\operatorname{QO},}\allowbreak\detokenize{\varphi\{b\}}} not carried to the members. \\
\phantomsection\label{entry:symbol-registry-06494}\(\ \varphi\{b\}\) & \texttt{\detokenize{\ \varphi\{b\}}} & quantifier-occurrence addresses and the subformula occurrence at an address symbol family member \(\ \varphi\{b\}\) & \hyperref[sec:p01-w02-p01-024-binding]{\ref*{sec:p01-w02-p01-024-binding}} & n/a & \hyperref[def:p01-w02-p01-024-addresses]{\ref*{def:p01-w02-p01-024-addresses}} & \thepage & Braces select a subformula occurrence; square brackets select a variable occurrence. Explicit displayed family; member types and common actual first-use exposed, not one atomic glyph. Family listing, not one atomic glyph. Author macro \texttt{\detokenize{\operatorname{QO},}\allowbreak\detokenize{\varphi\{b\}}} not carried to the members. \\
\phantomsection\label{entry:symbol-registry-06495}\(\operatorname{OV}\) & \texttt{\detokenize{\operatorname{OV}}} & occurring and free-occurring variable sets and the set of quantifier occurrences that could bind an occurrence symbol family member \(\operatorname{OV}\) & \hyperref[annex:p01-w02-p01-024-fv-occurrence]{\ref*{annex:p01-w02-p01-024-fv-occurrence}} & n/a & \hyperref[annex:p01-w02-p01-024-fv-occurrence]{\ref*{annex:p01-w02-p01-024-fv-occurrence}} & \thepage & Proof-local names. Explicit displayed family; member types and common actual first-use exposed, not one atomic glyph. Family listing, not one atomic glyph. Author macro \texttt{\detokenize{\operatorname{OV},}\allowbreak\detokenize{\operatorname{FO},}\allowbreak\detokenize{P\_\varphi}} not carried to the members. \\
\phantomsection\label{entry:symbol-registry-06496}\(\operatorname{FO}\) & \texttt{\detokenize{\operatorname{FO}}} & occurring and free-occurring variable sets and the set of quantifier occurrences that could bind an occurrence symbol family member \(\operatorname{FO}\) & \hyperref[annex:p01-w02-p01-024-fv-occurrence]{\ref*{annex:p01-w02-p01-024-fv-occurrence}} & n/a & \hyperref[annex:p01-w02-p01-024-fv-occurrence]{\ref*{annex:p01-w02-p01-024-fv-occurrence}} & \thepage & Proof-local names. Explicit displayed family; member types and common actual first-use exposed, not one atomic glyph. Family listing, not one atomic glyph. Author macro \texttt{\detokenize{\operatorname{OV},}\allowbreak\detokenize{\operatorname{FO},}\allowbreak\detokenize{P\_\varphi}} not carried to the members. \\
\phantomsection\label{entry:symbol-registry-06497}\(P_\varphi(a)\) & \texttt{\detokenize{P_\varphi(a)}} & occurring and free-occurring variable sets and the set of quantifier occurrences that could bind an occurrence symbol family member \(P_\varphi(a)\) & \hyperref[annex:p01-w02-p01-024-fv-occurrence]{\ref*{annex:p01-w02-p01-024-fv-occurrence}} & n/a & \hyperref[annex:p01-w02-p01-024-fv-occurrence]{\ref*{annex:p01-w02-p01-024-fv-occurrence}} & \thepage & Proof-local names. Explicit displayed family; member types and common actual first-use exposed, not one atomic glyph. Family listing, not one atomic glyph. Author macro \texttt{\detokenize{\operatorname{OV},}\allowbreak\detokenize{\operatorname{FO},}\allowbreak\detokenize{P\_\varphi}} not carried to the members. \\
\phantomsection\label{entry:symbol-registry-06498}\(\operatorname{Diag}^+(\mathcal M)\) & \texttt{\detokenize{\operatorname{Diag}^+}} & positive diagram & \hyperref[sec:p01-w02-p01-024-diagrams]{\ref*{sec:p01-w02-p01-024-diagrams}} & n/a & \hyperref[def:p01-w02-p01-024-positive-diagram]{\ref*{def:p01-w02-p01-024-positive-diagram}} & \thepage & Atomic L(M)-sentences true in the parameter expansion; the atomic diagram Diag(M) adds every true negated atomic sentence. \\
\phantomsection\label{entry:symbol-registry-06499}\(\operatorname{NAND}\) & \texttt{\detokenize{\operatorname{NAND}}} & binary connective symbols interpreted by rows 1110 and 1000 symbol family member \(\operatorname{NAND}\) & \hyperref[sec:p01-w02-p01-024-post]{\ref*{sec:p01-w02-p01-024-post}} & n/a & \hyperref[sec:p01-w02-p01-024-post]{\ref*{sec:p01-w02-p01-024-post}} & \thepage & Connectives of the formula sets studied in this section; the P1 registry section and the P2 row use the same names as eliminable metalinguistic abbreviations in their own scopes. Explicit displayed family; member types and common actual first-use exposed, not one atomic glyph. Family listing, not one atomic glyph. Author macro \texttt{\detokenize{\operatorname{NAND},}\allowbreak\detokenize{\operatorname{NOR}}} not carried to the members. \\
\phantomsection\label{entry:symbol-registry-06500}\(\operatorname{NOR}\) & \texttt{\detokenize{\operatorname{NOR}}} & binary connective symbols interpreted by rows 1110 and 1000 symbol family member \(\operatorname{NOR}\) & \hyperref[sec:p01-w02-p01-024-post]{\ref*{sec:p01-w02-p01-024-post}} & n/a & \hyperref[sec:p01-w02-p01-024-post]{\ref*{sec:p01-w02-p01-024-post}} & \thepage & Connectives of the formula sets studied in this section; the P1 registry section and the P2 row use the same names as eliminable metalinguistic abbreviations in their own scopes. Explicit displayed family; member types and common actual first-use exposed, not one atomic glyph. Family listing, not one atomic glyph. Author macro \texttt{\detokenize{\operatorname{NAND},}\allowbreak\detokenize{\operatorname{NOR}}} not carried to the members. \\
\phantomsection\label{entry:symbol-registry-06501}\(O\) & \texttt{\detokenize{O}} & positive-arity Boolean operations of all arities or of arity n and the projections symbol family member \(O\) & \hyperref[sec:p01-w02-p01-024-post]{\ref*{sec:p01-w02-p01-024-post}} & n/a & \hyperref[def:p01-w02-p01-024-clones]{\ref*{def:p01-w02-p01-024-clones}} & \thepage & Arity is at least 1 throughout. Explicit displayed family; member types and common actual first-use exposed, not one atomic glyph. Family listing, not one atomic glyph. Author macro \texttt{\detokenize{O,}\allowbreak\detokenize{O\_n,}\allowbreak\detokenize{e\^{}n\_i}} not carried to the members. \\
\phantomsection\label{entry:symbol-registry-06502}\(O_n\) & \texttt{\detokenize{O_n}} & positive-arity Boolean operations of all arities or of arity n and the projections symbol family member \(O_n\) & \hyperref[sec:p01-w02-p01-024-post]{\ref*{sec:p01-w02-p01-024-post}} & n/a & \hyperref[def:p01-w02-p01-024-clones]{\ref*{def:p01-w02-p01-024-clones}} & \thepage & Arity is at least 1 throughout. Explicit displayed family; member types and common actual first-use exposed, not one atomic glyph. Family listing, not one atomic glyph. Author macro \texttt{\detokenize{O,}\allowbreak\detokenize{O\_n,}\allowbreak\detokenize{e\^{}n\_i}} not carried to the members. \\
\phantomsection\label{entry:symbol-registry-06503}\(e^n_i\) & \texttt{\detokenize{e^n_i}} & positive-arity Boolean operations of all arities or of arity n and the projections symbol family member \(e^n_i\) & \hyperref[sec:p01-w02-p01-024-post]{\ref*{sec:p01-w02-p01-024-post}} & n/a & \hyperref[def:p01-w02-p01-024-clones]{\ref*{def:p01-w02-p01-024-clones}} & \thepage & Arity is at least 1 throughout. Explicit displayed family; member types and common actual first-use exposed, not one atomic glyph. Family listing, not one atomic glyph. Author macro \texttt{\detokenize{O,}\allowbreak\detokenize{O\_n,}\allowbreak\detokenize{e\^{}n\_i}} not carried to the members. \\
\phantomsection\label{entry:symbol-registry-06504}\(\langle C\rangle\) & \texttt{\detokenize{\langle C\rangle}} & clone generated by a set of positive-arity truth functions & \hyperref[sec:p01-w02-p01-024-post]{\ref*{sec:p01-w02-p01-024-post}} & n/a & \hyperref[def:p01-w02-p01-024-clones]{\ref*{def:p01-w02-p01-024-clones}} & \thepage & Same role as the P2 prop-syntax-semantics row, in a disjoint scope. \\
\phantomsection\label{entry:symbol-registry-06505}\(\mathsf T_0\) & \texttt{\detokenize{\mathsf T_0}} & Post maximal clones symbol family member \(\mathsf T_0\) & \hyperref[sec:p01-w02-p01-024-post]{\ref*{sec:p01-w02-p01-024-post}} & n/a & \hyperref[def:p01-w02-p01-024-clones]{\ref*{def:p01-w02-p01-024-clones}} & \thepage & 0-preserving, 1-preserving, self-dual, monotone and affine functions; same role as the P2 prop-syntax-semantics row, in a disjoint scope. Explicit displayed family; member types and common actual first-use exposed, not one atomic glyph. Family listing, not one atomic glyph. Author macro \texttt{\detokenize{\mathsf T\_0,}\allowbreak\detokenize{\mathsf T\_1,}\allowbreak\detokenize{\mathsf S,}\allowbreak\detokenize{\mathsf M,}\allowbreak\detokenize{\mathsf L}} not carried to the members. \\
\phantomsection\label{entry:symbol-registry-06506}\(\mathsf T_1\) & \texttt{\detokenize{\mathsf T_1}} & Post maximal clones symbol family member \(\mathsf T_1\) & \hyperref[sec:p01-w02-p01-024-post]{\ref*{sec:p01-w02-p01-024-post}} & n/a & \hyperref[def:p01-w02-p01-024-clones]{\ref*{def:p01-w02-p01-024-clones}} & \thepage & 0-preserving, 1-preserving, self-dual, monotone and affine functions; same role as the P2 prop-syntax-semantics row, in a disjoint scope. Explicit displayed family; member types and common actual first-use exposed, not one atomic glyph. Family listing, not one atomic glyph. Author macro \texttt{\detokenize{\mathsf T\_0,}\allowbreak\detokenize{\mathsf T\_1,}\allowbreak\detokenize{\mathsf S,}\allowbreak\detokenize{\mathsf M,}\allowbreak\detokenize{\mathsf L}} not carried to the members. \\
\phantomsection\label{entry:symbol-registry-06507}\(\mathsf S\) & \texttt{\detokenize{\mathsf S}} & Post maximal clones symbol family member \(\mathsf S\) & \hyperref[sec:p01-w02-p01-024-post]{\ref*{sec:p01-w02-p01-024-post}} & n/a & \hyperref[def:p01-w02-p01-024-clones]{\ref*{def:p01-w02-p01-024-clones}} & \thepage & 0-preserving, 1-preserving, self-dual, monotone and affine functions; same role as the P2 prop-syntax-semantics row, in a disjoint scope. Explicit displayed family; member types and common actual first-use exposed, not one atomic glyph. Family listing, not one atomic glyph. Author macro \texttt{\detokenize{\mathsf T\_0,}\allowbreak\detokenize{\mathsf T\_1,}\allowbreak\detokenize{\mathsf S,}\allowbreak\detokenize{\mathsf M,}\allowbreak\detokenize{\mathsf L}} not carried to the members. \\
\phantomsection\label{entry:symbol-registry-06508}\(\mathsf M\) & \texttt{\detokenize{\mathsf M}} & Post maximal clones symbol family member \(\mathsf M\) & \hyperref[sec:p01-w02-p01-024-post]{\ref*{sec:p01-w02-p01-024-post}} & n/a & \hyperref[def:p01-w02-p01-024-clones]{\ref*{def:p01-w02-p01-024-clones}} & \thepage & 0-preserving, 1-preserving, self-dual, monotone and affine functions; same role as the P2 prop-syntax-semantics row, in a disjoint scope. Explicit displayed family; member types and common actual first-use exposed, not one atomic glyph. Family listing, not one atomic glyph. Author macro \texttt{\detokenize{\mathsf T\_0,}\allowbreak\detokenize{\mathsf T\_1,}\allowbreak\detokenize{\mathsf S,}\allowbreak\detokenize{\mathsf M,}\allowbreak\detokenize{\mathsf L}} not carried to the members. \\
\phantomsection\label{entry:symbol-registry-06509}\(\mathsf L\) & \texttt{\detokenize{\mathsf L}} & Post maximal clones symbol family member \(\mathsf L\) & \hyperref[sec:p01-w02-p01-024-post]{\ref*{sec:p01-w02-p01-024-post}} & n/a & \hyperref[def:p01-w02-p01-024-clones]{\ref*{def:p01-w02-p01-024-clones}} & \thepage & 0-preserving, 1-preserving, self-dual, monotone and affine functions; same role as the P2 prop-syntax-semantics row, in a disjoint scope. Explicit displayed family; member types and common actual first-use exposed, not one atomic glyph. Family listing, not one atomic glyph. Author macro \texttt{\detokenize{\mathsf T\_0,}\allowbreak\detokenize{\mathsf T\_1,}\allowbreak\detokenize{\mathsf S,}\allowbreak\detokenize{\mathsf M,}\allowbreak\detokenize{\mathsf L}} not carried to the members. \\
\phantomsection\label{entry:symbol-registry-06510}\(\bar a^{\,c}\) & \texttt{\detokenize{\bar a^{\,c}}} & coordinatewise complement constant tuples and change of one coordinate symbol family member \(\bar a^{\,c}\) & \hyperref[sec:p01-w02-p01-024-post]{\ref*{sec:p01-w02-p01-024-post}} & n/a & \hyperref[def:p01-w02-p01-024-clones]{\ref*{def:p01-w02-p01-024-clones}} & \thepage & The sign \(\oplus\) is addition modulo 2 here, not linear-logic plus. Explicit displayed family; member types and common actual first-use exposed, not one atomic glyph. Family listing, not one atomic glyph. Author macro \texttt{\detokenize{\bar a\^{}{\,c},}\allowbreak\detokenize{\bar0,}\allowbreak\detokenize{\bar1,}\allowbreak\detokenize{\oplus}} not carried to the members. \\
\phantomsection\label{entry:symbol-registry-06511}\(\ \bar0\) & \texttt{\detokenize{\ \bar0}} & coordinatewise complement constant tuples and change of one coordinate symbol family member \(\ \bar0\) & \hyperref[sec:p01-w02-p01-024-post]{\ref*{sec:p01-w02-p01-024-post}} & n/a & \hyperref[def:p01-w02-p01-024-clones]{\ref*{def:p01-w02-p01-024-clones}} & \thepage & The sign \(\oplus\) is addition modulo 2 here, not linear-logic plus. Explicit displayed family; member types and common actual first-use exposed, not one atomic glyph. Family listing, not one atomic glyph. Author macro \texttt{\detokenize{\bar a\^{}{\,c},}\allowbreak\detokenize{\bar0,}\allowbreak\detokenize{\bar1,}\allowbreak\detokenize{\oplus}} not carried to the members. \\
\phantomsection\label{entry:symbol-registry-06512}\(\bar1\) & \texttt{\detokenize{\bar1}} & coordinatewise complement constant tuples and change of one coordinate symbol family member \(\bar1\) & \hyperref[sec:p01-w02-p01-024-post]{\ref*{sec:p01-w02-p01-024-post}} & n/a & \hyperref[def:p01-w02-p01-024-clones]{\ref*{def:p01-w02-p01-024-clones}} & \thepage & The sign \(\oplus\) is addition modulo 2 here, not linear-logic plus. Explicit displayed family; member types and common actual first-use exposed, not one atomic glyph. Family listing, not one atomic glyph. Author macro \texttt{\detokenize{\bar a\^{}{\,c},}\allowbreak\detokenize{\bar0,}\allowbreak\detokenize{\bar1,}\allowbreak\detokenize{\oplus}} not carried to the members. \\
\phantomsection\label{entry:symbol-registry-06513}\(\ \bar a\oplus e_i\) & \texttt{\detokenize{\ \bar a\oplus e_i}} & coordinatewise complement constant tuples and change of one coordinate symbol family member \(\ \bar a\oplus e_i\) & \hyperref[sec:p01-w02-p01-024-post]{\ref*{sec:p01-w02-p01-024-post}} & n/a & \hyperref[def:p01-w02-p01-024-clones]{\ref*{def:p01-w02-p01-024-clones}} & \thepage & The sign \(\oplus\) is addition modulo 2 here, not linear-logic plus. Explicit displayed family; member types and common actual first-use exposed, not one atomic glyph. Family listing, not one atomic glyph. Author macro \texttt{\detokenize{\bar a\^{}{\,c},}\allowbreak\detokenize{\bar0,}\allowbreak\detokenize{\bar1,}\allowbreak\detokenize{\oplus}} not carried to the members. \\
\phantomsection\label{entry:symbol-registry-06514}\(k_0\) & \texttt{\detokenize{k_0}} & unary constant functions ternary majority and the unary identity symbol family member \(k_0\) & \hyperref[annex:p01-w02-p01-024-post]{\ref*{annex:p01-w02-p01-024-post}} & n/a & \hyperref[annex:p01-w02-p01-024-post]{\ref*{annex:p01-w02-p01-024-post}} & \thepage & Proof-local witnesses. Explicit displayed family; member types and common actual first-use exposed, not one atomic glyph. Family listing, not one atomic glyph. Author macro \texttt{\detokenize{k\_0,}\allowbreak\detokenize{k\_1,}\allowbreak\detokenize{\operatorname{maj},}\allowbreak\detokenize{\mathrm{id}}} not carried to the members. \\
\phantomsection\label{entry:symbol-registry-06515}\(k_1\) & \texttt{\detokenize{k_1}} & unary constant functions ternary majority and the unary identity symbol family member \(k_1\) & \hyperref[annex:p01-w02-p01-024-post]{\ref*{annex:p01-w02-p01-024-post}} & n/a & \hyperref[annex:p01-w02-p01-024-post]{\ref*{annex:p01-w02-p01-024-post}} & \thepage & Proof-local witnesses. Explicit displayed family; member types and common actual first-use exposed, not one atomic glyph. Family listing, not one atomic glyph. Author macro \texttt{\detokenize{k\_0,}\allowbreak\detokenize{k\_1,}\allowbreak\detokenize{\operatorname{maj},}\allowbreak\detokenize{\mathrm{id}}} not carried to the members. \\
\phantomsection\label{entry:symbol-registry-06516}\(\operatorname{maj}\) & \texttt{\detokenize{\operatorname{maj}}} & unary constant functions ternary majority and the unary identity symbol family member \(\operatorname{maj}\) & \hyperref[annex:p01-w02-p01-024-post]{\ref*{annex:p01-w02-p01-024-post}} & n/a & \hyperref[annex:p01-w02-p01-024-post]{\ref*{annex:p01-w02-p01-024-post}} & \thepage & Proof-local witnesses. Explicit displayed family; member types and common actual first-use exposed, not one atomic glyph. Family listing, not one atomic glyph. Author macro \texttt{\detokenize{k\_0,}\allowbreak\detokenize{k\_1,}\allowbreak\detokenize{\operatorname{maj},}\allowbreak\detokenize{\mathrm{id}}} not carried to the members. \\
\phantomsection\label{entry:symbol-registry-06517}\(\mathrm{id}\) & \texttt{\detokenize{\mathrm{id}}} & unary constant functions ternary majority and the unary identity symbol family member \(\mathrm{id}\) & \hyperref[annex:p01-w02-p01-024-post]{\ref*{annex:p01-w02-p01-024-post}} & n/a & \hyperref[annex:p01-w02-p01-024-post]{\ref*{annex:p01-w02-p01-024-post}} & \thepage & Proof-local witnesses. Explicit displayed family; member types and common actual first-use exposed, not one atomic glyph. Family listing, not one atomic glyph. Author macro \texttt{\detokenize{k\_0,}\allowbreak\detokenize{k\_1,}\allowbreak\detokenize{\operatorname{maj},}\allowbreak\detokenize{\mathrm{id}}} not carried to the members. \\
\phantomsection\label{entry:symbol-registry-06518}\(C_1\) & \texttt{\detokenize{C_1}} & Post's names for the complete system and the five systems immediately included in it symbol family member \(C_1\) & \hyperref[sec:p01-w02-p01-024-post]{\ref*{sec:p01-w02-p01-024-post}} & n/a & \hyperref[sec:p01-w02-p01-024-post]{\ref*{sec:p01-w02-p01-024-post}} & \thepage & Source notation of Post 1941, sections 25 and 26. Explicit displayed family; member types and common actual first-use exposed, not one atomic glyph. Family listing, not one atomic glyph. Author macro \texttt{\detokenize{C\_1,}\allowbreak\detokenize{C\_2,}\allowbreak\detokenize{C\_3,}\allowbreak\detokenize{D\_3,}\allowbreak\detokenize{A\_1,}\allowbreak\detokenize{L\_1}} not carried to the members. \\
\phantomsection\label{entry:symbol-registry-06519}\(C_2\) & \texttt{\detokenize{C_2}} & Post's names for the complete system and the five systems immediately included in it symbol family member \(C_2\) & \hyperref[sec:p01-w02-p01-024-post]{\ref*{sec:p01-w02-p01-024-post}} & n/a & \hyperref[sec:p01-w02-p01-024-post]{\ref*{sec:p01-w02-p01-024-post}} & \thepage & Source notation of Post 1941, sections 25 and 26. Explicit displayed family; member types and common actual first-use exposed, not one atomic glyph. Family listing, not one atomic glyph. Author macro \texttt{\detokenize{C\_1,}\allowbreak\detokenize{C\_2,}\allowbreak\detokenize{C\_3,}\allowbreak\detokenize{D\_3,}\allowbreak\detokenize{A\_1,}\allowbreak\detokenize{L\_1}} not carried to the members. \\
\phantomsection\label{entry:symbol-registry-06520}\(C_3\) & \texttt{\detokenize{C_3}} & Post's names for the complete system and the five systems immediately included in it symbol family member \(C_3\) & \hyperref[sec:p01-w02-p01-024-post]{\ref*{sec:p01-w02-p01-024-post}} & n/a & \hyperref[sec:p01-w02-p01-024-post]{\ref*{sec:p01-w02-p01-024-post}} & \thepage & Source notation of Post 1941, sections 25 and 26. Explicit displayed family; member types and common actual first-use exposed, not one atomic glyph. Family listing, not one atomic glyph. Author macro \texttt{\detokenize{C\_1,}\allowbreak\detokenize{C\_2,}\allowbreak\detokenize{C\_3,}\allowbreak\detokenize{D\_3,}\allowbreak\detokenize{A\_1,}\allowbreak\detokenize{L\_1}} not carried to the members. \\
\phantomsection\label{entry:symbol-registry-06521}\(D_3\) & \texttt{\detokenize{D_3}} & Post's names for the complete system and the five systems immediately included in it symbol family member \(D_3\) & \hyperref[sec:p01-w02-p01-024-post]{\ref*{sec:p01-w02-p01-024-post}} & n/a & \hyperref[sec:p01-w02-p01-024-post]{\ref*{sec:p01-w02-p01-024-post}} & \thepage & Source notation of Post 1941, sections 25 and 26. Explicit displayed family; member types and common actual first-use exposed, not one atomic glyph. Family listing, not one atomic glyph. Author macro \texttt{\detokenize{C\_1,}\allowbreak\detokenize{C\_2,}\allowbreak\detokenize{C\_3,}\allowbreak\detokenize{D\_3,}\allowbreak\detokenize{A\_1,}\allowbreak\detokenize{L\_1}} not carried to the members. \\
\phantomsection\label{entry:symbol-registry-06522}\(A_1\) & \texttt{\detokenize{A_1}} & Post's names for the complete system and the five systems immediately included in it symbol family member \(A_1\) & \hyperref[sec:p01-w02-p01-024-post]{\ref*{sec:p01-w02-p01-024-post}} & n/a & \hyperref[sec:p01-w02-p01-024-post]{\ref*{sec:p01-w02-p01-024-post}} & \thepage & Source notation of Post 1941, sections 25 and 26. Explicit displayed family; member types and common actual first-use exposed, not one atomic glyph. Family listing, not one atomic glyph. Author macro \texttt{\detokenize{C\_1,}\allowbreak\detokenize{C\_2,}\allowbreak\detokenize{C\_3,}\allowbreak\detokenize{D\_3,}\allowbreak\detokenize{A\_1,}\allowbreak\detokenize{L\_1}} not carried to the members. \\
\phantomsection\label{entry:symbol-registry-06523}\(L_1\) & \texttt{\detokenize{L_1}} & Post's names for the complete system and the five systems immediately included in it symbol family member \(L_1\) & \hyperref[sec:p01-w02-p01-024-post]{\ref*{sec:p01-w02-p01-024-post}} & n/a & \hyperref[sec:p01-w02-p01-024-post]{\ref*{sec:p01-w02-p01-024-post}} & \thepage & Source notation of Post 1941, sections 25 and 26. Explicit displayed family; member types and common actual first-use exposed, not one atomic glyph. Family listing, not one atomic glyph. Author macro \texttt{\detokenize{C\_1,}\allowbreak\detokenize{C\_2,}\allowbreak\detokenize{C\_3,}\allowbreak\detokenize{D\_3,}\allowbreak\detokenize{A\_1,}\allowbreak\detokenize{L\_1}} not carried to the members. \\
\phantomsection\label{entry:symbol-registry-06524}\(\alpha\) & \texttt{\detokenize{\alpha}} & Post's class of a two-valued function whose associated first-order function (all arguments identified) is the identity & \hyperref[sec:p01-w02-p01-024-post]{\ref*{sec:p01-w02-p01-024-post}} & n/a & \hyperref[sec:p01-w02-p01-024-post]{\ref*{sec:p01-w02-p01-024-post}} & \thepage & Source notation of Post 1941, section 10, written as a prefix to '-function' (as in \(\alpha\)-function); disjoint from the ordinal and type metavariable scopes. \\
\phantomsection\label{entry:symbol-registry-06525}\(\beta\) & \texttt{\detokenize{\beta}} & Post's class of a two-valued function whose associated first-order function (all arguments identified) is the constant 1 & \hyperref[sec:p01-w02-p01-024-post]{\ref*{sec:p01-w02-p01-024-post}} & n/a & \hyperref[sec:p01-w02-p01-024-post]{\ref*{sec:p01-w02-p01-024-post}} & \thepage & Source notation of Post 1941, section 10, written as a prefix to '-function' (as in \(\alpha\)-function); disjoint from the ordinal and type metavariable scopes. \\
\phantomsection\label{entry:symbol-registry-06526}\(\gamma\) & \texttt{\detokenize{\gamma}} & Post's class of a two-valued function whose associated first-order function (all arguments identified) is the constant 0 & \hyperref[sec:p01-w02-p01-024-post]{\ref*{sec:p01-w02-p01-024-post}} & n/a & \hyperref[sec:p01-w02-p01-024-post]{\ref*{sec:p01-w02-p01-024-post}} & \thepage & Source notation of Post 1941, section 10, written as a prefix to '-function' (as in \(\alpha\)-function); disjoint from the ordinal and type metavariable scopes. \\
\phantomsection\label{entry:symbol-registry-06527}\(\delta\) & \texttt{\detokenize{\delta}} & Post's class of a two-valued function whose associated first-order function (all arguments identified) is the negation & \hyperref[sec:p01-w02-p01-024-post]{\ref*{sec:p01-w02-p01-024-post}} & n/a & \hyperref[sec:p01-w02-p01-024-post]{\ref*{sec:p01-w02-p01-024-post}} & \thepage & Source notation of Post 1941, section 10, written as a prefix to '-function' (as in \(\alpha\)-function); disjoint from the ordinal and type metavariable scopes. \\
\phantomsection\label{entry:symbol-registry-06551}\(\alpha\) & \texttt{\detokenize{\alpha}} & ordinal and cardinal metavariables symbol family member \(\alpha\) & \hyperref[sec:set-prerequisites-05]{\ref*{sec:set-prerequisites-05}}--\hyperref[sec:set-prerequisites-06]{\ref*{sec:set-prerequisites-06}} & n/a & \hyperref[def:ordinal-cardinal]{\ref*{def:ordinal-cardinal}} & \thepage & Ordinals compared by membership and cardinals as initial ordinals; Part 2 compactness and Lowenheim-Skolem constructions bind alpha, xi, kappa, lambda, mu locally in the same roles. Explicit displayed family; member types and common actual first-use exposed, not one atomic glyph. Family listing, not one atomic glyph. Author macro \texttt{\detokenize{\alpha,}\allowbreak\detokenize{\beta,}\allowbreak\detokenize{\kappa}} not carried to the members. \\
\phantomsection\label{entry:symbol-registry-06552}\(\beta;\ \kappa\) & \texttt{\detokenize{\beta;\ \kappa}} & ordinal and cardinal metavariables symbol family member \(\beta;\ \kappa\) & \hyperref[sec:set-prerequisites-05]{\ref*{sec:set-prerequisites-05}}--\hyperref[sec:set-prerequisites-06]{\ref*{sec:set-prerequisites-06}} & n/a & \hyperref[def:ordinal-cardinal]{\ref*{def:ordinal-cardinal}} & \thepage & Ordinals compared by membership and cardinals as initial ordinals; Part 2 compactness and Lowenheim-Skolem constructions bind alpha, xi, kappa, lambda, mu locally in the same roles. Explicit displayed family; member types and common actual first-use exposed, not one atomic glyph. Family listing, not one atomic glyph. Author macro \texttt{\detokenize{\alpha,}\allowbreak\detokenize{\beta,}\allowbreak\detokenize{\kappa}} not carried to the members. \\
\phantomsection\label{entry:symbol-registry-06554}\(\alpha\equiv\beta\) & \texttt{\detokenize{\equiv}} & source metalinguistic equivalence sign joining lines of a displayed propositional calculation & \hyperref[sec:writing-foundations-05]{\ref*{sec:writing-foundations-05}} & meta & \hyperref[sec:writing-foundations-05]{\ref*{sec:writing-foundations-05}} & \thepage & Asserts logical equivalence of consecutive displayed CPL formulas; a sign of the metalanguage, with the object-language biconditional written \(\leftrightarrow\). Holds iff \(\alpha\leftrightarrow\beta\) is true under every valuation (thm:p01-w02-p01-025-equiv-bridges). \\
\phantomsection\label{entry:symbol-registry-06557}\(\mathrm{Eq}_{\mathcal L}\) & \texttt{\detokenize{\mathrm{Eq}_{\mathcal L}}} & set of reflexivity and function/relation congruence axioms of a signature & \hyperref[sec:p01-w02-p01-025-equality]{\ref*{sec:p01-w02-p01-025-equality}} & n/a & \hyperref[def:p01-w02-p01-025-eq-axioms]{\ref*{def:p01-w02-p01-025-eq-axioms}} & \thepage & Contains every choice of distinct bound variables, hence every alpha-variant of its members. \\
\phantomsection\label{entry:symbol-registry-06558}\(\mathsf{ND}^{-}_{\mathsf{FOL}}\) & \texttt{\detokenize{\mathsf{ND}^{-}_{\mathsf{FOL}}}} & natural deduction for first-order logic without the rules =I and =E & \hyperref[sec:p01-w02-p01-025-equality]{\ref*{sec:p01-w02-p01-025-equality}} & n/a & \hyperref[def:p01-w02-p01-025-eq-axioms]{\ref*{def:p01-w02-p01-025-eq-axioms}} & \thepage & Same language as ND\_FOL; the identity sign is an ordinary binary relation symbol with no rules. \\
\phantomsection\label{entry:symbol-registry-06559}\(e\sigma\) & \texttt{\detokenize{e\sigma}} & application of a finite substitution sequence from left to right & \hyperref[sec:annex-p01-w02-p01-025-equality]{\ref*{sec:annex-p01-w02-p01-025-equality}} & n/a & \hyperref[annex:p01-w02-p01-025-axiom-replacement]{\ref*{annex:p01-w02-p01-025-axiom-replacement}} & \thepage & Proof-local notation in the annex; \(\sigma\cdot[s/z]\) extends the sequence. \\
\phantomsection\label{entry:symbol-registry-06564}\(\mathrm{id}_A\) & \texttt{\detokenize{\mathrm{id}_A}} & identity relation on a set & \hyperref[sec:p01-w02-p01-008-orders]{\ref*{sec:p01-w02-p01-008-orders}} & n/a & \hyperref[def:p01-w02-p01-008-order-kinds]{\ref*{def:p01-w02-p01-008-order-kinds}} & \thepage & The set of pairs (a,a) with a in A, formed by separation from A x A; used to pass between strict and non-strict orders. \\
\phantomsection\label{entry:symbol-registry-06565}\(S(x)\) & \texttt{\detokenize{S}} & successor of a set & \hyperref[sec:p01-w02-p01-008-naturals]{\ref*{sec:p01-w02-p01-008-naturals}} & n/a & \hyperref[def:p01-w02-p01-008-inductive]{\ref*{def:p01-w02-p01-008-inductive}} & \thepage & x union {x}, the source notation of the definition Infinity in the ZF chapter, here in the metalanguage; for an ordinal it is alpha+1. \\
\phantomsection\label{entry:symbol-registry-06566}\(\omega\) & \texttt{\detokenize{\omega}} & least inductive set & \hyperref[sec:p01-w02-p01-008-naturals]{\ref*{sec:p01-w02-p01-008-naturals}} & n/a & \hyperref[thm:p01-w02-p01-008-omega]{\ref*{thm:p01-w02-p01-008-omega}} & \thepage & Metatheoretic least inductive set, identified with the set N of natural numbers; same construction as the object-language omega of sec:zf-zfc-01. \\
\phantomsection\label{entry:symbol-registry-06567}\(m<n\) & \texttt{\detokenize{m<n}} & order of natural numbers symbol family member \(m<n\) & \hyperref[sec:p01-w02-p01-008-naturals]{\ref*{sec:p01-w02-p01-008-naturals}} & 2 & \hyperref[def:p01-w02-p01-008-naturals]{\ref*{def:p01-w02-p01-008-naturals}} & \thepage & m < n iff m is a member of n; the ordinal order of def:ordinal-cardinal restricted to N. The global P5-P19 canonical-ordering row does not overlap this scope. Explicit displayed family; member types and common actual first-use exposed, not one atomic glyph. Family listing, not one atomic glyph. Author macro \texttt{\detokenize{<,}\allowbreak\detokenize{\leq}} not carried to the members. \\
\phantomsection\label{entry:symbol-registry-06568}\(\ m\leq n\) & \texttt{\detokenize{\ m\leq n}} & order of natural numbers symbol family member \(\ m\leq n\) & \hyperref[sec:p01-w02-p01-008-naturals]{\ref*{sec:p01-w02-p01-008-naturals}} & 2 & \hyperref[def:p01-w02-p01-008-naturals]{\ref*{def:p01-w02-p01-008-naturals}} & \thepage & m < n iff m is a member of n; the ordinal order of def:ordinal-cardinal restricted to N. The global P5-P19 canonical-ordering row does not overlap this scope. Explicit displayed family; member types and common actual first-use exposed, not one atomic glyph. Family listing, not one atomic glyph. Author macro \texttt{\detokenize{<,}\allowbreak\detokenize{\leq}} not carried to the members. \\
\phantomsection\label{entry:symbol-registry-06569}\(n-1\) & \texttt{\detokenize{-1}} & predecessor of a nonzero natural number & \hyperref[sec:p01-w02-p01-008-naturals]{\ref*{sec:p01-w02-p01-008-naturals}} & n/a & \hyperref[def:p01-w02-p01-008-naturals]{\ref*{def:p01-w02-p01-008-naturals}} & \thepage & The unique m with m+1=n; defined only for n different from 0. No subtraction operation is introduced. \\
\phantomsection\label{entry:symbol-registry-06570}\(\operatorname{seg}(w)\) & \texttt{\detokenize{\operatorname{seg}}} & initial segment of a well-order & \hyperref[sec:p01-w02-p01-008-representation]{\ref*{sec:p01-w02-p01-008-representation}} & n/a & \hyperref[def:p01-w02-p01-008-isomorphism]{\ref*{def:p01-w02-p01-008-isomorphism}} & \thepage & Same notation and role as def:p01-w02-p01-011-readings item 4. \\
\phantomsection\label{entry:symbol-registry-06571}\(f:(W,\prec)\cong(V,<)\) & \texttt{\detokenize{\cong}} & order isomorphism of well-orders & \hyperref[sec:p01-w02-p01-008-representation]{\ref*{sec:p01-w02-p01-008-representation}} & meta & \hyperref[def:p01-w02-p01-008-isomorphism]{\ref*{def:p01-w02-p01-008-isomorphism}} & \thepage & Instance of the global isomorphism row: an order-preserving and order-reflecting bijection. \\
\phantomsection\label{entry:symbol-registry-06572}\(B^{+}\) & \texttt{\detokenize{^{+}}} & choice successor of a subset & \hyperref[sec:p01-w02-p01-008-wot]{\ref*{sec:p01-w02-p01-008-wot}} & n/a & \hyperref[def:p01-w02-p01-008-tower-method]{\ref*{def:p01-w02-p01-008-tower-method}} & \thepage & B union {c(A minus B)} for B different from A, and A itself for B=A, relative to a fixed choice function c; local to the tower construction. \\
\phantomsection\label{entry:symbol-registry-06573}\(\mathcal T\) & \texttt{\detokenize{\mathcal T}} & towers and the least tower symbol family member \(\mathcal T\) & \hyperref[sec:p01-w02-p01-008-wot]{\ref*{sec:p01-w02-p01-008-wot}} & n/a & \hyperref[def:p01-w02-p01-008-tower-method]{\ref*{def:p01-w02-p01-008-tower-method}} & \thepage & Families of subsets of A closed under the empty set, B\^{}{+} and arbitrary unions; T\_x is the union of the members of T\_0 omitting x. Local to the tower construction and its annex proof. Explicit displayed family; member types and common actual first-use exposed, not one atomic glyph. Family listing, not one atomic glyph. Author macro \texttt{\detokenize{\mathcal T,}\allowbreak\detokenize{\mathcal T\_0,}\allowbreak\detokenize{T\_x}} not carried to the members. \\
\phantomsection\label{entry:symbol-registry-06574}\(\mathcal T_0\) & \texttt{\detokenize{\mathcal T_0}} & towers and the least tower symbol family member \(\mathcal T_0\) & \hyperref[sec:p01-w02-p01-008-wot]{\ref*{sec:p01-w02-p01-008-wot}} & n/a & \hyperref[def:p01-w02-p01-008-tower-method]{\ref*{def:p01-w02-p01-008-tower-method}} & \thepage & Families of subsets of A closed under the empty set, B\^{}{+} and arbitrary unions; T\_x is the union of the members of T\_0 omitting x. Local to the tower construction and its annex proof. Explicit displayed family; member types and common actual first-use exposed, not one atomic glyph. Family listing, not one atomic glyph. Author macro \texttt{\detokenize{\mathcal T,}\allowbreak\detokenize{\mathcal T\_0,}\allowbreak\detokenize{T\_x}} not carried to the members. \\
\phantomsection\label{entry:symbol-registry-06575}\(T_x\) & \texttt{\detokenize{T_x}} & towers and the least tower symbol family member \(T_x\) & \hyperref[sec:p01-w02-p01-008-wot]{\ref*{sec:p01-w02-p01-008-wot}} & n/a & \hyperref[def:p01-w02-p01-008-tower-method]{\ref*{def:p01-w02-p01-008-tower-method}} & \thepage & Families of subsets of A closed under the empty set, B\^{}{+} and arbitrary unions; T\_x is the union of the members of T\_0 omitting x. Local to the tower construction and its annex proof. Explicit displayed family; member types and common actual first-use exposed, not one atomic glyph. Family listing, not one atomic glyph. Author macro \texttt{\detokenize{\mathcal T,}\allowbreak\detokenize{\mathcal T\_0,}\allowbreak\detokenize{T\_x}} not carried to the members. \\
\phantomsection\label{entry:symbol-registry-06576}\(I\) & \texttt{\detokenize{I}} & index set whose subsets are filtered & \hyperref[sec:p01-w02-p01-009-ultrafilter-lemma]{\ref*{sec:p01-w02-p01-009-ultrafilter-lemma}} & n/a & \hyperref[sec:p01-w02-p01-009-ultrafilter-lemma]{\ref*{sec:p01-w02-p01-009-ultrafilter-lemma}} & \thepage & Source notation of def:ultrafilter-lemma; a metalanguage set, possibly empty. Local typed role distinct from the P8 agent-index and phase sets. \\
\phantomsection\label{entry:symbol-registry-06577}\(\mathcal F\) & \texttt{\detokenize{\mathcal F}} & family of subsets satisfying (F1)--(F4) & \hyperref[sec:p01-w02-p01-009-ultrafilter-lemma]{\ref*{sec:p01-w02-p01-009-ultrafilter-lemma}} & n/a & \hyperref[sec:p01-w02-p01-009-ultrafilter-lemma]{\ref*{sec:p01-w02-p01-009-ultrafilter-lemma}} & \thepage & Source notation of def:ultrafilter-lemma; a proper set filter on I by thm:p01-w02-p01-009-terms. Overlaps the global row for function-symbol sets; the local type governs in this section. \\
\phantomsection\label{entry:symbol-registry-06578}\(\mathcal U\) & \texttt{\detokenize{\mathcal U}} & deciding extension of a family of subsets & \hyperref[sec:p01-w02-p01-009-ultrafilter-lemma]{\ref*{sec:p01-w02-p01-009-ultrafilter-lemma}} & n/a & \hyperref[sec:p01-w02-p01-009-ultrafilter-lemma]{\ref*{sec:p01-w02-p01-009-ultrafilter-lemma}} & \thepage & Source notation of def:ultrafilter-lemma; an ultrafilter on I by thm:p01-w02-p01-009-terms. Same glyph as the P6 measurability ultrafilter row, without any completeness requirement. \\
\phantomsection\label{entry:symbol-registry-06579}\(\mathcal E\) & \texttt{\detokenize{\mathcal E}} & given family of subsets to be extended & \hyperref[sec:p01-w02-p01-009-side-conditions]{\ref*{sec:p01-w02-p01-009-side-conditions}} & n/a & \hyperref[thm:p01-w02-p01-009-side-conditions]{\ref*{thm:p01-w02-p01-009-side-conditions}} & \thepage & Arbitrary subfamily of the power set of I; no closure conditions assumed. \\
\phantomsection\label{entry:symbol-registry-06580}\(\mathcal E^{\uparrow}\) & \texttt{\detokenize{\mathcal E^{\uparrow}}} & family generated by a family of subsets & \hyperref[sec:p01-w02-p01-009-side-conditions]{\ref*{sec:p01-w02-p01-009-side-conditions}} & n/a & \hyperref[thm:p01-w02-p01-009-side-conditions]{\ref*{thm:p01-w02-p01-009-side-conditions}} & \thepage & Supersets of finite intersections of members of E, the empty intersection being I. \\
\phantomsection\label{entry:symbol-registry-06581}\(\mathcal W\) & \texttt{\detokenize{\mathcal W}} & auxiliary set filter adjoining one subset & \hyperref[sec:annex-p01-w02-p01-009-ultrafilter-lemma]{\ref*{sec:annex-p01-w02-p01-009-ultrafilter-lemma}} & n/a & \hyperref[annex:p01-w02-p01-009-terms]{\ref*{annex:p01-w02-p01-009-terms}} & \thepage & Proof-local; supersets of Z intersected with X for Z in U. \\
\phantomsection\label{entry:symbol-registry-06582}\(\mathcal V\) & \texttt{\detokenize{\mathcal V}} & arbitrary family satisfying (F1)--(F4) containing E & \hyperref[sec:annex-p01-w02-p01-009-ultrafilter-lemma]{\ref*{sec:annex-p01-w02-p01-009-ultrafilter-lemma}} & n/a & \hyperref[annex:p01-w02-p01-009-side-conditions]{\ref*{annex:p01-w02-p01-009-side-conditions}} & \thepage & Proof-local; used for the leastness of the generated family. \\
\phantomsection\label{entry:symbol-registry-06583}\(T\) & \texttt{\detokenize{T}} & the theory ZF plus BPI plus the negation of AC & \hyperref[annex:p01-w02-p01-009-lemma]{\ref*{annex:p01-w02-p01-009-lemma}} & n/a & \hyperref[annex:p01-w02-p01-009-lemma]{\ref*{annex:p01-w02-p01-009-lemma}} & \thepage & Proof-local name for a set-theoretic theory; not an object theory under study. \\
\phantomsection\label{entry:symbol-registry-06584}\(\mathcal P(I)\) & \texttt{\detokenize{\mathcal P}} & power set of the index set I & \hyperref[sec:p01-w02-p01-009-ultrafilter-lemma]{\ref*{sec:p01-w02-p01-009-ultrafilter-lemma}} & n/a & \hyperref[sec:p01-w02-p01-009-ultrafilter-lemma]{\ref*{sec:p01-w02-p01-009-ultrafilter-lemma}} & \thepage & Renders the source's \texttt{\detokenize{\Pow}}(I) of def:ultrafilter-lemma; handbook.sty restricts \texttt{\detokenize{\Pow}} to restoration captures, so the body uses \(\mathcal P(I)\), which typesets the same glyph with the same typed role, as in sec:p01-boolean-prime-ideal-and-ultrafilters. \\
\phantomsection\label{entry:symbol-registry-06585}\(x\in a\) & \texttt{\detokenize{\in}} & membership: the primitive binary relation symbol of \(\Lang_{\in}\) & \hyperref[sec:p01-w02-p01-005-set-operations]{\ref*{sec:p01-w02-p01-005-set-operations}} & n/a & \hyperref[def:p01-w02-p01-005-membership]{\ref*{def:p01-w02-p01-005-membership}} & \thepage & Undefined primitive; every other set-theoretic notion of the chapter is defined from it and identity. For a class \(C=\{x:\varphi(x,\bar p)\}\), \(x\in C\) abbreviates \(\varphi(x,\bar p)\). \\
\phantomsection\label{entry:symbol-registry-06586}\(x\notin a\) & \texttt{\detokenize{\notin}} & non-membership & \hyperref[sec:p01-w02-p01-005-set-operations]{\ref*{sec:p01-w02-p01-005-set-operations}} & n/a & \hyperref[def:p01-w02-p01-005-membership]{\ref*{def:p01-w02-p01-005-membership}} & \thepage & Abbreviates \(\neg\,x\in a\). \\
\phantomsection\label{entry:symbol-registry-06587}\(\Lang_{\in}\) & \texttt{\detokenize{\Lang_{\in}}} & one-sorted first-order language of set theory with the single relation symbol \(\in\) & \hyperref[sec:p01-w02-p01-005-set-operations]{\ref*{sec:p01-w02-p01-005-set-operations}} & n/a & \hyperref[def:p01-w02-p01-005-membership]{\ref*{def:p01-w02-p01-005-membership}} & \thepage & Same language and notation as the inherited axiomatic set-theory chapter (sec:zf-zfc-01); equality is logical identity. \\
\phantomsection\label{entry:symbol-registry-06588}\(\varnothing\) & \texttt{\detokenize{\varnothing}} & empty set & \hyperref[sec:p01-w02-p01-005-set-operations]{\ref*{sec:p01-w02-p01-005-set-operations}} & n/a & \hyperref[sec:p01-w02-p01-005-operations-defs]{\ref*{sec:p01-w02-p01-005-operations-defs}} & \thepage & Unique witness of the empty-set axiom; same role as in sec:zf-zfc-01. Distinct from the iterated-jump notation \(\varnothing^{(n)}\). \\
\phantomsection\label{entry:symbol-registry-06589}\(\{a,b\}\) & \texttt{\detokenize{\{a,b\}}} & unordered pair and singleton symbol family member \(\{a,b\}\) & \hyperref[sec:p01-w02-p01-005-set-operations]{\ref*{sec:p01-w02-p01-005-set-operations}} & n/a & \hyperref[sec:p01-w02-p01-005-operations-defs]{\ref*{sec:p01-w02-p01-005-operations-defs}} & \thepage & Unique witness of pairing; \(\{a\}=\{a,a\}\). Same role as in sec:zf-zfc-01. Explicit displayed family; member types and common actual first-use exposed, not one atomic glyph. Family listing, not one atomic glyph. Author macro \texttt{\detokenize{\{a,b\},}\allowbreak\detokenize{\{a\}}} not carried to the members. \\
\phantomsection\label{entry:symbol-registry-06590}\(\{a\}\) & \texttt{\detokenize{\{a\}}} & unordered pair and singleton symbol family member \(\{a\}\) & \hyperref[sec:p01-w02-p01-005-set-operations]{\ref*{sec:p01-w02-p01-005-set-operations}} & n/a & \hyperref[sec:p01-w02-p01-005-operations-defs]{\ref*{sec:p01-w02-p01-005-operations-defs}} & \thepage & Unique witness of pairing; \(\{a\}=\{a,a\}\). Same role as in sec:zf-zfc-01. Explicit displayed family; member types and common actual first-use exposed, not one atomic glyph. Family listing, not one atomic glyph. Author macro \texttt{\detokenize{\{a,b\},}\allowbreak\detokenize{\{a\}}} not carried to the members. \\
\phantomsection\label{entry:symbol-registry-06591}\(\{a,b,c\}\) & \texttt{\detokenize{\{a,b,c\}}} & three-element listing & \hyperref[sec:p01-w02-p01-005-set-operations]{\ref*{sec:p01-w02-p01-005-set-operations}} & n/a & \hyperref[def:p01-w02-p01-005-operations]{\ref*{def:p01-w02-p01-005-operations}} & \thepage & Defined as \(\{a,b\}\cup\{c\}\). \\
\phantomsection\label{entry:symbol-registry-06592}\(\bigcup a\) & \texttt{\detokenize{\bigcup}} & union of a set & \hyperref[sec:p01-w02-p01-005-set-operations]{\ref*{sec:p01-w02-p01-005-set-operations}} & n/a & \hyperref[sec:p01-w02-p01-005-operations-defs]{\ref*{sec:p01-w02-p01-005-operations-defs}} & \thepage & Unique witness of the union axiom; same role as in sec:zf-zfc-01. \\
\phantomsection\label{entry:symbol-registry-06593}\(A\cup B\) & \texttt{\detokenize{\cup}} & binary union & \hyperref[sec:p01-w02-p01-005-set-operations]{\ref*{sec:p01-w02-p01-005-set-operations}} & n/a & \hyperref[def:p01-w02-p01-005-operations]{\ref*{def:p01-w02-p01-005-operations}} & \thepage & Defined as \(\bigcup\{A,B\}\); same role as the Boolean join of def:p01-bpi-filter-ideal on \(\mathcal P(S)\). \\
\phantomsection\label{entry:symbol-registry-06594}\(A\cap B\) & \texttt{\detokenize{\cap}} & binary intersection & \hyperref[sec:p01-w02-p01-005-set-operations]{\ref*{sec:p01-w02-p01-005-set-operations}} & n/a & \hyperref[def:p01-w02-p01-005-operations]{\ref*{def:p01-w02-p01-005-operations}} & \thepage & Defined as \(\{x\in A:x\in B\}\); same role as the Boolean meet of def:p01-bpi-filter-ideal on \(\mathcal P(S)\). \\
\phantomsection\label{entry:symbol-registry-06595}\(A\setminus B\) & \texttt{\detokenize{\setminus}} & relative complement (difference) & \hyperref[sec:p01-w02-p01-005-set-operations]{\ref*{sec:p01-w02-p01-005-set-operations}} & n/a & \hyperref[def:p01-w02-p01-005-operations]{\ref*{def:p01-w02-p01-005-operations}} & \thepage & Defined as \(\{x\in A:x\notin B\}\). No absolute complement exists (thm:p01-w02-p01-005-no-universal). \\
\phantomsection\label{entry:symbol-registry-06596}\(\bigcap F\) & \texttt{\detokenize{\bigcap}} & intersection of a nonempty set & \hyperref[sec:p01-w02-p01-005-set-operations]{\ref*{sec:p01-w02-p01-005-set-operations}} & n/a & \hyperref[def:p01-w02-p01-005-big-intersection]{\ref*{def:p01-w02-p01-005-big-intersection}} & \thepage & Defined only for \(F\neq\varnothing\) as \(\{x\in\bigcup F:\forall Y\,(Y\in F\to x\in Y)\}\); same role as in equivalence-closure-and-quotients. \\
\phantomsection\label{entry:symbol-registry-06597}\(\{x\in A:\varphi(x,\bar p)\}\) & \texttt{\detokenize{\{x\in A:\varphi\}}} & bounded set formation (separated set) & \hyperref[sec:p01-w02-p01-005-set-operations]{\ref*{sec:p01-w02-p01-005-set-operations}} & n/a & \hyperref[def:p01-w02-p01-005-bounded-formation]{\ref*{def:p01-w02-p01-005-bounded-formation}} & \thepage & Denotes a set only under the side conditions of def:p01-w02-p01-005-bounded-formation; contrast the class notation \(\{x:\varphi(x)\}\) of sec:set-prerequisites-01. \\
\phantomsection\label{entry:symbol-registry-06598}\(s\cup\{s\}\) & \texttt{\detokenize{s\cup\{s\}}} & successor of a set & \hyperref[sec:p01-w02-p01-005-set-operations]{\ref*{sec:p01-w02-p01-005-set-operations}} & n/a & \hyperref[def:p01-w02-p01-005-operations]{\ref*{def:p01-w02-p01-005-operations}} & \thepage & The operation written \(S(x)\) in the axiom of infinity of sec:zf-zfc-01. \\
\phantomsection\label{entry:symbol-registry-06599}\(e_0\) & \texttt{\detokenize{e_0}} & three named pure sets of the worked example symbol family member \(e_0\) & \hyperref[def:p01-w02-p01-005-example]{\ref*{def:p01-w02-p01-005-example}} & n/a & \hyperref[def:p01-w02-p01-005-example]{\ref*{def:p01-w02-p01-005-example}} & \thepage & \(e_0=\varnothing\), \(e_1=\{\varnothing\}\), \(e_2=\{\{\varnothing\}\}\); local names only. Explicit displayed family; member types and common actual first-use exposed, not one atomic glyph. Family listing, not one atomic glyph. Author macro \texttt{\detokenize{e\_0,}\allowbreak\detokenize{e\_1,}\allowbreak\detokenize{e\_2}} not carried to the members. \\
\phantomsection\label{entry:symbol-registry-06600}\(e_1\) & \texttt{\detokenize{e_1}} & three named pure sets of the worked example symbol family member \(e_1\) & \hyperref[def:p01-w02-p01-005-example]{\ref*{def:p01-w02-p01-005-example}} & n/a & \hyperref[def:p01-w02-p01-005-example]{\ref*{def:p01-w02-p01-005-example}} & \thepage & \(e_0=\varnothing\), \(e_1=\{\varnothing\}\), \(e_2=\{\{\varnothing\}\}\); local names only. Explicit displayed family; member types and common actual first-use exposed, not one atomic glyph. Family listing, not one atomic glyph. Author macro \texttt{\detokenize{e\_0,}\allowbreak\detokenize{e\_1,}\allowbreak\detokenize{e\_2}} not carried to the members. \\
\phantomsection\label{entry:symbol-registry-06601}\(e_2\) & \texttt{\detokenize{e_2}} & three named pure sets of the worked example symbol family member \(e_2\) & \hyperref[def:p01-w02-p01-005-example]{\ref*{def:p01-w02-p01-005-example}} & n/a & \hyperref[def:p01-w02-p01-005-example]{\ref*{def:p01-w02-p01-005-example}} & \thepage & \(e_0=\varnothing\), \(e_1=\{\varnothing\}\), \(e_2=\{\{\varnothing\}\}\); local names only. Explicit displayed family; member types and common actual first-use exposed, not one atomic glyph. Family listing, not one atomic glyph. Author macro \texttt{\detokenize{e\_0,}\allowbreak\detokenize{e\_1,}\allowbreak\detokenize{e\_2}} not carried to the members. \\
\phantomsection\label{entry:symbol-registry-06602}\(r_a\) & \texttt{\detokenize{r_a}} & separated set of non-self-members of a & \hyperref[thm:p01-w02-p01-005-no-universal]{\ref*{thm:p01-w02-p01-005-no-universal}} & n/a & \hyperref[thm:p01-w02-p01-005-no-universal]{\ref*{thm:p01-w02-p01-005-no-universal}} & \thepage & \(\{x\in a:x\notin x\}\); local name. \\
\phantomsection\label{entry:symbol-registry-06603}\(\pi(z,A,B)\) & \texttt{\detokenize{\pi}} & ordered-pair condition bounding the product & \hyperref[sec:annex-p01-w02-p01-005-set-operations]{\ref*{sec:annex-p01-w02-p01-005-set-operations}} & n/a & \hyperref[annex:p01-w02-p01-005-product-bound]{\ref*{annex:p01-w02-p01-005-product-bound}} & \thepage & Local formula name inside the annex proof only. \\
\phantomsection\label{entry:symbol-registry-06604}\(G_{\cup}\) & \texttt{\detokenize{G_{\cup}}} & graphs of union intersection and complement on \(\mathcal P(S)\) symbol family member \(G_{\cup}\) & \hyperref[sec:annex-p01-w02-p01-005-set-operations]{\ref*{sec:annex-p01-w02-p01-005-set-operations}} & n/a & \hyperref[annex:p01-w02-p01-005-subset-algebra]{\ref*{annex:p01-w02-p01-005-subset-algebra}} & \thepage & Local names inside the annex proof only. Explicit displayed family; member types and common actual first-use exposed, not one atomic glyph. Family listing, not one atomic glyph. Author macro \texttt{\detokenize{G\_{\cup},}\allowbreak\detokenize{G\_{\cap},}\allowbreak\detokenize{G\_{c}}} not carried to the members. \\
\phantomsection\label{entry:symbol-registry-06605}\(G_{\cap}\) & \texttt{\detokenize{G_{\cap}}} & graphs of union intersection and complement on \(\mathcal P(S)\) symbol family member \(G_{\cap}\) & \hyperref[sec:annex-p01-w02-p01-005-set-operations]{\ref*{sec:annex-p01-w02-p01-005-set-operations}} & n/a & \hyperref[annex:p01-w02-p01-005-subset-algebra]{\ref*{annex:p01-w02-p01-005-subset-algebra}} & \thepage & Local names inside the annex proof only. Explicit displayed family; member types and common actual first-use exposed, not one atomic glyph. Family listing, not one atomic glyph. Author macro \texttt{\detokenize{G\_{\cup},}\allowbreak\detokenize{G\_{\cap},}\allowbreak\detokenize{G\_{c}}} not carried to the members. \\
\phantomsection\label{entry:symbol-registry-06606}\(G_{c}\) & \texttt{\detokenize{G_{c}}} & graphs of union intersection and complement on \(\mathcal P(S)\) symbol family member \(G_{c}\) & \hyperref[sec:annex-p01-w02-p01-005-set-operations]{\ref*{sec:annex-p01-w02-p01-005-set-operations}} & n/a & \hyperref[annex:p01-w02-p01-005-subset-algebra]{\ref*{annex:p01-w02-p01-005-subset-algebra}} & \thepage & Local names inside the annex proof only. Explicit displayed family; member types and common actual first-use exposed, not one atomic glyph. Family listing, not one atomic glyph. Author macro \texttt{\detokenize{G\_{\cup},}\allowbreak\detokenize{G\_{\cap},}\allowbreak\detokenize{G\_{c}}} not carried to the members. \\
\phantomsection\label{entry:symbol-registry-06607}\(\mathcal D_{\mathcal S}\) & \texttt{\detokenize{\mathcal D_{\mathcal S}}} & class of permitted semantic data of a semantics & \hyperref[sec:p01-w02-p01-016-semantic-data]{\ref*{sec:p01-w02-p01-016-semantic-data}} & n/a & \hyperref[def:p01-w02-p01-016-datum]{\ref*{def:p01-w02-p01-016-datum}} & \thepage & Members are the complete tuples of the semantics \(\mathcal S\): a valuation, a structure-assignment pair, or a Kripke structure-world pair in the three instances. \\
\phantomsection\label{entry:symbol-registry-06608}\(d\) & \texttt{\detokenize{d}} & semantic datum & \hyperref[sec:p01-w02-p01-016-semantic-data]{\ref*{sec:p01-w02-p01-016-semantic-data}} & n/a & \hyperref[def:p01-w02-p01-016-datum]{\ref*{def:p01-w02-p01-016-datum}} & \thepage & A member of \(\mathcal D_{\mathcal S}\); local metavariable, not a domain element. \\
\phantomsection\label{entry:symbol-registry-06609}\(d\models_{\mathcal S}\varphi\) & \texttt{\detokenize{\entails_{\mathcal S}}} & satisfaction of a formula by a datum of a named semantics & \hyperref[sec:p01-w02-p01-016-semantic-data]{\ref*{sec:p01-w02-p01-016-semantic-data}} & meta & \hyperref[def:p01-w02-p01-016-satisfaction]{\ref*{def:p01-w02-p01-016-satisfaction}} & \thepage & Instance of the registered per-system pattern \(\models_{S}\); with a set operand, satisfaction of every member. The global \(v\models\varphi\), \(\mathcal M,s\models\varphi\) and \(\mathcal M,w\models A\) rows are its three instances. \\
\phantomsection\label{entry:symbol-registry-06610}\(\llbracket t\rrbracket^{\mathcal M}_s\) & \texttt{\detokenize{\sem{t}^{\mathcal M}_s}} & value of a term in a structure under an assignment & \hyperref[sec:p01-w02-p01-016-semantic-data]{\ref*{sec:p01-w02-p01-016-semantic-data}} & n/a & \hyperref[def:p01-w02-p01-016-satisfaction]{\ref*{def:p01-w02-p01-016-satisfaction}} & \thepage & Inherited source term-value notation, used here in the atomic satisfaction clauses. \\
\phantomsection\label{entry:symbol-registry-06611}\(u^{+}\) & \texttt{\detokenize{u^{+}}} & valuation extending a finite partial valuation by F & \hyperref[thm:p01-w02-p01-016-finite-check]{\ref*{thm:p01-w02-p01-016-finite-check}} & n/a & \hyperref[thm:p01-w02-p01-016-finite-check]{\ref*{thm:p01-w02-p01-016-finite-check}} & \thepage & \(u\colon P\to\mathbb B\) for finite \(P\subseteq\mathsf{Prop}\); \(u^{+}\) agrees with \(u\) on \(P\) and is \(\mathbf F\) elsewhere. \\
\phantomsection\label{entry:symbol-registry-06612}\(\mathbb B\) & \texttt{\detokenize{\mathbb B}} & source CPL truth-value set and values symbol family member \(\mathbb B\) & \hyperref[sec:p01-w02-p01-016-semantic-data]{\ref*{sec:p01-w02-p01-016-semantic-data}} & n/a & \hyperref[def:p01-w02-p01-016-datum]{\ref*{def:p01-w02-p01-016-datum}} & \thepage & Same typed role as the registered rows scoped to P2 prop-syntax-semantics; this row extends their scope to this section. Explicit displayed family; member types and common actual first-use exposed, not one atomic glyph. Family listing, not one atomic glyph. Author macro \texttt{\detokenize{\mathbb B,}\allowbreak\detokenize{\mathbf F,}\allowbreak\detokenize{\mathbf T}} not carried to the members. \\
\phantomsection\label{entry:symbol-registry-06613}\(\mathbf F\) & \texttt{\detokenize{\mathbf F}} & source CPL truth-value set and values symbol family member \(\mathbf F\) & \hyperref[sec:p01-w02-p01-016-semantic-data]{\ref*{sec:p01-w02-p01-016-semantic-data}} & n/a & \hyperref[def:p01-w02-p01-016-datum]{\ref*{def:p01-w02-p01-016-datum}} & \thepage & Same typed role as the registered rows scoped to P2 prop-syntax-semantics; this row extends their scope to this section. Explicit displayed family; member types and common actual first-use exposed, not one atomic glyph. Family listing, not one atomic glyph. Author macro \texttt{\detokenize{\mathbb B,}\allowbreak\detokenize{\mathbf F,}\allowbreak\detokenize{\mathbf T}} not carried to the members. \\
\phantomsection\label{entry:symbol-registry-06614}\(\mathbf T\) & \texttt{\detokenize{\mathbf T}} & source CPL truth-value set and values symbol family member \(\mathbf T\) & \hyperref[sec:p01-w02-p01-016-semantic-data]{\ref*{sec:p01-w02-p01-016-semantic-data}} & n/a & \hyperref[def:p01-w02-p01-016-datum]{\ref*{def:p01-w02-p01-016-datum}} & \thepage & Same typed role as the registered rows scoped to P2 prop-syntax-semantics; this row extends their scope to this section. Explicit displayed family; member types and common actual first-use exposed, not one atomic glyph. Family listing, not one atomic glyph. Author macro \texttt{\detokenize{\mathbb B,}\allowbreak\detokenize{\mathbf F,}\allowbreak\detokenize{\mathbf T}} not carried to the members. \\
\phantomsection\label{entry:symbol-registry-06615}\(\Gamma\models^{\ell}_{\mathcal C}A\) & \texttt{\detokenize{\entails^{\ell}_{\mathcal C}}} & local modal consequence over a frame class & \hyperref[thm:p01-w02-p01-016-local-global]{\ref*{thm:p01-w02-p01-016-local-global}} & meta & \hyperref[thm:p01-w02-p01-016-local-global]{\ref*{thm:p01-w02-p01-016-local-global}} & \thepage & Same typed role as the registered row scoped to P8 modal-temporal-epistemic; used in this section for the modal instance of consequence over whole tuples. \\
\phantomsection\label{entry:symbol-registry-06616}\(\mathcal M=(W,R,V)\) & \texttt{\detokenize{W,}\allowbreak\detokenize{R,}\allowbreak\detokenize{V}} & Kripke structure components & \hyperref[sec:p01-w02-p01-016-semantic-data]{\ref*{sec:p01-w02-p01-016-semantic-data}} & n/a & \hyperref[def:p01-w02-p01-016-datum]{\ref*{def:p01-w02-p01-016-datum}} & \thepage & Worlds, accessibility relation, valuation \(\mathsf{Prop}\to\mathcal P(W)\); same typed roles as the registered P8 rows. \(R\) is the accessibility relation only inside the modal clauses; elsewhere in the section \(R\) is a unary relation symbol. \\
\phantomsection\label{entry:symbol-registry-06617}\(\operatorname{Open}(D)\) & \texttt{\detokenize{\operatorname{Open}}} & set of formulas with an open assumption occurrence in a derivation & \hyperref[sec:p01-w02-p01-016-assumptions]{\ref*{sec:p01-w02-p01-016-assumptions}} & n/a & \hyperref[def:p01-w02-p01-016-assumption]{\ref*{def:p01-w02-p01-016-assumption}} & \thepage & Defined for natural-deduction derivations of ND\_CPL and ND\_FOL; open means undischarged. \\
\phantomsection\label{entry:symbol-registry-06618}\(\rho\) & \texttt{\detokenize{\rho}} & last rule application of a derivation and the formulas it discharges symbol family member \(\rho\) & \hyperref[thm:p01-w02-p01-016-open-bounds]{\ref*{thm:p01-w02-p01-016-open-bounds}} & n/a & \hyperref[thm:p01-w02-p01-016-open-bounds]{\ref*{thm:p01-w02-p01-016-open-bounds}} & \thepage & \(\Delta_\rho\) is the set of formulas of which \(\rho\) discharges at least one occurrence; it is a formula set, not a context. Explicit displayed family; member types and common actual first-use exposed, not one atomic glyph. Family listing, not one atomic glyph. Author macro \texttt{\detokenize{\rho,}\allowbreak\detokenize{\Delta\_\rho}} not carried to the members. \\
\phantomsection\label{entry:symbol-registry-06619}\(\Delta_\rho\) & \texttt{\detokenize{\Delta_\rho}} & last rule application of a derivation and the formulas it discharges symbol family member \(\Delta_\rho\) & \hyperref[thm:p01-w02-p01-016-open-bounds]{\ref*{thm:p01-w02-p01-016-open-bounds}} & n/a & \hyperref[thm:p01-w02-p01-016-open-bounds]{\ref*{thm:p01-w02-p01-016-open-bounds}} & \thepage & \(\Delta_\rho\) is the set of formulas of which \(\rho\) discharges at least one occurrence; it is a formula set, not a context. Explicit displayed family; member types and common actual first-use exposed, not one atomic glyph. Family listing, not one atomic glyph. Author macro \texttt{\detokenize{\rho,}\allowbreak\detokenize{\Delta\_\rho}} not carried to the members. \\
\phantomsection\label{entry:symbol-registry-06620}\(D\) & \texttt{\detokenize{D}} & natural-deduction derivations symbol family member \(D\) & \hyperref[sec:p01-w02-p01-016-assumptions]{\ref*{sec:p01-w02-p01-016-assumptions}} & n/a & \hyperref[def:p01-w02-p01-016-assumption]{\ref*{def:p01-w02-p01-016-assumption}} & \thepage & Local metavariables for derivation trees in the statements and annex proofs. Explicit displayed family; member types and common actual first-use exposed, not one atomic glyph. Family listing, not one atomic glyph. Author macro \texttt{\detokenize{D,}\allowbreak\detokenize{D',}\allowbreak\detokenize{D'',}\allowbreak\detokenize{D\_k,}\allowbreak\detokenize{A,}\allowbreak\detokenize{E}} not carried to the members. \\
\phantomsection\label{entry:symbol-registry-06621}\(D'\) & \texttt{\detokenize{D'}} & natural-deduction derivations symbol family member \(D'\) & \hyperref[sec:p01-w02-p01-016-assumptions]{\ref*{sec:p01-w02-p01-016-assumptions}} & n/a & \hyperref[def:p01-w02-p01-016-assumption]{\ref*{def:p01-w02-p01-016-assumption}} & \thepage & Local metavariables for derivation trees in the statements and annex proofs. Explicit displayed family; member types and common actual first-use exposed, not one atomic glyph. Family listing, not one atomic glyph. Author macro \texttt{\detokenize{D,}\allowbreak\detokenize{D',}\allowbreak\detokenize{D'',}\allowbreak\detokenize{D\_k,}\allowbreak\detokenize{A,}\allowbreak\detokenize{E}} not carried to the members. \\
\phantomsection\label{entry:symbol-registry-06622}\(D''\) & \texttt{\detokenize{D''}} & natural-deduction derivations symbol family member \(D''\) & \hyperref[sec:p01-w02-p01-016-assumptions]{\ref*{sec:p01-w02-p01-016-assumptions}} & n/a & \hyperref[def:p01-w02-p01-016-assumption]{\ref*{def:p01-w02-p01-016-assumption}} & \thepage & Local metavariables for derivation trees in the statements and annex proofs. Explicit displayed family; member types and common actual first-use exposed, not one atomic glyph. Family listing, not one atomic glyph. Author macro \texttt{\detokenize{D,}\allowbreak\detokenize{D',}\allowbreak\detokenize{D'',}\allowbreak\detokenize{D\_k,}\allowbreak\detokenize{A,}\allowbreak\detokenize{E}} not carried to the members. \\
\phantomsection\label{entry:symbol-registry-06623}\(D_k\) & \texttt{\detokenize{D_k}} & natural-deduction derivations symbol family member \(D_k\) & \hyperref[sec:p01-w02-p01-016-assumptions]{\ref*{sec:p01-w02-p01-016-assumptions}} & n/a & \hyperref[def:p01-w02-p01-016-assumption]{\ref*{def:p01-w02-p01-016-assumption}} & \thepage & Local metavariables for derivation trees in the statements and annex proofs. Explicit displayed family; member types and common actual first-use exposed, not one atomic glyph. Family listing, not one atomic glyph. Author macro \texttt{\detokenize{D,}\allowbreak\detokenize{D',}\allowbreak\detokenize{D'',}\allowbreak\detokenize{D\_k,}\allowbreak\detokenize{A,}\allowbreak\detokenize{E}} not carried to the members. \\
\phantomsection\label{entry:symbol-registry-06624}\(A\) & \texttt{\detokenize{A}} & natural-deduction derivations symbol family member \(A\) & \hyperref[sec:p01-w02-p01-016-assumptions]{\ref*{sec:p01-w02-p01-016-assumptions}} & n/a & \hyperref[def:p01-w02-p01-016-assumption]{\ref*{def:p01-w02-p01-016-assumption}} & \thepage & Local metavariables for derivation trees in the statements and annex proofs. Explicit displayed family; member types and common actual first-use exposed, not one atomic glyph. Family listing, not one atomic glyph. Author macro \texttt{\detokenize{D,}\allowbreak\detokenize{D',}\allowbreak\detokenize{D'',}\allowbreak\detokenize{D\_k,}\allowbreak\detokenize{A,}\allowbreak\detokenize{E}} not carried to the members. \\
\phantomsection\label{entry:symbol-registry-06625}\(E\) & \texttt{\detokenize{E}} & natural-deduction derivations symbol family member \(E\) & \hyperref[sec:p01-w02-p01-016-assumptions]{\ref*{sec:p01-w02-p01-016-assumptions}} & n/a & \hyperref[def:p01-w02-p01-016-assumption]{\ref*{def:p01-w02-p01-016-assumption}} & \thepage & Local metavariables for derivation trees in the statements and annex proofs. Explicit displayed family; member types and common actual first-use exposed, not one atomic glyph. Family listing, not one atomic glyph. Author macro \texttt{\detokenize{D,}\allowbreak\detokenize{D',}\allowbreak\detokenize{D'',}\allowbreak\detokenize{D\_k,}\allowbreak\detokenize{A,}\allowbreak\detokenize{E}} not carried to the members. \\
\phantomsection\label{entry:symbol-registry-06626}\(P\) & \texttt{\detokenize{P}} & specified propositional variables symbol family member \(P\) & \hyperref[def:p01-w02-p01-016-worked-discharge]{\ref*{def:p01-w02-p01-016-worked-discharge}} & n/a & \hyperref[def:p01-w02-p01-016-worked-discharge]{\ref*{def:p01-w02-p01-016-worked-discharge}} & \thepage & Members of \(\mathsf{Prop}\) under the inherited readable-example convention. Explicit displayed family; member types and common actual first-use exposed, not one atomic glyph. Family listing, not one atomic glyph. Author macro \texttt{\detokenize{P,}\allowbreak\detokenize{Q,}\allowbreak\detokenize{S}} not carried to the members. \\
\phantomsection\label{entry:symbol-registry-06627}\(Q\) & \texttt{\detokenize{Q}} & specified propositional variables symbol family member \(Q\) & \hyperref[def:p01-w02-p01-016-worked-discharge]{\ref*{def:p01-w02-p01-016-worked-discharge}} & n/a & \hyperref[def:p01-w02-p01-016-worked-discharge]{\ref*{def:p01-w02-p01-016-worked-discharge}} & \thepage & Members of \(\mathsf{Prop}\) under the inherited readable-example convention. Explicit displayed family; member types and common actual first-use exposed, not one atomic glyph. Family listing, not one atomic glyph. Author macro \texttt{\detokenize{P,}\allowbreak\detokenize{Q,}\allowbreak\detokenize{S}} not carried to the members. \\
\phantomsection\label{entry:symbol-registry-06628}\(S\) & \texttt{\detokenize{S}} & specified propositional variables symbol family member \(S\) & \hyperref[def:p01-w02-p01-016-worked-discharge]{\ref*{def:p01-w02-p01-016-worked-discharge}} & n/a & \hyperref[def:p01-w02-p01-016-worked-discharge]{\ref*{def:p01-w02-p01-016-worked-discharge}} & \thepage & Members of \(\mathsf{Prop}\) under the inherited readable-example convention. Explicit displayed family; member types and common actual first-use exposed, not one atomic glyph. Family listing, not one atomic glyph. Author macro \texttt{\detokenize{P,}\allowbreak\detokenize{Q,}\allowbreak\detokenize{S}} not carried to the members. \\
\phantomsection\label{entry:symbol-registry-06629}\(P\) & \texttt{\detokenize{P}} & unary relation symbols of a local example symbol family member \(P\) & \hyperref[def:p01-w02-p01-016-worked-models]{\ref*{def:p01-w02-p01-016-worked-models}} & n/a & \hyperref[def:p01-w02-p01-016-worked-models]{\ref*{def:p01-w02-p01-016-worked-models}} & \thepage & Object-language signature symbol local to the worked example, the local-global theorem, the closure-gap theorem and the generalization failure case. Explicit displayed family; member types and common actual first-use exposed, not one atomic glyph. Family listing, not one atomic glyph. Author macro \texttt{\detokenize{P,}\allowbreak\detokenize{Q}} not carried to the members. \\
\phantomsection\label{entry:symbol-registry-06630}\(Q\) & \texttt{\detokenize{Q}} & unary relation symbols of a local example symbol family member \(Q\) & \hyperref[def:p01-w02-p01-016-worked-models]{\ref*{def:p01-w02-p01-016-worked-models}} & n/a & \hyperref[def:p01-w02-p01-016-worked-models]{\ref*{def:p01-w02-p01-016-worked-models}} & \thepage & Object-language signature symbol local to the worked example, the local-global theorem, the closure-gap theorem and the generalization failure case. Explicit displayed family; member types and common actual first-use exposed, not one atomic glyph. Family listing, not one atomic glyph. Author macro \texttt{\detokenize{P,}\allowbreak\detokenize{Q}} not carried to the members. \\
\phantomsection\label{entry:symbol-registry-06631}\(c\) & \texttt{\detokenize{c}} & constant symbol of a local example & \hyperref[def:p01-w02-p01-016-worked-models]{\ref*{def:p01-w02-p01-016-worked-models}} & n/a & \hyperref[def:p01-w02-p01-016-worked-models]{\ref*{def:p01-w02-p01-016-worked-models}} & \thepage & Object-language signature symbol local to the worked example, the local-global theorem, the closure-gap theorem and the generalization failure case. \\
\phantomsection\label{entry:symbol-registry-06632}\(R\) & \texttt{\detokenize{R}} & unary relation symbol of local examples & \hyperref[thm:p01-w02-p01-016-local-global]{\ref*{thm:p01-w02-p01-016-local-global}}--\hyperref[thm:p01-w02-p01-016-closure-gap]{\ref*{thm:p01-w02-p01-016-closure-gap}} & n/a & \hyperref[thm:p01-w02-p01-016-local-global]{\ref*{thm:p01-w02-p01-016-local-global}} & \thepage & Object-language signature symbol local to the worked example, the local-global theorem, the closure-gap theorem and the generalization failure case. Used in the local-global theorem, the incomplete-tuple warning and the closure-gap theorem. \\
\phantomsection\label{entry:symbol-registry-06633}\(P\) & \texttt{\detokenize{P}} & unary relation symbol of the generalization failure case & \hyperref[thm:p01-w02-p01-016-eigen-failure]{\ref*{thm:p01-w02-p01-016-eigen-failure}} & n/a & \hyperref[thm:p01-w02-p01-016-eigen-failure]{\ref*{thm:p01-w02-p01-016-eigen-failure}} & \thepage & Object-language signature symbol local to the worked example, the local-global theorem, the closure-gap theorem and the generalization failure case. \\
\phantomsection\label{entry:symbol-registry-06634}\(\Gamma\vdash_{\mathcal C}\varphi\) & \texttt{\detokenize{\derives_{\mathcal C}}} & derivability in a specified calculus & \hyperref[sec:p01-w02-p01-019-consistency]{\ref*{sec:p01-w02-p01-019-consistency}} & meta & \hyperref[sec:p01-w02-p01-019-consistency]{\ref*{sec:p01-w02-p01-019-consistency}} & \thepage & Inherited notation of sec:levels-03 (definition Derivability); used here for an arbitrary calculus \(\mathcal C\) and for the test calculi \(\mathcal C_0,\mathcal C_{\mathrm{ex}},\mathcal C_P\). Same typed role as the registered per-system turnstile pattern. \\
\phantomsection\label{entry:symbol-registry-06635}\(\mathcal D_{\mathcal S}\) & \texttt{\detokenize{\mathcal D_{\mathcal S}}} & class of permitted semantic data of a semantics & \hyperref[sec:p01-w02-p01-019-consistency]{\ref*{sec:p01-w02-p01-019-consistency}} & n/a & \hyperref[sec:p01-w02-p01-019-consistency]{\ref*{sec:p01-w02-p01-019-consistency}} & \thepage & Same typed role as the row scoped to sec:p01-w02-p01-016-semantic-data; this row extends its scope to this section. \\
\phantomsection\label{entry:symbol-registry-06636}\(d\models_{\mathcal S}\varphi\) & \texttt{\detokenize{\entails_{\mathcal S}}} & satisfaction of a formula by a datum of a named semantics & \hyperref[sec:p01-w02-p01-019-consistency]{\ref*{sec:p01-w02-p01-019-consistency}} & meta & \hyperref[sec:p01-w02-p01-019-consistency]{\ref*{sec:p01-w02-p01-019-consistency}} & \thepage & Same typed role as the row scoped to sec:p01-w02-p01-016-semantic-data; with a set operand, satisfaction of every member. Also the consequence relation of \(\mathcal S\) when the left operand is a formula set. \\
\phantomsection\label{entry:symbol-registry-06637}\(F\) & \texttt{\detokenize{F}} & formula set of the language under discussion & \hyperref[sec:p01-w02-p01-019-consistency]{\ref*{sec:p01-w02-p01-019-consistency}} & n/a & \hyperref[sec:p01-w02-p01-019-standing]{\ref*{sec:p01-w02-p01-019-standing}} & \thepage & Nonempty and closed under \(\neg\) throughout the section; same role as in def:p01-w02-p01-016-datum. Not the truth value \(\mathbf F\). \\
\phantomsection\label{entry:symbol-registry-06638}\(R\) & \texttt{\detokenize{R}} & set of rule instances and the calculus it generates symbol family member \(R\) & \hyperref[def:p01-w02-p01-019-rule-calculus]{\ref*{def:p01-w02-p01-019-rule-calculus}} & n/a & \hyperref[def:p01-w02-p01-019-rule-calculus]{\ref*{def:p01-w02-p01-019-rule-calculus}} & \thepage & \(R\) is a set of pairs \((\Pi,\chi)\) with \(\Pi\) a finite formula set; local to the test-calculus construction. Explicit displayed family; member types and common actual first-use exposed, not one atomic glyph. Family listing, not one atomic glyph. Author macro \texttt{\detokenize{R,}\allowbreak\detokenize{\mathcal C\_R}} not carried to the members. \\
\phantomsection\label{entry:symbol-registry-06639}\(\ \mathcal C_R\) & \texttt{\detokenize{\ \mathcal C_R}} & set of rule instances and the calculus it generates symbol family member \(\ \mathcal C_R\) & \hyperref[def:p01-w02-p01-019-rule-calculus]{\ref*{def:p01-w02-p01-019-rule-calculus}} & n/a & \hyperref[def:p01-w02-p01-019-rule-calculus]{\ref*{def:p01-w02-p01-019-rule-calculus}} & \thepage & \(R\) is a set of pairs \((\Pi,\chi)\) with \(\Pi\) a finite formula set; local to the test-calculus construction. Explicit displayed family; member types and common actual first-use exposed, not one atomic glyph. Family listing, not one atomic glyph. Author macro \texttt{\detokenize{R,}\allowbreak\detokenize{\mathcal C\_R}} not carried to the members. \\
\phantomsection\label{entry:symbol-registry-06640}\(\Pi\) & \texttt{\detokenize{\Pi}} & finite premise set of a rule instance & \hyperref[def:p01-w02-p01-019-rule-calculus]{\ref*{def:p01-w02-p01-019-rule-calculus}} & n/a & \hyperref[def:p01-w02-p01-019-rule-calculus]{\ref*{def:p01-w02-p01-019-rule-calculus}} & \thepage & A finite formula set, consistent with the inherited normative standing of \(\Pi\) as a formula collection. \\
\phantomsection\label{entry:symbol-registry-06641}\(\mathcal C_0\) & \texttt{\detokenize{\mathcal C_0}} & three test calculi on the classical propositional language symbol family member \(\mathcal C_0\) & \hyperref[def:p01-w02-p01-019-three-calculi]{\ref*{def:p01-w02-p01-019-three-calculi}} & n/a & \hyperref[def:p01-w02-p01-019-three-calculi]{\ref*{def:p01-w02-p01-019-three-calculi}} & \thepage & Editorial counterexample calculi with rule sets \(\varnothing\), \(R_{\mathrm{ex}}\), \(R_P\); not proposed as logics. Explicit displayed family; member types and common actual first-use exposed, not one atomic glyph. Family listing, not one atomic glyph. Author macro \texttt{\detokenize{\mathcal C\_0,}\allowbreak\detokenize{\mathcal C\_{\mathrm{ex}},}\allowbreak\detokenize{\mathcal C\_P}} not carried to the members. \\
\phantomsection\label{entry:symbol-registry-06642}\(\mathcal C_{\mathrm{ex}}\) & \texttt{\detokenize{\mathcal C_{\mathrm{ex}}}} & three test calculi on the classical propositional language symbol family member \(\mathcal C_{\mathrm{ex}}\) & \hyperref[def:p01-w02-p01-019-three-calculi]{\ref*{def:p01-w02-p01-019-three-calculi}} & n/a & \hyperref[def:p01-w02-p01-019-three-calculi]{\ref*{def:p01-w02-p01-019-three-calculi}} & \thepage & Editorial counterexample calculi with rule sets \(\varnothing\), \(R_{\mathrm{ex}}\), \(R_P\); not proposed as logics. Explicit displayed family; member types and common actual first-use exposed, not one atomic glyph. Family listing, not one atomic glyph. Author macro \texttt{\detokenize{\mathcal C\_0,}\allowbreak\detokenize{\mathcal C\_{\mathrm{ex}},}\allowbreak\detokenize{\mathcal C\_P}} not carried to the members. \\
\phantomsection\label{entry:symbol-registry-06643}\(\mathcal C_P\) & \texttt{\detokenize{\mathcal C_P}} & three test calculi on the classical propositional language symbol family member \(\mathcal C_P\) & \hyperref[def:p01-w02-p01-019-three-calculi]{\ref*{def:p01-w02-p01-019-three-calculi}} & n/a & \hyperref[def:p01-w02-p01-019-three-calculi]{\ref*{def:p01-w02-p01-019-three-calculi}} & \thepage & Editorial counterexample calculi with rule sets \(\varnothing\), \(R_{\mathrm{ex}}\), \(R_P\); not proposed as logics. Explicit displayed family; member types and common actual first-use exposed, not one atomic glyph. Family listing, not one atomic glyph. Author macro \texttt{\detokenize{\mathcal C\_0,}\allowbreak\detokenize{\mathcal C\_{\mathrm{ex}},}\allowbreak\detokenize{\mathcal C\_P}} not carried to the members. \\
\phantomsection\label{entry:symbol-registry-06644}\(R_{\mathrm{ex}}\) & \texttt{\detokenize{R_{\mathrm{ex}}}} & rule-instance sets of the test calculi symbol family member \(R_{\mathrm{ex}}\) & \hyperref[def:p01-w02-p01-019-three-calculi]{\ref*{def:p01-w02-p01-019-three-calculi}} & n/a & \hyperref[def:p01-w02-p01-019-three-calculi]{\ref*{def:p01-w02-p01-019-three-calculi}} & \thepage & Local to the test-calculus definition. Explicit displayed family; member types and common actual first-use exposed, not one atomic glyph. Family listing, not one atomic glyph. Author macro \texttt{\detokenize{R\_{\mathrm{ex}},}\allowbreak\detokenize{R\_P}} not carried to the members. \\
\phantomsection\label{entry:symbol-registry-06645}\(R_P\) & \texttt{\detokenize{R_P}} & rule-instance sets of the test calculi symbol family member \(R_P\) & \hyperref[def:p01-w02-p01-019-three-calculi]{\ref*{def:p01-w02-p01-019-three-calculi}} & n/a & \hyperref[def:p01-w02-p01-019-three-calculi]{\ref*{def:p01-w02-p01-019-three-calculi}} & \thepage & Local to the test-calculus definition. Explicit displayed family; member types and common actual first-use exposed, not one atomic glyph. Family listing, not one atomic glyph. Author macro \texttt{\detokenize{R\_{\mathrm{ex}},}\allowbreak\detokenize{R\_P}} not carried to the members. \\
\phantomsection\label{entry:symbol-registry-06646}\(\Gamma_v\) & \texttt{\detokenize{\Gamma_v}} & diagram set of a valuation (variables true under v and negations of the others) & \hyperref[thm:p01-w02-p01-019-separation]{\ref*{thm:p01-w02-p01-019-separation}} & n/a & \hyperref[thm:p01-w02-p01-019-separation]{\ref*{thm:p01-w02-p01-019-separation}} & \thepage & Local to clause 4 of the separation corollary. \\
\phantomsection\label{entry:symbol-registry-06647}\(v_{\mathbf T}\) & \texttt{\detokenize{v_{\mathbf T}}} & constant valuations symbol family member \(v_{\mathbf T}\) & \hyperref[thm:p01-w02-p01-019-three-calculi]{\ref*{thm:p01-w02-p01-019-three-calculi}} & n/a & \hyperref[thm:p01-w02-p01-019-three-calculi]{\ref*{thm:p01-w02-p01-019-three-calculi}} & \thepage & Send every propositional variable to \(\mathbf T\) (respectively \(\mathbf F\)); used in annex proofs only. Explicit displayed family; member types and common actual first-use exposed, not one atomic glyph. Family listing, not one atomic glyph. Author macro \texttt{\detokenize{v\_{\mathbf T},}\allowbreak\detokenize{v\_{\mathbf F}}} not carried to the members. \\
\phantomsection\label{entry:symbol-registry-06648}\(v_{\mathbf F}\) & \texttt{\detokenize{v_{\mathbf F}}} & constant valuations symbol family member \(v_{\mathbf F}\) & \hyperref[thm:p01-w02-p01-019-three-calculi]{\ref*{thm:p01-w02-p01-019-three-calculi}} & n/a & \hyperref[thm:p01-w02-p01-019-three-calculi]{\ref*{thm:p01-w02-p01-019-three-calculi}} & \thepage & Send every propositional variable to \(\mathbf T\) (respectively \(\mathbf F\)); used in annex proofs only. Explicit displayed family; member types and common actual first-use exposed, not one atomic glyph. Family listing, not one atomic glyph. Author macro \texttt{\detokenize{v\_{\mathbf T},}\allowbreak\detokenize{v\_{\mathbf F}}} not carried to the members. \\
\phantomsection\label{entry:symbol-registry-06649}\(P\) & \texttt{\detokenize{P}} & specified propositional variables symbol family member \(P\) & \hyperref[sec:p01-w02-p01-019-consistency]{\ref*{sec:p01-w02-p01-019-consistency}} & n/a & \hyperref[def:p01-w02-p01-019-three-calculi]{\ref*{def:p01-w02-p01-019-three-calculi}} & \thepage & Members of \(\mathsf{Prop}\) under the inherited readable-example convention (ch. Propositional Syntax and Semantics, definition Propositional alphabet). \(S\) occurs only in the worked example. Explicit displayed family; member types and common actual first-use exposed, not one atomic glyph. Family listing, not one atomic glyph. Author macro \texttt{\detokenize{P,}\allowbreak\detokenize{Q,}\allowbreak\detokenize{S}} not carried to the members. \\
\phantomsection\label{entry:symbol-registry-06650}\(Q\) & \texttt{\detokenize{Q}} & specified propositional variables symbol family member \(Q\) & \hyperref[sec:p01-w02-p01-019-consistency]{\ref*{sec:p01-w02-p01-019-consistency}} & n/a & \hyperref[def:p01-w02-p01-019-three-calculi]{\ref*{def:p01-w02-p01-019-three-calculi}} & \thepage & Members of \(\mathsf{Prop}\) under the inherited readable-example convention (ch. Propositional Syntax and Semantics, definition Propositional alphabet). \(S\) occurs only in the worked example. Explicit displayed family; member types and common actual first-use exposed, not one atomic glyph. Family listing, not one atomic glyph. Author macro \texttt{\detokenize{P,}\allowbreak\detokenize{Q,}\allowbreak\detokenize{S}} not carried to the members. \\
\phantomsection\label{entry:symbol-registry-06651}\(S\) & \texttt{\detokenize{S}} & specified propositional variables symbol family member \(S\) & \hyperref[sec:p01-w02-p01-019-consistency]{\ref*{sec:p01-w02-p01-019-consistency}} & n/a & \hyperref[def:p01-w02-p01-019-three-calculi]{\ref*{def:p01-w02-p01-019-three-calculi}} & \thepage & Members of \(\mathsf{Prop}\) under the inherited readable-example convention (ch. Propositional Syntax and Semantics, definition Propositional alphabet). \(S\) occurs only in the worked example. Explicit displayed family; member types and common actual first-use exposed, not one atomic glyph. Family listing, not one atomic glyph. Author macro \texttt{\detokenize{P,}\allowbreak\detokenize{Q,}\allowbreak\detokenize{S}} not carried to the members. \\
\phantomsection\label{entry:symbol-registry-06652}\(\mathcal M\) & \texttt{\detokenize{\mathcal M}} & intuitionistic forcing at a world symbol family member \(\mathcal M\) & \hyperref[thm:p01-w02-p01-019-intuitionistic]{\ref*{thm:p01-w02-p01-019-intuitionistic}} & meta & \hyperref[thm:p01-w02-p01-019-intuitionistic]{\ref*{thm:p01-w02-p01-019-intuitionistic}} & \thepage & Inherited notation of ch:intuitionistic-logic, with \(A\) a formula and \(p\) a propositional variable under that chapter's metavariable convention; the registered single forcing macro. Explicit displayed family; member types and common actual first-use exposed, not one atomic glyph. Family listing, not one atomic glyph. Author macro \texttt{\detokenize{\forces}} not carried to the members. \\
\phantomsection\label{entry:symbol-registry-06653}\(w\Vdash A\) & \texttt{\detokenize{w\Vdash A}} & intuitionistic forcing at a world symbol family member \(w\Vdash A\) & \hyperref[thm:p01-w02-p01-019-intuitionistic]{\ref*{thm:p01-w02-p01-019-intuitionistic}} & meta & \hyperref[thm:p01-w02-p01-019-intuitionistic]{\ref*{thm:p01-w02-p01-019-intuitionistic}} & \thepage & Inherited notation of ch:intuitionistic-logic, with \(A\) a formula and \(p\) a propositional variable under that chapter's metavariable convention; the registered single forcing macro. Explicit displayed family; member types and common actual first-use exposed, not one atomic glyph. Family listing, not one atomic glyph. Author macro \texttt{\detokenize{\forces}} not carried to the members. \\
\phantomsection\label{entry:symbol-registry-06654}\(\Gamma\models_{\mathsf{IPL}}A\) & \texttt{\detokenize{\entails_{\mathsf{IPL}}}} & intuitionistic semantic consequence & \hyperref[thm:p01-w02-p01-019-intuitionistic]{\ref*{thm:p01-w02-p01-019-intuitionistic}} & meta & \hyperref[thm:p01-w02-p01-019-intuitionistic]{\ref*{thm:p01-w02-p01-019-intuitionistic}} & \thepage & Inherited notation of ch:intuitionistic-logic (definition Intuitionistic semantic consequence). \\
\end{longtable}
\endgroup
\Needspace{.4\textheight}
\subsection*{Stable part P2 scopes}
\begingroup\scriptsize\setlength{\tabcolsep}{3pt}
\begin{longtable}{@{}p{0.9in}p{1.3in}p{1.0in}p{0.5in}p{0.3in}p{0.4in}p{0.35in}p{1.65in}@{}}
\toprule
\textbf{Glyph} & \textbf{Macro} & \textbf{Role} & \textbf{Scope} & \textbf{Prec./form} & \textbf{First} & \textbf{Reg. page} & \textbf{Principal uses / notes} \\
\midrule
\endfirsthead
\toprule
\textbf{Glyph} & \textbf{Macro} & \textbf{Role} & \textbf{Scope} & \textbf{Prec./form} & \textbf{First} & \textbf{Reg. page} & \textbf{Principal uses / notes} \\
\midrule
\endhead
\midrule\multicolumn{8}{r}{\emph{continued}}\\\endfoot
\bottomrule\endlastfoot
\phantomsection\label{entry:symbol-registry-00013}\(\theta,\rho,\eta\) & \texttt{\detokenize{\theta,}\allowbreak\detokenize{\rho,}\allowbreak\detokenize{\eta}} & simultaneous term substitutions & Parts 2–12 from prop-deduction to mechanization & n/a & \hyperref[sec:prop-deduction-01]{\ref*{sec:prop-deduction-01}} & \thepage & Finite-support variable-to-term maps with declared action and composition. \\
\phantomsection\label{entry:symbol-registry-00045}\(\operatorname{ElDiag}(\mathcal M)\) & \texttt{\detokenize{\operatorname{ElDiag}}} & elementary diagram & Parts 2–3 from prop-deduction & n/a & \hyperref[sec:prop-deduction-01]{\ref*{sec:prop-deduction-01}} & \thepage & All parameter sentences true in \(\mathcal M\). \\
\phantomsection\label{entry:symbol-registry-00076}\(\models_{\mathrm{full}}\) & \texttt{\detokenize{\entails_{\mathrm{full}}}} & full second-order consequence & P2 prop-metatheory & meta & \hyperref[sec:prop-metatheory-01]{\ref*{sec:prop-metatheory-01}} & \thepage & Consequence over full standard structures. \\
\phantomsection\label{entry:symbol-registry-00077}\(\models_{\mathrm H}\) & \texttt{\detokenize{\entails_{\mathrm H}}} & Henkin consequence & P2 prop-metatheory & meta & \hyperref[sec:prop-metatheory-01]{\ref*{sec:prop-metatheory-01}} & \thepage & Consequence over Henkin structures. \\
\phantomsection\label{entry:symbol-registry-00121}\(v\models\varphi\) & \texttt{\detokenize{\entails}} & valuation satisfaction & P2 prop-syntax-semantics & meta & \hyperref[sec:prop-syntax-semantics-01]{\ref*{sec:prop-syntax-semantics-01}} & \thepage & The valuation's homomorphic extension designates \(\varphi\). \\
\phantomsection\label{entry:symbol-registry-00123}\(\mathsf{Prop}\) & \texttt{\detokenize{\mathsf{Prop}}} & set of propositional variables & P2 prop-syntax-semantics & n/a & \hyperref[sec:prop-syntax-semantics-01]{\ref*{sec:prop-syntax-semantics-01}} & \thepage & Exactly \(\{p_n:n\in\mathbb N\}\) in the fixed countable CPL language. \\
\phantomsection\label{entry:symbol-registry-00125}\(p_n; P,Q,R,S\) & \texttt{\detokenize{p_n; P,}\allowbreak\detokenize{Q,}\allowbreak\detokenize{R,}\allowbreak\detokenize{S}} & object-language propositional variables & P2 prop-syntax-semantics & n/a & \hyperref[sec:prop-syntax-semantics-01]{\ref*{sec:prop-syntax-semantics-01}} & \thepage & \(p_n\) is the indexed variable schema; uppercase letters denote specified members of \(\mathsf{Prop}\) in examples and diagrams, never formula metavariables. \\
\phantomsection\label{entry:symbol-registry-00126}\(\mathsf{Fm}_{\mathsf{CPL}}\) & \texttt{\detokenize{\mathsf{Fm}_{\mathsf{CPL}}}} & set of CPL formulas & P2 prop-syntax-semantics & n/a & \hyperref[sec:prop-syntax-semantics-01]{\ref*{sec:prop-syntax-semantics-01}} & \thepage & Least set generated by the declared propositional grammar. \\
\phantomsection\label{entry:symbol-registry-00127}\(\mathbb B\) & \texttt{\detokenize{\mathbb B}} & source CPL truth-value set & P2 prop-syntax-semantics & n/a & \hyperref[sec:prop-syntax-semantics-01]{\ref*{sec:prop-syntax-semantics-01}} & \thepage & The active source defines \(\{\mathbf F,\mathbf T\}\). Numerical Boolean valuations use the explicit coding \(\mathbf F\mapsto0\) and \(\mathbf T\mapsto1\); this is not literal equality of the two sets. \\
\phantomsection\label{entry:symbol-registry-00128}\(\mathbf F,\mathbf T\) & \texttt{\detokenize{\mathbf F,}\allowbreak\detokenize{\mathbf T}} & source CPL truth values & P2 prop-syntax-semantics & n/a & \hyperref[sec:prop-syntax-semantics-01]{\ref*{sec:prop-syntax-semantics-01}} & \thepage & False and true in the inherited recursive valuation clauses and the semantic substitution/replacement theorems. \\
\phantomsection\label{entry:symbol-registry-00129}\(\varphi\equiv\psi\) & \texttt{\detokenize{\equiv}} & source CPL semantic formula equivalence & P2 prop-syntax-semantics & meta & \hyperref[sec:prop-syntax-semantics-01]{\ref*{sec:prop-syntax-semantics-01}} & \thepage & Every source CPL valuation gives the formulas the same truth value; active in sec:prop-syntax-semantics-02 and sec:p02-propositional-substitution. It is a metalinguistic relation, not an object-language connective. \\
\phantomsection\label{entry:symbol-registry-00130}\(c(\varphi)\) & \texttt{\detokenize{c}} & formula complexity & P2 prop-syntax-semantics to prop-deduction & n/a & \hyperref[sec:prop-syntax-semantics-01]{\ref*{sec:prop-syntax-semantics-01}} & \thepage & Recursively counts constructor occurrences under the active propositional or first-order grammar. \\
\phantomsection\label{entry:symbol-registry-00131}\(\operatorname{Sub}(\varphi)\) & \texttt{\detokenize{\operatorname{Sub}}} & subformula-set operation & P2 prop-syntax-semantics & n/a & \hyperref[sec:prop-syntax-semantics-01]{\ref*{sec:prop-syntax-semantics-01}} & \thepage & Set of subformulas fixed by the recursive CPL definition. \\
\phantomsection\label{entry:symbol-registry-00132}\(\upsilon\) & \texttt{\detokenize{\upsilon}} & uniform propositional substitution & P2 prop-syntax-semantics & n/a & \hyperref[sec:prop-syntax-semantics-01]{\ref*{sec:prop-syntax-semantics-01}} & \thepage & A metalinguistic map \(\mathsf{Prop}\to\mathsf{Fm}_{\mathsf{CPL}}\); not the globally registered sentence metavariables \(\sigma,\tau\). \\
\phantomsection\label{entry:symbol-registry-00133}\(\varphi[\upsilon]\) & \texttt{\detokenize{[\upsilon]}} & uniform-substitution instance & P2 prop-syntax-semantics & meta & \hyperref[sec:prop-syntax-semantics-01]{\ref*{sec:prop-syntax-semantics-01}} & \thepage & Recursive action of the declared uniform substitution on a CPL formula. \\
\phantomsection\label{entry:symbol-registry-00134}\(\operatorname{Pos}(\varphi)\) & \texttt{\detokenize{\operatorname{Pos}}} & formula occurrence-position set & P2 prop-syntax-semantics & meta & \hyperref[sec:prop-syntax-semantics-01]{\ref*{sec:prop-syntax-semantics-01}} & \thepage & Finite parse-tree addresses over \(\{1,2\}\), defined in def:p02-selected-replacement; occurrences are distinct even when their subformula strings agree. \\
\phantomsection\label{entry:symbol-registry-00135}\(\varphi_a\) & \texttt{\detokenize{\varphi_a}} & subformula at a selected position & P2 prop-syntax-semantics & meta & \hyperref[sec:prop-syntax-semantics-01]{\ref*{sec:prop-syntax-semantics-01}} & \thepage & The subformula rooted at \(a\in\operatorname{Pos}(\varphi)\); the address is in the original formula tree. \\
\phantomsection\label{entry:symbol-registry-00136}\(\varphi[O\leftarrow\chi]\) & \texttt{\detokenize{[O\leftarrow\chi]}} & simultaneous selected-occurrence replacement & P2 prop-syntax-semantics & meta & \hyperref[sec:prop-syntax-semantics-01]{\ref*{sec:prop-syntax-semantics-01}} & \thepage & Replaces the chosen original occurrences in \(O\) by \(\chi\), with the exact admissibility conditions of def:p02-selected-replacement; inserted copies are not searched again. \\
\phantomsection\label{entry:symbol-registry-00137}\(\epsilon\) & \texttt{\detokenize{\epsilon}} & empty occurrence address & P2 prop-syntax-semantics & n/a & \hyperref[sec:prop-syntax-semantics-01]{\ref*{sec:prop-syntax-semantics-01}} & \thepage & The empty word over \(\{1,2\}\) denotes the root position in def:p02-selected-replacement; this is not an epsilon binder or an ordinal. \\
\phantomsection\label{entry:symbol-registry-00138}\(\overline L\) & \texttt{\detokenize{\overline L}} & literal complement & Legacy pre-batch4 P2.1-P7.8 reader meaning & meta & \hyperref[sec:prop-syntax-semantics-01]{\ref*{sec:prop-syntax-semantics-01}} & \thepage & Metalinguistic involution exchanging \(p_n\) and \(\neg p_n\); P7.8 reuses it on the extended-resolution sheet; disjoint from numeral bars.  Numeric chapter scope migrated to its frozen pre-batch4 source-labelled reader regions; glyph, role and precedence unchanged. \\
\phantomsection\label{entry:symbol-registry-00139}\(\operatorname{code}(L)\) & \texttt{\detokenize{\operatorname{code}}} & propositional-literal ordering code & P2 prop-syntax-semantics & n/a & \hyperref[sec:prop-syntax-semantics-01]{\ref*{sec:prop-syntax-semantics-01}} & \thepage & \(\operatorname{code}(p_n)=2n\) and \(\operatorname{code}(\neg p_n)=2n+1\); fixes one formula string for each clause. \\
\phantomsection\label{entry:symbol-registry-00140}\(\operatorname{form}_{\mathsf{CPL}}(C)\) & \texttt{\detokenize{\operatorname{form}_{\mathsf{CPL}}}} & formula interpretation of a clause & P2 prop-syntax-semantics & n/a & \hyperref[sec:prop-syntax-semantics-01]{\ref*{sec:prop-syntax-semantics-01}} & \thepage & Increasing-code disjunction of a finite clause; the empty clause maps to falsum. \\
\phantomsection\label{entry:symbol-registry-00141}\(\alpha^{\mathrm d}\) & \texttt{\detokenize{^{\mathrm d}}} & NNF dual operation & P2 prop-syntax-semantics & meta & \hyperref[sec:prop-syntax-semantics-01]{\ref*{sec:prop-syntax-semantics-01}} & \thepage & Metalinguistic involution fixing literals and exchanging conjunction with disjunction and falsum with verum. \\
\phantomsection\label{entry:symbol-registry-00142}\(\bar v\) & \texttt{\detokenize{\bar v}} & complemented valuation & P2 prop-syntax-semantics & n/a & \hyperref[sec:prop-syntax-semantics-01]{\ref*{sec:prop-syntax-semantics-01}} & \thepage & For the explicitly numerical Boolean presentation, \(\bar v(p)=1-v(p)\); with source-labelled values the same complement swaps \(\mathbf F\) and \(\mathbf T\). Used in propositional duality. \\
\phantomsection\label{entry:symbol-registry-00143}\(\langle C\rangle\) & \texttt{\detokenize{\langle C\rangle}} & clone generated by truth functions & P2 prop-syntax-semantics & meta & \hyperref[sec:prop-syntax-semantics-01]{\ref*{sec:prop-syntax-semantics-01}} & \thepage & Intersection of all positive-arity Boolean clones containing C. \\
\phantomsection\label{entry:symbol-registry-00144}\(\mathsf T_0,\mathsf T_1,\mathsf M,\mathsf L,\mathsf S\) & \texttt{\detokenize{\mathsf T_0,}\allowbreak\detokenize{\mathsf T_1,}\allowbreak\detokenize{\mathsf M,}\allowbreak\detokenize{\mathsf L,}\allowbreak\detokenize{\mathsf S}} & Post maximal clones & P2 prop-syntax-semantics & n/a & \hyperref[sec:prop-syntax-semantics-01]{\ref*{sec:prop-syntax-semantics-01}} & \thepage & Respectively 0-preserving, 1-preserving, monotone, affine, and self-dual positive-arity Boolean functions. \\
\phantomsection\label{entry:symbol-registry-00145}\(\mathbb F_2\) & \texttt{\detokenize{\mathbb F_2}} & two-element field & P2 prop-syntax-semantics & n/a & \hyperref[sec:prop-syntax-semantics-01]{\ref*{sec:prop-syntax-semantics-01}} & \thepage & Field used for affine Boolean functions and XOR constraints. \\
\phantomsection\label{entry:symbol-registry-00146}\(\mathrm{NAND}(\varphi,\psi),\mathrm{NOR}(\varphi,\psi)\) & \texttt{\detokenize{\mathrm{NAND},}\allowbreak\detokenize{\mathrm{NOR}}} & metalinguistic Boolean operator names & P2 prop-syntax-semantics & n/a & \hyperref[sec:prop-syntax-semantics-01]{\ref*{sec:prop-syntax-semantics-01}} & \thepage & Formula abbreviations for negated conjunction and negated disjunction; not active binary glyphs. \\
\phantomsection\label{entry:symbol-registry-00147}\(\mathsf{ND}_{\mathsf{CPL}},\mathsf{LK}_{\mathsf{CPL}},\mathsf H_{\mathsf{CPL}}\) & \texttt{\detokenize{\mathsf{ND}_{\mathsf{CPL}},}\allowbreak\detokenize{\mathsf{LK}_{\mathsf{CPL}},}\allowbreak\detokenize{\mathsf H_{\mathsf{CPL}}}} & CPL calculus identifiers & P2 prop-syntax-semantics & n/a & \hyperref[sec:prop-syntax-semantics-01]{\ref*{sec:prop-syntax-semantics-01}} & \thepage & Names the registered natural-deduction sequent and Hilbert calculi; derivability still carries the calculus as a turnstile subscript. \\
\phantomsection\label{entry:symbol-registry-00148}\(\operatorname{supp}(\Gamma)\) & \texttt{\detokenize{\operatorname{supp}}} & support of a finite multiset & Parts 2–7 from prop-syntax-semantics to lk-lj & n/a & \hyperref[sec:prop-syntax-semantics-01]{\ref*{sec:prop-syntax-semantics-01}} & \thepage & Formulas having positive multiplicity in the declared propositional or first-order finite-multiset context; reused in P7.1. \\
\phantomsection\label{entry:symbol-registry-00149}\(\operatorname{tr}_{\mathsf H}(\varphi)\) & \texttt{\detokenize{\operatorname{tr}_{\mathsf H}}} & full-language-to-Hilbert translation & P2 prop-syntax-semantics & n/a & \hyperref[sec:prop-syntax-semantics-01]{\ref*{sec:prop-syntax-semantics-01}} & \thepage & Metalinguistic recursive map into the negation-implication fragment; avoids the globally reserved sentence metavariable \(\tau\). \\
\phantomsection\label{entry:symbol-registry-00150}\(\mathsf P\) & \texttt{\detokenize{\mathsf P}} & deterministic polynomial-time language class & P2 prop-syntax-semantics & n/a & \hyperref[sec:prop-syntax-semantics-01]{\ref*{sec:prop-syntax-semantics-01}} & \thepage & Local propositional-complexity standing. \\
\phantomsection\label{entry:symbol-registry-00151}\(\mathsf{NP}\) & \texttt{\detokenize{\mathsf{NP}}} & nondeterministic polynomial-time language class & Parts 2–12 from prop-syntax-semantics to formal-practice & n/a & \hyperref[sec:prop-syntax-semantics-01]{\ref*{sec:prop-syntax-semantics-01}} & \thepage & Equivalently polynomially witnessed languages under the declared encoding; reused in P4.6 P7.8 and the P12.1 proof-certificate comparison. \\
\phantomsection\label{entry:symbol-registry-00152}\(\mathsf{coNP}\) & \texttt{\detokenize{\mathsf{coNP}}} & complements of NP languages & Legacy pre-batch4 P2.1-P7.8 reader meaning & n/a & \hyperref[sec:prop-syntax-semantics-01]{\ref*{sec:prop-syntax-semantics-01}} & \thepage & Decision-language class under the declared encoding; reused in P4.6 and P7.8.  Numeric chapter scope migrated to its frozen pre-batch4 source-labelled reader regions; glyph, role and precedence unchanged. \\
\phantomsection\label{entry:symbol-registry-00153}\(\mathsf{SAT}\) & \texttt{\detokenize{\mathsf{SAT}}} & satisfiable propositional encodings & P2 prop-syntax-semantics & n/a & \hyperref[sec:prop-syntax-semantics-01]{\ref*{sec:prop-syntax-semantics-01}} & \thepage & Malformed strings are excluded by the local encoding convention. \\
\phantomsection\label{entry:symbol-registry-00154}\(\mathsf{TAUT}\) & \texttt{\detokenize{\mathsf{TAUT}}} & valid propositional encodings & Legacy pre-batch4 P2.1-P7.8 reader meaning & n/a & \hyperref[sec:prop-syntax-semantics-01]{\ref*{sec:prop-syntax-semantics-01}} & \thepage & Malformed strings are excluded by the fixed encoding convention; reused in P4.6 and P7.8.  Numeric chapter scope migrated to its frozen pre-batch4 source-labelled reader regions; glyph, role and precedence unchanged. \\
\phantomsection\label{entry:symbol-registry-00165}\(\iota_n\) & \texttt{\detokenize{\iota_n}} & at-least-n sentence & P2 prop-deduction & n/a & \hyperref[sec:prop-deduction-01]{\ref*{sec:prop-deduction-01}} & \thepage & The first-order sentence asserting the existence of \(n\) pairwise distinct elements. \\
\phantomsection\label{entry:symbol-registry-00166}\(\Lang_R,\delta_R\) & \texttt{\detokenize{\Lang_R,}\allowbreak\detokenize{\delta_R}} & definitional-extension signature and axiom & P2 prop-deduction & n/a & \hyperref[sec:prop-deduction-01]{\ref*{sec:prop-deduction-01}} & \thepage & The signature obtained by adjoining relation symbol \(R\) and its exact defining biconditional axiom. \\
\phantomsection\label{entry:symbol-registry-00167}\(R^{\mathcal M}_{\varphi}\) & \texttt{\detokenize{R^{\mathcal M}_{\varphi}}} & relation defined by a formula in a structure & P2 prop-deduction & n/a & \hyperref[sec:prop-deduction-01]{\ref*{sec:prop-deduction-01}} & \thepage & The tuple relation cut out by \(\varphi\) in \(\mathcal M\); deliberately not \(R^*\), which is reserved for reflexive-transitive closure. \\
\phantomsection\label{entry:symbol-registry-00168}\(\mathsf{ND}_{\mathsf{FOL}},\mathsf{LK},\mathsf H_{\mathsf{FOL}}\) & \texttt{\detokenize{\mathsf{ND}_{\mathsf{FOL}},}\allowbreak\detokenize{\mathsf{LK},}\allowbreak\detokenize{\mathsf H_{\mathsf{FOL}}}} & first-order calculus identifiers & Parts 2–8 from prop-deduction to intuitionistic-logic & n/a & \hyperref[sec:prop-deduction-01]{\ref*{sec:prop-deduction-01}} & \thepage & Names the registered natural-deduction sequent and equality-free Hilbert presentations; canonical LK recurs in P7 and the P8.1 negative-translation statements; turnstiles retain exact subscripts. \\
\phantomsection\label{entry:symbol-registry-00169}\(\Lang_{\Gamma},\Lang_{\Gamma}^{*}\) & \texttt{\detokenize{\Lang_{\Gamma},}\allowbreak\detokenize{\Lang_{\Gamma}^{*}}} & active tableau signatures & P2 prop-deduction & n/a & \hyperref[sec:prop-deduction-01]{\ref*{sec:prop-deduction-01}} & \thepage & The finite root subsignature and its countable fresh-parameter expansion used by the fair first-order tableau construction. \\
\phantomsection\label{entry:symbol-registry-00170}\(B^{+}_{\mathrm{at}},\mathcal M_{B^{+}_{\mathrm{at}}}\) & \texttt{\detokenize{B^{+}_{\mathrm{at}},}\allowbreak\detokenize{\mathcal M_{B^{+}_{\mathrm{at}}}}} & positive atomic branch and tableau countermodel & P2 prop-deduction & n/a & \hyperref[sec:prop-deduction-01]{\ref*{sec:prop-deduction-01}} & \thepage & The positive ground atoms on a fair saturated branch and the induced active-language Herbrand structure. \\
\phantomsection\label{entry:symbol-registry-00171}\(\mathcal M_B\) & \texttt{\detokenize{\mathcal M_B}} & Herbrand structure determined by ground atoms & P2 prop-deduction & n/a & \hyperref[sec:prop-deduction-01]{\ref*{sec:prop-deduction-01}} & \thepage & Ground-term structure whose positive relation atoms are exactly the members of the declared Herbrand-base subset \(B\). \\
\phantomsection\label{entry:symbol-registry-00172}\(\operatorname{nnf}(\varphi),\operatorname{nnf}^{\neg}(\varphi)\) & \texttt{\detokenize{\operatorname{nnf},}\allowbreak\detokenize{\operatorname{nnf}^{\neg}}} & first-order NNF translations & P2 prop-deduction & n/a & \hyperref[sec:prop-deduction-01]{\ref*{sec:prop-deduction-01}} & \thepage & Simultaneously recursive maps equivalent to \(\varphi\) and \(\neg\varphi\), respectively. \\
\phantomsection\label{entry:symbol-registry-00173}\(\operatorname{rect}(\varphi,V),\operatorname{pren}(\varphi)\) & \texttt{\detokenize{\operatorname{rect},}\allowbreak\detokenize{\operatorname{pren}}} & rectification and prenex maps & P2 prop-deduction & n/a & \hyperref[sec:prop-deduction-01]{\ref*{sec:prop-deduction-01}} & \thepage & Metalinguistic constructions; the forbidden-variable parameter belongs only to rectification. \\
\phantomsection\label{entry:symbol-registry-00174}\(\operatorname{Sk}(\sigma)\) & \texttt{\detokenize{\operatorname{Sk}}} & Skolem form & P2 prop-deduction & n/a & \hyperref[sec:prop-deduction-01]{\ref*{sec:prop-deduction-01}} & \thepage & Universal Skolem sentence in the expanded signature; equisatisfiability and reduct directions retain their separate choice standing. \\
\phantomsection\label{entry:symbol-registry-00175}\(\operatorname{form}_{\Lang}(C),\operatorname{ucl}_{\Lang}(C)\) & \texttt{\detokenize{\operatorname{form}_{\Lang},}\allowbreak\detokenize{\operatorname{ucl}_{\Lang}}} & clause formula and universal closure & Parts 2–12 from prop-deduction to automated-reasoning & n/a & \hyperref[sec:prop-deduction-01]{\ref*{sec:prop-deduction-01}} & \thepage & Fixed literal-order disjunction for clause \(C\) and its universal closure over every occurring variable. \\
\phantomsection\label{entry:symbol-registry-00176}\(\operatorname{supp}_{\mathrm{var}}(\theta)\) & \texttt{\detokenize{\operatorname{supp}_{\mathrm{var}}}} & variable support of a term substitution & Parts 2–12 from prop-deduction to mechanization & n/a & \hyperref[sec:prop-deduction-01]{\ref*{sec:prop-deduction-01}} & \thepage & Variables not fixed by \(\theta\); the subscript prevents collision with finite-multiset support. \\
\phantomsection\label{entry:symbol-registry-00177}\(t\theta,\theta\eta\) & \texttt{\detokenize{t\theta,}\allowbreak\detokenize{\theta\eta}} & simultaneous substitution action and composition & Parts 2–12 from prop-deduction to mechanization & meta & \hyperref[sec:prop-deduction-01]{\ref*{sec:prop-deduction-01}} & \thepage & Action extends homomorphically and composition satisfies \((\theta\eta)(x)=(\theta(x))\eta\). \\
\phantomsection\label{entry:symbol-registry-00178}\(\theta\preceq_{\mathrm{inst}}\rho\) & \texttt{\detokenize{\preceq_{\mathrm{inst}}}} & substitution instance preorder & Parts 2–12 from prop-deduction to mechanization & meta & \hyperref[sec:prop-deduction-01]{\ref*{sec:prop-deduction-01}} & \thepage & Every instance of \(\rho\) is an instance of \(\theta\); equivalently \(\rho=\theta\eta\) for some substitution \(\eta\). \\
\phantomsection\label{entry:symbol-registry-00179}\(r\doteq t\) & \texttt{\detokenize{\doteq}} & unification equation & Parts 2–12 from prop-deduction to mechanization & meta & \hyperref[sec:prop-deduction-01]{\ref*{sec:prop-deduction-01}} & \thepage & Metalinguistic equation constraint processed by unification; not object-language identity. \\
\phantomsection\label{entry:symbol-registry-00180}\(H_{x,\varphi,c}\) & \texttt{\detokenize{H_{x,\varphi,c}}} & Henkin witness implication & P2 prop-deduction & n/a & \hyperref[sec:prop-deduction-01]{\ref*{sec:prop-deduction-01}} & \thepage & Metalinguistic abbreviation for \(\exists x\,\varphi\to\varphi[c/x]\) with fresh constant \(c\). \\
\phantomsection\label{entry:symbol-registry-00181}\(\Lang^H\) & \texttt{\detokenize{\Lang^H}} & Henkin constant expansion & P2 prop-deduction & n/a & \hyperref[sec:prop-deduction-01]{\ref*{sec:prop-deduction-01}} & \thepage & Signature obtained by adjoining the fresh constants used by the Henkin witness construction. \\
\phantomsection\label{entry:symbol-registry-00182}\(\sim_{\Delta},[t],\mathcal M_{\Delta},s_0\) & \texttt{\detokenize{\sim_{\Delta},}\allowbreak\detokenize{[t],}\allowbreak\detokenize{\mathcal M_{\Delta},}\allowbreak\detokenize{s_0}} & Henkin quotient-model notation & P2 prop-deduction & n/a & \hyperref[sec:prop-deduction-01]{\ref*{sec:prop-deduction-01}} & \thepage & Provable-equality relation term class quotient term structure and canonical variable assignment for a maximal consistent witness set \(\Delta\). \\
\phantomsection\label{entry:symbol-registry-00183}\(\Lang(\mathcal M),\mathcal M_{\mathcal M}\) & \texttt{\detokenize{\Lang(\mathcal M),}\allowbreak\detokenize{\mathcal M_{\mathcal M}}} & parameter signature and canonical expansion & Parts 2–3 from prop-deduction & n/a & \hyperref[sec:prop-deduction-01]{\ref*{sec:prop-deduction-01}} & \thepage & Adds one constant for each domain element and interprets every such constant by the element it names. \\
\phantomsection\label{entry:symbol-registry-00184}\(\mathcal N_{\mathrm{ar}}\) & \texttt{\detokenize{\mathcal N_{\mathrm{ar}}}} & subscripted alias for the standard first-order arithmetic structure & Parts 2–4 from prop-deduction & n/a & \hyperref[sec:prop-deduction-01]{\ref*{sec:prop-deduction-01}} & \thepage & Earlier separately authored material uses this explicit subscript; exact inherited Part V writes \(\mathcal N\) for the same structure in the fixed arithmetic signature. \\
\phantomsection\label{entry:symbol-registry-00185}\(\mathrm{FO}^{2},\mathrm{GF},\mathrm{EPR}\) & \texttt{\detokenize{\mathrm{FO}^{2},}\allowbreak\detokenize{\mathrm{GF},}\allowbreak\detokenize{\mathrm{EPR}}} & decidable first-order fragment identifiers & P2 prop-deduction & n/a & \hyperref[sec:prop-deduction-01]{\ref*{sec:prop-deduction-01}} & \thepage & Respectively the two-variable relational fragment guarded relational fragment and Bernays--Schoenfinkel effectively propositional class; equality and function scopes are row-local. \\
\phantomsection\label{entry:symbol-registry-00186}\(\mathsf{NEXPTIME},\mathsf{coNEXPTIME},\mathsf{EXPTIME},2\mathsf{EXPTIME}\) & \texttt{\detokenize{\mathsf{NEXPTIME},}\allowbreak\detokenize{\mathsf{coNEXPTIME},}\allowbreak\detokenize{\mathsf{EXPTIME},}\allowbreak\detokenize{2\mathsf{EXPTIME}}} & decision-problem complexity classes & Parts 2–4 from prop-deduction & n/a & \hyperref[sec:prop-deduction-01]{\ref*{sec:prop-deduction-01}} & \thepage & Every bound applies only to the named decision problem. \\
\phantomsection\label{entry:symbol-registry-00187}\(\exists^{*}\forall^{*}\) & \texttt{\detokenize{\exists^{*}\forall^{*}}} & finite quantifier-block pattern & P2 prop-deduction & n/a & \hyperref[sec:prop-deduction-01]{\ref*{sec:prop-deduction-01}} & \thepage & A finite existential block followed by a finite universal block; stars are metalinguistic repetition rather than object-language operators. \\
\phantomsection\label{entry:symbol-registry-00190}\(\Lang_2,\mathsf{RelVar}_n\) & \texttt{\detokenize{\Lang_2,}\allowbreak\detokenize{\mathsf{RelVar}_n}} & relational second-order signature and variable stocks & P2 prop-metatheory & n/a & \hyperref[sec:prop-metatheory-01]{\ref*{sec:prop-metatheory-01}} & \thepage & For every positive n the second component is the countably infinite stock of n-ary relation variables. \\
\phantomsection\label{entry:symbol-registry-00191}\(X^n,Y^n,U^n,F^{n+1}\) & \texttt{\detokenize{X^n,}\allowbreak\detokenize{Y^n,}\allowbreak\detokenize{U^n,}\allowbreak\detokenize{F^{n+1}}} & relation-variable metavariables & P2 prop-metatheory & n/a & \hyperref[sec:prop-metatheory-01]{\ref*{sec:prop-metatheory-01}} & \thepage & Superscripts record arity; Y and F carry freshness conditions and U is an admissible relation-variable or predicate-constant instance. \\
\phantomsection\label{entry:symbol-registry-00192}\(P^n\) & \texttt{\detokenize{P^n}} & predicate-constant metavariable & P2 prop-metatheory & n/a & \hyperref[sec:prop-metatheory-01]{\ref*{sec:prop-metatheory-01}} & \thepage & An n-ary nonlogical predicate constant in the active relational second-order signature. \\
\phantomsection\label{entry:symbol-registry-00193}\(\varphi[U^n/X^n]\) & \texttt{\detokenize{[U^n/}\allowbreak\detokenize{X^n]}} & capture-avoiding relation-variable substitution & P2 prop-metatheory & meta & \hyperref[sec:prop-metatheory-01]{\ref*{sec:prop-metatheory-01}} & \thepage & Metalinguistic substitution of a relation variable or predicate constant of the same arity; distinct from term substitution. \\
\phantomsection\label{entry:symbol-registry-00194}\(\mathsf{PA}_2\) & \texttt{\detokenize{\mathsf{PA}_2}} & finite full second-order Peano sentence & P2 prop-metatheory & n/a & \hyperref[sec:prop-metatheory-01]{\ref*{sec:prop-metatheory-01}} & \thepage & Definition def:p02-pa2 packages the three successor/induction axioms and four arithmetic recursion axioms into one closed sentence in the expanded signature; the corresponding finite axiom set has the same full models. \\
\phantomsection\label{entry:symbol-registry-00195}\(\mathsf D\) & \texttt{\detokenize{\mathsf D}} & effective calculus metavariable & P2 prop-metatheory & n/a & \hyperref[sec:prop-metatheory-01]{\ref*{sec:prop-metatheory-01}} & \thepage & A calculus whose finite derivations are computably checkable; local to the non-enumerability theorem. \\
\phantomsection\label{entry:symbol-registry-00196}\(\mathcal H,(\mathcal R_n)_{n\geq1}\) & \texttt{\detokenize{\mathcal H,}\allowbreak\detokenize{\mathcal R_n}} & relational Henkin structure and relation domains & P2 prop-metatheory & n/a & \hyperref[sec:prop-metatheory-01]{\ref*{sec:prop-metatheory-01}} & \thepage & The object domain is M and each nonempty relation domain is a selected subset of the corresponding powerset. \\
\phantomsection\label{entry:symbol-registry-00197}\(\mathsf H_{\mathsf{CC}}\) & \texttt{\detokenize{\mathsf H_{\mathsf{CC}}}} & matched comprehension-choice Henkin model class & P2 prop-metatheory & n/a & \hyperref[sec:prop-metatheory-01]{\ref*{sec:prop-metatheory-01}} & \thepage & Models satisfy exactly the eligible comprehension and graph-choice schemas; distinct from arbitrary Henkin structures. \\
\phantomsection\label{entry:symbol-registry-00198}\(\mathsf{Comp}_{\chi}^{n},\mathsf{Choice}_{\chi}^{n}\) & \texttt{\detokenize{\mathsf{Comp}_{\chi}^{n},}\allowbreak\detokenize{\mathsf{Choice}_{\chi}^{n}}} & matched second-order schema abbreviations & P2 prop-metatheory & n/a & \hyperref[sec:prop-metatheory-01]{\ref*{sec:prop-metatheory-01}} & \thepage & Metalinguistic names for the universally closed eligible formulas displayed on sheet:sol. \\
\phantomsection\label{entry:symbol-registry-00199}\(\Gamma\models_{\mathrm H}^{\mathsf{CC}}\varphi\) & \texttt{\detokenize{\entails_{\mathrm H}^{\mathsf{CC}}}} & matched Henkin consequence & P2 prop-metatheory & meta & \hyperref[sec:prop-metatheory-01]{\ref*{sec:prop-metatheory-01}} & \thepage & Truth preservation over exactly the matched comprehension-choice Henkin model class and matching assignments. \\
\phantomsection\label{entry:symbol-registry-00200}\(\operatorname{ucl}(\varphi)\) & \texttt{\detokenize{\operatorname{ucl}}} & formula-level universal closure & P2 prop-metatheory & n/a & \hyperref[sec:prop-metatheory-01]{\ref*{sec:prop-metatheory-01}} & \thepage & Universally closes exactly the displayed free object and relation parameters; distinct from the clause-language subscripted closure. \\
\phantomsection\label{entry:symbol-registry-00201}\(D,S_n,E_n,\widehat P\) & \texttt{\detokenize{D,}\allowbreak\detokenize{S_n,}\allowbreak\detokenize{E_n,}\allowbreak\detokenize{\widehat P}} & incidence-translation signature data & P2 prop-metatheory & n/a & \hyperref[sec:prop-metatheory-01]{\ref*{sec:prop-metatheory-01}} & \thepage & \(D\) is the object sort; \(S_n\) the n-ary relation sort; \(E_n\) incidence; and \(\widehat P\) the relation-sort constant representing predicate \(P\). \\
\phantomsection\label{entry:symbol-registry-00202}\((\cdot)^\sharp\) & \texttt{\detokenize{{}^\sharp}} & incidence translation on formulas contexts structures and assignments & P2 prop-metatheory & n/a & \hyperref[sec:prop-metatheory-01]{\ref*{sec:prop-metatheory-01}} & \thepage & A metalinguistic tagging and translation operation; the operand determines its eliminable clause. \\
\phantomsection\label{entry:symbol-registry-00203}\(T_{\mathsf{CC}}^\sharp\) & \texttt{\detokenize{T_{\mathsf{CC}}^\sharp}} & translated matched-schema theory & P2 prop-metatheory & n/a & \hyperref[sec:prop-metatheory-01]{\ref*{sec:prop-metatheory-01}} & \thepage & The many-sorted first-order translations of exactly all eligible comprehension and graph-choice instances. \\
\phantomsection\label{entry:symbol-registry-00204}\(\mathsf{MSFOL},\mathsf{SOL}_{\mathrm H}^{\mathsf{CC}}\) & \texttt{\detokenize{\mathsf{MSFOL},}\allowbreak\detokenize{\mathsf{SOL}_{\mathrm H}^{\mathsf{CC}}}} & translation-defined calculus identifiers & P2 prop-metatheory & n/a & \hyperref[sec:prop-metatheory-01]{\ref*{sec:prop-metatheory-01}} & \thepage & Eliminable metalinguistic derivability abbreviations; neither names an independently generated native rule set. \\
\phantomsection\label{entry:symbol-registry-00205}\(\operatorname{Reach}(a,b)\) & \texttt{\detokenize{\operatorname{Reach}}} & second-order reachability formula & P2 prop-metatheory & n/a & \hyperref[sec:prop-metatheory-01]{\ref*{sec:prop-metatheory-01}} & \thepage & The full-semantic closure formula for finite directed reachability; its Henkin converse can fail. \\
\phantomsection\label{entry:symbol-registry-00206}\(\alpha,\beta,\gamma\) & \texttt{\detokenize{\alpha,}\allowbreak\detokenize{\beta,}\allowbreak\detokenize{\gamma}} & type metavariables & P2 prop-metatheory & n/a & \hyperref[sec:prop-metatheory-01]{\ref*{sec:prop-metatheory-01}} & \thepage & Local to the simple-type grammar; disjoint from ordinal metavariables and the separate P7.4 type scope. \\
\phantomsection\label{entry:symbol-registry-00207}\(\iota,o\) & \texttt{\detokenize{\iota,}\allowbreak\detokenize{o}} & individual and truth-value base types & P2 prop-metatheory & n/a & \hyperref[sec:prop-metatheory-01]{\ref*{sec:prop-metatheory-01}} & \thepage & Base types of the representative simple type theory; iota here is a type and not the P2.4 description binder. \\
\phantomsection\label{entry:symbol-registry-00208}\(\alpha\to_{\mathrm{ty}}\beta\) & \texttt{\detokenize{\to_{\mathrm{ty}}}} & object-language function-type former & Legacy pre-batch4 P2.3-P7.6 reader meaning & 6 & \hyperref[sec:prop-metatheory-01]{\ref*{sec:prop-metatheory-01}} & \thepage & Right-associative type constructor in the active typed calculus; distinct from object-language implication.  Numeric chapter scope migrated to its frozen pre-batch4 source-labelled reader regions; glyph, role and precedence unchanged. \\
\phantomsection\label{entry:symbol-registry-00209}\(A,B,C,F,P,R\) & \texttt{\detokenize{A,}\allowbreak\detokenize{B,}\allowbreak\detokenize{C,}\allowbreak\detokenize{F,}\allowbreak\detokenize{P,}\allowbreak\detokenize{R}} & typed-term metavariables & P2 prop-metatheory & n/a & \hyperref[sec:prop-metatheory-01]{\ref*{sec:prop-metatheory-01}} & \thepage & Used only under local type declarations in P2.3; disjoint from P8 formula metavariables and algebraic carriers. \\
\phantomsection\label{entry:symbol-registry-00210}\(A:\alpha\) & \texttt{\detokenize{:}} & type ascription & Legacy pre-batch4 P2.3-P7.6 reader meaning & meta & \hyperref[sec:prop-metatheory-01]{\ref*{sec:prop-metatheory-01}} & \thepage & The active typed calculus fixes the displayed term and type categories; distinct from the P7.7 labelled-formula colon.  Numeric chapter scope migrated to its frozen pre-batch4 source-labelled reader regions; glyph, role and precedence unchanged. \\
\phantomsection\label{entry:symbol-registry-00211}\(\lambda x{:}\alpha.B\) & \texttt{\detokenize{\lambda}} & typed lambda abstraction & Legacy pre-batch4 P2.3-P7.6 reader meaning & 3 & \hyperref[sec:prop-metatheory-01]{\ref*{sec:prop-metatheory-01}} & \thepage & Object-language term former binding x in the displayed body; used by STLC System F and System T with capture-avoiding substitution.  Numeric chapter scope migrated to its frozen pre-batch4 source-labelled reader regions; glyph, role and precedence unchanged. \\
\phantomsection\label{entry:symbol-registry-00212}\(\mathsf{Eq}_{\alpha},\forall_{\alpha}\) & \texttt{\detokenize{\mathsf{Eq}_{\alpha},}\allowbreak\detokenize{\forall_{\alpha}}} & typed equality and universal-quantifier constants & P2 prop-metatheory & 3 & \hyperref[sec:prop-metatheory-01]{\ref*{sec:prop-metatheory-01}} & \thepage & Object-language typed constants; Eq-alpha is deliberately distinct from global metalinguistic alpha-equivalence. \\
\phantomsection\label{entry:symbol-registry-00213}\(D_\alpha\) & \texttt{\detokenize{D_\alpha}} & semantic domain at type alpha & P2 prop-metatheory & n/a & \hyperref[sec:prop-metatheory-01]{\ref*{sec:prop-metatheory-01}} & \thepage & Nonempty type domain; standard function types contain all functions and Henkin domains may select a closure-stable subset. \\
\phantomsection\label{entry:symbol-registry-00214}\(\lbrack\!\lbrack A\rbrack\!\rbrack_s^{\mathcal M}\) & \texttt{\detokenize{\lbrack\!\lbrack\cdot\rbrack\!\rbrack_s^{\mathcal M}}} & typed-term denotation & P2 prop-metatheory & n/a & \hyperref[sec:prop-metatheory-01]{\ref*{sec:prop-metatheory-01}} & \thepage & Recursive semantic value of a well-typed term under typed assignment s. \\
\phantomsection\label{entry:symbol-registry-00215}\(B[A/x]\) & \texttt{\detokenize{[A/}\allowbreak\detokenize{x]}} & capture-avoiding typed substitution & P2 prop-metatheory & meta & \hyperref[sec:prop-metatheory-01]{\ref*{sec:prop-metatheory-01}} & \thepage & Substitutes the typed term A for free occurrences of x in B after the required alpha-renaming. \\
\phantomsection\label{entry:symbol-registry-00216}\(\mathcal A\) & \texttt{\detokenize{\mathcal A}} & auxiliary translated structure & P2 prop-metatheory to fol-syntax-semantics & n/a & \hyperref[sec:prop-metatheory-01]{\ref*{sec:prop-metatheory-01}} & \thepage & A one-sorted or strict many-sorted model in the translation proofs; scope is disjoint from P8 algebraic structures. \\
\phantomsection\label{entry:symbol-registry-00217}\(\operatorname{ext}_n(r),s^{\mathrm{ext}}\) & \texttt{\detokenize{\operatorname{ext}_n,}\allowbreak\detokenize{s^{\mathrm{ext}}}} & extensional-collapse map and induced assignment & P2 prop-metatheory & n/a & \hyperref[sec:prop-metatheory-01]{\ref*{sec:prop-metatheory-01}} & \thepage & Maps a relation-sort object to its n-ary extension and transports relation-variable assignments through that map. \\
\phantomsection\label{entry:symbol-registry-00218}\(\Lang_{\mathrm{MS}},\mathsf{Sort},\operatorname{tp}\) & \texttt{\detokenize{\Lang_{\mathrm{MS}},}\allowbreak\detokenize{\mathsf{Sort},}\allowbreak\detokenize{\operatorname{tp}}} & strict many-sorted signature data & P2 fol-syntax-semantics & n/a & \hyperref[def:p02-fol-signature]{\ref*{def:p02-fol-signature}} & \thepage & The nonempty sort set and symbol-typing map determine the strict many-sorted language. \\
\phantomsection\label{entry:symbol-registry-00219}\(\mathsf s,\mathsf t,\mathsf s_i;x^{\mathsf s},M_{\mathsf s}\) & \texttt{\detokenize{\mathsf s,}\allowbreak\detokenize{\mathsf t,}\allowbreak\detokenize{\mathsf s_i;x^{\mathsf s},}\allowbreak\detokenize{M_{\mathsf s}}} & sort metavariables sorted variables and carriers & P2 fol-syntax-semantics & n/a & \hyperref[def:p02-fol-signature]{\ref*{def:p02-fol-signature}} & \thepage & Sans-serif sort letters avoid the canonical assignment s; strict carriers are nonempty and pairwise disjoint. \\
\phantomsection\label{entry:symbol-registry-00220}\(\Lang_{\mathrm{MS}}^{\flat},Q_{\mathsf s}\) & \texttt{\detokenize{\Lang_{\mathrm{MS}}^{\flat},}\allowbreak\detokenize{Q_{\mathsf s}}} & one-sorted reduct signature and sort predicates & P2 fol-syntax-semantics & n/a & \hyperref[def:p02-fol-signature]{\ref*{def:p02-fol-signature}} & \thepage & The flat signature forgets type declarations and \(Q_{\mathsf s}\) cuts out the component for sort \(\mathsf s\). \\
\phantomsection\label{entry:symbol-registry-00221}\(\operatorname{tr}_{\mathrm{MS}},\mathsf{Ax}_{\mathrm{sort}}\) & \texttt{\detokenize{\operatorname{tr}_{\mathrm{MS}},}\allowbreak\detokenize{\mathsf{Ax}_{\mathrm{sort}}}} & many-sorted guard translation and sort theory & P2 prop-metatheory to fol-syntax-semantics & meta & \hyperref[sec:prop-metatheory-01]{\ref*{sec:prop-metatheory-01}} & \thepage & The recursive translation guards bound variables; the theory asserts sort nonemptiness typing closure and disjointness. \\
\phantomsection\label{entry:symbol-registry-00222}\(\mathsf{VarGuard}_{\mathrm{MS}}(\Lang_{\mathrm{MS}})\) & \texttt{\detokenize{\mathsf{VarGuard}_{\mathrm{MS}}}} & total-assignment sort-guard set & P2 prop-metatheory to fol-syntax-semantics & meta & \hyperref[sec:prop-metatheory-01]{\ref*{sec:prop-metatheory-01}} & \thepage & Metalinguistic one-sorted atoms constraining every object and relation variable to its declared sort; nonredundant with bound-variable guards and sort axioms, and choice-free for total assignments. \\
\phantomsection\label{entry:symbol-registry-00223}\(D,D^{\mathrm{in}},D^{\mathrm{out}}\) & \texttt{\detokenize{D,}\allowbreak\detokenize{D^{\mathrm{in}},}\allowbreak\detokenize{D^{\mathrm{out}}}} & free-logic outer and inner domains & Parts 2–8 from fol-syntax-semantics to many-valued-paraconsistent & n/a & \hyperref[def:p02-fol-signature]{\ref*{def:p02-fol-signature}} & \thepage & Local outer-domain presentations; the inner existing domain is nonempty and contained in the outer domain. \\
\phantomsection\label{entry:symbol-registry-00224}\(E!(t)\) & \texttt{\detokenize{E!}} & free-logic existence predicate & Parts 2–8 from fol-syntax-semantics to many-valued-paraconsistent & 3 & \hyperref[def:p02-fol-signature]{\ref*{def:p02-fol-signature}} & \thepage & True exactly when the term value lies in the declared inner existing domain of the displayed free-logic semantics. \\
\phantomsection\label{entry:symbol-registry-00225}\(\iota x\,\varphi\) & \texttt{\detokenize{\iota}} & definite-description term former & Parts 2–8 from fol-syntax-semantics to many-valued-paraconsistent & 3 & \hyperref[def:p02-fol-signature]{\ref*{def:p02-fol-signature}} & \thepage & Binds x in the description formula; a local Russellian elimination or primitive partial-term policy fixes its standing. \\
\phantomsection\label{entry:symbol-registry-00226}\(L_{\omega_1\omega}\) & \texttt{\detokenize{L_{\omega_1\omega}}} & countably infinitary first-order language & P2 fol-syntax-semantics & n/a & \hyperref[def:p02-fol-signature]{\ref*{def:p02-fol-signature}} & \thepage & Allows countable Boolean combinations but only finite quantifier strings. \\
\phantomsection\label{entry:symbol-registry-00227}\(\bigwedge_{i\in I}\varphi_i\) & \texttt{\detokenize{\bigwedge}} & countable infinitary conjunction & P2 fol-syntax-semantics & 4 & \hyperref[def:p02-fol-signature]{\ref*{def:p02-fol-signature}} & \thepage & I is countable; the connective retains the ordinary conjunction precedence rank. \\
\phantomsection\label{entry:symbol-registry-00228}\(\bigvee_{i\in I}\varphi_i\) & \texttt{\detokenize{\bigvee}} & countable infinitary disjunction & P2 fol-syntax-semantics & 5 & \hyperref[def:p02-fol-signature]{\ref*{def:p02-fol-signature}} & \thepage & I is countable; the connective retains the ordinary disjunction precedence rank. \\
\phantomsection\label{entry:symbol-registry-00229}\(\mathsf{FOL}(Q)\) & \texttt{\detokenize{\mathsf{FOL}(Q)}} & generalized-quantifier extension placeholder & P2 fol-syntax-semantics & n/a & \hyperref[def:p02-fol-signature]{\ref*{def:p02-fol-signature}} & \thepage & Pointer notation only; Part 3 must define the supplied isomorphism-invariant quantifier family and its exact arity. \\
\phantomsection\label{entry:symbol-registry-00230}\(\Sigma^1_n,\Pi^1_n,\Delta^1_n,\boldsymbol\Sigma^1_n,\boldsymbol\Pi^1_n\) & \texttt{\detokenize{\Sigma^1_n,}\allowbreak\detokenize{\Pi^1_n,}\allowbreak\detokenize{\Delta^1_n,}\allowbreak\detokenize{\boldsymbol\Sigma^1_n,}\allowbreak\detokenize{\boldsymbol\Pi^1_n}} & lightface and boldface analytical hierarchy classes & Legacy pre-batch5 P2.4-P4.5 reader meaning & n/a & \hyperref[def:p02-fol-signature]{\ref*{def:p02-fol-signature}} & \thepage & Lightface classes have effective matrices and \(\Delta^1_n\) is their intersection; boldface classes allow real or set parameters and must not be identified with them.  Numeric chapter scope migrated to its frozen pre-batch5 source-labelled reader regions; glyph, role and precedence unchanged. \\
\phantomsection\label{entry:symbol-registry-00862}\(W,V\) & \texttt{\detokenize{W,}\allowbreak\detokenize{V}} & finite valuation domain and exact input-variable set & P2 prop-syntax-semantics & meta & \hyperref[sec:prop-syntax-semantics-01]{\ref*{sec:prop-syntax-semantics-01}} & \thepage & In sec:p02-tseitin-transformation a valuation on W assigns precisely W; V is the set of variables occurring in the input formula, before auxiliary variables are chosen. \\
\phantomsection\label{entry:symbol-registry-00863}\(G_\varphi\) & \texttt{\detokenize{G_\varphi}} & nonvariable occurrence positions & P2 prop-syntax-semantics & meta & \hyperref[sec:prop-syntax-semantics-01]{\ref*{sec:prop-syntax-semantics-01}} & \thepage & Positions whose rooted subformula is not a propositional variable, including bottom occurrences; defined in def:p02-tseitin-construction. \\
\phantomsection\label{entry:symbol-registry-00864}\(z_a,Z_\varphi\) & \texttt{\detokenize{z_a,}\allowbreak\detokenize{Z_\varphi}} & fresh occurrence variable and auxiliary-variable set & P2 prop-syntax-semantics & meta & \hyperref[sec:prop-syntax-semantics-01]{\ref*{sec:prop-syntax-semantics-01}} & \thepage & Distinct fresh variables are assigned to nonvariable occurrences; \(Z_\varphi\) is disjoint from the original variable set \(V\), and each \(z_a\) aliases a selected indexed propositional variable. \\
\phantomsection\label{entry:symbol-registry-00865}\(r_a\) & \texttt{\detokenize{r_a}} & variable representative of a formula occurrence & P2 prop-syntax-semantics & meta & \hyperref[sec:prop-syntax-semantics-01]{\ref*{sec:prop-syntax-semantics-01}} & \thepage & An original variable leaf represents itself; a nonvariable occurrence uses its fresh \(z_a\). The root representative is \(r_\epsilon\). \\
\phantomsection\label{entry:symbol-registry-00866}\(\mathcal D_\varphi\) & \texttt{\detokenize{\mathcal D_\varphi}} & unrooted defining clause set & P2 prop-syntax-semantics & meta & \hyperref[sec:prop-syntax-semantics-01]{\ref*{sec:prop-syntax-semantics-01}} & \thepage & All gate clauses in the Tseitin construction before the asserted root clause; every input assignment has a unique defining extension on \(V\cup Z_\varphi\). \\
\phantomsection\label{entry:symbol-registry-00867}\(S_\varphi\) & \texttt{\detokenize{S_\varphi}} & rooted Tseitin clause set & P2 prop-syntax-semantics & meta & \hyperref[sec:prop-syntax-semantics-01]{\ref*{sec:prop-syntax-semantics-01}} & \thepage & The defining clauses together with the unit clause containing the root representative; the root is necessary for the stated equisatisfiability theorem. \\
\phantomsection\label{entry:symbol-registry-00868}\(C_1,\ldots,C_m\) & \texttt{\detokenize{C_1,}\allowbreak\detokenize{\ldots,}\allowbreak\detokenize{C_m}} & ordered list of emitted Tseitin clauses & P2 prop-syntax-semantics & meta & \hyperref[sec:prop-syntax-semantics-01]{\ref*{sec:prop-syntax-semantics-01}} & \thepage & Node blocks follow preorder and displayed row order, with the root clause last; repeated emissions remain in this list although \(S_\varphi\) is a set. \\
\phantomsection\label{entry:symbol-registry-00869}\(T_\varphi\) & \texttt{\detokenize{T_\varphi}} & deterministic Tseitin CNF formula & P2 prop-syntax-semantics & meta & \hyperref[sec:prop-syntax-semantics-01]{\ref*{sec:prop-syntax-semantics-01}} & \thepage & The inherited right-associated conjunction of the ordered emitted clause interpretations; equisatisfiable with the input, without general logical equivalence on the enlarged vocabulary. \\
\phantomsection\label{entry:symbol-registry-00870}\(N,M,L;n_0,n_1,n_2\) & \texttt{\detokenize{N,}\allowbreak\detokenize{M,}\allowbreak\detokenize{L;n_0,}\allowbreak\detokenize{n_1,}\allowbreak\detokenize{n_2}} & Tseitin structural and encoding size parameters & P2 prop-syntax-semantics & meta & \hyperref[sec:prop-syntax-semantics-01]{\ref*{sec:prop-syntax-semantics-01}} & \thepage & \(N\) counts input-tree positions, \(M\) its nonvariable positions, \(L\) the primitive input bit length; \(n_0\) counts bottom occurrences, \(n_1\) negations and \(n_2\) binary connectives in annex:p2.prop.tseitin-size. \\
\phantomsection\label{entry:symbol-registry-01975}\(P\) & \texttt{\detokenize{P}} & countably infinite set of unary atoms & P2 syllogistic-inference & n/a & \hyperref[def:p02-syllogistic-language]{\ref*{def:p02-syllogistic-language}} & \thepage & No binary atoms or arbitrary Boolean term formation in the selected fragments. \\
\phantomsection\label{entry:symbol-registry-01976}\(p,q,o\) & \texttt{\detokenize{p,}\allowbreak\detokenize{q,}\allowbreak\detokenize{o}} & unary atom metavariables & P2 syllogistic-inference & n/a & \hyperref[def:p02-syllogistic-language]{\ref*{def:p02-syllogistic-language}} & \thepage & Atom substitutions send atoms to atoms. Named examples additionally declare t and a/b/c/d as atoms locally. \\
\phantomsection\label{entry:symbol-registry-01977}\(l,m,n\) & \texttt{\detokenize{l,}\allowbreak\detokenize{m,}\allowbreak\detokenize{n}} & unary literal metavariables & P2 syllogistic-inference & n/a & \hyperref[def:p02-syllogistic-language]{\ref*{def:p02-syllogistic-language}} & \thepage & Positive atoms or their complements; annex cardinality n is separately scoped. \\
\phantomsection\label{entry:symbol-registry-01978}\(\bar l\) & \texttt{\detokenize{\bar l}} & involutive complement of a unary literal & P2 syllogistic-inference & n/a & \hyperref[def:p02-syllogistic-language]{\ref*{def:p02-syllogistic-language}} & \thepage & The semantic complement is relative to the fixed nonempty universe. \\
\phantomsection\label{entry:symbol-registry-01979}\(\mathcal S,\mathcal S^\dagger\) & \texttt{\detokenize{\mathcal S,}\allowbreak\detokenize{\mathcal S^\dagger}} & sets of permitted syllogistic sentences & P2 syllogistic-inference & n/a & \hyperref[def:p02-syllogistic-language]{\ref*{def:p02-syllogistic-language}} & \thepage & Calligraphic language names distinguished from sans-serif rule-set names. \\
\phantomsection\label{entry:symbol-registry-01980}\(\forall(l,m),\exists(l,m)\) & \texttt{\detokenize{\forall(l,}\allowbreak\detokenize{m),}\allowbreak\detokenize{\exists(l,}\allowbreak\detokenize{m)}} & categorical sentence constructors & P2 syllogistic-inference & n/a & \hyperref[def:p02-syllogistic-language]{\ref*{def:p02-syllogistic-language}} & \thepage & Two literal arguments; not written bound-variable quantifiers and not formulas with arbitrary Boolean terms. \\
\phantomsection\label{entry:symbol-registry-01981}\(\mathfrak A,A\) & \texttt{\detokenize{\mathfrak A,}\allowbreak\detokenize{A}} & structure and its nonempty universe & P2 syllogistic-inference & n/a & \hyperref[def:p02-syllogistic-semantics]{\ref*{def:p02-syllogistic-semantics}} & \thepage & Predicate extensions may be empty. \\
\phantomsection\label{entry:symbol-registry-01982}\(l^{\mathfrak A}\) & \texttt{\detokenize{l^{\mathfrak A}}} & extension of a literal in a structure & P2 syllogistic-inference & n/a & \hyperref[def:p02-syllogistic-semantics]{\ref*{def:p02-syllogistic-semantics}} & \thepage & Annex rules proof explicitly abbreviates extensions by literal names in its first two soundness paragraphs only. \\
\phantomsection\label{entry:symbol-registry-01983}\(\Theta\models\theta\) & \texttt{\detokenize{\Theta\models\theta}} & semantic consequence in the named syllogistic fragment & P2 syllogistic-inference & n/a & \hyperref[def:p02-syllogistic-semantics]{\ref*{def:p02-syllogistic-semantics}} & \thepage & All structures have nonempty universes; the arbitrary premise set need not be finite. \\
\phantomsection\label{entry:symbol-registry-01984}\(\bar\varphi\) & \texttt{\detokenize{\bar\varphi}} & defined formula complement & P2 syllogistic-inference & n/a & \hyperref[def:p02-syllogistic-identifications]{\ref*{def:p02-syllogistic-identifications}} & \thepage & Swap sentence constructor and complement its second literal; this respects existential symmetry and universal contraposition. \\
\phantomsection\label{entry:symbol-registry-01985}\(\mathsf S,\mathsf S^\dagger\) & \texttt{\detokenize{\mathsf S,}\allowbreak\detokenize{\mathsf S^\dagger}} & finite source syllogistic rule systems & P2 syllogistic-inference & n/a & \hyperref[def:p02-syllogistic-derivation]{\ref*{def:p02-syllogistic-derivation}} & \thepage & Source schemas and their typed metavariables appear on the two exclusively named sheets. \\
\phantomsection\label{entry:symbol-registry-01986}\(\vdash_{\mathsf S},\vdash_{\mathsf S^\dagger}\) & \texttt{\detokenize{\vdash_{\mathsf S},}\allowbreak\detokenize{\vdash_{\mathsf S^\dagger}}} & finite direct derivability & P2 syllogistic-inference & n/a & \hyperref[def:p02-syllogistic-derivation]{\ref*{def:p02-syllogistic-derivation}} & \thepage & No assumption discharge or primitive reductio rule. \\
\phantomsection\label{entry:symbol-registry-01987}\(\Phi,\Psi\) & \texttt{\detokenize{\Phi,}\allowbreak\detokenize{\Psi}} & universal-premise set and existential-premise set & P2 syllogistic-inference & n/a & \hyperref[thm:p02-syllogistic-separation]{\ref*{thm:p02-syllogistic-separation}} & \thepage & Each is a subset of the named fragment. \\
\phantomsection\label{entry:symbol-registry-01988}\(V,V^\Phi\) & \texttt{\detokenize{V,}\allowbreak\detokenize{V^\Phi}} & literal set and its closure under BTA consequences & P2 syllogistic-inference & n/a & \hyperref[def:p02-syllogistic-literal-closure]{\ref*{def:p02-syllogistic-literal-closure}} & \thepage & Literal-set consistency/completeness differs from proof-system completeness. \\
\phantomsection\label{entry:symbol-registry-01989}\(\vdash_{\mathrm{BTA}}\) & \texttt{\detokenize{\vdash_{\mathrm{BTA}}}} & derivability using only dagger B/T/A rules & P2 syllogistic-inference & n/a & \hyperref[def:p02-syllogistic-literal-closure]{\ref*{def:p02-syllogistic-literal-closure}} & \thepage & No existential rule or reductio is used in this closure. \\
\phantomsection\label{entry:symbol-registry-01990}\(Q,n\) & \texttt{\detokenize{Q,}\allowbreak\detokenize{n}} & finite input atom set and its cardinality & P2 syllogistic saturation and annex:p2.syllogistic.saturation & n/a & \hyperref[def:p02-syllogistic-saturation]{\ref*{def:p02-syllogistic-saturation}} & \thepage & Here n is a natural number; literal metavariable n is used only in the earlier rule schemas. \\
\phantomsection\label{entry:symbol-registry-01991}\(a,b,r,u,v,w\) & \texttt{\detokenize{a,}\allowbreak\detokenize{b,}\allowbreak\detokenize{r,}\allowbreak\detokenize{u,}\allowbreak\detokenize{v,}\allowbreak\detokenize{w}} & locally declared literal metavariables in closure and completeness proofs & P2 syllogistic proof annex & n/a & \hyperref[annex:p2.syllogistic.extension]{\ref*{annex:p2.syllogistic.extension}} & \thepage & The existential completeness case explicitly declares a/b as literals; a in singleton-model paragraphs is a separately declared domain element. \\
\phantomsection\label{entry:symbol-registry-01992}\(L(x),M(x)\) & \texttt{\detokenize{L(x),}\allowbreak\detokenize{M(x)}} & first-order translations of two literals & P2 syllogistic-inference & n/a & \hyperref[thm:p02-syllogistic-semantic-identifications]{\ref*{thm:p02-syllogistic-semantic-identifications}} & \thepage & Each is p(x) or its ordinary negation, according to literal polarity. \\
\phantomsection\label{entry:symbol-registry-03764}\(h\) & \texttt{\detokenize{h}} & full equality-free formula to literal Hilbert-fragment translation & P2 Hilbert-FOL adequacy & n/a & \hyperref[def:p02-hfa-encoding]{\ref*{def:p02-hfa-encoding}} & \thepage & Memberwise action on premise sets; closed falsum anchor preserves free-variable conditions. \\
\phantomsection\label{entry:symbol-registry-03765}\(Q,T,F\) & \texttt{\detokenize{Q;T;F}} & chosen equality-free atom and closed tautology and its negation & P2 Hilbert-FOL adequacy encoding & n/a & \hyperref[def:p02-hfa-encoding]{\ref*{def:p02-hfa-encoding}} & \thepage & Formula abbreviations only; not predicate symbols added except in the explicitly declared atomless-signature expansion. \\
\phantomsection\label{entry:symbol-registry-05404}\(P\) & \texttt{\detokenize{P}} & Set of atomic symbols & P2 free Boolean syntax & object & \hyperref[def:p02-fbc-language]{\ref*{def:p02-fbc-language}} & \thepage & Arbitrary set; no enumeration or well-order is assumed \\
\phantomsection\label{entry:symbol-registry-05405}\(\operatorname{Form}(P)\) & \texttt{\detokenize{\operatorname{Form}(P)}} & Set of finite CPL formulas over P & P2 free Boolean syntax & object & \hyperref[def:p02-fbc-language]{\ref*{def:p02-fbc-language}} & \thepage & Nullary bottom is retained even when P is empty \\
\phantomsection\label{entry:symbol-registry-05406}\(\vdash_P\) & \texttt{\detokenize{\vdash_P}} & ND\_CPL derivability with permitted premises & P2 free Boolean syntax & relation & \hyperref[def:p02-fbc-language]{\ref*{def:p02-fbc-language}} & \thepage & Abbreviation of the existing arbitrary-alphabet relation; not a new connective \\
\phantomsection\label{entry:symbol-registry-05407}\(\operatorname{oa}(D)\) & \texttt{\detokenize{\operatorname{oa}(D)}} & Finite set of open-assumption formulas & P2 free Boolean syntax & object & \hyperref[def:p02-fbc-language]{\ref*{def:p02-fbc-language}} & \thepage & Distinct from the possibly infinite permitted premise set \\
\phantomsection\label{entry:symbol-registry-05408}\(\llbracket\alpha\rrbracket_f\) & \texttt{\detokenize{\llbracket\alpha\rrbracket_f}} & Boolean-algebra-valued recursive formula evaluation & P2 free Boolean evaluation & object & \hyperref[sec:p02-fbc-boolean-soundness]{\ref*{sec:p02-fbc-boolean-soundness}} & \thepage & Target may be trivial; implication is complement joined with consequent \\
\phantomsection\label{entry:symbol-registry-05409}\(\sim_P\) & \texttt{\detokenize{\sim_P}} & Provable-biconditional equivalence on formulas & P2 free Boolean quotient & relation & \hyperref[def:p02-fbc-quotient]{\ref*{def:p02-fbc-quotient}} & \thepage & Distinct from semantic equivalence and elementary equivalence \\
\phantomsection\label{entry:symbol-registry-05410}\([\alpha]_P\) & \texttt{\detokenize{[\alpha]_P}} & Syntactic provable-equivalence class & P2 free Boolean quotient & object & \hyperref[def:p02-fbc-quotient]{\ref*{def:p02-fbc-quotient}} & \thepage & No representative selector \\
\phantomsection\label{entry:symbol-registry-05411}\(L_P\) & \texttt{\detokenize{L_P}} & Syntactic free Boolean algebra & P2 free Boolean quotient & object & \hyperref[def:p02-fbc-quotient]{\ref*{def:p02-fbc-quotient}} & \thepage & Sethood and nontrivial Boolean laws are proved \\
\phantomsection\label{entry:symbol-registry-05412}\(0_L\) & \texttt{\detokenize{0_L}} & Quotient bottom & P2 free Boolean quotient & object & \hyperref[def:p02-fbc-quotient]{\ref*{def:p02-fbc-quotient}} & \thepage & Class of source bottom \\
\phantomsection\label{entry:symbol-registry-05413}\(1_L\) & \texttt{\detokenize{1_L}} & Quotient top & P2 free Boolean quotient & object & \hyperref[def:p02-fbc-quotient]{\ref*{def:p02-fbc-quotient}} & \thepage & Class of source topv equal to negation of bottom \\
\phantomsection\label{entry:symbol-registry-05414}\(\wedge\) & \texttt{\detokenize{\wedge}} & Boolean meet on algebra elements & P2 free Boolean quotient & infix & \hyperref[def:p02-fbc-quotient]{\ref*{def:p02-fbc-quotient}} & \thepage & Conjunction inside brackets retains its source logical type \\
\phantomsection\label{entry:symbol-registry-05415}\(\vee\) & \texttt{\detokenize{\vee}} & Boolean join on algebra elements & P2 free Boolean quotient & infix & \hyperref[def:p02-fbc-quotient]{\ref*{def:p02-fbc-quotient}} & \thepage & Disjunction inside brackets retains its source logical type \\
\phantomsection\label{entry:symbol-registry-05416}\(\neg\) & \texttt{\detokenize{\neg}} & Boolean complement on algebra elements & P2 free Boolean quotient & object & \hyperref[def:p02-fbc-quotient]{\ref*{def:p02-fbc-quotient}} & \thepage & Formula negation inside brackets retains the primitive syntax type \\
\phantomsection\label{entry:symbol-registry-05417}\(\eta_P\) & \texttt{\detokenize{\eta_P}} & Injection of atomic generators into the quotient & P2 free Boolean quotient & object & \hyperref[thm:p02-fbc-free]{\ref*{thm:p02-fbc-free}} & \thepage & Canonical map to classes; empty generators allowed \\
\phantomsection\label{entry:symbol-registry-05418}\(h_f\) & \texttt{\detokenize{h_f}} & Unique homomorphism extending an assignment & P2 free Boolean quotient & object & \hyperref[thm:p02-fbc-free]{\ref*{thm:p02-fbc-free}} & \thepage & Includes trivial target algebras \\
\phantomsection\label{entry:symbol-registry-05419}\(\leq_L\) & \texttt{\detokenize{\leq_L}} & Quotient lattice order & P2 free Boolean quotient & relation & \hyperref[thm:p02-fbc-free]{\ref*{thm:p02-fbc-free}} & \thepage & Equivalent to closed ND implication \\
\phantomsection\label{entry:symbol-registry-05420}\(F_\Gamma\) & \texttt{\detokenize{F_\Gamma}} & Filter generated by premise classes & P2 free Boolean premise filters & object & \hyperref[def:p02-fbc-premise-filter]{\ref*{def:p02-fbc-premise-filter}} & \thepage & Finite ordered witnesses include the empty sequence \\
\phantomsection\label{entry:symbol-registry-05421}\(H\) & \texttt{\detokenize{H}} & Explicit proper filter adjoining an excluded class complement & P2 free Boolean countervaluation & object & \hyperref[thm:p02-fbc-separation]{\ref*{thm:p02-fbc-separation}} & \thepage & Properness is proved before BPI extension \\
\phantomsection\label{entry:symbol-registry-05422}\(U\) & \texttt{\detokenize{U}} & One ultrafilter extending the explicit proper filter & P2 free Boolean countervaluation & object & \hyperref[thm:p02-fbc-separation]{\ref*{thm:p02-fbc-separation}} & \thepage & The sole unrestricted extension uses BPI \\
\phantomsection\label{entry:symbol-registry-05423}\(\chi_U\) & \texttt{\detokenize{\chi_U}} & Characteristic Boolean homomorphism to the two-element algebra & P2 free Boolean countervaluation & object & \hyperref[thm:p02-fbc-separation]{\ref*{thm:p02-fbc-separation}} & \thepage & Laws supplied by the proved P1 duality construction \\
\phantomsection\label{entry:symbol-registry-05424}\(v_U\) & \texttt{\detokenize{v_U}} & Total atomic assignment composed from the ultrafilter & P2 free Boolean countervaluation & object & \hyperref[thm:p02-fbc-separation]{\ref*{thm:p02-fbc-separation}} & \thepage & No separate choices for atoms \\
\phantomsection\label{entry:symbol-registry-05425}\(\widehat v_U\) & \texttt{\detokenize{\widehat v_U}} & Recursive two-valued formula evaluation & P2 free Boolean countervaluation & object & \hyperref[thm:p02-fbc-separation]{\ref*{thm:p02-fbc-separation}} & \thepage & All primitive constructors and defined top \\
\phantomsection\label{entry:symbol-registry-05426}\(\models_P\) & \texttt{\detokenize{\models_P}} & Two-valued semantic consequence over arbitrary atoms & P2 free Boolean countervaluation & relation & \hyperref[thm:p02-fbc-completeness]{\ref*{thm:p02-fbc-completeness}} & \thepage & Separate from derivability until both directions are proved \\
\phantomsection\label{entry:symbol-registry-06085}\ensuremath{\vdash} & \texttt{\detokenize{\vdash}} & Closed classical propositional assertion judgment & P2 classical rejection construction & n/a & \hyperref[def:p02-crf-judgments]{\ref*{def:p02-crf-judgments}} & \thepage & No open assumptions; positive derivations are the declared Hilbert calculus Original precedence descriptor 'judgment' preserved; n/a is not an invented rank. \\
\phantomsection\label{entry:symbol-registry-06086}\ensuremath{\dashv} & \texttt{\detokenize{\dashv}} & Formal rejection of classical validity & P2 classical rejection construction & n/a & \hyperref[def:p02-crf-judgments]{\ref*{def:p02-crf-judgments}} & \thepage & Not object-language negation or assertion of universal falsity Original precedence descriptor 'judgment' preserved; n/a is not an invented rank. \\
\phantomsection\label{entry:symbol-registry-06087}\ensuremath{A^{+} A^{-}} & \texttt{\detokenize{A^{+} A^{-}}} & Source signed metastatements of membership and nonmembership in a logic & P2 classical rejection construction & n/a & \hyperref[def:p02-crf-judgments]{\ref*{def:p02-crf-judgments}} & \thepage & Not an extension of the propositional formula grammar Original precedence descriptor 'metadata' preserved; n/a is not an invented rank. \\
\phantomsection\label{entry:symbol-registry-06088}\ensuremath{\sigma} & \texttt{\detokenize{\sigma}} & Simultaneous uniform substitution on propositional variables & P2 classical rejection construction & n/a & \hyperref[def:p02-crf-derivation]{\ref*{def:p02-crf-derivation}} & \thepage & Its rejection use is backward from an instance to the general formula Original precedence descriptor 'function' preserved; n/a is not an invented rank. \\
\phantomsection\label{entry:symbol-registry-06089}v w & \texttt{\detokenize{v w}} & Boolean valuations and substitution pullback & P2 classical rejection construction & n/a & \hyperref[thm:p02-crf-substitution]{\ref*{thm:p02-crf-substitution}} & \thepage & A negative line carries an existential countervaluation; different lines may use different valuations Original precedence descriptor 'function' preserved; n/a is not an invented rank. \\
\phantomsection\label{entry:symbol-registry-06090}T F & \texttt{\detokenize{T F}} & Constant-one and constant-zero formulas in the primitive Hilbert language & P2 classical rejection construction & n/a & \hyperref[def:p02-crf-constants]{\ref*{def:p02-crf-constants}} & \thepage & T is p to p; F is neg T; source p and not p is separately identified as the equivalent constant-zero formula Original precedence descriptor 'formula abbreviation' preserved; n/a is not an invented rank. \\
\phantomsection\label{entry:symbol-registry-06091}\ensuremath{\sigma_v} & \texttt{\detokenize{\sigma_v}} & Countervaluation-coded simultaneous formula substitution & P2 classical rejection construction & n/a & \hyperref[def:p02-crf-constants]{\ref*{def:p02-crf-constants}} & \thepage & Each variable is sent to T or F according to v; p need not be fresh Original precedence descriptor 'function' preserved; n/a is not an invented rank. \\
\phantomsection\label{entry:symbol-registry-06092}0 1 & \texttt{\detokenize{0 1}} & Boolean zero and one in the classical valuation range & P2 classical rejection construction & n/a & \hyperref[def:p02-crf-judgments]{\ref*{def:p02-crf-judgments}} & \thepage & Not object-language numerical constants; negation swaps the two values Original precedence descriptor 'semantic values' preserved; n/a is not an invented rank. \\
\phantomsection\label{entry:symbol-registry-06278}C & \texttt{\detokenize{C}} & finite set of first-order literals interpreted as one clause & P2 clausal composition & n/a & \hyperref[def:p02-clause-family-reading]{\ref*{def:p02-clause-family-reading}} & \thepage & The empty clause is false; not a formula string or a structure domain. \\
\phantomsection\label{entry:symbol-registry-06279}S & \texttt{\detokenize{S}} & finite family of universally interpreted first-order clauses & P2 clausal composition & n/a & \hyperref[def:p02-clause-family-reading]{\ref*{def:p02-clause-family-reading}} & \thepage & The empty family is true; universal closure is specified separately for each clause. \\
\phantomsection\label{entry:symbol-registry-06280}L & \texttt{\detokenize{L}} & first-order literal belonging to a finite clause & P2 clausal composition & n/a & \hyperref[def:p02-clause-family-reading]{\ref*{def:p02-clause-family-reading}} & \thepage & An atom or its negation, not the truth constants removed by the conversion. \\
\phantomsection\label{entry:symbol-registry-06281}D & \texttt{\detokenize{D}} & second finite clause in the distributive Cartesian-product construction & P2 clausal composition & n/a & \hyperref[def:p02-matrix-clause-construction]{\ref*{def:p02-matrix-clause-construction}} & \thepage & Union with the first clause forms one disjunctive output clause; not a domain or a truth-value designated set. \\
\phantomsection\label{entry:symbol-registry-06282}\ensuremath{\operatorname{Cl}} & \texttt{\detokenize{\operatorname{Cl}}} & recursive map from quantifier-free NNF matrices to finite clause families & P2 clausal composition & n/a & \hyperref[def:p02-matrix-clause-construction]{\ref*{def:p02-matrix-clause-construction}} & \thepage & Explicit editorial distributive construction; no new predicates and no linear-size claim. \\
\phantomsection\label{entry:symbol-registry-06283}\ensuremath{\mathcal L^*} & \texttt{\detokenize{\mathcal L^*}} & finite Skolem-symbol expansion of the input first-order signature & P2 clausal composition & n/a & \hyperref[thm:p02-fol-clausal-form]{\ref*{thm:p02-fol-clausal-form}} & \thepage & The new symbols are finite in number; the original signature may be infinite. \\
\phantomsection\label{entry:symbol-registry-06528}\(r\) & \texttt{\detokenize{r}} & term metavariables symbol family member \(r\) & P2 fol-syntax-semantics & n/a & \hyperref[sec:fol-syntax-semantics-01]{\ref*{sec:fol-syntax-semantics-01}} & \thepage & The chapter declaration adds u to the global term-metavariable row r,t; x,y,z stay object-language variables and s stays a variable assignment. Explicit displayed family; member types and common actual first-use exposed, not one atomic glyph. Family listing, not one atomic glyph. Author macro \texttt{\detokenize{r,}\allowbreak\detokenize{t,}\allowbreak\detokenize{u}} not carried to the members. \\
\phantomsection\label{entry:symbol-registry-06529}\(t\) & \texttt{\detokenize{t}} & term metavariables symbol family member \(t\) & P2 fol-syntax-semantics & n/a & \hyperref[sec:fol-syntax-semantics-01]{\ref*{sec:fol-syntax-semantics-01}} & \thepage & The chapter declaration adds u to the global term-metavariable row r,t; x,y,z stay object-language variables and s stays a variable assignment. Explicit displayed family; member types and common actual first-use exposed, not one atomic glyph. Family listing, not one atomic glyph. Author macro \texttt{\detokenize{r,}\allowbreak\detokenize{t,}\allowbreak\detokenize{u}} not carried to the members. \\
\phantomsection\label{entry:symbol-registry-06530}\(u\) & \texttt{\detokenize{u}} & term metavariables symbol family member \(u\) & P2 fol-syntax-semantics & n/a & \hyperref[sec:fol-syntax-semantics-01]{\ref*{sec:fol-syntax-semantics-01}} & \thepage & The chapter declaration adds u to the global term-metavariable row r,t; x,y,z stay object-language variables and s stays a variable assignment. Explicit displayed family; member types and common actual first-use exposed, not one atomic glyph. Family listing, not one atomic glyph. Author macro \texttt{\detokenize{r,}\allowbreak\detokenize{t,}\allowbreak\detokenize{u}} not carried to the members. \\
\phantomsection\label{entry:symbol-registry-06531}\(\alpha\) & \texttt{\detokenize{\alpha}} & propositional formula metavariables symbol family member \(\alpha\) & P2 prop-syntax-semantics & n/a & \hyperref[sec:p02-propositional-substitution]{\ref*{sec:p02-propositional-substitution}} & \thepage & Used without separate declaration in the substitution, functional-completeness and Tseitin passages of the propositional chapter; same reading as the global phi,psi,chi there. Explicit displayed family; member types and common actual first-use exposed, not one atomic glyph. Family listing, not one atomic glyph. Author macro \texttt{\detokenize{\alpha,}\allowbreak\detokenize{\beta,}\allowbreak\detokenize{\gamma}} not carried to the members. \\
\phantomsection\label{entry:symbol-registry-06532}\(\beta\) & \texttt{\detokenize{\beta}} & propositional formula metavariables symbol family member \(\beta\) & P2 prop-syntax-semantics & n/a & \hyperref[sec:p02-propositional-substitution]{\ref*{sec:p02-propositional-substitution}} & \thepage & Used without separate declaration in the substitution, functional-completeness and Tseitin passages of the propositional chapter; same reading as the global phi,psi,chi there. Explicit displayed family; member types and common actual first-use exposed, not one atomic glyph. Family listing, not one atomic glyph. Author macro \texttt{\detokenize{\alpha,}\allowbreak\detokenize{\beta,}\allowbreak\detokenize{\gamma}} not carried to the members. \\
\phantomsection\label{entry:symbol-registry-06533}\(\gamma\) & \texttt{\detokenize{\gamma}} & propositional formula metavariables symbol family member \(\gamma\) & P2 prop-syntax-semantics & n/a & \hyperref[sec:p02-propositional-substitution]{\ref*{sec:p02-propositional-substitution}} & \thepage & Used without separate declaration in the substitution, functional-completeness and Tseitin passages of the propositional chapter; same reading as the global phi,psi,chi there. Explicit displayed family; member types and common actual first-use exposed, not one atomic glyph. Family listing, not one atomic glyph. Author macro \texttt{\detokenize{\alpha,}\allowbreak\detokenize{\beta,}\allowbreak\detokenize{\gamma}} not carried to the members. \\
\phantomsection\label{entry:symbol-registry-06534}\(\alpha\) & \texttt{\detokenize{\alpha}} & propositional formula metavariables symbol family member \(\alpha\) & \hyperref[sec:p02-hilbert-cpl]{\ref*{sec:p02-hilbert-cpl}} & n/a & \hyperref[def:p02-hilbert-translation]{\ref*{def:p02-hilbert-translation}} & \thepage & Range over full-language CPL formulas in the clauses of tr\_H; gamma ranges over members of Gamma. Explicit displayed family; member types and common actual first-use exposed, not one atomic glyph. Family listing, not one atomic glyph. Author macro \texttt{\detokenize{\alpha,}\allowbreak\detokenize{\beta,}\allowbreak\detokenize{\gamma}} not carried to the members. \\
\phantomsection\label{entry:symbol-registry-06535}\(\beta\) & \texttt{\detokenize{\beta}} & propositional formula metavariables symbol family member \(\beta\) & \hyperref[sec:p02-hilbert-cpl]{\ref*{sec:p02-hilbert-cpl}} & n/a & \hyperref[def:p02-hilbert-translation]{\ref*{def:p02-hilbert-translation}} & \thepage & Range over full-language CPL formulas in the clauses of tr\_H; gamma ranges over members of Gamma. Explicit displayed family; member types and common actual first-use exposed, not one atomic glyph. Family listing, not one atomic glyph. Author macro \texttt{\detokenize{\alpha,}\allowbreak\detokenize{\beta,}\allowbreak\detokenize{\gamma}} not carried to the members. \\
\phantomsection\label{entry:symbol-registry-06536}\(\gamma\) & \texttt{\detokenize{\gamma}} & propositional formula metavariables symbol family member \(\gamma\) & \hyperref[sec:p02-hilbert-cpl]{\ref*{sec:p02-hilbert-cpl}} & n/a & \hyperref[def:p02-hilbert-translation]{\ref*{def:p02-hilbert-translation}} & \thepage & Range over full-language CPL formulas in the clauses of tr\_H; gamma ranges over members of Gamma. Explicit displayed family; member types and common actual first-use exposed, not one atomic glyph. Family listing, not one atomic glyph. Author macro \texttt{\detokenize{\alpha,}\allowbreak\detokenize{\beta,}\allowbreak\detokenize{\gamma}} not carried to the members. \\
\phantomsection\label{entry:symbol-registry-06537}\(\alpha\) & \texttt{\detokenize{\alpha}} & propositional formula metavariable & \hyperref[sec:p02-analytic-tableaux]{\ref*{sec:p02-analytic-tableaux}} & n/a & \hyperref[def:p02-cpl-tableau]{\ref*{def:p02-cpl-tableau}} & \thepage & Closure condition: a branch containing both alpha and its negation. \\
\phantomsection\label{entry:symbol-registry-06538}\(\alpha\) & \texttt{\detokenize{\alpha}} & NNF formula metavariables symbol family member \(\alpha\) & \hyperref[sec:p02-propositional-duality]{\ref*{sec:p02-propositional-duality}} & n/a & \hyperref[def:p02-nnf-dual]{\ref*{def:p02-nnf-dual}} & \thepage & Category symbol of the metalinguistic NNF grammar; alpha-d is the dual operation of the registry row in P2 prop-syntax-semantics. Explicit displayed family; member types and common actual first-use exposed, not one atomic glyph. Family listing, not one atomic glyph. Author macro \texttt{\detokenize{\alpha,}\allowbreak\detokenize{\beta}} not carried to the members. \\
\phantomsection\label{entry:symbol-registry-06539}\(\beta\) & \texttt{\detokenize{\beta}} & NNF formula metavariables symbol family member \(\beta\) & \hyperref[sec:p02-propositional-duality]{\ref*{sec:p02-propositional-duality}} & n/a & \hyperref[def:p02-nnf-dual]{\ref*{def:p02-nnf-dual}} & \thepage & Category symbol of the metalinguistic NNF grammar; alpha-d is the dual operation of the registry row in P2 prop-syntax-semantics. Explicit displayed family; member types and common actual first-use exposed, not one atomic glyph. Family listing, not one atomic glyph. Author macro \texttt{\detokenize{\alpha,}\allowbreak\detokenize{\beta}} not carried to the members. \\
\phantomsection\label{entry:symbol-registry-06540}\(\alpha\) & \texttt{\detokenize{\alpha}} & formula metavariables over Form(P) symbol family member \(\alpha\) & P2 free Boolean syntax & n/a & \hyperref[def:p02-fbc-language]{\ref*{def:p02-fbc-language}} & \thepage & Overlaps the region of the type-metavariable row recorded for P2 prop-metatheory, whose declaration is in sec:simple-type-theory-01; see the return file. Explicit displayed family; member types and common actual first-use exposed, not one atomic glyph. Family listing, not one atomic glyph. Author macro \texttt{\detokenize{\alpha,}\allowbreak\detokenize{\beta,}\allowbreak\detokenize{\gamma}} not carried to the members. \\
\phantomsection\label{entry:symbol-registry-06541}\(\beta\) & \texttt{\detokenize{\beta}} & formula metavariables over Form(P) symbol family member \(\beta\) & P2 free Boolean syntax & n/a & \hyperref[def:p02-fbc-language]{\ref*{def:p02-fbc-language}} & \thepage & Overlaps the region of the type-metavariable row recorded for P2 prop-metatheory, whose declaration is in sec:simple-type-theory-01; see the return file. Explicit displayed family; member types and common actual first-use exposed, not one atomic glyph. Family listing, not one atomic glyph. Author macro \texttt{\detokenize{\alpha,}\allowbreak\detokenize{\beta,}\allowbreak\detokenize{\gamma}} not carried to the members. \\
\phantomsection\label{entry:symbol-registry-06542}\(\gamma\) & \texttt{\detokenize{\gamma}} & formula metavariables over Form(P) symbol family member \(\gamma\) & P2 free Boolean syntax & n/a & \hyperref[def:p02-fbc-language]{\ref*{def:p02-fbc-language}} & \thepage & Overlaps the region of the type-metavariable row recorded for P2 prop-metatheory, whose declaration is in sec:simple-type-theory-01; see the return file. Explicit displayed family; member types and common actual first-use exposed, not one atomic glyph. Family listing, not one atomic glyph. Author macro \texttt{\detokenize{\alpha,}\allowbreak\detokenize{\beta,}\allowbreak\detokenize{\gamma}} not carried to the members. \\
\phantomsection\label{entry:symbol-registry-06543}\(\alpha\) & \texttt{\detokenize{\alpha}} & quantifier-free NNF matrix metavariables symbol family member \(\alpha\) & \hyperref[sec:p02-fol-clausal-pipeline]{\ref*{sec:p02-fol-clausal-pipeline}} & n/a & \hyperref[def:p02-matrix-clause-construction]{\ref*{def:p02-matrix-clause-construction}} & \thepage & Immediate subformulas in the recursion defining Cl. Explicit displayed family; member types and common actual first-use exposed, not one atomic glyph. Family listing, not one atomic glyph. Author macro \texttt{\detokenize{\alpha,}\allowbreak\detokenize{\beta}} not carried to the members. \\
\phantomsection\label{entry:symbol-registry-06544}\(\beta\) & \texttt{\detokenize{\beta}} & quantifier-free NNF matrix metavariables symbol family member \(\beta\) & \hyperref[sec:p02-fol-clausal-pipeline]{\ref*{sec:p02-fol-clausal-pipeline}} & n/a & \hyperref[def:p02-matrix-clause-construction]{\ref*{def:p02-matrix-clause-construction}} & \thepage & Immediate subformulas in the recursion defining Cl. Explicit displayed family; member types and common actual first-use exposed, not one atomic glyph. Family listing, not one atomic glyph. Author macro \texttt{\detokenize{\alpha,}\allowbreak\detokenize{\beta}} not carried to the members. \\
\phantomsection\label{entry:symbol-registry-06545}\(A\) & \texttt{\detokenize{A}} & formula metavariables of the negation-implication language symbol family member \(A\) & P2 classical rejection construction & n/a & \hyperref[def:p02-crf-judgments]{\ref*{def:p02-crf-judgments}} & \thepage & Schematic letters of the assertion, rejection and substitution clauses and of sheet:p02-classical-rejection. Explicit displayed family; member types and common actual first-use exposed, not one atomic glyph. Family listing, not one atomic glyph. Author macro \texttt{\detokenize{A,}\allowbreak\detokenize{B}} not carried to the members. \\
\phantomsection\label{entry:symbol-registry-06546}\(B\) & \texttt{\detokenize{B}} & formula metavariables of the negation-implication language symbol family member \(B\) & P2 classical rejection construction & n/a & \hyperref[def:p02-crf-judgments]{\ref*{def:p02-crf-judgments}} & \thepage & Schematic letters of the assertion, rejection and substitution clauses and of sheet:p02-classical-rejection. Explicit displayed family; member types and common actual first-use exposed, not one atomic glyph. Family listing, not one atomic glyph. Author macro \texttt{\detokenize{A,}\allowbreak\detokenize{B}} not carried to the members. \\
\phantomsection\label{entry:symbol-registry-06547}\(A\) & \texttt{\detokenize{A}} & Hilbert-fragment formula metavariables symbol family member \(A\) & \hyperref[sec:p02-hilbert-fol-adequacy]{\ref*{sec:p02-hilbert-fol-adequacy}} & n/a & \hyperref[def:p02-hfa-setting]{\ref*{def:p02-hfa-setting}} & \thepage & Equality-free formulas built with negation, implication and the two quantifiers. Explicit displayed family; member types and common actual first-use exposed, not one atomic glyph. Family listing, not one atomic glyph. Author macro \texttt{\detokenize{A,}\allowbreak\detokenize{B}} not carried to the members. \\
\phantomsection\label{entry:symbol-registry-06548}\(B\) & \texttt{\detokenize{B}} & Hilbert-fragment formula metavariables symbol family member \(B\) & \hyperref[sec:p02-hilbert-fol-adequacy]{\ref*{sec:p02-hilbert-fol-adequacy}} & n/a & \hyperref[def:p02-hfa-setting]{\ref*{def:p02-hfa-setting}} & \thepage & Equality-free formulas built with negation, implication and the two quantifiers. Explicit displayed family; member types and common actual first-use exposed, not one atomic glyph. Family listing, not one atomic glyph. Author macro \texttt{\detokenize{A,}\allowbreak\detokenize{B}} not carried to the members. \\
\phantomsection\label{entry:symbol-registry-06549}\(\Gamma\nvdash_{S}\varphi\) & \texttt{\detokenize{\notderives,}\allowbreak\detokenize{\nvdash}} & negated derivability: no derivation in the named calculus & P2 prop-metatheory & meta & \hyperref[sec:prop-metatheory-02]{\ref*{sec:prop-metatheory-02}} & \thepage & Negation of the global per-system turnstile row; used for consistency and in contrapositive completeness arguments; also written with an unsubscripted turnstile under a declared local abbreviation. \\
\phantomsection\label{entry:symbol-registry-06550}\(\Gamma\nvdash_{S}\varphi\) & \texttt{\detokenize{\notderives,}\allowbreak\detokenize{\nvdash}} & negated derivability: no derivation in the named calculus & \hyperref[sec:fol-metatheory-01]{\ref*{sec:fol-metatheory-01}}--\hyperref[sec:fol-metatheory-03]{\ref*{sec:fol-metatheory-03}} & meta & \hyperref[sec:fol-metatheory-01]{\ref*{sec:fol-metatheory-01}} & \thepage & First-order consistency and the contrapositive completeness argument; same reading as the propositional row. \\
\phantomsection\label{entry:symbol-registry-06553}\(\varphi\equiv\psi\) & \texttt{\detokenize{\equiv}} & source first-order logical equivalence of formulas & \hyperref[sec:fol-semantics-02]{\ref*{sec:fol-semantics-02}}--\hyperref[sec:fol-cardinality-normalforms-06]{\ref*{sec:fol-cardinality-normalforms-06}} & meta & \hyperref[sec:fol-semantics-02]{\ref*{sec:fol-semantics-02}} & \thepage & Holds iff every L-structure and assignment satisfy both formulas or neither; metalinguistic, not alpha-equivalence and not string identity; used in sec:fol-cardinality-normalforms-03. Same relation as the editorial \(\varphi\Leftrightarrow\psi\) on first-order formulas; typed apart from \(\mathcal M\equiv\mathcal N\) by its formula operands (def:p01-w02-p01-025-equivalence-roles). \\
\phantomsection\label{entry:symbol-registry-06556}\(\varphi\equiv_\alpha\psi\) & \texttt{\detokenize{\equiv_\alpha}} & source alpha-equivalence of first-order formulas & \hyperref[sec:fol-syntax-semantics-04]{\ref*{sec:fol-syntax-semantics-04}}--\hyperref[sec:fol-cardinality-normalforms-06]{\ref*{sec:fol-cardinality-normalforms-06}} & meta & \hyperref[sec:fol-syntax-semantics-04]{\ref*{sec:fol-syntax-semantics-04}} & \thepage & Least constructor congruence containing \(Qx\,\varphi\equiv_\alpha Qy\,(\varphi[y/x])\) for \(y\notin\operatorname{Var}(\varphi)\); the same relation as the registered \(=_{\alpha}\). Unrelated to the P3 Morley ordinal-indexed back-and-forth relation written \(\equiv_\alpha\) (def:p03-morley-approximants). \\
\phantomsection\label{entry:symbol-registry-06655}\(B_\Gamma\) & \texttt{\detokenize{B_\Gamma}} & root branch of a propositional tableau & \hyperref[sec:p02-w02-p02-014-tableau-metatheory]{\ref*{sec:p02-w02-p02-014-tableau-metatheory}} & n/a & \hyperref[def:p02-w02-p02-014-input]{\ref*{def:p02-w02-p02-014-input}} & \thepage & The list of the members of the finite set \(\Gamma\), each once, with no expansion marks. \\
\phantomsection\label{entry:symbol-registry-06656}\(B=(\langle\theta_1,\ldots,\theta_m\rangle,M)\) & \texttt{\detokenize{B,}\allowbreak\detokenize{\theta_i,}\allowbreak\detokenize{M}} & tableau branch as formula list with marked positions & \hyperref[sec:p02-w02-p02-014-tableau-metatheory]{\ref*{sec:p02-w02-p02-014-tableau-metatheory}} & n/a & \hyperref[def:p02-w02-p02-014-expandable]{\ref*{def:p02-w02-p02-014-expandable}} & \thepage & \(M\) is a set of positions holding expandable formulas; an occurrence is a position, so equal formulas at two positions are two occurrences. \(C\) names a child branch. \\
\phantomsection\label{entry:symbol-registry-06657}\(\operatorname{fm}(B)\) & \texttt{\detokenize{\operatorname{fm}}} & set of formulas on a branch & \hyperref[sec:p02-w02-p02-014-tableau-metatheory]{\ref*{sec:p02-w02-p02-014-tableau-metatheory}} & meta & \hyperref[def:p02-w02-p02-014-expandable]{\ref*{def:p02-w02-p02-014-expandable}} & \thepage & Chosen instead of \(|B|\) because the registry reserves \(|\cdot|\) for cardinality. \\
\phantomsection\label{entry:symbol-registry-06658}\(A\) & \texttt{\detokenize{A}} & alternative of a tableau rule and the separator between alternatives symbol family member \(A\) & \hyperref[sec:p02-w02-p02-014-tableau-metatheory]{\ref*{sec:p02-w02-p02-014-tableau-metatheory}} & n/a & \hyperref[def:p02-w02-p02-014-expandable]{\ref*{def:p02-w02-p02-014-expandable}} & \thepage & Linear rules have one alternative and branching rules two; the bar is the sheet's own separator on sheet:tableaux-cpl. Explicit displayed family; member types and common actual first-use exposed, not one atomic glyph. Family listing, not one atomic glyph. Author macro \texttt{\detokenize{A,}\allowbreak\detokenize{\mid}} not carried to the members. \\
\phantomsection\label{entry:symbol-registry-06659}\(\ \langle\ldots\rangle\mid\langle\ldots\rangle\) & \texttt{\detokenize{\ \langle\ldots\rangle\mid\langle\ldots\rangle}} & alternative of a tableau rule and the separator between alternatives symbol family member \(\ \langle\ldots\rangle\mid\langle\ldots\rangle\) & \hyperref[sec:p02-w02-p02-014-tableau-metatheory]{\ref*{sec:p02-w02-p02-014-tableau-metatheory}} & n/a & \hyperref[def:p02-w02-p02-014-expandable]{\ref*{def:p02-w02-p02-014-expandable}} & \thepage & Linear rules have one alternative and branching rules two; the bar is the sheet's own separator on sheet:tableaux-cpl. Explicit displayed family; member types and common actual first-use exposed, not one atomic glyph. Family listing, not one atomic glyph. Author macro \texttt{\detokenize{A,}\allowbreak\detokenize{\mid}} not carried to the members. \\
\phantomsection\label{entry:symbol-registry-06660}\(s\frown A\) & \texttt{\detokenize{\frown}} & list concatenation appending an alternative to a branch list & \hyperref[sec:p02-w02-p02-014-tableau-metatheory]{\ref*{sec:p02-w02-p02-014-tableau-metatheory}} & n/a & \hyperref[def:p02-w02-p02-014-expandable]{\ref*{def:p02-w02-p02-014-expandable}} & \thepage & Same glyph as the registered P6 stem concatenation \(s\mathbin{\frown}t\); same typed role (concatenation of finite sequences) in a disjoint scope. \\
\phantomsection\label{entry:symbol-registry-06661}\(T_0,\ldots,T_k\) & \texttt{\detokenize{T_j}} & expansion sequence of finite lists of branches & \hyperref[sec:p02-w02-p02-014-tableau-metatheory]{\ref*{sec:p02-w02-p02-014-tableau-metatheory}} & n/a & \hyperref[def:p02-w02-p02-014-expandable]{\ref*{def:p02-w02-p02-014-expandable}} & \thepage & A tableau rooted at \(\Gamma\) is the last member; \(T_0\) is the one-branch list of \(B_\Gamma\). Sequence glyph kept as one row; the ellipsis is not a member symbol. \\
\phantomsection\label{entry:symbol-registry-06662}\(\operatorname{sz}(\chi)\) & \texttt{\detokenize{\operatorname{sz}}} & propositional formula size & \hyperref[sec:p02-w02-p02-014-tableau-metatheory]{\ref*{sec:p02-w02-p02-014-tableau-metatheory}} & meta & \hyperref[sec:p02-w02-p02-014-input]{\ref*{sec:p02-w02-p02-014-input}} & \thepage & Number of occurrences of variables, falsum and connectives; parentheses are not counted. The P12 row uses the same macro for term-tree node count in its own scope. \\
\phantomsection\label{entry:symbol-registry-06663}\(\mu(B)\) & \texttt{\detokenize{\mu(B)}} & termination measure of a tableau branch and its root value symbol family member \(\mu(B)\) & \hyperref[sec:p02-w02-p02-014-tableau-metatheory]{\ref*{sec:p02-w02-p02-014-tableau-metatheory}} & meta & \hyperref[sec:p02-w02-p02-014-input]{\ref*{sec:p02-w02-p02-014-input}} & \thepage & Sum of \(3^{\operatorname{sz}}\) over unmarked positions holding expandable formulas; zero on an open branch exactly when it is saturated. Other registered \(\mu\) rows are in disjoint P8, P9, P15 and P17 scopes. Explicit displayed family; member types and common actual first-use exposed, not one atomic glyph. Family listing, not one atomic glyph. Author macro \texttt{\detokenize{\mu}} not carried to the members. \\
\phantomsection\label{entry:symbol-registry-06664}\(\mu_0\) & \texttt{\detokenize{\mu_0}} & termination measure of a tableau branch and its root value symbol family member \(\mu_0\) & \hyperref[sec:p02-w02-p02-014-tableau-metatheory]{\ref*{sec:p02-w02-p02-014-tableau-metatheory}} & meta & \hyperref[sec:p02-w02-p02-014-input]{\ref*{sec:p02-w02-p02-014-input}} & \thepage & Sum of \(3^{\operatorname{sz}}\) over unmarked positions holding expandable formulas; zero on an open branch exactly when it is saturated. Other registered \(\mu\) rows are in disjoint P8, P9, P15 and P17 scopes. Explicit displayed family; member types and common actual first-use exposed, not one atomic glyph. Family listing, not one atomic glyph. Author macro \texttt{\detokenize{\mu}} not carried to the members. \\
\phantomsection\label{entry:symbol-registry-06665}\(v_B\) & \texttt{\detokenize{v_B}} & branch valuation of a tableau branch & \hyperref[sec:p02-w02-p02-014-tableau-metatheory]{\ref*{sec:p02-w02-p02-014-tableau-metatheory}} & n/a & \hyperref[sec:p02-w02-p02-014-input]{\ref*{sec:p02-w02-p02-014-input}} & \thepage & True exactly at the propositional variables that are formulas on \(B\). The P15 and P19 rows using \(v_B\) are in disjoint scopes with different roles. \\
\phantomsection\label{entry:symbol-registry-06666}\(\mathrm{CT}_{\Gamma}\) & \texttt{\detokenize{\mathrm{CT}_{\Gamma}}} & set of closed terms of the parameter-extended root signature & \hyperref[sec:p02-w02-p02-014-fol-calculus]{\ref*{sec:p02-w02-p02-014-fol-calculus}} & n/a & \hyperref[def:p02-w02-p02-014-fol-setting]{\ref*{def:p02-w02-p02-014-fol-setting}} & \thepage & Closed \(\mathcal L_{\Gamma}^{*}\)-terms; countably infinite; enumerated as \(t_0,t_1,\ldots\). \(\mathcal L_{\Gamma}\), \(\mathcal L_{\Gamma}^{*}\) and the parameters \(a_n\) are those of def:p02-fol-tableau. \\
\phantomsection\label{entry:symbol-registry-06667}\(t_0,t_1,\ldots\) & \texttt{\detokenize{t_k}} & fixed computable enumeration of the closed root terms & \hyperref[sec:p02-w02-p02-014-fol-calculus]{\ref*{sec:p02-w02-p02-014-fol-calculus}} & n/a & \hyperref[def:p02-w02-p02-014-fol-setting]{\ref*{def:p02-w02-p02-014-fol-setting}} & \thepage & Every member of \(\mathrm{CT}_{\Gamma}\) occurs; the index \(k\) schedules gamma and reflexivity tasks in the systematic procedure. Sequence glyph kept as one row; the ellipsis is not a member symbol. \\
\phantomsection\label{entry:symbol-registry-06668}\(\mathsf T\) & \texttt{\detokenize{\mathsf T}} & the equality-free first-order tableau calculus and its extension by equality rules symbol family member \(\mathsf T\) & \hyperref[sec:p02-w02-p02-014-fol-calculus]{\ref*{sec:p02-w02-p02-014-fol-calculus}} & n/a & \hyperref[def:p02-w02-p02-014-fol-setting]{\ref*{def:p02-w02-p02-014-fol-setting}} & \thepage & \(\mathsf T\) is the calculus of sheet:tableaux-fol; \(\mathsf T^{=}\) adds reflexivity and replacement (def:p02-w02-p02-014-fol-equality). Explicit displayed family; member types and common actual first-use exposed, not one atomic glyph. Family listing, not one atomic glyph. Author macro \texttt{\detokenize{\mathsf T,}\allowbreak\detokenize{\mathsf T\^{}{=}}} not carried to the members. \\
\phantomsection\label{entry:symbol-registry-06669}\(\ \mathsf T^{=}\) & \texttt{\detokenize{\ \mathsf T^{=}}} & the equality-free first-order tableau calculus and its extension by equality rules symbol family member \(\ \mathsf T^{=}\) & \hyperref[sec:p02-w02-p02-014-fol-calculus]{\ref*{sec:p02-w02-p02-014-fol-calculus}} & n/a & \hyperref[def:p02-w02-p02-014-fol-setting]{\ref*{def:p02-w02-p02-014-fol-setting}} & \thepage & \(\mathsf T\) is the calculus of sheet:tableaux-fol; \(\mathsf T^{=}\) adds reflexivity and replacement (def:p02-w02-p02-014-fol-equality). Explicit displayed family; member types and common actual first-use exposed, not one atomic glyph. Family listing, not one atomic glyph. Author macro \texttt{\detokenize{\mathsf T,}\allowbreak\detokenize{\mathsf T\^{}{=}}} not carried to the members. \\
\phantomsection\label{entry:symbol-registry-06670}\(x_0\) & \texttt{\detokenize{x_0}} & the fixed variable marking the replaced occurrence in the replacement rule & \hyperref[def:p02-w02-p02-014-fol-equality]{\ref*{def:p02-w02-p02-014-fol-equality}} & n/a & \hyperref[def:p02-w02-p02-014-fol-equality]{\ref*{def:p02-w02-p02-014-fol-equality}} & \thepage & First variable of the fixed order on Var; occurs exactly once in the atomic formula \(\varphi\) and no other variable occurs. \\
\phantomsection\label{entry:symbol-registry-06671}\(D\) & \texttt{\detokenize{D}} & set of done tasks of a first-order branch & \hyperref[sec:p02-w02-p02-014-fol-calculus]{\ref*{sec:p02-w02-p02-014-fol-calculus}} & n/a & \hyperref[def:p02-w02-p02-014-fol-branch]{\ref*{def:p02-w02-p02-014-fol-branch}} & \thepage & Generalizes the marked-position set \(M\) of the propositional branches. \\
\phantomsection\label{entry:symbol-registry-06672}\(\langle i\rangle\) & \texttt{\detokenize{\langle i\rangle}} & tasks of a first-order branch symbol family member \(\langle i\rangle\) & \hyperref[sec:p02-w02-p02-014-fol-calculus]{\ref*{sec:p02-w02-p02-014-fol-calculus}} & n/a & \hyperref[def:p02-w02-p02-014-fol-branch]{\ref*{def:p02-w02-p02-014-fol-branch}} & \thepage & Propositional or delta expansion at position \(i\); gamma instance at term \(t\); reflexivity at \(t\); replacement of an occurrence in position \(j\) using the equation at position \(i\). Angle brackets also denote finite lists in the same section; the task forms are distinguished by their entries. Explicit displayed family; member types and common actual first-use exposed, not one atomic glyph. Family listing, not one atomic glyph. Author macro \texttt{\detokenize{\langle\ \rangle}} not carried to the members. \\
\phantomsection\label{entry:symbol-registry-06673}\(\ \langle i,t\rangle\) & \texttt{\detokenize{\ \langle i,}\allowbreak\detokenize{t\rangle}} & tasks of a first-order branch symbol family member \(\ \langle i,t\rangle\) & \hyperref[sec:p02-w02-p02-014-fol-calculus]{\ref*{sec:p02-w02-p02-014-fol-calculus}} & n/a & \hyperref[def:p02-w02-p02-014-fol-branch]{\ref*{def:p02-w02-p02-014-fol-branch}} & \thepage & Propositional or delta expansion at position \(i\); gamma instance at term \(t\); reflexivity at \(t\); replacement of an occurrence in position \(j\) using the equation at position \(i\). Angle brackets also denote finite lists in the same section; the task forms are distinguished by their entries. Explicit displayed family; member types and common actual first-use exposed, not one atomic glyph. Family listing, not one atomic glyph. Author macro \texttt{\detokenize{\langle\ \rangle}} not carried to the members. \\
\phantomsection\label{entry:symbol-registry-06674}\(\ \langle{=},t\rangle\) & \texttt{\detokenize{\ \langle{=},}\allowbreak\detokenize{t\rangle}} & tasks of a first-order branch symbol family member \(\ \langle{=},t\rangle\) & \hyperref[sec:p02-w02-p02-014-fol-calculus]{\ref*{sec:p02-w02-p02-014-fol-calculus}} & n/a & \hyperref[def:p02-w02-p02-014-fol-branch]{\ref*{def:p02-w02-p02-014-fol-branch}} & \thepage & Propositional or delta expansion at position \(i\); gamma instance at term \(t\); reflexivity at \(t\); replacement of an occurrence in position \(j\) using the equation at position \(i\). Angle brackets also denote finite lists in the same section; the task forms are distinguished by their entries. Explicit displayed family; member types and common actual first-use exposed, not one atomic glyph. Family listing, not one atomic glyph. Author macro \texttt{\detokenize{\langle\ \rangle}} not carried to the members. \\
\phantomsection\label{entry:symbol-registry-06675}\(\ \langle i,j,\varphi\rangle\) & \texttt{\detokenize{\ \langle i,}\allowbreak\detokenize{j,}\allowbreak\detokenize{\varphi\rangle}} & tasks of a first-order branch symbol family member \(\ \langle i,j,\varphi\rangle\) & \hyperref[sec:p02-w02-p02-014-fol-calculus]{\ref*{sec:p02-w02-p02-014-fol-calculus}} & n/a & \hyperref[def:p02-w02-p02-014-fol-branch]{\ref*{def:p02-w02-p02-014-fol-branch}} & \thepage & Propositional or delta expansion at position \(i\); gamma instance at term \(t\); reflexivity at \(t\); replacement of an occurrence in position \(j\) using the equation at position \(i\). Angle brackets also denote finite lists in the same section; the task forms are distinguished by their entries. Explicit displayed family; member types and common actual first-use exposed, not one atomic glyph. Family listing, not one atomic glyph. Author macro \texttt{\detokenize{\langle\ \rangle}} not carried to the members. \\
\phantomsection\label{entry:symbol-registry-06676}\(B_\omega\) & \texttt{\detokenize{B_\omega}} & limit branch of a path through an infinite expansion sequence symbol family member \(B_\omega\) & \hyperref[sec:p02-w02-p02-014-fol-calculus]{\ref*{sec:p02-w02-p02-014-fol-calculus}} & n/a & \hyperref[def:p02-w02-p02-014-fol-branch]{\ref*{def:p02-w02-p02-014-fol-branch}} & \thepage & Union of the lists and of the done sets along the path. Explicit displayed family; member types and common actual first-use exposed, not one atomic glyph. Family listing, not one atomic glyph. Author macro \texttt{\detokenize{B\_\omega}} not carried to the members. \\
\phantomsection\label{entry:symbol-registry-06677}\(\ (B_j)_{j\in\N}\) & \texttt{\detokenize{\ (B_j)_{j\in\N}}} & limit branch of a path through an infinite expansion sequence symbol family member \(\ (B_j)_{j\in\N}\) & \hyperref[sec:p02-w02-p02-014-fol-calculus]{\ref*{sec:p02-w02-p02-014-fol-calculus}} & n/a & \hyperref[def:p02-w02-p02-014-fol-branch]{\ref*{def:p02-w02-p02-014-fol-branch}} & \thepage & Union of the lists and of the done sets along the path. Explicit displayed family; member types and common actual first-use exposed, not one atomic glyph. Family listing, not one atomic glyph. Author macro \texttt{\detokenize{B\_\omega}} not carried to the members. \\
\phantomsection\label{entry:symbol-registry-06678}\(H\) & \texttt{\detokenize{H}} & Hintikka set of first-order sentences & \hyperref[sec:p02-w02-p02-014-fol-herbrand]{\ref*{sec:p02-w02-p02-014-fol-herbrand}} & n/a & \hyperref[def:p02-w02-p02-014-hintikka]{\ref*{def:p02-w02-p02-014-hintikka}} & \thepage & Closure conditions (H0)--(H4); the formula set of an open saturated branch is one. \\
\phantomsection\label{entry:symbol-registry-06679}\(\sim_H\) & \texttt{\detokenize{\sim_H}} & listed-equality relation on closed terms and its classes symbol family member \(\sim_H\) & \hyperref[sec:p02-w02-p02-014-fol-herbrand]{\ref*{sec:p02-w02-p02-014-fol-herbrand}} & n/a & \hyperref[def:p02-w02-p02-014-hintikka]{\ref*{def:p02-w02-p02-014-hintikka}} & \thepage & \(r\sim_H u\) iff \(r\) and \(u\) are the same term or \(r=u\in H\); identity when \(H\) is equality-free. Original precedence descriptor 'rel' retained; no rank inferred. Explicit displayed family; member types and common actual first-use exposed, not one atomic glyph. Family listing, not one atomic glyph. Author macro \texttt{\detokenize{\sim\_H,}\allowbreak\detokenize{[t]}} not carried to the members. \\
\phantomsection\label{entry:symbol-registry-06680}\(\ [t]\) & \texttt{\detokenize{\ [t]}} & listed-equality relation on closed terms and its classes symbol family member \(\ [t]\) & \hyperref[sec:p02-w02-p02-014-fol-herbrand]{\ref*{sec:p02-w02-p02-014-fol-herbrand}} & n/a & \hyperref[def:p02-w02-p02-014-hintikka]{\ref*{def:p02-w02-p02-014-hintikka}} & \thepage & \(r\sim_H u\) iff \(r\) and \(u\) are the same term or \(r=u\in H\); identity when \(H\) is equality-free. Original precedence descriptor 'rel' retained; no rank inferred. Explicit displayed family; member types and common actual first-use exposed, not one atomic glyph. Family listing, not one atomic glyph. Author macro \texttt{\detokenize{\sim\_H,}\allowbreak\detokenize{[t]}} not carried to the members. \\
\phantomsection\label{entry:symbol-registry-06681}\(\mathcal M_H\) & \texttt{\detokenize{\mathcal M_H}} & quotient Herbrand structure of a Hintikka set & \hyperref[sec:p02-w02-p02-014-fol-herbrand]{\ref*{sec:p02-w02-p02-014-fol-herbrand}} & n/a & \hyperref[def:p02-w02-p02-014-hintikka]{\ref*{def:p02-w02-p02-014-hintikka}} & \thepage & Domain \(\mathrm{CT}_{\Gamma}/{\sim_H}\); coincides with \(\mathcal M_{B^{+}_{\mathrm{at}}}\) of thm:p02-fol-tableau-adequacy when \(H\) is equality-free. \\
\phantomsection\label{entry:symbol-registry-06682}\(t^{\mathcal M}\) & \texttt{\detokenize{t^{\mathcal M}}} & value of a closed term in a structure & \hyperref[sec:p02-w02-p02-014-fol-calculus]{\ref*{sec:p02-w02-p02-014-fol-calculus}} & n/a & \hyperref[sec:p02-w02-p02-014-fol-calculus]{\ref*{sec:p02-w02-p02-014-fol-calculus}} & \thepage & Abbreviates \(\sem{t}^{\mathcal M}_s\), which does not depend on \(s\) for closed \(t\). \\
\phantomsection\label{entry:symbol-registry-06683}\(T^{(n)}\) & \texttt{\detokenize{T^{(n)}}} & stage tableau of the fair systematic procedure and its stage-\(n\) task list symbol family member \(T^{(n)}\) & \hyperref[sec:p02-w02-p02-014-fol-systematic]{\ref*{sec:p02-w02-p02-014-fol-systematic}} & n/a & \hyperref[def:p02-w02-p02-014-fol-procedure]{\ref*{def:p02-w02-p02-014-fol-procedure}} & \thepage & \(m_B\) is the list length of \(B\) at the start of the stage. Explicit displayed family; member types and common actual first-use exposed, not one atomic glyph. Family listing, not one atomic glyph. Author macro \texttt{\detokenize{T\^{}{(n)},}\allowbreak\detokenize{W\_n(B),}\allowbreak\detokenize{m\_B}} not carried to the members. \\
\phantomsection\label{entry:symbol-registry-06684}\(\ W_n(B)\) & \texttt{\detokenize{\ W_n(B)}} & stage tableau of the fair systematic procedure and its stage-\(n\) task list symbol family member \(\ W_n(B)\) & \hyperref[sec:p02-w02-p02-014-fol-systematic]{\ref*{sec:p02-w02-p02-014-fol-systematic}} & n/a & \hyperref[def:p02-w02-p02-014-fol-procedure]{\ref*{def:p02-w02-p02-014-fol-procedure}} & \thepage & \(m_B\) is the list length of \(B\) at the start of the stage. Explicit displayed family; member types and common actual first-use exposed, not one atomic glyph. Family listing, not one atomic glyph. Author macro \texttt{\detokenize{T\^{}{(n)},}\allowbreak\detokenize{W\_n(B),}\allowbreak\detokenize{m\_B}} not carried to the members. \\
\phantomsection\label{entry:symbol-registry-06685}\(\ m_B\) & \texttt{\detokenize{\ m_B}} & stage tableau of the fair systematic procedure and its stage-\(n\) task list symbol family member \(\ m_B\) & \hyperref[sec:p02-w02-p02-014-fol-systematic]{\ref*{sec:p02-w02-p02-014-fol-systematic}} & n/a & \hyperref[def:p02-w02-p02-014-fol-procedure]{\ref*{def:p02-w02-p02-014-fol-procedure}} & \thepage & \(m_B\) is the list length of \(B\) at the start of the stage. Explicit displayed family; member types and common actual first-use exposed, not one atomic glyph. Family listing, not one atomic glyph. Author macro \texttt{\detokenize{T\^{}{(n)},}\allowbreak\detokenize{W\_n(B),}\allowbreak\detokenize{m\_B}} not carried to the members. \\
\phantomsection\label{entry:symbol-registry-06686}\(\rho(\chi)\) & \texttt{\detokenize{\rho}} & induction rank for the Herbrand truth lemma & \hyperref[sec:annex-p02-w02-p02-014-fol]{\ref*{sec:annex-p02-w02-p02-014-fol}} & meta & \hyperref[annex:p02-w02-p02-014-herbrand-truth]{\ref*{annex:p02-w02-p02-014-herbrand-truth}} & \thepage & \(2c(\chi)\) for a negation and \(2c(\chi)+1\) otherwise, with \(c\) first-order complexity. \\
\phantomsection\label{entry:symbol-registry-06687}\(O_n\) & \texttt{\detokenize{O_n}} & set of open branches of the stage-\(n\) tableau & \hyperref[sec:annex-p02-w02-p02-014-fol]{\ref*{sec:annex-p02-w02-p02-014-fol}} & n/a & \hyperref[annex:p02-w02-p02-014-fol-semidecision]{\ref*{annex:p02-w02-p02-014-fol-semidecision}} & \thepage & Used in the infinite branch construction. \\
\phantomsection\label{entry:symbol-registry-06688}\(\theta_n\) & \texttt{\detokenize{\theta_n}} & disjunction of n conjunctions of distinct variable pairs & \hyperref[sec:p02-w02-p02-007-limitations]{\ref*{sec:p02-w02-p02-007-limitations}} & n/a & \hyperref[thm:p02-w02-p02-007-cnf-size]{\ref*{thm:p02-w02-p02-007-cnf-size}} & \thepage & Right-associated finite disjunction of P\_i and Q\_i conjunctions, i from 1 to n; local to the CNF size theorem. Distinct from the P15 uses of theta\_n, whose scopes do not overlap. \\
\phantomsection\label{entry:symbol-registry-06689}\(S_n\) & \texttt{\detokenize{S_n}} & choice functions, their clauses, the complementary choice and the falsifying valuation symbol family member \(S_n\) & \hyperref[sec:p02-w02-p02-007-limitations]{\ref*{sec:p02-w02-p02-007-limitations}} & n/a & \hyperref[thm:p02-w02-p02-007-cnf-size]{\ref*{thm:p02-w02-p02-007-cnf-size}} & \thepage & S\_n: functions s with s(i) in {P\_i,Q\_i}; K\_s = {s(1),...,s(n)} (a clause of positive literals); bar s(i) the other member of {P\_i,Q\_i}; v\_s the valuation false on s(i), true on bar s(i), false elsewhere. Proof-local notation. Explicit displayed family; member types and common actual first-use exposed, not one atomic glyph. Family listing, not one atomic glyph. Author macro \texttt{\detokenize{S\_n,}\allowbreak\detokenize{K\_s,}\allowbreak\detokenize{\bar s,}\allowbreak\detokenize{v\_s}} not carried to the members. \\
\phantomsection\label{entry:symbol-registry-06690}\(K_s\) & \texttt{\detokenize{K_s}} & choice functions, their clauses, the complementary choice and the falsifying valuation symbol family member \(K_s\) & \hyperref[sec:p02-w02-p02-007-limitations]{\ref*{sec:p02-w02-p02-007-limitations}} & n/a & \hyperref[thm:p02-w02-p02-007-cnf-size]{\ref*{thm:p02-w02-p02-007-cnf-size}} & \thepage & S\_n: functions s with s(i) in {P\_i,Q\_i}; K\_s = {s(1),...,s(n)} (a clause of positive literals); bar s(i) the other member of {P\_i,Q\_i}; v\_s the valuation false on s(i), true on bar s(i), false elsewhere. Proof-local notation. Explicit displayed family; member types and common actual first-use exposed, not one atomic glyph. Family listing, not one atomic glyph. Author macro \texttt{\detokenize{S\_n,}\allowbreak\detokenize{K\_s,}\allowbreak\detokenize{\bar s,}\allowbreak\detokenize{v\_s}} not carried to the members. \\
\phantomsection\label{entry:symbol-registry-06691}\(\bar s\) & \texttt{\detokenize{\bar s}} & choice functions, their clauses, the complementary choice and the falsifying valuation symbol family member \(\bar s\) & \hyperref[sec:p02-w02-p02-007-limitations]{\ref*{sec:p02-w02-p02-007-limitations}} & n/a & \hyperref[thm:p02-w02-p02-007-cnf-size]{\ref*{thm:p02-w02-p02-007-cnf-size}} & \thepage & S\_n: functions s with s(i) in {P\_i,Q\_i}; K\_s = {s(1),...,s(n)} (a clause of positive literals); bar s(i) the other member of {P\_i,Q\_i}; v\_s the valuation false on s(i), true on bar s(i), false elsewhere. Proof-local notation. Explicit displayed family; member types and common actual first-use exposed, not one atomic glyph. Family listing, not one atomic glyph. Author macro \texttt{\detokenize{S\_n,}\allowbreak\detokenize{K\_s,}\allowbreak\detokenize{\bar s,}\allowbreak\detokenize{v\_s}} not carried to the members. \\
\phantomsection\label{entry:symbol-registry-06692}\(v_s\) & \texttt{\detokenize{v_s}} & choice functions, their clauses, the complementary choice and the falsifying valuation symbol family member \(v_s\) & \hyperref[sec:p02-w02-p02-007-limitations]{\ref*{sec:p02-w02-p02-007-limitations}} & n/a & \hyperref[thm:p02-w02-p02-007-cnf-size]{\ref*{thm:p02-w02-p02-007-cnf-size}} & \thepage & S\_n: functions s with s(i) in {P\_i,Q\_i}; K\_s = {s(1),...,s(n)} (a clause of positive literals); bar s(i) the other member of {P\_i,Q\_i}; v\_s the valuation false on s(i), true on bar s(i), false elsewhere. Proof-local notation. Explicit displayed family; member types and common actual first-use exposed, not one atomic glyph. Family listing, not one atomic glyph. Author macro \texttt{\detokenize{S\_n,}\allowbreak\detokenize{K\_s,}\allowbreak\detokenize{\bar s,}\allowbreak\detokenize{v\_s}} not carried to the members. \\
\phantomsection\label{entry:symbol-registry-06693}\(A\sim B\) & \texttt{\detokenize{\sim}} & Kleene's object-language equivalence abbreviation & \hyperref[sec:p02-w02-p02-007-relationships]{\ref*{sec:p02-w02-p02-007-relationships}} & n/a & \hyperref[sec:p02-w02-p02-007-relationships]{\ref*{sec:p02-w02-p02-007-relationships}} & \thepage & Kleene 1952 section 26 p. 113: A ~ B abbreviates (A > B) \& (B > A) in his formal system; reported notation only, not used to state handbook results. \\
\phantomsection\label{entry:symbol-registry-06694}\(A\supset B\) & \texttt{\detokenize{\supset}} & Kleene's implication connective & \hyperref[sec:p02-w02-p02-007-relationships]{\ref*{sec:p02-w02-p02-007-relationships}} & n/a & \hyperref[sec:p02-w02-p02-007-relationships]{\ref*{sec:p02-w02-p02-007-relationships}} & \thepage & Kleene's notation for the conditional; corresponds to the handbook arrow. Reported notation only. \\
\phantomsection\label{entry:symbol-registry-06695}\(A\mathbin{\&}B\) & \texttt{\detokenize{\mathbin{\&}}} & Kleene's conjunction connective & \hyperref[sec:p02-w02-p02-007-relationships]{\ref*{sec:p02-w02-p02-007-relationships}} & n/a & \hyperref[sec:p02-w02-p02-007-relationships]{\ref*{sec:p02-w02-p02-007-relationships}} & \thepage & Kleene's notation for conjunction; corresponds to the handbook wedge. Reported notation only; distinct from the P8 additive and t-norm conjunction rows, whose scopes do not overlap. \\
\phantomsection\label{entry:symbol-registry-06696}\(\mathcal O_n\) & \texttt{\detokenize{\mathcal O_n}} & sets of n-ary and of all positive-arity Boolean operations symbol family member \(\mathcal O_n\) & \hyperref[sec:p02-w02-p02-008-post]{\ref*{sec:p02-w02-p02-008-post}} & n/a & \hyperref[def:p02-w02-p02-008-operations]{\ref*{def:p02-w02-p02-008-operations}} & \thepage & O is the union over n at least 1; no nullary operations. Same objects as the O\_n, O of the pending P1 file with label tag p01-w02-p01-024. Explicit displayed family; member types and common actual first-use exposed, not one atomic glyph. Family listing, not one atomic glyph. Author macro \texttt{\detokenize{\mathcal O\_n,}\allowbreak\detokenize{\mathcal O}} not carried to the members. \\
\phantomsection\label{entry:symbol-registry-06697}\(\mathcal O\) & \texttt{\detokenize{\mathcal O}} & sets of n-ary and of all positive-arity Boolean operations symbol family member \(\mathcal O\) & \hyperref[sec:p02-w02-p02-008-post]{\ref*{sec:p02-w02-p02-008-post}} & n/a & \hyperref[def:p02-w02-p02-008-operations]{\ref*{def:p02-w02-p02-008-operations}} & \thepage & O is the union over n at least 1; no nullary operations. Same objects as the O\_n, O of the pending P1 file with label tag p01-w02-p01-024. Explicit displayed family; member types and common actual first-use exposed, not one atomic glyph. Family listing, not one atomic glyph. Author macro \texttt{\detokenize{\mathcal O\_n,}\allowbreak\detokenize{\mathcal O}} not carried to the members. \\
\phantomsection\label{entry:symbol-registry-06698}\(e^n_i\) & \texttt{\detokenize{e^n_i}} & n-ary projection onto coordinate i & \hyperref[sec:p02-w02-p02-008-post]{\ref*{sec:p02-w02-p02-008-post}} & n/a & \hyperref[def:p02-w02-p02-008-operations]{\ref*{def:p02-w02-p02-008-operations}} & \thepage & Every clone contains all projections. \\
\phantomsection\label{entry:symbol-registry-06699}\(\kappa^n_b\) & \texttt{\detokenize{\kappa^n_b}} & n-ary constant Boolean operation with value b & \hyperref[sec:p02-w02-p02-008-post]{\ref*{sec:p02-w02-p02-008-post}} & n/a & \hyperref[def:p02-w02-p02-008-operations]{\ref*{def:p02-w02-p02-008-operations}} & \thepage & Unary kappa\^{}1\_0, kappa\^{}1\_1 represent nullary falsum and verum in a clone (def:p02-w02-p02-008-arity). The P15 row using kappa\_0, kappa\_1 has a disjoint scope. \\
\phantomsection\label{entry:symbol-registry-06700}\(f\circ(g_1,\ldots,g_m)\) & \texttt{\detokenize{\circ}} & composition of Boolean operations & \hyperref[sec:p02-w02-p02-008-post]{\ref*{sec:p02-w02-p02-008-post}} & meta & \hyperref[def:p02-w02-p02-008-operations]{\ref*{def:p02-w02-p02-008-operations}} & \thepage & f m-ary, all g\_j n-ary; result n-ary. \\
\phantomsection\label{entry:symbol-registry-06701}\(\langle C\rangle\) & \texttt{\detokenize{\langle C\rangle}} & clone generated by a set of Boolean operations & \hyperref[sec:p02-w02-p02-008-post]{\ref*{sec:p02-w02-p02-008-post}} & meta & \hyperref[def:p02-w02-p02-008-operations]{\ref*{def:p02-w02-p02-008-operations}} & \thepage & Same role as the existing symbols.csv row scoped 'P2 prop-syntax-semantics'; scope widening requested rather than a new role. \\
\phantomsection\label{entry:symbol-registry-06702}\(\mathsf T_0\) & \texttt{\detokenize{\mathsf T_0}} & Post's five maximal clones symbol family member \(\mathsf T_0\) & \hyperref[sec:p02-w02-p02-008-post]{\ref*{sec:p02-w02-p02-008-post}} & n/a & \hyperref[def:p02-w02-p02-008-five]{\ref*{def:p02-w02-p02-008-five}} & \thepage & 0-preserving, 1-preserving, self-dual, monotone, affine. Same role as the existing symbols.csv row scoped 'P2 prop-syntax-semantics'. Disjoint from the P2 syllogistic S, the P7 combinator S and the Lambek and log-space L rows by scope. Explicit displayed family; member types and common actual first-use exposed, not one atomic glyph. Family listing, not one atomic glyph. Author macro \texttt{\detokenize{\mathsf T\_0,}\allowbreak\detokenize{\mathsf T\_1,}\allowbreak\detokenize{\mathsf S,}\allowbreak\detokenize{\mathsf M,}\allowbreak\detokenize{\mathsf L}} not carried to the members. \\
\phantomsection\label{entry:symbol-registry-06703}\(\mathsf T_1\) & \texttt{\detokenize{\mathsf T_1}} & Post's five maximal clones symbol family member \(\mathsf T_1\) & \hyperref[sec:p02-w02-p02-008-post]{\ref*{sec:p02-w02-p02-008-post}} & n/a & \hyperref[def:p02-w02-p02-008-five]{\ref*{def:p02-w02-p02-008-five}} & \thepage & 0-preserving, 1-preserving, self-dual, monotone, affine. Same role as the existing symbols.csv row scoped 'P2 prop-syntax-semantics'. Disjoint from the P2 syllogistic S, the P7 combinator S and the Lambek and log-space L rows by scope. Explicit displayed family; member types and common actual first-use exposed, not one atomic glyph. Family listing, not one atomic glyph. Author macro \texttt{\detokenize{\mathsf T\_0,}\allowbreak\detokenize{\mathsf T\_1,}\allowbreak\detokenize{\mathsf S,}\allowbreak\detokenize{\mathsf M,}\allowbreak\detokenize{\mathsf L}} not carried to the members. \\
\phantomsection\label{entry:symbol-registry-06704}\(\mathsf S\) & \texttt{\detokenize{\mathsf S}} & Post's five maximal clones symbol family member \(\mathsf S\) & \hyperref[sec:p02-w02-p02-008-post]{\ref*{sec:p02-w02-p02-008-post}} & n/a & \hyperref[def:p02-w02-p02-008-five]{\ref*{def:p02-w02-p02-008-five}} & \thepage & 0-preserving, 1-preserving, self-dual, monotone, affine. Same role as the existing symbols.csv row scoped 'P2 prop-syntax-semantics'. Disjoint from the P2 syllogistic S, the P7 combinator S and the Lambek and log-space L rows by scope. Explicit displayed family; member types and common actual first-use exposed, not one atomic glyph. Family listing, not one atomic glyph. Author macro \texttt{\detokenize{\mathsf T\_0,}\allowbreak\detokenize{\mathsf T\_1,}\allowbreak\detokenize{\mathsf S,}\allowbreak\detokenize{\mathsf M,}\allowbreak\detokenize{\mathsf L}} not carried to the members. \\
\phantomsection\label{entry:symbol-registry-06705}\(\mathsf M\) & \texttt{\detokenize{\mathsf M}} & Post's five maximal clones symbol family member \(\mathsf M\) & \hyperref[sec:p02-w02-p02-008-post]{\ref*{sec:p02-w02-p02-008-post}} & n/a & \hyperref[def:p02-w02-p02-008-five]{\ref*{def:p02-w02-p02-008-five}} & \thepage & 0-preserving, 1-preserving, self-dual, monotone, affine. Same role as the existing symbols.csv row scoped 'P2 prop-syntax-semantics'. Disjoint from the P2 syllogistic S, the P7 combinator S and the Lambek and log-space L rows by scope. Explicit displayed family; member types and common actual first-use exposed, not one atomic glyph. Family listing, not one atomic glyph. Author macro \texttt{\detokenize{\mathsf T\_0,}\allowbreak\detokenize{\mathsf T\_1,}\allowbreak\detokenize{\mathsf S,}\allowbreak\detokenize{\mathsf M,}\allowbreak\detokenize{\mathsf L}} not carried to the members. \\
\phantomsection\label{entry:symbol-registry-06706}\(\mathsf L\) & \texttt{\detokenize{\mathsf L}} & Post's five maximal clones symbol family member \(\mathsf L\) & \hyperref[sec:p02-w02-p02-008-post]{\ref*{sec:p02-w02-p02-008-post}} & n/a & \hyperref[def:p02-w02-p02-008-five]{\ref*{def:p02-w02-p02-008-five}} & \thepage & 0-preserving, 1-preserving, self-dual, monotone, affine. Same role as the existing symbols.csv row scoped 'P2 prop-syntax-semantics'. Disjoint from the P2 syllogistic S, the P7 combinator S and the Lambek and log-space L rows by scope. Explicit displayed family; member types and common actual first-use exposed, not one atomic glyph. Family listing, not one atomic glyph. Author macro \texttt{\detokenize{\mathsf T\_0,}\allowbreak\detokenize{\mathsf T\_1,}\allowbreak\detokenize{\mathsf S,}\allowbreak\detokenize{\mathsf M,}\allowbreak\detokenize{\mathsf L}} not carried to the members. \\
\phantomsection\label{entry:symbol-registry-06707}\(C^+\) & \texttt{\detokenize{C^+}} & operation set of a connective set & \hyperref[sec:p02-w02-p02-008-post]{\ref*{sec:p02-w02-p02-008-post}} & n/a & \hyperref[def:p02-w02-p02-008-arity]{\ref*{def:p02-w02-p02-008-arity}} & \thepage & Positive-arity interpreting operations plus kappa\^{}1\_b for each nullary connective of value b. \\
\phantomsection\label{entry:symbol-registry-06708}\(x\oplus y\) & \texttt{\detokenize{\oplus}} & addition modulo 2 on the truth values 0 and 1 & \hyperref[sec:p02-w02-p02-008-post]{\ref*{sec:p02-w02-p02-008-post}} & n/a & \hyperref[def:p02-w02-p02-008-bits]{\ref*{def:p02-w02-p02-008-bits}} & \thepage & Also the binary Boolean operation of exclusive disjunction; P4, P8, P12, P16, P19 rows for oplus have disjoint scopes. \\
\phantomsection\label{entry:symbol-registry-06709}\(\bar a^{\,c}\) & \texttt{\detokenize{\bar a^{\,c}}} & complement of a bit tuple, coordinate update, constant tuples symbol family member \(\bar a^{\,c}\) & \hyperref[sec:p02-w02-p02-008-post]{\ref*{sec:p02-w02-p02-008-post}} & n/a & \hyperref[def:p02-w02-p02-008-bits]{\ref*{def:p02-w02-p02-008-bits}} & \thepage & Coordinatewise order on bit tuples is written with the ordinary less-or-equal sign. Explicit displayed family; member types and common actual first-use exposed, not one atomic glyph. Family listing, not one atomic glyph. Author macro \texttt{\detokenize{\bar a\^{}{\,c},}\allowbreak\detokenize{\bar a[i{:=}b],}\allowbreak\detokenize{\bar0,}\allowbreak\detokenize{\bar1}} not carried to the members. \\
\phantomsection\label{entry:symbol-registry-06710}\(\ \bar a[i{:=}b]\) & \texttt{\detokenize{\ \bar a[i{:=}b]}} & complement of a bit tuple, coordinate update, constant tuples symbol family member \(\ \bar a[i{:=}b]\) & \hyperref[sec:p02-w02-p02-008-post]{\ref*{sec:p02-w02-p02-008-post}} & n/a & \hyperref[def:p02-w02-p02-008-bits]{\ref*{def:p02-w02-p02-008-bits}} & \thepage & Coordinatewise order on bit tuples is written with the ordinary less-or-equal sign. Explicit displayed family; member types and common actual first-use exposed, not one atomic glyph. Family listing, not one atomic glyph. Author macro \texttt{\detokenize{\bar a\^{}{\,c},}\allowbreak\detokenize{\bar a[i{:=}b],}\allowbreak\detokenize{\bar0,}\allowbreak\detokenize{\bar1}} not carried to the members. \\
\phantomsection\label{entry:symbol-registry-06711}\(\ \bar0\) & \texttt{\detokenize{\ \bar0}} & complement of a bit tuple, coordinate update, constant tuples symbol family member \(\ \bar0\) & \hyperref[sec:p02-w02-p02-008-post]{\ref*{sec:p02-w02-p02-008-post}} & n/a & \hyperref[def:p02-w02-p02-008-bits]{\ref*{def:p02-w02-p02-008-bits}} & \thepage & Coordinatewise order on bit tuples is written with the ordinary less-or-equal sign. Explicit displayed family; member types and common actual first-use exposed, not one atomic glyph. Family listing, not one atomic glyph. Author macro \texttt{\detokenize{\bar a\^{}{\,c},}\allowbreak\detokenize{\bar a[i{:=}b],}\allowbreak\detokenize{\bar0,}\allowbreak\detokenize{\bar1}} not carried to the members. \\
\phantomsection\label{entry:symbol-registry-06712}\(\ \bar1\) & \texttt{\detokenize{\ \bar1}} & complement of a bit tuple, coordinate update, constant tuples symbol family member \(\ \bar1\) & \hyperref[sec:p02-w02-p02-008-post]{\ref*{sec:p02-w02-p02-008-post}} & n/a & \hyperref[def:p02-w02-p02-008-bits]{\ref*{def:p02-w02-p02-008-bits}} & \thepage & Coordinatewise order on bit tuples is written with the ordinary less-or-equal sign. Explicit displayed family; member types and common actual first-use exposed, not one atomic glyph. Family listing, not one atomic glyph. Author macro \texttt{\detokenize{\bar a\^{}{\,c},}\allowbreak\detokenize{\bar a[i{:=}b],}\allowbreak\detokenize{\bar0,}\allowbreak\detokenize{\bar1}} not carried to the members. \\
\phantomsection\label{entry:symbol-registry-06713}\(\operatorname{maj}\) & \texttt{\detokenize{\operatorname{maj}}} & ternary majority operation & \hyperref[sec:p02-w02-p02-008-post]{\ref*{sec:p02-w02-p02-008-post}} & n/a & \hyperref[def:p02-w02-p02-008-operations]{\ref*{def:p02-w02-p02-008-operations}} & \thepage & Value 1 iff at least two arguments are 1. \\
\phantomsection\label{entry:symbol-registry-06714}\(\mathrm{NAND}\) & \texttt{\detokenize{\mathrm{NAND}}} & binary Boolean operations and the connective symbols interpreted by them symbol family member \(\mathrm{NAND}\) & \hyperref[sec:p02-w02-p02-008-post]{\ref*{sec:p02-w02-p02-008-post}} & n/a & \hyperref[def:p02-w02-p02-008-operations]{\ref*{def:p02-w02-p02-008-operations}} & \thepage & NAND(x,y)=1+xy, NOR(x,y)=(1+x)(1+y) mod 2. The existing symbols.csv row (scope 'P2 prop-syntax-semantics') records them as formula abbreviations; here they are connectives of the connective sets studied. Explicit displayed family; member types and common actual first-use exposed, not one atomic glyph. Family listing, not one atomic glyph. Author macro \texttt{\detokenize{\mathrm{NAND},}\allowbreak\detokenize{\mathrm{NOR}}} not carried to the members. \\
\phantomsection\label{entry:symbol-registry-06715}\(\mathrm{NOR}\) & \texttt{\detokenize{\mathrm{NOR}}} & binary Boolean operations and the connective symbols interpreted by them symbol family member \(\mathrm{NOR}\) & \hyperref[sec:p02-w02-p02-008-post]{\ref*{sec:p02-w02-p02-008-post}} & n/a & \hyperref[def:p02-w02-p02-008-operations]{\ref*{def:p02-w02-p02-008-operations}} & \thepage & NAND(x,y)=1+xy, NOR(x,y)=(1+x)(1+y) mod 2. The existing symbols.csv row (scope 'P2 prop-syntax-semantics') records them as formula abbreviations; here they are connectives of the connective sets studied. Explicit displayed family; member types and common actual first-use exposed, not one atomic glyph. Family listing, not one atomic glyph. Author macro \texttt{\detokenize{\mathrm{NAND},}\allowbreak\detokenize{\mathrm{NOR}}} not carried to the members. \\
\phantomsection\label{entry:symbol-registry-06716}\(f^\Delta\) & \texttt{\detokenize{f^\Delta}} & diagonal of a Boolean operation & \hyperref[sec:annex-p02-w02-p02-008-post]{\ref*{sec:annex-p02-w02-p02-008-post}} & n/a & \hyperref[annex:p02-w02-p02-008-post]{\ref*{annex:p02-w02-p02-008-post}} & \thepage & f(x,...,x); proof-local. \\
\phantomsection\label{entry:symbol-registry-06717}\(\mathcal J\) & \texttt{\detokenize{\mathcal J}} & monomial set of an algebraic normal form and its sum symbol family member \(\mathcal J\) & \hyperref[sec:p02-w02-p02-008-post]{\ref*{sec:p02-w02-p02-008-post}} & n/a & \hyperref[thm:p02-w02-p02-008-anf]{\ref*{thm:p02-w02-p02-008-anf}} & \thepage & Sum mod 2 over J in the set of the product of a\_i over i in J. Explicit displayed family; member types and common actual first-use exposed, not one atomic glyph. Family listing, not one atomic glyph. Author macro \texttt{\detokenize{\mathcal J,}\allowbreak\detokenize{\Sigma\_{\mathcal J}}} not carried to the members. \\
\phantomsection\label{entry:symbol-registry-06718}\(\ \Sigma_{\mathcal J}\) & \texttt{\detokenize{\ \Sigma_{\mathcal J}}} & monomial set of an algebraic normal form and its sum symbol family member \(\ \Sigma_{\mathcal J}\) & \hyperref[sec:p02-w02-p02-008-post]{\ref*{sec:p02-w02-p02-008-post}} & n/a & \hyperref[thm:p02-w02-p02-008-anf]{\ref*{thm:p02-w02-p02-008-anf}} & \thepage & Sum mod 2 over J in the set of the product of a\_i over i in J. Explicit displayed family; member types and common actual first-use exposed, not one atomic glyph. Family listing, not one atomic glyph. Author macro \texttt{\detokenize{\mathcal J,}\allowbreak\detokenize{\Sigma\_{\mathcal J}}} not carried to the members. \\
\phantomsection\label{entry:symbol-registry-06719}\(\theta\) & \texttt{\detokenize{\theta}} & local names of three formulas over one ternary connective in a worked example symbol family member \(\theta\) & \hyperref[def:p02-w02-p02-009-worked]{\ref*{def:p02-w02-p02-009-worked}} & n/a & \hyperref[def:p02-w02-p02-009-worked]{\ref*{def:p02-w02-p02-009-worked}} & \thepage & Metavariables for formulas, local to the example only; theta replaces kappa there to avoid a clash with the constant operations kappa\^{}n\_b of sec:p02-w02-p02-008-post. Explicit displayed family; member types and common actual first-use exposed, not one atomic glyph. Family listing, not one atomic glyph. Author macro \texttt{\detokenize{\theta,}\allowbreak\detokenize{\iota}} not carried to the members. \\
\phantomsection\label{entry:symbol-registry-06720}\(\iota\) & \texttt{\detokenize{\iota}} & local names of three formulas over one ternary connective in a worked example symbol family member \(\iota\) & \hyperref[def:p02-w02-p02-009-worked]{\ref*{def:p02-w02-p02-009-worked}} & n/a & \hyperref[def:p02-w02-p02-009-worked]{\ref*{def:p02-w02-p02-009-worked}} & \thepage & Metavariables for formulas, local to the example only; theta replaces kappa there to avoid a clash with the constant operations kappa\^{}n\_b of sec:p02-w02-p02-008-post. Explicit displayed family; member types and common actual first-use exposed, not one atomic glyph. Family listing, not one atomic glyph. Author macro \texttt{\detokenize{\theta,}\allowbreak\detokenize{\iota}} not carried to the members. \\
\phantomsection\label{entry:symbol-registry-06721}\(\iota'\) & \texttt{\detokenize{\iota'}} & local names of three formulas over one ternary connective in a worked example symbol family member \(\iota'\) & \hyperref[def:p02-w02-p02-009-worked]{\ref*{def:p02-w02-p02-009-worked}} & n/a & \hyperref[def:p02-w02-p02-009-worked]{\ref*{def:p02-w02-p02-009-worked}} & \thepage & Metavariables for formulas, local to the example only; theta replaces kappa there to avoid a clash with the constant operations kappa\^{}n\_b of sec:p02-w02-p02-008-post. Explicit displayed family; member types and common actual first-use exposed, not one atomic glyph. Family listing, not one atomic glyph. Author macro \texttt{\detokenize{\theta,}\allowbreak\detokenize{\iota}} not carried to the members. \\
\phantomsection\label{entry:symbol-registry-06722}\(\upsilon_{\neg}\) & \texttt{\detokenize{\upsilon_{\neg}}} & complementing uniform substitution & \hyperref[sec:p02-w02-p02-010-duality-arbitrary]{\ref*{sec:p02-w02-p02-010-duality-arbitrary}} & n/a & \hyperref[def:p02-w02-p02-010-complement]{\ref*{def:p02-w02-p02-010-complement}} & \thepage & The uniform substitution with \(\upsilon_{\neg}(p_n)=\neg p_n\) for every n; an instance of the registered \(\upsilon\) of P2 prop-syntax-semantics, whose induced valuation is the complemented valuation \(\bar v\). \\
\phantomsection\label{entry:symbol-registry-06723}\(\mathsf{NNF}\) & \texttt{\detokenize{\mathsf{NNF}}} & set of NNF formulas & \hyperref[sec:annex-p02-w02-p02-010-duality-arbitrary]{\ref*{sec:annex-p02-w02-p02-010-duality-arbitrary}} & n/a & \hyperref[annex:p02-w02-p02-010-nnf-presentation]{\ref*{annex:p02-w02-p02-010-nnf-presentation}} & \thepage & Proof-local name for the least set of CPL formulas containing every literal, \(\bottom\) and \(\topv\) and closed under binary conjunction and disjunction; \(\mathsf{NNF}'\) is the proof-local subset used in the leastness argument. \\
\phantomsection\label{entry:symbol-registry-06724}\(q\mathbin{\mathsf a}p\) & \texttt{\detokenize{q\mathbin{\mathsf a}p}} & categorical patterns, predicate q first and subject p second, over distinct atoms symbol family member \(q\mathbin{\mathsf a}p\) & \hyperref[sec:p02-w02-p02-001-historical-contrast]{\ref*{sec:p02-w02-p02-001-historical-contrast}} & n/a & \hyperref[def:p02-w02-p02-001-patterns]{\ref*{def:p02-w02-p02-001-patterns}} & \thepage & Predicate-first order follows Boger's AaB; sans-serif constants distinguish them from atom names. Explicit displayed family; member types and common actual first-use exposed, not one atomic glyph. Family listing, not one atomic glyph. Author macro \texttt{\detokenize{q\mathbin{\mathsf a}p,}\allowbreak\detokenize{q\mathbin{\mathsf e}p,}\allowbreak\detokenize{q\mathbin{\mathsf i}p,}\allowbreak\detokenize{q\mathbin{\mathsf o}p}} not carried to the members. \\
\phantomsection\label{entry:symbol-registry-06725}\(q\mathbin{\mathsf e}p\) & \texttt{\detokenize{q\mathbin{\mathsf e}p}} & categorical patterns, predicate q first and subject p second, over distinct atoms symbol family member \(q\mathbin{\mathsf e}p\) & \hyperref[sec:p02-w02-p02-001-historical-contrast]{\ref*{sec:p02-w02-p02-001-historical-contrast}} & n/a & \hyperref[def:p02-w02-p02-001-patterns]{\ref*{def:p02-w02-p02-001-patterns}} & \thepage & Predicate-first order follows Boger's AaB; sans-serif constants distinguish them from atom names. Explicit displayed family; member types and common actual first-use exposed, not one atomic glyph. Family listing, not one atomic glyph. Author macro \texttt{\detokenize{q\mathbin{\mathsf a}p,}\allowbreak\detokenize{q\mathbin{\mathsf e}p,}\allowbreak\detokenize{q\mathbin{\mathsf i}p,}\allowbreak\detokenize{q\mathbin{\mathsf o}p}} not carried to the members. \\
\phantomsection\label{entry:symbol-registry-06726}\(q\mathbin{\mathsf i}p\) & \texttt{\detokenize{q\mathbin{\mathsf i}p}} & categorical patterns, predicate q first and subject p second, over distinct atoms symbol family member \(q\mathbin{\mathsf i}p\) & \hyperref[sec:p02-w02-p02-001-historical-contrast]{\ref*{sec:p02-w02-p02-001-historical-contrast}} & n/a & \hyperref[def:p02-w02-p02-001-patterns]{\ref*{def:p02-w02-p02-001-patterns}} & \thepage & Predicate-first order follows Boger's AaB; sans-serif constants distinguish them from atom names. Explicit displayed family; member types and common actual first-use exposed, not one atomic glyph. Family listing, not one atomic glyph. Author macro \texttt{\detokenize{q\mathbin{\mathsf a}p,}\allowbreak\detokenize{q\mathbin{\mathsf e}p,}\allowbreak\detokenize{q\mathbin{\mathsf i}p,}\allowbreak\detokenize{q\mathbin{\mathsf o}p}} not carried to the members. \\
\phantomsection\label{entry:symbol-registry-06727}\(q\mathbin{\mathsf o}p\) & \texttt{\detokenize{q\mathbin{\mathsf o}p}} & categorical patterns, predicate q first and subject p second, over distinct atoms symbol family member \(q\mathbin{\mathsf o}p\) & \hyperref[sec:p02-w02-p02-001-historical-contrast]{\ref*{sec:p02-w02-p02-001-historical-contrast}} & n/a & \hyperref[def:p02-w02-p02-001-patterns]{\ref*{def:p02-w02-p02-001-patterns}} & \thepage & Predicate-first order follows Boger's AaB; sans-serif constants distinguish them from atom names. Explicit displayed family; member types and common actual first-use exposed, not one atomic glyph. Family listing, not one atomic glyph. Author macro \texttt{\detokenize{q\mathbin{\mathsf a}p,}\allowbreak\detokenize{q\mathbin{\mathsf e}p,}\allowbreak\detokenize{q\mathbin{\mathsf i}p,}\allowbreak\detokenize{q\mathbin{\mathsf o}p}} not carried to the members. \\
\phantomsection\label{entry:symbol-registry-06728}\(\models_{\mathrm B}\) & \texttt{\detokenize{\models_{\mathrm B}}} & existence-committing comparison satisfaction and consequence for patterns & \hyperref[sec:p02-w02-p02-001-historical-contrast]{\ref*{sec:p02-w02-p02-001-historical-contrast}} & n/a & \hyperref[def:p02-w02-p02-001-patterns]{\ref*{def:p02-w02-p02-001-patterns}} & \thepage & Defined through alpha, iota, epsilon, omega of def:p13-boger-guarded-bridge; extensional comparison only. \\
\phantomsection\label{entry:symbol-registry-06729}\(\varphi^{\mathrm c}\) & \texttt{\detokenize{\varphi^{\mathrm c}}} & contradictory pattern (a/o and e/i exchanged, terms kept) & \hyperref[sec:p02-w02-p02-001-historical-contrast]{\ref*{sec:p02-w02-p02-001-historical-contrast}} & n/a & \hyperref[def:p02-w02-p02-001-patterns]{\ref*{def:p02-w02-p02-001-patterns}} & \thepage & Distinct from the formula complement bar on S-formulas; sigma maps one to the other. \\
\phantomsection\label{entry:symbol-registry-06730}\(\sigma\) & \texttt{\detokenize{\sigma}} & S-reading of a pattern as an S-formula & \hyperref[sec:p02-w02-p02-001-historical-contrast]{\ref*{sec:p02-w02-p02-001-historical-contrast}} & n/a & \hyperref[def:p02-w02-p02-001-patterns]{\ref*{def:p02-w02-p02-001-patterns}} & \thepage & Local map; unrelated to substitutions or signatures elsewhere. \\
\phantomsection\label{entry:symbol-registry-06731}\(\Gamma\) & \texttt{\detokenize{\Gamma}} & finite sets of patterns symbol family member \(\Gamma\) & \hyperref[sec:p02-w02-p02-001-historical-contrast]{\ref*{sec:p02-w02-p02-001-historical-contrast}} & n/a & \hyperref[def:p02-w02-p02-001-patterns]{\ref*{def:p02-w02-p02-001-patterns}} & \thepage & Delta is Gamma with the contradictory of the proposed conclusion added. Explicit displayed family; member types and common actual first-use exposed, not one atomic glyph. Family listing, not one atomic glyph. Author macro \texttt{\detokenize{\Gamma,}\allowbreak\detokenize{\Delta}} not carried to the members. \\
\phantomsection\label{entry:symbol-registry-06732}\(\Delta\) & \texttt{\detokenize{\Delta}} & finite sets of patterns symbol family member \(\Delta\) & \hyperref[sec:p02-w02-p02-001-historical-contrast]{\ref*{sec:p02-w02-p02-001-historical-contrast}} & n/a & \hyperref[def:p02-w02-p02-001-patterns]{\ref*{def:p02-w02-p02-001-patterns}} & \thepage & Delta is Gamma with the contradictory of the proposed conclusion added. Explicit displayed family; member types and common actual first-use exposed, not one atomic glyph. Family listing, not one atomic glyph. Author macro \texttt{\detokenize{\Gamma,}\allowbreak\detokenize{\Delta}} not carried to the members. \\
\phantomsection\label{entry:symbol-registry-06733}\(s\) & \texttt{\detokenize{s}} & selection for o-patterns & \hyperref[sec:p02-w02-p02-001-historical-contrast]{\ref*{sec:p02-w02-p02-001-historical-contrast}} & n/a & \hyperref[def:p02-w02-p02-001-selection]{\ref*{def:p02-w02-p02-001-selection}} & \thepage & s picks forall(p,p-bar) or exists(p,q-bar) for each o-pattern. \\
\phantomsection\label{entry:symbol-registry-06734}\(\tau_s\) & \texttt{\detokenize{\tau_s}} & translation of patterns into S-formula sets induced by a selection & \hyperref[sec:p02-w02-p02-001-historical-contrast]{\ref*{sec:p02-w02-p02-001-historical-contrast}} & n/a & \hyperref[def:p02-w02-p02-001-selection]{\ref*{def:p02-w02-p02-001-selection}} & \thepage & \(\tau_s(\Delta)\) is the union of the \(\tau_s(\psi)\) for \(\psi\in\Delta\). \\
\phantomsection\label{entry:symbol-registry-06735}\(\mathfrak M_1\) & \texttt{\detokenize{\mathfrak M_1}} & three named finite structures in one proof symbol family member \(\mathfrak M_1\) & \hyperref[annex:p02-w02-p02-001-o-inexpressible]{\ref*{annex:p02-w02-p02-001-o-inexpressible}} & n/a & \hyperref[annex:p02-w02-p02-001-o-inexpressible]{\ref*{annex:p02-w02-p02-001-o-inexpressible}} & \thepage & Proof-local. Explicit displayed family; member types and common actual first-use exposed, not one atomic glyph. Family listing, not one atomic glyph. Author macro \texttt{\detokenize{\mathfrak M\_1,}\allowbreak\detokenize{\mathfrak U,}\allowbreak\detokenize{\mathfrak M\_5}} not carried to the members. \\
\phantomsection\label{entry:symbol-registry-06736}\(\mathfrak U\) & \texttt{\detokenize{\mathfrak U}} & three named finite structures in one proof symbol family member \(\mathfrak U\) & \hyperref[annex:p02-w02-p02-001-o-inexpressible]{\ref*{annex:p02-w02-p02-001-o-inexpressible}} & n/a & \hyperref[annex:p02-w02-p02-001-o-inexpressible]{\ref*{annex:p02-w02-p02-001-o-inexpressible}} & \thepage & Proof-local. Explicit displayed family; member types and common actual first-use exposed, not one atomic glyph. Family listing, not one atomic glyph. Author macro \texttt{\detokenize{\mathfrak M\_1,}\allowbreak\detokenize{\mathfrak U,}\allowbreak\detokenize{\mathfrak M\_5}} not carried to the members. \\
\phantomsection\label{entry:symbol-registry-06737}\(\mathfrak M_5\) & \texttt{\detokenize{\mathfrak M_5}} & three named finite structures in one proof symbol family member \(\mathfrak M_5\) & \hyperref[annex:p02-w02-p02-001-o-inexpressible]{\ref*{annex:p02-w02-p02-001-o-inexpressible}} & n/a & \hyperref[annex:p02-w02-p02-001-o-inexpressible]{\ref*{annex:p02-w02-p02-001-o-inexpressible}} & \thepage & Proof-local. Explicit displayed family; member types and common actual first-use exposed, not one atomic glyph. Family listing, not one atomic glyph. Author macro \texttt{\detokenize{\mathfrak M\_1,}\allowbreak\detokenize{\mathfrak U,}\allowbreak\detokenize{\mathfrak M\_5}} not carried to the members. \\
\phantomsection\label{entry:symbol-registry-06738}\(XaZ\) & \texttt{\detokenize{XaZ}} & Boger's categorical sentences, predicate X first and subject Z second, as source notation symbol family member \(XaZ\) & \hyperref[sec:p02-w02-p02-002-fourteen]{\ref*{sec:p02-w02-p02-002-fourteen}} & n/a & \hyperref[def:p02-w02-p02-002-reading]{\ref*{def:p02-w02-p02-002-reading}} & \thepage & Source notation of Boger 2004 table 20; read as the patterns x a z etc. of def:p02-w02-p02-001-patterns. Explicit displayed family; member types and common actual first-use exposed, not one atomic glyph. Family listing, not one atomic glyph. Author macro \texttt{\detokenize{XaZ,}\allowbreak\detokenize{XeZ,}\allowbreak\detokenize{XiZ,}\allowbreak\detokenize{XoZ}} not carried to the members. \\
\phantomsection\label{entry:symbol-registry-06739}\(XeZ\) & \texttt{\detokenize{XeZ}} & Boger's categorical sentences, predicate X first and subject Z second, as source notation symbol family member \(XeZ\) & \hyperref[sec:p02-w02-p02-002-fourteen]{\ref*{sec:p02-w02-p02-002-fourteen}} & n/a & \hyperref[def:p02-w02-p02-002-reading]{\ref*{def:p02-w02-p02-002-reading}} & \thepage & Source notation of Boger 2004 table 20; read as the patterns x a z etc. of def:p02-w02-p02-001-patterns. Explicit displayed family; member types and common actual first-use exposed, not one atomic glyph. Family listing, not one atomic glyph. Author macro \texttt{\detokenize{XaZ,}\allowbreak\detokenize{XeZ,}\allowbreak\detokenize{XiZ,}\allowbreak\detokenize{XoZ}} not carried to the members. \\
\phantomsection\label{entry:symbol-registry-06740}\(XiZ\) & \texttt{\detokenize{XiZ}} & Boger's categorical sentences, predicate X first and subject Z second, as source notation symbol family member \(XiZ\) & \hyperref[sec:p02-w02-p02-002-fourteen]{\ref*{sec:p02-w02-p02-002-fourteen}} & n/a & \hyperref[def:p02-w02-p02-002-reading]{\ref*{def:p02-w02-p02-002-reading}} & \thepage & Source notation of Boger 2004 table 20; read as the patterns x a z etc. of def:p02-w02-p02-001-patterns. Explicit displayed family; member types and common actual first-use exposed, not one atomic glyph. Family listing, not one atomic glyph. Author macro \texttt{\detokenize{XaZ,}\allowbreak\detokenize{XeZ,}\allowbreak\detokenize{XiZ,}\allowbreak\detokenize{XoZ}} not carried to the members. \\
\phantomsection\label{entry:symbol-registry-06741}\(XoZ\) & \texttt{\detokenize{XoZ}} & Boger's categorical sentences, predicate X first and subject Z second, as source notation symbol family member \(XoZ\) & \hyperref[sec:p02-w02-p02-002-fourteen]{\ref*{sec:p02-w02-p02-002-fourteen}} & n/a & \hyperref[def:p02-w02-p02-002-reading]{\ref*{def:p02-w02-p02-002-reading}} & \thepage & Source notation of Boger 2004 table 20; read as the patterns x a z etc. of def:p02-w02-p02-001-patterns. Explicit displayed family; member types and common actual first-use exposed, not one atomic glyph. Family listing, not one atomic glyph. Author macro \texttt{\detokenize{XaZ,}\allowbreak\detokenize{XeZ,}\allowbreak\detokenize{XiZ,}\allowbreak\detokenize{XoZ}} not carried to the members. \\
\phantomsection\label{entry:symbol-registry-06742}\(a\) & \texttt{\detokenize{a}} & atoms naming Boger's term letters A,B,C,M,N,X,P,R,S symbol family member \(a\) & \hyperref[sec:p02-w02-p02-002-fourteen]{\ref*{sec:p02-w02-p02-002-fourteen}} & n/a & \hyperref[def:p02-w02-p02-002-reading]{\ref*{def:p02-w02-p02-002-reading}} & \thepage & Chapter-local; m and n are atoms here, not the literal metavariables of def:p02-syllogistic-language. Explicit displayed family; member types and common actual first-use exposed, not one atomic glyph. Family listing, not one atomic glyph. Author macro \texttt{\detokenize{a,}\allowbreak\detokenize{b,}\allowbreak\detokenize{c,}\allowbreak\detokenize{m,}\allowbreak\detokenize{n,}\allowbreak\detokenize{x,}\allowbreak\detokenize{p,}\allowbreak\detokenize{r,}\allowbreak\detokenize{s}} not carried to the members. \\
\phantomsection\label{entry:symbol-registry-06743}\(b\) & \texttt{\detokenize{b}} & atoms naming Boger's term letters A,B,C,M,N,X,P,R,S symbol family member \(b\) & \hyperref[sec:p02-w02-p02-002-fourteen]{\ref*{sec:p02-w02-p02-002-fourteen}} & n/a & \hyperref[def:p02-w02-p02-002-reading]{\ref*{def:p02-w02-p02-002-reading}} & \thepage & Chapter-local; m and n are atoms here, not the literal metavariables of def:p02-syllogistic-language. Explicit displayed family; member types and common actual first-use exposed, not one atomic glyph. Family listing, not one atomic glyph. Author macro \texttt{\detokenize{a,}\allowbreak\detokenize{b,}\allowbreak\detokenize{c,}\allowbreak\detokenize{m,}\allowbreak\detokenize{n,}\allowbreak\detokenize{x,}\allowbreak\detokenize{p,}\allowbreak\detokenize{r,}\allowbreak\detokenize{s}} not carried to the members. \\
\phantomsection\label{entry:symbol-registry-06744}\(c\) & \texttt{\detokenize{c}} & atoms naming Boger's term letters A,B,C,M,N,X,P,R,S symbol family member \(c\) & \hyperref[sec:p02-w02-p02-002-fourteen]{\ref*{sec:p02-w02-p02-002-fourteen}} & n/a & \hyperref[def:p02-w02-p02-002-reading]{\ref*{def:p02-w02-p02-002-reading}} & \thepage & Chapter-local; m and n are atoms here, not the literal metavariables of def:p02-syllogistic-language. Explicit displayed family; member types and common actual first-use exposed, not one atomic glyph. Family listing, not one atomic glyph. Author macro \texttt{\detokenize{a,}\allowbreak\detokenize{b,}\allowbreak\detokenize{c,}\allowbreak\detokenize{m,}\allowbreak\detokenize{n,}\allowbreak\detokenize{x,}\allowbreak\detokenize{p,}\allowbreak\detokenize{r,}\allowbreak\detokenize{s}} not carried to the members. \\
\phantomsection\label{entry:symbol-registry-06745}\(m\) & \texttt{\detokenize{m}} & atoms naming Boger's term letters A,B,C,M,N,X,P,R,S symbol family member \(m\) & \hyperref[sec:p02-w02-p02-002-fourteen]{\ref*{sec:p02-w02-p02-002-fourteen}} & n/a & \hyperref[def:p02-w02-p02-002-reading]{\ref*{def:p02-w02-p02-002-reading}} & \thepage & Chapter-local; m and n are atoms here, not the literal metavariables of def:p02-syllogistic-language. Explicit displayed family; member types and common actual first-use exposed, not one atomic glyph. Family listing, not one atomic glyph. Author macro \texttt{\detokenize{a,}\allowbreak\detokenize{b,}\allowbreak\detokenize{c,}\allowbreak\detokenize{m,}\allowbreak\detokenize{n,}\allowbreak\detokenize{x,}\allowbreak\detokenize{p,}\allowbreak\detokenize{r,}\allowbreak\detokenize{s}} not carried to the members. \\
\phantomsection\label{entry:symbol-registry-06746}\(n\) & \texttt{\detokenize{n}} & atoms naming Boger's term letters A,B,C,M,N,X,P,R,S symbol family member \(n\) & \hyperref[sec:p02-w02-p02-002-fourteen]{\ref*{sec:p02-w02-p02-002-fourteen}} & n/a & \hyperref[def:p02-w02-p02-002-reading]{\ref*{def:p02-w02-p02-002-reading}} & \thepage & Chapter-local; m and n are atoms here, not the literal metavariables of def:p02-syllogistic-language. Explicit displayed family; member types and common actual first-use exposed, not one atomic glyph. Family listing, not one atomic glyph. Author macro \texttt{\detokenize{a,}\allowbreak\detokenize{b,}\allowbreak\detokenize{c,}\allowbreak\detokenize{m,}\allowbreak\detokenize{n,}\allowbreak\detokenize{x,}\allowbreak\detokenize{p,}\allowbreak\detokenize{r,}\allowbreak\detokenize{s}} not carried to the members. \\
\phantomsection\label{entry:symbol-registry-06747}\(x\) & \texttt{\detokenize{x}} & atoms naming Boger's term letters A,B,C,M,N,X,P,R,S symbol family member \(x\) & \hyperref[sec:p02-w02-p02-002-fourteen]{\ref*{sec:p02-w02-p02-002-fourteen}} & n/a & \hyperref[def:p02-w02-p02-002-reading]{\ref*{def:p02-w02-p02-002-reading}} & \thepage & Chapter-local; m and n are atoms here, not the literal metavariables of def:p02-syllogistic-language. Explicit displayed family; member types and common actual first-use exposed, not one atomic glyph. Family listing, not one atomic glyph. Author macro \texttt{\detokenize{a,}\allowbreak\detokenize{b,}\allowbreak\detokenize{c,}\allowbreak\detokenize{m,}\allowbreak\detokenize{n,}\allowbreak\detokenize{x,}\allowbreak\detokenize{p,}\allowbreak\detokenize{r,}\allowbreak\detokenize{s}} not carried to the members. \\
\phantomsection\label{entry:symbol-registry-06748}\(p\) & \texttt{\detokenize{p}} & atoms naming Boger's term letters A,B,C,M,N,X,P,R,S symbol family member \(p\) & \hyperref[sec:p02-w02-p02-002-fourteen]{\ref*{sec:p02-w02-p02-002-fourteen}} & n/a & \hyperref[def:p02-w02-p02-002-reading]{\ref*{def:p02-w02-p02-002-reading}} & \thepage & Chapter-local; m and n are atoms here, not the literal metavariables of def:p02-syllogistic-language. Explicit displayed family; member types and common actual first-use exposed, not one atomic glyph. Family listing, not one atomic glyph. Author macro \texttt{\detokenize{a,}\allowbreak\detokenize{b,}\allowbreak\detokenize{c,}\allowbreak\detokenize{m,}\allowbreak\detokenize{n,}\allowbreak\detokenize{x,}\allowbreak\detokenize{p,}\allowbreak\detokenize{r,}\allowbreak\detokenize{s}} not carried to the members. \\
\phantomsection\label{entry:symbol-registry-06749}\(r\) & \texttt{\detokenize{r}} & atoms naming Boger's term letters A,B,C,M,N,X,P,R,S symbol family member \(r\) & \hyperref[sec:p02-w02-p02-002-fourteen]{\ref*{sec:p02-w02-p02-002-fourteen}} & n/a & \hyperref[def:p02-w02-p02-002-reading]{\ref*{def:p02-w02-p02-002-reading}} & \thepage & Chapter-local; m and n are atoms here, not the literal metavariables of def:p02-syllogistic-language. Explicit displayed family; member types and common actual first-use exposed, not one atomic glyph. Family listing, not one atomic glyph. Author macro \texttt{\detokenize{a,}\allowbreak\detokenize{b,}\allowbreak\detokenize{c,}\allowbreak\detokenize{m,}\allowbreak\detokenize{n,}\allowbreak\detokenize{x,}\allowbreak\detokenize{p,}\allowbreak\detokenize{r,}\allowbreak\detokenize{s}} not carried to the members. \\
\phantomsection\label{entry:symbol-registry-06750}\(s\) & \texttt{\detokenize{s}} & atoms naming Boger's term letters A,B,C,M,N,X,P,R,S symbol family member \(s\) & \hyperref[sec:p02-w02-p02-002-fourteen]{\ref*{sec:p02-w02-p02-002-fourteen}} & n/a & \hyperref[def:p02-w02-p02-002-reading]{\ref*{def:p02-w02-p02-002-reading}} & \thepage & Chapter-local; m and n are atoms here, not the literal metavariables of def:p02-syllogistic-language. Explicit displayed family; member types and common actual first-use exposed, not one atomic glyph. Family listing, not one atomic glyph. Author macro \texttt{\detokenize{a,}\allowbreak\detokenize{b,}\allowbreak\detokenize{c,}\allowbreak\detokenize{m,}\allowbreak\detokenize{n,}\allowbreak\detokenize{x,}\allowbreak\detokenize{p,}\allowbreak\detokenize{r,}\allowbreak\detokenize{s}} not carried to the members. \\
\phantomsection\label{entry:symbol-registry-06751}\(\Gamma_9\) & \texttt{\detokenize{\Gamma_9}} & S-readings of the premises and conclusions of Darapti and Felapton symbol family member \(\Gamma_9\) & \hyperref[thm:p02-w02-p02-002-third-figure]{\ref*{thm:p02-w02-p02-002-third-figure}} & n/a & \hyperref[thm:p02-w02-p02-002-third-figure]{\ref*{thm:p02-w02-p02-002-third-figure}} & \thepage & Result-local. Explicit displayed family; member types and common actual first-use exposed, not one atomic glyph. Family listing, not one atomic glyph. Author macro \texttt{\detokenize{\Gamma\_9,}\allowbreak\detokenize{\Gamma\_{10},}\allowbreak\detokenize{\theta\_9,}\allowbreak\detokenize{\theta\_{10}}} not carried to the members. \\
\phantomsection\label{entry:symbol-registry-06752}\(\Gamma_{10}\) & \texttt{\detokenize{\Gamma_{10}}} & S-readings of the premises and conclusions of Darapti and Felapton symbol family member \(\Gamma_{10}\) & \hyperref[thm:p02-w02-p02-002-third-figure]{\ref*{thm:p02-w02-p02-002-third-figure}} & n/a & \hyperref[thm:p02-w02-p02-002-third-figure]{\ref*{thm:p02-w02-p02-002-third-figure}} & \thepage & Result-local. Explicit displayed family; member types and common actual first-use exposed, not one atomic glyph. Family listing, not one atomic glyph. Author macro \texttt{\detokenize{\Gamma\_9,}\allowbreak\detokenize{\Gamma\_{10},}\allowbreak\detokenize{\theta\_9,}\allowbreak\detokenize{\theta\_{10}}} not carried to the members. \\
\phantomsection\label{entry:symbol-registry-06753}\(\theta_9\) & \texttt{\detokenize{\theta_9}} & S-readings of the premises and conclusions of Darapti and Felapton symbol family member \(\theta_9\) & \hyperref[thm:p02-w02-p02-002-third-figure]{\ref*{thm:p02-w02-p02-002-third-figure}} & n/a & \hyperref[thm:p02-w02-p02-002-third-figure]{\ref*{thm:p02-w02-p02-002-third-figure}} & \thepage & Result-local. Explicit displayed family; member types and common actual first-use exposed, not one atomic glyph. Family listing, not one atomic glyph. Author macro \texttt{\detokenize{\Gamma\_9,}\allowbreak\detokenize{\Gamma\_{10},}\allowbreak\detokenize{\theta\_9,}\allowbreak\detokenize{\theta\_{10}}} not carried to the members. \\
\phantomsection\label{entry:symbol-registry-06754}\(\theta_{10}\) & \texttt{\detokenize{\theta_{10}}} & S-readings of the premises and conclusions of Darapti and Felapton symbol family member \(\theta_{10}\) & \hyperref[thm:p02-w02-p02-002-third-figure]{\ref*{thm:p02-w02-p02-002-third-figure}} & n/a & \hyperref[thm:p02-w02-p02-002-third-figure]{\ref*{thm:p02-w02-p02-002-third-figure}} & \thepage & Result-local. Explicit displayed family; member types and common actual first-use exposed, not one atomic glyph. Family listing, not one atomic glyph. Author macro \texttt{\detokenize{\Gamma\_9,}\allowbreak\detokenize{\Gamma\_{10},}\allowbreak\detokenize{\theta\_9,}\allowbreak\detokenize{\theta\_{10}}} not carried to the members. \\
\phantomsection\label{entry:symbol-registry-06755}\(Z^\ast\) & \texttt{\detokenize{Z^\ast}} & extensions of the subject and predicate atoms of a sentence XyZ in a structure symbol family member \(Z^\ast\) & \hyperref[annex:p02-w02-p02-002-committing]{\ref*{annex:p02-w02-p02-002-committing}} & n/a & \hyperref[annex:p02-w02-p02-002-committing]{\ref*{annex:p02-w02-p02-002-committing}} & \thepage & Proof-local. Explicit displayed family; member types and common actual first-use exposed, not one atomic glyph. Family listing, not one atomic glyph. Author macro \texttt{\detokenize{Z\^{}\ast,}\allowbreak\detokenize{X\^{}\ast}} not carried to the members. \\
\phantomsection\label{entry:symbol-registry-06756}\(X^\ast\) & \texttt{\detokenize{X^\ast}} & extensions of the subject and predicate atoms of a sentence XyZ in a structure symbol family member \(X^\ast\) & \hyperref[annex:p02-w02-p02-002-committing]{\ref*{annex:p02-w02-p02-002-committing}} & n/a & \hyperref[annex:p02-w02-p02-002-committing]{\ref*{annex:p02-w02-p02-002-committing}} & \thepage & Proof-local. Explicit displayed family; member types and common actual first-use exposed, not one atomic glyph. Family listing, not one atomic glyph. Author macro \texttt{\detokenize{Z\^{}\ast,}\allowbreak\detokenize{X\^{}\ast}} not carried to the members. \\
\phantomsection\label{entry:symbol-registry-06757}\(\mathcal D=(D,H,V)\) & \texttt{\detokenize{\mathcal D}} & domino system (tile set and horizontal and vertical compatibility relations) & \hyperref[sec:p02-w02-p02-035-decision-procedures]{\ref*{sec:p02-w02-p02-035-decision-procedures}} & n/a & \hyperref[def:p02-w02-p02-035-domino]{\ref*{def:p02-w02-p02-035-domino}} & \thepage & Source notation of Boerger-Graedel-Gurevich 1997, Def. 6.1.1. H and V here are relations on tiles, not formulas; the formulas H(z,z') and V(z,z') of the monadic tiling sentence are named after them in the source. \\
\phantomsection\label{entry:symbol-registry-06758}\(U(s,t)\) & \texttt{\detokenize{U(s,}\allowbreak\detokenize{t)}} & torus \(\mathbb Z_s\times\mathbb Z_t\) & \hyperref[sec:p02-w02-p02-035-decision-procedures]{\ref*{sec:p02-w02-p02-035-decision-procedures}} & n/a & \hyperref[def:p02-w02-p02-035-domino]{\ref*{def:p02-w02-p02-035-domino}} & \thepage & Source notation, Defs. 6.1.1 and 6.1.5. \\
\phantomsection\label{entry:symbol-registry-06759}\(Z(t)\) & \texttt{\detokenize{Z(t)}} & square torus \(U(t,t)\) & \hyperref[sec:p02-w02-p02-035-decision-procedures]{\ref*{sec:p02-w02-p02-035-decision-procedures}} & n/a & \hyperref[def:p02-w02-p02-035-domino]{\ref*{def:p02-w02-p02-035-domino}} & \thepage & Source notation, Defs. 6.1.1 and 6.1.5. \\
\phantomsection\label{entry:symbol-registry-06760}\(\mathrm{DOMINO}(\mathcal D,2^n)\) & \texttt{\detokenize{\mathrm{DOMINO}}} & bounded torus domino problem (set of initial words admitting a tiling) & \hyperref[sec:p02-w02-p02-035-decision-procedures]{\ref*{sec:p02-w02-p02-035-decision-procedures}} & n/a & \hyperref[def:p02-w02-p02-035-domino]{\ref*{def:p02-w02-p02-035-domino}} & \thepage & Source notation, Def. 6.1.5, with T(n)=2\^{}n; a set of words, not a formula. \\
\phantomsection\label{entry:symbol-registry-06761}\(c(a)\) & \texttt{\detokenize{c(a)}} & colour of an element (its word of unary-predicate truth values) & \hyperref[sec:p02-w02-p02-035-monadic]{\ref*{sec:p02-w02-p02-035-monadic}} & n/a & \hyperref[def:p02-w02-p02-035-monadic]{\ref*{def:p02-w02-p02-035-monadic}} & \thepage & Source notation, proof of Prop. 6.2.1; a word in {0,1}\^{}m. Not a constant symbol. \\
\phantomsection\label{entry:symbol-registry-06762}\(\mathfrak A_{\mathrm{col}}\) & \texttt{\detokenize{\mathfrak A_{\mathrm{col}}}} & colour-collapse structure & \hyperref[sec:p02-w02-p02-035-monadic]{\ref*{sec:p02-w02-p02-035-monadic}} & n/a & \hyperref[thm:p02-w02-p02-035-monadic-small-model]{\ref*{thm:p02-w02-p02-035-monadic-small-model}} & \thepage & Handbook notation for this section. Blank macro filled with the glyph's TeX spelling. \\
\phantomsection\label{entry:symbol-registry-06763}\(\mathfrak B_S\) & \texttt{\detokenize{\mathfrak B_S}} & structure on a colour set S & \hyperref[sec:p02-w02-p02-035-monadic]{\ref*{sec:p02-w02-p02-035-monadic}} & n/a & \hyperref[thm:p02-w02-p02-035-monadic-small-model]{\ref*{thm:p02-w02-p02-035-monadic-small-model}} & \thepage & Handbook notation for this section. Blank macro filled with the glyph's TeX spelling. \\
\phantomsection\label{entry:symbol-registry-06764}\(X_i\) & \texttt{\detokenize{X_i}} & unary predicate symbols of the monadic tiling sentence symbol family member \(X_i\) & \hyperref[sec:p02-w02-p02-035-monadic]{\ref*{sec:p02-w02-p02-035-monadic}} & n/a & \hyperref[def:p02-w02-p02-035-monadic-tiling-sentence]{\ref*{def:p02-w02-p02-035-monadic-tiling-sentence}} & \thepage & Source notation, proof of Thm. 6.2.13. In the Bernays-Schoenfinkel lower bound P\_d is reused as a 2n-ary symbol, declared there. Explicit displayed family; member types and common actual first-use exposed, not one atomic glyph. Family listing, not one atomic glyph. Author macro \texttt{\detokenize{X\_i,}\allowbreak\detokenize{X\_i\^{}{*},}\allowbreak\detokenize{Y\_i,}\allowbreak\detokenize{Y\_i\^{}{*},}\allowbreak\detokenize{N\_i,}\allowbreak\detokenize{P\_d}} not carried to the members. \\
\phantomsection\label{entry:symbol-registry-06765}\(X_i^{*}\) & \texttt{\detokenize{X_i^{*}}} & unary predicate symbols of the monadic tiling sentence symbol family member \(X_i^{*}\) & \hyperref[sec:p02-w02-p02-035-monadic]{\ref*{sec:p02-w02-p02-035-monadic}} & n/a & \hyperref[def:p02-w02-p02-035-monadic-tiling-sentence]{\ref*{def:p02-w02-p02-035-monadic-tiling-sentence}} & \thepage & Source notation, proof of Thm. 6.2.13. In the Bernays-Schoenfinkel lower bound P\_d is reused as a 2n-ary symbol, declared there. Explicit displayed family; member types and common actual first-use exposed, not one atomic glyph. Family listing, not one atomic glyph. Author macro \texttt{\detokenize{X\_i,}\allowbreak\detokenize{X\_i\^{}{*},}\allowbreak\detokenize{Y\_i,}\allowbreak\detokenize{Y\_i\^{}{*},}\allowbreak\detokenize{N\_i,}\allowbreak\detokenize{P\_d}} not carried to the members. \\
\phantomsection\label{entry:symbol-registry-06766}\(Y_i\) & \texttt{\detokenize{Y_i}} & unary predicate symbols of the monadic tiling sentence symbol family member \(Y_i\) & \hyperref[sec:p02-w02-p02-035-monadic]{\ref*{sec:p02-w02-p02-035-monadic}} & n/a & \hyperref[def:p02-w02-p02-035-monadic-tiling-sentence]{\ref*{def:p02-w02-p02-035-monadic-tiling-sentence}} & \thepage & Source notation, proof of Thm. 6.2.13. In the Bernays-Schoenfinkel lower bound P\_d is reused as a 2n-ary symbol, declared there. Explicit displayed family; member types and common actual first-use exposed, not one atomic glyph. Family listing, not one atomic glyph. Author macro \texttt{\detokenize{X\_i,}\allowbreak\detokenize{X\_i\^{}{*},}\allowbreak\detokenize{Y\_i,}\allowbreak\detokenize{Y\_i\^{}{*},}\allowbreak\detokenize{N\_i,}\allowbreak\detokenize{P\_d}} not carried to the members. \\
\phantomsection\label{entry:symbol-registry-06767}\(Y_i^{*}\) & \texttt{\detokenize{Y_i^{*}}} & unary predicate symbols of the monadic tiling sentence symbol family member \(Y_i^{*}\) & \hyperref[sec:p02-w02-p02-035-monadic]{\ref*{sec:p02-w02-p02-035-monadic}} & n/a & \hyperref[def:p02-w02-p02-035-monadic-tiling-sentence]{\ref*{def:p02-w02-p02-035-monadic-tiling-sentence}} & \thepage & Source notation, proof of Thm. 6.2.13. In the Bernays-Schoenfinkel lower bound P\_d is reused as a 2n-ary symbol, declared there. Explicit displayed family; member types and common actual first-use exposed, not one atomic glyph. Family listing, not one atomic glyph. Author macro \texttt{\detokenize{X\_i,}\allowbreak\detokenize{X\_i\^{}{*},}\allowbreak\detokenize{Y\_i,}\allowbreak\detokenize{Y\_i\^{}{*},}\allowbreak\detokenize{N\_i,}\allowbreak\detokenize{P\_d}} not carried to the members. \\
\phantomsection\label{entry:symbol-registry-06768}\(N_i\) & \texttt{\detokenize{N_i}} & unary predicate symbols of the monadic tiling sentence symbol family member \(N_i\) & \hyperref[sec:p02-w02-p02-035-monadic]{\ref*{sec:p02-w02-p02-035-monadic}} & n/a & \hyperref[def:p02-w02-p02-035-monadic-tiling-sentence]{\ref*{def:p02-w02-p02-035-monadic-tiling-sentence}} & \thepage & Source notation, proof of Thm. 6.2.13. In the Bernays-Schoenfinkel lower bound P\_d is reused as a 2n-ary symbol, declared there. Explicit displayed family; member types and common actual first-use exposed, not one atomic glyph. Family listing, not one atomic glyph. Author macro \texttt{\detokenize{X\_i,}\allowbreak\detokenize{X\_i\^{}{*},}\allowbreak\detokenize{Y\_i,}\allowbreak\detokenize{Y\_i\^{}{*},}\allowbreak\detokenize{N\_i,}\allowbreak\detokenize{P\_d}} not carried to the members. \\
\phantomsection\label{entry:symbol-registry-06769}\(P_d\) & \texttt{\detokenize{P_d}} & unary predicate symbols of the monadic tiling sentence symbol family member \(P_d\) & \hyperref[sec:p02-w02-p02-035-monadic]{\ref*{sec:p02-w02-p02-035-monadic}} & n/a & \hyperref[def:p02-w02-p02-035-monadic-tiling-sentence]{\ref*{def:p02-w02-p02-035-monadic-tiling-sentence}} & \thepage & Source notation, proof of Thm. 6.2.13. In the Bernays-Schoenfinkel lower bound P\_d is reused as a 2n-ary symbol, declared there. Explicit displayed family; member types and common actual first-use exposed, not one atomic glyph. Family listing, not one atomic glyph. Author macro \texttt{\detokenize{X\_i,}\allowbreak\detokenize{X\_i\^{}{*},}\allowbreak\detokenize{Y\_i,}\allowbreak\detokenize{Y\_i\^{}{*},}\allowbreak\detokenize{N\_i,}\allowbreak\detokenize{P\_d}} not carried to the members. \\
\phantomsection\label{entry:symbol-registry-06770}\(\psi_w\) & \texttt{\detokenize{\psi_w}} & monadic tiling sentence for an initial word w & \hyperref[sec:p02-w02-p02-035-decision-procedures]{\ref*{sec:p02-w02-p02-035-decision-procedures}} & n/a & \hyperref[def:p02-w02-p02-035-monadic-tiling-sentence]{\ref*{def:p02-w02-p02-035-monadic-tiling-sentence}} & \thepage & Sentences of the object language built from w. Blank macro filled with the glyph's TeX spelling. \\
\phantomsection\label{entry:symbol-registry-06771}\(\chi_w\) & \texttt{\detokenize{\chi_w}} & Bernays-Schoenfinkel tiling sentence for an initial word w & \hyperref[sec:p02-w02-p02-035-decision-procedures]{\ref*{sec:p02-w02-p02-035-decision-procedures}} & n/a & \hyperref[def:p02-w02-p02-035-monadic-tiling-sentence]{\ref*{def:p02-w02-p02-035-monadic-tiling-sentence}} & \thepage & Sentences of the object language built from w. Blank macro filled with the glyph's TeX spelling. \\
\phantomsection\label{entry:symbol-registry-06772}\(\oplus\) & \texttt{\detokenize{\oplus}} & exclusive or (Boolean connective abbreviation) & \hyperref[sec:p02-w02-p02-035-decision-procedures]{\ref*{sec:p02-w02-p02-035-decision-procedures}} & n/a & \hyperref[def:p02-w02-p02-035-monadic-tiling-sentence]{\ref*{def:p02-w02-p02-035-monadic-tiling-sentence}} & \thepage & Abbreviation: A xor B is (A and not B) or (not A and B); source notation. \\
\phantomsection\label{entry:symbol-registry-06773}\(\varphi^{*}\) & \texttt{\detokenize{\varphi^{*}}} & Scott-class form of an FO2 sentence & \hyperref[sec:p02-w02-p02-035-fo2]{\ref*{sec:p02-w02-p02-035-fo2}} & n/a & \hyperref[thm:p02-w02-p02-035-scott]{\ref*{thm:p02-w02-p02-035-scott}} & \thepage & Source notation of Graedel-Kolaitis-Vardi 1997, Sect. 3. Not a dual or a star operation on formulas generally. \\
\phantomsection\label{entry:symbol-registry-06774}\(Q_\chi\) & \texttt{\detokenize{Q_\chi}} & defining relation symbol for a subformula occurrence chi & \hyperref[sec:p02-w02-p02-035-fo2]{\ref*{sec:p02-w02-p02-035-fo2}} & n/a & \hyperref[thm:p02-w02-p02-035-scott]{\ref*{thm:p02-w02-p02-035-scott}} & \thepage & Source notation (Q indexed by a subformula); arity 0, 1 or 2. \\
\phantomsection\label{entry:symbol-registry-06775}\(K\) & \texttt{\detokenize{K}} & set of kings & \hyperref[sec:p02-w02-p02-035-fo2]{\ref*{sec:p02-w02-p02-035-fo2}} & n/a & \hyperref[thm:p02-w02-p02-035-scott-model]{\ref*{thm:p02-w02-p02-035-scott-model}} & \thepage & Source notation, proof of Thm. 4.3. Local to that construction. Blank macro filled with the glyph's TeX spelling. \\
\phantomsection\label{entry:symbol-registry-06776}\(P\) & \texttt{\detokenize{P}} & set of realized 1-types & \hyperref[sec:p02-w02-p02-035-fo2]{\ref*{sec:p02-w02-p02-035-fo2}} & n/a & \hyperref[thm:p02-w02-p02-035-scott-model]{\ref*{thm:p02-w02-p02-035-scott-model}} & \thepage & Source notation, proof of Thm. 4.3. Local to that construction. Blank macro filled with the glyph's TeX spelling. \\
\phantomsection\label{entry:symbol-registry-06777}\(C\) & \texttt{\detokenize{C}} & royal court & \hyperref[sec:p02-w02-p02-035-fo2]{\ref*{sec:p02-w02-p02-035-fo2}} & n/a & \hyperref[thm:p02-w02-p02-035-scott-model]{\ref*{thm:p02-w02-p02-035-scott-model}} & \thepage & Source notation, proof of Thm. 4.3. Local to that construction. Blank macro filled with the glyph's TeX spelling. \\
\phantomsection\label{entry:symbol-registry-06778}\(G(\varphi)\) & \texttt{\detokenize{G(\varphi)}} & grounding of a Bernays-Schoenfinkel sentence & \hyperref[sec:p02-w02-p02-035-bsr]{\ref*{sec:p02-w02-p02-035-bsr}} & n/a & \hyperref[def:p02-w02-p02-035-bsr]{\ref*{def:p02-w02-p02-035-bsr}} & \thepage & Handbook notation; sets of ground quantifier-free formulas. Blank macro filled with the glyph's TeX spelling. \\
\phantomsection\label{entry:symbol-registry-06779}\(G_\equiv\) & \texttt{\detokenize{G_\equiv}} & grounding under an equality pattern & \hyperref[sec:p02-w02-p02-035-bsr]{\ref*{sec:p02-w02-p02-035-bsr}} & n/a & \hyperref[def:p02-w02-p02-035-bsr]{\ref*{def:p02-w02-p02-035-bsr}} & \thepage & Handbook notation; sets of ground quantifier-free formulas. Blank macro filled with the glyph's TeX spelling. \\
\phantomsection\label{entry:symbol-registry-06780}\(\mathrm{Hor}\) & \texttt{\detokenize{\mathrm{Hor}}} & quantifier-free neighbour formulas on bit tuples symbol family member \(\mathrm{Hor}\) & \hyperref[sec:p02-w02-p02-035-bsr]{\ref*{sec:p02-w02-p02-035-bsr}} & n/a & \hyperref[thm:p02-w02-p02-035-bsr-complexity]{\ref*{thm:p02-w02-p02-035-bsr-complexity}} & \thepage & Handbook notation, annex proof only. Explicit displayed family; member types and common actual first-use exposed, not one atomic glyph. Family listing, not one atomic glyph. Author macro \texttt{\detokenize{\mathrm{Hor},}\allowbreak\detokenize{\mathrm{Ver}}} not carried to the members. \\
\phantomsection\label{entry:symbol-registry-06781}\(\mathrm{Ver}\) & \texttt{\detokenize{\mathrm{Ver}}} & quantifier-free neighbour formulas on bit tuples symbol family member \(\mathrm{Ver}\) & \hyperref[sec:p02-w02-p02-035-bsr]{\ref*{sec:p02-w02-p02-035-bsr}} & n/a & \hyperref[thm:p02-w02-p02-035-bsr-complexity]{\ref*{thm:p02-w02-p02-035-bsr-complexity}} & \thepage & Handbook notation, annex proof only. Explicit displayed family; member types and common actual first-use exposed, not one atomic glyph. Family listing, not one atomic glyph. Author macro \texttt{\detokenize{\mathrm{Hor},}\allowbreak\detokenize{\mathrm{Ver}}} not carried to the members. \\
\phantomsection\label{entry:symbol-registry-06782}\(\mathrm{GF}\) & \texttt{\detokenize{\mathrm{GF}}} & guarded fragment (a set of formulas) & \hyperref[sec:p02-w02-p02-035-gf]{\ref*{sec:p02-w02-p02-035-gf}} & n/a & \hyperref[def:p02-w02-p02-035-gf]{\ref*{def:p02-w02-p02-035-gf}} & \thepage & Source notation of Graedel 1999, Def. 2.1 and Sect. 2. Blank macro filled with the glyph's TeX spelling. \\
\phantomsection\label{entry:symbol-registry-06783}\(\mathrm{GF}_k\) & \texttt{\detokenize{\mathrm{GF}_k}} & k-variable part of the guarded fragment (a set of formulas) & \hyperref[sec:p02-w02-p02-035-gf]{\ref*{sec:p02-w02-p02-035-gf}} & n/a & \hyperref[def:p02-w02-p02-035-gf]{\ref*{def:p02-w02-p02-035-gf}} & \thepage & Source notation of Graedel 1999, Def. 2.1 and Sect. 2. Blank macro filled with the glyph's TeX spelling. \\
\phantomsection\label{entry:symbol-registry-06784}\((\exists\bar y\,.\,\alpha)\psi\) & \texttt{\detokenize{(\exists\bar y\,}\allowbreak\detokenize{.\,}\allowbreak\detokenize{\alpha)\psi}} & guarded quantification (abbreviation of exists y(alpha and psi) and forall y(alpha implies psi)) symbol family member \((\exists\bar y\,.\,\alpha)\psi\) & \hyperref[sec:p02-w02-p02-035-gf]{\ref*{sec:p02-w02-p02-035-gf}} & n/a & \hyperref[def:p02-w02-p02-035-gf]{\ref*{def:p02-w02-p02-035-gf}} & \thepage & Source notation (Sect. 2); tuples written with bars here, bold in the source. Original precedence descriptor 'binds like a quantifier' retained; no rank inferred. Explicit displayed family; member types and common actual first-use exposed, not one atomic glyph. Family listing, not one atomic glyph. Author macro \texttt{\detokenize{(\exists\bar y\,}\allowbreak\detokenize{.\,}\allowbreak\detokenize{\alpha)\psi}} not carried to the members. \\
\phantomsection\label{entry:symbol-registry-06785}\((\forall\bar y\,.\,\alpha)\psi\) & \texttt{\detokenize{(\forall\bar y\,}\allowbreak\detokenize{.\,}\allowbreak\detokenize{\alpha)\psi}} & guarded quantification (abbreviation of exists y(alpha and psi) and forall y(alpha implies psi)) symbol family member \((\forall\bar y\,.\,\alpha)\psi\) & \hyperref[sec:p02-w02-p02-035-gf]{\ref*{sec:p02-w02-p02-035-gf}} & n/a & \hyperref[def:p02-w02-p02-035-gf]{\ref*{def:p02-w02-p02-035-gf}} & \thepage & Source notation (Sect. 2); tuples written with bars here, bold in the source. Original precedence descriptor 'binds like a quantifier' retained; no rank inferred. Explicit displayed family; member types and common actual first-use exposed, not one atomic glyph. Family listing, not one atomic glyph. Author macro \texttt{\detokenize{(\exists\bar y\,}\allowbreak\detokenize{.\,}\allowbreak\detokenize{\alpha)\psi}} not carried to the members. \\
\phantomsection\label{entry:symbol-registry-06786}\(\mathrm{At}_k(\psi)\) & \texttt{\detokenize{\mathrm{At}_k}} & set of atoms in x\_1..x\_k that are substitution instances of atoms of psi & \hyperref[sec:p02-w02-p02-035-gf]{\ref*{sec:p02-w02-p02-035-gf}} & n/a & \hyperref[def:p02-w02-p02-035-gf-types]{\ref*{def:p02-w02-p02-035-gf-types}} & \thepage & Handbook notation; restricts the source's full atomic types. \\
\phantomsection\label{entry:symbol-registry-06787}\(\mathrm{atp}^{\psi}_{\mathfrak A}(\bar a)\) & \texttt{\detokenize{\mathrm{atp}^{\psi}}} & psi-type of a tuple (set of At\_k literals true of it) & \hyperref[sec:p02-w02-p02-035-gf]{\ref*{sec:p02-w02-p02-035-gf}} & n/a & \hyperref[def:p02-w02-p02-035-gf-types]{\ref*{def:p02-w02-p02-035-gf-types}} & \thepage & Adapted from the source's atp (Def. 3.2); superscript psi marks the restriction. \\
\phantomsection\label{entry:symbol-registry-06788}\(t^\rho\) & \texttt{\detokenize{t^\rho}} & reduction of a type t along an index map rho & \hyperref[sec:p02-w02-p02-035-gf]{\ref*{sec:p02-w02-p02-035-gf}} & n/a & \hyperref[def:p02-w02-p02-035-gf-types]{\ref*{def:p02-w02-p02-035-gf-types}} & \thepage & Handbook notation for the source's reduction (Def. 3.2); annex only. \\
\phantomsection\label{entry:symbol-registry-06789}\((P,t_0,(\mathrm{ext}_i)_{i\in I})\) & \texttt{\detokenize{P,}\allowbreak\detokenize{t_0,}\allowbreak\detokenize{\mathrm{ext}_i}} & witness for satisfiability: type set; initial type; extension functions & \hyperref[sec:p02-w02-p02-035-gf]{\ref*{sec:p02-w02-p02-035-gf}} & n/a & \hyperref[def:p02-w02-p02-035-gf-types]{\ref*{def:p02-w02-p02-035-gf-types}} & \thepage & Source notation, Def. 3.3. P here is a set of types, unrelated to the predicates P\_d or P\_sigma. \\
\phantomsection\label{entry:symbol-registry-06790}\(T=(V,E)\) & \texttt{\detokenize{T=(V,}\allowbreak\detokenize{E)}} & tree of a k-tree & \hyperref[sec:p02-w02-p02-035-gf]{\ref*{sec:p02-w02-p02-035-gf}} & n/a & \hyperref[def:p02-w02-p02-035-gf-types]{\ref*{def:p02-w02-p02-035-gf-types}} & \thepage & Source notation, Def. 3.5 and proof of Thm. 3.7. T and E here are a tree and its edges; in the acceptance sentence T and E are relation symbols, as in the source. Blank macro filled with the glyph's TeX spelling. \\
\phantomsection\label{entry:symbol-registry-06791}\(F(v)\) & \texttt{\detokenize{F(v)}} & bag map of a k-tree & \hyperref[sec:p02-w02-p02-035-gf]{\ref*{sec:p02-w02-p02-035-gf}} & n/a & \hyperref[def:p02-w02-p02-035-gf-types]{\ref*{def:p02-w02-p02-035-gf-types}} & \thepage & Source notation, Def. 3.5 and proof of Thm. 3.7. T and E here are a tree and its edges; in the acceptance sentence T and E are relation symbols, as in the source. Blank macro filled with the glyph's TeX spelling. \\
\phantomsection\label{entry:symbol-registry-06792}\(\mathfrak F(v)\) & \texttt{\detokenize{\mathfrak F(v)}} & induced substructure on a bag & \hyperref[sec:p02-w02-p02-035-gf]{\ref*{sec:p02-w02-p02-035-gf}} & n/a & \hyperref[def:p02-w02-p02-035-gf-types]{\ref*{def:p02-w02-p02-035-gf-types}} & \thepage & Source notation, Def. 3.5 and proof of Thm. 3.7. T and E here are a tree and its edges; in the acceptance sentence T and E are relation symbols, as in the source. Blank macro filled with the glyph's TeX spelling. \\
\phantomsection\label{entry:symbol-registry-06793}\(\mathfrak A^{+}\) & \texttt{\detokenize{\mathfrak A^{+}}} & extension of a finite structure given by the imported extension theorem & \hyperref[sec:p02-w02-p02-035-gf]{\ref*{sec:p02-w02-p02-035-gf}} & n/a & \hyperref[def:p02-w02-p02-035-herwig-import]{\ref*{def:p02-w02-p02-035-herwig-import}} & \thepage & Source notation, Thm. 3.9. \\
\phantomsection\label{entry:symbol-registry-06794}\(P_j\) & \texttt{\detokenize{P_j}} & stages and fixed point of type elimination symbol family member \(P_j\) & \hyperref[sec:p02-w02-p02-035-gf]{\ref*{sec:p02-w02-p02-035-gf}} & n/a & \hyperref[def:p02-w02-p02-035-gf-procedure]{\ref*{def:p02-w02-p02-035-gf-procedure}} & \thepage & Handbook notation. Explicit displayed family; member types and common actual first-use exposed, not one atomic glyph. Family listing, not one atomic glyph. Author macro \texttt{\detokenize{P\_j,}\allowbreak\detokenize{P\^{}{*}}} not carried to the members. \\
\phantomsection\label{entry:symbol-registry-06795}\(P^{*}\) & \texttt{\detokenize{P^{*}}} & stages and fixed point of type elimination symbol family member \(P^{*}\) & \hyperref[sec:p02-w02-p02-035-gf]{\ref*{sec:p02-w02-p02-035-gf}} & n/a & \hyperref[def:p02-w02-p02-035-gf-procedure]{\ref*{def:p02-w02-p02-035-gf-procedure}} & \thepage & Handbook notation. Explicit displayed family; member types and common actual first-use exposed, not one atomic glyph. Family listing, not one atomic glyph. Author macro \texttt{\detokenize{P\_j,}\allowbreak\detokenize{P\^{}{*}}} not carried to the members. \\
\phantomsection\label{entry:symbol-registry-06796}\(\mathrm{ACCEPT}_{M,w}\) & \texttt{\detokenize{\mathrm{ACCEPT}_{M,w}}} & guarded sentence encoding acceptance of w by an alternating machine M & \hyperref[sec:p02-w02-p02-035-gf]{\ref*{sec:p02-w02-p02-035-gf}} & n/a & \hyperref[def:p02-w02-p02-035-gf-accept]{\ref*{def:p02-w02-p02-035-gf-accept}} & \thepage & Source notation (proof of Thm. 4.4); the handbook's sentence differs as stated in the text. \\
\phantomsection\label{entry:symbol-registry-06797}\(T\) & \texttt{\detokenize{T}} & relation symbol of the acceptance sentence: node & \hyperref[sec:p02-w02-p02-035-gf]{\ref*{sec:p02-w02-p02-035-gf}} & n/a & \hyperref[def:p02-w02-p02-035-gf-accept]{\ref*{def:p02-w02-p02-035-gf-accept}} & \thepage & Source notation, proof of Thm. 4.4 and p. 19 for N. Local to that definition and its annex proof. Blank macro filled with the glyph's TeX spelling. \\
\phantomsection\label{entry:symbol-registry-06798}\(R\) & \texttt{\detokenize{R}} & relation symbol of the acceptance sentence: root & \hyperref[sec:p02-w02-p02-035-gf]{\ref*{sec:p02-w02-p02-035-gf}} & n/a & \hyperref[def:p02-w02-p02-035-gf-accept]{\ref*{def:p02-w02-p02-035-gf-accept}} & \thepage & Source notation, proof of Thm. 4.4 and p. 19 for N. Local to that definition and its annex proof. Blank macro filled with the glyph's TeX spelling. \\
\phantomsection\label{entry:symbol-registry-06799}\(U\) & \texttt{\detokenize{U}} & relation symbol of the acceptance sentence: cell tuples & \hyperref[sec:p02-w02-p02-035-gf]{\ref*{sec:p02-w02-p02-035-gf}} & n/a & \hyperref[def:p02-w02-p02-035-gf-accept]{\ref*{def:p02-w02-p02-035-gf-accept}} & \thepage & Source notation, proof of Thm. 4.4 and p. 19 for N. Local to that definition and its annex proof. Blank macro filled with the glyph's TeX spelling. \\
\phantomsection\label{entry:symbol-registry-06800}\(Z\) & \texttt{\detokenize{Z}} & relation symbol of the acceptance sentence: head & \hyperref[sec:p02-w02-p02-035-gf]{\ref*{sec:p02-w02-p02-035-gf}} & n/a & \hyperref[def:p02-w02-p02-035-gf-accept]{\ref*{def:p02-w02-p02-035-gf-accept}} & \thepage & Source notation, proof of Thm. 4.4 and p. 19 for N. Local to that definition and its annex proof. Blank macro filled with the glyph's TeX spelling. \\
\phantomsection\label{entry:symbol-registry-06801}\(P_\sigma\) & \texttt{\detokenize{P_\sigma}} & relation symbol of the acceptance sentence: symbol & \hyperref[sec:p02-w02-p02-035-gf]{\ref*{sec:p02-w02-p02-035-gf}} & n/a & \hyperref[def:p02-w02-p02-035-gf-accept]{\ref*{def:p02-w02-p02-035-gf-accept}} & \thepage & Source notation, proof of Thm. 4.4 and p. 19 for N. Local to that definition and its annex proof. Blank macro filled with the glyph's TeX spelling. \\
\phantomsection\label{entry:symbol-registry-06802}\(Q_q\) & \texttt{\detokenize{Q_q}} & relation symbol of the acceptance sentence: state & \hyperref[sec:p02-w02-p02-035-gf]{\ref*{sec:p02-w02-p02-035-gf}} & n/a & \hyperref[def:p02-w02-p02-035-gf-accept]{\ref*{def:p02-w02-p02-035-gf-accept}} & \thepage & Source notation, proof of Thm. 4.4 and p. 19 for N. Local to that definition and its annex proof. Blank macro filled with the glyph's TeX spelling. \\
\phantomsection\label{entry:symbol-registry-06803}\(S_1\) & \texttt{\detokenize{S_1}} & relation symbols of the acceptance sentence: successors symbol family member \(S_1\) & \hyperref[sec:p02-w02-p02-035-gf]{\ref*{sec:p02-w02-p02-035-gf}} & n/a & \hyperref[def:p02-w02-p02-035-gf-accept]{\ref*{def:p02-w02-p02-035-gf-accept}} & \thepage & Source notation, proof of Thm. 4.4 and p. 19 for N. Local to that definition and its annex proof. Explicit displayed family; member types and common actual first-use exposed, not one atomic glyph. Blank macro filled with the glyph's TeX spelling. Family listing, not one atomic glyph. \\
\phantomsection\label{entry:symbol-registry-06804}\(S_2\) & \texttt{\detokenize{S_2}} & relation symbols of the acceptance sentence: successors symbol family member \(S_2\) & \hyperref[sec:p02-w02-p02-035-gf]{\ref*{sec:p02-w02-p02-035-gf}} & n/a & \hyperref[def:p02-w02-p02-035-gf-accept]{\ref*{def:p02-w02-p02-035-gf-accept}} & \thepage & Source notation, proof of Thm. 4.4 and p. 19 for N. Local to that definition and its annex proof. Explicit displayed family; member types and common actual first-use exposed, not one atomic glyph. Blank macro filled with the glyph's TeX spelling. Family listing, not one atomic glyph. \\
\phantomsection\label{entry:symbol-registry-06805}\(E\) & \texttt{\detokenize{E}} & relation symbol of the acceptance sentence: matching & \hyperref[sec:p02-w02-p02-035-gf]{\ref*{sec:p02-w02-p02-035-gf}} & n/a & \hyperref[def:p02-w02-p02-035-gf-accept]{\ref*{def:p02-w02-p02-035-gf-accept}} & \thepage & Source notation, proof of Thm. 4.4 and p. 19 for N. Local to that definition and its annex proof. Blank macro filled with the glyph's TeX spelling. \\
\phantomsection\label{entry:symbol-registry-06806}\(N\) & \texttt{\detokenize{N}} & relation symbol of the acceptance sentence: left marker & \hyperref[sec:p02-w02-p02-035-gf]{\ref*{sec:p02-w02-p02-035-gf}} & n/a & \hyperref[def:p02-w02-p02-035-gf-accept]{\ref*{def:p02-w02-p02-035-gf-accept}} & \thepage & Source notation, proof of Thm. 4.4 and p. 19 for N. Local to that definition and its annex proof. Blank macro filled with the glyph's TeX spelling. \\
\phantomsection\label{entry:symbol-registry-06807}\(\hat r(a,b)\) & \texttt{\detokenize{\hat r(a,}\allowbreak\detokenize{b)}} & tuple in \(\{a,b\}^s\) coding the number r & \hyperref[sec:p02-w02-p02-035-gf]{\ref*{sec:p02-w02-p02-035-gf}} & n/a & \hyperref[def:p02-w02-p02-035-gf-accept]{\ref*{def:p02-w02-p02-035-gf-accept}} & \thepage & Handbook notation; the source writes bold i. \\
\phantomsection\label{entry:symbol-registry-06808}\(\mathit{succ}\) & \texttt{\detokenize{\mathit{succ}}} & quantifier-free formulas comparing coded numbers symbol family member \(\mathit{succ}\) & \hyperref[sec:p02-w02-p02-035-gf]{\ref*{sec:p02-w02-p02-035-gf}} & n/a & \hyperref[def:p02-w02-p02-035-gf-accept]{\ref*{def:p02-w02-p02-035-gf-accept}} & \thepage & succ follows the source; diff is handbook notation. Not the arithmetic successor function. Explicit displayed family; member types and common actual first-use exposed, not one atomic glyph. Family listing, not one atomic glyph. Author macro \texttt{\detokenize{\mathit{succ},}\allowbreak\detokenize{\mathit{diff}}} not carried to the members. \\
\phantomsection\label{entry:symbol-registry-06809}\(\mathit{diff}\) & \texttt{\detokenize{\mathit{diff}}} & quantifier-free formulas comparing coded numbers symbol family member \(\mathit{diff}\) & \hyperref[sec:p02-w02-p02-035-gf]{\ref*{sec:p02-w02-p02-035-gf}} & n/a & \hyperref[def:p02-w02-p02-035-gf-accept]{\ref*{def:p02-w02-p02-035-gf-accept}} & \thepage & succ follows the source; diff is handbook notation. Not the arithmetic successor function. Explicit displayed family; member types and common actual first-use exposed, not one atomic glyph. Family listing, not one atomic glyph. Author macro \texttt{\detokenize{\mathit{succ},}\allowbreak\detokenize{\mathit{diff}}} not carried to the members. \\
\phantomsection\label{entry:symbol-registry-06810}\(\mathrm{EXIST}\) & \texttt{\detokenize{\mathrm{EXIST}}} & state disjunctions and transition formulas of the acceptance sentence symbol family member \(\mathrm{EXIST}\) & \hyperref[sec:p02-w02-p02-035-gf]{\ref*{sec:p02-w02-p02-035-gf}} & n/a & \hyperref[def:p02-w02-p02-035-gf-accept]{\ref*{def:p02-w02-p02-035-gf-accept}} & \thepage & Source notation, proof of Thm. 4.4. Explicit displayed family; member types and common actual first-use exposed, not one atomic glyph. Family listing, not one atomic glyph. Author macro \texttt{\detokenize{\mathrm{EXIST},}\allowbreak\detokenize{\mathrm{NEXT}\_i}} not carried to the members. \\
\phantomsection\label{entry:symbol-registry-06811}\(\mathrm{UNIV}\) & \texttt{\detokenize{\mathrm{UNIV}}} & state disjunctions and transition formulas of the acceptance sentence symbol family member \(\mathrm{UNIV}\) & \hyperref[sec:p02-w02-p02-035-gf]{\ref*{sec:p02-w02-p02-035-gf}} & n/a & \hyperref[def:p02-w02-p02-035-gf-accept]{\ref*{def:p02-w02-p02-035-gf-accept}} & \thepage & Source notation, proof of Thm. 4.4. Explicit displayed family; member types and common actual first-use exposed, not one atomic glyph. Family listing, not one atomic glyph. Author macro \texttt{\detokenize{\mathrm{EXIST},}\allowbreak\detokenize{\mathrm{NEXT}\_i}} not carried to the members. \\
\phantomsection\label{entry:symbol-registry-06812}\(\mathrm{REJECT}\) & \texttt{\detokenize{\mathrm{REJECT}}} & state disjunctions and transition formulas of the acceptance sentence symbol family member \(\mathrm{REJECT}\) & \hyperref[sec:p02-w02-p02-035-gf]{\ref*{sec:p02-w02-p02-035-gf}} & n/a & \hyperref[def:p02-w02-p02-035-gf-accept]{\ref*{def:p02-w02-p02-035-gf-accept}} & \thepage & Source notation, proof of Thm. 4.4. Explicit displayed family; member types and common actual first-use exposed, not one atomic glyph. Family listing, not one atomic glyph. Author macro \texttt{\detokenize{\mathrm{EXIST},}\allowbreak\detokenize{\mathrm{NEXT}\_i}} not carried to the members. \\
\phantomsection\label{entry:symbol-registry-06813}\(\mathrm{NEXT}_i\) & \texttt{\detokenize{\mathrm{NEXT}_i}} & state disjunctions and transition formulas of the acceptance sentence symbol family member \(\mathrm{NEXT}_i\) & \hyperref[sec:p02-w02-p02-035-gf]{\ref*{sec:p02-w02-p02-035-gf}} & n/a & \hyperref[def:p02-w02-p02-035-gf-accept]{\ref*{def:p02-w02-p02-035-gf-accept}} & \thepage & Source notation, proof of Thm. 4.4. Explicit displayed family; member types and common actual first-use exposed, not one atomic glyph. Family listing, not one atomic glyph. Author macro \texttt{\detokenize{\mathrm{EXIST},}\allowbreak\detokenize{\mathrm{NEXT}\_i}} not carried to the members. \\
\phantomsection\label{entry:symbol-registry-06814}\(\mathrm{NX}_i\) & \texttt{\detokenize{\mathrm{NX}_i}} & transition formula and unary symbols of the two-variable acceptance sentence symbol family member \(\mathrm{NX}_i\) & \hyperref[annex:p02-w02-p02-035-gf-hard]{\ref*{annex:p02-w02-p02-035-gf-hard}} & n/a & \hyperref[annex:p02-w02-p02-035-gf-hard]{\ref*{annex:p02-w02-p02-035-gf-hard}} & \thepage & Handbook notation, annex only. Explicit displayed family; member types and common actual first-use exposed, not one atomic glyph. Family listing, not one atomic glyph. Author macro \texttt{\detokenize{\mathrm{NX}\_i,}\allowbreak\detokenize{Z\_r,}\allowbreak\detokenize{P\_{\sigma,r}}} not carried to the members. \\
\phantomsection\label{entry:symbol-registry-06815}\(Z_r\) & \texttt{\detokenize{Z_r}} & transition formula and unary symbols of the two-variable acceptance sentence symbol family member \(Z_r\) & \hyperref[annex:p02-w02-p02-035-gf-hard]{\ref*{annex:p02-w02-p02-035-gf-hard}} & n/a & \hyperref[annex:p02-w02-p02-035-gf-hard]{\ref*{annex:p02-w02-p02-035-gf-hard}} & \thepage & Handbook notation, annex only. Explicit displayed family; member types and common actual first-use exposed, not one atomic glyph. Family listing, not one atomic glyph. Author macro \texttt{\detokenize{\mathrm{NX}\_i,}\allowbreak\detokenize{Z\_r,}\allowbreak\detokenize{P\_{\sigma,r}}} not carried to the members. \\
\phantomsection\label{entry:symbol-registry-06816}\(P_{\sigma,r}\) & \texttt{\detokenize{P_{\sigma,r}}} & transition formula and unary symbols of the two-variable acceptance sentence symbol family member \(P_{\sigma,r}\) & \hyperref[annex:p02-w02-p02-035-gf-hard]{\ref*{annex:p02-w02-p02-035-gf-hard}} & n/a & \hyperref[annex:p02-w02-p02-035-gf-hard]{\ref*{annex:p02-w02-p02-035-gf-hard}} & \thepage & Handbook notation, annex only. Explicit displayed family; member types and common actual first-use exposed, not one atomic glyph. Family listing, not one atomic glyph. Author macro \texttt{\detokenize{\mathrm{NX}\_i,}\allowbreak\detokenize{Z\_r,}\allowbreak\detokenize{P\_{\sigma,r}}} not carried to the members. \\
\phantomsection\label{entry:symbol-registry-06817}\(\kappa\) & \texttt{\detokenize{\kappa}} & place pattern of an atom (constant index or the placeholder *) and the relation symbol for the non-constant places symbol family member \(\kappa\) & \hyperref[annex:p02-w02-p02-035-gf-bounded-constants]{\ref*{annex:p02-w02-p02-035-gf-bounded-constants}} & n/a & \hyperref[annex:p02-w02-p02-035-gf-bounded-constants]{\ref*{annex:p02-w02-p02-035-gf-bounded-constants}} & \thepage & Handbook notation, annex only; * is a placeholder value, not an operator. Explicit displayed family; member types and common actual first-use exposed, not one atomic glyph. Family listing, not one atomic glyph. Author macro \texttt{\detokenize{R\_\kappa}} not carried to the members. \\
\phantomsection\label{entry:symbol-registry-06818}\(R_\kappa\) & \texttt{\detokenize{R_\kappa}} & place pattern of an atom (constant index or the placeholder *) and the relation symbol for the non-constant places symbol family member \(R_\kappa\) & \hyperref[annex:p02-w02-p02-035-gf-bounded-constants]{\ref*{annex:p02-w02-p02-035-gf-bounded-constants}} & n/a & \hyperref[annex:p02-w02-p02-035-gf-bounded-constants]{\ref*{annex:p02-w02-p02-035-gf-bounded-constants}} & \thepage & Handbook notation, annex only; * is a placeholder value, not an operator. Explicit displayed family; member types and common actual first-use exposed, not one atomic glyph. Family listing, not one atomic glyph. Author macro \texttt{\detokenize{R\_\kappa}} not carried to the members. \\
\phantomsection\label{entry:symbol-registry-06819}\(\ast\) & \texttt{\detokenize{\ast}} & place pattern of an atom (constant index or the placeholder *) and the relation symbol for the non-constant places symbol family member \(\ast\) & \hyperref[annex:p02-w02-p02-035-gf-bounded-constants]{\ref*{annex:p02-w02-p02-035-gf-bounded-constants}} & n/a & \hyperref[annex:p02-w02-p02-035-gf-bounded-constants]{\ref*{annex:p02-w02-p02-035-gf-bounded-constants}} & \thepage & Handbook notation, annex only; * is a placeholder value, not an operator. Explicit displayed family; member types and common actual first-use exposed, not one atomic glyph. Family listing, not one atomic glyph. Author macro \texttt{\detokenize{R\_\kappa}} not carried to the members. \\
\phantomsection\label{entry:symbol-registry-06820}\(\mathrm{sp}_\sigma(\theta)\) & \texttt{\detokenize{\mathrm{sp}_\sigma}} & split translation of a guarded formula under a map sigma from variables to {*;1..p} & \hyperref[annex:p02-w02-p02-035-gf-bounded-constants]{\ref*{annex:p02-w02-p02-035-gf-bounded-constants}} & n/a & \hyperref[annex:p02-w02-p02-035-gf-bounded-constants]{\ref*{annex:p02-w02-p02-035-gf-bounded-constants}} & \thepage & Handbook notation, annex only. \\
\phantomsection\label{entry:symbol-registry-06821}\(\mathit{val}\) & \texttt{\detokenize{\mathit{val}}} & truth-value map on ground patterns & \hyperref[annex:p02-w02-p02-035-gf-bounded-constants]{\ref*{annex:p02-w02-p02-035-gf-bounded-constants}} & n/a & \hyperref[annex:p02-w02-p02-035-gf-bounded-constants]{\ref*{annex:p02-w02-p02-035-gf-bounded-constants}} & \thepage & Handbook notation, annex only. Blank macro filled with the glyph's TeX spelling. \\
\phantomsection\label{entry:symbol-registry-06822}\(\varphi^{\mathit{val}}\) & \texttt{\detokenize{\varphi^{\mathit{val}}}} & split and simplified sentence & \hyperref[annex:p02-w02-p02-035-gf-bounded-constants]{\ref*{annex:p02-w02-p02-035-gf-bounded-constants}} & n/a & \hyperref[annex:p02-w02-p02-035-gf-bounded-constants]{\ref*{annex:p02-w02-p02-035-gf-bounded-constants}} & \thepage & Handbook notation, annex only. Blank macro filled with the glyph's TeX spelling. \\
\phantomsection\label{entry:symbol-registry-06823}\(\varphi^{\flat}\) & \texttt{\detokenize{\varphi^{\flat}}} & flattened sentence & \hyperref[annex:p02-w02-p02-035-gf-bounded-constants]{\ref*{annex:p02-w02-p02-035-gf-bounded-constants}} & n/a & \hyperref[annex:p02-w02-p02-035-gf-bounded-constants]{\ref*{annex:p02-w02-p02-035-gf-bounded-constants}} & \thepage & Handbook notation, annex only. Blank macro filled with the glyph's TeX spelling. \\
\phantomsection\label{entry:symbol-registry-06824}\(\langle\kappa,\bar u\rangle\) & \texttt{\detokenize{\langle\kappa;\bar u\rangle}} & tuple obtained by filling the * places of kappa with u and the other places with constant values & \hyperref[annex:p02-w02-p02-035-gf-bounded-constants]{\ref*{annex:p02-w02-p02-035-gf-bounded-constants}} & n/a & \hyperref[annex:p02-w02-p02-035-gf-bounded-constants]{\ref*{annex:p02-w02-p02-035-gf-bounded-constants}} & \thepage & Handbook notation, annex only; not a pairing function. \\
\phantomsection\label{entry:symbol-registry-06825}\(\mathrm{Sat}(X)\) & \texttt{\detokenize{\mathrm{Sat}(X)}} & satisfiable / valid / finitely satisfiable / infinity-axiom / unsatisfiable members of a class X of sentences symbol family member \(\mathrm{Sat}(X)\) & \hyperref[sec:p02-w02-p02-036-finite-validity]{\ref*{sec:p02-w02-p02-036-finite-validity}} & n/a & \hyperref[def:p02-w02-p02-036-classes]{\ref*{def:p02-w02-p02-036-classes}} & \thepage & Source notation of Boerger-Graedel-Gurevich 1997, Def. 2.1.27 (p. 33); X a class of formulas, X = FO written without argument, main case X = Sent(L) for a finite relational signature L. Values are sets of formula codes. Explicit displayed family; member types and common actual first-use exposed, not one atomic glyph. Family listing, not one atomic glyph. Author macro \texttt{\detokenize{\mathrm{Sat},}\allowbreak\detokenize{\mathrm{Val},}\allowbreak\detokenize{\mathrm{Fin\text{-}sat},}\allowbreak\detokenize{\mathrm{Inf\text{-}axioms},}\allowbreak\detokenize{\mathrm{Non\text{-}sat}}} not carried to the members. \\
\phantomsection\label{entry:symbol-registry-06826}\(\mathrm{Val}(X)\) & \texttt{\detokenize{\mathrm{Val}(X)}} & satisfiable / valid / finitely satisfiable / infinity-axiom / unsatisfiable members of a class X of sentences symbol family member \(\mathrm{Val}(X)\) & \hyperref[sec:p02-w02-p02-036-finite-validity]{\ref*{sec:p02-w02-p02-036-finite-validity}} & n/a & \hyperref[def:p02-w02-p02-036-classes]{\ref*{def:p02-w02-p02-036-classes}} & \thepage & Source notation of Boerger-Graedel-Gurevich 1997, Def. 2.1.27 (p. 33); X a class of formulas, X = FO written without argument, main case X = Sent(L) for a finite relational signature L. Values are sets of formula codes. Explicit displayed family; member types and common actual first-use exposed, not one atomic glyph. Family listing, not one atomic glyph. Author macro \texttt{\detokenize{\mathrm{Sat},}\allowbreak\detokenize{\mathrm{Val},}\allowbreak\detokenize{\mathrm{Fin\text{-}sat},}\allowbreak\detokenize{\mathrm{Inf\text{-}axioms},}\allowbreak\detokenize{\mathrm{Non\text{-}sat}}} not carried to the members. \\
\phantomsection\label{entry:symbol-registry-06827}\(\mathrm{Fin\text{-}sat}(X)\) & \texttt{\detokenize{\mathrm{Fin\text{-}sat}(X)}} & satisfiable / valid / finitely satisfiable / infinity-axiom / unsatisfiable members of a class X of sentences symbol family member \(\mathrm{Fin\text{-}sat}(X)\) & \hyperref[sec:p02-w02-p02-036-finite-validity]{\ref*{sec:p02-w02-p02-036-finite-validity}} & n/a & \hyperref[def:p02-w02-p02-036-classes]{\ref*{def:p02-w02-p02-036-classes}} & \thepage & Source notation of Boerger-Graedel-Gurevich 1997, Def. 2.1.27 (p. 33); X a class of formulas, X = FO written without argument, main case X = Sent(L) for a finite relational signature L. Values are sets of formula codes. Explicit displayed family; member types and common actual first-use exposed, not one atomic glyph. Family listing, not one atomic glyph. Author macro \texttt{\detokenize{\mathrm{Sat},}\allowbreak\detokenize{\mathrm{Val},}\allowbreak\detokenize{\mathrm{Fin\text{-}sat},}\allowbreak\detokenize{\mathrm{Inf\text{-}axioms},}\allowbreak\detokenize{\mathrm{Non\text{-}sat}}} not carried to the members. \\
\phantomsection\label{entry:symbol-registry-06828}\(\mathrm{Inf\text{-}axioms}(X)\) & \texttt{\detokenize{\mathrm{Inf\text{-}axioms}(X)}} & satisfiable / valid / finitely satisfiable / infinity-axiom / unsatisfiable members of a class X of sentences symbol family member \(\mathrm{Inf\text{-}axioms}(X)\) & \hyperref[sec:p02-w02-p02-036-finite-validity]{\ref*{sec:p02-w02-p02-036-finite-validity}} & n/a & \hyperref[def:p02-w02-p02-036-classes]{\ref*{def:p02-w02-p02-036-classes}} & \thepage & Source notation of Boerger-Graedel-Gurevich 1997, Def. 2.1.27 (p. 33); X a class of formulas, X = FO written without argument, main case X = Sent(L) for a finite relational signature L. Values are sets of formula codes. Explicit displayed family; member types and common actual first-use exposed, not one atomic glyph. Family listing, not one atomic glyph. Author macro \texttt{\detokenize{\mathrm{Sat},}\allowbreak\detokenize{\mathrm{Val},}\allowbreak\detokenize{\mathrm{Fin\text{-}sat},}\allowbreak\detokenize{\mathrm{Inf\text{-}axioms},}\allowbreak\detokenize{\mathrm{Non\text{-}sat}}} not carried to the members. \\
\phantomsection\label{entry:symbol-registry-06829}\(\mathrm{Non\text{-}sat}(X)\) & \texttt{\detokenize{\mathrm{Non\text{-}sat}(X)}} & satisfiable / valid / finitely satisfiable / infinity-axiom / unsatisfiable members of a class X of sentences symbol family member \(\mathrm{Non\text{-}sat}(X)\) & \hyperref[sec:p02-w02-p02-036-finite-validity]{\ref*{sec:p02-w02-p02-036-finite-validity}} & n/a & \hyperref[def:p02-w02-p02-036-classes]{\ref*{def:p02-w02-p02-036-classes}} & \thepage & Source notation of Boerger-Graedel-Gurevich 1997, Def. 2.1.27 (p. 33); X a class of formulas, X = FO written without argument, main case X = Sent(L) for a finite relational signature L. Values are sets of formula codes. Explicit displayed family; member types and common actual first-use exposed, not one atomic glyph. Family listing, not one atomic glyph. Author macro \texttt{\detokenize{\mathrm{Sat},}\allowbreak\detokenize{\mathrm{Val},}\allowbreak\detokenize{\mathrm{Fin\text{-}sat},}\allowbreak\detokenize{\mathrm{Inf\text{-}axioms},}\allowbreak\detokenize{\mathrm{Non\text{-}sat}}} not carried to the members. \\
\phantomsection\label{entry:symbol-registry-06830}\(\Lang_{\mathcal T}\) & \texttt{\detokenize{\Lang_{\mathcal T}}} & run signature of a Turing machine T & \hyperref[sec:p02-w02-p02-036-finite-validity]{\ref*{sec:p02-w02-p02-036-finite-validity}} & n/a & \hyperref[def:p02-w02-p02-036-run-sentence]{\ref*{def:p02-w02-p02-036-run-sentence}} & \thepage & Binary <, Hd, Cl\_a (a a tape symbol); unary St\_q (q a state). \\
\phantomsection\label{entry:symbol-registry-06831}\(\mathsf{St}_q\) & \texttt{\detokenize{\mathsf{St}_q}} & relation symbol of the run signature: state at a moment & \hyperref[sec:p02-w02-p02-036-finite-validity]{\ref*{sec:p02-w02-p02-036-finite-validity}} & n/a & \hyperref[def:p02-w02-p02-036-run-sentence]{\ref*{def:p02-w02-p02-036-run-sentence}} & \thepage & Object-language relation symbols; St\_q unary, Hd and Cl\_a binary. Blank macro filled with the glyph's TeX spelling. \\
\phantomsection\label{entry:symbol-registry-06832}\(\mathsf{Hd}\) & \texttt{\detokenize{\mathsf{Hd}}} & relation symbol of the run signature: head cell at a moment & \hyperref[sec:p02-w02-p02-036-finite-validity]{\ref*{sec:p02-w02-p02-036-finite-validity}} & n/a & \hyperref[def:p02-w02-p02-036-run-sentence]{\ref*{def:p02-w02-p02-036-run-sentence}} & \thepage & Object-language relation symbols; St\_q unary, Hd and Cl\_a binary. Blank macro filled with the glyph's TeX spelling. \\
\phantomsection\label{entry:symbol-registry-06833}\(\mathsf{Cl}_a\) & \texttt{\detokenize{\mathsf{Cl}_a}} & relation symbol of the run signature: cell content at a moment & \hyperref[sec:p02-w02-p02-036-finite-validity]{\ref*{sec:p02-w02-p02-036-finite-validity}} & n/a & \hyperref[def:p02-w02-p02-036-run-sentence]{\ref*{def:p02-w02-p02-036-run-sentence}} & \thepage & Object-language relation symbols; St\_q unary, Hd and Cl\_a binary. Blank macro filled with the glyph's TeX spelling. \\
\phantomsection\label{entry:symbol-registry-06834}\(\mathsf{Nx}\) & \texttt{\detokenize{\mathsf{Nx}}} & defined formula abbreviation: immediate successor & \hyperref[sec:p02-w02-p02-036-finite-validity]{\ref*{sec:p02-w02-p02-036-finite-validity}} & n/a & \hyperref[def:p02-w02-p02-036-run-sentence]{\ref*{def:p02-w02-p02-036-run-sentence}} & \thepage & Abbreviations of first-order formulas in the run signature; not new symbols of the signature. Blank macro filled with the glyph's TeX spelling. \\
\phantomsection\label{entry:symbol-registry-06835}\(\mathsf{Fst}\) & \texttt{\detokenize{\mathsf{Fst}}} & defined formula abbreviation: least element & \hyperref[sec:p02-w02-p02-036-finite-validity]{\ref*{sec:p02-w02-p02-036-finite-validity}} & n/a & \hyperref[def:p02-w02-p02-036-run-sentence]{\ref*{def:p02-w02-p02-036-run-sentence}} & \thepage & Abbreviations of first-order formulas in the run signature; not new symbols of the signature. Blank macro filled with the glyph's TeX spelling. \\
\phantomsection\label{entry:symbol-registry-06836}\(\mathsf{Mv}_d\) & \texttt{\detokenize{\mathsf{Mv}_d}} & defined formula abbreviation: head move in direction d & \hyperref[sec:p02-w02-p02-036-finite-validity]{\ref*{sec:p02-w02-p02-036-finite-validity}} & n/a & \hyperref[def:p02-w02-p02-036-run-sentence]{\ref*{def:p02-w02-p02-036-run-sentence}} & \thepage & Abbreviations of first-order formulas in the run signature; not new symbols of the signature. Blank macro filled with the glyph's TeX spelling. \\
\phantomsection\label{entry:symbol-registry-06837}\(\mathsf{In}_w\) & \texttt{\detokenize{\mathsf{In}_w}} & defined formula abbreviation: initial tape for input w & \hyperref[sec:p02-w02-p02-036-finite-validity]{\ref*{sec:p02-w02-p02-036-finite-validity}} & n/a & \hyperref[def:p02-w02-p02-036-run-sentence]{\ref*{def:p02-w02-p02-036-run-sentence}} & \thepage & Abbreviations of first-order formulas in the run signature; not new symbols of the signature. Blank macro filled with the glyph's TeX spelling. \\
\phantomsection\label{entry:symbol-registry-06838}\(\psi_{\mathcal T,w}\) & \texttt{\detokenize{\psi_{\mathcal T,w}}} & run sentence of machine T on input word w & \hyperref[sec:p02-w02-p02-036-finite-validity]{\ref*{sec:p02-w02-p02-036-finite-validity}} & n/a & \hyperref[def:p02-w02-p02-036-run-sentence]{\ref*{def:p02-w02-p02-036-run-sentence}} & \thepage & A sentence of the run signature; has a finite model iff T reaches its accepting state on w. \\
\phantomsection\label{entry:symbol-registry-06839}\(\pi_j(x,y)\) & \texttt{\detokenize{\pi_j}} & path formula of index j in the single binary relation E & \hyperref[def:p02-w02-p02-036-gadget]{\ref*{def:p02-w02-p02-036-gadget}} & n/a & \hyperref[def:p02-w02-p02-036-gadget]{\ref*{def:p02-w02-p02-036-gadget}} & \thepage & Formula abbreviation: there is a path of exactly j unlooped vertices from x to y. Scoped to the path translation; the global pi rows (proof code, projections, paths of transition structures) have other roles. \\
\phantomsection\label{entry:symbol-registry-06840}\(\varphi^{I}\) & \texttt{\detokenize{\varphi^{I}}} & path translation of a formula and of a sentence into the signature {E} symbol family member \(\varphi^{I}\) & \hyperref[sec:p02-w02-p02-036-finite-validity]{\ref*{sec:p02-w02-p02-036-finite-validity}} & n/a & \hyperref[def:p02-w02-p02-036-gadget]{\ref*{def:p02-w02-p02-036-gadget}} & \thepage & Postfix syntactic operations; varphi-sharp adds the conjunct that some vertex is looped. Explicit displayed family; member types and common actual first-use exposed, not one atomic glyph. Family listing, not one atomic glyph. Author macro \texttt{\detokenize{\^{}{I},}\allowbreak\detokenize{\^{}{\sharp}}} not carried to the members. \\
\phantomsection\label{entry:symbol-registry-06841}\(\varphi^{\sharp}\) & \texttt{\detokenize{\varphi^{\sharp}}} & path translation of a formula and of a sentence into the signature {E} symbol family member \(\varphi^{\sharp}\) & \hyperref[sec:p02-w02-p02-036-finite-validity]{\ref*{sec:p02-w02-p02-036-finite-validity}} & n/a & \hyperref[def:p02-w02-p02-036-gadget]{\ref*{def:p02-w02-p02-036-gadget}} & \thepage & Postfix syntactic operations; varphi-sharp adds the conjunct that some vertex is looped. Explicit displayed family; member types and common actual first-use exposed, not one atomic glyph. Family listing, not one atomic glyph. Author macro \texttt{\detokenize{\^{}{I},}\allowbreak\detokenize{\^{}{\sharp}}} not carried to the members. \\
\phantomsection\label{entry:symbol-registry-06842}\(D(\mathcal G)\) & \texttt{\detokenize{D(\mathcal G)}} & decoded structure of a graph with a loop & \hyperref[sec:p02-w02-p02-036-finite-validity]{\ref*{sec:p02-w02-p02-036-finite-validity}} & n/a & \hyperref[def:p02-w02-p02-036-gadget]{\ref*{def:p02-w02-p02-036-gadget}} & \thepage & Operations on structures; D(G(M)) = M (Lemma thm:p02-w02-p02-036-gadget). Blank macro filled with the glyph's TeX spelling. \\
\phantomsection\label{entry:symbol-registry-06843}\(G(\mathcal M)\) & \texttt{\detokenize{G(\mathcal M)}} & path graph of a structure & \hyperref[sec:p02-w02-p02-036-finite-validity]{\ref*{sec:p02-w02-p02-036-finite-validity}} & n/a & \hyperref[def:p02-w02-p02-036-gadget]{\ref*{def:p02-w02-p02-036-gadget}} & \thepage & Operations on structures; D(G(M)) = M (Lemma thm:p02-w02-p02-036-gadget). Blank macro filled with the glyph's TeX spelling. \\
\phantomsection\label{entry:symbol-registry-06844}\(c^f_i\) & \texttt{\detokenize{c^f_i}} & auxiliary vertex number i of the path for a fact f & \hyperref[sec:p02-w02-p02-036-finite-validity]{\ref*{sec:p02-w02-p02-036-finite-validity}} & n/a & \hyperref[def:p02-w02-p02-036-gadget]{\ref*{def:p02-w02-p02-036-gadget}} & \thepage & Fresh vertices of the path graph; unlooped. \\
\phantomsection\label{entry:symbol-registry-06845}\(\chi[R/E]\) & \texttt{\detokenize{[R/}\allowbreak\detokenize{E]}} & renaming of the relation symbol E to R in a sentence & \hyperref[sec:p02-w02-p02-036-finite-validity]{\ref*{sec:p02-w02-p02-036-finite-validity}} & n/a & \hyperref[thm:p02-w02-p02-036-transfer]{\ref*{thm:p02-w02-p02-036-transfer}} & \thepage & Symbol renaming, not term substitution; the global [t/x] row is capture-avoiding term substitution. \\
\phantomsection\label{entry:symbol-registry-06849}\(\prec_K\) & \texttt{\detokenize{\prec_K}} & the source's infinity axiom in one binary relation symbol K & \hyperref[sec:p02-w02-p02-036-finite-validity]{\ref*{sec:p02-w02-p02-036-finite-validity}} & n/a & \hyperref[def:p02-w02-p02-036-infinity-axiom]{\ref*{def:p02-w02-p02-036-infinity-axiom}} & \thepage & Source notation of Boerger-Graedel-Gurevich 1997 p. 33: forall x exists u forall y (not Kxx and Kxu and (Kyx -> Kyu)). A sentence name; prec is not used here as an order symbol. \\
\phantomsection\label{entry:symbol-registry-06850}\(X\) & \texttt{\detokenize{X}} & disjoint sets of natural numbers in the definition of recursive inseparability symbol family member \(X\) & \hyperref[def:p02-w02-p02-036-inseparable]{\ref*{def:p02-w02-p02-036-inseparable}} & n/a & \hyperref[def:p02-w02-p02-036-inseparable]{\ref*{def:p02-w02-p02-036-inseparable}} & \thepage & Bound letters of the source definition (BGG Def. 2.1.28); in proofs a separating set is written W. Explicit displayed family; member types and common actual first-use exposed, not one atomic glyph. Family listing, not one atomic glyph. Author macro \texttt{\detokenize{X,}\allowbreak\detokenize{Y}} not carried to the members. \\
\phantomsection\label{entry:symbol-registry-06851}\(Y\) & \texttt{\detokenize{Y}} & disjoint sets of natural numbers in the definition of recursive inseparability symbol family member \(Y\) & \hyperref[def:p02-w02-p02-036-inseparable]{\ref*{def:p02-w02-p02-036-inseparable}} & n/a & \hyperref[def:p02-w02-p02-036-inseparable]{\ref*{def:p02-w02-p02-036-inseparable}} & \thepage & Bound letters of the source definition (BGG Def. 2.1.28); in proofs a separating set is written W. Explicit displayed family; member types and common actual first-use exposed, not one atomic glyph. Family listing, not one atomic glyph. Author macro \texttt{\detokenize{X,}\allowbreak\detokenize{Y}} not carried to the members. \\
\phantomsection\label{entry:symbol-registry-06852}\(R\) & \texttt{\detokenize{R}} & recursive set separating X from Y in the definition of recursive inseparability & \hyperref[def:p02-w02-p02-036-inseparable]{\ref*{def:p02-w02-p02-036-inseparable}} & n/a & \hyperref[def:p02-w02-p02-036-inseparable]{\ref*{def:p02-w02-p02-036-inseparable}} & \thepage & Bound letters of the source definition (BGG Def. 2.1.28); in proofs a separating set is written W. \\
\phantomsection\label{entry:symbol-registry-06853}\(V\) & \texttt{\detokenize{V}} & set of valid program codes & \hyperref[sec:p02-w02-p02-036-inseparability]{\ref*{sec:p02-w02-p02-036-inseparability}} & n/a & \hyperref[def:p02-w02-p02-036-halting]{\ref*{def:p02-w02-p02-036-halting}} & \thepage & Program codes of sec:universal-reductions-01; M\_e is a Turing machine tuple, not a structure. Blank macro filled with the glyph's TeX spelling. \\
\phantomsection\label{entry:symbol-registry-06854}\(M_e\) & \texttt{\detokenize{M_e}} & machine decoded from a valid code e & \hyperref[sec:p02-w02-p02-036-inseparability]{\ref*{sec:p02-w02-p02-036-inseparability}} & n/a & \hyperref[def:p02-w02-p02-036-halting]{\ref*{def:p02-w02-p02-036-halting}} & \thepage & Program codes of sec:universal-reductions-01; M\_e is a Turing machine tuple, not a structure. Blank macro filled with the glyph's TeX spelling. \\
\phantomsection\label{entry:symbol-registry-06855}\(H_1\) & \texttt{\detokenize{H_1}} & codes of machines that reject / accept / never halt on the empty input symbol family member \(H_1\) & \hyperref[sec:p02-w02-p02-036-inseparability]{\ref*{sec:p02-w02-p02-036-inseparability}} & n/a & \hyperref[def:p02-w02-p02-036-halting]{\ref*{def:p02-w02-p02-036-halting}} & \thepage & Source notation of BGG p. 34 (H\_1, H\_2 and the non-halting set); the rejecting state plays the source stop state 1, the accepting state stop state 2. Sets of natural numbers; distinct from the binary relation symbol H of the reduction formula. Explicit displayed family; member types and common actual first-use exposed, not one atomic glyph. Family listing, not one atomic glyph. Author macro \texttt{\detokenize{H\_1,}\allowbreak\detokenize{H\_2,}\allowbreak\detokenize{H\_\infty}} not carried to the members. \\
\phantomsection\label{entry:symbol-registry-06856}\(H_2\) & \texttt{\detokenize{H_2}} & codes of machines that reject / accept / never halt on the empty input symbol family member \(H_2\) & \hyperref[sec:p02-w02-p02-036-inseparability]{\ref*{sec:p02-w02-p02-036-inseparability}} & n/a & \hyperref[def:p02-w02-p02-036-halting]{\ref*{def:p02-w02-p02-036-halting}} & \thepage & Source notation of BGG p. 34 (H\_1, H\_2 and the non-halting set); the rejecting state plays the source stop state 1, the accepting state stop state 2. Sets of natural numbers; distinct from the binary relation symbol H of the reduction formula. Explicit displayed family; member types and common actual first-use exposed, not one atomic glyph. Family listing, not one atomic glyph. Author macro \texttt{\detokenize{H\_1,}\allowbreak\detokenize{H\_2,}\allowbreak\detokenize{H\_\infty}} not carried to the members. \\
\phantomsection\label{entry:symbol-registry-06857}\(H_\infty\) & \texttt{\detokenize{H_\infty}} & codes of machines that reject / accept / never halt on the empty input symbol family member \(H_\infty\) & \hyperref[sec:p02-w02-p02-036-inseparability]{\ref*{sec:p02-w02-p02-036-inseparability}} & n/a & \hyperref[def:p02-w02-p02-036-halting]{\ref*{def:p02-w02-p02-036-halting}} & \thepage & Source notation of BGG p. 34 (H\_1, H\_2 and the non-halting set); the rejecting state plays the source stop state 1, the accepting state stop state 2. Sets of natural numbers; distinct from the binary relation symbol H of the reduction formula. Explicit displayed family; member types and common actual first-use exposed, not one atomic glyph. Family listing, not one atomic glyph. Author macro \texttt{\detokenize{H\_1,}\allowbreak\detokenize{H\_2,}\allowbreak\detokenize{H\_\infty}} not carried to the members. \\
\phantomsection\label{entry:symbol-registry-06858}\(\Lang_M^{\prec}\) & \texttt{\detokenize{\Lang_M^{\prec}}} & signature of the refined reduction formula of a machine M & \hyperref[def:p02-w02-p02-036-refined]{\ref*{def:p02-w02-p02-036-refined}} & n/a & \hyperref[def:p02-w02-p02-036-refined]{\ref*{def:p02-w02-p02-036-refined}} & \thepage & Binary K, S, H, T\_a; unary Z, I\_q. \\
\phantomsection\label{entry:symbol-registry-06859}\(K\) & \texttt{\detokenize{K}} & relation symbols of the reduction formula: order, successor, head position, tape letter, state, zero symbol family member \(K\) & \hyperref[def:p02-w02-p02-036-refined]{\ref*{def:p02-w02-p02-036-refined}} & n/a & \hyperref[def:p02-w02-p02-036-refined]{\ref*{def:p02-w02-p02-036-refined}} & \thepage & Source notation of BGG pp. 20-22 (H, T\_j, I\_i, S, K) with Z for the source constant 0 (as in BGG Exercise 2.1.31). Object-language relation symbols; K here is unrelated to the self-application set K of sec:universal-reductions-03. Explicit displayed family; member types and common actual first-use exposed, not one atomic glyph. Family listing, not one atomic glyph. Author macro \texttt{\detokenize{K,}\allowbreak\detokenize{S,}\allowbreak\detokenize{H,}\allowbreak\detokenize{T\_a,}\allowbreak\detokenize{I\_q,}\allowbreak\detokenize{Z}} not carried to the members. \\
\phantomsection\label{entry:symbol-registry-06860}\(S\) & \texttt{\detokenize{S}} & relation symbols of the reduction formula: order, successor, head position, tape letter, state, zero symbol family member \(S\) & \hyperref[def:p02-w02-p02-036-refined]{\ref*{def:p02-w02-p02-036-refined}} & n/a & \hyperref[def:p02-w02-p02-036-refined]{\ref*{def:p02-w02-p02-036-refined}} & \thepage & Source notation of BGG pp. 20-22 (H, T\_j, I\_i, S, K) with Z for the source constant 0 (as in BGG Exercise 2.1.31). Object-language relation symbols; K here is unrelated to the self-application set K of sec:universal-reductions-03. Explicit displayed family; member types and common actual first-use exposed, not one atomic glyph. Family listing, not one atomic glyph. Author macro \texttt{\detokenize{K,}\allowbreak\detokenize{S,}\allowbreak\detokenize{H,}\allowbreak\detokenize{T\_a,}\allowbreak\detokenize{I\_q,}\allowbreak\detokenize{Z}} not carried to the members. \\
\phantomsection\label{entry:symbol-registry-06861}\(H\) & \texttt{\detokenize{H}} & relation symbols of the reduction formula: order, successor, head position, tape letter, state, zero symbol family member \(H\) & \hyperref[def:p02-w02-p02-036-refined]{\ref*{def:p02-w02-p02-036-refined}} & n/a & \hyperref[def:p02-w02-p02-036-refined]{\ref*{def:p02-w02-p02-036-refined}} & \thepage & Source notation of BGG pp. 20-22 (H, T\_j, I\_i, S, K) with Z for the source constant 0 (as in BGG Exercise 2.1.31). Object-language relation symbols; K here is unrelated to the self-application set K of sec:universal-reductions-03. Explicit displayed family; member types and common actual first-use exposed, not one atomic glyph. Family listing, not one atomic glyph. Author macro \texttt{\detokenize{K,}\allowbreak\detokenize{S,}\allowbreak\detokenize{H,}\allowbreak\detokenize{T\_a,}\allowbreak\detokenize{I\_q,}\allowbreak\detokenize{Z}} not carried to the members. \\
\phantomsection\label{entry:symbol-registry-06862}\(T_a\) & \texttt{\detokenize{T_a}} & relation symbols of the reduction formula: order, successor, head position, tape letter, state, zero symbol family member \(T_a\) & \hyperref[def:p02-w02-p02-036-refined]{\ref*{def:p02-w02-p02-036-refined}} & n/a & \hyperref[def:p02-w02-p02-036-refined]{\ref*{def:p02-w02-p02-036-refined}} & \thepage & Source notation of BGG pp. 20-22 (H, T\_j, I\_i, S, K) with Z for the source constant 0 (as in BGG Exercise 2.1.31). Object-language relation symbols; K here is unrelated to the self-application set K of sec:universal-reductions-03. Explicit displayed family; member types and common actual first-use exposed, not one atomic glyph. Family listing, not one atomic glyph. Author macro \texttt{\detokenize{K,}\allowbreak\detokenize{S,}\allowbreak\detokenize{H,}\allowbreak\detokenize{T\_a,}\allowbreak\detokenize{I\_q,}\allowbreak\detokenize{Z}} not carried to the members. \\
\phantomsection\label{entry:symbol-registry-06863}\(I_q\) & \texttt{\detokenize{I_q}} & relation symbols of the reduction formula: order, successor, head position, tape letter, state, zero symbol family member \(I_q\) & \hyperref[def:p02-w02-p02-036-refined]{\ref*{def:p02-w02-p02-036-refined}} & n/a & \hyperref[def:p02-w02-p02-036-refined]{\ref*{def:p02-w02-p02-036-refined}} & \thepage & Source notation of BGG pp. 20-22 (H, T\_j, I\_i, S, K) with Z for the source constant 0 (as in BGG Exercise 2.1.31). Object-language relation symbols; K here is unrelated to the self-application set K of sec:universal-reductions-03. Explicit displayed family; member types and common actual first-use exposed, not one atomic glyph. Family listing, not one atomic glyph. Author macro \texttt{\detokenize{K,}\allowbreak\detokenize{S,}\allowbreak\detokenize{H,}\allowbreak\detokenize{T\_a,}\allowbreak\detokenize{I\_q,}\allowbreak\detokenize{Z}} not carried to the members. \\
\phantomsection\label{entry:symbol-registry-06864}\(Z\) & \texttt{\detokenize{Z}} & relation symbols of the reduction formula: order, successor, head position, tape letter, state, zero symbol family member \(Z\) & \hyperref[def:p02-w02-p02-036-refined]{\ref*{def:p02-w02-p02-036-refined}} & n/a & \hyperref[def:p02-w02-p02-036-refined]{\ref*{def:p02-w02-p02-036-refined}} & \thepage & Source notation of BGG pp. 20-22 (H, T\_j, I\_i, S, K) with Z for the source constant 0 (as in BGG Exercise 2.1.31). Object-language relation symbols; K here is unrelated to the self-application set K of sec:universal-reductions-03. Explicit displayed family; member types and common actual first-use exposed, not one atomic glyph. Family listing, not one atomic glyph. Author macro \texttt{\detokenize{K,}\allowbreak\detokenize{S,}\allowbreak\detokenize{H,}\allowbreak\detokenize{T\_a,}\allowbreak\detokenize{I\_q,}\allowbreak\detokenize{Z}} not carried to the members. \\
\phantomsection\label{entry:symbol-registry-06865}\(\mathrm{Act}(x)\) & \texttt{\detokenize{\mathrm{Act}}} & abbreviation: the disjunction of I\_q(x) over non-halting states q & \hyperref[def:p02-w02-p02-036-refined]{\ref*{def:p02-w02-p02-036-refined}} & n/a & \hyperref[def:p02-w02-p02-036-refined]{\ref*{def:p02-w02-p02-036-refined}} & \thepage & Formula abbreviation, not a symbol of the signature. \\
\phantomsection\label{entry:symbol-registry-06866}\(\mathrm{ORDER}^{-}\) & \texttt{\detokenize{\mathrm{ORDER}^{-}}} & named conjuncts of the reduction formula symbol family member \(\mathrm{ORDER}^{-}\) & \hyperref[def:p02-w02-p02-036-refined]{\ref*{def:p02-w02-p02-036-refined}} & n/a & \hyperref[def:p02-w02-p02-036-refined]{\ref*{def:p02-w02-p02-036-refined}} & \thepage & Source names (BGG pp. 20-22); ORDER\^{}- is the source ORDER without its successor-existence clause. Explicit displayed family; member types and common actual first-use exposed, not one atomic glyph. Family listing, not one atomic glyph. Author macro \texttt{\detokenize{\mathrm{ORDER}\^{}{-},}\allowbreak\detokenize{\mathrm{START},}\allowbreak\detokenize{\mathrm{STEP}\_M,}\allowbreak\detokenize{\mathrm{NONSTOP}}} not carried to the members. \\
\phantomsection\label{entry:symbol-registry-06867}\(\mathrm{START}\) & \texttt{\detokenize{\mathrm{START}}} & named conjuncts of the reduction formula symbol family member \(\mathrm{START}\) & \hyperref[def:p02-w02-p02-036-refined]{\ref*{def:p02-w02-p02-036-refined}} & n/a & \hyperref[def:p02-w02-p02-036-refined]{\ref*{def:p02-w02-p02-036-refined}} & \thepage & Source names (BGG pp. 20-22); ORDER\^{}- is the source ORDER without its successor-existence clause. Explicit displayed family; member types and common actual first-use exposed, not one atomic glyph. Family listing, not one atomic glyph. Author macro \texttt{\detokenize{\mathrm{ORDER}\^{}{-},}\allowbreak\detokenize{\mathrm{START},}\allowbreak\detokenize{\mathrm{STEP}\_M,}\allowbreak\detokenize{\mathrm{NONSTOP}}} not carried to the members. \\
\phantomsection\label{entry:symbol-registry-06868}\(\mathrm{STEP}_M\) & \texttt{\detokenize{\mathrm{STEP}_M}} & named conjuncts of the reduction formula symbol family member \(\mathrm{STEP}_M\) & \hyperref[def:p02-w02-p02-036-refined]{\ref*{def:p02-w02-p02-036-refined}} & n/a & \hyperref[def:p02-w02-p02-036-refined]{\ref*{def:p02-w02-p02-036-refined}} & \thepage & Source names (BGG pp. 20-22); ORDER\^{}- is the source ORDER without its successor-existence clause. Explicit displayed family; member types and common actual first-use exposed, not one atomic glyph. Family listing, not one atomic glyph. Author macro \texttt{\detokenize{\mathrm{ORDER}\^{}{-},}\allowbreak\detokenize{\mathrm{START},}\allowbreak\detokenize{\mathrm{STEP}\_M,}\allowbreak\detokenize{\mathrm{NONSTOP}}} not carried to the members. \\
\phantomsection\label{entry:symbol-registry-06869}\(\mathrm{NONSTOP}\) & \texttt{\detokenize{\mathrm{NONSTOP}}} & named conjuncts of the reduction formula symbol family member \(\mathrm{NONSTOP}\) & \hyperref[def:p02-w02-p02-036-refined]{\ref*{def:p02-w02-p02-036-refined}} & n/a & \hyperref[def:p02-w02-p02-036-refined]{\ref*{def:p02-w02-p02-036-refined}} & \thepage & Source names (BGG pp. 20-22); ORDER\^{}- is the source ORDER without its successor-existence clause. Explicit displayed family; member types and common actual first-use exposed, not one atomic glyph. Family listing, not one atomic glyph. Author macro \texttt{\detokenize{\mathrm{ORDER}\^{}{-},}\allowbreak\detokenize{\mathrm{START},}\allowbreak\detokenize{\mathrm{STEP}\_M,}\allowbreak\detokenize{\mathrm{NONSTOP}}} not carried to the members. \\
\phantomsection\label{entry:symbol-registry-06870}\(\prec_{K,M}\) & \texttt{\detokenize{\prec_{K,M}}} & relativized infinity axiom of a machine M & \hyperref[def:p02-w02-p02-036-refined]{\ref*{def:p02-w02-p02-036-refined}} & n/a & \hyperref[def:p02-w02-p02-036-refined]{\ref*{def:p02-w02-p02-036-refined}} & \thepage & Name from BGG p. 35; the relational formula here asserts S-successors for active moments and visited cells and S-predecessors for visited cells. \\
\phantomsection\label{entry:symbol-registry-06871}\(\rho_{M,\prec}\) & \texttt{\detokenize{\rho_{M,\prec}}} & refined reduction formula of M & \hyperref[def:p02-w02-p02-036-refined]{\ref*{def:p02-w02-p02-036-refined}} & n/a & \hyperref[def:p02-w02-p02-036-refined]{\ref*{def:p02-w02-p02-036-refined}} & \thepage & Source notation rho\_{M,prec} (BGG p. 34); sentences. The global rho row is a substitution metavariable; this is a named sentence in a scoped definition. Blank macro filled with the glyph's TeX spelling. \\
\phantomsection\label{entry:symbol-registry-06872}\(\rho^{\infty}_M\) & \texttt{\detokenize{\rho^{\infty}_M}} & variant of the refined reduction formula of M with unrestricted successors & \hyperref[def:p02-w02-p02-036-refined]{\ref*{def:p02-w02-p02-036-refined}} & n/a & \hyperref[def:p02-w02-p02-036-refined]{\ref*{def:p02-w02-p02-036-refined}} & \thepage & Source notation rho\_{M,prec} (BGG p. 34); sentences. The global rho row is a substitution metavariable; this is a named sentence in a scoped definition. Blank macro filled with the glyph's TeX spelling. \\
\phantomsection\label{entry:symbol-registry-06873}\(c_t=(q_t,B_t,h_t)\) & \texttt{\detokenize{c_t=(q_t,}\allowbreak\detokenize{B_t,}\allowbreak\detokenize{h_t)}} & configuration at step t with tape inscription B\_t & \hyperref[thm:p02-w02-p02-036-simulation]{\ref*{thm:p02-w02-p02-036-simulation}} & n/a & \hyperref[thm:p02-w02-p02-036-simulation]{\ref*{thm:p02-w02-p02-036-simulation}} & \thepage & B\_t replaces the configuration letter T of sec:turing-machines-01 to avoid a clash with the relation symbols T\_a. Blank macro filled with the glyph's TeX spelling. \\
\phantomsection\label{entry:symbol-registry-06874}\(\tau_t\) & \texttt{\detokenize{\tau_t}} & time element of a model & \hyperref[thm:p02-w02-p02-036-simulation]{\ref*{thm:p02-w02-p02-036-simulation}} & n/a & \hyperref[thm:p02-w02-p02-036-simulation]{\ref*{thm:p02-w02-p02-036-simulation}} & \thepage & B\_t replaces the configuration letter T of sec:turing-machines-01 to avoid a clash with the relation symbols T\_a. Blank macro filled with the glyph's TeX spelling. \\
\phantomsection\label{entry:symbol-registry-06875}\(\sigma_x\) & \texttt{\detokenize{\sigma_x}} & cell element of a model & \hyperref[thm:p02-w02-p02-036-simulation]{\ref*{thm:p02-w02-p02-036-simulation}} & n/a & \hyperref[thm:p02-w02-p02-036-simulation]{\ref*{thm:p02-w02-p02-036-simulation}} & \thepage & B\_t replaces the configuration letter T of sec:turing-machines-01 to avoid a clash with the relation symbols T\_a. Blank macro filled with the glyph's TeX spelling. \\
\phantomsection\label{entry:symbol-registry-06876}\(\varphi^{\exists}\) & \texttt{\detokenize{^{\exists}}} & existential closure of a formula & \hyperref[thm:p02-w02-p02-036-fmp]{\ref*{thm:p02-w02-p02-036-fmp}} & n/a & \hyperref[thm:p02-w02-p02-036-fmp]{\ref*{thm:p02-w02-p02-036-fmp}} & \thepage & Postfix syntactic operation, local to the proof in the annex. \\
\phantomsection\label{entry:symbol-registry-06877}\(\operatorname{voc}(\varphi)\) & \texttt{\detokenize{\operatorname{voc}}} & occurrence vocabulary of a formula & \hyperref[sec:p02-w02-p02-037-craig]{\ref*{sec:p02-w02-p02-037-craig}} & n/a & \hyperref[def:p02-w02-p02-037-vocabulary]{\ref*{def:p02-w02-p02-037-vocabulary}} & \thepage & The function (constants included) and relation symbols occurring in the formula, with kind and arity; equality and bottom are logical and never members. \\
\phantomsection\label{entry:symbol-registry-06878}\(\varphi^{\mathrm{c}}\) & \texttt{\detokenize{\varphi^{\mathrm{c}}}} & constant instance of a formula symbol family member \(\varphi^{\mathrm{c}}\) & \hyperref[sec:p02-w02-p02-037-craig]{\ref*{sec:p02-w02-p02-037-craig}} & n/a & \hyperref[def:p02-w02-p02-037-reduction]{\ref*{def:p02-w02-p02-037-reduction}} & \thepage & The sentence obtained by substituting fresh pairwise distinct constants for the free variables of either formula; not a complement. Explicit displayed family; member types and common actual first-use exposed, not one atomic glyph. Family listing, not one atomic glyph. Author macro \texttt{\detokenize{{}\^{}{\mathrm{c}}}} not carried to the members. \\
\phantomsection\label{entry:symbol-registry-06879}\(\psi^{\mathrm{c}}\) & \texttt{\detokenize{\psi^{\mathrm{c}}}} & constant instance of a formula symbol family member \(\psi^{\mathrm{c}}\) & \hyperref[sec:p02-w02-p02-037-craig]{\ref*{sec:p02-w02-p02-037-craig}} & n/a & \hyperref[def:p02-w02-p02-037-reduction]{\ref*{def:p02-w02-p02-037-reduction}} & \thepage & The sentence obtained by substituting fresh pairwise distinct constants for the free variables of either formula; not a complement. Explicit displayed family; member types and common actual first-use exposed, not one atomic glyph. Family listing, not one atomic glyph. Author macro \texttt{\detokenize{{}\^{}{\mathrm{c}}}} not carried to the members. \\
\phantomsection\label{entry:symbol-registry-06880}\(\Lang^{\mathrm{c}}\) & \texttt{\detokenize{\Lang^{\mathrm{c}}}} & signature expanded by the fresh constants for free variables & \hyperref[sec:annex-p02-w02-p02-037-craig-beth]{\ref*{sec:annex-p02-w02-p02-037-craig-beth}} & n/a & \hyperref[annex:p02-w02-p02-037-craig-formula]{\ref*{annex:p02-w02-p02-037-craig-formula}} & \thepage & The union of the two compatible signatures together with the pairwise distinct fresh constants c\_1 to c\_k; distinct from the P2 tableau signature with a star. \\
\phantomsection\label{entry:symbol-registry-06881}\(\eta^{\circ}\) & \texttt{\detokenize{{}^{\circ}}} & bound-variable renaming of a formula away from a finite variable set & \hyperref[sec:annex-p02-w02-p02-037-craig-beth]{\ref*{sec:annex-p02-w02-p02-037-craig-beth}} & n/a & \hyperref[annex:p02-w02-p02-037-craig-formula]{\ref*{annex:p02-w02-p02-037-craig-formula}} & \thepage & Alpha-equivalent to the input with the same free variables and vocabulary; local to the annex proof. Distinct from the P8 da Costa consistency operator. \\
\phantomsection\label{entry:symbol-registry-06882}\(\eta\{v/c\}\) & \texttt{\detokenize{\{v/c\}}} & replacement of a constant by a variable & \hyperref[sec:annex-p02-w02-p02-037-craig-beth]{\ref*{sec:annex-p02-w02-p02-037-craig-beth}} & n/a & \hyperref[annex:p02-w02-p02-037-craig-formula]{\ref*{annex:p02-w02-p02-037-craig-formula}} & \thepage & Replaces every occurrence of the constant c by v and leaves every quantifier unchanged; used only when v does not occur in the formula. Distinct from capture-avoiding substitution [t/x]. \\
\phantomsection\label{entry:symbol-registry-06883}\(\mathcal N^{\mathrm{ren}}\) & \texttt{\detokenize{{}^{\mathrm{ren}}}} & renamed-predicate structure & \hyperref[sec:p02-w02-p02-037-beth]{\ref*{sec:p02-w02-p02-037-beth}} & n/a & \hyperref[thm:p02-w02-p02-037-renaming]{\ref*{thm:p02-w02-p02-037-renaming}} & \thepage & The L(P)-structure with the domain and L-interpretations of an L(P,P')-structure N, interpreting P as P' is interpreted in N. \\
\phantomsection\label{entry:symbol-registry-06884}\(\alpha'\) & \texttt{\detokenize{{}'}} & uniform replacement of P by P' in a formula & \hyperref[sec:p02-w02-p02-037-beth]{\ref*{sec:p02-w02-p02-037-beth}} & n/a & \hyperref[thm:p02-w02-p02-037-renaming]{\ref*{thm:p02-w02-p02-037-renaming}} & \thepage & Same reading as T(P') in def:p02-beth-definability, applied to single formulas. \\
\phantomsection\label{entry:symbol-registry-06885}\(A\supset B\) & \texttt{\detokenize{\supset}} & Kleene's object-language implication in his predicate calculus & \hyperref[sec:p02-w02-p02-025-kleene]{\ref*{sec:p02-w02-p02-025-kleene}} & n/a & \hyperref[def:p02-w02-p02-025-kleene-postulates]{\ref*{def:p02-w02-p02-025-kleene-postulates}} & \thepage & Exact source notation of Kleene 1952, Section 19, p. 82. Same glyph as the registered P8 Fitting conditional and the P17 experiment comparison; those scopes are disjoint and the roles differ. Within this scope it plays the role that \(\to\) plays on sheet:hilbert-fol. Precedence (source descriptor): binds less tightly than \(\mathbin{\&}\) and \(\lor\). \\
\phantomsection\label{entry:symbol-registry-06886}\(A\mathbin{\&}B\) & \texttt{\detokenize{\mathbin{\&}}} & Kleene's object-language conjunction & \hyperref[sec:p02-w02-p02-025-kleene]{\ref*{sec:p02-w02-p02-025-kleene}} & n/a & \hyperref[def:p02-w02-p02-025-kleene-postulates]{\ref*{def:p02-w02-p02-025-kleene-postulates}} & \thepage & Exact source notation of Kleene 1952, Section 19, p. 82; primitive in Kleene's system, with postulates 3, 4a and 4b. Precedence (source descriptor): binds more tightly than \(\supset\). \\
\phantomsection\label{entry:symbol-registry-06887}\(A(x)\) & \texttt{\detokenize{A(x)}} & Kleene's notation for a formula and the result of substituting a term free for x symbol family member \(A(x)\) & \hyperref[sec:p02-w02-p02-025-kleene]{\ref*{sec:p02-w02-p02-025-kleene}} & n/a & \hyperref[def:p02-w02-p02-025-kleene-postulates]{\ref*{def:p02-w02-p02-025-kleene-postulates}} & \thepage & Source notation of Kleene 1952, Section 19, pp. 81-82, with the stipulation that t is free for x in A(x); not the handbook's total capture-avoiding operation \(\varphi[t/x]\). Explicit displayed family; member types and common actual first-use exposed, not one atomic glyph. Family listing, not one atomic glyph. Author macro \texttt{\detokenize{A(x),}\allowbreak\detokenize{A(t)}} not carried to the members. \\
\phantomsection\label{entry:symbol-registry-06888}\(\ A(t)\) & \texttt{\detokenize{\ A(t)}} & Kleene's notation for a formula and the result of substituting a term free for x symbol family member \(\ A(t)\) & \hyperref[sec:p02-w02-p02-025-kleene]{\ref*{sec:p02-w02-p02-025-kleene}} & n/a & \hyperref[def:p02-w02-p02-025-kleene-postulates]{\ref*{def:p02-w02-p02-025-kleene-postulates}} & \thepage & Source notation of Kleene 1952, Section 19, pp. 81-82, with the stipulation that t is free for x in A(x); not the handbook's total capture-avoiding operation \(\varphi[t/x]\). Explicit displayed family; member types and common actual first-use exposed, not one atomic glyph. Family listing, not one atomic glyph. Author macro \texttt{\detokenize{A(x),}\allowbreak\detokenize{A(t)}} not carried to the members. \\
\phantomsection\label{entry:symbol-registry-06889}\(D_1,\ldots,D_l\derives E\) & \texttt{\detokenize{\derives}} & Kleene's deducibility from a list of assumption formula occurrences & \hyperref[sec:p02-w02-p02-025-kleene]{\ref*{sec:p02-w02-p02-025-kleene}} & n/a & \hyperref[def:p02-w02-p02-025-kleene-dependence]{\ref*{def:p02-w02-p02-025-kleene-dependence}} & \thepage & Kleene 1952, Section 20, p. 87. Same glyph and the same typed role (derivability turnstile) as the global \(\derives\); the premise is a finite list of occurrences, not a set, and the relation belongs to Kleene's system, not sheet:hilbert-fol. Sequence glyph kept as one row; the ellipsis is not a member symbol. \\
\phantomsection\label{entry:symbol-registry-06890}\(\Gamma\) & \texttt{\detokenize{\Gamma}} & finite list of assumption occurrences with A appended last symbol family member \(\Gamma\) & \hyperref[sec:p02-w02-p02-025-kleene]{\ref*{sec:p02-w02-p02-025-kleene}} & n/a & \hyperref[def:p02-w02-p02-025-kleene-dependence]{\ref*{def:p02-w02-p02-025-kleene-dependence}} & \thepage & Kleene's list convention; the deduction theorem discharges the last occurrence. Explicit displayed family; member types and common actual first-use exposed, not one atomic glyph. Family listing, not one atomic glyph. Author macro \texttt{\detokenize{\Gamma,}\allowbreak\detokenize{A}} not carried to the members. \\
\phantomsection\label{entry:symbol-registry-06891}\(A\) & \texttt{\detokenize{A}} & finite list of assumption occurrences with A appended last symbol family member \(A\) & \hyperref[sec:p02-w02-p02-025-kleene]{\ref*{sec:p02-w02-p02-025-kleene}} & n/a & \hyperref[def:p02-w02-p02-025-kleene-dependence]{\ref*{def:p02-w02-p02-025-kleene-dependence}} & \thepage & Kleene's list convention; the deduction theorem discharges the last occurrence. Explicit displayed family; member types and common actual first-use exposed, not one atomic glyph. Family listing, not one atomic glyph. Author macro \texttt{\detokenize{\Gamma,}\allowbreak\detokenize{A}} not carried to the members. \\
\phantomsection\label{entry:symbol-registry-06892}\(\mathcal D_1\mathcal D_2\) & \texttt{\detokenize{\mathcal D_i}} & concatenation of dependency-annotated Hilbert derivations & \hyperref[sec:p02-w02-p02-025-kernels]{\ref*{sec:p02-w02-p02-025-kernels}} & n/a & \hyperref[thm:p02-w02-p02-025-support]{\ref*{thm:p02-w02-p02-025-support}} & \thepage & Juxtaposition with line references of the second derivation shifted; lines keep their supports. \\
\phantomsection\label{entry:symbol-registry-06893}\(E_i\) & \texttt{\detokenize{E_i}} & support of the transformed line in the deduction-theorem construction & \hyperref[sec:p02-w02-p02-025-kernels]{\ref*{sec:p02-w02-p02-025-kernels}} & n/a & \hyperref[thm:p02-w02-p02-025-discharge]{\ref*{thm:p02-w02-p02-025-discharge}} & \thepage & Finite subset of \(D_i\setminus\{\varphi\}\); the support notion is that of def:p02-hilbert-fol-dependencies. \\
\phantomsection\label{entry:symbol-registry-06894}\(\operatorname{ext}_{\mathcal M}(a)\) & \texttt{\detokenize{\operatorname{ext}_{\mathcal M}}} & external extension of an element of an L\_in-structure & \hyperref[sec:p02-w02-p02-033-skolem]{\ref*{sec:p02-w02-p02-033-skolem}} & n/a & \hyperref[def:p02-w02-p02-033-internal-countability]{\ref*{def:p02-w02-p02-033-internal-countability}} & \thepage & The ambient set of b in the domain with M satisfying b in a; a subset of the domain. Shares the operator name ext with the P2 extensional-collapse map and the P10 semilattice extension, in different scopes and roles. \\
\phantomsection\label{entry:symbol-registry-06895}\(\in^{\mathcal M}\) & \texttt{\detokenize{\in^{\mathcal M}}} & interpretation of the membership symbol in a structure & \hyperref[sec:p02-w02-p02-033-skolem]{\ref*{sec:p02-w02-p02-033-skolem}} & meta & \hyperref[def:p02-w02-p02-033-internal-countability]{\ref*{def:p02-w02-p02-033-internal-countability}} & \thepage & Instance of the general interpretation superscript R\^{}{M}; need not be the ambient membership relation. \\
\phantomsection\label{entry:symbol-registry-06896}\(\mathrm{Emp}(e)\) & \texttt{\detokenize{\mathrm{Emp}}} & named L\_in-formula abbreviation: e is empty & \hyperref[sec:p02-w02-p02-033-skolem]{\ref*{sec:p02-w02-p02-033-skolem}} & n/a & \hyperref[def:p02-w02-p02-033-internal-countability]{\ref*{def:p02-w02-p02-033-internal-countability}} & \thepage & Object-language abbreviations eliminating the empty set, S(x), omega and power set of sec:zf-zfc-01 into membership formulas. \\
\phantomsection\label{entry:symbol-registry-06897}\(\mathrm{Suc}(y,x)\) & \texttt{\detokenize{\mathrm{Suc}}} & named L\_in-formula abbreviation: y is the successor of x & \hyperref[sec:p02-w02-p02-033-skolem]{\ref*{sec:p02-w02-p02-033-skolem}} & n/a & \hyperref[def:p02-w02-p02-033-internal-countability]{\ref*{def:p02-w02-p02-033-internal-countability}} & \thepage & Object-language abbreviations eliminating the empty set, S(x), omega and power set of sec:zf-zfc-01 into membership formulas. \\
\phantomsection\label{entry:symbol-registry-06898}\(\mathrm{Ind}(v)\) & \texttt{\detokenize{\mathrm{Ind}}} & named L\_in-formula abbreviation: v is inductive & \hyperref[sec:p02-w02-p02-033-skolem]{\ref*{sec:p02-w02-p02-033-skolem}} & n/a & \hyperref[def:p02-w02-p02-033-internal-countability]{\ref*{def:p02-w02-p02-033-internal-countability}} & \thepage & Object-language abbreviations eliminating the empty set, S(x), omega and power set of sec:zf-zfc-01 into membership formulas. \\
\phantomsection\label{entry:symbol-registry-06899}\(\mathrm{Om}(w)\) & \texttt{\detokenize{\mathrm{Om}}} & named L\_in-formula abbreviation: w is the least inductive set & \hyperref[sec:p02-w02-p02-033-skolem]{\ref*{sec:p02-w02-p02-033-skolem}} & n/a & \hyperref[def:p02-w02-p02-033-internal-countability]{\ref*{def:p02-w02-p02-033-internal-countability}} & \thepage & Object-language abbreviations eliminating the empty set, S(x), omega and power set of sec:zf-zfc-01 into membership formulas. \\
\phantomsection\label{entry:symbol-registry-06900}\(\mathrm{Pw}(p,w)\) & \texttt{\detokenize{\mathrm{Pw}}} & named L\_in-formula abbreviation: p is the power set of w & \hyperref[sec:p02-w02-p02-033-skolem]{\ref*{sec:p02-w02-p02-033-skolem}} & n/a & \hyperref[def:p02-w02-p02-033-internal-countability]{\ref*{def:p02-w02-p02-033-internal-countability}} & \thepage & Object-language abbreviations eliminating the empty set, S(x), omega and power set of sec:zf-zfc-01 into membership formulas. \\
\phantomsection\label{entry:symbol-registry-06901}\(\mathrm{OP}(z,u,v)\) & \texttt{\detokenize{\mathrm{OP}}} & named L\_in-formula abbreviation for a Kuratowski pair & \hyperref[sec:p02-w02-p02-033-skolem]{\ref*{sec:p02-w02-p02-033-skolem}} & n/a & \hyperref[def:p02-w02-p02-033-internal-countability]{\ref*{def:p02-w02-p02-033-internal-countability}} & \thepage & OP(z,u,v) says z is the Kuratowski pair of u and v; In(u,v,g) says that pair belongs to g. \\
\phantomsection\label{entry:symbol-registry-06902}\(\mathrm{In}(u,v,g)\) & \texttt{\detokenize{\mathrm{In}}} & named L\_in-formula abbreviation for pair membership & \hyperref[sec:p02-w02-p02-033-skolem]{\ref*{sec:p02-w02-p02-033-skolem}} & n/a & \hyperref[def:p02-w02-p02-033-internal-countability]{\ref*{def:p02-w02-p02-033-internal-countability}} & \thepage & OP(z,u,v) says z is the Kuratowski pair of u and v; In(u,v,g) says that pair belongs to g. \\
\phantomsection\label{entry:symbol-registry-06903}\(\mathrm{Inj}(g,a,w)\) & \texttt{\detokenize{\mathrm{Inj}}} & named L\_in-formula: g injects a into w & \hyperref[sec:p02-w02-p02-033-skolem]{\ref*{sec:p02-w02-p02-033-skolem}} & n/a & \hyperref[def:p02-w02-p02-033-internal-countability]{\ref*{def:p02-w02-p02-033-internal-countability}} & \thepage & Ctbl(a) is the internal reading of the countability definition of sec:set-prerequisites-05 (injection into omega). \\
\phantomsection\label{entry:symbol-registry-06904}\(\mathrm{Ctbl}(a)\) & \texttt{\detokenize{\mathrm{Ctbl}}} & named L\_in-formula: a is internally countable & \hyperref[sec:p02-w02-p02-033-skolem]{\ref*{sec:p02-w02-p02-033-skolem}} & n/a & \hyperref[def:p02-w02-p02-033-internal-countability]{\ref*{def:p02-w02-p02-033-internal-countability}} & \thepage & Ctbl(a) is the internal reading of the countability definition of sec:set-prerequisites-05 (injection into omega). \\
\phantomsection\label{entry:symbol-registry-06905}\(\sigma_{\mathrm{unc}}\) & \texttt{\detokenize{\sigma_{\mathrm{unc}}}} & L\_in-sentence asserting that the power set of omega is not injectable into omega & \hyperref[sec:p02-w02-p02-033-skolem]{\ref*{sec:p02-w02-p02-033-skolem}} & n/a & \hyperref[def:p02-w02-p02-033-internal-countability]{\ref*{def:p02-w02-p02-033-internal-countability}} & \thepage & Proved from extensionality; power set; separation and infinity in thm:p02-w02-p02-033-cantor-internal. \\
\phantomsection\label{entry:symbol-registry-06906}\(\omega^{\mathcal M}\) & \texttt{\detokenize{\omega^{\mathcal M}}} & the unique element of a model satisfying Om over it & \hyperref[sec:p02-w02-p02-033-skolem]{\ref*{sec:p02-w02-p02-033-skolem}} & n/a & \hyperref[thm:p02-w02-p02-033-cantor-internal]{\ref*{thm:p02-w02-p02-033-cantor-internal}} & \thepage & Uniqueness by extensionality inside the model. \\
\phantomsection\label{entry:symbol-registry-06907}\(\mathcal P(\omega)^{\mathcal M}\) & \texttt{\detokenize{\mathcal P(\omega)^{\mathcal M}}} & the unique element of a model satisfying Pw over \(\omega^{\mathcal M}\) & \hyperref[sec:p02-w02-p02-033-skolem]{\ref*{sec:p02-w02-p02-033-skolem}} & n/a & \hyperref[thm:p02-w02-p02-033-cantor-internal]{\ref*{thm:p02-w02-p02-033-cantor-internal}} & \thepage & Uniqueness by extensionality inside the model. \\
\phantomsection\label{entry:symbol-registry-06908}\(R_g\) & \texttt{\detokenize{R_g}} & external relation coded by an element g of an L\_in-structure & \hyperref[thm:p02-w02-p02-033-skolem-resolution]{\ref*{thm:p02-w02-p02-033-skolem-resolution}} & n/a & \hyperref[thm:p02-w02-p02-033-skolem-resolution]{\ref*{thm:p02-w02-p02-033-skolem-resolution}} & \thepage & Set of pairs (b,c) with M satisfying In(b,c,g); R is also a generic relation-symbol metavariable elsewhere. \\
\phantomsection\label{entry:symbol-registry-06909}\(\operatorname{code}(s)\) & \texttt{\detokenize{\operatorname{code}}} & binary code of a finite 0-1 sequence & \hyperref[thm:p02-w02-p02-033-hypotheses-needed]{\ref*{thm:p02-w02-p02-033-hypotheses-needed}} & n/a & \hyperref[thm:p02-w02-p02-033-hypotheses-needed]{\ref*{thm:p02-w02-p02-033-hypotheses-needed}} & \thepage & 2 to the length plus the binary value; same operator name as other code maps registered in other scopes. \\
\phantomsection\label{entry:symbol-registry-06910}\(D\) & \texttt{\detokenize{D}} & set of new elements of the target outside the image of an elementary embedding & \hyperref[thm:p02-w02-p02-033-copy-over-image]{\ref*{thm:p02-w02-p02-033-copy-over-image}} & n/a & \hyperref[thm:p02-w02-p02-033-copy-over-image]{\ref*{thm:p02-w02-p02-033-copy-over-image}} & \thepage & Local names in the isomorphic-copy construction. \\
\phantomsection\label{entry:symbol-registry-06911}\(\pi\) & \texttt{\detokenize{\pi}} & transport isomorphism of the isomorphic copy & \hyperref[thm:p02-w02-p02-033-copy-over-image]{\ref*{thm:p02-w02-p02-033-copy-over-image}} & n/a & \hyperref[thm:p02-w02-p02-033-copy-over-image]{\ref*{thm:p02-w02-p02-033-copy-over-image}} & \thepage & Local names in the isomorphic-copy construction. \\
\phantomsection\label{entry:symbol-registry-06912}\(D'\) & \texttt{\detokenize{D'}} & fresh copy of the new-element set & \hyperref[annex:p02-w02-p02-033-copy]{\ref*{annex:p02-w02-p02-033-copy}} & n/a & \hyperref[annex:p02-w02-p02-033-copy]{\ref*{annex:p02-w02-p02-033-copy}} & \thepage & Local names in the isomorphic-copy construction. \\
\phantomsection\label{entry:symbol-registry-06913}\(r\) & \texttt{\detokenize{r}} & fresh index making the copy disjoint from the old domain & \hyperref[annex:p02-w02-p02-033-copy]{\ref*{annex:p02-w02-p02-033-copy}} & n/a & \hyperref[annex:p02-w02-p02-033-copy]{\ref*{annex:p02-w02-p02-033-copy}} & \thepage & Local names in the isomorphic-copy construction. \\
\phantomsection\label{entry:symbol-registry-06914}\(\operatorname{num}_{\mathcal M}\) & \texttt{\detokenize{\operatorname{num}_{\mathcal M}}} & numeral map of an L\_A-structure & \hyperref[sec:p02-w02-p02-033-nonstandard]{\ref*{sec:p02-w02-p02-033-nonstandard}} & n/a & \hyperref[def:p02-w02-p02-033-numeral-map]{\ref*{def:p02-w02-p02-033-numeral-map}} & \thepage & Sends n to the value of the numeral S\^{}n0 in M; its range is the set of standard elements. \\
\phantomsection\label{entry:symbol-registry-06915}\(<^{\mathcal M}\) & \texttt{\detokenize{<^{\mathcal M}}} & relation on an L\_A-structure defined by the abbreviation r<t & \hyperref[sec:p02-w02-p02-033-nonstandard]{\ref*{sec:p02-w02-p02-033-nonstandard}} & meta & \hyperref[def:p02-w02-p02-033-numeral-map]{\ref*{def:p02-w02-p02-033-numeral-map}} & \thepage & Pairs (b,b') with M satisfying exists z (b+Sz=b'); not a signature symbol of L\_A. \\
\phantomsection\label{entry:symbol-registry-06916}\(\theta_n\) & \texttt{\detokenize{\theta_n}} & L\_A-sentences describing the elements below and comparable with a numeral symbol family member \(\theta_n\) & \hyperref[def:p02-w02-p02-033-numeral-map]{\ref*{def:p02-w02-p02-033-numeral-map}} & n/a & \hyperref[def:p02-w02-p02-033-numeral-map]{\ref*{def:p02-w02-p02-033-numeral-map}} & \thepage & Local sentence names; theta\_n and tau\_n are registered with unrelated roles in P15 and P6 scopes. Explicit displayed family; member types and common actual first-use exposed, not one atomic glyph. Family listing, not one atomic glyph. Author macro \texttt{\detokenize{\theta\_n,}\allowbreak\detokenize{\tau\_n}} not carried to the members. \\
\phantomsection\label{entry:symbol-registry-06917}\(\tau_n\) & \texttt{\detokenize{\tau_n}} & L\_A-sentences describing the elements below and comparable with a numeral symbol family member \(\tau_n\) & \hyperref[def:p02-w02-p02-033-numeral-map]{\ref*{def:p02-w02-p02-033-numeral-map}} & n/a & \hyperref[def:p02-w02-p02-033-numeral-map]{\ref*{def:p02-w02-p02-033-numeral-map}} & \thepage & Local sentence names; theta\_n and tau\_n are registered with unrelated roles in P15 and P6 scopes. Explicit displayed family; member types and common actual first-use exposed, not one atomic glyph. Family listing, not one atomic glyph. Author macro \texttt{\detokenize{\theta\_n,}\allowbreak\detokenize{\tau\_n}} not carried to the members. \\
\phantomsection\label{entry:symbol-registry-06918}\(T_1\) & \texttt{\detokenize{T_1}} & compactness theories of the nonstandard-model examples symbol family member \(T_1\) & \hyperref[sec:p02-w02-p02-033-nonstandard]{\ref*{sec:p02-w02-p02-033-nonstandard}} & n/a & \hyperref[thm:p02-w02-p02-033-nonstandard-direct]{\ref*{thm:p02-w02-p02-033-nonstandard-direct}} & \thepage & T\_1: true arithmetic with c different from every numeral; T\_2: with c nonzero and divisible by every prime; T\_3: elementary diagram of the real ordered field with c above every real. Explicit displayed family; member types and common actual first-use exposed, not one atomic glyph. Family listing, not one atomic glyph. Author macro \texttt{\detokenize{T\_1,}\allowbreak\detokenize{T\_2,}\allowbreak\detokenize{T\_3}} not carried to the members. \\
\phantomsection\label{entry:symbol-registry-06919}\(T_2\) & \texttt{\detokenize{T_2}} & compactness theories of the nonstandard-model examples symbol family member \(T_2\) & \hyperref[sec:p02-w02-p02-033-nonstandard]{\ref*{sec:p02-w02-p02-033-nonstandard}} & n/a & \hyperref[thm:p02-w02-p02-033-nonstandard-direct]{\ref*{thm:p02-w02-p02-033-nonstandard-direct}} & \thepage & T\_1: true arithmetic with c different from every numeral; T\_2: with c nonzero and divisible by every prime; T\_3: elementary diagram of the real ordered field with c above every real. Explicit displayed family; member types and common actual first-use exposed, not one atomic glyph. Family listing, not one atomic glyph. Author macro \texttt{\detokenize{T\_1,}\allowbreak\detokenize{T\_2,}\allowbreak\detokenize{T\_3}} not carried to the members. \\
\phantomsection\label{entry:symbol-registry-06920}\(T_3\) & \texttt{\detokenize{T_3}} & compactness theories of the nonstandard-model examples symbol family member \(T_3\) & \hyperref[sec:p02-w02-p02-033-nonstandard]{\ref*{sec:p02-w02-p02-033-nonstandard}} & n/a & \hyperref[thm:p02-w02-p02-033-nonstandard-direct]{\ref*{thm:p02-w02-p02-033-nonstandard-direct}} & \thepage & T\_1: true arithmetic with c different from every numeral; T\_2: with c nonzero and divisible by every prime; T\_3: elementary diagram of the real ordered field with c above every real. Explicit displayed family; member types and common actual first-use exposed, not one atomic glyph. Family listing, not one atomic glyph. Author macro \texttt{\detokenize{T\_1,}\allowbreak\detokenize{T\_2,}\allowbreak\detokenize{T\_3}} not carried to the members. \\
\phantomsection\label{entry:symbol-registry-06921}\(m_\Delta\) & \texttt{\detokenize{m_\Delta}} & finite-satisfiability witness for a finite fragment Delta symbol family member \(m_\Delta\) & \hyperref[sec:p02-w02-p02-033-nonstandard]{\ref*{sec:p02-w02-p02-033-nonstandard}} & n/a & \hyperref[thm:p02-w02-p02-033-nonstandard-direct]{\ref*{thm:p02-w02-p02-033-nonstandard-direct}} & \thepage & The interpretation of the new constant c chosen for Delta; depends on Delta. Explicit displayed family; member types and common actual first-use exposed, not one atomic glyph. Family listing, not one atomic glyph. Author macro \texttt{\detokenize{m\_\Delta,}\allowbreak\detokenize{r\_\Delta}} not carried to the members. \\
\phantomsection\label{entry:symbol-registry-06922}\(r_\Delta\) & \texttt{\detokenize{r_\Delta}} & finite-satisfiability witness for a finite fragment Delta symbol family member \(r_\Delta\) & \hyperref[sec:p02-w02-p02-033-nonstandard]{\ref*{sec:p02-w02-p02-033-nonstandard}} & n/a & \hyperref[thm:p02-w02-p02-033-nonstandard-direct]{\ref*{thm:p02-w02-p02-033-nonstandard-direct}} & \thepage & The interpretation of the new constant c chosen for Delta; depends on Delta. Explicit displayed family; member types and common actual first-use exposed, not one atomic glyph. Family listing, not one atomic glyph. Author macro \texttt{\detokenize{m\_\Delta,}\allowbreak\detokenize{r\_\Delta}} not carried to the members. \\
\phantomsection\label{entry:symbol-registry-06923}\(t\mid u\) & \texttt{\detokenize{\mid}} & divisibility abbreviation in L\_A & \hyperref[sec:p02-w02-p02-033-nonstandard]{\ref*{sec:p02-w02-p02-033-nonstandard}} & meta & \hyperref[thm:p02-w02-p02-033-divisible]{\ref*{thm:p02-w02-p02-033-divisible}} & \thepage & Abbreviates exists z (t times z = u) with z the first variable not in t or u; the glyph is registered for conditional probability and dyadic obligation in other scopes. \\
\phantomsection\label{entry:symbol-registry-06924}\(\Lang_{\mathrm{of}}\) & \texttt{\detokenize{\Lang_{\mathrm{of}}}} & ordered-field signature and the real ordered field symbol family member \(\Lang_{\mathrm{of}}\) & \hyperref[sec:p02-w02-p02-033-nonstandard]{\ref*{sec:p02-w02-p02-033-nonstandard}} & n/a & \hyperref[thm:p02-w02-p02-033-infinitesimal]{\ref*{thm:p02-w02-p02-033-infinitesimal}} & \thepage & Signature {0,1,+,cdot,<}; structure (R,0,1,+,cdot,<). The calligraphic R also names the relation-symbol family of a signature globally; the subscript separates them. Explicit displayed family; member types and common actual first-use exposed, not one atomic glyph. Family listing, not one atomic glyph. Author macro \texttt{\detokenize{\Lang\_{\mathrm{of}},}\allowbreak\detokenize{\mathcal R\_{\mathrm{of}}}} not carried to the members. \\
\phantomsection\label{entry:symbol-registry-06925}\(\mathcal R_{\mathrm{of}}\) & \texttt{\detokenize{\mathcal R_{\mathrm{of}}}} & ordered-field signature and the real ordered field symbol family member \(\mathcal R_{\mathrm{of}}\) & \hyperref[sec:p02-w02-p02-033-nonstandard]{\ref*{sec:p02-w02-p02-033-nonstandard}} & n/a & \hyperref[thm:p02-w02-p02-033-infinitesimal]{\ref*{thm:p02-w02-p02-033-infinitesimal}} & \thepage & Signature {0,1,+,cdot,<}; structure (R,0,1,+,cdot,<). The calligraphic R also names the relation-symbol family of a signature globally; the subscript separates them. Explicit displayed family; member types and common actual first-use exposed, not one atomic glyph. Family listing, not one atomic glyph. Author macro \texttt{\detokenize{\Lang\_{\mathrm{of}},}\allowbreak\detokenize{\mathcal R\_{\mathrm{of}}}} not carried to the members. \\
\phantomsection\label{entry:symbol-registry-06926}\(\underline n\) & \texttt{\detokenize{\underline}} & closed L\_of-term for a natural number & \hyperref[sec:p02-w02-p02-033-nonstandard]{\ref*{sec:p02-w02-p02-033-nonstandard}} & n/a & \hyperref[thm:p02-w02-p02-033-infinitesimal]{\ref*{thm:p02-w02-p02-033-infinitesimal}} & \thepage & underline 0 = 0 and underline(n+1) = (underline n + 1); underline is registered with other roles in P10 and P12 scopes. \\
\phantomsection\label{entry:symbol-registry-06927}\(\alpha_1,\ldots,\alpha_5,\beta_1,\beta_2,\gamma_1,\gamma_2,\delta_a,\iota_\varphi,\sigma_{\mathcal F}\) & \texttt{\detokenize{\alpha_i,}\allowbreak\detokenize{\beta_i,}\allowbreak\detokenize{\gamma_i,}\allowbreak\detokenize{\delta_a,}\allowbreak\detokenize{\iota_\varphi,}\allowbreak\detokenize{\sigma_{\mathcal F}}} & local sentence names in proofs & \hyperref[annex:p02-w02-p02-033-numeral]{\ref*{annex:p02-w02-p02-033-numeral}} & n/a & \hyperref[annex:p02-w02-p02-033-numeral]{\ref*{annex:p02-w02-p02-033-numeral}} & \thepage & Used only inside the annex proofs of the nonstandard-model examples (numeral, infinitesimal and fails items). Sequence glyph kept as one row; the ellipsis is not a member symbol. \\
\phantomsection\label{entry:symbol-registry-06928}\(\llbracket\psi\rrbracket_v\) & \texttt{\detokenize{\llbracket\cdot\rrbracket_v}} & propositional value of a ground quantifier-free formula under an atom valuation & \hyperref[sec:p02-w02-p02-028-herbrand-structures]{\ref*{sec:p02-w02-p02-028-herbrand-structures}} & n/a & \hyperref[def:p02-w02-p02-028-atom-valuation]{\ref*{def:p02-w02-p02-028-atom-valuation}} & \thepage & Each distinct ground atom is one propositional variable; same bracket glyph as the registered semantic brackets, here indexed by an atom valuation. \\
\phantomsection\label{entry:symbol-registry-06929}\(v_{\mathcal N}\) & \texttt{\detokenize{v_{\mathcal N}}} & atom valuation induced by a structure or by a set of ground atoms symbol family member \(v_{\mathcal N}\) & \hyperref[sec:p02-w02-p02-028-herbrand-structures]{\ref*{sec:p02-w02-p02-028-herbrand-structures}} & n/a & \hyperref[def:p02-w02-p02-028-atom-valuation]{\ref*{def:p02-w02-p02-028-atom-valuation}} & \thepage & v\_N(A)=1 iff N satisfies A; v\_B is the characteristic function of B. Explicit displayed family; member types and common actual first-use exposed, not one atomic glyph. Family listing, not one atomic glyph. Author macro \texttt{\detokenize{v\_{\mathcal N},}\allowbreak\detokenize{v\_B}} not carried to the members. \\
\phantomsection\label{entry:symbol-registry-06930}\(v_B\) & \texttt{\detokenize{v_B}} & atom valuation induced by a structure or by a set of ground atoms symbol family member \(v_B\) & \hyperref[sec:p02-w02-p02-028-herbrand-structures]{\ref*{sec:p02-w02-p02-028-herbrand-structures}} & n/a & \hyperref[def:p02-w02-p02-028-atom-valuation]{\ref*{def:p02-w02-p02-028-atom-valuation}} & \thepage & v\_N(A)=1 iff N satisfies A; v\_B is the characteristic function of B. Explicit displayed family; member types and common actual first-use exposed, not one atomic glyph. Family listing, not one atomic glyph. Author macro \texttt{\detokenize{v\_{\mathcal N},}\allowbreak\detokenize{v\_B}} not carried to the members. \\
\phantomsection\label{entry:symbol-registry-06931}\(\operatorname{Hf}(\sigma)\) & \texttt{\detokenize{\operatorname{Hf}}} & Herbrand (dual-Skolem) form of a rectified prenex sentence & \hyperref[sec:p02-w02-p02-028-validity-forms]{\ref*{sec:p02-w02-p02-028-validity-forms}} & n/a & \hyperref[def:p02-w02-p02-028-herbrand-form]{\ref*{def:p02-w02-p02-028-herbrand-form}} & \thepage & Existential sentence in the expanded signature; universals replaced by new function terms of the preceding existential variables. \\
\phantomsection\label{entry:symbol-registry-06932}\(\theta^{\mathrm H}\) & \texttt{\detokenize{\theta^{\mathrm H}}} & Herbrand-form matrix & \hyperref[sec:p02-w02-p02-028-validity-forms]{\ref*{sec:p02-w02-p02-028-validity-forms}} & n/a & \hyperref[def:p02-w02-p02-028-herbrand-form]{\ref*{def:p02-w02-p02-028-herbrand-form}} & \thepage & theta here is the quantifier-free matrix of a prenex sentence, following the inherited Skolemization definition; it is not a substitution in this scope. Blank macro filled with the glyph's TeX spelling. \\
\phantomsection\label{entry:symbol-registry-06933}\(\Lang^{\mathrm H}\) & \texttt{\detokenize{\Lang^{\mathrm H}}} & expanded signature of the Herbrand form & \hyperref[sec:p02-w02-p02-028-validity-forms]{\ref*{sec:p02-w02-p02-028-validity-forms}} & n/a & \hyperref[def:p02-w02-p02-028-herbrand-form]{\ref*{def:p02-w02-p02-028-herbrand-form}} & \thepage & theta here is the quantifier-free matrix of a prenex sentence, following the inherited Skolemization definition; it is not a substitution in this scope. Blank macro filled with the glyph's TeX spelling. \\
\phantomsection\label{entry:symbol-registry-06934}\(h_i\) & \texttt{\detokenize{h_i}} & new Herbrand function symbol & \hyperref[sec:p02-w02-p02-028-validity-forms]{\ref*{sec:p02-w02-p02-028-validity-forms}} & n/a & \hyperref[def:p02-w02-p02-028-herbrand-form]{\ref*{def:p02-w02-p02-028-herbrand-form}} & \thepage & theta here is the quantifier-free matrix of a prenex sentence, following the inherited Skolemization definition; it is not a substitution in this scope. Blank macro filled with the glyph's TeX spelling. \\
\phantomsection\label{entry:symbol-registry-06935}\(\sigma^{\neg}\) & \texttt{\detokenize{\sigma^{\neg}}} & prenex dual of a prenex sentence (negation pushed through the prefix) & \hyperref[sec:p02-w02-p02-028-validity-forms]{\ref*{sec:p02-w02-p02-028-validity-forms}} & n/a & \hyperref[thm:p02-w02-p02-028-herbrand-skolem-duality]{\ref*{thm:p02-w02-p02-028-herbrand-skolem-duality}} & \thepage & Quantifiers exchanged and matrix negated; logically equivalent to the negation of sigma. \\
\phantomsection\label{entry:symbol-registry-06936}\(S_\sigma\) & \texttt{\detokenize{S_\sigma}} & clause family of the negated Herbrand matrix & \hyperref[sec:p02-w02-p02-028-clauses]{\ref*{sec:p02-w02-p02-028-clauses}} & n/a & \hyperref[thm:p02-w02-p02-028-validity-to-clauses]{\ref*{thm:p02-w02-p02-028-validity-to-clauses}} & \thepage & Cl(nnf(neg theta\^{}H)), read universally. \\
\phantomsection\label{entry:symbol-registry-06937}\(a\doteq b\) & \texttt{\detokenize{\doteq}} & Statman's ambiguous equation: either a=b or b=a & \hyperref[sec:p02-w02-p02-028-statman]{\ref*{sec:p02-w02-p02-028-statman}} & n/a & \hyperref[def:p02-w02-p02-028-st-K]{\ref*{def:p02-w02-p02-028-st-K}} & \thepage & Statman 1977 p.902 notation, kept as printed.  Overlaps the registered unification-equation role of \(\doteq\) (Parts 2--12); the local role is declared in def:p02-w02-p02-028-st-K and is confined to this section. \\
\phantomsection\label{entry:symbol-registry-06938}\(K\) & \texttt{\detokenize{K}} & Kreisel--Tait equation calculus and its substitution-only fragment symbol family member \(K\) & \hyperref[sec:p02-w02-p02-028-statman]{\ref*{sec:p02-w02-p02-028-statman}} & n/a & \hyperref[def:p02-w02-p02-028-st-K]{\ref*{def:p02-w02-p02-028-st-K}} & \thepage & Statman 1977 p.902; after p.904 K denotes the calculus without rules (3) and (4\_n). Explicit displayed family; member types and common actual first-use exposed, not one atomic glyph. Family listing, not one atomic glyph. Author macro \texttt{\detokenize{K,}\allowbreak\detokenize{H}} not carried to the members. \\
\phantomsection\label{entry:symbol-registry-06939}\(H\) & \texttt{\detokenize{H}} & Kreisel--Tait equation calculus and its substitution-only fragment symbol family member \(H\) & \hyperref[sec:p02-w02-p02-028-statman]{\ref*{sec:p02-w02-p02-028-statman}} & n/a & \hyperref[def:p02-w02-p02-028-st-K]{\ref*{def:p02-w02-p02-028-st-K}} & \thepage & Statman 1977 p.902; after p.904 K denotes the calculus without rules (3) and (4\_n). Explicit displayed family; member types and common actual first-use exposed, not one atomic glyph. Family listing, not one atomic glyph. Author macro \texttt{\detokenize{K,}\allowbreak\detokenize{H}} not carried to the members. \\
\phantomsection\label{entry:symbol-registry-06940}\(J\) & \texttt{\detokenize{J}} & Statman's axiom set (1)--(6\_n) for the structures with N and F & \hyperref[sec:p02-w02-p02-028-statman]{\ref*{sec:p02-w02-p02-028-statman}} & n/a & \hyperref[def:p02-w02-p02-028-st-J]{\ref*{def:p02-w02-p02-028-st-J}} & \thepage & Statman 1977 p.900. \\
\phantomsection\label{entry:symbol-registry-06941}\(s^n\) & \texttt{\detokenize{s^n}} & n-fold successor & \hyperref[sec:p02-w02-p02-028-statman]{\ref*{sec:p02-w02-p02-028-statman}} & n/a & \hyperref[def:p02-w02-p02-028-st-J]{\ref*{def:p02-w02-p02-028-st-J}} & \thepage & s\^{}1=s and s\^{}{k+1}=ss\^{}k (Statman 1977 p.900). \\
\phantomsection\label{entry:symbol-registry-06942}\(M\) & \texttt{\detokenize{M}} & measure on terms and its transfer along a substitution symbol family member \(M\) & \hyperref[sec:p02-w02-p02-028-statman]{\ref*{sec:p02-w02-p02-028-statman}} & n/a & \hyperref[def:p02-w02-p02-028-st-measure]{\ref*{def:p02-w02-p02-028-st-measure}} & \thepage & Statman 1977 p.901; M\^{}*(a)=M(theta a), or M(c(...,t/x\_t,...)) in the flattening proof. Explicit displayed family; member types and common actual first-use exposed, not one atomic glyph. Family listing, not one atomic glyph. Author macro \texttt{\detokenize{M,}\allowbreak\detokenize{M\^{}*}} not carried to the members. \\
\phantomsection\label{entry:symbol-registry-06943}\(M^*\) & \texttt{\detokenize{M^*}} & measure on terms and its transfer along a substitution symbol family member \(M^*\) & \hyperref[sec:p02-w02-p02-028-statman]{\ref*{sec:p02-w02-p02-028-statman}} & n/a & \hyperref[def:p02-w02-p02-028-st-measure]{\ref*{def:p02-w02-p02-028-st-measure}} & \thepage & Statman 1977 p.901; M\^{}*(a)=M(theta a), or M(c(...,t/x\_t,...)) in the flattening proof. Explicit displayed family; member types and common actual first-use exposed, not one atomic glyph. Family listing, not one atomic glyph. Author macro \texttt{\detokenize{M,}\allowbreak\detokenize{M\^{}*}} not carried to the members. \\
\phantomsection\label{entry:symbol-registry-06944}\(\operatorname{lh}\) & \texttt{\detokenize{\operatorname{lh}}} & length: number of symbol occurrences in a term or equation occurrences in a computation & \hyperref[sec:p02-w02-p02-028-statman]{\ref*{sec:p02-w02-p02-028-statman}} & n/a & \hyperref[def:p02-w02-p02-028-st-measure]{\ref*{def:p02-w02-p02-028-st-measure}} & \thepage & Statman 1977 pp.901--902 uses lh for both; the argument type disambiguates. \\
\phantomsection\label{entry:symbol-registry-06945}\([\cdots](a_1/x_1,\ldots,a_n/x_n)\) & \texttt{\detokenize{(a/}\allowbreak\detokenize{x)}} & Statman's simultaneous substitution notation & \hyperref[sec:p02-w02-p02-028-statman]{\ref*{sec:p02-w02-p02-028-statman}} & n/a & \hyperref[def:p02-w02-p02-028-st-terms]{\ref*{def:p02-w02-p02-028-st-terms}} & \thepage & Statman 1977 p.901; E(a/x) also marks the occurrences of x as the place of substitution in rule (1). \\
\phantomsection\label{entry:symbol-registry-06946}\(T_R\) & \texttt{\detokenize{T_R}} & computations obtained from T by Statman's transformations symbol family member \(T_R\) & \hyperref[sec:p02-w02-p02-028-statman]{\ref*{sec:p02-w02-p02-028-statman}} & n/a & \hyperref[def:p02-w02-p02-028-st-leftness]{\ref*{def:p02-w02-p02-028-st-leftness}} & \thepage & Statman 1977 pp.903--908: rightmost and leftmost reduction forms; M-direct, subterm and weak-subterm forms. Explicit displayed family; member types and common actual first-use exposed, not one atomic glyph. Family listing, not one atomic glyph. Author macro \texttt{\detokenize{T\_R,}\allowbreak\detokenize{T\_L,}\allowbreak\detokenize{T\_M,}\allowbreak\detokenize{T\_S,}\allowbreak\detokenize{T\_W}} not carried to the members. \\
\phantomsection\label{entry:symbol-registry-06947}\(T_L\) & \texttt{\detokenize{T_L}} & computations obtained from T by Statman's transformations symbol family member \(T_L\) & \hyperref[sec:p02-w02-p02-028-statman]{\ref*{sec:p02-w02-p02-028-statman}} & n/a & \hyperref[def:p02-w02-p02-028-st-leftness]{\ref*{def:p02-w02-p02-028-st-leftness}} & \thepage & Statman 1977 pp.903--908: rightmost and leftmost reduction forms; M-direct, subterm and weak-subterm forms. Explicit displayed family; member types and common actual first-use exposed, not one atomic glyph. Family listing, not one atomic glyph. Author macro \texttt{\detokenize{T\_R,}\allowbreak\detokenize{T\_L,}\allowbreak\detokenize{T\_M,}\allowbreak\detokenize{T\_S,}\allowbreak\detokenize{T\_W}} not carried to the members. \\
\phantomsection\label{entry:symbol-registry-06948}\(T_M\) & \texttt{\detokenize{T_M}} & computations obtained from T by Statman's transformations symbol family member \(T_M\) & \hyperref[sec:p02-w02-p02-028-statman]{\ref*{sec:p02-w02-p02-028-statman}} & n/a & \hyperref[def:p02-w02-p02-028-st-leftness]{\ref*{def:p02-w02-p02-028-st-leftness}} & \thepage & Statman 1977 pp.903--908: rightmost and leftmost reduction forms; M-direct, subterm and weak-subterm forms. Explicit displayed family; member types and common actual first-use exposed, not one atomic glyph. Family listing, not one atomic glyph. Author macro \texttt{\detokenize{T\_R,}\allowbreak\detokenize{T\_L,}\allowbreak\detokenize{T\_M,}\allowbreak\detokenize{T\_S,}\allowbreak\detokenize{T\_W}} not carried to the members. \\
\phantomsection\label{entry:symbol-registry-06949}\(T_S\) & \texttt{\detokenize{T_S}} & computations obtained from T by Statman's transformations symbol family member \(T_S\) & \hyperref[sec:p02-w02-p02-028-statman]{\ref*{sec:p02-w02-p02-028-statman}} & n/a & \hyperref[def:p02-w02-p02-028-st-leftness]{\ref*{def:p02-w02-p02-028-st-leftness}} & \thepage & Statman 1977 pp.903--908: rightmost and leftmost reduction forms; M-direct, subterm and weak-subterm forms. Explicit displayed family; member types and common actual first-use exposed, not one atomic glyph. Family listing, not one atomic glyph. Author macro \texttt{\detokenize{T\_R,}\allowbreak\detokenize{T\_L,}\allowbreak\detokenize{T\_M,}\allowbreak\detokenize{T\_S,}\allowbreak\detokenize{T\_W}} not carried to the members. \\
\phantomsection\label{entry:symbol-registry-06950}\(T_W\) & \texttt{\detokenize{T_W}} & computations obtained from T by Statman's transformations symbol family member \(T_W\) & \hyperref[sec:p02-w02-p02-028-statman]{\ref*{sec:p02-w02-p02-028-statman}} & n/a & \hyperref[def:p02-w02-p02-028-st-leftness]{\ref*{def:p02-w02-p02-028-st-leftness}} & \thepage & Statman 1977 pp.903--908: rightmost and leftmost reduction forms; M-direct, subterm and weak-subterm forms. Explicit displayed family; member types and common actual first-use exposed, not one atomic glyph. Family listing, not one atomic glyph. Author macro \texttt{\detokenize{T\_R,}\allowbreak\detokenize{T\_L,}\allowbreak\detokenize{T\_M,}\allowbreak\detokenize{T\_S,}\allowbreak\detokenize{T\_W}} not carried to the members. \\
\phantomsection\label{entry:symbol-registry-06951}\(S_n\) & \texttt{\detokenize{S_n}} & stage data of the identification procedure symbol family member \(S_n\) & \hyperref[def:p02-w02-p02-028-st-closure]{\ref*{def:p02-w02-p02-028-st-closure}} & n/a & \hyperref[def:p02-w02-p02-028-st-closure]{\ref*{def:p02-w02-p02-028-st-closure}} & \thepage & Statman 1977 pp.904--905; distinct from the example set S\_n of [2.15], whose scope is thm:p02-w02-p02-028-st-cycle. Explicit displayed family; member types and common actual first-use exposed, not one atomic glyph. Family listing, not one atomic glyph. Author macro \texttt{\detokenize{S\_n,}\allowbreak\detokenize{C\_n,}\allowbreak\detokenize{T\^{}*\_n,}\allowbreak\detokenize{x\_n,}\allowbreak\detokenize{y\_n}} not carried to the members. \\
\phantomsection\label{entry:symbol-registry-06952}\(C_n\) & \texttt{\detokenize{C_n}} & stage data of the identification procedure symbol family member \(C_n\) & \hyperref[def:p02-w02-p02-028-st-closure]{\ref*{def:p02-w02-p02-028-st-closure}} & n/a & \hyperref[def:p02-w02-p02-028-st-closure]{\ref*{def:p02-w02-p02-028-st-closure}} & \thepage & Statman 1977 pp.904--905; distinct from the example set S\_n of [2.15], whose scope is thm:p02-w02-p02-028-st-cycle. Explicit displayed family; member types and common actual first-use exposed, not one atomic glyph. Family listing, not one atomic glyph. Author macro \texttt{\detokenize{S\_n,}\allowbreak\detokenize{C\_n,}\allowbreak\detokenize{T\^{}*\_n,}\allowbreak\detokenize{x\_n,}\allowbreak\detokenize{y\_n}} not carried to the members. \\
\phantomsection\label{entry:symbol-registry-06953}\(T^*_n\) & \texttt{\detokenize{T^*_n}} & stage data of the identification procedure symbol family member \(T^*_n\) & \hyperref[def:p02-w02-p02-028-st-closure]{\ref*{def:p02-w02-p02-028-st-closure}} & n/a & \hyperref[def:p02-w02-p02-028-st-closure]{\ref*{def:p02-w02-p02-028-st-closure}} & \thepage & Statman 1977 pp.904--905; distinct from the example set S\_n of [2.15], whose scope is thm:p02-w02-p02-028-st-cycle. Explicit displayed family; member types and common actual first-use exposed, not one atomic glyph. Family listing, not one atomic glyph. Author macro \texttt{\detokenize{S\_n,}\allowbreak\detokenize{C\_n,}\allowbreak\detokenize{T\^{}*\_n,}\allowbreak\detokenize{x\_n,}\allowbreak\detokenize{y\_n}} not carried to the members. \\
\phantomsection\label{entry:symbol-registry-06954}\(x_n\) & \texttt{\detokenize{x_n}} & stage data of the identification procedure symbol family member \(x_n\) & \hyperref[def:p02-w02-p02-028-st-closure]{\ref*{def:p02-w02-p02-028-st-closure}} & n/a & \hyperref[def:p02-w02-p02-028-st-closure]{\ref*{def:p02-w02-p02-028-st-closure}} & \thepage & Statman 1977 pp.904--905; distinct from the example set S\_n of [2.15], whose scope is thm:p02-w02-p02-028-st-cycle. Explicit displayed family; member types and common actual first-use exposed, not one atomic glyph. Family listing, not one atomic glyph. Author macro \texttt{\detokenize{S\_n,}\allowbreak\detokenize{C\_n,}\allowbreak\detokenize{T\^{}*\_n,}\allowbreak\detokenize{x\_n,}\allowbreak\detokenize{y\_n}} not carried to the members. \\
\phantomsection\label{entry:symbol-registry-06955}\(y_n\) & \texttt{\detokenize{y_n}} & stage data of the identification procedure symbol family member \(y_n\) & \hyperref[def:p02-w02-p02-028-st-closure]{\ref*{def:p02-w02-p02-028-st-closure}} & n/a & \hyperref[def:p02-w02-p02-028-st-closure]{\ref*{def:p02-w02-p02-028-st-closure}} & \thepage & Statman 1977 pp.904--905; distinct from the example set S\_n of [2.15], whose scope is thm:p02-w02-p02-028-st-cycle. Explicit displayed family; member types and common actual first-use exposed, not one atomic glyph. Family listing, not one atomic glyph. Author macro \texttt{\detokenize{S\_n,}\allowbreak\detokenize{C\_n,}\allowbreak\detokenize{T\^{}*\_n,}\allowbreak\detokenize{x\_n,}\allowbreak\detokenize{y\_n}} not carried to the members. \\
\phantomsection\label{entry:symbol-registry-06956}\(\bigwedge S\to E\) & \texttt{\detokenize{\bigwedge S\to E}} & implication from the conjunction of a finite equation set & \hyperref[sec:p02-w02-p02-028-statman]{\ref*{sec:p02-w02-p02-028-statman}} & n/a & \hyperref[def:p02-w02-p02-028-st-analytic]{\ref*{def:p02-w02-p02-028-st-analytic}} & \thepage & Statman 1977 p.904; analytic when valid with satisfiable antecedent. \\
\phantomsection\label{entry:symbol-registry-06957}\(\approx\) & \texttt{\detokenize{\approx}} & fresh binary relation symbol read in place of logical equality & \hyperref[sec:annex-p02-w02-p02-028]{\ref*{sec:annex-p02-w02-p02-028}} & n/a & \hyperref[sec:annex-p02-w02-p02-028]{\ref*{sec:annex-p02-w02-p02-028}} & \thepage & Annex-local; used in the proofs of Herbrand's theorem with equality and of the decidability of quantifier-free validity.  Interpreted as identity when the equality axioms are discarded. \\
\phantomsection\label{entry:symbol-registry-06958}\(\varphi^{\approx}\) & \texttt{\detokenize{\varphi^{\approx}}} & the formula with every atom r=t replaced by \(r\approx t\) & \hyperref[sec:annex-p02-w02-p02-028]{\ref*{sec:annex-p02-w02-p02-028}} & n/a & \hyperref[sec:annex-p02-w02-p02-028]{\ref*{sec:annex-p02-w02-p02-028}} & \thepage & Annex-local; used in the proofs of Herbrand's theorem with equality and of the decidability of quantifier-free validity.  Interpreted as identity when the equality axioms are discarded. \\
\phantomsection\label{entry:symbol-registry-06959}\(\mathrm{Eq}\) & \texttt{\detokenize{\mathrm{Eq}}} & quantifier-free equality axioms for the symbols of a formula & \hyperref[annex:p02-w02-p02-028-st-herbrand-eq]{\ref*{annex:p02-w02-p02-028-st-herbrand-eq}} & n/a & \hyperref[annex:p02-w02-p02-028-st-herbrand-eq]{\ref*{annex:p02-w02-p02-028-st-herbrand-eq}} & \thepage & Annex-local to annex:p02-w02-p02-028-st-herbrand-eq. \\
\phantomsection\label{entry:symbol-registry-06960}\(\varepsilon\) & \texttt{\detokenize{\varepsilon}} & conjunction of the equality axioms with variables renamed apart & \hyperref[annex:p02-w02-p02-028-st-herbrand-eq]{\ref*{annex:p02-w02-p02-028-st-herbrand-eq}} & n/a & \hyperref[annex:p02-w02-p02-028-st-herbrand-eq]{\ref*{annex:p02-w02-p02-028-st-herbrand-eq}} & \thepage & Annex-local to annex:p02-w02-p02-028-st-herbrand-eq. \\
\phantomsection\label{entry:symbol-registry-06961}\(e|_\pi\) & \texttt{\detokenize{e|_\pi}} & subterm of an equation at a position & \hyperref[annex:p02-w02-p02-028-st-nesting]{\ref*{annex:p02-w02-p02-028-st-nesting}} & n/a & \hyperref[annex:p02-w02-p02-028-st-nesting]{\ref*{annex:p02-w02-p02-028-st-nesting}} & \thepage & Positions as in def:p02-w02-p02-028-st-positions (1 and 2 for the two sides, pi i for the i-th argument). \\
\phantomsection\label{entry:symbol-registry-06962}\(e[c]_\pi\) & \texttt{\detokenize{e[c]_\pi}} & result of putting a term at a position of an equation & \hyperref[annex:p02-w02-p02-028-st-nesting]{\ref*{annex:p02-w02-p02-028-st-nesting}} & n/a & \hyperref[annex:p02-w02-p02-028-st-nesting]{\ref*{annex:p02-w02-p02-028-st-nesting}} & \thepage & Positions as in def:p02-w02-p02-028-st-positions (1 and 2 for the two sides, pi i for the i-th argument). \\
\phantomsection\label{entry:symbol-registry-06963}\(i(F)\) & \texttt{\detokenize{i(F)}} & number of inferences of an H fragment & \hyperref[annex:p02-w02-p02-028-st-prop-3-5]{\ref*{annex:p02-w02-p02-028-st-prop-3-5}} & n/a & \hyperref[annex:p02-w02-p02-028-st-prop-3-5]{\ref*{annex:p02-w02-p02-028-st-prop-3-5}} & \thepage & Annex-local to the proof of [3.5]; p(R) plays the role of the chapter's p (Statman 1977 p.908). \\
\phantomsection\label{entry:symbol-registry-06964}\(U_0(R)\) & \texttt{\detokenize{U_0(R)}} & original premise subcomputation of a bad inference of form (2) & \hyperref[annex:p02-w02-p02-028-st-prop-3-5]{\ref*{annex:p02-w02-p02-028-st-prop-3-5}} & n/a & \hyperref[annex:p02-w02-p02-028-st-prop-3-5]{\ref*{annex:p02-w02-p02-028-st-prop-3-5}} & \thepage & Annex-local to the proof of [3.5]; p(R) plays the role of the chapter's p (Statman 1977 p.908). \\
\phantomsection\label{entry:symbol-registry-06965}\(p(R)\) & \texttt{\detokenize{p(R)}} & length of the original premise subcomputation of a bad inference & \hyperref[annex:p02-w02-p02-028-st-prop-3-5]{\ref*{annex:p02-w02-p02-028-st-prop-3-5}} & n/a & \hyperref[annex:p02-w02-p02-028-st-prop-3-5]{\ref*{annex:p02-w02-p02-028-st-prop-3-5}} & \thepage & Annex-local to the proof of [3.5]; p(R) plays the role of the chapter's p (Statman 1977 p.908). \\
\phantomsection\label{entry:symbol-registry-06966}\(N_R\) & \texttt{\detokenize{N_R}} & singular computation replacing the premise subcomputation of a bad inference & \hyperref[annex:p02-w02-p02-028-st-prop-3-5]{\ref*{annex:p02-w02-p02-028-st-prop-3-5}} & n/a & \hyperref[annex:p02-w02-p02-028-st-prop-3-5]{\ref*{annex:p02-w02-p02-028-st-prop-3-5}} & \thepage & Annex-local to the proof of [3.5]; p(R) plays the role of the chapter's p (Statman 1977 p.908). \\
\phantomsection\label{entry:symbol-registry-06967}\(\sigma_\kappa\) & \texttt{\detokenize{\sigma_\kappa}} & projection-anticipating substitution & \hyperref[annex:p02-w02-p02-028-st-prop-4-1-1]{\ref*{annex:p02-w02-p02-028-st-prop-4-1-1}} & n/a & \hyperref[annex:p02-w02-p02-028-st-prop-4-1-1]{\ref*{annex:p02-w02-p02-028-st-prop-4-1-1}} & \thepage & Annex-local to the proof of [4.1.1]; theta\_kappa is the chapter's theta\_i (Statman 1977 p.911). \\
\phantomsection\label{entry:symbol-registry-06968}\(\theta_\kappa\) & \texttt{\detokenize{\theta_\kappa}} & the matching non-projecting substitution & \hyperref[annex:p02-w02-p02-028-st-prop-4-1-1]{\ref*{annex:p02-w02-p02-028-st-prop-4-1-1}} & n/a & \hyperref[annex:p02-w02-p02-028-st-prop-4-1-1]{\ref*{annex:p02-w02-p02-028-st-prop-4-1-1}} & \thepage & Annex-local to the proof of [4.1.1]; theta\_kappa is the chapter's theta\_i (Statman 1977 p.911). \\
\phantomsection\label{entry:symbol-registry-06969}\(f_{l,J}\) & \texttt{\detokenize{f_{l,J}}} & new function variable for the argument places J of f\_l & \hyperref[annex:p02-w02-p02-028-st-prop-4-1-1]{\ref*{annex:p02-w02-p02-028-st-prop-4-1-1}} & n/a & \hyperref[annex:p02-w02-p02-028-st-prop-4-1-1]{\ref*{annex:p02-w02-p02-028-st-prop-4-1-1}} & \thepage & Annex-local to the proof of [4.1.1]; theta\_kappa is the chapter's theta\_i (Statman 1977 p.911). \\
\phantomsection\label{entry:symbol-registry-06970}\(\mathcal R\) & \texttt{\detokenize{\mathcal R}} & finite set of analytic rules & \hyperref[sec:p02-w02-p02-028-st-cases]{\ref*{sec:p02-w02-p02-028-st-cases}} & n/a & \hyperref[thm:p02-w02-p02-028-st-optimality]{\ref*{thm:p02-w02-p02-028-st-optimality}} & \thepage & The chapter (p.910) does not define analytic rules; the reading is stated in the reader text before thm:p02-w02-p02-028-st-optimality. Blank macro filled with the glyph's TeX spelling. \\
\phantomsection\label{entry:symbol-registry-06971}\(S_\rho\) & \texttt{\detokenize{S_\rho}} & premises of the rule rho & \hyperref[sec:p02-w02-p02-028-st-cases]{\ref*{sec:p02-w02-p02-028-st-cases}} & n/a & \hyperref[thm:p02-w02-p02-028-st-optimality]{\ref*{thm:p02-w02-p02-028-st-optimality}} & \thepage & The chapter (p.910) does not define analytic rules; the reading is stated in the reader text before thm:p02-w02-p02-028-st-optimality. Blank macro filled with the glyph's TeX spelling. \\
\phantomsection\label{entry:symbol-registry-06972}\(E_\rho\) & \texttt{\detokenize{E_\rho}} & conclusion of the rule rho & \hyperref[sec:p02-w02-p02-028-st-cases]{\ref*{sec:p02-w02-p02-028-st-cases}} & n/a & \hyperref[thm:p02-w02-p02-028-st-optimality]{\ref*{thm:p02-w02-p02-028-st-optimality}} & \thepage & The chapter (p.910) does not define analytic rules; the reading is stated in the reader text before thm:p02-w02-p02-028-st-optimality. Blank macro filled with the glyph's TeX spelling. \\
\phantomsection\label{entry:symbol-registry-06973}\(E\) & \texttt{\detokenize{E}} & antecedent of a formula G of the class L (the conjunction of the occurrences of the A\_i; I\_j; B\_k in a fixed bracketing) & \hyperref[def:p02-w02-p02-028-st-horn-cut-setting]{\ref*{def:p02-w02-p02-028-st-horn-cut-setting}} & n/a & \hyperref[def:p02-w02-p02-028-st-horn-cut-setting]{\ref*{def:p02-w02-p02-028-st-horn-cut-setting}} & \thepage & The occurrences are the leaves of E; theta E' for a conjunction E' inside E is a left formula. \\
\phantomsection\label{entry:symbol-registry-06974}\(\operatorname{par}(G)\) & \texttt{\detokenize{\operatorname{par}}} & variables of G other than x\_1;...;x\_n & \hyperref[def:p02-w02-p02-028-st-horn-cut-setting]{\ref*{def:p02-w02-p02-028-st-horn-cut-setting}} & n/a & \hyperref[def:p02-w02-p02-028-st-horn-cut-setting]{\ref*{def:p02-w02-p02-028-st-horn-cut-setting}} & \thepage & Fixed by every substitution theta for x\_1...x\_n; binders of a G-form avoid it. \\
\phantomsection\label{entry:symbol-registry-06975}\(\mathsf H(\Theta;\Gamma)\) & \texttt{\detokenize{\mathsf H(\Theta;\Gamma)}} & Horn system read off a finite set Theta of substitutions and a multiset Gamma of left formulas & \hyperref[def:p02-w02-p02-028-st-horn-system]{\ref*{def:p02-w02-p02-028-st-horn-system}} & n/a & \hyperref[def:p02-w02-p02-028-st-horn-system]{\ref*{def:p02-w02-p02-028-st-horn-system}} & \thepage & H(Theta;empty) is the system of Statman 1977 [1.3(ii)] restricted to Theta.  \(\mathsf H\) is also registered for a P8 past-tense operator; this role is confined to the cut-elimination route of sec:p02-w02-p02-028-st-cases. \\
\phantomsection\label{entry:symbol-registry-06976}\(\eta\) & \texttt{\detokenize{\eta}} & map from variables to terms & \hyperref[annex:p02-w02-p02-028-st-gform]{\ref*{annex:p02-w02-p02-028-st-gform}} & n/a & \hyperref[annex:p02-w02-p02-028-st-gform]{\ref*{annex:p02-w02-p02-028-st-gform}} & \thepage & Annex-local to annex:p02-w02-p02-028-st-gform (eta) and annex:p02-w02-p02-028-st-extraction (beta, beta\_1, beta\_2).  eta is not the solving substitution of the unification proof, which is scoped to its own annex item. \\
\phantomsection\label{entry:symbol-registry-06977}\(\beta\) & \texttt{\detokenize{\beta}} & target of a succedent & \hyperref[annex:p02-w02-p02-028-st-extraction]{\ref*{annex:p02-w02-p02-028-st-extraction}} & n/a & \hyperref[annex:p02-w02-p02-028-st-extraction]{\ref*{annex:p02-w02-p02-028-st-extraction}} & \thepage & Annex-local to annex:p02-w02-p02-028-st-gform (eta) and annex:p02-w02-p02-028-st-extraction (beta, beta\_1, beta\_2).  eta is not the solving substitution of the unification proof, which is scoped to its own annex item. \\
\phantomsection\label{entry:symbol-registry-06978}\(\mathrm{Var}_{\mathrm{MS}}\) & \texttt{\detokenize{\mathrm{Var}_{\mathrm{MS}}}} & set of all sorted variables of a strict many-sorted signature & \hyperref[sec:p02-w02-p02-038-many-sorted-reduction]{\ref*{sec:p02-w02-p02-038-many-sorted-reduction}} & n/a & \hyperref[sec:p02-w02-p02-038-clauses]{\ref*{sec:p02-w02-p02-038-clauses}} & \thepage & Union of the per-sort variable stocks of def:p02-many-sorted-signature; identified with its image under the fixed injection into the one-sorted variable supply. \\
\phantomsection\label{entry:symbol-registry-06979}\(\Delta\) & \texttt{\detokenize{\Delta}} & context: a set of sorted variables & \hyperref[sec:p02-w02-p02-038-many-sorted-reduction]{\ref*{sec:p02-w02-p02-038-many-sorted-reduction}} & n/a & \hyperref[def:p02-w02-p02-038-sorted-structure]{\ref*{def:p02-w02-p02-038-sorted-structure}} & \thepage & Local role only; Delta has many other roles elsewhere in the book (sets of formulas, diagrams). Declared before use here. \\
\phantomsection\label{entry:symbol-registry-06980}\(g\) & \texttt{\detokenize{g}} & contextual sorted assignment and its update symbol family member \(g\) & \hyperref[sec:p02-w02-p02-038-many-sorted-reduction]{\ref*{sec:p02-w02-p02-038-many-sorted-reduction}} & n/a & \hyperref[def:p02-w02-p02-038-sorted-structure]{\ref*{def:p02-w02-p02-038-sorted-structure}} & \thepage & A Delta-assignment sends each variable of the context into the carrier of its sort; total assignments are written s as in def:p02-many-sorted-structure. Explicit displayed family; member types and common actual first-use exposed, not one atomic glyph. Family listing, not one atomic glyph. Author macro \texttt{\detokenize{g,}\allowbreak\detokenize{g[v\mapsto b]}} not carried to the members. \\
\phantomsection\label{entry:symbol-registry-06981}\(\ g[v\mapsto b]\) & \texttt{\detokenize{\ g[v\mapsto b]}} & contextual sorted assignment and its update symbol family member \(\ g[v\mapsto b]\) & \hyperref[sec:p02-w02-p02-038-many-sorted-reduction]{\ref*{sec:p02-w02-p02-038-many-sorted-reduction}} & n/a & \hyperref[def:p02-w02-p02-038-sorted-structure]{\ref*{def:p02-w02-p02-038-sorted-structure}} & \thepage & A Delta-assignment sends each variable of the context into the carrier of its sort; total assignments are written s as in def:p02-many-sorted-structure. Explicit displayed family; member types and common actual first-use exposed, not one atomic glyph. Family listing, not one atomic glyph. Author macro \texttt{\detokenize{g,}\allowbreak\detokenize{g[v\mapsto b]}} not carried to the members. \\
\phantomsection\label{entry:symbol-registry-06982}\(\llbracket t\rrbracket^{\mathcal M}_{g}\) & \texttt{\detokenize{\llbracket t\rrbracket^{\mathcal M}_{g}}} & value of a sorted term under a contextual assignment & \hyperref[sec:p02-w02-p02-038-many-sorted-reduction]{\ref*{sec:p02-w02-p02-038-many-sorted-reduction}} & n/a & \hyperref[def:p02-w02-p02-038-satisfaction]{\ref*{def:p02-w02-p02-038-satisfaction}} & \thepage & Same bracket glyph as the categorical interpretation brackets of Part 10; here a set-theoretic term value. Scopes are disjoint. \\
\phantomsection\label{entry:symbol-registry-06983}\(\mathcal M\entails\varphi[g]\) & \texttt{\detokenize{\mathcal M\entails\varphi[g]}} & satisfaction of a sorted formula at a contextual assignment & \hyperref[sec:p02-w02-p02-038-many-sorted-reduction]{\ref*{sec:p02-w02-p02-038-many-sorted-reduction}} & n/a & \hyperref[def:p02-w02-p02-038-satisfaction]{\ref*{def:p02-w02-p02-038-satisfaction}} & \thepage & Requires FV(varphi) within the domain of g; for strict structures and total assignments it is M,s entails varphi of def:p02-many-sorted-structure. \\
\phantomsection\label{entry:symbol-registry-06984}\((\mathrm{Ne})\) & \texttt{\detokenize{(\mathrm{Ne})}} & names of the four sort-axiom families symbol family member \((\mathrm{Ne})\) & \hyperref[sec:p02-w02-p02-038-many-sorted-reduction]{\ref*{sec:p02-w02-p02-038-many-sorted-reduction}} & n/a & \hyperref[def:p02-w02-p02-038-axiom-families]{\ref*{def:p02-w02-p02-038-axiom-families}} & \thepage & Nonemptiness, constant typing, function closure, disjointness; together exactly Ax\_sort of def:p02-many-sorted-translation. Explicit displayed family; member types and common actual first-use exposed, not one atomic glyph. Family listing, not one atomic glyph. Author macro \texttt{\detokenize{(\mathrm{Ne}),}\allowbreak\detokenize{(\mathrm{Ct}),}\allowbreak\detokenize{(\mathrm{Cl}),}\allowbreak\detokenize{(\mathrm{Dj})}} not carried to the members. \\
\phantomsection\label{entry:symbol-registry-06985}\((\mathrm{Ct})\) & \texttt{\detokenize{(\mathrm{Ct})}} & names of the four sort-axiom families symbol family member \((\mathrm{Ct})\) & \hyperref[sec:p02-w02-p02-038-many-sorted-reduction]{\ref*{sec:p02-w02-p02-038-many-sorted-reduction}} & n/a & \hyperref[def:p02-w02-p02-038-axiom-families]{\ref*{def:p02-w02-p02-038-axiom-families}} & \thepage & Nonemptiness, constant typing, function closure, disjointness; together exactly Ax\_sort of def:p02-many-sorted-translation. Explicit displayed family; member types and common actual first-use exposed, not one atomic glyph. Family listing, not one atomic glyph. Author macro \texttt{\detokenize{(\mathrm{Ne}),}\allowbreak\detokenize{(\mathrm{Ct}),}\allowbreak\detokenize{(\mathrm{Cl}),}\allowbreak\detokenize{(\mathrm{Dj})}} not carried to the members. \\
\phantomsection\label{entry:symbol-registry-06986}\((\mathrm{Cl})\) & \texttt{\detokenize{(\mathrm{Cl})}} & names of the four sort-axiom families symbol family member \((\mathrm{Cl})\) & \hyperref[sec:p02-w02-p02-038-many-sorted-reduction]{\ref*{sec:p02-w02-p02-038-many-sorted-reduction}} & n/a & \hyperref[def:p02-w02-p02-038-axiom-families]{\ref*{def:p02-w02-p02-038-axiom-families}} & \thepage & Nonemptiness, constant typing, function closure, disjointness; together exactly Ax\_sort of def:p02-many-sorted-translation. Explicit displayed family; member types and common actual first-use exposed, not one atomic glyph. Family listing, not one atomic glyph. Author macro \texttt{\detokenize{(\mathrm{Ne}),}\allowbreak\detokenize{(\mathrm{Ct}),}\allowbreak\detokenize{(\mathrm{Cl}),}\allowbreak\detokenize{(\mathrm{Dj})}} not carried to the members. \\
\phantomsection\label{entry:symbol-registry-06987}\((\mathrm{Dj})\) & \texttt{\detokenize{(\mathrm{Dj})}} & names of the four sort-axiom families symbol family member \((\mathrm{Dj})\) & \hyperref[sec:p02-w02-p02-038-many-sorted-reduction]{\ref*{sec:p02-w02-p02-038-many-sorted-reduction}} & n/a & \hyperref[def:p02-w02-p02-038-axiom-families]{\ref*{def:p02-w02-p02-038-axiom-families}} & \thepage & Nonemptiness, constant typing, function closure, disjointness; together exactly Ax\_sort of def:p02-many-sorted-translation. Explicit displayed family; member types and common actual first-use exposed, not one atomic glyph. Family listing, not one atomic glyph. Author macro \texttt{\detokenize{(\mathrm{Ne}),}\allowbreak\detokenize{(\mathrm{Ct}),}\allowbreak\detokenize{(\mathrm{Cl}),}\allowbreak\detokenize{(\mathrm{Dj})}} not carried to the members. \\
\phantomsection\label{entry:symbol-registry-06988}\(\mathsf{Ax}^{\mathrm E}_{\mathrm{sort}}\) & \texttt{\detokenize{\mathsf{Ax}^{\mathrm E}_{\mathrm{sort}}}} & sort theories for the empty-sort and overlapping-sort reductions symbol family member \(\mathsf{Ax}^{\mathrm E}_{\mathrm{sort}}\) & \hyperref[sec:p02-w02-p02-038-many-sorted-reduction]{\ref*{sec:p02-w02-p02-038-many-sorted-reduction}} & meta & \hyperref[def:p02-w02-p02-038-axiom-families]{\ref*{def:p02-w02-p02-038-axiom-families}} & \thepage & Ax\_sort without the nonemptiness family, and Ax\_sort without the disjointness family, respectively. Explicit displayed family; member types and common actual first-use exposed, not one atomic glyph. Family listing, not one atomic glyph. Author macro \texttt{\detokenize{\mathsf{Ax}\^{}{\mathrm E}\_{\mathrm{sort}},}\allowbreak\detokenize{\mathsf{Ax}\^{}{\cup}\_{\mathrm{sort}}}} not carried to the members. \\
\phantomsection\label{entry:symbol-registry-06989}\(\mathsf{Ax}^{\cup}_{\mathrm{sort}}\) & \texttt{\detokenize{\mathsf{Ax}^{\cup}_{\mathrm{sort}}}} & sort theories for the empty-sort and overlapping-sort reductions symbol family member \(\mathsf{Ax}^{\cup}_{\mathrm{sort}}\) & \hyperref[sec:p02-w02-p02-038-many-sorted-reduction]{\ref*{sec:p02-w02-p02-038-many-sorted-reduction}} & meta & \hyperref[def:p02-w02-p02-038-axiom-families]{\ref*{def:p02-w02-p02-038-axiom-families}} & \thepage & Ax\_sort without the nonemptiness family, and Ax\_sort without the disjointness family, respectively. Explicit displayed family; member types and common actual first-use exposed, not one atomic glyph. Family listing, not one atomic glyph. Author macro \texttt{\detokenize{\mathsf{Ax}\^{}{\mathrm E}\_{\mathrm{sort}},}\allowbreak\detokenize{\mathsf{Ax}\^{}{\cup}\_{\mathrm{sort}}}} not carried to the members. \\
\phantomsection\label{entry:symbol-registry-06990}\(\mathsf{Guard}(\Delta)\) & \texttt{\detokenize{\mathsf{Guard}(\Delta)}} & context guard: sort atoms for the variables of a context & \hyperref[sec:p02-w02-p02-038-many-sorted-reduction]{\ref*{sec:p02-w02-p02-038-many-sorted-reduction}} & meta & \hyperref[def:p02-w02-p02-038-axiom-families]{\ref*{def:p02-w02-p02-038-axiom-families}} & \thepage & Guard of the set of all sorted variables is VarGuard\_MS of def:p02-ms-variable-guard. \\
\phantomsection\label{entry:symbol-registry-06991}\(\mathsf{Cov}(\Lang_{\mathrm{MS}})\) & \texttt{\detokenize{\mathsf{Cov}(\Lang_{\mathrm{MS}})}} & coverage axiom for finitely many sorts & \hyperref[sec:p02-w02-p02-038-many-sorted-reduction]{\ref*{sec:p02-w02-p02-038-many-sorted-reduction}} & n/a & \hyperref[def:p02-w02-p02-038-coverage]{\ref*{def:p02-w02-p02-038-coverage}} & \thepage & One-sorted sentence: every element satisfies some sort predicate; defined only when Sort is finite. \\
\phantomsection\label{entry:symbol-registry-06992}\(\mathcal A^{\downarrow}\) & \texttt{\detokenize{\mathcal A^{\downarrow}}} & restriction of a one-sorted model to its sort predicates & \hyperref[sec:p02-w02-p02-038-many-sorted-reduction]{\ref*{sec:p02-w02-p02-038-many-sorted-reduction}} & n/a & \hyperref[def:p02-w02-p02-038-restriction]{\ref*{def:p02-w02-p02-038-restriction}} & \thepage & Sorted structure with carriers Q\_s\^{}A and restricted interpretations. \\
\phantomsection\label{entry:symbol-registry-06993}\(\mathcal M^{h}\) & \texttt{\detokenize{\mathcal M^{h}}} & transport of a sorted structure along sort injections with fill functions symbol family member \(\mathcal M^{h}\) & \hyperref[sec:p02-w02-p02-038-many-sorted-reduction]{\ref*{sec:p02-w02-p02-038-many-sorted-reduction}} & n/a & \hyperref[def:p02-w02-p02-038-transport]{\ref*{def:p02-w02-p02-038-transport}} & \thepage & h\_s injects the carrier of sort s into the one-sorted domain; j\_f gives the value of f at ill-typed inputs. Explicit displayed family; member types and common actual first-use exposed, not one atomic glyph. Family listing, not one atomic glyph. Author macro \texttt{\detokenize{\mathcal M\^{}{h},}\allowbreak\detokenize{h\_{\mathsf s},}\allowbreak\detokenize{j\_f}} not carried to the members. \\
\phantomsection\label{entry:symbol-registry-06994}\(h_{\mathsf s}\) & \texttt{\detokenize{h_{\mathsf s}}} & transport of a sorted structure along sort injections with fill functions symbol family member \(h_{\mathsf s}\) & \hyperref[sec:p02-w02-p02-038-many-sorted-reduction]{\ref*{sec:p02-w02-p02-038-many-sorted-reduction}} & n/a & \hyperref[def:p02-w02-p02-038-transport]{\ref*{def:p02-w02-p02-038-transport}} & \thepage & h\_s injects the carrier of sort s into the one-sorted domain; j\_f gives the value of f at ill-typed inputs. Explicit displayed family; member types and common actual first-use exposed, not one atomic glyph. Family listing, not one atomic glyph. Author macro \texttt{\detokenize{\mathcal M\^{}{h},}\allowbreak\detokenize{h\_{\mathsf s},}\allowbreak\detokenize{j\_f}} not carried to the members. \\
\phantomsection\label{entry:symbol-registry-06995}\(j_f\) & \texttt{\detokenize{j_f}} & transport of a sorted structure along sort injections with fill functions symbol family member \(j_f\) & \hyperref[sec:p02-w02-p02-038-many-sorted-reduction]{\ref*{sec:p02-w02-p02-038-many-sorted-reduction}} & n/a & \hyperref[def:p02-w02-p02-038-transport]{\ref*{def:p02-w02-p02-038-transport}} & \thepage & h\_s injects the carrier of sort s into the one-sorted domain; j\_f gives the value of f at ill-typed inputs. Explicit displayed family; member types and common actual first-use exposed, not one atomic glyph. Family listing, not one atomic glyph. Author macro \texttt{\detokenize{\mathcal M\^{}{h},}\allowbreak\detokenize{h\_{\mathsf s},}\allowbreak\detokenize{j\_f}} not carried to the members. \\
\phantomsection\label{entry:symbol-registry-06996}\(\mathcal M^{\sharp}\) & \texttt{\detokenize{\mathcal M^{\sharp}}} & tagged one-sorted structure with junk point and its tagged assignment symbol family member \(\mathcal M^{\sharp}\) & \hyperref[sec:p02-w02-p02-038-many-sorted-reduction]{\ref*{sec:p02-w02-p02-038-many-sorted-reduction}} & n/a & \hyperref[def:p02-w02-p02-038-transport]{\ref*{def:p02-w02-p02-038-transport}} & \thepage & Junk point star is the pair (Sort,Sort). The sharp superscript also marks the incidence-translation assignment in thm:p02-sol-incidence-satisfaction; different scope and object. Explicit displayed family; member types and common actual first-use exposed, not one atomic glyph. Family listing, not one atomic glyph. Author macro \texttt{\detokenize{\mathcal M\^{}{\sharp},}\allowbreak\detokenize{s\^{}{\sharp},}\allowbreak\detokenize{\star}} not carried to the members. \\
\phantomsection\label{entry:symbol-registry-06997}\(\ s^{\sharp}\) & \texttt{\detokenize{\ s^{\sharp}}} & tagged one-sorted structure with junk point and its tagged assignment symbol family member \(\ s^{\sharp}\) & \hyperref[sec:p02-w02-p02-038-many-sorted-reduction]{\ref*{sec:p02-w02-p02-038-many-sorted-reduction}} & n/a & \hyperref[def:p02-w02-p02-038-transport]{\ref*{def:p02-w02-p02-038-transport}} & \thepage & Junk point star is the pair (Sort,Sort). The sharp superscript also marks the incidence-translation assignment in thm:p02-sol-incidence-satisfaction; different scope and object. Explicit displayed family; member types and common actual first-use exposed, not one atomic glyph. Family listing, not one atomic glyph. Author macro \texttt{\detokenize{\mathcal M\^{}{\sharp},}\allowbreak\detokenize{s\^{}{\sharp},}\allowbreak\detokenize{\star}} not carried to the members. \\
\phantomsection\label{entry:symbol-registry-06998}\(\ \star\) & \texttt{\detokenize{\ \star}} & tagged one-sorted structure with junk point and its tagged assignment symbol family member \(\ \star\) & \hyperref[sec:p02-w02-p02-038-many-sorted-reduction]{\ref*{sec:p02-w02-p02-038-many-sorted-reduction}} & n/a & \hyperref[def:p02-w02-p02-038-transport]{\ref*{def:p02-w02-p02-038-transport}} & \thepage & Junk point star is the pair (Sort,Sort). The sharp superscript also marks the incidence-translation assignment in thm:p02-sol-incidence-satisfaction; different scope and object. Explicit displayed family; member types and common actual first-use exposed, not one atomic glyph. Family listing, not one atomic glyph. Author macro \texttt{\detokenize{\mathcal M\^{}{\sharp},}\allowbreak\detokenize{s\^{}{\sharp},}\allowbreak\detokenize{\star}} not carried to the members. \\
\phantomsection\label{entry:symbol-registry-06999}\(\mathcal M^{\sharp}_{s}\) & \texttt{\detokenize{\mathcal M^{\sharp}_{s}}} & tagged one-sorted structure with defaults from a total assignment & \hyperref[sec:p02-w02-p02-038-many-sorted-reduction]{\ref*{sec:p02-w02-p02-038-many-sorted-reduction}} & n/a & \hyperref[def:p02-w02-p02-038-transport]{\ref*{def:p02-w02-p02-038-transport}} & \thepage & Used for the coverage variant; no junk point. \\
\phantomsection\label{entry:symbol-registry-07000}\(\mathcal M^{\cup}\) & \texttt{\detokenize{\mathcal M^{\cup}}} & plain-union one-sorted structure with junk point & \hyperref[sec:p02-w02-p02-038-many-sorted-reduction]{\ref*{sec:p02-w02-p02-038-many-sorted-reduction}} & n/a & \hyperref[def:p02-w02-p02-038-transport]{\ref*{def:p02-w02-p02-038-transport}} & \thepage & Domain is the union U of the carriers together with U itself as junk point. \\
\phantomsection\label{entry:symbol-registry-07001}\(\entails^{\Delta}_{\mathrm E}\) & \texttt{\detokenize{\entails^{\Delta}_{\mathrm E}}} & empty-sort consequence in a context and with total assignments symbol family member \(\entails^{\Delta}_{\mathrm E}\) & \hyperref[sec:p02-w02-p02-038-many-sorted-reduction]{\ref*{sec:p02-w02-p02-038-many-sorted-reduction}} & meta & \hyperref[def:p02-w02-p02-038-empty-semantics]{\ref*{def:p02-w02-p02-038-empty-semantics}} & \thepage & Quantifies over E-structures (disjoint carriers, possibly empty). Explicit displayed family; member types and common actual first-use exposed, not one atomic glyph. Family listing, not one atomic glyph. Author macro \texttt{\detokenize{\entails\^{}{\Delta}\_{\mathrm E},}\allowbreak\detokenize{\entails\^{}{\mathrm{tot}}\_{\mathrm E}}} not carried to the members. \\
\phantomsection\label{entry:symbol-registry-07002}\(\entails^{\mathrm{tot}}_{\mathrm E}\) & \texttt{\detokenize{\entails^{\mathrm{tot}}_{\mathrm E}}} & empty-sort consequence in a context and with total assignments symbol family member \(\entails^{\mathrm{tot}}_{\mathrm E}\) & \hyperref[sec:p02-w02-p02-038-many-sorted-reduction]{\ref*{sec:p02-w02-p02-038-many-sorted-reduction}} & meta & \hyperref[def:p02-w02-p02-038-empty-semantics]{\ref*{def:p02-w02-p02-038-empty-semantics}} & \thepage & Quantifies over E-structures (disjoint carriers, possibly empty). Explicit displayed family; member types and common actual first-use exposed, not one atomic glyph. Family listing, not one atomic glyph. Author macro \texttt{\detokenize{\entails\^{}{\Delta}\_{\mathrm E},}\allowbreak\detokenize{\entails\^{}{\mathrm{tot}}\_{\mathrm E}}} not carried to the members. \\
\phantomsection\label{entry:symbol-registry-07003}\(\entails_{\cup}\) & \texttt{\detokenize{\entails_{\cup}}} & overlapping-sort consequence & \hyperref[sec:p02-w02-p02-038-many-sorted-reduction]{\ref*{sec:p02-w02-p02-038-many-sorted-reduction}} & meta & \hyperref[def:p02-w02-p02-038-overlap-semantics]{\ref*{def:p02-w02-p02-038-overlap-semantics}} & \thepage & Quantifies over cup-structures (nonempty carriers, possibly overlapping) and total assignments; language includes cross-sort equality. \\
\phantomsection\label{entry:symbol-registry-07004}\(\mathcal M_0\) & \texttt{\detokenize{\mathcal M_0}} & proof-local witness structures symbol family member \(\mathcal M_0\) & \hyperref[sec:annex-p02-w02-p02-038-many-sorted-reduction]{\ref*{sec:annex-p02-w02-p02-038-many-sorted-reduction}} & n/a & \hyperref[annex:p02-w02-p02-038-inhabited]{\ref*{annex:p02-w02-p02-038-inhabited}} & \thepage & M\_0 empties exactly the sorts without variable-free terms; A\_0 is the covering one-sorted model on the set of sorts. Explicit displayed family; member types and common actual first-use exposed, not one atomic glyph. Family listing, not one atomic glyph. Author macro \texttt{\detokenize{\mathcal M\_0,}\allowbreak\detokenize{\mathcal A\_0}} not carried to the members. \\
\phantomsection\label{entry:symbol-registry-07005}\(\ \mathcal A_0\) & \texttt{\detokenize{\ \mathcal A_0}} & proof-local witness structures symbol family member \(\ \mathcal A_0\) & \hyperref[sec:annex-p02-w02-p02-038-many-sorted-reduction]{\ref*{sec:annex-p02-w02-p02-038-many-sorted-reduction}} & n/a & \hyperref[annex:p02-w02-p02-038-inhabited]{\ref*{annex:p02-w02-p02-038-inhabited}} & \thepage & M\_0 empties exactly the sorts without variable-free terms; A\_0 is the covering one-sorted model on the set of sorts. Explicit displayed family; member types and common actual first-use exposed, not one atomic glyph. Family listing, not one atomic glyph. Author macro \texttt{\detokenize{\mathcal M\_0,}\allowbreak\detokenize{\mathcal A\_0}} not carried to the members. \\
\phantomsection\label{entry:symbol-registry-07006}\(L_{\infty\omega}\) & \texttt{\detokenize{L_{\infty\omega}}} & unrestricted infinitary first-order language with set-sized conjunctions and disjunctions and finite quantifier strings & \hyperref[sec:p02-w02-p02-039-lomega-treatments]{\ref*{sec:p02-w02-p02-039-lomega-treatments}} & n/a & \hyperref[def:p02-w02-p02-039-notation]{\ref*{def:p02-w02-p02-039-notation}} & \thepage & Defined in def:p08-bvb501-potential-isomorphism; used here under that definition with the satisfaction clauses of def:p02-lomega-one-omega without the countability restriction. Not previously registered. \\
\phantomsection\label{entry:symbol-registry-07007}\(\delta_M\) & \texttt{\detokenize{\delta_M}} & characteristic \(L_{\infty\omega}\)-sentence of a structure M up to potential isomorphism & \hyperref[sec:p02-w02-p02-039-lomega-treatments]{\ref*{sec:p02-w02-p02-039-lomega-treatments}} & n/a & \hyperref[def:p02-w02-p02-039-notation]{\ref*{def:p02-w02-p02-039-notation}} & \thepage & The sentence of thm:p08-bvb501-potential-isomorphism. Overlaps with \(\delta_M\) in thm:p03-morley-scott-stabilization, where it is a countable stabilization ordinal; the countable-model sentence there is \(\sigma_M\). \\
\phantomsection\label{entry:symbol-registry-07008}\((\alpha,<)\) & \texttt{\detokenize{(\alpha,}\allowbreak\detokenize{<)}} & an ordinal as a structure in the signature with one binary relation symbol & \hyperref[sec:p02-w02-p02-039-lomega-treatments]{\ref*{sec:p02-w02-p02-039-lomega-treatments}} & n/a & \hyperref[def:p02-w02-p02-039-notation]{\ref*{def:p02-w02-p02-039-notation}} & \thepage & Domain the ordinals below alpha (alpha at least 1), relation the ordinal order. \\
\phantomsection\label{entry:symbol-registry-07009}\(F_\xi\) & \texttt{\detokenize{F_\xi}} & stage xi of the cumulative construction of \(L_{\omega_1\omega}\)-formulas over a fixed signature & \hyperref[sec:annex-p02-w02-p02-039-lomega-treatments]{\ref*{sec:annex-p02-w02-p02-039-lomega-treatments}} & n/a & \hyperref[annex:p02-w02-p02-039-formula-set]{\ref*{annex:p02-w02-p02-039-formula-set}} & \thepage & Proof-local; defined for xi at most omega\_1. F\_{omega\_1} is the set of all formulas. \\
\phantomsection\label{entry:symbol-registry-07010}\(T(\alpha)\) & \texttt{\detokenize{T(\alpha)}} & the set of \(L_{\omega_1\omega}\)-sentences true in the structure on alpha & \hyperref[sec:annex-p02-w02-p02-039-lomega-treatments]{\ref*{sec:annex-p02-w02-p02-039-lomega-treatments}} & n/a & \hyperref[annex:p02-w02-p02-039-proper-fragment]{\ref*{annex:p02-w02-p02-039-proper-fragment}} & \thepage & Proof-local; injective on the class A of ordinals whose characteristic sentence has an \(L_{\omega_1\omega}\) equivalent. \\
\phantomsection\label{entry:symbol-registry-07011}\(\alpha\mathrel{E}\beta\) & \texttt{\detokenize{\mathrel{E}}} & arbitrary equivalence relation on CPL formulas & \hyperref[sec:p02-w02-p02-003-exercise-link]{\ref*{sec:p02-w02-p02-003-exercise-link}} & meta & \hyperref[thm:p02-w02-p02-003-congruence-replacement]{\ref*{thm:p02-w02-p02-003-congruence-replacement}} & \thepage & Metavariable for a relation on \(\mathsf{Fm}_{\mathsf{CPL}}\); logical equivalence \(\equiv\) is one instance and equisatisfiability another. Chosen instead of \(\approx\) or \(\sim\), which the registry already assigns elsewhere. \\
\phantomsection\label{entry:symbol-registry-07012}\([\varphi]\) & \texttt{\detokenize{[\varphi]}} & class of a formula under semantic equivalence & \hyperref[sec:p02-w02-p02-004-equivalent-substitutions]{\ref*{sec:p02-w02-p02-004-equivalent-substitutions}} & n/a & \hyperref[sec:p02-w02-p02-004-equivalent-substitutions]{\ref*{sec:p02-w02-p02-004-equivalent-substitutions}} & \thepage & Equivalence class of a CPL formula modulo the semantic relation of Section sec:prop-syntax-semantics-02; local to this section and distinct from the Henkin term class [t]. \\
\phantomsection\label{entry:symbol-registry-07013}\(A\sim B\vdash C_A\sim C_B\) & \texttt{\detokenize{\sim,}\allowbreak\detokenize{\vdash,}\allowbreak\detokenize{C_A,}\allowbreak\detokenize{C_B}} & Kleene's replacement-theorem notation & \hyperref[sec:p02-w02-p02-004-equivalent-substitutions]{\ref*{sec:p02-w02-p02-004-equivalent-substitutions}} & n/a & \hyperref[sec:p02-w02-p02-004-equivalent-substitutions]{\ref*{sec:p02-w02-p02-004-equivalent-substitutions}} & \thepage & Exact source notation of Kleene 1952, section 26, Theorem 6, p. 116: tilde is Kleene's object-language equivalence connective, turnstile is his deducibility relation, C\_A is a formula with a specified occurrence of A and C\_B the result of replacing that occurrence by B. Used only when reporting that source; not the handbook's semantic equivalence. \\
\phantomsection\label{entry:symbol-registry-07014}\(\#(f)\) & \texttt{\detokenize{\#(f)}} & arity of a function or relation symbol of a first-order language (Barwise's notation) & \hyperref[sec:p02-w02-p02-005-barwise-setup]{\ref*{sec:p02-w02-p02-005-barwise-setup}} & n/a & \hyperref[def:p02-w02-p02-005-bw-language]{\ref*{def:p02-w02-p02-005-bw-language}} & \thepage & Positive integer; Barwise 1977 ch. A.1 p. 17. \\
\phantomsection\label{entry:symbol-registry-07015}\(x^{\mathfrak M}\) & \texttt{\detokenize{x^{\mathfrak M}}} & interpretation in the structure M of the symbol x of L & \hyperref[sec:p02-w02-p02-005-barwise-setup]{\ref*{sec:p02-w02-p02-005-barwise-setup}} & n/a & \hyperref[def:p02-w02-p02-005-bw-structure]{\ref*{def:p02-w02-p02-005-bw-structure}} & \thepage & Barwise writes F(x) for the same value; structure M = (M,F). \\
\phantomsection\label{entry:symbol-registry-07016}\(\mathfrak M\models\varphi[s]\) & \texttt{\detokenize{\models}} & satisfaction of a formula by an assignment in a structure & \hyperref[sec:p02-w02-p02-005-barwise-setup]{\ref*{sec:p02-w02-p02-005-barwise-setup}} & n/a & \hyperref[def:p02-w02-p02-005-bw-satisfaction]{\ref*{def:p02-w02-p02-005-bw-satisfaction}} & \thepage & Barwise's bracket notation for the assignment; same relation as the global satisfaction relation. \\
\phantomsection\label{entry:symbol-registry-07017}\(s(\tfrac{a}{v})\) & \texttt{\detokenize{s(\tfrac{a}{v})}} & assignment agreeing with s except that v is sent to a & \hyperref[sec:p02-w02-p02-005-barwise-setup]{\ref*{sec:p02-w02-p02-005-barwise-setup}} & n/a & \hyperref[def:p02-w02-p02-005-bw-satisfaction]{\ref*{def:p02-w02-p02-005-bw-satisfaction}} & \thepage & Barwise's notation. \\
\phantomsection\label{entry:symbol-registry-07018}\(t^{\mathfrak M}(s)\) & \texttt{\detokenize{t^{\mathfrak M}(s)}} & value of the term t at the assignment s & \hyperref[sec:p02-w02-p02-005-barwise-setup]{\ref*{sec:p02-w02-p02-005-barwise-setup}} & n/a & \hyperref[def:p02-w02-p02-005-bw-satisfaction]{\ref*{def:p02-w02-p02-005-bw-satisfaction}} & \thepage & Function from assignments to the domain. \\
\phantomsection\label{entry:symbol-registry-07019}\(\mathfrak M\equiv\mathfrak N\) & \texttt{\detokenize{\equiv}} & elementary equivalence of structures & \hyperref[sec:p02-w02-p02-005-barwise-setup]{\ref*{sec:p02-w02-p02-005-barwise-setup}} & n/a & \hyperref[def:p02-w02-p02-005-bw-satisfaction]{\ref*{def:p02-w02-p02-005-bw-satisfaction}} & \thepage & Same symbol is used for propositional equivalence of formulas in thm:p02-w02-p02-005-idiom-equivalences; the arguments (structures vs formulas) disambiguate. \\
\phantomsection\label{entry:symbol-registry-07020}\(nx\) & \texttt{\detokenize{nx}} & abbreviated terms: n-fold sum of x and of 1 symbol family member \(nx\) & \hyperref[sec:p02-w02-p02-005-barwise-realm]{\ref*{sec:p02-w02-p02-005-barwise-realm}} & n/a & \hyperref[def:p02-w02-p02-005-bw-groups]{\ref*{def:p02-w02-p02-005-bw-groups}} & \thepage & n at least 1; Barwise ch. A.1 pp. 8, 12. Explicit displayed family; member types and common actual first-use exposed, not one atomic glyph. Family listing, not one atomic glyph. Author macro \texttt{\detokenize{nx,}\allowbreak\detokenize{ n1}} not carried to the members. \\
\phantomsection\label{entry:symbol-registry-07021}\(\ n1\) & \texttt{\detokenize{\ n1}} & abbreviated terms: n-fold sum of x and of 1 symbol family member \(\ n1\) & \hyperref[sec:p02-w02-p02-005-barwise-realm]{\ref*{sec:p02-w02-p02-005-barwise-realm}} & n/a & \hyperref[def:p02-w02-p02-005-bw-groups]{\ref*{def:p02-w02-p02-005-bw-groups}} & \thepage & n at least 1; Barwise ch. A.1 pp. 8, 12. Explicit displayed family; member types and common actual first-use exposed, not one atomic glyph. Family listing, not one atomic glyph. Author macro \texttt{\detokenize{nx,}\allowbreak\detokenize{ n1}} not carried to the members. \\
\phantomsection\label{entry:symbol-registry-07022}\(\mathbb Z_p\) & \texttt{\detokenize{\mathbb Z_p}} & additive group (or ring) of integers modulo the prime p & \hyperref[sec:annex-p02-w02-p02-005-idioms]{\ref*{sec:annex-p02-w02-p02-005-idioms}} & n/a & \hyperref[annex:p02-w02-p02-005-bw-divisible]{\ref*{annex:p02-w02-p02-005-bw-divisible}} & \thepage & Barwise's notation. \\
\phantomsection\label{entry:symbol-registry-07023}\({}^*\mathbb R\) & \texttt{\detokenize{{}^*\mathbb R}} & non-Archimedean ordered field elementarily extending the reals with names & \hyperref[sec:p02-w02-p02-005-barwise-realm]{\ref*{sec:p02-w02-p02-005-barwise-realm}} & n/a & \hyperref[thm:p02-w02-p02-005-bw-nonstandard]{\ref*{thm:p02-w02-p02-005-bw-nonstandard}} & \thepage & Barwise's notation. \\
\phantomsection\label{entry:symbol-registry-07024}\(\underline r\) & \texttt{\detokenize{\underline r}} & constant symbol naming the real number r & \hyperref[sec:annex-p02-w02-p02-005-idioms]{\ref*{sec:annex-p02-w02-p02-005-idioms}} & n/a & \hyperref[annex:p02-w02-p02-005-bw-nonstandard]{\ref*{annex:p02-w02-p02-005-bw-nonstandard}} & \thepage & Proof-local; Barwise uses r as its own name. \\
\phantomsection\label{entry:symbol-registry-07025}\(N_\psi\) & \texttt{\detokenize{N_\psi}} & bound beyond which a sentence true in characteristic 0 holds in characteristic p & \hyperref[sec:p02-w02-p02-005-barwise-realm]{\ref*{sec:p02-w02-p02-005-barwise-realm}} & n/a & \hyperref[thm:p02-w02-p02-005-bw-characteristic]{\ref*{thm:p02-w02-p02-005-bw-characteristic}} & \thepage & Barwise's notation. \\
\phantomsection\label{entry:symbol-registry-07026}\(\mathbf t\) & \texttt{\detokenize{\mathbf t}} & the two truth values of a truth assignment (Barwise's notation) symbol family member \(\mathbf t\) & \hyperref[sec:p02-w02-p02-005-bw-prop]{\ref*{sec:p02-w02-p02-005-bw-prop}} & n/a & \hyperref[def:p02-w02-p02-005-bw-truth-assignment]{\ref*{def:p02-w02-p02-005-bw-truth-assignment}} & \thepage & Barwise 1977 ch. A.1 p. 24; the idiom section of the same file uses 0 and 1. Explicit displayed family; member types and common actual first-use exposed, not one atomic glyph. Family listing, not one atomic glyph. Author macro \texttt{\detokenize{\mathbf t,}\allowbreak\detokenize{ \mathbf f}} not carried to the members. \\
\phantomsection\label{entry:symbol-registry-07027}\(\ \mathbf f\) & \texttt{\detokenize{\ \mathbf f}} & the two truth values of a truth assignment (Barwise's notation) symbol family member \(\ \mathbf f\) & \hyperref[sec:p02-w02-p02-005-bw-prop]{\ref*{sec:p02-w02-p02-005-bw-prop}} & n/a & \hyperref[def:p02-w02-p02-005-bw-truth-assignment]{\ref*{def:p02-w02-p02-005-bw-truth-assignment}} & \thepage & Barwise 1977 ch. A.1 p. 24; the idiom section of the same file uses 0 and 1. Explicit displayed family; member types and common actual first-use exposed, not one atomic glyph. Family listing, not one atomic glyph. Author macro \texttt{\detokenize{\mathbf t,}\allowbreak\detokenize{ \mathbf f}} not carried to the members. \\
\phantomsection\label{entry:symbol-registry-07028}\(\bar v\) & \texttt{\detokenize{\bar v}} & extension of a truth assignment v to all propositional formulas & \hyperref[sec:p02-w02-p02-005-bw-prop]{\ref*{sec:p02-w02-p02-005-bw-prop}} & n/a & \hyperref[def:p02-w02-p02-005-bw-truth-assignment]{\ref*{def:p02-w02-p02-005-bw-truth-assignment}} & \thepage & Barwise 1977 ch. A.1 p. 24. \\
\phantomsection\label{entry:symbol-registry-07029}\(L(C)\) & \texttt{\detokenize{L(C)}} & witnessing expansion of the language L by the witnessing constants C & \hyperref[sec:p02-w02-p02-005-bw-henkin]{\ref*{sec:p02-w02-p02-005-bw-henkin}} & n/a & \hyperref[def:p02-w02-p02-005-bw-witnessing]{\ref*{def:p02-w02-p02-005-bw-witnessing}} & \thepage & With C\_n and L\_n = L union C\_n; Barwise p. 29. \\
\phantomsection\label{entry:symbol-registry-07030}\(c_{\varphi(v)}\) & \texttt{\detokenize{c_{\varphi(v)}}} & witnessing constant attached to the formula phi(v) & \hyperref[sec:p02-w02-p02-005-bw-henkin]{\ref*{sec:p02-w02-p02-005-bw-henkin}} & n/a & \hyperref[def:p02-w02-p02-005-bw-witnessing]{\ref*{def:p02-w02-p02-005-bw-witnessing}} & \thepage & Read for FV(phi) contained in {v}; Barwise requires exactly one free variable. \\
\phantomsection\label{entry:symbol-registry-07031}\(T_{\mathrm{Henkin}}\) & \texttt{\detokenize{T_{\mathrm{Henkin}}}} & set of Henkin axioms (types I and II) and quantifier axioms (types III and IV) & \hyperref[sec:p02-w02-p02-005-bw-henkin]{\ref*{sec:p02-w02-p02-005-bw-henkin}} & n/a & \hyperref[def:p02-w02-p02-005-bw-witnessing]{\ref*{def:p02-w02-p02-005-bw-witnessing}} & \thepage & Barwise Def. 4.6, p. 30. \\
\phantomsection\label{entry:symbol-registry-07032}\(\mathrm{Eq}\) & \texttt{\detokenize{\mathrm{Eq}}} & set of variable-free equality axioms of L(C) & \hyperref[sec:p02-w02-p02-005-bw-henkin]{\ref*{sec:p02-w02-p02-005-bw-henkin}} & n/a & \hyperref[def:p02-w02-p02-005-bw-witnessing]{\ref*{def:p02-w02-p02-005-bw-witnessing}} & \thepage & Barwise p. 31. \\
\phantomsection\label{entry:symbol-registry-07033}\(T\vdash\varphi\) & \texttt{\detokenize{\vdash}} & provability of a formula from a set of sentences in Barwise's system H & \hyperref[sec:p02-w02-p02-005-bw-hilbert]{\ref*{sec:p02-w02-p02-005-bw-hilbert}} & n/a & \hyperref[def:p02-w02-p02-005-bw-system-h]{\ref*{def:p02-w02-p02-005-bw-system-h}} & \thepage & Same glyph is the sequent arrow of G and G+ (def:p02-w02-p02-005-bw-sequents); left argument a set of sentences vs a finite set of formulas in a sequent disambiguates. Original precedence descriptor 'infix relation' retained; no rank inferred. \\
\phantomsection\label{entry:symbol-registry-07034}\(\Gamma\vdash\Delta\) & \texttt{\detokenize{\vdash}} & sequent of finite sets of formulas in Barwise's systems G and G+ & \hyperref[sec:p02-w02-p02-005-bw-gentzen]{\ref*{sec:p02-w02-p02-005-bw-gentzen}} & infix & \hyperref[def:p02-w02-p02-005-bw-sequents]{\ref*{def:p02-w02-p02-005-bw-sequents}} & \thepage & Barwise p. 37. \\
\phantomsection\label{entry:symbol-registry-07035}\(\mathbf G\) & \texttt{\detokenize{\mathbf G}} & Barwise's Gentzen-type systems without and with cut symbol family member \(\mathbf G\) & \hyperref[sec:p02-w02-p02-005-bw-gentzen]{\ref*{sec:p02-w02-p02-005-bw-gentzen}} & n/a & \hyperref[def:p02-w02-p02-005-bw-system-g]{\ref*{def:p02-w02-p02-005-bw-system-g}} & \thepage & Barwise prints G with a superscript plus (p. 39). Explicit displayed family; member types and common actual first-use exposed, not one atomic glyph. Family listing, not one atomic glyph. Author macro \texttt{\detokenize{\mathbf G,}\allowbreak\detokenize{ \mathbf G\^{}{+}}} not carried to the members. \\
\phantomsection\label{entry:symbol-registry-07036}\(\ \mathbf G^{+}\) & \texttt{\detokenize{\ \mathbf G^{+}}} & Barwise's Gentzen-type systems without and with cut symbol family member \(\ \mathbf G^{+}\) & \hyperref[sec:p02-w02-p02-005-bw-gentzen]{\ref*{sec:p02-w02-p02-005-bw-gentzen}} & n/a & \hyperref[def:p02-w02-p02-005-bw-system-g]{\ref*{def:p02-w02-p02-005-bw-system-g}} & \thepage & Barwise prints G with a superscript plus (p. 39). Explicit displayed family; member types and common actual first-use exposed, not one atomic glyph. Family listing, not one atomic glyph. Author macro \texttt{\detokenize{\mathbf G,}\allowbreak\detokenize{ \mathbf G\^{}{+}}} not carried to the members. \\
\phantomsection\label{entry:symbol-registry-07037}\(\mathrm{HF}(\mathfrak M)\) & \texttt{\detokenize{\mathrm{HF}(\mathfrak M)}} & structure of hereditarily finite sets over the structure M & \hyperref[sec:p02-w02-p02-005-bw-beyond]{\ref*{sec:p02-w02-p02-005-bw-beyond}} & n/a & \hyperref[def:p02-w02-p02-005-bw-weak-so]{\ref*{def:p02-w02-p02-005-bw-weak-so}} & \thepage & With HF\_n(M) and HF(M); Barwise 5.3, p. 43. \\
\phantomsection\label{entry:symbol-registry-07038}\(L^{*}\) & \texttt{\detokenize{L^{*}}} & two-sorted language of weak second-order logic over L with set variables and membership & \hyperref[sec:p02-w02-p02-005-bw-beyond]{\ref*{sec:p02-w02-p02-005-bw-beyond}} & n/a & \hyperref[def:p02-w02-p02-005-bw-weak-so]{\ref*{def:p02-w02-p02-005-bw-weak-so}} & \thepage & Barwise 5.3, p. 43. \\
\phantomsection\label{entry:symbol-registry-07039}\(L_{\omega_1\omega}\) & \texttt{\detokenize{L_{\omega_1\omega}}} & infinitary logic with countable conjunctions and disjunctions and finite quantifier strings & \hyperref[sec:p02-w02-p02-005-bw-beyond]{\ref*{sec:p02-w02-p02-005-bw-beyond}} & n/a & \hyperref[def:p02-w02-p02-005-bw-infinitary]{\ref*{def:p02-w02-p02-005-bw-infinitary}} & \thepage & Barwise 5.4, p. 43. \\
\phantomsection\label{entry:symbol-registry-07040}\(Q_0\) & \texttt{\detokenize{Q_0}} & quantifiers 'infinitely many' and 'uncountably many' symbol family member \(Q_0\) & \hyperref[sec:p02-w02-p02-005-bw-beyond]{\ref*{sec:p02-w02-p02-005-bw-beyond}} & n/a & \hyperref[def:p02-w02-p02-005-bw-infinitary]{\ref*{def:p02-w02-p02-005-bw-infinitary}} & \thepage & Barwise writes L(Q\_0) and L(Q) (p. 45); L(Q\_1) is a common later name for L(Q). Original precedence descriptor 'binder' retained; no rank inferred. Explicit displayed family; member types and common actual first-use exposed, not one atomic glyph. Family listing, not one atomic glyph. Author macro \texttt{\detokenize{Q\_0,}\allowbreak\detokenize{ Q}} not carried to the members. \\
\phantomsection\label{entry:symbol-registry-07041}\(\ Q\) & \texttt{\detokenize{\ Q}} & quantifiers 'infinitely many' and 'uncountably many' symbol family member \(\ Q\) & \hyperref[sec:p02-w02-p02-005-bw-beyond]{\ref*{sec:p02-w02-p02-005-bw-beyond}} & n/a & \hyperref[def:p02-w02-p02-005-bw-infinitary]{\ref*{def:p02-w02-p02-005-bw-infinitary}} & \thepage & Barwise writes L(Q\_0) and L(Q) (p. 45); L(Q\_1) is a common later name for L(Q). Original precedence descriptor 'binder' retained; no rank inferred. Explicit displayed family; member types and common actual first-use exposed, not one atomic glyph. Family listing, not one atomic glyph. Author macro \texttt{\detokenize{Q\_0,}\allowbreak\detokenize{ Q}} not carried to the members. \\
\phantomsection\label{entry:symbol-registry-07042}\(\mathrm{supp}(b)\) & \texttt{\detokenize{\mathrm{supp}(b)}} & set of atoms in the transitive closure of a hereditarily finite set b & \hyperref[sec:annex-p02-w02-p02-005-bw-s5]{\ref*{sec:annex-p02-w02-p02-005-bw-s5}} & n/a & \hyperref[annex:p02-w02-p02-005-bw-weak-so-ls]{\ref*{annex:p02-w02-p02-005-bw-weak-so-ls}} & \thepage & Proof-local. \\
\phantomsection\label{entry:symbol-registry-07043}\(\mathsf{CONS}\) & \texttt{\detokenize{\mathsf{CONS}}} & finite-premise consequence and equivalence languages symbol family member \(\mathsf{CONS}\) & \hyperref[sec:p02-w02-p02-018-decision-certificates]{\ref*{sec:p02-w02-p02-018-decision-certificates}} & n/a & \hyperref[def:p02-w02-p02-018-input]{\ref*{def:p02-w02-p02-018-input}} & \thepage & Sets of encoded bit strings under the declared token encoding; companions of the registered SAT and TAUT languages; malformed strings belong to neither. Explicit displayed family; member types and common actual first-use exposed, not one atomic glyph. Family listing, not one atomic glyph. Author macro \texttt{\detokenize{\mathsf{CONS},}\allowbreak\detokenize{\mathsf{EQV}}} not carried to the members. \\
\phantomsection\label{entry:symbol-registry-07044}\(\mathsf{EQV}\) & \texttt{\detokenize{\mathsf{EQV}}} & finite-premise consequence and equivalence languages symbol family member \(\mathsf{EQV}\) & \hyperref[sec:p02-w02-p02-018-decision-certificates]{\ref*{sec:p02-w02-p02-018-decision-certificates}} & n/a & \hyperref[def:p02-w02-p02-018-input]{\ref*{def:p02-w02-p02-018-input}} & \thepage & Sets of encoded bit strings under the declared token encoding; companions of the registered SAT and TAUT languages; malformed strings belong to neither. Explicit displayed family; member types and common actual first-use exposed, not one atomic glyph. Family listing, not one atomic glyph. Author macro \texttt{\detokenize{\mathsf{CONS},}\allowbreak\detokenize{\mathsf{EQV}}} not carried to the members. \\
\phantomsection\label{entry:symbol-registry-07045}\(m\) & \texttt{\detokenize{m}} & input bit length and number of distinct variables of an instance symbol family member \(m\) & \hyperref[sec:p02-w02-p02-018-decision-certificates]{\ref*{sec:p02-w02-p02-018-decision-certificates}} & n/a & \hyperref[def:p02-w02-p02-018-input]{\ref*{def:p02-w02-p02-018-input}} & \thepage & Size measures of the truth-table bound; k is at most m/12. Explicit displayed family; member types and common actual first-use exposed, not one atomic glyph. Family listing, not one atomic glyph. Author macro \texttt{\detokenize{m,}\allowbreak\detokenize{k}} not carried to the members. \\
\phantomsection\label{entry:symbol-registry-07046}\(k\) & \texttt{\detokenize{k}} & input bit length and number of distinct variables of an instance symbol family member \(k\) & \hyperref[sec:p02-w02-p02-018-decision-certificates]{\ref*{sec:p02-w02-p02-018-decision-certificates}} & n/a & \hyperref[def:p02-w02-p02-018-input]{\ref*{def:p02-w02-p02-018-input}} & \thepage & Size measures of the truth-table bound; k is at most m/12. Explicit displayed family; member types and common actual first-use exposed, not one atomic glyph. Family listing, not one atomic glyph. Author macro \texttt{\detokenize{m,}\allowbreak\detokenize{k}} not carried to the members. \\
\phantomsection\label{entry:symbol-registry-07047}\(\tau\) & \texttt{\detokenize{\tau}} & valuation table and the valuation it determines symbol family member \(\tau\) & \hyperref[sec:p02-w02-p02-018-decision-certificates]{\ref*{sec:p02-w02-p02-018-decision-certificates}} & n/a & \hyperref[def:p02-w02-p02-018-evaluator]{\ref*{def:p02-w02-p02-018-evaluator}} & \thepage & A finite list of numeral and truth-value pairs; the valuation is false off the table. Explicit displayed family; member types and common actual first-use exposed, not one atomic glyph. Family listing, not one atomic glyph. Author macro \texttt{\detokenize{\tau,}\allowbreak\detokenize{v\_\tau}} not carried to the members. \\
\phantomsection\label{entry:symbol-registry-07048}\(v_\tau\) & \texttt{\detokenize{v_\tau}} & valuation table and the valuation it determines symbol family member \(v_\tau\) & \hyperref[sec:p02-w02-p02-018-decision-certificates]{\ref*{sec:p02-w02-p02-018-decision-certificates}} & n/a & \hyperref[def:p02-w02-p02-018-evaluator]{\ref*{def:p02-w02-p02-018-evaluator}} & \thepage & A finite list of numeral and truth-value pairs; the valuation is false off the table. Explicit displayed family; member types and common actual first-use exposed, not one atomic glyph. Family listing, not one atomic glyph. Author macro \texttt{\detokenize{\tau,}\allowbreak\detokenize{v\_\tau}} not carried to the members. \\
\phantomsection\label{entry:symbol-registry-07049}\(\operatorname{red}(\sigma)\) & \texttt{\detokenize{\operatorname{red}}} & stack reduction of negation markers & \hyperref[annex:p02-w02-p02-018-evaluator]{\ref*{annex:p02-w02-p02-018-evaluator}} & n/a & \hyperref[annex:p02-w02-p02-018-evaluator]{\ref*{annex:p02-w02-p02-018-evaluator}} & \thepage & Repeatedly replaces a negation marker followed by a truth value by the opposite value; metalinguistic bookkeeping of the evaluator. \\
\phantomsection\label{entry:symbol-registry-07050}\(\Gamma_A\) & \texttt{\detokenize{\Gamma_A}} & premise set of the variables indexed by A & \hyperref[sec:p02-w02-p02-018-decision-certificates]{\ref*{sec:p02-w02-p02-018-decision-certificates}} & n/a & \hyperref[thm:p02-w02-p02-018-limits]{\ref*{thm:p02-w02-p02-018-limits}} & \thepage & For A a subset of the natural numbers, the set of p\_n with n in A. \\
\phantomsection\label{entry:symbol-registry-07051}\(\langle x,y\rangle\) & \texttt{\detokenize{\langle x,}\allowbreak\detokenize{y\rangle}} & bit-string pairing for verifiers & \hyperref[sec:p02-w02-p02-018-complexity]{\ref*{sec:p02-w02-p02-018-complexity}} & n/a & \hyperref[def:p02-w02-p02-018-verifier]{\ref*{def:p02-w02-p02-018-verifier}} & \thepage & Doubled bits of x, then 01, then y; injective and decoded in one scan. \\
\phantomsection\label{entry:symbol-registry-07052}\(L^{\mathsf c}\) & \texttt{\detokenize{L^{\mathsf c}}} & complement of a language in the set of bit strings & \hyperref[sec:p02-w02-p02-018-complexity]{\ref*{sec:p02-w02-p02-018-complexity}} & n/a & \hyperref[def:p02-w02-p02-018-verifier]{\ref*{def:p02-w02-p02-018-verifier}} & \thepage & The set of bit strings not in L; chosen instead of an overline because the overline is the registered literal complement. \\
\phantomsection\label{entry:symbol-registry-07053}\(h\) & \texttt{\detokenize{h}} & negation reduction map & \hyperref[sec:p02-w02-p02-018-complexity]{\ref*{sec:p02-w02-p02-018-complexity}} & n/a & \hyperref[thm:p02-w02-p02-018-negation]{\ref*{thm:p02-w02-p02-018-negation}} & \thepage & Maps a formula encoding to the encoding of its negation and malformed strings to the encoding of negated falsum. \\
\phantomsection\label{entry:symbol-registry-07054}\(\varphi_x\) & \texttt{\detokenize{\varphi_x}} & computation formula of a verifier on input x & \hyperref[sec:p02-w02-p02-018-complexity]{\ref*{sec:p02-w02-p02-018-complexity}} & n/a & \hyperref[thm:p02-w02-p02-018-cook-levin]{\ref*{thm:p02-w02-p02-018-cook-levin}} & \thepage & Satisfiable exactly when x is in the NP language. \\
\phantomsection\label{entry:symbol-registry-07055}\(S_{t,s}\) & \texttt{\detokenize{S_{t,s}}} & state head and cell variables of the computation formula symbol family member \(S_{t,s}\) & \hyperref[annex:p02-w02-p02-018-cook-levin]{\ref*{annex:p02-w02-p02-018-cook-levin}} & n/a & \hyperref[annex:p02-w02-p02-018-cook-levin]{\ref*{annex:p02-w02-p02-018-cook-levin}} & \thepage & Propositional variables with distinct mixed-radix indices; local to the proof. Explicit displayed family; member types and common actual first-use exposed, not one atomic glyph. Family listing, not one atomic glyph. Author macro \texttt{\detokenize{S\_{t,s},}\allowbreak\detokenize{H\_{t,j,i},}\allowbreak\detokenize{C\_{t,j,i,a}}} not carried to the members. \\
\phantomsection\label{entry:symbol-registry-07056}\(H_{t,j,i}\) & \texttt{\detokenize{H_{t,j,i}}} & state head and cell variables of the computation formula symbol family member \(H_{t,j,i}\) & \hyperref[annex:p02-w02-p02-018-cook-levin]{\ref*{annex:p02-w02-p02-018-cook-levin}} & n/a & \hyperref[annex:p02-w02-p02-018-cook-levin]{\ref*{annex:p02-w02-p02-018-cook-levin}} & \thepage & Propositional variables with distinct mixed-radix indices; local to the proof. Explicit displayed family; member types and common actual first-use exposed, not one atomic glyph. Family listing, not one atomic glyph. Author macro \texttt{\detokenize{S\_{t,s},}\allowbreak\detokenize{H\_{t,j,i},}\allowbreak\detokenize{C\_{t,j,i,a}}} not carried to the members. \\
\phantomsection\label{entry:symbol-registry-07057}\(C_{t,j,i,a}\) & \texttt{\detokenize{C_{t,j,i,a}}} & state head and cell variables of the computation formula symbol family member \(C_{t,j,i,a}\) & \hyperref[annex:p02-w02-p02-018-cook-levin]{\ref*{annex:p02-w02-p02-018-cook-levin}} & n/a & \hyperref[annex:p02-w02-p02-018-cook-levin]{\ref*{annex:p02-w02-p02-018-cook-levin}} & \thepage & Propositional variables with distinct mixed-radix indices; local to the proof. Explicit displayed family; member types and common actual first-use exposed, not one atomic glyph. Family listing, not one atomic glyph. Author macro \texttt{\detokenize{S\_{t,s},}\allowbreak\detokenize{H\_{t,j,i},}\allowbreak\detokenize{C\_{t,j,i,a}}} not carried to the members. \\
\phantomsection\label{entry:symbol-registry-07058}\(\operatorname{One}(\chi_1,\ldots,\chi_r)\) & \texttt{\detokenize{\operatorname{One}}} & exactly-one formula & \hyperref[annex:p02-w02-p02-018-cook-levin]{\ref*{annex:p02-w02-p02-018-cook-levin}} & n/a & \hyperref[annex:p02-w02-p02-018-cook-levin]{\ref*{annex:p02-w02-p02-018-cook-levin}} & \thepage & Disjunction of the arguments conjoined with all pairwise negated conjunctions. \\
\phantomsection\label{entry:symbol-registry-07059}\(\theta_1\) & \texttt{\detokenize{\theta_1}} & displayed most general unifiers of the annotated refutation symbol family member \(\theta_1\) & \hyperref[fig:p02-w02-p02-018-fol-dag]{\ref*{fig:p02-w02-p02-018-fol-dag}} & n/a & \hyperref[fig:p02-w02-p02-018-fol-dag]{\ref*{fig:p02-w02-p02-018-fol-dag}} & \thepage & Term substitutions in the sense of the first-order resolution section; theta\_3 is the identity. Explicit displayed family; member types and common actual first-use exposed, not one atomic glyph. Family listing, not one atomic glyph. Author macro \texttt{\detokenize{\theta\_i}} not carried to the members. \\
\phantomsection\label{entry:symbol-registry-07060}\(\theta_2\) & \texttt{\detokenize{\theta_2}} & displayed most general unifiers of the annotated refutation symbol family member \(\theta_2\) & \hyperref[fig:p02-w02-p02-018-fol-dag]{\ref*{fig:p02-w02-p02-018-fol-dag}} & n/a & \hyperref[fig:p02-w02-p02-018-fol-dag]{\ref*{fig:p02-w02-p02-018-fol-dag}} & \thepage & Term substitutions in the sense of the first-order resolution section; theta\_3 is the identity. Explicit displayed family; member types and common actual first-use exposed, not one atomic glyph. Family listing, not one atomic glyph. Author macro \texttt{\detokenize{\theta\_i}} not carried to the members. \\
\phantomsection\label{entry:symbol-registry-07061}\(\theta_3\) & \texttt{\detokenize{\theta_3}} & displayed most general unifiers of the annotated refutation symbol family member \(\theta_3\) & \hyperref[fig:p02-w02-p02-018-fol-dag]{\ref*{fig:p02-w02-p02-018-fol-dag}} & n/a & \hyperref[fig:p02-w02-p02-018-fol-dag]{\ref*{fig:p02-w02-p02-018-fol-dag}} & \thepage & Term substitutions in the sense of the first-order resolution section; theta\_3 is the identity. Explicit displayed family; member types and common actual first-use exposed, not one atomic glyph. Family listing, not one atomic glyph. Author macro \texttt{\detokenize{\theta\_i}} not carried to the members. \\
\phantomsection\label{entry:symbol-registry-07062}\(b'\) & \texttt{\detokenize{b'}} & opposite truth value & \hyperref[annex:p02-w02-p02-018-evaluator]{\ref*{annex:p02-w02-p02-018-evaluator}} & n/a & \hyperref[annex:p02-w02-p02-018-evaluator]{\ref*{annex:p02-w02-p02-018-evaluator}} & \thepage & The truth value other than b; local to the evaluator proof. \\
\phantomsection\label{entry:symbol-registry-07063}\(N\) & \texttt{\detokenize{N}} & number of literal occurrences of a clause list & \hyperref[sec:p02-w02-p02-019-sat-procedures]{\ref*{sec:p02-w02-p02-019-sat-procedures}} & n/a & \hyperref[def:p02-w02-p02-019-clause-input]{\ref*{def:p02-w02-p02-019-clause-input}} & \thepage & Same meaning as the N of def:p02-tractable-classes; size of a clause list is m+N. \\
\phantomsection\label{entry:symbol-registry-07064}\(m\) & \texttt{\detokenize{m}} & number of clauses and of occurring variables of a clause list symbol family member \(m\) & \hyperref[sec:p02-w02-p02-019-inputs]{\ref*{sec:p02-w02-p02-019-inputs}} & n/a & \hyperref[def:p02-w02-p02-019-clause-input]{\ref*{def:p02-w02-p02-019-clause-input}} & \thepage & Used with these roles in sec:p02-w02-p02-019-horn and sec:p02-w02-p02-019-two-sat and their annex items; in sec:p02-w02-p02-019-complexity n is instead the bit length of an input as in def:p02-complexity-standing. Explicit displayed family; member types and common actual first-use exposed, not one atomic glyph. Family listing, not one atomic glyph. Author macro \texttt{\detokenize{m,}\allowbreak\detokenize{n}} not carried to the members. \\
\phantomsection\label{entry:symbol-registry-07065}\(n\) & \texttt{\detokenize{n}} & number of clauses and of occurring variables of a clause list symbol family member \(n\) & \hyperref[sec:p02-w02-p02-019-inputs]{\ref*{sec:p02-w02-p02-019-inputs}} & n/a & \hyperref[def:p02-w02-p02-019-clause-input]{\ref*{def:p02-w02-p02-019-clause-input}} & \thepage & Used with these roles in sec:p02-w02-p02-019-horn and sec:p02-w02-p02-019-two-sat and their annex items; in sec:p02-w02-p02-019-complexity n is instead the bit length of an input as in def:p02-complexity-standing. Explicit displayed family; member types and common actual first-use exposed, not one atomic glyph. Family listing, not one atomic glyph. Author macro \texttt{\detokenize{m,}\allowbreak\detokenize{n}} not carried to the members. \\
\phantomsection\label{entry:symbol-registry-07066}\(A_S\) & \texttt{\detokenize{A_S}} & set of variables occurring in a clause list S & \hyperref[sec:p02-w02-p02-019-sat-procedures]{\ref*{sec:p02-w02-p02-019-sat-procedures}} & n/a & \hyperref[def:p02-w02-p02-019-clause-input]{\ref*{def:p02-w02-p02-019-clause-input}} & \thepage & Assumed renamed p\_1..p\_n in order of first occurrence. \\
\phantomsection\label{entry:symbol-registry-07067}\(B(C)\) & \texttt{\detokenize{B(C)}} & body and head of a Horn clause symbol family member \(B(C)\) & \hyperref[sec:p02-w02-p02-019-horn]{\ref*{sec:p02-w02-p02-019-horn}} & n/a & \hyperref[def:p02-w02-p02-019-horn-parts]{\ref*{def:p02-w02-p02-019-horn-parts}} & \thepage & B(C) is the set of atoms q with \(\neg q\in C\); the head is the unique positive literal of a definite clause. Explicit displayed family; member types and common actual first-use exposed, not one atomic glyph. Family listing, not one atomic glyph. Author macro \texttt{\detokenize{B,}\allowbreak\detokenize{\operatorname{hd}}} not carried to the members. \\
\phantomsection\label{entry:symbol-registry-07068}\(\operatorname{hd}(C)\) & \texttt{\detokenize{\operatorname{hd}(C)}} & body and head of a Horn clause symbol family member \(\operatorname{hd}(C)\) & \hyperref[sec:p02-w02-p02-019-horn]{\ref*{sec:p02-w02-p02-019-horn}} & n/a & \hyperref[def:p02-w02-p02-019-horn-parts]{\ref*{def:p02-w02-p02-019-horn-parts}} & \thepage & B(C) is the set of atoms q with \(\neg q\in C\); the head is the unique positive literal of a definite clause. Explicit displayed family; member types and common actual first-use exposed, not one atomic glyph. Family listing, not one atomic glyph. Author macro \texttt{\detokenize{B,}\allowbreak\detokenize{\operatorname{hd}}} not carried to the members. \\
\phantomsection\label{entry:symbol-registry-07069}\(v_X\) & \texttt{\detokenize{v_X}} & valuation making exactly the members of X true & \hyperref[sec:p02-w02-p02-019-horn]{\ref*{sec:p02-w02-p02-019-horn}} & n/a & \hyperref[def:p02-w02-p02-019-horn-parts]{\ref*{def:p02-w02-p02-019-horn-parts}} & \thepage & X a set of variables; v\_F for F the forward closure is the least model. \\
\phantomsection\label{entry:symbol-registry-07070}\(\operatorname{Fwd}(S)\) & \texttt{\detokenize{\operatorname{Fwd}}} & least forward closure of a Horn list & \hyperref[sec:p02-w02-p02-019-horn]{\ref*{sec:p02-w02-p02-019-horn}} & n/a & \hyperref[def:p02-w02-p02-019-forward-closure]{\ref*{def:p02-w02-p02-019-forward-closure}} & \thepage & Intersection of the subsets of A\_S closed under S; same role as S\_* of annex:p14.dc.horn for the circuit instance. \\
\phantomsection\label{entry:symbol-registry-07071}\(\mathit{mark}\) & \texttt{\detokenize{\mathit{mark}}} & arrays and queue of the forward-chaining procedure symbol family member \(\mathit{mark}\) & \hyperref[sec:p02-w02-p02-019-horn]{\ref*{sec:p02-w02-p02-019-horn}} & n/a & \hyperref[def:p02-w02-p02-019-horn-method]{\ref*{def:p02-w02-p02-019-horn-method}} & \thepage & Program variables; the array \(\mathit{hd}\) stores the index of \(\operatorname{hd}(C_j)\) or 0. M and D in annex:p02-w02-p02-019-hp are the marked and removed atom sets. Explicit displayed family; member types and common actual first-use exposed, not one atomic glyph. Family listing, not one atomic glyph. Author macro \texttt{\detokenize{\mathit{mark},}\allowbreak\detokenize{\mathit{occ},}\allowbreak\detokenize{\mathit{cnt},}\allowbreak\detokenize{\mathit{hd},}\allowbreak\detokenize{Q}} not carried to the members. \\
\phantomsection\label{entry:symbol-registry-07072}\(\mathit{occ}\) & \texttt{\detokenize{\mathit{occ}}} & arrays and queue of the forward-chaining procedure symbol family member \(\mathit{occ}\) & \hyperref[sec:p02-w02-p02-019-horn]{\ref*{sec:p02-w02-p02-019-horn}} & n/a & \hyperref[def:p02-w02-p02-019-horn-method]{\ref*{def:p02-w02-p02-019-horn-method}} & \thepage & Program variables; the array \(\mathit{hd}\) stores the index of \(\operatorname{hd}(C_j)\) or 0. M and D in annex:p02-w02-p02-019-hp are the marked and removed atom sets. Explicit displayed family; member types and common actual first-use exposed, not one atomic glyph. Family listing, not one atomic glyph. Author macro \texttt{\detokenize{\mathit{mark},}\allowbreak\detokenize{\mathit{occ},}\allowbreak\detokenize{\mathit{cnt},}\allowbreak\detokenize{\mathit{hd},}\allowbreak\detokenize{Q}} not carried to the members. \\
\phantomsection\label{entry:symbol-registry-07073}\(\mathit{cnt}\) & \texttt{\detokenize{\mathit{cnt}}} & arrays and queue of the forward-chaining procedure symbol family member \(\mathit{cnt}\) & \hyperref[sec:p02-w02-p02-019-horn]{\ref*{sec:p02-w02-p02-019-horn}} & n/a & \hyperref[def:p02-w02-p02-019-horn-method]{\ref*{def:p02-w02-p02-019-horn-method}} & \thepage & Program variables; the array \(\mathit{hd}\) stores the index of \(\operatorname{hd}(C_j)\) or 0. M and D in annex:p02-w02-p02-019-hp are the marked and removed atom sets. Explicit displayed family; member types and common actual first-use exposed, not one atomic glyph. Family listing, not one atomic glyph. Author macro \texttt{\detokenize{\mathit{mark},}\allowbreak\detokenize{\mathit{occ},}\allowbreak\detokenize{\mathit{cnt},}\allowbreak\detokenize{\mathit{hd},}\allowbreak\detokenize{Q}} not carried to the members. \\
\phantomsection\label{entry:symbol-registry-07074}\(\mathit{hd}\) & \texttt{\detokenize{\mathit{hd}}} & arrays and queue of the forward-chaining procedure symbol family member \(\mathit{hd}\) & \hyperref[sec:p02-w02-p02-019-horn]{\ref*{sec:p02-w02-p02-019-horn}} & n/a & \hyperref[def:p02-w02-p02-019-horn-method]{\ref*{def:p02-w02-p02-019-horn-method}} & \thepage & Program variables; the array \(\mathit{hd}\) stores the index of \(\operatorname{hd}(C_j)\) or 0. M and D in annex:p02-w02-p02-019-hp are the marked and removed atom sets. Explicit displayed family; member types and common actual first-use exposed, not one atomic glyph. Family listing, not one atomic glyph. Author macro \texttt{\detokenize{\mathit{mark},}\allowbreak\detokenize{\mathit{occ},}\allowbreak\detokenize{\mathit{cnt},}\allowbreak\detokenize{\mathit{hd},}\allowbreak\detokenize{Q}} not carried to the members. \\
\phantomsection\label{entry:symbol-registry-07075}\(Q\) & \texttt{\detokenize{Q}} & arrays and queue of the forward-chaining procedure symbol family member \(Q\) & \hyperref[sec:p02-w02-p02-019-horn]{\ref*{sec:p02-w02-p02-019-horn}} & n/a & \hyperref[def:p02-w02-p02-019-horn-method]{\ref*{def:p02-w02-p02-019-horn-method}} & \thepage & Program variables; the array \(\mathit{hd}\) stores the index of \(\operatorname{hd}(C_j)\) or 0. M and D in annex:p02-w02-p02-019-hp are the marked and removed atom sets. Explicit displayed family; member types and common actual first-use exposed, not one atomic glyph. Family listing, not one atomic glyph. Author macro \texttt{\detokenize{\mathit{mark},}\allowbreak\detokenize{\mathit{occ},}\allowbreak\detokenize{\mathit{cnt},}\allowbreak\detokenize{\mathit{hd},}\allowbreak\detokenize{Q}} not carried to the members. \\
\phantomsection\label{entry:symbol-registry-07076}\(G_S\) & \texttt{\detokenize{G_S}} & implication graph of a 2-CNF list and its reverse symbol family member \(G_S\) & \hyperref[sec:p02-w02-p02-019-two-sat]{\ref*{sec:p02-w02-p02-019-two-sat}} & n/a & \hyperref[def:p02-w02-p02-019-implication-graph]{\ref*{def:p02-w02-p02-019-implication-graph}} & \thepage & 2n literal vertices and 2m edges counted with repetition; G\^{}R reverses every edge. Explicit displayed family; member types and common actual first-use exposed, not one atomic glyph. Family listing, not one atomic glyph. Author macro \texttt{\detokenize{G\_S,}\allowbreak\detokenize{G\_S\^{}{\mathrm R}}} not carried to the members. \\
\phantomsection\label{entry:symbol-registry-07077}\(G_S^{\mathrm R}\) & \texttt{\detokenize{G_S^{\mathrm R}}} & implication graph of a 2-CNF list and its reverse symbol family member \(G_S^{\mathrm R}\) & \hyperref[sec:p02-w02-p02-019-two-sat]{\ref*{sec:p02-w02-p02-019-two-sat}} & n/a & \hyperref[def:p02-w02-p02-019-implication-graph]{\ref*{def:p02-w02-p02-019-implication-graph}} & \thepage & 2n literal vertices and 2m edges counted with repetition; G\^{}R reverses every edge. Explicit displayed family; member types and common actual first-use exposed, not one atomic glyph. Family listing, not one atomic glyph. Author macro \texttt{\detokenize{G\_S,}\allowbreak\detokenize{G\_S\^{}{\mathrm R}}} not carried to the members. \\
\phantomsection\label{entry:symbol-registry-07078}\(X\leadsto Y\) & \texttt{\detokenize{\leadsto}} & directed reachability by a path of length 0 or more & \hyperref[sec:p02-w02-p02-019-two-sat]{\ref*{sec:p02-w02-p02-019-two-sat}} & meta & \hyperref[def:p02-w02-p02-019-implication-graph]{\ref*{def:p02-w02-p02-019-implication-graph}} & \thepage & Same use as in annex:p04.dh.two-sat; distinct from the scoped P9 role-compatibility relation \(\leadsto_r\). \\
\phantomsection\label{entry:symbol-registry-07079}\(C(X)\) & \texttt{\detokenize{C(\cdot)}} & strong component of a vertex & \hyperref[sec:p02-w02-p02-019-two-sat]{\ref*{sec:p02-w02-p02-019-two-sat}} & n/a & \hyperref[def:p02-w02-p02-019-implication-graph]{\ref*{def:p02-w02-p02-019-implication-graph}} & \thepage & Applied to a vertex; bare C and C' in annex:p02-w02-p02-019-scc are components, while C and C\_j elsewhere in the section are clauses. Corresponds to C\_L of annex:p04.dh.two-sat. \\
\phantomsection\label{entry:symbol-registry-07080}\(\operatorname{idx}\) & \texttt{\detokenize{\operatorname{idx}}} & component number and finishing data of the two-pass method symbol family member \(\operatorname{idx}\) & \hyperref[sec:p02-w02-p02-019-two-sat]{\ref*{sec:p02-w02-p02-019-two-sat}} & n/a & \hyperref[def:p02-w02-p02-019-scc-method]{\ref*{def:p02-w02-p02-019-scc-method}} & \thepage & idx is a topological numbering of the condensation; f is the first-pass finishing number; F(C) its maximum on a component; fin the finishing list. h=idx in annex:p02-w02-p02-019-2sat. Explicit displayed family; member types and common actual first-use exposed, not one atomic glyph. Family listing, not one atomic glyph. Author macro \texttt{\detokenize{\operatorname{idx},}\allowbreak\detokenize{f,}\allowbreak\detokenize{F,}\allowbreak\detokenize{\mathit{fin}}} not carried to the members. \\
\phantomsection\label{entry:symbol-registry-07081}\(f(u)\) & \texttt{\detokenize{f(u)}} & component number and finishing data of the two-pass method symbol family member \(f(u)\) & \hyperref[sec:p02-w02-p02-019-two-sat]{\ref*{sec:p02-w02-p02-019-two-sat}} & n/a & \hyperref[def:p02-w02-p02-019-scc-method]{\ref*{def:p02-w02-p02-019-scc-method}} & \thepage & idx is a topological numbering of the condensation; f is the first-pass finishing number; F(C) its maximum on a component; fin the finishing list. h=idx in annex:p02-w02-p02-019-2sat. Explicit displayed family; member types and common actual first-use exposed, not one atomic glyph. Family listing, not one atomic glyph. Author macro \texttt{\detokenize{\operatorname{idx},}\allowbreak\detokenize{f,}\allowbreak\detokenize{F,}\allowbreak\detokenize{\mathit{fin}}} not carried to the members. \\
\phantomsection\label{entry:symbol-registry-07082}\(F(C)\) & \texttt{\detokenize{F(C)}} & component number and finishing data of the two-pass method symbol family member \(F(C)\) & \hyperref[sec:p02-w02-p02-019-two-sat]{\ref*{sec:p02-w02-p02-019-two-sat}} & n/a & \hyperref[def:p02-w02-p02-019-scc-method]{\ref*{def:p02-w02-p02-019-scc-method}} & \thepage & idx is a topological numbering of the condensation; f is the first-pass finishing number; F(C) its maximum on a component; fin the finishing list. h=idx in annex:p02-w02-p02-019-2sat. Explicit displayed family; member types and common actual first-use exposed, not one atomic glyph. Family listing, not one atomic glyph. Author macro \texttt{\detokenize{\operatorname{idx},}\allowbreak\detokenize{f,}\allowbreak\detokenize{F,}\allowbreak\detokenize{\mathit{fin}}} not carried to the members. \\
\phantomsection\label{entry:symbol-registry-07083}\(\mathit{fin}\) & \texttt{\detokenize{\mathit{fin}}} & component number and finishing data of the two-pass method symbol family member \(\mathit{fin}\) & \hyperref[sec:p02-w02-p02-019-two-sat]{\ref*{sec:p02-w02-p02-019-two-sat}} & n/a & \hyperref[def:p02-w02-p02-019-scc-method]{\ref*{def:p02-w02-p02-019-scc-method}} & \thepage & idx is a topological numbering of the condensation; f is the first-pass finishing number; F(C) its maximum on a component; fin the finishing list. h=idx in annex:p02-w02-p02-019-2sat. Explicit displayed family; member types and common actual first-use exposed, not one atomic glyph. Family listing, not one atomic glyph. Author macro \texttt{\detokenize{\operatorname{idx},}\allowbreak\detokenize{f,}\allowbreak\detokenize{F,}\allowbreak\detokenize{\mathit{fin}}} not carried to the members. \\
\phantomsection\label{entry:symbol-registry-07084}\(V\) & \texttt{\detokenize{V}} & numbers of vertices and edges of a directed graph symbol family member \(V\) & \hyperref[sec:p02-w02-p02-019-two-sat]{\ref*{sec:p02-w02-p02-019-two-sat}} & n/a & \hyperref[thm:p02-w02-p02-019-scc]{\ref*{thm:p02-w02-p02-019-scc}} & \thepage & Counts, not vertex or edge sets. Explicit displayed family; member types and common actual first-use exposed, not one atomic glyph. Family listing, not one atomic glyph. Author macro \texttt{\detokenize{V,}\allowbreak\detokenize{E}} not carried to the members. \\
\phantomsection\label{entry:symbol-registry-07085}\(E\) & \texttt{\detokenize{E}} & numbers of vertices and edges of a directed graph symbol family member \(E\) & \hyperref[sec:p02-w02-p02-019-two-sat]{\ref*{sec:p02-w02-p02-019-two-sat}} & n/a & \hyperref[thm:p02-w02-p02-019-scc]{\ref*{thm:p02-w02-p02-019-scc}} & \thepage & Counts, not vertex or edge sets. Explicit displayed family; member types and common actual first-use exposed, not one atomic glyph. Family listing, not one atomic glyph. Author macro \texttt{\detokenize{V,}\allowbreak\detokenize{E}} not carried to the members. \\
\phantomsection\label{entry:symbol-registry-07086}\(u^{\circ}\) & \texttt{\detokenize{u^{\circ}}} & truth-table assignment extended by false and the size parameters symbol family member \(u^{\circ}\) & \hyperref[sec:p02-w02-p02-019-complexity]{\ref*{sec:p02-w02-p02-019-complexity}} & n/a & \hyperref[def:p02-w02-p02-019-truth-table-method]{\ref*{def:p02-w02-p02-019-truth-table-method}} & \thepage & k distinct variables, ell total symbol occurrences, r formulas; \(\neg^j\chi\) in annex:p02-w02-p02-019-tt is chi preceded by j negations. Explicit displayed family; member types and common actual first-use exposed, not one atomic glyph. Family listing, not one atomic glyph. Author macro \texttt{\detokenize{u\^{}{\circ},}\allowbreak\detokenize{k,}\allowbreak\detokenize{\ell,}\allowbreak\detokenize{r}} not carried to the members. \\
\phantomsection\label{entry:symbol-registry-07087}\(k\) & \texttt{\detokenize{k}} & truth-table assignment extended by false and the size parameters symbol family member \(k\) & \hyperref[sec:p02-w02-p02-019-complexity]{\ref*{sec:p02-w02-p02-019-complexity}} & n/a & \hyperref[def:p02-w02-p02-019-truth-table-method]{\ref*{def:p02-w02-p02-019-truth-table-method}} & \thepage & k distinct variables, ell total symbol occurrences, r formulas; \(\neg^j\chi\) in annex:p02-w02-p02-019-tt is chi preceded by j negations. Explicit displayed family; member types and common actual first-use exposed, not one atomic glyph. Family listing, not one atomic glyph. Author macro \texttt{\detokenize{u\^{}{\circ},}\allowbreak\detokenize{k,}\allowbreak\detokenize{\ell,}\allowbreak\detokenize{r}} not carried to the members. \\
\phantomsection\label{entry:symbol-registry-07088}\(\ell\) & \texttt{\detokenize{\ell}} & truth-table assignment extended by false and the size parameters symbol family member \(\ell\) & \hyperref[sec:p02-w02-p02-019-complexity]{\ref*{sec:p02-w02-p02-019-complexity}} & n/a & \hyperref[def:p02-w02-p02-019-truth-table-method]{\ref*{def:p02-w02-p02-019-truth-table-method}} & \thepage & k distinct variables, ell total symbol occurrences, r formulas; \(\neg^j\chi\) in annex:p02-w02-p02-019-tt is chi preceded by j negations. Explicit displayed family; member types and common actual first-use exposed, not one atomic glyph. Family listing, not one atomic glyph. Author macro \texttt{\detokenize{u\^{}{\circ},}\allowbreak\detokenize{k,}\allowbreak\detokenize{\ell,}\allowbreak\detokenize{r}} not carried to the members. \\
\phantomsection\label{entry:symbol-registry-07089}\(r\) & \texttt{\detokenize{r}} & truth-table assignment extended by false and the size parameters symbol family member \(r\) & \hyperref[sec:p02-w02-p02-019-complexity]{\ref*{sec:p02-w02-p02-019-complexity}} & n/a & \hyperref[def:p02-w02-p02-019-truth-table-method]{\ref*{def:p02-w02-p02-019-truth-table-method}} & \thepage & k distinct variables, ell total symbol occurrences, r formulas; \(\neg^j\chi\) in annex:p02-w02-p02-019-tt is chi preceded by j negations. Explicit displayed family; member types and common actual first-use exposed, not one atomic glyph. Family listing, not one atomic glyph. Author macro \texttt{\detokenize{u\^{}{\circ},}\allowbreak\detokenize{k,}\allowbreak\detokenize{\ell,}\allowbreak\detokenize{r}} not carried to the members. \\
\phantomsection\label{entry:symbol-registry-07090}\(R(x,y)\) & \texttt{\detokenize{R(x,}\allowbreak\detokenize{y)}} & certificate predicate for SAT and its auxiliaries symbol family member \(R(x,y)\) & \hyperref[sec:p02-w02-p02-019-complexity]{\ref*{sec:p02-w02-p02-019-complexity}} & n/a & \hyperref[thm:p02-w02-p02-019-sat-in-np]{\ref*{thm:p02-w02-p02-019-sat-in-np}} & \thepage & R is the polynomial-time predicate of def:p02-complexity-standing; A is a marker character used only inside the recognizing machine. Explicit displayed family; member types and common actual first-use exposed, not one atomic glyph. Family listing, not one atomic glyph. Author macro \texttt{\detokenize{R,}\allowbreak\detokenize{v\_y,}\allowbreak\detokenize{p\_{a\_i},}\allowbreak\detokenize{\mathtt A}} not carried to the members. \\
\phantomsection\label{entry:symbol-registry-07091}\(v_y\) & \texttt{\detokenize{v_y}} & certificate predicate for SAT and its auxiliaries symbol family member \(v_y\) & \hyperref[sec:p02-w02-p02-019-complexity]{\ref*{sec:p02-w02-p02-019-complexity}} & n/a & \hyperref[thm:p02-w02-p02-019-sat-in-np]{\ref*{thm:p02-w02-p02-019-sat-in-np}} & \thepage & R is the polynomial-time predicate of def:p02-complexity-standing; A is a marker character used only inside the recognizing machine. Explicit displayed family; member types and common actual first-use exposed, not one atomic glyph. Family listing, not one atomic glyph. Author macro \texttt{\detokenize{R,}\allowbreak\detokenize{v\_y,}\allowbreak\detokenize{p\_{a\_i},}\allowbreak\detokenize{\mathtt A}} not carried to the members. \\
\phantomsection\label{entry:symbol-registry-07092}\(p_{a_i}\) & \texttt{\detokenize{p_{a_i}}} & certificate predicate for SAT and its auxiliaries symbol family member \(p_{a_i}\) & \hyperref[sec:p02-w02-p02-019-complexity]{\ref*{sec:p02-w02-p02-019-complexity}} & n/a & \hyperref[thm:p02-w02-p02-019-sat-in-np]{\ref*{thm:p02-w02-p02-019-sat-in-np}} & \thepage & R is the polynomial-time predicate of def:p02-complexity-standing; A is a marker character used only inside the recognizing machine. Explicit displayed family; member types and common actual first-use exposed, not one atomic glyph. Family listing, not one atomic glyph. Author macro \texttt{\detokenize{R,}\allowbreak\detokenize{v\_y,}\allowbreak\detokenize{p\_{a\_i},}\allowbreak\detokenize{\mathtt A}} not carried to the members. \\
\phantomsection\label{entry:symbol-registry-07093}\(\mathtt A\) & \texttt{\detokenize{\mathtt A}} & certificate predicate for SAT and its auxiliaries symbol family member \(\mathtt A\) & \hyperref[sec:p02-w02-p02-019-complexity]{\ref*{sec:p02-w02-p02-019-complexity}} & n/a & \hyperref[thm:p02-w02-p02-019-sat-in-np]{\ref*{thm:p02-w02-p02-019-sat-in-np}} & \thepage & R is the polynomial-time predicate of def:p02-complexity-standing; A is a marker character used only inside the recognizing machine. Explicit displayed family; member types and common actual first-use exposed, not one atomic glyph. Family listing, not one atomic glyph. Author macro \texttt{\detokenize{R,}\allowbreak\detokenize{v\_y,}\allowbreak\detokenize{p\_{a\_i},}\allowbreak\detokenize{\mathtt A}} not carried to the members. \\
\phantomsection\label{entry:symbol-registry-07094}\(\overline A\) & \texttt{\detokenize{\overline A}} & complement of a language in \(\{0,1\}^*\) & \hyperref[sec:p02-w02-p02-019-complexity]{\ref*{sec:p02-w02-p02-019-complexity}} & meta & \hyperref[thm:p02-w02-p02-019-reductions]{\ref*{thm:p02-w02-p02-019-reductions}} & \thepage & Applied to languages; the registered literal-complement bar applies to literals. Both meanings are complements, of different types. \\
\phantomsection\label{entry:symbol-registry-07095}\(\nu\) & \texttt{\detokenize{\nu}} & negation map on encodings and the two reductions between TAUT and the complement of SAT symbol family member \(\nu\) & \hyperref[sec:p02-w02-p02-019-complexity]{\ref*{sec:p02-w02-p02-019-complexity}} & n/a & \hyperref[thm:p02-w02-p02-019-negation-maps]{\ref*{thm:p02-w02-p02-019-negation-maps}} & \thepage & Disjoint in scope from the registered P6 P8 P12 and P3 uses of \(\nu\). Explicit displayed family; member types and common actual first-use exposed, not one atomic glyph. Family listing, not one atomic glyph. Author macro \texttt{\detokenize{\nu,}\allowbreak\detokenize{g\_1,}\allowbreak\detokenize{g\_2}} not carried to the members. \\
\phantomsection\label{entry:symbol-registry-07096}\(g_1\) & \texttt{\detokenize{g_1}} & negation map on encodings and the two reductions between TAUT and the complement of SAT symbol family member \(g_1\) & \hyperref[sec:p02-w02-p02-019-complexity]{\ref*{sec:p02-w02-p02-019-complexity}} & n/a & \hyperref[thm:p02-w02-p02-019-negation-maps]{\ref*{thm:p02-w02-p02-019-negation-maps}} & \thepage & Disjoint in scope from the registered P6 P8 P12 and P3 uses of \(\nu\). Explicit displayed family; member types and common actual first-use exposed, not one atomic glyph. Family listing, not one atomic glyph. Author macro \texttt{\detokenize{\nu,}\allowbreak\detokenize{g\_1,}\allowbreak\detokenize{g\_2}} not carried to the members. \\
\phantomsection\label{entry:symbol-registry-07097}\(g_2\) & \texttt{\detokenize{g_2}} & negation map on encodings and the two reductions between TAUT and the complement of SAT symbol family member \(g_2\) & \hyperref[sec:p02-w02-p02-019-complexity]{\ref*{sec:p02-w02-p02-019-complexity}} & n/a & \hyperref[thm:p02-w02-p02-019-negation-maps]{\ref*{thm:p02-w02-p02-019-negation-maps}} & \thepage & Disjoint in scope from the registered P6 P8 P12 and P3 uses of \(\nu\). Explicit displayed family; member types and common actual first-use exposed, not one atomic glyph. Family listing, not one atomic glyph. Author macro \texttt{\detokenize{\nu,}\allowbreak\detokenize{g\_1,}\allowbreak\detokenize{g\_2}} not carried to the members. \\
\phantomsection\label{entry:symbol-registry-07098}\(a_{ji}\) & \texttt{\detokenize{a_{ji}}} & coefficients right-hand sides and unknowns of a parity system over the two-element field symbol family member \(a_{ji}\) & \hyperref[sec:p02-w02-p02-019-xor]{\ref*{sec:p02-w02-p02-019-xor}} & n/a & \hyperref[def:p02-w02-p02-019-parity-input]{\ref*{def:p02-w02-p02-019-parity-input}} & \thepage & m equations and n variables as in def:p02-tractable-classes; x\_i is 1 iff p\_i is true. x\^{}* is the vector of the returned valuation v\^{}*. Explicit displayed family; member types and common actual first-use exposed, not one atomic glyph. Family listing, not one atomic glyph. Author macro \texttt{\detokenize{a\_{ji},}\allowbreak\detokenize{b\_j,}\allowbreak\detokenize{x\_i}} not carried to the members. \\
\phantomsection\label{entry:symbol-registry-07099}\(b_j\) & \texttt{\detokenize{b_j}} & coefficients right-hand sides and unknowns of a parity system over the two-element field symbol family member \(b_j\) & \hyperref[sec:p02-w02-p02-019-xor]{\ref*{sec:p02-w02-p02-019-xor}} & n/a & \hyperref[def:p02-w02-p02-019-parity-input]{\ref*{def:p02-w02-p02-019-parity-input}} & \thepage & m equations and n variables as in def:p02-tractable-classes; x\_i is 1 iff p\_i is true. x\^{}* is the vector of the returned valuation v\^{}*. Explicit displayed family; member types and common actual first-use exposed, not one atomic glyph. Family listing, not one atomic glyph. Author macro \texttt{\detokenize{a\_{ji},}\allowbreak\detokenize{b\_j,}\allowbreak\detokenize{x\_i}} not carried to the members. \\
\phantomsection\label{entry:symbol-registry-07100}\(x_i\) & \texttt{\detokenize{x_i}} & coefficients right-hand sides and unknowns of a parity system over the two-element field symbol family member \(x_i\) & \hyperref[sec:p02-w02-p02-019-xor]{\ref*{sec:p02-w02-p02-019-xor}} & n/a & \hyperref[def:p02-w02-p02-019-parity-input]{\ref*{def:p02-w02-p02-019-parity-input}} & \thepage & m equations and n variables as in def:p02-tractable-classes; x\_i is 1 iff p\_i is true. x\^{}* is the vector of the returned valuation v\^{}*. Explicit displayed family; member types and common actual first-use exposed, not one atomic glyph. Family listing, not one atomic glyph. Author macro \texttt{\detokenize{a\_{ji},}\allowbreak\detokenize{b\_j,}\allowbreak\detokenize{x\_i}} not carried to the members. \\
\phantomsection\label{entry:symbol-registry-07101}\(r\) & \texttt{\detokenize{r}} & current pivot row number of pivots current column and pivot columns of the elimination method symbol family member \(r\) & \hyperref[sec:p02-w02-p02-019-xor]{\ref*{sec:p02-w02-p02-019-xor}} & n/a & \hyperref[def:p02-w02-p02-019-gauss-method]{\ref*{def:p02-w02-p02-019-gauss-method}} & \thepage & Program variables of the method; r here is a row index, distinct from the formula count r of sec:p02-w02-p02-019-complexity; c is a column index. Explicit displayed family; member types and common actual first-use exposed, not one atomic glyph. Family listing, not one atomic glyph. Author macro \texttt{\detokenize{r,}\allowbreak\detokenize{t,}\allowbreak\detokenize{c,}\allowbreak\detokenize{c\_s}} not carried to the members. \\
\phantomsection\label{entry:symbol-registry-07102}\(t\) & \texttt{\detokenize{t}} & current pivot row number of pivots current column and pivot columns of the elimination method symbol family member \(t\) & \hyperref[sec:p02-w02-p02-019-xor]{\ref*{sec:p02-w02-p02-019-xor}} & n/a & \hyperref[def:p02-w02-p02-019-gauss-method]{\ref*{def:p02-w02-p02-019-gauss-method}} & \thepage & Program variables of the method; r here is a row index, distinct from the formula count r of sec:p02-w02-p02-019-complexity; c is a column index. Explicit displayed family; member types and common actual first-use exposed, not one atomic glyph. Family listing, not one atomic glyph. Author macro \texttt{\detokenize{r,}\allowbreak\detokenize{t,}\allowbreak\detokenize{c,}\allowbreak\detokenize{c\_s}} not carried to the members. \\
\phantomsection\label{entry:symbol-registry-07103}\(c\) & \texttt{\detokenize{c}} & current pivot row number of pivots current column and pivot columns of the elimination method symbol family member \(c\) & \hyperref[sec:p02-w02-p02-019-xor]{\ref*{sec:p02-w02-p02-019-xor}} & n/a & \hyperref[def:p02-w02-p02-019-gauss-method]{\ref*{def:p02-w02-p02-019-gauss-method}} & \thepage & Program variables of the method; r here is a row index, distinct from the formula count r of sec:p02-w02-p02-019-complexity; c is a column index. Explicit displayed family; member types and common actual first-use exposed, not one atomic glyph. Family listing, not one atomic glyph. Author macro \texttt{\detokenize{r,}\allowbreak\detokenize{t,}\allowbreak\detokenize{c,}\allowbreak\detokenize{c\_s}} not carried to the members. \\
\phantomsection\label{entry:symbol-registry-07104}\(c_s\) & \texttt{\detokenize{c_s}} & current pivot row number of pivots current column and pivot columns of the elimination method symbol family member \(c_s\) & \hyperref[sec:p02-w02-p02-019-xor]{\ref*{sec:p02-w02-p02-019-xor}} & n/a & \hyperref[def:p02-w02-p02-019-gauss-method]{\ref*{def:p02-w02-p02-019-gauss-method}} & \thepage & Program variables of the method; r here is a row index, distinct from the formula count r of sec:p02-w02-p02-019-complexity; c is a column index. Explicit displayed family; member types and common actual first-use exposed, not one atomic glyph. Family listing, not one atomic glyph. Author macro \texttt{\detokenize{r,}\allowbreak\detokenize{t,}\allowbreak\detokenize{c,}\allowbreak\detokenize{c\_s}} not carried to the members. \\
\phantomsection\label{entry:symbol-registry-07105}\(v^{*}\) & \texttt{\detokenize{v^{*}}} & valuation returned by the elimination method & \hyperref[sec:p02-w02-p02-019-xor]{\ref*{sec:p02-w02-p02-019-xor}} & n/a & \hyperref[def:p02-w02-p02-019-gauss-method]{\ref*{def:p02-w02-p02-019-gauss-method}} & \thepage & Pivot variables take the final right-hand side of their row; all others false. \\
\phantomsection\label{entry:symbol-registry-07106}\(P\) & \texttt{\detokenize{P}} & unary predicate symbol of the non-terminating resolution example & \hyperref[sec:p02-w02-p02-019-first-order-boundary]{\ref*{sec:p02-w02-p02-019-first-order-boundary}}--\hyperref[thm:p02-w02-p02-019-fol-resolution]{\ref*{thm:p02-w02-p02-019-fol-resolution}} & n/a & \hyperref[thm:p02-w02-p02-019-fol-resolution]{\ref*{thm:p02-w02-p02-019-fol-resolution}} & \thepage & Example symbols only, in the non-terminating resolution example after thm:p02-w02-p02-019-fol-resolution, as in the example after thm:resolution-lifting. \\
\phantomsection\label{entry:symbol-registry-07107}\(f\) & \texttt{\detokenize{f}} & unary function symbol of the non-terminating resolution example & \hyperref[sec:p02-w02-p02-019-first-order-boundary]{\ref*{sec:p02-w02-p02-019-first-order-boundary}}--\hyperref[thm:p02-w02-p02-019-fol-resolution]{\ref*{thm:p02-w02-p02-019-fol-resolution}} & n/a & \hyperref[thm:p02-w02-p02-019-fol-resolution]{\ref*{thm:p02-w02-p02-019-fol-resolution}} & \thepage & Example symbols only, in the non-terminating resolution example after thm:p02-w02-p02-019-fol-resolution, as in the example after thm:resolution-lifting. \\
\phantomsection\label{entry:symbol-registry-07108}\(a\) & \texttt{\detokenize{a}} & constant of the non-terminating resolution example & \hyperref[sec:p02-w02-p02-019-first-order-boundary]{\ref*{sec:p02-w02-p02-019-first-order-boundary}}--\hyperref[thm:p02-w02-p02-019-fol-resolution]{\ref*{thm:p02-w02-p02-019-fol-resolution}} & n/a & \hyperref[thm:p02-w02-p02-019-fol-resolution]{\ref*{thm:p02-w02-p02-019-fol-resolution}} & \thepage & Example symbols only, in the non-terminating resolution example after thm:p02-w02-p02-019-fol-resolution, as in the example after thm:resolution-lifting. \\
\phantomsection\label{entry:symbol-registry-07109}\(x\) & \texttt{\detokenize{x}} & first-order variable of the non-terminating resolution example & \hyperref[sec:p02-w02-p02-019-first-order-boundary]{\ref*{sec:p02-w02-p02-019-first-order-boundary}}--\hyperref[thm:p02-w02-p02-019-fol-resolution]{\ref*{thm:p02-w02-p02-019-fol-resolution}} & n/a & \hyperref[thm:p02-w02-p02-019-fol-resolution]{\ref*{thm:p02-w02-p02-019-fol-resolution}} & \thepage & Example symbols only, in the non-terminating resolution example after thm:p02-w02-p02-019-fol-resolution, as in the example after thm:resolution-lifting. x here is a first-order variable, not the parity unknown x\_i of sec:p02-w02-p02-019-xor. \\
\phantomsection\label{entry:symbol-registry-07110}\(\kappa_u\) & \texttt{\detokenize{\kappa_u}} & complete conjunction of X-literals describing an assignment u & \hyperref[sec:p02-w02-p02-015-interpolation-boundary]{\ref*{sec:p02-w02-p02-015-interpolation-boundary}} & n/a & \hyperref[def:p02-w02-p02-015-output]{\ref*{def:p02-w02-p02-015-output}} & \thepage & Right-associated conjunction, in increasing variable index, of x\_i or \(\neg x_i\) according as u(x\_i) is 1 or 0; equals \(\topv\) when X is empty. Exact reading of the complete conjunction of Method def:p02-interpolant-truth-table. \\
\phantomsection\label{entry:symbol-registry-07111}\(U_\varphi\) & \texttt{\detokenize{U_\varphi}} & selected family of assignments to the common variables & \hyperref[sec:p02-w02-p02-015-interpolation-boundary]{\ref*{sec:p02-w02-p02-015-interpolation-boundary}} & n/a & \hyperref[def:p02-w02-p02-015-output]{\ref*{def:p02-w02-p02-015-output}} & \thepage & The assignments X to {0,1} that are restrictions of valuations satisfying \(\varphi\); X is the common-variable set of the consequence under discussion. \\
\phantomsection\label{entry:symbol-registry-07112}\(u_\varnothing\) & \texttt{\detokenize{u_\varnothing}} & the empty assignment & \hyperref[sec:p02-w02-p02-015-interpolation-boundary]{\ref*{sec:p02-w02-p02-015-interpolation-boundary}} & n/a & \hyperref[def:p02-w02-p02-015-constants]{\ref*{def:p02-w02-p02-015-constants}} & \thepage & The unique function from the empty set to {0,1}; every valuation restricts to it. \\
\phantomsection\label{entry:symbol-registry-07113}\(\mathsf{Fm}_{\mathsf{CPL}}^{-}\) & \texttt{\detokenize{\mathsf{Fm}_{\mathsf{CPL}}^{-}}} & bottom-free propositional formulas & \hyperref[sec:p02-w02-p02-015-interpolation-boundary]{\ref*{sec:p02-w02-p02-015-interpolation-boundary}} & n/a & \hyperref[def:p02-w02-p02-015-bottom-free]{\ref*{def:p02-w02-p02-015-bottom-free}} & \thepage & Least set of strings containing every p\_n and closed under negation and the three binary constructors; the formulas of \(\mathsf{Fm}_{\mathsf{CPL}}\) in which \(\bottom\) does not occur. \\
\phantomsection\label{entry:symbol-registry-07114}\(\equiv_P\) & \texttt{\detokenize{\equiv_P}} & two-valued logical equivalence over arbitrary atoms & \hyperref[sec:p02-w02-p02-016-finite-input-scope]{\ref*{sec:p02-w02-p02-016-finite-input-scope}} & relation & \hyperref[def:p02-w02-p02-016-semantics]{\ref*{def:p02-w02-p02-016-semantics}} & \thepage & Same value under every valuation P to {0,1}; companion of the registered consequence relation \(\models_P\) (written \(\entails_P\)); reduces to source equivalence when P is Prop. \\
\phantomsection\label{entry:symbol-registry-07115}\(\mathsf r\) & \texttt{\detokenize{\mathsf r}} & injection of a finite atom set into Prop & \hyperref[sec:p02-w02-p02-016-finite-input-scope]{\ref*{sec:p02-w02-p02-016-finite-input-scope}} & n/a & \hyperref[def:p02-w02-p02-016-renaming]{\ref*{def:p02-w02-p02-016-renaming}} & \thepage & Domain a finite subset A of P; initial when a bijection onto p\_0 to p\_{k-1}. Sans-serif letter chosen to avoid the registered P2 uses of iota (base type, description binder, at-least-n sentence). \\
\phantomsection\label{entry:symbol-registry-07116}\(\mathsf r^{\star}\) & \texttt{\detokenize{{}^{\star}}} & induced translation Form(P) to Fm\_CPL & \hyperref[sec:p02-w02-p02-016-finite-input-scope]{\ref*{sec:p02-w02-p02-016-finite-input-scope}} & n/a & \hyperref[def:p02-w02-p02-016-renaming]{\ref*{def:p02-w02-p02-016-renaming}} & \thepage & Structural recursion: atoms of A to their r-images, other atoms to bottom, commutes with every constructor. Chosen instead of a sharp superscript, which is registered for the P2 incidence translation. \\
\phantomsection\label{entry:symbol-registry-07117}\(w^{\mathsf r}\) & \texttt{\detokenize{w^{\mathsf r}}} & pulled-back and pushed-forward valuations symbol family member \(w^{\mathsf r}\) & \hyperref[sec:annex-p02-w02-p02-016-finite-input-scope]{\ref*{sec:annex-p02-w02-p02-016-finite-input-scope}} & n/a & \hyperref[annex:p02-w02-p02-016-transfer]{\ref*{annex:p02-w02-p02-016-transfer}} & \thepage & w\^{}r: P to {0,1} reads w on r-star of each atom; v\_r: Prop to {T,F} is true exactly on r-images of atoms true under v. Explicit displayed family; member types and common actual first-use exposed, not one atomic glyph. Family listing, not one atomic glyph. Author macro \texttt{\detokenize{w\^{}{\mathsf r},}\allowbreak\detokenize{v\_{\mathsf r}}} not carried to the members. \\
\phantomsection\label{entry:symbol-registry-07118}\(v_{\mathsf r}\) & \texttt{\detokenize{v_{\mathsf r}}} & pulled-back and pushed-forward valuations symbol family member \(v_{\mathsf r}\) & \hyperref[sec:annex-p02-w02-p02-016-finite-input-scope]{\ref*{sec:annex-p02-w02-p02-016-finite-input-scope}} & n/a & \hyperref[annex:p02-w02-p02-016-transfer]{\ref*{annex:p02-w02-p02-016-transfer}} & \thepage & w\^{}r: P to {0,1} reads w on r-star of each atom; v\_r: Prop to {T,F} is true exactly on r-images of atoms true under v. Explicit displayed family; member types and common actual first-use exposed, not one atomic glyph. Family listing, not one atomic glyph. Author macro \texttt{\detokenize{w\^{}{\mathsf r},}\allowbreak\detokenize{v\_{\mathsf r}}} not carried to the members. \\
\phantomsection\label{entry:symbol-registry-07119}\(P_{\N}\) & \texttt{\detokenize{P_{\N}}} & the source propositional alphabet as a set of atoms & \hyperref[sec:p02-w02-p02-017-finite-input-scope]{\ref*{sec:p02-w02-p02-017-finite-input-scope}} & n/a & \hyperref[sec:p02-w02-p02-017-finite-input-scope]{\ref*{sec:p02-w02-p02-017-finite-input-scope}} & \thepage & Editorial name for the source set of variables p\_n (n in N) of sec:prop-syntax-semantics-01; the source macro is restoration-only. Form(P\_N) is the source formula set. \\
\phantomsection\label{entry:symbol-registry-07120}\(\alpha^{c}\) & \texttt{\detokenize{\alpha^{c}}} & formula obtained by renaming each atom q to p\_{c(q)} along an atom coding c & \hyperref[sec:p02-w02-p02-017-finite-input-scope]{\ref*{sec:p02-w02-p02-017-finite-input-scope}} & n/a & \hyperref[def:p02-w02-p02-017-coded-input]{\ref*{def:p02-w02-p02-017-coded-input}} & \thepage & Structural recursion; injective when c is injective. Superscript is the coding, not a power or complement. \\
\phantomsection\label{entry:symbol-registry-07121}\(s^{0}\) & \texttt{\detokenize{s^{0}}} & extension of a finite assignment s on A to a valuation on P by value 0 off A & \hyperref[thm:p02-w02-p02-017-finite-reduction]{\ref*{thm:p02-w02-p02-017-finite-reduction}} & n/a & \hyperref[thm:p02-w02-p02-017-finite-reduction]{\ref*{thm:p02-w02-p02-017-finite-reduction}} & \thepage & Defined by cases on membership in A; no choice. \\
\phantomsection\label{entry:symbol-registry-07122}\(\ulcorner\alpha\urcorner\) & \texttt{\detokenize{\ulcorner\cdot\urcorner}} & code number of a propositional formula: number with binary numeral 1 followed by the def:p02-sat-taut encoding & \hyperref[sec:p02-w02-p02-017-finite-input-scope]{\ref*{sec:p02-w02-p02-017-finite-input-scope}} & n/a & \hyperref[def:p02-w02-p02-017-coded-input]{\ref*{def:p02-w02-p02-017-coded-input}} & \thepage & Same glyph and the same 'fixed effective code' role family as the P5 Goedel-number row; different coding and scope. \\
\phantomsection\label{entry:symbol-registry-07123}\(\Gamma_x\) & \texttt{\detokenize{\Gamma_x}} & uniformly decidable premise family whose satisfiability is undecidable & P2 infinite-premise-undecidability & n/a & \hyperref[thm:p02-w02-p02-017-infinite-undecidable]{\ref*{thm:p02-w02-p02-017-infinite-undecidable}} & \thepage & theta\_y = not p\_0 and (p\_{y+1} or not p\_{y+1}); Gamma\_x = {p\_0} with theta\_y for witnesses y of x in K. \\
\phantomsection\label{entry:symbol-registry-07124}\(\theta_y\) & \texttt{\detokenize{\theta_y}} & padded witness formula of the premise family \(\Gamma_x\) & P2 infinite-premise-undecidability & n/a & \hyperref[annex:p02-w02-p02-017-infinite-undecidable]{\ref*{annex:p02-w02-p02-017-infinite-undecidable}} & \thepage & theta\_y = not p\_0 and (p\_{y+1} or not p\_{y+1}); Gamma\_x = {p\_0} with theta\_y for witnesses y of x in K. \\
\phantomsection\label{entry:symbol-registry-07125}\(\Delta\) & \texttt{\detokenize{\Delta}} & decidable satisfiable premise set without computable model & P2 infinite-premise-undecidability & n/a & \hyperref[thm:p02-w02-p02-017-no-computable-valuation]{\ref*{thm:p02-w02-p02-017-no-computable-valuation}} & \thepage & tau\_y is the tautology instance on p\_{2y+1}; delta\_{x,y} carries the sign of p\_{2x} fixed by U(y). \\
\phantomsection\label{entry:symbol-registry-07126}\(\delta_{x,y}\) & \texttt{\detokenize{\delta_{x,y}}} & member of the premise set \(\Delta\) carrying the sign of \(p_{2x}\) & P2 infinite-premise-undecidability & n/a & \hyperref[annex:p02-w02-p02-017-no-computable-valuation]{\ref*{annex:p02-w02-p02-017-no-computable-valuation}} & \thepage & tau\_y is the tautology instance on p\_{2y+1}; delta\_{x,y} carries the sign of p\_{2x} fixed by U(y). \\
\phantomsection\label{entry:symbol-registry-07127}\(\tau_y\) & \texttt{\detokenize{\tau_y}} & tautology instance on \(p_{2y+1}\) padding the members of \(\Delta\) & P2 infinite-premise-undecidability & n/a & \hyperref[annex:p02-w02-p02-017-no-computable-valuation]{\ref*{annex:p02-w02-p02-017-no-computable-valuation}} & \thepage & tau\_y is the tautology instance on p\_{2y+1}; delta\_{x,y} carries the sign of p\_{2x} fixed by U(y). \\
\phantomsection\label{entry:symbol-registry-07128}\(F\) & \texttt{\detokenize{F}} & local computable or partial computable functions and decidable relations in the annex constructions symbol family member \(F\) & \hyperref[sec:annex-p02-w02-p02-017-finite-input-scope]{\ref*{sec:annex-p02-w02-p02-017-finite-input-scope}} & n/a & \hyperref[sec:annex-p02-w02-p02-017-finite-input-scope]{\ref*{sec:annex-p02-w02-p02-017-finite-input-scope}} & \thepage & F(x) = varphi\_x(x); T(x,y) = T\_1(e*,x,y) from the cited normal form; W least-witness relation; B,C,D,M bounded-search functions; t,d code-number maps. Local to the annex section; F is not the filter F of sec:p02-free-boolean-completeness. Explicit displayed family; member types and common actual first-use exposed, not one atomic glyph. Family listing, not one atomic glyph. Author macro \texttt{\detokenize{F,}\allowbreak\detokenize{e\^{}{*},}\allowbreak\detokenize{T,}\allowbreak\detokenize{W,}\allowbreak\detokenize{N,}\allowbreak\detokenize{B,}\allowbreak\detokenize{C,}\allowbreak\detokenize{D,}\allowbreak\detokenize{M,}\allowbreak\detokenize{t,}\allowbreak\detokenize{d,}\allowbreak\detokenize{d\_0,}\allowbreak\detokenize{d\_1,}\allowbreak\detokenize{h}} not carried to the members. \\
\phantomsection\label{entry:symbol-registry-07129}\(e^{*}\) & \texttt{\detokenize{e^{*}}} & local computable or partial computable functions and decidable relations in the annex constructions symbol family member \(e^{*}\) & \hyperref[sec:annex-p02-w02-p02-017-finite-input-scope]{\ref*{sec:annex-p02-w02-p02-017-finite-input-scope}} & n/a & \hyperref[sec:annex-p02-w02-p02-017-finite-input-scope]{\ref*{sec:annex-p02-w02-p02-017-finite-input-scope}} & \thepage & F(x) = varphi\_x(x); T(x,y) = T\_1(e*,x,y) from the cited normal form; W least-witness relation; B,C,D,M bounded-search functions; t,d code-number maps. Local to the annex section; F is not the filter F of sec:p02-free-boolean-completeness. Explicit displayed family; member types and common actual first-use exposed, not one atomic glyph. Family listing, not one atomic glyph. Author macro \texttt{\detokenize{F,}\allowbreak\detokenize{e\^{}{*},}\allowbreak\detokenize{T,}\allowbreak\detokenize{W,}\allowbreak\detokenize{N,}\allowbreak\detokenize{B,}\allowbreak\detokenize{C,}\allowbreak\detokenize{D,}\allowbreak\detokenize{M,}\allowbreak\detokenize{t,}\allowbreak\detokenize{d,}\allowbreak\detokenize{d\_0,}\allowbreak\detokenize{d\_1,}\allowbreak\detokenize{h}} not carried to the members. \\
\phantomsection\label{entry:symbol-registry-07130}\(T\) & \texttt{\detokenize{T}} & local computable or partial computable functions and decidable relations in the annex constructions symbol family member \(T\) & \hyperref[sec:annex-p02-w02-p02-017-finite-input-scope]{\ref*{sec:annex-p02-w02-p02-017-finite-input-scope}} & n/a & \hyperref[sec:annex-p02-w02-p02-017-finite-input-scope]{\ref*{sec:annex-p02-w02-p02-017-finite-input-scope}} & \thepage & F(x) = varphi\_x(x); T(x,y) = T\_1(e*,x,y) from the cited normal form; W least-witness relation; B,C,D,M bounded-search functions; t,d code-number maps. Local to the annex section; F is not the filter F of sec:p02-free-boolean-completeness. Explicit displayed family; member types and common actual first-use exposed, not one atomic glyph. Family listing, not one atomic glyph. Author macro \texttt{\detokenize{F,}\allowbreak\detokenize{e\^{}{*},}\allowbreak\detokenize{T,}\allowbreak\detokenize{W,}\allowbreak\detokenize{N,}\allowbreak\detokenize{B,}\allowbreak\detokenize{C,}\allowbreak\detokenize{D,}\allowbreak\detokenize{M,}\allowbreak\detokenize{t,}\allowbreak\detokenize{d,}\allowbreak\detokenize{d\_0,}\allowbreak\detokenize{d\_1,}\allowbreak\detokenize{h}} not carried to the members. \\
\phantomsection\label{entry:symbol-registry-07131}\(W\) & \texttt{\detokenize{W}} & local computable or partial computable functions and decidable relations in the annex constructions symbol family member \(W\) & \hyperref[sec:annex-p02-w02-p02-017-finite-input-scope]{\ref*{sec:annex-p02-w02-p02-017-finite-input-scope}} & n/a & \hyperref[sec:annex-p02-w02-p02-017-finite-input-scope]{\ref*{sec:annex-p02-w02-p02-017-finite-input-scope}} & \thepage & F(x) = varphi\_x(x); T(x,y) = T\_1(e*,x,y) from the cited normal form; W least-witness relation; B,C,D,M bounded-search functions; t,d code-number maps. Local to the annex section; F is not the filter F of sec:p02-free-boolean-completeness. Explicit displayed family; member types and common actual first-use exposed, not one atomic glyph. Family listing, not one atomic glyph. Author macro \texttt{\detokenize{F,}\allowbreak\detokenize{e\^{}{*},}\allowbreak\detokenize{T,}\allowbreak\detokenize{W,}\allowbreak\detokenize{N,}\allowbreak\detokenize{B,}\allowbreak\detokenize{C,}\allowbreak\detokenize{D,}\allowbreak\detokenize{M,}\allowbreak\detokenize{t,}\allowbreak\detokenize{d,}\allowbreak\detokenize{d\_0,}\allowbreak\detokenize{d\_1,}\allowbreak\detokenize{h}} not carried to the members. \\
\phantomsection\label{entry:symbol-registry-07132}\(N\) & \texttt{\detokenize{N}} & local computable or partial computable functions and decidable relations in the annex constructions symbol family member \(N\) & \hyperref[sec:annex-p02-w02-p02-017-finite-input-scope]{\ref*{sec:annex-p02-w02-p02-017-finite-input-scope}} & n/a & \hyperref[sec:annex-p02-w02-p02-017-finite-input-scope]{\ref*{sec:annex-p02-w02-p02-017-finite-input-scope}} & \thepage & F(x) = varphi\_x(x); T(x,y) = T\_1(e*,x,y) from the cited normal form; W least-witness relation; B,C,D,M bounded-search functions; t,d code-number maps. Local to the annex section; F is not the filter F of sec:p02-free-boolean-completeness. Explicit displayed family; member types and common actual first-use exposed, not one atomic glyph. Family listing, not one atomic glyph. Author macro \texttt{\detokenize{F,}\allowbreak\detokenize{e\^{}{*},}\allowbreak\detokenize{T,}\allowbreak\detokenize{W,}\allowbreak\detokenize{N,}\allowbreak\detokenize{B,}\allowbreak\detokenize{C,}\allowbreak\detokenize{D,}\allowbreak\detokenize{M,}\allowbreak\detokenize{t,}\allowbreak\detokenize{d,}\allowbreak\detokenize{d\_0,}\allowbreak\detokenize{d\_1,}\allowbreak\detokenize{h}} not carried to the members. \\
\phantomsection\label{entry:symbol-registry-07133}\(B\) & \texttt{\detokenize{B}} & local computable or partial computable functions and decidable relations in the annex constructions symbol family member \(B\) & \hyperref[sec:annex-p02-w02-p02-017-finite-input-scope]{\ref*{sec:annex-p02-w02-p02-017-finite-input-scope}} & n/a & \hyperref[sec:annex-p02-w02-p02-017-finite-input-scope]{\ref*{sec:annex-p02-w02-p02-017-finite-input-scope}} & \thepage & F(x) = varphi\_x(x); T(x,y) = T\_1(e*,x,y) from the cited normal form; W least-witness relation; B,C,D,M bounded-search functions; t,d code-number maps. Local to the annex section; F is not the filter F of sec:p02-free-boolean-completeness. Explicit displayed family; member types and common actual first-use exposed, not one atomic glyph. Family listing, not one atomic glyph. Author macro \texttt{\detokenize{F,}\allowbreak\detokenize{e\^{}{*},}\allowbreak\detokenize{T,}\allowbreak\detokenize{W,}\allowbreak\detokenize{N,}\allowbreak\detokenize{B,}\allowbreak\detokenize{C,}\allowbreak\detokenize{D,}\allowbreak\detokenize{M,}\allowbreak\detokenize{t,}\allowbreak\detokenize{d,}\allowbreak\detokenize{d\_0,}\allowbreak\detokenize{d\_1,}\allowbreak\detokenize{h}} not carried to the members. \\
\phantomsection\label{entry:symbol-registry-07134}\(C\) & \texttt{\detokenize{C}} & local computable or partial computable functions and decidable relations in the annex constructions symbol family member \(C\) & \hyperref[sec:annex-p02-w02-p02-017-finite-input-scope]{\ref*{sec:annex-p02-w02-p02-017-finite-input-scope}} & n/a & \hyperref[sec:annex-p02-w02-p02-017-finite-input-scope]{\ref*{sec:annex-p02-w02-p02-017-finite-input-scope}} & \thepage & F(x) = varphi\_x(x); T(x,y) = T\_1(e*,x,y) from the cited normal form; W least-witness relation; B,C,D,M bounded-search functions; t,d code-number maps. Local to the annex section; F is not the filter F of sec:p02-free-boolean-completeness. Explicit displayed family; member types and common actual first-use exposed, not one atomic glyph. Family listing, not one atomic glyph. Author macro \texttt{\detokenize{F,}\allowbreak\detokenize{e\^{}{*},}\allowbreak\detokenize{T,}\allowbreak\detokenize{W,}\allowbreak\detokenize{N,}\allowbreak\detokenize{B,}\allowbreak\detokenize{C,}\allowbreak\detokenize{D,}\allowbreak\detokenize{M,}\allowbreak\detokenize{t,}\allowbreak\detokenize{d,}\allowbreak\detokenize{d\_0,}\allowbreak\detokenize{d\_1,}\allowbreak\detokenize{h}} not carried to the members. \\
\phantomsection\label{entry:symbol-registry-07135}\(D\) & \texttt{\detokenize{D}} & local computable or partial computable functions and decidable relations in the annex constructions symbol family member \(D\) & \hyperref[sec:annex-p02-w02-p02-017-finite-input-scope]{\ref*{sec:annex-p02-w02-p02-017-finite-input-scope}} & n/a & \hyperref[sec:annex-p02-w02-p02-017-finite-input-scope]{\ref*{sec:annex-p02-w02-p02-017-finite-input-scope}} & \thepage & F(x) = varphi\_x(x); T(x,y) = T\_1(e*,x,y) from the cited normal form; W least-witness relation; B,C,D,M bounded-search functions; t,d code-number maps. Local to the annex section; F is not the filter F of sec:p02-free-boolean-completeness. Explicit displayed family; member types and common actual first-use exposed, not one atomic glyph. Family listing, not one atomic glyph. Author macro \texttt{\detokenize{F,}\allowbreak\detokenize{e\^{}{*},}\allowbreak\detokenize{T,}\allowbreak\detokenize{W,}\allowbreak\detokenize{N,}\allowbreak\detokenize{B,}\allowbreak\detokenize{C,}\allowbreak\detokenize{D,}\allowbreak\detokenize{M,}\allowbreak\detokenize{t,}\allowbreak\detokenize{d,}\allowbreak\detokenize{d\_0,}\allowbreak\detokenize{d\_1,}\allowbreak\detokenize{h}} not carried to the members. \\
\phantomsection\label{entry:symbol-registry-07136}\(M\) & \texttt{\detokenize{M}} & local computable or partial computable functions and decidable relations in the annex constructions symbol family member \(M\) & \hyperref[sec:annex-p02-w02-p02-017-finite-input-scope]{\ref*{sec:annex-p02-w02-p02-017-finite-input-scope}} & n/a & \hyperref[sec:annex-p02-w02-p02-017-finite-input-scope]{\ref*{sec:annex-p02-w02-p02-017-finite-input-scope}} & \thepage & F(x) = varphi\_x(x); T(x,y) = T\_1(e*,x,y) from the cited normal form; W least-witness relation; B,C,D,M bounded-search functions; t,d code-number maps. Local to the annex section; F is not the filter F of sec:p02-free-boolean-completeness. Explicit displayed family; member types and common actual first-use exposed, not one atomic glyph. Family listing, not one atomic glyph. Author macro \texttt{\detokenize{F,}\allowbreak\detokenize{e\^{}{*},}\allowbreak\detokenize{T,}\allowbreak\detokenize{W,}\allowbreak\detokenize{N,}\allowbreak\detokenize{B,}\allowbreak\detokenize{C,}\allowbreak\detokenize{D,}\allowbreak\detokenize{M,}\allowbreak\detokenize{t,}\allowbreak\detokenize{d,}\allowbreak\detokenize{d\_0,}\allowbreak\detokenize{d\_1,}\allowbreak\detokenize{h}} not carried to the members. \\
\phantomsection\label{entry:symbol-registry-07137}\(t\) & \texttt{\detokenize{t}} & local computable or partial computable functions and decidable relations in the annex constructions symbol family member \(t\) & \hyperref[sec:annex-p02-w02-p02-017-finite-input-scope]{\ref*{sec:annex-p02-w02-p02-017-finite-input-scope}} & n/a & \hyperref[sec:annex-p02-w02-p02-017-finite-input-scope]{\ref*{sec:annex-p02-w02-p02-017-finite-input-scope}} & \thepage & F(x) = varphi\_x(x); T(x,y) = T\_1(e*,x,y) from the cited normal form; W least-witness relation; B,C,D,M bounded-search functions; t,d code-number maps. Local to the annex section; F is not the filter F of sec:p02-free-boolean-completeness. Explicit displayed family; member types and common actual first-use exposed, not one atomic glyph. Family listing, not one atomic glyph. Author macro \texttt{\detokenize{F,}\allowbreak\detokenize{e\^{}{*},}\allowbreak\detokenize{T,}\allowbreak\detokenize{W,}\allowbreak\detokenize{N,}\allowbreak\detokenize{B,}\allowbreak\detokenize{C,}\allowbreak\detokenize{D,}\allowbreak\detokenize{M,}\allowbreak\detokenize{t,}\allowbreak\detokenize{d,}\allowbreak\detokenize{d\_0,}\allowbreak\detokenize{d\_1,}\allowbreak\detokenize{h}} not carried to the members. \\
\phantomsection\label{entry:symbol-registry-07138}\(d\) & \texttt{\detokenize{d}} & local computable or partial computable functions and decidable relations in the annex constructions symbol family member \(d\) & \hyperref[sec:annex-p02-w02-p02-017-finite-input-scope]{\ref*{sec:annex-p02-w02-p02-017-finite-input-scope}} & n/a & \hyperref[sec:annex-p02-w02-p02-017-finite-input-scope]{\ref*{sec:annex-p02-w02-p02-017-finite-input-scope}} & \thepage & F(x) = varphi\_x(x); T(x,y) = T\_1(e*,x,y) from the cited normal form; W least-witness relation; B,C,D,M bounded-search functions; t,d code-number maps. Local to the annex section; F is not the filter F of sec:p02-free-boolean-completeness. Explicit displayed family; member types and common actual first-use exposed, not one atomic glyph. Family listing, not one atomic glyph. Author macro \texttt{\detokenize{F,}\allowbreak\detokenize{e\^{}{*},}\allowbreak\detokenize{T,}\allowbreak\detokenize{W,}\allowbreak\detokenize{N,}\allowbreak\detokenize{B,}\allowbreak\detokenize{C,}\allowbreak\detokenize{D,}\allowbreak\detokenize{M,}\allowbreak\detokenize{t,}\allowbreak\detokenize{d,}\allowbreak\detokenize{d\_0,}\allowbreak\detokenize{d\_1,}\allowbreak\detokenize{h}} not carried to the members. \\
\phantomsection\label{entry:symbol-registry-07139}\(d_0\) & \texttt{\detokenize{d_0}} & local computable or partial computable functions and decidable relations in the annex constructions symbol family member \(d_0\) & \hyperref[sec:annex-p02-w02-p02-017-finite-input-scope]{\ref*{sec:annex-p02-w02-p02-017-finite-input-scope}} & n/a & \hyperref[sec:annex-p02-w02-p02-017-finite-input-scope]{\ref*{sec:annex-p02-w02-p02-017-finite-input-scope}} & \thepage & F(x) = varphi\_x(x); T(x,y) = T\_1(e*,x,y) from the cited normal form; W least-witness relation; B,C,D,M bounded-search functions; t,d code-number maps. Local to the annex section; F is not the filter F of sec:p02-free-boolean-completeness. Explicit displayed family; member types and common actual first-use exposed, not one atomic glyph. Family listing, not one atomic glyph. Author macro \texttt{\detokenize{F,}\allowbreak\detokenize{e\^{}{*},}\allowbreak\detokenize{T,}\allowbreak\detokenize{W,}\allowbreak\detokenize{N,}\allowbreak\detokenize{B,}\allowbreak\detokenize{C,}\allowbreak\detokenize{D,}\allowbreak\detokenize{M,}\allowbreak\detokenize{t,}\allowbreak\detokenize{d,}\allowbreak\detokenize{d\_0,}\allowbreak\detokenize{d\_1,}\allowbreak\detokenize{h}} not carried to the members. \\
\phantomsection\label{entry:symbol-registry-07140}\(d_1\) & \texttt{\detokenize{d_1}} & local computable or partial computable functions and decidable relations in the annex constructions symbol family member \(d_1\) & \hyperref[sec:annex-p02-w02-p02-017-finite-input-scope]{\ref*{sec:annex-p02-w02-p02-017-finite-input-scope}} & n/a & \hyperref[sec:annex-p02-w02-p02-017-finite-input-scope]{\ref*{sec:annex-p02-w02-p02-017-finite-input-scope}} & \thepage & F(x) = varphi\_x(x); T(x,y) = T\_1(e*,x,y) from the cited normal form; W least-witness relation; B,C,D,M bounded-search functions; t,d code-number maps. Local to the annex section; F is not the filter F of sec:p02-free-boolean-completeness. Explicit displayed family; member types and common actual first-use exposed, not one atomic glyph. Family listing, not one atomic glyph. Author macro \texttt{\detokenize{F,}\allowbreak\detokenize{e\^{}{*},}\allowbreak\detokenize{T,}\allowbreak\detokenize{W,}\allowbreak\detokenize{N,}\allowbreak\detokenize{B,}\allowbreak\detokenize{C,}\allowbreak\detokenize{D,}\allowbreak\detokenize{M,}\allowbreak\detokenize{t,}\allowbreak\detokenize{d,}\allowbreak\detokenize{d\_0,}\allowbreak\detokenize{d\_1,}\allowbreak\detokenize{h}} not carried to the members. \\
\phantomsection\label{entry:symbol-registry-07141}\(h\) & \texttt{\detokenize{h}} & local computable or partial computable functions and decidable relations in the annex constructions symbol family member \(h\) & \hyperref[sec:annex-p02-w02-p02-017-finite-input-scope]{\ref*{sec:annex-p02-w02-p02-017-finite-input-scope}} & n/a & \hyperref[sec:annex-p02-w02-p02-017-finite-input-scope]{\ref*{sec:annex-p02-w02-p02-017-finite-input-scope}} & \thepage & F(x) = varphi\_x(x); T(x,y) = T\_1(e*,x,y) from the cited normal form; W least-witness relation; B,C,D,M bounded-search functions; t,d code-number maps. Local to the annex section; F is not the filter F of sec:p02-free-boolean-completeness. Explicit displayed family; member types and common actual first-use exposed, not one atomic glyph. Family listing, not one atomic glyph. Author macro \texttt{\detokenize{F,}\allowbreak\detokenize{e\^{}{*},}\allowbreak\detokenize{T,}\allowbreak\detokenize{W,}\allowbreak\detokenize{N,}\allowbreak\detokenize{B,}\allowbreak\detokenize{C,}\allowbreak\detokenize{D,}\allowbreak\detokenize{M,}\allowbreak\detokenize{t,}\allowbreak\detokenize{d,}\allowbreak\detokenize{d\_0,}\allowbreak\detokenize{d\_1,}\allowbreak\detokenize{h}} not carried to the members. \\
\phantomsection\label{entry:symbol-registry-07142}\(a\mathbin{\dot{-}}b\) & \texttt{\detokenize{a\mathbin{\dot{-}}b}} & cutoff subtraction predecessor zero-test and equality-test primitive recursive functions symbol family member \(a\mathbin{\dot{-}}b\) & \hyperref[sec:annex-p02-w02-p02-017-finite-input-scope]{\ref*{sec:annex-p02-w02-p02-017-finite-input-scope}} & n/a & \hyperref[sec:annex-p02-w02-p02-017-finite-input-scope]{\ref*{sec:annex-p02-w02-p02-017-finite-input-scope}} & \thepage & Numerical functions; the dotted minus has the P12 cutoff-subtraction meaning, not the P8 Brouwerian difference. Explicit displayed family; member types and common actual first-use exposed, not one atomic glyph. Family listing, not one atomic glyph. Author macro \texttt{\detokenize{\mathbin{\dot{-}},}\allowbreak\detokenize{\operatorname{pd},}\allowbreak\detokenize{\overline{\operatorname{sg}},}\allowbreak\detokenize{\operatorname{eq}}} not carried to the members. \\
\phantomsection\label{entry:symbol-registry-07143}\(\operatorname{pd}\) & \texttt{\detokenize{\operatorname{pd}}} & cutoff subtraction predecessor zero-test and equality-test primitive recursive functions symbol family member \(\operatorname{pd}\) & \hyperref[sec:annex-p02-w02-p02-017-finite-input-scope]{\ref*{sec:annex-p02-w02-p02-017-finite-input-scope}} & n/a & \hyperref[sec:annex-p02-w02-p02-017-finite-input-scope]{\ref*{sec:annex-p02-w02-p02-017-finite-input-scope}} & \thepage & Numerical functions; the dotted minus has the P12 cutoff-subtraction meaning, not the P8 Brouwerian difference. Explicit displayed family; member types and common actual first-use exposed, not one atomic glyph. Family listing, not one atomic glyph. Author macro \texttt{\detokenize{\mathbin{\dot{-}},}\allowbreak\detokenize{\operatorname{pd},}\allowbreak\detokenize{\overline{\operatorname{sg}},}\allowbreak\detokenize{\operatorname{eq}}} not carried to the members. \\
\phantomsection\label{entry:symbol-registry-07144}\(\overline{\operatorname{sg}}\) & \texttt{\detokenize{\overline{\operatorname{sg}}}} & cutoff subtraction predecessor zero-test and equality-test primitive recursive functions symbol family member \(\overline{\operatorname{sg}}\) & \hyperref[sec:annex-p02-w02-p02-017-finite-input-scope]{\ref*{sec:annex-p02-w02-p02-017-finite-input-scope}} & n/a & \hyperref[sec:annex-p02-w02-p02-017-finite-input-scope]{\ref*{sec:annex-p02-w02-p02-017-finite-input-scope}} & \thepage & Numerical functions; the dotted minus has the P12 cutoff-subtraction meaning, not the P8 Brouwerian difference. Explicit displayed family; member types and common actual first-use exposed, not one atomic glyph. Family listing, not one atomic glyph. Author macro \texttt{\detokenize{\mathbin{\dot{-}},}\allowbreak\detokenize{\operatorname{pd},}\allowbreak\detokenize{\overline{\operatorname{sg}},}\allowbreak\detokenize{\operatorname{eq}}} not carried to the members. \\
\phantomsection\label{entry:symbol-registry-07145}\(\operatorname{eq}\) & \texttt{\detokenize{\operatorname{eq}}} & cutoff subtraction predecessor zero-test and equality-test primitive recursive functions symbol family member \(\operatorname{eq}\) & \hyperref[sec:annex-p02-w02-p02-017-finite-input-scope]{\ref*{sec:annex-p02-w02-p02-017-finite-input-scope}} & n/a & \hyperref[sec:annex-p02-w02-p02-017-finite-input-scope]{\ref*{sec:annex-p02-w02-p02-017-finite-input-scope}} & \thepage & Numerical functions; the dotted minus has the P12 cutoff-subtraction meaning, not the P8 Brouwerian difference. Explicit displayed family; member types and common actual first-use exposed, not one atomic glyph. Family listing, not one atomic glyph. Author macro \texttt{\detokenize{\mathbin{\dot{-}},}\allowbreak\detokenize{\operatorname{pd},}\allowbreak\detokenize{\overline{\operatorname{sg}},}\allowbreak\detokenize{\operatorname{eq}}} not carried to the members. \\
\phantomsection\label{entry:symbol-registry-07146}\(\Gamma_{\mathcal A}\) & \texttt{\detokenize{\Gamma_{\mathcal A}}} & pair-selection premises for a family of disjoint pairs and the valuation/choice-function correspondence symbol family member \(\Gamma_{\mathcal A}\) & \hyperref[thm:p02-w02-p02-017-pairs]{\ref*{thm:p02-w02-p02-017-pairs}} & n/a & \hyperref[thm:p02-w02-p02-017-pairs]{\ref*{thm:p02-w02-p02-017-pairs}} & \thepage & A\_a is the member of the family containing a. Explicit displayed family; member types and common actual first-use exposed, not one atomic glyph. Family listing, not one atomic glyph. Author macro \texttt{\detokenize{\Gamma\_{\mathcal A},}\allowbreak\detokenize{A\_a,}\allowbreak\detokenize{g\_v,}\allowbreak\detokenize{v\_g}} not carried to the members. \\
\phantomsection\label{entry:symbol-registry-07147}\(A_a\) & \texttt{\detokenize{A_a}} & pair-selection premises for a family of disjoint pairs and the valuation/choice-function correspondence symbol family member \(A_a\) & \hyperref[thm:p02-w02-p02-017-pairs]{\ref*{thm:p02-w02-p02-017-pairs}} & n/a & \hyperref[thm:p02-w02-p02-017-pairs]{\ref*{thm:p02-w02-p02-017-pairs}} & \thepage & A\_a is the member of the family containing a. Explicit displayed family; member types and common actual first-use exposed, not one atomic glyph. Family listing, not one atomic glyph. Author macro \texttt{\detokenize{\Gamma\_{\mathcal A},}\allowbreak\detokenize{A\_a,}\allowbreak\detokenize{g\_v,}\allowbreak\detokenize{v\_g}} not carried to the members. \\
\phantomsection\label{entry:symbol-registry-07148}\(g_v\) & \texttt{\detokenize{g_v}} & pair-selection premises for a family of disjoint pairs and the valuation/choice-function correspondence symbol family member \(g_v\) & \hyperref[thm:p02-w02-p02-017-pairs]{\ref*{thm:p02-w02-p02-017-pairs}} & n/a & \hyperref[thm:p02-w02-p02-017-pairs]{\ref*{thm:p02-w02-p02-017-pairs}} & \thepage & A\_a is the member of the family containing a. Explicit displayed family; member types and common actual first-use exposed, not one atomic glyph. Family listing, not one atomic glyph. Author macro \texttt{\detokenize{\Gamma\_{\mathcal A},}\allowbreak\detokenize{A\_a,}\allowbreak\detokenize{g\_v,}\allowbreak\detokenize{v\_g}} not carried to the members. \\
\phantomsection\label{entry:symbol-registry-07149}\(v_g\) & \texttt{\detokenize{v_g}} & pair-selection premises for a family of disjoint pairs and the valuation/choice-function correspondence symbol family member \(v_g\) & \hyperref[thm:p02-w02-p02-017-pairs]{\ref*{thm:p02-w02-p02-017-pairs}} & n/a & \hyperref[thm:p02-w02-p02-017-pairs]{\ref*{thm:p02-w02-p02-017-pairs}} & \thepage & A\_a is the member of the family containing a. Explicit displayed family; member types and common actual first-use exposed, not one atomic glyph. Family listing, not one atomic glyph. Author macro \texttt{\detokenize{\Gamma\_{\mathcal A},}\allowbreak\detokenize{A\_a,}\allowbreak\detokenize{g\_v,}\allowbreak\detokenize{v\_g}} not carried to the members. \\
\phantomsection\label{entry:symbol-registry-07150}\(\sigma^{\bullet}\) & \texttt{\detokenize{\sigma^{\bullet}}} & translation from the P1 compactness language with top into Form(P) and the induced valuation symbol family member \(\sigma^{\bullet}\) & \hyperref[annex:p02-w02-p02-017-bpi-boundary]{\ref*{annex:p02-w02-p02-017-bpi-boundary}} & n/a & \hyperref[annex:p02-w02-p02-017-bpi-boundary]{\ref*{annex:p02-w02-p02-017-bpi-boundary}} & \thepage & top goes to not-bottom; v\_nu(p\_v) = nu(v). Explicit displayed family; member types and common actual first-use exposed, not one atomic glyph. Family listing, not one atomic glyph. Author macro \texttt{\detokenize{\sigma\^{}{\bullet},}\allowbreak\detokenize{v\_\nu}} not carried to the members. \\
\phantomsection\label{entry:symbol-registry-07151}\(v_\nu\) & \texttt{\detokenize{v_\nu}} & translation from the P1 compactness language with top into Form(P) and the induced valuation symbol family member \(v_\nu\) & \hyperref[annex:p02-w02-p02-017-bpi-boundary]{\ref*{annex:p02-w02-p02-017-bpi-boundary}} & n/a & \hyperref[annex:p02-w02-p02-017-bpi-boundary]{\ref*{annex:p02-w02-p02-017-bpi-boundary}} & \thepage & top goes to not-bottom; v\_nu(p\_v) = nu(v). Explicit displayed family; member types and common actual first-use exposed, not one atomic glyph. Family listing, not one atomic glyph. Author macro \texttt{\detokenize{\sigma\^{}{\bullet},}\allowbreak\detokenize{v\_\nu}} not carried to the members. \\
\phantomsection\label{entry:symbol-registry-07152}\(\iota\) & \texttt{\detokenize{\iota}} & Church's type symbol of individuals & \hyperref[def:p02-w02-p02-040-church-types]{\ref*{def:p02-w02-p02-040-church-types}} & n/a & \hyperref[def:p02-w02-p02-040-church-types]{\ref*{def:p02-w02-p02-040-church-types}} & \thepage & Type symbol (glyph cell previously read '\(\iota\) (type)'). Occurs only as a subscript; same glyph as the selection constant \(\iota_{\alpha(o\alpha)}\), distinguished by position (Church 1940 p. 57). Original precedence descriptor 'none' retained; no rank inferred. \\
\phantomsection\label{entry:symbol-registry-07153}\(\iota_{\alpha(o\alpha)}\) & \texttt{\detokenize{\iota}} & Church's selection (choice/description) constant of type \(\alpha(o\alpha)\) & \hyperref[def:p02-w02-p02-040-church-wff]{\ref*{def:p02-w02-p02-040-church-wff}} & n/a & \hyperref[def:p02-w02-p02-040-church-wff]{\ref*{def:p02-w02-p02-040-church-wff}} & \thepage & Logical constant governed by Axioms 9 and 11; translated to the nonlogical host constant \(\varepsilon_\alpha\). Original precedence descriptor 'none' retained; no rank inferred. \\
\phantomsection\label{entry:symbol-registry-07154}\(o\) & \texttt{\detokenize{o}} & Church's type symbol of propositions & \hyperref[def:p02-w02-p02-040-church-types]{\ref*{def:p02-w02-p02-040-church-types}} & n/a & \hyperref[def:p02-w02-p02-040-church-types]{\ref*{def:p02-w02-p02-040-church-types}} & \thepage & Same role as host type o. Original precedence descriptor 'none' retained; no rank inferred. \\
\phantomsection\label{entry:symbol-registry-07155}\((\alpha\beta)\) & \texttt{\detokenize{(\alpha\beta)}} & Church's function type: values in \(\alpha\) and arguments in \(\beta\) & \hyperref[def:p02-w02-p02-040-church-types]{\ref*{def:p02-w02-p02-040-church-types}} & n/a & \hyperref[def:p02-w02-p02-040-church-types]{\ref*{def:p02-w02-p02-040-church-types}} & \thepage & Value type first; host image \(\beta^\circ\to\alpha^\circ\). Original precedence descriptor 'association to the left' retained; no rank inferred. Blank macro filled with the glyph's TeX spelling. \\
\phantomsection\label{entry:symbol-registry-07156}\(\alpha'\) & \texttt{\detokenize{'}} & Church's numeral type \(((\alpha\alpha)(\alpha\alpha))\) & \hyperref[def:p02-w02-p02-040-church-types]{\ref*{def:p02-w02-p02-040-church-types}} & n/a & \hyperref[def:p02-w02-p02-040-church-types]{\ref*{def:p02-w02-p02-040-church-types}} & \thepage & Iterated: \(\alpha''\), \(\alpha'''\). Original precedence descriptor 'postfix' retained; no rank inferred. \\
\phantomsection\label{entry:symbol-registry-07157}\(N_{oo}\) & \texttt{\detokenize{N}} & Church's negation constant & \hyperref[def:p02-w02-p02-040-church-wff]{\ref*{def:p02-w02-p02-040-church-wff}} & n/a & \hyperref[def:p02-w02-p02-040-church-wff]{\ref*{def:p02-w02-p02-040-church-wff}} & \thepage & Same letter as the number predicate \(N_{o\alpha'}\); the subscript distinguishes them. Original precedence descriptor 'none' retained; no rank inferred. \\
\phantomsection\label{entry:symbol-registry-07158}\(N_{o\alpha'}\) & \texttt{\detokenize{N}} & Church's defined predicate 'is a natural number of type \(\alpha'\)' & \hyperref[def:p02-w02-p02-040-church-abbrev]{\ref*{def:p02-w02-p02-040-church-abbrev}} & n/a & \hyperref[def:p02-w02-p02-040-church-abbrev]{\ref*{def:p02-w02-p02-040-church-abbrev}} & \thepage & Abbreviation (Church 1940 p. 58). Original precedence descriptor 'none' retained; no rank inferred. \\
\phantomsection\label{entry:symbol-registry-07159}\(A_{ooo}\) & \texttt{\detokenize{A}} & Church's disjunction constant & \hyperref[def:p02-w02-p02-040-church-wff]{\ref*{def:p02-w02-p02-040-church-wff}} & n/a & \hyperref[def:p02-w02-p02-040-church-wff]{\ref*{def:p02-w02-p02-040-church-wff}} & \thepage & Bold \(\mathbf A\) is a syntactic variable for wffs; italic subscripted A is the constant. Original precedence descriptor 'none' retained; no rank inferred. \\
\phantomsection\label{entry:symbol-registry-07160}\(\Pi_{o(o\alpha)}\) & \texttt{\detokenize{\Pi}} & Church's universal-quantifier constant & \hyperref[def:p02-w02-p02-040-church-wff]{\ref*{def:p02-w02-p02-040-church-wff}} & n/a & \hyperref[def:p02-w02-p02-040-church-wff]{\ref*{def:p02-w02-p02-040-church-wff}} & \thepage & Host image \(\forall_{\alpha^\circ}\). Original precedence descriptor 'none' retained; no rank inferred. \\
\phantomsection\label{entry:symbol-registry-07161}\(\supset\) & \texttt{\detokenize{\supset}} & Church's defined material implication & \hyperref[def:p02-w02-p02-040-church-abbrev]{\ref*{def:p02-w02-p02-040-church-abbrev}} & n/a & \hyperref[def:p02-w02-p02-040-church-abbrev]{\ref*{def:p02-w02-p02-040-church-abbrev}} & \thepage & Source notation; host image \(\to\). Precedence (source descriptor): highest bracket rank (with \(\equiv\)). \\
\phantomsection\label{entry:symbol-registry-07162}\(\sim\) & \texttt{\detokenize{\sim}} & Church's negation (abbreviation of \(N_{oo}\)) & \hyperref[def:p02-w02-p02-040-church-abbrev]{\ref*{def:p02-w02-p02-040-church-abbrev}} & n/a & \hyperref[def:p02-w02-p02-040-church-abbrev]{\ref*{def:p02-w02-p02-040-church-abbrev}} & \thepage & Source notation. Original precedence descriptor 'lowest bracket rank' retained; no rank inferred. \\
\phantomsection\label{entry:symbol-registry-07163}\((\mathbf x_\alpha)\) & \texttt{\detokenize{(\mathbf x_\alpha)}} & Church's universal quantifier notation & \hyperref[def:p02-w02-p02-040-church-abbrev]{\ref*{def:p02-w02-p02-040-church-abbrev}} & n/a & \hyperref[def:p02-w02-p02-040-church-abbrev]{\ref*{def:p02-w02-p02-040-church-abbrev}} & \thepage & Abbreviates \(\Pi_{o(o\alpha)}(\lambda\mathbf x_\alpha\mathbf A_o)\). Original precedence descriptor 'second bracket rank' retained; no rank inferred. Blank macro filled with the glyph's TeX spelling. \\
\phantomsection\label{entry:symbol-registry-07164}\(\text{\rotatebox[origin=c]{180}{\ensuremath{\iota}}}\) & \texttt{\detokenize{\text{\rotatebox[origin=c]{180}{\ensuremath{\iota}}}}} & Church's description operator notation & \hyperref[def:p02-w02-p02-040-church-abbrev]{\ref*{def:p02-w02-p02-040-church-abbrev}} & n/a & \hyperref[def:p02-w02-p02-040-church-abbrev]{\ref*{def:p02-w02-p02-040-church-abbrev}} & \thepage & Turned iota, as typeset in the body. Abbreviates \(\iota_{\alpha(o\alpha)}(\lambda\mathbf x_\alpha\mathbf A_o)\). Original precedence descriptor 'second bracket rank' retained; no rank inferred. \\
\phantomsection\label{entry:symbol-registry-07165}\(Q_{o\alpha\alpha}\) & \texttt{\detokenize{Q}} & Church's defined Leibniz equality & \hyperref[def:p02-w02-p02-040-church-abbrev]{\ref*{def:p02-w02-p02-040-church-abbrev}} & n/a & \hyperref[def:p02-w02-p02-040-church-abbrev]{\ref*{def:p02-w02-p02-040-church-abbrev}} & \thepage & \([\mathbf A=\mathbf B]\) abbreviates \(Q\mathbf A\mathbf B\). Original precedence descriptor 'none' retained; no rank inferred. \\
\phantomsection\label{entry:symbol-registry-07166}\(I_{\alpha\alpha}\) & \texttt{\detokenize{I_{\alpha\alpha}}} & Church's identity combinator & \hyperref[def:p02-w02-p02-040-church-abbrev]{\ref*{def:p02-w02-p02-040-church-abbrev}} & n/a & \hyperref[def:p02-w02-p02-040-church-abbrev]{\ref*{def:p02-w02-p02-040-church-abbrev}} & \thepage & Abbreviations. Original precedence descriptor 'none' retained; no rank inferred. Blank macro filled with the glyph's TeX spelling. \\
\phantomsection\label{entry:symbol-registry-07167}\(K_{\alpha\beta\alpha}\) & \texttt{\detokenize{K_{\alpha\beta\alpha}}} & Church's constant-function combinator & \hyperref[def:p02-w02-p02-040-church-abbrev]{\ref*{def:p02-w02-p02-040-church-abbrev}} & n/a & \hyperref[def:p02-w02-p02-040-church-abbrev]{\ref*{def:p02-w02-p02-040-church-abbrev}} & \thepage & Abbreviations. Original precedence descriptor 'none' retained; no rank inferred. Blank macro filled with the glyph's TeX spelling. \\
\phantomsection\label{entry:symbol-registry-07168}\(0_{\alpha'}\) & \texttt{\detokenize{0_{\alpha'}}} & Church numeral zero at type \(\alpha'\) & \hyperref[def:p02-w02-p02-040-church-abbrev]{\ref*{def:p02-w02-p02-040-church-abbrev}} & n/a & \hyperref[def:p02-w02-p02-040-church-abbrev]{\ref*{def:p02-w02-p02-040-church-abbrev}} & \thepage & Abbreviation; likewise \(1_{\alpha'},2_{\alpha'},\ldots\). Not the arithmetic 0 of def:p02-pa2. Original precedence descriptor 'none' retained; no rank inferred. Blank macro filled with the glyph's TeX spelling. \\
\phantomsection\label{entry:symbol-registry-07169}\(S_{\alpha'\alpha'}\) & \texttt{\detokenize{S_{\alpha'\alpha'}}} & Church's successor at type \(\alpha'\) & \hyperref[def:p02-w02-p02-040-church-abbrev]{\ref*{def:p02-w02-p02-040-church-abbrev}} & n/a & \hyperref[def:p02-w02-p02-040-church-abbrev]{\ref*{def:p02-w02-p02-040-church-abbrev}} & \thepage & Abbreviation. Not the arithmetic S of def:p02-pa2. Original precedence descriptor 'none' retained; no rank inferred. Blank macro filled with the glyph's TeX spelling. \\
\phantomsection\label{entry:symbol-registry-07170}\(\omega_{\alpha''\alpha'\alpha'}\) & \texttt{\detokenize{\omega_{\alpha''\alpha'\alpha'}}} & Church's pairing combinator & \hyperref[def:p02-w02-p02-040-church-abbrev]{\ref*{def:p02-w02-p02-040-church-abbrev}} & n/a & \hyperref[def:p02-w02-p02-040-church-abbrev]{\ref*{def:p02-w02-p02-040-church-abbrev}} & \thepage & Abbreviation (Church 1940 p. 58). Original precedence descriptor 'none' retained; no rank inferred. Blank macro filled with the glyph's TeX spelling. \\
\phantomsection\label{entry:symbol-registry-07171}\(\langle\cdot,\cdot\rangle\) & \texttt{\detokenize{\langle\cdot,}\allowbreak\detokenize{\cdot\rangle}} & Church's pair notation: \(\langle\mathbf A_{\alpha'},\mathbf B_{\alpha'}\rangle\) abbreviates \(\omega_{\alpha''\alpha'\alpha'}\mathbf A_{\alpha'}\mathbf B_{\alpha'}\) & \hyperref[def:p02-w02-p02-040-church-abbrev]{\ref*{def:p02-w02-p02-040-church-abbrev}} & n/a & \hyperref[def:p02-w02-p02-040-church-abbrev]{\ref*{def:p02-w02-p02-040-church-abbrev}} & \thepage & Abbreviation (Church 1940 p. 58). Original precedence descriptor 'none' retained; no rank inferred. \\
\phantomsection\label{entry:symbol-registry-07172}\(P_{\alpha'\alpha'''}\) & \texttt{\detokenize{P_{\alpha'\alpha'''}}} & Kleene-style predecessor from type \(\alpha'''\) into type \(\alpha'\) & \hyperref[def:p02-w02-p02-040-church-abbrev]{\ref*{def:p02-w02-p02-040-church-abbrev}} & n/a & \hyperref[def:p02-w02-p02-040-church-abbrev]{\ref*{def:p02-w02-p02-040-church-abbrev}} & \thepage & Abbreviation (Church 1940 p. 58). Original precedence descriptor 'none' retained; no rank inferred. Blank macro filled with the glyph's TeX spelling. \\
\phantomsection\label{entry:symbol-registry-07173}\(P_{\alpha'\alpha'}\) & \texttt{\detokenize{P_{\alpha'\alpha'}}} & Church's predecessor within type \(\alpha'\) & \hyperref[def:p02-w02-p02-040-church-abbrev]{\ref*{def:p02-w02-p02-040-church-abbrev}} & n/a & \hyperref[def:p02-w02-p02-040-church-abbrev]{\ref*{def:p02-w02-p02-040-church-abbrev}} & \thepage & Abbreviation (Church 1940 p. 58). Original precedence descriptor 'none' retained; no rank inferred. Blank macro filled with the glyph's TeX spelling. \\
\phantomsection\label{entry:symbol-registry-07174}\(T_{\alpha''\alpha'}\) & \texttt{\detokenize{T_{\alpha''\alpha'}}} & Church's type-raising map on natural numbers from type \(\alpha'\) to type \(\alpha''\) & \hyperref[def:p02-w02-p02-040-church-abbrev]{\ref*{def:p02-w02-p02-040-church-abbrev}} & n/a & \hyperref[def:p02-w02-p02-040-church-abbrev]{\ref*{def:p02-w02-p02-040-church-abbrev}} & \thepage & Abbreviation (Church 1940 p. 58). Original precedence descriptor 'none' retained; no rank inferred. Blank macro filled with the glyph's TeX spelling. \\
\phantomsection\label{entry:symbol-registry-07175}\(\derives\) & \texttt{\detokenize{\derives}} & Church's derivability on assumptions with the substitution/generalization restriction & \hyperref[def:p02-w02-p02-040-church-axioms]{\ref*{def:p02-w02-p02-040-church-axioms}} & relation & \hyperref[def:p02-w02-p02-040-church-axioms]{\ref*{def:p02-w02-p02-040-church-axioms}} & \thepage & Local typed role; Church distinguishes it from Frege's assertion sign (p. 62). \\
\phantomsection\label{entry:symbol-registry-07176}\(\mathcal S(D,c)\) & \texttt{\detokenize{\mathcal S(D,}\allowbreak\detokenize{c)}} & Standard model of Church's system & \hyperref[def:p02-w02-p02-040-church-standard]{\ref*{def:p02-w02-p02-040-church-standard}} & n/a & \hyperref[def:p02-w02-p02-040-church-standard]{\ref*{def:p02-w02-p02-040-church-standard}} & \thepage & Handbook formalization of the source's informal interpretation. Original precedence descriptor 'none' retained; no rank inferred. Blank macro filled with the glyph's TeX spelling. \\
\phantomsection\label{entry:symbol-registry-07177}\(\mathcal D_\gamma\) & \texttt{\detokenize{\mathcal D_\gamma}} & Type domains of a standard model of Church's system & \hyperref[def:p02-w02-p02-040-church-standard]{\ref*{def:p02-w02-p02-040-church-standard}} & n/a & \hyperref[def:p02-w02-p02-040-church-standard]{\ref*{def:p02-w02-p02-040-church-standard}} & \thepage & Handbook formalization of the source's informal interpretation. Original precedence descriptor 'none' retained; no rank inferred. Blank macro filled with the glyph's TeX spelling. \\
\phantomsection\label{entry:symbol-registry-07178}\(\mathrm T\) & \texttt{\detokenize{\mathrm T}} & Truth value true of \(\mathcal D_o\) & \hyperref[def:p02-w02-p02-040-church-standard]{\ref*{def:p02-w02-p02-040-church-standard}} & n/a & \hyperref[def:p02-w02-p02-040-church-standard]{\ref*{def:p02-w02-p02-040-church-standard}} & \thepage & Identified with host 1 and 0 in the bridge. Original precedence descriptor 'none' retained; no rank inferred. Blank macro filled with the glyph's TeX spelling. \\
\phantomsection\label{entry:symbol-registry-07179}\(\mathrm F\) & \texttt{\detokenize{\mathrm F}} & Truth value false of \(\mathcal D_o\) & \hyperref[def:p02-w02-p02-040-church-standard]{\ref*{def:p02-w02-p02-040-church-standard}} & n/a & \hyperref[def:p02-w02-p02-040-church-standard]{\ref*{def:p02-w02-p02-040-church-standard}} & \thepage & Identified with host 1 and 0 in the bridge. Original precedence descriptor 'none' retained; no rank inferred. Blank macro filled with the glyph's TeX spelling. \\
\phantomsection\label{entry:symbol-registry-07180}\(\gamma^\circ\) & \texttt{\detokenize{^\circ}} & Church-to-host translation of types & \hyperref[def:p02-w02-p02-040-church-bridge]{\ref*{def:p02-w02-p02-040-church-bridge}} & n/a & \hyperref[def:p02-w02-p02-040-church-bridge]{\ref*{def:p02-w02-p02-040-church-bridge}} & \thepage & Metalinguistic map. Original precedence descriptor 'postfix' retained; no rank inferred. \\
\phantomsection\label{entry:symbol-registry-07181}\(\mathbf M^\circ\) & \texttt{\detokenize{^\circ}} & Church-to-host translation of wffs & \hyperref[def:p02-w02-p02-040-church-bridge]{\ref*{def:p02-w02-p02-040-church-bridge}} & n/a & \hyperref[def:p02-w02-p02-040-church-bridge]{\ref*{def:p02-w02-p02-040-church-bridge}} & \thepage & Metalinguistic map. Original precedence descriptor 'postfix' retained; no rank inferred. \\
\phantomsection\label{entry:symbol-registry-07182}\(\varepsilon_\alpha\) & \texttt{\detokenize{\varepsilon}} & Host nonlogical constant translating \(\iota_{\alpha(o\alpha)}\) & \hyperref[def:p02-w02-p02-040-church-bridge]{\ref*{def:p02-w02-p02-040-church-bridge}} & n/a & \hyperref[def:p02-w02-p02-040-church-bridge]{\ref*{def:p02-w02-p02-040-church-bridge}} & \thepage & Original precedence descriptor 'none' retained; no rank inferred. \\
\phantomsection\label{entry:symbol-registry-07183}\(L_{\mathcal F}(x)\) & \texttt{\detokenize{L}} & Rational cut of an element of a complete ordered field & \hyperref[thm:p02-w02-p02-040-reals-categorical]{\ref*{thm:p02-w02-p02-040-reals-categorical}} & n/a & \hyperref[thm:p02-w02-p02-040-reals-categorical]{\ref*{thm:p02-w02-p02-040-reals-categorical}} & \thepage & Original precedence descriptor 'none' retained; no rank inferred. \\
\phantomsection\label{entry:symbol-registry-07184}\(q_{\mathcal F}\) & \texttt{\detokenize{q_{\mathcal F}}} & Image of a rational in an ordered field & \hyperref[thm:p02-w02-p02-040-reals-categorical]{\ref*{thm:p02-w02-p02-040-reals-categorical}} & n/a & \hyperref[thm:p02-w02-p02-040-reals-categorical]{\ref*{thm:p02-w02-p02-040-reals-categorical}} & \thepage & Original precedence descriptor 'none' retained; no rank inferred. Blank macro filled with the glyph's TeX spelling. \\
\phantomsection\label{entry:symbol-registry-07185}\(\mathrm{Val}_2\) & \texttt{\detokenize{\mathrm{Val}_2}} & Code set of full second-order validities & \hyperref[def:p02-w02-p02-040-so-coding]{\ref*{def:p02-w02-p02-040-so-coding}} & n/a & \hyperref[def:p02-w02-p02-040-so-coding]{\ref*{def:p02-w02-p02-040-so-coding}} & \thepage & Original precedence descriptor 'none' retained; no rank inferred. Blank macro filled with the glyph's TeX spelling. \\
\phantomsection\label{entry:symbol-registry-07186}\(\mathrm{TA}\) & \texttt{\detokenize{\mathrm{TA}}} & Code set of true arithmetic sentences & \hyperref[def:p02-w02-p02-040-so-coding]{\ref*{def:p02-w02-p02-040-so-coding}} & n/a & \hyperref[def:p02-w02-p02-040-so-coding]{\ref*{def:p02-w02-p02-040-so-coding}} & \thepage & Original precedence descriptor 'none' retained; no rank inferred. Blank macro filled with the glyph's TeX spelling. \\
\phantomsection\label{entry:symbol-registry-07187}\(N_{15},\ldots,N_{18}\) & \texttt{\detokenize{N}} & Constructor codes for second-order arithmetic syntax & \hyperref[def:p02-w02-p02-040-so-coding]{\ref*{def:p02-w02-p02-040-so-coding}} & n/a & \hyperref[def:p02-w02-p02-040-so-coding]{\ref*{def:p02-w02-p02-040-so-coding}} & \thepage & Extends the constructor \(N_k\) of sec:arithmetic-coding-04; unrelated to Church's \(N_{oo}\). Original precedence descriptor 'none' retained; no rank inferred. Sequence glyph kept as one row; the ellipsis is not a member symbol. \\
\phantomsection\label{entry:symbol-registry-07188}\(\mathfrak s_{\alpha'\alpha'\alpha'}\) & \texttt{\detokenize{\mathfrak s}} & Rosser's addition combinator in Church's system & \hyperref[thm:p02-w02-p02-040-church-rosser]{\ref*{thm:p02-w02-p02-040-church-rosser}} & n/a & \hyperref[thm:p02-w02-p02-040-church-rosser]{\ref*{thm:p02-w02-p02-040-church-rosser}} & \thepage & Source glyphs are fraktur s and B (Church 1940 p. 68 n. 13); \([\mathbf A+\mathbf B]\) and \([\mathbf A\times\mathbf B]\) abbreviate their applications. Distinct from bold syntactic variable \(\mathbf B\). Original precedence descriptor 'none' retained; no rank inferred. \\
\phantomsection\label{entry:symbol-registry-07189}\(\mathfrak B_{\alpha'\alpha'\alpha'}\) & \texttt{\detokenize{\mathfrak B}} & Rosser's multiplication combinator in Church's system & \hyperref[thm:p02-w02-p02-040-church-rosser]{\ref*{thm:p02-w02-p02-040-church-rosser}} & n/a & \hyperref[thm:p02-w02-p02-040-church-rosser]{\ref*{thm:p02-w02-p02-040-church-rosser}} & \thepage & Source glyphs are fraktur s and B (Church 1940 p. 68 n. 13); \([\mathbf A+\mathbf B]\) and \([\mathbf A\times\mathbf B]\) abbreviate their applications. Distinct from bold syntactic variable \(\mathbf B\). Original precedence descriptor 'none' retained; no rank inferred. \\
\phantomsection\label{entry:symbol-registry-07190}\(+\) & \texttt{\detokenize{+}} & Rosser's defined addition at type \(\alpha'\) (Church) & \hyperref[thm:p02-w02-p02-040-church-rosser]{\ref*{thm:p02-w02-p02-040-church-rosser}} & n/a & \hyperref[thm:p02-w02-p02-040-church-rosser]{\ref*{thm:p02-w02-p02-040-church-rosser}} & \thepage & Bracketed abbreviation \([\mathbf A_{\alpha'}+\mathbf B_{\alpha'}]\); unrelated to the arithmetic + of def:p02-pa2. Original precedence descriptor 'none' retained; no rank inferred. \\
\phantomsection\label{entry:symbol-registry-07191}\(\times\) & \texttt{\detokenize{\times}} & Rosser's defined multiplication at type \(\alpha'\) (Church) & \hyperref[thm:p02-w02-p02-040-church-rosser]{\ref*{thm:p02-w02-p02-040-church-rosser}} & n/a & \hyperref[thm:p02-w02-p02-040-church-rosser]{\ref*{thm:p02-w02-p02-040-church-rosser}} & \thepage & Bracketed abbreviation \([\mathbf A_{\alpha'}\times\mathbf B_{\alpha'}]\); unrelated to the arithmetic multiplication of def:p02-pa2. Original precedence descriptor 'none' retained; no rank inferred. \\
\phantomsection\label{entry:symbol-registry-07192}\(\mathbf F_{\alpha'\alpha'}\) & \texttt{\detokenize{\mathbf F_{\alpha'\alpha'}}} & Church's primitive-recursion formula built from \(\mathbf A_{\alpha'}\) and \(\mathbf B_{\alpha'\alpha'\alpha'}\) & \hyperref[def:p02-w02-p02-040-church-recursor]{\ref*{def:p02-w02-p02-040-church-recursor}} & n/a & \hyperref[def:p02-w02-p02-040-church-recursor]{\ref*{def:p02-w02-p02-040-church-recursor}} & \thepage & Source notation (p. 68). Original precedence descriptor 'none' retained; no rank inferred. Blank macro filled with the glyph's TeX spelling. \\
\phantomsection\label{entry:symbol-registry-07193}\(26^\alpha,30^\alpha,31^\alpha,\ldots,43^\alpha\) & \texttt{\detokenize{26^\alpha,}\allowbreak\detokenize{30^\alpha,}\allowbreak\detokenize{31^\alpha,}\allowbreak\detokenize{\ldots,}\allowbreak\detokenize{43^\alpha}} & Church's theorem numbers with type superscripts & \hyperref[sec:p02-w02-p02-040-church]{\ref*{sec:p02-w02-p02-040-church}} & n/a & \hyperref[thm:p02-w02-p02-040-church-26-30]{\ref*{thm:p02-w02-p02-040-church-26-30}} & \thepage & Superscript is a type symbol, not an exponent (typical ambiguity). Original precedence descriptor 'none' retained; no rank inferred. Blank macro filled with the glyph's TeX spelling. Sequence glyph kept as one row; the ellipsis is not a member symbol. \\
\phantomsection\label{entry:symbol-registry-07194}\(\mathbf M^\flat\) & \texttt{\detokenize{^\flat}} & Abbreviation for \(\mathbf M\,S_{\alpha'\alpha'}0_{\alpha'}\) in the annex & \hyperref[sec:annex-p02-w02-p02-040]{\ref*{sec:annex-p02-w02-p02-040}} & n/a & \hyperref[sec:annex-p02-w02-p02-040]{\ref*{sec:annex-p02-w02-p02-040}} & \thepage & Handbook proof notation only; introduced in the unlabelled paragraph 'Conventions for the arithmetic derivations' preceding annex:p02-w02-p02-040-church-26-30. Original precedence descriptor 'postfix' retained; no rank inferred. \\
\phantomsection\label{entry:symbol-registry-07195}\(\derives_0\) & \texttt{\detokenize{\derives_0}} & Schematic derivability from (a1)--(a4) by modus ponens & \hyperref[annex:p02-w02-p02-040-church-prop-complete]{\ref*{annex:p02-w02-p02-040-church-prop-complete}} & relation & \hyperref[annex:p02-w02-p02-040-church-prop-complete]{\ref*{annex:p02-w02-p02-040-church-prop-complete}} & \thepage & Proof-local; distinct from Church's \(\derives\). \\
\phantomsection\label{entry:symbol-registry-07196}\(\mathrm{Fm}_2\) & \texttt{\detokenize{\mathrm{Fm}_2}} & Primitive-recursive test for second-order formula codes & \hyperref[annex:p02-w02-p02-040-no-effective-calculus]{\ref*{annex:p02-w02-p02-040-no-effective-calculus}} & n/a & \hyperref[annex:p02-w02-p02-040-no-effective-calculus]{\ref*{annex:p02-w02-p02-040-no-effective-calculus}} & \thepage & Outer tags 15-18 also occur as derivation, proof-code and alpha-record tags in sec:arithmetic-coding-04; disjoint by arity, shown in the annex. Original precedence descriptor 'none' retained; no rank inferred. Blank macro filled with the glyph's TeX spelling. \\
\phantomsection\label{entry:symbol-registry-07197}\(\mathrm{FV}_1\) & \texttt{\detokenize{\mathrm{FV}_1}} & Primitive-recursive list of free individual-variable indices of a second-order formula code & \hyperref[annex:p02-w02-p02-040-no-effective-calculus]{\ref*{annex:p02-w02-p02-040-no-effective-calculus}} & n/a & \hyperref[annex:p02-w02-p02-040-no-effective-calculus]{\ref*{annex:p02-w02-p02-040-no-effective-calculus}} & \thepage & Outer tags 15-18 also occur as derivation, proof-code and alpha-record tags in sec:arithmetic-coding-04; disjoint by arity, shown in the annex. Original precedence descriptor 'none' retained; no rank inferred. Blank macro filled with the glyph's TeX spelling. \\
\phantomsection\label{entry:symbol-registry-07198}\(\mathrm{FV}_2\) & \texttt{\detokenize{\mathrm{FV}_2}} & Primitive-recursive list of free relation-variable codes of a second-order formula code & \hyperref[annex:p02-w02-p02-040-no-effective-calculus]{\ref*{annex:p02-w02-p02-040-no-effective-calculus}} & n/a & \hyperref[annex:p02-w02-p02-040-no-effective-calculus]{\ref*{annex:p02-w02-p02-040-no-effective-calculus}} & \thepage & Outer tags 15-18 also occur as derivation, proof-code and alpha-record tags in sec:arithmetic-coding-04; disjoint by arity, shown in the annex. Original precedence descriptor 'none' retained; no rank inferred. Blank macro filled with the glyph's TeX spelling. \\
\phantomsection\label{entry:symbol-registry-07199}\(\mathrm{Sent}_2\) & \texttt{\detokenize{\mathrm{Sent}_2}} & Primitive-recursive test for second-order sentence codes & \hyperref[annex:p02-w02-p02-040-no-effective-calculus]{\ref*{annex:p02-w02-p02-040-no-effective-calculus}} & n/a & \hyperref[annex:p02-w02-p02-040-no-effective-calculus]{\ref*{annex:p02-w02-p02-040-no-effective-calculus}} & \thepage & Outer tags 15-18 also occur as derivation, proof-code and alpha-record tags in sec:arithmetic-coding-04; disjoint by arity, shown in the annex. Original precedence descriptor 'none' retained; no rank inferred. Blank macro filled with the glyph's TeX spelling. \\
\phantomsection\label{entry:symbol-registry-07200}\(\mathrm{RV}\) & \texttt{\detokenize{\mathrm{RV}}} & Test for a relation-variable code & \hyperref[annex:p02-w02-p02-040-no-effective-calculus]{\ref*{annex:p02-w02-p02-040-no-effective-calculus}} & n/a & \hyperref[annex:p02-w02-p02-040-no-effective-calculus]{\ref*{annex:p02-w02-p02-040-no-effective-calculus}} & \thepage & Outer tags 15-18 also occur as derivation, proof-code and alpha-record tags in sec:arithmetic-coding-04; disjoint by arity, shown in the annex. Original precedence descriptor 'none' retained; no rank inferred. Blank macro filled with the glyph's TeX spelling. \\
\phantomsection\label{entry:symbol-registry-07201}\(\mathcal T\) & \texttt{\detokenize{\mathcal T}} & Set of true propositions of a general model & \hyperref[def:p02-w02-p02-040-church-general]{\ref*{def:p02-w02-p02-040-church-general}} & n/a & \hyperref[def:p02-w02-p02-040-church-general]{\ref*{def:p02-w02-p02-040-church-general}} & \thepage & Subset of \(\mathcal D_o\); \(\mathcal D_o\) may have more than two members. Original precedence descriptor 'none' retained; no rank inferred. \\
\phantomsection\label{entry:symbol-registry-07202}\(n\) & \texttt{\detokenize{n}} & Denotations of N\_oo, A\_ooo, Pi and iota in a general model symbol family member \(n\) & \hyperref[def:p02-w02-p02-040-church-general]{\ref*{def:p02-w02-p02-040-church-general}} & n/a & \hyperref[def:p02-w02-p02-040-church-general]{\ref*{def:p02-w02-p02-040-church-general}} & \thepage & \(c_\alpha\) is also used in def:p02-w02-p02-040-church-standard for the choice functions of a standard model. Original precedence descriptor 'none' retained; no rank inferred. Explicit displayed family; member types and common actual first-use exposed, not one atomic glyph. Blank macro filled with the glyph's TeX spelling. Family listing, not one atomic glyph. \\
\phantomsection\label{entry:symbol-registry-07203}\(a\) & \texttt{\detokenize{a}} & Denotations of N\_oo, A\_ooo, Pi and iota in a general model symbol family member \(a\) & \hyperref[def:p02-w02-p02-040-church-general]{\ref*{def:p02-w02-p02-040-church-general}} & n/a & \hyperref[def:p02-w02-p02-040-church-general]{\ref*{def:p02-w02-p02-040-church-general}} & \thepage & \(c_\alpha\) is also used in def:p02-w02-p02-040-church-standard for the choice functions of a standard model. Original precedence descriptor 'none' retained; no rank inferred. Explicit displayed family; member types and common actual first-use exposed, not one atomic glyph. Blank macro filled with the glyph's TeX spelling. Family listing, not one atomic glyph. \\
\phantomsection\label{entry:symbol-registry-07204}\(\pi_\alpha\) & \texttt{\detokenize{\pi_\alpha}} & Denotations of N\_oo, A\_ooo, Pi and iota in a general model symbol family member \(\pi_\alpha\) & \hyperref[def:p02-w02-p02-040-church-general]{\ref*{def:p02-w02-p02-040-church-general}} & n/a & \hyperref[def:p02-w02-p02-040-church-general]{\ref*{def:p02-w02-p02-040-church-general}} & \thepage & \(c_\alpha\) is also used in def:p02-w02-p02-040-church-standard for the choice functions of a standard model. Original precedence descriptor 'none' retained; no rank inferred. Explicit displayed family; member types and common actual first-use exposed, not one atomic glyph. Blank macro filled with the glyph's TeX spelling. Family listing, not one atomic glyph. \\
\phantomsection\label{entry:symbol-registry-07205}\(c_\alpha\) & \texttt{\detokenize{c_\alpha}} & Denotations of N\_oo, A\_ooo, Pi and iota in a general model symbol family member \(c_\alpha\) & \hyperref[def:p02-w02-p02-040-church-general]{\ref*{def:p02-w02-p02-040-church-general}} & n/a & \hyperref[def:p02-w02-p02-040-church-general]{\ref*{def:p02-w02-p02-040-church-general}} & \thepage & \(c_\alpha\) is also used in def:p02-w02-p02-040-church-standard for the choice functions of a standard model. Original precedence descriptor 'none' retained; no rank inferred. Explicit displayed family; member types and common actual first-use exposed, not one atomic glyph. Blank macro filled with the glyph's TeX spelling. Family listing, not one atomic glyph. \\
\phantomsection\label{entry:symbol-registry-07206}\(\entails_{\mathrm G}\) & \texttt{\detokenize{\entails_{\mathrm G}}} & Local consequence over general models validating Axioms 7 and 8 & \hyperref[def:p02-w02-p02-040-church-general]{\ref*{def:p02-w02-p02-040-church-general}} & relation & \hyperref[def:p02-w02-p02-040-church-general]{\ref*{def:p02-w02-p02-040-church-general}} & \thepage & Holds assumption variables fixed in one assignment, matching Church's proof-on-assumptions restriction. \\
\phantomsection\label{entry:symbol-registry-07207}\(\Sigma_{\mathrm C}\) & \texttt{\detokenize{\Sigma_{\mathrm C}}} & Many-sorted signature of the encoding of Church's system & \hyperref[def:p02-w02-p02-040-church-encoding]{\ref*{def:p02-w02-p02-040-church-encoding}} & n/a & \hyperref[def:p02-w02-p02-040-church-encoding]{\ref*{def:p02-w02-p02-040-church-encoding}} & \thepage & Original precedence descriptor 'none' retained; no rank inferred. Blank macro filled with the glyph's TeX spelling. \\
\phantomsection\label{entry:symbol-registry-07208}\(\mathsf s_\gamma\) & \texttt{\detokenize{\mathsf s_\gamma}} & Sort of the encoding for type \(\gamma\) & \hyperref[def:p02-w02-p02-040-church-encoding]{\ref*{def:p02-w02-p02-040-church-encoding}} & n/a & \hyperref[def:p02-w02-p02-040-church-encoding]{\ref*{def:p02-w02-p02-040-church-encoding}} & \thepage & Original precedence descriptor 'none' retained; no rank inferred. Blank macro filled with the glyph's TeX spelling. \\
\phantomsection\label{entry:symbol-registry-07209}\(\mathsf{ap}_{\alpha\beta}\) & \texttt{\detokenize{\mathsf{ap}}} & First-order application function symbol & \hyperref[def:p02-w02-p02-040-church-encoding]{\ref*{def:p02-w02-p02-040-church-encoding}} & n/a & \hyperref[def:p02-w02-p02-040-church-encoding]{\ref*{def:p02-w02-p02-040-church-encoding}} & \thepage & Original precedence descriptor 'none' retained; no rank inferred. \\
\phantomsection\label{entry:symbol-registry-07210}\(\widehat N\) & \texttt{\detokenize{\widehat N}} & Encoding constants for Church's logical constants symbol family member \(\widehat N\) & \hyperref[def:p02-w02-p02-040-church-encoding]{\ref*{def:p02-w02-p02-040-church-encoding}} & n/a & \hyperref[def:p02-w02-p02-040-church-encoding]{\ref*{def:p02-w02-p02-040-church-encoding}} & \thepage & Hat also marks predicate-constant codes \(\widehat P\) in def:p02-sol-incidence-translation; scopes are disjoint. Original precedence descriptor 'none' retained; no rank inferred. Explicit displayed family; member types and common actual first-use exposed, not one atomic glyph. Family listing, not one atomic glyph. Author macro \texttt{\detokenize{\widehat}} not carried to the members. \\
\phantomsection\label{entry:symbol-registry-07211}\(\widehat A\) & \texttt{\detokenize{\widehat A}} & Encoding constants for Church's logical constants symbol family member \(\widehat A\) & \hyperref[def:p02-w02-p02-040-church-encoding]{\ref*{def:p02-w02-p02-040-church-encoding}} & n/a & \hyperref[def:p02-w02-p02-040-church-encoding]{\ref*{def:p02-w02-p02-040-church-encoding}} & \thepage & Hat also marks predicate-constant codes \(\widehat P\) in def:p02-sol-incidence-translation; scopes are disjoint. Original precedence descriptor 'none' retained; no rank inferred. Explicit displayed family; member types and common actual first-use exposed, not one atomic glyph. Family listing, not one atomic glyph. Author macro \texttt{\detokenize{\widehat}} not carried to the members. \\
\phantomsection\label{entry:symbol-registry-07212}\(\widehat\Pi_\alpha\) & \texttt{\detokenize{\widehat\Pi_\alpha}} & Encoding constants for Church's logical constants symbol family member \(\widehat\Pi_\alpha\) & \hyperref[def:p02-w02-p02-040-church-encoding]{\ref*{def:p02-w02-p02-040-church-encoding}} & n/a & \hyperref[def:p02-w02-p02-040-church-encoding]{\ref*{def:p02-w02-p02-040-church-encoding}} & \thepage & Hat also marks predicate-constant codes \(\widehat P\) in def:p02-sol-incidence-translation; scopes are disjoint. Original precedence descriptor 'none' retained; no rank inferred. Explicit displayed family; member types and common actual first-use exposed, not one atomic glyph. Family listing, not one atomic glyph. Author macro \texttt{\detokenize{\widehat}} not carried to the members. \\
\phantomsection\label{entry:symbol-registry-07213}\(\widehat\iota_\alpha\) & \texttt{\detokenize{\widehat\iota_\alpha}} & Encoding constants for Church's logical constants symbol family member \(\widehat\iota_\alpha\) & \hyperref[def:p02-w02-p02-040-church-encoding]{\ref*{def:p02-w02-p02-040-church-encoding}} & n/a & \hyperref[def:p02-w02-p02-040-church-encoding]{\ref*{def:p02-w02-p02-040-church-encoding}} & \thepage & Hat also marks predicate-constant codes \(\widehat P\) in def:p02-sol-incidence-translation; scopes are disjoint. Original precedence descriptor 'none' retained; no rank inferred. Explicit displayed family; member types and common actual first-use exposed, not one atomic glyph. Family listing, not one atomic glyph. Author macro \texttt{\detokenize{\widehat}} not carried to the members. \\
\phantomsection\label{entry:symbol-registry-07214}\(\ell_{\mathbf L}\) & \texttt{\detokenize{\ell}} & Function symbol naming the abstract L in the encoding & \hyperref[def:p02-w02-p02-040-church-encoding]{\ref*{def:p02-w02-p02-040-church-encoding}} & n/a & \hyperref[def:p02-w02-p02-040-church-encoding]{\ref*{def:p02-w02-p02-040-church-encoding}} & \thepage & One symbol per abstract; arguments are the free variables of L in enumeration order. Original precedence descriptor 'none' retained; no rank inferred. \\
\phantomsection\label{entry:symbol-registry-07215}\(\mathsf{True}\) & \texttt{\detokenize{\mathsf{True}}} & Truth predicate on the proposition sort & \hyperref[def:p02-w02-p02-040-church-encoding]{\ref*{def:p02-w02-p02-040-church-encoding}} & n/a & \hyperref[def:p02-w02-p02-040-church-encoding]{\ref*{def:p02-w02-p02-040-church-encoding}} & \thepage & Distinct from the arithmetic formula \(\mathsf{Tr}\) in annex:p02-w02-p02-040-no-effective-calculus. Original precedence descriptor 'none' retained; no rank inferred. \\
\phantomsection\label{entry:symbol-registry-07216}\(\mathbf M^\dagger\) & \texttt{\detokenize{^\dagger}} & First-order encoding of a Church wff & \hyperref[def:p02-w02-p02-040-church-encoding]{\ref*{def:p02-w02-p02-040-church-encoding}} & n/a & \hyperref[def:p02-w02-p02-040-church-encoding]{\ref*{def:p02-w02-p02-040-church-encoding}} & \thepage & Distinct from the proof-local \(\mathbf M^\flat\). Original precedence descriptor 'postfix' retained; no rank inferred. \\
\phantomsection\label{entry:symbol-registry-07217}\(\mathsf{Lb}_\gamma\) & \texttt{\detokenize{\mathsf{Lb}}} & First-order Leibniz formula at sort gamma & \hyperref[def:p02-w02-p02-040-church-encoding]{\ref*{def:p02-w02-p02-040-church-encoding}} & n/a & \hyperref[def:p02-w02-p02-040-church-encoding]{\ref*{def:p02-w02-p02-040-church-encoding}} & \thepage & Original precedence descriptor 'none' retained; no rank inferred. \\
\phantomsection\label{entry:symbol-registry-07218}\(\mathsf{Enc}\) & \texttt{\detokenize{\mathsf{Enc}}} & First-order theory encoding Church's system & \hyperref[def:p02-w02-p02-040-church-encoding]{\ref*{def:p02-w02-p02-040-church-encoding}} & n/a & \hyperref[def:p02-w02-p02-040-church-encoding]{\ref*{def:p02-w02-p02-040-church-encoding}} & \thepage & Original precedence descriptor 'none' retained; no rank inferred. \\
\phantomsection\label{entry:symbol-registry-07219}\(\mathcal M^\dagger\) & \texttt{\detokenize{^\dagger}} & Encoding of a general model as an Enc-model & \hyperref[thm:p02-w02-p02-040-church-encoding]{\ref*{thm:p02-w02-p02-040-church-encoding}} & n/a & \hyperref[thm:p02-w02-p02-040-church-encoding]{\ref*{thm:p02-w02-p02-040-church-encoding}} & \thepage & Original precedence descriptor 'postfix' retained; no rank inferred. \\
\phantomsection\label{entry:symbol-registry-07220}\(\mathfrak A^\natural\) & \texttt{\detokenize{^\natural}} & General model decoded from an Enc-model & \hyperref[thm:p02-w02-p02-040-church-encoding]{\ref*{thm:p02-w02-p02-040-church-encoding}} & n/a & \hyperref[thm:p02-w02-p02-040-church-encoding]{\ref*{thm:p02-w02-p02-040-church-encoding}} & \thepage & Original precedence descriptor 'postfix' retained; no rank inferred. \\
\phantomsection\label{entry:symbol-registry-07221}\(d_\gamma\) & \texttt{\detokenize{d_\gamma}} & Decoding maps of an Enc-model & \hyperref[thm:p02-w02-p02-040-church-encoding]{\ref*{thm:p02-w02-p02-040-church-encoding}} & n/a & \hyperref[thm:p02-w02-p02-040-church-encoding]{\ref*{thm:p02-w02-p02-040-church-encoding}} & \thepage & Blank macro filled with the glyph's TeX spelling. \\
\phantomsection\label{entry:symbol-registry-07222}\(\mathrel{\mathrm{cnv}}\) & \texttt{\detokenize{\mathrel{\mathrm{cnv}}}} & Convertibility of wffs by Rules I-III & \hyperref[annex:p02-w02-p02-040-church-general-completeness]{\ref*{annex:p02-w02-p02-040-church-general-completeness}} & relation & \hyperref[annex:p02-w02-p02-040-church-general-completeness]{\ref*{annex:p02-w02-p02-040-church-general-completeness}} & \thepage & Proof-local; chosen to avoid Church's material equivalence \(\equiv\) between wffs of type o. \\
\phantomsection\label{entry:symbol-registry-07223}\(\derives_{\mathrm{nat}}\) & \texttt{\detokenize{\derives_{\mathrm{nat}}}} & Native derivability on sheet:lk-eq plus the sheet:sol additions & \hyperref[def:p02-w02-p02-040-sol-native]{\ref*{def:p02-w02-p02-040-sol-native}} & relation & \hyperref[def:p02-w02-p02-040-sol-native]{\ref*{def:p02-w02-p02-040-sol-native}} & \thepage &  \\
\phantomsection\label{entry:symbol-registry-07224}\(\operatorname{ext}(U)\) & \texttt{\detokenize{\operatorname{ext}}} & Extension of a relation symbol in the canonical Henkin structure & \hyperref[annex:p02-w02-p02-040-sol-native]{\ref*{annex:p02-w02-p02-040-sol-native}} & n/a & \hyperref[annex:p02-w02-p02-040-sol-native]{\ref*{annex:p02-w02-p02-040-sol-native}} & \thepage & Same name as \(\operatorname{ext}_n\) in annex:p2.hol.8, different argument (symbols, not relation-sort elements). Original precedence descriptor 'none' retained; no rank inferred. \\
\phantomsection\label{entry:symbol-registry-07225}\(\mathbf Z^k\) & \texttt{\detokenize{\mathbf Z^k}} & initial recursion-term constructor: zero term of arity k & \hyperref[def:p02-w02-p02-041-terms]{\ref*{def:p02-w02-p02-041-terms}} & n/a & \hyperref[def:p02-w02-p02-041-terms]{\ref*{def:p02-w02-p02-041-terms}} & \thepage & Bold Z, S, P name syntactic terms; f\_rho gives the function. Bold S is distinct from the function symbol S of the arithmetic signature. Original precedence descriptor 'none' retained; no rank inferred. Blank macro filled with the glyph's TeX spelling. \\
\phantomsection\label{entry:symbol-registry-07226}\(\mathbf S\) & \texttt{\detokenize{\mathbf S}} & initial recursion-term constructor: successor term & \hyperref[def:p02-w02-p02-041-terms]{\ref*{def:p02-w02-p02-041-terms}} & n/a & \hyperref[def:p02-w02-p02-041-terms]{\ref*{def:p02-w02-p02-041-terms}} & \thepage & Bold Z, S, P name syntactic terms; f\_rho gives the function. Bold S is distinct from the function symbol S of the arithmetic signature. Original precedence descriptor 'none' retained; no rank inferred. Blank macro filled with the glyph's TeX spelling. \\
\phantomsection\label{entry:symbol-registry-07227}\(\mathbf P^k_i\) & \texttt{\detokenize{\mathbf P^k_i}} & initial recursion-term constructor: projection term & \hyperref[def:p02-w02-p02-041-terms]{\ref*{def:p02-w02-p02-041-terms}} & n/a & \hyperref[def:p02-w02-p02-041-terms]{\ref*{def:p02-w02-p02-041-terms}} & \thepage & Bold Z, S, P name syntactic terms; f\_rho gives the function. Bold S is distinct from the function symbol S of the arithmetic signature. Original precedence descriptor 'none' retained; no rank inferred. Blank macro filled with the glyph's TeX spelling. \\
\phantomsection\label{entry:symbol-registry-07228}\(\mathbf C^k(\gamma;\eta_1,\ldots,\eta_m)\) & \texttt{\detokenize{\mathbf C^k(\gamma;\eta_1,}\allowbreak\detokenize{\ldots,}\allowbreak\detokenize{\eta_m)}} & recursion-term constructor for composition & \hyperref[def:p02-w02-p02-041-terms]{\ref*{def:p02-w02-p02-041-terms}} & n/a & \hyperref[def:p02-w02-p02-041-terms]{\ref*{def:p02-w02-p02-041-terms}} & \thepage & Bold C, R, M; scoped. Bold B is Church's syntactic variable in def:p02-w02-p02-040-church-wff, not used here. Original precedence descriptor 'none' retained; no rank inferred. Blank macro filled with the glyph's TeX spelling. \\
\phantomsection\label{entry:symbol-registry-07229}\(\mathbf R(\gamma,\eta)\) & \texttt{\detokenize{\mathbf R(\gamma,}\allowbreak\detokenize{\eta)}} & recursion-term constructor for primitive recursion & \hyperref[def:p02-w02-p02-041-terms]{\ref*{def:p02-w02-p02-041-terms}} & n/a & \hyperref[def:p02-w02-p02-041-terms]{\ref*{def:p02-w02-p02-041-terms}} & \thepage & Bold C, R, M; scoped. Bold B is Church's syntactic variable in def:p02-w02-p02-040-church-wff, not used here. Original precedence descriptor 'none' retained; no rank inferred. Blank macro filled with the glyph's TeX spelling. \\
\phantomsection\label{entry:symbol-registry-07230}\(\mathbf M(\gamma)\) & \texttt{\detokenize{\mathbf M(\gamma)}} & recursion-term constructor for unbounded minimization & \hyperref[def:p02-w02-p02-041-terms]{\ref*{def:p02-w02-p02-041-terms}} & n/a & \hyperref[def:p02-w02-p02-041-terms]{\ref*{def:p02-w02-p02-041-terms}} & \thepage & Bold C, R, M; scoped. Bold B is Church's syntactic variable in def:p02-w02-p02-040-church-wff, not used here. Original precedence descriptor 'none' retained; no rank inferred. Blank macro filled with the glyph's TeX spelling. \\
\phantomsection\label{entry:symbol-registry-07231}\(N_{30},\ldots,N_{35}\) & \texttt{\detokenize{N_{k}}} & constructor tags of recursion-term codes & \hyperref[def:p02-w02-p02-041-terms]{\ref*{def:p02-w02-p02-041-terms}} & n/a & \hyperref[def:p02-w02-p02-041-terms]{\ref*{def:p02-w02-p02-041-terms}} & \thepage & Tags 30-35 are unused by sec:arithmetic-coding-04 (0-18) and by def:p02-w02-p02-040-so-coding (15-18). Original precedence descriptor 'none' retained; no rank inferred. Sequence glyph kept as one row; the ellipsis is not a member symbol. \\
\phantomsection\label{entry:symbol-registry-07232}\(\mathrm{ar}(\rho)\) & \texttt{\detokenize{\mathrm{ar}}} & arity of a recursion term; also the primitive recursive arity function on codes & \hyperref[def:p02-w02-p02-041-terms]{\ref*{def:p02-w02-p02-041-terms}} & n/a & \hyperref[def:p02-w02-p02-041-terms]{\ref*{def:p02-w02-p02-041-terms}} & \thepage & Original precedence descriptor 'function application' retained; no rank inferred. \\
\phantomsection\label{entry:symbol-registry-07233}\(f_\rho\) & \texttt{\detokenize{f_{\rho}}} & partial map denoted by a recursion term & \hyperref[def:p02-w02-p02-041-terms]{\ref*{def:p02-w02-p02-041-terms}} & n/a & \hyperref[def:p02-w02-p02-041-terms]{\ref*{def:p02-w02-p02-041-terms}} & \thepage & f alone is an isomorphism in thm:p02-w02-p02-041-iso; the subscript distinguishes. Original precedence descriptor 'none' retained; no rank inferred. \\
\phantomsection\label{entry:symbol-registry-07234}\(\mathrm{Gr}_\rho\) & \texttt{\detokenize{\mathrm{Gr}}} & graph of the partial map denoted by rho & \hyperref[def:p02-w02-p02-041-terms]{\ref*{def:p02-w02-p02-041-terms}} & n/a & \hyperref[def:p02-w02-p02-041-terms]{\ref*{def:p02-w02-p02-041-terms}} & \thepage & Original precedence descriptor 'none' retained; no rank inferred. \\
\phantomsection\label{entry:symbol-registry-07235}\(X_\sigma\) & \texttt{\detokenize{X_{\sigma}}} & the relation variable X\^{}{ar(sigma)+1} indexed by the code of sigma & \hyperref[def:p02-w02-p02-041-graph-sentence]{\ref*{def:p02-w02-p02-041-graph-sentence}} & n/a & \hyperref[def:p02-w02-p02-041-graph-sentence]{\ref*{def:p02-w02-p02-041-graph-sentence}} & \thepage & Original precedence descriptor 'none' retained; no rank inferred. \\
\phantomsection\label{entry:symbol-registry-07236}\(\mathsf{Def}_\sigma\) & \texttt{\detokenize{\mathsf{Def}}} & second-order graph condition of a recursion term & \hyperref[def:p02-w02-p02-041-graph-sentence]{\ref*{def:p02-w02-p02-041-graph-sentence}} & n/a & \hyperref[def:p02-w02-p02-041-graph-sentence]{\ref*{def:p02-w02-p02-041-graph-sentence}} & \thepage & Original precedence descriptor 'none' retained; no rank inferred. \\
\phantomsection\label{entry:symbol-registry-07237}\(\mathrm{Sub}(\rho)\) & \texttt{\detokenize{\mathrm{Sub}}} & subterm list of a recursion term (repetitions kept) & \hyperref[def:p02-w02-p02-041-graph-sentence]{\ref*{def:p02-w02-p02-041-graph-sentence}} & n/a & \hyperref[def:p02-w02-p02-041-graph-sentence]{\ref*{def:p02-w02-p02-041-graph-sentence}} & \thepage & Code-level counterpart Sb in annex:p02-w02-p02-041-effective. Original precedence descriptor 'function application' retained; no rank inferred. \\
\phantomsection\label{entry:symbol-registry-07238}\(\Theta_\rho\) & \texttt{\detokenize{\Theta}} & divergence sentence of a unary recursion term & \hyperref[def:p02-w02-p02-041-graph-sentence]{\ref*{def:p02-w02-p02-041-graph-sentence}} & n/a & \hyperref[def:p02-w02-p02-041-graph-sentence]{\ref*{def:p02-w02-p02-041-graph-sentence}} & \thepage & Original precedence descriptor 'none' retained; no rank inferred. \\
\phantomsection\label{entry:symbol-registry-07239}\(\mathrm{rd}\) & \texttt{\detokenize{\mathrm{rd}}} & primitive recursive map from term codes to codes of divergence sentences & \hyperref[thm:p02-w02-p02-041-effective]{\ref*{thm:p02-w02-p02-041-effective}} & n/a & \hyperref[thm:p02-w02-p02-041-effective]{\ref*{thm:p02-w02-p02-041-effective}} & \thepage & Named rd to avoid r of thm:p02-w02-p02-040-no-effective-calculus. Original precedence descriptor 'function application' retained; no rank inferred. \\
\phantomsection\label{entry:symbol-registry-07240}\(\mathrm{RT}\) & \texttt{\detokenize{\mathrm{RT}}} & primitive recursive test for recursion-term codes & \hyperref[thm:p02-w02-p02-041-effective]{\ref*{thm:p02-w02-p02-041-effective}} & n/a & \hyperref[thm:p02-w02-p02-041-effective]{\ref*{thm:p02-w02-p02-041-effective}} & \thepage & Original precedence descriptor 'function application' retained; no rank inferred. \\
\phantomsection\label{entry:symbol-registry-07241}\(\operatorname{ext}_n(r)\) & \texttt{\detokenize{\operatorname{ext}_n}} & set of n-tuples incident with a relation-sort object & \hyperref[def:p02-w02-p02-041-collapse]{\ref*{def:p02-w02-p02-041-collapse}} & n/a & \hyperref[def:p02-w02-p02-041-collapse]{\ref*{def:p02-w02-p02-041-collapse}} & \thepage & Same symbol and meaning as in the proof of thm:p02-sol-derived-adequacy (annex:p2.hol.8). Original precedence descriptor 'function application' retained; no rank inferred. \\
\phantomsection\label{entry:symbol-registry-07242}\(\mathcal H_{\mathrm{ext}}\) & \texttt{\detokenize{\mathcal H_{\mathrm{ext}}}} & extensional collapse of a strict incidence structure & \hyperref[def:p02-w02-p02-041-collapse]{\ref*{def:p02-w02-p02-041-collapse}} & n/a & \hyperref[def:p02-w02-p02-041-collapse]{\ref*{def:p02-w02-p02-041-collapse}} & \thepage & Original precedence descriptor 'none' retained; no rank inferred. \\
\phantomsection\label{entry:symbol-registry-07243}\(s^{\mathrm{ext}}\) & \texttt{\detokenize{s^{\mathrm{ext}}}} & decoded assignment of the extensional collapse & \hyperref[def:p02-w02-p02-041-collapse]{\ref*{def:p02-w02-p02-041-collapse}} & n/a & \hyperref[def:p02-w02-p02-041-collapse]{\ref*{def:p02-w02-p02-041-collapse}} & \thepage & Counterpart of the tagged assignment s\^{}sharp of def:p02-sol-incidence-translation. Original precedence descriptor 'none' retained; no rank inferred. \\
\phantomsection\label{entry:symbol-registry-07244}\(\bar k^{\beta}\) & \texttt{\detokenize{\bar k^{\beta}}} & iteration functional phi to phi\^{}k in the standard model of Church's system & \hyperref[annex:p02-w02-p02-041-church-30]{\ref*{annex:p02-w02-p02-041-church-30}} & n/a & \hyperref[annex:p02-w02-p02-041-church-30]{\ref*{annex:p02-w02-p02-041-church-30}} & \thepage & Type superscript added to the bar-n notation of annex:p02-w02-p02-040-church-soundness to distinguish levels o and o'. Original precedence descriptor 'none' retained; no rank inferred. \\
\end{longtable}
\endgroup
\Needspace{.4\textheight}
\subsection*{Stable part P3 scopes}
\begingroup\scriptsize\setlength{\tabcolsep}{3pt}
\begin{longtable}{@{}p{0.9in}p{1.3in}p{1.0in}p{0.5in}p{0.3in}p{0.4in}p{0.35in}p{1.65in}@{}}
\toprule
\textbf{Glyph} & \textbf{Macro} & \textbf{Role} & \textbf{Scope} & \textbf{Prec./form} & \textbf{First} & \textbf{Reg. page} & \textbf{Principal uses / notes} \\
\midrule
\endfirsthead
\toprule
\textbf{Glyph} & \textbf{Macro} & \textbf{Role} & \textbf{Scope} & \textbf{Prec./form} & \textbf{First} & \textbf{Reg. page} & \textbf{Principal uses / notes} \\
\midrule
\endhead
\midrule\multicolumn{8}{r}{\emph{continued}}\\\endfoot
\bottomrule\endlastfoot
\phantomsection\label{entry:symbol-registry-00044}\(\operatorname{Diag}(\mathcal M)\) & \texttt{\detokenize{\operatorname{Diag}}} & atomic diagram & P3 maps-elementarity to further-reading & n/a & \hyperref[sec:maps-elementarity-01]{\ref*{sec:maps-elementarity-01}} & \thepage & Atomic sentences with parameters true in \(\mathcal M\). \\
\phantomsection\label{entry:symbol-registry-00046}\(S_n(T)\) & \texttt{\detokenize{S_n(T)}} & type space & P3 maps-elementarity to further-reading & n/a & \hyperref[sec:types-ultraproducts-01]{\ref*{sec:types-ultraproducts-01}} & \thepage & Complete \(n\)-types consistent with \(T\) under the clopen topology. \\
\phantomsection\label{entry:symbol-registry-00047}\(\prod_{i\in I}\mathcal M_i/\mathcal U\) & \texttt{\detokenize{\prod_{i\in I}\mathcal M_i/}\allowbreak\detokenize{\mathcal U}} & ultraproduct & Legacy pre-batch5 P3.3 reader meaning & n/a & \hyperref[sec:qe-categoricity-01]{\ref*{sec:qe-categoricity-01}} & \thepage & Quotient of the direct product by equality on a set in the ultrafilter \(\mathcal U\).  Numeric chapter scope migrated to its frozen pre-batch5 source-labelled reader regions; glyph, role and precedence unchanged. \\
\phantomsection\label{entry:symbol-registry-00445}\(\PA\) & \texttt{\detokenize{\PA}} & Peano arithmetic & Legacy pre-batch4 P3-P7.5 reader meaning & n/a & \hyperref[sec:qe-categoricity-01]{\ref*{sec:qe-categoricity-01}} & \thepage & The fixed first-order theory defined in P5.1; earlier table and later metatheoretic uses denote that same theory.  Numeric chapter scope migrated to its frozen pre-batch4 source-labelled reader regions; glyph, role and precedence unchanged. \\
\phantomsection\label{entry:symbol-registry-02073}\(R:E^n\to\{+,-\}\) & \texttt{\detokenize{R}} & two-valued relation on a stated base & P3 fraisse-1954 relation presentation & n/a & \hyperref[def:p03-fraisse54-relation]{\ref*{def:p03-fraisse54-relation}} & \thepage & Positive finite arity; tuples may repeat entries; empty bases explicitly allowed. \\
\phantomsection\label{entry:symbol-registry-02074}\(+,-\) & \texttt{\detokenize{+,}\allowbreak\detokenize{-}} & positive and negative relation values & P3 fraisse-1954 relation presentation & n/a & \hyperref[def:p03-fraisse54-relation]{\ref*{def:p03-fraisse54-relation}} & \thepage & Signs are relation values rather than arithmetic operations. \\
\phantomsection\label{entry:symbol-registry-02075}\(A\ll B\) & \texttt{\detokenize{\ll}} & existence of an induced relational embedding & P3 fraisse-1954 relation presentation & relation & \hyperref[def:p03-fraisse54-age]{\ref*{def:p03-fraisse54-age}} & \thepage & Preserves and reflects every relation instance; not a cardinal comparison. \\
\phantomsection\label{entry:symbol-registry-02076}\(A\prec B\) & \texttt{\detokenize{\prec}} & embedding of every finite restriction & P3 fraisse-1954 relation presentation & relation & \hyperref[def:p03-fraisse54-age]{\ref*{def:p03-fraisse54-age}} & \thepage & Finite embeddability is distinct from one embedding of the whole structure. \\
\phantomsection\label{entry:symbol-registry-02077}\(\gamma_R\) & \texttt{\detokenize{\gamma_R}} & age of a relation & P3 fraisse-1954 relation presentation & n/a & \hyperref[def:p03-fraisse54-age]{\ref*{def:p03-fraisse54-age}} & \thepage & Isomorphism-closed class of finite induced restrictions. \\
\phantomsection\label{entry:symbol-registry-02078}\(\Gamma_R\) & \texttt{\detokenize{\Gamma_R}} & local class determined by an age & P3 fraisse-1954 relation presentation & n/a & \hyperref[def:p03-fraisse54-age]{\ref*{def:p03-fraisse54-age}} & \thepage & Contains all set-based relations whose finite restrictions belong to the specified age. \\
\phantomsection\label{entry:symbol-registry-02079}\(\mathrm C_1,\mathrm C_2,\mathrm C_3\) & \texttt{\detokenize{\mathrm C_1;\mathrm C_2;\mathrm C_3}} & restriction joint-embedding and chain-closure conditions & P3 fraisse-1954 relation presentation & n/a & \hyperref[def:p03-fraisse54-closures]{\ref*{def:p03-fraisse54-closures}} & \thepage & These are separately defined class conditions. \\
\phantomsection\label{entry:symbol-registry-02080}\(\mathcal F\) & \texttt{\detokenize{\mathcal F}} & proper filter on a set of maps & P3 fraisse-1954 projection filters & n/a & \hyperref[def:p03-fraisse54-projection]{\ref*{def:p03-fraisse54-projection}} & \thepage & Its underlying set is D raised to E with D nonempty. \\
\phantomsection\label{entry:symbol-registry-02081}\(\mathcal F^{-1}(R)\) & \texttt{\detokenize{\mathcal F^{-1}(R)}} & inverse relation determined by a projection filter & P3 fraisse-1954 projection filters & n/a & \hyperref[def:p03-fraisse54-projection]{\ref*{def:p03-fraisse54-projection}} & \thepage & Defined by membership of all tuple sign-sets; not a function inverse. \\
\phantomsection\label{entry:symbol-registry-02082}\(\mathcal F_H\) & \texttt{\detokenize{\mathcal F_H}} & restriction of a projection filter & P3 fraisse-1954 projection filters & n/a & \hyperref[def:p03-fraisse54-projection]{\ref*{def:p03-fraisse54-projection}} & \thepage & Direct-image filter under restriction of maps to H. \\
\phantomsection\label{entry:symbol-registry-02083}\(I_H\) & \texttt{\detokenize{I_H}} & set of maps injective on a finite subset & P3 fraisse-1954 projection filters & n/a & \hyperref[def:p03-fraisse54-projection]{\ref*{def:p03-fraisse54-projection}} & \thepage & Membership for every finite H defines the source normality condition. \\
\phantomsection\label{entry:symbol-registry-02084}\(a\sim_{\mathcal F}b\) & \texttt{\detokenize{\sim_{\mathcal F}}} & equality of images on a filter-large set of maps & P3 fraisse-1954 projection filters & relation & \hyperref[def:p03-fraisse54-filter-equality]{\ref*{def:p03-fraisse54-filter-equality}} & \thepage & An equivalence on the domain; normal filters identify no distinct elements. \\
\phantomsection\label{entry:symbol-registry-02085}\(\mathcal H\) & \texttt{\detokenize{\mathcal H}} & finite extension property of a relation & P3 fraisse-1954 homogeneity & n/a & \hyperref[def:p03-fraisse54-extension]{\ref*{def:p03-fraisse54-extension}} & \thepage & Embeddings of finite age members extend over every finite extension in the age. \\
\phantomsection\label{entry:symbol-registry-02086}\(\mathcal C\) & \texttt{\detokenize{\mathcal C}} & countable universality condition & P3 fraisse-1954 homogeneity & n/a & \hyperref[def:p03-fraisse54-universal]{\ref*{def:p03-fraisse54-universal}} & \thepage & One finite or countable member contains induced copies of all finite or countable members of the local class. \\
\phantomsection\label{entry:symbol-registry-02087}\(\mathcal D\) & \texttt{\detokenize{\mathcal D}} & ordinary amalgamation condition & P3 fraisse-1954 amalgamation & n/a & \hyperref[def:p03-fraisse54-amalgamation]{\ref*{def:p03-fraisse54-amalgamation}} & \thepage & Allows identifications of new points on opposite sides. \\
\phantomsection\label{entry:symbol-registry-02088}\(\mathcal D'\) & \texttt{\detokenize{\mathcal D'}} & extension-over-a-fixed-host form of amalgamation & P3 fraisse-1954 amalgamation & n/a & \hyperref[def:p03-fraisse54-amalgamation]{\ref*{def:p03-fraisse54-amalgamation}} & \thepage & Equivalent to ordinary amalgamation after isomorphic transport. \\
\phantomsection\label{entry:symbol-registry-02089}\(\overline{\mathcal D}\) & \texttt{\detokenize{\overline{\mathcal D}}} & literal common-extension condition & P3 fraisse-1954 amalgamation & n/a & \hyperref[def:p03-fraisse54-amalgamation]{\ref*{def:p03-fraisse54-amalgamation}} & \thepage & Retains distinct base elements; the corresponding strong amalgamation condition. \\
\phantomsection\label{entry:symbol-registry-02090}\(R_p\) & \texttt{\detokenize{R_p}} & finite stage in the fixed-arity construction & P3 fraisse-1954 original construction & n/a & \hyperref[annex:p3.fraisse54.original-construction]{\ref*{annex:p3.fraisse54.original-construction}} & \thepage & Stage p schedules all finite types of size at most p and the specified finite embeddings. \\
\phantomsection\label{entry:symbol-registry-02091}\(L_Q,L_R\) & \texttt{\detokenize{L_Q;L_R}} & specified uncountable order sums & P3 fraisse-1954 order examples & n/a & \hyperref[annex:p3.fraisse54.orders]{\ref*{annex:p3.fraisse54.orders}} & \thepage & Respectively omega1 rational blocks and a real prefix followed by omega1 blocks of one plus the reals. \\
\phantomsection\label{entry:symbol-registry-02092}\(B_i,M_s\) & \texttt{\detokenize{B_i;M_s}} & enumerated finite type and finite construction stage & P3 fraisse-1954 countable-signature bridge & n/a & \hyperref[annex:p3.fraisse54.countable-language]{\ref*{annex:p3.fraisse54.countable-language}} & \thepage & Stage s handles all listed types i at most s and all embeddings into the current finite stage. \\
\phantomsection\label{entry:symbol-registry-02093}\(P_i\) & \texttt{\detokenize{P_i}} & unary color predicate & P3 fraisse-1954 categoricity counterexample & n/a & \hyperref[annex:p3.fraisse54.categoricity-boundary]{\ref*{annex:p3.fraisse54.categoricity-boundary}} & \thepage & Countably many disjoint infinite color classes; an added uncolored element gives the contrasting model. \\
\phantomsection\label{entry:symbol-registry-02094}\(E\) & \texttt{\detokenize{E}} & irreflexive symmetric graph edge relation & P3 fraisse-1954 graph interface & n/a & \hyperref[def:p03-fraisse54-graph-axioms]{\ref*{def:p03-fraisse54-graph-axioms}} & \thepage & This local E is a binary relation rather than the base set used earlier. \\
\phantomsection\label{entry:symbol-registry-02095}\(\delta_m(\bar w)\) & \texttt{\detokenize{\delta_m}} & pairwise-distinctness formula for an m-tuple & P3 fraisse-1954 graph interface & n/a & \hyperref[def:p03-fraisse54-graph-axioms]{\ref*{def:p03-fraisse54-graph-axioms}} & \thepage & Empty conjunction is logical truth. \\
\phantomsection\label{entry:symbol-registry-02514}\(\Gamma=(A,B)\) & \texttt{\detokenize{\Gamma;A;B}} & simple undirected graph with stated vertex and edge sets & P3 Rado 1964 universal graphs and functions & n/a & \hyperref[def:p03-rado64-graph]{\ref*{def:p03-rado64-graph}} & \thepage & Edges are two-element subsets; a weak embedding preserves edges and an induced embedding also reflects them. \\
\phantomsection\label{entry:symbol-registry-02515}\([A]^r\) & \texttt{\detokenize{[A]^r}} & set of r-element subsets of the stated base & P3 Rado 1964 universal graphs and functions & n/a & \hyperref[def:p03-rado64-graph]{\ref*{def:p03-rado64-graph}} & \thepage & For a nonnegative integer r; uniform-complex simplices have r equal to the dimension plus one. \\
\phantomsection\label{entry:symbol-registry-02516}\(\varphi\) & \texttt{\detokenize{\varphi}} & injective representation map on vertices or arguments & P3 Rado 1964 universal graphs and functions & n/a & \hyperref[def:p03-rado64-graph]{\ref*{def:p03-rado64-graph}} & \thepage & The constructed maps are increasing; colours remain fixed. This local symbol denotes a function rather than a formula. \\
\phantomsection\label{entry:symbol-registry-02517}\(E_\Gamma\) & \texttt{\detokenize{E_\Gamma}} & edge relation associated with a stated simple graph & P3 Rado 1964 universal graphs and functions & relation & \hyperref[def:p03-rado64-graph]{\ref*{def:p03-rado64-graph}} & \thepage & Equivalent to membership of the unordered pair in the edge set. \\
\phantomsection\label{entry:symbol-registry-02518}\(d_k(t,j)\) & \texttt{\detokenize{d_k}} & jth base-k digit of a nonnegative integer & P3 Rado 1964 universal graphs and functions & n/a & \hyperref[def:p03-rado64-digits]{\ref*{def:p03-rado64-digits}} & \thepage & The base is an integer at least two; the digit lies from zero through k minus one. \\
\phantomsection\label{entry:symbol-registry-02519}\(R\) & \texttt{\detokenize{R}} & Rado's explicit binary-digit graph & P3 Rado 1964 universal graphs and functions & n/a & \hyperref[def:p03-rado64-digits]{\ref*{def:p03-rado64-digits}} & \thepage & Vertex set is the nonnegative integers; this graph object is distinct from a frame accessibility relation. \\
\phantomsection\label{entry:symbol-registry-02520}\(E_R\) & \texttt{\detokenize{E_R}} & symmetric irreflexive edge relation of the binary graph & P3 Rado 1964 universal graphs and functions & relation & \hyperref[def:p03-rado64-digits]{\ref*{def:p03-rado64-digits}} & \thepage & For x less than y its value is given by digit x of y; extend symmetrically. \\
\phantomsection\label{entry:symbol-registry-02521}\(\psi(\xi)\) & \texttt{\detokenize{\psi}} & source fresh high digit position at an integer construction stage & P3 Rado 1964 universal graphs and functions & n/a & \hyperref[thm:p03-rado64-graph-universal]{\ref*{thm:p03-rado64-graph-universal}} & \thepage & Chosen above preceding images and above the required coding positions in the later constructions. \\
\phantomsection\label{entry:symbol-registry-02522}\(k(\Gamma,x,m)\) & \texttt{\detokenize{k(\Gamma,}\allowbreak\detokenize{x,}\allowbreak\detokenize{m)}} & number of vertices reachable by a walk of exactly m edges & P3 Rado 1964 universal graphs and functions & n/a & \hyperref[def:p03-rado64-local]{\ref*{def:p03-rado64-local}} & \thepage & Walk vertices may repeat; local finiteness makes the count finite. \\
\phantomsection\label{entry:symbol-registry-02523}\(\tau\) & \texttt{\detokenize{\tau}} & source binomial code of an increasing tuple at fixed length & P3 Rado 1964 universal graphs and functions & n/a & \hyperref[def:p03-rado64-combinadic]{\ref*{def:p03-rado64-combinadic}} & \thepage & For each positive fixed length this is a colexicographic bijection with the nonnegative integers; no injectivity across lengths is asserted. \\
\phantomsection\label{entry:symbol-registry-02524}\(C=(A,B)\) & \texttt{\detokenize{C}} & uniform complex in the original incidence convention & P3 Rado 1964 universal graphs and functions & n/a & \hyperref[def:p03-rado64-complex]{\ref*{def:p03-rado64-complex}} & \thepage & No closure under taking faces is imposed; embeddings preserve and reflect all simplex membership. \\
\phantomsection\label{entry:symbol-registry-02525}\(C_\ell^*\) & \texttt{\detokenize{C_\ell^*}} & universal fixed-dimension binary-digit complex & P3 Rado 1964 universal graphs and functions & n/a & \hyperref[thm:p03-rado64-uniform-complex]{\ref*{thm:p03-rado64-uniform-complex}} & \thepage & The separate zero-dimensional case uses singleton parity. \\
\phantomsection\label{entry:symbol-registry-02526}\(\rho\) & \texttt{\detokenize{\rho}} & source binary code of an increasing tuple of arbitrary finite length & P3 Rado 1964 universal graphs and functions & n/a & \hyperref[def:p03-rado64-mixed-code]{\ref*{def:p03-rado64-mixed-code}} & \thepage & Equivalently the binary code of the finite set under its increasing enumeration; the empty tuple has code zero. \\
\phantomsection\label{entry:symbol-registry-02527}\(B_{\rm mix}^*\) & \texttt{\detokenize{B_{\rm mix}^*}} & simplex family of the universal mixed complex & P3 Rado 1964 universal graphs and functions & n/a & \hyperref[thm:p03-rado64-mixed-complex]{\ref*{thm:p03-rado64-mixed-complex}} & \thepage & Allows all nonempty finite sizes with singleton membership at digit zero; the empty simplex is excluded. \\
\phantomsection\label{entry:symbol-registry-02528}\(K_{k\ell},K_k\) & \texttt{\detokenize{K_{k\ell};K_k}} & fixed-length and all-finite-length colour-function classes & P3 Rado 1964 universal graphs and functions & n/a & \hyperref[def:p03-rado64-colour-functions]{\ref*{def:p03-rado64-colour-functions}} & \thepage & Arguments are strictly increasing finite tuples; this is not the later repeated-argument domain. \\
\phantomsection\label{entry:symbol-registry-02529}\(f^*\) & \texttt{\detokenize{f^*}} & universal colour function in its stated class & P3 Rado 1964 universal graphs and functions & n/a & \hyperref[def:p03-rado64-colour-functions]{\ref*{def:p03-rado64-colour-functions}} & \thepage & One injection represents a whole function; values in the colour set remain unchanged. \\
\phantomsection\label{entry:symbol-registry-02530}\(f_{k\ell}^*,f_k^*\) & \texttt{\detokenize{f_{k\ell}^*;f_k^*}} & universal fixed-length and mixed-length digit functions & P3 Rado 1964 universal graphs and functions & n/a & \hyperref[thm:p03-rado64-finite-colours]{\ref*{thm:p03-rado64-finite-colours}} & \thepage & The displayed formulas are the functions of the two parts of original Theorem 5. \\
\phantomsection\label{entry:symbol-registry-02531}\(f_0(A,t)\) & \texttt{\detokenize{f_0}} & source unrestricted fixed-length digit function & P3 Rado 1964 universal graphs and functions & n/a & \hyperref[annex:p03.rado64.finite-colours]{\ref*{annex:p03.rado64.finite-colours}} & \thepage & Before the universal function restricts all prefix entries to be below the final argument; the finite set is identified with its increasing enumeration. \\
\phantomsection\label{entry:symbol-registry-02532}\(\omega_n\) & \texttt{\detokenize{\omega_n}} & initial ordinal of cardinality aleph-n & P3 Rado 1964 universal graphs and functions & n/a & \hyperref[def:p03-rado64-transfinite-indices]{\ref*{def:p03-rado64-transfinite-indices}} & \thepage & Here n is an arbitrary ordinal index; this is not the proof-theoretic finite tower notation. \\
\phantomsection\label{entry:symbol-registry-02533}\(\mathrm{cf}(n)\) & \texttt{\detokenize{\mathrm{cf}}} & aleph index of the ordinary cofinality of aleph-n & P3 Rado 1964 universal graphs and functions & n/a & \hyperref[def:p03-rado64-transfinite-indices]{\ref*{def:p03-rado64-transfinite-indices}} & \thepage & Retains the original index-valued notation; ordinary cardinal cofinality of aleph-n is aleph indexed by this value. \\
\phantomsection\label{entry:symbol-registry-02534}\(n-1\) & \texttt{\detokenize{n-1}} & source predecessor convention for ordinal indices & P3 Rado 1964 universal graphs and functions & n/a & \hyperref[def:p03-rado64-transfinite-indices]{\ref*{def:p03-rado64-transfinite-indices}} & \thepage & Equals r when n equals r plus one; equals n at zero or a limit; not ordinal subtraction. \\
\phantomsection\label{entry:symbol-registry-02535}\(F(l,m,S)\) & \texttt{\detokenize{F}} & colour functions on successor-length sequences with a final maximum & P3 Rado 1964 universal graphs and functions & n/a & \hyperref[def:p03-rado64-transfinite-domain]{\ref*{def:p03-rado64-transfinite-domain}} & \thepage & The prefix length is below l and all prefix coordinates are at most the final coordinate below m; repetition of the maximum is permitted. \\
\phantomsection\label{entry:symbol-registry-02536}\(D_\eta\) & \texttt{\detokenize{D_\eta}} & predecessor set with one fresh maximum marker & P3 Rado 1964 universal graphs and functions & n/a & \hyperref[def:p03-rado64-marked-template]{\ref*{def:p03-rado64-marked-template}} & \thepage & The marker is adjoined disjointly even when the ordered domain contains another object written as a star. \\
\phantomsection\label{entry:symbol-registry-02537}\(*\) & \texttt{\detokenize{*}} & marked occurrence of a newly embedded maximum & P3 Rado 1964 universal graphs and functions & n/a & \hyperref[def:p03-rado64-marked-template]{\ref*{def:p03-rado64-marked-template}} & \thepage & Earlier images lie below the selected bound; all repeated occurrences of the new maximum use this marker. \\
\phantomsection\label{entry:symbol-registry-02538}\(h\) & \texttt{\detokenize{h}} & marked-prefix template with stated arity tags & P3 Rado 1964 marked templates & n/a & \hyperref[def:p03-rado64-marked-template]{\ref*{def:p03-rado64-marked-template}} & \thepage & Its domain is the disjoint union of all prefix powers below l; the empty prefix is included when l is positive. \\
\phantomsection\label{entry:symbol-registry-02539}\(h\) & \texttt{\detokenize{h}} & proposed universal host function defeated by diagonalization & P3 Rado 1964 nonuniversality proofs & n/a & \hyperref[annex:p03.rado64.order-type-boundary]{\ref*{annex:p03.rado64.order-type-boundary}} & \thepage & Local to the omega-plus-one and all-finite-repeated-arity counterexamples; distinct from a marked-prefix template. \\
\phantomsection\label{entry:symbol-registry-02540}\(\eta_z,h_z,u_i\) & \texttt{\detokenize{\eta_z;h_z;u_i}} & bound assigned template and encoded prefix coordinate & P3 Rado 1964 universal graphs and functions & n/a & \hyperref[def:p03-rado64-marked-template]{\ref*{def:p03-rado64-marked-template}} & \thepage & Defines the universal value for a tuple ending at z; replace each prefix coordinate at least eta-z by the marker. \\
\phantomsection\label{entry:symbol-registry-02541}\(\kappa\) & \texttt{\detokenize{\kappa}} & target cardinal aleph-n in the transfinite proof & P3 Rado 1964 universal graphs and functions & n/a & \hyperref[annex:p03.rado64.transfinite]{\ref*{annex:p03.rado64.transfinite}} & \thepage & Cardinals are identified with their initial ordinals in the proof; additional GCH assumptions remain those of Theorem 6. \\
\phantomsection\label{entry:symbol-registry-02542}\(P_\nu\) & \texttt{\detokenize{P_\nu}} & unbounded cell realizing one regular-case template & P3 Rado 1964 universal graphs and functions & n/a & \hyperref[annex:p03.rado64.transfinite]{\ref*{annex:p03.rado64.transfinite}} & \thepage & The cells partition the target into target-cardinality many cells each of target cardinality. \\
\phantomsection\label{entry:symbol-registry-02543}\(\kappa_m,X_m\) & \texttt{\detokenize{\kappa_m;X_m}} & successor block cardinal and corresponding ordered block & P3 Rado 1964 universal graphs and functions & n/a & \hyperref[annex:p03.rado64.transfinite]{\ref*{annex:p03.rado64.transfinite}} & \thepage & Kappa-m is aleph indexed by m plus two; X-m has the corresponding initial-ordinal order type. \\
\phantomsection\label{entry:symbol-registry-02544}\(X\) & \texttt{\detokenize{X}} & ordered sum of the singular-case blocks & P3 Rado 1964 universal graphs and functions & n/a & \hyperref[annex:p03.rado64.transfinite]{\ref*{annex:p03.rado64.transfinite}} & \thepage & The proof establishes its exact order type omega-n before transporting the function. \\
\phantomsection\label{entry:symbol-registry-02545}\(P_{m,\mu}\) & \texttt{\detokenize{P_{m,\mu}}} & unbounded cell for a marked template within its singular-case block & P3 Rado 1964 universal graphs and functions & n/a & \hyperref[annex:p03.rado64.transfinite]{\ref*{annex:p03.rado64.transfinite}} & \thepage & The representation preserves blocks and schedules all prefix requirements at the final maximum. \\
\phantomsection\label{entry:symbol-registry-03079}\(v_i\) & \texttt{\detokenize{v_i}} & indexed individual variable & P3 Vaught prime atomic & n/a & \hyperref[def:p03-vaught-standing]{\ref*{def:p03-vaught-standing}} & \thepage & The original countable variable alphabet; substitutions are capture avoiding. \\
\phantomsection\label{entry:symbol-registry-03080}\(F_n(T)\) & \texttt{\detokenize{F_n(T)}} & formulas with free variables among the first n slots & P3 Vaught prime atomic & n/a & \hyperref[def:p03-vaught-formula-algebra]{\ref*{def:p03-vaught-formula-algebra}} & \thepage & Also the source formula prealgebra with its equivalence relation. \\
\phantomsection\label{entry:symbol-registry-03081}\(\vdash_T\varphi\) & \texttt{\detokenize{\vdash_T}} & source validity of a formula's universal closure & P3 Vaught prime atomic & meta & \hyperref[def:p03-vaught-formula-algebra]{\ref*{def:p03-vaught-formula-algebra}} & \thepage & Explicit semantic standing in this source presentation rather than an undeclared proof-system turnstile. \\
\phantomsection\label{entry:symbol-registry-03082}\(\approx_T\) & \texttt{\detokenize{\approx_T}} & equivalence of formulas in every model of T & P3 Vaught prime atomic & relation & \hyperref[def:p03-vaught-formula-algebra]{\ref*{def:p03-vaught-formula-algebra}} & \thepage & Preserves the source's equivalence quotient; elementary equivalence of structures remains a different relation. \\
\phantomsection\label{entry:symbol-registry-03083}\(\mathbb B_n(T)\) & \texttt{\detokenize{\mathbb B_n(T)}} & Boolean quotient of the finite-variable formula prealgebra & P3 Vaught prime atomic & n/a & \hyperref[def:p03-vaught-formula-algebra]{\ref*{def:p03-vaught-formula-algebra}} & \thepage & Nonzero elements correspond to consistent formulas. \\
\phantomsection\label{entry:symbol-registry-03084}\(\alpha,\beta\) & \texttt{\detokenize{\alpha;\beta}} & formulas serving as atoms or consistent nonzero elements & P3 Vaught prime atomic & n/a & \hyperref[def:p03-vaught-formula-algebra]{\ref*{def:p03-vaught-formula-algebra}} & \thepage & Their particular consistency and atom assumptions are specified before each use. \\
\phantomsection\label{entry:symbol-registry-03085}\(\delta\) & \texttt{\detokenize{\delta}} & formula used in a consistency-refinement calculation & P3 Vaught prime atomic & n/a & \hyperref[def:p03-vaught-two-branch-construction]{\ref*{def:p03-vaught-two-branch-construction}} & \thepage & Chosen outside the complete type in the nonprincipal branch; quotient representatives are understood in algebra calculations. \\
\phantomsection\label{entry:symbol-registry-03086}\(P\) & \texttt{\detokenize{P}} & complete finite-tuple type as a dual prime ideal & P3 Vaught prime atomic & n/a & \hyperref[def:p03-vaught-prime-ideals]{\ref*{def:p03-vaught-prime-ideals}} & \thepage & The source ideal direction is the ultrafilter direction. \\
\phantomsection\label{entry:symbol-registry-03087}\(\mathcal P_n(T)\) & \texttt{\detokenize{\mathcal P_n(T)}} & set of complete dual prime ideals of the n-variable algebra & P3 Vaught prime atomic & n/a & \hyperref[def:p03-vaught-prime-ideals]{\ref*{def:p03-vaught-prime-ideals}} & \thepage & Identified with the existing complete-type space only by the proved algebra/type correspondence. \\
\phantomsection\label{entry:symbol-registry-03088}\(P(\mathcal M)\) & \texttt{\detokenize{P(\mathcal M)}} & tuples realizing every member of the complete type & P3 Vaught prime atomic & n/a & \hyperref[def:p03-vaught-prime-ideals]{\ref*{def:p03-vaught-prime-ideals}} & \thepage & Repeated coordinates are permitted; omission means this set is empty. \\
\phantomsection\label{entry:symbol-registry-03089}\(\overline\varphi\) & \texttt{\detokenize{\overline\varphi}} & sentence obtained by naming the free-variable slots & P3 Vaught prime atomic & n/a & \hyperref[def:p03-vaught-prime-ideals]{\ref*{def:p03-vaught-prime-ideals}} & \thepage & Bound variables are freshened when making or reversing the named instance. \\
\phantomsection\label{entry:symbol-registry-03090}\(\overline P\) & \texttt{\detokenize{\overline P}} & complete theory with a realizing tuple named & P3 Vaught prime atomic & n/a & \hyperref[def:p03-vaught-prime-ideals]{\ref*{def:p03-vaught-prime-ideals}} & \thepage & Consists of the base theory and the named type instances. \\
\phantomsection\label{entry:symbol-registry-03091}\([\varphi]\) & \texttt{\detokenize{[\varphi]}} & basic clopen set of types containing the formula & P3 Vaught prime atomic & n/a & \hyperref[thm:p03-vaught-algebra-types]{\ref*{thm:p03-vaught-algebra-types}} & \thepage & An empty set exactly when the formula is inconsistent. \\
\phantomsection\label{entry:symbol-registry-03092}\(C\) & \texttt{\detokenize{C}} & countable set of fresh Henkin constant symbols & P3 Vaught prime atomic & n/a & \hyperref[def:p03-vaught-henkin-data]{\ref*{def:p03-vaught-henkin-data}} & \thepage & Disjoint from the original language; distinct symbols may have equal values. \\
\phantomsection\label{entry:symbol-registry-03093}\(c_i\) & \texttt{\detokenize{c_i}} & fresh Henkin individual constant & P3 Vaught prime atomic & n/a & \hyperref[def:p03-vaught-henkin-data]{\ref*{def:p03-vaught-henkin-data}} & \thepage & The completed construction gives infinitely many names for each of their values. \\
\phantomsection\label{entry:symbol-registry-03094}\(\Gamma_s\) & \texttt{\detokenize{\Gamma_s}} & finite conjunction of accumulated stage sentences & P3 Vaught prime atomic & n/a & \hyperref[def:p03-vaught-henkin-data]{\ref*{def:p03-vaught-henkin-data}} & \thepage & The empty initial conjunction is true; each stage retains consistency. \\
\phantomsection\label{entry:symbol-registry-03095}\(\varphi_s(v_0)\) & \texttt{\detokenize{\varphi_s(v_0)}} & enumerated formula for a witness requirement & P3 Vaught prime atomic & n/a & \hyperref[def:p03-vaught-henkin-data]{\ref*{def:p03-vaught-henkin-data}} & \thepage & All expanded-language formulas with at most the displayed free variable occur. \\
\phantomsection\label{entry:symbol-registry-03096}\(\nu_s\) & \texttt{\detokenize{\nu_s}} & least index of a constant absent from the current formula and stage & P3 Vaught prime atomic & n/a & \hyperref[def:p03-vaught-henkin-data]{\ref*{def:p03-vaught-henkin-data}} & \thepage & Freshness is a syntactic condition on the actual finite sentences. \\
\phantomsection\label{entry:symbol-registry-03097}\(W_s\) & \texttt{\detokenize{W_s}} & Henkin existential-witness sentence & P3 Vaught prime atomic & n/a & \hyperref[def:p03-vaught-henkin-data]{\ref*{def:p03-vaught-henkin-data}} & \thepage & Retained on every atom and atomless branch. \\
\phantomsection\label{entry:symbol-registry-03098}\(P_j,r_j,J\) & \texttt{\detokenize{P_j;r_j;J}} & listed nonprincipal complete types with their finite arities and countable index set & P3 Vaught prime atomic & n/a & \hyperref[thm:p03-vaught-simultaneous-construction]{\ref*{thm:p03-vaught-simultaneous-construction}} & \thepage & Finite and empty lists are allowed; no nonprincipal zero-type exists for the standing complete T. \\
\phantomsection\label{entry:symbol-registry-03099}\(\Pi_s\) & \texttt{\detokenize{\Pi_s}} & stage task combining a type flag and distinct constant symbols & P3 Vaught prime atomic & n/a & \hyperref[def:p03-vaught-two-branch-construction]{\ref*{def:p03-vaught-two-branch-construction}} & \thepage & The empty flag selects only the atom or atomless tuple condition. \\
\phantomsection\label{entry:symbol-registry-03100}\(\bar d,\bar e\) & \texttt{\detokenize{\bar d;\bar e}} & task constant tuple and remaining private constant tuple & P3 Vaught prime atomic & n/a & \hyperref[def:p03-vaught-two-branch-construction]{\ref*{def:p03-vaught-two-branch-construction}} & \thepage & Syntactically disjoint name lists; their values need not be distinct. \\
\phantomsection\label{entry:symbol-registry-03101}\(\gamma_s\) & \texttt{\detokenize{\gamma_s}} & accumulated finite stage together with its new witness sentence & P3 Vaught prime atomic & n/a & \hyperref[def:p03-vaught-two-branch-construction]{\ref*{def:p03-vaught-two-branch-construction}} & \thepage & Its named-variable projection determines the next two-branch choice. \\
\phantomsection\label{entry:symbol-registry-03102}\(\theta_s\) & \texttt{\detokenize{\theta_s}} & formula obtained by replacing the task and private constants by variables & P3 Vaught prime atomic & n/a & \hyperref[def:p03-vaught-two-branch-construction]{\ref*{def:p03-vaught-two-branch-construction}} & \thepage & All bound variables are freshened before this replacement. \\
\phantomsection\label{entry:symbol-registry-03103}\(\beta_s\) & \texttt{\detokenize{\beta_s}} & existential projection to the task variable slots & P3 Vaught prime atomic & n/a & \hyperref[def:p03-vaught-two-branch-construction]{\ref*{def:p03-vaught-two-branch-construction}} & \thepage & Consistent projection; either contains an atom or is atomless. \\
\phantomsection\label{entry:symbol-registry-03104}\(T_1\) & \texttt{\detokenize{T_1}} & countable theory specified in the current source construction & P3 Vaught prime atomic & n/a & \hyperref[def:p03-vaught-relative-interpretation]{\ref*{def:p03-vaught-relative-interpretation}} & \thepage & In the simultaneous Henkin proof it is the new-constant base expansion; in section 4 it is the arbitrary target theory. \\
\phantomsection\label{entry:symbol-registry-03105}\(T_2\) & \texttt{\detokenize{T_2}} & base expansion with every chosen simultaneous stage sentence & \hyperref[annex:p3.vaught.simultaneous]{\ref*{annex:p3.vaught.simultaneous}} & n/a & \hyperref[annex:p3.vaught.simultaneous]{\ref*{annex:p3.vaught.simultaneous}} & \thepage & Consistency follows from all finite stages and compactness. \\
\phantomsection\label{entry:symbol-registry-03106}\(\kappa\) & \texttt{\detokenize{\kappa}} & infinite cardinal bounding the parameter-set size & P3 Vaught prime atomic & n/a & \hyperref[def:p03-vaught-saturation]{\ref*{def:p03-vaught-saturation}} & \thepage & The finite-parameter convention is kappa equal to aleph zero. \\
\phantomsection\label{entry:symbol-registry-03107}\(I\) & \texttt{\detokenize{I}} & relative interpretation on a definable domain & P3 Vaught prime atomic & n/a & \hyperref[def:p03-vaught-relative-interpretation]{\ref*{def:p03-vaught-relative-interpretation}} & \thepage & Uses actual identity on the domain and definable relations; no quotient interpretation is inserted. \\
\phantomsection\label{entry:symbol-registry-03108}\(D(x)\) & \texttt{\detokenize{D(x)}} & nonempty interpreted-domain formula & P3 Vaught prime atomic & n/a & \hyperref[def:p03-vaught-relative-interpretation]{\ref*{def:p03-vaught-relative-interpretation}} & \thepage & Quantifier translations guard the interpreted domain. \\
\phantomsection\label{entry:symbol-registry-03109}\(\rho_R(\bar x)\) & \texttt{\detokenize{\rho_R(\bar x)}} & formula defining an interpreted relation & P3 Vaught prime atomic & n/a & \hyperref[def:p03-vaught-relative-interpretation]{\ref*{def:p03-vaught-relative-interpretation}} & \thepage & Its truth is used only at tuples of interpreted-domain elements. \\
\phantomsection\label{entry:symbol-registry-03110}\(K_d(x)\) & \texttt{\detokenize{K_d(x)}} & formula selecting the unique interpreted distinguished element & P3 Vaught prime atomic & n/a & \hyperref[def:p03-vaught-relative-interpretation]{\ref*{def:p03-vaught-relative-interpretation}} & \thepage & Explicit existence and uniqueness within the interpreted domain. \\
\phantomsection\label{entry:symbol-registry-03111}\(\varphi^I\) & \texttt{\detokenize{\varphi^I}} & translated formula under the relative interpretation & P3 Vaught prime atomic & n/a & \hyperref[def:p03-vaught-relative-interpretation]{\ref*{def:p03-vaught-relative-interpretation}} & \thepage & Truth preservation is supplied by a complete formula induction. \\
\phantomsection\label{entry:symbol-registry-03112}\(I(\mathcal B)\) & \texttt{\detokenize{I(\mathcal B)}} & denotation of an interpretation in the target model & P3 Vaught prime atomic & n/a & \hyperref[def:p03-vaught-relative-interpretation]{\ref*{def:p03-vaught-relative-interpretation}} & \thepage & The interpretation requirement makes it an infinite model of the standing T. \\
\phantomsection\label{entry:symbol-registry-03113}\(\mathcal B^*\) & \texttt{\detokenize{\mathcal B^*}} & countable elementary extension meeting interpreted requirements & P3 Vaught prime atomic & n/a & \hyperref[thm:p03-vaught-interpreted-extension]{\ref*{thm:p03-vaught-interpreted-extension}} & \thepage & Constructed as an actual inclusion by relabeling the elementary-diagram copy. \\
\phantomsection\label{entry:symbol-registry-03114}\(\mathcal B_m\) & \texttt{\detokenize{\mathcal B_m}} & outer elementary-chain target stage & P3 Vaught prime atomic & n/a & \hyperref[annex:p3.vaught.saturated-existence]{\ref*{annex:p3.vaught.saturated-existence}} & \thepage & Repeats the inner filling operation so new finite tuples are also covered. \\
\phantomsection\label{entry:symbol-registry-03115}\(\mathcal C_s\) & \texttt{\detokenize{\mathcal C_s}} & inner elementary-chain stage for one fixed target's tasks & \hyperref[annex:p3.vaught.saturated-existence]{\ref*{annex:p3.vaught.saturated-existence}} & n/a & \hyperref[annex:p3.vaught.saturated-existence]{\ref*{annex:p3.vaught.saturated-existence}} & \thepage & Every old tuple requirement appears once in the countable inner enumeration. \\
\phantomsection\label{entry:symbol-registry-03116}\(\mathcal B_\infty\) & \texttt{\detokenize{\mathcal B_\infty}} & union of the outer elementary chain & \hyperref[annex:p3.vaught.saturated-existence]{\ref*{annex:p3.vaught.saturated-existence}} & n/a & \hyperref[annex:p3.vaught.saturated-existence]{\ref*{annex:p3.vaught.saturated-existence}} & \thepage & Its interpreted domains commute with the union and satisfy all one-coordinate extension requirements. \\
\phantomsection\label{entry:symbol-registry-03117}\(\chi_s\) & \texttt{\detokenize{\chi_s}} & nonzero node in a binary Boolean splitting tree & \hyperref[annex:p3.vaught.countable-types-atomistic]{\ref*{annex:p3.vaught.countable-types-atomistic}} & n/a & \hyperref[annex:p3.vaught.countable-types-atomistic]{\ref*{annex:p3.vaught.countable-types-atomistic}} & \thepage & Here s is a finite binary string; separate from the natural-number Henkin stage formula. \\
\phantomsection\label{entry:symbol-registry-03118}\(P_\eta\) & \texttt{\detokenize{P_\eta}} & complete type extending the chosen infinite binary branch & \hyperref[annex:p3.vaught.countable-types-atomistic]{\ref*{annex:p3.vaught.countable-types-atomistic}} & n/a & \hyperref[annex:p3.vaught.countable-types-atomistic]{\ref*{annex:p3.vaught.countable-types-atomistic}} & \thepage & Disjoint sibling formulas ensure different branches give different types. \\
\phantomsection\label{entry:symbol-registry-03119}\(T_{\rm dlo}\) & \texttt{\detokenize{T_{\rm dlo}}} & complete theory of dense linear order without endpoints & P3 Vaught prime atomic & n/a & \hyperref[thm:p03-vaught-dlo-example]{\ref*{thm:p03-vaught-dlo-example}} & \thepage & Its finite tuple types are exactly the equality and strict-order patterns. \\
\phantomsection\label{entry:symbol-registry-03120}\(T_{\mathbb Q}\) & \texttt{\detokenize{T_{\mathbb Q}}} & complete dense-order theory naming every rational & P3 Vaught prime atomic & n/a & \hyperref[thm:p03-vaught-dlo-example]{\ref*{thm:p03-vaught-dlo-example}} & \thepage & Countable naming changes the empty-parameter type space and destroys smallness. \\
\phantomsection\label{entry:symbol-registry-03121}\(T_{\rm disc}\) & \texttt{\detokenize{T_{\rm disc}}} & complete theory of the ordered natural numbers & P3 Vaught prime atomic & n/a & \hyperref[thm:p03-vaught-discrete-example]{\ref*{thm:p03-vaught-discrete-example}} & \thepage & The prime calculation uses unique finite-position formulas rather than an unsupplied decision procedure. \\
\phantomsection\label{entry:symbol-registry-03122}\(\gamma_k(x)\) & \texttt{\detokenize{\gamma_k(x)}} & formula saying exactly k predecessors & \hyperref[annex:p3.vaught.discrete-example]{\ref*{annex:p3.vaught.discrete-example}} & n/a & \hyperref[annex:p3.vaught.discrete-example]{\ref*{annex:p3.vaught.discrete-example}} & \thepage & Isolates the standard position k in the ordered natural model. \\
\phantomsection\label{entry:symbol-registry-03123}\(\delta_k(x)\) & \texttt{\detokenize{\delta_k(x)}} & formula saying at least k predecessors & \hyperref[annex:p3.vaught.discrete-example]{\ref*{annex:p3.vaught.discrete-example}} & n/a & \hyperref[annex:p3.vaught.discrete-example]{\ref*{annex:p3.vaught.discrete-example}} & \thepage & Their consistent partial type is nonprincipal and every completion is omitted by the prime model. \\
\phantomsection\label{entry:symbol-registry-03124}\(T_{\rm bits}\) & \texttt{\detokenize{T_{\rm bits}}} & complete independent-unary-predicate theory & P3 Vaught prime atomic & n/a & \hyperref[thm:p03-vaught-unary-obstruction]{\ref*{thm:p03-vaught-unary-obstruction}} & \thepage & Every finite Boolean predicate pattern has infinitely many realizations. \\
\phantomsection\label{entry:symbol-registry-03125}\(U_i\) & \texttt{\detokenize{U_i}} & indexed unary predicate & P3 Vaught prime atomic & n/a & \hyperref[thm:p03-vaught-unary-obstruction]{\ref*{thm:p03-vaught-unary-obstruction}} & \thepage & Changing a previously unmentioned predicate gives two consistent refinements. \\
\phantomsection\label{entry:symbol-registry-03126}\(G_f\) & \texttt{\detokenize{G_f}} & total single-valued graph relation replacing a function & P3 Vaught prime atomic & n/a & \hyperref[thm:p03-vaught-functional-bridge]{\ref*{thm:p03-vaught-functional-bridge}} & \thepage & Zero-arity functions remain individual constants. \\
\phantomsection\label{entry:symbol-registry-03127}\(E_t(\bar x,z)\) & \texttt{\detokenize{E_t(\bar x,}\allowbreak\detokenize{z)}} & term-value formula in the graph language & \hyperref[annex:p3.vaught.functional-bridge]{\ref*{annex:p3.vaught.functional-bridge}} & n/a & \hyperref[annex:p3.vaught.functional-bridge]{\ref*{annex:p3.vaught.functional-bridge}} & \thepage & Term induction proves existence and uniqueness of its selected value. \\
\phantomsection\label{entry:symbol-registry-03128}\(S^{\rm gr}\) & \texttt{\detokenize{S^{\rm gr}}} & translated theory together with all total graph axioms & \hyperref[annex:p3.vaught.functional-bridge]{\ref*{annex:p3.vaught.functional-bridge}} & n/a & \hyperref[annex:p3.vaught.functional-bridge]{\ref*{annex:p3.vaught.functional-bridge}} & \thepage & Its models correspond bijectively to the function-language models on the same domain. \\
\phantomsection\label{entry:symbol-registry-03129}\(p_j\) & \texttt{\detokenize{p_j}} & listed nonprincipal partial type relative to an arbitrary satisfiable theory & P3 Vaught prime atomic & n/a & \hyperref[thm:p03-vaught-partial-omission]{\ref*{thm:p03-vaught-partial-omission}} & \thepage & The broader omission theorem does not assume completeness or replace a partial type by an arbitrary completion. \\
\phantomsection\label{entry:symbol-registry-03130}\(L\) & \texttt{\detokenize{L}} & countable first-order language with logical identity & P3 categoricity finite types & n/a & \hyperref[def:p03-rn-standing]{\ref*{def:p03-rn-standing}} & \thepage & The main criterion allows constants and finite-arity functions; later boundaries explicitly change the standing. \\
\phantomsection\label{entry:symbol-registry-03131}\(T\) & \texttt{\detokenize{T}} & complete satisfiable infinite-model theory & P3 categoricity finite types & n/a & \hyperref[def:p03-rn-standing]{\ref*{def:p03-rn-standing}} & \thepage & Finite and incomplete cases state their changed hypotheses before use. \\
\phantomsection\label{entry:symbol-registry-03132}\(\bar x\) & \texttt{\detokenize{\bar x}} & ordered finite variable slots & P3 categoricity finite types & n/a & \hyperref[def:p03-rn-standing]{\ref*{def:p03-rn-standing}} & \thepage & Repeated tuple coordinates and zero slots are allowed. \\
\phantomsection\label{entry:symbol-registry-03133}\(n\) & \texttt{\detokenize{n}} & finite tuple arity & P3 categoricity finite types & n/a & \hyperref[def:p03-rn-standing]{\ref*{def:p03-rn-standing}} & \thepage & The criterion quantifies over every finite arity; there is no uniform numerical bound. \\
\phantomsection\label{entry:symbol-registry-03134}\(\operatorname{tp}^{\mathcal M}(\bar a)\) & \texttt{\detokenize{\operatorname{tp}}} & complete type of an interpreted tuple & P3 categoricity finite types & n/a & \hyperref[def:p03-rn-standing]{\ref*{def:p03-rn-standing}} & \thepage & The model superscript records where formula truth is evaluated. \\
\phantomsection\label{entry:symbol-registry-03135}\(\mathbb B_n(T)\) & \texttt{\detokenize{\mathbb B_n(T)}} & Boolean quotient of formulas modulo T-equivalence & P3 categoricity finite types & n/a & \hyperref[def:p03-rn-standing]{\ref*{def:p03-rn-standing}} & \thepage & Explicit reuse of the Vaught formula-algebra interface. \\
\phantomsection\label{entry:symbol-registry-03136}\(p\) & \texttt{\detokenize{p}} & complete finite-tuple type & P3 categoricity finite types & n/a & \hyperref[thm:p03-rn-type-realization]{\ref*{thm:p03-rn-type-realization}} & \thepage & Its consistency is a standing assumption before realization or isolation. \\
\phantomsection\label{entry:symbol-registry-03137}\(\theta(\bar x,y)\) & \texttt{\detokenize{\theta}} & consistent isolator of a joint type & P3 categoricity finite types & n/a & \hyperref[thm:p03-rn-isolated-extension]{\ref*{thm:p03-rn-isolated-extension}} & \thepage & Its existential projection permits extension over the fixed tuple. \\
\phantomsection\label{entry:symbol-registry-03138}\(q\) & \texttt{\detokenize{q}} & first-slot restriction of the isolated joint type & P3 categoricity finite types & n/a & \hyperref[annex:p3.rn.isolated-extension]{\ref*{annex:p3.rn.isolated-extension}} & \thepage & An auxiliary proof type not a new global parameter convention. \\
\phantomsection\label{entry:symbol-registry-03139}\(r(\bar y)\) & \texttt{\detokenize{r}} & complete type over a named finite parameter tuple & P3 categoricity finite types & n/a & \hyperref[annex:p3.rn.isolated-extension]{\ref*{annex:p3.rn.isolated-extension}} & \thepage & Consistency is with the full named theory; its joint empty-parameter type is separately formed. \\
\phantomsection\label{entry:symbol-registry-03140}\(d_i\) & \texttt{\detokenize{d_i}} & fresh coordinate constant symbols & P3 categoricity finite types & n/a & \hyperref[annex:p3.rn.type-realization]{\ref*{annex:p3.rn.type-realization}} & \thepage & Distinct symbols do not impose distinct interpretations. \\
\phantomsection\label{entry:symbol-registry-03141}\(\operatorname{Aut}(\mathcal M)\) & \texttt{\detokenize{\operatorname{Aut}}} & automorphism group acting on finite tuples & P3 categoricity finite types & n/a & \hyperref[thm:p03-rn-orbits]{\ref*{thm:p03-rn-orbits}} & \thepage & The prescribed tuple map is retained throughout back-and-forth. \\
\phantomsection\label{entry:symbol-registry-03142}\(m\) & \texttt{\detokenize{m}} & size of a finite nonempty model & P3 categoricity finite types & n/a & \hyperref[thm:p03-rn-finite-boundary]{\ref*{thm:p03-rn-finite-boundary}} & \thepage & This changed finite-theory case has no infinite model. \\
\phantomsection\label{entry:symbol-registry-03143}\(\rho_m\) & \texttt{\detokenize{\rho_m}} & exact-cardinality sentence for positive finite m & P3 categoricity finite types & n/a & \hyperref[annex:p3.rn.incomplete-boundary]{\ref*{annex:p3.rn.incomplete-boundary}} & \thepage & The different complete one-types witness failure without completeness. \\
\phantomsection\label{entry:symbol-registry-03144}\(c_i\) & \texttt{\detokenize{c_i}} & constant naming the natural number i in the uncountable-language example & P3 categoricity finite types & n/a & \hyperref[thm:p03-rn-uncountable-language]{\ref*{thm:p03-rn-uncountable-language}} & \thepage & All named constants are distinct; unnamed elements force continuum cardinality. \\
\phantomsection\label{entry:symbol-registry-03145}\(f_g\) & \texttt{\detokenize{f_g}} & unary symbol for every function from omega to omega & P3 categoricity finite types & n/a & \hyperref[thm:p03-rn-uncountable-language]{\ref*{thm:p03-rn-uncountable-language}} & \thepage & The complete full-function theory is an explicit uncountable-signature boundary. \\
\phantomsection\label{entry:symbol-registry-03146}\(e\) & \texttt{\detokenize{e}} & injection of finite binary strings into omega & P3 categoricity finite types & n/a & \hyperref[annex:p3.rn.uncountable-language]{\ref*{annex:p3.rn.uncountable-language}} & \thepage & Used to build a continuum family of functions with pairwise finite agreement. \\
\phantomsection\label{entry:symbol-registry-03147}\(\eta,\nu\) & \texttt{\detokenize{\eta;\nu}} & infinite binary sequences in the cardinality counterexample & P3 categoricity finite types & n/a & \hyperref[annex:p3.rn.uncountable-language]{\ref*{annex:p3.rn.uncountable-language}} & \thepage & The finite agreement lemma assumes they are different. \\
\phantomsection\label{entry:symbol-registry-03148}\(g_\eta\) & \texttt{\detokenize{g_\eta}} & function coding successive prefixes of eta & P3 categoricity finite types & n/a & \hyperref[annex:p3.rn.uncountable-language]{\ref*{annex:p3.rn.uncountable-language}} & \thepage & Its image at an unnamed point distinguishes it from every other coded function. \\
\phantomsection\label{entry:symbol-registry-03149}\(E_{\eta,\nu}\) & \texttt{\detokenize{E_{\eta,\nu}}} & finite agreement set of two coded functions & P3 categoricity finite types & n/a & \hyperref[annex:p3.rn.uncountable-language]{\ref*{annex:p3.rn.uncountable-language}} & \thepage & Each displayed disjunction is genuinely finite. \\
\phantomsection\label{entry:symbol-registry-03150}\(T_{\rm dlo}\) & \texttt{\detokenize{T_{\rm dlo}}} & dense linear order theory without endpoints & P3 categoricity finite types & n/a & \hyperref[def:p03-rn-dlo-data]{\ref*{def:p03-rn-dlo-data}} & \thepage & Only one strict-order relation and logical identity; no constants or extra predicates. \\
\phantomsection\label{entry:symbol-registry-03151}\((B_1,\ldots,B_k)\) & \texttt{\detokenize{B_i}} & ordered partition of the coordinate indices & P3 categoricity finite types & n/a & \hyperref[def:p03-rn-dlo-data]{\ref*{def:p03-rn-dlo-data}} & \thepage & The blocks are nonempty except for the single zero-arity empty pattern. \\
\phantomsection\label{entry:symbol-registry-03152}\(\delta_{(B_1,\ldots,B_k)}\) & \texttt{\detokenize{\delta}} & finite equality-and-order pattern formula & P3 categoricity finite types & n/a & \hyperref[def:p03-rn-dlo-data]{\ref*{def:p03-rn-dlo-data}} & \thepage & All patterns are realized and later proved to isolate complete types. \\
\phantomsection\label{entry:symbol-registry-03153}\({n\brace k}\) & \texttt{\detokenize{\brace}} & number of unordered k-block partitions of n indices & P3 categoricity finite types & n/a & \hyperref[def:p03-rn-dlo-data]{\ref*{def:p03-rn-dlo-data}} & \thepage & Zero-arity convention is explicit before the exact type-count sum. \\
\phantomsection\label{entry:symbol-registry-03154}\(a_i,b_i\) & \texttt{\detokenize{a_i;b_i}} & increasing matched points of a finite order map & P3 categoricity finite types & n/a & \hyperref[annex:p3.rn.dlo-extension]{\ref*{annex:p3.rn.dlo-extension}} & \thepage & Source and target points share exactly the required comparisons. \\
\phantomsection\label{entry:symbol-registry-03155}\(\mathcal Z\) & \texttt{\detokenize{\mathcal Z}} & integer successor-predecessor structure & P3 categoricity finite types & n/a & \hyperref[thm:p03-rn-function-boundary]{\ref*{thm:p03-rn-function-boundary}} & \thepage & A finite function signature with no nonempty finite closed substructures. \\
\phantomsection\label{entry:symbol-registry-03156}\(s,t\) & \texttt{\detokenize{s;t}} & successor and predecessor unary functions & P3 categoricity finite types & n/a & \hyperref[thm:p03-rn-function-boundary]{\ref*{thm:p03-rn-function-boundary}} & \thepage & Closure under both generates the whole integer domain. \\
\phantomsection\label{entry:symbol-registry-03157}\(s^k(x)\) & \texttt{\detokenize{s^k}} & iterated successor term & P3 categoricity finite types & n/a & \hyperref[thm:p03-rn-function-boundary]{\ref*{thm:p03-rn-function-boundary}} & \thepage & Infinitely many atomic distance formulas remain in the two-symbol signature. \\
\phantomsection\label{entry:symbol-registry-03158}\(I(T,\aleph_0)\) & \texttt{\detokenize{I(T,}\allowbreak\detokenize{\aleph_0)}} & countable-model isomorphism spectrum & \hyperref[sec:p03-vaught-never-two-presentation]{\ref*{sec:p03-vaught-never-two-presentation}} & n/a & \hyperref[def:p03-vaught-spectrum-notation]{\ref*{def:p03-vaught-spectrum-notation}} & \thepage & Finite models included under the explicitly stated countable convention. \\
\phantomsection\label{entry:symbol-registry-03159}\(F_m(T)\) & \texttt{\detokenize{F_m(T)}} & Boolean algebra of formulas modulo theory equivalence & \hyperref[sec:p03-vaught-never-two-presentation]{\ref*{sec:p03-vaught-never-two-presentation}} & n/a & \hyperref[def:p03-vaught-spectrum-notation]{\ref*{def:p03-vaught-spectrum-notation}} & \thepage & Matches the original formula-algebra role. \\
\phantomsection\label{entry:symbol-registry-03160}\(\mathfrak P_m(T)\) & \texttt{\detokenize{\mathfrak P_m(T)}} & Vaught dual prime ideals corresponding to complete types & \hyperref[sec:p03-vaught-never-two-presentation]{\ref*{sec:p03-vaught-never-two-presentation}} & n/a & \hyperref[def:p03-vaught-spectrum-notation]{\ref*{def:p03-vaught-spectrum-notation}} & \thepage & Original dual ideal terminology is preserved. \\
\phantomsection\label{entry:symbol-registry-03161}\(T_p\) & \texttt{\detokenize{T_p}} & complete theory after naming a realizing tuple & \hyperref[sec:p03-vaught-never-two-presentation]{\ref*{sec:p03-vaught-never-two-presentation}} & n/a & \hyperref[def:p03-vaught-spectrum-notation]{\ref*{def:p03-vaught-spectrum-notation}} & \thepage & Tuple entries may repeat; complete type decides equalities. \\
\phantomsection\label{entry:symbol-registry-03162}\(T_{r+2}\) & \texttt{\detokenize{T_{r+2}}} & colored Ehrenfeucht order theory & \hyperref[sec:p03-vaught-never-two-presentation]{\ref*{sec:p03-vaught-never-two-presentation}} & n/a & \hyperref[def:p03-vaught-colored-orders]{\ref*{def:p03-vaught-colored-orders}} & \thepage & Positive finite color count. \\
\phantomsection\label{entry:symbol-registry-03163}\(T'_{r+2}\) & \texttt{\detokenize{T'_{r+2}}} & finite-language quotient coding theory & \hyperref[sec:p03-vaught-never-two-presentation]{\ref*{sec:p03-vaught-never-two-presentation}} & n/a & \hyperref[def:p03-vaught-colored-orders]{\ref*{def:p03-vaught-colored-orders}} & \thepage & Countable language theory and finite language coding remain distinct. \\
\phantomsection\label{entry:symbol-registry-03164}\(C_i\) & \texttt{\detokenize{C_i}} & unique quotient class of positive finite size & \hyperref[sec:p03-vaught-never-two-presentation]{\ref*{sec:p03-vaught-never-two-presentation}} & n/a & \hyperref[def:p03-vaught-colored-orders]{\ref*{def:p03-vaught-colored-orders}} & \thepage & External quotient-point notation; no extra constants in the coding language. \\
\phantomsection\label{entry:symbol-registry-03165}\(SR(M),SR_p(M)\) & \texttt{\detokenize{SR(M),}\allowbreak\detokenize{SR_p(M)}} & unparameterized and parameterized Scott ranks & \hyperref[sec:p03-vaught-never-two-presentation]{\ref*{sec:p03-vaught-never-two-presentation}} & n/a & \hyperref[def:p03-vaught-scott-comparison]{\ref*{def:p03-vaught-scott-comparison}} & \thepage & Exact selected 2026 complexity convention. \\
\phantomsection\label{entry:symbol-registry-03166}\(\operatorname{vo}(\sigma)\) & \texttt{\detokenize{\operatorname{vo}(\sigma)}} & Vaught ordinal for an infinitary sentence & \hyperref[sec:p03-vaught-never-two-presentation]{\ref*{sec:p03-vaught-never-two-presentation}} & n/a & \hyperref[def:p03-vaught-scott-comparison]{\ref*{def:p03-vaught-scott-comparison}} & \thepage & Bounded ordinal comparison is not a general conjecture proof. \\
\phantomsection\label{entry:symbol-registry-03167}\(U_j\) & \texttt{\detokenize{U_j}} & dense unary color class & \hyperref[sec:p03-vaught-never-two-presentation]{\ref*{sec:p03-vaught-never-two-presentation}} & n/a & \hyperref[def:p03-vaught-colored-orders]{\ref*{def:p03-vaught-colored-orders}} & \thepage & Class-constant in the finite-language quotient coding. \\
\phantomsection\label{entry:symbol-registry-03168}\(R\) & \texttt{\detokenize{R}} & equivalence relation of coding fibers & \hyperref[sec:p03-vaught-never-two-presentation]{\ref*{sec:p03-vaught-never-two-presentation}} & n/a & \hyperref[def:p03-vaught-colored-orders]{\ref*{def:p03-vaught-colored-orders}} & \thepage & The order is constant on these fibers. \\
\phantomsection\label{entry:symbol-registry-03169}\(c_i\) & \texttt{\detokenize{c_i}} & increasing named order points & \hyperref[sec:p03-vaught-never-two-presentation]{\ref*{sec:p03-vaught-never-two-presentation}} & n/a & \hyperref[def:p03-vaught-colored-orders]{\ref*{def:p03-vaught-colored-orders}} & \thepage & Removed from the finite-language coding. \\
\phantomsection\label{entry:symbol-registry-03170}\(\Sigma^{\mathrm{in}}_\alpha,\Pi^{\mathrm{in}}_\alpha\) & \texttt{\detokenize{\Sigma^{\mathrm{in}}_\alpha,}\allowbreak\detokenize{\Pi^{\mathrm{in}}_\alpha}} & infinitary formula complexity classes & \hyperref[sec:p03-vaught-never-two-presentation]{\ref*{sec:p03-vaught-never-two-presentation}} & n/a & \hyperref[def:p03-vaught-scott-comparison]{\ref*{def:p03-vaught-scott-comparison}} & \thepage & Countable ordinals with explicitly stated base and recursion. \\
\phantomsection\label{entry:symbol-registry-03454}\(\Sigma^*\) & \texttt{\detokenize{\Sigma^*}} & Morley definitional relational expansion & P3 Morley1965 & n/a & \hyperref[def:p03-morley-standing]{\ref*{def:p03-morley-standing}} & \thepage & Not Skolemization; preserves original model universes and elementary maps. \\
\phantomsection\label{entry:symbol-registry-03455}\(T\) & \texttt{\detokenize{T}} & complete countable Morleyized infinite theory & P3 Morley1965 & n/a & \hyperref[def:p03-morley-standing]{\ref*{def:p03-morley-standing}} & \thepage & Standing throughout the source-native reconstruction. \\
\phantomsection\label{entry:symbol-registry-03456}\(\mathcal N(T)\) & \texttt{\detokenize{\mathcal N(T)}} & class of induced parameter subsystems & P3 Morley1965 & n/a & \hyperref[def:p03-morley-standing]{\ref*{def:p03-morley-standing}} & \thepage & Source writes a script class symbol; subsystems need not be models. \\
\phantomsection\label{entry:symbol-registry-03457}\(|A|\) & \texttt{\detokenize{|A|}} & universe of a parameter system & P3 Morley1965 & n/a & \hyperref[def:p03-morley-standing]{\ref*{def:p03-morley-standing}} & \thepage & Published source universe notation; cardinality is kappa(A) equal to kappa applied to this universe. \\
\phantomsection\label{entry:symbol-registry-03458}\(\kappa(X)\) & \texttt{\detokenize{\kappa(X)}} & cardinality operation on a set & P3 Morley1965 & n/a & \hyperref[def:p03-morley-standing]{\ref*{def:p03-morley-standing}} & \thepage & Published source notation; applies also to type spaces and extension fibers and ordinal sets. \\
\phantomsection\label{entry:symbol-registry-03459}\(\kappa(A)\) & \texttt{\detokenize{\kappa(A)}} & cardinality or power of a parameter system & P3 Morley1965 & n/a & \hyperref[def:p03-morley-standing]{\ref*{def:p03-morley-standing}} & \thepage & Published abbreviation kappa(A) equals kappa of the universe vertical-bar A vertical-bar; not a second vertical-bar glyph. \\
\phantomsection\label{entry:symbol-registry-03460}\(T(A)\) & \texttt{\detokenize{T(A)}} & complete theory with named parameter diagram & P3 Morley1965 & n/a & \hyperref[def:p03-morley-type-space]{\ref*{def:p03-morley-type-space}} & \thepage & Morleyization makes the induced diagram full elementary information. \\
\phantomsection\label{entry:symbol-registry-03461}\(F(A)\) & \texttt{\detokenize{F(A)}} & Boolean algebra of one-variable parameter formulas & P3 Morley1965 & n/a & \hyperref[def:p03-morley-type-space]{\ref*{def:p03-morley-type-space}} & \thepage & Quotient by equivalence in the named complete parameter theory. \\
\phantomsection\label{entry:symbol-registry-03462}\(S(A)\) & \texttt{\detokenize{S(A)}} & Stone space of complete one-element parameter types & P3 Morley1965 & n/a & \hyperref[def:p03-morley-type-space]{\ref*{def:p03-morley-type-space}} & \thepage & Not a model universe or a space of complete theories without parameter standing. \\
\phantomsection\label{entry:symbol-registry-03463}\(U_\varphi\) & \texttt{\detokenize{U_\varphi}} & basic clopen of types containing a formula & P3 Morley1965 & n/a & \hyperref[def:p03-morley-type-space]{\ref*{def:p03-morley-type-space}} & \thepage & Always a clopen subset of the specified type space. \\
\phantomsection\label{entry:symbol-registry-03464}\(f^*\) & \texttt{\detokenize{f^*}} & continuous onto restriction map dual to a monomorphism & P3 Morley1965 & n/a & \hyperref[def:p03-morley-type-space]{\ref*{def:p03-morley-type-space}} & \thepage & Contravariant; inverse-image notation is explicitly typed. \\
\phantomsection\label{entry:symbol-registry-03465}\(i_{AB}^*\) & \texttt{\detokenize{i_{AB}^*}} & restriction along an inclusion of parameter systems & P3 Morley1965 & n/a & \hyperref[def:p03-morley-type-space]{\ref*{def:p03-morley-type-space}} & \thepage & Maps S(B) to S(A). \\
\phantomsection\label{entry:symbol-registry-03466}\(S^\alpha(A)\) & \texttt{\detokenize{S^\alpha(A)}} & global extension-sensitive derivative of the type space & P3 Morley1965 & n/a & \hyperref[def:p03-morley-global-rank]{\ref*{def:p03-morley-global-rank}} & \thepage & Not the ordinary Cantor Bendixson derivative of a single space. \\
\phantomsection\label{entry:symbol-registry-03467}\(\operatorname{Tr}^\alpha(A)\) & \texttt{\detokenize{\operatorname{Tr}^\alpha(A)}} & points removed at ordinal rank alpha & P3 Morley1965 & n/a & \hyperref[def:p03-morley-global-rank]{\ref*{def:p03-morley-global-rank}} & \thepage & Defined by isolation of every tested extension fiber. \\
\phantomsection\label{entry:symbol-registry-03468}\(d(p)\) & \texttt{\detokenize{d(p)}} & uniform finite multiplicity of a rank-layer type fiber & P3 Morley1965 & n/a & \hyperref[thm:p03-morley-degree]{\ref*{thm:p03-morley-degree}} & \thepage & Positive integer; additive under extension; not dimension. \\
\phantomsection\label{entry:symbol-registry-03469}\(\alpha_T\) & \texttt{\detokenize{\alpha_T}} & first simultaneous stabilization ordinal & P3 Morley1965 & n/a & \hyperref[def:p03-morley-total-transcendence]{\ref*{def:p03-morley-total-transcendence}} & \thepage & The residual need not be empty unless total transcendence is assumed. \\
\phantomsection\label{entry:symbol-registry-03470}\(\operatorname{RM}(X)\) & \texttt{\detokenize{\operatorname{RM}(X)}} & definable splitting rank in finite arity & P3 Morley1965 & n/a & \hyperref[def:p03-morley-definable-rank]{\ref*{def:p03-morley-definable-rank}} & \thepage & Editorial bridge is proved and kept distinct from point-degree notation. \\
\phantomsection\label{entry:symbol-registry-03471}\(K(A)\) & \texttt{\detokenize{K(A)}} & field generated by the parameter subsystem & P3 Morley1965 & n/a & \hyperref[def:p03-morley-field-data]{\ref*{def:p03-morley-field-data}} & \thepage & Fixed-characteristic field example; not algebraic closure. \\
\phantomsection\label{entry:symbol-registry-03472}\(\mathrm{ACF}_p\) & \texttt{\detokenize{\mathrm{ACF}_p}} & complete algebraically closed field theory at fixed characteristic & P3 Morley1965 & n/a & \hyperref[def:p03-morley-field-data]{\ref*{def:p03-morley-field-data}} & \thepage & p is zero or a fixed prime. \\
\phantomsection\label{entry:symbol-registry-03473}\(R_0\) & \texttt{\detokenize{R_0}} & infinite unary base part in the tuple-code example & P3 Morley1965 & n/a & \hyperref[def:p03-morley-tuple-code-data]{\ref*{def:p03-morley-tuple-code-data}} & \thepage & Not the separately indexed predicates of the Cantor-space example. \\
\phantomsection\label{entry:symbol-registry-03474}\(R_1\) & \texttt{\detokenize{R_1}} & relation assigning a unique code to an ordered distinct tuple & P3 Morley1965 & n/a & \hyperref[def:p03-morley-tuple-code-data]{\ref*{def:p03-morley-tuple-code-data}} & \thepage & Arity n plus one with n at least one. \\
\phantomsection\label{entry:symbol-registry-03475}\(\overline B\) & \texttt{\detokenize{\overline B}} & closure under taking coordinates and forming tuple codes & P3 Morley1965 & n/a & \hyperref[def:p03-morley-tuple-code-data]{\ref*{def:p03-morley-tuple-code-data}} & \thepage & Finite definability closure in this particular example. \\
\phantomsection\label{entry:symbol-registry-03476}\(Y^{(\beta)}\) & \texttt{\detokenize{Y^{(\beta)}}} & ordinary Cantor Bendixson derivative of a closed Cantor subset & P3 Morley1965 & n/a & \hyperref[def:p03-morley-unary-data]{\ref*{def:p03-morley-unary-data}} & \thepage & Compared with global S derivatives only by the proved unary-example bridge. \\
\phantomsection\label{entry:symbol-registry-03477}\(\alpha_Y\) & \texttt{\detokenize{\alpha_Y}} & first stabilization ordinal of the ordinary derivative & P3 Morley1965 & n/a & \hyperref[def:p03-morley-unary-data]{\ref*{def:p03-morley-unary-data}} & \thepage & Source formula alpha T equals ordinal one plus alpha Y. \\
\phantomsection\label{entry:symbol-registry-03478}\(\Sigma^\#\) & \texttt{\detokenize{\Sigma^\#}} & universal conservative Skolem expansion & P3 Morley1965 & n/a & \hyperref[def:p03-morley-skolem-hull]{\ref*{def:p03-morley-skolem-hull}} & \thepage & Distinct from star relational expansion; combination star hash is explicit. \\
\phantomsection\label{entry:symbol-registry-03479}\(M(X,A)\) & \texttt{\detokenize{M(X,}\allowbreak\detokenize{A)}} & substructure generated by term values on the ordered generators & P3 Morley1965 & n/a & \hyperref[def:p03-morley-skolem-hull]{\ref*{def:p03-morley-skolem-hull}} & \thepage & Includes constant terms; universal Skolem axioms make it a model. \\
\phantomsection\label{entry:symbol-registry-03480}\(A_\xi\) & \texttt{\detokenize{A_\xi}} & continuous parameter-system chain & P3 Morley1965 & n/a & \hyperref[thm:p03-morley-chain-types]{\ref*{thm:p03-morley-chain-types}} & \thepage & Later prime-chain and omission-transfer uses specify their bases. \\
\phantomsection\label{entry:symbol-registry-03481}\(B_\xi\) & \texttt{\detokenize{B_\xi}} & compatible prime-model chain & P3 Morley1965 & n/a & \hyperref[thm:p03-morley-chain-primes]{\ref*{thm:p03-morley-chain-primes}} & \thepage & Continuity is a proved construction condition rather than automatic for independently selected primes. \\
\phantomsection\label{entry:symbol-registry-03482}\(\nu\) & \texttt{\detokenize{\nu}} & ordinal rank minimized or stabilized in chain arguments & P3 Morley1965 & n/a & \hyperref[thm:p03-morley-chain-types]{\ref*{thm:p03-morley-chain-types}} & \thepage & Not a cardinal or number of generators. \\
\phantomsection\label{entry:symbol-registry-03483}\(\kappa\) & \texttt{\detokenize{\kappa}} & cardinal variable & P3 Morley1965 & n/a & \hyperref[def:p03-morley-standing]{\ref*{def:p03-morley-standing}} & \thepage & Published bare cardinal variable; argument-bearing kappa is the separately declared cardinality operation. Uncountable and regular hypotheses are stated where needed. \\
\phantomsection\label{entry:symbol-registry-03484}\(A'\) & \texttt{\detokenize{A'}} & countable retained parameter subsystem witnessing omission & P3 Morley1965 & n/a & \hyperref[thm:p03-morley-countable-omission-witness]{\ref*{thm:p03-morley-countable-omission-witness}} & \thepage & Contained in the countable elementary model; not itself assumed to be a model. \\
\phantomsection\label{entry:symbol-registry-03485}\(q\) & \texttt{\detokenize{q}} & complete type over the retained subsystem omitted by the witness model & P3 Morley1965 & n/a & \hyperref[thm:p03-morley-countable-omission-witness]{\ref*{thm:p03-morley-countable-omission-witness}} & \thepage & The transferred omitted type remains over the same countable base. \\
\phantomsection\label{entry:symbol-registry-03486}\(V,W,R,\rho\) & \texttt{\detokenize{V;W;R;\rho}} & variables operations relations and arity function & P3 Lyndon 1959 preservation & n/a & \hyperref[def:p03-lyndon-language]{\ref*{def:p03-lyndon-language}} & \thepage & Pairwise disjoint symbol sets; V infinite; operation and relation arities are finite. \\
\phantomsection\label{entry:symbol-registry-03487}\(0,1\) & \texttt{\detokenize{0;1}} & false and true formula constants & P3 Lyndon 1959 preservation & object & \hyperref[def:p03-lyndon-language]{\ref*{def:p03-lyndon-language}} & \thepage & Source truth constants; 0 is allowed in the positive fragment and empty disjunction; 1 is empty conjunction. \\
\phantomsection\label{entry:symbol-registry-03488}\(e\) & \texttt{\detokenize{e}} & distinguished auxiliary binary relation later genuine identity & P3 Lyndon 1959 preservation & object & \hyperref[def:p03-lyndon-language]{\ref*{def:p03-lyndon-language}} & \thepage & Before quotienting this is a freely interpreted relation symbol; relational systems interpret it as identity. \\
\phantomsection\label{entry:symbol-registry-03489}\(Q,Q'\) & \texttt{\detokenize{Q;Q'}} & selected relation symbols and their disjoint primed copies & P3 Lyndon 1959 preservation & n/a & \hyperref[def:p03-lyndon-polarity]{\ref*{def:p03-lyndon-polarity}} & \thepage & Same respective arities; Q-positive permits negative atoms outside Q. \\
\phantomsection\label{entry:symbol-registry-03490}\(F'\) & \texttt{\detokenize{F'}} & simultaneous relation-symbol replacement in a formula & P3 Lyndon 1959 preservation & n/a & \hyperref[def:p03-lyndon-polarity]{\ref*{def:p03-lyndon-polarity}} & \thepage & Replace precisely symbols in Q by corresponding Q-prime symbols; do not prime the diagram constants or e\_A. \\
\phantomsection\label{entry:symbol-registry-03491}\(I_Q\) & \texttt{\detokenize{I_Q}} & relation inclusion theory & P3 Lyndon 1959 preservation & n/a & \hyperref[def:p03-lyndon-polarity]{\ref*{def:p03-lyndon-polarity}} & \thepage & One universally closed inclusion per selected relation symbol. \\
\phantomsection\label{entry:symbol-registry-03492}\(A,B,C,G\) & \texttt{\detokenize{A;B;C;G}} & local structures with domain understood from context & P3 Lyndon 1959 preservation & n/a & \hyperref[def:p03-lyndon-qmap]{\ref*{def:p03-lyndon-qmap}} & \thepage & Domain notation and indexed interpretations are declared together; these letters do not denote axioms. \\
\phantomsection\label{entry:symbol-registry-03493}\(h\) & \texttt{\detokenize{h}} & onto Q-map or onto homomorphism & P3 Lyndon 1959 preservation & n/a & \hyperref[def:p03-lyndon-qmap]{\ref*{def:p03-lyndon-qmap}} & \thepage & Preserves total operations and relations; reflects relations outside Q. \\
\phantomsection\label{entry:symbol-registry-03494}\(\Gamma^*,K^*\) & \texttt{\detokenize{\Gamma^*;K^*}} & model class and common sentence theory & P3 Lyndon 1959 preservation & n/a & \hyperref[def:p03-lyndon-qmap]{\ref*{def:p03-lyndon-qmap}} & \thepage & The input type determines the dual role; double star is logical or elementary closure respectively. \\
\phantomsection\label{entry:symbol-registry-03495}\(k\) & \texttt{\detokenize{k}} & map kernel or specified congruence & P3 Lyndon 1959 preservation & relation & \hyperref[thm:p03-lyndon-map-calculus]{\ref*{thm:p03-lyndon-map-calculus}} & \thepage & For a map h kernel means h(a)=h(a-prime); all auxiliary relation congruence conditions are needed for a zero-map quotient. \\
\phantomsection\label{entry:symbol-registry-03496}\(L_A,A_A,e_A,w_a\) & \texttt{\detokenize{L_A;A_A;e_A;w_a}} & expanded language named expansion fresh identity relation and constants & P3 Lyndon 1959 preservation & n/a & \hyperref[def:p03-lyndon-diagram]{\ref*{def:p03-lyndon-diagram}} & \thepage & e\_A is a new ordinary relation; its congruence quotient extracts the elementary extension; w\_a is the source named constant. \\
\phantomsection\label{entry:symbol-registry-03497}\(\bar C,\bar B,\bar h\) & \texttt{\detokenize{\bar C;\bar B;\bar h}} & identity quotients and induced map & P3 Lyndon 1959 preservation & n/a & \hyperref[thm:p03-lyndon-main]{\ref*{thm:p03-lyndon-main}} & \thepage & Quotient by the congruence interpreting e; genuine-identity projection is a {e}-map rather than an auxiliary empty-set-map. \\
\phantomsection\label{entry:symbol-registry-03498}\(H(K)\) & \texttt{\detokenize{H(K)}} & class of onto homomorphic images of a stated class & P3 Lyndon 1959 preservation & n/a & \hyperref[thm:p03-lyndon-image-class-boundary]{\ref*{thm:p03-lyndon-image-class-boundary}} & \thepage & Distinct from H the ambient relative elementary class. \\
\phantomsection\label{entry:symbol-registry-03499}\(H\) & \texttt{\detokenize{H}} & ambient elementary class in relative preservation & P3 Lyndon 1959 preservation & n/a & \hyperref[thm:p03-lyndon-relative]{\ref*{thm:p03-lyndon-relative}} & \thepage & Both endpoints of the relative maps must belong to H; H need not be image-closed. \\
\phantomsection\label{entry:symbol-registry-03500}\(S,S_n\) & \texttt{\detokenize{S;S_n}} & source rooted infinite-extension condition and finite chain conditions & P3 Lyndon 1959 preservation & n/a & \hyperref[def:p03-lyndon-chain-example]{\ref*{def:p03-lyndon-chain-example}} & \thepage & Initial edge is conjoined outside the implication in S; witnesses of S\_n need not be distinct. \\
\phantomsection\label{entry:symbol-registry-03501}\(s_0,s_1,s,t_i\) & \texttt{\detokenize{s_0;s_1;s;t_i}} & Skolem constants unary operation and iterated chain terms & P3 Lyndon 1959 preservation & n/a & \hyperref[thm:p03-lyndon-constructive-import]{\ref*{thm:p03-lyndon-constructive-import}} & \thepage & t\_0=s\_0 and t\_1=s\_1; t\_(i+1)=s(t\_i) only for i at least one. \\
\phantomsection\label{entry:symbol-registry-03502}\(\sigma,\pi_i\) & \texttt{\detokenize{\sigma;\pi_i}} & Skolem variable substitution and variable-disjoint copy renamings & P3 Lyndon 1959 preservation constructive import & n/a & \hyperref[def:p03-lyndon-constructive-operations]{\ref*{def:p03-lyndon-constructive-operations}} & \thepage & Skolem witness operations depend only on preceding universals; duplication prefix respects each original dependency. \\
\phantomsection\label{entry:symbol-registry-03503}\(x_{pq},y_r,u_h,s_r\) & \texttt{\detokenize{x_{pq};y_r;u_h;s_r}} & block universal variables existential witnesses fresh universals and Skolem operations & P3 Lyndon 1959 preservation constructive import & n/a & \hyperref[def:p03-lyndon-constructive-operations]{\ref*{def:p03-lyndon-constructive-operations}} & \thepage & Specialization terms for x\_(pq) may use only y\_s with s less than p; prefix dependency requirements are explicit. \\
\phantomsection\label{entry:symbol-registry-03504}\(\chi\) & \texttt{\detokenize{\chi}} & Boolean substitution on atomic formulas in source chain argument & P3 Lyndon 1959 preservation & n/a & \hyperref[thm:p03-lyndon-image-class-boundary]{\ref*{thm:p03-lyndon-image-class-boundary}} & \thepage & Keeps rooted-chain atoms and sends every other atom to 0; not a variable substitution. \\
\phantomsection\label{entry:symbol-registry-03505}\(\Sigma\) & \texttt{\detokenize{\Sigma}} & set of substitution instances of the Skolem matrix & P3 Lyndon 1959 preservation chain argument & n/a & \hyperref[thm:p03-lyndon-image-class-boundary]{\ref*{thm:p03-lyndon-image-class-boundary}} & \thepage & Formulas not sentences; source Boolean substitution and finite propositional compactness apply to this set. \\
\phantomsection\label{entry:symbol-registry-03506}\(\Sigma\) & \texttt{\detokenize{\Sigma}} & sentences true in all one-element models of T & P3 Lyndon 1959 preservation undecidability & n/a & \hyperref[thm:p03-lyndon-undecidable]{\ref*{thm:p03-lyndon-undecidable}} & \thepage & Recursive membership uses the finite axiomatization; distinct local use from source Skolem instance set in the chain annex. \\
\phantomsection\label{entry:symbol-registry-03708}\(L\) & \texttt{\detokenize{L}} & at most countable finite-arity first-order signature with identity & P3 Morley countable model dependencies & n/a & \hyperref[def:p03-morley-coding]{\ref*{def:p03-morley-coding}} & \thepage & Functions and constants are included in the coded interpretation data. \\
\phantomsection\label{entry:symbol-registry-03709}\(X_L\) & \texttt{\detokenize{X_L}} & space of structures labeled on omega & P3 Morley countable model dependencies & n/a & \hyperref[def:p03-morley-coding]{\ref*{def:p03-morley-coding}} & \thepage & Every code has countably infinite domain; finite components are kept separate. \\
\phantomsection\label{entry:symbol-registry-03710}\(X_{L,m}\) & \texttt{\detokenize{X_{L,m}}} & space of structures labeled on a fixed positive finite domain & P3 Morley countable model dependencies & n/a & \hyperref[def:p03-morley-coding]{\ref*{def:p03-morley-coding}} & \thepage & Even fixed finite domains can have continuum many isomorphism types for infinite signatures. \\
\phantomsection\label{entry:symbol-registry-03711}\(X_L^{\leq\omega}\) & \texttt{\detokenize{X_L^{\leq\omega}}} & countable topological sum of nonempty finite and infinite code components & P3 Morley countable model dependencies & n/a & \hyperref[def:p03-morley-coding]{\ref*{def:p03-morley-coding}} & \thepage & At most countable structures include the explicitly separate finite components. \\
\phantomsection\label{entry:symbol-registry-03712}\(\operatorname{arity}\) & \texttt{\detokenize{\operatorname{arity}}} & finite arity of a relation or positive-arity function symbol & P3 Morley countable model dependencies & n/a & \hyperref[def:p03-morley-coding]{\ref*{def:p03-morley-coding}} & \thepage & Constants and nullary relation coordinates have their explicit clauses. \\
\phantomsection\label{entry:symbol-registry-03713}\(I(K)\) & \texttt{\detokenize{I(K)}} & cardinality of a coded class modulo isomorphism & P3 Morley countable model dependencies & n/a & \hyperref[def:p03-morley-coding]{\ref*{def:p03-morley-coding}} & \thepage & This counts classes rather than labeled copies. \\
\phantomsection\label{entry:symbol-registry-03714}\(K_T,K_\sigma\) & \texttt{\detokenize{K_T;K_\sigma}} & fixed-omega-domain model classes & P3 Morley countable model dependencies & n/a & \hyperref[def:p03-morley-coding]{\ref*{def:p03-morley-coding}} & \thepage & No claim is made that an uncountable sentence family is a permitted countable conjunction. \\
\phantomsection\label{entry:symbol-registry-03715}\(L_{\omega_1,\omega}\) & \texttt{\detokenize{L_{\omega_1,\omega}}} & logic with countable connectives and finite quantifier strings & P3 Morley countable model dependencies & n/a & \hyperref[sec:p03-morley-codes]{\ref*{sec:p03-morley-codes}} & \thepage & Defined fully at def:p03-morley-infinitary before its truth and construction results. \\
\phantomsection\label{entry:symbol-registry-03716}\(\operatorname{qr}\) & \texttt{\detokenize{\operatorname{qr}}} & auxiliary countable infinitary quantifier rank & P3 Morley countable model dependencies & n/a & \hyperref[def:p03-morley-infinitary]{\ref*{def:p03-morley-infinitary}} & \thepage & This rank is not identified with the selected status source's Scott-rank conventions. \\
\phantomsection\label{entry:symbol-registry-03717}\(\equiv_\alpha\) & \texttt{\detokenize{\equiv_\alpha}} & pointed back-and-forth equivalence at a countable ordinal & P3 Morley countable model dependencies & relation & \hyperref[def:p03-morley-approximants]{\ref*{def:p03-morley-approximants}} & \thepage & Tuples have equal finite arity and may repeat coordinates. \\
\phantomsection\label{entry:symbol-registry-03718}\(E_\alpha(M,N)\) & \texttt{\detokenize{E_\alpha}} & empty-tuple back-and-forth equivalence & P3 Morley countable model dependencies & relation & \hyperref[def:p03-morley-approximants]{\ref*{def:p03-morley-approximants}} & \thepage & Each fixed level is Borel on omega-domain code pairs; the full ordinal intersection is uncountable. \\
\phantomsection\label{entry:symbol-registry-03719}\(\alpha,\beta,\lambda\) & \texttt{\detokenize{\alpha;\beta;\lambda}} & countable ordinals used in the recursion & P3 Morley countable model dependencies & n/a & \hyperref[def:p03-morley-approximants]{\ref*{def:p03-morley-approximants}} & \thepage & The limit clause uses only countably many preceding stages. \\
\phantomsection\label{entry:symbol-registry-03720}\(\bar a,\bar b\) & \texttt{\detokenize{\bar a;\bar b}} & equal-length interpreted finite tuples & P3 Morley countable model dependencies & n/a & \hyperref[def:p03-morley-approximants]{\ref*{def:p03-morley-approximants}} & \thepage & Their equality atoms make the underlying finite partial bijection well defined. \\
\phantomsection\label{entry:symbol-registry-03721}\(\theta_{\bar a}^{\alpha}\) & \texttt{\detokenize{\theta}} & infinitary formula attached to a tuple in one fixed countable source structure & P3 Morley countable model dependencies & n/a & \hyperref[def:p03-morley-scott-formulas]{\ref*{def:p03-morley-scott-formulas}} & \thepage & Its two successor clauses encode every forth and every back extension. \\
\phantomsection\label{entry:symbol-registry-03722}\(\delta_M\) & \texttt{\detokenize{\delta_M}} & countable simultaneous stabilization stage within one source structure & P3 Morley countable model dependencies & n/a & \hyperref[thm:p03-morley-scott-stabilization]{\ref*{thm:p03-morley-scott-stabilization}} & \thepage & The stage can depend on M and is not a uniform bound over all models. \\
\phantomsection\label{entry:symbol-registry-03723}\(\sigma_M\) & \texttt{\detokenize{\sigma_M}} & Scott sentence with explicit universal stability clauses & P3 Morley countable model dependencies & n/a & \hyperref[thm:p03-morley-scott-sentence]{\ref*{thm:p03-morley-scott-sentence}} & \thepage & Characterizes the source among nonempty at most countable models only. \\
\phantomsection\label{entry:symbol-registry-03724}\(S_\infty\) & \texttt{\detokenize{S_\infty}} & Polish group of permutations of omega & P3 Morley countable model dependencies & n/a & \hyperref[thm:p03-morley-isomorphism-complexity]{\ref*{thm:p03-morley-isomorphism-complexity}} & \thepage & Complete metric uses both forward and inverse coordinates. \\
\phantomsection\label{entry:symbol-registry-03725}\(\pi,\tau\) & \texttt{\detokenize{\pi;\tau}} & permutations used in the isomorphism witness space & P3 Morley countable model dependencies & n/a & \hyperref[annex:p3.morley.isomorphism-complexity]{\ref*{annex:p3.morley.isomorphism-complexity}} & \thepage & Their forward and inverse limits are checked to be mutual inverses. \\
\phantomsection\label{entry:symbol-registry-03726}\(d(u,v)\) & \texttt{\detokenize{d}} & complete product coordinate metric & P3 Morley countable model dependencies & n/a & \hyperref[annex:p3.morley.codes]{\ref*{annex:p3.morley.codes}} & \thepage & Coordinate stabilization and tail estimates prove completeness. \\
\phantomsection\label{entry:symbol-registry-03727}\(\mathbf1\) & \texttt{\detokenize{\mathbf1}} & indicator of a stated coordinate inequality & P3 Morley countable model dependencies & n/a & \hyperref[annex:p3.morley.codes]{\ref*{annex:p3.morley.codes}} & \thepage & The second metric later applies the same indicator to forward and inverse permutation coordinates. \\
\phantomsection\label{entry:symbol-registry-03728}\(C\) & \texttt{\detokenize{C}} & closed graph of permutation isomorphism witnesses & P3 Morley countable model dependencies & n/a & \hyperref[annex:p3.morley.isomorphism-complexity]{\ref*{annex:p3.morley.isomorphism-complexity}} & \thepage & Its projection is analytic; projection closedness or Borelness is not assumed. \\
\phantomsection\label{entry:symbol-registry-03729}\(U\) & \texttt{\detokenize{U}} & incomplete theory of uniform unary predicates & P3 Morley countable model dependencies & n/a & \hyperref[thm:p03-morley-perfect-example]{\ref*{thm:p03-morley-perfect-example}} & \thepage & The binary model family is not a many-model construction for a complete theory. \\
\phantomsection\label{entry:symbol-registry-03730}\(\eta\) & \texttt{\detokenize{\eta}} & infinite binary sequence indexing the perfect family & P3 Morley countable model dependencies & n/a & \hyperref[thm:p03-morley-perfect-example]{\ref*{thm:p03-morley-perfect-example}} & \thepage & Each bit specifies whether its predicate is everywhere or nowhere. \\
\phantomsection\label{entry:symbol-registry-03731}\(M_\eta\) & \texttt{\detokenize{M_\eta}} & omega-domain unary model associated with a branch & P3 Morley countable model dependencies & n/a & \hyperref[thm:p03-morley-perfect-example]{\ref*{thm:p03-morley-perfect-example}} & \thepage & Different branches are separated by an isomorphism-invariant predicate sentence. \\
\phantomsection\label{entry:symbol-registry-03732}\(P_i\) & \texttt{\detokenize{P_i}} & indexed unary predicate in the perfect-family example & P3 Morley countable model dependencies & n/a & \hyperref[thm:p03-morley-perfect-example]{\ref*{thm:p03-morley-perfect-example}} & \thepage & This is the example's countably infinite language rather than an unnamed parameter. \\
\phantomsection\label{entry:symbol-registry-03733}\(D\) & \texttt{\detokenize{D}} & strongly minimal definable set & P3 Baldwin Lachlan geometry & n/a & \hyperref[def:p03-bl-ambient]{\ref*{def:p03-bl-ambient}} & \thepage & Defined by delta over the finite tuple c; may be a proper subset of the model. \\
\phantomsection\label{entry:symbol-registry-03734}\(\bar c\) & \texttt{\detokenize{\bar c}} & finite tuple defining the strongly minimal set & P3 Baldwin Lachlan geometry & n/a & \hyperref[def:p03-bl-ambient]{\ref*{def:p03-bl-ambient}} & \thepage & Its coordinates are retained in every local closure base. \\
\phantomsection\label{entry:symbol-registry-03735}\(\operatorname{cl}_{\bar c}(X)\) & \texttt{\detokenize{\operatorname{cl}_{\bar c}}} & algebraic closure restricted to D with c in the base & P3 Baldwin Lachlan geometry & n/a & \hyperref[def:p03-bl-ambient]{\ref*{def:p03-bl-ambient}} & \thepage & Not a closure operation on the whole model unless D is its universe. \\
\phantomsection\label{entry:symbol-registry-03736}\(\dim_{\bar c}(D)\) & \texttt{\detokenize{\dim_{\bar c}}} & cardinality of a pregeometry basis & P3 Baldwin Lachlan geometry & n/a & \hyperref[thm:p03-bl-dimension]{\ref*{thm:p03-bl-dimension}} & \thepage & Distinct from ordinal Morley rank and finite Morley degree. \\
\phantomsection\label{entry:symbol-registry-03737}\(q\) & \texttt{\detokenize{q}} & uniform finite or finite complement bound & P3 Baldwin Lachlan geometry & n/a & \hyperref[thm:p03-bl-uniform]{\ref*{thm:p03-bl-uniform}} & \thepage & Chosen separately for each formula family. \\
\phantomsection\label{entry:symbol-registry-03738}\(R(x,y)\) & \texttt{\detokenize{R}} & definable bounded finite fibre relation & \hyperref[annex:p3.n-a-006.2.pregeometry]{\ref*{annex:p3.n-a-006.2.pregeometry}} & n/a & \hyperref[annex:p3.n-a-006.2.pregeometry]{\ref*{annex:p3.n-a-006.2.pregeometry}} & \thepage & Not a proof relation or a rank. \\
\phantomsection\label{entry:symbol-registry-03739}\(H\) & \texttt{\detokenize{H}} & definable set of points with cofinite reverse fibre & \hyperref[annex:p3.n-a-006.2.pregeometry]{\ref*{annex:p3.n-a-006.2.pregeometry}} & n/a & \hyperref[annex:p3.n-a-006.2.pregeometry]{\ref*{annex:p3.n-a-006.2.pregeometry}} & \thepage & Its definability uses the proved uniform bound. \\
\phantomsection\label{entry:symbol-registry-03740}\(J_i\) & \texttt{\detokenize{J_i}} & finite spanning set maintained in basis replacement & \hyperref[annex:p3.n-a-006.2.dimension]{\ref*{annex:p3.n-a-006.2.dimension}} & n/a & \hyperref[annex:p3.n-a-006.2.dimension]{\ref*{annex:p3.n-a-006.2.dimension}} & \thepage & Contains the first i independent elements and retains the original size. \\
\phantomsection\label{entry:symbol-registry-03741}\(n_0\) & \texttt{\detokenize{n_0}} & least finite dimension of an infinite model when such a dimension exists & P3 whole strongly minimal theory & n/a & \hyperref[thm:p03-bl-whole-models]{\ref*{thm:p03-bl-whole-models}} & \thepage & Undefined and unused in the alternative with no finite dimensional model. \\
\phantomsection\label{entry:symbol-registry-05427}\(L\) & \texttt{\detokenize{L}} & source first-order language & P3 quotient interpretations & n/a & \hyperref[def:p03-imaginaries-standing]{\ref*{def:p03-imaginaries-standing}} & \thepage & Set-sized and possibly many-sorted; named parameters are declared before uniform claims. \\
\phantomsection\label{entry:symbol-registry-05428}\(K\) & \texttt{\detokenize{K}} & target first-order language & P3 quotient interpretations & n/a & \hyperref[def:p03-imaginaries-standing]{\ref*{def:p03-imaginaries-standing}} & \thepage & Not a field here; the field example uses R and F locally. \\
\phantomsection\label{entry:symbol-registry-05429}\(T\) & \texttt{\detokenize{T}} & source theory & P3 quotient interpretations & n/a & \hyperref[def:p03-imaginaries-standing]{\ref*{def:p03-imaginaries-standing}} & \thepage & Uniform consequences are required in every model of this declared language. \\
\phantomsection\label{entry:symbol-registry-05430}\(n_s\) & \texttt{\detokenize{n_s}} & finite dimension of a representative block & P3 quotient interpretations & n/a & \hyperref[def:p03-imaginaries-standing]{\ref*{def:p03-imaginaries-standing}} & \thepage & May be zero; it is not the cardinality of the quotient sort. \\
\phantomsection\label{entry:symbol-registry-05431}\(D_s\) & \texttt{\detokenize{D_s}} & definable representative domain & P3 quotient interpretations & n/a & \hyperref[def:p03-imaginaries-interpretation]{\ref*{def:p03-imaginaries-interpretation}} & \thepage & Subset of a fixed finite product of typed home sorts. \\
\phantomsection\label{entry:symbol-registry-05432}\(E_s\) & \texttt{\detokenize{E_s}} & definable equality on representatives & P3 quotient interpretations & n/a & \hyperref[def:p03-imaginaries-interpretation]{\ref*{def:p03-imaginaries-interpretation}} & \thepage & Equivalence on D\_s; not actual tuple equality in general. \\
\phantomsection\label{entry:symbol-registry-05433}\(I(M)_s\) & \texttt{\detokenize{I(M)_s}} & interpreted quotient sort & P3 quotient interpretations & n/a & \hyperref[def:p03-imaginaries-interpretation]{\ref*{def:p03-imaginaries-interpretation}} & \thepage & Sort tags distinguish coincident underlying sets. \\
\phantomsection\label{entry:symbol-registry-05434}\(q_s\) & \texttt{\detokenize{q_s}} & set-theoretic quotient projection & P3 quotient interpretations & n/a & \hyperref[def:p03-imaginaries-interpretation]{\ref*{def:p03-imaginaries-interpretation}} & \thepage & Domain is D\_s; no representative section or partial source-language term is introduced. \\
\phantomsection\label{entry:symbol-registry-05435}\(G_f\) & \texttt{\detokenize{G_f}} & source formula defining a quotient function graph & P3 quotient interpretations & n/a & \hyperref[def:p03-imaginaries-interpretation]{\ref*{def:p03-imaginaries-interpretation}} & \thepage & Total and unique modulo the output equivalence; invariant in every block. \\
\phantomsection\label{entry:symbol-registry-05436}\(V_t^I\) & \texttt{\detokenize{V_t^I}} & source formula for representative term values & P3 quotient interpretations & n/a & \hyperref[def:p03-imaginaries-translation]{\ref*{def:p03-imaginaries-translation}} & \thepage & Term recursion includes functions and constants; intermediate output blocks are guarded. \\
\phantomsection\label{entry:symbol-registry-05437}\(\varphi^I\) & \texttt{\detokenize{\varphi^I}} & formula pullback through an interpretation & P3 quotient interpretations & n/a & \hyperref[def:p03-imaginaries-translation]{\ref*{def:p03-imaginaries-translation}} & \thepage & Equality is translated through E\_s; quantifiers use the representative domain guards. \\
\phantomsection\label{entry:symbol-registry-05438}\(h^I\) & \texttt{\detokenize{h^I}} & induced elementary quotient map & P3 quotient interpretations & n/a & \hyperref[thm:p03-imaginaries-elementarity]{\ref*{thm:p03-imaginaries-elementarity}} & \thepage & Defined on classes of coordinatewise images; not literal inclusion of successive classes. \\
\phantomsection\label{entry:symbol-registry-05439}\(I\circ J\) & \texttt{\detokenize{I\circ J}} & composition by flattened quotient data & P3 quotient interpretations & n/a & \hyperref[thm:p03-imaginaries-composition]{\ref*{thm:p03-imaginaries-composition}} & \thepage & Interpret J first and then I; comparison is with I(J(M)). \\
\phantomsection\label{entry:symbol-registry-05440}\(S_E\) & \texttt{\detokenize{S_E}} & tagged imaginary sort for an equivalence & P3 imaginary expansion & n/a & \hyperref[def:p03-imaginaries-eq-expansion]{\ref*{def:p03-imaginaries-eq-expansion}} & \thepage & The full representative product is the domain of its total projection. \\
\phantomsection\label{entry:symbol-registry-05441}\(M^{\mathrm{eq}}\) & \texttt{\detokenize{M^{\mathrm{eq}}}} & many-sorted imaginary expansion & P3 imaginary expansion & n/a & \hyperref[def:p03-imaginaries-eq-expansion]{\ref*{def:p03-imaginaries-eq-expansion}} & \thepage & Home sorts remain; quotient closure in expanded sorts is not elimination in original real sorts. \\
\phantomsection\label{entry:symbol-registry-05442}\(E^*\) & \texttt{\detokenize{E^*}} & whole-product extension of a restricted equivalence & P3 imaginary expansion & n/a & \hyperref[thm:p03-imaginaries-eq-pullback]{\ref*{thm:p03-imaginaries-eq-pullback}} & \thepage & Old classes within D; singleton classes outside D; no cross-region classes. \\
\phantomsection\label{entry:symbol-registry-05443}\(E_\varphi\) & \texttt{\detokenize{E_\varphi}} & equality of definable-family fibers & P3 canonical parameters & n/a & \hyperref[def:p03-imaginaries-canonical-parameter]{\ref*{def:p03-imaginaries-canonical-parameter}} & \thepage & Defined by universally quantified equality of the sets given by the same formula. \\
\phantomsection\label{entry:symbol-registry-05444}\(e_X\) & \texttt{\detokenize{e_X}} & imaginary canonical parameter of a family member & P3 canonical parameters & n/a & \hyperref[def:p03-imaginaries-canonical-parameter]{\ref*{def:p03-imaginaries-canonical-parameter}} & \thepage & Its sort and defining family are declared; it is not an unspecified real code. \\
\phantomsection\label{entry:symbol-registry-05445}\(\operatorname{dcl}^{\mathrm{eq}}\) & \texttt{\detokenize{\operatorname{dcl}^{\mathrm{eq}}}} & unique-solution definable closure & P3 imaginary coding & n/a & \hyperref[def:p03-imaginaries-real-codes]{\ref*{def:p03-imaginaries-real-codes}} & \thepage & Finite parameters and typed formulas; no nonsaturated automorphism converse is assumed. \\
\phantomsection\label{entry:symbol-registry-05446}\(\operatorname{acl}^{\mathrm{eq}}\) & \texttt{\detokenize{\operatorname{acl}^{\mathrm{eq}}}} & finite-solution algebraic closure & P3 imaginary coding & n/a & \hyperref[def:p03-imaginaries-real-codes]{\ref*{def:p03-imaginaries-real-codes}} & \thepage & Weak real codes may have finite ambiguity; every coordinate is algebraic. \\
\phantomsection\label{entry:symbol-registry-05447}\(\chi_E\) & \texttt{\detokenize{\chi_E}} & uniform real coding graph & P3 imaginary coding & n/a & \hyperref[def:p03-imaginaries-uniform-ei]{\ref*{def:p03-imaginaries-uniform-ei}} & \thepage & One formula for the equivalence presentation independent of model or class; zero output coordinates permitted. \\
\phantomsection\label{entry:symbol-registry-05448}\(s\) & \texttt{\detokenize{s}} & definable class-invariant selector & P3 imaginary coding & n/a & \hyperref[thm:p03-imaginaries-uniform-codes]{\ref*{thm:p03-imaginaries-uniform-codes}} & \thepage & This local function role is separate from sort index s in the interpretation definition. \\
\phantomsection\label{entry:symbol-registry-05449}\(R\) & \texttt{\detokenize{R}} & commutative integral-domain structure & P3 fraction-field quotient & n/a & \hyperref[thm:p03-imaginaries-fraction-field]{\ref*{thm:p03-imaginaries-fraction-field}} & \thepage & Not the equivalence relation R of the preceding counterexample; local scope is the fraction-field theorem. \\
\phantomsection\label{entry:symbol-registry-05450}\([a,b]\) & \texttt{\detokenize{[a,}\allowbreak\detokenize{b]}} & fraction-field equivalence class & P3 fraction-field quotient & n/a & \hyperref[thm:p03-imaginaries-fraction-field]{\ref*{thm:p03-imaginaries-fraction-field}} & \thepage & b is nonzero; cross multiplication defines equality. \\
\phantomsection\label{entry:symbol-registry-05451}\(F\) & \texttt{\detokenize{F}} & ordered field & P3 complexification & n/a & \hyperref[thm:p03-imaginaries-complexification]{\ref*{thm:p03-imaginaries-complexification}} & \thepage & Separate local role from the parameter-fiber equivalence F in the imaginary pullback proof. \\
\phantomsection\label{entry:symbol-registry-05847}\(L\) & \texttt{\detokenize{L}} & finitary many-sorted first-order signature & P3 FV products & n/a & \hyperref[def:p03-fv-standing]{\ref*{def:p03-fv-standing}} & \thepage & Set-sized; every factor sort is nonempty. \\
\phantomsection\label{entry:symbol-registry-05848}\(I\) & \texttt{\detokenize{I}} & set of indices & P3 FV products & n/a & \hyperref[def:p03-fv-standing]{\ref*{def:p03-fv-standing}} & \thepage & May be empty; the product then has one element in each sort. \\
\phantomsection\label{entry:symbol-registry-05849}\(M_i\) & \texttt{\detokenize{M_i}} & indexed factor structure & P3 FV products & n/a & \hyperref[def:p03-fv-product]{\ref*{def:p03-fv-product}} & \thepage & The coordinate domain is sort-specific. \\
\phantomsection\label{entry:symbol-registry-05850}\(\mathcal P(I)\) & \texttt{\detokenize{\mathcal P(I)}} & Boolean algebra of all subsets of the index set & P3 FV products & n/a & \hyperref[def:p03-fv-boolean-language]{\ref*{def:p03-fv-boolean-language}} & \thepage & Complement is relative to I; Boolean quantifiers range over subsets. \\
\phantomsection\label{entry:symbol-registry-05851}\([\![\psi(\bar a)]\!]\) & \texttt{\detokenize{[\![\psi(\bar a)]\!]}} & factor-formula truth set & P3 FV products & n/a & \hyperref[def:p03-fv-product]{\ref*{def:p03-fv-product}} & \thepage & Not the assertion that the product satisfies psi. \\
\phantomsection\label{entry:symbol-registry-05852}\(\Theta_\varphi\) & \texttt{\detokenize{\Theta_\varphi}} & first-order Boolean condition & P3 FV reduction & n/a & \hyperref[def:p03-fv-reduction]{\ref*{def:p03-fv-reduction}} & \thepage & Its logical negation is not Boolean complementation of one truth set. \\
\phantomsection\label{entry:symbol-registry-05853}\(\psi_j\) & \texttt{\detokenize{\psi_j}} & finite list of local formulas & P3 FV reduction & n/a & \hyperref[def:p03-fv-reduction]{\ref*{def:p03-fv-reduction}} & \thepage & Bound variables and Boolean variables are freshened capture-avoidantly. \\
\phantomsection\label{entry:symbol-registry-05854}\(\sigma\) & \texttt{\detokenize{\sigma}} & bit pattern in the finite local-formula list & P3 FV reduction & n/a & \hyperref[def:p03-fv-reduction]{\ref*{def:p03-fv-reduction}} & \thepage & All of {0,1}\^{}m is used including the unique pattern when m=0. \\
\phantomsection\label{entry:symbol-registry-05855}\(\delta_\sigma\) & \texttt{\detokenize{\delta_\sigma}} & complete signed local witness pattern & P3 FV reduction & n/a & \hyperref[def:p03-fv-reduction]{\ref*{def:p03-fv-reduction}} & \thepage & Conjunction of each psi\_j or its logical negation. \\
\phantomsection\label{entry:symbol-registry-05856}\(\eta_\sigma\) & \texttt{\detokenize{\eta_\sigma}} & local existence formula for a witness pattern & P3 FV reduction & n/a & \hyperref[def:p03-fv-reduction]{\ref*{def:p03-fv-reduction}} & \thepage & Existentially binds the declared witness sort. \\
\phantomsection\label{entry:symbol-registry-05857}\(Z_\sigma\) & \texttt{\detokenize{Z_\sigma}} & Boolean witness-partition cell & P3 FV reduction & n/a & \hyperref[def:p03-fv-reduction]{\ref*{def:p03-fv-reduction}} & \thepage & Zero cells are allowed; nonzero cells must lie below the local existence value. \\
\phantomsection\label{entry:symbol-registry-05858}\(\operatorname{Part}_F\) & \texttt{\detokenize{\operatorname{Part}_F}} & finite Boolean partition condition & P3 FV reduction & n/a & \hyperref[def:p03-fv-boolean-language]{\ref*{def:p03-fv-boolean-language}} & \thepage & Pairwise disjoint cells whose join is one; not an infinite completeness assumption. \\
\phantomsection\label{entry:symbol-registry-05859}\(W_i\) & \texttt{\detokenize{W_i}} & nonempty set of actual factor witnesses & P3 FV direct-product proof & n/a & \hyperref[thm:p03-fv-direct-reduction]{\ref*{thm:p03-fv-direct-reduction}} & \thepage & An indexed choice function is required in the full product; nonempty product domains alone do not suffice. \\
\phantomsection\label{entry:symbol-registry-05860}\(B\) & \texttt{\detokenize{B}} & external coefficient Boolean algebra & P3 finite-partition Boolean powers & n/a & \hyperref[def:p03-fv-partition-power]{\ref*{def:p03-fv-partition-power}} & \thepage & May be degenerate; not automatically the Boolean algebra of all idempotents of a ring. \\
\phantomsection\label{entry:symbol-registry-05861}\(A[B]_{\mathrm{part}}\) & \texttt{\detokenize{A[B]_{\mathrm{part}}}} & formal finite-partition Boolean power & P3 finite-partition Boolean powers & n/a & \hyperref[def:p03-fv-partition-power]{\ref*{def:p03-fv-partition-power}} & \thepage & Finite presentations modulo equal-label common refinements; finite witness choice suffices. \\
\phantomsection\label{entry:symbol-registry-05862}\((b_r,a_r)\) & \texttt{\detokenize{(b_r,}\allowbreak\detokenize{a_r)}} & Boolean coefficient and base-sort label & P3 finite-partition Boolean powers & n/a & \hyperref[def:p03-fv-partition-power]{\ref*{def:p03-fv-partition-power}} & \thepage & The coefficients form a finite partition; zero coefficients carry irrelevant labels. \\
\phantomsection\label{entry:symbol-registry-05863}\([\![\psi(\bar a)]\!]_B\) & \texttt{\detokenize{[\![\psi(\bar a)]\!]_B}} & factor-formula Boolean value from a finite refinement & P3 finite-partition Boolean powers & n/a & \hyperref[thm:p03-fv-power-structure]{\ref*{thm:p03-fv-power-structure}} & \thepage & Not the truth value of psi evaluated internally in the power. \\
\phantomsection\label{entry:symbol-registry-05864}\(X_B\) & \texttt{\detokenize{X_B}} & Stone space of the coefficient Boolean algebra & P3 Stone realization & n/a & \hyperref[thm:p03-fv-power-representations]{\ref*{thm:p03-fv-power-representations}} & \thepage & Uses the existing Stone representation theorem and its BPIT hypothesis; empty for degenerate B. \\
\phantomsection\label{entry:symbol-registry-05865}\(A[B]^*\) & \texttt{\detokenize{A[B]^*}} & source bounded continuous Boolean power & P3 source presentation IV section 5 & n/a & \hyperref[thm:p03-fv-power-representations]{\ref*{thm:p03-fv-power-representations}} & \thepage & Burris--Sankappanavar notation; maps from the compact Stone space to discrete A. \\
\phantomsection\label{entry:symbol-registry-05866}\(A^I\) & \texttt{\detokenize{A^I}} & full direct power & P3 elementary-preservation comparison & n/a & \hyperref[thm:p03-fv-power-elementarity]{\ref*{thm:p03-fv-power-elementarity}} & \thepage & Distinguished from its finite-range Boolean subpower and from the constant diagonal. \\
\phantomsection\label{entry:symbol-registry-05867}\(S\) & \texttt{\detokenize{S}} & subalgebra of a represented full function power & P3 intrinsic power criterion & n/a & \hyperref[thm:p03-fv-patching-characterization]{\ref*{thm:p03-fv-patching-characterization}} & \thepage & Constants plus clopen equalizers plus binary clopen patching characterize the bounded power. \\
\phantomsection\label{entry:symbol-registry-06105}\(L\) & \texttt{\detokenize{L}} & recursive base first-order language & P3 resplendence effective construction & n/a & \hyperref[def:p03-resplendence-coding]{\ref*{def:p03-resplendence-coding}} & \thepage & Recursive syntax and finite arities; no computable model diagram is assumed. \\
\phantomsection\label{entry:symbol-registry-06106}\(L^+\) & \texttt{\detokenize{L^+}} & recursive same-sort expansion language & P3 resplendence effective construction & n/a & \hyperref[def:p03-resplendence-coding]{\ref*{def:p03-resplendence-coding}} & \thepage & May have countably many new symbols; finite-sentence variants are separately named. \\
\phantomsection\label{entry:symbol-registry-06107}\(M\) & \texttt{\detokenize{M}} & fixed underlying model and domain & P3 resplendence effective construction & n/a & \hyperref[def:p03-resplendence-recursive-saturation]{\ref*{def:p03-resplendence-recursive-saturation}} & \thepage & Actual witnesses must be selected from this model; an elementary extension is not the desired fixed-domain expansion. \\
\phantomsection\label{entry:symbol-registry-06108}\(M^+\) & \texttt{\detokenize{M^+}} & expansion on exactly the original domain & P3 resplendence effective construction & n/a & \hyperref[def:p03-resplendence-expansion]{\ref*{def:p03-resplendence-expansion}} & \thepage & Old symbols and the designated finite parameter tuple are unchanged. \\
\phantomsection\label{entry:symbol-registry-06109}\(\bar a\) & \texttt{\detokenize{\bar a}} & finite named parameter tuple & P3 resplendence and definability & n/a & \hyperref[def:p03-resplendence-recursive-saturation]{\ref*{def:p03-resplendence-recursive-saturation}} & \thepage & External finite syntactic names do not make the domain computably numbered. \\
\phantomsection\label{entry:symbol-registry-06110}\(p(\bar x,\bar a)\) & \texttt{\detokenize{p(\bar x,}\allowbreak\detokenize{\bar a)}} & recursive or computably enumerable finite-block type & P3 resplendence effective construction & n/a & \hyperref[def:p03-resplendence-recursive-saturation]{\ref*{def:p03-resplendence-recursive-saturation}} & \thepage & Not necessarily complete; coding before parameter interpretation is explicit. \\
\phantomsection\label{entry:symbol-registry-06111}\(\tau_s\) & \texttt{\detokenize{\tau_s}} & canonical tautological padding tag & P3 c.e.-to-recursive proof & n/a & \hyperref[thm:p03-resplendence-padding]{\ref*{thm:p03-resplendence-padding}} & \thepage & Syntactic length encodes a finite enumeration stage; the tag is not simplified. \\
\phantomsection\label{entry:symbol-registry-06112}\(D\) & \texttt{\detokenize{D}} & full elementary diagram of the fixed model & P3 fixed-domain construction & n/a & \hyperref[thm:p03-resplendence-countable]{\ref*{thm:p03-resplendence-countable}} & \thepage & Not assumed recursive or c.e.; it is not included in the proof enumeration for the projection type. \\
\phantomsection\label{entry:symbol-registry-06113}\(S_i\) & \texttt{\detokenize{S_i}} & stage expansion requirement theory & P3 simultaneous Henkin/type schedule & n/a & \hyperref[thm:p03-resplendence-countable]{\ref*{thm:p03-resplendence-countable}} & \thepage & C.e. axioms and finite parameter footprint at each stage; the axiom set need not be finite. \\
\phantomsection\label{entry:symbol-registry-06114}\(\pi\) & \texttt{\detokenize{\pi}} & equality pattern relative to old parameters & P3 actual-witness lemma & n/a & \hyperref[thm:p03-resplendence-witness]{\ref*{thm:p03-resplendence-witness}} & \thepage & Separates coordinates fixed to old parameters from fresh distinct coordinate classes. \\
\phantomsection\label{entry:symbol-registry-06115}\(q(\bar y,\bar a)\) & \texttt{\detokenize{q(\bar y,}\allowbreak\detokenize{\bar a)}} & base-language consequence type & P3 actual-witness lemma & n/a & \hyperref[thm:p03-resplendence-witness]{\ref*{thm:p03-resplendence-witness}} & \thepage & C.e. consequences of effective expansion requirements; no complete-type or diagram decision oracle. \\
\phantomsection\label{entry:symbol-registry-06116}\(q(x;\bar a,\bar b,c)\) & \texttt{\detokenize{q(x;\bar a,}\allowbreak\detokenize{\bar b,}\allowbreak\detokenize{c)}} & recursive comparison type & P3 homogeneity proof & n/a & \hyperref[thm:p03-resplendence-homogeneity]{\ref*{thm:p03-resplendence-homogeneity}} & \thepage & Separate proof-local q role; implications compare actual parameter tuples rather than enumerating their true type. \\
\phantomsection\label{entry:symbol-registry-06117}\(h\) & \texttt{\detokenize{h}} & elementary base embedding used for transport & P3 fixed-domain Henkin proof & n/a & \hyperref[thm:p03-resplendence-countable]{\ref*{thm:p03-resplendence-countable}} & \thepage & The named image is proved an elementary expanded substructure before transporting the new interpretations. \\
\phantomsection\label{entry:symbol-registry-06118}\(R\) & \texttt{\detokenize{R}} & given relation whose definability is tested & P3 parameter-sensitive definability & n/a & \hyperref[thm:p03-resplendence-definability]{\ref*{thm:p03-resplendence-definability}} & \thepage & Recursive saturation of the given expansion (M,R) is an explicit extra hypothesis. \\
\phantomsection\label{entry:symbol-registry-06119}\(\theta_\sigma\) & \texttt{\detokenize{\theta_\sigma}} & finite signed base-formula cell & P3 parameter-sensitive definability & n/a & \hyperref[thm:p03-resplendence-definability]{\ref*{thm:p03-resplendence-definability}} & \thepage & The finite union of membership-positive cells is the actual defining formula. \\
\phantomsection\label{entry:symbol-registry-06120}\(A\) & \texttt{\detokenize{A}} & structure tested for atomic compactness & P3 atomic systems and retractions & n/a & \hyperref[def:p03-resplendence-atomic]{\ref*{def:p03-resplendence-atomic}} & \thepage & No recursive-language requirement in this scope. \\
\phantomsection\label{entry:symbol-registry-06121}\(X\) & \texttt{\detokenize{X}} & arbitrary set of unknown variables & P3 atomic systems and retractions & n/a & \hyperref[def:p03-resplendence-atomic]{\ref*{def:p03-resplendence-atomic}} & \thepage & Not necessarily finite or countable; each individual formula is finite. \\
\phantomsection\label{entry:symbol-registry-06122}\(\Phi(X)\) & \texttt{\detokenize{\Phi(X)}} & system of formulas with arbitrary parameters from A & P3 atomic systems and retractions & n/a & \hyperref[def:p03-resplendence-atomic]{\ref*{def:p03-resplendence-atomic}} & \thepage & Atomic and positive systems have distinct formula classes but the same full parameter/unknown scope. \\
\phantomsection\label{entry:symbol-registry-06123}\(r:B\to A\) & \texttt{\detokenize{r}} & retraction homomorphism fixing A & P3 atomic systems and retractions & n/a & \hyperref[def:p03-resplendence-purity]{\ref*{def:p03-resplendence-purity}} & \thepage & Onto but not asserted injective or elementary; universal-positive preservation uses surjectivity. \\
\phantomsection\label{entry:symbol-registry-06124}\(x_b\) & \texttt{\detokenize{x_b}} & unknown indexed by an extension element & P3 homomorphism-diagram construction & n/a & \hyperref[thm:p03-resplendence-atomic-retractions]{\ref*{thm:p03-resplendence-atomic-retractions}} & \thepage & Every function/relation constraint and every native x\_a equiv a fixed-parameter clause is included. \\
\phantomsection\label{entry:symbol-registry-06125}\(I\) & \texttt{\detokenize{I}} & nonempty reduced-power index set & P3 filter quotient construction & n/a & \hyperref[def:p03-resplendence-reduced-power]{\ref*{def:p03-resplendence-reduced-power}} & \thepage & In the converse I is the set of all finite subsystems including the empty one. \\
\phantomsection\label{entry:symbol-registry-06126}\(F\) & \texttt{\detokenize{F}} & proper filter of index subsets & P3 filter quotient construction & n/a & \hyperref[def:p03-resplendence-reduced-power]{\ref*{def:p03-resplendence-reduced-power}} & \thepage & No full first-order Los statement is claimed for a general proper filter. \\
\phantomsection\label{entry:symbol-registry-06127}\(u\sim_F v\) & \texttt{\detokenize{\sim_F}} & equality on a filter-large set & P3 filter quotient construction & n/a & \hyperref[def:p03-resplendence-reduced-power]{\ref*{def:p03-resplendence-reduced-power}} & \thepage & Defines quotient identity and representative independence. \\
\phantomsection\label{entry:symbol-registry-06128}\(C_v\) & \texttt{\detokenize{C_v}} & cone of finite subsystems extending v & P3 atomic reduced-power converse & n/a & \hyperref[thm:p03-resplendence-atomic-powers]{\ref*{thm:p03-resplendence-atomic-powers}} & \thepage & Finite intersections are cones for finite unions and are nonempty. \\
\phantomsection\label{entry:symbol-registry-06129}\(f_x\) & \texttt{\detokenize{f_x}} & coordinate function assembled from finite-system solutions & P3 atomic reduced-power converse & n/a & \hyperref[thm:p03-resplendence-atomic-powers]{\ref*{thm:p03-resplendence-atomic-powers}} & \thepage & Selected actual assignments supply a simultaneous quotient solution before retraction. \\
\phantomsection\label{entry:symbol-registry-06130}\(c_k\) & \texttt{\detokenize{c_k}} & constant naming the standard number k & P3 separating named-equality example & n/a & \hyperref[thm:p03-resplendence-examples]{\ref*{thm:p03-resplendence-examples}} & \thepage & All names form a recursive language; the omitted type has no extra parameters. \\
\phantomsection\label{entry:symbol-registry-06131}\(\equiv\) & \texttt{\detokenize{\equiv}} & Delhomme-native logical identity atom on same-sort terms & P3 Delhomme atomic presentation and homomorphism diagram & relation & \hyperref[def:p03-resplendence-atomic]{\ref*{def:p03-resplendence-atomic}} & \thepage & Binding selectors: def:p03-resplendence-atomic t\_1 equiv t\_2; thm:p03-resplendence-positive-maps left satisfaction atom; annex:p03-resplendence-positive-maps native-atom proof; annex:p03-resplendence-atomic-retractions x\_a equiv a and function-value identity atoms; annex:p03-resplendence-topological native diagonal constraint; annex:p03-resplendence-atomic-powers native quotient-atom clause. Both syntactic operands are L(A)-terms; interpreted operands are objects in one sort. Not elementary equivalence of structures. \\
\phantomsection\label{entry:symbol-registry-06132}\(=\) & \texttt{\detokenize{=}} & host logical identity atom used only in the satisfaction comparison & P3 native-to-host satisfaction bridge & relation & \hyperref[thm:p03-resplendence-positive-maps]{\ref*{thm:p03-resplendence-positive-maps}} & \thepage & Binding selectors: rightmost C,s models (t\_1=t\_2) in the theorem's bridge and its exact proof at annex:p03-resplendence-positive-maps. Both operands are same-sort terms. This bridge does not rewrite the source presentation or its homomorphism diagram. \\
\phantomsection\label{entry:symbol-registry-06133}\(=\) & \texttt{\detokenize{=}} & metamathematical identity of interpreted object values & P3 native-to-host satisfaction bridge & relation & \hyperref[thm:p03-resplendence-positive-maps]{\ref*{thm:p03-resplendence-positive-maps}} & \thepage & Binding selector: middle t\_1\^{}C(s)=t\_2\^{}C(s) clause of the bridge. Operands are interpreted objects not formulas or structures. Ordinary equals elsewhere remains unchanged. \\
\phantomsection\label{entry:symbol-registry-06134}\(t_j\) & \texttt{\detokenize{t_j}} & syntactic term operand of a logical identity atom & P3 Delhomme atomic definition and satisfaction bridge & n/a & \hyperref[def:p03-resplendence-atomic]{\ref*{def:p03-resplendence-atomic}} & \thepage & L(A)-terms j=1 or 2; all variables and named parameters receive their existing interpretations. \\
\phantomsection\label{entry:symbol-registry-06135}\(t_j^C(s)\) & \texttt{\detokenize{t_j^C(s)}} & object value of a term under assignment s & P3 satisfaction bridge & n/a & \hyperref[thm:p03-resplendence-positive-maps]{\ref*{thm:p03-resplendence-positive-maps}} & \thepage & The bridge proof evaluates variables parameters and function terms recursively before comparing the two values. \\
\phantomsection\label{entry:symbol-registry-06136}\(C\) & \texttt{\detokenize{C}} & structure interpreting the satisfaction bridge and typed contrast & P3 atomic-definition contrast and satisfaction bridge & n/a & \hyperref[def:p03-resplendence-atomic]{\ref*{def:p03-resplendence-atomic}} & \thepage & Binding selectors: structure C in C equiv D contrast and C,s models displays; not a Boolean coefficient algebra or an object term. \\
\phantomsection\label{entry:symbol-registry-06137}\(D\) & \texttt{\detokenize{D}} & structure operand of elementary-equivalence contrast & P3 atomic-definition typed contrast & n/a & \hyperref[def:p03-resplendence-atomic]{\ref*{def:p03-resplendence-atomic}} & \thepage & Binding selector: D in C equiv D only. The full elementary diagram D in the earlier Henkin proofs retains its separate declared role. \\
\phantomsection\label{entry:symbol-registry-06284}\(K\) & \texttt{\detokenize{K}} & isomorphism-closed model class & P3 invariant-class continuation & n/a & \hyperref[def:p03-lindstrom-standing]{\ref*{def:p03-lindstrom-standing}} & \thepage & One fixed relational vocabulary; set parameters permitted in the class predicate; only set-sized code families enter formulas. \\
\phantomsection\label{entry:symbol-registry-06285}\(L^*\) & \texttt{\detokenize{L^*}} & source abstract logic pair & P3 Kennedy-Vaananen presentation & n/a & \hyperref[def:p03-lindstrom-source-logic]{\ref*{def:p03-lindstrom-source-logic}} & \thepage & Native pair (Sigma,T); not identified with the stronger regular package from five operations alone. \\
\phantomsection\label{entry:symbol-registry-06286}\(\Sigma\) & \texttt{\detokenize{\Sigma}} & set or class of source sentences & P3 source-pair definition & n/a & \hyperref[def:p03-lindstrom-source-logic]{\ref*{def:p03-lindstrom-source-logic}} & \thepage & Syntactic carrier in the source pair; distinct from the projection-class operation Sigma(L*). \\
\phantomsection\label{entry:symbol-registry-06287}\(T\) & \texttt{\detokenize{T}} & source satisfaction relation & P3 source-pair definition & n/a & \hyperref[def:p03-lindstrom-source-logic]{\ref*{def:p03-lindstrom-source-logic}} & \thepage & Arguments are a sentence and a structure; not an axiom theory. \\
\phantomsection\label{entry:symbol-registry-06288}\(\operatorname{Mod}(\phi)\) & \texttt{\detokenize{\operatorname{Mod}(\phi)}} & model class of a source sentence & P3 source-pair definition & n/a & \hyperref[def:p03-lindstrom-source-logic]{\ref*{def:p03-lindstrom-source-logic}} & \thepage & Defined by T(phi,M) at the stated vocabulary. \\
\phantomsection\label{entry:symbol-registry-06289}\(\mathcal L^*\) & \texttt{\detokenize{\mathcal L^*}} & regular relational logic with finite occurrence & P3 compactness-LS characterization & n/a & \hyperref[def:p03-lindstrom-source-logic]{\ref*{def:p03-lindstrom-source-logic}} & \thepage & Includes relativization and fresh-constant abstraction as separate hypotheses; bridge to class-level operations is proved. \\
\phantomsection\label{entry:symbol-registry-06290}\(Q_K\) & \texttt{\detokenize{Q_K}} & generalized quantifier defined by a model class & P3 binary class-quantifier construction & n/a & \hyperref[thm:p03-lindstrom-generalized-class]{\ref*{thm:p03-lindstrom-generalized-class}} & \thepage & Matrix selects a binary relation on the unchanged carrier; defining the class does not establish compactness. Original precedence descriptor 'quantifier' retained; no numerical rank inferred. \\
\phantomsection\label{entry:symbol-registry-06291}\(L_{\kappa\lambda}\) & \texttt{\detokenize{L_{\kappa\lambda}}} & infinitary syntax with two cardinal bounds & P3 invariant-class coding & n/a & \hyperref[def:p03-lindstrom-infinitary]{\ref*{def:p03-lindstrom-infinitary}} & \thepage & Connective size below kappa and block length below lambda; finite support is not inferred. \\
\phantomsection\label{entry:symbol-registry-06292}\(\eta_\alpha(x)\) & \texttt{\detokenize{\eta_\alpha(x)}} & source predecessor-order-type formula & P3 ordinal-position code & n/a & \hyperref[def:p03-lindstrom-ordinal]{\ref*{def:p03-lindstrom-ordinal}} & \thepage & Actual transfinite recursion; object variable x; validity of the order-type test requires the linear-order guard. \\
\phantomsection\label{entry:symbol-registry-06293}\(\eta^0_\alpha\) & \texttt{\detokenize{\eta^0_\alpha}} & source whole-order-type sentence & P3 ordinal-position code & n/a & \hyperref[def:p03-lindstrom-ordinal]{\ref*{def:p03-lindstrom-ordinal}} & \thepage & Superscript zero is preserved; alpha=0 is false on a nonempty carrier. \\
\phantomsection\label{entry:symbol-registry-06294}\(\operatorname{LO}(<)\) & \texttt{\detokenize{\operatorname{LO}(<)}} & strict-linear-order guard & P3 ordinal-position code & n/a & \hyperref[def:p03-lindstrom-ordinal]{\ref*{def:p03-lindstrom-ordinal}} & \thepage & The source's linear-order regime is made an actual conjunct when expansions are arbitrary. \\
\phantomsection\label{entry:symbol-registry-06295}\(K_\lambda\) & \texttt{\detokenize{K_\lambda}} & cardinality-restricted model class & P3 cardinal-dependent definability & n/a & \hyperref[def:p03-lindstrom-cardinal-dependent]{\ref*{def:p03-lindstrom-cardinal-dependent}} & \thepage & Cardinal-relative witness pairs are distinguished from the complement among all nonempty base structures. \\
\phantomsection\label{entry:symbol-registry-06296}\(\Sigma(L^*)\) & \texttt{\detokenize{\Sigma(L^*)}} & projection-definable model classes & P3 semantic projection operation & n/a & \hyperref[def:p03-lindstrom-cardinal-dependent]{\ref*{def:p03-lindstrom-cardinal-dependent}} & \thepage & Forget additional interpretations on the same domain; not an object-language first-order relation quantifier. \\
\phantomsection\label{entry:symbol-registry-06297}\(\Delta(L^*)\) & \texttt{\detokenize{\Delta(L^*)}} & classes with both positive and complementary projections & P3 semantic projection operation & n/a & \hyperref[def:p03-lindstrom-cardinal-dependent]{\ref*{def:p03-lindstrom-cardinal-dependent}} & \thepage & Global complementarity and its admission conditions are explicit. \\
\phantomsection\label{entry:symbol-registry-06298}\(\theta_A\) & \texttt{\detokenize{\theta_A}} & complete enumerating signed-diagram sentence & P3 McGee-style cardinal code & n/a & \hyperref[thm:p03-lindstrom-mcgee-code]{\ref*{thm:p03-lindstrom-mcgee-code}} & \thepage & Uses a lambda-long variable block and exact domain cover; indexed by the actual code A. \\
\phantomsection\label{entry:symbol-registry-06299}\(f_A\) & \texttt{\detokenize{f_A}} & bijection from the fixed cardinal to the coded model & P3 relation-diagram construction & n/a & \hyperref[thm:p03-lindstrom-mcgee-code]{\ref*{thm:p03-lindstrom-mcgee-code}} & \thepage & Both sign-test arguments use this same bijection; no undefined f\_B is used. \\
\phantomsection\label{entry:symbol-registry-06300}\(\rho_{\alpha,\beta}\) & \texttt{\detokenize{\rho_{\alpha,\beta}}} & source positive or negative binary relation atom & P3 ordinal-expansion coding & n/a & \hyperref[thm:p03-lindstrom-delta-code]{\ref*{thm:p03-lindstrom-delta-code}} & \thepage & Sign chosen from A's complete relation table; variables are object operands. \\
\phantomsection\label{entry:symbol-registry-06301}\(\Phi_A\) & \texttt{\detokenize{\Phi_A}} & source relation diagram tested at ordinal positions & P3 Delta cardinal code & n/a & \hyperref[thm:p03-lindstrom-delta-code]{\ref*{thm:p03-lindstrom-delta-code}} & \thepage & Native finite-quantifier formula with set-sized conjunction; relation table and actual isomorphism bridge are proved. \\
\phantomsection\label{entry:symbol-registry-06302}\(\Theta_K\) & \texttt{\detokenize{\Theta_K}} & source disjunction over selected fixed-domain model codes & P3 Delta cardinal code & n/a & \hyperref[thm:p03-lindstrom-delta-code]{\ref*{thm:p03-lindstrom-delta-code}} & \thepage & Set-sized code set; no proper-class disjunction. \\
\phantomsection\label{entry:symbol-registry-06303}\(P\) & \texttt{\detokenize{P}} & auxiliary subset for the larger-cardinality complement branch & P3 global Delta complement construction & n/a & \hyperref[thm:p03-lindstrom-delta-code]{\ref*{thm:p03-lindstrom-delta-code}} & \thepage & Its induced order has type lambda-plus; not the predicate of the separate amalgamation counterexample. \\
\phantomsection\label{entry:symbol-registry-06304}\(\Sigma_1\) & \texttt{\detokenize{\Sigma_1}} & set-theoretic existential bounded-matrix definability & P3 syntax and satisfaction absoluteness & n/a & \hyperref[thm:p03-lindstrom-absolute-truth]{\ref*{thm:p03-lindstrom-absolute-truth}} & \thepage & Metadata and truth-table witnesses are sets; distinct from the projection-class operator. \\
\phantomsection\label{entry:symbol-registry-06305}\(\Delta_1\) & \texttt{\detokenize{\Delta_1}} & both truth and nontruth Sigma-one definability & P3 satisfaction absoluteness & n/a & \hyperref[thm:p03-lindstrom-absolute-truth]{\ref*{thm:p03-lindstrom-absolute-truth}} & \thepage & A semantic complexity class; not the Delta-extension operator or a blanket admission claim for all complement pairs. \\
\phantomsection\label{entry:symbol-registry-06306}\(Q_0\) & \texttt{\detokenize{Q_0}} & at-least-countably-infinite cardinality quantifier & P3 compactness assumption countercase & n/a & \hyperref[thm:p03-lindstrom-assumptions]{\ref*{thm:p03-lindstrom-assumptions}} & \thepage & Native Q-alpha convention at alpha zero; actual countable witness closure is proved. Original precedence descriptor 'quantifier' retained; no numerical rank inferred. \\
\phantomsection\label{entry:symbol-registry-06307}\(Q_1\) & \texttt{\detokenize{Q_1}} & uncountability quantifier & P3 LS assumption countercase & n/a & \hyperref[thm:p03-lindstrom-assumptions]{\ref*{thm:p03-lindstrom-assumptions}} & \thepage & Reuses the existing public semantics; only the stated LS failure is newly proved. Original precedence descriptor 'quantifier' retained; no numerical rank inferred. \\
\phantomsection\label{entry:symbol-registry-06846}\(e(\mathfrak A)\) & \texttt{\detokenize{e(\mathfrak A)}} & Fagin's word encoding of a finite structure and the set of encodings of a class symbol family member \(e(\mathfrak A)\) & \hyperref[def:p02-w02-p02-036-fagin-encoding]{\ref*{def:p02-w02-p02-036-fagin-encoding}} & n/a & \hyperref[def:p02-w02-p02-036-fagin-encoding]{\ref*{def:p02-w02-p02-036-fagin-encoding}} & \thepage & Source notation of Fagin 1974 pp. 52-53; universe {1,...,n}; same notation as the P9 presentation sec:p09-presentation-n-c-004. Explicit displayed family; member types and common actual first-use exposed, not one atomic glyph. Family listing, not one atomic glyph. Author macro \texttt{\detokenize{e(\cdot),}\allowbreak\detokenize{E(\cdot)}} not carried to the members. \\
\phantomsection\label{entry:symbol-registry-06847}\(E(A)\) & \texttt{\detokenize{E(A)}} & Fagin's word encoding of a finite structure and the set of encodings of a class symbol family member \(E(A)\) & \hyperref[def:p02-w02-p02-036-fagin-encoding]{\ref*{def:p02-w02-p02-036-fagin-encoding}} & n/a & \hyperref[def:p02-w02-p02-036-fagin-encoding]{\ref*{def:p02-w02-p02-036-fagin-encoding}} & \thepage & Source notation of Fagin 1974 pp. 52-53; universe {1,...,n}; same notation as the P9 presentation sec:p09-presentation-n-c-004. Explicit displayed family; member types and common actual first-use exposed, not one atomic glyph. Family listing, not one atomic glyph. Author macro \texttt{\detokenize{e(\cdot),}\allowbreak\detokenize{E(\cdot)}} not carried to the members. \\
\phantomsection\label{entry:symbol-registry-06848}\(\mathsf{NP}_1\) & \texttt{\detokenize{\mathsf{NP}_1}} & Fagin's class of sets recognized nondeterministically in exponential time & \hyperref[def:p02-w02-p02-036-fagin-encoding]{\ref*{def:p02-w02-p02-036-fagin-encoding}} & n/a & \hyperref[def:p02-w02-p02-036-fagin-encoding]{\ref*{def:p02-w02-p02-036-fagin-encoding}} & \thepage & Source notation of Fagin 1974 p. 51; also used in sec:p09-presentation-n-c-004. \\
\phantomsection\label{entry:symbol-registry-07506}\(\Omega\) & \texttt{\detokenize{\Omega}} & index set for the reduced-product presentation & P3 FV reduced products & n/a & \hyperref[def:p03-fv-reduced-standing]{\ref*{def:p03-fv-reduced-standing}} & \thepage & Nonempty set; distinct from the ideal I. \\
\phantomsection\label{entry:symbol-registry-07507}\(I\) & \texttt{\detokenize{I}} & proper ideal of subsets of the index set & P3 FV reduced products & n/a & \hyperref[def:p03-fv-reduced-standing]{\ref*{def:p03-fv-reduced-standing}} & \thepage & Native Ghasemi index/ideal roles retained; I in the earlier direct-product presentation was an index set. \\
\phantomsection\label{entry:symbol-registry-07508}\(F_I\) & \texttt{\detokenize{F_I}} & dual proper filter of the ideal & P3 FV reduced products & n/a & \hyperref[def:p03-fv-reduced-standing]{\ref*{def:p03-fv-reduced-standing}} & \thepage & A subset is large exactly when its complement belongs to I. \\
\phantomsection\label{entry:symbol-registry-07509}\(A_\gamma\) & \texttt{\detokenize{A_\gamma}} & factor structure at an index & P3 FV reduced products & n/a & \hyperref[def:p03-fv-reduced-standing]{\ref*{def:p03-fv-reduced-standing}} & \thepage & Every sort nonempty; common finitary many-sorted signature. \\
\phantomsection\label{entry:symbol-registry-07510}\(A_I\) & \texttt{\detokenize{A_I}} & ordinary two-valued reduced product over an ideal & P3 FV reduced products & n/a & \hyperref[def:p03-fv-reduced-standing]{\ref*{def:p03-fv-reduced-standing}} & \thepage & Native source letter used for the explicitly distinguished classical specialization; not its unrestricted real-valued metric structure. \\
\phantomsection\label{entry:symbol-registry-07511}\(\pi_I\) & \texttt{\detokenize{\pi_I}} & quotient map on product sort elements & P3 FV reduced products & n/a & \hyperref[def:p03-fv-reduced-standing]{\ref*{def:p03-fv-reduced-standing}} & \thepage & Applied coordinatewise to finite tuples. \\
\phantomsection\label{entry:symbol-registry-07512}\(\mathcal P(\Omega)/I\) & \texttt{\detokenize{\mathcal P(\Omega)/}\allowbreak\detokenize{I}} & quotient Boolean algebra of index subsets & P3 FV reduced products & n/a & \hyperref[def:p03-fv-reduced-standing]{\ref*{def:p03-fv-reduced-standing}} & \thepage & Symmetric differences in I identify sets; first-order Boolean quantifiers range over classes. \\
\phantomsection\label{entry:symbol-registry-07513}\([S]_I\) & \texttt{\detokenize{[S]_I}} & Boolean equivalence class of an index subset & P3 FV reduced products & n/a & \hyperref[def:p03-fv-reduced-standing]{\ref*{def:p03-fv-reduced-standing}} & \thepage & Distinct from the element quotient pi\_I and from an actual representative subset. \\
\phantomsection\label{entry:symbol-registry-07514}\(X_\psi\) & \texttt{\detokenize{X_\psi}} & factor-formula truth set for representatives & P3 FV reduced products & n/a & \hyperref[def:p03-fv-reduced-standing]{\ref*{def:p03-fv-reduced-standing}} & \thepage & Its class is invariant under equivalent parameter representatives. \\
\phantomsection\label{entry:symbol-registry-07515}\(Z_\sigma\) & \texttt{\detokenize{Z_\sigma}} & quotient-Boolean existential partition cell & P3 FV reduced composition & n/a & \hyperref[def:p03-fv-reduced-reduction]{\ref*{def:p03-fv-reduced-reduction}} & \thepage & Same recursive reduction syntax; now a class in P(Omega)/I rather than a subset in the earlier full-product presentation. \\
\phantomsection\label{entry:symbol-registry-07516}\(S_\sigma\) & \texttt{\detokenize{S_\sigma}} & actual lifted partition of the index set & P3 FV reduced composition proof & n/a & \hyperref[annex:p03-fv-reduced-reduction]{\ref*{annex:p03-fv-reduced-reduction}} & \thepage & Lifted sets partition Omega exactly and represent Z\_sigma. \\
\phantomsection\label{entry:symbol-registry-07517}\(E_\sigma\) & \texttt{\detokenize{E_\sigma}} & availability-failure subset of a lifted cell & P3 FV reduced composition proof & n/a & \hyperref[annex:p03-fv-reduced-reduction]{\ref*{annex:p03-fv-reduced-reduction}} & \thepage & Each belongs to I by the Boolean availability inequality. \\
\phantomsection\label{entry:symbol-registry-07518}\(E\) & \texttt{\detokenize{E}} & finite union of exceptional index subsets & P3 FV reduced composition proof & n/a & \hyperref[annex:p03-fv-reduced-reduction]{\ref*{annex:p03-fv-reduced-reduction}} & \thepage & Belongs to I; witness values may be arbitrary there. \\
\phantomsection\label{entry:symbol-registry-07519}\(h\) & \texttt{\detokenize{h}} & surjective Boolean-algebra homomorphism & P3 FV finite partition lifting & n/a & \hyperref[thm:p03-fv-finite-partition-lift]{\ref*{thm:p03-fv-finite-partition-lift}} & \thepage & Finite partition lifts use finite choices of preimages. \\
\phantomsection\label{entry:symbol-registry-07520}\(B^*\) & \texttt{\detokenize{B^*}} & compact Stone space of a Boolean algebra & P3 FV source bounded-power laws & n/a & \hyperref[def:p03-fv-source-power-standing]{\ref*{def:p03-fv-source-power-standing}} & \thepage & Burris--Sankappanavar native starred notation; empty for the trivial Boolean algebra. \\
\phantomsection\label{entry:symbol-registry-07521}\(C(X,A)\) & \texttt{\detokenize{C(X,}\allowbreak\detokenize{A)}} & algebra of continuous maps into a discrete base algebra & P3 FV source bounded-power laws & n/a & \hyperref[def:p03-fv-source-power-standing]{\ref*{def:p03-fv-source-power-standing}} & \thepage & Pointwise operations; finite range follows from compactness of the Stone domain. \\
\phantomsection\label{entry:symbol-registry-07522}\(N_b\) & \texttt{\detokenize{N_b}} & basic clopen of a coefficient Boolean element & P3 FV source bounded-power laws & n/a & \hyperref[def:p03-fv-source-power-standing]{\ref*{def:p03-fv-source-power-standing}} & \thepage & Source-native clopen consists of ultrafilters containing b. \\
\phantomsection\label{entry:symbol-registry-07523}\(X^*\) & \texttt{\detokenize{X^*}} & Boolean algebra of clopen subsets of a Boolean space & P3 FV source iterated-power law & n/a & \hyperref[def:p03-fv-stone-product-coefficients]{\ref*{def:p03-fv-stone-product-coefficients}} & \thepage & Oppositely typed star to B\^{}*; retained rather than renamed. \\
\phantomsection\label{entry:symbol-registry-07524}\(B_1*B_2\) & \texttt{\detokenize{B_1*B_2}} & clopen algebra of the Cartesian product of Stone spaces & P3 FV source iterated-power law & n/a & \hyperref[def:p03-fv-stone-product-coefficients]{\ref*{def:p03-fv-stone-product-coefficients}} & \thepage & Source IV section5 Exercise1; not the direct product B\_1 times B\_2. \\
\phantomsection\label{entry:symbol-registry-07525}\(F\) & \texttt{\detokenize{F}} & filter of a coefficient Boolean algebra & P3 FV source quotient-power law & n/a & \hyperref[def:p03-fv-power-filter-quotient]{\ref*{def:p03-fv-power-filter-quotient}} & \thepage & Source Exercise2 role permits the improper filter; distinct from the necessarily proper dual filter F\_I of the reduced-product ideal. \\
\phantomsection\label{entry:symbol-registry-07526}\(F^*\) & \texttt{\detokenize{F^*}} & closed Stone subspace of ultrafilters extending a coefficient filter & P3 FV source quotient-power law & n/a & \hyperref[def:p03-fv-power-filter-quotient]{\ref*{def:p03-fv-power-filter-quotient}} & \thepage & BPIT makes proper filters extendible; empty for the improper filter. \\
\phantomsection\label{entry:symbol-registry-07527}\(q_F\) & \texttt{\detokenize{q_F}} & coefficient Boolean quotient homomorphism & P3 FV source quotient-power law & n/a & \hyperref[def:p03-fv-power-filter-quotient]{\ref*{def:p03-fv-power-filter-quotient}} & \thepage & Sends each labeled partition coefficient to its quotient class. \\
\phantomsection\label{entry:symbol-registry-07528}\(B/F\) & \texttt{\detokenize{B/}\allowbreak\detokenize{F}} & Boolean quotient by filter agreement & P3 FV source quotient-power law & n/a & \hyperref[def:p03-fv-power-filter-quotient]{\ref*{def:p03-fv-power-filter-quotient}} & \thepage & b and c agree modulo F when their Boolean biconditional belongs to F. \\
\phantomsection\label{entry:symbol-registry-07529}\(\theta_F\) & \texttt{\detokenize{\theta_F}} & congruence on the bounded algebra power & P3 FV source quotient-power law & n/a & \hyperref[def:p03-fv-power-filter-quotient]{\ref*{def:p03-fv-power-filter-quotient}} & \thepage & Source-native equality on the closed subspace F\^{}*; kernel of coefficient transport. \\
\phantomsection\label{entry:symbol-registry-07530}\([\![a=b]\!]\) & \texttt{\detokenize{[\![a=b]\!]}} & clopen equalizer of bounded continuous-power elements & P3 FV source quotient-power law & n/a & \hyperref[def:p03-fv-power-filter-quotient]{\ref*{def:p03-fv-power-filter-quotient}} & \thepage & Source Exercise2 uses a subset of B\^{}*; formal presentation gives its corresponding Boolean element. \\
\phantomsection\label{entry:symbol-registry-07531}\(N_F\) & \texttt{\detokenize{N_F}} & Boolean ideal dual to a coefficient filter & P3 FV source quotient-power proof & n/a & \hyperref[annex:p03-fv-power-quotient]{\ref*{annex:p03-fv-power-quotient}} & \thepage & Defines and verifies the quotient Boolean congruence before coefficient transport. \\
\phantomsection\label{entry:symbol-registry-07532}\(\Delta(b,c)\) & \texttt{\detokenize{\Delta(b,}\allowbreak\detokenize{c)}} & Boolean symmetric-difference term & P3 FV source quotient-power proof & n/a & \hyperref[annex:p03-fv-power-quotient]{\ref*{annex:p03-fv-power-quotient}} & \thepage & Declared Boolean term used to prove the quotient equivalence and operation congruence; distinct from the factor pattern formulas delta\_sigma. \\
\phantomsection\label{entry:symbol-registry-07533}\(X\) & \texttt{\detokenize{X}} & nonempty topological point carrier & P3 topological correspondence language & n/a & \hyperref[def:p03-topological-carrier]{\ref*{def:p03-topological-carrier}} & \thepage & Not the coding space of ordinary structures. \\
\phantomsection\label{entry:symbol-registry-07534}\(\mathrm{Prop}\) & \texttt{\detokenize{\mathrm{Prop}}} & countably infinite set of propositional letters & P3 topological correspondence language & n/a & \hyperref[def:p03-topological-carrier]{\ref*{def:p03-topological-carrier}} & \thepage & Source2006 correspondence signature. \\
\phantomsection\label{entry:symbol-registry-07535}\(\tau\) & \texttt{\detokenize{\tau}} & full topology of a point carrier & P3 topological correspondence language & n/a & \hyperref[def:p03-topological-carrier]{\ref*{def:p03-topological-carrier}} & \thepage & Closed under arbitrary unions and finite intersections. \\
\phantomsection\label{entry:symbol-registry-07536}\(\sigma\) & \texttt{\detokenize{\sigma}} & basis of the point topology & P3 topological correspondence language & n/a & \hyperref[def:p03-topological-carrier]{\ref*{def:p03-topological-carrier}} & \thepage & October2006 source convention contains empty; distinguished from Boolean reduction sigma below. \\
\phantomsection\label{entry:symbol-registry-07537}\(\widehat\sigma\) & \texttt{\detokenize{\widehat\sigma}} & topology generated by a basis & P3 topological correspondence language & n/a & \hyperref[def:p03-topological-carrier]{\ref*{def:p03-topological-carrier}} & \thepage & May contain uncountably many opens even for a countable basis. \\
\phantomsection\label{entry:symbol-registry-07538}\(\nu\) & \texttt{\detokenize{\nu}} & assignment of point subsets to propositional letters & P3 topological correspondence language & n/a & \hyperref[def:p03-topological-carrier]{\ref*{def:p03-topological-carrier}} & \thepage & No openness assumption on predicate extensions. \\
\phantomsection\label{entry:symbol-registry-07539}\(\widehat M\) & \texttt{\detokenize{\widehat M}} & generated topological model of a basoid model & P3 topological correspondence language & n/a & \hyperref[def:p03-topological-carrier]{\ref*{def:p03-topological-carrier}} & \thepage & Same carrier and predicate interpretations; open-sort range changes. \\
\phantomsection\label{entry:symbol-registry-07540}\(L_2\) & \texttt{\detokenize{L_2}} & unrestricted two-sorted point and open language & P3 topological correspondence language & n/a & \hyperref[def:p03-topological-languages]{\ref*{def:p03-topological-languages}} & \thepage & Open equality is extensional; full-topology semantics differs from basis semantics. \\
\phantomsection\label{entry:symbol-registry-07541}\(L_t\) & \texttt{\detokenize{L_t}} & guarded membership-polarity topological language & P3 topological correspondence language & n/a & \hyperref[def:p03-topological-languages]{\ref*{def:p03-topological-languages}} & \thepage & Open equality is expanded before polarity testing; does not count equality as a monotone membership atom. \\
\phantomsection\label{entry:symbol-registry-07542}\(x\varepsilon U\) & \texttt{\detokenize{x\varepsilon U}} & point membership in an open parameter & P3 topological correspondence language & n/a & \hyperref[def:p03-topological-languages]{\ref*{def:p03-topological-languages}} & \thepage & Bound open range is topology or basis; free subset parameters can be fixed during the induction. \\
\phantomsection\label{entry:symbol-registry-07543}\(P_p\) & \texttt{\detokenize{P_p}} & unary point predicate corresponding to a modal letter & P3 topological correspondence language & n/a & \hyperref[def:p03-topological-languages]{\ref*{def:p03-topological-languages}} & \thepage & Interpretation is nu(p); not assumed open or closed. \\
\phantomsection\label{entry:symbol-registry-07544}\(\mathrm{Ext}\) & \texttt{\detokenize{\mathrm{Ext}}} & open-sort extensionality axiom & P3 topological basis axioms & n/a & \hyperref[def:p03-topological-basis-axioms]{\ref*{def:p03-topological-basis-axioms}} & \thepage & Part of the source four-sentence basoid presentation. \\
\phantomsection\label{entry:symbol-registry-07545}\(\mathrm{Union}\) & \texttt{\detokenize{\mathrm{Union}}} & basis cover axiom & P3 topological basis axioms & n/a & \hyperref[def:p03-topological-basis-axioms]{\ref*{def:p03-topological-basis-axioms}} & \thepage & Does not assert closure under all unions. \\
\phantomsection\label{entry:symbol-registry-07546}\(\mathrm{Empty}\) & \texttt{\detokenize{\mathrm{Empty}}} & empty basis-element axiom & P3 topological basis axioms & n/a & \hyperref[def:p03-topological-basis-axioms]{\ref*{def:p03-topological-basis-axioms}} & \thepage & The later ideal-box basis may require adjoining empty to meet this convention. \\
\phantomsection\label{entry:symbol-registry-07547}\(\mathrm{Bas}\) & \texttt{\detokenize{\mathrm{Bas}}} & local intersection refinement axiom & P3 topological basis axioms & n/a & \hyperref[def:p03-topological-basis-axioms]{\ref*{def:p03-topological-basis-axioms}} & \thepage & Point-dependent refinement is sufficient for a basis. \\
\phantomsection\label{entry:symbol-registry-07548}\(\mathrm{ST}_x\) & \texttt{\detokenize{\mathrm{ST}_x}} & standard modal-to-topological translation & P3 topological modal correspondence & n/a & \hyperref[def:p03-topological-modal-translation]{\ref*{def:p03-topological-modal-translation}} & \thepage & Box translates to a guarded open existential with a negative membership matrix. \\
\phantomsection\label{entry:symbol-registry-07549}\(I\) & \texttt{\detokenize{I}} & index set for topological ideal-products & P3 topological reduced ideal-products & n/a & \hyperref[def:p03-topological-rip]{\ref*{def:p03-topological-rip}} & \thepage & Source2001 index role; distinct from ideal I in the other reduced-product packet. \\
\phantomsection\label{entry:symbol-registry-07550}\(\Lambda\) & \texttt{\detokenize{\Lambda}} & ideal of allowed box support subsets & P3 topological reduced ideal-products & n/a & \hyperref[def:p03-topological-rip]{\ref*{def:p03-topological-rip}} & \thepage & May be improper to permit full boxes; not the element-equivalence filter. \\
\phantomsection\label{entry:symbol-registry-07551}\(\Psi\) & \texttt{\detokenize{\Psi}} & proper filter defining point and Boolean quotients & P3 topological reduced ideal-products & n/a & \hyperref[def:p03-topological-rip]{\ref*{def:p03-topological-rip}} & \thepage & Finite intersection of large agreement sets is used in the proof. \\
\phantomsection\label{entry:symbol-registry-07552}\(S_i\) & \texttt{\detokenize{S_i}} & factor point structure & P3 topological reduced ideal-products & n/a & \hyperref[def:p03-topological-rip]{\ref*{def:p03-topological-rip}} & \thepage & Nonempty carrier; operations need not be continuous. \\
\phantomsection\label{entry:symbol-registry-07553}\(O_i\) & \texttt{\detokenize{O_i}} & factor topology & P3 topological reduced ideal-products & n/a & \hyperref[def:p03-topological-rip]{\ref*{def:p03-topological-rip}} & \thepage & Contains empty and whole carrier. \\
\phantomsection\label{entry:symbol-registry-07554}\(B_\Lambda\) & \texttt{\detokenize{B_\Lambda}} & basis of ideal-supported boxes & P3 topological reduced ideal-products & n/a & \hyperref[def:p03-topological-rip]{\ref*{def:p03-topological-rip}} & \thepage & This ordinary basis need not contain empty before it is explicitly adjoined. \\
\phantomsection\label{entry:symbol-registry-07555}\(O_\Lambda\) & \texttt{\detokenize{O_\Lambda}} & generated topology on the full point product & P3 topological reduced ideal-products & n/a & \hyperref[def:p03-topological-rip]{\ref*{def:p03-topological-rip}} & \thepage & Finite support gives product topology; unrestricted support gives box topology. \\
\phantomsection\label{entry:symbol-registry-07556}\(q\) & \texttt{\detokenize{q}} & open point-quotient map & P3 topological reduced ideal-products & n/a & \hyperref[def:p03-topological-rip]{\ref*{def:p03-topological-rip}} & \thepage & Its image on a box is a basic quotient open. \\
\phantomsection\label{entry:symbol-registry-07557}\([f]\) & \texttt{\detokenize{[f]}} & filter equivalence class of a product point & P3 topological reduced ideal-products & n/a & \hyperref[def:p03-topological-rip]{\ref*{def:p03-topological-rip}} & \thepage & Point-sort value; not the Boolean index class [A]. \\
\phantomsection\label{entry:symbol-registry-07558}\(\pi_i\) & \texttt{\detokenize{\pi_i}} & coordinate projection of the full point product & P3 topological reduced ideal-products & n/a & \hyperref[def:p03-topological-rip]{\ref*{def:p03-topological-rip}} & \thepage & Applied to actual boxes to define canonical factor parameters. \\
\phantomsection\label{entry:symbol-registry-07559}\(\prod^\Lambda_\Psi S_i\) & \texttt{\detokenize{\prod^\Lambda_\Psi S_i}} & topological reduced ideal-product & P3 topological reduced ideal-products & n/a & \hyperref[def:p03-topological-rip]{\ref*{def:p03-topological-rip}} & \thepage & Carries quotient topology and ordinary filter-large relation interpretations. \\
\phantomsection\label{entry:symbol-registry-07560}\(O^\Lambda_\Psi\) & \texttt{\detokenize{O^\Lambda_\Psi}} & quotient topology of an ideal-supported product & P3 topological reduced ideal-products & n/a & \hyperref[def:p03-topological-rip]{\ref*{def:p03-topological-rip}} & \thepage & Basic image family generates this full topology. \\
\phantomsection\label{entry:symbol-registry-07561}\(B^\Lambda_\Psi\) & \texttt{\detokenize{B^\Lambda_\Psi}} & basic image family of the open quotient map & P3 topological reduced ideal-products & n/a & \hyperref[def:p03-topological-rip]{\ref*{def:p03-topological-rip}} & \thepage & Not identified with the collection of all quotient opens. \\
\phantomsection\label{entry:symbol-registry-07562}\(U_i\) & \texttt{\detokenize{U_i}} & canonical coordinate projection of an actual basic box & P3 topological reduced ideal-products & n/a & \hyperref[def:p03-topological-rip]{\ref*{def:p03-topological-rip}} & \thepage & All empty for an empty box; arbitrary written factors of an empty rectangle are not used as parameters. \\
\phantomsection\label{entry:symbol-registry-07563}\(\equiv\) & \texttt{\detokenize{\equiv}} & filter agreement congruence of index subsets & P3 topological Boolean index algebra & n/a & \hyperref[def:p03-topological-boolean-ideal]{\ref*{def:p03-topological-boolean-ideal}} & \thepage & Source-native role is not elementary equivalence of structures. \\
\phantomsection\label{entry:symbol-registry-07564}\([A]\) & \texttt{\detokenize{[A]}} & Boolean class of an index subset & P3 topological Boolean index algebra & n/a & \hyperref[def:p03-topological-boolean-ideal]{\ref*{def:p03-topological-boolean-ideal}} & \thepage & Distinct type from the point class [f] despite the source's bracket reuse. \\
\phantomsection\label{entry:symbol-registry-07565}\(\mathbf B\) & \texttt{\detokenize{\mathbf B}} & Boolean index quotient with distinguished support ideal & P3 topological Boolean index algebra & n/a & \hyperref[def:p03-topological-boolean-ideal]{\ref*{def:p03-topological-boolean-ideal}} & \thepage & Plain Boolean quotient alone loses the box support information. \\
\phantomsection\label{entry:symbol-registry-07566}\(+\) & \texttt{\detokenize{+}} & Boolean join of index classes & P3 topological Boolean index algebra & n/a & \hyperref[def:p03-topological-boolean-ideal]{\ref*{def:p03-topological-boolean-ideal}} & \thepage & Source-native operation; not a point-algebra operation. \\
\phantomsection\label{entry:symbol-registry-07567}\(\cdot\) & \texttt{\detokenize{\cdot}} & Boolean meet of index classes & P3 topological Boolean index algebra & n/a & \hyperref[def:p03-topological-boolean-ideal]{\ref*{def:p03-topological-boolean-ideal}} & \thepage & Source-native operation is induced by intersection. \\
\phantomsection\label{entry:symbol-registry-07568}\(u'\) & \texttt{\detokenize{u'}} & Boolean complement of an index class & P3 topological Boolean index algebra & n/a & \hyperref[def:p03-topological-boolean-ideal]{\ref*{def:p03-topological-boolean-ideal}} & \thepage & Complement in the declared quotient Boolean algebra. \\
\phantomsection\label{entry:symbol-registry-07569}\(\Lambda/{\equiv}\) & \texttt{\detokenize{\Lambda/}\allowbreak\detokenize{{\equiv}}} & image ideal in the Boolean index quotient & P3 topological Boolean index algebra & n/a & \hyperref[def:p03-topological-boolean-ideal]{\ref*{def:p03-topological-boolean-ideal}} & \thepage & Downward closure follows by replacing a smaller class by its intersection with a support representative. \\
\phantomsection\label{entry:symbol-registry-07570}\(\lambda\) & \texttt{\detokenize{\lambda}} & predicate for the support ideal & P3 topological Boolean index algebra & n/a & \hyperref[def:p03-topological-boolean-ideal]{\ref*{def:p03-topological-boolean-ideal}} & \thepage & Unary Boolean-sort predicate; not an infinitary join operator. \\
\phantomsection\label{entry:symbol-registry-07571}\(T_B\) & \texttt{\detokenize{T_B}} & Boolean algebra and ideal theory & P3 topological Boolean index algebra & n/a & \hyperref[def:p03-topological-boolean-ideal]{\ref*{def:p03-topological-boolean-ideal}} & \thepage & Native ideal axioms are displayed as first-order statements. \\
\phantomsection\label{entry:symbol-registry-07572}\(I_\psi\) & \texttt{\detokenize{I_\psi}} & factor truth set for point and canonical-box parameters & P3 topological Boolean index algebra & n/a & \hyperref[def:p03-topological-boolean-ideal]{\ref*{def:p03-topological-boolean-ideal}} & \thepage & Its class is input to the monotone Boolean reduction. \\
\phantomsection\label{entry:symbol-registry-07573}\(\varphi^Y\) & \texttt{\detokenize{\varphi^Y}} & whole-carrier replacement of a free open parameter & P3 topological Boolean reduction & n/a & \hyperref[thm:p03-topological-carrier-substitution]{\ref*{thm:p03-topological-carrier-substitution}} & \thepage & Only free memberships in Y are replaced by point reflexivity. \\
\phantomsection\label{entry:symbol-registry-07574}\(\sigma\) & \texttt{\detokenize{\sigma}} & increasing first-order Boolean reduction condition & P3 topological Boolean reduction & n/a & \hyperref[def:p03-topological-fv-recursion]{\ref*{def:p03-topological-fv-recursion}} & \thepage & Typed role distinct from the earlier topology basis sigma. \\
\phantomsection\label{entry:symbol-registry-07575}\(\tau\) & \texttt{\detokenize{\tau}} & increasing Boolean reduction condition of a matrix & P3 topological Boolean reduction & n/a & \hyperref[def:p03-topological-fv-recursion]{\ref*{def:p03-topological-fv-recursion}} & \thepage & Source-native recursive matrix condition; distinct from the earlier carrier topology tau. \\
\phantomsection\label{entry:symbol-registry-07576}\(\theta_j\) & \texttt{\detokenize{\theta_j}} & finite list of local matrix formulas & P3 topological Boolean reduction & n/a & \hyperref[def:p03-topological-fv-recursion]{\ref*{def:p03-topological-fv-recursion}} & \thepage & Free-open polarity retained in each local component. \\
\phantomsection\label{entry:symbol-registry-07577}\(z_k\) & \texttt{\detokenize{z_k}} & existential Boolean witness-class variable & P3 topological Boolean reduction & n/a & \hyperref[def:p03-topological-fv-recursion]{\ref*{def:p03-topological-fv-recursion}} & \thepage & Subset-conjunction constraints enforce intersections of witnesses; not an open neighborhood or product point. \\
\phantomsection\label{entry:symbol-registry-07578}\(s_k\) & \texttt{\detokenize{s_k}} & enumerated subset of the finite local-formula index set & P3 topological Boolean reduction & n/a & \hyperref[def:p03-topological-fv-recursion]{\ref*{def:p03-topological-fv-recursion}} & \thepage & All 2\^{}m subsets with singletons first and empty conjunction retained. \\
\phantomsection\label{entry:symbol-registry-07579}\(\eta_j\) & \texttt{\detokenize{\eta_j}} & whole-carrier local matrix formula & P3 topological Boolean reduction & n/a & \hyperref[def:p03-topological-fv-recursion]{\ref*{def:p03-topological-fv-recursion}} & \thepage & Records truth where the witness open is the entire factor carrier. \\
\phantomsection\label{entry:symbol-registry-07580}\(Z_k\) & \texttt{\detokenize{Z_k}} & representative of a Boolean witness class & P3 topological Boolean reduction proof & n/a & \hyperref[annex:p03-topological-fv]{\ref*{annex:p03-topological-fv}} & \thepage & Finite constraints are made exact on a common Psi-large set. \\
\phantomsection\label{entry:symbol-registry-07581}\(F\) & \texttt{\detokenize{F}} & common filter-large agreement set & P3 topological Boolean reduction proof & n/a & \hyperref[annex:p03-topological-fv]{\ref*{annex:p03-topological-fv}} & \thepage & Finite intersection of constraint-agreement sets. \\
\phantomsection\label{entry:symbol-registry-07582}\(L\) & \texttt{\detokenize{L}} & support set on which actual neighborhood witnesses are chosen & P3 topological Boolean reduction proof & n/a & \hyperref[annex:p03-topological-fv]{\ref*{annex:p03-topological-fv}} & \thepage & Union of the finitely many whole-carrier failures; belongs to Lambda. This set role differs from the point signature L. \\
\phantomsection\label{entry:symbol-registry-07583}\(\mathcal U\) & \texttt{\detokenize{\mathcal U}} & ultrafilter for the topological transfer corollary & P3 topological ultraproducts & n/a & \hyperref[thm:p03-topological-transfer]{\ref*{thm:p03-topological-transfer}} & \thepage & Full box support and generated quotient topology; not full unguarded L\_2 transfer. \\
\end{longtable}
\endgroup
\Needspace{.4\textheight}
\subsection*{Stable part P4 scopes}
\begingroup\scriptsize\setlength{\tabcolsep}{3pt}
\begin{longtable}{@{}p{0.9in}p{1.3in}p{1.0in}p{0.5in}p{0.3in}p{0.4in}p{0.35in}p{1.65in}@{}}
\toprule
\textbf{Glyph} & \textbf{Macro} & \textbf{Role} & \textbf{Scope} & \textbf{Prec./form} & \textbf{First} & \textbf{Reg. page} & \textbf{Principal uses / notes} \\
\midrule
\endfirsthead
\toprule
\textbf{Glyph} & \textbf{Macro} & \textbf{Role} & \textbf{Scope} & \textbf{Prec./form} & \textbf{First} & \textbf{Reg. page} & \textbf{Principal uses / notes} \\
\midrule
\endhead
\midrule\multicolumn{8}{r}{\emph{continued}}\\\endfoot
\bottomrule\endlastfoot
\phantomsection\label{entry:symbol-registry-00015}\(\mathcal A_{\mathrm{in}},\mathcal A_{\mathrm{tape}}\) & \texttt{\detokenize{\mathcal A_{\mathrm{in}},}\allowbreak\detokenize{\mathcal A_{\mathrm{tape}}}} & input and tape alphabets & P4 turing-machines & n/a & \hyperref[sec:turing-machines-01]{\ref*{sec:turing-machines-01}} & \thepage & No formula-context standing in this scope. \\
\phantomsection\label{entry:symbol-registry-00016}\(u,v,w\) & \texttt{\detokenize{u,}\allowbreak\detokenize{v,}\allowbreak\detokenize{w}} & finite words & P4 turing-machines & n/a & \hyperref[sec:turing-machines-01]{\ref*{sec:turing-machines-01}} & \thepage & Finite words over the declared alphabet; later P4 chapters use the letters numerically. \\
\phantomsection\label{entry:symbol-registry-00017}\(\Sigma^0_n,\Pi^0_n\) & \texttt{\detokenize{\Sigma^0_n,}\allowbreak\detokenize{\Pi^0_n}} & syntactic arithmetic formula classes & Legacy pre-batch4 P4-P7.5 reader meaning & n/a & \hyperref[sec:decision-hierarchy-01]{\ref*{sec:decision-hierarchy-01}} & \thepage & Classes fixed by the declared normalization and unbounded-quantifier alternation; P7.5 uses the same hierarchy in induction-fragment names.  Numeric chapter scope migrated to its frozen pre-batch4 source-labelled reader regions; glyph, role and precedence unchanged. \\
\phantomsection\label{entry:symbol-registry-00018}\(\boldsymbol{\Sigma}^0_n,\boldsymbol{\Pi}^0_n,\boldsymbol{\Delta}^0_n\) & \texttt{\detokenize{\boldsymbol{\Sigma}^0_n,}\allowbreak\detokenize{\boldsymbol{\Pi}^0_n,}\allowbreak\detokenize{\boldsymbol{\Delta}^0_n}} & arithmetical definability classes of relations & Parts 4–5 from turing-machines & n/a & \hyperref[sec:decision-hierarchy-01]{\ref*{sec:decision-hierarchy-01}} & \thepage & Bold symbols classify relations definable in \(\mathcal N_{\mathrm{ar}}\), not formula strings or Borel pointclasses. \\
\phantomsection\label{entry:symbol-registry-00049}\(\simeq\) & \texttt{\detokenize{\simeq}} & Kleene equality & Legacy pre-batch5 P4.2-P4.5 reader meaning & meta & \hyperref[sec:universal-reductions-01]{\ref*{sec:universal-reductions-01}} & \thepage & Both sides are defined with equal value or both are undefined.  Numeric chapter scope migrated to its frozen pre-batch5 source-labelled reader regions; glyph, role and precedence unchanged. \\
\phantomsection\label{entry:symbol-registry-00050}\(\leq_m\) & \texttt{\detokenize{\leq_m}} & many-one reducibility & Parts 4–5 from turing-machines & meta & \hyperref[sec:universal-reductions-01]{\ref*{sec:universal-reductions-01}} & \thepage & A total computable translator preserves membership. \\
\phantomsection\label{entry:symbol-registry-00051}\(\leq_T\) & \texttt{\detokenize{\leq_T}} & Turing reducibility & Parts 4–5 from turing-machines & meta & \hyperref[sec:p04-frontiers-computability]{\ref*{sec:p04-frontiers-computability}} & \thepage & An oracle machine for the right-hand set decides the left-hand set. \\
\phantomsection\label{entry:symbol-registry-00052}\(\varphi_e\) & \texttt{\detokenize{\varphi_e}} & partial computable function with index e & Parts 4–5 from turing-machines & n/a & \hyperref[sec:universal-reductions-01]{\ref*{sec:universal-reductions-01}} & \thepage & Oracle form \(\varphi^B_e\) is relative to \(B\); finite-string form \(\varphi_e^\tau\) and stage form \(\varphi^{B_s}_{e,s}\) are defined in the two priority proofs. \\
\phantomsection\label{entry:symbol-registry-00053}\(B'\) & \texttt{\detokenize{'}} & Turing jump & P4 frontiers-computability & n/a & \hyperref[sec:p04-frontiers-computability]{\ref*{sec:p04-frontiers-computability}} & \thepage & Postfix prime on an oracle set. \\
\phantomsection\label{entry:symbol-registry-00054}\(\varnothing^{(n)}\) & \texttt{\detokenize{\varnothing^{(n)}}} & iterated jump of the empty set & Parts 4–5 from frontiers-computability & n/a & \hyperref[sec:p04-frontiers-computability]{\ref*{sec:p04-frontiers-computability}} & \thepage & First defined with the general jump iteration. \\
\phantomsection\label{entry:symbol-registry-00056}\(\bar n\) & \texttt{\detokenize{\bar n}} & numeral & Parts 4–5 from decision-hierarchy & n/a & \hyperref[sec:decision-hierarchy-01]{\ref*{sec:decision-hierarchy-01}} & \thepage & Closed arithmetic term \(S^n0\). In the P5 representation construction, barred argument lists are typed tuples; an overline applied componentwise gives their numeral terms. \\
\phantomsection\label{entry:symbol-registry-00080}\(\leq\) & \texttt{\detokenize{\leq}} & object-language arithmetic order abbreviation & P4 decision-hierarchy & 2 & \hyperref[sec:decision-hierarchy-01]{\ref*{sec:decision-hierarchy-01}} & \thepage & Eliminable abbreviation \(r\leq t\defeq\exists z\,(z+r=t)\), with the least variable fresh for the displayed terms; marked bounded quantifiers inherit it. \\
\phantomsection\label{entry:symbol-registry-00081}\(\leq\) & \texttt{\detokenize{\leq}} & canonical ordering & Legacy pre-batch5 P4.4-P4.6 reader meaning & 2 & \hyperref[sec:p04-frontiers-computability]{\ref*{sec:p04-frontiers-computability}} & \thepage & Ambient order on degrees, natural numbers, and resource bounds; not arithmetic object-language syntax.  Numeric chapter scope migrated to its frozen pre-batch5 source-labelled reader regions; glyph, role and precedence unchanged. \\
\phantomsection\label{entry:symbol-registry-00084}\(a\leq_{\N}b\) & \texttt{\detokenize{\leq_{\N}}} & metalinguistic natural-number order & Parts 4–5 from decision-hierarchy to arithmetic-coding & 2 & \hyperref[sec:decision-hierarchy-01]{\ref*{sec:decision-hierarchy-01}} & \thepage & Used where bare \(\leq\) simultaneously denotes arithmetic-language syntax; operands are standard naturals or codes, never arithmetic terms. \\
\phantomsection\label{entry:symbol-registry-00086}\(<\) & \texttt{\detokenize{<}} & object-language arithmetic strict-order abbreviation & P4 decision-hierarchy & 2 & \hyperref[sec:decision-hierarchy-01]{\ref*{sec:decision-hierarchy-01}} & \thepage & Eliminable abbreviation \(r<t\defeq\exists z\,(r+Sz=t)\), with the least variable fresh for the displayed terms. \\
\phantomsection\label{entry:symbol-registry-00087}\(<\) & \texttt{\detokenize{<}} & canonical strict ordering & Parts 4–19 from frontiers-computability & 2 & \hyperref[sec:p04-frontiers-computability]{\ref*{sec:p04-frontiers-computability}} & \thepage & Ambient strict order on declared numerical and order-theoretic objects; not arithmetic object-language syntax. \\
\phantomsection\label{entry:symbol-registry-00088}\(a<_{\N}b\) & \texttt{\detokenize{<_{\N}}} & metalinguistic natural-number strict order & P4 decision-hierarchy & 2 & \hyperref[sec:decision-hierarchy-01]{\ref*{sec:decision-hierarchy-01}} & \thepage & Used where bare \(<\) simultaneously denotes arithmetic-language syntax; operands are standard naturals or codes, never arithmetic terms. \\
\phantomsection\label{entry:symbol-registry-00233}\(\vec a,\vec n\) & \texttt{\detokenize{\vec a,}\allowbreak\detokenize{\vec n}} & finite tuples of natural-number parameters & Parts 4–5 from turing-machines & n/a & \hyperref[sec:turing-machines-01]{\ref*{sec:turing-machines-01}} & \thepage & Vector arrows distinguish tuples from numeral bars. \\
\phantomsection\label{entry:symbol-registry-00234}\(\vec x,\vec y,\vec z,\vec u,\vec v\) & \texttt{\detokenize{\vec x,}\allowbreak\detokenize{\vec y,}\allowbreak\detokenize{\vec z,}\allowbreak\detokenize{\vec u,}\allowbreak\detokenize{\vec v}} & finite tuples of numerical variables & Parts 4–5 from turing-machines & n/a & \hyperref[sec:turing-machines-01]{\ref*{sec:turing-machines-01}} & \thepage & Vector arrows distinguish variable blocks from numeral bars. \\
\phantomsection\label{entry:symbol-registry-00235}\(\overline{\vec a}\) & \texttt{\detokenize{\overline{\vec a}}} & tuple of arithmetic numerals & Parts 4–5 from decision-hierarchy & n/a & \hyperref[sec:decision-hierarchy-01]{\ref*{sec:decision-hierarchy-01}} & \thepage & Componentwise numeral formation from a declared tuple of standard naturals. \\
\phantomsection\label{entry:symbol-registry-00236}\(\mathcal T\) & \texttt{\detokenize{\mathcal T}} & deterministic Turing-machine metavariable & P4 turing-machines to universal-reductions & n/a & \hyperref[sec:turing-machines-01]{\ref*{sec:turing-machines-01}} & \thepage & Machine under the fixed single-tape convention. \\
\phantomsection\label{entry:symbol-registry-00237}\(\sqcup\) & \texttt{\detokenize{\sqcup}} & blank tape symbol & P4 turing-machines to universal-reductions & n/a & \hyperref[sec:turing-machines-01]{\ref*{sec:turing-machines-01}} & \thepage & Distinguished tape-alphabet symbol. \\
\phantomsection\label{entry:symbol-registry-00238}\(\delta\) & \texttt{\detokenize{\delta}} & Turing-machine transition function & P4 turing-machines to universal-reductions & n/a & \hyperref[sec:turing-machines-01]{\ref*{sec:turing-machines-01}} & \thepage & Typed by the fixed machine tuple. \\
\phantomsection\label{entry:symbol-registry-00239}\(\longrightarrow_{\mathcal T}\) & \texttt{\detokenize{\longrightarrow_{\mathcal T}}} & one-step machine transition & P4 turing-machines to universal-reductions & meta & \hyperref[sec:turing-machines-01]{\ref*{sec:turing-machines-01}} & \thepage & Relation on configurations of the displayed machine. \\
\phantomsection\label{entry:symbol-registry-00240}\(\operatorname{bin}(n)\) & \texttt{\detokenize{\operatorname{bin}}} & canonical binary word & P4 turing-machines to universal-reductions & n/a & \hyperref[sec:turing-machines-01]{\ref*{sec:turing-machines-01}} & \thepage & Zero is represented by the one-symbol word zero. \\
\phantomsection\label{entry:symbol-registry-00241}\(\operatorname{bin}_m(i)\) & \texttt{\detokenize{\operatorname{bin}_m}} & fixed-width binary address & Legacy pre-batch5 P4.6 reader meaning & n/a & \hyperref[sec:p04-decision-problem-hardness]{\ref*{sec:p04-decision-problem-hardness}} & \thepage & Length-m expansion padded with leading zeros for succinct-instance bit addressing.  Numeric chapter scope migrated to its frozen pre-batch5 source-labelled reader regions; glyph, role and precedence unchanged. \\
\phantomsection\label{entry:symbol-registry-00242}\(\operatorname{pair}_{\N}(a,b)\) & \texttt{\detokenize{\operatorname{pair}_{\N}}} & numerical pairing function & P4 turing-machines to decision-hierarchy & n/a & \hyperref[sec:turing-machines-01]{\ref*{sec:turing-machines-01}} & \thepage & Distinct from the set-theoretic ordered pair. \\
\phantomsection\label{entry:symbol-registry-00243}\(\operatorname{code}_k(\vec n)\) & \texttt{\detokenize{\operatorname{code}_k}} & fixed-arity tuple code & P4 turing-machines to decision-hierarchy & n/a & \hyperref[sec:turing-machines-01]{\ref*{sec:turing-machines-01}} & \thepage & Defined recursively from the numerical pairing function. \\
\phantomsection\label{entry:symbol-registry-00244}\(\chi_A\) & \texttt{\detokenize{\chi_A}} & characteristic function of a set & P4 turing-machines to frontiers-computability & n/a & \hyperref[sec:turing-machines-01]{\ref*{sec:turing-machines-01}} & \thepage & Source-bound notation: one exactly on members of \(A\). \\
\phantomsection\label{entry:symbol-registry-00245}\(\chi_R\) & \texttt{\detokenize{\chi_R}} & characteristic function of a relation & P4 turing-machines to frontiers-computability & n/a & \hyperref[sec:turing-machines-01]{\ref*{sec:turing-machines-01}} & \thepage & Source-bound notation: one exactly on tuples in \(R\). \\
\phantomsection\label{entry:symbol-registry-00246}\(\mu g\) & \texttt{\detokenize{\mu}} & unbounded minimization & P4 turing-machines & n/a & \hyperref[sec:turing-machines-01]{\ref*{sec:turing-machines-01}} & \thepage & Definedness-sensitive least-zero operation; relational least-true use is the explicit zero-test abbreviation in P4.1. \\
\phantomsection\label{entry:symbol-registry-00247}\(\mathsf T_k(e,\vec x,y)\) & \texttt{\detokenize{\mathsf T_k}} & primitive-recursive accepting-computation-code relation & P4 turing-machines & n/a & \hyperref[sec:turing-machines-01]{\ref*{sec:turing-machines-01}} & \thepage & Kleene normal-form relation. \\
\phantomsection\label{entry:symbol-registry-00248}\(\mathsf U(y)\) & \texttt{\detokenize{\mathsf U}} & primitive-recursive normal-form output extractor & P4 turing-machines & n/a & \hyperref[sec:turing-machines-01]{\ref*{sec:turing-machines-01}} & \thepage & Distinct from the universal program evaluator. \\
\phantomsection\label{entry:symbol-registry-00249}\(U(e,x)\) & \texttt{\detokenize{U}} & universal partial computable function & P4 universal-reductions to frontiers-computability & n/a & \hyperref[sec:universal-reductions-01]{\ref*{sec:universal-reductions-01}} & \thepage & Exact source notation; oracle superscript \(U^B\) gives relativization. \\
\phantomsection\label{entry:symbol-registry-00250}\(\mathsf{par}_n^m\) & \texttt{\detokenize{\mathsf{par}_n^m}} & parameter-specialization compiler & P4 universal-reductions to frontiers-computability & n/a & \hyperref[sec:universal-reductions-01]{\ref*{sec:universal-reductions-01}} & \thepage & Total computable compiler for the fixed numbering. \\
\phantomsection\label{entry:symbol-registry-00251}\(K\) & \texttt{\detokenize{K}} & self-application domain set & P4 universal-reductions to frontiers-computability & n/a & \hyperref[sec:universal-reductions-01]{\ref*{sec:universal-reductions-01}} & \thepage & Indices \(e\) for which \(\varphi_e(e)\) is defined. \\
\phantomsection\label{entry:symbol-registry-00252}\(H\) & \texttt{\detokenize{H}} & binary application domain set & P4 universal-reductions & n/a & \hyperref[sec:universal-reductions-01]{\ref*{sec:universal-reductions-01}} & \thepage & Paired indices and inputs with defined computation. \\
\phantomsection\label{entry:symbol-registry-00253}\(\Lang_A\) & \texttt{\detokenize{\Lang_A}} & first-order language of arithmetic & Parts 4–5 from decision-hierarchy & n/a & \hyperref[sec:decision-hierarchy-01]{\ref*{sec:decision-hierarchy-01}} & \thepage & Signature containing zero successor addition and multiplication. \\
\phantomsection\label{entry:symbol-registry-00254}\(F(\vec x)\mathord\downarrow\) & \texttt{\detokenize{\downarrow}} & convergence of a partial numerical expression & Legacy pre-batch5 P4.3-P4.5 reader meaning & meta & \hyperref[sec:decision-hierarchy-01]{\ref*{sec:decision-hierarchy-01}} & \thepage & Metalinguistic definedness predicate.  Numeric chapter scope migrated to its frozen pre-batch5 source-labelled reader regions; glyph, role and precedence unchanged. \\
\phantomsection\label{entry:symbol-registry-00255}\(F(\vec x)\mathord\uparrow\) & \texttt{\detokenize{\uparrow}} & divergence of a partial numerical expression & Legacy pre-batch5 P4.3-P4.5 reader meaning & meta & \hyperref[sec:decision-hierarchy-01]{\ref*{sec:decision-hierarchy-01}} & \thepage & Metalinguistic undefinedness predicate.  Numeric chapter scope migrated to its frozen pre-batch5 source-labelled reader regions; glyph, role and precedence unchanged. \\
\phantomsection\label{entry:symbol-registry-00256}\(\mathsf P(x,y,z)\) & \texttt{\detokenize{\mathsf P}} & bounded arithmetic pairing-graph formula & P4 decision-hierarchy & n/a & \hyperref[sec:decision-hierarchy-01]{\ref*{sec:decision-hierarchy-01}} & \thepage & Scope is disjoint from complexity class \(\mathsf P\). \\
\phantomsection\label{entry:symbol-registry-00257}\(\mathsf B(w,i,v)\) & \texttt{\detokenize{\mathsf B}} & bounded Goedel-beta graph formula & P4 decision-hierarchy & n/a & \hyperref[sec:decision-hierarchy-01]{\ref*{sec:decision-hierarchy-01}} & \thepage & Object-language arithmetic formula. \\
\phantomsection\label{entry:symbol-registry-00258}\(\beta(w,i)\) & \texttt{\detokenize{\beta}} & Goedel beta function & P4 decision-hierarchy & n/a & \hyperref[sec:decision-hierarchy-01]{\ref*{sec:decision-hierarchy-01}} & \thepage & Numerical finite-sequence decoder. \\
\phantomsection\label{entry:symbol-registry-00259}\(\mathsf{Comp}(e,x,s)\) & \texttt{\detokenize{\mathsf{Comp}}} & bounded accepting-history formula & P4 decision-hierarchy & n/a & \hyperref[sec:decision-hierarchy-01]{\ref*{sec:decision-hierarchy-01}} & \thepage & The explicit H1/H2/H3/H4 conjunction in def:p04-computation-formula, with capture-avoiding expansion of its bounded helpers; thm:bounded-verification proves exact raw-program history and convergence equivalences in the standard model. \\
\phantomsection\label{entry:symbol-registry-00260}\(\mathsf{TOT}\) & \texttt{\detokenize{\mathsf{TOT}}} & totality index set & P4 decision-hierarchy & n/a & \hyperref[sec:decision-hierarchy-01]{\ref*{sec:decision-hierarchy-01}} & \thepage & Indices of programs accepting with valid output on every input. \\
\phantomsection\label{entry:symbol-registry-00261}\(\mathcal C\) & \texttt{\detokenize{\mathcal C}} & locally declared class of sets or encoded languages & Legacy pre-batch5 P4.3-P4.6 reader meaning & n/a & \hyperref[sec:decision-hierarchy-01]{\ref*{sec:decision-hierarchy-01}} & \thepage & The active class and member type are declared locally.  Numeric chapter scope migrated to its frozen pre-batch5 source-labelled reader regions; glyph, role and precedence unchanged. \\
\phantomsection\label{entry:symbol-registry-00262}\(N!\) & \texttt{\detokenize{!}} & factorial & P4 decision-hierarchy & 1 & \hyperref[sec:decision-hierarchy-01]{\ref*{sec:decision-hierarchy-01}} & \thepage & Scope disjoint from the Part-8 exponential. \\
\phantomsection\label{entry:symbol-registry-00263}\(B^{(n)}\) & \texttt{\detokenize{B^{(n)}}} & n-fold Turing jump & Parts 4–5 from frontiers-computability & n/a & \hyperref[sec:p04-frontiers-computability]{\ref*{sec:p04-frontiers-computability}} & \thepage & Iteration from an arbitrary oracle set. \\
\phantomsection\label{entry:symbol-registry-00264}\(\equiv_m\) & \texttt{\detokenize{\equiv_m}} & many-one equivalence & P4 frontiers-computability & meta & \hyperref[sec:p04-frontiers-computability]{\ref*{sec:p04-frontiers-computability}} & \thepage & Exact source notation for mutual many-one reducibility. \\
\phantomsection\label{entry:symbol-registry-00265}\(\sim_T\) & \texttt{\detokenize{\sim_T}} & Turing equivalence & P4 frontiers-computability & meta & \hyperref[sec:p04-frontiers-computability]{\ref*{sec:p04-frontiers-computability}} & \thepage & Mutual Turing reducibility. \\
\phantomsection\label{entry:symbol-registry-00266}\(\deg_T(A)\) & \texttt{\detokenize{\deg_T}} & Turing degree of a set & P4 frontiers-computability & n/a & \hyperref[sec:p04-frontiers-computability]{\ref*{sec:p04-frontiers-computability}} & \thepage & Equivalence class under Turing equivalence. \\
\phantomsection\label{entry:symbol-registry-00267}\([A]_T\) & \texttt{\detokenize{[A]_T}} & Turing-equivalence class & P4 frontiers-computability & n/a & \hyperref[sec:p04-frontiers-computability]{\ref*{sec:p04-frontiers-computability}} & \thepage & Alternate display of the degree of \(A\). \\
\phantomsection\label{entry:symbol-registry-00268}\(A\oplus B\) & \texttt{\detokenize{\oplus}} & computable join representative & P4 frontiers-computability & n/a & \hyperref[sec:p04-frontiers-computability]{\ref*{sec:p04-frontiers-computability}} & \thepage & Even-odd coded union. \\
\phantomsection\label{entry:symbol-registry-00269}\(\mathbf0\) & \texttt{\detokenize{\mathbf0}} & computable Turing degree & P4 frontiers-computability & n/a & \hyperref[sec:p04-frontiers-computability]{\ref*{sec:p04-frontiers-computability}} & \thepage & Degree of the empty set. \\
\phantomsection\label{entry:symbol-registry-00270}\(\mathbf0'\) & \texttt{\detokenize{\mathbf0'}} & jump of the computable degree & P4 frontiers-computability & n/a & \hyperref[sec:p04-frontiers-computability]{\ref*{sec:p04-frontiers-computability}} & \thepage & Degree of the empty-set jump. \\
\phantomsection\label{entry:symbol-registry-00271}\(R_e\) & \texttt{\detokenize{R_e}} & priority requirement against computing A from B & P4 frontiers-computability & n/a & \hyperref[sec:p04-frontiers-computability]{\ref*{sec:p04-frontiers-computability}} & \thepage & Finite-injury requirement. \\
\phantomsection\label{entry:symbol-registry-00272}\(S_e\) & \texttt{\detokenize{S_e}} & priority requirement against computing B from A & P4 frontiers-computability & n/a & \hyperref[sec:p04-frontiers-computability]{\ref*{sec:p04-frontiers-computability}} & \thepage & Finite-injury requirement. \\
\phantomsection\label{entry:symbol-registry-00273}\(W_e\) & \texttt{\detokenize{W_e}} & c.e. set induced by program index e & P4 frontiers-computability & n/a & \hyperref[sec:p04-frontiers-computability]{\ref*{sec:p04-frontiers-computability}} & \thepage & Domain under the acceptable numbering. \\
\phantomsection\label{entry:symbol-registry-00274}\(\leq_1\) & \texttt{\detokenize{\leq_1}} & one-one reducibility & P4 frontiers-computability & meta & \hyperref[sec:p04-frontiers-computability]{\ref*{sec:p04-frontiers-computability}} & \thepage & The translator is total computable and injective. \\
\phantomsection\label{entry:symbol-registry-00275}\(D_d\) & \texttt{\detokenize{D_d}} & finite set decoded from d & P4 frontiers-computability & n/a & \hyperref[sec:p04-frontiers-computability]{\ref*{sec:p04-frontiers-computability}} & \thepage & Positions of one-bits in \(d\). \\
\phantomsection\label{entry:symbol-registry-00276}\(I_{\mathcal C}\) & \texttt{\detokenize{I_{\mathcal C}}} & extensional index set & P4 frontiers-computability & n/a & \hyperref[sec:p04-frontiers-computability]{\ref*{sec:p04-frontiers-computability}} & \thepage & Indices \(e\) such that \(W_e\) belongs to the declared class. \\
\phantomsection\label{entry:symbol-registry-00277}\(\mathsf{Fin}_{\mathcal C}\) & \texttt{\detokenize{\mathsf{Fin}_{\mathcal C}}} & c.e. set of finite-witness codes & P4 frontiers-computability & n/a & \hyperref[sec:p04-frontiers-computability]{\ref*{sec:p04-frontiers-computability}} & \thepage & Rice-Shapiro finite positive witnesses. \\
\phantomsection\label{entry:symbol-registry-00278}\(\alpha_i\) & \texttt{\detokenize{\alpha_i}} & function variable in the analytical hierarchy & Legacy pre-batch5 P4.5 reader meaning & n/a & \hyperref[sec:p04-complexity]{\ref*{sec:p04-complexity}} & \thepage & Scope disjoint from ordinal and type metavariables.  Numeric chapter scope migrated to its frozen pre-batch5 source-labelled reader regions; glyph, role and precedence unchanged. \\
\phantomsection\label{entry:symbol-registry-00279}\(\mathcal O\) & \texttt{\detokenize{\mathcal O}} & Kleene system of ordinal notations & Legacy pre-batch5 P4.5 reader meaning & n/a & \hyperref[sec:p04-complexity]{\ref*{sec:p04-complexity}} & \thepage & The cited standard system fixed in the method entry.  Numeric chapter scope migrated to its frozen pre-batch5 source-labelled reader regions; glyph, role and precedence unchanged. \\
\phantomsection\label{entry:symbol-registry-00280}\(<_{\mathcal O}\) & \texttt{\detokenize{<_{\mathcal O}}} & Kleene notation order & Legacy pre-batch5 P4.5 reader meaning & 2 & \hyperref[sec:p04-complexity]{\ref*{sec:p04-complexity}} & \thepage & Order belonging to the fixed cited notation system.  Numeric chapter scope migrated to its frozen pre-batch5 source-labelled reader regions; glyph, role and precedence unchanged. \\
\phantomsection\label{entry:symbol-registry-00281}\(H_a\) & \texttt{\detokenize{H_a}} & effective jump-hierarchy set & Legacy pre-batch5 P4.5 reader meaning & n/a & \hyperref[sec:p04-complexity]{\ref*{sec:p04-complexity}} & \thepage & Indexed by a Kleene-O notation.  Numeric chapter scope migrated to its frozen pre-batch5 source-labelled reader regions; glyph, role and precedence unchanged. \\
\phantomsection\label{entry:symbol-registry-00282}\(\mathsf{HYP}\) & \texttt{\detokenize{\mathsf{HYP}}} & hyperarithmetical subsets of the naturals & Legacy pre-batch5 P4.5 reader meaning & n/a & \hyperref[sec:p04-complexity]{\ref*{sec:p04-complexity}} & \thepage & Lightface class.  Numeric chapter scope migrated to its frozen pre-batch5 source-labelled reader regions; glyph, role and precedence unchanged. \\
\phantomsection\label{entry:symbol-registry-00283}\(\mathsf{LO}\) & \texttt{\detokenize{\mathsf{LO}}} & valid computable linear-order codes & Legacy pre-batch5 P4.5 reader meaning & n/a & \hyperref[sec:p04-complexity]{\ref*{sec:p04-complexity}} & \thepage & Universe within which \(\mathsf{WF}\) and \(\mathsf{IF}\) are complementary.  Numeric chapter scope migrated to its frozen pre-batch5 source-labelled reader regions; glyph, role and precedence unchanged. \\
\phantomsection\label{entry:symbol-registry-00284}\(R_e\) & \texttt{\detokenize{R_e}} & computable relation decoded from e & Legacy pre-batch5 P4.5 reader meaning & n/a & \hyperref[sec:p04-complexity]{\ref*{sec:p04-complexity}} & \thepage & Scope disjoint from the P4.4 priority-requirement row.  Numeric chapter scope migrated to its frozen pre-batch5 source-labelled reader regions; glyph, role and precedence unchanged. \\
\phantomsection\label{entry:symbol-registry-00285}\(\mathsf{WF}\) & \texttt{\detokenize{\mathsf{WF}}} & codes of computable well-orders & Legacy pre-batch5 P4.5 reader meaning & n/a & \hyperref[sec:p04-complexity]{\ref*{sec:p04-complexity}} & \thepage & Subset of \(\mathsf{LO}\).  Numeric chapter scope migrated to its frozen pre-batch5 source-labelled reader regions; glyph, role and precedence unchanged. \\
\phantomsection\label{entry:symbol-registry-00286}\(\mathsf{IF}\) & \texttt{\detokenize{\mathsf{IF}}} & codes of ill-founded computable linear orders & Legacy pre-batch5 P4.5 reader meaning & n/a & \hyperref[sec:p04-complexity]{\ref*{sec:p04-complexity}} & \thepage & Complement of \(\mathsf{WF}\) within \(\mathsf{LO}\).  Numeric chapter scope migrated to its frozen pre-batch5 source-labelled reader regions; glyph, role and precedence unchanged. \\
\phantomsection\label{entry:symbol-registry-00287}\(\mathsf U^1_{\Sigma_n}\) & \texttt{\detokenize{\mathsf U^1_{\Sigma_n}}} & universal lightface Sigma-one-n relation & Legacy pre-batch5 P4.5 reader meaning & n/a & \hyperref[sec:p04-complexity]{\ref*{sec:p04-complexity}} & \thepage & Pairs a valid definition code with an input satisfying that definition.  Numeric chapter scope migrated to its frozen pre-batch5 source-labelled reader regions; glyph, role and precedence unchanged. \\
\phantomsection\label{entry:symbol-registry-00288}\(\mathsf U^1_{\Pi_n}\) & \texttt{\detokenize{\mathsf U^1_{\Pi_n}}} & universal lightface Pi-one-n relation & Legacy pre-batch5 P4.5 reader meaning & n/a & \hyperref[sec:p04-complexity]{\ref*{sec:p04-complexity}} & \thepage & Complementary universal relation under the fixed coding.  Numeric chapter scope migrated to its frozen pre-batch5 source-labelled reader regions; glyph, role and precedence unchanged. \\
\phantomsection\label{entry:symbol-registry-00289}\(\operatorname{len}(x)\) & \texttt{\detokenize{\operatorname{len}}} & bit length of a finite binary string & Legacy pre-batch5 P4.5-P6.2 reader meaning & n/a & \hyperref[sec:p04-complexity]{\ref*{sec:p04-complexity}} & \thepage & Replaces forbidden vertical-bar length notation.  Numeric chapter scope migrated to its frozen pre-batch5 source-labelled reader regions; glyph, role and precedence unchanged. \\
\phantomsection\label{entry:symbol-registry-00290}\(C_V(x)\) & \texttt{\detokenize{C_V}} & plain Kolmogorov complexity & Legacy pre-batch5 P4.5 reader meaning & n/a & \hyperref[sec:p04-complexity]{\ref*{sec:p04-complexity}} & \thepage & Relative to decompressor \(V\).  Numeric chapter scope migrated to its frozen pre-batch5 source-labelled reader regions; glyph, role and precedence unchanged. \\
\phantomsection\label{entry:symbol-registry-00291}\(K_V(x)\) & \texttt{\detokenize{K_V}} & prefix-free Kolmogorov complexity & Legacy pre-batch5 P4.5 reader meaning & n/a & \hyperref[sec:p04-complexity]{\ref*{sec:p04-complexity}} & \thepage & Relative to prefix-free decompressor \(V\).  Numeric chapter scope migrated to its frozen pre-batch5 source-labelled reader regions; glyph, role and precedence unchanged. \\
\phantomsection\label{entry:symbol-registry-00292}\(U_C\) & \texttt{\detokenize{U_C}} & optimal plain decompressor & Legacy pre-batch5 P4.5 reader meaning & n/a & \hyperref[sec:p04-complexity]{\ref*{sec:p04-complexity}} & \thepage & Distinguished from theory metavariables.  Numeric chapter scope migrated to its frozen pre-batch5 source-labelled reader regions; glyph, role and precedence unchanged. \\
\phantomsection\label{entry:symbol-registry-00293}\(U_K\) & \texttt{\detokenize{U_K}} & optimal prefix-free machine & Legacy pre-batch5 P4.5 reader meaning & n/a & \hyperref[sec:p04-complexity]{\ref*{sec:p04-complexity}} & \thepage & Used throughout the Chaitin entries instead of bare \(U\).  Numeric chapter scope migrated to its frozen pre-batch5 source-labelled reader regions; glyph, role and precedence unchanged. \\
\phantomsection\label{entry:symbol-registry-00294}\(\Omega_{U_K}\) & \texttt{\detokenize{\Omega_{U_K}}} & halting probability of the fixed optimal prefix-free machine & Legacy pre-batch5 P4.5 reader meaning & n/a & \hyperref[sec:p04-complexity]{\ref*{sec:p04-complexity}} & \thepage & Not attached to an arbitrary plain machine.  Numeric chapter scope migrated to its frozen pre-batch5 source-labelled reader regions; glyph, role and precedence unchanged. \\
\phantomsection\label{entry:symbol-registry-00295}\(\mathsf{Run}_{U_K}(p,x,t)\) & \texttt{\detokenize{\mathsf{Run}_{U_K}}} & bounded prefix-machine computation formula & Legacy pre-batch5 P4.5 reader meaning & n/a & \hyperref[sec:p04-complexity]{\ref*{sec:p04-complexity}} & \thepage & Fixed-machine/string analogue of the inherited Comp formula; the finite trace codings are primitively recursively intertranslatable.  Numeric chapter scope migrated to its frozen pre-batch5 source-labelled reader regions; glyph, role and precedence unchanged. \\
\phantomsection\label{entry:symbol-registry-00296}\(\operatorname{BB}_{\mathrm{step}}(n)\) & \texttt{\detokenize{\operatorname{BB}_{\mathrm{step}}}} & Busy Beaver step function & Legacy pre-batch5 P4.5 reader meaning & n/a & \hyperref[sec:p04-complexity]{\ref*{sec:p04-complexity}} & \thepage & Maximum halting transition count under the displayed machine convention.  Numeric chapter scope migrated to its frozen pre-batch5 source-labelled reader regions; glyph, role and precedence unchanged. \\
\phantomsection\label{entry:symbol-registry-00297}\(\operatorname{BB}_{1}(n)\) & \texttt{\detokenize{\operatorname{BB}_{1}}} & Busy Beaver one-count function & Legacy pre-batch5 P4.5 reader meaning & n/a & \hyperref[sec:p04-complexity]{\ref*{sec:p04-complexity}} & \thepage & Maximum number of one-filled cells at halt under the displayed convention.  Numeric chapter scope migrated to its frozen pre-batch5 source-labelled reader regions; glyph, role and precedence unchanged. \\
\phantomsection\label{entry:symbol-registry-00298}\(A\leq_m^pB\) & \texttt{\detokenize{\leq_m^p}} & polynomial-time many-one reducibility & Legacy pre-batch5 P4.6 reader meaning & meta & \hyperref[sec:p04-decision-problem-hardness]{\ref*{sec:p04-decision-problem-hardness}} & \thepage & A total deterministic polynomial-time translator preserves membership.  Numeric chapter scope migrated to its frozen pre-batch5 source-labelled reader regions; glyph, role and precedence unchanged. \\
\phantomsection\label{entry:symbol-registry-00299}\(\mathsf P\) & \texttt{\detokenize{\mathsf P}} & deterministic polynomial-time language class & Legacy pre-batch5 P4.6 reader meaning & n/a & \hyperref[sec:p04-decision-problem-hardness]{\ref*{sec:p04-decision-problem-hardness}} & \thepage & Scope disjoint from the P4.3 pairing formula.  Numeric chapter scope migrated to its frozen pre-batch5 source-labelled reader regions; glyph, role and precedence unchanged. \\
\phantomsection\label{entry:symbol-registry-00300}\(\mathsf{SAT}\) & \texttt{\detokenize{\mathsf{SAT}}} & satisfiable propositional encodings & Legacy pre-batch5 P4.6 reader meaning & n/a & \hyperref[sec:p04-decision-problem-hardness]{\ref*{sec:p04-decision-problem-hardness}} & \thepage & Cook-Levin target under the declared encoding.  Numeric chapter scope migrated to its frozen pre-batch5 source-labelled reader regions; glyph, role and precedence unchanged. \\
\phantomsection\label{entry:symbol-registry-00301}\(3\mathsf{SAT}\) & \texttt{\detokenize{3\mathsf{SAT}}} & 3-CNF satisfiability language & Legacy pre-batch5 P4.6 reader meaning & n/a & \hyperref[sec:p04-decision-problem-hardness]{\ref*{sec:p04-decision-problem-hardness}} & \thepage & Encodings of satisfiable 3-CNF formulas whose clauses contain at most three literals.  Numeric chapter scope migrated to its frozen pre-batch5 source-labelled reader regions; glyph, role and precedence unchanged. \\
\phantomsection\label{entry:symbol-registry-00302}\(\mathsf{DTIME}(f(n))\) & \texttt{\detokenize{\mathsf{DTIME}}} & deterministic time-class constructor & Legacy pre-batch5 P4.6 reader meaning & n/a & \hyperref[sec:p04-decision-problem-hardness]{\ref*{sec:p04-decision-problem-hardness}} & \thepage & Under the fixed machine convention.  Numeric chapter scope migrated to its frozen pre-batch5 source-labelled reader regions; glyph, role and precedence unchanged. \\
\phantomsection\label{entry:symbol-registry-00303}\(\mathsf{NTIME}(f(n))\) & \texttt{\detokenize{\mathsf{NTIME}}} & nondeterministic time-class constructor & Legacy pre-batch5 P4.6 reader meaning & n/a & \hyperref[sec:p04-decision-problem-hardness]{\ref*{sec:p04-decision-problem-hardness}} & \thepage & Under the fixed machine convention.  Numeric chapter scope migrated to its frozen pre-batch5 source-labelled reader regions; glyph, role and precedence unchanged. \\
\phantomsection\label{entry:symbol-registry-00304}\(\mathsf{SPACE}(f(n))\) & \texttt{\detokenize{\mathsf{SPACE}}} & deterministic space-class constructor & Legacy pre-batch5 P4.6 reader meaning & n/a & \hyperref[sec:p04-decision-problem-hardness]{\ref*{sec:p04-decision-problem-hardness}} & \thepage & Counts work-tape cells.  Numeric chapter scope migrated to its frozen pre-batch5 source-labelled reader regions; glyph, role and precedence unchanged. \\
\phantomsection\label{entry:symbol-registry-00305}\(\mathsf{NSPACE}(f(n))\) & \texttt{\detokenize{\mathsf{NSPACE}}} & nondeterministic space-class constructor & Legacy pre-batch5 P4.6 reader meaning & n/a & \hyperref[sec:p04-decision-problem-hardness]{\ref*{sec:p04-decision-problem-hardness}} & \thepage & Counts work-tape cells.  Numeric chapter scope migrated to its frozen pre-batch5 source-labelled reader regions; glyph, role and precedence unchanged. \\
\phantomsection\label{entry:symbol-registry-00306}\(\mathsf{PSPACE}\) & \texttt{\detokenize{\mathsf{PSPACE}}} & polynomial-space language class & Legacy pre-batch5 P4.6 reader meaning & n/a & \hyperref[sec:p04-decision-problem-hardness]{\ref*{sec:p04-decision-problem-hardness}} & \thepage & Decision-language class.  Numeric chapter scope migrated to its frozen pre-batch5 source-labelled reader regions; glyph, role and precedence unchanged. \\
\phantomsection\label{entry:symbol-registry-00307}\(\mathsf{NPSPACE}\) & \texttt{\detokenize{\mathsf{NPSPACE}}} & nondeterministic polynomial-space language class & Legacy pre-batch5 P4.6 reader meaning & n/a & \hyperref[sec:p04-decision-problem-hardness]{\ref*{sec:p04-decision-problem-hardness}} & \thepage & Decision-language class.  Numeric chapter scope migrated to its frozen pre-batch5 source-labelled reader regions; glyph, role and precedence unchanged. \\
\phantomsection\label{entry:symbol-registry-00308}\(\mathsf{EXPSPACE}\) & \texttt{\detokenize{\mathsf{EXPSPACE}}} & exponential-space language class & Legacy pre-batch5 P4.6 reader meaning & n/a & \hyperref[sec:p04-decision-problem-hardness]{\ref*{sec:p04-decision-problem-hardness}} & \thepage & Decision-language class.  Numeric chapter scope migrated to its frozen pre-batch5 source-labelled reader regions; glyph, role and precedence unchanged. \\
\phantomsection\label{entry:symbol-registry-00309}\(O(f(n))\) & \texttt{\detokenize{O}} & asymptotic upper bound & Legacy pre-batch5 P4.6 reader meaning & n/a & \hyperref[sec:p04-decision-problem-hardness]{\ref*{sec:p04-decision-problem-hardness}} & \thepage & Resource-bound asymptotics.  Numeric chapter scope migrated to its frozen pre-batch5 source-labelled reader regions; glyph, role and precedence unchanged. \\
\phantomsection\label{entry:symbol-registry-00310}\(o(g(n))\) & \texttt{\detokenize{o}} & strict asymptotic domination & Legacy pre-batch5 P4.6 reader meaning & n/a & \hyperref[sec:p04-decision-problem-hardness]{\ref*{sec:p04-decision-problem-hardness}} & \thepage & Hierarchy-gap hypothesis.  Numeric chapter scope migrated to its frozen pre-batch5 source-labelled reader regions; glyph, role and precedence unchanged. \\
\phantomsection\label{entry:symbol-registry-00311}\(\Sigma^p_k\) & \texttt{\detokenize{\Sigma^p_k}} & existential-start polynomial-hierarchy class & Legacy pre-batch5 P4.6 reader meaning & n/a & \hyperref[sec:p04-decision-problem-hardness]{\ref*{sec:p04-decision-problem-hardness}} & \thepage & Bounded alternating string quantifiers.  Numeric chapter scope migrated to its frozen pre-batch5 source-labelled reader regions; glyph, role and precedence unchanged. \\
\phantomsection\label{entry:symbol-registry-00312}\(\Pi^p_k\) & \texttt{\detokenize{\Pi^p_k}} & universal-start polynomial-hierarchy class & Legacy pre-batch5 P4.6 reader meaning & n/a & \hyperref[sec:p04-decision-problem-hardness]{\ref*{sec:p04-decision-problem-hardness}} & \thepage & Bounded alternating string quantifiers.  Numeric chapter scope migrated to its frozen pre-batch5 source-labelled reader regions; glyph, role and precedence unchanged. \\
\phantomsection\label{entry:symbol-registry-00313}\(\mathsf{PH}\) & \texttt{\detokenize{\mathsf{PH}}} & polynomial hierarchy & Legacy pre-batch5 P4.6 reader meaning & n/a & \hyperref[sec:p04-decision-problem-hardness]{\ref*{sec:p04-decision-problem-hardness}} & \thepage & Union of its \(\Sigma^p_k\) levels.  Numeric chapter scope migrated to its frozen pre-batch5 source-labelled reader regions; glyph, role and precedence unchanged. \\
\phantomsection\label{entry:symbol-registry-00314}\(\mathsf{TQBF}\) & \texttt{\detokenize{\mathsf{TQBF}}} & true closed quantified-Boolean-formula language & Legacy pre-batch5 P4.6 reader meaning & n/a & \hyperref[sec:p04-decision-problem-hardness]{\ref*{sec:p04-decision-problem-hardness}} & \thepage & Unbounded alternation.  Numeric chapter scope migrated to its frozen pre-batch5 source-labelled reader regions; glyph, role and precedence unchanged. \\
\phantomsection\label{entry:symbol-registry-00315}\(\Sigma_k\text{-}\mathsf{SAT}\) & \texttt{\detokenize{\Sigma_k\text{-}\mathsf{SAT}}} & existential-start fixed-alternation QBF language & Legacy pre-batch5 P4.6 reader meaning & n/a & \hyperref[sec:p04-decision-problem-hardness]{\ref*{sec:p04-decision-problem-hardness}} & \thepage & Fixed \(k\).  Numeric chapter scope migrated to its frozen pre-batch5 source-labelled reader regions; glyph, role and precedence unchanged. \\
\phantomsection\label{entry:symbol-registry-00316}\(\Pi_k\text{-}\mathsf{SAT}\) & \texttt{\detokenize{\Pi_k\text{-}\mathsf{SAT}}} & universal-start fixed-alternation QBF language & Legacy pre-batch5 P4.6 reader meaning & n/a & \hyperref[sec:p04-decision-problem-hardness]{\ref*{sec:p04-decision-problem-hardness}} & \thepage & Fixed \(k\).  Numeric chapter scope migrated to its frozen pre-batch5 source-labelled reader regions; glyph, role and precedence unchanged. \\
\phantomsection\label{entry:symbol-registry-00317}\(\mathfrak L\) & \texttt{\detokenize{\mathfrak L}} & logic metavariable in a capture statement & Legacy pre-batch5 P4.6 reader meaning & n/a & \hyperref[sec:p04-decision-problem-hardness]{\ref*{sec:p04-decision-problem-hardness}} & \thepage & Distinct from \(\Lang\), which globally denotes a signature.  Numeric chapter scope migrated to its frozen pre-batch5 source-labelled reader regions; glyph, role and precedence unchanged. \\
\phantomsection\label{entry:symbol-registry-00318}\(A\leq_m^{\log}B\) & \texttt{\detokenize{\leq_m^{\log}}} & logspace many-one reducibility & Legacy pre-batch5 P4.6 reader meaning & meta & \hyperref[sec:p04-decision-problem-hardness]{\ref*{sec:p04-decision-problem-hardness}} & \thepage & A deterministic logarithmic-work-space translator preserves membership.  Numeric chapter scope migrated to its frozen pre-batch5 source-labelled reader regions; glyph, role and precedence unchanged. \\
\phantomsection\label{entry:symbol-registry-00319}\(\mathsf{CVP}\) & \texttt{\detokenize{\mathsf{CVP}}} & circuit value problem & Legacy pre-batch5 P4.6 reader meaning & n/a & \hyperref[sec:p04-decision-problem-hardness]{\ref*{sec:p04-decision-problem-hardness}} & \thepage & Valid acyclic Boolean circuit encodings with output value one.  Numeric chapter scope migrated to its frozen pre-batch5 source-labelled reader regions; glyph, role and precedence unchanged. \\
\phantomsection\label{entry:symbol-registry-00320}\(\mathsf{SCVP}\) & \texttt{\detokenize{\mathsf{SCVP}}} & succinct circuit value problem & Legacy pre-batch5 P4.6 reader meaning & n/a & \hyperref[sec:p04-decision-problem-hardness]{\ref*{sec:p04-decision-problem-hardness}} & \thepage & Circuit-described exponentially large circuit value instances.  Numeric chapter scope migrated to its frozen pre-batch5 source-labelled reader regions; glyph, role and precedence unchanged. \\
\phantomsection\label{entry:symbol-registry-00321}\(\mathsf{STQBF}\) & \texttt{\detokenize{\mathsf{STQBF}}} & succinct true quantified Boolean formula problem & Legacy pre-batch5 P4.6 reader meaning & n/a & \hyperref[sec:p04-decision-problem-hardness]{\ref*{sec:p04-decision-problem-hardness}} & \thepage & Circuit-described exponentially long closed QBF instances.  Numeric chapter scope migrated to its frozen pre-batch5 source-labelled reader regions; glyph, role and precedence unchanged. \\
\phantomsection\label{entry:symbol-registry-00322}\([\operatorname{lfp}_{R,\bar x}\varphi](\bar t)\) & \texttt{\detokenize{\operatorname{lfp}}} & least fixed-point formula constructor & Legacy pre-batch5 P4.6-P9 reader meaning & n/a & \hyperref[sec:p04-decision-problem-hardness]{\ref*{sec:p04-decision-problem-hardness}} & \thepage & Requires positive occurrence of the relation variable and denotes the least monotone fixed point.  Numeric chapter scope migrated to its frozen pre-batch5 source-labelled reader regions; glyph, role and precedence unchanged. \\
\phantomsection\label{entry:symbol-registry-00871}\(n,b,m,k\) & \texttt{\detokenize{n,}\allowbreak\detokenize{b,}\allowbreak\detokenize{m,}\allowbreak\detokenize{k}} & normalized program serialization counts & P4 universal-reductions to decision-hierarchy & meta & \hyperref[sec:universal-reductions-01]{\ref*{sec:universal-reductions-01}} & \thepage & State count n and tape-alphabet size b are at least three; m=(n-2)b is the full row count and k=3+m the metadata-plus-row tuple length in def:p04-program-serialization. \\
\phantomsection\label{entry:symbol-registry-00872}\(A,\eta_a\) & \texttt{\detokenize{A,}\allowbreak\detokenize{\eta_a}} & input-alphabet bitmap and its digits & P4 universal-reductions to decision-hierarchy & meta & \hyperref[sec:universal-reductions-01]{\ref*{sec:universal-reductions-01}} & \thepage & The b-bit bitmap A records normalized input-alphabet membership; bit zero excludes blank and bits one and two include the two numerical digits. Exponentiation in the numerical coding description is metatheoretic. \\
\phantomsection\label{entry:symbol-registry-00873}\(\rho_i\) & \texttt{\detokenize{\rho_i}} & serialized five-field transition row & P4 universal-reductions to decision-hierarchy & meta & \hyperref[sec:universal-reductions-01]{\ref*{sec:universal-reductions-01}} & \thepage & The right-nested \(\operatorname{code}_5\) of source state, read symbol, target state, written symbol and direction, in exact lexicographic key order. \\
\phantomsection\label{entry:symbol-registry-00874}\(\varepsilon\in\{0,1\}\) & \texttt{\detokenize{\varepsilon}} & numerical transition-direction code & P4 universal-reductions to decision-hierarchy & meta & \hyperref[sec:universal-reductions-01]{\ref*{sec:universal-reductions-01}} & \thepage & Zero encodes L and one encodes R in the normalized program serialization; this is a numerical field rather than an empty-word or epsilon-binder symbol. \\
\phantomsection\label{entry:symbol-registry-00875}\(\mathsf E,\mathsf E^{(2)}\) & \texttt{\detokenize{\mathsf E,}\allowbreak\detokenize{\mathsf E^{(2)}}} & bounded tagged-entry formulas & P4 decision-hierarchy & meta & \hyperref[sec:decision-hierarchy-01]{\ref*{sec:decision-hierarchy-01}} & \thepage & E(s,t,j,v) uses one position; E2(s,t,a,b,v) flattens two positions. Their pairing indices have explicit arithmetic bounds and their unique values are at most s; each family tag has fixed positional arity. \\
\phantomsection\label{entry:symbol-registry-00876}\(s[t,j],s[t,a,b]\) & \texttt{\detokenize{s[t,}\allowbreak\detokenize{j],}\allowbreak\detokenize{s[t,}\allowbreak\detokenize{a,}\allowbreak\detokenize{b]}} & numerical tagged-entry notation & P4 decision-hierarchy & meta & \hyperref[sec:decision-hierarchy-01]{\ref*{sec:decision-hierarchy-01}} & \thepage & Metatheoretic beta lookups from def:p04-tagged-history-entry; each lookup inside a formula is replaced by a value variable and E or E2, never by a new arithmetic function symbol. \\
\phantomsection\label{entry:symbol-registry-00877}\(\mathsf{Lt},\mathsf{Le},\mathsf{Bit}\) & \texttt{\detokenize{\mathsf{Lt},}\allowbreak\detokenize{\mathsf{Le},}\allowbreak\detokenize{\mathsf{Bit}}} & bounded comparison and binary-digit formulas & P4 decision-hierarchy & meta & \hyperref[sec:decision-hierarchy-01]{\ref*{sec:decision-hierarchy-01}} & \thepage & Eliminable arithmetic formulas for strict comparison, non-strict comparison and membership in zero/one; matrix comparisons use marked bounded witnesses. \\
\phantomsection\label{entry:symbol-registry-00878}\(\mathsf H_1,\mathsf H_2,\mathsf H_3,\mathsf H_4\) & \texttt{\detokenize{\mathsf H_1,}\allowbreak\detokenize{\mathsf H_2,}\allowbreak\detokenize{\mathsf H_3,}\allowbreak\detokenize{\mathsf H_4}} & four bounded computation-certificate clauses & P4 decision-hierarchy & meta & \hyperref[sec:decision-hierarchy-01]{\ref*{sec:decision-hierarchy-01}} & \thepage & H1(x,s) checks canonical input; H2(e,s) exact program/table data; H3(s) configurations and transitions; H4(s) accepting canonical output. Their conjunction defines Comp. \\
\phantomsection\label{entry:symbol-registry-00879}\(\mathsf P_5,\mathsf R_5\) & \texttt{\detokenize{\mathsf P_5,}\allowbreak\detokenize{\mathsf R_5}} & five-field pairing and stored-row formulas & P4 decision-hierarchy & meta & \hyperref[sec:decision-hierarchy-01]{\ref*{sec:decision-hierarchy-01}} & \thepage & P5 unpacks a right-nested five-field tuple; R5 extracts row i from description position 3+i. Validity and field ranges require H2 and i<m; a padding value alone is not a transition. \\
\phantomsection\label{entry:symbol-registry-00880}\(\mathsf I,\mathsf J\) & \texttt{\detokenize{\mathsf I,}\allowbreak\detokenize{\mathsf J}} & input-cell and successor-prefix checks & P4 decision-hierarchy & meta & \hyperref[sec:decision-hierarchy-01]{\ref*{sec:decision-hierarchy-01}} & \thepage & I(x,s,L) checks all input digits/prefix bounds; J(s,lambda) checks the binary prefix recurrence. Both are eliminable helpers of H1. \\
\phantomsection\label{entry:symbol-registry-00881}\(\mathsf D,\mathsf M,\mathsf K\) & \texttt{\detokenize{\mathsf D,}\allowbreak\detokenize{\mathsf M,}\allowbreak\detokenize{\mathsf K}} & description-chain bitmap and table checks & P4 decision-hierarchy & meta & \hyperref[sec:decision-hierarchy-01]{\ref*{sec:decision-hierarchy-01}} & \thepage & Bounded H2 helpers: D checks the full right-nested decoding chain, M the finite-width bitmap with its tail values, and K every ordered nonhalting transition row. \\
\phantomsection\label{entry:symbol-registry-00882}\(\mathsf{Cfg},\mathsf{Init},\mathsf{Step}\) & \texttt{\detokenize{\mathsf{Cfg},}\allowbreak\detokenize{\mathsf{Init},}\allowbreak\detokenize{\mathsf{Step}}} & configuration initialization and transition formulas & P4 decision-hierarchy & meta & \hyperref[sec:decision-hierarchy-01]{\ref*{sec:decision-hierarchy-01}} & \thepage & Bounded H3 helpers enforce state/head/symbol ranges and tape prefixes, the canonical initial tape, and both head moves with exactly one overwritten cell and all other columns preserved. \\
\phantomsection\label{entry:symbol-registry-00883}\(\mathsf{OutCells},\mathsf{OutNext}\) & \texttt{\detokenize{\mathsf{OutCells},}\allowbreak\detokenize{\mathsf{OutNext}}} & canonical-output cell and prefix formulas & P4 decision-hierarchy & meta & \hyperref[sec:decision-hierarchy-01]{\ref*{sec:decision-hierarchy-01}} & \thepage & Bounded H4 helpers link each output bit to its normalized final-tape symbol and verify the complete binary prefix recurrence. \\
\phantomsection\label{entry:symbol-registry-00884}\(L,T,H,W,\ell,y\) & \texttt{\detokenize{L,}\allowbreak\detokenize{T,}\allowbreak\detokenize{H,}\allowbreak\detokenize{W,}\allowbreak\detokenize{\ell,}\allowbreak\detokenize{y}} & computation-history size origin and output fields & P4 decision-hierarchy & meta & \hyperref[sec:decision-hierarchy-01]{\ref*{sec:decision-hierarchy-01}} & \thepage & Input length L, transition count T, origin H=T, width W=L+2T+1, canonical output length ell and value y are header entries; shifted column j represents physical cell j-H. \\
\phantomsection\label{entry:symbol-registry-00885}\(q_t,h_t,a_{t,j},u_{t,j},z_t\) & \texttt{\detokenize{q_t,}\allowbreak\detokenize{h_t,}\allowbreak\detokenize{a_{t,j},}\allowbreak\detokenize{u_{t,j},}\allowbreak\detokenize{z_t}} & stored configuration and tape data & P4 decision-hierarchy & meta & \hyperref[sec:decision-hierarchy-01]{\ref*{sec:decision-hierarchy-01}} & \thepage & State, shifted head, tape symbol, base-b prefix value and complete fixed-width tape code, respectively, with family tags 7 through 11. \\
\phantomsection\label{entry:symbol-registry-00886}\(d_j,v_j,o_j,w_j,B_j\) & \texttt{\detokenize{d_j,}\allowbreak\detokenize{v_j,}\allowbreak\detokenize{o_j,}\allowbreak\detokenize{w_j,}\allowbreak\detokenize{B_j}} & input output and bitmap expansion data & P4 decision-hierarchy & meta & \hyperref[sec:decision-hierarchy-01]{\ref*{sec:decision-hierarchy-01}} & \thepage & Input bits/prefixes use tags 5/6, output bits/prefixes tags 12/13, and bitmap tail values tag 4. Canonical zero and positive-word conditions are explicit in H1 and H4. \\
\phantomsection\label{entry:symbol-registry-00887}\(\lambda,\mu,\nu\) & \texttt{\detokenize{\lambda,}\allowbreak\detokenize{\mu,}\allowbreak\detokenize{\nu}} & local position and nonhalting-state witnesses & P4 decision-hierarchy & meta & \hyperref[sec:decision-hierarchy-01]{\ref*{sec:decision-hierarchy-01}} & \thepage & Lambda denotes the last input, description or output position under separately bound scopes; mu is the last tape column and nu=n-2 the number of nonhalting states. Every use has an explicit arithmetic bound. \\
\phantomsection\label{entry:symbol-registry-03054}\(\beta(c,d,i)\) & \texttt{\detokenize{\beta}} & arithmetical finite-sequence remainder decoder & P4 delta-one hyperarithmetical construction & n/a & \hyperref[annex:p04.doh.arithmetic-graphs]{\ref*{annex:p04.doh.arithmetic-graphs}} & \thepage & The modulus is one plus (i+1)d; the CRT construction supplies codes for every finite sequence and the empty case is vacuous. \\
\phantomsection\label{entry:symbol-registry-03055}\(\bigoplus_i B_i\) & \texttt{\detokenize{\bigoplus}} & effective numerical join of sets & P4 delta-one hyperarithmetical construction & n/a & \hyperref[def:p04-doh-programs]{\ref*{def:p04-doh-programs}} & \thepage & The first coordinate of the fixed pairing chooses the oracle column. \\
\phantomsection\label{entry:symbol-registry-03056}\(\mathsf s(a)\) & \texttt{\detokenize{\mathsf s}} & successor-notation constructor & P4 delta-one hyperarithmetical construction & n/a & \hyperref[def:p04-doh-notations]{\ref*{def:p04-doh-notations}} & \thepage & The numerical code is two to the power a; this is not ordinal exponentiation. \\
\phantomsection\label{entry:symbol-registry-03057}\(\mathsf l(e)\) & \texttt{\detokenize{\mathsf l}} & limit-notation constructor & P4 delta-one hyperarithmetical construction & n/a & \hyperref[def:p04-doh-notations]{\ref*{def:p04-doh-notations}} & \thepage & The numerical code is three times five to the power e; the indexed sequence must be total and strictly increasing in notation order. \\
\phantomsection\label{entry:symbol-registry-03058}\(\leq_{\mathcal O}\) & \texttt{\detokenize{\leq_{\mathcal O}}} & reflexive notation predecessor relation & P4 delta-one hyperarithmetical construction & 2 & \hyperref[def:p04-doh-notations]{\ref*{def:p04-doh-notations}} & \thepage & Includes equality; unlike denoted ordinal comparison it retains the particular notation paths. \\
\phantomsection\label{entry:symbol-registry-03059}\(\nu_{\mathcal O}(a)\) & \texttt{\detokenize{\nu_{\mathcal O}}} & ordinal denotation of a valid notation & P4 delta-one hyperarithmetical construction & n/a & \hyperref[def:p04-doh-notations]{\ref*{def:p04-doh-notations}} & \thepage & No ordinal value is supplied on invalid numerical codes. \\
\phantomsection\label{entry:symbol-registry-03060}\(\operatorname{add}_{\mathcal O}(a,b)\) & \texttt{\detokenize{\operatorname{add}_{\mathcal O}}} & syntactic ordinal-notation addition compiler & P4 delta-one hyperarithmetical construction & n/a & \hyperref[def:p04-doh-addition]{\ref*{def:p04-doh-addition}} & \thepage & Total code compiler on all numerical inputs; validity and ordinal-addition claims hold on valid inputs only. \\
\phantomsection\label{entry:symbol-registry-03061}\(r(a,b)\) & \texttt{\detokenize{r}} & right-summand oracle-reduction index & P4 delta-one hyperarithmetical construction & n/a & \hyperref[thm:p04-doh-addition]{\ref*{thm:p04-doh-addition}} & \thepage & Uniformly reads the b hierarchy set from the hierarchy set at the compiled addition notation. \\
\phantomsection\label{entry:symbol-registry-03062}\(\mathsf C(z,i)\) & \texttt{\detokenize{\mathsf C}} & decidable numerical child relation & P4 delta-one hyperarithmetical construction & n/a & \hyperref[def:p04-doh-node-bounds]{\ref*{def:p04-doh-node-bounds}} & \thepage & No infinite child sequence is the promise used by the validity proof. \\
\phantomsection\label{entry:symbol-registry-03063}\(h(z)\) & \texttt{\detokenize{h}} & computable node-notation generator & P4 delta-one hyperarithmetical construction & n/a & \hyperref[def:p04-doh-node-bounds]{\ref*{def:p04-doh-node-bounds}} & \thepage & Top-level output builds two successors of a limit code without evaluating the child hierarchy. \\
\phantomsection\label{entry:symbol-registry-03064}\(a_{z,i}\) & \texttt{\detokenize{a_{z,i}}} & cumulative valid-notation sequence at a node & P4 delta-one hyperarithmetical construction & n/a & \hyperref[def:p04-doh-node-bounds]{\ref*{def:p04-doh-node-bounds}} & \thepage & The step adds a child notation to the successor of the preceding cumulative notation. \\
\phantomsection\label{entry:symbol-registry-03065}\(b_{z,i}\) & \texttt{\detokenize{b_{z,i}}} & child notation or base dummy & P4 delta-one hyperarithmetical construction & n/a & \hyperref[def:p04-doh-node-bounds]{\ref*{def:p04-doh-node-bounds}} & \thepage & Equals h(i) for an actual child and one otherwise. \\
\phantomsection\label{entry:symbol-registry-03066}\(e_z\) & \texttt{\detokenize{e_z}} & index of the cumulative node sequence & P4 delta-one hyperarithmetical construction & n/a & \hyperref[def:p04-doh-node-bounds]{\ref*{def:p04-doh-node-bounds}} & \thepage & Constructed as a finite program specialization; distinct from the ordinal denotation. \\
\phantomsection\label{entry:symbol-registry-03067}\(J_z\) & \texttt{\detokenize{J_z}} & node limit-join oracle & P4 delta-one hyperarithmetical construction & n/a & \hyperref[thm:p04-doh-node-bounds]{\ref*{thm:p04-doh-node-bounds}} & \thepage & The hierarchy set at the limit before the two successor jumps; column i+1 carries the oracle used for child i. \\
\phantomsection\label{entry:symbol-registry-03068}\(T_n,U_n\) & \texttt{\detokenize{T_n;U_n}} & recursive trees for complementary analytical definitions & P4 delta-one hyperarithmetical construction & n/a & \hyperref[def:p04-doh-alternation]{\ref*{def:p04-doh-alternation}} & \thepage & A single synchronized product forest combines all inputs; complementarity is a correctness promise rather than an effective test. \\
\phantomsection\label{entry:symbol-registry-03069}\(\vartheta(s,t)\) & \texttt{\detokenize{\vartheta}} & alternating well-founded product-tree bit & P4 delta-one hyperarithmetical construction & n/a & \hyperref[def:p04-doh-alternation]{\ref*{def:p04-doh-alternation}} & \thepage & Universal first-tree children and existential second-tree children; empty universal is true and empty existential is false. \\
\phantomsection\label{entry:symbol-registry-03070}\(\operatorname{node}(n,s,t)\) & \texttt{\detokenize{\operatorname{node}}} & numerical product-node code & P4 delta-one hyperarithmetical construction & n/a & \hyperref[def:p04-doh-alternation]{\ref*{def:p04-doh-alternation}} & \thepage & Positive tagged triple code reserves zero for the common root. \\
\phantomsection\label{entry:symbol-registry-03071}\(*\) & \texttt{\detokenize{*}} & common root of the product forest & P4 delta-one hyperarithmetical construction & n/a & \hyperref[def:p04-doh-alternation]{\ref*{def:p04-doh-alternation}} & \thepage & The numerical code zero; each input root is an immediate child. \\
\phantomsection\label{entry:symbol-registry-03072}\(P_z(m,k)\) & \texttt{\detokenize{P_z}} & oracle-decidable alternating matrix & P4 delta-one hyperarithmetical construction & n/a & \hyperref[annex:p04.doh.alternation]{\ref*{annex:p04.doh.alternation}} & \thepage & Child evaluation is simulated only when both child tests hold; well-founded induction proves totality of the matrix. \\
\phantomsection\label{entry:symbol-registry-03073}\(p_{z,m}\) & \texttt{\detokenize{p_{z,m}}} & first-jump existential-search program & P4 delta-one hyperarithmetical construction & n/a & \hyperref[annex:p04.doh.alternation]{\ref*{annex:p04.doh.alternation}} & \thepage & Halts exactly when some second-child choice gives the matrix value one. \\
\phantomsection\label{entry:symbol-registry-03074}\(q_z\) & \texttt{\detokenize{q_z}} & second-jump counterexample-search program & P4 delta-one hyperarithmetical construction & n/a & \hyperref[annex:p04.doh.alternation]{\ref*{annex:p04.doh.alternation}} & \thepage & With the first jump as oracle searches for a first-child choice with no successful second choice. \\
\phantomsection\label{entry:symbol-registry-03075}\(V_a\) & \texttt{\detokenize{V_a}} & recursive unfolding tree of a fixed valid notation & P4 delta-one hyperarithmetical construction & n/a & \hyperref[annex:p04.doh.hyp-to-delta]{\ref*{annex:p04.doh.hyp-to-delta}} & \thepage & Membership and labels are arithmetical uniformly in prospective codes; totality and well-foundedness are proved for valid a. \\
\phantomsection\label{entry:symbol-registry-03076}\(Y_u\) & \texttt{\detokenize{Y_u}} & hierarchy column at an unfolding-tree node & P4 delta-one hyperarithmetical construction & n/a & \hyperref[annex:p04.doh.hyp-to-delta]{\ref*{annex:p04.doh.hyp-to-delta}} & \thepage & A single subset of the natural numbers codes all these columns through pairing. \\
\phantomsection\label{entry:symbol-registry-03077}\(\operatorname{Hier}(a,Y)\) & \texttt{\detokenize{\operatorname{Hier}}} & arithmetical specification of the unfolded jump hierarchy & P4 delta-one hyperarithmetical construction & n/a & \hyperref[annex:p04.doh.hyp-to-delta]{\ref*{annex:p04.doh.hyp-to-delta}} & \thepage & Existence and uniqueness follow from well-founded set recursion; no set quantifier is hidden in its matrix. \\
\phantomsection\label{entry:symbol-registry-03078}\(s(p,t)\) & \texttt{\detokenize{s}} & computable program-parameter insertion & P4 delta-one hyperarithmetical construction & n/a & \hyperref[annex:p04.doh.program-tools]{\ref*{annex:p04.doh.program-tools}} & \thepage & Lowercase s is a code-specialization function; distinct from the sans-serif successor constructor. \\
\phantomsection\label{entry:symbol-registry-03309}\(\mathsf{NL}\) & \texttt{\detokenize{\mathsf{NL}}} & nondeterministic logarithmic-work-space language class & P04 decision hardness resources & n/a & \hyperref[def:p04-dh-nl]{\ref*{def:p04-dh-nl}} & \thepage & Read-only input; every computation branch halts; work space O(log(n+2)). \\
\phantomsection\label{entry:symbol-registry-03310}\(\mathsf{coNL}\) & \texttt{\detokenize{\mathsf{coNL}}} & complement class of NL on the whole binary input space & P04 decision hardness resources & n/a & \hyperref[def:p04-dh-nl]{\ref*{def:p04-dh-nl}} & \thepage & Not merely complement on syntactically valid graph encodings. \\
\phantomsection\label{entry:symbol-registry-03311}\(\mathsf{STCON}\) & \texttt{\detokenize{\mathsf{STCON}}} & valid explicit graph instances with directed s-to-t reachability & P04 decision hardness graphs & n/a & \hyperref[def:p04-dh-graphs]{\ref*{def:p04-dh-graphs}} & \thepage & Includes paths of length zero; malformed inputs are no instances. \\
\phantomsection\label{entry:symbol-registry-03312}\(\mathsf{NONREACH}\) & \texttt{\detokenize{\mathsf{NONREACH}}} & valid explicit graph instances without directed s-to-t reachability & P04 decision hardness graphs & n/a & \hyperref[def:p04-dh-graphs]{\ref*{def:p04-dh-graphs}} & \thepage & Both graph languages reject malformed inputs. \\
\phantomsection\label{entry:symbol-registry-03313}\(N\) & \texttt{\detokenize{N}} & explicit graph vertex count & P04 decision hardness graphs & n/a & \hyperref[def:p04-dh-graphs]{\ref*{def:p04-dh-graphs}} & \thepage & Positive count supplied in unary; matrix has N squared bits. \\
\phantomsection\label{entry:symbol-registry-03314}\(R_d\) & \texttt{\detokenize{R_d}} & vertices reachable from s within d directed edges & P04 decision hardness graphs & n/a & \hyperref[def:p04-dh-counting]{\ref*{def:p04-dh-counting}} & \thepage & Length-zero paths make R\_0 the singleton start vertex. \\
\phantomsection\label{entry:symbol-registry-03315}\(c_d\) & \texttt{\detokenize{c_d}} & exact cardinality of the bounded reachable set & P04 decision hardness graphs & n/a & \hyperref[def:p04-dh-counting]{\ref*{def:p04-dh-counting}} & \thepage & The checked inclusion census prevents omission of a reachable vertex. \\
\phantomsection\label{entry:symbol-registry-03316}\(\mathsf{2SAT}\) & \texttt{\detokenize{\mathsf{2SAT}}} & well-formed satisfiable explicit 2-CNF language & P04 decision hardness CNF & n/a & \hyperref[def:p04-dh-two-cnf]{\ref*{def:p04-dh-two-cnf}} & \thepage & Malformed inputs are no instances; empty clause lists are yes instances. \\
\phantomsection\label{entry:symbol-registry-03317}\(r\) & \texttt{\detokenize{r}} & declared Boolean variable count & P04 decision hardness CNF & n/a & \hyperref[def:p04-dh-two-cnf]{\ref*{def:p04-dh-two-cnf}} & \thepage & Unary declaration; unused variables allowed; r=0 allowed. \\
\phantomsection\label{entry:symbol-registry-03318}\(F_{G,s,t}\) & \texttt{\detokenize{F_{G,s,t}}} & edge-implication 2-CNF of the reachability reduction & P04 decision hardness CNF & n/a & \hyperref[sec:p04-dh-two-cnf]{\ref*{sec:p04-dh-two-cnf}} & \thepage & Satisfiable precisely for valid NONREACH inputs. \\
\phantomsection\label{entry:symbol-registry-03319}\(\Sigma_{\mathrm H}\) & \texttt{\detokenize{\Sigma_{\mathrm H}}} & fixed effective countable function-and-relation signature & P04 decision hardness validity & n/a & \hyperref[def:p04-dh-validity]{\ref*{def:p04-dh-validity}} & \thepage & Each encoded sentence uses only finitely many of its indexed symbols. \\
\phantomsection\label{entry:symbol-registry-03320}\(\mathsf{VALID}_{\mathrm H}\) & \texttt{\detokenize{\mathsf{VALID}_{\mathrm H}}} & unrestricted nonempty-structure first-order validity code set & P04 decision hardness validity & n/a & \hyperref[def:p04-dh-validity]{\ref*{def:p04-dh-validity}} & \thepage & Malformed codes and open formulas are no instances; not finite validity. \\
\phantomsection\label{entry:symbol-registry-03321}\(\mathsf{nil}\) & \texttt{\detokenize{\mathsf{nil}}} & empty-list constant in the machine compilation & P04 decision hardness validity & n/a & \hyperref[def:p04-dh-validity]{\ref*{def:p04-dh-validity}} & \thepage & Also ensures a nonempty free ground-term universe. \\
\phantomsection\label{entry:symbol-registry-03322}\(\mathsf{Reach}\) & \texttt{\detokenize{\mathsf{Reach}}} & unary object-language predicate on configuration terms & P04 decision hardness validity & n/a & \hyperref[def:p04-dh-validity]{\ref*{def:p04-dh-validity}} & \thepage & Not a newly postulated first-order transitive-closure operator. \\
\phantomsection\label{entry:symbol-registry-03323}\(\mathsf a_i\) & \texttt{\detokenize{\mathsf a_i}} & indexed unary constructor symbol family & P04 decision hardness validity & n/a & \hyperref[def:p04-dh-validity]{\ref*{def:p04-dh-validity}} & \thepage & Finite alphabet symbols are assigned to distinct family members. \\
\phantomsection\label{entry:symbol-registry-03324}\(\mathsf c_j\) & \texttt{\detokenize{\mathsf c_j}} & indexed binary constructor symbol family & P04 decision hardness validity & n/a & \hyperref[def:p04-dh-validity]{\ref*{def:p04-dh-validity}} & \thepage & Finite machine states are assigned to distinct family members. \\
\phantomsection\label{entry:symbol-registry-03325}\(\gamma(i)\) & \texttt{\detokenize{\gamma(i)}} & self-delimiting nonnegative index field & P04 decision hardness validity & n/a & \hyperref[def:p04-dh-validity]{\ref*{def:p04-dh-validity}} & \thepage & Unary digit-length prefix followed by the canonical binary string of i+1. \\
\phantomsection\label{entry:symbol-registry-03326}\(\operatorname{code}(u)\) & \texttt{\detokenize{\operatorname{code}(u)}} & natural-number code of the encoded binary sentence word & P04 decision hardness validity & n/a & \hyperref[def:p04-dh-validity]{\ref*{def:p04-dh-validity}} & \thepage & Equal to 2 to word length minus one plus its binary value; not program numbering. \\
\phantomsection\label{entry:symbol-registry-03327}\(\ell\) & \texttt{\detokenize{\ell}} & left-of-head finite tape list, nearest symbol first & P04 decision hardness validity & n/a & \hyperref[def:p04-dh-machine-encoding]{\ref*{def:p04-dh-machine-encoding}} & \thepage & Missing positions are blank; explicit blank entries are permitted. \\
\phantomsection\label{entry:symbol-registry-03328}\(\rho\) & \texttt{\detokenize{\rho}} & scanned-and-right finite tape list & P04 decision hardness validity & n/a & \hyperref[def:p04-dh-machine-encoding]{\ref*{def:p04-dh-machine-encoding}} & \thepage & Empty list scans a blank; nonempty list exposes its first symbol. \\
\phantomsection\label{entry:symbol-registry-03329}\([w]\) & \texttt{\detokenize{[w]}} & nested constructor term of a finite binary input word & P04 decision hardness validity & n/a & \hyperref[def:p04-dh-machine-encoding]{\ref*{def:p04-dh-machine-encoding}} & \thepage & The empty word is nil; the numeral for zero is the nonempty word 0. \\
\phantomsection\label{entry:symbol-registry-03330}\(T_M\) & \texttt{\detokenize{T_M}} & finite conjunction of closed Horn transition clauses & P04 decision hardness validity & n/a & \hyperref[def:p04-dh-horn-compilation]{\ref*{def:p04-dh-horn-compilation}} & \thepage & All clause variables are universally bound before conjunction. \\
\phantomsection\label{entry:symbol-registry-03331}\(I_{M,w}\) & \texttt{\detokenize{I_{M,w}}} & ground initial configuration Reach atom & P04 decision hardness validity & n/a & \hyperref[def:p04-dh-horn-compilation]{\ref*{def:p04-dh-horn-compilation}} & \thepage & Initial left list nil and right list the encoded input word. \\
\phantomsection\label{entry:symbol-registry-03332}\(H_M\) & \texttt{\detokenize{H_M}} & existential accepting-configuration sentence & P04 decision hardness validity & n/a & \hyperref[def:p04-dh-horn-compilation]{\ref*{def:p04-dh-horn-compilation}} & \thepage & Two list witnesses are arbitrary domain elements, including in the countermodel. \\
\phantomsection\label{entry:symbol-registry-03333}\(\Phi_{M,w}\) & \texttt{\detokenize{\Phi_{M,w}}} & closed machine-acceptance validity instance & P04 decision hardness validity & n/a & \hyperref[def:p04-dh-horn-compilation]{\ref*{def:p04-dh-horn-compilation}} & \thepage & Implication from transition conjunction and initial atom to accepting existence. \\
\phantomsection\label{entry:symbol-registry-03334}\(h(C)\) & \texttt{\detokenize{h(C)}} & topological rank of an implication-graph SCC & P04 decision hardness SCC proof & n/a & \hyperref[annex:p04.dh.two-sat]{\ref*{annex:p04.dh.two-sat}} & \thepage & Increases on edges between distinct SCCs; supplies the satisfying assignment. \\
\phantomsection\label{entry:symbol-registry-03335}\(C_L\) & \texttt{\detokenize{C_L}} & SCC of literal L in the implication graph & P04 decision hardness SCC proof & n/a & \hyperref[annex:p04.dh.two-sat]{\ref*{annex:p04.dh.two-sat}} & \thepage & Literal complementation reverses the condensation edges. \\
\phantomsection\label{entry:symbol-registry-03336}\(f_A\) & \texttt{\detokenize{f_A}} & total computable reduction for a fixed c.e. set A & P04 decision hardness reduction proof & n/a & \hyperref[annex:p04.dh.validity]{\ref*{annex:p04.dh.validity}} & \thepage & Compiles a fixed recognizer and canonical binary input without running to acceptance. \\
\phantomsection\label{entry:symbol-registry-03337}\(G_f\) & \texttt{\detokenize{G_f}} & relational graph predicate for a used function symbol & P04 decision hardness relational proof & n/a & \hyperref[annex:p04.dh.relational]{\ref*{annex:p04.dh.relational}} & \thepage & Binary for unary functions; ternary for binary functions. \\
\phantomsection\label{entry:symbol-registry-03338}\(J\) & \texttt{\detokenize{J}} & closed conjunction of totality and uniqueness graph axioms & P04 decision hardness relational proof & n/a & \hyperref[annex:p04.dh.relational]{\ref*{annex:p04.dh.relational}} & \thepage & Contains no constant-replacement parameter. \\
\phantomsection\label{entry:symbol-registry-03339}\(V_t(y)\) & \texttt{\detokenize{V_t(y)}} & first-order term-value translation formula & P04 decision hardness relational proof & n/a & \hyperref[annex:p04.dh.relational]{\ref*{annex:p04.dh.relational}} & \thepage & Generated variables are fresh; graph axioms imply existence and uniqueness of the value. \\
\phantomsection\label{entry:symbol-registry-03340}\(z_0\) & \texttt{\detokenize{z_0}} & fresh variable replacing the source constant & P04 decision hardness relational proof & n/a & \hyperref[annex:p04.dh.relational]{\ref*{annex:p04.dh.relational}} & \thepage & Universally bound outside the graph-axiom implication. \\
\phantomsection\label{entry:symbol-registry-03341}\(\Psi\) & \texttt{\detokenize{\Psi}} & closed purely relational validity translation & P04 decision hardness relational proof & n/a & \hyperref[annex:p04.dh.relational]{\ref*{annex:p04.dh.relational}} & \thepage & Equivalent unrestricted validity in both model directions. \\
\phantomsection\label{entry:symbol-registry-03342}\(\mathsf L\) & \texttt{\detokenize{\mathsf L}} & deterministic logarithmic-work-space language class & P04 decision hardness resources & n/a & \hyperref[def:p04-dh-nl]{\ref*{def:p04-dh-nl}} & \thepage & Distinct from the machine move code with the same glyph in its disjoint scope. \\
\phantomsection\label{entry:symbol-registry-03343}\(G\) & \texttt{\detokenize{G}} & explicit finite directed input graph & P04 decision hardness graphs & n/a & \hyperref[def:p04-dh-graphs]{\ref*{def:p04-dh-graphs}} & \thepage & Declared positive vertex count and full adjacency matrix. \\
\phantomsection\label{entry:symbol-registry-03344}\(s\) & \texttt{\detokenize{s}} & designated starting vertex index & P04 decision hardness graphs & n/a & \hyperref[def:p04-dh-graphs]{\ref*{def:p04-dh-graphs}} & \thepage & Binary index within the explicit graph; zero-length paths allowed. \\
\phantomsection\label{entry:symbol-registry-03345}\(t\) & \texttt{\detokenize{t}} & designated target vertex index & P04 decision hardness graphs & n/a & \hyperref[def:p04-dh-graphs]{\ref*{def:p04-dh-graphs}} & \thepage & May equal the starting vertex. \\
\phantomsection\label{entry:symbol-registry-03346}\(p_i\) & \texttt{\detokenize{p_i}} & indexed propositional variable of the declared 2-CNF instance & P04 decision hardness CNF & n/a & \hyperref[def:p04-dh-two-cnf]{\ref*{def:p04-dh-two-cnf}} & \thepage & Index lies between one and the unary declaration r. \\
\phantomsection\label{entry:symbol-registry-03347}\(\operatorname{value}_2(u)\) & \texttt{\detokenize{\operatorname{value}_2(u)}} & binary numerical value of a finite bit string & P04 decision hardness validity & n/a & \hyperref[def:p04-dh-validity]{\ref*{def:p04-dh-validity}} & \thepage & The empty word has value zero; leading bits remain significant in sentence coding. \\
\phantomsection\label{entry:symbol-registry-03348}\(Q\) & \texttt{\detokenize{Q}} & finite state set of the compiled recognizer & P04 decision hardness validity & n/a & \hyperref[def:p04-dh-machine-encoding]{\ref*{def:p04-dh-machine-encoding}} & \thepage & State constructors are injectively assigned by finite indexed list position. \\
\phantomsection\label{entry:symbol-registry-03349}\(\Gamma\) & \texttt{\detokenize{\Gamma}} & finite tape alphabet of the compiled recognizer & P04 decision hardness validity & n/a & \hyperref[def:p04-dh-machine-encoding]{\ref*{def:p04-dh-machine-encoding}} & \thepage & Contains distinct blank zero and one symbols; not a logical premise set here. \\
\phantomsection\label{entry:symbol-registry-03350}\(\delta\) & \texttt{\detokenize{\delta}} & partial deterministic machine transition table & P04 decision hardness validity & n/a & \hyperref[def:p04-dh-machine-encoding]{\ref*{def:p04-dh-machine-encoding}} & \thepage & No outgoing accepting-state transitions; missing entries do not accept. \\
\phantomsection\label{entry:symbol-registry-03351}\(q_0\) & \texttt{\detokenize{q_0}} & designated initial machine state & P04 decision hardness validity & n/a & \hyperref[def:p04-dh-machine-encoding]{\ref*{def:p04-dh-machine-encoding}} & \thepage & May already be the accepting state in the reduction presentation. \\
\phantomsection\label{entry:symbol-registry-03352}\(q_{\mathrm a}\) & \texttt{\detokenize{q_{\mathrm a}}} & halting accepting machine state & P04 decision hardness validity & n/a & \hyperref[def:p04-dh-machine-encoding]{\ref*{def:p04-dh-machine-encoding}} & \thepage & Acceptance rather than an arbitrary physical halt is the recognized event. \\
\phantomsection\label{entry:symbol-registry-03353}\(\mathsf a_a\) & \texttt{\detokenize{\mathsf a_a}} & unary term constructor assigned to tape symbol a & P04 decision hardness validity & n/a & \hyperref[def:p04-dh-machine-encoding]{\ref*{def:p04-dh-machine-encoding}} & \thepage & Abbreviation for the indexed family member assigned injectively to a. \\
\phantomsection\label{entry:symbol-registry-03354}\(\mathsf c_q\) & \texttt{\detokenize{\mathsf c_q}} & binary configuration constructor assigned to state q & P04 decision hardness validity & n/a & \hyperref[def:p04-dh-machine-encoding]{\ref*{def:p04-dh-machine-encoding}} & \thepage & Its two arguments are the relative-head tape lists. \\
\phantomsection\label{entry:symbol-registry-03355}\(\sqcup\) & \texttt{\detokenize{\sqcup}} & designated blank tape symbol & P04 decision hardness validity & n/a & \hyperref[def:p04-dh-horn-compilation]{\ref*{def:p04-dh-horn-compilation}} & \thepage & Not set or concept union; separate empty-right and empty-left clauses handle it. \\
\phantomsection\label{entry:symbol-registry-03356}\(\mathsf L\) & \texttt{\detokenize{\mathsf L}} & left machine move code & P04 decision hardness validity & n/a & \hyperref[def:p04-dh-machine-encoding]{\ref*{def:p04-dh-machine-encoding}} & \thepage & Distinct from the deterministic logarithmic-space class in the resources scope. \\
\phantomsection\label{entry:symbol-registry-03357}\(\mathsf R\) & \texttt{\detokenize{\mathsf R}} & right machine move code & P04 decision hardness validity & n/a & \hyperref[def:p04-dh-machine-encoding]{\ref*{def:p04-dh-machine-encoding}} & \thepage & Push the written symbol onto the left list and expose the right tail. \\
\phantomsection\label{entry:symbol-registry-03358}\(\mathsf N\) & \texttt{\detokenize{\mathsf N}} & stationary machine move code & P04 decision hardness validity & n/a & \hyperref[def:p04-dh-machine-encoding]{\ref*{def:p04-dh-machine-encoding}} & \thepage & Write without changing the represented head position. \\
\phantomsection\label{entry:symbol-registry-03359}\(D\) & \texttt{\detokenize{D}} & free ground-term countermodel domain & P04 decision hardness computation proof & n/a & \hyperref[annex:p04.dh.computation]{\ref*{annex:p04.dh.computation}} & \thepage & All finite ground terms over the declared signature; nonempty and unrestricted. \\
\phantomsection\label{entry:symbol-registry-03360}\(\widehat\Phi(z_0)\) & \texttt{\detokenize{\widehat\Phi(z_0)}} & relationalized formula with a constant-replacement parameter & P04 decision hardness relational proof & n/a & \hyperref[annex:p04.dh.relational]{\ref*{annex:p04.dh.relational}} & \thepage & Its truth corresponds to the functional formula under total unique graph interpretations. \\
\phantomsection\label{entry:symbol-registry-05967}\(\mathcal B=\mathbb N^{\mathbb N}\) & \texttt{\detokenize{\mathcal B}} & space of natural-valued infinite names & P4 represented spaces & object & \hyperref[def:p04-rs-baire]{\ref*{def:p04-rs-baire}} & \thepage & Baire topology on names; finite-alphabet delimiter coding is explicitly distinguished. \\
\phantomsection\label{entry:symbol-registry-05968}\(\mathbb N^{<\mathbb N}\) & \texttt{\detokenize{\mathbb N^{<\mathbb N}}} & finite natural sequences & P4 represented spaces & object & \hyperref[def:p04-rs-baire]{\ref*{def:p04-rs-baire}} & \thepage & Used as cylinder prefixes rather than infinite object names. \\
\phantomsection\label{entry:symbol-registry-05969}\([s]\) & \texttt{\detokenize{none}} & Baire cylinder extending finite sequence s & P4 represented spaces & object & \hyperref[def:p04-rs-baire]{\ref*{def:p04-rs-baire}} & \thepage & Not an interval of reals or a degree equivalence class. \\
\phantomsection\label{entry:symbol-registry-05970}\(p{\upharpoonright}k\) & \texttt{\detokenize{\upharpoonright}} & first k entries of a stream & P4 represented spaces & meta & \hyperref[def:p04-rs-baire]{\ref*{def:p04-rs-baire}} & \thepage & Finite information observed; no uniform precision modulus presumed. \\
\phantomsection\label{entry:symbol-registry-05971}\(\delta_X,\delta_Y\) & \texttt{\detokenize{\delta}} & partial surjective name interpretation & P4 represented spaces & object & \hyperref[def:p04-rs-representation]{\ref*{def:p04-rs-representation}} & \thepage & Every point has a name; invalid streams need not be recognizable. \\
\phantomsection\label{entry:symbol-registry-05972}\(E_X,E_Y,D_f\) & \texttt{\detokenize{E_X;E_Y;D_f}} & valid-name and promised-input domains & P4 represented spaces & object & \hyperref[def:p04-rs-realizer]{\ref*{def:p04-rs-realizer}} & \thepage & Subspace topologies; machine realizer may also be defined outside the promise. \\
\phantomsection\label{entry:symbol-registry-05973}\(G,F,H,T\) & \texttt{\detokenize{G;F;H;T}} & stream realizer or translator with locally stated role & P4 represented spaces & meta & \hyperref[def:p04-rs-realizer]{\ref*{def:p04-rs-realizer}} & \thepage & Decoded extensionality does not require identical output streams. \\
\phantomsection\label{entry:symbol-registry-05974}\(\leq_{\mathrm c}\) & \texttt{\detokenize{\leq_{\mathrm c}}} & computable translation reducibility of representations & P4 represented spaces & meta & \hyperref[def:p04-rs-reduction]{\ref*{def:p04-rs-reduction}} & \thepage & Direction is source names to target names; not Weihrauch reducibility of problems. \\
\phantomsection\label{entry:symbol-registry-05975}\(\leq_{\mathrm t}\) & \texttt{\detokenize{\leq_{\mathrm t}}} & continuous translation reducibility of representations & P4 represented spaces & meta & \hyperref[def:p04-rs-reduction]{\ref*{def:p04-rs-reduction}} & \thepage & Topological equivalence alone need not preserve computable points. \\
\phantomsection\label{entry:symbol-registry-05976}\(\delta_{\mathrm{nb}}\) & \texttt{\detokenize{\delta_{\mathrm{nb}}}} & all-basic-neighbourhood enumeration representation & P4 countably based T0 spaces & object & \hyperref[def:p04-rs-neighbourhood]{\ref*{def:p04-rs-neighbourhood}} & \thepage & Positive index-plus-one entries and zero padding; every containing basic open occurs. \\
\phantomsection\label{entry:symbol-registry-05977}\(\nu_{\mathbb Q},\zeta\) & \texttt{\detokenize{\nu;\zeta}} & rational decoder and signed-integer decoder & P4 real representation coding & meta & \hyperref[def:p04-rs-rational-code]{\ref*{def:p04-rs-rational-code}} & \thepage & Total effective rational decoder from numerical pairing; finite fractions have effective output codes. \\
\phantomsection\label{entry:symbol-registry-05978}\(\delta_{\mathrm C}\) & \texttt{\detokenize{\delta_{\mathrm C}}} & host pointwise-error fast Cauchy representation & P4 real computation & object & \hyperref[def:p04-rs-cauchy]{\ref*{def:p04-rs-cauchy}} & \thepage & Non-strict error to limit at most two to the minus n; not the strict 1995 pairwise clause. \\
\phantomsection\label{entry:symbol-registry-05979}\(\delta_{\mathrm I}\) & \texttt{\detokenize{\delta_{\mathrm I}}} & all-open-rational-interval enumeration representation & P4 real computation & object & \hyperref[sec:p04-rs-cauchy]{\ref*{sec:p04-rs-cauchy}} & \thepage & Distinct from the 1995 lower/upper bound sequence representation. \\
\phantomsection\label{entry:symbol-registry-05980}\(\rho_{\mathrm{pair}},\rho_{\mathrm{bounds}}\) & \texttt{\detokenize{\rho_{\mathrm{pair}};\rho_{\mathrm{bounds}}}} & decoded 1995 pairwise and bound naming clauses & P4 source-to-host bridge & object & \hyperref[sec:p04-rs-cauchy]{\ref*{sec:p04-rs-cauchy}} & \thepage & Strict pairwise inequality and supremum/infimum clauses remain explicit with two-way translations. \\
\phantomsection\label{entry:symbol-registry-05981}\(\lvert x\rvert\) & \texttt{\detokenize{\lvert;\rvert}} & real or rational absolute value & P4 represented real arithmetic & object & \hyperref[sec:p04-rs-arithmetic]{\ref*{sec:p04-rs-arithmetic}} & \thepage & Numerical magnitude; not cardinality or finite-sequence length. \\
\phantomsection\label{entry:symbol-registry-05982}\(B_x,B_y,L\) & \texttt{\detokenize{B_x;B_y;L}} & input-derived magnitude bounds and nonzero certificate & P4 real arithmetic algorithms & meta & \hyperref[sec:p04-rs-arithmetic]{\ref*{sec:p04-rs-arithmetic}} & \thepage & Known rational data calculated from names; no external bound silently assumed. \\
\phantomsection\label{entry:symbol-registry-05983}\(\delta_2,Z\) & \texttt{\detokenize{\delta_2;Z}} & discrete bit representation and zero-equality characteristic map & P4 decision obstruction & object & \hyperref[def:p04-rs-bit]{\ref*{def:p04-rs-bit}} & \thepage & Total zero decision is not semidecision; sign is computed only on the nonzero promise. \\
\phantomsection\label{entry:symbol-registry-05984}\(\oplus\) & \texttt{\detokenize{\oplus}} & bitwise exclusive-or in representation counterexample & P4 representation-equivalence annex & object & \hyperref[annex:p04-rs-invariance]{\ref*{annex:p04-rs-invariance}} & \thepage & Not degree join; a fixed noncomputable stream gives continuous translations only. \\
\phantomsection\label{entry:symbol-registry-06052}\(\leq\) & \texttt{\detokenize{\leq}} & canonical ordering & P4 represented-space numerical ordering & 2 & \hyperref[def:p04-rs-cauchy]{\ref*{def:p04-rs-cauchy}} & \thepage & Ordinary order on the declared numerical objects and accuracy bounds in the new represented-space chapter and its nine associated proofs; the original pre-batch5 ordering scopes deliberately do not include this inserted chapter. Representation reducibilities and arithmetic formula syntax have their own typed roles. \\
\phantomsection\label{entry:symbol-registry-06138}\(C[0,1],C(\mathbb R)\) & \texttt{\detokenize{C}} & spaces of continuous real functions with specified domain & P4 function operators & object & \hyperref[def:p04-fo-polygons]{\ref*{def:p04-fo-polygons}} & \thepage & Function names are additional input beyond scalar or pointwise value data. \\
\phantomsection\label{entry:symbol-registry-06139}\(\lVert f\rVert_\infty\) & \texttt{\detokenize{\lVert;\rVert}} & uniform function norm & P4 function operators & object & \hyperref[def:p04-fo-polygons]{\ref*{def:p04-fo-polygons}} & \thepage & Maximum absolute value on the compact unit interval. \\
\phantomsection\label{entry:symbol-registry-06140}\(P_n,Q_n\) & \texttt{\detokenize{P_n;Q_n}} & rational polygon approximants & P4 function operators & object & \hyperref[def:p04-fo-polygons]{\ref*{def:p04-fo-polygons}} & \thepage & Finite rational vertices and declared uniform error; not unnamed pointwise samples. \\
\phantomsection\label{entry:symbol-registry-06141}\(\gamma_{\mathrm P}\) & \texttt{\detokenize{\gamma_{\mathrm P}}} & host uniform-error polygon representation & P4 function operators & object & \hyperref[def:p04-fo-polygons]{\ref*{def:p04-fo-polygons}} & \thepage & Error to the named function at most two to the minus n. \\
\phantomsection\label{entry:symbol-registry-06142}\(\gamma_{\mathrm{pair}}\) & \texttt{\detokenize{\gamma_{\mathrm{pair}}}} & host strict-pairwise polygon representation on numerical-code streams & P4 source bridge & object & \hyperref[sec:p04-fo-names]{\ref*{sec:p04-fo-names}} & \thepage & Separately defined host bridge to native delta\_C symbol-stream function names; not a source-glyph replacement. Native function delta\_C and host scalar delta\_C retain their declared different operand types. \\
\phantomsection\label{entry:symbol-registry-06143}\(\gamma_{\mathrm B},\gamma_{\mathbb R}\) & \texttt{\detokenize{\gamma_{\mathrm B};\gamma_{\mathbb R}}} & local and global compact-range box representations & P4 function operators & object & \hyperref[def:p04-fo-boxes]{\ref*{def:p04-fo-boxes}} & \thepage & Lists every true rational compact-to-open range property; zero padding carries no property. \\
\phantomsection\label{entry:symbol-registry-06144}\(\mu,\mu'\) & \texttt{\detokenize{\mu;\mu'}} & function and derivative continuity moduli & P4 function operators & meta & \hyperref[thm:p04-fo-evaluation]{\ref*{thm:p04-fo-evaluation}} & \thepage & Function modulus is computed from a uniform name; derivative modulus is explicitly additional input. \\
\phantomsection\label{entry:symbol-registry-06145}\(J,\operatorname{Int}\) & \texttt{\detokenize{J;\operatorname{Int}}} & primitive function and integration operation & P4 function operators & object & \hyperref[thm:p04-fo-integration]{\ref*{thm:p04-fo-integration}} & \thepage & Scalar area and named function output remain distinct. \\
\phantomsection\label{entry:symbol-registry-06146}\(L,M,B,h,r\) & \texttt{\detokenize{L;M;B;h;r}} & derived slope or magnitude bounds and chosen mesh data & P4 function operator algorithms & meta & \hyperref[thm:p04-fo-evaluation]{\ref*{thm:p04-fo-evaluation}} & \thepage & Finite rational or integer data obtained from names; not free global information. \\
\phantomsection\label{entry:symbol-registry-06147}\(C^1[0,1],f'\) & \texttt{\detokenize{C^1;f'}} & continuously differentiable functions and derivative & P4 differentiation & object & \hyperref[thm:p04-fo-differentiation]{\ref*{thm:p04-fo-differentiation}} & \thepage & One-sided endpoint derivatives extend continuously; ordinary function names do not encode derivative accuracy. \\
\phantomsection\label{entry:symbol-registry-06148}\(\mathcal F_{\mathrm{nd}}\) & \texttt{\detokenize{\mathcal F_{\mathrm{nd}}}} & sign-changing functions with no nondegenerate zero interval & P4 zero relation & object & \hyperref[def:p04-fo-sign-change-class]{\ref*{def:p04-fo-sign-change-class}} & \thepage & Semantic promise; distinct from general sign change or uniqueness. \\
\phantomsection\label{entry:symbol-registry-06149}\(G_t,H_t\) & \texttt{\detokenize{G_t;H_t}} & source plateau and selector-counterexample polygon families & P4 zero finding & object & \hyperref[thm:p04-fo-general-zero]{\ref*{thm:p04-fo-general-zero}} & \thepage & The two families have different middle vertices and different evidential roles. \\
\phantomsection\label{entry:symbol-registry-06150}\(K,u,v\) & \texttt{\detokenize{K;u;v}} & excluded compact complement and candidate isolating interval endpoints & P4 unique-zero search & object & \hyperref[thm:p04-fo-unique-zero]{\ref*{thm:p04-fo-unique-zero}} & \thepage & Empty complement is vacuous; endpoints may lie outside the unit interval to handle boundary roots. \\
\phantomsection\label{entry:symbol-registry-06151}\(R_{\mathrm{nd}}\) & \texttt{\detokenize{R_{\mathrm{nd}}}} & name-realized zero relation on the promised class & P4 zero relation & object & \hyperref[thm:p04-fo-nonextensional]{\ref*{thm:p04-fo-nonextensional}} & \thepage & Different names can legitimately produce different zero objects. \\
\phantomsection\label{entry:symbol-registry-06152}\(U,T_j,F\) & \texttt{\detokenize{U;T_j;F}} & effective open cover rational tents and weighted function & P4 no-computable-zero construction & object & \hyperref[thm:p04-fo-no-computable-zero]{\ref*{thm:p04-fo-no-computable-zero}} & \thepage & Nonnegative weighted series with explicit uniform tail; complement supplies zeros. \\
\phantomsection\label{entry:symbol-registry-06153}\(\lvert f(x)\rvert\) & \texttt{\detokenize{\lvert;\rvert}} & absolute residual value & P4 function operators & object & \hyperref[thm:p04-fo-approximate-zero]{\ref*{thm:p04-fo-approximate-zero}} & \thepage & Small value does not itself bound distance to a zero or prove existence. \\
\phantomsection\label{entry:symbol-registry-06154}\(\Sigma^*,\Sigma^\omega\) & \texttt{\detokenize{\Sigma^*;\Sigma^\omega}} & native finite words and infinite source symbol streams & P4 function operators native and host bridge & object & \hyperref[def:p04-fo-native-notation]{\ref*{def:p04-fo-native-notation}} & \thepage & source=Weihrauch1995FunctionOperatorsFO; selectors=Definitions3.1/6.1/6.5 and AppendixC1; source-symbol streams are not Baire numerical-entry syntax. \\
\phantomsection\label{entry:symbol-registry-06155}\(\nu_{\mathrm{bin}}\) & \texttt{\detokenize{\nu_{\mathrm{bin}}}} & native finite binary-word natural notation & P4 function operators native and host bridge & meta & \hyperref[def:p04-fo-native-notation]{\ref*{def:p04-fo-native-notation}} & \thepage & source=Weihrauch1995FunctionOperatorsFO; selector=AppendixC1(1); input=Sigma*; output=N; canonical zero or leading-one word. \\
\phantomsection\label{entry:symbol-registry-06156}\(\nu_{\mathbb Q}\) & \texttt{\detokenize{\nu_{\mathbb Q}}} & native signed reduced rational-word notation & P4 shared rational notation with typed operands & meta & \hyperref[def:p04-fo-native-notation]{\ref*{def:p04-fo-native-notation}} & \thepage & source=Weihrauch1995FunctionOperatorsFO; selector=AppendixC1(2); input=Sigma*; output=Q; genuine typed overload with host numerical decoder N->Q at def:p04-rs-rational-code; not a fabricated disjoint namespace. \\
\phantomsection\label{entry:symbol-registry-06157}\(\bar u\) & \texttt{\detokenize{\bar}} & native rational value decoded from finite word u & P4 function operators native and host bridge & meta & \hyperref[def:p04-fo-native-notation]{\ref*{def:p04-fo-native-notation}} & \thepage & source=Weihrauch1995FunctionOperatorsFO; selector=AppendixC1; u belongs to dom(native nu\_Q). \\
\phantomsection\label{entry:symbol-registry-06158}\(\rho_C\) & \texttt{\detokenize{\rho_C}} & native scalar strict-pairwise Cauchy representation & P4 scalar and function representations with typed operands & object & \hyperref[def:p04-fo-native-notation]{\ref*{def:p04-fo-native-notation}} & \thepage & source=Weihrauch1995FunctionOperatorsFO; definition-selector=def:p04-fo-native-scalar; source-selector=Definition3.1; input=Sigma\^{}omega; output=R; limit plus strict pairwise clause; scalar native operator operands. \\
\phantomsection\label{entry:symbol-registry-06159}\(\delta_C\) & \texttt{\detokenize{\delta_C}} & native unit-interval continuous-function Cauchy representation & P4 scalar and function representations with typed operands & object & \hyperref[def:p04-fo-native-notation]{\ref*{def:p04-fo-native-notation}} & \thepage & source=Weihrauch1995FunctionOperatorsFO; definition-selector=def:p04-fo-native-functions; source-selector=Definition6.5(2); input=Sigma\^{}omega polygon-word blocks; output=C[0,1]; honest typed overload with host scalar delta\_C:Baire->R; typography is not a namespace separation. \\
\phantomsection\label{entry:symbol-registry-06160}\(\delta_{\mathbb R}\) & \texttt{\detokenize{\delta_{\mathbb R}}} & native global continuous-function box representation & P4 scalar and function representations with typed operands & object & \hyperref[def:p04-fo-native-notation]{\ref*{def:p04-fo-native-notation}} & \thepage & source=Weihrauch1995FunctionOperatorsFO; definition-selector=def:p04-fo-native-functions; source-selector=Definition6.1; input=Sigma\^{}omega four-rational-word records; output=C(R); every true property including vacuous reversed-domain boxes. \\
\phantomsection\label{entry:symbol-registry-06161}\(\alpha\) & \texttt{\detokenize{\alpha}} & native rational-polygon finite-word notation & P4 function operators native and host bridge & meta & \hyperref[def:p04-fo-native-functions]{\ref*{def:p04-fo-native-functions}} & \thepage & source=Weihrauch1995FunctionOperatorsFO; selector=Definition6.5(1); input=Sigma*; output=Pg; rational abscissae strictly increase from zero to one. \\
\phantomsection\label{entry:symbol-registry-06162}\(\#,\text{\textcent}\) & \texttt{\detokenize{\#;\textcent}} & native within-record and between-record separators & P4 function operators native and host bridge & object & \hyperref[def:p04-fo-native-notation]{\ref*{def:p04-fo-native-notation}} & \thepage & source=Weihrauch1995FunctionOperatorsFO; selectors=Definitions3.1/6.1/6.5; cent is distinct from hash; there is no native zero-padding record. \\
\phantomsection\label{entry:symbol-registry-06163}\(\operatorname{Diff}\) & \texttt{\detokenize{\operatorname{Diff}}} & native differentiation operator & P4 native operator signatures with typed operands & object & \hyperref[sec:p04-fo-differentiation]{\ref*{sec:p04-fo-differentiation}} & \thepage & source=Weihrauch1995FunctionOperatorsFO; selector=Theorem6.9; declaration-selector=def:p04-fo-native-operators; native function delta\_C names both input and output. \\
\phantomsection\label{entry:symbol-registry-06562}\(t\equiv a\pmod n\) & \texttt{\detokenize{\equiv}} & number-theoretic congruence of natural numbers modulo n & \hyperref[sec:decision-hierarchy-02]{\ref*{sec:decision-hierarchy-02}} & meta & \hyperref[sec:decision-hierarchy-02]{\ref*{sec:decision-hierarchy-02}} & \thepage & Holds iff n divides t-a; operands are numbers and the modulus is displayed by \texttt{\detokenize{\pmod}}; used once, in the proof of lem:goedel-beta with moduli 1+(i+1)m; not a formula of the arithmetic object language. Remainder bridge: thm:p01-w02-p01-025-equiv-bridges part 5 (def:p01-w02-p01-025-equivalence-roles item 6). \\
\phantomsection\label{entry:symbol-registry-06563}\(A\equiv_m B\) & \texttt{\detokenize{\equiv_m}} & many-one equivalence of sets of natural numbers & \hyperref[sec:decision-hierarchy-03]{\ref*{sec:decision-hierarchy-03}} & meta & \hyperref[sec:decision-hierarchy-03]{\ref*{sec:decision-hierarchy-03}} & \thepage & \(A\leq_m B\) and \(B\leq_m A\); same relation as the registered row for \(\equiv_m\), whose scope (P4 frontiers-computability) does not include rem:jump-convention (def:p01-w02-p01-025-equivalence-roles item 7). \\
\phantomsection\label{entry:symbol-registry-07474}\(\sigma_{\mathrm S}\) & \texttt{\detokenize{\sigma_{\mathrm S}}} & sampled-value representation of continuous functions & P4 real-analysis function information & object & \hyperref[def:p04-ca-sample-name]{\ref*{def:p04-ca-sample-name}} & \thepage & All rational arguments are sampled with scalar error control; no uniform function modulus is included. \\
\phantomsection\label{entry:symbol-registry-07475}\(q_i\) & \texttt{\detokenize{q_i}} & rational sample location and approximation table symbol family member \(q_i\) & P4 real-analysis function information & object & \hyperref[def:p04-ca-sample-name]{\ref*{def:p04-ca-sample-name}} & \thepage & Countable rational enumeration and paired table positions; not a scalar Cauchy sequence indexed only by precision. Explicit displayed family; member types and common actual first-use exposed, not one atomic glyph. Family listing, not one atomic glyph. Author macro \texttt{\detokenize{q\_i;r\_{i,n}}} not carried to the members. \\
\phantomsection\label{entry:symbol-registry-07476}\(r_{i,n}\) & \texttt{\detokenize{r_{i,n}}} & rational sample location and approximation table symbol family member \(r_{i,n}\) & P4 real-analysis function information & object & \hyperref[def:p04-ca-sample-name]{\ref*{def:p04-ca-sample-name}} & \thepage & Countable rational enumeration and paired table positions; not a scalar Cauchy sequence indexed only by precision. Explicit displayed family; member types and common actual first-use exposed, not one atomic glyph. Family listing, not one atomic glyph. Author macro \texttt{\detokenize{q\_i;r\_{i,n}}} not carried to the members. \\
\phantomsection\label{entry:symbol-registry-07477}\(\mu\) & \texttt{\detokenize{\mu}} & integer modulus supplied with a sampled-value name & P4 real-analysis function information & object & \hyperref[def:p04-ca-sample-name]{\ref*{def:p04-ca-sample-name}} & \thepage & Same function in both input components; continuous existence does not imply computable extraction from samples. \\
\phantomsection\label{entry:symbol-registry-07478}\(\theta\) & \texttt{\detokenize{\theta}} & basic and scaled C1 bump functions symbol family member \(\theta\) & P4 real-analysis derivative construction & object & \hyperref[def:p04-ca-bump]{\ref*{def:p04-ca-bump}} & \thepage & Actual native theta glyph and scaled presentation retained; corrected height formula is recorded outside reader. Explicit displayed family; member types and common actual first-use exposed, not one atomic glyph. Family listing, not one atomic glyph. Author macro \texttt{\detokenize{\theta;\theta\_{\alpha\beta}}} not carried to the members. \\
\phantomsection\label{entry:symbol-registry-07479}\(\theta_{\alpha\beta}\) & \texttt{\detokenize{\theta_{\alpha\beta}}} & basic and scaled C1 bump functions symbol family member \(\theta_{\alpha\beta}\) & P4 real-analysis derivative construction & object & \hyperref[def:p04-ca-bump]{\ref*{def:p04-ca-bump}} & \thepage & Actual native theta glyph and scaled presentation retained; corrected height formula is recorded outside reader. Explicit displayed family; member types and common actual first-use exposed, not one atomic glyph. Family listing, not one atomic glyph. Author macro \texttt{\detokenize{\theta;\theta\_{\alpha\beta}}} not carried to the members. \\
\phantomsection\label{entry:symbol-registry-07480}\(\lambda\) & \texttt{\detokenize{\lambda}} & positive maximum amplitude of the basic bump & P4 real-analysis derivative construction & object & \hyperref[def:p04-ca-bump]{\ref*{def:p04-ca-bump}} & \thepage & Exactly sixteen divided by the product of twenty-five and square root of five; not an eigenvalue or lambda binder. \\
\phantomsection\label{entry:symbol-registry-07481}\(\alpha_n\) & \texttt{\detokenize{\alpha_n}} & half-width and height at an enumerated centre symbol family member \(\alpha_n\) & P4 real-analysis derivative construction & object & \hyperref[def:p04-ca-coding-function]{\ref*{def:p04-ca-coding-function}} & \thepage & Defined for n in A with unique enumeration position k; height divided by width is two to minus n. Explicit displayed family; member types and common actual first-use exposed, not one atomic glyph. Family listing, not one atomic glyph. Author macro \texttt{\detokenize{\alpha\_n;\beta\_n}} not carried to the members. \\
\phantomsection\label{entry:symbol-registry-07482}\(\beta_n\) & \texttt{\detokenize{\beta_n}} & half-width and height at an enumerated centre symbol family member \(\beta_n\) & P4 real-analysis derivative construction & object & \hyperref[def:p04-ca-coding-function]{\ref*{def:p04-ca-coding-function}} & \thepage & Defined for n in A with unique enumeration position k; height divided by width is two to minus n. Explicit displayed family; member types and common actual first-use exposed, not one atomic glyph. Family listing, not one atomic glyph. Author macro \texttt{\detokenize{\alpha\_n;\beta\_n}} not carried to the members. \\
\phantomsection\label{entry:symbol-registry-07483}\(A\) & \texttt{\detokenize{A}} & undecidable enumerable set and total injective enumeration symbol family member \(A\) & P4 real-analysis derivative construction & meta & \hyperref[def:p04-ca-coding-function]{\ref*{def:p04-ca-coding-function}} & \thepage & Computing finite enumeration prefixes is not deciding arbitrary membership in the set. Explicit displayed family; member types and common actual first-use exposed, not one atomic glyph. Family listing, not one atomic glyph. Author macro \texttt{\detokenize{A;h}} not carried to the members. \\
\phantomsection\label{entry:symbol-registry-07484}\(h\) & \texttt{\detokenize{h}} & undecidable enumerable set and total injective enumeration symbol family member \(h\) & P4 real-analysis derivative construction & meta & \hyperref[def:p04-ca-coding-function]{\ref*{def:p04-ca-coding-function}} & \thepage & Computing finite enumeration prefixes is not deciding arbitrary membership in the set. Explicit displayed family; member types and common actual first-use exposed, not one atomic glyph. Family listing, not one atomic glyph. Author macro \texttt{\detokenize{A;h}} not carried to the members. \\
\phantomsection\label{entry:symbol-registry-07485}\(f\) & \texttt{\detokenize{f}} & bump-series function and finite partial sum symbol family member \(f\) & P4 real-analysis derivative construction & object & \hyperref[def:p04-ca-coding-function]{\ref*{def:p04-ca-coding-function}} & \thepage & Native sum through k equals M for g\_M; disjoint supports give a maximum omitted-term bound. Explicit displayed family; member types and common actual first-use exposed, not one atomic glyph. Family listing, not one atomic glyph. Author macro \texttt{\detokenize{f;g\_M}} not carried to the members. \\
\phantomsection\label{entry:symbol-registry-07486}\(g_M\) & \texttt{\detokenize{g_M}} & bump-series function and finite partial sum symbol family member \(g_M\) & P4 real-analysis derivative construction & object & \hyperref[def:p04-ca-coding-function]{\ref*{def:p04-ca-coding-function}} & \thepage & Native sum through k equals M for g\_M; disjoint supports give a maximum omitted-term bound. Explicit displayed family; member types and common actual first-use exposed, not one atomic glyph. Family listing, not one atomic glyph. Author macro \texttt{\detokenize{f;g\_M}} not carried to the members. \\
\phantomsection\label{entry:symbol-registry-07487}\(c\) & \texttt{\detokenize{c}} & real argument clamping map for the basic bump & P4 real-analysis derivative construction & object & \hyperref[sec:p04-ca-bumps]{\ref*{sec:p04-ca-bumps}} & \thepage & Defined by min and max before use; avoids deciding equality at a support endpoint. \\
\phantomsection\label{entry:symbol-registry-07488}\(K\) & \texttt{\detokenize{K}} & second-derivative bound exponent and derivative modulus symbol family member \(K\) & P4 real-analysis higher derivatives & object & \hyperref[thm:p04-ca-derivative-modulus]{\ref*{thm:p04-ca-derivative-modulus}} & \thepage & Existing bound for fixed-function computability differs from a supplied bound for a uniform realizer. Explicit displayed family; member types and common actual first-use exposed, not one atomic glyph. Family listing, not one atomic glyph. Author macro \texttt{\detokenize{K;\ell;\mu'}} not carried to the members. \\
\phantomsection\label{entry:symbol-registry-07489}\(\ell\) & \texttt{\detokenize{\ell}} & second-derivative bound exponent and derivative modulus symbol family member \(\ell\) & P4 real-analysis higher derivatives & object & \hyperref[thm:p04-ca-derivative-modulus]{\ref*{thm:p04-ca-derivative-modulus}} & \thepage & Existing bound for fixed-function computability differs from a supplied bound for a uniform realizer. Explicit displayed family; member types and common actual first-use exposed, not one atomic glyph. Family listing, not one atomic glyph. Author macro \texttt{\detokenize{K;\ell;\mu'}} not carried to the members. \\
\phantomsection\label{entry:symbol-registry-07490}\(\mu'\) & \texttt{\detokenize{\mu'}} & second-derivative bound exponent and derivative modulus symbol family member \(\mu'\) & P4 real-analysis higher derivatives & object & \hyperref[thm:p04-ca-derivative-modulus]{\ref*{thm:p04-ca-derivative-modulus}} & \thepage & Existing bound for fixed-function computability differs from a supplied bound for a uniform realizer. Explicit displayed family; member types and common actual first-use exposed, not one atomic glyph. Family listing, not one atomic glyph. Author macro \texttt{\detokenize{K;\ell;\mu'}} not carried to the members. \\
\phantomsection\label{entry:symbol-registry-07491}\(m\) & \texttt{\detokenize{m}} & supplied positive lower slope and unique scalar zero symbol family member \(m\) & P4 real-analysis root certificates & object & \hyperref[def:p04-ca-root-bound]{\ref*{def:p04-ca-root-bound}} & \thepage & The slope promise is additional information; uniqueness alone gives no positive lower slope. Explicit displayed family; member types and common actual first-use exposed, not one atomic glyph. Family listing, not one atomic glyph. Author macro \texttt{\detokenize{m;z}} not carried to the members. \\
\phantomsection\label{entry:symbol-registry-07492}\(z\) & \texttt{\detokenize{z}} & supplied positive lower slope and unique scalar zero symbol family member \(z\) & P4 real-analysis root certificates & object & \hyperref[def:p04-ca-root-bound]{\ref*{def:p04-ca-root-bound}} & \thepage & The slope promise is additional information; uniqueness alone gives no positive lower slope. Explicit displayed family; member types and common actual first-use exposed, not one atomic glyph. Family listing, not one atomic glyph. Author macro \texttt{\detokenize{m;z}} not carried to the members. \\
\phantomsection\label{entry:symbol-registry-07493}\(\epsilon\) & \texttt{\detokenize{\epsilon}} & uniform perturbation bound and an existing zero of the perturbed function symbol family member \(\epsilon\) & P4 real-analysis root certificates & object & \hyperref[thm:p04-ca-root-certificate]{\ref*{thm:p04-ca-root-certificate}} & \thepage & Perturbation estimate applies to every existing zero and does not promise that a zero survives. Explicit displayed family; member types and common actual first-use exposed, not one atomic glyph. Family listing, not one atomic glyph. Author macro \texttt{\detokenize{\epsilon;w}} not carried to the members. \\
\phantomsection\label{entry:symbol-registry-07494}\(w\) & \texttt{\detokenize{w}} & uniform perturbation bound and an existing zero of the perturbed function symbol family member \(w\) & P4 real-analysis root certificates & object & \hyperref[thm:p04-ca-root-certificate]{\ref*{thm:p04-ca-root-certificate}} & \thepage & Perturbation estimate applies to every existing zero and does not promise that a zero survives. Explicit displayed family; member types and common actual first-use exposed, not one atomic glyph. Family listing, not one atomic glyph. Author macro \texttt{\detokenize{\epsilon;w}} not carried to the members. \\
\phantomsection\label{entry:symbol-registry-07584}\(f:\subseteq X\rightrightarrows Y\) & \texttt{\detokenize{\subseteq;\rightrightarrows}} & partial multivalued problem with nonempty allowed-answer sets & P4 Weihrauch realizer interfaces & meta & \hyperref[def:p04-wc-multivalued]{\ref*{def:p04-wc-multivalued}} & \thepage & Native partial marker and double function arrow; not a total map into individual points. \\
\phantomsection\label{entry:symbol-registry-07585}\(F\vdash f\) & \texttt{\detokenize{\vdash}} & name-based realizer relation & P4 Weihrauch realizer interfaces & meta & \hyperref[def:p04-wc-multivalued]{\ref*{def:p04-wc-multivalued}} & \thepage & All promised input names and allowed output names; not a logical derivability turnstile. \\
\phantomsection\label{entry:symbol-registry-07586}\(H\) & \texttt{\detokenize{H}} & postprocessor and preprocessor in the 2009 convention symbol family member \(H\) & P4 Weihrauch realizer interfaces & meta & \hyperref[def:p04-wc-reduction]{\ref*{def:p04-wc-reduction}} & \thepage & Actual formula H paired id and GK; strong H receives no original input. Explicit displayed family; member types and common actual first-use exposed, not one atomic glyph. Family listing, not one atomic glyph. Author macro \texttt{\detokenize{H;K}} not carried to the members. \\
\phantomsection\label{entry:symbol-registry-07587}\(K\) & \texttt{\detokenize{K}} & postprocessor and preprocessor in the 2009 convention symbol family member \(K\) & P4 Weihrauch realizer interfaces & meta & \hyperref[def:p04-wc-reduction]{\ref*{def:p04-wc-reduction}} & \thepage & Actual formula H paired id and GK; strong H receives no original input. Explicit displayed family; member types and common actual first-use exposed, not one atomic glyph. Family listing, not one atomic glyph. Author macro \texttt{\detokenize{H;K}} not carried to the members. \\
\phantomsection\label{entry:symbol-registry-07588}\(H\) & \texttt{\detokenize{H}} & preprocessor and postprocessor in the 2010 convention symbol family member \(H\) & P4 Weihrauch realizer interfaces & meta & \hyperref[def:p04-wc-reduction]{\ref*{def:p04-wc-reduction}} & \thepage & Actual separate formula K paired id and GH; opposite operator names are explicitly declared. Explicit displayed family; member types and common actual first-use exposed, not one atomic glyph. Family listing, not one atomic glyph. Author macro \texttt{\detokenize{H;K}} not carried to the members. \\
\phantomsection\label{entry:symbol-registry-07589}\(K\) & \texttt{\detokenize{K}} & preprocessor and postprocessor in the 2010 convention symbol family member \(K\) & P4 Weihrauch realizer interfaces & meta & \hyperref[def:p04-wc-reduction]{\ref*{def:p04-wc-reduction}} & \thepage & Actual separate formula K paired id and GH; opposite operator names are explicitly declared. Explicit displayed family; member types and common actual first-use exposed, not one atomic glyph. Family listing, not one atomic glyph. Author macro \texttt{\detokenize{H;K}} not carried to the members. \\
\phantomsection\label{entry:symbol-registry-07590}\(\leq_{\mathrm W}\) & \texttt{\detokenize{\leq_{\mathrm W}}} & ordinary and strong Weihrauch reduction symbol family member \(\leq_{\mathrm W}\) & P4 Weihrauch realizer interfaces & meta & \hyperref[def:p04-wc-reduction]{\ref*{def:p04-wc-reduction}} & \thepage & Native upright subscripts; universal quantifier over oracle realizers not a favourable-oracle assumption. Explicit displayed family; member types and common actual first-use exposed, not one atomic glyph. Family listing, not one atomic glyph. Author macro \texttt{\detokenize{\leq\_{\mathrm W};\leq\_{\mathrm{sW}}}} not carried to the members. \\
\phantomsection\label{entry:symbol-registry-07591}\(\leq_{\mathrm{sW}}\) & \texttt{\detokenize{\leq_{\mathrm{sW}}}} & ordinary and strong Weihrauch reduction symbol family member \(\leq_{\mathrm{sW}}\) & P4 Weihrauch realizer interfaces & meta & \hyperref[def:p04-wc-reduction]{\ref*{def:p04-wc-reduction}} & \thepage & Native upright subscripts; universal quantifier over oracle realizers not a favourable-oracle assumption. Explicit displayed family; member types and common actual first-use exposed, not one atomic glyph. Family listing, not one atomic glyph. Author macro \texttt{\detokenize{\leq\_{\mathrm W};\leq\_{\mathrm{sW}}}} not carried to the members. \\
\phantomsection\label{entry:symbol-registry-07592}\(\equiv_{\mathrm W}\) & \texttt{\detokenize{\equiv_{\mathrm W}}} & mutual reduction strictness and incomparability symbol family member \(\equiv_{\mathrm W}\) & P4 Weihrauch realizer interfaces & meta & \hyperref[def:p04-wc-reduction]{\ref*{def:p04-wc-reduction}} & \thepage & Each claimed strict or incomparable comparison has its own missing-arrow proof. Explicit displayed family; member types and common actual first-use exposed, not one atomic glyph. Family listing, not one atomic glyph. Author macro \texttt{\detokenize{\equiv;\mid}} not carried to the members. \\
\phantomsection\label{entry:symbol-registry-07593}\(\equiv_{\mathrm{sW}}\) & \texttt{\detokenize{\equiv_{\mathrm{sW}}}} & mutual reduction strictness and incomparability symbol family member \(\equiv_{\mathrm{sW}}\) & P4 Weihrauch realizer interfaces & meta & \hyperref[def:p04-wc-reduction]{\ref*{def:p04-wc-reduction}} & \thepage & Each claimed strict or incomparable comparison has its own missing-arrow proof. Explicit displayed family; member types and common actual first-use exposed, not one atomic glyph. Family listing, not one atomic glyph. Author macro \texttt{\detokenize{\equiv;\mid}} not carried to the members. \\
\phantomsection\label{entry:symbol-registry-07594}\(<_{\mathrm W}\) & \texttt{\detokenize{<_{\mathrm W}}} & mutual reduction strictness and incomparability symbol family member \(<_{\mathrm W}\) & P4 Weihrauch realizer interfaces & meta & \hyperref[def:p04-wc-reduction]{\ref*{def:p04-wc-reduction}} & \thepage & Each claimed strict or incomparable comparison has its own missing-arrow proof. Explicit displayed family; member types and common actual first-use exposed, not one atomic glyph. Family listing, not one atomic glyph. Author macro \texttt{\detokenize{\equiv;\mid}} not carried to the members. \\
\phantomsection\label{entry:symbol-registry-07595}\(\mid_{\mathrm W}\) & \texttt{\detokenize{\mid_{\mathrm W}}} & mutual reduction strictness and incomparability symbol family member \(\mid_{\mathrm W}\) & P4 Weihrauch realizer interfaces & meta & \hyperref[def:p04-wc-reduction]{\ref*{def:p04-wc-reduction}} & \thepage & Each claimed strict or incomparable comparison has its own missing-arrow proof. Explicit displayed family; member types and common actual first-use exposed, not one atomic glyph. Family listing, not one atomic glyph. Author macro \texttt{\detokenize{\equiv;\mid}} not carried to the members. \\
\phantomsection\label{entry:symbol-registry-07596}\(\langle p,q\rangle\) & \texttt{\detokenize{\langle p,}\allowbreak\detokenize{q\rangle}} & binary and countable tupling of stream names symbol family member \(\langle p,q\rangle\) & P4 Weihrauch name products & meta & \hyperref[def:p04-wc-tupling]{\ref*{def:p04-wc-tupling}} & \thepage & Computable projections; not the numerical pairing of coordinate indices. Explicit displayed family; member types and common actual first-use exposed, not one atomic glyph. Family listing, not one atomic glyph. Author macro \texttt{\detokenize{\langle;\rangle}} not carried to the members. \\
\phantomsection\label{entry:symbol-registry-07597}\(\langle p_i\rangle_{i\in\mathbb N}\) & \texttt{\detokenize{\langle p_i\rangle_{i\in\mathbb N}}} & binary and countable tupling of stream names symbol family member \(\langle p_i\rangle_{i\in\mathbb N}\) & P4 Weihrauch name products & meta & \hyperref[def:p04-wc-tupling]{\ref*{def:p04-wc-tupling}} & \thepage & Computable projections; not the numerical pairing of coordinate indices. Explicit displayed family; member types and common actual first-use exposed, not one atomic glyph. Family listing, not one atomic glyph. Author macro \texttt{\detokenize{\langle;\rangle}} not carried to the members. \\
\phantomsection\label{entry:symbol-registry-07598}\(\langle i,j\rangle\) & \texttt{\detokenize{\langle;\rangle}} & effective numerical pairing of coordinate indices & P4 Weihrauch name products & meta & \hyperref[def:p04-wc-tupling]{\ref*{def:p04-wc-tupling}} & \thepage & Honest local typed overload with stream tupling in the same definition. \\
\phantomsection\label{entry:symbol-registry-07599}\(\mathcal C\) & \texttt{\detokenize{\mathcal C}} & Cantor binary-stream carrier & P4 Weihrauch choice & object & \hyperref[def:p04-wc-tupling]{\ref*{def:p04-wc-tupling}} & \thepage & Explicit shorthand for the source carrier zero-one to the natural-number power. \\
\phantomsection\label{entry:symbol-registry-07600}\(\mathrm{id}\times f\) & \texttt{\detokenize{\mathrm{id};\times}} & cylindrification retaining an independent Baire input & P4 Weihrauch cylinders & object & \hyperref[def:p04-wc-cylinder]{\ref*{def:p04-wc-cylinder}} & \thepage & Identity uses the identity representation; output recovers actual input digits. \\
\phantomsection\label{entry:symbol-registry-07601}\(\widehat f\) & \texttt{\detokenize{\widehat}} & countable parallelization of a multivalued problem & P4 Weihrauch parallelization & object & \hyperref[def:p04-wc-parallel]{\ref*{def:p04-wc-parallel}} & \thepage & All instances independently supplied; one allowed answer per coordinate. \\
\phantomsection\label{entry:symbol-registry-07602}\(f^*\) & \texttt{\detokenize{f^*}} & length-tagged finite parallelization & P4 Weihrauch parallelization & object & \hyperref[def:p04-wc-parallel]{\ref*{def:p04-wc-parallel}} & \thepage & Zero-length case included; ordinary postprocessing preserves the original sublist partition. \\
\phantomsection\label{entry:symbol-registry-07603}\(\psi_-^X\) & \texttt{\detokenize{\psi_-^X}} & negative closed-set representation & P4 Weihrauch choice & object & \hyperref[def:p04-wc-negative]{\ref*{def:p04-wc-negative}} & \thepage & Native discrete plus-one exclusion clause and explicitly coded Cantor exterior cylinders. \\
\phantomsection\label{entry:symbol-registry-07604}\(\mathcal A_-(X)\) & \texttt{\detokenize{\mathcal A_-(X)}} & represented space of negatively described closed sets & P4 Weihrauch choice & object & \hyperref[def:p04-wc-choice]{\ref*{def:p04-wc-choice}} & \thepage & Native subscript minus; no positive membership information is supplied. \\
\phantomsection\label{entry:symbol-registry-07605}\(\mathsf{C}_X\) & \texttt{\detokenize{\mathsf{C}_X}} & closed choice with locally specified carrier symbol family member \(\mathsf{C}_X\) & P4 Weihrauch choice & object & \hyperref[def:p04-wc-choice]{\ref*{def:p04-wc-choice}} & \thepage & Native sans-serif choice glyph; empty sets are outside the promise and point answers are not set names. Explicit displayed family; member types and common actual first-use exposed, not one atomic glyph. Family listing, not one atomic glyph. Author macro \texttt{\detokenize{\mathsf{C}}} not carried to the members. \\
\phantomsection\label{entry:symbol-registry-07606}\(\mathsf{C}_{\mathbb N}\) & \texttt{\detokenize{\mathsf{C}_{\mathbb N}}} & closed choice with locally specified carrier symbol family member \(\mathsf{C}_{\mathbb N}\) & P4 Weihrauch choice & object & \hyperref[def:p04-wc-choice]{\ref*{def:p04-wc-choice}} & \thepage & Native sans-serif choice glyph; empty sets are outside the promise and point answers are not set names. Explicit displayed family; member types and common actual first-use exposed, not one atomic glyph. Family listing, not one atomic glyph. Author macro \texttt{\detokenize{\mathsf{C}}} not carried to the members. \\
\phantomsection\label{entry:symbol-registry-07607}\(\mathsf{C}_{\{0,1\}}\) & \texttt{\detokenize{\mathsf{C}_{\{0,1\}}}} & closed choice with locally specified carrier symbol family member \(\mathsf{C}_{\{0,1\}}\) & P4 Weihrauch choice & object & \hyperref[def:p04-wc-choice]{\ref*{def:p04-wc-choice}} & \thepage & Native sans-serif choice glyph; empty sets are outside the promise and point answers are not set names. Explicit displayed family; member types and common actual first-use exposed, not one atomic glyph. Family listing, not one atomic glyph. Author macro \texttt{\detokenize{\mathsf{C}}} not carried to the members. \\
\phantomsection\label{entry:symbol-registry-07608}\(\mathsf{C}_{\mathcal C}\) & \texttt{\detokenize{\mathsf{C}_{\mathcal C}}} & closed choice with locally specified carrier symbol family member \(\mathsf{C}_{\mathcal C}\) & P4 Weihrauch choice & object & \hyperref[def:p04-wc-choice]{\ref*{def:p04-wc-choice}} & \thepage & Native sans-serif choice glyph; empty sets are outside the promise and point answers are not set names. Explicit displayed family; member types and common actual first-use exposed, not one atomic glyph. Family listing, not one atomic glyph. Author macro \texttt{\detokenize{\mathsf{C}}} not carried to the members. \\
\phantomsection\label{entry:symbol-registry-07609}\([w_j]\) & \texttt{\detokenize{none}} & binary cylinder extending a finite word & P4 Weihrauch Cantor coding & object & \hyperref[def:p04-wc-negative]{\ref*{def:p04-wc-negative}} & \thepage & Distinct typed argument from the path set bracket of a tree. \\
\phantomsection\label{entry:symbol-registry-07610}\(\mathsf{LPO}\) & \texttt{\detokenize{\mathsf{LPO}}} & native 2009 existence-of-a-zero coordinate decision & P4 Weihrauch omniscience & object & \hyperref[def:p04-wc-omniscience]{\ref*{def:p04-wc-omniscience}} & \thepage & Output zero iff some coordinate is zero; literal map differs from the zero-stream variant. \\
\phantomsection\label{entry:symbol-registry-07611}\(\mathsf{LPO}_0\) & \texttt{\detokenize{\mathsf{LPO}_0}} & alternative characteristic test of the all-zero stream & P4 Weihrauch omniscience & object & \hyperref[def:p04-wc-omniscience]{\ref*{def:p04-wc-omniscience}} & \thepage & Two explicit strong translations preserve the difference in literal input clauses. \\
\phantomsection\label{entry:symbol-registry-07612}\(\mathsf{LLPO}\) & \texttt{\detokenize{\mathsf{LLPO}}} & choice of a parity subsequence whose entries all vanish & P4 Weihrauch omniscience & object & \hyperref[def:p04-wc-omniscience]{\ref*{def:p04-wc-omniscience}} & \thepage & At most one nonzero input entry; all-zero input permits both answers. \\
\phantomsection\label{entry:symbol-registry-07613}\(\mathsf{MLPO}_n\) & \texttt{\detokenize{\mathsf{MLPO}_n}} & 2010 choice of an all-zero stream in a finite tuple & P4 Weihrauch omniscience & object & \hyperref[def:p04-wc-mlpo]{\ref*{def:p04-wc-mlpo}} & \thepage & One-based answer labels; at least one zero component and no single-marker restriction on other components. \\
\phantomsection\label{entry:symbol-registry-07614}\(\mathsf{LLPO}\) & \texttt{\detokenize{\mathsf{LLPO}}} & 2010 two-component MLPO presentation & P4 Weihrauch omniscience & object & \hyperref[def:p04-wc-mlpo]{\ref*{def:p04-wc-mlpo}} & \thepage & Same native glyph with a separately stated tuple domain and one-based answers; strongly equivalent not literally identical to the 2009 presentation. \\
\phantomsection\label{entry:symbol-registry-07615}\(\mathsf{LLPO}_{\mathbb R}\) & \texttt{\detokenize{\mathsf{LLPO}_{\mathbb R}}} & total real weak-comparison relation & P4 Weihrauch real comparison & object & \hyperref[def:p04-wc-real-llpo]{\ref*{def:p04-wc-real-llpo}} & \thepage & Zero permits both answers; the corrected parity translator is not a sign-equality decision. \\
\phantomsection\label{entry:symbol-registry-07616}\(\lim\) & \texttt{\detokenize{\lim}} & coordinatewise eventually constant Baire limit & P4 Weihrauch limits & object & \hyperref[def:p04-wc-limit-types]{\ref*{def:p04-wc-limit-types}} & \thepage & Coordinate stabilization without a supplied modulus. \\
\phantomsection\label{entry:symbol-registry-07617}\(\lim_{\mathbb N}\) & \texttt{\detokenize{\lim_{\mathbb N}}} & eventually constant scalar natural-number limit & P4 Weihrauch limits & object & \hyperref[def:p04-wc-limit-types]{\ref*{def:p04-wc-limit-types}} & \thepage & Separate output type from the source discrete topology on entire Baire streams. \\
\phantomsection\label{entry:symbol-registry-07618}\(\lim_\Delta\) & \texttt{\detokenize{\lim_\Delta}} & source whole-stream eventually constant Baire limit & P4 Weihrauch limits & object & \hyperref[def:p04-wc-limit-types]{\ref*{def:p04-wc-limit-types}} & \thepage & Ordinarily equivalent but not strongly reducible to discrete natural choice. \\
\phantomsection\label{entry:symbol-registry-07619}\(\delta_{\mathrm{Tr}}\) & \texttt{\detokenize{\delta_{\mathrm{Tr}}}} & characteristic representation and carrier of binary trees symbol family member \(\delta_{\mathrm{Tr}}\) & P4 Weihrauch binary trees & object & \hyperref[def:p04-wc-wkl]{\ref*{def:p04-wc-wkl}} & \thepage & Finite-node membership information rather than only a negative path-set description. Explicit displayed family; member types and common actual first-use exposed, not one atomic glyph. Family listing, not one atomic glyph. Author macro \texttt{\detokenize{\delta\_{\mathrm{Tr}};\mathrm{Tr}}} not carried to the members. \\
\phantomsection\label{entry:symbol-registry-07620}\(\mathrm{Tr}\) & \texttt{\detokenize{\mathrm{Tr}}} & characteristic representation and carrier of binary trees symbol family member \(\mathrm{Tr}\) & P4 Weihrauch binary trees & object & \hyperref[def:p04-wc-wkl]{\ref*{def:p04-wc-wkl}} & \thepage & Finite-node membership information rather than only a negative path-set description. Explicit displayed family; member types and common actual first-use exposed, not one atomic glyph. Family listing, not one atomic glyph. Author macro \texttt{\detokenize{\delta\_{\mathrm{Tr}};\mathrm{Tr}}} not carried to the members. \\
\phantomsection\label{entry:symbol-registry-07621}\([T]\) & \texttt{\detokenize{[T]}} & tree path set and path-selection problem symbol family member \([T]\) & P4 Weihrauch binary trees & object & \hyperref[def:p04-wc-wkl]{\ref*{def:p04-wc-wkl}} & \thepage & Native sans-serif problem glyph; classical nonempty path guarantee is not a computed path. Explicit displayed family; member types and common actual first-use exposed, not one atomic glyph. Family listing, not one atomic glyph. Author macro \texttt{\detokenize{\mathsf{WKL}}} not carried to the members. \\
\phantomsection\label{entry:symbol-registry-07622}\(\mathsf{WKL}\) & \texttt{\detokenize{\mathsf{WKL}}} & tree path set and path-selection problem symbol family member \(\mathsf{WKL}\) & P4 Weihrauch binary trees & object & \hyperref[def:p04-wc-wkl]{\ref*{def:p04-wc-wkl}} & \thepage & Native sans-serif problem glyph; classical nonempty path guarantee is not a computed path. Explicit displayed family; member types and common actual first-use exposed, not one atomic glyph. Family listing, not one atomic glyph. Author macro \texttt{\detokenize{\mathsf{WKL}}} not carried to the members. \\
\phantomsection\label{entry:symbol-registry-07623}\(q_{n,k,a}\) & \texttt{\detokenize{q_{n,k,a}}} & native-LPO query for a coordinate stable at a specified value & P4 Weihrauch limit annex & object & \hyperref[annex:p04-wc-parallel-limit]{\ref*{annex:p04-wc-parallel-limit}} & \thepage & Output one exactly when all future entries equal the candidate value; permits strong reconstruction without input. \\
\phantomsection\label{entry:symbol-registry-07624}\(q_s\) & \texttt{\detokenize{q_s}} & guarded single-marker LLPO query at a binary prefix & P4 Weihrauch Cantor-choice annex & object & \hyperref[annex:p04-wc-cantor-llpo]{\ref*{annex:p04-wc-cantor-llpo}} & \thepage & The first discovered empty child determines at most one marker even at an empty parent cylinder. \\
\phantomsection\label{entry:symbol-registry-07625}\(U_v\) & \texttt{\detokenize{U_v}} & open set of binary oracle inputs completing one output digit with value v & P4 Weihrauch weak-single annex & object & \hyperref[annex:p04-wc-weak-single]{\ref*{annex:p04-wc-weak-single}} & \thepage & Finite-computation cylinders and finite covers compute the unique identity-represented Baire digit. \\
\end{longtable}
\endgroup
\Needspace{.4\textheight}
\subsection*{Stable part P5 scopes}
\begingroup\scriptsize\setlength{\tabcolsep}{3pt}
\begin{longtable}{@{}p{0.9in}p{1.3in}p{1.0in}p{0.5in}p{0.3in}p{0.4in}p{0.35in}p{1.65in}@{}}
\toprule
\textbf{Glyph} & \textbf{Macro} & \textbf{Role} & \textbf{Scope} & \textbf{Prec./form} & \textbf{First} & \textbf{Reg. page} & \textbf{Principal uses / notes} \\
\midrule
\endfirsthead
\toprule
\textbf{Glyph} & \textbf{Macro} & \textbf{Role} & \textbf{Scope} & \textbf{Prec./form} & \textbf{First} & \textbf{Reg. page} & \textbf{Principal uses / notes} \\
\midrule
\endhead
\midrule\multicolumn{8}{r}{\emph{continued}}\\\endfoot
\bottomrule\endlastfoot
\phantomsection\label{entry:symbol-registry-00055}\(\ulcorner\varphi\urcorner\) & \texttt{\detokenize{\ulcorner\cdot\urcorner}} & Goedel number & P5 arithmetic-coding to reverse-mathematics & n/a & \hyperref[sec:arithmetic-coding-01]{\ref*{sec:arithmetic-coding-01}} & \thepage & Corner quotes around an expression denote its fixed effective code. \\
\phantomsection\label{entry:symbol-registry-00057}\(\mathsf{Proof}_T\) & \texttt{\detokenize{\mathsf{Proof}_T}} & proof predicate & P5 arithmetic-coding to reverse-mathematics & n/a & \hyperref[sec:arithmetic-coding-01]{\ref*{sec:arithmetic-coding-01}} & \thepage & First the explicitly checked numerical relation on certified proof codes; def:fixed-proof-predicate then uses the same notation for a chosen \(\Sigma^0_1\) formula numeralwise representing that relation in \(\mathsf Q\). Numerical correctness alone does not establish internal derivability conditions for that formula. \\
\phantomsection\label{entry:symbol-registry-00058}\(\mathsf{Prov}_T\) & \texttt{\detokenize{\mathsf{Prov}_T}} & provability predicate & P5 arithmetic-coding to reverse-mathematics & n/a & \hyperref[sec:arithmetic-coding-01]{\ref*{sec:arithmetic-coding-01}} & \thepage & For the earlier certified presentation this abbreviates \(\exists p\,\mathsf{Proof}_T(p,x)\). The later derivability and consistency chapter explicitly declares its chosen formula parameter; its internal hypotheses concern that choice. Arguments are formula codes. \\
\phantomsection\label{entry:symbol-registry-00059}\(\mathsf{Prov}_{\PA}\) & \texttt{\detokenize{\mathsf{Prov}_{\PA}}} & selected standard PA provability predicate & P5 arithmetic-coding to reverse-mathematics & n/a & \hyperref[sec:first-incompleteness-01]{\ref*{sec:first-incompleteness-01}} & \thepage & In the second-incompleteness chapter this is explicitly the literal selected \(\mathsf{Prv}^{19}_{\PA}\). The GL realization writes that selected formula directly. No internal identification with the earlier arbitrary numeralwise representative is assumed. \\
\phantomsection\label{entry:symbol-registry-00060}\(\operatorname{ProofRel}_T,\operatorname{RProofRel}_T,\operatorname{NegCode}\) & \texttt{\detokenize{\operatorname{ProofRel}_T,}\allowbreak\detokenize{\operatorname{RProofRel}_T,}\allowbreak\detokenize{\operatorname{NegCode}}} & metalinguistic primitive-recursive proof relations and negation-code function & P5 arithmetic-coding to reverse-mathematics & meta & \hyperref[sec:p05-independence-catalog]{\ref*{sec:p05-independence-catalog}} & \thepage & Relations and function on natural-number codes used before choosing object-language representatives; none is a term or predicate of the fixed arithmetic language. \\
\phantomsection\label{entry:symbol-registry-00061}\(\mathsf{RProof}_T\) & \texttt{\detokenize{\mathsf{RProof}_T}} & Rosser proof formula & P5 arithmetic-coding to reverse-mathematics & n/a & \hyperref[sec:p05-independence-catalog]{\ref*{sec:p05-independence-catalog}} & \thepage & A fixed Sigma-zero-one formula numeralwise representing the primitive-recursive Rosser proof relation. \\
\phantomsection\label{entry:symbol-registry-00062}\(\mathsf{RProv}_T\) & \texttt{\detokenize{\mathsf{RProv}_T}} & Rosser provability predicate & P5 arithmetic-coding to reverse-mathematics & n/a & \hyperref[sec:p05-independence-catalog]{\ref*{sec:p05-independence-catalog}} & \thepage & A proof of the coded sentence exists whose code is strictly smaller than every proof code for its negation; distinct from ordinary provability. \\
\phantomsection\label{entry:symbol-registry-00063}\(\mathsf{Con}(T)\) & \texttt{\detokenize{\mathsf{Con}}} & consistency sentence & Legacy pre-batch4 P5-P7.6 reader meaning & n/a & \hyperref[sec:first-incompleteness-01]{\ref*{sec:first-incompleteness-01}} & \thepage & Abbreviates \(\neg\mathsf{Prov}_T(\overline{\ulcorner\bottom_A\urcorner})\) for the specified formula parameter. Standard PA uses the selected tag-19 existential graph; P7 names its theory argument explicitly.  Numeric chapter scope migrated to its frozen pre-batch4 source-labelled reader regions; glyph, role and precedence unchanged. \\
\phantomsection\label{entry:symbol-registry-00064}\(\mathsf{Con}_{\mathsf{Proof}_T}(T)\) & \texttt{\detokenize{\mathsf{Con}_{\mathsf{Proof}_T}}} & presentation-indexed consistency sentence & P5 arithmetic-coding to reverse-mathematics & n/a & \hyperref[sec:p05-decidable-arithmetics]{\ref*{sec:p05-decidable-arithmetics}} & \thepage & Exposes the proof formula. It equals \(\mathsf{Con}(T)\) when the chosen provability formula is the existential closure of that proof formula. Standard PA uses the selected tag-19 pair; a certified representative retains its own indexed sentence. \\
\phantomsection\label{entry:symbol-registry-00078}\(\varepsilon_0\) & \texttt{\detokenize{\varepsilon_0}} & least epsilon ordinal & Parts 5–7 from arithmetic-coding & n/a & \hyperref[sec:p05-provability-logic]{\ref*{sec:p05-provability-logic}} & \thepage & Least fixed point of ordinal exponentiation base \(\omega\); calibrates PA in P7.5. \\
\phantomsection\label{entry:symbol-registry-00082}\(\leq\) & \texttt{\detokenize{\leq}} & object-language arithmetic order abbreviation & P5 arithmetic-coding & 2 & \hyperref[sec:arithmetic-coding-01]{\ref*{sec:arithmetic-coding-01}} & \thepage & The P4.3 abbreviation imported for the Q and Rosser chapters. \\
\phantomsection\label{entry:symbol-registry-00083}\(\leq\) & \texttt{\detokenize{\leq}} & canonical ordering & Parts 5–19 from first-incompleteness & 2 & \hyperref[sec:p06-source-map]{\ref*{sec:p06-source-map}} & \thepage & Ambient order on declared numerical and order-theoretic objects, including ordinals, cardinals, degrees, and frame orders. \\
\phantomsection\label{entry:symbol-registry-00091}\(\Box\) & \texttt{\detokenize{\Box}} & necessity operator & Parts 5–8 from arithmetic-coding to conditionals-deontic-free & 3 & \hyperref[sec:second-incompleteness-01]{\ref*{sec:second-incompleteness-01}} & \thepage & Prefix unary at the negation rank; modal and provability readings are fixed by the active system. \\
\phantomsection\label{entry:symbol-registry-00093}\(\Diamond\) & \texttt{\detokenize{\Diamond}} & possibility operator & Parts 5–8 from arithmetic-coding to conditionals-deontic-free & 3 & \hyperref[sec:second-incompleteness-01]{\ref*{sec:second-incompleteness-01}} & \thepage & Prefix unary at the negation rank; \(\neg\Box\neg\) where defined. \\
\phantomsection\label{entry:symbol-registry-00323}\(\operatorname{otp}(A,<)\) & \texttt{\detokenize{\operatorname{otp}}} & order type of a declared linear order & Parts 5–7 from provability-logic & n/a & \hyperref[sec:p05-provability-logic]{\ref*{sec:p05-provability-logic}} & \thepage & Metalinguistic isomorphism type of the displayed ordered set; not an ordinal unless the order is well founded. \\
\phantomsection\label{entry:symbol-registry-00324}\(\mathcal I,\varphi_{\mathrm{dom}},\psi^{\mathcal I}\) & \texttt{\detokenize{\mathcal I,}\allowbreak\detokenize{\varphi_{\mathrm{dom}},}\allowbreak\detokenize{\psi^{\mathcal I}}} & relative-interpretation data and formula translation & P5 provability-logic & n/a & \hyperref[sec:p05-provability-logic]{\ref*{sec:p05-provability-logic}} & \thepage & A one-dimensional parameter-free interpretation; the domain formula and recursive translation are fixed at def:p05-relative-interpretation. \\
\phantomsection\label{entry:symbol-registry-00325}\(E(x,y),F_f(\vec x,y)\) & \texttt{\detokenize{E,}\allowbreak\detokenize{F_f}} & interpreted equality and function-graph formulas & P5 provability-logic & n/a & \hyperref[sec:p05-provability-logic]{\ref*{sec:p05-provability-logic}} & \thepage & The target-language equality surrogate and total functional graph data fixed at def:p05-relative-interpretation; constants are nullary functions. \\
\phantomsection\label{entry:symbol-registry-00449}\(\mathsf Z_2\) & \texttt{\detokenize{\mathsf Z_2}} & full second-order arithmetic & Legacy pre-batch4 P5.7-P7.5 reader meaning & n/a & \hyperref[sec:p05-reverse-mathematics]{\ref*{sec:p05-reverse-mathematics}} & \thepage & RCA-zero augmented by full induction and full comprehension in the fixed two-sorted language; ambient theory for the displayed reverse-mathematics subsystems.  Numeric chapter scope migrated to its frozen pre-batch4 source-labelled reader regions; glyph, role and precedence unchanged. \\
\phantomsection\label{entry:symbol-registry-00450}\(\mathsf{RCA}_0\) & \texttt{\detokenize{\mathsf{RCA}_0}} & recursive-comprehension base theory & Legacy pre-batch4 P5.7-P7.5 reader meaning & n/a & \hyperref[sec:p05-reverse-mathematics]{\ref*{sec:p05-reverse-mathematics}} & \thepage & The exact two-sorted base with Q number axioms, set extensionality, Sigma-zero-one induction, and Delta-zero-one comprehension is fixed in Part V.  Numeric chapter scope migrated to its frozen pre-batch4 source-labelled reader regions; glyph, role and precedence unchanged. \\
\phantomsection\label{entry:symbol-registry-00451}\(\mathsf{WKL}_0\) & \texttt{\detokenize{\mathsf{WKL}_0}} & weak-Konig-lemma subsystem & Legacy pre-batch4 P5.7-P7.5 reader meaning & n/a & \hyperref[sec:p05-reverse-mathematics]{\ref*{sec:p05-reverse-mathematics}} & \thepage & RCA-zero plus the assertion that every infinite coded binary tree has an infinite path.  Numeric chapter scope migrated to its frozen pre-batch4 source-labelled reader regions; glyph, role and precedence unchanged. \\
\phantomsection\label{entry:symbol-registry-00452}\(\mathsf{ACA}_0\) & \texttt{\detokenize{\mathsf{ACA}_0}} & arithmetical-comprehension theory & Legacy pre-batch4 P5.7-P7.5 reader meaning & n/a & \hyperref[sec:p05-reverse-mathematics]{\ref*{sec:p05-reverse-mathematics}} & \thepage & RCA-zero plus comprehension for every arithmetical formula; its first-order consequences are exactly those of PA.  Numeric chapter scope migrated to its frozen pre-batch4 source-labelled reader regions; glyph, role and precedence unchanged. \\
\phantomsection\label{entry:symbol-registry-00453}\(\mathsf{ACA}\) & \texttt{\detokenize{\mathsf{ACA}}} & full-induction arithmetical-comprehension theory & P5 decidable-arithmetics & n/a & \hyperref[sec:p05-decidable-arithmetics]{\ref*{sec:p05-decidable-arithmetics}} & \thepage & ACA-zero strengthened by induction for every formula of the two-sorted language; used to calibrate the uniform Paris-Harrington principle. \\
\phantomsection\label{entry:symbol-registry-00454}\(\mathsf{ATR}_0\) & \texttt{\detokenize{\mathsf{ATR}_0}} & arithmetical-transfinite-recursion theory & Legacy pre-batch4 P5.7-P7.5 reader meaning & n/a & \hyperref[sec:p05-reverse-mathematics]{\ref*{sec:p05-reverse-mathematics}} & \thepage & RCA-zero plus the arithmetical-transfinite-recursion scheme along coded well-orders.  Numeric chapter scope migrated to its frozen pre-batch4 source-labelled reader regions; glyph, role and precedence unchanged. \\
\phantomsection\label{entry:symbol-registry-00455}\(\Pi^1_1\text{-}\mathsf{CA}_0\) & \texttt{\detokenize{\Pi^1_1\text{-}\mathsf{CA}_0}} & Pi-one-one comprehension theory & Legacy pre-batch4 P5.7-P7.5 reader meaning & n/a & \hyperref[sec:p05-reverse-mathematics]{\ref*{sec:p05-reverse-mathematics}} & \thepage & RCA-zero plus comprehension for formulas with universal set quantifiers over an arithmetical matrix.  Numeric chapter scope migrated to its frozen pre-batch4 source-labelled reader regions; glyph, role and precedence unchanged. \\
\phantomsection\label{entry:symbol-registry-00888}\(C(a,b)\) & \texttt{\detokenize{C}} & positive cons constructor for arithmetic syntax codes & P5 arithmetic-coding & meta & \hyperref[sec:arithmetic-coding-01]{\ref*{sec:arithmetic-coding-01}} & \thepage & In sec:arithmetic-coding-04, \(C(a,b)=1+\operatorname{pair}_{\N}(a,b)\); its two arguments are strictly smaller than its value. It is a numerical operation, not an arithmetic-language function symbol. \\
\phantomsection\label{entry:symbol-registry-00889}\([],[a_0,\ldots,a_{k-1}]\) & \texttt{\detokenize{[],}\allowbreak\detokenize{[a_0,}\allowbreak\detokenize{\ldots,}\allowbreak\detokenize{a_{k-1}]}} & finite-list codes for arithmetic syntax and derivations & P5 arithmetic-coding & meta & \hyperref[sec:arithmetic-coding-01]{\ref*{sec:arithmetic-coding-01}} & \thepage & Zero codes the empty list; a nonempty list uses nested positive cons. These governing list codes are distinct from the separate prime-power illustration and the Part IV computation-history witnesses. \\
\phantomsection\label{entry:symbol-registry-00890}\(h(c),t(c),t^j(c)\) & \texttt{\detokenize{h,}\allowbreak\detokenize{t,}\allowbreak\detokenize{t^j}} & list projections and iterated tail & P5 arithmetic-coding & meta & \hyperref[sec:arithmetic-coding-01]{\ref*{sec:arithmetic-coding-01}} & \thepage & For positive \(c\), head and tail invert \(C\); both have value zero at zero. The superscript on the tail denotes iteration. \\
\phantomsection\label{entry:symbol-registry-00891}\(\ell(c),(c)_j\) & \texttt{\detokenize{\ell,}\allowbreak\detokenize{(c)_j}} & length and indexed entry of a coded list & P5 arithmetic-coding & meta & \hyperref[sec:arithmetic-coding-01]{\ref*{sec:arithmetic-coding-01}} & \thepage & The least tail-iteration count reaching zero is the length. Entries use zero-based indices and return zero at or beyond the length. \\
\phantomsection\label{entry:symbol-registry-00892}\(N_k(a_0,\ldots,a_{r-1})\) & \texttt{\detokenize{N_k}} & tagged numerical constructor & P5 arithmetic-coding & meta & \hyperref[sec:arithmetic-coding-01]{\ref*{sec:arithmetic-coding-01}} & \thepage & Equals \(C(k,[a_0,\ldots,a_{r-1}])\). Tags zero through twelve encode arithmetic term/formula constructors with the exact listed arities; different record sorts have their own tags and validation. \\
\phantomsection\label{entry:symbol-registry-00893}\(N_{13}(a),N_{14}([a_0,\ldots,a_{k-1}])\) & \texttt{\detokenize{N_{13},}\allowbreak\detokenize{N_{14}}} & symbol and token-string codes & P5 arithmetic-coding & meta & \hyperref[sec:arithmetic-coding-01]{\ref*{sec:arithmetic-coding-01}} & \thepage & Symbol and string sorts are disjoint from tree codes. A formula's governing code is its expression tree; canonical spelling and parsing provide the proved conversions. \\
\phantomsection\label{entry:symbol-registry-00894}\(L_j(b),B(k,q)\) & \texttt{\detokenize{L_j,}\allowbreak\detokenize{B}} & list majorant and primitive-recursive parsing bound & P5 arithmetic-coding & meta & \hyperref[sec:arithmetic-coding-01]{\ref*{sec:arithmetic-coding-01}} & \thepage & The recurrences in sec:arithmetic-coding-04 bound a tree code for a string of \(k\) tokens bounded by \(q\). This is a numerical search bound, not a polynomial-time claim. \\
\phantomsection\label{entry:symbol-registry-00895}\(\nu(n)\) & \texttt{\detokenize{\nu}} & code of an arithmetic numeral & P5 arithmetic-coding & meta & \hyperref[sec:arithmetic-coding-01]{\ref*{sec:arithmetic-coding-01}} & \thepage & Defined by \(\nu(0)=N_1()\) and \(\nu(n+1)=N_2(\nu(n))\); it is the numerical code of \(\overline n\), not the numeral term itself. \\
\phantomsection\label{entry:symbol-registry-00896}\(N_{17}(i),N_{18}(j)\) & \texttt{\detokenize{N_{17},}\allowbreak\detokenize{N_{18}}} & free and bound occurrence comparison records & P5 arithmetic-coding & meta & \hyperref[sec:arithmetic-coding-01]{\ref*{sec:arithmetic-coding-01}} & \thepage & Alpha comparison retains the index \(i\) of a free variable and the zero-based nearest-binder stack position \(j\) of a bound occurrence. These are comparison records, not formula codes. \\
\phantomsection\label{entry:symbol-registry-00897}\(N_{11}(b),N_{12}(b)\) & \texttt{\detokenize{N_{11},}\allowbreak\detokenize{N_{12}}} & binder comparison records without variable names & P5 arithmetic-coding & meta & \hyperref[sec:arithmetic-coding-01]{\ref*{sec:arithmetic-coding-01}} & \thepage & Only in alpha comparison, the sole argument \(b\) is a matrix record. Formula-code quantifiers instead have two arguments, a binder index and a matrix code. \\
\phantomsection\label{entry:symbol-registry-00898}\(N_{15}(R)\) & \texttt{\detokenize{N_{15}}} & natural-deduction tree record & P5 arithmetic-coding & meta & \hyperref[sec:arithmetic-coding-01]{\ref*{sec:arithmetic-coding-01}} & \thepage & The nonempty list \(R\) contains postorder rows. Earlier-premise indices and the exact parent-count conditions encode one connected tree with its root last. \\
\phantomsection\label{entry:symbol-registry-00899}\([r,a,J,U,l]\) & \texttt{\detokenize{r,}\allowbreak\detokenize{a,}\allowbreak\detokenize{J,}\allowbreak\detokenize{U,}\allowbreak\detokenize{l}} & five fields of a derivation row & P5 arithmetic-coding & meta & \hyperref[sec:arithmetic-coding-01]{\ref*{sec:arithmetic-coding-01}} & \thepage & Rule tag, conclusion formula code, ordered premise-row list, parameter list and assumption marker, respectively. A marker is zero when unlabelled and \(i+1\) for label \(i\); it must be zero at other rules. \\
\phantomsection\label{entry:symbol-registry-00900}\(O_j,D_i^\theta(O)\) & \texttt{\detokenize{O_j,}\allowbreak\detokenize{D_i^\theta}} & open assumption occurrences and labelled discharge & P5 arithmetic-coding & meta & \hyperref[sec:arithmetic-coding-01]{\ref*{sec:arithmetic-coding-01}} & \thepage & The open list contains leaf indices, retaining occurrences and branch ownership. The discharge filter removes exactly leaves with marker \(i+1\) and literal formula \(\theta\), and every requested discharge must be nonvacuous. \\
\phantomsection\label{entry:symbol-registry-00901}\(e_T\) & \texttt{\detokenize{e_T}} & fixed domain index for theory axiom codes & P5 arithmetic-coding & meta & \hyperref[sec:arithmetic-coding-01]{\ref*{sec:arithmetic-coding-01}} & \thepage & In the fixed Part IV numbering, the program's domain is exactly the codes of axiom sentences in the c.e. set \(T\). This does not define \(T\) as its deductive closure. \\
\phantomsection\label{entry:symbol-registry-00902}\(V(e,a,s)\) & \texttt{\detokenize{V}} & numerical truth indicator for the fixed computation formula & P5 arithmetic-coding & meta & \hyperref[sec:arithmetic-coding-01]{\ref*{sec:arithmetic-coding-01}} & \thepage & The primitive-recursive zero/one value of the accepted bounded \(\mathsf{Comp}\) formula on numerical arguments. It is not a uniform truth predicate for arbitrary arithmetic formulas. \\
\phantomsection\label{entry:symbol-registry-00903}\(N_{16}(d,W)\) & \texttt{\detokenize{N_{16},}\allowbreak\detokenize{d,}\allowbreak\detokenize{W}} & certified theory-proof code & P5 arithmetic-coding & meta & \hyperref[sec:arithmetic-coding-01]{\ref*{sec:arithmetic-coding-01}} & \thepage & The derivation record \(d\) has a sentence root; \(W\) contains exactly one accepting-history witness per surviving open axiom leaf in its ordered list. The separately fixed finite-axiom convention instead uses \(W=0\) and direct membership. \\
\phantomsection\label{entry:symbol-registry-00934}\(d,f_d,\Gamma_d\) & \texttt{\detokenize{d,}\allowbreak\detokenize{f_d,}\allowbreak\detokenize{\Gamma_d}} & finite primitive-recursive description function and arithmetic graph & P5 arithmetic-coding & meta & \hyperref[sec:arithmetic-coding-01]{\ref*{sec:arithmetic-coding-01}} & \thepage & The description is a fixed finite metalinguistic tree with arities, its function is numerical, and its graph is the specified arithmetic formula. The index is not an arbitrary formula, machine index or arithmetic variable. \\
\phantomsection\label{entry:symbol-registry-00935}\(\bar x,\bar a,\overline{\bar a},\bar b\) & \texttt{\detokenize{\bar x,}\allowbreak\detokenize{\bar a,}\allowbreak\detokenize{\overline{\bar a},}\allowbreak\detokenize{\bar b}} & argument tuples and numeral substitution in representation & P5 arithmetic-coding & meta & \hyperref[sec:arithmetic-coding-01]{\ref*{sec:arithmetic-coding-01}} & \thepage & In the fixed-arity representation statements, barred argument lists are tuples. The overlined numerical tuple means its componentwise numeral terms; \(\bar b\) is the numeral for the single numerical output \(b\). \\
\phantomsection\label{entry:symbol-registry-00936}\(A,H,K,C;\Gamma_A,\Gamma_H,\Gamma_K,\Gamma_C\) & \texttt{\detokenize{A,}\allowbreak\detokenize{H,}\allowbreak\detokenize{K,}\allowbreak\detokenize{C,}\allowbreak\detokenize{\Gamma_A,}\allowbreak\detokenize{\Gamma_H,}\allowbreak\detokenize{\Gamma_K,}\allowbreak\detokenize{\Gamma_C}} & specified descriptions for addition triangular numbers pairing and cons & P5 arithmetic-coding & meta & \hyperref[sec:arithmetic-coding-01]{\ref*{sec:arithmetic-coding-01}} & \thepage & The displayed recursion/composition descriptions have the corresponding graph formulas. Their internal equations identify pairing with the existing \(\mathsf P\) graph and positive cons with its successor; function notation is eliminable. \\
\phantomsection\label{entry:symbol-registry-00937}\(T(u)\) & \texttt{\detokenize{T(u)}} & local triangular-value abbreviation in the PA proof & P5 arithmetic-coding & meta & \hyperref[sec:arithmetic-coding-01]{\ref*{sec:arithmetic-coding-01}} & \thepage & Abbreviates the unique natural value satisfying \(2T(u)=u(u+1)\), established directly by PA induction before the general graph theorem. It is unrelated to the theory parameter \(T\). \\
\phantomsection\label{entry:symbol-registry-00938}\(\mathsf{Cong}(r,v,d)\) & \texttt{\detokenize{\mathsf{Cong}}} & arithmetic congruence abbreviation in the extension proof & P5 arithmetic-coding & meta & \hyperref[sec:arithmetic-coding-01]{\ref*{sec:arithmetic-coding-01}} & \thepage & For positive \(d\), abbreviates \(\exists q\,\exists q'\,(r+dq=v+dq')\). Signed coefficients used in its proof are pairs of natural numbers with cross-addition equality. \\
\phantomsection\label{entry:symbol-registry-00939}\(\mathsf{Ev}_\varphi(\bar x,e)\) & \texttt{\detokenize{\mathsf{Ev}_\varphi}} & bounded evidence formula for a fixed Sigma-one formula & P5 arithmetic-coding & meta & \hyperref[sec:arithmetic-coding-01]{\ref*{sec:arithmetic-coding-01}} & \thepage & A structural construction for the separately fixed formula with fresh evidence variable. The original-sentence compiler now turns its evidence into selected tag-19 PA derivations. This is neither a uniform truth predicate nor an internal transfer to the earlier certified-proof representative. \\
\phantomsection\label{entry:symbol-registry-00940}\(\mathsf{mp}(p,q)\) & \texttt{\detokenize{\mathsf{mp}}} & numerical certified-proof modus-ponens constructor & P5 arithmetic-coding & meta & \hyperref[sec:arithmetic-coding-01]{\ref*{sec:arithmetic-coding-01}} & \thepage & The first argument proves an implication and the second its antecedent. The constructor places the antecedent tree first, shifts the implication tree, and concatenates the certificates in that order. Its proved correctness here is external, not PA D2. \\
\phantomsection\label{entry:symbol-registry-00941}\(+\!+\) & \texttt{\detokenize{+\!+}} & local coded-list concatenation in the numerical proof constructor & P5 arithmetic-coding & meta & \hyperref[sec:arithmetic-coding-01]{\ref*{sec:arithmetic-coding-01}} & \thepage & Used only in the displayed construction of the merged row and witness lists. It is the already defined numerical list operation, not arithmetic addition or a new arithmetic function symbol. \\
\phantomsection\label{entry:symbol-registry-00942}\(\rho(p,x),P_\rho(x),B_\sigma(e)\) & \texttt{\detokenize{\rho,}\allowbreak\detokenize{P_\rho,}\allowbreak\detokenize{B_\sigma}} & schematic proof formula existential closure and evidence builder & P5 arithmetic-coding & meta & \hyperref[sec:arithmetic-coding-01]{\ref*{sec:arithmetic-coding-01}} & \thepage & The earlier construction notation is realized by the later selected graph and its fixed original-sentence builders. D1, D2 and D3 concern that literal tag-19 PA predicate; no internal identification with an arbitrary numeralwise representative is inferred. \\
\phantomsection\label{entry:symbol-registry-00943}\(v_i,o,q_1,\ldots,q_7\) & \texttt{\detokenize{v_i,}\allowbreak\detokenize{o,}\allowbreak\detokenize{q_1,}\allowbreak\detokenize{\ldots,}\allowbreak\detokenize{q_7}} & variable zero and seven fixed arithmetic axiom codes & P5 arithmetic-coding & meta & \hyperref[sec:arithmetic-coding-01]{\ref*{sec:arithmetic-coding-01}} & \thepage & The direct PA presentation uses \(v_i=N_0(i)\), \(o=N_1()\), and the seven displayed constructor expressions for its literal arithmetic axioms. They are numbers coding expressions, not the expressions themselves. \\
\phantomsection\label{entry:symbol-registry-00944}\(F(m),U(P,c)\) & \texttt{\detokenize{F,}\allowbreak\detokenize{U}} & sorted free-index list and universal-closing constructor & P5 arithmetic-coding & meta & \hyperref[sec:arithmetic-coding-01]{\ref*{sec:arithmetic-coding-01}} & \thepage & The finite index scan gives increasing distinct free variables. Wrapping the reversed list gives the indicated outside-in universal prefix; the empty list leaves the formula code unchanged. \\
\phantomsection\label{entry:symbol-registry-00945}\(\mathsf{sub}(m,i,t)\) & \texttt{\detokenize{\mathsf{sub}}} & governing capture-avoiding substitution on codes & P5 arithmetic-coding & meta & \hyperref[sec:arithmetic-coding-01]{\ref*{sec:arithmetic-coding-01}} & \thepage & The previously specified deterministic syntax operation used by the induction recognizer; bound occurrences are retained and the source's least-fresh convention is preserved. \\
\phantomsection\label{entry:symbol-registry-00946}\(b(m,i),s(m,i),J(m,i),\mathsf{IndAx}(m,i)\) & \texttt{\detokenize{b,}\allowbreak\detokenize{s,}\allowbreak\detokenize{J,}\allowbreak\detokenize{\mathsf{IndAx}}} & base step core and closed induction-axiom constructors & P5 arithmetic-coding & meta & \hyperref[sec:arithmetic-coding-01]{\ref*{sec:arithmetic-coding-01}} & \thepage & The base substitutes zero, the step substitutes the successor of the induction variable, and the final constructor closes the remaining free parameters in increasing order. A nonformula matrix has output zero. \\
\phantomsection\label{entry:symbol-registry-00947}\(\mathsf{Ax}^{\rm d}_{\PA},\PA^{\rm d}\) & \texttt{\detokenize{\mathsf{Ax}^{\rm d}_{\PA},}\allowbreak\detokenize{\PA^{\rm d}}} & direct numerical axiom recognizer and canonical PA presentation & P5 arithmetic-coding & meta & \hyperref[sec:arithmetic-coding-01]{\ref*{sec:arithmetic-coding-01}} & \thepage & Recognizes exactly the seven fixed codes and reconstructed induction axioms. Noncanonical but derivable closures require actual ND derivations; they are not accepted through an implicit equivalence closure. \\
\phantomsection\label{entry:symbol-registry-00948}\(N_{19}(d,0),\mathsf{StdProof}_{\PA}(p,x)\) & \texttt{\detokenize{N_{19},}\allowbreak\detokenize{\mathsf{StdProof}_{\PA}}} & direct PA proof sort and numerical checker & P5 arithmetic-coding & meta & \hyperref[sec:arithmetic-coding-01]{\ref*{sec:arithmetic-coding-01}} & \thepage & Outer tag 19 and its empty second field distinguish this sort from certified tag 16. The checker keeps all twenty row rules and tests direct axiom membership at surviving open leaves; this notation does not select an arithmetic representative. \\
\phantomsection\label{entry:symbol-registry-00956}\(\mathsf{Prf}^{19}_{\PA},\mathsf{Prv}^{19}_{\PA}\) & \texttt{\detokenize{\mathsf{Prf}^{19}_{\PA},}\allowbreak\detokenize{\mathsf{Prv}^{19}_{\PA}}} & selected tag-19 arithmetic proof and provability formulas & P5 arithmetic-coding to reverse-mathematics & meta & \hyperref[sec:arithmetic-coding-01]{\ref*{sec:arithmetic-coding-01}} & \thepage & These are the actual fixed graph at output one and its existential closure. The completed D1-D3 and formalized Sigma-one results use this same selection, as do the standard PA consistency and GL applications. They are distinct from an arbitrary representative of the numerical checker. \\
\phantomsection\label{entry:symbol-registry-00957}\(d_{19},\chi_{19}\) & \texttt{\detokenize{d_{19},}\allowbreak\detokenize{\chi_{19}}} & fixed formation tree and direct checker characteristic description & P5 arithmetic-coding & meta & \hyperref[sec:arithmetic-coding-01]{\ref*{sec:arithmetic-coding-01}} & \thepage & The finite expansion conventions determine the description tree. Its numerical function is zero or one; its arithmetic graph is supplied by the fixed Gamma construction. \\
\phantomsection\label{entry:symbol-registry-00958}\(\mathsf{mp}_{19},d_{\rm mp}\) & \texttt{\detokenize{\mathsf{mp}_{19},}\allowbreak\detokenize{d_{\rm mp}}} & selected direct proof-joining function and its description & P5 arithmetic-coding & meta & \hyperref[sec:arithmetic-coding-01]{\ref*{sec:arithmetic-coding-01}} & \thepage & The implication proof is the first input and the antecedent proof the second. The output places the antecedent rows first, shifts the implication rows, and appends the actual implication-elimination row; failed input checks return zero. \\
\phantomsection\label{entry:symbol-registry-00959}\(D,Z,M,L,E,I,\mathsf{sg},\mathsf{lt}\) & \texttt{\detokenize{D,}\allowbreak\detokenize{Z,}\allowbreak\detokenize{M,}\allowbreak\detokenize{L,}\allowbreak\detokenize{E,}\allowbreak\detokenize{I,}\allowbreak\detokenize{\mathsf{sg},}\allowbreak\detokenize{\mathsf{lt}}} & local initial arithmetic description library & P5 arithmetic-coding & meta & \hyperref[sec:arithmetic-coding-01]{\ref*{sec:arithmetic-coding-01}} & \thepage & Within the specified-checker construction these names mean predecessor, zero test, multiplication, truncated subtraction, equality test, conditional selection, sign and strict-order test. Their displayed arities and recursion order govern every expansion. \\
\phantomsection\label{entry:symbol-registry-00960}\(h,t,T,\ell,e,V\) & \texttt{\detokenize{h,}\allowbreak\detokenize{t,}\allowbreak\detokenize{T,}\allowbreak\detokenize{\ell,}\allowbreak\detokenize{e,}\allowbreak\detokenize{V}} & specified raw-list decoding and reversal-fold descriptions & P5 arithmetic-coding & meta & \hyperref[sec:arithmetic-coding-01]{\ref*{sec:arithmetic-coding-01}} & \thepage & Here h and t decode head and tail; T iterates tail, ell is length, e is entry, and V is the three-argument reverse-prefix accumulator. Zero and out-of-range defaults are explicitly defined; T is not a theory parameter. \\
\phantomsection\label{entry:symbol-registry-00961}\(\mathsf{rev},\mathsf{cat},\mathsf{map},\mathsf{filter},\mathsf{range},\mathsf{mem},\mathsf{take}\) & \texttt{\detokenize{\mathsf{rev},}\allowbreak\detokenize{\mathsf{cat},}\allowbreak\detokenize{\mathsf{map},}\allowbreak\detokenize{\mathsf{filter},}\allowbreak\detokenize{\mathsf{range},}\allowbreak\detokenize{\mathsf{mem},}\allowbreak\detokenize{\mathsf{take}}} & specified finite list folds and collection descriptions & P5 arithmetic-coding & meta & \hyperref[sec:arithmetic-coding-01]{\ref*{sec:arithmetic-coding-01}} & \thepage & Each map or filter is separately fixed by its expression and parameters. These abbreviations expand into the declared folds and graphs; no uniform arithmetic evaluator for arbitrary functions is assumed. \\
\phantomsection\label{entry:symbol-registry-00962}\(\mathsf{node}_{k,r},k(c),a_j(c)\) & \texttt{\detokenize{\mathsf{node}_{k,r},}\allowbreak\detokenize{k(c),}\allowbreak\detokenize{a_j(c)}} & constructor tag arity test and argument projections & P5 arithmetic-coding & meta & \hyperref[sec:arithmetic-coding-01]{\ref*{sec:arithmetic-coding-01}} & \thepage & The node test requires positive code, exact tag and exact argument-list length. The projections have total raw-list defaults. Matching a constructor pattern includes its tag and arity tests. \\
\phantomsection\label{entry:symbol-registry-00963}\(\mathsf{tm},\mathsf{fm},\mathsf{nn},\mathsf{mx},\mathsf{av},\mathsf{fv}\) & \texttt{\detokenize{\mathsf{tm},}\allowbreak\detokenize{\mathsf{fm},}\allowbreak\detokenize{\mathsf{nn},}\allowbreak\detokenize{\mathsf{mx},}\allowbreak\detokenize{\mathsf{av},}\allowbreak\detokenize{\mathsf{fv}}} & local fields in the two specified syntax tables & P5 arithmetic-coding & meta & \hyperref[sec:arithmetic-coding-01]{\ref*{sec:arithmetic-coding-01}} & \thepage & The first row records term/formula bits, node count and greatest variable index. The second records all/free occurrence bits for a fixed index. Child fields come from the earlier-row accumulator. \\
\phantomsection\label{entry:symbol-registry-00964}\(\mathsf{Tm},\mathsf{Fm},n,m,\mathsf{AV},\mathsf{FV},\mathsf{Sent}\) & \texttt{\detokenize{\mathsf{Tm},}\allowbreak\detokenize{\mathsf{Fm},}\allowbreak\detokenize{n,}\allowbreak\detokenize{m,}\allowbreak\detokenize{\mathsf{AV},}\allowbreak\detokenize{\mathsf{FV},}\allowbreak\detokenize{\mathsf{Sent}}} & retrieved syntax-description outputs and sentence test & P5 arithmetic-coding & meta & \hyperref[sec:arithmetic-coding-01]{\ref*{sec:arithmetic-coding-01}} & \thepage & These outputs belong to the newly fixed graph descriptions. The occurrence tests distinguish free from all occurrences, including binders; the sentence test requires a formula and an empty free-index list. \\
\phantomsection\label{entry:symbol-registry-00965}\([b,D,V]\) & \texttt{\detokenize{[b,}\allowbreak\detokenize{D,}\allowbreak\detokenize{V]}} & bounded syntax-machine state & P5 arithmetic-coding & meta & \hyperref[sec:arithmetic-coding-01]{\ref*{sec:arithmetic-coding-01}} & \thepage & The state contains a success bit, pending-command list and value stack. In this local state notation D and V denote lists, not the earlier predecessor and fold descriptions. Malformed states and command failures enter the absorbing failed state. \\
\phantomsection\label{entry:symbol-registry-00966}\(\mathsf{sub},\mathsf{ar},\mathsf{ae}\) & \texttt{\detokenize{\mathsf{sub},}\allowbreak\detokenize{\mathsf{ar},}\allowbreak\detokenize{\mathsf{ae}}} & fixed substitution and alpha-comparison machine descriptions & P5 arithmetic-coding & meta & \hyperref[sec:arithmetic-coding-01]{\ref*{sec:arithmetic-coding-01}} & \thepage & The bounded machines reproduce the stated substitution and alpha-record operations. Alpha comparison is restricted to formula inputs. Their unchanged arguments, together with graph uniqueness, support the internal joining proof. \\
\phantomsection\label{entry:symbol-registry-00967}\(\mathsf{Ax}_{19},B(a)\) & \texttt{\detokenize{\mathsf{Ax}_{19},}\allowbreak\detokenize{B(a)}} & selected direct axiom test and induction-pattern test & P5 arithmetic-coding & meta & \hyperref[sec:arithmetic-coding-01]{\ref*{sec:arithmetic-coding-01}} & \thepage & B(a) combines prefix stripping, constructor/field tests and reconstruction equality for induction. Ax19 also tests the seven literal axioms and sentencehood; B(a) is local and differs from a table accumulator. \\
\phantomsection\label{entry:symbol-registry-00968}\(\mathsf{Rows},\mathsf{Tree},\mathsf{pc},\mathsf{root},\mathsf{open}\) & \texttt{\detokenize{\mathsf{Rows},}\allowbreak\detokenize{\mathsf{Tree},}\allowbreak\detokenize{\mathsf{pc},}\allowbreak\detokenize{\mathsf{root},}\allowbreak\detokenize{\mathsf{open}}} & specified row table parent counts and proof-record selectors & P5 arithmetic-coding & meta & \hyperref[sec:arithmetic-coding-01]{\ref*{sec:arithmetic-coding-01}} & \thepage & Rows stores a success bit and surviving open indices per row. Tree additionally requires one parent for each nonroot and zero for the last row. Root and open have zero defaults for the empty record. \\
\phantomsection\label{entry:symbol-registry-00969}\(r_j,f_R,\lambda_R,b_j,O_j\) & \texttt{\detokenize{r_j,}\allowbreak\detokenize{f_R,}\allowbreak\detokenize{\lambda_R,}\allowbreak\detokenize{b_j,}\allowbreak\detokenize{O_j}} & row projections and success/open-occurrence fields & P5 arithmetic-coding & meta & \hyperref[sec:arithmetic-coding-01]{\ref*{sec:arithmetic-coding-01}} & \thepage & The row formula and label lookups refer to the declared raw row list R. The second table field is a list of surviving assumption-row indices; it is not a set of formulas. \\
\phantomsection\label{entry:symbol-registry-00970}\(D_R^{i,\theta}\) & \texttt{\detokenize{D_R^{i,\theta}}} & formula-specific discharge filter on open occurrences & P5 arithmetic-coding & meta & \hyperref[sec:arithmetic-coding-01]{\ref*{sec:arithmetic-coding-01}} & \thepage & Removes exactly occurrences whose row label is Si and whose literal formula is theta. A strict reduction in list length is the separately required nonvacuity test. \\
\phantomsection\label{entry:symbol-registry-00971}\(\mathsf{sh},\mathsf{sr},\mathsf{SR}\) & \texttt{\detokenize{\mathsf{sh},}\allowbreak\detokenize{\mathsf{sr},}\allowbreak\detokenize{\mathsf{SR}}} & premise-index offset and row-block shift descriptions & P5 arithmetic-coding & meta & \hyperref[sec:arithmetic-coding-01]{\ref*{sec:arithmetic-coding-01}} & \thepage & The shift changes only premise references. Formula codes, term/formula parameters and labels remain unchanged; the internal row proof tracks the corresponding shifted open indices. \\
\phantomsection\label{entry:symbol-registry-00982}\(Q_7,e(Q_7,j)\) & \texttt{\detokenize{Q_7,}\allowbreak\detokenize{e(Q_7,}\allowbreak\detokenize{j)}} & fixed seven-axiom constant list and total counter selector & P5 arithmetic-coding & meta & \hyperref[sec:arithmetic-coding-01]{\ref*{sec:arithmetic-coding-01}} & \thepage & In def:p05-checker-axiom-description the seven literal axiom codes form this exact cons-list. Its selector is the already specified total entry function, with zero beyond the seven positions; it fixes the bounded-sum description. \\
\phantomsection\label{entry:symbol-registry-01025}\(v_i,o,s(t),a(t,u),b(t,u)\) & \texttt{\detokenize{v_i,}\allowbreak\detokenize{o,}\allowbreak\detokenize{s,}\allowbreak\detokenize{a,}\allowbreak\detokenize{b}} & term-code abbreviations for the numeral proof builders & P5 arithmetic-coding & meta & \hyperref[sec:arithmetic-coding-01]{\ref*{sec:arithmetic-coding-01}} & \thepage & Variable, zero, successor, addition and multiplication syntax constructors. The numerical operations retain their separately defined names A and M. \\
\phantomsection\label{entry:symbol-registry-01026}\(\mathsf{Eq},\mathsf{bc},\mathsf{Neg},\mathsf{All},\mathsf{Ex},\&,\Rightarrow\) & \texttt{\detokenize{\mathsf{Eq},}\allowbreak\detokenize{\mathsf{bc},}\allowbreak\detokenize{\mathsf{Neg},}\allowbreak\detokenize{\mathsf{All},}\allowbreak\detokenize{\mathsf{Ex},}\allowbreak\detokenize{\&,}\allowbreak\detokenize{\Rightarrow}} & formula-code constructors in the numeral builders & P5 arithmetic-coding & meta & \hyperref[sec:arithmetic-coding-01]{\ref*{sec:arithmetic-coding-01}} & \thepage & Local abbreviations for the fixed equality, bottom, negation, quantifier, conjunction and implication syntax tags. They construct codes; they do not add logical constants or inference rules. \\
\phantomsection\label{entry:symbol-registry-01027}\(\mathsf{entry}(R,j),\mathsf{nodes}(c),f(R)\) & \texttt{\detokenize{\mathsf{entry},}\allowbreak\detokenize{\mathsf{nodes},}\allowbreak\detokenize{f}} & list-entry syntax-node-count and raw-root projections & P5 arithmetic-coding & meta & \hyperref[sec:arithmetic-coding-01]{\ref*{sec:arithmetic-coding-01}} & \thepage & Reuse the previously fixed descriptions. A raw root can be open; it becomes a checked sentence proof only under the additional sentence and axiom-leaf conditions. \\
\phantomsection\label{entry:symbol-registry-01028}\(\mathsf{it}(t,j),\mathsf{pack}(R)\) & \texttt{\detokenize{\mathsf{it},}\allowbreak\detokenize{\mathsf{pack}}} & successor-string code and selected tag-19 proof packaging & P5 arithmetic-coding & meta & \hyperref[sec:arithmetic-coding-01]{\ref*{sec:arithmetic-coding-01}} & \thepage & Iteration starts at t and adds j successor constructors. Packing is exactly \(N_{19}(N_{15}(R),0)\); packaging alone does not establish validity. \\
\phantomsection\label{entry:symbol-registry-01029}\(\mathsf{nc},\mathsf{ns},\mathsf{cost}\) & \texttt{\detokenize{\mathsf{nc},}\allowbreak\detokenize{\mathsf{ns},}\allowbreak\detokenize{\mathsf{cost}}} & no-renaming substitution-table fields & P5 arithmetic-coding & meta & \hyperref[sec:arithmetic-coding-01]{\ref*{sec:arithmetic-coding-01}} & \thepage & The fields are success, substituted code and execution cost. Invalid expression or replacement types give three zero fields. Numeral invariance requires a valid replacement term. \\
\phantomsection\label{entry:symbol-registry-01030}\(\mathsf J_r(c,U;R_0,\ldots,R_{h-1}),l_j,t_j\) & \texttt{\detokenize{\mathsf J_r,}\allowbreak\detokenize{l_j,}\allowbreak\detokenize{t_j}} & fixed-arity raw proof join block lengths and offsets & P5 arithmetic-coding & meta & \hyperref[sec:arithmetic-coding-01]{\ref*{sec:arithmetic-coding-01}} & \thepage & Each premise tree is copied and shifted into its own interval before the final rule row. The arity is zero through three; validity requires nonempty successful premises and the actual last-row test. \\
\phantomsection\label{entry:symbol-registry-01031}\(\mathsf L(c,l),\mathsf R(t)\) & \texttt{\detokenize{\mathsf L,}\allowbreak\detokenize{\mathsf R}} & assumption and identity row templates & P5 arithmetic-coding & meta & \hyperref[sec:arithmetic-coding-01]{\ref*{sec:arithmetic-coding-01}} & \thepage & The stored assumption marker \(Si\) denotes label i; zero means unmarked. The identity template has no premise or open leaf. These names differ from numerical truncated subtraction and arbitrary row lists. \\
\phantomsection\label{entry:symbol-registry-01032}\(\mathsf{FE},\mathsf{FI},\mathsf{EI},\mathsf{EE}\) & \texttt{\detokenize{\mathsf{FE},}\allowbreak\detokenize{\mathsf{FI},}\allowbreak\detokenize{\mathsf{EI},}\allowbreak\detokenize{\mathsf{EE}}} & quantifier elimination and introduction row templates & P5 arithmetic-coding & meta & \hyperref[sec:arithmetic-coding-01]{\ref*{sec:arithmetic-coding-01}} & \thepage & Each carries the selected rule's actual binder, matrix, replacement or eigenvariable parameters. Validity includes the required free-variable or all-variable conditions and nonvacuous discharge. \\
\phantomsection\label{entry:symbol-registry-01033}\(\mathsf{II},\mathsf{NI},\mathsf B,\mathsf{And},\mathsf{MP}^{\rm row}\) & \texttt{\detokenize{\mathsf{II},}\allowbreak\detokenize{\mathsf{NI},}\allowbreak\detokenize{\mathsf B,}\allowbreak\detokenize{\mathsf{And},}\allowbreak\detokenize{\mathsf{MP}^{\rm row}}} & propositional raw-row construction templates & P5 arithmetic-coding & meta & \hyperref[sec:arithmetic-coding-01]{\ref*{sec:arithmetic-coding-01}} & \thepage & The templates retain exact premise order and discharge labels. Raw implication elimination can have an open formula root; sentence-proof modus ponens is its separately guarded packed instance. \\
\phantomsection\label{entry:symbol-registry-01034}\(\mathsf{EEq},\mathsf{Sym},\mathsf{Tr},\mathsf{Cg}_g\) & \texttt{\detokenize{\mathsf{EEq},}\allowbreak\detokenize{\mathsf{Sym},}\allowbreak\detokenize{\mathsf{Tr},}\allowbreak\detokenize{\mathsf{Cg}_g}} & equality elimination symmetry transitivity and context congruence & P5 arithmetic-coding & meta & \hyperref[sec:arithmetic-coding-01]{\ref*{sec:arithmetic-coding-01}} & \thepage & Use exact equality roots and fresh hole indices. Each permitted context has its own fixed constructor expression; repeated premises are copied as trees rather than shared rows. \\
\phantomsection\label{entry:symbol-registry-01035}\(\mathsf A_j,\mathsf A_Z,\mathsf A_L\) & \texttt{\detokenize{\mathsf A_j,}\allowbreak\detokenize{\mathsf A_Z,}\allowbreak\detokenize{\mathsf A_L}} & actual universal-elimination chains for arithmetic instances & P5 arithmetic-coding & meta & \hyperref[sec:arithmetic-coding-01]{\ref*{sec:arithmetic-coding-01}} & \thepage & Start at a literal arithmetic axiom or copy the complete derived addition proof, then append the stated elimination rows. A derived instance never becomes a new assumption axiom. \\
\phantomsection\label{entry:symbol-registry-01036}\(m_Z,m_L,I_Z,I_L,D_Z,D_L,\mathcal A_*\) & \texttt{\detokenize{m_Z,}\allowbreak\detokenize{m_L,}\allowbreak\detokenize{I_Z,}\allowbreak\detokenize{I_L,}\allowbreak\detokenize{D_Z,}\allowbreak\detokenize{D_L,}\allowbreak\detokenize{\mathcal A_*}} & fixed addition matrices induction leaves derived records and allowed leaves & P5 arithmetic-coding & meta & \hyperref[sec:arithmetic-coding-01]{\ref*{sec:arithmetic-coding-01}} & \thepage & The two explicit induction records use their canonical matrix/index pairs. The nine allowed unmarked leaves are the seven arithmetic constants and these two induction axioms, whose variable indices lie in zero and one. \\
\phantomsection\label{entry:symbol-registry-01037}\(R_+,R_\times,R_{\rm right},R_{\rm left}\) & \texttt{\detokenize{R_+,}\allowbreak\detokenize{R_\times,}\allowbreak\detokenize{R_{\rm right},}\allowbreak\detokenize{R_{\rm left}}} & numeral arithmetic and successor-string raw proof recurrences & P5 arithmetic-coding & meta & \hyperref[sec:arithmetic-coding-01]{\ref*{sec:arithmetic-coding-01}} & \thepage & Addition and multiplication normalize closed numerals. The two successor-string orientations additionally permit the specified open terms; those intermediate roots are not asserted to be sentence proofs. \\
\phantomsection\label{entry:symbol-registry-01038}\(P(t,k,b,R,j),\mathsf{Clash}(t,k,b,R)\) & \texttt{\detokenize{P,}\allowbreak\detokenize{\mathsf{Clash}}} & successor peeling recurrence and contradiction constructor & P5 arithmetic-coding & meta & \hyperref[sec:arithmetic-coding-01]{\ref*{sec:arithmetic-coding-01}} & \thepage & Peeling runs for the supplied bound b. Its success invariant requires the stated term domain, \(\mathsf{Tree}(R)=1\), \(k>b\) and the exact equality root; a matching root alone is insufficient. \\
\phantomsection\label{entry:symbol-registry-01039}\(\mathsf{Le}_i(t,u),\mathsf{Lt}_i(t,u)\) & \texttt{\detokenize{\mathsf{Le}_i,}\allowbreak\detokenize{\mathsf{Lt}_i}} & weak and strict comparison codes retaining the original binder & P5 arithmetic-coding & meta & \hyperref[sec:arithmetic-coding-01]{\ref*{sec:arithmetic-coding-01}} & \thepage & Weak comparison uses a witness on the left of addition; strict comparison uses its successor on the right. Numeral substitution preserves the supplied original binder instead of selecting a new alpha-equivalent code. \\
\phantomsection\label{entry:symbol-registry-01040}\(R_=,R_{\ne},R_{\leq}^{+},R_{\leq}^{-},R_<^{+},R_<^{-},R_{\neg\bot}\) & \texttt{\detokenize{R_=,}\allowbreak\detokenize{R_{\ne},}\allowbreak\detokenize{R_{\leq}^{+},}\allowbreak\detokenize{R_{\leq}^{-},}\allowbreak\detokenize{R_<^{+},}\allowbreak\detokenize{R_<^{-},}\allowbreak\detokenize{R_{\neg\bot}}} & signed numeral-literal raw proof builders & P5 arithmetic-coding & meta & \hyperref[sec:arithmetic-coding-01]{\ref*{sec:arithmetic-coding-01}} & \thepage & Each numerical success condition selects its exact literal root; failure gives zero where specified. Negated bottom has an explicit two-row proof with empty final open list. \\
\phantomsection\label{entry:symbol-registry-01041}\(t^\#,t^\nu,R_t,\lambda^\nu,d_\lambda\) & \texttt{\detokenize{t^\#,}\allowbreak\detokenize{t^\nu,}\allowbreak\detokenize{R_t,}\allowbreak\detokenize{\lambda^\nu,}\allowbreak\detokenize{d_\lambda}} & fixed-term value numeral instance evaluation proof and literal description & P5 arithmetic-coding & meta & \hyperref[sec:arithmetic-coding-01]{\ref*{sec:arithmetic-coding-01}} & \thepage & Each fixed term or literal has its separately constructed finite description. The arithmetic antecedent is the original formula at numerical parameters; it is not a code or an internal truth predicate. \\
\phantomsection\label{entry:symbol-registry-01042}\(\mathsf{Safe}(R,c)\) & \texttt{\detokenize{\mathsf{Safe}}} & internal successful-tree root and allowed-leaf invariant & P5 arithmetic-coding & meta & \hyperref[sec:arithmetic-coding-01]{\ref*{sec:arithmetic-coding-01}} & \thepage & Combines the exact tree test, root equation and unmarked membership in the fixed allowed-leaf list. Adding a sentence root yields the selected checker after packing. \\
\phantomsection\label{entry:symbol-registry-01043}\(\mathsf{step}^{,w},D,V\) & \texttt{\detokenize{\mathsf{step}^{,w},}\allowbreak\detokenize{D,}\allowbreak\detokenize{V}} & substitution-machine iteration continuation and value lists & P5 arithmetic-coding & meta & \hyperref[sec:arithmetic-coding-01]{\ref*{sec:arithmetic-coding-01}} & \thepage & The execution invariant quantifies over raw continuation and value lists. It returns the computed value before executing the continuation and fits inside the selected machine's fixed time bound. \\
\phantomsection\label{entry:symbol-registry-01044}\(G_<,G_{\leq}\) & \texttt{\detokenize{G_<,}\allowbreak\detokenize{G_{\leq}}} & marked strict and weak comparison guards in the bounded compiler & P5 arithmetic-coding & meta & \hyperref[sec:arithmetic-coding-01]{\ref*{sec:arithmetic-coding-01}} & \thepage & Retain the original comparison binder distinct from the outer quantified variable; the bound term excludes both variables where required. The code is the specified expanded comparison, not a new primitive atom. \\
\phantomsection\label{entry:symbol-registry-01045}\(c\mathbin{|}d\) & \texttt{\detokenize{\mathbin{|}}} & disjunction syntax-code abbreviation in bounded proof templates & P5 arithmetic-coding & meta & \hyperref[sec:arithmetic-coding-01]{\ref*{sec:arithmetic-coding-01}} & \thepage & Abbreviates \(N_9(c,d)\) in the local code construction. The arguments are formula codes, and this occurrence has no cardinality role. \\
\phantomsection\label{entry:symbol-registry-01046}\(\mathsf P_0,\mathsf P_1,\mathsf{OL},\mathsf{OR},\mathsf{OE}\) & \texttt{\detokenize{\mathsf P_0,}\allowbreak\detokenize{\mathsf P_1,}\allowbreak\detokenize{\mathsf{OL},}\allowbreak\detokenize{\mathsf{OR},}\allowbreak\detokenize{\mathsf{OE}}} & conjunction projection and disjunction raw-row templates & P5 arithmetic-coding & meta & \hyperref[sec:arithmetic-coding-01]{\ref*{sec:arithmetic-coding-01}} & \thepage & Use the selected tags and premise/parameter order. Projections return zero on a failed conjunction-root pattern; elimination requires actual surviving marked cases. \\
\phantomsection\label{entry:symbol-registry-01047}\(\mathsf{EF},\mathsf{RA},\mathsf{Keep},\mathsf{EM}\) & \texttt{\detokenize{\mathsf{EF},}\allowbreak\detokenize{\mathsf{RA},}\allowbreak\detokenize{\mathsf{Keep},}\allowbreak\detokenize{\mathsf{EM}}} & bottom elimination reductio assumption retention and excluded-middle records & P5 arithmetic-coding & meta & \hyperref[sec:arithmetic-coding-01]{\ref*{sec:arithmetic-coding-01}} & \thepage & Keep preserves an otherwise unused open assumption through conjunction and projection. EM is a specified copied-row construction with both negative-major occurrences discharged; it is not an additional axiom leaf. \\
\phantomsection\label{entry:symbol-registry-01048}\(\mathcal D(t,n)\) & \texttt{\detokenize{\mathcal D}} & finite equality-case disjunction code & P5 arithmetic-coding & meta & \hyperref[sec:arithmetic-coding-01]{\ref*{sec:arithmetic-coding-01}} & \thepage & The zero case is bottom; each successor appends the equality with the next numeral on the right. The argument t is a term code. \\
\phantomsection\label{entry:symbol-registry-01049}\(g_\epsilon,Z_\epsilon,S_\epsilon,\mathsf C_\epsilon,F_\epsilon\) & \texttt{\detokenize{g_\epsilon,}\allowbreak\detokenize{Z_\epsilon,}\allowbreak\detokenize{S_\epsilon,}\allowbreak\detokenize{\mathsf C_\epsilon,}\allowbreak\detokenize{F_\epsilon}} & two addition orientations axiom chains and sum-exhaustion recursions & P5 arithmetic-coding & meta & \hyperref[sec:arithmetic-coding-01]{\ref*{sec:arithmetic-coding-01}} & \thepage & The orientation is a fixed metalinguistic choice. Each recurrence has three distinct indices at least three, a specified stage and zero failure value; the reversed orientation copies the accepted left-addition proofs. \\
\phantomsection\label{entry:symbol-registry-01050}\(\mathsf{CS},\mathsf E_\diamond,N_<,N_{\leq}\) & \texttt{\detokenize{\mathsf{CS},}\allowbreak\detokenize{\mathsf E_\diamond,}\allowbreak\detokenize{N_<,}\allowbreak\detokenize{N_{\leq}}} & strict-sum and range-exhaustion records and case counts & P5 arithmetic-coding & meta & \hyperref[sec:arithmetic-coding-01]{\ref*{sec:arithmetic-coding-01}} & \thepage & Strict range has n cases and weak range has its successor number of cases. The zero strict case derives bottom; the zero weak case retains equality to zero. \\
\phantomsection\label{entry:symbol-registry-01051}\(\mathsf{Cases},\mathsf{CI},\mathsf U_\diamond\) & \texttt{\detokenize{\mathsf{Cases},}\allowbreak\detokenize{\mathsf{CI},}\allowbreak\detokenize{\mathsf U_\diamond}} & finite case elimination implication and quantified proof constructors & P5 arithmetic-coding & meta & \hyperref[sec:arithmetic-coding-01]{\ref*{sec:arithmetic-coding-01}} & \thepage & The numeral-instance proofs are premises of an explicit elimination tree. The final implication and universal introduction retain the original bound term and binder with all required discharge and freshness conditions. \\
\phantomsection\label{entry:symbol-registry-01052}\(b_\varphi,\varphi^\nu,R_\varphi^+,R_\varphi^-,d_\varphi^+,d_\varphi^-\) & \texttt{\detokenize{b_\varphi,}\allowbreak\detokenize{\varphi^\nu,}\allowbreak\detokenize{R_\varphi^+,}\allowbreak\detokenize{R_\varphi^-,}\allowbreak\detokenize{d_\varphi^+,}\allowbreak\detokenize{d_\varphi^-}} & fixed bounded-formula truth description literal code and signed proof descriptions & P5 arithmetic-coding & meta & \hyperref[sec:arithmetic-coding-01]{\ref*{sec:arithmetic-coding-01}} & \thepage & A separate finite description is built for each fixed formula. Its arithmetic truth assertion and its formula code have distinct roles; no internal uniform evaluator over formula descriptions is asserted. \\
\phantomsection\label{entry:symbol-registry-01096}\(\mathcal E_\theta,e_\theta\) & \texttt{\detokenize{\mathcal E_\theta,}\allowbreak\detokenize{e_\theta}} & literal-positive bounded evidence formula and its Boolean description & P5 arithmetic-coding & meta & \hyperref[sec:arithmetic-coding-01]{\ref*{sec:arithmetic-coding-01}} & \thepage & Each fixed positive formation has its own formula and finite description. Tuple projections place values in the actual free-index order, even when the evidence variable is not last in that order. \\
\phantomsection\label{entry:symbol-registry-01097}\(\pi_0,\pi_1,\beta_*\) & \texttt{\detokenize{\pi_0,}\allowbreak\detokenize{\pi_1,}\allowbreak\detokenize{\beta_*}} & original pairing inverses and bounded beta evidence selector & P5 arithmetic-coding & meta & \hyperref[sec:arithmetic-coding-01]{\ref*{sec:arithmetic-coding-01}} & \thepage & The pairing projections are \(h(Sz)\) and \(t(Sz)\), rather than head and tail of z itself. The bounded beta search retains the original formula and distinguishes zero values from its search status. \\
\phantomsection\label{entry:symbol-registry-01098}\(V_\theta,Q_\theta,d_\theta\) & \texttt{\detokenize{V_\theta,}\allowbreak\detokenize{Q_\theta,}\allowbreak\detokenize{d_\theta}} & raw positive evidence proof proof-list fold and guarded packed description & P5 arithmetic-coding & meta & \hyperref[sec:arithmetic-coding-01]{\ref*{sec:arithmetic-coding-01}} & \thepage & The fixed evidence guard precedes each constructor. The bounded universal collects actual child proof records and uses the accepted arbitrary-matrix quantified constructor. \\
\phantomsection\label{entry:symbol-registry-01099}\(\mathcal N^+,\mathcal N^-\) & \texttt{\detokenize{\mathcal N^+,}\allowbreak\detokenize{\mathcal N^-}} & signed fixed negation-normal-form operations & P5 arithmetic-coding & meta & \hyperref[sec:arithmetic-coding-01]{\ref*{sec:arithmetic-coding-01}} & \thepage & These are syntax operations on the fixed formula, distinct from the standard-model notation. The signed compiler consumes their evidence and returns the original signed root. \\
\phantomsection\label{entry:symbol-registry-01100}\(\mathcal E_\theta^\epsilon,e_\theta^\epsilon\) & \texttt{\detokenize{\mathcal E_\theta^\epsilon,}\allowbreak\detokenize{e_\theta^\epsilon}} & signed-normal-form evidence and its fixed Boolean description & P5 arithmetic-coding & meta & \hyperref[sec:arithmetic-coding-01]{\ref*{sec:arithmetic-coding-01}} & \thepage & The sign is a fixed positive or negative construction choice. Bounded-leaf Boolean identities and all permitted quantifier polarities are proved before the original-sentence conclusion. \\
\phantomsection\label{entry:symbol-registry-01101}\(\widehat V_\theta^\epsilon,\widehat Q_\theta^\epsilon,\widehat d_\theta^\epsilon\) & \texttt{\detokenize{\widehat V_\theta^\epsilon,}\allowbreak\detokenize{\widehat Q_\theta^\epsilon,}\allowbreak\detokenize{\widehat d_\theta^\epsilon}} & original-syntax signed proof proof-list and packed descriptions & P5 arithmetic-coding & meta & \hyperref[sec:arithmetic-coding-01]{\ref*{sec:arithmetic-coding-01}} & \thepage & Their roots are the original formula or its literal negation, including original guard and quantifier binders. They do not terminate at a merely equivalent normal-form code. \\
\phantomsection\label{entry:symbol-registry-01102}\(P_{19},g_{19},n_\varphi,g_\varphi,c_\varphi\) & \texttt{\detokenize{P_{19},}\allowbreak\detokenize{g_{19},}\allowbreak\detokenize{n_\varphi,}\allowbreak\detokenize{g_\varphi,}\allowbreak\detokenize{c_\varphi}} & fixed selected provability graph and substituted sentence codes & P5 arithmetic-coding & meta & \hyperref[sec:arithmetic-coding-01]{\ref*{sec:arithmetic-coding-01}} & \thepage & The original expanded graph is fixed before numeral inputs. The outer code \(c_\varphi\) is the literal selected provability sentence, distinct from the code \(n_\varphi\) of its argument sentence. \\
\phantomsection\label{entry:symbol-registry-01103}\(\sigma_\varphi,\gamma_\varphi,d^3_\varphi\) & \texttt{\detokenize{\sigma_\varphi,}\allowbreak\detokenize{\gamma_\varphi,}\allowbreak\detokenize{d^3_\varphi}} & literal provability sentence its open graph matrix and D3 evidence builder & P5 arithmetic-coding & meta & \hyperref[sec:arithmetic-coding-01]{\ref*{sec:arithmetic-coding-01}} & \thepage & One unary builder is specified for each fixed arithmetic sentence. Both provability occurrences in D3 use the same selected existential Gamma formula; no consistency assumption is introduced. \\
\phantomsection\label{entry:symbol-registry-01124}\(s_\theta^+,s_\theta^-,s_\varphi^+\) & \texttt{\detokenize{s_\theta^+,}\allowbreak\detokenize{s_\theta^-,}\allowbreak\detokenize{s_\varphi^+}} & fixed numeral-substitution and negated-root descriptions & P5 arithmetic-coding to reverse-mathematics & meta & \hyperref[sec:arithmetic-coding-01]{\ref*{sec:arithmetic-coding-01}} & \thepage & The positive description substitutes numeral codes in increasing free-index order into the original bounded formula, and the same construction extends to the original Sigma-one formula. The negative description appends the fixed N7 negation constructor; a closed input uses its constant original code. \\
\phantomsection\label{entry:symbol-registry-01125}\(E_\varphi(\vec x,e)\) & \texttt{\detokenize{E_\varphi}} & exact evidence alias in formalized Sigma-one completeness & P5 arithmetic-coding to reverse-mathematics & meta & \hyperref[sec:arithmetic-coding-01]{\ref*{sec:arithmetic-coding-01}} & \thepage & This is precisely the existing fixed evidence formula, including its marked normal form, fresh variables, conjunction associations and original comparison binders. It is not a newly selected equivalent evidence formula. \\
\phantomsection\label{entry:symbol-registry-04760}\(M_d\) & \texttt{\detokenize{M_d}} & fixed compiled numerical register program & local:sec:p05-representability & object & \hyperref[def:p05-rep-register-certificate]{\ref*{def:p05-rep-register-certificate}} & \thepage & Not governing Part IV numbering or an arithmetic symbol \\
\phantomsection\label{entry:symbol-registry-04761}\(\operatorname{Row}_d\) & \texttt{\detokenize{\operatorname{Row}_d}} & bounded formula extracting one register-program row & local:sec:p05-representability & relation & \hyperref[def:p05-rep-register-certificate]{\ref*{def:p05-rep-register-certificate}} & \thepage & Expands into actual tagged beta extraction formulas \\
\phantomsection\label{entry:symbol-registry-04762}\(\operatorname{Step}_d\) & \texttt{\detokenize{\operatorname{Step}_d}} & quantifier-free transition formula for a fixed program & local:sec:p05-representability & relation & \hyperref[def:p05-rep-register-certificate]{\ref*{def:p05-rep-register-certificate}} & \thepage & Finite disjunction over increment and zero-test/decrement instructions \\
\phantomsection\label{entry:symbol-registry-04763}\(\operatorname{Halt}_d\) & \texttt{\detokenize{\operatorname{Halt}_d}} & finite disjunction of fixed program halt labels & local:sec:p05-representability & relation & \hyperref[def:p05-rep-register-certificate]{\ref*{def:p05-rep-register-certificate}} & \thepage & Expanded quantifier-free formula not a program execution primitive \\
\phantomsection\label{entry:symbol-registry-04764}\(\operatorname{Label}_d\) & \texttt{\detokenize{\operatorname{Label}_d}} & finite disjunction of all fixed program instruction labels & local:sec:p05-representability & relation & \hyperref[def:p05-rep-register-certificate]{\ref*{def:p05-rep-register-certificate}} & \thepage & Includes halts and excludes unrecognized labels \\
\phantomsection\label{entry:symbol-registry-04765}\(\operatorname{Le}\) & \texttt{\detokenize{\operatorname{Le}}} & marked-bounded witness formula implying source comparison & local:sec:p05-representability & relation & \hyperref[def:p05-rep-register-certificate]{\ref*{def:p05-rep-register-certificate}} & \thepage & Standalone source comparison is not counted as atomic bounded formula \\
\phantomsection\label{entry:symbol-registry-04766}\(C_d\) & \texttt{\detokenize{C_d}} & bounded standard trace-certificate formula & local:sec:p05-representability & relation & \hyperref[def:p05-rep-register-certificate]{\ref*{def:p05-rep-register-certificate}} & \thepage & Separate from governing positive-cons C; d denotes fixed description \\
\phantomsection\label{entry:symbol-registry-04767}\(\delta_d\) & \texttt{\detokenize{\delta_d}} & bounded minimum-certificate matrix & local:sec:p05-representability & relation & \hyperref[def:p05-rep-least-graph]{\ref*{def:p05-rep-least-graph}} & \thepage & No Q normalization from a different Sigma1 representative assumed \\
\phantomsection\label{entry:symbol-registry-04768}\(\gamma_d\) & \texttt{\detokenize{\gamma_d}} & literal existential-prefix pointwise strong graph & local:sec:p05-representability & relation & \hyperref[def:p05-rep-least-graph]{\ref*{def:p05-rep-least-graph}} & \thepage & Universal output at each standard numeral input; not universal input-totality \\
\phantomsection\label{entry:symbol-registry-04769}\(\rho_R\) & \texttt{\detokenize{\rho_R}} & selected numeralwise relation representative & local:sec:p05-representability & relation & \hyperref[thm:p05-rep-numeralwise-relations]{\ref*{thm:p05-rep-numeralwise-relations}} & \thepage & Characteristic value one signifies truth \\
\phantomsection\label{entry:symbol-registry-04770}\(\kappa_T\) & \texttt{\detokenize{\kappa_T}} & bounded matrix of selected certified-proof representative & local:sec:p05-representability & relation & \hyperref[def:p05-rep-certified-presentation]{\ref*{def:p05-rep-certified-presentation}} & \thepage & Compiled from actual N16 checker with fixed e\_T \\
\phantomsection\label{entry:symbol-registry-04771}\(\mathsf{Proof}^{*}_T\) & \texttt{\detokenize{\mathsf{Proof}^{*}_T}} & selected Sigma1 certified-proof representative & local:sec:p05-representability & relation & \hyperref[def:p05-rep-certified-presentation]{\ref*{def:p05-rep-certified-presentation}} & \thepage & Selection not Q equivalence with all other numeralwise representatives \\
\phantomsection\label{entry:symbol-registry-04772}\(\mathsf{Prov}^{*}_T\) & \texttt{\detokenize{\mathsf{Prov}^{*}_T}} & provability formula for selected representative & local:sec:p05-representability & relation & \hyperref[def:p05-rep-certified-presentation]{\ref*{def:p05-rep-certified-presentation}} & \thepage & Literal existential proof and computation-certificate prefix \\
\phantomsection\label{entry:symbol-registry-04773}\(\tau\) & \texttt{\detokenize{\tau}} & fixed valid default code of arithmetic tautology & local:sec:p05-representability & object & \hyperref[def:p05-rep-diagonal-default]{\ref*{def:p05-rep-diagonal-default}} & \thepage & Code of 0=0; negation default is N7(tau) \\
\phantomsection\label{entry:symbol-registry-04774}\(D\) & \texttt{\detokenize{D}} & selected literal Sigma1 strong diagonal graph & local:sec:p05-representability & relation & \hyperref[thm:p05-rep-diagonal-graph]{\ref*{thm:p05-rep-diagonal-graph}} & \thepage & Distinguished variable x\_0 and tau default align with incompleteness fragment \\
\phantomsection\label{entry:symbol-registry-04775}\(J\) & \texttt{\detokenize{J}} & proof-local fully expanded successful-output existential & local:annex:p05-rep-strong-graph & relation & \hyperref[annex:p05-rep-strong-graph]{\ref*{annex:p05-rep-strong-graph}} & \thepage & Logical expansion for guard derivation; not itself asserted bounded in raw Lang\_A \\
\phantomsection\label{entry:symbol-registry-04776}\(d(n)\) & \texttt{\detokenize{d}} & governing numeral self-substitution code function & \hyperref[sec:p05-inc]{\ref*{sec:p05-inc}} & n/a & \hyperref[def:p05-inc-interface]{\ref*{def:p05-inc-interface}} & \thepage & Numerical primitive recursive function; malformed input gets a fixed sentence code; not an arithmetic function symbol. \\
\phantomsection\label{entry:symbol-registry-04777}\(\operatorname{neg}(n)\) & \texttt{\detokenize{\operatorname{neg}}} & governing formula negation-code constructor & \hyperref[sec:p05-inc]{\ref*{sec:p05-inc}} & n/a & \hyperref[def:p05-inc-interface]{\ref*{def:p05-inc-interface}} & \thepage & Numerical N7 constructor with malformed default; arithmetic uses a pointwise strong graph N not neg(a) as a term. \\
\phantomsection\label{entry:symbol-registry-04778}\(D(x,y),\delta_d(x,y,\vec s)\) & \texttt{\detokenize{D,}\allowbreak\detokenize{\delta_d}} & selected literal Sigma1 diagonal graph and bounded certificate matrix & \hyperref[sec:p05-inc]{\ref*{sec:p05-inc}} & n/a & \hyperref[def:p05-inc-interface]{\ref*{def:p05-inc-interface}} & \thepage & Separate numeralwise all-output uniqueness/existence interface; not Q-uniform totality. \\
\phantomsection\label{entry:symbol-registry-04779}\(N(x,y)\) & \texttt{\detokenize{N}} & selected arithmetic graph of the numerical negation-code function & \hyperref[sec:p05-inc]{\ref*{sec:p05-inc}} & n/a & \hyperref[def:p05-inc-interface]{\ref*{def:p05-inc-interface}} & \thepage & Pointwise strong graph at fixed standard inputs; distinct from the standard arithmetic structure mathcal N. \\
\phantomsection\label{entry:symbol-registry-04780}\(P_T(p,a),\operatorname{Pr}_T(a)\) & \texttt{\detokenize{P_T,}\allowbreak\detokenize{\operatorname{Pr}_T}} & chosen numeralwise Sigma1 certified-proof representative and existential provability & \hyperref[sec:p05-inc]{\ref*{sec:p05-inc}} & n/a & \hyperref[def:p05-inc-interface]{\ref*{def:p05-inc-interface}} & \thepage & Identifies with active Proof\_T and Prov\_T only by explicit formula selection; no equivalence of arbitrary representatives in Q. \\
\phantomsection\label{entry:symbol-registry-04781}\(r^-\) & \texttt{\detokenize{r^-}} & natural-number code of the negation of the constructed Rosser sentence & \hyperref[sec:p05-inc]{\ref*{sec:p05-inc}} & n/a & \hyperref[thm:p05-inc-rosser-chain]{\ref*{thm:p05-inc-rosser-chain}} & \thepage & r\^{}- is neg(r) in the metatheory; arithmetic parameter is its numeral. \\
\phantomsection\label{entry:symbol-registry-04782}\(C_D\) & \texttt{\detokenize{C_D}} & finite closed-sentence conjunction used in the actual ND consistency test & \hyperref[sec:p05-inc]{\ref*{sec:p05-inc}} & n/a & \hyperref[thm:p05-inc-finite-extension]{\ref*{thm:p05-inc-finite-extension}} & \thepage & Empty list gives 0=0; duplicate sentence entries allowed; no implicit Hilbert deduction theorem. \\
\end{longtable}
\endgroup
\Needspace{.4\textheight}
\subsection*{Stable part P6 scopes}
\begingroup\scriptsize\setlength{\tabcolsep}{3pt}
\begin{longtable}{@{}p{0.9in}p{1.3in}p{1.0in}p{0.5in}p{0.3in}p{0.4in}p{0.35in}p{1.65in}@{}}
\toprule
\textbf{Glyph} & \textbf{Macro} & \textbf{Role} & \textbf{Scope} & \textbf{Prec./form} & \textbf{First} & \textbf{Reg. page} & \textbf{Principal uses / notes} \\
\midrule
\endfirsthead
\toprule
\textbf{Glyph} & \textbf{Macro} & \textbf{Role} & \textbf{Scope} & \textbf{Prec./form} & \textbf{First} & \textbf{Reg. page} & \textbf{Principal uses / notes} \\
\midrule
\endhead
\midrule\multicolumn{8}{r}{\emph{continued}}\\\endfoot
\bottomrule\endlastfoot
\phantomsection\label{entry:symbol-registry-00019}\(\alpha,\beta\) & \texttt{\detokenize{\alpha,}\allowbreak\detokenize{\beta}} & ordinal metavariables & P6 zf-zfc to determinacy & n/a & \hyperref[sec:p06-source-map]{\ref*{sec:p06-source-map}} & \thepage & Ordinals of the ambient set theory; recurs under a local declaration in P7.5 ordinal analysis; disjoint from the type-metavariable scope. \\
\phantomsection\label{entry:symbol-registry-00065}\(V_\alpha\) & \texttt{\detokenize{V_\alpha}} & cumulative-hierarchy level & P6 zf-zfc to determinacy & n/a & \hyperref[sec:p06-source-map]{\ref*{sec:p06-source-map}} & \thepage & \(V=\bigcup_\alpha V_\alpha\) is class notation, registered with its level notation. \\
\phantomsection\label{entry:symbol-registry-00066}\(\operatorname{rk}(x)\) & \texttt{\detokenize{\rank}} & set-theoretic rank & P6 zf-zfc to determinacy & n/a & \hyperref[sec:p06-source-map]{\ref*{sec:p06-source-map}} & \thepage & Least \(\alpha\) with \(x\subseteq V_\alpha\). \\
\phantomsection\label{entry:symbol-registry-00067}\(L_\alpha\) & \texttt{\detokenize{L_\alpha}} & constructible level & P6 cardinals-choice & n/a & \hyperref[sec:p06-cardinals-choice]{\ref*{sec:p06-cardinals-choice}} & \thepage & The \(\alpha\)-th level of the constructible hierarchy. \\
\phantomsection\label{entry:symbol-registry-00068}\(\mathbb{P}\) & \texttt{\detokenize{\mathbb{P}}} & forcing poset & P6 cardinals-choice to constructibility-forcing & n/a & \hyperref[sec:p06-cardinals-choice]{\ref*{sec:p06-cardinals-choice}} & \thepage & \(q\leq p\) reads \(q\) is stronger, per the fixed forcing convention. \\
\phantomsection\label{entry:symbol-registry-00069}\(V^{\mathbb P}\) & \texttt{\detokenize{V^{\mathbb P}}} & class of P-names & P6 cardinals-choice & n/a & \hyperref[sec:p06-cardinals-choice]{\ref*{sec:p06-cardinals-choice}} & \thepage & Name metavariables \(\nu\) range over \(M\cap V^{\mathbb P}\). \\
\phantomsection\label{entry:symbol-registry-00070}\(\check x\) & \texttt{\detokenize{\check x}} & check name & P6 cardinals-choice & n/a & \hyperref[sec:p06-cardinals-choice]{\ref*{sec:p06-cardinals-choice}} & \thepage & Canonical forcing name of a ground-model set. \\
\phantomsection\label{entry:symbol-registry-00071}\(\nu^G\) & \texttt{\detokenize{\nu^G}} & name evaluation & P6 cardinals-choice & n/a & \hyperref[sec:p06-cardinals-choice]{\ref*{sec:p06-cardinals-choice}} & \thepage & Interpretation by the generic filter \(G\) of a locally declared forcing name \(\nu,\nu_f\), or \(\zeta\). \\
\phantomsection\label{entry:symbol-registry-00072}\(M[G]\) & \texttt{\detokenize{M[G]}} & generic extension & P6 cardinals-choice to infinitary-combinatorics & n/a & \hyperref[sec:p06-cardinals-choice]{\ref*{sec:p06-cardinals-choice}} & \thepage & The forcing extension obtained by adjoining the \(M\)-generic filter \(G\). \\
\phantomsection\label{entry:symbol-registry-00073}\(p\Vdash^{\mathrm{sem}}_{M,\mathbb P}\varphi\) & \texttt{\detokenize{\forces^{\mathrm{sem}}_{M,\mathbb P}}} & semantic set-theoretic forcing & P6 cardinals-choice & meta & \hyperref[sec:p06-cardinals-choice]{\ref*{sec:p06-cardinals-choice}} & \thepage & Exact external relation: every M-generic filter containing p yields an extension satisfying the formula at the interpreted names. \\
\phantomsection\label{entry:symbol-registry-00074}\(p\Vdash_{M,\mathbb P}\varphi\) & \texttt{\detokenize{\forces_{M,\mathbb P}}} & imported recursive forcing assertion & P6 cardinals-choice & meta & \hyperref[sec:p06-cardinals-choice]{\ref*{sec:p06-cardinals-choice}} & \thepage & Formula-by-formula recursive assertion imported by the cited forcing theorem; equivalent there to the separately registered semantic relation. Bare \(p\Vdash\varphi\) is its local abbreviation only under the renewed ccc/name-counting convention. \\
\phantomsection\label{entry:symbol-registry-00326}\(\nu,\nu_f,\zeta\) & \texttt{\detokenize{\nu,}\allowbreak\detokenize{\nu_f,}\allowbreak\detokenize{\zeta}} & forcing-name metavariables & P6 cardinals-choice & n/a & \hyperref[sec:p06-cardinals-choice]{\ref*{sec:p06-cardinals-choice}} & \thepage & Members of the active P-name hierarchy; subscripts identify a name by its intended value and never turn a name into that value. \\
\phantomsection\label{entry:symbol-registry-00327}\(\mathbb B\) & \texttt{\detokenize{\mathbb B}} & complete Boolean algebra used for Boolean-valued models & P6 cardinals-choice & n/a & \hyperref[sec:p06-cardinals-choice]{\ref*{sec:p06-cardinals-choice}} & \thepage & Scope is disjoint from the P2.1 Boolean truth-value set; completeness supplies all joins and meets used by the Boolean-valued recursion. \\
\phantomsection\label{entry:symbol-registry-00328}\(V^{\mathbb B}_\alpha,V^{\mathbb B}\) & \texttt{\detokenize{V^{\mathbb B}_\alpha,}\allowbreak\detokenize{V^{\mathbb B}}} & ambient Boolean-valued name hierarchy & P6 cardinals-choice & n/a & \hyperref[sec:p06-cardinals-choice]{\ref*{sec:p06-cardinals-choice}} & \thepage & Names are Boolean-valued functions on earlier levels for an ambient-complete algebra; distinct from the P-name hierarchy and the ground-model-computed variant. \\
\phantomsection\label{entry:symbol-registry-00329}\(\operatorname{Bool}_{\mathbb B}(\varphi)\) & \texttt{\detokenize{\operatorname{Bool}_{\mathbb B}}} & ambient Boolean truth value of a formula & P6 cardinals-choice & n/a & \hyperref[sec:p06-cardinals-choice]{\ref*{sec:p06-cardinals-choice}} & \thepage & Value in an ambient-complete Boolean algebra; this is not the two-valued satisfaction relation. \\
\phantomsection\label{entry:symbol-registry-00330}\((V^{\mathbb B})^M,\operatorname{Bool}_{\mathbb B}^M,S_{\varphi,\bar{\mathbf v}}^M\) & \texttt{\detokenize{(V^{\mathbb B})^M,}\allowbreak\detokenize{\operatorname{Bool}_{\mathbb B}^M,}\allowbreak\detokenize{S_{\varphi,\bar{\mathbf v}}^M}} & ground-model-computed Boolean names truth and attained-value set & P6 cardinals-choice & n/a & \hyperref[sec:p06-cardinals-choice]{\ref*{sec:p06-cardinals-choice}} & \thepage & All recursions, Separation, and joins are computed in M when M regards B as complete; ambient completeness is not assumed. \\
\phantomsection\label{entry:symbol-registry-00331}\(b^{\complement_{\mathbb B}}\) & \texttt{\detokenize{{}^{\complement_{\mathbb B}}}} & Boolean-algebra complement & P6 cardinals-choice & 1 & \hyperref[sec:p06-cardinals-choice]{\ref*{sec:p06-cardinals-choice}} & \thepage & Postfix operation in the declared algebra; the subscript prevents collision with object-language negation. \\
\phantomsection\label{entry:symbol-registry-00332}\(b\wedge_{\mathbb B}c\) & \texttt{\detokenize{\wedge_{\mathbb B}}} & Boolean-algebra meet & P6 cardinals-choice & 4 & \hyperref[sec:p06-cardinals-choice]{\ref*{sec:p06-cardinals-choice}} & \thepage & Operation in the declared algebra; distinct from the globally registered object-language conjunction. \\
\phantomsection\label{entry:symbol-registry-00333}\(b\vee_{\mathbb B}c\) & \texttt{\detokenize{\vee_{\mathbb B}}} & Boolean-algebra join & P6 cardinals-choice & 5 & \hyperref[sec:p06-cardinals-choice]{\ref*{sec:p06-cardinals-choice}} & \thepage & Operation in the declared algebra; distinct from the globally registered object-language disjunction. \\
\phantomsection\label{entry:symbol-registry-00334}\(\mathbf u,\mathbf v,\mathbf w\) & \texttt{\detokenize{\mathbf u,}\allowbreak\detokenize{\mathbf v,}\allowbreak\detokenize{\mathbf w}} & Boolean-valued name metavariables & P6 cardinals-choice & n/a & \hyperref[sec:p06-cardinals-choice]{\ref*{sec:p06-cardinals-choice}} & \thepage & Members of \(V^{\mathbb B}\); bold forms avoid the globally reserved valuation and term metavariables. \\
\phantomsection\label{entry:symbol-registry-00335}\(\mathbf u\sim_G\mathbf v\) & \texttt{\detokenize{\sim_G}} & generic-ultrafilter equivalence of Boolean names & P6 cardinals-choice & meta & \hyperref[sec:p06-cardinals-choice]{\ref*{sec:p06-cardinals-choice}} & \thepage & Defined by Boolean equality having value in \(G\); the quotient relation is not ordinary set equality. \\
\phantomsection\label{entry:symbol-registry-00336}\(\mathcal G,\mathfrak F\) & \texttt{\detokenize{\mathcal G,}\allowbreak\detokenize{\mathfrak F}} & permutation group and normal subgroup filter & P6 cardinals-choice & n/a & \hyperref[sec:p06-cardinals-choice]{\ref*{sec:p06-cardinals-choice}} & \thepage & Typed components of a permutation-model datum; \(\mathfrak F\) contains subgroups of \(\mathcal G\). \\
\phantomsection\label{entry:symbol-registry-00337}\(\operatorname{Sym}(A),\operatorname{Stab}_{\mathcal G}(x)\) & \texttt{\detokenize{\operatorname{Sym},}\allowbreak\detokenize{\operatorname{Stab}_{\mathcal G}}} & full permutation group and stabilizer & P6 cardinals-choice & n/a & \hyperref[sec:p06-cardinals-choice]{\ref*{sec:p06-cardinals-choice}} & \thepage & The stabilizer is a subgroup of the declared permutation group; support and symmetry are defined from it. \\
\phantomsection\label{entry:symbol-registry-00338}\(\omega^\omega\) & \texttt{\detokenize{\omega^\omega}} & ordinal exponentiation & P6 zf-zfc & n/a & \hyperref[sec:p06-source-map]{\ref*{sec:p06-source-map}} & \thepage & In the ordinal-arithmetic chapter this is the ordinal power obtained by right-recursion. \\
\phantomsection\label{entry:symbol-registry-00339}\(\omega^\omega\) & \texttt{\detokenize{\omega^\omega}} & Baire function space & P6 infinitary-combinatorics to large-cardinals-absoluteness & n/a & \hyperref[sec:p06-infinitary-combinatorics]{\ref*{sec:p06-infinitary-combinatorics}} & \thepage & The set of functions from \(\omega\) to \(\omega\); Chapter 6.7 adds the product topology. Scope is disjoint from the ordinal-power role. \\
\phantomsection\label{entry:symbol-registry-00340}\(\boldsymbol{\Sigma}^{0}_{\alpha},\boldsymbol{\Pi}^{0}_{\alpha},\boldsymbol{\Delta}^{0}_{\alpha}\) & \texttt{\detokenize{\boldsymbol{\Sigma}^{0}_{\alpha},}\allowbreak\detokenize{\boldsymbol{\Pi}^{0}_{\alpha},}\allowbreak\detokenize{\boldsymbol{\Delta}^{0}_{\alpha}}} & boldface Borel pointclasses & P6 cardinal-arithmetic & n/a & \hyperref[sec:p06-cardinal-arithmetic]{\ref*{sec:p06-cardinal-arithmetic}} & \thepage & Countable-ordinal Borel levels on a Polish space; distinct from the P4-P5 bold arithmetical-definability classes. \\
\phantomsection\label{entry:symbol-registry-00341}\(\boldsymbol{\Sigma}^{1}_{n},\boldsymbol{\Pi}^{1}_{n},\boldsymbol{\Delta}^{1}_{n}\) & \texttt{\detokenize{\boldsymbol{\Sigma}^{1}_{n},}\allowbreak\detokenize{\boldsymbol{\Pi}^{1}_{n},}\allowbreak\detokenize{\boldsymbol{\Delta}^{1}_{n}}} & boldface projective pointclasses & P6 infinitary-combinatorics to large-cardinals-absoluteness & n/a & \hyperref[sec:p06-infinitary-combinatorics]{\ref*{sec:p06-infinitary-combinatorics}} & \thepage & Real-parameter projective classes; distinct from lightface analytical classes and from the Levy hierarchy. \\
\phantomsection\label{entry:symbol-registry-00342}\(\operatorname{Bor}(X)\) & \texttt{\detokenize{\operatorname{Bor}}} & Borel sigma-algebra of a Polish space & P6 cardinal-arithmetic & n/a & \hyperref[sec:p06-cardinal-arithmetic]{\ref*{sec:p06-cardinal-arithmetic}} & \thepage & The least sigma-algebra containing the open subsets of \(X\). \\
\phantomsection\label{entry:symbol-registry-00343}\(\operatorname{crit}(j)\) & \texttt{\detokenize{\operatorname{crit}}} & critical point of an elementary embedding & P6 infinitary-combinatorics to large-cardinals-absoluteness & n/a & \hyperref[sec:p06-infinitary-combinatorics]{\ref*{sec:p06-infinitary-combinatorics}} & \thepage & Least ordinal moved by \(j\); subscripted embedding symbols such as \(j_U\) use the same operation. \\
\phantomsection\label{entry:symbol-registry-00344}\(G(A)\) & \texttt{\detokenize{G(A)}} & Gale--Stewart payoff game & P6 large-cardinals-absoluteness & n/a & \hyperref[sec:p06-large-cardinals-absoluteness]{\ref*{sec:p06-large-cardinals-absoluteness}} & \thepage & Length-omega perfect-information game whose payoff set is \(A\subseteq\omega^\omega\). \\
\phantomsection\label{entry:symbol-registry-00345}\(\mathsf{AD},\mathsf{PD}\) & \texttt{\detokenize{\mathsf{AD},}\allowbreak\detokenize{\mathsf{PD}}} & determinacy assertions & P6 large-cardinals-absoluteness & n/a & \hyperref[sec:p06-large-cardinals-absoluteness]{\ref*{sec:p06-large-cardinals-absoluteness}} & \thepage & AD quantifies over all payoff sets of reals; PD restricts to boldface projective payoff sets. \\
\phantomsection\label{entry:symbol-registry-00346}\(\mathsf{ZF},\mathsf{ZFC}\) & \texttt{\detokenize{\mathsf{ZF},}\allowbreak\detokenize{\mathsf{ZFC}}} & first-order set theories without and with Choice & P6 zf-zfc to determinacy & n/a & \hyperref[sec:p06-source-map]{\ref*{sec:p06-source-map}} & \thepage & ZF is the displayed axiom basis without Choice; ZFC adds the displayed object-language Axiom of Choice. \\
\phantomsection\label{entry:symbol-registry-00347}\(\mathsf Z,\mathsf{ZC},\mathsf{KP},\mathsf{KP}^{\mathcal P}\) & \texttt{\detokenize{\mathsf Z,}\allowbreak\detokenize{\mathsf{ZC},}\allowbreak\detokenize{\mathsf{KP},}\allowbreak\detokenize{\mathsf{KP}^{\mathcal P}}} & set-theory system abbreviations & P6 cumulative-hierarchy & n/a & \hyperref[sec:cumulative-hierarchy-01]{\ref*{sec:cumulative-hierarchy-01}} & \thepage & Z omits Replacement from the stated ZF basis; ZC adds Choice; KP is the stated admissible-set basis and its power-set variant is explicitly superscripted. \\
\phantomsection\label{entry:symbol-registry-00348}\(\mathsf{GB},\mathsf{NBG},\mathsf{MK},\mathsf{KM}\) & \texttt{\detokenize{\mathsf{GB},}\allowbreak\detokenize{\mathsf{NBG},}\allowbreak\detokenize{\mathsf{MK},}\allowbreak\detokenize{\mathsf{KM}}} & set-and-class theory abbreviations & P6 cumulative-hierarchy & n/a & \hyperref[sec:cumulative-hierarchy-01]{\ref*{sec:cumulative-hierarchy-01}} & \thepage & The comparison table fixes GB as choice-free NBG and treats MK and KM as names for the declared impredicative class theory convention. \\
\phantomsection\label{entry:symbol-registry-00349}\(\mathsf{NF},\mathsf{TTT}\) & \texttt{\detokenize{\mathsf{NF},}\allowbreak\detokenize{\mathsf{TTT}}} & New Foundations and Tangled Type Theory abbreviations & P6 cumulative-hierarchy & n/a & \hyperref[sec:cumulative-hierarchy-01]{\ref*{sec:cumulative-hierarchy-01}} & \thepage & Distinct theories linked only by the cited consistency reduction; neither is presented as a subtheory of ZF. \\
\phantomsection\label{entry:symbol-registry-00350}\(\mathsf{ZFA}\) & \texttt{\detokenize{\mathsf{ZFA}}} & set theory with atoms & P6 cardinals-choice & n/a & \hyperref[sec:p06-cardinals-choice]{\ref*{sec:p06-cardinals-choice}} & \thepage & Theory of the local permutation-model example; the cited transfer is required before drawing a pure-ZF conclusion. \\
\phantomsection\label{entry:symbol-registry-00351}\(\mathsf{AC}\) & \texttt{\detokenize{\mathsf{AC}}} & Axiom of Choice & P6 zf-zfc to determinacy & n/a & \hyperref[sec:p06-source-map]{\ref*{sec:p06-source-map}} & \thepage & Object-language set-theory axiom; occurrences in metatheoretic dependency tables name that exact principle. \\
\phantomsection\label{entry:symbol-registry-00352}\(\mathsf{DC},\mathsf{AC}_{\omega},\mathsf{AC}_{\omega}^{\mathrm{fin}},\mathsf{BPI}\) & \texttt{\detokenize{\mathsf{DC},}\allowbreak\detokenize{\mathsf{AC}_{\omega},}\allowbreak\detokenize{\mathsf{AC}_{\omega}^{\mathrm{fin}},}\allowbreak\detokenize{\mathsf{BPI}}} & restricted or weaker choice principles & P6 cumulative-hierarchy to constructibility-forcing & n/a & \hyperref[sec:cumulative-hierarchy-01]{\ref*{sec:cumulative-hierarchy-01}} & \thepage & Dependent Choice, Countable Choice, Countable Choice for finite sets, and the Boolean Prime Ideal Theorem; implication arrows are stated separately. \\
\phantomsection\label{entry:symbol-registry-00353}\(\kappa,\lambda,\mu\) & \texttt{\detokenize{\kappa,}\allowbreak\detokenize{\lambda,}\allowbreak\detokenize{\mu}} & ordinal or cardinal metavariables as explicitly quantified & P6 zf-zfc to determinacy & n/a & \hyperref[sec:p06-source-map]{\ref*{sec:p06-source-map}} & \thepage & The local quantifier fixes the standing: limit ordinals and conditional hierarchy markers occur before the cardinal-arithmetic chapters. \\
\phantomsection\label{entry:symbol-registry-00354}\(\gamma,\delta,\xi\) & \texttt{\detokenize{\gamma,}\allowbreak\detokenize{\delta,}\allowbreak\detokenize{\xi}} & additional ordinal metavariables & P6 zf-zfc to determinacy & n/a & \hyperref[sec:p06-source-map]{\ref*{sec:p06-source-map}} & \thepage & Ordinals in recursive, limit, club, forcing, and hierarchy clauses; complements the existing alpha-beta ordinal row. \\
\phantomsection\label{entry:symbol-registry-00355}\(\kappa^+\) & \texttt{\detokenize{{}^+}} & successor-cardinal operation & P6 cardinals-choice to large-cardinals-absoluteness & 1 & \hyperref[sec:p06-cardinals-choice]{\ref*{sec:p06-cardinals-choice}} & \thepage & Least cardinal strictly above kappa; distinct from ordinal successor alpha+1. \\
\phantomsection\label{entry:symbol-registry-00356}\(\kappa^\lambda,\lambda^{<\kappa}\) & \texttt{\detokenize{\kappa^\lambda,}\allowbreak\detokenize{\lambda^{<\kappa}}} & cardinal exponentiation and bounded-exponent supremum & P6 cardinals-choice to infinitary-combinatorics & n/a & \hyperref[sec:p06-cardinals-choice]{\ref*{sec:p06-cardinals-choice}} & \thepage & Cardinal function-space size and the supremum over exponents below kappa; the full entry is Chapter 6.5. \\
\phantomsection\label{entry:symbol-registry-00357}\(\aleph_\alpha,\beth_\alpha\) & \texttt{\detokenize{\aleph_\alpha,}\allowbreak\detokenize{\beth_\alpha}} & aleph and beth cardinal hierarchies & P6 cardinals-choice to large-cardinals-absoluteness & n/a & \hyperref[sec:p06-cardinals-choice]{\ref*{sec:p06-cardinals-choice}} & \thepage & Aleph enumerates infinite initial ordinals and beth iterates powerset cardinality; limit stages take suprema. \\
\phantomsection\label{entry:symbol-registry-00358}\(\operatorname{cf}(\alpha),\operatorname{cf}(P,\leq)\) & \texttt{\detokenize{\operatorname{cf}}} & cofinality of an ordinal or ordered set & P6 cardinals-choice to large-cardinals-absoluteness & n/a & \hyperref[sec:p06-cardinals-choice]{\ref*{sec:p06-cardinals-choice}} & \thepage & For an ordinal: least cofinal order type. For a nonempty ordered set: least cardinality of a cofinal family. Regularity is cf(kappa)=kappa. \\
\phantomsection\label{entry:symbol-registry-00359}\(\mathsf{CH},\mathsf{GCH}\) & \texttt{\detokenize{\mathsf{CH},}\allowbreak\detokenize{\mathsf{GCH}}} & continuum hypotheses & P6 cardinals-choice to infinitary-combinatorics & n/a & \hyperref[sec:p06-cardinals-choice]{\ref*{sec:p06-cardinals-choice}} & \thepage & CH is the aleph-one equation at the continuum; GCH is its all-aleph scheme. \\
\phantomsection\label{entry:symbol-registry-00360}\(\mathfrak c\) & \texttt{\detokenize{\mathfrak c}} & cardinality of the continuum & P6 constructibility-forcing & n/a & \hyperref[sec:constructibility-forcing-01]{\ref*{sec:constructibility-forcing-01}} & \thepage & Defined as two to the aleph-zero; used only to bound the Martin-axiom scheme. \\
\phantomsection\label{entry:symbol-registry-00361}\(\mathsf{Con}(\mathsf{ZFC})^\in\) & \texttt{\detokenize{\mathsf{Con}(\mathsf{ZFC})^\in}} & membership-language translation of the arithmetical ZFC-consistency sentence & P6 cardinals-choice & n/a & \hyperref[sec:p06-cardinals-choice]{\ref*{sec:p06-cardinals-choice}} & \thepage & Obtained by translating the fixed arithmetical no-proof sentence along the canonical interpretation of arithmetic by the finite ordinals; the separate countable-transitive-model hypothesis is stronger and must not be conflated with this sentence. \\
\phantomsection\label{entry:symbol-registry-00362}\(2^{<\omega},2^\omega\) & \texttt{\detokenize{2^{<\omega},}\allowbreak\detokenize{2^\omega}} & finite and infinite binary-sequence spaces & P6 cumulative-hierarchy to constructibility-forcing & n/a & \hyperref[sec:cumulative-hierarchy-01]{\ref*{sec:cumulative-hierarchy-01}} & \thepage & Finite strings and omega-sequences over two; the first carries prefix order and the second supplies branches or reals as locally declared. \\
\phantomsection\label{entry:symbol-registry-00363}\(\mathcal T\) & \texttt{\detokenize{\mathcal T}} & tree metavariable & P6 cumulative-hierarchy to constructibility-forcing & n/a & \hyperref[sec:cumulative-hierarchy-01]{\ref*{sec:cumulative-hierarchy-01}} & \thepage & A binary or set-theoretic tree as locally declared; scope is disjoint from the P4.1-P4.2 Turing-machine role. \\
\phantomsection\label{entry:symbol-registry-00364}\(\mathbf s\) & \texttt{\detokenize{\mathbf s}} & finite binary-sequence metavariable & P6 cumulative-hierarchy & n/a & \hyperref[sec:cumulative-hierarchy-01]{\ref*{sec:cumulative-hierarchy-01}} & \thepage & Node of the canonical binary tree; bold form avoids the globally reserved assignment metavariable s. \\
\phantomsection\label{entry:symbol-registry-00365}\(\operatorname{Fn}(I,J,{<}\delta)\) & \texttt{\detokenize{\operatorname{Fn}}} & bounded-size partial-function poset & P6 cardinals-choice & n/a & \hyperref[sec:p06-cardinals-choice]{\ref*{sec:p06-cardinals-choice}} & \thepage & Partial functions from I to J of size below delta, ordered by reverse inclusion. \\
\phantomsection\label{entry:symbol-registry-00366}\(\operatorname{Add}(\omega,\kappa)\) & \texttt{\detokenize{\operatorname{Add}}} & Cohen forcing adding kappa many reals & P6 cardinals-choice to constructibility-forcing & n/a & \hyperref[sec:p06-cardinals-choice]{\ref*{sec:p06-cardinals-choice}} & \thepage & The explicitly displayed Fn poset; preservation and continuum claims retain their separate hypotheses. \\
\phantomsection\label{entry:symbol-registry-00367}\(\operatorname{Coll}(\kappa,\lambda)\) & \texttt{\detokenize{\operatorname{Coll}}} & Levy collapse forcing & P6 cardinals-choice & n/a & \hyperref[sec:p06-cardinals-choice]{\ref*{sec:p06-cardinals-choice}} & \thepage & Partial functions from kappa to lambda of size below kappa under the catalog convention; collapse claims are cited. \\
\phantomsection\label{entry:symbol-registry-00368}\(1_{\mathbb P}\) & \texttt{\detokenize{1_{\mathbb P}}} & weakest forcing condition & P6 cardinals-choice & n/a & \hyperref[sec:p06-cardinals-choice]{\ref*{sec:p06-cardinals-choice}} & \thepage & Top element under the convention that q less-than-or-equal p means q is stronger. \\
\phantomsection\label{entry:symbol-registry-00369}\(\mathop{\triangle}_{\xi<\kappa}C_\xi\) & \texttt{\detokenize{\mathop{\triangle}}} & diagonal intersection & P6 constructibility-forcing & n/a & \hyperref[sec:constructibility-forcing-01]{\ref*{sec:constructibility-forcing-01}} & \thepage & Set of \(\alpha<\kappa\) belonging to every earlier \(C_\xi\); distinct from symmetric difference. \\
\phantomsection\label{entry:symbol-registry-00370}\([X]^\ell\) & \texttt{\detokenize{[X]^\ell}} & ell-element subsets of a set & P6 constructibility-forcing to infinitary-combinatorics & n/a & \hyperref[sec:constructibility-forcing-01]{\ref*{sec:constructibility-forcing-01}} & \thepage & Domain of finite-color partition relations; ell is a positive finite integer. \\
\phantomsection\label{entry:symbol-registry-00371}\(\kappa\longrightarrow(\lambda)^\ell_\mu\) & \texttt{\detokenize{\longrightarrow}} & partition relation & P6 constructibility-forcing to infinitary-combinatorics & meta & \hyperref[sec:constructibility-forcing-01]{\ref*{sec:constructibility-forcing-01}} & \thepage & Every mu-coloring of the ell-element subsets of kappa has a homogeneous subset of cardinality lambda; not metalinguistic implication. \\
\phantomsection\label{entry:symbol-registry-00372}\(\beth_n(\kappa)\) & \texttt{\detokenize{\beth_n(\kappa)}} & finite beth iteration above a cardinal & P6 constructibility-forcing & n/a & \hyperref[sec:constructibility-forcing-01]{\ref*{sec:constructibility-forcing-01}} & \thepage & Starts at kappa and iterates powerset cardinality n times; distinct from the global beth hierarchy indexed from aleph-zero. \\
\phantomsection\label{entry:symbol-registry-00373}\(\mathcal T_\alpha,\operatorname{ht}_{\mathcal T}(\mathbf t)\) & \texttt{\detokenize{\mathcal T_\alpha,}\allowbreak\detokenize{\operatorname{ht}_{\mathcal T}}} & tree level and node height & P6 constructibility-forcing & n/a & \hyperref[sec:constructibility-forcing-01]{\ref*{sec:constructibility-forcing-01}} & \thepage & Height is the predecessor order type and the alpha-level contains exactly nodes of that height. \\
\phantomsection\label{entry:symbol-registry-00374}\(\mathbf t\) & \texttt{\detokenize{\mathbf t}} & set-theoretic tree-node metavariable & P6 constructibility-forcing & n/a & \hyperref[sec:constructibility-forcing-01]{\ref*{sec:constructibility-forcing-01}} & \thepage & Node of the locally declared tree; bold form avoids the globally reserved term metavariable t. \\
\phantomsection\label{entry:symbol-registry-00375}\(\mathsf{SH}\) & \texttt{\detokenize{\mathsf{SH}}} & Suslin hypothesis & P6 constructibility-forcing & n/a & \hyperref[sec:constructibility-forcing-01]{\ref*{sec:constructibility-forcing-01}} & \thepage & Assertion that no Suslin tree exists under the chapter's tree convention. \\
\phantomsection\label{entry:symbol-registry-00376}\(\diamondsuit_{\omega_1}\) & \texttt{\detokenize{\diamondsuit_{\omega_1}}} & diamond principle at omega-one & P6 constructibility-forcing & n/a & \hyperref[sec:constructibility-forcing-01]{\ref*{sec:constructibility-forcing-01}} & \thepage & Stationary-guessing sequence defined in the combinatorial-principles entry. \\
\phantomsection\label{entry:symbol-registry-00377}\(\square_\kappa\) & \texttt{\detokenize{\square_\kappa}} & square principle & P6 constructibility-forcing & n/a & \hyperref[sec:constructibility-forcing-01]{\ref*{sec:constructibility-forcing-01}} & \thepage & Coherent club sequence with the explicit no-thread condition; distinct from the modal necessity glyph. \\
\phantomsection\label{entry:symbol-registry-00378}\(\mathsf{MA}_\kappa,\mathsf{MA}\) & \texttt{\detokenize{\mathsf{MA}_\kappa,}\allowbreak\detokenize{\mathsf{MA}}} & Martin axiom instance and scheme & P6 constructibility-forcing & n/a & \hyperref[sec:constructibility-forcing-01]{\ref*{sec:constructibility-forcing-01}} & \thepage & The instance meets at most kappa dense sets in a ccc poset; the scheme ranges below the continuum. \\
\phantomsection\label{entry:symbol-registry-00379}\(\Delta_0,\Sigma_n,\Pi_n,\Delta_n\) & \texttt{\detokenize{\Delta_0,}\allowbreak\detokenize{\Sigma_n,}\allowbreak\detokenize{\Pi_n,}\allowbreak\detokenize{\Delta_n}} & Levy-hierarchy formula classes & P6 infinitary-combinatorics & n/a & \hyperref[sec:p06-infinitary-combinatorics]{\ref*{sec:p06-infinitary-combinatorics}} & \thepage & Set-theoretic bounded and unbounded-quantifier classes; distinct from arithmetical and projective hierarchies. \\
\phantomsection\label{entry:symbol-registry-00380}\(\mathcal U\) & \texttt{\detokenize{\mathcal U}} & kappa-complete ultrafilter witnessing measurability & P6 infinitary-combinatorics & n/a & \hyperref[sec:p06-infinitary-combinatorics]{\ref*{sec:p06-infinitary-combinatorics}} & \thepage & Calligraphic form avoids the globally reserved bare theory metavariable U. \\
\phantomsection\label{entry:symbol-registry-00381}\(j:V\to M\) & \texttt{\detokenize{j:V\to M}} & elementary class embedding & P6 infinitary-combinatorics to large-cardinals-absoluteness & n/a & \hyperref[sec:p06-infinitary-combinatorics]{\ref*{sec:p06-infinitary-combinatorics}} & \thepage & Metalinguistic embedding read formula by formula; neither V nor M is added as a class object to ZFC. \\
\phantomsection\label{entry:symbol-registry-00382}\({}^{\lambda}M\subseteq M\) & \texttt{\detokenize{{}^{\lambda}M\subseteq M}} & closure under ambient lambda-sequences & P6 infinitary-combinatorics to large-cardinals-absoluteness & meta & \hyperref[sec:p06-infinitary-combinatorics]{\ref*{sec:p06-infinitary-combinatorics}} & \thepage & Every ambient function from lambda into M belongs to M; stronger than transitivity. \\
\phantomsection\label{entry:symbol-registry-00383}\(\operatorname{Ult}(V,\mathcal U)\) & \texttt{\detokenize{\operatorname{Ult}}} & universe ultrapower by a kappa-complete ultrafilter & P6 infinitary-combinatorics & n/a & \hyperref[sec:p06-infinitary-combinatorics]{\ref*{sec:p06-infinitary-combinatorics}} & \thepage & Class ultrapower whose well-foundedness and transitive collapse are cited. \\
\phantomsection\label{entry:symbol-registry-00384}\(\Lang_{\in}\) & \texttt{\detokenize{\Lang_{\in}}} & first-order membership signature & P6 zf-zfc to determinacy & n/a & \hyperref[sec:p06-source-map]{\ref*{sec:p06-source-map}} & \thepage & One-sorted signature with equality and membership as its sole nonlogical relation; classes remain metalinguistic. \\
\phantomsection\label{entry:symbol-registry-00385}\(\operatorname{Ord}\) & \texttt{\detokenize{\operatorname{Ord}}} & proper class of all ordinals & P6 zf-zfc to determinacy & n/a & \hyperref[sec:p06-source-map]{\ref*{sec:p06-source-map}} & \thepage & Metalinguistic class notation and never a set-valued union index. \\
\phantomsection\label{entry:symbol-registry-00386}\(\omega\) & \texttt{\detokenize{\omega}} & least infinite ordinal and von Neumann natural-number set & Legacy pre-batch4 P6-P7.5 reader meaning & n/a & \hyperref[sec:p06-source-map]{\ref*{sec:p06-source-map}} & \thepage & Least inductive set under the displayed ZF construction; P7.5 uses only the ordinal standing.  Numeric chapter scope migrated to its frozen pre-batch4 source-labelled reader regions; glyph, role and precedence unchanged. \\
\phantomsection\label{entry:symbol-registry-00387}\(V\) & \texttt{\detokenize{V}} & ambient universe of sets & P6 zf-zfc to determinacy & n/a & \hyperref[sec:p06-source-map]{\ref*{sec:p06-source-map}} & \thepage & Metalinguistic class abbreviation for the union of the cumulative hierarchy; not a set or a valuation. \\
\phantomsection\label{entry:symbol-registry-00388}\(\sup A\) & \texttt{\detokenize{\sup}} & ordinal supremum & Legacy pre-batch4 P6-P7.5 reader meaning & n/a & \hyperref[sec:p06-source-map]{\ref*{sec:p06-source-map}} & \thepage & Union of a set of ordinals and the least ordinal at least every member; extended to the P7.5 ordinal towers.  Numeric chapter scope migrated to its frozen pre-batch4 source-labelled reader regions; glyph, role and precedence unchanged. \\
\phantomsection\label{entry:symbol-registry-00389}\(\alpha+\beta,\alpha\cdot\beta,\alpha^\beta\) & \texttt{\detokenize{+,}\allowbreak\detokenize{\cdot,}\allowbreak\detokenize{{}^{}}} & ordinal arithmetic operations & Legacy pre-batch4 P6.1-P7.5 reader meaning & n/a & \hyperref[sec:p06-source-map]{\ref*{sec:p06-source-map}} & \thepage & Right-recursive ordinal sum product and exponentiation; distinct from cardinal arithmetic by the declared operands and chapter scope.  Numeric chapter scope migrated to its frozen pre-batch4 source-labelled reader regions; glyph, role and precedence unchanged. \\
\phantomsection\label{entry:symbol-registry-00390}\(A\approx B\) & \texttt{\detokenize{\approx}} & equipotence of sets & P6 cumulative-hierarchy & meta & \hyperref[sec:cumulative-hierarchy-01]{\ref*{sec:cumulative-hierarchy-01}} & \thepage & There is a bijection from A to B; distinct from isomorphism of structures and from formula equivalence. \\
\phantomsection\label{entry:symbol-registry-00391}\(h(A)\) & \texttt{\detokenize{h(A)}} & Hartogs number of a set & P6 cumulative-hierarchy & n/a & \hyperref[sec:cumulative-hierarchy-01]{\ref*{sec:cumulative-hierarchy-01}} & \thepage & Least ordinal admitting no injection into A; its existence is a ZF theorem. \\
\phantomsection\label{entry:symbol-registry-00392}\(\operatorname{Def}(A)\) & \texttt{\detokenize{\operatorname{Def}}} & definable subsets of a membership structure & P6 cardinals-choice & n/a & \hyperref[sec:p06-cardinals-choice]{\ref*{sec:p06-cardinals-choice}} & \thepage & The empty branch is \(\operatorname{Def}(\varnothing)=\{\varnothing\}\). For nonempty A, its members are subsets first-order definable over \((A,\in)\) with finitely many parameters from A; satisfaction is used only on the nonempty structure. \\
\phantomsection\label{entry:symbol-registry-00393}\(L\) & \texttt{\detokenize{L}} & constructible universe & P6 cardinals-choice to infinitary-combinatorics & n/a & \hyperref[sec:p06-cardinals-choice]{\ref*{sec:p06-cardinals-choice}} & \thepage & Proper class union of the constructible levels; distinct from the global signature macro and from a set-sized level. \\
\phantomsection\label{entry:symbol-registry-00447}\(\mathsf{PRA}\) & \texttt{\detokenize{\mathsf{PRA}}} & primitive recursive arithmetic & Legacy pre-batch4 P6-P7.5 reader meaning & n/a & \hyperref[sec:p06-infinitary-combinatorics]{\ref*{sec:p06-infinitary-combinatorics}} & \thepage & Fixed for the master strength map by primitive-recursive function symbols, defining equations, and quantifier-free induction; developed systematically in P7.5.  Numeric chapter scope migrated to its frozen pre-batch4 source-labelled reader regions; glyph, role and precedence unchanged. \\
\phantomsection\label{entry:symbol-registry-00448}\(\mathsf I\Sigma_n\) & \texttt{\detokenize{\mathsf I\Sigma_n}} & induction-fragment family & Legacy pre-batch4 P6-P7.5 reader meaning & n/a & \hyperref[sec:p06-infinitary-combinatorics]{\ref*{sec:p06-infinitary-combinatorics}} & \thepage & Robinson arithmetic plus induction exactly for Sigma-zero-n formulas, for n at least one; P6 first uses the n=1 conservativity calibration and P7.5 gives the systematic treatment.  Numeric chapter scope migrated to its frozen pre-batch4 source-labelled reader regions; glyph, role and precedence unchanged. \\
\phantomsection\label{entry:symbol-registry-00696}\(\prod\mathfrak a\) & \texttt{\detokenize{\prod}} & product space of regular cardinals & P6 proper-forcing-master & n/a & \hyperref[sec:p06-proper-forcing-master]{\ref*{sec:p06-proper-forcing-master}} & \thepage & Functions with f(kappa) below kappa at every coordinate in the nonempty set a; the cardinality of this space is the cardinal product. \\
\phantomsection\label{entry:symbol-registry-00697}\(f={}_Dg,\ f\leq_Dg,\ f<_Dg\) & \texttt{\detokenize{={}_D,}\allowbreak\detokenize{\leq_D,}\allowbreak\detokenize{<_D}} & equality and comparison modulo a filter & P6 proper-forcing-master & meta & \hyperref[sec:p06-proper-forcing-master]{\ref*{sec:p06-proper-forcing-master}} & \thepage & The set of coordinates satisfying the indicated comparison belongs to D; for an ultrafilter the quotient order is linear. \\
\phantomsection\label{entry:symbol-registry-00698}\(\operatorname{pcf}(\mathfrak a)\) & \texttt{\detokenize{\operatorname{pcf}}} & possible cofinalities of a product & P6 proper-forcing-master & n/a & \hyperref[sec:p06-proper-forcing-master]{\ref*{sec:p06-proper-forcing-master}} & \thepage & Cofinalities of the ultrafilter quotients of the product of the infinite regular cardinals in a; principal ultrafilters are included. \\
\phantomsection\label{entry:symbol-registry-00699}\(\omega_4\) & \texttt{\detokenize{\omega_4}} & fourth uncountable initial ordinal & P6 proper-forcing-master & n/a & \hyperref[sec:p06-proper-forcing-master]{\ref*{sec:p06-proper-forcing-master}} & \thepage & Initial ordinal of cardinality aleph-four; used as the aleph index in Shelah's bound at aleph-omega. \\
\phantomsection\label{entry:symbol-registry-01214}\(\operatorname{Tr}(x),\operatorname{Ord}_{\mathrm b}(x)\) & \texttt{\detokenize{\operatorname{Tr},}\allowbreak\detokenize{\operatorname{Ord}_{\mathrm b}}} & fixed bounded membership-language formulas & P6 cardinals-choice & meta & \hyperref[sec:p06-cardinals-choice]{\ref*{sec:p06-cardinals-choice}} & \thepage & Tr expresses transitivity. The bounded ordinal formula includes transitivity of each member and pairwise membership comparison; ambient ZF proves equivalence with the original ordinal definition. \\
\phantomsection\label{entry:symbol-registry-01215}\(\vartheta(f,y,\bar p),G,\operatorname{App}_G(\gamma,f),f_\gamma,F\) & \texttt{\detokenize{\vartheta,}\allowbreak\detokenize{G,}\allowbreak\detokenize{\operatorname{App}_G,}\allowbreak\detokenize{f_\gamma,}\allowbreak\detokenize{F}} & functional formula recursion operation approximations and class function & P6 cardinals-choice & meta & \hyperref[sec:p06-cardinals-choice]{\ref*{sec:p06-cardinals-choice}} & \thepage & The operation subscript records one fixed first-order formula with set parameters, not a class parameter. Every approximation and every restriction to a set of ordinals is a set function. \\
\phantomsection\label{entry:symbol-registry-01216}\(v_i,c,\mathcal C,\psi_i\) & \texttt{\detokenize{v_i,}\allowbreak\detokenize{c,}\allowbreak\detokenize{\mathcal C,}\allowbreak\detokenize{\psi_i}} & variables finite instruction codes their set and unfolded formulas & P6 cardinals-choice & meta & \hyperref[sec:p06-cardinals-choice]{\ref*{sec:p06-cardinals-choice}} & \thepage & Codes are nonempty finite lists of triple records with tags zero through eight. Variable-index fields are distinct from earlier-instruction position fields; sharing a prior instruction preserves its formula and truth value. \\
\phantomsection\label{entry:symbol-registry-01217}\(S_A,s[i\mapsto a],\operatorname{Sat}(A,c,s),\operatorname{Sat}_A\) & \texttt{\detokenize{S_A,}\allowbreak\detokenize{s[i\mapsto a],}\allowbreak\detokenize{\operatorname{Sat},}\allowbreak\detokenize{\operatorname{Sat}_A}} & assignment set update uniform satisfaction formula and set relation & P6 cardinals-choice & meta & \hyperref[sec:p06-cardinals-choice]{\ref*{sec:p06-cardinals-choice}} & \thepage & For nonempty A the assignment set is \(A^\omega\). The construction uses sets and one chosen member of that A, without Choice; it provides no satisfaction predicate for the universe and does not evaluate an empty structure. \\
\phantomsection\label{entry:symbol-registry-01218}\(B,t,T_k,r(k)\) & \texttt{\detokenize{B,}\allowbreak\detokenize{t,}\allowbreak\detokenize{T_k,}\allowbreak\detokenize{r(k)}} & record bounding set truth computation truth subsets and finite index shift & P6 cardinals-choice & meta & \hyperref[sec:p06-cardinals-choice]{\ref*{sec:p06-cardinals-choice}} & \thepage & The bounding set is a product of copies of omega. Each truth value is a subset of the assignment set; concatenation shifts only instruction-position references and leaves variable indices unchanged. \\
\phantomsection\label{entry:symbol-registry-01219}\(G_L,G_V,\Lambda(\alpha,z)\) & \texttt{\detokenize{G_L,}\allowbreak\detokenize{G_V,}\allowbreak\detokenize{\Lambda}} & hierarchy recursion operations and the formula for a level value & P6 cardinals-choice & meta & \hyperref[sec:p06-cardinals-choice]{\ref*{sec:p06-cardinals-choice}} & \thepage & The operations supply empty, successor and limit values by Def or Power Set and Union. The level predicate expands into a first-order formula with set quantifiers. \\
\phantomsection\label{entry:symbol-registry-01220}\(D,H,\tau,\mu\) & \texttt{\detokenize{D,}\allowbreak\detokenize{H,}\allowbreak\detokenize{\tau,}\allowbreak\detokenize{\mu}} & local disagreement sets replacement image and ordinal indices & P6 cardinals-choice & meta & \hyperref[sec:p06-cardinals-choice]{\ref*{sec:p06-cardinals-choice}} & \thepage & These are locally introduced set objects in the approximation and ordinal proofs. Replacement is applied only after the required functional uniqueness is established. \\
\phantomsection\label{entry:symbol-registry-01221}\(O,\nu\) & \texttt{\detokenize{O,}\allowbreak\detokenize{\nu}} & hypothetical set of all ordinals and its union & P6 cardinals-choice & meta & \hyperref[sec:p06-cardinals-choice]{\ref*{sec:p06-cardinals-choice}} & \thepage & The final proper-class argument forms O by Separation under the hypothetical set assumption on L, then uses its ordinal union and successor to derive a membership loop. \\
\phantomsection\label{entry:symbol-registry-05032}\(\mathbf P\) & \texttt{\detokenize{\mathbf P}} & complete forcing poset datum & P6 proper forcing master conditions & object & \hyperref[sec:p06-proper-forcing-master]{\ref*{sec:p06-proper-forcing-master}} & \thepage & Nonempty poset with greatest condition and stronger-is-smaller order \\
\phantomsection\label{entry:symbol-registry-05033}\(P\) & \texttt{\detokenize{P}} & underlying condition set & P6 proper forcing master conditions & object & \hyperref[sec:p06-proper-forcing-master]{\ref*{sec:p06-proper-forcing-master}} & \thepage & Component of the named forcing datum \\
\phantomsection\label{entry:symbol-registry-05034}\(\leq_P\) & \texttt{\detokenize{\leq_P}} & forcing order & P6 proper forcing master conditions & relation & \hyperref[sec:p06-proper-forcing-master]{\ref*{sec:p06-proper-forcing-master}} & \thepage & Left operand is stronger \\
\phantomsection\label{entry:symbol-registry-05035}\(1_P\) & \texttt{\detokenize{1_P}} & greatest weakest condition & P6 proper forcing master conditions & object & \hyperref[sec:p06-proper-forcing-master]{\ref*{sec:p06-proper-forcing-master}} & \thepage & Not an algebraic zero or numeral \\
\phantomsection\label{entry:symbol-registry-05036}\(\operatorname{tc}(\{x\})\) & \texttt{\detokenize{\operatorname{tc}(\{x\})}} & transitive closure containing x as member & P6 proper forcing master conditions & object & \hyperref[def:p06-pfm-hchi]{\ref*{def:p06-pfm-hchi}} & \thepage & Least transitive superset of the singleton \\
\phantomsection\label{entry:symbol-registry-05037}\(H_\chi\) & \texttt{\detokenize{H_\chi}} & hereditary-size set & P6 proper forcing master conditions & object & \hyperref[def:p06-pfm-hchi]{\ref*{def:p06-pfm-hchi}} & \thepage & Regular uncountable height \\
\phantomsection\label{entry:symbol-registry-05038}\(\chi\) & \texttt{\detokenize{\chi}} & regular cardinal height & P6 proper forcing master conditions & object & \hyperref[def:p06-pfm-hchi]{\ref*{def:p06-pfm-hchi}} & \thepage & Greater than two to the power kappa in the properness tests \\
\phantomsection\label{entry:symbol-registry-05039}\(\kappa\) & \texttt{\detokenize{\kappa}} & cardinal bound for hereditary size & P6 proper forcing master conditions & object & \hyperref[def:p06-pfm-hchi]{\ref*{def:p06-pfm-hchi}} & \thepage & Maximum of aleph one and size of transitive closure of the datum \\
\phantomsection\label{entry:symbol-registry-05040}\(N\) & \texttt{\detokenize{N}} & countable nontransitive elementary model domain & P6 proper forcing master conditions & object & \hyperref[def:p06-pfm-hchi]{\ref*{def:p06-pfm-hchi}} & \thepage & Membership structure induced from H chi; not the transitive ground model \\
\phantomsection\label{entry:symbol-registry-05041}\(\prec\) & \texttt{\detokenize{\prec}} & elementary substructure & P6 proper forcing master conditions & relation & \hyperref[def:p06-pfm-hchi]{\ref*{def:p06-pfm-hchi}} & \thepage & Fixed membership or declared expanded structure \\
\phantomsection\label{entry:symbol-registry-05042}\(\triangleleft\) & \texttt{\detokenize{\triangleleft}} & chosen well-order of ambient set & P6 proper forcing master conditions & relation & \hyperref[def:p06-pfm-hulls]{\ref*{def:p06-pfm-hulls}} & \thepage & External Choice parameter for witness selection \\
\phantomsection\label{entry:symbol-registry-05043}\(F_\varphi\) & \texttt{\detokenize{F_\varphi}} & finite-arity witness function & P6 proper forcing master conditions & object & \hyperref[def:p06-pfm-hulls]{\ref*{def:p06-pfm-hulls}} & \thepage & Least ambient witness or fixed default \\
\phantomsection\label{entry:symbol-registry-05044}\(A_n\) & \texttt{\detokenize{A_n}} & closure stage of seed & P6 proper forcing master conditions & object & \hyperref[def:p06-pfm-hulls]{\ref*{def:p06-pfm-hulls}} & \thepage & External natural-indexed construction \\
\phantomsection\label{entry:symbol-registry-05045}\(\operatorname{Hull}(A)\) & \texttt{\detokenize{\operatorname{Hull}(A)}} & countable witness hull & P6 proper forcing master conditions & object & \hyperref[def:p06-pfm-hulls]{\ref*{def:p06-pfm-hulls}} & \thepage & Union of all closure stages \\
\phantomsection\label{entry:symbol-registry-05046}\([H_\chi]^\omega\) & \texttt{\detokenize{[H_\chi]^\omega}} & countable subsets of ambient set & P6 proper forcing master conditions & object & \hyperref[thm:p06-pfm-hulls]{\ref*{thm:p06-pfm-hulls}} & \thepage & Club uses inclusion unboundedness and countable increasing-union closure \\
\phantomsection\label{entry:symbol-registry-05047}\(\delta_N\) & \texttt{\detokenize{\delta_N}} & countable ordinal trace & P6 proper forcing master conditions & object & \hyperref[thm:p06-pfm-hulls]{\ref*{thm:p06-pfm-hulls}} & \thepage & Equals N intersect ground omega one \\
\phantomsection\label{entry:symbol-registry-05048}\(X^\downarrow\) & \texttt{\detokenize{X^\downarrow}} & downward closure in poset & P6 proper forcing master conditions & object & \hyperref[def:p06-pfm-predensity]{\ref*{def:p06-pfm-predensity}} & \thepage & All conditions stronger than a member of X \\
\phantomsection\label{entry:symbol-registry-05049}\(\perp_P\) & \texttt{\detokenize{\perp_P}} & incompatibility & P6 proper forcing master conditions & relation & \hyperref[def:p06-pfm-predensity]{\ref*{def:p06-pfm-predensity}} & \thepage & No common strengthening \\
\phantomsection\label{entry:symbol-registry-05050}\(\dot\alpha\) & \texttt{\detokenize{\dot\alpha}} & forcing name for ground-bounded ordinal & P6 proper forcing master conditions & object & \hyperref[def:p06-pfm-semigeneric]{\ref*{def:p06-pfm-semigeneric}} & \thepage & Globally forced below check omega one \\
\phantomsection\label{entry:symbol-registry-05051}\(\check{\delta_N}\) & \texttt{\detokenize{\check{\delta_N}}} & canonical ground ordinal name & P6 proper forcing master conditions & object & \hyperref[def:p06-pfm-semigeneric]{\ref*{def:p06-pfm-semigeneric}} & \thepage & Not the extension omega-one ordinal \\
\phantomsection\label{entry:symbol-registry-05052}\(p\forces\varphi\) & \texttt{\detokenize{\forces}} & imported recursive forcing relation & P6 proper forcing master conditions & meta & \hyperref[sec:p06-pfm-semiproper]{\ref*{sec:p06-pfm-semiproper}} & \thepage & Relative to the fixed poset and declared transitive ground model; no new forcing theorem \\
\phantomsection\label{entry:symbol-registry-05053}\(T_\chi\) & \texttt{\detokenize{T_\chi}} & set of relevant ordinal names & P6 proper forcing master conditions & object & \hyperref[annex:p06.pfm.semigeneric]{\ref*{annex:p06.pfm.semigeneric}} & \thepage & External forcing predicate on names in H chi \\
\phantomsection\label{entry:symbol-registry-05054}\(A_\tau\) & \texttt{\detokenize{A_\tau}} & deciding maximal antichain & P6 proper forcing master conditions & object & \hyperref[annex:p06.pfm.semigeneric]{\ref*{annex:p06.pfm.semigeneric}} & \thepage & Chosen before the elementary model \\
\phantomsection\label{entry:symbol-registry-05055}\(f_\tau\) & \texttt{\detokenize{f_\tau}} & ordinal-value graph on deciding antichain & P6 proper forcing master conditions & object & \hyperref[annex:p06.pfm.semigeneric]{\ref*{annex:p06.pfm.semigeneric}} & \thepage & Set-valued witness operation in explicit language expansion \\
\phantomsection\label{entry:symbol-registry-05056}\(\tau_n\) & \texttt{\detokenize{\tau_n}} & explicit coordinate ordinal name & P6 proper forcing master conditions & object & \hyperref[annex:p06.pfm.collapse]{\ref*{annex:p06.pfm.collapse}} & \thepage & Finite collapse name with sequence ledger as model parameter \\
\phantomsection\label{entry:symbol-registry-05452}\(\forces^{\mathrm{rec}}_{M,P}\) & \texttt{\detokenize{forces}} & ground-computed density-closed recursive forcing & P6 recursive forcing and truth & relation & \hyperref[def:p06-ftr-atomic]{\ref*{def:p06-ftr-atomic}} & \thepage & Not the earlier external semantic relation; fixed-formula definition \\
\phantomsection\label{entry:symbol-registry-05453}\(\mathsf{Name}_P(\sigma)\) & \texttt{\detokenize{\mathsf{Name}_P(\sigma)}} & first-order hereditary name predicate & P6 recursive forcing and truth & n/a & \hyperref[def:p06-ftr-names]{\ref*{def:p06-ftr-names}} & \thepage & Ground set quantifiers; malformed entries fail Blank macro filled with exact unchanged glyph spelling; this is data, not a newly defined TeX macro. Original precedence descriptor 'predicate' preserved as semantic metadata; n/a is not an invented binding rank. \\
\phantomsection\label{entry:symbol-registry-05454}\(A_\sigma\) & \texttt{\detokenize{A_\sigma}} & set of immediate subnames & P6 recursive forcing and truth & object & \hyperref[def:p06-ftr-names]{\ref*{def:p06-ftr-names}} & \thepage & First coordinates of name--condition pairs Blank macro filled with exact unchanged glyph spelling; this is data, not a newly defined TeX macro. \\
\phantomsection\label{entry:symbol-registry-05455}\(S_\sigma\) & \texttt{\detokenize{S_\sigma}} & subname-closed set including sigma & P6 recursive forcing and truth & object & \hyperref[def:p06-ftr-names]{\ref*{def:p06-ftr-names}} & \thepage & Bounded by transitive closure of singleton sigma Blank macro filled with exact unchanged glyph spelling; this is data, not a newly defined TeX macro. \\
\phantomsection\label{entry:symbol-registry-05456}\(W_\alpha^P\) & \texttt{\detokenize{W_\alpha^P}} & ground-computed all-ordinal name stage & P6 recursive forcing and truth & object & \hyperref[def:p06-ftr-names]{\ref*{def:p06-ftr-names}} & \thepage & Stage zero contains the empty name; ranks at most alpha Blank macro filled with exact unchanged glyph spelling; this is data, not a newly defined TeX macro. \\
\phantomsection\label{entry:symbol-registry-05457}\(\rho_P(\sigma)\) & \texttt{\detokenize{\rho_P(\sigma)}} & intrinsic name rank & P6 recursive forcing and truth & object & \hyperref[def:p06-ftr-names]{\ref*{def:p06-ftr-names}} & \thepage & Supremum of child ranks plus one Blank macro filled with exact unchanged glyph spelling; this is data, not a newly defined TeX macro. \\
\phantomsection\label{entry:symbol-registry-05458}\(\mathcal P^M\) & \texttt{\detokenize{\mathcal P^M}} & power set computed by ground membership structure & P6 recursive forcing and truth & object & \hyperref[def:p06-ftr-names]{\ref*{def:p06-ftr-names}} & \thepage & Not the ambient collection of all subsets Blank macro filled with exact unchanged glyph spelling; this is data, not a newly defined TeX macro. \\
\phantomsection\label{entry:symbol-registry-05459}\(\mathsf{Incl}_P(p,\sigma,\tau)\) & \texttt{\detokenize{\mathsf{Incl}_P(p,}\allowbreak\detokenize{\sigma,}\allowbreak\detokenize{\tau)}} & auxiliary atomic inclusion test & P6 recursive forcing and truth & n/a & \hyperref[def:p06-ftr-atomic]{\ref*{def:p06-ftr-atomic}} & \thepage & Not a new primitive in the membership language Blank macro filled with exact unchanged glyph spelling; this is data, not a newly defined TeX macro. Original precedence descriptor 'predicate' preserved as semantic metadata; n/a is not an invented binding rank. \\
\phantomsection\label{entry:symbol-registry-05460}\(\mathsf{Frc}_{P,\varphi}(p,\bar\tau)\) & \texttt{\detokenize{\mathsf{Frc}_{P,\varphi}(p,}\allowbreak\detokenize{\bar\tau)}} & ground defining formula for one fixed standard formula & P6 recursive forcing and truth & n/a & \hyperref[def:p06-ftr-logical]{\ref*{def:p06-ftr-logical}} & \thepage & No universal internal truth predicate over formula codes Blank macro filled with exact unchanged glyph spelling; this is data, not a newly defined TeX macro. Original precedence descriptor 'predicate' preserved as semantic metadata; n/a is not an invented binding rank. \\
\phantomsection\label{entry:symbol-registry-05461}\((\max(\rho_P(\sigma),\rho_P(\tau)),\min(\rho_P(\sigma),\rho_P(\tau)))\) & \texttt{\detokenize{(\max(\rho_P(\sigma),}\allowbreak\detokenize{\rho_P(\tau)),}\allowbreak\detokenize{\min(\rho_P(\sigma),}\allowbreak\detokenize{\rho_P(\tau)))}} & lexicographic atomic recursion measure & P6 recursive forcing and truth & object & \hyperref[def:p06-ftr-atomic]{\ref*{def:p06-ftr-atomic}} & \thepage & Lowering either coordinate strictly lowers the sorted pair Blank macro filled with exact unchanged glyph spelling; this is data, not a newly defined TeX macro. \\
\phantomsection\label{entry:symbol-registry-05462}\(r,T,T_y,U,g\) & \texttt{\detokenize{r,}\allowbreak\detokenize{T,}\allowbreak\detokenize{T_y,}\allowbreak\detokenize{U,}\allowbreak\detokenize{g}} & physical-rank approximation and bounding sets & P6 forcing foundations annex & n/a & \hyperref[annex:p06.ftr.name-foundations]{\ref*{annex:p06.ftr.name-foundations}} & \thepage & Typed and quantified within the construction Blank macro filled with exact unchanged glyph spelling; this is data, not a newly defined TeX macro. Original precedence descriptor 'local working variables' preserved as semantic metadata; n/a is not an invented binding rank. \\
\phantomsection\label{entry:symbol-registry-05463}\(B_{\sigma,\tau},D_{\sigma,\tau}\) & \texttt{\detokenize{B_{\sigma,\tau},}\allowbreak\detokenize{D_{\sigma,\tau}}} & disagreement-witness set and equality-decision dense set & P6 atomic truth annex & object & \hyperref[annex:p06.ftr.atomic-truth]{\ref*{annex:p06.ftr.atomic-truth}} & \thepage & One ground dense set rather than infinitely many conditions Blank macro filled with exact unchanged glyph spelling; this is data, not a newly defined TeX macro. \\
\phantomsection\label{entry:symbol-registry-05464}\(U(\sigma,\tau),O(\sigma,\tau),L_\sigma\) & \texttt{\detokenize{U(\sigma,}\allowbreak\detokenize{\tau),}\allowbreak\detokenize{O(\sigma,}\allowbreak\detokenize{\tau),}\allowbreak\detokenize{L_\sigma}} & pair ordered-pair and exact-union name constructors & P6 ZFC transfer annex & object & \hyperref[annex:p06.ftr.zfc]{\ref*{annex:p06.ftr.zfc}} & \thepage & Ground set names with directly checked evaluations Blank macro filled with exact unchanged glyph spelling; this is data, not a newly defined TeX macro. \\
\phantomsection\label{entry:symbol-registry-05465}\(I,B,C,R,F,\delta_i,\beta,e\) & \texttt{\detokenize{I,}\allowbreak\detokenize{B,}\allowbreak\detokenize{C,}\allowbreak\detokenize{R,}\allowbreak\detokenize{F,}\allowbreak\detokenize{\delta_i,}\allowbreak\detokenize{\beta,}\allowbreak\detokenize{e}} & bounded candidate sets names and ordinal indices & P6 ZFC transfer annex & n/a & \hyperref[annex:p06.ftr.zfc]{\ref*{annex:p06.ftr.zfc}} & \thepage & Ground Collection Replacement Power Set and Choice constructions Blank macro filled with exact unchanged glyph spelling; this is data, not a newly defined TeX macro. Original precedence descriptor 'local working variables' preserved as semantic metadata; n/a is not an invented binding rank. \\
\phantomsection\label{entry:symbol-registry-05797}\(j(U),n(U),i(p)\) & \texttt{\detokenize{j(U),}\allowbreak\detokenize{n(U),}\allowbreak\detokenize{i(p)}} & regularization complement and condition-image maps & P6 Boolean forcing completion & n/a & \hyperref[def:p06-bfc-regular-open]{\ref*{def:p06-bfc-regular-open}} & \thepage & Downward-open subsets of the forcing poset Blank macro filled with exact unchanged glyph spelling; registry data only, not a new TeX command. Original precedence 'operator' retained as semantic descriptor; n/a is not an invented rank. \\
\phantomsection\label{entry:symbol-registry-05798}\(E_B(u,v),I_B(u,v)\) & \texttt{\detokenize{E_B(u,}\allowbreak\detokenize{v),}\allowbreak\detokenize{I_B(u,}\allowbreak\detokenize{v)}} & Boolean equality and membership values & P6 Boolean forcing consistency & n/a & \hyperref[def:p06-bfc-values]{\ref*{def:p06-bfc-values}} & \thepage & Values lie in the complete Boolean algebra Blank macro filled with exact unchanged glyph spelling; registry data only, not a new TeX command. Original precedence 'predicate value' retained as semantic descriptor; n/a is not an invented rank. \\
\phantomsection\label{entry:symbol-registry-05799}\(\lVert\varphi\rVert_B\) & \texttt{\detokenize{\lVert\varphi\rVert_B}} & fixed-formula Boolean truth value & P6 Boolean forcing consistency & object & \hyperref[def:p06-bfc-values]{\ref*{def:p06-bfc-values}} & \thepage & Not a norm or an all-formula-code truth predicate Blank macro filled with exact unchanged glyph spelling; registry data only, not a new TeX command. \\
\phantomsection\label{entry:symbol-registry-05800}\(N_B(u)\) & \texttt{\detokenize{N_B(u)}} & first-order function-style Boolean-name predicate & P6 Boolean forcing consistency & n/a & \hyperref[def:p06-bfc-values]{\ref*{def:p06-bfc-values}} & \thepage & Internal bounded rank construction Blank macro filled with exact unchanged glyph spelling; registry data only, not a new TeX command. Original precedence 'predicate' retained as semantic descriptor; n/a is not an invented rank. \\
\phantomsection\label{entry:symbol-registry-05801}\(\mathcal M=(D,E)\) & \texttt{\detokenize{\mathcal M=(D,}\allowbreak\detokenize{E)}} & possibly externally ill-founded membership structure & P6 Boolean quotient consistency & n/a & \hyperref[def:p06-bfc-ill-founded]{\ref*{def:p06-bfc-ill-founded}} & \thepage & Internal Foundation is not external well-foundedness Blank macro filled with exact unchanged glyph spelling; registry data only, not a new TeX command. Original precedence 'structure' retained as semantic descriptor; n/a is not an invented rank. \\
\phantomsection\label{entry:symbol-registry-05802}\(B^{\mathcal M},\mathcal D_b,v^{\mathcal M}_\varphi\) & \texttt{\detokenize{B^{\mathcal M},}\allowbreak\detokenize{\mathcal D_b,}\allowbreak\detokenize{v^{\mathcal M}_\varphi}} & interpreted algebra names and fixed-formula value function & P6 Boolean quotient consistency & object & \hyperref[def:p06-bfc-ill-founded]{\ref*{def:p06-bfc-ill-founded}} & \thepage & External algebra need not be complete Blank macro filled with exact unchanged glyph spelling; registry data only, not a new TeX command. \\
\phantomsection\label{entry:symbol-registry-05803}\(\sim_U,\in_U,\mathcal Q\) & \texttt{\detokenize{\sim_U,}\allowbreak\detokenize{\in_U,}\allowbreak\detokenize{\mathcal Q}} & equivalence quotient membership and actual quotient structure & P6 Boolean quotient consistency & relation & \hyperref[thm:p06-bfc-quotient]{\ref*{thm:p06-bfc-quotient}} & \thepage & No external recursive evaluation of internal names Blank macro filled with exact unchanged glyph spelling; registry data only, not a new TeX command. \\
\phantomsection\label{entry:symbol-registry-05868}\(\mathcal N\) & \texttt{\detokenize{\mathcal N}} & null ideal of arbitrary real subsets & P6 cardinal characteristics & object & \hyperref[def:p06-cc-standing]{\ref*{def:p06-cc-standing}} & \thepage & Root restored actual Fremlin source glyph. Ideal of arbitrary real subsets, explicitly not a structure or ground model. Original I\_null/I\_meager row is superseded evidence, not active notation. \\
\phantomsection\label{entry:symbol-registry-05869}\(\mathcal M\) & \texttt{\detokenize{\mathcal M}} & meager ideal of arbitrary real subsets & P6 cardinal characteristics & object & \hyperref[def:p06-cc-standing]{\ref*{def:p06-cc-standing}} & \thepage & Root restored actual Fremlin source glyph. Ideal of arbitrary real subsets, explicitly not a structure or ground model. Original I\_null/I\_meager row is superseded evidence, not active notation. \\
\phantomsection\label{entry:symbol-registry-05870}\(\operatorname{add}(I),\operatorname{cov}(I),\operatorname{non}(I),\operatorname{cof}(I)\) & \texttt{\detokenize{\operatorname{add}(I),}\allowbreak\detokenize{\operatorname{cov}(I),}\allowbreak\detokenize{\operatorname{non}(I),}\allowbreak\detokenize{\operatorname{cof}(I)}} & four ideal witness cardinals & P6 cardinal characteristics & n/a & \hyperref[def:p06-cc-ideal-four]{\ref*{def:p06-cc-ideal-four}} & \thepage & cof is ideal cofinality; Fremlin writes cf for this ideal invariant Blank macro filled with exact unchanged glyph spelling; registry data only, not a new TeX command. Original precedence 'operator' retained as semantic descriptor; n/a is not an invented rank. \\
\phantomsection\label{entry:symbol-registry-05871}\(\leq^*\) & \texttt{\detokenize{\leq^*}} & eventual domination on omega to omega & P6 cardinal characteristics & relation & \hyperref[def:p06-cc-bdsr]{\ref*{def:p06-cc-bdsr}} & \thepage & Finite exceptional coordinates allowed; not pointwise domination Blank macro filled with exact unchanged glyph spelling; registry data only, not a new TeX command. \\
\phantomsection\label{entry:symbol-registry-05872}\(\mathfrak b,\mathfrak d\) & \texttt{\detokenize{\mathfrak b,}\allowbreak\detokenize{\mathfrak d}} & least unbounded and dominating family sizes & P6 cardinal characteristics & n/a & \hyperref[def:p06-cc-bdsr]{\ref*{def:p06-cc-bdsr}} & \thepage & Separate existential and universal witness conditions Blank macro filled with exact unchanged glyph spelling; registry data only, not a new TeX command. Original precedence 'cardinal' retained as semantic descriptor; n/a is not an invented rank. \\
\phantomsection\label{entry:symbol-registry-05873}\(\mathfrak s,\mathfrak r\) & \texttt{\detokenize{\mathfrak s,}\allowbreak\detokenize{\mathfrak r}} & splitting and reaping family sizes & P6 cardinal characteristics & n/a & \hyperref[def:p06-cc-bdsr]{\ref*{def:p06-cc-bdsr}} & \thepage & Blass also calls r the unsplitting number Blank macro filled with exact unchanged glyph spelling; registry data only, not a new TeX command. Original precedence 'cardinal' retained as semantic descriptor; n/a is not an invented rank. \\
\phantomsection\label{entry:symbol-registry-05874}\(\mathfrak c\) & \texttt{\detokenize{\mathfrak c}} & cardinality of the continuum & P6 cardinal characteristics & n/a & \hyperref[def:p06-cc-standing]{\ref*{def:p06-cc-standing}} & \thepage & No characteristic-specific value inferred from c alone Root source-label reconciliation entry:p06-cc-standing -> def:p06-cc-standing; original intake preserved. Blank macro filled with exact unchanged glyph spelling; registry data only, not a new TeX command. Original precedence 'cardinal' retained as semantic descriptor; n/a is not an invented rank. \\
\phantomsection\label{entry:symbol-registry-05875}\(\Phi,\Psi\) & \texttt{\detokenize{\Phi,}\allowbreak\detokenize{\Psi}} & closed category witness and recovered eventual bound & P6 category coding & n/a & \hyperref[thm:p06-cc-category-coding]{\ref*{thm:p06-cc-category-coding}} & \thepage & Containment implication stated locally; not an inference rule Blank macro filled with exact unchanged glyph spelling; registry data only, not a new TeX command. Original precedence 'function' retained as semantic descriptor; n/a is not an invented rank. \\
\phantomsection\label{entry:symbol-registry-05876}\(f\in^*S\) & \texttt{\detokenize{f\in^*S}} & eventual finite-width localization & P6 cardinal characteristics & relation & \hyperref[def:p06-cc-slaloms]{\ref*{def:p06-cc-slaloms}} & \thepage & Every sufficiently late coordinate matches; infinitely often is insufficient Blank macro filled with exact unchanged glyph spelling; registry data only, not a new TeX command. \\
\phantomsection\label{entry:symbol-registry-05877}\(a,b,c,e\) & \texttt{\detokenize{a,}\allowbreak\detokenize{b,}\allowbreak\detokenize{c,}\allowbreak\detokenize{e}} & the four null localization maps & P6 null localization & n/a & \hyperref[thm:p06-cc-null-localization]{\ref*{thm:p06-cc-null-localization}} & \thepage & Local map letters; c is not the continuum cardinal mathfrak c Blank macro filled with exact unchanged glyph spelling; registry data only, not a new TeX command. Original precedence 'function' retained as semantic descriptor; n/a is not an invented rank. \\
\phantomsection\label{entry:symbol-registry-05878}\(u,v\) & \texttt{\detokenize{u,}\allowbreak\detokenize{v}} & meager to null containment transport maps & P6 localization bridge & n/a & \hyperref[thm:p06-cc-meager-null-transport]{\ref*{thm:p06-cc-meager-null-transport}} & \thepage & No inclusion of one ideal in the other asserted Blank macro filled with exact unchanged glyph spelling; registry data only, not a new TeX command. Original precedence 'function' retained as semantic descriptor; n/a is not an invented rank. \\
\phantomsection\label{entry:symbol-registry-05879}\(\operatorname{cf}(I)\) & \texttt{\detokenize{\operatorname{cf}}} & Fremlin ideal cofinal-base minimum & P6 cardinal characteristics & n/a & \hyperref[def:p06-cc-ideal-four]{\ref*{def:p06-cc-ideal-four}} & \thepage & Actual cf(I) source spelling restored beside cof(I); equality uses the SAME ideal and cofinal-base minimum. Ideal argument type differs from ordinal cf(alpha). Original grouped invariant row remains unchanged as evidence. \\
\phantomsection\label{entry:symbol-registry-06053}\(M,\kappa,U\) & \texttt{\detokenize{M,}\allowbreak\detokenize{\kappa,}\allowbreak\detokenize{U}} & transitive ground and its uncountable normal-measure cardinal and measure & P6 classical Prikry construction & object & \hyperref[def:p06-prikry-ground]{\ref*{def:p06-prikry-ground}} & \thepage & All completeness and measure membership claims are internal to M Blank macro filled with unchanged glyph spelling; registry data only, no new TeX command. \\
\phantomsection\label{entry:symbol-registry-06054}\(\triangle_{\xi<\kappa}A_\xi\) & \texttt{\detokenize{\triangle_{\xi<\kappa}A_\xi}} & diagonal intersection of measure-one sets & P6 Prikry normality & n/a & \hyperref[def:p06-prikry-normal-measure]{\ref*{def:p06-prikry-normal-measure}} & \thepage & At beta the index quantifier ranges only over xi below beta Blank macro filled with unchanged glyph spelling; registry data only, no new TeX command. Original precedence descriptor 'operator' preserved; n/a is not an invented rank. \\
\phantomsection\label{entry:symbol-registry-06055}\(P_U\) & \texttt{\detokenize{P_U}} & classical normal-measure Prikry poset & P6 classical Prikry construction & object & \hyperref[def:p06-prikry-poset]{\ref*{def:p06-prikry-poset}} & \thepage & Finite increasing stem and measure-one tail; not tree or extender forcing Blank macro filled with unchanged glyph spelling; registry data only, no new TeX command. \\
\phantomsection\label{entry:symbol-registry-06056}\((s_p,A_p)\) & \texttt{\detokenize{(s_p,}\allowbreak\detokenize{A_p)}} & stem and upper part of a condition & P6 classical Prikry construction & object & \hyperref[def:p06-prikry-poset]{\ref*{def:p06-prikry-poset}} & \thepage & Tail points are above every point of the stem Blank macro filled with unchanged glyph spelling; registry data only, no new TeX command. \\
\phantomsection\label{entry:symbol-registry-06057}\(\leq,\leq^*\) & \texttt{\detokenize{\leq,}\allowbreak\detokenize{\leq^*}} & full and direct Prikry extension orders & P6 classical Prikry construction & relation & \hyperref[def:p06-prikry-poset]{\ref*{def:p06-prikry-poset}} & \thepage & Smaller is stronger here; Gitik author draft uses opposite greater-is-stronger glyphs Blank macro filled with unchanged glyph spelling; registry data only, no new TeX command. \\
\phantomsection\label{entry:symbol-registry-06058}\(s\mathbin{\frown}t\) & \texttt{\detokenize{s\mathbin{\frown}t}} & concatenation of increasing stems & P6 classical Prikry construction & n/a & \hyperref[thm:p06-prikry-strong-dense]{\ref*{thm:p06-prikry-strong-dense}} & \thepage & All entries of t lie above s Blank macro filled with unchanged glyph spelling; registry data only, no new TeX command. Original precedence descriptor 'operator' preserved; n/a is not an invented rank. \\
\phantomsection\label{entry:symbol-registry-06059}\(C_G\) & \texttt{\detokenize{C_G}} & generic increasing omega-sequence and its graph & P6 classical Prikry construction & object & \hyperref[thm:p06-prikry-sequence]{\ref*{thm:p06-prikry-sequence}} & \thepage & Union of compatible stems; no external cardinality assertion about kappa Blank macro filled with unchanged glyph spelling; registry data only, no new TeX command. \\
\phantomsection\label{entry:symbol-registry-06060}\(G(C)\) & \texttt{\detokenize{G(C)}} & filter reconstructed from a sequence & P6 Prikry genericity & n/a & \hyperref[thm:p06-prikry-mathias]{\ref*{thm:p06-prikry-mathias}} & \thepage & Initial stem plus exact tail containment; eventual measure containment gives genericity Blank macro filled with unchanged glyph spelling; registry data only, no new TeX command. Original precedence descriptor 'function' preserved; n/a is not an invented rank. \\
\phantomsection\label{entry:symbol-registry-06061}\(\mathcal P(\lambda)^M,\mathcal P(\lambda)^{M[G]}\) & \texttt{\detokenize{\mathcal P(\lambda)^M,}\allowbreak\detokenize{\mathcal P(\lambda)^{M[G]}}} & ground and extension internal power sets & P6 Prikry preservation & object & \hyperref[thm:p06-prikry-no-bounded]{\ref*{thm:p06-prikry-no-bounded}} & \thepage & Equality asserted only for lambda below kappa Blank macro filled with unchanged glyph spelling; registry data only, no new TeX command. \\
\phantomsection\label{entry:symbol-registry-06062}\(\operatorname{cf}^M,\operatorname{cf}^{M[G]}\) & \texttt{\detokenize{\operatorname{cf}^M,}\allowbreak\detokenize{\operatorname{cf}^{M[G]}}} & internal ground and extension cofinalities & P6 Prikry preservation & n/a & \hyperref[thm:p06-prikry-preservation]{\ref*{thm:p06-prikry-preservation}} & \thepage & Only ground cofinality kappa changes in the displayed classification Blank macro filled with unchanged glyph spelling; registry data only, no new TeX command. Original precedence descriptor 'operator' preserved; n/a is not an invented rank. \\
\phantomsection\label{entry:symbol-registry-06063}\(P\) & \texttt{\detokenize{P}} & Gitik finite-subset-stem Prikry condition set & P6 Gitik native Prikry orders & object & \hyperref[def:p06-prikry-gitik-orders]{\ref*{def:p06-prikry-gitik-orders}} & \thepage & Distinct condition representation from sequence-stem P\_U; increasing enumeration is the proved bridge Blank macro filled with unchanged glyph spelling; registry data only, no new TeX command. \\
\phantomsection\label{entry:symbol-registry-06064}\(\langle p,A\rangle\) & \texttt{\detokenize{\langle p,}\allowbreak\detokenize{A\rangle}} & native finite subset and measure-one upper part & P6 Gitik native Prikry orders & object & \hyperref[def:p06-prikry-gitik-orders]{\ref*{def:p06-prikry-gitik-orders}} & \thepage & Finite subset p is not a sequence or a propositional formula Blank macro filled with unchanged glyph spelling; registry data only, no new TeX command. \\
\phantomsection\label{entry:symbol-registry-06065}\(\geq,\geq^*\) & \texttt{\detokenize{\geq,}\allowbreak\detokenize{\geq^*}} & native stronger and directly stronger relations on finite-set conditions & P6 Gitik native Prikry orders & relation & \hyperref[def:p06-prikry-gitik-orders]{\ref*{def:p06-prikry-gitik-orders}} & \thepage & Source Definitions1.2 and1.3 preserved with greater-is-stronger orientation Blank macro filled with unchanged glyph spelling; registry data only, no new TeX command. \\
\phantomsection\label{entry:symbol-registry-06066}\(\leq,\leq^*\) & \texttt{\detokenize{\leq,}\allowbreak\detokenize{\leq^*}} & native converse weak orders on finite-set conditions & P6 Gitik native Prikry orders & relation & \hyperref[def:p06-prikry-gitik-orders]{\ref*{def:p06-prikry-gitik-orders}} & \thepage & These native weak orders reverse the stronger relation; not the host smaller-is-stronger sequence-stem orders Blank macro filled with unchanged glyph spelling; registry data only, no new TeX command. \\
\phantomsection\label{entry:symbol-registry-06093}\(U,\kappa\) & \texttt{\detokenize{U,}\allowbreak\detokenize{\kappa}} & ground nonprincipal kappa-complete ultrafilter and its cardinal & P6 native tree-Prikry & object & \hyperref[def:p06-tree-prikry-ground]{\ref*{def:p06-tree-prikry-ground}} & \thepage & Normality is additional only in the final comparison theorem Blank macro filled with unchanged glyph spelling; registry data only, no new TeX command. \\
\phantomsection\label{entry:symbol-registry-06094}\(P\) & \texttt{\detokenize{P}} & native tree-Prikry condition set & P6 native tree-Prikry & object & \hyperref[def:p06-tree-prikry-native-orders]{\ref*{def:p06-tree-prikry-native-orders}} & \thepage & Typed pairs of trunks and trees; not the classical finite-set condition type Blank macro filled with unchanged glyph spelling; registry data only, no new TeX command. \\
\phantomsection\label{entry:symbol-registry-06095}\(\langle t,T\rangle\) & \texttt{\detokenize{\langle t,}\allowbreak\detokenize{T\rangle}} & native trunk and U-tree condition & P6 native tree-Prikry & object & \hyperref[def:p06-tree-prikry-native-orders]{\ref*{def:p06-tree-prikry-native-orders}} & \thepage & Source pair syntax retained Blank macro filled with unchanged glyph spelling; registry data only, no new TeX command. \\
\phantomsection\label{entry:symbol-registry-06096}\(\operatorname{Suc}_T(\eta),\operatorname{Lev}_n(T)\) & \texttt{\detokenize{\operatorname{Suc}_T(\eta),}\allowbreak\detokenize{\operatorname{Lev}_n(T)}} & native successor set and absolute tree level & P6 native tree-Prikry & n/a & \hyperref[def:p06-tree-prikry-trees]{\ref*{def:p06-tree-prikry-trees}} & \thepage & Each retained successor set is measure one Blank macro filled with unchanged glyph spelling; registry data only, no new TeX command. Original precedence descriptor 'operator' preserved; n/a is not an invented rank. \\
\phantomsection\label{entry:symbol-registry-06097}\((T)_r\) & \texttt{\detokenize{(T)_r}} & native rooted cone above a node & P6 native tree-Prikry & n/a & \hyperref[def:p06-tree-prikry-trees]{\ref*{def:p06-tree-prikry-trees}} & \thepage & Only extensions of r; no compulsory ancestors below r Blank macro filled with unchanged glyph spelling; registry data only, no new TeX command. Original precedence descriptor 'operator' preserved; n/a is not an invented rank. \\
\phantomsection\label{entry:symbol-registry-06098}\(\geq,\geq^*,\leq,\leq^*\) & \texttt{\detokenize{\geq,}\allowbreak\detokenize{\geq^*,}\allowbreak\detokenize{\leq,}\allowbreak\detokenize{\leq^*}} & native stronger and direct-stronger orders and their weak converses & P6 native tree-Prikry & relation & \hyperref[def:p06-tree-prikry-native-orders]{\ref*{def:p06-tree-prikry-native-orders}} & \thepage & Greater is stronger; not normalized into host smaller-is-stronger syntax Blank macro filled with unchanged glyph spelling; registry data only, no new TeX command. \\
\phantomsection\label{entry:symbol-registry-06099}\(\leq_h,\leq_h^*\) & \texttt{\detokenize{\leq_h,}\allowbreak\detokenize{\leq_h^*}} & explicit host order-dual and its direct order & P6 tree-Prikry host bridge & relation & \hyperref[thm:p06-tree-prikry-host]{\ref*{thm:p06-tree-prikry-host}} & \thepage & Same condition set and strengthening sets; native glyphs remain in use Blank macro filled with unchanged glyph spelling; registry data only, no new TeX command. \\
\phantomsection\label{entry:symbol-registry-06100}\(|t|\) & \texttt{\detokenize{|t|}} & source finite-stem length and finite cardinality & P6 native tree-Prikry & n/a & \hyperref[def:p06-tree-prikry-trees]{\ref*{def:p06-tree-prikry-trees}} & \thepage & For a finite sequence its function cardinality equals its finite domain length Blank macro filled with unchanged glyph spelling; registry data only, no new TeX command. Original precedence descriptor 'cardinal' preserved; n/a is not an invented rank. \\
\phantomsection\label{entry:symbol-registry-06101}\({}^{i}\sigma\) & \texttt{\detokenize{{}^{i}\sigma}} & native positive or negative forcing sentence sign & P6 native tree decision & object & \hyperref[thm:p06-tree-prikry-decision]{\ref*{thm:p06-tree-prikry-decision}} & \thepage & Source convention: sign zero is sigma and sign one is negation Blank macro filled with unchanged glyph spelling; registry data only, no new TeX command. \\
\phantomsection\label{entry:symbol-registry-06102}\(C,G(C)\) & \texttt{\detokenize{C,}\allowbreak\detokenize{G(C)}} & generic branch and its reconstructed tree filter & P6 tree-Prikry genericity & object & \hyperref[thm:p06-tree-prikry-branch]{\ref*{thm:p06-tree-prikry-branch}} & \thepage & Branch membership in each tree is explicit Blank macro filled with unchanged glyph spelling; registry data only, no new TeX command. \\
\phantomsection\label{entry:symbol-registry-06103}\(s\mapsto T_s\) & \texttt{\detokenize{s\mapsto T_s}} & ground assignment of a tree to every finite stem & P6 tree-Prikry family criterion & n/a & \hyperref[thm:p06-tree-prikry-family-criterion]{\ref*{thm:p06-tree-prikry-family-criterion}} & \thepage & Not replaced by individual eventual measure-set containment without normality Blank macro filled with unchanged glyph spelling; registry data only, no new TeX command. Original precedence descriptor 'function' preserved; n/a is not an invented rank. \\
\phantomsection\label{entry:symbol-registry-06104}\(\iota\) & \texttt{\detokenize{\iota}} & conditional dense classical-to-tree embedding & P6 normal tree comparison & n/a & \hyperref[thm:p06-tree-prikry-normal-comparison]{\ref*{thm:p06-tree-prikry-normal-comparison}} & \thepage & Normality supplies its direct density; native order correspondence proved Blank macro filled with unchanged glyph spelling; registry data only, no new TeX command. Original precedence descriptor 'function' preserved; n/a is not an invented rank. \\
\phantomsection\label{entry:symbol-registry-07245}\(\mathcal B(\dot Q)\) & \texttt{\detokenize{\mathcal B(\dot Q)}} & bounded mixed total condition-name set & P6 iteration names & object & \hyperref[thm:p06-iter-name-bound]{\ref*{thm:p06-iter-name-bound}} & \thepage & Not an unrestricted proper class Blank macro filled with the glyph's TeX spelling. \\
\phantomsection\label{entry:symbol-registry-07246}\(\operatorname{Mix}(A,\tau)\) & \texttt{\detokenize{\operatorname{Mix}(A,}\allowbreak\detokenize{\tau)}} & antichain mixing constructor & P6 iteration names & n/a & \hyperref[thm:p06-iter-name-bound]{\ref*{thm:p06-iter-name-bound}} & \thepage & Supports refine antichain and original supports Original precedence descriptor 'operator' retained; no rank inferred. Blank macro filled with the glyph's TeX spelling. \\
\phantomsection\label{entry:symbol-registry-07247}\(P*\dot Q\) & \texttt{\detokenize{P*\dot Q}} & raw two-step preorder & P6 iterations & object & \hyperref[def:p06-iter-two-step]{\ref*{def:p06-iter-two-step}} & \thepage & Karagila smaller-is-stronger convention retained Blank macro filled with the glyph's TeX spelling. \\
\phantomsection\label{entry:symbol-registry-07248}\(\approx\) & \texttt{\detokenize{\approx}} & mutual-extension equivalence & P6 iterations & relation & \hyperref[def:p06-iter-two-step]{\ref*{def:p06-iter-two-step}} & \thepage & Quotient partial order is distinguished from raw names Blank macro filled with the glyph's TeX spelling. \\
\phantomsection\label{entry:symbol-registry-07249}\(i\) & \texttt{\detokenize{i}} & unit embedding and projection symbol family member \(i\) & P6 iteration initial segments & n/a & \hyperref[thm:p06-iter-preorder]{\ref*{thm:p06-iter-preorder}} & \thepage & Order and antichain properties proved Original precedence descriptor 'function' retained; no rank inferred. Explicit displayed family; member types and common actual first-use exposed, not one atomic glyph. Blank macro filled with the glyph's TeX spelling. Family listing, not one atomic glyph. \\
\phantomsection\label{entry:symbol-registry-07250}\(\pi\) & \texttt{\detokenize{\pi}} & unit embedding and projection symbol family member \(\pi\) & P6 iteration initial segments & n/a & \hyperref[thm:p06-iter-preorder]{\ref*{thm:p06-iter-preorder}} & \thepage & Order and antichain properties proved Original precedence descriptor 'function' retained; no rank inferred. Explicit displayed family; member types and common actual first-use exposed, not one atomic glyph. Blank macro filled with the glyph's TeX spelling. Family listing, not one atomic glyph. \\
\phantomsection\label{entry:symbol-registry-07251}\(\tau^\downarrow\) & \texttt{\detokenize{\tau^\downarrow}} & P-name for second-stage name & P6 two-step names & object & \hyperref[thm:p06-iter-two-generics]{\ref*{thm:p06-iter-two-generics}} & \thepage & Double evaluation equals two-step evaluation Blank macro filled with the glyph's TeX spelling. \\
\phantomsection\label{entry:symbol-registry-07252}\(P_\alpha\) & \texttt{\detokenize{P_\alpha}} & stages and iterand names symbol family member \(P_\alpha\) & P6 support recursion & object & \hyperref[def:p06-iter-recursion]{\ref*{def:p06-iter-recursion}} & \thepage & Successor quotients and limit support explicit Explicit displayed family; member types and common actual first-use exposed, not one atomic glyph. Blank macro filled with the glyph's TeX spelling. Family listing, not one atomic glyph. \\
\phantomsection\label{entry:symbol-registry-07253}\(\dot Q_\xi\) & \texttt{\detokenize{\dot Q_\xi}} & stages and iterand names symbol family member \(\dot Q_\xi\) & P6 support recursion & object & \hyperref[def:p06-iter-recursion]{\ref*{def:p06-iter-recursion}} & \thepage & Successor quotients and limit support explicit Explicit displayed family; member types and common actual first-use exposed, not one atomic glyph. Blank macro filled with the glyph's TeX spelling. Family listing, not one atomic glyph. \\
\phantomsection\label{entry:symbol-registry-07254}\(\operatorname{supp}(p)\) & \texttt{\detokenize{\operatorname{supp}(p)}} & listed syntactic nonunit coordinates & P6 local support recursion & n/a & \hyperref[def:p06-iter-recursion]{\ref*{def:p06-iter-recursion}} & \thepage & Conditional nonunit names are not silently omitted Original precedence descriptor 'operator' retained; no rank inferred. Blank macro filled with the glyph's TeX spelling. \\
\phantomsection\label{entry:symbol-registry-07255}\(\pi_{\beta\alpha}\) & \texttt{\detokenize{\pi_{\beta\alpha}}} & coherent restrictions and complete embeddings symbol family member \(\pi_{\beta\alpha}\) & P6 iteration initial segments & n/a & \hyperref[thm:p06-iter-initial-segments]{\ref*{thm:p06-iter-initial-segments}} & \thepage & Maps are quotient invariant Original precedence descriptor 'function' retained; no rank inferred. Explicit displayed family; member types and common actual first-use exposed, not one atomic glyph. Blank macro filled with the glyph's TeX spelling. Family listing, not one atomic glyph. \\
\phantomsection\label{entry:symbol-registry-07256}\(i_{\alpha\beta}\) & \texttt{\detokenize{i_{\alpha\beta}}} & coherent restrictions and complete embeddings symbol family member \(i_{\alpha\beta}\) & P6 iteration initial segments & n/a & \hyperref[thm:p06-iter-initial-segments]{\ref*{thm:p06-iter-initial-segments}} & \thepage & Maps are quotient invariant Original precedence descriptor 'function' retained; no rank inferred. Explicit displayed family; member types and common actual first-use exposed, not one atomic glyph. Blank macro filled with the glyph's TeX spelling. Family listing, not one atomic glyph. \\
\phantomsection\label{entry:symbol-registry-07257}\(A_\theta\) & \texttt{\detokenize{A_\theta}} & unit and theta incompatible atoms & P6 support examples & object & \hyperref[thm:p06-iter-support-examples]{\ref*{thm:p06-iter-support-examples}} & \thepage & Individually countably closed and proper Blank macro filled with the glyph's TeX spelling. \\
\phantomsection\label{entry:symbol-registry-07258}\(\geq\) & \texttt{\detokenize{\geq}} & native Magidor full and direct orders with weak converses symbol family member \(\geq\) & P6 Gitik comparison & relation & \hyperref[def:p06-iter-gitik-native]{\ref*{def:p06-iter-gitik-native}} & \thepage & Greater is stronger; finite nondirect exceptions remain Explicit displayed family; member types and common actual first-use exposed, not one atomic glyph. Blank macro filled with the glyph's TeX spelling. Family listing, not one atomic glyph. \\
\phantomsection\label{entry:symbol-registry-07259}\(\geq^*\) & \texttt{\detokenize{\geq^*}} & native Magidor full and direct orders with weak converses symbol family member \(\geq^*\) & P6 Gitik comparison & relation & \hyperref[def:p06-iter-gitik-native]{\ref*{def:p06-iter-gitik-native}} & \thepage & Greater is stronger; finite nondirect exceptions remain Explicit displayed family; member types and common actual first-use exposed, not one atomic glyph. Blank macro filled with the glyph's TeX spelling. Family listing, not one atomic glyph. \\
\phantomsection\label{entry:symbol-registry-07260}\(\leq\) & \texttt{\detokenize{\leq}} & native Magidor full and direct orders with weak converses symbol family member \(\leq\) & P6 Gitik comparison & relation & \hyperref[def:p06-iter-gitik-native]{\ref*{def:p06-iter-gitik-native}} & \thepage & Greater is stronger; finite nondirect exceptions remain Explicit displayed family; member types and common actual first-use exposed, not one atomic glyph. Blank macro filled with the glyph's TeX spelling. Family listing, not one atomic glyph. \\
\phantomsection\label{entry:symbol-registry-07261}\(\leq^*\) & \texttt{\detokenize{\leq^*}} & native Magidor full and direct orders with weak converses symbol family member \(\leq^*\) & P6 Gitik comparison & relation & \hyperref[def:p06-iter-gitik-native]{\ref*{def:p06-iter-gitik-native}} & \thepage & Greater is stronger; finite nondirect exceptions remain Explicit displayed family; member types and common actual first-use exposed, not one atomic glyph. Blank macro filled with the glyph's TeX spelling. Family listing, not one atomic glyph. \\
\phantomsection\label{entry:symbol-registry-07262}\(\leq_g\) & \texttt{\detokenize{\leq_g}} & separate host orientation symbol family member \(\leq_g\) & P6 Gitik comparison & relation & \hyperref[thm:p06-iter-native-bridge]{\ref*{thm:p06-iter-native-bridge}} & \thepage & Does not replace native full support Explicit displayed family; member types and common actual first-use exposed, not one atomic glyph. Blank macro filled with the glyph's TeX spelling. Family listing, not one atomic glyph. \\
\phantomsection\label{entry:symbol-registry-07263}\(\leq_g^*\) & \texttt{\detokenize{\leq_g^*}} & separate host orientation symbol family member \(\leq_g^*\) & P6 Gitik comparison & relation & \hyperref[thm:p06-iter-native-bridge]{\ref*{thm:p06-iter-native-bridge}} & \thepage & Does not replace native full support Explicit displayed family; member types and common actual first-use exposed, not one atomic glyph. Blank macro filled with the glyph's TeX spelling. Family listing, not one atomic glyph. \\
\phantomsection\label{entry:symbol-registry-07463}\(N[G]\) & \texttt{\detokenize{N[G]}} & evaluated hereditary-small name model & P6 proper iteration & n/a & \hyperref[def:p06-pil-lift]{\ref*{def:p06-pil-lift}} & \thepage & Not N union G; countable elementary lifting and ground trace proved Original precedence descriptor 'structure' retained; no rank inferred. Blank macro filled with the glyph's TeX spelling. \\
\phantomsection\label{entry:symbol-registry-07464}\(R_\xi\) & \texttt{\detokenize{R_\xi}} & raw condition and quotient stage sets symbol family member \(R_\xi\) & P6 proper iteration & object & \hyperref[def:p06-pil-datum]{\ref*{def:p06-pil-datum}} & \thepage & Sections and quotient maps are in the full iteration datum Explicit displayed family; member types and common actual first-use exposed, not one atomic glyph. Blank macro filled with the glyph's TeX spelling. Family listing, not one atomic glyph. \\
\phantomsection\label{entry:symbol-registry-07465}\(P_\xi\) & \texttt{\detokenize{P_\xi}} & raw condition and quotient stage sets symbol family member \(P_\xi\) & P6 proper iteration & object & \hyperref[def:p06-pil-datum]{\ref*{def:p06-pil-datum}} & \thepage & Sections and quotient maps are in the full iteration datum Explicit displayed family; member types and common actual first-use exposed, not one atomic glyph. Blank macro filled with the glyph's TeX spelling. Family listing, not one atomic glyph. \\
\phantomsection\label{entry:symbol-registry-07466}\(T_{\alpha\beta}(G)\) & \texttt{\detokenize{T_{\alpha\beta}(G)}} & ground-represented initial-segment quotient & P6 proper iteration & object & \hyperref[def:p06-pil-tail]{\ref*{def:p06-pil-tail}} & \thepage & Conditions are old full-stage conditions with generic prefixes Blank macro filled with the glyph's TeX spelling. \\
\phantomsection\label{entry:symbol-registry-07467}\(\leq_s\) & \texttt{\detokenize{\leq_s}} & separative tail order and class map symbol family member \(\leq_s\) & P6 proper iteration & relation & \hyperref[def:p06-pil-tail]{\ref*{def:p06-pil-tail}} & \thepage & Generic-prefix identifications do not mean literal full-condition extension Explicit displayed family; member types and common actual first-use exposed, not one atomic glyph. Blank macro filled with the glyph's TeX spelling. Family listing, not one atomic glyph. \\
\phantomsection\label{entry:symbol-registry-07468}\(c_G\) & \texttt{\detokenize{c_G}} & separative tail order and class map symbol family member \(c_G\) & P6 proper iteration & relation & \hyperref[def:p06-pil-tail]{\ref*{def:p06-pil-tail}} & \thepage & Generic-prefix identifications do not mean literal full-condition extension Explicit displayed family; member types and common actual first-use exposed, not one atomic glyph. Blank macro filled with the glyph's TeX spelling. Family listing, not one atomic glyph. \\
\phantomsection\label{entry:symbol-registry-07469}\(\dot\rho\) & \texttt{\detokenize{\dot\rho}} & model-valued raw name and mixed coordinate name symbol family member \(\dot\rho\) & P6 proper iteration & object & \hyperref[thm:p06-pil-extension]{\ref*{thm:p06-pil-extension}} & \thepage & Name need not be in N; values are in N Explicit displayed family; member types and common actual first-use exposed, not one atomic glyph. Blank macro filled with the glyph's TeX spelling. Family listing, not one atomic glyph. \\
\phantomsection\label{entry:symbol-registry-07470}\(\widehat\rho_\xi\) & \texttt{\detokenize{\widehat\rho_\xi}} & model-valued raw name and mixed coordinate name symbol family member \(\widehat\rho_\xi\) & P6 proper iteration & object & \hyperref[thm:p06-pil-extension]{\ref*{thm:p06-pil-extension}} & \thepage & Name need not be in N; values are in N Explicit displayed family; member types and common actual first-use exposed, not one atomic glyph. Blank macro filled with the glyph's TeX spelling. Family listing, not one atomic glyph. \\
\phantomsection\label{entry:symbol-registry-07471}\(\delta=\sup(N\cap\beta)\) & \texttt{\detokenize{\delta=\sup(N\cap\beta)}} & model trace endpoint & P6 proper iteration & n/a & \hyperref[thm:p06-pil-extension]{\ref*{thm:p06-pil-extension}} & \thepage & May be below beta; final tail is unit beyond this endpoint Original precedence descriptor 'ordinal' retained; no rank inferred. Blank macro filled with the glyph's TeX spelling. \\
\phantomsection\label{entry:symbol-registry-07472}\(p_n\) & \texttt{\detokenize{p_n}} & coherent initial masters and descending model-condition names symbol family member \(p_n\) & P6 proper iteration & n/a & \hyperref[thm:p06-pil-extension]{\ref*{thm:p06-pil-extension}} & \thepage & No arbitrary proper lower bounds are taken Original precedence descriptor 'function' retained; no rank inferred. Explicit displayed family; member types and common actual first-use exposed, not one atomic glyph. Blank macro filled with the glyph's TeX spelling. Family listing, not one atomic glyph. \\
\phantomsection\label{entry:symbol-registry-07473}\(\rho_n\) & \texttt{\detokenize{\rho_n}} & coherent initial masters and descending model-condition names symbol family member \(\rho_n\) & P6 proper iteration & n/a & \hyperref[thm:p06-pil-extension]{\ref*{thm:p06-pil-extension}} & \thepage & No arbitrary proper lower bounds are taken Original precedence descriptor 'function' retained; no rank inferred. Explicit displayed family; member types and common actual first-use exposed, not one atomic glyph. Blank macro filled with the glyph's TeX spelling. Family listing, not one atomic glyph. \\
\phantomsection\label{entry:symbol-registry-07495}\(X=\omega\cup P\cup\mathcal P(P)\) & \texttt{\detokenize{X=\omega\cup P\cup\mathcal P(P)}} & small trace universe & P6 properness parameter invariance & object & \hyperref[def:p06-ppi-standing]{\ref*{def:p06-ppi-standing}} & \thepage & Includes actual dense subsets and all natural numbers; host construction not attributed as source-native glyph Blank macro filled with the glyph's TeX spelling. \\
\phantomsection\label{entry:symbol-registry-07496}\(\kappa\) & \texttt{\detokenize{\kappa}} & hereditary poset bound & P6 properness parameter invariance & n/a & \hyperref[def:p06-ppi-standing]{\ref*{def:p06-ppi-standing}} & \thepage & Maximum of aleph1 and the cardinality of the transitive closure of the poset tuple; exact existing book height convention Original precedence descriptor 'cardinality' retained; no rank inferred. Blank macro filled with the glyph's TeX spelling. \\
\phantomsection\label{entry:symbol-registry-07497}\(\chi\) & \texttt{\detokenize{\chi}} & tested and witness heights symbol family member \(\chi\) & P6 properness parameter invariance & n/a & \hyperref[def:p06-ppi-standing]{\ref*{def:p06-ppi-standing}} & \thepage & Both regular above 2\^{}kappa; Theta may contain a much larger parameter absent from Hchi Original precedence descriptor 'cardinality' retained; no rank inferred. Explicit displayed family; member types and common actual first-use exposed, not one atomic glyph. Blank macro filled with the glyph's TeX spelling. Family listing, not one atomic glyph. \\
\phantomsection\label{entry:symbol-registry-07498}\(\Theta\) & \texttt{\detokenize{\Theta}} & tested and witness heights symbol family member \(\Theta\) & P6 properness parameter invariance & n/a & \hyperref[def:p06-ppi-standing]{\ref*{def:p06-ppi-standing}} & \thepage & Both regular above 2\^{}kappa; Theta may contain a much larger parameter absent from Hchi Original precedence descriptor 'cardinality' retained; no rank inferred. Explicit displayed family; member types and common actual first-use exposed, not one atomic glyph. Blank macro filled with the glyph's TeX spelling. Family listing, not one atomic glyph. \\
\phantomsection\label{entry:symbol-registry-07499}\(\mathcal D_P\) & \texttt{\detokenize{\mathcal D_P}} & all dense subsets of the poset & P6 properness parameter invariance & object & \hyperref[def:p06-ppi-good-trace]{\ref*{def:p06-ppi-good-trace}} & \thepage & No selected enumeration substitutes for the full dense-set collection Blank macro filled with the glyph's TeX spelling. \\
\phantomsection\label{entry:symbol-registry-07500}\(\operatorname{Good}_P(Y)\) & \texttt{\detokenize{\operatorname{Good}_P(Y)}} & trace master-production predicate & P6 properness parameter invariance & n/a & \hyperref[def:p06-ppi-good-trace]{\ref*{def:p06-ppi-good-trace}} & \thepage & For each trace condition one master serves every trace dense set; master need not be in the trace Original precedence descriptor 'predicate' retained; no rank inferred. Blank macro filled with the glyph's TeX spelling. \\
\phantomsection\label{entry:symbol-registry-07501}\(g:[Z]^{<\omega}\to[Z]^{\leq\omega}\) & \texttt{\detokenize{g:[Z]^{<\omega}\to[Z]^{\leq\omega}}} & finite-input countable-output club code & P6 properness parameter invariance & n/a & \hyperref[thm:p06-ppi-club-codes]{\ref*{thm:p06-ppi-club-codes}} & \thepage & The input universe is infinite; individual codes are small even if their complete collection is not Original precedence descriptor 'function' retained; no rank inferred. Blank macro filled with the glyph's TeX spelling. \\
\phantomsection\label{entry:symbol-registry-07502}\(C_g\) & \texttt{\detokenize{C_g}} & club of countable code-closed sets & P6 properness parameter invariance & object & \hyperref[thm:p06-ppi-club-codes]{\ref*{thm:p06-ppi-club-codes}} & \thepage & Complete club coding construction supplied; no claim that arbitrary trace images are closed Blank macro filled with the glyph's TeX spelling. \\
\phantomsection\label{entry:symbol-registry-07503}\(h(S)\) & \texttt{\detokenize{h(S)}} & least closure under a countable-output code & P6 properness parameter invariance & n/a & \hyperref[thm:p06-ppi-trace-projection]{\ref*{thm:p06-ppi-trace-projection}} & \thepage & Finite character proved; this code hull is not assumed to be a Skolem hull until membership in the given elementary-model club is established Original precedence descriptor 'function' retained; no rank inferred. Blank macro filled with the glyph's TeX spelling. \\
\phantomsection\label{entry:symbol-registry-07504}\(g_N\) & \texttt{\detokenize{g_N}} & good code selected in the tested elementary model & P6 properness parameter invariance & n/a & \hyperref[thm:p06-ppi-all-model-recognition]{\ref*{thm:p06-ppi-all-model-recognition}} & \thepage & Elementarity selects a local witness; it need not be the original large-model witness code Original precedence descriptor 'function' retained; no rank inferred. Blank macro filled with the glyph's TeX spelling. \\
\phantomsection\label{entry:symbol-registry-07505}\(D_{\mathrm{at}}\) & \texttt{\detokenize{D_{\mathrm{at}}}} & dense set of atoms in the inheritance counterexample & P6 properness parameter invariance & object & \hyperref[thm:p06-ppi-no-downward-inheritance]{\ref*{thm:p06-ppi-no-downward-inheritance}} & \thepage & A larger-model master can fail the smaller-model test despite properness of the atomic poset Blank macro filled with the glyph's TeX spelling. \\
\phantomsection\label{entry:symbol-registry-07626}\(\|\mathbf A\|\) & \texttt{\detokenize{\|\mathbf A\|}} & least size of a relation-datum answer family & P6 Blass relational norms & n/a & \hyperref[def:p06-cc-relational-norm]{\ref*{def:p06-cc-relational-norm}} & \thepage & Source norm of a typed relation datum; not Euclidean norm or cardinality bars \\
\phantomsection\label{entry:symbol-registry-07627}\(\mathbf A\) & \texttt{\detokenize{\mathbf A}} & typed relation triples symbol family member \(\mathbf A\) & P6 Blass relational norms & n/a & \hyperref[def:p06-cc-relational-norm]{\ref*{def:p06-cc-relational-norm}} & \thepage & Family row for bold object names; B is declared in def:p06-cc-blass-morphism and C in def:p06-cc-sequential Explicit displayed family; member types and common actual first-use exposed, not one atomic glyph. Family listing, not one atomic glyph. Author macro \texttt{\detokenize{\mathbf A,}\allowbreak\detokenize{\mathbf B,}\allowbreak\detokenize{\mathbf C}} not carried to the members. \\
\phantomsection\label{entry:symbol-registry-07628}\(\mathbf B\) & \texttt{\detokenize{\mathbf B}} & typed relation triples symbol family member \(\mathbf B\) & P6 Blass relational norms & n/a & \hyperref[def:p06-cc-relational-norm]{\ref*{def:p06-cc-relational-norm}} & \thepage & Family row for bold object names; B is declared in def:p06-cc-blass-morphism and C in def:p06-cc-sequential Explicit displayed family; member types and common actual first-use exposed, not one atomic glyph. Family listing, not one atomic glyph. Author macro \texttt{\detokenize{\mathbf A,}\allowbreak\detokenize{\mathbf B,}\allowbreak\detokenize{\mathbf C}} not carried to the members. \\
\phantomsection\label{entry:symbol-registry-07629}\(\mathbf C\) & \texttt{\detokenize{\mathbf C}} & typed relation triples symbol family member \(\mathbf C\) & P6 Blass relational norms & n/a & \hyperref[def:p06-cc-relational-norm]{\ref*{def:p06-cc-relational-norm}} & \thepage & Family row for bold object names; B is declared in def:p06-cc-blass-morphism and C in def:p06-cc-sequential Explicit displayed family; member types and common actual first-use exposed, not one atomic glyph. Family listing, not one atomic glyph. Author macro \texttt{\detokenize{\mathbf A,}\allowbreak\detokenize{\mathbf B,}\allowbreak\detokenize{\mathbf C}} not carried to the members. \\
\phantomsection\label{entry:symbol-registry-07630}\(A_-\) & \texttt{\detokenize{A_-}} & negative question and positive answer domains symbol family member \(A_-\) & P6 Blass relational norms & n/a & \hyperref[def:p06-cc-relational-norm]{\ref*{def:p06-cc-relational-norm}} & \thepage & The analogous B and C domains have the same declared datum standing Explicit displayed family; member types and common actual first-use exposed, not one atomic glyph. Family listing, not one atomic glyph. Author macro \texttt{\detokenize{A\_-,}\allowbreak\detokenize{A\_+}} not carried to the members. \\
\phantomsection\label{entry:symbol-registry-07631}\(A_+\) & \texttt{\detokenize{A_+}} & negative question and positive answer domains symbol family member \(A_+\) & P6 Blass relational norms & n/a & \hyperref[def:p06-cc-relational-norm]{\ref*{def:p06-cc-relational-norm}} & \thepage & The analogous B and C domains have the same declared datum standing Explicit displayed family; member types and common actual first-use exposed, not one atomic glyph. Family listing, not one atomic glyph. Author macro \texttt{\detokenize{A\_-,}\allowbreak\detokenize{A\_+}} not carried to the members. \\
\phantomsection\label{entry:symbol-registry-07632}\(A\) & \texttt{\detokenize{A}} & binary relations in the corresponding typed triples symbol family member \(A\) & P6 Blass relational norms & n/a & \hyperref[def:p06-cc-relational-norm]{\ref*{def:p06-cc-relational-norm}} & \thepage & Italic binary-relation calls differ from bold datum names; family not three conflated object roles Explicit displayed family; member types and common actual first-use exposed, not one atomic glyph. Family listing, not one atomic glyph. Author macro \texttt{\detokenize{A,}\allowbreak\detokenize{B,}\allowbreak\detokenize{C}} not carried to the members. \\
\phantomsection\label{entry:symbol-registry-07633}\(B\) & \texttt{\detokenize{B}} & binary relations in the corresponding typed triples symbol family member \(B\) & P6 Blass relational norms & n/a & \hyperref[def:p06-cc-relational-norm]{\ref*{def:p06-cc-relational-norm}} & \thepage & Italic binary-relation calls differ from bold datum names; family not three conflated object roles Explicit displayed family; member types and common actual first-use exposed, not one atomic glyph. Family listing, not one atomic glyph. Author macro \texttt{\detokenize{A,}\allowbreak\detokenize{B,}\allowbreak\detokenize{C}} not carried to the members. \\
\phantomsection\label{entry:symbol-registry-07634}\(C\) & \texttt{\detokenize{C}} & binary relations in the corresponding typed triples symbol family member \(C\) & P6 Blass relational norms & n/a & \hyperref[def:p06-cc-relational-norm]{\ref*{def:p06-cc-relational-norm}} & \thepage & Italic binary-relation calls differ from bold datum names; family not three conflated object roles Explicit displayed family; member types and common actual first-use exposed, not one atomic glyph. Family listing, not one atomic glyph. Author macro \texttt{\detokenize{A,}\allowbreak\detokenize{B,}\allowbreak\detokenize{C}} not carried to the members. \\
\phantomsection\label{entry:symbol-registry-07635}\(\mathbf A^\perp\) & \texttt{\detokenize{\mathbf A^\perp}} & dual relation datum & P6 Blass relational norms & n/a & \hyperref[def:p06-cc-relational-norm]{\ref*{def:p06-cc-relational-norm}} & \thepage & Exchange the domains and complement the converse relation \\
\phantomsection\label{entry:symbol-registry-07636}\(\xi_-\) & \texttt{\detokenize{\xi_-}} & negative component from B questions to A questions & P6 Blass relational norms & n/a & \hyperref[def:p06-cc-blass-morphism]{\ref*{def:p06-cc-blass-morphism}} & \thepage & Source direction B\_minus to A\_minus \\
\phantomsection\label{entry:symbol-registry-07637}\(\xi_+\) & \texttt{\detokenize{\xi_+}} & positive component from A answers to B answers & P6 Blass relational norms & n/a & \hyperref[def:p06-cc-blass-morphism]{\ref*{def:p06-cc-blass-morphism}} & \thepage & Source direction A\_plus to B\_plus \\
\phantomsection\label{entry:symbol-registry-07638}\(\mathbf A;\mathbf B\) & \texttt{\detokenize{\mathbf A;\mathbf B}} & sequential composition of relation data & P6 Blass sequential composition & n/a & \hyperref[def:p06-cc-sequential]{\ref*{def:p06-cc-sequential}} & \thepage & The second task depends on the first answer via a full function space \\
\phantomsection\label{entry:symbol-registry-07639}\(R\) & \texttt{\detokenize{R}} & binary relation of the sequentially composed datum & P6 Blass sequential composition & n/a & \hyperref[def:p06-cc-sequential]{\ref*{def:p06-cc-sequential}} & \thepage & Local predicate on paired questions and paired answers; not a set of reals \\
\phantomsection\label{entry:symbol-registry-07640}\({}^{A_+}B_-\) & \texttt{\detokenize{{}^{A_+}B_-}} & all functions from the first answer domain to second question domain & P6 Blass sequential composition & n/a & \hyperref[def:p06-cc-sequential]{\ref*{def:p06-cc-sequential}} & \thepage & No unsupported standard Borel structure on the whole higher-type function space \\
\phantomsection\label{entry:symbol-registry-07641}\(\mathbf W\) & \texttt{\detokenize{\mathbf W}} & eventual-domination relation datum & P6 Blass relational norms & n/a & \hyperref[def:p06-cc-relational-norm]{\ref*{def:p06-cc-relational-norm}} & \thepage & Norm d and dual norm b; bold source object distinguished from italic relation \\
\phantomsection\label{entry:symbol-registry-07642}\(\mathcal P_\infty(\omega)\) & \texttt{\detokenize{\mathcal P_\infty(\omega)}} & infinite subsets of the natural numbers & P6 Blass relational norms & n/a & \hyperref[def:p06-cc-relational-norm]{\ref*{def:p06-cc-relational-norm}} & \thepage & Actual source glyph; explicitly equal to the standard family notation [omega]\^{}omega \\
\phantomsection\label{entry:symbol-registry-07643}\(\mu\) & \texttt{\detokenize{\mu}} & fair-coin product measure on Cantor space & P6 Cantor ideal and Cohen calculations & n/a & \hyperref[def:p06-ccsep-cantor]{\ref*{def:p06-ccsep-cantor}} & \thepage & Cylinder probability is 2 to the negative stem length \\
\phantomsection\label{entry:symbol-registry-07644}\([s]\) & \texttt{\detokenize{[s]}} & basic Cantor cylinder determined by a finite binary stem & P6 Cantor ideal and Cohen calculations & n/a & \hyperref[def:p06-ccsep-cantor]{\ref*{def:p06-ccsep-cantor}} & \thepage & Set of infinite binary extensions; not a quotient equivalence class \\
\phantomsection\label{entry:symbol-registry-07645}\(\mathcal B\) & \texttt{\detokenize{\mathcal B}} & meager ideal of arbitrary Cantor-space subsets & P6 Cantor ideal and Cohen calculations & n/a & \hyperref[def:p06-ccsep-cantor]{\ref*{def:p06-ccsep-cantor}} & \thepage & Blass 1994 source glyph; numerical values transferred to the real-line meager ideal by a proved bridge \\
\phantomsection\label{entry:symbol-registry-07646}\(\mathcal N\) & \texttt{\detokenize{\mathcal N}} & outer-null ideal of arbitrary Cantor-space subsets & P6 Cantor ideal and Cohen calculations & n/a & \hyperref[def:p06-ccsep-cantor]{\ref*{def:p06-ccsep-cantor}} & \thepage & Actual Bartoszynski source glyph with explicit Cantor domain; not literally the earlier real-line ideal; final value theorem explicitly transports back to the real-line domain \\
\phantomsection\label{entry:symbol-registry-07647}\(\pi\) & \texttt{\detokenize{\pi}} & binary-expansion map from Cantor space onto the unit interval & P6 Cantor real-space bridge & n/a & \hyperref[thm:p06-ccsep-domain-bridge]{\ref*{thm:p06-ccsep-domain-bridge}} & \thepage & Continuous surjection; measure images and preimages preserve outer nullity \\
\phantomsection\label{entry:symbol-registry-07648}\(h\) & \texttt{\detokenize{h}} & real-line to open-unit-interval measure-class homeomorphism & P6 Cantor real-space bridge & n/a & \hyperref[annex:p06.ccsep.domain-bridge]{\ref*{annex:p06.ccsep.domain-bridge}} & \thepage & Explicit rational formula; Lipschitz forward and locally Lipschitz inverse \\
\phantomsection\label{entry:symbol-registry-07649}\(\mathrm{CR}\) & \texttt{\detokenize{\mathrm{CR}}} & set of chopped reals & P6 Blass chopped category & n/a & \hyperref[def:p06-ccsep-chopped]{\ref*{def:p06-ccsep-chopped}} & \thepage & Pairs of binary real and consecutive nonempty finite-interval partition \\
\phantomsection\label{entry:symbol-registry-07650}\(\Pi\) & \texttt{\detokenize{\Pi}} & chopping partition of the natural numbers & P6 Blass chopped category & n/a & \hyperref[def:p06-ccsep-chopped]{\ref*{def:p06-ccsep-chopped}} & \thepage & Partition belongs to a chopped-real pair; not a forcing iteration or quantifier prefix \\
\phantomsection\label{entry:symbol-registry-07651}\(\operatorname{Match}(x,\Pi)\) & \texttt{\detokenize{\operatorname{Match}(x,}\allowbreak\detokenize{\Pi)}} & dense G-delta set of infinitely-often whole-block matches & P6 Blass chopped category & n/a & \hyperref[def:p06-ccsep-chopped]{\ref*{def:p06-ccsep-chopped}} & \thepage & Infinitely many matches; not eventual coordinate localization \\
\phantomsection\label{entry:symbol-registry-07652}\(\mathbf U\) & \texttt{\detokenize{\mathbf U}} & matching and engulfing relation data symbol family member \(\mathbf U\) & P6 Blass chopped category & n/a & \hyperref[def:p06-ccsep-category-data]{\ref*{def:p06-ccsep-category-data}} & \thepage & Bold object-family row; source U and V have explicitly different domains Explicit displayed family; member types and common actual first-use exposed, not one atomic glyph. Family listing, not one atomic glyph. Author macro \texttt{\detokenize{\mathbf U,}\allowbreak\detokenize{\mathbf V}} not carried to the members. \\
\phantomsection\label{entry:symbol-registry-07653}\(\mathbf V\) & \texttt{\detokenize{\mathbf V}} & matching and engulfing relation data symbol family member \(\mathbf V\) & P6 Blass chopped category & n/a & \hyperref[def:p06-ccsep-category-data]{\ref*{def:p06-ccsep-category-data}} & \thepage & Bold object-family row; source U and V have explicitly different domains Explicit displayed family; member types and common actual first-use exposed, not one atomic glyph. Family listing, not one atomic glyph. Author macro \texttt{\detokenize{\mathbf U,}\allowbreak\detokenize{\mathbf V}} not carried to the members. \\
\phantomsection\label{entry:symbol-registry-07654}\(U\) & \texttt{\detokenize{U}} & matching and engulfing binary predicates symbol family member \(U\) & P6 Blass chopped category & n/a & \hyperref[def:p06-ccsep-category-data]{\ref*{def:p06-ccsep-category-data}} & \thepage & Italic predicate-family row; distinct from the bold source triples and set-theoretic universe Explicit displayed family; member types and common actual first-use exposed, not one atomic glyph. Family listing, not one atomic glyph. Author macro \texttt{\detokenize{U,}\allowbreak\detokenize{V}} not carried to the members. \\
\phantomsection\label{entry:symbol-registry-07655}\(V\) & \texttt{\detokenize{V}} & matching and engulfing binary predicates symbol family member \(V\) & P6 Blass chopped category & n/a & \hyperref[def:p06-ccsep-category-data]{\ref*{def:p06-ccsep-category-data}} & \thepage & Italic predicate-family row; distinct from the bold source triples and set-theoretic universe Explicit displayed family; member types and common actual first-use exposed, not one atomic glyph. Family listing, not one atomic glyph. Author macro \texttt{\detokenize{U,}\allowbreak\detokenize{V}} not carried to the members. \\
\phantomsection\label{entry:symbol-registry-07656}\(\smile\) & \texttt{\detokenize{\smile}} & source superscript converse on the matching relation & P6 Blass chopped category & n/a & \hyperref[def:p06-ccsep-category-data]{\ref*{def:p06-ccsep-category-data}} & \thepage & The converse places the chopped real first and its matching binary real second \\
\phantomsection\label{entry:symbol-registry-07657}\(\beta\) & \texttt{\detokenize{\beta}} & answer-dependent search and reconstructed chopped answer symbol family member \(\beta\) & P6 Blass chopped category & n/a & \hyperref[annex:p06.ccsep.category-sequential]{\ref*{annex:p06.ccsep.category-sequential}} & \thepage & Beta is total with a zero branch on nonmatching inputs; its uncurried map is Borel; gamma is continuous Explicit displayed family; member types and common actual first-use exposed, not one atomic glyph. Family listing, not one atomic glyph. Author macro \texttt{\detokenize{\beta,}\allowbreak\detokenize{\gamma}} not carried to the members. \\
\phantomsection\label{entry:symbol-registry-07658}\(\gamma\) & \texttt{\detokenize{\gamma}} & answer-dependent search and reconstructed chopped answer symbol family member \(\gamma\) & P6 Blass chopped category & n/a & \hyperref[annex:p06.ccsep.category-sequential]{\ref*{annex:p06.ccsep.category-sequential}} & \thepage & Beta is total with a zero branch on nonmatching inputs; its uncurried map is Borel; gamma is continuous Explicit displayed family; member types and common actual first-use exposed, not one atomic glyph. Family listing, not one atomic glyph. Author macro \texttt{\detokenize{\beta,}\allowbreak\detokenize{\gamma}} not carried to the members. \\
\phantomsection\label{entry:symbol-registry-07659}\(\Theta_g\) & \texttt{\detokenize{\Theta_g}} & consecutive interval partition computed from an eventual bound & P6 Blass chopped category & n/a & \hyperref[annex:p06.ccsep.category-sequential]{\ref*{annex:p06.ccsep.category-sequential}} & \thepage & Endpoints recurse from zero with b\_next equal to one plus max of endpoint and g of endpoint \\
\phantomsection\label{entry:symbol-registry-07660}\(P_\kappa\) & \texttt{\detokenize{P_\kappa}} & binary finite-function Cohen forcing and its coordinate restriction symbol family member \(P_\kappa\) & P6 Cohen value calculations & n/a & \hyperref[def:p06-ccsep-ground]{\ref*{def:p06-ccsep-ground}} & \thepage & Reverse inclusion order; ground-model posets Explicit displayed family; member types and common actual first-use exposed, not one atomic glyph. Family listing, not one atomic glyph. Author macro \texttt{\detokenize{P\_\kappa,}\allowbreak\detokenize{P\_S}} not carried to the members. \\
\phantomsection\label{entry:symbol-registry-07661}\(P_S\) & \texttt{\detokenize{P_S}} & binary finite-function Cohen forcing and its coordinate restriction symbol family member \(P_S\) & P6 Cohen value calculations & n/a & \hyperref[def:p06-ccsep-ground]{\ref*{def:p06-ccsep-ground}} & \thepage & Reverse inclusion order; ground-model posets Explicit displayed family; member types and common actual first-use exposed, not one atomic glyph. Family listing, not one atomic glyph. Author macro \texttt{\detokenize{P\_\kappa,}\allowbreak\detokenize{P\_S}} not carried to the members. \\
\phantomsection\label{entry:symbol-registry-07662}\(G_S\) & \texttt{\detokenize{G_S}} & restricted generic filter on coordinate set S & P6 Cohen value calculations & n/a & \hyperref[def:p06-ccsep-ground]{\ref*{def:p06-ccsep-ground}} & \thepage & Restrictions of members of the supplied full generic; product factorization proves genericity \\
\phantomsection\label{entry:symbol-registry-07663}\(c_\alpha\) & \texttt{\detokenize{c_\alpha}} & binary coordinate Cohen real & P6 Cohen value calculations & n/a & \hyperref[def:p06-ccsep-ground]{\ref*{def:p06-ccsep-ground}} & \thepage & Generic union evaluated on the specified coordinate \\
\phantomsection\label{entry:symbol-registry-07664}\(f_c\) & \texttt{\detokenize{f_c}} & function enumerating positions of ones in a binary Cohen real & P6 Cohen value calculations & n/a & \hyperref[thm:p06-ccsep-fresh-obstructions]{\ref*{thm:p06-ccsep-fresh-obstructions}} & \thepage & Zero-based nth-one position; generic density establishes totality and non-domination \\
\phantomsection\label{entry:symbol-registry-07665}\(N_c\) & \texttt{\detokenize{N_c}} & Cohen-coded outer-null set of infinitely many matching blocks & P6 Cohen null calculation & n/a & \hyperref[def:p06-ccsep-cohen-null]{\ref*{def:p06-ccsep-cohen-null}} & \thepage & A subset of Cantor space; not a model or ideal \\
\phantomsection\label{entry:symbol-registry-07666}\(I_n\) & \texttt{\detokenize{I_n}} & nonempty block of length n plus one & P6 Cohen null calculation & n/a & \hyperref[def:p06-ccsep-cohen-null]{\ref*{def:p06-ccsep-cohen-null}} & \thepage & Fixed consecutive blocks support the explicit summable null cover \\
\phantomsection\label{entry:symbol-registry-07667}\(H_{q,n}\) & \texttt{\detokenize{H_{q,n}}} & ground dense open set constructed from a forcing condition and closed-set name & P6 countable forcing category preservation & n/a & \hyperref[annex:p06.ccsep.category-preservation]{\ref*{annex:p06.ccsep.category-preservation}} & \thepage & Union of cylinders forced to miss the nth named nowhere dense closed set \\
\phantomsection\label{entry:symbol-registry-07668}\(\mathbf S\) & \texttt{\detokenize{\mathbf S}} & local splitting relation datum & P6 Blass interval morphism & n/a & \hyperref[annex:p06.ccsep.morphism-duality]{\ref*{annex:p06.ccsep.morphism-duality}} & \thepage & Editorial proof-local name for the fully displayed source splitting triple; no new mathematical system \\
\end{longtable}
\endgroup
\Needspace{.4\textheight}
\subsection*{Stable part P7 scopes}
\begingroup\scriptsize\setlength{\tabcolsep}{3pt}
\begin{longtable}{@{}p{0.9in}p{1.3in}p{1.0in}p{0.5in}p{0.3in}p{0.4in}p{0.35in}p{1.65in}@{}}
\toprule
\textbf{Glyph} & \textbf{Macro} & \textbf{Role} & \textbf{Scope} & \textbf{Prec./form} & \textbf{First} & \textbf{Reg. page} & \textbf{Principal uses / notes} \\
\midrule
\endfirsthead
\toprule
\textbf{Glyph} & \textbf{Macro} & \textbf{Role} & \textbf{Scope} & \textbf{Prec./form} & \textbf{First} & \textbf{Reg. page} & \textbf{Principal uses / notes} \\
\midrule
\endhead
\midrule\multicolumn{8}{r}{\emph{continued}}\\\endfoot
\bottomrule\endlastfoot
\phantomsection\label{entry:symbol-registry-00020}\(\alpha,\beta,\gamma\) & \texttt{\detokenize{\alpha,}\allowbreak\detokenize{\beta,}\allowbreak\detokenize{\gamma}} & type metavariables & P7 normalization-curry-howard to lambda-calculi-and-type-theories & n/a & \hyperref[sec:normalization-curry-howard-01]{\ref*{sec:normalization-curry-howard-01}} & \thepage & Scope is disjoint from the P7.5 ordinal role and the P7.6 finite-type role. \\
\phantomsection\label{entry:symbol-registry-00394}\(\mathsf{LK}_{\mathrm d},\mathsf{LK}_{=}\) & \texttt{\detokenize{\mathsf{LK}_{\mathrm d},}\allowbreak\detokenize{\mathsf{LK}_{=}}} & defined-negation and equality sequent-calculus identifiers & P7 lk-lj to cut-elimination & n/a & \hyperref[sec:lk-lj-01]{\ref*{sec:lk-lj-01}} & \thepage & Exact registered presentations; neither identifier is an alias for canonical primitive-negation LK. \\
\phantomsection\label{entry:symbol-registry-00395}\(\mathsf{LJ}\) & \texttt{\detokenize{\mathsf{LJ}}} & intuitionistic sequent-calculus identifier & Parts 7–8 from lk-lj to intuitionistic-logic & n/a & \hyperref[sec:lk-lj-01]{\ref*{sec:lk-lj-01}} & \thepage & Exact single-succedent calculus fixed by sheet:lj; reused for IPL and equality-free IQC with every turnstile subscripted. \\
\phantomsection\label{entry:symbol-registry-00396}\(\mathsf L\) & \texttt{\detokenize{\mathsf L}} & associative Lambek-calculus identifier & Parts 7–8 from lk-lj to paraconsistent-and-nonmonotonic & n/a & \hyperref[sec:lk-lj-01]{\ref*{sec:lk-lj-01}} & \thepage & The cited associative ordered nonempty-antecedent Lambek calculus; \(\mathsf L^*\) and \(\mathsf L_1\) are source-named contrast extensions rather than aliases. \\
\phantomsection\label{entry:symbol-registry-00397}\(d(\varphi),d(\Gamma)\) & \texttt{\detokenize{d}} & defined-negation formula and context translation & P7 lk-lj & n/a & \hyperref[sec:lk-lj-01]{\ref*{sec:lk-lj-01}} & \thepage & Recursive equality-free translation from canonical LK syntax to the defined-negation language; context action is occurrencewise. \\
\phantomsection\label{entry:symbol-registry-00398}\(\bigwedge\Gamma\) & \texttt{\detokenize{\bigwedge}} & finite context conjunction & P7 lk-lj & 4 & \hyperref[sec:lk-lj-01]{\ref*{sec:lk-lj-01}} & \thepage & Metalinguistic fold of a finite context with empty value top; distinct from the P2.4 infinitary object-language connective. \\
\phantomsection\label{entry:symbol-registry-00399}\(\bigvee\Delta\) & \texttt{\detokenize{\bigvee}} & finite context disjunction & P7 lk-lj & 5 & \hyperref[sec:lk-lj-01]{\ref*{sec:lk-lj-01}} & \thepage & Metalinguistic fold of a finite context with empty value bottom; distinct from the P2.4 infinitary object-language connective. \\
\phantomsection\label{entry:symbol-registry-00400}\(D,D_1,D_2\) & \texttt{\detokenize{D,}\allowbreak\detokenize{D_1,}\allowbreak\detokenize{D_2}} & derivation metavariables & Legacy pre-batch4 P7.2-P7.6 reader meaning & n/a & \hyperref[def:p07-analytic-derivation]{\ref*{def:p07-analytic-derivation}} & \thepage & Finite derivation trees in the locally fixed calculus; distinct from the roman Dialectica superscript.  Numeric chapter scope migrated to its frozen pre-batch4 source-labelled reader regions; glyph, role and precedence unchanged. \\
\phantomsection\label{entry:symbol-registry-00401}\(\operatorname{rk}_{\mathrm{fm}}(\varphi)\) & \texttt{\detokenize{\operatorname{rk}_{\mathrm{fm}}}} & formula rank & P7 cut-elimination & n/a & \hyperref[def:p07-analytic-derivation]{\ref*{def:p07-analytic-derivation}} & \thepage & Number of logical symbols in the displayed formula; the subscript separates it from set-theoretic rank. \\
\phantomsection\label{entry:symbol-registry-00402}\(\operatorname{ht}(D)\) & \texttt{\detokenize{\operatorname{ht}}} & derivation height & P7 cut-elimination & n/a & \hyperref[def:p07-analytic-derivation]{\ref*{def:p07-analytic-derivation}} & \thepage & Greatest inference-step count on a leaf-to-root branch; distinct from tree-node height in P6.4. \\
\phantomsection\label{entry:symbol-registry-00403}\(\varphi^k\) & \texttt{\detokenize{{}^k}} & repeated Mix-formula occurrence & P7 lk-lj to cut-elimination & n/a & \hyperref[sec:lk-lj-01]{\ref*{sec:lk-lj-01}} & \thepage & Exactly k comma-separated occurrences on the registered Mix sheets; not exponentiation or iteration. \\
\phantomsection\label{entry:symbol-registry-00404}\(\mathsf{NJ}\) & \texttt{\detokenize{\mathsf{NJ}}} & intuitionistic natural-deduction calculus identifier & P7 lk-lj to normalization-curry-howard & n/a & \hyperref[sec:lk-lj-01]{\ref*{sec:lk-lj-01}} & \thepage & Exact calculus fixed by sheet:nj; excludes RAA, primitive negation, equality, and induction. \\
\phantomsection\label{entry:symbol-registry-00405}\(\longrightarrow_{\mathsf{NJ}}\) & \texttt{\detokenize{\longrightarrow_{\mathsf{NJ}}}} & one-step NJ derivation reduction & P7 normalization-curry-howard & meta & \hyperref[sec:normalization-curry-howard-01]{\ref*{sec:normalization-curry-howard-01}} & \thepage & Compatible closure of the exact principal reductions and permutation conversions fixed in the normalization entry. \\
\phantomsection\label{entry:symbol-registry-00406}\(\longrightarrow\) & \texttt{\detokenize{\longrightarrow}} & abstract reduction relation & P7 normalization-curry-howard & meta & \hyperref[sec:normalization-curry-howard-01]{\ref*{sec:normalization-curry-howard-01}} & \thepage & Locally declared reduction relation used only in the abstract normalization definitions; not implication or a sequent arrow. \\
\phantomsection\label{entry:symbol-registry-00407}\(\mathsf{Base},\mathsf b\) & \texttt{\detokenize{\mathsf{Base},}\allowbreak\detokenize{\mathsf b}} & base-type set and base-type metavariable & P7 normalization-curry-howard to lambda-calculi-and-type-theories & n/a & \hyperref[sec:normalization-curry-howard-01]{\ref*{sec:normalization-curry-howard-01}} & \thepage & Fixed generating set for the arrow-only STLC type grammar and one member of it. \\
\phantomsection\label{entry:symbol-registry-00408}\(\mathsf{STLC}\) & \texttt{\detokenize{\mathsf{STLC}}} & arrow-only simply typed lambda-calculus identifier & P7 normalization-curry-howard to lambda-calculi-and-type-theories & n/a & \hyperref[sec:normalization-curry-howard-01]{\ref*{sec:normalization-curry-howard-01}} & \thepage & Exact calculus fixed by sheet:stlc and used in every corresponding typing turnstile. \\
\phantomsection\label{entry:symbol-registry-00409}\(\Theta_{\mathrm{ty}}\) & \texttt{\detokenize{\Theta_{\mathrm{ty}}}} & term-variable typing context & Legacy pre-batch4 P7.3-P7.6 reader meaning & n/a & \hyperref[sec:normalization-curry-howard-01]{\ref*{sec:normalization-curry-howard-01}} & \thepage & Finite list of distinct term-variable declarations carrying well-formed types in the active typed calculus.  Numeric chapter scope migrated to its frozen pre-batch4 source-labelled reader regions; glyph, role and precedence unchanged. \\
\phantomsection\label{entry:symbol-registry-00410}\(M,N,P\) & \texttt{\detokenize{M,}\allowbreak\detokenize{N,}\allowbreak\detokenize{P}} & lambda-term metavariables & Legacy pre-batch4 P7.3-P7.6 reader meaning & n/a & \hyperref[sec:normalization-curry-howard-01]{\ref*{sec:normalization-curry-howard-01}} & \thepage & Terms of the locally fixed typed or untyped lambda grammar; the active calculus fixes the exact syntax.  Numeric chapter scope migrated to its frozen pre-batch4 source-labelled reader regions; glyph, role and precedence unchanged. \\
\phantomsection\label{entry:symbol-registry-00411}\(\Theta_{\mathrm{ty}}\derives_{\mathsf{STLC}}M{:}\alpha\) & \texttt{\detokenize{\derives_{\mathsf{STLC}}}} & STLC typing judgment & P7 normalization-curry-howard to lambda-calculi-and-type-theories & meta & \hyperref[sec:normalization-curry-howard-01]{\ref*{sec:normalization-curry-howard-01}} & \thepage & The exact arrow-only typing judgment of sheet:stlc; the context contains distinct term-variable declarations. \\
\phantomsection\label{entry:symbol-registry-00412}\(M[N/x]\) & \texttt{\detokenize{[N/}\allowbreak\detokenize{x]}} & capture-avoiding term substitution & Legacy pre-batch4 P7.3-P7.6 reader meaning & meta & \hyperref[sec:normalization-curry-howard-01]{\ref*{sec:normalization-curry-howard-01}} & \thepage & Substitutes term N for free occurrences of x in M after required alpha-renaming; distinct from formula substitution.  Numeric chapter scope migrated to its frozen pre-batch4 source-labelled reader regions; glyph, role and precedence unchanged. \\
\phantomsection\label{entry:symbol-registry-00413}\(\lambda x.M\) & \texttt{\detokenize{\lambda}} & untyped lambda abstraction & P7 lambda-calculi-and-type-theories & 3 & \hyperref[sec:p07-lambda-calculi-and-type-theories]{\ref*{sec:p07-lambda-calculi-and-type-theories}} & \thepage & Term former binding x in M; separate exact-scope role from typed abstraction. \\
\phantomsection\label{entry:symbol-registry-00414}\(M\longrightarrow_{\beta}N,\ M\longrightarrow_{\beta}^{*}N\) & \texttt{\detokenize{\longrightarrow_{\beta},}\allowbreak\detokenize{\longrightarrow_{\beta}^{*}}} & one-step beta reduction and finite reflexive-transitive closure & P7 normalization-curry-howard to lambda-calculi-and-type-theories & meta & \hyperref[sec:normalization-curry-howard-01]{\ref*{sec:normalization-curry-howard-01}} & \thepage & The active sheet fixes term beta reduction; the star denotes finite reflexive-transitive closure. \\
\phantomsection\label{entry:symbol-registry-00415}\(M=_{\beta}N\) & \texttt{\detokenize{=_{\beta}}} & beta conversion & P7 lambda-calculi-and-type-theories & meta & \hyperref[sec:p07-lambda-calculi-and-type-theories]{\ref*{sec:p07-lambda-calculi-and-type-theories}} & \thepage & Least congruence equivalence relation containing beta reduction; not elementary equivalence or formula equivalence. \\
\phantomsection\label{entry:symbol-registry-00416}\(M\mathrel{\mathsf{Par}}N\) & \texttt{\detokenize{\mathsf{Par}}} & parallel beta-reduction relation & P7 lambda-calculi-and-type-theories & meta & \hyperref[sec:p07-lambda-calculi-and-type-theories]{\ref*{sec:p07-lambda-calculi-and-type-theories}} & \thepage & Auxiliary relation on alpha-equivalence classes fixed by the four local clauses; its closure is compared with ordinary beta reduction. \\
\phantomsection\label{entry:symbol-registry-00417}\(\operatorname{dev}(M)\) & \texttt{\detokenize{\operatorname{dev}}} & complete development of a lambda term & P7 lambda-calculi-and-type-theories & n/a & \hyperref[sec:p07-lambda-calculi-and-type-theories]{\ref*{sec:p07-lambda-calculi-and-type-theories}} & \thepage & Metalinguistic recursion contracting all residuals exposed by the starting syntax; not a reduction relation. \\
\phantomsection\label{entry:symbol-registry-00418}\(\lambda f.\lambda x.f^n x\) & \texttt{\detokenize{\lambda f.\lambda x.f^n x}} & Church numeral encoding of n & P7 lambda-calculi-and-type-theories & n/a & \hyperref[sec:p07-lambda-calculi-and-type-theories]{\ref*{sec:p07-lambda-calculi-and-type-theories}} & \thepage & Closed untyped term used as the fixed numerical input-output encoding for lambda definability. \\
\phantomsection\label{entry:symbol-registry-00419}\(\mathsf S,\mathsf K,\mathsf{SK}\) & \texttt{\detokenize{\mathsf S,}\allowbreak\detokenize{\mathsf K,}\allowbreak\detokenize{\mathsf{SK}}} & S-K combinators and combinatory-calculus identifier & P7 lambda-calculi-and-type-theories & n/a & \hyperref[sec:p07-lambda-calculi-and-type-theories]{\ref*{sec:p07-lambda-calculi-and-type-theories}} & \thepage & Sheet:sk-combinatory fixes the two constants and exact calculus; scoped away from Post-clone and modal-system letters. \\
\phantomsection\label{entry:symbol-registry-00420}\(M\longrightarrow_{\mathsf{SK}}N\) & \texttt{\detokenize{\longrightarrow_{\mathsf{SK}}}} & one-step S-K contraction & P7 lambda-calculi-and-type-theories & meta & \hyperref[sec:p07-lambda-calculi-and-type-theories]{\ref*{sec:p07-lambda-calculi-and-type-theories}} & \thepage & Compatible contextual contraction fixed by sheet:sk-combinatory. \\
\phantomsection\label{entry:symbol-registry-00421}\(\Delta_{\mathrm{ty}}\) & \texttt{\detokenize{\Delta_{\mathrm{ty}}}} & System-F type-variable context & P7 lambda-calculi-and-type-theories & n/a & \hyperref[sec:p07-lambda-calculi-and-type-theories]{\ref*{sec:p07-lambda-calculi-and-type-theories}} & \thepage & Finite list of distinct type variables; distinct from formula and sequent contexts. \\
\phantomsection\label{entry:symbol-registry-00422}\(\mathsf F\) & \texttt{\detokenize{\mathsf F}} & Church-style System-F calculus identifier & P7 lambda-calculi-and-type-theories & n/a & \hyperref[sec:p07-lambda-calculi-and-type-theories]{\ref*{sec:p07-lambda-calculi-and-type-theories}} & \thepage & Exact identifier on sheet:system-f; its scope is disjoint from temporal eventually in P8.3. \\
\phantomsection\label{entry:symbol-registry-00423}\(\forall a.\alpha\) & \texttt{\detokenize{\forall}} & System-F universal type former & P7 lambda-calculi-and-type-theories & 3 & \hyperref[sec:p07-lambda-calculi-and-type-theories]{\ref*{sec:p07-lambda-calculi-and-type-theories}} & \thepage & Binds the type variable a in alpha; its syntactic category distinguishes it from the object-language formula quantifier. \\
\phantomsection\label{entry:symbol-registry-00424}\(\Lambda a.M\) & \texttt{\detokenize{\Lambda}} & System-F type abstraction & P7 lambda-calculi-and-type-theories & 3 & \hyperref[sec:p07-lambda-calculi-and-type-theories]{\ref*{sec:p07-lambda-calculi-and-type-theories}} & \thepage & Term former binding type variable a in M; distinct from ordinary term abstraction. \\
\phantomsection\label{entry:symbol-registry-00425}\(M[\alpha],\alpha[\beta/a],M[\alpha/a]\) & \texttt{\detokenize{[\alpha],}\allowbreak\detokenize{[\beta/}\allowbreak\detokenize{a],}\allowbreak\detokenize{[\alpha/}\allowbreak\detokenize{a]}} & System-F type application and capture-avoiding type substitution & P7 lambda-calculi-and-type-theories & meta & \hyperref[sec:p07-lambda-calculi-and-type-theories]{\ref*{sec:p07-lambda-calculi-and-type-theories}} & \thepage & The active syntactic category distinguishes term-level type application from substitution into a type or term. \\
\phantomsection\label{entry:symbol-registry-00426}\(\Delta_{\mathrm{ty}}\derives_{\mathsf F}\alpha\ \mathsf{type}\) & \texttt{\detokenize{\derives_{\mathsf F}}} & System-F type-formation judgment & P7 lambda-calculi-and-type-theories & meta & \hyperref[sec:p07-lambda-calculi-and-type-theories]{\ref*{sec:p07-lambda-calculi-and-type-theories}} & \thepage & Formation judgment for the least type set fixed by sheet:system-f. \\
\phantomsection\label{entry:symbol-registry-00427}\(\Delta_{\mathrm{ty}};\Theta_{\mathrm{ty}}\derives_{\mathsf F}M{:}\alpha\) & \texttt{\detokenize{\derives_{\mathsf F}}} & System-F term-typing judgment & P7 lambda-calculi-and-type-theories & meta & \hyperref[sec:p07-lambda-calculi-and-type-theories]{\ref*{sec:p07-lambda-calculi-and-type-theories}} & \thepage & Church-style typing judgment with separate type-variable and term-variable contexts. \\
\phantomsection\label{entry:symbol-registry-00428}\(\Gamma\;\mathsf{ctx}\) & \texttt{\detokenize{\mathsf{ctx}}} & MLTT context-formation judgment & P7 lambda-calculi-and-type-theories & meta & \hyperref[sec:p07-lambda-calculi-and-type-theories]{\ref*{sec:p07-lambda-calculi-and-type-theories}} & \thepage & Judgment of sheet:mltt-intensional asserting that the ordered dependent context Gamma is well formed. \\
\phantomsection\label{entry:symbol-registry-00429}\(\Gamma\vdash A\;\mathsf{type}\) & \texttt{\detokenize{\vdash,}\allowbreak\detokenize{\mathsf{type}}} & MLTT type-formation judgment & P7 lambda-calculi-and-type-theories & meta & \hyperref[sec:p07-lambda-calculi-and-type-theories]{\ref*{sec:p07-lambda-calculi-and-type-theories}} & \thepage & Type judgment in the selected intensional dependent calculus; distinct from System-F formation and formula derivability. \\
\phantomsection\label{entry:symbol-registry-00430}\(\Gamma\vdash a:A\) & \texttt{\detokenize{\vdash,}\allowbreak\detokenize{:}} & MLTT term-typing judgment & P7 lambda-calculi-and-type-theories & meta & \hyperref[sec:p07-lambda-calculi-and-type-theories]{\ref*{sec:p07-lambda-calculi-and-type-theories}} & \thepage & Dependent typing judgment on sheet:mltt-intensional; the colon is typed ascription rather than a labelled modal formula. \\
\phantomsection\label{entry:symbol-registry-00431}\(\Gamma\vdash A\equiv B\;\mathsf{type},\ \Gamma\vdash a\equiv b:A\) & \texttt{\detokenize{\equiv}} & MLTT judgmental type and term equality & P7 lambda-calculi-and-type-theories & meta & \hyperref[sec:p07-lambda-calculi-and-type-theories]{\ref*{sec:p07-lambda-calculi-and-type-theories}} & \thepage & Typed definitional equality generated by the displayed computations and congruence; equality reflection into identity types is absent. \\
\phantomsection\label{entry:symbol-registry-00432}\(\Pi_{x:A}B,\Sigma_{x:A}B\) & \texttt{\detokenize{\Pi,}\allowbreak\detokenize{\Sigma}} & MLTT dependent product and dependent sum type formers & P7 lambda-calculi-and-type-theories & 3 & \hyperref[sec:p07-lambda-calculi-and-type-theories]{\ref*{sec:p07-lambda-calculi-and-type-theories}} & \thepage & Bind x in B on sheet:mltt-intensional; distinct from analytical classes and sequent-context metavariables. \\
\phantomsection\label{entry:symbol-registry-00433}\(\operatorname{Id}_A(a,b)\) & \texttt{\detokenize{\operatorname{Id}}} & MLTT intensional identity type & P7 lambda-calculi-and-type-theories & n/a & \hyperref[sec:p07-lambda-calculi-and-type-theories]{\ref*{sec:p07-lambda-calculi-and-type-theories}} & \thepage & Proposition-as-type of identifications between terms a and b of A; no judgmental equality reflection is assumed. \\
\phantomsection\label{entry:symbol-registry-00434}\(\mathcal U_i\) & \texttt{\detokenize{\mathcal U_i}} & MLTT Russell universe at level i & P7 lambda-calculi-and-type-theories & n/a & \hyperref[sec:p07-lambda-calculi-and-type-theories]{\ref*{sec:p07-lambda-calculi-and-type-theories}} & \thepage & Cumulative predicative universe hierarchy on sheet:mltt-intensional; distinct from the P6 measurable-cardinal ultrafilter role of bare calligraphic U. \\
\phantomsection\label{entry:symbol-registry-00435}\(J_C(d,a,b,p)\) & \texttt{\detokenize{J_C}} & MLTT identity eliminator & P7 lambda-calculi-and-type-theories & n/a & \hyperref[sec:p07-lambda-calculi-and-type-theories]{\ref*{sec:p07-lambda-calculi-and-type-theories}} & \thepage & Typed eliminator for the displayed identity motive C with its judgmental refl computation rule. \\
\phantomsection\label{entry:symbol-registry-00436}\(\sigma,\tau:\Delta\to\Gamma\) & \texttt{\detokenize{\sigma,}\allowbreak\detokenize{\tau}} & MLTT substitutions with common domain and codomain & P7 lambda-calculi-and-type-theories & meta & \hyperref[sec:p07-lambda-calculi-and-type-theories]{\ref*{sec:p07-lambda-calculi-and-type-theories}} & \thepage & Capture-avoiding substitutions used by typed constructor congruence; equality premises and conclusions live in Delta. \\
\phantomsection\label{entry:symbol-registry-00437}\(\operatorname{Eq}_{\Delta}(J,J')\) & \texttt{\detokenize{\operatorname{Eq}_{\Delta}}} & multi-sorted MLTT equality premise & P7 lambda-calculi-and-type-theories & meta & \hyperref[sec:p07-lambda-calculi-and-type-theories]{\ref*{sec:p07-lambda-calculi-and-type-theories}} & \thepage & Eliminable metalinguistic abbreviation for type or term judgmental equality after transport to the common context Delta. \\
\phantomsection\label{entry:symbol-registry-00438}\(\mathsf 1_{\mathrm{ty}},\mathsf 0_{\mathrm{ty}}\) & \texttt{\detokenize{\mathsf 1_{\mathrm{ty}},}\allowbreak\detokenize{\mathsf 0_{\mathrm{ty}}}} & unit and empty types in the categorical STLC extension & P7 normalization-curry-howard & n/a & \hyperref[sec:normalization-curry-howard-01]{\ref*{sec:normalization-curry-howard-01}} & \thepage & Typed subscripts distinguish these formers from numerical constants and the computable Turing degree. \\
\phantomsection\label{entry:symbol-registry-00439}\(\alpha\mathbin{\times_{\mathrm{ty}}}\beta,\alpha\mathbin{+_{\mathrm{ty}}}\beta\) & \texttt{\detokenize{\times_{\mathrm{ty}},}\allowbreak\detokenize{+_{\mathrm{ty}}}} & product and sum type formers in the categorical STLC extension & P7 normalization-curry-howard & 3 & \hyperref[sec:normalization-curry-howard-01]{\ref*{sec:normalization-curry-howard-01}} & \thepage & Exact added type formers used to construct finite products and coproducts; distinct from arithmetic operations. \\
\phantomsection\label{entry:symbol-registry-00440}\(\star,\langle M,N\rangle,\pi_iM,\mathsf{inl}(M),\mathsf{inr}(M)\) & \texttt{\detokenize{\star,}\allowbreak\detokenize{\langle\ ,}\allowbreak\detokenize{\ \rangle,}\allowbreak\detokenize{\pi_i,}\allowbreak\detokenize{\mathsf{inl},}\allowbreak\detokenize{\mathsf{inr}}} & unit product and sum constructors for the categorical STLC extension & P7 normalization-curry-howard & n/a & \hyperref[sec:normalization-curry-howard-01]{\ref*{sec:normalization-curry-howard-01}} & \thepage & Term constructors governed by the exact typing table and beta-eta equations in the Curry-Howard entry. \\
\phantomsection\label{entry:symbol-registry-00441}\(\mathsf{case}(L;x.M;y.N),\mathsf{abort}_{\alpha}(M)\) & \texttt{\detokenize{\mathsf{case},}\allowbreak\detokenize{\mathsf{abort}_{\alpha}}} & sum and empty-type eliminators for the categorical STLC extension & P7 normalization-curry-howard & n/a & \hyperref[sec:normalization-curry-howard-01]{\ref*{sec:normalization-curry-howard-01}} & \thepage & Typed eliminators used only by the explicitly defined bicartesian extension. \\
\phantomsection\label{entry:symbol-registry-00442}\(r(A),h(c),o(D)\) & \texttt{\detokenize{r,}\allowbreak\detokenize{h,}\allowbreak\detokenize{o}} & formula rank Cut height and derivation ordinal measure & Legacy pre-batch4 P7.5 reader meaning & meta & \hyperref[sec:p07-arithmetic-proof-theory]{\ref*{sec:p07-arithmetic-proof-theory}} & \thepage & Metalinguistic measures for the finite-set propositional Cut strategy; o(D) is the Hessenberg natural sum displayed in its definition.  Numeric chapter scope migrated to its frozen pre-batch4 source-labelled reader regions; glyph, role and precedence unchanged. \\
\phantomsection\label{entry:symbol-registry-00443}\(\operatorname{Cut}_A(D_0,D_1),\operatorname{Cut}(D),A_c\) & \texttt{\detokenize{\operatorname{Cut}_A,}\allowbreak\detokenize{\operatorname{Cut},}\allowbreak\detokenize{A_c}} & Cut node Cut-node set and node-indexed cut formula & Legacy pre-batch4 P7.5 reader meaning & meta & \hyperref[sec:p07-arithmetic-proof-theory]{\ref*{sec:p07-arithmetic-proof-theory}} & \thepage & Metalinguistic notation for one Cut with premises D0 D1, the finite set of Cut nodes of D, and the formula attached to node c in the ordinal assignment.  Numeric chapter scope migrated to its frozen pre-batch4 source-labelled reader regions; glyph, role and precedence unchanged. \\
\phantomsection\label{entry:symbol-registry-00444}\(\mathop{\#}\) & \texttt{\detokenize{\#}} & Hessenberg natural sum & Legacy pre-batch4 P7.5 reader meaning & meta & \hyperref[sec:p07-arithmetic-proof-theory]{\ref*{sec:p07-arithmetic-proof-theory}} & \thepage & Commutative ordinal sum used only to combine the finitely many Cut-node contributions in o(D).  Numeric chapter scope migrated to its frozen pre-batch4 source-labelled reader regions; glyph, role and precedence unchanged. \\
\phantomsection\label{entry:symbol-registry-00446}\(\mathsf{HA}\) & \texttt{\detokenize{\mathsf{HA}}} & Heyting arithmetic & Legacy pre-batch4 P7.5-P7.6 reader meaning & n/a & \hyperref[sec:p07-arithmetic-proof-theory]{\ref*{sec:p07-arithmetic-proof-theory}} & \thepage & First-order arithmetic with full induction over intuitionistic logic; not a semantic notion of constructive truth.  Numeric chapter scope migrated to its frozen pre-batch4 source-labelled reader regions; glyph, role and precedence unchanged. \\
\phantomsection\label{entry:symbol-registry-00456}\(\alpha,\beta\) & \texttt{\detokenize{\alpha,}\allowbreak\detokenize{\beta}} & ordinal metavariables & Legacy pre-batch4 P7.5 reader meaning & n/a & \hyperref[sec:p07-arithmetic-proof-theory]{\ref*{sec:p07-arithmetic-proof-theory}} & \thepage & Ordinals in the local proof-theoretic calibration; disjoint from the P7.3-P7.4 and P7.6 type scopes.  Numeric chapter scope migrated to its frozen pre-batch4 source-labelled reader regions; glyph, role and precedence unchanged. \\
\phantomsection\label{entry:symbol-registry-00457}\(\gamma\) & \texttt{\detokenize{\gamma}} & ordinal metavariable & Legacy pre-batch4 P7.5 reader meaning & n/a & \hyperref[sec:p07-arithmetic-proof-theory]{\ref*{sec:p07-arithmetic-proof-theory}} & \thepage & Ordinal variable in the Veblen hierarchy and the definition of Gamma-zero; locally disjoint from formula and type metavariables.  Numeric chapter scope migrated to its frozen pre-batch4 source-labelled reader regions; glyph, role and precedence unchanged. \\
\phantomsection\label{entry:symbol-registry-00458}\(\omega^\omega\) & \texttt{\detokenize{\omega^\omega}} & ordinal exponentiation & Legacy pre-batch4 P7.5 reader meaning & n/a & \hyperref[sec:p07-arithmetic-proof-theory]{\ref*{sec:p07-arithmetic-proof-theory}} & \thepage & The ordinal power used in the PRA calibration; disjoint from the P6 Baire-space role.  Numeric chapter scope migrated to its frozen pre-batch4 source-labelled reader regions; glyph, role and precedence unchanged. \\
\phantomsection\label{entry:symbol-registry-00459}\(e\mapsto(D_e,\prec_e)\) & \texttt{\detokenize{\mapsto}} & effective notation-system enumeration & Legacy pre-batch4 P7.5 reader meaning & meta & \hyperref[sec:p07-arithmetic-proof-theory]{\ref*{sec:p07-arithmetic-proof-theory}} & \thepage & Maps each natural-number code to its recursive domain and ordering relation in the relative ordinal calibration.  Numeric chapter scope migrated to its frozen pre-batch4 source-labelled reader regions; glyph, role and precedence unchanged. \\
\phantomsection\label{entry:symbol-registry-00460}\(\mathcal C\) & \texttt{\detokenize{\mathcal C}} & formula class for transfinite induction & Legacy pre-batch4 P7.5 reader meaning & n/a & \hyperref[sec:p07-arithmetic-proof-theory]{\ref*{sec:p07-arithmetic-proof-theory}} & \thepage & Locally fixed family of formulas parameterizing the proof-theoretic calibration; disjoint from the P4 class scope.  Numeric chapter scope migrated to its frozen pre-batch4 source-labelled reader regions; glyph, role and precedence unchanged. \\
\phantomsection\label{entry:symbol-registry-00461}\(\Delta^0_1\) & \texttt{\detokenize{\Delta^0_1}} & Delta-zero-one arithmetic formula class & Legacy pre-batch4 P7.5 reader meaning & n/a & \hyperref[sec:p07-arithmetic-proof-theory]{\ref*{sec:p07-arithmetic-proof-theory}} & \thepage & Formulas both Sigma-zero-one and Pi-zero-one under the fixed arithmetic syntax; used in RCA-zero comprehension.  Numeric chapter scope migrated to its frozen pre-batch4 source-labelled reader regions; glyph, role and precedence unchanged. \\
\phantomsection\label{entry:symbol-registry-00462}\(\omega_n\) & \texttt{\detokenize{\omega_n}} & finite tower sequence below epsilon-zero & Legacy pre-batch4 P7.5 reader meaning & n/a & \hyperref[sec:p07-arithmetic-proof-theory]{\ref*{sec:p07-arithmetic-proof-theory}} & \thepage & Defined by omega-zero equals one and omega-(k+1) equals omega to the omega-k; used in the induction-fragment calibration.  Numeric chapter scope migrated to its frozen pre-batch4 source-labelled reader regions; glyph, role and precedence unchanged. \\
\phantomsection\label{entry:symbol-registry-00463}\(\Gamma_0\) & \texttt{\detokenize{\Gamma_0}} & Feferman--Schuette ordinal & Legacy pre-batch4 P7.5 reader meaning & n/a & \hyperref[sec:p07-arithmetic-proof-theory]{\ref*{sec:p07-arithmetic-proof-theory}} & \thepage & Least positive gamma satisfying Veb-gamma of zero equals gamma under the locally defined Veblen hierarchy.  Numeric chapter scope migrated to its frozen pre-batch4 source-labelled reader regions; glyph, role and precedence unchanged. \\
\phantomsection\label{entry:symbol-registry-00464}\(\operatorname{Veb}_{\alpha}(\beta)\) & \texttt{\detokenize{\operatorname{Veb}}} & Veblen hierarchy of normal ordinal functions & Legacy pre-batch4 P7.5 reader meaning & n/a & \hyperref[sec:p07-arithmetic-proof-theory]{\ref*{sec:p07-arithmetic-proof-theory}} & \thepage & Veb-zero is ordinal exponentiation by omega; later levels enumerate common fixed points of all earlier levels.  Numeric chapter scope migrated to its frozen pre-batch4 source-labelled reader regions; glyph, role and precedence unchanged. \\
\phantomsection\label{entry:symbol-registry-00465}\(\alpha,\beta\) & \texttt{\detokenize{\alpha,}\allowbreak\detokenize{\beta}} & finite-type metavariables & Legacy pre-batch4 P7.6 reader meaning & n/a & \hyperref[sec:p07-functional-interpretations]{\ref*{sec:p07-functional-interpretations}} & \thepage & Finite types generated from ground type zero by the registered type arrow; disjoint from P7.5 ordinals.  Numeric chapter scope migrated to its frozen pre-batch4 source-labelled reader regions; glyph, role and precedence unchanged. \\
\phantomsection\label{entry:symbol-registry-00466}\(\mathsf T_{\mathrm{tm}}\) & \texttt{\detokenize{\mathsf T_{\mathrm{tm}}}} & System-T term calculus identifier & Legacy pre-batch4 P7.6 reader meaning & n/a & \hyperref[sec:p07-functional-interpretations]{\ref*{sec:p07-functional-interpretations}} & \thepage & Exact finite-type term calculus fixed by sheet:system-t.  Numeric chapter scope migrated to its frozen pre-batch4 source-labelled reader regions; glyph, role and precedence unchanged. \\
\phantomsection\label{entry:symbol-registry-00467}\(\mathsf T_{\mathrm{qf}}\) & \texttt{\detokenize{\mathsf T_{\mathrm{qf}}}} & quantifier-free System-T target theory & Legacy pre-batch4 P7.6 reader meaning & n/a & \hyperref[sec:p07-functional-interpretations]{\ref*{sec:p07-functional-interpretations}} & \thepage & Separately named equality theory used as the Dialectica matrix target; not the term calculus alone.  Numeric chapter scope migrated to its frozen pre-batch4 source-labelled reader regions; glyph, role and precedence unchanged. \\
\phantomsection\label{entry:symbol-registry-00468}\(\Theta_{\mathrm{ty}}\derives_{\mathsf T_{\mathrm{tm}}}M{:}\alpha\) & \texttt{\detokenize{\derives_{\mathsf T_{\mathrm{tm}}}}} & System-T typing judgment & Legacy pre-batch4 P7.6 reader meaning & meta & \hyperref[sec:p07-functional-interpretations]{\ref*{sec:p07-functional-interpretations}} & \thepage & Typing judgment of sheet:system-t with the exact finite-type term calculus in the turnstile subscript.  Numeric chapter scope migrated to its frozen pre-batch4 source-labelled reader regions; glyph, role and precedence unchanged. \\
\phantomsection\label{entry:symbol-registry-00469}\(0\) & \texttt{\detokenize{0}} & System-T ground type and zero constant & Legacy pre-batch4 P7.6 reader meaning & n/a & \hyperref[sec:p07-functional-interpretations]{\ref*{sec:p07-functional-interpretations}} & \thepage & Syntactic category distinguishes the ground type from the term constant; both are fixed by sheet:system-t.  Numeric chapter scope migrated to its frozen pre-batch4 source-labelled reader regions; glyph, role and precedence unchanged. \\
\phantomsection\label{entry:symbol-registry-00470}\(S\) & \texttt{\detokenize{S}} & System-T successor constant & Legacy pre-batch4 P7.6 reader meaning & n/a & \hyperref[sec:p07-functional-interpretations]{\ref*{sec:p07-functional-interpretations}} & \thepage & Ground-type successor constant on sheet:system-t; scoped away from formula and Post-clone letters.  Numeric chapter scope migrated to its frozen pre-batch4 source-labelled reader regions; glyph, role and precedence unchanged. \\
\phantomsection\label{entry:symbol-registry-00471}\(R_\alpha\) & \texttt{\detokenize{R_\alpha}} & System-T finite-type recursor & Legacy pre-batch4 P7.6 reader meaning & n/a & \hyperref[sec:p07-functional-interpretations]{\ref*{sec:p07-functional-interpretations}} & \thepage & One primitive recursor at each finite type with argument order and computation equations fixed by sheet:system-t.  Numeric chapter scope migrated to its frozen pre-batch4 source-labelled reader regions; glyph, role and precedence unchanged. \\
\phantomsection\label{entry:symbol-registry-00472}\(M\longrightarrow_{\mathsf T_{\mathrm{tm}}}N\) & \texttt{\detokenize{\longrightarrow_{\mathsf T_{\mathrm{tm}}}}} & one-step System-T term reduction & Legacy pre-batch4 P7.6 reader meaning & meta & \hyperref[sec:p07-functional-interpretations]{\ref*{sec:p07-functional-interpretations}} & \thepage & Compatible reduction generated by beta contraction and the two registered recursor computation rules.  Numeric chapter scope migrated to its frozen pre-batch4 source-labelled reader regions; glyph, role and precedence unchanged. \\
\phantomsection\label{entry:symbol-registry-00473}\(\varphi^{\mathrm D}\) & \texttt{\detokenize{{}^{\mathrm D}}} & Dialectica translation of a formula & Legacy pre-batch4 P7.6 reader meaning & n/a & \hyperref[sec:p07-functional-interpretations]{\ref*{sec:p07-functional-interpretations}} & \thepage & Translation into an existential witness block, universal challenge block, and quantifier-free target matrix.  Numeric chapter scope migrated to its frozen pre-batch4 source-labelled reader regions; glyph, role and precedence unchanged. \\
\phantomsection\label{entry:symbol-registry-00474}\(\varphi_{\mathrm D}(\vec x,\vec y)\) & \texttt{\detokenize{{}_{\mathrm D}}} & Dialectica quantifier-free matrix & Legacy pre-batch4 P7.6 reader meaning & n/a & \hyperref[sec:p07-functional-interpretations]{\ref*{sec:p07-functional-interpretations}} & \thepage & Matrix associated with the locally fixed Dialectica translation; typed witness and challenge tuples remain explicit.  Numeric chapter scope migrated to its frozen pre-batch4 source-labelled reader regions; glyph, role and precedence unchanged. \\
\phantomsection\label{entry:symbol-registry-00475}\(\vec x,\vec y,\vec u,\vec v\) & \texttt{\detokenize{\vec x,}\allowbreak\detokenize{\vec y,}\allowbreak\detokenize{\vec u,}\allowbreak\detokenize{\vec v}} & Dialectica witness and challenge tuples & Legacy pre-batch4 P7.6 reader meaning & n/a & \hyperref[sec:p07-functional-interpretations]{\ref*{sec:p07-functional-interpretations}} & \thepage & Finite tuples of finite-type variables; distinct from the P4-P5 tuples of numerical variables.  Numeric chapter scope migrated to its frozen pre-batch4 source-labelled reader regions; glyph, role and precedence unchanged. \\
\phantomsection\label{entry:symbol-registry-00476}\(\vec f,\vec g,\vec t\) & \texttt{\detokenize{\vec f,}\allowbreak\detokenize{\vec g,}\allowbreak\detokenize{\vec t}} & finite tuples of System-T functionals and terms & Legacy pre-batch4 P7.6 reader meaning & n/a & \hyperref[sec:p07-functional-interpretations]{\ref*{sec:p07-functional-interpretations}} & \thepage & Typed functional witnesses in translation clauses and the extracted closed-term tuple in Dialectica soundness.  Numeric chapter scope migrated to its frozen pre-batch4 source-labelled reader regions; glyph, role and precedence unchanged. \\
\phantomsection\label{entry:symbol-registry-00477}\(\pi_0,\pi_1\) & \texttt{\detokenize{\pi_0,}\allowbreak\detokenize{\pi_1}} & primitive-recursive pairing projections & Legacy pre-batch4 P7.6 reader meaning & n/a & \hyperref[sec:p07-functional-interpretations]{\ref*{sec:p07-functional-interpretations}} & \thepage & First and second projections for the fixed numerical pairing function used in number realizability.  Numeric chapter scope migrated to its frozen pre-batch4 source-labelled reader regions; glyph, role and precedence unchanged. \\
\phantomsection\label{entry:symbol-registry-00478}\(\mathsf{ECT}_0\) & \texttt{\detokenize{\mathsf{ECT}_0}} & extended Church-thesis scheme & Legacy pre-batch4 P7.6 reader meaning & n/a & \hyperref[sec:p07-functional-interpretations]{\ref*{sec:p07-functional-interpretations}} & \thepage & Universal closure of the displayed almost-negative-domain scheme; not an unrestricted Church-thesis axiom.  Numeric chapter scope migrated to its frozen pre-batch4 source-labelled reader regions; glyph, role and precedence unchanged. \\
\phantomsection\label{entry:symbol-registry-00479}\(\mathsf{PC}_{\varepsilon}\) & \texttt{\detokenize{\mathsf{PC}_{\varepsilon}}} & equality-free epsilon-calculus identifier & Legacy pre-batch4 P7.7 reader meaning & n/a & \hyperref[sec:p07-structural-extensions]{\ref*{sec:p07-structural-extensions}} & \thepage & Exact Hilbert-FOL extension fixed by sheet:epsilon; primitive quantifiers are retained and equality is excluded.  Numeric chapter scope migrated to its frozen pre-batch4 source-labelled reader regions; glyph, role and precedence unchanged. \\
\phantomsection\label{entry:symbol-registry-00480}\(\mathsf{EC}\) & \texttt{\detokenize{\mathsf{EC}}} & epsilon-free quantifier-free fragment & Legacy pre-batch4 P7.7 reader meaning & n/a & \hyperref[sec:p07-structural-extensions]{\ref*{sec:p07-structural-extensions}} & \thepage & Eliminable metalinguistic abbreviation used in the first epsilon theorem; not a second epsilon system.  Numeric chapter scope migrated to its frozen pre-batch4 source-labelled reader regions; glyph, role and precedence unchanged. \\
\phantomsection\label{entry:symbol-registry-00481}\(\varepsilon x\,\varphi(x)\) & \texttt{\detokenize{\varepsilon}} & Hilbert epsilon term former & Legacy pre-batch4 P7.7 reader meaning & 3 & \hyperref[sec:p07-structural-extensions]{\ref*{sec:p07-structural-extensions}} & \thepage & Capture-avoiding term binder on the equality-free epsilon sheet; distinct from epsilon-zero.  Numeric chapter scope migrated to its frozen pre-batch4 source-labelled reader regions; glyph, role and precedence unchanged. \\
\phantomsection\label{entry:symbol-registry-00482}\(w{:}\varphi\) & \texttt{\detokenize{:}} & world-labelled formula & Legacy pre-batch4 P7.7 reader meaning & meta & \hyperref[sec:p07-structural-extensions]{\ref*{sec:p07-structural-extensions}} & \thepage & Labelled-sequent syntax only; distinct from type ascription and not a semantic satisfaction assertion.  Numeric chapter scope migrated to its frozen pre-batch4 source-labelled reader regions; glyph, role and precedence unchanged. \\
\phantomsection\label{entry:symbol-registry-00483}\(wRv\) & \texttt{\detokenize{R}} & relational atom in labelled-sequent syntax & Legacy pre-batch4 P7.7 reader meaning & n/a & \hyperref[sec:p07-structural-extensions]{\ref*{sec:p07-structural-extensions}} & \thepage & Syntactic frame-relation atom inside a labelled derivation language; distinct from semantic accessibility in P8.  Numeric chapter scope migrated to its frozen pre-batch4 source-labelled reader regions; glyph, role and precedence unchanged. \\
\phantomsection\label{entry:symbol-registry-00484}\(\mathsf P_{\mathrm{pf}},\mathsf Q_{\mathrm{pf}}\) & \texttt{\detokenize{\mathsf P_{\mathrm{pf}},}\allowbreak\detokenize{\mathsf Q_{\mathrm{pf}}}} & Cook-Reckhow proof-system metavariables & Legacy pre-batch4 P7.8 reader meaning & n/a & \hyperref[sec:p07-multiplicative-proof-nets]{\ref*{sec:p07-multiplicative-proof-nets}} & \thepage & Polynomial-time encoded surjections onto TAUT; subscripts prevent collision with complexity and formula letters.  Numeric chapter scope migrated to its frozen pre-batch4 source-labelled reader regions; glyph, role and precedence unchanged. \\
\phantomsection\label{entry:symbol-registry-00485}\(\operatorname{Prf}_{\mathsf P_{\mathrm{pf}}}(\pi,\varphi)\) & \texttt{\detokenize{\operatorname{Prf}_{\mathsf P_{\mathrm{pf}}}}} & encoded propositional proof relation & Legacy pre-batch4 P7.8 reader meaning & meta & \hyperref[sec:p07-multiplicative-proof-nets]{\ref*{sec:p07-multiplicative-proof-nets}} & \thepage & Polynomial-time decidable relation whose range is exactly TAUT under the fixed malformed-input convention.  Numeric chapter scope migrated to its frozen pre-batch4 source-labelled reader regions; glyph, role and precedence unchanged. \\
\phantomsection\label{entry:symbol-registry-00486}\(\pi\) & \texttt{\detokenize{\pi}} & encoded propositional proof metavariable & Legacy pre-batch4 P7.8 reader meaning & n/a & \hyperref[sec:p07-multiplicative-proof-nets]{\ref*{sec:p07-multiplicative-proof-nets}} & \thepage & Binary proof string for the fixed Cook-Reckhow encoding; no standing as a projection or permutation.  Numeric chapter scope migrated to its frozen pre-batch4 source-labelled reader regions; glyph, role and precedence unchanged. \\
\phantomsection\label{entry:symbol-registry-00487}\(\operatorname{size}(a)\) & \texttt{\detokenize{\operatorname{size}}} & bit length of an encoded proof or formula object & Legacy pre-batch4 P7.8 reader meaning & n/a & \hyperref[sec:p07-multiplicative-proof-nets]{\ref*{sec:p07-multiplicative-proof-nets}} & \thepage & Defined from the fixed binary encoding; deliberately not vertical-bar cardinality notation.  Numeric chapter scope migrated to its frozen pre-batch4 source-labelled reader regions; glyph, role and precedence unchanged. \\
\phantomsection\label{entry:symbol-registry-00488}\(\mathsf{PHP}^{n}_{n+1},x_{i,j}\) & \texttt{\detokenize{\mathsf{PHP}^{n}_{n+1},}\allowbreak\detokenize{x_{i,j}}} & pigeonhole-principle CNF family and its variables & Legacy pre-batch4 P7.8 reader meaning & n/a & \hyperref[sec:p07-multiplicative-proof-nets]{\ref*{sec:p07-multiplicative-proof-nets}} & \thepage & Exact weak PHP encoding with n+1 pigeons and n holes; no clauses forbid one pigeon from occupying two holes.  Numeric chapter scope migrated to its frozen pre-batch4 source-labelled reader regions; glyph, role and precedence unchanged. \\
\phantomsection\label{entry:symbol-registry-00489}\(\longrightarrow\) & \texttt{\detokenize{\longrightarrow}} & proof-system function arrow & Legacy pre-batch4 P7.8 reader meaning & meta & \hyperref[sec:p07-multiplicative-proof-nets]{\ref*{sec:p07-multiplicative-proof-nets}} & \thepage & Metalinguistic total-function arrow in the Cook-Reckhow definition; disjoint from the P7.3 abstract-reduction scope.  Numeric chapter scope migrated to its frozen pre-batch4 source-labelled reader regions; glyph, role and precedence unchanged. \\
\phantomsection\label{entry:symbol-registry-02119}\(\mathsf{Mar}\) & \texttt{\detokenize{\mathsf{Mar}}} & set of variance markings & P7 monotonicity-inference & n/a & \hyperref[def:p07-monotonicity-types]{\ref*{def:p07-monotonicity-types}} & \thepage & Source-local + - and dot marking roles. \\
\phantomsection\label{entry:symbol-registry-02120}\(+,-,\cdot\) & \texttt{\detokenize{+,}\allowbreak\detokenize{-,}\allowbreak\detokenize{\cdot}} & monotone antitone and unspecified markings & P7 monotonicity-inference & n/a & \hyperref[def:p07-monotonicity-types]{\ref*{def:p07-monotonicity-types}} & \thepage & Not truth values; dot does not assert nonmonotonicity. \\
\phantomsection\label{entry:symbol-registry-02121}\(\sqsubseteq\) & \texttt{\detokenize{\sqsubseteq}} & refinement order on markings & P7 monotonicity-inference & n/a & \hyperref[def:p07-monotonicity-types]{\ref*{def:p07-monotonicity-types}} & \thepage & Both signs refine dot. \\
\phantomsection\label{entry:symbol-registry-02122}\(\circ\) & \texttt{\detokenize{\circ}} & composition of occurrence polarities & P7 monotonicity-inference & n/a & \hyperref[def:p07-monotonicity-types]{\ref*{def:p07-monotonicity-types}} & \thepage & Also ordinary function composition when arguments are functions; both roles explicitly typed. \\
\phantomsection\label{entry:symbol-registry-02123}\(\sigma\xrightarrow m\tau\) & \texttt{\detokenize{\xrightarrow}} & marked function type constructor & P7 monotonicity-inference & n/a & \hyperref[def:p07-monotonicity-types]{\ref*{def:p07-monotonicity-types}} & \thepage & Greek letters are source-local type metavariables; m is a variance marking. \\
\phantomsection\label{entry:symbol-registry-02124}\(\preceq\) & \texttt{\detokenize{\preceq}} & admissible type coercion order & P7 monotonicity-inference & n/a & \hyperref[def:p07-monotonicity-types]{\ref*{def:p07-monotonicity-types}} & \thepage & Contravariant in arrow arguments and covariant in results and markings. \\
\phantomsection\label{entry:symbol-registry-02125}\(\uparrow\) & \texttt{\detokenize{\uparrow}} & partial same-argument compatibility relation on types & P7 monotonicity-inference & n/a & \hyperref[def:p07-monotonicity-types]{\ref*{def:p07-monotonicity-types}} & \thepage & Not divergence of a computation. \\
\phantomsection\label{entry:symbol-registry-02126}\(\vee\) & \texttt{\detokenize{\vee}} & join of markings or partial compatible type join & P7 monotonicity-inference & n/a & \hyperref[def:p07-monotonicity-types]{\ref*{def:p07-monotonicity-types}} & \thepage & Operands distinguish this local role from object-language disjunction. \\
\phantomsection\label{entry:symbol-registry-02127}\(\widehat\sigma\) & \texttt{\detokenize{\widehat}} & erasure of markings along the result spine & P7 monotonicity-inference & n/a & \hyperref[def:p07-monotonicity-types]{\ref*{def:p07-monotonicity-types}} & \thepage & Does not erase argument-type markings. \\
\phantomsection\label{entry:symbol-registry-02128}\(D_\tau\) & \texttt{\detokenize{D_\tau}} & semantic domain at a marked type & P7 monotonicity-inference & n/a & \hyperref[def:p07-monotonicity-prestructure]{\ref*{def:p07-monotonicity-prestructure}} & \thepage & A preordered set; arrow domains consist of actual functions. \\
\phantomsection\label{entry:symbol-registry-02129}\(\leq_\tau\) & \texttt{\detokenize{\leq}} & semantic preorder at type tau & P7 monotonicity-inference & n/a & \hyperref[def:p07-monotonicity-prestructure]{\ref*{def:p07-monotonicity-prestructure}} & \thepage & Pointwise at arrow types. \\
\phantomsection\label{entry:symbol-registry-02130}\(\pi_{\sigma,\tau}\) & \texttt{\detokenize{\pi}} & coherent semantic coercion & P7 monotonicity-inference & n/a & \hyperref[def:p07-monotonicity-prestructure]{\ref*{def:p07-monotonicity-prestructure}} & \thepage & Identity composition order preservation and application coherence are explicit. \\
\phantomsection\label{entry:symbol-registry-02131}\(\operatorname{tp}(M)\) & \texttt{\detokenize{\operatorname{tp}}} & intrinsic marked type of a term & P7 monotonicity-inference & n/a & \hyperref[def:p07-monotonicity-terms]{\ref*{def:p07-monotonicity-terms}} & \thepage & No implicit term subsumption rule. \\
\phantomsection\label{entry:symbol-registry-02132}\(\sem M_\varphi\) & \texttt{\detokenize{\sem}} & interpretation of a term at an assignment & P7 monotonicity-inference & n/a & \hyperref[def:p07-monotonicity-structures]{\ref*{def:p07-monotonicity-structures}} & \thepage & Promote to comparison types explicitly through pi. \\
\phantomsection\label{entry:symbol-registry-02133}\(\varphi_x^a\) & \texttt{\detokenize{\varphi}} & assignment with one variable replaced & P7 monotonicity-inference & n/a & \hyperref[def:p07-monotonicity-structures]{\ref*{def:p07-monotonicity-structures}} & \thepage & a belongs to the declared type of x. \\
\phantomsection\label{entry:symbol-registry-02134}\(M[s]\) & \texttt{\detokenize{M[s]}} & simultaneous capture-avoiding typed substitution & P7 monotonicity-inference & n/a & \hyperref[def:p07-monotonicity-substitution]{\ref*{def:p07-monotonicity-substitution}} & \thepage & Images may have smaller intrinsic types; binders retain their types. \\
\phantomsection\label{entry:symbol-registry-02135}\(\approx\) & \texttt{\detokenize{\approx}} & global alpha-beta-eta conversion with type-changing application & P7 monotonicity-inference & n/a & \hyperref[def:p07-monotonicity-substitution]{\ref*{def:p07-monotonicity-substitution}} & \thepage & Conversion coherence is proved by long elaboration at a fixed common comparison type. \\
\phantomsection\label{entry:symbol-registry-02136}\(\precsim_\tau\) & \texttt{\detokenize{\precsim}} & object assertion comparing terms at type tau & P7 monotonicity-inference & n/a & \hyperref[def:p07-monotonicity-consequence]{\ref*{def:p07-monotonicity-consequence}} & \thepage & Both intrinsic types lie below tau. \\
\phantomsection\label{entry:symbol-registry-02137}\(\models_{\mathrm{mon}}\) & \texttt{\detokenize{\models}} & consequence with universally interpreted open premises & P7 monotonicity-inference & n/a & \hyperref[def:p07-monotonicity-consequence]{\ref*{def:p07-monotonicity-consequence}} & \thepage & Not consequence at a single fixed assignment. \\
\phantomsection\label{entry:symbol-registry-02138}\(\vdash_{\mathrm{mon}}\) & \texttt{\detokenize{\vdash}} & finite derivation in the displayed monotonicity calculus & P7 monotonicity-inference & n/a & \hyperref[def:p07-monotonicity-consequence]{\ref*{def:p07-monotonicity-consequence}} & \thepage & Rule sheet p07-monotonicity; universal premises and type-changing comparisons. \\
\phantomsection\label{entry:symbol-registry-02139}\(\langle M\rangle\) & \texttt{\detokenize{\langle}} & term conversion class & P7 monotonicity-inference & n/a & \hyperref[def:p07-monotonicity-term-structure]{\ref*{def:p07-monotonicity-term-structure}} & \thepage & Global conversion class; its functional realization retains the comparison-type index as proved in thm:p07-monotonicity-term-model. \\
\phantomsection\label{entry:symbol-registry-02140}\(T_\tau\) & \texttt{\detokenize{T_\tau}} & classes with an intrinsic representative below tau & P7 monotonicity-inference & n/a & \hyperref[def:p07-monotonicity-term-structure]{\ref*{def:p07-monotonicity-term-structure}} & \thepage & The derivability preorder on these domains yields the indexed canonical model used in thm:p07-monotonicity-completeness. \\
\phantomsection\label{entry:symbol-registry-02141}\(\sqsubseteq_\tau\) & \texttt{\detokenize{\sqsubseteq}} & term-structure preorder & P7 monotonicity-inference & n/a & \hyperref[def:p07-monotonicity-term-structure]{\ref*{def:p07-monotonicity-term-structure}} & \thepage & Typed term domain distinguishes this role from marking refinement. \\
\phantomsection\label{entry:symbol-registry-02142}\(i_{\sigma,\tau}\) & \texttt{\detokenize{i}} & inclusion of term-class domains & P7 monotonicity-inference & n/a & \hyperref[def:p07-monotonicity-term-structure]{\ref*{def:p07-monotonicity-term-structure}} & \thepage & Order preservation is proved for every stated term structure and used in the completed canonical and occurrence-model constructions. \\
\phantomsection\label{entry:symbol-registry-02143}\(L_\tau(M)\) & \texttt{\detokenize{L_\tau}} & comparison-directed beta-normal eta-long elaboration & P7 monotonicity-inference & n/a & \hyperref[def:p07-monotonicity-long-form]{\ref*{def:p07-monotonicity-long-form}} & \thepage & Expected type tau and receiving argument types determine marked decoration. \\
\phantomsection\label{entry:symbol-registry-02144}\(\mathsf{Par}_{\beta\eta}\) & \texttt{\detokenize{\mathsf{Par}}} & ordinary unmarked parallel beta-eta reduction & P7 monotonicity-inference & n/a & \hyperref[thm:p07-monotonicity-ordinary-long]{\ref*{thm:p07-monotonicity-ordinary-long}} & \thepage & Auxiliary proof relation; not a claim of raw marked confluence. \\
\phantomsection\label{entry:symbol-registry-02145}\(\iota_\tau\) & \texttt{\detokenize{\iota}} & order isomorphism realizing a term class at comparison type tau & P7 monotonicity-inference & n/a & \hyperref[thm:p07-monotonicity-term-model]{\ref*{thm:p07-monotonicity-term-model}} & \thepage & The index must not be dropped when substituted representatives refine types. \\
\phantomsection\label{entry:symbol-registry-02146}\(F_a^{\sigma,\tau}\) & \texttt{\detokenize{F}} & actual function realizing a class at an argument and result type & P7 monotonicity-inference & n/a & \hyperref[thm:p07-monotonicity-term-model]{\ref*{thm:p07-monotonicity-term-model}} & \thepage & Same domain/codomain function across permitted variance markings. \\
\phantomsection\label{entry:symbol-registry-02147}\(\leq_{\mathrm{col}}\) & \texttt{\detokenize{\leq}} & signed pointwise order of distinguished free-variable colors & P7 monotonicity-inference & n/a & \hyperref[thm:p07-monotonicity-lyndon]{\ref*{thm:p07-monotonicity-lyndon}} & \thepage & Used only on equal color-erased typed trees in the occurrence-model proof. \\
\phantomsection\label{entry:symbol-registry-02148}\(\Theta,\Delta\) & \texttt{\detokenize{\Theta,}\allowbreak\detokenize{\Delta}} & finite source and target typing contexts & \hyperref[sec:p07-stlc-metatheory]{\ref*{sec:p07-stlc-metatheory}} & n/a & \hyperref[def:p07-stlc-meta-setting]{\ref*{def:p07-stlc-meta-setting}} & \thepage & Each context assigns one simple type to each distinct declared variable. \\
\phantomsection\label{entry:symbol-registry-02149}\(\rho\) & \texttt{\detokenize{\rho}} & simultaneous substitution of raw terms for free variables & \hyperref[sec:p07-stlc-metatheory]{\ref*{sec:p07-stlc-metatheory}} & n/a & \hyperref[def:p07-stlc-meta-setting]{\ref*{def:p07-stlc-meta-setting}} & \thepage & Images are inserted simultaneously and binders are renamed fresh. \\
\phantomsection\label{entry:symbol-registry-02150}\(\mathsf{SN}_\beta\) & \texttt{\detokenize{\mathsf{SN}_\beta}} & set of strongly beta-normalizing raw annotated terms & \hyperref[sec:p07-stlc-metatheory]{\ref*{sec:p07-stlc-metatheory}} & n/a & \hyperref[def:p07-stlc-reducibility]{\ref*{def:p07-stlc-reducibility}} & \thepage & The subscript specifies beta reduction; raw members need not carry typing derivations. \\
\phantomsection\label{entry:symbol-registry-02151}\(\mathcal R_\alpha\) & \texttt{\detokenize{\mathcal R_\alpha}} & reducibility predicate indexed by a finite simple type & \hyperref[sec:p07-stlc-metatheory]{\ref*{sec:p07-stlc-metatheory}} & n/a & \hyperref[def:p07-stlc-reducibility]{\ref*{def:p07-stlc-reducibility}} & \thepage & Defined on raw terms; the fundamental lemma connects the predicate with a typing derivation. \\
\phantomsection\label{entry:symbol-registry-02152}\(h(M)\) & \texttt{\detokenize{h}} & greatest beta-reduction length from a strongly normalizing raw term & \hyperref[annex:p7.stlc.strong-normalization]{\ref*{annex:p7.stlc.strong-normalization}} & n/a & \hyperref[annex:p7.stlc.strong-normalization]{\ref*{annex:p7.stlc.strong-normalization}} & \thepage & Existence follows from finite branching and absence of an infinite reduction; each immediate reduct has smaller height. \\
\phantomsection\label{entry:symbol-registry-02278}\(\mathcal T_{\mathsf F}\) & \texttt{\detokenize{\mathcal T_{\mathsf F}}} & set of all finite System F types & \hyperref[sec:p07-system-f-normalization]{\ref*{sec:p07-system-f-normalization}} & n/a & \hyperref[def:p07-system-f-sn-raw]{\ref*{def:p07-system-f-sn-raw}} & \thepage & Includes universally quantified types and open types \\
\phantomsection\label{entry:symbol-registry-02279}\(\mathcal M_{\mathsf F}\) & \texttt{\detokenize{\mathcal M_{\mathsf F}}} & set of raw Church-style System F terms & \hyperref[sec:p07-system-f-normalization]{\ref*{sec:p07-system-f-normalization}} & n/a & \hyperref[def:p07-system-f-sn-raw]{\ref*{def:p07-system-f-sn-raw}} & \thepage & Terms modulo alpha-renaming and without a required typing derivation \\
\phantomsection\label{entry:symbol-registry-02280}\(\to_{\mathsf F}\) & \texttt{\detokenize{\to_{\mathsf F}}} & one combined compatible term-beta or type-beta contraction & \hyperref[sec:p07-system-f-normalization]{\ref*{sec:p07-system-f-normalization}} & n/a & \hyperref[def:p07-system-f-sn-raw]{\ref*{def:p07-system-f-sn-raw}} & \thepage & Includes positions under both abstraction constructors \\
\phantomsection\label{entry:symbol-registry-02281}\(\to_{\mathsf F}^*\) & \texttt{\detokenize{\to_{\mathsf F}^*}} & finite reflexive transitive reduction closure & \hyperref[sec:p07-system-f-normalization]{\ref*{sec:p07-system-f-normalization}} & n/a & \hyperref[def:p07-system-f-sn-raw]{\ref*{def:p07-system-f-sn-raw}} & \thepage & Zero contractions permitted \\
\phantomsection\label{entry:symbol-registry-02282}\(\mathsf{SN}_{\mathsf F}\) & \texttt{\detokenize{\mathsf{SN}_{\mathsf F}}} & set of raw terms strongly normalizing for the combined relation & \hyperref[sec:p07-system-f-normalization]{\ref*{sec:p07-system-f-normalization}} & n/a & \hyperref[def:p07-system-f-sn-raw]{\ref*{def:p07-system-f-sn-raw}} & \thepage & Not restricted to typable or closed raw terms \\
\phantomsection\label{entry:symbol-registry-02283}\(h_{\mathsf F}(M)\) & \texttt{\detokenize{h_{\mathsf F}(M)}} & greatest combined reduction length from a strongly normalizing term & \hyperref[sec:p07-system-f-normalization]{\ref*{sec:p07-system-f-normalization}} & n/a & \hyperref[def:p07-system-f-sn-raw]{\ref*{def:p07-system-f-sn-raw}} & \thepage & Finite by finite branching and absence of an infinite path \\
\phantomsection\label{entry:symbol-registry-02284}\(\Delta,\Gamma\) & \texttt{\detokenize{\Delta,}\allowbreak\detokenize{\Gamma}} & type-variable and term-variable typing contexts & \hyperref[sec:p07-system-f-normalization]{\ref*{sec:p07-system-f-normalization}} & n/a & \hyperref[def:p07-system-f-sn-raw]{\ref*{def:p07-system-f-sn-raw}} & \thepage & Abbreviations for the exact sheet contexts \\
\phantomsection\label{entry:symbol-registry-02285}\(\tau\) & \texttt{\detokenize{\tau}} & simultaneous capture-avoiding type substitution & \hyperref[sec:p07-system-f-normalization]{\ref*{sec:p07-system-f-normalization}} & n/a & \hyperref[def:p07-system-f-sn-substitutions]{\ref*{def:p07-system-f-sn-substitutions}} & \thepage & Acts in annotations and type arguments before term substitution \\
\phantomsection\label{entry:symbol-registry-02286}\(\sigma\) & \texttt{\detokenize{\sigma}} & simultaneous capture-avoiding term substitution & \hyperref[sec:p07-system-f-normalization]{\ref*{sec:p07-system-f-normalization}} & n/a & \hyperref[def:p07-system-f-sn-substitutions]{\ref*{def:p07-system-f-sn-substitutions}} & \thepage & Inserted term images may be open raw terms \\
\phantomsection\label{entry:symbol-registry-02287}\(\operatorname{FV}_{\rm tm}\) & \texttt{\detokenize{\operatorname{FV}_{\rm tm}}} & free term-variable set & \hyperref[sec:p07-system-f-normalization]{\ref*{sec:p07-system-f-normalization}} & n/a & \hyperref[def:p07-system-f-sn-substitutions]{\ref*{def:p07-system-f-sn-substitutions}} & \thepage & Finite for each raw term \\
\phantomsection\label{entry:symbol-registry-02288}\(\operatorname{FV}_{\rm ty}\) & \texttt{\detokenize{\operatorname{FV}_{\rm ty}}} & free type-variable set including annotations and type arguments & \hyperref[sec:p07-system-f-normalization]{\ref*{sec:p07-system-f-normalization}} & n/a & \hyperref[def:p07-system-f-sn-substitutions]{\ref*{def:p07-system-f-sn-substitutions}} & \thepage & Both sorts of binders are respected \\
\phantomsection\label{entry:symbol-registry-02289}\(\mathfrak C_{\mathsf F}\) & \texttt{\detokenize{\mathfrak C_{\mathsf F}}} & set of all reducibility candidates & \hyperref[sec:p07-system-f-normalization]{\ref*{sec:p07-system-f-normalization}} & n/a & \hyperref[def:p07-system-f-sn-candidate]{\ref*{def:p07-system-f-sn-candidate}} & \thepage & Defined on one fixed raw-term set before interpreting types \\
\phantomsection\label{entry:symbol-registry-02290}\(R\Rightarrow S\) & \texttt{\detokenize{R\Rightarrow S}} & arrow constructor on reducibility candidates & \hyperref[sec:p07-system-f-normalization]{\ref*{sec:p07-system-f-normalization}} & n/a & \hyperref[thm:p07-system-f-sn-candidate-constructors]{\ref*{thm:p07-system-f-sn-candidate-constructors}} & \thepage & Metalinguistic operation on sets rather than an object-language connective \\
\phantomsection\label{entry:symbol-registry-02291}\(\mathcal U(H)\) & \texttt{\detokenize{\mathcal U(H)}} & universal candidate constructor indexed by types and candidates & \hyperref[sec:p07-system-f-normalization]{\ref*{sec:p07-system-f-normalization}} & n/a & \hyperref[thm:p07-system-f-sn-candidate-constructors]{\ref*{thm:p07-system-f-sn-candidate-constructors}} & \thepage & Quantifies over every finite System F type and every candidate \\
\phantomsection\label{entry:symbol-registry-02292}\(\rho\) & \texttt{\detokenize{\rho}} & candidate valuation of type variables & \hyperref[sec:p07-system-f-normalization]{\ref*{sec:p07-system-f-normalization}} & n/a & \hyperref[def:p07-system-f-sn-interpretation]{\ref*{def:p07-system-f-sn-interpretation}} & \thepage & Distinct from the term substitution used in the earlier STLC section \\
\phantomsection\label{entry:symbol-registry-02293}\(\mathcal R_\rho(\alpha)\) & \texttt{\detokenize{\mathcal R_\rho(\alpha)}} & interpretation of a type under a candidate valuation & \hyperref[sec:p07-system-f-normalization]{\ref*{sec:p07-system-f-normalization}} & n/a & \hyperref[def:p07-system-f-sn-interpretation]{\ref*{def:p07-system-f-sn-interpretation}} & \thepage & Structural recursion with candidate parameters at universal types \\
\phantomsection\label{entry:symbol-registry-02294}\(\iota\) & \texttt{\detokenize{\iota}} & polymorphic identity type in the worked example & \hyperref[sec:p07-system-f-sn-fundamental]{\ref*{sec:p07-system-f-sn-fundamental}} & n/a & \hyperref[def:p07-system-f-sn-example]{\ref*{def:p07-system-f-sn-example}} & \thepage & The example instantiates the identity at this universally quantified type \\
\phantomsection\label{entry:symbol-registry-02295}\(I\) & \texttt{\detokenize{I}} & explicitly typed polymorphic identity term & \hyperref[sec:p07-system-f-sn-fundamental]{\ref*{sec:p07-system-f-sn-fundamental}} & n/a & \hyperref[def:p07-system-f-sn-example]{\ref*{def:p07-system-f-sn-example}} & \thepage & Defined as a type abstraction followed by a term abstraction \\
\phantomsection\label{entry:symbol-registry-02296}\(\Omega\) & \texttt{\detokenize{\Omega}} & raw self-reducing term used to delimit the expansion lemma & \hyperref[sec:p07-system-f-sn-fundamental]{\ref*{sec:p07-system-f-sn-fundamental}} & n/a & \hyperref[def:p07-system-f-sn-side-condition]{\ref*{def:p07-system-f-sn-side-condition}} & \thepage & Not asserted typable in System F \\
\phantomsection\label{entry:symbol-registry-02297}\(\mathrm{SN}\) & \texttt{\detokenize{\mathrm{SN}}} & local proof abbreviation for the combined strong-normalization set & \hyperref[annex:p7.system-f.candidate-constructors]{\ref*{annex:p7.system-f.candidate-constructors}} & n/a & \hyperref[annex:p7.system-f.candidate-constructors]{\ref*{annex:p7.system-f.candidate-constructors}} & \thepage & Same set as mathsf SN with System F subscript in the body \\
\phantomsection\label{entry:symbol-registry-02468}\(\mathsf{NJ}_{\mathrm{nv}}\) & \texttt{\detokenize{\mathsf{NJ}_{\mathrm{nv}}}} & the existing nonvacuous finite-tree NJ presentation & \hyperref[sec:p07-natural-deduction-normalization]{\ref*{sec:p07-natural-deduction-normalization}} & meta & \hyperref[def:p07-nj-presentations]{\ref*{def:p07-nj-presentations}} & \thepage & Refers exactly to sheet:nj; the sheet itself is preserved. \\
\phantomsection\label{entry:symbol-registry-02469}\(\mathsf{NJ}_{\mathrm v}\) & \texttt{\detokenize{\mathsf{NJ}_{\mathrm v}}} & the distinct empty-discharge finite-tree NJ presentation & \hyperref[sec:p07-natural-deduction-normalization]{\ref*{sec:p07-natural-deduction-normalization}} & meta & \hyperref[def:p07-nj-presentations]{\ref*{def:p07-nj-presentations}} & \thepage & Sheet:p07-nj-vacuous; empty classes are allowed in implication introduction and both case-elimination rules. \\
\phantomsection\label{entry:symbol-registry-02470}\(\mathrm{OA}\) & \texttt{\detokenize{\mathrm{OA}}} & set of actual open assumption formulas of a derivation & \hyperref[sec:p07-natural-deduction-normalization]{\ref*{sec:p07-natural-deduction-normalization}} & meta & \hyperref[def:p07-nj-presentations]{\ref*{def:p07-nj-presentations}} & \thepage & An unused member of a supporting context is not an actual assumption occurrence. \\
\phantomsection\label{entry:symbol-registry-02471}\(M:A\) & \texttt{\detokenize{M:A}} & formula-annotated natural-deduction proof term & \hyperref[sec:p07-natural-deduction-normalization]{\ref*{sec:p07-natural-deduction-normalization}} & meta & \hyperref[def:p07-nj-proof-terms]{\ref*{def:p07-nj-proof-terms}} & \thepage & Proof-variable and individual-variable alphabets are distinct. \\
\phantomsection\label{entry:symbol-registry-02472}\(\pi_i\) & \texttt{\detokenize{\pi_i}} & conjunction projection on a proof term & \hyperref[sec:p07-natural-deduction-normalization]{\ref*{sec:p07-natural-deduction-normalization}} & meta & \hyperref[def:p07-nj-proof-terms]{\ref*{def:p07-nj-proof-terms}} & \thepage & i is one or two. \\
\phantomsection\label{entry:symbol-registry-02473}\(\iota_i\) & \texttt{\detokenize{\iota_i}} & disjunction injection on a proof term & \hyperref[sec:p07-natural-deduction-normalization]{\ref*{sec:p07-natural-deduction-normalization}} & meta & \hyperref[def:p07-nj-proof-terms]{\ref*{def:p07-nj-proof-terms}} & \thepage & The actual derivation retains its missing-summand formula annotation. \\
\phantomsection\label{entry:symbol-registry-02474}\(\mathsf{case}\) & \texttt{\detokenize{\mathsf{case}}} & disjunction elimination with two bound assumption classes & \hyperref[sec:p07-natural-deduction-normalization]{\ref*{sec:p07-natural-deduction-normalization}} & meta & \hyperref[def:p07-nj-proof-terms]{\ref*{def:p07-nj-proof-terms}} & \thepage & Fresh copies retain the external discharge ownership of free proof variables. \\
\phantomsection\label{entry:symbol-registry-02475}\(\mathsf{abort}\) & \texttt{\detokenize{\mathsf{abort}}} & bottom elimination to an annotated formula & \hyperref[sec:p07-natural-deduction-normalization]{\ref*{sec:p07-natural-deduction-normalization}} & meta & \hyperref[def:p07-nj-proof-terms]{\ref*{def:p07-nj-proof-terms}} & \thepage & There is no matching bottom introduction. \\
\phantomsection\label{entry:symbol-registry-02476}\(\Lambda\) & \texttt{\detokenize{\Lambda}} & universal proof introduction binding an individual eigenvariable & \hyperref[sec:p07-natural-deduction-normalization]{\ref*{sec:p07-natural-deduction-normalization}} & meta & \hyperref[def:p07-nj-proof-terms]{\ref*{def:p07-nj-proof-terms}} & \thepage & Individual substitution is distinct from proof substitution. \\
\phantomsection\label{entry:symbol-registry-02477}\(M[t]\) & \texttt{\detokenize{M[t]}} & universal proof elimination at an individual term & \hyperref[sec:p07-natural-deduction-normalization]{\ref*{sec:p07-natural-deduction-normalization}} & meta & \hyperref[def:p07-nj-proof-terms]{\ref*{def:p07-nj-proof-terms}} & \thepage & The actual conclusion is the total capture-avoiding formula instance. \\
\phantomsection\label{entry:symbol-registry-02478}\(\mathsf{pack}\) & \texttt{\detokenize{\mathsf{pack}}} & existential introduction with witness and proof & \hyperref[sec:p07-natural-deduction-normalization]{\ref*{sec:p07-natural-deduction-normalization}} & meta & \hyperref[def:p07-nj-proof-terms]{\ref*{def:p07-nj-proof-terms}} & \thepage & Actual terms retain the witness and existential formula. \\
\phantomsection\label{entry:symbol-registry-02479}\(\mathsf{unpack}\) & \texttt{\detokenize{\mathsf{unpack}}} & existential elimination binding an eigenvariable and proof class & \hyperref[sec:p07-natural-deduction-normalization]{\ref*{sec:p07-natural-deduction-normalization}} & meta & \hyperref[def:p07-nj-proof-terms]{\ref*{def:p07-nj-proof-terms}} & \thepage & Strict absence conditions are retained on the actual quantified proof. \\
\phantomsection\label{entry:symbol-registry-02480}\(E[-]\) & \texttt{\detokenize{E[-]}} & one-step elimination frame with hole in its major premise & \hyperref[sec:p07-natural-deduction-normalization]{\ref*{sec:p07-natural-deduction-normalization}} & meta & \hyperref[def:p07-nj-elimination-frames]{\ref*{def:p07-nj-elimination-frames}} & \thepage & Exactly seven frame forms; no minor-premise or introduction permutation is implicit. \\
\phantomsection\label{entry:symbol-registry-02481}\(\to_{\mathrm D}\) & \texttt{\detokenize{\to_{\mathrm D}}} & compatible principal conversion & \hyperref[sec:p07-natural-deduction-normalization]{\ref*{sec:p07-natural-deduction-normalization}} & meta & \hyperref[def:p07-nj-vacuous-normality]{\ref*{def:p07-nj-vacuous-normality}} & \thepage & Seven concrete introduction-elimination rows. \\
\phantomsection\label{entry:symbol-registry-02482}\(\to_{\mathrm P}\) & \texttt{\detokenize{\to_{\mathrm P}}} & compatible case or existential permutation & \hyperref[sec:p07-natural-deduction-normalization]{\ref*{sec:p07-natural-deduction-normalization}} & meta & \hyperref[def:p07-nj-vacuous-normality]{\ref*{def:p07-nj-vacuous-normality}} & \thepage & Fourteen concrete rows, with copied branch freshness. \\
\phantomsection\label{entry:symbol-registry-02483}\(\to_{\mathrm B}\) & \texttt{\detokenize{\to_{\mathrm B}}} & compatible bottom-absorption conversion & \hyperref[sec:p07-natural-deduction-normalization]{\ref*{sec:p07-natural-deduction-normalization}} & meta & \hyperref[def:p07-nj-vacuous-normality]{\ref*{def:p07-nj-vacuous-normality}} & \thepage & Seven concrete rows; the proved strategy postpones this class. \\
\phantomsection\label{entry:symbol-registry-02484}\(e\) & \texttt{\detokenize{e}} & termination-only erasure of individual and formula annotations & \hyperref[sec:p07-natural-deduction-normalization]{\ref*{sec:p07-natural-deduction-normalization}} & meta & \hyperref[def:p07-nj-erasure]{\ref*{def:p07-nj-erasure}} & \thepage & It is not an identification of actual quantified derivations. \\
\phantomsection\label{entry:symbol-registry-02485}\(\mathsf{ai}\) & \texttt{\detokenize{\mathsf{ai}}} & erased universal-introduction constructor & \hyperref[sec:p07-natural-deduction-normalization]{\ref*{sec:p07-natural-deduction-normalization}} & meta & \hyperref[def:p07-nj-erasure]{\ref*{def:p07-nj-erasure}} & \thepage & Retained so universal principal conversion has a positive beta simulation. \\
\phantomsection\label{entry:symbol-registry-02486}\(\mathsf{ae}\) & \texttt{\detokenize{\mathsf{ae}}} & erased universal-elimination constructor & \hyperref[sec:p07-natural-deduction-normalization]{\ref*{sec:p07-natural-deduction-normalization}} & meta & \hyperref[def:p07-nj-erasure]{\ref*{def:p07-nj-erasure}} & \thepage & Only its individual term is erased. \\
\phantomsection\label{entry:symbol-registry-02487}\(A^{\mathrm v}\) & \texttt{\detokenize{A^{\mathrm v}}} & value type in the arrow-only simulation & \hyperref[sec:p07-natural-deduction-normalization]{\ref*{sec:p07-natural-deduction-normalization}} & meta & \hyperref[def:p07-nj-cps-types]{\ref*{def:p07-nj-cps-types}} & \thepage & Distinct role from the v naming the empty-discharge source calculus. \\
\phantomsection\label{entry:symbol-registry-02488}\(A^{\mathrm c}\) & \texttt{\detokenize{A^{\mathrm c}}} & computation type in the arrow-only simulation & \hyperref[sec:p07-natural-deduction-normalization]{\ref*{sec:p07-natural-deduction-normalization}} & meta & \hyperref[def:p07-nj-cps-types]{\ref*{def:p07-nj-cps-types}} & \thepage & Defined as continuation-to-o applied twice. \\
\phantomsection\label{entry:symbol-registry-02489}\(\star\) & \texttt{\detokenize{\star}} & fresh target variable of individual-marker type & \hyperref[sec:p07-natural-deduction-normalization]{\ref*{sec:p07-natural-deduction-normalization}} & meta & \hyperref[def:p07-nj-cps-types]{\ref*{def:p07-nj-cps-types}} & \thepage & An open STLC marker; no witness is thereby added to the object logic. \\
\phantomsection\label{entry:symbol-registry-02490}\(\mathcal T\) & \texttt{\detokenize{\mathcal T}} & modified CPS translation of an erased proof term & \hyperref[sec:p07-natural-deduction-normalization]{\ref*{sec:p07-natural-deduction-normalization}} & meta & \hyperref[def:p07-nj-modified-cps]{\ref*{def:p07-nj-modified-cps}} & \thepage & A typed arrow-only term with an explicit continuation wrapper. \\
\phantomsection\label{entry:symbol-registry-02491}\(\mathcal C\) & \texttt{\detokenize{\mathcal C}} & CPS term with a supplied continuation & \hyperref[sec:p07-natural-deduction-normalization]{\ref*{sec:p07-natural-deduction-normalization}} & meta & \hyperref[def:p07-nj-modified-cps]{\ref*{def:p07-nj-modified-cps}} & \thepage & The bottom clause keeps its continuation inside an uncontracted administrative beta redex. \\
\phantomsection\label{entry:symbol-registry-02492}\(n\) & \texttt{\detokenize{n}} & positive auxiliary branching weight & \hyperref[sec:p07-natural-deduction-normalization]{\ref*{sec:p07-natural-deduction-normalization}} & meta & \hyperref[def:p07-nj-permutation-measure]{\ref*{def:p07-nj-permutation-measure}} & \thepage & Preserved by every permutation; corresponds to the source's hash-mark quantity. \\
\phantomsection\label{entry:symbol-registry-02493}\(w\) & \texttt{\detokenize{w}} & positive permutation-decrease measure & \hyperref[sec:p07-natural-deduction-normalization]{\ref*{sec:p07-natural-deduction-normalization}} & meta & \hyperref[def:p07-nj-permutation-measure]{\ref*{def:p07-nj-permutation-measure}} & \thepage & Accounts for copied auxiliary proof branches. \\
\phantomsection\label{entry:symbol-registry-02494}\(a_E,b_E\) & \texttt{\detokenize{a_E,}\allowbreak\detokenize{b_E}} & positive frame coefficients for the two arithmetic recurrences & \hyperref[sec:p07-natural-deduction-normalization]{\ref*{sec:p07-natural-deduction-normalization}} & meta & \hyperref[def:p07-nj-permutation-measure]{\ref*{def:p07-nj-permutation-measure}} & \thepage & Independent of the proof in the frame's hole. \\
\phantomsection\label{entry:symbol-registry-02495}\(h_\beta\) & \texttt{\detokenize{h_\beta}} & greatest beta-reduction length of a strongly normalizing target term & \hyperref[annex:p7.nj.vacuous-normalization]{\ref*{annex:p7.nj.vacuous-normalization}} & meta & \hyperref[annex:p7.nj.vacuous-normalization]{\ref*{annex:p7.nj.vacuous-normalization}} & \thepage & Finite by compatible STLC strong normalization and finite branching. \\
\phantomsection\label{entry:symbol-registry-02496}\(\mu\) & \texttt{\detokenize{\mu}} & lexicographic height and permutation-weight measure & \hyperref[annex:p7.nj.vacuous-normalization]{\ref*{annex:p7.nj.vacuous-normalization}} & meta & \hyperref[annex:p7.nj.vacuous-normalization]{\ref*{annex:p7.nj.vacuous-normalization}} & \thepage & Used for principal/permutation Phase I; Phase II instead decreases constructor count. \\
\phantomsection\label{entry:symbol-registry-02497}\(\Rrightarrow_{\mathsf F}\) & \texttt{\detokenize{\Rrightarrow_{\mathsf F}}} & auxiliary one-step parallel relation for combined System F beta reduction & \hyperref[sec:p07-system-f-confluence-and-checking]{\ref*{sec:p07-system-f-confluence-and-checking}} & n/a & \hyperref[def:p07-system-f-parallel]{\ref*{def:p07-system-f-parallel}} & \thepage & Inductively defined with reflexive cases but without a transitivity clause \\
\phantomsection\label{entry:symbol-registry-02498}\(M^\star\) & \texttt{\detokenize{M^\star}} & complete development of the redexes in the original raw term & \hyperref[sec:p07-system-f-confluence-and-checking]{\ref*{sec:p07-system-f-confluence-and-checking}} & n/a & \hyperref[def:p07-system-f-development]{\ref*{def:p07-system-f-development}} & \thepage & May retain newly exposed redexes and need not be normal \\
\phantomsection\label{entry:symbol-registry-02499}\(\mathsf{WF}(\Delta,\alpha)\) & \texttt{\detokenize{\mathsf{WF}(\Delta,}\allowbreak\detokenize{\alpha)}} & finite type-formation check in a type context & \hyperref[sec:p07-system-f-confluence-and-checking]{\ref*{sec:p07-system-f-confluence-and-checking}} & n/a & \hyperref[def:p07-system-f-checking-inputs]{\ref*{def:p07-system-f-checking-inputs}} & \thepage & Free names are checked against the type context \\
\phantomsection\label{entry:symbol-registry-02500}\(\mathsf{Synth}(\Delta,\Gamma,M)\) & \texttt{\detokenize{\mathsf{Synth}(\Delta,}\allowbreak\detokenize{\Gamma,}\allowbreak\detokenize{M)}} & syntax-directed synthesis for a fully explicit Church-style term & \hyperref[sec:p07-system-f-checking-algorithm]{\ref*{sec:p07-system-f-checking-algorithm}} & n/a & \hyperref[def:p07-system-f-synthesis]{\ref*{def:p07-system-f-synthesis}} & \thepage & Returns a type or the failure marker without implicit type instantiation \\
\phantomsection\label{entry:symbol-registry-02501}\(\mathsf{fail}\) & \texttt{\detokenize{\mathsf{fail}}} & failure result of explicit type synthesis & \hyperref[sec:p07-system-f-checking-algorithm]{\ref*{sec:p07-system-f-checking-algorithm}} & n/a & \hyperref[def:p07-system-f-synthesis]{\ref*{def:p07-system-f-synthesis}} & \thepage & Failure is not a System F type \\
\phantomsection\label{entry:symbol-registry-02502}\(A_{i,j}\) & \texttt{\detokenize{A_{i,j}}} & raw term at a vertex of the finite confluence grid & \hyperref[annex:p7.system-f.confluence]{\ref*{annex:p7.system-f.confluence}} & n/a & \hyperref[annex:p7.system-f.confluence]{\ref*{annex:p7.system-f.confluence}} & \thepage & Abstract finite joining construction independent of normalization \\
\phantomsection\label{entry:symbol-registry-02503}\([A]_\alpha\) & \texttt{\detokenize{[A]_\alpha}} & alpha-equivalence class of a literal formula & P7 intuitionistic cut elimination & n/a & \hyperref[def:p07-lj-alpha-proof]{\ref*{def:p07-lj-alpha-proof}} & \thepage & Auxiliary proof labels; the exact literal endsequent is recovered after normalization. \\
\phantomsection\label{entry:symbol-registry-02504}\(\rho(A)\) & \texttt{\detokenize{\rho}} & number of logical symbols in a formula & P7 intuitionistic cut elimination & n/a & \hyperref[def:p07-lj-alpha-proof]{\ref*{def:p07-lj-alpha-proof}} & \thepage & Atoms rank zero and bottom rank one; term symbols are not counted; defined negation expanded. \\
\phantomsection\label{entry:symbol-registry-02505}\(h(D)\) & \texttt{\detokenize{h}} & derivation height & P7 intuitionistic cut elimination & n/a & \hyperref[def:p07-lj-alpha-proof]{\ref*{def:p07-lj-alpha-proof}} & \thepage & Initial height zero and inference height one plus maximum premise height. \\
\phantomsection\label{entry:symbol-registry-02506}\(\Gamma\) & \texttt{\detokenize{\Gamma}} & supplied finite-multiset antecedent & P7 marked multicut & n/a & \hyperref[def:p07-lj-marked-multicut]{\ref*{def:p07-lj-marked-multicut}} & \thepage & Exactly one copy is retained in the multicut output; may contain the cut formula. \\
\phantomsection\label{entry:symbol-registry-02507}\(\Pi\) & \texttt{\detokenize{\Pi}} & unselected right-premise antecedent remainder & P7 marked multicut & n/a & \hyperref[def:p07-lj-marked-multicut]{\ref*{def:p07-lj-marked-multicut}} & \thepage & Every unselected occurrence remains even when equal to the cut formula. \\
\phantomsection\label{entry:symbol-registry-02508}\(A^n\) & \texttt{\detokenize{A^n}} & n individually selected antecedent occurrences & P7 marked multicut & n/a & \hyperref[def:p07-lj-marked-multicut]{\ref*{def:p07-lj-marked-multicut}} & \thepage & Occurrence bookkeeping; n zero is derived padding and sheet Mix still requires n at least one. \\
\phantomsection\label{entry:symbol-registry-02509}\(\Gamma^m\) & \texttt{\detokenize{\Gamma^m}} & m designated copies of a finite multiset & P7 intuitionistic context collapse & n/a & \hyperref[thm:p07-lj-context-collapse]{\ref*{thm:p07-lj-context-collapse}} & \thepage & Contraction targets supplied copies and preserves the other context. \\
\phantomsection\label{entry:symbol-registry-02510}\(D,E\) & \texttt{\detokenize{D;E}} & cut-free input derivation trees & P7 marked multicut & n/a & \hyperref[def:p07-lj-marked-multicut]{\ref*{def:p07-lj-marked-multicut}} & \thepage & No Cut or Mix nodes in either input of the local admissibility lemma. \\
\phantomsection\label{entry:symbol-registry-02511}\(\theta\) & \texttt{\detokenize{\theta}} & finite simultaneous substitution on free variables & P7 intuitionistic substitution & n/a & \hyperref[def:p07-lj-alpha-proof]{\ref*{def:p07-lj-alpha-proof}} & \thepage & Capture avoiding on alpha-classes; respects constructor rank and fresh eigenparameters. \\
\phantomsection\label{entry:symbol-registry-02512}\(z\) & \texttt{\detokenize{z}} & fresh eigenparameter or temporary substitution variable & P7 intuitionistic quantifier reductions & n/a & \hyperref[annex:p7.lj.alpha-substitution]{\ref*{annex:p7.lj.alpha-substitution}} & \thepage & Chosen outside the finite variables specified in each reduction. \\
\phantomsection\label{entry:symbol-registry-02513}\(t_0\) & \texttt{\detokenize{t_0}} & grounded witness term & P7 intuitionistic existence property & n/a & \hyperref[annex:p7.lj.existence-property]{\ref*{annex:p7.lj.existence-property}} & \thepage & Requires a closed existential theorem and an available ground term in the signature. \\
\phantomsection\label{entry:symbol-registry-02584}\(\operatorname{rk}_{\mathrm{fm}}(A)\) & \texttt{\detokenize{\operatorname{rk}_{\mathrm{fm}}}} & formula rank & P7 classical cut elimination & n/a & \hyperref[def:p07-classical-cut-domain]{\ref*{def:p07-classical-cut-domain}} & \thepage & Counts logical symbols after expanding defined negation; atoms rank zero and bottom rank one. \\
\phantomsection\label{entry:symbol-registry-02585}\(D,E\) & \texttt{\detokenize{D;E}} & finite derivation trees & P7 classical cut elimination & n/a & \hyperref[def:p07-classical-cut-domain]{\ref*{def:p07-classical-cut-domain}} & \thepage & The local Mix lemma takes cut-free inputs; the global elimination theorem also allows Cut and Mix in its original tree. \\
\phantomsection\label{entry:symbol-registry-02586}\(A^m\) & \texttt{\detokenize{A^m}} & m individually selected formula occurrences & P7 classical cut elimination & n/a & \hyperref[def:p07-classical-marked-mix]{\ref*{def:p07-classical-marked-mix}} & \thepage & An occurrence multiplicity in a multiset context; unselected equal formulas remain in the retained context. \\
\phantomsection\label{entry:symbol-registry-02587}\(A^n\) & \texttt{\detokenize{A^n}} & n individually selected formula occurrences & P7 classical cut elimination & n/a & \hyperref[def:p07-classical-marked-mix]{\ref*{def:p07-classical-marked-mix}} & \thepage & Selected antecedent copies in the second Mix premise; contraction marks both selected ancestors. \\
\phantomsection\label{entry:symbol-registry-02588}\(m,n\) & \texttt{\detokenize{m;n}} & selected occurrence counts & P7 classical cut elimination & n/a & \hyperref[def:p07-classical-marked-mix]{\ref*{def:p07-classical-marked-mix}} & \thepage & Nonnegative in the strengthened admissibility statement; the registered Mix rule requires both positive. \\
\phantomsection\label{entry:symbol-registry-02589}\(\Gamma,\Delta,\Pi,\Sigma\) & \texttt{\detokenize{\Gamma;\Delta;\Pi;\Sigma}} & retained finite multiset contexts & P7 classical cut elimination & n/a & \hyperref[def:p07-classical-marked-mix]{\ref*{def:p07-classical-marked-mix}} & \thepage & Every retained occurrence survives with its intended multiplicity even when its formula equals the eliminated formula. \\
\phantomsection\label{entry:symbol-registry-02590}\([A]\) & \texttt{\detokenize{[A]}} & alpha-equivalence class of a literal formula & P7 classical cut elimination & n/a & \hyperref[def:p07-classical-alpha-construction]{\ref*{def:p07-classical-alpha-construction}} & \thepage & Syntactic auxiliary proof label; equality of classes is not semantic equivalence. \\
\phantomsection\label{entry:symbol-registry-02591}\(\operatorname{rk}_{\mathrm{fm}}(A),\operatorname{ht}(D)+\operatorname{ht}(E)\) & \texttt{\detokenize{\operatorname{rk}_{\mathrm{fm}};\operatorname{ht}}} & lexicographic measure of a Mix recursion problem & P7 classical cut elimination & n/a & \hyperref[def:p07-classical-cut-construction]{\ref*{def:p07-classical-cut-construction}} & \thepage & Same-rank calls use an original immediate premise; principal auxiliary calls have smaller rank with unrestricted finite heights. \\
\phantomsection\label{entry:symbol-registry-02592}\(\sigma\) & \texttt{\detokenize{\sigma}} & finite simultaneous substitution on free variables & P7 classical cut elimination & n/a & \hyperref[annex:p07.classical.alpha-identity]{\ref*{annex:p07.classical.alpha-identity}} & \thepage & Range terms are inserted once; binders are freshened to avoid capture; the resulting alpha-class does not depend on fresh bound names. \\
\phantomsection\label{entry:symbol-registry-02593}\(C\sigma\) & \texttt{\detokenize{C\sigma}} & formula after a finite simultaneous substitution & P7 classical cut elimination & n/a & \hyperref[annex:p07.classical.alpha-identity]{\ref*{annex:p07.classical.alpha-identity}} & \thepage & Used only in the syntactic interchange and auxiliary substitution arguments. \\
\phantomsection\label{entry:symbol-registry-02594}\(S,T\) & \texttt{\detokenize{S;T}} & combined retained antecedent and succedent multisets & P7 classical cut elimination & n/a & \hyperref[annex:p07.classical.marked-mix]{\ref*{annex:p07.classical.marked-mix}} & \thepage & In principal reductions they combine the two retained context pairs; copied contexts are contracted to their original internal multiplicity. \\
\phantomsection\label{entry:symbol-registry-02595}\(L,R\) & \texttt{\detokenize{L;R}} & added antecedent and succedent context pair in a commuting case & P7 classical cut elimination & n/a & \hyperref[annex:p07.classical.marked-mix]{\ref*{annex:p07.classical.marked-mix}} & \thepage & The opposite input proof supplies this pair; split implication-left duplicates it and requires explicit contraction. \\
\phantomsection\label{entry:symbol-registry-02672}\(\mathcal C\) & \texttt{\detokenize{\mathcal C}} & finite binder-free S-K term domain & \hyperref[sec:p07-combinatory-abstraction]{\ref*{sec:p07-combinatory-abstraction}} & n/a & \hyperref[def:p07-combinatory-syntax]{\ref*{def:p07-combinatory-syntax}} & \thepage & Generated by variables two primitive constants and binary application. \\
\phantomsection\label{entry:symbol-registry-02673}\(\Lambda\) & \texttt{\detokenize{\Lambda}} & finite untyped lambda term domain modulo bound-variable renaming & \hyperref[sec:p07-combinatory-abstraction]{\ref*{sec:p07-combinatory-abstraction}} & n/a & \hyperref[def:p07-combinatory-syntax]{\ref*{def:p07-combinatory-syntax}} & \thepage & Distinct syntactic domain from combinatory terms. \\
\phantomsection\label{entry:symbol-registry-02674}\(\mathsf S\) & \texttt{\detokenize{\mathsf S}} & primitive substitution combinator & \hyperref[sec:p07-combinatory-abstraction]{\ref*{sec:p07-combinatory-abstraction}} & n/a & \hyperref[def:p07-combinatory-syntax]{\ref*{def:p07-combinatory-syntax}} & \thepage & Exactly the primitive S of sheet:sk-combinatory. \\
\phantomsection\label{entry:symbol-registry-02675}\(\mathsf K\) & \texttt{\detokenize{\mathsf K}} & primitive constant-function combinator & \hyperref[sec:p07-combinatory-abstraction]{\ref*{sec:p07-combinatory-abstraction}} & n/a & \hyperref[def:p07-combinatory-syntax]{\ref*{def:p07-combinatory-syntax}} & \thepage & Exactly the primitive K of sheet:sk-combinatory. \\
\phantomsection\label{entry:symbol-registry-02676}\(\mathsf I\) & \texttt{\detokenize{\mathsf I}} & abbreviation for the five-node S-K identity term & \hyperref[sec:p07-combinatory-abstraction]{\ref*{sec:p07-combinatory-abstraction}} & n/a & \hyperref[def:p07-combinatory-syntax]{\ref*{def:p07-combinatory-syntax}} & \thepage & Always expands to S K K; not an additional primitive or atom. \\
\phantomsection\label{entry:symbol-registry-02677}\(\FV(P)\) & \texttt{\detokenize{\FV}} & free-variable set of a term & \hyperref[sec:p07-combinatory-abstraction]{\ref*{sec:p07-combinatory-abstraction}} & n/a & \hyperref[def:p07-combinatory-syntax]{\ref*{def:p07-combinatory-syntax}} & \thepage & All combinatory variable occurrences are free; lambda binders remove their variables. \\
\phantomsection\label{entry:symbol-registry-02678}\(P[Q/x]\) & \texttt{\detokenize{[Q/}\allowbreak\detokenize{x]}} & tree substitution into a combinatory term & \hyperref[sec:p07-combinatory-abstraction]{\ref*{sec:p07-combinatory-abstraction}} & meta & \hyperref[def:p07-combinatory-syntax]{\ref*{def:p07-combinatory-syntax}} & \thepage & No binders occur in the receiving combinatory syntax. \\
\phantomsection\label{entry:symbol-registry-02679}\(\mathcal B_x(P)\) & \texttt{\detokenize{\mathcal B_x}} & host bracket abstraction with whole-subtree absent-variable shortcut & \hyperref[sec:p07-combinatory-abstraction]{\ref*{sec:p07-combinatory-abstraction}} & n/a & \hyperref[def:p07-combinatory-abstraction-algorithms]{\ref*{def:p07-combinatory-abstraction-algorithms}} & \thepage & Three disjoint cases; commutes literally with substitution under stated freshness. \\
\phantomsection\label{entry:symbol-registry-02680}\(\mathcal A_x(P)\) & \texttt{\detokenize{\mathcal A_x}} & atom-by-atom bracket abstraction from the selected Selinger construction & \hyperref[sec:p07-combinatory-abstraction]{\ref*{sec:p07-combinatory-abstraction}} & n/a & \hyperref[def:p07-combinatory-abstraction-algorithms]{\ref*{def:p07-combinatory-abstraction-algorithms}} & \thepage & Four cases; compound substitution commutes generally only extensionally. \\
\phantomsection\label{entry:symbol-registry-02681}\(\mathcal D\) & \texttt{\detokenize{\mathcal D}} & metavariable for an explicitly selected abstraction algorithm & \hyperref[sec:p07-combinatory-abstraction]{\ref*{sec:p07-combinatory-abstraction}} & meta & \hyperref[thm:p07-combinatory-substitution]{\ref*{thm:p07-combinatory-substitution}} & \thepage & Each theorem states its allowed algorithms; the size theorem also includes H. \\
\phantomsection\label{entry:symbol-registry-02682}\(\longrightarrow_{\mathsf{SK}}\) & \texttt{\detokenize{\longrightarrow_{\mathsf{SK}}}} & one compatible S-K contraction & \hyperref[sec:p07-combinatory-abstraction]{\ref*{sec:p07-combinatory-abstraction}} & relation & \hyperref[def:p07-combinatory-relations]{\ref*{def:p07-combinatory-relations}} & \thepage & Application contexts only; exactly the existing two primitive rules. \\
\phantomsection\label{entry:symbol-registry-02683}\(\longrightarrow_{\mathsf{SK}}^*\) & \texttt{\detokenize{\longrightarrow_{\mathsf{SK}}^*}} & reflexive transitive S-K reduction & \hyperref[sec:p07-combinatory-abstraction]{\ref*{sec:p07-combinatory-abstraction}} & relation & \hyperref[def:p07-combinatory-relations]{\ref*{def:p07-combinatory-relations}} & \thepage & Zero or finitely many weak contractions. \\
\phantomsection\label{entry:symbol-registry-02684}\(=_{\mathsf{SK}}\) & \texttt{\detokenize{=_{\mathsf{SK}}}} & weak S-K conversion & \hyperref[sec:p07-combinatory-abstraction]{\ref*{sec:p07-combinatory-abstraction}} & relation & \hyperref[def:p07-combinatory-relations]{\ref*{def:p07-combinatory-relations}} & \thepage & Equivalence closure of compatible S-K contraction. \\
\phantomsection\label{entry:symbol-registry-02685}\(=_\beta\) & \texttt{\detokenize{=_\beta}} & beta conversion of lambda terms & \hyperref[sec:p07-combinatory-abstraction]{\ref*{sec:p07-combinatory-abstraction}} & relation & \hyperref[def:p07-combinatory-relations]{\ref*{def:p07-combinatory-relations}} & \thepage & Contains conversion beneath lambda binders. \\
\phantomsection\label{entry:symbol-registry-02686}\(=_{\beta\eta}\) & \texttt{\detokenize{=_{\beta\eta}}} & beta-eta conversion of lambda terms & \hyperref[sec:p07-combinatory-abstraction]{\ref*{sec:p07-combinatory-abstraction}} & relation & \hyperref[def:p07-combinatory-relations]{\ref*{def:p07-combinatory-relations}} & \thepage & Eta requires absence of the contracted binder variable from the function term. \\
\phantomsection\label{entry:symbol-registry-02687}\(=_{\mathrm e}\) & \texttt{\detokenize{=_{\mathrm e}}} & extensional combinatory equality & \hyperref[sec:p07-combinatory-abstraction]{\ref*{sec:p07-combinatory-abstraction}} & relation & \hyperref[def:p07-combinatory-relations]{\ref*{def:p07-combinatory-relations}} & \thepage & Least application congruence with S-K equations and the fresh-variable extensionality rule. \\
\phantomsection\label{entry:symbol-registry-02688}\(E(P)\) & \texttt{\detokenize{E}} & expansion of combinatory terms into lambda terms & \hyperref[sec:p07-combinatory-abstraction]{\ref*{sec:p07-combinatory-abstraction}} & n/a & \hyperref[def:p07-combinatory-translations]{\ref*{def:p07-combinatory-translations}} & \thepage & Primitive S and K expand to the two specified closed lambda terms. \\
\phantomsection\label{entry:symbol-registry-02689}\(T_{\mathcal D}(M)\) & \texttt{\detokenize{T_{\mathcal D}}} & compilation of lambda syntax by the selected abstraction algorithm & \hyperref[sec:p07-combinatory-abstraction]{\ref*{sec:p07-combinatory-abstraction}} & n/a & \hyperref[def:p07-combinatory-translations]{\ref*{def:p07-combinatory-translations}} & \thepage & Algorithm parameter is explicit; removes inner binders first. \\
\phantomsection\label{entry:symbol-registry-02690}\(=_{\lambda\beta}\) & \texttt{\detokenize{=_{\lambda\beta}}} & beta equality transported through expansion E & \hyperref[sec:p07-combinatory-translations]{\ref*{sec:p07-combinatory-translations}} & relation & \hyperref[def:p07-combinatory-beta-presentation]{\ref*{def:p07-combinatory-beta-presentation}} & \thepage & Definition compares lambda expansions; distinct from weak S-K conversion. \\
\phantomsection\label{entry:symbol-registry-02691}\(\nu_n\) & \texttt{\detokenize{\nu_n}} & Church numeral lambda term for n & \hyperref[sec:p07-combinatory-translations]{\ref*{sec:p07-combinatory-translations}} & n/a & \hyperref[def:p07-combinatory-beta-presentation]{\ref*{def:p07-combinatory-beta-presentation}} & \thepage & Uses the existing Church numeral encoding. \\
\phantomsection\label{entry:symbol-registry-02692}\(c_n\) & \texttt{\detokenize{c_n}} & compiled Church numeral for a fixed chosen algorithm & \hyperref[sec:p07-combinatory-translations]{\ref*{sec:p07-combinatory-translations}} & n/a & \hyperref[def:p07-combinatory-beta-presentation]{\ref*{def:p07-combinatory-beta-presentation}} & \thepage & Observation is specified through transported beta equality. \\
\phantomsection\label{entry:symbol-registry-02693}\(\Rightarrow_{\mathsf{SK}}\) & \texttt{\detokenize{\Rightarrow_{\mathsf{SK}}}} & auxiliary parallel S-K reduction & \hyperref[sec:p07-combinatory-abstraction]{\ref*{sec:p07-combinatory-abstraction}} & relation & \hyperref[def:p07-combinatory-parallel]{\ref*{def:p07-combinatory-parallel}} & \thepage & Four generating clauses on the new sheet; used in the confluence proof. \\
\phantomsection\label{entry:symbol-registry-02694}\(d(P)\) & \texttt{\detokenize{d}} & complete parallel development of a combinatory term & \hyperref[annex:p07.combinatory.confluence]{\ref*{annex:p07.combinatory.confluence}} & n/a & \hyperref[annex:p07.combinatory.confluence]{\ref*{annex:p07.combinatory.confluence}} & \thepage & Four exhaustive root cases in the proof; not an algorithm for normalization. \\
\phantomsection\label{entry:symbol-registry-02695}\(\mathcal H_x(P)\) & \texttt{\detokenize{\mathcal H_x}} & eta-optimized host bracket abstraction & \hyperref[sec:p07-combinatory-extensional]{\ref*{sec:p07-combinatory-extensional}} & n/a & \hyperref[def:p07-combinatory-eta-algorithm]{\ref*{def:p07-combinatory-eta-algorithm}} & \thepage & Ordered extra U x case requires x absent from FV(U). \\
\phantomsection\label{entry:symbol-registry-02696}\(A\) & \texttt{\detokenize{A}} & nonempty carrier of an applicative structure & \hyperref[sec:p07-combinatory-algebras]{\ref*{sec:p07-combinatory-algebras}} & n/a & \hyperref[def:p07-combinatory-algebra]{\ref*{def:p07-combinatory-algebra}} & \thepage & Application is a total binary operation and need not be associative. \\
\phantomsection\label{entry:symbol-registry-02697}\(s,k\) & \texttt{\detokenize{s;k}} & designated S and K elements of a combinatory algebra & \hyperref[sec:p07-combinatory-algebras]{\ref*{sec:p07-combinatory-algebras}} & n/a & \hyperref[def:p07-combinatory-algebra]{\ref*{def:p07-combinatory-algebra}} & \thepage & Distinct from the later size function s(P). \\
\phantomsection\label{entry:symbol-registry-02698}\(\llbracket p\rrbracket_\rho\) & \texttt{\detokenize{\llbracket p\rrbracket_\rho}} & value of an applicative polynomial under an assignment & \hyperref[sec:p07-combinatory-algebras]{\ref*{sec:p07-combinatory-algebras}} & n/a & \hyperref[def:p07-combinatory-algebra]{\ref*{def:p07-combinatory-algebra}} & \thepage & Coefficient atoms retain their named algebra elements. \\
\phantomsection\label{entry:symbol-registry-02699}\(\rho\) & \texttt{\detokenize{\rho}} & assignment of polynomial variables to algebra elements & \hyperref[sec:p07-combinatory-algebras]{\ref*{sec:p07-combinatory-algebras}} & n/a & \hyperref[def:p07-combinatory-algebra]{\ref*{def:p07-combinatory-algebra}} & \thepage & The abstraction identity quantifies over every assignment and argument. \\
\phantomsection\label{entry:symbol-registry-02700}\(s(P)\) & \texttt{\detokenize{s(P)}} & syntax-tree node count & \hyperref[sec:p07-combinatory-size]{\ref*{sec:p07-combinatory-size}} & n/a & \hyperref[def:p07-combinatory-size]{\ref*{def:p07-combinatory-size}} & \thepage & Counts every leaf application and lambda binder; I is expanded to five nodes. \\
\phantomsection\label{entry:symbol-registry-02701}\(b(M)\) & \texttt{\detokenize{b(M)}} & number of lambda binders & \hyperref[sec:p07-combinatory-size]{\ref*{sec:p07-combinatory-size}} & n/a & \hyperref[def:p07-combinatory-size]{\ref*{def:p07-combinatory-size}} & \thepage & Used in the proved compilation upper bound; no lower-bound claim. \\
\phantomsection\label{entry:symbol-registry-03171}\(\Sigma_1,\Sigma_2\) & \texttt{\detokenize{\Sigma_1,}\allowbreak\detokenize{\Sigma_2}} & source finite input and output alphabets & P7 Cook Reckhow source & meta & \hyperref[def:p07-cr-encoding]{\ref*{def:p07-cr-encoding}} & \thepage & String alphabets, not sequent contexts. \\
\phantomsection\label{entry:symbol-registry-03172}\(|x|,|y|,|\pi|\) & \texttt{\detokenize{|x|,}\allowbreak\detokenize{|y|,}\allowbreak\detokenize{|\pi|}} & source encoded string length & P7 Cook Reckhow source & n/a & \hyperref[def:p07-cr-encoding]{\ref*{def:p07-cr-encoding}} & \thepage & Locally typed string measure; not set cardinality or proof-line count. \\
\phantomsection\label{entry:symbol-registry-03173}\(\langle y,w\rangle\) & \texttt{\detokenize{\langle y,}\allowbreak\detokenize{w\rangle}} & explicit paired conclusion and certificate code & P7 Cook Reckhow source & n/a & \hyperref[def:p07-cr-encoding]{\ref*{def:p07-cr-encoding}} & \thepage & Both fields are present in the proof-map input; no output enumeration is required. \\
\phantomsection\label{entry:symbol-registry-03174}\(L\) & \texttt{\detokenize{L}} & target string language for an abstract proof system & P7 Cook abstract proof languages & n/a & \hyperref[def:p07-cr-proof-map]{\ref*{def:p07-cr-proof-map}} & \thepage & The empty language has no total surjection from all strings. \\
\phantomsection\label{entry:symbol-registry-03175}\(f,h\) & \texttt{\detokenize{f,}\allowbreak\detokenize{h}} & source total polynomial-time proof maps & P7 Cook Reckhow source & meta & \hyperref[def:p07-cr-proof-map]{\ref*{def:p07-cr-proof-map}} & \thepage & Common target language is required for their simulation comparison. \\
\phantomsection\label{entry:symbol-registry-03176}\(R\) & \texttt{\detokenize{R}} & polynomial-time certificate verification relation & P7 Cook certificate relation & relation & \hyperref[thm:p07-cr-map-relation]{\ref*{thm:p07-cr-map-relation}} & \thepage & Range is existential in the certificate; the conclusion is explicit in its paired proof code. \\
\phantomsection\label{entry:symbol-registry-03177}\(f_R\) & \texttt{\detokenize{f_R}} & paired-code proof map constructed from a verifier & P7 Cook Reckhow source & meta & \hyperref[thm:p07-cr-map-relation]{\ref*{thm:p07-cr-map-relation}} & \thepage & Returns stored conclusion or the fixed default. \\
\phantomsection\label{entry:symbol-registry-03178}\(y_0\) & \texttt{\detokenize{y_0}} & fixed default member of a nonempty target language & P7 Cook Reckhow source & n/a & \hyperref[def:p07-cr-encoding]{\ref*{def:p07-cr-encoding}} & \thepage & Invalid proof strings are handled without leaving the target range. \\
\phantomsection\label{entry:symbol-registry-03179}\(g\) & \texttt{\detokenize{g}} & one total polynomial-time proof translation & P7 Cook Reckhow source & meta & \hyperref[def:p07-cr-source-simulation]{\ref*{def:p07-cr-source-simulation}} & \thepage & Identity holds on valid and invalid source codes. \\
\phantomsection\label{entry:symbol-registry-03180}\(s_f(y)\) & \texttt{\detokenize{s_f(y)}} & minimum encoded f-proof length of a target string & P7 Cook Reckhow source & n/a & \hyperref[def:p07-cr-source-simulation]{\ref*{def:p07-cr-source-simulation}} & \thepage & An existential length comparison does not supply a translation algorithm. \\
\phantomsection\label{entry:symbol-registry-03181}\(l(A)\) & \texttt{\detokenize{l(A)}} & source count of atom and nullary-connective occurrences & P7 Cook Reckhow source & n/a & \hyperref[def:p07-cr-source-measures]{\ref*{def:p07-cr-source-measures}} & \thepage & Counts occurrences rather than distinct atoms; excludes unary constructors and identifier bit lengths. \\
\phantomsection\label{entry:symbol-registry-03182}\(m(A)\) & \texttt{\detokenize{m(A)}} & all constructor nodes in a formula parse tree & P7 Cook Reckhow source & n/a & \hyperref[def:p07-cr-source-measures]{\ref*{def:p07-cr-source-measures}} & \thepage & Explicit measure for complete defining sets and unary-chain overhead. \\
\phantomsection\label{entry:symbol-registry-03183}\(\lambda(\pi)\) & \texttt{\detokenize{\lambda(\pi)}} & source derivation line count & P7 Cook Reckhow source & n/a & \hyperref[def:p07-cr-source-measures]{\ref*{def:p07-cr-source-measures}} & \thepage & A reusable sequence does not duplicate a line when it is referenced twice. \\
\phantomsection\label{entry:symbol-registry-03184}\(\rho(\pi)\) & \texttt{\detokenize{\rho(\pi)}} & maximum source atom-occurrence count of a derivation line & P7 Cook Reckhow source & n/a & \hyperref[def:p07-cr-source-measures]{\ref*{def:p07-cr-source-measures}} & \thepage & Not maximum encoded bit width when unary strings or long atom identifiers occur. \\
\phantomsection\label{entry:symbol-registry-03185}\(\kappa,\kappa'\) & \texttt{\detokenize{\kappa,}\allowbreak\detokenize{\kappa'}} & fixed finite adequate connective bases & P7 Cook Reckhow source & meta & \hyperref[def:p07-cr-source-frege]{\ref*{def:p07-cr-source-frege}} & \thepage & Arities at most two; source comparisons distinguish same and different formula languages. \\
\phantomsection\label{entry:symbol-registry-03186}\(\mathcal F,\mathcal F_1,\mathcal F_2\) & \texttt{\detokenize{\mathcal F,}\allowbreak\detokenize{\mathcal F_1,}\allowbreak\detokenize{\mathcal F_2}} & source implicationally complete finite Frege systems & P7 Cook Reckhow source & meta & \hyperref[def:p07-cr-source-frege]{\ref*{def:p07-cr-source-frege}} & \thepage & Rules are sound formula schemata and use simultaneous substitution. \\
\phantomsection\label{entry:symbol-registry-03187}\(R\) & \texttt{\detokenize{R}} & fixed schematic calculus rule & P7 Cook schematic rules & meta & \hyperref[annex:p7.cr.frege-comparison]{\ref*{annex:p7.cr.frege-comparison}} & \thepage & Rule template rather than the separately scoped verifier relation. \\
\phantomsection\label{entry:symbol-registry-03188}\(\mathcal N\) & \texttt{\detokenize{\mathcal N}} & source ordered-context natural-deduction system & P7 Cook Reckhow source & meta & \hyperref[def:p07-cr-source-nd]{\ref*{def:p07-cr-source-nd}} & \thepage & Proofs are reusable sequences of line records. \\
\phantomsection\label{entry:symbol-registry-03189}\(L\) & \texttt{\detokenize{L}} & source natural-deduction line record & P7 Cook ordered ND & n/a & \hyperref[def:p07-cr-source-nd]{\ref*{def:p07-cr-source-nd}} & \thepage & A finite ordered hypothesis sequence paired with its formula conclusion. \\
\phantomsection\label{entry:symbol-registry-03190}\(L^*,L^\#\) & \texttt{\detokenize{L^*,}\allowbreak\detokenize{L^\#}} & left and right associated formula representations of an ND line & P7 Cook Reckhow source & n/a & \hyperref[def:p07-cr-source-nd]{\ref*{def:p07-cr-source-nd}} & \thepage & Star is a source representation mark here, not relation closure. \\
\phantomsection\label{entry:symbol-registry-03191}\(e\mathcal F,s\mathcal F\) & \texttt{\detokenize{e\mathcal F,}\allowbreak\detokenize{s\mathcal F}} & extension and theorem-only substitution versions of a Frege system & P7 Cook Reckhow source & meta & \hyperref[def:p07-cr-source-extension]{\ref*{def:p07-cr-source-extension}} & \thepage & Extension variables occur in neither hypotheses nor final conclusion; substitution proofs have no hypotheses. \\
\phantomsection\label{entry:symbol-registry-03192}\(P_A,P_B,P_E\) & \texttt{\detokenize{P_A,}\allowbreak\detokenize{P_B,}\allowbreak\detokenize{P_E}} & proxy atoms associated to syntactic formula records & P7 Cook Reckhow source & object & \hyperref[def:p07-cr-gate-definitions]{\ref*{def:p07-cr-gate-definitions}} & \thepage & Identity on atomic records; fresh identifiers are allocated without storing an expanded formula in every name. \\
\phantomsection\label{entry:symbol-registry-03193}\(\operatorname{def}_{\kappa'}(A)\) & \texttt{\detokenize{\operatorname{def}_{\kappa'}(A)}} & source defining set evaluated in the target connective basis & P7 Cook Reckhow source & n/a & \hyperref[def:p07-cr-gate-definitions]{\ref*{def:p07-cr-gate-definitions}} & \thepage & One finite gate equation per distinct nonatomic subformula record. \\
\phantomsection\label{entry:symbol-registry-03194}\(L\) & \texttt{\detokenize{L}} & nondecreasing proof-line length bound & P7 Cook relative extension bounds & n/a & \hyperref[thm:p07-cr-ef-relative-bounds]{\ref*{thm:p07-cr-ef-relative-bounds}} & \thepage & Argument counts full formula constructors; at least the identity bound. \\
\phantomsection\label{entry:symbol-registry-03195}\(L^\uparrow(n)\) & \texttt{\detokenize{L^\uparrow(n)}} & prefix maximum of an arbitrary natural-number length bound & P7 Cook Reckhow source & n/a & \hyperref[annex:p7.cr.ef-relative-bounds]{\ref*{annex:p7.cr.ef-relative-bounds}} & \thepage & Makes monotonicity explicit without changing polynomial boundedness. \\
\phantomsection\label{entry:symbol-registry-03196}\(D_k\) & \texttt{\detokenize{D_k}} & left-associated conjunction of the first k extension declarations & P7 Cook Reckhow source & object & \hyperref[annex:p7.cr.substitution-simulation]{\ref*{annex:p7.cr.substitution-simulation}} & \thepage & Shared guard for theorem-only reverse-order pruning. \\
\phantomsection\label{entry:symbol-registry-03197}\(\mathcal P_n,A_n\) & \texttt{\detokenize{\mathcal P_n,}\allowbreak\detokenize{A_n}} & source n-pigeon and n-minus-one-hole clause family and its tautology & P7 Cook Reckhow source & n/a & \hyperref[def:p07-cr-source-php]{\ref*{def:p07-cr-source-php}} & \thepage & No functionality clauses; the source index shift relative to the existing CNF family is explicit. \\
\phantomsection\label{entry:symbol-registry-03198}\(B_{ij}\) & \texttt{\detokenize{B_{ij}}} & source smaller-pigeonhole substitution formula & P7 Cook Reckhow source & object & \hyperref[thm:p07-cr-php-compression]{\ref*{thm:p07-cr-php-compression}} & \thepage & Old direct hole or last-hole transfer through the last pigeon. \\
\phantomsection\label{entry:symbol-registry-03199}\(Q^k_{ij}\) & \texttt{\detokenize{Q^k_{ij}}} & source fresh stage-k pigeonhole proxy atom & P7 Cook Reckhow source & object & \hyperref[annex:p7.cr.php-compression]{\ref*{annex:p7.cr.php-compression}} & \thepage & Definitions depend only on the preceding stage; unshared expansion has three branches per level. \\
\phantomsection\label{entry:symbol-registry-03200}\(x_i,y_i\) & \texttt{\detokenize{x_i,}\allowbreak\detokenize{y_i}} & hole-one and hole-two atoms in the worked three-pigeon refutation & P7 Cook Reckhow source & object & \hyperref[annex:p7.cr.small-php]{\ref*{annex:p7.cr.small-php}} & \thepage & Aliases for P-i-one and P-i-two only in the complete finite example. \\
\phantomsection\label{entry:symbol-registry-03201}\(\sigma,\tau\) & \texttt{\detokenize{\sigma,}\allowbreak\detokenize{\tau}} & simultaneous source formula substitutions & P7 Cook formula substitutions & meta & \hyperref[annex:p7.cr.frege-comparison]{\ref*{annex:p7.cr.frege-comparison}} & \thepage & Source substitutions and their composition; these are functions on formulas, not sentence metavariables. \\
\phantomsection\label{entry:symbol-registry-03265}\(A_i\) & \texttt{\detokenize{A_i}} & host pigeon i row clause as a literal set & P7 Haken elementary setup & object & \hyperref[def:p07-haken-host-encoding]{\ref*{def:p07-haken-host-encoding}} & \thepage & Distinct from the entire formula family; no within-pigeon exclusivity input is added. \\
\phantomsection\label{entry:symbol-registry-03266}\(B_{ikj}\) & \texttt{\detokenize{B_{ikj}}} & host collision clause for pigeons i and k in hole j & P7 Haken elementary setup & object & \hyperref[def:p07-haken-host-encoding]{\ref*{def:p07-haken-host-encoding}} & \thepage & The set consists of the two negative occupancy literals. \\
\phantomsection\label{entry:symbol-registry-03267}\(\mathcal C_n\) & \texttt{\detokenize{\mathcal C_n}} & finite family of host critical full truth assignments & P7 Haken elementary setup & n/a & \hyperref[def:p07-haken-critical]{\ref*{def:p07-haken-critical}} & \thepage & Exactly n matched pigeons, all n holes, and one all-false unmatched row. \\
\phantomsection\label{entry:symbol-registry-03268}\(v_{u,h}\) & \texttt{\detokenize{v_{u,h}}} & critical assignment with unmatched pigeon u and bijection h & P7 Haken elementary setup & meta & \hyperref[def:p07-haken-critical]{\ref*{def:p07-haken-critical}} & \thepage & Full values on all occupancy variables; this is not a satisfying assignment of the complete CNF. \\
\phantomsection\label{entry:symbol-registry-03269}\(\rho_f\) & \texttt{\detokenize{\rho_f}} & partial truth assignment induced by an injective matching & P7 Haken elementary setup & meta & \hyperref[def:p07-haken-matching-restriction]{\ref*{def:p07-haken-matching-restriction}} & \thepage & Sets matched rows and occupied columns; leaves only the unmatched-row and free-column rectangle unset. \\
\phantomsection\label{entry:symbol-registry-03270}\(I,J\) & \texttt{\detokenize{I,}\allowbreak\detokenize{J}} & matched-pigeon domain and occupied-hole image & P7 Haken elementary setup & n/a & \hyperref[def:p07-haken-matching-restriction]{\ref*{def:p07-haken-matching-restriction}} & \thepage & f maps I injectively onto J, with equal finite cardinality r. \\
\phantomsection\label{entry:symbol-registry-03271}\(C\upharpoonright\rho\) & \texttt{\detokenize{C\upharpoonright\rho}} & clause restriction by partial truth values & P7 Haken elementary setup & meta & \hyperref[def:p07-haken-matching-restriction]{\ref*{def:p07-haken-matching-restriction}} & \thepage & Satisfied marker is separate from an empty residual clause; tracked proof clauses may be subclauses. \\
\phantomsection\label{entry:symbol-registry-03272}\(S\) & \texttt{\detokenize{S}} & designated route-barrier node set in a fixed proof DAG & P7 Haken elementary setup & n/a & \hyperref[def:p07-haken-first-charge]{\ref*{def:p07-haken-first-charge}} & \thepage & Every critical route must hit S before the conditional counting inequality applies. \\
\phantomsection\label{entry:symbol-registry-03273}\(b\) & \texttt{\detokenize{b}} & positive upper bound on critical assignments charged to one node & P7 Haken elementary setup & n/a & \hyperref[thm:p07-haken-charging]{\ref*{thm:p07-haken-charging}} & \thepage & A capacity bound is a hypothesis, not an imported unverified Haken constant. \\
\phantomsection\label{entry:symbol-registry-03274}\(\beta\) & \texttt{\detokenize{\beta}} & positive coefficient in a conditional exponential capacity bound & P7 Haken elementary setup & n/a & \hyperref[thm:p07-haken-charging]{\ref*{thm:p07-haken-charging}} & \thepage & No particular source-native value is claimed before the original paper is read. \\
\phantomsection\label{entry:symbol-registry-03275}\(N,W\) & \texttt{\detokenize{N,}\allowbreak\detokenize{W}} & proof-node count and total recorded literal occurrences & P7 Haken elementary setup & n/a & \hyperref[def:p07-haken-proof-measures]{\ref*{def:p07-haken-proof-measures}} & \thepage & These counts determine the direct bit-encoding bound along with variable and parent identifier lengths. \\
\phantomsection\label{entry:symbol-registry-03575}\(\mathsf{LK}_{=}\) & \texttt{\detokenize{\mathsf{LK}_{=}}} & the literal existing equality-LK calculus including Cut and Mix & \hyperref[sec:p07-equality-lk-metatheory]{\ref*{sec:p07-equality-lk-metatheory}} & meta & \hyperref[def:p07-eqlk-domain]{\ref*{def:p07-eqlk-domain}} & \thepage & Exactly sheet:lk-eq; no right replacement or reflexivity deletion is added. \\
\phantomsection\label{entry:symbol-registry-03576}\(D_\Delta\) & \texttt{\detokenize{D_\Delta}} & fixed finite-support succedent disjunction & \hyperref[sec:p07-equality-lk-metatheory]{\ref*{sec:p07-equality-lk-metatheory}} & meta & \hyperref[def:p07-eqlk-domain]{\ref*{def:p07-eqlk-domain}} & \thepage & Reuses def:p07-sequent-nd-translation; empty support denotes falsum. \\
\phantomsection\label{entry:symbol-registry-03577}\(a,b\) & \texttt{\detokenize{a;b}} & syntactically distinct constant symbols in the essential-Cut counterexample & P7 equality-LK metatheory & object & \hyperref[thm:p07-eqlk-cut-failure]{\ref*{thm:p07-eqlk-cut-failure}} & \thepage & Only this local signature uses these constants; the independent variable version works in every signature. \\
\phantomsection\label{entry:symbol-registry-03578}\(\pi\) & \texttt{\detokenize{\pi}} & finite transposition of the individual-variable alphabet in a proof-renaming construction & \hyperref[annex:p7.eqlk.nd]{\ref*{annex:p7.eqlk.nd}} & meta & \hyperref[annex:p7.eqlk.nd]{\ref*{annex:p7.eqlk.nd}} & \thepage & Renames bound and free variables; total substitution covariance is up to alpha-equivalence and is proved locally. \\
\phantomsection\label{entry:symbol-registry-04617}\(\mathcal A_i^n,\mathcal A_o^n,\mathcal A^n\) & \texttt{\detokenize{\mathcal A_i^n,}\allowbreak\detokenize{\mathcal A_o^n,}\allowbreak\detokenize{\mathcal A^n}} & indexed object logical-context and combined syntax classes & \hyperref[sec:p07-ch88]{\ref*{sec:p07-ch88}} & n/a & \hyperref[def:p07-ch88-syntax]{\ref*{def:p07-ch88-syntax}} & \thepage & Source one-star grammar at external depth n; objects include schemas and types. \\
\phantomsection\label{entry:symbol-registry-04618}\(\ast\) & \texttt{\detokenize{\ast}} & source untyped universe and empty environment term & \hyperref[sec:p07-ch88]{\ref*{sec:p07-ch88}} & n/a & \hyperref[def:p07-ch88-syntax]{\ref*{def:p07-ch88-syntax}} & \thepage & A logical-context atom with no source type; not a star self-type axiom. \\
\phantomsection\label{entry:symbol-registry-04619}\([x:A]B\) & \texttt{\detokenize{[x:A]B}} & source indexed product binder & \hyperref[sec:p07-ch88]{\ref*{sec:p07-ch88}} & n/a & \hyperref[def:p07-ch88-syntax]{\ref*{def:p07-ch88-syntax}} & \thepage & B has one added index; annotation A is outside the binder. \\
\phantomsection\label{entry:symbol-registry-04620}\((\lambda x:A)M\) & \texttt{\detokenize{(\lambda x:A)M}} & source indexed annotated abstraction & \hyperref[sec:p07-ch88]{\ref*{sec:p07-ch88}} & n/a & \hyperref[def:p07-ch88-syntax]{\ref*{def:p07-ch88-syntax}} & \thepage & The written binder name is concrete only; abstract variable occurrences are positive indices. \\
\phantomsection\label{entry:symbol-registry-04621}\(|\Gamma|\) & \texttt{\detokenize{|\Gamma|}} & source local environment length & \hyperref[sec:p07-ch88]{\ref*{sec:p07-ch88}} & n/a & \hyperref[def:p07-ch88-environments]{\ref*{def:p07-ch88-environments}} & \thepage & Source-bound length of a product-list environment; not cardinality or a proof-theoretic ordinal. \\
\phantomsection\label{entry:symbol-registry-04622}\(\Gamma;\Delta,\Gamma[x:A]\) & \texttt{\detokenize{\Gamma;\Delta,}\allowbreak\detokenize{\Gamma[x:A]}} & source context concatenation and fresh final declaration & \hyperref[sec:p07-ch88]{\ref*{sec:p07-ch88}} & n/a & \hyperref[def:p07-ch88-environments]{\ref*{def:p07-ch88-environments}} & \thepage & Delta remains a term in a type position; concatenation replaces final star. \\
\phantomsection\label{entry:symbol-registry-04623}\(\xi_i,M^+,\operatorname{add}(j,M)\) & \texttt{\detokenize{\xi_i,}\allowbreak\detokenize{M^+,}\allowbreak\detokenize{\operatorname{add}}} & corrected free-reference relocation & \hyperref[sec:p07-ch88]{\ref*{sec:p07-ch88}} & n/a & \hyperref[def:p07-ch88-index-operations]{\ref*{def:p07-ch88-index-operations}} & \thepage & One-based bound indices at most cutoff i stay fixed; annotations and bodies use different cutoffs. \\
\phantomsection\label{entry:symbol-registry-04624}\(\Gamma/k\) & \texttt{\detokenize{\Gamma/}\allowbreak\detokenize{k}} & source declaration-type lookup relocated to current scope & \hyperref[sec:p07-ch88]{\ref*{sec:p07-ch88}} & n/a & \hyperref[def:p07-ch88-index-operations]{\ref*{def:p07-ch88-index-operations}} & \thepage & Equals add(k,A\_k) with corrected relocation; not arithmetic division. \\
\phantomsection\label{entry:symbol-registry-04625}\([N/x]Q,\operatorname{sub}_i(N,Q)\) & \texttt{\detokenize{[N/}\allowbreak\detokenize{x]Q,}\allowbreak\detokenize{\operatorname{sub}_i}} & binder-removal object substitution & \hyperref[sec:p07-ch88]{\ref*{sec:p07-ch88}} & n/a & \hyperref[def:p07-ch88-index-operations]{\ref*{def:p07-ch88-index-operations}} & \thepage & Inserted object argument is relocated under inner binders and images are not recursively substituted. \\
\phantomsection\label{entry:symbol-registry-04626}\(\Gamma\vdash\Delta,\Gamma\vdash M:P\) & \texttt{\detokenize{\Gamma\vdash\Delta,}\allowbreak\detokenize{\Gamma\vdash M:P}} & source validity and object typing & \hyperref[sec:p07-ch88]{\ref*{sec:p07-ch88}} & n/a & \hyperref[def:p07-ch88-judgments]{\ref*{def:p07-ch88-judgments}} & \thepage & Logical-context validity is not a typing of star; exact full and restricted rule regimes are stated locally. \\
\phantomsection\label{entry:symbol-registry-04627}\(\Gamma\vdash M\cong N\) & \texttt{\detokenize{\Gamma\vdash M\cong N}} & source typed conversion judgment & \hyperref[sec:p07-ch88]{\ref*{sec:p07-ch88}} & n/a & \hyperref[def:p07-ch88-judgments]{\ref*{def:p07-ch88-judgments}} & \thepage & Corrected restricted logical guards and actual full replacements remain distinct. \\
\phantomsection\label{entry:symbol-registry-04628}\(\longrightarrow_\beta,\longrightarrow_\beta^+,\longrightarrow_\beta^*\) & \texttt{\detokenize{\longrightarrow_\beta}} & full indexed annotated beta and its path closures & \hyperref[sec:p07-ch88]{\ref*{sec:p07-ch88}} & n/a & \hyperref[def:p07-ch88-beta]{\ref*{def:p07-ch88-beta}} & \thepage & Every binder field is retained; separate from either typed conversion and from directed logical reduction. \\
\phantomsection\label{entry:symbol-registry-04629}\(\Gamma\vdash M\mathrel{\triangleright}N,\delta(M)\) & \texttt{\detokenize{\triangleright,}\allowbreak\detokenize{\delta}} & source directed logical reduction and logical rank & \hyperref[sec:p07-ch88]{\ref*{sec:p07-ch88}} & n/a & \hyperref[def:p07-ch88-logical-reduction]{\ref*{def:p07-ch88-logical-reduction}} & \thepage & The directed relation is transitive but not reflexive; it is not exactly one full beta step. \\
\phantomsection\label{entry:symbol-registry-04630}\(T_{\vec z},\widehat\Gamma\) & \texttt{\detokenize{T_{\vec z},}\allowbreak\detokenize{\widehat\Gamma}} & editorial indexed-to-named term and environment encoding & \hyperref[sec:p07-ch88]{\ref*{sec:p07-ch88}} & n/a & \hyperref[def:p07-ch88-named-encoding]{\ref*{def:p07-ch88-named-encoding}} & \thepage & Preserves both binder fields; source validity is encoded at target Box rather than by inventing a source Box. \\
\phantomsection\label{entry:symbol-registry-04631}\(\ast,\Box\) & \texttt{\detokenize{\ast,}\allowbreak\detokenize{\Box}} & target CC sorts in the declared translation only & \hyperref[sec:p07-ch88]{\ref*{sec:p07-ch88}} & n/a & \hyperref[def:p07-ch88-named-encoding]{\ref*{def:p07-ch88-named-encoding}} & \thepage & The target star has sort Box; the original one-star presentation instead has a separate validity judgment and no Box constant. \\
\phantomsection\label{entry:symbol-registry-04632}\(\alpha_\Gamma,j_\Gamma,\nu_\Gamma(M),\rho(M)\) & \texttt{\detokenize{\alpha_\Gamma,}\allowbreak\detokenize{j_\Gamma,}\allowbreak\detokenize{\nu_\Gamma,}\allowbreak\detokenize{\rho}} & source proof-program arity parameter injection erasure and closed interpretation & \hyperref[sec:p07-ch88]{\ref*{sec:p07-ch88}} & n/a & \hyperref[def:p07-ch88-stripping]{\ref*{def:p07-ch88-stripping}} & \thepage & Defined only for proof/functionality subjects with an object type; rho names the empty-environment pure image only; this erasure is not the annotation-retaining normalization bridge. \\
\phantomsection\label{entry:symbol-registry-04821}\(\Omega\) & \texttt{\detokenize{\Omega}} & inert marker of the original closed program model & \hyperref[sec:p07-ch88-realization]{\ref*{sec:p07-ch88-realization}} & n/a & \hyperref[def:p07-ch88-native-saturation]{\ref*{def:p07-ch88-native-saturation}} & \thepage & No reduction rule; not divergent self-application. \\
\phantomsection\label{entry:symbol-registry-04822}\(\mathcal Y\) & \texttt{\detokenize{\mathcal Y}} & all closed pure lambda programs with an inert marker & \hyperref[sec:p07-ch88-realization]{\ref*{sec:p07-ch88-realization}} & n/a & \hyperref[def:p07-ch88-native-saturation]{\ref*{def:p07-ch88-native-saturation}} & \thepage & Ambient programs may be non-normalizing; realization restricts their membership. \\
\phantomsection\label{entry:symbol-registry-04823}\(\mathcal U\) & \texttt{\detokenize{\mathcal U}} & source head-expansion saturated program subsets & \hyperref[sec:p07-ch88-realization]{\ref*{sec:p07-ch88-realization}} & n/a & \hyperref[def:p07-ch88-native-saturation]{\ref*{def:p07-ch88-native-saturation}} & \thepage & Not a universe sort or reduction-closed neutral candidate collection. \\
\phantomsection\label{entry:symbol-registry-04824}\(\Pi(A,F)\) & \texttt{\detokenize{\Pi(A,}\allowbreak\detokenize{F)}} & native dependent saturated product & \hyperref[sec:p07-ch88-realization]{\ref*{sec:p07-ch88-realization}} & n/a & \hyperref[def:p07-ch88-native-saturation]{\ref*{def:p07-ch88-native-saturation}} & \thepage & A saturated; F total from closed programs to saturated sets; index by the actual argument. \\
\phantomsection\label{entry:symbol-registry-04825}\(\varphi_q(M),\varphi_{q,\mathrm{env}}(\Gamma)\) & \texttt{\detokenize{\varphi_q,}\allowbreak\detokenize{\varphi_{q,\mathrm{env}}}} & functionality and separately qualified ambient Cartesian environment & \hyperref[sec:p07-ch88-realization]{\ref*{sec:p07-ch88-realization}} & n/a & \hyperref[def:p07-ch88-native-functionality]{\ref*{def:p07-ch88-native-functionality}} & \thepage & The full object-domain function space requires beta invariance; the restricted one does not. \\
\phantomsection\label{entry:symbol-registry-04826}\(\rho_{q,\Gamma}(M)(e)\) & \texttt{\detokenize{\rho_{q,\Gamma}}} & native constructor or stripped-program value in its actual typed mode & \hyperref[sec:p07-ch88-realization]{\ref*{sec:p07-ch88-realization}} & n/a & \hyperref[def:p07-ch88-native-values]{\ref*{def:p07-ch88-native-values}} & \thepage & Not annotation-retaining E; constructor equality is literal while full proof equality is beta. \\
\phantomsection\label{entry:symbol-registry-04827}\(D_q(\Gamma)\) & \texttt{\detokenize{D_q}} & dependent realized subenvironment & \hyperref[sec:p07-ch88-realization]{\ref*{sec:p07-ch88-realization}} & n/a & \hyperref[def:p07-ch88-native-realizers]{\ref*{def:p07-ch88-native-realizers}} & \thepage & Object coordinates inhabit their older-prefix type interpretation; constructor coordinates range over functionality. \\
\phantomsection\label{entry:symbol-registry-04899}\(\vdash_n\) & \texttt{\detokenize{\vdash_n}} & separate p112 normal-type variant validity and typing & \hyperref[sec:p07-ch88-checking]{\ref*{sec:p07-ch88-checking}} & n/a & \hyperref[def:p07-ch88-check-normal-variant]{\ref*{def:p07-ch88-check-normal-variant}} & \thepage & No equality/conversion; not raw-beta subject-preserving. \\
\phantomsection\label{entry:symbol-registry-04900}\(\mathcal N(E)\) & \texttt{\detokenize{\mathcal N}} & partial full annotated source beta normalizer & \hyperref[sec:p07-ch88-checking]{\ref*{sec:p07-ch88-checking}} & n/a & \hyperref[def:p07-ch88-check-normal-variant]{\ref*{def:p07-ch88-check-normal-variant}} & \thepage & Termination only on certified typed images; not a total raw or erased-program normalizer. \\
\phantomsection\label{entry:symbol-registry-04901}\(\mathrm{Sort},\mathrm{Valid},\mathrm{Infer}\) & \texttt{\detokenize{\mathrm{Sort},}\allowbreak\detokenize{\mathrm{Valid},}\allowbreak\detokenize{\mathrm{Infer}}} & structural operations for the actual normal-type checker & \hyperref[sec:p07-ch88-checking]{\ref*{sec:p07-ch88-checking}} & n/a & \hyperref[def:p07-ch88-check-algorithm]{\ref*{def:p07-ch88-check-algorithm}} & \thepage & Oldest-first context validation; proper original syntax recursion; normalize only certified synthesized types. \\
\phantomsection\label{entry:symbol-registry-05076}\(\eta_\beta(E)\) & \texttt{\detokenize{\eta_\beta}} & maximum full compatible beta path length of a strongly normalizing scoped term & \hyperref[sec:p07-ch88-object-product]{\ref*{sec:p07-ch88-object-product}} & n/a & \hyperref[def:p07-ch88-object-height]{\ref*{def:p07-ch88-object-height}} & \thepage & Includes annotations and bodies; local endpoint decrease does not bound commuted intermediate terms. \\
\phantomsection\label{entry:symbol-registry-05077}\(\operatorname{ht}(d)\) & \texttt{\detokenize{\operatorname{ht}}} & height of an original supplied finite source proof tree & \hyperref[sec:p07-ch88-object-product]{\ref*{sec:p07-ch88-object-product}} & n/a & \hyperref[def:p07-ch88-object-height]{\ref*{def:p07-ch88-object-height}} & \thepage & Certified transformation recurses only on original proper subproofs; never on inflated generated proofs. \\
\phantomsection\label{entry:symbol-registry-05466}\(V,E,s,t\) & \texttt{\detokenize{V,}\allowbreak\detokenize{E,}\allowbreak\detokenize{s,}\allowbreak\detokenize{t}} & finite partial incidence graph data & P7 multiplicative-proof-nets & n/a & \hyperref[def:p07-mll-partial-graph]{\ref*{def:p07-mll-partial-graph}} & \thepage & Vertices arcs and partial source/target maps; arcs are formula occurrences. \\
\phantomsection\label{entry:symbol-registry-05467}\(G_\sigma(N)\) & \texttt{\detokenize{G_\sigma(N)}} & undirected switched incidence graph & P7 multiplicative-proof-nets & n/a & \hyperref[def:p07-mll-switching]{\ref*{def:p07-mll-switching}} & \thepage & Par premise incidence severed; occurrence arc retained with fresh boundary. \\
\phantomsection\label{entry:symbol-registry-05468}\(D(\pi)\) & \texttt{\detokenize{D(\pi)}} & desequentialization map & P7 multiplicative-proof-nets & n/a & \hyperref[def:p07-mll-desequentialization]{\ref*{def:p07-mll-desequentialization}} & \thepage & Construction of incidence graph from the exact selected proof. \\
\phantomsection\label{entry:symbol-registry-05469}\(F^+,F^-\) & \texttt{\detokenize{F^+,}\allowbreak\detokenize{F^-}} & original Girard exit port states & P7 multiplicative-proof-nets & n/a & \hyperref[def:p07-mll-girard-trip]{\ref*{def:p07-mll-girard-trip}} & \thepage & Signs label upward/downward ports; neither is linear negation or focusing polarity. \\
\phantomsection\label{entry:symbol-registry-05470}\(eA\) & \texttt{\detokenize{eA}} & Girard tensor-relative empire & P7 multiplicative-proof-nets & n/a & \hyperref[def:p07-mll-girard-empire]{\ref*{def:p07-mll-girard-empire}} & \thepage & Intersection of both-gate trip intervals over all original settings; not a single switching component. \\
\phantomsection\label{entry:symbol-registry-05471}\(M_v\) & \texttt{\detokenize{M_v}} & minimal-bridge cycle set & P7 multiplicative-proof-nets & n/a & \hyperref[def:p07-mll-strong-path]{\ref*{def:p07-mll-strong-path}} & \thepage & Finite simple cycles based at v containing its conclusion and minimizing bridges. \\
\phantomsection\label{entry:symbol-registry-05472}\(\prec\) & \texttt{\detokenize{\prec}} & first bridge dependency on par vertices & P7 multiplicative-proof-nets & n/a & \hyperref[def:p07-mll-strong-path]{\ref*{def:p07-mll-strong-path}} & \thepage & Strictness is proved before maximal vertices are used. \\
\phantomsection\label{entry:symbol-registry-05473}\(d(N)\) & \texttt{\detokenize{d(N)}} & switching component count & P7 multiplicative-proof-nets & n/a & \hyperref[thm:p07-mll-desequentialization]{\ref*{thm:p07-mll-desequentialization}} & \thepage & Invariant on proved auxiliary proof images; empty graph has zero components. \\
\phantomsection\label{entry:symbol-registry-05474}\(\mathsf{cut}\) & \texttt{\detokenize{\mathsf{cut}}} & incidence cut vertex or local graph surgery pattern & P7 multiplicative-proof-nets & n/a & \hyperref[def:p07-mll-net-cut]{\ref*{def:p07-mll-net-cut}} & \thepage & Two dual premise arcs and no conclusion; pattern notation is not textual formula substitution. \\
\phantomsection\label{entry:symbol-registry-05623}\(I_\sigma(X)\) & \texttt{\detokenize{I_\sigma(X)}} & closed forward original-trip gate interval & P7 girard-1987-trips-and-empires & n/a & \hyperref[def:p07-girard-empire-interval]{\ref*{def:p07-girard-empire-interval}} & \thepage & From the plus gate to the minus gate in one full original trip; not an empire by itself. \\
\phantomsection\label{entry:symbol-registry-05624}\(\widehat E\) & \texttt{\detokenize{\widehat E}} & both-gate expansion of an occurrence set & P7 girard-1987-trips-and-empires & n/a & \hyperref[def:p07-girard-empire-interval]{\ref*{def:p07-girard-empire-interval}} & \thepage & Contains exactly both ports of each occurrence in E; set type differs from E. \\
\phantomsection\label{entry:symbol-registry-05625}\(h,g_L,g_R\) & \texttt{\detokenize{h,}\allowbreak\detokenize{g_L,}\allowbreak\detokenize{g_R}} & three-path endpoint and reconnection permutations & P7 girard-1987-trips-and-empires & n/a & \hyperref[thm:p07-girard-empire-local-blocks]{\ref*{thm:p07-girard-empire-local-blocks}} & \thepage & Permutations of three fixed path names; composition cycles count the reconnected trip orbits. \\
\phantomsection\label{entry:symbol-registry-05626}\(\dot\cup\) & \texttt{\detokenize{\dot\cup}} & disjoint union of occurrence sets & P7 girard-1987-trips-and-empires & n/a & \hyperref[thm:p07-girard-empire-joint]{\ref*{thm:p07-girard-empire-joint}} & \thepage & Disjointness is established by original tensor block intervals and root exclusion. \\
\phantomsection\label{entry:symbol-registry-05650}\(\equiv\) & \texttt{\detokenize{\equiv}} & Source judgmental equality & P7 HoTT source theory & relation & \hyperref[def:p07-hott-core]{\ref*{def:p07-hott-core}} & \thepage & Genuine source role; distinct from an identity type and from elementary equivalence elsewhere Blank macro filled with exact unchanged glyph spelling; registry data only, not a new TeX command. \\
\phantomsection\label{entry:symbol-registry-05651}\(x=_A y\) & \texttt{\detokenize{x=_A y}} & Identity type of identifications & P7 HoTT source theory & relation & \hyperref[def:p07-hott-core]{\ref*{def:p07-hott-core}} & \thepage & Typed by A; not equality reflection or proof irrelevance Blank macro filled with exact unchanged glyph spelling; registry data only, not a new TeX command. \\
\phantomsection\label{entry:symbol-registry-05652}\(p\cdot q\) & \texttt{\detokenize{p\cdot q}} & Path concatenation & P7 HoTT paths & n/a & \hyperref[def:p07-hott-path-operations]{\ref*{def:p07-hott-path-operations}} & \thepage & Traverse p then q; distinct from function composition and multiplication Blank macro filled with exact unchanged glyph spelling; registry data only, not a new TeX command. Original precedence 'operator' retained as semantic descriptor; n/a is not an invented rank. \\
\phantomsection\label{entry:symbol-registry-05653}\(p^{-1}\) & \texttt{\detokenize{p^{-1}}} & Inverse identification & P7 HoTT paths & n/a & \hyperref[def:p07-hott-path-operations]{\ref*{def:p07-hott-path-operations}} & \thepage & J-defined with reflexivity computation Blank macro filled with exact unchanged glyph spelling; registry data only, not a new TeX command. Original precedence 'operator' retained as semantic descriptor; n/a is not an invented rank. \\
\phantomsection\label{entry:symbol-registry-05654}\(\Omega(A,a)\) & \texttt{\detokenize{\Omega(A,}\allowbreak\detokenize{a)}} & Pointed loop type & P7 HoTT paths & object & \hyperref[def:p07-hott-path-operations]{\ref*{def:p07-hott-path-operations}} & \thepage & Higher loops iterate identity at the reflexivity basepoint Blank macro filled with exact unchanged glyph spelling; registry data only, not a new TeX command. \\
\phantomsection\label{entry:symbol-registry-05655}\(\mathsf{transport}^{P}(p,u)\) & \texttt{\detokenize{\mathsf{transport}^{P}(p,}\allowbreak\detokenize{u)}} & Dependent transport & P7 HoTT transport & n/a & \hyperref[def:p07-hott-transport]{\ref*{def:p07-hott-transport}} & \thepage & Target P(y); no cross-fiber ill-typed equality Blank macro filled with exact unchanged glyph spelling; registry data only, not a new TeX command. Original precedence 'function' retained as semantic descriptor; n/a is not an invented rank. \\
\phantomsection\label{entry:symbol-registry-05656}\(\mathsf{ap}_f(p)\) & \texttt{\detokenize{\mathsf{ap}_f(p)}} & Function action on paths & P7 HoTT transport & n/a & \hyperref[def:p07-hott-transport]{\ref*{def:p07-hott-transport}} & \thepage & Source-typed ordinary action Blank macro filled with exact unchanged glyph spelling; registry data only, not a new TeX command. Original precedence 'function' retained as semantic descriptor; n/a is not an invented rank. \\
\phantomsection\label{entry:symbol-registry-05657}\(\mathsf{apd}_f(p)\) & \texttt{\detokenize{\mathsf{apd}_f(p)}} & Dependent function action & P7 HoTT transport & n/a & \hyperref[def:p07-hott-transport]{\ref*{def:p07-hott-transport}} & \thepage & Path from transported value to target value Blank macro filled with exact unchanged glyph spelling; registry data only, not a new TeX command. Original precedence 'function' retained as semantic descriptor; n/a is not an invented rank. \\
\phantomsection\label{entry:symbol-registry-05658}\(f\sim g\) & \texttt{\detokenize{f\sim g}} & Pointwise homotopy type & P7 HoTT homotopies & relation & \hyperref[def:p07-hott-homotopy]{\ref*{def:p07-hott-homotopy}} & \thepage & Not automatically a path of functions Blank macro filled with exact unchanged glyph spelling; registry data only, not a new TeX command. \\
\phantomsection\label{entry:symbol-registry-05659}\(\mathsf{qinv}(f)\) & \texttt{\detokenize{\mathsf{qinv}(f)}} & Raw quasi-inverse data & P7 HoTT equivalences & object & \hyperref[def:p07-hott-equivalence]{\ref*{def:p07-hott-equivalence}} & \thepage & Not silently a mere proposition Blank macro filled with exact unchanged glyph spelling; registry data only, not a new TeX command. \\
\phantomsection\label{entry:symbol-registry-05660}\(\mathsf{ishae}(f)\) & \texttt{\detokenize{\mathsf{ishae}(f)}} & Half-adjoint equivalence data & P7 HoTT equivalences & object & \hyperref[def:p07-hott-equivalence]{\ref*{def:p07-hott-equivalence}} & \thepage & Includes an actual triangle path Blank macro filled with exact unchanged glyph spelling; registry data only, not a new TeX command. \\
\phantomsection\label{entry:symbol-registry-05661}\(\mathsf{linv}(f),\mathsf{rinv}(f)\) & \texttt{\detokenize{\mathsf{linv}(f),}\allowbreak\detokenize{\mathsf{rinv}(f)}} & Left and right inverse types & P7 HoTT equivalences & object & \hyperref[def:p07-hott-equivalence]{\ref*{def:p07-hott-equivalence}} & \thepage & Distinct inverse functions are permitted Blank macro filled with exact unchanged glyph spelling; registry data only, not a new TeX command. \\
\phantomsection\label{entry:symbol-registry-05662}\(\mathsf{biinv}(f)\) & \texttt{\detokenize{\mathsf{biinv}(f)}} & Bi-invertibility type & P7 HoTT equivalences & object & \hyperref[def:p07-hott-equivalence]{\ref*{def:p07-hott-equivalence}} & \thepage & Product of left and right inverse data Blank macro filled with exact unchanged glyph spelling; registry data only, not a new TeX command. \\
\phantomsection\label{entry:symbol-registry-05663}\(\mathsf{isequiv}(f)\) & \texttt{\detokenize{\mathsf{isequiv}(f)}} & Source chosen equivalence predicate & P7 HoTT equivalences & object & \hyperref[def:p07-hott-equivalence]{\ref*{def:p07-hott-equivalence}} & \thepage & Source Section4.5 chooses ishae; bridges are proved separately Blank macro filled with exact unchanged glyph spelling; registry data only, not a new TeX command. \\
\phantomsection\label{entry:symbol-registry-05664}\(A\simeq B\) & \texttt{\detokenize{A\simeq B}} & Type of specified equivalences & P7 HoTT equivalences & relation & \hyperref[def:p07-hott-equivalence]{\ref*{def:p07-hott-equivalence}} & \thepage & Not merely truncated existence of an equivalence Blank macro filled with exact unchanged glyph spelling; registry data only, not a new TeX command. \\
\phantomsection\label{entry:symbol-registry-05665}\(\mathsf{isContr}(A)\) & \texttt{\detokenize{\mathsf{isContr}(A)}} & Type-level contraction data & P7 HoTT contractibility & object & \hyperref[def:p07-hott-contractibility]{\ref*{def:p07-hott-contractibility}} & \thepage & Same source operator has a distinct map-typed use Blank macro filled with exact unchanged glyph spelling; registry data only, not a new TeX command. \\
\phantomsection\label{entry:symbol-registry-05666}\(\mathsf{isContr}(f)\) & \texttt{\detokenize{\mathsf{isContr}(f)}} & Contractibility of all map fibers & P7 HoTT map contractibility & object & \hyperref[def:p07-hott-fiber]{\ref*{def:p07-hott-fiber}} & \thepage & Source overload retained; f is a map not a type Blank macro filled with exact unchanged glyph spelling; registry data only, not a new TeX command. \\
\phantomsection\label{entry:symbol-registry-05667}\(\mathsf{isProp}(A),\mathsf{isSet}(A)\) & \texttt{\detokenize{\mathsf{isProp}(A),}\allowbreak\detokenize{\mathsf{isSet}(A)}} & Mere-proposition and set predicates & P7 HoTT levels & object & \hyperref[def:p07-hott-contractibility]{\ref*{def:p07-hott-contractibility}} & \thepage & Their proof standing is established under declared assumptions Blank macro filled with exact unchanged glyph spelling; registry data only, not a new TeX command. \\
\phantomsection\label{entry:symbol-registry-05668}\(\mathsf{fib}_f(b)\) & \texttt{\detokenize{\mathsf{fib}_f(b)}} & Homotopy fiber & P7 HoTT fibers & object & \hyperref[def:p07-hott-fiber]{\ref*{def:p07-hott-fiber}} & \thepage & Includes endpoint path evidence Blank macro filled with exact unchanged glyph spelling; registry data only, not a new TeX command. \\
\phantomsection\label{entry:symbol-registry-05669}\(\mathsf{happly},\mathsf{funext}\) & \texttt{\detokenize{\mathsf{happly},}\allowbreak\detokenize{\mathsf{funext}}} & Function-path comparison and assumed inverse & P7 HoTT function extensionality & n/a & \hyperref[def:p07-hott-funext]{\ref*{def:p07-hott-funext}} & \thepage & Inverse comparisons are propositional; funext is not derived from J Blank macro filled with exact unchanged glyph spelling; registry data only, not a new TeX command. Original precedence 'function' retained as semantic descriptor; n/a is not an invented rank. \\
\phantomsection\label{entry:symbol-registry-05670}\(\mathsf{idtoeqv},\mathsf{ua}\) & \texttt{\detokenize{\mathsf{idtoeqv},}\allowbreak\detokenize{\mathsf{ua}}} & Universe-path comparison and univalent inverse & P7 HoTT univalence & n/a & \hyperref[def:p07-hott-univalence]{\ref*{def:p07-hott-univalence}} & \thepage & Particular source map; no judgmental ua computation Blank macro filled with exact unchanged glyph spelling; registry data only, not a new TeX command. Original precedence 'function' retained as semantic descriptor; n/a is not an invented rank. \\
\phantomsection\label{entry:symbol-registry-05671}\(\|A\|\) & \texttt{\detokenize{\|A\|}} & Assumed propositional truncation & P7 HoTT truncation interface & object & \hyperref[def:p07-hott-propositional-truncation]{\ref*{def:p07-hott-propositional-truncation}} & \thepage & Construction belongs to separate HIT scope; elimination only into mere propositions Blank macro filled with exact unchanged glyph spelling; registry data only, not a new TeX command. \\
\phantomsection\label{entry:symbol-registry-05672}\(\mathsf{is}\mbox{-}n\mbox{-}\mathsf{type}(A)\) & \texttt{\detokenize{\mathsf{is}\mbox{-}n\mbox{-}\mathsf{type}(A)}} & Source recursive homotopy-level predicate & P7 HoTT levels & object & \hyperref[def:p07-hott-levels]{\ref*{def:p07-hott-levels}} & \thepage & Index begins at minus2; not a truncation constructor Blank macro filled with exact unchanged glyph spelling; registry data only, not a new TeX command. \\
\phantomsection\label{entry:symbol-registry-05673}\(\equiv_{\mathrm j}\) & \texttt{\detokenize{\equiv_{\mathrm j}}} & judgmental equality & P7 HIT principles & n/a & \hyperref[def:p07-hit-standing]{\ref*{def:p07-hit-standing}} & \thepage & Distinct from the identity type = and from elementary equivalence. \\
\phantomsection\label{entry:symbol-registry-05674}\(\operatorname{DPath}_P\) & \texttt{\detokenize{\operatorname{DPath}_P}} & transport-based dependent path type & P7 HIT principles & n/a & \hyperref[def:p07-hit-dependent-paths]{\ref*{def:p07-hit-dependent-paths}} & \thepage & Typed endpoints u:P(x) and v:P(y); not a path solely in the source fiber. \\
\phantomsection\label{entry:symbol-registry-05675}\(\operatorname{tr}^2_P\) & \texttt{\detokenize{\operatorname{tr}^2_P}} & action of a 2-path on transport & P7 HIT principles & n/a & \hyperref[def:p07-hit-dependent-paths]{\ref*{def:p07-hit-dependent-paths}} & \thepage & A path between two transported values in P(y). \\
\phantomsection\label{entry:symbol-registry-05676}\(\operatorname{DPath}^2_P\) & \texttt{\detokenize{\operatorname{DPath}^2_P}} & dependent 2-path type & P7 HIT principles & n/a & \hyperref[def:p07-hit-dependent-paths]{\ref*{def:p07-hit-dependent-paths}} & \thepage & Its two sides are parallel paths in the target fiber. \\
\phantomsection\label{entry:symbol-registry-05677}\(\gamma_P\) & \texttt{\detokenize{\gamma_P}} & propositional transport-composition comparison & P7 HIT path calculus & n/a & \hyperref[thm:p07-hit-path-calculus]{\ref*{thm:p07-hit-path-calculus}} & \thepage & Derived by identity elimination; not judgmental at arbitrary p and q. \\
\phantomsection\label{entry:symbol-registry-05678}\(\star\) & \texttt{\detokenize{\star}} & dependent-path composition & P7 HIT path calculus & n/a & \hyperref[thm:p07-hit-path-calculus]{\ref*{thm:p07-hit-path-calculus}} & \thepage & Uses transport composition then action on the first lift then the second lift. \\
\phantomsection\label{entry:symbol-registry-05679}\(\mu_f\) & \texttt{\detokenize{\mu_f}} & dependent-action composition comparison & P7 HIT path calculus & n/a & \hyperref[thm:p07-hit-path-calculus]{\ref*{thm:p07-hit-path-calculus}} & \thepage & Full path witness; no judgmental preservation assertion. \\
\phantomsection\label{entry:symbol-registry-05680}\(S^1\) & \texttt{\detokenize{S^1}} & circle with basepoint and generating loop & P7 HIT circle & n/a & \hyperref[def:p07-hit-circle]{\ref*{def:p07-hit-circle}} & \thepage & Added HIT principle; its paths retain ordinary identity elimination. \\
\phantomsection\label{entry:symbol-registry-05681}\(\mathrm{base}\) & \texttt{\detokenize{\mathrm{base}}} & circle or sphere point constructor & P7 HIT circle and spheres & n/a & \hyperref[def:p07-hit-circle]{\ref*{def:p07-hit-circle}} & \thepage & Each local type identifies which basepoint this source notation denotes. \\
\phantomsection\label{entry:symbol-registry-05682}\(\mathrm{loop}\) & \texttt{\detokenize{\mathrm{loop}}} & circle path constructor & P7 HIT circle & n/a & \hyperref[def:p07-hit-circle]{\ref*{def:p07-hit-circle}} & \thepage & Not judgmentally reflexivity; path computation is propositional. \\
\phantomsection\label{entry:symbol-registry-05683}\(I\) & \texttt{\detokenize{I}} & higher-inductive interval type & P7 HIT interval & n/a & \hyperref[def:p07-hit-interval]{\ref*{def:p07-hit-interval}} & \thepage & Not a cubical dimension object or face-formula algebra. \\
\phantomsection\label{entry:symbol-registry-05684}\(\mathrm{seg}\) & \texttt{\detokenize{\mathrm{seg}}} & interval endpoint identity constructor & P7 HIT interval & n/a & \hyperref[def:p07-hit-interval]{\ref*{def:p07-hit-interval}} & \thepage & Endpoint point computations are judgmental; its action computation is propositional. \\
\phantomsection\label{entry:symbol-registry-05685}\(\Sigma A\) & \texttt{\detokenize{\Sigma A}} & suspension of A & P7 HIT suspension & n/a & \hyperref[def:p07-hit-suspension]{\ref*{def:p07-hit-suspension}} & \thepage & Published notation distinct from a dependent-sum expression with binder. \\
\phantomsection\label{entry:symbol-registry-05686}\(\mathrm{merid}\) & \texttt{\detokenize{\mathrm{merid}}} & A-indexed suspension path constructor & P7 HIT suspension & n/a & \hyperref[def:p07-hit-suspension]{\ref*{def:p07-hit-suspension}} & \thepage & All paths in A also act on the constructor family; not merely component indexing. \\
\phantomsection\label{entry:symbol-registry-05687}\(\mathrm{surf}\) & \texttt{\detokenize{\mathrm{surf}}} & direct 2-sphere 2-path constructor & P7 HIT spheres & n/a & \hyperref[def:p07-hit-spheres]{\ref*{def:p07-hit-spheres}} & \thepage & Dependent computation uses the declared dependent 2-path convention. \\
\phantomsection\label{entry:symbol-registry-05688}\(\mathrm{glue}\) & \texttt{\detokenize{\mathrm{glue}}} & pushout path constructor & P7 HIT pushout & n/a & \hyperref[def:p07-hit-pushout]{\ref*{def:p07-hit-pushout}} & \thepage & Indexed by the actual span input with typed endpoint images. \\
\phantomsection\label{entry:symbol-registry-05689}\(\|A\|\) & \texttt{\detokenize{\|A\|}} & propositional truncation & P7 HIT truncations & n/a & \hyperref[def:p07-hit-truncations]{\ref*{def:p07-hit-truncations}} & \thepage & Recursive squash constructor; elimination restricted as stated. \\
\phantomsection\label{entry:symbol-registry-05690}\(\|A\|_0\) & \texttt{\detokenize{\|A\|_0}} & set truncation & P7 HIT truncations & n/a & \hyperref[def:p07-hit-truncations]{\ref*{def:p07-hit-truncations}} & \thepage & Direct constructor identifies parallel paths by a 2-path. \\
\phantomsection\label{entry:symbol-registry-05691}\(\mathrm{squash}\) & \texttt{\detokenize{\mathrm{squash}}} & propositional-truncation path constructor & P7 HIT truncations & n/a & \hyperref[def:p07-hit-truncations]{\ref*{def:p07-hit-truncations}} & \thepage & Its parameters range over the type being generated. \\
\phantomsection\label{entry:symbol-registry-05692}\(A/R\) & \texttt{\detokenize{A/}\allowbreak\detokenize{R}} & set-quotient HIT & P7 HIT quotient & n/a & \hyperref[def:p07-hit-set-quotient]{\ref*{def:p07-hit-set-quotient}} & \thepage & Mere relation need not be an equivalence; no global representative choice is supplied. \\
\phantomsection\label{entry:symbol-registry-05693}\(q\) & \texttt{\detokenize{q}} & set-quotient point projection & P7 HIT quotient & n/a & \hyperref[def:p07-hit-set-quotient]{\ref*{def:p07-hit-set-quotient}} & \thepage & Mere surjectivity is bracketed rather than a section of q. \\
\phantomsection\label{entry:symbol-registry-05694}\(\|A\|_n\) & \texttt{\detokenize{\|A\|_n}} & general hub-and-spoke truncation construction & P7 HIT general truncation & n/a & \hyperref[def:p07-hit-hub-truncation]{\ref*{def:p07-hit-hub-truncation}} & \thepage & The remaining sphere-adjunction and level-characterization proof duty is explicit. \\
\phantomsection\label{entry:symbol-registry-05695}\(\mathrm{hub}\) & \texttt{\detokenize{\mathrm{hub}}} & recursive hub point constructor & P7 HIT general truncation & n/a & \hyperref[def:p07-hit-hub-truncation]{\ref*{def:p07-hit-hub-truncation}} & \thepage & Argument is a sphere map into the type being generated. \\
\phantomsection\label{entry:symbol-registry-05696}\(\mathrm{spoke}\) & \texttt{\detokenize{\mathrm{spoke}}} & sphere-indexed path to the hub & P7 HIT general truncation & n/a & \hyperref[def:p07-hit-hub-truncation]{\ref*{def:p07-hit-hub-truncation}} & \thepage & Dependent raw elimination includes lifted recursive input and hub value. \\
\phantomsection\label{entry:symbol-registry-05697}I & \texttt{\detokenize{\mathbb I}} & Dimension sort and free de Morgan algebra & P07 CCHM2016 v1 & n/a & \hyperref[sec:p07-cchm2016-dimensions]{\ref*{sec:p07-cchm2016-dimensions}} & \thepage & Not an ordinary type or Boolean algebra; no new global macro Original precedence 'sort' retained as semantic descriptor; n/a is not an invented rank. \\
\phantomsection\label{entry:symbol-registry-05698}F & \texttt{\detokenize{\mathbb F}} & Face lattice & P07 CCHM2016 v1 & n/a & \hyperref[sec:p07-cchm2016-dimensions]{\ref*{sec:p07-cchm2016-dimensions}} & \thepage & Distinct zero and one from interval; no face complement Original precedence 'sort' retained as semantic descriptor; n/a is not an invented rank. \\
\phantomsection\label{entry:symbol-registry-05699}i j k & \texttt{\detokenize{literal names}} & Dimension names & P07 CCHM2016 v1 & n/a & \hyperref[sec:p07-cchm2016-dimensions]{\ref*{sec:p07-cchm2016-dimensions}} & \thepage & Freshness checked at each dimension binder Original precedence 'atomic' retained as semantic descriptor; n/a is not an invented rank. \\
\phantomsection\label{entry:symbol-registry-05700}1-r & \texttt{\detokenize{literal notation}} & De Morgan involution & P07 CCHM2016 v1 & n/a & \hyperref[sec:p07-cchm2016-dimensions]{\ref*{sec:p07-cchm2016-dimensions}} & \thepage & Not subtraction of ordinary terms Original precedence 'unary' retained as semantic descriptor; n/a is not an invented rank. \\
\phantomsection\label{entry:symbol-registry-05701}(i=0) (i=1) & \texttt{\detokenize{literal notation}} & Face generators & P07 CCHM2016 v1 & n/a & \hyperref[sec:p07-cchm2016-dimensions]{\ref*{sec:p07-cchm2016-dimensions}} & \thepage & Meet of opposite generators is zero; their join is not top Original precedence 'atomic' retained as semantic descriptor; n/a is not an invented rank. \\
\phantomsection\label{entry:symbol-registry-05702}(i/r) (x/a) & \texttt{\detokenize{literal notation}} & Source-direction capture-avoiding substitution & P07 CCHM2016 v1 & n/a & \hyperref[sec:p07-cchm2016-dimensions]{\ref*{sec:p07-cchm2016-dimensions}} & \thepage & First argument is replaced variable; second is replacement Original precedence 'postfix' retained as semantic descriptor; n/a is not an invented rank. \\
\phantomsection\label{entry:symbol-registry-05703}<i>t & \texttt{\detokenize{\langle i\rangle t}} & Dimension abstraction & P07 CCHM2016 v1 & n/a & \hyperref[sec:p07-cchm2016-paths]{\ref*{sec:p07-cchm2016-paths}} & \thepage & Not ordinary lambda abstraction Original precedence 'binder' retained as semantic descriptor; n/a is not an invented rank. \\
\phantomsection\label{entry:symbol-registry-05704}Path & \texttt{\detokenize{\mathsf{Path}}} & Nondependent path type of the core theory & P07 CCHM2016 v1 & n/a & \hyperref[sec:p07-cchm2016-paths]{\ref*{sec:p07-cchm2016-paths}} & \thepage & Dependent path extension in source Section 9 remains separate Original precedence 'prefix' retained as semantic descriptor; n/a is not an invented rank. \\
\phantomsection\label{entry:symbol-registry-05705}{}[phi->u] & \texttt{\detokenize{literal notation}} & Tube with stated extent or strict boundary abbreviation & P07 CCHM2016 v1 & n/a & \hyperref[sec:p07-cchm2016-composition]{\ref*{sec:p07-cchm2016-composition}} & \thepage & Not a new type; system formation requires a cover in its current context Original precedence 'bracketed' retained as semantic descriptor; n/a is not an invented rank. \\
\phantomsection\label{entry:symbol-registry-05706}comp & \texttt{\detokenize{\mathsf{comp}}} & Primitive composition & P07 CCHM2016 v1 & n/a & \hyperref[sec:p07-cchm2016-composition]{\ref*{sec:p07-cchm2016-composition}} & \thepage & Extent independent of bound direction; cap compatibility mandatory Original precedence 'prefix binder' retained as semantic descriptor; n/a is not an invented rank. \\
\phantomsection\label{entry:symbol-registry-05707}fill & \texttt{\detokenize{\mathsf{fill}}} & Derived filling using composition and meet connection & P07 CCHM2016 v1 & n/a & \hyperref[sec:p07-cchm2016-composition]{\ref*{sec:p07-cchm2016-composition}} & \thepage & No additional primitive or constant-transport regularity axiom Original precedence 'prefix binder' retained as semantic descriptor; n/a is not an invented rank. \\
\phantomsection\label{entry:symbol-registry-05708}transp & \texttt{\detokenize{\mathsf{transp}}} & Empty-extent composition & P07 CCHM2016 v1 & n/a & \hyperref[sec:p07-cchm2016-composition]{\ref*{sec:p07-cchm2016-composition}} & \thepage & Not generally judgmentally identity in a constant family Original precedence 'prefix binder' retained as semantic descriptor; n/a is not an invented rank. \\
\phantomsection\label{entry:symbol-registry-05709}isContr & \texttt{\detokenize{\mathsf{isContr}}} & Center plus contraction paths & P07 CCHM2016 v1 & n/a & \hyperref[sec:p07-cchm2016-equivalence]{\ref*{sec:p07-cchm2016-equivalence}} & \thepage & Path orientation from center to arbitrary point Original precedence 'prefix' retained as semantic descriptor; n/a is not an invented rank. \\
\phantomsection\label{entry:symbol-registry-05710}isEquiv & \texttt{\detokenize{\mathsf{isEquiv}}} & Fiber contractibility & P07 CCHM2016 v1 & n/a & \hyperref[sec:p07-cchm2016-equivalence]{\ref*{sec:p07-cchm2016-equivalence}} & \thepage & Fiber paths oriented y to f x Original precedence 'prefix' retained as semantic descriptor; n/a is not an invented rank. \\
\phantomsection\label{entry:symbol-registry-05711}equiv & \texttt{\detokenize{\mathsf{equiv}}} & uniform partial-fiber extension operation & P07 CCHM2016 v1 partial fiber extension & n/a & \hyperref[sec:p07-cchm2016-equivalence]{\ref*{sec:p07-cchm2016-equivalence}} & \thepage & Two source uses disambiguated by superscript and arguments Original precedence 'prefix' retained as semantic descriptor; n/a is not an invented rank. Root-authorized typed decomposition; original grouped intake row remains unchanged as evidence. \\
\phantomsection\label{entry:symbol-registry-05712}\ensuremath{equiv^i} & \texttt{\detokenize{\mathsf{equiv}^{i}}} & superscripted type-line equivalence construction & P07 CCHM2016 v1 type-line equivalence & n/a & \hyperref[sec:p07-cchm2016-universe]{\ref*{sec:p07-cchm2016-universe}} & \thepage & Two source uses disambiguated by superscript and arguments Original precedence 'prefix' retained as semantic descriptor; n/a is not an invented rank. Root-authorized typed decomposition; original grouped intake row remains unchanged as evidence. \\
\phantomsection\label{entry:symbol-registry-05713}pres & \texttt{\detokenize{\mathsf{pres}}} & Composition-preservation path & P07 CCHM2016 v1 & n/a & \hyperref[sec:p07-cchm2016-equivalence]{\ref*{sec:p07-cchm2016-equivalence}} & \thepage & Constant on the original tube extent Original precedence 'prefix binder' retained as semantic descriptor; n/a is not an invented rank. \\
\phantomsection\label{entry:symbol-registry-05714}Glue & \texttt{\detokenize{\mathsf{Glue}}} & gluing type formation operation & P07 CCHM2016 v1 & n/a & \hyperref[sec:p07-cchm2016-glue]{\ref*{sec:p07-cchm2016-glue}} & \thepage & No empty-Glue judgmental collapse Original precedence 'prefix' retained as semantic descriptor; n/a is not an invented rank. Root-authorized typed decomposition; original grouped intake row remains unchanged as evidence. \\
\phantomsection\label{entry:symbol-registry-05715}glue & \texttt{\detokenize{\mathsf{glue}}} & gluing term introduction operation & P07 CCHM2016 v1 & n/a & \hyperref[sec:p07-cchm2016-glue]{\ref*{sec:p07-cchm2016-glue}} & \thepage & No empty-Glue judgmental collapse Original precedence 'prefix' retained as semantic descriptor; n/a is not an invented rank. Root-authorized typed decomposition; original grouped intake row remains unchanged as evidence. \\
\phantomsection\label{entry:symbol-registry-05716}unglue & \texttt{\detokenize{\mathsf{unglue}}} & gluing term elimination operation & P07 CCHM2016 v1 & n/a & \hyperref[sec:p07-cchm2016-glue]{\ref*{sec:p07-cchm2016-glue}} & \thepage & No empty-Glue judgmental collapse Original precedence 'prefix' retained as semantic descriptor; n/a is not an invented rank. Root-authorized typed decomposition; original grouped intake row remains unchanged as evidence. \\
\phantomsection\label{entry:symbol-registry-05717}forall i & \texttt{\detokenize{\forall i}} & Largest direction-independent face below a face & P07 CCHM2016 v1 & n/a & \hyperref[sec:p07-cchm2016-glue]{\ref*{sec:p07-cchm2016-glue}} & \thepage & Face-lattice homomorphism; not ordinary universal quantification Original precedence 'binder' retained as semantic descriptor; n/a is not an invented rank. \\
\phantomsection\label{entry:symbol-registry-05718}U & \texttt{\detokenize{literal U}} & Source Russell universe & P07 CCHM2016 v1 & n/a & \hyperref[sec:p07-cchm2016-universe]{\ref*{sec:p07-cchm2016-universe}} & \thepage & No U:U; semantic model assumes two external Grothendieck universes Original precedence 'atomic' retained as semantic descriptor; n/a is not an invented rank. \\
\phantomsection\label{entry:symbol-registry-05719}eqToPath & \texttt{\detokenize{\mathsf{eqToPath}}} & Explicit Glue type-path construction & P07 CCHM2016 v1 & n/a & \hyperref[sec:p07-cchm2016-universe]{\ref*{sec:p07-cchm2016-universe}} & \thepage & Typed path does not alone discharge full univalence coherence Original precedence 'prefix' retained as semantic descriptor; n/a is not an invented rank. \\
\phantomsection\label{entry:symbol-registry-05880}\(\kappa\) & \texttt{\detokenize{\kappa}} & terminal CUT marker identity & P7 girard-1987-cut-trips-and-boxes & n/a & \hyperref[def:p07-girard-cut-marker]{\ref*{def:p07-girard-cut-marker}} & \thepage & Port-bearing symbol not an object-language formula; cannot be consumed or counted as an ordinary door. \\
\phantomsection\label{entry:symbol-registry-05881}\(\kappa^-,\kappa^+\) & \texttt{\detokenize{\kappa^-,}\allowbreak\detokenize{\kappa^+}} & two auxiliary CUT gates & P7 girard-1987-cut-trips-and-boxes & n/a & \hyperref[def:p07-girard-cut-marker]{\ref*{def:p07-girard-cut-marker}} & \thepage & Their insertion subdivides one direct cut crossing and does not change trip orbits. \\
\phantomsection\label{entry:symbol-registry-05882}\(p,c,s\) & \texttt{\detokenize{p,}\allowbreak\detokenize{c,}\allowbreak\detokenize{s}} & PN0 ordinary occurrence CUT and combined size counts & P7 girard-1987-cut-trips-and-boxes & n/a & \hyperref[def:p07-girard-cut-marker]{\ref*{def:p07-girard-cut-marker}} & \thepage & Source measure s=p+c; not link count or an arbitrary boxed normalization bound. \\
\phantomsection\label{entry:symbol-registry-05883}\(\rho\) & \texttt{\detokenize{\rho}} & one-cycle permutation of box door indices & P7 girard-1987-cut-trips-and-boxes & n/a & \hyperref[def:p07-girard-box-doors]{\ref*{def:p07-girard-box-doors}} & \thepage & All n-cycles are allowed; arbitrary permutations are not; n=1 uses identity. \\
\phantomsection\label{entry:symbol-registry-05884}\(\operatorname{Conc}(\beta)\) & \texttt{\detokenize{\operatorname{Conc}(\beta)}} & ordinary terminal formula-door list & P7 girard-1987-cut-trips-and-boxes & n/a & \hyperref[def:p07-girard-box-front-aux]{\ref*{def:p07-girard-box-front-aux}} & \thepage & CUT symbols excluded; actual inner-net certificate required by each formation card. \\
\phantomsection\label{entry:symbol-registry-05885}\(\Gamma_?\) & \texttt{\detokenize{\Gamma_?}} & quest-prefixed auxiliary occurrence context & P7 girard-1987-cut-trips-and-boxes & n/a & \hyperref[def:p07-girard-box-front-aux]{\ref*{def:p07-girard-box-front-aux}} & \thepage & Every member quest-prefixed in promotion; not an arbitrary context. \\
\phantomsection\label{entry:symbol-registry-05886}\(P_B,P_C,Q_1,Q_2\) & \texttt{\detokenize{P_B,}\allowbreak\detokenize{P_C,}\allowbreak\detokenize{Q_1,}\allowbreak\detokenize{Q_2}} & four exhaustive principal-cut trip slices & P7 girard-1987-cut-trips-and-boxes & n/a & \hyperref[def:p07-girard-cut-slices]{\ref*{def:p07-girard-cut-slices}} & \thepage & Disjoint outside paths; source principal contraction rejoins their order without deleting a slice. \\
\phantomsection\label{entry:symbol-registry-06164}\(C\) & \texttt{\detokenize{C}} & signed canonical carrier type & P7 2013 circle & n/a & \hyperref[def:p07-circle-integer-normalization]{\ref*{def:p07-circle-integer-normalization}} & \thepage & Type 1+N+N; zero and positive/negative successor tags are disjoint. \\
\phantomsection\label{entry:symbol-registry-06165}\(g\) & \texttt{\detokenize{g}} & pair normalization function & P7 2013 circle & n/a & \hyperref[def:p07-circle-integer-normalization]{\ref*{def:p07-circle-integer-normalization}} & \thepage & Typed N×N\ensuremath{\to}C; cancels common successors. \\
\phantomsection\label{entry:symbol-registry-06166}\(j\) & \texttt{\detokenize{j}} & canonical pair embedding & P7 2013 circle & n/a & \hyperref[def:p07-circle-integer-normalization]{\ref*{def:p07-circle-integer-normalization}} & \thepage & Typed C\ensuremath{\to}N×N; not judgmentally inverse at a free variable. \\
\phantomsection\label{entry:symbol-registry-06167}\(r\) & \texttt{\detokenize{r}} & idempotent pair canonicalizer & P7 2013 circle & n/a & \hyperref[def:p07-circle-integer-normalization]{\ref*{def:p07-circle-integer-normalization}} & \thepage & Typed N×N\ensuremath{\to}N×N and defined as jg. \\
\phantomsection\label{entry:symbol-registry-06168}\(\mathbb Z\) & \texttt{\detokenize{\mathbb Z}} & source canonical fixed-point integer type & P7 2013 circle & n/a & \hyperref[def:p07-circle-integer-normalization]{\ref*{def:p07-circle-integer-normalization}} & \thepage & Sigma x:N×N of r(x)=x; proved equivalent to the source quotient into sets. \\
\phantomsection\label{entry:symbol-registry-06169}\(E\) & \texttt{\detokenize{E}} & signed presentation comparison & P7 2013 circle & n/a & \hyperref[def:p07-circle-integer-normalization]{\ref*{def:p07-circle-integer-normalization}} & \thepage & Typed C\ensuremath{\to}Z; its fixed-point witness is defined by tag induction. \\
\phantomsection\label{entry:symbol-registry-06170}\(G\) & \texttt{\detokenize{G}} & inverse signed presentation comparison & P7 2013 circle & n/a & \hyperref[def:p07-circle-integer-normalization]{\ref*{def:p07-circle-integer-normalization}} & \thepage & Typed Z\ensuremath{\to}C; ignores the proposition-valued fixed-point proof. \\
\phantomsection\label{entry:symbol-registry-06171}\(\nu_z\) & \texttt{\detokenize{\nu}} & canonical presentation homotopy & P7 2013 circle & n/a & \hyperref[thm:p07-circle-integer-presentation]{\ref*{thm:p07-circle-integer-presentation}} & \thepage & Typed E(G(z))=z; arbitrary signed case proofs transport along this path. \\
\phantomsection\label{entry:symbol-registry-06172}\(\mathsf{succ}\) & \texttt{\detokenize{\mathsf{succ}}} & integer successor equivalence forward function & P7 2013 circle & n/a & \hyperref[def:p07-circle-integer-steps]{\ref*{def:p07-circle-integer-steps}} & \thepage & Typed Z\ensuremath{\to}Z; inverse pred; no universe judgmental univalence computation. \\
\phantomsection\label{entry:symbol-registry-06173}\(\mathsf{pred}\) & \texttt{\detokenize{\mathsf{pred}}} & integer predecessor inverse function & P7 2013 circle & n/a & \hyperref[def:p07-circle-integer-steps]{\ref*{def:p07-circle-integer-steps}} & \thepage & Typed Z\ensuremath{\to}Z; includes negative and zero boundary cases. \\
\phantomsection\label{entry:symbol-registry-06174}\(\ell^n\) & \texttt{\detokenize{\ell^n}} & signed loop power & P7 2013 circle & n/a & \hyperref[def:p07-circle-powers]{\ref*{def:p07-circle-powers}} & \thepage & Typed a=a for ell:a=a and n:Z; only zero power is asserted judgmentally reflexive. Original precedence descriptor 'exponent' preserved; n/a is not an invented rank. \\
\phantomsection\label{entry:symbol-registry-06175}\(\mathrm{code}\) & \texttt{\detokenize{\mathrm{code}}} & univalent circle family & P7 2013 circle & n/a & \hyperref[def:p07-circle-code]{\ref*{def:p07-circle-code}} & \thepage & Typed S1\ensuremath{\to}U; base fiber judgmentally Z; loop monodromy propositionally successor. \\
\phantomsection\label{entry:symbol-registry-06176}\(\mathrm{encode}_x\) & \texttt{\detokenize{\mathrm{encode}}} & transport encoder & P7 2013 circle & n/a & \hyperref[def:p07-circle-encode]{\ref*{def:p07-circle-encode}} & \thepage & Typed (base=x)\ensuremath{\to}code(x); sends p to transport(p,0). \\
\phantomsection\label{entry:symbol-registry-06177}\(\mathrm{decode}_x\) & \texttt{\detokenize{\mathrm{decode}}} & circle-inductively constructed decoder & P7 2013 circle & n/a & \hyperref[thm:p07-circle-decoder]{\ref*{thm:p07-circle-decoder}} & \thepage & Typed code(x)\ensuremath{\to}(base=x); supplied transported function lift uses predecessor powers and cancellation. \\
\phantomsection\label{entry:symbol-registry-06178}\(\eta_x\) & \texttt{\detokenize{\eta}} & decode-after-encode homotopy & P7 2013 circle & n/a & \hyperref[thm:p07-circle-decode-encode]{\ref*{thm:p07-circle-decode-encode}} & \thepage & Typed \ensuremath{\Pi}p:base=x decode\_x(encode\_x(p))=p; based J varies the endpoint. \\
\phantomsection\label{entry:symbol-registry-06179}\(\epsilon_x\) & \texttt{\detokenize{\epsilon}} & encode-after-decode homotopy & P7 2013 circle & n/a & \hyperref[thm:p07-circle-encode-decode]{\ref*{thm:p07-circle-encode-decode}} & \thepage & Typed \ensuremath{\Pi}c:code(x) encode\_x(decode\_x(c))=c; circle coherence follows from integer setness and funext. \\
\phantomsection\label{entry:symbol-registry-06308}\((R,r_0)\) & \texttt{\detokenize{(R,}\allowbreak\detokenize{r_0)}} & pointed predicate data & P7 identity-system subsection & n/a & \hyperref[def:p07-is-pointed]{\ref*{def:p07-is-pointed}} & \thepage & R:A\ensuremath{\to}U and r0:R(a0); distinct from the cover HIT R below. \\
\phantomsection\label{entry:symbol-registry-06309}\(\mathsf{ppmap}(R,S)\) & \texttt{\detokenize{\mathsf{ppmap}}} & pointed dependent mapping type & P7 identity-system subsection & n/a & \hyperref[def:p07-is-pointed]{\ref*{def:p07-is-pointed}} & \thepage & Sigma of \ensuremath{\Pi}b R(b)\ensuremath{\to}S(b) and u(a0,r0)=s0; pointing proof is retained. \\
\phantomsection\label{entry:symbol-registry-06310}\(D\) & \texttt{\detokenize{D}} & identity-system dependent motive & P7 identity-system subsection & n/a & \hyperref[def:p07-is-pointed]{\ref*{def:p07-is-pointed}} & \thepage & Typed \ensuremath{\Pi}b:A R(b)\ensuremath{\to}V; not the flattening equivalence family D. \\
\phantomsection\label{entry:symbol-registry-06311}\(e_b\) & \texttt{\detokenize{e_b}} & canonical pointed transport map & P7 identity-system subsection & n/a & \hyperref[thm:p07-is-characterization]{\ref*{thm:p07-is-characterization}} & \thepage & Typed (a0=b)\ensuremath{\to}R(b); p\ensuremath{\mapsto}transport\_R(p,r0). \\
\phantomsection\label{entry:symbol-registry-06312}\(R(a,b)\) & \texttt{\detokenize{R}} & two-sided identity-system family & P7 two-sided subsection & n/a & \hyperref[def:p07-is-two-sided]{\ref*{def:p07-is-two-sided}} & \thepage & Typed A\ensuremath{\to}A\ensuremath{\to}U with r0(a):R(a,a); eliminated r:R(a,b). \\
\phantomsection\label{entry:symbol-registry-06313}\(\mathsf c\) & \texttt{\detokenize{\mathsf c}} & coequalizer point constructor & P7 general flattening subsection & n/a & \hyperref[def:p07-flat-coequalizer]{\ref*{def:p07-flat-coequalizer}} & \thepage & Typed A\ensuremath{\to}W; distinct from cover point constructor Z\ensuremath{\to}R. \\
\phantomsection\label{entry:symbol-registry-06314}\(\mathsf p_b\) & \texttt{\detokenize{\mathsf p}} & coequalizer edge constructor & P7 general flattening subsection & n/a & \hyperref[def:p07-flat-coequalizer]{\ref*{def:p07-flat-coequalizer}} & \thepage & Typed c(fb)=c(gb) in W; not fb=gb in A. \\
\phantomsection\label{entry:symbol-registry-06315}\(C\) & \texttt{\detokenize{C}} & flattening fiber family & P7 general flattening subsection & n/a & \hyperref[def:p07-flat-coequalizer]{\ref*{def:p07-flat-coequalizer}} & \thepage & Typed A\ensuremath{\to}U; distinct from the preceding circle packet's signed carrier C. \\
\phantomsection\label{entry:symbol-registry-06316}\(D_b\) & \texttt{\detokenize{D}} & flattening equivalence family & P7 general flattening subsection & n/a & \hyperref[def:p07-flat-coequalizer]{\ref*{def:p07-flat-coequalizer}} & \thepage & Typed C(fb)\ensuremath{\simeq}C(gb); source-native overlap with the prior motive is explicitly typed. \\
\phantomsection\label{entry:symbol-registry-06317}\(P\) & \texttt{\detokenize{P}} & univalent coequalizer family & P7 general flattening subsection & n/a & \hyperref[def:p07-flat-coequalizer]{\ref*{def:p07-flat-coequalizer}} & \thepage & Typed W\ensuremath{\to}U; point computation judgmental and universe edge computation propositional. \\
\phantomsection\label{entry:symbol-registry-06318}\(\widetilde W\) & \texttt{\detokenize{\widetilde W}} & flattened coequalizer HIT & P7 general flattening subsection & n/a & \hyperref[def:p07-flat-coequalizer]{\ref*{def:p07-flat-coequalizer}} & \thepage & Point index a and fiber value; edge index b and source fiber value. \\
\phantomsection\label{entry:symbol-registry-06319}\([s,r]\) & \texttt{\detokenize{[s,}\allowbreak\detokenize{r]}} & decoded Sigma path & P7 flattening and circle identity systems & n/a & \hyperref[def:p07-flat-pairpaths]{\ref*{def:p07-flat-pairpaths}} & \thepage & s:x=x' and r:transport\_Y(s,u)=u'; neither coordinate is omitted. \\
\phantomsection\label{entry:symbol-registry-06320}\(H(s,u)\) & \texttt{\detokenize{H}} & horizontal Sigma lift & P7 flattening and circle identity systems & n/a & \hyperref[def:p07-flat-pairpaths]{\ref*{def:p07-flat-pairpaths}} & \thepage & Typed (x,u)=(x',transport\_Y(s,u)); second coordinate reflexivity. \\
\phantomsection\label{entry:symbol-registry-06321}\(V_x(r)\) & \texttt{\detokenize{V}} & vertical Sigma path & P7 flattening and circle identity systems & n/a & \hyperref[def:p07-flat-pairpaths]{\ref*{def:p07-flat-pairpaths}} & \thepage & Typed (x,u)=(x,u'); first coordinate reflexivity. \\
\phantomsection\label{entry:symbol-registry-06322}\(\mathcal K_s\) & \texttt{\detokenize{\mathcal K}} & calibrated dependent-function-path equivalence & P7 flattening and circle identity systems & n/a & \hyperref[def:p07-flat-pairpaths]{\ref*{def:p07-flat-pairpaths}} & \thepage & Chosen by J with happly/funext at reflexivity; both inverse comparisons retained. \\
\phantomsection\label{entry:symbol-registry-06323}\(\mathcal K^0_s\) & \texttt{\detokenize{\mathcal K^0}} & constant-codomain calibrated comparison & P7 flattening and circle-cover applications & n/a & \hyperref[def:p07-flat-pairpaths]{\ref*{def:p07-flat-pairpaths}} & \thepage & Components v(u)=v'(transport\_Y(s,u)); actual inverse counit controls recursor beta. \\
\phantomsection\label{entry:symbol-registry-06324}\(\theta_b(y)\) & \texttt{\detokenize{\theta}} & univalent fiber transport comparison & P7 general flattening subsection & n/a & \hyperref[thm:p07-flat-total-constructors]{\ref*{thm:p07-flat-total-constructors}} & \thepage & Typed transport\_P(p\_b,y)=D\_b(y); propositional. \\
\phantomsection\label{entry:symbol-registry-06325}\(\mathsf c'\) & \texttt{\detokenize{\mathsf c'}} & total-space point constructor & P7 general flattening subsection & n/a & \hyperref[thm:p07-flat-total-constructors]{\ref*{thm:p07-flat-total-constructors}} & \thepage & Typed \ensuremath{\Pi}a C(a)\ensuremath{\to}\ensuremath{\Sigma}w P(w); y\ensuremath{\mapsto}(c(a),y). \\
\phantomsection\label{entry:symbol-registry-06326}\(\mathsf p'_{b,y}\) & \texttt{\detokenize{\mathsf p'}} & total-space edge constructor & P7 general flattening subsection & n/a & \hyperref[thm:p07-flat-total-constructors]{\ref*{thm:p07-flat-total-constructors}} & \thepage & Decoded pair (p\_b,theta\_b(y)); monodromy retained. \\
\phantomsection\label{entry:symbol-registry-06327}\(R\) & \texttt{\detokenize{R}} & integer-indexed cover HIT & P7 circle-cover subsection & n/a & \hyperref[def:p07-flat-circle-cover]{\ref*{def:p07-flat-circle-cover}} & \thepage & Type with c:Z\ensuremath{\to}R and p\_n:c(n)=c(succ n); not the earlier predicate R:A\ensuremath{\to}U. \\
\phantomsection\label{entry:symbol-registry-06328}\(\mathsf c\) & \texttt{\detokenize{\mathsf c}} & integer-cover point constructor & P7 circle-cover subsection & n/a & \hyperref[def:p07-flat-circle-cover]{\ref*{def:p07-flat-circle-cover}} & \thepage & Typed Z\ensuremath{\to}R; distinct local role from A\ensuremath{\to}W. \\
\phantomsection\label{entry:symbol-registry-06329}\(\mathsf p_n\) & \texttt{\detokenize{\mathsf p}} & integer-cover edge constructor & P7 circle-cover subsection & n/a & \hyperref[def:p07-flat-circle-cover]{\ref*{def:p07-flat-circle-cover}} & \thepage & Typed c(n)=c(succ n); negative recursion uses the lower integer's edge. \\
\phantomsection\label{entry:symbol-registry-06330}\(q_n\) & \texttt{\detokenize{q}} & integer-cover contraction path & P7 circle-cover subsection & n/a & \hyperref[thm:p07-flat-circle-cover]{\ref*{thm:p07-flat-circle-cover}} & \thepage & Typed c(0)=c(n); signed recurrence and canonical-representative transport proved. \\
\phantomsection\label{entry:symbol-registry-06331}\(K_n\) & \texttt{\detokenize{K}} & signed winding total-space path & P7 native circle identity-system subsection & n/a & \hyperref[thm:p07-is-circle-pair-coherence]{\ref*{thm:p07-is-circle-pair-coherence}} & \thepage & Typed (base,0)=(base,n); pair(loop\^{}n,h\_n). \\
\phantomsection\label{entry:symbol-registry-06332}\(S_n\) & \texttt{\detokenize{S}} & successor total-space edge & P7 native circle identity-system subsection & n/a & \hyperref[thm:p07-is-circle-pair-coherence]{\ref*{thm:p07-is-circle-pair-coherence}} & \thepage & Typed (base,n)=(base,succ n); pair(loop,theta\_n). \\
\phantomsection\label{entry:symbol-registry-06333}\(h_n\) & \texttt{\detokenize{h}} & signed winding witness & P7 native circle identity-system subsection & n/a & \hyperref[thm:p07-is-circle-pair-coherence]{\ref*{thm:p07-is-circle-pair-coherence}} & \thepage & Typed transport\_code(loop\^{}n,0)=n; transported variants compared by integer setness. \\
\phantomsection\label{entry:symbol-registry-06334}\(\widehat D\) & \texttt{\detokenize{\widehat D}} & uncurried code-dependent motive & P7 native circle identity-system subsection & n/a & \hyperref[thm:p07-is-circle-native-induction]{\ref*{thm:p07-is-circle-native-induction}} & \thepage & Typed (\ensuremath{\Sigma}x:S1 code(x))\ensuremath{\to}V; transport along K\_n precedes circle lift construction. \\
\end{longtable}
\endgroup
\Needspace{.4\textheight}
\subsection*{Stable part P8 scopes}
\begingroup\scriptsize\setlength{\tabcolsep}{3pt}
\begin{longtable}{@{}p{0.9in}p{1.3in}p{1.0in}p{0.5in}p{0.3in}p{0.4in}p{0.35in}p{1.65in}@{}}
\toprule
\textbf{Glyph} & \textbf{Macro} & \textbf{Role} & \textbf{Scope} & \textbf{Prec./form} & \textbf{First} & \textbf{Reg. page} & \textbf{Principal uses / notes} \\
\midrule
\endfirsthead
\toprule
\textbf{Glyph} & \textbf{Macro} & \textbf{Role} & \textbf{Scope} & \textbf{Prec./form} & \textbf{First} & \textbf{Reg. page} & \textbf{Principal uses / notes} \\
\midrule
\endhead
\midrule\multicolumn{8}{r}{\emph{continued}}\\\endfoot
\bottomrule\endlastfoot
\phantomsection\label{entry:symbol-registry-00090}\(w\Vdash A\) & \texttt{\detokenize{\forces}} & Kripke forcing & P8 intuitionistic-logic & meta & \hyperref[sec:p08-source-map]{\ref*{sec:p08-source-map}} & \thepage & The negated partner is \(\nVdash\) (macro \texttt{\detokenize{\nforces}}); modal satisfaction continues to use \(\models\). \\
\phantomsection\label{entry:symbol-registry-00092}\(\Box A\) & \texttt{\detokenize{\Box}} & Anderson--Belnap relevance abbreviation & P8 paraconsistent-and-nonmonotonic & 3 & \hyperref[sec:p08-paraconsistent-and-nonmonotonic]{\ref*{sec:p08-paraconsistent-and-nonmonotonic}} & \thepage & Exact eliminable abbreviation \((A\to A)\to A\) on the E and R sheets; it is not a primitive modality. \\
\phantomsection\label{entry:symbol-registry-00094}\(\mathsf X\) & \texttt{\detokenize{\mathsf X}} & temporal next & Parts 8–9 from conditionals-deontic-free & 3 & \hyperref[sec:p08-conditionals-deontic-free]{\ref*{sec:p08-conditionals-deontic-free}} & \thepage & Prefix unary at the negation rank. \\
\phantomsection\label{entry:symbol-registry-00095}\(A\mathbin{\mathsf U}B\) & \texttt{\detokenize{\mathbin{\mathsf U}}} & temporal until & Parts 8–9 from conditionals-deontic-free & 3.5 & \hyperref[sec:p08-conditionals-deontic-free]{\ref*{sec:p08-conditionals-deontic-free}} & \thepage & Binds tighter than every binary connective and looser than the unary rank; licenses the U-unfolding axiom's written form. \\
\phantomsection\label{entry:symbol-registry-00096}\(\mathsf F\) & \texttt{\detokenize{\mathsf F}} & temporal eventually & Parts 8–9 from conditionals-deontic-free & 3 & \hyperref[sec:p08-conditionals-deontic-free]{\ref*{sec:p08-conditionals-deontic-free}} & \thepage & Defined \(\mathsf FA:\equiv\top\mathbin{\mathsf U}A\), prefix unary at the negation rank. \\
\phantomsection\label{entry:symbol-registry-00097}\(\mathsf G\) & \texttt{\detokenize{\mathsf G}} & temporal always & Parts 8–9 from conditionals-deontic-free & 3 & \hyperref[sec:p08-conditionals-deontic-free]{\ref*{sec:p08-conditionals-deontic-free}} & \thepage & Defined \(\mathsf GA:\equiv\neg\mathsf F\neg A\), prefix unary at the negation rank. \\
\phantomsection\label{entry:symbol-registry-00098}\(K_i\) & \texttt{\detokenize{K_i}} & knowledge operator (agent i) & P8 conditionals-deontic-free & 3 & \hyperref[sec:p08-conditionals-deontic-free]{\ref*{sec:p08-conditionals-deontic-free}} & \thepage & Prefix unary at the negation rank. \\
\phantomsection\label{entry:symbol-registry-00099}\(\sim_i\) & \texttt{\detokenize{\sim_i}} & epistemic accessibility & P8 conditionals-deontic-free & n/a & \hyperref[sec:p08-conditionals-deontic-free]{\ref*{sec:p08-conditionals-deontic-free}} & \thepage & Agent-indexed accessibility; an equivalence relation for knowledge readings. \\
\phantomsection\label{entry:symbol-registry-00100}\(E_G\) & \texttt{\detokenize{E_G}} & mutual knowledge & P8 conditionals-deontic-free & 3 & \hyperref[sec:p08-conditionals-deontic-free]{\ref*{sec:p08-conditionals-deontic-free}} & \thepage & \(E_GA:\equiv\bigwedge_{i\in G}K_iA\), prefix unary at the negation rank. \\
\phantomsection\label{entry:symbol-registry-00101}\(C_G\) & \texttt{\detokenize{C_G}} & common knowledge & P8 conditionals-deontic-free & 3 & \hyperref[sec:p08-conditionals-deontic-free]{\ref*{sec:p08-conditionals-deontic-free}} & \thepage & Prefix unary at the negation rank; satisfaction via reflexive-transitive reachability under the union of the \(\sim_i\). \\
\phantomsection\label{entry:symbol-registry-00102}\(\otimes\) & \texttt{\detokenize{\otimes}} & linear tensor and its homomorphic semantic interpretation & P8 paraconsistent-and-nonmonotonic & 4 & \hyperref[sec:p08-paraconsistent-and-nonmonotonic]{\ref*{sec:p08-paraconsistent-and-nonmonotonic}} & \thepage & Formula operands give the exact IMLL or CLL connective; the IMLL algebra and CLL phase clauses interpret that same connective with their source-standard tensor operation. \\
\phantomsection\label{entry:symbol-registry-00103}\(\multimap\) & \texttt{\detokenize{\multimap}} & linear implication and its residual interpretation & P8 paraconsistent-and-nonmonotonic & 6 & \hyperref[sec:p08-paraconsistent-and-nonmonotonic]{\ref*{sec:p08-paraconsistent-and-nonmonotonic}} & \thepage & Formula operands give IMLL or CLL linear implication; the algebraic and phase clauses interpret that same connective by the source-standard residual. \\
\phantomsection\label{entry:symbol-registry-00104}\(\mathbf 1\) & \texttt{\detokenize{\mathbf 1}} & multiplicative unit & P8 paraconsistent-and-nonmonotonic & n/a & \hyperref[sec:p08-paraconsistent-and-nonmonotonic]{\ref*{sec:p08-paraconsistent-and-nonmonotonic}} & \thepage & Syntactic tensor unit in IMLL and CLL; the ambient calculus distinguishes it from a semantic monoid identity. \\
\phantomsection\label{entry:symbol-registry-00105}\(!A\) & \texttt{\detokenize{!}} & bang exponential and its phase interpretation & P8 paraconsistent-and-nonmonotonic & 2.5 & \hyperref[sec:p08-paraconsistent-and-nonmonotonic]{\ref*{sec:p08-paraconsistent-and-nonmonotonic}} & \thepage & Prefix connective on the source-bounded IMLL extension and full CLL; the phase clause for \(!F\) interprets that same connective, and neither use is factorial. \\
\phantomsection\label{entry:symbol-registry-00106}\(W\) & \texttt{\detokenize{W}} & set of worlds & P8 intuitionistic-logic to rationality-2-6 & n/a & \hyperref[sec:p08-source-map]{\ref*{sec:p08-source-map}} & \thepage & Nonempty world carrier for intuitionistic, modal, temporal, epistemic, and broad Routley--Meyer semantics; world letters are declared locally. \\
\phantomsection\label{entry:symbol-registry-00107}\(R\) & \texttt{\detokenize{R}} & binary accessibility relation & P8 modal-temporal-epistemic to conditionals-deontic-free & n/a & \hyperref[sec:modal-temporal-epistemic-01]{\ref*{sec:modal-temporal-epistemic-01}} & \thepage & Frame relation of modal, temporal, and epistemic structures; scope is disjoint from the P8.6 ternary relation. \\
\phantomsection\label{entry:symbol-registry-00108}\(N\) & \texttt{\detokenize{N}} & normal worlds & P8 paraconsistent-and-nonmonotonic & n/a & \hyperref[sec:p08-paraconsistent-and-nonmonotonic]{\ref*{sec:p08-paraconsistent-and-nonmonotonic}} & \thepage & Nonempty distinguished subset of the broad inherited Routley--Meyer structure fixing its source-scoped validity notion. \\
\phantomsection\label{entry:symbol-registry-00109}\(R\) & \texttt{\detokenize{R}} & ternary accessibility relation & P8 paraconsistent-and-nonmonotonic & n/a & \hyperref[sec:p08-paraconsistent-and-nonmonotonic]{\ref*{sec:p08-paraconsistent-and-nonmonotonic}} & \thepage & Ternary relation in broad and matched Routley--Meyer structures; scope is disjoint from binary accessibility. \\
\phantomsection\label{entry:symbol-registry-00110}\(w^*\) & \texttt{\detokenize{^{*}}} & Routley star & P8 paraconsistent-and-nonmonotonic & n/a & \hyperref[sec:p08-paraconsistent-and-nonmonotonic]{\ref*{sec:p08-paraconsistent-and-nonmonotonic}} & \thepage & Postfix unary operation on relevance worlds used by the negation clause; distinct from closure stars and the extension marker in \(\mathsf L^*\). \\
\phantomsection\label{entry:symbol-registry-00111}\(\mathsf V\) & \texttt{\detokenize{\mathsf V}} & matrix value set & P8 relevance-linear & n/a & \hyperref[sec:relevance-linear-01]{\ref*{sec:relevance-linear-01}} & \thepage & Finite nonempty truth-value set for the generic matrix sheet; distinct from world valuation \(V\) and cumulative \(V_\alpha\). \\
\phantomsection\label{entry:symbol-registry-00112}\(D\) & \texttt{\detokenize{D}} & designated values & P8 relevance-linear & n/a & \hyperref[sec:relevance-linear-01]{\ref*{sec:relevance-linear-01}} & \thepage & Subset of the active matrix value set fixing matrix consequence. \\
\phantomsection\label{entry:symbol-registry-00113}\(f_c\) & \texttt{\detokenize{f_c}} & matrix operation for connective c & P8 relevance-linear & n/a & \hyperref[sec:relevance-linear-01]{\ref*{sec:relevance-linear-01}} & \thepage & Total n-ary operation on \(\mathsf V\) interpreting the declared connective \(c\). \\
\phantomsection\label{entry:symbol-registry-00115}\(\models^{\ell}\) & \texttt{\detokenize{\entails^{\ell}}} & local modal consequence & P8 modal-temporal-epistemic to conditionals-deontic-free & meta & \hyperref[sec:modal-temporal-epistemic-01]{\ref*{sec:modal-temporal-epistemic-01}} & \thepage & Truth preservation at one world in the declared modal or epistemic frame class. \\
\phantomsection\label{entry:symbol-registry-00116}\(\models^{g}\) & \texttt{\detokenize{\entails^{g}}} & global modal consequence & P8 modal-temporal-epistemic to conditionals-deontic-free & meta & \hyperref[sec:modal-temporal-epistemic-01]{\ref*{sec:modal-temporal-epistemic-01}} & \thepage & Truth preservation at all worlds in the declared modal or epistemic frame class. \\
\phantomsection\label{entry:symbol-registry-00117}\(\mathcal M,w\models A\) & \texttt{\detokenize{\entails}} & satisfaction at a world & P8 intuitionistic-logic to rationality-2-6 & meta & \hyperref[sec:modal-temporal-epistemic-01]{\ref*{sec:modal-temporal-epistemic-01}} & \thepage & Satisfaction at world \(w\) of the structure \(\mathcal M\). \\
\phantomsection\label{entry:symbol-registry-00118}\(A,B,C\) & \texttt{\detokenize{A,}\allowbreak\detokenize{B,}\allowbreak\detokenize{C}} & default formula metavariables in inherited Part-8 prose & P8 intuitionistic-logic to rationality-2-6 & n/a & \hyperref[sec:p08-source-map]{\ref*{sec:p08-source-map}} & \thepage & Default only where the inherited or cited presentation uses these letters.  A named source presentation's own locally declared syntax overrides this row; this row never licenses renaming source metavariables. \\
\phantomsection\label{entry:symbol-registry-00119}\(\mathcal A=(A,\leq,\otimes,1,\multimap)\) & \texttt{\detokenize{\mathcal A,}\allowbreak\detokenize{A,}\allowbreak\detokenize{\leq,}\allowbreak\detokenize{\otimes,}\allowbreak\detokenize{1,}\allowbreak\detokenize{\multimap}} & residuated pomonoid datum & P8 paraconsistent-and-nonmonotonic & n/a & \hyperref[sec:p08-paraconsistent-and-nonmonotonic]{\ref*{sec:p08-paraconsistent-and-nonmonotonic}} & \thepage & The inherited source tuple.  Within its definition A is locally typed as the algebra carrier, overriding the default formula-metavariable role. \\
\phantomsection\label{entry:symbol-registry-00120}\(V\) & \texttt{\detokenize{V}} & world valuation & P8 intuitionistic-logic to rationality-2-6 & n/a & \hyperref[sec:p08-source-map]{\ref*{sec:p08-source-map}} & \thepage & \(V:\mathsf{Prop}\to\mathcal P(W)\) (temporal structures: \(V:\mathbb N\to\mathcal P(\mathsf{Prop})\)); distinct from matrix \(\mathsf V\) and cumulative \(V_\alpha\). \\
\phantomsection\label{entry:symbol-registry-00124}\(\mathsf{Prop}\) & \texttt{\detokenize{\mathsf{Prop}}} & default propositional-atom set & P8 intuitionistic-logic to rationality-2-6 & n/a & \hyperref[sec:p08-source-map]{\ref*{sec:p08-source-map}} & \thepage & Used where the inherited or cited presentation declares this atom-set glyph.  A named source presentation may instead declare its own atom notation, which governs locally without house renaming. \\
\phantomsection\label{entry:symbol-registry-00490}\(\mathsf{IPL},\mathsf{KC},\mathsf{LC}\) & \texttt{\detokenize{\mathsf{IPL},}\allowbreak\detokenize{\mathsf{KC},}\allowbreak\detokenize{\mathsf{LC}}} & intuitionistic and intermediate logic identifiers & P8 intuitionistic-logic & n/a & \hyperref[sec:p08-source-map]{\ref*{sec:p08-source-map}} & \thepage & Exact IPL card and the named KC and Goedel--Dummett LC extensions; IQC is omitted because no IQC identifier is printed. \\
\phantomsection\label{entry:symbol-registry-00491}\(\Gamma\derives_{\mathsf{IPL}}A\) & \texttt{\detokenize{\derives_{\mathsf{IPL}}}} & finite-support IPL derivability & P8 intuitionistic-logic to modal-temporal-epistemic & meta & \hyperref[sec:p08-source-map]{\ref*{sec:p08-source-map}} & \thepage & A premise-set judgment witnessed by a finite multiset subcontext derivable on sheet:lj; P8.2 reuses it only in the exact IPL-to-S4 embedding theorem. \\
\phantomsection\label{entry:symbol-registry-00492}\(\Gamma\entails_{\mathsf{IPL}}A\) & \texttt{\detokenize{\entails_{\mathsf{IPL}}}} & IPL Kripke consequence & P8 intuitionistic-logic & meta & \hyperref[sec:p08-source-map]{\ref*{sec:p08-source-map}} & \thepage & Truth preservation at one world in every intuitionistic Kripke structure of the declared model class under all premises in Gamma. \\
\phantomsection\label{entry:symbol-registry-00493}\(\Gamma_0\) & \texttt{\detokenize{\Gamma_0}} & finite-support witness context & P8 intuitionistic-logic & n/a & \hyperref[sec:p08-source-map]{\ref*{sec:p08-source-map}} & \thepage & Finite multiset whose members lie in the ambient IPL premise set; disjoint from the P7.5 Feferman--Schuette ordinal. \\
\phantomsection\label{entry:symbol-registry-00494}\(G(A),A^{\mathrm{Ko}},A^{\mathrm N},A_{\mathrm{Ku}},A^{\mathrm{Ku}}\) & \texttt{\detokenize{G,}\allowbreak\detokenize{A^{\mathrm{Ko}},}\allowbreak\detokenize{A^{\mathrm N},}\allowbreak\detokenize{A_{\mathrm{Ku}},}\allowbreak\detokenize{A^{\mathrm{Ku}}}} & negative-translation maps & P8 intuitionistic-logic & meta & \hyperref[sec:p08-source-map]{\ref*{sec:p08-source-map}} & \thepage & Exact Glivenko Kolmogorov Goedel--Gentzen and inner/full Kuroda maps; superscripts act memberwise on premise contexts when printed there. \\
\phantomsection\label{entry:symbol-registry-00495}\(\mathcal H=(H,\wedge,\vee,0,1,\Rightarrow)\) & \texttt{\detokenize{\mathcal H,}\allowbreak\detokenize{H,}\allowbreak\detokenize{\wedge,}\allowbreak\detokenize{\vee,}\allowbreak\detokenize{0,}\allowbreak\detokenize{1,}\allowbreak\detokenize{\Rightarrow}} & Heyting algebra datum & P8 intuitionistic-logic & n/a & \hyperref[sec:p08-source-map]{\ref*{sec:p08-source-map}} & \thepage & The cited presentation's bounded distributive lattice with residual.  Local element operands type these glyphs as algebraic operations rather than object-language connectives or a sequent arrow. \\
\phantomsection\label{entry:symbol-registry-00496}\(a\wedge b,a\vee b,0,1,a\Rightarrow b,\neg a\) & \texttt{\detokenize{\wedge,}\allowbreak\detokenize{\vee,}\allowbreak\detokenize{0,}\allowbreak\detokenize{1,}\allowbreak\detokenize{\Rightarrow,}\allowbreak\detokenize{\neg}} & Heyting operations and bounds & P8 intuitionistic-logic & n/a & \hyperref[sec:p08-source-map]{\ref*{sec:p08-source-map}} & \thepage & Source operations on elements of the locally declared Heyting algebra; the operand type, not a house subscript, distinguishes them from formula connectives. \\
\phantomsection\label{entry:symbol-registry-00497}\(\nu\) & \texttt{\detokenize{\nu}} & Heyting or open-set formula valuation & P8 intuitionistic-logic & n/a & \hyperref[sec:p08-source-map]{\ref*{sec:p08-source-map}} & \thepage & Local semantic valuation extended homomorphically from propositional atoms; disjoint from the P6.3 forcing-name metavariable. \\
\phantomsection\label{entry:symbol-registry-00498}\(\operatorname{int}_{X}\) & \texttt{\detokenize{\operatorname{int}_{X}}} & topological interior operator & P8 intuitionistic-logic & meta & \hyperref[sec:p08-source-map]{\ref*{sec:p08-source-map}} & \thepage & Interior in the declared topological space X used by the open-set implication clause. \\
\phantomsection\label{entry:symbol-registry-00499}\(D_w\) & \texttt{\detokenize{D_w}} & world-indexed intuitionistic domain & P8 intuitionistic-logic & n/a & \hyperref[sec:p08-source-map]{\ref*{sec:p08-source-map}} & \thepage & Nonempty domain at world w in expanding-domain IQC semantics; monotone along the intuitionistic preorder. \\
\phantomsection\label{entry:symbol-registry-00500}\(K,KD,KT,KTB,KB\) & \texttt{\detokenize{K,}\allowbreak\detokenize{KD,}\allowbreak\detokenize{KT,}\allowbreak\detokenize{KTB,}\allowbreak\detokenize{KB}} & Chellas normal modal logic identifiers & P8 modal-temporal-epistemic & n/a & \hyperref[sec:modal-temporal-epistemic-01]{\ref*{sec:modal-temporal-epistemic-01}} & \thepage & Chellas writes KD KT and KTB for the serial reflexive and reflexive-symmetric extensions and writes KB for the symmetry-only extension.  The familiar D T and B system names occur only as Hughes--Cresswell aliases on the corresponding cards. \\
\phantomsection\label{entry:symbol-registry-00501}\(K4,KT4,S4,KD45\) & \texttt{\detokenize{K4,}\allowbreak\detokenize{KT4,}\allowbreak\detokenize{S4,}\allowbreak\detokenize{KD45}} & Chellas normal modal logic identifiers & P8 modal-temporal-epistemic & n/a & \hyperref[sec:modal-temporal-epistemic-01]{\ref*{sec:modal-temporal-epistemic-01}} & \thepage & Chellas's exact plain math-italic names; S4 is the source-named system KT4. \\
\phantomsection\label{entry:symbol-registry-00502}\(S4.2,S4.3,S5\) & \texttt{\detokenize{S4.2,}\allowbreak\detokenize{S4.3,}\allowbreak\detokenize{S5}} & Chellas normal modal logic identifiers & P8 modal-temporal-epistemic & n/a & \hyperref[sec:modal-temporal-epistemic-01]{\ref*{sec:modal-temporal-epistemic-01}} & \thepage & Chellas's exact plain math-italic names for the displayed extensions. \\
\phantomsection\label{entry:symbol-registry-00503}\(\mathsf{GL},\mathsf{Grz}\) & \texttt{\detokenize{\mathsf{GL},}\allowbreak\detokenize{\mathsf{Grz}}} & normal modal logic identifiers & P8 modal-temporal-epistemic & n/a & \hyperref[sec:modal-temporal-epistemic-01]{\ref*{sec:modal-temporal-epistemic-01}} & \thepage & Exact named provability and Grzegorczyk systems outside the ordinary modal-cube route. \\
\phantomsection\label{entry:symbol-registry-00504}\(E,EM,EC,EN,EMC,EMCN\) & \texttt{\detokenize{E,}\allowbreak\detokenize{EM,}\allowbreak\detokenize{EC,}\allowbreak\detokenize{EN,}\allowbreak\detokenize{EMC,}\allowbreak\detokenize{EMCN}} & Chellas compositional modal logic identifiers & P8 modal-temporal-epistemic & n/a & \hyperref[sec:modal-temporal-epistemic-01]{\ref*{sec:modal-temporal-epistemic-01}} & \thepage & Chellas's exact plain math-italic compositional system names.  M C and N remain schema labels rather than substitute system identifiers; normal K is EMCN. \\
\phantomsection\label{entry:symbol-registry-00505}\(\mathfrak L\) & \texttt{\detokenize{\mathfrak L}} & normal modal logic metavariable & P8 modal-temporal-epistemic & n/a & \hyperref[sec:modal-temporal-epistemic-01]{\ref*{sec:modal-temporal-epistemic-01}} & \thepage & Locally fixed normal modal logic in derivability canonical-model and frame-class constructions; disjoint from the P4.6 capture-logic use. \\
\phantomsection\label{entry:symbol-registry-00506}\(\mathcal C\) & \texttt{\detokenize{\mathcal C}} & modal frame-class metavariable & P8 modal-temporal-epistemic & n/a & \hyperref[sec:modal-temporal-epistemic-01]{\ref*{sec:modal-temporal-epistemic-01}} & \thepage & Locally declared class of modal frames; disjoint from the P4 set or language class and the P7 formula-class use. \\
\phantomsection\label{entry:symbol-registry-00507}\(\operatorname{ST}_x,P_p\) & \texttt{\detokenize{\operatorname{ST}_x,}\allowbreak\detokenize{P_p}} & modal standard translation and atom predicate & P8 modal-temporal-epistemic & meta & \hyperref[sec:modal-temporal-epistemic-01]{\ref*{sec:modal-temporal-epistemic-01}} & \thepage & ST at free world variable x maps each modal atom p to the corresponding unary first-order predicate \(P_p\). \\
\phantomsection\label{entry:symbol-registry-00508}\(A^{\Box}\) & \texttt{\detokenize{{}^{\Box}}} & Goedel--McKinsey--Tarski modal translation & P8 modal-temporal-epistemic & meta & \hyperref[sec:modal-temporal-epistemic-01]{\ref*{sec:modal-temporal-epistemic-01}} & \thepage & Recursive IPL-to-S4 formula translation; distinct from applying the object-language box to an already translated formula. \\
\phantomsection\label{entry:symbol-registry-00509}\(Z\) & \texttt{\detokenize{Z}} & modal bisimulation relation & P8 modal-temporal-epistemic & meta & \hyperref[sec:modal-temporal-epistemic-01]{\ref*{sec:modal-temporal-epistemic-01}} & \thepage & Relation between worlds of two pointed Kripke structures satisfying atom zig and zag clauses; disjoint from P8.3 fixpoint variables. \\
\phantomsection\label{entry:symbol-registry-00510}\(\Phi,\sim_{\Phi},[w]_{\Phi}\) & \texttt{\detokenize{\Phi,}\allowbreak\detokenize{\sim_{\Phi},}\allowbreak\detokenize{[\cdot]_{\Phi}}} & filtration set equivalence and quotient world & P8 modal-temporal-epistemic & meta & \hyperref[sec:modal-temporal-epistemic-01]{\ref*{sec:modal-temporal-epistemic-01}} & \thepage & Finite subformula-closed set Phi induces modal equivalence and its world classes in the declared filtration. \\
\phantomsection\label{entry:symbol-registry-00511}\(N:W\to\mathcal P(\mathcal P(W))\) & \texttt{\detokenize{N}} & neighborhood function & P8 modal-temporal-epistemic & n/a & \hyperref[sec:modal-temporal-epistemic-01]{\ref*{sec:modal-temporal-epistemic-01}} & \thepage & The source's world-indexed neighborhood assignment.  Its function type distinguishes it from system EN and the later Routley--Meyer set of normal worlds. \\
\phantomsection\label{entry:symbol-registry-00512}\(\lbrack\!\lbrack A\rbrack\!\rbrack_{\mathcal M}\) & \texttt{\detokenize{\lbrack\!\lbrack\cdot\rbrack\!\rbrack_{\mathcal M}}} & neighborhood truth-set abbreviation & P8 modal-temporal-epistemic & meta & \hyperref[sec:modal-temporal-epistemic-01]{\ref*{sec:modal-temporal-epistemic-01}} & \thepage & Set of worlds satisfying A in the declared neighborhood model; distinct from typed-term and P8.6 semantic brackets. \\
\phantomsection\label{entry:symbol-registry-00513}\(\mathcal A\) & \texttt{\detokenize{\mathcal A}} & admissible proposition family & P8 modal-temporal-epistemic & n/a & \hyperref[sec:modal-temporal-epistemic-01]{\ref*{sec:modal-temporal-epistemic-01}} & \thepage & Boolean-and-modal-operation closed proposition family in a general frame; disjoint from the P8.3 automaton and P8.6 pomonoid. \\
\phantomsection\label{entry:symbol-registry-00514}\(A\alpha,\mathcal F\alpha\) & \texttt{\detokenize{A,}\allowbreak\detokenize{\mathcal F}} & Davies--Humberstone actuality and fixedly operators & P8 modal-temporal-epistemic & 3 & \hyperref[sec:modal-temporal-epistemic-01]{\ref*{sec:modal-temporal-epistemic-01}} & \thepage & Exact source unary operators in the two-index S5 presentation: A resets the lower evaluation index to the upper actual-world index and script F universally varies the upper index. \\
\phantomsection\label{entry:symbol-registry-00515}\(\langle W,V\rangle\models_x^y\alpha\) & \texttt{\detokenize{\models_x^y}} & Davies--Humberstone two-index truth relation & P8 modal-temporal-epistemic & meta & \hyperref[sec:modal-temporal-epistemic-01]{\ref*{sec:modal-temporal-epistemic-01}} & \thepage & Exact source judgment: x is the lower evaluation-world index and y is the upper actual-world perspective. \\
\phantomsection\label{entry:symbol-registry-00516}\(\alpha\) & \texttt{\detokenize{\alpha}} & Davies--Humberstone formula metavariable & P8 modal-temporal-epistemic & n/a & \hyperref[sec:modal-temporal-epistemic-01]{\ref*{sec:modal-temporal-epistemic-01}} & \thepage & Ranges over formulas of the source two-index S5 language; disjoint from the CTL-star path-formula role in P8.3. \\
\phantomsection\label{entry:symbol-registry-00517}\(V(p,x)\) & \texttt{\detokenize{V}} & Davies--Humberstone atomic valuation & P8 modal-temporal-epistemic & n/a & \hyperref[sec:modal-temporal-epistemic-01]{\ref*{sec:modal-temporal-epistemic-01}} & \thepage & Exact source two-place truth-value notation for atom p at world x inside the model pair; locally typed apart from set-valued world valuations. \\
\phantomsection\label{entry:symbol-registry-00518}\(\mathcal M=\langle\mathcal G,\mathcal R,\mathcal D,\mathcal I\rangle\) & \texttt{\detokenize{\mathcal M,}\allowbreak\detokenize{\mathcal G,}\allowbreak\detokenize{\mathcal R,}\allowbreak\detokenize{\mathcal D,}\allowbreak\detokenize{\mathcal I}} & Fitting quantified-modal model datum & P8 modal-temporal-epistemic & n/a & \hyperref[sec:modal-temporal-epistemic-01]{\ref*{sec:modal-temporal-epistemic-01}} & \thepage & Exact source tuple of worlds accessibility world-domain assignment and world-relative interpretation; the local types govern overlapping glyphs. \\
\phantomsection\label{entry:symbol-registry-00519}\(\Gamma,\Delta\) & \texttt{\detokenize{\Gamma,}\allowbreak\detokenize{\Delta}} & Fitting quantified-modal worlds & P8 modal-temporal-epistemic & n/a & \hyperref[sec:modal-temporal-epistemic-01]{\ref*{sec:modal-temporal-epistemic-01}} & \thepage & Source world metavariables in the displayed Fitting model; their local element-of-G type overrides the default context role. \\
\phantomsection\label{entry:symbol-registry-00520}\(\Phi\) & \texttt{\detokenize{\Phi}} & Fitting quantified-modal formula metavariable & P8 modal-temporal-epistemic & n/a & \hyperref[sec:modal-temporal-epistemic-01]{\ref*{sec:modal-temporal-epistemic-01}} & \thepage & Ranges over formulas in the source presentation; locally typed apart from the filtration-set use of Phi. \\
\phantomsection\label{entry:symbol-registry-00521}\(\mathcal D(\Gamma),\mathcal I(P)(\Gamma),\mathcal I(E)(\Gamma)\) & \texttt{\detokenize{\mathcal D,}\allowbreak\detokenize{\mathcal I}} & Fitting local domain and predicate extensions & P8 modal-temporal-epistemic & n/a & \hyperref[sec:modal-temporal-epistemic-01]{\ref*{sec:modal-temporal-epistemic-01}} & \thepage & The source world-domain and world-relative interpretation notation; the special existence predicate satisfies I(E)(Gamma)=D(Gamma). \\
\phantomsection\label{entry:symbol-registry-00522}\((\forall^E x)\Phi,(\exists^E x)\Phi\) & \texttt{\detokenize{\forall^E,}\allowbreak\detokenize{\exists^E}} & Fitting actualist quantifier abbreviations & P8 modal-temporal-epistemic & 3 & \hyperref[sec:modal-temporal-epistemic-01]{\ref*{sec:modal-temporal-epistemic-01}} & \thepage & Exact source relativizations of the primitive possibilist quantifiers by the existence predicate E. \\
\phantomsection\label{entry:symbol-registry-00523}\(E(x)\supset\Phi\) & \texttt{\detokenize{\supset}} & Fitting object-language conditional & P8 modal-temporal-epistemic & 6 & \hyperref[sec:modal-temporal-epistemic-01]{\ref*{sec:modal-temporal-epistemic-01}} & \thepage & Exact source conditional inside the actualist universal-quantifier abbreviation; locally typed apart from the handbook's default implication glyph. \\
\phantomsection\label{entry:symbol-registry-00524}\(\mathcal M,\Gamma\Vdash_v\Phi\) & \texttt{\detokenize{\forces_v}} & Fitting quantified-modal truth judgment & P8 modal-temporal-epistemic & meta & \hyperref[sec:modal-temporal-epistemic-01]{\ref*{sec:modal-temporal-epistemic-01}} & \thepage & Exact source-shaped truth judgment at world Gamma under a world-independent valuation v. \\
\phantomsection\label{entry:symbol-registry-00525}\(\mathsf{LTL},\mathsf{CTL},\mathsf{CTL}^{*},\mathsf{PAL}\) & \texttt{\detokenize{\mathsf{LTL},}\allowbreak\detokenize{\mathsf{CTL},}\allowbreak\detokenize{\mathsf{CTL}^{*},}\allowbreak\detokenize{\mathsf{PAL}}} & temporal and dynamic calculus identifiers & P8 conditionals-deontic-free & n/a & \hyperref[sec:p08-conditionals-deontic-free]{\ref*{sec:p08-conditionals-deontic-free}} & \thepage & Exact registered LTL CTL full CTL-star and public-announcement presentations. \\
\phantomsection\label{entry:symbol-registry-00526}\(L_{\mu}\) & \texttt{\detokenize{L_{\mu}}} & Kozen modal fixed-point language and calculus identifier & P8 conditionals-deontic-free & n/a & \hyperref[sec:p08-conditionals-deontic-free]{\ref*{sec:p08-conditionals-deontic-free}} & \thepage & The exact italic identifier used by Kozen for the finitary propositional modal mu-calculus; the same identifier subscripts its equational derivability relation. \\
\phantomsection\label{entry:symbol-registry-00527}\(\mathcal K=(S,R,V),\pi\) & \texttt{\detokenize{\mathcal K,}\allowbreak\detokenize{S,}\allowbreak\detokenize{R,}\allowbreak\detokenize{V,}\allowbreak\detokenize{\pi}} & total transition structure and path & Parts 8–9 from conditionals-deontic-free & n/a & \hyperref[sec:p08-conditionals-deontic-free]{\ref*{sec:p08-conditionals-deontic-free}} & \thepage & K is a total transition structure and pi is an infinite path from the declared state; this path role is disjoint from the P7.8 proof-code metavariable. \\
\phantomsection\label{entry:symbol-registry-00528}\(\mathsf{S5}_{I},\mathsf{S5}_n\) & \texttt{\detokenize{\mathsf{S5}_{I},}\allowbreak\detokenize{\mathsf{S5}_n}} & multi-agent epistemic calculus identifiers & P8 conditionals-deontic-free & n/a & \hyperref[sec:p08-conditionals-deontic-free]{\ref*{sec:p08-conditionals-deontic-free}} & \thepage & Exact agent-indexed S5 system and its finite-cardinality alias. \\
\phantomsection\label{entry:symbol-registry-00529}\(\mathsf{KD45}_{I},\mathsf{KD45}_n\) & \texttt{\detokenize{\mathsf{KD45}_{I},}\allowbreak\detokenize{\mathsf{KD45}_n}} & multi-agent doxastic calculus identifiers & P8 conditionals-deontic-free & n/a & \hyperref[sec:p08-conditionals-deontic-free]{\ref*{sec:p08-conditionals-deontic-free}} & \thepage & Exact agent-indexed KD45 system and its finite-cardinality alias. \\
\phantomsection\label{entry:symbol-registry-00530}\(\mathsf{S5}_{I}^{C},\mathsf{S5}_{n}^{C}\) & \texttt{\detokenize{\mathsf{S5}_{I}^{C},}\allowbreak\detokenize{\mathsf{S5}_{n}^{C}}} & common-knowledge calculus identifiers & P8 conditionals-deontic-free & n/a & \hyperref[sec:p08-conditionals-deontic-free]{\ref*{sec:p08-conditionals-deontic-free}} & \thepage & Exact common-knowledge extension and its finite-cardinality alias. \\
\phantomsection\label{entry:symbol-registry-00531}\(I\) & \texttt{\detokenize{I}} & nonempty agent-index set & P8 conditionals-deontic-free & n/a & \hyperref[sec:p08-conditionals-deontic-free]{\ref*{sec:p08-conditionals-deontic-free}} & \thepage & Fixed set indexing the epistemic or doxastic modalities and relations; disjoint from the P8.6 idempotent exponential phase set. \\
\phantomsection\label{entry:symbol-registry-00532}\(\mathsf E,\mathsf A\) & \texttt{\detokenize{\mathsf E,}\allowbreak\detokenize{\mathsf A}} & CTL and CTL-star path quantifiers & Parts 8–9 from conditionals-deontic-free & 3 & \hyperref[sec:p08-conditionals-deontic-free]{\ref*{sec:p08-conditionals-deontic-free}} & \thepage & Source-defined existential and universal path quantifiers; scoped away from P8.2 non-normal E. \\
\phantomsection\label{entry:symbol-registry-00533}\(\mathsf{EX},\mathsf{AX},\mathsf{EF},\mathsf{AF}\) & \texttt{\detokenize{\mathsf{EX},}\allowbreak\detokenize{\mathsf{AX},}\allowbreak\detokenize{\mathsf{EF},}\allowbreak\detokenize{\mathsf{AF}}} & CTL next and eventuality operators & Parts 8–9 from conditionals-deontic-free & 3 & \hyperref[sec:p08-conditionals-deontic-free]{\ref*{sec:p08-conditionals-deontic-free}} & \thepage & Source-defined paired existential and universal next and eventuality operators. \\
\phantomsection\label{entry:symbol-registry-00534}\(Z,Y,\mu Z.A,\nu Z.A\) & \texttt{\detokenize{Z,}\allowbreak\detokenize{Y,}\allowbreak\detokenize{\mu,}\allowbreak\detokenize{\nu}} & modal fixed-point variables and binders & Parts 8–9 from conditionals-deontic-free & 3 & \hyperref[sec:p08-conditionals-deontic-free]{\ref*{sec:p08-conditionals-deontic-free}} & \thepage & Z and Y are Kozen's fixed-point variables; mu is the primitive least-fixed-point binder and nu is its capture-avoiding defined dual in the exact L-mu grammar. \\
\phantomsection\label{entry:symbol-registry-00535}\(\zeta,\operatorname{Ext}_{\zeta}(A)\) & \texttt{\detokenize{\zeta,}\allowbreak\detokenize{\operatorname{Ext}_{\zeta}}} & modal fixed-point environment and extension map & Parts 8–9 from conditionals-deontic-free & meta & \hyperref[sec:p08-conditionals-deontic-free]{\ref*{sec:p08-conditionals-deontic-free}} & \thepage & Zeta maps fixed-point variables to state sets and Ext gives the set denoted by A; the P8.3 role is scoped away from the P6.3 forcing-name use. \\
\phantomsection\label{entry:symbol-registry-00536}\(a,b,\langle a\rangle A,[a]A\) & \texttt{\detokenize{a,}\allowbreak\detokenize{b,}\allowbreak\detokenize{\langle a\rangle,}\allowbreak\detokenize{[a]}} & Kozen program constants and indexed modalities & P8 conditionals-deontic-free & 3 & \hyperref[sec:p08-conditionals-deontic-free]{\ref*{sec:p08-conditionals-deontic-free}} & \thepage & Program constants index Kozen's primitive diamond; the box is its exact defined dual.  Automaton input letters use sigma to keep this role distinct. \\
\phantomsection\label{entry:symbol-registry-00537}\(\mathcal I\) & \texttt{\detokenize{\mathcal I}} & Kozen-model interpretation & P8 conditionals-deontic-free & n/a & \hyperref[sec:p08-conditionals-deontic-free]{\ref*{sec:p08-conditionals-deontic-free}} & \thepage & Maps propositional constants to state sets and program constants to binary relations in the exact L-mu semantics; disjoint from the P5.4 relative-interpretation datum. \\
\phantomsection\label{entry:symbol-registry-00538}\(\alpha,\beta\) & \texttt{\detokenize{\alpha,}\allowbreak\detokenize{\beta}} & CTL-star path-formula metavariables & P8 conditionals-deontic-free & n/a & \hyperref[sec:p08-conditionals-deontic-free]{\ref*{sec:p08-conditionals-deontic-free}} & \thepage & Range over Reynolds's path-formula category after the mutually generated CTL-star state/path syntax is declared; disjoint from ordinal and type roles in earlier parts. \\
\phantomsection\label{entry:symbol-registry-00539}\(L,Q;b,b_0\) & \texttt{\detokenize{L,}\allowbreak\detokenize{Q;b,}\allowbreak\detokenize{b_0}} & Reynolds auxiliary atom data & P8 conditionals-deontic-free & n/a & \hyperref[sec:p08-conditionals-deontic-free]{\ref*{sec:p08-conditionals-deontic-free}} & \thepage & In the exact CTL-star AA side condition \(L\) and \(Q\) are disjoint atom sets and \(b\) and \(b_0\) are auxiliary atoms; \(Q\) is locally typed as an atom set, unlike the Buchi state-set role. \\
\phantomsection\label{entry:symbol-registry-00540}\(B_iA\) & \texttt{\detokenize{B_i}} & agent-indexed belief operator & P8 conditionals-deontic-free & 3 & \hyperref[sec:p08-conditionals-deontic-free]{\ref*{sec:p08-conditionals-deontic-free}} & \thepage & Doxastic modality for agent i on the registered \(\mathsf{KD45}_{I}\) sheet; distinct from factive \(K_i\) knowledge. \\
\phantomsection\label{entry:symbol-registry-00541}\(D_GA\) & \texttt{\detokenize{D_G}} & distributed-knowledge operator & P8 conditionals-deontic-free & 3 & \hyperref[sec:p08-conditionals-deontic-free]{\ref*{sec:p08-conditionals-deontic-free}} & \thepage & Truth at every successor of the intersection of the group relations; distinct from everyone-knows and common knowledge. \\
\phantomsection\label{entry:symbol-registry-00542}\(R_G^{E},R_G^{C},R_G^{D}\) & \texttt{\detokenize{R_G^{E},}\allowbreak\detokenize{R_G^{C},}\allowbreak\detokenize{R_G^{D}}} & group epistemic relations & P8 conditionals-deontic-free & n/a & \hyperref[sec:p08-conditionals-deontic-free]{\ref*{sec:p08-conditionals-deontic-free}} & \thepage & Union relation its reflexive-transitive closure and relation intersection for everyone common and distributed knowledge. \\
\phantomsection\label{entry:symbol-registry-00543}\(U\,A\,B,S\,A\,B\) & \texttt{\detokenize{U,}\allowbreak\detokenize{S}} & Kamp strict tense connectives & P8 conditionals-deontic-free & 3 & \hyperref[sec:p08-conditionals-deontic-free]{\ref*{sec:p08-conditionals-deontic-free}} & \thepage & Kamp's prefix source syntax: the first operand holds at the future or past witness and the second at every strictly intervening point. \\
\phantomsection\label{entry:symbol-registry-00544}\(\mathsf P,\mathsf H\) & \texttt{\detokenize{\mathsf P,}\allowbreak\detokenize{\mathsf H}} & unary past tense operators & P8 conditionals-deontic-free & 3 & \hyperref[sec:p08-conditionals-deontic-free]{\ref*{sec:p08-conditionals-deontic-free}} & \thepage & Prior-style sometime-past and always-past operators; their temporal standing is fixed locally. \\
\phantomsection\label{entry:symbol-registry-00545}\([!\chi]A,\langle!\chi\rangle A\) & \texttt{\detokenize{[!\chi],}\allowbreak\detokenize{\langle!\chi\rangle}} & public-announcement modalities & P8 conditionals-deontic-free & 3 & \hyperref[sec:p08-conditionals-deontic-free]{\ref*{sec:p08-conditionals-deontic-free}} & \thepage & Truthful-update box and its executable dual; a false box precondition is vacuously true. \\
\phantomsection\label{entry:symbol-registry-00546}\(\mathcal A=(Q,\Sigma,\delta,q_0,F)\) & \texttt{\detokenize{\mathcal A,}\allowbreak\detokenize{Q,}\allowbreak\detokenize{\Sigma,}\allowbreak\detokenize{\delta,}\allowbreak\detokenize{q_0,}\allowbreak\detokenize{F}} & Buchi automaton datum & Parts 8–9 from conditionals-deontic-free & n/a & \hyperref[sec:p08-conditionals-deontic-free]{\ref*{sec:p08-conditionals-deontic-free}} & \thepage & Finite state set alphabet set-valued transition function initial state and accepting set; a deterministic specialization has singleton transition sets. \\
\phantomsection\label{entry:symbol-registry-00547}\(\operatorname{Inf}(r)\) & \texttt{\detokenize{\operatorname{Inf}}} & infinitely recurring run states & Parts 8–9 from conditionals-deontic-free & meta & \hyperref[sec:p08-conditionals-deontic-free]{\ref*{sec:p08-conditionals-deontic-free}} & \thepage & States occurring infinitely often in run r; Buchi acceptance intersects F while a Rabin pair separately tests its E and F components. \\
\phantomsection\label{entry:symbol-registry-00548}\(\mathbf K_3,\mathsf{LP},\mathbf{FDE}\) & \texttt{\detokenize{\mathbf K_3,}\allowbreak\detokenize{\mathsf{LP},}\allowbreak\detokenize{\mathbf{FDE}}} & three-valued and four-valued logic identifiers & P8 relevance-linear & n/a & \hyperref[sec:relevance-linear-01]{\ref*{sec:relevance-linear-01}} & \thepage & The inherited presentation controls bold K3 and sans-serif LP; Omori--Wansing controls bold upright FDE.  The weak-Kleene matrix remains unnamed because Kleene's presentation supplies no symbolic system identifier. \\
\phantomsection\label{entry:symbol-registry-00549}\(\text{\L}_n\) & \texttt{\detokenize{\text{\L}_n}} & finite Lukasiewicz chain & P8 relevance-linear & n/a & \hyperref[sec:relevance-linear-01]{\ref*{sec:relevance-linear-01}} & \thepage & Cignoli--D'Ottaviano--Mundici's stroked-L name for the n-element chain itself; the corresponding n-valued system is described in source-bound prose. \\
\phantomsection\label{entry:symbol-registry-00550}\(\text{\L}\) & \texttt{\detokenize{\text{\L}}} & standard infinite-valued Lukasiewicz logic identifier & P8 relevance-linear & n/a & \hyperref[sec:relevance-linear-01]{\ref*{sec:relevance-linear-01}} & \thepage & Hajek's exact stroked-L identifier for the infinite-valued logic; its standard value algebra is the unit interval. \\
\phantomsection\label{entry:symbol-registry-00551}\(\mathrm{MTL}\) & \texttt{\detokenize{\mathrm{MTL}}} & monoidal t-norm logic identifier & P8 relevance-linear & n/a & \hyperref[sec:relevance-linear-01]{\ref*{sec:relevance-linear-01}} & \thepage & Esteva--Godo's upright Roman system identifier; the source controls both its letters and typography. \\
\phantomsection\label{entry:symbol-registry-00552}\(\mathrm{BL},G,\Pi\) & \texttt{\detokenize{\mathrm{BL},}\allowbreak\detokenize{G,}\allowbreak\detokenize{\Pi}} & basic Goedel and Product fuzzy-logic identifiers & P8 relevance-linear & n/a & \hyperref[sec:relevance-linear-01]{\ref*{sec:relevance-linear-01}} & \thepage & Hajek's source identifiers: upright Roman BL, plain italic G, and plain Pi; local types distinguish G and Pi from metavariables. \\
\phantomsection\label{entry:symbol-registry-00553}\(\mathrm{BD}\) & \texttt{\detokenize{\mathrm{BD}}} & Belnap--Dunn four-valued algebra & P8 relevance-linear & n/a & \hyperref[sec:relevance-linear-01]{\ref*{sec:relevance-linear-01}} & \thepage & Omori--Wansing's upright source identifier for the four-valued algebra, represented here by the exact positive/negative support pairs. \\
\phantomsection\label{entry:symbol-registry-00554}\(x^+,x^-\) & \texttt{\detokenize{{}^+,}\allowbreak\detokenize{{}^-}} & positive and negative support coordinates & P8 relevance-linear & n/a & \hyperref[sec:relevance-linear-01]{\ref*{sec:relevance-linear-01}} & \thepage & Boolean support indicators of an FDE value; the two coordinates are required to distinguish both from neither. \\
\phantomsection\label{entry:symbol-registry-00555}\(A{:}a\) & \texttt{\detokenize{A{:}a}} & matrix value-labelled formula & P8 relevance-linear & meta & \hyperref[sec:relevance-linear-01]{\ref*{sec:relevance-linear-01}} & \thepage & Tableau occurrence asserting that formula A has exactly matrix value a; distinct from term typing and world labels. \\
\phantomsection\label{entry:symbol-registry-00556}\(f,n,b,t\) & \texttt{\detokenize{f,}\allowbreak\detokenize{n,}\allowbreak\detokenize{b,}\allowbreak\detokenize{t}} & Belnap--Dunn value names & P8 relevance-linear & n/a & \hyperref[sec:relevance-linear-01]{\ref*{sec:relevance-linear-01}} & \thepage & Omori--Wansing's source letters for false-only neither both and true-only values in the BD algebra. \\
\phantomsection\label{entry:symbol-registry-00557}\(\leq_t,\leq_k\) & \texttt{\detokenize{\leq_t,}\allowbreak\detokenize{\leq_k}} & Belnap--Dunn truth and knowledge orders & P8 relevance-linear & meta & \hyperref[sec:relevance-linear-01]{\ref*{sec:relevance-linear-01}} & \thepage & Omori--Wansing's source subscripts for the distinct truth and knowledge orders on BD; neither is a consequence relation. \\
\phantomsection\label{entry:symbol-registry-00558}\(x\Rightarrow y\) & \texttt{\detokenize{\Rightarrow}} & t-norm residuum and Lukasiewicz implication & P8 relevance-linear & 6 & \hyperref[sec:relevance-linear-01]{\ref*{sec:relevance-linear-01}} & \thepage & Hajek's plain source glyph.  For the Lukasiewicz t-norm it is min of one and one-minus-x-plus-y; numeric operands type it as a truth operation rather than a sequent arrow. \\
\phantomsection\label{entry:symbol-registry-00559}\(x\mathbin{\ast_{\text{\L}}}y,x\mathbin{\ast_{\Pi}}y\) & \texttt{\detokenize{\mathbin{\ast_{\text{\L}}},}\allowbreak\detokenize{\mathbin{\ast_{\Pi}}}} & Lukasiewicz and Product t-norms & P8 relevance-linear & 4 & \hyperref[sec:relevance-linear-01]{\ref*{sec:relevance-linear-01}} & \thepage & Hajek's source subscripts use stroked L for the Lukasiewicz max-zero-sum-minus-one t-norm and Pi for the Product multiplication t-norm. \\
\phantomsection\label{entry:symbol-registry-00560}\(x\ast y\) & \texttt{\detokenize{\ast}} & t-norm operation & P8 relevance-linear & 4 & \hyperref[sec:relevance-linear-01]{\ref*{sec:relevance-linear-01}} & \thepage & Hajek's generic associative commutative monotone binary operation on the unit interval with unit one. \\
\phantomsection\label{entry:symbol-registry-00561}\(\mathbf E,\mathbf R,\mathbf{RM}\) & \texttt{\detokenize{\mathbf E,}\allowbreak\detokenize{\mathbf R,}\allowbreak\detokenize{\mathbf{RM}}} & relevance-system identifiers & P8 paraconsistent-and-nonmonotonic & n/a & \hyperref[sec:p08-paraconsistent-and-nonmonotonic]{\ref*{sec:p08-paraconsistent-and-nonmonotonic}} & \thepage & Bold Roman system typography reproduced from Dunn--Restall's Anderson--Belnap convention; locally typed as system identifiers. \\
\phantomsection\label{entry:symbol-registry-00562}\(\mathbf E_{\to},\mathbf R_{\to}\) & \texttt{\detokenize{\mathbf E_{\to},}\allowbreak\detokenize{\mathbf R_{\to}}} & pure implicational relevance-fragment identifiers & P8 paraconsistent-and-nonmonotonic & n/a & \hyperref[sec:p08-paraconsistent-and-nonmonotonic]{\ref*{sec:p08-paraconsistent-and-nonmonotonic}} & \thepage & Dunn--Restall's source subscripts for the pure implicational fragments of bold Roman E and R. \\
\phantomsection\label{entry:symbol-registry-00563}\(\gamma\) & \texttt{\detokenize{\gamma}} & Ackermann theoremhood metarule & P8 paraconsistent-and-nonmonotonic & meta & \hyperref[sec:p08-paraconsistent-and-nonmonotonic]{\ref*{sec:p08-paraconsistent-and-nonmonotonic}} & \thepage & From theoremhood of A and of not-A-or-B infer theoremhood of B; admissible for E and R but not a primitive or local-consequence rule. \\
\phantomsection\label{entry:symbol-registry-00564}\(\mathfrak F\) & \texttt{\detokenize{\mathfrak F}} & declared frame datum & P8 modal-temporal-epistemic to paraconsistent-and-nonmonotonic & n/a & \hyperref[sec:modal-temporal-epistemic-01]{\ref*{sec:modal-temporal-epistemic-01}} & \thepage & The locally declared tuple fixes whether this is a modal Kripke frame a broad Routley--Meyer frame or the pointed Routley--Meyer frame used by the cited presentation; disjoint from the P6.3 normal-subgroup-filter use. \\
\phantomsection\label{entry:symbol-registry-00565}\(a\leq b\) & \texttt{\detokenize{\leq}} & derived relevance-frame preorder & P8 paraconsistent-and-nonmonotonic & meta & \hyperref[sec:p08-paraconsistent-and-nonmonotonic]{\ref*{sec:p08-paraconsistent-and-nonmonotonic}} & \thepage & Metalinguistic abbreviation for R(0,a,b) in the displayed pointed frame; distinct from the algebraic orders later in the chapter. \\
\phantomsection\label{entry:symbol-registry-00566}\(R^2(ab)cd,R^2a(bc)d\) & \texttt{\detokenize{R^2}} & composite ternary-frame relations & P8 paraconsistent-and-nonmonotonic & meta & \hyperref[sec:p08-paraconsistent-and-nonmonotonic]{\ref*{sec:p08-paraconsistent-and-nonmonotonic}} & \thepage & The first two exact existential composites in the cited E R and RM frame presentation; parenthesization fixes the composition. \\
\phantomsection\label{entry:symbol-registry-00567}\(R^2b(ac)d,R^2abbc\) & \texttt{\detokenize{R^2}} & composite ternary-frame relations & P8 paraconsistent-and-nonmonotonic & meta & \hyperref[sec:p08-paraconsistent-and-nonmonotonic]{\ref*{sec:p08-paraconsistent-and-nonmonotonic}} & \thepage & The remaining exact existential composites in the cited E R and RM frame presentation; parenthesization fixes the composition. \\
\phantomsection\label{entry:symbol-registry-00568}\(\mathsf{IMLL},\mathsf{CLL}\) & \texttt{\detokenize{\mathsf{IMLL},}\allowbreak\detokenize{\mathsf{CLL}}} & linear-calculus identifiers & P8 paraconsistent-and-nonmonotonic & n/a & \hyperref[sec:p08-paraconsistent-and-nonmonotonic]{\ref*{sec:p08-paraconsistent-and-nonmonotonic}} & \thepage & Named multiplicative intuitionistic core and distinct full one-sided classical calculus.  The inherited four-rule bang extension is unnamed and therefore has no fabricated identifier. \\
\phantomsection\label{entry:symbol-registry-00569}\(A\bindnasrepma B\) & \texttt{\detokenize{\bindnasrepma}} & multiplicative disjunction (par) & P8 paraconsistent-and-nonmonotonic & 4 & \hyperref[sec:p08-paraconsistent-and-nonmonotonic]{\ref*{sec:p08-paraconsistent-and-nonmonotonic}} & \thepage & CLL connective dual to tensor; shares the disjunction rank only in this exact scope. \\
\phantomsection\label{entry:symbol-registry-00570}\(A\mathbin{\&}B\) & \texttt{\detokenize{\mathbin{\&}}} & additive conjunction (with) & P8 paraconsistent-and-nonmonotonic & 4 & \hyperref[sec:p08-paraconsistent-and-nonmonotonic]{\ref*{sec:p08-paraconsistent-and-nonmonotonic}} & \thepage & CLL additive conjunction; distinct from metalinguistic ampersands. \\
\phantomsection\label{entry:symbol-registry-00571}\(A\oplus B\) & \texttt{\detokenize{\oplus}} & additive disjunction (plus) & P8 paraconsistent-and-nonmonotonic & 5 & \hyperref[sec:p08-paraconsistent-and-nonmonotonic]{\ref*{sec:p08-paraconsistent-and-nonmonotonic}} & \thepage & CLL additive disjunction; exact P8.6 scope is disjoint from the P4.4 computable join role. \\
\phantomsection\label{entry:symbol-registry-00572}\(0\) & \texttt{\detokenize{0}} & CLL additive false unit & P8 paraconsistent-and-nonmonotonic & n/a & \hyperref[sec:p08-paraconsistent-and-nonmonotonic]{\ref*{sec:p08-paraconsistent-and-nonmonotonic}} & \thepage & Additive unit dual to top in the displayed CLL grammar; not the relevance base point or an ordinary numeral. \\
\phantomsection\label{entry:symbol-registry-00573}\(?A\) & \texttt{\detokenize{?}} & why-not exponential & P8 paraconsistent-and-nonmonotonic & 2.5 & \hyperref[sec:p08-paraconsistent-and-nonmonotonic]{\ref*{sec:p08-paraconsistent-and-nonmonotonic}} & \thepage & CLL exponential dual to bang; prefix unary and the only side supporting the displayed structural rules. \\
\phantomsection\label{entry:symbol-registry-00574}\(A^\perp\) & \texttt{\detokenize{{}^\perp}} & formula-level linear negation & P8 paraconsistent-and-nonmonotonic & 2.5 & \hyperref[sec:p08-paraconsistent-and-nonmonotonic]{\ref*{sec:p08-paraconsistent-and-nonmonotonic}} & \thepage & Involutive De Morgan operation on CLL formulas; distinct from phase orthogonality on subsets. \\
\phantomsection\label{entry:symbol-registry-00575}\(X^\perp\) & \texttt{\detokenize{{}^\perp}} & phase orthogonal of a subset & P8 paraconsistent-and-nonmonotonic & meta & \hyperref[sec:p08-paraconsistent-and-nonmonotonic]{\ref*{sec:p08-paraconsistent-and-nonmonotonic}} & \thepage & Girard's source notation for subset orthogonality relative to the phase-space pole.  Subset operands distinguish it from formula linear negation without altering the glyph. \\
\phantomsection\label{entry:symbol-registry-00576}\(\perp\subseteq M\) & \texttt{\detokenize{\perp}} & Girard phase-space pole & P8 paraconsistent-and-nonmonotonic & n/a & \hyperref[sec:p08-paraconsistent-and-nonmonotonic]{\ref*{sec:p08-paraconsistent-and-nonmonotonic}} & \thepage & Girard's exact perpendicular glyph, locally typed as the distinguished subset of the phase monoid rather than formula linear negation. \\
\phantomsection\label{entry:symbol-registry-00577}\(\sem A\) & \texttt{\detokenize{\sem}} & IMLL algebraic interpretation & P8 paraconsistent-and-nonmonotonic & meta & \hyperref[sec:p08-paraconsistent-and-nonmonotonic]{\ref*{sec:p08-paraconsistent-and-nonmonotonic}} & \thepage & The inherited source's semantic value of formula A in the declared commutative residuated partially ordered monoid. \\
\phantomsection\label{entry:symbol-registry-00578}\(A^*\) & \texttt{\detokenize{{}^*}} & Girard phase interpretation of a CLL formula & P8 paraconsistent-and-nonmonotonic & meta & \hyperref[sec:p08-paraconsistent-and-nonmonotonic]{\ref*{sec:p08-paraconsistent-and-nonmonotonic}} & \thepage & Fact assigned to formula A in Girard's source notation; phase validity is membership of the monoid unit 1 in A-star. \\
\phantomsection\label{entry:symbol-registry-00579}\(I\) & \texttt{\detokenize{I}} & idempotent exponential phase set & P8 paraconsistent-and-nonmonotonic & n/a & \hyperref[sec:p08-paraconsistent-and-nonmonotonic]{\ref*{sec:p08-paraconsistent-and-nonmonotonic}} & \thepage & Idempotent members of the tensor-unit fact used in the displayed full-CLL exponential clauses. \\
\phantomsection\label{entry:symbol-registry-00580}\(A\mathbin{\cdot}B\) & \texttt{\detokenize{\mathbin{\cdot}}} & Lambek product & P8 paraconsistent-and-nonmonotonic & 4 & \hyperref[sec:p08-paraconsistent-and-nonmonotonic]{\ref*{sec:p08-paraconsistent-and-nonmonotonic}} & \thepage & Associative unit-free product in Lambek L; operand type distinguishes it from phase and arithmetic multiplication. \\
\phantomsection\label{entry:symbol-registry-00581}\(A\backslash B,B/A\) & \texttt{\detokenize{\backslash,}\allowbreak\detokenize{/}} & Lambek residuals & P8 paraconsistent-and-nonmonotonic & 6 & \hyperref[sec:p08-paraconsistent-and-nonmonotonic]{\ref*{sec:p08-paraconsistent-and-nonmonotonic}} & \thepage & Left and right divisions of the ordered Lambek product; neither admits exchange weakening or contraction. \\
\phantomsection\label{entry:symbol-registry-00582}\((M,\perp);mn;1\) & \texttt{\detokenize{M,}\allowbreak\detokenize{\perp,}\allowbreak\detokenize{mn,}\allowbreak\detokenize{1}} & Girard phase-space datum and monoid notation & P8 paraconsistent-and-nonmonotonic & n/a & \hyperref[sec:p08-paraconsistent-and-nonmonotonic]{\ref*{sec:p08-paraconsistent-and-nonmonotonic}} & \thepage & Girard's phase space is the pair of a commutative monoid M and perpendicular pole; multiplication is juxtaposition and the monoid unit is 1. \\
\phantomsection\label{entry:symbol-registry-00583}\(X,Y\) & \texttt{\detokenize{X,}\allowbreak\detokenize{Y}} & phase subsets and facts & P8 paraconsistent-and-nonmonotonic & n/a & \hyperref[sec:p08-paraconsistent-and-nonmonotonic]{\ref*{sec:p08-paraconsistent-and-nonmonotonic}} & \thepage & Girard uses the same source metavariables for arbitrary phase subsets and for biorthogonally closed facts; the locally stated closure property supplies the fact role. \\
\phantomsection\label{entry:symbol-registry-00584}\(\mathsf L^*,\mathsf L_1\) & \texttt{\detokenize{\mathsf L^*,}\allowbreak\detokenize{\mathsf L_1}} & Lambek extension identifiers & P8 paraconsistent-and-nonmonotonic & n/a & \hyperref[sec:p08-paraconsistent-and-nonmonotonic]{\ref*{sec:p08-paraconsistent-and-nonmonotonic}} & \thepage & Contrast names for empty-antecedent and unit extensions; neither is the associative unit-free Lambek-L calculus displayed on its card. \\
\phantomsection\label{entry:symbol-registry-00585}\(\lbrack A\rbrack\) & \texttt{\detokenize{\lbrack\cdot\rbrack}} & IMLL Lindenbaum quotient class & P8 paraconsistent-and-nonmonotonic & meta & \hyperref[sec:p08-paraconsistent-and-nonmonotonic]{\ref*{sec:p08-paraconsistent-and-nonmonotonic}} & \thepage & Equivalence class of formula A in the annex sketch's syntactic algebra; distinct from semantic double brackets. \\
\phantomsection\label{entry:symbol-registry-00586}\(\Phi\) & \texttt{\detokenize{\Phi}} & ordered Lambek context metavariable & P8 paraconsistent-and-nonmonotonic & n/a & \hyperref[sec:p08-paraconsistent-and-nonmonotonic]{\ref*{sec:p08-paraconsistent-and-nonmonotonic}} & \thepage & Possibly empty context segment in rule notation subject to every displayed antecedent remaining nonempty. \\
\phantomsection\label{entry:symbol-registry-00587}\(i,j\) & \texttt{\detokenize{i,}\allowbreak\detokenize{j}} & hybrid-logic nominals & P8 conditionals-deontic-free & n/a & \hyperref[sec:p08-conditionals-deontic-free]{\ref*{sec:p08-conditionals-deontic-free}} & \thepage & Object-language atoms constrained to name exactly one world each; disjoint from metalinguistic world labels. \\
\phantomsection\label{entry:symbol-registry-00588}\(@_iA\) & \texttt{\detokenize{@_i}} & hybrid satisfaction operator & P8 conditionals-deontic-free & 3 & \hyperref[sec:p08-conditionals-deontic-free]{\ref*{sec:p08-conditionals-deontic-free}} & \thepage & Evaluates A at the unique world named by nominal i; prefix unary and internal to the object language. \\
\phantomsection\label{entry:symbol-registry-00589}\(@_i\Diamond j\) & \texttt{\detokenize{@_i\Diamond}} & hybrid tableau accessibility atom & P8 conditionals-deontic-free & meta & \hyperref[sec:p08-conditionals-deontic-free]{\ref*{sec:p08-conditionals-deontic-free}} & \thepage & Branch notation asserting that the world named by j is accessible from the world named by i. \\
\phantomsection\label{entry:symbol-registry-00590}\(\mathcal I(c)(\Gamma)\) & \texttt{\detokenize{\mathcal I(c)}} & world-relative interpretation of a quantified-modal constant & P8 modal-temporal-epistemic & n/a & \hyperref[sec:modal-temporal-epistemic-01]{\ref*{sec:modal-temporal-epistemic-01}} & \thepage & The value of constant c at Fitting world Gamma; rigidity means equality of these values at every pair of worlds and is independent of inner-domain existence. \\
\phantomsection\label{entry:symbol-registry-00591}\(\mathsf f,\mathsf u,\mathsf t\) & \texttt{\detokenize{\mathsf f,}\allowbreak\detokenize{\mathsf u,}\allowbreak\detokenize{\mathsf t}} & neutral free-logic truth values & P8 many-valued-paraconsistent & n/a & \hyperref[sec:many-valued-paraconsistent-01]{\ref*{sec:many-valued-paraconsistent-01}} & \thepage & Strong-Kleene false undefined and true values for the partial-term variant; only true is designated. \\
\phantomsection\label{entry:symbol-registry-00592}\(\models_{\mathrm{neg}},\models_{\mathrm{pos}},\models_{\mathrm{neu}}\) & \texttt{\detokenize{\models_{\mathrm{neg}},}\allowbreak\detokenize{\models_{\mathrm{pos}},}\allowbreak\detokenize{\models_{\mathrm{neu}}}} & free-logic variant consequence relations & P8 many-valued-paraconsistent & meta & \hyperref[sec:many-valued-paraconsistent-01]{\ref*{sec:many-valued-paraconsistent-01}} & \thepage & Truth preservation in the negative and positive two-valued rows and designated-t preservation in the neutral strong-Kleene row; the subscript fixes the atomic and partial-term policy. \\
\phantomsection\label{entry:symbol-registry-00593}\(\mathsf{SDL},O A,\mathsf P_OA\) & \texttt{\detokenize{\mathsf{SDL},}\allowbreak\detokenize{O,}\allowbreak\detokenize{\mathsf P_O}} & standard deontic identifier obligation and permission & P8 many-valued-paraconsistent & 3 & \hyperref[sec:many-valued-paraconsistent-01]{\ref*{sec:many-valued-paraconsistent-01}} & \thepage & SDL is KD under the deontic reading; permission is the eliminable abbreviation \(\mathsf P_OA\defeq\neg O\neg A\). \\
\phantomsection\label{entry:symbol-registry-00594}\(O(A\mid B)\) & \texttt{\detokenize{O(\cdot\mid\cdot)}} & dyadic obligation & P8 many-valued-paraconsistent & 6 & \hyperref[sec:many-valued-paraconsistent-01]{\ref*{sec:many-valued-paraconsistent-01}} & \thepage & Binary modality evaluated at the best accessible B-worlds; it is not the monadic formula \(O(B\to A)\). \\
\phantomsection\label{entry:symbol-registry-00595}\(\preceq_w,\operatorname{Best}_w(B)\) & \texttt{\detokenize{\preceq_w,}\allowbreak\detokenize{\operatorname{Best}_w}} & world-relative ideality preorder and best-world selector & P8 many-valued-paraconsistent & n/a & \hyperref[sec:many-valued-paraconsistent-01]{\ref*{sec:many-valued-paraconsistent-01}} & \thepage & A smooth total preorder on accessible worlds and the explicitly defined set of its minimal accessible B-worlds; the set is empty when there are no candidates. \\
\phantomsection\label{entry:symbol-registry-00596}\(A>B\) & \texttt{\detokenize{>}} & conditional connective & P8 many-valued-paraconsistent & 6 & \hyperref[sec:many-valued-paraconsistent-01]{\ref*{sec:many-valued-paraconsistent-01}} & \thepage & Formula operands type this glyph as the primitive CK or Stalnaker conditional; numerical operands retain the strict-order role. \\
\phantomsection\label{entry:symbol-registry-00597}\(A\Box\!\!\to B\) & \texttt{\detokenize{\Box\!\!\to}} & Lewis variably strict conditional & P8 many-valued-paraconsistent & 6 & \hyperref[sec:many-valued-paraconsistent-01]{\ref*{sec:many-valued-paraconsistent-01}} & \thepage & Evaluated by nested similarity spheres; distinct from material implication and the unique-selection conditional. \\
\phantomsection\label{entry:symbol-registry-00598}\(f(w,X),F_w(X)\) & \texttt{\detokenize{f,}\allowbreak\detokenize{F_w}} & singular and set-valued conditional selection functions & P8 many-valued-paraconsistent & n/a & \hyperref[sec:many-valued-paraconsistent-01]{\ref*{sec:many-valued-paraconsistent-01}} & \thepage & Stalnaker f selects one member of a nonempty proposition; CK \(F_w\) returns an arbitrary set and need not be contained in X. \\
\phantomsection\label{entry:symbol-registry-00599}\(\mathcal S_w,U_w\) & \texttt{\detokenize{\mathcal S_w,}\allowbreak\detokenize{U_w}} & Lewis sphere system and accessible union & P8 many-valued-paraconsistent & n/a & \hyperref[sec:many-valued-paraconsistent-01]{\ref*{sec:many-valued-paraconsistent-01}} & \thepage & The inclusion-nested union/intersection-closed sphere family at w and its union; no least antecedent-permitting sphere is assumed. \\
\phantomsection\label{entry:symbol-registry-00600}\(\mathsf{CK},\mathsf V,\mathsf{VW},\mathsf{VC},\mathsf{C1},\mathsf{C2}\) & \texttt{\detokenize{\mathsf{CK},}\allowbreak\detokenize{\mathsf V,}\allowbreak\detokenize{\mathsf{VW},}\allowbreak\detokenize{\mathsf{VC},}\allowbreak\detokenize{\mathsf{C1},}\allowbreak\detokenize{\mathsf{C2}}} & conditional-logic identifiers & P8 many-valued-paraconsistent & n/a & \hyperref[sec:many-valued-paraconsistent-01]{\ref*{sec:many-valued-paraconsistent-01}} & \thepage & Basic conditional logic and the separately stated Lewis and Stalnaker extensions; VC is also called C1 in the cited nomenclature. \\
\phantomsection\label{entry:symbol-registry-00601}\(P(B\mid A)\) & \texttt{\detokenize{P(\cdot\mid\cdot)}} & conditional probability & P8 many-valued-paraconsistent & meta & \hyperref[sec:many-valued-paraconsistent-01]{\ref*{sec:many-valued-paraconsistent-01}} & \thepage & Defined when \(P(A)>0\); probability P is distinct from the deontic permission abbreviation \(\mathsf P_O\). \\
\phantomsection\label{entry:symbol-registry-00602}\(A\mathbin{\&}B\) & \texttt{\detokenize{\mathbin{\&}}} & strong t-norm conjunction & P8 relevance-linear & 4 & \hyperref[sec:relevance-linear-01]{\ref*{sec:relevance-linear-01}} & \thepage & Primitive multiplicative conjunction of the MTL and BL-family sheets; distinct from weak lattice conjunction and CLL additive with. \\
\phantomsection\label{entry:symbol-registry-00603}\(\mathsf L_n,\mathsf L_{\infty}\) & \texttt{\detokenize{\mathsf L_n,}\allowbreak\detokenize{\mathsf L_{\infty}}} & finite and infinite Lukasiewicz system identifiers & P8 relevance-linear & n/a & \hyperref[sec:relevance-linear-01]{\ref*{sec:relevance-linear-01}} & \thepage & Editorial system-card identifiers; the source symbols \(\text{\L}_n\) and \(\text{\L}\) name the finite chain and standard infinite logic. \\
\phantomsection\label{entry:symbol-registry-00604}\(\derives_{\mathrm{MTL}},\derives_{\mathrm{BL}},\derives_G,\derives_{\Pi}\) & \texttt{\detokenize{\derives_{\mathrm{MTL}},}\allowbreak\detokenize{\derives_{\mathrm{BL}},}\allowbreak\detokenize{\derives_G,}\allowbreak\detokenize{\derives_{\Pi}}} & t-norm Hilbert derivability relations & P8 relevance-linear & meta & \hyperref[sec:relevance-linear-01]{\ref*{sec:relevance-linear-01}} & \thepage & The subscript fixes one exact sheet; their standard completeness strengths are stated separately. \\
\phantomsection\label{entry:symbol-registry-00605}\(\mathsf C_1,\derives_{\mathsf C_1}\) & \texttt{\detokenize{\mathsf C_1,}\allowbreak\detokenize{\derives_{\mathsf C_1}}} & da Costa C1 calculus and derivability & P8 further-reading & meta & \hyperref[sec:p08-further-reading]{\ref*{sec:p08-further-reading}} & \thepage & The source-fixed propositional calculus on sheet:hilbert-da-costa-c1; no claim is transferred to every member of the da Costa hierarchy. \\
\phantomsection\label{entry:symbol-registry-00606}\(A^\circ\) & \texttt{\detokenize{{}^\circ}} & da Costa C1 consistency formula & P8 further-reading & 2.5 & \hyperref[sec:p08-further-reading]{\ref*{sec:p08-further-reading}} & \thepage & Eliminable abbreviation \(\neg(A\land\neg A)\) expressing formula-level well-behavedness in C1; not a matrix truth value or a metalinguistic theory-consistency claim. \\
\phantomsection\label{entry:symbol-registry-00607}\(\mathrel{|\!\sim}\) & \texttt{\detokenize{|\!\sim}} & defeasible consequence & P8 further-reading & meta & \hyperref[sec:p08-further-reading]{\ref*{sec:p08-further-reading}} & \thepage & A relation between information states or formulas whose conclusions may be defeated by additional premises. \\
\phantomsection\label{entry:symbol-registry-00608}\((S,\ell,\prec),\mu(X)\) & \texttt{\detokenize{S,}\allowbreak\detokenize{\ell,}\allowbreak\detokenize{\prec,}\allowbreak\detokenize{\mu}} & preferential-model datum and minimal-state operator & P8 further-reading & meta & \hyperref[sec:p08-further-reading]{\ref*{sec:p08-further-reading}} & \thepage & The labeling map gives classical valuations; \(\mu(X)\) selects the \(\prec\)-minimal states of X. \\
\phantomsection\label{entry:symbol-registry-00609}\(K*A,K+A,K-A\) & \texttt{\detokenize{*,}\allowbreak\detokenize{+,}\allowbreak\detokenize{-}} & AGM revision expansion and contraction & P8 further-reading & meta & \hyperref[sec:p08-further-reading]{\ref*{sec:p08-further-reading}} & \thepage & Operations on the deductively closed belief set K; their standing is fixed by the AGM postulates and Levi and Harper identities. \\
\phantomsection\label{entry:symbol-registry-00610}\(\preceq_K,\min([A],\preceq_K)\) & \texttt{\detokenize{\preceq_K,}\allowbreak\detokenize{\min}} & AGM faithful preorder and minimal input valuations & P8 further-reading & meta & \hyperref[sec:p08-further-reading]{\ref*{sec:p08-further-reading}} & \thepage & A total preorder assigned to prior belief set K; the minimum selects its most preferred A-valuations. \\
\phantomsection\label{entry:symbol-registry-00611}\(\frac{\alpha:\beta_1,\ldots,\beta_n}{\gamma}\) & \texttt{\detokenize{:,}\allowbreak\detokenize{/}} & Reiter default & P8 further-reading & meta & \hyperref[sec:p08-further-reading]{\ref*{sec:p08-further-reading}} & \thepage & Prerequisite alpha justifications beta-i and conclusion gamma; this fraction is not an ordinary inference rule. \\
\phantomsection\label{entry:symbol-registry-00612}\((W,D),\Gamma_E,U_k\) & \texttt{\detokenize{W,}\allowbreak\detokenize{D,}\allowbreak\detokenize{\Gamma_E,}\allowbreak\detokenize{U_k}} & default theory grounded closure and stages & P8 further-reading & meta & \hyperref[sec:p08-further-reading]{\ref*{sec:p08-further-reading}} & \thepage & Gamma-E is the least closure grounded through the stages U-k while justification consistency is tested against candidate E. \\
\phantomsection\label{entry:symbol-registry-00613}\(T(\bar P,\bar Q,\bar Z),\bar P'\leq\bar P\) & \texttt{\detokenize{T,}\allowbreak\detokenize{\bar P,}\allowbreak\detokenize{\bar Q,}\allowbreak\detokenize{\bar Z,}\allowbreak\detokenize{\leq}} & circumscription theory and predicate-tuple inclusion & P8 further-reading & meta & \hyperref[sec:p08-further-reading]{\ref*{sec:p08-further-reading}} & \thepage & P is minimized Q is fixed and Z may vary; tuple inclusion abbreviates componentwise universal closure. \\
\phantomsection\label{entry:symbol-registry-00614}\(LA,\operatorname{Th}\) & \texttt{\detokenize{L,}\allowbreak\detokenize{\operatorname{Th}}} & autoepistemic belief and classical closure & P8 further-reading & 3 & \hyperref[sec:p08-further-reading]{\ref*{sec:p08-further-reading}} & \thepage & L is the belief operator; Th is ordinary classical propositional closure over the autoepistemic language and adds no modal axiom. \\
\phantomsection\label{entry:symbol-registry-01909}\(\mathsf{Fm}_{\Box}\) & \texttt{\detokenize{\mathsf{Fm}_{\Box}}} & set of finite modal formulas & P8 modal canonical extensions & n/a & \hyperref[def:p08-mce-language]{\ref*{def:p08-mce-language}} & \thepage & Existing countable propositional alphabet with primitive Boolean connectives and box. \\
\phantomsection\label{entry:symbol-registry-01910}\(L\) & \texttt{\detokenize{L}} & selected normal Hilbert system & P8 modal canonical extensions & n/a & \hyperref[def:p08-mce-language]{\ref*{def:p08-mce-language}} & \thepage & Ranges over KD KTB and K4 with their existing sheets. \\
\phantomsection\label{entry:symbol-registry-01911}\(\mathcal C_L\) & \texttt{\detokenize{\mathcal C_L}} & class of nonempty frames for L & P8 modal canonical extensions & n/a & \hyperref[def:p08-mce-frame-classes]{\ref*{def:p08-mce-frame-classes}} & \thepage & Serial for KD; reflexive symmetric for KTB; transitive for K4. \\
\phantomsection\label{entry:symbol-registry-01912}\(W_L\) & \texttt{\detokenize{W_L}} & set of maximal L-consistent theories & P8 modal canonical extensions & n/a & \hyperref[def:p08-mce-canonical-model]{\ref*{def:p08-mce-canonical-model}} & \thepage & Subset of the powerset of the countable formula set; may itself be uncountable. \\
\phantomsection\label{entry:symbol-registry-01913}\(U_L(w)\) & \texttt{\detokenize{U_L(w)}} & contents of formulas boxed at canonical world w & P8 modal canonical extensions & n/a & \hyperref[def:p08-mce-canonical-model]{\ref*{def:p08-mce-canonical-model}} & \thepage & The set of C for which Box C belongs to w. \\
\phantomsection\label{entry:symbol-registry-01914}\(R_L\) & \texttt{\detokenize{R_L}} & canonical accessibility relation & P8 modal canonical extensions & n/a & \hyperref[def:p08-mce-canonical-model]{\ref*{def:p08-mce-canonical-model}} & \thepage & \(wR_Lv\) iff \(U_L(w)\) is contained in \(v\); distinct from derivability. \\
\phantomsection\label{entry:symbol-registry-01915}\(V_L\) & \texttt{\detokenize{V_L}} & canonical atomic valuation & P8 modal canonical extensions & n/a & \hyperref[def:p08-mce-canonical-model]{\ref*{def:p08-mce-canonical-model}} & \thepage & Maps an atom to all maximal theories containing it. \\
\phantomsection\label{entry:symbol-registry-01916}\(M_L\) & \texttt{\detokenize{M_L}} & canonical Kripke structure & P8 modal canonical extensions & n/a & \hyperref[def:p08-mce-canonical-model]{\ref*{def:p08-mce-canonical-model}} & \thepage & Triple of canonical worlds relation and valuation; nonemptiness proved separately. \\
\phantomsection\label{entry:symbol-registry-01917}\(A_n\) & \texttt{\detokenize{A_n}} & enumerated modal formula & P8 modal canonical extensions annex & n/a & \hyperref[annex:p8.mce.3]{\ref*{annex:p8.mce.3}} & \thepage & Fixed enumeration indexed by natural numbers. \\
\phantomsection\label{entry:symbol-registry-01918}\(X_n\) & \texttt{\detokenize{X_n}} & Lindenbaum extension stage & P8 modal canonical extensions annex & n/a & \hyperref[annex:p8.mce.3]{\ref*{annex:p8.mce.3}} & \thepage & Increasing consistent sets beginning with the entire arbitrary premise set. \\
\phantomsection\label{entry:symbol-registry-02440}\(W,s,t\) & \texttt{\detokenize{W,}\allowbreak\detokenize{s,}\allowbreak\detokenize{t}} & world universe and information states & P8 inquisitive-propositional & n/a & \hyperref[def:p08-inq-states]{\ref*{def:p08-inq-states}} & \thepage & States are subsets of W; t subset s means informational enhancement. W may be empty. \\
\phantomsection\label{entry:symbol-registry-02441}\(\mathcal A\) & \texttt{\detokenize{\mathcal A}} & nonempty downward-closed family of information states & P8 inquisitive-propositional & n/a & \hyperref[def:p08-inq-propositions]{\ref*{def:p08-inq-propositions}} & \thepage & An inquisitive proposition contains the empty state. \\
\phantomsection\label{entry:symbol-registry-02442}\(T^\downarrow\) & \texttt{\detokenize{T^\downarrow}} & downward closure of a family of states & P8 inquisitive-propositional & n/a & \hyperref[def:p08-inq-propositions]{\ref*{def:p08-inq-propositions}} & \thepage & Generator closure; not downward closure of a logical theory. \\
\phantomsection\label{entry:symbol-registry-02443}\(\operatorname{Alt}(\mathcal A)\) & \texttt{\detokenize{\operatorname{Alt}}} & inclusion-maximal supporting states & P8 inquisitive-propositional & n/a & \hyperref[def:p08-inq-propositions]{\ref*{def:p08-inq-propositions}} & \thepage & Alternatives need not exist for arbitrary infinite inquisitive propositions; all InqB formulas have finite generators. \\
\phantomsection\label{entry:symbol-registry-02444}\([\varphi]_M\) & \texttt{\detokenize{[\varphi]_M}} & support proposition of a formula in a model & P8 inquisitive-propositional & n/a & \hyperref[def:p08-inq-general-entailment]{\ref*{def:p08-inq-general-entailment}} & \thepage & Set of states; distinct from the truth set and a proof-equivalence class. \\
\phantomsection\label{entry:symbol-registry-02445}\(|\varphi|_M\) & \texttt{\detokenize{|\varphi|_M}} & source truth-set notation & P8 inquisitive-propositional & n/a & \hyperref[def:p08-inq-general-entailment]{\ref*{def:p08-inq-general-entailment}} & \thepage & Typed source notation for worlds supporting the singleton formula; not cardinality. \\
\phantomsection\label{entry:symbol-registry-02446}\(\models_s,\models_M\) & \texttt{\detokenize{\models_s,}\allowbreak\detokenize{\models_M}} & support entailment in a context or whole model & P8 inquisitive-propositional & n/a & \hyperref[def:p08-inq-general-entailment]{\ref*{def:p08-inq-general-entailment}} & \thepage & Contextual entailment quantifies over all substates; global entailment quantifies over models as well. \\
\phantomsection\label{entry:symbol-registry-02447}\(\equiv_{\mathrm{InqB}}\) & \texttt{\detokenize{\equiv_{\mathrm{InqB}}}} & mutual support entailment & P8 inquisitive-propositional & n/a & \hyperref[def:p08-inq-general-entailment]{\ref*{def:p08-inq-general-entailment}} & \thepage & Locally typed formula equivalence; not elementary equivalence of structures. \\
\phantomsection\label{entry:symbol-registry-02448}\(P\) & \texttt{\detokenize{P}} & countable set of propositional atoms & P8 inquisitive expressiveness & n/a & \hyperref[def:p08-inq-language]{\ref*{def:p08-inq-language}} & \thepage & The atom set is used in the language, valuation and finite-profile declarations. Proof P has its own local natural-deduction scope. \\
\phantomsection\label{entry:symbol-registry-02449}\(L_c\) & \texttt{\detokenize{L_c}} & syntactically classical formula fragment & P8 inquisitive-propositional & n/a & \hyperref[def:p08-inq-resolutions]{\ref*{def:p08-inq-resolutions}} & \thepage & The literal \(L_c\) first appears as the codomain of a resolution selection, where it is explicitly declared to be the classical fragment. Classical means inquisitive-disjunction-free after expanding defined connectives. \\
\phantomsection\label{entry:symbol-registry-02450}\(\mathbin{\vee\!\vee}\) & \texttt{\detokenize{\mathbin{\vee\!\vee}}} & inquisitive disjunction & P8 inquisitive-propositional & n/a & \hyperref[def:p08-inq-language]{\ref*{def:p08-inq-language}} & \thepage & Doubled-vee composite reproduces the source connective without defining a shared style macro; support requires one whole-state disjunct. \\
\phantomsection\label{entry:symbol-registry-02451}\(\lor\) & \texttt{\detokenize{\lor}} & defined classical disjunction & P8 inquisitive-propositional & n/a & \hyperref[def:p08-inq-language]{\ref*{def:p08-inq-language}} & \thepage & Defined by negation and conjunction; always flat even when operands are questions. \\
\phantomsection\label{entry:symbol-registry-02452}\(?\varphi\) & \texttt{\detokenize{?\varphi}} & question-mark operator & P8 inquisitive-propositional & n/a & \hyperref[def:p08-inq-language]{\ref*{def:p08-inq-language}} & \thepage & Defined as inquisitive disjunction of phi and its negation; a polar question is formed from a statement. \\
\phantomsection\label{entry:symbol-registry-02453}\(M,s\models\varphi\) & \texttt{\detokenize{M,}\allowbreak\detokenize{s\models\varphi}} & support relation & P8 inquisitive-propositional & n/a & \hyperref[def:p08-inq-general-entailment]{\ref*{def:p08-inq-general-entailment}} & \thepage & World truth is singleton support; empty state supports every formula. \\
\phantomsection\label{entry:symbol-registry-02454}\(\varphi^{\mathrm{cl}}\) & \texttt{\detokenize{\varphi^{\mathrm{cl}}}} & classical variant and presupposition formula & P8 inquisitive-propositional & n/a & \hyperref[def:p08-inq-flatness]{\ref*{def:p08-inq-flatness}} & \thepage & Replace each inquisitive disjunction by classical disjunction; preserves truth but may lose support information. \\
\phantomsection\label{entry:symbol-registry-02455}\(R(\varphi),R(\Phi)\) & \texttt{\detokenize{R(\varphi),}\allowbreak\detokenize{R(\Phi)}} & syntactic resolutions of a formula or premise set & P8 inquisitive-propositional & n/a & \hyperref[def:p08-inq-resolutions]{\ref*{def:p08-inq-resolutions}} & \thepage & Each formula has finitely many nonempty classical resolutions; a set resolution selects one for each premise. \\
\phantomsection\label{entry:symbol-registry-02456}\(\gamma_f\) & \texttt{\detokenize{\gamma_f}} & classical conjunction encoding a resolution-transfer function & P8 inquisitive-propositional & n/a & \hyperref[def:p08-inq-resolutions]{\ref*{def:p08-inq-resolutions}} & \thepage & Conjoins alpha implies f(alpha) for every antecedent resolution. \\
\phantomsection\label{entry:symbol-registry-02457}\(f\) & \texttt{\detokenize{f}} & map between resolution sets; a dependence function when the stated entailment condition holds & P8 inquisitive natural deduction & n/a & \hyperref[def:p08-inq-resolutions]{\ref*{def:p08-inq-resolutions}} & \thepage & The recursive implication-resolution clause introduces arbitrary \(f:R(\varphi)\to R(\psi)\). Its additional semantic dependence condition is defined later at def:p08-inq-logical-dependence. The finite code-table content label f is outside this scope. \\
\phantomsection\label{entry:symbol-registry-02458}\(P:\Phi\vdash_{\mathrm{InqB}}\psi\) & \texttt{\detokenize{\vdash_{\mathrm{InqB}}}} & labelled natural-deduction proof and derivability & P8 inquisitive natural deduction & n/a & \hyperref[def:p08-inq-proofs]{\ref*{def:p08-inq-proofs}} & \thepage & Here P denotes a proof rather than the earlier atom set; discharge is source-specific and may be vacuous. \\
\phantomsection\label{entry:symbol-registry-02459}\(F_P\) & \texttt{\detokenize{F_P}} & resolution algorithm extracted from a proof & P8 inquisitive natural deduction & n/a & \hyperref[thm:p08-inq-proof-resolution]{\ref*{thm:p08-inq-proof-resolution}} & \thepage & Returns a classical proof of one conclusion resolution from supplied assumption resolutions. \\
\phantomsection\label{entry:symbol-registry-02460}\(M^c,W^c,V^c\) & \texttt{\detokenize{M^c,}\allowbreak\detokenize{W^c,}\allowbreak\detokenize{V^c}} & canonical model on complete classical theories & P8 inquisitive completeness & n/a & \hyperref[def:p08-inq-canonical-model]{\ref*{def:p08-inq-canonical-model}} & \thepage & The empty intersection of worlds is all classical formulas. \\
\phantomsection\label{entry:symbol-registry-02461}\(\Omega_Q,\chi_v,\chi_A\) & \texttt{\detokenize{\Omega_Q,}\allowbreak\detokenize{\chi_v,}\allowbreak\detokenize{\chi_A}} & finite valuation space and characteristic classical formulas & P8 inquisitive expressiveness & n/a & \hyperref[def:p08-inq-profiles]{\ref*{def:p08-inq-profiles}} & \thepage & Empty conjunction is top; empty disjunction is bottom; Q may be empty. \\
\phantomsection\label{entry:symbol-registry-02462}\(d(\varphi)\) & \texttt{\detokenize{d(\varphi)}} & coherence degree & P8 inquisitive coherence & n/a & \hyperref[def:p08-inq-coherence]{\ref*{def:p08-inq-coherence}} & \thepage & Least positive bound on the size of substates needed to detect failure of support. \\
\phantomsection\label{entry:symbol-registry-02463}\(D(m)\) & \texttt{\detokenize{D(m)}} & Dedekind antichain count for the subset lattice of an m-element set & P8 inquisitive expressiveness & n/a & \hyperref[thm:p08-inq-expressive]{\ref*{thm:p08-inq-expressive}} & \thepage & The InqB class count over \(n\) atoms is \(D(2^n)-1\); the empty antichain is excluded. \\
\phantomsection\label{entry:symbol-registry-02464}\(\mathsf{InqC}\) & \texttt{\detokenize{\mathsf{InqC}}} & polar-operator inquisitive propositional fragment & P8 inquisitive coherence & n/a & \hyperref[thm:p08-inq-coherence]{\ref*{thm:p08-inq-coherence}} & \thepage & Primitives bottom/conjunction/implication/question mark; not the full doubled-vee language. \\
\phantomsection\label{entry:symbol-registry-02465}\(\otimes\) & \texttt{\detokenize{\otimes}} & tensor disjunction on information states & P8 inquisitive tensor extension & n/a & \hyperref[def:p08-inq-tensor]{\ref*{def:p08-inq-tensor}} & \thepage & A state splits as a union of operand-supporting states; distinct from the P8 linear-logic tensor. \\
\phantomsection\label{entry:symbol-registry-02466}\(R_i,C_i,A_i\) & \texttt{\detokenize{R_i,}\allowbreak\detokenize{C_i,}\allowbreak\detokenize{A_i}} & row/column/digit-occurrence alternatives in the code example & P8 inquisitive finite code example & n/a & \hyperref[def:p08-inq-code-table]{\ref*{def:p08-inq-code-table}} & \thepage & These are sets of code worlds; not syntactic resolutions R(phi). \\
\phantomsection\label{entry:symbol-registry-02467}\(P_0,P_1,Q_0,Q_1\) & \texttt{\detokenize{P_0,}\allowbreak\detokenize{P_1,}\allowbreak\detokenize{Q_0,}\allowbreak\detokenize{Q_1}} & truth-value blocks in the four-world example & P8 inquisitive finite four-world example & n/a & \hyperref[def:p08-inq-four-world-calculations]{\ref*{def:p08-inq-four-world-calculations}} & \thepage & First coordinate gives p and second q. \\
\phantomsection\label{entry:symbol-registry-02596}\(\mathcal F\) & \texttt{\detokenize{\mathcal F}} & set of finite IPL formulas over the specified atom set & \hyperref[sec:p08-ipl-kripke-completeness]{\ref*{sec:p08-ipl-kripke-completeness}} & n/a & \hyperref[def:p08-ipl-completeness-setting]{\ref*{def:p08-ipl-completeness-setting}} & \thepage & Primitive bottom, conjunction, disjunction and implication; negation abbreviates implication to bottom. \\
\phantomsection\label{entry:symbol-registry-02597}\(\Gamma\vdash_{\mathsf{IPL}}A\) & \texttt{\detokenize{\derives}} & finite-support propositional LJ derivability from a formula set & \hyperref[sec:p08-ipl-kripke-completeness]{\ref*{sec:p08-ipl-kripke-completeness}} & meta & \hyperref[def:p08-ipl-completeness-setting]{\ref*{def:p08-ipl-completeness-setting}} & \thepage & A finite multiset over Gamma witnesses the relation; arbitrary premise sets remain permitted. \\
\phantomsection\label{entry:symbol-registry-02598}\(\Gamma_0\Rightarrow A\) & \texttt{\detokenize{\Rightarrow}} & single-succedent LJ sequent with finite-multiset antecedent & \hyperref[sec:p08-ipl-kripke-completeness]{\ref*{sec:p08-ipl-kripke-completeness}} & meta & \hyperref[def:p08-ipl-completeness-setting]{\ref*{def:p08-ipl-completeness-setting}} & \thepage & Exactly one formula in the succedent; exchange is implicit in multisets. \\
\phantomsection\label{entry:symbol-registry-02599}\(u,v,w\) & \texttt{\detokenize{u,}\allowbreak\detokenize{v,}\allowbreak\detokenize{w}} & worlds of the specified intuitionistic Kripke model & \hyperref[sec:p08-ipl-kripke-completeness]{\ref*{sec:p08-ipl-kripke-completeness}} & n/a & \hyperref[def:p08-ipl-completeness-semantics]{\ref*{def:p08-ipl-completeness-semantics}} & \thepage & Declared worlds; the finite-word convention belongs to a disjoint P4 scope. \\
\phantomsection\label{entry:symbol-registry-02600}\(\leq\) & \texttt{\detokenize{\leq}} & reflexive transitive persistence order on worlds & \hyperref[sec:p08-ipl-kripke-completeness]{\ref*{sec:p08-ipl-kripke-completeness}} & relation & \hyperref[def:p08-ipl-completeness-semantics]{\ref*{def:p08-ipl-completeness-semantics}} & \thepage & Original theorem uses partial orders; the explicitly compared variant permits preorders. \\
\phantomsection\label{entry:symbol-registry-02601}\(\mathcal M,w\Vdash A\) & \texttt{\detokenize{\forces}} & intuitionistic Kripke forcing of a formula at a world & \hyperref[sec:p08-ipl-kripke-completeness]{\ref*{sec:p08-ipl-kripke-completeness}} & meta & \hyperref[def:p08-ipl-completeness-semantics]{\ref*{def:p08-ipl-completeness-semantics}} & \thepage & Defined implication quantifies over all future worlds; no world forces bottom. \\
\phantomsection\label{entry:symbol-registry-02602}\(\mathcal M,w\Vdash\Gamma\) & \texttt{\detokenize{\forces}} & forcing every member of a formula set at the same world & \hyperref[sec:p08-ipl-kripke-completeness]{\ref*{sec:p08-ipl-kripke-completeness}} & meta & \hyperref[def:p08-ipl-completeness-semantics]{\ref*{def:p08-ipl-completeness-semantics}} & \thepage & Pointwise extension to arbitrary premise sets. \\
\phantomsection\label{entry:symbol-registry-02603}\(\Gamma\models_{\mathsf{IPL}}A\) & \texttt{\detokenize{\entails}} & local consequence over all persistent partially ordered Kripke models & \hyperref[sec:p08-ipl-kripke-completeness]{\ref*{sec:p08-ipl-kripke-completeness}} & meta & \hyperref[def:p08-ipl-completeness-semantics]{\ref*{def:p08-ipl-completeness-semantics}} & \thepage & Preorder and rooted-frame variants are named separately and proved equivalent. \\
\phantomsection\label{entry:symbol-registry-02604}\(\operatorname{Cn}_{\mathsf{IPL}}\) & \texttt{\detokenize{\operatorname{Cn}}} & deductive closure on subsets of the IPL formula universe & \hyperref[sec:p08-ipl-prime-extension]{\ref*{sec:p08-ipl-prime-extension}} & n/a & \hyperref[def:p08-ipl-prime-theory]{\ref*{def:p08-ipl-prime-theory}} & \thepage & Extensive, monotone and idempotent by finite-support substitution of LJ derivations. \\
\phantomsection\label{entry:symbol-registry-02605}\(S\) & \texttt{\detokenize{S}} & formula set tested for avoidance of the target C & \hyperref[sec:p08-ipl-prime-extension]{\ref*{sec:p08-ipl-prime-extension}} & n/a & \hyperref[def:p08-ipl-prime-theory]{\ref*{def:p08-ipl-prime-theory}} & \thepage & A set avoids C when it does not derive C; maximality ranges over all extensions inside the fixed formula universe. \\
\phantomsection\label{entry:symbol-registry-02606}\(D_n\) & \texttt{\detokenize{D_n}} & nth formula in a supplied enumeration & \hyperref[thm:p08-ipl-enumerated-extension]{\ref*{thm:p08-ipl-enumerated-extension}} & n/a & \hyperref[thm:p08-ipl-enumerated-extension]{\ref*{thm:p08-ipl-enumerated-extension}} & \thepage & The enumeration is assumed as data; no effective derivability test is claimed. \\
\phantomsection\label{entry:symbol-registry-02607}\(S_n\) & \texttt{\detokenize{S_n}} & nth target-avoiding formula set in the enumerated construction & \hyperref[thm:p08-ipl-enumerated-extension]{\ref*{thm:p08-ipl-enumerated-extension}} & n/a & \hyperref[thm:p08-ipl-enumerated-extension]{\ref*{thm:p08-ipl-enumerated-extension}} & \thepage & Each successor is uniquely specified by a classical derivability case distinction. \\
\phantomsection\label{entry:symbol-registry-02608}\(\mathcal Q,\mathcal C,S^*\) & \texttt{\detokenize{\mathcal Q,}\allowbreak\detokenize{\mathcal C,}\allowbreak\detokenize{S^*}} & poset of target-avoiding extensions, an inclusion chain, and its union & \hyperref[annex:p8.ipl.complete.prime-extension]{\ref*{annex:p8.ipl.complete.prime-extension}} & n/a & \hyperref[annex:p8.ipl.complete.prime-extension]{\ref*{annex:p8.ipl.complete.prime-extension}} & \thepage & Proof-local Zorn construction; the empty chain has Gamma as upper bound. \\
\phantomsection\label{entry:symbol-registry-02609}\(W_{T_0},V_{T_0},\mathcal M_{T_0}\) & \texttt{\detokenize{W_{T_0},}\allowbreak\detokenize{V_{T_0},}\allowbreak\detokenize{\mathcal M_{T_0}}} & world set, atom valuation and canonical model above a prime root theory & \hyperref[sec:p08-ipl-canonical-countermodel]{\ref*{sec:p08-ipl-canonical-countermodel}} & n/a & \hyperref[def:p08-ipl-canonical-cone]{\ref*{def:p08-ipl-canonical-cone}} & \thepage & Contains all prime extensions of the root; later worlds need not omit the root's target formula. \\
\phantomsection\label{entry:symbol-registry-02610}\(\Sigma\) & \texttt{\detokenize{\Sigma}} & finite subformula-closed set used for filtration & \hyperref[sec:p08-ipl-finite-filtration]{\ref*{sec:p08-ipl-finite-filtration}} & n/a & \hyperref[def:p08-ipl-filtration]{\ref*{def:p08-ipl-filtration}} & \thepage & A finite set of formulas, not a multiset sequent antecedent. \\
\phantomsection\label{entry:symbol-registry-02611}\(\tau_\Sigma(w)\) & \texttt{\detokenize{\tau_\Sigma}} & the set of formulas from Sigma forced at w & \hyperref[sec:p08-ipl-finite-filtration]{\ref*{sec:p08-ipl-finite-filtration}} & n/a & \hyperref[def:p08-ipl-filtration]{\ref*{def:p08-ipl-filtration}} & \thepage & Truth profile is a subset of Sigma; tau is not a sentence metavariable in this scope. \\
\phantomsection\label{entry:symbol-registry-02612}\(w\sim_\Sigma v\) & \texttt{\detokenize{\sim_\Sigma}} & equality of finite subformula truth profiles & \hyperref[sec:p08-ipl-finite-filtration]{\ref*{sec:p08-ipl-finite-filtration}} & relation & \hyperref[def:p08-ipl-filtration]{\ref*{def:p08-ipl-filtration}} & \thepage & Equivalence relation used to define the filtration quotient. \\
\phantomsection\label{entry:symbol-registry-02613}\([w]_\Sigma\) & \texttt{\detokenize{[w]_\Sigma}} & equivalence class of a world under equal Sigma truth profiles & \hyperref[sec:p08-ipl-finite-filtration]{\ref*{sec:p08-ipl-finite-filtration}} & n/a & \hyperref[def:p08-ipl-filtration]{\ref*{def:p08-ipl-filtration}} & \thepage & This denotes a set of worlds; it is not a proposition or formula equivalence class. \\
\phantomsection\label{entry:symbol-registry-02614}\(\leq_\Sigma\) & \texttt{\detokenize{\leq_\Sigma}} & partial order on filtered worlds defined by inclusion of truth profiles & \hyperref[sec:p08-ipl-finite-filtration]{\ref*{sec:p08-ipl-finite-filtration}} & relation & \hyperref[def:p08-ipl-filtration]{\ref*{def:p08-ipl-filtration}} & \thepage & Comparison includes implication formulas in Sigma and may add pairs absent from the original frame. \\
\phantomsection\label{entry:symbol-registry-02615}\(W_\Sigma,V_\Sigma,\mathcal M_\Sigma\) & \texttt{\detokenize{W_\Sigma,}\allowbreak\detokenize{V_\Sigma,}\allowbreak\detokenize{\mathcal M_\Sigma}} & filtered world set, persistent atom valuation and finite model & \hyperref[sec:p08-ipl-finite-filtration]{\ref*{sec:p08-ipl-finite-filtration}} & n/a & \hyperref[def:p08-ipl-filtration]{\ref*{def:p08-ipl-filtration}} & \thepage & Atoms outside Sigma have the empty valuation; the model is defined for the entire language. \\
\phantomsection\label{entry:symbol-registry-02616}\(m,P_0\) & \texttt{\detokenize{m,}\allowbreak\detokenize{P_0}} & number of input subformulas and finite set of input atoms & \hyperref[annex:p8.ipl.complete.decision]{\ref*{annex:p8.ipl.complete.decision}} & n/a & \hyperref[annex:p8.ipl.complete.decision]{\ref*{annex:p8.ipl.complete.decision}} & \thepage & Local finite model search parameters; m is a natural number and P0 a set of atoms. \\
\phantomsection\label{entry:symbol-registry-02617}\(r,s\) & \texttt{\detokenize{r,}\allowbreak\detokenize{s}} & root and strict successor in the two-world countermodel & \hyperref[def:p08-ipl-completeness-two-worlds]{\ref*{def:p08-ipl-completeness-two-worlds}} & n/a & \hyperref[def:p08-ipl-completeness-two-worlds]{\ref*{def:p08-ipl-completeness-two-worlds}} & \thepage & Local worlds, not terms or assignments; reflexive loops are included. \\
\phantomsection\label{entry:symbol-registry-02618}\(t\) & \texttt{\detokenize{t}} & propositional atom in the implication-composition example & \hyperref[def:p08-ipl-completeness-persistent-premise]{\ref*{def:p08-ipl-completeness-persistent-premise}} & n/a & \hyperref[def:p08-ipl-completeness-persistent-premise]{\ref*{def:p08-ipl-completeness-persistent-premise}} & \thepage & Explicit local atom role overrides the global term-metavariable default. \\
\phantomsection\label{entry:symbol-registry-02619}\(\vdash_L\) & \texttt{\detokenize{\vdash_L}} & no-premise derivability in the registered LTL calculus & P8 linear temporal completeness & n/a & \hyperref[def:p08-ltl-completeness-language]{\ref*{def:p08-ltl-completeness-language}} & \thepage & All derivations are finite; both temporal rules are theorem-only. \\
\phantomsection\label{entry:symbol-registry-02620}\(\mathsf U_s\) & \texttt{\detokenize{\mathsf U_s}} & strict future until operator & P8 linear temporal completeness & n/a & \hyperref[thm:p08-ltl-strict-bridge]{\ref*{thm:p08-ltl-strict-bridge}} & \thepage & Witness strictly later; the left operand is required only strictly between endpoints. \\
\phantomsection\label{entry:symbol-registry-02621}\(\models_\omega\) & \texttt{\detokenize{\models_\omega}} & local semantic premise consequence over infinite words & P8 linear temporal completeness & n/a & \hyperref[def:p08-ltl-local-consequence]{\ref*{def:p08-ltl-local-consequence}} & \thepage & Every premise holds at the evaluation position; no global-premise convention is introduced. \\
\phantomsection\label{entry:symbol-registry-02622}\(\vdash_L^{\rm fin}\) & \texttt{\detokenize{\vdash_L^{\rm fin}}} & finitely supported local premise derivability & P8 linear temporal completeness & n/a & \hyperref[def:p08-ltl-local-consequence]{\ref*{def:p08-ltl-local-consequence}} & \thepage & Defined through a provable implication from a finite conjunction. \\
\phantomsection\label{entry:symbol-registry-02623}\(C=\operatorname{Sub}(\varphi)\) & \texttt{\detokenize{\operatorname{Sub}}} & finite syntactic subformula set & P8 linear temporal completeness & n/a & \hyperref[def:p08-ltl-consistent-type]{\ref*{def:p08-ltl-consistent-type}} & \thepage & Abbreviations are expanded before forming this set. \\
\phantomsection\label{entry:symbol-registry-02624}\(m,n\) & \texttt{\detokenize{m,}\allowbreak\detokenize{n}} & subformula count and expanded formula length & P8 linear temporal completeness & n/a & \hyperref[def:p08-ltl-consistent-type]{\ref*{def:p08-ltl-consistent-type}} & \thepage & Here n is input length; indices of word positions remain bound variables in the semantic clauses. \\
\phantomsection\label{entry:symbol-registry-02625}\(t\) & \texttt{\detokenize{t}} & Boolean truth record on the subformula set & P8 linear temporal completeness & n/a & \hyperref[def:p08-ltl-consistent-type]{\ref*{def:p08-ltl-consistent-type}} & \thepage & Maps each subformula to zero or one; Hilbert consistency and Boolean coherence are different restrictions. \\
\phantomsection\label{entry:symbol-registry-02626}\(D^1,D^0\) & \texttt{\detokenize{D^1,}\allowbreak\detokenize{D^0}} & positive and negative signed formula literals & P8 linear temporal completeness & n/a & \hyperref[def:p08-ltl-consistent-type]{\ref*{def:p08-ltl-consistent-type}} & \thepage & The superscript is a truth bit; it is not numerical exponentiation. \\
\phantomsection\label{entry:symbol-registry-02627}\(\chi_t\) & \texttt{\detokenize{\chi_t}} & characteristic signed conjunction of a truth record & P8 linear temporal completeness & n/a & \hyperref[def:p08-ltl-consistent-type]{\ref*{def:p08-ltl-consistent-type}} & \thepage & Fixed finite ordering; it records every subformula truth bit. \\
\phantomsection\label{entry:symbol-registry-02628}\(Q_C\) & \texttt{\detokenize{Q_C}} & set of Hilbert-consistent finite truth records & P8 linear temporal completeness & n/a & \hyperref[def:p08-ltl-consistent-type]{\ref*{def:p08-ltl-consistent-type}} & \thepage & Used for the existence proof; its consistency definition is not the decision algorithm. \\
\phantomsection\label{entry:symbol-registry-02629}\(R_C\) & \texttt{\detokenize{R_C}} & consistency-based successor relation between truth records & P8 linear temporal completeness & n/a & \hyperref[def:p08-ltl-consistent-type]{\ref*{def:p08-ltl-consistent-type}} & \thepage & An edge says that the conjunction of the current characteristic formula and next successor characteristic formula is not refutable. \\
\phantomsection\label{entry:symbol-registry-02630}\((\Sigma,Q,I,\Delta,F)\) & \texttt{\detokenize{(\Sigma,}\allowbreak\detokenize{Q,}\allowbreak\detokenize{I,}\allowbreak\detokenize{\Delta,}\allowbreak\detokenize{F)}} & finite alphabet states permitted starts transition function and Buchi accepting set & P8 linear temporal completeness & n/a & \hyperref[def:p08-ltl-buchi-convention]{\ref*{def:p08-ltl-buchi-convention}} & \thepage & The current state describes the letter consumed by its outgoing transition. \\
\phantomsection\label{entry:symbol-registry-02631}\(Q_b\) & \texttt{\detokenize{Q_b}} & all Boolean-coherent subformula truth records & P8 linear temporal completeness & n/a & \hyperref[def:p08-ltl-local-automaton]{\ref*{def:p08-ltl-local-automaton}} & \thepage & Effectively tested without a Hilbert-consistency oracle. \\
\phantomsection\label{entry:symbol-registry-02632}\(P\) & \texttt{\detokenize{P}} & finite set of atoms in the input formula & P8 linear temporal completeness & n/a & \hyperref[def:p08-ltl-local-automaton]{\ref*{def:p08-ltl-local-automaton}} & \thepage & Its powerset is the alphabet of valuation letters. \\
\phantomsection\label{entry:symbol-registry-02633}\(E\) & \texttt{\detokenize{E}} & effective one-step coherence relation on truth records & P8 linear temporal completeness & n/a & \hyperref[def:p08-ltl-local-automaton]{\ref*{def:p08-ltl-local-automaton}} & \thepage & Tests both next equality and non-strict until unfolding. \\
\phantomsection\label{entry:symbol-registry-02634}\(\mathcal U,r\) & \texttt{\detokenize{\mathcal U,}\allowbreak\detokenize{r}} & set and count of until subformulas & P8 linear temporal completeness & n/a & \hyperref[def:p08-ltl-local-automaton]{\ref*{def:p08-ltl-local-automaton}} & \thepage & Every member has a separate finite-witness obligation. \\
\phantomsection\label{entry:symbol-registry-02635}\(D_t\) & \texttt{\detokenize{D_t}} & until obligations true but not currently witnessed in a record & P8 linear temporal completeness & n/a & \hyperref[def:p08-ltl-local-automaton]{\ref*{def:p08-ltl-local-automaton}} & \thepage & These obligations are collected when a pending batch is empty. \\
\phantomsection\label{entry:symbol-registry-02636}\(H_t\) & \texttt{\detokenize{H_t}} & until obligations whose right operands hold currently & P8 linear temporal completeness & n/a & \hyperref[def:p08-ltl-local-automaton]{\ref*{def:p08-ltl-local-automaton}} & \thepage & Their pending copies may be removed in the next batch. \\
\phantomsection\label{entry:symbol-registry-02637}\(O,O'\) & \texttt{\detokenize{O,}\allowbreak\detokenize{O'}} & current and next pending obligation sets & P8 linear temporal completeness & n/a & \hyperref[def:p08-ltl-local-automaton]{\ref*{def:p08-ltl-local-automaton}} & \thepage & An accepting run must have an empty batch infinitely often; a single initial empty batch is insufficient. \\
\phantomsection\label{entry:symbol-registry-02638}\(\mathcal A_\varphi\) & \texttt{\detokenize{\mathcal A_\varphi}} & ordinary Buchi automaton for the formula & P8 linear temporal completeness & n/a & \hyperref[def:p08-ltl-local-automaton]{\ref*{def:p08-ltl-local-automaton}} & \thepage & Product of local truth records and pending obligation batches. \\
\phantomsection\label{entry:symbol-registry-02639}\(F_1,\ldots,F_h\) & \texttt{\detokenize{F_1,}\allowbreak\detokenize{\ldots,}\allowbreak\detokenize{F_h}} & finite list of generalized Buchi accepting sets & P8 linear temporal completeness & n/a & \hyperref[def:p08-ltl-generalized-buchi]{\ref*{def:p08-ltl-generalized-buchi}} & \thepage & Each set must recur; an empty list imposes no acceptance condition on an existing infinite run. \\
\phantomsection\label{entry:symbol-registry-02640}\(F_U\) & \texttt{\detokenize{F_U}} & accepting set for one until subformula & P8 linear temporal completeness & n/a & \hyperref[def:p08-ltl-generalized-buchi]{\ref*{def:p08-ltl-generalized-buchi}} & \thepage & Contains records where that until is false or its right operand is true. \\
\phantomsection\label{entry:symbol-registry-02641}\(\mathcal G_\varphi\) & \texttt{\detokenize{\mathcal G_\varphi}} & generalized Buchi local-type automaton & P8 linear temporal completeness & n/a & \hyperref[def:p08-ltl-generalized-buchi]{\ref*{def:p08-ltl-generalized-buchi}} & \thepage & Same letter-matching edges as the effective local automaton. \\
\phantomsection\label{entry:symbol-registry-02642}\(j\) & \texttt{\detokenize{j}} & next required accepting-set index in the counter conversion & P8 linear temporal completeness & n/a & \hyperref[def:p08-ltl-buchi-counter]{\ref*{def:p08-ltl-buchi-counter}} & \thepage & Its transition update tests the current first-coordinate state; acceptance also requires actual last-set membership. \\
\phantomsection\label{entry:symbol-registry-02643}\(F^*\) & \texttt{\detokenize{F^*}} & ordinary accepting set after generalized Buchi conversion & P8 linear temporal completeness & n/a & \hyperref[def:p08-ltl-buchi-counter]{\ref*{def:p08-ltl-buchi-counter}} & \thepage & A last-counter state alone is insufficient; its first coordinate must lie in the last required set. \\
\phantomsection\label{entry:symbol-registry-02644}\(\to^*,\to^+\) & \texttt{\detokenize{\to^*,}\allowbreak\detokenize{\to^+}} & zero-or-more and positive-length graph reachability & P8 linear temporal completeness & n/a & \hyperref[def:p08-ltl-automaton-graph]{\ref*{def:p08-ltl-automaton-graph}} & \thepage & The positive return prevents false acceptance of a loopless singleton. \\
\phantomsection\label{entry:symbol-registry-02645}\(N\) & \texttt{\detokenize{N}} & number of automaton states for lasso length bounds & P8 linear temporal completeness & n/a & \hyperref[thm:p08-ltl-accepting-cycle]{\ref*{thm:p08-ltl-accepting-cycle}} & \thepage & A shortest reachable prefix has at most N minus one edges and a shortest positive return at most N. \\
\phantomsection\label{entry:symbol-registry-02646}\(\operatorname{len}(u)\) & \texttt{\detokenize{\operatorname{len}}} & finite word length and written formula length & P8 linear temporal completeness & n/a & \hyperref[thm:p08-ltl-accepting-cycle]{\ref*{thm:p08-ltl-accepting-cycle}} & \thepage & Its operand determines the domain; the cardinality bars retain their set-cardinality meaning. \\
\phantomsection\label{entry:symbol-registry-02647}\(K=(W,w_0,R,\ell)\) & \texttt{\detokenize{K=(W,}\allowbreak\detokenize{w_0,}\allowbreak\detokenize{R,}\allowbreak\detokenize{\ell)}} & finite total labelled transition structure with a single initial state & P8 linear temporal completeness & n/a & \hyperref[def:p08-ltl-product-checking]{\ref*{def:p08-ltl-product-checking}} & \thepage & Universal path checking uses the automaton for the syntactic negation. \\
\phantomsection\label{entry:symbol-registry-02648}\(\lVert K\rVert\) & \texttt{\detokenize{\lVert K\rVert}} & complete encoded finite-model input length & P8 linear temporal completeness & n/a & \hyperref[def:p08-ltl-product-checking]{\ref*{def:p08-ltl-product-checking}} & \thepage & Includes state edge and label lists and is at least two; it is not a vector norm. \\
\phantomsection\label{entry:symbol-registry-02649}\(\operatorname{Reach}(u,v,k)\) & \texttt{\detokenize{\operatorname{Reach}}} & bounded graph reachability recursion & P8 linear temporal completeness annex & n/a & \hyperref[annex:p08.ltl.polynomial-space]{\ref*{annex:p08.ltl.polynomial-space}} & \thepage & Tests paths of length at most two to the power k including zero-length paths. \\
\phantomsection\label{entry:symbol-registry-02650}\(b,b'\) & \texttt{\detokenize{b,}\allowbreak\detokenize{b'}} & encoded automaton and product state bit lengths & P8 linear temporal completeness annex & n/a & \hyperref[annex:p08.ltl.polynomial-space]{\ref*{annex:p08.ltl.polynomial-space}} & \thepage & The deterministic recursive decision procedure uses quadratic workspace in the product encoding length. \\
\phantomsection\label{entry:symbol-registry-02651}\(\mathcal D_M\) & \texttt{\detokenize{\mathcal D_M}} & cell alphabet for the fixed polynomial-space machine & P8 linear temporal completeness & n/a & \hyperref[def:p08-ltl-configuration-encoding]{\ref*{def:p08-ltl-configuration-encoding}} & \thepage & Boundary symbols ordinary tape symbols and state-symbol pairs; distinct from the automaton transition function. \\
\phantomsection\label{entry:symbol-registry-02652}\(s\) & \texttt{\detokenize{s}} & polynomially bounded machine-configuration width & P8 linear temporal completeness & n/a & \hyperref[def:p08-ltl-configuration-encoding]{\ref*{def:p08-ltl-configuration-encoding}} & \thepage & Includes two permanent boundary cells and is at least three. \\
\phantomsection\label{entry:symbol-registry-02653}\(\Psi_{M,w}\) & \texttt{\detokenize{\Psi_{M,w}}} & temporal formula encoding a machine computation on an input word & P8 linear temporal completeness & n/a & \hyperref[def:p08-ltl-configuration-encoding]{\ref*{def:p08-ltl-configuration-encoding}} & \thepage & Stationary accepting and rejecting configurations give the required infinite history. \\
\phantomsection\label{entry:symbol-registry-02654}\(C_i,P_d,H,H_{\rm acc}\) & \texttt{\detokenize{C_i,}\allowbreak\detokenize{P_d,}\allowbreak\detokenize{H,}\allowbreak\detokenize{H_{\rm acc}}} & configuration phases cell-symbol atoms head and accepting-head tests & P8 linear temporal completeness & n/a & \hyperref[def:p08-ltl-configuration-encoding]{\ref*{def:p08-ltl-configuration-encoding}} & \thepage & Exactly one phase cell symbol and configuration head are enforced. \\
\phantomsection\label{entry:symbol-registry-02655}\(\operatorname{One}\) & \texttt{\detokenize{\operatorname{One}}} & finite exactly-one Boolean abbreviation & P8 linear temporal completeness & n/a & \hyperref[def:p08-ltl-configuration-encoding]{\ref*{def:p08-ltl-configuration-encoding}} & \thepage & Defined on a nonempty list; all uses expand to the registered Boolean language. \\
\phantomsection\label{entry:symbol-registry-02656}\(f\) & \texttt{\detokenize{f}} & local next-cell update function of the fixed machine & P8 linear temporal completeness & n/a & \hyperref[def:p08-ltl-configuration-encoding]{\ref*{def:p08-ltl-configuration-encoding}} & \thepage & Its finite table and treatment of malformed triples are supplied in the hardness proof. \\
\phantomsection\label{entry:symbol-registry-02657}\(T\) & \texttt{\detokenize{T}} & clock-block translation of temporal formulas & P8 linear temporal completeness & n/a & \hyperref[def:p08-ltl-block-translation]{\ref*{def:p08-ltl-block-translation}} & \thepage & Truth is compared only at phase-zero positions; next advances a full encoded block. \\
\phantomsection\label{entry:symbol-registry-02658}\(w_*,c,b_*\) & \texttt{\detokenize{w_*,}\allowbreak\detokenize{c,}\allowbreak\detokenize{b_*}} & fresh model start phase-zero marker and valuation-bit atom & P8 linear temporal completeness & n/a & \hyperref[def:p08-ltl-block-translation]{\ref*{def:p08-ltl-block-translation}} & \thepage & The initial next step skips the fresh start and preserves the first encoded valuation. \\
\phantomsection\label{entry:symbol-registry-03020}\(P_n\) & \texttt{\detokenize{P_n}} & source atomic sentence indexed by a natural number & P8 Chellas standard systems & object & \hyperref[def:p08-chellas-source-language]{\ref*{def:p08-chellas-source-language}} & \thepage & Primitive atom; source overload with valuation set is explicitly typed in the model definition. \\
\phantomsection\label{entry:symbol-registry-03021}\(P_n\) & \texttt{\detokenize{P_n}} & source valuation set assigned to the nth atom & P8 Chellas standard systems & n/a & \hyperref[def:p08-chellas-standard-model]{\ref*{def:p08-chellas-standard-model}} & \thepage & In membership expressions this is P(n); not a sentence. Source uses the same printed glyph in these two types. \\
\phantomsection\label{entry:symbol-registry-03022}\(P\) & \texttt{\detokenize{P}} & source valuation mapping from natural numbers to world sets & P8 Chellas standard systems & n/a & \hyperref[def:p08-chellas-standard-model]{\ref*{def:p08-chellas-standard-model}} & \thepage & The model component; P\_n is its nth truth set. \\
\phantomsection\label{entry:symbol-registry-03023}\(\operatorname{Sn}(A)\) & \texttt{\detokenize{\operatorname{Sn}}} & source finite subsentence set & P8 Chellas standard systems & n/a & \hyperref[def:p08-chellas-source-language]{\ref*{def:p08-chellas-source-language}} & \thepage & Contains A and every recursively occurring matrix; syntax has both primitive modalities. \\
\phantomsection\label{entry:symbol-registry-03024}\(\Sigma\) & \texttt{\detokenize{\Sigma}} & source modal system as a set of theorem sentences & P8 Chellas standard systems & meta & \hyperref[def:p08-chellas-source-consequence]{\ref*{def:p08-chellas-source-consequence}} & \thepage & Normality is a closure property; Sigma need not be effectively axiomatized in the general canonical lemma. \\
\phantomsection\label{entry:symbol-registry-03025}\(\vdash_\Sigma\) & \texttt{\detokenize{\vdash_\Sigma}} & source theoremhood and finitely supported premise deducibility & P8 Chellas standard systems & relation & \hyperref[def:p08-chellas-source-consequence]{\ref*{def:p08-chellas-source-consequence}} & \thepage & An absent premise set denotes membership in Sigma; premises use a finite theorem implication. \\
\phantomsection\label{entry:symbol-registry-03026}\(\Gamma\) & \texttt{\detokenize{\Gamma}} & arbitrary source local premise set & P8 Chellas source consequence & meta & \hyperref[def:p08-chellas-source-consequence]{\ref*{def:p08-chellas-source-consequence}} & \thepage & No global-premise necessitation; filtration uses its separately bound sentence-set role. \\
\phantomsection\label{entry:symbol-registry-03027}\(\Gamma\) & \texttt{\detokenize{\Gamma}} & source subsentence-closed filtration sentence set & P8 Chellas source filtration & meta & \hyperref[def:p08-chellas-source-filtration]{\ref*{def:p08-chellas-source-filtration}} & \thepage & May be syntactically infinite when modally closed but must be logically finite to yield a finite quotient. \\
\phantomsection\label{entry:symbol-registry-03028}\(\alpha,\beta,\gamma,\delta,\epsilon\) & \texttt{\detokenize{\alpha,}\allowbreak\detokenize{\beta,}\allowbreak\detokenize{\gamma,}\allowbreak\detokenize{\delta,}\allowbreak\detokenize{\epsilon}} & source model worlds and canonical-world maximal theories & P8 Chellas standard systems & n/a & \hyperref[def:p08-chellas-standard-model]{\ref*{def:p08-chellas-standard-model}} & \thepage & World type is fixed by the model; canonical worlds are particular sets of sentences. \\
\phantomsection\label{entry:symbol-registry-03029}\(\models_\alpha^{\mathcal M}\) & \texttt{\detokenize{\models_\alpha^{\mathcal M}}} & source truth of a sentence at a world of a standard model & P8 Chellas standard systems & relation & \hyperref[def:p08-chellas-standard-model]{\ref*{def:p08-chellas-standard-model}} & \thepage & Classical truth clauses and both primitive modality clauses. \\
\phantomsection\label{entry:symbol-registry-03030}\(\models_C\) & \texttt{\detokenize{\models_C}} & source validity throughout a class of models & P8 Chellas standard systems & relation & \hyperref[def:p08-chellas-standard-model]{\ref*{def:p08-chellas-standard-model}} & \thepage & Truth at every world in every model; determination concerns theoremhood. \\
\phantomsection\label{entry:symbol-registry-03031}\(\models_C^\ell\) & \texttt{\detokenize{\models_C^\ell}} & source local premise consequence on a model class & P8 Chellas standard systems & relation & \hyperref[def:p08-chellas-standard-model]{\ref*{def:p08-chellas-standard-model}} & \thepage & One evaluation world satisfies every premise; not global premise truth. \\
\phantomsection\label{entry:symbol-registry-03032}\(\|A\|^{\mathcal M}\) & \texttt{\detokenize{\|A\|^{\mathcal M}}} & source truth set of a sentence in a model & P8 Chellas standard systems & n/a & \hyperref[def:p08-chellas-standard-model]{\ref*{def:p08-chellas-standard-model}} & \thepage & A subset of the model world set. \\
\phantomsection\label{entry:symbol-registry-03033}\(W_\Sigma\) & \texttt{\detokenize{W_\Sigma}} & all maximal source Sigma-consistent theories & P8 Chellas standard systems & n/a & \hyperref[def:p08-chellas-proper-canonical]{\ref*{def:p08-chellas-proper-canonical}} & \thepage & Nonempty when Sigma is consistent; the source general canonical definition and proper construction are distinguished. \\
\phantomsection\label{entry:symbol-registry-03034}\(R_\Sigma\) & \texttt{\detokenize{R_\Sigma}} & proper source canonical accessibility relation & P8 Chellas standard systems & relation & \hyperref[def:p08-chellas-proper-canonical]{\ref*{def:p08-chellas-proper-canonical}} & \thepage & Boxed-content containment; constructed relation rather than an assumed truth lemma. \\
\phantomsection\label{entry:symbol-registry-03035}\(R^j\) & \texttt{\detokenize{R^j}} & exact j-step source accessibility path relation & P8 Chellas standard systems & relation & \hyperref[def:p08-chellas-g-family]{\ref*{def:p08-chellas-g-family}} & \thepage & Zero is identity; exact path length not reflexive-transitive reachability. \\
\phantomsection\label{entry:symbol-registry-03036}\(G^{k,l,m,n}\) & \texttt{\detokenize{G^{k,l,m,n}}} & source four-index modal schema & P8 Chellas standard systems & object & \hyperref[def:p08-chellas-g-family]{\ref*{def:p08-chellas-g-family}} & \thepage & Diamond-k Box-l implies Box-m Diamond-n; preserve index order and zero cases. \\
\phantomsection\label{entry:symbol-registry-03037}\(W^\alpha,R^\alpha,P^\alpha\) & \texttt{\detokenize{W^\alpha,}\allowbreak\detokenize{R^\alpha,}\allowbreak\detokenize{P^\alpha}} & source generated restriction components & P8 Chellas standard systems & n/a & \hyperref[def:p08-chellas-generated-model]{\ref*{def:p08-chellas-generated-model}} & \thepage & Reachable world subset and restricted relation and valuation. \\
\phantomsection\label{entry:symbol-registry-03038}\(\mathcal M^\alpha\) & \texttt{\detokenize{\mathcal M^\alpha}} & source model generated by a world & P8 Chellas standard systems & n/a & \hyperref[def:p08-chellas-generated-model]{\ref*{def:p08-chellas-generated-model}} & \thepage & All finite descendants including the root; no outgoing edge is discarded. \\
\phantomsection\label{entry:symbol-registry-03039}\(\sim_\Gamma\) & \texttt{\detokenize{\sim_\Gamma}} & agreement on every source filtration sentence & P8 Chellas standard systems & relation & \hyperref[def:p08-chellas-source-filtration]{\ref*{def:p08-chellas-source-filtration}} & \thepage & Equivalence relation on worlds in one fixed original model. \\
\phantomsection\label{entry:symbol-registry-03040}\([\alpha],[X]\) & \texttt{\detokenize{[\alpha],}\allowbreak\detokenize{[X]}} & source quotient world and quotient image of a world set & P8 Chellas standard systems & n/a & \hyperref[def:p08-chellas-source-filtration]{\ref*{def:p08-chellas-source-filtration}} & \thepage & X image consists of the equivalence classes with a member in X. \\
\phantomsection\label{entry:symbol-registry-03041}\(W^*,R^*,P^*,\mathcal M^*\) & \texttt{\detokenize{W^*,}\allowbreak\detokenize{R^*,}\allowbreak\detokenize{P^*,}\allowbreak\detokenize{\mathcal M^*}} & source filtration components and quotient model & P8 Chellas standard systems & n/a & \hyperref[def:p08-chellas-source-filtration]{\ref*{def:p08-chellas-source-filtration}} & \thepage & Source star marks a filtration object; it is not relation reflexive-transitive closure. \\
\phantomsection\label{entry:symbol-registry-03042}\(c_1,c_2,c_3,c_4\) & \texttt{\detokenize{c_1,}\allowbreak\detokenize{c_2,}\allowbreak\detokenize{c_3,}\allowbreak\detokenize{c_4}} & source paired filtration constraints on two worlds & P8 Chellas standard systems & meta & \hyperref[def:p08-chellas-filtration-conditions]{\ref*{def:p08-chellas-filtration-conditions}} & \thepage & c2 reverses c1; c4 reverses c3; corrected symmetric transitive construction retains all four. \\
\phantomsection\label{entry:symbol-registry-03043}\(U_j(\alpha)\) & \texttt{\detokenize{U_j(\alpha)}} & contents of all j-times-boxed sentences in a canonical world & P8 Chellas standard systems & n/a & \hyperref[annex:p8.chellas.g-determination]{\ref*{annex:p8.chellas.g-determination}} & \thepage & Proof-local j-indexed set; j=0 is the world itself. \\
\phantomsection\label{entry:symbol-registry-03044}\(B_A(\alpha),D_A(\alpha)\) & \texttt{\detokenize{B_A(\alpha),}\allowbreak\detokenize{D_A(\alpha)}} & proof-local truth predicates for boxed and diamond sentences & P8 Chellas standard systems & meta & \hyperref[annex:p8.chellas.paired-filtrations]{\ref*{annex:p8.chellas.paired-filtrations}} & \thepage & Used only for modalities whose sentences occur in the filtration set. \\
\phantomsection\label{entry:symbol-registry-03045}\(e\) & \texttt{\detokenize{e}} & recursive source-to-sheet formula translation & P8 Chellas standard systems & meta & \hyperref[annex:p8.chellas.normal-interface]{\ref*{annex:p8.chellas.normal-interface}} & \thepage & Eliminates primitive diamond through the source duality theorem and maps uppercase atoms to sheet atoms. \\
\phantomsection\label{entry:symbol-registry-03202}\(t(A)\) & \texttt{\detokenize{t(A)}} & translation eliminating source primitive possibility into the host defined dual & P8 Chellas neighborhood & object & \hyperref[def:p08-cn-language]{\ref*{def:p08-cn-language}} & \thepage & The translation fixes atoms constants and boxes and commutes with Boolean connectives. \\
\phantomsection\label{entry:symbol-registry-03203}\(\Gamma\vdash_L A\) & \texttt{\detokenize{\Gamma\vdash_L A}} & local finite-sequence consequence with only theorem-level modal replacement & P8 Chellas neighborhood & relation & \hyperref[def:p08-cn-language]{\ref*{def:p08-cn-language}} & \thepage & Premises are true at one world; every RE equivalence needs a separate proof without those premises. \\
\phantomsection\label{entry:symbol-registry-03204}\(A^*,\mu^*\) & \texttt{\detokenize{A^*,}\allowbreak\detokenize{\mu^*}} & source formula dual and dual modality word & P8 Chellas neighborhood & object & \hyperref[def:p08-cn-dual]{\ref*{def:p08-cn-dual}} & \thepage & The formula dual negates atoms and swaps Boolean duals; the word dual swaps box and diamond while retaining negation. \\
\phantomsection\label{entry:symbol-registry-03205}\(\|A\|_{\mathcal M}\) & \texttt{\detokenize{\|A\|_{\mathcal M}}} & formula truth set in a neighborhood model & P8 Chellas neighborhood & object & \hyperref[def:p08-cn-models]{\ref*{def:p08-cn-models}} & \thepage & Source truth sets are propositions in the powerset of worlds. \\
\phantomsection\label{entry:symbol-registry-03206}\(f(X),d(X)\) & \texttt{\detokenize{f(X),}\allowbreak\detokenize{d(X)}} & powerset necessity operation and its complement dual & P8 Chellas neighborhood & object & \hyperref[def:p08-cn-models]{\ref*{def:p08-cn-models}} & \thepage & The operation determines neighborhood membership and conversely; iteration uses composition. \\
\phantomsection\label{entry:symbol-registry-03207}\(\mathcal B^+,\mathcal B^-,\mathcal B^\#\) & \texttt{\detokenize{\mathcal B^+,}\allowbreak\detokenize{\mathcal B^-,}\allowbreak\detokenize{\mathcal B^\#}} & supplementation nonempty finite intersection closure and quasi-filtering of a family & P8 Chellas neighborhood & object & \hyperref[def:p08-cn-closures]{\ref*{def:p08-cn-closures}} & \thepage & Intersection arity is positive; supplementation and this intersection closure commute. \\
\phantomsection\label{entry:symbol-registry-03208}\((M,C,N)\) & \texttt{\detokenize{(M,}\allowbreak\detokenize{C,}\allowbreak\detokenize{N)}} & Chellas axiom-schema identifiers for three independent neighborhood properties & P8 Chellas neighborhood & object & \hyperref[def:p08-cn-systems]{\ref*{def:p08-cn-systems}} & \thepage & These names identify schemas; the source also uses M as a system name with its separately declared type. \\
\phantomsection\label{entry:symbol-registry-03209}\((M,R,K)\) & \texttt{\detokenize{(M,}\allowbreak\detokenize{R,}\allowbreak\detokenize{K)}} & source smallest monotonic regular and normal systems & P8 Chellas neighborhood & object & \hyperref[sec:p08-n-b-006-03]{\ref*{sec:p08-n-b-006-03}} & \thepage & Proved names are M equals EM; R equals EMC; K equals EMCN. System names are distinct typed roles from same-letter schemas. \\
\phantomsection\label{entry:symbol-registry-03210}\(L_S,S\) & \texttt{\detokenize{L_S,}\allowbreak\detokenize{S}} & classical extension selected by a subset of the M C N schema names & P8 Chellas neighborhood & object & \hyperref[def:p08-cn-systems]{\ref*{def:p08-cn-systems}} & \thepage & The selected family has eight systems; it is not the full lattice of modal logics. \\
\phantomsection\label{entry:symbol-registry-03211}\((RM,RR,RK)\) & \texttt{\detokenize{(RM,}\allowbreak\detokenize{RR,}\allowbreak\detokenize{RK)}} & source implication binary regular and finite normal theorem rules & P8 Chellas neighborhood & object & \hyperref[def:p08-cn-systems]{\ref*{def:p08-cn-systems}} & \thepage & The normal rule includes zero antecedents; all rules apply to theorems. \\
\phantomsection\label{entry:symbol-registry-03212}\(W_L\) & \texttt{\detokenize{W_L}} & world set of maximal consistent theories in a classical system & P8 Chellas neighborhood & object & \hyperref[def:p08-cn-proof-sets]{\ref*{def:p08-cn-proof-sets}} & \thepage & Existence uses the countable enumeration and local deduction argument; base systems are proved consistent. \\
\phantomsection\label{entry:symbol-registry-03213}\(|A|_L\) & \texttt{\detokenize{|A|_L}} & Chellas published proof set consisting of maximal worlds containing A & P8 Chellas neighborhood & object & \hyperref[def:p08-cn-proof-sets]{\ref*{def:p08-cn-proof-sets}} & \thepage & This source-local bar notation is a proposition set not cardinality; it differs from the logical equivalence class. \\
\phantomsection\label{entry:symbol-registry-03214}\(D_L\) & \texttt{\detokenize{D_L}} & family of all formula proof sets in the canonical world set & P8 Chellas neighborhood & object & \hyperref[def:p08-cn-proof-sets]{\ref*{def:p08-cn-proof-sets}} & \thepage & Non-proof subsets are those outside this family; changing their membership may affect structural frame conditions. \\
\phantomsection\label{entry:symbol-registry-03215}\(N^{\min}(u),N^{\max}(u)\) & \texttt{\detokenize{N^{\min}(u),}\allowbreak\detokenize{N^{\max}(u)}} & smallest and largest canonical neighborhood choices & P8 Chellas neighborhood & object & \hyperref[def:p08-cn-proof-sets]{\ref*{def:p08-cn-proof-sets}} & \thepage & The smallest includes forced proof sets only; the largest additionally includes all non-proof subsets. \\
\phantomsection\label{entry:symbol-registry-03216}\(N_R(w)\) & \texttt{\detokenize{N_R(w)}} & principal neighborhood family of an accessibility relation & P8 Chellas neighborhood & object & \hyperref[def:p08-cn-augmentation]{\ref*{def:p08-cn-augmentation}} & \thepage & Contains exactly the supersets of the successor set; an empty successor set gives all subsets. \\
\phantomsection\label{entry:symbol-registry-03217}\(\bigcap N(w)\) & \texttt{\detokenize{\bigcap N(w)}} & core of a world's neighborhood family & P8 Chellas neighborhood & object & \hyperref[def:p08-cn-augmentation]{\ref*{def:p08-cn-augmentation}} & \thepage & Empty-family intersection is the whole world domain; an augmented family contains its core and every superset. \\
\phantomsection\label{entry:symbol-registry-03218}\(F,W_F,\pi\) & \texttt{\detokenize{F,}\allowbreak\detokenize{W_F,}\allowbreak\detokenize{\pi}} & finite subformula observation set quotient worlds and quotient map & P8 Chellas neighborhood & object & \hyperref[def:p08-cn-filtration]{\ref*{def:p08-cn-filtration}} & \thepage & A required unit is retained by adding the boxed truth observation; all observed proposition sets are saturated. \\
\phantomsection\label{entry:symbol-registry-03219}\((W,N,P)\) & \texttt{\detokenize{(W,}\allowbreak\detokenize{N,}\allowbreak\detokenize{P)}} & the 1980 source minimal model with a sequence of atom valuations & P8 Chellas neighborhood & object & \hyperref[def:p08-cn-models]{\ref*{def:p08-cn-models}} & \thepage & Source P at index n is the host valuation of atom p at index n; the source permits empty world sets. \\
\phantomsection\label{entry:symbol-registry-03220}\(\ast\) & \texttt{\detokenize{\ast}} & arbitrary unary modal operation on a Boolean algebra & P8 Chellas neighborhood & object & \hyperref[def:p08-cn-algebra]{\ref*{def:p08-cn-algebra}} & \thepage & This semantic unary operation is distinct from the formula and modality-word dual transformations. \\
\phantomsection\label{entry:symbol-registry-03221}\([A]_L\) & \texttt{\detokenize{[A]_L}} & logical equivalence class in the source Lindenbaum modal algebra construction & P8 Chellas neighborhood & object & \hyperref[def:p08-cn-algebra]{\ref*{def:p08-cn-algebra}} & \thepage & Equivalent exactly when the formula biconditional is a theorem; RE makes its boxed class well defined. \\
\phantomsection\label{entry:symbol-registry-03222}\(B_F,r,\ast_F\) & \texttt{\detokenize{B_F,}\allowbreak\detokenize{r,}\allowbreak\detokenize{\ast_F}} & finite Boolean subalgebra greatest-lower-bound projection and projected modal operation & P8 Chellas neighborhood & object & \hyperref[thm:p08-cn-algebra-finite]{\ref*{thm:p08-cn-algebra-finite}} & \thepage & Only represented modal outputs need exact agreement; projection preserves each M C N inequality. \\
\phantomsection\label{entry:symbol-registry-03223}\(f^n,d^n\) & \texttt{\detokenize{f^n,}\allowbreak\detokenize{d^n}} & finite iterates of the necessity and possibility operations & P8 Chellas neighborhood & object & \hyperref[thm:p08-cn-other-conditions]{\ref*{thm:p08-cn-other-conditions}} & \thepage & The zeroth iterate is identity; the diamond iterate is complement then the box iterate then complement. \\
\phantomsection\label{entry:symbol-registry-03224}\((d,t,b,iv,v)\) & \texttt{\detokenize{(d,}\allowbreak\detokenize{t,}\allowbreak\detokenize{b,}\allowbreak\detokenize{iv,}\allowbreak\detokenize{v)}} & source neighborhood-property names corresponding to schemas D T B 4 5 & P8 Chellas neighborhood & object & \hyperref[def:p08-cn-other-conditions]{\ref*{def:p08-cn-other-conditions}} & \thepage & These are conditions on proposition families; no undeclared relational accessibility properties are substituted. \\
\phantomsection\label{entry:symbol-registry-03225}\(u^\circ,u_n\) & \texttt{\detokenize{u^\circ,}\allowbreak\detokenize{u_n}} & maximal source theories in the explicit smallest-E5 counterexample & \hyperref[annex:p8.cn.nonproof]{\ref*{annex:p8.cn.nonproof}} & object & \hyperref[annex:p8.cn.nonproof]{\ref*{annex:p8.cn.nonproof}} & \thepage & The first lacks boxed truth; the indexed identity theories are pairwise distinct and define the diagonal non-proof set. \\
\phantomsection\label{entry:symbol-registry-03226}\(\mathcal C\) & \texttt{\detokenize{\mathcal C}} & finite complement-closed proposition family used in the ET5 non-normality witness & \hyperref[annex:p8.cn.t5-principles]{\ref*{annex:p8.cn.t5-principles}} & object & \hyperref[annex:p8.cn.t5-principles]{\ref*{annex:p8.cn.t5-principles}} & \thepage & Its induced operation fixes these six sets and sends every other set to empty. \\
\phantomsection\label{entry:symbol-registry-03361}\(\mathrm{PROP}\) & \texttt{\detokenize{\mathrm{PROP}}} & proposition-symbol signature & \hyperref[sec:p08-src-025-501]{\ref*{sec:p08-src-025-501}} & n/a & \hyperref[def:p08-bvb501-language]{\ref*{def:p08-bvb501-language}} & \thepage & Arbitrary set; effectiveness and cardinal hypotheses are stated where needed. \\
\phantomsection\label{entry:symbol-registry-03362}\(\mathrm{MOD}\) & \texttt{\detokenize{\mathrm{MOD}}} & modality-name signature & \hyperref[sec:p08-src-025-501]{\ref*{sec:p08-src-025-501}} & n/a & \hyperref[def:p08-bvb501-language]{\ref*{def:p08-bvb501-language}} & \thepage & Each name selects a binary relation. \\
\phantomsection\label{entry:symbol-registry-03363}\(R_m\) & \texttt{\detokenize{R_m}} & accessibility relation for one named modality & \hyperref[sec:p08-src-025-501]{\ref*{sec:p08-src-025-501}} & n/a & \hyperref[def:p08-bvb501-language]{\ref*{def:p08-bvb501-language}} & \thepage & No implicit transitivity or functionality. \\
\phantomsection\label{entry:symbol-registry-03364}\(\operatorname{md}(\varphi)\) & \texttt{\detokenize{\operatorname{md}(\varphi)}} & modal nesting depth & \hyperref[sec:p08-src-025-501]{\ref*{sec:p08-src-025-501}} & n/a & \hyperref[def:p08-bvb501-language]{\ref*{def:p08-bvb501-language}} & \thepage & Maximum on Boolean combinations. \\
\phantomsection\label{entry:symbol-registry-03365}\(\langle m\rangle,[m]\) & \texttt{\detokenize{\langle m\rangle,}\allowbreak\detokenize{[m]}} & indexed diamond and box & \hyperref[sec:p08-src-025-501]{\ref*{sec:p08-src-025-501}} & n/a & \hyperref[def:p08-bvb501-language]{\ref*{def:p08-bvb501-language}} & \thepage & Relational semantic duals. \\
\phantomsection\label{entry:symbol-registry-03366}\(\operatorname{ST}_x\) & \texttt{\detokenize{\operatorname{ST}_x}} & first-order standard translation at a displayed state variable & \hyperref[sec:p08-src-025-501]{\ref*{sec:p08-src-025-501}} & n/a & \hyperref[def:p08-bvb501-standard-translation]{\ref*{def:p08-bvb501-standard-translation}} & \thepage & The output has at most the displayed free variable. \\
\phantomsection\label{entry:symbol-registry-03367}\(P_p\) & \texttt{\detokenize{P_p}} & unary predicate interpreting a proposition letter & \hyperref[sec:p08-src-025-501]{\ref*{sec:p08-src-025-501}} & n/a & \hyperref[def:p08-bvb501-standard-translation]{\ref*{def:p08-bvb501-standard-translation}} & \thepage & Interpreted by the valuation extension. \\
\phantomsection\label{entry:symbol-registry-03368}\(E\) & \texttt{\detokenize{E}} & bisimulation relation between two modal domains & \hyperref[sec:p08-src-025-501]{\ref*{sec:p08-src-025-501}} & n/a & \hyperref[def:p08-bvb501-bisimulation]{\ref*{def:p08-bvb501-bisimulation}} & \thepage & Existence is pointed or unpointed as explicitly stated; no totality requirement. \\
\phantomsection\label{entry:symbol-registry-03369}\(\sim\) & \texttt{\detokenize{\sim}} & largest self-bisimulation equivalence & \hyperref[sec:p08-src-025-501]{\ref*{sec:p08-src-025-501}} & n/a & \hyperref[def:p08-bvb501-contraction]{\ref*{def:p08-bvb501-contraction}} & \thepage & Unions and compositions justify the equivalence property. \\
\phantomsection\label{entry:symbol-registry-03370}\([w]\) & \texttt{\detokenize{[w]}} & bisimulation quotient class & \hyperref[sec:p08-src-025-501]{\ref*{sec:p08-src-025-501}} & n/a & \hyperref[def:p08-bvb501-contraction]{\ref*{def:p08-bvb501-contraction}} & \thepage & Quotient point distinct from world cardinality. \\
\phantomsection\label{entry:symbol-registry-03371}\(\lvert w\rvert\) & \texttt{\detokenize{\lvert w\rvert}} & original source notation for a bisimulation class & \hyperref[sec:p08-src-025-501]{\ref*{sec:p08-src-025-501}} & n/a & \hyperref[def:p08-bvb501-contraction]{\ref*{def:p08-bvb501-contraction}} & \thepage & Explicitly source-local class notation; not cardinality in this declaration. \\
\phantomsection\label{entry:symbol-registry-03372}\(I\) & \texttt{\detokenize{I}} & potential-isomorphism family of finite maps & \hyperref[sec:p08-bvb501-first-order-comparison]{\ref*{sec:p08-bvb501-first-order-comparison}} & n/a & \hyperref[def:p08-bvb501-potential-isomorphism]{\ref*{def:p08-bvb501-potential-isomorphism}} & \thepage & Closed under the stated forth and back extensions. \\
\phantomsection\label{entry:symbol-registry-03373}\(L_{\infty\omega},L_{\omega_1\omega}\) & \texttt{\detokenize{L_{\infty\omega},}\allowbreak\detokenize{L_{\omega_1\omega}}} & first-order infinitary languages & \hyperref[sec:p08-bvb501-first-order-comparison]{\ref*{sec:p08-bvb501-first-order-comparison}} & n/a & \hyperref[def:p08-bvb501-potential-isomorphism]{\ref*{def:p08-bvb501-potential-isomorphism}} & \thepage & Set-sized versus countable connectives; finite free-variable and quantifier conditions. \\
\phantomsection\label{entry:symbol-registry-03374}\(L_{\infty\omega}\) & \texttt{\detokenize{L_{\infty\omega}}} & source notation for infinitary modal language & \hyperref[sec:p08-bvb501-comparison-games]{\ref*{sec:p08-bvb501-comparison-games}} & n/a & \hyperref[def:p08-bvb501-infinitary-modal-language]{\ref*{def:p08-bvb501-infinitary-modal-language}} & \thepage & Declared modal role has no world variables; distinct from first-order infinitary syntax. \\
\phantomsection\label{entry:symbol-registry-03375}\(\delta_M\) & \texttt{\detokenize{\delta_M}} & characteristic sentence for potential isomorphism & \hyperref[sec:p08-src-025-501]{\ref*{sec:p08-src-025-501}} & n/a & \hyperref[thm:p08-bvb501-potential-isomorphism]{\ref*{thm:p08-bvb501-potential-isomorphism}} & \thepage & Characterizes isomorphism only within countable structures when model and signature are countable. \\
\phantomsection\label{entry:symbol-registry-03376}\(g\) & \texttt{\detokenize{g}} & assignment as a state of an assignment-reset modal model & \hyperref[sec:p08-src-025-501]{\ref*{sec:p08-src-025-501}} & n/a & \hyperref[def:p08-bvb501-assignment-model]{\ref*{def:p08-bvb501-assignment-model}} & \thepage & Each variable-labelled transition resets only one coordinate. \\
\phantomsection\label{entry:symbol-registry-03377}\([\![\sigma]\!]\) & \texttt{\detokenize{[\![\sigma]\!]}} & modal formula truth set & \hyperref[sec:p08-src-025-501]{\ref*{sec:p08-src-025-501}} & n/a & \hyperref[def:p08-bvb501-labelling]{\ref*{def:p08-bvb501-labelling}} & \thepage & Computed on finite domains in the algorithm; the same semantic notation applies to arbitrary models in later proofs. \\
\phantomsection\label{entry:symbol-registry-03378}\(s(\varphi)\) & \texttt{\detokenize{s(\varphi)}} & number of distinct subformulas & \hyperref[sec:p08-bvb501-computation]{\ref*{sec:p08-bvb501-computation}} & n/a & \hyperref[thm:p08-bvb501-model-checking]{\ref*{thm:p08-bvb501-model-checking}} & \thepage & Counts the original primitive box and diamond grammar without a duality expansion. \\
\phantomsection\label{entry:symbol-registry-03379}\(U_{\exists}\) & \texttt{\detokenize{U_{\exists}}} & existential finite-path Until & \hyperref[sec:p08-bvb501-model-checking]{\ref*{sec:p08-bvb501-model-checking}} & n/a & \hyperref[def:p08-bvb501-until-check]{\ref*{def:p08-bvb501-until-check}} & \thepage & Distinct from the universally quantified source-target interval clause on branching graphs. \\
\phantomsection\label{entry:symbol-registry-03380}\(\sim_\Sigma\) & \texttt{\detokenize{\sim_\Sigma}} & agreement on a subformula-closed set & \hyperref[sec:p08-bvb501-filtration]{\ref*{sec:p08-bvb501-filtration}} & n/a & \hyperref[def:p08-bvb501-paired-filtration]{\ref*{def:p08-bvb501-paired-filtration}} & \thepage & Not the earlier largest self-bisimulation relation. \\
\phantomsection\label{entry:symbol-registry-03381}\([w]_\Sigma\) & \texttt{\detokenize{[w]_\Sigma}} & filtration class & \hyperref[sec:p08-bvb501-filtration]{\ref*{sec:p08-bvb501-filtration}} & n/a & \hyperref[def:p08-bvb501-paired-filtration]{\ref*{def:p08-bvb501-paired-filtration}} & \thepage & Truth-pattern quotient rather than a bisimulation quotient. \\
\phantomsection\label{entry:symbol-registry-03382}\(\lvert w\rvert_\Sigma\) & \texttt{\detokenize{\lvert w\rvert_\Sigma}} & original source filtration-class notation & \hyperref[sec:p08-bvb501-filtration]{\ref*{sec:p08-bvb501-filtration}} & n/a & \hyperref[def:p08-bvb501-paired-filtration]{\ref*{def:p08-bvb501-paired-filtration}} & \thepage & Explicit local class role rather than cardinality. \\
\phantomsection\label{entry:symbol-registry-03383}\(R_m^s,R_m^l,R_m^f\) & \texttt{\detokenize{R_m^s,}\allowbreak\detokenize{R_m^l,}\allowbreak\detokenize{R_m^f}} & minimum maximum and intervening filtration relations & \hyperref[sec:p08-bvb501-filtration]{\ref*{sec:p08-bvb501-filtration}} & n/a & \hyperref[def:p08-bvb501-paired-filtration]{\ref*{def:p08-bvb501-paired-filtration}} & \thepage & Maximum includes both primitive box and primitive diamond constraints. \\
\phantomsection\label{entry:symbol-registry-03384}\(l(X)\) & \texttt{\detokenize{l(X)}} & universal-successor subset operator & \hyperref[sec:p08-bvb501-frame-definability]{\ref*{sec:p08-bvb501-frame-definability}} & n/a & \hyperref[def:p08-bvb501-ultrafilter-extension]{\ref*{def:p08-bvb501-ultrafilter-extension}} & \thepage & Preserves finite intersections and the whole domain. \\
\phantomsection\label{entry:symbol-registry-03385}\(\operatorname{uf}(W)\) & \texttt{\detokenize{\operatorname{uf}(W)}} & domain of proper ultrafilters & \hyperref[sec:p08-bvb501-frame-definability]{\ref*{sec:p08-bvb501-frame-definability}} & n/a & \hyperref[def:p08-bvb501-ultrafilter-extension]{\ref*{def:p08-bvb501-ultrafilter-extension}} & \thepage & Each member decides a subset against its complement. \\
\phantomsection\label{entry:symbol-registry-03386}\(\operatorname{ue}(F)\) & \texttt{\detokenize{\operatorname{ue}(F)}} & ultrafilter extension of a frame & \hyperref[sec:p08-bvb501-frame-definability]{\ref*{sec:p08-bvb501-frame-definability}} & n/a & \hyperref[def:p08-bvb501-ultrafilter-extension]{\ref*{def:p08-bvb501-ultrafilter-extension}} & \thepage & Validity reflection has direction from extension to original. \\
\phantomsection\label{entry:symbol-registry-03387}\(R^{\mathrm{ue}}\) & \texttt{\detokenize{R^{\mathrm{ue}}}} & accessibility between ultrafilters & \hyperref[sec:p08-bvb501-frame-definability]{\ref*{sec:p08-bvb501-frame-definability}} & n/a & \hyperref[def:p08-bvb501-ultrafilter-extension]{\ref*{def:p08-bvb501-ultrafilter-extension}} & \thepage & Defined by the l-preimage filter condition. \\
\phantomsection\label{entry:symbol-registry-03388}\(R_x\) & \texttt{\detokenize{R_x}} & edges whose departure is a successor of the fixed source & \hyperref[sec:p08-bvb501-fixed-correspondence]{\ref*{sec:p08-bvb501-fixed-correspondence}} & n/a & \hyperref[thm:p08-bvb501-minimal-fixed-point]{\ref*{thm:p08-bvb501-minimal-fixed-point}} & \thepage & Restricted return relation; reduces to ordinary return under transitivity. \\
\phantomsection\label{entry:symbol-registry-03389}\(E,A\) & \texttt{\detokenize{E,}\allowbreak\detokenize{A}} & universal diamond and box & \hyperref[sec:p08-bvb501-universal]{\ref*{sec:p08-bvb501-universal}} & n/a & \hyperref[def:p08-bvb501-universal]{\ref*{def:p08-bvb501-universal}} & \thepage & Fixed interpretation by the universal relation; E is a modality here rather than the earlier bisimulation relation. \\
\phantomsection\label{entry:symbol-registry-03390}\(D\) & \texttt{\detokenize{D}} & difference diamond & \hyperref[sec:p08-bvb501-universal]{\ref*{sec:p08-bvb501-universal}} & n/a & \hyperref[def:p08-bvb501-universal]{\ref*{def:p08-bvb501-universal}} & \thepage & Fixed interpretation by inequality; not freely chosen accessibility. \\
\phantomsection\label{entry:symbol-registry-03391}\(R^*,R^+\) & \texttt{\detokenize{R^*,}\allowbreak\detokenize{R^+}} & reflexive-transitive and transitive accessibility closures & \hyperref[sec:p08-src-025-501]{\ref*{sec:p08-src-025-501}} & n/a & \hyperref[def:p08-bvb501-until-check]{\ref*{def:p08-bvb501-until-check}} & \thepage & Finite-path reachability; not infinitary closure. \\
\phantomsection\label{entry:symbol-registry-03392}\(i,j,k\) & \texttt{\detokenize{i,}\allowbreak\detokenize{j,}\allowbreak\detokenize{k}} & hybrid nominals with singleton extensions & \hyperref[sec:p08-bvb501-hybrid]{\ref*{sec:p08-bvb501-hybrid}} & n/a & \hyperref[def:p08-bvb501-hybrid]{\ref*{def:p08-bvb501-hybrid}} & \thepage & Different nominals may name the same point. \\
\phantomsection\label{entry:symbol-registry-03393}\(d_i\) & \texttt{\detokenize{d_i}} & denotation of a hybrid nominal & \hyperref[sec:p08-bvb501-hybrid]{\ref*{sec:p08-bvb501-hybrid}} & n/a & \hyperref[def:p08-bvb501-hybrid]{\ref*{def:p08-bvb501-hybrid}} & \thepage & Related by named closure even when unreachable from the selected point. \\
\phantomsection\label{entry:symbol-registry-03394}\(@_i\) & \texttt{\detokenize{@_i}} & hybrid satisfaction operator & \hyperref[sec:p08-bvb501-hybrid]{\ref*{sec:p08-bvb501-hybrid}} & n/a & \hyperref[def:p08-bvb501-hybrid]{\ref*{def:p08-bvb501-hybrid}} & \thepage & Truth is evaluated at the nominal denotation. \\
\phantomsection\label{entry:symbol-registry-03395}\(U,S\) & \texttt{\detokenize{U,}\allowbreak\detokenize{S}} & strict relational interval Until and Since & \hyperref[sec:p08-bvb501-temporal]{\ref*{sec:p08-bvb501-temporal}} & n/a & \hyperref[def:p08-bvb501-until-since]{\ref*{def:p08-bvb501-until-since}} & \thepage & Distinct from finite-path Until used by the queue algorithm. \\
\phantomsection\label{entry:symbol-registry-03396}\(C_w\) & \texttt{\detokenize{C_w}} & reference-world relative preference preorder & \hyperref[sec:p08-bvb501-conditional]{\ref*{sec:p08-bvb501-conditional}} & n/a & \hyperref[def:p08-bvb501-conditional]{\ref*{def:p08-bvb501-conditional}} & \thepage & First compared point is at least as preferred as the second; comparability is additional. \\
\phantomsection\label{entry:symbol-registry-03397}\(>\) & \texttt{\detokenize{>}} & conditional connective & \hyperref[sec:p08-bvb501-conditional]{\ref*{sec:p08-bvb501-conditional}} & n/a & \hyperref[def:p08-bvb501-conditional]{\ref*{def:p08-bvb501-conditional}} & \thepage & Contextual conditional role rather than numeric greater-than. \\
\phantomsection\label{entry:symbol-registry-03398}\(R_w^c\) & \texttt{\detokenize{R_w^c}} & auxiliary inward accessibility with reference fixed & \hyperref[sec:p08-bvb501-conditional]{\ref*{sec:p08-bvb501-conditional}} & n/a & \hyperref[thm:p08-bvb501-conditional-clause]{\ref*{thm:p08-bvb501-conditional-clause}} & \thepage & Recomputed at a different reference point; not a global translation of nested conditionals. \\
\phantomsection\label{entry:symbol-registry-03399}\(G,\gamma\) & \texttt{\detokenize{G,}\allowbreak\detokenize{\gamma}} & atomic and pair-covering composite guards & \hyperref[sec:p08-bvb501-guarded]{\ref*{sec:p08-bvb501-guarded}} & n/a & \hyperref[def:p08-bvb501-guarded]{\ref*{def:p08-bvb501-guarded}} & \thepage & All matrix free variables must be covered. \\
\phantomsection\label{entry:symbol-registry-03400}\(\pi,\rho\) & \texttt{\detokenize{\pi,}\allowbreak\detokenize{\rho}} & PDL program expressions & \hyperref[sec:p08-bvb501-pdl]{\ref*{sec:p08-bvb501-pdl}} & n/a & \hyperref[def:p08-bvb501-pdl]{\ref*{def:p08-bvb501-pdl}} & \thepage & Programs and formulas are defined mutually because of tests. \\
\phantomsection\label{entry:symbol-registry-03401}\(\chi?\) & \texttt{\detokenize{\chi?}} & PDL test program & \hyperref[sec:p08-bvb501-pdl]{\ref*{sec:p08-bvb501-pdl}} & n/a & \hyperref[def:p08-bvb501-pdl]{\ref*{def:p08-bvb501-pdl}} & \thepage & Identity restricted to the formula truth set; valuation dependent. \\
\phantomsection\label{entry:symbol-registry-03402}\(\cup,;,{}^*\) & \texttt{\detokenize{\cup,}\allowbreak\detokenize{;,}\allowbreak\detokenize{{}^*}} & PDL choice composition and finite iteration constructors & \hyperref[sec:p08-bvb501-pdl]{\ref*{sec:p08-bvb501-pdl}} & n/a & \hyperref[def:p08-bvb501-pdl]{\ref*{def:p08-bvb501-pdl}} & \thepage & Composition executes left first; iteration includes zero steps. \\
\phantomsection\label{entry:symbol-registry-03403}\(R_\pi\) & \texttt{\detokenize{R_\pi}} & program transition relation & \hyperref[sec:p08-bvb501-pdl]{\ref*{sec:p08-bvb501-pdl}} & n/a & \hyperref[def:p08-bvb501-pdl]{\ref*{def:p08-bvb501-pdl}} & \thepage & Composite relations follow regular-model equations. \\
\phantomsection\label{entry:symbol-registry-03404}\(\mu p.,\nu p.\) & \texttt{\detokenize{\mu p.,}\allowbreak\detokenize{\nu p.}} & least and greatest propositional fixed-point binders & \hyperref[sec:p08-bvb501-mu]{\ref*{sec:p08-bvb501-mu}} & n/a & \hyperref[def:p08-bvb501-mu]{\ref*{def:p08-bvb501-mu}} & \thepage & Free bound-variable occurrences must be positive; substitutions are hygienic. \\
\phantomsection\label{entry:symbol-registry-03405}\(V[p\leftarrow U]\) & \texttt{\detokenize{V[p\leftarrow U]}} & changed propositional valuation extended to formula truth sets & \hyperref[sec:p08-bvb501-mu]{\ref*{sec:p08-bvb501-mu}} & n/a & \hyperref[def:p08-bvb501-mu]{\ref*{def:p08-bvb501-mu}} & \thepage & Distinct from individual assignment update in first-order modal semantics. \\
\phantomsection\label{entry:symbol-registry-03406}\(F\) & \texttt{\detokenize{F}} & powerset operator induced by a fixed-point body & \hyperref[sec:p08-bvb501-mu]{\ref*{sec:p08-bvb501-mu}} & n/a & \hyperref[def:p08-bvb501-mu]{\ref*{def:p08-bvb501-mu}} & \thepage & Positivity establishes monotonicity before fixed-point existence is used. \\
\phantomsection\label{entry:symbol-registry-03407}\(T,T_\pi\) & \texttt{\detokenize{T,}\allowbreak\detokenize{T_\pi}} & PDL formula and continuation translations into the modal fixed-point calculus & \hyperref[sec:p08-bvb501-mu]{\ref*{sec:p08-bvb501-mu}} & n/a & \hyperref[thm:p08-bvb501-mu-constructions]{\ref*{thm:p08-bvb501-mu-constructions}} & \thepage & All tests and composite program modalities are eliminated. \\
\phantomsection\label{entry:symbol-registry-03408}\(\otimes\) & \texttt{\detokenize{\otimes}} & normal modal fusion & \hyperref[sec:p08-bvb501-combinations]{\ref*{sec:p08-bvb501-combinations}} & n/a & \hyperref[def:p08-bvb501-fusion-product]{\ref*{def:p08-bvb501-fusion-product}} & \thepage & Disjoint modality alphabets with joint-language uniform substitution; not linear tensor. \\
\phantomsection\label{entry:symbol-registry-03409}\(R_h,R_v\) & \texttt{\detokenize{R_h,}\allowbreak\detokenize{R_v}} & horizontal and vertical product accessibility & \hyperref[sec:p08-bvb501-combinations]{\ref*{sec:p08-bvb501-combinations}} & n/a & \hyperref[def:p08-bvb501-fusion-product]{\ref*{def:p08-bvb501-fusion-product}} & \thepage & One coordinate moves while the other stays fixed. \\
\phantomsection\label{entry:symbol-registry-03410}\(I\) & \texttt{\detokenize{I}} & epistemic accessibility relation for the interaction theorem & \hyperref[sec:p08-bvb501-combinations]{\ref*{sec:p08-bvb501-combinations}} & n/a & \hyperref[thm:p08-bvb501-interactions]{\ref*{thm:p08-bvb501-interactions}} & \thepage & Local binary-relation role distinct from the earlier family of partial isomorphisms. \\
\phantomsection\label{entry:symbol-registry-03411}\(D\) & \texttt{\detokenize{D}} & outer individual domain & \hyperref[sec:p08-bvb501-first-order-modal]{\ref*{sec:p08-bvb501-first-order-modal}} & n/a & \hyperref[def:p08-bvb501-varying-domain]{\ref*{def:p08-bvb501-varying-domain}} & \thepage & Nonempty carrier of assignments; distinct from difference modality in its earlier scope. \\
\phantomsection\label{entry:symbol-registry-03412}\(\delta_w\) & \texttt{\detokenize{\delta_w}} & world-specific quantifier domain & \hyperref[sec:p08-bvb501-first-order-modal]{\ref*{sec:p08-bvb501-first-order-modal}} & n/a & \hyperref[def:p08-bvb501-varying-domain]{\ref*{def:p08-bvb501-varying-domain}} & \thepage & Subset of outer D; source convention allows it to be empty. \\
\phantomsection\label{entry:symbol-registry-03413}\(V_w\) & \texttt{\detokenize{V_w}} & world-specific interpretation of individual predicates & \hyperref[sec:p08-bvb501-first-order-modal]{\ref*{sec:p08-bvb501-first-order-modal}} & n/a & \hyperref[def:p08-bvb501-varying-domain]{\ref*{def:p08-bvb501-varying-domain}} & \thepage & Defined on all of D to the appropriate finite power. \\
\phantomsection\label{entry:symbol-registry-03414}\(g[x\mapsto d]\) & \texttt{\detokenize{g[x\mapsto d]}} & individual assignment update & \hyperref[sec:p08-bvb501-first-order-modal]{\ref*{sec:p08-bvb501-first-order-modal}} & n/a & \hyperref[def:p08-bvb501-varying-domain]{\ref*{def:p08-bvb501-varying-domain}} & \thepage & Quantification remains at the current world. \\
\phantomsection\label{entry:symbol-registry-03415}\(\operatorname{Rel}(\chi,p)\) & \texttt{\detokenize{\operatorname{Rel}(\chi,}\allowbreak\detokenize{p)}} & abstract-language relativization to a proposition-defined submodel & \hyperref[sec:p08-bvb501-abstract]{\ref*{sec:p08-bvb501-abstract}} & n/a & \hyperref[def:p08-bvb501-abstract]{\ref*{def:p08-bvb501-abstract}} & \thepage & Requires the evaluation point to satisfy p. \\
\phantomsection\label{entry:symbol-registry-03416}\(B,+,\times,-,1,0\) & \texttt{\detokenize{B,}\allowbreak\detokenize{+,}\allowbreak\detokenize{\times,}\allowbreak\detokenize{-,}\allowbreak\detokenize{1,}\allowbreak\detokenize{0}} & Boolean algebra carrier and operations & \hyperref[sec:p08-bvb501-algebra]{\ref*{sec:p08-bvb501-algebra}} & n/a & \hyperref[def:p08-bvb501-bao]{\ref*{def:p08-bvb501-bao}} & \thepage & Meet with 1 returns the argument; operators are algebraic rather than numerical in this scope. \\
\phantomsection\label{entry:symbol-registry-03417}\(m\) & \texttt{\detokenize{m}} & normal additive Boolean diamond operator & \hyperref[sec:p08-bvb501-algebra]{\ref*{sec:p08-bvb501-algebra}} & n/a & \hyperref[def:p08-bvb501-bao]{\ref*{def:p08-bvb501-bao}} & \thepage & Preserves finite joins and zero. \\
\phantomsection\label{entry:symbol-registry-03418}\(\sim_L\) & \texttt{\detokenize{\sim_L}} & provable formula-equivalence congruence & \hyperref[sec:p08-bvb501-algebra]{\ref*{sec:p08-bvb501-algebra}} & n/a & \hyperref[thm:p08-bvb501-algebra-completeness]{\ref*{thm:p08-bvb501-algebra-completeness}} & \thepage & Used for the Lindenbaum algebra rather than a state bisimulation. \\
\phantomsection\label{entry:symbol-registry-03419}\(R\) & \texttt{\detokenize{R}} & neighborhood point-to-subset relation & \hyperref[sec:p08-bvb501-neighborhood]{\ref*{sec:p08-bvb501-neighborhood}} & n/a & \hyperref[def:p08-bvb501-neighborhood]{\ref*{def:p08-bvb501-neighborhood}} & \thepage & Subset of W times its powerset; explicitly not world-to-world accessibility. \\
\phantomsection\label{entry:symbol-registry-03420}\(N_w\) & \texttt{\detokenize{N_w}} & arbitrary neighborhood family of subsets & \hyperref[sec:p08-bvb501-neighborhood]{\ref*{sec:p08-bvb501-neighborhood}} & n/a & \hyperref[def:p08-bvb501-neighborhood]{\ref*{def:p08-bvb501-neighborhood}} & \thepage & Relational embedding uses all supersets of the successor set. \\
\phantomsection\label{entry:symbol-registry-03421}\(\tau\) & \texttt{\detokenize{\tau}} & topology of open subsets & \hyperref[sec:p08-bvb501-topology]{\ref*{sec:p08-bvb501-topology}} & n/a & \hyperref[def:p08-bvb501-topology]{\ref*{def:p08-bvb501-topology}} & \thepage & Contains empty and whole sets; closed under finite intersection and arbitrary union. \\
\phantomsection\label{entry:symbol-registry-03422}\(\operatorname{int},\operatorname{cl}\) & \texttt{\detokenize{\operatorname{int},}\allowbreak\detokenize{\operatorname{cl}}} & topological interior and closure operations & \hyperref[sec:p08-bvb501-topology]{\ref*{sec:p08-bvb501-topology}} & n/a & \hyperref[thm:p08-bvb501-topological-definability]{\ref*{thm:p08-bvb501-topological-definability}} & \thepage & Interpret box and diamond respectively. \\
\phantomsection\label{entry:symbol-registry-03423}\(\operatorname{pre}_a,\operatorname{box}_a\) & \texttt{\detokenize{\operatorname{pre}_a,}\allowbreak\detokenize{\operatorname{box}_a}} & existential and universal successor-set operators & \hyperref[sec:p08-bvb501-mu]{\ref*{sec:p08-bvb501-mu}} & n/a & \hyperref[def:p08-bvb501-mu]{\ref*{def:p08-bvb501-mu}} & \thepage & Used for least and greatest fixed-point constructions. \\
\phantomsection\label{entry:symbol-registry-03424}\(\mathsf{At},\mathcal L(\mathsf{At})\) & \texttt{\detokenize{\mathsf{At},}\allowbreak\detokenize{\mathcal L(\mathsf{At})}} & source alphabet and its classical propositional modal language & P8 Pacuit neighborhood & object & \hyperref[def:p08-pn-language-model]{\ref*{def:p08-pn-language-model}} & \thepage & Finite or countable alphabet; constants are explicitly included when the alphabet is empty. \\
\phantomsection\label{entry:symbol-registry-03425}\(\wp(W)\) & \texttt{\detokenize{\wp(W)}} & Pacuit source powerset notation & P8 Pacuit neighborhood & object & \hyperref[def:p08-pn-language-model]{\ref*{def:p08-pn-language-model}} & \thepage & The source-local proposition sets are subsets of the nonempty world domain. \\
\phantomsection\label{entry:symbol-registry-03426}\([![\varphi]!]_{\mathcal M}\) & \texttt{\detokenize{[![\varphi]!]_{\mathcal M}}} & formula truth set in the source neighborhood model & P8 Pacuit neighborhood & object & \hyperref[def:p08-pn-language-model]{\ref*{def:p08-pn-language-model}} & \thepage & This brackets convention is kept separately from Chellas's double-bar truth sets. \\
\phantomsection\label{entry:symbol-registry-03427}\(N:W\to\wp(\wp(W))\) & \texttt{\detokenize{N}} & source neighborhood assignment of proposition families & P8 Pacuit neighborhood & object & \hyperref[def:p08-pn-language-model]{\ref*{def:p08-pn-language-model}} & \thepage & This N is not the source N axiom name or the EN system name. \\
\phantomsection\label{entry:symbol-registry-03428}\(\langle\ ]\varphi\) & \texttt{\detokenize{\langle\ ]\varphi}} & source alternative modality with a neighborhood contained in the truth set & P8 Pacuit neighborhood & n/a & \hyperref[def:p08-pn-language-model]{\ref*{def:p08-pn-language-model}} & \thepage & Unary formula constructor; no separate numerical precedence rank is imposed here. It agrees with exact-membership box only on supplemented models. \\
\phantomsection\label{entry:symbol-registry-03429}\((E,EM,EC,EN,EMC,ECM,EMCN,EK,K)\) & \texttt{\detokenize{E,}\allowbreak\detokenize{EM,}\allowbreak\detokenize{EC,}\allowbreak\detokenize{EN,}\allowbreak\detokenize{EMC,}\allowbreak\detokenize{ECM,}\allowbreak\detokenize{EMCN,}\allowbreak\detokenize{EK,}\allowbreak\detokenize{K}} & source theorem-set system names & P8 Pacuit neighborhood & object & \hyperref[def:p08-pn-logics]{\ref*{def:p08-pn-logics}} & \thepage & ECM is EMC; EMCN is K; EK is E with distribution and without necessitation. \\
\phantomsection\label{entry:symbol-registry-03430}\((M,C,N)\) & \texttt{\detokenize{M,}\allowbreak\detokenize{C,}\allowbreak\detokenize{N}} & source axiom-schema identifiers for supplementation intersection and unit & P8 Pacuit neighborhood & object & \hyperref[thm:p08-pn-correspondence]{\ref*{thm:p08-pn-correspondence}} & \thepage & These are schemas with distinct roles from the M model or N neighborhood assignment. \\
\phantomsection\label{entry:symbol-registry-03431}\(\Gamma\vdash_L\varphi\) & \texttt{\detokenize{\Gamma\vdash_L\varphi}} & source local finite-implication deduction relation & P8 Pacuit neighborhood & relation & \hyperref[def:p08-pn-local-deduction]{\ref*{def:p08-pn-local-deduction}} & \thepage & Modal inferences require separate theorem proofs; an empty premise conjunction means theoremhood. \\
\phantomsection\label{entry:symbol-registry-03432}\(\models^{\ell}_{\mathscr F}\) & \texttt{\detokenize{\models^{\ell}_{\mathscr F}}} & local semantic consequence on a frame class & P8 Pacuit neighborhood & relation & \hyperref[def:p08-pn-local-deduction]{\ref*{def:p08-pn-local-deduction}} & \thepage & Premises are true at one world; distinguished from model-global and frame-global consequence. \\
\phantomsection\label{entry:symbol-registry-03433}\(M_L\) & \texttt{\detokenize{M_L}} & source world set of maximal L-consistent theories & P8 Pacuit neighborhood & object & \hyperref[def:p08-pn-proof-sets]{\ref*{def:p08-pn-proof-sets}} & \thepage & Source ML is not an unconditionally uncountable set; consistency ensures it is nonempty. \\
\phantomsection\label{entry:symbol-registry-03434}\(|\varphi|_L\) & \texttt{\detokenize{|\varphi|_L}} & source proof set consisting of maximal worlds containing the formula & P8 Pacuit neighborhood & object & \hyperref[def:p08-pn-proof-sets]{\ref*{def:p08-pn-proof-sets}} & \thepage & Bars are the published typed proposition-set notation and not numerical cardinality. \\
\phantomsection\label{entry:symbol-registry-03435}\(P_L\) & \texttt{\detokenize{P_L}} & family of all formula proof sets & P8 Pacuit neighborhood & object & \hyperref[def:p08-pn-proof-sets]{\ref*{def:p08-pn-proof-sets}} & \thepage & Canonical neighborhoods can differ on subsets outside this Boolean algebra. \\
\phantomsection\label{entry:symbol-registry-03436}\(N_L^{\min},N_L^{\max}\) & \texttt{\detokenize{N_L^{\min},}\allowbreak\detokenize{N_L^{\max}}} & minimal and maximal canonical neighborhood functions & P8 Pacuit neighborhood & object & \hyperref[def:p08-pn-canonical]{\ref*{def:p08-pn-canonical}} & \thepage & The maximum adds all non-proof subsets and no rejected proof sets. \\
\phantomsection\label{entry:symbol-registry-03437}\(\mathcal U^{\rm mon}\) & \texttt{\detokenize{\mathcal U^{\rm mon}}} & source supplementation of a proposition family & P8 Pacuit neighborhood & object & \hyperref[def:p08-pn-supplementation]{\ref*{def:p08-pn-supplementation}} & \thepage & A member contains some original member; an empty family remains empty. \\
\phantomsection\label{entry:symbol-registry-03438}\(N^R,R_N\) & \texttt{\detokenize{N^R,}\allowbreak\detokenize{R_N}} & principal relational neighborhoods and the relation induced by augmented cores & P8 Pacuit neighborhood & object & \hyperref[thm:p08-pn-augmentation]{\ref*{thm:p08-pn-augmentation}} & \thepage & Pointwise conversion requires augmentation; filter completeness alone does not establish a principal representation. \\
\phantomsection\label{entry:symbol-registry-03439}\(\bigcap N(w)\) & \texttt{\detokenize{\bigcap N(w)}} & source core of a neighborhood family & P8 Pacuit neighborhood & object & \hyperref[def:p08-pn-closure-properties]{\ref*{def:p08-pn-closure-properties}} & \thepage & The empty-family intersection is W; an augmented family contains its core. \\
\phantomsection\label{entry:symbol-registry-03440}\(W^N\) & \texttt{\detokenize{W^N}} & normal-world subset in the source impossible-world relational model & P8 Pacuit neighborhood & object & \hyperref[def:p08-pn-impossible-worlds]{\ref*{def:p08-pn-impossible-worlds}} & \thepage & Outside this subset box is always false; Boolean truth remains classical. \\
\phantomsection\label{entry:symbol-registry-03441}\(B_\Gamma\) & \texttt{\detokenize{B_\Gamma}} & unboxed members of a maximal theory used in the relational canonical construction & P8 Pacuit neighborhood & object & \hyperref[annex:p8.pn.augmentation]{\ref*{annex:p8.pn.augmentation}} & \thepage & Finite conjunction consistency supplies successors; the zero-member case uses unit or necessitation explicitly. \\
\phantomsection\label{entry:symbol-registry-03442}\(q,[w]_\Sigma,[X]_\Sigma\) & \texttt{\detokenize{q,}\allowbreak\detokenize{[w]_\Sigma,}\allowbreak\detokenize{[X]_\Sigma}} & source filtration quotient map world classes and subset images & P8 Pacuit neighborhood & object & \hyperref[def:p08-pn-filtration]{\ref*{def:p08-pn-filtration}} & \thepage & Tracked truth sets are saturated; arbitrary images are not presumed to preserve intersections. \\
\phantomsection\label{entry:symbol-registry-03443}\(\mathcal B([w])\) & \texttt{\detokenize{\mathcal B([w])}} & finite generating family of tracked accepted quotient truth sets & P8 Pacuit neighborhood & object & \hyperref[def:p08-pn-filtration]{\ref*{def:p08-pn-filtration}} & \thepage & Positive intersections and supplementation are applied only for chosen C and M; the unit is added separately for N. \\
\phantomsection\label{entry:symbol-registry-03444}\(\Sigma_\varphi,D,\Gamma_D,W_\varphi\) & \texttt{\detokenize{\Sigma_\varphi,}\allowbreak\detokenize{D,}\allowbreak\detokenize{\Gamma_D,}\allowbreak\detokenize{W_\varphi}} & subformula observations consistent descriptions their selected maximal extensions and finite worlds & P8 Pacuit neighborhood & object & \hyperref[def:p08-pn-finite-canonical]{\ref*{def:p08-pn-finite-canonical}} & \thepage & Inconsistent descriptions have no representative; the construction is not a consistency decision algorithm. \\
\phantomsection\label{entry:symbol-registry-03445}\(\lambda_D,\lambda_Y\) & \texttt{\detokenize{\lambda_D,}\allowbreak\detokenize{\lambda_Y}} & Boolean labels of a description world and a subset of finite worlds & P8 Pacuit neighborhood & object & \hyperref[def:p08-pn-finite-canonical]{\ref*{def:p08-pn-finite-canonical}} & \thepage & The empty subset label is a contradiction; labels are Boolean combinations of the subformula observations. \\
\phantomsection\label{entry:symbol-registry-03446}\(|\psi|_\varphi,N_\varphi\) & \texttt{\detokenize{|\psi|_\varphi,}\allowbreak\detokenize{N_\varphi}} & finite canonical membership set and label-based neighborhood function & P8 Pacuit neighborhood & object & \hyperref[def:p08-pn-finite-canonical]{\ref*{def:p08-pn-finite-canonical}} & \thepage & Agreement is proved on observations their Boolean closure and boxes of that closure; not every arbitrary iterated formula. \\
\phantomsection\label{entry:symbol-registry-03447}\(\mathcal A,m_N\) & \texttt{\detokenize{\mathcal A,}\allowbreak\detokenize{m_N}} & general-frame admissible Boolean algebra and neighborhood membership operation & P8 Pacuit neighborhood & object & \hyperref[def:p08-pn-general-frames]{\ref*{def:p08-pn-general-frames}} & \thepage & The canonical algebra is PL; admissible valuation completeness is distinct from unrestricted ordinary-frame completeness. \\
\phantomsection\label{entry:symbol-registry-03448}\(\Delta_n,Z,X,Y\) & \texttt{\detokenize{\Delta_n,}\allowbreak\detokenize{Z,}\allowbreak\detokenize{X,}\allowbreak\detokenize{Y}} & distinct chain-root theories diagonal non-proof complement witness and its union with a proof set & \hyperref[annex:p8.pn.minimal-not-monotone]{\ref*{annex:p8.pn.minimal-not-monotone}} & object & \hyperref[annex:p8.pn.minimal-not-monotone]{\ref*{annex:p8.pn.minimal-not-monotone}} & \thepage & The diagonal is wholly outside X; this condition is needed for the Boolean intersection argument. \\
\phantomsection\label{entry:symbol-registry-03449}\(A_i,B_j,A_I,r\) & \texttt{\detokenize{A_i,}\allowbreak\detokenize{B_j,}\allowbreak\detokenize{A_I,}\allowbreak\detokenize{r}} & old-component matrix truth sets positive antecedent intersections and the fresh countermodel root & \hyperref[annex:p8.pn.one-step]{\ref*{annex:p8.pn.one-step}} & object & \hyperref[annex:p8.pn.one-step]{\ref*{annex:p8.pn.one-step}} & \thepage & Root families are defined using already evaluated old components; they do not depend circularly on root truth values. \\
\phantomsection\label{entry:symbol-registry-03579}\(\mathsf{MPL}\) & \texttt{\detokenize{\mathsf{MPL}}} & selected defined-negation minimal propositional calculus & P8 minimal-and-intermediate-presentations & n/a & \hyperref[def:p08-mi-language]{\ref*{def:p08-mi-language}} & \thepage & Empty-discharge ND with no falsity elimination; not an unchanged Johansson primitive-negation presentation. \\
\phantomsection\label{entry:symbol-registry-03580}\(\mathsf H_0\) & \texttt{\detokenize{\mathsf H_0}} & common full-language intuitionistic Hilbert base & P8 minimal-and-intermediate-presentations & n/a & \hyperref[def:p08-mi-language]{\ref*{def:p08-mi-language}} & \thepage & Nine schemas and modus ponens; delete only the KC or LC extra schema from their sheets. \\
\phantomsection\label{entry:symbol-registry-03581}\(\mathsf H_{\mathsf{KC}}\) & \texttt{\detokenize{\mathsf H_{\mathsf{KC}}}} & modern KC Hilbert presentation & P8 minimal-and-intermediate-presentations & n/a & \hyperref[def:p08-mi-language]{\ref*{def:p08-mi-language}} & \thepage & Common Hilbert base plus weak excluded middle; no substitution rule on arbitrary premise lines. \\
\phantomsection\label{entry:symbol-registry-03582}\(\mathsf H_{\mathsf{LC}}\) & \texttt{\detokenize{\mathsf H_{\mathsf{LC}}}} & modern LC Hilbert presentation & P8 minimal-and-intermediate-presentations & n/a & \hyperref[def:p08-mi-language]{\ref*{def:p08-mi-language}} & \thepage & Common Hilbert base plus prelinearity; source matrix retained separately. \\
\phantomsection\label{entry:symbol-registry-03583}\(F_{\bot}\) & \texttt{\detokenize{F_{\bot}}} & persistent upset of worlds forcing minimal falsum & P8 minimal-and-intermediate-presentations & n/a & \hyperref[def:p08-mi-models]{\ref*{def:p08-mi-models}} & \thepage & May be nonempty; for KC and LC it is fixed empty. \\
\phantomsection\label{entry:symbol-registry-03584}\(f\) & \texttt{\detokenize{f}} & distinguished new atom replacing minimal falsum & P8 minimal positive-language bridge & n/a & \hyperref[thm:p08-mi-minimal-falsum-recovery]{\ref*{thm:p08-mi-minimal-falsum-recovery}} & \thepage & Chosen outside the original atom set; not the new-bottom algebra element of the separate algebra definition. \\
\phantomsection\label{entry:symbol-registry-03585}\(\alpha,\alpha^*,f\) & \texttt{\detokenize{alpha;alpha-star;f}} & finite Heyting algebra and its new-bottom extension and new least element & P8 new-bottom algebra construction & n/a & \hyperref[def:p08-mi-new-bottom-algebra]{\ref*{def:p08-mi-new-bottom-algebra}} & \thepage & Algebraic roles only; alpha does not denote a type and f does not denote the distinguished syntactic atom in this scope. \\
\phantomsection\label{entry:symbol-registry-03586}\(0,1,\Rightarrow\) & \texttt{\detokenize{0;1;Rightarrow}} & bottom and top and Godel-chain conditional operation & P8 finite numerical-chain bridge & n/a & \hyperref[thm:p08-mi-finite-chain-bridge]{\ref*{thm:p08-mi-finite-chain-bridge}} & \thepage & Rightarrow is an algebraic operation and not a sequent arrow in this theorem. \\
\phantomsection\label{entry:symbol-registry-03587}\(\derives_{\mathsf{MPL}},\derives_{\mathsf H_0},\derives_{\mathsf H_{\mathsf{KC}}},\derives_{\mathsf H_{\mathsf{LC}}}\) & \texttt{\detokenize{\derives}} & derivability in the four exact named propositional presentations & P8 minimal-and-intermediate-presentations & n/a & \hyperref[def:p08-mi-language]{\ref*{def:p08-mi-language}} & \thepage & Distinct proof relations with the language and premise conventions stated in the body and sheets. \\
\phantomsection\label{entry:symbol-registry-03588}\((V,D,A_i,I)\) & \texttt{\detokenize{V,}\allowbreak\detokenize{D,}\allowbreak\detokenize{A_i,}\allowbreak\detokenize{I}} & source virtual possible actual individual domains and reference indices & P8 Scott Advice & object & \hyperref[def:p08-sa-domains]{\ref*{def:p08-sa-domains}} & \thepage & D is nonempty; Ai is a subset of D and may be empty; pointed comparisons use nonempty I. \\
\phantomsection\label{entry:symbol-registry-03589}\(\bar a\) & \texttt{\detokenize{\bar a}} & source rigid individual constant for a virtual or possible individual & P8 Scott Advice & object & \hyperref[def:p08-sa-domains]{\ref*{def:p08-sa-domains}} & \thepage & This bar is not a numeral bar; its whole concept is the constant function with value a. \\
\phantomsection\label{entry:symbol-registry-03590}\(\forall_{\cdot},\exists_{\cdot}\) & \texttt{\detokenize{\forall_{\cdot},}\allowbreak\detokenize{\exists_{\cdot}}} & source dotted actual-domain quantifiers & P8 Scott Advice & n/a & \hyperref[def:p08-sa-domains]{\ref*{def:p08-sa-domains}} & \thepage & At index i these quantifiers range over Ai; ordinary variables and quantifiers range over fixed D. \\
\phantomsection\label{entry:symbol-registry-03591}\(\widetilde A\) & \texttt{\detokenize{\widetilde A}} & source index-varying actuality predicate & P8 Scott Advice & n/a & \hyperref[def:p08-sa-domains]{\ref*{def:p08-sa-domains}} & \thepage & True only when the term is defined and its current value is in Ai. \\
\phantomsection\label{entry:symbol-registry-03592}\(|x|,\operatorname{Card}(\tau)\) & \texttt{\detokenize{|x|,}\allowbreak\detokenize{\operatorname{Card}(\tau)}} & source virtual cardinal equipollence class and its image predicate & P8 Scott Advice & object & \hyperref[def:p08-sa-virtuals]{\ref*{def:p08-sa-virtuals}} & \thepage & This class-theoretic illustration is conditional on the chosen class system; not numerical cardinality of worlds. \\
\phantomsection\label{entry:symbol-registry-03593}\(g[D]\) & \texttt{\detokenize{g[D]}} & image of the source's indexing map for restricted virtual quantification & P8 Scott Advice & object & \hyperref[def:p08-sa-virtuals]{\ref*{def:p08-sa-virtuals}} & \thepage & Image quantification does not expand ordinary D-valued assignments to all virtual objects. \\
\phantomsection\label{entry:symbol-registry-03594}\(\|\Phi\|,\|\Phi\|_i,2^I\) & \texttt{\detokenize{\|\Phi\|,}\allowbreak\detokenize{\|\Phi\|_i,}\allowbreak\detokenize{2^I}} & source whole propositional concept its point value and its type & P8 Scott Advice & object & \hyperref[def:p08-sa-concepts]{\ref*{def:p08-sa-concepts}} & \thepage & 2 is the two truth values; a whole truth concept is not its extension at a single index. \\
\phantomsection\label{entry:symbol-registry-03595}\(V^{(I)},V^I\) & \texttt{\detokenize{V^{(I)},}\allowbreak\detokenize{V^I}} & source partial and total individual-concept function spaces & P8 Scott Advice & object & \hyperref[def:p08-sa-concepts]{\ref*{def:p08-sa-concepts}} & \thepage & The parenthesized exponent denotes partial functions; undefinedness is not a virtual null individual. \\
\phantomsection\label{entry:symbol-registry-03596}\(=,\equiv,\approx\) & \texttt{\detokenize{=,}\allowbreak\detokenize{\equiv,}\allowbreak\detokenize{\approx}} & source local extension equality whole-concept equality and incidence & P8 Scott Advice & relation & \hyperref[def:p08-sa-concepts]{\ref*{def:p08-sa-concepts}} & \thepage & Equiv is this source-local intensional equality and not elementary equivalence; incidence is separately defined later. \\
\phantomsection\label{entry:symbol-registry-03597}\(\mathsf I,\mathsf I_{\cdot}\) & \texttt{\detokenize{\mathsf I,}\allowbreak\detokenize{\mathsf I_{\cdot}}} & source unique possible and actual individual description binders & P8 Scott Advice & n/a & \hyperref[def:p08-sa-concepts]{\ref*{def:p08-sa-concepts}} & \thepage & An improper description is undefined; the dotted binder is reduced by an actual existential test. \\
\phantomsection\label{entry:symbol-registry-03598}\(\ast\) & \texttt{\detokenize{\ast}} & source totally undefined individual concept & P8 Scott Advice & object & \hyperref[def:p08-sa-concepts]{\ref*{def:p08-sa-concepts}} & \thepage & The value of a contradictory description is a partial function with empty domain and not a chosen element of V. \\
\phantomsection\label{entry:symbol-registry-03599}\(\bar R,\widetilde R,R_i\) & \texttt{\detokenize{\bar R,}\allowbreak\detokenize{\widetilde R,}\allowbreak\detokenize{R_i}} & source strict fixed and varying extensional predicates and their relations & P8 Scott Advice & n/a & \hyperref[def:p08-sa-operators]{\ref*{def:p08-sa-operators}} & \thepage & Both arguments must have defined values; equality has the distinct both-undefined convention. \\
\phantomsection\label{entry:symbol-registry-03600}\(\boxed{\mathscr P},\mathscr P_i\) & \texttt{\detokenize{\boxed{\mathscr P},}\allowbreak\detokenize{\mathscr P_i}} & source general monadic propositional operator and its subset-family representation & P8 Scott Advice & n/a & \hyperref[def:p08-sa-operators]{\ref*{def:p08-sa-operators}} & \thepage & An arbitrary function from truth concepts to truth concepts is exactly represented by one family per index. \\
\phantomsection\label{entry:symbol-registry-03601}\(\Box,\boxed\rho,\boxed{\leftrightarrow}\) & \texttt{\detokenize{\Box,}\allowbreak\detokenize{\boxed\rho,}\allowbreak\detokenize{\boxed{\leftrightarrow}}} & source universal relational and real-line progressive modalities & P8 Scott Advice & n/a & \hyperref[def:p08-sa-operators]{\ref*{def:p08-sa-operators}} & \thepage & Universal box quantifies all indices; relational box tests rho successors; progressive is ordinary topological interior. \\
\phantomsection\label{entry:symbol-registry-03602}\(\circ,t_i\) & \texttt{\detokenize{\circ,}\allowbreak\detokenize{t_i}} & source index composition and individual transformation family & P8 Scott Advice & object & \hyperref[def:p08-sa-operators]{\ref*{def:p08-sa-operators}} & \thepage & Source orientation is tj(ti(a)) equals t at i-circle-j of a; its modal example tests all i-circle-j. \\
\phantomsection\label{entry:symbol-registry-03603}\(\$\) & \texttt{\detokenize{\$}} & source general variable-binding term or formula maker placeholder & P8 Scott Advice & n/a & \hyperref[def:p08-sa-binding]{\ref*{def:p08-sa-binding}} & \thepage & Its arguments are D-indexed whole concept functions; the codomain distinguishes term from formula makers. \\
\phantomsection\label{entry:symbol-registry-03604}\(f,F\) & \texttt{\detokenize{f,}\allowbreak\detokenize{F}} & source whole individual-concept and propositional-concept argument functions of the binder & P8 Scott Advice binding & object & \hyperref[def:p08-sa-binding]{\ref*{def:p08-sa-binding}} & \thepage & These source functions have domain D; the earlier editorial bridge F has its separately declared type from 2I to 2I. \\
\phantomsection\label{entry:symbol-registry-03605}\(F:2^I\to2^I\) & \texttt{\detokenize{F}} & editorial arbitrary truth-concept operation in the exact source representation proof & P8 Scott Advice neighborhood bridge & object & \hyperref[thm:p08-sa-neighborhood-bridge]{\ref*{thm:p08-sa-neighborhood-bridge}} & \thepage & This F is locally typed separately from the source binder argument function with domain D. \\
\phantomsection\label{entry:symbol-registry-03606}\(\mathfrak A\) & \texttt{\detokenize{\mathfrak A}} & source complete interpretation with domains indices and logical types & P8 Scott Advice & object & \hyperref[def:p08-sa-binding]{\ref*{def:p08-sa-binding}} & \thepage & Logical validity varies interpretations and D-valued free-variable assignments before evaluating at each index. \\
\phantomsection\label{entry:symbol-registry-03607}\(S^I,a_i,(\|\tau\|_i)_i\) & \texttt{\detokenize{S^I,}\allowbreak\detokenize{a_i,}\allowbreak\detokenize{(\|\tau\|_i)_i}} & source history space individual state and diagonal state of a denoted history & P8 Scott Advice & object & \hyperref[def:p08-sa-incidence]{\ref*{def:p08-sa-incidence}} & \thepage & There are two function layers: individuals are total histories while term concepts are partial functions selecting histories. \\
\phantomsection\label{entry:symbol-registry-03608}\(C_i,\widetilde C,H_i(a)\) & \texttt{\detokenize{C_i,}\allowbreak\detokenize{\widetilde C,}\allowbreak\detokenize{H_i(a)}} & source current counterpart relation predicate and counterpart profile & P8 Scott Advice & relation & \hyperref[def:p08-sa-counterparts]{\ref*{def:p08-sa-counterparts}} & \thepage & The counterpart-profile incidence is a separate construction from state incidence; its one-direction implication has a full countermodel to the converse. \\
\phantomsection\label{entry:symbol-registry-03609}\(\mathfrak A=\langle A,F,R\rangle\) & \texttt{\detokenize{\mathfrak A;A;F;R}} & possible pragmatic interpretation & P8 Montague Synthese 1970 ordinary pragmatics & n/a & \hyperref[def:p08-montague-pragmatic-language]{\ref*{def:p08-montague-pragmatic-language}} & \thepage & A assigns actuality sets; F assigns individual/predicate intensions; R assigns arbitrary operator relations. \\
\phantomsection\label{entry:symbol-registry-03610}\(I=DA,A_i,U\) & \texttt{\detokenize{I;DA;A_i;U}} & index set actuality sets and union of possible individuals & P8 Montague Synthese 1970 & n/a & \hyperref[def:p08-montague-pragmatic-language]{\ref*{def:p08-montague-pragmatic-language}} & \thepage & Source DA is domain of A; quantifiers range over U not only A\_i; empty U is permitted but contributes no countermodels to sentence consequence. \\
\phantomsection\label{entry:symbol-registry-03611}\(E\) & \texttt{\detokenize{E}} & logical actuality predicate & P8 Montague Synthese 1970 & object & \hyperref[def:p08-montague-pragmatic-language]{\ref*{def:p08-montague-pragmatic-language}} & \thepage & E[zeta] tests membership of the denoted possible individual in A\_i; not the modern congruential modal system E. \\
\phantomsection\label{entry:symbol-registry-03612}\(R_N\) & \texttt{\detokenize{R_N}} & relation between an index and a proposition truth set & P8 Montague Synthese 1970 ordinary pragmatics & relation & \hyperref[def:p08-montague-pragmatic-language]{\ref*{def:p08-montague-pragmatic-language}} & \thepage & Subset of I times power set of I; no closure conditions are imposed by a possible interpretation. \\
\phantomsection\label{entry:symbol-registry-03613}\(\Lambda,\mathrm F,\mathrm T\) & \texttt{\detokenize{\Lambda;\mathrm F;\mathrm T}} & empty tuple and zero-place relation truth values & P8 Montague Synthese 1970 & n/a & \hyperref[def:p08-montague-empty-relation]{\ref*{def:p08-montague-empty-relation}} & \thepage & Lambda is the empty sequence/set; false is Lambda and true is its singleton; Roman F is distinct from the interpretation function F. \\
\phantomsection\label{entry:symbol-registry-03614}\(\operatorname{Ext}_{i,\mathfrak A}(\xi)\) & \texttt{\detokenize{\operatorname{Ext}}} & indexed extension of a term or formula & P8 Montague Synthese 1970 & n/a & \hyperref[def:p08-montague-pragmatic-extension]{\ref*{def:p08-montague-pragmatic-extension}} & \thepage & Terms give functions on the appropriate assignment set; formulas give its subsets. \\
\phantomsection\label{entry:symbol-registry-03615}\(\operatorname{Int}_{\mathfrak A}(\xi)\) & \texttt{\detokenize{\operatorname{Int}}} & index-to-extension intension & P8 Montague Synthese 1970 & n/a & \hyperref[def:p08-montague-pragmatic-extension]{\ref*{def:p08-montague-pragmatic-extension}} & \thepage & At each index returns Ext; no separate hierarchy of indirect intensions is introduced. \\
\phantomsection\label{entry:symbol-registry-03616}\(x[n:=y]\) & \texttt{\detokenize{x[n:=y]}} & replacement of an ordinary assignment coordinate & P8 Montague Synthese 1970 ordinary and extended pragmatics & n/a & \hyperref[def:p08-montague-pragmatic-extension]{\ref*{def:p08-montague-pragmatic-extension}} & \thepage & The source uses function-replacement subscripts/superscripts; replacement is explicit here. \\
\phantomsection\label{entry:symbol-registry-03617}\(G_{n,k}\) & \texttt{\detokenize{G_{n,k}}} & k-place predicate variable numbered n & P8 Montague Synthese 1970 intensional logic & n/a & \hyperref[def:p08-montague-intensional-language]{\ref*{def:p08-montague-intensional-language}} & \thepage & k is a nonnegative arity; variables denote predicate intensions rather than index-local relations. \\
\phantomsection\label{entry:symbol-registry-03618}\(\langle s_0,\ldots,s_{r-1}\rangle\) & \texttt{\detokenize{\langle s_0,}\allowbreak\detokenize{\ldots,}\allowbreak\detokenize{s_{r-1}\rangle}} & predicate-constant argument type sequence & P8 Montague Synthese 1970 intensional logic & n/a & \hyperref[def:p08-montague-intensional-language]{\ref*{def:p08-montague-intensional-language}} & \thepage & Minus one is an individual argument; k at least zero is a predicate-intension argument of arity k; type <0> is unary over propositions. \\
\phantomsection\label{entry:symbol-registry-03619}\(\mathsf T G\phi\) & \texttt{\detokenize{\mathsf T}} & unique predicate description with empty default & P8 Montague Synthese 1970 intensional logic & n/a & \hyperref[def:p08-montague-intensional-language]{\ref*{def:p08-montague-intensional-language}} & \thepage & Binds predicate variable G in phi; no individual description is admitted; Roman truth value T is a different role. \\
\phantomsection\label{entry:symbol-registry-03620}\(\mathcal P_k,D_k\) & \texttt{\detokenize{\mathcal P_k;D_k}} & full predicate-intension and typed variable domains & P8 Montague Synthese 1970 intensional logic & n/a & \hyperref[def:p08-montague-intensional-interpretation]{\ref*{def:p08-montague-intensional-interpretation}} & \thepage & P\_k consists of all functions I to power set of U\^{}k; D\_minus-one=U; these are not arbitrary Henkin subdomains. \\
\phantomsection\label{entry:symbol-registry-03621}\(S,x_{n,k}\) & \texttt{\detokenize{S;x_{n,k}}} & all typed assignment systems and their coordinates & P8 Montague Synthese 1970 intensional logic & n/a & \hyperref[def:p08-montague-intensional-interpretation]{\ref*{def:p08-montague-intensional-interpretation}} & \thepage & Coordinate (n,-1) is individual v\_n; coordinate (n,k) is predicate G\_(n,k) for nonnegative k. \\
\phantomsection\label{entry:symbol-registry-03622}\(x[(n,k):=Y]\) & \texttt{\detokenize{x[(n,}\allowbreak\detokenize{k):=Y]}} & typed coordinate replacement & P8 Montague Synthese 1970 intensional logic & n/a & \hyperref[def:p08-montague-intensional-extension]{\ref*{def:p08-montague-intensional-extension}} & \thepage & Y belongs to the correct typed coordinate domain. \\
\phantomsection\label{entry:symbol-registry-03623}\(\llbracket\xi\rrbracket_i^x\) & \texttt{\detokenize{\llbracket\xi\rrbracket_i^x}} & term extension evaluated on a typed assignment & P8 Montague Synthese 1970 intensional logic & n/a & \hyperref[def:p08-montague-intensional-interpretation]{\ref*{def:p08-montague-intensional-interpretation}} & \thepage & Abbreviates Ext\_(i,A)(xi)(x); a predicate term still denotes the whole I-indexed intension. \\
\phantomsection\label{entry:symbol-registry-03624}\(O_k\) & \texttt{\detokenize{O_k}} & empty predicate intension of arity k & P8 Montague Synthese 1970 intensional logic & n/a & \hyperref[def:p08-montague-intensional-interpretation]{\ref*{def:p08-montague-intensional-interpretation}} & \thepage & Value is empty relation at every index; graph is I times singleton of empty set; default for failed description uniqueness. \\
\phantomsection\label{entry:symbol-registry-03625}\(\widehat{\bar u}\phi,\widehat\phi\) & \texttt{\detokenize{\widehat{\bar u}\phi;\widehat\phi}} & source predicate/property/proposition names & P8 Montague Synthese 1970 intensional logic & n/a & \hyperref[def:p08-montague-predicate-names]{\ref*{def:p08-montague-predicate-names}} & \thepage & Descriptions uniquely defining the formula's full indexed predicate; empty variable list gives the proposition that phi. \\
\phantomsection\label{entry:symbol-registry-03626}\(N'\) & \texttt{\detokenize{N'}} & predicate of propositions corresponding to a pragmatic operator & P8 Montague Synthese 1970 reduction & n/a & \hyperref[thm:p08-montague-pragmatic-reduction]{\ref*{thm:p08-montague-pragmatic-reduction}} & \thepage & Type <0>; its local extension accepts exactly the propositions whose truth sets are related by R\_N to the current index. \\
\phantomsection\label{entry:symbol-registry-03627}\(B,\mathit{Se}\) & \texttt{\detokenize{B;\mathit{Se}}} & empirical belief and seeming predicate constants & P8 Montague Synthese 1970 attitude examples & n/a & \hyperref[def:p08-montague-belief-seems]{\ref*{def:p08-montague-belief-seems}} & \thepage & Types <-1,0> and <-1,1,-1>; no logical psychological axioms imposed. \\
\phantomsection\label{entry:symbol-registry-03628}\(\langle m,n,p\rangle,Y_{q,x}\) & \texttt{\detokenize{\langle m,}\allowbreak\detokenize{n,}\allowbreak\detokenize{p\rangle;Y_{q,x}}} & extended operator degree and bound-variable operand predicate & P8 Montague Synthese 1970 extended pragmatics & n/a & \hyperref[def:p08-montague-extended-pragmatics]{\ref*{def:p08-montague-extended-pragmatics}} & \thepage & m bound individual variables; n individual arguments; p formula arguments generating m-place intensions under fixed remaining assignment. \\
\phantomsection\label{entry:symbol-registry-03629}\(\mathcal N_N(i)\) & \texttt{\detokenize{\mathcal N_N}} & modern neighborhood family obtained as a section of R\_N & P8 Montague Synthese 1970 neighborhood comparison & n/a & \hyperref[thm:p08-montague-neighborhood-bridge]{\ref*{thm:p08-montague-neighborhood-bridge}} & \thepage & J belongs iff (i,J) belongs to R\_N; modern closure systems are not attributed as Montague's axiomatizations. \\
\phantomsection\label{entry:symbol-registry-03630}\(\operatorname{Eq}_k,\operatorname{Empty}_k\) & \texttt{\detokenize{\operatorname{Eq}_k;\operatorname{Empty}_k}} & eliminable equality and emptiness formulas for predicate intensions & P8 Montague Synthese 1970 description-elimination proof & n/a & \hyperref[thm:p08-montague-descriptions]{\ref*{thm:p08-montague-descriptions}} & \thepage & Defined by universal-index agreement on all individual tuples; k=0 has no individual quantifiers. \\
\phantomsection\label{entry:symbol-registry-03631}\(\operatorname{Unique}_\psi,\operatorname{Select}_\psi\) & \texttt{\detokenize{\operatorname{Unique}_\psi;\operatorname{Select}_\psi}} & description-elimination abbreviations & P8 Montague Synthese 1970 description-elimination proof & n/a & \hyperref[thm:p08-montague-descriptions]{\ref*{thm:p08-montague-descriptions}} & \thepage & Explicitly defined formulas selecting unique satisfying predicate or the empty default; not new semantic primitives. \\
\phantomsection\label{entry:symbol-registry-03754}\(U,D\) & \texttt{\detokenize{U,}\allowbreak\detokenize{D}} & outer domain and nonempty existing inner domain & P8 derived negative-free calculus & n/a & \hyperref[def:p08-negative-models]{\ref*{def:p08-negative-models}} & \thepage & D is a subset of U; both are nonempty. \\
\phantomsection\label{entry:symbol-registry-03755}\(X\) & \texttt{\detokenize{X}} & set of free parameters required to have existing values & P8 derived negative-free calculus & n/a & \hyperref[def:p08-negative-models]{\ref*{def:p08-negative-models}} & \thepage & Other free variables remain outer-valued; the parameter context records the convention explicitly. \\
\phantomsection\label{entry:symbol-registry-03756}\(\models_{-,X}\) & \texttt{\detokenize{\models_{-,X}}} & negative semantic consequence under existing-parameter requirements & P8 derived negative-free calculus & n/a & \hyperref[def:p08-negative-models]{\ref*{def:p08-negative-models}} & \thepage & Sentence consequence is independent of X. \\
\phantomsection\label{entry:symbol-registry-03757}\(E\) & \texttt{\detokenize{E}} & fresh target unary predicate coding the existing domain & P8 derived negative-free target language & n/a & \hyperref[def:p08-negative-translation]{\ref*{def:p08-negative-translation}} & \thepage & Not an ordinary source predicate; E! is the source existence atom. \\
\phantomsection\label{entry:symbol-registry-03758}\(V_t(z)\) & \texttt{\detokenize{V_t(z)}} & target formula expressing the defined outer value of a source term & P8 derived negative-free target language & n/a & \hyperref[def:p08-negative-translation]{\ref*{def:p08-negative-translation}} & \thepage & Can have no satisfying value for a failed description; is functional when it has a value. \\
\phantomsection\label{entry:symbol-registry-03759}\(\delta\) & \texttt{\detokenize{\delta}} & finite simultaneous term-value and formula translation & P8 derived negative-free target language & n/a & \hyperref[def:p08-negative-translation]{\ref*{def:p08-negative-translation}} & \thepage & Delta on formulas is recursively defined; it does not map raw strict equality to unguarded target equality. \\
\phantomsection\label{entry:symbol-registry-03760}\(\Theta,\mathcal E_X\) & \texttt{\detokenize{\Theta,}\allowbreak\detokenize{\mathcal E_X}} & closed target nonemptiness theory and target parameter-guard assumptions & P8 derived negative-free target language & n/a & \hyperref[def:p08-negative-translation]{\ref*{def:p08-negative-translation}} & \thepage & Theta is the singleton exists z E(z); E\_X contains one guard for each parameter name. \\
\phantomsection\label{entry:symbol-registry-03761}\(\vdash^{\mathrm{tr}}_{-,X}\) & \texttt{\detokenize{\vdash^{\mathrm{tr}}_{-,X}}} & negative derivability by a finite registered classical target certificate & P8 derived negative-free calculus & n/a & \hyperref[def:p08-negative-calculus]{\ref*{def:p08-negative-calculus}} & \thepage & Native rows are derived interfaces; no semantic-validity oracle is an inference premise. \\
\phantomsection\label{entry:symbol-registry-03762}\(\operatorname{val}_g\) & \texttt{\detokenize{\operatorname{val}_g}} & partial source term evaluation under an outer assignment & P8 derived negative-free calculus & n/a & \hyperref[thm:p08-negative-substitution]{\ref*{thm:p08-negative-substitution}} & \thepage & Equality of partial values includes two undefined results. \\
\phantomsection\label{entry:symbol-registry-03763}\(\dashv\vdash^{\mathrm{tr}}_{-,X}\) & \texttt{\detokenize{\dashv\vdash^{\mathrm{tr}}_{-,X}}} & pair of finite negative derivability implications between displayed single formulas & P8 derived negative-free calculus & n/a & \hyperref[thm:p08-negative-equality-existence]{\ref*{thm:p08-negative-equality-existence}} & \thepage & This is explicit shorthand for two directions; it does not identify a formula with a semantic or target-equality object. \\
\phantomsection\label{entry:symbol-registry-03868}\((G,K,R)\) & \texttt{\detokenize{G,}\allowbreak\detokenize{K,}\allowbreak\detokenize{R}} & Considerations pointed model structure with distinguished world and reflexive accessibility & P8 Considerations supporting comparison & object & \hyperref[def:p08-kc-model]{\ref*{def:p08-kc-model}} & \thepage & K here is a world set not the later normal-logic system name; G belongs to nonempty K. \\
\phantomsection\label{entry:symbol-registry-03869}\((M,S4,S5)\) & \texttt{\detokenize{M,}\allowbreak\detokenize{S4,}\allowbreak\detokenize{S5}} & source modal-base system names in the reprint & P8 Considerations supporting comparison & object & \hyperref[def:p08-kc-model]{\ref*{def:p08-kc-model}} & \thepage & Source M is reflexive normal T and not Chellas's least monotonic M or Pacuit EM; the symmetric extension is called the Brouwersche system in the source. \\
\phantomsection\label{entry:symbol-registry-03870}\(\psi(H),U\) & \texttt{\detokenize{\psi(H),}\allowbreak\detokenize{U}} & local existing-individual domain and union of all local domains & P8 Considerations supporting comparison & object & \hyperref[def:p08-kc-model]{\ref*{def:p08-kc-model}} & \thepage & Local domains may be empty and have no implied inclusion along edges; U is the source's chosen assignment space. \\
\phantomsection\label{entry:symbol-registry-03871}\(\varphi(P^n,H),\varphi_s(A,H)\) & \texttt{\detokenize{\varphi(P^n,}\allowbreak\detokenize{H),}\allowbreak\detokenize{\varphi_s(A,}\allowbreak\detokenize{H)}} & predicate extension and formula truth under an explicit free-variable assignment & P8 Considerations supporting comparison & object & \hyperref[def:p08-kc-model]{\ref*{def:p08-kc-model}} & \thepage & Positive arity extensions lie in U to the n even for absent arguments; zero arity is a truth value. \\
\phantomsection\label{entry:symbol-registry-03872}\((x),\exists x\) & \texttt{\detokenize{(x),}\allowbreak\detokenize{\exists x}} & source local-domain universal binder and its Boolean-dual existential & P8 Considerations supporting comparison & object & \hyperref[def:p08-kc-model]{\ref*{def:p08-kc-model}} & \thepage & The source binder ranges over psi(H); it is translated to Fitting's actualist quantifier rather than his primitive possibilist quantifier. \\
\phantomsection\label{entry:symbol-registry-03873}\(s:\operatorname{FV}(A)\to U\) & \texttt{\detokenize{s:\operatorname{FV}(A)\to U}} & explicit source free-variable assignment parameter & P8 Considerations supporting comparison & object & \hyperref[def:p08-kc-model]{\ref*{def:p08-kc-model}} & \thepage & Assignments stay fixed through modal movement; a closed formula has an empty assignment even if U is empty. \\
\phantomsection\label{entry:symbol-registry-03874}\((\mathrm{BF},\mathrm{CBF})\) & \texttt{\detokenize{\mathrm{BF},}\allowbreak\detokenize{\mathrm{CBF}}} & Barcan and converse Barcan directions under the source binder & P8 Considerations supporting comparison & object & \hyperref[def:p08-kc-barcan]{\ref*{def:p08-kc-barcan}} & \thepage & BF is local universal box implies box local universal; CBF is its reverse. \\
\phantomsection\label{entry:symbol-registry-03875}\(E(x),x=y\) & \texttt{\detokenize{E(x),}\allowbreak\detokenize{x=y}} & source optional existence predicate and U-value identity extension & P8 Considerations supporting comparison & relation & \hyperref[def:p08-kc-existence]{\ref*{def:p08-kc-existence}} & \thepage & Existence is local-domain membership; identity compares assigned U-values even when absent locally. \\
\phantomsection\label{entry:symbol-registry-03876}\(t(A)\) & \texttt{\detokenize{t(A)}} & editorial translation replacing the source binder by Fitting's E-guarded quantifier & P8 Considerations supporting comparison & object & \hyperref[thm:p08-kc-fitting-bridge]{\ref*{thm:p08-kc-fitting-bridge}} & \thepage & The translation requires nonempty local domains to remain inside the currently declared Fitting consumer class. \\
\phantomsection\label{entry:symbol-registry-03877}\((c,d)\) & \texttt{\detokenize{c,}\allowbreak\detokenize{d}} & explicitly added later-model individual names in the designation comparison & P8 Considerations designation extension & object & \hyperref[thm:p08-kc-designation]{\ref*{thm:p08-kc-designation}} & \thepage & The supporting source syntax is variable-only; this separate extension distinguishes fixed designation from existence and nonrigid name substitution. \\
\phantomsection\label{entry:symbol-registry-04036}\(M=(W,(\sim_i)_{i\in I},V)\) & \texttt{\detokenize{M}} & finite epistemic model & p08 rationality-2-6 & n/a & \hyperref[def:p08-rat26-epistemic-update]{\ref*{def:p08-rat26-epistemic-update}} & \thepage & Equivalence accessibility; restricted valuation after truthful announcement \\
\phantomsection\label{entry:symbol-registry-04037}\(K_i\) & \texttt{\detokenize{K_i}} & knowledge operator & p08 rationality-2-6 & n/a & \hyperref[def:p08-rat26-epistemic-update]{\ref*{def:p08-rat26-epistemic-update}} & \thepage & All accessible worlds; existing P8 convention \\
\phantomsection\label{entry:symbol-registry-04038}\(C\) & \texttt{\detokenize{C}} & common knowledge operator & p08 epistemic update & n/a & \hyperref[def:p08-rat26-epistemic-update]{\ref*{def:p08-rat26-epistemic-update}} & \thepage & All finite accessibility paths; not injury cost in later payoff section \\
\phantomsection\label{entry:symbol-registry-04039}\([!\phi]\) & \texttt{\detokenize{[!\phi]}} & truthful public announcement box & p08 rationality-2-6 & n/a & \hyperref[def:p08-rat26-epistemic-update]{\ref*{def:p08-rat26-epistemic-update}} & \thepage & Vacuous at worlds failing the precondition \\
\phantomsection\label{entry:symbol-registry-04040}\(\preceq\) & \texttt{\detokenize{\preceq}} & total plausibility preorder & p08 soft information & n/a & \hyperref[def:p08-rat26-soft-update]{\ref*{def:p08-rat26-soft-update}} & \thepage & Earlier is more plausible; finite domain \\
\phantomsection\label{entry:symbol-registry-04041}\(B\) & \texttt{\detokenize{B}} & belief operator & p08 soft information & n/a & \hyperref[def:p08-rat26-soft-update]{\ref*{def:p08-rat26-soft-update}} & \thepage & Truth at best worlds; not necessarily factive \\
\phantomsection\label{entry:symbol-registry-04042}\(B^\chi\) & \texttt{\detokenize{B^\chi}} & conditional belief & p08 soft information & n/a & \hyperref[def:p08-rat26-soft-update]{\ref*{def:p08-rat26-soft-update}} & \thepage & Best chi worlds; empty condition vacuous \\
\phantomsection\label{entry:symbol-registry-04043}\(\uparrow\phi\) & \texttt{\detokenize{\uparrow\phi}} & conservative upgrade & p08 soft information & n/a & \hyperref[def:p08-rat26-soft-update]{\ref*{def:p08-rat26-soft-update}} & \thepage & Promote best phi worlds only \\
\phantomsection\label{entry:symbol-registry-04044}\(\Uparrow\phi\) & \texttt{\detokenize{\Uparrow\phi}} & radical upgrade & p08 soft information & n/a & \hyperref[def:p08-rat26-soft-update]{\ref*{def:p08-rat26-soft-update}} & \thepage & All phi worlds before complement; preserve within-class ranks \\
\phantomsection\label{entry:symbol-registry-04045}\(V,C\) & \texttt{\detokenize{V,}\allowbreak\detokenize{C}} & resource value and injury cost & p08 Hawk-Dove & n/a & \hyperref[def:p08-rat26-hawk-dove]{\ref*{def:p08-rat26-hawk-dove}} & \thepage & Positive payoff parameters; local scope distinct from valuation and common knowledge \\
\phantomsection\label{entry:symbol-registry-04046}\(u(x,z)\) & \texttt{\detokenize{u}} & expected payoff of random pairing & p08 Hawk-Dove & n/a & \hyperref[def:p08-rat26-hawk-dove]{\ref*{def:p08-rat26-hawk-dove}} & \thepage & Row Hawk probability x against Hawk frequency z \\
\phantomsection\label{entry:symbol-registry-04047}\(N(s)\) & \texttt{\detokenize{N}} & outgoing consultation neighborhood & p08 influence dynamics & n/a & \hyperref[def:p08-rat26-influence]{\ref*{def:p08-rat26-influence}} & \thepage & Arrow from consulting agent to consulted agent \\
\phantomsection\label{entry:symbol-registry-04048}\(T(X)\) & \texttt{\detokenize{T}} & synchronous influence update & p08 influence dynamics & n/a & \hyperref[def:p08-rat26-influence]{\ref*{def:p08-rat26-influence}} & \thepage & All consultation neighbors must be active \\
\phantomsection\label{entry:symbol-registry-04049}\(\eta\) & \texttt{\detokenize{\eta}} & positive fixed gradient step & p08 learning countercase & n/a & \hyperref[thm:p08-rat26-learning-limits]{\ref*{thm:p08-rat26-learning-limits}} & \thepage & Not a graph multiplicity encoding \\
\phantomsection\label{entry:symbol-registry-05887}\(+\) & \texttt{\detokenize{+}} & Lattice join & P8 1946 Brouwerian algebra & n/a & \hyperref[def:p08-mt-brouwerian]{\ref*{def:p08-mt-brouwerian}} & \thepage & Source lattice operation; dual to modern Heyting meet Blank macro filled with exact unchanged glyph spelling; registry data only, not a new TeX command. Original precedence 'operator' retained as semantic descriptor; n/a is not an invented rank. \\
\phantomsection\label{entry:symbol-registry-05888}\(\cdot\) & \texttt{\detokenize{\cdot}} & Lattice meet & P8 1946 Brouwerian algebra & n/a & \hyperref[def:p08-mt-brouwerian]{\ref*{def:p08-mt-brouwerian}} & \thepage & Source lattice operation; dual to modern Heyting join Blank macro filled with exact unchanged glyph spelling; registry data only, not a new TeX command. Original precedence 'operator' retained as semantic descriptor; n/a is not an invented rank. \\
\phantomsection\label{entry:symbol-registry-05889}\(\mathbin{\dot{-}}\) & \texttt{\detokenize{\mathbin{\dot{-}}}} & Brouwerian difference & P8 1946 closed-element convention & n/a & \hyperref[def:p08-mt-brouwerian]{\ref*{def:p08-mt-brouwerian}} & \thepage & Typed TeX transcription of binary source difference; least c with a less than or equal to b+c; not unreversed Heyting implication Blank macro filled with exact unchanged glyph spelling; registry data only, not a new TeX command. Original precedence 'operator' retained as semantic descriptor; n/a is not an invented rank. \\
\phantomsection\label{entry:symbol-registry-05890}\(\mathord{\sim}a\) & \texttt{\detokenize{\mathord{\sim}a}} & Brouwerian complement & P8 1946 Brouwerian algebra & n/a & \hyperref[def:p08-mt-brouwerian]{\ref*{def:p08-mt-brouwerian}} & \thepage & 1 difference a; not Boolean complement; exact printed glyph requires retained-image collation Blank macro filled with exact unchanged glyph spelling; registry data only, not a new TeX command. Original precedence 'unary' retained as semantic descriptor; n/a is not an invented rank. \\
\phantomsection\label{entry:symbol-registry-05891}\(-a\) & \texttt{\detokenize{-a}} & Boolean complement & P8 closure algebra & n/a & \hyperref[def:p08-mt-closure]{\ref*{def:p08-mt-closure}} & \thepage & Relative complement in Boolean carrier; distinct from Brouwerian complement Blank macro filled with exact unchanged glyph spelling; registry data only, not a new TeX command. Original precedence 'unary' retained as semantic descriptor; n/a is not an invented rank. \\
\phantomsection\label{entry:symbol-registry-05892}\(C\) & \texttt{\detokenize{C}} & Closure operator & P8 Boolean closure algebra & n/a & \hyperref[def:p08-mt-closure]{\ref*{def:p08-mt-closure}} & \thepage & Extensive idempotent finite-join-preserving operation; actual topological closure in powerset examples Blank macro filled with exact unchanged glyph spelling; registry data only, not a new TeX command. Original precedence 'unary' retained as semantic descriptor; n/a is not an invented rank. \\
\phantomsection\label{entry:symbol-registry-05893}\(I\) & \texttt{\detokenize{I}} & Interior operator & P8 Boolean closure algebra & n/a & \hyperref[def:p08-mt-closure]{\ref*{def:p08-mt-closure}} & \thepage & I x equals minus C minus x Blank macro filled with exact unchanged glyph spelling; registry data only, not a new TeX command. Original precedence 'unary' retained as semantic descriptor; n/a is not an invented rank. \\
\phantomsection\label{entry:symbol-registry-05894}\(0_A\) & \texttt{\detokenize{0_A}} & Least closed-convention bound & P8 Brouwerian versus Heyting duality & n/a & \hyperref[thm:p08-mt-dual]{\ref*{thm:p08-mt-dual}} & \thepage & Equals modern Heyting top; designated closed-convention formula value Root source-label reconciliation lem:p08-mt-dual -> thm:p08-mt-dual; original intake preserved. Blank macro filled with exact unchanged glyph spelling; registry data only, not a new TeX command. Original precedence 'constant' retained as semantic descriptor; n/a is not an invented rank. \\
\phantomsection\label{entry:symbol-registry-05895}\(1_A\) & \texttt{\detokenize{1_A}} & Greatest closed-convention bound & P8 Brouwerian versus Heyting duality & n/a & \hyperref[thm:p08-mt-dual]{\ref*{thm:p08-mt-dual}} & \thepage & Equals modern Heyting bottom; do not confuse with designated formula truth Root source-label reconciliation lem:p08-mt-dual -> thm:p08-mt-dual; original intake preserved. Blank macro filled with exact unchanged glyph spelling; registry data only, not a new TeX command. Original precedence 'constant' retained as semantic descriptor; n/a is not an invented rank. \\
\phantomsection\label{entry:symbol-registry-05896}\(\Rightarrow_H\) & \texttt{\detokenize{\Rightarrow_H}} & Heyting residual & P8 modern dual interpretation & n/a & \hyperref[thm:p08-mt-dual]{\ref*{thm:p08-mt-dual}} & \thepage & b implies a is the value of source a difference b Root source-label reconciliation lem:p08-mt-dual -> thm:p08-mt-dual; original intake preserved. Blank macro filled with exact unchanged glyph spelling; registry data only, not a new TeX command. Original precedence 'operator' retained as semantic descriptor; n/a is not an invented rank. \\
\phantomsection\label{entry:symbol-registry-05897}\(K\) & \texttt{\detokenize{K}} & Boolean-generated extension carrier & P8 minimal closure extension & n/a & \hyperref[thm:p08-mt-extension]{\ref*{thm:p08-mt-extension}} & \thepage & Finite normal-form Boolean algebra; separate typed role from the real Cantor set below Blank macro filled with exact unchanged glyph spelling; registry data only, not a new TeX command. Original precedence 'set' retained as semantic descriptor; n/a is not an invented rank. \\
\phantomsection\label{entry:symbol-registry-05898}\(A_a\) & \texttt{\detokenize{A_a}} & Relativized Brouwerian algebra & P8 source relativization & n/a & \hyperref[def:p08-mt-connected]{\ref*{def:p08-mt-connected}} & \thepage & Carrier x greater than or equal to a; zero a; difference a+(x difference y) Blank macro filled with exact unchanged glyph spelling; registry data only, not a new TeX command. Original precedence 'constructor' retained as semantic descriptor; n/a is not an invented rank. \\
\phantomsection\label{entry:symbol-registry-05899}\(J(L)\) & \texttt{\detokenize{J(L)}} & Nonzero join-irreducibles & P8 finite representation & n/a & \hyperref[def:p08-mt-finite-space]{\ref*{def:p08-mt-finite-space}} & \thepage & Finite ordered carrier; empty for the trivial lattice Blank macro filled with exact unchanged glyph spelling; registry data only, not a new TeX command. Original precedence 'set' retained as semantic descriptor; n/a is not an invented rank. \\
\phantomsection\label{entry:symbol-registry-05900}\(e(a)\) & \texttt{\detokenize{e(a)}} & Closed-set lattice representation & P8 finite representation & n/a & \hyperref[def:p08-mt-finite-space]{\ref*{def:p08-mt-finite-space}} & \thepage & Join-irreducibles below a; all downsets represented Blank macro filled with exact unchanged glyph spelling; registry data only, not a new TeX command. Original precedence 'function' retained as semantic descriptor; n/a is not an invented rank. \\
\phantomsection\label{entry:symbol-registry-05901}\(\downarrow S\) & \texttt{\detokenize{\downarrow S}} & Downward closure & P8 upper Alexandrov topology & n/a & \hyperref[def:p08-mt-finite-space]{\ref*{def:p08-mt-finite-space}} & \thepage & Actual closure of S; contains all lower points of members of S Blank macro filled with exact unchanged glyph spelling; registry data only, not a new TeX command. Original precedence 'operator' retained as semantic descriptor; n/a is not an invented rank. \\
\phantomsection\label{entry:symbol-registry-05902}\(\uparrow w\) & \texttt{\detokenize{\uparrow w}} & Principal upset & P8 upper Alexandrov topology & n/a & \hyperref[thm:p08-mt-formulas]{\ref*{thm:p08-mt-formulas}} & \thepage & Least open neighborhood; implication quantifies over its points Root source-label reconciliation lem:p08-mt-formulas -> thm:p08-mt-formulas; original intake preserved. Blank macro filled with exact unchanged glyph spelling; registry data only, not a new TeX command. Original precedence 'operator' retained as semantic descriptor; n/a is not an invented rank. \\
\phantomsection\label{entry:symbol-registry-05903}\(t_B\) & \texttt{\detokenize{t_B}} & Closed-convention formula term & P8 formula bridge & n/a & \hyperref[def:p08-mt-formula-term]{\ref*{def:p08-mt-formula-term}} & \thepage & Conjunction becomes join; disjunction becomes meet; implication reverses difference operands Blank macro filled with exact unchanged glyph spelling; registry data only, not a new TeX command. Original precedence 'term' retained as semantic descriptor; n/a is not an invented rank. \\
\phantomsection\label{entry:symbol-registry-05904}\(K\subseteq[0,1]\) & \texttt{\detokenize{K\subseteq[0,}\allowbreak\detokenize{1]}} & Middle-thirds Cantor set & P8 real-line reconstruction & n/a & \hyperref[thm:p08-mt-cantor]{\ref*{thm:p08-mt-cantor}} & \thepage & Typed subset of reals; not the Boolean extension carrier; generation coloring local at every point Root source-label reconciliation lem:p08-mt-cantor -> thm:p08-mt-cantor; original intake preserved. Blank macro filled with exact unchanged glyph spelling; registry data only, not a new TeX command. Original precedence 'set' retained as semantic descriptor; n/a is not an invented rank. \\
\phantomsection\label{entry:symbol-registry-05905}\(f:X\to Y\) & \texttt{\detokenize{f:X\to Y}} & Interior map & P8 topological formula preservation & n/a & \hyperref[thm:p08-mt-pullback]{\ref*{thm:p08-mt-pullback}} & \thepage & Continuous and open; surjectivity additionally required for validity reflection Root source-label reconciliation lem:p08-mt-pullback -> thm:p08-mt-pullback; original intake preserved. Blank macro filled with exact unchanged glyph spelling; registry data only, not a new TeX command. Original precedence 'function' retained as semantic descriptor; n/a is not an invented rank. \\
\phantomsection\label{entry:symbol-registry-05906}\(\mathbb R^d\) & \texttt{\detokenize{\mathbb R^d}} & Positive-dimensional Euclidean space & P8 reconstructed completeness & n/a & \hyperref[thm:p08-mt-topological-completeness]{\ref*{thm:p08-mt-topological-completeness}} & \thepage & d is a positive integer; singleton R\^{}0 excluded by concrete classical countercase Blank macro filled with exact unchanged glyph spelling; registry data only, not a new TeX command. Original precedence 'space' retained as semantic descriptor; n/a is not an invented rank. \\
\phantomsection\label{entry:symbol-registry-05907}\(f_A:L^n\to L\) & \texttt{\detokenize{f_A:L^n\to L}} & Uniform algebraic function at one algebra & P8 1946 algebraic-function presentation & n/a & \hyperref[def:p08-mt-source-functions]{\ref*{def:p08-mt-source-functions}} & \thepage & The family ranges over every Brouwerian algebra; equality on one A is not global equality of uniform functions Blank macro filled with exact unchanged glyph spelling; registry data only, not a new TeX command. Original precedence 'function' retained as semantic descriptor; n/a is not an invented rank. \\
\phantomsection\label{entry:symbol-registry-07264}\(L\) & \texttt{\detokenize{L}} & Set of finitary classical formulas & P8 KLM source framework & n/a & \hyperref[def:p08-klm-language]{\ref*{def:p08-klm-language}} & \thepage & Not a language with nested defeasible assertions Original precedence descriptor 'set' retained; no rank inferred. Blank macro filled with the glyph's TeX spelling. \\
\phantomsection\label{entry:symbol-registry-07265}\(U\) & \texttt{\detokenize{U}} & Fixed universe of classical worlds & P8 KLM source framework & n/a & \hyperref[def:p08-klm-language]{\ref*{def:p08-klm-language}} & \thepage & May be background-theory restricted; compactness explicitly assumed Original precedence descriptor 'set' retained; no rank inferred. Blank macro filled with the glyph's TeX spelling. \\
\phantomsection\label{entry:symbol-registry-07266}\(m\models\alpha\) & \texttt{\detokenize{m\models\alpha}} & Classical truth at a world & P8 KLM source worlds & relation & \hyperref[def:p08-klm-language]{\ref*{def:p08-klm-language}} & \thepage & Ordinary Boolean connectives; distinct from preferential assertion Blank macro filled with the glyph's TeX spelling. \\
\phantomsection\label{entry:symbol-registry-07267}\(\models\alpha\) & \texttt{\detokenize{\models\alpha}} & Validity relative to fixed U & P8 KLM source classical guards & relation & \hyperref[def:p08-klm-language]{\ref*{def:p08-klm-language}} & \thepage & Not necessarily substitution-invariant under background constraints Blank macro filled with the glyph's TeX spelling. \\
\phantomsection\label{entry:symbol-registry-07268}\(\alpha\models\beta\) & \texttt{\detokenize{\alpha\models\beta}} & Material semantic consequence relative to U & P8 KLM source classical guards & relation & \hyperref[def:p08-klm-language]{\ref*{def:p08-klm-language}} & \thepage & Abbreviates validity of alpha material-implies beta Blank macro filled with the glyph's TeX spelling. \\
\phantomsection\label{entry:symbol-registry-07269}\(\alpha\sim\beta\) & \texttt{\detokenize{\alpha\sim\beta}} & Flat conditional assertion and abstract relation & P8 KLM C/P syntax & relation & \hyperref[def:p08-klm-assertions]{\ref*{def:p08-klm-assertions}} & \thepage & Published source notation; a pair of formulas in the metalanguage, not a nestable connective Blank macro filled with the glyph's TeX spelling. \\
\phantomsection\label{entry:symbol-registry-07270}\(K\vdash_P\alpha\sim\beta\) & \texttt{\detokenize{K\vdash_P\alpha\sim\beta}} & Finite P derivability from an assertion base & P8 KLM calculus & n/a & \hyperref[def:p08-klm-assertions]{\ref*{def:p08-klm-assertions}} & \thepage & K may be infinite; classical-validity premises are external fixed guards Original precedence descriptor 'judgment' retained; no rank inferred. Blank macro filled with the glyph's TeX spelling. \\
\phantomsection\label{entry:symbol-registry-07271}\(S\) & \texttt{\detokenize{S}} & Set of states & P8 preferential models & n/a & \hyperref[def:p08-klm-model]{\ref*{def:p08-klm-model}} & \thepage & Copies permitted; can be empty; not required to equal the world universe Original precedence descriptor 'set' retained; no rank inferred. Blank macro filled with the glyph's TeX spelling. \\
\phantomsection\label{entry:symbol-registry-07272}\(l:S\to U\) & \texttt{\detokenize{l:S\to U}} & Singleton-world preferential label map & P8 preferential model type & n/a & \hyperref[def:p08-klm-model]{\ref*{def:p08-klm-model}} & \thepage & Need not be injective or onto Original precedence descriptor 'function' retained; no rank inferred. Blank macro filled with the glyph's TeX spelling. \\
\phantomsection\label{entry:symbol-registry-07273}\(\prec\) & \texttt{\detokenize{\prec}} & Strict preference on states & P8 preferential model type & relation & \hyperref[def:p08-klm-model]{\ref*{def:p08-klm-model}} & \thepage & Preferred state is on the left; irreflexive and transitive; not the formula preorder Blank macro filled with the glyph's TeX spelling. \\
\phantomsection\label{entry:symbol-registry-07274}\(\widehat\alpha\) & \texttt{\detokenize{\widehat\alpha}} & Formula-definable state set & P8 preferential semantics & n/a & \hyperref[def:p08-klm-model]{\ref*{def:p08-klm-model}} & \thepage & Label-world satisfies alpha; smoothness required on these sets Original precedence descriptor 'set' retained; no rank inferred. Blank macro filled with the glyph's TeX spelling. \\
\phantomsection\label{entry:symbol-registry-07275}\(\mu(X)\) & \texttt{\detokenize{\mu(X)}} & Minimal states of X & P8 preferential semantics & n/a & \hyperref[def:p08-klm-model]{\ref*{def:p08-klm-model}} & \thepage & All states with no predecessor in X; not one selected minimum or a probability mass Original precedence descriptor 'operator' retained; no rank inferred. Blank macro filled with the glyph's TeX spelling. \\
\phantomsection\label{entry:symbol-registry-07276}\(\sim_W\) & \texttt{\detokenize{\sim_W}} & Model-defined preferential consequence & P8 preferential semantics & relation & \hyperref[def:p08-klm-model]{\ref*{def:p08-klm-model}} & \thepage & All minimal antecedent states satisfy consequent; linked to abstract relation by representation Blank macro filled with the glyph's TeX spelling. \\
\phantomsection\label{entry:symbol-registry-07277}\(C(\alpha)\) & \texttt{\detokenize{C(\alpha)}} & Full conditional consequence theory & P8 normal-world witness & n/a & \hyperref[def:p08-klm-normal]{\ref*{def:p08-klm-normal}} & \thepage & Can be infinite; finite satisfiability and compactness are actually proved/used Original precedence descriptor 'set' retained; no rank inferred. Blank macro filled with the glyph's TeX spelling. \\
\phantomsection\label{entry:symbol-registry-07278}\(N(\alpha)\) & \texttt{\detokenize{N(\alpha)}} & All alpha-normal worlds & P8 normal-world witness & n/a & \hyperref[def:p08-klm-normal]{\ref*{def:p08-klm-normal}} & \thepage & Definition does not presume existence; witnesses supplied by proved compactness lemma Original precedence descriptor 'set' retained; no rank inferred. Blank macro filled with the glyph's TeX spelling. \\
\phantomsection\label{entry:symbol-registry-07279}\(\alpha\leq\beta\) & \texttt{\detokenize{\alpha\leq\beta}} & Not-less-ordinary formula preorder & P8 KLM source representation & relation & \hyperref[def:p08-klm-ordinary]{\ref*{def:p08-klm-ordinary}} & \thepage & Defined by alpha disjunction beta conditionally entails alpha; not material implication or state preference Blank macro filled with the glyph's TeX spelling. \\
\phantomsection\label{entry:symbol-registry-07280}\((m,\alpha)\) & \texttt{\detokenize{(m,}\allowbreak\detokenize{\alpha)}} & Canonical copied state & P8 canonical preferential representation & n/a & \hyperref[def:p08-klm-canonical]{\ref*{def:p08-klm-canonical}} & \thepage & m is alpha-normal; formula coordinate must not be erased Original precedence descriptor 'state' retained; no rank inferred. Blank macro filled with the glyph's TeX spelling. \\
\phantomsection\label{entry:symbol-registry-07281}\(l:S\to\mathcal P(U)\setminus\{\varnothing\}\) & \texttt{\detokenize{l:S\to\mathcal P(U)\setminus\{\varnothing\}}} & Set-of-worlds cumulative label map & P8 local C-versus-P comparison & n/a & \hyperref[thm:p08-klm-cumulative-separation]{\ref*{thm:p08-klm-cumulative-separation}} & \thepage & Different genuine source model type from singleton-world preferential labels Original precedence descriptor 'function' retained; no rank inferred. Blank macro filled with the glyph's TeX spelling. \\
\phantomsection\label{entry:symbol-registry-07282}\(s\equiv\alpha\) & \texttt{\detokenize{s\equiv\alpha}} & Cumulative state satisfaction & P8 local C-versus-P comparison & relation & \hyperref[thm:p08-klm-cumulative-separation]{\ref*{thm:p08-klm-cumulative-separation}} & \thepage & Every label-world satisfies alpha; not identity-type or logical equivalence of formulas Blank macro filled with the glyph's TeX spelling. \\
\phantomsection\label{entry:symbol-registry-07669}\(\mathcal L\) & \texttt{\detokenize{\mathcal L}} & Finitary classical propositional language & P8 Lehmann Magidor source language & n/a & \hyperref[def:p08-rational-language]{\ref*{def:p08-rational-language}} & \thepage & Full propositional language; logically finite means finitely many atoms, not finitely many syntactic formulas Original precedence descriptor 'set' retained; no rank inferred. Blank macro filled with the glyph's TeX spelling. \\
\phantomsection\label{entry:symbol-registry-07670}\(\mathcal U\) & \texttt{\detokenize{\mathcal U}} & All Boolean valuations of the atom set & P8 Lehmann Magidor source language & n/a & \hyperref[def:p08-rational-language]{\ref*{def:p08-rational-language}} & \thepage & Source1992 full-universe presentation; distinct from background-restricted KLM U Original precedence descriptor 'set' retained; no rank inferred. Blank macro filled with the glyph's TeX spelling. \\
\phantomsection\label{entry:symbol-registry-07671}\(\alpha\mathrel{|\!\sim}\beta\) & \texttt{\detokenize{\alpha\mathrel{|\!\sim}\beta}} & Flat conditional assertion and abstract rational relation & P8 Lehmann Magidor assertions & relation & \hyperref[def:p08-rational-language]{\ref*{def:p08-rational-language}} & \thepage & Source hooked-bar glyph; distinct from the unbarred quotient equivalence alpha sim beta Blank macro filled with the glyph's TeX spelling. \\
\phantomsection\label{entry:symbol-registry-07672}\(\alpha\mathrel{|\!\not\sim}\beta\) & \texttt{\detokenize{\alpha\mathrel{|\!\not\sim}\beta}} & Actual nonmembership of an assertion & P8 Lehmann Magidor assertions & relation & \hyperref[def:p08-rational-language]{\ref*{def:p08-rational-language}} & \thepage & Not object-language negation or unsuccessful proof search Blank macro filled with the glyph's TeX spelling. \\
\phantomsection\label{entry:symbol-registry-07673}\(V\) & \texttt{\detokenize{V}} & Set of ranked-model states & P8 Lehmann Magidor ranked models & n/a & \hyperref[def:p08-rational-ranked]{\ref*{def:p08-rational-ranked}} & \thepage & Copies permitted; not the consistent-formula set S Original precedence descriptor 'set' retained; no rank inferred. Blank macro filled with the glyph's TeX spelling. \\
\phantomsection\label{entry:symbol-registry-07674}\(l\) & \texttt{\detokenize{l}} & State-to-world label map & P8 Lehmann Magidor ranked models & n/a & \hyperref[def:p08-rational-ranked]{\ref*{def:p08-rational-ranked}} & \thepage & Domain V and codomain mathcal U; need not be injective or onto Original precedence descriptor 'function' retained; no rank inferred. Blank macro filled with the glyph's TeX spelling. \\
\phantomsection\label{entry:symbol-registry-07675}\(\prec\) & \texttt{\detokenize{\prec}} & Strict modular preference between states & P8 Lehmann Magidor ranked models & relation & \hyperref[def:p08-rational-ranked]{\ref*{def:p08-rational-ranked}} & \thepage & Lower ranked state on left; distinct typed role from preference between whole rational relations Blank macro filled with the glyph's TeX spelling. \\
\phantomsection\label{entry:symbol-registry-07676}\(\Omega\) & \texttt{\detokenize{\Omega}} & Totally ordered rank codomain & P8 Lehmann Magidor ranked models & n/a & \hyperref[def:p08-rational-ranked]{\ref*{def:p08-rational-ranked}} & \thepage & Need not be well ordered or ordinal Original precedence descriptor 'set' retained; no rank inferred. Blank macro filled with the glyph's TeX spelling. \\
\phantomsection\label{entry:symbol-registry-07677}\(r\) & \texttt{\detokenize{r}} & State ranking map into Omega & P8 Lehmann Magidor ranked models & n/a & \hyperref[def:p08-rational-ranked]{\ref*{def:p08-rational-ranked}} & \thepage & Strict rank comparison induces preference Original precedence descriptor 'function' retained; no rank inferred. Blank macro filled with the glyph's TeX spelling. \\
\phantomsection\label{entry:symbol-registry-07678}\(s\equiv\alpha\) & \texttt{\detokenize{s\equiv\alpha}} & Ranked state satisfies formula through its label & P8 Lehmann Magidor ranked models & relation & \hyperref[def:p08-rational-ranked]{\ref*{def:p08-rational-ranked}} & \thepage & Source state/formula satisfaction; not logical equivalence or identity Blank macro filled with the glyph's TeX spelling. \\
\phantomsection\label{entry:symbol-registry-07679}\(\widehat\alpha\) & \texttt{\detokenize{\widehat\alpha}} & Formula-definable ranked-state set & P8 Lehmann Magidor ranked models & n/a & \hyperref[def:p08-rational-ranked]{\ref*{def:p08-rational-ranked}} & \thepage & Smoothness required on each such set Original precedence descriptor 'set' retained; no rank inferred. Blank macro filled with the glyph's TeX spelling. \\
\phantomsection\label{entry:symbol-registry-07680}\(\mathrel{|\!\sim}_W\) & \texttt{\detokenize{\mathrel{|\!\sim}_W}} & Conditional consequence defined by ranked minima & P8 Lehmann Magidor ranked models & relation & \hyperref[def:p08-rational-ranked]{\ref*{def:p08-rational-ranked}} & \thepage & All minimal antecedent states; vacuous for empty antecedent Blank macro filled with the glyph's TeX spelling. \\
\phantomsection\label{entry:symbol-registry-07681}\(S\) & \texttt{\detokenize{S}} & Set of formulas consistent for the given rational relation & P8 rational formula quotient & n/a & \hyperref[def:p08-rational-normal]{\ref*{def:p08-rational-normal}} & \thepage & Native source local type differs from KLM state-set S; consistency means no assertion to bottom Original precedence descriptor 'set' retained; no rank inferred. Blank macro filled with the glyph's TeX spelling. \\
\phantomsection\label{entry:symbol-registry-07682}\(\mathcal R\) & \texttt{\detokenize{\mathcal R}} & Not-more-exceptional comparison on consistent formulas & P8 rational formula quotient & relation & \hyperref[def:p08-rational-normal]{\ref*{def:p08-rational-normal}} & \thepage & alpha disjunction beta does not conditionally entail not alpha; total preorder actually proved Blank macro filled with the glyph's TeX spelling. \\
\phantomsection\label{entry:symbol-registry-07683}\(\alpha\sim\beta\) & \texttt{\detokenize{\alpha\sim\beta}} & Equivalence of exception classes of consistent formulas & P8 rational formula quotient & relation & \hyperref[def:p08-rational-normal]{\ref*{def:p08-rational-normal}} & \thepage & Published native unbarred glyph; not the conditional assertion with hooked bar Blank macro filled with the glyph's TeX spelling. \\
\phantomsection\label{entry:symbol-registry-07684}\(\overline\alpha\) & \texttt{\detokenize{\overline\alpha}} & Equivalence class of a consistent formula & P8 rational formula quotient & n/a & \hyperref[def:p08-rational-normal]{\ref*{def:p08-rational-normal}} & \thepage & Only defined for alpha in S Original precedence descriptor 'class' retained; no rank inferred. Blank macro filled with the glyph's TeX spelling. \\
\phantomsection\label{entry:symbol-registry-07685}\(E\) & \texttt{\detokenize{E}} & Set of exception-equivalence classes & P8 rational formula quotient & n/a & \hyperref[def:p08-rational-normal]{\ref*{def:p08-rational-normal}} & \thepage & Different source type from exception-removal operator E(K) Original precedence descriptor 'set' retained; no rank inferred. Blank macro filled with the glyph's TeX spelling. \\
\phantomsection\label{entry:symbol-registry-07686}\(\overline\alpha\leq\overline\beta\) & \texttt{\detokenize{\overline\alpha\leq\overline\beta}} & Total class order induced by mathcal R & P8 rational formula quotient & relation & \hyperref[def:p08-rational-normal]{\ref*{def:p08-rational-normal}} & \thepage & Not material implication or the earlier preferential formula preorder Blank macro filled with the glyph's TeX spelling. \\
\phantomsection\label{entry:symbol-registry-07687}\(\overline\alpha<\overline\beta\) & \texttt{\detokenize{\overline\alpha<\overline\beta}} & Strict order on consistent-formula classes & P8 rational formula quotient & relation & \hyperref[def:p08-rational-normal]{\ref*{def:p08-rational-normal}} & \thepage & Distinct from source closure-comparison abbreviation alpha less beta between formulas in a rational relation Blank macro filled with the glyph's TeX spelling. \\
\phantomsection\label{entry:symbol-registry-07688}\(\langle m,\alpha\rangle\) & \texttt{\detokenize{\langle m,}\allowbreak\detokenize{\alpha\rangle}} & Canonical normal-world and formula pair & P8 rational copied state construction & n/a & \hyperref[def:p08-rational-canonical]{\ref*{def:p08-rational-canonical}} & \thepage & m normal for alpha; every such pair retained; rank is overline alpha Original precedence descriptor 'state' retained; no rank inferred. Blank macro filled with the glyph's TeX spelling. \\
\phantomsection\label{entry:symbol-registry-07689}\(K^p\) & \texttt{\detokenize{K^p}} & Preferential closure of a conditional base & P8 rational closure ranks & n/a & \hyperref[def:p08-rational-exception-ranks]{\ref*{def:p08-rational-exception-ranks}} & \thepage & Not rational closure or all-ranked-model consequence identified without proof Original precedence descriptor 'set' retained; no rank inferred. Blank macro filled with the glyph's TeX spelling. \\
\phantomsection\label{entry:symbol-registry-07690}\(E(K)\) & \texttt{\detokenize{E(K)}} & Operator retaining defaults with exceptional antecedents & P8 rational closure ranks & n/a & \hyperref[def:p08-rational-exception-ranks]{\ref*{def:p08-rational-exception-ranks}} & \thepage & Native source overload distinct from set E of formula classes Original precedence descriptor 'operator' retained; no rank inferred. Blank macro filled with the glyph's TeX spelling. \\
\phantomsection\label{entry:symbol-registry-07691}\(C_\tau\) & \texttt{\detokenize{C_\tau}} & Decreasing transfinite sequence of default subbases & P8 rational closure ranks & n/a & \hyperref[def:p08-rational-exception-ranks]{\ref*{def:p08-rational-exception-ranks}} & \thepage & Successor E(C\_tau), limit intersection; finite base stabilizes by n strict removals Original precedence descriptor 'set' retained; no rank inferred. Blank macro filled with the glyph's TeX spelling. \\
\phantomsection\label{entry:symbol-registry-07692}\(\widetilde{\mathcal A}\) & \texttt{\detokenize{\widetilde{\mathcal A}}} & Material counterpart of one flat assertion & P8 rational materialization & n/a & \hyperref[def:p08-rational-exception-ranks]{\ref*{def:p08-rational-exception-ranks}} & \thepage & alpha material-implies beta; not conditional connective identification Original precedence descriptor 'formula' retained; no rank inferred. Blank macro filled with the glyph's TeX spelling. \\
\phantomsection\label{entry:symbol-registry-07693}\(\widetilde K\) & \texttt{\detokenize{\widetilde K}} & Set of material counterparts of a conditional base & P8 rational materialization & n/a & \hyperref[def:p08-rational-exception-ranks]{\ref*{def:p08-rational-exception-ranks}} & \thepage & Classical implication set; exact top-level exception test proved Original precedence descriptor 'set' retained; no rank inferred. Blank macro filled with the glyph's TeX spelling. \\
\phantomsection\label{entry:symbol-registry-07694}\(\alpha<\beta\) & \texttt{\detokenize{\alpha<\beta}} & Conditional assertion abbreviation in closure comparison & P8 rational extension selection & relation & \hyperref[def:p08-rational-abstract-closure]{\ref*{def:p08-rational-abstract-closure}} & \thepage & Native source alpha disjunction beta conditionally entails not beta; distinguish class-order type, especially inconsistent formulas Blank macro filled with the glyph's TeX spelling. \\
\phantomsection\label{entry:symbol-registry-07695}\(K_0\prec K_1\) & \texttt{\detokenize{K_0\prec K_1}} & Preference between whole rational consequence relations & P8 rational extension selection & relation & \hyperref[def:p08-rational-abstract-closure]{\ref*{def:p08-rational-abstract-closure}} & \thepage & Native source typed overload; two quantified conditions; not state preference Blank macro filled with the glyph's TeX spelling. \\
\phantomsection\label{entry:symbol-registry-07696}\(\overline K\) & \texttt{\detokenize{\overline K}} & Rational closure selected among rational extensions & P8 rational extension selection & n/a & \hyperref[def:p08-rational-abstract-closure]{\ref*{def:p08-rational-abstract-closure}} & \thepage & Conditional existence; full admissible-base selection proof supplied separately from finite realization Original precedence descriptor 'set' retained; no rank inferred. Blank macro filled with the glyph's TeX spelling. \\
\phantomsection\label{entry:symbol-registry-07697}\(h(\alpha)\) & \texttt{\detokenize{h(\alpha)}} & Exception rank of a formula for the fixed base & P8 rational exception rank proofs and finite realization & n/a & \hyperref[def:p08-rational-finite-realization]{\ref*{def:p08-rational-finite-realization}} & \thepage & Explicit editorial name for source rank; finite values for finite K and ordinal values in general admissible-base proof; no-rank marker infinity ordered above all occurring ranks Original precedence descriptor 'ordinal' retained; no rank inferred. Blank macro filled with the glyph's TeX spelling. \\
\phantomsection\label{entry:symbol-registry-07698}\(\infty\) & \texttt{\detokenize{\infty}} & No-rank marker for formula exception ranks & P8 finite rational ranked realization & n/a & \hyperref[def:p08-rational-finite-realization]{\ref*{def:p08-rational-finite-realization}} & \thepage & Not an occupied state rank; no-rank antecedents separately vacuous Original precedence descriptor 'marker' retained; no rank inferred. Blank macro filled with the glyph's TeX spelling. \\
\phantomsection\label{entry:symbol-registry-07699}\(\mathcal U_*\) & \texttt{\detokenize{\mathcal U_*}} & Valuations satisfying the stable material subbase & P8 finite rational ranked realization & n/a & \hyperref[def:p08-rational-finite-realization]{\ref*{def:p08-rational-finite-realization}} & \thepage & Exactly possible labels; omitted worlds are not assigned rank infinity Original precedence descriptor 'set' retained; no rank inferred. Blank macro filled with the glyph's TeX spelling. \\
\phantomsection\label{entry:symbol-registry-07700}\(k(m)\) & \texttt{\detokenize{k(m)}} & Least material stage satisfied by a permitted valuation & P8 finite rational ranked realization & n/a & \hyperref[def:p08-rational-finite-realization]{\ref*{def:p08-rational-finite-realization}} & \thepage & Concrete finite layer ranking, not an assumed abstract formula rank Original precedence descriptor 'integer' retained; no rank inferred. Blank macro filled with the glyph's TeX spelling. \\
\phantomsection\label{entry:symbol-registry-07701}\(W_K\) & \texttt{\detokenize{W_K}} & Complete valuation ranked model of the finite base & P8 finite rational ranked realization & n/a & \hyperref[def:p08-rational-finite-realization]{\ref*{def:p08-rational-finite-realization}} & \thepage & One state per permitted world; identity labeling; finite number of layers even over infinite vocabulary Original precedence descriptor 'model' retained; no rank inferred. Blank macro filled with the glyph's TeX spelling. \\
\phantomsection\label{entry:symbol-registry-07702}\(f(m)\) & \texttt{\detokenize{f(m)}} & Normalized least-copy rank in a compared complete finite-vocabulary model & P8 finite rational pointwise comparison & n/a & \hyperref[thm:p08-rational-pointwise-minimum]{\ref*{thm:p08-rational-pointwise-minimum}} & \thepage & Finite occupied layers renumbered from zero; distinct from propositional fly atom f in the penguin example Original precedence descriptor 'integer' retained; no rank inferred. Blank macro filled with the glyph's TeX spelling. \\
\phantomsection\label{entry:symbol-registry-07703}\(x\approx y\) & \texttt{\detokenize{x\approx y}} & Incomparability equivalence in the modular-order proof & P8 modularity quotient proof & relation & \hyperref[annex:p08.rational.modularity]{\ref*{annex:p08.rational.modularity}} & \thepage & Explicitly defined inside the proof; not an object-language connective Blank macro filled with the glyph's TeX spelling. \\
\phantomsection\label{entry:symbol-registry-07704}\(T\) & \texttt{\detokenize{T}} & Rank-criterion relation in the admissible selection proof & P8 admissible rational selection proof & relation & \hyperref[annex:p08.rational.selection]{\ref*{annex:p08.rational.selection}} & \thepage & Local proof binding; not the normal-world theory T of the separate witness proof Blank macro filled with the glyph's TeX spelling. \\
\phantomsection\label{entry:symbol-registry-07705}\(H\) & \texttt{\detokenize{H}} & Ranked formulas with exception rank at least the first difference ordinal & P8 admissible rational selection proof & n/a & \hyperref[annex:p08.rational.selection]{\ref*{annex:p08.rational.selection}} & \thepage & No-rank formulas excluded; their assertions agree by admissibility Original precedence descriptor 'set' retained; no rank inferred. Blank macro filled with the glyph's TeX spelling. \\
\phantomsection\label{entry:symbol-registry-07706}\(\lambda(\varphi)\) & \texttt{\detokenize{\lambda(\varphi)}} & Least layer of a consistent formula in a competing ranked relation & P8 admissible rational selection proof & n/a & \hyperref[annex:p08.rational.selection]{\ref*{annex:p08.rational.selection}} & \thepage & Defined only for formulas with states in R; codomain can be non-well-ordered; no minimum case proved explicitly Original precedence descriptor 'rank' retained; no rank inferred. Blank macro filled with the glyph's TeX spelling. \\
\phantomsection\label{entry:symbol-registry-07707}\(\Gamma\) & \texttt{\detokenize{\Gamma}} & Finite disjunction of persistent default antecedents & P8 finite admissibility proof & n/a & \hyperref[annex:p08.rational.admissibility]{\ref*{annex:p08.rational.admissibility}} & \thepage & Bottom for the empty stable base; smoothness applies to this formula-definable union Original precedence descriptor 'formula' retained; no rank inferred. Blank macro filled with the glyph's TeX spelling. \\
\end{longtable}
\endgroup
\Needspace{.4\textheight}
\subsection*{Stable part P9 scopes}
\begingroup\scriptsize\setlength{\tabcolsep}{3pt}
\begin{longtable}{@{}p{0.9in}p{1.3in}p{1.0in}p{0.5in}p{0.3in}p{0.4in}p{0.35in}p{1.65in}@{}}
\toprule
\textbf{Glyph} & \textbf{Macro} & \textbf{Role} & \textbf{Scope} & \textbf{Prec./form} & \textbf{First} & \textbf{Reg. page} & \textbf{Principal uses / notes} \\
\midrule
\endfirsthead
\toprule
\textbf{Glyph} & \textbf{Macro} & \textbf{Role} & \textbf{Scope} & \textbf{Prec./form} & \textbf{First} & \textbf{Reg. page} & \textbf{Principal uses / notes} \\
\midrule
\endhead
\midrule\multicolumn{8}{r}{\emph{continued}}\\\endfoot
\bottomrule\endlastfoot
\phantomsection\label{entry:symbol-registry-00641}\(\tau_{\mathrm{fin}}\) & \texttt{\detokenize{\tau_{\mathrm{fin}}}} & fixed finite relational signature & P9 logic-in-computer-science & n/a & \hyperref[sec:p09-logic-in-computer-science]{\ref*{sec:p09-logic-in-computer-science}} & \thepage & The signature relative to which finite satisfiability and finite validity are defined; it is not the varying database signature of a combined-complexity input. \\
\phantomsection\label{entry:symbol-registry-00642}\(\mathcal{ALC}\) & \texttt{\detokenize{\mathcal{ALC}}} & description-logic language and system identifier & P9 legacy unit 3 description-logics & n/a & \hyperref[sec:p09-description-logics]{\ref*{sec:p09-description-logics}} & \thepage & The exact Boolean concept language with qualified role restrictions from Definition def:p09-alc; extensions in the description-logic table have different identifiers. \\
\phantomsection\label{entry:symbol-registry-00643}\(I=(\Delta^I,\mathord{\cdot}^I)\) & \texttt{\detokenize{I,}\allowbreak\detokenize{\Delta^I,}\allowbreak\detokenize{\mathord{\cdot}^I}} & description-logic interpretation datum & P9 legacy unit 3 description-logics & n/a & \hyperref[sec:p09-description-logics]{\ref*{sec:p09-description-logics}} & \thepage & Delta-I is the nonempty individual domain and dot-I interprets concept role and individual names with the types fixed in Definition def:p09-alc. \\
\phantomsection\label{entry:symbol-registry-00644}\(\mathcal T,\mathcal A\) & \texttt{\detokenize{\mathcal T,}\allowbreak\detokenize{\mathcal A}} & ALC TBox and ABox & P9 legacy unit 3 description-logics & n/a & \hyperref[sec:p09-description-logics]{\ref*{sec:p09-description-logics}} & \thepage & Finite sets respectively of concept inclusions and concept or role assertions; script A is locally typed apart from automata. \\
\phantomsection\label{entry:symbol-registry-00645}\(C\sqsubseteq D\) & \texttt{\detokenize{\sqsubseteq}} & description-logic concept inclusion & P9 legacy unit 3 description-logics & 6 & \hyperref[sec:p09-description-logics]{\ref*{sec:p09-description-logics}} & \thepage & An object-language TBox axiom interpreted by subset inclusion of concept extensions; not metalinguistic formula consequence. \\
\phantomsection\label{entry:symbol-registry-00646}\(\exists r.C,\forall r.C\) & \texttt{\detokenize{\exists r.,}\allowbreak\detokenize{\forall r.}} & ALC role restrictions & P9 legacy unit 3 description-logics & 3 & \hyperref[sec:p09-description-logics]{\ref*{sec:p09-description-logics}} & \thepage & Qualified existential and universal concept constructors interpreted over the declared role relation. \\
\phantomsection\label{entry:symbol-registry-00647}\(\operatorname{nnf},\operatorname{cl}\) & \texttt{\detokenize{\operatorname{nnf},}\allowbreak\detokenize{\operatorname{cl}}} & ALC normalization and finite closure operations & P9 legacy unit 3 description-logics & meta & \hyperref[sec:p09-description-logics]{\ref*{sec:p09-description-logics}} & \thepage & Metalinguistic operations used by blocking completion and type elimination; they add no concept constructors. \\
\phantomsection\label{entry:symbol-registry-00648}\(\mathsf{Glob}(\mathcal T)\) & \texttt{\detokenize{\mathsf{Glob}}} & ALC tableau global-GCI abbreviation & P9 legacy unit 3 description-logics & meta & \hyperref[sec:p09-description-logics]{\ref*{sec:p09-description-logics}} & \thepage & Metalinguistic abbreviation for the NNF internalizations of every GCI in the active TBox; used in body initialization and on sheet:tableaux-alc. \\
\phantomsection\label{entry:symbol-registry-00649}\(\mathcal L(x),\mathcal L(y)\subseteq\mathcal L(x)\) & \texttt{\detokenize{\mathcal L,}\allowbreak\detokenize{\subseteq}} & ALC completion label and subset-blocking comparison & P9 legacy unit 3 description-logics & meta & \hyperref[sec:p09-description-logics]{\ref*{sec:p09-description-logics}} & \thepage & Labels are subsets of the finite closure; a nonroot y is directly blocked by ancestor x exactly under the displayed subset comparison, re-evaluated after label growth. \\
\phantomsection\label{entry:symbol-registry-00650}\(\widehat y,b(y)\) & \texttt{\detokenize{\widehat y,}\allowbreak\detokenize{b}} & ALC retained-edge target and direct blocker & P9 legacy unit 3 description-logics & meta & \hyperref[sec:p09-description-logics]{\ref*{sec:p09-description-logics}} & \thepage & Defined only for a target y of an edge retained from an unblocked source; widehat-y is y when unblocked and otherwise a directly blocking unblocked ancestor b(y). \\
\phantomsection\label{entry:symbol-registry-00651}\(\mathsf Y A\) & \texttt{\detokenize{\mathsf Y}} & temporal previous operator & P9 legacy unit 4 temporal-automata & 3 & \hyperref[sec:p09-temporal-automata]{\ref*{sec:p09-temporal-automata}} & \thepage & Prefix unary past operator on pointed words; false at position zero. \\
\phantomsection\label{entry:symbol-registry-00652}\(A\mathbin{\mathsf S}B\) & \texttt{\detokenize{\mathbin{\mathsf S}}} & temporal since operator & P9 legacy unit 4 temporal-automata & 3.5 & \hyperref[sec:p09-temporal-automata]{\ref*{sec:p09-temporal-automata}} & \thepage & Past-time binary operator with a witness at or before the current position and A at every strictly intervening position. \\
\phantomsection\label{entry:symbol-registry-00653}\((Q,\Sigma,\delta,q_0,\Omega)\) & \texttt{\detokenize{Q,}\allowbreak\detokenize{\Sigma,}\allowbreak\detokenize{\delta,}\allowbreak\detokenize{q_0,}\allowbreak\detokenize{\Omega}} & deterministic parity word-automaton datum & P9 legacy unit 4 temporal-automata & n/a & \hyperref[sec:p09-temporal-automata]{\ref*{sec:p09-temporal-automata}} & \thepage & Omega assigns a finite priority to each state; the least priority occurring infinitely often determines acceptance. \\
\phantomsection\label{entry:symbol-registry-00654}\((E_i,F_i)_{i<k}\) & \texttt{\detokenize{E_i,}\allowbreak\detokenize{F_i}} & Rabin acceptance pairs & P9 legacy unit 4 temporal-automata & n/a & \hyperref[sec:p09-temporal-automata]{\ref*{sec:p09-temporal-automata}} & \thepage & A run accepts when one pair has no infinitely recurring E-state and at least one infinitely recurring F-state. \\
\phantomsection\label{entry:symbol-registry-00655}\(\mathsf{AP}(A),\ell_A(s)\) & \texttt{\detokenize{\mathsf{AP},}\allowbreak\detokenize{\ell_A}} & formula alphabet and induced state-label map & P9 legacy unit 4 temporal-automata & meta & \hyperref[sec:p09-temporal-automata]{\ref*{sec:p09-temporal-automata}} & \thepage & AP(A) is the finite set of proposition letters occurring in A; ell-A maps a state to the subset true there under the predicate-indexed valuation V. \\
\phantomsection\label{entry:symbol-registry-00656}\(\Sigma=\bigsqcup_{n\geq0}\Sigma_n,T_\Sigma\) & \texttt{\detokenize{\Sigma,}\allowbreak\detokenize{\Sigma_n,}\allowbreak\detokenize{T_\Sigma}} & finite ranked alphabet and finite tree set & P9 legacy unit 4 temporal-automata & meta & \hyperref[sec:p09-temporal-automata]{\ref*{sec:p09-temporal-automata}} & \thepage & Every symbol in Sigma-n has arity n and T-Sigma is the least finite-term set generated according to those ranks. \\
\phantomsection\label{entry:symbol-registry-00657}\(\mathcal A=(Q,\Sigma,\Delta,F),L(\mathcal A)\) & \texttt{\detokenize{\mathcal A,}\allowbreak\detokenize{Q,}\allowbreak\detokenize{\Sigma,}\allowbreak\detokenize{\Delta,}\allowbreak\detokenize{F,}\allowbreak\detokenize{L}} & bottom-up finite-tree automaton and accepted language & P9 legacy unit 4 temporal-automata & meta & \hyperref[sec:p09-temporal-automata]{\ref*{sec:p09-temporal-automata}} & \thepage & Delta contains ranked bottom-up transitions and acceptance requires a run whose root state lies in F; distinct from infinite-word and Rabin infinite-tree tuples. \\
\phantomsection\label{entry:symbol-registry-00658}\(\operatorname{Pre}_{\exists},\operatorname{Pre}_{\forall}\) & \texttt{\detokenize{\operatorname{Pre}_{\exists},}\allowbreak\detokenize{\operatorname{Pre}_{\forall}}} & existential and universal predecessor maps & P9 legacy unit 5 model-checking & meta & \hyperref[sec:p09-model-checking]{\ref*{sec:p09-model-checking}} & \thepage & Algorithmic maps from state sets to state sets in a finite transition structure; not CTL path quantifiers. \\
\phantomsection\label{entry:symbol-registry-00659}\(L(A)\) & \texttt{\detokenize{L}} & CTL labelling truth set & P9 legacy unit 5 model-checking & meta & \hyperref[sec:p09-model-checking]{\ref*{sec:p09-model-checking}} & \thepage & The set of states assigned to subformula A by bottom-up CTL labelling; locally typed apart from languages and well-founded bounds. \\
\phantomsection\label{entry:symbol-registry-00660}\(\operatorname{gfp}(F)\) & \texttt{\detokenize{\operatorname{gfp}}} & greatest fixed point of a monotone map & P9 legacy units 5–7 model-checking to logic-programming & meta & \hyperref[sec:p09-model-checking]{\ref*{sec:p09-model-checking}} & \thepage & The greatest set fixed by the displayed monotone operator; CTL formulas continue to use the object-language nu binder. \\
\phantomsection\label{entry:symbol-registry-00661}\(\Sigma_n^{N\mu},\Pi_n^{N\mu};\Sigma_n^{S\mu},\Pi_n^{S\mu};\Sigma_n^{EL\mu},\Pi_n^{EL\mu}\) & \texttt{\detokenize{\Sigma_n^{N\mu},}\allowbreak\detokenize{\Pi_n^{N\mu},}\allowbreak\detokenize{\Sigma_n^{S\mu},}\allowbreak\detokenize{\Pi_n^{S\mu},}\allowbreak\detokenize{\Sigma_n^{EL\mu},}\allowbreak\detokenize{\Pi_n^{EL\mu}}} & Niwiński simple and Emerson--Lei modal-mu alternation classes & P9 legacy unit 5 model-checking & meta & \hyperref[sec:p09-model-checking]{\ref*{sec:p09-model-checking}} & \thepage & Source-specific positive-form syntactic classes; the variants differ by capture-avoiding substitution closure, closed substitution only, or no substitution closure. \\
\phantomsection\label{entry:symbol-registry-00662}\(M=(I(V),T(V,V'),F(V))\) & \texttt{\detokenize{M,}\allowbreak\detokenize{I,}\allowbreak\detokenize{T,}\allowbreak\detokenize{F,}\allowbreak\detokenize{V,}\allowbreak\detokenize{V'}} & Boolean finite transition-system presentation & P9 legacy unit 5 model-checking & meta & \hyperref[sec:p09-model-checking]{\ref*{sec:p09-model-checking}} & \thepage & I and F denote initial and bad states and total T relates current to next state variables in the McMillan interpolation procedure. \\
\phantomsection\label{entry:symbol-registry-00663}\(W_i,R,P,\operatorname{rd}(M,F)\) & \texttt{\detokenize{W_i,}\allowbreak\detokenize{R,}\allowbreak\detokenize{P,}\allowbreak\detokenize{\operatorname{rd}}} & state copies reachable overapproximation interpolant and reverse depth & P9 legacy unit 5 model-checking & meta & \hyperref[sec:p09-model-checking]{\ref*{sec:p09-model-checking}} & \thepage & W-i are renamed state blocks R overapproximates reachable states P is the one-state interpolant and rd is the maximum finite shortest distance to the bad set. \\
\phantomsection\label{entry:symbol-registry-00664}\(\mathcal D=(S,P,L),\mathcal M=(S,\mathsf{Act},P,L)\) & \texttt{\detokenize{\mathcal D,}\allowbreak\detokenize{\mathcal M,}\allowbreak\detokenize{S,}\allowbreak\detokenize{P,}\allowbreak\detokenize{L,}\allowbreak\detokenize{\mathsf{Act}}} & finite DTMC and MDP data & P9 legacy unit 5 model-checking & n/a & \hyperref[sec:p09-model-checking]{\ref*{sec:p09-model-checking}} & \thepage & P is row-stochastic for the chain and action-indexed for the MDP; only a scheduler-resolved MDP induces one path measure. \\
\phantomsection\label{entry:symbol-registry-00665}\(\mathsf{Act}(s),\sigma:\operatorname{Paths}_{\mathrm{fin}}(\mathcal M)\to\mathsf{Act},\operatorname{Cyl}(\pi)\) & \texttt{\detokenize{\mathsf{Act},}\allowbreak\detokenize{\sigma,}\allowbreak\detokenize{\operatorname{Paths}_{\mathrm{fin}},}\allowbreak\detokenize{\operatorname{Cyl}}} & enabled actions scheduler and finite-path cylinder & P9 legacy unit 5 model-checking & meta & \hyperref[sec:p09-model-checking]{\ref*{sec:p09-model-checking}} & \thepage & Act-s is nonempty sigma chooses an enabled action from each finite history and Cyl-pi is the set of infinite paths extending pi; PCTL extrema range over exactly this scheduler class. \\
\phantomsection\label{entry:symbol-registry-00666}\(\Pr_s,\Pr_s^{\min},\Pr_s^{\max}\) & \texttt{\detokenize{\Pr_s,}\allowbreak\detokenize{\Pr_s^{\min},}\allowbreak\detokenize{\Pr_s^{\max}}} & path and extremal scheduler probabilities & P9 legacy unit 5 model-checking & meta & \hyperref[sec:p09-model-checking]{\ref*{sec:p09-model-checking}} & \thepage & The undecorated measure belongs to a DTMC or CTMC or to an MDP after fixing a scheduler; min and max quantify over MDP schedulers. \\
\phantomsection\label{entry:symbol-registry-00667}\(\mathsf P_{\bowtie q}[\psi],\mathsf P^{\min}_{\bowtie q}[\psi],\mathsf P^{\max}_{\bowtie q}[\psi]\) & \texttt{\detokenize{\mathsf P_{\bowtie q},}\allowbreak\detokenize{\mathsf P^{\min}_{\bowtie q},}\allowbreak\detokenize{\mathsf P^{\max}_{\bowtie q}}} & probability-threshold modalities & P9 legacy unit 5 model-checking & 3 & \hyperref[sec:p09-model-checking]{\ref*{sec:p09-model-checking}} & \thepage & The undecorated form uses the active DTMC or CTMC path measure; min and max forms take the MDP infimum and supremum over the declared scheduler class. \\
\phantomsection\label{entry:symbol-registry-00668}\(\mathsf U^{\leq k},\mathsf X^{I},\mathsf U^{I}\) & \texttt{\detokenize{\mathsf U^{\leq k},}\allowbreak\detokenize{\mathsf X^{I},}\allowbreak\detokenize{\mathsf U^{I}}} & bounded-step and real-time temporal operators & P9 legacy unit 5 model-checking & 3.5 & \hyperref[sec:p09-model-checking]{\ref*{sec:p09-model-checking}} & \thepage & The first bounds DTMC or MDP steps by natural k; X-I tests the first CTMC transition time and U-I bounds elapsed time by a declared real interval I. \\
\phantomsection\label{entry:symbol-registry-00669}\(\mathcal C=(S,R,L),R(s,t),E(s),P(s,t),\pi[i],\pi@t,\tau_i\) & \texttt{\detokenize{\mathcal C,}\allowbreak\detokenize{S,}\allowbreak\detokenize{R,}\allowbreak\detokenize{L,}\allowbreak\detokenize{E,}\allowbreak\detokenize{P,}\allowbreak\detokenize{\pi,}\allowbreak\detokenize{\tau}} & finite CTMC rates embedded jumps and timed paths & P9 legacy unit 5 model-checking & n/a & \hyperref[sec:p09-model-checking]{\ref*{sec:p09-model-checking}} & \thepage & The nonnegative rational rate map is typed on all state pairs and may contain positive self-loops; E sums the full row P is R divided by E at nonabsorbing states and tau-i is the residence time before path position i-plus-one. Pi-at-t uses the locally stated endpoint convention. \\
\phantomsection\label{entry:symbol-registry-00670}\(\mathsf S_{\bowtie q}[\Phi]\) & \texttt{\detokenize{\mathsf S_{\bowtie q}}} & CSL long-run threshold modality & P9 legacy unit 5 model-checking & 3 & \hyperref[sec:p09-model-checking]{\ref*{sec:p09-model-checking}} & \thepage & Compares the starting-state-dependent limiting probability of Phi-states with q; not the temporal since connective. \\
\phantomsection\label{entry:symbol-registry-00671}\((W,\mathcal F,\mu,\pi),w(\varphi)\) & \texttt{\detokenize{W,}\allowbreak\detokenize{\mathcal F,}\allowbreak\detokenize{\mu,}\allowbreak\detokenize{\pi,}\allowbreak\detokenize{w}} & atemporal probability structure and likelihood term & P9 legacy unit 5 model-checking & meta & \hyperref[sec:p09-model-checking]{\ref*{sec:p09-model-checking}} & \thepage & Mu is the event measure and w-phi its value on the phi-event; pi is locally the world-valuation map and not the temporal path metavariable. \\
\phantomsection\label{entry:symbol-registry-00672}\(\mathsf{PDL}\) & \texttt{\detokenize{\mathsf{PDL}}} & propositional dynamic logic and calculus identifier & P9 legacy unit 6 program-logics & n/a & \hyperref[sec:p09-program-logics]{\ref*{sec:p09-program-logics}} & \thepage & The regular-program modal system whose exact grammar semantics and Hilbert sheet are registered in Part IX. \\
\phantomsection\label{entry:symbol-registry-00673}\(\alpha;\beta,\alpha\cup\beta,\alpha^*,A?\) & \texttt{\detokenize{;,}\allowbreak\detokenize{\cup,}\allowbreak\detokenize{{}^*,}\allowbreak\detokenize{?}} & PDL program constructors & P9 legacy unit 6 program-logics & n/a & \hyperref[sec:p09-program-logics]{\ref*{sec:p09-program-logics}} & \thepage & Relational composition choice reflexive-transitive closure and identity test in the mutually generated PDL grammar. \\
\phantomsection\label{entry:symbol-registry-00674}\([\alpha]A\) & \texttt{\detokenize{[\alpha]}} & PDL universal program modality & P9 legacy unit 6 program-logics & 3 & \hyperref[sec:p09-program-logics]{\ref*{sec:p09-program-logics}} & \thepage & True when every alpha-successor satisfies A; the index is a regular program rather than a modal frame label. \\
\phantomsection\label{entry:symbol-registry-00675}\(\derives_{\mathsf{PDL}}A\) & \texttt{\detokenize{\derives_{\mathsf{PDL}}}} & PDL theoremhood judgment & P9 legacy unit 6 program-logics & meta & \hyperref[sec:p09-program-logics]{\ref*{sec:p09-program-logics}} & \thepage & Finite no-premise derivability on sheet:hilbert-pdl. \\
\phantomsection\label{entry:symbol-registry-00676}\(\{P\}\,C\,\{Q\}\) & \texttt{\detokenize{\{\mathord{\cdot}\}\,}\allowbreak\detokenize{\mathord{\cdot}\,}\allowbreak\detokenize{\{\mathord{\cdot}\}}} & Hoare partial or total correctness triple & P9 legacy unit 6 program-logics & meta & \hyperref[sec:p09-program-logics]{\ref*{sec:p09-program-logics}} & \thepage & The active correctness tag determines whether termination is required; braces do not denote sets in this role. \\
\phantomsection\label{entry:symbol-registry-00677}\(\langle C,s\rangle\Downarrow t\) & \texttt{\detokenize{\Downarrow}} & deterministic while-language big-step evaluation & P9 legacy unit 6 program-logics & meta & \hyperref[sec:p09-program-logics]{\ref*{sec:p09-program-logics}} & \thepage & Least evaluation relation from a command and initial integer state to a terminating final state. \\
\phantomsection\label{entry:symbol-registry-00678}\(\operatorname{wp}(C,Q),\operatorname{wlp}(C,Q)\) & \texttt{\detokenize{\operatorname{wp},}\allowbreak\detokenize{\operatorname{wlp}}} & weakest and weakest liberal preconditions & P9 legacy unit 6 program-logics & meta & \hyperref[sec:p09-program-logics]{\ref*{sec:p09-program-logics}} & \thepage & Predicate transformers distinguished by whether termination is required. \\
\phantomsection\label{entry:symbol-registry-00679}\(\derives_{\mathsf{par}},\derives_{\mathsf{tot}}\) & \texttt{\detokenize{\derives_{\mathsf{par}},}\allowbreak\detokenize{\derives_{\mathsf{tot}}}} & partial and total Hoare derivability tags & P9 legacy unit 6 program-logics & meta & \hyperref[sec:p09-program-logics]{\ref*{sec:p09-program-logics}} & \thepage & Judgments on sheet:hoare-while; only the total tag licenses the decreasing-variant loop rule. \\
\phantomsection\label{entry:symbol-registry-00680}\(h_1\perp h_2,h_1\uplus h_2\) & \texttt{\detokenize{\perp,}\allowbreak\detokenize{\uplus}} & disjoint finite heaps and their union & P9 legacy unit 6 program-logics & meta & \hyperref[sec:p09-program-logics]{\ref*{sec:p09-program-logics}} & \thepage & Perpendicular denotes disjoint domains and uplus the defined union of compatible finite partial heaps. \\
\phantomsection\label{entry:symbol-registry-00681}\(P*Q\) & \texttt{\detokenize{*}} & separating conjunction & P9 legacy unit 6 program-logics & 4 & \hyperref[sec:p09-program-logics]{\ref*{sec:p09-program-logics}} & \thepage & Assertion connective splitting one heap into disjoint subheaps satisfying P and Q; distinct from ordinary conjunction and arithmetic multiplication. \\
\phantomsection\label{entry:symbol-registry-00682}\(\mathit{HB}_P,T_P,M_P\) & \texttt{\detokenize{\mathit{HB}_P,}\allowbreak\detokenize{T_P,}\allowbreak\detokenize{M_P}} & Herbrand base immediate-consequence map and least model & P9 legacy unit 7 logic-programming & meta & \hyperref[sec:p09-logic-programming]{\ref*{sec:p09-logic-programming}} & \thepage & HB-P is a set of ground atoms T-P is its monotone consequence operator and M-P is the least fixed point. \\
\phantomsection\label{entry:symbol-registry-00683}\(\leftarrow A_1,\ldots,A_n\) & \texttt{\detokenize{\leftarrow}} & SLD goal & P9 legacy unit 7 logic-programming & meta & \hyperref[sec:p09-logic-programming]{\ref*{sec:p09-logic-programming}} & \thepage & Finite ordered goal sequence; an empty sequence is refutation success rather than a clause with an omitted head. \\
\phantomsection\label{entry:symbol-registry-00684}\(\xRightarrow{\theta}\) & \texttt{\detokenize{\xRightarrow{\theta}}} & SLD transition labelled by an mgu & P9 legacy unit 7 logic-programming & meta & \hyperref[sec:p09-logic-programming]{\ref*{sec:p09-logic-programming}} & \thepage & Applies the displayed most general unifier to the entire resolvent on sheet:sld-resolution. \\
\phantomsection\label{entry:symbol-registry-00685}\(\mathsf{not}\ A\) & \texttt{\detokenize{\mathsf{not}}} & logic-programming default negation & P9 legacy unit 7 logic-programming & 3 & \hyperref[sec:p09-logic-programming]{\ref*{sec:p09-logic-programming}} & \thepage & Ground operational negation by finite failure in SLDNF; not classical object-language negation. \\
\phantomsection\label{entry:symbol-registry-00686}\(\operatorname{Comp}(P)\) & \texttt{\detokenize{\operatorname{Comp}}} & Clark completion of a normal program & P9 legacy unit 7 logic-programming & meta & \hyperref[sec:p09-logic-programming]{\ref*{sec:p09-logic-programming}} & \thepage & Predicate-defining biconditionals plus the explicitly stated signature-relative free-term equality axioms; distinct from least and stable models. \\
\phantomsection\label{entry:symbol-registry-00687}\(P^M\) & \texttt{\detokenize{P^M}} & Gelfond-Lifschitz reduct & P9 legacy unit 7 logic-programming & meta & \hyperref[sec:p09-logic-programming]{\ref*{sec:p09-logic-programming}} & \thepage & Positive program obtained relative to atom set M by deleting defeated rules and remaining default-negated literals. \\
\phantomsection\label{entry:symbol-registry-00688}\(\Gamma_P,L,G\) & \texttt{\detokenize{\Gamma_P,}\allowbreak\detokenize{L,}\allowbreak\detokenize{G}} & stable operator and well-founded approximation bounds & P9 legacy unit 7 logic-programming & meta & \hyperref[sec:p09-logic-programming]{\ref*{sec:p09-logic-programming}} & \thepage & Gamma-P maps M to the least model of P-superscript-M; L and G are the least and greatest fixed points of its square. \\
\phantomsection\label{entry:symbol-registry-00689}\(q(\bar x),D_q\) & \texttt{\detokenize{q,}\allowbreak\detokenize{D_q}} & conjunctive query and canonical database & P9 legacy unit 7 logic-programming & meta & \hyperref[sec:p09-logic-programming]{\ref*{sec:p09-logic-programming}} & \thepage & The canonical finite structure D-q carries the query variables constants atoms and distinguished free tuple. \\
\phantomsection\label{entry:symbol-registry-00690}\(q_1\subseteq q_2\) & \texttt{\detokenize{\subseteq}} & conjunctive-query containment & P9 legacy unit 7 logic-programming & 2 & \hyperref[sec:p09-logic-programming]{\ref*{sec:p09-logic-programming}} & \thepage & Every answer to q1 is an answer to q2 on every finite database; distinct from set inclusion inside one instance. \\
\phantomsection\label{entry:symbol-registry-00691}\(\operatorname{CSP}(B)\) & \texttt{\detokenize{\operatorname{CSP}}} & fixed-template constraint-satisfaction problem & P9 legacy unit 7 logic-programming & n/a & \hyperref[sec:p09-logic-programming]{\ref*{sec:p09-logic-programming}} & \thepage & B is the fixed finite relational target and the input asks for a homomorphism from A to B. \\
\phantomsection\label{entry:symbol-registry-00692}\(\operatorname{Step}_{M,x},\operatorname{Eq},\operatorname{Sel}_b\) & \texttt{\detokenize{\operatorname{Step},}\allowbreak\detokenize{\operatorname{Eq},}\allowbreak\detokenize{\operatorname{Sel}}} & QBF reduction transition equality and selector formulas & P9 legacy unit 8 complexity-boundaries & meta & \hyperref[sec:p09-complexity-boundaries]{\ref*{sec:p09-complexity-boundaries}} & \thepage & Metalinguistic polynomial-size Boolean formulas on configuration blocks; the selector permits one textual recursive occurrence. \\
\phantomsection\label{entry:symbol-registry-00693}\(R_i(u,v)\) & \texttt{\detokenize{R_i}} & bounded configuration-reachability QBF family & P9 legacy unit 8 complexity-boundaries & meta & \hyperref[sec:p09-complexity-boundaries]{\ref*{sec:p09-complexity-boundaries}} & \thepage & True exactly when v is reachable from u within two-to-the-i steps; a metalinguistic formula family rather than a new object-language predicate. \\
\phantomsection\label{entry:symbol-registry-00694}\(c_{\mathrm{init}}(x),c_{\mathrm{acc}}\) & \texttt{\detokenize{c_{\mathrm{init}},}\allowbreak\detokenize{c_{\mathrm{acc}}}} & initial and canonical accepting configuration blocks & P9 legacy unit 8 complexity-boundaries & meta & \hyperref[sec:p09-complexity-boundaries]{\ref*{sec:p09-complexity-boundaries}} & \thepage & Boolean constant blocks substituted into the final closed QBF; not first-order individual constants. \\
\phantomsection\label{entry:symbol-registry-00695}\(\Phi_{M,x}\) & \texttt{\detokenize{\Phi_{M,x}}} & closed QBF produced by the PSPACE reduction & P9 legacy unit 8 complexity-boundaries & n/a & \hyperref[sec:p09-complexity-boundaries]{\ref*{sec:p09-complexity-boundaries}} & \thepage & The polynomial-size formula R-m applied to the initial and canonical accepting configuration blocks. \\
\phantomsection\label{entry:symbol-registry-00700}\(\bar E\) & \texttt{\detokenize{\bar}} & NNF complement of an ALC concept & P9 legacy unit 3 description-logics & meta & \hyperref[sec:p09-description-logics]{\ref*{sec:p09-description-logics}} & \thepage & The recursive involution complements literals and constants and exchanges Boolean and modal duals; its extension is the complement in the interpretation domain. \\
\phantomsection\label{entry:symbol-registry-00701}\(\mathcal C,m\) & \texttt{\detokenize{\mathcal C,}\allowbreak\detokenize{m}} & ALC input closure and its cardinality & P9 legacy unit 3 description-logics & meta & \hyperref[sec:p09-description-logics]{\ref*{sec:p09-description-logics}} & \thepage & Finite set of normalized input concepts their sub-concepts and NNF complements; m is linear in the total input length. \\
\phantomsection\label{entry:symbol-registry-00702}\(t,\operatorname{tp}_I(x)\) & \texttt{\detokenize{t,}\allowbreak\detokenize{\operatorname{tp}}} & complete ALC concept type and realized type & P9 legacy unit 3 description-logics & meta & \hyperref[sec:p09-description-logics]{\ref*{sec:p09-description-logics}} & \thepage & A type is a subset of the fixed closure satisfying Boolean completeness and every internalized GCI; the realized type records truth at x in I. \\
\phantomsection\label{entry:symbol-registry-00703}\(t\leadsto_r u\) & \texttt{\detokenize{\leadsto_r}} & role compatibility of ALC types & P9 legacy unit 3 description-logics & meta & \hyperref[sec:p09-description-logics]{\ref*{sec:p09-description-logics}} & \thepage & Every universal r-restriction in t has its filler in u; this relation supplies all role edges in the canonical model. \\
\phantomsection\label{entry:symbol-registry-00704}\(S_i,S\) & \texttt{\detokenize{S_i,}\allowbreak\detokenize{S}} & surviving ALC type sets & P9 legacy unit 3 description-logics & meta & \hyperref[sec:p09-description-logics]{\ref*{sec:p09-description-logics}} & \thepage & S-zero contains all complete types; each stage removes types lacking existential witnesses; S is the stable survivor set. \\
\phantomsection\label{entry:symbol-registry-00705}\(\mathsf Q_q,\mathsf H_i,\mathsf A_{i,\gamma}\) & \texttt{\detokenize{\mathsf Q,}\allowbreak\detokenize{\mathsf H,}\allowbreak\detokenize{\mathsf A}} & configuration concept names in the ALC hardness reduction & P9 legacy unit 3 description-logics & meta & \hyperref[sec:p09-description-logics]{\ref*{sec:p09-description-logics}} & \thepage & Unary concept names record machine control head position and one tape symbol per cell; their extensions are constrained by finite general inclusions. \\
\phantomsection\label{entry:symbol-registry-00706}\(\mathsf B_\ell,Z,L_\ell\) & \texttt{\detokenize{\mathsf B,}\allowbreak\detokenize{Z,}\allowbreak\detokenize{L}} & binary countdown concepts in the ALC hardness reduction & P9 legacy unit 3 description-logics & meta & \hyperref[sec:p09-description-logics]{\ref*{sec:p09-description-logics}} & \thepage & B-ell records one rank bit; Z and L-ell abbreviate conjunctions saying all bits or all lower bits are zero. \\
\phantomsection\label{entry:symbol-registry-00707}\(\rho(x)\) & \texttt{\detokenize{\rho}} & nonnegative rank of an ALC reduction model element & P9 legacy unit 3 description-logics & meta & \hyperref[sec:p09-description-logics]{\ref*{sec:p09-description-logics}} & \thepage & Integer obtained from the concept-membership bits at x; every transition role edge decreases this value by exactly one. \\
\phantomsection\label{entry:symbol-registry-00708}\(C_{M,w},\mathcal T_{M,s}\) & \texttt{\detokenize{C_{M,w},}\allowbreak\detokenize{\mathcal T_{M,s}}} & ALC concept and TBox encoding alternating acceptance & P9 legacy unit 3 description-logics & meta & \hyperref[sec:p09-description-logics]{\ref*{sec:p09-description-logics}} & \thepage & The concept fixes the input configuration and maximal binary rank; the finite TBox enforces configurations branching transitions and countdown. \\
\phantomsection\label{entry:symbol-registry-00904}\(\mathsf{Lin}\) & \texttt{\detokenize{\mathsf{Lin}}} & Kamp class of nonempty strict linear orders & P9 legacy unit 4 temporal-automata & meta & \hyperref[sec:p09-temporal-automata]{\ref*{sec:p09-temporal-automata}} & \thepage & The III.1 normal form allows arbitrary unary interpretations, endpoints and empty gaps. Density and order completeness are not assumed. \\
\phantomsection\label{entry:symbol-registry-00905}\(x\doteq y\) & \texttt{\detokenize{\doteq}} & Kamp coincidence abbreviation & P9 legacy unit 4 temporal-automata & 2 & \hyperref[sec:p09-temporal-automata]{\ref*{sec:p09-temporal-automata}} & \thepage & In the equality-free source language, abbreviates \(\neg(x<y)\land\neg(y<x)\). It expresses equality of assigned points over \(\mathsf{Lin}\), with a different local role from a unification constraint. \\
\phantomsection\label{entry:symbol-registry-00906}\(\top_x,\bot_x\) & \texttt{\detokenize{\top_x,}\allowbreak\detokenize{\bot_x}} & Kamp starred constant slots & P9 legacy unit 4 temporal-automata & meta & \hyperref[sec:p09-temporal-automata]{\ref*{sec:p09-temporal-automata}} & \thepage & The formulas \(\neg(x<x)\) and \(x<x\) retain the variable \(x\). They belong to the starred rendition's atomic component class, not Kamp's literal unextended unary-literal class. \\
\phantomsection\label{entry:symbol-registry-00907}\((B_1,\ldots,B_s),(y_1,\ldots,y_s)\) & \texttt{\detokenize{B_k,}\allowbreak\detokenize{y_k,}\allowbreak\detokenize{s}} & ordered coincidence blocks and essential representatives & P9 legacy unit 4 temporal-automata & meta & \hyperref[sec:p09-temporal-automata]{\ref*{sec:p09-temporal-automata}} & \thepage & The blocks partition the finite nonempty free-variable set in assigned-point order; each representative is the first variable name in its block. Representative order need not agree with variable-name order. \\
\phantomsection\label{entry:symbol-registry-00908}\(\pi,l(\pi)\) & \texttt{\detokenize{\pi,}\allowbreak\detokenize{l}} & Kamp standard temporal prefix and its length & P9 legacy unit 4 temporal-automata & meta & \hyperref[sec:p09-temporal-automata]{\ref*{sec:p09-temporal-automata}} & \thepage & The prefix records adjacent strict comparisons and representative equalities; \(l(\pi)=s\) counts essential variables. A single named variable uses self-coincidence. \\
\phantomsection\label{entry:symbol-registry-00909}\(\operatorname{At}(x),T^K_x,F^K_x\) & \texttt{\detokenize{\operatorname{At},}\allowbreak\detokenize{T^K_x,}\allowbreak\detokenize{F^K_x}} & Kamp literal atomic components and source constants & P9 legacy unit 4 temporal-automata & meta & \hyperref[sec:p09-temporal-automata]{\ref*{sec:p09-temporal-automata}} & \thepage & Nonempty conjunctions of unary literals and the specific tautology \(P_0(x)\lor\neg P_0(x)\). This tautology and \(P_0(x)\land\neg P_0(x)\) supply the source's true and false components. \\
\phantomsection\label{entry:symbol-registry-00910}\(\operatorname{Bef}_n(x),\operatorname{Aft}_n(x),\operatorname{Bet}_n(x,y)\) & \texttt{\detokenize{\operatorname{Bef}_n,}\allowbreak\detokenize{\operatorname{Aft}_n,}\allowbreak\detokenize{\operatorname{Bet}_n}} & finite stages of Kamp regional component classes & P9 legacy unit 4 temporal-automata & meta & \hyperref[sec:p09-temporal-automata]{\ref*{sec:p09-temporal-automata}} & \thepage & Stage zero is empty; each next stage retains its predecessor and adds the stated guarded existential-conjunction and universal-disjunction constructors. The source-index regularization is identified in its provenance note. \\
\phantomsection\label{entry:symbol-registry-00911}\(\operatorname{Bef}(x),\operatorname{Aft}(x),\operatorname{Bet}(x,y)\) & \texttt{\detokenize{\operatorname{Bef},}\allowbreak\detokenize{\operatorname{Aft},}\allowbreak\detokenize{\operatorname{Bet}}} & Kamp first-order regional component classes & P9 legacy unit 4 temporal-automata & meta & \hyperref[sec:p09-temporal-automata]{\ref*{sec:p09-temporal-automata}} & \thepage & Unions of finite stages, with only the displayed endpoints free. They describe the before, after and open-between regions and are not themselves temporal operators. \\
\phantomsection\label{entry:symbol-registry-00912}\(\operatorname{At}^{*},\operatorname{Bef}^{*},\operatorname{Aft}^{*},\operatorname{Bet}^{*}\) & \texttt{\detokenize{\operatorname{At}^{*},}\allowbreak\detokenize{\operatorname{Bef}^{*},}\allowbreak\detokenize{\operatorname{Aft}^{*},}\allowbreak\detokenize{\operatorname{Bet}^{*}}} & starred component classes for support preservation & P9 legacy unit 4 temporal-automata & meta & \hyperref[sec:p09-temporal-automata]{\ref*{sec:p09-temporal-automata}} & \thepage & Repeat the regional recursion with the order-defined constant slots. The III.1 annex suppresses stars only within its explicitly declared starred proof scope. \\
\phantomsection\label{entry:symbol-registry-00913}\(\varepsilon(N^*_{\varphi})\) & \texttt{\detokenize{\varepsilon}} & Kamp source-grammar conversion & P9 legacy unit 4 temporal-automata & meta & \hyperref[sec:p09-temporal-automata]{\ref*{sec:p09-temporal-automata}} & \thepage & Replaces distinguished starred constants by the source constants at the same variable and commutes with regional constructors. It preserves truth and free variables; its constant encodings can introduce \(P_0\). \\
\phantomsection\label{entry:symbol-registry-00914}\(\mathcal C_\pi\) & \texttt{\detokenize{\mathcal C_\pi}} & permitted local components of a Kamp standard clause & P9 legacy unit 4 temporal-automata & meta & \hyperref[sec:p09-temporal-automata]{\ref*{sec:p09-temporal-automata}} & \thepage & Atomic components at representatives, between components at consecutive representatives, and before/after components at the two extremes. The annex's declared starred scope changes all these component classes together. \\
\phantomsection\label{entry:symbol-registry-00915}\(N_\varphi,N^*_{\varphi}\) & \texttt{\detokenize{N_\varphi,}\allowbreak\detokenize{N^*_{\varphi}}} & source-standard and starred normal forms & P9 legacy unit 4 temporal-automata & meta & \hyperref[sec:p09-temporal-automata]{\ref*{sec:p09-temporal-automata}} & \thepage & Effectively constructed III.1 forms on exactly the input free variables. Unary-predicate support preservation applies to the starred form; the source form is its conversion under \(\varepsilon\). \\
\phantomsection\label{entry:symbol-registry-00916}\(H_{\rho,\beta}\) & \texttt{\detokenize{H_{\rho,\beta}}} & Kamp universal-step position group & P9 legacy unit 4 temporal-automata & meta & \hyperref[sec:p09-temporal-automata]{\ref*{sec:p09-temporal-automata}} & \thepage & Disjunction of local conjunctions sharing residual prefix \(\rho\) and position guard \(\beta\). All \(2s+1\) positions are retained; an absent group has the appropriately typed false component. \\
\phantomsection\label{entry:symbol-registry-00917}\(\widehat R\) & \texttt{\detokenize{\widehat R}} & Kamp predicate-to-proposition correspondence & P9 legacy unit 4 temporal-automata & meta & \hyperref[sec:p09-temporal-automata]{\ref*{sec:p09-temporal-automata}} & \thepage & For unary predicate letter \(R=P_i\), the corresponding temporal proposition is \(\widehat R=q_i\). Named-point expression uses the fixed translation \(D\) and capture-avoiding current-variable renaming. \\
\phantomsection\label{entry:symbol-registry-00918}\(\operatorname{Bp}_n(x,y),\operatorname{Bf}_n(x,y)\) & \texttt{\detokenize{\operatorname{Bp}_n,}\allowbreak\detokenize{\operatorname{Bf}_n}} & Kamp binary basic-pattern stages & P9 legacy unit 4 temporal-automata & meta & \hyperref[sec:p09-temporal-automata]{\ref*{sec:p09-temporal-automata}} & \thepage & Exact stages indexed by \(n\geq1\), with two free endpoints and \(n\) interval pieces. The binary base has no strict endpoint guard; these stages do not retain all earlier stages. \\
\phantomsection\label{entry:symbol-registry-00919}\(\operatorname{Bp}_n(w),\operatorname{Bf}_n(w),\operatorname{Bp}(w),\operatorname{Bf}(w)\) & \texttt{\detokenize{\operatorname{Bp},}\allowbreak\detokenize{\operatorname{Bf}}} & Kamp guarded one-variable basic patterns & P9 legacy unit 4 temporal-automata & meta & \hyperref[sec:p09-temporal-automata]{\ref*{sec:p09-temporal-automata}} & \thepage & Existential past or future closures supply a strict inequality between the two binary endpoints. The unindexed classes are unions over \(n\geq1\), distinct from the III.1 regional classes. \\
\phantomsection\label{entry:symbol-registry-00920}\(\sigma(\phi)=(A_0,B_1,A_1,\ldots,B_n,A_n)\) & \texttt{\detokenize{\sigma}} & Kamp basic-pattern predicate occurrence word & P9 legacy unit 4 temporal-automata & meta & \hyperref[sec:p09-temporal-automata]{\ref*{sec:p09-temporal-automata}} & \thepage & The ordered \(2n+1\) occurrences retain repeated predicate letters. This is the source's sequence \(S(\phi)\), renamed to distinguish it from strict since. \\
\phantomsection\label{entry:symbol-registry-00921}\(a_0<\cdots<a_n\) & \texttt{\detokenize{a_0,}\allowbreak\detokenize{\ldots,}\allowbreak\detokenize{a_n}} & realization of a Kamp predicate word & P9 legacy unit 4 temporal-automata & meta & \hyperref[sec:p09-temporal-automata]{\ref*{sec:p09-temporal-automata}} & \thepage & Each point satisfies its endpoint letter and each open interval satisfies its interval letter throughout. Empty intervals impose no interval-predicate restriction, while every displayed strict point remains required. \\
\phantomsection\label{entry:symbol-registry-00922}\(L_i,R_i\) & \texttt{\detokenize{L_i,}\allowbreak\detokenize{R_i}} & explicit temporal recurrences for Kamp basic patterns & P9 legacy unit 4 temporal-automata & meta & \hyperref[sec:p09-temporal-automata]{\ref*{sec:p09-temporal-automata}} & \thepage & The past recurrence grows by strict since from \(L_0\); the future recurrence grows by strict until backward from \(R_n\). Their first temporal operands name witness conditions, and their second operands name interval conditions. \\
\phantomsection\label{entry:symbol-registry-00995}\(H',G',P',F'\) & \texttt{\detokenize{H',}\allowbreak\detokenize{G',}\allowbreak\detokenize{P',}\allowbreak\detokenize{F'}} & Kamp derived strict interval and limit operations & P9 legacy unit 4 temporal-automata & meta & \hyperref[sec:p09-temporal-automata]{\ref*{sec:p09-temporal-automata}} & \thepage & The displayed definitions expand through the actual S/U witness-first conventions. P' and F' require a satisfying point in every strict past or future adjoining interval; they are true when the respective cut domain is empty. \\
\phantomsection\label{entry:symbol-registry-00996}\(\partial^-\theta,\partial^+\theta\) & \texttt{\detokenize{\partial^-,}\allowbreak\detokenize{\partial^+}} & first-order strict past and future limit formula operations & P9 legacy unit 4 temporal-automata & meta & \hyperref[sec:p09-temporal-automata]{\ref*{sec:p09-temporal-automata}} & \thepage & Fresh quantified points witness theta strictly between every cut and the current point. These symbols are formula operations, not analytic derivatives, and impose no current-point truth requirement. \\
\phantomsection\label{entry:symbol-registry-00997}\(\operatorname{Lp}_n(w),\operatorname{Lf}_n(w)\) & \texttt{\detokenize{\operatorname{Lp}_n,}\allowbreak\detokenize{\operatorname{Lf}_n}} & Kamp limiting-pattern stages under interval-count indexing & P9 legacy unit 4 temporal-automata & meta & \hyperref[sec:p09-temporal-automata]{\ref*{sec:p09-temporal-automata}} & \thepage & Stage zero is a one-letter witness limit. Positive stage n uses an n-interval basic pattern realized wholly inside every adjoining interval. Stages do not retain earlier stages, and no Bp\_0 is introduced. \\
\phantomsection\label{entry:symbol-registry-00998}\(\operatorname{Lp},\operatorname{Lf}\) & \texttt{\detokenize{\operatorname{Lp},}\allowbreak\detokenize{\operatorname{Lf}}} & unions of the limiting past and future classes & P9 legacy unit 4 temporal-automata & meta & \hyperref[sec:p09-temporal-automata]{\ref*{sec:p09-temporal-automata}} & \thepage & Take all finite stages and the declared choices of current free variable. The lemma proves expressibility over Lin, without density, endpointlessness or completeness of the order. \\
\phantomsection\label{entry:symbol-registry-00999}\(E^-_\sigma,E^+_\sigma\) & \texttt{\detokenize{E^-_\sigma,}\allowbreak\detokenize{E^+_\sigma}} & basic temporal expressions ending or beginning at the current point & P9 legacy unit 4 temporal-automata & meta & \hyperref[sec:p09-temporal-automata]{\ref*{sec:p09-temporal-automata}} & \thepage & For positive interval count these are the already proved Lemma 1 expressions L\_n and R\_0. At zero both are the sole proposition letter of the word. \\
\phantomsection\label{entry:symbol-registry-01000}\(\sigma^-,\sigma^+\) & \texttt{\detokenize{\sigma^-,}\allowbreak\detokenize{\sigma^+}} & suffix and prefix words for the limiting recurrence & P9 legacy unit 4 temporal-automata & meta & \hyperref[sec:p09-temporal-automata]{\ref*{sec:p09-temporal-automata}} & \thepage & The minus word deletes the first interval and its first endpoint; the plus word deletes the last interval and its last endpoint. At one interval each becomes a one-letter word. \\
\phantomsection\label{entry:symbol-registry-01001}\(\Lambda^-_\sigma,\Lambda^+_\sigma\) & \texttt{\detokenize{\Lambda^-_\sigma,}\allowbreak\detokenize{\Lambda^+_\sigma}} & explicit temporal expressions for finite limiting patterns & P9 legacy unit 4 temporal-automata & meta & \hyperref[sec:p09-temporal-automata]{\ref*{sec:p09-temporal-automata}} & \thepage & Each recurrence conjoins a strict limit of the whole basic expression with the limit for the shorter suffix or prefix. The second conjunct excludes provisional chains whose remote endpoint crosses the current point. \\
\phantomsection\label{entry:symbol-registry-01002}\(C^-_\sigma(t),C^+_\sigma(t)\) & \texttt{\detokenize{C^-_\sigma,}\allowbreak\detokenize{C^+_\sigma}} & proof-local finite chain conditions inside every adjoining interval & P9 legacy unit 4 temporal-automata & meta & \hyperref[sec:p09-temporal-automata]{\ref*{sec:p09-temporal-automata}} & \thepage & The conditions quantify all strict cuts and require a complete realization of the fixed word inside the intervening interval. They are equivalent to the corresponding limiting-class formulas. \\
\phantomsection\label{entry:symbol-registry-01003}\(\operatorname{CurF}_A(w),\tau_A\) & \texttt{\detokenize{\operatorname{CurF}_A,}\allowbreak\detokenize{\tau_A}} & handbook current-point comparison and its proposition tautology & P9 legacy unit 4 temporal-automata & meta & \hyperref[sec:p09-temporal-automata]{\ref*{sec:p09-temporal-automata}} & \thepage & The comparison places A at the current point and requires an intervening point at every future cut. It is an additional handbook example, distinct from Kamp's actual witness-point future base. Tau\_A is the displayed tautology using the same proposition letter. \\
\phantomsection\label{entry:symbol-registry-01004}\(\operatorname{psub}_{P_i}(\phi,\psi),[\phi]^\psi_{P_i}\) & \texttt{\detokenize{\operatorname{psub}_{P_i},}\allowbreak\detokenize{[\phi]^\psi_{P_i}}} & Kamp proper substitution for a unary predicate & P9 legacy unit 4 temporal-automata & meta & \hyperref[sec:p09-temporal-automata]{\ref*{sec:p09-temporal-automata}} & \thepage & The replacement has exactly one free parameter. Source-selected bound names and a direct fresh-name implementation are distinguished up to alphabetic variation; original target atoms are visited once. \\
\phantomsection\label{entry:symbol-registry-01005}\(\eta[\zeta/q_i]\) & \texttt{\detokenize{\eta[\zeta/}\allowbreak\detokenize{q_i]}} & single-pass temporal proposition substitution & P9 legacy unit 4 temporal-automata & meta & \hyperref[sec:p09-temporal-automata]{\ref*{sec:p09-temporal-automata}} & \thepage & Structural recursion preserves negation, conjunction and the two ordered S/U operands. The inserted replacement is not revisited in the same substitution step. \\
\phantomsection\label{entry:symbol-registry-01006}\(\mathcal G(w),\operatorname{E}_m(w),\operatorname{E}(w),\operatorname{E}\) & \texttt{\detokenize{\mathcal G,}\allowbreak\detokenize{\operatorname{E}_m,}\allowbreak\detokenize{\operatorname{E}}} & Kamp five-generator family and its finite least closure & P9 legacy unit 4 temporal-automata & meta & \hyperref[sec:p09-temporal-automata]{\ref*{sec:p09-temporal-automata}} & \thepage & The retained stages add Boolean constructors and proper predicate substitution with operands in the same stage. This is a syntactic finite closure, not arbitrary quantification or semantic equivalence. \\
\phantomsection\label{entry:symbol-registry-01007}\(g[x\mapsto a]\) & \texttt{\detokenize{g[x\mapsto a]}} & assignment update in the proper-substitution proof & P9 legacy unit 4 temporal-automata & meta & \hyperref[sec:p09-temporal-automata]{\ref*{sec:p09-temporal-automata}} & \thepage & Changes the assignment at the named individual variable while preserving all other values. It is distinct from predicate substitution and from changing a structure's unary interpretation. \\
\phantomsection\label{entry:symbol-registry-01008}\(A,B,\mathcal M_i^A,\mathcal M_i^B\) & \texttt{\detokenize{A,}\allowbreak\detokenize{B,}\allowbreak\detokenize{\mathcal M_i^A,}\allowbreak\detokenize{\mathcal M_i^B}} & replacement truth sets and modified unary interpretations & P9 legacy unit 4 temporal-automata & meta & \hyperref[sec:p09-temporal-automata]{\ref*{sec:p09-temporal-automata}} & \thepage & The proof evaluates each replacement in the original structure and changes only \(P_i\) in the modified structure. A and B here are truth sets, not the unary letters in the later mixed example. \\
\phantomsection\label{entry:symbol-registry-01053}\(\operatorname{Ep}(x,y),\operatorname{Ef}(x,y)\) & \texttt{\detokenize{\operatorname{Ep},}\allowbreak\detokenize{\operatorname{Ef}}} & Kamp elementary past and future interval formula classes & P9 legacy unit 4 temporal-automata & meta & \hyperref[sec:p09-temporal-automata]{\ref*{sec:p09-temporal-automata}} & \thepage & Obtained by finite proper predicate substitutions into the respective basic class. Ordered substitution histories preserve shared and repeated predicate occurrences. \\
\phantomsection\label{entry:symbol-registry-01054}\(W,V,\theta(v),\mathcal B_p(W;x,y),\mathcal B_f(W;x,y)\) & \texttt{\detokenize{W,}\allowbreak\detokenize{V,}\allowbreak\detokenize{\theta(v),}\allowbreak\detokenize{\mathcal B_p,}\allowbreak\detokenize{\mathcal B_f}} & elementary labelled words parameter instances and binary pattern recursions & P9 legacy unit 4 temporal-automata & meta & \hyperref[sec:p09-temporal-automata]{\ref*{sec:p09-temporal-automata}} & \thepage & The words alternate point and open-interval labels from the one-variable elementary class. Parameter instantiation avoids capture; the two recursions have their respective source orientations. \\
\phantomsection\label{entry:symbol-registry-01055}\(I(x,y)\) & \texttt{\detokenize{I(x,}\allowbreak\detokenize{y)}} & Kamp property of guarded interval representation on complete linear orders & P9 legacy unit 4 temporal-automata & meta & \hyperref[sec:p09-temporal-automata]{\ref*{sec:p09-temporal-automata}} & \thepage & A formula with no other free variables is equivalent under x less than y to a finite elementary past disjunction on Com. This does not promote its given representation to Lin. \\
\phantomsection\label{entry:symbol-registry-01056}\(e_i,f_j,b_i,d_j,\mathcal T\) & \texttt{\detokenize{e_i,}\allowbreak\detokenize{f_j,}\allowbreak\detokenize{b_i,}\allowbreak\detokenize{d_j,}\allowbreak\detokenize{\mathcal T}} & distinct point and interval occurrence tags and their finite set & P9 legacy unit 4 temporal-automata & meta & \hyperref[sec:p09-temporal-automata]{\ref*{sec:p09-temporal-automata}} & \thepage & Tags record occurrences rather than distinct formulas. Shared or repeated formula labels keep separate positions and constraints. \\
\phantomsection\label{entry:symbol-registry-01057}\(F,\sim_F,\mathcal F_{m,n}\) & \texttt{\detokenize{F,}\allowbreak\detokenize{\sim_F,}\allowbreak\detokenize{\mathcal F_{m,n}}} & finite fusion preorder coincidence and admissible fusion family & P9 legacy unit 4 temporal-automata & meta & \hyperref[sec:p09-temporal-automata]{\ref*{sec:p09-temporal-automata}} & \thepage & Each input sequence remains strictly ordered between its distinct positions; the two first endpoints coincide and the two last endpoints coincide. Cross-sequence interior coincidences are permitted. \\
\phantomsection\label{entry:symbol-registry-01058}\(C_k,r_F,h_F\) & \texttt{\detokenize{C_k,}\allowbreak\detokenize{r_F,}\allowbreak\detokenize{h_F}} & ordered fusion classes their count and point-position map & P9 legacy unit 4 temporal-automata & meta & \hyperref[sec:p09-temporal-automata]{\ref*{sec:p09-temporal-automata}} & \thepage & The class count equals the source fusion length. Point position k corresponds to alternating slot \(2k-1\), rather than to slot k. \\
\phantomsection\label{entry:symbol-registry-01059}\(f_F\) & \texttt{\detokenize{f_F}} & map from occurrence tags to required alternating slots & P9 legacy unit 4 temporal-automata & meta & \hyperref[sec:p09-temporal-automata]{\ref*{sec:p09-temporal-automata}} & \thepage & A point tag requires its singleton odd slot. An interval tag requires every slot strictly between its endpoint slots, including inserted odd point slots; unlisted slots impose no negative condition. \\
\phantomsection\label{entry:symbol-registry-01060}\(\rho_{F,s},\chi_F\) & \texttt{\detokenize{\rho_{F,s},}\allowbreak\detokenize{\chi_F}} & slot conjunction and elementary formula for one fusion & P9 legacy unit 4 temporal-automata & meta & \hyperref[sec:p09-temporal-automata]{\ref*{sec:p09-temporal-automata}} & \thepage & The full input conjunction corresponds to the finite disjunction over admissible fusions. A particular strict fusion need not follow from the input conjunction. \\
\phantomsection\label{entry:symbol-registry-01061}\(A_i,B_i,C_i,D_i\) & \texttt{\detokenize{A_i,}\allowbreak\detokenize{B_i,}\allowbreak\detokenize{C_i,}\allowbreak\detokenize{D_i}} & elementary formula labels in the three-case fusion example & P9 legacy unit 4 temporal-automata & meta & \hyperref[sec:p09-temporal-automata]{\ref*{sec:p09-temporal-automata}} & \thepage & These are one-variable elementary formulas in the declared parameter, not an extra family of primitive predicate symbols. \\
\phantomsection\label{entry:symbol-registry-01104}\(W\mathbin{\star_\eta}V\) & \texttt{\detokenize{\star_\eta}} & elementary interval-word concatenation with a midpoint label & P9 legacy unit 4 temporal-automata & meta & \hyperref[sec:p09-temporal-automata]{\ref*{sec:p09-temporal-automata}} & \thepage & Joins two positive-interval words and conjoins the midpoint condition with both former endpoint labels. Every original gap label remains; the operation is not an object-language connective. \\
\phantomsection\label{entry:symbol-registry-01105}\(a,b,\eta\) & \texttt{\detokenize{a,}\allowbreak\detokenize{b,}\allowbreak\detokenize{\eta}} & input interval counts and elementary midpoint formula & P9 legacy unit 4 temporal-automata & meta & \hyperref[sec:p09-temporal-automata]{\ref*{sec:p09-temporal-automata}} & \thepage & The two counts are positive in the word operation. The midpoint label is in the one-variable elementary class and is evaluated in the whole structure. \\
\phantomsection\label{entry:symbol-registry-01106}\(\theta(z),\Phi(x,z),\Psi(z,y),\Gamma(x,y)\) & \texttt{\detokenize{\theta,}\allowbreak\detokenize{\Phi,}\allowbreak\detokenize{\Psi,}\allowbreak\detokenize{\Gamma}} & finite midpoint and side aggregates and their guarded existential & P9 legacy unit 4 temporal-automata & meta & \hyperref[sec:p09-temporal-automata]{\ref*{sec:p09-temporal-automata}} & \thepage & All side conditions use the same strict interior witness. Empty families have explicit tautology conventions, while the all-empty aggregate still requires a strict midpoint. \\
\phantomsection\label{entry:symbol-registry-01107}\(\mathcal L,\mathcal R,W_\rho,V_\sigma\) & \texttt{\detokenize{\mathcal L,}\allowbreak\detokenize{\mathcal R,}\allowbreak\detokenize{W_\rho,}\allowbreak\detokenize{V_\sigma}} & finite elementary representation lists and their labelled words & P9 legacy unit 4 temporal-automata & meta & \hyperref[sec:p09-temporal-automata]{\ref*{sec:p09-temporal-automata}} & \thepage & Every side-conjunction disjunct and fusion is retained, then every final left/right disjunct pair is joined. Given interval-property representations remain on Com. \\
\phantomsection\label{entry:symbol-registry-01108}\(Q_0,Q',Q'',\pi,\pi'\) & \texttt{\detokenize{Q_0,}\allowbreak\detokenize{Q',}\allowbreak\detokenize{Q'',}\allowbreak\detokenize{\pi,}\allowbreak\detokenize{\pi'}} & fresh midpoint and endpoint placeholders and basic source concatenations & P9 legacy unit 4 temporal-automata & meta & \hyperref[sec:p09-temporal-automata]{\ref*{sec:p09-temporal-automata}} & \thepage & The first right placeholder is removed before concatenation. Proper substitution for the remaining midpoint inserts all three point labels and retains its self-target once; temporary letters disappear when labels are restored. \\
\phantomsection\label{entry:symbol-registry-01126}\(\phi(r,s),\phi^{\exists<}(y)\) & \texttt{\detokenize{\phi(r,}\allowbreak\detokenize{s),}\allowbreak\detokenize{\phi^{\exists<}}} & semantic basic-pattern assertion and guarded past closure & P9 legacy unit 4 temporal-automata & meta & \hyperref[sec:p09-temporal-automata]{\ref*{sec:p09-temporal-automata}} & \thepage & The first expression abbreviates truth under the displayed endpoint assignment. The second is a formula with a fresh existential initial endpoint strictly before y; that endpoint is not restricted by the separately fixed a. \\
\phantomsection\label{entry:symbol-registry-01127}\(a,b,c,d\) & \texttt{\detokenize{a,}\allowbreak\detokenize{b,}\allowbreak\detokenize{c,}\allowbreak\detokenize{d}} & four ordered points in Kamp's fixed-initial-point lemma & P9 legacy unit 4 temporal-automata & meta & \hyperref[sec:p09-temporal-automata]{\ref*{sec:p09-temporal-automata}} & \thepage & They satisfy \(a<b<c<d\) in a strict linear order. Realizations from a to b and a to d use the identical basic pattern and predicate interpretation; density and completeness are not required. \\
\phantomsection\label{entry:symbol-registry-01128}\(\mathbf v,\mathbf w^-,\mathbf w^+,\mathbf w,\mathbf r\) & \texttt{\detokenize{\mathbf v,}\allowbreak\detokenize{\mathbf w^-,}\allowbreak\detokenize{\mathbf w^+,}\allowbreak\detokenize{\mathbf w,}\allowbreak\detokenize{\mathbf r}} & realizing sequences and their joined sequence & P9 legacy unit 4 temporal-automata & meta & \hyperref[sec:p09-temporal-automata]{\ref*{sec:p09-temporal-automata}} & \thepage & Every sequence has exactly n+1 points for the one fixed n-interval word. The chosen old prefix and past-realization suffix meet at an explicitly justified switch, preserving all point and open-gap labels. \\
\phantomsection\label{entry:symbol-registry-01129}\(k,J,K,j\) & \texttt{\detokenize{k,}\allowbreak\detokenize{J,}\allowbreak\detokenize{K,}\allowbreak\detokenize{j}} & switch indices and finite comparison sets & P9 legacy unit 4 temporal-automata & meta & \hyperref[sec:p09-temporal-automata]{\ref*{sec:p09-temporal-automata}} & \thepage & The fixed-initial-point proof finds a switch in indices zero through n-1. Its finite comparison sets establish existence; endpoint switches and equality of initial points are included. \\
\phantomsection\label{entry:symbol-registry-01142}\(\ell(W),\ell(\phi)\) & \texttt{\detokenize{\ell(W),}\allowbreak\detokenize{\ell(\phi)}} & occurrence length of a basic word and its pattern & P9 legacy unit 4 temporal-automata & meta & \hyperref[sec:p09-temporal-automata]{\ref*{sec:p09-temporal-automata}} & \thepage & The length counts every point and gap occurrence with repetition. An n-interval word has length \(2n+1\); the residual restriction compares this occurrence length. \\
\phantomsection\label{entry:symbol-registry-01143}\(\mathcal C_n,\mathrm p,\mathrm g,c\) & \texttt{\detokenize{\mathcal C_n,}\allowbreak\detokenize{\mathrm p,}\allowbreak\detokenize{\mathrm g,}\allowbreak\detokenize{c}} & finite ordered cut set point and gap tags and a cut index & P9 legacy unit 4 temporal-automata & meta & \hyperref[sec:p09-temporal-automata]{\ref*{sec:p09-temporal-automata}} & \thepage & There are n gap cuts and n-1 interior-point cuts. The point and gap tags distinguish positions by their occurrence role even when a predicate letter is repeated. \\
\phantomsection\label{entry:symbol-registry-01144}\(W_c^L,W_c^R,L_c,R_c\) & \texttt{\detokenize{W_c^L,}\allowbreak\detokenize{W_c^R,}\allowbreak\detokenize{L_c,}\allowbreak\detokenize{R_c}} & cut words and their basic formulas & P9 legacy unit 4 temporal-automata & meta & \hyperref[sec:p09-temporal-automata]{\ref*{sec:p09-temporal-automata}} & \thepage & Each legal cut has positive-stage left and right basic patterns. A divided gap puts its predicate on the shared point and both remaining open gaps; both factorization directions retain that condition. \\
\phantomsection\label{entry:symbol-registry-01145}\(J,\alpha\) & \texttt{\detokenize{J,}\allowbreak\detokenize{\alpha}} & cut subset and complete Boolean truth assignment & P9 legacy unit 4 temporal-automata & meta & \hyperref[sec:p09-temporal-automata]{\ref*{sec:p09-temporal-automata}} & \thepage & All subsets of the finite cut set are used. A subset records the fixed-side truths while its residual disjunction permits different cuts at different universally tested points. \\
\phantomsection\label{entry:symbol-registry-01146}\(\Sigma_J^L,\Sigma_J^R,\rho_J^L,\rho_J^R\) & \texttt{\detokenize{\Sigma_J^L,}\allowbreak\detokenize{\Sigma_J^R,}\allowbreak\detokenize{\rho_J^L,}\allowbreak\detokenize{\rho_J^R}} & fixed-side conjunctions and varying-side disjunctions in both orientations & P9 legacy unit 4 temporal-automata & meta & \hyperref[sec:p09-temporal-automata]{\ref*{sec:p09-temporal-automata}} & \thepage & Empty conjunctions are true and empty disjunctions false. These are explicit finite aggregates; the corrected right-interior statement does not require each aggregate to be one basic pattern. \\
\phantomsection\label{entry:symbol-registry-01147}\(\Omega_\alpha,\widetilde\Omega_\alpha\) & \texttt{\detokenize{\Omega_\alpha,}\allowbreak\detokenize{\widetilde\Omega_\alpha}} & complete truth-pattern formulas for the fixed side & P9 legacy unit 4 temporal-automata & meta & \hyperref[sec:p09-temporal-automata]{\ref*{sec:p09-temporal-automata}} & \thepage & Positive and negative truth tests are first retained together. Negative tests are removed only in the proved complete existential and disjunctive construction, with its separate converse. \\
\phantomsection\label{entry:symbol-registry-01148}\(a,d,b,a_i,l_i,r_i,z\) & \texttt{\detokenize{a,}\allowbreak\detokenize{d,}\allowbreak\detokenize{b,}\allowbreak\detokenize{a_i,}\allowbreak\detokenize{l_i,}\allowbreak\detokenize{r_i,}\allowbreak\detokenize{z}} & endpoints realization entries and tested cut point & P9 legacy unit 4 temporal-automata & meta & \hyperref[sec:p09-temporal-automata]{\ref*{sec:p09-temporal-automata}} & \thepage & These are local domain points in the cut proofs, distinct from object-language variable names and integer cut indices. Every joined realization retains the original strict endpoint order. \\
\phantomsection\label{entry:symbol-registry-01191}\(F,G,G^Q,Q\) & \texttt{\detokenize{F,}\allowbreak\detokenize{G,}\allowbreak\detokenize{G^Q,}\allowbreak\detokenize{Q}} & finite formula families their freshened family and an added unary atom & P9 legacy unit 4 temporal-automata & meta & \hyperref[sec:p09-temporal-automata]{\ref*{sec:p09-temporal-automata}} & \thepage & Each whole At member receives one fresh atom, preserving the special tautology case. The added atom is later replaced by the declared falsehood; its addition preserves the comparison predecessors when maximal length exceeds one. \\
\phantomsection\label{entry:symbol-registry-01192}\(D_F,\mathcal U_F\) & \texttt{\detokenize{D_F,}\allowbreak\detokenize{\mathcal U_F}} & family disjunction and its universal interior formula & P9 legacy unit 4 temporal-automata & meta & \hyperref[sec:p09-temporal-automata]{\ref*{sec:p09-temporal-automata}} & \thepage & Empty disjunction uses the declared falsehood. The universal guard is strictly between the original free endpoints, whose occurrences inside every member remain fixed. \\
\phantomsection\label{entry:symbol-registry-01193}\(\ell,M(F),c(F),\prec\) & \texttt{\detokenize{\ell,}\allowbreak\detokenize{M(F),}\allowbreak\detokenize{c(F),}\allowbreak\detokenize{\prec}} & formula length maximum length longest-member count and family comparison & P9 legacy unit 4 temporal-automata & meta & \hyperref[sec:p09-temporal-automata]{\ref*{sec:p09-temporal-automata}} & \thepage & At members have length one and basic lengths are occurrence lengths. The empty family has maximum and count zero. Comparison is lexicographic in maximum and count; identical formulas are not counted as separate set members. \\
\phantomsection\label{entry:symbol-registry-01194}\(\rho_j,n_j,L_{j,c},R_{j,c}\) & \texttt{\detokenize{\rho_j,}\allowbreak\detokenize{n_j,}\allowbreak\detokenize{L_{j,c},}\allowbreak\detokenize{R_{j,c}}} & crossing basic patterns interval counts and indexed cut formulas & P9 legacy unit 4 temporal-automata & meta & \hyperref[sec:p09-temporal-automata]{\ref*{sec:p09-temporal-automata}} & \thepage & Every original crossing pattern retains its own cut index. Each contributes no longer residual and at most one of its original length; all fixed-side conjuncts remain in the cut output. \\
\phantomsection\label{entry:symbol-registry-01195}\(\mathcal J,\mathbf J,S_{\mathbf J}^L,S_{\mathbf J}^R,R_{\mathbf J}^L,R_{\mathbf J}^R\) & \texttt{\detokenize{\mathcal J,}\allowbreak\detokenize{\mathbf J,}\allowbreak\detokenize{S_{\mathbf J}^L,}\allowbreak\detokenize{S_{\mathbf J}^R,}\allowbreak\detokenize{R_{\mathbf J}^L,}\allowbreak\detokenize{R_{\mathbf J}^R}} & finite product of cut subsets and both aggregate orientations & P9 legacy unit 4 temporal-automata & meta & \hyperref[sec:p09-temporal-automata]{\ref*{sec:p09-temporal-automata}} & \thepage & The product has one empty tuple when there are no crossing patterns. Every tuple is enumerated, with true empty conjunctions and false empty disjunctions; a varying universal witness may use a different residual disjunct. \\
\phantomsection\label{entry:symbol-registry-01196}\(C_W^+,C_W^-\) & \texttt{\detokenize{C_W^+,}\allowbreak\detokenize{C_W^-}} & accumulation formulas adjacent to the fixed endpoint & P9 legacy unit 4 temporal-automata & meta & \hyperref[sec:p09-temporal-automata]{\ref*{sec:p09-temporal-automata}} & \thepage & The original endpoint stays fixed inside the basic pattern. Their elementary-representative equivalences are guarded by ordered endpoints; they do not assert remote-endpoint accumulation. \\
\phantomsection\label{entry:symbol-registry-01197}\(\lambda_W^+,\lambda_W^-\) & \texttt{\detokenize{\lambda_W^+,}\allowbreak\detokenize{\lambda_W^-}} & elementary limiting labels for a suffix or prefix & P9 legacy unit 4 temporal-automata & meta & \hyperref[sec:p09-temporal-automata]{\ref*{sec:p09-temporal-automata}} & \thepage & For one interval the label is the corresponding one-letter limiting formula. For more intervals it is the exact positive-stage limiting suffix or prefix, lying wholly inside each required adjoining interval. \\
\phantomsection\label{entry:symbol-registry-01198}\(V_W^+,V_W^-\) & \texttt{\detokenize{V_W^+,}\allowbreak\detokenize{V_W^-}} & two-interval elementary representatives of fixed-endpoint accumulation & P9 legacy unit 4 temporal-automata & meta & \hyperref[sec:p09-temporal-automata]{\ref*{sec:p09-temporal-automata}} & \thepage & Each word retains the original endpoint atom and the adjoining gap predicate. Both directions of its guarded equivalence are proved for linear orders. \\
\phantomsection\label{entry:symbol-registry-01199}\(H(t),P(t),R_{a,b},D(t)\) & \texttt{\detokenize{H(t),}\allowbreak\detokenize{P(t),}\allowbreak\detokenize{R_{a,b},}\allowbreak\detokenize{D(t)}} & fixed-endpoint truth guarded closure its interior truth set and remaining disjunction & P9 legacy unit 4 temporal-automata & meta & \hyperref[sec:p09-temporal-automata]{\ref*{sec:p09-temporal-automata}} & \thepage & These semantic assignment schemata have explicit formula versions. The closure may choose any endpoint satisfying its strict guard; the truth set retains the original fixed endpoint. D retains both original endpoints. \\
\phantomsection\label{entry:symbol-registry-01200}\(\mathcal V_D,\mathcal V_{D\lor P}\) & \texttt{\detokenize{\mathcal V_D,}\allowbreak\detokenize{\mathcal V_{D\lor P}}} & universal interval abbreviations with original formula parameters retained & P9 legacy unit 4 temporal-automata & meta & \hyperref[sec:p09-temporal-automata]{\ref*{sec:p09-temporal-automata}} & \thepage & Changing the displayed interval changes only its outer guard. It does not substitute new endpoints inside the remaining disjunction or the distinguished truth formula. \\
\phantomsection\label{entry:symbol-registry-01201}\(l,h,\lambda,\mu\) & \texttt{\detokenize{l,}\allowbreak\detokenize{h,}\allowbreak\detokenize{\lambda,}\allowbreak\detokenize{\mu}} & lower and upper truth-set bounds and their three boundary types & P9 legacy unit 4 temporal-automata & meta & \hyperref[sec:p09-temporal-automata]{\ref*{sec:p09-temporal-automata}} & \thepage & Type zero is an outer endpoint, one an attained interior extremum, and two an unattained interior extremum. The nonempty truth set has unique bounds on complete orders; singleton and empty cases are separate. \\
\phantomsection\label{entry:symbol-registry-01202}\(\Gamma_i,E_j,K_i,Z_i\) & \texttt{\detokenize{\Gamma_i,}\allowbreak\detokenize{E_j,}\allowbreak\detokenize{K_i,}\allowbreak\detokenize{Z_i}} & exact boundary descriptors missing-point tests matrices and guarded branches & P9 legacy unit 4 temporal-automata & meta & \hyperref[sec:p09-temporal-automata]{\ref*{sec:p09-temporal-automata}} & \thepage & Every branch retains its bracketing descriptor, outside-interval tests and tests at unattained interior bounds. Conditional branch equivalence holds on linear orders; completeness supplies the exhaustive eleven-way classification, not an elementary representation by itself. \\
\phantomsection\label{entry:symbol-registry-01203}\(<^{\mathrm{op}},W^{\mathrm{rev}},\phi^{\mathrm{rev}}\) & \texttt{\detokenize{<^{\mathrm{op}},}\allowbreak\detokenize{W^{\mathrm{rev}},}\allowbreak\detokenize{\phi^{\mathrm{rev}}}} & opposite strict order reversed occurrence word and its basic pattern & P9 legacy unit 4 temporal-automata & meta & \hyperref[sec:p09-temporal-automata]{\ref*{sec:p09-temporal-automata}} & \thepage & Reversing a realizing sequence retains the corresponding point labels and the same underlying gap sets. This gives the explicit opposite-orientation fixed-endpoint interpolation proof. \\
\phantomsection\label{entry:symbol-registry-01222}\(\theta^{\mathrm d},W^{\mathrm d},\mathcal M^{\mathrm{op}}\) & \texttt{\detokenize{\theta^{\mathrm d},}\allowbreak\detokenize{W^{\mathrm d},}\allowbreak\detokenize{\mathcal M^{\mathrm{op}}}} & formula and labelled-word duals and opposite structure & P9 legacy unit 4 temporal-automata & meta & \hyperref[sec:p09-temporal-automata]{\ref*{sec:p09-temporal-automata}} & \thepage & Every order atom is reversed, with all Boolean and quantifier structure retained. Word reversal also reverses each elementary label as a formula, yielding an actual elementary presentation of the dual. \\
\phantomsection\label{entry:symbol-registry-01223}\(N_B^+,N_B^-,e_W^+,e_W^-\) & \texttt{\detokenize{N_B^+,}\allowbreak\detokenize{N_B^-,}\allowbreak\detokenize{e_W^+,}\allowbreak\detokenize{e_W^-}} & one-sided neighborhood and global unary accumulation labels & P9 legacy unit 4 temporal-automata & meta & \hyperref[sec:p09-temporal-automata]{\ref*{sec:p09-temporal-automata}} & \thepage & Neighborhood formulas require an actual earlier or later point. The global accumulation labels include the first/last-point vacuity cases explicitly and have proved E presentations. \\
\phantomsection\label{entry:symbol-registry-01224}\(B_W^+,B_W^-\) & \texttt{\detokenize{B_W^+,}\allowbreak\detokenize{B_W^-}} & remote cofinal and coinitial realization conditions & P9 legacy unit 4 temporal-automata & meta & \hyperref[sec:p09-temporal-automata]{\ref*{sec:p09-temporal-automata}} & \thepage & Their quantifiers include an endpoint instance forcing a realization, so adjacent outer endpoints do not make the conditions vacuously true. \\
\phantomsection\label{entry:symbol-registry-01225}\(p_{W,k}^-,p_{W,k}^+,V_{W,k}^+,V_{W,k}^-\) & \texttt{\detokenize{p_{W,k}^-,}\allowbreak\detokenize{p_{W,k}^+,}\allowbreak\detokenize{V_{W,k}^+,}\allowbreak\detokenize{V_{W,k}^-}} & limiting suffix and prefix labels and remote candidate words & P9 legacy unit 4 temporal-automata & meta & \hyperref[sec:p09-temporal-automata]{\ref*{sec:p09-temporal-automata}} & \thepage & All n candidate cuts are retained. The limiting one-letter and positive-stage cases are distinguished, and every label keeps its whole-structure meaning. \\
\phantomsection\label{entry:symbol-registry-01226}\(I_j\) & \texttt{\detokenize{I_j}} & cofinal-position condition for realizing sequences & P9 legacy unit 4 temporal-automata & meta & \hyperref[sec:p09-temporal-automata]{\ref*{sec:p09-temporal-automata}} & \thepage & The least cofinal position gives a common upper bound on the preceding positions, not a common prefix. The selected old gap and later witness gaps together cover the full stretched interval. \\
\phantomsection\label{entry:symbol-registry-01227}\(A_W^+\) & \texttt{\detokenize{A_W^+}} & fixed-initial pattern approaching an interior point from its right & P9 legacy unit 4 temporal-automata & meta & \hyperref[sec:p09-temporal-automata]{\ref*{sec:p09-temporal-automata}} & \thepage & Its elementary formula includes the no-future case and every successful cut. Order duality supplies the corresponding fixed-final construction from the left. \\
\phantomsection\label{entry:symbol-registry-01228}\(a,r,s,b,\eta,Q,G\) & \texttt{\detokenize{a,}\allowbreak\detokenize{r,}\allowbreak\detokenize{s,}\allowbreak\detokenize{b,}\allowbreak\detokenize{\eta,}\allowbreak\detokenize{Q,}\allowbreak\detokenize{G}} & localization endpoints elementary label placeholder and smaller family & P9 legacy unit 4 temporal-automata & meta & \hyperref[sec:p09-temporal-automata]{\ref*{sec:p09-temporal-automata}} & \thepage & Fixed prefix and suffix conditions remain outside the local universal formula. The fresh unary atom has length one; its elementary replacement is substituted only after applying the strictly smaller-family induction. \\
\phantomsection\label{entry:symbol-registry-01229}\(Y_i\) & \texttt{\detokenize{Y_i}} & positive supported branches for the full universal-family construction & P9 legacy unit 4 temporal-automata & meta & \hyperref[sec:p09-temporal-automata]{\ref*{sec:p09-temporal-automata}} & \thepage & Every exact boundary descriptor supplies its positive support, and each supported branch implies the target. All outside-region and unattained-boundary D tests and both necessary brackets are retained. \\
\phantomsection\label{entry:symbol-registry-01230}\(x_0,\ldots,x_m\) & \texttt{\detokenize{x_0,}\allowbreak\detokenize{\ldots,}\allowbreak\detokenize{x_m}} & ordered assembly vertices carrying point labels and interval edges & P9 legacy unit 4 temporal-automata & meta & \hyperref[sec:p09-temporal-automata]{\ref*{sec:p09-temporal-automata}} & \thepage & Every edge is cut at every intermediate named vertex before adjacent conditions are fused and concatenated. The existential assembly preserves all point labels and all supplied interval conditions. \\
\phantomsection\label{entry:symbol-registry-01231}\(\chi,\Phi,\Psi,U\) & \texttt{\detokenize{\chi,}\allowbreak\detokenize{\Phi,}\allowbreak\detokenize{\Psi,}\allowbreak\detokenize{U}} & finite formula families and their aggregate universal condition & P9 legacy unit 4 temporal-automata & meta & \hyperref[sec:p09-temporal-automata]{\ref*{sec:p09-temporal-automata}} & \thepage & The Lemma 9 construction retains all three families and guards every I representation; U is a local first-order aggregate, not the temporal until constructor. \\
\phantomsection\label{entry:symbol-registry-01232}\(\mathcal D(w),\phi_i(y_i),Q_i,\xi_i,\xi'\) & \texttt{\detokenize{\mathcal D(w),}\allowbreak\detokenize{\phi_i,}\allowbreak\detokenize{Q_i,}\allowbreak\detokenize{\xi_i,}\allowbreak\detokenize{\xi'}} & primitive regional formulas and maximal proper occurrence replacement & P9 legacy unit 4 temporal-automata & meta & \hyperref[sec:p09-temporal-automata]{\ref*{sec:p09-temporal-automata}} & \thepage & The source primitive class is not asserted closed under disjunction. Fresh predicates name whole-structure unary truths and are eliminated by the proved Com substitution argument. \\
\phantomsection\label{entry:symbol-registry-01233}\(\mathcal C,\Xi(w),\eta,D(\eta)\) & \texttt{\detokenize{\mathcal C,}\allowbreak\detokenize{\Xi,}\allowbreak\detokenize{\eta,}\allowbreak\detokenize{D(\eta)}} & literal temporal image class and the completed one-variable translation & P9 legacy unit 4 temporal-automata & meta & \hyperref[sec:p09-temporal-automata]{\ref*{sec:p09-temporal-automata}} & \thepage & The final output is a temporal TL\_1 formula and its literal D-image. Arbitrary intermediate elementary formulas are not identified with members of that image. \\
\phantomsection\label{entry:symbol-registry-01234}\(Q,R,q,r,\tau\) & \texttt{\detokenize{Q,}\allowbreak\detokenize{R,}\allowbreak\detokenize{q,}\allowbreak\detokenize{r,}\allowbreak\detokenize{\tau}} & unary and corresponding propositional letters and a tautology in the regional example & P9 legacy unit 4 temporal-automata & meta & \hyperref[sec:p09-temporal-automata]{\ref*{sec:p09-temporal-automata}} & \thepage & The example explicitly treats the first-point vacuity of a universal past condition. Temporary predicate letters do not survive the final substitution. \\
\phantomsection\label{entry:symbol-registry-01266}\(\#a,w,\operatorname{cons},\operatorname{nil}\) & \texttt{\detokenize{\#a,}\allowbreak\detokenize{w,}\allowbreak\detokenize{\operatorname{cons},}\allowbreak\detokenize{\operatorname{nil}}} & numerical syntax codes token weight and list constructors & P9 legacy unit 4 temporal-automata & meta & \hyperref[sec:p09-temporal-automata]{\ref*{sec:p09-temporal-automata}} & \thepage & The quantitative Kamp construction defines prefix and frame tags, parsing, malformed-input defaults and all bounded arithmetic routines explicitly. \\
\phantomsection\label{entry:symbol-registry-01267}\(R,A,f^{[k]},\mu,b_i,B,h,c,z,s_*,u_i,L_8,L\) & \texttt{\detokenize{R,}\allowbreak\detokenize{A,}\allowbreak\detokenize{f^{[k]},}\allowbreak\detokenize{\mu,}\allowbreak\detokenize{b_i,}\allowbreak\detokenize{B,}\allowbreak\detokenize{h,}\allowbreak\detokenize{c,}\allowbreak\detokenize{z,}\allowbreak\detokenize{s_*,}\allowbreak\detokenize{u_i,}\allowbreak\detokenize{L_8,}\allowbreak\detokenize{L}} & local expansion envelopes structural family mass and complete recursion bounds & P9 legacy unit 4 temporal-automata & meta & \hyperref[sec:p09-temporal-automata]{\ref*{sec:p09-temporal-automata}} & \thepage & The mass counts every retained word and At tree. Level bounds control whole-tree size and returned expansion, rather than only a decreasing branch. \\
\phantomsection\label{entry:symbol-registry-01268}\(J,S,C,V,W,n_i,T,\operatorname{Exp}\) & \texttt{\detokenize{J,}\allowbreak\detokenize{S,}\allowbreak\detokenize{C,}\allowbreak\detokenize{V,}\allowbreak\detokenize{W,}\allowbreak\detokenize{n_i,}\allowbreak\detokenize{T,}\allowbreak\detokenize{\operatorname{Exp}}} & outer compilation envelopes final clock and total numerical translation & P9 legacy unit 4 temporal-automata & meta & \hyperref[sec:p09-temporal-automata]{\ref*{sec:p09-temporal-automata}} & \thepage & These are mathematical primitive-recursive bounds, not practical runtime estimates. Positive-arity schema padding retains all required letters; the nullary exception is separately proved. \\
\phantomsection\label{entry:symbol-registry-01318}\(\mathcal R,\mathfrak M_{\mathcal R},\mathsf{TM}_{\mathcal R},\llbracket\eta\rrbracket_V\) & \texttt{\detokenize{\mathcal R,}\allowbreak\detokenize{\mathfrak M_{\mathcal R},}\allowbreak\detokenize{\mathsf{TM}_{\mathcal R},}\allowbreak\detokenize{\llbracket\eta\rrbracket_V}} & real order all source-definable unary operations finite sentential language and truth-set interpretation & P9 legacy unit 4 temporal-automata & meta & \hyperref[sec:p09-temporal-automata]{\ref*{sec:p09-temporal-automata}} & \thepage & The unary operations have no extra real or predicate parameters. Every occurring seed predicate is restricted as stated by the corresponding field theorem. \\
\phantomsection\label{entry:symbol-registry-01319}\(\operatorname{DU}(A),\mathcal F(A),\mathcal B(A_1,\ldots,A_n),A_I,E,H,d_-,d_+\) & \texttt{\detokenize{\operatorname{DU},}\allowbreak\detokenize{\mathcal F,}\allowbreak\detokenize{\mathcal B,}\allowbreak\detokenize{A_I,}\allowbreak\detokenize{E,}\allowbreak\detokenize{H,}\allowbreak\detokenize{d_-,}\allowbreak\detokenize{d_+}} & source discrete-unbounded condition four-choice family finite generated field even/odd sets and nearest points & P9 legacy unit 4 temporal-automata & meta & \hyperref[sec:p09-temporal-automata]{\ref*{sec:p09-temporal-automata}} & \thepage & The empty proposition is treated separately. Discreteness is required at every real point, and both nearest-point directions are proved before use. \\
\phantomsection\label{entry:symbol-registry-01320}\(J,\mathsf{TM}_J,I,d_n,\mathcal B_I,h\) & \texttt{\detokenize{J,}\allowbreak\detokenize{\mathsf{TM}_J,}\allowbreak\detokenize{I,}\allowbreak\detokenize{d_n,}\allowbreak\detokenize{\mathcal B_I,}\allowbreak\detokenize{h}} & dense order full unary language convex interval markers relative field and finite partial-map index shift & P9 legacy unit 4 temporal-automata & meta & \hyperref[sec:p09-temporal-automata]{\ref*{sec:p09-temporal-automata}} & \thepage & The dense-order argument permits endpoints and noncomplete orders and fixes the selected exterior points while mapping inside the convex interval. \\
\phantomsection\label{entry:symbol-registry-01321}\(\mathsf{In},\pi_-,\pi_+,\iota_-,\iota_+,\mathcal S,\mathsf P_1,\mathsf F_1,\mathsf S_1,\mathsf U_1\) & \texttt{\detokenize{\mathsf{In},}\allowbreak\detokenize{\pi_-,}\allowbreak\detokenize{\pi_+,}\allowbreak\detokenize{\iota_-,}\allowbreak\detokenize{\iota_+,}\allowbreak\detokenize{\mathcal S,}\allowbreak\detokenize{\mathsf P_1,}\allowbreak\detokenize{\mathsf F_1,}\allowbreak\detokenize{\mathsf S_1,}\allowbreak\detokenize{\mathsf U_1}} & integer order defining formulas six-operation basis neighbor and isolation operations & P9 legacy unit 4 temporal-automata & meta & \hyperref[sec:p09-temporal-automata]{\ref*{sec:p09-temporal-automata}} & \thepage & The subscript-one temporal symbols are unary. They are distinct from the binary strict SINCE/UNTIL operations and from first-order predicate letters. \\
\phantomsection\label{entry:symbol-registry-01322}\(C,B_P,B_F,\eta_1,\eta_2,\delta,N_\delta,J_\delta,\eta_\delta,G_F,G_P,\zeta_1,\zeta_2\) & \texttt{\detokenize{C,}\allowbreak\detokenize{B_P,}\allowbreak\detokenize{B_F,}\allowbreak\detokenize{\eta_1,}\allowbreak\detokenize{\eta_2,}\allowbreak\detokenize{\delta,}\allowbreak\detokenize{N_\delta,}\allowbreak\detokenize{J_\delta,}\allowbreak\detokenize{\eta_\delta,}\allowbreak\detokenize{G_F,}\allowbreak\detokenize{G_P,}\allowbreak\detokenize{\zeta_1,}\allowbreak\detokenize{\zeta_2}} & continuation and failure formulas corrected markers directional abbreviations source guards and integer target formulas & P9 legacy unit 4 temporal-automata & meta & \hyperref[sec:p09-temporal-automata]{\ref*{sec:p09-temporal-automata}} & \thepage & The future search uses the backward-inspecting marker. Both target expressions retain every source disjunct, with the proved neighbor shifts and no current-point substitution. \\
\phantomsection\label{entry:symbol-registry-01323}\(\mathcal S'',R_r,V,W,\mathcal B_4\) & \texttt{\detokenize{\mathcal S'',}\allowbreak\detokenize{R_r,}\allowbreak\detokenize{V,}\allowbreak\detokenize{W,}\allowbreak\detokenize{\mathcal B_4}} & restricted integer basis modulo-four residues even and odd sets and the eight-set field & P9 legacy unit 4 temporal-automata & meta & \hyperref[sec:p09-temporal-automata]{\ref*{sec:p09-temporal-automata}} & \thepage & The family includes the complement of residue zero modulo four. The complete operation table and formula induction prove irredundancy of immediate past and future. \\
\phantomsection\label{entry:symbol-registry-01324}\(\mathsf{Ra},(a,b)_{\mathbb Q},\alpha,b_n,Q,Q',\mathcal M,\mathcal M',\beta,\varphi\) & \texttt{\detokenize{\mathsf{Ra},}\allowbreak\detokenize{(a,}\allowbreak\detokenize{b)_{\mathbb Q},}\allowbreak\detokenize{\alpha,}\allowbreak\detokenize{b_n,}\allowbreak\detokenize{Q,}\allowbreak\detokenize{Q',}\allowbreak\detokenize{\mathcal M,}\allowbreak\detokenize{\mathcal M',}\allowbreak\detokenize{\beta,}\allowbreak\detokenize{\varphi}} & rational order rational intervals irrational cut reciprocal endpoints interpretations and distinguishing formulas & P9 legacy unit 4 temporal-automata & meta & \hyperref[sec:p09-temporal-automata]{\ref*{sec:p09-temporal-automata}} & \thepage & The inserted bands have ordered endpoints. Beta detects the accumulating obstruction with rational witnesses and a nonvacuity clause; its negation gives the required separation at current point three. \\
\phantomsection\label{entry:symbol-registry-01346}\(\mathsf{TL}_0,\mathsf P,\mathsf F,\rho_0,\rho,\mathsf{PR}_0,\mathsf{PR},T_*,x^-,x^+\) & \texttt{\detokenize{\mathsf{TL}_0,}\allowbreak\detokenize{\mathsf P,}\allowbreak\detokenize{\mathsf F,}\allowbreak\detokenize{\rho_0,}\allowbreak\detokenize{\rho,}\allowbreak\detokenize{\mathsf{PR}_0,}\allowbreak\detokenize{\mathsf{PR},}\allowbreak\detokenize{T_*,}\allowbreak\detokenize{x^-,}\allowbreak\detokenize{x^+}} & strict past/future language progressive proposals two-sided points and neighbors & P9 legacy unit 4 temporal-automata & meta & \hyperref[sec:p09-temporal-automata]{\ref*{sec:p09-temporal-automata}} & \thepage & The revised progressive includes present truth and both endpoint predicates. Nonexpressibility concerns the smaller past/future basis. \\
\phantomsection\label{entry:symbol-registry-01347}\(d_-,d_+,a_n,b_n,I,I_+,I_-,H_m^+,H_m^-,A,A_i^+,A_i^-,A_{ij},C_{mn},J^{op},\eta^*,N\) & \texttt{\detokenize{d_-,}\allowbreak\detokenize{d_+,}\allowbreak\detokenize{a_n,}\allowbreak\detokenize{b_n,}\allowbreak\detokenize{I,}\allowbreak\detokenize{I_+,}\allowbreak\detokenize{I_-,}\allowbreak\detokenize{H_m^+,}\allowbreak\detokenize{H_m^-,}\allowbreak\detokenize{A,}\allowbreak\detokenize{A_i^+,}\allowbreak\detokenize{A_i^-,}\allowbreak\detokenize{A_{ij},}\allowbreak\detokenize{C_{mn},}\allowbreak\detokenize{J^{op},}\allowbreak\detokenize{\eta^*,}\allowbreak\detokenize{N}} & directional depths markers convex domains tails inputs windows order reversal syntax dual and finite bound & P9 legacy unit 4 temporal-automata & meta & \hyperref[sec:p09-temporal-automata]{\ref*{sec:p09-temporal-automata}} & \thepage & The general infinite-order proof selects two finite chains. No cofinality or two infinite sequence orientations are assumed. \\
\phantomsection\label{entry:symbol-registry-01348}\(\operatorname{At}_G,\operatorname{Bet}_{G,n},\operatorname{Bef}_{G,n},\operatorname{Aft}_{G,n},g_{ab},B_j,y_j,\pi\) & \texttt{\detokenize{\operatorname{At}_G,}\allowbreak\detokenize{\operatorname{Bet}_{G,n},}\allowbreak\detokenize{\operatorname{Bef}_{G,n},}\allowbreak\detokenize{\operatorname{Aft}_{G,n},}\allowbreak\detokenize{g_{ab},}\allowbreak\detokenize{B_j,}\allowbreak\detokenize{y_j,}\allowbreak\detokenize{\pi}} & GPSS signed point and regional grammars endpoint guards equality blocks representatives and order type & P9 legacy unit 4 temporal-automata & meta & \hyperref[sec:p09-temporal-automata]{\ref*{sec:p09-temporal-automata}} & \thepage & These syntactic classes differ from Kamp's normal-form grammar. Formal infinite endpoints are not model points. \\
\phantomsection\label{entry:symbol-registry-01349}\((\cdot)^*,\mathsf U_s,F,G,\delta,z_i,\theta_i,\varphi_i,\operatorname{Inf},H_i,D_p\) & \texttt{\detokenize{(\cdot)^*,}\allowbreak\detokenize{\mathsf U_s,}\allowbreak\detokenize{F,}\allowbreak\detokenize{G,}\allowbreak\detokenize{\delta,}\allowbreak\detokenize{z_i,}\allowbreak\detokenize{\theta_i,}\allowbreak\detokenize{\varphi_i,}\allowbreak\detokenize{\operatorname{Inf},}\allowbreak\detokenize{H_i,}\allowbreak\detokenize{D_p}} & future translation decomposition witnesses point and gap labels cofinal condition suffix histories and merge example & P9 legacy unit 4 temporal-automata & meta & \hyperref[sec:p09-temporal-automata]{\ref*{sec:p09-temporal-automata}} & \thepage & Labels are future-only L(U) translations. Cofinality translates as F psi AND G F psi. \\
\phantomsection\label{entry:symbol-registry-01367}\(W,r,\operatorname{First}_A,\operatorname{Last}_A,\operatorname{Test}_A,\operatorname{Def},\Lambda_i,\mu_i,V_m,V^j,W_j,\ell,T_W,E_W,B(r),C_i\) & \texttt{\detokenize{W,}\allowbreak\detokenize{r,}\allowbreak\detokenize{\operatorname{First}_A,}\allowbreak\detokenize{\operatorname{Last}_A,}\allowbreak\detokenize{\operatorname{Test}_A,}\allowbreak\detokenize{\operatorname{Def},}\allowbreak\detokenize{\Lambda_i,}\allowbreak\detokenize{\mu_i,}\allowbreak\detokenize{V_m,}\allowbreak\detokenize{V^j,}\allowbreak\detokenize{W_j,}\allowbreak\detokenize{\ell,}\allowbreak\detokenize{T_W,}\allowbreak\detokenize{E_W,}\allowbreak\detokenize{B(r),}\allowbreak\detokenize{C_i}} & GPSS word rank partial event selectors statuses future continuation labels canonical values determining decision lists bounds and compiled cases & P9 legacy unit 4 temporal-automata & meta & \hyperref[sec:p09-temporal-automata]{\ref*{sec:p09-temporal-automata}} & \thepage & Undefined is not a model point. All auxiliary events are eliminated; the node bound does not bound expanded label syntax. \\
\phantomsection\label{entry:symbol-registry-01368}\(e_{\mathrm{last}},L_E,w,C_n,I\) & \texttt{\detokenize{e_{\mathrm{last}},}\allowbreak\detokenize{L_E,}\allowbreak\detokenize{w,}\allowbreak\detokenize{C_n,}\allowbreak\detokenize{I}} & last-point truth value unfolded successor label binary interpretation closed cylinders and infinitely-many-failures class & P9 legacy unit 4 temporal-automata & meta & \hyperref[sec:p09-temporal-automata]{\ref*{sec:p09-temporal-automata}} & \thepage & The formal-left-endpoint countertheorem holds in the unchanged future-label decomposition grammar. It does not refute the actual-lower-endpoint discrete future theorem. \\
\phantomsection\label{entry:symbol-registry-01379}\(N^+,\partial^+,K^+,E_W^+,\Gamma_W,W^-,\alpha,\operatorname{Inf}_A(s,t),\operatorname{Test}_A(s)\) & \texttt{\detokenize{N^+,}\allowbreak\detokenize{\partial^+,}\allowbreak\detokenize{K^+,}\allowbreak\detokenize{E_W^+,}\allowbreak\detokenize{\Gamma_W,}\allowbreak\detokenize{W^-,}\allowbreak\detokenize{\alpha,}\allowbreak\detokenize{\operatorname{Inf}_A,}\allowbreak\detokenize{\operatorname{Test}_A}} & right neighborhood vacuous and nonvacuous limits word realization accumulation suffix lower-endpoint label and partial infimum events & P9 legacy unit 4 temporal-automata & meta & \hyperref[sec:p09-temporal-automata]{\ref*{sec:p09-temporal-automata}} & \thepage & These finite formulas and events retain their local proven scopes. They do not assert U-only future completeness on arbitrary complete orders. \\
\phantomsection\label{entry:symbol-registry-01380}\(V_{\Theta;\Phi},c,h,j,k,R,D_A,\operatorname{Adj},\mathrm{FIN},\mathrm{INF}\) & \texttt{\detokenize{V_{\Theta;\Phi},}\allowbreak\detokenize{c,}\allowbreak\detokenize{h,}\allowbreak\detokenize{j,}\allowbreak\detokenize{k,}\allowbreak\detokenize{R,}\allowbreak\detokenize{D_A,}\allowbreak\detokenize{\operatorname{Adj},}\allowbreak\detokenize{\mathrm{FIN},}\allowbreak\detokenize{\mathrm{INF}}} & bounded positive-chain tests infimum values recursion indices word decisions initial-gap formula adjacency and binary interpretation classes & P9 legacy unit 4 temporal-automata & meta & \hyperref[sec:p09-temporal-automata]{\ref*{sec:p09-temporal-automata}} & \thepage & The positive-chain recursion excludes the general final-gap condition. The endpoint-free bounded class has a proved complete-order obstruction. \\
\phantomsection\label{entry:symbol-registry-01395}\(\mathsf{FTL}_d,O_k,\alpha,\beta,\rho,\equiv_n,L_n(x,t),\phi_n,\phi_n^P\) & \texttt{\detokenize{\mathsf{FTL}_d,}\allowbreak\detokenize{O_k,}\allowbreak\detokenize{\alpha,}\allowbreak\detokenize{\beta,}\allowbreak\detokenize{\rho,}\allowbreak\detokenize{\equiv_n,}\allowbreak\detokenize{L_n,}\allowbreak\detokenize{\phi_n,}\allowbreak\detokenize{\phi_n^P}} & bounded-depth future truth-table language ordinal orders padding game equivalence relativized limit hierarchy and separators & P9 legacy unit 4 temporal-automata & meta & \hyperref[sec:p09-temporal-automata]{\ref*{sec:p09-temporal-automata}} & \thepage & The depth bound applies to each modality truth table, not to temporal nesting. The reference point and positive limit-level witness remove initial-point vacuity. \\
\phantomsection\label{entry:symbol-registry-01396}\(C_k,P_k,K^-,\mathcal A,\top,\bot,P,\neg P\) & \texttt{\detokenize{C_k,}\allowbreak\detokenize{P_k,}\allowbreak\detokenize{K^-,}\allowbreak\detokenize{\mathcal A,}\allowbreak\detokenize{\top,}\allowbreak\detokenize{\bot,}\allowbreak\detokenize{P,}\allowbreak\detokenize{\neg P}} & marked real models past-limit operation full and guarded languages and four invariant truth sets & P9 legacy unit 4 temporal-automata & meta & \hyperref[sec:p09-temporal-automata]{\ref*{sec:p09-temporal-automata}} & \thepage & The past limit is vacuously true at a first point. Its guarded fragment has future semantics even though an unguarded past-limit formula need not. \\
\phantomsection\label{entry:symbol-registry-01397}\(\mathcal A,[W](a,b),Q_W,\operatorname{All}_B,\operatorname{UC}_W,S_{(a,b)},\Sigma_W\) & \texttt{\detokenize{\mathcal A,}\allowbreak\detokenize{[W],}\allowbreak\detokenize{Q_W,}\allowbreak\detokenize{\operatorname{All}_B,}\allowbreak\detokenize{\operatorname{UC}_W,}\allowbreak\detokenize{S_{(a,b)},}\allowbreak\detokenize{\Sigma_W}} & canonical guarded-basis labels interval words positive-chain decision uniform-gap handler failure set and retained boundary descriptor & P9 legacy unit 4 temporal-automata & meta & \hyperref[sec:p09-temporal-automata]{\ref*{sec:p09-temporal-automata}} & \thepage & Both actual endpoints have labels. A boundary variable is retained across all conjuncts and is never replaced independently by another witness. \\
\phantomsection\label{entry:symbol-registry-01398}\(P_i,S_i,L_i,R_i,N_W,E_1,E_2,E_3,D_i,H_i,M_W,T_i,g\) & \texttt{\detokenize{P_i,}\allowbreak\detokenize{S_i,}\allowbreak\detokenize{L_i,}\allowbreak\detokenize{R_i,}\allowbreak\detokenize{N_W,}\allowbreak\detokenize{E_1,}\allowbreak\detokenize{E_2,}\allowbreak\detokenize{E_3,}\allowbreak\detokenize{D_i,}\allowbreak\detokenize{H_i,}\allowbreak\detokenize{M_W,}\allowbreak\detokenize{T_i,}\allowbreak\detokenize{g}} & point and gap cut words complement recursion branches compiled lists ray formulas and initial guard normalizer & P9 legacy unit 4 temporal-automata & meta & \hyperref[sec:p09-temporal-automata]{\ref*{sec:p09-temporal-automata}} & \thepage & All ordinary calls strictly decrease interval count. Extreme same-length cases use the separate uniform-gap handler or are proved impossible. \\
\phantomsection\label{entry:symbol-registry-01405}\(L(S),L(U),A^*(x),(\cdot)^d,\mathcal M^{op},T_F^{dc},T_P^{dc},\operatorname{Sep}_{dc},C_i,\beta_i,\gamma_i\) & \texttt{\detokenize{L(S),}\allowbreak\detokenize{L(U),}\allowbreak\detokenize{A^*,}\allowbreak\detokenize{(\cdot)^d,}\allowbreak\detokenize{\mathcal M^{op},}\allowbreak\detokenize{T_F^{dc},}\allowbreak\detokenize{T_P^{dc},}\allowbreak\detokenize{\operatorname{Sep}_{dc},}\allowbreak\detokenize{C_i,}\allowbreak\detokenize{\beta_i,}\allowbreak\detokenize{\gamma_i}} & pure temporal fragments first-order translation typed reversals opposite model discrete translators separation and regional clauses & P9 legacy unit 4 temporal-automata & meta & \hyperref[sec:p09-temporal-automata]{\ref*{sec:p09-temporal-automata}} & \thepage & Pure is an operator restriction with current-point atoms allowed. Reversal exchanges U and S without swapping their infix operand positions. \\
\phantomsection\label{entry:symbol-registry-01406}\(\mathcal A_F,\mathcal A_P,K_v^+,T_F^{com},T_P^{com},\Xi\) & \texttt{\detokenize{\mathcal A_F,}\allowbreak\detokenize{\mathcal A_P,}\allowbreak\detokenize{K_v^+,}\allowbreak\detokenize{T_F^{com},}\allowbreak\detokenize{T_P^{com},}\allowbreak\detokenize{\Xi}} & guarded future and past languages vacuous right limit complete-order translators and mixed separation counterformula & P9 legacy unit 4 temporal-automata & meta & \hyperref[sec:p09-temporal-automata]{\ref*{sec:p09-temporal-automata}} & \thepage & The vacuous right limit is true at a last point. Guarded separation on complete orders is distinct from pure-S/pure-U separation, which the real countermodels refute. \\
\phantomsection\label{entry:symbol-registry-01425}\((L,R),g,\gamma^\pm,\Gamma_P,\operatorname{Run}_P,\mathsf U_{\mathrm{gap}},\mathsf S_{\mathrm{gap}},\mathcal T_{\mathrm{gap}}\) & \texttt{\detokenize{(L,}\allowbreak\detokenize{R),}\allowbreak\detokenize{g,}\allowbreak\detokenize{\gamma^\pm,}\allowbreak\detokenize{\Gamma_P,}\allowbreak\detokenize{\operatorname{Run}_P,}\allowbreak\detokenize{\mathsf U_{\mathrm{gap}},}\allowbreak\detokenize{\mathsf S_{\mathrm{gap}},}\allowbreak\detokenize{\mathcal T_{\mathrm{gap}}}} & proper gap point-gap comparisons unique run gap formula gap temporal operators and full four-operator language & P9 legacy unit 4 temporal-automata & meta & \hyperref[sec:p09-temporal-automata]{\ref*{sec:p09-temporal-automata}} & \thepage & A proper gap has no point at which a label could be evaluated. The defined 2018 gap operators are distinct from ordinary strict until and since and from undeclared 1980 symbols. \\
\phantomsection\label{entry:symbol-registry-01426}\(\operatorname{Part},O,U_*,F_i,UG,OG_D,CG_D\) & \texttt{\detokenize{\operatorname{Part},}\allowbreak\detokenize{O,}\allowbreak\detokenize{U_*,}\allowbreak\detokenize{F_i,}\allowbreak\detokenize{UG,}\allowbreak\detokenize{OG_D,}\allowbreak\detokenize{CG_D}} & labelled convex partition singleton flags two-cell operator prefix recurrence common-gap guard and open and closed gap-capped translations & P9 legacy unit 4 temporal-automata & meta & \hyperref[sec:p09-temporal-automata]{\ref*{sec:p09-temporal-automata}} & \thepage & Conjoined gap-side formulas retain one common missing cut. A final cell approaching a proper gap cannot be a forced singleton. \\
\phantomsection\label{entry:symbol-registry-01427}\(E,\operatorname{Occ}_B,\operatorname{All}_B,W,\Pi_W,\Pi_{i:j},\operatorname{Inf}_B,\operatorname{Sup}_B,A_m,O_m,L_W,R_W\) & \texttt{\detokenize{E,}\allowbreak\detokenize{\operatorname{Occ}_B,}\allowbreak\detokenize{\operatorname{All}_B,}\allowbreak\detokenize{W,}\allowbreak\detokenize{\Pi_W,}\allowbreak\detokenize{\Pi_{i:j},}\allowbreak\detokenize{\operatorname{Inf}_B,}\allowbreak\detokenize{\operatorname{Sup}_B,}\allowbreak\detokenize{A_m,}\allowbreak\detokenize{O_m,}\allowbreak\detokenize{L_W,}\allowbreak\detokenize{R_W}} & emptiness occurrence uniform label flagged partition word shifted subword boundary descriptors plain chain and its complement and independent prefix and suffix handlers & P9 legacy unit 4 temporal-automata & meta & \hyperref[sec:p09-temporal-automata]{\ref*{sec:p09-temporal-automata}} & \thepage & The uniform handlers are proved independently from plain-chain complement. Forced first and last singleton cells are resolved before uniform-cell extension. \\
\phantomsection\label{entry:symbol-registry-01428}\(A_i,M_W,N_W,I_W,J_W,G_W,V_W,B_W,T_W,Q_W,\lambda_i,\rho_j,Z_W\) & \texttt{\detokenize{A_i,}\allowbreak\detokenize{M_W,}\allowbreak\detokenize{N_W,}\allowbreak\detokenize{I_W,}\allowbreak\detokenize{J_W,}\allowbreak\detokenize{G_W,}\allowbreak\detokenize{V_W,}\allowbreak\detokenize{B_W,}\allowbreak\detokenize{T_W,}\allowbreak\detokenize{Q_W,}\allowbreak\detokenize{\lambda_i,}\allowbreak\detokenize{\rho_j,}\allowbreak\detokenize{Z_W}} & fixed-point partition cut interior condition guarded complement actual and missing boundary descriptors final-label common-gap test gap-side partitions and full cut complement & P9 legacy unit 4 temporal-automata & meta & \hyperref[sec:p09-temporal-automata]{\ref*{sec:p09-temporal-automata}} & \thepage & Every recursive partition call decreases cell count. Equal-gap cuts include within-cell crossings only for unflagged cells and use the actual final label. \\
\phantomsection\label{entry:symbol-registry-03632}\(S,T\) & \texttt{\detokenize{S;T}} & disjoint finite similarity types & P9 Fagin 1974 generalized spectra & n/a & \hyperref[def:p09-fagin74-spectrum]{\ref*{def:p09-fagin74-spectrum}} & \thepage & Theorem 6 restricts the input type S to positive-arity relations; the auxiliary type T may include functions before graph conversion. \\
\phantomsection\label{entry:symbol-registry-03633}\(\operatorname{Fin}(S)\) & \texttt{\detokenize{\operatorname{Fin}}} & class of nonempty finite structures of the specified type & P9 Fagin 1974 generalized spectra & n/a & \hyperref[def:p09-fagin74-spectrum]{\ref*{def:p09-fagin74-spectrum}} & \thepage & Empty vocabulary is allowed; empty universe is not. \\
\phantomsection\label{entry:symbol-registry-03634}\(\operatorname{Mod}_S\exists T\sigma\) & \texttt{\detokenize{\operatorname{Mod};\exists T}} & finite reduct class admitting the specified same-universe expansion & P9 Fagin 1974 generalized spectra & n/a & \hyperref[def:p09-fagin74-spectrum]{\ref*{def:p09-fagin74-spectrum}} & \thepage & The second-order existential prefix binds all auxiliary relation and function symbols. \\
\phantomsection\label{entry:symbol-registry-03635}\(P_i,m_i,b_i\) & \texttt{\detokenize{P_i;m_i;b_i}} & input relation its positive arity and characteristic word & P9 Fagin 1974 generalized spectra & n/a & \hyperref[def:p09-fagin74-encoding]{\ref*{def:p09-fagin74-encoding}} & \thepage & Both type order and tuple lexicographic order are fixed by the encoding. \\
\phantomsection\label{entry:symbol-registry-03636}\(e(\mathcal A),E(\mathcal C)\) & \texttt{\detokenize{e;E}} & original encoding and encoded language of a finite-structure class & P9 Fagin 1974 generalized spectra & n/a & \hyperref[def:p09-fagin74-encoding]{\ref*{def:p09-fagin74-encoding}} & \thepage & Binary size prefix and relation-table words; this is not an arbitrary or succinct encoding. \\
\phantomsection\label{entry:symbol-registry-03637}\(n,\ell\) & \texttt{\detokenize{n;\ell}} & universe cardinality and encoded word length & P9 Fagin 1974 generalized spectra & n/a & \hyperref[def:p09-fagin74-encoding]{\ref*{def:p09-fagin74-encoding}} & \thepage & They are polynomially related only when the fixed positive-arity input type is nonempty. \\
\phantomsection\label{entry:symbol-registry-03638}\(NP,NP_1\) & \texttt{\detokenize{NP;NP_1}} & original polynomial and single-exponential nondeterministic time classes & P9 Fagin 1974 generalized spectra & n/a & \hyperref[def:p09-fagin74-encoding]{\ref*{def:p09-fagin74-encoding}} & \thepage & NP\_1 uses some time bound 2 raised to k times input length; it is not NP on the original binary numerical input. \\
\phantomsection\label{entry:symbol-registry-03639}\(F_f,V_t\) & \texttt{\detokenize{F_f;V_t}} & auxiliary function graph and recursive term-value formula & P9 Fagin 1974 generalized spectra & n/a & \hyperref[annex:p9.n-c-004.1.theorem6]{\ref*{annex:p9.n-c-004.1.theorem6}} & \thepage & Graph coding enforces totality and single-valuedness; constants use unary singleton graphs. \\
\phantomsection\label{entry:symbol-registry-03640}\(F,f\) & \texttt{\detokenize{F;f}} & pair relation and fixed-point-free involution witness & P9 Fagin 1974 generalized spectra & n/a & \hyperref[def:p09-fagin74-even-witness]{\ref*{def:p09-fagin74-even-witness}} & \thepage & F is a symmetric irreflexive total single-valued relation; f is its corresponding unary function. \\
\phantomsection\label{entry:symbol-registry-03641}\(<,c_1,c_2\) & \texttt{\detokenize{<;c_1;c_2}} & existentially chosen strict linear order and its endpoints & P9 Fagin 1974 tableau & n/a & \hyperref[annex:p9.n-c-004.1.theorem6]{\ref*{annex:p9.n-c-004.1.theorem6}} & \thepage & This order is an auxiliary witness and is not required in the input. \\
\phantomsection\label{entry:symbol-registry-03642}\(S\) & \texttt{\detokenize{S}} & source cyclic successor function of the auxiliary order & P9 Fagin 1974 tableau & n/a & \hyperref[annex:p9.n-c-004.1.theorem6]{\ref*{annex:p9.n-c-004.1.theorem6}} & \thepage & Typed function role distinct from source S denoting the fixed similarity type; last element maps to first. \\
\phantomsection\label{entry:symbol-registry-03643}\(S_1\) & \texttt{\detokenize{S_1}} & source cyclic lexicographic successor relation on k-tuples & P9 Fagin 1974 tableau & relation & \hyperref[annex:p9.n-c-004.1.theorem6]{\ref*{annex:p9.n-c-004.1.theorem6}} & \thepage & Its final-to-initial pair is excluded for time transitions and guarded against at tape boundaries. \\
\phantomsection\label{entry:symbol-registry-03644}\(\mathbf a,\mathbf z,N\) & \texttt{\detokenize{\mathbf a;\mathbf z;N}} & first and last tuple addresses and tuple-address count & P9 Fagin 1974 tableau & n/a & \hyperref[annex:p9.n-c-004.1.theorem6]{\ref*{annex:p9.n-c-004.1.theorem6}} & \thepage & N equals n raised to the fixed tableau exponent k; slack separates the accepted run from the last address. \\
\phantomsection\label{entry:symbol-registry-03645}\(D(\bar t,\bar u)\) & \texttt{\detokenize{D}} & nonwrapping consecutive-time predicate & P9 Fagin 1974 tableau & relation & \hyperref[annex:p9.n-c-004.1.theorem6]{\ref*{annex:p9.n-c-004.1.theorem6}} & \thepage & Abbreviates tuple successor and the exclusion of the final time. \\
\phantomsection\label{entry:symbol-registry-03646}\(K,\Gamma,\delta\) & \texttt{\detokenize{K;\Gamma;\delta}} & fixed machine state set tape alphabet and finite transition table & P9 Fagin 1974 tableau & n/a & \hyperref[annex:p9.n-c-004.1.theorem6]{\ref*{annex:p9.n-c-004.1.theorem6}} & \thepage & Machine moves select one whole common transition for all tapes. \\
\phantomsection\label{entry:symbol-registry-03647}\(q,v_i,H_i\) & \texttt{\detokenize{q;v_i;H_i}} & source state function tape-symbol function and unique-head relation & P9 Fagin 1974 tableau & n/a & \hyperref[annex:p9.n-c-004.1.theorem6]{\ref*{annex:p9.n-c-004.1.theorem6}} & \thepage & Arity k for q and 2k for the tape and head symbols; state and alphabet range constraints are explicit. \\
\phantomsection\label{entry:symbol-registry-03648}\(\kappa_b,\gamma_e\) & \texttt{\detokenize{\kappa_b;\gamma_e}} & successor-term codes of a state or tape symbol & P9 Fagin 1974 tableau & n/a & \hyperref[annex:p9.n-c-004.1.theorem6]{\ref*{annex:p9.n-c-004.1.theorem6}} & \thepage & Code lists are individually injective but may overlap across their different typed roles. \\
\phantomsection\label{entry:symbol-registry-03649}\(0,1,B\) & \texttt{\detokenize{0;1;B}} & source tape-symbol codes for zero one and blank & P9 Fagin 1974 tableau & n/a & \hyperref[annex:p9.n-c-004.1.theorem6]{\ref*{annex:p9.n-c-004.1.theorem6}} & \thepage & These stand for the first three auxiliary order elements rather than arithmetic constants in the input structure. \\
\phantomsection\label{entry:symbol-registry-03650}\(q_0,q_A,q_R\) & \texttt{\detokenize{q_0;q_A;q_R}} & initial accepting and rejecting state codes & P9 Fagin 1974 tableau & n/a & \hyperref[annex:p9.n-c-004.1.theorem6]{\ref*{annex:p9.n-c-004.1.theorem6}} & \thepage & Final configurations are frozen through the remaining time rows. \\
\phantomsection\label{entry:symbol-registry-03651}\(G(x,y),C(x,y)\) & \texttt{\detokenize{G;C}} & source binary-digit function and explicit reconstruction carry condition & P9 Fagin 1974 tableau & n/a & \hyperref[annex:p9.n-c-004.1.theorem6]{\ref*{annex:p9.n-c-004.1.theorem6}} & \thepage & y runs from least significant bit upward; C says all lower positions are one. \\
\phantomsection\label{entry:symbol-registry-03652}\(\Phi_s\) & \texttt{\detokenize{\Phi_s}} & first-order formula for one complete allowed machine move & P9 Fagin 1974 tableau & n/a & \hyperref[annex:p9.n-c-004.1.theorem6]{\ref*{annex:p9.n-c-004.1.theorem6}} & \thepage & State update writes unchanged cells and both head-boundary guards are included. \\
\phantomsection\label{entry:symbol-registry-03653}\(\chi,\Delta_{\mathcal D}\) & \texttt{\detokenize{\chi;\Delta_{\mathcal D}}} & finite exceptional-size disjunction and complete finite relational diagram & P9 Fagin 1974 generalized spectra & n/a & \hyperref[annex:p9.n-c-004.1.theorem6]{\ref*{annex:p9.n-c-004.1.theorem6}} & \thepage & These supply the source's finite-modification step on the original same universe. \\
\phantomsection\label{entry:symbol-registry-03654}\(v\) & \texttt{\detokenize{v}} & finite variable-slot bound used for pure-set equality-pattern invariance & P9 Fagin 1974 unordered comparison & n/a & \hyperref[annex:p9.n-c-004.1.unordered-boundary]{\ref*{annex:p9.n-c-004.1.unordered-boundary}} & \thepage & A scalar variable bound distinct from indexed source tape functions v\_i; includes tuple binders parameters and all alpha-renamed individual variables. \\
\phantomsection\label{entry:symbol-registry-03655}\(Y\) & \texttt{\detokenize{Y}} & fixed finite positive-arity relational vocabulary & \hyperref[sec:p09-n-c-002-02]{\ref*{sec:p09-n-c-002-02}} & n/a & \hyperref[def:p09-fagin1976-probabilities]{\ref*{def:p09-fagin1976-probabilities}} & \thepage & No constants or functions in the zero-one theorem; auxiliary names and counterexample vocabularies are separately declared. \\
\phantomsection\label{entry:symbol-registry-03656}\([n]\) & \texttt{\detokenize{[n]}} & labelled universe {1,...,n} & \hyperref[sec:p09-n-c-002-02]{\ref*{sec:p09-n-c-002-02}} & n/a & \hyperref[def:p09-fagin1976-probabilities]{\ref*{def:p09-fagin1976-probabilities}} & \thepage & n is a positive integer; differs by indexing only from the adjacent P3 convention. \\
\phantomsection\label{entry:symbol-registry-03657}\(A_n(Y)\) & \texttt{\detokenize{A_n}} & set of labelled Y-structures & \hyperref[sec:p09-n-c-002-02]{\ref*{sec:p09-n-c-002-02}} & n/a & \hyperref[def:p09-fagin1976-probabilities]{\ref*{def:p09-fagin1976-probabilities}} & \thepage & Every ordered relation tuple including repetitions is independently fair. \\
\phantomsection\label{entry:symbol-registry-03658}\(a_n(Y)\) & \texttt{\detokenize{a_n}} & cardinality of the set of labelled Y-structures & \hyperref[sec:p09-n-c-002-02]{\ref*{sec:p09-n-c-002-02}} & n/a & \hyperref[def:p09-fagin1976-probabilities]{\ref*{def:p09-fagin1976-probabilities}} & \thepage & 2\^{}{sum\_i k\_i n\^{}i}. \\
\phantomsection\label{entry:symbol-registry-03659}\(B_n(Y)\) & \texttt{\detokenize{B_n}} & chosen set of isomorphism-type representatives & \hyperref[sec:p09-n-c-002-02]{\ref*{sec:p09-n-c-002-02}} & n/a & \hyperref[def:p09-fagin1976-probabilities]{\ref*{def:p09-fagin1976-probabilities}} & \thepage & One representative per isomorphism type; choice does not affect sentence truth. \\
\phantomsection\label{entry:symbol-registry-03660}\(b_n(Y)\) & \texttt{\detokenize{b_n}} & number of Y-isomorphism types of size n & \hyperref[sec:p09-n-c-002-02]{\ref*{sec:p09-n-c-002-02}} & n/a & \hyperref[def:p09-fagin1976-probabilities]{\ref*{def:p09-fagin1976-probabilities}} & \thepage & Cardinality of B\_n(Y); the unlabelled measure is uniform on these types. \\
\phantomsection\label{entry:symbol-registry-03661}\(\mu_n^Y\) & \texttt{\detokenize{\mu_n^Y}} & uniform probability measure on subsets of A\_n(Y) & \hyperref[sec:p09-n-c-002-02]{\ref*{sec:p09-n-c-002-02}} & n/a & \hyperref[def:p09-fagin1976-probabilities]{\ref*{def:p09-fagin1976-probabilities}} & \thepage & Distinct from sentence probability p\_n and final-page nonuniform mixture. \\
\phantomsection\label{entry:symbol-registry-03662}\(p_n(\sigma)\) & \texttt{\detokenize{p_n}} & labelled finite-model probability of a sentence & \hyperref[sec:p09-n-c-002-02]{\ref*{sec:p09-n-c-002-02}} & n/a & \hyperref[def:p09-fagin1976-probabilities]{\ref*{def:p09-fagin1976-probabilities}} & \thepage & Unaffected by unused relation symbols; conditional probability is separately normalized. \\
\phantomsection\label{entry:symbol-registry-03663}\(v_n^Y(\sigma)\) & \texttt{\detokenize{v_n^Y}} & unlabelled sentence probability under uniform isomorphism types & \hyperref[sec:p09-n-c-002-02]{\ref*{sec:p09-n-c-002-02}} & n/a & \hyperref[def:p09-fagin1976-probabilities]{\ref*{def:p09-fagin1976-probabilities}} & \thepage & Vocabulary superscript retained since unused symbols can affect the measure. \\
\phantomsection\label{entry:symbol-registry-03664}\(M(\bar x)\) & \texttt{\detokenize{M}} & complete relational open description on parameter variables & \hyperref[sec:p09-n-c-002-02]{\ref*{sec:p09-n-c-002-02}} & n/a & \hyperref[def:p09-fagin1976-extension]{\ref*{def:p09-fagin1976-extension}} & \thepage & No equality atoms; repeated variable tuples included; distinctness lives in axiom guards. \\
\phantomsection\label{entry:symbol-registry-03665}\(N(\bar x,y)\) & \texttt{\detokenize{N}} & complete open description extending M by one fresh variable & \hyperref[sec:p09-n-c-002-02]{\ref*{sec:p09-n-c-002-02}} & n/a & \hyperref[def:p09-fagin1976-extension]{\ref*{def:p09-fagin1976-extension}} & \thepage & Retains every conjunct of M and decides every new relational atom. \\
\phantomsection\label{entry:symbol-registry-03666}\(e_{M,N}\) & \texttt{\detokenize{e_{M,N}}} & one-point extension sentence & \hyperref[sec:p09-n-c-002-02]{\ref*{sec:p09-n-c-002-02}} & n/a & \hyperref[def:p09-fagin1976-extension]{\ref*{def:p09-fagin1976-extension}} & \thepage & Outer distinctness guard and fresh witness inequalities are part of the sentence. \\
\phantomsection\label{entry:symbol-registry-03667}\(T\) & \texttt{\detokenize{T}} & theory of all relational one-point extension sentences & \hyperref[sec:p09-n-c-002-02]{\ref*{sec:p09-n-c-002-02}} & n/a & \hyperref[def:p09-fagin1976-extension]{\ref*{def:p09-fagin1976-extension}} & \thepage & A sentence set, not a distribution or a countable-model name. \\
\phantomsection\label{entry:symbol-registry-03668}\(k_i\) & \texttt{\detokenize{k_i}} & number of i-ary symbols in Y & \hyperref[sec:p09-n-c-002-02]{\ref*{sec:p09-n-c-002-02}} & n/a & \hyperref[def:p09-fagin1976-probabilities]{\ref*{def:p09-fagin1976-probabilities}} & \thepage & Only finitely many k\_i are nonzero. \\
\phantomsection\label{entry:symbol-registry-03669}\(d_Y(m)\) & \texttt{\detokenize{d_Y(m)}} & number of fresh relational atoms involving one new point & \hyperref[sec:p09-n-c-002-02]{\ref*{sec:p09-n-c-002-02}} & n/a & \hyperref[thm:p09-fagin1976-extension-estimate]{\ref*{thm:p09-fagin1976-extension-estimate}} & \thepage & sum\_i k\_i((m+1)\^{}i-m\^{}i), including repeated entries. \\
\phantomsection\label{entry:symbol-registry-03670}\(q\) & \texttt{\detokenize{q}} & success probability of a fixed fresh candidate extension & \hyperref[sec:p09-n-c-002-02]{\ref*{sec:p09-n-c-002-02}} & n/a & \hyperref[thm:p09-fagin1976-extension-estimate]{\ref*{thm:p09-fagin1976-extension-estimate}} & \thepage & q=2\^{}{-d\_Y(m)}; graph q=p\^{}u(1-p)\^{}v is separately declared in the graph theorem. \\
\phantomsection\label{entry:symbol-registry-03671}\(\Sigma\) & \texttt{\detokenize{\Sigma}} & canonical full-relation projection-function infinite theory for spectrum proof & P9 Fagin1976 spectrum argument & n/a & \hyperref[thm:p09-fagin1976-finite-fragments]{\ref*{thm:p09-fagin1976-finite-fragments}} & \thepage & Not the relational extension theory T; source second spectrum proof. \\
\phantomsection\label{entry:symbol-registry-03672}\(\tau_1\) & \texttt{\detokenize{\tau_1}} & symmetry and irreflexivity conditioning sentence & \hyperref[sec:p09-n-c-002-02]{\ref*{sec:p09-n-c-002-02}} & n/a & \hyperref[thm:p09-fagin1976-graph]{\ref*{thm:p09-fagin1976-graph}} & \thepage & Conditional unrestricted binary measure equals fair simple-graph sampling. \\
\phantomsection\label{entry:symbol-registry-03673}\(\tau_2\) & \texttt{\detokenize{\tau_2}} & tournament conditioning sentence & \hyperref[sec:p09-n-c-002-02]{\ref*{sec:p09-n-c-002-02}} & n/a & \hyperref[thm:p09-fagin1976-graph]{\ref*{thm:p09-fagin1976-graph}} & \thepage & One direction per unordered pair and no loops. \\
\phantomsection\label{entry:symbol-registry-03674}\(\tau_3\) & \texttt{\detokenize{\tau_3}} & conditioning sentence making U uniformly true or uniformly false & \hyperref[sec:p09-n-c-002-02]{\ref*{sec:p09-n-c-002-02}} & n/a & \hyperref[thm:p09-fagin1976-boundaries]{\ref*{thm:p09-fagin1976-boundaries}} & \thepage & Conditional sentence probabilities are eventually 0, 1/2, or 1. \\
\phantomsection\label{entry:symbol-registry-03675}\(\alpha\) & \texttt{\detokenize{\alpha}} & mixture mass on the full relation & \hyperref[sec:p09-n-c-002-02]{\ref*{sec:p09-n-c-002-02}} & n/a & \hyperref[thm:p09-fagin1976-mixed-measure]{\ref*{thm:p09-fagin1976-mixed-measure}} & \thepage & 0<=alpha<1; source statement-only limit claim retained. \\
\phantomsection\label{entry:symbol-registry-03676}\(f(t)\) & \texttt{\detokenize{f}} & auxiliary weight function for the mixed relation measure & \hyperref[sec:p09-n-c-002-02]{\ref*{sec:p09-n-c-002-02}} & n/a & \hyperref[thm:p09-fagin1976-mixed-measure]{\ref*{thm:p09-fagin1976-mixed-measure}} & \thepage & f(t)=(1-alpha)p\^{}t+alpha. \\
\phantomsection\label{entry:symbol-registry-03677}\(\beta\) & \texttt{\detokenize{\beta}} & fixed lower proportion used to count many nonisomorphic models & \hyperref[sec:p09-n-c-002-02]{\ref*{sec:p09-n-c-002-02}} & n/a & \hyperref[thm:p09-fagin1976-finite-fragments]{\ref*{thm:p09-fagin1976-finite-fragments}} & \thepage & 0<beta<1; unrelated to the mixture mass alpha. \\
\phantomsection\label{entry:symbol-registry-03678}\(p\) & \texttt{\detokenize{p}} & fixed probability of an independent edge or ordered relation fact & \hyperref[sec:p09-n-c-002-02]{\ref*{sec:p09-n-c-002-02}} & n/a & \hyperref[thm:p09-fagin1976-graph]{\ref*{thm:p09-fagin1976-graph}} & \thepage & Fixed in (0,1); also used in the separately defined final-page Bernoulli mixture. Not the sentence-probability function p\_n. \\
\phantomsection\label{entry:symbol-registry-03679}\(p(n)\) & \texttt{\detokenize{p(n)}} & varying graph-edge probability in the boundary example & \hyperref[def:p09-fagin1976-outside-hypotheses]{\ref*{def:p09-fagin1976-outside-hypotheses}} & n/a & \hyperref[def:p09-fagin1976-outside-hypotheses]{\ref*{def:p09-fagin1976-outside-hypotheses}} & \thepage & p(n)=1/n\^{}2 is outside the fixed-density graph theorem. \\
\phantomsection\label{entry:symbol-registry-03680}\(Q\) & \texttt{\detokenize{Q}} & source graph predicate after replacement of unrestricted binary P & \hyperref[thm:p09-fagin1976-graph]{\ref*{thm:p09-fagin1976-graph}} & n/a & \hyperref[thm:p09-fagin1976-graph]{\ref*{thm:p09-fagin1976-graph}} & \thepage & An unordered-pair graph interpretation; the displayed E language is explicitly the same graph language with its edge symbol renamed. \\
\phantomsection\label{entry:symbol-registry-03681}\(\sigma^Q\) & \texttt{\detokenize{\sigma^Q}} & sentence obtained by replacing every P atom with the corresponding Q atom & \hyperref[thm:p09-fagin1976-graph]{\ref*{thm:p09-fagin1976-graph}} & n/a & \hyperref[thm:p09-fagin1976-graph]{\ref*{thm:p09-fagin1976-graph}} & \thepage & Source conditional-to-graph translation preserves equality and all logical operations. \\
\phantomsection\label{entry:symbol-registry-03682}\(s\) & \texttt{\detokenize{s}} & number 2\^{}h of unary one-point types & \hyperref[thm:p09-fagin1976-monadic]{\ref*{thm:p09-fagin1976-monadic}} & n/a & \hyperref[thm:p09-fagin1976-monadic]{\ref*{thm:p09-fagin1976-monadic}} & \thepage & Not the final-page count of false ordered-pair facts. \\
\phantomsection\label{entry:symbol-registry-03683}\(s\) & \texttt{\detokenize{s}} & number of false ordered-pair facts in one structure & \hyperref[thm:p09-fagin1976-mixed-measure]{\ref*{thm:p09-fagin1976-mixed-measure}} & n/a & \hyperref[thm:p09-fagin1976-mixed-measure]{\ref*{thm:p09-fagin1976-mixed-measure}} & \thepage & r+s=n\^{}2; disjoint local role from monadic type count. \\
\phantomsection\label{entry:symbol-registry-03684}\(\exists!y\) & \texttt{\detokenize{\exists!}} & unique-existence abbreviation in the pairing sentences & \hyperref[thm:p09-fagin1976-boundaries]{\ref*{thm:p09-fagin1976-boundaries}} & n/a & \hyperref[thm:p09-fagin1976-boundaries]{\ref*{thm:p09-fagin1976-boundaries}} & \thepage & Eliminable first-order abbreviation with a fresh universally quantified comparison variable. \\
\phantomsection\label{entry:symbol-registry-03878}\(B=(D,R_1,\ldots,R_k)\) & \texttt{\detokenize{B;D;R_i}} & source finite relational database domain and input relations & P9 Vardi 1982 query complexity & n/a & \hyperref[def:p09-vardi82-query]{\ref*{def:p09-vardi82-query}} & \thepage & D is a finite subset of the natural numbers; relation ranks are positive. Renaming preserves the query answer. \\
\phantomsection\label{entry:symbol-registry-03879}\(\bar a,b\) & \texttt{\detokenize{\bar a;b}} & source database rank vector and output rank & P9 Vardi 1982 query complexity & n/a & \hyperref[def:p09-vardi82-query]{\ref*{def:p09-vardi82-query}} & \thepage & For Boolean queries the output rank is zero; each fixed expression fixes its output rank. \\
\phantomsection\label{entry:symbol-registry-03880}\(Q_e,\mu,C(L)\) & \texttt{\detokenize{Q_e;\mu;C}} & meaning of an expression meaning function and definable query class & P9 Vardi 1982 query complexity & n/a & \hyperref[def:p09-vardi82-query]{\ref*{def:p09-vardi82-query}} & \thepage & General source queries are partial functions; logical expressions in this unit define total queries. \\
\phantomsection\label{entry:symbol-registry-03881}\(\operatorname{Gr}(Q),\operatorname{Gr}_L(B)\) & \texttt{\detokenize{\operatorname{Gr}}} & query answer-membership graph and fixed-database expression graph & P9 Vardi 1982 query complexity & n/a & \hyperref[def:p09-vardi82-query]{\ref*{def:p09-vardi82-query}} & \thepage & The graph does not distinguish undefined input from a defined empty answer. \\
\phantomsection\label{entry:symbol-registry-03882}\(Q\mathcal K,Q\mathcal K_{\rm tot}\) & \texttt{\detokenize{Q\mathcal K;Q\mathcal K_{\rm tot}}} & source graph-complexity query class and its explicitly qualified total subclass & P9 Vardi 1982 query complexity & n/a & \hyperref[def:p09-vardi82-axes]{\ref*{def:p09-vardi82-axes}} & \thepage & The total subclass prevents a false partial-function capture equality. \\
\phantomsection\label{entry:symbol-registry-03883}\(F^-,F,TF,YF,SF\) & \texttt{\detokenize{F^-;F;TF;YF;SF}} & source quantifier-free first-order closure fixed-point and existential-relation languages & P9 Vardi 1982 query complexity & n/a & \hyperref[def:p09-vardi82-languages]{\ref*{def:p09-vardi82-languages}} & \thepage & TF and YF here initially add specified outer operators; they are not assumed to be modern full FO(TC) or FO(LFP). \\
\phantomsection\label{entry:symbol-registry-03884}\(\Upsilon R.\bar x.\varphi\) & \texttt{\detokenize{\Upsilon}} & source positive least-fixed-point expression & P9 Vardi 1982 query complexity & object & \hyperref[def:p09-vardi82-languages]{\ref*{def:p09-vardi82-languages}} & \thepage & The bound relation rank equals the output-tuple length; the first-order matrix is positive in that relation. \\
\phantomsection\label{entry:symbol-registry-03885}\(\tau.(x,y).\varphi\) & \texttt{\detokenize{\tau}} & source positive-length transitive closure of a binary first-order query & P9 Vardi 1982 query complexity & object & \hyperref[def:p09-vardi82-languages]{\ref*{def:p09-vardi82-languages}} & \thepage & Not the general later parameterized nested transitive-closure grammar. \\
\phantomsection\label{entry:symbol-registry-03886}\(YF^+,Y\) & \texttt{\detokenize{YF^+;Y}} & source equality-restriction projection extension and separately logically closed fixed-point language & P9 Vardi 1982 query complexity & n/a & \hyperref[def:p09-vardi82-languages]{\ref*{def:p09-vardi82-languages}} & \thepage & The source's C(Y)=C(YF-plus) is a cited normal-form identity; it is not silently built into the initial syntax. \\
\phantomsection\label{entry:symbol-registry-03887}\(\Phi_B,X_j\) & \texttt{\detokenize{\Phi_B;X_j}} & positive operator and its finite approximation stages & P9 Vardi 1982 query complexity & n/a & \hyperref[def:p09-vardi82-languages]{\ref*{def:p09-vardi82-languages}} & \thepage & Starts from the empty relation; at most domain-size raised to rank strict increases. \\
\phantomsection\label{entry:symbol-registry-03888}\(B_0\) & \texttt{\detokenize{B_0}} & fixed two-element source database for expression hardness & P9 Vardi 1982 query complexity & n/a & \hyperref[thm:p09-vardi82-theorem5]{\ref*{thm:p09-vardi82-theorem5}} & \thepage & Domain zero and one with a unary predicate naming one; the database stays fixed while expression arity grows. \\
\phantomsection\label{entry:symbol-registry-03889}\(PTIME,EXPTIME\) & \texttt{\detokenize{PTIME;EXPTIME}} & source polynomial-time data and exponential-time expression classes & P9 Vardi 1982 query complexity & n/a & \hyperref[thm:p09-vardi82-theorem5]{\ref*{thm:p09-vardi82-theorem5}} & \thepage & Theorem 5 completeness uses logspace reductions and two different input families. \\
\phantomsection\label{entry:symbol-registry-03890}\(QPTIME_{\rm tot},QNPTIME_{\rm tot}\) & \texttt{\detokenize{QPTIME_{\rm tot};QNPTIME_{\rm tot}}} & qualified total-query recognition classes for the capture display & P9 Vardi 1982 query complexity & n/a & \hyperref[thm:p09-vardi82-remark2]{\ref*{thm:p09-vardi82-remark2}} & \thepage & Total-query or graph-equivalent semantics is explicit; domain-definedness of arbitrary partial functions is not captured. \\
\phantomsection\label{entry:symbol-registry-03891}\(\leq\) & \texttt{\detokenize{\leq}} & source distinguished input linear-order relation & P9 Vardi 1982 ordered capture & relation & \hyperref[thm:p09-vardi82-remark2]{\ref*{thm:p09-vardi82-remark2}} & \thepage & The original printed capture uses the weak-order glyph; strict order is derived by inequality. \\
\phantomsection\label{entry:symbol-registry-03892}\(EVEN\) & \texttt{\detokenize{EVEN}} & source Boolean domain-parity query & P9 Vardi 1982 unordered boundary & n/a & \hyperref[def:p09-vardi82-even]{\ref*{def:p09-vardi82-even}} & \thepage & Logspace computable but not expressible in unordered YF-plus. \\
\phantomsection\label{entry:symbol-registry-03893}\(A,P,X\) & \texttt{\detokenize{A;P;X}} & path-system initial points binary-premise triples and accessible relation & P9 Vardi 1982 data-hardness sketch & n/a & \hyperref[annex:p9.n-c-004.2.theorem5]{\ref*{annex:p9.n-c-004.2.theorem5}} & \thepage & P is ternary; both premises must already be accessible. External PTIME-completeness is retained as cited. \\
\phantomsection\label{entry:symbol-registry-03894}\(d,s,N,r,w\) & \texttt{\detokenize{d;s;N;r;w}} & domain cardinality expression size combined size relation rank and variable-slot bound & P9 Vardi 1982 iteration lemma & n/a & \hyperref[annex:p9.n-c-004.2.iteration]{\ref*{annex:p9.n-c-004.2.iteration}} & \thepage & Their roles are local and declared; fixed rank on the data axis is not fixed on the expression axis. \\
\phantomsection\label{entry:symbol-registry-03895}\(k,w,h,c\) & \texttt{\detokenize{k;w;h;c}} & fixed time-cell address width token-code width relation-block padding width and input-word exponent & P9 Vardi 1982 ordered tableau & n/a & \hyperref[annex:p9.n-c-004.2.ordered]{\ref*{annex:p9.n-c-004.2.ordered}} & \thepage & This local w is a code width rather than the earlier lemma's individual-variable bound; k is not the source database's number of relations. \\
\phantomsection\label{entry:symbol-registry-03896}\(\bar d\) & \texttt{\detokenize{\bar d}} & candidate query tuple and the prefix separating its computation & P9 Vardi 1982 ordered tableau & n/a & \hyperref[annex:p9.n-c-004.2.ordered]{\ref*{annex:p9.n-c-004.2.ordered}} & \thepage & These are bound coordinates of the one tableau relation; the scalar d separately denotes domain cardinality. \\
\phantomsection\label{entry:symbol-registry-03897}\(\bar a\) & \texttt{\detokenize{\bar a}} & finite tuple coding one complete cell token & P9 Vardi 1982 ordered tableau & n/a & \hyperref[annex:p9.n-c-004.2.ordered]{\ref*{annex:p9.n-c-004.2.ordered}} & \thepage & This local element tuple is distinct from the source's integer-valued database rank vector at its database-definition scope. \\
\phantomsection\label{entry:symbol-registry-03898}\(\operatorname{Min},\operatorname{Max},\operatorname{Succ}_k\) & \texttt{\detokenize{\operatorname{Min};\operatorname{Max};\operatorname{Succ}_k}} & first-order endpoint tests and nonwrapping tuple successor & P9 Vardi 1982 ordered reconstruction & relation & \hyperref[annex:p9.n-c-004.2.ordered]{\ref*{annex:p9.n-c-004.2.ordered}} & \thepage & Defined from the input order without adding endpoint constants or a BIT predicate. \\
\phantomsection\label{entry:symbol-registry-03899}\(\Omega,f\) & \texttt{\detokenize{\Omega;f}} & finite full-cell token alphabet and deterministic radius-one update map & P9 Vardi 1982 ordered reconstruction & n/a & \hyperref[annex:p9.n-c-004.2.ordered]{\ref*{annex:p9.n-c-004.2.ordered}} & \thepage & Complete cell tokens encode blanks zero bits and absent heads as positive facts. \\
\phantomsection\label{entry:symbol-registry-03900}\(T,T_a\) & \texttt{\detokenize{T;T_a}} & single positive tableau relation and code-selected computational membership abbreviation & P9 Vardi 1982 ordered reconstruction & relation & \hyperref[annex:p9.n-c-004.2.ordered]{\ref*{annex:p9.n-c-004.2.ordered}} & \thepage & Candidate tuple prefixes are bound coordinates; T\_a is positive in the one evolving T relation. \\
\phantomsection\label{entry:symbol-registry-03901}\(f_0,f_1\) & \texttt{\detokenize{f_0;f_1}} & tuple flag coordinates separating computation from output signals & P9 Vardi 1982 ordered reconstruction & n/a & \hyperref[annex:p9.n-c-004.2.ordered]{\ref*{annex:p9.n-c-004.2.ordered}} & \thepage & Distinct endpoint flags select computation; equality selects answer signals using literal YF-plus projection. These are individual variables not function symbols. \\
\phantomsection\label{entry:symbol-registry-03902}\(\epsilon,\psi\) & \texttt{\detokenize{\epsilon;\psi}} & empty-domain Boolean truth case and the constructed positive operator matrix & P9 Vardi 1982 ordered reconstruction & n/a & \hyperref[annex:p9.n-c-004.2.ordered]{\ref*{annex:p9.n-c-004.2.ordered}} & \thepage & The Boolean case is combined in Y and passed to YF-plus by the source-cited normal-form identity. \\
\phantomsection\label{entry:symbol-registry-04013}\(\tau\) & \texttt{\detokenize{\tau}} & fixed finite similarity type of relations and constants & \hyperref[sec:p09-n-c-004-03]{\ref*{sec:p09-n-c-004-03}} & n/a & \hyperref[def:p09-immerman1986-query]{\ref*{def:p09-immerman1986-query}} & \thepage & At least c\_0 in the source normal form; order is separately required for capture. \\
\phantomsection\label{entry:symbol-registry-04014}\(B\) & \texttt{\detokenize{B}} & finite database of the fixed similarity type & \hyperref[sec:p09-n-c-004-03]{\ref*{sec:p09-n-c-004-03}} & n/a & \hyperref[def:p09-immerman1986-query]{\ref*{def:p09-immerman1986-query}} & \thepage & Relations and constants interpreted on its finite nonempty universe. \\
\phantomsection\label{entry:symbol-registry-04015}\(U\) & \texttt{\detokenize{U}} & universe of the finite database & \hyperref[sec:p09-n-c-004-03]{\ref*{sec:p09-n-c-004-03}} & n/a & \hyperref[def:p09-immerman1986-query]{\ref*{def:p09-immerman1986-query}} & \thepage & Quantifiers range over U; tuple coordinates also use its elements. \\
\phantomsection\label{entry:symbol-registry-04016}\(n\) & \texttt{\detokenize{n}} & cardinality of the database universe & \hyperref[sec:p09-n-c-004-03]{\ref*{sec:p09-n-c-004-03}} & n/a & \hyperref[def:p09-immerman1986-query]{\ref*{def:p09-immerman1986-query}} & \thepage & n=|U|; explicit dense encodings only. \\
\phantomsection\label{entry:symbol-registry-04017}\(\gamma^B\) & \texttt{\detokenize{\gamma^B}} & answer relation of a fixed query on B & \hyperref[sec:p09-n-c-004-03]{\ref*{sec:p09-n-c-004-03}} & n/a & \hyperref[def:p09-immerman1986-query]{\ref*{def:p09-immerman1986-query}} & \thepage & Free variable assignment is explicit in the defining satisfaction clause. \\
\phantomsection\label{entry:symbol-registry-04018}\(\mathcal L(\tau)\) & \texttt{\detokenize{\mathcal L}} & source domain relational calculus & \hyperref[sec:p09-n-c-004-03]{\ref*{sec:p09-n-c-004-03}} & n/a & \hyperref[def:p09-immerman1986-query]{\ref*{def:p09-immerman1986-query}} & \thepage & First-order logic with equality, relation and constant symbols. \\
\phantomsection\label{entry:symbol-registry-04019}\(\mathcal L_\mu(\tau)\) & \texttt{\detokenize{\mathcal L_\mu}} & source positive least fixed-point query language & \hyperref[sec:p09-n-c-004-03]{\ref*{sec:p09-n-c-004-03}} & n/a & \hyperref[def:p09-immerman1986-lfp]{\ref*{def:p09-immerman1986-lfp}} & \thepage & Closed under ordinary Boolean operations, quantification and admitted nested fixed points. \\
\phantomsection\label{entry:symbol-registry-04020}\(\mathrm{QPTIME}\) & \texttt{\detokenize{\mathrm{QPTIME}}} & class of fixed polynomial-time query membership graphs & \hyperref[sec:p09-n-c-004-03]{\ref*{sec:p09-n-c-004-03}} & n/a & \hyperref[def:p09-immerman1986-query]{\ref*{def:p09-immerman1986-query}} & \thepage & Not formula-as-input or combined evaluation complexity. \\
\phantomsection\label{entry:symbol-registry-04021}\(\Gamma_{\varphi,\bar b}^B\) & \texttt{\detokenize{\Gamma}} & monotone operator on relation interpretations for fixed parameter assignment & \hyperref[sec:p09-n-c-004-03]{\ref*{sec:p09-n-c-004-03}} & n/a & \hyperref[def:p09-immerman1986-lfp]{\ref*{def:p09-immerman1986-lfp}} & \thepage & Defined on subsets of U\^{}r; positivity ensures monotonicity. \\
\phantomsection\label{entry:symbol-registry-04022}\(I_i\) & \texttt{\detokenize{I_i}} & ith iterate from the empty relation & \hyperref[def:p09-immerman1986-lfp]{\ref*{def:p09-immerman1986-lfp}} & n/a & \hyperref[def:p09-immerman1986-lfp]{\ref*{def:p09-immerman1986-lfp}} & \thepage & Ascending from bottom even when the operator is not inflationary on every relation. \\
\phantomsection\label{entry:symbol-registry-04023}\((\bar t:\mu R(\bar x))\varphi\) & \texttt{\detokenize{\mu}} & source least-fixed-point membership binder & \hyperref[sec:p09-n-c-004-03]{\ref*{sec:p09-n-c-004-03}} & n/a & \hyperref[def:p09-immerman1986-lfp]{\ref*{def:p09-immerman1986-lfp}} & \thepage & R has positive arity |x| under the constant-proviso convention and occurs positively; t is the queried tuple of terms. \\
\phantomsection\label{entry:symbol-registry-04024}\(\Sigma_\omega\) & \texttt{\detokenize{\Sigma_\omega}} & source one-fixed-point hierarchy level & \hyperref[thm:p09-immerman1986-capture]{\ref*{thm:p09-immerman1986-capture}} & n/a & \hyperref[thm:p09-immerman1986-capture]{\ref*{thm:p09-immerman1986-capture}} & \thepage & One outer mu with first-order defining matrix; not arithmetical-hierarchy notation. \\
\phantomsection\label{entry:symbol-registry-04025}\(\operatorname{Lex}_k\) & \texttt{\detokenize{\operatorname{Lex}_k}} & first-order strict lexicographic tuple order & \hyperref[sec:p09-n-c-004-03]{\ref*{sec:p09-n-c-004-03}} & relation & \hyperref[thm:p09-immerman1986-initial]{\ref*{thm:p09-immerman1986-initial}} & \thepage & Finite k fixed by the query/machine; no added arithmetic functions. \\
\phantomsection\label{entry:symbol-registry-04026}\(\operatorname{Succ}_k\) & \texttt{\detokenize{\operatorname{Succ}_k}} & first-order immediate successor on the tuple order & \hyperref[sec:p09-n-c-004-03]{\ref*{sec:p09-n-c-004-03}} & relation & \hyperref[thm:p09-immerman1986-initial]{\ref*{thm:p09-immerman1986-initial}} & \thepage & Absence of an intermediate tuple; p and t superscripts distinguish coordinate orders. \\
\phantomsection\label{entry:symbol-registry-04027}\(M_0(\bar p,a)\) & \texttt{\detokenize{M_0}} & first-order initial tape-letter formula & \hyperref[sec:p09-n-c-004-03]{\ref*{sec:p09-n-c-004-03}} & n/a & \hyperref[thm:p09-immerman1986-initial]{\ref*{thm:p09-immerman1986-initial}} & \thepage & Source Lemma 3.1 binary layout preserved; general decoding explicitly completed. \\
\phantomsection\label{entry:symbol-registry-04028}\(\delta_M\) & \texttt{\detokenize{\delta_M}} & finite set of local neighborhood-to-letter quadruples & \hyperref[sec:p09-n-c-004-03]{\ref*{sec:p09-n-c-004-03}} & n/a & \hyperref[thm:p09-immerman1986-tableau]{\ref*{thm:p09-immerman1986-tableau}} & \thepage & Encodes deterministic one-tape transitions, including specified endpoint behavior. \\
\phantomsection\label{entry:symbol-registry-04029}\(\Delta_M\) & \texttt{\detokenize{\Delta_M}} & first-order finite-disjunction expansion of the source local transition table & \hyperref[sec:p09-n-c-004-03]{\ref*{sec:p09-n-c-004-03}} & n/a & \hyperref[thm:p09-immerman1986-tableau]{\ref*{thm:p09-immerman1986-tableau}} & \thepage & No recursively bound relation occurs negated in the tableau update. \\
\phantomsection\label{entry:symbol-registry-04030}\(C_M\) & \texttt{\detokenize{C_M}} & least-fixed-point cell/time/letter relation for a normalized machine & \hyperref[sec:p09-n-c-004-03]{\ref*{sec:p09-n-c-004-03}} & n/a & \hyperref[thm:p09-immerman1986-tableau]{\ref*{thm:p09-immerman1986-tableau}} & \thepage & Arity 2k+1; one letter per tape cell and time in the complete tableau. \\
\phantomsection\label{entry:symbol-registry-04031}\(q_f\) & \texttt{\detokenize{q_f}} & designated stationary accepting letter at first tape cell & \hyperref[sec:p09-n-c-004-03]{\ref*{sec:p09-n-c-004-03}} & n/a & \hyperref[thm:p09-immerman1986-tableau]{\ref*{thm:p09-immerman1986-tableau}} & \thepage & Machine returns to that cell before terminal padding; coded by c\_0 for source normal form. \\
\phantomsection\label{entry:symbol-registry-04032}\(\#\) & \texttt{\detokenize{\#}} & exterior tape-neighbor letter & \hyperref[sec:p09-n-c-004-03]{\ref*{sec:p09-n-c-004-03}} & n/a & \hyperref[thm:p09-immerman1986-tableau]{\ref*{thm:p09-immerman1986-tableau}} & \thepage & Supplies missing endpoint neighbors; not a wraparound tape cell. \\
\phantomsection\label{entry:symbol-registry-04033}\(<_p\) & \texttt{\detokenize{<_p}} & position order with c\_0 first & P9 Immerman1986 coordinate recoding & relation & \hyperref[thm:p09-immerman1986-capture]{\ref*{thm:p09-immerman1986-capture}} & \thepage & First-order recoding of the given order; the input order is retained in its own table. \\
\phantomsection\label{entry:symbol-registry-04034}\(<_t\) & \texttt{\detokenize{<_t}} & time order with c\_0 last & P9 Immerman1986 coordinate recoding & relation & \hyperref[thm:p09-immerman1986-capture]{\ref*{thm:p09-immerman1986-capture}} & \thepage & Different coordinate role from the position order. \\
\phantomsection\label{entry:symbol-registry-04035}\(\operatorname{Code}_b(v)\) & \texttt{\detokenize{\operatorname{Code}_b}} & first-order test for a designated letter-code element & \hyperref[sec:p09-n-c-004-03]{\ref*{sec:p09-n-c-004-03}} & relation & \hyperref[thm:p09-immerman1986-tableau]{\ref*{thm:p09-immerman1986-tableau}} & \thepage & Eliminates displayed code equalities; accepting code may be c\_0, and unavailable fixed-rank tests are false. \\
\phantomsection\label{entry:symbol-registry-04234}\(\Pi\) & \texttt{\detokenize{\Pi}} & source ground normal program & P9 Gelfond-Lifschitz 1988 & n/a & \hyperref[def:p09-gl88-language]{\ref*{def:p09-gl88-language}} & \thepage & Possibly infinite after all ground instances; one atomic head and finite body. \\
\phantomsection\label{entry:symbol-registry-04235}\(A,L_i,B,H\) & \texttt{\detokenize{A;L_i;B;H}} & source head atoms body literals and atoms inside negative body conditions & P9 Gelfond-Lifschitz 1988 program syntax & n/a & \hyperref[def:p09-gl88-language]{\ref*{def:p09-gl88-language}} & \thepage & L\_i is a rule literal; L without an index later is the autoepistemic operator. A later objective theory A has a separately declared local type. \\
\phantomsection\label{entry:symbol-registry-04236}\(\neg B\) & \texttt{\detokenize{\neg}} & source negative body condition with classical implication reading and candidate-dependent default use & P9 Gelfond-Lifschitz 1988 program syntax & n/a & \hyperref[def:p09-gl88-language]{\ref*{def:p09-gl88-language}} & \thepage & Not a strong-negation head literal. The source symbol is preserved instead of replaced by house not notation. \\
\phantomsection\label{entry:symbol-registry-04237}\(At,M,N\) & \texttt{\detokenize{At;M;N}} & grounded occurring atom set candidate stable set and comparison Herbrand interpretation & P9 Gelfond-Lifschitz 1988 & n/a & \hyperref[def:p09-gl88-language]{\ref*{def:p09-gl88-language}} & \thepage & At can be infinite and differs from the larger Herbrand base. M and N are sets of atoms not theories of formulas. \\
\phantomsection\label{entry:symbol-registry-04238}\(CM(\Pi)\) & \texttt{\detokenize{CM}} & source unique canonical Herbrand model & P9 Gelfond-Lifschitz 1988 & n/a & \hyperref[def:p09-gl88-reduct]{\ref*{def:p09-gl88-reduct}} & \thepage & Only designated by the original stable model semantics when there is exactly one stable model. \\
\phantomsection\label{entry:symbol-registry-04239}\(\Pi_M,S_\Pi(M)\) & \texttt{\detokenize{\Pi_M;S_\Pi}} & source candidate reduct and its positive least model & P9 Gelfond-Lifschitz 1988 & n/a & \hyperref[def:p09-gl88-reduct]{\ref*{def:p09-gl88-reduct}} & \thepage & Retains original subscript and S operator; maps candidates antimonotonically. \\
\phantomsection\label{entry:symbol-registry-04240}\(Q,T_Q,X_n,M_Q\) & \texttt{\detokenize{Q;T_Q;X_n;M_Q}} & editorial positive program consequence operator finite-stage chain and least model & P9 Gelfond-Lifschitz 1988 positive closure & n/a & \hyperref[thm:p09-gl88-positive-closure]{\ref*{thm:p09-gl88-positive-closure}} & \thepage & Individual rule bodies are finite even if the rule set is infinite; omega union gives the least model. \\
\phantomsection\label{entry:symbol-registry-04241}\(D,D'\) & \texttt{\detokenize{D;D'}} & sets of source example input facts & P9 Gelfond-Lifschitz 1988 examples & n/a & \hyperref[def:p09-gl88-one-model]{\ref*{def:p09-gl88-one-model}} & \thepage & D contains p(1 2); D-prime also contains p(2 1). Not domains or default theories. \\
\phantomsection\label{entry:symbol-registry-04242}\(L\) & \texttt{\detokenize{L}} & source autoepistemic belief-forming operator & P9 Gelfond-Lifschitz 1988 autoepistemic translation & n/a & \hyperref[def:p09-gl88-expansion]{\ref*{def:p09-gl88-expansion}} & \thepage & Modal atoms LF are independent letters for ordinary propositional consequence; no modal axioms or necessitation are added. \\
\phantomsection\label{entry:symbol-registry-04243}\(A,E,F\) & \texttt{\detokenize{A;E;F}} & objective or translated theory expansion candidate and arbitrary autoepistemic formula & P9 Gelfond-Lifschitz 1988 autoepistemic translation & n/a & \hyperref[def:p09-gl88-expansion]{\ref*{def:p09-gl88-expansion}} & \thepage & This A is a set of formulas rather than the earlier head atom; F ranges over all finite formulas including modal ones. \\
\phantomsection\label{entry:symbol-registry-04244}\(th(X)\) & \texttt{\detokenize{th}} & source ordinary propositional consequence closure & P9 Gelfond-Lifschitz 1988 autoepistemic translation & n/a & \hyperref[def:p09-gl88-expansion]{\ref*{def:p09-gl88-expansion}} & \thepage & Objective kernel is consequence from the objective theory not all truths of its least Herbrand valuation. \\
\phantomsection\label{entry:symbol-registry-04245}\(I(\Pi)\) & \texttt{\detokenize{I}} & source rulewise autoepistemic translation & P9 Gelfond-Lifschitz 1988 autoepistemic translation & n/a & \hyperref[def:p09-gl88-expansion]{\ref*{def:p09-gl88-expansion}} & \thepage & Replaces each negative body B condition by neg L B; ground atoms become propositional letters. \\
\phantomsection\label{entry:symbol-registry-04246}\(\Delta(E)\) & \texttt{\detokenize{\Delta}} & editorial abbreviation for all positive and negative introspective assumptions & P9 Gelfond-Lifschitz 1988 expansion proofs & n/a & \hyperref[def:p09-gl88-expansion]{\ref*{def:p09-gl88-expansion}} & \thepage & Includes LF for F in E and neg LF for F outside E over the entire formula language; not just program atoms. \\
\phantomsection\label{entry:symbol-registry-04247}\(r_A,E_A\) & \texttt{\detokenize{r_A;E_A}} & editorial finite-modal-depth objective reduction and constructed expansion & P9 Gelfond-Lifschitz 1988 objective-expansion proof & n/a & \hyperref[annex:p9.n-c-023.1.objective-expansion]{\ref*{annex:p9.n-c-023.1.objective-expansion}} & \thepage & Recurses on proper subformulas; no fixed-point existence is assumed in the construction. \\
\phantomsection\label{entry:symbol-registry-04248}\(\top,\bot\) & \texttt{\detokenize{\top;\bot}} & objective Boolean truth and contradiction abbreviations & P9 Gelfond-Lifschitz 1988 objective-expansion proof & n/a & \hyperref[annex:p9.n-c-023.1.objective-expansion]{\ref*{annex:p9.n-c-023.1.objective-expansion}} & \thepage & May be replaced by a fixed objective tautology and contradiction. \\
\phantomsection\label{entry:symbol-registry-04249}\(v,\widehat v,d,G\) & \texttt{\detokenize{v;\widehat v;d;G}} & objective valuation its modal-skeleton extension modal depth and modal operand & P9 Gelfond-Lifschitz 1988 objective-expansion proof & n/a & \hyperref[annex:p9.n-c-023.1.objective-expansion]{\ref*{annex:p9.n-c-023.1.objective-expansion}} & \thepage & Modal assignments are independent Boolean values for syntactically distinct LF letters. \\
\phantomsection\label{entry:symbol-registry-04250}\(w,d,k\) & \texttt{\detokenize{w;d;k}} & fresh reduction witness atom constant count and variable count & P9 Gelfond-Lifschitz 1988 task boundary & n/a & \hyperref[sec:p09-n-c-023-01]{\ref*{sec:p09-n-c-023-01}} & \thepage & d and k here describe grounding size; local proof-depth d has a different declared scope. \\
\phantomsection\label{entry:symbol-registry-05057}\(A\) & \texttt{\detokenize{A}} & finite-control timed automaton with declared invariant convention & local:sec:p09-timed-automata-regions-and-zones & object & \hyperref[def:p09-timed-model]{\ref*{def:p09-timed-model}} & \thepage & Not a formula or generic hybrid automaton \\
\phantomsection\label{entry:symbol-registry-05058}\(C\) & \texttt{\detokenize{C}} & finite set of rate-one real-valued clocks & local:sec:p09-timed-automata-regions-and-zones & object & \hyperref[def:p09-timed-model]{\ref*{def:p09-timed-model}} & \thepage & One global elapsed time; no stopped or drifting clocks \\
\phantomsection\label{entry:symbol-registry-05059}\(v[r]\) & \texttt{\detokenize{v[r]}} & simultaneous zero-reset valuation & local:sec:p09-timed-automata-regions-and-zones & meta & \hyperref[def:p09-timed-model]{\ref*{def:p09-timed-model}} & \thepage & Reset image is not zero substitution in every old constraint \\
\phantomsection\label{entry:symbol-registry-05060}\(M_x\) & \texttt{\detokenize{M_x}} & nonnegative integer comparison bound for a clock & local:sec:p09-timed-automata-regions-and-zones & object & \hyperref[def:p09-timed-region-equivalence]{\ref*{def:p09-timed-region-equivalence}} & \thepage & Includes invariant constants; unused clock default zero \\
\phantomsection\label{entry:symbol-registry-05061}\(\langle t\rangle\) & \texttt{\detokenize{\langle t\rangle}} & fractional part of a nonnegative real & local:sec:p09-timed-automata-regions-and-zones & object & \hyperref[def:p09-timed-region-equivalence]{\ref*{def:p09-timed-region-equivalence}} & \thepage & Unary angle notation distinct from tuple delimiters \\
\phantomsection\label{entry:symbol-registry-05062}\(\sim_M\) & \texttt{\detokenize{\sim_M}} & explicit bounded-cell region equivalence & local:sec:p09-timed-automata-regions-and-zones & relation & \hyperref[def:p09-timed-region-equivalence]{\ref*{def:p09-timed-region-equivalence}} & \thepage & Matching over-limit status with floor and zero flags and weak fractional ordering \\
\phantomsection\label{entry:symbol-registry-05063}\(\mathcal D\) & \texttt{\detokenize{\mathcal D}} & finite set of tested diagonal atoms & local:sec:p09-timed-automata-regions-and-zones & object & \hyperref[def:p09-timed-diagonal-refinement]{\ref*{def:p09-timed-diagonal-refinement}} & \thepage & Truth-vector refinement preserves diagonal tests \\
\phantomsection\label{entry:symbol-registry-05064}\(\sim_{M,\mathcal D}\) & \texttt{\detokenize{\sim_{M,\mathcal D}}} & diagonal-sensitive finite region refinement & local:sec:p09-timed-automata-regions-and-zones & relation & \hyperref[def:p09-timed-diagonal-refinement]{\ref*{def:p09-timed-diagonal-refinement}} & \thepage & A refined region need not be one convex zone \\
\phantomsection\label{entry:symbol-registry-05065}\(D\) & \texttt{\detokenize{D}} & strict difference-bound matrix & local:sec:p09-timed-automata-regions-and-zones & object & \hyperref[def:p09-timed-dbm]{\ref*{def:p09-timed-dbm}} & \thepage & Entry ij bounds v\_i minus v\_j; nonnegativity rows included \\
\phantomsection\label{entry:symbol-registry-05066}\(x_0\) & \texttt{\detokenize{x_0}} & fixed zero coordinate of a DBM & local:sec:p09-timed-automata-regions-and-zones & object & \hyperref[def:p09-timed-dbm]{\ref*{def:p09-timed-dbm}} & \thepage & Not another evolving clock \\
\phantomsection\label{entry:symbol-registry-05067}\(\preceq_b\) & \texttt{\detokenize{\preceq_b}} & no-weaker-than order on strict bounds & local:sec:p09-timed-automata-regions-and-zones & relation & \hyperref[def:p09-timed-dbm]{\ref*{def:p09-timed-dbm}} & \thepage & At equal constant strict is tighter than nonstrict \\
\phantomsection\label{entry:symbol-registry-05068}\(\prec_b\) & \texttt{\detokenize{\prec_b}} & strictly tighter order on strict bounds & local:sec:p09-timed-automata-regions-and-zones & relation & \hyperref[def:p09-timed-dbm]{\ref*{def:p09-timed-dbm}} & \thepage & Zero strict diagonal is inconsistent \\
\phantomsection\label{entry:symbol-registry-05069}\(\oplus_b\) & \texttt{\detokenize{\oplus_b}} & sum of bounds with propagated strictness & local:sec:p09-timed-automata-regions-and-zones & infix & \hyperref[def:p09-timed-dbm]{\ref*{def:p09-timed-dbm}} & \thepage & Infinity absorbing; sum strict iff either summand strict \\
\phantomsection\label{entry:symbol-registry-05070}\(Z(D)\) & \texttt{\detokenize{Z(D)}} & nonnegative valuation set of a DBM & local:sec:p09-timed-automata-regions-and-zones & object & \hyperref[def:p09-timed-dbm]{\ref*{def:p09-timed-dbm}} & \thepage & Empty zone is separately propagated \\
\phantomsection\label{entry:symbol-registry-05071}\(Z^\uparrow\) & \texttt{\detokenize{Z^\uparrow}} & common nonnegative time-elapse image & local:sec:p09-timed-automata-regions-and-zones & meta & \hyperref[thm:p09-timed-dbm-operations]{\ref*{thm:p09-timed-dbm-operations}} & \thepage & All clock differences preserved during delay \\
\phantomsection\label{entry:symbol-registry-05072}\(\operatorname{Post}_e\) & \texttt{\detokenize{\operatorname{Post}_e}} & exact delay/action successor-set operator & local:sec:p09-timed-zone-reachability & meta & \hyperref[def:p09-timed-zone-post]{\ref*{def:p09-timed-zone-post}} & \thepage & Endpoint invariant test justified for declared convex conjunctive class \\
\phantomsection\label{entry:symbol-registry-05073}\(\alpha\) & \texttt{\detokenize{\alpha}} & finite region-union saturation operator & local:sec:p09-timed-zone-reachability & meta & \hyperref[def:p09-timed-region-saturation]{\ref*{def:p09-timed-region-saturation}} & \thepage & Not one DBM or an unproved practical LU extrapolation \\
\phantomsection\label{entry:symbol-registry-05074}\(\triangleleft\) & \texttt{\detokenize{\triangleleft}} & strict or nonstrict upper-bound comparison flag & local:sec:p09-timed-automata-regions-and-zones & relation & \hyperref[def:p09-timed-dbm]{\ref*{def:p09-timed-dbm}} & \thepage & Means less-than or less-than-or-equal as selected by bound \\
\phantomsection\label{entry:symbol-registry-05075}\(\triangleright\) & \texttt{\detokenize{\triangleright}} & reversed lower-bound comparison flag in projection proof & local:annex:p09-timed-dbm-operations & relation & \hyperref[annex:p09-timed-dbm-operations]{\ref*{annex:p09-timed-dbm-operations}} & \thepage & Means greater-than for strict and greater-than-or-equal for nonstrict \\
\phantomsection\label{entry:symbol-registry-05208}\(B\) & \texttt{\detokenize{B}} & all positive closed judgements & p09-negative-premises & n/a & \hyperref[def:p09-np-language]{\ref*{def:p09-np-language}} & \thepage & Disjoint action triples and unary predicate judgements \\
\phantomsection\label{entry:symbol-registry-05209}\(D(\eta)\) & \texttt{\detokenize{D}} & positive deniers & p09-negative-premises & n/a & \hyperref[def:p09-np-language]{\ref*{def:p09-np-language}} & \thepage & Action absence excludes all closed targets; singleton target negatives extend only proof language \\
\phantomsection\label{entry:symbol-registry-05210}\(\Gamma_T(M)\) & \texttt{\detokenize{\Gamma_T}} & least positive reduct relation & p09-negative-premises & n/a & \hyperref[def:p09-np-models]{\ref*{def:p09-np-models}} & \thepage & Antitone operator; not a generic search-failure test \\
\phantomsection\label{entry:symbol-registry-05211}\(L_*,U_*\) & \texttt{\detokenize{L_*,}\allowbreak\detokenize{U_*}} & certain and possible well-founded bounds & p09-negative-premises & n/a & \hyperref[def:p09-np-three-valued]{\ref*{def:p09-np-three-valued}} & \thepage & Unknown is upper minus lower; source uses disjoint true and unknown sets \\
\phantomsection\label{entry:symbol-registry-05212}\(S:B\to\kappa\) & \texttt{\detokenize{S}} & ordinal stratification & p09-negative-premises & n/a & \hyperref[def:p09-np-stratification]{\ref*{def:p09-np-stratification}} & \thepage & Positive weak and every negative-denier strict rank condition \\
\phantomsection\label{entry:symbol-registry-05213}\(\lambda\) & \texttt{\detokenize{\lambda}} & infinite regular positive-premise bound & p09-negative-premises & n/a & \hyperref[def:p09-np-stratification]{\ref*{def:p09-np-stratification}} & \thepage & Strictly greater than every positive premise-set size \\
\phantomsection\label{entry:symbol-registry-05214}\(t\checkmark\) & \texttt{\detokenize{\checkmark}} & successful termination predicate & p09-negative-premises & n/a & \hyperref[thm:p09-np-processes]{\ref*{thm:p09-np-processes}} & \thepage & Not absence of action transitions \\
\phantomsection\label{entry:symbol-registry-05215}\(\sigma_d,\sigma\) & \texttt{\detokenize{\sigma_d,}\allowbreak\detokenize{\sigma}} & delay operator and time action & p09-negative-premises & n/a & \hyperref[thm:p09-np-processes]{\ref*{thm:p09-np-processes}} & \thepage & Time is outside the ordinary action alphabet \\
\phantomsection\label{entry:symbol-registry-05399}\(\mathsf{LI}\) & \texttt{\detokenize{\mathsf{LI}}} & declared local-input predicate subclass & p09-sos-formats & n/a & \hyperref[def:p09-sos-local]{\ref*{def:p09-sos-local}} & \thepage & Not all panth; native nonrecursive GSOS is a literal special case \\
\phantomsection\label{entry:symbol-registry-05400}\(\sim\) & \texttt{\detokenize{\sim}} & largest predicate-respecting strong bisimulation & p09-sos-formats & n/a & \hyperref[def:p09-sos-bisimulation]{\ref*{def:p09-sos-bisimulation}} & \thepage & Predicates are not discarded \\
\phantomsection\label{entry:symbol-registry-05401}\(K\) & \texttt{\detokenize{K}} & least compatible closure of bisimilarity & p09-sos-formats & n/a & \hyperref[thm:p09-sos-congruence]{\ref*{thm:p09-sos-congruence}} & \thepage & Constructed by finite stages; not assumed transitive \\
\phantomsection\label{entry:symbol-registry-05402}\(C_0,C_1\) & \texttt{\detokenize{C_0,}\allowbreak\detokenize{C_1}} & old and extended closed-term universes & p09-sos-formats & n/a & \hyperref[def:p09-sos-conservative]{\ref*{def:p09-sos-conservative}} & \thepage & Free finite first-order syntax; matching arities in the intersection \\
\phantomsection\label{entry:symbol-registry-05403}\(T_0\oplus T_1\) & \texttt{\detokenize{\oplus}} & union TSS on compatible union signature & p09-sos-formats & n/a & \hyperref[def:p09-sos-conservative]{\ref*{def:p09-sos-conservative}} & \thepage & Not the disjoint TRS-sum convention from Part12 \\
\phantomsection\label{entry:symbol-registry-05908}F & \texttt{\detokenize{F}} & Finite ranked alphabet with at least one constant & p09 finite-tree-algebra & n/a & \hyperref[sec:p09-tree-algebra-input]{\ref*{sec:p09-tree-algebra-input}} & \thepage & TATA source role; distinct from final-state notation in the generic handbook tuple Original precedence 'atomic' retained as semantic descriptor; n/a is not an invented rank. \\
\phantomsection\label{entry:symbol-registry-05909}\ensuremath{F_n} & \texttt{\detokenize{F_n}} & Symbols of arity n & p09 finite-tree-algebra & n/a & \hyperref[sec:p09-tree-algebra-input]{\ref*{sec:p09-tree-algebra-input}} & \thepage & Ordered arities; F\_0 is nonempty in this selected source treatment Original precedence 'atomic' retained as semantic descriptor; n/a is not an invented rank. \\
\phantomsection\label{entry:symbol-registry-05910}T(F) & \texttt{\detokenize{T(F)}} & Ground term algebra carrier & p09 finite-tree-algebra & n/a & \hyperref[sec:p09-tree-algebra-input]{\ref*{sec:p09-tree-algebra-input}} & \thepage & Finite ordered trees; not the infinite successor structure Original precedence 'atomic' retained as semantic descriptor; n/a is not an invented rank. \\
\phantomsection\label{entry:symbol-registry-05911}C(F) & \texttt{\detokenize{C(F)}} & One-hole ground contexts & p09 finite-tree-algebra & n/a & \hyperref[sec:p09-tree-algebra-input]{\ref*{sec:p09-tree-algebra-input}} & \thepage & Exactly one hole occurrence and ground side branches Original precedence 'atomic' retained as semantic descriptor; n/a is not an invented rank. \\
\phantomsection\label{entry:symbol-registry-05912}C[t] & \texttt{\detokenize{C[t]}} & Context plugging operation & p09 finite-tree-algebra & n/a & \hyperref[sec:p09-tree-algebra-input]{\ref*{sec:p09-tree-algebra-input}} & \thepage & Also context composition when the inserted object has one hole Original precedence 'application' retained as semantic descriptor; n/a is not an invented rank. \\
\phantomsection\label{entry:symbol-registry-05913}x & \texttt{\detokenize{x}} & Distinguished single context hole & p09 finite-tree-algebra contexts & n/a & \hyperref[sec:p09-tree-algebra-input]{\ref*{sec:p09-tree-algebra-input}} & \thepage & Not an independently varying node marker in this local use Original precedence 'atomic' retained as semantic descriptor; n/a is not an invented rank. \\
\phantomsection\label{entry:symbol-registry-05914}A & \texttt{\detokenize{A}} & F-algebra with total operations & p09 finite-tree-algebra & n/a & \hyperref[def:p09-finite-f-algebra]{\ref*{def:p09-finite-f-algebra}} & \thepage & No attribution to unheld original1968algebra conventions Original precedence 'atomic' retained as semantic descriptor; n/a is not an invented rank. \\
\phantomsection\label{entry:symbol-registry-05915}|A| & \texttt{\detokenize{|A|}} & Carrier of an F-algebra in the source algebra notation & p09 finite-tree-algebra & n/a & \hyperref[def:p09-finite-f-algebra]{\ref*{def:p09-finite-f-algebra}} & \thepage & Published typed carrier role explicitly distinguished from cardinality; cardinality written operatorname card Original precedence 'atomic' retained as semantic descriptor; n/a is not an invented rank. \\
\phantomsection\label{entry:symbol-registry-05916}A' & \texttt{\detokenize{A'}} & Designated accepting subset of an algebra carrier & p09 finite-tree-algebra & n/a & \hyperref[def:p09-finite-f-algebra]{\ref*{def:p09-finite-f-algebra}} & \thepage & Recognition data rather than extra signature operation Original precedence 'atomic' retained as semantic descriptor; n/a is not an invented rank. \\
\phantomsection\label{entry:symbol-registry-05917}\ensuremath{f^A} & \texttt{\detokenize{f^A}} & Total algebra operation for an n-ary ranked symbol & p09 finite-tree-algebra & n/a & \hyperref[def:p09-finite-f-algebra]{\ref*{def:p09-finite-f-algebra}} & \thepage & Nullary case is a designated carrier element Original precedence 'application' retained as semantic descriptor; n/a is not an invented rank. \\
\phantomsection\label{entry:symbol-registry-05918}\ensuremath{\varphi_A} & \texttt{\detokenize{\varphi_A}} & Unique ground-term evaluation homomorphism & p09 finite-tree-algebra & n/a & \hyperref[thm:p09-tree-algebra-evaluation]{\ref*{thm:p09-tree-algebra-evaluation}} & \thepage & Its image is the constant-generated accessible subalgebra Original precedence 'application' retained as semantic descriptor; n/a is not an invented rank. \\
\phantomsection\label{entry:symbol-registry-05919}\ensuremath{\sim} & \texttt{\detokenize{\sim}} & Arbitrary equivalence or congruence on ground terms & p09 finite-tree-algebra & relation & \hyperref[def:p09-tree-context-congruence]{\ref*{def:p09-tree-context-congruence}} & \thepage & Operator compatibility required separately \\
\phantomsection\label{entry:symbol-registry-05920}\ensuremath{\equiv_L} & \texttt{\detokenize{\equiv_L}} & Source syntactic context congruence of a tree language & p09 finite-tree-algebra & relation & \hyperref[def:p09-tree-context-congruence]{\ref*{def:p09-tree-context-congruence}} & \thepage & Published syntax retained with explicit local type; not elementary equivalence \\
\phantomsection\label{entry:symbol-registry-05921}{}[t] & \texttt{\detokenize{[t]}} & Congruence class of a ground term & p09 finite-tree-algebra & n/a & \hyperref[thm:p09-tree-algebra-recognition]{\ref*{thm:p09-tree-algebra-recognition}} & \thepage & Relation fixed in each use Original precedence 'atomic' retained as semantic descriptor; n/a is not an invented rank. \\
\phantomsection\label{entry:symbol-registry-05922}\ensuremath{Q_{\min}} & \texttt{\detokenize{Q_{\min}}} & Finite carrier of the canonical context quotient & p09 finite-tree-algebra & n/a & \hyperref[thm:p09-tree-algebra-recognition]{\ref*{thm:p09-tree-algebra-recognition}} & \thepage & Nonempty since source alphabet has a constant Original precedence 'atomic' retained as semantic descriptor; n/a is not an invented rank. \\
\phantomsection\label{entry:symbol-registry-05923}\ensuremath{Q_{\min,f}} & \texttt{\detokenize{Q_{\min,f}}} & Designated classes of accepted terms in the canonical quotient & p09 finite-tree-algebra & n/a & \hyperref[thm:p09-tree-algebra-recognition]{\ref*{thm:p09-tree-algebra-recognition}} & \thepage & Membership is well defined by saturation Original precedence 'atomic' retained as semantic descriptor; n/a is not an invented rank. \\
\phantomsection\label{entry:symbol-registry-05924}\ensuremath{\pi} & \texttt{\detokenize{\pi}} & Canonical homomorphism from reachable carrier to syntactic quotient & p09 finite-tree-algebra & n/a & \hyperref[thm:p09-tree-algebra-minimal]{\ref*{thm:p09-tree-algebra-minimal}} & \thepage & Not a track-erasure projection Original precedence 'application' retained as semantic descriptor; n/a is not an invented rank. \\
\phantomsection\label{entry:symbol-registry-05925}\ensuremath{P_j} & \texttt{\detokenize{P_j}} & Round-j equivalence relation in partition refinement & p09 finite-tree-algebra & relation & \hyperref[sec:p09-tree-algebra-refinement]{\ref*{sec:p09-tree-algebra-refinement}} & \thepage & Uses old relation throughout each simultaneous round \\
\phantomsection\label{entry:symbol-registry-05926}\ensuremath{\delta_f} & \texttt{\detokenize{\delta_f}} & Total deterministic bottom-up transition operation & p09 finite-tree-algebra & n/a & \hyperref[sec:p09-tree-algebra-refinement]{\ref*{sec:p09-tree-algebra-refinement}} & \thepage & Same ranked operation type as f\^{}A Original precedence 'application' retained as semantic descriptor; n/a is not an invented rank. \\
\phantomsection\label{entry:symbol-registry-05927}\ensuremath{C^\delta} & \texttt{\detokenize{C^\delta}} & Total state function induced by a ground context & p09 finite-tree-algebra annex & n/a & \hyperref[annex:p09.tree-algebra.refinement]{\ref*{annex:p09.tree-algebra.refinement}} & \thepage & Hole-state evaluation; ground side children have witnesses by reducedness Original precedence 'application' retained as semantic descriptor; n/a is not an invented rank. \\
\phantomsection\label{entry:symbol-registry-05928}\ensuremath{\tau} & \texttt{\detokenize{\tau}} & Elementary one-hole state operation & p09 finite-tree-algebra annex & n/a & \hyperref[annex:p09.tree-algebra.refinement]{\ref*{annex:p09.tree-algebra.refinement}} & \thepage & One coordinate varies; all other states fixed Original precedence 'application' retained as semantic descriptor; n/a is not an invented rank. \\
\phantomsection\label{entry:symbol-registry-05929}R & \texttt{\detokenize{R}} & Reachable generated subalgebra of a recognition algebra & p09 finite-tree-algebra & n/a & \hyperref[thm:p09-tree-algebra-minimal]{\ref*{thm:p09-tree-algebra-minimal}} & \thepage & Closed under all total operations Original precedence 'atomic' retained as semantic descriptor; n/a is not an invented rank. \\
\phantomsection\label{entry:symbol-registry-07283}\(V\) & \texttt{\detokenize{V}} & finite module or splicing alphabet & p09-recombination-development & n/a & \hyperref[def:p09-recombination-splicing]{\ref*{def:p09-recombination-splicing}} & \thepage & Source-local alphabets; no inference that real-parameter module set is finite Original precedence descriptor 'atomic' retained; no rank inferred. Actual displayed reader TeX form, not the raw ASCII alias. Blank macro filled with the glyph's TeX spelling. \\
\phantomsection\label{entry:symbol-registry-07284}\(\#\) & \texttt{\detokenize{\#}} & first and second recognition-cut separator & p09-splicing & n/a & \hyperref[def:p09-recombination-splicing]{\ref*{def:p09-recombination-splicing}} & \thepage & Literal source rule syntax; not plant line-width command Original precedence descriptor 'atomic' retained; no rank inferred. Actual displayed reader TeX form, not the raw ASCII alias. Blank macro filled with the glyph's TeX spelling. \\
\phantomsection\label{entry:symbol-registry-07285}\(\$\) & \texttt{\detokenize{\$}} & separator between the two splicing sites & p09-splicing & n/a & \hyperref[def:p09-recombination-splicing]{\ref*{def:p09-recombination-splicing}} & \thepage & Literal source rule syntax; not mathematical-mode delimiter in the interpretation Original precedence descriptor 'atomic' retained; no rank inferred. Actual displayed reader TeX form, not the raw ASCII alias. Blank macro filled with the glyph's TeX spelling. \\
\phantomsection\label{entry:symbol-registry-07286}\(u_i\) & \texttt{\detokenize{u_i}} & four rule handles & p09-splicing & n/a & \hyperref[def:p09-recombination-splicing]{\ref*{def:p09-recombination-splicing}} & \thepage & Handles may be empty; sites are u\_1u\_2 and u\_3u\_4 Original precedence descriptor 'atomic' retained; no rank inferred. Actual displayed reader TeX form, not the raw ASCII alias. Blank macro filled with the glyph's TeX spelling. \\
\phantomsection\label{entry:symbol-registry-07287}\(\sigma\) & \texttt{\detokenize{\sigma}} & H scheme and its language operation & p09-splicing & n/a & \hyperref[def:p09-recombination-splicing]{\ref*{def:p09-recombination-splicing}} & \thepage & Locally rendered \(\sigma\); separate operation on ordered pair versus input set Original precedence descriptor 'application' retained; no rank inferred. Actual displayed reader TeX form, not the raw ASCII alias. Blank macro filled with the glyph's TeX spelling. Mechanical TeX fix: math command in the notes cell wrapped in inline math. \\
\phantomsection\label{entry:symbol-registry-07288}\(\vdash_r\) & \texttt{\detokenize{\vdash_r}} & oriented splicing relation & p09-splicing & relation & \hyperref[def:p09-recombination-splicing]{\ref*{def:p09-recombination-splicing}} & \thepage & Locally rendered \(\vdash_r\); not logical derivability and not complementary-product semantics Actual displayed reader TeX form, not the raw ASCII alias. Blank macro filled with the glyph's TeX spelling. Mechanical TeX fix: math command in the notes cell wrapped in inline math. \\
\phantomsection\label{entry:symbol-registry-07289}\(A/B\) & \texttt{\detokenize{A/}\allowbreak\detokenize{B}} & right language quotient & p09-splicing & n/a & \hyperref[sec:p09-recombination-splicing]{\ref*{sec:p09-recombination-splicing}} & \thepage & Set of p with pq in A for some q in B; slash is not parametric numerical division Original precedence descriptor 'application' retained; no rank inferred. Actual displayed reader TeX form, not the raw ASCII alias. Blank macro filled with the glyph's TeX spelling. \\
\phantomsection\label{entry:symbol-registry-07290}\(B\backslash A\) & \texttt{\detokenize{B\backslash A}} & left language quotient & p09-splicing & n/a & \hyperref[sec:p09-recombination-splicing]{\ref*{sec:p09-recombination-splicing}} & \thepage & Set of s with qs in A for some q in B Original precedence descriptor 'application' retained; no rank inferred. Actual displayed reader TeX form, not the raw ASCII alias. Blank macro filled with the glyph's TeX spelling. \\
\phantomsection\label{entry:symbol-registry-07291}\(\sigma^*\) & \texttt{\detokenize{\sigma^*}} & iterated least H-scheme closure & p09-splicing & n/a & \hyperref[sec:p09-recombination-iteration]{\ref*{sec:p09-recombination-iteration}} & \thepage & Locally rendered \(\sigma^*\); input availability preserved; not consuming finite multiset Original precedence descriptor 'application' retained; no rank inferred. Actual displayed reader TeX form, not the raw ASCII alias. Blank macro filled with the glyph's TeX spelling. Mechanical TeX fix: math command in the notes cell wrapped in inline math. \\
\phantomsection\label{entry:symbol-registry-07292}\(V_B\) & \texttt{\detokenize{V_B}} & module alphabet extended by branch brackets & p09-development & n/a & \hyperref[sec:p09-development-architecture]{\ref*{sec:p09-development-architecture}} & \thepage & Brackets are not module labels; ordered well-nested representation Original precedence descriptor 'atomic' retained; no rank inferred. Actual displayed reader TeX form, not the raw ASCII alias. Blank macro filled with the glyph's TeX spelling. \\
\phantomsection\label{entry:symbol-registry-07293}\(P\) & \texttt{\detokenize{P}} & turtle position vector & \hyperref[sec:p09-development-architecture]{\ref*{sec:p09-development-architecture}} & n/a & \hyperref[sec:p09-development-architecture]{\ref*{sec:p09-development-architecture}} & \thepage & Source-local U denotes up vector; require H cross L equals U Original precedence descriptor 'atomic' retained; no rank inferred. Blank macro filled with the glyph's TeX spelling. Family listing, not one atomic glyph. Scope narrowed to the heading of first use (turtle interpretation); production-set P is declared in sec:p09-development-parametric. \\
\phantomsection\label{entry:symbol-registry-07294}\(H\) & \texttt{\detokenize{H}} & orthonormal heading vector & \hyperref[sec:p09-development-architecture]{\ref*{sec:p09-development-architecture}} & n/a & \hyperref[sec:p09-development-architecture]{\ref*{sec:p09-development-architecture}} & \thepage & Source-local U denotes up vector; require H cross L equals U Original precedence descriptor 'atomic' retained; no rank inferred. Blank macro filled with the glyph's TeX spelling. Family listing, not one atomic glyph. Scope narrowed to the heading of first use (turtle interpretation); production-set P is declared in sec:p09-development-parametric. \\
\phantomsection\label{entry:symbol-registry-07295}\(L\) & \texttt{\detokenize{L}} & orthonormal left vector & \hyperref[sec:p09-development-architecture]{\ref*{sec:p09-development-architecture}} & n/a & \hyperref[sec:p09-development-architecture]{\ref*{sec:p09-development-architecture}} & \thepage & Source-local U denotes up vector; require H cross L equals U Original precedence descriptor 'atomic' retained; no rank inferred. Blank macro filled with the glyph's TeX spelling. Family listing, not one atomic glyph. Scope narrowed to the heading of first use (turtle interpretation); production-set P is declared in sec:p09-development-parametric. \\
\phantomsection\label{entry:symbol-registry-07296}\(U\) & \texttt{\detokenize{U}} & orthonormal up vector & \hyperref[sec:p09-development-architecture]{\ref*{sec:p09-development-architecture}} & n/a & \hyperref[sec:p09-development-architecture]{\ref*{sec:p09-development-architecture}} & \thepage & Source-local U denotes up vector; require H cross L equals U Original precedence descriptor 'atomic' retained; no rank inferred. Blank macro filled with the glyph's TeX spelling. Family listing, not one atomic glyph. Scope narrowed to the heading of first use (turtle interpretation); production-set P is declared in sec:p09-development-parametric. \\
\phantomsection\label{entry:symbol-registry-07297}\(M\) & \texttt{\detokenize{M}} & parametric module set V times R star & p09-development & n/a & \hyperref[def:p09-development-parametric]{\ref*{def:p09-development-parametric}} & \thepage & Locally rendered V times \(\mathbb R^*\); finite letter alphabet does not make M finite Original precedence descriptor 'atomic' retained; no rank inferred. Actual displayed reader TeX form, not the raw ASCII alias. Blank macro filled with the glyph's TeX spelling. Mechanical TeX fix: math command in the notes cell wrapped in inline math. \\
\phantomsection\label{entry:symbol-registry-07298}\(G\) & \texttt{\detokenize{G}} & parametric OL-system tuple & p09-development & n/a & \hyperref[def:p09-development-parametric]{\ref*{def:p09-development-parametric}} & \thepage & Source-local model; not a sequential grammar tuple Original precedence descriptor 'atomic' retained; no rank inferred. Actual displayed reader TeX form, not the raw ASCII alias. Blank macro filled with the glyph's TeX spelling. \\
\phantomsection\label{entry:symbol-registry-07299}\(E\) & \texttt{\detokenize{E}} & set of formal parameter names & p09-development & n/a & \hyperref[def:p09-development-parametric]{\ref*{def:p09-development-parametric}} & \thepage & Distinct from expectation operator \(\mathbb E\) Original precedence descriptor 'atomic' retained; no rank inferred. Actual displayed reader TeX form, not the raw ASCII alias. Blank macro filled with the glyph's TeX spelling. Mechanical TeX fix: math command in the notes cell wrapped in inline math. \\
\phantomsection\label{entry:symbol-registry-07300}\(\omega\) & \texttt{\detokenize{\omega}} & nonempty parametric axiom & p09-development & n/a & \hyperref[def:p09-development-parametric]{\ref*{def:p09-development-parametric}} & \thepage & Locally rendered \(\omega\) Original precedence descriptor 'atomic' retained; no rank inferred. Actual displayed reader TeX form, not the raw ASCII alias. Blank macro filled with the glyph's TeX spelling. Mechanical TeX fix: math command in the notes cell wrapped in inline math. \\
\phantomsection\label{entry:symbol-registry-07301}\(P\) & \texttt{\detokenize{P}} & finite production set & P9 development production sets & n/a & \hyperref[def:p09-development-parametric]{\ref*{def:p09-development-parametric}} & \thepage & Source-local P differs from turtle position P in its explicitly separate typed paragraph Original precedence descriptor 'atomic' retained; no rank inferred. Actual displayed reader TeX form, not the raw ASCII alias. Blank macro filled with the glyph's TeX spelling. Production set P: declared in sec:p09-development-parametric, used through sec:p09-development-stochastic and its annex proof; a named scope covers those regions without overlapping the turtle P of sec:p09-development-architecture. \\
\phantomsection\label{entry:symbol-registry-07302}\(h\) & \texttt{\detokenize{h}} & morphism induced by nonparametric DOL successors & p09-development & n/a & \hyperref[thm:p09-development-morphic-expressiveness]{\ref*{thm:p09-development-morphic-expressiveness}} & \thepage & Editorial finite-alphabet subcase; not a claim that parametric module space is finite Original precedence descriptor 'application' retained; no rank inferred. Actual displayed reader TeX form, not the raw ASCII alias. Blank macro filled with the glyph's TeX spelling. \\
\phantomsection\label{entry:symbol-registry-07303}\(\pi\) & \texttt{\detokenize{\pi}} & stochastic production-factor map & p09-development-stochastic & n/a & \hyperref[sec:p09-development-stochastic]{\ref*{sec:p09-development-stochastic}} & \thepage & Locally rendered \(\pi\); distinguished from circle constant in crown-area formula by type Original precedence descriptor 'application' retained; no rank inferred. Actual displayed reader TeX form, not the raw ASCII alias. Blank macro filled with the glyph's TeX spelling. Mechanical TeX fix: math command in the notes cell wrapped in inline math. \\
\phantomsection\label{entry:symbol-registry-07304}\(N_k\) & \texttt{\detokenize{N_k}} & current apex count & p09-development-stochastic & n/a & \hyperref[thm:p09-development-stochastic-kernel]{\ref*{thm:p09-development-stochastic-kernel}} & \thepage & Expectation result uses marginals; variance additionally requires conditional independence Original precedence descriptor 'atomic' retained; no rank inferred. Blank macro filled with the glyph's TeX spelling. Family listing, not one atomic glyph. \\
\phantomsection\label{entry:symbol-registry-07305}\(p_k\) & \texttt{\detokenize{p_k}} & conditional binary branching probability & p09-development-stochastic & n/a & \hyperref[thm:p09-development-stochastic-kernel]{\ref*{thm:p09-development-stochastic-kernel}} & \thepage & Expectation result uses marginals; variance additionally requires conditional independence Original precedence descriptor 'atomic' retained; no rank inferred. Blank macro filled with the glyph's TeX spelling. Family listing, not one atomic glyph. \\
\phantomsection\label{entry:symbol-registry-07306}\(\sigma_{\min}\) & \texttt{\detokenize{\sigma_{\min}}} & minimum exposed leaf area per viable shoot & p09-development-stochastic & n/a & \hyperref[sec:p09-development-stochastic]{\ref*{sec:p09-development-stochastic}} & \thepage & Locally rendered \(\sigma_{\min}\); not the H-scheme operation Original precedence descriptor 'atomic' retained; no rank inferred. Actual displayed reader TeX form, not the raw ASCII alias. Blank macro filled with the glyph's TeX spelling. Mechanical TeX fix: math command in the notes cell wrapped in inline math. \\
\phantomsection\label{entry:symbol-registry-07307}\(\lambda\) & \texttt{\detokenize{\lambda}} & Source-native empty word in oriented splicing and finite-word development & P9 recombination development native words & meta & \hyperref[def:p09-recombination-splicing]{\ref*{def:p09-recombination-splicing}} & \thepage & HeadPaunPixton1997RecombinationIntake source-local word notation restored; not a term binder. Original epsilon intake remains immutable evidence and is superseded for this active displayed role. \\
\phantomsection\label{entry:symbol-registry-07360}\(\lambda\) & \texttt{\detokenize{\lambda}} & empty word & p09-recombination-iteration & n/a & \hyperref[sec:p09-recombination-iterated-closures]{\ref*{sec:p09-recombination-iterated-closures}} & \thepage & Published source-local notation; first packet uses epsilon and states the comparison Original precedence descriptor 'atomic' retained; no rank inferred. Blank macro filled with the glyph's TeX spelling. \\
\phantomsection\label{entry:symbol-registry-07361}\(|w|\) & \texttt{\detokenize{|w|}} & word length & p09-recombination-iteration & n/a & \hyperref[sec:p09-recombination-iterated-closures]{\ref*{sec:p09-recombination-iterated-closures}} & \thepage & Published source-local typed role; not global cardinality notation Original precedence descriptor 'application' retained; no rank inferred. Blank macro filled with the glyph's TeX spelling. \\
\phantomsection\label{entry:symbol-registry-07362}\(\sigma\) & \texttt{\detokenize{\sigma}} & oriented H scheme and set-splicing operation & p09-recombination-iteration & n/a & \hyperref[sec:p09-recombination-iterated-closures]{\ref*{sec:p09-recombination-iterated-closures}} & \thepage & Finite rule set in this theorem; inputs remain reusable Original precedence descriptor 'application' retained; no rank inferred. Blank macro filled with the glyph's TeX spelling. \\
\phantomsection\label{entry:symbol-registry-07363}\(\alpha_r\) & \texttt{\detokenize{\alpha_r}} & left recognition site u\_1u\_2 & p09-recombination-iteration & n/a & \hyperref[sec:p09-recombination-iterated-closures]{\ref*{sec:p09-recombination-iterated-closures}} & \thepage & Source-local role; may be empty Original precedence descriptor 'atomic' retained; no rank inferred. Blank macro filled with the glyph's TeX spelling. \\
\phantomsection\label{entry:symbol-registry-07364}\(\gamma_r\) & \texttt{\detokenize{\gamma_r}} & right recognition site u\_3u\_4 & p09-recombination-iteration & n/a & \hyperref[sec:p09-recombination-iterated-closures]{\ref*{sec:p09-recombination-iterated-closures}} & \thepage & Source-local role; may be empty Original precedence descriptor 'atomic' retained; no rank inferred. Blank macro filled with the glyph's TeX spelling. \\
\phantomsection\label{entry:symbol-registry-07365}\(\beta_r\) & \texttt{\detokenize{\beta_r}} & recombinant handle word u\_1u\_4 & p09-recombination-iteration & n/a & \hyperref[sec:p09-recombination-iterated-closures]{\ref*{sec:p09-recombination-iterated-closures}} & \thepage & Source-local role; not the whole output Original precedence descriptor 'atomic' retained; no rank inferred. Blank macro filled with the glyph's TeX spelling. \\
\phantomsection\label{entry:symbol-registry-07366}\(FA\) & \texttt{\detokenize{FA}} & language family with stated closure properties & p09-recombination-iteration & n/a & \hyperref[thm:p09-recombination-full-afl-iteration]{\ref*{thm:p09-recombination-full-afl-iteration}} & \thepage & Contains every regular language; regular quotients and regular-control substitution explicitly assumed Original precedence descriptor 'atomic' retained; no rank inferred. Blank macro filled with the glyph's TeX spelling. \\
\phantomsection\label{entry:symbol-registry-07367}\(\mathopen{(\!(}_r\) & \texttt{\detokenize{\mathopen{(\!(}_r}} & left-nest opening token & p09-recombination-iteration & n/a & \hyperref[sec:p09-recombination-iterated-closures]{\ref*{sec:p09-recombination-iterated-closures}} & \thepage & Published doubled-round glyph; closes with square closing token of same label Original precedence descriptor 'atomic' retained; no rank inferred. Blank macro filled with the glyph's TeX spelling. \\
\phantomsection\label{entry:symbol-registry-07368}\(\mathclose{)\!)}_r\) & \texttt{\detokenize{\mathclose{)\!)}_r}} & right-nest closing token & p09-recombination-iteration & n/a & \hyperref[sec:p09-recombination-iterated-closures]{\ref*{sec:p09-recombination-iterated-closures}} & \thepage & Published doubled-round glyph; opens with square opening token of same label Original precedence descriptor 'atomic' retained; no rank inferred. Blank macro filled with the glyph's TeX spelling. \\
\phantomsection\label{entry:symbol-registry-07369}\([_r\) & \texttt{\detokenize{[_r}} & right-nest opening token & p09-recombination-iteration & n/a & \hyperref[sec:p09-recombination-iterated-closures]{\ref*{sec:p09-recombination-iterated-closures}} & \thepage & Not a plant branch delimiter in this scoped history alphabet Original precedence descriptor 'atomic' retained; no rank inferred. Blank macro filled with the glyph's TeX spelling. \\
\phantomsection\label{entry:symbol-registry-07370}\(]_r\) & \texttt{\detokenize{]_r}} & left-nest closing token & p09-recombination-iteration & n/a & \hyperref[sec:p09-recombination-iterated-closures]{\ref*{sec:p09-recombination-iterated-closures}} & \thepage & Not ordinary square-to-square matching Original precedence descriptor 'atomic' retained; no rank inferred. Blank macro filled with the glyph's TeX spelling. \\
\phantomsection\label{entry:symbol-registry-07371}\(N\) & \texttt{\detokenize{N}} & well-nested history-word domain & p09-recombination-iteration & n/a & \hyperref[sec:p09-recombination-iterated-closures]{\ref*{sec:p09-recombination-iterated-closures}} & \thepage & Not natural numbers; balanced factors are essential to justified splicing Original precedence descriptor 'atomic' retained; no rank inferred. Blank macro filled with the glyph's TeX spelling. \\
\phantomsection\label{entry:symbol-registry-07372}\(E\) & \texttt{\detokenize{E}} & whole-nest erasure on N & p09-recombination-iteration & n/a & \hyperref[sec:p09-recombination-iterated-closures]{\ref*{sec:p09-recombination-iterated-closures}} & \thepage & Deletes maximal nests including contents; not merely bracket erasure or a free-monoid morphism on every token string Original precedence descriptor 'application' retained; no rank inferred. Blank macro filled with the glyph's TeX spelling. \\
\phantomsection\label{entry:symbol-registry-07373}\(J\) & \texttt{\detokenize{J}} & least justified-splicing history language & p09-recombination-iteration & n/a & \hyperref[sec:p09-recombination-iterated-closures]{\ref*{sec:p09-recombination-iterated-closures}} & \thepage & Not assumed to belong to FA Original precedence descriptor 'atomic' retained; no rank inferred. Blank macro filled with the glyph's TeX spelling. \\
\phantomsection\label{entry:symbol-registry-07374}\(\$\) & \texttt{\detokenize{\$}} & separator within a splicing-rule word & p09-recombination-iteration & n/a & \hyperref[sec:p09-recombination-iterated-closures]{\ref*{sec:p09-recombination-iterated-closures}} & \thepage & Published role inherited from ordinary H scheme Original precedence descriptor 'atomic' retained; no rank inferred. Blank macro filled with the glyph's TeX spelling. \\
\phantomsection\label{entry:symbol-registry-07375}\(\$\) & \texttt{\detokenize{\$}} & dummy index denoting original input-word birth & p09-recombination-iteration & n/a & \hyperref[sec:p09-recombination-iterated-closures]{\ref*{sec:p09-recombination-iterated-closures}} & \thepage & Second published typed role explicitly declared; dummy index cannot equal a rule word Original precedence descriptor 'atomic' retained; no rank inferred. Blank macro filled with the glyph's TeX spelling. \\
\phantomsection\label{entry:symbol-registry-07376}\(L_r\) & \texttt{\detokenize{L_r}} & base language L at dummy index or singleton beta\_r at a rule & p09-recombination-iteration & n/a & \hyperref[sec:p09-recombination-iterated-closures]{\ref*{sec:p09-recombination-iterated-closures}} & \thepage & Birth-language role; not the indexed quotient cell Original precedence descriptor 'application' retained; no rank inferred. Blank macro filled with the glyph's TeX spelling. \\
\phantomsection\label{entry:symbol-registry-07377}\(F_r(s,t)\) & \texttt{\detokenize{F_r(s,}\allowbreak\detokenize{t)}} & bounded-gap and terminal suffix-transition language & p09-recombination-iteration & n/a & \hyperref[sec:p09-recombination-iterated-closures]{\ref*{sec:p09-recombination-iterated-closures}} & \thepage & Defined by actual history occurrence; finite semantic tables do not by themselves give a terminating discovery algorithm Original precedence descriptor 'application' retained; no rank inferred. Blank macro filled with the glyph's TeX spelling. \\
\phantomsection\label{entry:symbol-registry-07378}\(S_r(s)\) & \texttt{\detokenize{S_r(s)}} & regular suffix language from finite rule-index paths & p09-recombination-iteration & n/a & \hyperref[sec:p09-recombination-iterated-closures]{\ref*{sec:p09-recombination-iterated-closures}} & \thepage & Dummy entry corrected to singleton empty word; literal source empty-set reading recorded in provenance Original precedence descriptor 'application' retained; no rank inferred. Blank macro filled with the glyph's TeX spelling. \\
\phantomsection\label{entry:symbol-registry-07379}\(P_r(s)\) & \texttt{\detokenize{P_r(s)}} & regular prefix language defined by explicit reversal & p09-recombination-iteration & n/a & \hyperref[sec:p09-recombination-iterated-closures]{\ref*{sec:p09-recombination-iterated-closures}} & \thepage & Dummy entry singleton empty word; no reversal closure of FA is assumed Original precedence descriptor 'application' retained; no rank inferred. Blank macro filled with the glyph's TeX spelling. \\
\phantomsection\label{entry:symbol-registry-07380}\(L_r(s,t)\) & \texttt{\detokenize{L_r(s,}\allowbreak\detokenize{t)}} & double-quotient visible-block language & p09-recombination-iteration & n/a & \hyperref[sec:p09-recombination-iterated-closures]{\ref*{sec:p09-recombination-iterated-closures}} & \thepage & Uses P\_r(s) on left and S\_r(t) on right; source's repeated s-index recorded in correction provenance Original precedence descriptor 'application' retained; no rank inferred. Blank macro filled with the glyph's TeX spelling. \\
\phantomsection\label{entry:symbol-registry-07381}\(K\) & \texttt{\detokenize{K}} & regular-control substitution language of compatible visible pieces & p09-recombination-iteration & n/a & \hyperref[sec:p09-recombination-iterated-closures]{\ref*{sec:p09-recombination-iterated-closures}} & \thepage & Both inclusions with E(J) proved; not assumed saturation Original precedence descriptor 'atomic' retained; no rank inferred. Blank macro filled with the glyph's TeX spelling. \\
\phantomsection\label{entry:symbol-registry-07382}\(\rho\) & \texttt{\detokenize{\rho}} & word-reversal and bracket-rule isomorphism & p09-recombination-iteration & n/a & \hyperref[annex:p09.recombination-full-afl-iteration]{\ref*{annex:p09.recombination-full-afl-iteration}} & \thepage & Reverses order and swaps left/right tokens; transformed rule r\^{}rho reverses and exchanges handles Original precedence descriptor 'application' retained; no rank inferred. Blank macro filled with the glyph's TeX spelling. \\
\phantomsection\label{entry:symbol-registry-07383}\(\operatorname{rad}(\sigma)\) & \texttt{\detokenize{\operatorname{rad}(\sigma)}} & maximum single-handle length & p09-recombination-iteration & n/a & \hyperref[sec:p09-recombination-iterated-closures]{\ref*{sec:p09-recombination-iterated-closures}} & \thepage & Empty rule set assigned radius zero; site lengths at most twice radius Original precedence descriptor 'application' retained; no rank inferred. Blank macro filled with the glyph's TeX spelling. \\
\phantomsection\label{entry:symbol-registry-07384}\(L_k\) & \texttt{\detokenize{L_k}} & radius-witness and wrapper-language symbol family member \(L_k\) & p09-recombination-iteration & n/a & \hyperref[thm:p09-recombination-radius-hierarchy]{\ref*{thm:p09-recombination-radius-hierarchy}} & \thepage & L\_k has fixed 4k b-count; wrapper endpoint uses L=a+ rather than asserting the complete general lemma Original precedence descriptor 'atomic' retained; no rank inferred. Blank macro filled with the glyph's TeX spelling. Family listing, not one atomic glyph. \\
\phantomsection\label{entry:symbol-registry-07385}\(W_1\) & \texttt{\detokenize{W_1}} & radius-witness and wrapper-language symbol family member \(W_1\) & p09-recombination-iteration & n/a & \hyperref[thm:p09-recombination-radius-hierarchy]{\ref*{thm:p09-recombination-radius-hierarchy}} & \thepage & L\_k has fixed 4k b-count; wrapper endpoint uses L=a+ rather than asserting the complete general lemma Original precedence descriptor 'atomic' retained; no rank inferred. Blank macro filled with the glyph's TeX spelling. Family listing, not one atomic glyph. \\
\phantomsection\label{entry:symbol-registry-07386}\(W_2\) & \texttt{\detokenize{W_2}} & radius-witness and wrapper-language symbol family member \(W_2\) & p09-recombination-iteration & n/a & \hyperref[thm:p09-recombination-radius-hierarchy]{\ref*{thm:p09-recombination-radius-hierarchy}} & \thepage & L\_k has fixed 4k b-count; wrapper endpoint uses L=a+ rather than asserting the complete general lemma Original precedence descriptor 'atomic' retained; no rank inferred. Blank macro filled with the glyph's TeX spelling. Family listing, not one atomic glyph. \\
\phantomsection\label{entry:symbol-registry-07387}\(W\) & \texttt{\detokenize{W}} & radius-witness and wrapper-language symbol family member \(W\) & p09-recombination-iteration & n/a & \hyperref[thm:p09-recombination-radius-hierarchy]{\ref*{thm:p09-recombination-radius-hierarchy}} & \thepage & L\_k has fixed 4k b-count; wrapper endpoint uses L=a+ rather than asserting the complete general lemma Original precedence descriptor 'atomic' retained; no rank inferred. Blank macro filled with the glyph's TeX spelling. Family listing, not one atomic glyph. \\
\phantomsection\label{entry:symbol-registry-07388}\(M\) & \texttt{\detokenize{M}} & consuming multiset with integer or infinite multiplicities & p09-recombination-iteration-multiset & n/a & \hyperref[sec:p09-recombination-iterated-closures]{\ref*{sec:p09-recombination-iterated-closures}} & \thepage & Distinct from reusable set semantics; published next-section event consumes two occurrences and emits both complementary products Original precedence descriptor 'application' retained; no rank inferred. Blank macro filled with the glyph's TeX spelling. \\
\phantomsection\label{entry:symbol-registry-07389}\(T\) & \texttt{\detokenize{T}} & single nonterminal of the explicit grammar for N & p09-recombination-iteration-nest-proof & n/a & \hyperref[annex:p09.recombination-nest-domain]{\ref*{annex:p09.recombination-nest-domain}} & \thepage & Editorial proof grammar; not a terminal alphabet or source rule-index set Original precedence descriptor 'atomic' retained; no rank inferred. Blank macro filled with the glyph's TeX spelling. \\
\phantomsection\label{entry:symbol-registry-07708}\(G=(N,T,S,P)\) & \texttt{\detokenize{G=(N,}\allowbreak\detokenize{T,}\allowbreak\detokenize{S,}\allowbreak\detokenize{P)}} & input finite type-0 grammar & p09-grammar-normalization & n/a & \hyperref[def:p09-short-grammar-source]{\ref*{def:p09-short-grammar-source}} & \thepage & Grouped structural record; original terminal alphabet T is unchanged by compiler Original precedence descriptor 'atomic' retained; no rank inferred. Blank macro filled with the glyph's TeX spelling. \\
\phantomsection\label{entry:symbol-registry-07709}\(\lambda\) & \texttt{\detokenize{\lambda}} & empty word & p09-grammar-normalization & n/a & \hyperref[def:p09-short-grammar-source]{\ref*{def:p09-short-grammar-source}} & \thepage & Literal empty-word preservation; not only equality modulo epsilon Original precedence descriptor 'atomic' retained; no rank inferred. Blank macro filled with the glyph's TeX spelling. \\
\phantomsection\label{entry:symbol-registry-07710}\(|w|\) & \texttt{\detokenize{|w|}} & word length & p09-grammar-normalization & n/a & \hyperref[def:p09-short-grammar-source]{\ref*{def:p09-short-grammar-source}} & \thepage & Source-local role retained in the host construction Original precedence descriptor 'application' retained; no rank inferred. Blank macro filled with the glyph's TeX spelling. \\
\phantomsection\label{entry:symbol-registry-07711}\(A_a\) & \texttt{\detokenize{A_a}} & one fresh nonterminal encoding original symbol a & p09-grammar-normalization & n/a & \hyperref[def:p09-short-grammar-source]{\ref*{def:p09-short-grammar-source}} & \thepage & New host compiler name; both original terminals and nonterminals encoded; not a two-letter word Original precedence descriptor 'atomic' retained; no rank inferred. Blank macro filled with the glyph's TeX spelling. \\
\phantomsection\label{entry:symbol-registry-07712}\(\widehat w\) & \texttt{\detokenize{\widehat w}} & letter-to-letter encoded word & p09-grammar-normalization & n/a & \hyperref[def:p09-short-grammar-source]{\ref*{def:p09-short-grammar-source}} & \thepage & New host word map sends each a to A\_a; no erasing observation is used Original precedence descriptor 'application' retained; no rank inferred. Blank macro filled with the glyph's TeX spelling. \\
\phantomsection\label{entry:symbol-registry-07713}\(r\) & \texttt{\detokenize{r}} & production label and its two finite words and lengths symbol family member \(r\) & p09-grammar-normalization & n/a & \hyperref[def:p09-short-grammar-compiler]{\ref*{def:p09-short-grammar-compiler}} & \thepage & Grouped host indexing family; each range explicitly finite Original precedence descriptor 'atomic' retained; no rank inferred. Explicit displayed family; member types and common actual first-use exposed, not one atomic glyph. Blank macro filled with the glyph's TeX spelling. Family listing, not one atomic glyph. \\
\phantomsection\label{entry:symbol-registry-07714}\(u_r\) & \texttt{\detokenize{u_r}} & production label and its two finite words and lengths symbol family member \(u_r\) & p09-grammar-normalization & n/a & \hyperref[def:p09-short-grammar-compiler]{\ref*{def:p09-short-grammar-compiler}} & \thepage & Grouped host indexing family; each range explicitly finite Original precedence descriptor 'atomic' retained; no rank inferred. Explicit displayed family; member types and common actual first-use exposed, not one atomic glyph. Blank macro filled with the glyph's TeX spelling. Family listing, not one atomic glyph. \\
\phantomsection\label{entry:symbol-registry-07715}\(v_r\) & \texttt{\detokenize{v_r}} & production label and its two finite words and lengths symbol family member \(v_r\) & p09-grammar-normalization & n/a & \hyperref[def:p09-short-grammar-compiler]{\ref*{def:p09-short-grammar-compiler}} & \thepage & Grouped host indexing family; each range explicitly finite Original precedence descriptor 'atomic' retained; no rank inferred. Explicit displayed family; member types and common actual first-use exposed, not one atomic glyph. Blank macro filled with the glyph's TeX spelling. Family listing, not one atomic glyph. \\
\phantomsection\label{entry:symbol-registry-07716}\(m_r\) & \texttt{\detokenize{m_r}} & production label and its two finite words and lengths symbol family member \(m_r\) & p09-grammar-normalization & n/a & \hyperref[def:p09-short-grammar-compiler]{\ref*{def:p09-short-grammar-compiler}} & \thepage & Grouped host indexing family; each range explicitly finite Original precedence descriptor 'atomic' retained; no rank inferred. Explicit displayed family; member types and common actual first-use exposed, not one atomic glyph. Blank macro filled with the glyph's TeX spelling. Family listing, not one atomic glyph. \\
\phantomsection\label{entry:symbol-registry-07717}\(n_r\) & \texttt{\detokenize{n_r}} & production label and its two finite words and lengths symbol family member \(n_r\) & p09-grammar-normalization & n/a & \hyperref[def:p09-short-grammar-compiler]{\ref*{def:p09-short-grammar-compiler}} & \thepage & Grouped host indexing family; each range explicitly finite Original precedence descriptor 'atomic' retained; no rank inferred. Explicit displayed family; member types and common actual first-use exposed, not one atomic glyph. Blank macro filled with the glyph's TeX spelling. Family listing, not one atomic glyph. \\
\phantomsection\label{entry:symbol-registry-07718}\(S^\dagger\) & \texttt{\detokenize{S^\dagger}} & compiled start symbol and boundary and initialization symbols symbol family member \(S^\dagger\) & p09-grammar-normalization & n/a & \hyperref[def:p09-short-grammar-compiler]{\ref*{def:p09-short-grammar-compiler}} & \thepage & Pairwise fresh nonterminals; grouped names are not atomic glyph counts Original precedence descriptor 'atomic' retained; no rank inferred. Explicit displayed family; member types and common actual first-use exposed, not one atomic glyph. Blank macro filled with the glyph's TeX spelling. Family listing, not one atomic glyph. \\
\phantomsection\label{entry:symbol-registry-07719}\(L\) & \texttt{\detokenize{L}} & compiled start symbol and boundary and initialization symbols symbol family member \(L\) & p09-grammar-normalization & n/a & \hyperref[def:p09-short-grammar-compiler]{\ref*{def:p09-short-grammar-compiler}} & \thepage & Pairwise fresh nonterminals; grouped names are not atomic glyph counts Original precedence descriptor 'atomic' retained; no rank inferred. Explicit displayed family; member types and common actual first-use exposed, not one atomic glyph. Blank macro filled with the glyph's TeX spelling. Family listing, not one atomic glyph. \\
\phantomsection\label{entry:symbol-registry-07720}\(R\) & \texttt{\detokenize{R}} & compiled start symbol and boundary and initialization symbols symbol family member \(R\) & p09-grammar-normalization & n/a & \hyperref[def:p09-short-grammar-compiler]{\ref*{def:p09-short-grammar-compiler}} & \thepage & Pairwise fresh nonterminals; grouped names are not atomic glyph counts Original precedence descriptor 'atomic' retained; no rank inferred. Explicit displayed family; member types and common actual first-use exposed, not one atomic glyph. Blank macro filled with the glyph's TeX spelling. Family listing, not one atomic glyph. \\
\phantomsection\label{entry:symbol-registry-07721}\(I\) & \texttt{\detokenize{I}} & compiled start symbol and boundary and initialization symbols symbol family member \(I\) & p09-grammar-normalization & n/a & \hyperref[def:p09-short-grammar-compiler]{\ref*{def:p09-short-grammar-compiler}} & \thepage & Pairwise fresh nonterminals; grouped names are not atomic glyph counts Original precedence descriptor 'atomic' retained; no rank inferred. Explicit displayed family; member types and common actual first-use exposed, not one atomic glyph. Blank macro filled with the glyph's TeX spelling. Family listing, not one atomic glyph. \\
\phantomsection\label{entry:symbol-registry-07722}\(J\) & \texttt{\detokenize{J}} & compiled start symbol and boundary and initialization symbols symbol family member \(J\) & p09-grammar-normalization & n/a & \hyperref[def:p09-short-grammar-compiler]{\ref*{def:p09-short-grammar-compiler}} & \thepage & Pairwise fresh nonterminals; grouped names are not atomic glyph counts Original precedence descriptor 'atomic' retained; no rank inferred. Explicit displayed family; member types and common actual first-use exposed, not one atomic glyph. Blank macro filled with the glyph's TeX spelling. Family listing, not one atomic glyph. \\
\phantomsection\label{entry:symbol-registry-07723}\(Q\) & \texttt{\detokenize{Q}} & unique ready-mode cursor nonterminal & p09-grammar-normalization & n/a & \hyperref[def:p09-short-grammar-compiler]{\ref*{def:p09-short-grammar-compiler}} & \thepage & Not quantifier Q of SSAT or question Q; actual compiler-local standing Original precedence descriptor 'atomic' retained; no rank inferred. Blank macro filled with the glyph's TeX spelling. \\
\phantomsection\label{entry:symbol-registry-07724}\(Q_{r,i}\) & \texttt{\detokenize{Q_{r,i}}} & unique match-state nonterminal after i consumed letters & p09-grammar-normalization & n/a & \hyperref[def:p09-short-grammar-compiler]{\ref*{def:p09-short-grammar-compiler}} & \thepage & One indexed letter; wrong match is stuck and cannot restore ready mode Original precedence descriptor 'atomic' retained; no rank inferred. Blank macro filled with the glyph's TeX spelling. \\
\phantomsection\label{entry:symbol-registry-07725}\(E_{r,j}\) & \texttt{\detokenize{E_{r,j}}} & unique emission-state nonterminal after j emitted coded letters & p09-grammar-normalization & n/a & \hyperref[def:p09-short-grammar-compiler]{\ref*{def:p09-short-grammar-compiler}} & \thepage & One indexed letter; no other-mode production can interleave Original precedence descriptor 'atomic' retained; no rank inferred. Blank macro filled with the glyph's TeX spelling. \\
\phantomsection\label{entry:symbol-registry-07726}\(D\) & \texttt{\detokenize{D}} & unique cleanup-mode nonterminal & p09-grammar-normalization & n/a & \hyperref[def:p09-short-grammar-compiler]{\ref*{def:p09-short-grammar-compiler}} & \thepage & Not data symbol or domain; moves only across encoded original terminals Original precedence descriptor 'atomic' retained; no rank inferred. Blank macro filled with the glyph's TeX spelling. \\
\phantomsection\label{entry:symbol-registry-07727}\(G^\dagger\) & \texttt{\detokenize{G^\dagger}} & compiled grammar and its finite nonterminal and production sets symbol family member \(G^\dagger\) & p09-grammar-normalization & n/a & \hyperref[def:p09-short-grammar-compiler]{\ref*{def:p09-short-grammar-compiler}} & \thepage & Grouped host output record; original T retained literally Original precedence descriptor 'atomic' retained; no rank inferred. Explicit displayed family; member types and common actual first-use exposed, not one atomic glyph. Blank macro filled with the glyph's TeX spelling. Family listing, not one atomic glyph. \\
\phantomsection\label{entry:symbol-registry-07728}\(N^\dagger\) & \texttt{\detokenize{N^\dagger}} & compiled grammar and its finite nonterminal and production sets symbol family member \(N^\dagger\) & p09-grammar-normalization & n/a & \hyperref[def:p09-short-grammar-compiler]{\ref*{def:p09-short-grammar-compiler}} & \thepage & Grouped host output record; original T retained literally Original precedence descriptor 'atomic' retained; no rank inferred. Explicit displayed family; member types and common actual first-use exposed, not one atomic glyph. Blank macro filled with the glyph's TeX spelling. Family listing, not one atomic glyph. \\
\phantomsection\label{entry:symbol-registry-07729}\(P^\dagger\) & \texttt{\detokenize{P^\dagger}} & compiled grammar and its finite nonterminal and production sets symbol family member \(P^\dagger\) & p09-grammar-normalization & n/a & \hyperref[def:p09-short-grammar-compiler]{\ref*{def:p09-short-grammar-compiler}} & \thepage & Grouped host output record; original T retained literally Original precedence descriptor 'atomic' retained; no rank inferred. Explicit displayed family; member types and common actual first-use exposed, not one atomic glyph. Blank macro filled with the glyph's TeX spelling. Family listing, not one atomic glyph. \\
\phantomsection\label{entry:symbol-registry-07730}\(\Longrightarrow_P^*\) & \texttt{\detokenize{\Longrightarrow_P^*}} & ordinary context replacement derivation of input grammar & p09-grammar-normalization & relation & \hyperref[def:p09-short-grammar-source]{\ref*{def:p09-short-grammar-source}} & \thepage & Typed grammar derivation; distinct from consuming multiset event relation Blank macro filled with the glyph's TeX spelling. \\
\phantomsection\label{entry:symbol-registry-07731}\(\sigma\) & \texttt{\detokenize{\sigma}} & finite consuming H scheme obtained from compiled grammar & p09-grammar-normalization-consequence & n/a & \hyperref[thm:p09-consuming-universality-short-compiler]{\ref*{thm:p09-consuming-universality-short-compiler}} & \thepage & Exact repaired prior consuming construction retained; no supply convention changed Original precedence descriptor 'application' retained; no rank inferred. Blank macro filled with the glyph's TeX spelling. \\
\phantomsection\label{entry:symbol-registry-07732}\(M\) & \texttt{\detokenize{M}} & finite-support initial multiset with infinite helper multiplicities & p09-grammar-normalization-consequence & n/a & \hyperref[thm:p09-consuming-universality-short-compiler]{\ref*{thm:p09-consuming-universality-short-compiler}} & \thepage & Finite support is not finite total population; terminal extraction T* remains necessary Original precedence descriptor 'application' retained; no rank inferred. Blank macro filled with the glyph's TeX spelling. \\
\phantomsection\label{entry:symbol-registry-07733}\(\lambda\) & \texttt{\detokenize{\lambda}} & empty word & p09-recombination-computation & n/a & \hyperref[sec:p09-recombination-computation]{\ref*{sec:p09-recombination-computation}} & \thepage & Published source-local notation; literal empty-word cases are treated Original precedence descriptor 'atomic' retained; no rank inferred. Blank macro filled with the glyph's TeX spelling. \\
\phantomsection\label{entry:symbol-registry-07734}\(|w|\) & \texttt{\detokenize{|w|}} & word length & p09-recombination-computation & n/a & \hyperref[sec:p09-recombination-computation]{\ref*{sec:p09-recombination-computation}} & \thepage & Source-local type distinct from cardinality role Original precedence descriptor 'application' retained; no rank inferred. Blank macro filled with the glyph's TeX spelling. \\
\phantomsection\label{entry:symbol-registry-07735}\(G=(N,T,S,P)\) & \texttt{\detokenize{G=(N,}\allowbreak\detokenize{T,}\allowbreak\detokenize{S,}\allowbreak\detokenize{P)}} & type-0 or right-linear grammar with explicitly distinct hypotheses & p09-recombination-computation & n/a & \hyperref[def:p09-recombination-regular-simulation]{\ref*{def:p09-recombination-regular-simulation}} & \thepage & Grouped record; not four atomic glyph counts; each grammar regime redeclared before use Original precedence descriptor 'atomic' retained; no rank inferred. Blank macro filled with the glyph's TeX spelling. \\
\phantomsection\label{entry:symbol-registry-07736}\(U\) & \texttt{\detokenize{U}} & finite working alphabet & p09-recombination-computation & n/a & \hyperref[def:p09-recombination-regular-simulation]{\ref*{def:p09-recombination-regular-simulation}} & \thepage & First simulation includes B; consuming simulation explicitly redeclares U=N union T Original precedence descriptor 'atomic' retained; no rank inferred. Blank macro filled with the glyph's TeX spelling. \\
\phantomsection\label{entry:symbol-registry-07737}\(X\) & \texttt{\detokenize{X}} & fresh markers of the regular-rule set construction symbol family member \(X\) & P9 recombination regular-simulation markers & n/a & \hyperref[def:p09-recombination-regular-simulation]{\ref*{def:p09-recombination-regular-simulation}} & \thepage & Grouped published marker family; distinct from grammar nonterminals by standing Original precedence descriptor 'atomic' retained; no rank inferred. Explicit displayed family; member types and common actual first-use exposed, not one atomic glyph. Blank macro filled with the glyph's TeX spelling. Family listing, not one atomic glyph. Scope narrowed from the whole computation subsection to the regular-rule simulation where the markers X, X', B, Y, Z are declared and used (paragraph of def:p09-recombination-regular-simulation and annex:p09.recombination-regular-simulation); the whole-subsection scope overlapped the separate consuming-construction marker Y. \\
\phantomsection\label{entry:symbol-registry-07738}\(X'\) & \texttt{\detokenize{X'}} & fresh markers of the regular-rule set construction symbol family member \(X'\) & P9 recombination regular-simulation markers & n/a & \hyperref[def:p09-recombination-regular-simulation]{\ref*{def:p09-recombination-regular-simulation}} & \thepage & Grouped published marker family; distinct from grammar nonterminals by standing Original precedence descriptor 'atomic' retained; no rank inferred. Explicit displayed family; member types and common actual first-use exposed, not one atomic glyph. Blank macro filled with the glyph's TeX spelling. Family listing, not one atomic glyph. Scope narrowed from the whole computation subsection to the regular-rule simulation where the markers X, X', B, Y, Z are declared and used (paragraph of def:p09-recombination-regular-simulation and annex:p09.recombination-regular-simulation); the whole-subsection scope overlapped the separate consuming-construction marker Y. \\
\phantomsection\label{entry:symbol-registry-07739}\(B\) & \texttt{\detokenize{B}} & fresh markers of the regular-rule set construction symbol family member \(B\) & P9 recombination regular-simulation markers & n/a & \hyperref[def:p09-recombination-regular-simulation]{\ref*{def:p09-recombination-regular-simulation}} & \thepage & Grouped published marker family; distinct from grammar nonterminals by standing Original precedence descriptor 'atomic' retained; no rank inferred. Explicit displayed family; member types and common actual first-use exposed, not one atomic glyph. Blank macro filled with the glyph's TeX spelling. Family listing, not one atomic glyph. Scope narrowed from the whole computation subsection to the regular-rule simulation where the markers X, X', B, Y, Z are declared and used (paragraph of def:p09-recombination-regular-simulation and annex:p09.recombination-regular-simulation); the whole-subsection scope overlapped the separate consuming-construction marker Y. \\
\phantomsection\label{entry:symbol-registry-07740}\(Y\) & \texttt{\detokenize{Y}} & fresh markers of the regular-rule set construction symbol family member \(Y\) & P9 recombination regular-simulation markers & n/a & \hyperref[def:p09-recombination-regular-simulation]{\ref*{def:p09-recombination-regular-simulation}} & \thepage & Grouped published marker family; distinct from grammar nonterminals by standing Original precedence descriptor 'atomic' retained; no rank inferred. Explicit displayed family; member types and common actual first-use exposed, not one atomic glyph. Blank macro filled with the glyph's TeX spelling. Family listing, not one atomic glyph. Scope narrowed from the whole computation subsection to the regular-rule simulation where the markers X, X', B, Y, Z are declared and used (paragraph of def:p09-recombination-regular-simulation and annex:p09.recombination-regular-simulation); the whole-subsection scope overlapped the separate consuming-construction marker Y. \\
\phantomsection\label{entry:symbol-registry-07741}\(Z\) & \texttt{\detokenize{Z}} & fresh markers of the regular-rule set construction symbol family member \(Z\) & P9 recombination regular-simulation markers & n/a & \hyperref[def:p09-recombination-regular-simulation]{\ref*{def:p09-recombination-regular-simulation}} & \thepage & Grouped published marker family; distinct from grammar nonterminals by standing Original precedence descriptor 'atomic' retained; no rank inferred. Explicit displayed family; member types and common actual first-use exposed, not one atomic glyph. Blank macro filled with the glyph's TeX spelling. Family listing, not one atomic glyph. Scope narrowed from the whole computation subsection to the regular-rule simulation where the markers X, X', B, Y, Z are declared and used (paragraph of def:p09-recombination-regular-simulation and annex:p09.recombination-regular-simulation); the whole-subsection scope overlapped the separate consuming-construction marker Y. \\
\phantomsection\label{entry:symbol-registry-07742}\(Y_\alpha\) & \texttt{\detokenize{Y_\alpha}} & single tagged symbol storing a removed letter & p09-recombination-computation & n/a & \hyperref[def:p09-recombination-regular-simulation]{\ref*{def:p09-recombination-regular-simulation}} & \thepage & Not the two-letter word Y alpha Original precedence descriptor 'atomic' retained; no rank inferred. Blank macro filled with the glyph's TeX spelling. \\
\phantomsection\label{entry:symbol-registry-07743}\(\sigma\) & \texttt{\detokenize{\sigma}} & H scheme and its splicing operations & p09-recombination-computation & n/a & \hyperref[def:p09-recombination-regular-simulation]{\ref*{def:p09-recombination-regular-simulation}} & \thepage & Reusable-set and consuming-multiset argument types explicitly separated Original precedence descriptor 'application' retained; no rank inferred. Blank macro filled with the glyph's TeX spelling. \\
\phantomsection\label{entry:symbol-registry-07744}\(L_0\) & \texttt{\detokenize{L_0}} & finite seed language & p09-recombination-computation & n/a & \hyperref[def:p09-recombination-regular-simulation]{\ref*{def:p09-recombination-regular-simulation}} & \thepage & Separate regular-simulation and transition-coding seeds redeclared Original precedence descriptor 'atomic' retained; no rank inferred. Blank macro filled with the glyph's TeX spelling. \\
\phantomsection\label{entry:symbol-registry-07745}\(L_1\) & \texttt{\detokenize{L_1}} & wrapper language family symbol family member \(L_1\) & p09-recombination-computation-wrapper & n/a & \hyperref[def:p09-recombination-wrapper]{\ref*{def:p09-recombination-wrapper}} & \thepage & Grouped source names; terminal extraction cannot imply unrestricted generated equality Original precedence descriptor 'atomic' retained; no rank inferred. Explicit displayed family; member types and common actual first-use exposed, not one atomic glyph. Blank macro filled with the glyph's TeX spelling. Family listing, not one atomic glyph. \\
\phantomsection\label{entry:symbol-registry-07746}\(L_2\) & \texttt{\detokenize{L_2}} & wrapper language family symbol family member \(L_2\) & p09-recombination-computation-wrapper & n/a & \hyperref[def:p09-recombination-wrapper]{\ref*{def:p09-recombination-wrapper}} & \thepage & Grouped source names; terminal extraction cannot imply unrestricted generated equality Original precedence descriptor 'atomic' retained; no rank inferred. Explicit displayed family; member types and common actual first-use exposed, not one atomic glyph. Blank macro filled with the glyph's TeX spelling. Family listing, not one atomic glyph. \\
\phantomsection\label{entry:symbol-registry-07747}\(L'\) & \texttt{\detokenize{L'}} & wrapper language family symbol family member \(L'\) & p09-recombination-computation-wrapper & n/a & \hyperref[def:p09-recombination-wrapper]{\ref*{def:p09-recombination-wrapper}} & \thepage & Grouped source names; terminal extraction cannot imply unrestricted generated equality Original precedence descriptor 'atomic' retained; no rank inferred. Explicit displayed family; member types and common actual first-use exposed, not one atomic glyph. Blank macro filled with the glyph's TeX spelling. Family listing, not one atomic glyph. \\
\phantomsection\label{entry:symbol-registry-07748}\(h\) & \texttt{\detokenize{h}} & endpoint-erasing monoid morphism restricted to cV*d & p09-recombination-computation-wrapper & n/a & \hyperref[def:p09-recombination-wrapper]{\ref*{def:p09-recombination-wrapper}} & \thepage & Source-local type; erases exactly two endpoints on the stated restricted domain Original precedence descriptor 'application' retained; no rank inferred. Blank macro filled with the glyph's TeX spelling. \\
\phantomsection\label{entry:symbol-registry-07749}\([X,a,Y]\) & \texttt{\detokenize{[X,}\allowbreak\detokenize{a,}\allowbreak\detokenize{Y]}} & single transition and final-transition letters symbol family member \([X,a,Y]\) & p09-recombination-computation-coding & n/a & \hyperref[def:p09-recombination-transition-coding]{\ref*{def:p09-recombination-transition-coding}} & \thepage & Each tuple-looking printed name is one letter; not a three-letter word Original precedence descriptor 'atomic' retained; no rank inferred. Explicit displayed family; member types and common actual first-use exposed, not one atomic glyph. Blank macro filled with the glyph's TeX spelling. Family listing, not one atomic glyph. \\
\phantomsection\label{entry:symbol-registry-07750}\([X,a,*]\) & \texttt{\detokenize{[X,}\allowbreak\detokenize{a,}\allowbreak\detokenize{*]}} & single transition and final-transition letters symbol family member \([X,a,*]\) & p09-recombination-computation-coding & n/a & \hyperref[def:p09-recombination-transition-coding]{\ref*{def:p09-recombination-transition-coding}} & \thepage & Each tuple-looking printed name is one letter; not a three-letter word Original precedence descriptor 'atomic' retained; no rank inferred. Explicit displayed family; member types and common actual first-use exposed, not one atomic glyph. Blank macro filled with the glyph's TeX spelling. Family listing, not one atomic glyph. \\
\phantomsection\label{entry:symbol-registry-07751}\(h\) & \texttt{\detokenize{h}} & letter-to-letter coding from transition alphabet V to T & p09-recombination-computation-coding & n/a & \hyperref[def:p09-recombination-transition-coding]{\ref*{def:p09-recombination-transition-coding}} & \thepage & Distinct domain/type from endpoint erasure; local declaration preserves source glyph Original precedence descriptor 'application' retained; no rank inferred. Blank macro filled with the glyph's TeX spelling. \\
\phantomsection\label{entry:symbol-registry-07752}\(M\) & \texttt{\detokenize{M}} & consuming multiset with natural or infinite multiplicity & p09-recombination-computation-consuming & n/a & \hyperref[def:p09-recombination-consuming-simulation]{\ref*{def:p09-recombination-consuming-simulation}} & \thepage & Single active occurrence; reservoirs infinite; both complementary outputs counted Original precedence descriptor 'application' retained; no rank inferred. Blank macro filled with the glyph's TeX spelling. \\
\phantomsection\label{entry:symbol-registry-07753}\(X_1\) & \texttt{\detokenize{X_1}} & fresh consuming-construction marker family symbol family member \(X_1\) & P9 recombination consuming-simulation markers & n/a & \hyperref[def:p09-recombination-consuming-simulation]{\ref*{def:p09-recombination-consuming-simulation}} & \thepage & Grouped source marker family; two-copy boundary blocks are words Original precedence descriptor 'atomic' retained; no rank inferred. Explicit displayed family; member types and common actual first-use exposed, not one atomic glyph. Blank macro filled with the glyph's TeX spelling. Family listing, not one atomic glyph. Scope narrowed to the consuming construction where the markers X\_1, X\_2, Y, Z\_1, Z\_2 are declared and used (paragraph of def:p09-recombination-consuming-simulation, its annex proof, and the short-compiler composition proof that builds the start word with them); the previous scope reached the regular-rule simulation's distinct marker Y. \\
\phantomsection\label{entry:symbol-registry-07754}\(X_2\) & \texttt{\detokenize{X_2}} & fresh consuming-construction marker family symbol family member \(X_2\) & P9 recombination consuming-simulation markers & n/a & \hyperref[def:p09-recombination-consuming-simulation]{\ref*{def:p09-recombination-consuming-simulation}} & \thepage & Grouped source marker family; two-copy boundary blocks are words Original precedence descriptor 'atomic' retained; no rank inferred. Explicit displayed family; member types and common actual first-use exposed, not one atomic glyph. Blank macro filled with the glyph's TeX spelling. Family listing, not one atomic glyph. Scope narrowed to the consuming construction where the markers X\_1, X\_2, Y, Z\_1, Z\_2 are declared and used (paragraph of def:p09-recombination-consuming-simulation, its annex proof, and the short-compiler composition proof that builds the start word with them); the previous scope reached the regular-rule simulation's distinct marker Y. \\
\phantomsection\label{entry:symbol-registry-07755}\(Y\) & \texttt{\detokenize{Y}} & fresh consuming-construction marker family symbol family member \(Y\) & P9 recombination consuming-simulation markers & n/a & \hyperref[def:p09-recombination-consuming-simulation]{\ref*{def:p09-recombination-consuming-simulation}} & \thepage & Grouped source marker family; two-copy boundary blocks are words Original precedence descriptor 'atomic' retained; no rank inferred. Explicit displayed family; member types and common actual first-use exposed, not one atomic glyph. Blank macro filled with the glyph's TeX spelling. Family listing, not one atomic glyph. Scope narrowed to the consuming construction where the markers X\_1, X\_2, Y, Z\_1, Z\_2 are declared and used (paragraph of def:p09-recombination-consuming-simulation, its annex proof, and the short-compiler composition proof that builds the start word with them); the previous scope reached the regular-rule simulation's distinct marker Y. \\
\phantomsection\label{entry:symbol-registry-07756}\(Z_1\) & \texttt{\detokenize{Z_1}} & fresh consuming-construction marker family symbol family member \(Z_1\) & P9 recombination consuming-simulation markers & n/a & \hyperref[def:p09-recombination-consuming-simulation]{\ref*{def:p09-recombination-consuming-simulation}} & \thepage & Grouped source marker family; two-copy boundary blocks are words Original precedence descriptor 'atomic' retained; no rank inferred. Explicit displayed family; member types and common actual first-use exposed, not one atomic glyph. Blank macro filled with the glyph's TeX spelling. Family listing, not one atomic glyph. Scope narrowed to the consuming construction where the markers X\_1, X\_2, Y, Z\_1, Z\_2 are declared and used (paragraph of def:p09-recombination-consuming-simulation, its annex proof, and the short-compiler composition proof that builds the start word with them); the previous scope reached the regular-rule simulation's distinct marker Y. \\
\phantomsection\label{entry:symbol-registry-07757}\(Z_2\) & \texttt{\detokenize{Z_2}} & fresh consuming-construction marker family symbol family member \(Z_2\) & P9 recombination consuming-simulation markers & n/a & \hyperref[def:p09-recombination-consuming-simulation]{\ref*{def:p09-recombination-consuming-simulation}} & \thepage & Grouped source marker family; two-copy boundary blocks are words Original precedence descriptor 'atomic' retained; no rank inferred. Explicit displayed family; member types and common actual first-use exposed, not one atomic glyph. Blank macro filled with the glyph's TeX spelling. Family listing, not one atomic glyph. Scope narrowed to the consuming construction where the markers X\_1, X\_2, Y, Z\_1, Z\_2 are declared and used (paragraph of def:p09-recombination-consuming-simulation, its annex proof, and the short-compiler composition proof that builds the start word with them); the previous scope reached the regular-rule simulation's distinct marker Y. \\
\phantomsection\label{entry:symbol-registry-07758}\((r)\) & \texttt{\detokenize{(r)}} & single grammar-production tag letters symbol family member \((r)\) & p09-recombination-computation-consuming & n/a & \hyperref[def:p09-recombination-consuming-simulation]{\ref*{def:p09-recombination-consuming-simulation}} & \thepage & Distinct from punctuation token sequences and from history nesting brackets Original precedence descriptor 'atomic' retained; no rank inferred. Explicit displayed family; member types and common actual first-use exposed, not one atomic glyph. Blank macro filled with the glyph's TeX spelling. Family listing, not one atomic glyph. \\
\phantomsection\label{entry:symbol-registry-07759}\([r]\) & \texttt{\detokenize{[r]}} & single grammar-production tag letters symbol family member \([r]\) & p09-recombination-computation-consuming & n/a & \hyperref[def:p09-recombination-consuming-simulation]{\ref*{def:p09-recombination-consuming-simulation}} & \thepage & Distinct from punctuation token sequences and from history nesting brackets Original precedence descriptor 'atomic' retained; no rank inferred. Explicit displayed family; member types and common actual first-use exposed, not one atomic glyph. Blank macro filled with the glyph's TeX spelling. Family listing, not one atomic glyph. \\
\phantomsection\label{entry:symbol-registry-07760}\(K\) & \texttt{\detokenize{K}} & finite injective production-label set & p09-recombination-computation-consuming & n/a & \hyperref[def:p09-recombination-consuming-simulation]{\ref*{def:p09-recombination-consuming-simulation}} & \thepage & Not compatible-history language K of the earlier packet Original precedence descriptor 'atomic' retained; no rank inferred. Blank macro filled with the glyph's TeX spelling. \\
\phantomsection\label{entry:symbol-registry-07761}\(\delta_1\) & \texttt{\detokenize{\delta_1}} & left-context guard letters symbol family member \(\delta_1\) & p09-recombination-computation-consuming & n/a & \hyperref[def:p09-recombination-consuming-simulation]{\ref*{def:p09-recombination-consuming-simulation}} & \thepage & Typed U or X1; guards distinguish live fragments from reservoirs Original precedence descriptor 'atomic' retained; no rank inferred. Explicit displayed family; member types and common actual first-use exposed, not one atomic glyph. Blank macro filled with the glyph's TeX spelling. Family listing, not one atomic glyph. \\
\phantomsection\label{entry:symbol-registry-07762}\(\delta_2\) & \texttt{\detokenize{\delta_2}} & left-context guard letters symbol family member \(\delta_2\) & p09-recombination-computation-consuming & n/a & \hyperref[def:p09-recombination-consuming-simulation]{\ref*{def:p09-recombination-consuming-simulation}} & \thepage & Typed U or X1; guards distinguish live fragments from reservoirs Original precedence descriptor 'atomic' retained; no rank inferred. Explicit displayed family; member types and common actual first-use exposed, not one atomic glyph. Blank macro filled with the glyph's TeX spelling. Family listing, not one atomic glyph. \\
\phantomsection\label{entry:symbol-registry-07763}\(\delta\) & \texttt{\detokenize{\delta}} & left-context guard letters symbol family member \(\delta\) & p09-recombination-computation-consuming & n/a & \hyperref[def:p09-recombination-consuming-simulation]{\ref*{def:p09-recombination-consuming-simulation}} & \thepage & Typed U or X1; guards distinguish live fragments from reservoirs Original precedence descriptor 'atomic' retained; no rank inferred. Explicit displayed family; member types and common actual first-use exposed, not one atomic glyph. Blank macro filled with the glyph's TeX spelling. Family listing, not one atomic glyph. \\
\phantomsection\label{entry:symbol-registry-07764}\(\beta_1\) & \texttt{\detokenize{\beta_1}} & right-context guard letters symbol family member \(\beta_1\) & p09-recombination-computation-consuming & n/a & \hyperref[def:p09-recombination-consuming-simulation]{\ref*{def:p09-recombination-consuming-simulation}} & \thepage & Published group 5 retains beta3 and beta1 in U; repaired group adds only exact X2 squared boundary Original precedence descriptor 'atomic' retained; no rank inferred. Explicit displayed family; member types and common actual first-use exposed, not one atomic glyph. Blank macro filled with the glyph's TeX spelling. Family listing, not one atomic glyph. \\
\phantomsection\label{entry:symbol-registry-07765}\(\beta_2\) & \texttt{\detokenize{\beta_2}} & right-context guard letters symbol family member \(\beta_2\) & p09-recombination-computation-consuming & n/a & \hyperref[def:p09-recombination-consuming-simulation]{\ref*{def:p09-recombination-consuming-simulation}} & \thepage & Published group 5 retains beta3 and beta1 in U; repaired group adds only exact X2 squared boundary Original precedence descriptor 'atomic' retained; no rank inferred. Explicit displayed family; member types and common actual first-use exposed, not one atomic glyph. Blank macro filled with the glyph's TeX spelling. Family listing, not one atomic glyph. \\
\phantomsection\label{entry:symbol-registry-07766}\(\beta_3\) & \texttt{\detokenize{\beta_3}} & right-context guard letters symbol family member \(\beta_3\) & p09-recombination-computation-consuming & n/a & \hyperref[def:p09-recombination-consuming-simulation]{\ref*{def:p09-recombination-consuming-simulation}} & \thepage & Published group 5 retains beta3 and beta1 in U; repaired group adds only exact X2 squared boundary Original precedence descriptor 'atomic' retained; no rank inferred. Explicit displayed family; member types and common actual first-use exposed, not one atomic glyph. Blank macro filled with the glyph's TeX spelling. Family listing, not one atomic glyph. \\
\phantomsection\label{entry:symbol-registry-07767}\(\beta\) & \texttt{\detokenize{\beta}} & right-context guard letters symbol family member \(\beta\) & p09-recombination-computation-consuming & n/a & \hyperref[def:p09-recombination-consuming-simulation]{\ref*{def:p09-recombination-consuming-simulation}} & \thepage & Published group 5 retains beta3 and beta1 in U; repaired group adds only exact X2 squared boundary Original precedence descriptor 'atomic' retained; no rank inferred. Explicit displayed family; member types and common actual first-use exposed, not one atomic glyph. Blank macro filled with the glyph's TeX spelling. Family listing, not one atomic glyph. \\
\phantomsection\label{entry:symbol-registry-07768}\(\Longrightarrow_\sigma\) & \texttt{\detokenize{\Longrightarrow_\sigma}} & one consuming two-output event & p09-recombination-computation-consuming & relation & \hyperref[sec:p09-recombination-computation]{\ref*{sec:p09-recombination-computation}} & \thepage & Source double-arrow relation; actual multiplicity precondition retained Blank macro filled with the glyph's TeX spelling. \\
\phantomsection\label{entry:symbol-registry-07769}\(C(p,q)\) & \texttt{\detokenize{C(p,}\allowbreak\detokenize{q)}} & editorial abbreviation for completed word X1 squared p Y q X2 squared & p09-recombination-computation-consuming-proof & n/a & \hyperref[annex:p09.recombination-consuming-simulation]{\ref*{annex:p09.recombination-consuming-simulation}} & \thepage & Explicit proof abbreviation; not claimed to be a published source name Original precedence descriptor 'application' retained; no rank inferred. Blank macro filled with the glyph's TeX spelling. \\
\phantomsection\label{entry:symbol-registry-07770}\(\operatorname{cnt}(M)\) & \texttt{\detokenize{\operatorname{cnt}(M)}} & finite population size & p09-recombination-computation-population & n/a & \hyperref[def:p09-recombination-finite-population]{\ref*{def:p09-recombination-finite-population}} & \thepage & Editorial conservation quantity needed for empty-word state space Original precedence descriptor 'application' retained; no rank inferred. Blank macro filled with the glyph's TeX spelling. \\
\phantomsection\label{entry:symbol-registry-07771}\(\operatorname{wht}(M)\) & \texttt{\detokenize{\operatorname{wht}(M)}} & total word-length weight & p09-recombination-computation-population & n/a & \hyperref[def:p09-recombination-finite-population]{\ref*{def:p09-recombination-finite-population}} & \thepage & Published weight role; theorem also requires finite counts Original precedence descriptor 'application' retained; no rank inferred. Blank macro filled with the glyph's TeX spelling. \\
\phantomsection\label{entry:symbol-registry-07772}\(U_{n,W}\) & \texttt{\detokenize{U_{n,W}}} & finite count-and-weight constrained multiset state space & p09-recombination-computation-population & n/a & \hyperref[def:p09-recombination-finite-population]{\ref*{def:p09-recombination-finite-population}} & \thepage & Editorial repaired vertex set; not alphabet U of grammar constructions Original precedence descriptor 'atomic' retained; no rank inferred. Blank macro filled with the glyph's TeX spelling. \\
\phantomsection\label{entry:symbol-registry-07773}\(\Gamma(\sigma,M)\) & \texttt{\detokenize{\Gamma(\sigma,}\allowbreak\detokenize{M)}} & directed multiset evolution graph & p09-recombination-computation-population & n/a & \hyperref[def:p09-recombination-finite-population]{\ref*{def:p09-recombination-finite-population}} & \thepage & Source graph name with corrected count-constrained vertex domain and explicit effective rule-membership hypothesis Original precedence descriptor 'application' retained; no rank inferred. Blank macro filled with the glyph's TeX spelling. \\
\end{longtable}
\endgroup
\Needspace{.4\textheight}
\subsection*{Stable part P10 scopes}
\begingroup\scriptsize\setlength{\tabcolsep}{3pt}
\begin{longtable}{@{}p{0.9in}p{1.3in}p{1.0in}p{0.5in}p{0.3in}p{0.4in}p{0.35in}p{1.65in}@{}}
\toprule
\textbf{Glyph} & \textbf{Macro} & \textbf{Role} & \textbf{Scope} & \textbf{Prec./form} & \textbf{First} & \textbf{Reg. page} & \textbf{Principal uses / notes} \\
\midrule
\endfirsthead
\toprule
\textbf{Glyph} & \textbf{Macro} & \textbf{Role} & \textbf{Scope} & \textbf{Prec./form} & \textbf{First} & \textbf{Reg. page} & \textbf{Principal uses / notes} \\
\midrule
\endhead
\midrule\multicolumn{8}{r}{\emph{continued}}\\\endfoot
\bottomrule\endlastfoot
\phantomsection\label{entry:symbol-registry-00709}\(T_\tau(X)\) & \texttt{\detokenize{T_\tau}} & term algebra on a generator set & P10 algebraic-and-categorical-logic & meta & \hyperref[sec:p10-algebraic-and-categorical-logic]{\ref*{sec:p10-algebraic-and-categorical-logic}} & \thepage & Finite terms over the finitary signature tau and variables X; a nonempty algebra only when variables or nullary symbols exist. \\
\phantomsection\label{entry:symbol-registry-00710}\(t^A(v)\) & \texttt{\detokenize{t^A(v)}} & evaluation of an algebraic term & P10 algebraic-and-categorical-logic & meta & \hyperref[sec:p10-algebraic-and-categorical-logic]{\ref*{sec:p10-algebraic-and-categorical-logic}} & \thepage & Term value in A under the variable assignment v; every assignment extends uniquely to a term-algebra homomorphism. \\
\phantomsection\label{entry:symbol-registry-00711}\(t\sigma,\widehat\sigma\) & \texttt{\detokenize{\widehat\sigma}} & simultaneous algebraic term substitution & P10 algebraic-and-categorical-logic & meta & \hyperref[sec:p10-algebraic-and-categorical-logic]{\ref*{sec:p10-algebraic-and-categorical-logic}} & \thepage & The map on variables extends homomorphically to terms; the extension may have a different generator alphabet as codomain. \\
\phantomsection\label{entry:symbol-registry-00712}\(\Gamma\models_{\mathrm{alg}}s\approx t\) & \texttt{\detokenize{\models_{\mathrm{alg}}}} & universal identity consequence & P10 algebraic-and-categorical-logic & meta & \hyperref[sec:p10-algebraic-and-categorical-logic]{\ref*{sec:p10-algebraic-and-categorical-logic}} & \thepage & Every nonempty algebra validating all premise identities under all valuations validates the conclusion identity under all valuations. \\
\phantomsection\label{entry:symbol-registry-00713}\(\Gamma\vdash_{\mathrm{eq}}s\approx t\) & \texttt{\detokenize{\vdash_{\mathrm{eq}}}} & finite equational derivability & P10 algebraic-and-categorical-logic & meta & \hyperref[sec:p10-algebraic-and-categorical-logic]{\ref*{sec:p10-algebraic-and-categorical-logic}} & \thepage & Derivability from identity premises by the complete registered reflexivity symmetry transitivity compatibility and substitution rules. \\
\phantomsection\label{entry:symbol-registry-00714}\(\theta_\Gamma\) & \texttt{\detokenize{\theta_\Gamma}} & congruence of provable identities & P10 algebraic-and-categorical-logic & meta & \hyperref[sec:p10-algebraic-and-categorical-logic]{\ref*{sec:p10-algebraic-and-categorical-logic}} & \thepage & The least fully invariant congruence on the nonempty term algebra containing all premise equation pairs. \\
\phantomsection\label{entry:symbol-registry-00715}\(F_{\mathcal V}(X),\eta_X\) & \texttt{\detokenize{F_{\mathcal V},}\allowbreak\detokenize{\eta_X}} & free algebra and its generator map & P10 algebraic-and-categorical-logic & meta & \hyperref[sec:p10-algebraic-and-categorical-logic]{\ref*{sec:p10-algebraic-and-categorical-logic}} & \thepage & The map from X has a unique homomorphic extension into every member of the variety; it need not be injective in a trivial variety. \\
\phantomsection\label{entry:symbol-registry-00716}\(\theta_{\mathcal V,X}\) & \texttt{\detokenize{\theta_{\mathcal V,X}}} & semantic congruence defining a relative free algebra & P10 algebraic-and-categorical-logic & meta & \hyperref[sec:p10-algebraic-and-categorical-logic]{\ref*{sec:p10-algebraic-and-categorical-logic}} & \thepage & Two X-terms are related exactly when all evaluations into all algebras of the variety identify them. \\
\phantomsection\label{entry:symbol-registry-00717}\(\operatorname{Con}A,\Delta_A,\nabla_A\) & \texttt{\detokenize{\operatorname{Con},}\allowbreak\detokenize{\Delta,}\allowbreak\detokenize{\nabla}} & congruence lattice and its bounds & P10 algebraic-and-categorical-logic & meta & \hyperref[sec:p10-algebraic-and-categorical-logic]{\ref*{sec:p10-algebraic-and-categorical-logic}} & \thepage & Congruences ordered by inclusion; Delta is equality and nabla is the universal relation; meets are intersections. \\
\phantomsection\label{entry:symbol-registry-00718}\(\operatorname{Cg}_A(a,b)\) & \texttt{\detokenize{\operatorname{Cg}}} & principal congruence generated by a pair & P10 algebraic-and-categorical-logic & meta & \hyperref[sec:p10-algebraic-and-categorical-logic]{\ref*{sec:p10-algebraic-and-categorical-logic}} & \thepage & Intersection of every congruence on A containing the pair (a,b). \\
\phantomsection\label{entry:symbol-registry-00719}\(\ker h\) & \texttt{\detokenize{\ker}} & kernel congruence of an algebra homomorphism & P10 algebraic-and-categorical-logic & meta & \hyperref[sec:p10-algebraic-and-categorical-logic]{\ref*{sec:p10-algebraic-and-categorical-logic}} & \thepage & Pairs of domain elements with equal h-images; the kernel of a product map is the intersection of its coordinate kernels. \\
\phantomsection\label{entry:symbol-registry-00720}\(\phi/\theta\) & \texttt{\detokenize{/}} & quotient congruence above a kernel & P10 algebraic-and-categorical-logic & meta & \hyperref[sec:p10-algebraic-and-categorical-logic]{\ref*{sec:p10-algebraic-and-categorical-logic}} & \thepage & For theta contained in phi this relation on A/theta relates exactly the classes of phi-related representatives. \\
\phantomsection\label{entry:symbol-registry-00721}\(\mu\) & \texttt{\detokenize{\mu}} & monolith of a nontrivial algebra & P10 subdirect monolith & meta & \hyperref[def:p10-subdirect-irreducible]{\ref*{def:p10-subdirect-irreducible}} & \thepage & Least congruence strictly above equality when it exists; this condition characterizes nontrivial subdirect irreducibility. \\
\phantomsection\label{entry:symbol-registry-00722}\(\theta_{a,b}\) & \texttt{\detokenize{\theta_{a,b}}} & maximal congruence separating a chosen pair & P10 algebraic-and-categorical-logic & meta & \hyperref[sec:p10-algebraic-and-categorical-logic]{\ref*{sec:p10-algebraic-and-categorical-logic}} & \thepage & Maximal among congruences omitting (a,b); the quotient is subdirectly irreducible and still separates the chosen elements. \\
\phantomsection\label{entry:symbol-registry-00723}\(F[a]\) & \texttt{\detokenize{F[a]}} & Boolean filter generated by adjoining one element & P10 legacy unit 2 stone-duality & meta & \hyperref[sec:p10-stone-duality]{\ref*{sec:p10-stone-duality}} & \thepage & Elements above f meet a for some f in F; used in ultrafilter dichotomy and avoiding-element extension. \\
\phantomsection\label{entry:symbol-registry-00724}\(X_B,\widehat a\) & \texttt{\detokenize{X_B,}\allowbreak\detokenize{\widehat}} & Stone space and its basic clopen sets & P10 legacy unit 2 stone-duality & meta & \hyperref[sec:p10-stone-duality]{\ref*{sec:p10-stone-duality}} & \thepage & X-B consists of the ultrafilters of B; hat-a consists of those containing a; trivial B gives the empty space. \\
\phantomsection\label{entry:symbol-registry-00725}\(\operatorname{Clop}(X)\) & \texttt{\detokenize{\operatorname{Clop}}} & Boolean algebra of clopen subsets & P10 legacy unit 2 stone-duality & meta & \hyperref[sec:p10-stone-duality]{\ref*{sec:p10-stone-duality}} & \thepage & Intersection union and relative complement act on the clopens of X; the bounds are the empty set and X. \\
\phantomsection\label{entry:symbol-registry-00726}\(\beta_B\) & \texttt{\detokenize{\beta_B}} & Boolean representation isomorphism & P10 legacy unit 2 stone-duality & meta & \hyperref[sec:p10-stone-duality]{\ref*{sec:p10-stone-duality}} & \thepage & Maps a Boolean element a to its basic clopen hat-a in the Stone space X-B. \\
\phantomsection\label{entry:symbol-registry-00727}\(\varepsilon_X\) & \texttt{\detokenize{\varepsilon_X}} & Stone-space reconstruction homeomorphism & P10 legacy unit 2 stone-duality & meta & \hyperref[sec:p10-stone-duality]{\ref*{sec:p10-stone-duality}} & \thepage & Maps x to the ultrafilter of all clopen subsets containing x. \\
\phantomsection\label{entry:symbol-registry-00728}\(X(h)\) & \texttt{\detokenize{X(h)}} & continuous dual of a Boolean homomorphism & P10 legacy unit 2 stone-duality & meta & \hyperref[sec:p10-stone-duality]{\ref*{sec:p10-stone-duality}} & \thepage & For h from B to C sends an ultrafilter of C to its inverse image under h, an ultrafilter of B. \\
\phantomsection\label{entry:symbol-registry-00729}\(\operatorname{Clop}(g)\) & \texttt{\detokenize{\operatorname{Clop}(g)}} & Boolean dual of a continuous map & P10 legacy unit 2 stone-duality & meta & \hyperref[sec:p10-stone-duality]{\ref*{sec:p10-stone-duality}} & \thepage & For g from X to Y sends a clopen subset of Y to its inverse image in X. \\
\phantomsection\label{entry:symbol-registry-00730}\(\mathsf{BA},\mathsf{Stone}\) & \texttt{\detokenize{\mathsf{BA},}\allowbreak\detokenize{\mathsf{Stone}}} & categories in Stone duality & P10 legacy unit 2 stone-duality & meta & \hyperref[sec:p10-stone-duality]{\ref*{sec:p10-stone-duality}} & \thepage & Boolean algebras with homomorphisms preserving both bounds; compact Hausdorff spaces with clopen bases and continuous maps; trivial and empty objects are allowed. \\
\phantomsection\label{entry:symbol-registry-00731}\(K=2^{\mathbb N}\) & \texttt{\detokenize{K,}\allowbreak\detokenize{2^{\mathbb N}}} & Cantor space of binary sequences & P10 legacy unit 2 stone-duality & meta & \hyperref[sec:p10-stone-duality]{\ref*{sec:p10-stone-duality}} & \thepage & The countable product of the discrete two-element space with coordinates starting at zero. \\
\phantomsection\label{entry:symbol-registry-00732}\([s],x|_n\) & \texttt{\detokenize{[s],}\allowbreak\detokenize{|_n}} & binary cylinder and length-n initial segment & P10 legacy unit 2 stone-duality & meta & \hyperref[sec:p10-stone-duality]{\ref*{sec:p10-stone-duality}} & \thepage & The cylinder contains precisely the sequences extending the finite word s; the empty word gives all Cantor space. \\
\phantomsection\label{entry:symbol-registry-00733}\(\mathcal B_n,\mathcal B_{\rm cyl}\) & \texttt{\detokenize{\mathcal B_n,}\allowbreak\detokenize{\mathcal B_{\rm cyl}}} & finite cylinder algebras and their union & P10 legacy unit 2 stone-duality & meta & \hyperref[sec:p10-stone-duality]{\ref*{sec:p10-stone-duality}} & \thepage & B-n contains all unions of length-n cylinders; the union is the clopen algebra of Cantor space and the direct limit of the inclusions. \\
\phantomsection\label{entry:symbol-registry-00734}\(H:K\to X_B\) & \texttt{\detokenize{H}} & explicit Cantor-to-Stone homeomorphism & P10 legacy unit 2 stone-duality & meta & \hyperref[sec:p10-stone-duality]{\ref*{sec:p10-stone-duality}} & \thepage & Using a Boolean isomorphism alpha from B to the cylinder algebra sends x to the elements b for which x belongs to alpha(b). \\
\phantomsection\label{entry:symbol-registry-00735}\(U_p\) & \texttt{\detokenize{U_p}} & principal ultrafilter of a Boolean atom & P10 legacy unit 2 stone-duality & meta & \hyperref[sec:p10-stone-duality]{\ref*{sec:p10-stone-duality}} & \thepage & All elements above atom p; its singleton is the basic clopen hat-p and hence an isolated Stone point. \\
\phantomsection\label{entry:symbol-registry-00736}\(\Lambda,n_f\) & \texttt{\detokenize{\Lambda,}\allowbreak\detokenize{n_f}} & finitary Boolean-operator signature and arity & P10 legacy unit 3 nonclassical-algebras & meta & \hyperref[sec:p10-nonclassical-algebras]{\ref*{sec:p10-nonclassical-algebras}} & \thepage & A set of additional symbols with one finite arity per symbol; nullary symbols designate arbitrary Boolean elements. \\
\phantomsection\label{entry:symbol-registry-00737}\(f^B\) & \texttt{\detokenize{f^B}} & abstract Boolean operator & P10 legacy unit 3 nonclassical-algebras & meta & \hyperref[sec:p10-nonclassical-algebras]{\ref*{sec:p10-nonclassical-algebras}} & \thepage & For positive arity preserves zero and finite joins in each coordinate; the conditions are vacuous for nullary symbols. \\
\phantomsection\label{entry:symbol-registry-00738}\(\mathfrak F^+\) & \texttt{\detokenize{\mathfrak F^+}} & full complex algebra of a relational frame & P10 legacy unit 3 nonclassical-algebras & meta & \hyperref[sec:p10-nonclassical-algebras]{\ref*{sec:p10-nonclassical-algebras}} & \thepage & Power-set Boolean algebra with existential successor-tuple operations; the carrier may be the power set of the empty frame. \\
\phantomsection\label{entry:symbol-registry-00739}\(R_f(U,V_1,\ldots,V_n)\) & \texttt{\detokenize{R_f}} & canonical relation of a finitary Boolean operator & P10 legacy unit 3 nonclassical-algebras & meta & \hyperref[sec:p10-nonclassical-algebras]{\ref*{sec:p10-nonclassical-algebras}} & \thepage & Every tuple of elements from the successor ultrafilters has its operator value in U; for arity zero this is a unary predicate. \\
\phantomsection\label{entry:symbol-registry-00740}\(\mathbf B^\bullet\) & \texttt{\detokenize{\mathbf B^\bullet}} & ultrafilter frame of a Boolean algebra with operators & P10 legacy unit 3 nonclassical-algebras & meta & \hyperref[sec:p10-nonclassical-algebras]{\ref*{sec:p10-nonclassical-algebras}} & \thepage & The set of Boolean ultrafilters with all canonical operator relations. \\
\phantomsection\label{entry:symbol-registry-00741}\(\mathbf B^\sigma\) & \texttt{\detokenize{\mathbf B^\sigma}} & canonical extension of a Boolean algebra with operators & P10 legacy unit 3 nonclassical-algebras & meta & \hyperref[sec:p10-nonclassical-algebras]{\ref*{sec:p10-nonclassical-algebras}} & \thepage & The full complex algebra of the ultrafilter frame; the Boolean embedding is dense and compact. \\
\phantomsection\label{entry:symbol-registry-00742}\(\mathbf B_*\) & \texttt{\detokenize{\mathbf B_*}} & descriptive general frame of a Boolean algebra with operators & P10 legacy unit 3 nonclassical-algebras & meta & \hyperref[sec:p10-nonclassical-algebras]{\ref*{sec:p10-nonclassical-algebras}} & \thepage & The canonical relations on the Stone space together with precisely the basic-clopen image as admissible algebra. \\
\phantomsection\label{entry:symbol-registry-00743}\(\mathcal G^+,\mathcal A\) & \texttt{\detokenize{\mathcal G^+,}\allowbreak\detokenize{\mathcal A}} & admissible algebra of a general frame & P10 legacy unit 3 nonclassical-algebras & meta & \hyperref[sec:p10-nonclassical-algebras]{\ref*{sec:p10-nonclassical-algebras}} & \thepage & The carrier is the Boolean family of admissible sets closed under complex operators. \\
\phantomsection\label{entry:symbol-registry-00744}\(h_*\) & \texttt{\detokenize{h_*}} & bounded morphism dual to a Boolean-operator homomorphism & P10 legacy unit 3 nonclassical-algebras & meta & \hyperref[sec:p10-nonclassical-algebras]{\ref*{sec:p10-nonclassical-algebras}} & \thepage & Inverse image on ultrafilters; its back condition follows from the simultaneous filter-extension lemma. \\
\phantomsection\label{entry:symbol-registry-00745}\(p_U=\bigwedge e[U]\) & \texttt{\detokenize{p_U}} & atom in a dense compact Boolean completion & P10 legacy unit 3 nonclassical-algebras & meta & \hyperref[sec:p10-nonclassical-algebras]{\ref*{sec:p10-nonclassical-algebras}} & \thepage & The meet of the images of all elements of an ultrafilter U; compactness makes it nonzero and density makes it an atom. \\
\phantomsection\label{entry:symbol-registry-00746}\(\Box_A,\Diamond_A\) & \texttt{\detokenize{\Box_A,}\allowbreak\detokenize{\Diamond_A}} & dual operations of a normal modal algebra & P10 legacy unit 3 nonclassical-algebras & meta & \hyperref[sec:p10-nonclassical-algebras]{\ref*{sec:p10-nonclassical-algebras}} & \thepage & Box preserves finite meets and one; diamond is its Boolean dual and preserves finite joins and zero. \\
\phantomsection\label{entry:symbol-registry-00747}\(\varphi\equiv_L\psi,[\varphi]_L\) & \texttt{\detokenize{\equiv_L,}\allowbreak\detokenize{[\cdot]_L}} & normal-modal provable equivalence and its classes & P10 legacy unit 3 nonclassical-algebras & meta & \hyperref[sec:p10-nonclassical-algebras]{\ref*{sec:p10-nonclassical-algebras}} & \thepage & Equivalence means theoremhood of the biconditional in L and is a congruence for all formula constructors. \\
\phantomsection\label{entry:symbol-registry-00748}\(F_L\) & \texttt{\detokenize{F_L}} & normal modal Lindenbaum algebra & P10 legacy unit 3 nonclassical-algebras & meta & \hyperref[sec:p10-nonclassical-algebras]{\ref*{sec:p10-nonclassical-algebras}} & \thepage & The formula algebra modulo provable equivalence; arbitrary valuations are handled with finite representative selections per formula. \\
\phantomsection\label{entry:symbol-registry-00749}\(\operatorname{Alg}(L)\) & \texttt{\detokenize{\operatorname{Alg}}} & algebra class validating a normal modal logic & P10 legacy unit 3 nonclassical-algebras & meta & \hyperref[sec:p10-nonclassical-algebras]{\ref*{sec:p10-nonclassical-algebras}} & \thepage & Normal modal algebras validating every member of L under all algebra valuations. \\
\phantomsection\label{entry:symbol-registry-00750}\(\mathfrak G_L,\mathcal A_L\) & \texttt{\detokenize{\mathfrak G_L,}\allowbreak\detokenize{\mathcal A_L}} & canonical general frame and its admissible sets & P10 legacy unit 3 nonclassical-algebras & meta & \hyperref[sec:p10-nonclassical-algebras]{\ref*{sec:p10-nonclassical-algebras}} & \thepage & Ultrafilters of the modal Lindenbaum algebra with the canonical relation and Stone clopens as admissible variable values. \\
\phantomsection\label{entry:symbol-registry-00751}\(\mathfrak F_L,X_L,R_L\) & \texttt{\detokenize{\mathfrak F_L,}\allowbreak\detokenize{X_L,}\allowbreak\detokenize{R_L}} & underlying canonical relational frame & P10 legacy unit 3 nonclassical-algebras & meta & \hyperref[sec:p10-nonclassical-algebras]{\ref*{sec:p10-nonclassical-algebras}} & \thepage & Same worlds and relation as the canonical general frame; ordinary valuations may use every subset of X-L. \\
\phantomsection\label{entry:symbol-registry-00752}\(\operatorname{Gen}(L),\operatorname{Fr}(L)\) & \texttt{\detokenize{\operatorname{Gen},}\allowbreak\detokenize{\operatorname{Fr}}} & nonempty validating general-frame and ordinary-frame classes & P10 legacy unit 3 nonclassical-algebras & meta & \hyperref[sec:p10-nonclassical-algebras]{\ref*{sec:p10-nonclassical-algebras}} & \thepage & General frames validate L using admissible valuations; ordinary frames validate it using all subset valuations. \\
\phantomsection\label{entry:symbol-registry-00753}\(L(Y)\) & \texttt{\detokenize{L(Y)}} & schematic extension of a modal logic to another variable set & P10 legacy unit 3 nonclassical-algebras & meta & \hyperref[sec:p10-nonclassical-algebras]{\ref*{sec:p10-nonclassical-algebras}} & \thepage & The least normal logic over Y containing all substituted instances of the original logic; used to state all-variable-set canonicity. \\
\phantomsection\label{entry:symbol-registry-00754}\(\operatorname{Fm}_\Lambda(X)\) & \texttt{\detokenize{\operatorname{Fm}_\Lambda}} & formulas for a finitary modal signature & P10 legacy unit 3 nonclassical-algebras & meta & \hyperref[sec:p10-nonclassical-algebras]{\ref*{sec:p10-nonclassical-algebras}} & \thepage & Finite Boolean formulas over X with one constructor of arity n-f for every symbol f in Lambda. \\
\phantomsection\label{entry:symbol-registry-00755}\(f_i(\chi)\) & \texttt{\detokenize{f_i}} & one-coordinate modal formula context & P10 legacy unit 3 nonclassical-algebras & meta & \hyperref[sec:p10-nonclassical-algebras]{\ref*{sec:p10-nonclassical-algebras}} & \thepage & All other arguments are fixed; used in theorem-level monotonicity and the normality and additivity axiom schemes. \\
\phantomsection\label{entry:symbol-registry-00756}\(F_{\Lambda,L}\) & \texttt{\detokenize{F_{\Lambda,L}}} & Lindenbaum algebra of a finitary normal logic & P10 legacy unit 3 nonclassical-algebras & meta & \hyperref[sec:p10-nonclassical-algebras]{\ref*{sec:p10-nonclassical-algebras}} & \thepage & The formula quotient by L-provable equivalence with all finite-arity and nullary operators interpreted on classes. \\
\phantomsection\label{entry:symbol-registry-00757}\(\operatorname{Alg}_\Lambda(L)\) & \texttt{\detokenize{\operatorname{Alg}_\Lambda}} & Boolean-operator algebras validating a finitary logic & P10 legacy unit 3 nonclassical-algebras & meta & \hyperref[sec:p10-nonclassical-algebras]{\ref*{sec:p10-nonclassical-algebras}} & \thepage & All Lambda-BAOs validating L under every algebra valuation, including trivial algebras. \\
\phantomsection\label{entry:symbol-registry-00758}\(\operatorname{Ult}(A)\) & \texttt{\detokenize{\operatorname{Ult}}} & ultrafilters of the Boolean reduct & P10 legacy unit 3 nonclassical-algebras & meta & \hyperref[sec:p10-nonclassical-algebras]{\ref*{sec:p10-nonclassical-algebras}} & \thepage & The set of proper maximal Boolean filters, used as the canonical frame carrier of A. \\
\phantomsection\label{entry:symbol-registry-00759}\(\mathcal D_L,W_L\) & \texttt{\detokenize{\mathcal D_L,}\allowbreak\detokenize{W_L}} & descriptive canonical frame of a finitary normal logic & P10 legacy unit 3 nonclassical-algebras & meta & \hyperref[sec:p10-nonclassical-algebras]{\ref*{sec:p10-nonclassical-algebras}} & \thepage & Worlds are ultrafilters of F-Lambda-L; the relations are the full canonical operator relations and the admissible sets are Stone clopens. \\
\phantomsection\label{entry:symbol-registry-00760}\(\operatorname{DGen}_\Lambda(L)\) & \texttt{\detokenize{\operatorname{DGen}_\Lambda}} & descriptive general frames validating a finitary logic & P10 legacy unit 3 nonclassical-algebras & meta & \hyperref[sec:p10-nonclassical-algebras]{\ref*{sec:p10-nonclassical-algebras}} & \thepage & The nonempty descriptive Lambda-general frames validating every theorem under all admissible valuations. \\
\phantomsection\label{entry:symbol-registry-00761}\(v_0(p)=\widehat{[p]_L}\) & \texttt{\detokenize{v_0}} & canonical modal valuation & P10 legacy unit 3 nonclassical-algebras & meta & \hyperref[sec:p10-nonclassical-algebras]{\ref*{sec:p10-nonclassical-algebras}} & \thepage & Assigns each proposition variable the basic clopen determined by its formula class. \\
\phantomsection\label{entry:symbol-registry-00762}\(X_D,\widehat a\) & \texttt{\detokenize{X_D,}\allowbreak\detokenize{\widehat}} & Priestley spectrum and its basic clopen upsets & P10 legacy unit 3 nonclassical-algebras & meta & \hyperref[sec:p10-nonclassical-algebras]{\ref*{sec:p10-nonclassical-algebras}} & \thepage & Proper prime lattice filters ordered by inclusion; hat-a consists of those containing a, and the topology also uses their complements. \\
\phantomsection\label{entry:symbol-registry-00763}\(\operatorname{ClopUp}(X)\) & \texttt{\detokenize{\operatorname{ClopUp}}} & lattice of clopen upsets & P10 legacy unit 3 nonclassical-algebras & meta & \hyperref[sec:p10-nonclassical-algebras]{\ref*{sec:p10-nonclassical-algebras}} & \thepage & Bounded distributive lattice for a Priestley space; carries the residual complement-of-downclosure operation for an Esakia space. \\
\phantomsection\label{entry:symbol-registry-00764}\(\eta_D\) & \texttt{\detokenize{\eta_D}} & Priestley lattice representation & P10 legacy unit 3 nonclassical-algebras & meta & \hyperref[sec:p10-nonclassical-algebras]{\ref*{sec:p10-nonclassical-algebras}} & \thepage & Maps each lattice element to the clopen upset of prime filters containing it. \\
\phantomsection\label{entry:symbol-registry-00765}\(\varepsilon_X:X\to X_{\operatorname{ClopUp}(X)}\) & \texttt{\detokenize{\varepsilon_X}} & ordered-space reconstruction by prime filters & P10 legacy unit 3 nonclassical-algebras & meta & \hyperref[sec:p10-nonclassical-algebras]{\ref*{sec:p10-nonclassical-algebras}} & \thepage & Maps a point to its membership filter of clopen upsets; an order isomorphism and homeomorphism. \\
\phantomsection\label{entry:symbol-registry-00766}\(\Pi_D,p_a\) & \texttt{\detokenize{\Pi_D,}\allowbreak\detokenize{p_a}} & classical encoding of prime lattice filters & P10 legacy unit 3 nonclassical-algebras & meta & \hyperref[sec:p10-nonclassical-algebras]{\ref*{sec:p10-nonclassical-algebras}} & \thepage & The proposition p-a records membership of a; Pi-D imposes only bounded-lattice homomorphism laws to two, not Heyting implication preservation. \\
\phantomsection\label{entry:symbol-registry-00767}\(\mathsf{DC}\) & \texttt{\detokenize{\mathsf{DC}}} & dependent choice for a prescribed start & P10 legacy unit 3 nonclassical-algebras & meta & \hyperref[sec:p10-nonclassical-algebras]{\ref*{sec:p10-nonclassical-algebras}} & \thepage & A serial relation on a nonempty set admits an infinite path beginning at a prescribed point; used in the arbitrary-poset dense-set lemma. \\
\phantomsection\label{entry:symbol-registry-00768}\((p_n),\mathcal D\) & \texttt{\detokenize{p_n,}\allowbreak\detokenize{\mathcal D}} & descending conditions and dense requirements & P10 legacy unit 3 nonclassical-algebras & meta & \hyperref[sec:p10-nonclassical-algebras]{\ref*{sec:p10-nonclassical-algebras}} & \thepage & Rasiowa-Sikorski construction has p-n-plus-one stronger than p-n and in D-n; its upward closure is generic for the given family. \\
\phantomsection\label{entry:symbol-registry-00769}\(a\Rightarrow b\) & \texttt{\detokenize{\Rightarrow}} & Heyting residual on algebra elements & P10 legacy unit 3 nonclassical-algebras & meta & \hyperref[sec:p10-nonclassical-algebras]{\ref*{sec:p10-nonclassical-algebras}} & \thepage & The greatest c with c meet a below b; source notation a-arrow-b denotes this same operation when operands are algebra elements. \\
\phantomsection\label{entry:symbol-registry-00770}\(\mathcal O(X)\) & \texttt{\detokenize{\mathcal O}} & Heyting algebra of open subsets & P10 legacy unit 3 nonclassical-algebras & meta & \hyperref[sec:p10-nonclassical-algebras]{\ref*{sec:p10-nonclassical-algebras}} & \thepage & Finite meets are intersections and implication is the interior of the union of the first operand's complement with the second. \\
\phantomsection\label{entry:symbol-registry-00771}\(\operatorname{Up}(W)\) & \texttt{\detokenize{\operatorname{Up}}} & Heyting algebra of upsets of a preorder & P10 legacy unit 3 nonclassical-algebras & meta & \hyperref[sec:p10-nonclassical-algebras]{\ref*{sec:p10-nonclassical-algebras}} & \thepage & Implication holds at w when every larger point in the antecedent upset belongs to the consequent upset. \\
\phantomsection\label{entry:symbol-registry-00772}\(H_{\mathsf{IPC}},[\varphi]\) & \texttt{\detokenize{H_{\mathsf{IPC}},}\allowbreak\detokenize{[\cdot]}} & IPC Lindenbaum Heyting algebra & P10 legacy unit 3 nonclassical-algebras & meta & \hyperref[sec:p10-nonclassical-algebras]{\ref*{sec:p10-nonclassical-algebras}} & \thepage & Formula classes under mutual IPL derivability, ordered by one-premise derivability. \\
\phantomsection\label{entry:symbol-registry-00773}\(\operatorname{Pf}(H)\) & \texttt{\detokenize{\operatorname{Pf}}} & prime-filter order of a Heyting algebra & P10 legacy unit 3 nonclassical-algebras & meta & \hyperref[sec:p10-nonclassical-algebras]{\ref*{sec:p10-nonclassical-algebras}} & \thepage & All proper prime lattice filters ordered by inclusion; filters need not be maximal. \\
\phantomsection\label{entry:symbol-registry-00774}\(\beta_H(a)\) & \texttt{\detokenize{\beta_H}} & Heyting embedding into upsets of prime filters & P10 legacy unit 3 nonclassical-algebras & meta & \hyperref[sec:p10-nonclassical-algebras]{\ref*{sec:p10-nonclassical-algebras}} & \thepage & The set of prime filters containing a; the implication witness proves preservation of the residual. \\
\phantomsection\label{entry:symbol-registry-00775}\(\downarrow S,\uparrow x\) & \texttt{\detokenize{\downarrow,}\allowbreak\detokenize{\uparrow}} & downward closure and principal upset & P10 legacy unit 3 nonclassical-algebras & meta & \hyperref[sec:p10-nonclassical-algebras]{\ref*{sec:p10-nonclassical-algebras}} & \thepage & Order closures in a Priestley or Esakia space; the downclosure condition applies to every clopen subset. \\
\phantomsection\label{entry:symbol-registry-00776}\(g[\uparrow x]=\uparrow g(x)\) & \texttt{\detokenize{g[\uparrow x]}} & Esakia order back condition & P10 legacy unit 3 nonclassical-algebras & meta & \hyperref[sec:p10-nonclassical-algebras]{\ref*{sec:p10-nonclassical-algebras}} & \thepage & Together with continuity and monotonicity this characterizes Esakia morphisms and ensures implication-preserving inverse images. \\
\phantomsection\label{entry:symbol-registry-00777}\(\mathsf{HA},\mathsf{BA}\) & \texttt{\detokenize{\mathsf{HA},}\allowbreak\detokenize{\mathsf{BA}}} & Heyting and Boolean classes in the common residual signature & P10 legacy unit 3 nonclassical-algebras & meta & \hyperref[sec:p10-nonclassical-algebras]{\ref*{sec:p10-nonclassical-algebras}} & \thepage & Classes of algebras for the intermediate-logic correspondence; HA here is not Heyting arithmetic from Part VII. \\
\phantomsection\label{entry:symbol-registry-00778}\(H_L\) & \texttt{\detokenize{H_L}} & Lindenbaum Heyting algebra of an intermediate logic & P10 legacy unit 3 nonclassical-algebras & meta & \hyperref[sec:p10-nonclassical-algebras]{\ref*{sec:p10-nonclassical-algebras}} & \thepage & The quotient by L-provable biconditional, with formula implication inducing the algebraic residual. \\
\phantomsection\label{entry:symbol-registry-00779}\(\mathsf V(L),\operatorname{Log}(\mathcal V)\) & \texttt{\detokenize{\mathsf V,}\allowbreak\detokenize{\operatorname{Log}}} & order-reversing logic-variety assignments & P10 legacy unit 3 nonclassical-algebras & meta & \hyperref[sec:p10-nonclassical-algebras]{\ref*{sec:p10-nonclassical-algebras}} & \thepage & Heyting algebras validating L and formulas valid in V; inverse on the intervals IPC-to-CPL and BA-to-HA. \\
\phantomsection\label{entry:symbol-registry-00780}\(s^\flat,\varphi^{\rm alg}\) & \texttt{\detokenize{{}^\flat,}\allowbreak\detokenize{{}^{\rm alg}}} & translation between residual terms and formulas & P10 legacy unit 3 nonclassical-algebras & meta & \hyperref[sec:p10-nonclassical-algebras]{\ref*{sec:p10-nonclassical-algebras}} & \thepage & Converts algebraic zero one and residual into bottom top and implication or conversely; identities correspond to biconditionals with value one. \\
\phantomsection\label{entry:symbol-registry-00781}\(\mathcal C(A,B)\) & \texttt{\detokenize{\mathcal C(A,}\allowbreak\detokenize{B)}} & categorical hom-set & P10 legacy units 7–9 lawvere-theories to doctrine-quantifiers & meta & \hyperref[sec:p10-lawvere-theories]{\ref*{sec:p10-lawvere-theories}} & \thepage & Arrows from object A to object B; local smallness places each such set in the fixed universe. \\
\phantomsection\label{entry:symbol-registry-00782}\(\mathcal C^{\mathrm{op}}\) & \texttt{\detokenize{{}^{\mathrm{op}}}} & opposite category & P10 legacy units 7–9 lawvere-theories to doctrine-quantifiers & meta & \hyperref[sec:p10-lawvere-theories]{\ref*{sec:p10-lawvere-theories}} & \thepage & Reverses the arrows and composition order; a contravariant functor is an ordinary functor out of this category. \\
\phantomsection\label{entry:symbol-registry-00783}\(1_A,1_{\mathcal C}\) & \texttt{\detokenize{1_A,}\allowbreak\detokenize{1_{\mathcal C}}} & identity arrow and identity functor & P10 legacy units 7–9 lawvere-theories to doctrine-quantifiers & meta & \hyperref[sec:p10-lawvere-theories]{\ref*{sec:p10-lawvere-theories}} & \thepage & The subscript states the categorical level: an object for an arrow and a category for its identity functor. \\
\phantomsection\label{entry:symbol-registry-00784}\(\operatorname{Nat}(F,G)\) & \texttt{\detokenize{\operatorname{Nat}}} & family of natural transformations & P10 legacy units 7–9 lawvere-theories to doctrine-quantifiers & meta & \hyperref[sec:p10-lawvere-theories]{\ref*{sec:p10-lawvere-theories}} & \thepage & Component families in the containing ambient universe; Yoneda gives a small set of parameters for transformations from a representable. \\
\phantomsection\label{entry:symbol-registry-00785}\(F\dashv G,\Phi_{A,B}\) & \texttt{\detokenize{\dashv,}\allowbreak\detokenize{\Phi_{A,B}}} & adjunction and hom-set transposition & P10 legacy units 7–9 lawvere-theories to doctrine-quantifiers & meta & \hyperref[sec:p10-lawvere-theories]{\ref*{sec:p10-lawvere-theories}} & \thepage & Phi sends an arrow from F(A) to B to an arrow from A to G(B); its bijections are natural in both arguments. \\
\phantomsection\label{entry:symbol-registry-00786}\(\eta,\varepsilon\) & \texttt{\detokenize{\eta,}\allowbreak\detokenize{\varepsilon}} & unit and counit of an adjunction & P10 legacy units 7–9 lawvere-theories to doctrine-quantifiers & meta & \hyperref[sec:p10-lawvere-theories]{\ref*{sec:p10-lawvere-theories}} & \thepage & Natural transformations from the identity to GF and from FG to the identity; their components satisfy the two triangle identities. \\
\phantomsection\label{entry:symbol-registry-00787}\(\pi_i,\iota_i,\langle f_i\rangle_i\) & \texttt{\detokenize{\pi_i,}\allowbreak\detokenize{\iota_i,}\allowbreak\detokenize{\langle\rangle}} & product projections coproduct injections and pairing & P10 legacy units 7–9 lawvere-theories to doctrine-quantifiers & meta & \hyperref[sec:p10-lawvere-theories]{\ref*{sec:p10-lawvere-theories}} & \thepage & Domains and codomains distinguish projections from injections; pairing is the unique arrow with its specified product components. \\
\phantomsection\label{entry:symbol-registry-00788}\(f\times g\) & \texttt{\detokenize{\times}} & product of categorical arrows & P10 legacy units 7–9 lawvere-theories to doctrine-quantifiers & meta & \hyperref[sec:p10-lawvere-theories]{\ref*{sec:p10-lawvere-theories}} & \thepage & The arrow with components f after the first projection and g after the second; respects identity and composition. \\
\phantomsection\label{entry:symbol-registry-00789}\(B^A,\operatorname{ev}_{A,B},\Lambda(f)\) & \texttt{\detokenize{{}^{A},}\allowbreak\detokenize{\operatorname{ev},}\allowbreak\detokenize{\Lambda}} & exponential object evaluation and transpose & P10 legacy units 7–9 lawvere-theories to doctrine-quantifiers & meta & \hyperref[sec:p10-lawvere-theories]{\ref*{sec:p10-lawvere-theories}} & \thepage & Evaluation has domain \(B^A\times A\); transposition sends \(f:X\times A\to B\) to the unique corresponding arrow \(X\to B^A\). \\
\phantomsection\label{entry:symbol-registry-00790}\(\mathcal S_{\beta\eta}\) & \texttt{\detokenize{\mathcal S_{\beta\eta}}} & category of arrow-only typed contexts & P10 legacy unit 8 cartesian-lambda & meta & \hyperref[sec:p10-cartesian-lambda]{\ref*{sec:p10-cartesian-lambda}} & \thepage & Objects are finite type lists and arrows are typed term tuples modulo componentwise beta-eta equality. \\
\phantomsection\label{entry:symbol-registry-00791}\(\Theta\vdash_{\beta\eta}M=N:\alpha\) & \texttt{\detokenize{\vdash_{\beta\eta}}} & typed beta-eta equational consequence & P10 legacy unit 8 cartesian-lambda & meta & \hyperref[sec:p10-cartesian-lambda]{\ref*{sec:p10-cartesian-lambda}} & \thepage & The equational metatheory includes alpha-renaming and typed beta and eta schemas; the registered reduction relation remains beta reduction. \\
\phantomsection\label{entry:symbol-registry-00792}\(\sigma:\Theta\to\Psi,M[\sigma]\) & \texttt{\detokenize{\sigma,}\allowbreak\detokenize{M[\sigma]}} & typed simultaneous substitution and its action & P10 legacy unit 8 cartesian-lambda & meta & \hyperref[sec:p10-cartesian-lambda]{\ref*{sec:p10-cartesian-lambda}} & \thepage & A tuple of Theta-terms with the types declared in Psi; substitutes into a Psi-term without resubstituting into inserted terms. \\
\phantomsection\label{entry:symbol-registry-00793}\(c_\Gamma(\beta),\Delta^\Gamma\) & \texttt{\detokenize{c_\Gamma,}\allowbreak\detokenize{\Delta^\Gamma}} & curried type and context exponential & P10 legacy unit 8 cartesian-lambda & meta & \hyperref[sec:p10-cartesian-lambda]{\ref*{sec:p10-cartesian-lambda}} & \thepage & Right-associated arrows over Gamma's type list; the exponential list has one such curried type per component of Delta. \\
\phantomsection\label{entry:symbol-registry-00794}\(\llbracket\Gamma\rrbracket,\llbracket M\rrbracket\) & \texttt{\detokenize{\llbracket\rrbracket}} & categorical interpretation of a context or term & P10 legacy unit 8 cartesian-lambda & meta & \hyperref[sec:p10-cartesian-lambda]{\ref*{sec:p10-cartesian-lambda}} & \thepage & Context products begin with the chosen terminal object; a term denotes an arrow from its context object to its interpreted type. \\
\phantomsection\label{entry:symbol-registry-00795}\(\underline A,A_T,P_\Gamma\) & \texttt{\detokenize{\underline,}\allowbreak\detokenize{A_T,}\allowbreak\detokenize{P_\Gamma}} & named type interpreted type and context object & P10 legacy unit 8 cartesian-lambda & meta & \hyperref[sec:p10-cartesian-lambda]{\ref*{sec:p10-cartesian-lambda}} & \thepage & The named base type denotes \(A\); \(A_T\) interprets the type \(T\), and \(P_\Gamma\) is the recursively chosen product with leading terminal factor. \\
\phantomsection\label{entry:symbol-registry-00796}\(c_f,u,\mathsf{pair}_{A,B}\) & \texttt{\detokenize{c_f,}\allowbreak\detokenize{u,}\allowbreak\detokenize{\mathsf{pair}}} & named arrow terminal and pairing constants & P10 legacy unit 8 cartesian-lambda & meta & \hyperref[sec:p10-cartesian-lambda]{\ref*{sec:p10-cartesian-lambda}} & \thepage & Closed constants in the expanded internal language; their declared applications denote the corresponding arrows and product pairing. \\
\phantomsection\label{entry:symbol-registry-00797}\(e_{A,B},l_{A,B}\) & \texttt{\detokenize{e_{A,B},}\allowbreak\detokenize{l_{A,B}}} & comparison constants for a named exponential & P10 legacy unit 8 cartesian-lambda & meta & \hyperref[sec:p10-cartesian-lambda]{\ref*{sec:p10-cartesian-lambda}} & \thepage & Compare the base type naming \(B^A\) with the arrow type between the named types for \(A\) and \(B\); both unary interpretations are identities. \\
\phantomsection\label{entry:symbol-registry-00798}\(\operatorname{Diag}(\mathcal C),=_{\mathcal C}\) & \texttt{\detokenize{\operatorname{Diag},}\allowbreak\detokenize{=_{\mathcal C}}} & full true diagram and its equational theory & P10 legacy unit 8 cartesian-lambda & meta & \hyperref[sec:p10-cartesian-lambda]{\ref*{sec:p10-cartesian-lambda}} & \thepage & Contains every true typed equation in every finite context of the named language; no effective enumeration is asserted. \\
\phantomsection\label{entry:symbol-registry-00799}\(\kappa_T,\iota_T\) & \texttt{\detokenize{\kappa_T,}\allowbreak\detokenize{\iota_T}} & recursive type comparison terms & P10 legacy unit 8 cartesian-lambda & meta & \hyperref[sec:p10-cartesian-lambda]{\ref*{sec:p10-cartesian-lambda}} & \thepage & Closed inverse comparisons between T and the named type for its interpreted object; applications preserve denotation. \\
\phantomsection\label{entry:symbol-registry-00800}\(\operatorname{pack}_\Gamma,\operatorname{unpack}_\Gamma\) & \texttt{\detokenize{\operatorname{pack},}\allowbreak\detokenize{\operatorname{unpack}}} & inverse context packing arrows & P10 legacy unit 8 cartesian-lambda & meta & \hyperref[sec:p10-cartesian-lambda]{\ref*{sec:p10-cartesian-lambda}} & \thepage & Compare \(\Gamma\) with the singleton named context for \(P_\Gamma\); the packing term and its context arrow have different codomains. \\
\phantomsection\label{entry:symbol-registry-00801}\(\mathcal S_{\mathcal C},E,J\) & \texttt{\detokenize{\mathcal S_{\mathcal C},}\allowbreak\detokenize{E,}\allowbreak\detokenize{J}} & internal context category evaluation and quasi-inverse & P10 legacy unit 8 cartesian-lambda & meta & \hyperref[sec:p10-cartesian-lambda]{\ref*{sec:p10-cartesian-lambda}} & \thepage & Evaluation sends \(\Gamma\) to \(P_\Gamma\); \(J\) names objects and arrows, with \(EJ(A)=1\times A\) for the chosen product. \\
\phantomsection\label{entry:symbol-registry-00802}\(j_{\Gamma,\Delta},\chi_{\Gamma,\Delta}\) & \texttt{\detokenize{j_{\Gamma,\Delta},}\allowbreak\detokenize{\chi_{\Gamma,\Delta}}} & product and exponential comparison isomorphisms & P10 legacy unit 8 cartesian-lambda & meta & \hyperref[sec:p10-cartesian-lambda]{\ref*{sec:p10-cartesian-lambda}} & \thepage & The product comparison uses interpreted block projections; the exponential comparison transposes interpreted evaluation after the inverse product comparison. \\
\phantomsection\label{entry:symbol-registry-00803}\(\mathbb F,j:\mathbb F^{\mathrm{op}}\to T\) & \texttt{\detokenize{\mathbb F,}\allowbreak\detokenize{j}} & finite ordinals and the structural theory functor & P10 legacy unit 7 lawvere-theories & meta & \hyperref[sec:p10-lawvere-theories]{\ref*{sec:p10-lawvere-theories}} & \thepage & The identity-on-objects functor preserves finite products; a finite-set function \(m\to n\) selects an arrow \(n\to m\) in the theory. \\
\phantomsection\label{entry:symbol-registry-00804}\(p_i^n,\langle f_i\rangle_{i<n},!_n\) & \texttt{\detokenize{p_i^n,}\allowbreak\detokenize{\langle\rangle,}\allowbreak\detokenize{!_n}} & chosen powers and tuples in a Lawvere theory & P10 legacy unit 7 lawvere-theories & meta & \hyperref[sec:p10-lawvere-theories]{\ref*{sec:p10-lawvere-theories}} & \thepage & The arrows \(p_i^n:n\to1\) exhibit \(n=1^n\), with \(p_0^1=1_1\); the empty tuple is the unique arrow \(n\to0\). \\
\phantomsection\label{entry:symbol-registry-00805}\((k,s,t),E\) & \texttt{\detokenize{(k,}\allowbreak\detokenize{s,}\allowbreak\detokenize{t),}\allowbreak\detokenize{E}} & equations with explicit finite contexts & P10 legacy unit 7 lawvere-theories & meta & \hyperref[sec:p10-lawvere-theories]{\ref*{sec:p10-lawvere-theories}} & \thepage & Each triple asserts equality of its two terms under every assignment to the declared \(k\)-variable context. \\
\phantomsection\label{entry:symbol-registry-00806}\(\operatorname{Alg}_{\mathrm{all}}(\Sigma,E),\operatorname{Alg}_{\mathrm{ne}}(\Sigma,E)\) & \texttt{\detokenize{\operatorname{Alg}_{\mathrm{all}},}\allowbreak\detokenize{\operatorname{Alg}_{\mathrm{ne}}}} & all-carrier and nonempty-carrier equational categories & P10 legacy unit 7 lawvere-theories & meta & \hyperref[sec:p10-lawvere-theories]{\ref*{sec:p10-lawvere-theories}} & \thepage & Both use carriers in the fixed universe and all operation-preserving functions; the latter is the full subcategory with nonempty carrier. \\
\phantomsection\label{entry:symbol-registry-00807}\(\equiv_E^n\) & \texttt{\detokenize{\equiv_E^n}} & equational congruence in a finite context & P10 legacy unit 7 lawvere-theories & meta & \hyperref[sec:p10-lawvere-theories]{\ref*{sec:p10-lawvere-theories}} & \thepage & The least family containing the stated equations and closed under operation compatibility and simultaneous substitutions between all finite contexts. \\
\phantomsection\label{entry:symbol-registry-00808}\(T_{\Sigma,E}(n,m)\) & \texttt{\detokenize{T_{\Sigma,E}}} & term theory of an equational presentation & P10 legacy unit 7 lawvere-theories & meta & \hyperref[sec:p10-lawvere-theories]{\ref*{sec:p10-lawvere-theories}} & \thepage & Its arrows are \(m\)-tuples of congruence classes of \(n\)-variable terms; composition is substitution and object \(0\) is terminal. \\
\phantomsection\label{entry:symbol-registry-00809}\(M_A,M_h\) & \texttt{\detokenize{M_A,}\allowbreak\detokenize{M_h}} & evaluation model and homomorphism transformation & P10 legacy unit 7 lawvere-theories & meta & \hyperref[sec:p10-lawvere-theories]{\ref*{sec:p10-lawvere-theories}} & \thepage & The model assigns \(A^n\) to \(n\) and evaluates term tuples; a homomorphism \(h\) gives components \(h^n\). \\
\phantomsection\label{entry:symbol-registry-00810}\(\operatorname{Mod}(T),U\) & \texttt{\detokenize{\operatorname{Mod},}\allowbreak\detokenize{U}} & Lawvere models and the underlying-set functor & P10 legacy unit 7 lawvere-theories & meta & \hyperref[sec:p10-lawvere-theories]{\ref*{sec:p10-lawvere-theories}} & \thepage & Finite-product-preserving functors to Set and all natural transformations; \(U(M)=M(1)\) and \(U(\alpha)=\alpha_1\). \\
\phantomsection\label{entry:symbol-registry-00811}\(c_n^M\) & \texttt{\detokenize{c_n^M}} & canonical model power comparison & P10 legacy unit 7 lawvere-theories & meta & \hyperref[sec:p10-lawvere-theories]{\ref*{sec:p10-lawvere-theories}} & \thepage & The bijection \(M(n)\to M(1)^n\) has coordinates \(M(p_i^n)\); the comparison at one uses singleton tuples. \\
\phantomsection\label{entry:symbol-registry-00812}\(\bar t,\bar\sigma\) & \texttt{\detokenize{\bar t,}\allowbreak\detokenize{\bar\sigma}} & one-component term and basic-operation arrows & P10 legacy unit 7 lawvere-theories & meta & \hyperref[sec:p10-lawvere-theories]{\ref*{sec:p10-lawvere-theories}} & \thepage & An \(n\)-variable term determines an arrow \(n\to1\); a symbol of arity \(k\) gives its basic term arrow \(k\to1\). \\
\phantomsection\label{entry:symbol-registry-00813}\(\Sigma_{\mathcal T},E_{\mathcal T},\sigma_f\) & \texttt{\detokenize{\Sigma_{\mathcal T},}\allowbreak\detokenize{E_{\mathcal T},}\allowbreak\detokenize{\sigma_f}} & named equational presentation of a theory & P10 legacy unit 7 lawvere-theories & meta & \hyperref[sec:p10-lawvere-theories]{\ref*{sec:p10-lawvere-theories}} & \thepage & One operation symbol for each arrow to the generator; projection and substitution equations recover precisely the original arrows. \\
\phantomsection\label{entry:symbol-registry-00814}\(\llbracket t\rrbracket_n\) & \texttt{\detokenize{\llbracket\rrbracket_n}} & categorical evaluation of a named algebraic term & P10 legacy unit 7 lawvere-theories & meta & \hyperref[sec:p10-lawvere-theories]{\ref*{sec:p10-lawvere-theories}} & \thepage & An arrow \(n\to1\) in the given theory; substitution of terms is composition with their tuple of interpreted arrows. \\
\phantomsection\label{entry:symbol-registry-00815}\(\varepsilon_{\mathcal T},\rho\) & \texttt{\detokenize{\varepsilon_{\mathcal T},}\allowbreak\detokenize{\rho}} & isomorphism recovering a Lawvere theory & P10 legacy unit 7 lawvere-theories & meta & \hyperref[sec:p10-lawvere-theories]{\ref*{sec:p10-lawvere-theories}} & \thepage & Evaluation from the named term theory is identity on objects; \(\rho\) denotes this fixed isomorphism when transporting free models. \\
\phantomsection\label{entry:symbol-registry-00816}\(H^*,\widehat H\) & \texttt{\detokenize{H^*,}\allowbreak\detokenize{\widehat H}} & restriction of models and translation of named terms & P10 legacy unit 7 lawvere-theories & meta & \hyperref[sec:p10-lawvere-theories]{\ref*{sec:p10-lawvere-theories}} & \thepage & A theory morphism induces precomposition on models; its named-symbol translation commutes with the recovery isomorphisms. \\
\phantomsection\label{entry:symbol-registry-00817}\(R_n,\eta^n\) & \texttt{\detokenize{R_n,}\allowbreak\detokenize{\eta^n}} & representable finite free model and its generators & P10 legacy unit 7 lawvere-theories & meta & \hyperref[sec:p10-lawvere-theories]{\ref*{sec:p10-lawvere-theories}} & \thepage & The model \(R_n=T(n,-)\) is free on \(n\); its generator indexed by \(i\) is \(p_i^n:n\to1\). \\
\phantomsection\label{entry:symbol-registry-00818}\(\theta_X,Q_X\) & \texttt{\detokenize{\theta_X,}\allowbreak\detokenize{Q_X}} & semantic term congruence and free equational algebra & P10 legacy unit 7 lawvere-theories & meta & \hyperref[sec:p10-lawvere-theories]{\ref*{sec:p10-lawvere-theories}} & \thepage & The congruence requires equality under every evaluation into an all-carrier model; its quotient includes the empty term case. \\
\phantomsection\label{entry:symbol-registry-00819}\(\mathcal F_{\mathbb T},\eta_X\) & \texttt{\detokenize{\mathcal F_{\mathbb T},}\allowbreak\detokenize{\eta_X}} & free Lawvere-model functor and generator map & P10 legacy unit 7 lawvere-theories & meta & \hyperref[sec:p10-lawvere-theories]{\ref*{sec:p10-lawvere-theories}} & \thepage & The free model has values \(Q_X^n\) after the fixed theory isomorphism; the generator maps \(x\) to the singleton tuple \(([x])\). \\
\phantomsection\label{entry:symbol-registry-00820}\(P(I),\top_I,\bot_I\) & \texttt{\detokenize{P,}\allowbreak\detokenize{\top_I,}\allowbreak\detokenize{\bot_I}} & predicate fibre and its bounds & P10 legacy unit 9 doctrine-quantifiers & meta & \hyperref[sec:p10-doctrine-quantifiers]{\ref*{sec:p10-doctrine-quantifiers}} & \thepage & A small Heyting algebra of predicates in context object \(I\); its order is semantic entailment in that context. \\
\phantomsection\label{entry:symbol-registry-00821}\(\mathsf{Heyt}\) & \texttt{\detokenize{\mathsf{Heyt}}} & category of Heyting algebras and homomorphisms & P10 legacy unit 9 doctrine-quantifiers & meta & \hyperref[sec:p10-doctrine-quantifiers]{\ref*{sec:p10-doctrine-quantifiers}} & \thepage & Maps preserve both bounds, finite meets and joins, and Heyting implication; reindexing takes values in this category. \\
\phantomsection\label{entry:symbol-registry-00822}\(f^*:P(I)\to P(J)\) & \texttt{\detokenize{f^*}} & predicate reindexing along a context map & P10 legacy unit 9 doctrine-quantifiers & meta & \hyperref[sec:p10-doctrine-quantifiers]{\ref*{sec:p10-doctrine-quantifiers}} & \thepage & For \(f:J\to I\), the contravariant functor sends predicates in \(I\) to predicates in \(J\), preserving every Heyting operation. \\
\phantomsection\label{entry:symbol-registry-00823}\(\pi_{I,A},f\times1_A\) & \texttt{\detokenize{\pi_{I,A},}\allowbreak\detokenize{f\times1_A}} & context projection and parameter substitution & P10 legacy unit 9 doctrine-quantifiers & meta & \hyperref[sec:p10-doctrine-quantifiers]{\ref*{sec:p10-doctrine-quantifiers}} & \thepage & The projection \(I\times A\to I\) forgets one variable; a parameter map lifts to the displayed product substitution square. \\
\phantomsection\label{entry:symbol-registry-00824}\(\exists_\pi,\forall_\pi\) & \texttt{\detokenize{\exists_\pi,}\allowbreak\detokenize{\forall_\pi}} & projection quantifier adjoints & P10 legacy unit 9 doctrine-quantifiers & meta & \hyperref[sec:p10-doctrine-quantifiers]{\ref*{sec:p10-doctrine-quantifiers}} & \thepage & The maps from \(P(I\times A)\) to \(P(I)\) are respectively left and right adjoints to projection reindexing. \\
\phantomsection\label{entry:symbol-registry-00825}\(\delta_A\) & \texttt{\detokenize{\delta_A}} & internal equality predicate & P10 legacy unit 9 doctrine-quantifiers & meta & \hyperref[sec:p10-doctrine-quantifiers]{\ref*{sec:p10-doctrine-quantifiers}} & \thepage & An element of \(P(A\times A)\) satisfying the full parameterized diagonal adjunction; no subobject representation is assumed. \\
\phantomsection\label{entry:symbol-registry-00826}\(e_{I,A},p_{I,A},q_{I,A},r_{I,A}\) & \texttt{\detokenize{e_{I,A},}\allowbreak\detokenize{p_{I,A},}\allowbreak\detokenize{q_{I,A},}\allowbreak\detokenize{r_{I,A}}} & maps for equality with parameters & P10 legacy unit 9 doctrine-quantifiers & meta & \hyperref[sec:p10-doctrine-quantifiers]{\ref*{sec:p10-doctrine-quantifiers}} & \thepage & Use \((I\times A)\times A\): duplicate the final variable, select either copy with its parameters, or select the equality pair. \\
\phantomsection\label{entry:symbol-registry-00827}\(L_e(\alpha)\) & \texttt{\detokenize{L_e}} & left adjoint for parameterized equality & P10 legacy unit 9 doctrine-quantifiers & meta & \hyperref[sec:p10-doctrine-quantifiers]{\ref*{sec:p10-doctrine-quantifiers}} & \thepage & Defined by \(p^*\alpha\wedge r^*\delta_A\), with \(L_e\dashv e^*\); its full order adjunction gives equality elimination. \\
\phantomsection\label{entry:symbol-registry-00828}\(\Delta_A,\tau_A\) & \texttt{\detokenize{\Delta_A,}\allowbreak\detokenize{\tau_A}} & diagonal and exchange of an equality pair & P10 legacy unit 9 doctrine-quantifiers & meta & \hyperref[sec:p10-doctrine-quantifiers]{\ref*{sec:p10-doctrine-quantifiers}} & \thepage & The arrows \(A\to A\times A\) and \(A\times A\to A\times A\) used to state reflexivity and symmetry. \\
\phantomsection\label{entry:symbol-registry-00829}\(\exists_{\pi_{1,A}}\top_{1\times A}=\top_1\) & \texttt{\detokenize{\exists_{\pi_{1,A}}}} & internal inhabitance of a sort & P10 legacy unit 9 doctrine-quantifiers & meta & \hyperref[sec:p10-doctrine-quantifiers]{\ref*{sec:p10-doctrine-quantifiers}} & \thepage & Equivalent by Beck--Chevalley to the corresponding condition in every parameter context; it implies reflection of predicate order, not of base-arrow equality. \\
\phantomsection\label{entry:symbol-registry-00830}\(P(X)=\mathcal P(X)\) & \texttt{\detokenize{\mathcal P}} & powerset instance of a predicate fibre & P10 legacy unit 9 doctrine-quantifiers & meta & \hyperref[sec:p10-doctrine-quantifiers]{\ref*{sec:p10-doctrine-quantifiers}} & \thepage & Boolean algebra of all subsets of a small set \(X\), including the one-element fibre for empty \(X\). \\
\phantomsection\label{entry:symbol-registry-00831}\(U\Rightarrow_X V\) & \texttt{\detokenize{\Rightarrow_X}} & Boolean implication between subsets & P10 legacy unit 9 doctrine-quantifiers & meta & \hyperref[sec:p10-doctrine-quantifiers]{\ref*{sec:p10-doctrine-quantifiers}} & \thepage & The subset \((X\setminus U)\cup V\) is the Heyting residual in the powerset fibre. \\
\phantomsection\label{entry:symbol-registry-00832}\(\exists_f,\forall_f\) & \texttt{\detokenize{\exists_f,}\allowbreak\detokenize{\forall_f}} & powerset quantifiers along an arbitrary function & P10 legacy unit 9 doctrine-quantifiers & meta & \hyperref[sec:p10-doctrine-quantifiers]{\ref*{sec:p10-doctrine-quantifiers}} & \thepage & Direct image and the condition that the entire fibre belongs to the predicate; empty fibres satisfy the universal condition. \\
\phantomsection\label{entry:symbol-registry-00833}\(\Gamma,\Phi\) & \texttt{\detokenize{\Gamma,}\allowbreak\detokenize{\Phi}} & variable context and hypothesis list in Pitts's calculus & P10 legacy unit 9 doctrine-quantifiers & meta & \hyperref[sec:p10-doctrine-quantifiers]{\ref*{sec:p10-doctrine-quantifiers}} & \thepage & The first is a finite list of distinct sorted variables; the second is a finite list of formulas in that context. \\
\phantomsection\label{entry:symbol-registry-00834}\(M:\sigma[\Gamma],M=N:\sigma[\Gamma]\) & \texttt{\detokenize{M:\sigma,}\allowbreak\detokenize{M=N:\sigma}} & term formation and external term-equation judgments & P10 legacy unit 9 doctrine-quantifiers & meta & \hyperref[sec:p10-doctrine-quantifiers]{\ref*{sec:p10-doctrine-quantifiers}} & \thepage & Formation is well-sortedness in the retained context; an equation is interpreted as equality of term arrows. \\
\phantomsection\label{entry:symbol-registry-00835}\(\Phi\vdash\psi[\Gamma]\) & \texttt{\detokenize{\vdash}} & contextual predicate sequent & P10 legacy unit 9 doctrine-quantifiers & meta & \hyperref[sec:p10-doctrine-quantifiers]{\ref*{sec:p10-doctrine-quantifiers}} & \thepage & The conjunction of the finite assumptions is below the conclusion in \(P(I_\Gamma)\); unused variable declarations are retained. \\
\phantomsection\label{entry:symbol-registry-00836}\(M=_\sigma N\) & \texttt{\detokenize{=_\sigma}} & equality as a first-order formula & P10 legacy unit 9 doctrine-quantifiers & meta & \hyperref[sec:p10-doctrine-quantifiers]{\ref*{sec:p10-doctrine-quantifiers}} & \thepage & Reindexes \(\delta_{A_\sigma}\) along the pair of term arrows; its truth need not reflect external equality of those arrows. \\
\phantomsection\label{entry:symbol-registry-00837}\(\varphi\mathbin{\&}\psi,\varphi\Rightarrow\psi,\mathsf{false}\) & \texttt{\detokenize{\&,}\allowbreak\detokenize{\Rightarrow,}\allowbreak\detokenize{\mathsf{false}}} & Pitts's propositional formula constructors & P10 legacy unit 9 doctrine-quantifiers & meta & \hyperref[sec:p10-doctrine-quantifiers]{\ref*{sec:p10-doctrine-quantifiers}} & \thepage & Conjunction, implication, and falsity at the formula level; their denotations are the corresponding Heyting operations. \\
\phantomsection\label{entry:symbol-registry-00838}\(\forall_\sigma(\varphi),\exists_\sigma(\varphi)\) & \texttt{\detokenize{\forall_\sigma,}\allowbreak\detokenize{\exists_\sigma}} & quantifiers on unary formula abstractions & P10 legacy unit 9 doctrine-quantifiers & meta & \hyperref[sec:p10-doctrine-quantifiers]{\ref*{sec:p10-doctrine-quantifiers}} & \thepage & Pitts's printed quantifier notation; the body \(\varphi(x)\) is interpreted in the context extended by \(x:\sigma\). \\
\phantomsection\label{entry:symbol-registry-00839}\((x)e,\varphi(M)\) & \texttt{\detokenize{(x)e,}\allowbreak\detokenize{\varphi(M)}} & metalevel formula abstraction and instantiation & P10 legacy unit 9 doctrine-quantifiers & meta & \hyperref[sec:p10-doctrine-quantifiers]{\ref*{sec:p10-doctrine-quantifiers}} & \thepage & If \(\varphi=(x)e\), then \(\varphi(M)=e[M/x]\); this is binding notation, not an additional first-order connective. \\
\phantomsection\label{entry:symbol-registry-00840}\(I_\Gamma\) & \texttt{\detokenize{I_\Gamma}} & interpreted finite variable context & P10 legacy unit 9 doctrine-quantifiers & meta & \hyperref[sec:p10-doctrine-quantifiers]{\ref*{sec:p10-doctrine-quantifiers}} & \thepage & A chosen product starting with terminal object \(1\), extended by one sort object for each declaration. \\
\phantomsection\label{entry:symbol-registry-00841}\(\llbracket M\rrbracket_\Gamma,\llbracket\varphi\rrbracket_\Gamma\) & \texttt{\detokenize{\llbracket\rrbracket_\Gamma}} & contextual term and predicate interpretations & P10 legacy unit 9 doctrine-quantifiers & meta & \hyperref[sec:p10-doctrine-quantifiers]{\ref*{sec:p10-doctrine-quantifiers}} & \thepage & A term denotes an arrow from \(I_\Gamma\) to its sort object; a formula denotes an element of \(P(I_\Gamma)\). \\
\phantomsection\label{entry:symbol-registry-00842}\(s:I_\Delta\to I_\Gamma\) & \texttt{\detokenize{s}} & interpreted simultaneous term substitution & P10 legacy unit 9 doctrine-quantifiers & meta & \hyperref[sec:p10-doctrine-quantifiers]{\ref*{sec:p10-doctrine-quantifiers}} & \thepage & Its components interpret one substituting term per variable of \(\Gamma\); formula substitution is reindexing along this map. \\
\phantomsection\label{entry:symbol-registry-00843}\(\mathcal U\) & \texttt{\detokenize{\mathcal U}} & universe used for categorical size bounds & P10 legacy units 7–9 lawvere-theories to doctrine-quantifiers & meta & \hyperref[sec:p10-lawvere-theories]{\ref*{sec:p10-lawvere-theories}} & \thepage & May be the ambient collection of all sets; a set-sized relativization has the explicitly stated Grothendieck closures, with existence treated as a hypothesis. \\
\phantomsection\label{entry:symbol-registry-00844}\(\mathrm{ND},\mathrm{ADJ}\) & \texttt{\detokenize{\mathrm{ND},}\allowbreak\detokenize{\mathrm{ADJ}}} & local abbreviations for two contextual rule presentations & P10 legacy unit 9 doctrine-quantifiers & meta & \hyperref[sec:p10-doctrine-quantifiers]{\ref*{sec:p10-doctrine-quantifiers}} & \thepage & The same theory and structural/equational basis, with either the natural-deduction logical rules or every direction of Pitts's adjoint bi-rules. \\
\phantomsection\label{entry:symbol-registry-00845}\(\mathrm{Eq0}\) & \texttt{\detokenize{\mathrm{Eq0}}} & parameter-free equality bi-rule & P10 legacy unit 9 doctrine-quantifiers & meta & \hyperref[sec:p10-doctrine-quantifiers]{\ref*{sec:p10-doctrine-quantifiers}} & \thepage & Both directions of Pitts's Remark 5.6.2; the displayed variable contexts have no hidden parameter prefix. \\
\phantomsection\label{entry:symbol-registry-00846}\((x)(x')e,\psi(M,N)\) & \texttt{\detokenize{(x)(x')e,}\allowbreak\detokenize{\psi(M,}\allowbreak\detokenize{N)}} & binary formula abstraction and simultaneous instantiation & P10 legacy unit 9 doctrine-quantifiers & meta & \hyperref[sec:p10-doctrine-quantifiers]{\ref*{sec:p10-doctrine-quantifiers}} & \thepage & Application replaces the two displayed argument variables simultaneously in the body, avoiding capture. \\
\phantomsection\label{entry:symbol-registry-00847}\(\widehat{\mathcal C}\) & \texttt{\detokenize{\widehat{\mathcal C}}} & presheaf category on a small base & P10 legacy unit 9 doctrine-quantifiers & meta & \hyperref[sec:p10-doctrine-quantifiers]{\ref*{sec:p10-doctrine-quantifiers}} & \thepage & The category of functors \(\mathcal C^{\mathrm{op}}\to\mathsf{Set}\) and natural transformations, under the declared size convention. \\
\phantomsection\label{entry:symbol-registry-00848}\(\mathrm y c,\mathrm y(u)\) & \texttt{\detokenize{\mathrm y}} & representable presheaf and the Yoneda map & P10 legacy unit 9 doctrine-quantifiers & meta & \hyperref[sec:p10-doctrine-quantifiers]{\ref*{sec:p10-doctrine-quantifiers}} & \thepage & The value at \(d\) is \(\mathcal C(d,c)\); an arrow between representing objects acts by postcomposition. \\
\phantomsection\label{entry:symbol-registry-00849}\(X(u),\alpha_c\) & \texttt{\detokenize{X(u),}\allowbreak\detokenize{\alpha_c}} & presheaf restriction and natural-transformation component & P10 legacy unit 9 doctrine-quantifiers & meta & \hyperref[sec:p10-doctrine-quantifiers]{\ref*{sec:p10-doctrine-quantifiers}} & \thepage & For \(u:d\to c\), restriction maps \(X(c)\) to \(X(d)\); the component \(\alpha_c\) maps \(X(c)\) to \(Y(c)\). \\
\phantomsection\label{entry:symbol-registry-00850}\(\operatorname{Sub}(X)\) & \texttt{\detokenize{\operatorname{Sub}}} & ordered set of subpresheaves & P10 legacy unit 9 doctrine-quantifiers & meta & \hyperref[sec:p10-doctrine-quantifiers]{\ref*{sec:p10-doctrine-quantifiers}} & \thepage & Restriction-stable families of component subsets; canonical images represent categorical subobjects and their factorization order. \\
\phantomsection\label{entry:symbol-registry-00851}\(0(c)=\varnothing,\ 1(c)=\{*\}\) & \texttt{\detokenize{0,}\allowbreak\detokenize{1}} & constant initial and terminal presheaves & P10 legacy unit 9 doctrine-quantifiers & meta & \hyperref[sec:p10-doctrine-quantifiers]{\ref*{sec:p10-doctrine-quantifiers}} & \thepage & These denote presheaves; numeric object names of a particular base category have their separately declared base-object role. \\
\phantomsection\label{entry:symbol-registry-00852}\(\operatorname{Im}(m),q,i\) & \texttt{\detokenize{\operatorname{Im},}\allowbreak\detokenize{q,}\allowbreak\detokenize{i}} & componentwise image and its factorization & P10 legacy unit 9 doctrine-quantifiers & meta & \hyperref[sec:p10-doctrine-quantifiers]{\ref*{sec:p10-doctrine-quantifiers}} & \thepage & The natural map factors through its image by a componentwise surjection followed by an inclusion; a monic map is isomorphic to that inclusion. \\
\phantomsection\label{entry:symbol-registry-00853}\(\Omega,\mathsf{true}:1\to\Omega\) & \texttt{\detokenize{\Omega,}\allowbreak\detokenize{\mathsf{true}}} & subobject classifier and truth map & P10 legacy unit 9 doctrine-quantifiers & meta & \hyperref[sec:p10-doctrine-quantifiers]{\ref*{sec:p10-doctrine-quantifiers}} & \thepage & Every monomorphism is a pullback along its unique characteristic arrow; the definition retains the full universal property. \\
\phantomsection\label{entry:symbol-registry-00854}\((Y^X)(c)\) & \texttt{\detokenize{Y^X}} & component of a presheaf exponential & P10 legacy unit 9 doctrine-quantifiers & meta & \hyperref[sec:p10-doctrine-quantifiers]{\ref*{sec:p10-doctrine-quantifiers}} & \thepage & The set \(\operatorname{Nat}(\mathrm y c\times X,Y)\), with restriction by precomposition; it need not be the pointwise external function set. \\
\phantomsection\label{entry:symbol-registry-00855}\(\operatorname{ev}_c,\Lambda h\) & \texttt{\detokenize{\operatorname{ev}_c,}\allowbreak\detokenize{\Lambda}} & presheaf evaluation and transposition & P10 legacy unit 9 doctrine-quantifiers & meta & \hyperref[sec:p10-doctrine-quantifiers]{\ref*{sec:p10-doctrine-quantifiers}} & \thepage & Evaluation uses the representing identity arrow; transposition uses every restriction of the parameter element. \\
\phantomsection\label{entry:symbol-registry-00856}\(\Omega(c),\Omega(u),t(c)\) & \texttt{\detokenize{\Omega(c),}\allowbreak\detokenize{\Omega(u),}\allowbreak\detokenize{t(c)}} & sieves their restrictions and the maximal sieve & P10 legacy unit 9 doctrine-quantifiers & meta & \hyperref[sec:p10-doctrine-quantifiers]{\ref*{sec:p10-doctrine-quantifiers}} & \thepage & A sieve is a set of arrows into \(c\) closed under precomposition; the maximal sieve contains the identity. \\
\phantomsection\label{entry:symbol-registry-00857}\(\chi_A:X\to\Omega\) & \texttt{\detokenize{\chi_A}} & characteristic map of a subpresheaf & P10 legacy unit 9 doctrine-quantifiers & meta & \hyperref[sec:p10-doctrine-quantifiers]{\ref*{sec:p10-doctrine-quantifiers}} & \thepage & At \(x\in X(c)\), it is the sieve of arrows \(u:d\to c\) whose restriction \(X(u)(x)\) belongs to \(A(d)\). \\
\phantomsection\label{entry:symbol-registry-00858}\(P(X)=\Omega^X\) & \texttt{\detokenize{P(X)}} & power object of a presheaf & P10 legacy unit 9 doctrine-quantifiers & meta & \hyperref[sec:p10-doctrine-quantifiers]{\ref*{sec:p10-doctrine-quantifiers}} & \thepage & Here \(X\) is a presheaf and \(P(X)\) is an object of \(\widehat{\mathcal C}\), distinct from the earlier powerset-fibre usage for sets. \\
\phantomsection\label{entry:symbol-registry-00859}\(\in_X\hookrightarrow P(X)\times X\) & \texttt{\detokenize{\in_X}} & membership subobject of a power object & P10 legacy unit 9 doctrine-quantifiers & meta & \hyperref[sec:p10-doctrine-quantifiers]{\ref*{sec:p10-doctrine-quantifiers}} & \thepage & The pullback of truth along evaluation; under stage classification membership tests whether \((1_c,x)\) lies in the represented subpresheaf. \\
\phantomsection\label{entry:symbol-registry-00860}\(S|u\) & \texttt{\detokenize{S|u}} & restriction of a classified power-object element & P10 legacy unit 9 doctrine-quantifiers & meta & \hyperref[sec:p10-doctrine-quantifiers]{\ref*{sec:p10-doctrine-quantifiers}} & \thepage & The pullback of \(S\subseteq\mathrm y c\times X\) along \(\mathrm y u\times1_X\). \\
\phantomsection\label{entry:symbol-registry-00861}\(S_z\) & \texttt{\detokenize{S_z}} & subpresheaf determined by a relation and a parameter element & P10 legacy unit 9 doctrine-quantifiers & meta & \hyperref[sec:p10-doctrine-quantifiers]{\ref*{sec:p10-doctrine-quantifiers}} & \thepage & For \(z\in Z(c)\), its component consists of pairs \((u,x)\) with \((Z(u)(z),x)\) in the original relation. \\
\phantomsection\label{entry:symbol-registry-04251}\(x\cdot u\) & \texttt{\detokenize{\cdot}} & restriction of a presheaf element along u:d to c & p10/n-d-023-source-03 & n/a & \hyperref[sec:p10-n-d-023-source-03]{\ref*{sec:p10-n-d-023-source-03}} & \thepage & Source notation x.u equals host X(u)(x); x at c and output at d. \\
\phantomsection\label{entry:symbol-registry-04252}\(\Rightarrow\) & \texttt{\detokenize{\Rightarrow}} & Heyting implication of subobjects or internal truth values & p10/n-d-023-source-03 & n/a & \hyperref[thm:p10-mlm-presheaf-predicates]{\ref*{thm:p10-mlm-presheaf-predicates}} & \thepage & Typed operation on Sub(X) or on P(X); not material implication. \\
\phantomsection\label{entry:symbol-registry-04253}\(f^*\) & \texttt{\detokenize{{}^*}} & inverse-image operation on external subobject posets & p10/n-d-023-source-03 & n/a & \hyperref[thm:p10-mlm-presheaf-predicates]{\ref*{thm:p10-mlm-presheaf-predicates}} & \thepage & Source f inverse superscript; distinguish from internal P(f). \\
\phantomsection\label{entry:symbol-registry-04254}\(\exists_f\) & \texttt{\detokenize{\exists}} & left order adjoint on external subobject posets & p10/n-d-023-source-03 & n/a & \hyperref[thm:p10-mlm-presheaf-predicates]{\ref*{thm:p10-mlm-presheaf-predicates}} & \thepage & Input subobject of X; output subobject of Y for f:X to Y. \\
\phantomsection\label{entry:symbol-registry-04255}\(\forall_f\) & \texttt{\detokenize{\forall}} & right order adjoint on external subobject posets & p10/n-d-023-source-03 & n/a & \hyperref[thm:p10-mlm-presheaf-predicates]{\ref*{thm:p10-mlm-presheaf-predicates}} & \thepage & All-arrow formula restricts target y to stage d before comparison. \\
\phantomsection\label{entry:symbol-registry-04256}\(P(f)\) & \texttt{\detokenize{P}} & internal inverse-image map P(Y) to P(X) & p10/n-d-023-source-03 & n/a & \hyperref[def:p10-mlm-internal-operations]{\ref*{def:p10-mlm-internal-operations}} & \thepage & Contravariant power-object operation; not an external poset arrow. \\
\phantomsection\label{entry:symbol-registry-04257}\(\exists_f\) & \texttt{\detokenize{\exists}} & internal power-object map P(X) to P(Y) & p10/n-d-023-source-03 & n/a & \hyperref[def:p10-mlm-internal-operations]{\ref*{def:p10-mlm-internal-operations}} & \thepage & Same source glyph; typed role declared separately from external adjoint. \\
\phantomsection\label{entry:symbol-registry-04258}\(\forall_f\) & \texttt{\detokenize{\forall}} & internal power-object map P(X) to P(Y) & p10/n-d-023-source-03 & n/a & \hyperref[def:p10-mlm-internal-operations]{\ref*{def:p10-mlm-internal-operations}} & \thepage & Same source glyph; typed role declared separately from external adjoint. \\
\phantomsection\label{entry:symbol-registry-04259}\([\phi]\) & \texttt{\detokenize{[\phi]}} & extension subobject of a formula in its declared context & p10/n-d-023-source-03 & n/a & \hyperref[def:p10-mlm-language]{\ref*{def:p10-mlm-language}} & \thepage & Editorial notation for source set-builder extension. \\
\phantomsection\label{entry:symbol-registry-04260}\(\delta_X\) & \texttt{\detokenize{\delta}} & characteristic map of the diagonal X to X times X & p10/n-d-023-source-03 & n/a & \hyperref[def:p10-mlm-language]{\ref*{def:p10-mlm-language}} & \thepage & Interprets typed equality; distinct from the diagonal arrow itself. \\
\phantomsection\label{entry:symbol-registry-04261}\(\lambda x.t\) & \texttt{\detokenize{\lambda}} & exponential transposition in the internal language & p10/n-d-023-source-03 & n/a & \hyperref[def:p10-mlm-language]{\ref*{def:p10-mlm-language}} & \thepage & Source notation; capture-avoiding binding with all parameter types retained. \\
\phantomsection\label{entry:symbol-registry-04262}\(\{x:X\mid\phi\}\) & \texttt{\detokenize{\mid}} & predicate-valued comprehension term G to P(X) & p10/n-d-023-source-03 & n/a & \hyperref[def:p10-mlm-language]{\ref*{def:p10-mlm-language}} & \thepage & Separate typed role from the extension object using the same braces. \\
\phantomsection\label{entry:symbol-registry-04263}\(\{x:X\mid\phi\}\) & \texttt{\detokenize{\mid}} & extension subobject of G times X & p10/n-d-023-source-03 & n/a & \hyperref[def:p10-mlm-language]{\ref*{def:p10-mlm-language}} & \thepage & Separate typed role from the predicate-valued term using the same braces. \\
\phantomsection\label{entry:symbol-registry-04264}\(U\forces\phi(\alpha)\) & \texttt{\detokenize{\forces}} & generalized-element factorization forcing in an elementary topos & p10/n-d-023-source-03 & relation & \hyperref[def:p10-mlm-generalized-forcing]{\ref*{def:p10-mlm-generalized-forcing}} & \thepage & U is a topos object; disjunction and existence are epimorphically local. \\
\phantomsection\label{entry:symbol-registry-04265}\(c\forces\phi(\gamma)\) & \texttt{\detokenize{\forces}} & site-stage or presheaf-stage semantic membership forcing & p10/n-d-023-source-03 & relation & \hyperref[def:p10-mlm-site-setup]{\ref*{def:p10-mlm-site-setup}} & \thepage & c is a base object; site and trivial-topology clauses distinguished. \\
\phantomsection\label{entry:symbol-registry-04266}\(J(c)\) & \texttt{\detokenize{J}} & set of covering sieves on c & p10/n-d-023-source-03 & n/a & \hyperref[def:p10-mlm-site-setup]{\ref*{def:p10-mlm-site-setup}} & \thepage & Empty covering sieve permits bottom at a site stage. \\
\phantomsection\label{entry:symbol-registry-04267}\(a(\mathrm yc)\) & \texttt{\detokenize{a}} & associated sheaf of a representable & p10/n-d-023-source-03 & n/a & \hyperref[def:p10-mlm-site-setup]{\ref*{def:p10-mlm-site-setup}} & \thepage & Already established associated-sheaf functor; construction remains outside scope. \\
\phantomsection\label{entry:symbol-registry-04268}\(\mathcal L_1\) & \texttt{\detokenize{\mathcal L}} & Lambek-Scott1980 pure product/power type theory & p10/n-d-023-source-03 & n/a & \hyperref[def:p10-ls1980-language]{\ref*{def:p10-ls1980-language}} & \thepage & Full source includes N and Peano; nonarithmetic restriction and applied extension are explicitly separate. \\
\phantomsection\label{entry:symbol-registry-04269}\(p\vdash_Xq\) & \texttt{\detokenize{\vdash}} & source variable-set-indexed derivability relation & p10/n-d-023-source-03 & relation & \hyperref[def:p10-ls1980-contextual-validity]{\ref*{def:p10-ls1980-contextual-validity}} & \thepage & Not generalized forcing; X is a declared typed variable set, possibly infinite. \\
\phantomsection\label{entry:symbol-registry-04270}\(p\models_Xq\) & \texttt{\detokenize{\models}} & semantic entailment for all generalized assignments on X & p10/n-d-023-source-03 & relation & \hyperref[def:p10-ls1980-contextual-validity]{\ref*{def:p10-ls1980-contextual-validity}} & \thepage & Finite X means subobject inclusion in its full context; infinite X does not assume an infinite product. \\
\phantomsection\label{entry:symbol-registry-04271}\(a=_Ab\) & \texttt{\detokenize{=}} & source Leibniz-defined equality formula at type A & p10/n-d-023-source-03 & relation & \hyperref[def:p10-ls1980-language]{\ref*{def:p10-ls1980-language}} & \thepage & Abbreviates universal PA membership equivalence; proven to classify the diagonal, not made primitive. \\
\phantomsection\label{entry:symbol-registry-04272}\(*\) & \texttt{\detokenize{*}} & source closed unit term of type 1 & p10/n-d-023-source-03 & n/a & \hyperref[def:p10-ls1980-language]{\ref*{def:p10-ls1980-language}} & \thepage & Distinct from the meta-level symbol for subobject pullback. \\
\phantomsection\label{entry:symbol-registry-04273}\(\operatorname{Fun}(A,B)\) & \texttt{\detokenize{\operatorname{Fun}}} & subobject of P(A times B) of total single-valued relations & p10/n-d-023-source-03 & n/a & \hyperref[thm:p10-ls1980-exponential-graph-bridge]{\ref*{thm:p10-ls1980-exponential-graph-bridge}} & \thepage & Represents B\^{}A by a proved parameter-natural graph correspondence; not a source primitive function type. \\
\phantomsection\label{entry:symbol-registry-04274}\(N\) & \texttt{\detokenize{N}} & optional source natural-number type interpreted as an NNO & p10/n-d-023-source-03 & n/a & \hyperref[def:p10-ls1980-language]{\ref*{def:p10-ls1980-language}} & \thepage & Absent from the explicitly nonarithmetic fragment. \\
\phantomsection\label{entry:symbol-registry-04275}\(S\) & \texttt{\detokenize{S}} & optional source successor term constructor & p10/n-d-023-source-03 & n/a & \hyperref[def:p10-ls1980-language]{\ref*{def:p10-ls1980-language}} & \thepage & Interpreted as NNO successor; Peano validation derived from recursion universal property. \\
\phantomsection\label{entry:symbol-registry-04513}\(T\) & \texttt{\detokenize{T}} & original source subobject classifier and internal truth-value object & p10/n-d-023-source-01 & n/a & \hyperref[def:p10-lawvere-sheaves-topos]{\ref*{def:p10-lawvere-sheaves-topos}} & \thepage & Identified with host Omega only in the concrete presheaf bridge; source T is not a theory here. \\
\phantomsection\label{entry:symbol-registry-04514}\(\mathsf{true}\) & \texttt{\detokenize{\mathsf{true}}} & classifier truth map 1 to T & p10/n-d-023-source-01 & n/a & \hyperref[def:p10-lawvere-sheaves-topos]{\ref*{def:p10-lawvere-sheaves-topos}} & \thepage & Subobjects are its unique characteristic pullbacks. \\
\phantomsection\label{entry:symbol-registry-04515}\(\{X/\varphi\}\) & \texttt{\detokenize{\{X/\varphi\}}} & source subobject of X classified by predicate phi:X to T & p10/n-d-023-source-01 & n/a & \hyperref[def:p10-lawvere-sheaves-topos]{\ref*{def:p10-lawvere-sheaves-topos}} & \thepage & Original source notation; not a quotient object. \\
\phantomsection\label{entry:symbol-registry-04516}\(T^X\) & \texttt{\detokenize{{}^X}} & internal predicate or power object of X & p10/n-d-023-source-01 & n/a & \hyperref[def:p10-lawvere-sheaves-topos]{\ref*{def:p10-lawvere-sheaves-topos}} & \thepage & Classifier exponential; use internal order for adjunctions. \\
\phantomsection\label{entry:symbol-registry-04517}\(T^f\) & \texttt{\detokenize{{}^f}} & precomposition map from T\^{}Y to T\^{}X for f:X to Y & p10/n-d-023-source-01 & n/a & \hyperref[def:p10-lawvere-sheaves-topos]{\ref*{def:p10-lawvere-sheaves-topos}} & \thepage & Contravariant exponent action; not a pointwise function-set replacement. \\
\phantomsection\label{entry:symbol-registry-04518}\(\exists_f\) & \texttt{\detokenize{\exists}} & internal predicate-object left order adjoint of T\^{}f & p10/n-d-023-source-01 & n/a & \hyperref[def:p10-lawvere-sheaves-topos]{\ref*{def:p10-lawvere-sheaves-topos}} & \thepage & Map T\^{}X to T\^{}Y; distinguished from its induced global subobject-poset operation. \\
\phantomsection\label{entry:symbol-registry-04519}\(\forall_f\) & \texttt{\detokenize{\forall}} & internal predicate-object right order adjoint of T\^{}f & p10/n-d-023-source-01 & n/a & \hyperref[def:p10-lawvere-sheaves-topos]{\ref*{def:p10-lawvere-sheaves-topos}} & \thepage & Map T\^{}X to T\^{}Y; internal completeness is not external infinite completeness. \\
\phantomsection\label{entry:symbol-registry-04520}\(f^*\) & \texttt{\detokenize{{}^*}} & pullback functor on families or its restriction to subobjects & p10/n-d-023-source-01 & n/a & \hyperref[def:p10-lawvere-sheaves-families]{\ref*{def:p10-lawvere-sheaves-families}} & \thepage & For f:X to Y maps E/Y to E/X; typed separately from a geometric inverse image. \\
\phantomsection\label{entry:symbol-registry-04521}\(\Pi_f\) & \texttt{\detokenize{\Pi}} & dependent-product right adjoint on slices & p10/n-d-023-source-01 & n/a & \hyperref[def:p10-lawvere-sheaves-families]{\ref*{def:p10-lawvere-sheaves-families}} & \thepage & Maps E/X to E/Y; includes arbitrary families and not only predicates. \\
\phantomsection\label{entry:symbol-registry-04522}\(\sigma_X\) & \texttt{\detokenize{\sigma}} & support characteristic predicate of a family E to X & p10/n-d-023-source-01 & n/a & \hyperref[def:p10-lawvere-sheaves-families]{\ref*{def:p10-lawvere-sheaves-families}} & \thepage & Takes the image subobject; values in E(X,T). \\
\phantomsection\label{entry:symbol-registry-04523}\(\widehat C\) & \texttt{\detokenize{\widehat}} & internal category of contravariant actions on C-indexed families & p10/n-d-023-source-01 & n/a & \hyperref[def:p10-lawvere-sheaves-families]{\ref*{def:p10-lawvere-sheaves-families}} & \thepage & Source C-op E; not solely an external Set-valued presheaf category. \\
\phantomsection\label{entry:symbol-registry-04524}\(f_0\) & \texttt{\detokenize{f_0}} & source left adjoint to restriction along an internal functor & p10/n-d-023-source-01 & n/a & \hyperref[def:p10-lawvere-sheaves-families]{\ref*{def:p10-lawvere-sheaves-families}} & \thepage & From C-hat to C-prime-hat. \\
\phantomsection\label{entry:symbol-registry-04525}\(f_*\) & \texttt{\detokenize{f_*}} & source right adjoint to restriction along an internal functor & p10/n-d-023-source-01 & n/a & \hyperref[def:p10-lawvere-sheaves-families]{\ref*{def:p10-lawvere-sheaves-families}} & \thepage & From C-hat to C-prime-hat; not geometric unless separately stated. \\
\phantomsection\label{entry:symbol-registry-04526}\(\Omega\) & \texttt{\detokenize{\Omega}} & source category encoding an object with one endomap & p10/n-d-023-source-01/recursion & n/a & \hyperref[sec:p10-n-d-023-source-01]{\ref*{sec:p10-n-d-023-source-01}} & \thepage & In E\^{}Omega only; source truth object remains T. \\
\phantomsection\label{entry:symbol-registry-04527}\(\omega\) & \texttt{\detokenize{\omega}} & source recursion object and finite-coordinate index & p10/n-d-023-source-01 & n/a & \hyperref[sec:p10-n-d-023-source-01]{\ref*{sec:p10-n-d-023-source-01}} & \thepage & Additional omega condition; not assumed merely from the original topos definition. \\
\phantomsection\label{entry:symbol-registry-04528}\(u_*\) & \texttt{\detokenize{u_*}} & geometric direct-image functor with exact left adjoint & p10/n-d-023-source-01/morphisms & n/a & \hyperref[def:p10-lawvere-sheaves-morphisms]{\ref*{def:p10-lawvere-sheaves-morphisms}} & \thepage & Direction E-prime to E-double-prime. \\
\phantomsection\label{entry:symbol-registry-04529}\(u^*\) & \texttt{\detokenize{u^*}} & geometric inverse-image exact left adjoint & p10/n-d-023-source-01/morphisms & n/a & \hyperref[def:p10-lawvere-sheaves-morphisms]{\ref*{def:p10-lawvere-sheaves-morphisms}} & \thepage & Not automatically logical or exponential-preserving. \\
\phantomsection\label{entry:symbol-registry-04530}\(j\) & \texttt{\detokenize{j}} & internal topology on T & p10/n-d-023-source-01/sheaves & n/a & \hyperref[def:p10-lawvere-sheaves-topology]{\ref*{def:p10-lawvere-sheaves-topology}} & \thepage & Idempotent and preserves truth and conjunction; defines universal closure. \\
\phantomsection\label{entry:symbol-registry-04531}\(T_j\) & \texttt{\detokenize{T_j}} & image of topology j as a subobject of T & p10/n-d-023-source-01/sheaves & n/a & \hyperref[def:p10-lawvere-sheaves-topology]{\ref*{def:p10-lawvere-sheaves-topology}} & \thepage & Equals Eq(j,id) by the split-idempotent proof; dense-extension property now proved editorially. \\
\phantomsection\label{entry:symbol-registry-04532}\(u_*\) & \texttt{\detokenize{u_*}} & finite-limit-preserving base functor in generalized-ultraproduct datum & p10/n-d-023-source-01/ultraproduct & n/a & \hyperref[def:p10-lawvere-sheaves-ultraproduct]{\ref*{def:p10-lawvere-sheaves-ultraproduct}} & \thepage & Not required to have a geometric inverse image; same source glyph but different explicitly typed standing. \\
\phantomsection\label{entry:symbol-registry-04533}\(h\) & \texttt{\detokenize{h}} & internal Heyting homomorphism u\_*(T) to T\_0 & p10/n-d-023-source-01/ultraproduct & n/a & \hyperref[def:p10-lawvere-sheaves-ultraproduct]{\ref*{def:p10-lawvere-sheaves-ultraproduct}} & \thepage & Its top-value condition determines the inverted monos. \\
\phantomsection\label{entry:symbol-registry-04534}\(\mathcal E_h\) & \texttt{\detokenize{\mathcal E_h}} & source category localization inverting the specified monos & p10/n-d-023-source-01/ultraproduct & n/a & \hyperref[def:p10-lawvere-sheaves-ultraproduct]{\ref*{def:p10-lawvere-sheaves-ultraproduct}} & \thepage & Announced topos/logical-morphism theorem; not an ordinary single product quotient. \\
\phantomsection\label{entry:symbol-registry-04535}\(\mathsf{BVM}/\mathcal S\) & \texttt{\detokenize{\mathsf{BVM}}} & source bicategory of Boolean-valued topoi over abstract-set base S & p10/n-d-023-source-01/bvm & n/a & \hyperref[def:p10-lawvere-sheaves-bvm]{\ref*{def:p10-lawvere-sheaves-bvm}} & \thepage & Boolean includes the Skolem condition as well as classical internal logic. \\
\phantomsection\label{entry:symbol-registry-04536}\(\mathsf{CBA}(\mathcal S)\) & \texttt{\detokenize{\mathsf{CBA}}} & source category of S-complete Boolean algebra objects & p10/n-d-023-source-01/bvm & n/a & \hyperref[def:p10-lawvere-sheaves-bvm]{\ref*{def:p10-lawvere-sheaves-bvm}} & \thepage & Equivalence to BVM/S is announced; missing morphism/functor proof remains owed. \\
\phantomsection\label{entry:symbol-registry-04537}\(P\) & \texttt{\detokenize{P}} & poset of finite partial maps omega times I to 2 & p10/lawvere-ch-chain-poset-reuse & n/a & \hyperref[thm:p10-lawvere-sheaves-ch]{\ref*{thm:p10-lawvere-sheaves-ch}} & \thepage & Ordered by extension; distinct from the spectrum membership object. \\
\phantomsection\label{entry:symbol-registry-04538}\(\widetilde P\) & \texttt{\detokenize{\widetilde P}} & double-negation sheaf topos of P S & p10/lawvere-ch-chain-poset-reuse & n/a & \hyperref[thm:p10-lawvere-sheaves-ch]{\ref*{thm:p10-lawvere-sheaves-ch}} & \thepage & Source CH forcing stages use q greater than or equal to p and cofinal locality. \\
\phantomsection\label{entry:symbol-registry-04539}\(\operatorname{Epi}(X,Y)\) & \texttt{\detokenize{\operatorname{Epi}}} & object of internally parametrized epimorphisms from X to Y & p10/n-d-023-source-01/ch & n/a & \hyperref[thm:p10-lawvere-sheaves-ch]{\ref*{thm:p10-lawvere-sheaves-ch}} & \thepage & Subobject of Y\^{}X defined by full-image predicate; not just its global sections. \\
\phantomsection\label{entry:symbol-registry-04540}\(<\) & \texttt{\detokenize{<}} & source cardinal comparison: mono from smaller object to larger exists but no epi in that same direction & p10/n-d-023-source-01/ch & relation & \hyperref[thm:p10-lawvere-sheaves-ch]{\ref*{thm:p10-lawvere-sheaves-ch}} & \thepage & Exact source use in 2\^{}omega<I and the sheaf inequality; ground and new power objects differ. \\
\phantomsection\label{entry:symbol-registry-04541}\(P\) & \texttt{\detokenize{P}} & membership subobject of X times multiplicative monoid A & p10/n-d-023-source-01/spectrum & n/a & \hyperref[def:p10-lawvere-sheaves-prime]{\ref*{def:p10-lawvere-sheaves-prime}} & \thepage & Source reuse of glyph; separate from CH poset P. \\
\phantomsection\label{entry:symbol-registry-04542}\(F\) & \texttt{\detokenize{F}} & order/incidence relation object on basis indices A & p10/n-d-023-source-01/spectrum & n/a & \hyperref[def:p10-lawvere-sheaves-prime]{\ref*{def:p10-lawvere-sheaves-prime}} & \thepage & Gives a multiplicative subobject of the constant monoid presheaf. \\
\phantomsection\label{entry:symbol-registry-04543}\([f\in x]\) & \texttt{\detokenize{[f\in x]}} & truth-value predicate for membership in source prime subobject x & p10/n-d-023-source-01/spectrum & n/a & \hyperref[def:p10-lawvere-sheaves-prime]{\ref*{def:p10-lawvere-sheaves-prime}} & \thepage & Prime is multiplicative subobject not an ideal; join may hold only locally. \\
\phantomsection\label{entry:symbol-registry-04544}\(\widetilde A\) & \texttt{\detokenize{\widetilde A}} & reflected fraction ring in intermediate spectrum sheaf category & p10/n-d-023-source-01/spectrum & n/a & \hyperref[thm:p10-lawvere-sheaves-spectrum]{\ref*{thm:p10-lawvere-sheaves-spectrum}} & \thepage & Locality and universal property announced with full proof still owed. \\
\phantomsection\label{entry:symbol-registry-04545}\(\operatorname{Spec}(A)\) & \texttt{\detokenize{\operatorname{Spec}}} & global spectrum of ringed topos (E,A) & p10/n-d-023-source-01/spectrum & n/a & \hyperref[thm:p10-lawvere-sheaves-spectrum]{\ref*{thm:p10-lawvere-sheaves-spectrum}} & \thepage & Source intrinsic construction; not an unqualified prime-ideal set. \\
\phantomsection\label{entry:symbol-registry-04546}\(D_0,D_1,D\) & \texttt{\detokenize{D}} & successive parameter objects for supported single-valued total fiber graphs & p10/n-d-023-source-01/functional-graphs & n/a & \hyperref[def:p10-lawvere-sheaves-functional-graphs]{\ref*{def:p10-lawvere-sheaves-functional-graphs}} & \thepage & Constructed by power objects and finite equalizers; unrelated to the later internal category D. \\
\phantomsection\label{entry:symbol-registry-04547}\(\delta_W\) & \texttt{\detokenize{\delta}} & classifier of the diagonal of W & p10/n-d-023-source-01/functional-graphs & n/a & \hyperref[def:p10-lawvere-sheaves-functional-graphs]{\ref*{def:p10-lawvere-sheaves-functional-graphs}} & \thepage & Internal equality predicate; not a diagonal arrow. \\
\phantomsection\label{entry:symbol-registry-04548}\(\Lambda_f\) & \texttt{\detokenize{\Lambda}} & dependent transpose of a map over X to a map into Pi\_f over Y & p10/n-d-023-source-01/functional-graphs & n/a & \hyperref[thm:p10-lawvere-sheaves-dependent-product]{\ref*{thm:p10-lawvere-sheaves-dependent-product}} & \thepage & Built by the classified total graph; both inverse identities are proved. \\
\phantomsection\label{entry:symbol-registry-04549}\(\varepsilon_p\) & \texttt{\detokenize{\varepsilon}} & dependent-product counit evaluated on D fiber-product X & p10/n-d-023-source-01/functional-graphs & n/a & \hyperref[thm:p10-lawvere-sheaves-dependent-product]{\ref*{thm:p10-lawvere-sheaves-dependent-product}} & \thepage & Evaluation is extracted from equality of actual graph-domain monos; no chosen epi section. \\
\phantomsection\label{entry:symbol-registry-04550}\(\Rightarrow\) & \texttt{\detokenize{\Rightarrow}} & internal Heyting implication classifying the meet-order equalizer & p10/n-d-023-source-01/predicate-adjoints & n/a & \hyperref[thm:p10-lawvere-sheaves-predicate-adjoints]{\ref*{thm:p10-lawvere-sheaves-predicate-adjoints}} & \thepage & Residuates subobject intersection; not Boolean material implication. \\
\phantomsection\label{entry:symbol-registry-04551}\(Q\) & \texttt{\detokenize{Q}} & effective quotient obtained as image of an equivalence-relation row map & p10/n-d-023-source-01/effective-quotients & n/a & \hyperref[thm:p10-lawvere-sheaves-effective-quotients]{\ref*{thm:p10-lawvere-sheaves-effective-quotients}} & \thepage & Its epi has exactly the given relation as kernel pair; no quotient exactness assumed. \\
\phantomsection\label{entry:symbol-registry-04552}\(e\cdot a\) & \texttt{\detokenize{\cdot}} & restriction of an internal presheaf along a:s(a) to t(a) & p10/n-d-023-source-01/internal-actions & n/a & \hyperref[def:p10-lawvere-sheaves-internal-actions]{\ref*{def:p10-lawvere-sheaves-internal-actions}} & \thepage & e lies over t(a); result lies over s(a); editorial generalized-element notation. \\
\phantomsection\label{entry:symbol-registry-04553}\(U_C\) & \texttt{\detokenize{U_C}} & forgetful functor from internal C-presheaves to E/C\_0 & p10/n-d-023-source-01/internal-actions & n/a & \hyperref[def:p10-lawvere-sheaves-internal-actions]{\ref*{def:p10-lawvere-sheaves-internal-actions}} & \thepage & Creates finite limits and finite colimits; monomorphisms detected through its free adjoint. \\
\phantomsection\label{entry:symbol-registry-04554}\(L_C\) & \texttt{\detokenize{L_C}} & free action left adjoint of U\_C & p10/n-d-023-source-01/internal-actions & n/a & \hyperref[thm:p10-lawvere-sheaves-internal-presheaf-topos]{\ref*{thm:p10-lawvere-sheaves-internal-presheaf-topos}} & \thepage & Underlying Sigma\_s t-star B with precomposition. \\
\phantomsection\label{entry:symbol-registry-04555}\(R_C\) & \texttt{\detokenize{R_C}} & cofree action right adjoint of U\_C & p10/n-d-023-source-01/internal-actions & n/a & \hyperref[thm:p10-lawvere-sheaves-internal-presheaf-topos]{\ref*{thm:p10-lawvere-sheaves-internal-presheaf-topos}} & \thepage & Underlying Pi\_t s-star B; dependent assignments on incoming arrows. \\
\phantomsection\label{entry:symbol-registry-04556}\(\Sigma_s\) & \texttt{\detokenize{\Sigma}} & composition of a family structure map with s & p10/n-d-023-source-01/internal-actions & n/a & \hyperref[thm:p10-lawvere-sheaves-internal-presheaf-topos]{\ref*{thm:p10-lawvere-sheaves-internal-presheaf-topos}} & \thepage & Slice left adjoint to pullback; not predicate existential quantification. \\
\phantomsection\label{entry:symbol-registry-04557}\(T_C\) & \texttt{\detokenize{T_C}} & classifier of the internal C-presheaf topos & p10/n-d-023-source-01/internal-actions & n/a & \hyperref[thm:p10-lawvere-sheaves-internal-presheaf-topos]{\ref*{thm:p10-lawvere-sheaves-internal-presheaf-topos}} & \thepage & Internal incoming-arrow sieve predicates; distinct from base classifier T. \\
\phantomsection\label{entry:symbol-registry-04558}\(F_0\) & \texttt{\detokenize{F_0}} & object component C\_0 to D\_0 of an internal functor & p10/n-d-023-source-01/internal-kan-restriction & n/a & \hyperref[thm:p10-lawvere-sheaves-internal-kan]{\ref*{thm:p10-lawvere-sheaves-internal-kan}} & \thepage & Typed separately from F\_0 A denoting the left extension of a presheaf. \\
\phantomsection\label{entry:symbol-registry-04559}\(F_1\) & \texttt{\detokenize{F_1}} & arrow component C\_1 to D\_1 of an internal functor & p10/n-d-023-source-01/internal-kan-restriction & n/a & \hyperref[thm:p10-lawvere-sheaves-internal-kan]{\ref*{thm:p10-lawvere-sheaves-internal-kan}} & \thepage & Respects source target identities and composition. \\
\phantomsection\label{entry:symbol-registry-04560}\(F^*\) & \texttt{\detokenize{{}^*}} & restriction of internal D-presheaves to C-presheaves & p10/n-d-023-source-01/internal-kan & n/a & \hyperref[thm:p10-lawvere-sheaves-internal-kan]{\ref*{thm:p10-lawvere-sheaves-internal-kan}} & \thepage & Underlying family pullback along object component; action uses arrow component. \\
\phantomsection\label{entry:symbol-registry-04561}\(F_0 A\) & \texttt{\detokenize{F_0}} & left extension of a C-presheaf along the internal functor F & p10/n-d-023-source-01/internal-kan-left & n/a & \hyperref[thm:p10-lawvere-sheaves-internal-kan]{\ref*{thm:p10-lawvere-sheaves-internal-kan}} & \thepage & Finite parallel-pair coequalizer; presheaf argument distinguishes it from object component F\_0. \\
\phantomsection\label{entry:symbol-registry-04562}\(F_* A\) & \texttt{\detokenize{F_*}} & right extension of a C-presheaf along the internal functor F & p10/n-d-023-source-01/internal-kan-right & n/a & \hyperref[thm:p10-lawvere-sheaves-internal-kan]{\ref*{thm:p10-lawvere-sheaves-internal-kan}} & \thepage & Compatible dependent assignments indexed by b:F(c) to d. \\
\phantomsection\label{entry:symbol-registry-04563}\(B_0,B_1\) & \texttt{\detokenize{B}} & generator and relation families for the left extension & p10/n-d-023-source-01/internal-kan-left & n/a & \hyperref[thm:p10-lawvere-sheaves-internal-kan]{\ref*{thm:p10-lawvere-sheaves-internal-kan}} & \thepage & Both are constructed D-presheaves; their coequalizer uses the proved finite colimit package. \\
\phantomsection\label{entry:symbol-registry-04564}\(R_X\) & \texttt{\detokenize{R_X}} & closure of the diagonal of X as an internal equivalence relation & p10/n-d-023-source-01/associated-sheaf & n/a & \hyperref[def:p10-lawvere-sheaves-associated-construction]{\ref*{def:p10-lawvere-sheaves-associated-construction}} & \thepage & Classifier j delta\_X; exact kernel pair of the closed-singleton map. \\
\phantomsection\label{entry:symbol-registry-04565}\(\kappa_X\) & \texttt{\detokenize{\kappa}} & closed-singleton map X to T\_j\^{}X & p10/n-d-023-source-01/associated-sheaf & n/a & \hyperref[def:p10-lawvere-sheaves-associated-construction]{\ref*{def:p10-lawvere-sheaves-associated-construction}} & \thepage & Editorial companion to the source singleton map; distinct from source support sigma\_X. \\
\phantomsection\label{entry:symbol-registry-04566}\(S_X\) & \texttt{\detokenize{S_X}} & separated image of the closed-singleton map & p10/n-d-023-source-01/associated-sheaf & n/a & \hyperref[def:p10-lawvere-sheaves-associated-construction]{\ref*{def:p10-lawvere-sheaves-associated-construction}} & \thepage & Subobject of the sheaf T\_j\^{}X; separated quotient universal property proved. \\
\phantomsection\label{entry:symbol-registry-04567}\(q_X\) & \texttt{\detokenize{q_X}} & effective quotient of the closed diagonal onto S\_X & p10/n-d-023-source-01/associated-sheaf & n/a & \hyperref[def:p10-lawvere-sheaves-associated-construction]{\ref*{def:p10-lawvere-sheaves-associated-construction}} & \thepage & Constructed image epi; its kernel pair is R\_X. \\
\phantomsection\label{entry:symbol-registry-04568}\(a_j\) & \texttt{\detokenize{a_j}} & associated-sheaf reflector built by closing S\_X in T\_j\^{}X & p10/n-d-023-source-01/associated-sheaf & n/a & \hyperref[def:p10-lawvere-sheaves-associated-construction]{\ref*{def:p10-lawvere-sheaves-associated-construction}} & \thepage & Universal property functor laws arbitrary-pullback left exactness and classifier package all proved. \\
\phantomsection\label{entry:symbol-registry-04569}\(\eta_X\) & \texttt{\detokenize{\eta}} & unit X to its constructed associated sheaf a\_jX & p10/n-d-023-source-01/associated-sheaf & n/a & \hyperref[def:p10-lawvere-sheaves-associated-construction]{\ref*{def:p10-lawvere-sheaves-associated-construction}} & \thepage & Dense image and dense relative diagonal; no assumed local-isomorphism property. \\
\phantomsection\label{entry:symbol-registry-04570}\(\mathcal E_j\) & \texttt{\detokenize{\mathcal E_j}} & full topos of j-sheaves & p10/n-d-023-source-01/sheaf-topos & n/a & \hyperref[thm:p10-lawvere-sheaves-sheaf-topos]{\ref*{thm:p10-lawvere-sheaves-sheaf-topos}} & \thepage & Finite limits and exponentials ambient; finite colimits reflected; classifier T\_j. \\
\phantomsection\label{entry:symbol-registry-04571}\(i_j\) & \texttt{\detokenize{i_j}} & full-and-faithful inclusion of j-sheaves into E & p10/n-d-023-source-01/sheaf-topos & n/a & \hyperref[thm:p10-lawvere-sheaves-sheaf-topos]{\ref*{thm:p10-lawvere-sheaves-sheaf-topos}} & \thepage & Right adjoint of the proved left-exact reflector. \\
\phantomsection\label{entry:symbol-registry-04572}\(\alpha_X\) & \texttt{\detokenize{\alpha}} & unit X to R L X of the given geometric-morphism adjunction & p10/n-d-023-source-01/geometric-factorization & n/a & \hyperref[thm:p10-lawvere-sheaves-factorization]{\ref*{thm:p10-lawvere-sheaves-factorization}} & \thepage & Not the separately constructed sheafification unit eta\_X. \\
\phantomsection\label{entry:symbol-registry-04573}\(L_X\) & \texttt{\detokenize{L_X}} & inverse-image operation on subobjects induced by the geometric inverse image L & p10/n-d-023-source-01/geometric-factorization & n/a & \hyperref[thm:p10-lawvere-sheaves-factorization]{\ref*{thm:p10-lawvere-sheaves-factorization}} & \thepage & Sub(X) to Sub(LX); preserves finite meets. \\
\phantomsection\label{entry:symbol-registry-04574}\(G_X\) & \texttt{\detokenize{G_X}} & right order adjoint on subobjects using pullback of R along alpha\_X & p10/n-d-023-source-01/geometric-factorization & n/a & \hyperref[thm:p10-lawvere-sheaves-factorization]{\ref*{thm:p10-lawvere-sheaves-factorization}} & \thepage & Actual adjunction L\_X left adjoint G\_X is proved by transpose and pullback. \\
\phantomsection\label{entry:symbol-registry-04575}\(c_X\) & \texttt{\detokenize{c_X}} & universal closure G\_X L\_X induced by the geometric morphism & p10/n-d-023-source-01/geometric-factorization & n/a & \hyperref[thm:p10-lawvere-sheaves-factorization]{\ref*{thm:p10-lawvere-sheaves-factorization}} & \thepage & Idempotence follows from the order adjunction; pullback stability from unit naturality and finite-limit preservation. \\
\phantomsection\label{entry:symbol-registry-04576}\(j_u\) & \texttt{\detokenize{j_u}} & topology classifying c\_T(true) for the geometric morphism u & p10/n-d-023-source-01/geometric-factorization & n/a & \hyperref[thm:p10-lawvere-sheaves-factorization]{\ref*{thm:p10-lawvere-sheaves-factorization}} & \thepage & Gives the conservative and full-sheaf-embedding factors; full-and-faithful morphisms recovered up to equivalence. \\
\phantomsection\label{entry:symbol-registry-04577}\(\neg\) & \texttt{\detokenize{\neg}} & internal Heyting predicate negation & p10/n-d-023-source-01/double-negation & n/a & \hyperref[thm:p10-lawvere-sheaves-double-negation]{\ref*{thm:p10-lawvere-sheaves-double-negation}} & \thepage & Residuates to bottom; not assumed Boolean in the ambient topos. \\
\phantomsection\label{entry:symbol-registry-04578}\(U_a\) & \texttt{\detokenize{U_a}} & basis member classified by the pairing at index a & p10/n-d-023-source-01/godement & n/a & \hyperref[def:p10-lawvere-sheaves-godement-basis]{\ref*{def:p10-lawvere-sheaves-godement-basis}} & \thepage & Subobject of X; generalized indices handled at every parameter base. \\
\phantomsection\label{entry:symbol-registry-04579}\(P\) & \texttt{\detokenize{P}} & basis incidence subobject of X times A & p10/n-d-023-source-01/godement & n/a & \hyperref[def:p10-lawvere-sheaves-godement-basis]{\ref*{def:p10-lawvere-sheaves-godement-basis}} & \thepage & Cover/refinement criteria are actual epi projections of internal objects. \\
\phantomsection\label{entry:symbol-registry-04580}\(A\) & \texttt{\detokenize{A}} & thin internal category of basis indices ordered by subobject inclusion & p10/n-d-023-source-01/godement & n/a & \hyperref[def:p10-lawvere-sheaves-godement-basis]{\ref*{def:p10-lawvere-sheaves-godement-basis}} & \thepage & Preorder if duplicate labels occur; source calls it a poset. \\
\phantomsection\label{entry:symbol-registry-04581}\(\leq\) & \texttt{\detokenize{\leq}} & basis inclusion order U\_a contained in U\_b & p10/n-d-023-source-01/godement & relation & \hyperref[def:p10-lawvere-sheaves-godement-basis]{\ref*{def:p10-lawvere-sheaves-godement-basis}} & \thepage & Defined by the internal universally quantified predicate; not a chosen global refinement. \\
\phantomsection\label{entry:symbol-registry-04582}\(L\) & \texttt{\detokenize{L}} & finite-coequalizer Godement realization functor & p10/n-d-023-source-01/godement & n/a & \hyperref[thm:p10-lawvere-sheaves-godement]{\ref*{thm:p10-lawvere-sheaves-godement}} & \thepage & Presheaves on A to E/X; source lim notation. \\
\phantomsection\label{entry:symbol-registry-04583}\(\mathsf{lim}\) & \texttt{\detokenize{\mathsf{lim}}} & source name for the Godement realization left adjoint & p10/n-d-023-source-01/godement & n/a & \hyperref[thm:p10-lawvere-sheaves-godement]{\ref*{thm:p10-lawvere-sheaves-godement}} & \thepage & Not an assumed external infinite limit; reconstructed by internal generators and relations. \\
\phantomsection\label{entry:symbol-registry-04584}\(\Gamma\) & \texttt{\detokenize{\Gamma}} & dependent section right adjoint of Godement realization & p10/n-d-023-source-01/godement & n/a & \hyperref[thm:p10-lawvere-sheaves-godement]{\ref*{thm:p10-lawvere-sheaves-godement}} & \thepage & E/X to internal A-presheaves; restriction by basis inclusions. \\
\phantomsection\label{entry:symbol-registry-04585}\(H_0,H_1\) & \texttt{\detokenize{H}} & representative and relation families for Godement germs & p10/n-d-023-source-01/godement & n/a & \hyperref[thm:p10-lawvere-sheaves-godement]{\ref*{thm:p10-lawvere-sheaves-godement}} & \thepage & Finite pullback objects over X; their finite coequalizer defines L. \\
\phantomsection\label{entry:symbol-registry-04586}\(G\) & \texttt{\detokenize{G}} & left-exact Godement cotriple L Gamma on E/X & p10/n-d-023-source-01/godement-cotriple & n/a & \hyperref[thm:p10-lawvere-sheaves-godement]{\ref*{thm:p10-lawvere-sheaves-godement}} & \thepage & Coalgebra comparison proved through finite equalizers and the established conservative image factor. \\
\phantomsection\label{entry:symbol-registry-04587}\(\zeta\) & \texttt{\detokenize{\zeta}} & unit of the Godement realization-section adjunction & p10/n-d-023-source-01/godement-cotriple & n/a & \hyperref[thm:p10-lawvere-sheaves-godement]{\ref*{thm:p10-lawvere-sheaves-godement}} & \thepage & Distinct by type from associated-sheaf unit eta. \\
\phantomsection\label{entry:symbol-registry-04588}\(\theta\) & \texttt{\detokenize{\theta}} & coalgebra structure map M to L Gamma M & p10/n-d-023-source-01/godement-cotriple & n/a & \hyperref[thm:p10-lawvere-sheaves-godement]{\ref*{thm:p10-lawvere-sheaves-godement}} & \thepage & Counit and coassociativity equations are explicitly declared. \\
\phantomsection\label{entry:symbol-registry-04589}\(\delta_M\) & \texttt{\detokenize{\delta}} & comultiplication of Godement cotriple at M & p10/n-d-023-source-01/godement-cotriple & n/a & \hyperref[thm:p10-lawvere-sheaves-godement]{\ref*{thm:p10-lawvere-sheaves-godement}} & \thepage & L of the adjunction unit; different typed role from equality classifier delta\_W. \\
\phantomsection\label{entry:symbol-registry-04642}\(\theta_h\) & \texttt{\detokenize{\theta}} & u\_*/h induced global Heyting truth map & p10/lawvere-localization & n/a & \hyperref[def:p10-lawvere-localization-filter]{\ref*{def:p10-lawvere-localization-filter}} & \thepage & Not asserted to classify u\_*(U). \\
\phantomsection\label{entry:symbol-registry-04643}\(\mathcal F_h\) & \texttt{\detokenize{\mathcal F}} & full arbitrary admitted subterminal filter & p10/lawvere-localization & n/a & \hyperref[def:p10-lawvere-localization-filter]{\ref*{def:p10-lawvere-localization-filter}} & \thepage & Proper for nondegenerate E\_0; improper boundary explicitly proved. \\
\phantomsection\label{entry:symbol-registry-04644}\(\mathcal W_h\) & \texttt{\detokenize{\mathcal W}} & exact source selected monomorphism class & p10/lawvere-localization & n/a & \hyperref[def:p10-lawvere-localization-filter]{\ref*{def:p10-lawvere-localization-filter}} & \thepage & Uniform cofinal denominator equivalence proved. \\
\phantomsection\label{entry:symbol-registry-04645}\(\gamma_m\) & \texttt{\detokenize{\gamma}} & universal truth subterminal of mono m & p10/lawvere-localization & n/a & \hyperref[def:p10-lawvere-localization-filter]{\ref*{def:p10-lawvere-localization-filter}} & \thepage & Same predicate as source sigma\_1 Pi(m). \\
\phantomsection\label{entry:symbol-registry-04646}\(i_{U,X}\) & \texttt{\detokenize{i}} & uniform restriction U times X to X & p10/lawvere-localization & n/a & \hyperref[thm:p10-lawvere-localization-denominators]{\ref*{thm:p10-lawvere-localization-denominators}} & \thepage & Source denominator; admitted U need not be finite or principal. \\
\phantomsection\label{entry:symbol-registry-04647}\([U,f]\) & \texttt{\detokenize{[U,}\allowbreak\detokenize{f]}} & fraction germ with numerator f:U times X to Y & p10/lawvere-localization & n/a & \hyperref[def:p10-lawvere-localization-fractions]{\ref*{def:p10-lawvere-localization-fractions}} & \thepage & Equality is further admitted restriction. \\
\phantomsection\label{entry:symbol-registry-04648}\(\sim\) & \texttt{\detokenize{\sim}} & equivalence of fraction representatives & p10/lawvere-localization & relation & \hyperref[def:p10-lawvere-localization-fractions]{\ref*{def:p10-lawvere-localization-fractions}} & \thepage & Actual equality relation not elementary equivalence. \\
\phantomsection\label{entry:symbol-registry-04649}\(f|W\) & \texttt{\detokenize{|}} & restriction of numerator to W times domain & p10/lawvere-localization & n/a & \hyperref[def:p10-lawvere-localization-fractions]{\ref*{def:p10-lawvere-localization-fractions}} & \thepage & W below its original denominator. \\
\phantomsection\label{entry:symbol-registry-04650}\(Q\) & \texttt{\detokenize{Q}} & canonical logical localization functor & p10/lawvere-localization & n/a & \hyperref[def:p10-lawvere-localization-fractions]{\ref*{def:p10-lawvere-localization-fractions}} & \thepage & Not a claimed geometric or closed-base descent of u\_*. \\
\phantomsection\label{entry:symbol-registry-04651}\(\mathcal E/\mathcal F_h\) & \texttt{\detokenize{\mathcal E/}\allowbreak\detokenize{\mathcal F_h}} & ordinary fraction category & p10/lawvere-localization & n/a & \hyperref[def:p10-lawvere-localization-fractions]{\ref*{def:p10-lawvere-localization-fractions}} & \thepage & Exact source localization by its universal property; not an object slice. \\
\phantomsection\label{entry:symbol-registry-04652}\(\bar H\) & \texttt{\detokenize{\bar H}} & unique universal extension of denominator-inverting functor & p10/lawvere-localization & n/a & \hyperref[thm:p10-lawvere-localization-fractions]{\ref*{thm:p10-lawvere-localization-fractions}} & \thepage & Acts by numerator followed by inverse denominator; transformations extend. \\
\phantomsection\label{entry:symbol-registry-04796}\(\mathcal S\) & \texttt{\detokenize{\mathcal S}} & source abstract-set base & p10/lawvere-bvm & object & \hyperref[thm:p10-lawvere-bvm-base]{\ref*{thm:p10-lawvere-bvm-base}} & \thepage & Not an ordinary ZFC universe. \\
\phantomsection\label{entry:symbol-registry-04797}\(\mathcal P_+I\) & \texttt{\detokenize{\mathcal P}} & nonempty predicates on a base object & p10/lawvere-bvm & object & \hyperref[thm:p10-lawvere-bvm-base]{\ref*{thm:p10-lawvere-bvm-base}} & \thepage & Membership epi is split. \\
\phantomsection\label{entry:symbol-registry-04798}\(B_{\mathcal E}\) & \texttt{\detokenize{B}} & complete Boolean algebra p\_*T & p10/lawvere-bvm & object & \hyperref[thm:p10-lawvere-bvm-presentation]{\ref*{thm:p10-lawvere-bvm-presentation}} & \thepage & Completeness internally S-indexed. \\
\phantomsection\label{entry:symbol-registry-04799}\(\bigvee\) & \texttt{\detokenize{\bigvee}} & internally indexed least upper bound & p10/lawvere-bvm & meta & \hyperref[def:p10-lawvere-bvm-matrices]{\ref*{def:p10-lawvere-bvm-matrices}} & \thepage & Actual power-object operation. \\
\phantomsection\label{entry:symbol-registry-04800}\(\operatorname{Mat}(B)\) & \texttt{\detokenize{\operatorname}} & Boolean-matrix presentation category & p10/lawvere-bvm & object & \hyperref[def:p10-lawvere-bvm-matrices]{\ref*{def:p10-lawvere-bvm-matrices}} & \thepage & Maps saturated functional relations. \\
\phantomsection\label{entry:symbol-registry-04801}\((I,b,e)\) & \texttt{\detokenize{(I,}\allowbreak\detokenize{b,}\allowbreak\detokenize{e)}} & index weights and equality & p10/lawvere-bvm & object & \hyperref[def:p10-lawvere-bvm-matrices]{\ref*{def:p10-lawvere-bvm-matrices}} & \thepage & Zero weights and empty indices allowed. \\
\phantomsection\label{entry:symbol-registry-04802}\(r_{ij}\) & \texttt{\detokenize{r}} & truth matrix of a map graph & p10/lawvere-bvm & object & \hyperref[def:p10-lawvere-bvm-matrices]{\ref*{def:p10-lawvere-bvm-matrices}} & \thepage & Single-valued modulo target equality. \\
\phantomsection\label{entry:symbol-registry-04803}\(\partial X\) & \texttt{\detokenize{\partial}} & subsingleton predicates on X & p10/lawvere-bvm & object & \hyperref[thm:p10-lawvere-bvm-presentation]{\ref*{thm:p10-lawvere-bvm-presentation}} & \thepage & Partial elements not global total points. \\
\phantomsection\label{entry:symbol-registry-04804}\(D_X\) & \texttt{\detokenize{D}} & canonical partial-element cover & p10/lawvere-bvm & object & \hyperref[thm:p10-lawvere-bvm-presentation]{\ref*{thm:p10-lawvere-bvm-presentation}} & \thepage & Evaluation epi proved. \\
\phantomsection\label{entry:symbol-registry-04805}\(Q\) & \texttt{\detokenize{Q}} & effective-quotient matrix realization & p10/lawvere-bvm & object & \hyperref[thm:p10-lawvere-bvm-presentation]{\ref*{thm:p10-lawvere-bvm-presentation}} & \thepage & Distinct from localization Q. \\
\phantomsection\label{entry:symbol-registry-04806}\(M\) & \texttt{\detokenize{M}} & canonical partial-element matrix functor & p10/lawvere-bvm & object & \hyperref[thm:p10-lawvere-bvm-presentation]{\ref*{thm:p10-lawvere-bvm-presentation}} & \thepage & Actual inverse comparisons. \\
\phantomsection\label{entry:symbol-registry-04807}\(\operatorname{Sh}_{\mathcal S}(B)\) & \texttt{\detokenize{\operatorname}} & internal join-coverage sheaves & p10/lawvere-bvm & object & \hyperref[thm:p10-lawvere-bvm-realization]{\ref*{thm:p10-lawvere-bvm-realization}} & \thepage & Uses established internal reflector. \\
\phantomsection\label{entry:symbol-registry-04808}\(t_{ij}\) & \texttt{\detokenize{t}} & well-ordered Boolean disjointification & p10/lawvere-bvm & object & \hyperref[thm:p10-lawvere-bvm-realization]{\ref*{thm:p10-lawvere-bvm-realization}} & \thepage & No global sheaf witness selected. \\
\phantomsection\label{entry:symbol-registry-04809}\(\mathsf{BVM}^{\leftarrow}(\mathcal S)\) & \texttt{\detokenize{\mathsf}} & inverse-image-directed bicategory & p10/lawvere-bvm & object & \hyperref[def:p10-lawvere-bvm-cell-regime]{\ref*{def:p10-lawvere-bvm-cell-regime}} & \thepage & Editorial specified cell regime. \\
\phantomsection\label{entry:symbol-registry-04810}\(T_\alpha\) & \texttt{\detokenize{T}} & entrywise homomorphism transport & p10/lawvere-bvm & object & \hyperref[thm:p10-lawvere-bvm-correspondence]{\ref*{thm:p10-lawvere-bvm-correspondence}} & \thepage & Not truth classifier T. \\
\phantomsection\label{entry:symbol-registry-04811}\(R_\alpha\) & \texttt{\detokenize{R}} & sheaf right adjoint by precomposition & p10/lawvere-bvm & object & \hyperref[thm:p10-lawvere-bvm-correspondence]{\ref*{thm:p10-lawvere-bvm-correspondence}} & \thepage & Units counits triangles proved. \\
\phantomsection\label{entry:symbol-registry-04812}\(\alpha_L\) & \texttt{\detokenize{\alpha}} & homomorphism of a geometric inverse image & p10/lawvere-bvm & object & \hyperref[thm:p10-lawvere-bvm-correspondence]{\ref*{thm:p10-lawvere-bvm-correspondence}} & \thepage & Opposite conventional geometric direction. \\
\phantomsection\label{entry:symbol-registry-04902}\(\partial X\) & \texttt{\detokenize{\partial}} & subsingleton predicates for base completion & p10/lawvere-closed-base & object & \hyperref[def:p10-lawvere-closed-base-partials]{\ref*{def:p10-lawvere-closed-base-partials}} & \thepage & Partial members at arbitrary generalized bases. \\
\phantomsection\label{entry:symbol-registry-04903}\(D_X\) & \texttt{\detokenize{D}} & universal monic partial-member domain & p10/lawvere-closed-base & object & \hyperref[def:p10-lawvere-closed-base-partials]{\ref*{def:p10-lawvere-closed-base-partials}} & \thepage & Not the BVM partial-element cover of the same local spelling. \\
\phantomsection\label{entry:symbol-registry-04904}\(d_X\) & \texttt{\detokenize{d}} & truth-valued domain map on partial members & p10/lawvere-closed-base & object & \hyperref[def:p10-lawvere-closed-base-partials]{\ref*{def:p10-lawvere-closed-base-partials}} & \thepage & Classifier before applying u\_* and h. \\
\phantomsection\label{entry:symbol-registry-04905}\(e_X\) & \texttt{\detokenize{e}} & common-member domain truth & p10/lawvere-closed-base & object & \hyperref[def:p10-lawvere-closed-base-partials]{\ref*{def:p10-lawvere-closed-base-partials}} & \thepage & Not full equality of partial predicates. \\
\phantomsection\label{entry:symbol-registry-04906}\(\widehat d_X\) & \texttt{\detokenize{\widehat}} & h of u\_* domain truth & p10/lawvere-closed-base & object & \hyperref[def:p10-lawvere-closed-base-partials]{\ref*{def:p10-lawvere-closed-base-partials}} & \thepage & Not the classifier of u\_*D\_X. \\
\phantomsection\label{entry:symbol-registry-04907}\(\widehat e_X\) & \texttt{\detokenize{\widehat}} & h of u\_* common-member truth & p10/lawvere-closed-base & object & \hyperref[def:p10-lawvere-closed-base-partials]{\ref*{def:p10-lawvere-closed-base-partials}} & \thepage & New base equivalence predicate. \\
\phantomsection\label{entry:symbol-registry-04908}\(A_X\) & \texttt{\detokenize{A}} & h-total partial members in the base & p10/lawvere-closed-base & object & \hyperref[def:p10-lawvere-closed-base-partials]{\ref*{def:p10-lawvere-closed-base-partials}} & \thepage & Full truth fiber of hatted domain. \\
\phantomsection\label{entry:symbol-registry-04909}\(R_X\) & \texttt{\detokenize{R}} & effective equality on h-total partial members & p10/lawvere-closed-base & relation & \hyperref[def:p10-lawvere-closed-base-partials]{\ref*{def:p10-lawvere-closed-base-partials}} & \thepage & Actual equivalence relation proved. \\
\phantomsection\label{entry:symbol-registry-04910}\(H_h\) & \texttt{\detokenize{H}} & new finite-limit base functor & p10/lawvere-closed-base & object & \hyperref[thm:p10-lawvere-closed-base-partials]{\ref*{thm:p10-lawvere-closed-base-partials}} & \thepage & Not a descent of original u\_*. \\
\phantomsection\label{entry:symbol-registry-04911}\(\eta_X\) & \texttt{\detokenize{\eta}} & singleton map from u\_*X to H\_hX & p10/lawvere-closed-base & object & \hyperref[thm:p10-lawvere-closed-base-partials]{\ref*{thm:p10-lawvere-closed-base-partials}} & \thepage & Natural comparison not an assumed adjunction unit. \\
\phantomsection\label{entry:symbol-registry-04912}\(\tau\) & \texttt{\detokenize{\tau}} & truth map H\_hT to T\_0 extending h & p10/lawvere-closed-base & object & \hyperref[thm:p10-lawvere-closed-base-partials]{\ref*{thm:p10-lawvere-closed-base-partials}} & \thepage & Not asserted invertible or logical. \\
\phantomsection\label{entry:symbol-registry-04913}\(\rho_m\) & \texttt{\detokenize{\rho}} & restriction of a partial member to a mono & p10/lawvere-closed-base & object & \hyperref[thm:p10-lawvere-closed-base-partials]{\ref*{thm:p10-lawvere-closed-base-partials}} & \thepage & h-total when exact source denominator is selected. \\
\phantomsection\label{entry:symbol-registry-04914}\(\overline H_h\) & \texttt{\detokenize{\overline}} & unique localization extension of H\_h & p10/lawvere-closed-base & object & \hyperref[thm:p10-lawvere-closed-base-partials]{\ref*{thm:p10-lawvere-closed-base-partials}} & \thepage & Actual fraction formula. \\
\phantomsection\label{entry:symbol-registry-04915}\(\kappa_{X,Y}\) & \texttt{\detokenize{\kappa}} & product-induced closed comparison & p10/lawvere-closed-base & object & \hyperref[thm:p10-lawvere-closed-base-closed]{\ref*{thm:p10-lawvere-closed-base-closed}} & \thepage & No strong or normal assertion. \\
\phantomsection\label{entry:symbol-registry-04916}\(\mathcal C_h(X,Y)\) & \texttt{\detokenize{\mathcal C}} & base hom object H\_h of localized exponential & p10/lawvere-closed-base & object & \hyperref[thm:p10-lawvere-closed-base-closed]{\ref*{thm:p10-lawvere-closed-base-closed}} & \thepage & Global underlying maps need not equal ordinary fraction maps. \\
\phantomsection\label{entry:symbol-registry-04917}\(\mathcal U\) & \texttt{\detokenize{\mathcal U}} & boundedly constructed proper Boolean deciding filter & p10/lawvere-closed-base & object & \hyperref[thm:p10-lawvere-closed-base-collapse]{\ref*{thm:p10-lawvere-closed-base-collapse}} & \thepage & Choice of well-order allowed; no complete-homomorphism assertion. \\
\phantomsection\label{entry:symbol-registry-04918}\(F[b]\) & \texttt{\detokenize{F[b]}} & explicit proper-filter extension by a Boolean value & p10/lawvere-closed-base & object & \hyperref[thm:p10-lawvere-closed-base-collapse]{\ref*{thm:p10-lawvere-closed-base-collapse}} & \thepage & Generated from one added element. \\
\phantomsection\label{entry:symbol-registry-04919}\(F_b\) & \texttt{\detokenize{F}} & filter at a bounded well-ordered history stage & p10/lawvere-closed-base & object & \hyperref[thm:p10-lawvere-closed-base-collapse]{\ref*{thm:p10-lawvere-closed-base-collapse}} & \thepage & History graphs lie in power objects not external ordinal recursion. \\
\phantomsection\label{entry:symbol-registry-04955}\(\operatorname{Fin}(D)\) & \texttt{\detokenize{\operatorname}} & least singleton and binary-union closed finite-subset object & p10/lawvere-ch-chain & meta & \hyperref[def:p10-lawvere-ch-chain-poset]{\ref*{def:p10-lawvere-ch-chain-poset}} & \thepage & Constructed in power objects; finite enumeration equivalence proved in this base. \\
\phantomsection\label{entry:symbol-registry-04956}\([n]\) & \texttt{\detokenize{[n]}} & natural initial segment of length n & p10/lawvere-ch-chain & object & \hyperref[def:p10-lawvere-ch-chain-poset]{\ref*{def:p10-lawvere-ch-chain-poset}} & \thepage & Internal NNO arithmetic not external ordinals. \\
\phantomsection\label{entry:symbol-registry-04957}\(D\) & \texttt{\detokenize{D}} & ground coordinate object omega times I & p10/lawvere-ch-chain & object & \hyperref[def:p10-lawvere-ch-chain-poset]{\ref*{def:p10-lawvere-ch-chain-poset}} & \thepage & Distinct from deciding subobjects D\_x. \\
\phantomsection\label{entry:symbol-registry-04958}\(\succeq\) & \texttt{\detokenize{\succeq}} & extension toward stronger conditions & p10/lawvere-ch-chain & relation & \hyperref[def:p10-lawvere-ch-chain-poset]{\ref*{def:p10-lawvere-ch-chain-poset}} & \thepage & Source q greater than or equal to p; site arrow q to p. \\
\phantomsection\label{entry:symbol-registry-04959}\(\mathcal M\) & \texttt{\detokenize{\mathcal M}} & hereditary chain-union-closed subset property & p10/lawvere-ch-chain/annex-cardinal & object & \hyperref[annex:p10.lawvere-ch-chain.cardinal-tools]{\ref*{annex:p10.lawvere-ch-chain.cardinal-tools}} & \thepage & Bounded greedy history proves maximalization. \\
\phantomsection\label{entry:symbol-registry-04960}\(t_n\) & \texttt{\detokenize{t}} & triangular natural-number block endpoints & p10/lawvere-ch-chain/annex-cardinal & object & \hyperref[annex:p10.lawvere-ch-chain.cardinal-tools]{\ref*{annex:p10.lawvere-ch-chain.cardinal-tools}} & \thepage & Actual recursive pairing bijection omega squared to omega. \\
\phantomsection\label{entry:symbol-registry-04961}\(\mathcal A\) & \texttt{\detokenize{\mathcal A}} & family of distinct finite supports & p10/lawvere-ch-chain/annex-ccc & object & \hyperref[annex:p10.lawvere-ch-chain.ccc]{\ref*{annex:p10.lawvere-ch-chain.ccc}} & \thepage & Internal uncountability and delta-system induction. \\
\phantomsection\label{entry:symbol-registry-04962}\(D_x\) & \texttt{\detokenize{D}} & bounded deciding condition and value pairs & p10/lawvere-ch-chain/annex-no-epi & object & \hyperref[annex:p10.lawvere-ch-chain.no-epi]{\ref*{annex:p10.lawvere-ch-chain.no-epi}} & \thepage & Constructed from internal action restriction evaluation and equality. \\
\phantomsection\label{entry:symbol-registry-04963}\(S_x\) & \texttt{\detokenize{S}} & countable possible output values for a ground input & p10/lawvere-ch-chain/annex-no-epi & object & \hyperref[annex:p10.lawvere-ch-chain.no-epi]{\ref*{annex:p10.lawvere-ch-chain.no-epi}} & \thepage & Image of a maximal deciding antichain. \\
\phantomsection\label{entry:symbol-registry-04964}\(\alpha\) & \texttt{\detokenize{\alpha}} & descended generic sequence injection & p10/lawvere-ch-chain & object & \hyperref[thm:p10-lawvere-ch-chain-generic]{\ref*{thm:p10-lawvere-ch-chain-generic}} & \thepage & Dense partial evaluation reflected then transposed. \\
\phantomsection\label{entry:symbol-registry-04965}\(C\) & \texttt{\detokenize{C}} & ground sequence power object 2 to omega & p10/lawvere-ch-chain & object & \hyperref[thm:p10-lawvere-ch-chain-comparisons]{\ref*{thm:p10-lawvere-ch-chain-comparisons}} & \thepage & Not the newly formed sheaf sequence object. \\
\phantomsection\label{entry:symbol-registry-04966}\(\kappa\) & \texttt{\detokenize{\kappa}} & transported-evaluation ground-to-new sequence comparison & p10/lawvere-ch-chain/annex-comparisons & object & \hyperref[annex:p10.lawvere-ch-chain.comparisons]{\ref*{annex:p10.lawvere-ch-chain.comparisons}} & \thepage & Monicity proved without assuming inverse image logical. \\
\phantomsection\label{entry:symbol-registry-04967}\(r\) & \texttt{\detokenize{r}} & complement retraction from new sequences to u*I & p10/lawvere-ch-chain/annex-comparisons & object & \hyperref[annex:p10.lawvere-ch-chain.comparisons]{\ref*{annex:p10.lawvere-ch-chain.comparisons}} & \thepage & Different local role from a condition extension; every parametrized epi is excluded. \\
\phantomsection\label{entry:symbol-registry-04968}\(\operatorname{Epi}\) & \texttt{\detokenize{\operatorname}} & full internally parametrized epi predicate object & p10/lawvere-ch-chain & meta & \hyperref[thm:p10-lawvere-ch-chain-no-epi]{\ref*{thm:p10-lawvere-ch-chain-no-epi}} & \thepage & Zero object assertion not merely absence of a global map. \\
\phantomsection\label{entry:symbol-registry-05475}\(\sqsubseteq\) & \texttt{\detokenize{\sqsubseteq}} & information order & p10/domain-theory & relation & \hyperref[def:p10-domain-order-convention]{\ref*{def:p10-domain-order-convention}} & \thepage & Not numerical order. \\
\phantomsection\label{entry:symbol-registry-05476}\(\bigsqcup^{\uparrow}A\) & \texttt{\detokenize{\bigsqcup}} & nonempty directed supremum & p10/domain-theory & n/a & \hyperref[def:p10-domain-order-convention]{\ref*{def:p10-domain-order-convention}} & \thepage & Not only countable-chain convergence. Original precedence descriptor 'operator' preserved as semantic metadata; n/a is not an invented binding rank. \\
\phantomsection\label{entry:symbol-registry-05477}\(\bot_D\) & \texttt{\detokenize{\bot}} & least information element & p10/domain-theory & n/a & \hyperref[def:p10-domain-order-convention]{\ref*{def:p10-domain-order-convention}} & \thepage & Lifting adds a new bottom. Original precedence descriptor 'constant' preserved as semantic metadata; n/a is not an invented binding rank. \\
\phantomsection\label{entry:symbol-registry-05478}\(x\ll y\) & \texttt{\detokenize{\ll}} & way-below approximation & p10/domain-theory & relation & \hyperref[def:p10-domain-way-below]{\ref*{def:p10-domain-way-below}} & \thepage & Tests all directed subsets. \\
\phantomsection\label{entry:symbol-registry-05479}\(K(D)\) & \texttt{\detokenize{K}} & compact elements & p10/domain-theory & object & \hyperref[def:p10-domain-way-below]{\ref*{def:p10-domain-way-below}} & \thepage & Self-approximating elements. \\
\phantomsection\label{entry:symbol-registry-05480}\(B_x\) & \texttt{\detokenize{B}} & basis approximants to x & p10/domain-theory & object & \hyperref[def:p10-domain-way-below]{\ref*{def:p10-domain-way-below}} & \thepage & Basis intersected with way-below predecessors. \\
\phantomsection\label{entry:symbol-registry-05481}\(\mathord{\downarrow\!\downarrow}x\) & \texttt{\detokenize{\downarrow}} & way-below predecessors & p10/domain-theory & n/a & \hyperref[def:p10-domain-way-below]{\ref*{def:p10-domain-way-below}} & \thepage & Not ordinary down-set. Original precedence descriptor 'operator' preserved as semantic metadata; n/a is not an invented binding rank. \\
\phantomsection\label{entry:symbol-registry-05482}\(\mathord{\uparrow\!\uparrow}x\) & \texttt{\detokenize{\uparrow}} & points approximated by x & p10/domain-theory & n/a & \hyperref[def:p10-domain-way-below]{\ref*{def:p10-domain-way-below}} & \thepage & Scott-open under continuous-domain hypotheses. Original precedence descriptor 'operator' preserved as semantic metadata; n/a is not an invented binding rank. \\
\phantomsection\label{entry:symbol-registry-05483}\([D\longrightarrow E]\) & \texttt{\detokenize{\longrightarrow}} & continuous function dcpo & p10/domain-theory & object & \hyperref[def:p10-domain-continuous-map]{\ref*{def:p10-domain-continuous-map}} & \thepage & Pointwise information order. \\
\phantomsection\label{entry:symbol-registry-05484}\([D\mathrel{\mathop{\longrightarrow}^{\bot!}}E]\) & \texttt{\detokenize{\longrightarrow}} & strict continuous function dcpo & p10/domain-theory & object & \hyperref[def:p10-domain-continuous-map]{\ref*{def:p10-domain-continuous-map}} & \thepage & Bottom-preserving maps. \\
\phantomsection\label{entry:symbol-registry-05485}\(\mathsf{DCPO}_{\bot!}\) & \texttt{\detokenize{\mathsf{DCPO}}} & strict pointed dcpo category & p10/domain-theory & n/a & \hyperref[def:p10-domain-continuous-map]{\ref*{def:p10-domain-continuous-map}} & \thepage & Distinct from unrestricted and non-strict pointed categories. Original precedence descriptor 'category' preserved as semantic metadata; n/a is not an invented binding rank. \\
\phantomsection\label{entry:symbol-registry-05486}\(\prec\) & \texttt{\detokenize{\prec}} & abstract-basis relation & p10/domain-theory & relation & \hyperref[def:p10-domain-abstract-basis]{\ref*{def:p10-domain-abstract-basis}} & \thepage & Transitive with finite interpolation; not necessarily reflexive. \\
\phantomsection\label{entry:symbol-registry-05487}\(\operatorname{Idl}(B,\prec)\) & \texttt{\detokenize{\operatorname{Idl}}} & round-ideal completion & p10/domain-theory & object & \hyperref[def:p10-domain-abstract-basis]{\ref*{def:p10-domain-abstract-basis}} & \thepage & Nonempty lower directed round ideals. \\
\phantomsection\label{entry:symbol-registry-05488}\(i(b)\) & \texttt{\detokenize{i}} & basis lower-ideal map & p10/domain-theory & n/a & \hyperref[def:p10-domain-abstract-basis]{\ref*{def:p10-domain-abstract-basis}} & \thepage & Need not contain b. Original precedence descriptor 'map' preserved as semantic metadata; n/a is not an invented binding rank. \\
\phantomsection\label{entry:symbol-registry-05489}\(\sigma D\) & \texttt{\detokenize{\sigma}} & Scott topology & p10/domain-theory & object & \hyperref[def:p10-domain-scott-topology]{\ref*{def:p10-domain-scott-topology}} & \thepage & Upper sets inaccessible by directed suprema. \\
\phantomsection\label{entry:symbol-registry-05490}\(\operatorname{fix}(f)\) & \texttt{\detokenize{\operatorname{fix}}} & least fixed-point operator & p10/domain-theory & n/a & \hyperref[thm:p10-domain-least-fixpoint]{\ref*{thm:p10-domain-least-fixpoint}} & \thepage & Its continuity is separately proved. Original precedence descriptor 'operator' preserved as semantic metadata; n/a is not an invented binding rank. \\
\phantomsection\label{entry:symbol-registry-05491}\(\Phi\) & \texttt{\detokenize{\Phi}} & factorial partial-graph functional & p10/domain-theory & n/a & \hyperref[def:p10-domain-recursive-factorial]{\ref*{def:p10-domain-recursive-factorial}} & \thepage & Finite input observation for each output. Original precedence descriptor 'map' preserved as semantic metadata; n/a is not an invented binding rank. \\
\phantomsection\label{entry:symbol-registry-05492}\(I\) & \texttt{\detokenize{I}} & one-point dcpo & p10/domain-theory & object & \hyperref[def:p10-domain-type-constructors]{\ref*{def:p10-domain-type-constructors}} & \thepage & I\_bot is the two-element smash unit. \\
\phantomsection\label{entry:symbol-registry-05493}\(\operatorname{up}(d)\) & \texttt{\detokenize{\operatorname{up}}} & lifting injection & p10/domain-theory & n/a & \hyperref[def:p10-domain-type-constructors]{\ref*{def:p10-domain-type-constructors}} & \thepage & Old and new bottoms remain distinct. Original precedence descriptor 'map' preserved as semantic metadata; n/a is not an invented binding rank. \\
\phantomsection\label{entry:symbol-registry-05494}\(D\oplus E\) & \texttt{\detokenize{\oplus}} & coalesced sum & p10/domain-theory & n/a & \hyperref[def:p10-domain-type-constructors]{\ref*{def:p10-domain-type-constructors}} & \thepage & Strict-map coproduct; not MV addition. Original precedence descriptor 'constructor' preserved as semantic metadata; n/a is not an invented binding rank. \\
\phantomsection\label{entry:symbol-registry-05495}\(D\otimes E\) & \texttt{\detokenize{\otimes}} & smash product & p10/domain-theory & n/a & \hyperref[def:p10-domain-type-constructors]{\ref*{def:p10-domain-type-constructors}} & \thepage & Bottom-axis pairs collapsed. Original precedence descriptor 'constructor' preserved as semantic metadata; n/a is not an invented binding rank. \\
\phantomsection\label{entry:symbol-registry-05496}\(a\mathbin{\searrow}b\) & \texttt{\detokenize{\searrow}} & compact step function & p10/domain-theory & n/a & \hyperref[thm:p10-domain-algebraic-functions]{\ref*{thm:p10-domain-algebraic-functions}} & \thepage & Finite compatible joins give the exact compact basis. Original precedence descriptor 'constructor' preserved as semantic metadata; n/a is not an invented binding rank. \\
\phantomsection\label{entry:symbol-registry-05497}\(e,p\) & \texttt{\detokenize{e,}\allowbreak\detokenize{p}} & embedding-projection pair & p10/domain-theory & n/a & \hyperref[def:p10-domain-expanding-sequence]{\ref*{def:p10-domain-expanding-sequence}} & \thepage & pe identity and ep below identity. Original precedence descriptor 'maps' preserved as semantic metadata; n/a is not an invented binding rank. \\
\phantomsection\label{entry:symbol-registry-05498}\(e_{mn},p_{nm}\) & \texttt{\detokenize{e,}\allowbreak\detokenize{p}} & stage embedding and projection & p10/domain-theory & n/a & \hyperref[def:p10-domain-expanding-sequence]{\ref*{def:p10-domain-expanding-sequence}} & \thepage & e\_mn maps D\_n to D\_m. Original precedence descriptor 'maps' preserved as semantic metadata; n/a is not an invented binding rank. \\
\phantomsection\label{entry:symbol-registry-05499}\(e_n,p_n\) & \texttt{\detokenize{e,}\allowbreak\detokenize{p}} & bilimit embedding and projection & p10/domain-theory & n/a & \hyperref[thm:p10-domain-bilimit]{\ref*{thm:p10-domain-bilimit}} & \thepage & Composite idempotents increase to identity. Original precedence descriptor 'maps' preserved as semantic metadata; n/a is not an invented binding rank. \\
\phantomsection\label{entry:symbol-registry-05500}\(F,\widehat F,G\) & \texttt{\detokenize{F,}\allowbreak\detokenize{G}} & mixed and induced embedding functors & p10/domain-theory & n/a & \hyperref[def:p10-domain-continuous-functor]{\ref*{def:p10-domain-continuous-functor}} & \thepage & Local continuity and bilimit preservation proved, not presumed. Original precedence descriptor 'functors' preserved as semantic metadata; n/a is not an invented binding rank. \\
\phantomsection\label{entry:symbol-registry-05501}\(\operatorname{fold},\operatorname{unfold}\) & \texttt{\detokenize{\operatorname{fold},}\allowbreak\detokenize{\operatorname{unfold}}} & recursive comparison isomorphisms & p10/domain-theory & n/a & \hyperref[thm:p10-domain-recursive-solution]{\ref*{thm:p10-domain-recursive-solution}} & \thepage & Actual typed mediator suprema. Original precedence descriptor 'maps' preserved as semantic metadata; n/a is not an invented binding rank. \\
\phantomsection\label{entry:symbol-registry-05502}\(\epsilon\) & \texttt{\detokenize{\epsilon}} & empty partial word & p10/domain-theory & n/a & \hyperref[thm:p10-domain-recursive-examples]{\ref*{thm:p10-domain-recursive-examples}} & \thepage & Bottom prefix, not termination tag or ordinal. Original precedence descriptor 'constant' preserved as semantic metadata; n/a is not an invented binding rank. \\
\phantomsection\label{entry:symbol-registry-05627}\(\mathcal D,\mathcal D_{\bot!}\) & \texttt{\detokenize{\mathcal D}} & unrestricted and strict domain categories & p10/domain-invariance & n/a & \hyperref[def:p10-domain-invariance-category]{\ref*{def:p10-domain-invariance-category}} & \thepage & Contains I and actual expanding-sequence bilimits; strict-only functors not silently extended. Original precedence descriptor 'category' preserved as semantic metadata; n/a is not an invented binding rank. \\
\phantomsection\label{entry:symbol-registry-05628}\(i_D,\operatorname{fold},\operatorname{unfold}\) & \texttt{\detokenize{i_D,}\allowbreak\detokenize{\operatorname{fold},}\allowbreak\detokenize{\operatorname{unfold}}} & canonical solution structure comparisons & p10/domain-invariance & n/a & \hyperref[def:p10-domain-invariance-algebra]{\ref*{def:p10-domain-invariance-algebra}} & \thepage & Canonical means the one-point seed I. Original precedence descriptor 'maps' preserved as semantic metadata; n/a is not an invented binding rank. \\
\phantomsection\label{entry:symbol-registry-05629}\(q_n\) & \texttt{\detokenize{q}} & canonical approximating idempotents & p10/domain-invariance & n/a & \hyperref[def:p10-domain-invariance-algebra]{\ref*{def:p10-domain-invariance-algebra}} & \thepage & e\_n p\_n with directed supremum identity. Original precedence descriptor 'maps' preserved as semantic metadata; n/a is not an invented binding rank. \\
\phantomsection\label{entry:symbol-registry-05630}\(\phi_{C,A}\) & \texttt{\detokenize{\phi}} & covariant homomorphism functional & p10/domain-invariance & n/a & \hyperref[thm:p10-domain-invariance-least-hom]{\ref*{thm:p10-domain-invariance-least-hom}} & \thepage & Typed hom dcpo action a F(u) i inverse. Original precedence descriptor 'map' preserved as semantic metadata; n/a is not an invented binding rank. \\
\phantomsection\label{entry:symbol-registry-05631}\(h_{C,A}\) & \texttt{\detokenize{h}} & least covariant algebra comparison & p10/domain-invariance & n/a & \hyperref[thm:p10-domain-invariance-least-hom]{\ref*{thm:p10-domain-invariance-least-hom}} & \thepage & Composition is independently derived. Original precedence descriptor 'map' preserved as semantic metadata; n/a is not an invented binding rank. \\
\phantomsection\label{entry:symbol-registry-05632}\(\mu_A\) & \texttt{\detokenize{\mu}} & least solution endomorphism & p10/domain-invariance & n/a & \hyperref[thm:p10-domain-invariance-least-hom]{\ref*{thm:p10-domain-invariance-least-hom}} & \thepage & Strict idempotent below identity; not assumed identity. Original precedence descriptor 'map' preserved as semantic metadata; n/a is not an invented binding rank. \\
\phantomsection\label{entry:symbol-registry-05633}\(\theta_D\) & \texttt{\detokenize{\theta}} & canonical endomap functional & p10/domain-invariance & n/a & \hyperref[thm:p10-domain-invariance-covariant]{\ref*{thm:p10-domain-invariance-covariant}} & \thepage & Its iterates are actual q\_n; mixed case uses both slots. Original precedence descriptor 'map' preserved as semantic metadata; n/a is not an invented binding rank. \\
\phantomsection\label{entry:symbol-registry-05634}\(\Phi_{B,A}\) & \texttt{\detokenize{\Phi}} & paired mixed comparison functional & p10/domain-invariance & n/a & \hyperref[thm:p10-domain-invariance-least-pairs]{\ref*{thm:p10-domain-invariance-least-pairs}} & \thepage & Forward and backward maps are differently typed. Original precedence descriptor 'map' preserved as semantic metadata; n/a is not an invented binding rank. \\
\phantomsection\label{entry:symbol-registry-05635}\(\Psi\) & \texttt{\detokenize{\Psi}} & free mixed comparison functional & p10/domain-invariance & n/a & \hyperref[thm:p10-domain-invariance-free-mixed]{\ref*{thm:p10-domain-invariance-free-mixed}} & \thepage & Product of strict hom dcpos. Original precedence descriptor 'map' preserved as semantic metadata; n/a is not an invented binding rank. \\
\phantomsection\label{entry:symbol-registry-05636}\(M(A),M(u)\) & \texttt{\detokenize{M}} & parameter fixed-point functor & p10/domain-invariance & n/a & \hyperref[thm:p10-domain-invariance-parameters]{\ref*{thm:p10-domain-invariance-parameters}} & \thepage & Actual object and stagewise map construction. Original precedence descriptor 'functor' preserved as semantic metadata; n/a is not an invented binding rank. \\
\phantomsection\label{entry:symbol-registry-05637}\(X_n(A),u_n\) & \texttt{\detokenize{X,}\allowbreak\detokenize{u}} & parameter stage objects and maps & p10/domain-invariance & n/a & \hyperref[thm:p10-domain-invariance-parameters]{\ref*{thm:p10-domain-invariance-parameters}} & \thepage & u\_0 identity on I; u\_(n+1)=F(u,u\_n). Original precedence descriptor 'objects and maps' preserved as semantic metadata; n/a is not an invented binding rank. \\
\phantomsection\label{entry:symbol-registry-05638}\(\alpha_A\) & \texttt{\detokenize{\alpha}} & natural parameter fold & p10/domain-invariance & n/a & \hyperref[thm:p10-domain-invariance-parameters]{\ref*{thm:p10-domain-invariance-parameters}} & \thepage & Naturality proved without adjoints to arbitrary M(u). Original precedence descriptor 'map' preserved as semantic metadata; n/a is not an invented binding rank. \\
\phantomsection\label{entry:symbol-registry-05639}\(M(P,Q),r_n,s_n\) & \texttt{\detokenize{M,}\allowbreak\detokenize{r,}\allowbreak\detokenize{s}} & mixed parameter objects and stage maps & p10/domain-invariance & n/a & \hyperref[thm:p10-domain-invariance-mixed-parameters]{\ref*{thm:p10-domain-invariance-mixed-parameters}} & \thepage & Four embedding/projection compatibilities derived. Original precedence descriptor 'objects and maps' preserved as semantic metadata; n/a is not an invented binding rank. \\
\phantomsection\label{entry:symbol-registry-05640}\(P,Q,H,K\) & \texttt{\detokenize{P,}\allowbreak\detokenize{Q,}\allowbreak\detokenize{H,}\allowbreak\detokenize{K}} & simultaneous and nested domain functors & p10/domain-invariance & n/a & \hyperref[thm:p10-domain-invariance-bekic]{\ref*{thm:p10-domain-invariance-bekic}} & \thepage & Strict categorical Bekić distinct from value elimination. Original precedence descriptor 'functors' preserved as semantic metadata; n/a is not an invented binding rank. \\
\phantomsection\label{entry:symbol-registry-05641}\(T(R),J_u\) & \texttt{\detokenize{T,}\allowbreak\detokenize{J}} & resumption functor and branch action & p10/domain-invariance & n/a & \hyperref[thm:p10-domain-invariance-resumptions]{\ref*{thm:p10-domain-invariance-resumptions}} & \thepage & Functor maps strict; computations inside its object need not be strict in state. Original precedence descriptor 'functor and map' preserved as semantic metadata; n/a is not an invented binding rank. \\
\phantomsection\label{entry:symbol-registry-05642}\(\operatorname{return}_s,\operatorname{step}(s,r),\Omega_s\) & \texttt{\detokenize{\operatorname{return},}\allowbreak\detokenize{\operatorname{step},}\allowbreak\detokenize{\Omega}} & resumption observations and infinite iteration & p10/domain-invariance & n/a & \hyperref[thm:p10-domain-invariance-resumptions]{\ref*{thm:p10-domain-invariance-resumptions}} & \thepage & Return and continuation tags distinct. Original precedence descriptor 'elements and map' preserved as semantic metadata; n/a is not an invented binding rank. \\
\phantomsection\label{entry:symbol-registry-05643}\(\mathcal R_n(D)\) & \texttt{\detokenize{\mathcal R}} & admissible finite-arity relations & p10/domain-invariance & object & \hyperref[def:p10-domain-invariance-relations]{\ref*{def:p10-domain-invariance-relations}} & \thepage & Bottom tuple and omega-chain closure. \\
\phantomsection\label{entry:symbol-registry-05644}\(u^{\mathrm{rel}}R\) & \texttt{\detokenize{\mathrm{rel}}} & coordinatewise relation direct image & p10/domain-invariance & n/a & \hyperref[def:p10-domain-invariance-relations]{\ref*{def:p10-domain-invariance-relations}} & \thepage & Does not assert every image is itself admissible. Original precedence descriptor 'operation' preserved as semantic metadata; n/a is not an invented binding rank. \\
\phantomsection\label{entry:symbol-registry-05645}\(F^{\mathrm{rel}}(R,S)\) & \texttt{\detokenize{F^{\mathrm{rel}}}} & typed admissible relation action & p10/domain-invariance & n/a & \hyperref[def:p10-domain-invariance-relations]{\ref*{def:p10-domain-invariance-relations}} & \thepage & Negative antitone; positive monotone; not application to relation objects. Original precedence descriptor 'operation' preserved as semantic metadata; n/a is not an invented binding rank. \\
\phantomsection\label{entry:symbol-registry-05646}\([R\longrightarrow S]\) & \texttt{\detokenize{\longrightarrow}} & logical or extensional function relation action & p10/domain-invariance & relation & \hyperref[thm:p10-domain-invariance-actions]{\ref*{thm:p10-domain-invariance-actions}} & \thepage & Unary preservation versus binary pointwise action. \\
\phantomsection\label{entry:symbol-registry-05647}\(R_{\bot}\) & \texttt{\detokenize{R_{\bot}}} & admissible lifting action & p10/domain-invariance & relation & \hyperref[thm:p10-domain-invariance-actions]{\ref*{thm:p10-domain-invariance-actions}} & \thepage & Binary carrier (D\_bot)\^{}2; new bottom and old injections explicit. \\
\phantomsection\label{entry:symbol-registry-05648}\(R\) & \texttt{\detokenize{R}} & admissible simulation & p10/domain-invariance & relation & \hyperref[thm:p10-domain-invariance-induction-coinduction]{\ref*{thm:p10-domain-invariance-induction-coinduction}} & \thepage & Unfold compatibility required for coinduction. \\
\phantomsection\label{entry:symbol-registry-05649}\(\infty\) & \texttt{\detokenize{\infty}} & nonterminating set-coalgebra value & p10/domain-invariance & n/a & \hyperref[thm:p10-domain-invariance-failures]{\ref*{thm:p10-domain-invariance-failures}} & \thepage & Not a universal domain bottom or top convention. Original precedence descriptor 'element' preserved as semantic metadata; n/a is not an invented binding rank. \\
\phantomsection\label{entry:symbol-registry-06335}\(s\) & \texttt{\detokenize{s}} & continuous section D to E & P10 continuous retractions & n/a & \hyperref[def:p10-dc-retraction]{\ref*{def:p10-dc-retraction}} & \thepage & Source-native s; rs=1. Typed map, not a variable assignment; actual section/retraction declaration precedes use. \\
\phantomsection\label{entry:symbol-registry-06336}\(r\) & \texttt{\detokenize{r}} & continuous retraction E to D & P10 continuous retractions & n/a & \hyperref[def:p10-dc-retraction]{\ref*{def:p10-dc-retraction}} & \thepage & Source-native r; not assumed to preserve every way-below pair. \\
\phantomsection\label{entry:symbol-registry-06337}\(f\) & \texttt{\detokenize{f}} & continuous idempotent endomap & P10 continuous retractions & n/a & \hyperref[def:p10-dc-retraction]{\ref*{def:p10-dc-retraction}} & \thepage & The idempotent endomap is f:E->E with f squared=f; argument/object types stated. \\
\phantomsection\label{entry:symbol-registry-06338}\(\operatorname{im}(f)\) & \texttt{\detokenize{\operatorname{im}}} & fixed-point image dcpo of a continuous idempotent & P10 continuous retractions & n/a & \hyperref[def:p10-dc-retraction]{\ref*{def:p10-dc-retraction}} & \thepage & Inherited order and parent directed suprema; not an arbitrary set image with an invented domain structure. \\
\phantomsection\label{entry:symbol-registry-06339}\(\operatorname{Idl}(B,\sqsubseteq)\) & \texttt{\detokenize{\operatorname{Idl}}} & algebraic ideal cover using reflexive information order & P10 algebraic retract cover & n/a & \hyperref[thm:p10-dc-algebraic-cover]{\ref*{thm:p10-dc-algebraic-cover}} & \thepage & Explicitly distinct from the round ideal completion using nonreflexive way-below; same source construction with its relation parameter stated. \\
\phantomsection\label{entry:symbol-registry-06340}\(S\) & \texttt{\detokenize{S}} & section on the continuous function dcpo & P10 function-domain retract transport & n/a & \hyperref[thm:p10-dc-function-retract]{\ref*{thm:p10-dc-function-retract}} & \thepage & S(h)=s\_E h r\_D; actual function-domain map, not a set of states. \\
\phantomsection\label{entry:symbol-registry-06341}\(R\) & \texttt{\detokenize{R}} & retraction on the continuous function dcpo & P10 function-domain retract transport & n/a & \hyperref[thm:p10-dc-function-retract]{\ref*{thm:p10-dc-function-retract}} & \thepage & R(k)=r\_E k s\_D and RS=1; projection inequality only under e-p hypotheses. \\
\phantomsection\label{entry:symbol-registry-06342}\(\mathcal P_f(B)\) & \texttt{\detokenize{\mathcal P_f}} & nonempty finite subsets of the chosen domain basis & P10 Plotkin finite basis & n/a & \hyperref[def:p10-dc-plotkin-basis]{\ref*{def:p10-dc-plotkin-basis}} & \thepage & No empty finite choice or nullary union is included. \\
\phantomsection\label{entry:symbol-registry-06343}\(\ll_{\mathrm{EM}}\) & \texttt{\detokenize{\ll_{\mathrm{EM}}}} & two-sided Egli-Milner approximation on finite basis sets & P10 Plotkin finite basis & n/a & \hyperref[def:p10-dc-plotkin-basis]{\ref*{def:p10-dc-plotkin-basis}} & \thepage & Both forall-exists clauses retained; not set inclusion or one-sided Hoare/Smyth approximation. \\
\phantomsection\label{entry:symbol-registry-06344}\(P^{\mathrm P}(D)\) & \texttt{\detokenize{P^{\mathrm P}}} & free Plotkin powerdomain of a continuous domain & P10 Plotkin construction & n/a & \hyperref[def:p10-dc-plotkin-basis]{\ref*{def:p10-dc-plotkin-basis}} & \thepage & Source powerdomain notation; actual round-ideal carrier and complete free universal property. \\
\phantomsection\label{entry:symbol-registry-06345}\(i(M)\) & \texttt{\detokenize{i}} & principal round ideal of the Egli-Milner basis & P10 Plotkin finite basis & n/a & \hyperref[def:p10-dc-plotkin-basis]{\ref*{def:p10-dc-plotkin-basis}} & \thepage & Same principal-round-ideal construction applied to the finite-set abstract basis; M need not belong to i(M) in the continuous case. \\
\phantomsection\label{entry:symbol-registry-06346}\(\eta_D\) & \texttt{\detokenize{\eta_D}} & continuous generator map into the free Plotkin algebra & P10 Plotkin construction & n/a & \hyperref[def:p10-dc-plotkin-basis]{\ref*{def:p10-dc-plotkin-basis}} & \thepage & Finite sets whose every element approximates x; not an assumed literal singleton of outcomes. \\
\phantomsection\label{entry:symbol-registry-06347}\(\cup\) & \texttt{\detokenize{\cup}} & formal continuous semilattice union & P10 Plotkin typed union & n/a & \hyperref[def:p10-dc-plotkin-basis]{\ref*{def:p10-dc-plotkin-basis}} & \thepage & SAME actual declaration/operation/proof regions as finite-set union: operands are powerdomain ideals or dcpo-semilattice elements. Not an information-order join. Preserve genuine typed overload, not fake disjoint scopes. \\
\phantomsection\label{entry:symbol-registry-06348}\(\cup\) & \texttt{\detokenize{\cup}} & ordinary union of finite basis sets & P10 Plotkin typed union & n/a & \hyperref[def:p10-dc-plotkin-basis]{\ref*{def:p10-dc-plotkin-basis}} & \thepage & SAME source-labelled regions also use M cup N as set union. Operands are finite basis sets; the original declaration explicitly contrasts this with formal ideal union. Canonical semantic reuse should be considered before appending a duplicate set-union identity. \\
\phantomsection\label{entry:symbol-registry-06349}\(\cup_A\) & \texttt{\detokenize{\cup_A}} & formal union in a target dcpo-semilattice & P10 Plotkin universal property & n/a & \hyperref[thm:p10-dc-free-powerdomain]{\ref*{thm:p10-dc-free-powerdomain}} & \thepage & Target operation is jointly Scott-continuous/commutative/associative/idempotent; no inflationary or deflationary information-order axiom imposed. \\
\phantomsection\label{entry:symbol-registry-06350}\(\operatorname{ext}(g)\) & \texttt{\detokenize{\operatorname{ext}}} & unique continuous semilattice homomorphic extension & P10 Plotkin universal property & n/a & \hyperref[thm:p10-dc-free-powerdomain]{\ref*{thm:p10-dc-free-powerdomain}} & \thepage & Defined by the directed finite-choice supremum and characterized uniquely on generators. \\
\phantomsection\label{entry:symbol-registry-06351}\(\tau_g(M)\) & \texttt{\detokenize{\tau_g}} & finite formal evaluation of a basis choice & P10 Plotkin universal property & n/a & \hyperref[annex:p10.dc.free-powerdomain]{\ref*{annex:p10.dc.free-powerdomain}} & \thepage & Actual proof declaration; repeated pairs preserve this value by idempotence, allowing monotonicity for both EM clauses. \\
\phantomsection\label{entry:symbol-registry-06352}\(\mu_D\) & \texttt{\detokenize{\mu_D}} & free Plotkin-algebra flattening & P10 Plotkin functor and monad & n/a & \hyperref[thm:p10-dc-powerdomain-functor]{\ref*{thm:p10-dc-powerdomain-functor}} & \thepage & Domain P\^{}P(P\^{}P(D)); actual extension of identity, naturality and three monad laws proved. Not the least recursive-solution endomorphism mu\_A. \\
\phantomsection\label{entry:symbol-registry-06353}\(\sqsubseteq_{\mathrm{EM}}\) & \texttt{\detokenize{\sqsubseteq_{\mathrm{EM}}}} & two-sided information refinement of algebraic finite choices & P10 Plotkin finite example & n/a & \hyperref[thm:p10-dc-choice-example]{\ref*{thm:p10-dc-choice-example}} & \thepage & Algebraic specialization of way-below lifting on compact basis; seven finite flat-domain outcome classes explicitly constructed. \\
\phantomsection\label{entry:symbol-registry-06354}\(b_n\) & \texttt{\detokenize{b_n}} & numbered compact basis element of an effective algebraic source domain & P10 effective algebraic choice & n/a & \hyperref[def:p10-dc-effective]{\ref*{def:p10-dc-effective}} & \thepage & Explicit fixed numbering and recursive compact order are object data; repeated names allowed. \\
\phantomsection\label{entry:symbol-registry-06355}\(c_m\) & \texttt{\detokenize{c_m}} & numbered compact basis element of an effective algebraic target domain & P10 effective algebraic choice & n/a & \hyperref[def:p10-dc-effective]{\ref*{def:p10-dc-effective}} & \thepage & Same compact-basis kind as b\_n with actual target operand type; no invented canonical representative of quotient names. \\
\phantomsection\label{entry:symbol-registry-06356}\(I_x\) & \texttt{\detokenize{I_x}} & recursively enumerable compact-name set of a computable point & P10 effective algebraic choice & n/a & \hyperref[def:p10-dc-effective]{\ref*{def:p10-dc-effective}} & \thepage & All n with b\_n below x; recursive enumerability not a promise of decidable point membership. \\
\phantomsection\label{entry:symbol-registry-06357}\(R_f\) & \texttt{\detokenize{R_f}} & recursively enumerable compact approximation matrix of an effective continuous map & P10 effective algebraic choice & n/a & \hyperref[def:p10-dc-effective]{\ref*{def:p10-dc-effective}} & \thepage & Pairs n m with c\_m below f(b\_n); explicit algebraic specialization of source approximable mappings. Composition and the Plotkin lift are proved by finite witnesses/dovetailing. \\
\phantomsection\label{entry:symbol-registry-07410}\(\mathsf{num}\) & \texttt{\detokenize{\mathsf{num}}} & Gunter PCF numerical ground type & P10 Gunter PCF language & n/a & \hyperref[def:p10-oa-language]{\ref*{def:p10-oa-language}} & \thepage & Source keyword num retained; not renamed to nat. \\
\phantomsection\label{entry:symbol-registry-07411}\(\mathsf{bool}\) & \texttt{\detokenize{\mathsf{bool}}} & Gunter PCF boolean ground type & P10 Gunter PCF language & n/a & \hyperref[def:p10-oa-language]{\ref*{def:p10-oa-language}} & \thepage & Exactly two boolean observations plus semantic bottom. \\
\phantomsection\label{entry:symbol-registry-07412}\(\mu\) & \texttt{\detokenize{\mu}} & Gunter PCF typed recursive term binder & P10 Gunter PCF language & n/a & \hyperref[def:p10-oa-language]{\ref*{def:p10-oa-language}} & \thepage & mu x:t.M is explicitly typed as syntax, not a probability mean or semantic monad multiplication; native glyph preserved. \\
\phantomsection\label{entry:symbol-registry-07413}\(H\) & \texttt{\detokenize{H}} & finite PCF variable type assignment & P10 Gunter PCF language & n/a & \hyperref[def:p10-oa-language]{\ref*{def:p10-oa-language}} & \thepage & Distinct variable names and explicit bound type annotations; source H retained. \\
\phantomsection\label{entry:symbol-registry-07414}\([N/x]M\) & \texttt{\detokenize{[N/}\allowbreak\detokenize{x]M}} & capture-avoiding PCF term substitution & P10 Gunter PCF language & n/a & \hyperref[def:p10-oa-language]{\ref*{def:p10-oa-language}} & \thepage & Both lambda and mu binders; unlike program contexts substitution does not capture free variables. \\
\phantomsection\label{entry:symbol-registry-07415}\(\bar n\) & \texttt{\detokenize{\bar n}} & PCF numeral succ iterated n times on zero & P10 Gunter PCF language & n/a & \hyperref[def:p10-oa-language]{\ref*{def:p10-oa-language}} & \thepage & Declared abbreviation of source numeral values, not a new primitive constant. \\
\phantomsection\label{entry:symbol-registry-07416}\(\operatorname{pred}\) & \texttt{\detokenize{\operatorname{pred}}} & Gunter PCF predecessor with pred zero equal zero & P10 Gunter PCF language & n/a & \hyperref[def:p10-oa-language]{\ref*{def:p10-oa-language}} & \thepage & Source rule differs from original Plotkin1977 arithmetic predecessor at zero; operational/model convention explicit. \\
\phantomsection\label{entry:symbol-registry-07417}\(\operatorname{zero?}\) & \texttt{\detokenize{\operatorname{zero?}}} & PCF strict numerical zero-test & P10 Gunter PCF language & n/a & \hyperref[def:p10-oa-language]{\ref*{def:p10-oa-language}} & \thepage & Returns true on zero and false on positive numeral; bottom input has bottom semantic result. \\
\phantomsection\label{entry:symbol-registry-07418}\(\to\) & \texttt{\detokenize{\to}} & PCF simple arrow type constructor & P10 Gunter PCF typed arrows & n/a & \hyperref[def:p10-oa-language]{\ref*{def:p10-oa-language}} & \thepage & SAME regions also use source arrow for small steps. Type operands s t distinguish the type constructor; no fake disjoint-file namespace. \\
\phantomsection\label{entry:symbol-registry-07419}\(\to\) & \texttt{\detokenize{\to}} & weak call-by-name PCF operational step & P10 Gunter PCF typed arrows & n/a & \hyperref[def:p10-oa-evaluation]{\ref*{def:p10-oa-evaluation}} & \thepage & Term operands M N distinguish operational steps from simple-type arrows and semantic maps; declaration before literal rule input. \\
\phantomsection\label{entry:symbol-registry-07420}\(E\) & \texttt{\detokenize{E}} & typed weak PCF evaluation context & P10 Gunter PCF evaluation & n/a & \hyperref[def:p10-oa-evaluation]{\ref*{def:p10-oa-evaluation}} & \thepage & Exactly the source operator/guard/unary-operation positions; no argument or under-lambda position. \\
\phantomsection\label{entry:symbol-registry-07421}\(\Downarrow\) & \texttt{\detokenize{\Downarrow}} & finite call-by-name PCF natural evaluation & P10 Gunter PCF evaluation & n/a & \hyperref[def:p10-oa-evaluation]{\ref*{def:p10-oa-evaluation}} & \thepage & Least relation generated by the source-natural rules; equivalent to finite weak small-step evaluation, not a denotational order. \\
\phantomsection\label{entry:symbol-registry-07422}\(\Uparrow\) & \texttt{\detokenize{\Uparrow}} & PCF program has no finite ground-value evaluation & P10 Gunter PCF evaluation & n/a & \hyperref[def:p10-oa-evaluation]{\ref*{def:p10-oa-evaluation}} & \thepage & Deterministic typed progress proves an infinite reduction sequence; this does not decide nontermination. \\
\phantomsection\label{entry:symbol-registry-07423}\(C\llbracket t\rrbracket\) & \texttt{\detokenize{C\llbracket t\rrbracket}} & Gunter continuous bc-domain interpretation of a PCF type & P10 Gunter PCF interpretation & n/a & \hyperref[def:p10-oa-interpretation]{\ref*{def:p10-oa-interpretation}} & \thepage & Source C interpretation distinct from stable D interpretation; D\_t is explicitly declared abbreviation, not source/version aliasing. \\
\phantomsection\label{entry:symbol-registry-07424}\(\llbracket M\rrbracket_\rho\) & \texttt{\detokenize{\llbracket M\rrbracket_\rho}} & PCF term denotation in a typed domain environment & P10 Gunter PCF interpretation & n/a & \hyperref[def:p10-oa-interpretation]{\ref*{def:p10-oa-interpretation}} & \thepage & Abbreviates source C[[H triangle M:t]]rho, retaining the chosen language/type/environment. \\
\phantomsection\label{entry:symbol-registry-07425}\(\mathcal R_t\) & \texttt{\detokenize{\mathcal R_t}} & typed semantic-element-to-PCF-term adequacy relation & P10 PCF adequacy logical relation & n/a & \hyperref[def:p10-oa-logical-relation]{\ref*{def:p10-oa-logical-relation}} & \thepage & Ground observation reflection and all-related-arguments arrow clause; not equality of two denotations or Plotkin's verbatim Comp predicate. \\
\phantomsection\label{entry:symbol-registry-07426}\(\sqsubseteq_o\) & \texttt{\detokenize{\sqsubseteq_o}} & Gunter PCF contextual operational approximation & P10 Gunter PCF contextual observations & n/a & \hyperref[def:p10-oa-contextual]{\ref*{def:p10-oa-contextual}} & \thepage & All typed closing ground contexts; their declared binders may capture inserted free variables. \\
\phantomsection\label{entry:symbol-registry-07427}\(\simeq_o\) & \texttt{\detokenize{\simeq_o}} & Gunter PCF contextual operational equivalence & P10 Gunter PCF contextual observations & n/a & \hyperref[def:p10-oa-contextual]{\ref*{def:p10-oa-contextual}} & \thepage & Two operational approximation directions; no full-abstraction converse asserted. \\
\phantomsection\label{entry:symbol-registry-07428}\(\Omega_t\) & \texttt{\detokenize{\Omega_t}} & typed divergent PCF term mu x:t.x & P10 PCF strategy boundaries & n/a & \hyperref[thm:p10-oa-strategy-boundaries]{\ref*{thm:p10-oa-strategy-boundaries}} & \thepage & Ground divergence equals bottom; at arrow type a terminating constant-bottom abstraction can have this same denotation. \\
\phantomsection\label{entry:symbol-registry-07429}\(K\) & \texttt{\detokenize{K}} & recursive PCF countdown example computing strict numerical identity & P10 PCF recursive example & n/a & \hyperref[thm:p10-oa-countdown]{\ref*{thm:p10-oa-countdown}} & \thepage & Editorial worked abbreviation of actual source syntax, not a new primitive or source-named constant; unevaluated numerical arguments are covered. \\
\phantomsection\label{entry:symbol-registry-07430}\(Q_t\) & \texttt{\detokenize{Q_t}} & ternary semantic sequentiality relation at PCF type t & P10 PCF sequentiality relations & n/a & \hyperref[def:p10-oa-sequential-relations]{\ref*{def:p10-oa-sequential-relations}} & \thepage & Editorial relation with first/second bottom or diagonal ground clause; arrow lift and full semantic fundamental lemma supplied. \\
\phantomsection\label{entry:symbol-registry-07431}\(A_t\) & \texttt{\detokenize{A_t}} & semantic elements with a Q\_t-related diagonal triple & P10 PCF sequentiality relations & n/a & \hyperref[def:p10-oa-sequential-relations]{\ref*{def:p10-oa-sequential-relations}} & \thepage & Selected semantic subset, not asserted to be a new full domain category or precisely the definable elements. \\
\phantomsection\label{entry:symbol-registry-07432}\(B_t\) & \texttt{\detokenize{B_t}} & binary logical relation on selected sequential semantic elements & P10 PCF sequentiality relations & n/a & \hyperref[def:p10-oa-sequential-relations]{\ref*{def:p10-oa-sequential-relations}} & \thepage & Ground equality and arrow relation on all related argument pairs; directed closure and all constructors proved. \\
\phantomsection\label{entry:symbol-registry-07433}\(T\) & \texttt{\detokenize{T}} & Gunter's first higher-type PCF full-abstraction counterexample term & P10 PCF source T F counterexample & n/a & \hyperref[thm:p10-oa-not-fully-abstract]{\ref*{thm:p10-oa-not-fully-abstract}} & \thepage & Native source Table9 name retained; all sequential arguments yield bottom but por yields true. \\
\phantomsection\label{entry:symbol-registry-07434}\(F\) & \texttt{\detokenize{F}} & Gunter's second higher-type PCF full-abstraction counterexample term & P10 PCF source T F counterexample & n/a & \hyperref[thm:p10-oa-not-fully-abstract]{\ref*{thm:p10-oa-not-fully-abstract}} & \thepage & Native source Table9 name retained; same source term except innermost false result, not generic semantic functional F. \\
\phantomsection\label{entry:symbol-registry-07435}\(\operatorname{por}\) & \texttt{\detokenize{\operatorname{por}}} & continuous parallel disjunction separating T and F in the full model & P10 PCF source T F counterexample & n/a & \hyperref[thm:p10-oa-not-fully-abstract]{\ref*{thm:p10-oa-not-fully-abstract}} & \thepage & Source Table10 all nine outcomes preserved; not added as a primitive of this sequential PCF. \\
\phantomsection\label{entry:symbol-registry-07436}\(\operatorname{succ}\) & \texttt{\detokenize{\operatorname{succ}}} & Gunter PCF numerical successor term constructor & P10 Gunter PCF language & n/a & \hyperref[def:p10-oa-language]{\ref*{def:p10-oa-language}} & \thepage & Strict numerical semantic operation; source constructor and numeral-value grammar preserved. \\
\phantomsection\label{entry:symbol-registry-07774}\(\mathcal F\) & \texttt{\detokenize{\mathcal F}} & source signature and symbols of specified finite arity symbol family member \(\mathcal F\) & P10 HSP source conventions & meta & \hyperref[def:p10-hsp-conventions]{\ref*{def:p10-hsp-conventions}} & \thepage & Native Burris2012 II1.2; same signature as host tau; not a free-algebra carrier. Explicit displayed family; member types and common actual first-use exposed, not one atomic glyph. Family listing, not one atomic glyph. Author macro \texttt{\detokenize{\mathcal F;\mathcal F\_n;\mathcal F\_0}} not carried to the members. \\
\phantomsection\label{entry:symbol-registry-07775}\(\mathcal F_n\) & \texttt{\detokenize{\mathcal F_n}} & source signature and symbols of specified finite arity symbol family member \(\mathcal F_n\) & P10 HSP source conventions & meta & \hyperref[def:p10-hsp-conventions]{\ref*{def:p10-hsp-conventions}} & \thepage & Native Burris2012 II1.2; same signature as host tau; not a free-algebra carrier. Explicit displayed family; member types and common actual first-use exposed, not one atomic glyph. Family listing, not one atomic glyph. Author macro \texttt{\detokenize{\mathcal F;\mathcal F\_n;\mathcal F\_0}} not carried to the members. \\
\phantomsection\label{entry:symbol-registry-07776}\(\mathcal F_0\) & \texttt{\detokenize{\mathcal F_0}} & source signature and symbols of specified finite arity symbol family member \(\mathcal F_0\) & P10 HSP source conventions & meta & \hyperref[def:p10-hsp-conventions]{\ref*{def:p10-hsp-conventions}} & \thepage & Native Burris2012 II1.2; same signature as host tau; not a free-algebra carrier. Explicit displayed family; member types and common actual first-use exposed, not one atomic glyph. Family listing, not one atomic glyph. Author macro \texttt{\detokenize{\mathcal F;\mathcal F\_n;\mathcal F\_0}} not carried to the members. \\
\phantomsection\label{entry:symbol-registry-07777}\(\mathbf A\) & \texttt{\detokenize{\mathbf A}} & algebra structure and its nonempty carrier symbol family member \(\mathbf A\) & P10 HSP constructions & object & \hyperref[def:p10-hsp-conventions]{\ref*{def:p10-hsp-conventions}} & \thepage & Native structure/carrier distinction from Burris2012 II1.3; operations total and nullary symbols included. Explicit displayed family; member types and common actual first-use exposed, not one atomic glyph. Family listing, not one atomic glyph. Author macro \texttt{\detokenize{\mathbf A;A}} not carried to the members. \\
\phantomsection\label{entry:symbol-registry-07778}\(A\) & \texttt{\detokenize{A}} & algebra structure and its nonempty carrier symbol family member \(A\) & P10 HSP constructions & object & \hyperref[def:p10-hsp-conventions]{\ref*{def:p10-hsp-conventions}} & \thepage & Native structure/carrier distinction from Burris2012 II1.3; operations total and nullary symbols included. Explicit displayed family; member types and common actual first-use exposed, not one atomic glyph. Family listing, not one atomic glyph. Author macro \texttt{\detokenize{\mathbf A;A}} not carried to the members. \\
\phantomsection\label{entry:symbol-registry-07779}\(T(X)\) & \texttt{\detokenize{T(X)}} & term set and term algebra over variables symbol family member \(T(X)\) & P10 HSP source conventions & object & \hyperref[def:p10-hsp-conventions]{\ref*{def:p10-hsp-conventions}} & \thepage & Burris2012 II10.1 and10.4; existence iff X union F0 is nonempty. Explicit displayed family; member types and common actual first-use exposed, not one atomic glyph. Family listing, not one atomic glyph. Author macro \texttt{\detokenize{T;\mathbf T}} not carried to the members. \\
\phantomsection\label{entry:symbol-registry-07780}\(\mathbf T(X)\) & \texttt{\detokenize{\mathbf T(X)}} & term set and term algebra over variables symbol family member \(\mathbf T(X)\) & P10 HSP source conventions & object & \hyperref[def:p10-hsp-conventions]{\ref*{def:p10-hsp-conventions}} & \thepage & Burris2012 II10.1 and10.4; existence iff X union F0 is nonempty. Explicit displayed family; member types and common actual first-use exposed, not one atomic glyph. Family listing, not one atomic glyph. Author macro \texttt{\detokenize{T;\mathbf T}} not carried to the members. \\
\phantomsection\label{entry:symbol-registry-07781}\(I\) & \texttt{\detokenize{I}} & source class operators and least nonempty variety symbol family member \(I\) & P10 Burris source conventions & meta & \hyperref[def:p10-hsp-conventions]{\ref*{def:p10-hsp-conventions}} & \thepage & Burris2012 II9.1–9.4; native P uses nonempty families; V has the nonempty-class convention. Explicit displayed family; member types and common actual first-use exposed, not one atomic glyph. Family listing, not one atomic glyph. Author macro \texttt{\detokenize{I;S;H;P;V}} not carried to the members. \\
\phantomsection\label{entry:symbol-registry-07782}\(S\) & \texttt{\detokenize{S}} & source class operators and least nonempty variety symbol family member \(S\) & P10 Burris source conventions & meta & \hyperref[def:p10-hsp-conventions]{\ref*{def:p10-hsp-conventions}} & \thepage & Burris2012 II9.1–9.4; native P uses nonempty families; V has the nonempty-class convention. Explicit displayed family; member types and common actual first-use exposed, not one atomic glyph. Family listing, not one atomic glyph. Author macro \texttt{\detokenize{I;S;H;P;V}} not carried to the members. \\
\phantomsection\label{entry:symbol-registry-07783}\(H\) & \texttt{\detokenize{H}} & source class operators and least nonempty variety symbol family member \(H\) & P10 Burris source conventions & meta & \hyperref[def:p10-hsp-conventions]{\ref*{def:p10-hsp-conventions}} & \thepage & Burris2012 II9.1–9.4; native P uses nonempty families; V has the nonempty-class convention. Explicit displayed family; member types and common actual first-use exposed, not one atomic glyph. Family listing, not one atomic glyph. Author macro \texttt{\detokenize{I;S;H;P;V}} not carried to the members. \\
\phantomsection\label{entry:symbol-registry-07784}\(P\) & \texttt{\detokenize{P}} & source class operators and least nonempty variety symbol family member \(P\) & P10 Burris source conventions & meta & \hyperref[def:p10-hsp-conventions]{\ref*{def:p10-hsp-conventions}} & \thepage & Burris2012 II9.1–9.4; native P uses nonempty families; V has the nonempty-class convention. Explicit displayed family; member types and common actual first-use exposed, not one atomic glyph. Family listing, not one atomic glyph. Author macro \texttt{\detokenize{I;S;H;P;V}} not carried to the members. \\
\phantomsection\label{entry:symbol-registry-07785}\(V\) & \texttt{\detokenize{V}} & source class operators and least nonempty variety symbol family member \(V\) & P10 Burris source conventions & meta & \hyperref[def:p10-hsp-conventions]{\ref*{def:p10-hsp-conventions}} & \thepage & Burris2012 II9.1–9.4; native P uses nonempty families; V has the nonempty-class convention. Explicit displayed family; member types and common actual first-use exposed, not one atomic glyph. Family listing, not one atomic glyph. Author macro \texttt{\detokenize{I;S;H;P;V}} not carried to the members. \\
\phantomsection\label{entry:symbol-registry-07786}\(\mathsf H\) & \texttt{\detokenize{\mathsf H}} & host homomorphic-image subalgebra and product operators symbol family member \(\mathsf H\) & P10 HSP host convention & meta & \hyperref[def:p10-hsp-conventions]{\ref*{def:p10-hsp-conventions}} & \thepage & Existing reader P also permits empty products; native-source and host product conventions are explicitly bridged rather than aliased. Explicit displayed family; member types and common actual first-use exposed, not one atomic glyph. Family listing, not one atomic glyph. Author macro \texttt{\detokenize{\mathsf H;\mathsf S;\mathsf P}} not carried to the members. \\
\phantomsection\label{entry:symbol-registry-07787}\(\mathsf S\) & \texttt{\detokenize{\mathsf S}} & host homomorphic-image subalgebra and product operators symbol family member \(\mathsf S\) & P10 HSP host convention & meta & \hyperref[def:p10-hsp-conventions]{\ref*{def:p10-hsp-conventions}} & \thepage & Existing reader P also permits empty products; native-source and host product conventions are explicitly bridged rather than aliased. Explicit displayed family; member types and common actual first-use exposed, not one atomic glyph. Family listing, not one atomic glyph. Author macro \texttt{\detokenize{\mathsf H;\mathsf S;\mathsf P}} not carried to the members. \\
\phantomsection\label{entry:symbol-registry-07788}\(\mathsf P\) & \texttt{\detokenize{\mathsf P}} & host homomorphic-image subalgebra and product operators symbol family member \(\mathsf P\) & P10 HSP host convention & meta & \hyperref[def:p10-hsp-conventions]{\ref*{def:p10-hsp-conventions}} & \thepage & Existing reader P also permits empty products; native-source and host product conventions are explicitly bridged rather than aliased. Explicit displayed family; member types and common actual first-use exposed, not one atomic glyph. Family listing, not one atomic glyph. Author macro \texttt{\detokenize{\mathsf H;\mathsf S;\mathsf P}} not carried to the members. \\
\phantomsection\label{entry:symbol-registry-07789}\(W\) & \texttt{\detokenize{W}} & identity alphabet generator alphabet and enlarged deduction alphabet symbol family member \(W\) & P10 HSP syntax & meta & \hyperref[def:p10-hsp-conventions]{\ref*{def:p10-hsp-conventions}} & \thepage & W countably infinite handles each finite variable support; X may be arbitrary set; U includes X and all variables of a law basis. Explicit displayed family; member types and common actual first-use exposed, not one atomic glyph. Family listing, not one atomic glyph. Author macro \texttt{\detokenize{W;X;U}} not carried to the members. \\
\phantomsection\label{entry:symbol-registry-07790}\(X\) & \texttt{\detokenize{X}} & identity alphabet generator alphabet and enlarged deduction alphabet symbol family member \(X\) & P10 HSP syntax & meta & \hyperref[def:p10-hsp-conventions]{\ref*{def:p10-hsp-conventions}} & \thepage & W countably infinite handles each finite variable support; X may be arbitrary set; U includes X and all variables of a law basis. Explicit displayed family; member types and common actual first-use exposed, not one atomic glyph. Family listing, not one atomic glyph. Author macro \texttt{\detokenize{W;X;U}} not carried to the members. \\
\phantomsection\label{entry:symbol-registry-07791}\(U\) & \texttt{\detokenize{U}} & identity alphabet generator alphabet and enlarged deduction alphabet symbol family member \(U\) & P10 HSP syntax & meta & \hyperref[def:p10-hsp-conventions]{\ref*{def:p10-hsp-conventions}} & \thepage & W countably infinite handles each finite variable support; X may be arbitrary set; U includes X and all variables of a law basis. Explicit displayed family; member types and common actual first-use exposed, not one atomic glyph. Family listing, not one atomic glyph. Author macro \texttt{\detokenize{W;X;U}} not carried to the members. \\
\phantomsection\label{entry:symbol-registry-07792}\(\operatorname{Id}_K(X)\) & \texttt{\detokenize{\operatorname{Id}_K(X)}} & common identity set and its model class symbol family member \(\operatorname{Id}_K(X)\) & P10 HSP semantics & meta & \hyperref[def:p10-hsp-conventions]{\ref*{def:p10-hsp-conventions}} & \thepage & Burris2012 II11.1 and11.7; definable proper classes still yield a set of identities for fixed X. Explicit displayed family; member types and common actual first-use exposed, not one atomic glyph. Family listing, not one atomic glyph. Author macro \texttt{\detokenize{\operatorname{Id};M;K;\Sigma}} not carried to the members. \\
\phantomsection\label{entry:symbol-registry-07793}\(M(\Sigma)\) & \texttt{\detokenize{M(\Sigma)}} & common identity set and its model class symbol family member \(M(\Sigma)\) & P10 HSP semantics & meta & \hyperref[def:p10-hsp-conventions]{\ref*{def:p10-hsp-conventions}} & \thepage & Burris2012 II11.1 and11.7; definable proper classes still yield a set of identities for fixed X. Explicit displayed family; member types and common actual first-use exposed, not one atomic glyph. Family listing, not one atomic glyph. Author macro \texttt{\detokenize{\operatorname{Id};M;K;\Sigma}} not carried to the members. \\
\phantomsection\label{entry:symbol-registry-07794}\(\Phi_K(X)\) & \texttt{\detokenize{\Phi_K(X)}} & set of relative quotient kernels and one congruence symbol family member \(\Phi_K(X)\) & P10 source relative-free construction & meta & \hyperref[def:p10-hsp-relative-free]{\ref*{def:p10-hsp-relative-free}} & \thepage & Burris2012 II10.9; each quotient belongs to IS(K); subset of Con of the set-sized term algebra. Explicit displayed family; member types and common actual first-use exposed, not one atomic glyph. Family listing, not one atomic glyph. Author macro \texttt{\detokenize{\Phi\_K;\phi}} not carried to the members. \\
\phantomsection\label{entry:symbol-registry-07795}\(\phi\) & \texttt{\detokenize{\phi}} & set of relative quotient kernels and one congruence symbol family member \(\phi\) & P10 source relative-free construction & meta & \hyperref[def:p10-hsp-relative-free]{\ref*{def:p10-hsp-relative-free}} & \thepage & Burris2012 II10.9; each quotient belongs to IS(K); subset of Con of the set-sized term algebra. Explicit displayed family; member types and common actual first-use exposed, not one atomic glyph. Family listing, not one atomic glyph. Author macro \texttt{\detokenize{\Phi\_K;\phi}} not carried to the members. \\
\phantomsection\label{entry:symbol-registry-07796}\(\theta_K(X)\) & \texttt{\detokenize{\theta_K}} & intersection of source quotient kernels & P10 source relative-free construction & meta & \hyperref[def:p10-hsp-relative-free]{\ref*{def:p10-hsp-relative-free}} & \thepage & Burris2012 II10.9; empty intersection is universal congruence; equivalent to all evaluation kernels by HSP.4. \\
\phantomsection\label{entry:symbol-registry-07797}\(\mathbf F_K(\bar X)\) & \texttt{\detokenize{\mathbf F_K(\bar X)}} & relative-free algebra generator map and classes symbol family member \(\mathbf F_K(\bar X)\) & P10 source relative-free construction & object & \hyperref[def:p10-hsp-relative-free]{\ref*{def:p10-hsp-relative-free}} & \thepage & Burris2012 II10.9–10.12; generator map may collapse in a trivial class; map formulation distinguishes original X from source class set barX. Explicit displayed family; member types and common actual first-use exposed, not one atomic glyph. Family listing, not one atomic glyph. Author macro \texttt{\detokenize{\mathbf F\_K;\eta\_X;\bar X;\bar p}} not carried to the members. \\
\phantomsection\label{entry:symbol-registry-07798}\(\eta_X\) & \texttt{\detokenize{\eta_X}} & relative-free algebra generator map and classes symbol family member \(\eta_X\) & P10 source relative-free construction & object & \hyperref[def:p10-hsp-relative-free]{\ref*{def:p10-hsp-relative-free}} & \thepage & Burris2012 II10.9–10.12; generator map may collapse in a trivial class; map formulation distinguishes original X from source class set barX. Explicit displayed family; member types and common actual first-use exposed, not one atomic glyph. Family listing, not one atomic glyph. Author macro \texttt{\detokenize{\mathbf F\_K;\eta\_X;\bar X;\bar p}} not carried to the members. \\
\phantomsection\label{entry:symbol-registry-07799}\(\bar X\) & \texttt{\detokenize{\bar X}} & relative-free algebra generator map and classes symbol family member \(\bar X\) & P10 source relative-free construction & object & \hyperref[def:p10-hsp-relative-free]{\ref*{def:p10-hsp-relative-free}} & \thepage & Burris2012 II10.9–10.12; generator map may collapse in a trivial class; map formulation distinguishes original X from source class set barX. Explicit displayed family; member types and common actual first-use exposed, not one atomic glyph. Family listing, not one atomic glyph. Author macro \texttt{\detokenize{\mathbf F\_K;\eta\_X;\bar X;\bar p}} not carried to the members. \\
\phantomsection\label{entry:symbol-registry-07800}\(\bar p\) & \texttt{\detokenize{\bar p}} & relative-free algebra generator map and classes symbol family member \(\bar p\) & P10 source relative-free construction & object & \hyperref[def:p10-hsp-relative-free]{\ref*{def:p10-hsp-relative-free}} & \thepage & Burris2012 II10.9–10.12; generator map may collapse in a trivial class; map formulation distinguishes original X from source class set barX. Explicit displayed family; member types and common actual first-use exposed, not one atomic glyph. Family listing, not one atomic glyph. Author macro \texttt{\detokenize{\mathbf F\_K;\eta\_X;\bar X;\bar p}} not carried to the members. \\
\phantomsection\label{entry:symbol-registry-07801}\(q_\phi\) & \texttt{\detokenize{q_\phi}} & quotient map chosen embedding and composite witness homomorphism symbol family member \(q_\phi\) & P10 separating product & meta & \hyperref[thm:p10-hsp-witness-product]{\ref*{thm:p10-hsp-witness-product}} & \thepage & Typed construction HSP.5; kernels of h\_phi are exactly phi; witnesses selected over a set after collection. Explicit displayed family; member types and common actual first-use exposed, not one atomic glyph. Family listing, not one atomic glyph. Author macro \texttt{\detokenize{q\_\phi;j\_\phi;h\_\phi}} not carried to the members. \\
\phantomsection\label{entry:symbol-registry-07802}\(j_\phi\) & \texttt{\detokenize{j_\phi}} & quotient map chosen embedding and composite witness homomorphism symbol family member \(j_\phi\) & P10 separating product & meta & \hyperref[thm:p10-hsp-witness-product]{\ref*{thm:p10-hsp-witness-product}} & \thepage & Typed construction HSP.5; kernels of h\_phi are exactly phi; witnesses selected over a set after collection. Explicit displayed family; member types and common actual first-use exposed, not one atomic glyph. Family listing, not one atomic glyph. Author macro \texttt{\detokenize{q\_\phi;j\_\phi;h\_\phi}} not carried to the members. \\
\phantomsection\label{entry:symbol-registry-07803}\(h_\phi\) & \texttt{\detokenize{h_\phi}} & quotient map chosen embedding and composite witness homomorphism symbol family member \(h_\phi\) & P10 separating product & meta & \hyperref[thm:p10-hsp-witness-product]{\ref*{thm:p10-hsp-witness-product}} & \thepage & Typed construction HSP.5; kernels of h\_phi are exactly phi; witnesses selected over a set after collection. Explicit displayed family; member types and common actual first-use exposed, not one atomic glyph. Family listing, not one atomic glyph. Author macro \texttt{\detokenize{q\_\phi;j\_\phi;h\_\phi}} not carried to the members. \\
\phantomsection\label{entry:symbol-registry-07804}\(d\) & \texttt{\detokenize{d}} & embedding of relative-free quotient and its product image symbol family member \(d\) & P10 separating product & meta & \hyperref[thm:p10-hsp-witness-product]{\ref*{thm:p10-hsp-witness-product}} & \thepage & Actual image subalgebra is generated by diagonal generator tuples; bold B is not the original law basis B. Explicit displayed family; member types and common actual first-use exposed, not one atomic glyph. Family listing, not one atomic glyph. Author macro \texttt{\detokenize{d;\mathbf B}} not carried to the members. \\
\phantomsection\label{entry:symbol-registry-07805}\(\mathbf B\) & \texttt{\detokenize{\mathbf B}} & embedding of relative-free quotient and its product image symbol family member \(\mathbf B\) & P10 separating product & meta & \hyperref[thm:p10-hsp-witness-product]{\ref*{thm:p10-hsp-witness-product}} & \thepage & Actual image subalgebra is generated by diagonal generator tuples; bold B is not the original law basis B. Explicit displayed family; member types and common actual first-use exposed, not one atomic glyph. Family listing, not one atomic glyph. Author macro \texttt{\detokenize{d;\mathbf B}} not carried to the members. \\
\phantomsection\label{entry:symbol-registry-07806}\(e_A\) & \texttt{\detokenize{e_A}} & target evaluation quotient surjection and image-subalgebra surjection symbol family member \(e_A\) & P10 HSP target representation & meta & \hyperref[def:p10-hsp-target-map]{\ref*{def:p10-hsp-target-map}} & \thepage & HSP.6–HSP.7; well-definedness uses theta inclusion in evaluation kernel; b\_A=r\_A composed with inverse embedding. Explicit displayed family; member types and common actual first-use exposed, not one atomic glyph. Family listing, not one atomic glyph. Author macro \texttt{\detokenize{e\_A;r\_A;b\_A}} not carried to the members. \\
\phantomsection\label{entry:symbol-registry-07807}\(r_A\) & \texttt{\detokenize{r_A}} & target evaluation quotient surjection and image-subalgebra surjection symbol family member \(r_A\) & P10 HSP target representation & meta & \hyperref[def:p10-hsp-target-map]{\ref*{def:p10-hsp-target-map}} & \thepage & HSP.6–HSP.7; well-definedness uses theta inclusion in evaluation kernel; b\_A=r\_A composed with inverse embedding. Explicit displayed family; member types and common actual first-use exposed, not one atomic glyph. Family listing, not one atomic glyph. Author macro \texttt{\detokenize{e\_A;r\_A;b\_A}} not carried to the members. \\
\phantomsection\label{entry:symbol-registry-07808}\(b_A\) & \texttt{\detokenize{b_A}} & target evaluation quotient surjection and image-subalgebra surjection symbol family member \(b_A\) & P10 HSP target representation & meta & \hyperref[def:p10-hsp-target-map]{\ref*{def:p10-hsp-target-map}} & \thepage & HSP.6–HSP.7; well-definedness uses theta inclusion in evaluation kernel; b\_A=r\_A composed with inverse embedding. Explicit displayed family; member types and common actual first-use exposed, not one atomic glyph. Family listing, not one atomic glyph. Author macro \texttt{\detokenize{e\_A;r\_A;b\_A}} not carried to the members. \\
\phantomsection\label{entry:symbol-registry-07809}\(x_a\) & \texttt{\detokenize{x_a}} & tagged generator for a target carrier element & P10 HSP target representation & object & \hyperref[def:p10-hsp-target-map]{\ref*{def:p10-hsp-target-map}} & \thepage & Nonempty target carrier gives a nonempty term alphabet independently of nullary operations. \\
\phantomsection\label{entry:symbol-registry-07810}\(\mathfrak A\) & \texttt{\detokenize{\mathfrak A}} & 1935 algebra basis primitive-symbol count valuations generated factors and diagonal elements symbol family member \(\mathfrak A\) & P10 original all-valuations route & object & \hyperref[def:p10-hsp-original-route]{\ref*{def:p10-hsp-original-route}} & \thepage & Original Birkhoff1935 pp440–441 product construction; modern valuation-index notation is declared; basis is a set not a proper-class product. Explicit displayed family; member types and common actual first-use exposed, not one atomic glyph. Family listing, not one atomic glyph. Author macro \texttt{\detokenize{\mathfrak A;n;\xi;S\_{i,\xi};S\^{}*;z\_x}} not carried to the members. \\
\phantomsection\label{entry:symbol-registry-07811}\(n\) & \texttt{\detokenize{n}} & 1935 algebra basis primitive-symbol count valuations generated factors and diagonal elements symbol family member \(n\) & P10 original all-valuations route & object & \hyperref[def:p10-hsp-original-route]{\ref*{def:p10-hsp-original-route}} & \thepage & Original Birkhoff1935 pp440–441 product construction; modern valuation-index notation is declared; basis is a set not a proper-class product. Explicit displayed family; member types and common actual first-use exposed, not one atomic glyph. Family listing, not one atomic glyph. Author macro \texttt{\detokenize{\mathfrak A;n;\xi;S\_{i,\xi};S\^{}*;z\_x}} not carried to the members. \\
\phantomsection\label{entry:symbol-registry-07812}\(\xi\) & \texttt{\detokenize{\xi}} & 1935 algebra basis primitive-symbol count valuations generated factors and diagonal elements symbol family member \(\xi\) & P10 original all-valuations route & object & \hyperref[def:p10-hsp-original-route]{\ref*{def:p10-hsp-original-route}} & \thepage & Original Birkhoff1935 pp440–441 product construction; modern valuation-index notation is declared; basis is a set not a proper-class product. Explicit displayed family; member types and common actual first-use exposed, not one atomic glyph. Family listing, not one atomic glyph. Author macro \texttt{\detokenize{\mathfrak A;n;\xi;S\_{i,\xi};S\^{}*;z\_x}} not carried to the members. \\
\phantomsection\label{entry:symbol-registry-07813}\(S_{i,\xi}\) & \texttt{\detokenize{S_{i,\xi}}} & 1935 algebra basis primitive-symbol count valuations generated factors and diagonal elements symbol family member \(S_{i,\xi}\) & P10 original all-valuations route & object & \hyperref[def:p10-hsp-original-route]{\ref*{def:p10-hsp-original-route}} & \thepage & Original Birkhoff1935 pp440–441 product construction; modern valuation-index notation is declared; basis is a set not a proper-class product. Explicit displayed family; member types and common actual first-use exposed, not one atomic glyph. Family listing, not one atomic glyph. Author macro \texttt{\detokenize{\mathfrak A;n;\xi;S\_{i,\xi};S\^{}*;z\_x}} not carried to the members. \\
\phantomsection\label{entry:symbol-registry-07814}\(S^*\) & \texttt{\detokenize{S^*}} & 1935 algebra basis primitive-symbol count valuations generated factors and diagonal elements symbol family member \(S^*\) & P10 original all-valuations route & object & \hyperref[def:p10-hsp-original-route]{\ref*{def:p10-hsp-original-route}} & \thepage & Original Birkhoff1935 pp440–441 product construction; modern valuation-index notation is declared; basis is a set not a proper-class product. Explicit displayed family; member types and common actual first-use exposed, not one atomic glyph. Family listing, not one atomic glyph. Author macro \texttt{\detokenize{\mathfrak A;n;\xi;S\_{i,\xi};S\^{}*;z\_x}} not carried to the members. \\
\phantomsection\label{entry:symbol-registry-07815}\(z_x\) & \texttt{\detokenize{z_x}} & 1935 algebra basis primitive-symbol count valuations generated factors and diagonal elements symbol family member \(z_x\) & P10 original all-valuations route & object & \hyperref[def:p10-hsp-original-route]{\ref*{def:p10-hsp-original-route}} & \thepage & Original Birkhoff1935 pp440–441 product construction; modern valuation-index notation is declared; basis is a set not a proper-class product. Explicit displayed family; member types and common actual first-use exposed, not one atomic glyph. Family listing, not one atomic glyph. Author macro \texttt{\detokenize{\mathfrak A;n;\xi;S\_{i,\xi};S\^{}*;z\_x}} not carried to the members. \\
\phantomsection\label{entry:symbol-registry-07816}\(F(\mathfrak A,n)\) & \texttt{\detokenize{F(\mathfrak A,}\allowbreak\detokenize{n)}} & original finite-generator free algebras and declared set-alphabet bridge symbol family member \(F(\mathfrak A,n)\) & P10 original all-valuations route & object & \hyperref[def:p10-hsp-original-route]{\ref*{def:p10-hsp-original-route}} & \thepage & Original equation10-1 and preceding construction; F signature in the book has a different declared type; original ordinal-rank scope is not claimed proved by finite terms. Explicit displayed family; member types and common actual first-use exposed, not one atomic glyph. Family listing, not one atomic glyph. Author macro \texttt{\detokenize{F;\mathbf F}} not carried to the members. \\
\phantomsection\label{entry:symbol-registry-07817}\(F(B,n)\) & \texttt{\detokenize{F(B,}\allowbreak\detokenize{n)}} & original finite-generator free algebras and declared set-alphabet bridge symbol family member \(F(B,n)\) & P10 original all-valuations route & object & \hyperref[def:p10-hsp-original-route]{\ref*{def:p10-hsp-original-route}} & \thepage & Original equation10-1 and preceding construction; F signature in the book has a different declared type; original ordinal-rank scope is not claimed proved by finite terms. Explicit displayed family; member types and common actual first-use exposed, not one atomic glyph. Family listing, not one atomic glyph. Author macro \texttt{\detokenize{F;\mathbf F}} not carried to the members. \\
\phantomsection\label{entry:symbol-registry-07818}\(\mathbf F(\mathfrak A,X)\) & \texttt{\detokenize{\mathbf F(\mathfrak A,}\allowbreak\detokenize{X)}} & original finite-generator free algebras and declared set-alphabet bridge symbol family member \(\mathbf F(\mathfrak A,X)\) & P10 original all-valuations route & object & \hyperref[def:p10-hsp-original-route]{\ref*{def:p10-hsp-original-route}} & \thepage & Original equation10-1 and preceding construction; F signature in the book has a different declared type; original ordinal-rank scope is not claimed proved by finite terms. Explicit displayed family; member types and common actual first-use exposed, not one atomic glyph. Family listing, not one atomic glyph. Author macro \texttt{\detokenize{F;\mathbf F}} not carried to the members. \\
\phantomsection\label{entry:symbol-registry-07819}\(R_{B,X}\) & \texttt{\detokenize{R_{B,X}}} & restricted deductive congruence and enlarged-alphabet deductive congruence symbol family member \(R_{B,X}\) & P10 modern deduction bridge & meta & \hyperref[thm:p10-hsp-original-product]{\ref*{thm:p10-hsp-original-product}} & \thepage & Restriction to T(X)\^{}2 permits a law basis with more variables than the chosen generator set; equality to semantic relative-free kernel follows from existing completeness. Explicit displayed family; member types and common actual first-use exposed, not one atomic glyph. Family listing, not one atomic glyph. Author macro \texttt{\detokenize{R\_{B,X};\theta\_{B,U}}} not carried to the members. \\
\phantomsection\label{entry:symbol-registry-07820}\(\theta_{B,U}\) & \texttt{\detokenize{\theta_{B,U}}} & restricted deductive congruence and enlarged-alphabet deductive congruence symbol family member \(\theta_{B,U}\) & P10 modern deduction bridge & meta & \hyperref[thm:p10-hsp-original-product]{\ref*{thm:p10-hsp-original-product}} & \thepage & Restriction to T(X)\^{}2 permits a law basis with more variables than the chosen generator set; equality to semantic relative-free kernel follows from existing completeness. Explicit displayed family; member types and common actual first-use exposed, not one atomic glyph. Family listing, not one atomic glyph. Author macro \texttt{\detokenize{R\_{B,X};\theta\_{B,U}}} not carried to the members. \\
\phantomsection\label{entry:symbol-registry-07821}\(\Sigma\) & \texttt{\detokenize{\Sigma}} & original1935 species of uniform operators & P10 original all-valuations route & meta & \hyperref[def:p10-hsp-original-route]{\ref*{def:p10-hsp-original-route}} & \thepage & Source Birkhoff1935 pp439–441; this is not the book's identity set Sigma used in M(Sigma). \\
\phantomsection\label{entry:symbol-registry-07822}\(\mathfrak F(\mathfrak A)\) & \texttt{\detokenize{\mathfrak F(\mathfrak A)}} & original generated algebra family or family determined by laws symbol family member \(\mathfrak F(\mathfrak A)\) & P10 original family correspondence & meta & \hyperref[def:p10-hsp-original-route]{\ref*{def:p10-hsp-original-route}} & \thepage & Source Definitions4–5 and Theorems9–10 pp440–441; algebra-basis and law-basis operands are explicitly typed source overloads. Explicit displayed family; member types and common actual first-use exposed, not one atomic glyph. Family listing, not one atomic glyph. Author macro \texttt{\detokenize{\mathfrak F}} not carried to the members. \\
\phantomsection\label{entry:symbol-registry-07823}\(\mathfrak F(B)\) & \texttt{\detokenize{\mathfrak F(B)}} & original generated algebra family or family determined by laws symbol family member \(\mathfrak F(B)\) & P10 original family correspondence & meta & \hyperref[def:p10-hsp-original-route]{\ref*{def:p10-hsp-original-route}} & \thepage & Source Definitions4–5 and Theorems9–10 pp440–441; algebra-basis and law-basis operands are explicitly typed source overloads. Explicit displayed family; member types and common actual first-use exposed, not one atomic glyph. Family listing, not one atomic glyph. Author macro \texttt{\detokenize{\mathfrak F}} not carried to the members. \\
\phantomsection\label{entry:symbol-registry-07824}\(\mathfrak F\) & \texttt{\detokenize{\mathfrak F}} & original generated algebra family or family determined by laws symbol family member \(\mathfrak F\) & P10 original family correspondence & meta & \hyperref[def:p10-hsp-original-route]{\ref*{def:p10-hsp-original-route}} & \thepage & Source Definitions4–5 and Theorems9–10 pp440–441; algebra-basis and law-basis operands are explicitly typed source overloads. Explicit displayed family; member types and common actual first-use exposed, not one atomic glyph. Family listing, not one atomic glyph. Author macro \texttt{\detokenize{\mathfrak F}} not carried to the members. \\
\phantomsection\label{entry:symbol-registry-07825}\(\Phi(\mathfrak A)\) & \texttt{\detokenize{\Phi(\mathfrak A)}} & original common-law family or deductive law family symbol family member \(\Phi(\mathfrak A)\) & P10 original family correspondence & meta & \hyperref[def:p10-hsp-original-route]{\ref*{def:p10-hsp-original-route}} & \thepage & Source Definitions3–5 and Theorems9–10 pp439–441; these are not the book's Phi\_K(X) set of congruence kernels. Explicit displayed family; member types and common actual first-use exposed, not one atomic glyph. Family listing, not one atomic glyph. Author macro \texttt{\detokenize{\Phi}} not carried to the members. \\
\phantomsection\label{entry:symbol-registry-07826}\(\Phi(B)\) & \texttt{\detokenize{\Phi(B)}} & original common-law family or deductive law family symbol family member \(\Phi(B)\) & P10 original family correspondence & meta & \hyperref[def:p10-hsp-original-route]{\ref*{def:p10-hsp-original-route}} & \thepage & Source Definitions3–5 and Theorems9–10 pp439–441; these are not the book's Phi\_K(X) set of congruence kernels. Explicit displayed family; member types and common actual first-use exposed, not one atomic glyph. Family listing, not one atomic glyph. Author macro \texttt{\detokenize{\Phi}} not carried to the members. \\
\phantomsection\label{entry:symbol-registry-07827}\(\Phi\) & \texttt{\detokenize{\Phi}} & original common-law family or deductive law family symbol family member \(\Phi\) & P10 original family correspondence & meta & \hyperref[def:p10-hsp-original-route]{\ref*{def:p10-hsp-original-route}} & \thepage & Source Definitions3–5 and Theorems9–10 pp439–441; these are not the book's Phi\_K(X) set of congruence kernels. Explicit displayed family; member types and common actual first-use exposed, not one atomic glyph. Family listing, not one atomic glyph. Author macro \texttt{\detokenize{\Phi}} not carried to the members. \\
\phantomsection\label{entry:symbol-registry-07828}\(\sim\) & \texttt{\detokenize{\sim}} & original algebra isomorphism marker & P10 original equation10-1 & meta & \hyperref[thm:p10-hsp-original-product]{\ref*{thm:p10-hsp-original-product}} & \thepage & Verified on original printed441; not a similarity relation or an equality of literal algebra carriers. \\
\phantomsection\label{entry:symbol-registry-07829}\(a_i\) & \texttt{\detokenize{a_i}} & finite basis carrier size and generator count symbol family member \(a_i\) & P10 original cardinality bound & meta & \hyperref[thm:p10-hsp-original-product]{\ref*{thm:p10-hsp-original-product}} & \thepage & Original p441 Corollary2; HSP.12 gives its finite-basis case from all a\_i\^{}m valuations rather than confusing product size with image size. Explicit displayed family; member types and common actual first-use exposed, not one atomic glyph. Family listing, not one atomic glyph. Author macro \texttt{\detokenize{a\_i;m}} not carried to the members. \\
\phantomsection\label{entry:symbol-registry-07830}\(m\) & \texttt{\detokenize{m}} & finite basis carrier size and generator count symbol family member \(m\) & P10 original cardinality bound & meta & \hyperref[thm:p10-hsp-original-product]{\ref*{thm:p10-hsp-original-product}} & \thepage & Original p441 Corollary2; HSP.12 gives its finite-basis case from all a\_i\^{}m valuations rather than confusing product size with image size. Explicit displayed family; member types and common actual first-use exposed, not one atomic glyph. Family listing, not one atomic glyph. Author macro \texttt{\detokenize{a\_i;m}} not carried to the members. \\
\phantomsection\label{entry:symbol-registry-07831}\(\ker h\) & \texttt{\detokenize{\ker h}} & homomorphism kernel congruence set and generated subalgebra symbol family member \(\ker h\) & P10 HSP constructions & meta & \hyperref[def:p10-hsp-relative-free]{\ref*{def:p10-hsp-relative-free}} & \thepage & Existing algebraic meanings retained; source Phi set contains quotient kernels not arbitrary equivalences. Explicit displayed family; member types and common actual first-use exposed, not one atomic glyph. Family listing, not one atomic glyph. Author macro \texttt{\detokenize{\ker;\operatorname{Con};\operatorname{Sg}}} not carried to the members. \\
\phantomsection\label{entry:symbol-registry-07832}\(\operatorname{Con}\) & \texttt{\detokenize{\operatorname{Con}}} & homomorphism kernel congruence set and generated subalgebra symbol family member \(\operatorname{Con}\) & P10 HSP constructions & meta & \hyperref[def:p10-hsp-relative-free]{\ref*{def:p10-hsp-relative-free}} & \thepage & Existing algebraic meanings retained; source Phi set contains quotient kernels not arbitrary equivalences. Explicit displayed family; member types and common actual first-use exposed, not one atomic glyph. Family listing, not one atomic glyph. Author macro \texttt{\detokenize{\ker;\operatorname{Con};\operatorname{Sg}}} not carried to the members. \\
\phantomsection\label{entry:symbol-registry-07833}\(\operatorname{Sg}\) & \texttt{\detokenize{\operatorname{Sg}}} & homomorphism kernel congruence set and generated subalgebra symbol family member \(\operatorname{Sg}\) & P10 HSP constructions & meta & \hyperref[def:p10-hsp-relative-free]{\ref*{def:p10-hsp-relative-free}} & \thepage & Existing algebraic meanings retained; source Phi set contains quotient kernels not arbitrary equivalences. Explicit displayed family; member types and common actual first-use exposed, not one atomic glyph. Family listing, not one atomic glyph. Author macro \texttt{\detokenize{\ker;\operatorname{Con};\operatorname{Sg}}} not carried to the members. \\
\phantomsection\label{entry:symbol-registry-07834}\(K(D)\) & \texttt{\detokenize{K}} & compact elements of a dI-domain & P10 stable-domain foundations & n/a & \hyperref[def:p10-sd-domains]{\ref*{def:p10-sd-domains}} & \thepage & Same compact-element construction as earlier domain theory; semantic reuse appropriate, with finite-principal-ideal hypothesis now stated. \\
\phantomsection\label{entry:symbol-registry-07835}\(x\mathrel{\uparrow}y\) & \texttt{\detokenize{\uparrow}} & compatibility by a common upper bound & P10 stable-domain foundations & n/a & \hyperref[def:p10-sd-domains]{\ref*{def:p10-sd-domains}} & \thepage & Declared shorthand for the source's common-upper-bound condition, not a claimed literal source glyph, numerical divergence or upward closure; operands x y are domain points. \\
\phantomsection\label{entry:symbol-registry-07836}\(\sqcap\) & \texttt{\detokenize{\sqcap}} & nonempty infimum in the declared dI domain order & P10 stable-domain foundations & n/a & \hyperref[def:p10-sd-domains]{\ref*{def:p10-sd-domains}} & \thepage & Meets exist; stable-map preservation requires compatible inputs. \\
\phantomsection\label{entry:symbol-registry-07837}\(\sqcup\) & \texttt{\detokenize{\sqcup}} & bounded join in the declared dI domain order & P10 stable-domain foundations & n/a & \hyperref[def:p10-sd-domains]{\ref*{def:p10-sd-domains}} & \thepage & No unbounded join is assumed; distributive-law hypothesis is compatibility of y z. \\
\phantomsection\label{entry:symbol-registry-07838}\([D\longrightarrow_s E]\) & \texttt{\detokenize{\longrightarrow_s}} & stable-function domain with stable order & P10 stable function domains & n/a & \hyperref[def:p10-sd-stable-order]{\ref*{def:p10-sd-stable-order}} & \thepage & Actual function carrier and finite-cut compact basis proved; not full continuous function space with renamed order. \\
\phantomsection\label{entry:symbol-registry-07839}\(\sqsubseteq_s\) & \texttt{\detokenize{\sqsubseteq_s}} & source stable order on stable functions & P10 stable function domains & n/a & \hyperref[def:p10-sd-stable-order]{\ref*{def:p10-sd-stable-order}} & \thepage & For x below y, f(x)=f(y) meet g(x); stronger than pointwise order. \\
\phantomsection\label{entry:symbol-registry-07840}\(\sqsubseteq\) & \texttt{\detokenize{\sqsubseteq}} & ordinary pointwise comparison of stable maps & P10 stable function domains & n/a & \hyperref[def:p10-sd-stable-order]{\ref*{def:p10-sd-stable-order}} & \thepage & Explicitly contrasted with stable order; ordinary domain-point information order retains its declared operand type. \\
\phantomsection\label{entry:symbol-registry-07841}\(M(f,x,b)\) & \texttt{\detokenize{M}} & least stable input modulus for compact output information & P10 stability moduli & n/a & \hyperref[def:p10-sd-stable-order]{\ref*{def:p10-sd-stable-order}} & \thepage & Least d below x with b below f(d), not f(d) below b; compact algebraic source specialization. \\
\phantomsection\label{entry:symbol-registry-07842}\(f_{a,b}\) & \texttt{\detokenize{f_{a,b}}} & finite input/output cut approximant of a stable function & P10 stable finite cuts & n/a & \hyperref[def:p10-sd-finite-cuts]{\ref*{def:p10-sd-finite-cuts}} & \thepage & Editorial full construction f(x meet a) meet b; actual compactness/algebraicity/property I proof, not attributed as printed source formula. \\
\phantomsection\label{entry:symbol-registry-07843}\(\sqcap_s\) & \texttt{\detokenize{\sqcap_s}} & meet in the stable function order & P10 stable function domains & n/a & \hyperref[thm:p10-sd-ordered-functions]{\ref*{thm:p10-sd-ordered-functions}} & \thepage & Pointwise only for a pair with common stable upper bound; p/q example has different stable and pointwise meets. \\
\phantomsection\label{entry:symbol-registry-07844}\(\downarrow_s h\) & \texttt{\detokenize{\downarrow_s}} & stable principal ideal of a stable function & P10 stable finite cuts & n/a & \hyperref[thm:p10-sd-function-domain]{\ref*{thm:p10-sd-function-domain}} & \thepage & Finite for compact h, established by actual finite restriction tables. \\
\phantomsection\label{entry:symbol-registry-07845}\(\operatorname{fix}\) & \texttt{\detokenize{\operatorname{fix}}} & least fixed-point operator as a stable morphism & P10 stable fixed point & n/a & \hyperref[thm:p10-sd-stable-fix]{\ref*{thm:p10-sd-stable-fix}} & \thepage & Same least-element fixed point construction, with new stable-order continuity/compatible-meet property independently proved. \\
\phantomsection\label{entry:symbol-registry-07846}\(S_t\) & \texttt{\detokenize{S_t}} & abbreviation for Gunter's stable PCF type interpretation D[[t]] & P10 stable PCF interpretation & n/a & \hyperref[def:p10-sd-pcf-model]{\ref*{def:p10-sd-pcf-model}} & \thepage & Not the continuous C interpretation; explicit source interpretation abbreviation, not title/version equivalence. \\
\phantomsection\label{entry:symbol-registry-07847}\(D\llbracket t\rrbracket\) & \texttt{\detokenize{D\llbracket t\rrbracket}} & Gunter stable dI-domain type interpretation & P10 stable PCF interpretation & n/a & \hyperref[def:p10-sd-pcf-model]{\ref*{def:p10-sd-pcf-model}} & \thepage & Native D interpretation distinguished from C; higher arrows have stable carrier/order. \\
\phantomsection\label{entry:symbol-registry-07848}\(T\) & \texttt{\detokenize{T}} & source flat boolean domain carrier & P10 stable PCF interpretation & n/a & \hyperref[def:p10-sd-pcf-model]{\ref*{def:p10-sd-pcf-model}} & \thepage & Actual domain declaration before use; not the earlier closed PCF full-abstraction counterexample term T. \\
\phantomsection\label{entry:symbol-registry-07849}\(\mathcal R^s_t\) & \texttt{\detokenize{\mathcal R^s_t}} & stable-carrier semantic-to-PCF-term adequacy relation & P10 stable PCF adequacy & n/a & \hyperref[def:p10-sd-relation]{\ref*{def:p10-sd-relation}} & \thepage & Actual reader declaration before use; explicitly reconstructed on stable type carriers, not silently the earlier continuous-carrier relation. \\
\phantomsection\label{entry:symbol-registry-07850}\(p\) & \texttt{\detokenize{p}} & constant-true stable boolean function & P10 stable source p q r counterexample & n/a & \hyperref[thm:p10-sd-counterexample]{\ref*{thm:p10-sd-counterexample}} & \thepage & SAME source theorem also uses p as a closed syntax abbreviation; mathematical function operands declare the semantic role. \\
\phantomsection\label{entry:symbol-registry-07851}\(p\) & \texttt{\detokenize{p}} & closed PCF abbreviation for constant-true abstraction & P10 stable source p q r counterexample & n/a & \hyperref[thm:p10-sd-counterexample]{\ref*{thm:p10-sd-counterexample}} & \thepage & SAME actual theorem/proof region, term p=lambda x.true; real syntax/denotation typing, not fictional disjoint filenames. \\
\phantomsection\label{entry:symbol-registry-07852}\(q\) & \texttt{\detokenize{q}} & strict-true stable boolean function & P10 stable source p q r counterexample & n/a & \hyperref[thm:p10-sd-counterexample]{\ref*{thm:p10-sd-counterexample}} & \thepage & Maps bottom to bottom and both booleans to true; pointwise below p but stable incompatible. \\
\phantomsection\label{entry:symbol-registry-07853}\(q\) & \texttt{\detokenize{q}} & closed PCF abbreviation for guard-strict true abstraction & P10 stable source p q r counterexample & n/a & \hyperref[thm:p10-sd-counterexample]{\ref*{thm:p10-sd-counterexample}} & \thepage & SAME actual theorem/proof region as semantic q; source term if x then true else true, preserving guard strictness. \\
\phantomsection\label{entry:symbol-registry-07854}\(r\) & \texttt{\detokenize{r}} & stable higher-type separator on the boolean function domain & P10 stable source p q r counterexample & n/a & \hyperref[thm:p10-sd-counterexample]{\ref*{thm:p10-sd-counterexample}} & \thepage & True at q, false at p, bottom elsewhere; monotonicity/continuity/stability proved, not added as PCF primitive. \\
\phantomsection\label{entry:symbol-registry-07855}\(M\) & \texttt{\detokenize{M}} & first source PCF counterexample term for the stable model & P10 stable source p q r counterexample & n/a & \hyperref[thm:p10-sd-counterexample]{\ref*{thm:p10-sd-counterexample}} & \thepage & Native source M name retained; explicit typed closed term, distinct from generic term metavariable role. \\
\phantomsection\label{entry:symbol-registry-07856}\(N\) & \texttt{\detokenize{N}} & second source PCF counterexample term for the stable model & P10 stable source p q r counterexample & n/a & \hyperref[thm:p10-sd-counterexample]{\ref*{thm:p10-sd-counterexample}} & \thepage & Native source N name retained; innermost false-result counterpart of M. \\
\end{longtable}
\endgroup
\Needspace{.4\textheight}
\subsection*{Stable part P11 scopes}
\begingroup\scriptsize\setlength{\tabcolsep}{3pt}
\begin{longtable}{@{}p{0.9in}p{1.3in}p{1.0in}p{0.5in}p{0.3in}p{0.4in}p{0.35in}p{1.65in}@{}}
\toprule
\textbf{Glyph} & \textbf{Macro} & \textbf{Role} & \textbf{Scope} & \textbf{Prec./form} & \textbf{First} & \textbf{Reg. page} & \textbf{Principal uses / notes} \\
\midrule
\endfirsthead
\toprule
\textbf{Glyph} & \textbf{Macro} & \textbf{Role} & \textbf{Scope} & \textbf{Prec./form} & \textbf{First} & \textbf{Reg. page} & \textbf{Principal uses / notes} \\
\midrule
\endhead
\midrule\multicolumn{8}{r}{\emph{continued}}\\\endfoot
\bottomrule\endlastfoot
\phantomsection\label{entry:symbol-registry-01874}\(F\) & \texttt{\detokenize{F}} & schematic action or attitude specification & \hyperref[sec:p11-reasons-and-rationality]{\ref*{sec:p11-reasons-and-rationality}} & meta & \hyperref[def:p11-reasons-response]{\ref*{def:p11-reasons-response}} & \thepage & Broome's schematic response content: an act or an attitude or arrangement of attitudes. Not a predicate symbol in a supplied formal calculus. Source: printed 131; local discussion assumes combined reasons do not require both F and not-F. The declaration-bound chapter scope remains stable when preceding chapters are inserted. \\
\phantomsection\label{entry:symbol-registry-02835}\(c\) & \texttt{\detokenize{c}} & injective code map on sentences & P11 Kripke truth fixed points & n/a & \hyperref[def:p11-ktf-base]{\ref*{def:p11-ktf-base}} & \thepage & Maps closed formulas into the fixed base domain; sentence codes are distinguished from names denoting them. \\
\phantomsection\label{entry:symbol-registry-02836}\(C\) & \texttt{\detokenize{C}} & set of sentence codes & P11 Kripke truth fixed points & n/a & \hyperref[def:p11-ktf-base]{\ref*{def:p11-ktf-base}} & \thepage & Image of the sentence code map. Noncodes enter the negative truth extension after jump 1. \\
\phantomsection\label{entry:symbol-registry-02837}\(T\) & \texttt{\detokenize{T}} & object-language unary partial truth predicate & P11 Kripke truth fixed points & n/a & \hyperref[def:p11-ktf-base]{\ref*{def:p11-ktf-base}} & \thepage & Ordinary predicates remain bivalent. At stage 0 this predicate is undefined on the whole domain. \\
\phantomsection\label{entry:symbol-registry-02838}\((S_1,S_2)\) & \texttt{\detokenize{(S_1,}\allowbreak\detokenize{S_2)}} & partial positive and negative truth extensions & P11 Kripke truth fixed points & n/a & \hyperref[def:p11-ktf-states]{\ref*{def:p11-ktf-states}} & \thepage & The two coordinates are disjoint subsets of the fixed domain. \\
\phantomsection\label{entry:symbol-registry-02839}\(\mathcal P\) & \texttt{\detokenize{\mathcal P}} & ordered set of consistent partial truth states & P11 Kripke truth fixed points & n/a & \hyperref[def:p11-ktf-states]{\ref*{def:p11-ktf-states}} & \thepage & This state space is not a lattice. Powerset notation applied to a set later has its usual distinct typed role. \\
\phantomsection\label{entry:symbol-registry-02840}\(\preceq\) & \texttt{\detokenize{\preceq}} & information extension relation between partial states & P11 Kripke truth fixed points & meta & \hyperref[def:p11-ktf-states]{\ref*{def:p11-ktf-states}} & \thepage & Coordinatewise inclusion; opposite signs are incomparable. \\
\phantomsection\label{entry:symbol-registry-02841}\(\bot_T\) & \texttt{\detokenize{\bot_T}} & wholly undefined truth state & P11 Kripke truth fixed points & n/a & \hyperref[def:p11-ktf-states]{\ref*{def:p11-ktf-states}} & \thepage & Both truth coordinates are empty; the base interpretation is fixed. \\
\phantomsection\label{entry:symbol-registry-02842}\(\mathbf u\) & \texttt{\detokenize{\mathbf u}} & undefined semantic value & P11 Kripke truth fixed points & n/a & \hyperref[def:p11-ktf-states]{\ref*{def:p11-ktf-states}} & \thepage & A gap in the interpretation, not a numerical probability or truth degree. \\
\phantomsection\label{entry:symbol-registry-02843}\(\leq_i\) & \texttt{\detokenize{\leq_i}} & information order on semantic values & P11 Kripke truth fixed points & meta & \hyperref[def:p11-ktf-states]{\ref*{def:p11-ktf-states}} & \thepage & Undefined is below each defined value; false and true are incomparable. \\
\phantomsection\label{entry:symbol-registry-02844}\(v_S(A,g)\) & \texttt{\detokenize{v_S(A,}\allowbreak\detokenize{g)}} & Strong-Kleene formula evaluation & P11 Kripke truth fixed points & n/a & \hyperref[def:p11-ktf-evaluation]{\ref*{def:p11-ktf-evaluation}} & \thepage & Sentences are independent of assignments. Terms are evaluated in the fixed base structure. \\
\phantomsection\label{entry:symbol-registry-02845}\(\Phi\) & \texttt{\detokenize{\Phi}} & Strong-Kleene truth jump & P11 Kripke truth fixed points & n/a & \hyperref[def:p11-ktf-jump]{\ref*{def:p11-ktf-jump}} & \thepage & Outputs true sentence codes positively and false sentence codes plus all noncodes negatively. \\
\phantomsection\label{entry:symbol-registry-02846}\(S^\alpha\) & \texttt{\detokenize{S^\alpha}} & ordinal iteration from an inflationary seed & P11 Kripke truth fixed points & n/a & \hyperref[def:p11-ktf-iteration]{\ref*{def:p11-ktf-iteration}} & \thepage & Successors use the jump and nonzero limits use compatible coordinate unions. \\
\phantomsection\label{entry:symbol-registry-02847}\(\operatorname{Cl}(S)\) & \texttt{\detokenize{\operatorname{Cl}(S)}} & least fixed extension of an inflationary seed & P11 Kripke truth fixed points & n/a & \hyperref[def:p11-ktf-iteration]{\ref*{def:p11-ktf-iteration}} & \thepage & Existence and countable stabilization are proved for countable sentence syntax. \\
\phantomsection\label{entry:symbol-registry-02848}\(K^\alpha\) & \texttt{\detokenize{K^\alpha}} & bottom iteration of the truth jump & P11 Kripke truth fixed points & n/a & \hyperref[def:p11-ktf-iteration]{\ref*{def:p11-ktf-iteration}} & \thepage & Its stage zero leaves even noncode truth atoms undefined. \\
\phantomsection\label{entry:symbol-registry-02849}\(K\) & \texttt{\detokenize{K}} & least Strong-Kleene truth fixed point & P11 Kripke truth fixed points & n/a & \hyperref[def:p11-ktf-iteration]{\ref*{def:p11-ktf-iteration}} & \thepage & Grounded values are defined values here; larger points can decide ungrounded sentences. \\
\phantomsection\label{entry:symbol-registry-02850}\(d\) & \texttt{\detokenize{d}} & total syntactic diagonal-substitution function & P11 Kripke truth fixed points & n/a & \hyperref[def:p11-ktf-arithmetic]{\ref*{def:p11-ktf-arithmetic}} & \thepage & A declared arithmetic function symbol acting on word codes; maps other inputs to zero. Domain objects denoted by an italic letter in assignment clauses have their ordinary local role. \\
\phantomsection\label{entry:symbol-registry-02851}\(n\) & \texttt{\detokenize{n}} & total negated-sentence-code function & P11 Kripke truth fixed points & n/a & \hyperref[def:p11-ktf-arithmetic]{\ref*{def:p11-ktf-arithmetic}} & \thepage & Maps a sentence code to its negation's code, and a non-sentence input to zero. Numerical indices use other letters. \\
\phantomsection\label{entry:symbol-registry-02852}\(q\) & \texttt{\detokenize{q}} & total recursive list of a finite truth chain's codes & P11 Kripke truth fixed points & n/a & \hyperref[def:p11-ktf-arithmetic]{\ref*{def:p11-ktf-arithmetic}} & \thepage & Reserved in the arithmetic vocabulary and determined by the syntactic recursion in the omega example. \\
\phantomsection\label{entry:symbol-registry-02853}\(\lambda\) & \texttt{\detokenize{\lambda}} & explicit liar sentence & P11 Kripke truth fixed points & n/a & \hyperref[def:p11-ktf-diagonal]{\ref*{def:p11-ktf-diagonal}} & \thepage & The displayed sentence negates the truth atom whose term denotes its own code. Limit ordinals in the general recursion are schematic indices. \\
\phantomsection\label{entry:symbol-registry-02854}\(\tau\) & \texttt{\detokenize{\tau}} & explicit truth-teller sentence & P11 Kripke truth fixed points & n/a & \hyperref[def:p11-ktf-diagonal]{\ref*{def:p11-ktf-diagonal}} & \thepage & Ungrounded at the least point and assigned either value by different larger fixed points. \\
\phantomsection\label{entry:symbol-registry-02855}\(\sigma\) & \texttt{\detokenize{\sigma}} & universal sentence about an unbounded finite truth chain & P11 Kripke truth fixed points & n/a & \hyperref[def:p11-ktf-omega-example]{\ref*{def:p11-ktf-omega-example}} & \thepage & First evaluated true at the first limit stage; its code enters one successor later. \\
\phantomsection\label{entry:symbol-registry-02856}\(\rho\) & \texttt{\detokenize{\rho}} & self-supporting intrinsically true sentence & P11 Kripke truth fixed points & n/a & \hyperref[def:p11-ktf-intrinsic-example]{\ref*{def:p11-ktf-intrinsic-example}} & \thepage & Says that it or its negation is true; ungrounded but assigned true in an intrinsic fixed point. \\
\phantomsection\label{entry:symbol-registry-02857}\(\chi\) & \texttt{\detokenize{\chi}} & truth-teller excluded-middle sentence & P11 Kripke truth fixed points & n/a & \hyperref[def:p11-ktf-intrinsic-example]{\ref*{def:p11-ktf-intrinsic-example}} & \thepage & Never false at a Strong-Kleene fixed point yet defined at no intrinsic point. \\
\phantomsection\label{entry:symbol-registry-02858}\(H\) & \texttt{\detokenize{H}} & parameterized universal falsity assertion & P11 Kripke truth fixed points & n/a & \hyperref[def:p11-ktf-empirical]{\ref*{def:p11-ktf-empirical}} & \thepage & Held syntax is grounded or paradoxical according to the bivalent empirical predicate's extension. \\
\phantomsection\label{entry:symbol-registry-02859}\(A_r\) & \texttt{\detokenize{A_r}} & bivalent finite-level sentence-code guard & P11 Kripke truth fixed points & n/a & \hyperref[def:p11-ktf-tarski-guards]{\ref*{def:p11-ktf-tarski-guards}} & \thepage & A syntactic base predicate; subscripted guards are distinct from schematic formulas. \\
\phantomsection\label{entry:symbol-registry-02860}\(T_r\) & \texttt{\detokenize{T_r}} & guarded truth predicate for a finite language level & P11 Kripke truth fixed points & n/a & \hyperref[def:p11-ktf-tarski-guards]{\ref*{def:p11-ktf-tarski-guards}} & \thepage & Defined as the truth predicate conjoined with the corresponding sentence-code guard. \\
\phantomsection\label{entry:symbol-registry-02861}\(\iota\) & \texttt{\detokenize{\iota}} & fixed injection coding finite domain sequences & P11 Kripke truth fixed points & n/a & \hyperref[def:p11-ktf-satisfaction]{\ref*{def:p11-ktf-satisfaction}} & \thepage & Extra resource for satisfaction; no sequence-code assumption was needed for truth on sentences. \\
\phantomsection\label{entry:symbol-registry-02862}\(c_f\) & \texttt{\detokenize{c_f}} & injective formula code map for the satisfaction language & P11 Kripke truth fixed points & n/a & \hyperref[def:p11-ktf-satisfaction]{\ref*{def:p11-ktf-satisfaction}} & \thepage & Includes formulas containing the added binary satisfaction predicate. \\
\phantomsection\label{entry:symbol-registry-02863}\(\mathsf{Sat}\) & \texttt{\detokenize{\mathsf{Sat}}} & binary object-language partial satisfaction predicate & P11 Kripke truth fixed points & n/a & \hyperref[def:p11-ktf-satisfaction]{\ref*{def:p11-ktf-satisfaction}} & \thepage & Arguments name a coded finite assignment and a formula; invalid pairs become false after the first jump. \\
\phantomsection\label{entry:symbol-registry-02864}\(\Psi\) & \texttt{\detokenize{\Psi}} & Strong-Kleene satisfaction jump & P11 Kripke truth fixed points & n/a & \hyperref[def:p11-ktf-satisfaction]{\ref*{def:p11-ktf-satisfaction}} & \thepage & Operates on partial subsets of the square of the base domain. \\
\phantomsection\label{entry:symbol-registry-02865}\(W\) & \texttt{\detokenize{W}} & least closed set of signed sentences & P11 Kripke truth fixed points & n/a & \hyperref[def:p11-ktf-signed]{\ref*{def:p11-ktf-signed}} & \thepage & Includes base signs, equal-denotation replacement, forcing cases and truth-prefix introduction; proved equal to the least point's defined sentence values. \\
\phantomsection\label{entry:symbol-registry-02866}\(v^{\mathrm{wk}}_S\) & \texttt{\detokenize{v^{\mathrm{wk}}_S}} & weak-Kleene evaluation & P11 Kripke truth fixed points & n/a & \hyperref[def:p11-ktf-weak]{\ref*{def:p11-ktf-weak}} & \thepage & A gap in an operand or instance propagates to the compound or quantifier. \\
\phantomsection\label{entry:symbol-registry-02867}\(v^{\mathrm{sv}}_S\) & \texttt{\detokenize{v^{\mathrm{sv}}_S}} & unrestricted supervaluation & P11 Kripke truth fixed points & n/a & \hyperref[def:p11-ktf-weak]{\ref*{def:p11-ktf-weak}} & \thepage & Unanimous classical evaluation over all total extensions of the partial state. \\
\phantomsection\label{entry:symbol-registry-02868}\(\Phi_{\mathrm{wk}}\) & \texttt{\detokenize{\Phi_{\mathrm{wk}}}} & weak-Kleene truth jump & P11 Kripke truth fixed points & n/a & \hyperref[def:p11-ktf-weak]{\ref*{def:p11-ktf-weak}} & \thepage & Distinct operation with the same sentence-code and noncode staging clauses. \\
\phantomsection\label{entry:symbol-registry-02869}\(\Phi_{\mathrm{sv}}\) & \texttt{\detokenize{\Phi_{\mathrm{sv}}}} & unrestricted supervaluation truth jump & P11 Kripke truth fixed points & n/a & \hyperref[def:p11-ktf-weak]{\ref*{def:p11-ktf-weak}} & \thepage & Distinct operation whose least point can define more sentence values than the Strong-Kleene least point. \\
\phantomsection\label{entry:symbol-registry-02870}\(\widehat F\) & \texttt{\detokenize{\widehat F}} & closed-off total extension of a partial truth fixed point & P11 Kripke truth fixed points & n/a & \hyperref[def:p11-ktf-closed]{\ref*{def:p11-ktf-closed}} & \thepage & Retains only the former positive extension as its total truth predicate; it is not itself a fixed point. \\
\phantomsection\label{entry:symbol-registry-02871}\(\widehat v_F\) & \texttt{\detokenize{\widehat v_F}} & classical evaluation after closing off & P11 Kripke truth fixed points & n/a & \hyperref[def:p11-ktf-closed]{\ref*{def:p11-ktf-closed}} & \thepage & Distinguishes the classical truth set of the new model from membership in the old positive truth extension. \\
\phantomsection\label{entry:symbol-registry-02872}\(\omega_1^{\mathrm{CK}}\) & \texttt{\detokenize{\omega_1^{\mathrm{CK}}}} & first nonrecursive ordinal & P11 Kripke truth fixed points & n/a & \hyperref[def:p11-ktf-effective]{\ref*{def:p11-ktf-effective}} & \thepage & Effective arithmetic classification, distinct from the arbitrary-syntax countable bound. \\
\phantomsection\label{entry:symbol-registry-02873}\(\kappa\) & \texttt{\detokenize{\kappa}} & ungrounded sentence defined at some but not all maximal points & P11 Kripke truth fixed points & n/a & \hyperref[def:p11-ktf-maximal]{\ref*{def:p11-ktf-maximal}} & \thepage & The truth-teller disjoined with its negation conjoined with the liar. \\
\phantomsection\label{entry:symbol-registry-02874}\(H_1,H_2\) & \texttt{\detokenize{H_1,}\allowbreak\detokenize{H_2}} & mutually dependent universal falsity assertions & P11 Kripke truth fixed points & n/a & \hyperref[def:p11-ktf-empirical]{\ref*{def:p11-ktf-empirical}} & \thepage & Their ground ranks depend on an ordinary true assertion in one predicate's extension. \\
\phantomsection\label{entry:symbol-registry-02875}\(K_{\mathrm{wk}},K_{\mathrm{sv}}\) & \texttt{\detokenize{K_{\mathrm{wk}},}\allowbreak\detokenize{K_{\mathrm{sv}}}} & least fixed points of the weak and supervaluation jumps & P11 Kripke truth fixed points & n/a & \hyperref[def:p11-ktf-weak]{\ref*{def:p11-ktf-weak}} & \thepage & Scheme-specific least points; the underlying truth predicate is not renamed. \\
\phantomsection\label{entry:symbol-registry-02876}\(Q^*\) & \texttt{\detokenize{Q^*}} & least satisfaction fixed point & P11 Kripke truth fixed points & n/a & \hyperref[thm:p11-ktf-satisfaction]{\ref*{thm:p11-ktf-satisfaction}} & \thepage & Its closure bound counts satisfaction arguments rather than only formula codes. \\
\phantomsection\label{entry:symbol-registry-02877}\(\mathsf{Sent}\) & \texttt{\detokenize{\mathsf{Sent}}} & bivalent sentence-code recognition predicate & P11 Kripke truth fixed points & n/a & \hyperref[def:p11-ktf-restricted]{\ref*{def:p11-ktf-restricted}} & \thepage & Provides the guard on the restricted-completion disjunction. \\
\phantomsection\label{entry:symbol-registry-02878}\(\mathsf{PQT}\) & \texttt{\detokenize{\mathsf{PQT}}} & bivalent recognition of classical first-order theorem codes & P11 Kripke truth fixed points & n/a & \hyperref[thm:p11-ktf-schemes]{\ref*{thm:p11-ktf-schemes}} & \thepage & Its guarded universal truth assertion is settled at stage 1 of the unrestricted supervaluation iteration. \\
\phantomsection\label{entry:symbol-registry-02879}\(\mathcal T_n\) & \texttt{\detokenize{\mathcal T_n}} & uniform recursive family of finite-sequence trees & P11 Kripke truth fixed points & n/a & \hyperref[def:p11-ktf-recursive-trees]{\ref*{def:p11-ktf-recursive-trees}} & \thepage & Rooted trees are prefix closed and a branch consists of all prefixes of an infinite sequence. \\
\phantomsection\label{entry:symbol-registry-02880}\(h_{\mathcal T}(s)\) & \texttt{\detokenize{h_{\mathcal T}(s)}} & height of a well-founded tree node & P11 Kripke truth fixed points & n/a & \hyperref[def:p11-ktf-recursive-trees]{\ref*{def:p11-ktf-recursive-trees}} & \thepage & The supremum of immediate child heights plus one; leaves have height zero. \\
\phantomsection\label{entry:symbol-registry-02881}\(\psi_z\) & \texttt{\detokenize{\psi_z}} & actual diagonal truth sentence for a forest node & P11 Kripke truth fixed points & n/a & \hyperref[def:p11-ktf-tree-formulas]{\ref*{def:p11-ktf-tree-formulas}} & \thepage & Its code map is computable and injective, and its evaluation rank equals its well-founded node height. \\
\phantomsection\label{entry:symbol-registry-02882}\(\operatorname{Closed}_{\mathrm{SK}}(X)\) & \texttt{\detokenize{\operatorname{Closed}_{\mathrm{SK}}(X)}} & arithmetical signed-closure condition & P11 Kripke truth fixed points & n/a & \hyperref[annex:p11.ktf.pi-upper]{\ref*{annex:p11.ktf.pi-upper}} & \thepage & The set variable codes signs; every closure premise is specified using number quantifiers. \\
\phantomsection\label{entry:symbol-registry-02883}\(G(X)\) & \texttt{\detokenize{G(X)}} & numerical supervaluation jump on the ambient full powerset & P11 Kripke truth fixed points & n/a & \hyperref[annex:p11.ktf.pi-upper]{\ref*{annex:p11.ktf.pi-upper}} & \thepage & Extends the signed-state operation to arbitrary subsets for the pre-fixed-set argument; the actual bottom iteration stays consistent. \\
\phantomsection\label{entry:symbol-registry-02884}\(\operatorname{Val}(R,V)\) & \texttt{\detokenize{\operatorname{Val}(R,}\allowbreak\detokenize{V)}} & arithmetical verification of a classical valuation set & P11 Kripke truth fixed points & n/a & \hyperref[annex:p11.ktf.pi-upper]{\ref*{annex:p11.ktf.pi-upper}} & \thepage & Two set arguments encode a total truth-predicate interpretation and its ordinary classical sentence truth set. \\
\phantomsection\label{entry:symbol-registry-02885}\(\operatorname{Pre}(X,Y)\) & \texttt{\detokenize{\operatorname{Pre}(X,}\allowbreak\detokenize{Y)}} & arithmetical packed countermodel condition for a pre-fixed set & P11 Kripke truth fixed points & n/a & \hyperref[annex:p11.ktf.pi-upper]{\ref*{annex:p11.ktf.pi-upper}} & \thepage & The second set packs countermodel and valuation witnesses for all missing signed conclusions. \\
\phantomsection\label{entry:symbol-registry-02886}\(L_N^r,L_C^r\) & \texttt{\detokenize{L_N^r,}\allowbreak\detokenize{L_C^r}} & noncumulative and cumulative auxiliary Tarski languages & P11 Kripke truth fixed points & n/a & \hyperref[def:p11-ktf-two-finite-hierarchies]{\ref*{def:p11-ktf-two-finite-hierarchies}} & \thepage & Classical finite language levels with separate predicate tokens and numerical code operations. \\
\phantomsection\label{entry:symbol-registry-02887}\(\mathsf N_r,\mathsf C_r\) & \texttt{\detokenize{\mathsf N_r,}\allowbreak\detokenize{\mathsf C_r}} & classical truth predicates for the two preceding auxiliary levels & P11 Kripke truth fixed points & n/a & \hyperref[def:p11-ktf-two-finite-hierarchies]{\ref*{def:p11-ktf-two-finite-hierarchies}} & \thepage & Each predicate is false on every input that is not a true sentence code of its own preceding language. \\
\phantomsection\label{entry:symbol-registry-02888}\(u_r,v_r\) & \texttt{\detokenize{u_r,}\allowbreak\detokenize{v_r}} & effective translations between the finite auxiliary hierarchies & P11 Kripke truth fixed points & n/a & \hyperref[thm:p11-ktf-finite-equivalence]{\ref*{thm:p11-ktf-finite-equivalence}} & \thepage & Preserve free variables and classical values under all assignments; they are not asserted to be syntactic inverses. \\
\phantomsection\label{entry:symbol-registry-03046}\(\mathsf U(e,z)\) & \texttt{\detokenize{\mathsf U(e,}\allowbreak\detokenize{z)}} & universal analytical definition-index relation & P11 Kripke truth fixed points & n/a & \hyperref[def:p11-ktf-pi-universal]{\ref*{def:p11-ktf-pi-universal}} & \thepage & Uses separately coded second-order definitions and a numerical instance argument; invalid definition codes have no instances. \\
\phantomsection\label{entry:symbol-registry-03047}\(\mathcal S_{e,z}\) & \texttt{\detokenize{\mathcal S_{e,z}}} & recursive tree representing a universal analytical instance & P11 Kripke truth fixed points & n/a & \hyperref[def:p11-ktf-pi-universal]{\ref*{def:p11-ktf-pi-universal}} & \thepage & Its Kleene-Brouwer order is a recursive well-order when the instance holds. \\
\phantomsection\label{entry:symbol-registry-03048}\(\rho(e,z)\) & \texttt{\detokenize{\rho(e,}\allowbreak\detokenize{z)}} & recursive ordinal type of the analytical comparison tree order & P11 Kripke truth fixed points & n/a & \hyperref[def:p11-ktf-pi-universal]{\ref*{def:p11-ktf-pi-universal}} & \thepage & Defined when the corresponding universal instance holds; distinct typed use from the self-supporting sentence. \\
\phantomsection\label{entry:symbol-registry-03049}\(\operatorname{Below}(p,q)\) & \texttt{\detokenize{\operatorname{Below}(p,}\allowbreak\detokenize{q)}} & strict analytical norm comparison with specified outside-domain behaviour & P11 Kripke truth fixed points & n/a & \hyperref[def:p11-ktf-pi-universal]{\ref*{def:p11-ktf-pi-universal}} & \thepage & No initial-segment embedding from the second order into the first well-order; the second argument may be outside the well-founded domain. \\
\phantomsection\label{entry:symbol-registry-03050}\(\Gamma\) & \texttt{\detokenize{\Gamma}} & monotone numerical induction operator & P11 Kripke truth fixed points & n/a & \hyperref[thm:p11-ktf-effective-induction]{\ref*{thm:p11-ktf-effective-induction}} & \thepage & Requires a verified effective compiler on analytical definitions; monotonicity alone does not supply that hypothesis. \\
\phantomsection\label{entry:symbol-registry-03051}\(\operatorname{diag}_{\Pi}\) & \texttt{\detokenize{\operatorname{diag}_{\Pi}}} & effective self-index substitution on analytical formula codes & P11 Kripke truth fixed points & n/a & \hyperref[annex:p11.ktf.effective-induction]{\ref*{annex:p11.ktf.effective-induction}} & \thepage & Replaces the first designated free number variable by the numeral for its own formula code; invalid inputs return a fixed false definition code. \\
\phantomsection\label{entry:symbol-registry-03052}\(\Omega\) & \texttt{\detokenize{\Omega}} & positive arithmetic signed-rule closure operator & P11 Kripke truth fixed points & n/a & \hyperref[annex:p11.ktf.ck-exact]{\ref*{annex:p11.ktf.ck-exact}} & \thepage & Retains the input signs and adds base and forcing-rule conclusions; its stages are compared with the original truth stages. \\
\phantomsection\label{entry:symbol-registry-03053}\(\vartheta_z\) & \texttt{\detokenize{\vartheta_z}} & alternating product-tree comparison sentence & P11 Kripke truth fixed points & n/a & \hyperref[def:p11-ktf-product-comparison]{\ref*{def:p11-ktf-product-comparison}} & \thepage & Uses an outer universal over first-tree children and an inner existential over second-tree children; total at the least point when one subtree is well founded. \\
\phantomsection\label{entry:symbol-registry-03234}\(\mathcal A\) & \texttt{\detokenize{\mathcal A}} & acceptable finite-signature relational structure & P11 acceptable-structure truth & n/a & \hyperref[def:p11-ktf-acceptable]{\ref*{def:p11-ktf-acceptable}} & \thepage & Infinite domain with elementary internal numerical and finite-sequence resources; all definitions permit finitely many domain parameters. \\
\phantomsection\label{entry:symbol-registry-03235}\(\mathsf N_{\mathcal A}\) & \texttt{\detokenize{\mathsf N_{\mathcal A}}} & elementary standard numerical copy inside the acceptable domain & P11 acceptable-structure truth & n/a & \hyperref[def:p11-ktf-acceptable]{\ref*{def:p11-ktf-acceptable}} & \thepage & Its order is isomorphic to ordinary natural numbers; no nonstandard numerical copy is substituted. \\
\phantomsection\label{entry:symbol-registry-03236}\(\leq_{\mathcal A}\) & \texttt{\detokenize{\leq_{\mathcal A}}} & elementary non-strict numerical order in the structure & P11 acceptable-structure truth & n/a & \hyperref[def:p11-ktf-acceptable]{\ref*{def:p11-ktf-acceptable}} & \thepage & The native acceptable coding scheme uses the usual non-strict order on its standard numerical copy. \\
\phantomsection\label{entry:symbol-registry-03237}\(\iota\) & \texttt{\detokenize{\iota}} & injective code of every finite domain tuple & P11 acceptable-structure truth & n/a & \hyperref[def:p11-ktf-acceptable]{\ref*{def:p11-ktf-acceptable}} & \thepage & Includes the empty tuple; each fixed-arity tuple graph is elementary from sequence length and projections. \\
\phantomsection\label{entry:symbol-registry-03238}\(0_{\mathcal A},1_{\mathcal A},\ldots\) & \texttt{\detokenize{0_{\mathcal A},}\allowbreak\detokenize{1_{\mathcal A}}} & internal numerical domain markers & P11 acceptable-structure truth & n/a & \hyperref[def:p11-ktf-acceptable]{\ref*{def:p11-ktf-acceptable}} & \thepage & These domain objects and constructor tags are not semantic truth values; the empty tuple code need not be the numerical zero. \\
\phantomsection\label{entry:symbol-registry-03239}\(\mathsf{Seq}_{\mathcal A}\) & \texttt{\detokenize{\mathsf{Seq}_{\mathcal A}}} & elementary image predicate for finite tuple codes & P11 acceptable-structure truth & n/a & \hyperref[def:p11-ktf-acceptable]{\ref*{def:p11-ktf-acceptable}} & \thepage & The coding injection has all finite domain tuples as its input; arbitrary domain objects need not be sequence codes. \\
\phantomsection\label{entry:symbol-registry-03240}\(\ell_{\mathcal A}\) & \texttt{\detokenize{\ell_{\mathcal A}}} & elementary internal sequence length & P11 acceptable-structure truth & n/a & \hyperref[def:p11-ktf-acceptable]{\ref*{def:p11-ktf-acceptable}} & \thepage & Length is an object of the internal numerical copy; a nonsequence has the default length zero. \\
\phantomsection\label{entry:symbol-registry-03241}\(q_{\mathcal A}\) & \texttt{\detokenize{q_{\mathcal A}}} & elementary finite-sequence projection with numerical-zero default & P11 acceptable-structure truth & n/a & \hyperref[def:p11-ktf-acceptable]{\ref*{def:p11-ktf-acceptable}} & \thepage & Indices are one based; nonsequences and out-of-range projections return the internal numerical zero. \\
\phantomsection\label{entry:symbol-registry-03242}\(\operatorname{Ind}(\mathcal A)\) & \texttt{\detokenize{\operatorname{Ind}(\mathcal A)}} & parameter sections of positive elementary least fixed points & P11 acceptable-structure truth & n/a & \hyperref[def:p11-ktf-general-classes]{\ref*{def:p11-ktf-general-classes}} & \thepage & Positivity concerns the new inductive predicate only; ordinary predicates may appear negatively. \\
\phantomsection\label{entry:symbol-registry-03243}\(\operatorname{HE}(\mathcal A)\) & \texttt{\detokenize{\operatorname{HE}(\mathcal A)}} & inductive and coinductive relations on the structure & P11 acceptable-structure truth & n/a & \hyperref[def:p11-ktf-general-classes]{\ref*{def:p11-ktf-general-classes}} & \thepage & A relation and its complement must both be inductive with their permitted finite individual parameters. \\
\phantomsection\label{entry:symbol-registry-03244}\(\Pi^1_1(\mathcal A)\) & \texttt{\detokenize{\Pi^1_1(\mathcal A)}} & universal-relation definability over the fixed structure & P11 acceptable-structure truth & n/a & \hyperref[def:p11-ktf-general-classes]{\ref*{def:p11-ktf-general-classes}} & \thepage & Finite universal blocks of finite-arity relation variables followed by elementary matrices; not numerical lightface definability by assumption. \\
\phantomsection\label{entry:symbol-registry-03245}\(L_{\mathcal A}\) & \texttt{\detokenize{L_{\mathcal A}}} & name-rich ordinary language of the acceptable structure & P11 acceptable-structure truth & n/a & \hyperref[def:p11-ktf-general-language]{\ref*{def:p11-ktf-general-language}} & \thepage & Has a separate constant for every domain object; may be uncountable although formulas are finite. \\
\phantomsection\label{entry:symbol-registry-03246}\(\dot a\) & \texttt{\detokenize{\dot a}} & object-language constant naming a domain object & P11 acceptable-structure truth & n/a & \hyperref[def:p11-ktf-general-language]{\ref*{def:p11-ktf-general-language}} & \thepage & Constant-code payloads are opaque data, never recursively parsed as expressions. \\
\phantomsection\label{entry:symbol-registry-03247}\(K_{\mathcal A}^{\mathrm K},K_{\mathcal A}^{\mathrm V}\) & \texttt{\detokenize{K_{\mathcal A}^{\mathrm K},}\allowbreak\detokenize{K_{\mathcal A}^{\mathrm V}}} & least partial truth points in the general naming regime & P11 acceptable-structure truth & n/a & \hyperref[def:p11-ktf-general-language]{\ref*{def:p11-ktf-general-language}} & \thepage & Original wholly undefined truth-predicate stage zero; the second point uses unrestricted extensions of both coordinates. \\
\phantomsection\label{entry:symbol-registry-03248}\(\mathsf{Cl}_{\mathcal A}\) & \texttt{\detokenize{\mathsf{Cl}_{\mathcal A}}} & elementary total finite-assignment closing graph & P11 acceptable-structure truth & n/a & \hyperref[thm:p11-ktf-internal-syntax]{\ref*{thm:p11-ktf-internal-syntax}} & \thepage & Graph is proved by finite parsing and substitution histories; invalid inputs return a fixed false sentence code. \\
\phantomsection\label{entry:symbol-registry-03249}\(\mathsf V(e,s,b)\) & \texttt{\detokenize{\mathsf V(e,}\allowbreak\detokenize{s,}\allowbreak\detokenize{b)}} & positive signed formula-evaluation induction & P11 acceptable-structure truth & n/a & \hyperref[def:p11-ktf-elementary-interpreter]{\ref*{def:p11-ktf-elementary-interpreter}} & \thepage & Domain marker one means true and zero false; the original 1974 classical Val uses the opposite flag convention. \\
\phantomsection\label{entry:symbol-registry-03250}\(\mathsf Z(u,b)\) & \texttt{\detokenize{\mathsf Z(u,}\allowbreak\detokenize{b)}} & positive signed truth-coordinate induction & P11 acceptable-structure truth & n/a & \hyperref[def:p11-ktf-elementary-interpreter]{\ref*{def:p11-ktf-elementary-interpreter}} & \thepage & Copies evaluated sentence signs and supplies false nonsentence signs; packed with the formula interpreter into one positive relation. \\
\phantomsection\label{entry:symbol-registry-03251}\(\mathsf{Den}(t,s,a)\) & \texttt{\detokenize{\mathsf{Den}(t,}\allowbreak\detokenize{s,}\allowbreak\detokenize{a)}} & elementary term-denotation graph & P11 acceptable-structure truth & n/a & \hyperref[annex:p11.ktf.internal-syntax]{\ref*{annex:p11.ktf.internal-syntax}} & \thepage & Terms are variables and constants; finite assignments have numerical-zero defaults. \\
\phantomsection\label{entry:symbol-registry-03252}\(\mathsf{Upd}(s,i,a,s')\) & \texttt{\detokenize{\mathsf{Upd}(s,}\allowbreak\detokenize{i,}\allowbreak\detokenize{a,}\allowbreak\detokenize{s')}} & elementary unique assignment-update graph & P11 acceptable-structure truth & n/a & \hyperref[annex:p11.ktf.internal-syntax]{\ref*{annex:p11.ktf.internal-syntax}} & \thepage & Output length is the maximum of old length and the changed positive numerical index. \\
\phantomsection\label{entry:symbol-registry-03253}\(\mathsf{Sent}_{\mathcal A}\) & \texttt{\detokenize{\mathsf{Sent}_{\mathcal A}}} & elementary internally parsed sentence-code predicate & P11 acceptable-structure truth & n/a & \hyperref[annex:p11.ktf.internal-induction]{\ref*{annex:p11.ktf.internal-induction}} & \thepage & Finite postorder certificate and empty root free-index list establish the predicate; cyclic constructor equations alone do not suffice. \\
\phantomsection\label{entry:symbol-registry-03254}\(G(k,\vec z,u)\) & \texttt{\detokenize{G(k,}\allowbreak\detokenize{\vec z,}\allowbreak\detokenize{u)}} & elementary prospective-template closing graph & P11 acceptable-structure truth & n/a & \hyperref[annex:p11.ktf.acceptable-ind]{\ref*{annex:p11.ktf.acceptable-ind}} & \thepage & Closes template k under the assignment beginning with k itself; supplies actual diagonal sentences without defining the desired least relation as a premise. \\
\phantomsection\label{entry:symbol-registry-03255}\(P,Q\) & \texttt{\detokenize{P,}\allowbreak\detokenize{Q}} & one-sign elementary-control games in a selective shuffle & P11 acceptable-structure truth & n/a & \hyperref[thm:p11-ktf-selective-shuffle]{\ref*{thm:p11-ktf-selective-shuffle}} & \thepage & P only accepts or diverges; Q only rejects or diverges; the alternating state retains disjoint component registers. \\
\phantomsection\label{entry:symbol-registry-03256}\(\infty\) & \texttt{\detokenize{\infty}} & unresolved-game rank comparison marker & P11 acceptable-structure truth & n/a & \hyperref[def:p11-ktf-elementary-games]{\ref*{def:p11-ktf-elementary-games}} & \thepage & Placed above all proper resolution ranks for comparisons; not an ordinal stage or a semantic truth value. \\
\phantomsection\label{entry:symbol-registry-03257}\(\mathsf{CV}_{\mathcal A}(R,V)\) & \texttt{\detokenize{\mathsf{CV}_{\mathcal A}(R,}\allowbreak\detokenize{V)}} & first-order matrix verifying a complete classical formula valuation & P11 acceptable-structure truth & n/a & \hyperref[def:p11-ktf-general-classical-valuation]{\ref*{def:p11-ktf-general-classical-valuation}} & \thepage & Unary R interprets T and ternary V evaluates all coded formulas and finite assignments; not absolute elementary satisfaction. \\
\phantomsection\label{entry:symbol-registry-03258}\(\tau_a\) & \texttt{\detokenize{\tau_a}} & independent name-rich family of internally coded truth-tellers & P11 acceptable-structure truth & n/a & \hyperref[def:p11-ktf-general-tellers]{\ref*{def:p11-ktf-general-tellers}} & \thepage & Distinct code for each domain payload; all remain undefined at the least unrestricted-supervaluation point. \\
\phantomsection\label{entry:symbol-registry-03259}\(\mathsf{Ext}(X,R)\) & \texttt{\detokenize{\mathsf{Ext}(X,}\allowbreak\detokenize{R)}} & first-order two-coordinate classical-extension condition & P11 acceptable-structure truth & n/a & \hyperref[annex:p11.ktf.general-sv]{\ref*{annex:p11.ktf.general-sv}} & \thepage & Signed X may be inconsistent in the ambient powerset; then no total extension exists. \\
\phantomsection\label{entry:symbol-registry-03260}\(\mathsf{Pre}_{\mathcal A}(X,Y)\) & \texttt{\detokenize{\mathsf{Pre}_{\mathcal A}(X,}\allowbreak\detokenize{Y)}} & first-order packed general pre-fixed countermodel matrix & P11 acceptable-structure truth & n/a & \hyperref[annex:p11.ktf.general-sv]{\ref*{annex:p11.ktf.general-sv}} & \thepage & Six-ary Y packs countermodels and classical valuations indexed by missing signed codes; equivalence to pre-fixedness explicitly uses choice. \\
\phantomsection\label{entry:symbol-registry-03261}\(p_j(\vec z)\) & \texttt{\detokenize{p_j(\vec z)}} & disjoint elementary relation-payload tuple code & P11 acceptable-structure truth & n/a & \hyperref[annex:p11.ktf.general-sv]{\ref*{annex:p11.ktf.general-sv}} & \thepage & Distinct internal tags encode finite-arity relation tuples for the independent truth-teller family, including zero arity. \\
\phantomsection\label{entry:symbol-registry-03262}\(\theta_{\vec a},\theta_z\) & \texttt{\detokenize{\theta_{\vec a},}\allowbreak\detokenize{\theta_z}} & actual generic diagonal induction and game sentences & P11 acceptable-structure truth & n/a & \hyperref[annex:p11.ktf.acceptable-ind]{\ref*{annex:p11.ktf.acceptable-ind}} & \thepage & Named prospective-template parameter is closed to its own fixed code; retained argument tautologies ensure distinct payload strings. \\
\phantomsection\label{entry:symbol-registry-03263}\(H\) & \texttt{\detokenize{H}} & ambient typed signed-powerset supervaluation operator & P11 acceptable-structure truth & n/a & \hyperref[annex:p11.ktf.general-sv]{\ref*{annex:p11.ktf.general-sv}} & \thepage & May take inconsistent inputs in the complete ambient lattice; its consistent bottom iteration is the original partial truth iteration. \\
\phantomsection\label{entry:symbol-registry-03264}\(u_r,v_r\) & \texttt{\detokenize{u_r,}\allowbreak\detokenize{v_r}} & elementary total translations between finite acceptable-structure Tarski languages & P11 acceptable-structure truth & n/a & \hyperref[annex:p11.ktf.general-finite-hierarchy]{\ref*{annex:p11.ktf.general-finite-hierarchy}} & \thepage & Graphs manipulate parsed formulas with opaque domain constants; invalid inputs return the variable-term code rather than an assumed numerical noncode. \\
\phantomsection\label{entry:symbol-registry-03450}\(U\) & \texttt{\detokenize{U}} & real-valued prospect representation in the formal-unification survey & P11 practical-and-theoretical-rationality & meta & \hyperref[thm:p11-rat22-numerical-order]{\ref*{thm:p11-rat22-numerical-order}} & \thepage & Editorial name for the survey's utility function; higher values exactly represent strict preferences and equal values exactly represent indifference. No existence or expected-utility interpretation follows from this row. \\
\phantomsection\label{entry:symbol-registry-03451}\(A,B,C\) & \texttt{\detokenize{A,}\allowbreak\detokenize{B,}\allowbreak\detokenize{C}} & prospects in a scalar preference representation & P11 practical-and-theoretical-rationality & meta & \hyperref[thm:p11-rat22-numerical-order]{\ref*{thm:p11-rat22-numerical-order}} & \thepage & Wedgwood's source roles for A B C at printed 140; these are decision prospects rather than object-language formulas or propositions. \\
\phantomsection\label{entry:symbol-registry-03452}\(p\) & \texttt{\detokenize{p}} & schematic proposition assessed as a belief content & P11 practical-and-theoretical-rationality & meta & \hyperref[def:p11-rat22-local-goals]{\ref*{def:p11-rat22-local-goals}} & \thepage & Used in the compelling-evidence case and the graded-correctness proposal; not a probability mass or a prospect. \\
\phantomsection\label{entry:symbol-registry-03453}\(P\) & \texttt{\detokenize{P}} & supposed rational probability function for an agent at a time & P11 practical-and-theoretical-rationality & meta & \hyperref[def:p11-rat22-correctness]{\ref*{def:p11-rat22-correctness}} & \thepage & Wedgwood's symbol at printed 143; its existence uniqueness approximation metric and expected-correctness bridge are not proved in the selected chapter. Not automatically the agent's actual credences or a preference-derived probability. \\
\phantomsection\label{entry:symbol-registry-03511}\(L_{\mathrm e},L_{\mathrm e,T}\) & \texttt{\detokenize{L_{\mathrm e},}\allowbreak\detokenize{L_{\mathrm e,T}}} & fixed empirical arithmetic base and its unary truth expansion & P11 recursive truth hierarchies & n/a & \hyperref[def:p11-ktf-recursive-hierarchy-base]{\ref*{def:p11-ktf-recursive-hierarchy-base}} & \thepage & Standard arithmetic and a finite ordinary predicate signature; syntax is held fixed while empirical interpretations vary. \\
\phantomsection\label{entry:symbol-registry-03512}\(R_1,\ldots,R_l\) & \texttt{\detokenize{R_1,}\allowbreak\detokenize{\ldots,}\allowbreak\detokenize{R_l}} & finite ordinary empirical predicate vocabulary & P11 recursive truth hierarchies & n/a & \hyperref[def:p11-ktf-recursive-hierarchy-base]{\ref*{def:p11-ktf-recursive-hierarchy-base}} & \thepage & Predicate arities are fixed; their extensions need not be computable. \\
\phantomsection\label{entry:symbol-registry-03513}\(\mathfrak B\) & \texttt{\detokenize{\mathfrak B}} & specified nonempty empirical interpretation family & P11 recursive truth hierarchies & n/a & \hyperref[def:p11-ktf-recursive-hierarchy-base]{\ref*{def:p11-ktf-recursive-hierarchy-base}} & \thepage & A set of interpretations on the fixed numerical domain; a relative safety claim does not silently quantify over a larger family. \\
\phantomsection\label{entry:symbol-registry-03514}\(c_T\) & \texttt{\detokenize{c_T}} & computable finite-word code for the original empirical truth language & P11 recursive truth hierarchies & n/a & \hyperref[def:p11-ktf-recursive-hierarchy-base]{\ref*{def:p11-ktf-recursive-hierarchy-base}} & \thepage & Separate explicit coding regime with zero a nonword code; all ordinary syntax operations have arithmetic graphs. \\
\phantomsection\label{entry:symbol-registry-03515}\(S_{\mathcal M}^{\alpha}\) & \texttt{\detokenize{S_{\mathcal M}^{\alpha}}} & original bottom Strong-Kleene stage over an empirical interpretation & P11 recursive truth hierarchies & n/a & \hyperref[def:p11-ktf-recursive-hierarchy-base]{\ref*{def:p11-ktf-recursive-hierarchy-base}} & \thepage & Both truth coordinates start empty; nonsentence inputs become negative after jump one. \\
\phantomsection\label{entry:symbol-registry-03516}\(W=(B,\prec_W)\) & \texttt{\detokenize{W=(B,}\allowbreak\detokenize{\prec_W)}} & recursive well-order on a recursive numerical domain & P11 recursive truth hierarchies & n/a & \hyperref[def:p11-ktf-recursive-hierarchy-base]{\ref*{def:p11-ktf-recursive-hierarchy-base}} & \thepage & Zero is the least element; the order relation is false outside its domain and is not the branching order on all ordinal notations. \\
\phantomsection\label{entry:symbol-registry-03517}\(\alpha_W(b)\) & \texttt{\detokenize{\alpha_W(b)}} & ordinal rank of an index in the fixed recursive well-order & P11 recursive truth hierarchies & n/a & \hyperref[def:p11-ktf-recursive-hierarchy-base]{\ref*{def:p11-ktf-recursive-hierarchy-base}} & \thepage & Order type of the strict predecessor segment; each represented ordinal is recursive and below the Church--Kleene ordinal. \\
\phantomsection\label{entry:symbol-registry-03518}\(\mathsf{Tr}_b\) & \texttt{\detokenize{\mathsf{Tr}_b}} & classical truth token at a nonzero completed hierarchy index & P11 recursive truth hierarchies & n/a & \hyperref[def:p11-ktf-completed-hierarchy]{\ref*{def:p11-ktf-completed-hierarchy}} & \thepage & True exactly on true sentences of the strictly lower indexed language and false on every other number; also introduced at limit indices. \\
\phantomsection\label{entry:symbol-registry-03519}\(\mathcal H_{<b},\mathcal H_b\) & \texttt{\detokenize{\mathcal H_{<b},}\allowbreak\detokenize{\mathcal H_b}} & prior union language and completed language at an indexed level & P11 recursive truth hierarchies & n/a & \hyperref[def:p11-ktf-completed-hierarchy]{\ref*{def:p11-ktf-completed-hierarchy}} & \thepage & The completed language includes its new truth token; a formula remains a finite word even at an infinite level. \\
\phantomsection\label{entry:symbol-registry-03520}\(c_W\) & \texttt{\detokenize{c_W}} & tagged finite-word code for the indexed auxiliary language & P11 recursive truth hierarchies & n/a & \hyperref[def:p11-ktf-completed-hierarchy]{\ref*{def:p11-ktf-completed-hierarchy}} & \thepage & Numerical predicate indices form part of finite token encodings; syntax guards use the fixed recursive order. \\
\phantomsection\label{entry:symbol-registry-03521}\(\delta_W\) & \texttt{\detokenize{\delta_W}} & uniform capture-free hierarchy embedding into unary partial truth & P11 recursive truth hierarchies & n/a & \hyperref[thm:p11-ktf-transfinite-embedding]{\ref*{thm:p11-ktf-transfinite-embedding}} & \thepage & One actual computable translation for the whole indexed language; lower token guards and self-template closing implement it. \\
\phantomsection\label{entry:symbol-registry-03522}\(\mathsf P_b,\mathsf N_b\) & \texttt{\detokenize{\mathsf P_b,}\allowbreak\detokenize{\mathsf N_b}} & classical auxiliary formulas defining signed original truth-stage coordinates & P11 recursive truth hierarchies & n/a & \hyperref[thm:p11-ktf-stage-compiler]{\ref*{thm:p11-ktf-stage-compiler}} & \thepage & Their total finite-word compiler is constructed first; their exact equality to stage coordinates is a later proved result rather than a primitive premise. \\
\phantomsection\label{entry:symbol-registry-03523}\(C_W^+(b,F),C_W^-(b,F)\) & \texttt{\detokenize{C_W^+(b,}\allowbreak\detokenize{F),}\allowbreak\detokenize{C_W^-(b,}\allowbreak\detokenize{F)}} & positive and negative finite-word formula compilations & P11 recursive truth hierarchies & n/a & \hyperref[thm:p11-ktf-stage-compiler]{\ref*{thm:p11-ktf-stage-compiler}} & \thepage & Classical conditions test the corresponding Strong-Kleene signs at the represented stage; both are false for a gap. \\
\phantomsection\label{entry:symbol-registry-03524}\(\mathsf{Sent}_T\) & \texttt{\detokenize{\mathsf{Sent}_T}} & arithmetic original-language sentence-code guard & P11 recursive truth hierarchies & n/a & \hyperref[thm:p11-ktf-stage-compiler]{\ref*{thm:p11-ktf-stage-compiler}} & \thepage & Validates actual finite original words independently of every empirical predicate interpretation. \\
\phantomsection\label{entry:symbol-registry-03525}\(G_e(\sigma,d,x,y)\) & \texttt{\detokenize{G_e(\sigma,}\allowbreak\detokenize{d,}\allowbreak\detokenize{x,}\allowbreak\detokenize{y)}} & arithmetic graph of the self-indexed signed word compiler & P11 recursive truth hierarchies & n/a & \hyperref[thm:p11-ktf-stage-compiler]{\ref*{thm:p11-ktf-stage-compiler}} & \thepage & The actual program e is total because its equal-index factory programme terminates by finite parse recursion. \\
\phantomsection\label{entry:symbol-registry-03526}\(\sigma\) & \texttt{\detokenize{\sigma}} & numerical selector of a signed compilation output & P11 recursive truth hierarchies & n/a & \hyperref[thm:p11-ktf-stage-compiler]{\ref*{thm:p11-ktf-stage-compiler}} & \thepage & One selects the positive word and zero the negative word; this local programme input is distinct from the earlier omega-rank sentence of the same glyph. \\
\phantomsection\label{entry:symbol-registry-03527}\(\mathsf{Sent}_{<W}(b,e)\) & \texttt{\detokenize{\mathsf{Sent}_{<W}(b,}\allowbreak\detokenize{e)}} & uniform arithmetic strict-lower auxiliary sentence guard & P11 recursive truth hierarchies & n/a & \hyperref[annex:p11.ktf.transfinite-embedding]{\ref*{annex:p11.ktf.transfinite-embedding}} & \thepage & Checks every finitely occurring indexed predicate token against the predecessor relation. \\
\phantomsection\label{entry:symbol-registry-03528}\(\operatorname{tr}(k,e)\) & \texttt{\detokenize{\operatorname{tr}(k,}\allowbreak\detokenize{e)}} & total prospective-template word translation & P11 recursive truth hierarchies & n/a & \hyperref[annex:p11.ktf.transfinite-embedding]{\ref*{annex:p11.ktf.transfinite-embedding}} & \thepage & Copies the prospective template without recursively retranslating it; invalid inputs return a fixed false base word. \\
\phantomsection\label{entry:symbol-registry-03529}\(E(k,e,y)\) & \texttt{\detokenize{E(k,}\allowbreak\detokenize{e,}\allowbreak\detokenize{y)}} & arithmetic unique-output graph of prospective-template translation & P11 recursive truth hierarchies & n/a & \hyperref[annex:p11.ktf.transfinite-embedding]{\ref*{annex:p11.ktf.transfinite-embedding}} & \thepage & Graph construction precedes interpretation of the original or auxiliary truth predicates. \\
\phantomsection\label{entry:symbol-registry-03530}\(Q(k,b,x)\) & \texttt{\detokenize{Q(k,}\allowbreak\detokenize{b,}\allowbreak\detokenize{x)}} & three-slot guarded original truth-language embedding template & P11 recursive truth hierarchies & n/a & \hyperref[annex:p11.ktf.transfinite-embedding]{\ref*{annex:p11.ktf.transfinite-embedding}} & \thepage & Checks the index and lower-language syntax before reading T at the translated word code. \\
\phantomsection\label{entry:symbol-registry-03531}\(q,D_b\) & \texttt{\detokenize{q,}\allowbreak\detokenize{D_b}} & fixed embedding-template code and its closed indexed instances & P11 recursive truth hierarchies & n/a & \hyperref[annex:p11.ktf.transfinite-embedding]{\ref*{annex:p11.ktf.transfinite-embedding}} & \thepage & q is the actual code of Q; D\_b substitutes q and b into its designated slots and retains the numerical input variable. \\
\phantomsection\label{entry:symbol-registry-03532}\(p(i,b),n(i,b)\) & \texttt{\detokenize{p(i,}\allowbreak\detokenize{b),}\allowbreak\detokenize{n(i,}\allowbreak\detokenize{b)}} & total factories for positive and negative coordinate formula words & P11 recursive truth hierarchies & n/a & \hyperref[annex:p11.ktf.stage-compiler]{\ref*{annex:p11.ktf.stage-compiler}} & \thepage & The prospective programme index i is inserted only as an arithmetic graph word; no programme execution or truth-stage query is needed. \\
\phantomsection\label{entry:symbol-registry-03533}\(f(i),e\) & \texttt{\detokenize{f(i),}\allowbreak\detokenize{e}} & compiler programme-index factory and its self-reference index & P11 recursive truth hierarchies & n/a & \hyperref[annex:p11.ktf.stage-compiler]{\ref*{annex:p11.ktf.stage-compiler}} & \thepage & The independent P4 programme self-reference lemma gives equal computed functions; totality is proved by finite syntax recursion. \\
\phantomsection\label{entry:symbol-registry-03534}\(r_{\mathcal M}(F,g)\) & \texttt{\detokenize{r_{\mathcal M}(F,}\allowbreak\detokenize{g)}} & first defined evaluation rank for a grounded formula instance & P11 recursive truth hierarchies & n/a & \hyperref[def:p11-ktf-intrinsic-profile]{\ref*{def:p11-ktf-intrinsic-profile}} & \thepage & Assignments on the finitely many free variables are numerical tuples; this rank is countable in each fixed interpretation. \\
\phantomsection\label{entry:symbol-registry-03535}\(\ell_{\mathfrak B}(F)\) & \texttt{\detokenize{\ell_{\mathfrak B}(F)}} & intrinsic uniform evaluation bound relative to an empirical family & P11 recursive truth hierarchies & n/a & \hyperref[def:p11-ktf-intrinsic-profile]{\ref*{def:p11-ktf-intrinsic-profile}} & \thepage & Defined only for safe formulas as the supremum of their first ranks; it is at most omega1 and need not be represented by a recursive well-order. \\
\phantomsection\label{entry:symbol-registry-03536}\(\sim_{\mathfrak B}\) & \texttt{\detokenize{\sim_{\mathfrak B}}} & semantic equivalence over every empirical interpretation and assignment & P11 recursive truth hierarchies & meta & \hyperref[def:p11-ktf-intrinsic-profile]{\ref*{def:p11-ktf-intrinsic-profile}} & \thepage & Compares least-fixed-point values of safe original formulas and classical values of auxiliary formulas; ranks are not invariants of this relation. \\
\phantomsection\label{entry:symbol-registry-03537}\(\omega_1\) & \texttt{\detokenize{\omega_1}} & first uncountable ordinal bounding all relative numerical evaluation ranks & P11 recursive truth hierarchies & n/a & \hyperref[def:p11-ktf-intrinsic-profile]{\ref*{def:p11-ktf-intrinsic-profile}} & \thepage & The supremum over an arbitrary interpretation family can equal it even though every individual rank is countable; distinct from the Church--Kleene ordinal. \\
\phantomsection\label{entry:symbol-registry-03538}\(\mathfrak B_{\mathrm{wf}}\) & \texttt{\detokenize{\mathfrak B_{\mathrm{wf}}}} & family of all well-founded empirical numerical tree interpretations & P11 recursive truth hierarchies & n/a & \hyperref[thm:p11-ktf-intrinsic-scope]{\ref*{thm:p11-ktf-intrinsic-scope}} & \thepage & Well-foundedness is a consequential restriction on allowed empirical facts; adding a branch interpretation destroys the example sentence's safety. \\
\phantomsection\label{entry:symbol-registry-03539}\(\eta,\eta_s\) & \texttt{\detokenize{\eta,}\allowbreak\detokenize{\eta_s}} & one fixed safe root sentence and its internally closed node family & P11 recursive truth hierarchies & n/a & \hyperref[thm:p11-ktf-intrinsic-scope]{\ref*{thm:p11-ktf-intrinsic-scope}} & \thepage & The root ranges over no model-dependent syntax; its evaluation ranks across the well-founded family are cofinal in omega1. \\
\phantomsection\label{entry:symbol-registry-03540}\(G(k,t,u)\) & \texttt{\detokenize{G(k,}\allowbreak\detokenize{t,}\allowbreak\detokenize{u)}} & arithmetic prospective-template closing graph for empirical tree sentences & P11 recursive truth hierarchies & n/a & \hyperref[annex:p11.ktf.intrinsic-scope]{\ref*{annex:p11.ktf.intrinsic-scope}} & \thepage & Closes two designated template slots at its own numerical code and a node numeral; does not define the empirical child relation arithmetically. \\
\phantomsection\label{entry:symbol-registry-03541}\(\preceq_W\) & \texttt{\detokenize{\preceq_W}} & nonstrict indexed hierarchy order & P11 recursive truth hierarchies & meta & \hyperref[def:p11-ktf-recursive-hierarchy-base]{\ref*{def:p11-ktf-recursive-hierarchy-base}} & \thepage & Requires both indices in the recursive domain and then either equality or strict predecessorhood. \\
\phantomsection\label{entry:symbol-registry-03542}\(d_W,D_W\) & \texttt{\detokenize{d_W,}\allowbreak\detokenize{D_W}} & total auxiliary diagonal word map and its arithmetic graph & P11 recursive truth hierarchies & n/a & \hyperref[annex:p11.ktf.limit-separation]{\ref*{annex:p11.ktf.limit-separation}} & \thepage & Finite-word self-substitution in the indexed auxiliary language supports the actual truth-undefinability proof; invalid inputs return a fixed noncode. \\
\phantomsection\label{entry:symbol-registry-03543}\(\kappa,\kappa^+\) & \texttt{\detokenize{\kappa,}\allowbreak\detokenize{\kappa^+}} & fixed acceptable code-domain cardinal and its successor cardinal ordinal & P11 semantic-order truth hierarchies & n/a & \hyperref[def:p11-ktf-semantic-order-base]{\ref*{def:p11-ktf-semantic-order-base}} & \thepage & Metatheory is classical set theory with choice; individual truth iterations stabilize below the successor and a codable index order has smaller order type. \\
\phantomsection\label{entry:symbol-registry-03544}\(I,o,\prec_I\) & \texttt{\detokenize{I,}\allowbreak\detokenize{o,}\allowbreak\detokenize{\prec_I}} & codable index subset with a named least element and strict semantic well-order & P11 semantic-order truth hierarchies & n/a & \hyperref[def:p11-ktf-semantic-order-base]{\ref*{def:p11-ktf-semantic-order-base}} & \thepage & Order predicates must be elementary in the base or explicitly supplied as an ordinary finite expansion; their extensions need not be decidable. \\
\phantomsection\label{entry:symbol-registry-03545}\(\preceq_I\) & \texttt{\detokenize{\preceq_I}} & nonstrict semantic index order & P11 semantic-order truth hierarchies & meta & \hyperref[def:p11-ktf-semantic-order-base]{\ref*{def:p11-ktf-semantic-order-base}} & \thepage & Requires both indices in the supplied index subset and either equality or strict predecessorhood. \\
\phantomsection\label{entry:symbol-registry-03546}\(\beta_I(b)\) & \texttt{\detokenize{\beta_I(b)}} & set-ordinal rank of an opaque semantic-order index & P11 semantic-order truth hierarchies & n/a & \hyperref[def:p11-ktf-semantic-order-base]{\ref*{def:p11-ktf-semantic-order-base}} & \thepage & Order type of its strict predecessor segment; not restricted to recursive ordinal presentations. \\
\phantomsection\label{entry:symbol-registry-03547}\(\mathfrak F\) & \texttt{\detokenize{\mathfrak F}} & held-coding family of ordinary acceptable-base interpretations & P11 semantic-order truth hierarchies & n/a & \hyperref[def:p11-ktf-semantic-order-base]{\ref*{def:p11-ktf-semantic-order-base}} & \thepage & Domain coding names and index order remain fixed for the common-bound theorem; other ordinary predicates may vary while preserving the specified coding formulas. \\
\phantomsection\label{entry:symbol-registry-03548}\(L_*,L_{*,T}\) & \texttt{\detokenize{L_*,}\allowbreak\detokenize{L_{*,T}}} & name-rich ordinary semantic-order base and its unary partial-truth expansion & P11 semantic-order truth hierarchies & n/a & \hyperref[def:p11-ktf-semantic-order-base]{\ref*{def:p11-ktf-semantic-order-base}} & \thepage & The additional ordinary index predicates are counted as an expansion when not already elementary in the original base. \\
\phantomsection\label{entry:symbol-registry-03549}\(\mathsf{Tr}^{I}_b\) & \texttt{\detokenize{\mathsf{Tr}^{I}_b}} & classical truth token with opaque index payload & P11 semantic-order truth hierarchies & n/a & \hyperref[def:p11-ktf-semantic-order-hierarchy]{\ref*{def:p11-ktf-semantic-order-hierarchy}} & \thepage & Assigned full truth of the strict-lower finite-expression language at every nonzero index including limits; false on all other domain objects. \\
\phantomsection\label{entry:symbol-registry-03550}\(\mathcal J_{<b},\mathcal J_b\) & \texttt{\detokenize{\mathcal J_{<b},}\allowbreak\detokenize{\mathcal J_b}} & strict-lower language and its completed opaque-index hierarchy level & P11 semantic-order truth hierarchies & n/a & \hyperref[def:p11-ktf-semantic-order-hierarchy]{\ref*{def:p11-ktf-semantic-order-hierarchy}} & \thepage & Both keep the same ordinary name-rich base; the completed level adds its new full truth token. \\
\phantomsection\label{entry:symbol-registry-03551}\(c_I\) & \texttt{\detokenize{c_I}} & opaque-index auxiliary formula coding by the acceptable tuple scheme & P11 semantic-order truth hierarchies & n/a & \hyperref[def:p11-ktf-semantic-order-hierarchy]{\ref*{def:p11-ktf-semantic-order-hierarchy}} & \thepage & The indexed prime has constructor tag eight with index and term-code payloads; indices and constant payloads are not recursively parsed. \\
\phantomsection\label{entry:symbol-registry-03552}\(\theta\) & \texttt{\detokenize{\theta}} & set order type of the supplied semantic index order & P11 semantic-order truth hierarchies & n/a & \hyperref[def:p11-ktf-semantic-order-hierarchy]{\ref*{def:p11-ktf-semantic-order-hierarchy}} & \thepage & Its cardinality is at most that of the fixed code domain; it is not a proper-class all-ordinal hierarchy. \\
\phantomsection\label{entry:symbol-registry-03553}\(\mathsf{Sent}^{I}_{<},\mathsf{Sent}^{I}_{\leq}\) & \texttt{\detokenize{\mathsf{Sent}^{I}_{<},}\allowbreak\detokenize{\mathsf{Sent}^{I}_{\leq}}} & elementary order-dependent lower-language sentence legality guards & P11 semantic-order truth hierarchies & n/a & \hyperref[thm:p11-ktf-semantic-order-syntax]{\ref*{thm:p11-ktf-semantic-order-syntax}} & \thepage & Finite postorder certificates check each indexed prime by ordinary semantic order relations; raw parsing does not need to decide those extensions. \\
\phantomsection\label{entry:symbol-registry-03554}\(E_I(k,e,y)\) & \texttt{\detokenize{E_I(k,}\allowbreak\detokenize{e,}\allowbreak\detokenize{y)}} & elementary total graph of prospective opaque-index embedding translation & P11 semantic-order truth hierarchies & n/a & \hyperref[thm:p11-ktf-semantic-order-syntax]{\ref*{thm:p11-ktf-semantic-order-syntax}} & \thepage & Copies a three-slot original template at an indexed prime using named prospective code and index; does not translate the inserted copy again. \\
\phantomsection\label{entry:symbol-registry-03555}\(\mathsf C_I(k,\sigma,b,e,y)\) & \texttt{\detokenize{\mathsf C_I(k,}\allowbreak\detokenize{\sigma,}\allowbreak\detokenize{b,}\allowbreak\detokenize{e,}\allowbreak\detokenize{y)}} & elementary total graph of prospective signed condition compilation & P11 semantic-order truth hierarchies & n/a & \hyperref[thm:p11-ktf-semantic-order-syntax]{\ref*{thm:p11-ktf-semantic-order-syntax}} & \thepage & Finite parse certificates verify paired outputs and literal copies of a prospective ordinary relation-template code; no stage relation or numerical self-index theorem is assumed. \\
\phantomsection\label{entry:symbol-registry-03556}\(K_k\) & \texttt{\detokenize{K_k}} & syntactic instance of a prospective ordinary five-slot relation template & P11 semantic-order truth hierarchies & n/a & \hyperref[thm:p11-ktf-semantic-order-syntax]{\ref*{thm:p11-ktf-semantic-order-syntax}} & \thepage & This denotes a copied formula word with its own code named in the first slot; it is not an added semantic relation symbol. \\
\phantomsection\label{entry:symbol-registry-03557}\(p_I(k,b,x),n_I(k,b,x)\) & \texttt{\detokenize{p_I(k,}\allowbreak\detokenize{b,}\allowbreak\detokenize{x),}\allowbreak\detokenize{n_I(k,}\allowbreak\detokenize{b,}\allowbreak\detokenize{x)}} & opaque-index signed coordinate factory words & P11 semantic-order truth hierarchies & n/a & \hyperref[thm:p11-ktf-semantic-order-syntax]{\ref*{thm:p11-ktf-semantic-order-syntax}} & \thepage & Formed for every domain index using named payloads and ordinary predecessor guard words; the least index uses the constantly false pair. \\
\phantomsection\label{entry:symbol-registry-03558}\(\mathsf{Sent}_{*,T}\) & \texttt{\detokenize{\mathsf{Sent}_{*,T}}} & elementary original name-rich sentence-code guard in the expanded base & P11 semantic-order truth hierarchies & n/a & \hyperref[thm:p11-ktf-semantic-order-syntax]{\ref*{thm:p11-ktf-semantic-order-syntax}} & \thepage & Distinguishes original sentence objects from nonsentences before the signed coordinate tests. \\
\phantomsection\label{entry:symbol-registry-03559}\(\Delta_I\) & \texttt{\detokenize{\Delta_I}} & elementary-code hierarchy embedding for semantic well-orders & P11 semantic-order truth hierarchies & n/a & \hyperref[thm:p11-ktf-semantic-order-embedding]{\ref*{thm:p11-ktf-semantic-order-embedding}} & \thepage & A single prospective-template translation preserves values under every assignment at the ordinal rank of its supplied index. \\
\phantomsection\label{entry:symbol-registry-03560}\(h\) & \texttt{\detokenize{h}} & named actual ordinary formula code of the signed compiler graph & P11 semantic-order truth hierarchies & n/a & \hyperref[thm:p11-ktf-semantic-order-stages]{\ref*{thm:p11-ktf-semantic-order-stages}} & \thepage & Selected after the graph construction; copying its formula at its own named code is finite syntax and does not recursively interpret that formula. \\
\phantomsection\label{entry:symbol-registry-03561}\(P^I_b,N^I_b\) & \texttt{\detokenize{P^I_b,}\allowbreak\detokenize{N^I_b}} & actual classical formulas for signed original truth-stage coordinates & P11 semantic-order truth hierarchies & n/a & \hyperref[thm:p11-ktf-semantic-order-stages]{\ref*{thm:p11-ktf-semantic-order-stages}} & \thepage & The exact stage equality is proved by well-order induction after construction of their total prospective compiler. \\
\phantomsection\label{entry:symbol-registry-03562}\(U_b^+(F),U_b^-(F)\) & \texttt{\detokenize{U_b^+(F),}\allowbreak\detokenize{U_b^-(F)}} & actual positive and negative auxiliary condition outputs & P11 semantic-order truth hierarchies & n/a & \hyperref[thm:p11-ktf-semantic-order-stages]{\ref*{thm:p11-ktf-semantic-order-stages}} & \thepage & True exactly for the corresponding Strong-Kleene formula sign at the represented stage; both false for a gap. \\
\phantomsection\label{entry:symbol-registry-03563}\(\ell_{\mathfrak F}(F)\) & \texttt{\detokenize{\ell_{\mathfrak F}(F)}} & supremum intrinsic evaluation bound over a held acceptable interpretation family & P11 semantic-order truth hierarchies & n/a & \hyperref[thm:p11-ktf-semantic-order-converse]{\ref*{thm:p11-ktf-semantic-order-converse}} & \thepage & Uses finite domain assignments and the previously explicit relative safety convention; does not claim the historical outline supplied that convention. \\
\phantomsection\label{entry:symbol-registry-03564}\(D_I,\lambda_I\) & \texttt{\detokenize{D_I,}\allowbreak\detokenize{\lambda_I}} & elementary auxiliary diagonal closing graph and its actual liar sentence & P11 semantic-order truth hierarchies & n/a & \hyperref[annex:p11.ktf.semantic-order-embedding]{\ref*{annex:p11.ktf.semantic-order-embedding}} & \thepage & Establishes undefinability of full truth in its own prior completed-hierarchy language using finite opaque-name syntax. \\
\phantomsection\label{entry:symbol-registry-03565}\(Q_I,q_I\) & \texttt{\detokenize{Q_I,}\allowbreak\detokenize{q_I}} & three-slot original embedding template and its named code & P11 semantic-order truth hierarchies & n/a & \hyperref[annex:p11.ktf.semantic-order-embedding]{\ref*{annex:p11.ktf.semantic-order-embedding}} & \thepage & Lower-language semantic guards precede the unary truth query at the unique prospective translation output. \\
\phantomsection\label{entry:symbol-registry-03566}\(\sigma\) & \texttt{\detokenize{\sigma}} & internal-domain selector marker for an opaque signed compilation & P11 semantic-order truth hierarchies & n/a & \hyperref[thm:p11-ktf-semantic-order-syntax]{\ref*{thm:p11-ktf-semantic-order-syntax}} & \thepage & Its allowed values are the named internal markers one\_A and zero\_A; these domain objects are distinct from semantic truth values and from the earlier omega-rank sentence. \\
\phantomsection\label{entry:symbol-registry-03567}\(S_{\mathcal M}^{\alpha}\) & \texttt{\detokenize{S_{\mathcal M}^{\alpha}}} & original Strong-Kleene stages over a held acceptable-base interpretation & P11 semantic-order truth hierarchies & n/a & \hyperref[def:p11-ktf-semantic-order-base]{\ref*{def:p11-ktf-semantic-order-base}} & \thepage & Uses the full name-rich code domain rather than numerical-only syntax; each held interpretation has its own ordinal stage sequence. \\
\phantomsection\label{entry:symbol-registry-03568}\(D_b\) & \texttt{\detokenize{D_b}} & opaque-index instance of the actual original embedding template & P11 semantic-order truth hierarchies & n/a & \hyperref[annex:p11.ktf.semantic-order-embedding]{\ref*{annex:p11.ktf.semantic-order-embedding}} & \thepage & Closes the named template code and index while retaining its domain input argument; its value bounds use the supplied semantic-order rank. \\
\phantomsection\label{entry:symbol-registry-03569}\(c_0,c_1\) & \texttt{\detokenize{c_0,}\allowbreak\detokenize{c_1}} & compatible source and ordinary-expansion target formula codes & P11 semantic-order truth hierarchies & n/a & \hyperref[thm:p11-ktf-expansion-transport]{\ref*{thm:p11-ktf-expansion-transport}} & \thepage & The held tuple scheme and old predicate indices are retained; the target adds finitely many ordinary predicate cases. \\
\phantomsection\label{entry:symbol-registry-03570}\(\Lambda\) & \texttt{\detokenize{\Lambda}} & stage-preserving source-to-expansion formula recoding & P11 semantic-order truth hierarchies & n/a & \hyperref[thm:p11-ktf-expansion-transport]{\ref*{thm:p11-ktf-expansion-transport}} & \thepage & Routes old truth atoms through the actual translated source code graph; does not assume unchanged strings keep their old truth values after vocabulary expansion. \\
\phantomsection\label{entry:symbol-registry-03571}\(\mathsf R(k,e,y),\mu_R(a)\) & \texttt{\detokenize{\mathsf R(k,}\allowbreak\detokenize{e,}\allowbreak\detokenize{y),}\allowbreak\detokenize{\mu_R(a)}} & total elementary prospective transport compiler graph and its selected output map & P11 semantic-order truth hierarchies & n/a & \hyperref[thm:p11-ktf-expansion-transport]{\ref*{thm:p11-ktf-expansion-transport}} & \thepage & Valid source sentences map to their actual translated sentence codes; source nonsentences map to target nonsentences. \\
\phantomsection\label{entry:symbol-registry-03572}\(q_R\) & \texttt{\detokenize{q_R}} & actual named target code of the transport compiler graph formula & P11 semantic-order truth hierarchies & n/a & \hyperref[thm:p11-ktf-expansion-transport]{\ref*{thm:p11-ktf-expansion-transport}} & \thepage & Selected after construction of the prospective compiler; its copied formula has exactly three retained free slots. \\
\phantomsection\label{entry:symbol-registry-03573}\(S^{\alpha}_{\mathrm{old}},S^{\alpha}_{\mathrm{new}}\) & \texttt{\detokenize{S^{\alpha}_{\mathrm{old}},}\allowbreak\detokenize{S^{\alpha}_{\mathrm{new}}}} & source and ordinary-expansion bottom truth-stage sequences & P11 semantic-order truth hierarchies & n/a & \hyperref[thm:p11-ktf-expansion-transport]{\ref*{thm:p11-ktf-expansion-transport}} & \thepage & They may stabilize at different stages; the translated source formula has the same value at every ordinal stage under matched ordinary interpretations. \\
\phantomsection\label{entry:symbol-registry-03574}\(a_*\) & \texttt{\detokenize{a_*}} & genuine target nonsentence default for malformed source formula inputs & P11 semantic-order truth hierarchies & n/a & \hyperref[annex:p11.ktf.expansion-transport]{\ref*{annex:p11.ktf.expansion-transport}} & \thepage & Chosen as the variable-term code rather than a false-sentence code so that negative noncodes enter at the original first jump. \\
\phantomsection\label{entry:symbol-registry-04085}\(\mathrm{Coop},\mathrm{Def}\) & \texttt{\detokenize{\mathrm{Coop},}\allowbreak\detokenize{\mathrm{Def}}} & actions in a stipulated two-player individual/joint contrast & P11 opening rationality frameworks & n/a & \hyperref[def:p11-ri-two-agent-model]{\ref*{def:p11-ri-two-agent-model}} & \thepage & The payoffs and best-response criterion are explicit model assumptions; no participant choices or universal rationality norm are asserted. \\
\phantomsection\label{entry:symbol-registry-04086}\(\mathcal C\) & \texttt{\detokenize{\mathcal C}} & set of the four displayed card objects & P11 rationality-framework card task & n/a & \hyperref[def:p11-ri-card-task]{\ref*{def:p11-ri-card-task}} & \thepage & The labels A D 4 and 7 name the visible-face cards; the proof checks this closed display, not an unobserved card population. \\
\phantomsection\label{entry:symbol-registry-04087}\(\mathsf V(k),\mathsf E(k)\) & \texttt{\detokenize{\mathsf V(k),}\allowbreak\detokenize{\mathsf E(k)}} & vowel-letter and even-number predicates on a card & P11 rationality-framework card task & n/a & \hyperref[def:p11-ri-card-task]{\ref*{def:p11-ri-card-task}} & \thepage & The rule's only countercase is a vowel paired with an odd number; these are predicates on the card's two respective faces. \\
\phantomsection\label{entry:symbol-registry-04088}\(t,n,\Omega,E_i\) & \texttt{\detokenize{t,}\allowbreak\detokenize{n,}\allowbreak\detokenize{\Omega,}\allowbreak\detokenize{E_i}} & threshold and finite uniform-event construction & P11 rationality-framework threshold example & n/a & \hyperref[def:p11-ri-threshold-model]{\ref*{def:p11-ri-threshold-model}} & \thepage & 0<=t<1; n is the selected number of equiprobable outcomes; E\_i omits exactly outcome i. \\
\phantomsection\label{entry:symbol-registry-04089}\(\mathbb P\) & \texttt{\detokenize{\mathbb P}} & probability measure on specified finite events & P11 rationality-framework threshold swan and urn examples & n/a & \hyperref[def:p11-ri-threshold-model]{\ref*{def:p11-ri-threshold-model}} & \thepage & Measures belong to separate stipulated models; none is an elicited population distribution. \\
\phantomsection\label{entry:symbol-registry-04090}\(\Gamma,N,D\) & \texttt{\detokenize{\Gamma,}\allowbreak\detokenize{N,}\allowbreak\detokenize{D}} & descriptive premise set and fresh normative and descriptive atoms & P11 rationality-framework no-bridge example & n/a & \hyperref[def:p11-ri-no-bridge-language]{\ref*{def:p11-ri-no-bridge-language}} & \thepage & N is an ordinary fresh propositional atom with no bridge axiom; not a deontic operator. \\
\phantomsection\label{entry:symbol-registry-04091}\(W_1,W_2\) & \texttt{\detokenize{W_1,}\allowbreak\detokenize{W_2}} & white-swan observation events in a stipulated joint law & P11 rationality-framework norm illustration & n/a & \hyperref[sec:p11-ri-normative-descriptive]{\ref*{sec:p11-ri-normative-descriptive}} & \thepage & The negatively associated law demonstrates a probability-axiom boundary; no empirical swan frequencies. \\
\phantomsection\label{entry:symbol-registry-04092}\(R,W,B\) & \texttt{\detokenize{R,}\allowbreak\detokenize{W,}\allowbreak\detokenize{B}} & red white and black color events & P11 rationality-framework introduction Ellsberg variant & n/a & \hyperref[def:p11-ri-ellsberg-model]{\ref*{def:p11-ri-ellsberg-model}} & \thepage & Colors retain the introduction's labeling; disjoint exhaustive events with unknown white/black split. \\
\phantomsection\label{entry:symbol-registry-04093}\(H,L,u,\Delta u\) & \texttt{\detokenize{H,}\allowbreak\detokenize{L,}\allowbreak\detokenize{u,}\allowbreak\detokenize{\Delta u}} & common consequences utility and positive utility difference & P11 rationality-framework introduction Ellsberg variant & n/a & \hyperref[def:p11-ri-ellsberg-model]{\ref*{def:p11-ri-ellsberg-model}} & \thepage & State-independent common consequences; Delta u is used in the annex proof only. \\
\phantomsection\label{entry:symbol-registry-04094}\(a,b,c,\succ_i,\succ_M\) & \texttt{\detokenize{a,}\allowbreak\detokenize{b,}\allowbreak\detokenize{c,}\allowbreak\detokenize{\succ_i,}\allowbreak\detokenize{\succ_M}} & alternatives individual strict rankings and pairwise majority & P11 rationality-framework social aggregation example & n/a & \hyperref[def:p11-ri-majority-model]{\ref*{def:p11-ri-majority-model}} & \thepage & Ordinal comparisons only; majority requires at least two of the three voters. \\
\phantomsection\label{entry:symbol-registry-04101}\(N,\chi^2,p\) & \texttt{\detokenize{N,}\allowbreak\detokenize{\chi^2,}\allowbreak\detokenize{p}} & reported response-unit count chi-square statistic and significance probability & \hyperref[sec:p11-ri-individual-social]{\ref*{sec:p11-ri-individual-social}} & n/a & \hyperref[def:p11-ri-collaborative-statistics]{\ref*{def:p11-ri-collaborative-statistics}} & \thepage & MoshmanGeil1998: N=52 means32 control responses plus20 group decisions; N=20 compares the ten groups in each interactive condition. Group members' postdiscussion responses are not independent group trials. Statistical p is not an object-language atom or a probability of the mechanism. \\
\phantomsection\label{entry:symbol-registry-04321}\(A,B\) & \texttt{\detokenize{A,}\allowbreak\detokenize{B}} & players in the source's two-person dilemma & P11 structural-rationality-10-5 & meta & \hyperref[def:p11-rat105-source-pd]{\ref*{def:p11-rat105-source-pd}} & \thepage & Source player names at printed 626; not propositions or action labels. \\
\phantomsection\label{entry:symbol-registry-04322}\(C,D\) & \texttt{\detokenize{C,}\allowbreak\detokenize{D}} & cooperation and defection action labels & P11 structural-rationality-10-5 & meta & \hyperref[def:p11-rat105-source-pd]{\ref*{def:p11-rat105-source-pd}} & \thepage & Source actions at printed 626; full utility interpretation is explicit. \\
\phantomsection\label{entry:symbol-registry-04323}\(q\) & \texttt{\detokenize{q}} & opponent cooperation probability in the local dilemma calculations & P11 structural-rationality-10-5 & meta & \hyperref[thm:p11-rat105-pd]{\ref*{thm:p11-rat105-pd}} & \thepage & Editorial probability in the closed unit interval; not an object-language proposition. \\
\phantomsection\label{entry:symbol-registry-04324}\(Z\) & \texttt{\detokenize{Z}} & finite nonempty set of fully described consequences & P11 structural-rationality-10-5 & meta & \hyperref[thm:p11-rat105-value]{\ref*{thm:p11-rat105-value}} & \thepage & Includes the four joint action profiles in the explicitly editorial structural-value example. \\
\phantomsection\label{entry:symbol-registry-04325}\(z\) & \texttt{\detokenize{z}} & consequence belonging to the supplied finite set & P11 structural-rationality-10-5 & meta & \hyperref[thm:p11-rat105-value]{\ref*{thm:p11-rat105-value}} & \thepage & Dummy consequence in finite expectation sums. \\
\phantomsection\label{entry:symbol-registry-04326}\(v,w\) & \texttt{\detokenize{v,}\allowbreak\detokenize{w}} & supplied consequence-value functions and positive affine transformation & P11 structural-rationality-10-5 & meta & \hyperref[thm:p11-rat105-value]{\ref*{thm:p11-rat105-value}} & \thepage & Structural content of a value is separately specified and not inferred from representability. \\
\phantomsection\label{entry:symbol-registry-04327}\(V_v,V_w\) & \texttt{\detokenize{V_v,}\allowbreak\detokenize{V_w}} & expected-value functions on supplied lottery laws & P11 structural-rationality-10-5 & meta & \hyperref[thm:p11-rat105-value]{\ref*{thm:p11-rat105-value}} & \thepage & Finite sums; source's subjective-value idea made explicit locally without claiming a general Ramsey existence theorem. \\
\phantomsection\label{entry:symbol-registry-04328}\(p,r,s,p',r'\) & \texttt{\detokenize{p,}\allowbreak\detokenize{r,}\allowbreak\detokenize{s,}\allowbreak\detokenize{p',}\allowbreak\detokenize{r'}} & lotteries on the finite consequence domain & P11 structural-rationality-10-5 & meta & \hyperref[thm:p11-rat105-value]{\ref*{thm:p11-rat105-value}} & \thepage & Probability vectors with nonnegative masses summing to one; not propositions. \\
\phantomsection\label{entry:symbol-registry-04329}\(a,b\) & \texttt{\detokenize{a,}\allowbreak\detokenize{b}} & positive affine slope and arbitrary translation & P11 structural-rationality-10-5 & meta & \hyperref[thm:p11-rat105-value]{\ref*{thm:p11-rat105-value}} & \thepage & Real a is strictly positive; b is any real. Distinct from the source's players A and B. \\
\phantomsection\label{entry:symbol-registry-04330}\(\alpha\) & \texttt{\detokenize{\alpha}} & common mixture weight for preference independence & P11 structural-rationality-10-5 & meta & \hyperref[thm:p11-rat105-value]{\ref*{thm:p11-rat105-value}} & \thepage & Strictly between zero and one in the independence biconditional; not event independence. \\
\phantomsection\label{entry:symbol-registry-04331}\(M,\delta,\varepsilon\) & \texttt{\detokenize{M,}\allowbreak\detokenize{\delta,}\allowbreak\detokenize{\varepsilon}} & finite value bound strict-comparison margin and continuity radius & P11 structural-rationality-10-5 & meta & \hyperref[thm:p11-rat105-value]{\ref*{thm:p11-rat105-value}} & \thepage & M is the positive sum of one plus each squared consequence value; delta is a positive value gap; epsilon is delta divided by four M. \\
\phantomsection\label{entry:symbol-registry-04332}\(d\) & \texttt{\detokenize{d}} & relative Euclidean distance on lottery probability vectors & P11 structural-rationality-10-5 & meta & \hyperref[thm:p11-rat105-value]{\ref*{thm:p11-rat105-value}} & \thepage & Uses the exact P17 continuity interface; not the defection action D. \\
\phantomsection\label{entry:symbol-registry-04333}\(t\) & \texttt{\detokenize{t}} & own cooperation mixing probability in the editorial value example & P11 structural-rationality-10-5 & meta & \hyperref[thm:p11-rat105-value]{\ref*{thm:p11-rat105-value}} & \thepage & Closed unit interval; independent draws yield structural expected value t q. \\
\phantomsection\label{entry:symbol-registry-04334}\(A,B\) & \texttt{\detokenize{A,}\allowbreak\detokenize{B}} & sure program and risky lottery in the source disease example & P11 adaptationism disease model & object & \hyperref[def:p11-rat106-disease-model]{\ref*{def:p11-rat106-disease-model}} & \thepage & A saves200 certainly; B saves600 with probability1/3 and0 otherwise. Distinct from propositions or event labels. \\
\phantomsection\label{entry:symbol-registry-04335}\(u,u_k\) & \texttt{\detokenize{u,}\allowbreak\detokenize{u_k}} & consequence utility and specified threshold indicator utility & P11 adaptationism disease model & meta & \hyperref[def:p11-rat106-disease-model]{\ref*{def:p11-rat106-disease-model}} & \thepage & Editorial utility symbols make the source's omitted expected-utility assumption explicit; not evolutionary fitness. \\
\phantomsection\label{entry:symbol-registry-04336}\(k,x\) & \texttt{\detokenize{k,}\allowbreak\detokenize{x}} & need threshold and number of lives saved & P11 adaptationism disease model & n/a & \hyperref[def:p11-rat106-disease-model]{\ref*{def:p11-rat106-disease-model}} & \thepage & Editorial threshold0<=k<=600 and consequence argument x. \\
\phantomsection\label{entry:symbol-registry-04337}\(P,Q\) & \texttt{\detokenize{P,}\allowbreak\detokenize{Q}} & propositional schematic letters in a classical material conditional & P11 adaptationism conditional model & object & \hyperref[thm:p11-rat106-conditional]{\ref*{thm:p11-rat106-conditional}} & \thepage & Published roles at637 and639-640; these are propositions rather than the one-shot base rate P of the separate analytic model. \\
\phantomsection\label{entry:symbol-registry-04338}\(P\) & \texttt{\detokenize{P}} & one-shot dyad base probability in the dependent analytic model & P11 adaptationism reciprocity model & n/a & \hyperref[def:p11-rat106-analytic]{\ref*{def:p11-rat106-analytic}} & \thepage & Published symbol in Delton2011SI3-4;0<P<1. Not the conditional antecedent or an epistemic attitude. \\
\phantomsection\label{entry:symbol-registry-04339}\(b,c\) & \texttt{\detokenize{b,}\allowbreak\detokenize{c}} & within-round delivered benefit and personal cost & P11 adaptationism reciprocity model & n/a & \hyperref[sec:p11-rat106-reciprocity]{\ref*{sec:p11-rat106-reciprocity}} & \thepage & Published supplement roles with b>c>0; simulations setc=1. Benefit/cost numbers are mapped into model fitness not moral worth. \\
\phantomsection\label{entry:symbol-registry-04340}\(w\) & \texttt{\detokenize{w}} & probability a repeated dyad continues to the next round & P11 adaptationism reciprocity model & n/a & \hyperref[sec:p11-rat106-reciprocity]{\ref*{sec:p11-rat106-reciprocity}} & \thepage & 0<=w<1; theoretical w=0 endpoint distinguished from simulated values. \\
\phantomsection\label{entry:symbol-registry-04341}\(d\) & \texttt{\detokenize{d}} & distance between the cue distributions' means in standard-deviation units & P11 adaptationism reciprocity model & n/a & \hyperref[sec:p11-rat106-reciprocity]{\ref*{sec:p11-rat106-reciprocity}} & \thepage & Published SI1 symbol; tested values1,2,3. \\
\phantomsection\label{entry:symbol-registry-04342}\(\Theta_{\mathrm{OS}},\Theta_R\) & \texttt{\detokenize{\Theta_{\mathrm{OS}},}\allowbreak\detokenize{\Theta_R}} & probabilities of a one-shot belief conditional on actual one-shot or repeated dyad & P11 adaptationism reciprocity model & n/a & \hyperref[def:p11-rat106-analytic]{\ref*{def:p11-rat106-analytic}} & \thepage & Published SI3-4 roles; not posterior probabilities of dyad type given a particular cue. \\
\phantomsection\label{entry:symbol-registry-04343}\(\alpha_i,\beta_i\) & \texttt{\detokenize{\alpha_i,}\allowbreak\detokenize{\beta_i}} & cooperation probabilities given one-shot or repeated belief for player i & P11 adaptationism reciprocity model & n/a & \hyperref[def:p11-rat106-analytic]{\ref*{def:p11-rat106-analytic}} & \thepage & Published SI3-4 values in[0,1]; conditional classification-to-action mappings rather than cue densities. \\
\phantomsection\label{entry:symbol-registry-04344}\(V_1\) & \texttt{\detokenize{V_1}} & player one's ex ante expected payoff in the fixed-partner analytic model & P11 adaptationism reciprocity model & n/a & \hyperref[def:p11-rat106-analytic]{\ref*{def:p11-rat106-analytic}} & \thepage & Published SI4 payoff; the best response is not an evolutionary convergence claim. \\
\phantomsection\label{entry:symbol-registry-04345}\(s_i,r_i\) & \texttt{\detokenize{s_i,}\allowbreak\detokenize{r_i}} & conditional cooperation probabilities given actual one-shot and repeated dyad types & P11 adaptationism reciprocity model & n/a & \hyperref[thm:p11-rat106-best-response]{\ref*{thm:p11-rat106-best-response}} & \thepage & Editorial abbreviations of the source's published belief-conditioned mixtures used only in the algebraic S1 reconstruction. \\
\phantomsection\label{entry:symbol-registry-04346}\(i\) & \texttt{\detokenize{i}} & player index1 or2 in a dyad & P11 adaptationism reciprocity model & meta & \hyperref[def:p11-rat106-analytic]{\ref*{def:p11-rat106-analytic}} & \thepage & Published subscript role; fixed partner2 when optimizing player1. \\
\phantomsection\label{entry:symbol-registry-04347}\(a,t\) & \texttt{\detokenize{a,}\allowbreak\detokenize{t}} & affine utility slope and intercept in the lottery proof & P11 adaptationism disease model & n/a & \hyperref[annex:p11.rat106-lottery]{\ref*{annex:p11.rat106-lottery}} & \thepage & Editorial real coefficients in the proof's u(x)=ax+t representation. Not the source's dyadic action probabilities. \\
\phantomsection\label{entry:symbol-registry-04348}\(y,\widetilde u\) & \texttt{\detokenize{y,}\allowbreak\detokenize{\widetilde u}} & death-count argument and transported utility in the frame proof & P11 adaptationism disease model & meta & \hyperref[annex:p11.rat106-lottery]{\ref*{annex:p11.rat106-lottery}} & \thepage & Editorial y=600-x and transported utility of death outcomes; same consequence evaluation rather than a new preference. \\
\phantomsection\label{entry:symbol-registry-04349}\(n\) & \texttt{\detokenize{n}} & positive round index in the geometric-series payoff calculation & P11 adaptationism reciprocity model & meta & \hyperref[annex:p11.rat106-best-response]{\ref*{annex:p11.rat106-best-response}} & \thepage & Editorial positive integer for the probability of reaching round n; not an agent sample size. \\
\phantomsection\label{entry:symbol-registry-04590}\(T,R\) & \texttt{\detokenize{T,}\allowbreak\detokenize{R}} & fixed qualifying-case predicate and interpreted legal-consequence-judgment predicate & P11.rationality-11-3.fixed-rule & meta & \hyperref[def:p11-rat113-fixed-rule]{\ref*{def:p11-rat113-fixed-rule}} & \thepage & R is a truth-assessable proposition in the declared reconstruction; not physical enforcement or a bare command. \\
\phantomsection\label{entry:symbol-registry-04591}\(x,a\) & \texttt{\detokenize{x,}\allowbreak\detokenize{a}} & case variable and named case constant & P11.rationality-11-3.fixed-rule & object & \hyperref[def:p11-rat113-fixed-rule]{\ref*{def:p11-rat113-fixed-rule}} & \thepage & Nonempty domain of cases with the chosen rule interpretation fixed. \\
\phantomsection\label{entry:symbol-registry-04592}\(T_1,\ldots,T_n,T'_1\) & \texttt{\detokenize{T_1,}\allowbreak\detokenize{\ldots,}\allowbreak\detokenize{T_n,}\allowbreak\detokenize{T'_1}} & source rule's interpreted condition predicates and replacement condition & P11.rationality-11-3.analogy & meta & \hyperref[def:p11-rat113-analogy]{\ref*{def:p11-rat113-analogy}} & \thepage & Similarity alone does not define a logical inclusion or an authoritative extended rule. \\
\phantomsection\label{entry:symbol-registry-04593}\(T,T'\) & \texttt{\detokenize{T,}\allowbreak\detokenize{T'}} & original and comparison-case aggregate predicates & P11.rationality-11-3.analogy & meta & \hyperref[def:p11-rat113-analogy]{\ref*{def:p11-rat113-analogy}} & \thepage & T and T' are explicitly defined conjunctive condition predicates. \\
\phantomsection\label{entry:symbol-registry-04594}\(n,i\) & \texttt{\detokenize{n,}\allowbreak\detokenize{i}} & number of source conditions and conjunction index & P11.rationality-11-3.analogy & meta & \hyperref[def:p11-rat113-analogy]{\ref*{def:p11-rat113-analogy}} & \thepage & n>=1; remaining-condition conjunction is true at n=1. \\
\phantomsection\label{entry:symbol-registry-04595}\(\bigwedge\) & \texttt{\detokenize{\bigwedge}} & finite conjunction of interpreted conditions & P11.rationality-11-3.analogy & infix & \hyperref[def:p11-rat113-analogy]{\ref*{def:p11-rat113-analogy}} & \thepage & Empty finite conjunction is explicitly true. \\
\phantomsection\label{entry:symbol-registry-04596}\(d\) & \texttt{\detokenize{d}} & single domain object in formal countermodels & P11.rationality-11-3.annex & object & \hyperref[annex:p11.rat113-detachment]{\ref*{annex:p11.rat113-detachment}} & \thepage & The countermodels test formal consequence not facts about a legal system. \\
\phantomsection\label{entry:symbol-registry-04597}\(\Gamma,\Delta\) & \texttt{\detokenize{\Gamma,}\allowbreak\detokenize{\Delta}} & classical premise sets before and after adding premises & P11.rationality-11-3.annex & meta & \hyperref[annex:p11.rat113-detachment]{\ref*{annex:p11.rat113-detachment}} & \thepage & Monotonicity holds for unchanged premises and consequence relation; not a model of legal defeasibility. \\
\phantomsection\label{entry:symbol-registry-04598}\(R\) & \texttt{\detokenize{R}} & domain of practical rational assessment of actions & P11 rationality-and-morality & meta & \hyperref[sec:p11-rationality-12-1]{\ref*{sec:p11-rationality-12-1}} & \thepage & Published source domain abbreviation at681; neither a preference relation nor the adaptive-function standard of neighboring10.6. \\
\phantomsection\label{entry:symbol-registry-04599}\(M\) & \texttt{\detokenize{M}} & domain of moral assessment of actions & P11 rationality-and-morality & meta & \hyperref[sec:p11-rationality-12-1]{\ref*{sec:p11-rationality-12-1}} & \thepage & Published source abbreviation at681; moral rightness here means permission and does not imply unique obligation. \\
\phantomsection\label{entry:symbol-registry-04600}\(p\) & \texttt{\detokenize{p}} & propositional content of an intrinsic desire in the RM-desire classification & P11 rationality-and-morality & object & \hyperref[def:p11-rat121-rm-desires]{\ref*{def:p11-rat121-rm-desires}} & \thepage & Published figure12.1.2 role; not a probability or an expected-utility lottery. \\
\phantomsection\label{entry:symbol-registry-04601}\(\mathrm{mpi}\) & \texttt{\detokenize{\mathrm{mpi}}} & source abbreviation for morally positive in itself & P11 rationality-and-morality & meta & \hyperref[def:p11-rat121-rm-desires]{\ref*{def:p11-rat121-rm-desires}} & \thepage & Published figure12.1.2 classification abbreviation; what qualifies depends on the adopted morality. \\
\phantomsection\label{entry:symbol-registry-04602}\(t\) & \texttt{\detokenize{t}} & time of the acting person-stage and its contemporaneous desires & P11 rationality-and-morality & meta & \hyperref[sec:p11-rat121-desires]{\ref*{sec:p11-rat121-desires}} & \thepage & Published source time index at683; no intertemporal aggregation formula is implied. \\
\phantomsection\label{entry:symbol-registry-04603}\(k\) & \texttt{\detokenize{k}} & permitted factor weighting own welfare against others' welfare & P11 rationality-and-morality prerogatives & n/a & \hyperref[def:p11-rat121-prerogatives]{\ref*{def:p11-rat121-prerogatives}} & \thepage & Published source686 says k>1; independent of cost cap l and welfare floor m. \\
\phantomsection\label{entry:symbol-registry-04604}\(l\) & \texttt{\detokenize{l}} & permitted upper limit on units of own welfare surrendered & P11 rationality-and-morality prerogatives & n/a & \hyperref[def:p11-rat121-prerogatives]{\ref*{def:p11-rat121-prerogatives}} & \thepage & Published source686 says l>0; units require an adopted welfare scale. \\
\phantomsection\label{entry:symbol-registry-04605}\(m\) & \texttt{\detokenize{m}} & own-welfare floor allowed by a threshold prerogative & P11 rationality-and-morality prerogatives & n/a & \hyperref[def:p11-rat121-prerogatives]{\ref*{def:p11-rat121-prerogatives}} & \thepage & Published source686-687 threshold; not a morality variable or a universal normative constant. \\
\phantomsection\label{entry:symbol-registry-04606}\(H\) & \texttt{\detokenize{H}} & independently assignable hypothesis atom in the local Boolean test schema & P11 scientific-rationality-14-1 & object & \hyperref[thm:p11-rat141-test-logic]{\ref*{thm:p11-rat141-test-logic}} & \thepage & Not a probability event in this schema; later probability interfaces retain their own declared event roles. \\
\phantomsection\label{entry:symbol-registry-04607}\(A\) & \texttt{\detokenize{A}} & independently assignable auxiliary-premise atom in the local Boolean test schema & P11 scientific-rationality-14-1 & object & \hyperref[thm:p11-rat141-test-logic]{\ref*{thm:p11-rat141-test-logic}} & \thepage & Represents a conjunction of test assumptions; not an act or a group agent. \\
\phantomsection\label{entry:symbol-registry-04608}\(E\) & \texttt{\detokenize{E}} & independently assignable predicted-observation atom in the local Boolean test schema & P11 scientific-rationality-14-1 & object & \hyperref[thm:p11-rat141-test-logic]{\ref*{thm:p11-rat141-test-logic}} & \thepage & False E refutes the tested H-and-A conjunction; true E alone does not verify H. Not an expectation operator. \\
\phantomsection\label{entry:symbol-registry-04609}\(H\) & \texttt{\detokenize{H}} & descriptive hypothesis being accepted or rejected & P11 value-freedom of science & meta & \hyperref[def:p11-rat142-risk-profile]{\ref*{def:p11-rat142-risk-profile}} & \thepage & False positive means acceptance of a false H in the source's terminology; mapping to statistical null rejection must be declared. \\
\phantomsection\label{entry:symbol-registry-04610}\(\alpha\) & \texttt{\detokenize{\alpha}} & selected frequentist test level & P11 value-freedom of science & object & \hyperref[def:p11-rat142-risk-profile]{\ref*{def:p11-rat142-risk-profile}} & \thepage & Uses the existing P15 test definition; the mathematical error bound does not select an ethically legitimate evidential convention. \\
\phantomsection\label{entry:symbol-registry-04659}\(D\) & \texttt{\detokenize{D}} & finite set of candidate die face counts & P11 moral learning dice & object & \hyperref[def:p11-r123-dice]{\ref*{def:p11-r123-dice}} & \thepage & Candidates are 4 6 8 and 10 under an explicitly uniform prior. \\
\phantomsection\label{entry:symbol-registry-04660}\((d,e)\) & \texttt{\detokenize{d;e}} & candidate die face counts & P11 moral learning dice & object & \hyperref[def:p11-r123-dice]{\ref*{def:p11-r123-dice}} & \thepage & Distinct from dilemma or ethical-response codes. \\
\phantomsection\label{entry:symbol-registry-04661}\(n\) & \texttt{\detokenize{n}} & number of independent retained-die rolls & P11 moral learning dice & object & \hyperref[def:p11-r123-dice]{\ref*{def:p11-r123-dice}} & \thepage & Empty and incompatible data boundaries are explicit. \\
\phantomsection\label{entry:symbol-registry-04662}\(H_d\) & \texttt{\detokenize{H_d}} & selected-die hypothesis & P11 moral learning dice & meta & \hyperref[def:p11-r123-dice]{\ref*{def:p11-r123-dice}} & \thepage & Not a moral principle or a participant classification. \\
\phantomsection\label{entry:symbol-registry-04663}\(x\) & \texttt{\detokenize{x}} & ordered observed roll sequence & P11 moral learning dice & object & \hyperref[def:p11-r123-dice]{\ref*{def:p11-r123-dice}} & \thepage & Uses the exact ten outcomes printed in the 2021 chapter. \\
\phantomsection\label{entry:symbol-registry-04664}\(L_d(x)\) & \texttt{\detokenize{L_d(x)}} & likelihood of the observed sequence under a die hypothesis & P11 moral learning dice & object & \hyperref[def:p11-r123-dice]{\ref*{def:p11-r123-dice}} & \thepage & Zero when any observation lies outside the candidate face set. \\
\phantomsection\label{entry:symbol-registry-04665}\(p_d(x)\) & \texttt{\detokenize{p_d(x)}} & posterior probability of a die hypothesis & P11 moral learning dice & object & \hyperref[def:p11-r123-dice]{\ref*{def:p11-r123-dice}} & \thepage & Requires positive marginal probability of the observation. \\
\phantomsection\label{entry:symbol-registry-04666}\(\Pr\) & \texttt{\detokenize{\Pr}} & probability under the declared finite dice model & P11 moral learning dice & meta & \hyperref[def:p11-r123-dice]{\ref*{def:p11-r123-dice}} & \thepage & Not a moral correctness score or estimated human learning parameter. \\
\phantomsection\label{entry:symbol-registry-04667}\((t,F,r,p)\) & \texttt{\detokenize{t;F;r;p}} & reported test statistics correlation and p-value & P11 moral judgment original records & meta & \hyperref[sec:p11-r123-kahane]{\ref*{sec:p11-r123-kahane}} & \thepage & Statistical roles differ from die-size posterior probabilities and from response-time variables. \\
\phantomsection\label{entry:symbol-registry-04713}\(\mathrm{worm}\) & \texttt{\detokenize{\mathrm{worm}}} & unary worm predicate in the source example & P11 normative systems logic & object & \hyperref[def:p11-nsc-logic-model]{\ref*{def:p11-nsc-logic-model}} & \thepage & Actual source vocabulary under classical interpretation. \\
\phantomsection\label{entry:symbol-registry-04714}\(\mathrm{warble}\) & \texttt{\detokenize{\mathrm{warble}}} & unary warbling predicate & P11 normative systems logic & object & \hyperref[def:p11-nsc-logic-model]{\ref*{def:p11-nsc-logic-model}} & \thepage & Distinct from the later event H. \\
\phantomsection\label{entry:symbol-registry-04715}\(\mathrm{Albert}\) & \texttt{\detokenize{\mathrm{Albert}}} & constant naming the example object & P11 normative systems logic & object & \hyperref[def:p11-nsc-logic-model]{\ref*{def:p11-nsc-logic-model}} & \thepage & Not a universally quantified variable. \\
\phantomsection\label{entry:symbol-registry-04716}\(x\) & \texttt{\detokenize{x}} & object variable bound by the universal quantifier & P11 normative systems logic & object & \hyperref[def:p11-nsc-logic-model]{\ref*{def:p11-nsc-logic-model}} & \thepage & Not the complexity section's encoded input string. \\
\phantomsection\label{entry:symbol-registry-04717}\(a\) & \texttt{\detokenize{a}} & denotation of Albert in a model & P11 normative systems logic & object & \hyperref[def:p11-nsc-logic-model]{\ref*{def:p11-nsc-logic-model}} & \thepage & The repaired positive example denotes5. \\
\phantomsection\label{entry:symbol-registry-04718}\(S_{\rm worm}\) & \texttt{\detokenize{S_{\rm worm}}} & extension of the worm predicate & P11 normative systems logic & object & \hyperref[def:p11-nsc-logic-model]{\ref*{def:p11-nsc-logic-model}} & \thepage & Subset of the declared model domain. \\
\phantomsection\label{entry:symbol-registry-04719}\(S_{\rm warble}\) & \texttt{\detokenize{S_{\rm warble}}} & extension of the warbling predicate & P11 normative systems logic & object & \hyperref[def:p11-nsc-logic-model]{\ref*{def:p11-nsc-logic-model}} & \thepage & Contains Albert in the repaired positive model. \\
\phantomsection\label{entry:symbol-registry-04720}\(\subset\) & \texttt{\detokenize{\subset}} & inclusion allowing equality in the source's local notation & P11 normative systems logic & relation & \hyperref[def:p11-nsc-logic-model]{\ref*{def:p11-nsc-logic-model}} & \thepage & The source explicitly permits equality; not undeclared strict inclusion. \\
\phantomsection\label{entry:symbol-registry-04721}\(\&\) & \texttt{\detokenize{\&}} & Boolean conjunction in the source notation & P11 normative systems logic & infix & \hyperref[def:p11-nsc-logic-model]{\ref*{def:p11-nsc-logic-model}} & \thepage & Not an alignment separator in the mathematical content. \\
\phantomsection\label{entry:symbol-registry-04722}\(\forall\) & \texttt{\detokenize{\forall}} & universal quantifier & P11 normative systems logic & meta & \hyperref[def:p11-nsc-logic-model]{\ref*{def:p11-nsc-logic-model}} & \thepage & Nonempty domain semantics. \\
\phantomsection\label{entry:symbol-registry-04723}\(\to\) & \texttt{\detokenize{\to}} & material conditional & P11 normative systems logic & infix & \hyperref[def:p11-nsc-logic-model]{\ref*{def:p11-nsc-logic-model}} & \thepage & Not every ordinary-language conditional interpretation. \\
\phantomsection\label{entry:symbol-registry-04724}\(\neg\) & \texttt{\detokenize{\neg}} & classical negation & P11 normative systems logic & meta & \hyperref[def:p11-nsc-logic-model]{\ref*{def:p11-nsc-logic-model}} & \thepage & The refutation equivalence uses classical reductio. \\
\phantomsection\label{entry:symbol-registry-04725}\(\Gamma\) & \texttt{\detokenize{\Gamma}} & premise set in the consequence test & P11 normative systems logic & object & \hyperref[def:p11-nsc-logic-model]{\ref*{def:p11-nsc-logic-model}} & \thepage & Static premises do not choose a revision policy. \\
\phantomsection\label{entry:symbol-registry-04726}\(\varphi\) & \texttt{\detokenize{\varphi}} & warbling conclusion and generic sentence & P11 normative systems logic & object & \hyperref[def:p11-nsc-logic-model]{\ref*{def:p11-nsc-logic-model}} & \thepage & The particular example precedes the general classical assertion. \\
\phantomsection\label{entry:symbol-registry-04727}\(\Omega\) & \texttt{\detokenize{\Omega}} & finite probability sample space & P11 normative systems probability & object & \hyperref[def:p11-nsc-probability-model]{\ref*{def:p11-nsc-probability-model}} & \thepage & Anchor sec:p11-nsc-probability; coherence is not calibration. \\
\phantomsection\label{entry:symbol-registry-04728}\(\Pr\) & \texttt{\detokenize{\Pr}} & probability measure in the credence model & P11 normative systems probability & n/a & \hyperref[def:p11-nsc-probability-model]{\ref*{def:p11-nsc-probability-model}} & \thepage & Interpretation is distinct from validity of the identities. \\
\phantomsection\label{entry:symbol-registry-04729}\(A,B\) & \texttt{\detokenize{A,}\allowbreak\detokenize{B}} & event arguments in Bayes' formula & P11 normative systems probability & object & \hyperref[def:p11-nsc-probability-model]{\ref*{def:p11-nsc-probability-model}} & \thepage & Not the later menu sets. \\
\phantomsection\label{entry:symbol-registry-04730}\(\Pr(A,B)\) & \texttt{\detokenize{\Pr(A,}\allowbreak\detokenize{B)}} & joint intersection probability & P11 normative systems probability & n/a & \hyperref[def:p11-nsc-probability-model]{\ref*{def:p11-nsc-probability-model}} & \thepage & The source's comma notation is explicitly declared. \\
\phantomsection\label{entry:symbol-registry-04731}\(\Pr(A\mid B)\) & \texttt{\detokenize{\Pr(A\mid B)}} & conditional event probability & P11 normative systems probability & n/a & \hyperref[def:p11-nsc-probability-model]{\ref*{def:p11-nsc-probability-model}} & \thepage & Defined only when the conditioning event has positive probability. \\
\phantomsection\label{entry:symbol-registry-04732}\(A^c,B^c\) & \texttt{\detokenize{A^c,}\allowbreak\detokenize{B^c}} & event complements in the Bayes counterexample & P11 normative systems probability & object & \hyperref[annex:p11.nsc.bayes]{\ref*{annex:p11.nsc.bayes}} & \thepage & All conditioning denominators in the counterexample are positive. \\
\phantomsection\label{entry:symbol-registry-04733}\(H\) & \texttt{\detokenize{H}} & warbling event in the toy case & P11 normative systems worked case & object & \hyperref[def:p11-nsc-worked-model]{\ref*{def:p11-nsc-worked-model}} & \thepage & Anchor sec:p11-nsc-worked; not certainty of the premises. \\
\phantomsection\label{entry:symbol-registry-04734}\(E\) & \texttt{\detokenize{E}} & fallible report event & P11 normative systems worked case & object & \hyperref[def:p11-nsc-worked-model]{\ref*{def:p11-nsc-worked-model}} & \thepage & Report occurrence is not its asserted truth. \\
\phantomsection\label{entry:symbol-registry-04735}\(H^c,E^c\) & \texttt{\detokenize{H^c,}\allowbreak\detokenize{E^c}} & complementary events in the toy law & P11 normative systems worked case & object & \hyperref[def:p11-nsc-worked-model]{\ref*{def:p11-nsc-worked-model}} & \thepage & Four joint masses sum to one. \\
\phantomsection\label{entry:symbol-registry-04736}\(c\) & \texttt{\detokenize{c}} & nonnegative utility cost of the perfect test & P11 normative systems worked case & object & \hyperref[def:p11-nsc-worked-model]{\ref*{def:p11-nsc-worked-model}} & \thepage & Stipulated value rather than observed cost. \\
\phantomsection\label{entry:symbol-registry-04737}\(P,M,D\) & \texttt{\detokenize{P,}\allowbreak\detokenize{M,}\allowbreak\detokenize{D}} & poke meal and decoy alternatives & P11 normative systems choice & object & \hyperref[sec:p11-nsc-choice]{\ref*{sec:p11-nsc-choice}} & \thepage & Alternatives rather than propositions or probabilities. \\
\phantomsection\label{entry:symbol-registry-04738}\(X\) & \texttt{\detokenize{X}} & finite option set & P11 normative systems choice & object & \hyperref[def:p11-nsc-menu-model]{\ref*{def:p11-nsc-menu-model}} & \thepage & All nonempty menus belong to the domain. \\
\phantomsection\label{entry:symbol-registry-04739}\(x,y\) & \texttt{\detokenize{x,}\allowbreak\detokenize{y}} & options compared in menus and revealed preference & P11 normative systems choice & object & \hyperref[def:p11-nsc-menu-model]{\ref*{def:p11-nsc-menu-model}} & \thepage & Not first-order quantified objects or encoded verifier inputs. \\
\phantomsection\label{entry:symbol-registry-04740}\(C\) & \texttt{\detokenize{C}} & nonempty choice correspondence & P11 normative systems choice & n/a & \hyperref[def:p11-nsc-menu-model]{\ref*{def:p11-nsc-menu-model}} & \thepage & Retains ties; explicit contraction and expansion conditions. \\
\phantomsection\label{entry:symbol-registry-04741}\(A,B\) & \texttt{\detokenize{A,}\allowbreak\detokenize{B}} & nonempty menus & P11 normative systems choice & object & \hyperref[def:p11-nsc-menu-model]{\ref*{def:p11-nsc-menu-model}} & \thepage & Not the earlier event arguments. \\
\phantomsection\label{entry:symbol-registry-04742}\(\succeq_C\) & \texttt{\detokenize{\succeq_C}} & revealed weak comparison & P11 normative systems choice & relation & \hyperref[def:p11-nsc-menu-model]{\ref*{def:p11-nsc-menu-model}} & \thepage & Transitivity is an additional condition. \\
\phantomsection\label{entry:symbol-registry-04743}\(u\) & \texttt{\detokenize{u}} & ordinal real representation & P11 normative systems choice & n/a & \hyperref[def:p11-nsc-menu-model]{\ref*{def:p11-nsc-menu-model}} & \thepage & Increasing ordinal transformations need not preserve expectations. \\
\phantomsection\label{entry:symbol-registry-04744}\(\operatorname{argmax}\) & \texttt{\detokenize{\operatorname{argmax}}} & set of maximizing menu alternatives & P11 normative systems choice & n/a & \hyperref[thm:p11-nsc-menu-representation]{\ref*{thm:p11-nsc-menu-representation}} & \thepage & All tied maximizers are retained. \\
\phantomsection\label{entry:symbol-registry-04745}\(\Delta_1,\Delta_2\) & \texttt{\detokenize{\Delta_1,}\allowbreak\detokenize{\Delta_2}} & positive transaction fees & P11 normative systems money pump & object & \hyperref[thm:p11-nsc-pump]{\ref*{thm:p11-nsc-pump}} & \thepage & Anchor sec:p11-nsc-pump; acceptance is a premise. \\
\phantomsection\label{entry:symbol-registry-04746}\(k\) & \texttt{\detokenize{k}} & number of affordable accepted rounds & P11 normative systems money pump & object & \hyperref[thm:p11-nsc-pump]{\ref*{thm:p11-nsc-pump}} & \thepage & No unbounded repetition is guaranteed. \\
\phantomsection\label{entry:symbol-registry-04747}\(w\) & \texttt{\detokenize{w}} & current wealth & P11 normative systems money pump & object & \hyperref[annex:p11.nsc.pump]{\ref*{annex:p11.nsc.pump}} & \thepage & Not the SAT certificate with the same glyph. \\
\phantomsection\label{entry:symbol-registry-04748}\(q\) & \texttt{\detokenize{q}} & positive lottery component weight & P11 normative systems lotteries & object & \hyperref[sec:p11-nsc-lotteries]{\ref*{sec:p11-nsc-lotteries}} & \thepage & Zero weight gives identical lotteries. \\
\phantomsection\label{entry:symbol-registry-04749}\(s\) & \texttt{\detokenize{s}} & lottery interpolation weight & P11 normative systems lotteries & object & \hyperref[sec:p11-nsc-lotteries]{\ref*{sec:p11-nsc-lotteries}} & \thepage & Full weak-order and continuity assumptions required. \\
\phantomsection\label{entry:symbol-registry-04750}\(R\) & \texttt{\detokenize{R}} & intermediate prize in the source example & P11 normative systems lotteries & object & \hyperref[sec:p11-nsc-lotteries]{\ref*{sec:p11-nsc-lotteries}} & \thepage & Not a revealed comparison relation. \\
\phantomsection\label{entry:symbol-registry-04751}\(\mathsf{SAT}\) & \texttt{\detokenize{\mathsf{SAT}}} & Boolean satisfiability language & P11 normative systems complexity & object & \hyperref[def:p11-nsc-sat-task]{\ref*{def:p11-nsc-sat-task}} & \thepage & Anchor sec:p11-nsc-limits; differs from UNSAT and entailment. \\
\phantomsection\label{entry:symbol-registry-04752}\(x,n\) & \texttt{\detokenize{x,}\allowbreak\detokenize{n}} & encoded input and its length & P11 normative systems complexity & object & \hyperref[def:p11-nsc-sat-task]{\ref*{def:p11-nsc-sat-task}} & \thepage & Not a first-order object variable. \\
\phantomsection\label{entry:symbol-registry-04753}\(w\) & \texttt{\detokenize{w}} & verifier certificate & P11 normative systems complexity & object & \hyperref[def:p11-nsc-sat-task]{\ref*{def:p11-nsc-sat-task}} & \thepage & Not wealth in the pump proof. \\
\phantomsection\label{entry:symbol-registry-04754}\(V\) & \texttt{\detokenize{V}} & fixed deterministic verifier & P11 normative systems complexity & object & \hyperref[def:p11-nsc-sat-task]{\ref*{def:p11-nsc-sat-task}} & \thepage & Finite local transition model used in the full construction. \\
\phantomsection\label{entry:symbol-registry-04755}\(p(n)\) & \texttt{\detokenize{p(n)}} & polynomial certificate-length bound & P11 normative systems complexity & n/a & \hyperref[def:p11-nsc-sat-task]{\ref*{def:p11-nsc-sat-task}} & \thepage & Not a probability or significance level. \\
\phantomsection\label{entry:symbol-registry-04756}\(T,L,b\) & \texttt{\detokenize{T,}\allowbreak\detokenize{L,}\allowbreak\detokenize{b}} & time bound tape length and code width & P11 normative systems complexity & object & \hyperref[def:p11-nsc-sat-task]{\ref*{def:p11-nsc-sat-task}} & \thepage & Finite simulation has polynomial circuit size. \\
\phantomsection\label{entry:symbol-registry-04757}\(\mathsf F\) & \texttt{\detokenize{\mathsf F}} & total local update function & P11 normative systems complexity & n/a & \hyperref[def:p11-nsc-sat-task]{\ref*{def:p11-nsc-sat-task}} & \thepage & Fixed truth table; illegal codes are not reached from the legal initialization. \\
\phantomsection\label{entry:symbol-registry-04758}\(z,a_g,b_g\) & \texttt{\detokenize{z,}\allowbreak\detokenize{a_g,}\allowbreak\detokenize{b_g}} & Boolean gate-output and input variables & P11 normative systems complexity & object & \hyperref[def:p11-nsc-sat-task]{\ref*{def:p11-nsc-sat-task}} & \thepage & Clause equivalences proved explicitly. \\
\phantomsection\label{entry:symbol-registry-04759}\(\lor\) & \texttt{\detokenize{\lor}} & Boolean disjunction in clauses & P11 normative systems complexity & infix & \hyperref[thm:p11-nsc-sat-complete]{\ref*{thm:p11-nsc-sat-complete}} & \thepage & No editorial inference-rule display is introduced. \\
\end{longtable}
\endgroup
\Needspace{.4\textheight}
\subsection*{Stable part P12 scopes}
\begingroup\scriptsize\setlength{\tabcolsep}{3pt}
\begin{longtable}{@{}p{0.9in}p{1.3in}p{1.0in}p{0.5in}p{0.3in}p{0.4in}p{0.35in}p{1.65in}@{}}
\toprule
\textbf{Glyph} & \textbf{Macro} & \textbf{Role} & \textbf{Scope} & \textbf{Prec./form} & \textbf{First} & \textbf{Reg. page} & \textbf{Principal uses / notes} \\
\midrule
\endfirsthead
\toprule
\textbf{Glyph} & \textbf{Macro} & \textbf{Role} & \textbf{Scope} & \textbf{Prec./form} & \textbf{First} & \textbf{Reg. page} & \textbf{Principal uses / notes} \\
\midrule
\endhead
\midrule\multicolumn{8}{r}{\emph{continued}}\\\endfoot
\bottomrule\endlastfoot
\phantomsection\label{entry:symbol-registry-00615}\(\mathsf P\) & \texttt{\detokenize{\mathsf P}} & deterministic polynomial-time language class & P12 formal-practice & n/a & \hyperref[sec:formal-practice-01]{\ref*{sec:formal-practice-01}} & \thepage & Complexity class reused in the P12 SAT and certificate comparison; disjoint from the P4.3 pairing-graph formula. \\
\phantomsection\label{entry:symbol-registry-00616}\(\mathsf{SAT}\) & \texttt{\detokenize{\mathsf{SAT}}} & satisfiable propositional encodings & P12 formal-practice & n/a & \hyperref[sec:formal-practice-01]{\ref*{sec:formal-practice-01}} & \thepage & SAT instances under the exact clause-set and encoding conventions fixed in the automation chapter. \\
\phantomsection\label{entry:symbol-registry-00617}\(\mathsf{DPLL}\) & \texttt{\detokenize{\mathsf{DPLL}}} & Davis--Putnam--Logemann--Loveland method identifier & P12 formal-practice & n/a & \hyperref[sec:formal-practice-01]{\ref*{sec:formal-practice-01}} & \thepage & Exact recursive finite-clause search method defined before its correctness theorem and trace. \\
\phantomsection\label{entry:symbol-registry-00618}\(\#\mathsf{SAT}\) & \texttt{\detokenize{\#\mathsf{SAT}}} & satisfying-assignment counting function & P12 formal-practice & n/a & \hyperref[sec:formal-practice-01]{\ref*{sec:formal-practice-01}} & \thepage & Maps a propositional encoding to its number of satisfying total assignments; not a decision language. \\
\phantomsection\label{entry:symbol-registry-00619}\(\#\mathsf P\) & \texttt{\detokenize{\#\mathsf P}} & counting complexity class & P12 formal-practice & n/a & \hyperref[sec:formal-practice-01]{\ref*{sec:formal-practice-01}} & \thepage & Functions counting accepting witnesses of a polynomial-time decidable relation. \\
\phantomsection\label{entry:symbol-registry-00620}\(\mathsf{MaxSAT}\) & \texttt{\detokenize{\mathsf{MaxSAT}}} & maximum satisfiability optimization problem & P12 formal-practice & n/a & \hyperref[sec:formal-practice-01]{\ref*{sec:formal-practice-01}} & \thepage & Optimization problem maximizing the number or declared weight of satisfied clauses under the local variant. \\
\phantomsection\label{entry:symbol-registry-00621}\(\mathsf{SMT}\) & \texttt{\detokenize{\mathsf{SMT}}} & satisfiability-modulo-theories identifier & P12 automated-reasoning & n/a & \hyperref[sec:automated-reasoning-01]{\ref*{sec:automated-reasoning-01}} & \thepage & Decision problem for quantifier-free formulas interpreted in a fixed background theory. \\
\phantomsection\label{entry:symbol-registry-00622}\(\mathsf{EUF}\) & \texttt{\detokenize{\mathsf{EUF}}} & equality with uninterpreted functions theory identifier & P12 automated-reasoning & n/a & \hyperref[sec:automated-reasoning-01]{\ref*{sec:automated-reasoning-01}} & \thepage & Quantifier-free equality theory whose function symbols have no axioms beyond equality congruence. \\
\phantomsection\label{entry:symbol-registry-00623}\(M\) & \texttt{\detokenize{M}} & Slagle ground-interpretation metavariable & P12 mechanization & n/a & \hyperref[sec:mechanization-01]{\ref*{sec:mechanization-01}} & \thepage & The fixed ground interpretation in semantic resolution; exact source role, disjoint from the P7 term-metavariable scope. \\
\phantomsection\label{entry:symbol-registry-00624}\(N,E_1,\ldots,E_q\) & \texttt{\detokenize{N,}\allowbreak\detokenize{E_i}} & semantic-clash nucleus and electron clauses & P12 mechanization & n/a & \hyperref[sec:mechanization-01]{\ref*{sec:mechanization-01}} & \thepage & Exact Slagle roles: \(N\) is the nucleus and the \(E_i\) are electrons tested against the fixed ground interpretation \(M\). \\
\phantomsection\label{entry:symbol-registry-00625}\(R_1,\ldots,R_{q+1}\) & \texttt{\detokenize{R_i}} & semantic-clash intermediate resolvents & P12 mechanization & n/a & \hyperref[sec:mechanization-01]{\ref*{sec:mechanization-01}} & \thepage & Exact Slagle resolvent chain with \(R_1=N\) and final false \(M\)-resolvent \(R_{q+1}\); distinct from frame relations. \\
\phantomsection\label{entry:symbol-registry-00626}\(\mathsf{LCF}\) & \texttt{\detokenize{\mathsf{LCF}}} & LCF trusted-kernel architecture identifier & P12 research-workflow & n/a & \hyperref[sec:research-workflow-01]{\ref*{sec:research-workflow-01}} & \thepage & Interactive-prover architecture in which an abstract theorem type can be constructed only by a small trusted kernel. \\
\phantomsection\label{entry:symbol-registry-00627}\(\mathsf{PTS}\) & \texttt{\detokenize{\mathsf{PTS}}} & pure type system identifier & P12 research-workflow & n/a & \hyperref[sec:research-workflow-01]{\ref*{sec:research-workflow-01}} & \thepage & Parameterized typed calculus determined by its sorts axioms and allowed product triples. \\
\phantomsection\label{entry:symbol-registry-00628}\(S=(\mathcal S,\mathcal A,\mathcal R);\lambda S\) & \texttt{\detokenize{S,}\allowbreak\detokenize{\mathcal S,}\allowbreak\detokenize{\mathcal A,}\allowbreak\detokenize{\mathcal R,}\allowbreak\detokenize{\lambda S}} & Barendregt pure-type-system specification and calculus & P12 research-workflow & n/a & \hyperref[sec:research-workflow-01]{\ref*{sec:research-workflow-01}} & \thepage & The exact source datum consists of sorts sort-typing axioms and dependent-product triples; lambda S names the calculus it determines. \\
\phantomsection\label{entry:symbol-registry-00629}\(\ast,\Box;(s_1,s_2)\) & \texttt{\detokenize{\ast,}\allowbreak\detokenize{\Box;(s_1,}\allowbreak\detokenize{s_2)}} & Barendregt lambda-cube sorts and product-sort pair & P12 research-workflow & n/a & \hyperref[sec:research-workflow-01]{\ref*{sec:research-workflow-01}} & \thepage & Exact source notation: the pair abbreviates the triple (s1,s2,s2); the locally declared PTS role distinguishes Box and star from modal and algebraic uses. \\
\phantomsection\label{entry:symbol-registry-00630}\(\lambda{\to},\lambda\mathrm P\) & \texttt{\detokenize{\lambda{\to},}\allowbreak\detokenize{\lambda\mathrm P}} & Barendregt lambda-cube calculus identifiers & P12 research-workflow & n/a & \hyperref[sec:research-workflow-01]{\ref*{sec:research-workflow-01}} & \thepage & Exact source identifiers for the base simply typed vertex and the types-on-terms vertex. \\
\phantomsection\label{entry:symbol-registry-00631}\(\lambda 2,\lambda\mathrm P2\) & \texttt{\detokenize{\lambda 2,}\allowbreak\detokenize{\lambda\mathrm P2}} & Barendregt lambda-cube calculus identifiers & P12 research-workflow & n/a & \hyperref[sec:research-workflow-01]{\ref*{sec:research-workflow-01}} & \thepage & Exact source identifiers for the terms-on-types vertex and its types-on-terms extension. \\
\phantomsection\label{entry:symbol-registry-00632}\(\lambda\underline{\omega},\lambda\mathrm P\underline{\omega}\) & \texttt{\detokenize{\lambda\underline{\omega},}\allowbreak\detokenize{\lambda\mathrm P\underline{\omega}}} & Barendregt lambda-cube calculus identifiers & P12 research-workflow & n/a & \hyperref[sec:research-workflow-01]{\ref*{sec:research-workflow-01}} & \thepage & Exact source identifiers for the types-on-types vertex and its types-on-terms extension. \\
\phantomsection\label{entry:symbol-registry-00633}\(\lambda\omega,\lambda\mathrm P\omega=\lambda\mathrm C\) & \texttt{\detokenize{\lambda\omega,}\allowbreak\detokenize{\lambda\mathrm P\omega,}\allowbreak\detokenize{\lambda\mathrm C}} & Barendregt lambda-cube calculus identifiers & P12 research-workflow & n/a & \hyperref[sec:research-workflow-01]{\ref*{sec:research-workflow-01}} & \thepage & Exact source identifiers for the terms-on-types plus types-on-types vertex and the all-three-axes vertex; lambda P omega equals lambda C. \\
\phantomsection\label{entry:symbol-registry-00634}\(\mathsf{CIC}\) & \texttt{\detokenize{\mathsf{CIC}}} & Calculus of Inductive Constructions identifier & P12 research-workflow & n/a & \hyperref[sec:research-workflow-01]{\ref*{sec:research-workflow-01}} & \thepage & Calculus of Constructions extended by the separately described inductive families and eliminators. \\
\phantomsection\label{entry:symbol-registry-00635}\(\alpha,\beta,\gamma\) & \texttt{\detokenize{\alpha,}\allowbreak\detokenize{\beta,}\allowbreak\detokenize{\gamma}} & type metavariables & P12 research-workflow & n/a & \hyperref[sec:research-workflow-01]{\ref*{sec:research-workflow-01}} & \thepage & Types in the reused STLC and lambda-cube examples; this scope is disjoint from ordinal and path-formula roles. \\
\phantomsection\label{entry:symbol-registry-00636}\(\alpha\to_{\mathrm{ty}}\beta\) & \texttt{\detokenize{\to_{\mathrm{ty}}}} & object-language function-type former & P12 research-workflow & 6 & \hyperref[sec:research-workflow-01]{\ref*{sec:research-workflow-01}} & \thepage & Right-associative arrow of the exact reused STLC syntax. \\
\phantomsection\label{entry:symbol-registry-00637}\(M{:}\alpha\) & \texttt{\detokenize{:}} & type ascription & P12 research-workflow & meta & \hyperref[sec:research-workflow-01]{\ref*{sec:research-workflow-01}} & \thepage & The local STLC example fixes the term and type categories; distinct from labelled formulas. \\
\phantomsection\label{entry:symbol-registry-00638}\(\Theta_{\mathrm{ty}}\) & \texttt{\detokenize{\Theta_{\mathrm{ty}}}} & term-variable typing context & P12 research-workflow & n/a & \hyperref[sec:research-workflow-01]{\ref*{sec:research-workflow-01}} & \thepage & Finite list of distinct term-variable declarations in the reused STLC example. \\
\phantomsection\label{entry:symbol-registry-00639}\(\mathsf{STLC}\) & \texttt{\detokenize{\mathsf{STLC}}} & arrow-only simply typed lambda-calculus identifier & P12 research-workflow & n/a & \hyperref[sec:research-workflow-01]{\ref*{sec:research-workflow-01}} & \thepage & Exact calculus reused from sheet:stlc in the mechanization example. \\
\phantomsection\label{entry:symbol-registry-00640}\(\Theta_{\mathrm{ty}}\derives_{\mathsf{STLC}}M{:}\alpha\) & \texttt{\detokenize{\derives_{\mathsf{STLC}}}} & STLC typing judgment & P12 research-workflow & meta & \hyperref[sec:research-workflow-01]{\ref*{sec:research-workflow-01}} & \thepage & Exact reused arrow-only typing judgment; no new P12 calculus is introduced. \\
\phantomsection\label{entry:symbol-registry-01919}\(\Sigma\) & \texttt{\detokenize{\Sigma}} & declared first-order signature & \hyperref[sec:p12-small-clause-standard]{\ref*{sec:p12-small-clause-standard}} & meta & \hyperref[def:p12-small-language]{\ref*{def:p12-small-language}} & \thepage & Fresh Skolem symbols lie outside the current signature for fixed-structure expansion. \\
\phantomsection\label{entry:symbol-registry-01920}\(\mathcal F\) & \texttt{\detokenize{\mathcal F}} & function-symbol alphabet with finite arities & \hyperref[sec:p12-small-clause-standard]{\ref*{sec:p12-small-clause-standard}} & meta & \hyperref[def:p12-small-language]{\ref*{def:p12-small-language}} & \thepage & Separate from relation-symbol alphabet and variable names. \\
\phantomsection\label{entry:symbol-registry-01921}\(\mathcal R\) & \texttt{\detokenize{\mathcal R}} & predicate-symbol alphabet with finite arities & \hyperref[sec:p12-small-clause-standard]{\ref*{sec:p12-small-clause-standard}} & meta & \hyperref[def:p12-small-language]{\ref*{def:p12-small-language}} & \thepage & Calligraphic alphabet; plain R is an example predicate. \\
\phantomsection\label{entry:symbol-registry-01922}\(X\) & \texttt{\detokenize{X}} & infinite variable-name supply & \hyperref[sec:p12-small-clause-standard]{\ref*{sec:p12-small-clause-standard}} & meta & \hyperref[def:p12-small-language]{\ref*{def:p12-small-language}} & \thepage & Finite formulas use finite subsets; fresh binder blocks are explicitly allocated. \\
\phantomsection\label{entry:symbol-registry-01923}\(\mathfrak A\) & \texttt{\detokenize{\mathfrak A}} & first-order structure & \hyperref[sec:p12-small-clause-standard]{\ref*{sec:p12-small-clause-standard}} & meta & \hyperref[def:p12-small-language]{\ref*{def:p12-small-language}} & \thepage & Nonempty domain and interpretation of every declared signature symbol. \\
\phantomsection\label{entry:symbol-registry-01924}\(D\) & \texttt{\detokenize{D}} & nonempty interpretation domain & \hyperref[sec:p12-small-clause-standard]{\ref*{sec:p12-small-clause-standard}} & meta & \hyperref[def:p12-small-language]{\ref*{def:p12-small-language}} & \thepage & The letter D also denotes a clause only in the explicitly typed clause-reduction contexts. \\
\phantomsection\label{entry:symbol-registry-01925}\(\nu\) & \texttt{\detokenize{\nu}} & variable valuation & \hyperref[sec:p12-small-clause-standard]{\ref*{sec:p12-small-clause-standard}} & meta & \hyperref[def:p12-small-language]{\ref*{def:p12-small-language}} & \thepage & A map from variable names to the nonempty domain. \\
\phantomsection\label{entry:symbol-registry-01926}\(\nu[x/d]\) & \texttt{\detokenize{\nu[x/}\allowbreak\detokenize{d]}} & valuation update & \hyperref[sec:p12-small-clause-standard]{\ref*{sec:p12-small-clause-standard}} & meta & \hyperref[def:p12-small-language]{\ref*{def:p12-small-language}} & \thepage & Changes the designated variable coordinate only. \\
\phantomsection\label{entry:symbol-registry-01927}\(t^{\mathfrak A,\nu}\) & \texttt{\detokenize{t^{\mathfrak A,\nu}}} & term evaluation & \hyperref[sec:p12-small-clause-standard]{\ref*{sec:p12-small-clause-standard}} & meta & \hyperref[def:p12-small-language]{\ref*{def:p12-small-language}} & \thepage & Homomorphic application of the declared function interpretations. \\
\phantomsection\label{entry:symbol-registry-01928}\(\equiv_\Sigma\) & \texttt{\detokenize{\equiv_\Sigma}} & pointwise first-order semantic equivalence & \hyperref[sec:p12-small-clause-standard]{\ref*{sec:p12-small-clause-standard}} & meta & \hyperref[def:p12-small-language]{\ref*{def:p12-small-language}} & \thepage & Bare equiv in transformation chains abbreviates this local typed relation. \\
\phantomsection\label{entry:symbol-registry-01929}\(\operatorname{free}\) & \texttt{\detokenize{\operatorname{free}}} & free-variable-name set & \hyperref[sec:p12-small-clause-standard]{\ref*{sec:p12-small-clause-standard}} & meta & \hyperref[def:p12-small-language]{\ref*{def:p12-small-language}} & \thepage & Names rather than binder occurrences; union and binding deletion are explicit. \\
\phantomsection\label{entry:symbol-registry-01930}\(\operatorname{bound}\) & \texttt{\detokenize{\operatorname{bound}}} & bound-variable-name set & \hyperref[sec:p12-small-clause-standard]{\ref*{sec:p12-small-clause-standard}} & meta & \hyperref[def:p12-small-language]{\ref*{def:p12-small-language}} & \thepage & A repeated binder name occurs once in this set. \\
\phantomsection\label{entry:symbol-registry-01931}\(\operatorname{vars}\) & \texttt{\detokenize{\operatorname{vars}}} & all occurring variable names & \hyperref[sec:p12-small-clause-standard]{\ref*{sec:p12-small-clause-standard}} & meta & \hyperref[def:p12-small-language]{\ref*{def:p12-small-language}} & \thepage & Union of free names and bound names. \\
\phantomsection\label{entry:symbol-registry-01932}\(\Phi^\exists\) & \texttt{\detokenize{\Phi^\exists}} & existential closure of an open formula & \hyperref[sec:p12-small-clause-standard]{\ref*{sec:p12-small-clause-standard}} & meta & \hyperref[thm:p12-small-sat-closure]{\ref*{thm:p12-small-sat-closure}} & \thepage & Closes exactly the finite free-variable list. \\
\phantomsection\label{entry:symbol-registry-01933}\(\pi\) & \texttt{\detokenize{\pi}} & formula position & \hyperref[sec:p12-small-clause-standard]{\ref*{sec:p12-small-clause-standard}} & meta & \hyperref[def:p12-small-positions]{\ref*{def:p12-small-positions}} & \thepage & Finite child-index word with empty-word root. \\
\phantomsection\label{entry:symbol-registry-01934}\(\Phi|_\pi\) & \texttt{\detokenize{\Phi|_\pi}} & subformula at a position & \hyperref[sec:p12-small-clause-standard]{\ref*{sec:p12-small-clause-standard}} & meta & \hyperref[def:p12-small-positions]{\ref*{def:p12-small-positions}} & \thepage & Formula children are distinguished from term children. \\
\phantomsection\label{entry:symbol-registry-01935}\(\Phi[\pi/\Psi]\) & \texttt{\detokenize{\Phi[\pi/}\allowbreak\detokenize{\Psi]}} & replacement at a selected occurrence & \hyperref[sec:p12-small-clause-standard]{\ref*{sec:p12-small-clause-standard}} & meta & \hyperref[def:p12-small-positions]{\ref*{def:p12-small-positions}} & \thepage & Surrounding binders remain in place. \\
\phantomsection\label{entry:symbol-registry-01936}\(\operatorname{pol}\) & \texttt{\detokenize{\operatorname{pol}}} & occurrence polarity & \hyperref[sec:p12-small-clause-standard]{\ref*{sec:p12-small-clause-standard}} & meta & \hyperref[def:p12-small-positions]{\ref*{def:p12-small-positions}} & \thepage & Positive one, negative one, or zero below a biconditional. \\
\phantomsection\label{entry:symbol-registry-01937}\(\sigma\) & \texttt{\detokenize{\sigma}} & finite-support term substitution & \hyperref[sec:p12-small-clause-standard]{\ref*{sec:p12-small-clause-standard}} & meta & \hyperref[def:p12-small-positions]{\ref*{def:p12-small-positions}} & \thepage & Simultaneous substitution; raw variable renaming is defined separately. \\
\phantomsection\label{entry:symbol-registry-01938}\(b(A)\) & \texttt{\detokenize{b(A)}} & number of Boolean and-or nodes & \hyperref[sec:p12-small-clause-standard]{\ref*{sec:p12-small-clause-standard}} & meta & \hyperref[thm:p12-small-miniscope]{\ref*{thm:p12-small-miniscope}} & \thepage & Used inside the nonnegative quantifier-scope termination measure. \\
\phantomsection\label{entry:symbol-registry-01939}\(\rho\) & \texttt{\detokenize{\rho}} & injective map on the occurring variable set & \hyperref[sec:p12-small-clause-standard]{\ref*{sec:p12-small-clause-standard}} & meta & \hyperref[def:p12-small-raw-renaming]{\ref*{def:p12-small-raw-renaming}} & \thepage & No injective identity-extension to the whole alphabet is required. \\
\phantomsection\label{entry:symbol-registry-01940}\(\operatorname{raw}_\rho\) & \texttt{\detokenize{\operatorname{raw}_\rho}} & raw replacement of variables and binder names & \hyperref[sec:p12-small-clause-standard]{\ref*{sec:p12-small-clause-standard}} & meta & \hyperref[def:p12-small-raw-renaming]{\ref*{def:p12-small-raw-renaming}} & \thepage & Semantic valuation lemma covers both quantifiers. \\
\phantomsection\label{entry:symbol-registry-01941}\(nq\) & \texttt{\detokenize{nq}} & quantifier-occurrence count & \hyperref[sec:p12-small-clause-standard]{\ref*{sec:p12-small-clause-standard}} & meta & \hyperref[def:p12-small-rectification]{\ref*{def:p12-small-rectification}} & \thepage & Counts binders with multiplicity. \\
\phantomsection\label{entry:symbol-registry-01942}\(\operatorname{ren}\) & \texttt{\detokenize{\operatorname{ren}}} & indexed fresh-binder rectification & \hyperref[sec:p12-small-clause-standard]{\ref*{sec:p12-small-clause-standard}} & meta & \hyperref[def:p12-small-rectification]{\ref*{def:p12-small-rectification}} & \thepage & Disjoint finite name blocks prove freshness and allocation. \\
\phantomsection\label{entry:symbol-registry-01943}\(C\sigma\subseteq D\) & \texttt{\detokenize{C\sigma\subseteq D}} & literal-multiset inclusion & \hyperref[def:p12-small-clause-reductions]{\ref*{def:p12-small-clause-reductions}} & meta & \hyperref[def:p12-small-clause-reductions]{\ref*{def:p12-small-clause-reductions}} & \thepage & Includes multiplicity; C and D in this definition are clauses. \\
\phantomsection\label{entry:symbol-registry-01944}\(\llbracket A\rrbracket_\nu\) & \texttt{\detokenize{\llbracket A\rrbracket_\nu}} & Boolean truth value at a valuation & \hyperref[annex:p12-small-clause-standard]{\ref*{annex:p12-small-clause-standard}} & meta & \hyperref[annex:p12.small.context]{\ref*{annex:p12.small.context}} & \thepage & One exactly when the formula is satisfied; zero otherwise. \\
\phantomsection\label{entry:symbol-registry-01945}\(W_{\bar d}\) & \texttt{\detokenize{W_{\bar d}}} & fiber of existential witnesses & \hyperref[annex:p12.small.skolem]{\ref*{annex:p12.small.skolem}} & meta & \hyperref[annex:p12.small.skolem]{\ref*{annex:p12.small.skolem}} & \thepage & Indexed by free-parameter tuples; Choice selects from nonempty fibers. \\
\phantomsection\label{entry:symbol-registry-01946}\(\nu_\sigma\) & \texttt{\detokenize{\nu_\sigma}} & valuation induced by evaluating a substitution & \hyperref[annex:p12.small.clause.simplify]{\ref*{annex:p12.small.clause.simplify}} & meta & \hyperref[annex:p12.small.clause.simplify]{\ref*{annex:p12.small.clause.simplify}} & \thepage & Provides the universal clause-instance semantics. \\
\phantomsection\label{entry:symbol-registry-01947}\(\Sigma\) & \texttt{\detokenize{\Sigma}} & declared first-order signature & \hyperref[sec:p12-opt-small-clauses]{\ref*{sec:p12-opt-small-clauses}} & meta & \hyperref[def:p12-opt-interface]{\ref*{def:p12-opt-interface}} & \thepage & Fresh symbols lie outside this signature for fixed-structure expansion. \\
\phantomsection\label{entry:symbol-registry-01948}\(D\) & \texttt{\detokenize{D}} & nonempty first-order domain & \hyperref[sec:p12-opt-small-clauses]{\ref*{sec:p12-opt-small-clauses}} & meta & \hyperref[def:p12-opt-interface]{\ref*{def:p12-opt-interface}} & \thepage & Arbitrary set; no finite-domain restriction. \\
\phantomsection\label{entry:symbol-registry-01949}\(\operatorname{FV}\) & \texttt{\detokenize{\operatorname{FV}}} & free-variable set & \hyperref[sec:p12-opt-small-clauses]{\ref*{sec:p12-opt-small-clauses}} & meta & \hyperref[def:p12-opt-interface]{\ref*{def:p12-opt-interface}} & \thepage & Same local role as companion free notation. \\
\phantomsection\label{entry:symbol-registry-01950}\(C[A]\) & \texttt{\detokenize{C[A]}} & one-occurrence formula context & \hyperref[sec:p12-opt-small-clauses]{\ref*{sec:p12-opt-small-clauses}} & meta & \hyperref[def:p12-opt-polarity]{\ref*{def:p12-opt-polarity}} & \thepage & Surrounding binders remain. \\
\phantomsection\label{entry:symbol-registry-01951}\(\operatorname{pol}\) & \texttt{\detokenize{\operatorname{pol}}} & occurrence polarity & \hyperref[sec:p12-opt-small-clauses]{\ref*{sec:p12-opt-small-clauses}} & meta & \hyperref[def:p12-opt-polarity]{\ref*{def:p12-opt-polarity}} & \thepage & Positive one, negative one or zero. \\
\phantomsection\label{entry:symbol-registry-01952}\(\pi\) & \texttt{\detokenize{\pi}} & formula occurrence position & \hyperref[sec:p12-opt-small-clauses]{\ref*{sec:p12-opt-small-clauses}} & meta & \hyperref[def:p12-opt-renaming]{\ref*{def:p12-opt-renaming}} & \thepage & Finite child-index word. \\
\phantomsection\label{entry:symbol-registry-01953}\(R\) & \texttt{\detokenize{R}} & fresh auxiliary predicate & \hyperref[sec:p12-opt-small-clauses]{\ref*{sec:p12-opt-small-clauses}} & meta & \hyperref[def:p12-opt-renaming]{\ref*{def:p12-opt-renaming}} & \thepage & All free variables are arguments; example relation R has a separately declared role. \\
\phantomsection\label{entry:symbol-registry-01954}\(R\) & \texttt{\detokenize{R}} & original binary relation in the two witness examples & \hyperref[def:p12-opt-skolem-comparisons]{\ref*{def:p12-opt-skolem-comparisons}} & meta & \hyperref[def:p12-opt-skolem-comparisons]{\ref*{def:p12-opt-skolem-comparisons}} & \thepage & Explicit nested source rebinding within the two examples; this is not the fresh auxiliary predicate of the naming construction. \\
\phantomsection\label{entry:symbol-registry-01955}\(D_\pi\) & \texttt{\detokenize{D_\pi}} & defining sentence & \hyperref[sec:p12-opt-small-clauses]{\ref*{sec:p12-opt-small-clauses}} & meta & \hyperref[def:p12-opt-renaming]{\ref*{def:p12-opt-renaming}} & \thepage & One implication or full biconditional according to polarity. \\
\phantomsection\label{entry:symbol-registry-01956}\(p\) & \texttt{\detokenize{p}} & unreduced positive clause count & \hyperref[sec:p12-opt-small-clauses]{\ref*{sec:p12-opt-small-clauses}} & meta & \hyperref[def:p12-opt-counts]{\ref*{def:p12-opt-counts}} & \thepage & Counts multiplicity after initial constant handling. \\
\phantomsection\label{entry:symbol-registry-01957}\(q\) & \texttt{\detokenize{q}} & unreduced negative clause count & \hyperref[sec:p12-opt-small-clauses]{\ref*{sec:p12-opt-small-clauses}} & meta & \hyperref[def:p12-opt-counts]{\ref*{def:p12-opt-counts}} & \thepage & q(A)=p(not A); translates source barred p. \\
\phantomsection\label{entry:symbol-registry-01958}\(N\) & \texttt{\detokenize{N}} & finite input tree size & \hyperref[sec:p12-opt-small-clauses]{\ref*{sec:p12-opt-small-clauses}} & meta & \hyperref[thm:p12-opt-counts]{\ref*{thm:p12-opt-counts}} & \thepage & Include explicit first-order symbols for output-size claims. \\
\phantomsection\label{entry:symbol-registry-01959}\(a_\pi\) & \texttt{\detokenize{a_\pi}} & positive count coefficient & \hyperref[sec:p12-opt-small-clauses]{\ref*{sec:p12-opt-small-clauses}} & meta & \hyperref[def:p12-opt-coefficients]{\ref*{def:p12-opt-coefficients}} & \thepage & Nonnegative integer including zero. \\
\phantomsection\label{entry:symbol-registry-01960}\(b_\pi\) & \texttt{\detokenize{b_\pi}} & negative count coefficient & \hyperref[sec:p12-opt-small-clauses]{\ref*{sec:p12-opt-small-clauses}} & meta & \hyperref[def:p12-opt-coefficients]{\ref*{def:p12-opt-coefficients}} & \thepage & Nonnegative integer propagated from root. \\
\phantomsection\label{entry:symbol-registry-01961}\(c_\pi\) & \texttt{\detokenize{c_\pi}} & count remainder & \hyperref[sec:p12-opt-small-clauses]{\ref*{sec:p12-opt-small-clauses}} & meta & \hyperref[thm:p12-opt-coefficients]{\ref*{thm:p12-opt-coefficients}} & \thepage & Independent of formula inserted at designated occurrence. \\
\phantomsection\label{entry:symbol-registry-01962}\(\langle v\rangle_3\) & \texttt{\detokenize{\langle v\rangle_3}} & saturation at three & \hyperref[sec:p12-opt-small-clauses]{\ref*{sec:p12-opt-small-clauses}} & meta & \hyperref[def:p12-opt-capped]{\ref*{def:p12-opt-capped}} & \thepage & Minimum of nonnegative integer v and three. \\
\phantomsection\label{entry:symbol-registry-01963}\(f_i\) & \texttt{\detokenize{f_i}} & fresh witness function & \hyperref[sec:p12-opt-small-clauses]{\ref*{sec:p12-opt-small-clauses}} & meta & \hyperref[def:p12-opt-skolem-modes]{\ref*{def:p12-opt-skolem-modes}} & \thepage & Arity equals declared dependency-list length. \\
\phantomsection\label{entry:symbol-registry-01964}\(\sigma\) & \texttt{\detokenize{\sigma}} & capture-free substitution & \hyperref[sec:p12-opt-small-clauses]{\ref*{sec:p12-opt-small-clauses}} & meta & \hyperref[thm:p12-opt-sharing]{\ref*{thm:p12-opt-sharing}} & \thepage & Instance matching here; later simultaneous witness substitution. \\
\phantomsection\label{entry:symbol-registry-01965}\(E\) & \texttt{\detokenize{E}} & existential-conjunction subformula & \hyperref[sec:p12-opt-small-clauses]{\ref*{sec:p12-opt-small-clauses}} & meta & \hyperref[def:p12-opt-skolem-modes]{\ref*{def:p12-opt-skolem-modes}} & \thepage & Positive occurrence in rectified NNF. \\
\phantomsection\label{entry:symbol-registry-01966}\(Z_i\) & \texttt{\detokenize{Z_i}} & ordered free-variable block & \hyperref[sec:p12-opt-small-clauses]{\ref*{sec:p12-opt-small-clauses}} & meta & \hyperref[def:p12-opt-strong-skolem]{\ref*{def:p12-opt-strong-skolem}} & \thepage & Disjoint blocks; empty blocks allowed. \\
\phantomsection\label{entry:symbol-registry-01967}\(W_i\) & \texttt{\detokenize{W_i}} & fresh universal replacement block & \hyperref[sec:p12-opt-small-clauses]{\ref*{sec:p12-opt-small-clauses}} & meta & \hyperref[def:p12-opt-strong-skolem]{\ref*{def:p12-opt-strong-skolem}} & \thepage & Same length as corresponding Z block. \\
\phantomsection\label{entry:symbol-registry-01968}\(\Xi\) & \texttt{\detokenize{\Xi}} & strong-Skolemized replacement & \hyperref[sec:p12-opt-small-clauses]{\ref*{sec:p12-opt-small-clauses}} & meta & \hyperref[def:p12-opt-strong-skolem]{\ref*{def:p12-opt-strong-skolem}} & \thepage & Conjunction with universally quantified suffix blocks. \\
\phantomsection\label{entry:symbol-registry-01969}\(h\) & \texttt{\detokenize{h}} & selector or fiber retraction & \hyperref[sec:p12-opt-small-clauses]{\ref*{sec:p12-opt-small-clauses}} & meta & \hyperref[def:p12-opt-selectors]{\ref*{def:p12-opt-selectors}} & \thepage & Onto nonempty fiber; arbitrary total values at empty fibers. \\
\phantomsection\label{entry:symbol-registry-01970}\(\operatorname{pr}_j\) & \texttt{\detokenize{\operatorname{pr}_j}} & tuple projection & \hyperref[sec:p12-opt-small-clauses]{\ref*{sec:p12-opt-small-clauses}} & meta & \hyperref[def:p12-opt-selectors]{\ref*{def:p12-opt-selectors}} & \thepage & Returns coordinate j. \\
\phantomsection\label{entry:symbol-registry-01971}\(\Theta\) & \texttt{\detokenize{\Theta}} & retained background sentence & \hyperref[sec:p12-opt-small-clauses]{\ref*{sec:p12-opt-small-clauses}} & meta & \hyperref[def:p12-opt-atom-definition]{\ref*{def:p12-opt-atom-definition}} & \thepage & Implies the definition and any globally used guard. \\
\phantomsection\label{entry:symbol-registry-01972}\(H\) & \texttt{\detokenize{H}} & conditional-definition guard & \hyperref[def:p12-opt-atom-definition]{\ref*{def:p12-opt-atom-definition}} & meta & \hyperref[def:p12-opt-atom-definition]{\ref*{def:p12-opt-atom-definition}} & \thepage & Its free variables are among the head parameters; distinct from the later BDD high successor. \\
\phantomsection\label{entry:symbol-registry-01973}\(v_i\) & \texttt{\detokenize{v_i}} & ordered Boolean variable & \hyperref[sec:p12-opt-small-clauses]{\ref*{sec:p12-opt-small-clauses}} & meta & \hyperref[def:p12-opt-bdd]{\ref*{def:p12-opt-bdd}} & \thepage & Indices increase along paths. \\
\phantomsection\label{entry:symbol-registry-01974}\(L,H\) & \texttt{\detokenize{L,}\allowbreak\detokenize{H}} & low and high successors & \hyperref[sec:p12-opt-small-clauses]{\ref*{sec:p12-opt-small-clauses}} & meta & \hyperref[def:p12-opt-bdd]{\ref*{def:p12-opt-bdd}} & \thepage & Local graph role; H also denotes an earlier definition guard. \\
\phantomsection\label{entry:symbol-registry-02096}\(\beta,p\) & \texttt{\detokenize{\beta,}\allowbreak\detokenize{p}} & even radix and significand precision & P12 numerical-error-and-inference & meta & \hyperref[def:p12-numerical-format]{\ref*{def:p12-numerical-format}} & \thepage & Integers with beta at least two and p at least two. \\
\phantomsection\label{entry:symbol-registry-02097}\(e,e_{\min},e_{\max}\) & \texttt{\detokenize{e,}\allowbreak\detokenize{e_{\min},}\allowbreak\detokenize{e_{\max}}} & exponent and finite integer exponent limits & P12 numerical-error-and-inference & meta & \hyperref[def:p12-numerical-format]{\ref*{def:p12-numerical-format}} & \thepage & Unrestricted integer exponents are used only in the idealized interior proof model. \\
\phantomsection\label{entry:symbol-registry-02098}\(d_j\) & \texttt{\detokenize{d_j}} & significand digits & P12 numerical-error-and-inference & meta & \hyperref[def:p12-numerical-format]{\ref*{def:p12-numerical-format}} & \thepage & Leading digit is nonzero in a normal value. \\
\phantomsection\label{entry:symbol-registry-02099}\(R\) & \texttt{\detokenize{R}} & round to a nearest format value with the declared even-midpoint rule & P12 numerical-error-and-inference & meta & \hyperref[def:p12-numerical-rounding]{\ref*{def:p12-numerical-rounding}} & \thepage & Zero is fixed; finite implementation transfer requires the explicit range assumptions. \\
\phantomsection\label{entry:symbol-registry-02100}\(u\) & \texttt{\detokenize{u}} & unit roundoff beta to the power one minus p divided by two & P12 numerical-error-and-inference & meta & \hyperref[def:p12-numerical-rounding]{\ref*{def:p12-numerical-rounding}} & \thepage & Goldberg's machine epsilon; distinct from the next-value gap 2u at one. \\
\phantomsection\label{entry:symbol-registry-02101}\(\operatorname{ulp}(z)\) & \texttt{\detokenize{\operatorname{ulp}(z)}} & last-place grid spacing in the binade of nonzero z & P12 numerical-error-and-inference & meta & \hyperref[def:p12-numerical-rounding]{\ref*{def:p12-numerical-rounding}} & \thepage & The value selecting the binade is explicit; no relative-error claim at zero. \\
\phantomsection\label{entry:symbol-registry-02102}\(\widehat z\) & \texttt{\detokenize{\widehat z}} & approximation to an explicitly stated real target & P12 numerical-error-and-inference & meta & \hyperref[def:p12-numerical-rounding]{\ref*{def:p12-numerical-rounding}} & \thepage & A stored input need not equal a measured target. \\
\phantomsection\label{entry:symbol-registry-02103}\(g\) & \texttt{\detokenize{g}} & number of guard digits in aligned addition or subtraction & P12 numerical-error-and-inference & meta & \hyperref[def:p12-numerical-guard]{\ref*{def:p12-numerical-guard}} & \thepage & No-guard g=0 examples are named separately. \\
\phantomsection\label{entry:symbol-registry-02104}\(\oplus,\ominus\) & \texttt{\detokenize{\oplus,}\allowbreak\detokenize{\ominus}} & declared computed addition and subtraction & P12 numerical-error-and-inference & infix & \hyperref[def:p12-numerical-guard]{\ref*{def:p12-numerical-guard}} & \thepage & At least one guard digit unless explicitly stated; exact rounding is a stronger separately declared procedure. \\
\phantomsection\label{entry:symbol-registry-02105}\(\otimes,\oslash\) & \texttt{\detokenize{\otimes,}\allowbreak\detokenize{\oslash}} & nearest-rounded multiplication and division & P12 numerical-error-and-inference & infix & \hyperref[def:p12-numerical-guard]{\ref*{def:p12-numerical-guard}} & \thepage & Division requires a nonzero denominator. These local arithmetic roles are distinct from logical connectives. \\
\phantomsection\label{entry:symbol-registry-02106}\(\delta_i,a_i\) & \texttt{\detokenize{\delta_i,}\allowbreak\detokenize{a_i}} & multiplicative perturbations and their nonnegative magnitude bounds & P12 numerical-error-and-inference & meta & \hyperref[thm:p12-numerical-product-errors]{\ref*{thm:p12-numerical-product-errors}} & \thepage & The product bound retains all higher-order terms. \\
\phantomsection\label{entry:symbol-registry-02107}\(a,b,c,d\) & \texttt{\detokenize{a,}\allowbreak\detokenize{b,}\allowbreak\detokenize{c,}\allowbreak\detokenize{d}} & quadratic coefficients and discriminant; ordered triangle sides in the named later section & P12 numerical-error-and-inference & meta & \hyperref[thm:p12-numerical-quadratic]{\ref*{thm:p12-numerical-quadratic}} & \thepage & The two uses are explicitly reintroduced in nonoverlapping sections. \\
\phantomsection\label{entry:symbol-registry-02108}\(q,Q\) & \texttt{\detokenize{q,}\allowbreak\detokenize{Q}} & exact and computed ordered-side radicands & P12 numerical-error-and-inference triangle section & meta & \hyperref[def:p12-numerical-triangle]{\ref*{def:p12-numerical-triangle}} & \thepage & Multiplication order and four inner parentheses are specified. \\
\phantomsection\label{entry:symbol-registry-02109}\(\operatorname{SQRT},\operatorname{LN}\) & \texttt{\detokenize{\operatorname{SQRT},}\allowbreak\detokenize{\operatorname{LN}}} & assumed half-ulp accurate library square root and natural logarithm & P12 numerical-error-and-inference & meta & \hyperref[def:p12-numerical-triangle]{\ref*{def:p12-numerical-triangle}} & \thepage & No specific library midpoint direction is needed; the logarithm assumption is introduced in def:p12-numerical-log-increment. \\
\phantomsection\label{entry:symbol-registry-02110}\(\mu\) & \texttt{\detokenize{\mu}} & continuous logarithmic ratio ln(1+x)/x with value one at zero & P12 numerical-error-and-inference & meta & \hyperref[thm:p12-numerical-mu]{\ref*{thm:p12-numerical-mu}} & \thepage & Integral representation and derivative bounds are proved. \\
\phantomsection\label{entry:symbol-registry-02111}\(H_k\) & \texttt{\detokenize{H_k}} & coarse rounding obtained by discarding k low radix digits with integer midpoint parity & P12 numerical-error-and-inference & meta & \hyperref[def:p12-numerical-coarse-rounding]{\ref*{def:p12-numerical-coarse-rounding}} & \thepage & Definition explicitly includes a carry and a one-digit result. \\
\phantomsection\label{entry:symbol-registry-02112}\(T,S,h,l\) & \texttt{\detokenize{T,}\allowbreak\detokenize{S,}\allowbreak\detokenize{h,}\allowbreak\detokenize{l}} & split algorithm intermediates and high/low components & P12 numerical-error-and-inference exact-rounding section & meta & \hyperref[thm:p12-numerical-split-rounding]{\ref*{thm:p12-numerical-split-rounding}} & \thepage & All split operations are exactly rounded. \\
\phantomsection\label{entry:symbol-registry-02113}\(P,B,V,X\) & \texttt{\detokenize{P,}\allowbreak\detokenize{B,}\allowbreak\detokenize{V,}\allowbreak\detokenize{X}} & integer lattice quantities in the complete split proof & P12 numerical-error-and-inference annex split proof & meta & \hyperref[annex:p12.numerical.split-rounding]{\ref*{annex:p12.numerical.split-rounding}} & \thepage & \(P=\beta^p\); \(B=\beta^k\); normalized positive integer \(X\); rounded product \(BV\). These are introduced together in this proof. \\
\phantomsection\label{entry:symbol-registry-02114}\(\operatorname{nint}\) & \texttt{\detokenize{\operatorname{nint}}} & nearest-integer rounding with ties to even & P12 numerical coarse rounding and split proof & meta & \hyperref[def:p12-numerical-coarse-rounding]{\ref*{def:p12-numerical-coarse-rounding}} & \thepage & Definition def:p12-numerical-coarse-rounding and its use in annex:p12.numerical.split-rounding; translation through a midpoint uses the proved parity condition. \\
\phantomsection\label{entry:symbol-registry-02115}\(E,E_v,E_w\) & \texttt{\detokenize{E,}\allowbreak\detokenize{E_v,}\allowbreak\detokenize{E_w}} & proved absolute-error radii for numerical certificates & P12 numerical-error-and-inference certification section & meta & \hyperref[thm:p12-numerical-certification]{\ref*{thm:p12-numerical-certification}} & \thepage & Nonnegative; interval overlap cannot certify strict order. \\
\phantomsection\label{entry:symbol-registry-02116}\(\rho,\eta,q_0,a_0,t_-,t_+\) & \texttt{\detokenize{\rho,}\allowbreak\detokenize{\eta,}\allowbreak\detokenize{q_0,}\allowbreak\detokenize{a_0,}\allowbreak\detokenize{t_-,}\allowbreak\detokenize{t_+}} & rounding displacements and adjacent-cell quantities in the no-drift proof & P12 numerical-error-and-inference annex no-drift proof & meta & \hyperref[annex:p12.numerical.no-drift]{\ref*{annex:p12.numerical.no-drift}} & \thepage & The chapter's u remains unit roundoff; q0 and a0 are the large and small adjacent gaps at a radix power. These local quantities are introduced in the proof. \\
\phantomsection\label{entry:symbol-registry-02117}\(m,n,m',n',q\) & \texttt{\detokenize{m,}\allowbreak\detokenize{n,}\allowbreak\detokenize{m',}\allowbreak\detokenize{n',}\allowbreak\detokenize{q}} & integer operands and scaled operands/quotient in the radix round-trip identity & P12 numerical-error-and-inference round-trip subsection & meta & \hyperref[thm:p12-numerical-round-trip]{\ref*{thm:p12-numerical-round-trip}} & \thepage & The pair m,n are integers; scaled n has the proved two-radix-power form. \\
\phantomsection\label{entry:symbol-registry-02118}\(a_*,b_*,c_*,\eta_a,\eta_b,\eta_c\) & \texttt{\detokenize{a_*,}\allowbreak\detokenize{b_*,}\allowbreak\detokenize{c_*,}\allowbreak\detokenize{\eta_a,}\allowbreak\detokenize{\eta_b,}\allowbreak\detokenize{\eta_c}} & target coefficients and bounds on their discrepancy from stored coefficients & P12 numerical-error-and-inference certification section & meta & \hyperref[thm:p12-numerical-data-budget]{\ref*{thm:p12-numerical-data-budget}} & \thepage & Input uncertainty is separate from arithmetic roundoff. \\
\phantomsection\label{entry:symbol-registry-02157}\(I\) & \texttt{\detokenize{I}} & weighted search instance with hidden graph and revealed information & \hyperref[sec:p12-best-first-search-and-heuristic-guarantees]{\ref*{sec:p12-best-first-search-and-heuristic-guarantees}} & meta & \hyperref[def:p12-hsearch-instance]{\ref*{def:p12-hsearch-instance}} & \thepage & Mathematical descriptor; the whole adjacency structure is not supplied to comparison algorithms. \\
\phantomsection\label{entry:symbol-registry-02158}\(G\) & \texttt{\detokenize{G}} & directed state graph & \hyperref[sec:p12-best-first-search-and-heuristic-guarantees]{\ref*{sec:p12-best-first-search-and-heuristic-guarantees}} & meta & \hyperref[def:p12-hsearch-instance]{\ref*{def:p12-hsearch-instance}} & \thepage & Distinct from its path-unfolding tree. \\
\phantomsection\label{entry:symbol-registry-02159}\(\Gamma\) & \texttt{\detokenize{\Gamma}} & set of goal states & \hyperref[sec:p12-best-first-search-and-heuristic-guarantees]{\ref*{sec:p12-best-first-search-and-heuristic-guarantees}} & meta & \hyperref[def:p12-hsearch-instance]{\ref*{def:p12-hsearch-instance}} & \thepage & A state's goal flag is revealed on generation. \\
\phantomsection\label{entry:symbol-registry-02160}\(c\) & \texttt{\detokenize{c}} & nonnegative edge cost & \hyperref[sec:p12-best-first-search-and-heuristic-guarantees]{\ref*{sec:p12-best-first-search-and-heuristic-guarantees}} & meta & \hyperref[def:p12-hsearch-instance]{\ref*{def:p12-hsearch-instance}} & \thepage & A positive uniform lower bound is imposed only in stated infinite-domain theorems. \\
\phantomsection\label{entry:symbol-registry-02161}\(g(P)\) & \texttt{\detokenize{g(P)}} & cost of a persistent path record & \hyperref[sec:p12-best-first-search-and-heuristic-guarantees]{\ref*{sec:p12-best-first-search-and-heuristic-guarantees}} & meta & \hyperref[def:p12-hsearch-instance]{\ref*{def:p12-hsearch-instance}} & \thepage & Typed by path argument; additive edge sum. \\
\phantomsection\label{entry:symbol-registry-02162}\(g(v)\) & \texttt{\detokenize{g(v)}} & current best discovered cost label of a state & \hyperref[sec:p12-best-first-search-and-heuristic-guarantees]{\ref*{sec:p12-best-first-search-and-heuristic-guarantees}} & meta & \hyperref[def:p12-hsearch-algorithm]{\ref*{def:p12-hsearch-algorithm}} & \thepage & Typed by state argument; updates strictly decrease the label. \\
\phantomsection\label{entry:symbol-registry-02163}\(g^*\) & \texttt{\detokenize{g^*}} & least start-to-state cost or infimum & \hyperref[sec:p12-best-first-search-and-heuristic-guarantees]{\ref*{sec:p12-best-first-search-and-heuristic-guarantees}} & meta & \hyperref[def:p12-hsearch-instance]{\ref*{def:p12-hsearch-instance}} & \thepage & Actual graph cost; not generally the current label. \\
\phantomsection\label{entry:symbol-registry-02164}\(h\) & \texttt{\detokenize{h}} & nonnegative normalized heuristic estimate & \hyperref[sec:p12-best-first-search-and-heuristic-guarantees]{\ref*{sec:p12-best-first-search-and-heuristic-guarantees}} & meta & \hyperref[def:p12-hsearch-instance]{\ref*{def:p12-hsearch-instance}} & \thepage & Admissibility and consistency are separate properties. \\
\phantomsection\label{entry:symbol-registry-02165}\(h^*\) & \texttt{\detokenize{h^*}} & least remaining goal cost or infinity & \hyperref[sec:p12-best-first-search-and-heuristic-guarantees]{\ref*{sec:p12-best-first-search-and-heuristic-guarantees}} & meta & \hyperref[def:p12-hsearch-instance]{\ref*{def:p12-hsearch-instance}} & \thepage & Infinity is permitted at dead ends. \\
\phantomsection\label{entry:symbol-registry-02166}\(C^*\) & \texttt{\detokenize{C^*}} & optimal cost for the currently specified solution-cost functional & \hyperref[sec:p12-best-first-search-and-heuristic-guarantees]{\ref*{sec:p12-best-first-search-and-heuristic-guarantees}} & meta & \hyperref[def:p12-hsearch-instance]{\ref*{def:p12-hsearch-instance}} & \thepage & Initially additive h*(s); the generalized block declares the minimum under C(P). Attainment is stated before use. \\
\phantomsection\label{entry:symbol-registry-02167}\(W\) & \texttt{\detokenize{W}} & persistent witnessing path record & \hyperref[sec:p12-best-first-search-and-heuristic-guarantees]{\ref*{sec:p12-best-first-search-and-heuristic-guarantees}} & meta & \hyperref[def:p12-hsearch-algorithm]{\ref*{def:p12-hsearch-algorithm}} & \thepage & Old witnesses do not silently change when labels are updated. \\
\phantomsection\label{entry:symbol-registry-02168}\(f\) & \texttt{\detokenize{f}} & A-star selection score g plus h & \hyperref[sec:p12-best-first-search-and-heuristic-guarantees]{\ref*{sec:p12-best-first-search-and-heuristic-guarantees}} & meta & \hyperref[def:p12-hsearch-algorithm]{\ref*{def:p12-hsearch-algorithm}} & \thepage & Not a path cost at arbitrary non-goals. \\
\phantomsection\label{entry:symbol-registry-02169}\(\rho\) & \texttt{\detokenize{\rho}} & map to a relaxed search problem & \hyperref[sec:p12-best-first-search-and-heuristic-guarantees]{\ref*{sec:p12-best-first-search-and-heuristic-guarantees}} & meta & \hyperref[def:p12-hsearch-relaxation]{\ref*{def:p12-hsearch-relaxation}} & \thepage & Must preserve goals and simulate each edge at no greater cost. \\
\phantomsection\label{entry:symbol-registry-02170}\(d_R(x)\) & \texttt{\detokenize{d_R(x)}} & exact relaxed remaining cost as a unary function & \hyperref[def:p12-hsearch-relaxation]{\ref*{def:p12-hsearch-relaxation}} & meta & \hyperref[def:p12-hsearch-relaxation]{\ref*{def:p12-hsearch-relaxation}} & \thepage & An arbitrary relaxed feasible solution is not this value; distinguish the later binary revealed-graph distance. \\
\phantomsection\label{entry:symbol-registry-02171}\(E_A(I)\) & \texttt{\detokenize{E_A(I)}} & set of distinct non-goal states expanded & \hyperref[sec:p12-best-first-search-and-heuristic-guarantees]{\ref*{sec:p12-best-first-search-and-heuristic-guarantees}} & meta & \hyperref[def:p12-hsearch-expansion-count]{\ref*{def:p12-hsearch-expansion-count}} & \thepage & Excludes goal selection without successor generation and does not count re-expansions twice. \\
\phantomsection\label{entry:symbol-registry-02172}\(S(I)\) & \texttt{\detokenize{S(I)}} & states reachable by a strict optimal-cost contour & \hyperref[sec:p12-best-first-search-and-heuristic-guarantees]{\ref*{sec:p12-best-first-search-and-heuristic-guarantees}} & meta & \hyperref[def:p12-hsearch-sure-set]{\ref*{def:p12-hsearch-sure-set}} & \thepage & Initial non-goal state included; goals excluded. \\
\phantomsection\label{entry:symbol-registry-02173}\(F(P)\) & \texttt{\detokenize{F(P)}} & maximum g-plus-h prefix score of a path record & \hyperref[sec:p12-best-first-search-and-heuristic-guarantees]{\ref*{sec:p12-best-first-search-and-heuristic-guarantees}} & meta & \hyperref[def:p12-hsearch-pathmax]{\ref*{def:p12-hsearch-pathmax}} & \thepage & Used for A-double-star selection; duplicate dominance still uses g. \\
\phantomsection\label{entry:symbol-registry-02174}\(V\) & \texttt{\detokenize{V}} & set of graph states & \hyperref[sec:p12-best-first-search-and-heuristic-guarantees]{\ref*{sec:p12-best-first-search-and-heuristic-guarantees}} & meta & \hyperref[def:p12-hsearch-instance]{\ref*{def:p12-hsearch-instance}} & \thepage & The finite-graph theorems assume this set finite. \\
\phantomsection\label{entry:symbol-registry-02175}\(E\) & \texttt{\detokenize{E}} & directed graph edge relation & \hyperref[sec:p12-best-first-search-and-heuristic-guarantees]{\ref*{sec:p12-best-first-search-and-heuristic-guarantees}} & meta & \hyperref[def:p12-hsearch-instance]{\ref*{def:p12-hsearch-instance}} & \thepage & Typed differently from the expansion-set notation \(E_A(I)\). \\
\phantomsection\label{entry:symbol-registry-02176}\(s\) & \texttt{\detokenize{s}} & initial search state & \hyperref[sec:p12-best-first-search-and-heuristic-guarantees]{\ref*{sec:p12-best-first-search-and-heuristic-guarantees}} & meta & \hyperref[def:p12-hsearch-instance]{\ref*{def:p12-hsearch-instance}} & \thepage & A single initial seed is used for oracle-expansion comparisons. \\
\phantomsection\label{entry:symbol-registry-02177}\(\mathrm{OPEN}\) & \texttt{\detokenize{\mathrm{OPEN}}} & current unexpanded or reopened state labels & \hyperref[sec:p12-best-first-search-and-heuristic-guarantees]{\ref*{sec:p12-best-first-search-and-heuristic-guarantees}} & meta & \hyperref[def:p12-hsearch-algorithm]{\ref*{def:p12-hsearch-algorithm}} & \thepage & Obsolete heap entries are not members of this current-label set. \\
\phantomsection\label{entry:symbol-registry-02178}\(\mathrm{CLOSED}\) & \texttt{\detokenize{\mathrm{CLOSED}}} & states processed at their current label & \hyperref[sec:p12-best-first-search-and-heuristic-guarantees]{\ref*{sec:p12-best-first-search-and-heuristic-guarantees}} & meta & \hyperref[def:p12-hsearch-algorithm]{\ref*{def:p12-hsearch-algorithm}} & \thepage & Strictly improving labels reopen a state. \\
\phantomsection\label{entry:symbol-registry-02179}\(\delta\) & \texttt{\detokenize{\delta}} & positive uniform lower bound on edge costs & \hyperref[sec:p12-best-first-search-and-heuristic-guarantees]{\ref*{sec:p12-best-first-search-and-heuristic-guarantees}} & meta & \hyperref[thm:p12-hsearch-infinite]{\ref*{thm:p12-hsearch-infinite}} & \thepage & Not imposed in the finite nonnegative-cost results. \\
\phantomsection\label{entry:symbol-registry-02180}\(\tau\) & \texttt{\detokenize{\tau}} & goal-preferring tie policy & \hyperref[sec:p12-best-first-search-and-heuristic-guarantees]{\ref*{sec:p12-best-first-search-and-heuristic-guarantees}} & meta & \hyperref[def:p12-hsearch-comparison]{\ref*{def:p12-hsearch-comparison}} & \thepage & Its quantifier position is part of every comparison. \\
\phantomsection\label{entry:symbol-registry-02181}\(\mathcal I_{\rm adm}\) & \texttt{\detokenize{\mathcal I_{\rm adm}}} & domain of admissible solvable finite instances & \hyperref[sec:p12-best-first-search-and-heuristic-guarantees]{\ref*{sec:p12-best-first-search-and-heuristic-guarantees}} & meta & \hyperref[def:p12-hsearch-domains]{\ref*{def:p12-hsearch-domains}} & \thepage & Entire-domain correctness defines competitor membership. \\
\phantomsection\label{entry:symbol-registry-02182}\(\mathcal I_{\rm con}\) & \texttt{\detokenize{\mathcal I_{\rm con}}} & consistent subdomain of admissible finite instances & \hyperref[sec:p12-best-first-search-and-heuristic-guarantees]{\ref*{sec:p12-best-first-search-and-heuristic-guarantees}} & meta & \hyperref[def:p12-hsearch-domains]{\ref*{def:p12-hsearch-domains}} & \thepage & Goal normalization is retained. \\
\phantomsection\label{entry:symbol-registry-02183}\(\mathcal A_{\rm adm}\) & \texttt{\detokenize{\mathcal A_{\rm adm}}} & algorithms optimal on all admissible finite instances & \hyperref[sec:p12-best-first-search-and-heuristic-guarantees]{\ref*{sec:p12-best-first-search-and-heuristic-guarantees}} & meta & \hyperref[def:p12-hsearch-domains]{\ref*{def:p12-hsearch-domains}} & \thepage & No ungenerated-graph oracle is available. \\
\phantomsection\label{entry:symbol-registry-02184}\(\mathcal A_{\rm bf}\) & \texttt{\detokenize{\mathcal A_{\rm bf}}} & normalized admissible path-only best-first class & \hyperref[sec:p12-best-first-search-and-heuristic-guarantees]{\ref*{sec:p12-best-first-search-and-heuristic-guarantees}} & meta & \hyperref[def:p12-hsearch-bestfirst-class]{\ref*{def:p12-hsearch-bestfirst-class}} & \thepage & Goal scores equal actual path cost; duplicate merging requires continuation-safe dominance. \\
\phantomsection\label{entry:symbol-registry-02185}\(b(P)\) & \texttt{\detokenize{b(P)}} & normalized competitor path score & \hyperref[sec:p12-best-first-search-and-heuristic-guarantees]{\ref*{sec:p12-best-first-search-and-heuristic-guarantees}} & meta & \hyperref[def:p12-hsearch-bestfirst-class]{\ref*{def:p12-hsearch-bestfirst-class}} & \thepage & Typed by path argument; distinct from the example state named b. \\
\phantomsection\label{entry:symbol-registry-02186}\(H\) & \texttt{\detokenize{H}} & two-input competitor scoring function & \hyperref[sec:p12-best-first-search-and-heuristic-guarantees]{\ref*{sec:p12-best-first-search-and-heuristic-guarantees}} & meta & \hyperref[thm:p12-hsearch-two-input-score]{\ref*{thm:p12-hsearch-two-input-score}} & \thepage & Inputs are current path cost and endpoint heuristic. \\
\phantomsection\label{entry:symbol-registry-02187}\(G_e\) & \texttt{\detokenize{G_e}} & finite graph of exposed successor lists & \hyperref[annex:p12-heuristic-search]{\ref*{annex:p12-heuristic-search}} & meta & \hyperref[annex:p12.hsearch-tradeoff]{\ref*{annex:p12.hsearch-tradeoff}} & \thepage & Unexpanded non-goals have no outgoing edges in the adversarial completion. \\
\phantomsection\label{entry:symbol-registry-02188}\(d_e\) & \texttt{\detokenize{d_e}} & shortest-path distance in the exposed graph & \hyperref[annex:p12-heuristic-search]{\ref*{annex:p12-heuristic-search}} & meta & \hyperref[annex:p12.hsearch-tradeoff]{\ref*{annex:p12.hsearch-tradeoff}} & \thepage & Equality to current labels is established only after all needed relaxations. \\
\phantomsection\label{entry:symbol-registry-02189}\(P\) & \texttt{\detokenize{P}} & finite start-rooted path record & \hyperref[sec:p12-best-first-search-and-heuristic-guarantees]{\ref*{sec:p12-best-first-search-and-heuristic-guarantees}} & meta & \hyperref[def:p12-hsearch-instance]{\ref*{def:p12-hsearch-instance}} & \thepage & A graph state and a path occurrence are different objects. Accepted additive-cost witnesses are simple. \\
\phantomsection\label{entry:symbol-registry-02190}\(P_{s,u}\) & \texttt{\detokenize{P_{s,u}}} & prefix ending at a specified path occurrence & \hyperref[sec:p12-best-first-search-and-heuristic-guarantees]{\ref*{sec:p12-best-first-search-and-heuristic-guarantees}} & meta & \hyperref[def:p12-hsearch-instance]{\ref*{def:p12-hsearch-instance}} & \thepage & Occurrences distinguish repeated appearances of a state in generalized records. \\
\phantomsection\label{entry:symbol-registry-02191}\(h_R\) & \texttt{\detokenize{h_R}} & heuristic induced by an exact relaxed remaining cost & \hyperref[def:p12-hsearch-relaxation]{\ref*{def:p12-hsearch-relaxation}} & meta & \hyperref[def:p12-hsearch-relaxation]{\ref*{def:p12-hsearch-relaxation}} & \thepage & \(h_R(v)=d_R(\rho(v))\); goal and edge simulations supply its certificate. \\
\phantomsection\label{entry:symbol-registry-02192}\(q\) & \texttt{\detokenize{q}} & tentative additive cost in an A-star edge relaxation & \hyperref[def:p12-hsearch-algorithm]{\ref*{def:p12-hsearch-algorithm}} & meta & \hyperref[def:p12-hsearch-algorithm]{\ref*{def:p12-hsearch-algorithm}} & \thepage & Local scalar proposal g(u)+c(u,v), distinct from the later function q(P). \\
\phantomsection\label{entry:symbol-registry-02193}\(A^\dagger\) & \texttt{\detokenize{A^\dagger}} & fully relaxed revealed-graph path-barrier algorithm & \hyperref[sec:p12-hsearch-cost-closure]{\ref*{sec:p12-hsearch-cost-closure}} & meta & \hyperref[def:p12-hsearch-saturated]{\ref*{def:p12-hsearch-saturated}} & \thepage & Explicit strengthened procedure, not an assertion that pending-reopening A** already performs cost closure. \\
\phantomsection\label{entry:symbol-registry-02194}\(R\) & \texttt{\detokenize{R}} & currently revealed directed graph & \hyperref[sec:p12-hsearch-cost-closure]{\ref*{sec:p12-hsearch-cost-closure}} & meta & \hyperref[def:p12-hsearch-saturated]{\ref*{def:p12-hsearch-saturated}} & \thepage & Contains all revealed successor lists and generated vertices; initially one vertex and no edges. \\
\phantomsection\label{entry:symbol-registry-02195}\(D\) & \texttt{\detokenize{D}} & set of non-goals with revealed successor lists & \hyperref[sec:p12-hsearch-cost-closure]{\ref*{sec:p12-hsearch-cost-closure}} & meta & \hyperref[def:p12-hsearch-saturated]{\ref*{def:p12-hsearch-saturated}} & \thepage & Cached-edge propagation through D is not an additional distinct-state expansion. \\
\phantomsection\label{entry:symbol-registry-02196}\(U\) & \texttt{\detokenize{U}} & generated frontier states outside D & \hyperref[sec:p12-hsearch-cost-closure]{\ref*{sec:p12-hsearch-cost-closure}} & meta & \hyperref[def:p12-hsearch-saturated]{\ref*{def:p12-hsearch-saturated}} & \thepage & Goals remain frontier candidates until return. \\
\phantomsection\label{entry:symbol-registry-02197}\(d_R(u,v)\) & \texttt{\detokenize{d_R(u,}\allowbreak\detokenize{v)}} & binary shortest-path distance in the revealed graph R & \hyperref[sec:p12-hsearch-cost-closure]{\ref*{sec:p12-hsearch-cost-closure}} & meta & \hyperref[thm:p12-hsearch-saturated-invariant]{\ref*{thm:p12-hsearch-saturated-invariant}} & \thepage & Different argument type from the earlier unary relaxed goal-distance function. \\
\phantomsection\label{entry:symbol-registry-02198}\(m\) & \texttt{\detokenize{m}} & barrier of the next newly selected state & \hyperref[sec:p12-hsearch-cost-closure]{\ref*{sec:p12-hsearch-cost-closure}} & meta & \hyperref[thm:p12-hsearch-saturated-invariant]{\ref*{thm:p12-hsearch-saturated-invariant}} & \thepage & The positive-cost proof separately handles g<m and g=m. \\
\phantomsection\label{entry:symbol-registry-02199}\(U_m\) & \texttt{\detokenize{U_m}} & flat minimum-barrier frontier with cost equal to m & \hyperref[annex:p12.hsearch-positive-type-three]{\ref*{annex:p12.hsearch-positive-type-three}} & meta & \hyperref[annex:p12.hsearch-positive-type-three]{\ref*{annex:p12.hsearch-positive-type-three}} & \thepage & All members have heuristic zero under the stated positive-domain tie rule. \\
\phantomsection\label{entry:symbol-registry-02200}\(q(P)\) & \texttt{\detokenize{q(P)}} & general finite path evaluation score & \hyperref[sec:p12-hsearch-generalized]{\ref*{sec:p12-hsearch-generalized}} & meta & \hyperref[def:p12-hsearch-general]{\ref*{def:p12-hsearch-general}} & \thepage & Each evaluation terminates; merging requires continuation-safe order preservation. \\
\phantomsection\label{entry:symbol-registry-02201}\(C(P)\) & \texttt{\detokenize{C(P)}} & specified possibly nonadditive solution cost & \hyperref[sec:p12-hsearch-generalized]{\ref*{sec:p12-hsearch-generalized}} & meta & \hyperref[def:p12-hsearch-general]{\ref*{def:p12-hsearch-general}} & \thepage & No relation to the path score is assumed until a strictly increasing goal-score transform is specified. \\
\phantomsection\label{entry:symbol-registry-02202}\(B_q(P)\) & \texttt{\detokenize{B_q(P)}} & largest noninitial score on a goal record & \hyperref[sec:p12-hsearch-generalized]{\ref*{sec:p12-hsearch-generalized}} & meta & \hyperref[def:p12-hsearch-general]{\ref*{def:p12-hsearch-general}} & \thepage & The forced initial selection is excluded from this barrier. \\
\phantomsection\label{entry:symbol-registry-02203}\(\mu\) & \texttt{\detokenize{\mu}} & minimum attained noninitial goal-path barrier & \hyperref[sec:p12-hsearch-generalized]{\ref*{sec:p12-hsearch-generalized}} & meta & \hyperref[def:p12-hsearch-general]{\ref*{def:p12-hsearch-general}} & \thepage & Assumed finite and attained; gives the effective post-root expansion threshold. \\
\phantomsection\label{entry:symbol-registry-02204}\(M\) & \texttt{\detokenize{M}} & minmax including the forced initial score & \hyperref[sec:p12-hsearch-generalized]{\ref*{sec:p12-hsearch-generalized}} & meta & \hyperref[def:p12-hsearch-general]{\ref*{def:p12-hsearch-general}} & \thepage & M=max(q((s)),mu); retaining this source quantity does not make it the sharper expansion threshold. \\
\phantomsection\label{entry:symbol-registry-02205}\(\psi\) & \texttt{\detokenize{\psi}} & strictly increasing map from solution cost to goal score & \hyperref[sec:p12-hsearch-generalized]{\ref*{sec:p12-hsearch-generalized}} & meta & \hyperref[thm:p12-hsearch-general-barrier]{\ref*{thm:p12-hsearch-general-barrier}} & \thepage & The inverse bound is used only where its inverse value exists. \\
\phantomsection\label{entry:symbol-registry-02206}\(p(P)\) & \texttt{\detokenize{p(P)}} & path-dependent modulation of heuristic weight & \hyperref[thm:p12-hsearch-weighted]{\ref*{thm:p12-hsearch-weighted}} & meta & \hyperref[thm:p12-hsearch-weighted]{\ref*{thm:p12-hsearch-weighted}} & \thepage & Only its upper bound along an optimal path is used for the stated return-quality ratio. \\
\phantomsection\label{entry:symbol-registry-02207}\(\varepsilon\) & \texttt{\detokenize{\varepsilon}} & nonnegative relative-error factor in the weighted score & \hyperref[thm:p12-hsearch-weighted]{\ref*{thm:p12-hsearch-weighted}} & meta & \hyperref[thm:p12-hsearch-weighted]{\ref*{thm:p12-hsearch-weighted}} & \thepage & Adversarial proofs separately quantify positive local perturbations of the same glyph. \\
\phantomsection\label{entry:symbol-registry-02208}\(\beta\) & \texttt{\detokenize{\beta}} & multiplicative upper bound on heuristic error along an optimal path & \hyperref[thm:p12-hsearch-weighted]{\ref*{thm:p12-hsearch-weighted}} & meta & \hyperref[thm:p12-hsearch-weighted]{\ref*{thm:p12-hsearch-weighted}} & \thepage & May exceed one; the separate weighted admissibility condition remains explicit. \\
\phantomsection\label{entry:symbol-registry-02209}\(w\) & \texttt{\detokenize{w}} & convex weighting parameter of g and h & \hyperref[thm:p12-hsearch-weighted]{\ref*{thm:p12-hsearch-weighted}} & meta & \hyperref[thm:p12-hsearch-weighted]{\ref*{thm:p12-hsearch-weighted}} & \thepage & Zero is allowed; one is excluded because normalized goal scores must preserve strict cost order. \\
\phantomsection\label{entry:symbol-registry-02210}\(\phi\) & \texttt{\detokenize{\phi}} & increasing concave bijection of nonnegative reals fixing zero & \hyperref[thm:p12-hsearch-weighted]{\ref*{thm:p12-hsearch-weighted}} & meta & \hyperref[thm:p12-hsearch-weighted]{\ref*{thm:p12-hsearch-weighted}} & \thepage & Its inverse supplies the nonlinear score construction. \\
\phantomsection\label{entry:symbol-registry-02211}\(r\) & \texttt{\detokenize{r}} & inverse function of the concave score generator phi & \hyperref[annex:p12.hsearch-weighted]{\ref*{annex:p12.hsearch-weighted}} & meta & \hyperref[annex:p12.hsearch-weighted]{\ref*{annex:p12.hsearch-weighted}} & \thepage & Local function role; not the ordered-tree vertex r or right-side barrier r(n). \\
\phantomsection\label{entry:symbol-registry-02212}\(L_i\) & \texttt{\detokenize{L_i}} & maximum future score on the returned record beyond position i & \hyperref[thm:p12-hsearch-returned-path]{\ref*{thm:p12-hsearch-returned-path}} & meta & \hyperref[thm:p12-hsearch-returned-path]{\ref*{thm:p12-hsearch-returned-path}} & \thepage & Evaluations on divergent continuations retain the whole start-rooted prefix. \\
\phantomsection\label{entry:symbol-registry-02213}\(T\) & \texttt{\detokenize{T}} & finite rooted tree with ordered children and terminal goals & \hyperref[sec:p12-hsearch-ordered-tree]{\ref*{sec:p12-hsearch-ordered-tree}} & meta & \hyperref[def:p12-hsearch-tree-convention]{\ref*{def:p12-hsearch-tree-convention}} & \thepage & Pure-leftmost tie priority applies only in this section. \\
\phantomsection\label{entry:symbol-registry-02214}\(q(u)\) & \texttt{\detokenize{q(u)}} & score at the unique root-to-u record in the ordered tree & \hyperref[sec:p12-hsearch-ordered-tree]{\ref*{sec:p12-hsearch-ordered-tree}} & meta & \hyperref[def:p12-hsearch-tree-convention]{\ref*{def:p12-hsearch-tree-convention}} & \thepage & Abbreviation of q(P) available because the tree path is unique. \\
\phantomsection\label{entry:symbol-registry-02215}\(m(r)\) & \texttt{\detokenize{m(r)}} & minimum goal-path maximum inside a subtree rooted at r & \hyperref[sec:p12-hsearch-ordered-tree]{\ref*{sec:p12-hsearch-ordered-tree}} & meta & \hyperref[def:p12-hsearch-tree-convention]{\ref*{def:p12-hsearch-tree-convention}} & \thepage & Both endpoints included; infinity denotes a subtree with no goal. \\
\phantomsection\label{entry:symbol-registry-02216}\(\ell(n)\) & \texttt{\detokenize{\ell(n)}} & minimum barrier of eligible left side subtrees & \hyperref[sec:p12-hsearch-ordered-tree]{\ref*{sec:p12-hsearch-ordered-tree}} & meta & \hyperref[def:p12-hsearch-tree-convention]{\ref*{def:p12-hsearch-tree-convention}} & \thepage & Side branches occur at strict ancestors of n, with empty minimum infinity. \\
\phantomsection\label{entry:symbol-registry-02217}\(r(n)\) & \texttt{\detokenize{r(n)}} & minimum barrier of eligible right side subtrees & \hyperref[sec:p12-hsearch-ordered-tree]{\ref*{sec:p12-hsearch-ordered-tree}} & meta & \hyperref[def:p12-hsearch-tree-convention]{\ref*{def:p12-hsearch-tree-convention}} & \thepage & Function notation differs from the locally bound subtree root r. \\
\phantomsection\label{entry:symbol-registry-02218}\(D\) & \texttt{\detokenize{D}} & largest strict-contour score in the nonpathological trade proof & \hyperref[annex:p12.hsearch-tradeoff]{\ref*{annex:p12.hsearch-tradeoff}} & meta & \hyperref[annex:p12.hsearch-tradeoff]{\ref*{annex:p12.hsearch-tradeoff}} & \thepage & Local real scalar, distinct from the expanded-set D of the saturated algorithm. \\
\phantomsection\label{entry:symbol-registry-02219}\(Q\) & \texttt{\detokenize{Q}} & maximum A-star prefix score on a competing-score witness path & \hyperref[annex:p12.hsearch-bestfirst-forcing]{\ref*{annex:p12.hsearch-bestfirst-forcing}} & meta & \hyperref[annex:p12.hsearch-bestfirst-forcing]{\ref*{annex:p12.hsearch-bestfirst-forcing}} & \thepage & Includes the initial score under the explicit h(s)<C* premise. \\
\phantomsection\label{entry:symbol-registry-02220}\(Q_b\) & \texttt{\detokenize{Q_b}} & largest noninitial competitor score on that witness path & \hyperref[annex:p12.hsearch-bestfirst-forcing]{\ref*{annex:p12.hsearch-bestfirst-forcing}} & meta & \hyperref[annex:p12.hsearch-bestfirst-forcing]{\ref*{annex:p12.hsearch-bestfirst-forcing}} & \thepage & The two-goal counterexample places a goal cost strictly between \(Q\) and \(Q_b\). \\
\phantomsection\label{entry:symbol-registry-02221}\(M_P\) & \texttt{\detokenize{M_P}} & largest A-star score along one fixed solution path & \hyperref[annex:p12.hsearch-infinite]{\ref*{annex:p12.hsearch-infinite}} & meta & \hyperref[annex:p12.hsearch-infinite]{\ref*{annex:p12.hsearch-infinite}} & \thepage & A finite upper bound used even when the heuristic overestimates. \\
\phantomsection\label{entry:symbol-registry-02222}\(K\) & \texttt{\detokenize{K}} & cost of a cheapest goal path in the currently revealed graph & \hyperref[annex:p12.hsearch-pending-domination]{\ref*{annex:p12.hsearch-pending-domination}} & meta & \hyperref[annex:p12.hsearch-pending-domination]{\ref*{annex:p12.hsearch-pending-domination}} & \thepage & The early-return certificate compares this feasible cost with the simulated OPEN barrier. \\
\phantomsection\label{entry:symbol-registry-02223}\(\varepsilon\) & \texttt{\detokenize{\varepsilon}} & strictly positive appended-edge slack & \hyperref[annex:p12.hsearch-consistent-efficiency]{\ref*{annex:p12.hsearch-consistent-efficiency}} & meta & \hyperref[annex:p12.hsearch-consistent-efficiency]{\ref*{annex:p12.hsearch-consistent-efficiency}} & \thepage & Chosen below the remaining strict cost gap. \\
\phantomsection\label{entry:symbol-registry-02224}\(\varepsilon\) & \texttt{\detokenize{\varepsilon}} & strictly positive cost gap in an exposed-graph completion & \hyperref[annex:p12.hsearch-tradeoff]{\ref*{annex:p12.hsearch-tradeoff}} & meta & \hyperref[annex:p12.hsearch-tradeoff]{\ref*{annex:p12.hsearch-tradeoff}} & \thepage & Chosen between zero and C*-D. \\
\phantomsection\label{entry:symbol-registry-02225}\(\varepsilon\) & \texttt{\detokenize{\varepsilon}} & strictly positive perturbation for costwise comparison & \hyperref[annex:p12.hsearch-cost-compatibility]{\ref*{annex:p12.hsearch-cost-compatibility}} & meta & \hyperref[annex:p12.hsearch-cost-compatibility]{\ref*{annex:p12.hsearch-cost-compatibility}} & \thepage & The modified instance need not retain heuristic admissibility. \\
\phantomsection\label{entry:symbol-registry-02226}\(\varepsilon\) & \texttt{\detokenize{\varepsilon}} & small strictly positive terminal-edge cost for a flat frontier group & \hyperref[annex:p12.hsearch-positive-type-three]{\ref*{annex:p12.hsearch-positive-type-three}} & meta & \hyperref[annex:p12.hsearch-positive-type-three]{\ref*{annex:p12.hsearch-positive-type-three}} & \thepage & Chosen below every finite positive gap to the current barrier. \\
\phantomsection\label{entry:symbol-registry-02227}\(\varepsilon'\) & \texttt{\detokenize{\varepsilon'}} & smaller positive terminal-edge perturbation & \hyperref[annex:p12.hsearch-positive-type-three]{\ref*{annex:p12.hsearch-positive-type-three}} & meta & \hyperref[annex:p12.hsearch-positive-type-three]{\ref*{annex:p12.hsearch-positive-type-three}} & \thepage & Used to force inspection of every still unqueried member of the flat group. \\
\phantomsection\label{entry:symbol-registry-02228}\(\varepsilon\) & \texttt{\detokenize{\varepsilon}} & positive perturbation keeping a comparison instance inside the positive-cost domain & \hyperref[annex:p12.hsearch-positive-domain-extension]{\ref*{annex:p12.hsearch-positive-domain-extension}} & meta & \hyperref[annex:p12.hsearch-positive-domain-extension]{\ref*{annex:p12.hsearch-positive-domain-extension}} & \thepage & Its bound preserves the strict inequalities in the finite counterexample. \\
\phantomsection\label{entry:symbol-registry-02229}\(X_d\) & \texttt{\detokenize{X_d}} & domain of candidate unknowns for fixed data & \hyperref[sec:p12-heuristic-transformations-and-backward-reasoning]{\ref*{sec:p12-heuristic-transformations-and-backward-reasoning}} & meta & \hyperref[def:p12-heuristic-problem]{\ref*{def:p12-heuristic-problem}} & \thepage & The domain is part of the original problem and survives transformations. \\
\phantomsection\label{entry:symbol-registry-02230}\(C_d\) & \texttt{\detokenize{C_d}} & condition on a candidate with fixed data & \hyperref[sec:p12-heuristic-transformations-and-backward-reasoning]{\ref*{sec:p12-heuristic-transformations-and-backward-reasoning}} & meta & \hyperref[def:p12-heuristic-problem]{\ref*{def:p12-heuristic-problem}} & \thepage & Includes all original side conditions. \\
\phantomsection\label{entry:symbol-registry-02231}\(S_d\) & \texttt{\detokenize{S_d}} & original solution set & \hyperref[sec:p12-heuristic-transformations-and-backward-reasoning]{\ref*{sec:p12-heuristic-transformations-and-backward-reasoning}} & meta & \hyperref[def:p12-heuristic-problem]{\ref*{def:p12-heuristic-problem}} & \thepage & One solution and all solutions are distinct requested outputs. \\
\phantomsection\label{entry:symbol-registry-02232}\(S_B\) & \texttt{\detokenize{S_B}} & solution set of condition B in the fixed original domain & \hyperref[sec:p12-heuristic-transformations-and-backward-reasoning]{\ref*{sec:p12-heuristic-transformations-and-backward-reasoning}} & meta & \hyperref[def:p12-heuristic-transformation]{\ref*{def:p12-heuristic-transformation}} & \thepage & Subscripts \(C,N,H,E\) and \(C_i\) instantiate this notation. \\
\phantomsection\label{entry:symbol-registry-02233}\(\Gamma\) & \texttt{\detokenize{\Gamma}} & fixed proof assumptions & \hyperref[sec:p12-heuristic-transformations-and-backward-reasoning]{\ref*{sec:p12-heuristic-transformations-and-backward-reasoning}} & meta & \hyperref[def:p12-heuristic-problem]{\ref*{def:p12-heuristic-problem}} & \thepage & Temporary discovery assumptions are separately discharged. \\
\phantomsection\label{entry:symbol-registry-02234}\(T\) & \texttt{\detokenize{T}} & target assertion & \hyperref[sec:p12-heuristic-transformations-and-backward-reasoning]{\ref*{sec:p12-heuristic-transformations-and-backward-reasoning}} & meta & \hyperref[def:p12-heuristic-problem]{\ref*{def:p12-heuristic-problem}} & \thepage & The proof system and meaning of the assertion are specified. \\
\phantomsection\label{entry:symbol-registry-02235}\(N\) & \texttt{\detokenize{N}} & necessary transformed condition & \hyperref[def:p12-heuristic-transformation]{\ref*{def:p12-heuristic-transformation}} & meta & \hyperref[def:p12-heuristic-transformation]{\ref*{def:p12-heuristic-transformation}} & \thepage & Original solutions satisfy it; candidates need a reverse check. \\
\phantomsection\label{entry:symbol-registry-02236}\(H\) & \texttt{\detokenize{H}} & sufficient transformed condition & \hyperref[def:p12-heuristic-transformation]{\ref*{def:p12-heuristic-transformation}} & meta & \hyperref[def:p12-heuristic-transformation]{\ref*{def:p12-heuristic-transformation}} & \thepage & This local predicate is distinct from the longest-altitude H in the triangle example. \\
\phantomsection\label{entry:symbol-registry-02237}\(E\) & \texttt{\detokenize{E}} & equivalent transformed condition & \hyperref[def:p12-heuristic-transformation]{\ref*{def:p12-heuristic-transformation}} & meta & \hyperref[def:p12-heuristic-transformation]{\ref*{def:p12-heuristic-transformation}} & \thepage & Both implication directions must be established. \\
\phantomsection\label{entry:symbol-registry-02238}\(F\) & \texttt{\detokenize{F}} & auxiliary-solution reconstruction map & \hyperref[thm:p12-heuristic-solution-relations]{\ref*{thm:p12-heuristic-solution-relations}} & meta & \hyperref[thm:p12-heuristic-solution-relations]{\ref*{thm:p12-heuristic-solution-relations}} & \thepage & A map back to original candidates with proved soundness; all-solutions transfer needs surjectivity. \\
\phantomsection\label{entry:symbol-registry-02239}\(G\) & \texttt{\detokenize{G}} & right inverse witnessing full solution coverage & \hyperref[thm:p12-heuristic-solution-relations]{\ref*{thm:p12-heuristic-solution-relations}} & meta & \hyperref[thm:p12-heuristic-solution-relations]{\ref*{thm:p12-heuristic-solution-relations}} & \thepage & Distinct from the goal set G in transition reasoning. \\
\phantomsection\label{entry:symbol-registry-02240}\(\operatorname{Pre}_R\) & \texttt{\detokenize{\operatorname{Pre}_R}} & existential predecessor operation & \hyperref[sec:p12-heuristic-transformations-and-backward-reasoning]{\ref*{sec:p12-heuristic-transformations-and-backward-reasoning}} & meta & \hyperref[def:p12-heuristic-predecessor]{\ref*{def:p12-heuristic-predecessor}} & \thepage & Stores witnesses for legal forward transitions. \\
\phantomsection\label{entry:symbol-registry-02241}\(R\) & \texttt{\detokenize{R}} & legal one-step transition relation & \hyperref[def:p12-heuristic-predecessor]{\ref*{def:p12-heuristic-predecessor}} & meta & \hyperref[def:p12-heuristic-predecessor]{\ref*{def:p12-heuristic-predecessor}} & \thepage & Local to state/path arguments; not the circumradius R. \\
\phantomsection\label{entry:symbol-registry-02242}\(G_k\) & \texttt{\detokenize{G_k}} & states reaching a goal in exactly k legal steps & \hyperref[thm:p12-heuristic-backward-path]{\ref*{thm:p12-heuristic-backward-path}} & meta & \hyperref[thm:p12-heuristic-backward-path]{\ref*{thm:p12-heuristic-backward-path}} & \thepage & \(G_0\) is the goal set; iteration uses existential predecessors. \\
\phantomsection\label{entry:symbol-registry-02243}\((L,S)\) & \texttt{\detokenize{(L,}\allowbreak\detokenize{S)}} & amounts in nine- and four-quart containers & \hyperref[thm:p12-heuristic-pails]{\ref*{thm:p12-heuristic-pails}} & meta & \hyperref[thm:p12-heuristic-pails]{\ref*{thm:p12-heuristic-pails}} & \thepage & First coordinate is the larger container. \\
\phantomsection\label{entry:symbol-registry-02244}\(p_N\) & \texttt{\detokenize{p_N}} & polynomial separating finite samples from a universal identity & \hyperref[thm:p12-heuristic-generalization]{\ref*{thm:p12-heuristic-generalization}} & meta & \hyperref[thm:p12-heuristic-generalization]{\ref*{thm:p12-heuristic-generalization}} & \thepage & Product of x-j for j from one through N. \\
\phantomsection\label{entry:symbol-registry-02245}\(r\) & \texttt{\detokenize{r}} & inradius of the specified Euclidean triangle & \hyperref[thm:p12-heuristic-triangle-counterexample]{\ref*{thm:p12-heuristic-triangle-counterexample}} & meta & \hyperref[thm:p12-heuristic-triangle-counterexample]{\ref*{thm:p12-heuristic-triangle-counterexample}} & \thepage & Positive for the stated nondegenerate triangles. \\
\phantomsection\label{entry:symbol-registry-02246}\(R\) & \texttt{\detokenize{R}} & circumradius of the specified Euclidean triangle & \hyperref[thm:p12-heuristic-triangle-counterexample]{\ref*{thm:p12-heuristic-triangle-counterexample}} & meta & \hyperref[thm:p12-heuristic-triangle-counterexample]{\ref*{thm:p12-heuristic-triangle-counterexample}} & \thepage & Distinct local geometric role from the transition relation. \\
\phantomsection\label{entry:symbol-registry-02247}\(H\) & \texttt{\detokenize{H}} & longest of the triangle's three altitudes & \hyperref[thm:p12-heuristic-triangle-counterexample]{\ref*{thm:p12-heuristic-triangle-counterexample}} & meta & \hyperref[thm:p12-heuristic-triangle-counterexample]{\ref*{thm:p12-heuristic-triangle-counterexample}} & \thepage & Not necessarily the altitude to the displayed base. \\
\phantomsection\label{entry:symbol-registry-02248}\(d\) & \texttt{\detokenize{d}} & initial relative position vector & \hyperref[thm:p12-heuristic-ships]{\ref*{thm:p12-heuristic-ships}} & meta & \hyperref[thm:p12-heuristic-ships]{\ref*{thm:p12-heuristic-ships}} & \thepage & Local definition d=a-b; distinct from the generic data record d. \\
\phantomsection\label{entry:symbol-registry-02249}\(w\) & \texttt{\detokenize{w}} & relative velocity vector & \hyperref[thm:p12-heuristic-ships]{\ref*{thm:p12-heuristic-ships}} & meta & \hyperref[thm:p12-heuristic-ships]{\ref*{thm:p12-heuristic-ships}} & \thepage & The zero-relative-velocity case is separate. \\
\phantomsection\label{entry:symbol-registry-02250}\(t_*\) & \texttt{\detokenize{t_*}} & unrestricted minimizing time & \hyperref[thm:p12-heuristic-ships]{\ref*{thm:p12-heuristic-ships}} & meta & \hyperref[thm:p12-heuristic-ships]{\ref*{thm:p12-heuristic-ships}} & \thepage & For future times the candidate is \(\max(0,t_*)\). \\
\phantomsection\label{entry:symbol-registry-02251}\(t|_p,t[u]_p\) & \texttt{\detokenize{t|_p,}\allowbreak\detokenize{t[u]_p}} & subterm selection and replacement at a valid position & \hyperref[sec:p12-rewriting]{\ref*{sec:p12-rewriting}} & meta & \hyperref[def:p12-rewrite-positions]{\ref*{def:p12-rewrite-positions}} & \thepage & Positions are finite words of child indices; replacement changes one selected occurrence. \\
\phantomsection\label{entry:symbol-registry-02252}\(\epsilon,p,q\) & \texttt{\detokenize{\epsilon,}\allowbreak\detokenize{p,}\allowbreak\detokenize{q}} & empty root position and term-position words & \hyperref[sec:p12-rewriting]{\ref*{sec:p12-rewriting}} & meta & \hyperref[def:p12-rewrite-positions]{\ref*{def:p12-rewrite-positions}} & \thepage & Prefix comparison separates disjoint and nested redexes. \\
\phantomsection\label{entry:symbol-registry-02253}\(\downarrow_R\) & \texttt{\detokenize{\downarrow_R}} & joinability by finite forward reductions & \hyperref[sec:p12-rewriting]{\ref*{sec:p12-rewriting}} & relation & \hyperref[def:p12-rewrite-positions]{\ref*{def:p12-rewrite-positions}} & \thepage & Both terms reach one common descendant; zero steps are permitted. \\
\phantomsection\label{entry:symbol-registry-02254}\(\theta,\delta\) & \texttt{\detokenize{\theta,}\allowbreak\detokenize{\delta}} & most general unifier and residual substitution in critical-overlap factorization & \hyperref[sec:p12-rewriting]{\ref*{sec:p12-rewriting}} & meta & \hyperref[def:p12-critical-pair]{\ref*{def:p12-critical-pair}} & \thepage & Rule variables are renamed apart; the selected MGU fixes variables outside the unified terms. \\
\phantomsection\label{entry:symbol-registry-02255}\(\sigma,\sigma'\) & \texttt{\detokenize{\sigma,}\allowbreak\detokenize{\sigma'}} & substitutions before and after changing one variable image & \hyperref[sec:p12-rewriting]{\ref*{sec:p12-rewriting}} & meta & \hyperref[thm:p12-rewrite-closure]{\ref*{thm:p12-rewrite-closure}} & \thepage & The synchronization lemma includes repeated occurrences and an erased variable. \\
\phantomsection\label{entry:symbol-registry-02256}\(\leftrightarrow_R^*\) & \texttt{\detokenize{\leftrightarrow_R^*}} & finite conversion by forward and backward rewrite steps & \hyperref[sec:p12-rewriting]{\ref*{sec:p12-rewriting}} & relation & \hyperref[thm:p12-rewrite-normal-forms]{\ref*{thm:p12-rewrite-normal-forms}} & \thepage & Distinct from joinability before confluence is established. \\
\phantomsection\label{entry:symbol-registry-02257}\(\operatorname{nf}_R(t)\) & \texttt{\detokenize{\operatorname{nf}_R(t)}} & unique normal form of a term under terminating confluent rewriting & \hyperref[sec:p12-rewriting]{\ref*{sec:p12-rewriting}} & meta & \hyperref[thm:p12-rewrite-normal-forms]{\ref*{thm:p12-rewrite-normal-forms}} & \thepage & Existence and uniqueness precede use as a decision result. \\
\phantomsection\label{entry:symbol-registry-02258}\(\equiv_Q\) & \texttt{\detokenize{\equiv_Q}} & equivalence closed under contexts and substitutions generated by equations Q & \hyperref[sec:p12-rewriting]{\ref*{sec:p12-rewriting}} & relation & \hyperref[def:p12-completion]{\ref*{def:p12-completion}} & \thepage & This source-local congruence notation is not elementary equivalence of first-order structures. \\
\phantomsection\label{entry:symbol-registry-02259}\(\rho\) & \texttt{\detokenize{\rho}} & combined matching substitution for renamed rule variables & \hyperref[annex:p12.rewriting.critical-pairs]{\ref*{annex:p12.rewriting.critical-pairs}} & meta & \hyperref[annex:p12.rewriting.critical-pairs]{\ref*{annex:p12.rewriting.critical-pairs}} & \thepage & Factorization through the MGU is needed on both full rule-variable sets. \\
\phantomsection\label{entry:symbol-registry-02260}\(\rho\) & \texttt{\detokenize{\rho}} & assignment of variables to natural numbers for a termination interpretation & \hyperref[annex:p12.rewriting.example]{\ref*{annex:p12.rewriting.example}} & meta & \hyperref[annex:p12.rewriting.example]{\ref*{annex:p12.rewriting.example}} & \thepage & This separate local role is not the matching substitution in the critical-pair proof. \\
\phantomsection\label{entry:symbol-registry-02261}\(\llbracket t\rrbracket_\rho\) & \texttt{\detokenize{\llbracket t\rrbracket_\rho}} & natural-number evaluation of a term in the worked ordering & \hyperref[annex:p12.rewriting.example]{\ref*{annex:p12.rewriting.example}} & meta & \hyperref[annex:p12.rewriting.example]{\ref*{annex:p12.rewriting.example}} & \thepage & The interpretation fixes three constants and one strictly increasing unary operation. \\
\phantomsection\label{entry:symbol-registry-02262}\(x\) & \texttt{\detokenize{x}} & specified query condition for a model-to-target comparison & P12 simulation-results-as-evidence error-chain section & meta & \hyperref[def:p12-simulation-error-chain]{\ref*{def:p12-simulation-error-chain}} & \thepage & The target and every model/referent comparison use the same declared conditions. \\
\phantomsection\label{entry:symbol-registry-02263}\(Y(x)\) & \texttt{\detokenize{Y(x)}} & quantity of interest in the target system & P12 simulation-results-as-evidence & meta & \hyperref[def:p12-simulation-error-chain]{\ref*{def:p12-simulation-error-chain}} & \thepage & An empirical target is distinct from the exact mathematical model. \\
\phantomsection\label{entry:symbol-registry-02264}\(m,\theta_*,\widehat\theta\) & \texttt{\detokenize{m,}\allowbreak\detokenize{\theta_*,}\allowbreak\detokenize{\widehat\theta}} & exact model and target/supplied inputs or parameters & P12 simulation-results-as-evidence error-chain section & meta & \hyperref[def:p12-simulation-error-chain]{\ref*{def:p12-simulation-error-chain}} & \thepage & Their distinction identifies the input-error contribution. \\
\phantomsection\label{entry:symbol-registry-02265}\(m_h,\widehat m_h\) & \texttt{\detokenize{m_h,}\allowbreak\detokenize{\widehat m_h}} & discretized mathematical calculation and computed output & P12 simulation-results-as-evidence error-chain section & meta & \hyperref[def:p12-simulation-error-chain]{\ref*{def:p12-simulation-error-chain}} & \thepage & h is a discretization parameter; hats denote computed values. \\
\phantomsection\label{entry:symbol-registry-02266}\(E_{\rm model},E_{\rm input},E_{\rm disc},E_{\rm arith}\) & \texttt{\detokenize{E_{\rm model},}\allowbreak\detokenize{E_{\rm input},}\allowbreak\detokenize{E_{\rm disc},}\allowbreak\detokenize{E_{\rm arith}}} & nonnegative bounds on four explicitly named differences & P12 simulation-results-as-evidence & meta & \hyperref[def:p12-simulation-error-chain]{\ref*{def:p12-simulation-error-chain}} & \thepage & Worst-case premises; unavailable quantities are not assigned zero. \\
\phantomsection\label{entry:symbol-registry-02267}\(r,E_r,E_n,\tau\) & \texttt{\detokenize{r,}\allowbreak\detokenize{E_r,}\allowbreak\detokenize{E_n,}\allowbreak\detokenize{\tau}} & referent value, referent and numerical error radii, and declared discrepancy tolerance & P12 simulation-results-as-evidence comparison section & meta & \hyperref[thm:p12-simulation-reference-comparison]{\ref*{thm:p12-simulation-reference-comparison}} & \thepage & The same units and query condition are required. \\
\phantomsection\label{entry:symbol-registry-02268}\(q,q_0,\lambda,t\) & \texttt{\detokenize{q,}\allowbreak\detokenize{q_0,}\allowbreak\detokenize{\lambda,}\allowbreak\detokenize{t}} & decaying quantity, initial amount, nonnegative rate, and time & P12 simulation-results-as-evidence decay section & meta & \hyperref[def:p12-simulation-decay]{\ref*{def:p12-simulation-decay}} & \thepage & lambda has reciprocal-time units; lambda t is dimensionless. \\
\phantomsection\label{entry:symbol-registry-02269}\(h,N,T,a\) & \texttt{\detokenize{h,}\allowbreak\detokenize{N,}\allowbreak\detokenize{T,}\allowbreak\detokenize{a}} & time step, number of steps, final time and dimensionless step rate & P12 simulation-results-as-evidence decay section & meta & \hyperref[def:p12-simulation-decay]{\ref*{def:p12-simulation-decay}} & \thepage & T=Nh and a=lambda h; step count N is a nonnegative integer. \\
\phantomsection\label{entry:symbol-registry-02270}\(v_j,\widehat v_j\) & \texttt{\detokenize{v_j,}\allowbreak\detokenize{\widehat v_j}} & exact-real Euler iterate and computed iterate & P12 simulation-results-as-evidence decay section & meta & \hyperref[def:p12-simulation-decay]{\ref*{def:p12-simulation-decay}} & \thepage & The computed recurrence is separately specified in thm:p12-simulation-arithmetic-propagation. \\
\phantomsection\label{entry:symbol-registry-02271}\(\delta_j,\rho_j,\rho,E_0\) & \texttt{\detokenize{\delta_j,}\allowbreak\detokenize{\rho_j,}\allowbreak\detokenize{\rho,}\allowbreak\detokenize{E_0}} & additive arithmetic perturbations and their bounds & P12 simulation-results-as-evidence arithmetic-propagation section & meta & \hyperref[thm:p12-simulation-arithmetic-propagation]{\ref*{thm:p12-simulation-arithmetic-propagation}} & \thepage & Coefficient rounding is included in the actual perturbation budget; no independence assumption. \\
\phantomsection\label{entry:symbol-registry-02272}\(Q,\Lambda\) & \texttt{\detokenize{Q,}\allowbreak\detokenize{\Lambda}} & upper bounds for initial quantity and decay rate in a sensitivity rectangle & P12 simulation-results-as-evidence decay sensitivity section & meta & \hyperref[thm:p12-simulation-decay-inputs]{\ref*{thm:p12-simulation-decay-inputs}} & \thepage & These letters are local bounds, not probability distributions. \\
\phantomsection\label{entry:symbol-registry-02273}\(\widetilde q_0,\widetilde\lambda,\widetilde t\) & \texttt{\detokenize{\widetilde q_0,}\allowbreak\detokenize{\widetilde\lambda,}\allowbreak\detokenize{\widetilde t}} & simultaneously perturbed inputs in the same declared rectangle & P12 simulation-results-as-evidence decay sensitivity section & meta & \hyperref[thm:p12-simulation-decay-inputs]{\ref*{thm:p12-simulation-decay-inputs}} & \thepage & The bound applies to joint changes within that rectangle. \\
\phantomsection\label{entry:symbol-registry-02274}\(B,\widehat B,\widehat\lambda,y,e\) & \texttt{\detokenize{B,}\allowbreak\detokenize{\widehat B,}\allowbreak\detokenize{\widehat\lambda,}\allowbreak\detokenize{y,}\allowbreak\detokenize{e}} & BoxBOD asymptote and reported fit parameters, response and residual-model error term & P12 simulation-results-as-evidence BoxBOD section & meta & \hyperref[def:p12-simulation-boxbod]{\ref*{def:p12-simulation-boxbod}} & \thepage & B and y use mg/l; lambda uses reciprocal days; e in the statistical model is not Euler's constant. \\
\phantomsection\label{entry:symbol-registry-02275}\(t_i,s,f,P,g\) & \texttt{\detokenize{t_i,}\allowbreak\detokenize{s,}\allowbreak\detokenize{f,}\allowbreak\detokenize{P,}\allowbreak\detokenize{g}} & comparison points, new point, baseline function, vanishing polynomial and alternative function & P12 simulation-results-as-evidence finite-comparison section & meta & \hyperref[thm:p12-simulation-finite-comparisons]{\ref*{thm:p12-simulation-finite-comparisons}} & \thepage & P and g are introduced in the associated annex proof. \\
\phantomsection\label{entry:symbol-registry-02276}\(q_1,q_2,\mu\) & \texttt{\detokenize{q_1,}\allowbreak\detokenize{q_2,}\allowbreak\detokenize{\mu}} & two decay trajectories and the post-change positive rate & P12 simulation-results-as-evidence transfer section & meta & \hyperref[thm:p12-simulation-regime-change]{\ref*{thm:p12-simulation-regime-change}} & \thepage & The two rates differ; both trajectories agree at and before T. \\
\phantomsection\label{entry:symbol-registry-02277}\(A,B\) & \texttt{\detokenize{A,}\allowbreak\detokenize{B}} & two ordinal assessment profiles in the counterexample & P12 simulation-results-as-evidence ordinal-score annex & meta & \hyperref[annex:p12.simulation.ordinal-scores]{\ref*{annex:p12.simulation.ordinal-scores}} & \thepage & Local profile labels; distinct from BoxBOD's asymptote parameter. \\
\phantomsection\label{entry:symbol-registry-02428}\(>_\Sigma\) & \texttt{\detokenize{>_\Sigma}} & precedence on ranked function symbols & \hyperref[sec:p12-reduction-orderings]{\ref*{sec:p12-reduction-orderings}} & relation & \hyperref[def:p12-ro-data]{\ref*{def:p12-ro-data}} & \thepage & Partial for the selected LPO; strict total for the historical KBO; variables are outside this precedence. \\
\phantomsection\label{entry:symbol-registry-02429}\(\pi_f\) & \texttt{\detokenize{\pi_f}} & fixed permutation specifying lexicographic argument status & \hyperref[sec:p12-reduction-orderings]{\ref*{sec:p12-reduction-orderings}} & meta & \hyperref[def:p12-ro-data]{\ref*{def:p12-ro-data}} & \thepage & The original KBO uses the left-to-right status; LPO permits the declared permutations. \\
\phantomsection\label{entry:symbol-registry-02430}\(\operatorname{occ}_x(t)\) & \texttt{\detokenize{\operatorname{occ}_x(t)}} & number of variable occurrences in a finite term & \hyperref[sec:p12-reduction-orderings]{\ref*{sec:p12-reduction-orderings}} & meta & \hyperref[def:p12-ro-data]{\ref*{def:p12-ro-data}} & \thepage & Counts are multiplicities, not merely membership in a variable set. \\
\phantomsection\label{entry:symbol-registry-02431}\(\lvert t\rvert\) & \texttt{\detokenize{\lvert t\rvert}} & number of nodes in a finite term & \hyperref[sec:p12-reduction-orderings]{\ref*{sec:p12-reduction-orderings}} & meta & \hyperref[def:p12-ro-data]{\ref*{def:p12-ro-data}} & \thepage & Explicit local syntactic-size use of bars; not a claim that terms are sets. \\
\phantomsection\label{entry:symbol-registry-02432}\(>_{\rm lpo},\geq_{\rm lpo}\) & \texttt{\detokenize{>_{\rm lpo},}\allowbreak\detokenize{\geq_{\rm lpo}}} & strict lexicographic path order and its reflexive closure & \hyperref[sec:p12-reduction-orderings]{\ref*{sec:p12-reduction-orderings}} & relation & \hyperref[def:p12-ro-lpo]{\ref*{def:p12-ro-lpo}} & \thepage & The root comparison requires domination of every target argument. \\
\phantomsection\label{entry:symbol-registry-02433}\(\unlhd\) & \texttt{\detokenize{\unlhd}} & ordered homeomorphic embedding of finite ranked trees & \hyperref[sec:p12-reduction-orderings]{\ref*{sec:p12-reduction-orderings}} & relation & \hyperref[def:p12-ro-embedding]{\ref*{def:p12-ro-embedding}} & \thepage & Coupling keeps the same root and child coordinates; diving enters a child. \\
\phantomsection\label{entry:symbol-registry-02434}\(w(f),w(t),w_0\) & \texttt{\detokenize{w(f),}\allowbreak\detokenize{w(t),}\allowbreak\detokenize{w_0}} & symbol weight term weight and minimum nullary weight & \hyperref[sec:p12-reduction-orderings]{\ref*{sec:p12-reduction-orderings}} & meta & \hyperref[def:p12-ro-kbo-weights]{\ref*{def:p12-ro-kbo-weights}} & \thepage & All term weights are positive; zero weights remain allowed at the declared nonnullary symbols. \\
\phantomsection\label{entry:symbol-registry-02435}\(z^k(t)\) & \texttt{\detokenize{z^k(t)}} & iteration of the exceptional greatest zero-weight unary symbol & \hyperref[sec:p12-reduction-orderings]{\ref*{sec:p12-reduction-orderings}} & meta & \hyperref[def:p12-ro-kbo-weights]{\ref*{def:p12-ro-kbo-weights}} & \thepage & The zeroth iterate is the term itself; this is not exponentiation in an algebra. \\
\phantomsection\label{entry:symbol-registry-02436}\(>_{\rm kb70}\) & \texttt{\detokenize{>_{\rm kb70}}} & the exact historical 1970 Knuth-Bendix comparison & \hyperref[sec:p12-reduction-orderings]{\ref*{sec:p12-reduction-orderings}} & relation & \hyperref[def:p12-ro-kbo]{\ref*{def:p12-ro-kbo}} & \thepage & The equal-weight branch requires equality of every variable count. \\
\phantomsection\label{entry:symbol-registry-02437}\(n_j\) & \texttt{\detokenize{n_j}} & number of nodes of arity j in a ground tree & \hyperref[annex:p12.ro.kbo-termination]{\ref*{annex:p12.ro.kbo-termination}} & meta & \hyperref[annex:p12.ro.kbo-termination]{\ref*{annex:p12.ro.kbo-termination}} & \thepage & Used in the edge-count identity and the bounded-core argument. \\
\phantomsection\label{entry:symbol-registry-02438}\(d_x\) & \texttt{\detokenize{d_x}} & difference of variable occurrence counts between compared terms & \hyperref[annex:p12.ro.kbo-substitution]{\ref*{annex:p12.ro.kbo-substitution}} & meta & \hyperref[annex:p12.ro.kbo-substitution]{\ref*{annex:p12.ro.kbo-substitution}} & \thepage & Nonnegative in the strict-weight case and zero in the equal-weight case. \\
\phantomsection\label{entry:symbol-registry-02439}\(\mathsf A,\mathsf s,\mathsf0\) & \texttt{\detokenize{\mathsf A,}\allowbreak\detokenize{\mathsf s,}\allowbreak\detokenize{\mathsf0}} & Ackermann successor and zero function symbols in the orientation example & \hyperref[sec:p12-reduction-orderings]{\ref*{sec:p12-reduction-orderings}} & meta & \hyperref[thm:p12-ro-comparison-examples]{\ref*{thm:p12-ro-comparison-examples}} & \thepage & Arities are two one and zero respectively. \\
\phantomsection\label{entry:symbol-registry-02659}\(>\) & \texttt{\detokenize{>}} & fixed decidable reduction ordering for an entire completion run & \hyperref[sec:p12-fair-completion]{\ref*{sec:p12-fair-completion}} & relation & \hyperref[def:p12-fc-common-order]{\ref*{def:p12-fc-common-order}} & \thepage & Every generated rule decreases this one substitution-stable and context-stable ordering. \\
\phantomsection\label{entry:symbol-registry-02660}\(E_i\) & \texttt{\detokenize{E_i}} & finite set of pending equations at iteration i & \hyperref[sec:p12-fair-completion]{\ref*{sec:p12-fair-completion}} & object & \hyperref[def:p12-fc-state]{\ref*{def:p12-fc-state}} & \thepage & Initial equation set is E\_0; marking phases generate the next finite critical-equation set. \\
\phantomsection\label{entry:symbol-registry-02661}\(\mathcal R_i\) & \texttt{\detokenize{\mathcal R_i}} & finite labeled and marked rule-record set at iteration i & \hyperref[sec:p12-fair-completion]{\ref*{sec:p12-fair-completion}} & object & \hyperref[def:p12-fc-state]{\ref*{def:p12-fc-state}} & \thepage & Each record has unique label left side right side and mark. \\
\phantomsection\label{entry:symbol-registry-02662}\(R_i\) & \texttt{\detokenize{R_i}} & current rewriting system after forgetting labels and marks & \hyperref[sec:p12-fair-completion]{\ref*{sec:p12-fair-completion}} & object & \hyperref[def:p12-fc-state]{\ref*{def:p12-fc-state}} & \thepage & Right sides may change while surviving labels retain their left sides and marks. \\
\phantomsection\label{entry:symbol-registry-02663}\(c_i\) & \texttt{\detokenize{c_i}} & counter for fresh increasing rule labels & \hyperref[sec:p12-fair-completion]{\ref*{sec:p12-fair-completion}} & object & \hyperref[def:p12-fc-state]{\ref*{def:p12-fc-state}} & \thepage & The next new label is c\_i plus one; deleted labels are never reused. \\
\phantomsection\label{entry:symbol-registry-02664}\((k,\ell,r,m)\) & \texttt{\detokenize{(k,}\allowbreak\detokenize{\ell,}\allowbreak\detokenize{r,}\allowbreak\detokenize{m)}} & completion rule record & \hyperref[sec:p12-fair-completion]{\ref*{sec:p12-fair-completion}} & object & \hyperref[def:p12-fc-state]{\ref*{def:p12-fc-state}} & \thepage & Natural-number label k; term endpoints ell and r; marked or unmarked state m. \\
\phantomsection\label{entry:symbol-registry-02665}\(K\) & \texttt{\detokenize{K}} & labels whose current left sides collapse under the new rule & \hyperref[def:p12-fc-equation-phase]{\ref*{def:p12-fc-equation-phase}} & object & \hyperref[def:p12-fc-equation-phase]{\ref*{def:p12-fc-equation-phase}} & \thepage & This finite label set is local to one insertion macrostep. \\
\phantomsection\label{entry:symbol-registry-02666}\(S_i\) & \texttt{\detokenize{S_i}} & pre-insertion rewriting system together with the new rule & \hyperref[def:p12-fc-equation-phase]{\ref*{def:p12-fc-equation-phase}} & object & \hyperref[def:p12-fc-equation-phase]{\ref*{def:p12-fc-equation-phase}} & \thepage & Retained right sides are normalized under this precise historical rule set. \\
\phantomsection\label{entry:symbol-registry-02667}\(R_*\) & \texttt{\detokenize{R_*}} & union of all current rule versions throughout a run & \hyperref[sec:p12-fair-completion]{\ref*{sec:p12-fair-completion}} & object & \hyperref[def:p12-fc-fairness]{\ref*{def:p12-fc-fairness}} & \thepage & Includes versions of deleted rules and earlier right sides; it need not be computably finite. \\
\phantomsection\label{entry:symbol-registry-02668}\(R_\infty\) & \texttt{\detokenize{R_\infty}} & persistent final versions of never-deleted rule labels & \hyperref[sec:p12-fair-completion]{\ref*{sec:p12-fair-completion}} & object & \hyperref[def:p12-fc-fairness]{\ref*{def:p12-fc-fairness}} & \thepage & Well-founded right-side decrease proves final versions exist; the limit may be infinite. \\
\phantomsection\label{entry:symbol-registry-02669}\(\gg\) & \texttt{\detokenize{\gg}} & strict pattern containment relation & \hyperref[sec:p12-fair-completion]{\ref*{sec:p12-fair-completion}} & relation & \hyperref[def:p12-fc-containment]{\ref*{def:p12-fc-containment}} & \thepage & The left term contains an instance of the right pattern but the converse containment fails; different from the reduction ordering. \\
\phantomsection\label{entry:symbol-registry-02670}\(\to_*,\to_\infty\) & \texttt{\detokenize{\to_*,}\allowbreak\detokenize{\to_\infty}} & historical-union and persistent-limit one-step reduction abbreviations & \hyperref[annex:p12.fc.limit]{\ref*{annex:p12.fc.limit}} & relation & \hyperref[annex:p12.fc.limit]{\ref*{annex:p12.fc.limit}} & \thepage & Proof-local abbreviations; a superscript star forms the reflexive transitive closure. \\
\phantomsection\label{entry:symbol-registry-02671}\((\lvert a\rvert,\lvert a\rvert-\#\operatorname{Var}(a))\) & \texttt{\detokenize{(\lvert a\rvert,}\allowbreak\detokenize{\lvert a\rvert-\#\operatorname{Var}(a))}} & well-founded syntactic rank for strict pattern containment & \hyperref[annex:p12.fc.containment]{\ref*{annex:p12.fc.containment}} & object & \hyperref[annex:p12.fc.containment]{\ref*{annex:p12.fc.containment}} & \thepage & Bars count tree nodes and the variable symbol counts distinct occurring variables; both rank coordinates are nonnegative. \\
\phantomsection\label{entry:symbol-registry-03766}\(\operatorname{Sol}(E)\) & \texttt{\detokenize{\operatorname{Sol}}} & set of all finite-support syntactic solutions of a finite equation multiset & \hyperref[sec:p12-unification-extraction-adequacy]{\ref*{sec:p12-unification-extraction-adequacy}} & meta & \hyperref[def:p12-ue-substitutions]{\ref*{def:p12-ue-substitutions}} & \thepage & Global variable-map equality is used in factorization; irrelevant input variables are not silently quotiented. \\
\phantomsection\label{entry:symbol-registry-03767}\(\theta\eta\) & \texttt{\detokenize{\theta\eta}} & simultaneous substitution composition with theta acting first & \hyperref[sec:p12-unification-extraction-adequacy]{\ref*{sec:p12-unification-extraction-adequacy}} & meta & \hyperref[def:p12-ue-substitutions]{\ref*{def:p12-ue-substitutions}} & \thepage & Reuses the exact inherited P2 orientation; inserted images are not recursively reapplied. \\
\phantomsection\label{entry:symbol-registry-03768}\(\operatorname{iso}(E)\) & \texttt{\detokenize{\operatorname{iso}}} & variables with a single isolated left binding in the literal sheet state & \hyperref[sec:p12-unification-extraction-adequacy]{\ref*{sec:p12-unification-extraction-adequacy}} & meta & \hyperref[def:p12-ue-measure]{\ref*{def:p12-ue-measure}} & \thepage & Permanent isolated bindings define the first termination coordinate; no accumulator is added to the sheet. \\
\phantomsection\label{entry:symbol-registry-03769}\(\operatorname{uv}(E),\operatorname{sz}(E),\operatorname{bad}(E)\) & \texttt{\detokenize{\operatorname{uv};\operatorname{sz};\operatorname{bad}}} & three natural-number termination counters on the equation state & \hyperref[sec:p12-unification-extraction-adequacy]{\ref*{sec:p12-unification-extraction-adequacy}} & meta & \hyperref[def:p12-ue-measure]{\ref*{def:p12-ue-measure}} & \thepage & Unisolated variables; total term-tree nodes with multiplicity; nonvariable-left and variable-right equations. \\
\phantomsection\label{entry:symbol-registry-03770}\(\operatorname{nodes}(t)\) & \texttt{\detokenize{\operatorname{nodes}}} & number of nodes in a finite term tree & \hyperref[sec:p12-unification-extraction-adequacy]{\ref*{sec:p12-unification-extraction-adequacy}} & meta & \hyperref[def:p12-ue-measure]{\ref*{def:p12-ue-measure}} & \thepage & Counts variable and constant leaves and every function constructor; does not use cardinality bars for syntax size. \\
\phantomsection\label{entry:symbol-registry-03771}\(\#\) & \texttt{\detokenize{\#}} & finite count of variables or equation occurrences in a termination coordinate & \hyperref[sec:p12-unification-extraction-adequacy]{\ref*{sec:p12-unification-extraction-adequacy}} & meta & \hyperref[def:p12-ue-measure]{\ref*{def:p12-ue-measure}} & \thepage & Equation counts retain multiset multiplicity; variable counts count the finite set once. \\
\phantomsection\label{entry:symbol-registry-03772}\(\mu\) & \texttt{\detokenize{\mu}} & finite-support most general substitution and the subsequently extracted answer & \hyperref[sec:p12-unification-extraction-adequacy]{\ref*{sec:p12-unification-extraction-adequacy}} & meta & \hyperref[def:p12-ue-substitutions]{\ref*{def:p12-ue-substitutions}} & \thepage & The abstract type is declared first; terminal extraction reads all literal solved bindings simultaneously and is identity elsewhere. \\
\phantomsection\label{entry:symbol-registry-03773}\(D,V_0\) & \texttt{\detokenize{D;V_0}} & eliminated variable set and finite original variable set in the accumulator proof & \hyperref[annex:p12.ue.accumulator]{\ref*{annex:p12.ue.accumulator}} & meta & \hyperref[annex:p12.ue.accumulator]{\ref*{annex:p12.ue.accumulator}} & \thepage & Proof bookkeeping only; the actual inherited algorithm state remains E and theta. \\
\phantomsection\label{entry:symbol-registry-03774}\(\pi\) & \texttt{\detokenize{\pi}} & finite-support permutation comparing two most general solutions & P12 unification extraction & meta & \hyperref[thm:p12-ue-renaming]{\ref*{thm:p12-ue-renaming}} & \thepage & Named body-and-matching-annex scope recorded in p12-unification-extraction-source-note.md; distinct from an arbitrary factor substitution. \\
\phantomsection\label{entry:symbol-registry-04218}\(U\) & \texttt{\detokenize{U}} & set-diagram universe & P12 rationality-13-1 set and spider interpretation & n/a & \hyperref[def:p12-r131-region-reading]{\ref*{def:p12-r131-region-reading}} & \thepage & Fixed common semantic universe; geometry does not determine its cardinality. Nonempty domain is stated for comparison with ordinary monadic FOL. \\
\phantomsection\label{entry:symbol-registry-04219}\(P,Q,R,A,B\) & \texttt{\detokenize{P;Q;R;A;B}} & set extensions named by contour labels & P12 rationality-13-1 set examples & n/a & \hyperref[def:p12-r131-free-ride]{\ref*{def:p12-r131-free-ride}} & \thepage & P/Q/R Euler and cardinality examples; A/B Speedith example. The corresponding labels in L are syntax while Phi maps them to these extensions. \\
\phantomsection\label{entry:symbol-registry-04220}\(d=(L,Z,\mathrm{ShZ},S,\eta)\) & \texttt{\detokenize{d;L;Z;\mathrm{ShZ};S;\eta}} & Speedith unitary diagram and typed syntax fields & P12 rationality-13-1 spider syntax and rule fragment & n/a & \hyperref[def:p12-r131-spider-syntax]{\ref*{def:p12-r131-spider-syntax}} & \thepage & L finite contour-label set; Z present Boolean zones; ShZ shaded subset; S finite spider set; eta maps spiders to nonempty habitats. d is not the block-world individual in the later distinct example scope. \\
\phantomsection\label{entry:symbol-registry-04221}\(z=(\mathrm{in},\mathrm{out})\) & \texttt{\detokenize{z;\mathrm{in};\mathrm{out}}} & Boolean membership zone syntax & P12 rationality-13-1 spider semantics & n/a & \hyperref[def:p12-r131-spider-syntax]{\ref*{def:p12-r131-spider-syntax}} & \thepage & Inside/outside label sets partition L. Missing Boolean zones have empty interpretations; no diagram-area nonemptiness inference is licensed. \\
\phantomsection\label{entry:symbol-registry-04222}\(I=(U,\Phi)\) & \texttt{\detokenize{I;\Phi}} & Speedith interpretation tuple and contour extension map & P12 rationality-13-1 spider semantics & n/a & \hyperref[def:p12-r131-spider-semantics]{\ref*{def:p12-r131-spider-semantics}} & \thepage & Phi maps contour labels to subsets of U. I is an interpretation not the categorical I mood used in psychological syllogistics. \\
\phantomsection\label{entry:symbol-registry-04223}\(\zeta_I(z),\chi_I(h)\) & \texttt{\detokenize{\zeta;\chi}} & zone and habitat interpretation functions & P12 rationality-13-1 spider semantics & n/a & \hyperref[def:p12-r131-spider-semantics]{\ref*{def:p12-r131-spider-semantics}} & \thepage & zeta maps a zone to its membership intersection; chi maps a habitat region to the union of zone extensions. These retain the primary source notation and local types. \\
\phantomsection\label{entry:symbol-registry-04224}\(\Sigma:S\to U\) & \texttt{\detokenize{\Sigma}} & injective existential spider witness map & P12 rationality-13-1 spider semantics & n/a & \hyperref[def:p12-r131-spider-semantics]{\ref*{def:p12-r131-spider-semantics}} & \thepage & Different spiders have distinct witnesses; each unitary assertion existentially selects its map. Capital Sigma here is not a theory or summation sign. \\
\phantomsection\label{entry:symbol-registry-04225}\(I\models_{\rm sd}d\) & \texttt{\detokenize{\models_{\rm sd}}} & truth of a spider diagram in an interpretation & P12 rationality-13-1 spider semantics & n/a & \hyperref[def:p12-r131-spider-semantics]{\ref*{def:p12-r131-spider-semantics}} & \thepage & Indexed semantic truth under the declared Speedith conditions. Consequence quantifies over interpretations; no claim of psychological use follows. \\
\phantomsection\label{entry:symbol-registry-04226}\(s_i,t_i,u_i\) & \texttt{\detokenize{s_i;t_i;u_i}} & spider handles and explicitly bound existential witnesses & P12 rationality-13-1 cardinality and Speedith examples & n/a & \hyperref[def:p12-r131-spider-cardinality]{\ref*{def:p12-r131-spider-cardinality}} & \thepage & Handles in drawn syntax are not rigid individual constants. When inside an existential set formula they name that formula's bound witness elements. \\
\phantomsection\label{entry:symbol-registry-04227}\(d[s:h]\) & \texttt{\detokenize{d[s:h]}} & metalinguistic replacement of one spider habitat & P12 rationality-13-1 selected Speedith rule sheet & n/a & \hyperref[sheet:p12-r131-speedith-fragment]{\ref*{sheet:p12-r131-speedith-fragment}} & \thepage & Only eta(s) changes; all other tuple fields unchanged. The notation abbreviates actual AddFeet and SplitSpider transformations, not a new complete calculus. \\
\phantomsection\label{entry:symbol-registry-04228}\(c,d\) & \texttt{\detokenize{c;d}} & block-world individual constants & P12 rationality-13-1 Hyperproof example & n/a & \hyperref[thm:p12-r131-hyperproof]{\ref*{thm:p12-r131-hyperproof}} & \thepage & Local object-language constants in Dodec/Small/SameShape; distinct role from unitary diagram d in the spider-section scope. \\
\phantomsection\label{entry:symbol-registry-04229}\(\operatorname{Dodec},\operatorname{Small},\operatorname{SameShape}\) & \texttt{\detokenize{\operatorname{Dodec};\operatorname{Small};\operatorname{SameShape}}} & block-world shape and size predicates & P12 rationality-13-1 Hyperproof example & n/a & \hyperref[thm:p12-r131-hyperproof]{\ref*{thm:p12-r131-hyperproof}} & \thepage & Dodec and Small unary; SameShape binary under the task shape convention. A displayed unknown property is not silently an observation. \\
\phantomsection\label{entry:symbol-registry-04230}\(D_n,L_k\) & \texttt{\detokenize{D_n;L_k}} & finite dot square and disjoint ell layers & P12 rationality-13-1 Diamond example & n/a & \hyperref[thm:p12-r131-odd-sum]{\ref*{thm:p12-r131-odd-sum}} & \thepage & D\_n={1,...,n} squared and L\_k=D\_k minus D\_(k-1); not contour-label set L. Empty square D\_0 is the verified recursive base. \\
\phantomsection\label{entry:symbol-registry-04231}\(x,o\) & \texttt{\detokenize{x;o}} & Venn and Peirce existential or empty-region marks & P12 rationality-13-1 Venn discussion & n/a & \hyperref[sec:p12-rationality-13-1]{\ref*{sec:p12-rationality-13-1}} & \thepage & x-sequences give one witness in an alternative habitat; Peirce o-mark is historical emptiness notation. These are diagram marks rather than rigid constants or general quantified variable x in annex set arguments. \\
\phantomsection\label{entry:symbol-registry-04232}\(h,r\) & \texttt{\detokenize{h;r}} & sets of zones used as habitats or rule regions & P12 rationality-13-1 selected Speedith rule sheet & n/a & \hyperref[sheet:p12-r131-speedith-fragment]{\ref*{sheet:p12-r131-speedith-fragment}} & \thepage & h is an added-foot region; r is a nonempty proper habitat part for SplitSpider or a shaded subregion for RemoveShading. Each rule states its local subset side conditions. \\
\phantomsection\label{entry:symbol-registry-04233}\(i,j,n,k\) & \texttt{\detokenize{i;j;n;k}} & bound cell coordinates and natural-number size or layer indices & P12 rationality-13-1 checkerboard and ell proofs & n/a & \hyperref[thm:p12-r131-checkerboard]{\ref*{thm:p12-r131-checkerboard}} & \thepage & i/j are finite cell coordinates; n is a nonnegative square-size input and k a positive layer index. No measured participant count or population parameter is denoted. \\
\phantomsection\label{entry:symbol-registry-04397}\(\Gamma\vdash_T M:A\) & \texttt{\detokenize{\Gamma\vdash_T M:A}} & source type-theoretic judgment & p12/har2001-c18 & n/a & \hyperref[def:p12-har-checking]{\ref*{def:p12-har-checking}} & \thepage & Explicit context term and type; distinct from unqualified logical theoremhood. \\
\phantomsection\label{entry:symbol-registry-04398}\([A],[p]\) & \texttt{\detokenize{[A],}\allowbreak\detokenize{[p]}} & source translation into types and proof terms & p12/har2001-c18 & n/a & \hyperref[def:p12-har-checking]{\ref*{def:p12-har-checking}} & \thepage & Brackets denote a declared encoding, not its adequacy proof. \\
\phantomsection\label{entry:symbol-registry-04399}\(\mathrm{TCP},\mathrm{TSP},\mathrm{TIP}\) & \texttt{\detokenize{\mathrm{TCP},}\allowbreak\detokenize{\mathrm{TSP},}\allowbreak\detokenize{\mathrm{TIP}}} & checking synthesis and inhabitation task names & p12/har2001-c18 & n/a & \hyperref[def:p12-har-three-problems]{\ref*{def:p12-har-three-problems}} & \thepage & Metalinguistic problem names, not calculi. \\
\phantomsection\label{entry:symbol-registry-04400}\(\supset\) & \texttt{\detokenize{\supset}} & source logical implication & p12/har2001-c18 & n/a & \hyperref[def:p12-har-propositions-types]{\ref*{def:p12-har-propositions-types}} & \thepage & Direct source reading mirrors the arrow type. \\
\phantomsection\label{entry:symbol-registry-04401}\(\mathbin{\&}\) & \texttt{\detokenize{\mathbin{\&}}} & source conjunction & p12/har2001-c18 & n/a & \hyperref[def:p12-har-propositions-types]{\ref*{def:p12-har-propositions-types}} & \thepage & Published local syntax retained. \\
\phantomsection\label{entry:symbol-registry-04402}\(\cup\) & \texttt{\detokenize{\cup}} & source disjoint-union type & p12/har2001-c18 & n/a & \hyperref[def:p12-har-propositions-types]{\ref*{def:p12-har-propositions-types}} & \thepage & Tagged union of types, not ordinary untagged set union. \\
\phantomsection\label{entry:symbol-registry-04403}\(\varnothing\) & \texttt{\detokenize{\varnothing}} & source empty type & p12/har2001-c18 & n/a & \hyperref[def:p12-har-propositions-types]{\ref*{def:p12-har-propositions-types}} & \thepage & Type with no constructors in the constructive interpretation. \\
\phantomsection\label{entry:symbol-registry-04404}\(\Pi x:A.B\) & \texttt{\detokenize{\Pi x:A.B}} & source dependent-product type & p12/har2001-c18 & n/a & \hyperref[def:p12-har-propositions-types]{\ref*{def:p12-har-propositions-types}} & \thepage & B may depend on x; the sort combination must be allowed. \\
\phantomsection\label{entry:symbol-registry-04405}\(\Sigma x:A.B\) & \texttt{\detokenize{\Sigma x:A.B}} & source dependent-sum type & p12/har2001-c18 & n/a & \hyperref[def:p12-har-propositions-types]{\ref*{def:p12-har-propositions-types}} & \thepage & Dependent pairs; projection into data is regime-specific. \\
\phantomsection\label{entry:symbol-registry-04406}\(\mathrm{Set},\mathrm{Prop}\) & \texttt{\detokenize{\mathrm{Set},}\allowbreak\detokenize{\mathrm{Prop}}} & source domain and proposition sorts & p12/har2001-c18 & n/a & \hyperref[def:p12-har-relation-terms]{\ref*{def:p12-har-relation-terms}} & \thepage & Local source roles, not modern assistant universes. \\
\phantomsection\label{entry:symbol-registry-04407}\(\mathrm{Rel},\mathrm{AntiSym},\mathrm{Irrefl}\) & \texttt{\detokenize{\mathrm{Rel},}\allowbreak\detokenize{\mathrm{AntiSym},}\allowbreak\detokenize{\mathrm{Irrefl}}} & source defined relation operators & p12/har2001-c18 & n/a & \hyperref[def:p12-har-relation-terms]{\ref*{def:p12-har-relation-terms}} & \thepage & AntiSym names the source asymmetric-relation formula, not ordinary antisymmetry. \\
\phantomsection\label{entry:symbol-registry-04408}\(\mathrm{prop},\mathrm{type},T\) & \texttt{\detokenize{\mathrm{prop},}\allowbreak\detokenize{\mathrm{type},}\allowbreak\detokenize{T}} & framework proposition-name type host type sort and proof family & p12/har2001-c18 & n/a & \hyperref[def:p12-har-framework]{\ref*{def:p12-har-framework}} & \thepage & T maps proposition names to proof types; prop differs from direct Prop. \\
\phantomsection\label{entry:symbol-registry-04409}\(\Sigma_{\mathrm{PROP}},\Sigma_{\mathrm{PRED}},\Sigma_{\mathrm{lambda}}\) & \texttt{\detokenize{\Sigma_{\mathrm{PROP}},}\allowbreak\detokenize{\Sigma_{\mathrm{PRED}},}\allowbreak\detokenize{\Sigma_{\mathrm{lambda}}}} & source framework signatures & p12/har2001-c18 & n/a & \hyperref[def:p12-har-prop-signature]{\ref*{def:p12-har-prop-signature}} & \thepage & These Sigma glyphs denote signatures, not dependent sums. \\
\phantomsection\label{entry:symbol-registry-04410}\(\mathrm{imp\_intr},\mathrm{imp\_el}\) & \texttt{\detokenize{\mathrm{imp\_intr},}\allowbreak\detokenize{\mathrm{imp\_el}}} & framework implication proof constants & p12/har2001-c18 & n/a & \hyperref[def:p12-har-prop-signature]{\ref*{def:p12-har-prop-signature}} & \thepage & Declared object inference constants. \\
\phantomsection\label{entry:symbol-registry-04411}\(\mathrm{forall\_intr},\mathrm{forall\_elim}\) & \texttt{\detokenize{\mathrm{forall\_intr},}\allowbreak\detokenize{\mathrm{forall\_elim}}} & framework universal proof constants & p12/har2001-c18 & n/a & \hyperref[def:p12-har-pred-signature]{\ref*{def:p12-har-pred-signature}} & \thepage & Binding uses host abstraction. \\
\phantomsection\label{entry:symbol-registry-04412}\(\mathrm{app},\mathrm{abs},\mathrm{eq}\) & \texttt{\detokenize{\mathrm{app},}\allowbreak\detokenize{\mathrm{abs},}\allowbreak\detokenize{\mathrm{eq}}} & untyped-lambda framework constructors and equality family & p12/har2001-c18 & n/a & \hyperref[def:p12-har-lambda-signature]{\ref*{def:p12-har-lambda-signature}} & \thepage & Object application abstraction and beta equality are encoded constants, not host definitional equality. \\
\phantomsection\label{entry:symbol-registry-04413}\(\to_\beta,\to_\delta,\to_\iota\) & \texttt{\detokenize{\to_\beta,}\allowbreak\detokenize{\to_\delta,}\allowbreak\detokenize{\to_\iota}} & source beta unfolding and recursor reductions & p12/har2001-c18 & n/a & \hyperref[def:p12-har-relation-terms]{\ref*{def:p12-har-relation-terms}} & \thepage & No universal normalization assertion follows. \\
\phantomsection\label{entry:symbol-registry-04414}\(=_{\beta\iota\delta}\) & \texttt{\detokenize{=_{\beta\iota\delta}}} & source definitional convertibility & p12/har2001-c18 & n/a & \hyperref[def:p12-har-conversion]{\ref*{def:p12-har-conversion}} & \thepage & Distinct from book equality and extensional reflection. \\
\phantomsection\label{entry:symbol-registry-04415}\(\overline n\) & \texttt{\detokenize{\overline n}} & source numeral encoding & p12/har2001-c18 & n/a & \hyperref[def:p12-har-plus-conversion]{\ref*{def:p12-har-plus-conversion}} & \thepage & The term (S\^{}+)\^{}n0, not metalevel n. \\
\phantomsection\label{entry:symbol-registry-04416}\(\mathrm{Rec},S^+\) & \texttt{\detokenize{\mathrm{Rec},}\allowbreak\detokenize{S^+}} & source primitive recursor and successor & p12/har2001-c18 & n/a & \hyperref[def:p12-har-functions]{\ref*{def:p12-har-functions}} & \thepage & Recursor computation uses the specified iota clauses. \\
\phantomsection\label{entry:symbol-registry-04417}\(\pi_1,\pi_2\) & \texttt{\detokenize{\pi_1,}\allowbreak\detokenize{\pi_2}} & source dependent-sum projections & p12/har2001-c18 & n/a & \hyperref[thm:p12-har-sigma-extraction]{\ref*{thm:p12-har-sigma-extraction}} & \thepage & The second projection's type depends on the first. \\
\phantomsection\label{entry:symbol-registry-04418}\(=_A\) & \texttt{\detokenize{=_A}} & source book equality or declared setoid equality & p12/har2001-c18 & n/a & \hyperref[def:p12-har-equalities]{\ref*{def:p12-har-equalities}} & \thepage & Each choice is declared before use and differs from definitional conversion. \\
\phantomsection\label{entry:symbol-registry-04419}\(\mathrm{Eq}_A,\mathrm{Id}_A\) & \texttt{\detokenize{\mathrm{Eq}_A,}\allowbreak\detokenize{\mathrm{Id}_A}} & source Prop-valued equality and Set-valued identity & p12/har2001-c18 & n/a & \hyperref[def:p12-har-book-equality]{\ref*{def:p12-har-book-equality}} & \thepage & Different presentations and elimination permissions. \\
\phantomsection\label{entry:symbol-registry-04420}\([A,=_A]\) & \texttt{\detokenize{[A,}\allowbreak\detokenize{=_A]}} & source setoid & p12/har2001-c18 & n/a & \hyperref[def:p12-har-setoids]{\ref*{def:p12-har-setoids}} & \thepage & Carrier relation and proofs of equivalence; not literally the carrier type. \\
\phantomsection\label{entry:symbol-registry-04421}\(A\to^s B\) & \texttt{\detokenize{A\to^s B}} & source setoid function-space carrier & p12/har2001-c18 & n/a & \hyperref[def:p12-har-setoids]{\ref*{def:p12-har-setoids}} & \thepage & Maps paired with preservation proofs; equality ignores those proofs. \\
\phantomsection\label{entry:symbol-registry-04422}\(c:A\subseteq B\) & \texttt{\detokenize{c:A\subseteq B}} & source marked coercion map & p12/har2001-c18 & n/a & \hyperref[def:p12-har-coercions]{\ref*{def:p12-har-coercions}} & \thepage & An explicit term inserted by elaboration, not object set inclusion. \\
\phantomsection\label{entry:symbol-registry-04423}\(\iota P\) & \texttt{\detokenize{\iota P}} & source Church description operator & p12/har2001-c18 & n/a & \hyperref[def:p12-har-functions]{\ref*{def:p12-har-functions}} & \thepage & Predicate description term in its selected regime; distinct from the iota recursor-reduction subscript. \\
\phantomsection\label{entry:symbol-registry-04424}\(D,\Omega,\mathrm{Term}_\sigma\) & \texttt{\detokenize{D,}\allowbreak\detokenize{\Omega,}\allowbreak\detokenize{\mathrm{Term}_\sigma}} & source HOL domains proposition domain and domain-indexed term class & p12/har2001-c18 & n/a & \hyperref[def:p12-har-hol-language]{\ref*{def:p12-har-hol-language}} & \thepage & Omega is the logical domain of propositions, distinct from a PTS sort. \\
\phantomsection\label{entry:symbol-registry-04425}\(\Rightarrow\) & \texttt{\detokenize{\Rightarrow}} & source HOL object implication & p12/har2001-c18 & n/a & \hyperref[def:p12-har-hol-language]{\ref*{def:p12-har-hol-language}} & \thepage & Within the source HOL language this is a connective, not an unqualified metalevel implication. \\
\phantomsection\label{entry:symbol-registry-04426}\(x^\sigma,c_i^\sigma\) & \texttt{\detokenize{x^\sigma,}\allowbreak\detokenize{c_i^\sigma}} & source domain-annotated HOL variables and constants & p12/har2001-c18 & n/a & \hyperref[def:p12-har-hol-language]{\ref*{def:p12-har-hol-language}} & \thepage & Superscripts record object domains, not exponentiation. \\
\phantomsection\label{entry:symbol-registry-04427}\(\mathrm{Type},\mathrm{Type}'\) & \texttt{\detokenize{\mathrm{Type},}\allowbreak\detokenize{\mathrm{Type}'}} & source lambdaHOL data sort and untyped top sort & p12/har2001-c18 & n/a & \hyperref[def:p12-har-lambda-hol]{\ref*{def:p12-har-lambda-hol}} & \thepage & Prop:Type and Type:Type'; no type for Type' and no delta rule in this formal calculus. \\
\phantomsection\label{entry:symbol-registry-04428}\(\mathcal T,\mathcal S,\mathcal A,\mathcal R\) & \texttt{\detokenize{\mathcal T,}\allowbreak\detokenize{\mathcal S,}\allowbreak\detokenize{\mathcal A,}\allowbreak\detokenize{\mathcal R}} & source pseudo terms sorts axioms and product triples & p12/har2001-c18 & n/a & \hyperref[def:p12-har-lambda-hol]{\ref*{def:p12-har-lambda-hol}} & \thepage & The PTS specification preserves the source pair/triple convention. \\
\phantomsection\label{entry:symbol-registry-04429}\(\mathrm{Type}_\Gamma(M),\mathrm{Ok}(\Gamma)\) & \texttt{\detokenize{\mathrm{Type}_\Gamma(M),}\allowbreak\detokenize{\mathrm{Ok}(\Gamma)}} & source synthesis and context-validation functions & p12/har2001-c18 & n/a & \hyperref[def:p12-har-type-algorithm]{\ref*{def:p12-har-type-algorithm}} & \thepage & Type as a function name differs from the sort Type; false is a failure token and Ok true is a Boolean result, not the reflection proof constant true. \\
\phantomsection\label{entry:symbol-registry-04430}\(\vdash^{\mathrm{tc}}\) & \texttt{\detokenize{\vdash^{\mathrm{tc}}}} & source syntax-directed typing and context judgments & p12/har2001-c18 & n/a & \hyperref[def:p12-har-syntax-directed]{\ref*{def:p12-har-syntax-directed}} & \thepage & Modified source surface, distinct from the ordinary typing judgment. \\
\phantomsection\label{entry:symbol-registry-04431}\(\mathrm{Type}^s\) & \texttt{\detokenize{\mathrm{Type}^s}} & source PRED sort of first-order-domain sorts & p12/har2001-c18 & n/a & \hyperref[def:p12-har-pts-specifications]{\ref*{def:p12-har-pts-specifications}} & \thepage & Set:Type\^{}s in PRED and PREDomega; not lambdaHOL Type'. \\
\phantomsection\label{entry:symbol-registry-04432}\(\lambda\mathrm{PRED},\lambda\mathrm{PRED}\omega,\mathrm{CC},\lambda\ast\) & \texttt{\detokenize{\lambda\mathrm{PRED},}\allowbreak\detokenize{\lambda\mathrm{PRED}\omega,}\allowbreak\detokenize{\mathrm{CC},}\allowbreak\detokenize{\lambda\ast}} & source distinct PTS identifiers & p12/har2001-c18 & n/a & \hyperref[def:p12-har-pts-specifications]{\ref*{def:p12-har-pts-specifications}} & \thepage & Do not merge PRED's higher-order function domains with predicate quantification, or CC with CIC. \\
\phantomsection\label{entry:symbol-registry-04433}\(\leadsto,T^*\) & \texttt{\detokenize{\leadsto,}\allowbreak\detokenize{T^*}} & editorial parallel-beta proof relation and complete development & p12/har2001-c18 & n/a & \hyperref[thm:p12-har-pseudo-confluence]{\ref*{thm:p12-har-pseudo-confluence}} & \thepage & Annex proof devices for the source beta grammar; not a new admitted typing rule or object calculus. \\
\phantomsection\label{entry:symbol-registry-04434}\(f:\mathcal S\to\mathcal S'\) & \texttt{\detokenize{f:\mathcal S\to\mathcal S'}} & source metalevel PTS sort map & p12/har2001-c18/pts-morphism & n/a & \hyperref[def:p12-har-pts-morphism]{\ref*{def:p12-har-pts-morphism}} & \thepage & Homomorphic extension to pseudo terms and judgments; distinct local standing from earlier object functions named f. \\
\phantomsection\label{entry:symbol-registry-04435}\(g\) & \texttt{\detokenize{g}} & source proposed reverse judgment translation & p12/har2001-c18/pts-inverse & n/a & \hyperref[def:p12-har-pts-inverse]{\ref*{def:p12-har-pts-inverse}} & \thepage & Published sort clauses plus result-sort overrides; fails on the stated counterexample and is not treated as a valid inverse morphism. \\
\phantomsection\label{entry:symbol-registry-04436}\(V_D,V_L,V_P,D,O\) & \texttt{\detokenize{V_D,}\allowbreak\detokenize{V_L,}\allowbreak\detokenize{V_P,}\allowbreak\detokenize{D,}\allowbreak\detokenize{O}} & auxiliary variable classes and domain/object grammars & p12/har2001-c18/disambiguation & n/a & \hyperref[thm:p12-har-disambiguation]{\ref*{thm:p12-har-disambiguation}} & \thepage & Explicit proof-device syntax for the source stratification; D here denotes a domain expression class, not the earlier HOL domain inventory. \\
\phantomsection\label{entry:symbol-registry-04437}\(h,\sigma(D)\) & \texttt{\detokenize{h,}\allowbreak\detokenize{\sigma(D)}} & editorial raw transport and target-kind selector & p12/har2001-c18/cc-transport & n/a & \hyperref[thm:p12-har-hol-cc-transport]{\ref*{thm:p12-har-hol-cc-transport}} & \thepage & The result-sort override is judgment-dependent; h is not a PTS sort morphism. \\
\phantomsection\label{entry:symbol-registry-04438}\(A_1+A_2\) & \texttt{\detokenize{A_1+A_2}} & source sum-type former & p12/har2001-c18/extensions & n/a & \hyperref[def:p12-har-product-sum-extensions]{\ref*{def:p12-har-product-sum-extensions}} & \thepage & Distinct from arithmetic addition; the stated encoding is at Prop. \\
\phantomsection\label{entry:symbol-registry-04439}\(\circ\) & \texttt{\detokenize{\circ}} & source group binary operation & p12/har2001-c18/group-record & n/a & \hyperref[def:p12-har-group-record]{\ref*{def:p12-har-group-record}} & \thepage & Locally typed A-to-A-to-A, not substitution or function composition. \\
\phantomsection\label{entry:symbol-registry-04440}\(\mu,\sigma_j^i(\mu)\) & \texttt{\detokenize{\mu,}\allowbreak\detokenize{\sigma_j^i(\mu)}} & source fresh inductive type and strictly positive constructor argument scheme & p12/har2001-c18/inductives & n/a & \hyperref[def:p12-har-inductive-scheme]{\ref*{def:p12-har-inductive-scheme}} & \thepage & Domains of each argument scheme are mu-free; later dependent-mu example is separately declared. \\
\phantomsection\label{entry:symbol-registry-04441}\(\mathcal B_i(C),\mathcal B_i(P)\) & \texttt{\detokenize{\mathcal B_i(C),}\allowbreak\detokenize{\mathcal B_i(P)}} & metalevel generated branch types & p12/har2001-c18/inductive-scheme & n/a & \hyperref[def:p12-har-inductive-scheme]{\ref*{def:p12-har-inductive-scheme}} & \thepage & Explicit reconstruction device for original arguments and recursive value/proof families; not new object type formers. \\
\phantomsection\label{entry:symbol-registry-04442}\(\mathrm{nat},0,S\) & \texttt{\detokenize{\mathrm{nat},}\allowbreak\detokenize{0,}\allowbreak\detokenize{S}} & source fixed natural inductive type and constructors & p12/har2001-c18/inductives & n/a & \hyperref[def:p12-har-nat-list-tree]{\ref*{def:p12-har-nat-list-tree}} & \thepage & Distinct declaration from the earlier atomic HOL domain named natural numbers. \\
\phantomsection\label{entry:symbol-registry-04443}\(\mathrm{List},\mathrm{Nil},\mathrm{Cons}\) & \texttt{\detokenize{\mathrm{List},}\allowbreak\detokenize{\mathrm{Nil},}\allowbreak\detokenize{\mathrm{Cons}}} & source fixed-D list type and constructors & p12/har2001-c18/inductives & n/a & \hyperref[def:p12-har-nat-list-tree]{\ref*{def:p12-har-nat-list-tree}} & \thepage & The later List-D family is a separately proposed parametric extension. \\
\phantomsection\label{entry:symbol-registry-04444}\(\mathrm{Tree},\mathrm{Leaf},\mathrm{Join}\) & \texttt{\detokenize{\mathrm{Tree},}\allowbreak\detokenize{\mathrm{Leaf},}\allowbreak\detokenize{\mathrm{Join}}} & source countably branching well-founded tree type and constructors & p12/har2001-c18/inductives & n/a & \hyperref[def:p12-har-nat-list-tree]{\ref*{def:p12-har-nat-list-tree}} & \thepage & Join supplies a nat-indexed child family; infinite node count is not an infinite descending branch. \\
\phantomsection\label{entry:symbol-registry-04445}\(\mathrm{Rec}_\mu,\to_\iota\) & \texttt{\detokenize{\mathrm{Rec}_\mu,}\allowbreak\detokenize{\to_\iota}} & source instantiated inductive recursor and constructor computation & p12/har2001-c18/inductives & n/a & \hyperref[def:p12-har-inductive-scheme]{\ref*{def:p12-har-inductive-scheme}} & \thepage & Its target/predicate is chosen by the displayed rule, not an unlicensed polymorphic base constant type. \\
\phantomsection\label{entry:symbol-registry-04446}\(\mathrm{In},\mathrm{Eq},\mathrm{Refl}\) & \texttt{\detokenize{\mathrm{In},}\allowbreak\detokenize{\mathrm{Eq},}\allowbreak\detokenize{\mathrm{Refl}}} & source dependent constructor and inductive equality constructors & p12/har2001-c18/inductive-variants & n/a & \hyperref[def:p12-har-dependent-inductives]{\ref*{def:p12-har-dependent-inductives}} & \thepage & In belongs to the separately added Prop/Type regime; Eq eliminates at endpoint-dependent but proof-independent Prop conclusions. \\
\phantomsection\label{entry:symbol-registry-04447}\(K,C,O,\kappa(A)\) & \texttt{\detokenize{K,}\allowbreak\detokenize{C,}\allowbreak\detokenize{O,}\allowbreak\detokenize{\kappa(A)}} & CC proof classification and finite kind shape & p12/har2001-c18/cc-normalization & n/a & \hyperref[def:p12-har-cc-normalization-method]{\ref*{def:p12-har-cc-normalization-method}} & \thepage & Proof-device syntax for exact CC, distinct from HOL domain stratification. \\
\phantomsection\label{entry:symbol-registry-04448}\(\mathsf{SAT},\mathcal C,R\) & \texttt{\detokenize{\mathsf{SAT},}\allowbreak\detokenize{\mathcal C,}\allowbreak\detokenize{R}} & primary key-redex saturation and declared candidate adaptation & p12/har2001-c18/cc-normalization & n/a & \hyperref[def:p12-har-cc-normalization-method]{\ref*{def:p12-har-cc-normalization-method}} & \thepage & The two definitions are explicitly distinguished. \\
\phantomsection\label{entry:symbol-registry-04449}\(V,\xi,\rho\) & \texttt{\detokenize{V,}\allowbreak\detokenize{\xi,}\allowbreak\detokenize{\rho}} & kind-domain sets semantic constructor valuation and raw-term valuation & p12/har2001-c18/cc-normalization & n/a & \hyperref[def:p12-har-cc-normalization-method]{\ref*{def:p12-har-cc-normalization-method}} & \thepage & Semantic values and raw terms have different standing. \\
\phantomsection\label{entry:symbol-registry-04450}\(E,E_\rho,[\cdot]_\xi,d,\mathrm K\) & \texttt{\detokenize{E,}\allowbreak\detokenize{E_\rho,}\allowbreak\detokenize{[\cdot]_\xi,}\allowbreak\detokenize{d,}\allowbreak\detokenize{\mathrm K}} & annotation-retaining CC proof interpretations inert constant and K combinator & p12/har2001-c18/cc-normalization & n/a & \hyperref[def:p12-har-cc-normalization-method]{\ref*{def:p12-har-cc-normalization-method}} & \thepage & Inert d is not the object-language doubling function. \\
\phantomsection\label{entry:symbol-registry-04451}\(\eta(P)\) & \texttt{\detokenize{\eta(P)}} & finite maximum raw beta reduction length & p12/har2001-c18/cc-normalization & n/a & \hyperref[thm:p12-har-cc-normalization]{\ref*{thm:p12-har-cc-normalization}} & \thepage & Defined for SN raw terms, not eta conversion. \\
\phantomsection\label{entry:symbol-registry-04452}\(R\Rightarrow S\) & \texttt{\detokenize{R\Rightarrow S}} & CC proof candidate function space & p12/har2001-c18/cc-normalization & n/a & \hyperref[thm:p12-har-cc-normalization]{\ref*{thm:p12-har-cc-normalization}} & \thepage & A metalevel candidate construction, distinct from the source HOL object implication declared elsewhere. \\
\phantomsection\label{entry:symbol-registry-04453}\ensuremath{\texttt{(x:A)B},\texttt{[x:A]M}} & \texttt{\detokenize{(x:A)B,}\allowbreak\detokenize{[x:A]M}} & historical Coq product and abstraction input syntax & p12/har2001-c18/development & n/a & \hyperref[def:p12-har-tactic-workflow]{\ref*{def:p12-har-tactic-workflow}} & \thepage & Old parser conventions and the version6.3 example are not modern syntax. \\
\phantomsection\label{entry:symbol-registry-04454}\(\mathrm{DN},\mathrm{TT},\mathrm{FF}\) & \texttt{\detokenize{\mathrm{DN},}\allowbreak\detokenize{\mathrm{TT},}\allowbreak\detokenize{\mathrm{FF}}} & source development classical axiom truth and falsehood & p12/har2001-c18/development & n/a & \hyperref[def:p12-har-prime-development]{\ref*{def:p12-har-prime-development}} & \thepage & DN remains an explicit proof assumption. \\
\phantomsection\label{entry:symbol-registry-04455}\(\mathrm{div},\mathrm{propdiv},\mathrm{prime},\mathrm{primediv},\mathrm{HPD}\) & \texttt{\detokenize{\mathrm{div},}\allowbreak\detokenize{\mathrm{propdiv},}\allowbreak\detokenize{\mathrm{prime},}\allowbreak\detokenize{\mathrm{primediv},}\allowbreak\detokenize{\mathrm{HPD}}} & source historical number-theory definitions & p12/har2001-c18/development & n/a & \hyperref[def:p12-har-prime-development]{\ref*{def:p12-har-prime-development}} & \thepage & Divisor is the first argument; definitions do not prove the end theorem. \\
\phantomsection\label{entry:symbol-registry-04456}\(G,E_\nu,\mathrm{KPrime},F\) & \texttt{\detokenize{G,}\allowbreak\detokenize{E_\nu,}\allowbreak\detokenize{\mathrm{KPrime},}\allowbreak\detokenize{F}} & source internal transformations interpreters and characteristic/numerical terms & p12/har2001-c18/autarkic & n/a & \hyperref[def:p12-har-autarkic]{\ref*{def:p12-har-autarkic}} & \thepage & Object-language algorithms, distinct from the normalization proof's metalevel E. \\
\phantomsection\label{entry:symbol-registry-04457}\(\Gamma\) & \texttt{\detokenize{\Gamma}} & finite source equational-unification problem & P12 Baader-Snyder equational unification & meta & \hyperref[def:p12-eu-equational-problem]{\ref*{def:p12-eu-equational-problem}} & \thepage & Body and matching annex only; not a sequent context. \\
\phantomsection\label{entry:symbol-registry-04458}\(\mathcal F,\mathcal V,\mathcal X\) & \texttt{\detokenize{\mathcal F;\mathcal V;\mathcal X}} & ranked signature variable alphabet and finite comparison domain & P12 Baader-Snyder equational unification & meta & \hyperref[def:p12-eu-equational-problem]{\ref*{def:p12-eu-equational-problem}} & \thepage & X is fixed as Var(Gamma) before solution comparison. \\
\phantomsection\label{entry:symbol-registry-04459}\(=_{E}\) & \texttt{\detokenize{=_{E}}} & least substitution-closed congruence generated by E & P12 Baader-Snyder equational unification & relation & \hyperref[def:p12-eu-equational-problem]{\ref*{def:p12-eu-equational-problem}} & \thepage & Distinct from literal equality and the symmetric instance relation. \\
\phantomsection\label{entry:symbol-registry-04460}\(\approx\) & \texttt{\detokenize{\approx}} & identity used to generate an equational theory & P12 Baader-Snyder equational unification & relation & \hyperref[def:p12-eu-equational-problem]{\ref*{def:p12-eu-equational-problem}} & \thepage & Source identity notation; not elementary equivalence. \\
\phantomsection\label{entry:symbol-registry-04461}\(=^{?}_{E}\) & \texttt{\detokenize{=^{?}_{E}}} & input equation to be solved modulo E & P12 Baader-Snyder equational unification & relation & \hyperref[def:p12-eu-equational-problem]{\ref*{def:p12-eu-equational-problem}} & \thepage & Question-mark annotation denotes a unification problem. \\
\phantomsection\label{entry:symbol-registry-04462}\(\mathcal U_E(\Gamma)\) & \texttt{\detokenize{\mathcal U_E}} & all finite-support E-unifiers of the input & P12 Baader-Snyder equational unification & meta & \hyperref[def:p12-eu-equational-problem]{\ref*{def:p12-eu-equational-problem}} & \thepage & Source calligraphic U; finite trees; fresh variables permitted; no rational trees. \\
\phantomsection\label{entry:symbol-registry-04463}\(\mathcal C\) & \texttt{\detokenize{\mathcal C}} & complete set of E-unifiers & \hyperref[def:p12-eu-generality]{\ref*{def:p12-eu-generality}} & meta & \hyperref[def:p12-eu-generality]{\ref*{def:p12-eu-generality}} & \thepage & Source calligraphic complete-set variable; distinct from theory C and later set of constants. \\
\phantomsection\label{entry:symbol-registry-04464}\(\mathrel{\underline{\lessdot}}_E^{\mathcal X}\) & \texttt{\detokenize{\mathrel{\underline{\lessdot}}_E^{\mathcal X}}} & dotted source instance comparison modulo E on X & P12 Baader-Snyder equational unification & relation & \hyperref[def:p12-eu-generality]{\ref*{def:p12-eu-generality}} & \thepage & Source p470 detail visually confirms dot inside less-than angle and lower horizontal line; typography still needs manuscript visual checking. \\
\phantomsection\label{entry:symbol-registry-04465}\(\doteq_E^{\mathcal X}\) & \texttt{\detokenize{\doteq_E^{\mathcal X}}} & mutual instance comparison modulo E on X & P12 Baader-Snyder equational unification & relation & \hyperref[def:p12-eu-generality]{\ref*{def:p12-eu-generality}} & \thepage & Source dotted equality; not an input equation or equivalence under a variable permutation. \\
\phantomsection\label{entry:symbol-registry-04466}\(1,\omega,\infty,0\) & \texttt{\detokenize{1;\omega;\infty;0}} & ordered unification types unitary finitary infinitary zero & P12 Baader-Snyder equational unification & meta & \hyperref[def:p12-eu-types]{\ref*{def:p12-eu-types}} & \thepage & Zero means no minimal complete set not absence of solutions. \\
\phantomsection\label{entry:symbol-registry-04467}\(C_{\mathcal F}(E)\) & \texttt{\detokenize{C_{\mathcal F}(E)}} & category of finitely generated E-free algebras & P12 Baader-Snyder equational unification & meta & \hyperref[def:p12-eu-free-algebra]{\ref*{def:p12-eu-free-algebra}} & \thepage & Arrows are homomorphisms; MGU factor is not required unique. \\
\phantomsection\label{entry:symbol-registry-04468}\(<\) & \texttt{\detokenize{<}} & linear ordering inducing constant restrictions & P12 Baader-Snyder equational unification & relation & \hyperref[def:p12-eu-lcr]{\ref*{def:p12-eu-lcr}} & \thepage & x<c forbids c in the image of x; different from numeric or reduction orders. \\
\phantomsection\label{entry:symbol-registry-04469}\(|t|_q\) & \texttt{\detokenize{|t|_q}} & source generator-occurrence count in an ACU term & P12 Baader-Snyder equational unification & meta & \hyperref[def:p12-eu-three-theories]{\ref*{def:p12-eu-three-theories}} & \thepage & Noncardinality source use explicitly retained locally. \\
\phantomsection\label{entry:symbol-registry-04470}\(\operatorname{coef}_q(t)\) & \texttt{\detokenize{\operatorname{coef}}} & typed coordinate in the declared free algebra & P12 Baader-Snyder equational unification & meta & \hyperref[def:p12-eu-three-theories]{\ref*{def:p12-eu-three-theories}} & \thepage & Editorial coordinate notation; domains N Boolean and Z explicitly distinguished. \\
\phantomsection\label{entry:symbol-registry-04471}\(BS\) & \texttt{\detokenize{BS}} & Boolean semiring with OR addition and AND multiplication & P12 Baader-Snyder equational unification & meta & \hyperref[def:p12-eu-three-theories]{\ref*{def:p12-eu-three-theories}} & \thepage & Not a Boolean ring. \\
\phantomsection\label{entry:symbol-registry-04472}\(U(M_1,M_2)\) & \texttt{\detokenize{U(M_1,}\allowbreak\detokenize{M_2)}} & right semimodule of homogeneous solution vectors & P12 Baader-Snyder equational unification & meta & \hyperref[def:p12-eu-homogeneous-construction]{\ref*{def:p12-eu-homogeneous-construction}} & \thepage & Vectors used as columns in multiplication; source describes ambient n-tuples as row vectors on p511. \\
\phantomsection\label{entry:symbol-registry-04473}\(P_c\) & \texttt{\detokenize{P_c}} & finite particular-solution set for one fixed free constant & P12 Baader-Snyder equational unification & meta & \hyperref[def:p12-eu-constant-construction]{\ref*{def:p12-eu-constant-construction}} & \thepage & Each particular solution plus homogeneous solutions covers its inhomogeneous set. \\
\phantomsection\label{entry:symbol-registry-04474}\(S_E\) & \texttt{\detokenize{S_E}} & one-generator free-algebra endomorphism semiring & P12 Baader-Snyder equational unification & meta & \hyperref[def:p12-eu-semiring]{\ref*{def:p12-eu-semiring}} & \thepage & Multiplication is source-order composition and need not be commutative. \\
\phantomsection\label{entry:symbol-registry-04475}\(M_\sigma\) & \texttt{\detokenize{M_\sigma}} & matrix representing a free-algebra arrow & P12 Baader-Snyder equational unification & meta & \hyperref[def:p12-eu-semiring]{\ref*{def:p12-eu-semiring}} & \thepage & Source generators index rows; target generators index columns. \\
\phantomsection\label{entry:symbol-registry-04476}\(*\) & \texttt{\detokenize{*}} & implicit commutative monoid operation for endomorphism addition & P12 Baader-Snyder equational unification & meta & \hyperref[def:p12-eu-semiring]{\ref*{def:p12-eu-semiring}} & \thepage & Not necessarily an operation symbol in the original signature. \\
\phantomsection\label{entry:symbol-registry-04477}\(\mathcal U,\mathcal G,\mathcal B,\mathcal N\) & \texttt{\detokenize{\mathcal U;\mathcal G;\mathcal B;\mathcal N}} & source syntactic arbitrary-equational basic and standard narrowing calculi & P12 Baader-Snyder equational unification & meta & \hyperref[def:p12-eu-gbn-states]{\ref*{def:p12-eu-gbn-states}} & \thepage & Distinct state shapes and allowed positions; U is not the solution-set symbol. \\
\phantomsection\label{entry:symbol-registry-04478}\(P;C;S\) & \texttt{\detokenize{P;C;S}} & skeleton multiset syntactic constraints and solved bindings & P12 Baader-Snyder equational unification & meta & \hyperref[def:p12-eu-gbn-states]{\ref*{def:p12-eu-gbn-states}} & \thepage & B elimination leaves P uninstantiated; N instantiates P. \\
\phantomsection\label{entry:symbol-registry-04479}\(\sigma_S,[\sigma]\) & \texttt{\detokenize{\sigma_S;[\sigma]}} & solved-binding simultaneous map and an idempotent map's solved equations & P12 Baader-Snyder equational unification & meta & \hyperref[def:p12-eu-gbn-states]{\ref*{def:p12-eu-gbn-states}} & \thepage & Not a retained-equation-sheet accumulator; inserted images act once. \\
\phantomsection\label{entry:symbol-registry-04480}\(e[u]\) & \texttt{\detokenize{e[u]}} & equation with distinguished subterm occurrence & P12 Baader-Snyder equational unification & meta & \hyperref[def:p12-eu-gbn-states]{\ref*{def:p12-eu-gbn-states}} & \thepage & Source lowercase e context role; not ACU's unit constant. \\
\phantomsection\label{entry:symbol-registry-04481}\(F_j,F_\omega\) & \texttt{\detokenize{F_j;F_\omega}} & explicit unrestricted editorial critical closure & P12 Baader-Snyder equational unification & meta & \hyperref[def:p12-eu-abstract-completion]{\ref*{def:p12-eu-abstract-completion}} & \thepage & Not identified with the source ordered E\_omega construction. \\
\phantomsection\label{entry:symbol-registry-04482}\(\operatorname{Gr}^{\succ}(F_\omega)\) & \texttt{\detokenize{\operatorname{Gr}^{\succ}}} & decreasing ground identity instances in the declared parameter extension & P12 Baader-Snyder equational unification & meta & \hyperref[def:p12-eu-abstract-completion]{\ref*{def:p12-eu-abstract-completion}} & \thepage & May be infinite; conceptual normalization not an effective algorithm claim. \\
\phantomsection\label{entry:symbol-registry-04483}\(\succ\) & \texttt{\detokenize{\succ}} & well-founded monotone reduction ordering total on ground terms & P12 Baader-Snyder equational unification & relation & \hyperref[def:p12-eu-abstract-completion]{\ref*{def:p12-eu-abstract-completion}} & \thepage & Ground-domain extension declared before lifting. \\
\phantomsection\label{entry:symbol-registry-04484}\(\operatorname{Cr}^{\succ}(E),\operatorname{cp}(E)\) & \texttt{\detokenize{\operatorname{Cr}^{\succ};\operatorname{cp}}} & cross-referred source ordered-critical-pair operator & P12 Baader-Snyder equational unification & meta & \hyperref[def:p12-eu-ordered-cr]{\ref*{def:p12-eu-ordered-cr}} & \thepage & Some-substitution two-filter admission; not all overlaps followed by orientation. \\
\phantomsection\label{entry:symbol-registry-04485}\(\preceq\) & \texttt{\detokenize{\preceq}} & nonstrict comparison in the ordered-overlap filters & P12 Baader-Snyder equational unification & relation & \hyperref[def:p12-eu-ordered-cr]{\ref*{def:p12-eu-ordered-cr}} & \thepage & a preceq b iff a=b or b succ a; source p569 writes <=. \\
\phantomsection\label{entry:symbol-registry-04486}\(E_j,E_\omega\) & \texttt{\detokenize{E_j;E_\omega}} & source ordered critical closure & P12 Baader-Snyder equational unification & meta & \hyperref[def:p12-eu-ordered-cr]{\ref*{def:p12-eu-ordered-cr}} & \thepage & Distinct from the editorial unrestricted F\_j and F\_omega. \\
\phantomsection\label{entry:symbol-registry-04487}\(M,i,n_1,n_2,n_3\) & \texttt{\detokenize{M;i;n_1;n_2;n_3}} & completeness induction counters for eager refinement & P12 Baader-Snyder equational unification & meta & \hyperref[def:p12-eu-eager-strategy]{\ref*{def:p12-eu-eager-strategy}} & \thepage & Ground skeleton multiset; nonvariable skeleton symbols; C variable/node/orientation counters. B and N lexicographic orders stated separately. \\
\phantomsection\label{entry:symbol-registry-04488}\(\gamma\) & \texttt{\detokenize{\gamma}} & fresh-only literal matching map in Reduce & P12 Baader-Snyder equational unification & meta & \hyperref[def:p12-eu-eager-strategy]{\ref*{def:p12-eu-eager-strategy}} & \thepage & l gamma=u sigma\_S; isolated output bindings differ between literal and cleaned variants. \\
\phantomsection\label{entry:symbol-registry-04489}\(\prec_{\rm red}\) & \texttt{\detokenize{\prec_{\rm red}}} & source redex-position ordering extending proper-subterm priority & P12 Baader-Snyder equational unification & relation & \hyperref[def:p12-eu-redex-ordered-strategy]{\ref*{def:p12-eu-redex-ordered-strategy}} & \thepage & Finite retained-context abstraction pass and lowest-index rule test declared. \\
\phantomsection\label{entry:symbol-registry-04490}\(\twoheadrightarrow_\alpha\) & \texttt{\detokenize{\twoheadrightarrow_\alpha}} & reflexive transitive reduction closure & p12-src-027-501 & n/a & \hyperref[def:p12-klop1987-ars]{\ref*{def:p12-klop1987-ars}} & \thepage & Klop1987 double-arrow source convention; not a second one-step relation \\
\phantomsection\label{entry:symbol-registry-04491}\(=_\alpha\) & \texttt{\detokenize{=_\alpha}} & conversion relation & p12-src-027-501 & n/a & \hyperref[def:p12-klop1987-ars]{\ref*{def:p12-klop1987-ars}} & \thepage & Equivalence generated by the declared reduction; ordinary equality remains identity elsewhere \\
\phantomsection\label{entry:symbol-registry-04492}\(|a|\) & \texttt{\detokenize{|a|}} & natural-number increasing rank & p12-src-027-501 & n/a & \hyperref[def:p12-klop1987-properties]{\ref*{def:p12-klop1987-properties}} & \thepage & Source-specific bars are the Inc rank; not cardinality \\
\phantomsection\label{entry:symbol-registry-04493}\(\mathcal G(a)\) & \texttt{\detokenize{\mathcal G(a)}} & closed reduction graph & p12-src-027-501 & n/a & \hyperref[def:p12-klop1987-subars]{\ref*{def:p12-klop1987-subars}} & \thepage & All descendants of a with restricted reduction \\
\phantomsection\label{entry:symbol-registry-04494}\(t^\sigma\) & \texttt{\detokenize{t^\sigma}} & homomorphic substitution action & p12-src-027-501 & n/a & \hyperref[def:p12-klop1987-terms]{\ref*{def:p12-klop1987-terms}} & \thepage & Report source notation; finite terms \\
\phantomsection\label{entry:symbol-registry-04495}\(R_1\oplus R_2\) & \texttt{\detokenize{\oplus}} & disjoint-signature sum of TRSs & p12-src-027-501 & n/a & \hyperref[def:p12-klop1987-direct-sum]{\ref*{def:p12-klop1987-direct-sum}} & \thepage & Nonvariable alphabets renamed apart; not arbitrary union \\
\phantomsection\label{entry:symbol-registry-04496}\(\mathcal T^*\) & \texttt{\detokenize{\mathcal T^*}} & finite trees permitting marker-labelled nodes & p12-src-027-501 & n/a & \hyperref[def:p12-klop1987-star-order]{\ref*{def:p12-klop1987-star-order}} & \thepage & Star on the tree family is marking permission; not reduction closure \\
\phantomsection\label{entry:symbol-registry-04497}\(t^*\) & \texttt{\detokenize{t^*}} & root-marker operation & p12-src-027-501 & n/a & \hyperref[def:p12-klop1987-star-order]{\ref*{def:p12-klop1987-star-order}} & \thepage & Local to the marked-tree order; not a normal-form operator \\
\phantomsection\label{entry:symbol-registry-04498}\(\to_{\parallel}\) & \texttt{\detokenize{\to_{\parallel}}} & finite disjoint-redex parallel contraction & p12-src-027-501 & n/a & \hyperref[def:p12-klop1987-regular]{\ref*{def:p12-klop1987-regular}} & \thepage & Empty set expressly allowed; distinct from nested full substitution \\
\phantomsection\label{entry:symbol-registry-04499}\(\mathbb F\) & \texttt{\detokenize{\mathbb F}} & computable reduction strategy & p12-src-027-501 & n/a & \hyperref[def:p12-klop1987-strategy]{\ref*{def:p12-klop1987-strategy}} & \thepage & Fixes normal forms; chooses one or positive finitely many steps otherwise \\
\phantomsection\label{entry:symbol-registry-04500}\(I_t\) & \texttt{\detokenize{I_t}} & parameterized combinatory identity term & p12-src-027-501 & n/a & \hyperref[thm:p12-klop1987-omega]{\ref*{thm:p12-klop1987-omega}} & \thepage & Source p35 defines I\_t=SKt; acts as identity when given an argument \\
\phantomsection\label{entry:symbol-registry-04501}\(\Omega_t\) & \texttt{\detokenize{\Omega_t}} & source outermost cofinality counterexample term & p12-src-027-501 & n/a & \hyperref[thm:p12-klop1987-omega]{\ref*{thm:p12-klop1987-omega}} & \thepage & Source p35 defines SI\_tI\_t(SI\_tI\_t) \\
\phantomsection\label{entry:symbol-registry-04502}\(\Lambda\) & \texttt{\detokenize{\Lambda}} & unary categorical abstraction operator & p12-src-027-501 & n/a & \hyperref[def:p12-klop1987-weak-categorical]{\ref*{def:p12-klop1987-weak-categorical}} & \thepage & Source Table8 operator; not quantifier or lambda binder \\
\phantomsection\label{entry:symbol-registry-04503}\(\mathsf{App}\) & \texttt{\detokenize{\mathsf{App}}} & categorical application constant & p12-src-027-501 & n/a & \hyperref[def:p12-klop1987-weak-categorical]{\ref*{def:p12-klop1987-weak-categorical}} & \thepage & Distinct from binary underlying application operator Ap \\
\phantomsection\label{entry:symbol-registry-04504}\(\rho\) & \texttt{\detokenize{\rho}} & finite or infinite reduction sequence & p12-src-027-501 & n/a & \hyperref[def:p12-klop1987-fair-traces]{\ref*{def:p12-klop1987-fair-traces}} & \thepage & Editorial residual-projection proof notation; not a term substitution \\
\phantomsection\label{entry:symbol-registry-04505}\(U_i\) & \texttt{\detokenize{U_i}} & finite disjoint set of a fixed redex's residuals at stage i & p12-src-027-501 & n/a & \hyperref[def:p12-klop1987-outermost-fair]{\ref*{def:p12-klop1987-outermost-fair}} & \thepage & Original-origin labels are not newly created redexes \\
\phantomsection\label{entry:symbol-registry-04506}\(\mathsf a,\mathsf w,\mathsf i,\mathsf k\) & \texttt{\detokenize{\mathsf a;\mathsf w;\mathsf i;\mathsf k}} & auxiliary constant codes in Omega invariant proof & p12-src-027-501 & n/a & \hyperref[thm:p12-klop1987-omega]{\ref*{thm:p12-klop1987-omega}} & \thepage & Editorial only; interpreted as SKI and S(SKI)(SKI) and I and K; original CL is unchanged \\
\phantomsection\label{entry:symbol-registry-04507}\(V\) & \texttt{\detokenize{V}} & auxiliary finite value grammar for CL invariant & p12-src-027-501 & n/a & \hyperref[thm:p12-klop1987-omega]{\ref*{thm:p12-klop1987-omega}} & \thepage & Editorial grammar class in the Omega proof; not the original variable set \\
\phantomsection\label{entry:symbol-registry-04508}\(\to_3\) & \texttt{\detokenize{\to_3}} & finite1-then-finite2 composite relation & p12-src-027-501-ars & n/a & \hyperref[def:p12-klop1987-ars-exercises]{\ref*{def:p12-klop1987-ars-exercises}} & \thepage & Source exercise7 composite; both component arrows are reflexive transitive closures \\
\phantomsection\label{entry:symbol-registry-04509}\(I\) & \texttt{\detokenize{I}} & nullary identity combinator & p12-src-027-501-applicative & n/a & \hyperref[def:p12-klop1987-applicative]{\ref*{def:p12-klop1987-applicative}} & \thepage & CL Table2-4 constant; not the unary group inverse \\
\phantomsection\label{entry:symbol-registry-04510}\(I\) & \texttt{\detokenize{I}} & unary group inverse operator & p12-src-027-501-group & n/a & \hyperref[def:p12-klop1987-group-signature]{\ref*{def:p12-klop1987-group-signature}} & \thepage & Group Tables10 and12; fixed arity1 distinct from CL constant \\
\phantomsection\label{entry:symbol-registry-04511}\(e\) & \texttt{\detokenize{e}} & nullary group identity symbol & p12-src-027-501-group & n/a & \hyperref[def:p12-klop1987-group-signature]{\ref*{def:p12-klop1987-group-signature}} & \thepage & Exact source group signature \\
\phantomsection\label{entry:symbol-registry-04512}\(.\) & \texttt{\detokenize{\mathbin{.}}} & binary noncommutative group product & p12-src-027-501-group & n/a & \hyperref[def:p12-klop1987-group-signature]{\ref*{def:p12-klop1987-group-signature}} & \thepage & Parentheses remain explicit; associativity is oriented not assumed during matching \\
\phantomsection\label{entry:symbol-registry-04633}\(R_E\) & \texttt{\detokenize{R_E}} & source subset of reduced oriented ground identity instances & P12 equational hereditary fragments & meta & \hyperref[def:p12-eu-h-ground-instances]{\ref*{def:p12-eu-h-ground-instances}} & \thepage & Generating maps restricted to identity variables; reduced with respect to full Gr\^{}succ(E). \\
\phantomsection\label{entry:symbol-registry-04634}\(m,m_{R_E},n_1,n_2,n_3\) & \texttt{\detokenize{m;m_{R_E};n_1;n_2;n_3}} & printed minimum-basic-proof tuple and explicit reduced-proof minimum & P12 equational hereditary fragments & meta & \hyperref[def:p12-eu-h-source-measure]{\ref*{def:p12-eu-h-source-measure}} & \thepage & Full LP/matching blocks decrease reduced minimum; raw LP need not decrease printed tuple. \\
\phantomsection\label{entry:symbol-registry-04635}\(q\theta\rho\) & \texttt{\detokenize{q\theta\rho}} & updated source skeleton with first-then-second finite substitution composition & P12 equational hereditary fragments & meta & \hyperref[def:p12-eu-h-basic-proof]{\ref*{def:p12-eu-h-basic-proof}} & \thepage & Fresh rule pattern remains uninstantiated in skeleton. \\
\phantomsection\label{entry:symbol-registry-04636}\(\mathcal T\) & \texttt{\detokenize{\mathcal T}} & finite nonempty unordered natural-labelled trees & p12-klop1987-marked-trees & n/a & \hyperref[def:p12-klop1987-mt-objects]{\ref*{def:p12-klop1987-mt-objects}} & \thepage & Unranked source objects; numerical labels may have any finite number of children \\
\phantomsection\label{entry:symbol-registry-04637}\(\simeq\) & \texttt{\detokenize{\simeq}} & isomorphism by finite sibling permutations & p12-klop1987-marked-trees & n/a & \hyperref[def:p12-klop1987-mt-objects]{\ref*{def:p12-klop1987-mt-objects}} & \thepage & Equivalence on planar representations; tree objects are its classes \\
\phantomsection\label{entry:symbol-registry-04638}\(\sqsubseteq\) & \texttt{\detokenize{\sqsubseteq}} & source ancestry-order node embedding & p12-klop1987-marked-trees & n/a & \hyperref[def:p12-klop1987-mt-embedding]{\ref*{def:p12-klop1987-mt-embedding}} & \thepage & Preserves and reflects ancestry and increases labels; not silently meet preserving \\
\phantomsection\label{entry:symbol-registry-04639}\(\preccurlyeq_{\rm pl}\) & \texttt{\detokenize{\preccurlyeq_{\rm pl}}} & stronger planar diving/coupling tree embedding & p12-klop1987-marked-trees & n/a & \hyperref[def:p12-klop1987-mt-embedding]{\ref*{def:p12-klop1987-mt-embedding}} & \thepage & Auxiliary proof relation implies source embedding after forgetting planar order \\
\phantomsection\label{entry:symbol-registry-04640}\(\rightsquigarrow\) & \texttt{\detokenize{\rightsquigarrow}} & all-context auxiliary marker reduction & p12-klop1987-marked-trees & n/a & \hyperref[def:p12-klop1987-mt-objects]{\ref*{def:p12-klop1987-mt-objects}} & \thepage & Source marker relation is a subrelation; auxiliary enlargement is explicit \\
\phantomsection\label{entry:symbol-registry-04641}\(|t|\) & \texttt{\detokenize{|t|}} & node count for norm-indexed source example & p12-klop1987-marked-trees & n/a & \hyperref[thm:p12-klop1987-mt-indexed]{\ref*{thm:p12-klop1987-mt-indexed}} & \thepage & Counts operators and variables once; not the earlier Inc rank or cardinality \\
\phantomsection\label{entry:symbol-registry-04653}\(\mathrm{type},\ast,\Box\) & \texttt{\detokenize{\mathrm{type},}\allowbreak\detokenize{\ast,}\allowbreak\detokenize{\Box}} & weak LF kind and two-sort PTS view & p12/har-framework-faithfulness & n/a & \hyperref[def:p12-har-fw-host]{\ref*{def:p12-har-fw-host}} & \thepage & Only object-domain product pairs; distinct from direct lambdaHOL. \\
\phantomsection\label{entry:symbol-registry-04654}\(\mathcal D\) & \texttt{\detokenize{\mathcal D}} & normal-name source reification & p12/har-framework-faithfulness & n/a & \hyperref[def:p12-har-fw-contexts]{\ref*{def:p12-har-fw-contexts}} & \thepage & Pointwise clauses do not assert object eta. \\
\phantomsection\label{entry:symbol-registry-04655}\(\mathrm{nf}_\beta\) & \texttt{\detokenize{\mathrm{nf}_\beta}} & typed LF beta normalization & p12/har-framework-faithfulness & n/a & \hyperref[def:p12-har-fw-contexts]{\ref*{def:p12-har-fw-contexts}} & \thepage & No assumption of beta conversion of every term to eta-long form. \\
\phantomsection\label{entry:symbol-registry-04656}\(r,j,d\) & \texttt{\detokenize{r,}\allowbreak\detokenize{j,}\allowbreak\detokenize{d}} & rational reduction inclusion and positive gcd & p12/har-framework-rationals & n/a & \hyperref[thm:p12-har-rational-normalization]{\ref*{thm:p12-har-rational-normalization}} & \thepage & Unique numerical representative need not have a unique proof witness. \\
\phantomsection\label{entry:symbol-registry-04783}\(\Sigma_1,\Sigma_2\) & \texttt{\detokenize{\Sigma_1;\Sigma_2}} & disjoint ranked nonvariable signatures & p12-klop1987-direct-sum & n/a & \hyperref[def:p12-klop1987-ds-signatures]{\ref*{def:p12-klop1987-ds-signatures}} & \thepage & Includes constants; sharing application is not allowed \\
\phantomsection\label{entry:symbol-registry-04784}\(C[\mathbf A]\) & \texttt{\detokenize{C[\mathbf A]}} & maximal one-colour cap with principal-alien occurrences & p12-klop1987-direct-sum & n/a & \hyperref[def:p12-klop1987-ds-layers]{\ref*{def:p12-klop1987-ds-layers}} & \thepage & Ordinary variables remain neutral leaves \\
\phantomsection\label{entry:symbol-registry-04785}\(H[\mathbf A]\) & \texttt{\detokenize{H[\mathbf A]}} & proof frontier cap with neutral variables also holed & p12-klop1987-direct-sum & n/a & \hyperref[def:p12-klop1987-ds-layers]{\ref*{def:p12-klop1987-ds-layers}} & \thepage & Separate occurrence presentation; not silently the principal-alien tuple \\
\phantomsection\label{entry:symbol-registry-04786}\(\operatorname{rank}(t)\) & \texttt{\detokenize{\operatorname{rank}(t)}} & alternating-signature layer rank & p12-klop1987-direct-sum & n/a & \hyperref[def:p12-klop1987-ds-layers]{\ref*{def:p12-klop1987-ds-layers}} & \thepage & Minimum1 including a variable; not node count \\
\phantomsection\label{entry:symbol-registry-04787}\(\mathbf A\propto\mathbf B\) & \texttt{\detokenize{\propto}} & preservation of literal equality among frontier indices & p12-klop1987-direct-sum & n/a & \hyperref[def:p12-klop1987-ds-coherence]{\ref*{def:p12-klop1987-ds-coherence}} & \thepage & Original equal values must update equally; distinct values may merge \\
\phantomsection\label{entry:symbol-registry-04788}\(d(t)\) & \texttt{\detokenize{d(t)}} & maximum depth of minimal nonpreserved special layers & p12-klop1987-direct-sum & n/a & \hyperref[def:p12-klop1987-ds-depth]{\ref*{def:p12-klop1987-ds-depth}} & \thepage & Semantic depth for a confluence proof; not an ordinary termination measure \\
\phantomsection\label{entry:symbol-registry-04789}\(\Rightarrow_{\rm d}\) & \texttt{\detokenize{\Rightarrow_{\rm d}}} & parallel deepest-layer deletion by finite ordinary reductions & p12-klop1987-direct-sum & n/a & \hyperref[def:p12-klop1987-ds-depth]{\ref*{def:p12-klop1987-ds-depth}} & \thepage & Not a single ordinary rewrite and not assumed computable \\
\phantomsection\label{entry:symbol-registry-04790}\(D,d_c\) & \texttt{\detokenize{D;d_c}} & signed variable-count matrix and fixed-constant offset column & P12 pure AC fragments & meta & \hyperref[def:p12-ac-input]{\ref*{def:p12-ac-input}} & \thepage & D=M\_L-M\_R; d\_c=u\_R,c-u\_L,c; constants are fixed. \\
\phantomsection\label{entry:symbol-registry-04791}\(H(A)\) & \texttt{\detokenize{H(A)}} & computed minimal nonzero homogeneous integer generators & P12 pure AC fragments & meta & \hyperref[def:p12-ac-system-basis]{\ref*{def:p12-ac-system-basis}} & \thepage & Finite rational-ray and bounded integer remainder algorithm with full proof. \\
\phantomsection\label{entry:symbol-registry-04792}\(r_j,B_i\) & \texttt{\detokenize{r_j;B_i}} & primitive positive rational cone rays and finite coordinate bounds & P12 pure AC fragments & meta & \hyperref[def:p12-ac-system-basis]{\ref*{def:p12-ac-system-basis}} & \thepage & Rays alone do not necessarily span integer solutions. \\
\phantomsection\label{entry:symbol-registry-04793}\(\widetilde D_c,P_c\) & \texttt{\detokenize{\widetilde D_c;P_c}} & augmented homogeneous matrix and exact minimal affine particular set & P12 pure AC fragments & meta & \hyperref[def:p12-ac-particulars]{\ref*{def:p12-ac-particulars}} & \thepage & Last coordinate1 generators; last coordinate0 gives homogeneous residuals. \\
\phantomsection\label{entry:symbol-registry-04794}\(\sigma_{p,I}\) & \texttt{\detokenize{\sigma_{p,I}}} & unit-free positive-row answer selected by particular columns and basis subset & P12 pure AC fragments & meta & \hyperref[def:p12-ac-extraction]{\ref*{def:p12-ac-extraction}} & \thepage & No parameter receives empty AC term; I can be empty only if fixed constants cover all rows. \\
\phantomsection\label{entry:symbol-registry-04795}\(K,L\) & \texttt{\detokenize{K;L}} & source single-equation balance bound and maximum pairwise lcm & P12 pure AC fragments & meta & \hyperref[def:p12-ac-single-bound]{\ref*{def:p12-ac-single-bound}} & \thepage & Positive coefficients and both sides present; zero/cancelled columns treated separately. \\
\phantomsection\label{entry:symbol-registry-04813}\(\ast\) & \texttt{\detokenize{\ast}} & self-typed sort & p12/har-lambda-star & n/a & \hyperref[def:p12-har-star-specification]{\ref*{def:p12-har-star-specification}} & \thepage & Exact one-sort axiom/triple. \\
\phantomsection\label{entry:symbol-registry-04814}\(\ast,\Box,\triangle\) & \texttt{\detokenize{\ast,}\allowbreak\detokenize{\Box,}\allowbreak\detokenize{\triangle}} & source hierarchy sorts & p12/har-hurkens & n/a & \hyperref[def:p12-har-hurkens-source]{\ref*{def:p12-har-hurkens-source}} & \thepage & Actual four pairs; U extra pair separate. \\
\phantomsection\label{entry:symbol-registry-04815}\(\wp X\) & \texttt{\detokenize{\wp X}} & source predicate-space type & p12/har-hurkens & n/a & \hyperref[def:p12-har-hurkens-source]{\ref*{def:p12-har-hurkens-source}} & \thepage & X-to-star, not primitive powerset. \\
\phantomsection\label{entry:symbol-registry-04816}\(\mathcal U,\tau t,\sigma s\) & \texttt{\detokenize{\mathcal U,}\allowbreak\detokenize{\tau t,}\allowbreak\detokenize{\sigma s}} & source universe and explicit recipes & p12/har-hurkens & n/a & \hyperref[def:p12-har-hurkens-data]{\ref*{def:p12-har-hurkens-data}} & \thepage & No equality is postulated. \\
\phantomsection\label{entry:symbol-registry-04817}\(I,K,W,q_p\) & \texttt{\detokenize{I,}\allowbreak\detokenize{K,}\allowbreak\detokenize{W,}\allowbreak\detokenize{q_p}} & proof-index abbreviations & p12/har-hurkens & n/a & \hyperref[def:p12-har-hurkens-proof]{\ref*{def:p12-har-hurkens-proof}} & \thepage & No additional axioms. \\
\phantomsection\label{entry:symbol-registry-04818}\(\Delta,\Omega\) & \texttt{\detokenize{\Delta,}\allowbreak\detokenize{\Omega}} & diagonal predicate and computed object & p12/har-hurkens & n/a & \hyperref[def:p12-har-hurkens-proof]{\ref*{def:p12-har-hurkens-proof}} & \thepage & Roles distinct from upper sort. \\
\phantomsection\label{entry:symbol-registry-04819}\(0,1,2,3\) & \texttt{\detokenize{0,}\allowbreak\detokenize{1,}\allowbreak\detokenize{2,}\allowbreak\detokenize{3}} & numeric proof-variable labels & p12/har-hurkens & n/a & \hyperref[def:p12-har-hurkens-proof]{\ref*{def:p12-har-hurkens-proof}} & \thepage & Not arithmetic numerals. \\
\phantomsection\label{entry:symbol-registry-04820}\(R_0,M_0,L_0,D_0,D_\ast\) & \texttt{\detokenize{R_0,}\allowbreak\detokenize{M_0,}\allowbreak\detokenize{L_0,}\allowbreak\detokenize{D_0,}\allowbreak\detokenize{D_\ast}} & factored proof and transported witness & p12/har-lambda-star & n/a & \hyperref[def:p12-har-hurkens-proof]{\ref*{def:p12-har-hurkens-proof}} & \thepage & All application types and conversions supplied. \\
\phantomsection\label{entry:symbol-registry-04863}\(\Pi\) & \texttt{\detokenize{\Pi}} & partition of shared input variables & p12-equational-combination & object & \hyperref[def:p12-comb-branch]{\ref*{def:p12-comb-branch}} & \thepage & Source signature-identification branch; not a probability measure \\
\phantomsection\label{entry:symbol-registry-04864}\(L\) & \texttt{\detokenize{L}} & signature owner labelling of partition representatives & p12-equational-combination & object & \hyperref[def:p12-comb-branch]{\ref*{def:p12-comb-branch}} & \thepage & Nonshared variables have their unique block owner \\
\phantomsection\label{entry:symbol-registry-04865}\(\sigma_\Pi\) & \texttt{\detokenize{\sigma_\Pi}} & identification substitution & p12-equational-combination & meta & \hyperref[def:p12-comb-branch]{\ref*{def:p12-comb-branch}} & \thepage & Maps every shared-variable block to its chosen representative \\
\phantomsection\label{entry:symbol-registry-04866}\(F_i(S)\) & \texttt{\detokenize{F_i(S)}} & pure Ei free algebra on named generators S & p12-equational-combination & object & \hyperref[def:p12-comb-tower]{\ref*{def:p12-comb-tower}} & \thepage & Generators are algebraic atoms not assumed signature constants \\
\phantomsection\label{entry:symbol-registry-04867}\(N_i^k\) & \texttt{\detokenize{N_i^k}} & stage k nongenerator classes of colour i & p12-equational-combination & object & \hyperref[def:p12-comb-tower]{\ref*{def:p12-comb-tower}} & \thepage & Editorial equational tower; old nodes identified by proved inclusions \\
\phantomsection\label{entry:symbol-registry-04868}\(nu\) & \texttt{\detokenize{\nu}} & evaluation into the combined free algebra & p12-equational-combination & object & \hyperref[def:p12-comb-tower]{\ref*{def:p12-comb-tower}} & \thepage & Not the narrowing accumulator or a ground normal-form algorithm \\
\phantomsection\label{entry:symbol-registry-04869}\(e_i\) & \texttt{\detokenize{e_i}} & bijection from combined carrier to a pure reduct's free algebra & p12-equational-combination & object & \hyperref[def:p12-comb-tower]{\ref*{def:p12-comb-tower}} & \thepage & Used to transport Sigma\_i operations \\
\phantomsection\label{entry:symbol-registry-04870}\(Merge\) & \texttt{\detokenize{\operatorname{Merge}_{L,<}}} & finite ordered assembly of local LCR answers & p12-equational-combination & object & \hyperref[def:p12-comb-merge]{\ref*{def:p12-comb-merge}} & \thepage & Not arbitrary cyclic substitution composition \\
\phantomsection\label{entry:symbol-registry-04871}\(VarConst\) & \texttt{\detokenize{\operatorname{VarConst}}} & set of variable and designated free-constant names in a term & p12-equational-combination & object & \hyperref[def:p12-comb-domain]{\ref*{def:p12-comb-domain}} & \thepage & Occurrence restriction not cardinality \\
\phantomsection\label{entry:symbol-registry-04872}\(0\) & \texttt{\detokenize{0}} & group identity e in the explicitly additive AG counterexample & \hyperref[def:p12-comb-ag-filter-failure]{\ref*{def:p12-comb-ag-filter-failure}} & object & \hyperref[def:p12-comb-ag-filter-failure]{\ref*{def:p12-comb-ag-filter-failure}} & \thepage & Local group element not a coordinate scalar or a new free constant \\
\phantomsection\label{entry:symbol-registry-04873}\(X\) & \texttt{\detokenize{X}} & variables shared by at least two purified component blocks & p12-equational-combination & object & \hyperref[def:p12-comb-branch]{\ref*{def:p12-comb-branch}} & \thepage & Source role distinct from original-variable comparison scope in the pure AC subsection \\
\phantomsection\label{entry:symbol-registry-04874}\(X_Pi\) & \texttt{\detokenize{X_\Pi}} & representatives of shared-variable partition classes & p12-equational-combination & object & \hyperref[def:p12-comb-branch]{\ref*{def:p12-comb-branch}} & \thepage & Ordered and signature-labelled in each combination branch \\
\phantomsection\label{entry:symbol-registry-04875}\(S_i^k\) & \texttt{\detokenize{S_i^k}} & stage k pure free-algebra generator set & p12-equational-combination & object & \hyperref[def:p12-comb-tower]{\ref*{def:p12-comb-tower}} & \thepage & Contains Z and earlier foreign-colour nodes \\
\phantomsection\label{entry:symbol-registry-04876}\(D\) & \texttt{\detokenize{D}} & combined free-product carrier & p12-equational-combination & object & \hyperref[def:p12-comb-tower]{\ref*{def:p12-comb-tower}} & \thepage & Local carrier not the AC coefficient matrix D or a distributivity theory \\
\phantomsection\label{entry:symbol-registry-04877}\(Z\) & \texttt{\detokenize{Z}} & countably infinite generator set for the completeness model & p12-equational-combination & object & \hyperref[def:p12-comb-tower]{\ref*{def:p12-comb-tower}} & \thepage & Not integers; the integer domain is written mathbb Z \\
\phantomsection\label{entry:symbol-registry-04878}\(P_i\) & \texttt{\detokenize{P_i}} & disjoint fresh range-parameter set of local answer i & p12-equational-combination & object & \hyperref[def:p12-comb-merge]{\ref*{def:p12-comb-merge}} & \thepage & Not a partition block P\_j in the source's identification definition \\
\phantomsection\label{entry:symbol-registry-04879}\(B_m\) & \texttt{\detokenize{B_m}} & number of partitions of an m-element set & p12-equational-combination & object & \hyperref[thm:p12-comb-decision]{\ref*{thm:p12-comb-decision}} & \thepage & Finite branch bound defined in reciprocal proof \\
\phantomsection\label{entry:symbol-registry-04880}\(\Sigma_M\) & \texttt{\detokenize{\Sigma_M}} & finite reachable ranked numerical signature & p12-klop1987-ordered-markers & n/a & \hyperref[thm:p12-klop1987-om-bound]{\ref*{thm:p12-klop1987-om-bound}} & \thepage & Labels at most initial M; not a common global bound on all inputs \\
\phantomsection\label{entry:symbol-registry-04881}\(\bot\) & \texttt{\detokenize{\bot}} & fresh semantic nullary bottom symbol & p12-klop1987-ordered-markers & n/a & \hyperref[def:p12-klop1987-om-interpretation]{\ref*{def:p12-klop1987-om-interpretation}} & \thepage & Below every numerical label; not an actual marked-game node or rule \\
\phantomsection\label{entry:symbol-registry-04882}\(U_M\) & \texttt{\detokenize{U_M}} & ground semantic term universe & p12-klop1987-ordered-markers & n/a & \hyperref[def:p12-klop1987-om-interpretation]{\ref*{def:p12-klop1987-om-interpretation}} & \thepage & Finite ranked alphabet Sigma\_M plus fresh bottom; unbounded term depth and size \\
\phantomsection\label{entry:symbol-registry-04883}\(D(v)\) & \texttt{\detokenize{D}} & weak principal LPO lower set & p12-klop1987-ordered-markers & n/a & \hyperref[def:p12-klop1987-om-interpretation]{\ref*{def:p12-klop1987-om-interpretation}} & \thepage & Contains v and all strictly smaller semantic terms \\
\phantomsection\label{entry:symbol-registry-04884}\(D^-(v)\) & \texttt{\detokenize{D^-}} & strict principal LPO lower set & p12-klop1987-ordered-markers & n/a & \hyperref[def:p12-klop1987-om-interpretation]{\ref*{def:p12-klop1987-om-interpretation}} & \thepage & Contains semantic bottom for every actual source-headed application \\
\phantomsection\label{entry:symbol-registry-04885}\(J(t)\) & \texttt{\detokenize{J}} & nonempty lower-set interpretation of a marked tree & p12-klop1987-ordered-markers & n/a & \hyperref[def:p12-klop1987-om-interpretation]{\ref*{def:p12-klop1987-om-interpretation}} & \thepage & Union over argument choices; marked roots use strict lower sets \\
\phantomsection\label{entry:symbol-registry-04886}\(\to_{\rm lex}\) & \texttt{\detokenize{\to_{\rm lex}}} & source fixed-arity ordered marker relation & p12-klop1987-ordered-markers & n/a & \hyperref[def:p12-klop1987-om-game]{\ref*{def:p12-klop1987-om-game}} & \thepage & Five operations with plain-root context closure and well-typed root marking \\
\phantomsection\label{entry:symbol-registry-04887}\(\rightsquigarrow_{\rm lex}\) & \texttt{\detokenize{\rightsquigarrow_{\rm lex}}} & all-context ordered auxiliary marker relation & p12-klop1987-ordered-markers & n/a & \hyperref[def:p12-klop1987-om-game]{\ref*{def:p12-klop1987-om-game}} & \thepage & Explicit enlargement including marked contexts; only plain endpoint positive restriction proved well-founded \\
\phantomsection\label{entry:symbol-registry-04888}\(O,S,\mathsf{one},\mathsf{two}\) & \texttt{\detokenize{O,}\allowbreak\detokenize{S,}\allowbreak\detokenize{\mathsf{one},}\allowbreak\detokenize{\mathsf{two}}} & historical nat constructors and numerals & p12/har-prime-development & n/a & \hyperref[def:p12-har-prime-arithmetic]{\ref*{def:p12-har-prime-arithmetic}} & \thepage & Coq6.3 nat:Set; O is constructor letter not a new proof variable. \\
\phantomsection\label{entry:symbol-registry-04889}\(\mathsf{nat\_ind},\mathsf{nat\_rec}\) & \texttt{\detokenize{\mathsf{nat\_ind},}\allowbreak\detokenize{\mathsf{nat\_rec}}} & historical propositional induction and instantiated data recursion & p12/har-prime-development & n/a & \hyperref[def:p12-har-prime-regime]{\ref*{def:p12-har-prime-regime}} & \thepage & Not identified with unchanged beta-only lambdaHOL or a top-sort-polymorphic constant. \\
\phantomsection\label{entry:symbol-registry-04890}\(\mathsf{eq},\mathsf{refl\_equal}\) & \texttt{\detokenize{\mathsf{eq},}\allowbreak\detokenize{\mathsf{refl\_equal}}} & historical Set-indexed primitive propositional equality & p12/har-prime-development & n/a & \hyperref[def:p12-har-prime-regime]{\ref*{def:p12-har-prime-regime}} & \thepage & Only generated Prop transport used; no equality reflection. \\
\phantomsection\label{entry:symbol-registry-04891}\(\mathsf{DN}\) & \texttt{\detokenize{\mathsf{DN}}} & explicit double-negation axiom & p12/har-prime-development & n/a & \hyperref[def:p12-har-prime-regime]{\ref*{def:p12-har-prime-regime}} & \thepage & Used for tnd on existence of proper divisor; no classical data extraction. \\
\phantomsection\label{entry:symbol-registry-04892}\(\mathsf{Is\_suc},\mathsf{pred}\) & \texttt{\detokenize{\mathsf{Is\_suc},}\allowbreak\detokenize{\mathsf{pred}}} & predicate-code case and data predecessor & p12/har-prime-development & n/a & \hyperref[def:p12-har-prime-arithmetic]{\ref*{def:p12-har-prime-arithmetic}} & \thepage & Both scrutinize nat data; no Prop evidence eliminated into Set. \\
\phantomsection\label{entry:symbol-registry-04893}\(\mathsf{leseq},\mathsf{lthan},\le,<,\mathsf{before}\) & \texttt{\detokenize{\mathsf{leseq},}\allowbreak\detokenize{\mathsf{lthan},}\allowbreak\detokenize{\le,}\allowbreak\detokenize{<,}\allowbreak\detokenize{\mathsf{before}}} & historical inductive order and induction predicate & p12/har-prime-development & relation & \hyperref[def:p12-har-prime-arithmetic]{\ref*{def:p12-har-prime-arithmetic}} & \thepage & n<m means S n<=m; cv\_ind is proved from nat\_ind. \\
\phantomsection\label{entry:symbol-registry-04894}\(\mathsf{plus},\mathsf{times},+,nm\) & \texttt{\detokenize{\mathsf{plus},}\allowbreak\detokenize{\mathsf{times},}\allowbreak\detokenize{+,}\allowbreak\detokenize{nm}} & source first-input arithmetic recursion and abbreviations & p12/har-prime-development & object & \hyperref[def:p12-har-prime-arithmetic]{\ref*{def:p12-har-prime-arithmetic}} & \thepage & Semiring and order identities proved; definitions not treated as their theorems. \\
\phantomsection\label{entry:symbol-registry-04895}\(\mathsf{div},\mathsf{propdiv},\mathsf{prime},\mathsf{primediv},\mathsf{HPD}\) & \texttt{\detokenize{\mathsf{div},}\allowbreak\detokenize{\mathsf{propdiv},}\allowbreak\detokenize{\mathsf{prime},}\allowbreak\detokenize{\mathsf{primediv},}\allowbreak\detokenize{\mathsf{HPD}}} & source divisor and primality predicates & p12/har-prime-development & n/a & \hyperref[def:p12-har-prime-arithmetic]{\ref*{def:p12-har-prime-arithmetic}} & \thepage & div d n is Prop existential coefficient x with times x d=n. \\
\phantomsection\label{entry:symbol-registry-04896}\(a\mathbin{\dot{-}}b\) & \texttt{\detokenize{\mathbin{\dot{-}}}} & explicit cutoff-subtraction data function & p12/har-prime-development & object & \hyperref[def:p12-har-prime-arithmetic]{\ref*{def:p12-har-prime-arithmetic}} & \thepage & Constructed local divisibility coefficient; zero multiplier separately handled. \\
\phantomsection\label{entry:symbol-registry-04897}\(\mathsf{fac}\) & \texttt{\detokenize{\mathsf{fac}}} & factorial data recursion for source Euclid continuation & p12/har-prime-development & n/a & \hyperref[def:p12-har-prime-arithmetic]{\ref*{def:p12-har-prime-arithmetic}} & \thepage & Fac O=one; factorial bound and divisor property separately proved. \\
\phantomsection\label{entry:symbol-registry-04898}\(B k\ a\) & \texttt{\detokenize{B k\ a}} & finite-family sum-bound data fold & p12/har-prime-development & n/a & \hyperref[thm:p12-har-prime-euclid]{\ref*{thm:p12-har-prime-euclid}} & \thepage & Author-supplied typed finite-family representation; no first-order sequence coding inferred. \\
\phantomsection\label{entry:symbol-registry-04940}\(\operatorname{Inst}\) & \texttt{\detokenize{\operatorname{Inst}}} & compilation of a reachable basic state to its instantiated ordinary state & p12-equational-narrowing-transport & meta & \hyperref[def:p12-en-compile]{\ref*{def:p12-en-compile}} & \thepage & Exact P sigma\_S; identical C and S; not normalization \\
\phantomsection\label{entry:symbol-registry-04941}\(A\) & \texttt{\detokenize{A}} & solved substitution sigma\_S of the current state & p12-equational-narrowing-transport & object & \hyperref[def:p12-en-variants]{\ref*{def:p12-en-variants}} & \thepage & Source-order composition; distinct from associative theory A \\
\phantomsection\label{entry:symbol-registry-04942}\(D_S\) & \texttt{\detokenize{D_S}} & domain of the current solved substitution & p12-equational-narrowing-transport & object & \hyperref[def:p12-en-variants]{\ref*{def:p12-en-variants}} & \thepage & Absent from literal constraints and solved images \\
\phantomsection\label{entry:symbol-registry-04943}\(\mathcal L\) & \texttt{\detokenize{\mathcal L}} & fixed family of fresh LP identities used in both representations & p12-equational-narrowing-transport & object & \hyperref[def:p12-en-variants]{\ref*{def:p12-en-variants}} & \thepage & Forward R and inverse-identity families are separate choices \\
\phantomsection\label{entry:symbol-registry-04944}\(\widehat A\) & \texttt{\detokenize{\widehat A}} & terminal image-normalized answer substitution & p12-equational-narrowing-transport & object & \hyperref[thm:p12-en-terminal-normalization]{\ref*{thm:p12-en-terminal-normalization}} & \thepage & Not a replacement for the active literal S accumulator \\
\phantomsection\label{entry:symbol-registry-04945}\(0\) & \texttt{\detokenize{0}} & free constant in the convergent infinite-narrowing example & \hyperref[thm:p12-en-infinite-search]{\ref*{thm:p12-en-infinite-search}} & object & \hyperref[thm:p12-en-infinite-search]{\ref*{thm:p12-en-infinite-search}} & \thepage & Term atom distinct from a; its integer interpretation is one \\
\phantomsection\label{entry:symbol-registry-04946}\(Q\) & \texttt{\detokenize{Q}} & finite literal equation problem selected for the syntactic helper & p12-equational-narrowing-transport & object & \hyperref[def:p12-en-source-composites]{\ref*{def:p12-en-source-composites}} & \thepage & MGU computed on e sigma\_S or l sigma\_S=u sigma\_S; not equational unification \\
\phantomsection\label{entry:symbol-registry-04947}\(\lhd\) & \texttt{\detokenize{\lhd}} & binary PA left merge & p12-klop1987-pa & n/a & \hyperref[def:p12-klop1987-pa-regime]{\ref*{def:p12-klop1987-pa-regime}} & \thepage & Typographic surrogate for source left-merge glyph; not an order relation \\
\phantomsection\label{entry:symbol-registry-04948}\(R_3\) & \texttt{\detokenize{R_3}} & exact three ordinary PA source rules & p12-klop1987-pa & n/a & \hyperref[def:p12-klop1987-pa-regime]{\ref*{def:p12-klop1987-pa-regime}} & \thepage & Atomic a is a constant; all substitutions and contexts allowed \\
\phantomsection\label{entry:symbol-registry-04949}\(R_7\) & \texttt{\detokenize{R_7}} & explicit seven-rule ordinary PA presentation & p12-klop1987-pa & n/a & \hyperref[def:p12-klop1987-pa-regime]{\ref*{def:p12-klop1987-pa-regime}} & \thepage & Adds atomic left merge and oriented A3/A4/A5; not full ACP or rewriting modulo choice equations \\
\phantomsection\label{entry:symbol-registry-04950}\(\nu(t)\) & \texttt{\detokenize{\nu}} & positive repaired max/sum norm & p12-klop1987-pa & n/a & \hyperref[def:p12-klop1987-pa-stable-norm]{\ref*{def:p12-klop1987-pa-stable-norm}} & \thepage & Unit action and variable leaves; maximum for choice and sum for three other binary operators \\
\phantomsection\label{entry:symbol-registry-04951}\(E(t)\) & \texttt{\detokenize{E}} & canonical equality-indexed annotation & p12-klop1987-pa & n/a & \hyperref[def:p12-klop1987-pa-stable-norm]{\ref*{def:p12-klop1987-pa-stable-norm}} & \thepage & Merge tag is sum of actual argument norms; not ordinary substitution-homomorphic \\
\phantomsection\label{entry:symbol-registry-04952}\(P_n,L_n\) & \texttt{\detokenize{P_n,}\allowbreak\detokenize{L_n}} & ranked parallel and left-merge operators & p12-klop1987-pa & n/a & \hyperref[def:p12-klop1987-pa-stable-norm]{\ref*{def:p12-klop1987-pa-stable-norm}} & \thepage & Fixed arity two and n at least2; level first then parallel over left merge \\
\phantomsection\label{entry:symbol-registry-04953}\(\widehat C_w\) & \texttt{\detokenize{\widehat C_w}} & common annotated context at fixed hole norm & p12-klop1987-pa & n/a & \hyperref[thm:p12-klop1987-pa-norm]{\ref*{thm:p12-klop1987-pa-norm}} & \thepage & Only norm-equal endpoints use the same literal context \\
\phantomsection\label{entry:symbol-registry-04954}\(Q\) & \texttt{\detokenize{Q}} & isolated embedded-left-side one-rule TRS & p12-klop1987-pa & n/a & \hyperref[thm:p12-klop1987-pa-puel-example]{\ref*{thm:p12-klop1987-pa-puel-example}} & \thepage & No rule for p or s; active-origin proof is not Puel's general extension \\
\phantomsection\label{entry:symbol-registry-05078}\ensuremath{H_q} & \texttt{\detokenize{H_q}} & proof-level evaluation sending all variables to one original ground term & p12-equational-basic-completeness & meta & \hyperref[def:p12-eb-ground-map]{\ref*{def:p12-eb-ground-map}} & \thepage & Not an infinite-support substitution permitted by the unification chapter; each rewrite instance uses a finite restriction \\
\phantomsection\label{entry:symbol-registry-05079}\ensuremath{theta_0} & \texttt{\detokenize{\theta_0}} & initial finite-support image-normalized open witness & p12-equational-basic-completeness & meta & \hyperref[thm:p12-eb-initial-normalization]{\ref*{thm:p12-eb-initial-normalization}} & \thepage & R-equal to the target; no ground constants added and no active S rewritten \\
\phantomsection\label{entry:symbol-registry-05080}len & \texttt{\detokenize{\operatorname{len}}} & number of nodes in a unary term for the explicit reduction ordering & \hyperref[annex:p12.eb.progress-countercase]{\ref*{annex:p12.eb.progress-countercase}} & meta & \hyperref[thm:p12-eb-progress-countercase]{\ref*{thm:p12-eb-progress-countercase}} & \thepage & Includes the terminal leaf; not a cardinality glyph \\
\phantomsection\label{entry:symbol-registry-05081}\(\mathsf{eqb},\mathsf{choose}\) & \texttt{\detokenize{\mathsf{eqb},}\allowbreak\detokenize{\mathsf{choose}}} & computed equality code and data selector & p12/har-autarkic-f & meta & \hyperref[def:p12-har-autarkic-eqb]{\ref*{def:p12-har-autarkic-eqb}} & \thepage & Fixed-target nat\_rec; total O/one iff proof without DN. \\
\phantomsection\label{entry:symbol-registry-05082}\(Q(r),F\) & \texttt{\detokenize{Q(r),}\allowbreak\detokenize{F}} & square function and square-root program & p12/har-autarkic-f & meta & \hyperref[def:p12-har-autarkic-square-root]{\ref*{def:p12-har-autarkic-square-root}} & \thepage & Universal bound and finite 45-to-6 reduction; not kernel replay. \\
\phantomsection\label{entry:symbol-registry-05083}\(\ulcorner k\urcorner_{\mathsf{nat}}\) & \texttt{\detokenize{\ulcorner k\urcorner_{\mathsf{nat}}}} & unary numeral with editorial sort tag & p12/har-autarkic-f & n/a & \hyperref[def:p12-har-autarkic-square-root]{\ref*{def:p12-har-autarkic-square-root}} & \thepage & S\^{}k O data; not syntax quotation or numeric proof labels. \\
\phantomsection\label{entry:symbol-registry-05084}\(T,Z,U,\mathsf{Var},\mathsf{Add},\mathsf{Mul}\) & \texttt{\detokenize{T,}\allowbreak\detokenize{Z,}\allowbreak\detokenize{U,}\allowbreak\detokenize{\mathsf{Var},}\allowbreak\detokenize{\mathsf{Add},}\allowbreak\detokenize{\mathsf{Mul}}} & positive arithmetic syntax constructors & p12/har-autarkic-syntax & n/a & \hyperref[def:p12-har-autarkic-syntax]{\ref*{def:p12-har-autarkic-syntax}} & \thepage & Nat/+/*/0/1 and nat-coded names; not arbitrary lambda syntax. \\
\phantomsection\label{entry:symbol-registry-05085}\(\mathsf{fold}_T^C\) & \texttt{\detokenize{\mathsf{fold}_T^C}} & fixed-target syntax recursion & p12/har-autarkic-syntax & meta & \hyperref[def:p12-har-autarkic-syntax]{\ref*{def:p12-har-autarkic-syntax}} & \thepage & Five curried branches/computations; not top-sort-polymorphic constant. \\
\phantomsection\label{entry:symbol-registry-05086}\(A,0_A,1_A,+_A,\times_A\) & \texttt{\detokenize{A,}\allowbreak\detokenize{0_A,}\allowbreak\detokenize{1_A,}\allowbreak\detokenize{+_A,}\allowbreak\detokenize{\times_A}} & fixed carrier and interpretation parameters & p12/har-autarkic-interpretation & n/a & \hyperref[def:p12-har-autarkic-interpreter]{\ref*{def:p12-har-autarkic-interpreter}} & \thepage & Only typing needed for quote/substitution laws; Nat instance explicit. \\
\phantomsection\label{entry:symbol-registry-05087}\(E_\nu^A,E_\nu\) & \texttt{\detokenize{E_\nu^A,}\allowbreak\detokenize{E_\nu}} & actual valuation-sensitive interpreter & p12/har-autarkic-interpretation & meta & \hyperref[def:p12-har-autarkic-interpreter]{\ref*{def:p12-har-autarkic-interpreter}} & \thepage & Five branches and open laws; distinct from CC raw interpretation. \\
\phantomsection\label{entry:symbol-registry-05088}\(\ulcorner p\urcorner_T\) & \texttt{\detokenize{\ulcorner p\urcorner_T}} & arithmetic quotation with editorial sort tag & p12/har-autarkic-syntax & n/a & \hyperref[def:p12-har-autarkic-interpreter]{\ref*{def:p12-har-autarkic-interpreter}} & \thepage & Five grammar clauses; name indices encoded as nat data. \\
\phantomsection\label{entry:symbol-registry-05089}\(\mathsf{sel}_A,\nu_\Gamma\) & \texttt{\detokenize{\mathsf{sel}_A,}\allowbreak\detokenize{\nu_\Gamma}} & selector and finite valuation extension & p12/har-autarkic-interpretation & meta & \hyperref[def:p12-har-autarkic-interpreter]{\ref*{def:p12-har-autarkic-interpreter}} & \thepage & Distinct literal indices and default zero; assigned-index reductions without choice. \\
\phantomsection\label{entry:symbol-registry-05090}\(\mathsf{sub}_\sigma,\nu_\sigma\) & \texttt{\detokenize{\mathsf{sub}_\sigma,}\allowbreak\detokenize{\nu_\sigma}} & syntax substitution and composed valuation & p12/har-autarkic-interpretation & meta & \hyperref[def:p12-har-autarkic-interpreter]{\ref*{def:p12-har-autarkic-interpreter}} & \thepage & Universal five-branch substitution proof; no extensionality. \\
\phantomsection\label{entry:symbol-registry-05091}\(\mathsf{Bit}\) & \texttt{\detokenize{\mathsf{Bit}}} & Prop code-range predicate & p12/har-kprime & meta & \hyperref[def:p12-har-kprime-bits]{\ref*{def:p12-har-kprime-bits}} & \thepage & Datum equals O or one; not a datatype consumed to select data. \\
\phantomsection\label{entry:symbol-registry-05092}\(\mathsf{andb},\mathsf{orb},\mathsf{notb}\) & \texttt{\detokenize{\mathsf{andb},}\allowbreak\detokenize{\mathsf{orb},}\allowbreak\detokenize{\mathsf{notb}}} & actual 0/1 code operations & p12/har-kprime & meta & \hyperref[def:p12-har-kprime-bits]{\ref*{def:p12-har-kprime-bits}} & \thepage & Fixed nat choose terms; exact range premises and all positive/negative cases proved. \\
\phantomsection\label{entry:symbol-registry-05093}\(l_0,l_S,\mathsf{leb},\mathsf{ltb}\) & \texttt{\detokenize{l_0,}\allowbreak\detokenize{l_S,}\allowbreak\detokenize{\mathsf{leb},}\allowbreak\detokenize{\mathsf{ltb}}} & actual computed source order codes & p12/har-kprime & meta & \hyperref[def:p12-har-kprime-order]{\ref*{def:p12-har-kprime-order}} & \thepage & Nat-to-nat outer target and nat inner target; match leseq/lthan with both iff directions. \\
\phantomsection\label{entry:symbol-registry-05094}\(\mathsf{Any}\) & \texttt{\detokenize{\mathsf{Any}}} & inclusive existential data fold & p12/har-kprime & meta & \hyperref[def:p12-har-kprime-any]{\ref*{def:p12-har-kprime-any}} & \thepage & Scans O through N including endpoints; range/true/false/universal specifications constructed. \\
\phantomsection\label{entry:symbol-registry-05095}\(\mathsf{Coeff},\mathsf{Proper},\mathsf{HasProper}\) & \texttt{\detokenize{\mathsf{Coeff},}\allowbreak\detokenize{\mathsf{Proper},}\allowbreak\detokenize{\mathsf{HasProper}}} & nested coefficient and divisor data scans & p12/har-kprime & meta & \hyperref[def:p12-har-kprime-program]{\ref*{def:p12-har-kprime-program}} & \thepage & Both bounds n; coefficient-left times; no witness extraction from Prop. \\
\phantomsection\label{entry:symbol-registry-05096}\(\mathsf{KPrime}\) & \texttt{\detokenize{\mathsf{KPrime}}} & direct bounded primality program & p12/har-kprime & meta & \hyperref[def:p12-har-kprime-program]{\ref*{def:p12-har-kprime-program}} & \thepage & Source prime characteristic/range/negative proofs; distinct from compiler/Pocklington variants. \\
\phantomsection\label{entry:symbol-registry-05097}\(\overline{k}\) & \texttt{\detokenize{\overline{k}}} & local unary constructor numeral & p12/har-kprime & n/a & \hyperref[thm:p12-har-kprime-correct]{\ref*{thm:p12-har-kprime-correct}} & \thepage & S\^{}k O; no primitive decimal constants; closed0/1/4/61 reductions described mathematically. \\
\phantomsection\label{entry:symbol-registry-05098}\(W,[],i::w\) & \texttt{\detokenize{W,}\allowbreak\detokenize{[],}\allowbreak\detokenize{i::w}} & positive finite variable-word datatype & p12/har-polynomial-normalizer & n/a & \hyperref[def:p12-har-poly-data]{\ref*{def:p12-har-poly-data}} & \thepage & Nat indices sorted nondecreasing; repeated indices retained. \\
\phantomsection\label{entry:symbol-registry-05099}\(R,\mathsf{row}\ c\ w\) & \texttt{\detokenize{R,}\allowbreak\detokenize{\mathsf{row}}} & positive coefficient/word row datatype & p12/har-polynomial-normalizer & n/a & \hyperref[def:p12-har-poly-data]{\ref*{def:p12-har-poly-data}} & \thepage & Nat coefficients; constructor cases explicitly generated. \\
\phantomsection\label{entry:symbol-registry-05100}\(P,\epsilon,r\mathbin{\#}p\) & \texttt{\detokenize{P,}\allowbreak\detokenize{\epsilon,}\allowbreak\detokenize{\#}} & positive finite sparse-polynomial datatype & p12/har-polynomial-normalizer & n/a & \hyperref[def:p12-har-poly-data]{\ref*{def:p12-har-poly-data}} & \thepage & Strict descending lex keys and nonzero coefficients in Good outputs. \\
\phantomsection\label{entry:symbol-registry-05101}\(I_C,\mathrm{Rec}_W,\mathrm{Rec}_R,\mathrm{Rec}_P\) & \texttt{\detokenize{I_C,}\allowbreak\detokenize{\mathrm{Rec}_W,}\allowbreak\detokenize{\mathrm{Rec}_R,}\allowbreak\detokenize{\mathrm{Rec}_P}} & actual fixed-target selection/finite recursion & p12/har-polynomial-normalizer & meta & \hyperref[def:p12-har-poly-data]{\ref*{def:p12-har-poly-data}} & \thepage & All branch types/computations supplied; no top-sort-polymorphic base constant or assumed list library. \\
\phantomsection\label{entry:symbol-registry-05102}\(\mathsf{weq},\mathsf{wlt},<_{\rm lex}\) & \texttt{\detokenize{\mathsf{weq},}\allowbreak\detokenize{\mathsf{wlt},}\allowbreak\detokenize{<_{\rm lex}}} & computed word equality/strict order and relation & p12/har-polynomial-normalizer & relation & \hyperref[def:p12-har-poly-word-code]{\ref*{def:p12-har-poly-word-code}} & \thepage & All code cases plus primitive word equality and total/transitive lex proof. \\
\phantomsection\label{entry:symbol-registry-05103}\(\mathsf{MW}\) & \texttt{\detokenize{\mathsf{MW}}} & nested primitive monomial-word merge & p12/har-polynomial-normalizer & meta & \hyperref[def:p12-har-poly-word-code]{\ref*{def:p12-har-poly-word-code}} & \thepage & Both recursion arguments justified; membership/sortedness/value proofs explicit. \\
\phantomsection\label{entry:symbol-registry-05104}\(\mathrm{Sorted}_W,\mathrm{Good}\) & \texttt{\detokenize{\mathrm{Sorted}_W,}\allowbreak\detokenize{\mathrm{Good}}} & Prop sparse-data invariants & p12/har-polynomial-normalizer & n/a & \hyperref[def:p12-har-poly-invariants]{\ref*{def:p12-har-poly-invariants}} & \thepage & Evidence is not a data-program argument. \\
\phantomsection\label{entry:symbol-registry-05105}\(\chi,V_\nu,B_\nu,L_\nu\) & \texttt{\detokenize{\chi,}\allowbreak\detokenize{V_\nu,}\allowbreak\detokenize{B_\nu,}\allowbreak\detokenize{L_\nu}} & actual coefficient/word/row/polynomial interpretations & p12/har-polynomial-normalizer & meta & \hyperref[def:p12-har-poly-semantics]{\ref*{def:p12-har-poly-semantics}} & \thepage & Fixed commutative-semiring parameters; Nat instance uses supplied laws; not arbitrary uninterpreted operations. \\
\phantomsection\label{entry:symbol-registry-05106}\(\mathsf{mk},\mathsf{Ins},\mathsf{Sum},\mathsf{Scale},\mathsf{Prod}\) & \texttt{\detokenize{\mathsf{mk},}\allowbreak\detokenize{\mathsf{Ins},}\allowbreak\detokenize{\mathsf{Sum},}\allowbreak\detokenize{\mathsf{Scale},}\allowbreak\detokenize{\mathsf{Prod}}} & zero-aware polynomial data operations & p12/har-polynomial-normalizer & meta & \hyperref[def:p12-har-poly-insert]{\ref*{def:p12-har-poly-insert}} & \thepage & Every combination/ordering/zero branch and universal value proof supplied. \\
\phantomsection\label{entry:symbol-registry-05107}\(N\) & \texttt{\detokenize{N}} & actual five-branch arithmetic normalization fold & p12/har-polynomial-normalizer & meta & \hyperref[def:p12-har-poly-normalizer]{\ref*{def:p12-har-poly-normalizer}} & \thepage & T-to-P; Good output and all source-grammar value cases proved. \\
\phantomsection\label{entry:symbol-registry-05108}\(\mathsf{lit},Q_W,Q_R,Q_P,G\) & \texttt{\detokenize{\mathsf{lit},}\allowbreak\detokenize{Q_W,}\allowbreak\detokenize{Q_R,}\allowbreak\detokenize{Q_P,}\allowbreak\detokenize{G}} & actual quotation-back and symbolic normalizer & p12/har-polynomial-normalizer & meta & \hyperref[def:p12-har-poly-quotation]{\ref*{def:p12-har-poly-quotation}} & \thepage & Right-associated sums/no trailing zero or unit; coefficient one omitted; book retraction/idempotence and example supplied. \\
\phantomsection\label{entry:symbol-registry-05304}\(Z_k,\mathsf{Succ},P_i^k\) & \texttt{\detokenize{Z_k,}\allowbreak\detokenize{\mathsf{Succ},}\allowbreak\detokenize{P_i^k}} & ranked initial-function prescriptions & p12/har-primitive-recursive-compiler & meta & \hyperref[def:p12-har-prcompiler-prescriptions]{\ref*{def:p12-har-prcompiler-prescriptions}} & \thepage & Editorial metalevel finite syntax; zeroary zero allowed; projection only k>=1 and 1<=i<=k; not an internal source datatype. \\
\phantomsection\label{entry:symbol-registry-05305}\(\mathsf{Comp}_k,\mathsf{Rec}_k\) & \texttt{\detokenize{\mathsf{Comp}_k,}\allowbreak\detokenize{\mathsf{Rec}_k}} & ranked composition and primitive-recursion prescriptions & p12/har-primitive-recursive-compiler & meta & \hyperref[def:p12-har-prcompiler-prescriptions]{\ref*{def:p12-har-prcompiler-prescriptions}} & \thepage & Composition permits zero arguments; primitive recursion's last argument convention explicit; actual source nat:Set templates. \\
\phantomsection\label{entry:symbol-registry-05306}\(T_k,F_k\) & \texttt{\detokenize{T_k,}\allowbreak\detokenize{F_k}} & curried function and formula-family arity types & p12/har-primitive-recursive-compiler & meta & \hyperref[def:p12-har-prcompiler-prescriptions]{\ref*{def:p12-har-prcompiler-prescriptions}} & \thepage & T0=nat and F0=form; no assumed zeroary unit/product datatype. \\
\phantomsection\label{entry:symbol-registry-05307}\(\mathcal C\) & \texttt{\detokenize{\mathcal C}} & external finite prescription syntax compiler & p12/har-primitive-recursive-compiler & meta & \hyperref[def:p12-har-prcompiler-prescriptions]{\ref*{def:p12-har-prcompiler-prescriptions}} & \thepage & Every initial/composition/recursion template supplied; typing and numerical representation proved; not internal Prop reification. \\
\phantomsection\label{entry:symbol-registry-05308}\(\langle e\rangle,\overline a\) & \texttt{\detokenize{\langle e\rangle,}\allowbreak\detokenize{\overline a}} & independent numerical meaning and unary constructor numeral & p12/har-primitive-recursive-compiler & meta & \hyperref[def:p12-har-prcompiler-denotation]{\ref*{def:p12-har-prcompiler-denotation}} & \thepage & Metalevel denotation independently recursive; numeral is S\^{}a O; universal reduction and standard valuations proved. \\
\phantomsection\label{entry:symbol-registry-05309}\(\mathsf{bool},\mathsf{true},\mathsf{false}\) & \texttt{\detokenize{\mathsf{bool},}\allowbreak\detokenize{\mathsf{true},}\allowbreak\detokenize{\mathsf{false}}} & source inductive boolean data type and constructors & p12/har-primitive-recursive-compiler & meta & \hyperref[def:p12-har-prcompiler-booleans]{\ref*{def:p12-har-prcompiler-booleans}} & \thepage & bool:Set from author preprint Example5; distinct from nat O/one codes. \\
\phantomsection\label{entry:symbol-registry-05310}\(\mathsf{if}_C,\mathsf{istrue},\mathsf{code}\) & \texttt{\detokenize{\mathsf{if}_C,}\allowbreak\detokenize{\mathsf{istrue},}\allowbreak\detokenize{\mathsf{code}}} & actual boolean selector and Prop/nat interpretation bridges & p12/har-primitive-recursive-compiler & meta & \hyperref[def:p12-har-prcompiler-booleans]{\ref*{def:p12-har-prcompiler-booleans}} & \thepage & Generated bool recursor; Prop:Type code target admitted; both positive/negative nat code equivalences derived. \\
\phantomsection\label{entry:symbol-registry-05311}\(\mathsf{b\_not},\mathsf{b\_and},\mathsf{b\_or},\mathsf{b\_imp}\) & \texttt{\detokenize{\mathsf{b\_not},}\allowbreak\detokenize{\mathsf{b\_and},}\allowbreak\detokenize{\mathsf{b\_or},}\allowbreak\detokenize{\mathsf{b\_imp}}} & source boolean connective programs & p12/har-primitive-recursive-compiler & meta & \hyperref[def:p12-har-prcompiler-booleans]{\ref*{def:p12-har-prcompiler-booleans}} & \thepage & All truth-constructor cases and constructive implication equivalence proved. \\
\phantomsection\label{entry:symbol-registry-05312}\(q_0,e_0,d,\mathsf{b\_le},\mathsf{b\_eq},\mathsf{b\_lt},\mathsf{b\_ge},\mathsf{b\_gt}\) & \texttt{\detokenize{q_0,}\allowbreak\detokenize{e_0,}\allowbreak\detokenize{d,}\allowbreak\detokenize{\mathsf{b\_le},}\allowbreak\detokenize{\mathsf{b\_eq},}\allowbreak\detokenize{\mathsf{b\_lt},}\allowbreak\detokenize{\mathsf{b\_ge},}\allowbreak\detokenize{\mathsf{b\_gt}}} & actual fixed-target atomic comparison programs & p12/har-primitive-recursive-compiler & meta & \hyperref[def:p12-har-prcompiler-booleans]{\ref*{def:p12-har-prcompiler-booleans}} & \thepage & Outer nat-to-bool and inner bool recursions; five relation equivalences relative to historical equality/order. \\
\phantomsection\label{entry:symbol-registry-05313}\(\mathsf{b\_all},\mathsf{b\_ex}\) & \texttt{\detokenize{\mathsf{b\_all},}\allowbreak\detokenize{\mathsf{b\_ex}}} & source exclusive-bound data folds & p12/har-primitive-recursive-compiler & meta & \hyperref[def:p12-har-prcompiler-bounds]{\ref*{def:p12-har-prcompiler-bounds}} & \thepage & Scan x<b with endpoint b-1; descending combination; not earlier inclusive Any; both iff directions constructed. \\
\phantomsection\label{entry:symbol-registry-05314}\(\mathsf{form}\) & \texttt{\detokenize{\mathsf{form}}} & source internal formula datatype & p12/har-primitive-recursive-compiler & meta & \hyperref[def:p12-har-prcompiler-form]{\ref*{def:p12-har-prcompiler-form}} & \thepage & Eleven constructors; atom fields are already host nat terms; quantifier family nat->form and generated family induction retained. \\
\phantomsection\label{entry:symbol-registry-05315}\(\mathsf{f\_lt},\mathsf{f\_le},\mathsf{f\_eq},\mathsf{f\_ge},\mathsf{f\_gt},\mathsf{f\_not},\mathsf{f\_and},\mathsf{f\_or},\mathsf{f\_imp},\mathsf{f\_all},\mathsf{f\_ex}\) & \texttt{\detokenize{\mathsf{f\_lt},}\allowbreak\detokenize{\mathsf{f\_le},}\allowbreak\detokenize{\mathsf{f\_eq},}\allowbreak\detokenize{\mathsf{f\_ge},}\allowbreak\detokenize{\mathsf{f\_gt},}\allowbreak\detokenize{\mathsf{f\_not},}\allowbreak\detokenize{\mathsf{f\_and},}\allowbreak\detokenize{\mathsf{f\_or},}\allowbreak\detokenize{\mathsf{f\_imp},}\allowbreak\detokenize{\mathsf{f\_all},}\allowbreak\detokenize{\mathsf{f\_ex}}} & exact source formula constructors & p12/har-primitive-recursive-compiler & meta & \hyperref[def:p12-har-prcompiler-form]{\ref*{def:p12-har-prcompiler-form}} & \thepage & Source Definition7 pp18-19; distinct typed signatures and all eleven fold branches supplied. \\
\phantomsection\label{entry:symbol-registry-05316}\(\mathsf{Rec}_{\mathsf{form}}^C,\mathsf{fold}_C\) & \texttt{\detokenize{\mathsf{Rec}_{\mathsf{form}}^C,}\allowbreak\detokenize{\mathsf{fold}_C}} & source-generated fixed-target formula recursor & p12/har-primitive-recursive-compiler & meta & \hyperref[def:p12-har-prcompiler-form]{\ref*{def:p12-har-prcompiler-form}} & \thepage & Fold is that primitive applied to eleven branch terms, not an extra assumed helper; family recursive result lambda x.fold(hx). \\
\phantomsection\label{entry:symbol-registry-05317}\(\mathcal P,\mathsf{comp},\llbracket-\rrbracket,\llparenthesis-\rrparenthesis\) & \texttt{\detokenize{\mathcal P,}\allowbreak\detokenize{\mathsf{comp},}\allowbreak\detokenize{\llbracket-\rrbracket,}\allowbreak\detokenize{\llparenthesis-\rrparenthesis}} & internal Prop and bool interpreters with source local aliases & p12/har-primitive-recursive-compiler & meta & \hyperref[def:p12-har-prcompiler-interpreters]{\ref*{def:p12-har-prcompiler-interpreters}} & \thepage & Exact source Definition10/11 interpretations; full comp recursor term; P code target Prop:Type. \\
\phantomsection\label{entry:symbol-registry-05318}\(\mathsf{ok},K_p\) & \texttt{\detokenize{\mathsf{ok},}\allowbreak\detokenize{K_p}} & universal internal adequacy proof and computed characteristic instances & p12/har-primitive-recursive-compiler & meta & \hyperref[thm:p12-har-prcompiler-ok]{\ref*{thm:p12-har-prcompiler-ok}} & \thepage & Form induction in all valid open contexts; source bool interpreter then actual nat code bridge; no inverse Prop reifier. \\
\phantomsection\label{entry:symbol-registry-05319}\(R_e,p_e,K_e\) & \texttt{\detokenize{R_e,}\allowbreak\detokenize{p_e,}\allowbreak\detokenize{K_e}} & explicit arbitrary-prescription predicate instance & p12/har-primitive-recursive-compiler & meta & \hyperref[thm:p12-har-prcompiler-predicate]{\ref*{thm:p12-har-prcompiler-predicate}} & \thepage & Equality-to-one construction for every ranked prescription; complete finite 0/1/>=2 output proof; internal open adequacy distinct from standard numerical semantics. \\
\phantomsection\label{entry:symbol-registry-05320}\(\mathsf{f\_Divides},\mathsf{f\_Prime}\) & \texttt{\detokenize{\mathsf{f\_Divides},}\allowbreak\detokenize{\mathsf{f\_Prime}}} & source bounded-division and primality formula-valued inputs & p12/har-primitive-recursive-compiler & meta & \hyperref[def:p12-har-prcompiler-source-prime]{\ref*{def:p12-har-prcompiler-source-prime}} & \thepage & Exact Example9 terms; coefficient k<S m and divisor d<n; generic adequacy applies; no name-based identity with separate unbounded-div/proper-divisor predicate or direct KPrime. \\
\phantomsection\label{entry:symbol-registry-05573}\(X\) & \texttt{\detokenize{X}} & finite propositional variable universe for one input and trace & \hyperref[sec:p12-drat-certificates-and-checking]{\ref*{sec:p12-drat-certificates-and-checking}} & meta & \hyperref[def:p12-drat-finite-cnf]{\ref*{def:p12-drat-finite-cnf}} & \thepage & Contains the union of input and finite trace variables. \\
\phantomsection\label{entry:symbol-registry-05574}\(x\) & \texttt{\detokenize{x}} & propositional variable & \hyperref[sec:p12-drat-certificates-and-checking]{\ref*{sec:p12-drat-certificates-and-checking}} & object & \hyperref[def:p12-drat-finite-cnf]{\ref*{def:p12-drat-finite-cnf}} & \thepage & Indexed x\_i and named p a b c d share this type. \\
\phantomsection\label{entry:symbol-registry-05575}\(p\) & \texttt{\detokenize{p}} & propositional variable in boundary examples & \hyperref[sec:p12-drat-certificates-and-checking]{\ref*{sec:p12-drat-certificates-and-checking}} & object & \hyperref[def:p12-drat-trace]{\ref*{def:p12-drat-trace}} & \thepage & Not a proof identifier or integer code. \\
\phantomsection\label{entry:symbol-registry-05576}\(a\) & \texttt{\detokenize{a}} & propositional variable in the 2014 examples & \hyperref[sec:p12-drat-certificates-and-checking]{\ref*{sec:p12-drat-certificates-and-checking}} & object & \hyperref[def:p12-drat-nonrup-rat]{\ref*{def:p12-drat-nonrup-rat}} & \thepage & The native DIMACS example explicitly decodes 1 through 4 as a b c d. \\
\phantomsection\label{entry:symbol-registry-05577}\(b\) & \texttt{\detokenize{b}} & propositional variable in the 2014 examples & \hyperref[sec:p12-drat-certificates-and-checking]{\ref*{sec:p12-drat-certificates-and-checking}} & object & \hyperref[def:p12-drat-nonrup-rat]{\ref*{def:p12-drat-nonrup-rat}} & \thepage & Same variable role as source Example 1. \\
\phantomsection\label{entry:symbol-registry-05578}\(c\) & \texttt{\detokenize{c}} & propositional variable in the 2014 examples & \hyperref[sec:p12-drat-certificates-and-checking]{\ref*{sec:p12-drat-certificates-and-checking}} & object & \hyperref[def:p12-drat-nonrup-rat]{\ref*{def:p12-drat-nonrup-rat}} & \thepage & Not a textual comment prefix. \\
\phantomsection\label{entry:symbol-registry-05579}\(d\) & \texttt{\detokenize{d}} & propositional variable decoding DIMACS index four & \hyperref[sec:p12-drat-certificates-and-checking]{\ref*{sec:p12-drat-certificates-and-checking}} & object & \hyperref[def:p12-drat-native-refutation]{\ref*{def:p12-drat-native-refutation}} & \thepage & Not the textual deletion prefix. \\
\phantomsection\label{entry:symbol-registry-05580}\(l\) & \texttt{\detokenize{l}} & abstract signed propositional literal & \hyperref[sec:p12-drat-certificates-and-checking]{\ref*{sec:p12-drat-certificates-and-checking}} & object & \hyperref[def:p12-drat-finite-cnf]{\ref*{def:p12-drat-finite-cnf}} & \thepage & Binary serialization explicitly declares a separate local integer-code role. \\
\phantomsection\label{entry:symbol-registry-05581}\(\bar l\) & \texttt{\detokenize{\bar l}} & complement of a signed propositional literal & \hyperref[sec:p12-drat-certificates-and-checking]{\ref*{sec:p12-drat-certificates-and-checking}} & meta & \hyperref[def:p12-drat-finite-cnf]{\ref*{def:p12-drat-finite-cnf}} & \thepage & Source complement notation retained; an involution. \\
\phantomsection\label{entry:symbol-registry-05582}\(C\) & \texttt{\detokenize{C}} & finite propositional clause as a literal set & \hyperref[sec:p12-drat-certificates-and-checking]{\ref*{sec:p12-drat-certificates-and-checking}} & object & \hyperref[def:p12-drat-finite-cnf]{\ref*{def:p12-drat-finite-cnf}} & \thepage & Semantic core allows tautologies; the 2016 parser separately forbids them. \\
\phantomsection\label{entry:symbol-registry-05583}\(D\) & \texttt{\detokenize{D}} & database clause including an opposing-pivot clause & \hyperref[sec:p12-drat-certificates-and-checking]{\ref*{sec:p12-drat-certificates-and-checking}} & object & \hyperref[def:p12-drat-up]{\ref*{def:p12-drat-up}} & \thepage & D\_i names the exact ten decoded 2014 example clauses. \\
\phantomsection\label{entry:symbol-registry-05584}\(E\) & \texttt{\detokenize{E}} & clause extending C in RUP monotonicity & \hyperref[sec:p12-drat-certificates-and-checking]{\ref*{sec:p12-drat-certificates-and-checking}} & object & \hyperref[thm:p12-drat-ala-rup]{\ref*{thm:p12-drat-ala-rup}} & \thepage & Specialized to a pivot-retaining test clause in the bridge proof. \\
\phantomsection\label{entry:symbol-registry-05585}\(F\) & \texttt{\detokenize{F}} & finite CNF clause-set database & \hyperref[sec:p12-drat-certificates-and-checking]{\ref*{sec:p12-drat-certificates-and-checking}} & object & \hyperref[def:p12-drat-finite-cnf]{\ref*{def:p12-drat-finite-cnf}} & \thepage & 2014 set interpretation; the 2016 multiset role is declared separately. \\
\phantomsection\label{entry:symbol-registry-05586}\(G\) & \texttt{\detokenize{G}} & finite CNF database extending F & \hyperref[sec:p12-drat-certificates-and-checking]{\ref*{sec:p12-drat-certificates-and-checking}} & object & \hyperref[thm:p12-drat-ala-rup]{\ref*{thm:p12-drat-ala-rup}} & \thepage & Used for RUP database monotonicity. \\
\phantomsection\label{entry:symbol-registry-05587}\(\models\) & \texttt{\detokenize{\models}} & valuation satisfaction of a propositional clause or CNF & \hyperref[sec:p12-drat-certificates-and-checking]{\ref*{sec:p12-drat-certificates-and-checking}} & relation & \hyperref[def:p12-drat-finite-cnf]{\ref*{def:p12-drat-finite-cnf}} & \thepage & Right argument type specifies clause or database. \\
\phantomsection\label{entry:symbol-registry-05588}\(\operatorname{Sat}(F)\) & \texttt{\detokenize{\operatorname{Sat}(F)}} & existence of a Boolean model of a typed CNF database & \hyperref[sec:p12-drat-certificates-and-checking]{\ref*{sec:p12-drat-certificates-and-checking}} & meta & \hyperref[def:p12-drat-finite-cnf]{\ref*{def:p12-drat-finite-cnf}} & \thepage & Deletion preserves only the forward implication. \\
\phantomsection\label{entry:symbol-registry-05589}\(\varnothing\) & \texttt{\detokenize{\varnothing}} & empty literal-set clause & \hyperref[def:p12-drat-finite-cnf]{\ref*{def:p12-drat-finite-cnf}} & object & \hyperref[def:p12-drat-finite-cnf]{\ref*{def:p12-drat-finite-cnf}} & \thepage & False-clause type; genuine shared glyph retained. \\
\phantomsection\label{entry:symbol-registry-05590}\(\varnothing\) & \texttt{\detokenize{\varnothing}} & empty clause-set database & \hyperref[def:p12-drat-finite-cnf]{\ref*{def:p12-drat-finite-cnf}} & object & \hyperref[def:p12-drat-finite-cnf]{\ref*{def:p12-drat-finite-cnf}} & \thepage & True-conjunction type; distinct from the empty clause. \\
\phantomsection\label{entry:symbol-registry-05591}\(\tau\) & \texttt{\detokenize{\tau}} & consistent ordered unit-propagation trail & \hyperref[sec:p12-drat-certificates-and-checking]{\ref*{sec:p12-drat-certificates-and-checking}} & meta & \hyperref[def:p12-drat-up]{\ref*{def:p12-drat-up}} & \thepage & No solver decision levels or backjump state. \\
\phantomsection\label{entry:symbol-registry-05592}\(A\) & \texttt{\detokenize{A}} & set of simultaneously true assumption literals & \hyperref[sec:p12-drat-certificates-and-checking]{\ref*{sec:p12-drat-certificates-and-checking}} & object & \hyperref[def:p12-drat-up]{\ref*{def:p12-drat-up}} & \thepage & Complementary assumptions conflict immediately. \\
\phantomsection\label{entry:symbol-registry-05593}\(\operatorname{res}_{\tau}(D)\) & \texttt{\detokenize{\operatorname{res}_{\tau}(D)}} & residual of a clause not already satisfied by the trail & \hyperref[sec:p12-drat-certificates-and-checking]{\ref*{sec:p12-drat-certificates-and-checking}} & meta & \hyperref[def:p12-drat-up]{\ref*{def:p12-drat-up}} & \thepage & False literals are deleted only after checking satisfaction. \\
\phantomsection\label{entry:symbol-registry-05594}\(\operatorname{UP}(F,A)\) & \texttt{\detokenize{\operatorname{UP}(F,}\allowbreak\detokenize{A)}} & outcome of the finite unit-conflict test & \hyperref[sec:p12-drat-certificates-and-checking]{\ref*{sec:p12-drat-certificates-and-checking}} & meta & \hyperref[def:p12-drat-up]{\ref*{def:p12-drat-up}} & \thepage & Saturation without conflict is not a SAT claim. \\
\phantomsection\label{entry:symbol-registry-05595}\(\bot\) & \texttt{\detokenize{\bot}} & unit-conflict outcome token & \hyperref[sec:p12-drat-certificates-and-checking]{\ref*{sec:p12-drat-certificates-and-checking}} & object & \hyperref[def:p12-drat-up]{\ref*{def:p12-drat-up}} & \thepage & Not an object-language constant or an empty database. \\
\phantomsection\label{entry:symbol-registry-05596}\(\operatorname{ALA}(F,C)\) & \texttt{\detokenize{\operatorname{ALA}(F,}\allowbreak\detokenize{C)}} & 2014 least asymmetric-literal-addition closure & \hyperref[sec:p12-drat-certificates-and-checking]{\ref*{sec:p12-drat-certificates-and-checking}} & meta & \hyperref[def:p12-drat-ala-rup]{\ref*{def:p12-drat-ala-rup}} & \thepage & Empty input clauses need the explicit terminal conflict case. \\
\phantomsection\label{entry:symbol-registry-05597}\(\bar C\) & \texttt{\detokenize{\bar C}} & complementary unit assumptions for a clause & \hyperref[sec:p12-drat-certificates-and-checking]{\ref*{sec:p12-drat-certificates-and-checking}} & meta & \hyperref[def:p12-drat-ala-rup]{\ref*{def:p12-drat-ala-rup}} & \thepage & Conjunctive assumptions rather than their disjunction. \\
\phantomsection\label{entry:symbol-registry-05598}\(\operatorname{RUP}(F,C)\) & \texttt{\detokenize{\operatorname{RUP}(F,}\allowbreak\detokenize{C)}} & reverse-unit-propagation redundancy predicate & \hyperref[sec:p12-drat-certificates-and-checking]{\ref*{sec:p12-drat-certificates-and-checking}} & meta & \hyperref[def:p12-drat-ala-rup]{\ref*{def:p12-drat-ala-rup}} & \thepage & Checked against the pre-addition database. \\
\phantomsection\label{entry:symbol-registry-05599}\(\operatorname{Res}_l(C,D)\) & \texttt{\detokenize{\operatorname{Res}_l(C,}\allowbreak\detokenize{D)}} & ordinary resolvent removing the designated pivot pair & \hyperref[sec:p12-drat-certificates-and-checking]{\ref*{sec:p12-drat-certificates-and-checking}} & meta & \hyperref[def:p12-drat-rat]{\ref*{def:p12-drat-rat}} & \thepage & The bridge also needs the opposing clause active in F. \\
\phantomsection\label{entry:symbol-registry-05600}\(\operatorname{RAT}(F,C,l)\) & \texttt{\detokenize{\operatorname{RAT}(F,}\allowbreak\detokenize{C,}\allowbreak\detokenize{l)}} & 2014 universal pivot-retaining RAT redundancy predicate & \hyperref[sec:p12-drat-certificates-and-checking]{\ref*{sec:p12-drat-certificates-and-checking}} & meta & \hyperref[def:p12-drat-rat]{\ref*{def:p12-drat-rat}} & \thepage & All opposing active clauses are tested; some pivot must exist. \\
\phantomsection\label{entry:symbol-registry-05601}\(P\) & \texttt{\detokenize{P}} & ordered finite addition-deletion certificate trace & \hyperref[sec:p12-drat-certificates-and-checking]{\ref*{sec:p12-drat-certificates-and-checking}} & meta & \hyperref[def:p12-drat-trace]{\ref*{def:p12-drat-trace}} & \thepage & Not a proposition letter or projection set. \\
\phantomsection\label{entry:symbol-registry-05602}\(L_i\) & \texttt{\detokenize{L_i}} & clause recorded at certificate step i & \hyperref[sec:p12-drat-certificates-and-checking]{\ref*{sec:p12-drat-certificates-and-checking}} & object & \hyperref[def:p12-drat-trace]{\ref*{def:p12-drat-trace}} & \thepage & First literal remains designated by the ordered serialization. \\
\phantomsection\label{entry:symbol-registry-05603}\(F^i_P\) & \texttt{\detokenize{F^i_P}} & 2014 clause-set database after prefix i of trace P & \hyperref[sec:p12-drat-certificates-and-checking]{\ref*{sec:p12-drat-certificates-and-checking}} & object & \hyperref[def:p12-drat-trace]{\ref*{def:p12-drat-trace}} & \thepage & Source superscript-subscript update notation retained. \\
\phantomsection\label{entry:symbol-registry-05604}\(s\) & \texttt{\detokenize{s}} & variable-prefix assignment in the finite construction & \hyperref[sec:p12-drat-certificates-and-checking]{\ref*{sec:p12-drat-certificates-and-checking}} & meta & \hyperref[def:p12-drat-truth-tree]{\ref*{def:p12-drat-truth-tree}} & \thepage & Full and empty prefixes are included. \\
\phantomsection\label{entry:symbol-registry-05605}\(B_s\) & \texttt{\detokenize{B_s}} & blocking clause for extensions of prefix assignment s & \hyperref[sec:p12-drat-certificates-and-checking]{\ref*{sec:p12-drat-certificates-and-checking}} & object & \hyperref[def:p12-drat-truth-tree]{\ref*{def:p12-drat-truth-tree}} & \thepage & Opposite literals; empty prefix gives the empty clause. \\
\phantomsection\label{entry:symbol-registry-05606}\(i\) & \texttt{\detokenize{i}} & nonnegative certificate-prefix index & \hyperref[sec:p12-drat-certificates-and-checking]{\ref*{sec:p12-drat-certificates-and-checking}} & meta & \hyperref[def:p12-drat-trace]{\ref*{def:p12-drat-trace}} & \thepage & Zero denotes the input database. \\
\phantomsection\label{entry:symbol-registry-05607}\(k\) & \texttt{\detokenize{k}} & finite trail length or displayed finite clause arity & \hyperref[sec:p12-drat-certificates-and-checking]{\ref*{sec:p12-drat-certificates-and-checking}} & meta & \hyperref[def:p12-drat-up]{\ref*{def:p12-drat-up}} & \thepage & Later prefix length is also a nonnegative finite index. \\
\phantomsection\label{entry:symbol-registry-05608}\(n\) & \texttt{\detokenize{n}} & number of input variables in the completeness construction & \hyperref[sec:p12-drat-certificates-and-checking]{\ref*{sec:p12-drat-certificates-and-checking}} & meta & \hyperref[def:p12-drat-truth-tree]{\ref*{def:p12-drat-truth-tree}} & \thepage & Includes n equals zero. \\
\phantomsection\label{entry:symbol-registry-05609}\(l\) & \texttt{\detokenize{l}} & nonzero signed integer literal code in 2016 serialization & \hyperref[sec:p12-drat-binary-2016]{\ref*{sec:p12-drat-binary-2016}} & object & \hyperref[sec:p12-drat-binary-2016]{\ref*{sec:p12-drat-binary-2016}} & \thepage & Explicit local retyping needs a root-owned override of the abstract-literal role. \\
\phantomsection\label{entry:symbol-registry-05610}\(F\) & \texttt{\detokenize{F}} & CNF multiset in the separately dated 2016 interpretation & \hyperref[def:p12-drat-format-2016]{\ref*{def:p12-drat-format-2016}} & object & \hyperref[def:p12-drat-format-2016]{\ref*{def:p12-drat-format-2016}} & \thepage & Single-copy deletion; not identified with 2014 set trace behavior. \\
\phantomsection\label{entry:symbol-registry-05611}\(K_j\) & \texttt{\detokenize{K_j}} & finite stage of the least ALA clause closure & \hyperref[annex:p12-drat-ala-rup]{\ref*{annex:p12-drat-ala-rup}} & object & \hyperref[annex:p12-drat-ala-rup]{\ref*{annex:p12-drat-ala-rup}} & \thepage & Index j is a nonnegative round number. \\
\phantomsection\label{entry:symbol-registry-05612}\(K_*\) & \texttt{\detokenize{K_*}} & least ALA-closed literal set containing the seed & \hyperref[annex:p12-drat-ala-rup]{\ref*{annex:p12-drat-ala-rup}} & object & \hyperref[annex:p12-drat-ala-rup]{\ref*{annex:p12-drat-ala-rup}} & \thepage & Computed over a finite two-polarity universe. \\
\phantomsection\label{entry:symbol-registry-05613}\(K\) & \texttt{\detokenize{K}} & arbitrary ALA-closed literal set containing the seed & \hyperref[annex:p12-drat-ala-rup]{\ref*{annex:p12-drat-ala-rup}} & object & \hyperref[annex:p12-drat-ala-rup]{\ref*{annex:p12-drat-ala-rup}} & \thepage & Used to prove leastness. \\
\phantomsection\label{entry:symbol-registry-05614}\(j\) & \texttt{\detokenize{j}} & nonnegative ALA closure-round index & \hyperref[annex:p12-drat-ala-rup]{\ref*{annex:p12-drat-ala-rup}} & meta & \hyperref[annex:p12-drat-ala-rup]{\ref*{annex:p12-drat-ala-rup}} & \thepage & No ordinal or transfinite step is used. \\
\phantomsection\label{entry:symbol-registry-05615}\(T\) & \texttt{\detokenize{T}} & true-literal set of a saturated nonconflicting trail & \hyperref[annex:p12-drat-ala-rup]{\ref*{annex:p12-drat-ala-rup}} & object & \hyperref[annex:p12-drat-ala-rup]{\ref*{annex:p12-drat-ala-rup}} & \thepage & Consistent by its trail type. \\
\phantomsection\label{entry:symbol-registry-05616}\(\bar T\) & \texttt{\detokenize{\bar T}} & literal set made false by the saturated trail & \hyperref[annex:p12-drat-ala-rup]{\ref*{annex:p12-drat-ala-rup}} & object & \hyperref[annex:p12-drat-ala-rup]{\ref*{annex:p12-drat-ala-rup}} & \thepage & Shown closed under ALA. \\
\phantomsection\label{entry:symbol-registry-05617}\(R\) & \texttt{\detokenize{R}} & ordinary resolvent clause in the bridge proof & \hyperref[annex:p12-drat-resolvent-test]{\ref*{annex:p12-drat-resolvent-test}} & object & \hyperref[annex:p12-drat-resolvent-test]{\ref*{annex:p12-drat-resolvent-test}} & \thepage & R equals Res\_l(C,D). \\
\phantomsection\label{entry:symbol-registry-05618}\(E\) & \texttt{\detokenize{E}} & pivot-retaining source test clause in the bridge proof & \hyperref[annex:p12-drat-resolvent-test]{\ref*{annex:p12-drat-resolvent-test}} & object & \hyperref[annex:p12-drat-resolvent-test]{\ref*{annex:p12-drat-resolvent-test}} & \thepage & E equals R union the singleton pivot. \\
\phantomsection\label{entry:symbol-registry-05619}\(K_*\) & \texttt{\detokenize{K_*}} & least ALA closure of the ordinary resolvent & \hyperref[annex:p12-drat-resolvent-test]{\ref*{annex:p12-drat-resolvent-test}} & object & \hyperref[annex:p12-drat-resolvent-test]{\ref*{annex:p12-drat-resolvent-test}} & \thepage & The active opposing clause licenses the retained pivot. \\
\phantomsection\label{entry:symbol-registry-05620}\(w\) & \texttt{\detokenize{w}} & pivot-repaired total Boolean valuation & \hyperref[annex:p12-drat-rat-preservation]{\ref*{annex:p12-drat-rat-preservation}} & meta & \hyperref[annex:p12-drat-rat-preservation]{\ref*{annex:p12-drat-rat-preservation}} & \thepage & Agrees with v off the pivot variable. \\
\phantomsection\label{entry:symbol-registry-05621}\(\operatorname{Var}(F)\) & \texttt{\detokenize{\operatorname{Var}(F)}} & finite variable set occurring in the input CNF & \hyperref[annex:p12-drat-rat-preservation]{\ref*{annex:p12-drat-rat-preservation}} & meta & \hyperref[annex:p12-drat-rat-preservation]{\ref*{annex:p12-drat-rat-preservation}} & \thepage & Extend its valuations to the combined trace-variable universe. \\
\phantomsection\label{entry:symbol-registry-05622}\(\varnothing\) & \texttt{\detokenize{\varnothing}} & empty variable-prefix assignment & \hyperref[def:p12-drat-truth-tree]{\ref*{def:p12-drat-truth-tree}} & object & \hyperref[def:p12-drat-truth-tree]{\ref*{def:p12-drat-truth-tree}} & \thepage & The subscript of B\_varnothing is an empty assignment; the resulting clause is the separately typed empty clause. \\
\phantomsection\label{entry:symbol-registry-05804}\(X\) & \texttt{\detokenize{X}} & fixed finite propositional variable set & \hyperref[sec:p12-cdcl-uip-and-watched-literals]{\ref*{sec:p12-cdcl-uip-and-watched-literals}} & meta & \hyperref[def:p12-cdcl-uip-state]{\ref*{def:p12-cdcl-uip-state}} & \thepage & No fresh variables in the proved core. \\
\phantomsection\label{entry:symbol-registry-05805}\(\varphi\) & \texttt{\detokenize{\varphi}} & active CNF clause-set database & \hyperref[sec:p12-cdcl-uip-and-watched-literals]{\ref*{sec:p12-cdcl-uip-and-watched-literals}} & object & \hyperref[def:p12-cdcl-uip-state]{\ref*{def:p12-cdcl-uip-state}} & \thepage & Older source notation retained. \\
\phantomsection\label{entry:symbol-registry-05806}\(\varphi_0\) & \texttt{\detokenize{\varphi_0}} & immutable input CNF database & \hyperref[sec:p12-cdcl-uip-and-watched-literals]{\ref*{sec:p12-cdcl-uip-and-watched-literals}} & object & \hyperref[def:p12-cdcl-uip-state]{\ref*{def:p12-cdcl-uip-state}} & \thepage & Original clauses are retained. \\
\phantomsection\label{entry:symbol-registry-05807}\(\omega\) & \texttt{\detokenize{\omega}} & finite literal-set clause & \hyperref[sec:p12-cdcl-uip-and-watched-literals]{\ref*{sec:p12-cdcl-uip-and-watched-literals}} & object & \hyperref[def:p12-cdcl-uip-state]{\ref*{def:p12-cdcl-uip-state}} & \thepage & Omega\_i names source example clauses. \\
\phantomsection\label{entry:symbol-registry-05808}\(\nu\) & \texttt{\detokenize{\nu}} & partial Boolean assignment & \hyperref[sec:p12-cdcl-uip-and-watched-literals]{\ref*{sec:p12-cdcl-uip-and-watched-literals}} & meta & \hyperref[def:p12-cdcl-uip-state]{\ref*{def:p12-cdcl-uip-state}} & \thepage & Map to zero u and one; not a three-valued logic. \\
\phantomsection\label{entry:symbol-registry-05809}\(u\) & \texttt{\detokenize{u}} & unassigned marker in the older source convention & \hyperref[def:p12-cdcl-uip-state]{\ref*{def:p12-cdcl-uip-state}} & object & \hyperref[def:p12-cdcl-uip-state]{\ref*{def:p12-cdcl-uip-state}} & \thepage & Later failure example explicitly declares a variable role. \\
\phantomsection\label{entry:symbol-registry-05810}\(l^\nu\) & \texttt{\detokenize{l^\nu}} & partial evaluation of a signed literal & \hyperref[def:p12-cdcl-uip-state]{\ref*{def:p12-cdcl-uip-state}} & meta & \hyperref[def:p12-cdcl-uip-state]{\ref*{def:p12-cdcl-uip-state}} & \thepage & Negation leaves the unassigned marker fixed. \\
\phantomsection\label{entry:symbol-registry-05811}\(T\) & \texttt{\detokenize{T}} & consistent ordered true-literal trail & \hyperref[sec:p12-cdcl-uip-and-watched-literals]{\ref*{sec:p12-cdcl-uip-and-watched-literals}} & meta & \hyperref[def:p12-cdcl-uip-state]{\ref*{def:p12-cdcl-uip-state}} & \thepage & Decision-level blocks. \\
\phantomsection\label{entry:symbol-registry-05812}\(d\) & \texttt{\detokenize{d}} & current nonnegative decision level & \hyperref[sec:p12-cdcl-uip-and-watched-literals]{\ref*{sec:p12-cdcl-uip-and-watched-literals}} & meta & \hyperref[def:p12-cdcl-uip-state]{\ref*{def:p12-cdcl-uip-state}} & \thepage & Root zero; one decision per positive block. \\
\phantomsection\label{entry:symbol-registry-05813}\(\delta(x)\) & \texttt{\detokenize{\delta(x)}} & level of a variable assignment & \hyperref[sec:p12-cdcl-uip-and-watched-literals]{\ref*{sec:p12-cdcl-uip-and-watched-literals}} & meta & \hyperref[def:p12-cdcl-uip-state]{\ref*{def:p12-cdcl-uip-state}} & \thepage & Minus one when unassigned; two source level policies distinguished. \\
\phantomsection\label{entry:symbol-registry-05814}\(\delta(l)\) & \texttt{\detokenize{\delta(l)}} & level of a signed literal's variable & \hyperref[sec:p12-cdcl-uip-and-watched-literals]{\ref*{sec:p12-cdcl-uip-and-watched-literals}} & meta & \hyperref[def:p12-cdcl-uip-analyze]{\ref*{def:p12-cdcl-uip-analyze}} & \thepage & Independent of polarity. \\
\phantomsection\label{entry:symbol-registry-05815}\(\alpha(x)\) & \texttt{\detokenize{\alpha(x)}} & retained antecedent of a propagated variable & \hyperref[sec:p12-cdcl-uip-and-watched-literals]{\ref*{sec:p12-cdcl-uip-and-watched-literals}} & meta & \hyperref[def:p12-cdcl-uip-state]{\ref*{def:p12-cdcl-uip-state}} & \thepage & Alpha(p) means reason of the literal's variable. \\
\phantomsection\label{entry:symbol-registry-05816}\(\mathrm{NIL}\) & \texttt{\detokenize{\mathrm{NIL}}} & absence-of-reason marker & \hyperref[sec:p12-cdcl-uip-and-watched-literals]{\ref*{sec:p12-cdcl-uip-and-watched-literals}} & object & \hyperref[def:p12-cdcl-uip-state]{\ref*{def:p12-cdcl-uip-state}} & \thepage & Not an empty clause. \\
\phantomsection\label{entry:symbol-registry-05817}\(Q\) & \texttt{\detokenize{Q}} & queue of unprocessed assignment watch events & \hyperref[sec:p12-cdcl-uip-and-watched-literals]{\ref*{sec:p12-cdcl-uip-and-watched-literals}} & meta & \hyperref[def:p12-cdcl-uip-state]{\ref*{def:p12-cdcl-uip-state}} & \thepage & In-progress list references also remain pending. \\
\phantomsection\label{entry:symbol-registry-05818}\(p\) & \texttt{\detokenize{p}} & true signed assignment literal & \hyperref[sec:p12-cdcl-uip-and-watched-literals]{\ref*{sec:p12-cdcl-uip-and-watched-literals}} & object & \hyperref[def:p12-cdcl-uip-state]{\ref*{def:p12-cdcl-uip-state}} & \thepage & Generic enqueue literal; scoped examples use a variable. \\
\phantomsection\label{entry:symbol-registry-05819}\(x\) & \texttt{\detokenize{x}} & propositional variable & \hyperref[sec:p12-cdcl-uip-and-watched-literals]{\ref*{sec:p12-cdcl-uip-and-watched-literals}} & object & \hyperref[def:p12-cdcl-uip-state]{\ref*{def:p12-cdcl-uip-state}} & \thepage & Indexed variables preserve source names. \\
\phantomsection\label{entry:symbol-registry-05820}\(\kappa\) & \texttt{\detokenize{\kappa}} & implication-graph conflict node & \hyperref[sec:p12-cdcl-uip-and-watched-literals]{\ref*{sec:p12-cdcl-uip-and-watched-literals}} & object & \hyperref[def:p12-cdcl-uip-graph]{\ref*{def:p12-cdcl-uip-graph}} & \thepage & Not a cardinal or variable. \\
\phantomsection\label{entry:symbol-registry-05821}\(\omega_L\) & \texttt{\detokenize{\omega_L}} & temporary and final learned clause & \hyperref[sec:p12-cdcl-uip-and-watched-literals]{\ref*{sec:p12-cdcl-uip-and-watched-literals}} & object & \hyperref[def:p12-cdcl-uip-analyze]{\ref*{def:p12-cdcl-uip-analyze}} & \thepage & Input consequence false on the conflict trail. \\
\phantomsection\label{entry:symbol-registry-05822}\(a\) & \texttt{\detokenize{a}} & asserting literal of a learned clause & \hyperref[sec:p12-cdcl-uip-and-watched-literals]{\ref*{sec:p12-cdcl-uip-and-watched-literals}} & object & \hyperref[def:p12-cdcl-uip-analyze]{\ref*{def:p12-cdcl-uip-analyze}} & \thepage & Its complement is the UIP. \\
\phantomsection\label{entry:symbol-registry-05823}\(b\) & \texttt{\detokenize{b}} & maximal nonasserting level and backjump target & \hyperref[sec:p12-cdcl-uip-and-watched-literals]{\ref*{sec:p12-cdcl-uip-and-watched-literals}} & meta & \hyperref[def:p12-cdcl-uip-analyze]{\ref*{def:p12-cdcl-uip-analyze}} & \thepage & Zero for a learned unit. \\
\phantomsection\label{entry:symbol-registry-05824}\(w_1(\omega)\) & \texttt{\detokenize{w_1(\omega)}} & first watched clause position & \hyperref[sec:p12-cdcl-uip-and-watched-literals]{\ref*{sec:p12-cdcl-uip-and-watched-literals}} & meta & \hyperref[def:p12-cdcl-uip-watches]{\ref*{def:p12-cdcl-uip-watches}} & \thepage & Distinct watched positions. \\
\phantomsection\label{entry:symbol-registry-05825}\(w_2(\omega)\) & \texttt{\detokenize{w_2(\omega)}} & second watched clause position & \hyperref[sec:p12-cdcl-uip-and-watched-literals]{\ref*{sec:p12-cdcl-uip-and-watched-literals}} & meta & \hyperref[def:p12-cdcl-uip-watches]{\ref*{def:p12-cdcl-uip-watches}} & \thepage & Maximal nonasserting level on insertion. \\
\phantomsection\label{entry:symbol-registry-05826}\(W(l)\) & \texttt{\detokenize{W(l)}} & live clauses watching literal l & \hyperref[sec:p12-cdcl-uip-and-watched-literals]{\ref*{sec:p12-cdcl-uip-and-watched-literals}} & meta & \hyperref[def:p12-cdcl-uip-watches]{\ref*{def:p12-cdcl-uip-watches}} & \thepage & False-literal versus true-complement indexing distinguished. \\
\phantomsection\label{entry:symbol-registry-05827}\(q\) & \texttt{\detokenize{q}} & other watched signed literal & \hyperref[def:p12-cdcl-uip-watch-update]{\ref*{def:p12-cdcl-uip-watch-update}} & object & \hyperref[def:p12-cdcl-uip-watch-update]{\ref*{def:p12-cdcl-uip-watch-update}} & \thepage & True false or unassigned. \\
\phantomsection\label{entry:symbol-registry-05828}\(u\) & \texttt{\detokenize{u}} & propositional variable in the failure example & \hyperref[def:p12-cdcl-uip-watch-failure]{\ref*{def:p12-cdcl-uip-watch-failure}} & object & \hyperref[def:p12-cdcl-uip-watch-failure]{\ref*{def:p12-cdcl-uip-watch-failure}} & \thepage & Explicit retyping; not the unassigned marker. \\
\phantomsection\label{entry:symbol-registry-05829}\(p\) & \texttt{\detokenize{p}} & propositional variable in the failure example & \hyperref[def:p12-cdcl-uip-watch-failure]{\ref*{def:p12-cdcl-uip-watch-failure}} & object & \hyperref[def:p12-cdcl-uip-watch-failure]{\ref*{def:p12-cdcl-uip-watch-failure}} & \thepage & Not the generic signed enqueue literal. \\
\phantomsection\label{entry:symbol-registry-05830}\(q\) & \texttt{\detokenize{q}} & propositional variable in the failure example & \hyperref[def:p12-cdcl-uip-watch-failure]{\ref*{def:p12-cdcl-uip-watch-failure}} & object & \hyperref[def:p12-cdcl-uip-watch-failure]{\ref*{def:p12-cdcl-uip-watch-failure}} & \thepage & Not the other-watch parameter. \\
\phantomsection\label{entry:symbol-registry-05831}\(r\) & \texttt{\detokenize{r}} & propositional variable in the failure example & \hyperref[def:p12-cdcl-uip-watch-failure]{\ref*{def:p12-cdcl-uip-watch-failure}} & object & \hyperref[def:p12-cdcl-uip-watch-failure]{\ref*{def:p12-cdcl-uip-watch-failure}} & \thepage & Negative literal forced at level two. \\
\phantomsection\label{entry:symbol-registry-05832}\(s\) & \texttt{\detokenize{s}} & propositional variable in the failure example & \hyperref[def:p12-cdcl-uip-watch-failure]{\ref*{def:p12-cdcl-uip-watch-failure}} & object & \hyperref[def:p12-cdcl-uip-watch-failure]{\ref*{def:p12-cdcl-uip-watch-failure}} & \thepage & Worked resolution pivot. \\
\phantomsection\label{entry:symbol-registry-05833}\(\Delta\) & \texttt{\detokenize{\Delta}} & positive lazy activity-bump increment & \hyperref[sec:p12-cdcl-uip-and-watched-literals]{\ref*{sec:p12-cdcl-uip-and-watched-literals}} & meta & \hyperref[def:p12-cdcl-uip-activity]{\ref*{def:p12-cdcl-uip-activity}} & \thepage & Rescaled with all scores. \\
\phantomsection\label{entry:symbol-registry-05834}\(\rho\) & \texttt{\detokenize{\rho}} & positive decay factor below one & \hyperref[sec:p12-cdcl-uip-and-watched-literals]{\ref*{sec:p12-cdcl-uip-and-watched-literals}} & meta & \hyperref[thm:p12-cdcl-uip-activity-scale]{\ref*{thm:p12-cdcl-uip-activity-scale}} & \thepage & Exact-real theorem only. \\
\phantomsection\label{entry:symbol-registry-05835}\(l\) & \texttt{\detokenize{l}} & signed consequence literal in look-ahead & \hyperref[def:p12-cdcl-uip-lookahead]{\ref*{def:p12-cdcl-uip-lookahead}} & object & \hyperref[def:p12-cdcl-uip-lookahead]{\ref*{def:p12-cdcl-uip-lookahead}} & \thepage & Prefix-relative conclusion. \\
\phantomsection\label{entry:symbol-registry-05836}\(A\) & \texttt{\detokenize{A}} & nonpivot positive-parent clause remainder & \hyperref[annex:p12-cdcl-uip-local-sound]{\ref*{annex:p12-cdcl-uip-local-sound}} & object & \hyperref[annex:p12-cdcl-uip-local-sound]{\ref*{annex:p12-cdcl-uip-local-sound}} & \thepage & Literal-set clause. \\
\phantomsection\label{entry:symbol-registry-05837}\(B\) & \texttt{\detokenize{B}} & nonpivot negative-parent clause remainder & \hyperref[annex:p12-cdcl-uip-local-sound]{\ref*{annex:p12-cdcl-uip-local-sound}} & object & \hyperref[annex:p12-cdcl-uip-local-sound]{\ref*{annex:p12-cdcl-uip-local-sound}} & \thepage & Union with A is the resolvent. \\
\phantomsection\label{entry:symbol-registry-05838}\(z\) & \texttt{\detokenize{z}} & candidate downstream UIP node & \hyperref[annex:p12-cdcl-uip-analysis]{\ref*{annex:p12-cdcl-uip-analysis}} & object & \hyperref[annex:p12-cdcl-uip-analysis]{\ref*{annex:p12-cdcl-uip-analysis}} & \thepage & Acyclic trail order. \\
\phantomsection\label{entry:symbol-registry-05839}\(r\) & \texttt{\detokenize{r}} & nonfalse replacement watched literal & \hyperref[annex:p12-cdcl-uip-watch-correct]{\ref*{annex:p12-cdcl-uip-watch-correct}} & object & \hyperref[annex:p12-cdcl-uip-watch-correct]{\ref*{annex:p12-cdcl-uip-watch-correct}} & \thepage & Not the body's example variable. \\
\phantomsection\label{entry:symbol-registry-05840}\(s_t(x)\) & \texttt{\detokenize{s_t(x)}} & eager score after t decays & \hyperref[annex:p12-cdcl-uip-activity-scale]{\ref*{annex:p12-cdcl-uip-activity-scale}} & meta & \hyperref[annex:p12-cdcl-uip-activity-scale]{\ref*{annex:p12-cdcl-uip-activity-scale}} & \thepage & Ideal exact-real score. \\
\phantomsection\label{entry:symbol-registry-05841}\(a_t(x)\) & \texttt{\detokenize{a_t(x)}} & stored lazy score after t decays & \hyperref[annex:p12-cdcl-uip-activity-scale]{\ref*{annex:p12-cdcl-uip-activity-scale}} & meta & \hyperref[annex:p12-cdcl-uip-activity-scale]{\ref*{annex:p12-cdcl-uip-activity-scale}} & \thepage & Not an asserting literal. \\
\phantomsection\label{entry:symbol-registry-05842}\(h_t(x)\) & \texttt{\detokenize{h_t(x)}} & common step-t bump vector & \hyperref[annex:p12-cdcl-uip-activity-scale]{\ref*{annex:p12-cdcl-uip-activity-scale}} & meta & \hyperref[annex:p12-cdcl-uip-activity-scale]{\ref*{annex:p12-cdcl-uip-activity-scale}} & \thepage & Identical arbitrary bump sequences. \\
\phantomsection\label{entry:symbol-registry-05843}\(\Delta_t\) & \texttt{\detokenize{\Delta_t}} & lazy bump increment at update t & \hyperref[annex:p12-cdcl-uip-activity-scale]{\ref*{annex:p12-cdcl-uip-activity-scale}} & meta & \hyperref[annex:p12-cdcl-uip-activity-scale]{\ref*{annex:p12-cdcl-uip-activity-scale}} & \thepage & Positive, before common rescaling. \\
\phantomsection\label{entry:symbol-registry-05844}\(t\) & \texttt{\detokenize{t}} & nonnegative activity-update index & \hyperref[annex:p12-cdcl-uip-activity-scale]{\ref*{annex:p12-cdcl-uip-activity-scale}} & meta & \hyperref[annex:p12-cdcl-uip-activity-scale]{\ref*{annex:p12-cdcl-uip-activity-scale}} & \thepage & Finite induction index. \\
\phantomsection\label{entry:symbol-registry-05845}\(p_i\) & \texttt{\detokenize{p_i}} & true prefix literal in a global consequence clause & \hyperref[annex:p12-cdcl-uip-lookahead-sound]{\ref*{annex:p12-cdcl-uip-lookahead-sound}} & object & \hyperref[annex:p12-cdcl-uip-lookahead-sound]{\ref*{annex:p12-cdcl-uip-lookahead-sound}} & \thepage & Negation retained until discharged. \\
\phantomsection\label{entry:symbol-registry-05846}\(\neg p\) & \texttt{\detokenize{\neg p}} & sign-flipping complement of a signed propositional literal & \hyperref[sec:p12-cdcl-uip-and-watched-literals]{\ref*{sec:p12-cdcl-uip-and-watched-literals}} & meta & \hyperref[def:p12-cdcl-uip-state]{\ref*{def:p12-cdcl-uip-state}} & \thepage & Explicit literal-complement role; complement of negative x is x rather than a new double-negation formula. \\
\phantomsection\label{entry:symbol-registry-05999}\(\varphi\) & \texttt{\detokenize{\varphi}} & original GRASP CNF clause database & \hyperref[sec:p12-grasp-and-minisat-originals]{\ref*{sec:p12-grasp-and-minisat-originals}} & object & \hyperref[def:p12-grasp-native-state]{\ref*{def:p12-grasp-native-state}} & \thepage & Publication notation retained; not a Boolean-constraint class instance. \\
\phantomsection\label{entry:symbol-registry-06000}\(A\) & \texttt{\detokenize{A}} & partial assignment as variable-value pairs & \hyperref[sec:p12-grasp-and-minisat-originals]{\ref*{sec:p12-grasp-and-minisat-originals}} & meta & \hyperref[def:p12-grasp-native-state]{\ref*{def:p12-grasp-native-state}} & \thepage & Can include fixed initial query assumptions. \\
\phantomsection\label{entry:symbol-registry-06001}\(\nu(x)\) & \texttt{\detokenize{\nu(x)}} & assigned Boolean value of a variable & \hyperref[sec:p12-grasp-and-minisat-originals]{\ref*{sec:p12-grasp-and-minisat-originals}} & meta & \hyperref[def:p12-grasp-native-state]{\ref*{def:p12-grasp-native-state}} & \thepage & The GRASP source's value notation. \\
\phantomsection\label{entry:symbol-registry-06002}\(\delta(x)\) & \texttt{\detokenize{\delta(x)}} & level annotation of an assigned variable & \hyperref[sec:p12-grasp-and-minisat-originals]{\ref*{sec:p12-grasp-and-minisat-originals}} & meta & \hyperref[def:p12-grasp-native-state]{\ref*{def:p12-grasp-native-state}} & \thepage & Source reason-max Eq.1 and current-frame Fig.3/FDA convention remain separate. \\
\phantomsection\label{entry:symbol-registry-06003}\(X\) & \texttt{\detokenize{X}} & GRASP unassigned value marker denoting both Boolean possibilities & \hyperref[def:p12-grasp-native-state]{\ref*{def:p12-grasp-native-state}} & object & \hyperref[def:p12-grasp-native-state]{\ref*{def:p12-grasp-native-state}} & \thepage & Published X={0,1}; genuinely distinct from a finite variable universe in the other CDCL presentation. \\
\phantomsection\label{entry:symbol-registry-06004}\(x\) & \texttt{\detokenize{x}} & Boolean variable & \hyperref[sec:p12-grasp-and-minisat-originals]{\ref*{sec:p12-grasp-and-minisat-originals}} & object & \hyperref[def:p12-grasp-native-state]{\ref*{def:p12-grasp-native-state}} & \thepage & Indexed x\_1 through x\_13 preserve the original graph's names. \\
\phantomsection\label{entry:symbol-registry-06005}\(x^0\) & \texttt{\detokenize{x^0}} & positive source literal opposite assignment x=0 & \hyperref[sec:p12-grasp-and-minisat-originals]{\ref*{sec:p12-grasp-and-minisat-originals}} & object & \hyperref[def:p12-grasp-native-state]{\ref*{def:p12-grasp-native-state}} & \thepage & Published exponent is a value index not arithmetic exponentiation. \\
\phantomsection\label{entry:symbol-registry-06006}\(x^1\) & \texttt{\detokenize{x^1}} & negative source literal opposite assignment x=1 & \hyperref[sec:p12-grasp-and-minisat-originals]{\ref*{sec:p12-grasp-and-minisat-originals}} & object & \hyperref[def:p12-grasp-native-state]{\ref*{def:p12-grasp-native-state}} & \thepage & x\^{}nu(x) is false under that assignment. \\
\phantomsection\label{entry:symbol-registry-06007}\(+\) & \texttt{\detokenize{+}} & source clause disjunction & \hyperref[sec:p12-grasp-and-minisat-originals]{\ref*{sec:p12-grasp-and-minisat-originals}} & infix & \hyperref[def:p12-grasp-native-state]{\ref*{def:p12-grasp-native-state}} & \thepage & Published glyph retained; not numerical addition. \\
\phantomsection\label{entry:symbol-registry-06008}\(\omega\) & \texttt{\detokenize{\omega}} & selected clause or unit reason & \hyperref[sec:p12-grasp-and-minisat-originals]{\ref*{sec:p12-grasp-and-minisat-originals}} & object & \hyperref[def:p12-grasp-native-state]{\ref*{def:p12-grasp-native-state}} & \thepage & Omega\_i indexes the original nine-clause example. \\
\phantomsection\label{entry:symbol-registry-06009}\(A^\omega(x)\) & \texttt{\detokenize{A^\omega(x)}} & antecedent assignments falsifying other literals of a unit reason & \hyperref[sec:p12-grasp-and-minisat-originals]{\ref*{sec:p12-grasp-and-minisat-originals}} & meta & \hyperref[def:p12-grasp-native-state]{\ref*{def:p12-grasp-native-state}} & \thepage & Assignments rather than a clause object. \\
\phantomsection\label{entry:symbol-registry-06010}\(I\) & \texttt{\detokenize{I}} & finite implication graph of assignments and conflicts & \hyperref[sec:p12-grasp-and-minisat-originals]{\ref*{sec:p12-grasp-and-minisat-originals}} & meta & \hyperref[def:p12-grasp-native-state]{\ref*{def:p12-grasp-native-state}} & \thepage & Edges point from earlier antecedents to consequences. \\
\phantomsection\label{entry:symbol-registry-06011}\(\kappa\) & \texttt{\detokenize{\kappa}} & GRASP conflict vertex & \hyperref[sec:p12-grasp-and-minisat-originals]{\ref*{sec:p12-grasp-and-minisat-originals}} & object & \hyperref[def:p12-grasp-native-state]{\ref*{def:p12-grasp-native-state}} & \thepage & Not a Boolean variable. \\
\phantomsection\label{entry:symbol-registry-06012}\(d\) & \texttt{\detokenize{d}} & current recursive search frame level & \hyperref[sec:p12-grasp-and-minisat-originals]{\ref*{sec:p12-grasp-and-minisat-originals}} & meta & \hyperref[def:p12-grasp-native-state]{\ref*{def:p12-grasp-native-state}} & \thepage & Root convention explicitly completed rather than silently aliasing the printed wrapper. \\
\phantomsection\label{entry:symbol-registry-06013}\(\Lambda(x)\) & \texttt{\detokenize{\Lambda(x)}} & older-level reason assignments & \hyperref[sec:p12-grasp-and-minisat-originals]{\ref*{sec:p12-grasp-and-minisat-originals}} & meta & \hyperref[def:p12-grasp-native-state]{\ref*{def:p12-grasp-native-state}} & \thepage & Retained as leaves in current-level cause expansion. \\
\phantomsection\label{entry:symbol-registry-06014}\(\Sigma(x)\) & \texttt{\detokenize{\Sigma(x)}} & current-level reason assignments & \hyperref[sec:p12-grasp-and-minisat-originals]{\ref*{sec:p12-grasp-and-minisat-originals}} & meta & \hyperref[def:p12-grasp-native-state]{\ref*{def:p12-grasp-native-state}} & \thepage & Expanded backward in the acyclic graph. \\
\phantomsection\label{entry:symbol-registry-06015}\(\operatorname{causes\_of}(x)\) & \texttt{\detokenize{\operatorname{causes\_of}(x)}} & standard GRASP backward cause-assignment construction & \hyperref[sec:p12-grasp-and-minisat-originals]{\ref*{sec:p12-grasp-and-minisat-originals}} & meta & \hyperref[def:p12-grasp-causes]{\ref*{def:p12-grasp-causes}} & \thepage & Source Eq.3 reconstruction with explicit conflict/root boundary. \\
\phantomsection\label{entry:symbol-registry-06016}\(A^{\omega_C}(\kappa)\) & \texttt{\detokenize{A^{\omega_C}(\kappa)}} & sufficient assignment support of a conflict & \hyperref[sec:p12-grasp-and-minisat-originals]{\ref*{sec:p12-grasp-and-minisat-originals}} & meta & \hyperref[def:p12-grasp-causes]{\ref*{def:p12-grasp-causes}} & \thepage & Includes expanded FDA reasons rather than treating them as decisions. \\
\phantomsection\label{entry:symbol-registry-06017}\(\omega_C(\kappa)\) & \texttt{\detokenize{\omega_C(\kappa)}} & conflict-induced implicate negating the cause assignment & \hyperref[sec:p12-grasp-and-minisat-originals]{\ref*{sec:p12-grasp-and-minisat-originals}} & object & \hyperref[def:p12-grasp-causes]{\ref*{def:p12-grasp-causes}} & \thepage & Input consequence; not a RAT redundancy clause. \\
\phantomsection\label{entry:symbol-registry-06018}\(\sum\) & \texttt{\detokenize{\sum}} & finite disjunction in the GRASP source clause formula & \hyperref[sec:p12-grasp-and-minisat-originals]{\ref*{sec:p12-grasp-and-minisat-originals}} & meta & \hyperref[def:p12-grasp-causes]{\ref*{def:p12-grasp-causes}} & \thepage & Not a numerical sum; empty disjunction is the empty clause. \\
\phantomsection\label{entry:symbol-registry-06019}\(\beta\) & \texttt{\detokenize{\beta}} & highest-cause-level CDB target & \hyperref[sec:p12-grasp-and-minisat-originals]{\ref*{sec:p12-grasp-and-minisat-originals}} & meta & \hyperref[def:p12-grasp-standard-engines]{\ref*{def:p12-grasp-standard-engines}} & \thepage & Distinct from the second-highest assertion level in later first-UIP backjumping. \\
\phantomsection\label{entry:symbol-registry-06020}\(A_0\) & \texttt{\detokenize{A_0}} & fixed initial query assignment & \hyperref[sec:p12-grasp-and-minisat-originals]{\ref*{sec:p12-grasp-and-minisat-originals}} & meta & \hyperref[thm:p12-grasp-standard-search]{\ref*{thm:p12-grasp-standard-search}} & \thepage & Kept at root; success and failure are relative to extensions of this assignment. \\
\phantomsection\label{entry:symbol-registry-06021}\(u_i\) & \texttt{\detokenize{u_i}} & UIP vertex indexed in forward decision-to-conflict order & \hyperref[sec:p12-grasp-and-minisat-originals]{\ref*{sec:p12-grasp-and-minisat-originals}} & object & \hyperref[def:p12-grasp-uip-links]{\ref*{def:p12-grasp-uip-links}} & \thepage & Corrected endpoint order is explicitly stated and proved. \\
\phantomsection\label{entry:symbol-registry-06022}\(\operatorname{causes\_of}(z,u_i)\) & \texttt{\detokenize{\operatorname{causes\_of}(z,}\allowbreak\detokenize{u_i)}} & UIP-relative backward support of downstream z with stop u\_i & \hyperref[sec:p12-grasp-and-minisat-originals]{\ref*{sec:p12-grasp-and-minisat-originals}} & meta & \hyperref[def:p12-grasp-uip-links]{\ref*{def:p12-grasp-uip-links}} & \thepage & Two-argument source operation; endpoint roles are not interchanged. \\
\phantomsection\label{entry:symbol-registry-06023}\(z\) & \texttt{\detokenize{z}} & downstream graph vertex in a UIP-relative support & P12 GRASP UIP construction & object & \hyperref[def:p12-grasp-uip-links]{\ref*{def:p12-grasp-uip-links}} & \thepage & Can be an assignment or the conflict node as explicitly typed; separate from Boolean variable z. Root actual declaration-region reconciliation: Downstream graph-vertex role belongs to the actual UIP method and reciprocal proof, not every Boolean z. \\
\phantomsection\label{entry:symbol-registry-06024}\(P(\varphi)\) & \texttt{\detokenize{P(\varphi)}} & complete product of all prime implicate clauses & \hyperref[sec:p12-grasp-and-minisat-originals]{\ref*{sec:p12-grasp-and-minisat-originals}} & meta & \hyperref[def:p12-grasp-complete-product]{\ref*{def:p12-grasp-complete-product}} & \thepage & Empty clause is the unique prime implicate of false; product of true is empty. \\
\phantomsection\label{entry:symbol-registry-06025}\(\psi\) & \texttt{\detokenize{\psi}} & clause containing a prime subclause & \hyperref[def:p12-grasp-complete-product]{\ref*{def:p12-grasp-complete-product}} & object & \hyperref[def:p12-grasp-complete-product]{\ref*{def:p12-grasp-complete-product}} & \thepage & Correct implication goes from subclause to containing clause. \\
\phantomsection\label{entry:symbol-registry-06026}\(\bot\) & \texttt{\detokenize{\bot}} & MiniSat undefined lifted-value marker & \hyperref[def:p12-minisat-original-interface]{\ref*{def:p12-minisat-original-interface}} & object & \hyperref[def:p12-minisat-original-interface]{\ref*{def:p12-minisat-original-interface}} & \thepage & Published role retained; not Boolean false or a refutation outcome. \\
\phantomsection\label{entry:symbol-registry-06027}\(\mathrm{True}_{\bot}\) & \texttt{\detokenize{\mathrm{True}_{\bot}}} & lifted defined Boolean true & \hyperref[def:p12-minisat-original-interface]{\ref*{def:p12-minisat-original-interface}} & object & \hyperref[def:p12-minisat-original-interface]{\ref*{def:p12-minisat-original-interface}} & \thepage & Distinct from undefined. \\
\phantomsection\label{entry:symbol-registry-06028}\(\mathrm{False}_{\bot}\) & \texttt{\detokenize{\mathrm{False}_{\bot}}} & lifted defined Boolean false & \hyperref[def:p12-minisat-original-interface]{\ref*{def:p12-minisat-original-interface}} & object & \hyperref[def:p12-minisat-original-interface]{\ref*{def:p12-minisat-original-interface}} & \thepage & Search false is query UNSAT; undefined is a budget exit. \\
\phantomsection\label{entry:symbol-registry-06029}\(\bot_{lit}\) & \texttt{\detokenize{\bot_{lit}}} & undefined literal sentinel requesting a conflict reason & \hyperref[sec:p12-grasp-and-minisat-originals]{\ref*{sec:p12-grasp-and-minisat-originals}} & object & \hyperref[def:p12-minisat-original-callbacks]{\ref*{def:p12-minisat-original-callbacks}} & \thepage & Not a constraint literal. \\
\phantomsection\label{entry:symbol-registry-06030}\(K\) & \texttt{\detokenize{K}} & permanent Boolean predicate constraint & \hyperref[sec:p12-grasp-and-minisat-originals]{\ref*{sec:p12-grasp-and-minisat-originals}} & object & \hyperref[def:p12-minisat-original-reasons]{\ref*{def:p12-minisat-original-reasons}} & \thepage & Need not have a clause syntax. \\
\phantomsection\label{entry:symbol-registry-06031}\(R\) & \texttt{\detokenize{R}} & finite conjunction-support set of true reason assignments & \hyperref[sec:p12-grasp-and-minisat-originals]{\ref*{sec:p12-grasp-and-minisat-originals}} & meta & \hyperref[def:p12-minisat-original-reasons]{\ref*{def:p12-minisat-original-reasons}} & \thepage & Complements become a learned clause; not already a disjunction. \\
\phantomsection\label{entry:symbol-registry-06032}\(p\) & \texttt{\detokenize{p}} & propagated signed assignment literal & \hyperref[sec:p12-grasp-and-minisat-originals]{\ref*{sec:p12-grasp-and-minisat-originals}} & object & \hyperref[def:p12-minisat-original-callbacks]{\ref*{def:p12-minisat-original-callbacks}} & \thepage & Complement operation flips literal sign. \\
\phantomsection\label{entry:symbol-registry-06033}\(F\) & \texttt{\detokenize{F}} & permanent Boolean constraint conjunction & \hyperref[sec:p12-grasp-and-minisat-originals]{\ref*{sec:p12-grasp-and-minisat-originals}} & object & \hyperref[thm:p12-minisat-original-query]{\ref*{thm:p12-minisat-original-query}} & \thepage & Distinct from GRASP phi clause database. \\
\phantomsection\label{entry:symbol-registry-06034}\(G\) & \texttt{\detokenize{G}} & conjunction of permanent root facts entailed by F & \hyperref[sec:p12-grasp-and-minisat-originals]{\ref*{sec:p12-grasp-and-minisat-originals}} & object & \hyperref[thm:p12-minisat-original-query]{\ref*{thm:p12-minisat-original-query}} & \thepage & Remains part of represented semantics after simplification. \\
\phantomsection\label{entry:symbol-registry-06035}\(F'\) & \texttt{\detokenize{F'}} & remaining relatively simplified constraint conjunction & \hyperref[sec:p12-grasp-and-minisat-originals]{\ref*{sec:p12-grasp-and-minisat-originals}} & object & \hyperref[thm:p12-minisat-original-query]{\ref*{thm:p12-minisat-original-query}} & \thepage & G and F' together preserve F; F' alone need not. \\
\phantomsection\label{entry:symbol-registry-06036}\(H\) & \texttt{\detokenize{H}} & temporary query-assumption conjunction & \hyperref[sec:p12-grasp-and-minisat-originals]{\ref*{sec:p12-grasp-and-minisat-originals}} & object & \hyperref[thm:p12-minisat-original-query]{\ref*{thm:p12-minisat-original-query}} & \thepage & Not necessarily entailed by permanent constraints. \\
\phantomsection\label{entry:symbol-registry-06037}\(B\) & \texttt{\detokenize{B}} & consistent partial assignment in the prime-product proof & \hyperref[annex:p12-grasp-complete-product]{\ref*{annex:p12-grasp-complete-product}} & meta & \hyperref[annex:p12-grasp-complete-product]{\ref*{annex:p12-grasp-complete-product}} & \thepage & Its blocking clause contains opposite literals. \\
\phantomsection\label{entry:symbol-registry-06038}\(v\) & \texttt{\detokenize{v}} & chosen Boolean value in the prime-product proof & \hyperref[annex:p12-grasp-complete-product]{\ref*{annex:p12-grasp-complete-product}} & object & \hyperref[annex:p12-grasp-complete-product]{\ref*{annex:p12-grasp-complete-product}} & \thepage & Zero or one; not a valuation function. \\
\phantomsection\label{entry:symbol-registry-06039}\ensuremath{\texttt{var}} & \texttt{\detokenize{\texttt{var}}} & native MiniSat variable-index type & \hyperref[def:p12-minisat-original-interface]{\ref*{def:p12-minisat-original-interface}} & meta & \hyperref[def:p12-minisat-original-interface]{\ref*{def:p12-minisat-original-interface}} & \thepage & Integer alias with undefined variable sentinel. \\
\phantomsection\label{entry:symbol-registry-06040}\ensuremath{\texttt{lit}} & \texttt{\detokenize{\texttt{lit}}} & native MiniSat signed-literal type & \hyperref[def:p12-minisat-original-interface]{\ref*{def:p12-minisat-original-interface}} & meta & \hyperref[def:p12-minisat-original-interface]{\ref*{def:p12-minisat-original-interface}} & \thepage & Has an undefined sentinel distinct from real literals. \\
\phantomsection\label{entry:symbol-registry-06041}\ensuremath{\texttt{lbool}} & \texttt{\detokenize{\texttt{lbool}}} & native MiniSat lifted Boolean type & \hyperref[def:p12-minisat-original-interface]{\ref*{def:p12-minisat-original-interface}} & meta & \hyperref[def:p12-minisat-original-interface]{\ref*{def:p12-minisat-original-interface}} & \thepage & False True and undefined remain separate outcomes. \\
\phantomsection\label{entry:symbol-registry-06042}\ensuremath{\texttt{Vec}} & \texttt{\detokenize{\texttt{Vec}}} & native extensible vector type constructor & \hyperref[def:p12-minisat-original-interface]{\ref*{def:p12-minisat-original-interface}} & meta & \hyperref[def:p12-minisat-original-interface]{\ref*{def:p12-minisat-original-interface}} & \thepage & Transfer ownership semantics are an explicit implementation dependency. \\
\phantomsection\label{entry:symbol-registry-06043}\ensuremath{\texttt{Queue}} & \texttt{\detokenize{\texttt{Queue}}} & native event-queue type constructor & \hyperref[def:p12-minisat-original-interface]{\ref*{def:p12-minisat-original-interface}} & meta & \hyperref[def:p12-minisat-original-interface]{\ref*{def:p12-minisat-original-interface}} & \thepage & Distinct from clause databases and reason vectors. \\
\phantomsection\label{entry:symbol-registry-06044}\(y\) & \texttt{\detokenize{y}} & Boolean variable in original source clause examples & \hyperref[sec:p12-grasp-and-minisat-originals]{\ref*{sec:p12-grasp-and-minisat-originals}} & object & \hyperref[def:p12-grasp-native-state]{\ref*{def:p12-grasp-native-state}} & \thepage & Literal values determine antecedent assignments. \\
\phantomsection\label{entry:symbol-registry-06045}\(z\) & \texttt{\detokenize{z}} & Boolean variable in original source clause examples & \hyperref[sec:p12-grasp-and-minisat-originals]{\ref*{sec:p12-grasp-and-minisat-originals}} & object & \hyperref[def:p12-grasp-native-state]{\ref*{def:p12-grasp-native-state}} & \thepage & Separate locally declared graph-placeholder role in the UIP construction. \\
\phantomsection\label{entry:symbol-registry-06046}\(v\) & \texttt{\detokenize{v}} & assigned Boolean value in the GRASP FDA condition & \hyperref[def:p12-grasp-fda-cdb]{\ref*{def:p12-grasp-fda-cdb}} & object & \hyperref[def:p12-grasp-fda-cdb]{\ref*{def:p12-grasp-fda-cdb}} & \thepage & Zero or one; not a variable in this local condition. \\
\phantomsection\label{entry:symbol-registry-06047}\(u\) & \texttt{\detokenize{u}} & Boolean variable in the original MiniSat reason example & \hyperref[def:p12-minisat-original-reason-example]{\ref*{def:p12-minisat-original-reason-example}} & object & \hyperref[def:p12-minisat-original-reason-example]{\ref*{def:p12-minisat-original-reason-example}} & \thepage & Explicit source-example type declaration. \\
\phantomsection\label{entry:symbol-registry-06048}\(v\) & \texttt{\detokenize{v}} & Boolean variable in the original MiniSat reason example & \hyperref[def:p12-minisat-original-reason-example]{\ref*{def:p12-minisat-original-reason-example}} & object & \hyperref[def:p12-minisat-original-reason-example]{\ref*{def:p12-minisat-original-reason-example}} & \thepage & Not the earlier FDA assigned-value parameter. \\
\phantomsection\label{entry:symbol-registry-06049}\(\varnothing\) & \texttt{\detokenize{\varnothing}} & empty antecedent-assignment set & \hyperref[sec:p12-grasp-and-minisat-originals]{\ref*{sec:p12-grasp-and-minisat-originals}} & object & \hyperref[def:p12-grasp-causes]{\ref*{def:p12-grasp-causes}} & \thepage & Conflict with empty causes has a separately explained empty-clause result. \\
\phantomsection\label{entry:symbol-registry-06050}\(k\) & \texttt{\detokenize{k}} & number of UIP vertices in forward list & P12 GRASP UIP construction & meta & \hyperref[def:p12-grasp-uip-links]{\ref*{def:p12-grasp-uip-links}} & \thepage & Separate from the fixed clause-width threshold in the red/green policy. Root actual declaration-region reconciliation: UIP-count declaration and its reciprocal argument share this exact source construction. \\
\phantomsection\label{entry:symbol-registry-06051}\(k\) & \texttt{\detokenize{k}} & fixed permanent-green clause-width threshold & \hyperref[def:p12-grasp-retention-policy]{\ref*{def:p12-grasp-retention-policy}} & meta & \hyperref[def:p12-grasp-retention-policy]{\ref*{def:p12-grasp-retention-policy}} & \thepage & Source paragraph declares this integer parameter separately; root binds the UIP-list local override. Root actual declaration-region reconciliation: Actual new retention-policy method encloses the source green-width paragraph; width is not the UIP-count parameter. \\
\phantomsection\label{entry:symbol-registry-06243}\(\phi\) & \texttt{\detokenize{\phi}} & closed PCNF quantified Boolean input & \hyperref[sec:p12-q-resolution-and-quantified-certificates]{\ref*{sec:p12-q-resolution-and-quantified-certificates}} & object & \hyperref[def:p12-qres-pcnf]{\ref*{def:p12-qres-pcnf}} & \thepage & Prefix and matrix are both part of the input. \\
\phantomsection\label{entry:symbol-registry-06244}\(\Pi\) & \texttt{\detokenize{\Pi}} & fixed finite quantifier prefix & \hyperref[sec:p12-q-resolution-and-quantified-certificates]{\ref*{sec:p12-q-resolution-and-quantified-certificates}} & meta & \hyperref[def:p12-qres-pcnf]{\ref*{def:p12-qres-pcnf}} & \thepage & Immutable during proof checking; native source glyph retained. \\
\phantomsection\label{entry:symbol-registry-06245}\(\psi\) & \texttt{\detokenize{\psi}} & finite CNF matrix as clause set & \hyperref[sec:p12-q-resolution-and-quantified-certificates]{\ref*{sec:p12-q-resolution-and-quantified-certificates}} & object & \hyperref[def:p12-qres-pcnf]{\ref*{def:p12-qres-pcnf}} & \thepage & The empty matrix is true. \\
\phantomsection\label{entry:symbol-registry-06246}\(Q_i\) & \texttt{\detokenize{Q_i}} & quantifier tag of block i & \hyperref[sec:p12-q-resolution-and-quantified-certificates]{\ref*{sec:p12-q-resolution-and-quantified-certificates}} & meta & \hyperref[def:p12-qres-pcnf]{\ref*{def:p12-qres-pcnf}} & \thepage & Existential or universal; no free matrix variables. \\
\phantomsection\label{entry:symbol-registry-06247}\(X_i\) & \texttt{\detokenize{X_i}} & finite disjoint variable block & \hyperref[sec:p12-q-resolution-and-quantified-certificates]{\ref*{sec:p12-q-resolution-and-quantified-certificates}} & object & \hyperref[def:p12-qres-pcnf]{\ref*{def:p12-qres-pcnf}} & \thepage & Order inside an equal-quantifier block may be chosen for construction. \\
\phantomsection\label{entry:symbol-registry-06248}\(x\) & \texttt{\detokenize{x}} & Boolean variable & \hyperref[sec:p12-q-resolution-and-quantified-certificates]{\ref*{sec:p12-q-resolution-and-quantified-certificates}} & object & \hyperref[def:p12-qres-pcnf]{\ref*{def:p12-qres-pcnf}} & \thepage & Indexed x\_1 x\_2 x\_3 x\_4 preserve source example names. \\
\phantomsection\label{entry:symbol-registry-06249}\(l\) & \texttt{\detokenize{l}} & signed Boolean literal & \hyperref[sec:p12-q-resolution-and-quantified-certificates]{\ref*{sec:p12-q-resolution-and-quantified-certificates}} & object & \hyperref[def:p12-qres-pcnf]{\ref*{def:p12-qres-pcnf}} & \thepage & Complement flips signs without new double-negation formulas. \\
\phantomsection\label{entry:symbol-registry-06250}\(\bar l\) & \texttt{\detokenize{\bar l}} & complementary signed literal & \hyperref[sec:p12-q-resolution-and-quantified-certificates]{\ref*{sec:p12-q-resolution-and-quantified-certificates}} & meta & \hyperref[def:p12-qres-pcnf]{\ref*{def:p12-qres-pcnf}} & \thepage & An involution; source complement notation retained. \\
\phantomsection\label{entry:symbol-registry-06251}\(\varnothing\) & \texttt{\detokenize{\varnothing}} & empty literal-set clause & \hyperref[sec:p12-q-resolution-and-quantified-certificates]{\ref*{sec:p12-q-resolution-and-quantified-certificates}} & object & \hyperref[def:p12-qres-pcnf]{\ref*{def:p12-qres-pcnf}} & \thepage & False under every complete valuation; not an empty CNF database. \\
\phantomsection\label{entry:symbol-registry-06252}\(\bot\) & \texttt{\detokenize{\bot}} & empty-clause false constant in this PCNF source & \hyperref[sec:p12-q-resolution-and-quantified-certificates]{\ref*{sec:p12-q-resolution-and-quantified-certificates}} & object & \hyperref[def:p12-qres-pcnf]{\ref*{def:p12-qres-pcnf}} & \thepage & Not the undefined lifted value used by MiniSat in a separate chapter. \\
\phantomsection\label{entry:symbol-registry-06253}\(\operatorname{lv}(x)\) & \texttt{\detokenize{\operatorname{lv}(x)}} & prefix block level of a Boolean variable & \hyperref[sec:p12-q-resolution-and-quantified-certificates]{\ref*{sec:p12-q-resolution-and-quantified-certificates}} & meta & \hyperref[def:p12-qres-pcnf]{\ref*{def:p12-qres-pcnf}} & \thepage & Extends to signed literals through their underlying variable. \\
\phantomsection\label{entry:symbol-registry-06254}\(\operatorname{quant}(\Pi,x)\) & \texttt{\detokenize{\operatorname{quant}(\Pi,}\allowbreak\detokenize{x)}} & quantifier tag of a prefix variable or literal & \hyperref[sec:p12-q-resolution-and-quantified-certificates]{\ref*{sec:p12-q-resolution-and-quantified-certificates}} & meta & \hyperref[def:p12-qres-pcnf]{\ref*{def:p12-qres-pcnf}} & \thepage & Arguments may be variables or signed literals at the same bound position. \\
\phantomsection\label{entry:symbol-registry-06255}\(\leq_\Pi\) & \texttt{\detokenize{\leq_\Pi}} & comparison of prefix block levels & \hyperref[sec:p12-q-resolution-and-quantified-certificates]{\ref*{sec:p12-q-resolution-and-quantified-certificates}} & relation & \hyperref[def:p12-qres-pcnf]{\ref*{def:p12-qres-pcnf}} & \thepage & Opposite quantifier types cannot share a block. \\
\phantomsection\label{entry:symbol-registry-06256}\(\operatorname{var}(l)\) & \texttt{\detokenize{\operatorname{var}(l)}} & underlying variable of a signed literal & \hyperref[sec:p12-q-resolution-and-quantified-certificates]{\ref*{sec:p12-q-resolution-and-quantified-certificates}} & meta & \hyperref[def:p12-qres-pcnf]{\ref*{def:p12-qres-pcnf}} & \thepage & Native preliminaries argument type. \\
\phantomsection\label{entry:symbol-registry-06257}\(\operatorname{var}(\psi)\) & \texttt{\detokenize{\operatorname{var}(\psi)}} & set of occurring matrix variables & \hyperref[sec:p12-q-resolution-and-quantified-certificates]{\ref*{sec:p12-q-resolution-and-quantified-certificates}} & meta & \hyperref[def:p12-qres-pcnf]{\ref*{def:p12-qres-pcnf}} & \thepage & Typed formula argument distinct from the single-literal argument. \\
\phantomsection\label{entry:symbol-registry-06258}\(\operatorname{vars}(\Pi)\) & \texttt{\detokenize{\operatorname{vars}(\Pi)}} & union of prefix variable blocks in the native rule guards & \hyperref[sec:p12-q-resolution-and-quantified-certificates]{\ref*{sec:p12-q-resolution-and-quantified-certificates}} & meta & \hyperref[def:p12-qres-pcnf]{\ref*{def:p12-qres-pcnf}} & \thepage & Source Fig.2.1 spelling retained rather than silently replacing it by var. \\
\phantomsection\label{entry:symbol-registry-06259}\(e\) & \texttt{\detokenize{e}} & existential Boolean variable & \hyperref[sec:p12-q-resolution-and-quantified-certificates]{\ref*{sec:p12-q-resolution-and-quantified-certificates}} & object & \hyperref[def:p12-qres-pcnf]{\ref*{def:p12-qres-pcnf}} & \thepage & May depend only on earlier universal variables. \\
\phantomsection\label{entry:symbol-registry-06260}\(u\) & \texttt{\detokenize{u}} & universal Boolean variable & \hyperref[sec:p12-q-resolution-and-quantified-certificates]{\ref*{sec:p12-q-resolution-and-quantified-certificates}} & object & \hyperref[def:p12-qres-pcnf]{\ref*{def:p12-qres-pcnf}} & \thepage & Its counterstrategy may depend only on earlier existentials. \\
\phantomsection\label{entry:symbol-registry-06261}\(f_e\) & \texttt{\detokenize{f_e}} & legal existential Skolem strategy function & \hyperref[sec:p12-q-resolution-and-quantified-certificates]{\ref*{sec:p12-q-resolution-and-quantified-certificates}} & meta & \hyperref[def:p12-qres-strategies]{\ref*{def:p12-qres-strategies}} & \thepage & Declared preceding-universal argument list; empty list is constant. \\
\phantomsection\label{entry:symbol-registry-06262}\(g_u\) & \texttt{\detokenize{g_u}} & legal universal Herbrand strategy function & \hyperref[sec:p12-q-resolution-and-quantified-certificates]{\ref*{sec:p12-q-resolution-and-quantified-certificates}} & meta & \hyperref[def:p12-qres-strategies]{\ref*{def:p12-qres-strategies}} & \thepage & Declared preceding-existential argument list. \\
\phantomsection\label{entry:symbol-registry-06263}\(S\) & \texttt{\detokenize{S}} & family of legal existential strategy functions & \hyperref[sec:p12-q-resolution-and-quantified-certificates]{\ref*{sec:p12-q-resolution-and-quantified-certificates}} & meta & \hyperref[def:p12-qres-strategy-consequence]{\ref*{def:p12-qres-strategy-consequence}} & \thepage & Clause satisfaction is required over every universal assignment. \\
\phantomsection\label{entry:symbol-registry-06264}\(C\) & \texttt{\detokenize{C}} & finite Boolean clause or rule conclusion & \hyperref[sec:p12-q-resolution-and-quantified-certificates]{\ref*{sec:p12-q-resolution-and-quantified-certificates}} & object & \hyperref[def:p12-qres-strategy-consequence]{\ref*{def:p12-qres-strategy-consequence}} & \thepage & All accepted proof lines are non-tautological. \\
\phantomsection\label{entry:symbol-registry-06265}\(C_1\) & \texttt{\detokenize{C_1}} & clause remainder in a resolution parent & \hyperref[sheet:p12-q-resolution]{\ref*{sheet:p12-q-resolution}} & object & \hyperref[sheet:p12-q-resolution]{\ref*{sheet:p12-q-resolution}} & \thepage & Printed decomposition guards retained; completeness chooses canonical removed-pivot remainder. \\
\phantomsection\label{entry:symbol-registry-06266}\(C_2\) & \texttt{\detokenize{C_2}} & clause remainder in the other resolution parent & \hyperref[sheet:p12-q-resolution]{\ref*{sheet:p12-q-resolution}} & object & \hyperref[sheet:p12-q-resolution]{\ref*{sheet:p12-q-resolution}} & \thepage & Not a separate quantified matrix. \\
\phantomsection\label{entry:symbol-registry-06267}\(p\) & \texttt{\detokenize{p}} & existential pivot literal in common Q-Res & \hyperref[sheet:p12-q-resolution]{\ref*{sheet:p12-q-resolution}} & object & \hyperref[sheet:p12-q-resolution]{\ref*{sheet:p12-q-resolution}} & \thepage & Universal-pivot variants have separate rules and are not credited here. \\
\phantomsection\label{entry:symbol-registry-06268}\(C_i\) & \texttt{\detokenize{C_i}} & derived clause in the native blocked-reduction example & \hyperref[def:p12-qres-native-derivation]{\ref*{def:p12-qres-native-derivation}} & object & \hyperref[def:p12-qres-native-derivation]{\ref*{def:p12-qres-native-derivation}} & \thepage & Indices zero one two denote the actual source derivation. \\
\phantomsection\label{entry:symbol-registry-06269}\(y_5\) & \texttt{\detokenize{y_5}} & universal variable in the source example & \hyperref[def:p12-qres-native-derivation]{\ref*{def:p12-qres-native-derivation}} & object & \hyperref[def:p12-qres-native-derivation]{\ref*{def:p12-qres-native-derivation}} & \thepage & Initially blocked by later x\_2 and then reducible. \\
\phantomsection\label{entry:symbol-registry-06270}\(A\) & \texttt{\detokenize{A}} & positive-pivot parent remainder in innermost elimination & \hyperref[sec:p12-q-resolution-and-quantified-certificates]{\ref*{sec:p12-q-resolution-and-quantified-certificates}} & object & \hyperref[def:p12-qres-elimination]{\ref*{def:p12-qres-elimination}} & \thepage & Canonical remainder contains neither pivot polarity. \\
\phantomsection\label{entry:symbol-registry-06271}\(B\) & \texttt{\detokenize{B}} & negative-pivot parent remainder in innermost elimination & \hyperref[sec:p12-q-resolution-and-quantified-certificates]{\ref*{sec:p12-q-resolution-and-quantified-certificates}} & object & \hyperref[def:p12-qres-elimination]{\ref*{def:p12-qres-elimination}} & \thepage & Not the partial assignment called B in older chapters. \\
\phantomsection\label{entry:symbol-registry-06272}\(a\) & \texttt{\detokenize{a}} & outer existential variable in the alternating certificate & \hyperref[def:p12-qres-alternating-certificate]{\ref*{def:p12-qres-alternating-certificate}} & object & \hyperref[def:p12-qres-alternating-certificate]{\ref*{def:p12-qres-alternating-certificate}} & \thepage & Cannot depend on later universal u. \\
\phantomsection\label{entry:symbol-registry-06273}\(b\) & \texttt{\detokenize{b}} & inner existential variable in the alternating certificate & \hyperref[def:p12-qres-alternating-certificate]{\ref*{def:p12-qres-alternating-certificate}} & object & \hyperref[def:p12-qres-alternating-certificate]{\ref*{def:p12-qres-alternating-certificate}} & \thepage & Its elimination permits subsequent legal universal reductions. \\
\phantomsection\label{entry:symbol-registry-06274}\(y\) & \texttt{\detokenize{y}} & existential variable of the illegal-rule counterexample & \hyperref[def:p12-qres-illegal-reduction]{\ref*{def:p12-qres-illegal-reduction}} & object & \hyperref[def:p12-qres-illegal-reduction]{\ref*{def:p12-qres-illegal-reduction}} & \thepage & Legal function y=x with prefix forall x exists y only. \\
\phantomsection\label{entry:symbol-registry-06275}\(\operatorname{ur}_{\Pi}(C)\) & \texttt{\detokenize{\operatorname{ur}_{\Pi}(C)}} & full trailing-universal reduction of a clause & \hyperref[sec:p12-q-resolution-and-quantified-certificates]{\ref*{sec:p12-q-resolution-and-quantified-certificates}} & meta & \hyperref[def:p12-qres-full-reduction]{\ref*{def:p12-qres-full-reduction}} & \thepage & 2013 full-reduction convention is separately defined and bridged. \\
\phantomsection\label{entry:symbol-registry-06276}\(v\) & \texttt{\detokenize{v}} & complete valuation of earlier variables in elimination proof & P12 QRes finite-domain valuations & meta & \hyperref[def:p12-qres-pcnf]{\ref*{def:p12-qres-pcnf}} & \thepage & Quantifier-elimination equality is established for every such valuation. Root-authorized declaration/type reconciliation: Actual finite-domain valuation declaration now occurs in the PCNF definition BEFORE proof use; the raw earlier annex first-use remains historical evidence. \\
\phantomsection\label{entry:symbol-registry-06277}\(D\) & \texttt{\detokenize{D}} & universal-parent remainder in elimination proof & \hyperref[annex:p12-qres-elimination-step]{\ref*{annex:p12-qres-elimination-step}} & object & \hyperref[annex:p12-qres-elimination-step]{\ref*{annex:p12-qres-elimination-step}} & \thepage & Contains neither polarity of the eliminated universal. \\
\phantomsection\label{entry:symbol-registry-06560}\(\varphi\equiv_\alpha\psi\) & \texttt{\detokenize{\equiv_\alpha}} & source alpha-equivalence of first-order formulas in the alpha rule side condition & \hyperref[sec:rule-sheets-01]{\ref*{sec:rule-sheets-01}} & meta & \hyperref[sec:rule-sheets-01]{\ref*{sec:rule-sheets-01}} & \thepage & Same role as the row scoped to sec:fol-syntax-semantics-04; side condition of the bound-variable conversion rule on the rule sheet (def:p01-w02-p01-025-equivalence-roles item 4). \\
\phantomsection\label{entry:symbol-registry-06561}\(\varphi\equiv\psi\) & \texttt{\detokenize{\equiv}} & schematic register entry: logical equivalence relative to a declared semantic class & \hyperref[sec:latex-style-01]{\ref*{sec:latex-style-01}}--\hyperref[sec:registered-notation-01]{\ref*{sec:registered-notation-01}} & meta & \hyperref[sec:latex-style-01]{\ref*{sec:latex-style-01}} & \thepage & Style rule of sec:latex-style-01 contrasting the metalinguistic relation with the object-language biconditional, and the notation-table entry of app:notation (the chapter preamble before sec:registered-notation-01), available after a semantic class is declared; instances are the CPL row and the first-order rows (def:p01-w02-p01-025-equivalence-roles item 5). No other use of the glyph lies in the range. \\
\end{longtable}
\endgroup
\Needspace{.4\textheight}
\subsection*{Stable part P13 scopes}
\begingroup\scriptsize\setlength{\tabcolsep}{3pt}
\begin{longtable}{@{}p{0.9in}p{1.3in}p{1.0in}p{0.5in}p{0.3in}p{0.4in}p{0.35in}p{1.65in}@{}}
\toprule
\textbf{Glyph} & \textbf{Macro} & \textbf{Role} & \textbf{Scope} & \textbf{Prec./form} & \textbf{First} & \textbf{Reg. page} & \textbf{Principal uses / notes} \\
\midrule
\endfirsthead
\toprule
\textbf{Glyph} & \textbf{Macro} & \textbf{Role} & \textbf{Scope} & \textbf{Prec./form} & \textbf{First} & \textbf{Reg. page} & \textbf{Principal uses / notes} \\
\midrule
\endhead
\midrule\multicolumn{8}{r}{\emph{continued}}\\\endfoot
\bottomrule\endlastfoot
\phantomsection\label{entry:symbol-registry-04050}\(r,w,p,q\) & \texttt{\detokenize{r;w;p;q}} & Boolean atoms in declared historical-example checks & P13 descriptive-normative formal benchmarks & n/a & \hyperref[thm:p13-r11-logical-benchmarks]{\ref*{thm:p13-r11-logical-benchmarks}} & \thepage & rain/wet and generic implication operands; r is not a first-order term here. \\
\phantomsection\label{entry:symbol-registry-04051}\(D,N\) & \texttt{\detokenize{D;N}} & descriptive and normative statement atoms in the stipulated is-ought example & P13 descriptive-normative formal benchmarks & n/a & \hyperref[thm:p13-r11-logical-benchmarks]{\ref*{thm:p13-r11-logical-benchmarks}} & \thepage & N is uninterpreted in the Boolean proof; no deontic operator or normative axiom has been added. \\
\phantomsection\label{entry:symbol-registry-04052}\(A,B,C\) & \texttt{\detokenize{A;B;C}} & term extensions in the review's syllogism countermodels & P13 Evans rationality debate syllogisms & n/a & \hyperref[def:p13-r12-syllogisms]{\ref*{def:p13-r12-syllogisms}} & \thepage & All terms are nonempty in the content countermodel; validity is distinct from conclusion believability. \\
\phantomsection\label{entry:symbol-registry-04053}\(u,v\) & \texttt{\detokenize{u;v}} & two distinct objects in the explicit countermodels & P13 Evans rationality debate syllogisms & n/a & \hyperref[def:p13-r12-syllogisms]{\ref*{def:p13-r12-syllogisms}} & \thepage & These are formal witnesses not study participants. \\
\phantomsection\label{entry:symbol-registry-04054}\(p,q\) & \texttt{\detokenize{p;q}} & vowel and even predicates for the material selection rule & P13 Evans rationality debate selection task & n/a & \hyperref[def:p13-r12-selection]{\ref*{def:p13-r12-selection}} & \thepage & The false case is p and not q; advance selection and adaptive stopping have different stages. \\
\phantomsection\label{entry:symbol-registry-04055}\(E,P,4,7\) & \texttt{\detokenize{E;P;4;7}} & visible card-face labels in the source's selection example & P13 Evans rationality debate selection task & n/a & \hyperref[def:p13-r12-selection]{\ref*{def:p13-r12-selection}} & \thepage & P is the displayed consonant here and is not a probability function; the cards also have independently unknown opposite faces. \\
\phantomsection\label{entry:symbol-registry-04056}\(T,F\) & \texttt{\detokenize{T;F}} & bank-teller and feminist events in the Linda interpretation example & P13 Evans rationality debate interpretation & n/a & \hyperref[sec:p13-r12-interpretation]{\ref*{sec:p13-r12-interpretation}} & \thepage & Nested conjunction comparison differs from a comparison against teller and not feminist. \\
\phantomsection\label{entry:symbol-registry-04057}\(P\) & \texttt{\detokenize{P}} & probability measure in the explicit toy distributions & P13 Evans rationality debate probability comparisons & n/a & \hyperref[sec:p13-rationality-1-2]{\ref*{sec:p13-rationality-1-2}} & \thepage & The numbers are stipulated mathematical examples not participant estimates or historical frequencies. \\
\phantomsection\label{entry:symbol-registry-04058}\(p,q\) & \texttt{\detokenize{p;q}} & antecedent and consequent of the review's material-versus-suppositional contrast & P13 Evans rationality debate conditional semantics & n/a & \hyperref[sec:p13-r12-conditional-semantics]{\ref*{sec:p13-r12-conditional-semantics}} & \thepage & Locally typed propositions distinct from the selection-task predicate role; the displayed political example is a stipulated valuation rather than current news. \\
\phantomsection\label{entry:symbol-registry-05761}\(X_1,X_2\) & \texttt{\detokenize{X_1,}\allowbreak\detokenize{X_2}} & source absolute mode symbols & p13-street-2004-modal-tables & n/a & \hyperref[sec:p13-street-2004-modal-tables]{\ref*{sec:p13-street-2004-modal-tables}} & \thepage & General and special absolute modes; not object-language variables. \\
\phantomsection\label{entry:symbol-registry-05762}\(A\) & \texttt{\detokenize{A}} & source perpetual mode symbol & p13-street-2004-modal-tables & n/a & \hyperref[sec:p13-street-2004-modal-tables]{\ref*{sec:p13-street-2004-modal-tables}} & \thepage & Typed by position as mode before predicate and subject; distinct from predicate A in the rendition. \\
\phantomsection\label{entry:symbol-registry-05763}\(A\) & \texttt{\detokenize{A}} & predicate in the source temporal-modal rendition & p13-street-2004-modal-tables & n/a & \hyperref[sec:p13-street-2004-modal-tables]{\ref*{sec:p13-street-2004-modal-tables}} & \thepage & Actual body426-434 declares A,B predicates while the SAME modal-table region retains A as a mode. Argument/position type distinguishes readings; no artificial disjoint scope. \\
\phantomsection\label{entry:symbol-registry-05764}\(B\) & \texttt{\detokenize{B}} & predicate in the source temporal-modal rendition & p13-street-2004-modal-tables & n/a & \hyperref[sec:p13-street-2004-modal-tables]{\ref*{sec:p13-street-2004-modal-tables}} & \thepage & Actual body426-434 declares A,B predicates while the SAME modal-table region retains A as a mode. Argument/position type distinguishes readings; no artificial disjoint scope. \\
\phantomsection\label{entry:symbol-registry-05765}\(M_1,M_2,L\) & \texttt{\detokenize{M_1,}\allowbreak\detokenize{M_2,}\allowbreak\detokenize{L}} & possible and necessary mode symbols & p13-street-2004-modal-tables & n/a & \hyperref[sec:p13-street-2004-modal-tables]{\ref*{sec:p13-street-2004-modal-tables}} & \thepage & Mode triples ordered major minor conclusion; not sets/model metavariables. \\
\phantomsection\label{entry:symbol-registry-05766}\(w\) & \texttt{\detokenize{w}} & descriptional reading marker & p13-street-2004-modal-tables & n/a & \hyperref[sec:p13-street-2004-modal-tables]{\ref*{sec:p13-street-2004-modal-tables}} & \thepage & Superscript on a mode; no exponentiation interpretation. \\
\phantomsection\label{entry:symbol-registry-05767}\(a,e,i,o\) & \texttt{\detokenize{a,}\allowbreak\detokenize{e,}\allowbreak\detokenize{i,}\allowbreak\detokenize{o}} & quantity and quality markers & p13-street-2004-modal-tables & n/a & \hyperref[sec:p13-street-2004-modal-tables]{\ref*{sec:p13-street-2004-modal-tables}} & \thepage & Universal affirmative/negative and particular affirmative/negative. \\
\phantomsection\label{entry:symbol-registry-05768}\(\not\equiv\) & \texttt{\detokenize{\not\equiv}} & source contradictory pairing & p13-street-2004-modal-tables & n/a & \hyperref[sec:p13-street-2004-modal-tables]{\ref*{sec:p13-street-2004-modal-tables}} & \thepage & Local contradiction-table role; not ordinary nonequivalence. \\
\phantomsection\label{entry:symbol-registry-05769}\(R_t(Q_x)\) & \texttt{\detokenize{R_t}} & realization at time & p13-street-2004-modal-tables & n/a & \hyperref[sec:p13-street-2004-modal-tables]{\ref*{sec:p13-street-2004-modal-tables}} & \thepage & Individual x and time t; not an accessibility relation. \\
\phantomsection\label{entry:symbol-registry-05770}\(T,S\) & \texttt{\detokenize{T,}\allowbreak\detokenize{S}} & specified time constants & p13-street-2004-modal-tables & n/a & \hyperref[sec:p13-street-2004-modal-tables]{\ref*{sec:p13-street-2004-modal-tables}} & \thepage & Source temporal and spread instantiations; not sets or a truth predicate. \\
\phantomsection\label{entry:symbol-registry-05771}\(\exists Q_x,\forall Q_x\) & \texttt{\detokenize{\exists,}\allowbreak\detokenize{\forall}} & abbreviated temporal operators & p13-street-2004-modal-tables & n/a & \hyperref[sec:p13-street-2004-modal-tables]{\ref*{sec:p13-street-2004-modal-tables}} & \thepage & Quantify time through R\_t; distinct from the outer object binder forall x. \\
\phantomsection\label{entry:symbol-registry-05772}\(\supset\) & \texttt{\detokenize{\supset}} & material conditional in rendition & p13-street-2004-modal-tables & n/a & \hyperref[sec:p13-street-2004-modal-tables]{\ref*{sec:p13-street-2004-modal-tables}} & \thepage & Not set inclusion; confined to the reported modern temporal-modal rendition. \\
\phantomsection\label{entry:symbol-registry-05773}A B & \texttt{\detokenize{A B}} & Predicate and subject schematic placeholders in Boger's predicate-first patterns & p13 src063-0103 & n/a & \hyperref[sec:p13-boger-categorical-grammar]{\ref*{sec:p13-boger-categorical-grammar}} & \thepage & Not bound or free individual variables; reverse term order when comparing Johnson Original precedence 'atomic' retained as semantic descriptor; n/a is not an invented rank. \\
\phantomsection\label{entry:symbol-registry-05774}a & \texttt{\detokenize{a}} & Universal attributive categorical relation in Boger's AaB & p13 src063-0103 & n/a & \hyperref[sec:p13-boger-categorical-grammar]{\ref*{sec:p13-boger-categorical-grammar}} & \thepage & Published reconstruction's lower-case logical constant Original precedence 'atomic' retained as semantic descriptor; n/a is not an invented rank. \\
\phantomsection\label{entry:symbol-registry-05775}e & \texttt{\detokenize{e}} & Universal privative categorical relation in Boger's AeB & p13 src063-0103 & n/a & \hyperref[sec:p13-boger-categorical-grammar]{\ref*{sec:p13-boger-categorical-grammar}} & \thepage & Published reconstruction's lower-case logical constant Original precedence 'atomic' retained as semantic descriptor; n/a is not an invented rank. \\
\phantomsection\label{entry:symbol-registry-05776}i & \texttt{\detokenize{i}} & Partial attributive categorical relation in Boger's AiB & p13 src063-0103 & n/a & \hyperref[sec:p13-boger-categorical-grammar]{\ref*{sec:p13-boger-categorical-grammar}} & \thepage & Determinate versus indeterminate some remains explicitly distinguished Original precedence 'atomic' retained as semantic descriptor; n/a is not an invented rank. \\
\phantomsection\label{entry:symbol-registry-05777}o & \texttt{\detokenize{o}} & Partial privative categorical relation in Boger's AoB & p13 src063-0103 & n/a & \hyperref[sec:p13-boger-categorical-grammar]{\ref*{sec:p13-boger-categorical-grammar}} & \thepage & Not an unqualified empty-class existential Original precedence 'atomic' retained as semantic descriptor; n/a is not an invented rank. \\
\phantomsection\label{entry:symbol-registry-05778}\ensuremath{\alpha} & \texttt{\detokenize{\alpha}} & Guarded universal Boolean comparison condition & p13 src063-0103 & n/a & \hyperref[def:p13-boger-guarded-bridge]{\ref*{def:p13-boger-guarded-bridge}} & \thepage & Editorial modern bridge; not historical ontology Original precedence 'atomic' retained as semantic descriptor; n/a is not an invented rank. \\
\phantomsection\label{entry:symbol-registry-05779}\ensuremath{\iota} & \texttt{\detokenize{\iota}} & Positive particular Boolean comparison condition & p13 src063-0103 & n/a & \hyperref[def:p13-boger-guarded-bridge]{\ref*{def:p13-boger-guarded-bridge}} & \thepage & Nonempty intersection Original precedence 'atomic' retained as semantic descriptor; n/a is not an invented rank. \\
\phantomsection\label{entry:symbol-registry-05780}\ensuremath{\varepsilon} & \texttt{\detokenize{\varepsilon}} & Negative universal Boolean comparison condition & p13 src063-0103 & n/a & \hyperref[def:p13-boger-guarded-bridge]{\ref*{def:p13-boger-guarded-bridge}} & \thepage & Negation of iota Original precedence 'atomic' retained as semantic descriptor; n/a is not an invented rank. \\
\phantomsection\label{entry:symbol-registry-05781}\ensuremath{\omega} & \texttt{\detokenize{\omega}} & Negative particular comparison condition & p13 src063-0103 & n/a & \hyperref[def:p13-boger-guarded-bridge]{\ref*{def:p13-boger-guarded-bridge}} & \thepage & Negation of guarded alpha; true for empty subject Original precedence 'atomic' retained as semantic descriptor; n/a is not an invented rank. \\
\phantomsection\label{entry:symbol-registry-05782}S P & \texttt{\detokenize{S P}} & Subject and predicate extensions in the editorial bridge & p13 src063-0103 & n/a & \hyperref[def:p13-boger-guarded-bridge]{\ref*{def:p13-boger-guarded-bridge}} & \thepage & Distinct from historical schematic letter positions Original precedence 'atomic' retained as semantic descriptor; n/a is not an invented rank. \\
\phantomsection\label{entry:symbol-registry-05783}A & \texttt{\detokenize{A}} & Universal affirmative binary categorical operator in Johnson's subject-first notation & p13 src063-0104 & n/a & \hyperref[sec:p13-johnson-assertoric]{\ref*{sec:p13-johnson-assertoric}} & \thepage & Aab means all a are b Original precedence 'atomic' retained as semantic descriptor; n/a is not an invented rank. \\
\phantomsection\label{entry:symbol-registry-05784}I & \texttt{\detokenize{I}} & Particular affirmative binary categorical operator in Johnson's notation & p13 src063-0104 & n/a & \hyperref[sec:p13-johnson-assertoric]{\ref*{sec:p13-johnson-assertoric}} & \thepage & Iaa in LA carries term nonemptiness Original precedence 'atomic' retained as semantic descriptor; n/a is not an invented rank. \\
\phantomsection\label{entry:symbol-registry-05785}E & \texttt{\detokenize{E}} & Defined negative universal categorical operator & p13 src063-0104 & n/a & \hyperref[sec:p13-johnson-assertoric]{\ref*{sec:p13-johnson-assertoric}} & \thepage & Eab is neg Iab Original precedence 'atomic' retained as semantic descriptor; n/a is not an invented rank. \\
\phantomsection\label{entry:symbol-registry-05786}O & \texttt{\detokenize{O}} & Defined negative particular categorical operator & p13 src063-0104 & n/a & \hyperref[sec:p13-johnson-assertoric]{\ref*{sec:p13-johnson-assertoric}} & \thepage & Oab is neg Aab Original precedence 'atomic' retained as semantic descriptor; n/a is not an invented rank. \\
\phantomsection\label{entry:symbol-registry-05787}a b c & \texttt{\detokenize{a b c}} & Term variables or interpreted term extensions in Johnson's reconstruction & p13 src063-0104 & n/a & \hyperref[sec:p13-johnson-assertoric]{\ref*{sec:p13-johnson-assertoric}} & \thepage & Formula letters retain original subject-first syntax Original precedence 'atomic' retained as semantic descriptor; n/a is not an invented rank. \\
\phantomsection\label{entry:symbol-registry-05788}\ensuremath{\vdash} & \texttt{\detokenize{\vdash}} & System-relative assertion judgment & p13 src063-0104 & n/a & \hyperref[sec:p13-johnson-assertoric]{\ref*{sec:p13-johnson-assertoric}} & \thepage & Not object-language implication or empirical acceptance Original precedence 'judgment' retained as semantic descriptor; n/a is not an invented rank. \\
\phantomsection\label{entry:symbol-registry-05789}\ensuremath{\dashv} & \texttt{\detokenize{\dashv}} & System-relative rejection judgment & p13 src063-0104 & n/a & \hyperref[sec:p13-johnson-assertoric]{\ref*{sec:p13-johnson-assertoric}} & \thepage & Not assertion of a formula's negation Original precedence 'judgment' retained as semantic descriptor; n/a is not an invented rank. \\
\phantomsection\label{entry:symbol-registry-05790}L & \texttt{\detokenize{L}} & Necessity operator in the four-valued LM matrix & p13 src063-0104 LM comparison & n/a & \hyperref[sec:p13-johnson-lm]{\ref*{sec:p13-johnson-lm}} & \thepage & Not unrestricted normal modal necessity Original precedence 'unary' retained as semantic descriptor; n/a is not an invented rank. \\
\phantomsection\label{entry:symbol-registry-05791}M & \texttt{\detokenize{M}} & Possibility operator in the four-valued LM matrix & p13 src063-0104 LM comparison & n/a & \hyperref[sec:p13-johnson-lm]{\ref*{sec:p13-johnson-lm}} & \thepage & Table fixed by four values and one designated value Original precedence 'unary' retained as semantic descriptor; n/a is not an invented rank. \\
\phantomsection\label{entry:symbol-registry-05792}X L & \texttt{\detokenize{X L}} & Assertoric and necessary premise-conclusion types in mood combinations & p13 src063-0104 modal classification & n/a & \hyperref[sec:p13-johnson-lm]{\ref*{sec:p13-johnson-lm}} & \thepage & XLL is an ordered modal combination; not an object-language formula Original precedence 'metadata' retained as semantic descriptor; n/a is not an invented rank. \\
\phantomsection\label{entry:symbol-registry-05793}\ensuremath{\lambda} & \texttt{\detokenize{\lambda}} & Necessary complex-term constructor in Geach's G-system & p13 src063-0104 G comparison & n/a & \hyperref[sec:p13-johnson-geach]{\ref*{sec:p13-johnson-geach}} & \thepage & lambda a subset a Original precedence 'unary' retained as semantic descriptor; n/a is not an invented rank. \\
\phantomsection\label{entry:symbol-registry-05794}\ensuremath{\mu} & \texttt{\detokenize{\mu}} & Possible complex-term constructor in Geach's G-system & p13 src063-0104 G comparison & n/a & \hyperref[sec:p13-johnson-geach]{\ref*{sec:p13-johnson-geach}} & \thepage & a subset mu a Original precedence 'unary' retained as semantic descriptor; n/a is not an invented rank. \\
\phantomsection\label{entry:symbol-registry-05795}W D & \texttt{\detokenize{W D}} & World set and constant individual domain of the declared S5 countermodel & p13 src063-0104 modal proof & n/a & \hyperref[annex:p13.johnson.modal]{\ref*{annex:p13.johnson.modal}} & \thepage & Frame assumptions are explicit Original precedence 'atomic' retained as semantic descriptor; n/a is not an invented rank. \\
\phantomsection\label{entry:symbol-registry-05796}\ensuremath{\Rightarrow} & \texttt{\detokenize{\Rightarrow}} & Classical implication between Boolean bits in the four-value representation proof & p13 src063-0104 LM annex & n/a & \hyperref[annex:p13.johnson.lm]{\ref*{annex:p13.johnson.lm}} & \thepage & Distinct local type from formula implication; coordinates explicitly defined Original precedence 'binary' retained as semantic descriptor; n/a is not an invented rank. \\
\phantomsection\label{entry:symbol-registry-05985}P M S & \texttt{\detokenize{P M S}} & Major middle and minor schematic term places in Boger's figure patterns & p13 src063-0103-deduction & n/a & \hyperref[sec:p13-boger-deduction-formation]{\ref*{sec:p13-boger-deduction-formation}} & \thepage & Predicate-first; no individual quantification Original precedence 'atomic' retained as semantic descriptor; n/a is not an invented rank. \\
\phantomsection\label{entry:symbol-registry-05986}x y z & \texttt{\detokenize{x y z}} & Metalinguistic places for the four categorical constants & p13 src063-0103-deduction & n/a & \hyperref[sec:p13-boger-deduction-formation]{\ref*{sec:p13-boger-deduction-formation}} & \thepage & Not bound or free individual variables Original precedence 'atomic' retained as semantic descriptor; n/a is not an invented rank. \\
\phantomsection\label{entry:symbol-registry-05987}\ensuremath{c^{*}} & \texttt{\detokenize{c^{*}}} & Categorical contradictory of a sentence in editorial deduction bookkeeping & p13 src063-0103-deduction & n/a & \hyperref[def:p13-boger-deduction-bookkeeping]{\ref*{def:p13-boger-deduction-bookkeeping}} & \thepage & a exchanges with o and e with i; no new object-language negation Original precedence 'unary metalinguistic' retained as semantic descriptor; n/a is not an invented rank. \\
\phantomsection\label{entry:symbol-registry-05988}\ensuremath{\Gamma} & \texttt{\detokenize{\Gamma}} & Finite categorical premise context in editorial bookkeeping & p13 src063-0103-deduction & n/a & \hyperref[def:p13-boger-deduction-bookkeeping]{\ref*{def:p13-boger-deduction-bookkeeping}} & \thepage & Not Johnson's axiom or a historical symbol Original precedence 'metalinguistic set' retained as semantic descriptor; n/a is not an invented rank. \\
\phantomsection\label{entry:symbol-registry-05989}\ensuremath{\vdash_{\mathrm{full}}} & \texttt{\detokenize{\vdash_{\mathrm{full}}}} & Finite derivability with Boger's full fourteen-rule inventory & p13 src063-0103-deduction & n/a & \hyperref[def:p13-boger-deduction-bookkeeping]{\ref*{def:p13-boger-deduction-bookkeeping}} & \thepage & Editorial notation; distinct from Johnson assertion/rejection Original precedence 'judgment' retained as semantic descriptor; n/a is not an invented rank. \\
\phantomsection\label{entry:symbol-registry-05990}\ensuremath{\vdash_{\mathrm{red}}} & \texttt{\detokenize{\vdash_{\mathrm{red}}}} & Finite derivability with only Barbara and Celarent as two-premise rules & p13 src063-0103-deduction & n/a & \hyperref[def:p13-boger-reduced-system]{\ref*{def:p13-boger-reduced-system}} & \thepage & Conversions repetition opposition and indirect deduction remain Original precedence 'judgment' retained as semantic descriptor; n/a is not an invented rank. \\
\phantomsection\label{entry:symbol-registry-05991}\ensuremath{\vdash} & \texttt{\detokenize{\vdash}} & System-parametric deduction judgment in the categorical indirect-rule sheet & p13 p13-boger-deduction sheet & n/a & \hyperref[rule:p13-boger-indirect]{\ref*{rule:p13-boger-indirect}} & \thepage & Full or reduced bookkeeping; not Johnson rejection or object-language implication Original precedence 'judgment' retained as semantic descriptor; n/a is not an invented rank. \\
\phantomsection\label{entry:symbol-registry-05992}A B C; M N X; P R S & \texttt{\detokenize{A B C; M N X; P R S}} & Schematic term places in the original three figure-specific rule blocks & p13 src063-0103-deduction & n/a & \hyperref[sec:p13-boger-deduction-rules]{\ref*{sec:p13-boger-deduction-rules}} & \thepage & Original predicate-first formulas retained; letters are not individual variables Original precedence 'atomic' retained as semantic descriptor; n/a is not an invented rank. \\
\phantomsection\label{entry:symbol-registry-05993}a e i o & \texttt{\detokenize{a e i o}} & Four categorical belonging constants in Boger's predicate-first syntax & p13 src063-0103-deduction & n/a & \hyperref[sec:p13-boger-deduction]{\ref*{sec:p13-boger-deduction}} & \thepage & Same roles as the prior original packet; include existent-subject and indeterminate-some policy Original precedence 'atomic' retained as semantic descriptor; n/a is not an invented rank. \\
\phantomsection\label{entry:symbol-registry-06067}A E I O & \texttt{\detokenize{A E I O}} & Subject-first categorical operators in Johnson's native calculus & p13 src063-0104-mccall & n/a & \hyperref[def:p13-johnson-mccall-grammar]{\ref*{def:p13-johnson-mccall-grammar}} & \thepage & E is neg I; O is neg A; not Boger's predicate-first constants Original precedence descriptor 'atomic' preserved; n/a is not an invented rank. \\
\phantomsection\label{entry:symbol-registry-06068}L & \texttt{\detokenize{L}} & Restricted categorical necessity operator & p13 src063-0104-mccall & n/a & \hyperref[def:p13-johnson-mccall-grammar]{\ref*{def:p13-johnson-mccall-grammar}} & \thepage & Lp only for categorical expression p; LLAab not well formed Original precedence descriptor 'unary' preserved; n/a is not an invented rank. \\
\phantomsection\label{entry:symbol-registry-06069}M & \texttt{\detokenize{M}} & Defined categorical possibility neg L neg p & p13 src063-0104-mccall & n/a & \hyperref[def:p13-johnson-mccall-grammar]{\ref*{def:p13-johnson-mccall-grammar}} & \thepage & Not the earlier four-valued LM modal matrix Original precedence descriptor 'unary' preserved; n/a is not an invented rank. \\
\phantomsection\label{entry:symbol-registry-06070}X & \texttt{\detokenize{X}} & Assertoric mode position in a three-letter mood code & p13 src063-0104-mccall & n/a & \hyperref[sec:p13-johnson-mccall]{\ref*{sec:p13-johnson-mccall}} & \thepage & Not an object-language operator Original precedence descriptor 'metadata' preserved; n/a is not an invented rank. \\
\phantomsection\label{entry:symbol-registry-06071}\ensuremath{\vdash} & \texttt{\detokenize{\vdash}} & System-relative native assertion judgment & p13 src063-0104-mccall & n/a & \hyperref[sec:p13-johnson-mccall]{\ref*{sec:p13-johnson-mccall}} & \thepage & Distinct from rejection and from Boger's deduction bookkeeping Original precedence descriptor 'judgment' preserved; n/a is not an invented rank. \\
\phantomsection\label{entry:symbol-registry-06072}\ensuremath{\dashv} & \texttt{\detokenize{\dashv}} & System-relative native rejection judgment & p13 src063-0104-mccall & n/a & \hyperref[sec:p13-johnson-mccall]{\ref*{sec:p13-johnson-mccall}} & \thepage & Rejection of p is not assertion of neg p Original precedence descriptor 'judgment' preserved; n/a is not an invented rank. \\
\phantomsection\label{entry:symbol-registry-06073}\ensuremath{n^{+} q^{+} n^{-} q^{-}} & \texttt{\detokenize{n^{+} q^{+} n^{-} q^{-}}} & Four object regions assigned to a term in native semantics & p13 src063-0104-mccall & n/a & \hyperref[def:p13-johnson-jmodel]{\ref*{def:p13-johnson-jmodel}} & \thepage & Pairwise disjoint partition; n-plus nonemptiness is the J-model condition Original precedence descriptor 'set valued function' preserved; n/a is not an invented rank. \\
\phantomsection\label{entry:symbol-registry-06074}+(x) & \texttt{\detokenize{+(x)}} & Actual positive term extension n-plus union q-plus & p13 src063-0104-mccall & n/a & \hyperref[def:p13-johnson-jmodel]{\ref*{def:p13-johnson-jmodel}} & \thepage & Source notation retained; not unary arithmetic plus Original precedence descriptor 'set valued abbreviation' preserved; n/a is not an invented rank. \\
\phantomsection\label{entry:symbol-registry-06075}o & \texttt{\detokenize{\mathbin{o}}} & Nonempty overlap of two sets & p13 src063-0104-mccall & n/a & \hyperref[def:p13-johnson-jmodel]{\ref*{def:p13-johnson-jmodel}} & \thepage & Native semantic notation; not categorical O or Boger's lower-case o Original precedence descriptor 'binary relation' preserved; n/a is not an invented rank. \\
\phantomsection\label{entry:symbol-registry-06076}W & \texttt{\detokenize{W}} & Native object domain & p13 src063-0104-mccall & n/a & \hyperref[def:p13-johnson-jmodel]{\ref*{def:p13-johnson-jmodel}} & \thepage & Not a set of modal worlds Original precedence descriptor 'set' preserved; n/a is not an invented rank. \\
\phantomsection\label{entry:symbol-registry-06077}Q & \texttt{\detokenize{Q}} & Primitive restricted categorical contingency operator & p13 src063-0104-mccall & n/a & \hyperref[sec:p13-johnson-mccall-contingency]{\ref*{sec:p13-johnson-mccall-contingency}} & \thepage & Not defined as unrestricted Mp and M neg p; no nesting Original precedence descriptor 'unary' preserved; n/a is not an invented rank. \\
\phantomsection\label{entry:symbol-registry-06078}q(x) & \texttt{\detokenize{q(x)}} & Contingent term region q-plus union q-minus & p13 src063-0104-mccall & n/a & \hyperref[def:p13-johnson-qmodel]{\ref*{def:p13-johnson-qmodel}} & \thepage & Distinct from schematic propositional letter q in rule displays Original precedence descriptor 'set valued abbreviation' preserved; n/a is not an invented rank. \\
\phantomsection\label{entry:symbol-registry-06079}\ensuremath{\sigma_2 \sigma_3 \sigma_4} & \texttt{\detokenize{\sigma_2 \sigma_3 \sigma_4}} & Three literal rejected-formula abbreviations & p13 src063-0104-mccall & n/a & \hyperref[sec:p13-johnson-mccall-rejection]{\ref*{sec:p13-johnson-mccall-rejection}} & \thepage & Not assertion axioms; source R3 final consequent Aac Original precedence descriptor 'formula metaname' preserved; n/a is not an invented rank. \\
\phantomsection\label{entry:symbol-registry-06080}{}[A] [I] & \texttt{\detokenize{[A] [I]}} & Source tally places for A-or-E and I-or-O quality choices & p13 src063-0104-mccall tally & n/a & \hyperref[sec:p13-johnson-mccall-limits]{\ref*{sec:p13-johnson-mccall-limits}} & \thepage & No new categorical operators; grouped counts are source claims Original precedence descriptor 'metadata' preserved; n/a is not an invented rank. \\
\phantomsection\label{entry:symbol-registry-06081}BT3 BT4 BT5 BQ2 BQ3 BQ4 BQ5 BQ6 BQ7 BQ8 & \texttt{\detokenize{BT3 BT4 BT5 BQ2 BQ3 BQ4 BQ5 BQ6 BQ7 BQ8}} & Named universally quantified native model conditions & p13 src063-0104-mccall & n/a & \hyperref[sec:p13-johnson-mccall-t3-chain]{\ref*{sec:p13-johnson-mccall-t3-chain}} & \thepage & BQ7 printed repeated overlap retained; ordinary Q-model nonemptiness explicit Original precedence descriptor 'semantic condition names' preserved; n/a is not an invented rank. \\
\phantomsection\label{entry:symbol-registry-06082}\ensuremath{\mathcal M_i} & \texttt{\detokenize{\mathcal M_i}} & Named source term-region structures and their row permutations & p13 src063-0104-mccall & n/a & \hyperref[sec:p13-johnson-mccall-countermodels]{\ref*{sec:p13-johnson-mccall-countermodels}} & \thepage & Model1 to Model4 are J-models; Model5 to Model7 checked as completed Q-models Original precedence descriptor 'model metaname' preserved; n/a is not an invented rank. \\
\phantomsection\label{entry:symbol-registry-06083}C[p] & \texttt{\detokenize{C[p]}} & Well-formed formula context for declared definition and DN operations & p13 p13-johnson-mccall sheet & n/a & \hyperref[rule:p13-johnson-mccall-def]{\ref*{rule:p13-johnson-mccall-def}} & \thepage & Categorical p restriction and well-formedness prevent arbitrary nested modality Original precedence descriptor 'metalinguistic context' preserved; n/a is not an invented rank. \\
\phantomsection\label{entry:symbol-registry-06084}\ensuremath{\sigma} & \texttt{\detokenize{\sigma}} & Uniform map on term variables in the native rule sheet & p13 p13-johnson-mccall sheet & n/a & \hyperref[rule:p13-johnson-mccall-us]{\ref*{rule:p13-johnson-mccall-us}} & \thepage & Distinct from the sigma-indexed rejected-formula names in the reader Original precedence descriptor 'metalinguistic substitution' preserved; n/a is not an invented rank. \\
\phantomsection\label{entry:symbol-registry-06180}P Q & \texttt{\detokenize{literal proposition metavariables}} & Respondent commitments in a modern dialogue reconstruction & P13 Ganeri debate & n/a & \hyperref[sec:p13-ganeri-debate]{\ref*{sec:p13-ganeri-debate}} & \thepage & Not ancient object-language letters Original precedence descriptor 'atomic' preserved; n/a is not an invented rank. \\
\phantomsection\label{entry:symbol-registry-06181}F G & \texttt{\detokenize{literal predicate metavariables}} & Reason and target in the Nyaya reconstruction & P13 Ganeri bridge and number & n/a & \hyperref[sec:p13-ganeri-debate]{\ref*{sec:p13-ganeri-debate}} & \thepage & Roles are stated locally; not reversed in universal instantiation Original precedence descriptor 'atomic' preserved; n/a is not an invented rank. \\
\phantomsection\label{entry:symbol-registry-06182}a b p & \texttt{\detokenize{literal site names}} & Object or site of a reason & P13 Ganeri reasons & n/a & \hyperref[sec:p13-ganeri-reasons]{\ref*{sec:p13-ganeri-reasons}} & \thepage & p is a site here but a formula metavariable in the standpoint section Original precedence descriptor 'atomic' preserved; n/a is not an invented rank. \\
\phantomsection\label{entry:symbol-registry-06183}preceq & \texttt{\detokenize{\preceq}} & Abstract pervasion preorder & P13 Ganeri tarka & relation & \hyperref[sec:p13-ganeri-debate]{\ref*{sec:p13-ganeri-debate}} & \thepage & Reflexive and transitive; no empirical law is established by notation \\
\phantomsection\label{entry:symbol-registry-06184}H S & \texttt{\detokenize{literal predicates}} & Reason and target properties & P13 Ganeri Buddhist reasons & n/a & \hyperref[sec:p13-ganeri-reasons]{\ref*{sec:p13-ganeri-reasons}} & \thepage & Positive comparison is relative to S Original precedence descriptor 'atomic' preserved; n/a is not an invented rank. \\
\phantomsection\label{entry:symbol-registry-06185}sp vp & \texttt{\detokenize{\mathsf{sp};\mathsf{vp}}} & Positive and negative comparison classes & P13 Ganeri wheels & n/a & \hyperref[sec:p13-ganeri-reasons]{\ref*{sec:p13-ganeri-reasons}} & \thepage & They need not include the disputed site Original precedence descriptor 'set' preserved; n/a is not an invented rank. \\
\phantomsection\label{entry:symbol-registry-06186}+ - +/- 0 & \texttt{\detokenize{literal status signs}} & All; none; some but not all; empty comparison class & P13 Ganeri wheels & n/a & \hyperref[sec:p13-ganeri-reasons]{\ref*{sec:p13-ganeri-reasons}} & \thepage & 0 is an extra empty-range status rather than another truth value Original precedence descriptor 'atomic' preserved; n/a is not an invented rank. \\
\phantomsection\label{entry:symbol-registry-06187}sigma:r & \texttt{\detokenize{\sigma:r}} & Assertion at a standpoint & P13 Ganeri modern assertion kernel & relation & \hyperref[sec:p13-ganeri-standpoints]{\ref*{sec:p13-ganeri-standpoints}} & \thepage & Distinct from absolute truth and from union-theory membership \\
\phantomsection\label{entry:symbol-registry-06188}Sigma & \texttt{\detokenize{\Sigma}} & Nonempty family of standpoints & P13 Ganeri modern assertion kernel & n/a & \hyperref[sec:p13-ganeri-standpoints]{\ref*{sec:p13-ganeri-standpoints}} & \thepage & Discourse; no claim that it exhausts reality Original precedence descriptor 'set' preserved; n/a is not an invented rank. \\
\phantomsection\label{entry:symbol-registry-06189}nabla & \texttt{\detokenize{\nabla}} & Assertion at some standpoint & P13 Ganeri modern assertion kernel & n/a & \hyperref[sec:p13-ganeri-standpoints]{\ref*{sec:p13-ganeri-standpoints}} & \thepage & Existential abbreviation; no global adjunction Original precedence descriptor 'unary' preserved; n/a is not an invented rank. \\
\phantomsection\label{entry:symbol-registry-06190}t f u & \texttt{\detokenize{literal support markers}} & Affirmative; negative; neutral local assertion states represented somewhere & P13 Ganeri standpoint patterns & n/a & \hyperref[sec:p13-ganeri-standpoints]{\ref*{sec:p13-ganeri-standpoints}} & \thepage & Marker predicates can overlap globally; seven nonempty support sets Original precedence descriptor 'atomic' preserved; n/a is not an invented rank. \\
\phantomsection\label{entry:symbol-registry-06191}Cn & \texttt{\detokenize{\operatorname{Cn}}} & Classical semantic consequence closure & P13 Ganeri local proof model & n/a & \hyperref[thm:p13-ganeri-standpoints]{\ref*{thm:p13-ganeri-standpoints}} & \thepage & Modern explanatory model; not an original Sanskrit calculus Original precedence descriptor 'unary' preserved; n/a is not an invented rank. \\
\phantomsection\label{entry:symbol-registry-06192}U I Q M S & \texttt{\detokenize{\mathsf U;\mathsf I;\mathsf Q;\mathsf M;\mathsf S}} & Universal; individuator; quality; motion; substance graph types & P13 Ganeri finite graph repair & n/a & \hyperref[sec:p13-ganeri-navya]{\ref*{sec:p13-ganeri-navya}} & \thepage & S here is not the target predicate of the Buddhist section Original precedence descriptor 'type' preserved; n/a is not an invented rank. \\
\phantomsection\label{entry:symbol-registry-06193}Talpha & \texttt{\detokenize{T\alpha}} & Truth at a specified location in the higher-absence reconstruction & P13 Ganeri absence & n/a & \hyperref[sec:p13-ganeri-navya]{\ref*{sec:p13-ganeri-navya}} & \thepage & External denial of T differs from object-level absence Original precedence descriptor 'unary' preserved; n/a is not an invented rank. \\
\phantomsection\label{entry:symbol-registry-06194}neg & \texttt{\detokenize{\neg}} & object-level absence constructor & P13 Ganeri absence & n/a & \hyperref[sec:p13-ganeri-navya]{\ref*{sec:p13-ganeri-navya}} & \thepage & Typing stated explicitly; no unrestricted explosion or contraposition imported Original precedence descriptor 'unary' preserved; n/a is not an invented rank. Root-authorized typed split in the SAME actual Navya declaration/rule/proof regions; original combined row remains evidence. \\
\phantomsection\label{entry:symbol-registry-06195}neg & \texttt{\detokenize{\neg}} & external classical denial of location truth & P13 Ganeri absence & n/a & \hyperref[sec:p13-ganeri-navya]{\ref*{sec:p13-ganeri-navya}} & \thepage & Typing stated explicitly; no unrestricted explosion or contraposition imported Original precedence descriptor 'unary' preserved; n/a is not an invented rank. Root-authorized typed split in the SAME actual Navya declaration/rule/proof regions; original combined row remains evidence. \\
\phantomsection\label{entry:symbol-registry-06196}bn & \texttt{\detokenize{b_n}} & Truth indicator of the nth iterated absence & P13 Ganeri local absence model & n/a & \hyperref[thm:p13-ganeri-absence]{\ref*{thm:p13-ganeri-absence}} & \thepage & The initial pair excludes 0 0; recurrence starts at n=1 Original precedence descriptor 'atomic' preserved; n/a is not an invented rank. \\
\phantomsection\label{entry:symbol-registry-06197}C n & \texttt{\detokenize{literal set and numerical node}} & Collective target and finite cardinality & P13 Ganeri number & n/a & \hyperref[sec:p13-ganeri-navya]{\ref*{sec:p13-ganeri-navya}} & \thepage & A set of distinct nodes rather than a duplicate-counting tuple Original precedence descriptor 'atomic' preserved; n/a is not an invented rank. \\
\phantomsection\label{entry:symbol-registry-07308}\(Phi\) & \texttt{\detokenize{\Phi}} & Syncategorematic placeholder or its absence & P13 Latin supposition & n/a & \hyperref[sec:p13-latin-supposition]{\ref*{sec:p13-latin-supposition}} & \thepage & Distinct later role as proposal history Original precedence descriptor 'schema' retained; no rank inferred. \\
\phantomsection\label{entry:symbol-registry-07309}\(A B\) & \texttt{\detokenize{literal term metavariables}} & Terms whose supposition is classified & P13 Latin supposition & n/a & \hyperref[sec:p13-latin-supposition]{\ref*{sec:p13-latin-supposition}} & \thepage & Not consequence-section formula metavariables Original precedence descriptor 'atomic' retained; no rank inferred. \\
\phantomsection\label{entry:symbol-registry-07310}\(S P\) & \texttt{\detokenize{literal predicate and extension names}} & Finite local model of categorical descents & P13 finite descent proof & n/a & \hyperref[thm:p13-latin-finite-descents]{\ref*{thm:p13-latin-finite-descents}} & \thepage & Classical finite nonempty extensions; not a complete historical semantics Original precedence descriptor 'atomic' retained; no rank inferred. \\
\phantomsection\label{entry:symbol-registry-07311}\(Det Dist Conf\) & \texttt{\detokenize{\operatorname{Det};\operatorname{Dist};\operatorname{Conf}}} & Determinate; confused-and-distributive; merely-confused modes & P13 supposition negation & n/a & \hyperref[sec:p13-latin-supposition]{\ref*{sec:p13-latin-supposition}} & \thepage & Novaes's modern notation for attributed variants Original precedence descriptor 'operator' retained; no rank inferred. \\
\phantomsection\label{entry:symbol-registry-07312}\(P P-star\) & \texttt{\detokenize{P;P^*}} & Original and negated propositions containing classified terms & P13 supposition negation & n/a & \hyperref[sec:p13-latin-supposition]{\ref*{sec:p13-latin-supposition}} & \thepage & Distinct from finite extension P Original precedence descriptor 'atomic' retained; no rank inferred. \\
\phantomsection\label{entry:symbol-registry-07313}\(<sim A>\) & \texttt{\detokenize{\langle\sim,}\allowbreak\detokenize{A\rangle}} & Term follows added negation in the specified proposition & P13 supposition negation & relation & \hyperref[sec:p13-latin-supposition]{\ref*{sec:p13-latin-supposition}} & \thepage & Placement condition; not a truth-functional conjunction \\
\phantomsection\label{entry:symbol-registry-07314}\(Rightarrow\) & \texttt{\detokenize{\Rightarrow}} & Source modern consequence representation & P13 consequence operations & relation & \hyperref[sec:p13-latin-consequences]{\ref*{sec:p13-latin-consequences}} & \thepage & Not the object conditional or a claim of ancient symbolism \\
\phantomsection\label{entry:symbol-registry-07315}\(nRightarrow\) & \texttt{\detokenize{\nRightarrow}} & Countervaluation-backed non-consequence & P13 consequence operations & relation & \hyperref[sec:p13-latin-consequences]{\ref*{sec:p13-latin-consequences}} & \thepage & Negative principles distinguished despite duplicated source display \\
\phantomsection\label{entry:symbol-registry-07316}\(A B C D\) & \texttt{\detokenize{literal formula metavariables}} & Antecedents and consequents & P13 consequence operations & n/a & \hyperref[sec:p13-latin-consequences]{\ref*{sec:p13-latin-consequences}} & \thepage & Distinct from terms of supposition section Original precedence descriptor 'atomic' retained; no rank inferred. \\
\phantomsection\label{entry:symbol-registry-07317}\(comma\) & \texttt{\detokenize{literal comma}} & Simultaneous premises or joint conclusions & P13 source-derived cuts & n/a & \hyperref[sec:p13-latin-consequences]{\ref*{sec:p13-latin-consequences}} & \thepage & Right comma is not the disjunctive LK comma Original precedence descriptor 'structural' retained; no rank inferred. \\
\phantomsection\label{entry:symbol-registry-07318}\(T bot\) & \texttt{\detokenize{T;\bot}} & True and false throughout the declared valuation class & P13 consequence operations & n/a & \hyperref[sec:p13-latin-consequences]{\ref*{sec:p13-latin-consequences}} & \thepage & Relative to W; restricted case does not give global necessity Original precedence descriptor 'atomic' retained; no rank inferred. \\
\phantomsection\label{entry:symbol-registry-07319}\(W\) & \texttt{\detokenize{literal W}} & Fixed valuation class for a local reconstruction & P13 consequence proofs & n/a & \hyperref[sec:p13-latin-consequences]{\ref*{sec:p13-latin-consequences}} & \thepage & Absolute versus restricted-case scope is explicit Original precedence descriptor 'set' retained; no rank inferred. \\
\phantomsection\label{entry:symbol-registry-07320}\(Vdash\) & \texttt{\detokenize{\Vdash}} & Semantic following from the specified proposition or premise base & P13 obligation games & relation & \hyperref[sec:p13-latin-burley-obligations]{\ref*{sec:p13-latin-burley-obligations}} & \thepage & Source glyph is preserved; not set-theoretic forcing \\
\phantomsection\label{entry:symbol-registry-07321}\(nVdash\) & \texttt{\detokenize{\nVdash}} & Failure of the specified semantic following & P13 obligation games & relation & \hyperref[sec:p13-latin-burley-obligations]{\ref*{sec:p13-latin-burley-obligations}} & \thepage & A countermodel rather than object-language negation \\
\phantomsection\label{entry:symbol-registry-07322}\(Ob\) & \texttt{\detokenize{\operatorname{Ob}}} & Reconstructed obligation-game tuple & P13 obligation games & n/a & \hyperref[sec:p13-latin-burley-obligations]{\ref*{sec:p13-latin-burley-obligations}} & \thepage & Burley and Swyneshed tuples have distinct data Original precedence descriptor 'tuple' retained; no rank inferred. \\
\phantomsection\label{entry:symbol-registry-07323}\(KC\) & \texttt{\detokenize{K_C}} & Consistent fixed shared information and casus & P13 Burley game & n/a & \hyperref[sec:p13-latin-burley-obligations]{\ref*{sec:p13-latin-burley-obligations}} & \thepage & Not globally united with the counterfactual positum Original precedence descriptor 'set' retained; no rank inferred. \\
\phantomsection\label{entry:symbol-registry-07324}\(Phi\) & \texttt{\detokenize{\Phi}} & Ordered proposal-occurrence history & P13 obligation games & n/a & \hyperref[sec:p13-latin-burley-obligations]{\ref*{sec:p13-latin-burley-obligations}} & \thepage & Allows repetitions; distinct from supposition placeholder Original precedence descriptor 'sequence' retained; no rank inferred. \\
\phantomsection\label{entry:symbol-registry-07325}\(Gamma\) & \texttt{\detokenize{\Gamma_n}} & Cumulative concession and denied-literal base & P13 Burley game & n/a & \hyperref[sec:p13-latin-burley-obligations]{\ref*{sec:p13-latin-burley-obligations}} & \thepage & Doubt adds no literal Original precedence descriptor 'set' retained; no rank inferred. \\
\phantomsection\label{entry:symbol-registry-07326}\(R\) & \texttt{\detokenize{R(\phi_n)}} & Reply selector with current state implicitly indexed & P13 obligation games & n/a & \hyperref[sec:p13-latin-burley-obligations]{\ref*{sec:p13-latin-burley-obligations}} & \thepage & Values 1 0 ? are replies rather than proposition truth values Original precedence descriptor 'function' retained; no rank inferred. \\
\phantomsection\label{entry:symbol-registry-07327}\(UM\) & \texttt{\detokenize{\mathrm{UM}_n;\mathrm{UM}_{\phi}}} & Classes of models satisfying a base or proposition & P13 Burley model restrictions & n/a & \hyperref[sec:p13-latin-burley-obligations]{\ref*{sec:p13-latin-burley-obligations}} & \thepage & Increasing base can strictly reduce model class Original precedence descriptor 'set' retained; no rank inferred. \\
\phantomsection\label{entry:symbol-registry-07328}\(Sigma Sn\) & \texttt{\detokenize{\Sigma;S_n}} & Sequence of consistent possibly incomplete knowledge states & P13 Swyneshed game & n/a & \hyperref[sec:p13-latin-swyneshed-obligations]{\ref*{sec:p13-latin-swyneshed-obligations}} & \thepage & State may change; reply log does not determine relevance Original precedence descriptor 'sequence' retained; no rank inferred. \\
\phantomsection\label{entry:symbol-registry-07329}\(I iota\) & \texttt{\detokenize{I;\iota_n}} & Ordered reply history and individual reply records & P13 Swyneshed game & n/a & \hyperref[sec:p13-latin-swyneshed-obligations]{\ref*{sec:p13-latin-swyneshed-obligations}} & \thepage & Proposition index can recur at several reply indices Original precedence descriptor 'sequence' retained; no rank inferred. \\
\phantomsection\label{entry:symbol-registry-07330}\(gamma kappa\) & \texttt{\detokenize{\gamma;\kappa}} & Reply values in source's proposed history test & P13 Swyneshed admission discussion & n/a & \hyperref[sec:p13-latin-swyneshed-obligations]{\ref*{sec:p13-latin-swyneshed-obligations}} & \thepage & Finite variation is not equivalent to modal contingency Original precedence descriptor 'atomic' retained; no rank inferred. \\
\phantomsection\label{entry:symbol-registry-07331}\(Delta Pi\) & \texttt{\detokenize{\Delta_{\phi_0};\Pi_{\phi_0}}} & Pertinent and impertinent proposition classes & P13 Swyneshed game & n/a & \hyperref[sec:p13-latin-swyneshed-obligations]{\ref*{sec:p13-latin-swyneshed-obligations}} & \thepage & Fixed positum; pragmatic repeat exception is separate Original precedence descriptor 'set' retained; no rank inferred. \\
\phantomsection\label{entry:symbol-registry-07332}\(L\) & \texttt{\detokenize{\mathcal L}} & Fixed universe of proposals & P13 Swyneshed partition & n/a & \hyperref[sec:p13-latin-swyneshed-obligations]{\ref*{sec:p13-latin-swyneshed-obligations}} & \thepage & Makes the source set-builder's intended domain explicit Original precedence descriptor 'set' retained; no rank inferred. \\
\phantomsection\label{entry:symbol-registry-07333}\(Cn\) & \texttt{\detokenize{\operatorname{Cn}}} & Classical semantic consequence closure & P13 partial knowledge example & n/a & \hyperref[thm:p13-latin-partial-knowledge]{\ref*{thm:p13-latin-partial-knowledge}} & \thepage & Knowledge of disjunction need not select a known disjunct Original precedence descriptor 'unary' retained; no rank inferred. \\
\phantomsection\label{entry:symbol-registry-07334}\(ampersand\) & \texttt{\detokenize{\mathbin{\&}}} & Object conjunction in the source consequence display & P13 consequence operations & n/a & \hyperref[sheet:p13-latin-consequence-operations]{\ref*{sheet:p13-latin-consequence-operations}} & \thepage & Distinct from nominal collective predication Original precedence descriptor 'binary' retained; no rank inferred. \\
\phantomsection\label{entry:symbol-registry-07857}\(E_t C_t T_t\) & \texttt{\detokenize{literal indexed propositions}} & Token realization; obtaining content; gated truth & P13 Wilks token comparison & n/a & \hyperref[thm:p13-abelard-token]{\ref*{thm:p13-abelard-token}} & \thepage & Explicit diagnostic rather than full historical truth definition Original precedence descriptor 'atomic' retained; no rank inferred. \\
\phantomsection\label{entry:symbol-registry-07858}\(A I E O\) & \texttt{\detokenize{literal form labels}} & Existence-loaded affirmative and extinctive oppositions & P13 Abelard opposition & n/a & \hyperref[sec:p13-abelard-opposition]{\ref*{sec:p13-abelard-opposition}} & \thepage & Not proposition variables of hypothetical taxonomy Original precedence descriptor 'form' retained; no rank inferred. \\
\phantomsection\label{entry:symbol-registry-07859}\(O_sep\) & \texttt{\detokenize{O_{\mathrm{sep}}}} & Declared subject-existence particular negative & P13 Abelard extension comparison & n/a & \hyperref[sec:p13-abelard-opposition]{\ref*{sec:p13-abelard-opposition}} & \thepage & Does not convert unrestrictedly Original precedence descriptor 'form' retained; no rank inferred. \\
\phantomsection\label{entry:symbol-registry-07860}\(S P\) & \texttt{\detokenize{literal extensions}} & Subject and predicate ranges & P13 opposition and syllogistic comparisons & n/a & \hyperref[sec:p13-abelard-opposition]{\ref*{sec:p13-abelard-opposition}} & \thepage & Empty classes and existence conditions explicit Original precedence descriptor 'set' retained; no rank inferred. \\
\phantomsection\label{entry:symbol-registry-07861}\(sim\) & \texttt{\detokenize{\sim}} & Extinctive negation or declared classical diagnostic negation & P13 Abelard local models & n/a & \hyperref[sec:p13-abelard-opposition]{\ref*{sec:p13-abelard-opposition}} & \thepage & Does not identify separative and extinctive original operations Original precedence descriptor 'unary' retained; no rank inferred. \\
\phantomsection\label{entry:symbol-registry-07862}\(Box Diamond\) & \texttt{\detokenize{\Box;\Diamond}} & Modern world-operator comparison for source modal cases & P13 Wilks modal diagnostic & n/a & \hyperref[thm:p13-abelard-modal]{\ref*{thm:p13-abelard-modal}} & \thepage & Not original Abelard symbolic notation or full nature semantics Original precedence descriptor 'unary' retained; no rank inferred. \\
\phantomsection\label{entry:symbol-registry-07863}\(W W_C\) & \texttt{\detokenize{W;W_C}} & Global and condition-restricted state ranges & P13 modal and nesting diagnostic & n/a & \hyperref[thm:p13-abelard-modal]{\ref*{thm:p13-abelard-modal}} & \thepage & Restriction nonempty where stipulated Original precedence descriptor 'set' retained; no rank inferred. \\
\phantomsection\label{entry:symbol-registry-07864}\(D\) & \texttt{\detokenize{literal domain}} & Constant domain for modal/quantifier comparison & P13 modal comparison & n/a & \hyperref[thm:p13-abelard-modal]{\ref*{thm:p13-abelard-modal}} & \thepage & Distinct from conjunction D in Alberic argument Original precedence descriptor 'domain' retained; no rank inferred. \\
\phantomsection\label{entry:symbol-registry-07865}\(M P S J G E V\) & \texttt{\detokenize{literal predicates or extensions}} & Middle and term extensions for stated syllogistic cases & P13 syllogistic comparisons & n/a & \hyperref[thm:p13-abelard-syllogisms]{\ref*{thm:p13-abelard-syllogisms}} & \thepage & Roles declared per example rather than merged into one ontology Original precedence descriptor 'atomic' retained; no rank inferred. \\
\phantomsection\label{entry:symbol-registry-07866}\(CurrentBoy PastBoy PastYoung CurrentOld\) & \texttt{\detokenize{literal predicate names}} & Distinct time-indexed middle/example extensions & P13 tense countermodel & n/a & \hyperref[annex:p13.abelard.middle]{\ref*{annex:p13.abelard.middle}} & \thepage & Not one unchanged middle Original precedence descriptor 'atomic' retained; no rank inferred. \\
\phantomsection\label{entry:symbol-registry-07867}\(supset\) & \texttt{\detokenize{\supset}} & Accepted conditional judgment in the source's relevance/connexive discussion & P13 topical and school arguments & relation & \hyperref[sec:p13-abelard-topics]{\ref*{sec:p13-abelard-topics}} & \thepage & Not silently material implication \\
\phantomsection\label{entry:symbol-registry-07868}\(to\) & \texttt{\detokenize{\to}} & Material implication in explicitly qualified classical comparisons & P13 diagnostic models & n/a & \hyperref[thm:p13-abelard-middle]{\ref*{thm:p13-abelard-middle}} & \thepage & Distinct from source superset conditional Original precedence descriptor 'binary' retained; no rank inferred. \\
\phantomsection\label{entry:symbol-registry-07869}\(p q X D\) & \texttt{\detokenize{literal formula metavariables}} & Source opposite-topic and Alberic combinations & P13 conditional diagnostic & n/a & \hyperref[thm:p13-abelard-diagnostic]{\ref*{thm:p13-abelard-diagnostic}} & \thepage & D=p and not q in Alberic proof; X generic self-negation candidate Original precedence descriptor 'atomic' retained; no rank inferred. \\
\phantomsection\label{entry:symbol-registry-07870}\(K_W\) & \texttt{\detokenize{K_W}} & All-state material connection used as necessity-only diagnostic & P13 nested-connective model & n/a & \hyperref[sec:p13-abelard-nesting]{\ref*{sec:p13-abelard-nesting}} & \thepage & Omits original containment and connexive conditions Original precedence descriptor 'binary' retained; no rank inferred. \\
\phantomsection\label{entry:symbol-registry-07871}\(D_W\) & \texttt{\detokenize{D_W}} & Stronger diagnostic disjunction via K\_W of negated first input & P13 nested-connective model & n/a & \hyperref[sec:p13-abelard-nesting]{\ref*{sec:p13-abelard-nesting}} & \thepage & Global truth value; ordinary inner disjunction remains pointwise Original precedence descriptor 'binary' retained; no rank inferred. \\
\phantomsection\label{entry:symbol-registry-07872}\(lor\) & \texttt{\detokenize{\lor}} & Source disjunction or explicitly specified diagnostic disjunction & P13 nested and hypothetical cases & n/a & \hyperref[sec:p13-abelard-nesting]{\ref*{sec:p13-abelard-nesting}} & \thepage & Strong sentential and weak predicative roles distinguished Original precedence descriptor 'binary' retained; no rank inferred. \\
\phantomsection\label{entry:symbol-registry-07873}\(P Q\) & \texttt{\detokenize{literal proposition metavariables}} & Arbitrary contradiction and conclusion of explosion chain & P13 Parvipontanian argument & n/a & \hyperref[thm:p13-abelard-explosion]{\ref*{thm:p13-abelard-explosion}} & \thepage & No full propositional calculus admission Original precedence descriptor 'atomic' retained; no rank inferred. \\
\phantomsection\label{entry:symbol-registry-07874}\(alpha beta gamma\) & \texttt{\detokenize{\alpha;\beta;\gamma}} & Signed literals in the exhaustive simple-case proof & P13 hypothetical comparison & n/a & \hyperref[annex:p13.abelard.hypothetical-cases]{\ref*{annex:p13.abelard.hypothetical-cases}} & \thepage & Double-negation normalization is classical and local Original precedence descriptor 'atomic' retained; no rank inferred. \\
\phantomsection\label{entry:symbol-registry-07875}\(sim\) & \texttt{\detokenize{\sim}} & Classical negation in declared modern literal comparisons & P13 early Latin formal comparisons & n/a & \hyperref[thm:p13-early-latin-modes]{\ref*{thm:p13-early-latin-modes}} & \thepage & Not automatically the negation of a historical term connection Original precedence descriptor 'unary' retained; no rank inferred. \\
\phantomsection\label{entry:symbol-registry-07876}\(ampersand\) & \texttt{\detokenize{\mathbin{\&}}} & Conjunction in the displayed source translations & P13 early Latin literal modes & n/a & \hyperref[thm:p13-early-latin-modes]{\ref*{thm:p13-early-latin-modes}} & \thepage & Shared truth valuation at a stated propositional level Original precedence descriptor 'binary' retained; no rank inferred. \\
\phantomsection\label{entry:symbol-registry-07877}\(to\) & \texttt{\detokenize{\to}} & Material implication in the local formal comparison & P13 early Latin formal comparisons & n/a & \hyperref[thm:p13-early-latin-modes]{\ref*{thm:p13-early-latin-modes}} & \thepage & Does not silently define historical si or cum Original precedence descriptor 'binary' retained; no rank inferred. \\
\phantomsection\label{entry:symbol-registry-07878}\(p q\) & \texttt{\detokenize{literal proposition metavariables}} & Whole proposition contents in modern mode comparisons & P13 early Latin literal modes & n/a & \hyperref[thm:p13-early-latin-modes]{\ref*{thm:p13-early-latin-modes}} & \thepage & Source history is distinct from this stipulated propositional interpretation Original precedence descriptor 'atomic' retained; no rank inferred. \\
\phantomsection\label{entry:symbol-registry-07879}\(E N M T C F\) & \texttt{\detokenize{literal proposition abbreviations}} & Existence; numerical oneness; commonness; thought alone; strict correspondence; empty false thought & P13 Boethius universals objection & n/a & \hyperref[sec:p13-early-latin-boethius]{\ref*{sec:p13-early-latin-boethius}} & \thepage & One candidate universal; philosophical premise truth is not proved by the Boolean derivation Original precedence descriptor 'atomic' retained; no rank inferred. \\
\phantomsection\label{entry:symbol-registry-07880}\(Q\) & \texttt{\detokenize{literal proposition abbreviation}} & Prescription to abandon inquiry & P13 universals proof & n/a & \hyperref[annex:p13.early-latin.universals]{\ref*{annex:p13.early-latin.universals}} & \thepage & Requires an additional normative bridge Original precedence descriptor 'atomic' retained; no rank inferred. \\
\phantomsection\label{entry:symbol-registry-07881}\(A B C\) & \texttt{\detokenize{literal extension metavariables}} & Term extensions in the categorical inclusion model & P13 early Latin categorical comparison & n/a & \hyperref[thm:p13-early-latin-categorical]{\ref*{thm:p13-early-latin-categorical}} & \thepage & Distinct from proposition abbreviations C and from historical predicate-letter positions Original precedence descriptor 'set' retained; no rank inferred. \\
\phantomsection\label{entry:symbol-registry-07882}\(A B\) & \texttt{\detokenize{literal predicate metavariables}} & Predicates in the hypothetical term-relation comparison & P13 early Latin hypothetical comparison & n/a & \hyperref[thm:p13-early-latin-contraries]{\ref*{thm:p13-early-latin-contraries}} & \thepage & Historical term letters retained with modern interpretation explicitly stated Original precedence descriptor 'atomic' retained; no rank inferred. \\
\phantomsection\label{entry:symbol-registry-07883}\(a b\) & \texttt{\detokenize{literal object names}} & Fixed subjects and countermodel objects & P13 early Latin term comparisons & n/a & \hyperref[thm:p13-early-latin-contraries]{\ref*{thm:p13-early-latin-contraries}} & \thepage & Site-specific consequence is different from existence of an unspecified witness Original precedence descriptor 'atomic' retained; no rank inferred. \\
\phantomsection\label{entry:symbol-registry-07884}\(D\) & \texttt{\detokenize{literal domain symbol}} & Range of applicability for immediate contraries & P13 hypothetical extension model & n/a & \hyperref[thm:p13-early-latin-contraries]{\ref*{thm:p13-early-latin-contraries}} & \thepage & Exhaustiveness and disjointness are separate conditions Original precedence descriptor 'set' retained; no rank inferred. \\
\phantomsection\label{entry:symbol-registry-07885}\(Part\) & \texttt{\detokenize{\operatorname{Part}}} & Declared whole-part relation & P13 topical warrant comparison & relation & \hyperref[thm:p13-early-latin-topical-warrant]{\ref*{thm:p13-early-latin-topical-warrant}} & \thepage & Part relation alone does not transfer arbitrary properties \\
\phantomsection\label{entry:symbol-registry-07886}\(P\) & \texttt{\detokenize{literal predicate}} & Relevant property in the topical argument & P13 topical warrant comparison & n/a & \hyperref[thm:p13-early-latin-topical-warrant]{\ref*{thm:p13-early-latin-topical-warrant}} & \thepage & The maximal proposition supplies the additional transfer warrant Original precedence descriptor 'atomic' retained; no rank inferred. \\
\phantomsection\label{entry:symbol-registry-07887}\(h w\) & \texttt{\detokenize{literal object names}} & Part and whole of the source-based formal argument & P13 topical warrant comparison & n/a & \hyperref[thm:p13-early-latin-topical-warrant]{\ref*{thm:p13-early-latin-topical-warrant}} & \thepage & Not metaphysical existence theorems about providence Original precedence descriptor 'atomic' retained; no rank inferred. \\
\phantomsection\label{entry:symbol-registry-07888}\(Box Diamond\) & \texttt{\detokenize{\Box;\Diamond}} & Necessity and possibility in the explicit modern frame comparison & P13 early Latin modal scope & n/a & \hyperref[thm:p13-early-latin-modal-scope]{\ref*{thm:p13-early-latin-modal-scope}} & \thepage & Not original Damian or Anselm symbolic notation Original precedence descriptor 'unary' retained; no rank inferred. \\
\phantomsection\label{entry:symbol-registry-07889}\(W R\) & \texttt{\detokenize{literal frame symbols}} & Worlds and accessibility relation & P13 early Latin modal proof & n/a & \hyperref[annex:p13.early-latin.modal-scope]{\ref*{annex:p13.early-latin.modal-scope}} & \thepage & Reflexivity is needed for actual detachment from a boxed conditional Original precedence descriptor 'frame' retained; no rank inferred. \\
\phantomsection\label{entry:symbol-registry-07890}\(F P\) & \texttt{\detokenize{literal proposition metavariables}} & Earlier condition and event in modal dependency comparison & P13 prescience scope model & n/a & \hyperref[thm:p13-early-latin-modal-scope]{\ref*{thm:p13-early-latin-modal-scope}} & \thepage & Distinct from universals false-thought F and topical predicate P Original precedence descriptor 'atomic' retained; no rank inferred. \\
\phantomsection\label{entry:symbol-registry-07891}\(G M S\) & \texttt{\detokenize{literal predicates}} & Grammaticus; man; substance in the declared categorical comparison & P13 Anselm denominatives & n/a & \hyperref[thm:p13-early-latin-grammaticus]{\ref*{thm:p13-early-latin-grammaticus}} & \thepage & M here is not the universals commonness proposition Original precedence descriptor 'atomic' retained; no rank inferred. \\
\phantomsection\label{entry:symbol-registry-07892}\(t q s\) & \texttt{\detokenize{literal model object names}} & Word; quality; substance & P13 typed semantic coexistence model & n/a & \hyperref[annex:p13.early-latin.grammaticus]{\ref*{annex:p13.early-latin.grammaticus}} & \thepage & Relations connect distinct semantic types Original precedence descriptor 'atomic' retained; no rank inferred. \\
\phantomsection\label{entry:symbol-registry-07893}\(signifies appellates\) & \texttt{\detokenize{literal relation names}} & Distinct relations from words to their significates and appellates & P13 Anselm semantic-level model & relation & \hyperref[thm:p13-early-latin-grammaticus]{\ref*{thm:p13-early-latin-grammaticus}} & \thepage & A compatibility witness rather than a complete formal semantics of Latin \\
\phantomsection\label{entry:symbol-registry-07894}\(P H C B D SC SH\) & \texttt{\detokenize{literal catalogue prefixes}} & Source commentary-family identifiers & P13 turn-of-twelfth-century source corpus & n/a & \hyperref[sec:p13-early-latin-turn-twelfth]{\ref*{sec:p13-early-latin-turn-twelfth}} & \thepage & Not logical variables; manuscript literal labels and attribution uncertainty retained Original precedence descriptor 'identifier' retained; no rank inferred. \\
\end{longtable}
\endgroup
\Needspace{.4\textheight}
\subsection*{Stable part P14 scopes}
\begingroup\scriptsize\setlength{\tabcolsep}{3pt}
\begin{longtable}{@{}p{0.9in}p{1.3in}p{1.0in}p{0.5in}p{0.3in}p{0.4in}p{0.35in}p{1.65in}@{}}
\toprule
\textbf{Glyph} & \textbf{Macro} & \textbf{Role} & \textbf{Scope} & \textbf{Prec./form} & \textbf{First} & \textbf{Reg. page} & \textbf{Principal uses / notes} \\
\midrule
\endfirsthead
\toprule
\textbf{Glyph} & \textbf{Macro} & \textbf{Role} & \textbf{Scope} & \textbf{Prec./form} & \textbf{First} & \textbf{Reg. page} & \textbf{Principal uses / notes} \\
\midrule
\endhead
\midrule\multicolumn{8}{r}{\emph{continued}}\\\endfoot
\bottomrule\endlastfoot
\phantomsection\label{entry:symbol-registry-02546}\(W\) & \texttt{\detokenize{W}} & finite outcome space & P14 finite event-weight comparison & n/a & \hyperref[def:p14-evidence-uncertainty-data]{\ref*{def:p14-evidence-uncertainty-data}} & \thepage & Finite and nonempty; the same event space underlies each comparison. \\
\phantomsection\label{entry:symbol-registry-02547}\(A,B\) & \texttt{\detokenize{A,}\allowbreak\detokenize{B}} & events in the finite outcome space & P14 finite event-weight comparison & n/a & \hyperref[def:p14-evidence-uncertainty-data]{\ref*{def:p14-evidence-uncertainty-data}} & \thepage & Subsets of the declared outcome space; subset masses also use the bound subset letter B. \\
\phantomsection\label{entry:symbol-registry-02548}\(A^c\) & \texttt{\detokenize{A^c}} & event complement & P14 finite event-weight comparison & n/a & \hyperref[def:p14-evidence-uncertainty-data]{\ref*{def:p14-evidence-uncertainty-data}} & \thepage & Complement within the declared finite outcome space. \\
\phantomsection\label{entry:symbol-registry-02549}\(p\) & \texttt{\detokenize{p}} & probability mass on outcomes & P14 finite event-weight comparison & n/a & \hyperref[def:p14-evidence-uncertainty-data]{\ref*{def:p14-evidence-uncertainty-data}} & \thepage & Nonnegative outcome weights with total one. \\
\phantomsection\label{entry:symbol-registry-02550}\(P_p\) & \texttt{\detokenize{P_p}} & event probability induced by a mass & P14 finite event-weight comparison & n/a & \hyperref[def:p14-evidence-uncertainty-data]{\ref*{def:p14-evidence-uncertainty-data}} & \thepage & Adds the mass on members of an event. \\
\phantomsection\label{entry:symbol-registry-02551}\(\mathcal C\) & \texttt{\detokenize{\mathcal C}} & nonempty family of probability masses & P14 finite event-weight comparison & n/a & \hyperref[def:p14-evidence-uncertainty-data]{\ref*{def:p14-evidence-uncertainty-data}} & \thepage & No distribution over the family is implied. \\
\phantomsection\label{entry:symbol-registry-02552}\(\underline P_{\mathcal C},\overline P_{\mathcal C}\) & \texttt{\detokenize{\underline P_{\mathcal C},}\allowbreak\detokenize{\overline P_{\mathcal C}}} & lower and upper event envelopes & P14 finite event-weight comparison & n/a & \hyperref[def:p14-evidence-uncertainty-data]{\ref*{def:p14-evidence-uncertainty-data}} & \thepage & Infimum and supremum across the retained family. \\
\phantomsection\label{entry:symbol-registry-02553}\(\pi\) & \texttt{\detokenize{\pi}} & normalized singleton possibility scores & P14 finite event-weight comparison & n/a & \hyperref[def:p14-evidence-uncertainty-data]{\ref*{def:p14-evidence-uncertainty-data}} & \thepage & Values lie in the unit interval and their maximum is one. \\
\phantomsection\label{entry:symbol-registry-02554}\(\operatorname{Poss}_{\pi}\) & \texttt{\detokenize{\operatorname{Poss}_{\pi}}} & finite possibility measure & P14 finite event-weight comparison & n/a & \hyperref[def:p14-evidence-uncertainty-data]{\ref*{def:p14-evidence-uncertainty-data}} & \thepage & Maximum over the event; empty-event value is zero. \\
\phantomsection\label{entry:symbol-registry-02555}\(\operatorname{Nec}_{\pi}\) & \texttt{\detokenize{\operatorname{Nec}_{\pi}}} & dual finite necessity measure & P14 finite event-weight comparison & n/a & \hyperref[def:p14-evidence-uncertainty-data]{\ref*{def:p14-evidence-uncertainty-data}} & \thepage & One minus possibility of the complement. \\
\phantomsection\label{entry:symbol-registry-02556}\(m\) & \texttt{\detokenize{m}} & normalized mass assignment to subsets & P14 finite event-weight comparison & n/a & \hyperref[def:p14-evidence-uncertainty-data]{\ref*{def:p14-evidence-uncertainty-data}} & \thepage & Empty subset has zero mass; masses over all subsets total one. \\
\phantomsection\label{entry:symbol-registry-02557}\(\operatorname{Bel}_m,\operatorname{Pl}_m\) & \texttt{\detokenize{\operatorname{Bel}_m,}\allowbreak\detokenize{\operatorname{Pl}_m}} & mass-based belief and plausibility functions & P14 finite event-weight comparison & n/a & \hyperref[def:p14-evidence-uncertainty-data]{\ref*{def:p14-evidence-uncertainty-data}} & \thepage & Containment and nonempty-intersection sums over subset masses. \\
\phantomsection\label{entry:symbol-registry-02558}\(q_B\) & \texttt{\detokenize{q_B}} & probability allocation within a positive-mass subset & P14 finite event-weight comparison & n/a & \hyperref[def:p14-evidence-uncertainty-data]{\ref*{def:p14-evidence-uncertainty-data}} & \thepage & A probability mass on the subset extended by zero outside it. \\
\phantomsection\label{entry:symbol-registry-02559}\(\mathcal C_m\) & \texttt{\detokenize{\mathcal C_m}} & family of masses induced by subset allocations & P14 finite event-weight comparison & n/a & \hyperref[def:p14-evidence-uncertainty-data]{\ref*{def:p14-evidence-uncertainty-data}} & \thepage & Every permitted allocation contributes its mixture of subset probabilities. \\
\phantomsection\label{entry:symbol-registry-02952}\(D\) & \texttt{\detokenize{D}} & finite first-order domain & P14 empty-predicate countermodel & n/a & \hyperref[sec:p14-counter-empty-predicate]{\ref*{sec:p14-counter-empty-predicate}} & \thepage & The stipulated singleton universe; the predicate extension can still be empty. \\
\phantomsection\label{entry:symbol-registry-02953}\(W\) & \texttt{\detokenize{W}} & finite modal world set & P14 modal countermodels & n/a & \hyperref[sec:p14-counter-modal]{\ref*{sec:p14-counter-modal}} & \thepage & Each row declares its complete world set. \\
\phantomsection\label{entry:symbol-registry-02954}\(R\) & \texttt{\detokenize{R}} & modal accessibility relation & P14 modal countermodels & n/a & \hyperref[sec:p14-counter-modal]{\ref*{sec:p14-counter-modal}} & \thepage & Exactly the listed ordered pairs; omitted edges are absent. \\
\phantomsection\label{entry:symbol-registry-02955}\(V(p)\) & \texttt{\detokenize{V(p)}} & extension of the modal atom at worlds & P14 modal countermodels & n/a & \hyperref[sec:p14-counter-modal]{\ref*{sec:p14-counter-modal}} & \thepage & The worlds where the declared propositional atom is true. \\
\phantomsection\label{entry:symbol-registry-02956}\(W,R,V(p)\) & \texttt{\detokenize{W,}\allowbreak\detokenize{R,}\allowbreak\detokenize{V(p)}} & finite local-necessitation model data & P14 local necessitation countermodel & n/a & \hyperref[sec:p14-counter-local-necessitation]{\ref*{sec:p14-counter-local-necessitation}} & \thepage & The two-world full relation and one-world atomic extension. \\
\phantomsection\label{entry:symbol-registry-02957}\(K\) & \texttt{\detokenize{K}} & self-application domain set & P14 diagonal domain countermodel & n/a & \hyperref[sec:p14-counter-decidability]{\ref*{sec:p14-counter-decidability}} & \thepage & Defined output retains the referenced program numbering's accepting-with-valid-output convention. \\
\phantomsection\label{entry:symbol-registry-02958}\(h\) & \texttt{\detokenize{h}} & purported total characteristic function of the domain set & P14 diagonal domain countermodel & n/a & \hyperref[sec:p14-counter-decidability]{\ref*{sec:p14-counter-decidability}} & \thepage & One means membership; both possible answers at the diagonal index contradict the assumption. \\
\phantomsection\label{entry:symbol-registry-02959}\(d\) & \texttt{\detokenize{d}} & partial diagonal procedure & P14 diagonal domain countermodel & n/a & \hyperref[sec:p14-counter-decidability]{\ref*{sec:p14-counter-decidability}} & \thepage & Returns zero on a negative characteristic answer and diverges otherwise. \\
\phantomsection\label{entry:symbol-registry-02960}\(P\) & \texttt{\detokenize{P}} & probability measure in a stipulated counterexample & P14 probability countermodels & n/a & \hyperref[sec:p14-counter-independence]{\ref*{sec:p14-counter-independence}} & \thepage & Each numerical example states its own probability model; no shared physical population law is implied. \\
\phantomsection\label{entry:symbol-registry-02961}\(A,B,C\) & \texttt{\detokenize{A,}\allowbreak\detokenize{B,}\allowbreak\detokenize{C}} & binary-coordinate and agreement events & P14 independence countermodel & n/a & \hyperref[sec:p14-counter-independence]{\ref*{sec:p14-counter-independence}} & \thepage & First coordinate one; second coordinate one; the two coordinates agree. \\
\phantomsection\label{entry:symbol-registry-02962}\(C(0),C(1)\) & \texttt{\detokenize{C(0),}\allowbreak\detokenize{C(1)}} & credible sets indexed by the Bernoulli observation & P14 credible coverage countermodel & n/a & \hyperref[sec:p14-counter-coverage]{\ref*{sec:p14-counter-coverage}} & \thepage & Equal-tail posterior intervals after the stated zero or one outcome. \\
\phantomsection\label{entry:symbol-registry-02963}\(K\) & \texttt{\detokenize{K}} & number of independent candidate feature columns & P14 selected predictor countermodel & n/a & \hyperref[sec:p14-counter-selection]{\ref*{sec:p14-counter-selection}} & \thepage & A positive integer candidate count; distinct from the program-domain set. \\
\phantomsection\label{entry:symbol-registry-02964}\(h_j\) & \texttt{\detokenize{h_j}} & feature-coordinate predictor & P14 selected predictor countermodel & n/a & \hyperref[sec:p14-counter-selection]{\ref*{sec:p14-counter-selection}} & \thepage & Predictor j reports its fair binary feature against the constant-zero label. \\
\phantomsection\label{entry:symbol-registry-02965}\(X,Y,Z\) & \texttt{\detokenize{X,}\allowbreak\detokenize{Y,}\allowbreak\detokenize{Z}} & binary endogenous random variables & P14 causal countermodels & n/a & \hyperref[sec:p14-counter-causal]{\ref*{sec:p14-counter-causal}} & \thepage & The explicitly given acyclic structural equations determine them. \\
\phantomsection\label{entry:symbol-registry-02966}\(U_Z,U_X,U_Y\) & \texttt{\detokenize{U_Z,}\allowbreak\detokenize{U_X,}\allowbreak\detokenize{U_Y}} & mutually independent binary structural disturbances & P14 causal countermodels & n/a & \hyperref[sec:p14-counter-causal]{\ref*{sec:p14-counter-causal}} & \thepage & The first bit is fair; each other bit is one with probability one quarter in the common-cause model. \\
\phantomsection\label{entry:symbol-registry-02967}\(D\) & \texttt{\detokenize{D}} & independent binary disturbance in the comparison causal model & P14 causal countermodels & n/a & \hyperref[sec:p14-counter-causal]{\ref*{sec:p14-counter-causal}} & \thepage & The second model has probability three eighths of disturbance one. \\
\phantomsection\label{entry:symbol-registry-02968}\(X,Y\) & \texttt{\detokenize{X,}\allowbreak\detokenize{Y}} & input and squared-output random variables & P14 nonlinear uncertainty countermodel & n/a & \hyperref[sec:p14-counter-linearization]{\ref*{sec:p14-counter-linearization}} & \thepage & The input is uniform on the specified symmetric interval; the output is its square. \\
\phantomsection\label{entry:symbol-registry-02969}\(p,q\) & \texttt{\detokenize{p,}\allowbreak\detokenize{q}} & extreme lottery and certain-middle lottery & P14 lottery utility countermodels & n/a & \hyperref[sec:p14-counter-utility]{\ref*{sec:p14-counter-utility}} & \thepage & These are probability lotteries over the three prizes; distinct from propositional atoms. \\
\phantomsection\label{entry:symbol-registry-02970}\(\ell,c,h\) & \texttt{\detokenize{\ell,}\allowbreak\detokenize{c,}\allowbreak\detokenize{h}} & ordered sure consequences & P14 lottery utility countermodels & n/a & \hyperref[sec:p14-counter-utility]{\ref*{sec:p14-counter-utility}} & \thepage & The same ordinal order supports two different calibrated lottery preferences. \\
\phantomsection\label{entry:symbol-registry-02971}\(u\) & \texttt{\detokenize{u}} & consequence utility on three prizes & P14 lottery utility countermodels & n/a & \hyperref[sec:p14-counter-utility]{\ref*{sec:p14-counter-utility}} & \thepage & The middle utility is supplied as two fifths or three fifths in the two comparisons. \\
\phantomsection\label{entry:symbol-registry-02982}\(\Box,\Diamond\) & \texttt{\detokenize{\Box,}\allowbreak\detokenize{\Diamond}} & normal modal necessity and possibility constructors & P14 proof translation comparisons & n/a & \hyperref[sec:p14-translation-standard-interface]{\ref*{sec:p14-translation-standard-interface}} & \thepage & The model regime and Boolean syntax are stated before the standard and IPL translations. \\
\phantomsection\label{entry:symbol-registry-02983}\(\operatorname{ST}_x\) & \texttt{\detokenize{\operatorname{ST}_x}} & modal-to-first-order formula translation at a world variable & P14 standard translation interface & n/a & \hyperref[sec:p14-translation-standard-interface]{\ref*{sec:p14-translation-standard-interface}} & \thepage & The linked exact recursion uses fresh object variables and unary valuation predicates. \\
\phantomsection\label{entry:symbol-registry-02984}\(P_p\) & \texttt{\detokenize{P_p}} & unary first-order predicate encoding a modal atom & P14 standard translation interface & n/a & \hyperref[sec:p14-translation-standard-interface]{\ref*{sec:p14-translation-standard-interface}} & \thepage & Distinct predicate symbols correspond to distinct propositional atoms. \\
\phantomsection\label{entry:symbol-registry-02985}\(R\) & \texttt{\detokenize{R}} & binary accessibility relation in the paired model presentations & P14 standard translation interface & n/a & \hyperref[sec:p14-translation-standard-interface]{\ref*{sec:p14-translation-standard-interface}} & \thepage & The modal edge relation and its first-order interpretation are the same subset of the squared domain. \\
\phantomsection\label{entry:symbol-registry-02986}\(\mathfrak A_{\mathcal M}\) & \texttt{\detokenize{\mathfrak A_{\mathcal M}}} & first-order encoding of a nonempty Kripke structure & P14 standard translation interface & n/a & \hyperref[def:p14-standard-translation-models]{\ref*{def:p14-standard-translation-models}} & \thepage & Shares the world domain and accessibility; unary predicate extensions are the atomic valuation. \\
\phantomsection\label{entry:symbol-registry-02987}\(\mathcal M_{\mathfrak A}\) & \texttt{\detokenize{\mathcal M_{\mathfrak A}}} & Kripke decoding of a nonempty first-order frame structure & P14 standard translation interface & n/a & \hyperref[def:p14-standard-translation-models]{\ref*{def:p14-standard-translation-models}} & \thepage & Uses the declared binary relation and unary valuation predicates. \\
\phantomsection\label{entry:symbol-registry-02988}\(\operatorname{ST}_x[\Gamma]\) & \texttt{\detokenize{\operatorname{ST}_x[\Gamma]}} & memberwise translated premise set & P14 standard translation interface & n/a & \hyperref[def:p14-standard-translation-models]{\ref*{def:p14-standard-translation-models}} & \thepage & All translated premises use the same displayed free world variable. \\
\phantomsection\label{entry:symbol-registry-02989}\(A^N,A^{\mathrm{Ko}},A^{\mathrm{Ku}},A^{\Box}\) & \texttt{\detokenize{A^N,}\allowbreak\detokenize{A^{\mathrm{Ko}},}\allowbreak\detokenize{A^{\mathrm{Ku}},}\allowbreak\detokenize{A^{\Box}}} & formula images under the separately defined proof translations & P14 proof translation comparisons & n/a & \hyperref[sec:p14-translation-embeddings]{\ref*{sec:p14-translation-embeddings}} & \thepage & These denote formulas; the referenced source definitions determine each recursion. \\
\phantomsection\label{entry:symbol-registry-02990}\(\Gamma^N,\Gamma^{\mathrm{Ko}},\Gamma^{\mathrm{Ku}}\) & \texttt{\detokenize{\Gamma^N,}\allowbreak\detokenize{\Gamma^{\mathrm{Ko}},}\allowbreak\detokenize{\Gamma^{\mathrm{Ku}}}} & memberwise translated finite formula contexts & P14 proof translation comparisons & n/a & \hyperref[sec:p14-translation-embeddings]{\ref*{sec:p14-translation-embeddings}} & \thepage & Glivenko's exact input-query row instead retains its untranslated premise context. \\
\phantomsection\label{entry:symbol-registry-02991}\(\widehat a\) & \texttt{\detokenize{\widehat a}} & Stone clopen associated with a Boolean element & P14 algebraic representation comparisons & n/a & \hyperref[sec:p14-translation-frame-algebra]{\ref*{sec:p14-translation-frame-algebra}} & \thepage & The ultrafilters containing the given element; not every subset is an image clopen. \\
\phantomsection\label{entry:symbol-registry-02992}\(\beta_B\) & \texttt{\detokenize{\beta_B}} & Boolean or Boolean-operator Stone embedding & P14 algebraic representation comparisons & n/a & \hyperref[sec:p14-translation-frame-algebra]{\ref*{sec:p14-translation-frame-algebra}} & \thepage & The row specifies whether the target is the clopen algebra or the full canonical complex algebra. \\
\phantomsection\label{entry:symbol-registry-02993}\(\Box U,\Diamond U\) & \texttt{\detokenize{\Box U,}\allowbreak\detokenize{\Diamond U}} & values of the lifted modal operators on subsets & P14 algebraic representation comparisons & n/a & \hyperref[sec:p14-translation-frame-algebra]{\ref*{sec:p14-translation-frame-algebra}} & \thepage & The operand is a subset of the world carrier; these are the explicitly defined complex-algebra interpretations of the source modal operators. \\
\phantomsection\label{entry:symbol-registry-02994}\(X_B\) & \texttt{\detokenize{X_B}} & Stone space of Boolean ultrafilters & P14 algebraic representation comparisons & n/a & \hyperref[sec:p14-translation-frame-algebra]{\ref*{sec:p14-translation-frame-algebra}} & \thepage & The source representation includes trivial algebras with empty dual space. \\
\phantomsection\label{entry:symbol-registry-02995}\(\operatorname{Clop}\) & \texttt{\detokenize{\operatorname{Clop}}} & clopen-algebra construction and inverse-image functor & P14 algebraic representation comparisons & n/a & \hyperref[sec:p14-translation-frame-algebra]{\ref*{sec:p14-translation-frame-algebra}} & \thepage & On continuous maps it reverses direction. \\
\phantomsection\label{entry:symbol-registry-02996}\(D_{\alpha\to_{\mathrm{ty}}\beta}\) & \texttt{\detokenize{D_{\alpha\to_{\mathrm{ty}}\beta}}} & arrow-type function domain & P14 model regime comparisons & n/a & \hyperref[sec:p14-translation-model-regimes]{\ref*{sec:p14-translation-model-regimes}} & \thepage & Full function space in the standard row; a lambda-closed selected subset in the Henkin row. \\
\phantomsection\label{entry:symbol-registry-02997}\(S\) & \texttt{\detokenize{S}} & bimeasurable bijection of probability observation coordinates & P14 uncertainty representation comparisons & n/a & \hyperref[sec:p14-translation-definability-invariance]{\ref*{sec:p14-translation-definability-invariance}} & \thepage & The law and reference measure are transported with the likelihood; a new flat prior is a different operation. \\
\phantomsection\label{entry:symbol-registry-02998}\(\mathcal C_m\) & \texttt{\detokenize{\mathcal C_m}} & family of additive allocations from a subset mass assignment & P14 uncertainty representation comparisons & n/a & \hyperref[sec:p14-translation-definability-invariance]{\ref*{sec:p14-translation-definability-invariance}} & \thepage & The exact construction and attained event envelopes remain in the linked Chapter 287 proof. \\
\phantomsection\label{entry:symbol-registry-02999}\(\mathcal C\) & \texttt{\detokenize{\mathcal C}} & finite retained prior list & P14 uncertainty representation comparisons & n/a & \hyperref[sec:p14-translation-definability-invariance]{\ref*{sec:p14-translation-definability-invariance}} & \thepage & Taking its convex hull preserves minimum expected utility of every supplied act. \\
\phantomsection\label{entry:symbol-registry-03000}\(M\) & \texttt{\detokenize{M}} & two-element object domain in the selected/full comparison & P14 selected reachability comparison & n/a & \hyperref[def:p14-selected-reachability-model]{\ref*{def:p14-selected-reachability-model}} & \thepage & The same domain is used in both structures. \\
\phantomsection\label{entry:symbol-registry-03001}\(\mathcal R_n\) & \texttt{\detokenize{\mathcal R_n}} & explicit selected relation domain at each arity & P14 selected reachability comparison & n/a & \hyperref[def:p14-selected-reachability-model]{\ref*{def:p14-selected-reachability-model}} & \thepage & Unary domain contains only the full unary relation; the binary domain also admits the empty edge relation. \\
\phantomsection\label{entry:symbol-registry-03002}\(\mathcal H,\mathcal F\) & \texttt{\detokenize{\mathcal H,}\allowbreak\detokenize{\mathcal F}} & selected-domain and full-domain second-order structures & P14 selected reachability comparison & n/a & \hyperref[def:p14-selected-reachability-model]{\ref*{def:p14-selected-reachability-model}} & \thepage & The same object constants and empty graph relation are retained. \\
\phantomsection\label{entry:symbol-registry-03003}\(a,b\) & \texttt{\detokenize{a,}\allowbreak\detokenize{b}} & object constants in the selected/full reachability comparison & P14 selected reachability comparison & n/a & \hyperref[def:p14-selected-reachability-model]{\ref*{def:p14-selected-reachability-model}} & \thepage & Interpretations are zero and one respectively. \\
\phantomsection\label{entry:symbol-registry-03004}\(E\) & \texttt{\detokenize{E}} & binary predicate constant for the empty directed graph & P14 selected reachability comparison & n/a & \hyperref[def:p14-selected-reachability-model]{\ref*{def:p14-selected-reachability-model}} & \thepage & Its empty extension belongs to the binary selected relation domain. \\
\phantomsection\label{entry:symbol-registry-03227}\(L_{\mathrm{bit}}\) & \texttt{\detokenize{L_{\mathrm{bit}}}} & total encoded input bit length & P14 decidability inputs & n/a & \hyperref[def:p14-dc-instance]{\ref*{def:p14-dc-instance}} & \thepage & Syntax nodes and unit-cost arithmetic operations are separate measures. \\
\phantomsection\label{entry:symbol-registry-03228}\(n\) & \texttt{\detokenize{n}} & number of circuit gates & P14 Horn reduction & n/a & \hyperref[annex:p14.dc.horn]{\ref*{annex:p14.dc.horn}} & \thepage & Topological list positions serve as normalized gate indices; distinct from formula lengths in cited table rows. \\
\phantomsection\label{entry:symbol-registry-03229}\(p_i^+\) & \texttt{\detokenize{p_i^+}} & true-value rail atom for gate i & P14 Horn reduction & n/a & \hyperref[annex:p14.dc.horn]{\ref*{annex:p14.dc.horn}} & \thepage & A fresh Horn propositional atom; membership in the least closure represents the computed gate value. \\
\phantomsection\label{entry:symbol-registry-03230}\(p_i^-\) & \texttt{\detokenize{p_i^-}} & false-value rail atom for gate i & P14 Horn reduction & n/a & \hyperref[annex:p14.dc.horn]{\ref*{annex:p14.dc.horn}} & \thepage & Not a formula negation; the final negative unit forbids this output rail. \\
\phantomsection\label{entry:symbol-registry-03231}\(H^+\) & \texttt{\detokenize{H^+}} & positive Horn clause family of the circuit reduction & P14 Horn reduction & n/a & \hyperref[annex:p14.dc.horn]{\ref*{annex:p14.dc.horn}} & \thepage & Excludes the single negative output constraint. \\
\phantomsection\label{entry:symbol-registry-03232}\(S_*\) & \texttt{\detokenize{S_*}} & least forward closure of the positive Horn clauses & P14 Horn reduction & n/a & \hyperref[annex:p14.dc.horn]{\ref*{annex:p14.dc.horn}} & \thepage & Contained in every model; itself the least positive-clause model. \\
\phantomsection\label{entry:symbol-registry-03233}\(o\) & \texttt{\detokenize{o}} & designated circuit output gate index & P14 Horn reduction & n/a & \hyperref[annex:p14.dc.horn]{\ref*{annex:p14.dc.horn}} & \thepage & The emitted negative unit is the negation of the false rail p\_o\^{}-. \\
\phantomsection\label{entry:symbol-registry-03303}\(\operatorname{Con}_{\mathrm{ext}}(T)\) & \texttt{\detokenize{\operatorname{Con}_{\mathrm{ext}}}} & external proof-calculus-relative consistency assertion & P14 strength comparisons & n/a & \hyperref[def:p14-strength-relations]{\ref*{def:p14-strength-relations}} & \thepage & No finite derivation of falsity; distinct from the chosen internal arithmetic consistency formula. \\
\phantomsection\label{entry:symbol-registry-03304}\(\mathcal C\) & \texttt{\detokenize{\mathcal C}} & old-language sentence class for conservation & P14 strength comparisons & n/a & \hyperref[def:p14-strength-relations]{\ref*{def:p14-strength-relations}} & \thepage & The class restriction is part of the comparison and cannot be dropped. \\
\phantomsection\label{entry:symbol-registry-03305}\(\mathsf B\) & \texttt{\detokenize{\mathsf B}} & formal arithmetic base for a coded proof-transfer implication & P14 strength comparisons & n/a & \hyperref[thm:p14-strength-proof-transfer]{\ref*{thm:p14-strength-proof-transfer}} & \thepage & First-order reasoning proves the consistency implication from separately proved totality and contradiction-preservation formulas. \\
\phantomsection\label{entry:symbol-registry-03306}\(P_U(p,x),P_V(q,x)\) & \texttt{\detokenize{P_U;P_V}} & selected proof-predicate formulas & P14 strength comparisons & n/a & \hyperref[thm:p14-strength-proof-transfer]{\ref*{thm:p14-strength-proof-transfer}} & \thepage & The second argument is the code of the conclusion in that presentation. \\
\phantomsection\label{entry:symbol-registry-03307}\(c_U,c_V\) & \texttt{\detokenize{c_U;c_V}} & respective numerical falsity codes & P14 strength comparisons & n/a & \hyperref[thm:p14-strength-proof-transfer]{\ref*{thm:p14-strength-proof-transfer}} & \thepage & Bars denote their fixed standard numerals; source and target codes need not be equal. \\
\phantomsection\label{entry:symbol-registry-03308}\(F(p,q)\) & \texttt{\detokenize{F}} & formula coding the proposed proof-transformation graph & P14 strength comparisons & n/a & \hyperref[thm:p14-strength-proof-transfer]{\ref*{thm:p14-strength-proof-transfer}} & \thepage & Totality and contradiction preservation are premises proved in the named base; output uniqueness is not needed by the transfer lemma. \\
\phantomsection\label{entry:symbol-registry-03507}\(\Gamma,\Delta\) & \texttt{\detokenize{\Gamma,}\allowbreak\detokenize{\Delta}} & formula-premise sets or exact sheet-indexed resource contexts & P14 nonclassical modal and resource sheet comparisons & n/a & \hyperref[sec:p14-nonclassical-and-specialized-calculi]{\ref*{sec:p14-nonclassical-and-specialized-calculi}} & \thepage & The opening declares the family-specific type; ordered sequences and multisets are not identified. \\
\phantomsection\label{entry:symbol-registry-03508}\(\Gamma,\Phi\) & \texttt{\detokenize{\Gamma,}\allowbreak\detokenize{\Phi}} & distinct sorted-variable context and finite formula-hypothesis list & P14 nonclassical categorical sheet comparisons & n/a & \hyperref[sec:p14-nonclassical-specialized]{\ref*{sec:p14-nonclassical-specialized}} & \thepage & Pitts's source roles are separate from the modal-premise and resource-context roles. \\
\phantomsection\label{entry:symbol-registry-03509}\(\Box\) & \texttt{\detokenize{\Box}} & primitive normal modal formula constructor & P14 nonclassical modal and conditional sheet comparisons & n/a & \hyperref[sec:p14-nonclassical-modal-temporal]{\ref*{sec:p14-nonclassical-modal-temporal}} & \thepage & In relevance comparisons Box instead retains the source eliminable implication abbreviation. \\
\phantomsection\label{entry:symbol-registry-03510}\(\Box\) & \texttt{\detokenize{\Box}} & relevant-entailment formula abbreviation & P14 nonclassical relevance sheet comparisons & n/a & \hyperref[sec:p14-nonclassical-matrices-relevance]{\ref*{sec:p14-nonclassical-matrices-relevance}} & \thepage & The chapter explicitly distinguishes this from the primitive normal-modal box and links the exact E sheet. \\
\end{longtable}
\endgroup
\Needspace{.4\textheight}
\subsection*{Stable part P15 scopes}
\begingroup\scriptsize\setlength{\tabcolsep}{3pt}
\begin{longtable}{@{}p{0.9in}p{1.3in}p{1.0in}p{0.5in}p{0.3in}p{0.4in}p{0.35in}p{1.65in}@{}}
\toprule
\textbf{Glyph} & \textbf{Macro} & \textbf{Role} & \textbf{Scope} & \textbf{Prec./form} & \textbf{First} & \textbf{Reg. page} & \textbf{Principal uses / notes} \\
\midrule
\endfirsthead
\toprule
\textbf{Glyph} & \textbf{Macro} & \textbf{Role} & \textbf{Scope} & \textbf{Prec./form} & \textbf{First} & \textbf{Reg. page} & \textbf{Principal uses / notes} \\
\midrule
\endhead
\midrule\multicolumn{8}{r}{\emph{continued}}\\\endfoot
\bottomrule\endlastfoot
\phantomsection\label{entry:symbol-registry-00923}\(\Omega,\omega,2^\Omega\) & \texttt{\detokenize{\Omega,}\allowbreak\detokenize{\omega,}\allowbreak\detokenize{2^\Omega}} & finite outcome space outcome and event set & P15 finite-probability & meta & \hyperref[sec:p15-finite-probability]{\ref*{sec:p15-finite-probability}} & \thepage & The outcome space is finite and nonempty, and all its subsets are events. This first probability unit uses the full power set rather than a general sigma-algebra. \\
\phantomsection\label{entry:symbol-registry-00924}\(p,\mathbb P,A^c\) & \texttt{\detokenize{p,}\allowbreak\detokenize{\mathbb P,}\allowbreak\detokenize{A^c}} & finite atom mass event measure and relative complement & P15 finite-probability & meta & \hyperref[sec:p15-finite-probability]{\ref*{sec:p15-finite-probability}} & \thepage & Masses \(p(\omega)\) sum to one; \(\mathbb P(A)\) is their sum over \(A\), and \(A^c=\Omega\setminus A\). A zero mass is permitted; almost-sure truth need not hold at every outcome. \\
\phantomsection\label{entry:symbol-registry-00925}\(X:\Omega\to S,p_X,p_{X,Y}\) & \texttt{\detokenize{X,}\allowbreak\detokenize{p_X,}\allowbreak\detokenize{p_{X,Y}}} & finite random variable and marginal or joint laws & P15 finite-probability & meta & \hyperref[sec:p15-finite-probability]{\ref*{sec:p15-finite-probability}} & \thepage & A variable is a function on outcomes; its law is the induced mass function on its finite value set. Joint masses concern pairs of values, including unattainable pairs of mass zero. \\
\phantomsection\label{entry:symbol-registry-00926}\(\{X\in D\},\{X=x\},f(X)\) & \texttt{\detokenize{\{X\in D\},}\allowbreak\detokenize{\{X=x\},}\allowbreak\detokenize{f(X)}} & inverse-image events and composition of random variables & P15 finite-probability & meta & \hyperref[sec:p15-finite-probability]{\ref*{sec:p15-finite-probability}} & \thepage & The braces denote subsets of the declared outcome space. The expression \(f(X)\) denotes \(f\circ X\); grouping over its fibers computes its law or expectation. \\
\phantomsection\label{entry:symbol-registry-00927}\(\mathbb P(A\mid B),\mathbb P_B,p_B\) & \texttt{\detokenize{\mathbb P(\cdot\mid B),}\allowbreak\detokenize{\mathbb P_B,}\allowbreak\detokenize{p_B}} & positive-event conditional probability measure and masses & P15 finite-probability & meta & \hyperref[sec:p15-finite-probability]{\ref*{sec:p15-finite-probability}} & \thepage & Defined only when \(\mathbb P(B)>0\). The conditional measure retains the same event space, assigns zero mass outside \(B\), and renormalizes the masses inside it. \\
\phantomsection\label{entry:symbol-registry-00928}\(H_i,E,\pi_i,J,\ell_i\) & \texttt{\detokenize{H_i,}\allowbreak\detokenize{E,}\allowbreak\detokenize{\pi_i,}\allowbreak\detokenize{J,}\allowbreak\detokenize{\ell_i}} & hypothesis partition evidence priors and likelihoods & P15 finite-probability & meta & \hyperref[sec:p15-finite-probability]{\ref*{sec:p15-finite-probability}} & \thepage & The hypotheses form a finite partition; \(J\) contains precisely the positive-prior indices. Likelihood \(\ell_i=\mathbb P(E\mid H_i)\) is defined for \(i\in J\); null blocks enter unconditioned sums with zero contribution. \\
\phantomsection\label{entry:symbol-registry-00929}\(q_1/q_0,\pi/(1-\pi)\) & \texttt{\detokenize{q_1/}\allowbreak\detokenize{q_0,}\allowbreak\detokenize{\pi/}\allowbreak\detokenize{(1-\pi)}} & finite likelihood ratio and prior odds & P15 finite-probability & meta & \hyperref[sec:p15-finite-probability]{\ref*{sec:p15-finite-probability}} & \thepage & The two-hypothesis odds identity assumes \(0<\pi<1\) and \(q_0>0\). Likelihood ratio and posterior probability have different meanings; no ordinary real ratio is assigned a zero denominator. \\
\phantomsection\label{entry:symbol-registry-00930}\(p_2(h,u,v),\pi_h,r_h\) & \texttt{\detokenize{p_2,}\allowbreak\detokenize{\pi_h,}\allowbreak\detokenize{r_h}} & stipulated finite two-alert product model & P15 finite-probability & meta & \hyperref[sec:p15-finite-probability]{\ref*{sec:p15-finite-probability}} & \thepage & The three binary coordinates encode hypothesis and two alerts, with mass \(\pi_h r_h(u)r_h(v)\). The product specifies conditional independence under each hypothesis; one-alert rates alone do not establish that premise. \\
\phantomsection\label{entry:symbol-registry-00931}\(\mathbb E[X],\mathbf1_A\) & \texttt{\detokenize{\mathbb E,}\allowbreak\detokenize{\mathbf1_A}} & finite real expectation and event indicator & P15 finite-probability & meta & \hyperref[sec:p15-finite-probability]{\ref*{sec:p15-finite-probability}} & \thepage & Expectation is the finite mass-weighted sum. The indicator is one on its event and zero outside it. Linearity needs no independence; factorization of product expectation requires its stated hypothesis. \\
\phantomsection\label{entry:symbol-registry-00932}\(\mathbb E[X\mid B]\) & \texttt{\detokenize{\mathbb E[\cdot\mid B]}} & conditional mean given a positive event & P15 finite-probability & meta & \hyperref[sec:p15-finite-probability]{\ref*{sec:p15-finite-probability}} & \thepage & A real number computed using the conditional masses, with \(\mathbb P(B)>0\). It differs in type from conditional expectation given a random variable. \\
\phantomsection\label{entry:symbol-registry-00933}\(\mathbb E[X\mid Y],g(Y)\) & \texttt{\detokenize{\mathbb E[\cdot\mid Y],}\allowbreak\detokenize{g(Y)}} & finite conditional-expectation random variable & P15 finite-probability & meta & \hyperref[sec:p15-finite-probability]{\ref*{sec:p15-finite-probability}} & \thepage & A function of the finite-valued variable \(Y\), using conditional means on its positive-mass fibers. Values chosen on null fibers affect it only on a null event; they do not define null-event conditional probabilities. \\
\phantomsection\label{entry:symbol-registry-01062}\(\Omega,\omega,2^\Omega,A^c\) & \texttt{\detokenize{\Omega,}\allowbreak\detokenize{\omega,}\allowbreak\detokenize{2^\Omega,}\allowbreak\detokenize{A^c}} & outcome space outcomes full power set and relative complement & P15 measurable-probability & meta & \hyperref[sec:p15-measurable-probability]{\ref*{sec:p15-measurable-probability}} & \thepage & The outcome set is nonempty. A complement is taken in that fixed outcome space; the full power set need not be the chosen measurable event collection. \\
\phantomsection\label{entry:symbol-registry-01063}\(\mathcal F,\mathcal H,\mathcal C,\sigma(\mathcal C)\) & \texttt{\detokenize{\mathcal F,}\allowbreak\detokenize{\mathcal H,}\allowbreak\detokenize{\mathcal C,}\allowbreak\detokenize{\sigma}} & event sigma-algebras generators and generated sigma-algebra & P15 measurable-probability & meta & \hyperref[sec:p15-measurable-probability]{\ref*{sec:p15-measurable-probability}} & \thepage & Generation is intersection of all sigma-algebras containing the specified collection. Finite union closure alone describes an algebra and does not imply countable union closure. \\
\phantomsection\label{entry:symbol-registry-01064}\(\mu,\mathbb P,[0,\infty]\) & \texttt{\detokenize{\mu,}\allowbreak\detokenize{\mathbb P,}\allowbreak\detokenize{[0,}\allowbreak\detokenize{\infty]}} & measure probability measure and nonnegative extended codomain & P15 measurable-probability & meta & \hyperref[sec:p15-measurable-probability]{\ref*{sec:p15-measurable-probability}} & \thepage & Measures may take infinity; probability measures have total mass one. The finite-measure condition is retained whenever subtraction or continuity from above requires it. \\
\phantomsection\label{entry:symbol-registry-01065}\(\sum_{i\in I}a_i\) & \texttt{\detokenize{\sum}} & nonnegative countable sum defined by finite-subset suprema & P15 measurable-probability & meta & \hyperref[sec:p15-measurable-probability]{\ref*{sec:p15-measurable-probability}} & \thepage & Terms lie in the nonnegative extended reals and the index set is countable. Grouping includes empty sets and infinite values; no signed-series rearrangement is asserted. \\
\phantomsection\label{entry:symbol-registry-01066}\(p,\mu_p\) & \texttt{\detokenize{p,}\allowbreak\detokenize{\mu_p}} & countable atomic weights and their induced all-subset measure & P15 measurable-probability & meta & \hyperref[sec:p15-measurable-probability]{\ref*{sec:p15-measurable-probability}} & \thepage & The event measure is the nonnegative sum of its point weights. Total weight one gives a probability measure, and singleton probabilities uniquely recover those weights. \\
\phantomsection\label{entry:symbol-registry-01067}\(\mathbb N_+,A_m,Y,V,D_k\) & \texttt{\detokenize{\mathbb N_+,}\allowbreak\detokenize{A_m,}\allowbreak\detokenize{Y,}\allowbreak\detokenize{V,}\allowbreak\detokenize{D_k}} & positive-integer model tails observations and large-jump sets & P15 measurable-probability & meta & \hyperref[sec:p15-measurable-probability]{\ref*{sec:p15-measurable-probability}} & \thepage & The waiting example uses weights \(2^{-n}\), tail events, parity and reciprocal observations. In the separate CDF proof, \(D_k\) denotes points with mass at least \(1/k\), not a waiting-time event. \\
\phantomsection\label{entry:symbol-registry-01068}\(\mathcal B(\mathbb R^d),\mathbb Q,d\) & \texttt{\detokenize{\mathcal B,}\allowbreak\detokenize{\mathbb Q,}\allowbreak\detokenize{d}} & Borel sigma-algebra rational endpoints and positive dimension & P15 measurable-probability & meta & \hyperref[sec:p15-measurable-probability]{\ref*{sec:p15-measurable-probability}} & \thepage & Rational open rectangles form the countable generating collection in each fixed positive finite dimension. Real and rational closed rays generate the one-dimensional Borel sets. \\
\phantomsection\label{entry:symbol-registry-01069}\(X,(S,\mathcal S),X^{-1}(B)\) & \texttt{\detokenize{X,}\allowbreak\detokenize{S,}\allowbreak\detokenize{\mathcal S,}\allowbreak\detokenize{X^{-1}}} & measurable random element target space and inverse-image event & P15 measurable-probability & meta & \hyperref[sec:p15-measurable-probability]{\ref*{sec:p15-measurable-probability}} & \thepage & The inverse-image notation does not require an invertible function. Measurability is relative to the declared source and target sigma-algebras. \\
\phantomsection\label{entry:symbol-registry-01070}\(\sigma(X),f(X)\) & \texttt{\detokenize{\sigma(X),}\allowbreak\detokenize{f(X)}} & information generated by a random element and measurable composition & P15 measurable-probability & meta & \hyperref[sec:p15-measurable-probability]{\ref*{sec:p15-measurable-probability}} & \thepage & The target sigma-algebra is fixed when forming sigma(X). Composition means \(f\circ X\); its measurability requires the declared measurable maps. \\
\phantomsection\label{entry:symbol-registry-01071}\(\mathbf1_A,c_j,s_n\) & \texttt{\detokenize{\mathbf1_A,}\allowbreak\detokenize{c_j,}\allowbreak\detokenize{s_n}} & event indicators simple coefficients and dyadic approximants & P15 measurable-probability & meta & \hyperref[sec:p15-measurable-probability]{\ref*{sec:p15-measurable-probability}} & \thepage & Simple functions are finite-valued and measurable. The nonnegative dyadic sequence increases pointwise to the extended variable and stays finite even at outcomes where that variable is infinite. \\
\phantomsection\label{entry:symbol-registry-01072}\(\overline{\mathbb R},Z_n,\inf,\sup,\liminf,\limsup\) & \texttt{\detokenize{\overline{\mathbb R},}\allowbreak\detokenize{Z_n,}\allowbreak\detokenize{\inf,}\allowbreak\detokenize{\sup,}\allowbreak\detokenize{\liminf,}\allowbreak\detokenize{\limsup}} & extended random variables and pointwise order limits & P15 measurable-probability & meta & \hyperref[sec:p15-measurable-probability]{\ref*{sec:p15-measurable-probability}} & \thepage & The ordered codomain includes both infinities. The real convergence event uses finite differences; no difference of equal infinite limits is formed. \\
\phantomsection\label{entry:symbol-registry-01073}\(\mathbb P_X,p_X\) & \texttt{\detokenize{\mathbb P_X,}\allowbreak\detokenize{p_X}} & pushforward law and singleton point mass & P15 measurable-probability & meta & \hyperref[sec:p15-measurable-probability]{\ref*{sec:p15-measurable-probability}} & \thepage & The law is a measure on measurable value sets. A point mass is defined only when its singleton is measurable; neither object is silently identified with a density. \\
\phantomsection\label{entry:symbol-registry-01074}\(\mathcal L,\mathcal L_0,\mathcal G_A,\nu_1,\nu_2\) & \texttt{\detokenize{\mathcal L,}\allowbreak\detokenize{\mathcal L_0,}\allowbreak\detokenize{\mathcal G_A,}\allowbreak\detokenize{\nu_1,}\allowbreak\detokenize{\nu_2}} & lambda-systems intersection families and compared probability measures & P15 measurable-probability & meta & \hyperref[sec:p15-measurable-probability]{\ref*{sec:p15-measurable-probability}} & \thepage & The two-stage pi-lambda proof supplies equality on the generated sigma-algebra. Its probability-measure subtraction is finite; uniqueness supplies no measure-existence theorem. \\
\phantomsection\label{entry:symbol-registry-01075}\(F_X(a),F_X(a-)\) & \texttt{\detokenize{F_X,}\allowbreak\detokenize{F_X(a-)}} & cumulative distribution and its left limit & P15 measurable-probability & meta & \hyperref[sec:p15-measurable-probability]{\ref*{sec:p15-measurable-probability}} & \thepage & The cumulative function uses the weak lower-tail event. Its difference from the left limit is the point mass; right continuity and both infinite endpoint limits are proved for an already given law. \\
\phantomsection\label{entry:symbol-registry-01109}\(\mathcal S,\mathcal A(\mathcal S),m,\overline m,\mu_0\) & \texttt{\detokenize{\mathcal S,}\allowbreak\detokenize{\mathcal A,}\allowbreak\detokenize{m,}\allowbreak\detokenize{\overline m,}\allowbreak\detokenize{\mu_0}} & semialgebra generated algebra elementary masses and premeasure & P15 measure-extension & meta & \hyperref[sec:p15-measure-extension]{\ref*{sec:p15-measure-extension}} & \thepage & The finite representation of an algebra set is proved independent of its partition. Countable additivity of a premeasure is asserted when the union remains in its algebra. \\
\phantomsection\label{entry:symbol-registry-01110}\(q,\mathcal M_q,T,E\) & \texttt{\detokenize{q,}\allowbreak\detokenize{\mathcal M_q,}\allowbreak\detokenize{T,}\allowbreak\detokenize{E}} & outer measure splitting domain test set and splitting set & P15 measure-extension & meta & \hyperref[sec:p15-measure-extension]{\ref*{sec:p15-measure-extension}} & \thepage & The outer measure is defined on all subsets. Its splitting criterion selects the complete measurable domain, with all infinite-value and null-subset cases included. \\
\phantomsection\label{entry:symbol-registry-01111}\(\mu^*,A_n,\mu\) & \texttt{\detokenize{\mu^*,}\allowbreak\detokenize{A_n,}\allowbreak\detokenize{\mu}} & cover-cost outer measure algebra covers and measure extension & P15 measure-extension & meta & \hyperref[sec:p15-measure-extension]{\ref*{sec:p15-measure-extension}} & \thepage & The outer measure is an infimum over countable algebra covers. Existence does not require sigma-finiteness; the stated uniqueness does. Restriction to the generated sigma-algebra need not be complete. \\
\phantomsection\label{entry:symbol-registry-01112}\(E_n,H_n,\mathcal L_n\) & \texttt{\detokenize{E_n,}\allowbreak\detokenize{H_n,}\allowbreak\detokenize{\mathcal L_n}} & finite-measure covers increasing exhaustion and equality collections & P15 measure-extension & meta & \hyperref[sec:p15-measure-extension]{\ref*{sec:p15-measure-extension}} & \thepage & These are the local sets used in sigma-finite uniqueness. Restrictions are compared on finite pieces before passage to their increasing union. \\
\phantomsection\label{entry:symbol-registry-01113}\(I(a,b),\mathcal S_{\mathbb R},F,m_F,\mu_F\) & \texttt{\detokenize{I(a,}\allowbreak\detokenize{b),}\allowbreak\detokenize{\mathcal S_{\mathbb R},}\allowbreak\detokenize{F,}\allowbreak\detokenize{m_F,}\allowbreak\detokenize{\mu_F}} & half-open real intervals and Stieltjes assignments and measures & P15 measure-extension & meta & \hyperref[sec:p15-measure-extension]{\ref*{sec:p15-measure-extension}} & \thepage & Endpoints may be infinite but interval members are real. F is finite at each real argument and right continuous; its allowed endpoint increments never subtract equal infinities. \\
\phantomsection\label{entry:symbol-registry-01114}\(J_i,K_n,\delta,\eta_i\) & \texttt{\detokenize{J_i,}\allowbreak\detokenize{K_n,}\allowbreak\detokenize{\delta,}\allowbreak\detokenize{\eta_i}} & clipped intervals truncations and endpoint adjustments in the covering proof & P15 measure-extension & meta & \hyperref[sec:p15-measure-extension]{\ref*{sec:p15-measure-extension}} & \thepage & The proof clips unbounded covering pieces before choosing right-continuity enlargements and a compact subcover. Infinite target endpoints are handled by explicit truncation. \\
\phantomsection\label{entry:symbol-registry-01115}\(\lambda,B+t,\Omega_U,\mathcal F_U,\mathbb U,U\) & \texttt{\detokenize{\lambda,}\allowbreak\detokenize{B+t,}\allowbreak\detokenize{\Omega_U,}\allowbreak\detokenize{\mathcal F_U,}\allowbreak\detokenize{\mathbb U,}\allowbreak\detokenize{U}} & Borel length translated sets and the uniform unit-interval model & P15 measure-extension & meta & \hyperref[sec:p15-measure-extension]{\ref*{sec:p15-measure-extension}} & \thepage & Length is constructed from the Stieltjes function F(x)=x. The uniform space has domain (0,1), its trace Borel sigma-algebra, normalized length and the identity observation. \\
\phantomsection\label{entry:symbol-registry-01116}\(Q_F(u),L_u,F(x-)\) & \texttt{\detokenize{Q_F,}\allowbreak\detokenize{L_u,}\allowbreak\detokenize{F(x-)}} & supremum quantile strict sublevel set and cumulative left limit & P15 measure-extension & meta & \hyperref[sec:p15-measure-extension]{\ref*{sec:p15-measure-extension}} & \thepage & The quantile domain is strictly between zero and one. The exact ray equivalence is proved using right continuity; jump fibers may differ at their lower singleton endpoint and flat portions are handled explicitly. \\
\phantomsection\label{entry:symbol-registry-01130}\(\int f\,d\mu,\int_A f\,d\mu\) & \texttt{\detokenize{\int f\,}\allowbreak\detokenize{d\mu,}\allowbreak\detokenize{\int_A f\,}\allowbreak\detokenize{d\mu}} & measurable integral and its restriction to an event & P15 integration-expectation & meta & \hyperref[sec:p15-integration-expectation]{\ref*{sec:p15-integration-expectation}} & \thepage & Nonnegative integrals may be infinite. Restriction equals f on the measurable event and zero outside it, including at infinite original values. A signed integral is defined only when its two part integrals are not both infinite. \\
\phantomsection\label{entry:symbol-registry-01131}\(s,t,a_j,A_j,\nu_f\) & \texttt{\detokenize{s,}\allowbreak\detokenize{t,}\allowbreak\detokenize{a_j,}\allowbreak\detokenize{A_j,}\allowbreak\detokenize{\nu_f}} & simple approximants coefficients measurable parts and integral measure & P15 integration-expectation & meta & \hyperref[sec:p15-integration-expectation]{\ref*{sec:p15-integration-expectation}} & \thepage & Nonnegative simple integrals are independent of their finite partition. The convention zero times infinity equals zero includes zero-valued regions of infinite measure. For nonnegative measurable f the function \(\nu_f(A)=\int_A f\,d\mu\) is a measure. \\
\phantomsection\label{entry:symbol-registry-01132}\(f^+,f^-,|f|\) & \texttt{\detokenize{f^+,}\allowbreak\detokenize{f^-,}\allowbreak\detokenize{|f|}} & positive and negative parts and function absolute value & P15 integration-expectation & meta & \hyperref[sec:p15-integration-expectation]{\ref*{sec:p15-integration-expectation}} & \thepage & Integrability means a finite integral of the absolute value. Integrable extended functions have measurable real representatives equal almost everywhere; linearity uses those representatives to avoid undefined pointwise sums. \\
\phantomsection\label{entry:symbol-registry-01133}\(\mu\text{-a.e.}\) & \texttt{\detokenize{\mu\text{-a.e.}}} & almost-everywhere qualification for a specified measure & P15 integration-expectation & meta & \hyperref[sec:p15-integration-expectation]{\ref*{sec:p15-integration-expectation}} & \thepage & The exceptional set is contained in a measurable null set. Functions and replacements must be measurable; the measure space is not implicitly assumed complete. \\
\phantomsection\label{entry:symbol-registry-01134}\(\|f\|_p,\|f\|_\infty,L^1(\mu)\) & \texttt{\detokenize{\|f\|_p,}\allowbreak\detokenize{\|f\|_\infty,}\allowbreak\detokenize{L^1}} & integral sizes essential supremum and integrable equivalence classes & P15 integration-expectation & meta & \hyperref[sec:p15-integration-expectation]{\ref*{sec:p15-integration-expectation}} & \thepage & The finite-p size is the integral power expression; the infinity size is the infimum of almost-everywhere bounds, with empty infimum infinity. L1 identifies integrable real functions equal almost everywhere. \\
\phantomsection\label{entry:symbol-registry-01135}\(f_n\to f\text{ in measure},\|f_n-f\|_1\to0\) & \texttt{\detokenize{\text{in measure},}\allowbreak\detokenize{\text{in }L^1}} & convergence under the specified measure and in integral distance & P15 integration-expectation & meta & \hyperref[sec:p15-integration-expectation]{\ref*{sec:p15-integration-expectation}} & \thepage & Convergence in measure uses vanishing measure of every positive-error event. Bounded convergence in this chapter requires a common finite-measure support; on probability spaces convergence in measure is convergence in probability. \\
\phantomsection\label{entry:symbol-registry-01136}\(p,q,\varphi,t,c\) & \texttt{\detokenize{p,}\allowbreak\detokenize{q,}\allowbreak\detokenize{\varphi,}\allowbreak\detokenize{t,}\allowbreak\detokenize{c}} & Holder conjugate exponents convex function mixture weight and support point & P15 integration-expectation & meta & \hyperref[sec:p15-integration-expectation]{\ref*{sec:p15-integration-expectation}} & \thepage & Holder permits both endpoint exponents with reciprocal infinity zero. Jensen uses a finite convex function on all real numbers, a normalized measure, and the stated integrability hypotheses; its supporting affine function is proved. \\
\phantomsection\label{entry:symbol-registry-01137}\(\mathbb E[X],\mathbb E[X;A],m,\operatorname{Var}(X)\) & \texttt{\detokenize{\mathbb E[X],}\allowbreak\detokenize{\mathbb E[X;A],}\allowbreak\detokenize{m,}\allowbreak\detokenize{\operatorname{Var}}} & general expectation restricted expectation finite mean and variance & P15 integration-expectation & meta & \hyperref[sec:p15-integration-expectation]{\ref*{sec:p15-integration-expectation}} & \thepage & Expectations use the integral existence conventions. Restricted expectation is defined here for nonnegative or integrable variables. A finite second moment implies a finite mean and justifies the stated variance identities. \\
\phantomsection\label{entry:symbol-registry-01138}\(\mathbb P_X,h,X\) & \texttt{\detokenize{\mathbb P_X,}\allowbreak\detokenize{h,}\allowbreak\detokenize{X}} & law of a random element measurable observation and underlying element & P15 integration-expectation & meta & \hyperref[sec:p15-integration-expectation]{\ref*{sec:p15-integration-expectation}} & \thepage & Integrating the observation over the law equals integrating its composite on the original space. The nonnegative and integrable signed cases are proved separately, with equivalence of the absolute-integrability tests. \\
\phantomsection\label{entry:symbol-registry-01139}\(c_A,r\) & \texttt{\detokenize{c_A,}\allowbreak\detokenize{r}} & nonnegative event infimum and positive probability-bound threshold & P15 integration-expectation & meta & \hyperref[sec:p15-integration-expectation]{\ref*{sec:p15-integration-expectation}} & \thepage & For a nonempty Borel event set the finite nonnegative observable has finite infimum \(c_A\). The empty event is treated separately. Markov and variance bounds divide only by a strictly positive finite threshold. \\
\phantomsection\label{entry:symbol-registry-01140}\(E_n,\theta,a_n,b_n,\lfloor\cdot\rfloor\) & \texttt{\detokenize{E_n,}\allowbreak\detokenize{\theta,}\allowbreak\detokenize{a_n,}\allowbreak\detokenize{b_n,}\allowbreak\detokenize{\lfloor\cdot\rfloor}} & monotone comparison sets finite-measure exhaustion and rounding approximants & P15 integration-expectation & meta & \hyperref[sec:p15-integration-expectation]{\ref*{sec:p15-integration-expectation}} & \thepage & In the monotone proof \(E_n\) is the set where \(f_n\geq\theta s\), for \(0<\theta<1\), and need not have finite measure. In the separate source comparison it is a finite-measure exhaustion; that comparison covers arbitrary nonnegative measurable functions and permits an infinite limiting integral. \\
\phantomsection\label{entry:symbol-registry-01141}\(g,h,K,Y^{(M)},\varepsilon_M,C_M\) & \texttt{\detokenize{g,}\allowbreak\detokenize{h,}\allowbreak\detokenize{K,}\allowbreak\detokenize{Y^{(M)},}\allowbreak\detokenize{\varepsilon_M,}\allowbreak\detokenize{C_M}} & moment-control functions bound truncated variable tail ratio and compact bound & P15 integration-expectation & meta & \hyperref[sec:p15-integration-expectation]{\ref*{sec:p15-integration-expectation}} & \thepage & The stronger-moment argument proves integrability before comparing expectations. Truncations use thresholds with no atom at the limiting absolute value; compact bounds and a vanishing tail ratio control the error uniformly. \\
\phantomsection\label{entry:symbol-registry-01149}\(X,Y,\mathcal A,\mathcal B\) & \texttt{\detokenize{X,}\allowbreak\detokenize{Y,}\allowbreak\detokenize{\mathcal A,}\allowbreak\detokenize{\mathcal B}} & factor sets and their sigma-algebras & P15 product-measures & meta & \hyperref[sec:p15-product-measures]{\ref*{sec:p15-product-measures}} & \thepage & The two measurable spaces are declared separately. Their sigma-algebras need not be complete and their measures need not be finite. \\
\phantomsection\label{entry:symbol-registry-01150}\(\mathcal S,\mathcal A\times\mathcal B,\mathcal A\otimes\mathcal B\) & \texttt{\detokenize{\mathcal S,}\allowbreak\detokenize{\mathcal A\times\mathcal B,}\allowbreak\detokenize{\mathcal A\otimes\mathcal B}} & rectangle family and generated product sigma-algebra & P15 product-measures & meta & \hyperref[sec:p15-product-measures]{\ref*{sec:p15-product-measures}} & \thepage & The source's sigma-algebra product notation is retained and the tensor symbol is an explicit alias. This domain is generated by measurable rectangles and is not silently completed. \\
\phantomsection\label{entry:symbol-registry-01151}\(A\times B,m(A\times B)\) & \texttt{\detokenize{A\times B,}\allowbreak\detokenize{m(A\times B)}} & measurable rectangle and its preliminary nonnegative mass & P15 product-measures & meta & \hyperref[sec:p15-product-measures]{\ref*{sec:p15-product-measures}} & \thepage & The preliminary mass is the product of the two factor masses, with zero times infinity zero. Its partition and countable-additivity properties are proved before extension. \\
\phantomsection\label{entry:symbol-registry-01152}\(\mu_1,\mu_2,\mu_1\times\mu_2,\rho\) & \texttt{\detokenize{\mu_1,}\allowbreak\detokenize{\mu_2,}\allowbreak\detokenize{\mu_1\times\mu_2,}\allowbreak\detokenize{\rho}} & factor measures selected product measure and its local alias & P15 product-measures & meta & \hyperref[sec:p15-product-measures]{\ref*{sec:p15-product-measures}} & \thepage & Existence follows from the specified cover construction without sigma-finiteness. The uniqueness and integral-exchange theorems retain their explicit sigma-finite hypotheses. \\
\phantomsection\label{entry:symbol-registry-01153}\(E_x,E^y,f_x,f^y\) & \texttt{\detokenize{E_x,}\allowbreak\detokenize{E^y,}\allowbreak\detokenize{f_x,}\allowbreak\detokenize{f^y}} & set and function sections in the two coordinate directions & P15 product-measures & meta & \hyperref[sec:p15-product-measures]{\ref*{sec:p15-product-measures}} & \thepage & A section fixes the indicated coordinate and leaves the other as argument. Product measurability implies measurability of every section; separate measurability alone is insufficient. \\
\phantomsection\label{entry:symbol-registry-01154}\(\pi_X,\pi_Y\) & \texttt{\detokenize{\pi_X,}\allowbreak\detokenize{\pi_Y}} & coordinate projection maps & P15 product-measures & meta & \hyperref[sec:p15-product-measures]{\ref*{sec:p15-product-measures}} & \thepage & A map into the product is measurable exactly when both coordinate maps are measurable for the generated product sigma-algebra. \\
\phantomsection\label{entry:symbol-registry-01155}\(\lambda^n\) & \texttt{\detokenize{\lambda^n}} & finite product Borel Lebesgue measure & P15 product-measures & meta & \hyperref[sec:p15-product-measures]{\ref*{sec:p15-product-measures}} & \thepage & The measure lives on the Euclidean Borel sigma-algebra. Finite rectangle lengths, grouping and coordinate permutations are justified by the product construction and uniqueness. \\
\phantomsection\label{entry:symbol-registry-01156}\(T_1f,T_2f\) & \texttt{\detokenize{T_1f,}\allowbreak\detokenize{T_2f}} & nonnegative extended section-integral functions & P15 product-measures & meta & \hyperref[sec:p15-product-measures]{\ref*{sec:p15-product-measures}} & \thepage & Tonelli defines these at every factor point, permitting infinity. Their measurability and both iterated equalities are proved from indicators, simple functions and monotone convergence. \\
\phantomsection\label{entry:symbol-registry-01157}\(H_1,H_2,N_1,N_2\) & \texttt{\detokenize{H_1,}\allowbreak\detokenize{H_2,}\allowbreak\detokenize{N_1,}\allowbreak\detokenize{N_2}} & absolute-section integrals and their infinite-value exceptional sets & P15 product-measures & meta & \hyperref[sec:p15-product-measures]{\ref*{sec:p15-product-measures}} & \thepage & For an integrable product function these exceptional sets are measurable and null. A section at an exceptional point need not be integrable. \\
\phantomsection\label{entry:symbol-registry-01158}\(I_1f,I_2f\) & \texttt{\detokenize{I_1f,}\allowbreak\detokenize{I_2f}} & real measurable signed section-integral versions & P15 product-measures & meta & \hyperref[sec:p15-product-measures]{\ref*{sec:p15-product-measures}} & \thepage & The versions use the actual signed section integral off the specified null sets and zero on them. No undefined infinity-minus-infinity value is assigned to an exceptional section. \\
\phantomsection\label{entry:symbol-registry-01159}\(F_n,G_n,H_n,F_{i,k}\) & \texttt{\detokenize{F_n,}\allowbreak\detokenize{G_n,}\allowbreak\detokenize{H_n,}\allowbreak\detokenize{F_{i,k}}} & finite-measure factor covers and product rectangles & P15 product-measures & meta & \hyperref[sec:p15-product-measures]{\ref*{sec:p15-product-measures}} & \thepage & These local exhaustions allow finite subtraction inside a box before increasing limits. Their role differs from the absolute-section functions with similar letters. \\
\phantomsection\label{entry:symbol-registry-01160}\(g_{n,E},\mathcal L_n\) & \texttt{\detokenize{g_{n,E},}\allowbreak\detokenize{\mathcal L_n}} & finite-box section-measure function and comparison lambda-system & P15 product-measures & meta & \hyperref[sec:p15-product-measures]{\ref*{sec:p15-product-measures}} & \thepage & The lambda-system proof establishes the section-integral identity on a finite box and then extends it by the actual pi-lambda theorem and increasing exhaustions. \\
\phantomsection\label{entry:symbol-registry-01161}\(P_i,Q_i,\widetilde P_i,\widetilde Q_i\) & \texttt{\detokenize{P_i,}\allowbreak\detokenize{Q_i,}\allowbreak\detokenize{\widetilde P_i,}\allowbreak\detokenize{\widetilde Q_i}} & positive and negative section integrals and finite versions & P15 product-measures & meta & \hyperref[sec:p15-product-measures]{\ref*{sec:p15-product-measures}} & \thepage & The signed proof replaces their infinite values only on the specified measurable null set before forming a difference. These are functions, not probability measures or Kamp predicate letters. \\
\phantomsection\label{entry:symbol-registry-01162}\(\phi,c_+,c_-\) & \texttt{\detokenize{\phi,}\allowbreak\detokenize{c_+,}\allowbreak\detokenize{c_-}} & exponential difference and positive compact bounds in a counterexample & P15 product-measures & meta & \hyperref[sec:p15-product-measures]{\ref*{sec:p15-product-measures}} & \thepage & The bounds separately prove that both nonnegative product parts have infinite integral. Two existing iterated values therefore do not define an integrable signed product function. \\
\phantomsection\label{entry:symbol-registry-01163}\(\omega_1,b,<'\) & \texttt{\detokenize{\omega_1,}\allowbreak\detokenize{b,}\allowbreak\detokenize{<'}} & countable-ordinal domain assumed bijection and transported order & P15 product-measures & meta & \hyperref[sec:p15-product-measures]{\ref*{sec:p15-product-measures}} & \thepage & The separate-measurability counterexample explicitly retains its choice and continuum-hypothesis assumptions. The order is transported by that bijection and is not asserted product measurable. \\
\phantomsection\label{entry:symbol-registry-01164}\(\kappa,F(n,y),H(y),N,S\) & \texttt{\detokenize{\kappa,}\allowbreak\detokenize{F(n,}\allowbreak\detokenize{y),}\allowbreak\detokenize{H(y),}\allowbreak\detokenize{N,}\allowbreak\detokenize{S}} & counting measure sequence encoding absolute sum exceptional set and signed sum & P15 product-measures & meta & \hyperref[sec:p15-product-measures]{\ref*{sec:p15-product-measures}} & \thepage & The series theorem first establishes absolute product integrability. The signed sum is a measurable version defined after a.e. absolute convergence, and the numerical integral series is absolutely convergent. \\
\phantomsection\label{entry:symbol-registry-01165}\(D_g\) & \texttt{\detokenize{D_g}} & measurable strict subgraph for layer-cake integration & P15 product-measures & meta & \hyperref[sec:p15-product-measures]{\ref*{sec:p15-product-measures}} & \thepage & Its countable-rectangle description proves product measurability before Tonelli is used. Section lengths include zero, finite and infinite function values. \\
\phantomsection\label{entry:symbol-registry-01235}\(A_i,I,J,K\) & \texttt{\detokenize{A_i,}\allowbreak\detokenize{I,}\allowbreak\detokenize{J,}\allowbreak\detokenize{K}} & indexed event families and finite test index sets & P15 independent-random-variables & meta & \hyperref[sec:p15-independent-random-variables]{\ref*{sec:p15-independent-random-variables}} & \thepage & Independence is tested on every finite subfamily; K selects the factors different from the whole space, including the empty test. \\
\phantomsection\label{entry:symbol-registry-01236}\(\mathcal C_i,\overline{\mathcal C}_i,\mathcal F_i,\mathcal L_D\) & \texttt{\detokenize{\mathcal C_i,}\allowbreak\detokenize{\overline{\mathcal C}_i,}\allowbreak\detokenize{\mathcal F_i,}\allowbreak\detokenize{\mathcal L_D}} & event generators their whole-space extensions generated sigma-algebras and a proof lambda-system & P15 independent-random-variables & meta & \hyperref[sec:p15-independent-random-variables]{\ref*{sec:p15-independent-random-variables}} & \thepage & The extension proof starts with pi-systems and adds the whole space before extending one coordinate at a time. \\
\phantomsection\label{entry:symbol-registry-01237}\(I_j,\mathcal G_j,X_{i,j},f_i\) & \texttt{\detokenize{I_j,}\allowbreak\detokenize{\mathcal G_j,}\allowbreak\detokenize{X_{i,j},}\allowbreak\detokenize{f_i}} & disjoint block index sets generated sigma-algebras and measurable block outputs & P15 independent-random-variables & meta & \hyperref[sec:p15-independent-random-variables]{\ref*{sec:p15-independent-random-variables}} & \thepage & Blocks may be finite or countable, with independence still quantified over finite subfamilies. Overlapping blocks do not inherit the conclusion. \\
\phantomsection\label{entry:symbol-registry-01238}\(X_i,(S_i,\mathcal S_i),t_i,\mu_i,\mu,\nu\) & \texttt{\detokenize{X_i,}\allowbreak\detokenize{(S_i,}\allowbreak\detokenize{\mathcal S_i),}\allowbreak\detokenize{t_i,}\allowbreak\detokenize{\mu_i,}\allowbreak\detokenize{\mu,}\allowbreak\detokenize{\nu}} & random elements target measurable spaces thresholds and marginal laws & P15 independent-random-variables & meta & \hyperref[sec:p15-independent-random-variables]{\ref*{sec:p15-independent-random-variables}} & \thepage & Here mu and nu are measures. Their scalar mean uses in the moment section are separately declared; a joint law is a product exactly under the proved independence hypotheses. \\
\phantomsection\label{entry:symbol-registry-01239}\(\operatorname{Cov},\operatorname{Corr},\mu_X,\mu_Y,a_i,S_n,C\) & \texttt{\detokenize{\operatorname{Cov},}\allowbreak\detokenize{\operatorname{Corr},}\allowbreak\detokenize{\mu_X,}\allowbreak\detokenize{\mu_Y,}\allowbreak\detokenize{a_i,}\allowbreak\detokenize{S_n,}\allowbreak\detokenize{C}} & covariance correlation means coefficients partial sums and a variance bound & P15 independent-random-variables & meta & \hyperref[sec:p15-independent-random-variables]{\ref*{sec:p15-independent-random-variables}} & \thepage & Correlation requires two positive finite variances. The weak law here uses pairwise uncorrelated variables and the stated finite second moments. \\
\phantomsection\label{entry:symbol-registry-01240}\(Y_n(u),q_n(u),C_k,\varepsilon_j,\lfloor\cdot\rfloor\) & \texttt{\detokenize{Y_n(u),}\allowbreak\detokenize{q_n(u),}\allowbreak\detokenize{C_k,}\allowbreak\detokenize{\varepsilon_j,}\allowbreak\detokenize{\lfloor\cdot\rfloor}} & dyadic bit variables floor recursion prefix cells and specified bit patterns & P15 independent-random-variables & meta & \hyperref[sec:p15-independent-random-variables]{\ref*{sec:p15-independent-random-variables}} & \thepage & The construction uses the actual uniform space on [0,1), specifies dyadic endpoints and proves every finite bit pattern; it does not assume an infinite product theorem. \\
\phantomsection\label{entry:symbol-registry-01241}\(f,g,h,p,\eta,\lambda,\lambda^d,\delta_c\) & \texttt{\detokenize{f,}\allowbreak\detokenize{g,}\allowbreak\detokenize{h,}\allowbreak\detokenize{p,}\allowbreak\detokenize{\eta,}\allowbreak\detokenize{\lambda,}\allowbreak\detokenize{\lambda^d,}\allowbreak\detokenize{\delta_c}} & densities probability laws Lebesgue measures and Dirac masses & P15 convolution-gamma-normal & meta & \hyperref[sec:p15-convolution-gamma-normal]{\ref*{sec:p15-convolution-gamma-normal}} & \thepage & Density versions may differ on a Lebesgue null set. Dirac masses are measures and are not represented by a pointwise infinite density. \\
\phantomsection\label{entry:symbol-registry-01242}\(T(x)=cx+d,c_i\) & \texttt{\detokenize{T,}\allowbreak\detokenize{c,}\allowbreak\detokenize{d,}\allowbreak\detokenize{c_i}} & affine transformation and factor normalization constants & P15 convolution-gamma-normal & meta & \hyperref[sec:p15-convolution-gamma-normal]{\ref*{sec:p15-convolution-gamma-normal}} & \thepage & The density transformation treats positive and negative nonzero slopes; zero slope is a Dirac law. Each factor integral c\_i is proved positive and finite. \\
\phantomsection\label{entry:symbol-registry-01243}\(\mu*\nu,F*G,f*g,dF,dG\) & \texttt{\detokenize{\mu*\nu,}\allowbreak\detokenize{F*G,}\allowbreak\detokenize{f*g,}\allowbreak\detokenize{dF,}\allowbreak\detokenize{dG}} & convolution of measures distribution functions and densities and Stieltjes integration & P15 convolution-gamma-normal & meta & \hyperref[sec:p15-convolution-gamma-normal]{\ref*{sec:p15-convolution-gamma-normal}} & \thepage & Each convolution has its declared input type. The Stieltjes integrators denote probability measures associated with the CDFs. \\
\phantomsection\label{entry:symbol-registry-01244}\(\Gamma(\alpha),g_{\alpha,\lambda},\operatorname{gamma}(\alpha,\lambda)\) & \texttt{\detokenize{\Gamma,}\allowbreak\detokenize{g_{\alpha,\lambda},}\allowbreak\detokenize{\operatorname{gamma}}} & gamma function density and positive shape-rate law & P15 convolution-gamma-normal & meta & \hyperref[sec:p15-convolution-gamma-normal]{\ref*{sec:p15-convolution-gamma-normal}} & \thepage & The scalar lambda here is a strictly positive rate, not the Lebesgue measure used in the density preliminaries. Both gamma-sum source routes retain a common rate. \\
\phantomsection\label{entry:symbol-registry-01245}\(B(\alpha,\beta),\operatorname{beta}(\alpha,\beta),c_{\alpha,\beta}\) & \texttt{\detokenize{B,}\allowbreak\detokenize{\operatorname{beta},}\allowbreak\detokenize{c_{\alpha,\beta}}} & beta integral normalized law and gamma convolution constant & P15 convolution-gamma-normal & meta & \hyperref[sec:p15-convolution-gamma-normal]{\ref*{sec:p15-convolution-gamma-normal}} & \thepage & Both shape parameters are positive. Endpoint density values are explicitly selected null-set versions; the power formula applies on the open support. \\
\phantomsection\label{entry:symbol-registry-01246}\(\phi_{\mu,a},\operatorname{normal}(\mu,a),J,v\) & \texttt{\detokenize{\phi_{\mu,a},}\allowbreak\detokenize{\operatorname{normal},}\allowbreak\detokenize{J,}\allowbreak\detokenize{v}} & normal density mean-variance law Gaussian half-integral and completed-square variance & P15 convolution-gamma-normal & meta & \hyperref[sec:p15-convolution-gamma-normal]{\ref*{sec:p15-convolution-gamma-normal}} & \thepage & The scalar mu is a mean and a is variance. The zero-variance law is separately defined as a Dirac mass; positive-variance formulas never divide by zero. \\
\phantomsection\label{entry:symbol-registry-01273}\(\mathbb N_+,\Omega_\infty,\pi_n,X_i,\mathcal S,\mathcal A,\mathcal C,\mathcal F_\infty\) & \texttt{\detokenize{\mathbb N_+,}\allowbreak\detokenize{\Omega_\infty,}\allowbreak\detokenize{\pi_n,}\allowbreak\detokenize{X_i,}\allowbreak\detokenize{\mathcal S,}\allowbreak\detokenize{\mathcal A,}\allowbreak\detokenize{\mathcal C,}\allowbreak\detokenize{\mathcal F_\infty}} & positive coordinate indices real sequence space projections cylinder collections and product sigma-algebra & P15 countable-product-measures & meta & \hyperref[sec:p15-countable-product-measures]{\ref*{sec:p15-countable-product-measures}} & \thepage & The half-open cylinder semialgebra, its finite-union algebra and the larger Borel-cylinder algebra are explicitly distinguished. \\
\phantomsection\label{entry:symbol-registry-01274}\(\mu_n,p_0,m_n,B_n^*,K_n,C_n,D_n,D_n^*\) & \texttt{\detokenize{\mu_n,}\allowbreak\detokenize{p_0,}\allowbreak\detokenize{m_n,}\allowbreak\detokenize{B_n^*,}\allowbreak\detokenize{K_n,}\allowbreak\detokenize{C_n,}\allowbreak\detokenize{D_n,}\allowbreak\detokenize{D_n^*}} & consistent finite laws coherent cylinder mass growing prefix dimensions and compact approximation sets & P15 countable-product-measures & meta & \hyperref[sec:p15-countable-product-measures]{\ref*{sec:p15-countable-product-measures}} & \thepage & p\_0 is proved finitely additive before its continuity argument. Compact bases lie in finite-dimensional spaces; their cylinders still have unrestricted later coordinates. \\
\phantomsection\label{entry:symbol-registry-01275}\(\delta,r_i(j),s_k,\theta_i,\eta_i,E_i\) & \texttt{\detokenize{\delta,}\allowbreak\detokenize{r_i(j),}\allowbreak\detokenize{s_k,}\allowbreak\detokenize{\theta_i,}\allowbreak\detokenize{\eta_i,}\allowbreak\detokenize{E_i}} & hypothetical positive mass limit diagonal subsequences coordinate limits marginal laws and events & P15 countable-product-measures & meta & \hyperref[sec:p15-countable-product-measures]{\ref*{sec:p15-countable-product-measures}} & \thepage & The diagonal indices increase and converge on each fixed coordinate. The event probability for all coordinates is derived only after countable additivity is established. \\
\phantomsection\label{entry:symbol-registry-01276}\(Q,\delta(x,y),b_k(t),\theta(j,k),\phi_0,E_0,D_j,D\) & \texttt{\detokenize{Q,}\allowbreak\detokenize{\delta,}\allowbreak\detokenize{b_k,}\allowbreak\detokenize{\theta,}\allowbreak\detokenize{\phi_0,}\allowbreak\detokenize{E_0,}\allowbreak\detokenize{D_j,}\allowbreak\detokenize{D}} & unit-interval sequence space product metric canonical binary digits diagonal code image and inverse & P15 countable-product-measures & meta & \hyperref[sec:p15-countable-product-measures]{\ref*{sec:p15-countable-product-measures}} & \thepage & Every decoded row must have infinitely many zeros. The formula gives the exact Borel image, includes zero, excludes one and resolves all dyadic ambiguity. \\
\phantomsection\label{entry:symbol-registry-01277}\(M,\rho,S,H,q_j,d,\psi,C,R,K\) & \texttt{\detokenize{M,}\allowbreak\detokenize{\rho,}\allowbreak\detokenize{S,}\allowbreak\detokenize{H,}\allowbreak\detokenize{q_j,}\allowbreak\detokenize{d,}\allowbreak\detokenize{\psi,}\allowbreak\detokenize{C,}\allowbreak\detokenize{R,}\allowbreak\detokenize{K}} & complete metric space Borel subset closure dense sequence bounded metric and distance-coordinate image & P15 countable-product-measures & meta & \hyperref[sec:p15-countable-product-measures]{\ref*{sec:p15-countable-product-measures}} & \thepage & The inverse R is constructed on the explicitly Borel image C using a Cauchy sequence in the complete closure H. K is the Borel image of the possibly nonclosed S. \\
\phantomsection\label{entry:symbol-registry-01278}\(\phi_i,E_i,E^{(n)},\Omega_E,T,\nu_n\) & \texttt{\detokenize{\phi_i,}\allowbreak\detokenize{E_i,}\allowbreak\detokenize{E^{(n)},}\allowbreak\detokenize{\Omega_E,}\allowbreak\detokenize{T,}\allowbreak\detokenize{\nu_n}} & real codings Borel images full-measure coordinate subset inverse transport and original finite laws & P15 countable-product-measures & meta & \hyperref[sec:p15-countable-product-measures]{\ref*{sec:p15-countable-product-measures}} & \thepage & Borel images make the countable concentration set measurable. The transported extension is for the declared coordinate product sigma-algebra. \\
\phantomsection\label{entry:symbol-registry-01285}\(\alpha,\alpha^+,\alpha^-,\lvert\alpha\rvert\) & \texttt{\detokenize{\alpha,}\allowbreak\detokenize{\alpha^+,}\allowbreak\detokenize{\alpha^-,}\allowbreak\detokenize{\lvert\alpha\rvert}} & signed measure Jordan components and total variation measure & P15 radon-nikodym & meta & \hyperref[sec:p15-radon-nikodym]{\ref*{sec:p15-radon-nikodym}} & \thepage & The signed convention permits positive infinity and excludes negative infinity. The negative Jordan measure is finite; the bars around alpha name a measure, while bars around its scalar values denote absolute value. \\
\phantomsection\label{entry:symbol-registry-01286}\(A,B,F_k,E_k,n_k,c\) & \texttt{\detokenize{A,}\allowbreak\detokenize{B,}\allowbreak\detokenize{F_k,}\allowbreak\detokenize{E_k,}\allowbreak\detokenize{n_k,}\allowbreak\detokenize{c}} & positive and negative sets removal remainders selected positive subsets reciprocal thresholds and negative infimum & P15 radon-nikodym & meta & \hyperref[sec:p15-radon-nikodym]{\ref*{sec:p15-radon-nikodym}} & \thepage & A zero signed value need not make a set signed-null. The negative-subset proof uses least positive reciprocal thresholds and accounts for the infinite remainder. \\
\phantomsection\label{entry:symbol-registry-01287}\(\mu\perp\nu,\nu\ll\mu,\nu_{\mathrm r},\nu_{\mathrm s},d\nu/d\mu\) & \texttt{\detokenize{\perp,}\allowbreak\detokenize{\ll,}\allowbreak\detokenize{\nu_{\mathrm r},}\allowbreak\detokenize{\nu_{\mathrm s},}\allowbreak\detokenize{d\nu/}\allowbreak\detokenize{d\mu}} & mutual singularity absolute continuity regular and singular parts and density & P15 radon-nikodym & meta & \hyperref[sec:p15-radon-nikodym]{\ref*{sec:p15-radon-nikodym}} & \thepage & The base measure and original sigma-algebra are fixed. Density versions agree almost everywhere for that base; reciprocal and chain identities carry their specified null-set qualifications. \\
\phantomsection\label{entry:symbol-registry-01288}\(\mathcal G,\kappa,g_n,h_n,h,D_m,A_m,\sigma_m\) & \texttt{\detokenize{\mathcal G,}\allowbreak\detokenize{\kappa,}\allowbreak\detokenize{g_n,}\allowbreak\detokenize{h_n,}\allowbreak\detokenize{h,}\allowbreak\detokenize{D_m,}\allowbreak\detokenize{A_m,}\allowbreak\detokenize{\sigma_m}} & dominated-density family maximal integral approximations common finite partition and singular residuals & P15 radon-nikodym & meta & \hyperref[sec:p15-radon-nikodym]{\ref*{sec:p15-radon-nikodym}} & \thepage & The common partition makes both given measures finite on each piece. The global decomposition is assembled by nonnegative sums and does not subtract infinite totals. \\
\phantomsection\label{entry:symbol-registry-01289}\(H,F,G,\tau,\mathcal A_I,U,E\triangle U\) & \texttt{\detokenize{H,}\allowbreak\detokenize{F,}\allowbreak\detokenize{G,}\allowbreak\detokenize{\tau,}\allowbreak\detokenize{\mathcal A_I,}\allowbreak\detokenize{U,}\allowbreak\detokenize{\triangle}} & signed cumulative function positive cumulative parts variation restriction interval algebra and approximation error & P15 radon-nikodym & meta & \hyperref[sec:p15-radon-nikodym]{\ref*{sec:p15-radon-nikodym}} & \thepage & The Stieltjes converse includes right continuity and zero limits at negative infinity. Variation on (a,b] excludes the initial jump at a. \\
\phantomsection\label{entry:symbol-registry-01290}\(b,c_b,I_b,C_n,C,B_j,X,\eta\) & \texttt{\detokenize{b,}\allowbreak\detokenize{c_b,}\allowbreak\detokenize{I_b,}\allowbreak\detokenize{C_n,}\allowbreak\detokenize{C,}\allowbreak\detokenize{B_j,}\allowbreak\detokenize{X,}\allowbreak\detokenize{\eta}} & Cantor digit words interval endpoints stage intervals stages fair bits series and singular law & P15 radon-nikodym & meta & \hyperref[sec:p15-radon-nikodym]{\ref*{sec:p15-radon-nikodym}} & \thepage & Each stage interval has law mass 2 to the power minus its word length. All singleton masses vanish while the Cantor support has Lebesgue measure zero. \\
\phantomsection\label{entry:symbol-registry-01291}\(\limsup A_n,\liminf A_n,\mathrm{i.o.},N_n,N_\infty,p_j,t_n,n_k,T_k\) & \texttt{\detokenize{\limsup,}\allowbreak\detokenize{\liminf,}\allowbreak\detokenize{\mathrm{i.o.},}\allowbreak\detokenize{N_n,}\allowbreak\detokenize{N_\infty,}\allowbreak\detokenize{p_j,}\allowbreak\detokenize{t_n,}\allowbreak\detokenize{n_k,}\allowbreak\detokenize{T_k}} & event limits occurrence counts expected counts deterministic subsequence and subsequence counts & P15 limsup-events-and-zero-one-laws & meta & \hyperref[sec:p15-limsup-events-and-zero-one-laws]{\ref*{sec:p15-limsup-events-and-zero-one-laws}} & \thepage & The second Borel--Cantelli argument uses mutual independence; the separate count refinement has its explicitly proved pairwise-independence hypotheses. \\
\phantomsection\label{entry:symbol-registry-01292}\(\mathcal H_n,\mathcal H_\infty,\mathcal F'_n,\mathcal T,\mathcal E\) & \texttt{\detokenize{\mathcal H_n,}\allowbreak\detokenize{\mathcal H_\infty,}\allowbreak\detokenize{\mathcal F'_n,}\allowbreak\detokenize{\mathcal T,}\allowbreak\detokenize{\mathcal E}} & finite-information full observation future tail and finite-permutation event fields & P15 limsup-events-and-zero-one-laws & meta & \hyperref[sec:p15-limsup-events-and-zero-one-laws]{\ref*{sec:p15-limsup-events-and-zero-one-laws}} & \thepage & Tail and finite-permutation fields are distinct collections of events even when both are probability-trivial under the specified independent identically distributed law. \\
\phantomsection\label{entry:symbol-registry-01293}\(A\triangle B,d_n,\varepsilon_n,\pi_n,r_m,R_n^{(m)}\) & \texttt{\detokenize{\triangle,}\allowbreak\detokenize{d_n,}\allowbreak\detokenize{\varepsilon_n,}\allowbreak\detokenize{\pi_n,}\allowbreak\detokenize{r_m,}\allowbreak\detokenize{R_n^{(m)}}} & symmetric difference approximation errors prefix projection replacement and tail sums & P15 limsup-events-and-zero-one-laws & meta & \hyperref[sec:p15-limsup-events-and-zero-one-laws]{\ref*{sec:p15-limsup-events-and-zero-one-laws}} & \thepage & The approximation theorem supplies a sequence indexed by the actual number of observed coordinates, with exact measurable prefix representations. \\
\phantomsection\label{entry:symbol-registry-01294}\(\pi,\pi^{(n)},T_\pi,B_n,A'_n,p(q),U,H_n\) & \texttt{\detokenize{\pi,}\allowbreak\detokenize{\pi^{(n)},}\allowbreak\detokenize{T_\pi,}\allowbreak\detokenize{B_n,}\allowbreak\detokenize{A'_n,}\allowbreak\detokenize{p(q),}\allowbreak\detokenize{U,}\allowbreak\detokenize{H_n}} & finite permutations block swaps prefix events rational-cut probabilities and a harmonic scalar & P15 limsup-events-and-zero-one-laws & meta & \hyperref[sec:p15-limsup-events-and-zero-one-laws]{\ref*{sec:p15-limsup-events-and-zero-one-laws}} & \thepage & Permutation law invariance requires identical independent coordinate laws. H\_n is a scalar in an example, distinct from the calligraphic information field. \\
\phantomsection\label{entry:symbol-registry-01295}\((\Omega,\mathcal F_0,P),\mathcal G,\mathcal H,C_{\mathcal G}X\) & \texttt{\detokenize{\Omega,}\allowbreak\detokenize{\mathcal F_0,}\allowbreak\detokenize{P,}\allowbreak\detokenize{\mathcal G,}\allowbreak\detokenize{\mathcal H,}\allowbreak\detokenize{C_{\mathcal G}}} & ambient probability space conditioning sigma-algebras and a conditional-expectation version & P15 conditional-expectation & meta & \hyperref[sec:p15-conditional-expectation]{\ref*{sec:p15-conditional-expectation}} & \thepage & Versions are measurable for the stated conditioning field and agree only almost everywhere; local equality on a conditioning event does not imply equality outside it. \\
\phantomsection\label{entry:symbol-registry-01296}\(P(A\mid\mathcal G),E(X\mid Y),h(Y),\mu_Y\) & \texttt{\detokenize{P(A\mid\mathcal G),}\allowbreak\detokenize{E(X\mid Y),}\allowbreak\detokenize{h(Y),}\allowbreak\detokenize{\mu_Y}} & conditional probabilities observation conditioning scalar factorization and observation law & P15 conditional-expectation & meta & \hyperref[sec:p15-conditional-expectation]{\ref*{sec:p15-conditional-expectation}} & \thepage & A chosen value h(y) is not a ratio conditional on a null event. Its almost-everywhere uniqueness is with respect to the law of Y. \\
\phantomsection\label{entry:symbol-registry-01297}\(Z_n,A_{mn},Y_n,\Omega_0\) & \texttt{\detokenize{Z_n,}\allowbreak\detokenize{A_{mn},}\allowbreak\detokenize{Y_n,}\allowbreak\detokenize{\Omega_0}} & conditional truncation versions bounded comparison events tail infima and common good sets & P15 conditional-expectation & meta & \hyperref[sec:p15-conditional-expectation]{\ref*{sec:p15-conditional-expectation}} & \thepage & Order comparisons bound the finite-side output by m before subtracting event integrals. Countably many output exceptions are combined in the conditioning sigma-algebra. \\
\phantomsection\label{entry:symbol-registry-01298}\(\varphi,L_x,L',L_\theta,a_x,b_x,\theta\) & \texttt{\detokenize{\varphi,}\allowbreak\detokenize{L_x,}\allowbreak\detokenize{L',}\allowbreak\detokenize{L_\theta,}\allowbreak\detokenize{a_x,}\allowbreak\detokenize{b_x,}\allowbreak\detokenize{\theta}} & convex function supporting lines and rational-minorant interpolation & P15 conditional-expectation & meta & \hyperref[sec:p15-conditional-expectation]{\ref*{sec:p15-conditional-expectation}} & \thepage & The source's countable rational-coefficient minorants represent every nonaffine finite convex function. An irrational-slope affine function is treated separately by linearity. \\
\phantomsection\label{entry:symbol-registry-01299}\(L^p(\mathcal G),\|\cdot\|_p,\|\cdot\|_\infty,\langle X,Y\rangle\) & \texttt{\detokenize{L^p,}\allowbreak\detokenize{\|\cdot\|_p,}\allowbreak\detokenize{\|\cdot\|_\infty,}\allowbreak\detokenize{\langle X,}\allowbreak\detokenize{Y\rangle}} & almost-everywhere function classes norms essential bound and L2 inner product & P15 conditional-expectation & meta & \hyperref[sec:p15-conditional-expectation]{\ref*{sec:p15-conditional-expectation}} & \thepage & The norm laws and conditional contraction are proved for the displayed p range. Completeness and the closed-subspace assertion are explicitly proved for L2. \\
\phantomsection\label{entry:symbol-registry-01300}\(M,Z,D_j,H,n_j,\operatorname{Var}(X\mid\mathcal G)\) & \texttt{\detokenize{M,}\allowbreak\detokenize{Z,}\allowbreak\detokenize{D_j,}\allowbreak\detokenize{H,}\allowbreak\detokenize{n_j,}\allowbreak\detokenize{\operatorname{Var}(X\mid\mathcal G)}} & conditional predictor measurable multiplier Cauchy increments absolute-series bound subsequence and conditional variance & P15 conditional-expectation & meta & \hyperref[sec:p15-conditional-expectation]{\ref*{sec:p15-conditional-expectation}} & \thepage & The predictor minimizes square error. A finite-error competitor is proved square integrable; arbitrary prediction losses are not covered by that claim. \\
\phantomsection\label{entry:symbol-registry-01301}\(m(y),a(y),D,h(y),\phi(s,t),A_i,B_i,D_i\) & \texttt{\detokenize{m,}\allowbreak\detokenize{a,}\allowbreak\detokenize{D,}\allowbreak\detokenize{h,}\allowbreak\detokenize{\phi,}\allowbreak\detokenize{A_i,}\allowbreak\detokenize{B_i,}\allowbreak\detokenize{D_i}} & marginal and absolute section integrals admissible domain parameter integrand and lifted information cells & P15 conditional-expectation & meta & \hyperref[sec:p15-conditional-expectation]{\ref*{sec:p15-conditional-expectation}} & \thepage & Zero or infinite marginal densities and nonintegrable sections are handled on their explicit null domains. The scalar lift never assumes the observation image measurable. \\
\phantomsection\label{entry:symbol-registry-01302}\(N,S,\mu,\sigma^2,A_+,A_-,A_0\) & \texttt{\detokenize{N,}\allowbreak\detokenize{S,}\allowbreak\detokenize{\mu,}\allowbreak\detokenize{\sigma^2,}\allowbreak\detokenize{A_+,}\allowbreak\detokenize{A_-,}\allowbreak\detokenize{A_0}} & independent random count random sum moments and shifted-sign comparison events & P15 conditional-expectation & meta & \hyperref[sec:p15-conditional-expectation]{\ref*{sec:p15-conditional-expectation}} & \thepage & The random-sum proof establishes its second moment before conditioning. The same-law result for merely integrable variables uses countably many rational shifted thresholds. \\
\phantomsection\label{entry:symbol-registry-01303}\(K(y,A),K(\omega,A),K(y,dx)\) & \texttt{\detokenize{K(y,}\allowbreak\detokenize{A),}\allowbreak\detokenize{K(\omega,}\allowbreak\detokenize{A),}\allowbreak\detokenize{K(y,}\allowbreak\detokenize{dx)}} & probability kernels on observation or conditioning spaces and their integration notation & P15 regular-conditional-laws & meta & \hyperref[sec:p15-regular-conditional-laws]{\ref*{sec:p15-regular-conditional-laws}} & \thepage & The second argument is a measurable target event. The constructed kernels are probability measures at every parameter, with a measurable default on the exceptional domain. \\
\phantomsection\label{entry:symbol-registry-01304}\(g_q,H_q,F_\omega,F_y,\Omega_0,T_0,\nu_\omega,\nu_y\) & \texttt{\detokenize{g_q,}\allowbreak\detokenize{H_q,}\allowbreak\detokenize{F_\omega,}\allowbreak\detokenize{F_y,}\allowbreak\detokenize{\Omega_0,}\allowbreak\detokenize{T_0,}\allowbreak\detokenize{\nu_\omega,}\allowbreak\detokenize{\nu_y}} & rational conditional probabilities their observation lifts CDF families common good domains and Stieltjes laws & P15 regular-conditional-laws & meta & \hyperref[sec:p15-regular-conditional-laws]{\ref*{sec:p15-regular-conditional-laws}} & \thepage & Monotonicity, right continuity at rational arguments and both infinite-endpoint limits hold on one countably defined full-measure domain. \\
\phantomsection\label{entry:symbol-registry-01305}\(\phi,E,s_0,\Omega_1,T_1,\mu_Y\) & \texttt{\detokenize{\phi,}\allowbreak\detokenize{E,}\allowbreak\detokenize{s_0,}\allowbreak\detokenize{\Omega_1,}\allowbreak\detokenize{T_1,}\allowbreak\detokenize{\mu_Y}} & Borel real coding image fixed default point concentration domains and observation law & P15 regular-conditional-laws & meta & \hyperref[sec:p15-regular-conditional-laws]{\ref*{sec:p15-regular-conditional-laws}} & \thepage & The target image is proved Borel. Concentration on it is proved before transporting the real kernel back; the conditioning value space itself may be arbitrary measurable. \\
\phantomsection\label{entry:symbol-registry-01306}\(E_y,\mathcal D,Q,h(y,x),m(y),D,x_0\) & \texttt{\detokenize{E_y,}\allowbreak\detokenize{\mathcal D,}\allowbreak\detokenize{Q,}\allowbreak\detokenize{h,}\allowbreak\detokenize{m,}\allowbreak\detokenize{D,}\allowbreak\detokenize{x_0}} & event sections generating class kernel-induced joint measure parameter integrand and density-kernel domain & P15 regular-conditional-laws & meta & \hyperref[sec:p15-regular-conditional-laws]{\ref*{sec:p15-regular-conditional-laws}} & \thepage & Kernel integration supplies a joint probability on the product sigma-algebra. Signed inner integrals use measurable zero assignments on their absolute-integrability exceptional sets. \\
\phantomsection\label{entry:symbol-registry-01307}\(\Theta,\Phi,\Theta',\Psi,r(\theta,\phi),X,Y,Z,\varepsilon\) & \texttt{\detokenize{\Theta,}\allowbreak\detokenize{\Phi,}\allowbreak\detokenize{\Theta',}\allowbreak\detokenize{\Psi,}\allowbreak\detokenize{r,}\allowbreak\detokenize{X,}\allowbreak\detokenize{Y,}\allowbreak\detokenize{Z,}\allowbreak\detokenize{\varepsilon}} & spherical coordinate variables parametrization uniform-square observations and shrinking-event size & P15 regular-conditional-laws & meta & \hyperref[sec:p15-regular-conditional-laws]{\ref*{sec:p15-regular-conditional-laws}} & \thepage & The alternative great-circle angle has density absolute cosine divided by four. The two square-coordinate observations approach the same null event through distinct positive-probability event families. \\
\phantomsection\label{entry:symbol-registry-01325}\(S_n,S_0,S_\infty,C_X,w_M\) & \texttt{\detokenize{S_n,}\allowbreak\detokenize{S_0,}\allowbreak\detokenize{S_\infty,}\allowbreak\detokenize{C_X,}\allowbreak\detokenize{w_M}} & series partial sums finite limit convergence event and tail oscillations & P15 independent-random-series & meta & \hyperref[sec:p15-independent-random-series]{\ref*{sec:p15-independent-random-series}} & \thepage & Convergence is the countable Cauchy event and a tail event. A finite initial modification does not change it; the zero-one assertion uses mutual independence. \\
\phantomsection\label{entry:symbol-registry-01326}\(A_k,T_k,V,U_j,M,W,v,D,q_m\) & \texttt{\detokenize{A_k,}\allowbreak\detokenize{T_k,}\allowbreak\detokenize{V,}\allowbreak\detokenize{U_j,}\allowbreak\detokenize{M,}\allowbreak\detokenize{W,}\allowbreak\detokenize{v,}\allowbreak\detokenize{D,}\allowbreak\detokenize{q_m}} & first-crossing events future increments variances bounded summands and oscillation controls & P15 independent-random-series & meta & \hyperref[sec:p15-independent-random-series]{\ref*{sec:p15-independent-random-series}} & \thepage & The fourth-moment lower bound supplies an elementary bounded converse. Local uses of M as an index and as a uniform bound are distinguished in their respective proofs. \\
\phantomsection\label{entry:symbol-registry-01327}\(\widehat Y_n,\widetilde Y_n,Z_n,\Omega^2\) & \texttt{\detokenize{\widehat Y_n,}\allowbreak\detokenize{\widetilde Y_n,}\allowbreak\detokenize{Z_n,}\allowbreak\detokenize{\Omega^2}} & constructed independent copies their differences and the finite product probability space & P15 independent-random-series & meta & \hyperref[sec:p15-independent-random-series]{\ref*{sec:p15-independent-random-series}} & \thepage & The proof constructs the copies and combined-family independence, rather than assuming an infinite product. Difference variances are doubled and their joint law is swap symmetric. \\
\phantomsection\label{entry:symbol-registry-01328}\(A,Y_n^{(A)},p_n^{(A)},\mu_n^{(A)},v_n^{(A)}\) & \texttt{\detokenize{A,}\allowbreak\detokenize{Y_n^{(A)},}\allowbreak\detokenize{p_n^{(A)},}\allowbreak\detokenize{\mu_n^{(A)},}\allowbreak\detokenize{v_n^{(A)}}} & positive deletion cutoff truncated variables tail probabilities means and variances & P15 independent-random-series & meta & \hyperref[sec:p15-independent-random-series]{\ref*{sec:p15-independent-random-series}} & \thepage & The strict and weak cutoff inequalities are fixed explicitly. Three-series necessity imposes no original moment assumption; the full criterion is independent of the chosen positive cutoff. \\
\phantomsection\label{entry:symbol-registry-01329}\(\mathrm{LF},U_{n,m},G,\Phi,\sigma,R_n,c_n,Z_n,V_n,d_n\) & \texttt{\detokenize{\mathrm{LF},}\allowbreak\detokenize{U_{n,m},}\allowbreak\detokenize{G,}\allowbreak\detokenize{\Phi,}\allowbreak\detokenize{\sigma,}\allowbreak\detokenize{R_n,}\allowbreak\detokenize{c_n,}\allowbreak\detokenize{Z_n,}\allowbreak\detokenize{V_n,}\allowbreak\detokenize{d_n}} & stated Lindeberg-Feller premise triangular array normal law normalization and deterministic drift & P15 independent-random-series & meta & \hyperref[sec:p15-independent-random-series]{\ref*{sec:p15-independent-random-series}} & \thepage & The full source necessity route is explicitly conditional on the displayed LF premise. The separate elementary necessity proof already establishes three-series without that still-owed CLT proof. \\
\phantomsection\label{entry:symbol-registry-01330}\(\varepsilon_n,\varepsilon,\alpha,p,a_n\) & \texttt{\detokenize{\varepsilon_n,}\allowbreak\detokenize{\varepsilon,}\allowbreak\detokenize{\alpha,}\allowbreak\detokenize{p,}\allowbreak\detokenize{a_n}} & fair signs common sign power exponents and alternating numerical partial sums & P15 independent-random-series & meta & \hyperref[sec:p15-independent-random-series]{\ref*{sec:p15-independent-random-series}} & \thepage & The signed-power example includes every real exponent and both absolute and conditional convergence thresholds. Common-sign and noncentered examples expose the independence and centering requirements. \\
\phantomsection\label{entry:symbol-registry-01331}\(X_n,S_n,\overline X_n,\mu,Y_k,T_n,R,W,t,y,u_k,u\) & \texttt{\detokenize{X_n,}\allowbreak\detokenize{S_n,}\allowbreak\detokenize{\overline X_n,}\allowbreak\detokenize{\mu,}\allowbreak\detokenize{Y_k,}\allowbreak\detokenize{T_n,}\allowbreak\detokenize{R,}\allowbreak\detokenize{W,}\allowbreak\detokenize{t,}\allowbreak\detokenize{y,}\allowbreak\detokenize{u_k,}\allowbreak\detokenize{u}} & samples averages deletion truncations finite modification bound and tail variables & P15 strong-laws-of-large-numbers & meta & \hyperref[sec:p15-strong-laws-of-large-numbers]{\ref*{sec:p15-strong-laws-of-large-numbers}} & \thepage & Pairwise and mutual independence belong to separate strong-law statements. Truncation estimates need only the common law. \\
\phantomsection\label{entry:symbol-registry-01332}\(\alpha,k(j),j_0,E_\alpha,E,\alpha_\ell,X_n^+,X_n^-\) & \texttt{\detokenize{\alpha,}\allowbreak\detokenize{k(j),}\allowbreak\detokenize{j_0,}\allowbreak\detokenize{E_\alpha,}\allowbreak\detokenize{E,}\allowbreak\detokenize{\alpha_\ell,}\allowbreak\detokenize{X_n^+,}\allowbreak\detokenize{X_n^-}} & geometric ratio floor indices common events countable interpolation ratios and signed parts & P15 strong-laws-of-large-numbers & meta & \hyperref[sec:p15-strong-laws-of-large-numbers]{\ref*{sec:p15-strong-laws-of-large-numbers}} & \thepage & Repeated floor indices are retained. The countable ratios approach one on one probability-one event. \\
\phantomsection\label{entry:symbol-registry-01333}\(a_n,a_0,x_n,b_n,b,q_{n,m},H_M,\widetilde a_j,\widetilde x_j,\mu_n,Z_n\) & \texttt{\detokenize{a_n,}\allowbreak\detokenize{a_0,}\allowbreak\detokenize{x_n,}\allowbreak\detokenize{b_n,}\allowbreak\detokenize{b,}\allowbreak\detokenize{q_{n,m},}\allowbreak\detokenize{H_M,}\allowbreak\detokenize{\widetilde a_j,}\allowbreak\detokenize{\widetilde x_j,}\allowbreak\detokenize{\mu_n,}\allowbreak\detokenize{Z_n}} & Kronecker weights weighted sums normalized weights finite-head errors shifted sequences and centered variables & P15 strong-laws-of-large-numbers & meta & \hyperref[sec:p15-strong-laws-of-large-numbers]{\ref*{sec:p15-strong-laws-of-large-numbers}} & \thepage & The auxiliary zero weight is not a denominator. The eventual-positive extension retains the finite prefix. \\
\phantomsection\label{entry:symbol-registry-01334}\(\nu,g,B,U_k,f,U_j,V_j,\varepsilon,A_n,H_j\) & \texttt{\detokenize{\nu,}\allowbreak\detokenize{g,}\allowbreak\detokenize{B,}\allowbreak\detokenize{U_k,}\allowbreak\detokenize{f,}\allowbreak\detokenize{U_j,}\allowbreak\detokenize{V_j,}\allowbreak\detokenize{\varepsilon,}\allowbreak\detokenize{A_n,}\allowbreak\detokenize{H_j}} & observation laws functions events uniform inputs densities fair signs rare jumps and harmonic sums & P15 strong-laws-of-large-numbers & meta & \hyperref[sec:p15-strong-laws-of-large-numbers]{\ref*{sec:p15-strong-laws-of-large-numbers}} & \thepage & Uniform inputs and fair signs have distinct locally declared roles. The three-sign blocks have pairwise but not mutual independence. \\
\phantomsection\label{entry:symbol-registry-01335}\(X_i^{[M]},S_n^{[M]},Z_n,Z,J_m,K_m,B_n,a_n,H\) & \texttt{\detokenize{X_i^{[M]},}\allowbreak\detokenize{S_n^{[M]},}\allowbreak\detokenize{Z_n,}\allowbreak\detokenize{Z,}\allowbreak\detokenize{J_m,}\allowbreak\detokenize{K_m,}\allowbreak\detokenize{B_n,}\allowbreak\detokenize{a_n,}\allowbreak\detokenize{H}} & extended-mean clipping random-index limits finite indices extended first hits and Bernoulli events & P15 strong-law-applications-chapter & meta & \hyperref[sec:p15-strong-law-applications-chapter]{\ref*{sec:p15-strong-law-applications-chapter}} & \thepage & Undefined expectations and infinite limits are distinguished. Random indices are finite where the theorem requires this. \\
\phantomsection\label{entry:symbol-registry-01336}\(T_n,T_0,\mathbb N_0,N(t),\mu,\Omega_0\) & \texttt{\detokenize{T_n,}\allowbreak\detokenize{T_0,}\allowbreak\detokenize{\mathbb N_0,}\allowbreak\detokenize{N(t),}\allowbreak\detokenize{\mu,}\allowbreak\detokenize{\Omega_0}} & renewal epochs nonnegative indices last tied count extended mean and common event & P15 strong-law-applications-chapter & meta & \hyperref[sec:p15-strong-law-applications-chapter]{\ref*{sec:p15-strong-law-applications-chapter}} & \thepage & Zero waits are allowed with nonzero probability of a positive wait. The all-zero law is separate. \\
\phantomsection\label{entry:symbol-registry-01337}\(\nu,\nu_n,F,F_n,D_n,x_{j,k},G_n,N_k,n_k,\Omega_1\) & \texttt{\detokenize{\nu,}\allowbreak\detokenize{\nu_n,}\allowbreak\detokenize{F,}\allowbreak\detokenize{F_n,}\allowbreak\detokenize{D_n,}\allowbreak\detokenize{x_{j,k},}\allowbreak\detokenize{G_n,}\allowbreak\detokenize{N_k,}\allowbreak\detokenize{n_k,}\allowbreak\detokenize{\Omega_1}} & laws empirical measures CDFs rational supremum quantile grids thresholds and common event & P15 strong-law-applications-chapter & meta & \hyperref[sec:p15-strong-law-applications-chapter]{\ref*{sec:p15-strong-law-applications-chapter}} & \thepage & Both left limits and actual grid points are included. Repeated quantiles and both tails have explicit cases. \\
\phantomsection\label{entry:symbol-registry-01338}\(\mathcal A,\mathcal A_+,m,p,w,W_n,\pi_n,I_0,H,\sigma_I^2,L_n,T_n(\epsilon),\eta_n,\delta,C,\bot\) & \texttt{\detokenize{\mathcal A,}\allowbreak\detokenize{\mathcal A_+,}\allowbreak\detokenize{m,}\allowbreak\detokenize{p,}\allowbreak\detokenize{w,}\allowbreak\detokenize{W_n,}\allowbreak\detokenize{\pi_n,}\allowbreak\detokenize{I_0,}\allowbreak\detokenize{H,}\allowbreak\detokenize{\sigma_I^2,}\allowbreak\detokenize{L_n,}\allowbreak\detokenize{T_n,}\allowbreak\detokenize{\eta_n,}\allowbreak\detokenize{\delta,}\allowbreak\detokenize{C,}\allowbreak\detokenize{\bot}} & finite alphabet support words information entropy typical sets and error budgets & P15 strong-law-applications-chapter & meta & \hyperref[sec:p15-strong-law-applications-chapter]{\ref*{sec:p15-strong-law-applications-chapter}} & \thepage & Logarithms are natural. Zero probabilities use the positive support. The bottom glyph is an alphabet symbol in the example. \\
\phantomsection\label{entry:symbol-registry-01339}\((\Omega,\mathcal F,P),\mathcal F_n,\mathcal G_n,\mathcal F_n^X,\Delta X_n,\mathcal H_n\) & \texttt{\detokenize{\Omega,}\allowbreak\detokenize{\mathcal F,}\allowbreak\detokenize{P,}\allowbreak\detokenize{\mathcal F_n,}\allowbreak\detokenize{\mathcal G_n,}\allowbreak\detokenize{\mathcal F_n^X,}\allowbreak\detokenize{\Delta X_n,}\allowbreak\detokenize{\mathcal H_n}} & probability space information filtrations natural filtration increments and independent-product filtration & P15 martingales-and-convergence & meta & \hyperref[sec:p15-martingales-and-convergence]{\ref*{sec:p15-martingales-and-convergence}} & \thepage & Natural information includes X0. Product inputs and random-walk increments have separate filtrations. \\
\phantomsection\label{entry:symbol-registry-01340}\(\xi_i,\eta_i,S_n,R_n,c,\mu,\sigma^2,M_n,\varphi(\theta),\theta\) & \texttt{\detokenize{\xi_i,}\allowbreak\detokenize{\eta_i,}\allowbreak\detokenize{S_n,}\allowbreak\detokenize{R_n,}\allowbreak\detokenize{c,}\allowbreak\detokenize{\mu,}\allowbreak\detokenize{\sigma^2,}\allowbreak\detokenize{M_n,}\allowbreak\detokenize{\varphi,}\allowbreak\detokenize{\theta}} & increments walks initial value mean variance product and exponential normalizer & P15 martingales-and-convergence & meta & \hyperref[sec:p15-martingales-and-convergence]{\ref*{sec:p15-martingales-and-convergence}} & \thepage & Theta is real and the normalizer is finite and positive. Shifted and unshifted exponential martingales have distinct initial values. \\
\phantomsection\label{entry:symbol-registry-01341}\(H_i,K_i,T,S,\mathcal F_T,X^T,T_j,U_n,U_\infty,a,b,Y_n\) & \texttt{\detokenize{H_i,}\allowbreak\detokenize{K_i,}\allowbreak\detokenize{T,}\allowbreak\detokenize{S,}\allowbreak\detokenize{\mathcal F_T,}\allowbreak\detokenize{X^T,}\allowbreak\detokenize{T_j,}\allowbreak\detokenize{U_n,}\allowbreak\detokenize{U_\infty,}\allowbreak\detokenize{a,}\allowbreak\detokenize{b,}\allowbreak\detokenize{Y_n}} & predictable holdings stopping times stopped information entrance and exit times upcrossings and clipped process & P15 martingales-and-convergence & meta & \hyperref[sec:p15-martingales-and-convergence]{\ref*{sec:p15-martingales-and-convergence}} & \thepage & T0=-1 is only an initial marker. Literal stopped-field inclusion assumes pointwise time order. \\
\phantomsection\label{entry:symbol-registry-01342}\(C,L,X_\infty,B_k,\varepsilon_k,D,T_K,c_K,Z_i^{(K)}\) & \texttt{\detokenize{C,}\allowbreak\detokenize{L,}\allowbreak\detokenize{X_\infty,}\allowbreak\detokenize{B_k,}\allowbreak\detokenize{\varepsilon_k,}\allowbreak\detokenize{D,}\allowbreak\detokenize{T_K,}\allowbreak\detokenize{c_K,}\allowbreak\detokenize{Z_i^{(K)}}} & moment bound limit Bernoulli inputs signs upper increment bound localization and logarithmic truncations & P15 martingales-and-convergence & meta & \hyperref[sec:p15-martingales-and-convergence]{\ref*{sec:p15-martingales-and-convergence}} & \thepage & Limit integrability and exclusion of both infinite limits are proved. Log zero has value negative infinity. \\
\phantomsection\label{entry:symbol-registry-01343}\(A_n,Z_n,X_n^{(1)},X_n^{(2)},Y_n,Z_n^{(j)},G_n,p\) & \texttt{\detokenize{A_n,}\allowbreak\detokenize{Z_n,}\allowbreak\detokenize{X_n^{(1)},}\allowbreak\detokenize{X_n^{(2)},}\allowbreak\detokenize{Y_n,}\allowbreak\detokenize{Z_n^{(j)},}\allowbreak\detokenize{G_n,}\allowbreak\detokenize{p}} & compensation products switched processes finite switch approximants stopped gains and success probability & P15 martingales-and-convergence & meta & \hyperref[sec:p15-martingales-and-convergence]{\ref*{sec:p15-martingales-and-convergence}} & \thepage & Each switching construction declares its local Y and Z. Dubins has 0<a<b and a positive integer crossing count. \\
\phantomsection\label{entry:symbol-registry-01344}\(S_n,S_n^{(c)},a_n,a_0,W,p,q,y,Y_m,T_n,\mu_m,I_j,d_m\) & \texttt{\detokenize{S_n,}\allowbreak\detokenize{S_n^{(c)},}\allowbreak\detokenize{a_n,}\allowbreak\detokenize{a_0,}\allowbreak\detokenize{W,}\allowbreak\detokenize{p,}\allowbreak\detokenize{q,}\allowbreak\detokenize{y,}\allowbreak\detokenize{Y_m,}\allowbreak\detokenize{T_n,}\allowbreak\detokenize{\mu_m,}\allowbreak\detokenize{I_j,}\allowbreak\detokenize{d_m}} & rate sums denominators powers truncations means power intervals and tail moments & P15 strong-law-rates-and-normalization & meta & \hyperref[sec:p15-strong-law-rates-and-normalization]{\ref*{sec:p15-strong-law-rates-and-normalization}} & \thepage & The three MZ ranges below two have separate hypotheses. Feller retains strict cutoffs and ratio monotonicity. \\
\phantomsection\label{entry:symbol-registry-01345}\(\sigma,\sigma^2,\epsilon,V_n,\Delta_n,b_n,q_n,Q,\nu,A_m,N,B,C,\beta,d,u,f,N_k,N_0,A_k\) & \texttt{\detokenize{\sigma,}\allowbreak\detokenize{\sigma^2,}\allowbreak\detokenize{\epsilon,}\allowbreak\detokenize{V_n,}\allowbreak\detokenize{\Delta_n,}\allowbreak\detokenize{b_n,}\allowbreak\detokenize{q_n,}\allowbreak\detokenize{Q,}\allowbreak\detokenize{\nu,}\allowbreak\detokenize{A_m,}\allowbreak\detokenize{N,}\allowbreak\detokenize{B,}\allowbreak\detokenize{C,}\allowbreak\detokenize{\beta,}\allowbreak\detokenize{d,}\allowbreak\detokenize{u,}\allowbreak\detokenize{f,}\allowbreak\detokenize{N_k,}\allowbreak\detokenize{N_0,}\allowbreak\detokenize{A_k}} & variance rate parameters Feller ratios tails absolute-value law interval bounds and staircase construction & P15 strong-law-rates-and-normalization & meta & \hyperref[sec:p15-strong-law-rates-and-normalization]{\ref*{sec:p15-strong-law-rates-and-normalization}} & \thepage & Feller intervals and staircase heights use A in separate local roles. Increasing denominators alone are insufficient. \\
\phantomsection\label{entry:symbol-registry-01357}\(I,X_i,M,\epsilon,\delta,C,\Phi,a_j,\log^+,\phi_M,\psi_M\) & \texttt{\detokenize{I,}\allowbreak\detokenize{X_i,}\allowbreak\detokenize{M,}\allowbreak\detokenize{\epsilon,}\allowbreak\detokenize{\delta,}\allowbreak\detokenize{C,}\allowbreak\detokenize{\Phi,}\allowbreak\detokenize{a_j,}\allowbreak\detokenize{\log^+,}\allowbreak\detokenize{\phi_M,}\allowbreak\detokenize{\psi_M}} & uniform-integrability family tail threshold error budget moment bound superlinear function levels clipping and continuous cutoff & P15 uniform-integrability-and-conditional-limits & meta & \hyperref[sec:p15-uniform-integrability-and-conditional-limits]{\ref*{sec:p15-uniform-integrability-and-conditional-limits}} & \thepage & Empty nonnegative suprema are zero. The convex converse is an explicit locally finite positive-part sum. \\
\phantomsection\label{entry:symbol-registry-01358}\(Z,\mathcal G,\mathcal F_n,\mathcal F_\infty,X_\infty,M_n,W_N,Y_n,Y\) & \texttt{\detokenize{Z,}\allowbreak\detokenize{\mathcal G,}\allowbreak\detokenize{\mathcal F_n,}\allowbreak\detokenize{\mathcal F_\infty,}\allowbreak\detokenize{X_\infty,}\allowbreak\detokenize{M_n,}\allowbreak\detokenize{W_N,}\allowbreak\detokenize{Y_n,}\allowbreak\detokenize{Y}} & closing variable information sigma-algebras terminal limit conditional martingale and oscillation envelope & P15 uniform-integrability-and-conditional-limits & meta & \hyperref[sec:p15-uniform-integrability-and-conditional-limits]{\ref*{sec:p15-uniform-integrability-and-conditional-limits}} & \thepage & The observable limit of a closing variable is its conditional expectation on the generated information. Exceptional-set versions are constructed measurably. \\
\phantomsection\label{entry:symbol-registry-01359}\(B_n,C_n,\varepsilon,I_{k,n},\mathcal D_n,f,K,g,s_m,t_m,l,r,m\) & \texttt{\detokenize{B_n,}\allowbreak\detokenize{C_n,}\allowbreak\detokenize{\varepsilon,}\allowbreak\detokenize{I_{k,n},}\allowbreak\detokenize{\mathcal D_n,}\allowbreak\detokenize{f,}\allowbreak\detokenize{K,}\allowbreak\detokenize{g,}\allowbreak\detokenize{s_m,}\allowbreak\detokenize{t_m,}\allowbreak\detokenize{l,}\allowbreak\detokenize{r,}\allowbreak\detokenize{m}} & counterexample Bernoulli variables sign dyadic atoms information functions Lipschitz bound Borel versions and interval endpoints & P15 uniform-integrability-and-conditional-limits & meta & \hyperref[sec:p15-uniform-integrability-and-conditional-limits]{\ref*{sec:p15-uniform-integrability-and-conditional-limits}} & \thepage & Lipschitz functions on [0,1) are extended to 1 before terminal slopes are used. Completed measurable functions use proved Borel versions. \\
\phantomsection\label{entry:symbol-registry-01360}\(\xi_i,\mu,\theta,D,\delta_K,Z_n,Z_0,\xi_{n,j},p_0,\mathcal H_n,\alpha,\beta,U_n\) & \texttt{\detokenize{\xi_i,}\allowbreak\detokenize{\mu,}\allowbreak\detokenize{\theta,}\allowbreak\detokenize{D,}\allowbreak\detokenize{\delta_K,}\allowbreak\detokenize{Z_n,}\allowbreak\detokenize{Z_0,}\allowbreak\detokenize{\xi_{n,j},}\allowbreak\detokenize{p_0,}\allowbreak\detokenize{\mathcal H_n,}\allowbreak\detokenize{\alpha,}\allowbreak\detokenize{\beta,}\allowbreak\detokenize{U_n}} & additive noise mean prior parameter hitting event escape bound offspring population law generation information and polarization inputs & P15 uniform-integrability-and-conditional-limits & meta & \hyperref[sec:p15-uniform-integrability-and-conditional-limits]{\ref*{sec:p15-uniform-integrability-and-conditional-limits}} & \thepage & The branching dichotomy needs positive zero-offspring probability but no offspring mean. The posterior theorem uses the specified additive model. \\
\phantomsection\label{entry:symbol-registry-01361}\(\mathcal G_n,\mathcal G_\infty,Y_n,Y_\infty,X_{-n},\mathcal F_{-n},U_n(a,b)\) & \texttt{\detokenize{\mathcal G_n,}\allowbreak\detokenize{\mathcal G_\infty,}\allowbreak\detokenize{Y_n,}\allowbreak\detokenize{Y_\infty,}\allowbreak\detokenize{X_{-n},}\allowbreak\detokenize{\mathcal F_{-n},}\allowbreak\detokenize{U_n(a,}\allowbreak\detokenize{b)}} & decreasing information reverse-time process intersection limit source negative indices and reversed-string upcrossings & P15 uniform-integrability-and-conditional-limits & meta & \hyperref[sec:p15-uniform-integrability-and-conditional-limits]{\ref*{sec:p15-uniform-integrability-and-conditional-limits}} & \thepage & The reverse-time tower direction is explicit. Finite reversed strings have increasing filtrations; their maximal upcrossing counts grow when an initial observation is added. \\
\phantomsection\label{entry:symbol-registry-01362}\(X_{n,k},S_n,\mu_i,\mu_n,b_n,\sigma_n^2,\overline X_{n,k},\overline S_n,a_n,u_n,w_n\) & \texttt{\detokenize{X_{n,k},}\allowbreak\detokenize{S_n,}\allowbreak\detokenize{\mu_i,}\allowbreak\detokenize{\mu_n,}\allowbreak\detokenize{b_n,}\allowbreak\detokenize{\sigma_n^2,}\allowbreak\detokenize{\overline X_{n,k},}\allowbreak\detokenize{\overline S_n,}\allowbreak\detokenize{a_n,}\allowbreak\detokenize{u_n,}\allowbreak\detokenize{w_n}} & weak-law rows sums means scales variances truncations centers and error terms & P15 weak-laws-and-triangular-arrays & meta & \hyperref[sec:p15-weak-laws-and-triangular-arrays]{\ref*{sec:p15-weak-laws-and-triangular-arrays}} & \thepage & Triangular arrays require independence within each row; dependence between different rows is permitted. \\
\phantomsection\label{entry:symbol-registry-01363}\(f_n,M_f,\omega_f,\lambda_n,A_{n,\epsilon},M_n,R_n\) & \texttt{\detokenize{f_n,}\allowbreak\detokenize{M_f,}\allowbreak\detokenize{\omega_f,}\allowbreak\detokenize{\lambda_n,}\allowbreak\detokenize{A_{n,\epsilon},}\allowbreak\detokenize{M_n,}\allowbreak\detokenize{R_n}} & Bernstein polynomial approximation bound modulus product volume annulus mean square radius and radius & P15 weak-laws-and-triangular-arrays & meta & \hyperref[sec:p15-weak-laws-and-triangular-arrays]{\ref*{sec:p15-weak-laws-and-triangular-arrays}} & \thepage & The polynomial estimate is uniform on the closed unit interval. The cube example states its probability and radial normalization. \\
\phantomsection\label{entry:symbol-registry-01364}\(\tau_k^n,D_{n,k},T_n,p_{n,k},G,q,\pi,i_k,j_k,I_{n,k},K_n,H_n,v_n\) & \texttt{\detokenize{\tau_k^n,}\allowbreak\detokenize{D_{n,k},}\allowbreak\detokenize{T_n,}\allowbreak\detokenize{p_{n,k},}\allowbreak\detokenize{G,}\allowbreak\detokenize{q,}\allowbreak\detokenize{\pi,}\allowbreak\detokenize{i_k,}\allowbreak\detokenize{j_k,}\allowbreak\detokenize{I_{n,k},}\allowbreak\detokenize{K_n,}\allowbreak\detokenize{H_n,}\allowbreak\detokenize{v_n}} & coupon hits durations geometric laws permutation construction closing indicators cycle counts and harmonic moments & P15 weak-laws-and-triangular-arrays & meta & \hyperref[sec:p15-weak-laws-and-triangular-arrays]{\ref*{sec:p15-weak-laws-and-triangular-arrays}} & \thepage & Null exceptional coupon paths have finite versions. The full permutation construction proves mutual independence of closing indicators. \\
\phantomsection\label{entry:symbol-registry-01365}\(A_i,N_n,r_n,q_n,s_n,T_k,A_k,A_k^+,A_k^-,R_k\) & \texttt{\detokenize{A_i,}\allowbreak\detokenize{N_n,}\allowbreak\detokenize{r_n,}\allowbreak\detokenize{q_n,}\allowbreak\detokenize{s_n,}\allowbreak\detokenize{T_k,}\allowbreak\detokenize{A_k,}\allowbreak\detokenize{A_k^+,}\allowbreak\detokenize{A_k^-,}\allowbreak\detokenize{R_k}} & occupancy events counts ball load probabilities and Levy crossing sums events and future increments & P15 weak-laws-and-triangular-arrays & meta & \hyperref[sec:p15-weak-laws-and-triangular-arrays]{\ref*{sec:p15-weak-laws-and-triangular-arrays}} & \thepage & Occupancy and maximal-inequality notation have separate local declarations. Both strict and non-strict crossing bounds are proved. \\
\phantomsection\label{entry:symbol-registry-01366}\(\widehat X_i,\widetilde X_i,\widehat S_n,\widetilde S_n,Q,Y_i,\eta,W,c_n,g,r\) & \texttt{\detokenize{\widehat X_i,}\allowbreak\detokenize{\widetilde X_i,}\allowbreak\detokenize{\widehat S_n,}\allowbreak\detokenize{\widetilde S_n,}\allowbreak\detokenize{Q,}\allowbreak\detokenize{Y_i,}\allowbreak\detokenize{\eta,}\allowbreak\detokenize{W,}\allowbreak\detokenize{c_n,}\allowbreak\detokenize{g,}\allowbreak\detokenize{r}} & finite-product copies product measure symmetric differences tail scales original absolute input centers and moment exponent & P15 weak-laws-and-triangular-arrays & meta & \hyperref[sec:p15-weak-laws-and-triangular-arrays]{\ref*{sec:p15-weak-laws-and-triangular-arrays}} & \thepage & The necessity proof distinguishes original and product probabilities. The lower-moment consequence requires 0<r<1. \\
\phantomsection\label{entry:symbol-registry-01381}\(F(x),F(x-),F\{dx\},-F,{}^\circ X,{}^\circ F,{}^\circ S,m,a,Q,p\) & \texttt{\detokenize{F,}\allowbreak\detokenize{F(x-),}\allowbreak\detokenize{F\{dx\},}\allowbreak\detokenize{-F,}\allowbreak\detokenize{{}^\circ X,}\allowbreak\detokenize{{}^\circ F,}\allowbreak\detokenize{{}^\circ S,}\allowbreak\detokenize{m,}\allowbreak\detokenize{a,}\allowbreak\detokenize{Q,}\allowbreak\detokenize{p}} & Feller CDF left limit measure reflection symmetrization medians radius product probability and tail bound & P15 feller-weak-law-presentation & meta & \hyperref[sec:p15-feller-weak-law-presentation]{\ref*{sec:p15-feller-weak-law-presentation}} & \thepage & Minus F is a reflected probability law. Reflection retains the left limit at atoms, and the independent product measure is distinguished from the original measure. \\
\phantomsection\label{entry:symbol-registry-01382}\(Y_i,S,J,M,R,G,q_>,q_\geq\) & \texttt{\detokenize{Y_i,}\allowbreak\detokenize{S,}\allowbreak\detokenize{J,}\allowbreak\detokenize{M,}\allowbreak\detokenize{R,}\allowbreak\detokenize{G,}\allowbreak\detokenize{q_>,}\allowbreak\detokenize{q_\geq}} & symmetric summands sum first-largest index selected signed term remainder common CDF and exact absolute tails & P15 feller-weak-law-presentation & meta & \hyperref[sec:p15-feller-weak-law-presentation]{\ref*{sec:p15-feller-weak-law-presentation}} & \thepage & The selected term and remainder need not be independent. The proof uses their joint sign symmetry and explicitly retains zero remainders. \\
\phantomsection\label{entry:symbol-registry-01383}\(a,b,u,M,h,x_j,\Delta F_j,\Delta u_j,A_N,B_N\) & \texttt{\detokenize{a,}\allowbreak\detokenize{b,}\allowbreak\detokenize{u,}\allowbreak\detokenize{M,}\allowbreak\detokenize{h,}\allowbreak\detokenize{x_j,}\allowbreak\detokenize{\Delta F_j,}\allowbreak\detokenize{\Delta u_j,}\allowbreak\detokenize{A_N,}\allowbreak\detokenize{B_N}} & Stieltjes interval endpoints smooth integrand derivative bound mesh partition increments and approximation sums & P15 feller-weak-law-presentation & meta & \hyperref[sec:p15-feller-weak-law-presentation]{\ref*{sec:p15-feller-weak-law-presentation}} & \thepage & The two finite telescoping sums retain their actual interval endpoints and both error bounds. \\
\phantomsection\label{entry:symbol-registry-01384}\(\tau(t),\sigma(t),\sigma_h(t),m(t),m_h(t),W,r(t),s,X_k',S_n',m_n',\mu_n,c_n,\lambda,h\) & \texttt{\detokenize{\tau,}\allowbreak\detokenize{\sigma,}\allowbreak\detokenize{\sigma_h,}\allowbreak\detokenize{m,}\allowbreak\detokenize{m_h,}\allowbreak\detokenize{W,}\allowbreak\detokenize{r,}\allowbreak\detokenize{s,}\allowbreak\detokenize{X_k',}\allowbreak\detokenize{S_n',}\allowbreak\detokenize{m_n',}\allowbreak\detokenize{\mu_n,}\allowbreak\detokenize{c_n,}\allowbreak\detokenize{\lambda,}\allowbreak\detokenize{h}} & Feller tail function closed and half-open truncated moments absolute input cutoff clipped sums centers and example density parameters & P15 feller-weak-law-presentation & meta & \hyperref[sec:p15-feller-weak-law-presentation]{\ref*{sec:p15-feller-weak-law-presentation}} & \thepage & The atom at minus t distinguishes closed and half-open formulas. The truncation probability bound uses the single-summand coefficient 1/(n epsilon squared). \\
\phantomsection\label{entry:symbol-registry-01385}\(X,X_k,F,S_n,q(s),\mu(s),v(s),\mu_<(s),q_{\geq}(s),\delta(s),\rho(s),a_n,s_0\) & \texttt{\detokenize{X,}\allowbreak\detokenize{X_k,}\allowbreak\detokenize{F,}\allowbreak\detokenize{S_n,}\allowbreak\detokenize{q,}\allowbreak\detokenize{\mu,}\allowbreak\detokenize{v,}\allowbreak\detokenize{\mu_<,}\allowbreak\detokenize{q_{\geq},}\allowbreak\detokenize{\delta,}\allowbreak\detokenize{\rho,}\allowbreak\detokenize{a_n,}\allowbreak\detokenize{s_0}} & positive-gain laws sums strict and closed tails truncated moments reciprocal ratios normalizers and positive anchor & P15 infinite-mean-weak-normalization & meta & \hyperref[sec:p15-infinite-mean-weak-normalization]{\ref*{sec:p15-infinite-mean-weak-normalization}} & \thepage & The source's strictly positive domain and the proved nonnegative extension are distinct. Nonzero normalizers giving relative limit one are eventually positive. \\
\phantomsection\label{entry:symbol-registry-01386}\(N_0,s_n,L,c,K(n),m(n),b_n,\overline X_{n,k},a_n,\ell(s)\) & \texttt{\detokenize{N_0,}\allowbreak\detokenize{s_n,}\allowbreak\detokenize{L,}\allowbreak\detokenize{c,}\allowbreak\detokenize{K,}\allowbreak\detokenize{m,}\allowbreak\detokenize{b_n,}\allowbreak\detokenize{\overline X_{n,k},}\allowbreak\detokenize{a_n,}\allowbreak\detokenize{\ell}} & anchored crossing normalizers slowly varying function dilation and Petersburg truncation data & P15 infinite-mean-weak-normalization & meta & \hyperref[sec:p15-infinite-mean-weak-normalization]{\ref*{sec:p15-infinite-mean-weak-normalization}} & \thepage & The exact normalizer equation uses the closed truncated mean and a positive lower anchor. The Petersburg center a\_n=n m(n) is local to its array calculation. \\
\phantomsection\label{entry:symbol-registry-01387}\(f_a,a,b,\log_2,\log\) & \texttt{\detokenize{f_a,}\allowbreak\detokenize{a,}\allowbreak\detokenize{b,}\allowbreak\detokenize{\log_2,}\allowbreak\detokenize{\log}} & Cauchy density scales base-two and natural logarithms & P15 infinite-mean-weak-normalization & meta & \hyperref[sec:p15-infinite-mean-weak-normalization]{\ref*{sec:p15-infinite-mean-weak-normalization}} & \thepage & Cauchy convolution is integrated on finite symmetric intervals before limits. Petersburg logarithms are explicitly base two. \\
\phantomsection\label{entry:symbol-registry-01388}\(\overline X_n,X_n^*,X_\infty^*,\|Y\|_p,\|Y\|_\infty,S,T,L,A,k,H_j\) & \texttt{\detokenize{\overline X_n,}\allowbreak\detokenize{X_n^*,}\allowbreak\detokenize{X_\infty^*,}\allowbreak\detokenize{\|Y\|_p,}\allowbreak\detokenize{\|Y\|_\infty,}\allowbreak\detokenize{S,}\allowbreak\detokenize{T,}\allowbreak\detokenize{L,}\allowbreak\detokenize{A,}\allowbreak\detokenize{k,}\allowbreak\detokenize{H_j}} & positive and absolute maxima norms stopping times glued time stopped event deterministic horizon and predictable holdings & P15 martingale-maximal-inequalities & meta & \hyperref[sec:p15-martingale-maximal-inequalities]{\ref*{sec:p15-martingale-maximal-inequalities}} & \thepage & Time zero is included. Gluing requires pointwise time order; the deterministic horizon bound may hold almost surely. \\
\phantomsection\label{entry:symbol-registry-01389}\(\lambda,M,Z,p,q,C,\xi_i,\sigma_i^2,S_n,v_n,T_0,T_m\) & \texttt{\detokenize{\lambda,}\allowbreak\detokenize{M,}\allowbreak\detokenize{Z,}\allowbreak\detokenize{p,}\allowbreak\detokenize{q,}\allowbreak\detokenize{C,}\allowbreak\detokenize{\xi_i,}\allowbreak\detokenize{\sigma_i^2,}\allowbreak\detokenize{S_n,}\allowbreak\detokenize{v_n,}\allowbreak\detokenize{T_0,}\allowbreak\detokenize{T_m}} & maximal threshold truncation cap truncated maximum conjugate exponents moment bound increments variances sums and hitting times & P15 martingale-maximal-inequalities & meta & \hyperref[sec:p15-martingale-maximal-inequalities]{\ref*{sec:p15-martingale-maximal-inequalities}} & \thepage & The Holder coefficient is p/(p-1). Division treats the zero truncated moment separately; the fair-walk reciprocal tail is at integer levels. \\
\phantomsection\label{entry:symbol-registry-01390}\(Y,\Delta X_j,\Delta Y_j,H_j,K_j,\widetilde X_n,\widetilde Y_n,b_j,\mathcal H,\mathcal G\) & \texttt{\detokenize{Y,}\allowbreak\detokenize{\Delta X_j,}\allowbreak\detokenize{\Delta Y_j,}\allowbreak\detokenize{H_j,}\allowbreak\detokenize{K_j,}\allowbreak\detokenize{\widetilde X_n,}\allowbreak\detokenize{\widetilde Y_n,}\allowbreak\detokenize{b_j,}\allowbreak\detokenize{\mathcal H,}\allowbreak\detokenize{\mathcal G}} & earlier-information variable increments predictable covariance weights transformed processes normalizing weights and nested fields & P15 martingale-maximal-inequalities & meta & \hyperref[sec:p15-martingale-maximal-inequalities]{\ref*{sec:p15-martingale-maximal-inequalities}} & \thepage & Both covariance processes use the same filtration. Orthogonality and weighted isometry do not assert independent increments. \\
\phantomsection\label{entry:symbol-registry-01391}\(Q_j,K,x,v,c,a,a_0,\log^+\) & \texttt{\detokenize{Q_j,}\allowbreak\detokenize{K,}\allowbreak\detokenize{x,}\allowbreak\detokenize{v,}\allowbreak\detokenize{c,}\allowbreak\detokenize{a,}\allowbreak\detokenize{a_0,}\allowbreak\detokenize{\log^+}} & compensated square overshoot bound radius variance optimizing shift entropy input and positive scale & P15 martingale-maximal-inequalities & meta & \hyperref[sec:p15-martingale-maximal-inequalities]{\ref*{sec:p15-martingale-maximal-inequalities}} & \thepage & Small-ball exit is strict. The entropy input a in its proof and the deterministic scale a in the final formula are separately declared; log-plus zero is zero. \\
\phantomsection\label{entry:symbol-registry-01392}\(\mathcal T,\varnothing,v,vi,|v|_{\mathrm{len}},\xi_v,V_n,Z_n,\mathcal H_n,\mathcal F_n^{\mathrm{row}}\) & \texttt{\detokenize{\mathcal T,}\allowbreak\detokenize{\varnothing,}\allowbreak\detokenize{v,}\allowbreak\detokenize{vi,}\allowbreak\detokenize{|v|_{\mathrm{len}},}\allowbreak\detokenize{\xi_v,}\allowbreak\detokenize{V_n,}\allowbreak\detokenize{Z_n,}\allowbreak\detokenize{\mathcal H_n,}\allowbreak\detokenize{\mathcal F_n^{\mathrm{row}}}} & genealogical labels empty root concatenation length offspring realized generations counts and tree/source-row information & P15 martingale-maximal-inequalities & meta & \hyperref[sec:p15-martingale-maximal-inequalities]{\ref*{sec:p15-martingale-maximal-inequalities}} & \thepage & Genealogical and row constructions have the same generation law and separately proved martingale identities. Length bars are explicitly sequence length. \\
\phantomsection\label{entry:symbol-registry-01393}\(p_k,\mu,f(s),W_n,W,D,\rho,\theta_n,\sigma^2,Z_n^{(i)},W_n^{(i)},W^{(i)},z\) & \texttt{\detokenize{p_k,}\allowbreak\detokenize{\mu,}\allowbreak\detokenize{f,}\allowbreak\detokenize{W_n,}\allowbreak\detokenize{W,}\allowbreak\detokenize{D,}\allowbreak\detokenize{\rho,}\allowbreak\detokenize{\theta_n,}\allowbreak\detokenize{\sigma^2,}\allowbreak\detokenize{Z_n^{(i)},}\allowbreak\detokenize{W_n^{(i)},}\allowbreak\detokenize{W^{(i)},}\allowbreak\detokenize{z}} & offspring masses mean generating function normalized population limit extinction probabilities variance descendant families and initial size & P15 martingale-maximal-inequalities & meta & \hyperref[sec:p15-martingale-maximal-inequalities]{\ref*{sec:p15-martingale-maximal-inequalities}} & \thepage & Normalization requires positive mu. The zero-mean case is treated separately, and the positive-limit theorem does not assume the unproved general logarithmic-moment criterion. \\
\phantomsection\label{entry:symbol-registry-01394}\(A_n,Y_n,k_n,\alpha_n,\varepsilon_1,\varepsilon_2,\xi_1,\xi_2\) & \texttt{\detokenize{A_n,}\allowbreak\detokenize{Y_n,}\allowbreak\detokenize{k_n,}\allowbreak\detokenize{\alpha_n,}\allowbreak\detokenize{\varepsilon_1,}\allowbreak\detokenize{\varepsilon_2,}\allowbreak\detokenize{\xi_1,}\allowbreak\detokenize{\xi_2}} & nested-event countermodel dyadic cut-cell data fair signs and dependent martingale differences & P15 martingale-maximal-inequalities & meta & \hyperref[sec:p15-martingale-maximal-inequalities]{\ref*{sec:p15-martingale-maximal-inequalities}} & \thepage & The negative submartingale has constant pth moments but no Lp limit. The bounded dyadic martingale's essential-supremum error remains 2/3. \\
\phantomsection\label{entry:symbol-registry-01401}\(\mu_n,F_n,X_n,\Rightarrow,\stackrel{d}{=},Q_F,F^{-1},a_F(u),b_F(u),Y_n,E\) & \texttt{\detokenize{\mu_n,}\allowbreak\detokenize{F_n,}\allowbreak\detokenize{X_n,}\allowbreak\detokenize{\Rightarrow,}\allowbreak\detokenize{\stackrel{d}{=},}\allowbreak\detokenize{Q_F,}\allowbreak\detokenize{F^{-1},}\allowbreak\detokenize{a_F,}\allowbreak\detokenize{b_F,}\allowbreak\detokenize{Y_n,}\allowbreak\detokenize{E}} & laws CDFs weak convergence equality in law inverse conventions plateau endpoints coupled variables and exceptional set & P15 weak-convergence-of-laws & meta & \hyperref[sec:p15-weak-convergence-of-laws]{\ref*{sec:p15-weak-convergence-of-laws}} & \thepage & The source strict inverse uses u in (0,1). Coupling preserves the marginal laws and need not preserve an original joint law. \\
\phantomsection\label{entry:symbol-registry-01402}\(g,g_{x,\epsilon},D_g,\overline A,A^\circ,\partial A,\Rightarrow_v,\beta,\gamma,G,n_k,\nu\) & \texttt{\detokenize{g,}\allowbreak\detokenize{g_{x,\epsilon},}\allowbreak\detokenize{D_g,}\allowbreak\detokenize{\overline A,}\allowbreak\detokenize{A^\circ,}\allowbreak\detokenize{\partial A,}\allowbreak\detokenize{\Rightarrow_v,}\allowbreak\detokenize{\beta,}\allowbreak\detokenize{\gamma,}\allowbreak\detokenize{G,}\allowbreak\detokenize{n_k,}\allowbreak\detokenize{\nu}} & test maps continuous ramps discontinuity set closure interior boundary source vague-CDF limit offsets diagonal and remaining measure & P15 weak-convergence-of-laws & meta & \hyperref[sec:p15-weak-convergence-of-laws]{\ref*{sec:p15-weak-convergence-of-laws}} & \thepage & A defective limit F has remaining-measure CDF F minus beta. The negative endpoint offset is escaped mass rather than part of that CDF. \\
\phantomsection\label{entry:symbol-registry-01403}\(M_\epsilon,R,\varphi,C,L,\rho,\alpha,\|\cdot\|_{\mathrm{TV}}\) & \texttt{\detokenize{M_\epsilon,}\allowbreak\detokenize{R,}\allowbreak\detokenize{\varphi,}\allowbreak\detokenize{C,}\allowbreak\detokenize{L,}\allowbreak\detokenize{\rho,}\allowbreak\detokenize{\alpha,}\allowbreak\detokenize{\|\cdot\|_{\mathrm{TV}}}} & tightness radii coercive function integral bound threshold Levy metric Ky Fan metric and event-supremum total variation & P15 weak-convergence-of-laws & meta & \hyperref[sec:p15-weak-convergence-of-laws]{\ref*{sec:p15-weak-convergence-of-laws}} & \thepage & Ky Fan distance is a metric on almost-sure equivalence classes. Its defining infimum is proved attainable. \\
\phantomsection\label{entry:symbol-registry-01404}\(\phi,\Phi,k_n,\ell,C_n,X_p,T_N,V_{n+1},Y_n\) & \texttt{\detokenize{\phi,}\allowbreak\detokenize{\Phi,}\allowbreak\detokenize{k_n,}\allowbreak\detokenize{\ell,}\allowbreak\detokenize{C_n,}\allowbreak\detokenize{X_p,}\allowbreak\detokenize{T_N,}\allowbreak\detokenize{V_{n+1},}\allowbreak\detokenize{Y_n}} & normal density CDF product-row length limit central coefficient geometric law birthday time median and scaled median & P15 weak-convergence-of-laws & meta & \hyperref[sec:p15-weak-convergence-of-laws]{\ref*{sec:p15-weak-convergence-of-laws}} & \thepage & The central-binomial and normal-limit arguments here are independent of Stirling. The median keeps its separate density and Scheffe route. \\
\phantomsection\label{entry:symbol-registry-01407}\(n!,\binom nk,\varepsilon_i,S_N,\mathcal L_N,B_N,p,q,\phi,\Phi\) & \texttt{\detokenize{n!,}\allowbreak\detokenize{\binom nk,}\allowbreak\detokenize{\varepsilon_i,}\allowbreak\detokenize{S_N,}\allowbreak\detokenize{\mathcal L_N,}\allowbreak\detokenize{B_N,}\allowbreak\detokenize{p,}\allowbreak\detokenize{q,}\allowbreak\detokenize{\phi,}\allowbreak\detokenize{\Phi}} & factorials binomial coefficients fair signs sums admissible lattice binomial law parameters and normal density and CDF & P15 de-moivre-laplace-and-stirling & meta & \hyperref[sec:p15-de-moivre-laplace-and-stirling]{\ref*{sec:p15-de-moivre-laplace-and-stirling}} & \thepage & The local approximation retains parity and actual lattice spacing. Fixed binomial parameters are distinct from a changing triangular-row parameter. \\
\phantomsection\label{entry:symbol-registry-01408}\(r_n,d_n,\Delta_n,C,A,c,u,H(u),R_{n,k},x_{n,k},K_A\) & \texttt{\detokenize{r_n,}\allowbreak\detokenize{d_n,}\allowbreak\detokenize{\Delta_n,}\allowbreak\detokenize{C,}\allowbreak\detokenize{A,}\allowbreak\detokenize{c,}\allowbreak\detokenize{u,}\allowbreak\detokenize{H(u),}\allowbreak\detokenize{R_{n,k},}\allowbreak\detokenize{x_{n,k},}\allowbreak\detokenize{K_A}} & Stirling remainder logarithmic differences normalization constants entropy expansion local remainder standardized point and bounded-window constant & P15 de-moivre-laplace-and-stirling & meta & \hyperref[sec:p15-de-moivre-laplace-and-stirling]{\ref*{sec:p15-de-moivre-laplace-and-stirling}} & \thepage & The Stirling constant and both remainder bounds are proved before they are used for the normal approximation. Repeated scalar letters have their displayed local proof scopes. \\
\phantomsection\label{entry:symbol-registry-01409}\(\ell,y,b(k;N,p),r,L_N,U_N,\lambda,T_n,\gamma(a),I(a),\sigma_N,D(t\Vert p),M_p,A_n,p_N\) & \texttt{\detokenize{\ell,}\allowbreak\detokenize{y,}\allowbreak\detokenize{b,}\allowbreak\detokenize{r,}\allowbreak\detokenize{L_N,}\allowbreak\detokenize{U_N,}\allowbreak\detokenize{\lambda,}\allowbreak\detokenize{T_n,}\allowbreak\detokenize{\gamma,}\allowbreak\detokenize{I,}\allowbreak\detokenize{\sigma_N,}\allowbreak\detokenize{D,}\allowbreak\detokenize{M_p,}\allowbreak\detokenize{A_n,}\allowbreak\detokenize{p_N}} & row parameters binomial masses ratio rounded endpoints Poisson parameter and count tail-rate functions standard deviation divergence entropy remainder and moving parameter & P15 de-moivre-laplace-and-stirling & meta & \hyperref[sec:p15-de-moivre-laplace-and-stirling]{\ref*{sec:p15-de-moivre-laplace-and-stirling}} & \thepage & Endpoint rounding is explicit. Degenerate binomial parameters and moving-parameter counterexamples retain their separate domains. \\
\phantomsection\label{entry:symbol-registry-01410}\(Y_n,d_n,D_n,M_n,\mathcal F_n,K,Z_n,U_n,V_n,C_n^U,C_n^V\) & \texttt{\detokenize{Y_n,}\allowbreak\detokenize{d_n,}\allowbreak\detokenize{D_n,}\allowbreak\detokenize{M_n,}\allowbreak\detokenize{\mathcal F_n,}\allowbreak\detokenize{K,}\allowbreak\detokenize{Z_n,}\allowbreak\detokenize{U_n,}\allowbreak\detokenize{V_n,}\allowbreak\detokenize{C_n^U,}\allowbreak\detokenize{C_n^V}} & adapted process conditional drift predictable sum martingale filtration deterministic increment bound and additive-error processes & P15 predictable-quadratic-variation & meta & \hyperref[sec:p15-predictable-quadratic-variation]{\ref*{sec:p15-predictable-quadratic-variation}} & \thepage & The predictable drift starts at zero. Error allowances are known before the next step and their almost-sure finite sum does not imply an integrable total. \\
\phantomsection\label{entry:symbol-registry-01411}\(X_n,\Delta X_n,a_n,A_n,[X]_n,T,H_n,T_c,A_\infty,X^*\) & \texttt{\detokenize{X_n,}\allowbreak\detokenize{\Delta X_n,}\allowbreak\detokenize{a_n,}\allowbreak\detokenize{A_n,}\allowbreak\detokenize{[X]_n,}\allowbreak\detokenize{T,}\allowbreak\detokenize{H_n,}\allowbreak\detokenize{T_c,}\allowbreak\detokenize{A_\infty,}\allowbreak\detokenize{X^*}} & centered square-integrable martingale increments conditional variance predictable and realized variation stopping time transform pre-overshoot stop and maximal process & P15 predictable-quadratic-variation & meta & \hyperref[sec:p15-predictable-quadratic-variation]{\ref*{sec:p15-predictable-quadratic-variation}} & \thepage & Predictable A and realized square variation are different processes. T\_c observes A\_(n+1) at time n and excludes the variance overshoot. \\
\phantomsection\label{entry:symbol-registry-01412}\(f,I,\varepsilon,f_\varepsilon,B_n,p_n,C_n,P_n,V_n\) & \texttt{\detokenize{f,}\allowbreak\detokenize{I,}\allowbreak\detokenize{\varepsilon,}\allowbreak\detokenize{f_\varepsilon,}\allowbreak\detokenize{B_n,}\allowbreak\detokenize{p_n,}\allowbreak\detokenize{C_n,}\allowbreak\detokenize{P_n,}\allowbreak\detokenize{V_n}} & normalizing function reciprocal-square integral logarithmic rate adaptive events conditional probabilities counts compensator and count-martingale variance & P15 predictable-quadratic-variation & meta & \hyperref[sec:p15-predictable-quadratic-variation]{\ref*{sec:p15-predictable-quadratic-variation}} & \thepage & The normalizing conclusion is restricted to the infinite variance-clock event. The logarithmic function is defined at every nonnegative argument. Count variance and compensator can have different finiteness regimes. \\
\phantomsection\label{entry:symbol-registry-01413}\(g,r,a,b,G_n,R_n,D_n,t_n,Y_n,\Theta,U_n\) & \texttt{\detokenize{g,}\allowbreak\detokenize{r,}\allowbreak\detokenize{a,}\allowbreak\detokenize{b,}\allowbreak\detokenize{G_n,}\allowbreak\detokenize{R_n,}\allowbreak\detokenize{D_n,}\allowbreak\detokenize{t_n,}\allowbreak\detokenize{Y_n,}\allowbreak\detokenize{\Theta,}\allowbreak\detokenize{U_n}} & Friedman urn initial weights reinforcement constants color weights green draws total weight centered fraction random propensity and independent uniforms & P15 predictable-quadratic-variation & meta & \hyperref[sec:p15-predictable-quadratic-variation]{\ref*{sec:p15-predictable-quadratic-variation}} & \thepage & Urn D is a draw count in its local scope. The zero same-color reinforcement case has its own proof. The Pólya boundary has zero opposite-color reinforcement. \\
\phantomsection\label{entry:symbol-registry-01414}\(S_n,N,H,\xi_n,\zeta_n,\eta_n,\epsilon_n,\pi\) & \texttt{\detokenize{S_n,}\allowbreak\detokenize{N,}\allowbreak\detokenize{H,}\allowbreak\detokenize{\xi_n,}\allowbreak\detokenize{\zeta_n,}\allowbreak\detokenize{\eta_n,}\allowbreak\detokenize{\epsilon_n,}\allowbreak\detokenize{\pi}} & fair walk initially known horizon first hit sparse increments reset values reset switches reset probabilities and three-point law & P15 predictable-quadratic-variation & meta & \hyperref[sec:p15-predictable-quadratic-variation]{\ref*{sec:p15-predictable-quadratic-variation}} & \thepage & Samplewise boundedness does not supply a deterministic increment bound. A finite square-root stopping moment need not give a finite first stopping moment. \\
\phantomsection\label{entry:symbol-registry-01415}\(\mathbb C,\Re z,\Im z,\overline z,|z|,e^{it},\varphi_X,\mu\) & \texttt{\detokenize{\mathbb C,}\allowbreak\detokenize{\Re,}\allowbreak\detokenize{\Im,}\allowbreak\detokenize{\overline z,}\allowbreak\detokenize{|z|,}\allowbreak\detokenize{e^{it},}\allowbreak\detokenize{\varphi_X,}\allowbreak\detokenize{\mu}} & complex numbers components conjugate modulus exponential characteristic function and probability law & P15 characteristic-functions-and-inversion & meta & \hyperref[sec:p15-characteristic-functions-and-inversion]{\ref*{sec:p15-characteristic-functions-and-inversion}} & \thepage & Complex modulus has the declared numerical role here. Characteristic function varphi is distinct from normal density phi. \\
\phantomsection\label{entry:symbol-registry-01416}\(I,f,g,u,v,q(s),E,\lambda_j\) & \texttt{\detokenize{I,}\allowbreak\detokenize{f,}\allowbreak\detokenize{g,}\allowbreak\detokenize{u,}\allowbreak\detokenize{v,}\allowbreak\detokenize{q,}\allowbreak\detokenize{E,}\allowbreak\detokenize{\lambda_j}} & parameter interval integrand derivative endpoints integrated derivative support and mixture weights & P15 characteristic-functions-and-inversion & meta & \hyperref[sec:p15-characteristic-functions-and-inversion]{\ref*{sec:p15-characteristic-functions-and-inversion}} & \thepage & Differentiation on an arbitrary measure space is reduced to a sigma-finite derivative support. The derivative is identically zero off that support. \\
\phantomsection\label{entry:symbol-registry-01417}\(\operatorname{sinc},S(A),R(\theta,T),K_{a,b},\operatorname{sgn},f_P,b+h\mathbb Z\) & \texttt{\detokenize{\operatorname{sinc},}\allowbreak\detokenize{S,}\allowbreak\detokenize{R,}\allowbreak\detokenize{K,}\allowbreak\detokenize{\operatorname{sgn},}\allowbreak\detokenize{f_P,}\allowbreak\detokenize{b+h\mathbb Z}} & continuous sinc extension Dirichlet integral scaled sine integral inversion kernel sign Polya density and affine lattice & P15 characteristic-functions-and-inversion & meta & \hyperref[sec:p15-characteristic-functions-and-inversion]{\ref*{sec:p15-characteristic-functions-and-inversion}} & \thepage & Finite inversion retains both endpoint half-atoms. Affine-lattice inversion retains its phase and the infinite integral is a symmetric improper limit. \\
\phantomsection\label{entry:symbol-registry-01418}\(\mu_n,\varphi_n,Q(h),c,n\) & \texttt{\detokenize{\mu_n,}\allowbreak\detokenize{\varphi_n,}\allowbreak\detokenize{Q,}\allowbreak\detokenize{c,}\allowbreak\detokenize{n}} & law sequence transform sequence real symmetric difference quotient finite one-sided limit and derivative or remainder order & P15 characteristic-functions-and-inversion & meta & \hyperref[sec:p15-characteristic-functions-and-inversion]{\ref*{sec:p15-characteristic-functions-and-inversion}} & \thepage & Continuity at zero supplies tightness before a limiting probability law is asserted. The second-order expansion requires only a finite second moment. \\
\phantomsection\label{entry:symbol-registry-01429}\(\mu,\nu,c,\mathcal F_n,\mu_n,\nu_n,X_n,X,\sim\) & \texttt{\detokenize{\mu,}\allowbreak\detokenize{\nu,}\allowbreak\detokenize{c,}\allowbreak\detokenize{\mathcal F_n,}\allowbreak\detokenize{\mu_n,}\allowbreak\detokenize{\nu_n,}\allowbreak\detokenize{X_n,}\allowbreak\detokenize{X,}\allowbreak\detokenize{\sim}} & finite target measure probability reference target mass generating filtration restrictions density martingale limsup and mutual absolute continuity & P15 density-martingales & meta & \hyperref[sec:p15-density-martingales]{\ref*{sec:p15-density-martingales}} & \thepage & The finite-time density versions are common to both measures up to their local null sets. Almost-sure assertions retain their stated probability measure. \\
\phantomsection\label{entry:symbol-registry-01430}\(\rho,\rho_n,Y_n,Z_n,Y,Z,W,\mu_r,\mu_s,I_{k,n},K_n,I_n(x_0),\theta\) & \texttt{\detokenize{\rho,}\allowbreak\detokenize{\rho_n,}\allowbreak\detokenize{Y_n,}\allowbreak\detokenize{Z_n,}\allowbreak\detokenize{Y,}\allowbreak\detokenize{Z,}\allowbreak\detokenize{W,}\allowbreak\detokenize{\mu_r,}\allowbreak\detokenize{\mu_s,}\allowbreak\detokenize{I_{k,n},}\allowbreak\detokenize{K_n,}\allowbreak\detokenize{I_n,}\allowbreak\detokenize{\theta}} & normalized mixture restrictions bounded densities limits reciprocal on positive set regular and singular components partition cells and atom mixing weight & P15 density-martingales & meta & \hyperref[sec:p15-density-martingales]{\ref*{sec:p15-density-martingales}} & \thepage & The mixture is normalized by c plus one and includes zero target mass. Singular mass is carried by the infinite-density event under the target measure. \\
\phantomsection\label{entry:symbol-registry-01431}\(\alpha_m,\beta_m,\xi_m,q_m,h_m,H_n,H_\infty,V_n,R_n,z_m,C\) & \texttt{\detokenize{\alpha_m,}\allowbreak\detokenize{\beta_m,}\allowbreak\detokenize{\xi_m,}\allowbreak\detokenize{q_m,}\allowbreak\detokenize{h_m,}\allowbreak\detokenize{H_n,}\allowbreak\detokenize{H_\infty,}\allowbreak\detokenize{V_n,}\allowbreak\detokenize{R_n,}\allowbreak\detokenize{z_m,}\allowbreak\detokenize{C}} & coordinate laws projections density square-root affinity finite and infinite affinity products root likelihood reciprocal likelihood logarithmic terms and finite-series convergence event & P15 density-martingales & meta & \hyperref[sec:p15-density-martingales]{\ref*{sec:p15-density-martingales}} & \thepage & The affinity criterion permits zero coordinate densities. The separate logarithmic tail argument requires equivalent coordinate laws and uses both reference and target probabilities. \\
\phantomsection\label{entry:symbol-registry-01432}\(a_m,b_m,\sigma_m,\tau_m,D,v,p_m,q_m,\varepsilon_m\) & \texttt{\detokenize{a_m,}\allowbreak\detokenize{b_m,}\allowbreak\detokenize{\sigma_m,}\allowbreak\detokenize{\tau_m,}\allowbreak\detokenize{D,}\allowbreak\detokenize{v,}\allowbreak\detokenize{p_m,}\allowbreak\detokenize{q_m,}\allowbreak\detokenize{\varepsilon_m}} & normal means standard deviations completed-square parameters Bernoulli success probabilities and coordinate perturbations & P15 density-martingales & meta & \hyperref[sec:p15-density-martingales]{\ref*{sec:p15-density-martingales}} & \thepage & Bernoulli q is a scalar success probability within that theorem. Its earlier density-function role belongs to the product-law construction. \\
\phantomsection\label{entry:symbol-registry-01437}\(g,r,c,B_j,S_n,G_n,R_n,X_n,\mathcal F_n,\alpha,\beta,t,(a)_k\) & \texttt{\detokenize{g,}\allowbreak\detokenize{r,}\allowbreak\detokenize{c,}\allowbreak\detokenize{B_j,}\allowbreak\detokenize{S_n,}\allowbreak\detokenize{G_n,}\allowbreak\detokenize{R_n,}\allowbreak\detokenize{X_n,}\allowbreak\detokenize{\mathcal F_n,}\allowbreak\detokenize{\alpha,}\allowbreak\detokenize{\beta,}\allowbreak\detokenize{t,}\allowbreak\detokenize{(a)_k}} & initial color weights reinforcement draw indicators count color weights fraction draw information beta shapes their sum and rising products & P15 polya-urns & meta & \hyperref[sec:p15-polya-urns]{\ref*{sec:p15-polya-urns}} & \thepage & Time n is after n completed draws. The positive-reinforcement formulas use alpha equals g divided by c and beta equals r divided by c; zero reinforcement is defined separately. \\
\phantomsection\label{entry:symbol-registry-01438}\(U_j,V_j,\Theta,C_j,L(b),X_\infty,\sim\) & \texttt{\detokenize{U_j,}\allowbreak\detokenize{V_j,}\allowbreak\detokenize{\Theta,}\allowbreak\detokenize{C_j,}\allowbreak\detokenize{L,}\allowbreak\detokenize{X_\infty,}\allowbreak\detokenize{\sim}} & independent uniform inputs beta mixing parameter mixture draw sequence measurable limiting-frequency map urn limit and distribution notation & P15 polya-urns & meta & \hyperref[sec:p15-polya-urns]{\ref*{sec:p15-polya-urns}} & \thepage & The complete path law carries the joint frequency representation. Parameter-dependent strong-law events are integrated as measurable sections. The local sim relation assigns a law to a random variable. \\
\phantomsection\label{entry:symbol-registry-01439}\(a,b,T,K_{n,m},k,l,p,\ell\) & \texttt{\detokenize{a,}\allowbreak\detokenize{b,}\allowbreak\detokenize{T,}\allowbreak\detokenize{K_{n,m},}\allowbreak\detokenize{k,}\allowbreak\detokenize{l,}\allowbreak\detokenize{p,}\allowbreak\detokenize{\ell}} & updated beta shapes total shape future count mixed-moment orders initial green proportion and endpoint interpolant & P15 polya-urns & meta & \hyperref[sec:p15-polya-urns]{\ref*{sec:p15-polya-urns}} & \thepage & Updated scalar b is distinct from a binary word in the earlier word-law scope. Predictive future counts retain their common random parameter and are beta-binomial. \\
\phantomsection\label{entry:symbol-registry-01440}\(H_U,H_F,O,\pi_U,\pi_F\) & \texttt{\detokenize{H_U,}\allowbreak\detokenize{H_F,}\allowbreak\detokenize{O,}\allowbreak\detokenize{\pi_U,}\allowbreak\detokenize{\pi_F}} & urn and fixed-fair models prior odds and model probabilities & P15 polya-urns & meta & \hyperref[sec:p15-polya-urns]{\ref*{sec:p15-polya-urns}} & \thepage & The Bayes factor is the ratio of the probabilities of the same observed data under the two models. A posterior model probability is distinct from a within-model success prediction. \\
\phantomsection\label{entry:symbol-registry-01441}\(p_m,A_m,D_n,h_n,T_N\) & \texttt{\detokenize{p_m,}\allowbreak\detokenize{A_m,}\allowbreak\detokenize{D_n,}\allowbreak\detokenize{h_n,}\allowbreak\detokenize{T_N}} & event probabilities occurrence events finite survival events scalar hazards and no-occurrence tail events & P15 density-and-product-exercises & meta & \hyperref[sec:p15-density-and-product-exercises]{\ref*{sec:p15-density-and-product-exercises}} & \thepage & The event-product proof retains the Borel-Cantelli route. Hazards are defined only while survival has positive probability and require no event independence. \\
\phantomsection\label{entry:symbol-registry-01442}\(\mathcal P_n,\mathcal F_n,I,J,W,X_n,Y_n,M,H_j,U_N(a,b),G_N\) & \texttt{\detokenize{\mathcal P_n,}\allowbreak\detokenize{\mathcal F_n,}\allowbreak\detokenize{I,}\allowbreak\detokenize{J,}\allowbreak\detokenize{W,}\allowbreak\detokenize{X_n,}\allowbreak\detokenize{Y_n,}\allowbreak\detokenize{M,}\allowbreak\detokenize{H_j,}\allowbreak\detokenize{U_N,}\allowbreak\detokenize{G_N}} & refining partitions their sigma-fields parent and child cells integrand ratios bounded process bound predictable holdings upcrossing count and trading sum & P15 density-and-product-exercises & meta & \hyperref[sec:p15-density-and-product-exercises]{\ref*{sec:p15-density-and-product-exercises}} & \thepage & Every conditional average is computed by finite-cell integration. The elementary upcrossing proof does not assume the general density theorem. \\
\phantomsection\label{entry:symbol-registry-01443}\(E_j,\mu,\nu,c,\rho,Y_n,Z_n,Y,Z,f,D_i,\lambda_i,f_i,N,A_{i,k}\) & \texttt{\detokenize{E_j,}\allowbreak\detokenize{\mu,}\allowbreak\detokenize{\nu,}\allowbreak\detokenize{c,}\allowbreak\detokenize{\rho,}\allowbreak\detokenize{Y_n,}\allowbreak\detokenize{Z_n,}\allowbreak\detokenize{Y,}\allowbreak\detokenize{Z,}\allowbreak\detokenize{f,}\allowbreak\detokenize{D_i,}\allowbreak\detokenize{\lambda_i,}\allowbreak\detokenize{f_i,}\allowbreak\detokenize{N,}\allowbreak\detokenize{A_{i,k}}} & generating sets target and reference measures target mass mixture bounded ratios limiting ratios density disjoint finite blocks normalized references local densities and uniqueness null sets & P15 density-and-product-exercises & meta & \hyperref[sec:p15-density-and-product-exercises]{\ref*{sec:p15-density-and-product-exercises}} & \thepage & The sigma-finite assembly rescales each block density by the reference block mass. Both nonfinite density sets are removed before subtraction in the uniqueness proof. \\
\phantomsection\label{entry:symbol-registry-01444}\(\alpha_m,\beta_m,a_m,\epsilon,F,\eta\) & \texttt{\detokenize{\alpha_m,}\allowbreak\detokenize{\beta_m,}\allowbreak\detokenize{a_m,}\allowbreak\detokenize{\epsilon,}\allowbreak\detokenize{F,}\allowbreak\detokenize{\eta}} & Bernoulli success probabilities affinity common interior bound finitely-many-successes event and common submeasure & P15 density-and-product-exercises & meta & \hyperref[sec:p15-density-and-product-exercises]{\ref*{sec:p15-density-and-product-exercises}} & \thepage & The endpoint criterion states local support separately from the affinity product. Square-summable differences without interior bounds can still give singular sequence laws. \\
\phantomsection\label{entry:symbol-registry-01451}\(I,J,T,\lambda,\chi_R,Q_F,x_{m,j},\mu_m,F_m,g,h\) & \texttt{\detokenize{I,}\allowbreak\detokenize{J,}\allowbreak\detokenize{T,}\allowbreak\detokenize{\lambda,}\allowbreak\detokenize{\chi_R,}\allowbreak\detokenize{Q_F,}\allowbreak\detokenize{x_{m,j},}\allowbreak\detokenize{\mu_m,}\allowbreak\detokenize{F_m,}\allowbreak\detokenize{g,}\allowbreak\detokenize{h}} & intervals monotone substitution Lebesgue measure continuous cutoffs quantile midpoint atoms approximating laws and controlling functions & P15 characteristic-function-criteria-and-moments & meta & \hyperref[sec:p15-characteristic-function-criteria-and-moments]{\ref*{sec:p15-characteristic-function-criteria-and-moments}} & \thepage & The midpoint approximation uses exactly m atoms. Weak integration retains the tail bounds needed for unbounded moment transfer. \\
\phantomsection\label{entry:symbol-registry-01452}\(L,\varphi_X,\Psi_n,z,w,f_n,(u)_+,S_{m,j},Z_n,N_j,A_n,D_n,a\) & \texttt{\detokenize{L,}\allowbreak\detokenize{\varphi_X,}\allowbreak\detokenize{\Psi_n,}\allowbreak\detokenize{z,}\allowbreak\detokenize{w,}\allowbreak\detokenize{f_n,}\allowbreak\detokenize{(u)_+,}\allowbreak\detokenize{S_{m,j},}\allowbreak\detokenize{Z_n,}\allowbreak\detokenize{N_j,}\allowbreak\detokenize{A_n,}\allowbreak\detokenize{D_n,}\allowbreak\detokenize{a}} & local complex logarithm transforms complex arguments uniform-sum densities positive parts partial series probability-Cauchy variables indices and Pitman centering functions & P15 characteristic-function-criteria-and-moments & meta & \hyperref[sec:p15-characteristic-function-criteria-and-moments]{\ref*{sec:p15-characteristic-function-criteria-and-moments}} & \thepage & The local logarithm has a controlled branch. Pitman's derivative sign is positive i times the real centering. \\
\phantomsection\label{entry:symbol-registry-01453}\(d,\eta,\nu,s,c,\nu_c,\beta,c_n,\alpha,\psi,A,\chi,v,\varphi_2\) & \texttt{\detokenize{d,}\allowbreak\detokenize{\eta,}\allowbreak\detokenize{\nu,}\allowbreak\detokenize{s,}\allowbreak\detokenize{c,}\allowbreak\detokenize{\nu_c,}\allowbreak\detokenize{\beta,}\allowbreak\detokenize{c_n,}\allowbreak\detokenize{\alpha,}\allowbreak\detokenize{\psi,}\allowbreak\detokenize{A,}\allowbreak\detokenize{\chi,}\allowbreak\detokenize{v,}\allowbreak\detokenize{\varphi_2}} & convex right derivative Stieltjes and weighted measures positive scale tail atom binomial coefficients stable parameters equality set cutoff and transform perturbation & P15 characteristic-function-criteria-and-moments & meta & \hyperref[sec:p15-characteristic-function-criteria-and-moments]{\ref*{sec:p15-characteristic-function-criteria-and-moments}} & \thepage & The convex measure may be infinite near zero. The finite-support approximation and direct mixture routes retain their separate constructions. \\
\phantomsection\label{entry:symbol-registry-01454}\(m_k,\nu_k,\Lambda,r,\theta,p,q,f_0,f_a,a_k,c_a,Y_a,\lambda,c_\lambda\) & \texttt{\detokenize{m_k,}\allowbreak\detokenize{\nu_k,}\allowbreak\detokenize{\Lambda,}\allowbreak\detokenize{r,}\allowbreak\detokenize{\theta,}\allowbreak\detokenize{p,}\allowbreak\detokenize{q,}\allowbreak\detokenize{f_0,}\allowbreak\detokenize{f_a,}\allowbreak\detokenize{a_k,}\allowbreak\detokenize{c_a,}\allowbreak\detokenize{Y_a,}\allowbreak\detokenize{\lambda,}\allowbreak\detokenize{c_\lambda}} & moments absolute moments essential bound growth and expansion parameters complex gamma variables lognormal families and counterexample normalizers & P15 characteristic-function-criteria-and-moments & meta & \hyperref[sec:p15-characteristic-function-criteria-and-moments]{\ref*{sec:p15-characteristic-function-criteria-and-moments}} & \thepage & Complex powers use the specified right-half-plane branch. Continuous and discrete lognormal families retain their exact parameter identifications. \\
\phantomsection\label{entry:symbol-registry-01455}\(M_n,a_n,B,x_N,y_N,L_N,C,\epsilon\) & \texttt{\detokenize{M_n,}\allowbreak\detokenize{a_n,}\allowbreak\detokenize{B,}\allowbreak\detokenize{x_N,}\allowbreak\detokenize{y_N,}\allowbreak\detokenize{L_N,}\allowbreak\detokenize{C,}\allowbreak\detokenize{\epsilon}} & quasi-analytic derivative bounds ratios envelope interval points lengths generalized-gamma normalization and independent symmetric sign & P15 characteristic-function-criteria-and-moments & meta & \hyperref[sec:p15-characteristic-function-criteria-and-moments]{\ref*{sec:p15-characteristic-function-criteria-and-moments}} & \thepage & The derivative-envelope lemma is proved before Carleman's criterion. The symmetric-square correspondence retains the atom at zero. \\
\phantomsection\label{entry:symbol-registry-01456}\(\phi_0,\Phi,\chi,\mu,\sigma,\rho,k_n,X_{n,j},S_n,v_{n,j},v_n,L_n(\epsilon),O_P,o_P\) & \texttt{\detokenize{\phi_0,}\allowbreak\detokenize{\Phi,}\allowbreak\detokenize{\chi,}\allowbreak\detokenize{\mu,}\allowbreak\detokenize{\sigma,}\allowbreak\detokenize{\rho,}\allowbreak\detokenize{k_n,}\allowbreak\detokenize{X_{n,j},}\allowbreak\detokenize{S_n,}\allowbreak\detokenize{v_{n,j},}\allowbreak\detokenize{v_n,}\allowbreak\detokenize{L_n,}\allowbreak\detokenize{O_P,}\allowbreak\detokenize{o_P}} & normal density CDF and variable moments finite row lengths row variables sums variances Lindeberg tail and probability-order notation & P15 central-limit-theorems & meta & \hyperref[sec:p15-central-limit-theorems]{\ref*{sec:p15-central-limit-theorems}} & \thepage & Rows require independence within each row; no independence between rows is imposed. All normalizing variances are positive where divided by. \\
\phantomsection\label{entry:symbol-registry-01457}\(z_j,w_j,\theta,b,c_n,c,\gamma,Z_n,R_n,A_n,B_n,a,d\) & \texttt{\detokenize{z_j,}\allowbreak\detokenize{w_j,}\allowbreak\detokenize{\theta,}\allowbreak\detokenize{b,}\allowbreak\detokenize{c_n,}\allowbreak\detokenize{c,}\allowbreak\detokenize{\gamma,}\allowbreak\detokenize{Z_n,}\allowbreak\detokenize{R_n,}\allowbreak\detokenize{A_n,}\allowbreak\detokenize{B_n,}\allowbreak\detokenize{a,}\allowbreak\detokenize{d}} & complex-product comparison variables perturbations random errors deterministic limits and uniform law distance & P15 central-limit-theorems & meta & \hyperref[sec:p15-central-limit-theorems]{\ref*{sec:p15-central-limit-theorems}} & \thepage & The source complex-product argument is retained separately from the preceding logarithm proof. Perturbation variables may be dependent. \\
\phantomsection\label{entry:symbol-registry-01458}\(D,Z_\lambda,\xi_j,X'_j,Y_j,U_j,V_j,R_j,\varepsilon_j,\tau_A^2,N_n,a_n,T_m,N(t),D_t,M_t\) & \texttt{\detokenize{D,}\allowbreak\detokenize{Z_\lambda,}\allowbreak\detokenize{\xi_j,}\allowbreak\detokenize{X'_j,}\allowbreak\detokenize{Y_j,}\allowbreak\detokenize{U_j,}\allowbreak\detokenize{V_j,}\allowbreak\detokenize{R_j,}\allowbreak\detokenize{\varepsilon_j,}\allowbreak\detokenize{\tau_A^2,}\allowbreak\detokenize{N_n,}\allowbreak\detokenize{a_n,}\allowbreak\detokenize{T_m,}\allowbreak\detokenize{N(t),}\allowbreak\detokenize{D_t,}\allowbreak\detokenize{M_t}} & worked counts Poisson variable sign inputs symmetrized copies magnitudes conditional signs truncated variance random indices and renewal quantities & P15 central-limit-theorems & meta & \hyperref[sec:p15-central-limit-theorems]{\ref*{sec:p15-central-limit-theorems}} & \thepage & Random-index results do not assume index independence. The source pairwise-sign recursion retains consistent dyadic indices. \\
\phantomsection\label{entry:symbol-registry-01459}\(N_r,\epsilon_n,Y_{n,j},Z_{n,j},R_{n,j},U_n,q_{n,j},q_n,r_l,i_l,b_n\) & \texttt{\detokenize{N_r,}\allowbreak\detokenize{\epsilon_n,}\allowbreak\detokenize{Y_{n,j},}\allowbreak\detokenize{Z_{n,j},}\allowbreak\detokenize{R_{n,j},}\allowbreak\detokenize{U_n,}\allowbreak\detokenize{q_{n,j},}\allowbreak\detokenize{q_n,}\allowbreak\detokenize{r_l,}\allowbreak\detokenize{i_l,}\allowbreak\detokenize{b_n}} & diagonal truncation indices levels centered retained and discarded rows and moment multiplicity parameters & P15 central-limit-theorems & meta & \hyperref[sec:p15-central-limit-theorems]{\ref*{sec:p15-central-limit-theorems}} & \thepage & All collision patterns are included in the moment proof. Centering and discarded-tail estimates are explicit. \\
\phantomsection\label{entry:symbol-registry-01460}\(B_j,H_n,V_j,A_j,A,Y_j,c_n,U_n,\beta,u_{n,j},f_t,Q\) & \texttt{\detokenize{B_j,}\allowbreak\detokenize{H_n,}\allowbreak\detokenize{V_j,}\allowbreak\detokenize{A_j,}\allowbreak\detokenize{A,}\allowbreak\detokenize{Y_j,}\allowbreak\detokenize{c_n,}\allowbreak\detokenize{U_n,}\allowbreak\detokenize{\beta,}\allowbreak\detokenize{u_{n,j},}\allowbreak\detokenize{f_t,}\allowbreak\detokenize{Q}} & record indicators harmonic sums ranks cycle variables three-series cutoffs and sparse-jump parameters & P15 central-limit-theorems & meta & \hyperref[sec:p15-central-limit-theorems]{\ref*{sec:p15-central-limit-theorems}} & \thepage & The three-series source route now uses a proved Lindeberg-Feller premise. All sparse-jump parameter regimes remain distinct. \\
\phantomsection\label{entry:symbol-registry-01461}\(p,\vartheta(x),C,A(x),P(x),\delta,s,P_n,E_n,\delta_p,g,g_n,\ell_n,\alpha_n,B_p,V_n,b_n,a_n,q\) & \texttt{\detokenize{p,}\allowbreak\detokenize{\vartheta,}\allowbreak\detokenize{C,}\allowbreak\detokenize{A,}\allowbreak\detokenize{P,}\allowbreak\detokenize{\delta,}\allowbreak\detokenize{s,}\allowbreak\detokenize{P_n,}\allowbreak\detokenize{E_n,}\allowbreak\detokenize{\delta_p,}\allowbreak\detokenize{g,}\allowbreak\detokenize{g_n,}\allowbreak\detokenize{\ell_n,}\allowbreak\detokenize{\alpha_n,}\allowbreak\detokenize{B_p,}\allowbreak\detokenize{V_n,}\allowbreak\detokenize{b_n,}\allowbreak\detokenize{a_n,}\allowbreak\detokenize{q}} & prime variables elementary prime sums uniform integer law divisibility indicators truncations and independent comparison parameters & P15 central-limit-theorems & meta & \hyperref[sec:p15-central-limit-theorems]{\ref*{sec:p15-central-limit-theorems}} & \thepage & Divisibility indicators under finite uniform sampling are not independent. The independent comparison has success probabilities one divided by the prime. \\
\phantomsection\label{entry:symbol-registry-01462}\(L,h_L,H_L,\omega_L,F,G,D,\eta,\eta_L,\lambda,x_0,e,q,K_i,k_i,\kappa_i,\varphi_F,\varphi_G\) & \texttt{\detokenize{L,}\allowbreak\detokenize{h_L,}\allowbreak\detokenize{H_L,}\allowbreak\detokenize{\omega_L,}\allowbreak\detokenize{F,}\allowbreak\detokenize{G,}\allowbreak\detokenize{D,}\allowbreak\detokenize{\eta,}\allowbreak\detokenize{\eta_L,}\allowbreak\detokenize{\lambda,}\allowbreak\detokenize{x_0,}\allowbreak\detokenize{e,}\allowbreak\detokenize{q,}\allowbreak\detokenize{K_i,}\allowbreak\detokenize{k_i,}\allowbreak\detokenize{\kappa_i,}\allowbreak\detokenize{\varphi_F,}\allowbreak\detokenize{\varphi_G}} & smoothing scales density law tail mass CDF discrepancy approximating points and inversion kernels and transforms & P15 central-limit-theorems & meta & \hyperref[sec:p15-central-limit-theorems]{\ref*{sec:p15-central-limit-theorems}} & \thepage & The smoothing law has infinite absolute first moment. Its application uses the proved integrable-quotient inversion premise. \\
\phantomsection\label{entry:symbol-registry-01463}\(S,\mathcal S,P,T_\pi,\mathcal E_n,\mathcal E,\mathcal T,I_{n,k},(n)_k,A_{n,k},D_{n,m}\) & \texttt{\detokenize{S,}\allowbreak\detokenize{\mathcal S,}\allowbreak\detokenize{P,}\allowbreak\detokenize{T_\pi,}\allowbreak\detokenize{\mathcal E_n,}\allowbreak\detokenize{\mathcal E,}\allowbreak\detokenize{\mathcal T,}\allowbreak\detokenize{I_{n,k},}\allowbreak\detokenize{(n)_k,}\allowbreak\detokenize{A_{n,k},}\allowbreak\detokenize{D_{n,m}}} & coordinate space law finite permutations invariant fields tail field ordered distinct indices falling products symmetrizations and deleted-block averages & P15 exchangeability-and-definetti & meta & \hyperref[sec:p15-exchangeability-and-definetti]{\ref*{sec:p15-exchangeability-and-definetti}} & \thepage & The law invariance and event invariance are distinct definitions. The local product is falling. Deletion gives an almost-sure tail-measurable version. \\
\phantomsection\label{entry:symbol-registry-01464}\(V,Y,\mathcal G,\mathcal H,\mathcal H_n,S_n,\xi_j,G,Q_m,Y_r,\tau,\alpha,\beta\) & \texttt{\detokenize{V,}\allowbreak\detokenize{Y,}\allowbreak\detokenize{\mathcal G,}\allowbreak\detokenize{\mathcal H,}\allowbreak\detokenize{\mathcal H_n,}\allowbreak\detokenize{S_n,}\allowbreak\detokenize{\xi_j,}\allowbreak\detokenize{G,}\allowbreak\detokenize{Q_m,}\allowbreak\detokenize{Y_r,}\allowbreak\detokenize{\tau,}\allowbreak\detokenize{\alpha,}\allowbreak\detokenize{\beta}} & conditional-projection variables partial-sum fields ballot sums and event conditional final-sum law reverse martingale stop and integer vote totals & P15 exchangeability-and-definetti & meta & \hyperref[sec:p15-exchangeability-and-definetti]{\ref*{sec:p15-exchangeability-and-definetti}} & \thepage & The ballot proof conditions before stopping so the source's moment-free integer hypothesis is retained. \\
\phantomsection\label{entry:symbol-registry-01465}\(\mathscr P(S),e_B,\mathcal C,K_\omega,\Lambda,J,\mathcal G^P,\sigma(K),h_A\) & \texttt{\detokenize{\mathscr P(S),}\allowbreak\detokenize{e_B,}\allowbreak\detokenize{\mathcal C,}\allowbreak\detokenize{K_\omega,}\allowbreak\detokenize{\Lambda,}\allowbreak\detokenize{J,}\allowbreak\detokenize{\mathcal G^P,}\allowbreak\detokenize{\sigma(K),}\allowbreak\detokenize{h_A}} & probability-law space evaluations determining class directing kernel mixing law full product kernel null-event saturation and invariant-event probability map & P15 exchangeability-and-definetti & meta & \hyperref[sec:p15-exchangeability-and-definetti]{\ref*{sec:p15-exchangeability-and-definetti}} & \thepage & The probability-valued representation has a standard Borel target. Invariant and tail information agree only modulo the stated ambient measurable null events. \\
\phantomsection\label{entry:symbol-registry-01466}\(m_f,\widehat\mu_n,\widehat F_n,F_\omega,\Theta,\lambda,H_{m,N}^{(k)},I_j,J_j,D,\lambda_N\) & \texttt{\detokenize{m_f,}\allowbreak\detokenize{\widehat\mu_n,}\allowbreak\detokenize{\widehat F_n,}\allowbreak\detokenize{F_\omega,}\allowbreak\detokenize{\Theta,}\allowbreak\detokenize{\lambda,}\allowbreak\detokenize{H_{m,N}^{(k)},}\allowbreak\detokenize{I_j,}\allowbreak\detokenize{J_j,}\allowbreak\detokenize{D,}\allowbreak\detokenize{\lambda_N}} & conditional means empirical laws and CDFs binary mixing parameter without-replacement laws coupled indices collision-free event and frequency laws & P15 exchangeability-and-definetti & meta & \hyperref[sec:p15-exchangeability-and-definetti]{\ref*{sec:p15-exchangeability-and-definetti}} & \thepage & One rational determining class supports empirical weak convergence. The finite-urn bound uses event-supremum total variation and requires k no greater than N. \\
\phantomsection\label{entry:symbol-registry-01467}\(M,M_b,g_b,Z^{(b)},Y_n^{(b)},a_n,Z_n,W,R_N\) & \texttt{\detokenize{M,}\allowbreak\detokenize{M_b,}\allowbreak\detokenize{g_b,}\allowbreak\detokenize{Z^{(b)},}\allowbreak\detokenize{Y_n^{(b)},}\allowbreak\detokenize{a_n,}\allowbreak\detokenize{Z_n,}\allowbreak\detokenize{W,}\allowbreak\detokenize{R_N}} & conditional means clipped variables reverse-power approximations shrinking atom boundary changing integrands dominator and decreasing oscillation envelope & P15 exchangeability-and-definetti & meta & \hyperref[sec:p15-exchangeability-and-definetti]{\ref*{sec:p15-exchangeability-and-definetti}} & \thepage & L2 clipping licenses the unbounded moment product. Finite Lp reverse convergence does not extend to L-infinity. \\
\phantomsection\label{entry:symbol-registry-01477}\(X,X_j,X',S_n,q_Z,H_Z,W_Z\) & \texttt{\detokenize{X,}\allowbreak\detokenize{X_j,}\allowbreak\detokenize{X',}\allowbreak\detokenize{S_n,}\allowbreak\detokenize{q_Z,}\allowbreak\detokenize{H_Z,}\allowbreak\detokenize{W_Z}} & input law independent copy sum tail probability truncated second moment and integrated-tail function & P15 normal-domains-of-attraction & meta & \hyperref[sec:p15-normal-domains-of-attraction]{\ref*{sec:p15-normal-domains-of-attraction}} & \thepage & Truncated moments use the stated strict or weak cutoff; H\_Z(x-) is distinguished at atoms. \\
\phantomsection\label{entry:symbol-registry-01478}\(\mu,Y,V,g,b_n,a_n,d_n\) & \texttt{\detokenize{\mu,}\allowbreak\detokenize{Y,}\allowbreak\detokenize{V,}\allowbreak\detokenize{g,}\allowbreak\detokenize{b_n,}\allowbreak\detokenize{a_n,}\allowbreak\detokenize{d_n}} & finite mean centered input truncated variance ratio constructed scaling and arbitrary scaling and centering & P15 normal-domains-of-attraction & meta & \hyperref[sec:p15-normal-domains-of-attraction]{\ref*{sec:p15-normal-domains-of-attraction}} & \thepage & The constructed b\_n is supplied in the normalizer theorem. Arbitrary normalizing constants are not assumed to equal it. \\
\phantomsection\label{entry:symbol-registry-01479}\(\chi,\psi,T_{nj},m_n,U_{nj},M_n(r),K_A,\delta\) & \texttt{\detokenize{\chi,}\allowbreak\detokenize{\psi,}\allowbreak\detokenize{T_{nj},}\allowbreak\detokenize{m_n,}\allowbreak\detokenize{U_{nj},}\allowbreak\detokenize{M_n(r),}\allowbreak\detokenize{K_A,}\allowbreak\detokenize{\delta}} & normal law symmetric transform truncated row variables row means centered variables maximal tail and cutoff controls & P15 normal-domains-of-attraction & meta & \hyperref[sec:p15-normal-domains-of-attraction]{\ref*{sec:p15-normal-domains-of-attraction}} & \thepage & Normal domains include infinite-variance laws satisfying the proved slow-variation condition. \\
\phantomsection\label{entry:symbol-registry-01480}\(m(C),U(C),\mathcal H,C_n,\delta_n,N_k,M_k,B_n\) & \texttt{\detokenize{m(C),}\allowbreak\detokenize{U(C),}\allowbreak\detokenize{\mathcal H,}\allowbreak\detokenize{C_n,}\allowbreak\detokenize{\delta_n,}\allowbreak\detokenize{N_k,}\allowbreak\detokenize{M_k,}\allowbreak\detokenize{B_n}} & historical truncated mean variance cutoff ratio diagonal scales and source normalization & P15 normal-domains-of-attraction & meta & \hyperref[sec:p15-normal-domains-of-attraction]{\ref*{sec:p15-normal-domains-of-attraction}} & \thepage & The exact 1954 cutoff route is distinct from the broader general-array and canonical-measure developments. \\
\phantomsection\label{entry:symbol-registry-01481}\(\theta,v_{n,\theta},u,\omega_u,L,B_nu,\Delta,\Delta_h,m_N\) & \texttt{\detokenize{\theta,}\allowbreak\detokenize{v_{n,\theta},}\allowbreak\detokenize{u,}\allowbreak\detokenize{\omega_u,}\allowbreak\detokenize{L,}\allowbreak\detokenize{B_nu,}\allowbreak\detokenize{\Delta,}\allowbreak\detokenize{\Delta_h,}\allowbreak\detokenize{m_N}} & kernel parameter variance test modulus Lipschitz constant Bernstein approximation forward and normalized differences and truncated grid size & P15 feller-bernstein-hausdorff & meta & \hyperref[sec:p15-feller-bernstein-hausdorff]{\ref*{sec:p15-feller-bernstein-hausdorff}} & \thepage & Normalized differences include h to the negative order. In the exact grid theorem m\_N=floor(aN)<N. \\
\phantomsection\label{entry:symbol-registry-01482}\(a_r,a_{n,r},c_k,w_{n,j},E_n,L,u_{t,\varepsilon},U_\varepsilon,F,\mu\) & \texttt{\detokenize{a_r,}\allowbreak\detokenize{a_{n,r},}\allowbreak\detokenize{c_k,}\allowbreak\detokenize{w_{n,j},}\allowbreak\detokenize{E_n,}\allowbreak\detokenize{L,}\allowbreak\detokenize{u_{t,\varepsilon},}\allowbreak\detokenize{U_\varepsilon,}\allowbreak\detokenize{F,}\allowbreak\detokenize{\mu}} & positive series coefficients moment sequence Hausdorff weights discrete expectation positive functional ramp ramp value CDF and representing measure & P15 feller-bernstein-hausdorff & meta & \hyperref[sec:p15-feller-bernstein-hausdorff]{\ref*{sec:p15-feller-bernstein-hausdorff}} & \thepage & L here is a functional. The CDF may have mass at zero and inversion is asserted at continuity points. \\
\phantomsection\label{entry:symbol-registry-01483}\(t_0,a,b,m_{2n},p_{k,n},\Theta,\lambda,\alpha,\beta,\Delta_d\) & \texttt{\detokenize{t_0,}\allowbreak\detokenize{a,}\allowbreak\detokenize{b,}\allowbreak\detokenize{m_{2n},}\allowbreak\detokenize{p_{k,n},}\allowbreak\detokenize{\Theta,}\allowbreak\detokenize{\lambda,}\allowbreak\detokenize{\alpha,}\allowbreak\detokenize{\beta,}\allowbreak\detokenize{\Delta_d}} & support point local interval or mixture parameters even moments word probability directing parameter mixing law beta parameters and probability simplex & P15 feller-bernstein-hausdorff & meta & \hyperref[sec:p15-feller-bernstein-hausdorff]{\ref*{sec:p15-feller-bernstein-hausdorff}} & \thepage & Interval parameters and lognormal perturbation parameters are declared separately. Binary word probabilities precede count multiplicities. \\
\phantomsection\label{entry:symbol-registry-01484}\(Z_j,\overline Z,Y_j,\Lambda,q_\mu,\phi,c(t)\) & \texttt{\detokenize{Z_j,}\allowbreak\detokenize{\overline Z,}\allowbreak\detokenize{Y_j,}\allowbreak\detokenize{\Lambda,}\allowbreak\detokenize{q_\mu,}\allowbreak\detokenize{\phi,}\allowbreak\detokenize{c(t)}} & independent input sample mean residual mixture law component transform input transform and squared transform modulus & P15 feller-bernstein-hausdorff & meta & \hyperref[sec:p15-feller-bernstein-hausdorff]{\ref*{sec:p15-feller-bernstein-hausdorff}} & \thepage & The residual obstruction requires n at least 3 and a nondegenerate input law; it assumes no moments. The function c(t) is distinct from urn reinforcement. \\
\phantomsection\label{entry:symbol-registry-01485}\(\mathcal F_n,\mathcal F_T,L,M,T,X_\infty,X_T,Y_M,Z_n,R_n\) & \texttt{\detokenize{\mathcal F_n,}\allowbreak\detokenize{\mathcal F_T,}\allowbreak\detokenize{L,}\allowbreak\detokenize{M,}\allowbreak\detokenize{T,}\allowbreak\detokenize{X_\infty,}\allowbreak\detokenize{X_T,}\allowbreak\detokenize{Y_M,}\allowbreak\detokenize{Z_n,}\allowbreak\detokenize{R_n}} & filtrations stopping-time sigma-algebra ordered sampling times terminal and sampled values stopped processes and residual term & P15 optional-stopping-and-random-walks & meta & \hyperref[sec:p15-optional-stopping-and-random-walks]{\ref*{sec:p15-optional-stopping-and-random-walks}} & \thepage & Infinite-time values are assigned through the explicitly proved limits. Ordered stopping times can be infinite where stated. \\
\phantomsection\label{entry:symbol-registry-01486}\(K,\lambda,A,A_k,D,B,\xi_j,x,S_n,\mu,\sigma^2,\mathbb P_x,\mathbb E_x\) & \texttt{\detokenize{K,}\allowbreak\detokenize{\lambda,}\allowbreak\detokenize{A,}\allowbreak\detokenize{A_k,}\allowbreak\detokenize{D,}\allowbreak\detokenize{B,}\allowbreak\detokenize{\xi_j,}\allowbreak\detokenize{x,}\allowbreak\detokenize{S_n,}\allowbreak\detokenize{\mu,}\allowbreak\detokenize{\sigma^2,}\allowbreak\detokenize{\mathbb P_x,}\allowbreak\detokenize{\mathbb E_x}} & uniform-integrability cutoffs events majorants bounded increments starting state walk drift variance and started law and expectation & P15 optional-stopping-and-random-walks & meta & \hyperref[sec:p15-optional-stopping-and-random-walks]{\ref*{sec:p15-optional-stopping-and-random-walks}} & \thepage & Wald second-moment identities establish the required L2 limit before taking expectations. \\
\phantomsection\label{entry:symbol-registry-01487}\(T_a,\tau_{ab},G_T(s),\varphi(\theta),p,q,r,d,\rho,V_0,\eta_j,R_T\) & \texttt{\detokenize{T_a,}\allowbreak\detokenize{\tau_{ab},}\allowbreak\detokenize{G_T(s),}\allowbreak\detokenize{\varphi(\theta),}\allowbreak\detokenize{p,}\allowbreak\detokenize{q,}\allowbreak\detokenize{r,}\allowbreak\detokenize{d,}\allowbreak\detokenize{\rho,}\allowbreak\detokenize{V_0,}\allowbreak\detokenize{\eta_j,}\allowbreak\detokenize{R_T}} & level and interval hitting times passage generating function exponential transform walk probabilities ratios drifts stopped variance and increment variables & P15 optional-stopping-and-random-walks & meta & \hyperref[sec:p15-optional-stopping-and-random-walks]{\ref*{sec:p15-optional-stopping-and-random-walks}} & \thepage & The exponential transform is locally declared separately from a characteristic function. Defective hitting times retain their finite-time indicator. \\
\phantomsection\label{entry:symbol-registry-01488}\(f(s),C_k,u,\alpha,\beta,\theta_0,\delta,p_\delta,\zeta_n,c,G_n\) & \texttt{\detokenize{f(s),}\allowbreak\detokenize{C_k,}\allowbreak\detokenize{u,}\allowbreak\detokenize{\alpha,}\allowbreak\detokenize{\beta,}\allowbreak\detokenize{\theta_0,}\allowbreak\detokenize{\delta,}\allowbreak\detokenize{p_\delta,}\allowbreak\detokenize{\zeta_n,}\allowbreak\detokenize{c,}\allowbreak\detokenize{G_n}} & passage series Catalan numbers exponential variable roots integrability parameters insurance claims premium and doubling-game gain & P15 optional-stopping-and-random-walks & meta & \hyperref[sec:p15-optional-stopping-and-random-walks]{\ref*{sec:p15-optional-stopping-and-random-walks}} & \thepage & Capital and step assumptions are explicit in the insurance and skip-free formulas; exponential overshoots require the proved integrability bound. \\
\phantomsection\label{entry:symbol-registry-01489}\(s_j,s'_j,K,N_{n,z},p_{n,z},S_n,\xi_j,u_n,T_j,R,R_m,f_{2n}\) & \texttt{\detokenize{s_j,}\allowbreak\detokenize{s'_j,}\allowbreak\detokenize{K,}\allowbreak\detokenize{N_{n,z},}\allowbreak\detokenize{p_{n,z},}\allowbreak\detokenize{S_n,}\allowbreak\detokenize{\xi_j,}\allowbreak\detokenize{u_n,}\allowbreak\detokenize{T_j,}\allowbreak\detokenize{R,}\allowbreak\detokenize{R_m,}\allowbreak\detokenize{f_{2n}}} & path and reflection first zero counts fair masses walk increments return probabilities level hit first positive return successive returns and return masses & P15 reflection-and-ballot-counts & meta & \hyperref[sec:p15-reflection-and-ballot-counts]{\ref*{sec:p15-reflection-and-ballot-counts}} & \thepage & The inverse reflection starts at the correct positive endpoint. Return times and level-hit times have distinct finite laws. \\
\phantomsection\label{entry:symbol-registry-01490}\(\alpha,\beta,h,(E,\mathcal E),P,\mathbb P_a,X_n,\mathcal F_m,\tau_a^+,f_n,v_n,h_r,V(s),F(s)\) & \texttt{\detokenize{\alpha,}\allowbreak\detokenize{\beta,}\allowbreak\detokenize{h,}\allowbreak\detokenize{(E,}\allowbreak\detokenize{\mathcal E),}\allowbreak\detokenize{P,}\allowbreak\detokenize{\mathbb P_a,}\allowbreak\detokenize{X_n,}\allowbreak\detokenize{\mathcal F_m,}\allowbreak\detokenize{\tau_a^+,}\allowbreak\detokenize{f_n,}\allowbreak\detokenize{v_n,}\allowbreak\detokenize{h_r,}\allowbreak\detokenize{V(s),}\allowbreak\detokenize{F(s)}} & vote counts lead measurable state space kernel started Markov law filtration positive return masses transition iteration and renewal series & P15 reflection-and-ballot-counts & meta & \hyperref[sec:p15-reflection-and-ballot-counts]{\ref*{sec:p15-reflection-and-ballot-counts}} & \thepage & The singleton return state must be measurable. The renewal exercise uses a time-homogeneous Markov premise and first positive return. \\
\phantomsection\label{entry:symbol-registry-01491}\(L_{2n},A_{2n},\beta_{k,n},h,H,a_n,b_n,h_n,\delta\) & \texttt{\detokenize{L_{2n},}\allowbreak\detokenize{A_{2n},}\allowbreak\detokenize{\beta_{k,n},}\allowbreak\detokenize{h,}\allowbreak\detokenize{H,}\allowbreak\detokenize{a_n,}\allowbreak\detokenize{b_n,}\allowbreak\detokenize{h_n,}\allowbreak\detokenize{\delta}} & last zero positive-side edge occupation occupation probabilities arcsine density and CDF rounded endpoints step density and cutoff & P15 reflection-and-ballot-counts & meta & \hyperref[sec:p15-reflection-and-ballot-counts]{\ref*{sec:p15-reflection-and-ballot-counts}} & \thepage & Occupation counts edges including crossings adjacent to zero. Equality with the last-zero distribution is equality in law. \\
\phantomsection\label{entry:symbol-registry-01492}\(g,Y,\lambda,J(a),U_n,U,Z_n,R^{(j)},T_1^{(j)},b_m,w_j,W(z)\) & \texttt{\detokenize{g,}\allowbreak\detokenize{Y,}\allowbreak\detokenize{\lambda,}\allowbreak\detokenize{J(a),}\allowbreak\detokenize{U_n,}\allowbreak\detokenize{U,}\allowbreak\detokenize{Z_n,}\allowbreak\detokenize{R^{(j)},}\allowbreak\detokenize{T_1^{(j)},}\allowbreak\detokenize{b_m,}\allowbreak\detokenize{w_j,}\allowbreak\detokenize{W(z)}} & half-stable density limit variable Laplace parameter Gaussian integral compact transform variables nonnegative variables passage copies expansion term and finite mass generating series & P15 reflection-and-ballot-counts & meta & \hyperref[sec:p15-reflection-and-ballot-counts]{\ref*{sec:p15-reflection-and-ballot-counts}} & \thepage & The return and level-hit sums satisfy R\_m equal in law to m+T\_m and share the stated square-scale limit. \\
\phantomsection\label{entry:symbol-registry-01513}\(M_n,F,\phi,\Phi,\overline\Phi,\alpha,\beta,H_\alpha,H_\beta^-,G,u,b_n,\theta\) & \texttt{\detokenize{M_n,}\allowbreak\detokenize{F,}\allowbreak\detokenize{\phi,}\allowbreak\detokenize{\Phi,}\allowbreak\detokenize{\overline\Phi,}\allowbreak\detokenize{\alpha,}\allowbreak\detokenize{\beta,}\allowbreak\detokenize{H_\alpha,}\allowbreak\detokenize{H_\beta^-,}\allowbreak\detokenize{G,}\allowbreak\detokenize{u,}\allowbreak\detokenize{b_n,}\allowbreak\detokenize{\theta}} & maximum laws normal density and tails extreme-value shape parameters limit CDFs quantile level and Gaussian normalization & P15 weak-convergence-exercises & meta & \hyperref[sec:p15-weak-convergence-exercises]{\ref*{sec:p15-weak-convergence-exercises}} & \thepage & The Gaussian tail quantile is positive only from n=3. The three model limits do not assert the general domain-of-attraction classification. \\
\phantomsection\label{entry:symbol-registry-01514}\(E_n,f_n,\ell,F_n,\|\mu-\nu\|_{\mathrm{TV}}\) & \texttt{\detokenize{E_n,}\allowbreak\detokenize{f_n,}\allowbreak\detokenize{\ell,}\allowbreak\detokenize{F_n,}\allowbreak\detokenize{\|\mu-\nu\|_{\mathrm{TV}}}} & dyadic half-cells density cell length CDF and event-supremum total variation & P15 weak-convergence-exercises & meta & \hyperref[sec:p15-weak-convergence-exercises]{\ref*{sec:p15-weak-convergence-exercises}} & \thepage & The density values at both endpoints are fixed explicitly. Weak convergence coexists with constant total-variation distance one half. \\
\phantomsection\label{entry:symbol-registry-01515}\(\alpha(X,Y),\beta(X,Y),a,D,h(t)\) & \texttt{\detokenize{\alpha(X,}\allowbreak\detokenize{Y),}\allowbreak\detokenize{\beta(X,}\allowbreak\detokenize{Y),}\allowbreak\detokenize{a,}\allowbreak\detokenize{D,}\allowbreak\detokenize{h(t)}} & Ky Fan metric bounded expectation metric attained distance absolute difference and bounded increasing transform & P15 weak-convergence-exercises & meta & \hyperref[sec:p15-weak-convergence-exercises]{\ref*{sec:p15-weak-convergence-exercises}} & \thepage & Both metrics identify variables modulo almost-sure equality and require no moments. \\
\phantomsection\label{entry:symbol-registry-01516}\(\lambda_d,B_d,\mathbb S^{d-1},v_d,O,C(A),\sigma_d,\operatorname{area}_r,\pi\) & \texttt{\detokenize{\lambda_d,}\allowbreak\detokenize{B_d,}\allowbreak\detokenize{\mathbb S^{d-1},}\allowbreak\detokenize{v_d,}\allowbreak\detokenize{O,}\allowbreak\detokenize{C(A),}\allowbreak\detokenize{\sigma_d,}\allowbreak\detokenize{\operatorname{area}_r,}\allowbreak\detokenize{\pi}} & product volume unit ball sphere ball volume orthogonal map cones surface probability shell area and radial projection & P15 weak-convergence-exercises & meta & \hyperref[sec:p15-weak-convergence-exercises]{\ref*{sec:p15-weak-convergence-exercises}} & \thepage & Surface probability is defined geometrically before the Gaussian representation. The polar formula is proved from radial slabs. \\
\phantomsection\label{entry:symbol-registry-01517}\(X^{(n)},Y_j,D_n,Z^{(n)},R,U\) & \texttt{\detokenize{X^{(n)},}\allowbreak\detokenize{Y_j,}\allowbreak\detokenize{D_n,}\allowbreak\detokenize{Z^{(n)},}\allowbreak\detokenize{R,}\allowbreak\detokenize{U}} & surface vector independent Gaussian coordinates empirical squared norm coupled vector radius and direction & P15 weak-convergence-exercises & meta & \hyperref[sec:p15-weak-convergence-exercises]{\ref*{sec:p15-weak-convergence-exercises}} & \thepage & Almost-sure convergence belongs to the constructed coupling; the exercise concludes convergence of the original coordinate laws. \\
\phantomsection\label{entry:symbol-registry-01518}\(Y_n,\alpha,\beta,p,D(z),c_\varepsilon,K,y\) & \texttt{\detokenize{Y_n,}\allowbreak\detokenize{\alpha,}\allowbreak\detokenize{\beta,}\allowbreak\detokenize{p,}\allowbreak\detokenize{D(z),}\allowbreak\detokenize{c_\varepsilon,}\allowbreak\detokenize{K,}\allowbreak\detokenize{y}} & nonnegative moment variables exponent ratio convexity gap deviation constant fourth-moment bound and tail threshold & P15 weak-convergence-exercises & meta & \hyperref[sec:p15-weak-convergence-exercises]{\ref*{sec:p15-weak-convergence-exercises}} & \thepage & The exponents satisfy 0<alpha<beta. Negative thresholds and impossible K<1 inputs are handled separately. \\
\phantomsection\label{entry:symbol-registry-01531}\(\mathsf{At},\Omega,D_v,M_v,[S],\mathcal C(S)\) & \texttt{\detokenize{\mathsf{At},}\allowbreak\detokenize{\Omega,}\allowbreak\detokenize{D_v,}\allowbreak\detokenize{M_v,}\allowbreak\detokenize{[S],}\allowbreak\detokenize{\mathcal C(S)}} & finite ground atoms truth assignments state descriptions minimal clauses truth sets and content normal forms & P15 systematic-power & meta & \hyperref[sec:p15-systematic-power]{\ref*{sec:p15-systematic-power}} & \thepage & Empty conjunction is true and empty disjunction false; no domain closure is inferred from the names. \\
\phantomsection\label{entry:symbol-registry-01532}\(T,S,K,K'_1,K'_2,K_T,T^\flat,s_0\) & \texttt{\detokenize{T,}\allowbreak\detokenize{S,}\allowbreak\detokenize{K,}\allowbreak\detokenize{K'_1,}\allowbreak\detokenize{K'_2,}\allowbreak\detokenize{K_T,}\allowbreak\detokenize{T^\flat,}\allowbreak\detokenize{s_0}} & theory sentence data sentence finite data set content partition theory-entailed minimal clauses ground projection and original systematic score & P15 systematic-power & meta & \hyperref[sec:p15-systematic-power]{\ref*{sec:p15-systematic-power}} & \thepage & T may be quantified or inconsistent; the uniform singular formula uses the ground projection when needed. \\
\phantomsection\label{entry:symbol-registry-01533}\(r_1,g_1,r,g,p_r,s_r,\mathsf D\) & \texttt{\detokenize{r_1,}\allowbreak\detokenize{g_1,}\allowbreak\detokenize{r,}\allowbreak\detokenize{g,}\allowbreak\detokenize{p_r,}\allowbreak\detokenize{s_r,}\allowbreak\detokenize{\mathsf D}} & uniform and general range-content measures conditional logical probability generalized systematic power and Boolean dual transformation & P15 systematic-power & meta & \hyperref[sec:p15-systematic-power]{\ref*{sec:p15-systematic-power}} & \thepage & The p denominator requires positive range and the s denominator positive content; undefined inner factors are not multiplied by zero. \\
\phantomsection\label{entry:symbol-registry-01534}\(w_{00},w_{01},w_{10},w_{11},\theta,n\) & \texttt{\detokenize{w_{00},}\allowbreak\detokenize{w_{01},}\allowbreak\detokenize{w_{10},}\allowbreak\detokenize{w_{11},}\allowbreak\detokenize{\theta,}\allowbreak\detokenize{n}} & two-atom state weights latent probability parameter and number of ground atoms & P15 systematic-power & meta & \hyperref[sec:p15-systematic-power]{\ref*{sec:p15-systematic-power}} & \thepage & The explicit latent mixture is not identified with Carnap's separate measure. \\
\phantomsection\label{entry:symbol-registry-01540}\(\Theta,\mathcal T,\mathcal X,\mathcal A,\nu,\Pi,f_\theta,P_\theta,Q\) & \texttt{\detokenize{\Theta,}\allowbreak\detokenize{\mathcal T,}\allowbreak\detokenize{\mathcal X,}\allowbreak\detokenize{\mathcal A,}\allowbreak\detokenize{\nu,}\allowbreak\detokenize{\Pi,}\allowbreak\detokenize{f_\theta,}\allowbreak\detokenize{P_\theta,}\allowbreak\detokenize{Q}} & parameter and observation spaces reference measure prior sampling densities fixed-parameter sampling laws and joint prior law & P15 bayesian-updating-prediction & meta & \hyperref[sec:p15-bayesian-updating-prediction]{\ref*{sec:p15-bayesian-updating-prediction}} & \thepage & Likelihoods are finite jointly measurable and normalized under a sigma-finite observation measure. \\
\phantomsection\label{entry:symbol-registry-01541}\(m,G,\Pi_x\) & \texttt{\detokenize{m,}\allowbreak\detokenize{G,}\allowbreak\detokenize{\Pi_x}} & prior predictive density regular observation set and posterior kernel & P15 bayesian-updating-prediction & meta & \hyperref[sec:p15-bayesian-updating-prediction]{\ref*{sec:p15-bayesian-updating-prediction}} & \thepage & The ratio formula uses 0<m(x)<infinity; the chosen kernel equals the prior on the predictive-null exceptional set. \\
\phantomsection\label{entry:symbol-registry-01542}\(H,p,\ell_H,\ell_{H^c},B_H,Z_1,Z_2\) & \texttt{\detokenize{H,}\allowbreak\detokenize{p,}\allowbreak\detokenize{\ell_H,}\allowbreak\detokenize{\ell_{H^c},}\allowbreak\detokenize{B_H,}\allowbreak\detokenize{Z_1,}\allowbreak\detokenize{Z_2}} & parameter event prior event probability conditional-prior likelihoods odds multiplier and two inspection reports & P15 bayesian-updating-prediction & meta & \hyperref[sec:p15-bayesian-updating-prediction]{\ref*{sec:p15-bayesian-updating-prediction}} & \thepage & Conditional independence of reports and copying a report are different joint models. \\
\phantomsection\label{entry:symbol-registry-01543}\(L_n,m_n,\Pi_i,r_i,x_{1:n}\) & \texttt{\detokenize{L_n,}\allowbreak\detokenize{m_n,}\allowbreak\detokenize{\Pi_i,}\allowbreak\detokenize{r_i,}\allowbreak\detokenize{x_{1:n}}} & sequential joint likelihood predictive density posterior after a prefix predictive factor and observed path & P15 bayesian-updating-prediction & meta & \hyperref[sec:p15-bayesian-updating-prediction]{\ref*{sec:p15-bayesian-updating-prediction}} & \thepage & Every observed prefix in the update formula has a finite positive normalizer. \\
\phantomsection\label{entry:symbol-registry-01544}\(R,\widehat P_x,a,b,Y\) & \texttt{\detokenize{R,}\allowbreak\detokenize{\widehat P_x,}\allowbreak\detokenize{a,}\allowbreak\detokenize{b,}\allowbreak\detokenize{Y}} & future conditional kernel posterior predictive kernel conditional mean variance and future variable & P15 bayesian-updating-prediction & meta & \hyperref[sec:p15-bayesian-updating-prediction]{\ref*{sec:p15-bayesian-updating-prediction}} & \thepage & This kernel R is distinct from the local contamination prior also denoted R. \\
\phantomsection\label{entry:symbol-registry-01545}\(\alpha,\beta,a,b,t,k,n,r,J_r,K_n,B\) & \texttt{\detokenize{\alpha,}\allowbreak\detokenize{\beta,}\allowbreak\detokenize{a,}\allowbreak\detokenize{b,}\allowbreak\detokenize{t,}\allowbreak\detokenize{k,}\allowbreak\detokenize{n,}\allowbreak\detokenize{r,}\allowbreak\detokenize{J_r,}\allowbreak\detokenize{K_n,}\allowbreak\detokenize{B}} & beta shapes updated shapes total shape observed successes sample size future sample size future and observed counts and beta integral & P15 bayesian-updating-prediction & meta & \hyperref[sec:p15-bayesian-updating-prediction]{\ref*{sec:p15-bayesian-updating-prediction}} & \thepage & Shape parameters are positive; word and count likelihoods differ by their stated binomial coefficient. \\
\phantomsection\label{entry:symbol-registry-01546}\(\mu_0,\tau^2,\sigma^2,\bar x,S,s^2,u,v,\phi_V\) & \texttt{\detokenize{\mu_0,}\allowbreak\detokenize{\tau^2,}\allowbreak\detokenize{\sigma^2,}\allowbreak\detokenize{\bar x,}\allowbreak\detokenize{S,}\allowbreak\detokenize{s^2,}\allowbreak\detokenize{u,}\allowbreak\detokenize{v,}\allowbreak\detokenize{\phi_V}} & normal prior mean prior and sampling variances observed mean residual sum mean variance posterior mean variance and centered normal density & P15 bayesian-updating-prediction & meta & \hyperref[sec:p15-bayesian-updating-prediction]{\ref*{sec:p15-bayesian-updating-prediction}} & \thepage & Normal laws use variance as the second parameter; all density constants are retained. \\
\phantomsection\label{entry:symbol-registry-01547}\(g,g_\#\Pi_x,\eta,w,\Pi^w,m^w,\varepsilon,\lambda_x\) & \texttt{\detokenize{g,}\allowbreak\detokenize{g_\#\Pi_x,}\allowbreak\detokenize{\eta,}\allowbreak\detokenize{w,}\allowbreak\detokenize{\Pi^w,}\allowbreak\detokenize{m^w,}\allowbreak\detokenize{\varepsilon,}\allowbreak\detokenize{\lambda_x}} & parameter map posterior pushforward transformed quantity prior weight reweighted prior predictive density contamination fraction and posterior fraction & P15 bayesian-updating-prediction & meta & \hyperref[sec:p15-bayesian-updating-prediction]{\ref*{sec:p15-bayesian-updating-prediction}} & \thepage & Transporting a prior and choosing a new uniform prior are different operations. \\
\phantomsection\label{entry:symbol-registry-01548}\(C(x),c(\theta),\alpha,\Phi,z_{1-\alpha/2}\) & \texttt{\detokenize{C(x),}\allowbreak\detokenize{c(\theta),}\allowbreak\detokenize{\alpha,}\allowbreak\detokenize{\Phi,}\allowbreak\detokenize{z_{1-\alpha/2}}} & credible-set procedure fixed-parameter coverage tail probability normal CDF and quantile & P15 bayesian-updating-prediction & meta & \hyperref[sec:p15-bayesian-updating-prediction]{\ref*{sec:p15-bayesian-updating-prediction}} & \thepage & The graph of C is product measurable; posterior content calibrates prior-average coverage. \\
\phantomsection\label{entry:symbol-registry-01549}\(\mathcal D,L,\rho_x,c_{\rm FP},c_{\rm FN},\mathcal L_p(q)\) & \texttt{\detokenize{\mathcal D,}\allowbreak\detokenize{L,}\allowbreak\detokenize{\rho_x,}\allowbreak\detokenize{c_{\rm FP},}\allowbreak\detokenize{c_{\rm FN},}\allowbreak\detokenize{\mathcal L_p(q)}} & actions loss posterior risk false-positive and false-negative losses and expected logarithmic loss & P15 bayesian-updating-prediction & meta & \hyperref[sec:p15-bayesian-updating-prediction]{\ref*{sec:p15-bayesian-updating-prediction}} & \thepage & Finite-action integrated risk and pointwise squared-error minimization have distinct hypotheses. \\
\phantomsection\label{entry:symbol-registry-01550}\(z,c(z),c_\theta(z),L_z,N,a_i\) & \texttt{\detokenize{z,}\allowbreak\detokenize{c(z),}\allowbreak\detokenize{c_\theta(z),}\allowbreak\detokenize{L_z,}\allowbreak\detokenize{N,}\allowbreak\detokenize{a_i}} & record parameter-independent and parameter-dependent recording factors record likelihood first-success time and stopping probabilities & P15 bayesian-updating-prediction & meta & \hyperref[sec:p15-bayesian-updating-prediction]{\ref*{sec:p15-bayesian-updating-prediction}} & \thepage & Cancelling an observed-path factor does not assert a frequentist optional-stopping guarantee. \\
\phantomsection\label{entry:symbol-registry-01551}\(\mathcal K,K,H_k\) & \texttt{\detokenize{\mathcal K,}\allowbreak\detokenize{K,}\allowbreak\detokenize{H_k}} & finite model index set and model labels & P15 bayesian-model-comparison & meta & \hyperref[def:p15-model-comparison-experiment]{\ref*{def:p15-model-comparison-experiment}} & \thepage & Local context: sec:p15-bayesian-model-comparison. K+1 declared models; no nesting requirement. \\
\phantomsection\label{entry:symbol-registry-01552}\((\mathcal X,\mathcal A,\nu)\) & \texttt{\detokenize{\mathcal X,}\allowbreak\detokenize{\mathcal A,}\allowbreak\detokenize{\nu}} & common data space sigma algebra and sigma-finite reference measure & P15 bayesian-model-comparison & meta & \hyperref[def:p15-model-comparison-experiment]{\ref*{def:p15-model-comparison-experiment}} & \thepage & Local context: sec:p15-bayesian-model-comparison. The same observed quantity and dominating data measure are used across models. \\
\phantomsection\label{entry:symbol-registry-01553}\(\Theta_k,\mathcal T_k,\Pi_k\) & \texttt{\detokenize{\Theta_k,}\allowbreak\detokenize{\mathcal T_k,}\allowbreak\detokenize{\Pi_k}} & model-specific parameter space sigma algebra and proper prior & P15 bayesian-model-comparison & meta & \hyperref[def:p15-model-comparison-experiment]{\ref*{def:p15-model-comparison-experiment}} & \thepage & Local context: sec:p15-bayesian-model-comparison. Proper within-model priors are separate from prior mass on model labels. \\
\phantomsection\label{entry:symbol-registry-01554}\(f_k(x\mid\theta),m_k(x)\) & \texttt{\detokenize{f_k,}\allowbreak\detokenize{m_k}} & conditional data density and model marginal likelihood & P15 bayesian-model-comparison & meta & \hyperref[def:p15-model-comparison-experiment]{\ref*{def:p15-model-comparison-experiment}} & \thepage & Local context: sec:p15-bayesian-model-comparison. Joint likelihood measurability; density versions have almost-everywhere qualifications. \\
\phantomsection\label{entry:symbol-registry-01555}\(Q_k,Q_\alpha,e(x)\) & \texttt{\detokenize{Q_k,}\allowbreak\detokenize{Q_\alpha,}\allowbreak\detokenize{e}} & model predictive laws and prior-weighted predictive mixture density & P15 bayesian-model-comparison & meta & \hyperref[def:p15-model-comparison-experiment]{\ref*{def:p15-model-comparison-experiment}} & \thepage & Local context: sec:p15-bayesian-model-comparison. Zero prior terms are excluded before multiplying a possibly infinite density. \\
\phantomsection\label{entry:symbol-registry-01556}\(\alpha_k,w_k(x)\) & \texttt{\detokenize{\alpha_k,}\allowbreak\detokenize{w_k}} & prior and posterior model probabilities & P15 bayesian-model-comparison & meta & \hyperref[def:p15-model-comparison-experiment]{\ref*{def:p15-model-comparison-experiment}} & \thepage & Local context: sec:p15-bayesian-model-comparison. Posterior weights use the default alpha outside the predictive-full comparison domain. \\
\phantomsection\label{entry:symbol-registry-01557}\(\mathcal G\) & \texttt{\detokenize{\mathcal G}} & positive finite data comparison domain & P15 bayesian-model-comparison & meta & \hyperref[def:p15-model-comparison-experiment]{\ref*{def:p15-model-comparison-experiment}} & \thepage & Local context: sec:p15-bayesian-model-comparison. Not a sigma algebra here; all model marginal likelihoods finite and predictive evidence positive. \\
\phantomsection\label{entry:symbol-registry-01558}\(B_{ij}\) & \texttt{\detokenize{B_{ij}}} & Bayes factor with numerator model i and denominator model j & P15 bayesian-model-comparison & meta & \hyperref[def:p15-model-comparison-experiment]{\ref*{def:p15-model-comparison-experiment}} & \thepage & Local context: sec:p15-bayesian-model-comparison. Defined on G when the denominator marginal likelihood is positive; no ordinary value for 0/0. \\
\phantomsection\label{entry:symbol-registry-01559}\(Q_*,r_k\) & \texttt{\detokenize{Q_*,}\allowbreak\detokenize{r_k}} & equal-weight model predictive mixture and intrinsic Radon-Nikodym densities & P15 bayesian-model-comparison & meta & \hyperref[thm:p15-model-density-invariance]{\ref*{thm:p15-model-density-invariance}} & \thepage & Local context: sec:p15-model-density-invariance. r\_k=dQ\_k/dQ\_*; distinct from the scalar prior-scale ratio r. \\
\phantomsection\label{entry:symbol-registry-01560}\(h,T\) & \texttt{\detokenize{h,}\allowbreak\detokenize{T}} & data reference-measure multiplier and parameter-coordinate bijection & P15 bayesian-model-comparison & meta & \hyperref[thm:p15-model-density-invariance]{\ref*{thm:p15-model-density-invariance}} & \thepage & Local context: sec:p15-model-density-invariance. h is positive finite almost everywhere; T transports both likelihood and prior. \\
\phantomsection\label{entry:symbol-registry-01561}\(\Psi_k,\mathcal N_k,\psi,\lambda\) & \texttt{\detokenize{\Psi_k,}\allowbreak\detokenize{\mathcal N_k,}\allowbreak\detokenize{\psi,}\allowbreak\detokenize{\lambda}} & interest and nuisance parameter coordinates and spaces & P15 bayesian-model-comparison & meta & \hyperref[thm:p15-nuisance-integration]{\ref*{thm:p15-nuisance-integration}} & \thepage & Local context: sec:p15-nuisance-model-integration. Lambda as a nuisance coordinate is distinct from Lebesgue measure notation elsewhere. \\
\phantomsection\label{entry:symbol-registry-01562}\(\Lambda_k,\Pi_k^\psi,\bar f_k\) & \texttt{\detokenize{\Lambda_k,}\allowbreak\detokenize{\Pi_k^\psi,}\allowbreak\detokenize{\bar f_k}} & conditional nuisance-prior kernel marginal interest prior and integrated nuisance likelihood & P15 bayesian-model-comparison & meta & \hyperref[thm:p15-nuisance-integration]{\ref*{thm:p15-nuisance-integration}} & \thepage & Local context: sec:p15-nuisance-model-integration. The kernel is part of the declared prior; no automatic equality between models. \\
\phantomsection\label{entry:symbol-registry-01563}\(\theta_0,\mu_1,s^2,\tau^2\) & \texttt{\detokenize{\theta_0,}\allowbreak\detokenize{\mu_1,}\allowbreak\detokenize{s^2,}\allowbreak\detokenize{\tau^2}} & normal null value alternative prior center sampling variance and prior variance & P15 bayesian-model-comparison & meta & \hyperref[thm:p15-normal-point-null-factor]{\ref*{thm:p15-normal-point-null-factor}} & \thepage & Local context: sec:p15-normal-bayes-factor. Both variances positive; normal notation uses variance as its second parameter. \\
\phantomsection\label{entry:symbol-registry-01564}\(z,r\) & \texttt{\detokenize{z,}\allowbreak\detokenize{r}} & standardized normal observation and prior-to-sampling variance ratio & P15 bayesian-model-comparison & meta & \hyperref[thm:p15-normal-point-null-factor]{\ref*{thm:p15-normal-point-null-factor}} & \thepage & Local context: sec:p15-normal-bayes-factor. z=(x-theta\_0)/s and r=tau\^{}2/s\^{}2; z is real and |z| denotes absolute value. \\
\phantomsection\label{entry:symbol-registry-01565}\(\bar X,n,\sigma^2,\delta\) & \texttt{\detokenize{\bar X,}\allowbreak\detokenize{n,}\allowbreak\detokenize{\sigma^2,}\allowbreak\detokenize{\delta}} & sample mean size known observation variance and fixed departure & P15 bayesian-model-comparison & meta & \hyperref[thm:p15-normal-point-null-factor]{\ref*{thm:p15-normal-point-null-factor}} & \thepage & Local context: sec:p15-normal-bayes-factor. Fixed standardized and fixed nonzero-departure asymptotic paths are different. \\
\phantomsection\label{entry:symbol-registry-01566}\(\mathcal Z,T_k,R_k,R\) & \texttt{\detokenize{\mathcal Z,}\allowbreak\detokenize{T_k,}\allowbreak\detokenize{R_k,}\allowbreak\detokenize{R}} & common quantity space parameter-to-quantity maps and conditional/averaged posterior kernels & P15 bayesian-model-comparison & meta & \hyperref[def:p15-model-average-law]{\ref*{def:p15-model-average-law}} & \thepage & Local context: sec:p15-posterior-model-averaging. A common physical quantity is required; R\_k is not the nuisance-prior kernel Lambda\_k. \\
\phantomsection\label{entry:symbol-registry-01567}\(a_k,v_k,a\) & \texttt{\detokenize{a_k,}\allowbreak\detokenize{v_k,}\allowbreak\detokenize{a}} & component posterior means variances and mixture mean & P15 bayesian-model-comparison & meta & \hyperref[thm:p15-model-average-moments]{\ref*{thm:p15-model-average-moments}} & \thepage & Local context: sec:p15-posterior-model-averaging. Finite second moments are required only for positive posterior-weight components. \\
\phantomsection\label{entry:symbol-registry-01568}\(\widetilde\Pi_k,Z_k,c_k\) & \texttt{\detokenize{\widetilde\Pi_k,}\allowbreak\detokenize{Z_k,}\allowbreak\detokenize{c_k}} & improper parameter measure its integrated likelihood and arbitrary scaling constant & P15 bayesian-model-comparison & meta & \hyperref[thm:p15-improper-model-priors]{\ref*{thm:p15-improper-model-priors}} & \thepage & Local context: sec:p15-model-prior-conventions. Z\_k is not a proper prior predictive density. \\
\phantomsection\label{entry:symbol-registry-01569}\(\beta,\psi,R_a\) & \texttt{\detokenize{\beta,}\allowbreak\detokenize{\psi,}\allowbreak\detokenize{R_a}} & proper-prior null-constraint example coordinates and conditional versions & P15 bayesian-model-comparison & meta & \hyperref[thm:p15-improper-model-priors]{\ref*{thm:p15-improper-model-priors}} & \thepage & Local context: annex:p15.improper-model-priors. Independent uniform coordinates; version at psi=1/2 can be changed on that prior-null fiber. \\
\phantomsection\label{entry:symbol-registry-01570}\(p_{k\ell},a_{k\ell},t_{k\ell},t_*\) & \texttt{\detokenize{p_{k\ell},}\allowbreak\detokenize{a_{k\ell},}\allowbreak\detokenize{t_{k\ell},}\allowbreak\detokenize{t_*}} & finite parameter-prior masses weighted likelihood terms log terms and their maximum & P15 bayesian-model-comparison & meta & \hyperref[thm:p15-finite-model-evaluation]{\ref*{thm:p15-finite-model-evaluation}} & \thepage & Local context: sec:p15-finite-model-evaluation. Only finite positive terms are logged; zero terms contribute zero and are omitted. \\
\phantomsection\label{entry:symbol-registry-01571}\(\Phi,p_{\rm tail}\) & \texttt{\detokenize{\Phi,}\allowbreak\detokenize{p_{\rm tail}}} & standard normal CDF and two-sided null-tail probability & P15 bayesian-model-comparison & meta & \hyperref[thm:p15-normal-significance-comparison]{\ref*{thm:p15-normal-significance-comparison}} & \thepage & Local context: sec:p15-bayes-factor-significance. Tail calibration under H\_0 is distinct from posterior model probability. \\
\phantomsection\label{entry:symbol-registry-01572}\(S,S_\#\nu\) & \texttt{\detokenize{S,}\allowbreak\detokenize{S_\#\nu}} & observation-coordinate bijection and pushed-forward reference measure & P15 bayesian-model-comparison & meta & \hyperref[thm:p15-model-density-invariance]{\ref*{thm:p15-model-density-invariance}} & \thepage & Local context: sec:p15-model-density-invariance. Defined on measurable sets by (S\_\#nu)(A)=nu(S\^{}-1 A); model likelihoods are transported with the data. \\
\phantomsection\label{entry:symbol-registry-01573}\(\mathsf X,\mathscr X,\Theta,\Theta_0,\Theta_1,P_\theta,E_\theta\) & \texttt{\detokenize{\mathsf X,}\allowbreak\detokenize{\mathscr X,}\allowbreak\detokenize{\Theta,}\allowbreak\detokenize{\Theta_0,}\allowbreak\detokenize{\Theta_1,}\allowbreak\detokenize{P_\theta,}\allowbreak\detokenize{E_\theta}} & sample space sigma field parameter and hypothesis sets sampling laws and expectations & P15 hypothesis-tests-confidence & meta & \hyperref[sec:p15-hypothesis-tests-confidence]{\ref*{sec:p15-hypothesis-tests-confidence}} & \thepage & Tests and confidence procedures refer to repeated data under fixed-parameter laws. \\
\phantomsection\label{entry:symbol-registry-01574}\(\varphi,\psi,\varphi^*,\beta_\varphi,\alpha,U,\lambda_{01}\) & \texttt{\detokenize{\varphi,}\allowbreak\detokenize{\psi,}\allowbreak\detokenize{\varphi^*,}\allowbreak\detokenize{\beta_\varphi,}\allowbreak\detokenize{\alpha,}\allowbreak\detokenize{U,}\allowbreak\detokenize{\lambda_{01}}} & test competing test optimal test power significance level and auxiliary uniform randomization & P15 hypothesis-tests-confidence & meta & \hyperref[sec:p15-hypothesis-tests-confidence]{\ref*{sec:p15-hypothesis-tests-confidence}} & \thepage & The source's confidence coefficient alpha is renamed q or one minus alpha in this chapter. \\
\phantomsection\label{entry:symbol-registry-01575}\(\mu,f_0,f_1,L,c,\gamma\) & \texttt{\detokenize{\mu,}\allowbreak\detokenize{f_0,}\allowbreak\detokenize{f_1,}\allowbreak\detokenize{L,}\allowbreak\detokenize{c,}\allowbreak\detokenize{\gamma}} & dominating measure null and alternative densities extended likelihood ratio threshold and tie randomization & P15 hypothesis-tests-confidence & meta & \hyperref[sec:p15-hypothesis-tests-confidence]{\ref*{sec:p15-hypothesis-tests-confidence}} & \thepage & The ratio is infinite where only the null density is zero; threshold atoms and significance endpoints are explicit. \\
\phantomsection\label{entry:symbol-registry-01576}\(K,n,p,p_0,b_{n,p}(k),\pi,\pi_r\) & \texttt{\detokenize{K,}\allowbreak\detokenize{n,}\allowbreak\detokenize{p,}\allowbreak\detokenize{p_0,}\allowbreak\detokenize{b_{n,p}(k),}\allowbreak\detokenize{\pi,}\allowbreak\detokenize{\pi_r}} & binomial count sample size success probabilities mass function ordinary and randomized p-values & P15 hypothesis-tests-confidence & meta & \hyperref[sec:p15-hypothesis-tests-confidence]{\ref*{sec:p15-hypothesis-tests-confidence}} & \thepage & Discrete inclusive tails are conservative; independent randomization gives exact uniformity. \\
\phantomsection\label{entry:symbol-registry-01577}\(\tau,\Lambda,t,A_t,C(x,u),S(x),\ell,u,q\) & \texttt{\detokenize{\tau,}\allowbreak\detokenize{\Lambda,}\allowbreak\detokenize{t,}\allowbreak\detokenize{A_t,}\allowbreak\detokenize{C(x,}\allowbreak\detokenize{u),}\allowbreak\detokenize{S(x),}\allowbreak\detokenize{\ell,}\allowbreak\detokenize{u,}\allowbreak\detokenize{q}} & target parameter map target space target value acceptance regions confidence sets sections endpoints and confidence coefficient & P15 hypothesis-tests-confidence & meta & \hyperref[sec:p15-hypothesis-tests-confidence]{\ref*{sec:p15-hypothesis-tests-confidence}} & \thepage & Target-map inversion retains nuisance parameters and requires joint measurability. \\
\phantomsection\label{entry:symbol-registry-01578}\(d,T_{n,p},F_{n,p},\ell_n,u_n,Q_{a,b},I_p(a,b)\) & \texttt{\detokenize{d,}\allowbreak\detokenize{T_{n,p},}\allowbreak\detokenize{F_{n,p},}\allowbreak\detokenize{\ell_n,}\allowbreak\detokenize{u_n,}\allowbreak\detokenize{Q_{a,b},}\allowbreak\detokenize{I_p(a,}\allowbreak\detokenize{b)}} & half-tail size binomial tails confidence limits beta quantiles and incomplete beta integral & P15 hypothesis-tests-confidence & meta & \hyperref[sec:p15-hypothesis-tests-confidence]{\ref*{sec:p15-hypothesis-tests-confidence}} & \thepage & Here d=alpha/2; zero beta shapes are excluded and endpoint limits are given separately. \\
\phantomsection\label{entry:symbol-registry-01579}\(d,\ell,q,\eta,s,k,h\) & \texttt{\detokenize{d,}\allowbreak\detokenize{\ell,}\allowbreak\detokenize{q,}\allowbreak\detokenize{\eta,}\allowbreak\detokenize{s,}\allowbreak\detokenize{k,}\allowbreak\detokenize{h}} & local variation dimension boundary functions perturbation variation parameter multiplier and signed density integrand & P15 hypothesis-tests-confidence & meta & \hyperref[sec:p15-hypothesis-tests-confidence]{\ref*{sec:p15-hypothesis-tests-confidence}} & \thepage & These roles are confined to the smooth boundary-variation proof and distinct from the confidence and binomial roles. \\
\phantomsection\label{entry:symbol-registry-01587}\(X,\mathcal A,Y,H,f,h,P,S,V(S)\) & \texttt{\detokenize{X,}\allowbreak\detokenize{\mathcal A,}\allowbreak\detokenize{Y,}\allowbreak\detokenize{H,}\allowbreak\detokenize{f,}\allowbreak\detokenize{h,}\allowbreak\detokenize{P,}\allowbreak\detokenize{S,}\allowbreak\detokenize{V(S)}} & input space sigma field labels hypothesis class target candidate input law labelled sample and version space & P15 fisher-estimation & meta & \hyperref[sec:p15-fisher-estimation]{\ref*{sec:p15-fisher-estimation}} & \thepage & The class is fixed before sampling; labels are noiseless and samples independent with the same input distribution. \\
\phantomsection\label{entry:symbol-registry-01588}\(R_P(h,f),\widehat R_S(h),m,r,\varepsilon,\delta\) & \texttt{\detokenize{R_P(h,}\allowbreak\detokenize{f),}\allowbreak\detokenize{\widehat R_S(h),}\allowbreak\detokenize{m,}\allowbreak\detokenize{r,}\allowbreak\detokenize{\varepsilon,}\allowbreak\detokenize{\delta}} & generalization and empirical errors sample size finite class size error tolerance and sample failure tolerance & P15 fisher-estimation & meta & \hyperref[sec:p15-fisher-estimation]{\ref*{sec:p15-fisher-estimation}} & \thepage & The probability of a bad selected hypothesis is distinct from that hypothesis's error on a fresh input. \\
\phantomsection\label{entry:symbol-registry-01589}\(\mathcal C,\ell(h),L,n,c,a,q\) & \texttt{\detokenize{\mathcal C,}\allowbreak\detokenize{\ell(h),}\allowbreak\detokenize{L,}\allowbreak\detokenize{n,}\allowbreak\detokenize{c,}\allowbreak\detokenize{a,}\allowbreak\detokenize{q}} & fixed binary codes minimum length length limit target size bound compression exponents and minus log one minus epsilon & P15 fisher-estimation & meta & \hyperref[sec:p15-fisher-estimation]{\ref*{sec:p15-fisher-estimation}} & \thepage & Arbitrary variable-length and prefix-free code counts are separate; c is also a local individual code in the Kraft proof. \\
\phantomsection\label{entry:symbol-registry-01590}\(b(S),b_+(S),w_h,H_L,t_h\) & \texttt{\detokenize{b(S),}\allowbreak\detokenize{b_+(S),}\allowbreak\detokenize{w_h,}\allowbreak\detokenize{H_L,}\allowbreak\detokenize{t_h}} & shortest consistent length its positive modification predetermined hypothesis weights bounded-code class and hypothesis error threshold & P15 fisher-estimation & meta & \hyperref[sec:p15-fisher-estimation]{\ref*{sec:p15-fisher-estimation}} & \thepage & The weighted union allocation is not a posterior over hypotheses; weights and decoding are fixed before the sample. \\
\phantomsection\label{entry:symbol-registry-01591}\(h_t,O(S),h_S,N,F\) & \texttt{\detokenize{h_t,}\allowbreak\detokenize{O(S),}\allowbreak\detokenize{h_S,}\allowbreak\detokenize{N,}\allowbreak\detokenize{F}} & worked threshold hypotheses observed-input set memorizing hypothesis finite domain size and random target used in a lower-bound proof & P15 fisher-estimation & meta & \hyperref[sec:p15-fisher-estimation]{\ref*{sec:p15-fisher-estimation}} & \thepage & Averaging over target labels proves existence of a fixed difficult target; it does not change the upper-bound model. \\
\phantomsection\label{entry:symbol-registry-01592}\(n,M,d\) & \texttt{\detokenize{n,}\allowbreak\detokenize{M,}\allowbreak\detokenize{d}} & information index simulation count and fixed parameter dimension & P15 routed chapter bayes-factor-computation & meta & \hyperref[def:p15-laplace-family]{\ref*{def:p15-laplace-family}} & \thepage & Local context: sec:p15-bayes-factor-computation. The asymptotic sample index and Monte Carlo count are distinct. \\
\phantomsection\label{entry:symbol-registry-01593}\(\Theta_n,\theta_n,h_n,a_n,I_n,\widehat I_n\) & \texttt{\detokenize{\Theta_n,}\allowbreak\detokenize{\theta_n,}\allowbreak\detokenize{h_n,}\allowbreak\detokenize{a_n,}\allowbreak\detokenize{I_n,}\allowbreak\detokenize{\widehat I_n}} & Laplace domain mode exponent amplitude integral and approximation & P15 routed chapter bayes-factor-computation & meta & \hyperref[def:p15-laplace-family]{\ref*{def:p15-laplace-family}} & \thepage & Local context: sec:p15-laplace-hypotheses. Uniform local and tail conditions are explicitly enumerated. \\
\phantomsection\label{entry:symbol-registry-01594}\(J_n,c,C,A,a_-,a_+,\rho,c_0,A_T,\eta\) & \texttt{\detokenize{J_n,}\allowbreak\detokenize{c,}\allowbreak\detokenize{C,}\allowbreak\detokenize{A,}\allowbreak\detokenize{a_-,}\allowbreak\detokenize{a_+,}\allowbreak\detokenize{\rho,}\allowbreak\detokenize{c_0,}\allowbreak\detokenize{A_T,}\allowbreak\detokenize{\eta}} & scaled curvature and fixed Laplace bounds & P15 routed chapter bayes-factor-computation & meta & \hyperref[def:p15-laplace-family]{\ref*{def:p15-laplace-family}} & \thepage & Local context: sec:p15-laplace-hypotheses. All constants independent of n; derivative norms are multilinear operator norms. \\
\phantomsection\label{entry:symbol-registry-01595}\(q_n,\ell_n,\widetilde\theta_n,\widehat\theta_n,H_n,F_n,\pi\) & \texttt{\detokenize{q_n,}\allowbreak\detokenize{\ell_n,}\allowbreak\detokenize{\widetilde\theta_n,}\allowbreak\detokenize{\widehat\theta_n,}\allowbreak\detokenize{H_n,}\allowbreak\detokenize{F_n,}\allowbreak\detokenize{\pi}} & posterior and likelihood logs modes Hessians and proper prior density & P15 routed chapter bayes-factor-computation & meta & \hyperref[thm:p15-laplace-mode-variants]{\ref*{thm:p15-laplace-mode-variants}} & \thepage & Local context: sec:p15-laplace-mode-variants. Each variant must meet its own exponent/amplitude hypotheses. \\
\phantomsection\label{entry:symbol-registry-01596}\(\widetilde J_n,\delta_n\) & \texttt{\detokenize{\widetilde J_n,}\allowbreak\detokenize{\delta_n}} & symmetric approximate scaled information and operator-norm error & P15 routed chapter bayes-factor-computation & meta & \hyperref[thm:p15-laplace-curvature-replacement]{\ref*{thm:p15-laplace-curvature-replacement}} & \thepage & Local context: sec:p15-laplace-mode-variants. Matrix replacement rate is a premise, not assumed from its name. \\
\phantomsection\label{entry:symbol-registry-01597}\(c_{k,n},\mathrm{BIC}_k\) & \texttt{\detokenize{c_{k,n},}\allowbreak\detokenize{\mathrm{BIC}_k}} & bounded log-evidence correction and fixed-dimensional BIC score & P15 routed chapter bayes-factor-computation & meta & \hyperref[thm:p15-bic-leading-order]{\ref*{thm:p15-bic-leading-order}} & \thepage & Local context: sec:p15-bic-consequence. Discarded constant-order terms need not yield small relative factor error. \\
\phantomsection\label{entry:symbol-registry-01598}\(L,\Pi,I,T,P\) & \texttt{\detokenize{L,}\allowbreak\detokenize{\Pi,}\allowbreak\detokenize{I,}\allowbreak\detokenize{T,}\allowbreak\detokenize{P}} & likelihood proper prior evidence unnormalized posterior measure and posterior probability & P15 routed chapter bayes-factor-computation & meta & \hyperref[thm:p15-direct-evidence-mc]{\ref*{thm:p15-direct-evidence-mc}} & \thepage & Local context: sec:p15-evidence-monte-carlo. T=L Pi and P=T/I; I is positive finite. \\
\phantomsection\label{entry:symbol-registry-01599}\(Q,V,W,\widehat I_M,\widetilde I_M\) & \texttt{\detokenize{Q,}\allowbreak\detokenize{V,}\allowbreak\detokenize{W,}\allowbreak\detokenize{\widehat I_M,}\allowbreak\detokenize{\widetilde I_M}} & proposal target derivative prior derivative ordinary and normalized importance estimates & P15 routed chapter bayes-factor-computation & meta & \hyperref[thm:p15-importance-evidence]{\ref*{thm:p15-importance-evidence}} & \thepage & Local context: sec:p15-evidence-monte-carlo. Ordinary requires T<<Q; normalized requires Pi<<Q. \\
\phantomsection\label{entry:symbol-registry-01600}\(G,H,A_M,s\) & \texttt{\detokenize{G,}\allowbreak\detokenize{H,}\allowbreak\detokenize{A_M,}\allowbreak\detokenize{s}} & auxiliary probability generalized harmonic weight reciprocal mean and positive-likelihood prior mass & P15 routed chapter bayes-factor-computation & meta & \hyperref[thm:p15-harmonic-evidence]{\ref*{thm:p15-harmonic-evidence}} & \thepage & Local context: sec:p15-harmonic-evidence. G<<T; H=dG/dT; A\_M estimates 1/I. \\
\phantomsection\label{entry:symbol-registry-01601}\(v,v_{\rm SN},v_H\) & \texttt{\detokenize{v,}\allowbreak\detokenize{v_{\rm SN},}\allowbreak\detokenize{v_H}} & variance parameters in independent-simulation normal limits & P15 routed chapter bayes-factor-computation & meta & \hyperref[thm:p15-evidence-mc-clt]{\ref*{thm:p15-evidence-mc-clt}} & \thepage & Local context: sec:p15-harmonic-evidence. All are finite nonnegative under the stated theorem; no MCMC conclusion. \\
\phantomsection\label{entry:symbol-registry-01602}\(\alpha,\theta,l,L,L_n,S,\bar X,\beta,q,q_n\) & \texttt{\detokenize{\alpha,}\allowbreak\detokenize{\theta,}\allowbreak\detokenize{l,}\allowbreak\detokenize{L,}\allowbreak\detokenize{L_n,}\allowbreak\detokenize{S,}\allowbreak\detokenize{\bar X,}\allowbreak\detokenize{\beta,}\allowbreak\detokenize{q,}\allowbreak\detokenize{q_n}} & source confidence coefficient uniform scale minimum maxima sum mean and interval cutoffs & P15 occam-learning & meta & \hyperref[sec:p15-occam-learning]{\ref*{sec:p15-occam-learning}} & \thepage & Alpha is coverage here; the testing size in the preceding chapter has a different local convention. \\
\phantomsection\label{entry:symbol-registry-01603}\(A_j^{\rm U},I_j^{\rm U},D_2,D_3\) & \texttt{\detokenize{A_j^{\rm U},}\allowbreak\detokenize{I_j^{\rm U},}\allowbreak\detokenize{D_2,}\allowbreak\detokenize{D_3}} & three uniform acceptance systems their inverses and sum/max interval lengths & P15 occam-learning & meta & \hyperref[sec:p15-occam-learning]{\ref*{sec:p15-occam-learning}} & \thepage & The sum-strip calibration is corrected and boundary sets are included for exact closed inversion. \\
\phantomsection\label{entry:symbol-registry-01604}\(T,V,m_{\theta,\eta},h_\theta,\theta_1,\theta_2,\rho\) & \texttt{\detokenize{T,}\allowbreak\detokenize{V,}\allowbreak\detokenize{m_{\theta,\eta},}\allowbreak\detokenize{h_\theta,}\allowbreak\detokenize{\theta_1,}\allowbreak\detokenize{\theta_2,}\allowbreak\detokenize{\rho}} & conditional coordinates marginal and conditional densities source scale nuisance parameter and dimensionless nuisance value & P15 occam-learning & meta & \hyperref[sec:p15-occam-learning]{\ref*{sec:p15-occam-learning}} & \thepage & Coordinates and density factorization are explicit; the source triangle retains -1<rho<=2. \\
\phantomsection\label{entry:symbol-registry-01605}\(t,s,\eta,b,A_b^{\rm T},I_b^{\rm T},I_*^{\rm T},A_{\rm sw}^{\rm T}\) & \texttt{\detokenize{t,}\allowbreak\detokenize{s,}\allowbreak\detokenize{\eta,}\allowbreak\detokenize{b,}\allowbreak\detokenize{A_b^{\rm T},}\allowbreak\detokenize{I_b^{\rm T},}\allowbreak\detokenize{I_*^{\rm T},}\allowbreak\detokenize{A_{\rm sw}^{\rm T}}} & candidate and true scales nuisance placement parameter fixed-placement optimum and switched triangle systems & P15 occam-learning & meta & \hyperref[sec:p15-occam-learning]{\ref*{sec:p15-occam-learning}} & \thepage & The switched system includes a null-boundary completion and every positive-observation case. \\
\phantomsection\label{entry:symbol-registry-01606}\(Q_I,\operatorname{len},g,w,\Phi,z_\alpha,\nu\) & \texttt{\detokenize{Q_I,}\allowbreak\detokenize{\operatorname{len},}\allowbreak\detokenize{g,}\allowbreak\detokenize{w,}\allowbreak\detokenize{\Phi,}\allowbreak\detokenize{z_\alpha,}\allowbreak\detokenize{\nu}} & false-coverage function length L1 derivative score normal CDF quantile and dominating measure & P15 occam-learning & meta & \hyperref[sec:p15-occam-learning]{\ref*{sec:p15-occam-learning}} & \thepage & Global shortest and unbiased criteria are distinguished from local derivative and curvature requirements. \\
\phantomsection\label{entry:symbol-registry-01610}\(X_i,a_i,b_i,\mu_i,S_k,V_n,\lambda\) & \texttt{\detokenize{X_i,}\allowbreak\detokenize{a_i,}\allowbreak\detokenize{b_i,}\allowbreak\detokenize{\mu_i,}\allowbreak\detokenize{S_k,}\allowbreak\detokenize{V_n,}\allowbreak\detokenize{\lambda}} & bounded variables deterministic ranges means centered sums squared-range total and exponential parameter & P15 vc-dimension-and-uniform-learning & meta & \hyperref[sec:p15-vc-dimension-and-uniform-learning]{\ref*{sec:p15-vc-dimension-and-uniform-learning}} & \thepage & The independent theorem requires mutual independence; zero-width terms and signed weights are explicit. \\
\phantomsection\label{entry:symbol-registry-01611}\(D(q\Vert p),g(p),G(t,p),H(z),K_p(s)\) & \texttt{\detokenize{D(q\Vert p),}\allowbreak\detokenize{g(p),}\allowbreak\detokenize{G(t,}\allowbreak\detokenize{p),}\allowbreak\detokenize{H(z),}\allowbreak\detokenize{K_p(s)}} & binary divergence exact mean-dependent quadratic coefficient and proof functions for the exponential refinements & P15 vc-dimension-and-uniform-learning & meta & \hyperref[sec:p15-vc-dimension-and-uniform-learning]{\ref*{sec:p15-vc-dimension-and-uniform-learning}} & \thepage & Divergence has continuous endpoint values; the p-dependent coefficient is proved on its full stated interval. \\
\phantomsection\label{entry:symbol-registry-01612}\(\mathcal F_i,D_i,M_k\) & \texttt{\detokenize{\mathcal F_i,}\allowbreak\detokenize{D_i,}\allowbreak\detokenize{M_k}} & filtration centered adapted differences and martingale partial sums & P15 vc-dimension-and-uniform-learning & meta & \hyperref[sec:p15-vc-dimension-and-uniform-learning]{\ref*{sec:p15-vc-dimension-and-uniform-learning}} & \thepage & Conditional means are zero and range endpoints deterministic; maximal and time-uniform guarantees have their stated horizons. \\
\phantomsection\label{entry:symbol-registry-01613}\(\mathcal Z,\mathcal B,P,\mathcal H,\ell,R,\widehat R_n,r,b_n,\widehat h,\eta\) & \texttt{\detokenize{\mathcal Z,}\allowbreak\detokenize{\mathcal B,}\allowbreak\detokenize{P,}\allowbreak\detokenize{\mathcal H,}\allowbreak\detokenize{\ell,}\allowbreak\detokenize{R,}\allowbreak\detokenize{\widehat R_n,}\allowbreak\detokenize{r,}\allowbreak\detokenize{b_n,}\allowbreak\detokenize{\widehat h,}\allowbreak\detokenize{\eta}} & record space law hypothesis class bounded loss true and empirical risks class size uniform radius selected hypothesis and optimization tolerance & P15 vc-dimension-and-uniform-learning & meta & \hyperref[sec:p15-vc-dimension-and-uniform-learning]{\ref*{sec:p15-vc-dimension-and-uniform-learning}} & \thepage & The class is fixed before sampling and may have no perfect predictor; selection is covered by a simultaneous event. \\
\phantomsection\label{entry:symbol-registry-01614}\(w_h,b_h,\delta,T,K,m,R_j,\widehat R_j\) & \texttt{\detokenize{w_h,}\allowbreak\detokenize{b_h,}\allowbreak\detokenize{\delta,}\allowbreak\detokenize{T,}\allowbreak\detokenize{K,}\allowbreak\detokenize{m,}\allowbreak\detokenize{R_j,}\allowbreak\detokenize{\widehat R_j}} & fixed hypothesis weights penalties failure probability training object validation candidate count validation size and risks & P15 vc-dimension-and-uniform-learning & meta & \hyperref[sec:p15-vc-dimension-and-uniform-learning]{\ref*{sec:p15-vc-dimension-and-uniform-learning}} & \thepage & Validation observations are independent of the full training object; countable penalized selection with eta=0 needs attainment. \\
\phantomsection\label{entry:symbol-registry-01615}\(B_{ij},\mathbb F_2,Z,a,X_a,t_n,\tau\) & \texttt{\detokenize{B_{ij},}\allowbreak\detokenize{\mathbb F_2,}\allowbreak\detokenize{Z,}\allowbreak\detokenize{a,}\allowbreak\detokenize{X_a,}\allowbreak\detokenize{t_n,}\allowbreak\detokenize{\tau}} & feature matrix two-element field fair-vector parity labels binary variables simultaneous deviation radius and stopping time & P15 vc-dimension-and-uniform-learning & meta & \hyperref[sec:p15-vc-dimension-and-uniform-learning]{\ref*{sec:p15-vc-dimension-and-uniform-learning}} & \thepage & The fifteen-variable construction is pairwise independent but not mutually independent. \\
\phantomsection\label{entry:symbol-registry-01616}\(A_1,A_2,A_3,J(t),t',c\) & \texttt{\detokenize{A_1,}\allowbreak\detokenize{A_2,}\allowbreak\detokenize{A_3,}\allowbreak\detokenize{J(t),}\allowbreak\detokenize{t',}\allowbreak\detokenize{c}} & three common-range exponential bounds divergence numerator reflected argument and reflection center & P15 vc-dimension-and-uniform-learning & meta & \hyperref[sec:p15-vc-dimension-and-uniform-learning]{\ref*{sec:p15-vc-dimension-and-uniform-learning}} & \thepage & J extends to -p<t<1-p so the reflection argument remains defined when t' is negative. \\
\phantomsection\label{entry:symbol-registry-01617}\(W,v,u,\phi,\Phi,\xi_{n,i},v_n\) & \texttt{\detokenize{W,}\allowbreak\detokenize{v,}\allowbreak\detokenize{u,}\allowbreak\detokenize{\phi,}\allowbreak\detokenize{\Phi,}\allowbreak\detokenize{\xi_{n,i},}\allowbreak\detokenize{v_n}} & centered moment benchmark variance positive threshold normal density and CDF centered row variables and variance sums & P15 vc-dimension-and-uniform-learning & meta & \hyperref[sec:p15-vc-dimension-and-uniform-learning]{\ref*{sec:p15-vc-dimension-and-uniform-learning}} & \thepage & The source Lyapunov comparison is asymptotic; the exponential and moment inequalities hold at finite sample sizes. \\
\phantomsection\label{entry:symbol-registry-01618}\(X_i,\overline X,b,v_i,\sigma^2\) & \texttt{\detokenize{X_i,}\allowbreak\detokenize{\overline X,}\allowbreak\detokenize{b,}\allowbreak\detokenize{v_i,}\allowbreak\detokenize{\sigma^2}} & centered variables sample mean upper bound individual and averaged variances & P15 routed chapter variance-and-dependent-concentration & meta & \hyperref[def:p15-upper-variance-setting]{\ref*{def:p15-upper-variance-setting}} & \thepage & Local context: sec:p15-upper-variance. Mutual independence is required; no lower bound is imposed. \\
\phantomsection\label{entry:symbol-registry-01619}\(M_{b,v}(h),h\) & \texttt{\detokenize{M_{b,v}(h),}\allowbreak\detokenize{h}} & two-point exponential moment envelope and positive exponential parameter & P15 routed chapter variance-and-dependent-concentration & meta & \hyperref[thm:p15-upper-moment-extremal]{\ref*{thm:p15-upper-moment-extremal}} & \thepage & Local context: sec:p15-upper-variance. v is a nonnegative moment parameter; b is positive. \\
\phantomsection\label{entry:symbol-registry-01620}\(f_c(u),c,u\) & \texttt{\detokenize{f_c(u),}\allowbreak\detokenize{c,}\allowbreak\detokenize{u}} & log moment as a function of normalized variance & P15 routed chapter variance-and-dependent-concentration & meta & \hyperref[thm:p15-variance-moment-concavity]{\ref*{thm:p15-variance-moment-concavity}} & \thepage & Local context: sec:p15-upper-variance. c=bh and u=v/b\^{}2 in the moment application. \\
\phantomsection\label{entry:symbol-registry-01621}\(B_1,B_2,\lambda,\tau,h_1,h_2,h_3\) & \texttt{\detokenize{B_1,}\allowbreak\detokenize{B_2,}\allowbreak\detokenize{\lambda,}\allowbreak\detokenize{\tau,}\allowbreak\detokenize{h_1,}\allowbreak\detokenize{h_2,}\allowbreak\detokenize{h_3}} & exact exponential bound Bennett bound normalized parameters and comparison exponents & P15 routed chapter variance-and-dependent-concentration & meta & \hyperref[thm:p15-variance-tail-chain]{\ref*{thm:p15-variance-tail-chain}} & \thepage & Local context: sec:p15-upper-variance. lambda=bt/sigma squared and tau=nt/b; h\_j are functions not simulation counts. \\
\phantomsection\label{entry:symbol-registry-01622}\(v,w,\rho,\phi\) & \texttt{\detokenize{v,}\allowbreak\detokenize{w,}\allowbreak\detokenize{\rho,}\allowbreak\detokenize{\phi}} & normalized ratio coordinates series and exact log-ratio correction & P15 routed chapter variance-and-dependent-concentration & meta & \hyperref[thm:p15-variance-bound-comparison]{\ref*{thm:p15-variance-bound-comparison}} & \thepage & Local context: thm:p15-variance-bound-comparison. Here v=bt/(bt+sigma squared), distinct from the earlier one-variable moment parameter. \\
\phantomsection\label{entry:symbol-registry-01623}\(T,T_j,p_j\) & \texttt{\detokenize{T,}\allowbreak\detokenize{T_j,}\allowbreak\detokenize{p_j}} & pointwise convex average components and deterministic weights & P15 routed chapter variance-and-dependent-concentration & meta & \hyperref[thm:p15-concentration-convex-mixture]{\ref*{thm:p15-concentration-convex-mixture}} & \thepage & Local context: sec:p15-dependent-blocks. Components may be dependent; zero weights are omitted. \\
\phantomsection\label{entry:symbol-registry-01624}\((n)_r,U,g,k,V_\pi\) & \texttt{\detokenize{(n)_r,}\allowbreak\detokenize{U,}\allowbreak\detokenize{g,}\allowbreak\detokenize{k,}\allowbreak\detokenize{V_\pi}} & falling factorial ordered U statistic kernel effective block count and permutation block average & P15 routed chapter variance-and-dependent-concentration & meta & \hyperref[def:p15-ordered-u-statistic]{\ref*{def:p15-ordered-u-statistic}} & \thepage & Local context: sec:p15-dependent-blocks. g need not be symmetric; n>=r>=1 and k=floor(n/r). \\
\phantomsection\label{entry:symbol-registry-01625}\(m,n,r,s,V_{\pi,\eta}\) & \texttt{\detokenize{m,}\allowbreak\detokenize{n,}\allowbreak\detokenize{r,}\allowbreak\detokenize{s,}\allowbreak\detokenize{V_{\pi,\eta}}} & sample sizes kernel orders and two-sample permutation average & P15 routed chapter variance-and-dependent-concentration & meta & \hyperref[thm:p15-two-sample-u-concentration]{\ref*{thm:p15-two-sample-u-concentration}} & \thepage & Local context: thm:p15-two-sample-u-concentration. k=min(floor(m/r),floor(n/s)); input independence is joint across both samples. \\
\phantomsection\label{entry:symbol-registry-01626}\(W,g^*,\mathcal P,\beta_{\mathcal P}\) & \texttt{\detokenize{W,}\allowbreak\detokenize{g^*,}\allowbreak\detokenize{\mathcal P,}\allowbreak\detokenize{\beta_{\mathcal P}}} & repeated-index average equivalent U kernel equality partition and block label & P15 routed chapter variance-and-dependent-concentration & meta & \hyperref[thm:p15-repeated-index-concentration]{\ref*{thm:p15-repeated-index-concentration}} & \thepage & Local context: sec:p15-dependent-blocks. Partition weights are (n)\_{number of blocks}/n\^{}r. \\
\phantomsection\label{entry:symbol-registry-01627}\(G,F_n,\omega^2\) & \texttt{\detokenize{G,}\allowbreak\detokenize{F_n,}\allowbreak\detokenize{\omega^2}} & continuous reference cdf empirical cdf and Cramer-von Mises integral & P15 routed chapter variance-and-dependent-concentration & meta & \hyperref[thm:p15-cramer-von-mises-kernel]{\ref*{thm:p15-cramer-von-mises-kernel}} & \thepage & Local context: thm:p15-cramer-von-mises-kernel. The concentration statement does not impose the null distribution G on inputs. \\
\phantomsection\label{entry:symbol-registry-01628}\(Y_i,r,I_i,n_i,S_i,S\) & \texttt{\detokenize{Y_i,}\allowbreak\detokenize{r,}\allowbreak\detokenize{I_i,}\allowbreak\detokenize{n_i,}\allowbreak\detokenize{S_i,}\allowbreak\detokenize{S}} & finite-range dependent sequence separation size residues residue counts and sums & P15 routed chapter variance-and-dependent-concentration & meta & \hyperref[def:p15-finite-range-dependence]{\ref*{def:p15-finite-range-dependence}} & \thepage & Local context: sec:p15-dependent-blocks. The definition is independence of separated vectors, not pairwise independence. \\
\phantomsection\label{entry:symbol-registry-01629}\(c_j,N,\mu,\sigma_C^2,X_i,Y_i\) & \texttt{\detokenize{c_j,}\allowbreak\detokenize{N,}\allowbreak\detokenize{\mu,}\allowbreak\detokenize{\sigma_C^2,}\allowbreak\detokenize{X_i,}\allowbreak\detokenize{Y_i}} & population values size mean variance and without/with replacement samples & P15 routed chapter variance-and-dependent-concentration & meta & \hyperref[def:p15-finite-population-sampling]{\ref*{def:p15-finite-population-sampling}} & \thepage & Local context: sec:p15-finite-population-concentration. Sampling distinct indices allows repeated population values; samples are unweighted. \\
\phantomsection\label{entry:symbol-registry-01630}\(\overline g,p_{\mathcal P},r_\ell\) & \texttt{\detokenize{\overline g,}\allowbreak\detokenize{p_{\mathcal P},}\allowbreak\detokenize{r_\ell}} & symmetrized comparison kernel partition weights and block multiplicities & P15 routed chapter variance-and-dependent-concentration & meta & \hyperref[thm:p15-without-replacement-convex]{\ref*{thm:p15-without-replacement-convex}} & \thepage & Local context: annex:p15.without-replacement-convex. The coefficients and their linear-mean identity are explicitly proved. \\
\phantomsection\label{entry:symbol-registry-01631}\(S_j,A_j\) & \texttt{\detokenize{S_j,}\allowbreak\detokenize{A_j}} & independent centered partial sum and first-crossing event & P15 routed chapter variance-and-dependent-concentration & meta & \hyperref[thm:p15-variance-maximal]{\ref*{thm:p15-variance-maximal}} & \thepage & Local context: thm:p15-variance-maximal. The variance-sensitive maximum uses the terminal exponential moment with independent increments. \\
\phantomsection\label{entry:symbol-registry-01632}\(Q,c_j,\alpha_j,J,W,a\) & \texttt{\detokenize{Q,}\allowbreak\detokenize{c_j,}\allowbreak\detokenize{\alpha_j,}\allowbreak\detokenize{J,}\allowbreak\detokenize{W,}\allowbreak\detokenize{a}} & finite exponential majorant coefficients parameters term count variable and threshold & P15 routed chapter variance-and-dependent-concentration & meta & \hyperref[thm:p15-variance-optimality-boundary]{\ref*{thm:p15-variance-optimality-boundary}} & \thepage & Local context: thm:p15-variance-optimality-boundary. Every exponential moment must be finite; signed coefficients retain their inequality direction. \\
\phantomsection\label{entry:symbol-registry-01633}\(n,\nu,P_{\mu,\sigma},E_{\mu,\sigma},\bar X,V,S,e,T_m\) & \texttt{\detokenize{n,}\allowbreak\detokenize{\nu,}\allowbreak\detokenize{P_{\mu,\sigma},}\allowbreak\detokenize{E_{\mu,\sigma},}\allowbreak\detokenize{\bar X,}\allowbreak\detokenize{V,}\allowbreak\detokenize{S,}\allowbreak\detokenize{e,}\allowbreak\detokenize{T_m}} & sample size residual degrees of freedom normal experiment mean residual sum sample standard deviation standard error and Student statistic & Legacy pre-batch7 P15.46 reader meaning & meta & \hyperref[sec:p15-neyman-confidence-constructions]{\ref*{sec:p15-neyman-confidence-constructions}} & \thepage & n>=2 and sigma>0; source1937 s equals the standard error e rather than S.  Numeric chapter scope migrated to its frozen pre-batch7 source-labelled reader regions; glyph, role and precedence unchanged. \\
\phantomsection\label{entry:symbol-registry-01634}\(\chi_k^2,O,b,d,v,\omega_n,h_n\) & \texttt{\detokenize{\chi_k^2,}\allowbreak\detokenize{O,}\allowbreak\detokenize{b,}\allowbreak\detokenize{d,}\allowbreak\detokenize{v,}\allowbreak\detokenize{\omega_n,}\allowbreak\detokenize{h_n}} & chi-square law Householder matrix coordinate vectors sphere probability and coordinate density & Legacy pre-batch7 P15.46 reader meaning & meta & \hyperref[sec:p15-neyman-confidence-constructions]{\ref*{sec:p15-neyman-confidence-constructions}} & \thepage & The gamma convention uses rate; orthogonal radius-direction and ratio independence are proved before the pivot law.  Numeric chapter scope migrated to its frozen pre-batch7 source-labelled reader regions; glyph, role and precedence unchanged. \\
\phantomsection\label{entry:symbol-registry-01635}\(f_\nu,F_\nu,t_{\nu,p},C_-,C_+,C_0,\alpha,\varepsilon\) & \texttt{\detokenize{f_\nu,}\allowbreak\detokenize{F_\nu,}\allowbreak\detokenize{t_{\nu,p},}\allowbreak\detokenize{C_-,}\allowbreak\detokenize{C_+,}\allowbreak\detokenize{C_0,}\allowbreak\detokenize{\alpha,}\allowbreak\detokenize{\varepsilon}} & Student density CDF quantile confidence sets coverage coefficient and rejection probability & Legacy pre-batch7 P15.46 reader meaning & meta & \hyperref[sec:p15-neyman-confidence-constructions]{\ref*{sec:p15-neyman-confidence-constructions}} & \thepage & Alpha is confidence and epsilon=1-alpha; unscaled inequalities define zero-residual outcomes.  Numeric chapter scope migrated to its frozen pre-batch7 source-labelled reader regions; glyph, role and precedence unchanged. \\
\phantomsection\label{entry:symbol-registry-01636}\(\varphi_m^+,\varphi_m^-,\varphi_m^0,\beta_\varphi,Q_m,U_m,Z_m,H_m\) & \texttt{\detokenize{\varphi_m^+,}\allowbreak\detokenize{\varphi_m^-,}\allowbreak\detokenize{\varphi_m^0,}\allowbreak\detokenize{\beta_\varphi,}\allowbreak\detokenize{Q_m,}\allowbreak\detokenize{U_m,}\allowbreak\detokenize{Z_m,}\allowbreak\detokenize{H_m}} & one- and two-sided tests power null radius and direction quantities and score sum & Legacy pre-batch7 P15.46 reader meaning & meta & \hyperref[sec:p15-neyman-confidence-constructions]{\ref*{sec:p15-neyman-confidence-constructions}} & \thepage & One-sided optimum is among similar tests; two-sided global unbiased optimum is proved separately.  Numeric chapter scope migrated to its frozen pre-batch7 source-labelled reader regions; glyph, role and precedence unchanged. \\
\phantomsection\label{entry:symbol-registry-01637}\(\delta,a(q),M(a),A_\varphi,B_\varphi,c_+,c,r,\ell_a\) & \texttt{\detokenize{\delta,}\allowbreak\detokenize{a(q),}\allowbreak\detokenize{M(a),}\allowbreak\detokenize{A_\varphi,}\allowbreak\detokenize{B_\varphi,}\allowbreak\detokenize{c_+,}\allowbreak\detokenize{c,}\allowbreak\detokenize{r,}\allowbreak\detokenize{\ell_a}} & mean shift spherical tilt normalizer conditional angular mass and score cutoffs selector and exponential chord & Legacy pre-batch7 P15.46 reader meaning & meta & \hyperref[sec:p15-neyman-confidence-constructions]{\ref*{sec:p15-neyman-confidence-constructions}} & \thepage & Conditional sphere constraints hold almost everywhere with the specified common null-set qualification.  Numeric chapter scope migrated to its frozen pre-batch7 source-labelled reader regions; glyph, role and precedence unchanged. \\
\phantomsection\label{entry:symbol-registry-01638}\(a(x,m),v,\varphi_i,\varphi_{ij},s_{33},z_{33}\) & \texttt{\detokenize{a(x,}\allowbreak\detokenize{m),}\allowbreak\detokenize{v,}\allowbreak\detokenize{\varphi_i,}\allowbreak\detokenize{\varphi_{ij},}\allowbreak\detokenize{s_{33},}\allowbreak\detokenize{z_{33}}} & randomized interval inclusion source variance coordinate derivative scores and original1933 scale/statistic & Legacy pre-batch7 P15.46 reader meaning & meta & \hyperref[sec:p15-neyman-confidence-constructions]{\ref*{sec:p15-neyman-confidence-constructions}} & \thepage & These source-local coordinates are distinct from the posterior and modern sample-standard-deviation conventions.  Numeric chapter scope migrated to its frozen pre-batch7 source-labelled reader regions; glyph, role and precedence unchanged. \\
\phantomsection\label{entry:symbol-registry-01642}\(\mathcal C_i,\varphi,A(\varphi),B(\varphi)\) & \texttt{\detokenize{\mathcal C_i,}\allowbreak\detokenize{\varphi,}\allowbreak\detokenize{A(\varphi),}\allowbreak\detokenize{B(\varphi)}} & neighborhoods test worst-case size and worst-case type-II error & P15 routed chapter robust-hypothesis-testing & meta & \hyperref[def:p15-robust-test-experiment]{\ref*{def:p15-robust-test-experiment}} & \thepage & Local scope sec:p15-robust-hypothesis-testing. Sample laws are R tensor n with one common law R; A and B later denote sequential boundaries only in their own section. \\
\phantomsection\label{entry:symbol-registry-01643}\(L_i,w_i,\lambda_i,\alpha\) & \texttt{\detokenize{L_i,}\allowbreak\detokenize{w_i,}\allowbreak\detokenize{\lambda_i,}\allowbreak\detokenize{\alpha}} & error losses risk weights hypothesis priors and test level & P15 routed chapter robust-hypothesis-testing & meta & \hyperref[def:p15-robust-test-experiment]{\ref*{def:p15-robust-test-experiment}} & \thepage & Local scope sec:p15-robust-hypothesis-testing. w\_i=lambda\_i L\_i in the source weighted problem; alpha is not a contamination fraction. \\
\phantomsection\label{entry:symbol-registry-01644}\(\varepsilon_i,a_i,H,\mathcal V_i,\varepsilon,\delta\) & \texttt{\detokenize{\varepsilon_i,}\allowbreak\detokenize{a_i,}\allowbreak\detokenize{H,}\allowbreak\detokenize{\mathcal V_i,}\allowbreak\detokenize{\varepsilon,}\allowbreak\detokenize{\delta}} & contamination fractions retained masses arbitrary contamination and variation neighborhoods/radii & P15 routed chapter robust-hypothesis-testing & meta & \hyperref[def:p15-huber-neighborhoods]{\ref*{def:p15-huber-neighborhoods}} & \thepage & Local scope sec:p15-robust-neighborhoods. a\_i=1-epsilon\_i; the source variation radius epsilon is twice the sup-event radius delta. \\
\phantomsection\label{entry:symbol-registry-01645}\(d_{\rm TV},\|\cdot\|_{\rm var},m,d\) & \texttt{\detokenize{d_{\rm TV},}\allowbreak\detokenize{\|\cdot\|_{\rm var},}\allowbreak\detokenize{m,}\allowbreak\detokenize{d}} & sup-event distance total variation norm domination-envelope mass and nominal distance & P15 routed chapter robust-hypothesis-testing & meta & \hyperref[thm:p15-robust-overlap]{\ref*{thm:p15-robust-overlap}} & \thepage & Local scope sec:p15-robust-neighborhoods. Contamination overlap iff m<=1; variation overlap iff epsilon>=d. \\
\phantomsection\label{entry:symbol-registry-01646}\(\mu,p_i,r,F,G,u,v,q_i,Q_i,b,c',c'',\ell\) & \texttt{\detokenize{\mu,}\allowbreak\detokenize{p_i,}\allowbreak\detokenize{r,}\allowbreak\detokenize{F,}\allowbreak\detokenize{G,}\allowbreak\detokenize{u,}\allowbreak\detokenize{v,}\allowbreak\detokenize{q_i,}\allowbreak\detokenize{Q_i,}\allowbreak\detokenize{b,}\allowbreak\detokenize{c',}\allowbreak\detokenize{c'',}\allowbreak\detokenize{\ell}} & dominating measure densities nominal ratio normalization functions clipping parameters least-favorable pair and robust score & P15 routed chapter robust-hypothesis-testing & meta & \hyperref[thm:p15-huber-contamination-pair]{\ref*{thm:p15-huber-contamination-pair}} & \thepage & Local scope sec:p15-contamination-clipping. b=a\_1/a\_0; c'=u and c''=1/v; the score is explicitly assigned on common nominal null sets. \\
\phantomsection\label{entry:symbol-registry-01647}\(s,k_-,k_+\) & \texttt{\detokenize{s,}\allowbreak\detokenize{k_-,}\allowbreak\detokenize{k_+}} & sum density and variation mass-transfer thresholds & P15 routed chapter robust-hypothesis-testing & meta & \hyperref[thm:p15-huber-tv-pair]{\ref*{thm:p15-huber-tv-pair}} & \thepage & Local scope thm:p15-huber-tv-pair. s=p\_0+p\_1; lower and upper transfer integrals both equal epsilon/2. \\
\phantomsection\label{entry:symbol-registry-01648}\(T_n,c,\kappa\) & \texttt{\detokenize{T_n,}\allowbreak\detokenize{c,}\allowbreak\detokenize{\kappa}} & product score decision threshold and constant boundary randomization & P15 routed chapter robust-hypothesis-testing & meta & \hyperref[thm:p15-robust-product-order]{\ref*{thm:p15-robust-product-order}} & \thepage & Local scope sec:p15-robust-minimax. No multiplication of zero by infinity; all-zero-radius cases use joint product densities. \\
\phantomsection\label{entry:symbol-registry-01649}\(q_F,U,Z\) & \texttt{\detokenize{q_F,}\allowbreak\detokenize{U,}\allowbreak\detokenize{Z}} & inverse distribution function uniform coupling variable and transformed score & P15 routed chapter robust-hypothesis-testing & meta & \hyperref[thm:p15-robust-product-order]{\ref*{thm:p15-robust-product-order}} & \thepage & Local scope annex:p15.robust.product.order. Coupling is on real score distributions, not a disintegration on the observation space. \\
\phantomsection\label{entry:symbol-registry-01650}\(e,e_*,D,\Phi\) & \texttt{\detokenize{e,}\allowbreak\detokenize{e_*,}\allowbreak\detokenize{D,}\allowbreak\detokenize{\Phi}} & equal contamination critical radius normal-example normalizer and standard normal cdf & P15 routed chapter robust-hypothesis-testing & meta & \hyperref[thm:p15-robust-sensitivity]{\ref*{thm:p15-robust-sensitivity}} & \thepage & Local scope sec:p15-robust-examples. D is local to the normal example; the finite example uses observation labels a,b,c. \\
\phantomsection\label{entry:symbol-registry-01651}\(\varepsilon_{\rm N},\zeta,u_0\) & \texttt{\detokenize{\varepsilon_{\rm N},}\allowbreak\detokenize{\zeta,}\allowbreak\detokenize{u_0}} & normal-example contamination fraction overshoot remainder and lower essential ratio endpoint & P15 routed chapter robust-hypothesis-testing & meta & \hyperref[thm:p15-robust-testing-calculation]{\ref*{thm:p15-robust-testing-calculation}} & \thepage & Local scope sec:p15-robust-hypothesis-testing. The contamination fraction is distinct from the exponential base; zeta is local to the stopping proof. \\
\phantomsection\label{entry:symbol-registry-01652}\(p_\theta,r_*,S,\theta_0,\theta_1\) & \texttt{\detokenize{p_\theta,}\allowbreak\detokenize{r_*,}\allowbreak\detokenize{S,}\allowbreak\detokenize{\theta_0,}\allowbreak\detokenize{\theta_1}} & ordered density family monotone ratio extension endpoint support and separated parameter boundaries & P15 routed chapter robust-hypothesis-testing & meta & \hyperref[thm:p15-robust-mlr-extension]{\ref*{thm:p15-robust-mlr-extension}} & \thepage & Local scope thm:p15-robust-mlr-extension. The global cross-product condition and common dominating measure are explicit; supports may vary. \\
\phantomsection\label{entry:symbol-registry-01653}\(\gamma,A,B,N,\eta,S_n\) & \texttt{\detokenize{\gamma,}\allowbreak\detokenize{A,}\allowbreak\detokenize{B,}\allowbreak\detokenize{N,}\allowbreak\detokenize{\eta,}\allowbreak\detokenize{S_n}} & log score lower/upper stopping boundaries stopping time affinity and partial log sum & P15 routed chapter robust-hypothesis-testing & meta & \hyperref[thm:p15-robust-sequential-errors]{\ref*{thm:p15-robust-sequential-errors}} & \thepage & Local scope sec:p15-robust-source-extensions. eta=int sqrt(q\_0 q\_1)<1 proves finite expected termination, including infinite-score endpoints. \\
\phantomsection\label{entry:symbol-registry-01654}\(\ell_-,\ell_+,C,\kappa_0,\kappa_1\) & \texttt{\detokenize{\ell_-,}\allowbreak\detokenize{\ell_+,}\allowbreak\detokenize{C,}\allowbreak\detokenize{\kappa_0,}\allowbreak\detokenize{\kappa_1}} & positive finite score bounds overshoot bound and least-favorable drift magnitudes & P15 routed chapter robust-hypothesis-testing & meta & \hyperref[thm:p15-robust-sequential-asymptotic]{\ref*{thm:p15-robust-sequential-asymptotic}} & \thepage & Local scope thm:p15-robust-sequential-asymptotic. Bounded-score premises are additional for the uniform expected-size asymptotic. \\
\phantomsection\label{entry:symbol-registry-01697}\(\mathcal H_i,Q_i,\ell,\gamma,N_{a,b}\) & \texttt{\detokenize{\mathcal H_i,}\allowbreak\detokenize{Q_i,}\allowbreak\detokenize{\ell,}\allowbreak\detokenize{\gamma,}\allowbreak\detokenize{N_{a,b}}} & unchanged neighborhoods least-favorable pair score log score and two-boundary stop & P15 routed chapter robust-sequential-endpoints & meta & \hyperref[def:p15-sequential-endpoint-setting]{\ref*{def:p15-sequential-endpoint-setting}} & \thepage & Local scope sec:p15-robust-sequential-endpoints. A fixed R supplies independent observations on the original arbitrary measurable space. \\
\phantomsection\label{entry:symbol-registry-01698}\(Y_i,Z_i,t_i,u_i,D_i,s_i\) & \texttt{\detokenize{Y_i,}\allowbreak\detokenize{Z_i,}\allowbreak\detokenize{t_i,}\allowbreak\detokenize{u_i,}\allowbreak\detokenize{D_i,}\allowbreak\detokenize{s_i}} & oriented score least-favorable score law boundaries entropy and singular mass & P15 routed chapter robust-sequential-endpoints & meta & \hyperref[def:p15-sequential-endpoint-setting]{\ref*{def:p15-sequential-endpoint-setting}} & \thepage & Local scope sec:p15-robust-sequential-endpoints. Null boundaries are (a,b); alternative boundaries are (b,a); D\_i may be infinite. \\
\phantomsection\label{entry:symbol-registry-01699}\(z^+,z^-\) & \texttt{\detokenize{z^+,}\allowbreak\detokenize{z^-}} & positive and negative parts of an extended score & P15 routed chapter robust-sequential-endpoints & meta & \hyperref[def:p15-sequential-endpoint-setting]{\ref*{def:p15-sequential-endpoint-setting}} & \thepage & Local scope sec:p15-robust-sequential-endpoints. Negative parts are integrable. \\
\phantomsection\label{entry:symbol-registry-01700}\(\tau_t,m(t),d_M\) & \texttt{\detokenize{\tau_t,}\allowbreak\detokenize{m(t),}\allowbreak\detokenize{d_M}} & one-sided passage time mean and upper-truncated score mean & P15 routed chapter robust-sequential-endpoints & meta & \hyperref[def:p15-one-sided-hitting]{\ref*{def:p15-one-sided-hitting}} & \thepage & Local scope sec:p15-robust-sequential-endpoints. tau\_0=m(0)=0; a finite positive truncated drift exists. \\
\phantomsection\label{entry:symbol-registry-01701}\(S_n,W,v\) & \texttt{\detokenize{S_n,}\allowbreak\detokenize{W,}\allowbreak\detokenize{v}} & oriented partial sum wrong-exit event and overshoot magnitude & P15 routed chapter robust-sequential-endpoints & meta & \hyperref[thm:p15-hitting-restart]{\ref*{thm:p15-hitting-restart}} & \thepage & Local scope thm:p15-hitting-restart. v=-S\_N on the wrong-exit event. \\
\phantomsection\label{entry:symbol-registry-01702}\(\epsilon_M(u),\delta_M(p),p,K\) & \texttt{\detokenize{\epsilon_M(u),}\allowbreak\detokenize{\delta_M(p),}\allowbreak\detokenize{p,}\allowbreak\detokenize{K}} & boundary/error comparison remainders actual error and overshoot split & P15 routed chapter robust-sequential-endpoints & meta & \hyperref[thm:p15-two-boundary-passage-comparison]{\ref*{thm:p15-two-boundary-passage-comparison}} & \thepage & Local scope thm:p15-two-boundary-passage-comparison. Uniform in the correct boundary; delta\_M(0)=0. \\
\phantomsection\label{entry:symbol-registry-01703}\(\mathcal M_i(\alpha,\beta),T,\Delta\) & \texttt{\detokenize{\mathcal M_i(\alpha,}\allowbreak\detokenize{\beta),}\allowbreak\detokenize{T,}\allowbreak\detokenize{\Delta}} & minimax expected cost competing stop and terminal decision & P15 routed chapter robust-sequential-endpoints & meta & \hyperref[def:p15-sequential-procedure-cost]{\ref*{def:p15-sequential-procedure-cost}} & \thepage & Local scope sec:p15-sequential-cost-minimax. A common independent random seed is allowed. \\
\phantomsection\label{entry:symbol-registry-01704}\(g_i,\rho,L_r,C_i,T_\infty\) & \texttt{\detokenize{g_i,}\allowbreak\detokenize{\rho,}\allowbreak\detokenize{L_r,}\allowbreak\detokenize{C_i,}\allowbreak\detokenize{T_\infty}} & stopped densities seed law history density singular event and first singular observation & P15 routed chapter robust-sequential-endpoints & meta & \hyperref[thm:p15-stopped-binary-information]{\ref*{thm:p15-stopped-binary-information}} & \thepage & Local scope sec:p15-sequential-cost-minimax. Stopped entropy and singular cost bounds are derived. \\
\phantomsection\label{entry:symbol-registry-01705}\(k,L,P,Q\) & \texttt{\detokenize{k,}\allowbreak\detokenize{L,}\allowbreak\detokenize{P,}\allowbreak\detokenize{Q}} & normalizing series negative score magnitude and countable example laws & P15 routed chapter robust-sequential-endpoints & meta & \hyperref[thm:p15-infinite-entropy-finite-scores]{\ref*{thm:p15-infinite-entropy-finite-scores}} & \thepage & Local scope sec:p15-endpoint-examples. Full support and finite scores coexist with infinite forward entropy. \\
\phantomsection\label{entry:symbol-registry-01706}\(c,r\) & \texttt{\detokenize{c,}\allowbreak\detokenize{r}} & negative log-score magnitude and zero-run stopping length & P15 routed chapter robust-sequential-endpoints & meta & \hyperref[thm:p15-singular-contamination-stopping]{\ref*{thm:p15-singular-contamination-stopping}} & \thepage & Local scope thm:p15-singular-contamination-stopping. c=log(5/3) and r=ceil(a/c). \\
\phantomsection\label{entry:symbol-registry-01722}\(\rho,Q_n,\mathcal T_n,T_n,\psi_-,\psi_+,\psi\) & \texttt{\detokenize{\rho,}\allowbreak\detokenize{Q_n,}\allowbreak\detokenize{\mathcal T_n,}\allowbreak\detokenize{T_n,}\allowbreak\detokenize{\psi_-,}\allowbreak\detokenize{\psi_+,}\allowbreak\detokenize{\psi}} & convex residual criterion empirical sum minimizing interval selection and derivative versions & P15 routed chapter huber-robust-location & meta & \hyperref[def:p15-convex-location-criterion]{\ref*{def:p15-convex-location-criterion}} & \thepage & Local scope sec:p15-convex-location-estimation. The midpoint selector is measurable and translation equivariant; a general selector need not preserve equivariance. \\
\phantomsection\label{entry:symbol-registry-01723}\(\lambda,c,B(t),B,a,a_-,a_+\) & \texttt{\detokenize{\lambda,}\allowbreak\detokenize{c,}\allowbreak\detokenize{B(t),}\allowbreak\detokenize{B,}\allowbreak\detokenize{a,}\allowbreak\detokenize{a_-,}\allowbreak\detokenize{a_+}} & population score center second moment and population derivative magnitudes & P15 routed chapter huber-robust-location & meta & \hyperref[thm:p15-convex-location-consistency]{\ref*{thm:p15-convex-location-consistency}} & \thepage & Local scope sec:p15-convex-location-estimation. One-sided slopes determine a joined-half-normal limit when unequal. \\
\phantomsection\label{entry:symbol-registry-01724}\(k,\rho_k,\psi_k,T_{n,k},\sigma\) & \texttt{\detokenize{k,}\allowbreak\detokenize{\rho_k,}\allowbreak\detokenize{\psi_k,}\allowbreak\detokenize{T_{n,k},}\allowbreak\detokenize{\sigma}} & positive clipping threshold criterion score selected estimate and known scale & P15 routed chapter huber-robust-location & meta & \hyperref[def:p15-huber-location-score]{\ref*{def:p15-huber-location-score}} & \thepage & Local scope sec:p15-huber-location-computation. k=0 is a normalized limiting median procedure, not literal minimization of a constant criterion. \\
\phantomsection\label{entry:symbol-registry-01725}\(I,n_-,n_+\) & \texttt{\detokenize{I,}\allowbreak\detokenize{n_-,}\allowbreak\detokenize{n_+}} & active observations and saturated counts in a score interval & P15 routed chapter huber-robust-location & meta & \hyperref[thm:p15-huber-location-computation]{\ref*{thm:p15-huber-location-computation}} & \thepage & Local scope sec:p15-huber-location-computation. This I is an index set; information is I(f) or I\_* below. \\
\phantomsection\label{entry:symbol-registry-01726}\(\phi,\Phi,\varepsilon,H,\mathcal C_\varepsilon,\xi\) & \texttt{\detokenize{\phi,}\allowbreak\detokenize{\Phi,}\allowbreak\detokenize{\varepsilon,}\allowbreak\detokenize{H,}\allowbreak\detokenize{\mathcal C_\varepsilon,}\allowbreak\detokenize{\xi}} & normal density/cdf contamination fraction symmetric contaminant class and location & P15 routed chapter huber-robust-location & meta & \hyperref[def:p15-huber-symmetric-location-model]{\ref*{def:p15-huber-symmetric-location-model}} & \thepage & Local scope sec:p15-huber-normal-neighborhood. The prototype has known normal scale one; H may have atoms or infinite moments. \\
\phantomsection\label{entry:symbol-registry-01727}\(\mathcal R_\varepsilon,C\) & \texttt{\detokenize{\mathcal R_\varepsilon,}\allowbreak\detokenize{C}} & limiting truncated-risk criterion and truncation level & P15 routed chapter huber-robust-location & meta & \hyperref[def:p15-huber-symmetric-location-model]{\ref*{def:p15-huber-symmetric-location-model}} & \thepage & Local scope sec:p15-huber-normal-neighborhood. The order is sup\_F after limit\_C after liminf\_n; no interchange is claimed. \\
\phantomsection\label{entry:symbol-registry-01728}\(A_F(k),A(k),B(k),V_k(F),V^*(k,\varepsilon)\) & \texttt{\detokenize{A_F(k),}\allowbreak\detokenize{A(k),}\allowbreak\detokenize{B(k),}\allowbreak\detokenize{V_k(F),}\allowbreak\detokenize{V^*(k,}\allowbreak\detokenize{\varepsilon)}} & population slope normal score moments actual and worst asymptotic variances & P15 routed chapter huber-robust-location & meta & \hyperref[thm:p15-huber-normal-variance]{\ref*{thm:p15-huber-normal-variance}} & \thepage & Local scope sec:p15-huber-normal-neighborhood. Endpoint atoms carry half weight for symmetric F. \\
\phantomsection\label{entry:symbol-registry-01729}\(f_*,F_*,I_*,I(f),a_\psi\) & \texttt{\detokenize{f_*,}\allowbreak\detokenize{F_*,}\allowbreak\detokenize{I_*,}\allowbreak\detokenize{I(f),}\allowbreak\detokenize{a_\psi}} & least-favorable density/law information and competitor population slope & P15 routed chapter huber-robust-location & meta & \hyperref[thm:p15-huber-least-favorable-normal]{\ref*{thm:p15-huber-least-favorable-normal}} & \thepage & Local scope sec:p15-huber-normal-neighborhood. Density uniqueness uses the positive normal-component lower bound. \\
\phantomsection\label{entry:symbol-registry-01730}\(w,\pi_a,\widehat T,H\) & \texttt{\detokenize{w,}\allowbreak\detokenize{\pi_a,}\allowbreak\detokenize{\widehat T,}\allowbreak\detokenize{H}} & cosine-squared prior rescaled prior clipped estimator and joint score & P15 routed chapter huber-robust-location & meta & \hyperref[thm:p15-location-integrated-score-bound]{\ref*{thm:p15-location-integrated-score-bound}} & \thepage & Local scope annex:p15.location.integrated.score.bound. H is local joint score here, distinct from the contamination law in the body. \\
\phantomsection\label{entry:symbol-registry-01731}\(c(F),c_*,m_k\) & \texttt{\detokenize{c(F),}\allowbreak\detokenize{c_*,}\allowbreak\detokenize{m_k}} & asymmetric contamination displacement maximal bias and normal population score & P15 routed chapter huber-robust-location & meta & \hyperref[thm:p15-huber-asymmetric-bias]{\ref*{thm:p15-huber-asymmetric-bias}} & \thepage & Local scope sec:p15-huber-location-boundaries. The displayed root is exact; the linear expression is explicitly an approximation. \\
\phantomsection\label{entry:symbol-registry-01732}\(\mathcal M_{\leq1},\mathcal D,C_c,C_c^1,C_0\) & \texttt{\detokenize{\mathcal M_{\leq1},}\allowbreak\detokenize{\mathcal D,}\allowbreak\detokenize{C_c,}\allowbreak\detokenize{C_c^1,}\allowbreak\detokenize{C_0}} & subprobability domain test scores compact support classes and functions vanishing at infinity & P15 routed chapter huber-information-game & meta & \hyperref[def:p15-information-game-domain]{\ref*{def:p15-information-game-domain}} & \thepage & Local scope sec:p15-information-game-foundations. Vague topology uses compactly supported continuous tests. \\
\phantomsection\label{entry:symbol-registry-01733}\(A_\psi,B_\psi,J,K,I\) & \texttt{\detokenize{A_\psi,}\allowbreak\detokenize{B_\psi,}\allowbreak\detokenize{J,}\allowbreak\detokenize{K,}\allowbreak\detokenize{I}} & derivative moment score second moment affine payoff ratio payoff and information & P15 routed chapter huber-information-game & meta & \hyperref[def:p15-information-game-domain]{\ref*{def:p15-information-game-domain}} & \thepage & Local scope sec:p15-information-game-foundations. Zero denominators and the zero measure have explicit conventions. \\
\phantomsection\label{entry:symbol-registry-01734}\(C,\Pi,\overline F,\alpha\) & \texttt{\detokenize{C,}\allowbreak\detokenize{\Pi,}\allowbreak\detokenize{\overline F,}\allowbreak\detokenize{\alpha}} & convex compact model class mixed strategy barycenter and game value & P15 routed chapter huber-information-game & meta & \hyperref[thm:p15-information-payoff-bridge]{\ref*{thm:p15-information-payoff-bridge}} & \thepage & Local scope sec:p15-information-game-foundations. Pi is a probability on the compact space of measures. \\
\phantomsection\label{entry:symbol-registry-01735}\(S,a_g,b_g,r_g\) & \texttt{\detokenize{S,}\allowbreak\detokenize{a_g,}\allowbreak\detokenize{b_g,}\allowbreak\detokenize{r_g}} & positive support efficient-score cross moment second moment and new-support information & P15 routed chapter huber-information-game & meta & \hyperref[thm:p15-information-first-variation]{\ref*{thm:p15-information-first-variation}} & \thepage & Local scope sec:p15-information-support-correction. The term r\_g cannot be dropped without domination. \\
\phantomsection\label{entry:symbol-registry-01736}\(a_\psi,b_\psi,K_{\rm d},\chi\) & \texttt{\detokenize{a_\psi,}\allowbreak\detokenize{b_\psi,}\allowbreak\detokenize{K_{\rm d},}\allowbreak\detokenize{\chi}} & extended-score moments density payoff and specified weak derivative representative & P15 routed chapter huber-information-game & meta & \hyperref[def:p15-extended-information-score]{\ref*{def:p15-extended-information-score}} & \thepage & Local scope sec:p15-information-support-correction. Values of chi on atoms are part of a strategy. \\
\phantomsection\label{entry:symbol-registry-01737}\(C_1,d_F\) & \texttt{\detokenize{C_1,}\allowbreak\detokenize{d_F}} & finite-information members and actual population location slope & P15 routed chapter huber-information-game & meta & \hyperref[thm:p15-information-closure-extension]{\ref*{thm:p15-information-closure-extension}} & \thepage & Local scope sec:p15-information-support-correction. Density of C1 is explicit. \\
\phantomsection\label{entry:symbol-registry-01738}\(g,h,h_k,q_k,Z(k),t_0,t_1\) & \texttt{\detokenize{g,}\allowbreak\detokenize{h,}\allowbreak\detokenize{h_k,}\allowbreak\detokenize{q_k,}\allowbreak\detokenize{Z(k),}\allowbreak\detokenize{t_0,}\allowbreak\detokenize{t_1}} & prototype potential clipped envelope unnormalized density normalizer and tangent endpoints & P15 routed chapter huber-information-game & meta & \hyperref[def:p15-logconcave-prototype]{\ref*{def:p15-logconcave-prototype}} & \thepage & Local scope sec:p15-huber-general-prototype. Infinite endpoints have omitted tail terms. \\
\phantomsection\label{entry:symbol-registry-01739}\(k,\psi_k,f_*,F_*,I_*,\mathcal C_k^0\) & \texttt{\detokenize{k,}\allowbreak\detokenize{\psi_k,}\allowbreak\detokenize{f_*,}\allowbreak\detokenize{F_*,}\allowbreak\detokenize{I_*,}\allowbreak\detokenize{\mathcal C_k^0}} & clipping threshold efficient score least-favorable density law information and centered class & P15 routed chapter huber-information-game & meta & \hyperref[thm:p15-huber-general-density]{\ref*{thm:p15-huber-general-density}} & \thepage & Local scope sec:p15-huber-general-prototype. The asymmetric centered class uses the fixed constructed score. \\
\phantomsection\label{entry:symbol-registry-01740}\(U_n,a_n,b_n\) & \texttt{\detokenize{U_n,}\allowbreak\detokenize{a_n,}\allowbreak\detokenize{b_n}} & sample score and endpoints of its zero interval & P15 routed chapter huber-information-game & meta & \hyperref[sec:p15-huber-general-prototype]{\ref*{sec:p15-huber-general-prototype}} & \thepage & Local scope sec:p15-huber-general-prototype. Separate strict-sign bisections locate both endpoints for midpoint selection. \\
\phantomsection\label{entry:symbol-registry-01778}\(x_{(j)},f_0\) & \texttt{\detokenize{x_{(j)},}\allowbreak\detokenize{f_0}} & ordered observation and local median density & P15 routed chapter huber-remaining-methods & meta & \hyperref[thm:p15-huber-mean-median]{\ref*{thm:p15-huber-mean-median}} & \thepage & Local scope sec:p15-huber-median-trimming. No finite observation variance is required for the median statement. \\
\phantomsection\label{entry:symbol-registry-01779}\(\rho_k^{\rm tr},Q_n^{\rm tr},\psi_k^{\rm tr},A,D\) & \texttt{\detokenize{\rho_k^{\rm tr},}\allowbreak\detokenize{Q_n^{\rm tr},}\allowbreak\detokenize{\psi_k^{\rm tr},}\allowbreak\detokenize{A,}\allowbreak\detokenize{D}} & capped quadratic empirical criterion score active set and population slope quantity & P15 routed chapter huber-remaining-methods & meta & \hyperref[def:p15-huber-hard-trimming]{\ref*{def:p15-huber-hard-trimming}} & \thepage & Local scope sec:p15-huber-median-trimming. The active set A is distinct from the normal score moment A(k). \\
\phantomsection\label{entry:symbol-registry-01780}\(c_\varepsilon,c_{\rm flat}\) & \texttt{\detokenize{c_\varepsilon,}\allowbreak\detokenize{c_{\rm flat}}} & exact contamination median bias and historical flat-density lower bound & P15 routed chapter huber-remaining-methods & meta & \hyperref[thm:p15-huber-exact-median-bias]{\ref*{thm:p15-huber-exact-median-bias}} & \thepage & Local scope sec:p15-huber-median-bias. The flat value is a first-order approximation and a smaller lower bound. \\
\phantomsection\label{entry:symbol-registry-01781}\(A(k),B(k),W(k,\varepsilon),\varepsilon_{\min}\) & \texttt{\detokenize{A(k),}\allowbreak\detokenize{B(k),}\allowbreak\detokenize{W(k,}\allowbreak\detokenize{\varepsilon),}\allowbreak\detokenize{\varepsilon_{\min}}} & normal score moments variance bound and historical tuning fraction & P15 routed chapter huber-remaining-methods & meta & \hyperref[sec:p15-huber-table-one]{\ref*{sec:p15-huber-table-one}} & \thepage & Local scope sec:p15-huber-table-one. Variance entries concern square-root-n times the estimate. \\
\phantomsection\label{entry:symbol-registry-01782}\(B_F,m,\mathcal K_\varepsilon\) & \texttt{\detokenize{B_F,}\allowbreak\detokenize{m,}\allowbreak\detokenize{\mathcal K_\varepsilon}} & completed symmetric cumulative function real-line mass and CDF neighborhood & P15 routed chapter huber-remaining-methods & meta & \hyperref[def:p15-huber-cdf-neighborhood]{\ref*{def:p15-huber-cdf-neighborhood}} & \thepage & Local scope sec:p15-huber-cdf-neighborhood. Missing mass is split equally between the two infinities. \\
\phantomsection\label{entry:symbol-registry-01783}\(a,b,c,z,E(a),\varepsilon_*,f_*,\psi_*,\chi_*,I_*\) & \texttt{\detokenize{a,}\allowbreak\detokenize{b,}\allowbreak\detokenize{c,}\allowbreak\detokenize{z,}\allowbreak\detokenize{E(a),}\allowbreak\detokenize{\varepsilon_*,}\allowbreak\detokenize{f_*,}\allowbreak\detokenize{\psi_*,}\allowbreak\detokenize{\chi_*,}\allowbreak\detokenize{I_*}} & CDF construction joins core parameter tuning function crossing density score derivative and information & P15 routed chapter huber-remaining-methods & meta & \hyperref[thm:p15-huber-cdf-construction]{\ref*{thm:p15-huber-cdf-construction}} & \thepage & Local scope sec:p15-huber-cdf-neighborhood. Use actual one-sided derivative averages when a and b merge. \\
\phantomsection\label{entry:symbol-registry-01784}\(Y,\tau,g_Y,G_Y,\chi_c,u_-,u_+,q,\ell,u,j\) & \texttt{\detokenize{Y,}\allowbreak\detokenize{\tau,}\allowbreak\detokenize{g_Y,}\allowbreak\detokenize{G_Y,}\allowbreak\detokenize{\chi_c,}\allowbreak\detokenize{u_-,}\allowbreak\detokenize{u_+,}\allowbreak\detokenize{q,}\allowbreak\detokenize{\ell,}\allowbreak\detokenize{u,}\allowbreak\detokenize{j}} & log squared observation log scale transformed density cumulative function score tangents cutoff caps and nonzero count & P15 routed chapter huber-remaining-methods & meta & \hyperref[def:p15-huber-scale-construction]{\ref*{def:p15-huber-scale-construction}} & \thepage & Local scope sec:p15-huber-log-scale. Zero observations receive only the separately defined algebraic extension. G\_Y is the transformed normal cumulative function. \\
\phantomsection\label{entry:symbol-registry-01785}\(S_n,S_\infty,S_n',q',I_Y\) & \texttt{\detokenize{S_n,}\allowbreak\detokenize{S_\infty,}\allowbreak\detokenize{S_n',}\allowbreak\detokenize{q',}\allowbreak\detokenize{I_Y}} & raw scale normal limit calibrated scale calibrated cutoff and log-scale information & P15 routed chapter huber-remaining-methods & meta & \hyperref[thm:p15-huber-scale-model]{\ref*{thm:p15-huber-scale-model}} & \thepage & Local scope sec:p15-huber-log-scale. Normal calibration and strict-tail calibration have different targets. \\
\phantomsection\label{entry:symbol-registry-01786}\(T,S,k,\beta,m_1,m_2,m_3,\overline x_2\) & \texttt{\detokenize{T,}\allowbreak\detokenize{S,}\allowbreak\detokenize{k,}\allowbreak\detokenize{\beta,}\allowbreak\detokenize{m_1,}\allowbreak\detokenize{m_2,}\allowbreak\detokenize{m_3,}\allowbreak\detokenize{\overline x_2}} & joint estimates tuning moment level partition counts and interior mean & P15 routed chapter huber-remaining-methods & meta & \hyperref[def:p15-huber-three-proposals]{\ref*{def:p15-huber-three-proposals}} & \thepage & Local scope sec:p15-huber-joint-proposals. Finite existence requires positive beta and the exact multiplicity condition. \\
\phantomsection\label{entry:symbol-registry-01787}\(S_0,q,q_{\max},r,V,\widehat v_n\) & \texttt{\detokenize{S_0,}\allowbreak\detokenize{q,}\allowbreak\detokenize{q_{\max},}\allowbreak\detokenize{r,}\allowbreak\detokenize{V,}\allowbreak\detokenize{\widehat v_n}} & population scale effective cutoff maximal cutoff moment ratio limit variance and proposed variance estimate & P15 routed chapter huber-remaining-methods & meta & \hyperref[thm:p15-huber-joint-regular-limit]{\ref*{thm:p15-huber-joint-regular-limit}} & \thepage & Local scope sec:p15-huber-joint-proposals. Joint normality here excludes atoms at the population cutoff. \\
\phantomsection\label{entry:symbol-registry-01788}\(r,\rho_r,Q_{n,r},h_t,U_n,p,R_n,D\) & \texttt{\detokenize{r,}\allowbreak\detokenize{\rho_r,}\allowbreak\detokenize{Q_{n,r},}\allowbreak\detokenize{h_t,}\allowbreak\detokenize{U_n,}\allowbreak\detokenize{p,}\allowbreak\detokenize{R_n,}\allowbreak\detokenize{D}} & subset order criterion kernel normalized score first projection remainder and population slope & P15 routed chapter huber-remaining-methods & meta & \hyperref[def:p15-huber-mr-estimator]{\ref*{def:p15-huber-mr-estimator}} & \thepage & Local scope sec:p15-huber-mr. This r is a fixed subset size and differs from the joint moment ratio. \\
\phantomsection\label{entry:symbol-registry-01789}\(Z_{ij},H(t),U_n(t),p_t\) & \texttt{\detokenize{Z_{ij},}\allowbreak\detokenize{H(t),}\allowbreak\detokenize{U_n(t),}\allowbreak\detokenize{p_t}} & distinct-pair average population and empirical CDF and centered projection & P15 routed chapter huber-remaining-methods & meta & \hyperref[thm:p15-huber-pair-median]{\ref*{thm:p15-huber-pair-median}} & \thepage & Local scope sec:p15-huber-mr. The overlapping pair averages are dependent; diagonal pairs are omitted. \\
\phantomsection\label{entry:symbol-registry-01790}\(k,\beta,\psi,\rho,L_n,F,m_0,S_0\) & \texttt{\detokenize{k,}\allowbreak\detokenize{\beta,}\allowbreak\detokenize{\psi,}\allowbreak\detokenize{\rho,}\allowbreak\detokenize{L_n,}\allowbreak\detokenize{F,}\allowbreak\detokenize{m_0,}\allowbreak\detokenize{S_0}} & clipping target moment score potential criterion law zero atom and centered scale & P15 routed chapter huber-joint-endpoints & meta & \hyperref[def:p15-huber-joint-boundary-model]{\ref*{def:p15-huber-joint-boundary-model}} & \thepage & Local scope sec:p15-huber-joint-identification. Criterion differences are centered before expectation. \\
\phantomsection\label{entry:symbol-registry-01791}\(U,q,A,C,p,\eta^2\) & \texttt{\detokenize{U,}\allowbreak\detokenize{q,}\allowbreak\detokenize{A,}\allowbreak\detokenize{C,}\allowbreak\detokenize{p,}\allowbreak\detokenize{\eta^2}} & standardized observation cutoff interior moments each clipping atom and second score variance & P15 routed chapter huber-joint-endpoints & meta & \hyperref[def:p15-huber-joint-directional-constants]{\ref*{def:p15-huber-joint-directional-constants}} & \thepage & Local scope sec:p15-huber-joint-directional. A and C are positive under identification. \\
\phantomsection\label{entry:symbol-registry-01792}\(h,a,b,x_+,v_+,v_-,\Gamma,\mathcal D\) & \texttt{\detokenize{h,}\allowbreak\detokenize{a,}\allowbreak\detokenize{b,}\allowbreak\detokenize{x_+,}\allowbreak\detokenize{v_+,}\allowbreak\detokenize{v_-,}\allowbreak\detokenize{\Gamma,}\allowbreak\detokenize{\mathcal D}} & local parameters positive part boundary directions convex potential and directional map & P15 routed chapter huber-joint-endpoints & meta & \hyperref[def:p15-huber-joint-directional-constants]{\ref*{def:p15-huber-joint-directional-constants}} & \thepage & Local scope sec:p15-huber-joint-directional. The directional map is nonlinear when p is positive. \\
\phantomsection\label{entry:symbol-registry-01793}\(e_+,e_-,M_e\) & \texttt{\detokenize{e_+,}\allowbreak\detokenize{e_-,}\allowbreak\detokenize{M_e}} & cone activation indicators and positive definite matrix piece & P15 routed chapter huber-joint-endpoints & meta & \hyperref[thm:p15-huber-joint-cone-inverse]{\ref*{thm:p15-huber-joint-cone-inverse}} & \thepage & Local scope sec:p15-huber-joint-directional. Boundary inverse values agree. \\
\phantomsection\label{entry:symbol-registry-01794}\(H_n,H,Z,Z_1,Z_2,f_H\) & \texttt{\detokenize{H_n,}\allowbreak\detokenize{H,}\allowbreak\detokenize{Z,}\allowbreak\detokenize{Z_1,}\allowbreak\detokenize{Z_2,}\allowbreak\detokenize{f_H}} & joint normalized error limiting vector Gaussian driver and limit density & P15 routed chapter huber-joint-endpoints & meta & \hyperref[thm:p15-huber-joint-atomic-limit]{\ref*{thm:p15-huber-joint-atomic-limit}} & \thepage & Local scope sec:p15-huber-joint-directional. H is a random vector; the contaminating law in the bound is H\_c. \\
\phantomsection\label{entry:symbol-registry-01795}\(R,\widehat v_n,H_{\rm c}\) & \texttt{\detokenize{R,}\allowbreak\detokenize{\widehat v_n,}\allowbreak\detokenize{H_{\rm c}}} & truncation proposed variance estimate and contaminating law & P15 routed chapter huber-joint-endpoints & meta & \hyperref[thm:p15-huber-joint-law-and-variance]{\ref*{thm:p15-huber-joint-law-and-variance}} & \thepage & Local scope sec:p15-huber-joint-directional. The proposed variance estimate has a nonconstant boundary-law limit. \\
\phantomsection\label{entry:symbol-registry-01796}\(r,w,s_1,s_2,\theta_-,\theta_+,F_c,J\) & \texttt{\detokenize{r,}\allowbreak\detokenize{w,}\allowbreak\detokenize{s_1,}\allowbreak\detokenize{s_2,}\allowbreak\detokenize{\theta_-,}\allowbreak\detokenize{\theta_+,}\allowbreak\detokenize{F_c,}\allowbreak\detokenize{J}} & worked inverse coordinates scales angular endpoints and integral primitives & P15 routed chapter huber-joint-endpoints & meta & \hyperref[thm:p15-huber-joint-worked-atomic-law]{\ref*{thm:p15-huber-joint-worked-atomic-law}} & \thepage & Local scope annex:p15.huber.joint.worked.atomic.law. The variance is evaluated from an exact angular integral. \\
\phantomsection\label{entry:symbol-registry-01797}\(\Delta,\alpha,d,\ell,u,v,a_0\) & \texttt{\detokenize{\Delta,}\allowbreak\detokenize{\alpha,}\allowbreak\detokenize{d,}\allowbreak\detokenize{\ell,}\allowbreak\detokenize{u,}\allowbreak\detokenize{v,}\allowbreak\detokenize{a_0}} & determinant side-cone coefficients clipping slopes comparison variance and normal weight & P15 routed chapter huber-joint-endpoints & meta & \hyperref[thm:p15-huber-joint-normal-contamination-counterexample]{\ref*{thm:p15-huber-joint-normal-contamination-counterexample}} & \thepage & Local scope annex:p15.huber.joint.normal.contamination.counterexample. Tail and moment arguments prove a non-Gaussian marginal. \\
\phantomsection\label{entry:symbol-registry-01798}\(\mathcal F_q,S_n,\sigma F,\mathcal R_q\) & \texttt{\detokenize{\mathcal F_q,}\allowbreak\detokenize{S_n,}\allowbreak\detokenize{\sigma F,}\allowbreak\detokenize{\mathcal R_q}} & strict-tail scale class competing estimate scaled law and ordered risk & P15 routed chapter huber-joint-endpoints & meta & \hyperref[def:p15-huber-scale-sign-class]{\ref*{def:p15-huber-scale-sign-class}} & \thepage & Local scope sec:p15-huber-scale-signs. Competing estimators may use all signs and magnitudes. \\
\phantomsection\label{entry:symbol-registry-01799}\(f_*,g_*,s,T_{n,s},I_Y\) & \texttt{\detokenize{f_*,}\allowbreak\detokenize{g_*,}\allowbreak\detokenize{s,}\allowbreak\detokenize{T_{n,s},}\allowbreak\detokenize{I_Y}} & least-favorable densities sign pattern conditional estimator and information & P15 routed chapter huber-joint-endpoints & meta & \hyperref[thm:p15-huber-scale-sign-minimax]{\ref*{thm:p15-huber-scale-sign-minimax}} & \thepage & Local scope sec:p15-huber-scale-signs. Sign conditioning uses the actual least-favorable law. \\
\phantomsection\label{entry:symbol-registry-01800}\(\psi_\ell,A(\ell),B(\ell)\) & \texttt{\detokenize{\psi_\ell,}\allowbreak\detokenize{A(\ell),}\allowbreak\detokenize{B(\ell)}} & varying-cutoff score and normal score moments & P15 routed chapter huber-joint-endpoints & meta & \hyperref[def:p15-huber-varying-cutoff-notation]{\ref*{def:p15-huber-varying-cutoff-notation}} & \thepage & Local scope sec:p15-huber-joint-directional. A(ell) is distinct from the bare interior probability A. \\
\phantomsection\label{entry:symbol-registry-02012}\(H\) & \texttt{\detokenize{H}} & Bayesian model & \hyperref[sec:p15-posterior-predictive-model-criticism]{\ref*{sec:p15-posterior-predictive-model-criticism}} & n/a & \hyperref[def:p15-ppc-model]{\ref*{def:p15-ppc-model}} & \thepage & Includes sampling specification and prior \\
\phantomsection\label{entry:symbol-registry-02013}\(\Theta\) & \texttt{\detokenize{\Theta}} & parameter space & \hyperref[sec:p15-posterior-predictive-model-criticism]{\ref*{sec:p15-posterior-predictive-model-criticism}} & n/a & \hyperref[def:p15-ppc-model]{\ref*{def:p15-ppc-model}} & \thepage & Standard Borel space \\
\phantomsection\label{entry:symbol-registry-02014}\(\vartheta\) & \texttt{\detokenize{\vartheta}} & random parameter under the proper joint prior model & \hyperref[sec:p15-ppc-calibration]{\ref*{sec:p15-ppc-calibration}} & n/a & \hyperref[thm:p15-ppc-pivot]{\ref*{thm:p15-ppc-pivot}} & \thepage & Distinct from the parameter space \\
\phantomsection\label{entry:symbol-registry-02015}\(\Pi_y\) & \texttt{\detokenize{\Pi_y}} & posterior probability kernel at observed y & \hyperref[sec:p15-posterior-predictive-model-criticism]{\ref*{sec:p15-posterior-predictive-model-criticism}} & n/a & \hyperref[def:p15-ppc-model]{\ref*{def:p15-ppc-model}} & \thepage & Specified version with finite positive normalization \\
\phantomsection\label{entry:symbol-registry-02016}\(R_a\) & \texttt{\detokenize{R_a}} & replication kernel given auxiliary value a & \hyperref[sec:p15-posterior-predictive-model-criticism]{\ref*{sec:p15-posterior-predictive-model-criticism}} & n/a & \hyperref[def:p15-ppc-model]{\ref*{def:p15-ppc-model}} & \thepage & Fixed design and random auxiliary conditioning distinguished \\
\phantomsection\label{entry:symbol-registry-02017}\(A\) & \texttt{\detokenize{A}} & auxiliary statistic & \hyperref[sec:p15-ppc-construction]{\ref*{sec:p15-ppc-construction}} & n/a & \hyperref[def:p15-ppc-model]{\ref*{def:p15-ppc-model}} & \thepage & Not the separately typed detector operator in the tomography case \\
\phantomsection\label{entry:symbol-registry-02018}\(D\) & \texttt{\detokenize{D}} & realized discrepancy function & \hyperref[sec:p15-posterior-predictive-model-criticism]{\ref*{sec:p15-posterior-predictive-model-criticism}} & n/a & \hyperref[def:p15-ppc-discrepancy]{\ref*{def:p15-ppc-discrepancy}} & \thepage & Data and parameter arguments retained \\
\phantomsection\label{entry:symbol-registry-02019}\(q_D\) & \texttt{\detokenize{q_D}} & conditional discrepancy tail at fixed parameter & \hyperref[sec:p15-posterior-predictive-model-criticism]{\ref*{sec:p15-posterior-predictive-model-criticism}} & n/a & \hyperref[def:p15-ppc-discrepancy]{\ref*{def:p15-ppc-discrepancy}} & \thepage & Inclusive upper tail \\
\phantomsection\label{entry:symbol-registry-02020}\(p_B\) & \texttt{\detokenize{p_B}} & posterior predictive discrepancy probability & \hyperref[sec:p15-posterior-predictive-model-criticism]{\ref*{sec:p15-posterior-predictive-model-criticism}} & n/a & \hyperref[def:p15-ppc-discrepancy]{\ref*{def:p15-ppc-discrepancy}} & \thepage & Not posterior model probability \\
\phantomsection\label{entry:symbol-registry-02021}\(Q_y\) & \texttt{\detokenize{Q_y}} & joint posterior parameter and replication law & \hyperref[sec:p15-posterior-predictive-model-criticism]{\ref*{sec:p15-posterior-predictive-model-criticism}} & n/a & \hyperref[thm:p15-ppc-joint]{\ref*{thm:p15-ppc-joint}} & \thepage & Same parameter in observed and replicated discrepancies \\
\phantomsection\label{entry:symbol-registry-02022}\(D_{\min}\) & \texttt{\detokenize{D_{\min}}} & minimum or infimum discrepancy statistic & \hyperref[sec:p15-posterior-predictive-model-criticism]{\ref*{sec:p15-posterior-predictive-model-criticism}} & n/a & \hyperref[def:p15-ppc-reduced-discrepancies]{\ref*{def:p15-ppc-reduced-discrepancies}} & \thepage & Existence and measurability required \\
\phantomsection\label{entry:symbol-registry-02023}\(D_{\rm avg}\) & \texttt{\detokenize{D_{\rm avg}}} & posterior-average discrepancy statistic & \hyperref[sec:p15-posterior-predictive-model-criticism]{\ref*{sec:p15-posterior-predictive-model-criticism}} & n/a & \hyperref[def:p15-ppc-reduced-discrepancies]{\ref*{def:p15-ppc-reduced-discrepancies}} & \thepage & Replicated value refits its own posterior \\
\phantomsection\label{entry:symbol-registry-02024}\(\widehat p\) & \texttt{\detokenize{\widehat p}} & Monte Carlo average of discrepancy indicators & \hyperref[sec:p15-ppc-computation]{\ref*{sec:p15-ppc-computation}} & n/a & \hyperref[thm:p15-ppc-simulation]{\ref*{thm:p15-ppc-simulation}} & \thepage & Independent and correlated cases distinguished \\
\phantomsection\label{entry:symbol-registry-02025}\(\widetilde p\) & \texttt{\detokenize{\widetilde p}} & Monte Carlo average of conditional tails & \hyperref[sec:p15-ppc-computation]{\ref*{sec:p15-ppc-computation}} & n/a & \hyperref[thm:p15-ppc-simulation]{\ref*{thm:p15-ppc-simulation}} & \thepage & Conditional averaging removes replication variance \\
\phantomsection\label{entry:symbol-registry-02026}\(\gamma_h\) & \texttt{\detokenize{\gamma_h}} & lag h covariance of stationary indicators & \hyperref[sec:p15-ppc-computation]{\ref*{sec:p15-ppc-computation}} & n/a & \hyperref[thm:p15-ppc-simulation]{\ref*{thm:p15-ppc-simulation}} & \thepage & No mixing claim inferred from stationarity \\
\phantomsection\label{entry:symbol-registry-02027}\(\Phi\) & \texttt{\detokenize{\Phi}} & standard normal cumulative distribution function & \hyperref[sec:p15-posterior-predictive-model-criticism]{\ref*{sec:p15-posterior-predictive-model-criticism}} & n/a & \hyperref[thm:p15-ppc-normal-mean]{\ref*{thm:p15-ppc-normal-mean}} & \thepage & Not a logical formula in this treatment \\
\phantomsection\label{entry:symbol-registry-02028}\(X^2\) & \texttt{\detokenize{X^2}} & Pearson squared-residual discrepancy & \hyperref[sec:p15-posterior-predictive-model-criticism]{\ref*{sec:p15-posterior-predictive-model-criticism}} & n/a & \hyperref[sec:p15-ppc-linear]{\ref*{sec:p15-ppc-linear}} & \thepage & The name alone does not establish a chi-square distribution \\
\phantomsection\label{entry:symbol-registry-02029}\(B\) & \texttt{\detokenize{B}} & full-rank Gaussian design matrix & \hyperref[sec:p15-ppc-linear]{\ref*{sec:p15-ppc-linear}} & n/a & \hyperref[thm:p15-ppc-linear]{\ref*{thm:p15-ppc-linear}} & \thepage & Distinct from measurable set arguments of kernels \\
\phantomsection\label{entry:symbol-registry-02030}\(P\) & \texttt{\detokenize{P}} & orthogonal projection onto the design column space & \hyperref[sec:p15-ppc-linear]{\ref*{sec:p15-ppc-linear}} & n/a & \hyperref[thm:p15-ppc-linear]{\ref*{thm:p15-ppc-linear}} & \thepage & Matrix rather than a probability measure in this local theorem \\
\phantomsection\label{entry:symbol-registry-02031}\(r(y)\) & \texttt{\detokenize{r(y)}} & squared orthogonal residual norm & \hyperref[sec:p15-ppc-linear]{\ref*{sec:p15-ppc-linear}} & n/a & \hyperref[thm:p15-ppc-linear]{\ref*{thm:p15-ppc-linear}} & \thepage & Not the detector correction vector in the later application \\
\phantomsection\label{entry:symbol-registry-02032}\(A\) & \texttt{\detokenize{A}} & nonnegative detector operator & \hyperref[sec:p15-ppc-constrained-applications]{\ref*{sec:p15-ppc-constrained-applications}} & n/a & \hyperref[sec:p15-ppc-constrained-applications]{\ref*{sec:p15-ppc-constrained-applications}} & \thepage & Local tomography role \\
\phantomsection\label{entry:symbol-registry-02033}\(C\) & \texttt{\detokenize{C}} & increasing convex nonnegative rate region & \hyperref[sec:p15-ppc-constrained-applications]{\ref*{sec:p15-ppc-constrained-applications}} & n/a & \hyperref[thm:p15-ppc-mortality-posterior]{\ref*{thm:p15-ppc-mortality-posterior}} & \thepage & Bounded probability-rate variant stated separately \\
\phantomsection\label{entry:symbol-registry-02034}\(n_{habc}\) & \texttt{\detokenize{n_{habc}}} & augmented class and cell counts & \hyperref[sec:p15-ppc-latent-classes]{\ref*{sec:p15-ppc-latent-classes}} & n/a & \hyperref[def:p15-ppc-multinomial-deviance]{\ref*{def:p15-ppc-multinomial-deviance}} & \thepage & Plus signs in subscripts denote margins in the proof \\
\phantomsection\label{entry:symbol-registry-02035}\(\kappa_i\) & \texttt{\detokenize{\kappa_i}} & positive Dirichlet parameter & \hyperref[sec:p15-ppc-latent-classes]{\ref*{sec:p15-ppc-latent-classes}} & n/a & \hyperref[def:p15-ppc-dirichlet]{\ref*{def:p15-ppc-dirichlet}} & \thepage & Each probability block has its own parameter vector \\
\phantomsection\label{entry:symbol-registry-02036}\(M^{\rm mix}_y\) & \texttt{\detokenize{M^{\rm mix}_y}} & mixed predictive data law & \hyperref[sec:p15-ppc-replication-choices]{\ref*{sec:p15-ppc-replication-choices}} & n/a & \hyperref[sec:p15-ppc-replication-choices]{\ref*{sec:p15-ppc-replication-choices}} & \thepage & Fresh group parameter with shared posterior hyperparameter \\
\phantomsection\label{entry:symbol-registry-02037}\(G_t\) & \texttt{\detokenize{G_t}} & conditional distribution function of the absolute maximum & \hyperref[sec:p15-ppc-replication-choices]{\ref*{sec:p15-ppc-replication-choices}} & n/a & \hyperref[thm:p15-ppc-normal-maximum]{\ref*{thm:p15-ppc-normal-maximum}} & \thepage & Defined for t nonnegative \\
\phantomsection\label{entry:symbol-registry-02055}\(X,m,\Sigma\) & \texttt{\detokenize{X,}\allowbreak\detokenize{m,}\allowbreak\detokenize{\Sigma}} & input uncertainty vector its mean and covariance matrix & \hyperref[sec:p15-measurement-model]{\ref*{sec:p15-measurement-model}} & n/a & \hyperref[def:p15-measurement-model]{\ref*{def:p15-measurement-model}} & \thepage & Random uncertainty representation is distinct from a fixed physical quantity. \\
\phantomsection\label{entry:symbol-registry-02056}\(x,U,u_i,u_{ij},r_{ij}\) & \texttt{\detokenize{x,}\allowbreak\detokenize{U,}\allowbreak\detokenize{u_i,}\allowbreak\detokenize{u_{ij},}\allowbreak\detokenize{r_{ij}}} & input estimates assessed covariance standard uncertainties covariances and defined correlations & \hyperref[sec:p15-measurement-model]{\ref*{sec:p15-measurement-model}} & n/a & \hyperref[def:p15-measurement-model]{\ref*{def:p15-measurement-model}} & \thepage & Correlation quotient used only with positive denominator; U is a covariance matrix in this section. \\
\phantomsection\label{entry:symbol-registry-02057}\(Y,f,J,c\) & \texttt{\detokenize{Y,}\allowbreak\detokenize{f,}\allowbreak\detokenize{J,}\allowbreak\detokenize{c}} & output random quantity measurement map Jacobian and scalar sensitivities & \hyperref[sec:p15-measurement-model]{\ref*{sec:p15-measurement-model}}--\hyperref[sec:p15-measurement-nonlinear]{\ref*{sec:p15-measurement-nonlinear}} & n/a & \hyperref[def:p15-measurement-model]{\ref*{def:p15-measurement-model}} & \thepage & Coordinate units are fixed before derivatives and matrix norms are used. \\
\phantomsection\label{entry:symbol-registry-02058}\(P,V,R_0,\alpha,t,t_0\) & \texttt{\detokenize{P,}\allowbreak\detokenize{V,}\allowbreak\detokenize{R_0,}\allowbreak\detokenize{\alpha,}\allowbreak\detokenize{t,}\allowbreak\detokenize{t_0}} & power voltage reference resistance temperature coefficient actual and reference temperature & \hyperref[sec:p15-measurement-model]{\ref*{sec:p15-measurement-model}} & n/a & \hyperref[def:p15-measurement-power]{\ref*{def:p15-measurement-power}} & \thepage & A separate physical model; P is not a probability in this example. \\
\phantomsection\label{entry:symbol-registry-02059}\(Q_k,\bar Q,S,C\) & \texttt{\detokenize{Q_k,}\allowbreak\detokenize{\bar Q,}\allowbreak\detokenize{S,}\allowbreak\detokenize{C}} & independent replicate vectors sample mean sample covariance and population covariance & \hyperref[sec:p15-measurement-input-evaluation]{\ref*{sec:p15-measurement-input-evaluation}} & n/a & \hyperref[thm:p15-measurement-sample-covariance]{\ref*{thm:p15-measurement-sample-covariance}} & \thepage & S estimates covariance of observations; S/n estimates covariance of mean. \\
\phantomsection\label{entry:symbol-registry-02060}\(a,c,\beta,\Phi\) & \texttt{\detokenize{a,}\allowbreak\detokenize{c,}\allowbreak\detokenize{\beta,}\allowbreak\detokenize{\Phi}} & half width center trapezoid top ratio and uniform phase & \hyperref[sec:p15-measurement-input-evaluation]{\ref*{sec:p15-measurement-input-evaluation}} & n/a & \hyperref[thm:p15-measurement-bounded-distributions]{\ref*{thm:p15-measurement-bounded-distributions}} & \thepage & Phi is uniform phase in this theorem; normal CDF Phi is explicitly invoked in the later conversion. \\
\phantomsection\label{entry:symbol-registry-02061}\(A,B,m,b_-,b_+,\lambda,Z,q_\lambda\) & \texttt{\detokenize{A,}\allowbreak\detokenize{B,}\allowbreak\detokenize{m,}\allowbreak\detokenize{b_-,}\allowbreak\detokenize{b_+,}\allowbreak\detokenize{\lambda,}\allowbreak\detokenize{Z,}\allowbreak\detokenize{q_\lambda}} & support endpoints prescribed mean offsets tilt normalizer and density & \hyperref[sec:p15-measurement-input-evaluation]{\ref*{sec:p15-measurement-input-evaluation}} & n/a & \hyperref[thm:p15-measurement-asymmetric-model]{\ref*{thm:p15-measurement-asymmetric-model}} & \thepage & Distinct from offset and scale random variables in the later correction theorem. \\
\phantomsection\label{entry:symbol-registry-02062}\(A,B,Q,W,v_A,v_B,v_Q\) & \texttt{\detokenize{A,}\allowbreak\detokenize{B,}\allowbreak\detokenize{Q,}\allowbreak\detokenize{W,}\allowbreak\detokenize{v_A,}\allowbreak\detokenize{v_B,}\allowbreak\detokenize{v_Q}} & common offset scale averaged input output and respective variances & \hyperref[sec:p15-measurement-linear]{\ref*{sec:p15-measurement-linear}} & n/a & \hyperref[thm:p15-measurement-common-correction]{\ref*{thm:p15-measurement-common-correction}} & \thepage & Independence is explicit; common offset does not shrink with replicate count. \\
\phantomsection\label{entry:symbol-registry-02063}\(L,R,M,r,\delta\) & \texttt{\detokenize{L,}\allowbreak\detokenize{R,}\allowbreak\detokenize{M,}\allowbreak\detokenize{r,}\allowbreak\detokenize{\delta}} & linear surrogate remainder Hessian bound remainder norm bound and displacement & \hyperref[sec:p15-measurement-nonlinear]{\ref*{sec:p15-measurement-nonlinear}} & n/a & \hyperref[thm:p15-measurement-taylor-bound]{\ref*{thm:p15-measurement-taylor-bound}} & \thepage & R is a Taylor remainder here; resistance R is local to electrical example. \\
\phantomsection\label{entry:symbol-registry-02064}\(g,H,T,Z_i,A,B,C,\varepsilon,F_\varepsilon\) & \texttt{\detokenize{g,}\allowbreak\detokenize{H,}\allowbreak\detokenize{T,}\allowbreak\detokenize{Z_i,}\allowbreak\detokenize{A,}\allowbreak\detokenize{B,}\allowbreak\detokenize{C,}\allowbreak\detokenize{\varepsilon,}\allowbreak\detokenize{F_\varepsilon}} & linear quadratic cubic coefficients independent normal variables polynomial pieces scale and surrogate & \hyperref[sec:p15-measurement-nonlinear]{\ref*{sec:p15-measurement-nonlinear}} & n/a & \hyperref[thm:p15-measurement-gaussian-expansion]{\ref*{thm:p15-measurement-gaussian-expansion}} & \thepage & Fully symmetric cubic tensor; exact polynomial identity. \\
\phantomsection\label{entry:symbol-registry-02065}\(u_c,U,k,\Phi,\nu,\nu_{\rm eff}\) & \texttt{\detokenize{u_c,}\allowbreak\detokenize{U,}\allowbreak\detokenize{k,}\allowbreak\detokenize{\Phi,}\allowbreak\detokenize{\nu,}\allowbreak\detokenize{\nu_{\rm eff}}} & combined standard uncertainty expanded uncertainty coverage factor normal CDF and degrees of freedom & \hyperref[sec:p15-measurement-coverage]{\ref*{sec:p15-measurement-coverage}} & n/a & \hyperref[def:p15-measurement-expanded]{\ref*{def:p15-measurement-expanded}} & \thepage & U is now a scalar expanded uncertainty, not the earlier covariance matrix. \\
\phantomsection\label{entry:symbol-registry-02066}\(V_i,v_i,V,v,\chi_\nu^2\) & \texttt{\detokenize{V_i,}\allowbreak\detokenize{v_i,}\allowbreak\detokenize{V,}\allowbreak\detokenize{v,}\allowbreak\detokenize{\chi_\nu^2}} & variance estimators means sums and chi-square variable & \hyperref[sec:p15-measurement-coverage]{\ref*{sec:p15-measurement-coverage}} & n/a & \hyperref[thm:p15-measurement-effective-df]{\ref*{thm:p15-measurement-effective-df}} & \thepage & Matching two moments does not identify a distribution. \\
\phantomsection\label{entry:symbol-registry-02067}\(l,l_{\rm S},d,\alpha,\alpha_{\rm S},\theta,\theta_{\rm S}\) & \texttt{\detokenize{l,}\allowbreak\detokenize{l_{\rm S},}\allowbreak\detokenize{d,}\allowbreak\detokenize{\alpha,}\allowbreak\detokenize{\alpha_{\rm S},}\allowbreak\detokenize{\theta,}\allowbreak\detokenize{\theta_{\rm S}}} & gauge length standard length comparison and thermal parameters & \hyperref[sec:p15-measurement-gauge]{\ref*{sec:p15-measurement-gauge}} & n/a & \hyperref[sec:p15-measurement-gauge]{\ref*{sec:p15-measurement-gauge}} & \thepage & Lengths refer to 20 degrees Celsius; thermal deviations refer to the same reference. \\
\phantomsection\label{entry:symbol-registry-02068}\(A,D,L,\Theta,B,E_0,G,g_0\) & \texttt{\detokenize{A,}\allowbreak\detokenize{D,}\allowbreak\detokenize{L,}\allowbreak\detokenize{\Theta,}\allowbreak\detokenize{B,}\allowbreak\detokenize{E_0,}\allowbreak\detokenize{G,}\allowbreak\detokenize{g_0}} & coefficient difference temperature difference nominal length random temperatures coefficient additive error and gauge surrogate & \hyperref[sec:p15-measurement-gauge]{\ref*{sec:p15-measurement-gauge}} & n/a & \hyperref[thm:p15-measurement-gauge-products]{\ref*{thm:p15-measurement-gauge-products}} & \thepage & Independent-product model is explicitly stated rather than inferred from zero covariance. \\
\phantomsection\label{entry:symbol-registry-02069}\(V,I,\phi,R,X_{\rm r},Z,J,U\) & \texttt{\detokenize{V,}\allowbreak\detokenize{I,}\allowbreak\detokenize{\phi,}\allowbreak\detokenize{R,}\allowbreak\detokenize{X_{\rm r},}\allowbreak\detokenize{Z,}\allowbreak\detokenize{J,}\allowbreak\detokenize{U}} & voltage current phase resistance reactance magnitude Jacobian and covariance of input mean & \hyperref[sec:p15-measurement-impedance]{\ref*{sec:p15-measurement-impedance}} & n/a & \hyperref[sec:p15-measurement-impedance]{\ref*{sec:p15-measurement-impedance}} & \thepage & Source reactance \(X\) is written \(X_{\rm r}\) to distinguish it from the generic input vector; current converted to amperes. \\
\phantomsection\label{entry:symbol-registry-02070}\(q_k,\bar q,U_{\rm mean},H\) & \texttt{\detokenize{q_k,}\allowbreak\detokenize{\bar q,}\allowbreak\detokenize{U_{\rm mean},}\allowbreak\detokenize{H}} & fixed paired rows their mean covariance-of-mean estimate and Hessian & \hyperref[sec:p15-measurement-impedance]{\ref*{sec:p15-measurement-impedance}} & n/a & \hyperref[thm:p15-measurement-transform-average]{\ref*{thm:p15-measurement-transform-average}} & \thepage & Finite averaging correction has factor (n-1)/2 for this covariance normalization. \\
\phantomsection\label{entry:symbol-registry-02071}\(\delta_{ij}\) & \texttt{\detokenize{\delta_{ij}}} & Kronecker indicator for equal indices & \hyperref[annex:p15.measurement.gaussian-expansion]{\ref*{annex:p15.measurement.gaussian-expansion}} & n/a & \hyperref[annex:p15.measurement.gaussian-expansion]{\ref*{annex:p15.measurement.gaussian-expansion}} & \thepage & Used in the explicitly proved independent Gaussian fourth-moment identity. \\
\phantomsection\label{entry:symbol-registry-02072}\(h(p),D(p\Vert q)\) & \texttt{\detokenize{h(p),}\allowbreak\detokenize{D(p\Vert q)}} & differential entropy and relative entropy for densities on a fixed interval & \hyperref[sec:p15-measurement-input-evaluation]{\ref*{sec:p15-measurement-input-evaluation}} & n/a & \hyperref[thm:p15-measurement-asymmetric-model]{\ref*{thm:p15-measurement-asymmetric-model}} & \thepage & Natural logarithms and a fixed numerical coordinate; optimality proved directly from nonnegative relative entropy. \\
\phantomsection\label{entry:symbol-registry-02153}\(\eta_j\) & \texttt{\detokenize{\eta_j}} & nonnegative scaled-rate increment coefficient & \hyperref[sec:p15-ppc-constrained-applications]{\ref*{sec:p15-ppc-constrained-applications}} & n/a & \hyperref[thm:p15-ppc-mortality-fit]{\ref*{thm:p15-ppc-mortality-fit}} & \thepage & The first coefficient is strictly positive for the printed positive counts \\
\phantomsection\label{entry:symbol-registry-02154}\((x)_+\) & \texttt{\detokenize{(x)_+}} & positive part of a real number & \hyperref[sec:p15-ppc-constrained-applications]{\ref*{sec:p15-ppc-constrained-applications}} & n/a & \hyperref[thm:p15-ppc-mortality-fit]{\ref*{thm:p15-ppc-mortality-fit}} & \thepage & Defined as the maximum of x and zero \\
\phantomsection\label{entry:symbol-registry-02155}\(L\) & \texttt{\detokenize{L}} & negative Poisson log-likelihood objective & \hyperref[sec:p15-ppc-constrained-applications]{\ref*{sec:p15-ppc-constrained-applications}} & n/a & \hyperref[thm:p15-ppc-mortality-fit]{\ref*{thm:p15-ppc-mortality-fit}} & \thepage & Data-only additive constants omitted \\
\phantomsection\label{entry:symbol-registry-02156}\(P\) & \texttt{\detokenize{P}} & Pearson discrepancy objective & \hyperref[sec:p15-ppc-constrained-applications]{\ref*{sec:p15-ppc-constrained-applications}} & n/a & \hyperref[thm:p15-ppc-mortality-fit]{\ref*{thm:p15-ppc-mortality-fit}} & \thepage & Distinct from the orthogonal projection in the earlier linear-model section \\
\phantomsection\label{entry:symbol-registry-02315}\([q],d_i,x_i,s\) & \texttt{\detokenize{[q],}\allowbreak\detokenize{d_i,}\allowbreak\detokenize{x_i,}\allowbreak\detokenize{s}} & dimension vector input column numerical value and logarithmic unit change & \hyperref[sec:p15-dimensional-quantities]{\ref*{sec:p15-dimensional-quantities}} & n/a & \hyperref[def:p15-dimensional-coordinate]{\ref*{def:p15-dimensional-coordinate}} & \thepage & Positive numerical quantities and multiplicative unit changes. \\
\phantomsection\label{entry:symbol-registry-02316}\(D,a,x^a,r,n,k\) & \texttt{\detokenize{D,}\allowbreak\detokenize{a,}\allowbreak\detokenize{x^a,}\allowbreak\detokenize{r,}\allowbreak\detokenize{n,}\allowbreak\detokenize{k}} & dimension matrix exponent vector monomial rank and matrix sizes & \hyperref[sec:p15-dimensional-groups]{\ref*{sec:p15-dimensional-groups}} & n/a & \hyperref[def:p15-dimensional-matrix]{\ref*{def:p15-dimensional-matrix}} & \thepage & Independent dimensions determine rank. \\
\phantomsection\label{entry:symbol-registry-02317}\(b_j,B,\pi_j,\pi,\Psi\) & \texttt{\detokenize{b_j,}\allowbreak\detokenize{B,}\allowbreak\detokenize{\pi_j,}\allowbreak\detokenize{\pi,}\allowbreak\detokenize{\Psi}} & nullspace basis exponent-row matrix products and factorization & \hyperref[sec:p15-dimensional-pi]{\ref*{sec:p15-dimensional-pi}} & n/a & \hyperref[thm:p15-dimensional-pi-factorization]{\ref*{thm:p15-dimensional-pi-factorization}} & \thepage & The empty product target is allowed. \\
\phantomsection\label{entry:symbol-registry-02318}\(\odot\) & \texttt{\detokenize{\odot}} & componentwise multiplication of positive coordinate vectors & \hyperref[sec:p15-dimensional-pi]{\ref*{sec:p15-dimensional-pi}} & n/a & \hyperref[thm:p15-dimensional-pi-factorization]{\ref*{thm:p15-dimensional-pi-factorization}} & \thepage & Defined in the positive-domain unit-change factorization. \\
\phantomsection\label{entry:symbol-registry-02319}\(w\) & \texttt{\detokenize{w}} & dimension vector of a covariant numerical output & \hyperref[sec:p15-dimensional-pi]{\ref*{sec:p15-dimensional-pi}} & n/a & \hyperref[thm:p15-dimensional-covariant-output]{\ref*{thm:p15-dimensional-covariant-output}} & \thepage & The output transforms by the stated scalar exponential under a logarithmic unit change. \\
\phantomsection\label{entry:symbol-registry-02320}\(K,\alpha,\beta,p_j,a_{ji},u_j\) & \texttt{\detokenize{K,}\allowbreak\detokenize{\alpha,}\allowbreak\detokenize{\beta,}\allowbreak\detokenize{p_j,}\allowbreak\detokenize{a_{ji},}\allowbreak\detokenize{u_j}} & fixed coefficient exponents and intermediate monomials & \hyperref[sec:p15-dimensional-scaling]{\ref*{sec:p15-dimensional-scaling}} & n/a & \hyperref[def:p15-dimensional-scaling-law]{\ref*{def:p15-dimensional-scaling-law}} & \thepage & Coefficients may carry dimensions. \\
\phantomsection\label{entry:symbol-registry-02321}\(\xi,\eta,x_0,y_0\) & \texttt{\detokenize{\xi,}\allowbreak\detokenize{\eta,}\allowbreak\detokenize{x_0,}\allowbreak\detokenize{y_0}} & log coordinates and reference values & \hyperref[sec:p15-dimensional-scaling]{\ref*{sec:p15-dimensional-scaling}} & n/a & \hyperref[thm:p15-dimensional-log-slope]{\ref*{thm:p15-dimensional-log-slope}} & \thepage & Logarithms apply to numerical ratios. \\
\phantomsection\label{entry:symbol-registry-02322}\(r,v,\omega,T,G,M,g,l,\theta\) & \texttt{\detokenize{r,}\allowbreak\detokenize{v,}\allowbreak\detokenize{\omega,}\allowbreak\detokenize{T,}\allowbreak\detokenize{G,}\allowbreak\detokenize{M,}\allowbreak\detokenize{g,}\allowbreak\detokenize{l,}\allowbreak\detokenize{\theta}} & mechanical radius speed frequency period gravity parameters length and angle & \hyperref[sec:p15-dimensional-mechanical]{\ref*{sec:p15-dimensional-mechanical}}--\hyperref[sec:p15-dimensional-pendulum]{\ref*{sec:p15-dimensional-pendulum}} & n/a & \hyperref[thm:p15-dimensional-mechanical-models]{\ref*{thm:p15-dimensional-mechanical-models}} & \thepage & Explicit geometric and force models. \\
\phantomsection\label{entry:symbol-registry-02323}\(F_{\rm d},c_{\rm d},\rho,A,D_{\rm b},\nu,\mathrm{Re}\) & \texttt{\detokenize{F_{\rm d},}\allowbreak\detokenize{c_{\rm d},}\allowbreak\detokenize{\rho,}\allowbreak\detokenize{A,}\allowbreak\detokenize{D_{\rm b},}\allowbreak\detokenize{\nu,}\allowbreak\detokenize{\mathrm{Re}}} & drag force coefficient density area body length viscosity and Reynolds number & \hyperref[sec:p15-dimensional-drag]{\ref*{sec:p15-dimensional-drag}} & n/a & \hyperref[sec:p15-dimensional-drag]{\ref*{sec:p15-dimensional-drag}} & \thepage & Shape and flow regime remain part of the model. \\
\phantomsection\label{entry:symbol-registry-02324}\(E,A,B,p,q,v_*,v_E,v_P\) & \texttt{\detokenize{E,}\allowbreak\detokenize{A,}\allowbreak\detokenize{B,}\allowbreak\detokenize{p,}\allowbreak\detokenize{q,}\allowbreak\detokenize{v_*,}\allowbreak\detokenize{v_E,}\allowbreak\detokenize{v_P}} & two-power objective coefficients exponents and minimizing speeds & \hyperref[sec:p15-dimensional-optimization]{\ref*{sec:p15-dimensional-optimization}} & n/a & \hyperref[thm:p15-dimensional-competing-powers]{\ref*{thm:p15-dimensional-competing-powers}} & \thepage & Energy per distance and power are distinct objectives. \\
\phantomsection\label{entry:symbol-registry-02325}\(C,C',c_{\rm d},L,A_{\rm wing},A_{\rm cs},\beta,\mathcal E_{\rm fuel}\) & \texttt{\detokenize{C,}\allowbreak\detokenize{C',}\allowbreak\detokenize{c_{\rm d},}\allowbreak\detokenize{L,}\allowbreak\detokenize{A_{\rm wing},}\allowbreak\detokenize{A_{\rm cs},}\allowbreak\detokenize{\beta,}\allowbreak\detokenize{\mathcal E_{\rm fuel}}} & drag conventions span reference areas fuel fraction and usable energy per mass & \hyperref[sec:p15-dimensional-optimization]{\ref*{sec:p15-dimensional-optimization}} & n/a & \hyperref[sec:p15-dimensional-optimization]{\ref*{sec:p15-dimensional-optimization}} & \thepage & Area and one-half conventions are preserved. \\
\phantomsection\label{entry:symbol-registry-02326}\(P_0,M,n,k,N_b,[z^n],A_b(t)\) & \texttt{\detokenize{P_0,}\allowbreak\detokenize{M,}\allowbreak\detokenize{n,}\allowbreak\detokenize{k,}\allowbreak\detokenize{N_b,}\allowbreak\detokenize{[z^n],}\allowbreak\detokenize{A_b(t)}} & noncollision probability category/person counts occupancy threshold coefficient extraction and assignment counts & \hyperref[sec:p15-dimensional-birthday]{\ref*{sec:p15-dimensional-birthday}} & n/a & \hyperref[thm:p15-dimensional-birthday]{\ref*{thm:p15-dimensional-birthday}} & \thepage & Uniform independent birthday assignments do not make pair events jointly independent. \\
\phantomsection\label{entry:symbol-registry-02327}\(Q,l_i,u_i,\widehat x_i,\delta_i,\eta_i\) & \texttt{\detokenize{Q,}\allowbreak\detokenize{l_i,}\allowbreak\detokenize{u_i,}\allowbreak\detokenize{\widehat x_i,}\allowbreak\detokenize{\delta_i,}\allowbreak\detokenize{\eta_i}} & product estimate bounds nominal values logarithmic errors and error bounds & \hyperref[sec:p15-dimensional-estimates]{\ref*{sec:p15-dimensional-estimates}} & n/a & \hyperref[thm:p15-dimensional-product-bounds]{\ref*{thm:p15-dimensional-product-bounds}} & \thepage & A rectangular range is not a probability model. \\
\phantomsection\label{entry:symbol-registry-02328}\(C(n),M(n),h,B\) & \texttt{\detokenize{C(n),}\allowbreak\detokenize{M(n),}\allowbreak\detokenize{h,}\allowbreak\detokenize{B}} & abstract recursive costs depth and numeral block base & \hyperref[sec:p15-dimensional-source-tasks]{\ref*{sec:p15-dimensional-source-tasks}} & n/a & \hyperref[thm:p15-dimensional-source-extensions]{\ref*{thm:p15-dimensional-source-extensions}} & \thepage & Concrete implementations require their own operation accounting. \\
\phantomsection\label{entry:symbol-registry-02329}\(C,m,x_0,I,r,f\) & \texttt{\detokenize{C,}\allowbreak\detokenize{m,}\allowbreak\detokenize{x_0,}\allowbreak\detokenize{I,}\allowbreak\detokenize{r,}\allowbreak\detokenize{f}} & quartic coefficient mass amplitude inertia rolling radius and friction & \hyperref[sec:p15-dimensional-source-tasks]{\ref*{sec:p15-dimensional-source-tasks}} & n/a & \hyperref[thm:p15-dimensional-source-extensions]{\ref*{thm:p15-dimensional-source-extensions}} & \thepage & Quartic motion uses a restoring equation; rolling assumes no slip. \\
\phantomsection\label{entry:symbol-registry-02330}\(\epsilon,\sigma,A,F,L,T\) & \texttt{\detokenize{\epsilon,}\allowbreak\detokenize{\sigma,}\allowbreak\detokenize{A,}\allowbreak\detokenize{F,}\allowbreak\detokenize{L,}\allowbreak\detokenize{T}} & emissivity radiation coefficient albedo flux luminosity and absolute temperature & \hyperref[sec:p15-dimensional-mechanical]{\ref*{sec:p15-dimensional-mechanical}} & n/a & \hyperref[sec:p15-dimensional-mechanical]{\ref*{sec:p15-dimensional-mechanical}} & \thepage & The radiative-balance relation is a declared physical model. \\
\phantomsection\label{entry:symbol-registry-02331}\(\kappa_i,d_i,A_i\) & \texttt{\detokenize{\kappa_i,}\allowbreak\detokenize{d_i,}\allowbreak\detokenize{A_i}} & conductivity thickness and area in a thermal-resistance comparison & \hyperref[sec:p15-dimensional-source-tasks]{\ref*{sec:p15-dimensional-source-tasks}} & n/a & \hyperref[sec:p15-dimensional-source-tasks]{\ref*{sec:p15-dimensional-source-tasks}} & \thepage & Conductivity ratios cannot be omitted when comparing unlike materials. \\
\phantomsection\label{entry:symbol-registry-02332}\(\lambda_2\) & \texttt{\detokenize{\lambda_2}} & planar Lebesgue measure & \hyperref[annex:p15.dimensional.mechanical-models]{\ref*{annex:p15.dimensional.mechanical-models}} & n/a & \hyperref[annex:p15.dimensional.mechanical-models]{\ref*{annex:p15.dimensional.mechanical-models}} & \thepage & Area scaling is stated on Borel sets of positive finite area. \\
\phantomsection\label{entry:symbol-registry-02333}\(\mathcal X\) & \texttt{\detokenize{\mathcal X}} & declared feasible input domain & P15 global-sensitivity-and-robustness & n/a & \hyperref[def:p15-global-model]{\ref*{def:p15-global-model}} & \thepage & May differ from the support under one selected joint probability law. \\
\phantomsection\label{entry:symbol-registry-02334}\(P\) & \texttt{\detokenize{P}} & joint model-input probability law & P15 global-sensitivity-and-robustness & n/a & \hyperref[def:p15-global-model]{\ref*{def:p15-global-model}} & \thepage & Dependence and ranges belong to the sensitivity claim. \\
\phantomsection\label{entry:symbol-registry-02335}\(b\) & \texttt{\detokenize{b}} & fixed baseline input vector & P15 global-sensitivity-and-robustness OAT sections & n/a & \hyperref[def:p15-global-model]{\ref*{def:p15-global-model}} & \thepage & Regression b below is a separately typed coefficient vector. \\
\phantomsection\label{entry:symbol-registry-02336}\(A_b\) & \texttt{\detokenize{A_b}} & union of fixed-baseline coordinate-line designs & P15 global-sensitivity-and-robustness & n/a & \hyperref[def:p15-global-oat]{\ref*{def:p15-global-oat}} & \thepage & All but at most one coordinate equals its baseline value. \\
\phantomsection\label{entry:symbol-registry-02337}\(r_d\) & \texttt{\detokenize{r_d}} & inscribed-ball to unit-cube volume ratio & P15 global-sensitivity-and-robustness & n/a & \hyperref[thm:p15-global-ball-bound]{\ref*{thm:p15-global-ball-bound}} & \thepage & Containing-region ratio; not actual sample volume or a method accuracy probability. \\
\phantomsection\label{entry:symbol-registry-02338}\(\chi_u\) & \texttt{\detokenize{\chi_u}} & corner contrast character & P15 global-sensitivity-and-robustness & n/a & \hyperref[def:p15-global-factorial]{\ref*{def:p15-global-factorial}} & \thepage & Product of coded signs indexed by u. \\
\phantomsection\label{entry:symbol-registry-02339}\(\beta_u\) & \texttt{\detokenize{\beta_u}} & coefficient in the complete two-level corner expansion & P15 global-sensitivity-and-robustness & n/a & \hyperref[def:p15-global-factorial]{\ref*{def:p15-global-factorial}} & \thepage & Corner effects use the uniform design measure; not automatically the model-input law. \\
\phantomsection\label{entry:symbol-registry-02340}\(Z_i,W\) & \texttt{\detokenize{Z_i,}\allowbreak\detokenize{W}} & standardized predictors and standardized scalar response & P15 global-sensitivity-and-robustness & n/a & \hyperref[def:p15-global-regression]{\ref*{def:p15-global-regression}} & \thepage & Positive variances required. \\
\phantomsection\label{entry:symbol-registry-02341}\(\Sigma\) & \texttt{\detokenize{\Sigma}} & covariance matrix of standardized predictors & P15 global-sensitivity-and-robustness & n/a & \hyperref[def:p15-global-regression]{\ref*{def:p15-global-regression}} & \thepage & Population or empirical matrix as explicitly specified. \\
\phantomsection\label{entry:symbol-registry-02342}\(b\) & \texttt{\detokenize{b}} & least-squares coefficient vector & P15 global-sensitivity-and-robustness regression section & n/a & \hyperref[def:p15-global-regression]{\ref*{def:p15-global-regression}} & \thepage & Not the OAT baseline; uniqueness requires a nonsingular covariance matrix. \\
\phantomsection\label{entry:symbol-registry-02343}\(R^2\) & \texttt{\detokenize{R^2}} & fraction explained by the specified linear projection & P15 global-sensitivity-and-robustness & n/a & \hyperref[def:p15-global-regression]{\ref*{def:p15-global-regression}} & \thepage & Not a universal scalar measure of nonlinearity. \\
\phantomsection\label{entry:symbol-registry-02344}\(\operatorname{EE}_i\) & \texttt{\detokenize{\operatorname{EE}}} & finite elementary effect in specified design coordinates & P15 global-sensitivity-and-robustness & n/a & \hyperref[def:p15-global-elementary-effects]{\ref*{def:p15-global-elementary-effects}} & \thepage & Signed nonzero step and feasible endpoint required. \\
\phantomsection\label{entry:symbol-registry-02345}\(\overline e_i,\overline{|e_i|},s_{e_i}^2\) & \texttt{\detokenize{\overline}} & signed mean absolute-effect mean and sample spread & P15 global-sensitivity-and-robustness & n/a & \hyperref[def:p15-global-elementary-effects]{\ref*{def:p15-global-elementary-effects}} & \thepage & Descriptive summaries of the actual trajectory design. \\
\phantomsection\label{entry:symbol-registry-02346}\(I\) & \texttt{\detokenize{I}} & finite set of factor indices & P15 global-sensitivity-and-robustness variance sections & n/a & \hyperref[def:p15-global-variance-setting]{\ref*{def:p15-global-variance-setting}} & \thepage & The identity matrix I in the regression subsection has its separately typed matrix role. \\
\phantomsection\label{entry:symbol-registry-02347}\(\mathcal G_u\) & \texttt{\detokenize{\mathcal G_u}} & information from coordinates in u & P15 global-sensitivity-and-robustness & n/a & \hyperref[def:p15-global-variance-setting]{\ref*{def:p15-global-variance-setting}} & \thepage & Sigma-algebra generated by the indicated real input coordinates. \\
\phantomsection\label{entry:symbol-registry-02348}\(P_u\) & \texttt{\detokenize{P_u}} & conditional expectation projection onto coordinate information & P15 global-sensitivity-and-robustness & n/a & \hyperref[def:p15-global-variance-setting]{\ref*{def:p15-global-variance-setting}} & \thepage & Distinct from the joint probability law P; coordinate projections commute under the product law. \\
\phantomsection\label{entry:symbol-registry-02349}\(Y_u\) & \texttt{\detokenize{Y_u}} & orthogonal functional component for a coordinate subset & P15 global-sensitivity-and-robustness & n/a & \hyperref[def:p15-global-variance-setting]{\ref*{def:p15-global-variance-setting}} & \thepage & Defined by finite subset inversion; independence governs the orthogonal decomposition. \\
\phantomsection\label{entry:symbol-registry-02350}\(V,V_u\) & \texttt{\detokenize{V,}\allowbreak\detokenize{V_u}} & total output variance and component variances & P15 global-sensitivity-and-robustness & n/a & \hyperref[def:p15-global-variance-setting]{\ref*{def:p15-global-variance-setting}} & \thepage & The normalization V must be strictly positive. \\
\phantomsection\label{entry:symbol-registry-02351}\(S_i,S_{T_i}\) & \texttt{\detokenize{S_i,}\allowbreak\detokenize{S_{T_i}}} & first-order and total variance-based sensitivity indices & P15 global-sensitivity-and-robustness & n/a & \hyperref[def:p15-global-indices]{\ref*{def:p15-global-indices}} & \thepage & Conditional definitions remain meaningful under dependence but independent-input identities may fail. \\
\phantomsection\label{entry:symbol-registry-02352}\(A,B,C_i\) & \texttt{\detokenize{A,}\allowbreak\detokenize{B,}\allowbreak\detokenize{C_i}} & independent input copies and hybrid coordinate vector & P15 global-sensitivity-and-robustness & n/a & \hyperref[thm:p15-global-pick-freeze]{\ref*{thm:p15-global-pick-freeze}} & \thepage & \(C_i\) uses coordinate \(i\) from \(A\) and the other coordinates from \(B\). \\
\phantomsection\label{entry:symbol-registry-02353}\(\widehat V,\widehat S_i,\widehat S_{T_i}\) & \texttt{\detokenize{\widehat}} & sample variance and paired sensitivity estimates & P15 global-sensitivity-and-robustness & n/a & \hyperref[thm:p15-global-pick-freeze]{\ref*{thm:p15-global-pick-freeze}} & \thepage & Ratios used only with positive estimated variance; consistency does not imply finite-sample unbiasedness. \\
\phantomsection\label{entry:symbol-registry-02354}\(G,D\) & \texttt{\detokenize{G,}\allowbreak\detokenize{D}} & the two explicitly defined Saltelli-Annoni toy responses & P15 global-sensitivity-and-robustness source examples & n/a & \hyperref[thm:p15-global-source-toys]{\ref*{thm:p15-global-source-toys}} & \thepage & G uses the printed source experiment's uniform [-1,1] law. \\
\phantomsection\label{entry:symbol-registry-02355}\(m_i,v_i\) & \texttt{\detokenize{m_i,}\allowbreak\detokenize{v_i}} & mean and variance of a product-response factor & P15 global-sensitivity-and-robustness source examples & n/a & \hyperref[thm:p15-global-source-toys]{\ref*{thm:p15-global-source-toys}} & \thepage & The final ranking example separately uses \(v_i\) as input variance. \\
\phantomsection\label{entry:symbol-registry-02356}\(\lambda_i,N_i\) & \texttt{\detokenize{\lambda_i,}\allowbreak\detokenize{N_i}} & rate and concentration in the reported source application & P15 global-sensitivity-and-robustness source examples & n/a & \hyperref[sec:p15-global-rate-chain]{\ref*{sec:p15-global-rate-chain}} & \thepage & Evidence-only application record; rates have reciprocal-time units. \\
\phantomsection\label{entry:symbol-registry-02357}\(U_1,U_2\) & \texttt{\detokenize{U_1,}\allowbreak\detokenize{U_2}} & independent exogenous binary variables in the causal counterexample & P15 global-sensitivity-and-robustness causal boundary & n/a & \hyperref[annex:p15.global.causal-boundary]{\ref*{annex:p15.global.causal-boundary}} & \thepage & Distinct from output variance V and the standardized regression response W. \\
\phantomsection\label{entry:symbol-registry-02358}\(X\in\mathbb R^{N\times p}\) & \texttt{\detokenize{X}} & analysis matrix & P15 factorial inference & n/a & \hyperref[def:p15-factorial-linear-model]{\ref*{def:p15-factorial-linear-model}} & \thepage & Includes every fitted intercept treatment interaction and block column; distinct from a raw setting table. \\
\phantomsection\label{entry:symbol-registry-02359}\(\beta\) & \texttt{\detokenize{\beta}} & coefficient vector & P15 factorial inference & n/a & \hyperref[def:p15-factorial-linear-model]{\ref*{def:p15-factorial-linear-model}} & \thepage & Parameterization and admitted response model must be specified. \\
\phantomsection\label{entry:symbol-registry-02360}\(c^{\mathsf T}\beta\) & \texttt{\detokenize{c}} & coefficient-space contrast weights & P15 factorial estimability & n/a & \hyperref[def:p15-factorial-linear-model]{\ref*{def:p15-factorial-linear-model}} & \thepage & Estimable exactly when c belongs to the row space of X. \\
\phantomsection\label{entry:symbol-registry-02361}\(a^{\mathsf T}Y\) & \texttt{\detokenize{a}} & observation-space linear estimator weights & P15 factorial estimability & n/a & \hyperref[def:p15-factorial-linear-model]{\ref*{def:p15-factorial-linear-model}} & \thepage & Unbiasedness for all coefficients is X transpose a equals c. \\
\phantomsection\label{entry:symbol-registry-02362}\(\epsilon\) & \texttt{\detokenize{\epsilon}} & mean-zero response error vector & P15 factorial inference & n/a & \hyperref[def:p15-factorial-linear-model]{\ref*{def:p15-factorial-linear-model}} & \thepage & Common variance and Gaussian conditions are separate assumptions. \\
\phantomsection\label{entry:symbol-registry-02363}\(\mu_{ab}\) & \texttt{\detokenize{\mu_{ab}}} & four factorial cell means & P15 factorial two-by-two & n/a & \hyperref[def:p15-factorial-coding]{\ref*{def:p15-factorial-coding}} & \thepage & First binary coordinate is checking instruction; second is numerical aid. \\
\phantomsection\label{entry:symbol-registry-02364}\(\Delta_A(b)\) & \texttt{\detokenize{\Delta_A}} & conditional checking difference & P15 factorial two-by-two & n/a & \hyperref[def:p15-factorial-coding]{\ref*{def:p15-factorial-coding}} & \thepage & Equal-weight average is a separately defined main effect. \\
\phantomsection\label{entry:symbol-registry-02365}\(D_{AB}\) & \texttt{\detokenize{D_{AB}}} & difference of differences & P15 factorial two-by-two & n/a & \hyperref[def:p15-factorial-coding]{\ref*{def:p15-factorial-coding}} & \thepage & Factorial interaction effect is half this quantity; coded regression coefficient is one quarter. \\
\phantomsection\label{entry:symbol-registry-02366}\(H_2\) & \texttt{\detokenize{H_2}} & four-row saturated two-factor analysis matrix & P15 factorial two-by-two & n/a & \hyperref[thm:p15-factorial-two-by-two]{\ref*{thm:p15-factorial-two-by-two}} & \thepage & Cell order00 10 01 11; columns intercept x z xz. \\
\phantomsection\label{entry:symbol-registry-02367}\(a_*\) & \texttt{\detokenize{a_*}} & minimum-norm unbiased contrast weights & P15 factorial uncertainty & n/a & \hyperref[thm:p15-factorial-linear-uncertainty]{\ref*{thm:p15-factorial-linear-uncertainty}} & \thepage & Unique solution in the column space of X. \\
\phantomsection\label{entry:symbol-registry-02368}\(P_X\) & \texttt{\detokenize{P_X}} & orthogonal projection onto the fitted-mean space & P15 factorial uncertainty & n/a & \hyperref[thm:p15-factorial-linear-uncertainty]{\ref*{thm:p15-factorial-linear-uncertainty}} & \thepage & Projection rank is rank X. \\
\phantomsection\label{entry:symbol-registry-02369}\(\mathrm{SSE}\) & \texttt{\detokenize{\mathrm{SSE}}} & residual sum of squares & P15 factorial uncertainty & n/a & \hyperref[thm:p15-factorial-linear-uncertainty]{\ref*{thm:p15-factorial-linear-uncertainty}} & \thepage & Includes model lack of fit unless a correct mean model is assumed. \\
\phantomsection\label{entry:symbol-registry-02370}\(s^2\) & \texttt{\detokenize{s^2}} & residual variance estimate & P15 factorial uncertainty & n/a & \hyperref[thm:p15-factorial-linear-uncertainty]{\ref*{thm:p15-factorial-linear-uncertainty}} & \thepage & Defined only when N exceeds rank X. \\
\phantomsection\label{entry:symbol-registry-02371}\(\eta\) & \texttt{\detokenize{\eta}} & noncoverage probability & P15 factorial uncertainty & n/a & \hyperref[thm:p15-factorial-linear-uncertainty]{\ref*{thm:p15-factorial-linear-uncertainty}} & \thepage & The displayed normal-model interval has coverage1 minus eta. \\
\phantomsection\label{entry:symbol-registry-02372}\(F_{q,\nu}\) & \texttt{\detokenize{F}} & ratio law for a joint nested-model test & P15 factorial uncertainty & n/a & \hyperref[thm:p15-factorial-joint-test]{\ref*{thm:p15-factorial-joint-test}} & \thepage & Ratio of independent chi-squared variables divided by their positive degrees of freedom. \\
\phantomsection\label{entry:symbol-registry-02373}\(\mathrm{SS}_{\rm pure}\) & \texttt{\detokenize{\mathrm{SS}_{\rm pure}}} & within-setting pure-error sum of squares & P15 factorial replication & n/a & \hyperref[thm:p15-factorial-replication]{\ref*{thm:p15-factorial-replication}} & \thepage & Independent replicates are distinct experimental units. \\
\phantomsection\label{entry:symbol-registry-02374}\(U_i,e_{ij}\) & \texttt{\detokenize{U_i;e_{ij}}} & unit and repeated-measurement error components & P15 factorial replication & n/a & \hyperref[thm:p15-factorial-replication]{\ref*{thm:p15-factorial-replication}} & \thepage & The shared unit variance is not divided by the number of repeated measurements. \\
\phantomsection\label{entry:symbol-registry-02375}\(\tau_t,\gamma_b\) & \texttt{\detokenize{\tau_t;\gamma_b}} & centered treatment and block effects & P15 factorial blocking & n/a & \hyperref[sec:p15-factorial-randomized-blocks]{\ref*{sec:p15-factorial-randomized-blocks}} & \thepage & Additive model excludes arbitrary treatment-by-block interactions. \\
\phantomsection\label{entry:symbol-registry-02376}\(s,k\) & \texttt{\detokenize{s;k}} & levels and total number of factors in an orthogonal square & P15 factorial square designs & n/a & \hyperref[def:p15-factorial-orthogonal-squares]{\ref*{def:p15-factorial-orthogonal-squares}} & \thepage & The local s is a level count rather than a standard deviation. \\
\phantomsection\label{entry:symbol-registry-02377}\(\theta_j(a)\) & \texttt{\detokenize{\theta_j}} & centered effect of one level of a square factor & P15 factorial square designs & n/a & \hyperref[def:p15-factorial-orthogonal-squares]{\ref*{def:p15-factorial-orthogonal-squares}} & \thepage & Each level function sums to zero over its s levels. \\
\phantomsection\label{entry:symbol-registry-02378}\(\chi_S(x)\) & \texttt{\detokenize{\chi_S}} & product column indexed by a factor subset & P15 factorial full and fractional designs & n/a & \hyperref[def:p15-factorial-walsh]{\ref*{def:p15-factorial-walsh}} & \thepage & Empty subset gives the intercept; multiplication uses symmetric difference. \\
\phantomsection\label{entry:symbol-registry-02379}\(I\) & \texttt{\detokenize{I}} & all-ones intercept column or empty factor word & P15 factorial word notation & n/a & \hyperref[def:p15-factorial-walsh]{\ref*{def:p15-factorial-walsh}} & \thepage & Not the letter labeling a physical factor; letter factor notation omits I. \\
\phantomsection\label{entry:symbol-registry-02380}\(\mathcal G\) & \texttt{\detokenize{\mathcal G}} & defining word group of a regular fraction & P15 factorial fractions & n/a & \hyperref[def:p15-factorial-fraction]{\ref*{def:p15-factorial-fraction}} & \thepage & Generated by p independent binary words; has2 raised to p members. \\
\phantomsection\label{entry:symbol-registry-02381}\(s(H)\) & \texttt{\detokenize{s(H)}} & sign attached to a defining word & P15 factorial fractions & n/a & \hyperref[def:p15-factorial-fraction]{\ref*{def:p15-factorial-fraction}} & \thepage & Product of the selected generator signs; distinct from residual standard deviation. \\
\phantomsection\label{entry:symbol-registry-02382}\(S\mathbin{\triangle}T\) & \texttt{\detokenize{\triangle}} & symmetric difference of factor words & P15 factorial fractions & infix & \hyperref[def:p15-factorial-fraction]{\ref*{def:p15-factorial-fraction}} & \thepage & Corresponds to multiplication of sign columns. \\
\phantomsection\label{entry:symbol-registry-02383}\(\pi_i(t),\pi_{ij}(t,u)\) & \texttt{\detokenize{\pi_i;\pi_{ij}}} & marginal and joint assignment probabilities & P15 factorial causal interpretation & n/a & \hyperref[def:p15-factorial-assignment]{\ref*{def:p15-factorial-assignment}} & \thepage & Use the actual restricted or complete randomization distribution. \\
\phantomsection\label{entry:symbol-registry-02384}\(Y_i(t)\) & \texttt{\detokenize{Y_i}} & potential response of a recruited unit & P15 factorial causal interpretation & n/a & \hyperref[def:p15-factorial-assignment]{\ref*{def:p15-factorial-assignment}} & \thepage & Stable treatment versions consistency and no relevant interference are explicit. \\
\phantomsection\label{entry:symbol-registry-02385}\(Z_i\) & \texttt{\detokenize{Z_i}} & randomized assigned treatment cell & P15 factorial causal interpretation & n/a & \hyperref[def:p15-factorial-assignment]{\ref*{def:p15-factorial-assignment}} & \thepage & Assignment randomization is distinct from population sampling. \\
\phantomsection\label{entry:symbol-registry-02386}\(\tau,\widehat\tau\) & \texttt{\detokenize{\tau;\widehat\tau}} & finite assignment target and inverse-probability estimator & P15 factorial causal interpretation & n/a & \hyperref[thm:p15-factorial-assignment]{\ref*{thm:p15-factorial-assignment}} & \thepage & Target weights and every positive marginal assignment probability are specified. \\
\phantomsection\label{entry:symbol-registry-02400}\(\approx_{\mathrm{est}}\) & \texttt{\detokenize{\approx_{\mathrm{est}}}} & heuristic estimate without an asserted numerical error bound & \hyperref[sec:p15-approximate-quantitative-inference]{\ref*{sec:p15-approximate-quantitative-inference}} & relation & \hyperref[def:p15-approximation-relations]{\ref*{def:p15-approximation-relations}} & \thepage & Separates the source's comparable-scale notation from asymptotic equivalence. \\
\phantomsection\label{entry:symbol-registry-02401}\(\sim\) & \texttt{\detokenize{\sim}} & ratio tending to one in a declared limit & \hyperref[sec:p15-approximate-quantitative-inference]{\ref*{sec:p15-approximate-quantitative-inference}} & relation & \hyperref[def:p15-approximation-relations]{\ref*{def:p15-approximation-relations}} & \thepage & Comparison denominator is eventually nonzero; held-fixed variables are specified. \\
\phantomsection\label{entry:symbol-registry-02402}\(O\) & \texttt{\detokenize{O}} & bounded asymptotic remainder order & \hyperref[sec:p15-approximate-quantitative-inference]{\ref*{sec:p15-approximate-quantitative-inference}} & meta & \hyperref[def:p15-approximation-relations]{\ref*{def:p15-approximation-relations}} & \thepage & Existence of constants is distinct from a supplied numerical bound. \\
\phantomsection\label{entry:symbol-registry-02403}\(L,U\) & \texttt{\detokenize{L,}\allowbreak\detokenize{U}} & lower and upper bounds for an exact target & \hyperref[thm:p15-approximation-certificates]{\ref*{thm:p15-approximation-certificates}} & meta & \hyperref[thm:p15-approximation-certificates]{\ref*{thm:p15-approximation-certificates}} & \thepage & Not strict bounds unless expressly stated. \\
\phantomsection\label{entry:symbol-registry-02404}\(c\) & \texttt{\detokenize{c}} & decision threshold & \hyperref[thm:p15-approximation-certificates]{\ref*{thm:p15-approximation-certificates}} & meta & \hyperref[thm:p15-approximation-certificates]{\ref*{thm:p15-approximation-certificates}} & \thepage & Distinct from characteristic graph height and physical speed in later local models. \\
\phantomsection\label{entry:symbol-registry-02405}\(E_f,E_g\) & \texttt{\detokenize{E_f,}\allowbreak\detokenize{E_g}} & absolute error budgets & \hyperref[thm:p15-approximation-certificates]{\ref*{thm:p15-approximation-certificates}} & meta & \hyperref[thm:p15-approximation-certificates]{\ref*{thm:p15-approximation-certificates}} & \thepage & Bounds for two separate approximations. \\
\phantomsection\label{entry:symbol-registry-02406}\(Q_h\) & \texttt{\detokenize{Q_h}} & nearest rounding on a logarithmic grid & \hyperref[def:p15-approximation-log-round]{\ref*{def:p15-approximation-log-round}} & n/a & \hyperref[def:p15-approximation-log-round]{\ref*{def:p15-approximation-log-round}} & \thepage & Base \(b>1\) and exponent spacing \(h>0\); input positive. \\
\phantomsection\label{entry:symbol-registry-02407}\(b,h\) & \texttt{\detokenize{b,}\allowbreak\detokenize{h}} & logarithmic-grid base and exponent spacing & \hyperref[def:p15-approximation-log-round]{\ref*{def:p15-approximation-log-round}} & meta & \hyperref[def:p15-approximation-log-round]{\ref*{def:p15-approximation-log-round}} & \thepage & Both are dimensionless parameters; \(b>1\) and \(h>0\). \\
\phantomsection\label{entry:symbol-registry-02408}\(w_i\) & \texttt{\detokenize{w_i}} & nonnegative averaging weights summing to one & \hyperref[thm:p15-approximation-lumping]{\ref*{thm:p15-approximation-lumping}} & meta & \hyperref[thm:p15-approximation-lumping]{\ref*{thm:p15-approximation-lumping}} & \thepage & Finite weighted population, not assumed random sampling. \\
\phantomsection\label{entry:symbol-registry-02409}\(\mu\) & \texttt{\detokenize{\mu}} & weighted input mean & \hyperref[thm:p15-approximation-lumping]{\ref*{thm:p15-approximation-lumping}} & meta & \hyperref[thm:p15-approximation-lumping]{\ref*{thm:p15-approximation-lumping}} & \thepage & Mean of inputs, not automatically mean of nonlinear outputs. \\
\phantomsection\label{entry:symbol-registry-02410}\(M\) & \texttt{\detokenize{M}} & uniform bound on the magnitude of a second derivative & \hyperref[thm:p15-approximation-lumping]{\ref*{thm:p15-approximation-lumping}} & meta & \hyperref[thm:p15-approximation-lumping]{\ref*{thm:p15-approximation-lumping}} & \thepage & Nonlinear lumping bound. \\
\phantomsection\label{entry:symbol-registry-02411}\(M_0\) & \texttt{\detokenize{M_0}} & total mass of the velocity slabs & \hyperref[thm:p15-approximation-energy-average]{\ref*{thm:p15-approximation-energy-average}} & meta & \hyperref[thm:p15-approximation-energy-average]{\ref*{thm:p15-approximation-energy-average}} & \thepage & Sum of positive component masses. \\
\phantomsection\label{entry:symbol-registry-02412}\(\overline v\) & \texttt{\detokenize{\overline v}} & momentum-preserving common velocity & \hyperref[thm:p15-approximation-energy-average]{\ref*{thm:p15-approximation-energy-average}} & meta & \hyperref[thm:p15-approximation-energy-average]{\ref*{thm:p15-approximation-energy-average}} & \thepage & Mass-weighted velocity. \\
\phantomsection\label{entry:symbol-registry-02413}\(\alpha,\eta,e\) & \texttt{\detokenize{\alpha,}\allowbreak\detokenize{\eta,}\allowbreak\detokenize{e}} & muscle fraction, conversion fraction and muscle-specific energy & \hyperref[def:p15-approximation-jumps]{\ref*{def:p15-approximation-jumps}} & meta & \hyperref[def:p15-approximation-jumps]{\ref*{def:p15-approximation-jumps}} & \thepage & Local jump-model parameters; no equality across species is asserted. \\
\phantomsection\label{entry:symbol-registry-02414}\(k\) & \texttt{\detokenize{k}} & positive restoring-force coefficient & \hyperref[thm:p15-approximation-oscillators]{\ref*{thm:p15-approximation-oscillators}} & meta & \hyperref[thm:p15-approximation-oscillators]{\ref*{thm:p15-approximation-oscillators}} & \thepage & Local oscillator role; distinct from later angular wavenumber. \\
\phantomsection\label{entry:symbol-registry-02415}\(n\) & \texttt{\detokenize{n}} & positive restoring-force exponent & \hyperref[thm:p15-approximation-oscillators]{\ref*{thm:p15-approximation-oscillators}} & meta & \hyperref[thm:p15-approximation-oscillators]{\ref*{thm:p15-approximation-oscillators}} & \thepage & Force is \(-k|x|^{n-1}x\), defined as zero at \(x=0\); \(n\) need not be an integer. \\
\phantomsection\label{entry:symbol-registry-02416}\(V\) & \texttt{\detokenize{V}} & terminal speed in a quadratic-drag model & \hyperref[thm:p15-approximation-fall]{\ref*{thm:p15-approximation-fall}} & meta & \hyperref[thm:p15-approximation-fall]{\ref*{thm:p15-approximation-fall}} & \thepage & Positive constant; distinct from potential-energy notation in its annex proof. \\
\phantomsection\label{entry:symbol-registry-02417}\(I\) & \texttt{\detokenize{I}} & moment of inertia about the rolling axis & \hyperref[thm:p15-approximation-mechanical-cases]{\ref*{thm:p15-approximation-mechanical-cases}} & meta & \hyperref[thm:p15-approximation-mechanical-cases]{\ref*{thm:p15-approximation-mechanical-cases}} & \thepage & The rolling radius can differ from the largest geometric radius. \\
\phantomsection\label{entry:symbol-registry-02418}\(\mathrm{Re}\) & \texttt{\detokenize{\mathrm{Re}}} & diameter-based Reynolds number & \hyperref[def:p15-approximation-drag]{\ref*{def:p15-approximation-drag}} & meta & \hyperref[def:p15-approximation-drag]{\ref*{def:p15-approximation-drag}} & \thepage & Equals \(2rv/\nu\) for the sphere examples. \\
\phantomsection\label{entry:symbol-registry-02419}\(C_d\) & \texttt{\detokenize{C_d}} & drag coefficient & \hyperref[def:p15-approximation-drag]{\ref*{def:p15-approximation-drag}} & meta & \hyperref[def:p15-approximation-drag]{\ref*{def:p15-approximation-drag}} & \thepage & Defined with \(\tfrac12\rho\pi r^2v^2\) as denominator. \\
\phantomsection\label{entry:symbol-registry-02420}\(r_0,E_0\) & \texttt{\detokenize{r_0,}\allowbreak\detokenize{E_0}} & minimizing radius and minimum of the stipulated two-term energy & \hyperref[thm:p15-approximation-competing-energies]{\ref*{thm:p15-approximation-competing-energies}} & meta & \hyperref[thm:p15-approximation-competing-energies]{\ref*{thm:p15-approximation-competing-energies}} & \thepage & Exact values within \(E(r)=A/r^2-B/r\). \\
\phantomsection\label{entry:symbol-registry-02421}\(\Delta\) & \texttt{\detokenize{\Delta}} & energy above the stipulated potential minimum & \hyperref[thm:p15-approximation-competing-energies]{\ref*{thm:p15-approximation-competing-energies}} & meta & \hyperref[thm:p15-approximation-competing-energies]{\ref*{thm:p15-approximation-competing-energies}} & \thepage & Strictly between zero and the magnitude of the minimum energy. \\
\phantomsection\label{entry:symbol-registry-02422}\(\phi\) & \texttt{\detokenize{\phi}} & velocity potential of the linear surface-wave model & \hyperref[def:p15-approximation-water-model]{\ref*{def:p15-approximation-water-model}} & meta & \hyperref[def:p15-approximation-water-model]{\ref*{def:p15-approximation-water-model}} & \thepage & Incompressible irrotational flow in fixed depth. \\
\phantomsection\label{entry:symbol-registry-02423}\(\eta\) & \texttt{\detokenize{\eta}} & linear surface elevation & \hyperref[def:p15-approximation-water-model]{\ref*{def:p15-approximation-water-model}} & meta & \hyperref[def:p15-approximation-water-model]{\ref*{def:p15-approximation-water-model}} & \thepage & Distinct from the earlier jump conversion fraction. \\
\phantomsection\label{entry:symbol-registry-02424}\(\gamma\) & \texttt{\detokenize{\gamma}} & nonnegative surface tension & \hyperref[def:p15-approximation-water-model]{\ref*{def:p15-approximation-water-model}} & meta & \hyperref[def:p15-approximation-water-model]{\ref*{def:p15-approximation-water-model}} & \thepage & Capillary-dominant ratios require \(\gamma>0\). \\
\phantomsection\label{entry:symbol-registry-02425}\(k,\ell\) & \texttt{\detokenize{k,}\allowbreak\detokenize{\ell}} & angular wavenumber and reduced wavelength & \hyperref[thm:p15-approximation-water-regimes]{\ref*{thm:p15-approximation-water-regimes}} & meta & \hyperref[thm:p15-approximation-water-regimes]{\ref*{thm:p15-approximation-water-regimes}} & \thepage & \(k=2\pi/\lambda\) and \(\ell=1/k\). \\
\phantomsection\label{entry:symbol-registry-02426}\(\delta,\epsilon\) & \texttt{\detokenize{\delta,}\allowbreak\detokenize{\epsilon}} & restoring-term and depth-factor error budgets & \hyperref[thm:p15-approximation-water-regimes]{\ref*{thm:p15-approximation-water-regimes}} & meta & \hyperref[thm:p15-approximation-water-regimes]{\ref*{thm:p15-approximation-water-regimes}} & \thepage & Nonnegative bounds for the selected corner approximation. \\
\phantomsection\label{entry:symbol-registry-02427}\(\varepsilon\) & \texttt{\detokenize{\varepsilon}} & small positive parameter in the nonuniform-limit example & \hyperref[thm:p15-approximation-nonuniform]{\ref*{thm:p15-approximation-nonuniform}} & meta & \hyperref[thm:p15-approximation-nonuniform]{\ref*{thm:p15-approximation-nonuniform}} & \thepage & \(0<\varepsilon\leq1\); the moving input \(t=\varepsilon\) remains in the domain. \\
\phantomsection\label{entry:symbol-registry-02702}\(H_i\) & \texttt{\detokenize{H_i}} & indexed basic null hypothesis & P15 multiple testing & n/a & \hyperref[def:p15-mt-family]{\ref*{def:p15-mt-family}} & \thepage & Asserts membership in its specified null parameter set; the family is fixed before observing data. \\
\phantomsection\label{entry:symbol-registry-02703}\(\omega_i\) & \texttt{\detokenize{\omega_i}} & null parameter set for the indexed hypothesis & P15 multiple testing & n/a & \hyperref[def:p15-mt-family]{\ref*{def:p15-mt-family}} & \thepage & May overlap other null sets and include nuisance parameters. \\
\phantomsection\label{entry:symbol-registry-02704}\(n\) & \texttt{\detokenize{n}} & number of hypotheses in the fixed family & P15 multiple testing & n/a & \hyperref[def:p15-mt-family]{\ref*{def:p15-mt-family}} & \thepage & Positive integer for the main rule; the empty-family convention returns no rejections. \\
\phantomsection\label{entry:symbol-registry-02705}\(I_0(\theta)\) & \texttt{\detokenize{I_0(\theta)}} & indices of hypotheses true at the actual full parameter & P15 multiple testing & n/a & \hyperref[def:p15-mt-family]{\ref*{def:p15-mt-family}} & \thepage & Not an observed input to the rejection procedure. \\
\phantomsection\label{entry:symbol-registry-02706}\(m\) & \texttt{\detokenize{m}} & number of actual true hypotheses & P15 multiple testing & n/a & \hyperref[def:p15-mt-family]{\ref*{def:p15-mt-family}} & \thepage & Zero gives an empty family-error event; otherwise used in the first-true-rejection proof. \\
\phantomsection\label{entry:symbol-registry-02707}\(\mathcal R(X)\) & \texttt{\detokenize{\mathcal R(X)}} & measurable basic rejection set & P15 multiple testing & n/a & \hyperref[def:p15-mt-family]{\ref*{def:p15-mt-family}} & \thepage & A subset of the fixed hypothesis indices depending on the observed data. \\
\phantomsection\label{entry:symbol-registry-02708}\(C_i\) & \texttt{\detokenize{C_i}} & critical region for rejection of the indexed hypothesis & P15 multiple testing & n/a & \hyperref[def:p15-mt-family]{\ref*{def:p15-mt-family}} & \thepage & An event in the sample space; distinct from the lower-case positive weight. \\
\phantomsection\label{entry:symbol-registry-02709}\(\operatorname{FWER}_\theta\) & \texttt{\detokenize{\operatorname{FWER}_\theta}} & probability of at least one rejection of a true hypothesis & P15 multiple testing & n/a & \hyperref[def:p15-mt-family]{\ref*{def:p15-mt-family}} & \thepage & Evaluated at each actual configuration rather than only the complete null. \\
\phantomsection\label{entry:symbol-registry-02710}\(\alpha\) & \texttt{\detokenize{\alpha}} & chosen family significance level & P15 multiple testing & n/a & \hyperref[def:p15-mt-family]{\ref*{def:p15-mt-family}} & \thepage & Between zero and one; distinct from the source critical-level function with subscript. \\
\phantomsection\label{entry:symbol-registry-02711}\(R_i\) & \texttt{\detokenize{R_i}} & marginally valid obtained level or p-value & P15 multiple testing & n/a & \hyperref[def:p15-mt-obtained-levels]{\ref*{def:p15-mt-obtained-levels}} & \thepage & The null inequality quantifies over the entire parameter including other hypotheses and nuisance values. \\
\phantomsection\label{entry:symbol-registry-02712}\(Y_i\) & \texttt{\detokenize{Y_i}} & test statistic with large values favoring rejection & P15 multiple testing & n/a & \hyperref[def:p15-mt-obtained-levels]{\ref*{def:p15-mt-obtained-levels}} & \thepage & Dependence across these statistics is unrestricted for ordinary and weighted Holm. \\
\phantomsection\label{entry:symbol-registry-02713}\(\alpha_i(y)\) & \texttt{\detokenize{\alpha_i(y)}} & inclusive upper-tail supremum over the indexed null & P15 multiple testing & n/a & \hyperref[def:p15-mt-obtained-levels]{\ref*{def:p15-mt-obtained-levels}} & \thepage & Nonincreasing measurable function used in Holm's original obtained-level construction. \\
\phantomsection\label{entry:symbol-registry-02714}\(R^{(j)}\) & \texttt{\detokenize{R^{(j)}}} & raw obtained level in the ordinary sorted order & P15 multiple testing & n/a & \hyperref[def:p15-mt-holm-rule]{\ref*{def:p15-mt-holm-rule}} & \thepage & Uses the unweighted obtained-level ordering; the weighted method instead tests the explicitly carried original index. \\
\phantomsection\label{entry:symbol-registry-02715}\(i_j\) & \texttt{\detokenize{i_j}} & original index carried by the currently specified sorting order & P15 multiple testing & n/a & \hyperref[def:p15-mt-holm-rule]{\ref*{def:p15-mt-holm-rule}} & \thepage & Ordinary ordering sorts obtained levels; the weighted method explicitly redefines its ordering by ratios. \\
\phantomsection\label{entry:symbol-registry-02716}\(H^{(j)}\) & \texttt{\detokenize{H^{(j)}}} & hypothesis carried by the current sorting position & P15 multiple testing & n/a & \hyperref[def:p15-mt-holm-rule]{\ref*{def:p15-mt-holm-rule}} & \thepage & The ordinary and weighted methods each specify their ordering before using the position. \\
\phantomsection\label{entry:symbol-registry-02717}\(\mathcal R_{\rm H}\) & \texttt{\detokenize{\mathcal R_{\rm H}}} & ordinary Holm rejection set & P15 multiple testing & n/a & \hyperref[def:p15-mt-holm-rule]{\ref*{def:p15-mt-holm-rule}} & \thepage & Only the passing prefix is rejected; stop at the first failed comparison. \\
\phantomsection\label{entry:symbol-registry-02718}\(\mathcal R_{\rm B},\mathcal R_{\rm U}\) & \texttt{\detokenize{\mathcal R_{\rm B};\mathcal R_{\rm U}}} & fixed Bonferroni and unadjusted rejection sets & P15 multiple testing & n/a & \hyperref[def:p15-mt-comparison-rules]{\ref*{def:p15-mt-comparison-rules}} & \thepage & Use the same family and individual obtained levels as the Holm comparison. \\
\phantomsection\label{entry:symbol-registry-02719}\(\widetilde R^{(r)}\) & \texttt{\detokenize{\widetilde R^{(r)}}} & Holm adjusted level assigned to a ranked hypothesis & P15 multiple testing & n/a & \hyperref[def:p15-mt-adjusted]{\ref*{def:p15-mt-adjusted}} & \thepage & The capped running maximum retains all previous stopping conditions. \\
\phantomsection\label{entry:symbol-registry-02720}\(H_I\) & \texttt{\detokenize{H_I}} & secondary intersection hypothesis & P15 multiple testing & n/a & \hyperref[def:p15-mt-intersection-policy]{\ref*{def:p15-mt-intersection-policy}} & \thepage & The stated secondary rejection policy rejects an intersection when any basic member is rejected. \\
\phantomsection\label{entry:symbol-registry-02721}\(\gamma,\gamma_0\) & \texttt{\detokenize{\gamma;\gamma_0}} & scalar directional parameter and common null boundary & P15 multiple testing & n/a & \hyperref[thm:p15-mt-directional]{\ref*{thm:p15-mt-directional}} & \thepage & The paired opposite-tail argument requires a boundary law present in both null sets. \\
\phantomsection\label{entry:symbol-registry-02722}\(\mathcal C(x)\) & \texttt{\detokenize{\mathcal C(x)}} & all-details-nonrejection inverted candidate set & P15 multiple testing & n/a & \hyperref[thm:p15-mt-confidence-identical]{\ref*{thm:p15-mt-confidence-identical}} & \thepage & The equality with Bonferroni concerns this specific confidence construction. \\
\phantomsection\label{entry:symbol-registry-02723}\(\Phi\) & \texttt{\detokenize{\Phi}} & standard normal cumulative distribution function & P15 multiple testing & n/a & \hyperref[def:p15-mt-normal-power]{\ref*{def:p15-mt-normal-power}} & \thepage & The exact normal-model power calculation uses this CDF and its quantiles. \\
\phantomsection\label{entry:symbol-registry-02724}\(b_k\) & \texttt{\detokenize{b_k}} & ordered critical-level bin boundary & P15 multiple testing & n/a & \hyperref[thm:p15-mt-normal-power]{\ref*{thm:p15-mt-normal-power}} & \thepage & Zero initial boundary followed by ten Holm thresholds for the source normal example. \\
\phantomsection\label{entry:symbol-registry-02725}\(F_\mu(u)\) & \texttt{\detokenize{F_\mu(u)}} & CDF of an obtained level under a specified normal mean & P15 multiple testing & n/a & \hyperref[thm:p15-mt-normal-power]{\ref*{thm:p15-mt-normal-power}} & \thepage & Defined with zero and one endpoint values in addition to its open-interval formula. \\
\phantomsection\label{entry:symbol-registry-02726}\(c_i\) & \texttt{\detokenize{c_i}} & positive fixed hypothesis-priority weight & P15 multiple testing & n/a & \hyperref[def:p15-mt-weighted]{\ref*{def:p15-mt-weighted}} & \thepage & Chosen before the observed levels; common rescaling changes no decision. \\
\phantomsection\label{entry:symbol-registry-02727}\(S_i\) & \texttt{\detokenize{S_i}} & obtained level divided by its positive weight & P15 multiple testing & n/a & \hyperref[def:p15-mt-weighted]{\ref*{def:p15-mt-weighted}} & \thepage & These ratios determine the weighted ordering. \\
\phantomsection\label{entry:symbol-registry-02728}\(S^{(j)}\) & \texttt{\detokenize{S^{(j)}}} & ratio in the weighted sorted order & P15 multiple testing & n/a & \hyperref[def:p15-mt-weighted]{\ref*{def:p15-mt-weighted}} & \thepage & Equal ratio blocks have whole-block decisions. \\
\phantomsection\label{entry:symbol-registry-02729}\(c^{(j)}\) & \texttt{\detokenize{c^{(j)}}} & weight carried by the current weighted sorting position & P15 multiple testing & n/a & \hyperref[def:p15-mt-weighted]{\ref*{def:p15-mt-weighted}} & \thepage & The remaining sum changes after each accepted rejection step. \\
\phantomsection\label{entry:symbol-registry-02730}\(W_j\) & \texttt{\detokenize{W_j}} & total weight of the hypotheses remaining at step j & P15 multiple testing & n/a & \hyperref[def:p15-mt-weighted]{\ref*{def:p15-mt-weighted}} & \thepage & Includes all still-present true hypotheses at the first true rejection. \\
\phantomsection\label{entry:symbol-registry-02731}\(\mathcal R_{\rm W}\) & \texttt{\detokenize{\mathcal R_{\rm W}}} & fixed-weight Holm rejection set & P15 multiple testing & n/a & \hyperref[def:p15-mt-weighted]{\ref*{def:p15-mt-weighted}} & \thepage & Does not universally dominate unweighted Holm. \\
\phantomsection\label{entry:symbol-registry-02732}\(C_0\) & \texttt{\detokenize{C_0}} & sum of weights on the actual true-null indices & P15 multiple testing & n/a & \hyperref[thm:p15-mt-weighted]{\ref*{thm:p15-mt-weighted}} & \thepage & Positive only in the nonempty true-null case; distinct from a sample-space critical region. \\
\phantomsection\label{entry:symbol-registry-02733}\(d_k\) & \texttt{\detokenize{d_k}} & independence-based critical level with k hypotheses remaining & P15 multiple testing & n/a & \hyperref[def:p15-mt-independent]{\ref*{def:p15-mt-independent}} & \thepage & Its family guarantee requires joint independence of the true obtained levels at each actual configuration. \\
\phantomsection\label{entry:symbol-registry-02734}\(N_g,N_{\rm tot},\nu\) & \texttt{\detokenize{N_g;N_{\rm tot};\nu}} & normal-group observation counts and pooled residual degrees of freedom & P15 multiple testing & n/a & \hyperref[def:p15-mt-control-model]{\ref*{def:p15-mt-control-model}} & \thepage & The treatment-control model permits unequal group sizes with a common variance. \\
\phantomsection\label{entry:symbol-registry-02735}\(s_{\rm p}\) & \texttt{\detokenize{s_{\rm p}}} & pooled residual standard deviation & P15 multiple testing & n/a & \hyperref[def:p15-mt-control-model]{\ref*{def:p15-mt-control-model}} & \thepage & The obtained levels are explicitly defined at zero scale as well. \\
\phantomsection\label{entry:symbol-registry-02736}\(T_i\) & \texttt{\detokenize{T_i}} & Student statistic comparing treatment i with the common control & P15 multiple testing & n/a & \hyperref[def:p15-mt-control-model]{\ref*{def:p15-mt-control-model}} & \thepage & Tests share their control and pooled variance; marginal calibration suffices. \\
\phantomsection\label{entry:symbol-registry-02737}\(F_\nu\) & \texttt{\detokenize{F_\nu}} & Student cumulative distribution function with nu degrees of freedom & P15 multiple testing & n/a & \hyperref[def:p15-mt-control-model]{\ref*{def:p15-mt-control-model}} & \thepage & Uses the existing Student-law definition rather than a normal approximation. \\
\phantomsection\label{entry:symbol-registry-02738}\(Z_i\) & \texttt{\detokenize{Z_i}} & conditioning information for the indexed marginal calibration & P15 conditional multiple-testing calibration & n/a & \hyperref[thm:p15-mt-conditional-calibration]{\ref*{thm:p15-mt-conditional-calibration}} & \thepage & Each marginal test may use different conditioning information. \\
\phantomsection\label{entry:symbol-registry-02739}\(q_{g,k}\) & \texttt{\detokenize{q_{g,k}}} & conditional probability that a remaining group level falls in bin k & \hyperref[annex:p15.mt.normal-power]{\ref*{annex:p15.mt.normal-power}} & n/a & \hyperref[annex:p15.mt.normal-power]{\ref*{annex:p15.mt.normal-power}} & \thepage & The denominator conditions on being above the previous bin boundary. \\
\phantomsection\label{entry:symbol-registry-02740}\(K_k(a,t)\) & \texttt{\detokenize{K_k(a,}\allowbreak\detokenize{t)}} & finite conditional bin-count transition probability & \hyperref[annex:p15.mt.normal-power]{\ref*{annex:p15.mt.normal-power}} & n/a & \hyperref[annex:p15.mt.normal-power]{\ref*{annex:p15.mt.normal-power}} & \thepage & Product of three binomial count probabilities in the stated independent normal example. \\
\phantomsection\label{entry:symbol-registry-02741}\(D_k(a)\) & \texttt{\detokenize{D_k(a)}} & surviving probability mass at a placed-count state & \hyperref[annex:p15.mt.normal-power]{\ref*{annex:p15.mt.normal-power}} & n/a & \hyperref[annex:p15.mt.normal-power]{\ref*{annex:p15.mt.normal-power}} & \thepage & Requires no earlier stopping failure and fewer than two designated rejections. \\
\phantomsection\label{entry:symbol-registry-02742}\(B_k\) & \texttt{\detokenize{B_k}} & absorbed probability that both designated tests have been rejected & \hyperref[annex:p15.mt.normal-power]{\ref*{annex:p15.mt.normal-power}} & n/a & \hyperref[annex:p15.mt.normal-power]{\ref*{annex:p15.mt.normal-power}} & \thepage & The final exact finite sum is numerically evaluated using normal CDF values. \\
\phantomsection\label{entry:symbol-registry-02743}\(\lambda_i\) & \texttt{\detokenize{\lambda_i}} & standardized treatment-control mean difference & \hyperref[annex:p15.mt.control-model]{\ref*{annex:p15.mt.control-model}} & n/a & \hyperref[annex:p15.mt.control-model]{\ref*{annex:p15.mt.control-model}} & \thepage & Nonpositive under the one-sided comparison null. \\
\phantomsection\label{entry:symbol-registry-02744}\(\varepsilon_i\) & \texttt{\detokenize{\varepsilon_i}} & nonnegative marginal calibration-error bound & P15 multiple testing & n/a & \hyperref[thm:p15-mt-calibration-error]{\ref*{thm:p15-mt-calibration-error}} & \thepage & The error contributions remain explicit in the family guarantee. \\
\phantomsection\label{entry:symbol-registry-02745}\(M_g\) & \texttt{\detokenize{M_g}} & number of tested levels in a normal-mean group & \hyperref[annex:p15.mt.normal-power]{\ref*{annex:p15.mt.normal-power}} & n/a & \hyperref[annex:p15.mt.normal-power]{\ref*{annex:p15.mt.normal-power}} & \thepage & Group totals four at mean zero and four at mean six and two at mean three; distinct from observation sample sizes. \\
\phantomsection\label{entry:symbol-registry-02746}\(Z_i\) & \texttt{\detokenize{Z_i}} & standardized centered treatment-control mean contrast & \hyperref[annex:p15.mt.control-model]{\ref*{annex:p15.mt.control-model}} & n/a & \hyperref[annex:p15.mt.control-model]{\ref*{annex:p15.mt.control-model}} & \thepage & A standard normal variable independent of the pooled residual sum under the stated normal model. \\
\phantomsection\label{entry:symbol-registry-02747}\(P_i\) & \texttt{\detokenize{P_i}} & specified one-observation probability law & P15 sequential evidence & n/a & \hyperref[def:p15-se-model]{\ref*{def:p15-se-model}} & \thepage & Under the corresponding hypothesis observations have its product law. \\
\phantomsection\label{entry:symbol-registry-02748}\(H_i\) & \texttt{\detokenize{H_i}} & simple statistical hypothesis & P15 sequential evidence & n/a & \hyperref[def:p15-se-model]{\ref*{def:p15-se-model}} & \thepage & The two specified laws exhaust the simple-hypothesis comparison. \\
\phantomsection\label{entry:symbol-registry-02749}\(f_i\) & \texttt{\detokenize{f_i}} & observation density against the common measure & P15 sequential evidence & n/a & \hyperref[def:p15-se-model]{\ref*{def:p15-se-model}} & \thepage & Finite-score third-law calculations require positive finite density versions under that law. \\
\phantomsection\label{entry:symbol-registry-02750}\(p_{im}\) & \texttt{\detokenize{p_{im}}} & finite-history joint density & P15 sequential evidence & n/a & \hyperref[def:p15-se-model]{\ref*{def:p15-se-model}} & \thepage & Product of the observation densities in the independent model; the dependent extension uses actual joint densities. \\
\phantomsection\label{entry:symbol-registry-02751}\(L_m\) & \texttt{\detokenize{L_m}} & observed-history likelihood ratio & P15 sequential evidence & n/a & \hyperref[def:p15-se-model]{\ref*{def:p15-se-model}} & \thepage & One-sided zero densities yield zero or infinity and terminate the implemented test. \\
\phantomsection\label{entry:symbol-registry-02752}\(\mathcal F_m\) & \texttt{\detokenize{\mathcal F_m}} & observed information at stage m & P15 sequential evidence & n/a & \hyperref[def:p15-se-model]{\ref*{def:p15-se-model}} & \thepage & An explicit independent randomization seed may be included at the initial stage. \\
\phantomsection\label{entry:symbol-registry-02753}\(z_j\) & \texttt{\detokenize{z_j}} & one-observation log likelihood score & P15 sequential evidence & n/a & \hyperref[def:p15-se-model]{\ref*{def:p15-se-model}} & \thepage & Finite almost surely for equivalent laws; singular observations stop before any opposite infinity is added. \\
\phantomsection\label{entry:symbol-registry-02754}\(Z_m\) & \texttt{\detokenize{Z_m}} & cumulative log likelihood score & P15 sequential evidence & n/a & \hyperref[def:p15-se-model]{\ref*{def:p15-se-model}} & \thepage & The sum is used only with finite scores or before an immediate singular decision. \\
\phantomsection\label{entry:symbol-registry-02755}\(R_m^0,R_m^1,R_m\) & \texttt{\detokenize{R_m^0,}\allowbreak\detokenize{R_m^1,}\allowbreak\detokenize{R_m}} & partition into two terminal regions and continuation & P15 sequential evidence & n/a & \hyperref[def:p15-se-procedure]{\ref*{def:p15-se-procedure}} & \thepage & Regions are measurable in the observed finite history. \\
\phantomsection\label{entry:symbol-registry-02756}\(T\) & \texttt{\detokenize{T}} & first terminal stage & P15 sequential evidence & n/a & \hyperref[def:p15-se-procedure]{\ref*{def:p15-se-procedure}} & \thepage & Possibly infinite until termination is proved; equality at either boundary stops. \\
\phantomsection\label{entry:symbol-registry-02757}\(D\) & \texttt{\detokenize{D}} & terminal binary decision & P15 sequential evidence & n/a & \hyperref[def:p15-se-procedure]{\ref*{def:p15-se-procedure}} & \thepage & Zero chooses the null and one chooses the alternative. \\
\phantomsection\label{entry:symbol-registry-02758}\(Q_i\) & \texttt{\detokenize{Q_i}} & terminal decision event & P15 sequential evidence & n/a & \hyperref[def:p15-se-procedure]{\ref*{def:p15-se-procedure}} & \thepage & Here these are subsets of stopped histories; their local type differs from a probability law. \\
\phantomsection\label{entry:symbol-registry-02759}\(\alpha',\beta'\) & \texttt{\detokenize{\alpha',}\allowbreak\detokenize{\beta'}} & actual error probabilities of the chosen procedure & P15 sequential evidence & n/a & \hyperref[def:p15-se-procedure]{\ref*{def:p15-se-procedure}} & \thepage & They need not equal the prescribed targets. \\
\phantomsection\label{entry:symbol-registry-02760}\(\alpha,\beta\) & \texttt{\detokenize{\alpha,}\allowbreak\detokenize{\beta}} & prescribed type-I and type-II error targets & P15 sequential evidence & n/a & \hyperref[def:p15-se-procedure]{\ref*{def:p15-se-procedure}} & \thepage & Nominal and conservative calibrations have different exact guarantees. \\
\phantomsection\label{entry:symbol-registry-02761}\(g_0,g_1\) & \texttt{\detokenize{g_0,}\allowbreak\detokenize{g_1}} & prior probabilities of the exhaustive simple hypotheses & P15 sequential evidence & n/a & \hyperref[thm:p15-se-posterior]{\ref*{thm:p15-se-posterior}} & \thepage & They are positive and sum to one. \\
\phantomsection\label{entry:symbol-registry-02762}\(d_0,d_1\) & \texttt{\detokenize{d_0,}\allowbreak\detokenize{d_1}} & posterior correct-decision thresholds & P15 sequential evidence & n/a & \hyperref[thm:p15-se-posterior]{\ref*{thm:p15-se-posterior}} & \thepage & Each exceeds one half; these are posterior levels rather than frequentist errors. \\
\phantomsection\label{entry:symbol-registry-02763}\(A,B\) & \texttt{\detokenize{A,}\allowbreak\detokenize{B}} & upper and lower likelihood-ratio boundaries & P15 sequential evidence & n/a & \hyperref[def:p15-se-sprt]{\ref*{def:p15-se-sprt}} & \thepage & The original general rule has zero less than B less than A; the practical calibration scope has B below one and A above one. \\
\phantomsection\label{entry:symbol-registry-02764}\(a,b\) & \texttt{\detokenize{a,}\allowbreak\detokenize{b}} & logarithms of the upper and lower boundaries & P15 sequential evidence & n/a & \hyperref[def:p15-se-sprt]{\ref*{def:p15-se-sprt}} & \thepage & The upper log boundary is positive and the lower one negative in the declared practical calibration scope. \\
\phantomsection\label{entry:symbol-registry-02765}\(\eta_0\) & \texttt{\detokenize{\eta_0}} & affinity of the two observation laws & P15 sequential evidence & n/a & \hyperref[thm:p15-se-termination]{\ref*{thm:p15-se-termination}} & \thepage & Strictly below one for distinct laws; separate from the source exponential excess constant. \\
\phantomsection\label{entry:symbol-registry-02766}\(\mu\) & \texttt{\detokenize{\mu}} & mean one-observation score under the selected law & P15 sequential evidence & n/a & \hyperref[def:p15-se-overshoot]{\ref*{def:p15-se-overshoot}} & \thepage & Division by this mean is used only when it is nonzero. \\
\phantomsection\label{entry:symbol-registry-02767}\(\gamma\) & \texttt{\detokenize{\gamma}} & actual upper-terminal probability under the selected law & P15 sequential evidence & n/a & \hyperref[def:p15-se-overshoot]{\ref*{def:p15-se-overshoot}} & \thepage & Under a third law it is the operating characteristic at that law. \\
\phantomsection\label{entry:symbol-registry-02768}\(\sigma_z^2\) & \texttt{\detokenize{\sigma_z^2}} & finite positive score variance in the centered stopping identity & P15 sequential evidence & n/a & \hyperref[thm:p15-se-zero-drift]{\ref*{thm:p15-se-zero-drift}} & \thepage & The centered score has second moment equal to this variance. \\
\phantomsection\label{entry:symbol-registry-02769}\(r_-,r_+\) & \texttt{\detokenize{r_-,}\allowbreak\detokenize{r_+}} & nonnegative lower and upper boundary overshoots & P15 sequential evidence & n/a & \hyperref[def:p15-se-overshoot]{\ref*{def:p15-se-overshoot}} & \thepage & Defined on their respective terminal decision events. \\
\phantomsection\label{entry:symbol-registry-02770}\(c_-,c_+\) & \texttt{\detokenize{c_-,}\allowbreak\detokenize{c_+}} & conditional mean lower and upper overshoots & P15 sequential evidence & n/a & \hyperref[def:p15-se-overshoot]{\ref*{def:p15-se-overshoot}} & \thepage & A zero-probability terminal event contributes zero. \\
\phantomsection\label{entry:symbol-registry-02771}\(\xi,\xi'\) & \texttt{\detokenize{\xi,}\allowbreak\detokenize{\xi'}} & upper mean excess and signed lower mean excess of one score & P15 sequential evidence & n/a & \hyperref[def:p15-se-overshoot]{\ref*{def:p15-se-overshoot}} & \thepage & The prime retains the source negative signed convention; finite extrema are explicit assumptions. \\
\phantomsection\label{entry:symbol-registry-02772}\(\varphi(t)\) & \texttt{\detokenize{\varphi(t)}} & exponential moment of one score & P15 sequential evidence & n/a & \hyperref[def:p15-se-exponential-root]{\ref*{def:p15-se-exponential-root}} & \thepage & For complex t existence means absolute integrability of the real-part exponential. \\
\phantomsection\label{entry:symbol-registry-02773}\(h\) & \texttt{\detokenize{h}} & nonzero real exponential root & P15 sequential evidence & n/a & \hyperref[def:p15-se-exponential-root]{\ref*{def:p15-se-exponential-root}} & \thepage & A root is assumed or proved for the specific family; zero drift may leave only the zero root. \\
\phantomsection\label{entry:symbol-registry-02774}\(\eta,\delta\) & \texttt{\detokenize{\eta,}\allowbreak\detokenize{\delta}} & one-step multiplicative exponential excess extrema & P15 sequential evidence & n/a & \hyperref[def:p15-se-exponential-root]{\ref*{def:p15-se-exponential-root}} & \thepage & These quantities bound terminal exponential overshoot; they are not the affinity. \\
\phantomsection\label{entry:symbol-registry-02775}\(\bar\gamma\) & \texttt{\detokenize{\bar\gamma}} & operating-characteristic boundary approximation & P15 sequential evidence & n/a & \hyperref[thm:p15-se-oc]{\ref*{thm:p15-se-oc}} & \thepage & Obtained by neglecting both exponential overshoots; its zero-root expression is an algebraic limit. \\
\phantomsection\label{entry:symbol-registry-02776}\(g_1,g_2\) & \texttt{\detokenize{g_1,}\allowbreak\detokenize{g_2}} & integer magnitudes of the extreme lattice score increments & P15 sequential evidence & n/a & \hyperref[thm:p15-se-lattice]{\ref*{thm:p15-se-lattice}} & \thepage & Both endpoint masses are positive in the lattice construction. \\
\phantomsection\label{entry:symbol-registry-02777}\(a_0,b_0\) & \texttt{\detokenize{a_0,}\allowbreak\detokenize{b_0}} & upper ceiling and lower floor lattice boundary levels & P15 sequential evidence & n/a & \hyperref[thm:p15-se-lattice]{\ref*{thm:p15-se-lattice}} & \thepage & The rounded boundaries encode the unchanged real score exit criterion. \\
\phantomsection\label{entry:symbol-registry-02778}\(q_j\) & \texttt{\detokenize{q_j}} & probability of one integer score value & P15 sequential evidence & n/a & \hyperref[thm:p15-se-lattice]{\ref*{thm:p15-se-lattice}} & \thepage & A finite probability mass sequence used in the root equation. \\
\phantomsection\label{entry:symbol-registry-02779}\(c_j\) & \texttt{\detokenize{c_j}} & possible terminal lattice score value & P15 sequential evidence & n/a & \hyperref[thm:p15-se-lattice]{\ref*{thm:p15-se-lattice}} & \thepage & The listed lower and upper exit bands include all overshoots. \\
\phantomsection\label{entry:symbol-registry-02780}\(u_i\) & \texttt{\detokenize{u_i}} & nonzero complex root of the lattice equation & P15 sequential evidence & n/a & \hyperref[thm:p15-se-lattice]{\ref*{thm:p15-se-lattice}} & \thepage & Multiplicity is retained; a repeated row can make the source determinant singular. \\
\phantomsection\label{entry:symbol-registry-02781}\(\mathcal M,\mathcal M_j\) & \texttt{\detokenize{\mathcal M,}\allowbreak\detokenize{\mathcal M_j}} & root coefficient matrix and its column replacement & P15 sequential evidence & n/a & \hyperref[thm:p15-se-lattice]{\ref*{thm:p15-se-lattice}} & \thepage & Cramers rule is asserted only for a nonzero determinant. \\
\phantomsection\label{entry:symbol-registry-02782}\(J\) & \texttt{\detokenize{J}} & finite set of continuing lattice sums & P15 sequential evidence & n/a & \hyperref[thm:p15-se-finite-recursion]{\ref*{thm:p15-se-finite-recursion}} & \thepage & It contains the starting sum zero. \\
\phantomsection\label{entry:symbol-registry-02783}\(P_x,E_x,\gamma_x\) & \texttt{\detokenize{P_x,}\allowbreak\detokenize{E_x,}\allowbreak\detokenize{\gamma_x}} & probability expectation and upper-exit probability from starting sum x & P15 sequential evidence & n/a & \hyperref[thm:p15-se-finite-recursion]{\ref*{thm:p15-se-finite-recursion}} & \thepage & The start is a specified continuing lattice state. \\
\phantomsection\label{entry:symbol-registry-02784}\(K_J\) & \texttt{\detokenize{K_J}} & substochastic continuing-state transition matrix & P15 sequential evidence & n/a & \hyperref[thm:p15-se-finite-recursion]{\ref*{thm:p15-se-finite-recursion}} & \thepage & Its absorption series converges by a uniform block exit probability. \\
\phantomsection\label{entry:symbol-registry-02785}\(v_+,v_-\) & \texttt{\detokenize{v_+,}\allowbreak\detokenize{v_-}} & one-step upper and lower exit probability vectors & P15 sequential evidence & n/a & \hyperref[thm:p15-se-finite-recursion]{\ref*{thm:p15-se-finite-recursion}} & \thepage & Together with the continuing matrix they determine exact absorption and stopping-time transforms. \\
\phantomsection\label{entry:symbol-registry-02786}\(z_+,z_-\) & \texttt{\detokenize{z_+,}\allowbreak\detokenize{z_-}} & positive and negative Bernoulli score values & P15 sequential evidence & n/a & \hyperref[thm:p15-se-binomial]{\ref*{thm:p15-se-binomial}} & \thepage & Both values are finite for interior and separated Bernoulli parameters. \\
\phantomsection\label{entry:symbol-registry-02787}\(p_0,p_1,p\) & \texttt{\detokenize{p_0,}\allowbreak\detokenize{p_1,}\allowbreak\detokenize{p}} & Bernoulli null alternative and third-law parameters & P15 sequential evidence & n/a & \hyperref[thm:p15-se-binomial]{\ref*{thm:p15-se-binomial}} & \thepage & The two hypotheses satisfy zero less than p0 less than p1 less than one. \\
\phantomsection\label{entry:symbol-registry-02788}\(d,\theta\) & \texttt{\detokenize{d,}\allowbreak\detokenize{\theta}} & normal hypothesis half-separation and third-law mean & P15 sequential evidence & n/a & \hyperref[thm:p15-se-normal]{\ref*{thm:p15-se-normal}} & \thepage & The hypothesis means are minus d and d with variance one. \\
\phantomsection\label{entry:symbol-registry-02789}\(\phi,\Phi\) & \texttt{\detokenize{\phi,}\allowbreak\detokenize{\Phi}} & standard normal density and distribution function & P15 sequential evidence & n/a & \hyperref[thm:p15-se-normal]{\ref*{thm:p15-se-normal}} & \thepage & Distinct from the score exponential moment written varphi. \\
\phantomsection\label{entry:symbol-registry-02790}\(q_u\) & \texttt{\detokenize{q_u}} & standard normal u-quantile & P15 sequential evidence & n/a & \hyperref[def:p15-se-normal-saving]{\ref*{def:p15-se-normal-saving}} & \thepage & This indexed role is a quantile; the lattice probability masses use index j. \\
\phantomsection\label{entry:symbol-registry-02791}\(N_{\rm fixed}\) & \texttt{\detokenize{N_{\rm fixed}}} & minimum fixed normal sample size for the prescribed errors & P15 sequential evidence & n/a & \hyperref[def:p15-se-normal-saving]{\ref*{def:p15-se-normal-saving}} & \thepage & Integer ceiling is retained; the separate sequential saving formula neglects overshoot. \\
\phantomsection\label{entry:symbol-registry-02792}\(\Psi_+,\Psi_-\) & \texttt{\detokenize{\Psi_+,}\allowbreak\detokenize{\Psi_-}} & exact conditional stopping-time exponential transforms & P15 sequential evidence & n/a & \hyperref[def:p15-se-transform-approximation]{\ref*{def:p15-se-transform-approximation}} & \thepage & Their convergence domains are explicit; restriction to an imaginary argument gives characteristic functions. \\
\phantomsection\label{entry:symbol-registry-02793}\(\Psi\) & \texttt{\detokenize{\Psi}} & exact unconditional stopping-time exponential transform & P15 sequential evidence & n/a & \hyperref[def:p15-se-transform-approximation]{\ref*{def:p15-se-transform-approximation}} & \thepage & Its mixture weights are the actual terminal probabilities. \\
\phantomsection\label{entry:symbol-registry-02794}\(\widehat\Psi_+,\widehat\Psi_-\) & \texttt{\detokenize{\widehat\Psi_+,}\allowbreak\detokenize{\widehat\Psi_-}} & boundary approximations to the conditional transforms & P15 sequential evidence & n/a & \hyperref[def:p15-se-transform-approximation]{\ref*{def:p15-se-transform-approximation}} & \thepage & Use the approximated operating characteristic consistently for normalization at zero. \\
\phantomsection\label{entry:symbol-registry-02795}\(\widehat\Psi\) & \texttt{\detokenize{\widehat\Psi}} & boundary approximation to the unconditional transform & P15 sequential evidence & n/a & \hyperref[def:p15-se-transform-approximation]{\ref*{def:p15-se-transform-approximation}} & \thepage & The solved two-root expression is not a discrete stopping-time error guarantee. \\
\phantomsection\label{entry:symbol-registry-02796}\(\tau\) & \texttt{\detokenize{\tau}} & stopping-time transform argument & P15 sequential evidence & n/a & \hyperref[def:p15-se-transform-approximation]{\ref*{def:p15-se-transform-approximation}} & \thepage & Real arguments give exponential transforms and imaginary arguments give characteristic functions. \\
\phantomsection\label{entry:symbol-registry-02797}\(t_1,t_2\) & \texttt{\detokenize{t_1,}\allowbreak\detokenize{t_2}} & selected roots for the two-boundary transform system & P15 sequential evidence & n/a & \hyperref[def:p15-se-transform-approximation]{\ref*{def:p15-se-transform-approximation}} & \thepage & Local branches tend to zero and the nonzero score exponential root. \\
\phantomsection\label{entry:symbol-registry-02798}\(\mathcal D\) & \texttt{\detokenize{\mathcal D}} & two-root transform coefficient determinant & P15 sequential evidence & n/a & \hyperref[def:p15-se-transform-approximation]{\ref*{def:p15-se-transform-approximation}} & \thepage & At a zero denominator a justified limit is required. \\
\phantomsection\label{entry:symbol-registry-02799}\(f_c\) & \texttt{\detokenize{f_c}} & continuous one-sided stopping-time approximation density & P15 sequential evidence & n/a & \hyperref[def:p15-se-transform-approximation]{\ref*{def:p15-se-transform-approximation}} & \thepage & Its normalization and moments are proved; no unproved discrete-to-continuous error bound is claimed. \\
\phantomsection\label{entry:symbol-registry-02800}\(m,c\) & \texttt{\detokenize{m,}\allowbreak\detokenize{c}} & scaled continuous stopping size and its positive density parameter & P15 sequential evidence & n/a & \hyperref[def:p15-se-transform-approximation]{\ref*{def:p15-se-transform-approximation}} & \thepage & Here m is the declared continuous scaling variable rather than a finite inspection-stage index. \\
\phantomsection\label{entry:symbol-registry-02801}\(D_N\) & \texttt{\detokenize{D_N}} & forced finite-horizon terminal decision & P15 sequential evidence & n/a & \hyperref[thm:p15-se-truncation]{\ref*{thm:p15-se-truncation}} & \thepage & The zero-score tie accepts the null after the original continuation condition. \\
\phantomsection\label{entry:symbol-registry-02802}\(\mu_i,\sigma_i^2\) & \texttt{\detokenize{\mu_i,}\allowbreak\detokenize{\sigma_i^2}} & normal score mean and variance under the specified hypothesis & P15 sequential evidence & n/a & \hyperref[thm:p15-se-deadline]{\ref*{thm:p15-se-deadline}} & \thepage & The exact finite-horizon score has mean N times mu and variance N times the single-score variance. \\
\phantomsection\label{entry:symbol-registry-02803}\(\alpha_N,\beta_N\) & \texttt{\detokenize{\alpha_N,}\allowbreak\detokenize{\beta_N}} & actual errors of the capped procedure & P15 sequential evidence & n/a & \hyperref[thm:p15-se-truncation]{\ref*{thm:p15-se-truncation}} & \thepage & Bounds add a new terminal interval probability to the actual uncapped errors. \\
\phantomsection\label{entry:symbol-registry-02804}\(U_j\) & \texttt{\detokenize{U_j}} & independent uniform result for a fixed one-observation inspection & P15 sequential evidence & n/a & \hyperref[thm:p15-se-repeated-fixed]{\ref*{thm:p15-se-repeated-fixed}} & \thepage & Independence is specific to the calculated repeated-inspection counterexample. \\
\phantomsection\label{entry:symbol-registry-02805}\(M\) & \texttt{\detokenize{M}} & number of repeated fixed-sample inspections & P15 sequential evidence & n/a & \hyperref[thm:p15-se-repeated-fixed]{\ref*{thm:p15-se-repeated-fixed}} & \thepage & An integer horizon; the root coefficient matrix uses a calligraphic symbol. \\
\phantomsection\label{entry:symbol-registry-02806}\(u,v\) & \texttt{\detokenize{u,}\allowbreak\detokenize{v}} & actual errors of a competing sequential procedure & P15 sequential evidence & n/a & \hyperref[thm:p15-se-information]{\ref*{thm:p15-se-information}} & \thepage & The information lower bounds use these actual errors at the same strength. \\
\phantomsection\label{entry:symbol-registry-02807}\(I_i\) & \texttt{\detokenize{I_i}} & binary information numerator for the chosen test & P15 sequential evidence & n/a & \hyperref[thm:p15-se-no-overshoot]{\ref*{thm:p15-se-no-overshoot}} & \thepage & The bounded-excess theorem compares expected cost with its universal information lower bound. \\
\phantomsection\label{entry:symbol-registry-02808}\(C\) & \texttt{\detokenize{C}} & absolute bound on a single score increment & P15 sequential evidence & n/a & \hyperref[thm:p15-se-no-overshoot]{\ref*{thm:p15-se-no-overshoot}} & \thepage & In this result it bounds each terminal conditional score interval width. \\
\phantomsection\label{entry:symbol-registry-02809}\(G,H\) & \texttt{\detokenize{G,}\allowbreak\detokenize{H}} & normal upper-tail probability and hazard ratio & P15 sequential evidence annex & n/a & \hyperref[annex:p15.se.normal]{\ref*{annex:p15.se.normal}} & \thepage & Proof-local functions used to derive both excess extrema and tail-ratio monotonicity. \\
\phantomsection\label{entry:symbol-registry-02810}\(F(s)\) & \texttt{\detokenize{F(s)}} & auxiliary Gaussian integral for the density normalization & P15 sequential evidence annex & n/a & \hyperref[annex:p15.se.continuous-density]{\ref*{annex:p15.se.continuous-density}} & \thepage & A proof-local integral; it is not an observation information field. \\
\phantomsection\label{entry:symbol-registry-02811}\(|\varphi(t)|\) & \texttt{\detokenize{|\varphi(t)|}} & complex modulus of the score exponential transform & P15 sequential evidence & n/a & \hyperref[thm:p15-se-transform]{\ref*{thm:p15-se-transform}} & \thepage & Typed complex modulus in the stopped-transform premise; no cardinality role is intended. \\
\phantomsection\label{entry:symbol-registry-02812}\(|z_1|\) & \texttt{\detokenize{|z_1|}} & real absolute value of the observation score & P15 sequential evidence & n/a & \hyperref[thm:p15-se-mean]{\ref*{thm:p15-se-mean}} & \thepage & This typed absolute value supplies the integrability premise for Walds stopped-sum identity. \\
\phantomsection\label{entry:symbol-registry-02813}\(\Omega\) & \texttt{\detokenize{\Omega}} & nonempty outcome space & P15 imprecise-probability lower-prevision treatment & n/a & \hyperref[def:p15-ip-domain]{\ref*{def:p15-ip-domain}} & \thepage & The general space need not be finite or countable. \\
\phantomsection\label{entry:symbol-registry-02814}\(\mathcal L(\Omega)\) & \texttt{\detokenize{\mathcal L}} & all bounded real-valued gambles on the outcome space & P15 imprecise-probability lower-prevision treatment & n/a & \hyperref[def:p15-ip-domain]{\ref*{def:p15-ip-domain}} & \thepage & Payoffs use a specified linear utility scale. \\
\phantomsection\label{entry:symbol-registry-02815}\(\Sigma\) & \texttt{\detokenize{\Sigma}} & event algebra on the stated outcome space & P15 imprecise-probability lower-prevision treatment & n/a & \hyperref[def:p15-ip-domain]{\ref*{def:p15-ip-domain}} & \thepage & A sigma algebra is an additional assumption in the countable-additivity theorem. \\
\phantomsection\label{entry:symbol-registry-02816}\(B(\Sigma),B\) & \texttt{\detokenize{B(\Sigma);B}} & uniform closure of the finite-valued simple gambles on the algebra & P15 imprecise-probability lower-prevision treatment & n/a & \hyperref[def:p15-ip-domain]{\ref*{def:p15-ip-domain}} & \thepage & For the power-set algebra this is all bounded gambles; no replacement of the algebra by its generated sigma algebra is made. \\
\phantomsection\label{entry:symbol-registry-02817}\(\mathcal K\) & \texttt{\detokenize{\mathcal K}} & domain of assessed gambles & P15 imprecise-probability lower-prevision treatment & n/a & \hyperref[def:p15-ip-domain]{\ref*{def:p15-ip-domain}} & \thepage & The domain is an arbitrary subset of the full gamble space and need not be a vector space. \\
\phantomsection\label{entry:symbol-registry-02818}\(I_A\) & \texttt{\detokenize{I_A}} & indicator gamble of an event in the algebra & P15 imprecise-probability lower-prevision treatment & n/a & \hyperref[def:p15-ip-domain]{\ref*{def:p15-ip-domain}} & \thepage & An event outside the stated algebra is not silently assigned probability zero. \\
\phantomsection\label{entry:symbol-registry-02819}\(\underline P\) & \texttt{\detokenize{\underline P}} & assessed supremum buying-price functional & P15 imprecise-probability lower-prevision treatment & n/a & \hyperref[def:p15-ip-prices]{\ref*{def:p15-ip-prices}} & \thepage & The actual acceptance of a gamble at its supremum price is not determined by this number. \\
\phantomsection\label{entry:symbol-registry-02820}\(\overline P\) & \texttt{\detokenize{\overline P}} & infimum selling-price functional conjugate to the lower assessment & P15 imprecise-probability lower-prevision treatment & n/a & \hyperref[def:p15-ip-prices]{\ref*{def:p15-ip-prices}} & \thepage & The conjugate domain is the negative of the assessed domain. \\
\phantomsection\label{entry:symbol-registry-02821}\(G(f)\) & \texttt{\detokenize{G}} & centered assessed buying gain & P15 imprecise-probability lower-prevision treatment & n/a & \hyperref[def:p15-ip-coherence]{\ref*{def:p15-ip-coherence}} & \thepage & G of f is f minus its lower assessed price; portfolio criteria use pointwise suprema. \\
\phantomsection\label{entry:symbol-registry-02822}\(\underline E\) & \texttt{\detokenize{\underline E}} & natural extension of the lower assessment & P15 imprecise-probability lower-prevision treatment & n/a & \hyperref[def:p15-ip-natural-extension]{\ref*{def:p15-ip-natural-extension}} & \thepage & Finite exactly in the sure-loss-avoiding case; least coherent dominating full-domain functional. \\
\phantomsection\label{entry:symbol-registry-02823}\(\overline E\) & \texttt{\detokenize{\overline E}} & conjugate upper natural extension & P15 imprecise-probability lower-prevision treatment & n/a & \hyperref[thm:p15-ip-dual]{\ref*{thm:p15-ip-dual}} & \thepage & Equals the largest expectation among the compatible finitely additive models. \\
\phantomsection\label{entry:symbol-registry-02824}\(I\) & \texttt{\detokenize{I}} & full-domain lower-prevision functional in the coherence characterization & P15 imprecise-probability lower-prevision treatment & n/a & \hyperref[thm:p15-ip-linear-domain]{\ref*{thm:p15-ip-linear-domain}} & \thepage & This I is a functional rather than an event indicator; indicator notation always has an event subscript. \\
\phantomsection\label{entry:symbol-registry-02825}\(L\) & \texttt{\detokenize{L}} & positive normalized linear prevision & P15 imprecise-probability lower-prevision treatment & n/a & \hyperref[def:p15-ip-compatible]{\ref*{def:p15-ip-compatible}} & \thepage & Corresponds to a finitely additive probability on the selected event algebra. \\
\phantomsection\label{entry:symbol-registry-02826}\(\mathcal M(\underline P)\) & \texttt{\detokenize{\mathcal M}} & closed convex set of all compatible linear previsions & P15 imprecise-probability lower-prevision treatment & n/a & \hyperref[def:p15-ip-compatible]{\ref*{def:p15-ip-compatible}} & \thepage & Contains every precise model dominating the assessments; it is nonempty exactly when sure loss is avoided. \\
\phantomsection\label{entry:symbol-registry-02827}\(C,D\) & \texttt{\detokenize{C;D}} & full and reduced six-permutation prior hulls in the event-information example & P15 imprecise-probability event-information example & n/a & \hyperref[thm:p15-ip-event-separation]{\ref*{thm:p15-ip-event-separation}} & \thepage & The hulls have equal event envelopes but different expectations for the specified utility vector. \\
\phantomsection\label{entry:symbol-registry-02828}\(\pi_i,B_i,m(B_i)\) & \texttt{\detokenize{\pi_i;B_i;m}} & ordered possibility score nested focal set and its difference mass & P15 imprecise-probability finite possibility bridge & n/a & \hyperref[thm:p15-ip-possibility-bridge]{\ref*{thm:p15-ip-possibility-bridge}} & \thepage & The bridge uses the P14 definitions; it does not duplicate the finite comparison table. \\
\phantomsection\label{entry:symbol-registry-02829}\(Q\) & \texttt{\detokenize{Q}} & conditional positive normalized linear prevision supported on the observed event & P15 imprecise-probability lower-prevision treatment & n/a & \hyperref[def:p15-ip-conditional]{\ref*{def:p15-ip-conditional}} & \thepage & Bayes-compatible pairs retain the zero-mass branch rather than dividing by zero. \\
\phantomsection\label{entry:symbol-registry-02830}\(Q_L\) & \texttt{\detokenize{Q_L}} & Bayesian conditional associated with a positive-mass compatible model & P15 imprecise-probability lower-prevision treatment & n/a & \hyperref[annex:p15.ip.conditional-models]{\ref*{annex:p15.ip.conditional-models}} & \thepage & Its numerator is the expectation of the event-indicated gamble and its denominator is the positive event probability. \\
\phantomsection\label{entry:symbol-registry-02831}\(C(\cdot\mid A)\) & \texttt{\detokenize{C}} & arbitrary separately and jointly coherent conditional lower prevision & P15 imprecise-probability lower-prevision treatment & n/a & \hyperref[thm:p15-ip-conditional-bounds]{\ref*{thm:p15-ip-conditional-bounds}} & \thepage & The full-domain one-event portfolio criterion forces the stated generalized-Bayes identity. \\
\phantomsection\label{entry:symbol-registry-02832}\(\underline R(g\mid A)\) & \texttt{\detokenize{\underline R}} & regular conditional lower prevision & P15 imprecise-probability lower-prevision treatment & n/a & \hyperref[def:p15-ip-updates]{\ref*{def:p15-ip-updates}} & \thepage & Uses positive-mass compatible models; the all-zero-mass case has an explicit vacuous convention. \\
\phantomsection\label{entry:symbol-registry-02833}\(F(q)\) & \texttt{\detokenize{F}} & generalized-Bayes root functional of a candidate conditional price & P15 imprecise-probability Bayes-root proof & n/a & \hyperref[annex:p15.ip.generalized-bayes]{\ref*{annex:p15.ip.generalized-bayes}} & \thepage & Unique root requires positive lower event probability; at zero lower and positive upper probability the root set is a half-line. \\
\phantomsection\label{entry:symbol-registry-02834}\(X,Y\) & \texttt{\detokenize{X;Y}} & coordinates of the declared binary product outcome space & P15 imprecise-probability independence example & n/a & \hyperref[thm:p15-ip-independence-boundary]{\ref*{thm:p15-ip-independence-boundary}} & \thepage & The strong-product and conditional assertions refer to this specified full product space. \\
\phantomsection\label{entry:symbol-registry-02889}\(k\) & \texttt{\detokenize{k}} & number of study estimates & P15 evidence synthesis & n/a & \hyperref[def:p15-synthesis-record]{\ref*{def:p15-synthesis-record}} & \thepage & At least one for a review; dispersion and residual methods require at least two. \\
\phantomsection\label{entry:symbol-registry-02890}\(Y_i\) & \texttt{\detokenize{Y_i}} & study-level effect estimator & P15 evidence synthesis & n/a & \hyperref[def:p15-synthesis-record]{\ref*{def:p15-synthesis-record}} & \thepage & The recorded effect scale and analysis model belong to the estimate. \\
\phantomsection\label{entry:symbol-registry-02891}\(\theta_i\) & \texttt{\detokenize{\theta_i}} & study-specific estimand & P15 evidence synthesis & n/a & \hyperref[def:p15-synthesis-record]{\ref*{def:p15-synthesis-record}} & \thepage & Not a measured estimate or automatically a causal effect. \\
\phantomsection\label{entry:symbol-registry-02892}\(\theta_a\) & \texttt{\detokenize{\theta_a}} & declared finite weighted estimand & P15 evidence synthesis & n/a & \hyperref[def:p15-synthesis-record]{\ref*{def:p15-synthesis-record}} & \thepage & Deterministic normalized weights define this target. \\
\phantomsection\label{entry:symbol-registry-02893}\(a_i\) & \texttt{\detokenize{a_i}} & deterministic normalized synthesis weight & P15 synthesis weight model & n/a & \hyperref[def:p15-synthesis-record]{\ref*{def:p15-synthesis-record}} & \thepage & Separate from binary event-count notation in the table section. \\
\phantomsection\label{entry:symbol-registry-02894}\(\mu\) & \texttt{\detokenize{\mu}} & mean of modelled study effects & P15 evidence synthesis & n/a & \hyperref[def:p15-synthesis-record]{\ref*{def:p15-synthesis-record}} & \thepage & Confidence for the mean differs from prediction for a future effect. \\
\phantomsection\label{entry:symbol-registry-02895}\(\tau^2\) & \texttt{\detokenize{\tau^2}} & between-study effect variance & P15 evidence synthesis & n/a & \hyperref[def:p15-synthesis-record]{\ref*{def:p15-synthesis-record}} & \thepage & Known in exact results; estimated in practical calculations. \\
\phantomsection\label{entry:symbol-registry-02896}\(\Sigma\) & \texttt{\detokenize{\Sigma}} & covariance matrix of the estimate vector & P15 evidence synthesis & n/a & \hyperref[def:p15-synthesis-linear-model]{\ref*{def:p15-synthesis-linear-model}} & \thepage & Off-diagonal entries retain shared data and participants. \\
\phantomsection\label{entry:symbol-registry-02897}\(v_i\) & \texttt{\detokenize{v_i}} & known actual within-study sampling variance & P15 evidence synthesis & n/a & \hyperref[def:p15-synthesis-linear-model]{\ref*{def:p15-synthesis-linear-model}} & \thepage & Positive in exact normal and whitening arguments. \\
\phantomsection\label{entry:symbol-registry-02898}\(s_i^2\) & \texttt{\detokenize{s_i^2}} & estimated within-study sampling variance & P15 evidence synthesis & n/a & \hyperref[def:p15-synthesis-linear-model]{\ref*{def:p15-synthesis-linear-model}} & \thepage & Replacing known variances by these values does not inherit exact guarantees. \\
\phantomsection\label{entry:symbol-registry-02899}\(w_i\) & \texttt{\detokenize{w_i}} & inverse sampling-variance weight & P15 synthesis weight model & n/a & \hyperref[def:p15-synthesis-linear-model]{\ref*{def:p15-synthesis-linear-model}} & \thepage & Deterministic in the exact common-effect theorem. \\
\phantomsection\label{entry:symbol-registry-02900}\(W\) & \texttt{\detokenize{W}} & total inverse-variance weight & P15 synthesis weight model & n/a & \hyperref[def:p15-synthesis-linear-model]{\ref*{def:p15-synthesis-linear-model}} & \thepage & Positive sum over included study estimates. \\
\phantomsection\label{entry:symbol-registry-02901}\(a_*\) & \texttt{\detokenize{a_*}} & minimum-variance normalized linear weight vector & P15 synthesis weight model & n/a & \hyperref[thm:p15-synthesis-linear]{\ref*{thm:p15-synthesis-linear}} & \thepage & Uses an invertible positive definite covariance matrix. \\
\phantomsection\label{entry:symbol-registry-02902}\(A_i(Y)\) & \texttt{\detokenize{A_i(Y)}} & data-dependent normalized weight & P15 synthesis weight model & n/a & \hyperref[thm:p15-synthesis-estimated-weights]{\ref*{thm:p15-synthesis-estimated-weights}} & \thepage & Separate from conditional binary event-count variable. \\
\phantomsection\label{entry:symbol-registry-02903}\((a_i,b_i,c_i,d_i)\) & \texttt{\detokenize{a_i;b_i;c_i;d_i}} & binary table event and non-event counts & P15 synthesis binary tables & n/a & \hyperref[def:p15-synthesis-binary-tables]{\ref*{def:p15-synthesis-binary-tables}} & \thepage & Treatment event and nonevent followed by comparator event and nonevent. \\
\phantomsection\label{entry:symbol-registry-02904}\((n_{Ti},n_{Ci},n_i,m_i)\) & \texttt{\detokenize{n_{Ti};n_{Ci};n_i;m_i}} & arm totals study total and total events & P15 synthesis binary tables & n/a & \hyperref[def:p15-synthesis-binary-tables]{\ref*{def:p15-synthesis-binary-tables}} & \thepage & Each arm is nonempty. \\
\phantomsection\label{entry:symbol-registry-02905}\(\delta\) & \texttt{\detokenize{\delta}} & common conditional log odds ratio & P15 synthesis binary tables & n/a & \hyperref[def:p15-synthesis-binary-tables]{\ref*{def:p15-synthesis-binary-tables}} & \thepage & Independent Bernoulli arms with interior event probabilities. \\
\phantomsection\label{entry:symbol-registry-02906}\(A_i\) & \texttt{\detokenize{A_i}} & conditional treatment event-count variable & P15 synthesis binary tables & n/a & \hyperref[def:p15-synthesis-binary-tables]{\ref*{def:p15-synthesis-binary-tables}} & \thepage & Conditioning retains fixed arm sizes and total events. \\
\phantomsection\label{entry:symbol-registry-02907}\(Z_i(\delta)\) & \texttt{\detokenize{Z_i(\delta)}} & finite conditional table normalizer & P15 synthesis binary tables & n/a & \hyperref[def:p15-synthesis-binary-tables]{\ref*{def:p15-synthesis-binary-tables}} & \thepage & Sum over the feasible treatment event counts. \\
\phantomsection\label{entry:symbol-registry-02908}\((E_i,V_i)\) & \texttt{\detokenize{E_i;V_i}} & conditional null event-count mean and variance & P15 synthesis binary tables & n/a & \hyperref[def:p15-synthesis-binary-tables]{\ref*{def:p15-synthesis-binary-tables}} & \thepage & Information is zero for no-events tables. \\
\phantomsection\label{entry:symbol-registry-02909}\((U(\delta),I(\delta))\) & \texttt{\detokenize{U(\delta);I(\delta)}} & conditional log likelihood score and information & P15 synthesis binary tables & n/a & \hyperref[thm:p15-synthesis-score]{\ref*{thm:p15-synthesis-score}} & \thepage & Finite normalizer derivatives supply the exact identities. \\
\phantomsection\label{entry:symbol-registry-02910}\((R,S)\) & \texttt{\detokenize{R;S}} & table-weighted ratio numerator and denominator & P15 synthesis binary tables & n/a & \hyperref[thm:p15-synthesis-table-delta]{\ref*{thm:p15-synthesis-table-delta}} & \thepage & Distinct totals for odds-ratio and risk-ratio estimators. \\
\phantomsection\label{entry:symbol-registry-02911}\(\widehat V_f\) & \texttt{\detokenize{\widehat V_f}} & large-arm delta variance calculation & P15 synthesis binary tables & n/a & \hyperref[thm:p15-synthesis-table-delta]{\ref*{thm:p15-synthesis-table-delta}} & \thepage & Approximate independent-arm uncertainty with fixed number of studies. \\
\phantomsection\label{entry:symbol-registry-02912}\(\Theta_i\) & \texttt{\detokenize{\Theta_i}} & latent random study effect & P15 evidence synthesis & n/a & \hyperref[def:p15-synthesis-random-model]{\ref*{def:p15-synthesis-random-model}} & \thepage & Independent normal random effects in the declared model. \\
\phantomsection\label{entry:symbol-registry-02913}\(u\) & \texttt{\detokenize{u}} & candidate between-study variance & P15 evidence synthesis & n/a & \hyperref[def:p15-synthesis-random-model]{\ref*{def:p15-synthesis-random-model}} & \thepage & Nonnegative real parameter distinct from its estimate. \\
\phantomsection\label{entry:symbol-registry-02914}\(w_i(u)\) & \texttt{\detokenize{w_i(u)}} & inverse total-variance weight & P15 evidence synthesis & n/a & \hyperref[def:p15-synthesis-random-model]{\ref*{def:p15-synthesis-random-model}} & \thepage & Reciprocal of known within-study variance plus candidate heterogeneity. \\
\phantomsection\label{entry:symbol-registry-02915}\(W(u)\) & \texttt{\detokenize{W(u)}} & total candidate inverse-variance weight & P15 evidence synthesis & n/a & \hyperref[def:p15-synthesis-random-model]{\ref*{def:p15-synthesis-random-model}} & \thepage & Sum of the candidate weights. \\
\phantomsection\label{entry:symbol-registry-02916}\(M(u)\) & \texttt{\detokenize{M(u)}} & candidate weighted pooled mean & P15 evidence synthesis & n/a & \hyperref[def:p15-synthesis-random-model]{\ref*{def:p15-synthesis-random-model}} & \thepage & Fixed weights when u is specified; data-dependent when u is estimated. \\
\phantomsection\label{entry:symbol-registry-02917}\(Q(u)\) & \texttt{\detokenize{Q(u)}} & candidate weighted residual dispersion & P15 evidence synthesis & n/a & \hyperref[def:p15-synthesis-random-model]{\ref*{def:p15-synthesis-random-model}} & \thepage & Chi-squared only at the true variance under the exact normal model. \\
\phantomsection\label{entry:symbol-registry-02918}\(\Theta_*\) & \texttt{\detokenize{\Theta_*}} & future underlying effect & P15 evidence synthesis & n/a & \hyperref[thm:p15-synthesis-random-known]{\ref*{thm:p15-synthesis-random-known}} & \thepage & Independent draw from the same modelled effect distribution. \\
\phantomsection\label{entry:symbol-registry-02919}\(C\) & \texttt{\detokenize{C}} & moment dispersion normalization constant & P15 evidence synthesis & n/a & \hyperref[thm:p15-synthesis-q]{\ref*{thm:p15-synthesis-q}} & \thepage & Strictly positive for at least two positive known weights. \\
\phantomsection\label{entry:symbol-registry-02920}\(\widehat\tau_{\rm DL}^2\) & \texttt{\detokenize{\widehat\tau_{\rm DL}^2}} & truncated moment heterogeneity estimate & P15 evidence synthesis & n/a & \hyperref[thm:p15-synthesis-q]{\ref*{thm:p15-synthesis-q}} & \thepage & Nonnegative truncation is not an unbiasedness result. \\
\phantomsection\label{entry:symbol-registry-02921}\(\ell_R(u)\) & \texttt{\detokenize{\ell_R(u)}} & restricted normal log likelihood & P15 evidence synthesis & n/a & \hyperref[thm:p15-synthesis-reml]{\ref*{thm:p15-synthesis-reml}} & \thepage & Mean-eliminating contrast likelihood with the log total-weight term. \\
\phantomsection\label{entry:symbol-registry-02922}\((I_Q^2,I_\tau^2)\) & \texttt{\detokenize{I_Q^2;I_\tau^2}} & Q-based and tau-based inconsistency percentages & P15 evidence synthesis & n/a & \hyperref[def:p15-synthesis-inconsistency]{\ref*{def:p15-synthesis-inconsistency}} & \thepage & Equal with DL but not necessarily with REML. \\
\phantomsection\label{entry:symbol-registry-02923}\(\widetilde v\) & \texttt{\detokenize{\widetilde v}} & typical within-study variance for inconsistency & P15 evidence synthesis & n/a & \hyperref[def:p15-synthesis-inconsistency]{\ref*{def:p15-synthesis-inconsistency}} & \thepage & The ratio of degrees of freedom to dispersion normalization. \\
\phantomsection\label{entry:symbol-registry-02924}\((q_H,s_H^2,s_M^2)\) & \texttt{\detokenize{q_H;s_H^2;s_M^2}} & HKSJ residual factor adjusted variance and Wald variance & P15 evidence synthesis & n/a & \hyperref[def:p15-synthesis-hksj]{\ref*{def:p15-synthesis-hksj}} & \thepage & Their exactness requires separate model assumptions. \\
\phantomsection\label{entry:symbol-registry-02925}\(\Pi_{\rm syn}\) & \texttt{\detokenize{\Pi_{\rm syn}}} & proper probability measure for synthesis parameters & P15 evidence synthesis & n/a & \hyperref[thm:p15-synthesis-bayes]{\ref*{thm:p15-synthesis-bayes}} & \thepage & Prior finite moments guarantee the displayed posterior moments. \\
\phantomsection\label{entry:symbol-registry-02926}\((W_g,M_g)\) & \texttt{\detokenize{W_g;M_g}} & subgroup total weight and pooled effect & P15 evidence synthesis & n/a & \hyperref[thm:p15-synthesis-subgroups]{\ref*{thm:p15-synthesis-subgroups}} & \thepage & Independent group contrasts require disjoint independent evidence. \\
\phantomsection\label{entry:symbol-registry-02927}\(R\) & \texttt{\detokenize{R}} & observation or missingness indicator & P15 synthesis missing outcomes & n/a & \hyperref[def:p15-synthesis-missing]{\ref*{def:p15-synthesis-missing}} & \thepage & Separate from table-ratio totals and conditioned on observed data in MAR. \\
\phantomsection\label{entry:symbol-registry-02928}\((N,n,m,e)\) & \texttt{\detokenize{N;n;m;e}} & allocated observed missing and observed-event participant counts & P15 synthesis missing outcomes & n/a & \hyperref[thm:p15-synthesis-missing-bounds]{\ref*{thm:p15-synthesis-missing-bounds}} & \thepage & Full-sample denominators retained in deterministic binary bounds. \\
\phantomsection\label{entry:symbol-registry-02929}\((M_{-j},B_i)\) & \texttt{\detokenize{M_{-j};B_i}} & deleted fixed-weight mean and bounded study bias magnitudes & P15 evidence synthesis & n/a & \hyperref[thm:p15-synthesis-sensitivity]{\ref*{thm:p15-synthesis-sensitivity}} & \thepage & Re-estimated heterogeneity requires a fresh deletion calculation. \\
\phantomsection\label{entry:symbol-registry-02930}\((p_T,p_C)\) & \texttt{\detokenize{p_T;p_C}} & treatment and comparator event probabilities & P15 evidence synthesis & n/a & \hyperref[def:p15-synthesis-effect-scales]{\ref*{def:p15-synthesis-effect-scales}} & \thepage & Same defined outcome and horizon in an effect comparison. \\
\phantomsection\label{entry:symbol-registry-02931}\(z_p\) & \texttt{\detokenize{z_p}} & standard-normal quantile & P15 evidence synthesis & n/a & \hyperref[thm:p15-synthesis-linear]{\ref*{thm:p15-synthesis-linear}} & \thepage & Exact interval multiplier only under its stated normal model. \\
\phantomsection\label{entry:symbol-registry-02932}\(q_{p,d}\) & \texttt{\detokenize{q_{p,d}}} & chi-squared quantile & P15 evidence synthesis & n/a & \hyperref[thm:p15-synthesis-q]{\ref*{thm:p15-synthesis-q}} & \thepage & Positive integer degrees of freedom in the exact known-variance Q set. \\
\phantomsection\label{entry:symbol-registry-02933}\(t_{k-1,1-\alpha/2}\) & \texttt{\detokenize{t_{k-1,1-\alpha/2}}} & Student two-sided interval quantile & P15 evidence synthesis & n/a & \hyperref[def:p15-synthesis-hksj]{\ref*{def:p15-synthesis-hksj}} & \thepage & Student degrees of freedom are number of studies minus one. \\
\phantomsection\label{entry:symbol-registry-02934}\(\lambda_{gi}\) & \texttt{\detokenize{\lambda_{gi}}} & positive limiting arm sample-size fraction & P15 synthesis binary tables & n/a & \hyperref[thm:p15-synthesis-table-delta]{\ref*{thm:p15-synthesis-table-delta}} & \thepage & Fixed-study large-arm asymptotic regime. \\
\phantomsection\label{entry:symbol-registry-02935}\(\widehat p_{gi}\) & \texttt{\detokenize{\widehat p_{gi}}} & observed arm event proportion & P15 synthesis binary tables & n/a & \hyperref[thm:p15-synthesis-table-delta]{\ref*{thm:p15-synthesis-table-delta}} & \thepage & The arm-indexed T and C forms define the delta plug-in variance. \\
\phantomsection\label{entry:symbol-registry-02936}\(r\) & \texttt{\detokenize{r}} & correlation of two effect estimates & P15 synthesis dependence & n/a & \hyperref[thm:p15-synthesis-dependence]{\ref*{thm:p15-synthesis-dependence}} & \thepage & Not the earlier common risk-ratio parameter; its typed role is declared in this result. \\
\phantomsection\label{entry:symbol-registry-02937}\(c\) & \texttt{\detokenize{c}} & common unknown variance multiplier in fixed-shape normal sampling & P15 fixed-shape synthesis & n/a & \hyperref[thm:p15-synthesis-fixed-shape-t]{\ref*{thm:p15-synthesis-fixed-shape-t}} & \thepage & Distinct from the earlier affine measurement scale factor. \\
\phantomsection\label{entry:symbol-registry-02938}\(C_i\) & \texttt{\detokenize{C_i}} & comparator event-count variable & P15 synthesis binary tables & n/a & \hyperref[annex:p15.synthesis.score]{\ref*{annex:p15.synthesis.score}} & \thepage & Conditioned together with treatment count on the total event margin. \\
\phantomsection\label{entry:symbol-registry-02939}\(\widehat\delta_P\) & \texttt{\detokenize{\widehat\delta_P}} & conditional binary Peto one-step log odds estimate & P15 synthesis binary tables & n/a & \hyperref[def:p15-synthesis-binary-tables]{\ref*{def:p15-synthesis-binary-tables}} & \thepage & Score divided by null information with approximate uncertainty. \\
\phantomsection\label{entry:symbol-registry-02940}\(Y,\mathcal Y\) & \texttt{\detokenize{Y,}\allowbreak\detokenize{\mathcal Y}} & numerical study result and its finite support in the stipulated selection model & P15 missing evidence & object & \hyperref[def:p15-missing-selection]{\ref*{def:p15-missing-selection}} & \thepage & The model describes available-result distributions; it does not by itself identify a true treatment effect. \\
\phantomsection\label{entry:symbol-registry-02941}\(R\) & \texttt{\detokenize{R}} & indicator that the particular result is available to the synthesis & P15 missing evidence & object & \hyperref[def:p15-missing-selection]{\ref*{def:p15-missing-selection}} & \thepage & A result-level availability indicator; distinct from participant-level missingness within a study. \\
\phantomsection\label{entry:symbol-registry-02942}\(p(y),s(y),q(y),\rho\) & \texttt{\detokenize{p(y),}\allowbreak\detokenize{s(y),}\allowbreak\detokenize{q(y),}\allowbreak\detokenize{\rho}} & full result law availability probability observed conditional law and overall availability mass & P15 missing evidence & object & \hyperref[def:p15-missing-selection]{\ref*{def:p15-missing-selection}} & \thepage & Positive overall availability is required; inverse recovery additionally needs specified positive selection probabilities. \\
\phantomsection\label{entry:symbol-registry-02943}\(G,G^c,\lambda\) & \texttt{\detokenize{G,}\allowbreak\detokenize{G^c,}\allowbreak\detokenize{\lambda}} & inception cohort complementary studies and cohort weight in a declared full target & P15 missing evidence & object & \hyperref[thm:p15-missing-cohort]{\ref*{thm:p15-missing-cohort}} & \thepage & Closed-frame study identification is distinct from result availability and transport to excluded research. \\
\phantomsection\label{entry:symbol-registry-02944}\(b_i,\theta_G,\theta_{G^c}\) & \texttt{\detokenize{b_i,}\allowbreak\detokenize{\theta_G,}\allowbreak\detokenize{\theta_{G^c}}} & predeclared full-target weights and normalized within-cohort complementary effects & P15 missing evidence & object & \hyperref[thm:p15-missing-cohort]{\ref*{thm:p15-missing-cohort}} & \thepage & Positive fixed weights sum to one; the complement mean is used only when its total mass is positive. \\
\phantomsection\label{entry:symbol-registry-02945}\(\checkmark,\sim,?,X\) & \texttt{\detokenize{\checkmark,}\allowbreak\detokenize{\sim,}\allowbreak\detokenize{?,}\allowbreak\detokenize{X}} & four typed availability classifications in the selected ROB-ME results matrix & P15 missing evidence & object & \hyperref[def:p15-missing-results-matrix]{\ref*{def:p15-missing-results-matrix}} & \thepage & Available result; unrelated absence; eligible-result generation unclear; absence likely related to findings. \\
\phantomsection\label{entry:symbol-registry-02946}\(W_o,M_o\) & \texttt{\detokenize{W_o,}\allowbreak\detokenize{M_o}} & total fixed available-result weight and its available weighted mean & P15 missing evidence & object & \hyperref[def:p15-missing-weight-bounds]{\ref*{def:p15-missing-weight-bounds}} & \thepage & The total available weight is positive. \\
\phantomsection\label{entry:symbol-registry-02947}\(w_j,L_j,U_j,W_m\) & \texttt{\detokenize{w_j,}\allowbreak\detokenize{L_j,}\allowbreak\detokenize{U_j,}\allowbreak\detokenize{W_m}} & specified missing-result weights effect-range limits and total missing weight & P15 missing evidence & object & \hyperref[def:p15-missing-weight-bounds]{\ref*{def:p15-missing-weight-bounds}} & \thepage & Assumed sensitivity inputs; participant counts alone do not determine these inverse-variance weights. \\
\phantomsection\label{entry:symbol-registry-02948}\(M_m,d\) & \texttt{\detokenize{M_m,}\allowbreak\detokenize{d}} & weighted missing-result mean and an explicit decision threshold & P15 missing evidence & object & \hyperref[thm:p15-missing-weight-bounds]{\ref*{thm:p15-missing-weight-bounds}} & \thepage & The missing mean is used only when total missing weight is positive; the threshold comparison is deterministic. \\
\phantomsection\label{entry:symbol-registry-02949}\(v_L,v_S,y_L,y_S,M(u)\) & \texttt{\detokenize{v_L,}\allowbreak\detokenize{v_S,}\allowbreak\detokenize{y_L,}\allowbreak\detokenize{y_S,}\allowbreak\detokenize{M(u)}} & two study variances fixed effects and the summary under a common heterogeneity allowance & P15 missing evidence & object & \hyperref[thm:p15-missing-small-weight]{\ref*{thm:p15-missing-small-weight}} & \thepage & The less precise variance exceeds the more precise variance; increasing the allowance changes relative weights. \\
\phantomsection\label{entry:symbol-registry-02950}\(\epsilon\) & \texttt{\detokenize{\epsilon}} & positive rounding half-width in the funnel compatibility proof & \hyperref[annex:p15.missing.funnel]{\ref*{annex:p15.missing.funnel}} & object & \hyperref[annex:p15.missing.funnel]{\ref*{annex:p15.missing.funnel}} & \thepage & Chosen below one fortieth so displayed intervals remain distinct and the selected small-study intervals exceed the reporting threshold. \\
\phantomsection\label{entry:symbol-registry-02951}\(\eta,a\) & \texttt{\detokenize{\eta,}\allowbreak\detokenize{a}} & reported-result mass and an unreported value in the zero-selection boundary & \hyperref[annex:p15.missing.selection]{\ref*{annex:p15.missing.selection}} & object & \hyperref[annex:p15.missing.selection]{\ref*{annex:p15.missing.selection}} & \thepage & The observed result law remains fixed while the full mean changes. \\
\phantomsection\label{entry:symbol-registry-03742}\(\Omega,\mathcal F\) & \texttt{\detokenize{\Omega,}\allowbreak\detokenize{\mathcal F}} & possibility space and algebra of events & P15 subjective-probability-and-dynamics & meta & \hyperref[def:p15-rat41-probabilism]{\ref*{def:p15-rat41-probabilism}} & \thepage & Source uses Omega for possibilities; algebra closure is explicit. The source does not assume countable additivity. \\
\phantomsection\label{entry:symbol-registry-03743}\(P,C,C_{\rm old},C_{\rm new},C_t\) & \texttt{\detokenize{P,}\allowbreak\detokenize{C,}\allowbreak\detokenize{C_{\rm old},}\allowbreak\detokenize{C_{\rm new},}\allowbreak\detokenize{C_t}} & probability and interpreted credence functions with stage or time indices & P15 subjective-probability-and-dynamics & meta & \hyperref[def:p15-rat41-probabilism]{\ref*{def:p15-rat41-probabilism}} & \thepage & P can be objectively or subjectively interpreted; C denotes credence. Stage-indexing does not itself justify an update norm. \\
\phantomsection\label{entry:symbol-registry-03744}\(X,Y,A,E,F,E_i\) & \texttt{\detokenize{X,}\allowbreak\detokenize{Y,}\allowbreak\detokenize{A,}\allowbreak\detokenize{E,}\allowbreak\detokenize{F,}\allowbreak\detokenize{E_i}} & events or proposition events in the source probability framework & P15 subjective-probability-and-dynamics & meta & \hyperref[def:p15-rat41-probabilism]{\ref*{def:p15-rat41-probabilism}} & \thepage & E\_i are partition cells where so declared; F is the first column event in the Jeffrey example. X in the Principal Principle is a chance proposition. These are not random-variable names in this local treatment. \\
\phantomsection\label{entry:symbol-registry-03745}\(A_j,m,v(w),v_j(w),c,K\) & \texttt{\detokenize{A_j,}\allowbreak\detokenize{m,}\allowbreak\detokenize{v,}\allowbreak\detokenize{v_j,}\allowbreak\detokenize{c,}\allowbreak\detokenize{K}} & finite assessed event family truth vectors credence vector and coherent convex representation set & P15 subjective-probability-and-dynamics & meta & \hyperref[def:p15-rat41-brier]{\ref*{def:p15-rat41-brier}} & \thepage & Editorial finite realization of the source Brier-dominance argument; K is the convex hull of world truth vectors. Coordinates need not sum to one unless events form a partition. \\
\phantomsection\label{entry:symbol-registry-03746}\(S(c,w)\) & \texttt{\detokenize{S}} & finite summed Brier inaccuracy & P15 subjective-probability-and-dynamics & meta & \hyperref[def:p15-rat41-brier]{\ref*{def:p15-rat41-brier}} & \thepage & Lower values are better; the score is squared Euclidean distance from a truth vector. \\
\phantomsection\label{entry:symbol-registry-03747}\(q,x,r_i,n\) & \texttt{\detokenize{q,}\allowbreak\detokenize{x,}\allowbreak\detokenize{r_i,}\allowbreak\detokenize{n}} & specified event probability report value new partition weights and finite partition size & P15 subjective-probability-and-dynamics & meta & \hyperref[thm:p15-rat41-brier]{\ref*{thm:p15-rat41-brier}} & \thepage & q and x are scalars in the properness identity; q is explicitly a probability vector in the later repair example. r\_i sum to one; n is a positive finite count or a later-time increment when stated. \\
\phantomsection\label{entry:symbol-registry-03748}\(H(p),I(p),p_i\) & \texttt{\detokenize{H,}\allowbreak\detokenize{I,}\allowbreak\detokenize{p_i}} & finite entropy source negative-entropy information convention and probability coordinates & P15 subjective-probability-and-dynamics & meta & \hyperref[thm:p15-rat41-indifference]{\ref*{thm:p15-rat41-indifference}} & \thepage & H is negative sum p\_i log2 p\_i; source I is its negative. Zero-log-zero convention is explicit; no coordinate-invariant continuous prior is asserted. \\
\phantomsection\label{entry:symbol-registry-03749}\(Q,Q^*,q_j,p_j,D(Q,P)\) & \texttt{\detokenize{Q,}\allowbreak\detokenize{Q^*,}\allowbreak\detokenize{q_j,}\allowbreak\detokenize{p_j,}\allowbreak\detokenize{D}} & new probability mass constrained Jeffrey law coordinates and oriented relative entropy & P15 subjective-probability-and-dynamics & meta & \hyperref[thm:p15-rat41-minimal-update]{\ref*{thm:p15-rat41-minimal-update}} & \thepage & D uses new probability as first argument and old probability as second; support failure gives infinity. Q* is unique KL and Hellinger minimizer but not necessarily unique total-variation minimizer. \\
\phantomsection\label{entry:symbol-registry-03750}\(y,c_s,s,\|c-y\|\) & \texttt{\detokenize{y,}\allowbreak\detokenize{c_s,}\allowbreak\detokenize{s,}\allowbreak\detokenize{\|\cdot\|}} & Euclidean projection segment parameter and Euclidean norm for finite coherence repair & P15 subjective-probability-and-dynamics & meta & \hyperref[thm:p15-rat41-graded]{\ref*{thm:p15-rat41-graded}} & \thepage & y is the unique projection onto K; 0<s<=1. Squared distance and distance are distinguished; stake normalization uses unsquared Euclidean norm. \\
\phantomsection\label{entry:symbol-registry-03751}\(h,z,u,b,s_0,\bar c\) & \texttt{\detokenize{h,}\allowbreak\detokenize{z,}\allowbreak\detokenize{u,}\allowbreak\detokenize{b,}\allowbreak\detokenize{s_0,}\allowbreak\detokenize{\bar c}} & finite proof displacement coherent comparison segment parameter comparison or stake vector and normalization of a positive reference vector & P15 subjective-probability-and-dynamics & meta & \hyperref[annex:p15.rat41.brier]{\ref*{annex:p15.rat41.brier}} & \thepage & h=c-y; z in K; b is a comparison credence vector in the Brier proof and explicitly a unit-bounded stake vector in the graded proof. s0=0.7 and barred c=(3/7,4/7) only in the explicitly extended relative-entropy repair, not a source-exact numerical correction. \\
\phantomsection\label{entry:symbol-registry-03752}\(w,W,\mathbf1_X\) & \texttt{\detokenize{w,}\allowbreak\detokenize{W,}\allowbreak\detokenize{\mathbf1_X}} & world argument world random element and event indicator & P15 subjective-probability-and-dynamics & meta & \hyperref[def:p15-rat41-brier]{\ref*{def:p15-rat41-brier}} & \thepage & W has the finite probability masses chosen in the Brier proof. The indicator equals one on X and zero elsewhere; minus that indicator is the corrected zero-price semi-Dutch-book payoff. \\
\phantomsection\label{entry:symbol-registry-03753}\(t,k,a\) & \texttt{\detokenize{t,}\allowbreak\detokenize{k,}\allowbreak\detokenize{a}} & time integer threshold index and explicitly local scalar or probability-vector parameter & P15 subjective-probability-and-dynamics & meta & \hyperref[thm:p15-rat41-positive-atoms]{\ref*{thm:p15-rat41-positive-atoms}} & \thepage & t is time in the deference formulas; k is a positive integer reciprocal-mass threshold in the singleton proof. a is a scalar complement-pair coordinate in the repair theorem and explicitly a probability vector in the finite Gibbs proof. These typed local uses are not logical individual variables. \\
\phantomsection\label{entry:symbol-registry-03775}\(A\mathbin{\&}B,A\lor B,\sim A\) & \texttt{\detokenize{\mathbin{\&};\lor;\sim}} & source Boolean conjunction disjunction and negation & P15 Rationality 2021 chapter 4.4 & n/a & \hyperref[def:p15-rat44-events]{\ref*{def:p15-rat44-events}} & \thepage & Events/propositions retain Boolean truth conditions; probability values are not fuzzy truth values. \\
\phantomsection\label{entry:symbol-registry-03776}\(B\mid A\) & \texttt{\detokenize{B\mid A}} & nonnested conditional event & P15 Rationality 2021 chapter 4.4 & n/a & \hyperref[def:p15-rat44-events]{\ref*{def:p15-rat44-events}} & \thepage & True on A and B; false on A and not B; void outside A; antecedent must be possible even if its assessed probability is zero. \\
\phantomsection\label{entry:symbol-registry-03777}\(\Omega\) & \texttt{\detokenize{\Omega}} & finite logically possible constituent space & P15 Rationality 2021 chapter 4.4 & n/a & \hyperref[def:p15-rat44-events]{\ref*{def:p15-rat44-events}} & \thepage & Nonempty; logical dependencies remove impossible constituents; numerical probabilities may give possible constituents zero weight. \\
\phantomsection\label{entry:symbol-registry-03778}\(p_i,s_i,G_s,H_s\) & \texttt{\detokenize{p_i;s_i;G_s;H_s}} & conditional prices signed stakes portfolio gain and active conditioning union & P15 Rationality 2021 chapter 4.4 & n/a & \hyperref[def:p15-rat44-coherence]{\ref*{def:p15-rat44-coherence}} & \thepage & Two-sided precise prices; active union uses nonzero stakes; gain vanishes for a called-off conditional bet. \\
\phantomsection\label{entry:symbol-registry-03779}\(\mathbf1_E\) & \texttt{\detokenize{\mathbf1_E}} & event indicator & P15 Rationality 2021 chapter 4.4 & n/a & \hyperref[def:p15-rat44-coherence]{\ref*{def:p15-rat44-coherence}} & \thepage & Value one on E and zero elsewhere; conditional gain is 1\_H times the centered event indicator. \\
\phantomsection\label{entry:symbol-registry-03780}\(Q,\operatorname{conv}Q\) & \texttt{\detokenize{Q;\operatorname{conv}}} & possible binary evaluation vectors and their convex hull & P15 Rationality 2021 chapter 4.4 & n/a & \hyperref[thm:p15-rat44-coherence-method]{\ref*{thm:p15-rat44-coherence-method}} & \thepage & Constituent weights represent coherent unconditional assessments; impossible truth patterns must not be reintroduced. \\
\phantomsection\label{entry:symbol-registry-03781}\(\lambda_\omega,\mu\) & \texttt{\detokenize{\lambda_\omega;\mu}} & normalized constituent weights or recorded coherent-conditioning stratum & P15 Rationality 2021 chapter 4.4 proof construction & n/a & \hyperref[thm:p15-rat44-coherence-method]{\ref*{thm:p15-rat44-coherence-method}} & \thepage & Later strata determine prices only on conditioning events null at every earlier stratum; not an ordinary positive-mass assumption. \\
\phantomsection\label{entry:symbol-registry-03782}\(p(E\mid H),p(E)\) & \texttt{\detokenize{p(E\mid H);p(E)}} & primitive conditional and unconditional assessments & P15 Rationality 2021 chapter 4.4 & n/a & \hyperref[def:p15-rat44-coherence]{\ref*{def:p15-rat44-coherence}} & \thepage & Unconditional p(E) is p(E|top); a ratio identity follows only when p(H)>0; zero probability is not impossibility. \\
\phantomsection\label{entry:symbol-registry-03783}\(h\) & \texttt{\detokenize{h}} & high-assessment threshold for the selected connexive reading & P15 Rationality 2021 chapter 4.4 connexivity & n/a & \hyperref[def:p15-rat44-connexive-reading]{\ref*{def:p15-rat44-connexive-reading}} & \thepage & Strictly greater than one-half; narrow-scope negation changes the conditional's consequent; outer implications concern acceptance predicates. \\
\phantomsection\label{entry:symbol-registry-03784}\(x_L,x_U,y_L,y_U\) & \texttt{\detokenize{x_L;x_U;y_L;y_U}} & premise interval endpoints & P15 Rationality 2021 chapter 4.4 propagation & n/a & \hyperref[thm:p15-rat44-propagation]{\ref*{thm:p15-rat44-propagation}} & \thepage & Closed constraints on coherent precise selections; not automatically a full-domain lower-prevision assessment. \\
\phantomsection\label{entry:symbol-registry-03785}\(x_1,x_2,x_3,x_4\) & \texttt{\detokenize{x_1;x_2;x_3;x_4}} & four ordinary joint case probabilities in PTTT & P15 Rationality 2021 chapter 4.4 empirical comparison & n/a & \hyperref[sec:p15-rationality-4-4]{\ref*{sec:p15-rationality-4-4}} & \thepage & Nonnegative and sum to one; conditional-ratio task prediction requires x\_1+x\_2>0; chapter's approximately 80 percent concerns late-task responses. \\
\phantomsection\label{entry:symbol-registry-03903}\(g\) & \texttt{\detokenize{g}} & normalized monotone event-confidence capacity & P15 rationality4.7 confidence interface & n/a & \hyperref[def:p15-r47-capacity]{\ref*{def:p15-r47-capacity}} & \thepage & Not a many-valued truth function or an assumed additive probability. \\
\phantomsection\label{entry:symbol-registry-03904}\(N,\Pi,\pi,\iota\) & \texttt{\detokenize{N,}\allowbreak\detokenize{\Pi,}\allowbreak\detokenize{\pi,}\allowbreak\detokenize{\iota}} & necessity and possibility set functions with singleton possibility and exclusion strengths & P15 rationality4.7 possibility interface & n/a & \hyperref[thm:p15-r47-threshold]{\ref*{thm:p15-r47-threshold}} & \thepage & Finite representation; empty minima and maxima are explicitly fixed. \\
\phantomsection\label{entry:symbol-registry-03905}\(\kappa\) & \texttt{\detokenize{\kappa}} & normalized negative finite-state ranking & P15 rationality4.7 ranking comparison & n/a & \hyperref[def:p15-r47-ranking]{\ref*{def:p15-r47-ranking}} & \thepage & Static image is Pi(A)=2\^{}-kappa(A) and N(A)=1-2\^{}-kappa(complement); not a probability conditional update. \\
\phantomsection\label{entry:symbol-registry-03906}\(\Box^\lambda p\) & \texttt{\detokenize{\Box^\lambda p}} & graded meta-atom asserting a lower necessity bound on an objective formula & P15 rationality4.7 graded acceptance & n/a & \hyperref[def:p15-r47-gpl]{\ref*{def:p15-r47-gpl}} & \thepage & Outer Boolean language has no objective atoms or nested modalities; positive necessity grades only. \\
\phantomsection\label{entry:symbol-registry-03907}\(s(\lambda)\) & \texttt{\detokenize{s(\lambda)}} & next lower grade in the finite necessity scale & P15 rationality4.7 graded acceptance & n/a & \hyperref[def:p15-r47-gpl]{\ref*{def:p15-r47-gpl}} & \thepage & Predecessor shift fixes the strict-to-nonstrict dual bound; not an accessibility successor. \\
\phantomsection\label{entry:symbol-registry-03908}\(N_\lambda,\Pi_\lambda\) & \texttt{\detokenize{N_\lambda,}\allowbreak\detokenize{\Pi_\lambda}} & author2017 uniform-grid necessity and possibility-bound meta-atom notation & P15 rationality4.7 GPL source comparison & n/a & \hyperref[def:p15-r47-gpl]{\ref*{def:p15-r47-gpl}} & \thepage & Declared separately from the2021 boxed notation and its threshold labels. \\
\phantomsection\label{entry:symbol-registry-03909}\(m,\operatorname{Bel},\operatorname{Pl}\) & \texttt{\detokenize{m,}\allowbreak\detokenize{\operatorname{Bel},}\allowbreak\detokenize{\operatorname{Pl}}} & finite nonempty-focal mass and its containment/intersection event functions & P15 rationality4.7 evidence interface & n/a & \hyperref[def:p15-r47-mass]{\ref*{def:p15-r47-mass}} & \thepage & A mass is on information sets; it does not allocate a point within each focal set. \\
\phantomsection\label{entry:symbol-registry-03910}\(m_1\otimes m_2,K\) & \texttt{\detokenize{m_1\otimes m_2,}\allowbreak\detokenize{K}} & normalized independent-random-set combination and its nonconflict probability & P15 rationality4.7 evidence combination & n/a & \hyperref[def:p15-r47-combination]{\ref*{def:p15-r47-combination}} & \thepage & Combination requires K>0 and fixes empty mass separately. \\
\phantomsection\label{entry:symbol-registry-03911}\(\operatorname{BetP}\) & \texttt{\detokenize{\operatorname{BetP}}} & finite pignistic probability projection & P15 rationality4.7 decision comparison & n/a & \hyperref[thm:p15-r47-pignistic]{\ref*{thm:p15-r47-pignistic}} & \thepage & Not an injective recovery of an original belief function. \\
\phantomsection\label{entry:symbol-registry-03912}\(LP,UP,\mathcal P_{\mathcal G}\) & \texttt{\detokenize{LP,}\allowbreak\detokenize{UP,}\allowbreak\detokenize{\mathcal P_{\mathcal G}}} & lower and upper gamble previsions and compatible precise extensions & P15 rationality4.7 betting comparison & n/a & \hyperref[sec:p15-r47-lower-previsions]{\ref*{sec:p15-r47-lower-previsions}} & \thepage & Uses the actual bounded-gamble P15 interface; no latent precise probability is assumed. \\
\phantomsection\label{entry:symbol-registry-03913}\(\operatorname{lik}\) & \texttt{\detokenize{\operatorname{lik}}} & likelihood across supplied parameter models and its maxitive extension & P15 rationality4.7 likelihood comparison & n/a & \hyperref[thm:p15-r47-likelihood]{\ref*{thm:p15-r47-likelihood}} & \thepage & Not posterior parameter probability; positive maximum is required for normalization. \\
\phantomsection\label{entry:symbol-registry-03914}\(F,n,g(A\mid B)\) & \texttt{\detokenize{F,}\allowbreak\detokenize{n,}\allowbreak\detokenize{g(A\mid B)}} & conjunction composition and complement transform for source-reported conditional confidence & P15 rationality4.7 Cox comparison & n/a & \hyperref[def:p15-r47-cox]{\ref*{def:p15-r47-cox}} & \thepage & These are reported functional postulates and regularity conditions; no universal Cox theorem is promoted from them. \\
\phantomsection\label{entry:symbol-registry-03915}\(W,\mathcal P(W)\) & \texttt{\detokenize{W;\mathcal P}} & possibility space and full proposition algebra & P15 Ranking Theory objects & n/a & \hyperref[def:p15-rat53-ranking]{\ref*{def:p15-rat53-ranking}} & \thepage & Nonempty W; full intersection closure requires finite W or the stated world-minimum representation. \\
\phantomsection\label{entry:symbol-registry-03916}\(\neg A,A\land B,A\lor B\) & \texttt{\detokenize{\neg;\land;\lor}} & complement intersection and union of proposition events & P15 Ranking Theory objects & n/a & \hyperref[def:p15-rat53-ranking]{\ref*{def:p15-rat53-ranking}} & \thepage & The source uses logical notation for subsets. \\
\phantomsection\label{entry:symbol-registry-03917}\(\kappa\) & \texttt{\detokenize{\kappa}} & source negative event-ranking measuring disbelief & P15 Ranking Theory & n/a & \hyperref[def:p15-rat53-ranking]{\ref*{def:p15-rat53-ranking}} & \thepage & Values are nonnegative integers or infinity; zero means not disbelieved rather than automatically believed. \\
\phantomsection\label{entry:symbol-registry-03918}\(\beta\) & \texttt{\detokenize{\beta}} & source positive rank of a proposition & P15 Ranking Theory & n/a & \hyperref[def:p15-rat53-ranking]{\ref*{def:p15-rat53-ranking}} & \thepage & Negative rank of the complement; minimum law for conjunction. \\
\phantomsection\label{entry:symbol-registry-03919}\(\tau\) & \texttt{\detokenize{\tau}} & source signed two-sided rank & P15 Ranking Theory & n/a & \hyperref[def:p15-rat53-ranking]{\ref*{def:p15-rat53-ranking}} & \thepage & Complement disbelief minus event disbelief; no infinity-minus-infinity under the negation law. \\
\phantomsection\label{entry:symbol-registry-03920}\(\mathcal B_z,z\) & \texttt{\detokenize{\mathcal B_z;z}} & threshold belief set and finite nonnegative threshold & P15 Ranking Theory objects & n/a & \hyperref[def:p15-rat53-ranking]{\ref*{def:p15-rat53-ranking}} & \thepage & Belief is strict complement rank greater than z. \\
\phantomsection\label{entry:symbol-registry-03921}\(r(w),M_z\) & \texttt{\detokenize{r;M_z}} & world rank and worlds no higher than the threshold & P15 Ranking Theory world representation & n/a & \hyperref[thm:p15-rat53-rank-laws]{\ref*{thm:p15-rat53-rank-laws}} & \thepage & r is singleton rank; M\_z is the nonempty common core under the finite or stronger representation condition. \\
\phantomsection\label{entry:symbol-registry-03922}\(\kappa(B\mid A),\beta(B\mid A),\tau(B\mid A)\) & \texttt{\detokenize{\kappa(\cdot\mid\cdot);\beta(\cdot\mid\cdot);\tau(\cdot\mid\cdot)}} & conditional disbelief positive belief and signed ranks & P15 Ranking Theory conditioning & n/a & \hyperref[def:p15-rat53-conditionals]{\ref*{def:p15-rat53-conditionals}} & \thepage & A must have finite negative rank; it need not have rank zero. \\
\phantomsection\label{entry:symbol-registry-03923}\(B,P,F\) & \texttt{\detokenize{B;P;F}} & bird penguin and flight propositions & P15 Ranking Theory Tweety example & n/a & \hyperref[tab:p15-rat53-tweety]{\ref*{tab:p15-rat53-tweety}} & \thepage & Local events; P is not a probability measure here. \\
\phantomsection\label{entry:symbol-registry-03924}\(\kappa',\kappa''\) & \texttt{\detokenize{\kappa';\kappa''}} & Tweety states after bird and penguin evidence & P15 Ranking Theory Tweety example & n/a & \hyperref[tab:p15-rat53-tweety]{\ref*{tab:p15-rat53-tweety}} & \thepage & All eight ranks in each source table retained. \\
\phantomsection\label{entry:symbol-registry-03925}\(\kappa_E^{a,b},a,b,n\) & \texttt{\detokenize{\kappa_E^{a,b};a;b;n}} & two-cell update new cell ranks and target firmness & P15 Ranking Theory learning & n/a & \hyperref[def:p15-rat53-soft-update]{\ref*{def:p15-rat53-soft-update}} & \thepage & Both old cells finite; minimum target zero; n is not an unqualified signed-rank increment. \\
\phantomsection\label{entry:symbol-registry-03926}\(\Pi,N\) & \texttt{\detokenize{\Pi;N}} & static finite possibility and necessity transforms & P15 Ranking Theory comparison & n/a & \hyperref[thm:p15-rat53-possibility]{\ref*{thm:p15-rat53-possibility}} & \thepage & Pi(A)=2 to minus kappa(A); N(A)=1 minus Pi(complement A). \\
\phantomsection\label{entry:symbol-registry-03927}\(P_\varepsilon,\varepsilon\) & \texttt{\detokenize{P_\varepsilon;\varepsilon}} & editorial finite probability family and small positive parameter & P15 Ranking Theory probability bridge & n/a & \hyperref[thm:p15-rat53-possibility]{\ref*{thm:p15-rat53-possibility}} & \thepage & Weights proportional to epsilon raised to world rank; a limiting-exponent bridge not one unique probability. \\
\phantomsection\label{entry:symbol-registry-03928}\(A\mathrel{|\!\sim}B\) & \texttt{\detokenize{\mathrel{|\!\sim}}} & ranked nonmonotonic consequence & P15 Ranking Theory AI interface & relation & \hyperref[sec:p15-rationality-5-3]{\ref*{sec:p15-rationality-5-3}} & \thepage & Finite-rank A requires kappa(A and not B) greater than kappa(A). \\
\phantomsection\label{entry:symbol-registry-03960}\(P,P_A\) & \texttt{\detokenize{P,}\allowbreak\detokenize{P_A}} & credal probability and positive-support conditioning & P15 evidential-relevance & meta & \hyperref[def:p15-r43-polarity]{\ref*{def:p15-r43-polarity}} & \thepage & Finite-additive event model; P\_A requires P(A)>0. Neutral relevance at a null event is not a defined ratio conditional. \\
\phantomsection\label{entry:symbol-registry-03961}\(A,B,C,E,H,D,G,J,E_i,H_i\) & \texttt{\detokenize{A,}\allowbreak\detokenize{B,}\allowbreak\detokenize{C,}\allowbreak\detokenize{E,}\allowbreak\detokenize{H,}\allowbreak\detokenize{D,}\allowbreak\detokenize{G,}\allowbreak\detokenize{J,}\allowbreak\detokenize{E_i,}\allowbreak\detokenize{H_i}} & proposition events with explicitly assigned source target and observation roles & P15 evidential-relevance & meta & \hyperref[def:p15-r43-polarity]{\ref*{def:p15-r43-polarity}} & \thepage & Juxtaposition AB means intersection/conjunction; not random-variable multiplication. D names a data event only in the old-evidence section, distinct from the subscript D measure type. J is the disjoint added target in the characterization proof. \\
\phantomsection\label{entry:symbol-registry-03962}\(\Delta_P(A,B)\) & \texttt{\detokenize{\Delta_P}} & event covariance expression for polarity & P15 evidential-relevance & meta & \hyperref[def:p15-r43-polarity]{\ref*{def:p15-r43-polarity}} & \thepage & Editorial abbreviation P(AB)-P(A)P(B); not a posterior change or a source-defined extra measure. \\
\phantomsection\label{entry:symbol-registry-03963}\(\mathrm{pos}_P,\mathrm{neg}_P,\perp_P,\mathrm{expos}_P,\mathrm{exneg}_P\) & \texttt{\detokenize{pos,}\allowbreak\detokenize{neg,}\allowbreak\detokenize{\perp,}\allowbreak\detokenize{expos,}\allowbreak\detokenize{exneg}} & Merin's signed qualitative and extreme relevance relations & P15 evidential-relevance & meta & \hyperref[def:p15-r43-polarity]{\ref*{def:p15-r43-polarity}} & \thepage & Extreme positivity requires positive evidence mass, prior target below one and conditional target one. Conditional irrelevance uses both defined strata for plus/minus notation. \\
\phantomsection\label{entry:symbol-registry-03964}\(\Yright,R_T(A\Yright B)_P,\iota_T\) & \texttt{\detokenize{\Yright,}\allowbreak\detokenize{R_T,}\allowbreak\detokenize{\iota_T}} & source-target signpost relevance function and neutral value & P15 evidential-relevance & meta & \hyperref[def:p15-r43-measures]{\ref*{def:p15-r43-measures}} & \thepage & Source's forked separator is supplied by existing stmaryrd; it is not logical entailment. Amount symmetry is separate from polarity symmetry. \\
\phantomsection\label{entry:symbol-registry-03965}\(R_D,R_Q,R_L,R_{QL},R_{LL},R_{KO},R_Z,R_C\) & \texttt{\detokenize{R_D,}\allowbreak\detokenize{R_Q,}\allowbreak\detokenize{R_L,}\allowbreak\detokenize{R_{QL},}\allowbreak\detokenize{R_{LL},}\allowbreak\detokenize{R_{KO},}\allowbreak\detokenize{R_Z,}\allowbreak\detokenize{R_C}} & source measure types difference quotient likelihood logarithms bounded transforms and covariance & P15 evidential-relevance & meta & \hyperref[def:p15-r43-measures]{\ref*{def:p15-r43-measures}} & \thepage & Boundary completions and log domains are explicit. Natural logarithms are used; KO equals tanh of half log L, not half log LL. \\
\phantomsection\label{entry:symbol-registry-03966}\(e,h,q,a,b,p,a_i,b_i,c_i,d_i,\delta\) & \texttt{\detokenize{e,}\allowbreak\detokenize{h,}\allowbreak\detokenize{q,}\allowbreak\detokenize{a,}\allowbreak\detokenize{b,}\allowbreak\detokenize{p,}\allowbreak\detokenize{a_i,}\allowbreak\detokenize{b_i,}\allowbreak\detokenize{c_i,}\allowbreak\detokenize{d_i,}\allowbreak\detokenize{\delta}} & locally typed probability marginal posterior likelihood stratum and covariance scalars & P15 evidential-relevance & meta & \hyperref[def:p15-r43-measures]{\ref*{def:p15-r43-measures}} & \thepage & Local clauses define each role: q is posterior; a,b are likelihoods; p and stratum-indexed variables belong to screening. The annex delta abbreviates the event covariance. \\
\phantomsection\label{entry:symbol-registry-03967}\(O(H),O(x),\Lambda(x),f_1,f_0,s(x)\) & \texttt{\detokenize{O,}\allowbreak\detokenize{\Lambda,}\allowbreak\detokenize{f_1,}\allowbreak\detokenize{f_0,}\allowbreak\detokenize{s}} & odds binary-experiment likelihood ratio masses and supported mixture mass & P15 evidential-relevance & meta & \hyperref[thm:p15-r43-measures]{\ref*{thm:p15-r43-measures}} & \thepage & O(x)=x/(1-x) is scalar odds with declared ordered endpoints; O(H) abbreviates O(P(H)). Lambda is extended-valued at one-sided support and common-zero points are excluded; sufficient-statistic theorem is for two simple laws on a finite data space. \\
\phantomsection\label{entry:symbol-registry-03968}\(T,F,G,f_D,f_Q,f_L,\bowtie\) & \texttt{\detokenize{T,}\allowbreak\detokenize{F,}\allowbreak\detokenize{G,}\allowbreak\detokenize{f_D,}\allowbreak\detokenize{f_Q,}\allowbreak\detokenize{f_L,}\allowbreak\detokenize{\bowtie}} & measure type formal joint-table functions isotonic family transforms and comparison metavariable & P15 evidential-relevance & meta & \hyperref[thm:p15-r43-characterization]{\ref*{thm:p15-r43-characterization}} & \thepage & Bowtie is one of less equal greater. F in the characterization proof is a scalar function, while Linda's F is separately a feminist event; T in the sufficiency proof is a statistic rather than a measure type. \\
\phantomsection\label{entry:symbol-registry-03969}\(l_i,S,T,A,B,\alpha,\dot\vee\) & \texttt{\detokenize{l_i,}\allowbreak\detokenize{S,}\allowbreak\detokenize{T,}\allowbreak\detokenize{A,}\allowbreak\detokenize{B,}\allowbreak\detokenize{\alpha,}\allowbreak\detokenize{\dot\vee}} & compound-evidence likelihood ratios sum/product proof scalars mixture weight and exclusive disjunction & P15 evidential-relevance & meta & \hyperref[thm:p15-r43-composition]{\ref*{thm:p15-r43-composition}} & \thepage & The final annex uses S,T for likelihood sums and A,B for likelihood products, explicitly distinct from proposition events. Mixture weights differ for D/Q and L. \\
\phantomsection\label{entry:symbol-registry-03970}\(x,d,g,r,\lambda,u,v,t,g_i,k_t\) & \texttt{\detokenize{x,}\allowbreak\detokenize{d,}\allowbreak\detokenize{g,}\allowbreak\detokenize{r,}\allowbreak\detokenize{\lambda,}\allowbreak\detokenize{u,}\allowbreak\detokenize{v,}\allowbreak\detokenize{t,}\allowbreak\detokenize{g_i,}\allowbreak\detokenize{k_t}} & locally typed proof coordinates difference increments ratios likelihoods statistic fibers and conditional mass factors & P15 evidential-relevance & meta & \hyperref[annex:p15.r43.measures]{\ref*{annex:p15.r43.measures}} & \thepage & x is joint-event mass in the measures proof and a data point in the sufficiency proof; g=h2-h1 is an added disjoint target's probability; lambda is the extended likelihood ratio; t is a statistic value and k\_t its conditional mass kernel. Roles are declared within each proof, not shared logical-variable meanings. \\
\phantomsection\label{entry:symbol-registry-05109}\(\tau\) & \texttt{\detokenize{\tau}} & estimand function & P15 Fisher property distinctions & object & \hyperref[def:p15-fe-properties]{\ref*{def:p15-fe-properties}} & \thepage & Target equality differs from equality of sampling laws \\
\phantomsection\label{entry:symbol-registry-05110}\(T\) & \texttt{\detokenize{T}} & statistic & P15 Fisher property distinctions & object & \hyperref[def:p15-fe-properties]{\ref*{def:p15-fe-properties}} & \thepage & Finite conditional-law criterion \\
\phantomsection\label{entry:symbol-registry-05111}\(n\) & \texttt{\detokenize{n}} & fixed sample size & P15 Fisher binomial argument & object & \hyperref[def:p15-fe-binomial]{\ref*{def:p15-fe-binomial}} & \thepage & Positive integer shared by count and ordered data \\
\phantomsection\label{entry:symbol-registry-05112}\(Z_i\) & \texttt{\detokenize{Z_i}} & Bernoulli observations & P15 Fisher binomial argument & object & \hyperref[def:p15-fe-binomial]{\ref*{def:p15-fe-binomial}} & \thepage & Independence is an explicit sampling premise \\
\phantomsection\label{entry:symbol-registry-05113}\(K\) & \texttt{\detokenize{K}} & success count & P15 Fisher binomial argument & object & \hyperref[def:p15-fe-binomial]{\ref*{def:p15-fe-binomial}} & \thepage & Count statistic differs from source summation S \\
\phantomsection\label{entry:symbol-registry-05114}\(x,y\) & \texttt{\detokenize{x;y}} & observed success and failure counts & P15 Fisher binomial argument & object & \hyperref[def:p15-fe-binomial]{\ref*{def:p15-fe-binomial}} & \thepage & Source notation retained with x+y=n \\
\phantomsection\label{entry:symbol-registry-05115}\(p\) & \texttt{\detokenize{p}} & success parameter & P15 Fisher binomial argument & object & \hyperref[def:p15-fe-binomial]{\ref*{def:p15-fe-binomial}} & \thepage & Closed and open domains are different cases \\
\phantomsection\label{entry:symbol-registry-05116}\(\widehat p\) & \texttt{\detokenize{\widehat p}} & maximum likelihood proportion & P15 Fisher binomial argument & object & \hyperref[thm:p15-fe-binomial]{\ref*{thm:p15-fe-binomial}} & \thepage & Endpoint existence depends on the parameter space \\
\phantomsection\label{entry:symbol-registry-05117}\(\mathcal L\) & \texttt{\detokenize{\mathcal L}} & likelihood function & P15 Fisher binomial argument & object & \hyperref[def:p15-fe-binomial]{\ref*{def:p15-fe-binomial}} & \thepage & Editorial raw likelihood not a parameter density \\
\phantomsection\label{entry:symbol-registry-05118}\(L\) & \texttt{\detokenize{L}} & log likelihood & P15 Fisher binomial argument & object & \hyperref[def:p15-fe-binomial]{\ref*{def:p15-fe-binomial}} & \thepage & Source p331 log-likelihood role retained \\
\phantomsection\label{entry:symbol-registry-05119}\(S\) & \texttt{\detokenize{S}} & source summation notation & P15 Fisher binomial argument & meta & \hyperref[def:p15-fe-binomial]{\ref*{def:p15-fe-binomial}} & \thepage & Source L=S(log f) is not the count statistic K \\
\phantomsection\label{entry:symbol-registry-05120}\(\theta\) & \texttt{\detokenize{\theta}} & sine coordinate & P15 Fisher coordinate argument & object & \hyperref[thm:p15-fe-coordinates]{\ref*{thm:p15-fe-coordinates}} & \thepage & Source principal monotone coordinate interval \\
\phantomsection\label{entry:symbol-registry-05121}\(p\) & \texttt{\detokenize{p}} & original probability coordinate & P15 Fisher coordinate argument & object & \hyperref[thm:p15-fe-coordinates]{\ref*{thm:p15-fe-coordinates}} & \thepage & Flat priors in two coordinates are different measures \\
\phantomsection\label{entry:symbol-registry-05122}\(\eta\) & \texttt{\detokenize{\eta}} & finite-family parameter & P15 Fisher finite regularity & object & \hyperref[def:p15-fe-finite-regularity]{\ref*{def:p15-fe-finite-regularity}} & \thepage & Open interval and positive common support \\
\phantomsection\label{entry:symbol-registry-05123}\(q_\eta\) & \texttt{\detokenize{q_\eta}} & sampling masses & P15 Fisher finite regularity & object & \hyperref[def:p15-fe-finite-regularity]{\ref*{def:p15-fe-finite-regularity}} & \thepage & Finite differentiated normalization \\
\phantomsection\label{entry:symbol-registry-05124}\(U_\eta\) & \texttt{\detokenize{U_\eta}} & score function & P15 Fisher finite regularity & object & \hyperref[def:p15-fe-finite-regularity]{\ref*{def:p15-fe-finite-regularity}} & \thepage & Derivative of log mass \\
\phantomsection\label{entry:symbol-registry-05125}\(J\) & \texttt{\detokenize{J}} & expected score information & P15 Fisher finite regularity & object & \hyperref[def:p15-fe-finite-regularity]{\ref*{def:p15-fe-finite-regularity}} & \thepage & Zero information is retained without division \\
\phantomsection\label{entry:symbol-registry-05126}\(T\) & \texttt{\detokenize{T}} & fixed real statistic & P15 Fisher finite regularity & object & \hyperref[def:p15-fe-finite-regularity]{\ref*{def:p15-fe-finite-regularity}} & \thepage & Parameter-independent statistic \\
\phantomsection\label{entry:symbol-registry-05127}\(m\) & \texttt{\detokenize{m}} & statistic expectation function & P15 Fisher finite regularity & object & \hyperref[def:p15-fe-finite-regularity]{\ref*{def:p15-fe-finite-regularity}} & \thepage & Distinct from the Cauchy location parameter \\
\phantomsection\label{entry:symbol-registry-05128}\(\eta\) & \texttt{\detokenize{\eta}} & countermodel parameter & P15 Fisher finite countermodels & object & \hyperref[thm:p15-fe-failures]{\ref*{thm:p15-fe-failures}} & \thepage & Positive and signed extensions differ in identification \\
\phantomsection\label{entry:symbol-registry-05129}\(q_\eta\) & \texttt{\detokenize{q_\eta}} & three-outcome masses & P15 Fisher finite countermodels & object & \hyperref[thm:p15-fe-failures]{\ref*{thm:p15-fe-failures}} & \thepage & Constructed model not measured effects \\
\phantomsection\label{entry:symbol-registry-05130}\(T\) & \texttt{\detokenize{T}} & maximum likelihood statistic & P15 Fisher finite countermodels & object & \hyperref[thm:p15-fe-failures]{\ref*{thm:p15-fe-failures}} & \thepage & Equal maximizing points do not imply sufficient fibers \\
\phantomsection\label{entry:symbol-registry-05131}\(B\) & \texttt{\detokenize{B}} & shared Bernoulli variable & P15 Fisher finite countermodels & object & \hyperref[thm:p15-fe-failures]{\ref*{thm:p15-fe-failures}} & \thepage & Repeated random value not independent trials \\
\phantomsection\label{entry:symbol-registry-05132}\(m\) & \texttt{\detokenize{m}} & Cauchy location & P15 Fisher Cauchy location & object & \hyperref[thm:p15-fe-cauchy]{\ref*{thm:p15-fe-cauchy}} & \thepage & Source notation not a finite population mean \\
\phantomsection\label{entry:symbol-registry-05133}\(X_i\) & \texttt{\detokenize{X_i}} & Cauchy observations & P15 Fisher Cauchy location & object & \hyperref[thm:p15-fe-cauchy]{\ref*{thm:p15-fe-cauchy}} & \thepage & Independent unit-scale location sample \\
\phantomsection\label{entry:symbol-registry-05134}\(M_n\) & \texttt{\detokenize{M_n}} & middle order statistic & P15 Fisher Cauchy location & object & \hyperref[thm:p15-fe-cauchy]{\ref*{thm:p15-fe-cauchy}} & \thepage & Odd n=2r+1 only \\
\phantomsection\label{entry:symbol-registry-05135}\(r\) & \texttt{\detokenize{r}} & median order index & P15 Fisher Cauchy location & object & \hyperref[thm:p15-fe-cauchy]{\ref*{thm:p15-fe-cauchy}} & \thepage & Nonnegative integer \\
\phantomsection\label{entry:symbol-registry-05136}\(\delta_\epsilon\) & \texttt{\detokenize{\delta_\epsilon}} & Cauchy probability displacement & P15 Fisher Cauchy location & object & \hyperref[thm:p15-fe-cauchy]{\ref*{thm:p15-fe-cauchy}} & \thepage & Positive gap in finite deviation bound \\
\phantomsection\label{entry:symbol-registry-05137}\(N\) & \texttt{\detokenize{N}} & observed sequence length & p15-rmdl-source & object & \hyperref[def:p15-rmdl-lengths]{\ref*{def:p15-rmdl-lengths}} & \thepage & Positive integer; not parameter dimension \\
\phantomsection\label{entry:symbol-registry-05138}\(D(p,q)\) & \texttt{\detokenize{D(p,}\allowbreak\detokenize{q)}} & equal-mass divergence in nats & p15-rmdl-gibbs & meta & \hyperref[thm:p15-rmdl-gibbs]{\ref*{thm:p15-rmdl-gibbs}} & \thepage & Summable negative part and possible positive infinity \\
\phantomsection\label{entry:symbol-registry-05139}\(H_2(p)\) & \texttt{\detokenize{H_2(p)}} & discrete entropy in bits & p15-rmdl-gibbs & meta & \hyperref[thm:p15-rmdl-gibbs]{\ref*{thm:p15-rmdl-gibbs}} & \thepage & Fixed mass-function domain \\
\phantomsection\label{entry:symbol-registry-05140}\(S\) & \texttt{\detokenize{S}} & multinomial type-class cardinality & p15-rmdl-types & object & \hyperref[thm:p15-rmdl-types]{\ref*{thm:p15-rmdl-types}} & \thepage & Number of strings with fixed counts \\
\phantomsection\label{entry:symbol-registry-05141}\(s\) & \texttt{\detokenize{s}} & integer model structure & p15-rmdl-source & object & \hyperref[sec:p15-rmdl-structure]{\ref*{sec:p15-rmdl-structure}} & \thepage & Includes declared recurrence orders \\
\phantomsection\label{entry:symbol-registry-05142}\(\theta\) & \texttt{\detokenize{\theta}} & complete model description & p15-rmdl-source & object & \hyperref[sec:p15-rmdl-structure]{\ref*{sec:p15-rmdl-structure}} & \thepage & Structure and numerical parameters \\
\phantomsection\label{entry:symbol-registry-05143}\(\xi\) & \texttt{\detokenize{\xi}} & numerical parameter vector & p15-rmdl-source & object & \hyperref[sec:p15-rmdl-structure]{\ref*{sec:p15-rmdl-structure}} & \thepage & System and innovation-distribution parameters \\
\phantomsection\label{entry:symbol-registry-05144}\(\varepsilon\) & \texttt{\detokenize{\varepsilon}} & innovation coding resolution & p15-rmdl-source & object & \hyperref[sec:p15-rmdl-structure]{\ref*{sec:p15-rmdl-structure}} & \thepage & Not the Occam error tolerance \\
\phantomsection\label{entry:symbol-registry-05145}\(\delta_i\) & \texttt{\detokenize{\delta_i}} & parameter precision width & p15-rmdl-source & object & \hyperref[sec:p15-rmdl-structure]{\ref*{sec:p15-rmdl-structure}} & \thepage & Not the Occam failure-probability tolerance \\
\phantomsection\label{entry:symbol-registry-05146}\(L(e\mid\theta)\) & \texttt{\detokenize{L(e\mid\theta)}} & conditional residual-length approximation & p15-rmdl-source & meta & \hyperref[sec:p15-rmdl-structure]{\ref*{sec:p15-rmdl-structure}} & \thepage & Negative log small-cell density; not posterior probability \\
\phantomsection\label{entry:symbol-registry-05147}\(\lvert\xi_i\rvert\) & \texttt{\detokenize{\lvert\xi_i\rvert}} & numerical magnitude & p15-rmdl-source & object & \hyperref[sec:p15-rmdl-structure]{\ref*{sec:p15-rmdl-structure}} & \thepage & Explicit source-local absolute value; not cardinality \\
\phantomsection\label{entry:symbol-registry-05148}\(n\) & \texttt{\detokenize{n}} & autoregressive order & p15-rmdl-gaussian & object & \hyperref[def:p15-rmdl-arma]{\ref*{def:p15-rmdl-arma}} & \thepage & Nonnegative integer \\
\phantomsection\label{entry:symbol-registry-05149}\(m\) & \texttt{\detokenize{m}} & moving-average order & p15-rmdl-gaussian & object & \hyperref[def:p15-rmdl-arma]{\ref*{def:p15-rmdl-arma}} & \thepage & Nonnegative integer; not Occam sample count \\
\phantomsection\label{entry:symbol-registry-05150}\(\eta\) & \texttt{\detokenize{\eta}} & ARMA coefficient vector & p15-rmdl-gaussian & object & \hyperref[def:p15-rmdl-arma]{\ref*{def:p15-rmdl-arma}} & \thepage & Active nonzero coordinates in the local calculation \\
\phantomsection\label{entry:symbol-registry-05151}\(r\) & \texttt{\detokenize{r}} & innovation variance & p15-rmdl-gaussian & object & \hyperref[def:p15-rmdl-arma]{\ref*{def:p15-rmdl-arma}} & \thepage & Strictly positive Gaussian parameter \\
\phantomsection\label{entry:symbol-registry-05152}\(E_N(\eta)\) & \texttt{\detokenize{E_N(\eta)}} & residual sum of squares & p15-rmdl-gaussian & meta & \hyperref[def:p15-rmdl-arma]{\ref*{def:p15-rmdl-arma}} & \thepage & Recurrence-defined finite data quantity \\
\phantomsection\label{entry:symbol-registry-05153}\(\rho(\eta)\) & \texttt{\detokenize{\rho(\eta)}} & description-surrogate variance estimate & p15-rmdl-gaussian & meta & \hyperref[thm:p15-rmdl-variance]{\ref*{thm:p15-rmdl-variance}} & \thepage & Residual sum divided by N+2; not ordinary ML denominator N \\
\phantomsection\label{entry:symbol-registry-05154}\(\kappa_{n,m}\) & \texttt{\detokenize{\kappa_{n,m}}} & approximate structure cost & p15-rmdl-gaussian & meta & \hyperref[def:p15-rmdl-arma]{\ref*{def:p15-rmdl-arma}} & \thepage & Log product of shifted orders \\
\phantomsection\label{entry:symbol-registry-05155}\(K_N\) & \texttt{\detokenize{K_N}} & common observation-resolution cost & p15-rmdl-gaussian & meta & \hyperref[def:p15-rmdl-arma]{\ref*{def:p15-rmdl-arma}} & \thepage & Fixed units and resolution across candidates \\
\phantomsection\label{entry:symbol-registry-05156}\(L_\delta\) & \texttt{\detokenize{L_\delta}} & Gaussian local description surrogate & p15-rmdl-gaussian & meta & \hyperref[def:p15-rmdl-arma]{\ref*{def:p15-rmdl-arma}} & \thepage & Not an exact integer prefix-code length \\
\phantomsection\label{entry:symbol-registry-05157}\(\Phi\) & \texttt{\detokenize{\Phi}} & controlled local rounding cost & p15-rmdl-rounding & meta & \hyperref[thm:p15-rmdl-precision]{\ref*{thm:p15-rmdl-precision}} & \thepage & Third derivatives controlled on the entire box \\
\phantomsection\label{entry:symbol-registry-05158}\(h_i\) & \texttt{\detokenize{h_i}} & positive Hessian diagonal & p15-rmdl-rounding & object & \hyperref[thm:p15-rmdl-precision]{\ref*{thm:p15-rmdl-precision}} & \thepage & At the stationary fit; not a learning hypothesis \\
\phantomsection\label{entry:symbol-registry-05159}\(A_i\) & \texttt{\detokenize{A_i}} & centered independent rounding error & p15-rmdl-rounding & object & \hyperref[thm:p15-rmdl-precision]{\ref*{thm:p15-rmdl-precision}} & \thepage & Stipulated uniform law; not automatic deterministic-rounding behavior \\
\phantomsection\label{entry:symbol-registry-05160}\(G_i\) & \texttt{\detokenize{G_i}} & log residual-variance curvature & p15-rmdl-rounding & object & \hyperref[thm:p15-rmdl-precision]{\ref*{thm:p15-rmdl-precision}} & \thepage & Positive only in the declared approximation regime \\
\phantomsection\label{entry:symbol-registry-05161}\(Q(\delta)\) & \texttt{\detokenize{Q(\delta)}} & constant-retaining rounding surrogate & p15-rmdl-rounding & meta & \hyperref[thm:p15-rmdl-precision]{\ref*{thm:p15-rmdl-precision}} & \thepage & Fit and curvature held fixed \\
\phantomsection\label{entry:symbol-registry-05162}\(U_R\) & \texttt{\detokenize{U_R}} & published1978 approximate score & p15-rmdl-rounding & meta & \hyperref[thm:p15-rmdl-precision]{\ref*{thm:p15-rmdl-precision}} & \thepage & Not exact code length or integrated Bayesian evidence \\
\phantomsection\label{entry:symbol-registry-05163}\(d\) & \texttt{\detokenize{d}} & active numerical dimension & p15-rmdl-rounding & object & \hyperref[thm:p15-rmdl-precision]{\ref*{thm:p15-rmdl-precision}} & \thepage & Includes variance; fixed in the leading limit \\
\phantomsection\label{entry:symbol-registry-05164}\(L_{\omega P}\) & \texttt{\detokenize{L}} & source finitary probability-quantifier language & p15/probability-quantifiers & object & \hyperref[def:p15-pq-language]{\ref*{def:p15-pq-language}} & \thepage & First-order logic augmented; not the unavailable1985 chapter language. \\
\phantomsection\label{entry:symbol-registry-05165}\(M_n\) & \texttt{\detokenize{M}} & all relational structures on a fixed n-element universe & p15/probability-quantifiers & object & \hyperref[def:p15-pq-models]{\ref*{def:p15-pq-models}} & \thepage & Finite nonempty relational vocabulary; no functions. \\
\phantomsection\label{entry:symbol-registry-05166}\(\mu_n\) & \texttt{\detokenize{\mu}} & measure on random structures & p15/probability-quantifiers & meta & \hyperref[def:p15-pq-models]{\ref*{def:p15-pq-models}} & \thepage & Independent ground relation events with p\_R; distinct from tuple measure. \\
\phantomsection\label{entry:symbol-registry-05167}\(p_R\) & \texttt{\detokenize{p}} & ground relation-event Bernoulli parameter & p15/probability-quantifiers & object & \hyperref[def:p15-pq-models]{\ref*{def:p15-pq-models}} & \thepage & Strictly between zero and one; unbiased value one half. \\
\phantomsection\label{entry:symbol-registry-05168}\(\lambda_A^\ell\) & \texttt{\detokenize{\lambda}} & uniform counting measure on tuples inside one fixed model & p15/probability-quantifiers & meta & \hyperref[def:p15-pq-models]{\ref*{def:p15-pq-models}} & \thepage & All subsets measurable; coordinates sampled with replacement. \\
\phantomsection\label{entry:symbol-registry-05169}\(\exists^{\ge r}\) & \texttt{\detokenize{\exists}} & source probability quantifier binding a tuple & p15/probability-quantifiers & meta & \hyperref[def:p15-pq-language]{\ref*{def:p15-pq-language}} & \thepage & Real threshold strictly between zero and one; ordinary quantifiers retained. \\
\phantomsection\label{entry:symbol-registry-05170}\(\bar x,\bar y\) & \texttt{\detokenize{\bar}} & disjoint syntactically distinct free and bound variable lists & p15/probability-quantifiers & object & \hyperref[def:p15-pq-language]{\ref*{def:p15-pq-language}} & \thepage & Assigned values may coincide; capture-avoiding substitution. \\
\phantomsection\label{entry:symbol-registry-05171}\(\tau_n\) & \texttt{\detokenize{\tau}} & exact size-indexed FO counting translation & p15/probability-quantifiers & meta & \hyperref[thm:p15-pq-translation]{\ref*{thm:p15-pq-translation}} & \thepage & Effective at rational thresholds; no uniform elimination claim. \\
\phantomsection\label{entry:symbol-registry-05172}\(k=\lceil rn^\ell\rceil\) & \texttt{\detokenize{\lceil}} & integer number of distinct tuple witnesses & p15/probability-quantifiers & meta & \hyperref[thm:p15-pq-translation]{\ref*{thm:p15-pq-translation}} & \thepage & Distinct tuples not necessarily distinct coordinates within tuples. \\
\phantomsection\label{entry:symbol-registry-05173}\(\eta\) & \texttt{\detokenize{\eta}} & binary half-row counterexample sentence & p15/probability-quantifiers & object & \hyperref[thm:p15-pq-limit-failures]{\ref*{thm:p15-pq-limit-failures}} & \thepage & Unrestricted ordered relation atoms independent including diagonal. \\
\phantomsection\label{entry:symbol-registry-05174}\(q_m\) & \texttt{\detokenize{q}} & central binomial atom probability & p15/probability-quantifiers/annex & object & \hyperref[thm:p15-pq-limit-failures]{\ref*{thm:p15-pq-limit-failures}} & \thepage & Defines the exact even-size tail; square-root-exponent repair. \\
\phantomsection\label{entry:symbol-registry-05221}\(A\) & \texttt{\detokenize{A}} & Antecedent circumstance name & P15 Whewell Mill source schemata & object & \hyperref[def:p15-wm-agreement]{\ref*{def:p15-wm-agreement}} & \thepage & Mill's capital-letter scheme; juxtaposition lists circumstances and is not multiplication \\
\phantomsection\label{entry:symbol-registry-05222}\(B\) & \texttt{\detokenize{B}} & Antecedent circumstance name & P15 Whewell Mill source schemata & object & \hyperref[def:p15-wm-agreement]{\ref*{def:p15-wm-agreement}} & \thepage & Distinct circumstance; positive contexts and the editorial subset illustration retain this role \\
\phantomsection\label{entry:symbol-registry-05223}\(C\) & \texttt{\detokenize{C}} & Antecedent circumstance name & P15 Whewell Mill source schemata & object & \hyperref[def:p15-wm-agreement]{\ref*{def:p15-wm-agreement}} & \thepage & Distinct circumstance; not a credence function in this scope \\
\phantomsection\label{entry:symbol-registry-05224}\(D\) & \texttt{\detokenize{D}} & Antecedent circumstance name & P15 Whewell Mill source schemata & object & \hyperref[def:p15-wm-agreement]{\ref*{def:p15-wm-agreement}} & \thepage & Distinct circumstance; not a dataset in this scope \\
\phantomsection\label{entry:symbol-registry-05225}\(E\) & \texttt{\detokenize{E}} & Antecedent circumstance name & P15 Whewell Mill source schemata & object & \hyperref[def:p15-wm-agreement]{\ref*{def:p15-wm-agreement}} & \thepage & Distinct circumstance; not an expectation operator \\
\phantomsection\label{entry:symbol-registry-05226}\(a\) & \texttt{\detokenize{a}} & Consequent phenomenon & P15 Whewell Mill source schemata & object & \hyperref[def:p15-wm-agreement]{\ref*{def:p15-wm-agreement}} & \thepage & Occurrence in a context is distinguished from a numeral or coefficient \\
\phantomsection\label{entry:symbol-registry-05227}\(b\) & \texttt{\detokenize{b}} & Consequent phenomenon & P15 Whewell Mill source schemata & object & \hyperref[def:p15-wm-agreement]{\ref*{def:p15-wm-agreement}} & \thepage & Mill's lower-case consequent scheme \\
\phantomsection\label{entry:symbol-registry-05228}\(c\) & \texttt{\detokenize{c}} & Consequent phenomenon & P15 Whewell Mill source schemata & object & \hyperref[def:p15-wm-agreement]{\ref*{def:p15-wm-agreement}} & \thepage & Distinct from the real polynomial coefficient in the separate finite-continuation scope \\
\phantomsection\label{entry:symbol-registry-05229}\(d\) & \texttt{\detokenize{d}} & Consequent phenomenon & P15 Whewell Mill source schemata & object & \hyperref[def:p15-wm-agreement]{\ref*{def:p15-wm-agreement}} & \thepage & Mill's lower-case consequent paired with circumstance D \\
\phantomsection\label{entry:symbol-registry-05230}\(e\) & \texttt{\detokenize{e}} & Consequent phenomenon & P15 Whewell Mill source schemata & object & \hyperref[def:p15-wm-agreement]{\ref*{def:p15-wm-agreement}} & \thepage & Mill's lower-case consequent paired with circumstance E \\
\phantomsection\label{entry:symbol-registry-05231}\(U\) & \texttt{\detokenize{U}} & Fair exogenous binary variable & P15 Whewell Mill causal illustration & object & \hyperref[thm:p15-wm-causal-ambiguity]{\ref*{thm:p15-wm-causal-ambiguity}} & \thepage & Editorial variable; source does not present this equation-replacement model \\
\phantomsection\label{entry:symbol-registry-05232}\(X\) & \texttt{\detokenize{X}} & Binary assigned variable & P15 Whewell Mill causal illustration & object & \hyperref[thm:p15-wm-causal-ambiguity]{\ref*{thm:p15-wm-causal-ambiguity}} & \thepage & The equation for this variable is replaced by zero in the declared intervention \\
\phantomsection\label{entry:symbol-registry-05233}\(Y\) & \texttt{\detokenize{Y}} & Binary outcome variable & P15 Whewell Mill causal illustration & object & \hyperref[thm:p15-wm-causal-ambiguity]{\ref*{thm:p15-wm-causal-ambiguity}} & \thepage & Observational and replaced-assignment outcome laws differ in standing \\
\phantomsection\label{entry:symbol-registry-05234}\(m\) & \texttt{\detokenize{m}} & Positive number of exact observations & P15 Whewell Mill finite continuation & object & \hyperref[thm:p15-wm-finite-continuation]{\ref*{thm:p15-wm-finite-continuation}} & \thepage & The edge case m=1 uses empty products equal to one \\
\phantomsection\label{entry:symbol-registry-05235}\(x_i\) & \texttt{\detokenize{x_i}} & Distinct real observed abscissa & P15 Whewell Mill finite continuation & object & \hyperref[thm:p15-wm-finite-continuation]{\ref*{thm:p15-wm-finite-continuation}} & \thepage & Exact stipulated records; not recovered historical measurements \\
\phantomsection\label{entry:symbol-registry-05236}\(y_i\) & \texttt{\detokenize{y_i}} & Real observed ordinate & P15 Whewell Mill finite continuation & object & \hyperref[thm:p15-wm-finite-continuation]{\ref*{thm:p15-wm-finite-continuation}} & \thepage & Arbitrary real value at the indexed observed abscissa \\
\phantomsection\label{entry:symbol-registry-05237}\(x\) & \texttt{\detokenize{x}} & Real polynomial argument & P15 Whewell Mill finite continuation & object & \hyperref[thm:p15-wm-finite-continuation]{\ref*{thm:p15-wm-finite-continuation}} & \thepage & Dummy argument; distinct from binary capital X in the separate causal illustration \\
\phantomsection\label{entry:symbol-registry-05238}\(p\) & \texttt{\detokenize{p}} & Interpolating real polynomial & P15 Whewell Mill finite continuation & object & \hyperref[thm:p15-wm-finite-continuation]{\ref*{thm:p15-wm-finite-continuation}} & \thepage & Not a probability function \\
\phantomsection\label{entry:symbol-registry-05239}\(p_c\) & \texttt{\detokenize{p_c}} & One-parameter family of fitting polynomials & P15 Whewell Mill finite continuation & object & \hyperref[thm:p15-wm-finite-continuation]{\ref*{thm:p15-wm-finite-continuation}} & \thepage & Each real coefficient gives the same exact finite fitted values \\
\phantomsection\label{entry:symbol-registry-05240}\(c\) & \texttt{\detokenize{c}} & Real perturbation coefficient & P15 Whewell Mill finite continuation & object & \hyperref[thm:p15-wm-finite-continuation]{\ref*{thm:p15-wm-finite-continuation}} & \thepage & Distinct scoped role from Mill's lower-case consequent phenomenon \\
\phantomsection\label{entry:symbol-registry-05241}\(t\) & \texttt{\detokenize{t}} & Unobserved real evaluation point & P15 Whewell Mill finite continuation & object & \hyperref[thm:p15-wm-finite-continuation]{\ref*{thm:p15-wm-finite-continuation}} & \thepage & Must differ from every observed abscissa \\
\phantomsection\label{entry:symbol-registry-05242}\(R\) & \texttt{\detokenize{R}} & Nonzero product at the unobserved point & P15 Whewell Mill finite continuation & object & \hyperref[annex:p15.wm.finite-continuation]{\ref*{annex:p15.wm.finite-continuation}} & \thepage & Annex-local abbreviation used to solve for the coefficient \\
\phantomsection\label{entry:symbol-registry-05243}\(z\) & \texttt{\detokenize{z}} & Prescribed real continuation value & P15 Whewell Mill finite continuation & object & \hyperref[annex:p15.wm.finite-continuation]{\ref*{annex:p15.wm.finite-continuation}} & \thepage & Annex-local arbitrary target value \\
\phantomsection\label{entry:symbol-registry-05244}\(i\) & \texttt{\detokenize{i}} & Observed-point summation and product index & P15 Whewell Mill finite continuation & object & \hyperref[thm:p15-wm-finite-continuation]{\ref*{thm:p15-wm-finite-continuation}} & \thepage & Ranges from one through m \\
\phantomsection\label{entry:symbol-registry-05245}\(j\) & \texttt{\detokenize{j}} & Observed-point product index excluding i & P15 Whewell Mill finite continuation & object & \hyperref[thm:p15-wm-finite-continuation]{\ref*{thm:p15-wm-finite-continuation}} & \thepage & The stated exclusion makes all denominators nonzero \\
\phantomsection\label{entry:symbol-registry-05246}\(k\) & \texttt{\detokenize{k}} & Fixed observed-point evaluation index & P15 Whewell Mill finite continuation & object & \hyperref[annex:p15.wm.finite-continuation]{\ref*{annex:p15.wm.finite-continuation}} & \thepage & Annex-local index ranging from one through m \\
\phantomsection\label{entry:symbol-registry-05247}\(\mathsf{O}_t\) & \texttt{\detokenize{\mathsf{O}_t}} & examined before fixed cutoff t & p15-goodman-projectibility & object & \hyperref[def:p15-goodman-grue]{\ref*{def:p15-goodman-grue}} & \thepage & Historical observation condition; not a colour-change operator \\
\phantomsection\label{entry:symbol-registry-05248}\(\mathsf{E}\) & \texttt{\detokenize{\mathsf{E}}} & emerald predicate & p15-goodman-projectibility & object & \hyperref[def:p15-goodman-grue]{\ref*{def:p15-goodman-grue}} & \thepage & Object domain U \\
\phantomsection\label{entry:symbol-registry-05249}\(\mathsf{G}\) & \texttt{\detokenize{\mathsf{G}}} & green predicate & p15-goodman-projectibility & object & \hyperref[def:p15-goodman-grue]{\ref*{def:p15-goodman-grue}} & \thepage & No exhaustive colour partition assumed \\
\phantomsection\label{entry:symbol-registry-05250}\(\mathsf{B}\) & \texttt{\detokenize{\mathsf{B}}} & blue predicate & p15-goodman-projectibility & object & \hyperref[def:p15-goodman-grue]{\ref*{def:p15-goodman-grue}} & \thepage & Disjointness used only for specified incompatibility \\
\phantomsection\label{entry:symbol-registry-05251}\(\mathsf{Q}_t\) & \texttt{\detokenize{\mathsf{Q}_t}} & grue predicate at fixed cutoff & p15-goodman-projectibility & object & \hyperref[def:p15-goodman-grue]{\ref*{def:p15-goodman-grue}} & \thepage & Depends on observation history \\
\phantomsection\label{entry:symbol-registry-05252}\(\mathsf{R}_t\) & \texttt{\detokenize{\mathsf{R}_t}} & common-domain bleen counterpart & p15-goodman-projectibility & object & \hyperref[def:p15-goodman-bleen]{\ref*{def:p15-goodman-bleen}} & \thepage & Explicit formal extension of source illustration \\
\phantomsection\label{entry:symbol-registry-05253}\(\mathsf{A}_t\) & \texttt{\detokenize{\mathsf{A}_t}} & emerose-style antecedent & p15-goodman-projectibility & object & \hyperref[def:p15-goodman-antecedent]{\ref*{def:p15-goodman-antecedent}} & \thepage & Examined emerald or unexamined rose \\
\phantomsection\label{entry:symbol-registry-05254}\(\mathsf{V}\) & \texttt{\detokenize{\mathsf{V}}} & rose predicate & p15-goodman-projectibility & object & \hyperref[def:p15-goodman-antecedent]{\ref*{def:p15-goodman-antecedent}} & \thepage & Distinct from emerald predicate \\
\phantomsection\label{entry:symbol-registry-05255}\(C\) & \texttt{\detokenize{C}} & abstract confirmation relation on formulas & p15-goodman-projectibility & object & \hyperref[thm:p15-goodman-confirmation-collapse]{\ref*{thm:p15-goodman-confirmation-collapse}} & \thepage & Not a probability or a claim of normative adequacy \\
\phantomsection\label{entry:symbol-registry-05256}\(X\) & \texttt{\detokenize{X}} & ungrouped measurement & P15 Fisher midpoint experiment & object & \hyperref[def:p15-fg-grid]{\ref*{def:p15-fg-grid}} & \thepage & Density and absolute-moment hypotheses are explicit \\
\phantomsection\label{entry:symbol-registry-05257}\(h\) & \texttt{\detokenize{h}} & physical bin width & P15 Fisher midpoint experiment & object & \hyperref[def:p15-fg-grid]{\ref*{def:p15-fg-grid}} & \thepage & Positive fixed measurement-unit width \\
\phantomsection\label{entry:symbol-registry-05258}\(\vartheta\) & \texttt{\detokenize{\vartheta}} & grid phase & P15 Fisher midpoint experiment & object & \hyperref[def:p15-fg-grid]{\ref*{def:p15-fg-grid}} & \thepage & Fixed phase is not its uniform average \\
\phantomsection\label{entry:symbol-registry-05259}\(c_j\) & \texttt{\detokenize{c_j}} & lattice midpoint & P15 Fisher midpoint experiment & object & \hyperref[def:p15-fg-grid]{\ref*{def:p15-fg-grid}} & \thepage & Half-open class endpoint convention \\
\phantomsection\label{entry:symbol-registry-05260}\(B_j\) & \texttt{\detokenize{B_j}} & recorded class & P15 Fisher midpoint experiment & object & \hyperref[def:p15-fg-grid]{\ref*{def:p15-fg-grid}} & \thepage & Partition of the real line \\
\phantomsection\label{entry:symbol-registry-05261}\(Y_\vartheta\) & \texttt{\detokenize{Y_\vartheta}} & recorded midpoint & P15 Fisher midpoint experiment & object & \hyperref[def:p15-fg-grid]{\ref*{def:p15-fg-grid}} & \thepage & Deterministic function of measurement at fixed phase \\
\phantomsection\label{entry:symbol-registry-05262}\(G_k\) & \texttt{\detokenize{G_k}} & fixed-grid raw moment & P15 Fisher midpoint experiment & object & \hyperref[def:p15-fg-grid]{\ref*{def:p15-fg-grid}} & \thepage & Finite absolute kth moment required \\
\phantomsection\label{entry:symbol-registry-05263}\(M_k\) & \texttt{\detokenize{M_k}} & ungrouped raw moment & P15 Fisher midpoint experiment & object & \hyperref[def:p15-fg-grid]{\ref*{def:p15-fg-grid}} & \thepage & Same fixed origin for all corrections \\
\phantomsection\label{entry:symbol-registry-05264}\(\overline G_k\) & \texttt{\detokenize{\overline G_k}} & phase-averaged raw moment & P15 Fisher midpoint experiment & object & \hyperref[def:p15-fg-grid]{\ref*{def:p15-fg-grid}} & \thepage & Bar is phase integration not sample averaging \\
\phantomsection\label{entry:symbol-registry-05265}\(U\) & \texttt{\detokenize{U}} & independent uniform error & P15 Fisher phase-average correction & object & \hyperref[thm:p15-fg-phase]{\ref*{thm:p15-fg-phase}} & \thepage & Exact representation only after phase averaging \\
\phantomsection\label{entry:symbol-registry-05266}\(\sigma\) & \texttt{\detokenize{\sigma}} & normal standard deviation & P15 Fisher normal Fourier correction & object & \hyperref[thm:p15-fg-fourier]{\ref*{thm:p15-fg-fourier}} & \thepage & Positive scale in measurement units \\
\phantomsection\label{entry:symbol-registry-05267}\(H\) & \texttt{\detokenize{H}} & bin-mass Fourier transform & P15 Fisher normal Fourier correction & object & \hyperref[thm:p15-fg-fourier]{\ref*{thm:p15-fg-fourier}} & \thepage & Continuous value h at zero \\
\phantomsection\label{entry:symbol-registry-05268}\(\omega_s\) & \texttt{\detokenize{\omega_s}} & lattice Fourier frequency & P15 Fisher normal Fourier correction & object & \hyperref[thm:p15-fg-fourier]{\ref*{thm:p15-fg-fourier}} & \thepage & Equals two pi s over physical width \\
\phantomsection\label{entry:symbol-registry-05269}\(C_{k,s}\) & \texttt{\detokenize{C_{k,s}}} & complex moment coefficient & P15 Fisher normal Fourier correction & object & \hyperref[thm:p15-fg-fourier]{\ref*{thm:p15-fg-fourier}} & \thepage & Exact derivative formula \\
\phantomsection\label{entry:symbol-registry-05270}\(i\) & \texttt{\detokenize{i}} & imaginary unit & P15 Fisher normal Fourier correction & object & \hyperref[thm:p15-fg-fourier]{\ref*{thm:p15-fg-fourier}} & \thepage & Not an observation index \\
\phantomsection\label{entry:symbol-registry-05271}\(c\) & \texttt{\detokenize{c}} & Gaussian harmonic decay constant & P15 Fisher normal Fourier correction & object & \hyperref[thm:p15-fg-fourier]{\ref*{thm:p15-fg-fourier}} & \thepage & Not a lattice midpoint c subscript j \\
\phantomsection\label{entry:symbol-registry-05272}\(m\) & \texttt{\detokenize{m}} & harmonic truncation index & P15 Fisher normal Fourier correction & object & \hyperref[thm:p15-fg-fourier]{\ref*{thm:p15-fg-fourier}} & \thepage & Not cell expectation m subscript j \\
\phantomsection\label{entry:symbol-registry-05273}\(\epsilon\) & \texttt{\detokenize{\epsilon}} & source Fourier width coordinate & P15 Fisher normal Fourier correction & object & \hyperref[thm:p15-fg-fourier]{\ref*{thm:p15-fg-fourier}} & \thepage & Fisher uses physical width two pi epsilon \\
\phantomsection\label{entry:symbol-registry-05274}\(n\) & \texttt{\detokenize{n}} & independent repetition count & P15 Fisher moment precision & object & \hyperref[thm:p15-fg-moments]{\ref*{thm:p15-fg-moments}} & \thepage & Not a cell count n subscript j \\
\phantomsection\label{entry:symbol-registry-05275}\(p_j\) & \texttt{\detokenize{p_j}} & normalized cell probability & P15 Fisher cell likelihood & object & \hyperref[def:p15-fg-cells]{\ref*{def:p15-fg-cells}} & \thepage & Cell probabilities sum to one \\
\phantomsection\label{entry:symbol-registry-05276}\(\eta\) & \texttt{\detokenize{\eta}} & cell-model parameter & P15 Fisher cell likelihood & object & \hyperref[def:p15-fg-cells]{\ref*{def:p15-fg-cells}} & \thepage & Likelihood parameter not parameter probability \\
\phantomsection\label{entry:symbol-registry-05277}\(N\) & \texttt{\detokenize{N}} & total recorded count & P15 Fisher cell likelihood & object & \hyperref[def:p15-fg-cells]{\ref*{def:p15-fg-cells}} & \thepage & Positive sample size for discrepancy formulas \\
\phantomsection\label{entry:symbol-registry-05278}\(n_j\) & \texttt{\detokenize{n_j}} & observed cell count & P15 Fisher cell likelihood & object & \hyperref[def:p15-fg-cells]{\ref*{def:p15-fg-cells}} & \thepage & Zero-count log term has continuous convention \\
\phantomsection\label{entry:symbol-registry-05279}\(m_j\) & \texttt{\detokenize{m_j}} & expected cell count & P15 Fisher cell likelihood & object & \hyperref[def:p15-fg-cells]{\ref*{def:p15-fg-cells}} & \thepage & Equals N times p subscript j \\
\phantomsection\label{entry:symbol-registry-05280}\(D\) & \texttt{\detokenize{D}} & negative loglikelihood discrepancy & P15 Fisher cell likelihood & object & \hyperref[def:p15-fg-cells]{\ref*{def:p15-fg-cells}} & \thepage & Fisher L on printed357 has reversed likelihood sign \\
\phantomsection\label{entry:symbol-registry-05281}\(Q\) & \texttt{\detokenize{Q}} & Pearson discrepancy & P15 Fisher cell likelihood & object & \hyperref[def:p15-fg-cells]{\ref*{def:p15-fg-cells}} & \thepage & Exact objective differs from D \\
\phantomsection\label{entry:symbol-registry-05282}\(\delta_j\) & \texttt{\detokenize{\delta_j}} & relative count residual & P15 Fisher Pearson approximation & object & \hyperref[thm:p15-fg-deviance]{\ref*{thm:p15-fg-deviance}} & \thepage & Observed minus expected divided by expected \\
\phantomsection\label{entry:symbol-registry-05283}\(\rho\) & \texttt{\detokenize{\rho}} & uniform residual bound & P15 Fisher Pearson approximation & object & \hyperref[thm:p15-fg-deviance]{\ref*{thm:p15-fg-deviance}} & \thepage & Strictly below one \\
\phantomsection\label{entry:symbol-registry-05284}\(\mu\) & \texttt{\detokenize{\mu}} & normal location & P15 Fisher normal class experiment & object & \hyperref[def:p15-fg-normal]{\ref*{def:p15-fg-normal}} & \thepage & Endpoint probabilities rather than midpoint density \\
\phantomsection\label{entry:symbol-registry-05285}\(\phi\) & \texttt{\detokenize{\phi}} & standard normal density & P15 Fisher normal class experiment & object & \hyperref[def:p15-fg-normal]{\ref*{def:p15-fg-normal}} & \thepage & Gaussian derivatives and moments proved in annex \\
\phantomsection\label{entry:symbol-registry-05286}\(\Phi\) & \texttt{\detokenize{\Phi}} & standard normal distribution function & P15 Fisher normal class experiment & object & \hyperref[def:p15-fg-normal]{\ref*{def:p15-fg-normal}} & \thepage & Endpoint derivative is phi \\
\phantomsection\label{entry:symbol-registry-05287}\(b_j\) & \texttt{\detokenize{b_j}} & fixed physical cell endpoint & P15 Fisher normal class experiment & object & \hyperref[def:p15-fg-normal]{\ref*{def:p15-fg-normal}} & \thepage & Not standardized or parameter-dependent endpoint \\
\phantomsection\label{entry:symbol-registry-05288}\(z_j\) & \texttt{\detokenize{z_j}} & standardized endpoint & P15 Fisher normal class experiment & object & \hyperref[def:p15-fg-normal]{\ref*{def:p15-fg-normal}} & \thepage & Equals endpoint minus location divided by scale \\
\phantomsection\label{entry:symbol-registry-05289}\(S\) & \texttt{\detokenize{S}} & normal parameter score & P15 Fisher recorded normal information & object & \hyperref[thm:p15-fg-information]{\ref*{thm:p15-fg-information}} & \thepage & Conditional expectation yields recorded-class score \\
\phantomsection\label{entry:symbol-registry-05290}\(I_{\rm full}\) & \texttt{\detokenize{I_{\rm full}}} & ungrouped score information & P15 Fisher recorded normal information & object & \hyperref[thm:p15-fg-information]{\ref*{thm:p15-fg-information}} & \thepage & Scalar parameter with other coordinate held fixed \\
\phantomsection\label{entry:symbol-registry-05291}\(I_{\rm class}\) & \texttt{\detokenize{I_{\rm class}}} & class score information & P15 Fisher recorded normal information & object & \hyperref[thm:p15-fg-information]{\ref*{thm:p15-fg-information}} & \thepage & Exact fixed-partition information \\
\phantomsection\label{entry:symbol-registry-05292}\(a\) & \texttt{\detokenize{a}} & dimensionless normal bin width & P15 Fisher averaged information expansion & object & \hyperref[thm:p15-fg-expansion]{\ref*{thm:p15-fg-expansion}} & \thepage & Equals h over sigma not source earlier physical a \\
\phantomsection\label{entry:symbol-registry-05293}\(z\) & \texttt{\detokenize{z}} & standardized recorded midpoint & P15 Fisher averaged information expansion & object & \hyperref[thm:p15-fg-expansion]{\ref*{thm:p15-fg-expansion}} & \thepage & Centered at true location in this experiment \\
\phantomsection\label{entry:symbol-registry-05294}\(v_a\) & \texttt{\detokenize{v_a}} & averaged normal bin density & P15 Fisher averaged information expansion & object & \hyperref[thm:p15-fg-expansion]{\ref*{thm:p15-fg-expansion}} & \thepage & Exact phase-randomized density not every fixed-grid mass \\
\phantomsection\label{entry:symbol-registry-05295}\(J_\mu\) & \texttt{\detokenize{J_\mu}} & dimensionless location information & P15 Fisher averaged information expansion & object & \hyperref[thm:p15-fg-expansion]{\ref*{thm:p15-fg-expansion}} & \thepage & Physical information multiplied by sigma squared \\
\phantomsection\label{entry:symbol-registry-05296}\(J_\sigma\) & \texttt{\detokenize{J_\sigma}} & dimensionless scale information & P15 Fisher averaged information expansion & object & \hyperref[thm:p15-fg-expansion]{\ref*{thm:p15-fg-expansion}} & \thepage & Separate from estimator sampling variance \\
\phantomsection\label{entry:symbol-registry-05297}\(O(a^{10})\) & \texttt{\detokenize{O(a^{10})}} & controlled asymptotic remainder & P15 Fisher averaged information expansion & meta & \hyperref[thm:p15-fg-expansion]{\ref*{thm:p15-fg-expansion}} & \thepage & Compact width interval dominated integral bound \\
\phantomsection\label{entry:symbol-registry-05298}\(R\) & \texttt{\detokenize{R}} & bin density divided by Gaussian density & P15 Fisher expansion proof & object & \hyperref[annex:p15.fg.expansion]{\ref*{annex:p15.fg.expansion}} & \thepage & Jensen lower bound controls reciprocal derivatives \\
\phantomsection\label{entry:symbol-registry-05299}\(H_r\) & \texttt{\detokenize{H_r}} & probabilists Hermite polynomial & P15 Fisher expansion proof & object & \hyperref[annex:p15.fg.expansion]{\ref*{annex:p15.fg.expansion}} & \thepage & Not Fourier transform H in another scope \\
\phantomsection\label{entry:symbol-registry-05300}\(r_{2j}\) & \texttt{\detokenize{r_{2j}}} & even bin-density coefficient & P15 Fisher expansion proof & object & \hyperref[annex:p15.fg.expansion]{\ref*{annex:p15.fg.expansion}} & \thepage & Finite explicit Hermite coefficient \\
\phantomsection\label{entry:symbol-registry-05301}\(\ell_{2j}\) & \texttt{\detokenize{\ell_{2j}}} & even log-bin-density coefficient & P15 Fisher expansion proof & object & \hyperref[annex:p15.fg.expansion]{\ref*{annex:p15.fg.expansion}} & \thepage & Finite log series algebra \\
\phantomsection\label{entry:symbol-registry-05302}\(q\) & \texttt{\detokenize{q}} & standardized location logarithmic derivative & P15 Fisher expansion proof & object & \hyperref[annex:p15.fg.expansion]{\ref*{annex:p15.fg.expansion}} & \thepage & Odd in standardized midpoint \\
\phantomsection\label{entry:symbol-registry-05303}\(t\) & \texttt{\detokenize{t}} & standardized scale score factor & P15 Fisher expansion proof & object & \hyperref[annex:p15.fg.expansion]{\ref*{annex:p15.fg.expansion}} & \thepage & Even in standardized midpoint \\
\phantomsection\label{entry:symbol-registry-05321}\(\mathcal H\) & \texttt{\detokenize{\mathcal H}} & domain of formulated universal hypotheses & p15-goodman-entrenchment & object & \hyperref[def:p15-goodman-time-status]{\ref*{def:p15-goodman-time-status}} & \thepage & Not limited to hypotheses already projected at assessment time \\
\phantomsection\label{entry:symbol-registry-05322}\(T_\tau\) & \texttt{\detokenize{T_\tau}} & cases determined true at assessment time & p15-goodman-entrenchment & meta & \hyperref[def:p15-goodman-time-status]{\ref*{def:p15-goodman-time-status}} & \thepage & Epistemic record; conditional hypothesis instances are antecedent cases \\
\phantomsection\label{entry:symbol-registry-05323}\(F_\tau\) & \texttt{\detokenize{F_\tau}} & cases determined false at assessment time & p15-goodman-entrenchment & meta & \hyperref[def:p15-goodman-time-status]{\ref*{def:p15-goodman-time-status}} & \thepage & Not every actual false case has been determined \\
\phantomsection\label{entry:symbol-registry-05324}\(I_\tau\) & \texttt{\detokenize{I_\tau}} & instances not yet determined at assessment time & p15-goodman-entrenchment & meta & \hyperref[def:p15-goodman-time-status]{\ref*{def:p15-goodman-time-status}} & \thepage & May concern past events \\
\phantomsection\label{entry:symbol-registry-05325}\(\mathcal E_\tau\) & \texttt{\detokenize{\mathcal E_\tau}} & supported unviolated unexhausted competitors & p15-goodman-entrenchment & object & \hyperref[def:p15-goodman-time-status]{\ref*{def:p15-goodman-time-status}} & \thepage & Editorial abbreviation for the source's standing restriction \\
\phantomsection\label{entry:symbol-registry-05326}\(\succ_\tau\) & \texttt{\detokenize{\succ_\tau}} & better entrenchment of hypotheses by typed antecedent/consequent comparisons & p15-goodman-entrenchment & relation & \hyperref[def:p15-goodman-competitors]{\ref*{def:p15-goodman-competitors}} & \thepage & Coarse supplied comparison; not a sum or universal numerical ranking \\
\phantomsection\label{entry:symbol-registry-05327}\(\operatorname{Conf}_\tau\) & \texttt{\detokenize{\operatorname{Conf}_\tau}} & source conflict with target-ascription and nonconsequence conditions & p15-goodman-entrenchment & relation & \hyperref[def:p15-goodman-competitors]{\ref*{def:p15-goodman-competitors}} & \thepage & Eligible competitors; not arbitrary syntactic inequality \\
\phantomsection\label{entry:symbol-registry-05328}\(\operatorname{Ov}_\tau\) & \texttt{\detokenize{\operatorname{Ov}_\tau}} & printed short binary override relation & p15-goodman-entrenchment & relation & \hyperref[def:p15-goodman-short-override]{\ref*{def:p15-goodman-short-override}} & \thepage & Overrider must have no still-better eligible conflict \\
\phantomsection\label{entry:symbol-registry-05329}\(\operatorname{O}_\tau\) & \texttt{\detokenize{\operatorname{O}_\tau}} & overridden hypothesis predicate & p15-goodman-entrenchment & meta & \hyperref[def:p15-goodman-short-override]{\ref*{def:p15-goodman-short-override}} & \thepage & Not the observation-history predicate O\_t of the preceding chapter \\
\phantomsection\label{entry:symbol-registry-05330}\(\operatorname{P}_\tau\) & \texttt{\detokenize{\operatorname{P}_\tau}} & presumptively projectible hypothesis predicate & p15-goodman-entrenchment & meta & \hyperref[def:p15-goodman-short-override]{\ref*{def:p15-goodman-short-override}} & \thepage & Not a probability \\
\phantomsection\label{entry:symbol-registry-05331}\(\operatorname{N}_\tau\) & \texttt{\detokenize{\operatorname{N}_\tau}} & presumptively nonprojectible hypothesis predicate & p15-goodman-entrenchment & meta & \hyperref[def:p15-goodman-short-override]{\ref*{def:p15-goodman-short-override}} & \thepage & Unresolved nonoverridden eligible conflict \\
\phantomsection\label{entry:symbol-registry-05332}\(\operatorname{U}_\tau\) & \texttt{\detokenize{\operatorname{U}_\tau}} & presumptively unprojectible hypothesis predicate & p15-goodman-entrenchment & meta & \hyperref[def:p15-goodman-short-override]{\ref*{def:p15-goodman-short-override}} & \thepage & Ineligibility or overridden status \\
\phantomsection\label{entry:symbol-registry-05333}\(B\) & \texttt{\detokenize{B}} & consequent predicate in abstract hypothesis comparison & p15-goodman-override & object & \hyperref[def:p15-goodman-competitors]{\ref*{def:p15-goodman-competitors}} & \thepage & Local predicate role; distinct from the later bag name \\
\phantomsection\label{entry:symbol-registry-05334}\(B\) & \texttt{\detokenize{B}} & particular bag of marbles & p15-goodman-comparative and p15-goodman-negative-overhypotheses & object & \hyperref[def:p15-goodman-positive-overhypothesis]{\ref*{def:p15-goodman-positive-overhypothesis}} & \thepage & Source bag name; not the earlier abstract consequent predicate \\
\phantomsection\label{entry:symbol-registry-05335}\(G\) & \texttt{\detokenize{G}} & uniform-colour positive overhypothesis on bagfuls in stack S & p15-goodman-comparative and p15-goodman-negative-overhypotheses & object & \hyperref[def:p15-goodman-positive-overhypothesis]{\ref*{def:p15-goodman-positive-overhypothesis}} & \thepage & Not the green predicate of the preceding chapter \\
\phantomsection\label{entry:symbol-registry-05344}\(D_E(\bar x)\) & \texttt{\detokenize{D}} & complete equality diagram of parameter values & p15/probability-quantifier-elimination & meta & \hyperref[def:p15-pe-types]{\ref*{def:p15-pe-types}} & \thepage & Coordinates in one equivalence class coincide; syntactically distinct variables may have equal values. \\
\phantomsection\label{entry:symbol-registry-05345}\(\beta(\bar x)\) & \texttt{\detokenize{\beta}} & consistent complete quantifier-free parameter type & p15/probability-quantifier-elimination & meta & \hyperref[def:p15-pe-types]{\ref*{def:p15-pe-types}} & \thepage & Specifies equalities and every relation atom on its variable list. \\
\phantomsection\label{entry:symbol-registry-05346}\(\bar x_E\) & \texttt{\detokenize{x}} & parameter tuple with equality classes collapsed & p15/probability-quantifier-elimination & object & \hyperref[def:p15-pe-types]{\ref*{def:p15-pe-types}} & \thepage & One new representative variable per class; repeated at its original positions. \\
\phantomsection\label{entry:symbol-registry-05347}\(q_\beta\) & \texttt{\detokenize{q}} & conditional fresh-scope critical value & p15/probability-quantifier-elimination & object & \hyperref[def:p15-pe-types]{\ref*{def:p15-pe-types}} & \thepage & Actual limit constructed by finite atom-key truth-pattern sum. \\
\phantomsection\label{entry:symbol-registry-05348}\(\varphi^0\) & \texttt{\detokenize{\varphi}} & noncritical quantifier-free content & p15/probability-quantifier-elimination & object & \hyperref[def:p15-pe-content]{\ref*{def:p15-pe-content}} & \thepage & Probability content uses strict comparison; ordinary existence includes parameter witnesses. \\
\phantomsection\label{entry:symbol-registry-05349}\(P_0\) & \texttt{\detokenize{P}} & finite threshold or threshold-sequence inventory & p15/probability-quantifier-elimination & object & \hyperref[thm:p15-pe-strong]{\ref*{thm:p15-pe-strong}} & \thepage & Common-class scope also requires a finite variable budget. \\
\phantomsection\label{entry:symbol-registry-05350}\(C\) & \texttt{\detokenize{C}} & common class of finite structures & p15/probability-quantifier-elimination & object & \hyperref[def:p15-pe-ae]{\ref*{def:p15-pe-ae}} & \thepage & One class for a declared fixed-budget fragment; not all finite structures. \\
\phantomsection\label{entry:symbol-registry-05351}\(\pi\) & \texttt{\detokenize{\pi}} & parametric universal constraint & p15/probability-quantifier-elimination & object & \hyperref[def:p15-pe-parametric]{\ref*{def:p15-pe-parametric}} & \thepage & Every variable occurs in every literal; models of all positive finite sizes required. \\
\phantomsection\label{entry:symbol-registry-05352}\(q_\beta^\pi\) & \texttt{\detokenize{q}} & critical value under the conditioned support-block law & p15/probability-quantifier-elimination & object & \hyperref[def:p15-pe-parametric]{\ref*{def:p15-pe-parametric}} & \thepage & Finite normalized allowed-configuration sum; atoms within a block need not be independent. \\
\phantomsection\label{entry:symbol-registry-05353}\(\varphi^\pi\) & \texttt{\detokenize{\varphi}} & parametric quantifier-free content & p15/probability-quantifier-elimination & object & \hyperref[def:p15-pe-parametric]{\ref*{def:p15-pe-parametric}} & \thepage & Only complete types consistent with the constraint are used. \\
\phantomsection\label{entry:symbol-registry-05354}\(q_{\beta,n}\) & \texttt{\detokenize{q}} & size-dependent conditional scope mean & p15/probability-quantifier-elimination & object & \hyperref[def:p15-pe-sequences]{\ref*{def:p15-pe-sequences}} & \thepage & May not converge; criticality uses both liminf comparisons. \\
\phantomsection\label{entry:symbol-registry-05355}\(r(n)\) & \texttt{\detokenize{r}} & size-dependent probability threshold & p15/probability-quantifier-elimination & meta & \hyperref[def:p15-pe-sequences]{\ref*{def:p15-pe-sequences}} & \thepage & Every value strictly between zero and one; noncriticality requires eventual separation by type. \\
\phantomsection\label{entry:symbol-registry-05379}\(\mathcal C\) & \texttt{\detokenize{\mathcal C}} & fixed measurable event class & p15-vcl-events & object & \hyperref[def:p15-vcl-event-class]{\ref*{def:p15-vcl-event-class}} & \thepage & Nonempty and fixed before iid sampling \\
\phantomsection\label{entry:symbol-registry-05380}\(\mathcal C\upharpoonright S\) & \texttt{\detokenize{\mathcal C\upharpoonright S}} & finite trace family & p15-vcl-capacity & meta & \hyperref[def:p15-vcl-event-class]{\ref*{def:p15-vcl-event-class}} & \thepage & Intersection family on a finite set; no restriction on arbitrary functions is meant \\
\phantomsection\label{entry:symbol-registry-05381}\(\Pi_{\mathcal C}(t)\) & \texttt{\detokenize{\Pi_{\mathcal C}(t)}} & worst sampled trace count & p15-vcl-capacity & meta & \hyperref[def:p15-vcl-event-class]{\ref*{def:p15-vcl-event-class}} & \thepage & Ordered tuples permit repetitions; value at zero is one \\
\phantomsection\label{entry:symbol-registry-05382}\(N(F,n)\) & \texttt{\detokenize{N(F,}\allowbreak\detokenize{n)}} & source trace-growth notation & p15-vcl-source-dictionary & meta & \hyperref[def:p15-vcl-event-class]{\ref*{def:p15-vcl-event-class}} & \thepage & 2008 author notation retained in an explicit dictionary \\
\phantomsection\label{entry:symbol-registry-05383}\(d\) & \texttt{\detokenize{d}} & VC dimension & p15-vcl-capacity & object & \hyperref[thm:p15-vcl-sauer]{\ref*{thm:p15-vcl-sauer}} & \thepage & May be zero or infinite; polynomial expression only when one through t \\
\phantomsection\label{entry:symbol-registry-05384}\(P_m(C)\) & \texttt{\detokenize{P_m(C)}} & empirical event frequency & p15-vcl-events & meta & \hyperref[def:p15-vcl-event-class]{\ref*{def:p15-vcl-event-class}} & \thepage & Defined only for positive sample size \\
\phantomsection\label{entry:symbol-registry-05385}\(\Delta_m\) & \texttt{\detokenize{\Delta_m}} & uniform empirical discrepancy & p15-vcl-events & meta & \hyperref[def:p15-vcl-event-class]{\ref*{def:p15-vcl-event-class}} & \thepage & Measurability proved using pointwise countable core \\
\phantomsection\label{entry:symbol-registry-05386}\(\mathcal C_0\) & \texttt{\detokenize{\mathcal C_0}} & pointwise countable event core & p15-vcl-events & object & \hyperref[def:p15-vcl-countable-core]{\ref*{def:p15-vcl-countable-core}} & \thepage & Approximation of each indicator rather than assumed uniform convergence \\
\phantomsection\label{entry:symbol-registry-05387}\(H_0\) & \texttt{\detokenize{H_0}} & pointwise countable predictor core & p15-vcl-classification & object & \hyperref[def:p15-vcl-classification]{\ref*{def:p15-vcl-classification}} & \thepage & Enumeration used for a measurable first empirical minimizer \\
\phantomsection\label{entry:symbol-registry-05388}\(\mathcal E_H\) & \texttt{\detokenize{\mathcal E_H}} & binary error-event class & p15-vcl-classification & object & \hyperref[def:p15-vcl-classification]{\ref*{def:p15-vcl-classification}} & \thepage & Error patterns correspond to prediction patterns by fixed-label XOR \\
\phantomsection\label{entry:symbol-registry-05389}\(\Pi_H(t)\) & \texttt{\detokenize{\Pi_H(t)}} & prediction trace growth & p15-vcl-classification & meta & \hyperref[def:p15-vcl-classification]{\ref*{def:p15-vcl-classification}} & \thepage & Positive prediction events; same growth as binary error events \\
\phantomsection\label{entry:symbol-registry-05390}\(R(h)\) & \texttt{\detokenize{R(h)}} & population zero-one risk & p15-vcl-classification & meta & \hyperref[def:p15-vcl-classification]{\ref*{def:p15-vcl-classification}} & \thepage & Joint input-label law can be noisy \\
\phantomsection\label{entry:symbol-registry-05391}\(\widehat R_m(h)\) & \texttt{\detokenize{\widehat R_m(h)}} & empirical zero-one risk & p15-vcl-classification & meta & \hyperref[def:p15-vcl-classification]{\ref*{def:p15-vcl-classification}} & \thepage & Positive sample size and a predeclared predictor class \\
\phantomsection\label{entry:symbol-registry-05392}\(R_H^*\) & \texttt{\detokenize{R_H^*}} & infimal class risk & p15-vcl-classification & meta & \hyperref[def:p15-vcl-classification]{\ref*{def:p15-vcl-classification}} & \thepage & Infimum need not be attained; infimum zero is not realizability \\
\phantomsection\label{entry:symbol-registry-05393}\(\widehat h\) & \texttt{\detokenize{\widehat h}} & data-selected predictor & p15-vcl-classification & meta & \hyperref[thm:p15-vcl-erm]{\ref*{thm:p15-vcl-erm}} & \thepage & Mathematical measurable selection is not a terminating effective search claim \\
\phantomsection\label{entry:symbol-registry-05394}\(\eta\) & \texttt{\detokenize{\eta}} & empirical minimization tolerance & p15-vcl-classification & object & \hyperref[thm:p15-vcl-erm]{\ref*{thm:p15-vcl-erm}} & \thepage & Nonnegative additive empirical-risk tolerance \\
\phantomsection\label{entry:symbol-registry-05395}\(\varepsilon\) & \texttt{\detokenize{\varepsilon}} & deviation or error tolerance & p15-vcl-events & object & \hyperref[thm:p15-vcl-core]{\ref*{thm:p15-vcl-core}} & \thepage & Uniform bound uses a tolerance strictly between zero and one \\
\phantomsection\label{entry:symbol-registry-05396}\(\delta\) & \texttt{\detokenize{\delta}} & failure-probability allowance & p15-vcl-classification & object & \hyperref[thm:p15-vcl-samples]{\ref*{thm:p15-vcl-samples}} & \thepage & Strictly between zero and one \\
\phantomsection\label{entry:symbol-registry-05397}\(\sigma_i\) & \texttt{\detokenize{\sigma_i}} & independent fair pair-swap sign & p15-vcl-split & object & \hyperref[annex:p15.vcl.uniform]{\ref*{annex:p15.vcl.uniform}} & \thepage & Sign controls pooled sample orientation rather than original data-dependent hypothesis selection \\
\phantomsection\label{entry:symbol-registry-05398}\(v\) & \texttt{\detokenize{v}} & independent validation sample size & p15-vcl-classification & object & \hyperref[thm:p15-vcl-erm]{\ref*{thm:p15-vcl-erm}} & \thepage & Positive integer and same future sampling law \\
\phantomsection\label{entry:symbol-registry-05503}\(E\) & \texttt{\detokenize{E}} & Total stated observation report & P15 Hempel confirmation & object & \hyperref[def:p15-hempel-reports]{\ref*{def:p15-hempel-reports}} & \thepage & Editorial report metavariable; original PartII uses B; not an event or expectation Blank macro filled with exact unchanged glyph spelling; this is data, not a newly defined TeX macro. \\
\phantomsection\label{entry:symbol-registry-05504}\(H\) & \texttt{\detokenize{H}} & Hypothesis sentence & P15 Hempel confirmation & object & \hyperref[def:p15-hempel-reports]{\ref*{def:p15-hempel-reports}} & \thepage & May contain arbitrary first-order individual quantifier nesting in the stated identity-free relational language Blank macro filled with exact unchanged glyph spelling; this is data, not a newly defined TeX macro. \\
\phantomsection\label{entry:symbol-registry-05505}\(C(E,H)\) & \texttt{\detokenize{C(E,}\allowbreak\detokenize{H)}} & Nonquantitative confirmation relation & P15 Hempel confirmation & relation & \hyperref[def:p15-hempel-reports]{\ref*{def:p15-hempel-reports}} & \thepage & Not a probability or support increment; later satisfaction construction has explicit named-report boundaries Blank macro filled with exact unchanged glyph spelling; this is data, not a newly defined TeX macro. \\
\phantomsection\label{entry:symbol-registry-05506}\(\mathsf R\) & \texttt{\detokenize{\mathsf R}} & Unary raven predicate & P15 Hempel raven cases & object & \hyperref[thm:p15-hempel-nicod-spelling]{\ref*{thm:p15-hempel-nicod-spelling}} & \thepage & Distinct from the binary R in the separate ordered-pair case Blank macro filled with exact unchanged glyph spelling; this is data, not a newly defined TeX macro. \\
\phantomsection\label{entry:symbol-registry-05507}\(\mathsf B\) & \texttt{\detokenize{\mathsf B}} & Unary black predicate & P15 Hempel raven cases & object & \hyperref[thm:p15-hempel-nicod-spelling]{\ref*{thm:p15-hempel-nicod-spelling}} & \thepage & Not the original PartII report metavariable B Blank macro filled with exact unchanged glyph spelling; this is data, not a newly defined TeX macro. \\
\phantomsection\label{entry:symbol-registry-05508}\(H_R\) & \texttt{\detokenize{H_R}} & Universal black-raven hypothesis & P15 Hempel raven cases & object & \hyperref[thm:p15-hempel-nicod-spelling]{\ref*{thm:p15-hempel-nicod-spelling}} & \thepage & No existential import is built into the universal conditional Blank macro filled with exact unchanged glyph spelling; this is data, not a newly defined TeX macro. \\
\phantomsection\label{entry:symbol-registry-05509}\(H_N\) & \texttt{\detokenize{H_N}} & Contrapositive raven hypothesis & P15 Hempel raven cases & object & \hyperref[thm:p15-hempel-nicod-spelling]{\ref*{thm:p15-hempel-nicod-spelling}} & \thepage & Logically equivalent to H\_R before existential conjuncts are added Blank macro filled with exact unchanged glyph spelling; this is data, not a newly defined TeX macro. \\
\phantomsection\label{entry:symbol-registry-05510}\(R\) & \texttt{\detokenize{R}} & Binary relational predicate & P15 Hempel ordered-pair counterexample & object & \hyperref[thm:p15-hempel-multiple-instance]{\ref*{thm:p15-hempel-multiple-instance}} & \thepage & Source footnote's relational counterexample; not a unary raven predicate Blank macro filled with exact unchanged glyph spelling; this is data, not a newly defined TeX macro. \\
\phantomsection\label{entry:symbol-registry-05511}\(R_1,R_2\) & \texttt{\detokenize{R_1,}\allowbreak\detokenize{R_2}} & Binary observational relations & P15 Hempel prediction counterexample & object & \hyperref[thm:p15-hempel-prediction-gap]{\ref*{thm:p15-hempel-prediction-gap}} & \thepage & Fresh-object construction isolates finite ground observational conservativity Blank macro filled with exact unchanged glyph spelling; this is data, not a newly defined TeX macro. \\
\phantomsection\label{entry:symbol-registry-05512}\(E_1,E_2\) & \texttt{\detokenize{E_1,}\allowbreak\detokenize{E_2}} & Partition of a finite observation report & P15 Hempel prediction criterion & object & \hyperref[def:p15-hempel-prediction]{\ref*{def:p15-hempel-prediction}} & \thepage & Disjoint report subsets; E2 is nonempty; not probability events Blank macro filled with exact unchanged glyph spelling; this is data, not a newly defined TeX macro. \\
\phantomsection\label{entry:symbol-registry-05513}\(\Gamma\) & \texttt{\detokenize{\Gamma}} & Class of confirmed hypothesis sentences & P15 Hempel consequence and consistency & object & \hyperref[def:p15-hempel-conditions]{\ref*{def:p15-hempel-conditions}} & \thepage & Joint compatibility with the report is distinguished from fitting each hypothesis separately Blank macro filled with exact unchanged glyph spelling; this is data, not a newly defined TeX macro. \\
\phantomsection\label{entry:symbol-registry-05514}\(D\) & \texttt{\detokenize{D}} & Finite nonempty class of individual names & P15 Hempel development & object & \hyperref[def:p15-hempel-development]{\ref*{def:p15-hempel-development}} & \thepage & Domain coverage is a premise of the development lemma; names need not denote distinct objects there Blank macro filled with exact unchanged glyph spelling; this is data, not a newly defined TeX macro. \\
\phantomsection\label{entry:symbol-registry-05515}\(\operatorname{dev}_D(H)\) & \texttt{\detokenize{\operatorname{dev}_D(H)}} & Finite development of a hypothesis & P15 Hempel development & n/a & \hyperref[def:p15-hempel-development]{\ref*{def:p15-hempel-development}} & \thepage & Capture-avoiding quantified expansion; not an assertion that the actual scientific domain is finite Blank macro filled with exact unchanged glyph spelling; this is data, not a newly defined TeX macro. Original precedence descriptor 'function' preserved as semantic metadata; n/a is not an invented binding rank. \\
\phantomsection\label{entry:symbol-registry-05516}\(D(E)\) & \texttt{\detokenize{D(E)}} & Names in essential ground atoms of a report & P15 Hempel corrected report support & n/a & \hyperref[def:p15-hempel-essential]{\ref*{def:p15-hempel-essential}} & \thepage & Essential dependence is defined by the Boolean function; empty support is outside the proved fragment Blank macro filled with exact unchanged glyph spelling; this is data, not a newly defined TeX macro. Original precedence descriptor 'function' preserved as semantic metadata; n/a is not an invented binding rank. \\
\phantomsection\label{entry:symbol-registry-05517}\(\Delta_E\) & \texttt{\detokenize{\Delta_E}} & Class directly confirmed by a report & P15 Hempel satisfaction construction & object & \hyperref[def:p15-hempel-satisfaction-core]{\ref*{def:p15-hempel-satisfaction-core}} & \thepage & Editorial class symbol implementing source9.1a; general confirmation is its semantic consequence closure Blank macro filled with exact unchanged glyph spelling; this is data, not a newly defined TeX macro. \\
\phantomsection\label{entry:symbol-registry-05518}\(L\) & \texttt{\detokenize{L}} & Binary liking relation & P15 Hempel quantified confirmation cases & object & \hyperref[thm:p15-hempel-likes]{\ref*{thm:p15-hempel-likes}} & \thepage & Finite named-class confirmation can differ from truth over a larger domain Blank macro filled with exact unchanged glyph spelling; this is data, not a newly defined TeX macro. \\
\phantomsection\label{entry:symbol-registry-05519}\(f_0,f_1\) & \texttt{\detokenize{f_0,}\allowbreak\detokenize{f_1}} & Competing functions fitting two values & P15 Hempel joint-fit illustration & n/a & \hyperref[thm:p15-hempel-joint-consistency]{\ref*{thm:p15-hempel-joint-consistency}} & \thepage & Independent mathematical illustration of source curve-fit objection; outside the restricted relational confirmation language Blank macro filled with exact unchanged glyph spelling; this is data, not a newly defined TeX macro. Original precedence descriptor 'function' preserved as semantic metadata; n/a is not an invented binding rank. \\
\phantomsection\label{entry:symbol-registry-05520}N & \texttt{\detokenize{N}} & number of named individuals & P15 Carnap finite language & n/a & \hyperref[def:p15-carnap-state-range]{\ref*{def:p15-carnap-state-range}} & \thepage & Positive integer; extending N keeps the primitive vocabulary fixed Blank macro filled with exact unchanged glyph spelling; this is data, not a newly defined TeX macro. Original precedence descriptor 'atomic' preserved as semantic metadata; n/a is not an invented binding rank. \\
\phantomsection\label{entry:symbol-registry-05521}q & \texttt{\detokenize{q}} & number of independent primitive unary predicates & P15 Carnap monadic language & n/a & \hyperref[def:p15-carnap-state-range]{\ref*{def:p15-carnap-state-range}} & \thepage & Positive integer; not a confirmation value Blank macro filled with exact unchanged glyph spelling; this is data, not a newly defined TeX macro. Original precedence descriptor 'atomic' preserved as semantic metadata; n/a is not an invented binding rank. \\
\phantomsection\label{entry:symbol-registry-05522}\ensuremath{Omega_N} & \texttt{\detokenize{Omega_N}} & finite set of complete state descriptions & P15 Carnap finite language & n/a & \hyperref[def:p15-carnap-state-range]{\ref*{def:p15-carnap-state-range}} & \thepage & Every primitive ground instance is decided Blank macro filled with exact unchanged glyph spelling; this is data, not a newly defined TeX macro. Original precedence descriptor 'atomic' preserved as semantic metadata; n/a is not an invented binding rank. \\
\phantomsection\label{entry:symbol-registry-05523}\ensuremath{R_N(h)} & \texttt{\detokenize{R_N(h)}} & range of a sentence & P15 Carnap finite language & n/a & \hyperref[def:p15-carnap-state-range]{\ref*{def:p15-carnap-state-range}} & \thepage & Subset of Omega\_N determined by satisfaction Blank macro filled with exact unchanged glyph spelling; this is data, not a newly defined TeX macro. Original precedence descriptor 'function application' preserved as semantic metadata; n/a is not an invented binding rank. \\
\phantomsection\label{entry:symbol-registry-05524}k & \texttt{\detokenize{k}} & number of Q-properties & P15 Carnap monadic language & n/a & \hyperref[def:p15-carnap-structure-width]{\ref*{def:p15-carnap-structure-width}} & \thepage & k=2\^{}q; widths refer to this fixed division Blank macro filled with exact unchanged glyph spelling; this is data, not a newly defined TeX macro. Original precedence descriptor 'atomic' preserved as semantic metadata; n/a is not an invented binding rank. \\
\phantomsection\label{entry:symbol-registry-05525}w(M) & \texttt{\detokenize{w(M)}} & logical width of molecular property & P15 Carnap monadic language & n/a & \hyperref[def:p15-carnap-structure-width]{\ref*{def:p15-carnap-structure-width}} & \thepage & Number of Q-properties; not an empirical frequency Blank macro filled with exact unchanged glyph spelling; this is data, not a newly defined TeX macro. Original precedence descriptor 'function application' preserved as semantic metadata; n/a is not an invented binding rank. \\
\phantomsection\label{entry:symbol-registry-05526}\ensuremath{r_N} & \texttt{\detokenize{r_N}} & number of individual-renaming orbits & P15 Carnap finite language & n/a & \hyperref[thm:p15-carnap-language-counts]{\ref*{thm:p15-carnap-language-counts}} & \thepage & Monadic closed form uses binomial N+k-1 choose k-1 Blank macro filled with exact unchanged glyph spelling; this is data, not a newly defined TeX macro. Original precedence descriptor 'atomic' preserved as semantic metadata; n/a is not an invented binding rank. \\
\phantomsection\label{entry:symbol-registry-05527}\ensuremath{m_N} & \texttt{\detokenize{m_N}} & regular measure on states and ranges & P15 Carnap finite language & n/a & \hyperref[def:p15-carnap-regular-confirmation]{\ref*{def:p15-carnap-regular-confirmation}} & \thepage & Strictly positive state masses sum to one Blank macro filled with exact unchanged glyph spelling; this is data, not a newly defined TeX macro. Original precedence descriptor 'function application' preserved as semantic metadata; n/a is not an invented binding rank. \\
\phantomsection\label{entry:symbol-registry-05528}\ensuremath{c_N(h;e)} & \texttt{\detokenize{c_N(h;e)}} & confirmation of hypothesis on evidence & P15 Carnap finite language & n/a & \hyperref[def:p15-carnap-regular-confirmation]{\ref*{def:p15-carnap-regular-confirmation}} & \thepage & Displayed with comma in LaTeX; undefined for empty evidence range Blank macro filled with exact unchanged glyph spelling; this is data, not a newly defined TeX macro. Original precedence descriptor 'function application' preserved as semantic metadata; n/a is not an invented binding rank. \\
\phantomsection\label{entry:symbol-registry-05529}\ensuremath{m_N star} & \texttt{\detokenize{m_N star}} & equal structure-weight measure & P15 Carnap 1950 reconstruction & n/a & \hyperref[def:p15-carnap-star]{\ref*{def:p15-carnap-star}} & \thepage & Equal orbit mass and equal mass within each orbit Blank macro filled with exact unchanged glyph spelling; this is data, not a newly defined TeX macro. Original precedence descriptor 'function application' preserved as semantic metadata; n/a is not an invented binding rank. \\
\phantomsection\label{entry:symbol-registry-05530}\ensuremath{m_N dagger} & \texttt{\detokenize{m_N dagger}} & equal state-weight alternative & P15 Carnap contrast & n/a & \hyperref[def:p15-carnap-star]{\ref*{def:p15-carnap-star}} & \thepage & Not Carnap c star; no learning about an unmentioned individual Blank macro filled with exact unchanged glyph spelling; this is data, not a newly defined TeX macro. Original precedence descriptor 'function application' preserved as semantic metadata; n/a is not an invented binding rank. \\
\phantomsection\label{entry:symbol-registry-05531}x rising t & \texttt{\detokenize{x rising t}} & rising factorial & P15 Carnap finite calculations & n/a & \hyperref[thm:p15-carnap-universal]{\ref*{thm:p15-carnap-universal}} & \thepage & Empty product equals one Blank macro filled with exact unchanged glyph spelling; this is data, not a newly defined TeX macro. Original precedence descriptor 'function application' preserved as semantic metadata; n/a is not an invented binding rank. \\
\phantomsection\label{entry:symbol-registry-05532}\ensuremath{alpha_i} & \texttt{\detokenize{alpha_i}} & positive prior category weight & P15 distinct later-family reconstruction & n/a & \hyperref[def:p15-carnap-history-family]{\ref*{def:p15-carnap-history-family}} & \thepage & Does not replace the 1950 selected source Blank macro filled with exact unchanged glyph spelling; this is data, not a newly defined TeX macro. Original precedence descriptor 'atomic' preserved as semantic metadata; n/a is not an invented binding rank. \\
\phantomsection\label{entry:symbol-registry-05533}H,A,I,E & \texttt{\detokenize{H,}\allowbreak\detokenize{A,}\allowbreak\detokenize{I,}\allowbreak\detokenize{E}} & target theory; auxiliaries; application conditions; predicted observation & p15-popper-test-package & n/a & \hyperref[def:p15-popper-test-record]{\ref*{def:p15-popper-test-record}} & \thepage & Classical formulas; independently assignable atoms only for the nonentailment counterexamples Original precedence descriptor 'formula' preserved as semantic metadata; n/a is not an invented binding rank. \\
\phantomsection\label{entry:symbol-registry-05534}chi & \texttt{\detokenize{\chi}} & object-language conditional for the tested package & p15-popper-test-package & n/a & \hyperref[thm:p15-popper-package]{\ref*{thm:p15-popper-package}} & \thepage & Not the metalinguistic assertion of semantic consequence Original precedence descriptor 'formula' preserved as semantic metadata; n/a is not an invented binding rank. \\
\phantomsection\label{entry:symbol-registry-05535}U & \texttt{\detokenize{U}} & nonempty specimen domain & p15-popper-thread & n/a & \hyperref[def:p15-popper-thread]{\ref*{def:p15-popper-thread}} & \thepage & One fixed unit for the positive real-valued load and threshold functions Original precedence descriptor 'set' preserved as semantic metadata; n/a is not an invented binding rank. \\
\phantomsection\label{entry:symbol-registry-05536}L,S & \texttt{\detokenize{L,}\allowbreak\detokenize{S}} & actual load and breaking-strength threshold functions & p15-popper-thread & n/a & \hyperref[def:p15-popper-thread]{\ref*{def:p15-popper-thread}} & \thepage & Values are positive real numbers; deterministic idealization only Original precedence descriptor 'function' preserved as semantic metadata; n/a is not an invented binding rank. \\
\phantomsection\label{entry:symbol-registry-05537}\ensuremath{r_L,r_S} & \texttt{\detokenize{r_L,}\allowbreak\detokenize{r_S}} & reported load and reported strength functions & p15-popper-thread & n/a & \hyperref[def:p15-popper-thread]{\ref*{def:p15-popper-thread}} & \thepage & Calibration is an auxiliary equality; a reading alone does not establish the actual value Original precedence descriptor 'function' preserved as semantic metadata; n/a is not an invented binding rank. \\
\phantomsection\label{entry:symbol-registry-05538}Br,R & \texttt{\detokenize{\mathsf{Br},}\allowbreak\detokenize{\mathsf R}} & actual-breaking and detector-reported-breaking predicates & p15-popper-thread & n/a & \hyperref[def:p15-popper-thread]{\ref*{def:p15-popper-thread}} & \thepage & Detector equivalence belongs to the tested auxiliary Original precedence descriptor 'atomic' preserved as semantic metadata; n/a is not an invented binding rank. \\
\phantomsection\label{entry:symbol-registry-05539}c & \texttt{\detokenize{c}} & selected specimen constant & p15-popper-thread & n/a & \hyperref[def:p15-popper-thread]{\ref*{def:p15-popper-thread}} & \thepage & Same reported values in three explicit diagnostic models Original precedence descriptor 'term' preserved as semantic metadata; n/a is not an invented binding rank. \\
\phantomsection\label{entry:symbol-registry-05540}Rv,Wh & \texttt{\detokenize{\mathsf{Rv},}\allowbreak\detokenize{\mathsf{Wh}}} & raven and white predicates & p15-popper-asymmetry & n/a & \hyperref[thm:p15-popper-finite-record]{\ref*{thm:p15-popper-finite-record}} & \thepage & Empirical predicate identification is not supplied by the formal construction Original precedence descriptor 'atomic' preserved as semantic metadata; n/a is not an invented binding rank. \\
\phantomsection\label{entry:symbol-registry-05541}\ensuremath{D_n,H_rw} & \texttt{\detokenize{D_n,}\allowbreak\detokenize{H_{\rm rw}}} & finite named-observation record and universal raven hypothesis & p15-popper-asymmetry & n/a & \hyperref[thm:p15-popper-finite-record]{\ref*{thm:p15-popper-finite-record}} & \thepage & n may be zero; structures remain nonempty Original precedence descriptor 'formula' preserved as semantic metadata; n/a is not an invented binding rank. \\
\phantomsection\label{entry:symbol-registry-05542}\ensuremath{a_i,b} & \texttt{\detokenize{a_i,}\allowbreak\detokenize{b}} & observed object constants and an additional witness constant & p15-popper-asymmetry & n/a & \hyperref[thm:p15-popper-finite-record]{\ref*{thm:p15-popper-finite-record}} & \thepage & Distinctness is specified for observed names; final bounded equivalence does not need unique names Original precedence descriptor 'term' preserved as semantic metadata; n/a is not an invented binding rank. \\
\phantomsection\label{entry:symbol-registry-05543}n & \texttt{\detokenize{n}} & finite observation count or binary-sequence length & p15-popper-asymmetry and p15-popper-corroboration & n/a & \hyperref[thm:p15-popper-finite-record]{\ref*{thm:p15-popper-finite-record}} & \thepage & n>=0 for observation record; n>=1 for uniform binary experiment Original precedence descriptor 'integer' preserved as semantic metadata; n/a is not an invented binding rank. \\
\phantomsection\label{entry:symbol-registry-05544}P & \texttt{\detokenize{P}} & unary predicate in existential-search countermodels & p15-popper-asymmetry & n/a & \hyperref[thm:p15-popper-finite-record]{\ref*{thm:p15-popper-finite-record}} & \thepage & Distinct from the probability measure in the later section Original precedence descriptor 'atomic' preserved as semantic metadata; n/a is not an invented binding rank. \\
\phantomsection\label{entry:symbol-registry-05545}K,B & \texttt{\detokenize{\mathcal K,}\allowbreak\detokenize{\mathcal B}} & nonempty interpretation class and admitted test-statement set & p15-popper-falsifiers & n/a & \hyperref[def:p15-popper-relative-falsifiers]{\ref*{def:p15-popper-relative-falsifiers}} & \thepage & Every admitted test individually has a model in K; no exhaustive observation-language claim Original precedence descriptor 'class' preserved as semantic metadata; n/a is not an invented binding rank. \\
\phantomsection\label{entry:symbol-registry-05546}\ensuremath{T,T_1,T_2} & \texttt{\detokenize{T,}\allowbreak\detokenize{T_1,}\allowbreak\detokenize{T_2}} & theory sentences in a fixed observational comparison & p15-popper-falsifiers & n/a & \hyperref[def:p15-popper-relative-falsifiers]{\ref*{def:p15-popper-relative-falsifiers}} & \thepage & Entailment and compatibility are relative to the same K Original precedence descriptor 'formula' preserved as semantic metadata; n/a is not an invented binding rank. \\
\phantomsection\label{entry:symbol-registry-05547}\ensuremath{F_KB(T)} & \texttt{\detokenize{F_{\mathcal K,\mathcal B}(T)}} & set of admitted tests incompatible with T in K & p15-popper-falsifiers & n/a & \hyperref[def:p15-popper-relative-falsifiers]{\ref*{def:p15-popper-relative-falsifiers}} & \thepage & Potential falsifiers, not an accepted joint data record or a numerical severity score Original precedence descriptor 'set' preserved as semantic metadata; n/a is not an invented binding rank. \\
\phantomsection\label{entry:symbol-registry-05548}Omega,F,P & \texttt{\detokenize{\Omega,}\allowbreak\detokenize{\mathcal F,}\allowbreak\detokenize{\mathbb P}} & sample space; measurable-event sigma algebra; probability measure & p15-popper-corroboration & n/a & \hyperref[thm:p15-popper-probability-order]{\ref*{thm:p15-popper-probability-order}} & \thepage & Specified law only; not a posterior or corroboration score by definition Original precedence descriptor 'probability space' preserved as semantic metadata; n/a is not an invented binding rank. \\
\phantomsection\label{entry:symbol-registry-05549}\ensuremath{V_1,V_2,D} & \texttt{\detokenize{V_1,}\allowbreak\detokenize{V_2,}\allowbreak\detokenize{D}} & nested measurable hypothesis-world events and evidence event & p15-popper-corroboration & n/a & \hyperref[thm:p15-popper-probability-order]{\ref*{thm:p15-popper-probability-order}} & \thepage & P(D)>0 is required by ratio conditioning; nesting differs from exclusive-model comparison Original precedence descriptor 'event' preserved as semantic metadata; n/a is not an invented binding rank. \\
\phantomsection\label{entry:symbol-registry-05550}p,q,r & \texttt{\detokenize{p,}\allowbreak\detokenize{q,}\allowbreak\detokenize{r}} & independent Boolean atoms in observational-order counterexamples & p15-popper-falsifiers annex & n/a & \hyperref[annex:p15.popper.falsifier-order]{\ref*{annex:p15.popper.falsifier-order}} & \thepage & Deliberately incomplete admitted vocabulary; not an assertion about all scientific basic statements Original precedence descriptor 'atomic' preserved as semantic metadata; n/a is not an invented binding rank. \\
\phantomsection\label{entry:symbol-registry-05720}MC & \texttt{\detokenize{MC}} & purely comparative confirmation relation & P15 Carnap continuation & n/a & \hyperref[def:p15-carnap-comparative]{\ref*{def:p15-carnap-comparative}} & \thepage & Agrees with every regular finite confirmation function Blank macro filled with exact unchanged glyph spelling; registry data only, not a new TeX command. Original precedence 'function application' retained as semantic descriptor; n/a is not an invented rank. \\
\phantomsection\label{entry:symbol-registry-05721}Eq & \texttt{\detokenize{Eq}} & mutual comparative confirmation & P15 Carnap continuation & n/a & \hyperref[def:p15-carnap-comparative]{\ref*{def:p15-carnap-comparative}} & \thepage & MC holds in both directions Blank macro filled with exact unchanged glyph spelling; registry data only, not a new TeX command. Original precedence 'function application' retained as semantic descriptor; n/a is not an invented rank. \\
\phantomsection\label{entry:symbol-registry-05722}Gr & \texttt{\detokenize{Gr}} & strict comparative confirmation relation & P15 Carnap continuation & n/a & \hyperref[def:p15-carnap-comparative]{\ref*{def:p15-carnap-comparative}} & \thepage & MC holds only in the first direction Blank macro filled with exact unchanged glyph spelling; registry data only, not a new TeX command. Original precedence 'function application' retained as semantic descriptor; n/a is not an invented rank. \\
\phantomsection\label{entry:symbol-registry-05723}\ensuremath{Q_e} & \texttt{\detokenize{Q_e}} & relevance quotient & P15 Carnap continuation & n/a & \hyperref[def:p15-carnap-relevance]{\ref*{def:p15-carnap-relevance}} & \thepage & Requires positive displayed denominators Blank macro filled with exact unchanged glyph spelling; registry data only, not a new TeX command. Original precedence 'function application' retained as semantic descriptor; n/a is not an invented rank. \\
\phantomsection\label{entry:symbol-registry-05724}r(i;h;e) & \texttt{\detokenize{r(i;h;e)}} & signed relevance measure & P15 Carnap continuation & n/a & \hyperref[def:p15-carnap-relevance]{\ref*{def:p15-carnap-relevance}} & \thepage & Commas in manuscript; not a nonnegative measure Blank macro filled with exact unchanged glyph spelling; registry data only, not a new TeX command. Original precedence 'function application' retained as semantic descriptor; n/a is not an invented rank. \\
\phantomsection\label{entry:symbol-registry-05725}a & \texttt{\detokenize{a}} & scalar mass m(eih) & P15 Carnap relevance determinant & n/a & \hyperref[thm:p15-carnap-relevance-core]{\ref*{thm:p15-carnap-relevance-core}} & \thepage & Locally scoped c cell is not the confirmation function Blank macro filled with exact unchanged glyph spelling; registry data only, not a new TeX command. Original precedence 'atomic' retained as semantic descriptor; n/a is not an invented rank. Root-authorized typed decomposition; original grouped intake row remains unchanged as evidence. \\
\phantomsection\label{entry:symbol-registry-05726}b & \texttt{\detokenize{b}} & scalar mass m(ei not-h) & P15 Carnap relevance determinant & n/a & \hyperref[thm:p15-carnap-relevance-core]{\ref*{thm:p15-carnap-relevance-core}} & \thepage & Locally scoped c cell is not the confirmation function Blank macro filled with exact unchanged glyph spelling; registry data only, not a new TeX command. Original precedence 'atomic' retained as semantic descriptor; n/a is not an invented rank. Root-authorized typed decomposition; original grouped intake row remains unchanged as evidence. \\
\phantomsection\label{entry:symbol-registry-05727}c & \texttt{\detokenize{c}} & scalar mass m(e not-i h) & P15 Carnap relevance determinant & n/a & \hyperref[thm:p15-carnap-relevance-core]{\ref*{thm:p15-carnap-relevance-core}} & \thepage & Locally scoped c cell is not the confirmation function Blank macro filled with exact unchanged glyph spelling; registry data only, not a new TeX command. Original precedence 'atomic' retained as semantic descriptor; n/a is not an invented rank. Root-authorized typed decomposition; original grouped intake row remains unchanged as evidence. \\
\phantomsection\label{entry:symbol-registry-05728}d & \texttt{\detokenize{d}} & scalar mass m(e not-i not-h) & P15 Carnap relevance determinant & n/a & \hyperref[thm:p15-carnap-relevance-core]{\ref*{thm:p15-carnap-relevance-core}} & \thepage & Locally scoped c cell is not the confirmation function Blank macro filled with exact unchanged glyph spelling; registry data only, not a new TeX command. Original precedence 'atomic' retained as semantic descriptor; n/a is not an invented rank. Root-authorized typed decomposition; original grouped intake row remains unchanged as evidence. \\
\phantomsection\label{entry:symbol-registry-05729}c & \texttt{\detokenize{c}} & binary confirmation function & P15 Carnap continuation & n/a & \hyperref[sec:p15-carnap-confirmation-relevance-and-reliability]{\ref*{sec:p15-carnap-confirmation-relevance-and-reliability}} & \thepage & Binary c(h,e)=m(e and h)/m(e), distinct from the scalar c cell by argument type in the SAME theorem. Root-authorized continuation role; source Carnap1950Probability. \\
\phantomsection\label{entry:symbol-registry-05730}\ensuremath{x_l} & \texttt{\detokenize{x_l}} & signed relevance of observation cell l & P15 Carnap seventy-five-sign classification & n/a & \hyperref[thm:p15-carnap-relevance-sign-cases]{\ref*{thm:p15-carnap-relevance-sign-cases}} & \thepage & Four values sum to zero Blank macro filled with exact unchanged glyph spelling; registry data only, not a new TeX command. Original precedence 'atomic' retained as semantic descriptor; n/a is not an invented rank. \\
\phantomsection\label{entry:symbol-registry-05731}C prime & \texttt{\detokenize{C prime}} & pure logical confirming-evidence candidate & P15 Carnap classificatory comparison & n/a & \hyperref[sec:p15-carnap-classificatory]{\ref*{sec:p15-carnap-classificatory}} & \thepage & Not accepted by Carnap as an adequate finished explicatum Blank macro filled with exact unchanged glyph spelling; registry data only, not a new TeX command. Original precedence 'function application' retained as semantic descriptor; n/a is not an invented rank. \\
\phantomsection\label{entry:symbol-registry-05732}\ensuremath{e_c(f;u;e)} & \texttt{\detokenize{e_c(f;u;e)}} & c-mean estimate function & P15 Carnap general estimation & n/a & \hyperref[def:p15-carnap-c-mean]{\ref*{def:p15-carnap-c-mean}} & \thepage & Function and argument identify unknown quantity Blank macro filled with exact unchanged glyph spelling; registry data only, not a new TeX command. Original precedence 'function application' retained as semantic descriptor; n/a is not an invented rank. \\
\phantomsection\label{entry:symbol-registry-05733}\ensuremath{mu_e} & \texttt{\detokenize{mu_e}} & chosen evidence-relative distribution of unknown quantity & P15 Carnap general estimation & n/a & \hyperref[def:p15-carnap-c-mean]{\ref*{def:p15-carnap-c-mean}} & \thepage & General-language law is an explicit assumption Blank macro filled with exact unchanged glyph spelling; registry data only, not a new TeX command. Original precedence 'function application' retained as semantic descriptor; n/a is not an invented rank. \\
\phantomsection\label{entry:symbol-registry-05734}\ensuremath{F_e} & \texttt{\detokenize{F_e}} & cumulative confirmation distribution & P15 Carnap general estimation & n/a & \hyperref[def:p15-carnap-c-mean]{\ref*{def:p15-carnap-c-mean}} & \thepage & No density assumed automatically Blank macro filled with exact unchanged glyph spelling; registry data only, not a new TeX command. Original precedence 'function application' retained as semantic descriptor; n/a is not an invented rank. \\
\phantomsection\label{entry:symbol-registry-05735}X hat e & \texttt{\detokenize{X hat e}} & estimate of unknown quantity X & P15 Carnap general estimation & n/a & \hyperref[def:p15-carnap-c-mean]{\ref*{def:p15-carnap-c-mean}} & \thepage & Not the actual or necessarily a possible value Blank macro filled with exact unchanged glyph spelling; registry data only, not a new TeX command. Original precedence 'atomic' retained as semantic descriptor; n/a is not an invented rank. \\
\phantomsection\label{entry:symbol-registry-05736}V & \texttt{\detokenize{V}} & estimate minus actual quantity & P15 Carnap reliability & n/a & \hyperref[def:p15-carnap-estimated-error]{\ref*{def:p15-carnap-estimated-error}} & \thepage & Carnap's error sign convention Blank macro filled with exact unchanged glyph spelling; registry data only, not a new TeX command. Original precedence 'atomic' retained as semantic descriptor; n/a is not an invented rank. \\
\phantomsection\label{entry:symbol-registry-05737}\ensuremath{J_c} & \texttt{\detokenize{J_c}} & estimated standard error & P15 Carnap reliability & n/a & \hyperref[def:p15-carnap-estimated-error]{\ref*{def:p15-carnap-estimated-error}} & \thepage & Square is the c-mean squared error Blank macro filled with exact unchanged glyph spelling; registry data only, not a new TeX command. Original precedence 'function application' retained as semantic descriptor; n/a is not an invented rank. \\
\phantomsection\label{entry:symbol-registry-05738}gamma & \texttt{\detokenize{gamma}} & distinct-pair predictive indicator covariance & P15 Carnap predictive reliability & n/a & \hyperref[thm:p15-carnap-frequency-reliability]{\ref*{thm:p15-carnap-frequency-reliability}} & \thepage & Nonnegative only under stated extendability condition Blank macro filled with exact unchanged glyph spelling; registry data only, not a new TeX command. Original precedence 'atomic' retained as semantic descriptor; n/a is not an invented rank. \\
\phantomsection\label{entry:symbol-registry-05930}P(A;B) & \texttt{\detokenize{P(A;B)}} & reference-first relative probability & P15 Reichenbach & n/a & \hyperref[def:p15-reichenbach-probability-implication]{\ref*{def:p15-reichenbach-probability-implication}} & \thepage & Commas in reader; not joint probability Blank macro filled with exact unchanged glyph spelling; registry data only, not a new TeX command. Original precedence 'function application' retained as semantic descriptor; n/a is not an invented rank. \\
\phantomsection\label{entry:symbol-registry-05931}\ensuremath{A superset_p B} & \texttt{\detokenize{A superset_p B}} & probability implication & P15 Reichenbach & n/a & \hyperref[def:p15-reichenbach-probability-implication]{\ref*{def:p15-reichenbach-probability-implication}} & \thepage & Distinct from material implication Blank macro filled with exact unchanged glyph spelling; registry data only, not a new TeX command. Original precedence 'relational' retained as semantic descriptor; n/a is not an invented rank. \\
\phantomsection\label{entry:symbol-registry-05932}\ensuremath{F_n(A;B)} & \texttt{\detokenize{F_n(A;B)}} & finite selected frequency & P15 Reichenbach & n/a & \hyperref[def:p15-reichenbach-probability-implication]{\ref*{def:p15-reichenbach-probability-implication}} & \thepage & Nonzero reference count required Blank macro filled with exact unchanged glyph spelling; registry data only, not a new TeX command. Original precedence 'function application' retained as semantic descriptor; n/a is not an invented rank. \\
\phantomsection\label{entry:symbol-registry-05933}\ensuremath{x_i y_i} & \texttt{\detokenize{x_i y_i}} & ordered coordinated events & P15 Reichenbach & n/a & \hyperref[def:p15-reichenbach-probability-implication]{\ref*{def:p15-reichenbach-probability-implication}} & \thepage & Pairing and ordering specified Blank macro filled with exact unchanged glyph spelling; registry data only, not a new TeX command. Original precedence 'atomic' retained as semantic descriptor; n/a is not an invented rank. \\
\phantomsection\label{entry:symbol-registry-05934}bar A & \texttt{\detokenize{bar A}} & complement or typed empty-class predicate & P15 original schemata & n/a & \hyperref[sec:p15-reichenbach-relational-calculus]{\ref*{sec:p15-reichenbach-relational-calculus}} & \thepage & Parenthesized emptiness role distinguished Blank macro filled with exact unchanged glyph spelling; registry data only, not a new TeX command. Original precedence 'unary' retained as semantic descriptor; n/a is not an invented rank. \\
\phantomsection\label{entry:symbol-registry-05935}u v w & \texttt{\detokenize{u v w}} & marginal and joint reconstruction parameters & P15 coupling theorem & n/a & \hyperref[thm:p15-reichenbach-coupling]{\ref*{thm:p15-reichenbach-coupling}} & \thepage & Original French table p q u separately retained Blank macro filled with exact unchanged glyph spelling; registry data only, not a new TeX command. Original precedence 'atomic' retained as semantic descriptor; n/a is not an invented rank. \\
\phantomsection\label{entry:symbol-registry-05936}r & \texttt{\detokenize{r}} & positive-reference coupling & P15 coupling theorem & n/a & \hyperref[thm:p15-reichenbach-coupling]{\ref*{thm:p15-reichenbach-coupling}} & \thepage & r=w/u for u positive Blank macro filled with exact unchanged glyph spelling; registry data only, not a new TeX command. Original precedence 'atomic' retained as semantic descriptor; n/a is not an invented rank. \\
\phantomsection\label{entry:symbol-registry-05937}lambda k & \texttt{\detokenize{lambda k}} & division length and phase & P15 sequence classification & n/a & \hyperref[sec:p15-reichenbach-normality]{\ref*{sec:p15-reichenbach-normality}} & \thepage & Not a Carnap continuum parameter Blank macro filled with exact unchanged glyph spelling; registry data only, not a new TeX command. Original precedence 'atomic' retained as semantic descriptor; n/a is not an invented rank. \\
\phantomsection\label{entry:symbol-registry-05938}N & \texttt{\detokenize{N}} & extensional necessity & P15 modal table & n/a & \hyperref[thm:p15-reichenbach-modal-division]{\ref*{thm:p15-reichenbach-modal-division}} & \thepage & All indexed propositions true Blank macro filled with exact unchanged glyph spelling; registry data only, not a new TeX command. Original precedence 'atomic' retained as semantic descriptor; n/a is not an invented rank. \\
\phantomsection\label{entry:symbol-registry-05939}Pos & \texttt{\detokenize{Pos}} & extensional possibility & P15 modal table & n/a & \hyperref[thm:p15-reichenbach-modal-division]{\ref*{thm:p15-reichenbach-modal-division}} & \thepage & Both kinds occur; excludes necessity Blank macro filled with exact unchanged glyph spelling; registry data only, not a new TeX command. Original precedence 'atomic' retained as semantic descriptor; n/a is not an invented rank. \\
\phantomsection\label{entry:symbol-registry-05940}J calligraphic & \texttt{\detokenize{J calligraphic}} & extensional impossibility & P15 modal table & n/a & \hyperref[thm:p15-reichenbach-modal-division]{\ref*{thm:p15-reichenbach-modal-division}} & \thepage & All indexed propositions false Blank macro filled with exact unchanged glyph spelling; registry data only, not a new TeX command. Original precedence 'atomic' retained as semantic descriptor; n/a is not an invented rank. \\
\phantomsection\label{entry:symbol-registry-05941}tau & \texttt{\detokenize{tau}} & division threshold & P15 threshold theorem & n/a & \hyperref[thm:p15-reichenbach-modal-division]{\ref*{thm:p15-reichenbach-modal-division}} & \thepage & No Boolean preservation claim Blank macro filled with exact unchanged glyph spelling; registry data only, not a new TeX command. Original precedence 'atomic' retained as semantic descriptor; n/a is not an invented rank. \\
\phantomsection\label{entry:symbol-registry-05942}\ensuremath{f_n} & \texttt{\detokenize{f_n}} & empirical binary frequency & P15 straight rule & n/a & \hyperref[thm:p15-reichenbach-straight-rule-boundary]{\ref*{thm:p15-reichenbach-straight-rule-boundary}} & \thepage & Continuation posit separate Blank macro filled with exact unchanged glyph spelling; registry data only, not a new TeX command. Original precedence 'atomic' retained as semantic descriptor; n/a is not an invented rank. \\
\phantomsection\label{entry:symbol-registry-05943}\ensuremath{clip_[0;1]} & \texttt{\detokenize{clip_[0;1]}} & bounded rival-rule projection & P15 editorial convergence example & n/a & \hyperref[thm:p15-reichenbach-straight-rule-boundary]{\ref*{thm:p15-reichenbach-straight-rule-boundary}} & \thepage & Not a source inference rule Blank macro filled with exact unchanged glyph spelling; registry data only, not a new TeX command. Original precedence 'function application' retained as semantic descriptor; n/a is not an invented rank. \\
\phantomsection\label{entry:symbol-registry-05944}\ensuremath{g_c} & \texttt{\detokenize{g_c}} & alternative exact-fitting law & P15 historical failure case & n/a & \hyperref[thm:p15-reichenbach-finite-fit]{\ref*{thm:p15-reichenbach-finite-fit}} & \thepage & Finite real-input assumptions Blank macro filled with exact unchanged glyph spelling; registry data only, not a new TeX command. Original precedence 'function application' retained as semantic descriptor; n/a is not an invented rank. \\
\phantomsection\label{entry:symbol-registry-06358}\(\ssatrandom^{p}\) & \texttt{\detokenize{\ssatrandom^{p}}} & native rotated-R independent weighted Boolean random quantifier & P15 stochastic satisfiability & n/a & \hyperref[def:p15-ssat-input]{\ref*{def:p15-ssat-input}} & \thepage & Native Roman R rotated180 degrees, with parameter p for the declared fixed independent Bernoulli extension. Fair source operator is the separately recorded ssatrandom without annotation. Original upright R\^{}p readable-alias row remains immutable raw evidence; transformed by root's verified source-glyph repair. Original precedence descriptor 'quantifier' retained; no numerical rank inferred. \\
\phantomsection\label{entry:symbol-registry-06359}\(\Pi\) & \texttt{\detokenize{\Pi}} & ordered quantifier prefix & P15 stochastic satisfiability & n/a & \hyperref[def:p15-ssat-input]{\ref*{def:p15-ssat-input}} & \thepage & Information boundary part of input Actual displayed reader TeX spelling; original ASCII intake remains raw evidence. Original precedence descriptor 'prefix' retained; no numerical rank inferred. \\
\phantomsection\label{entry:symbol-registry-06360}\(\phi|_{x=b}\) & \texttt{\detokenize{\phi|_{x=b}}} & Boolean CNF restriction & P15 stochastic satisfiability & n/a & \hyperref[def:p15-ssat-input]{\ref*{def:p15-ssat-input}} & \thepage & Retains relative order of remaining prefix Actual displayed reader TeX spelling; original ASCII intake remains raw evidence. Original precedence descriptor 'restriction' retained; no numerical rank inferred. \\
\phantomsection\label{entry:symbol-registry-06361}\(\bar x\) & \texttt{\detokenize{\bar x}} & literal complement & P15 stochastic satisfiability & n/a & \hyperref[def:p15-ssat-input]{\ref*{def:p15-ssat-input}} & \thepage & Boolean negation not probabilistic complementation Actual displayed reader TeX spelling; original ASCII intake remains raw evidence. Original precedence descriptor 'unary' retained; no numerical rank inferred. \\
\phantomsection\label{entry:symbol-registry-06362}\(V(\Pi,\phi)\) & \texttt{\detokenize{V(\Pi,}\allowbreak\detokenize{\phi)}} & stochastic success value & P15 stochastic satisfiability & n/a & \hyperref[def:p15-ssat-value]{\ref*{def:p15-ssat-value}} & \thepage & Max min and weighted recurrence; not null-event conditional Actual displayed reader TeX spelling; original ASCII intake remains raw evidence. Original precedence descriptor 'function application' retained; no numerical rank inferred. \\
\phantomsection\label{entry:symbol-registry-06363}\(\theta\) & \texttt{\detokenize{\theta}} & inclusive satisfaction threshold & P15 stochastic satisfiability & n/a & \hyperref[def:p15-ssat-input]{\ref*{def:p15-ssat-input}} & \thepage & Decision compares V>=theta Actual displayed reader TeX spelling; original ASCII intake remains raw evidence. Original precedence descriptor 'scalar' retained; no numerical rank inferred. \\
\phantomsection\label{entry:symbol-registry-06364}\(\Downarrow\) & \texttt{\detokenize{\Downarrow}} & value-evaluation judgment & P15 stochastic satisfiability & relation & \hyperref[def:p15-ssat-evaluation-judgment]{\ref*{def:p15-ssat-evaluation-judgment}} & \thepage & Finite canonical certificate interpretation Actual displayed reader TeX spelling; original ASCII intake remains raw evidence. Real judgment definition before the literal sheet; no inferred definition from value alone. \\
\phantomsection\label{entry:symbol-registry-06365}\(\operatorname{Eval}\) & \texttt{\detokenize{\operatorname{Eval}}} & corrected threshold evaluator & P15 stochastic satisfiability & n/a & \hyperref[def:p15-ssat-window-evaluator]{\ref*{def:p15-ssat-window-evaluator}} & \thepage & Strict high early return supports closed-window exactness Actual displayed reader TeX spelling; original ASCII intake remains raw evidence. Original precedence descriptor 'function application' retained; no numerical rank inferred. \\
\phantomsection\label{entry:symbol-registry-06366}\(L,H\) & \texttt{\detokenize{L,}\allowbreak\detokenize{H}} & threshold-window endpoint symbol family & P15 stochastic satisfiability & n/a & \hyperref[def:p15-ssat-window-contract]{\ref*{def:p15-ssat-window-contract}} & \thepage & Recursive windows may be outside probability interval Family listing, not one atomic glyph: L lower threshold-window endpoint; H upper threshold-window endpoint. Original precedence descriptor 'scalar' retained; no numerical rank inferred. \\
\phantomsection\label{entry:symbol-registry-06367}\(p_b\) & \texttt{\detokenize{p_b}} & weight of forced random unit outcome & P15 stochastic satisfiability & n/a & \hyperref[thm:p15-ssat-unit-pure]{\ref*{thm:p15-ssat-unit-pure}} & \thepage & Zero weight returns zero without division Actual displayed reader TeX spelling; original ASCII intake remains raw evidence. Original precedence descriptor 'scalar' retained; no numerical rank inferred. \\
\phantomsection\label{entry:symbol-registry-06368}\(\mu\) & \texttt{\detokenize{\mu}} & product mass on random assignments & P15 stochastic satisfiability & n/a & \hyperref[thm:p15-ssat-projections]{\ref*{thm:p15-ssat-projections}} & \thepage & Independent specified draws Actual displayed reader TeX spelling; original ASCII intake remains raw evidence. Original precedence descriptor 'function application' retained; no numerical rank inferred. \\
\phantomsection\label{entry:symbol-registry-06369}\(C_{i,u}\) & \texttt{\detokenize{C_{i,u}}} & Bayesian CPT row coin & P15 stochastic satisfiability & n/a & \hyperref[def:p15-ssat-row-coin-encoding]{\ref*{def:p15-ssat-row-coin-encoding}} & \thepage & Independent coins selected by deterministic parent patterns Actual displayed reader TeX spelling; original ASCII intake remains raw evidence. Indexed symbol family at the stated indices, not the former hyphenated ASCII alias. Original precedence descriptor 'Boolean variable' retained; no numerical rank inferred. \\
\phantomsection\label{entry:symbol-registry-06370}\(\Psi\) & \texttt{\detokenize{\Psi}} & deterministic network and gate constraints & P15 stochastic satisfiability & n/a & \hyperref[def:p15-ssat-row-coin-encoding]{\ref*{def:p15-ssat-row-coin-encoding}} & \thepage & Final existential auxiliaries have unique extension Actual displayed reader TeX spelling; original ASCII intake remains raw evidence. Original precedence descriptor 'formula' retained; no numerical rank inferred. \\
\phantomsection\label{entry:symbol-registry-06371}\(\widehat v_s\) & \texttt{\detokenize{\widehat v_s}} & empirical complete-policy value & P15 stochastic satisfiability & n/a & \hyperref[def:p15-ssat-sample-policy]{\ref*{def:p15-ssat-sample-policy}} & \thepage & Sample joint history distribution retained Actual displayed reader TeX spelling; original ASCII intake remains raw evidence. Indexed symbol family at the stated indices, not the former hyphenated ASCII alias. Original precedence descriptor 'scalar' retained; no numerical rank inferred. \\
\phantomsection\label{entry:symbol-registry-06372}\(D\) & \texttt{\detokenize{D}} & number of complete policy table entries & P15 stochastic satisfiability & n/a & \hyperref[def:p15-ssat-sample-policy]{\ref*{def:p15-ssat-sample-policy}} & \thepage & Can be exponential in prefix length Actual displayed reader TeX spelling; original ASCII intake remains raw evidence. Original precedence descriptor 'scalar' retained; no numerical rank inferred. \\
\phantomsection\label{entry:symbol-registry-06373}\(P\) & \texttt{\detokenize{P}} & number of labelled policies & P15 stochastic satisfiability & n/a & \hyperref[def:p15-ssat-sample-policy]{\ref*{def:p15-ssat-sample-policy}} & \thepage & 2\^{}D; distinct from probability operator Actual displayed reader TeX spelling; original ASCII intake remains raw evidence. Original precedence descriptor 'scalar' retained; no numerical rank inferred. \\
\phantomsection\label{entry:symbol-registry-06374}\(N_{h,b}\) & \texttt{\detokenize{N_{h,b}}} & history and final-bit sample multiplicity & P15 stochastic satisfiability & n/a & \hyperref[thm:p15-ssat-sampling-overfit]{\ref*{thm:p15-ssat-sampling-overfit}} & \thepage & Majority chosen per earlier history Actual displayed reader TeX spelling; original ASCII intake remains raw evidence. Indexed symbol family at the stated indices, not the former hyphenated ASCII alias. Original precedence descriptor 'scalar' retained; no numerical rank inferred. \\
\phantomsection\label{entry:symbol-registry-06375}\(F_m^{k,n}\) & \texttt{\detokenize{F_m^{k,n}}} & fixed-clause random instance law & P15 stochastic satisfiability & n/a & \hyperref[sec:p15-ssat-random-instances]{\ref*{sec:p15-ssat-random-instances}} & \thepage & Distinct variables within clauses; not inner random quantification Actual displayed reader TeX spelling; original ASCII intake remains raw evidence. Indexed symbol family at the stated indices, not the former hyphenated ASCII alias. Original precedence descriptor 'distribution' retained; no numerical rank inferred. \\
\phantomsection\label{entry:symbol-registry-06376}\(Z\) & \texttt{\detokenize{Z}} & number of satisfying assignments & P15 stochastic satisfiability & n/a & \hyperref[thm:p15-ssat-random-formula-bound]{\ref*{thm:p15-ssat-random-formula-bound}} & \thepage & Randomness here is formula generation Actual displayed reader TeX spelling; original ASCII intake remains raw evidence. Original precedence descriptor 'random variable' retained; no numerical rank inferred. \\
\phantomsection\label{entry:symbol-registry-06377}\(\ssatrandom\) & \texttt{\detokenize{\ssatrandom}} & native fair randomized Boolean quantifier rendered as a Roman R rotated180 degrees & P15 SSAT native syntax & n/a & \hyperref[def:p15-ssat-native]{\ref*{def:p15-ssat-native}} & \thepage & Actual source glyph inspected on original printedp4; not upright R. \\
\phantomsection\label{entry:symbol-registry-06378}\(Q\) & \texttt{\detokenize{Q}} & native ordered-variable map to existential universal or fair random quantifiers & P15 SSAT native syntax & n/a & \hyperref[def:p15-ssat-native]{\ref*{def:p15-ssat-native}} & \thepage & Renumbering transports the map by the same order-preserving variable bijection. \\
\phantomsection\label{entry:symbol-registry-06379}\(\operatorname{val}\) & \texttt{\detokenize{\operatorname{val}}} & native fair Extended SSAT matrix-and-quantifier-map value & P15 SSAT native syntax & n/a & \hyperref[def:p15-ssat-native]{\ref*{def:p15-ssat-native}} & \thepage & Explicit correspondence with the weighted prefix evaluator is proved at the fair subcase. \\
\phantomsection\label{entry:symbol-registry-06380}\(\phi\upharpoonright_{x=b}\) & \texttt{\detokenize{\phi\upharpoonright_{x=b}}} & native Boolean restriction including order-preserving renumbering & P15 SSAT native syntax & n/a & \hyperref[def:p15-ssat-native]{\ref*{def:p15-ssat-native}} & \thepage & The companion cofactor notation retains old names; values agree under the proved variable correspondence. \\
\phantomsection\label{entry:symbol-registry-06381}\(\boldsymbol t\) & \texttt{\detokenize{\boldsymbol t}} & ordered vector of the declared row-coin probabilities & P15 stochastic satisfiability & n/a & \hyperref[def:p15-ssat-row-coin-encoding]{\ref*{def:p15-ssat-row-coin-encoding}} & \thepage & Defined in the same chosen order as the row coins before the weighted block is used. \\
\phantomsection\label{entry:symbol-registry-07437}\(A\) & \texttt{\detokenize{A}} & finite nonempty alphabet and finite-word universe symbol family member \(A\) & P15 formal identification & object & \hyperref[def:p15-flt-names]{\ref*{def:p15-flt-names}} & \thepage & Gold Sigma A and host A* denote the same word set; empty word is an ordinary word. Explicit displayed family; member types and common actual first-use exposed, not one atomic glyph. Family listing, not one atomic glyph. Author macro \texttt{\detokenize{A;\Sigma;\mathcal W}} not carried to the members. \\
\phantomsection\label{entry:symbol-registry-07438}\(\Sigma A\) & \texttt{\detokenize{\Sigma A}} & finite nonempty alphabet and finite-word universe symbol family member \(\Sigma A\) & P15 formal identification & object & \hyperref[def:p15-flt-names]{\ref*{def:p15-flt-names}} & \thepage & Gold Sigma A and host A* denote the same word set; empty word is an ordinary word. Explicit displayed family; member types and common actual first-use exposed, not one atomic glyph. Family listing, not one atomic glyph. Author macro \texttt{\detokenize{A;\Sigma;\mathcal W}} not carried to the members. \\
\phantomsection\label{entry:symbol-registry-07439}\(\mathcal W\) & \texttt{\detokenize{\mathcal W}} & finite nonempty alphabet and finite-word universe symbol family member \(\mathcal W\) & P15 formal identification & object & \hyperref[def:p15-flt-names]{\ref*{def:p15-flt-names}} & \thepage & Gold Sigma A and host A* denote the same word set; empty word is an ordinary word. Explicit displayed family; member types and common actual first-use exposed, not one atomic glyph. Family listing, not one atomic glyph. Author macro \texttt{\detokenize{A;\Sigma;\mathcal W}} not carried to the members. \\
\phantomsection\label{entry:symbol-registry-07440}\(b\) & \texttt{\detokenize{b}} & effective word enumeration and empty word symbol family member \(b\) & P15 formal identification & meta & \hyperref[def:p15-flt-names]{\ref*{def:p15-flt-names}} & \thepage & Word/integer coding is primitive recursive for the source primitive-recursive classes. Explicit displayed family; member types and common actual first-use exposed, not one atomic glyph. Family listing, not one atomic glyph. Author macro \texttt{\detokenize{b;\varepsilon}} not carried to the members. \\
\phantomsection\label{entry:symbol-registry-07441}\(\varepsilon\) & \texttt{\detokenize{\varepsilon}} & effective word enumeration and empty word symbol family member \(\varepsilon\) & P15 formal identification & meta & \hyperref[def:p15-flt-names]{\ref*{def:p15-flt-names}} & \thepage & Word/integer coding is primitive recursive for the source primitive-recursive classes. Explicit displayed family; member types and common actual first-use exposed, not one atomic glyph. Family listing, not one atomic glyph. Author macro \texttt{\detokenize{b;\varepsilon}} not carried to the members. \\
\phantomsection\label{entry:symbol-registry-07442}\(\Gamma\) & \texttt{\detokenize{\Gamma}} & hypothesis naming map and natural hypothesis code symbol family member \(\Gamma\) & P15 formal identification & meta & \hyperref[def:p15-flt-names]{\ref*{def:p15-flt-names}} & \thepage & Map can be partial and have synonyms; no uniform membership test is presumed. Explicit displayed family; member types and common actual first-use exposed, not one atomic glyph. Family listing, not one atomic glyph. Author macro \texttt{\detokenize{\Gamma;e}} not carried to the members. \\
\phantomsection\label{entry:symbol-registry-07443}\(e\) & \texttt{\detokenize{e}} & hypothesis naming map and natural hypothesis code symbol family member \(e\) & P15 formal identification & meta & \hyperref[def:p15-flt-names]{\ref*{def:p15-flt-names}} & \thepage & Map can be partial and have synonyms; no uniform membership test is presumed. Explicit displayed family; member types and common actual first-use exposed, not one atomic glyph. Family listing, not one atomic glyph. Author macro \texttt{\detokenize{\Gamma;e}} not carried to the members. \\
\phantomsection\label{entry:symbol-registry-07444}\(T\) & \texttt{\detokenize{T}} & positive presentation and host pause symbol symbol family member \(T\) & P15 formal identification & object & \hyperref[def:p15-flt-data]{\ref*{def:p15-flt-data}} & \thepage & Pause is not a word; source pure texts are separated from the empty-language host extension. Explicit displayed family; member types and common actual first-use exposed, not one atomic glyph. Family listing, not one atomic glyph. Author macro \texttt{\detokenize{T;\#}} not carried to the members. \\
\phantomsection\label{entry:symbol-registry-07445}\(\#\) & \texttt{\detokenize{\#}} & positive presentation and host pause symbol symbol family member \(\#\) & P15 formal identification & object & \hyperref[def:p15-flt-data]{\ref*{def:p15-flt-data}} & \thepage & Pause is not a word; source pure texts are separated from the empty-language host extension. Explicit displayed family; member types and common actual first-use exposed, not one atomic glyph. Family listing, not one atomic glyph. Author macro \texttt{\detokenize{T;\#}} not carried to the members. \\
\phantomsection\label{entry:symbol-registry-07446}\(G\) & \texttt{\detokenize{G}} & prefix learner guesses observations and abstention symbol family member \(G\) & P15 formal identification & meta & \hyperref[def:p15-flt-ex]{\ref*{def:p15-flt-ex}} & \thepage & Gold native g\_t=G(i\_1...i\_t); syntactic stability and semantic correctness are distinct. Explicit displayed family; member types and common actual first-use exposed, not one atomic glyph. Family listing, not one atomic glyph. Author macro \texttt{\detokenize{G;g\_t;i\_t;\bot}} not carried to the members. \\
\phantomsection\label{entry:symbol-registry-07447}\(g_t\) & \texttt{\detokenize{g_t}} & prefix learner guesses observations and abstention symbol family member \(g_t\) & P15 formal identification & meta & \hyperref[def:p15-flt-ex]{\ref*{def:p15-flt-ex}} & \thepage & Gold native g\_t=G(i\_1...i\_t); syntactic stability and semantic correctness are distinct. Explicit displayed family; member types and common actual first-use exposed, not one atomic glyph. Family listing, not one atomic glyph. Author macro \texttt{\detokenize{G;g\_t;i\_t;\bot}} not carried to the members. \\
\phantomsection\label{entry:symbol-registry-07448}\(i_t\) & \texttt{\detokenize{i_t}} & prefix learner guesses observations and abstention symbol family member \(i_t\) & P15 formal identification & meta & \hyperref[def:p15-flt-ex]{\ref*{def:p15-flt-ex}} & \thepage & Gold native g\_t=G(i\_1...i\_t); syntactic stability and semantic correctness are distinct. Explicit displayed family; member types and common actual first-use exposed, not one atomic glyph. Family listing, not one atomic glyph. Author macro \texttt{\detokenize{G;g\_t;i\_t;\bot}} not carried to the members. \\
\phantomsection\label{entry:symbol-registry-07449}\(\bot\) & \texttt{\detokenize{\bot}} & prefix learner guesses observations and abstention symbol family member \(\bot\) & P15 formal identification & meta & \hyperref[def:p15-flt-ex]{\ref*{def:p15-flt-ex}} & \thepage & Gold native g\_t=G(i\_1...i\_t); syntactic stability and semantic correctness are distinct. Explicit displayed family; member types and common actual first-use exposed, not one atomic glyph. Family listing, not one atomic glyph. Author macro \texttt{\detokenize{G;g\_t;i\_t;\bot}} not carried to the members. \\
\phantomsection\label{entry:symbol-registry-07450}\(T{\upharpoonright}n\) & \texttt{\detokenize{\upharpoonright}} & finite observed prefix & P15 formal identification & meta & \hyperref[def:p15-flt-ex]{\ref*{def:p15-flt-ex}} & \thepage & Length counts observations including host pauses. \\
\phantomsection\label{entry:symbol-registry-07451}\(D\) & \texttt{\detokenize{D}} & total joint membership decider and indexed family symbol family member \(D\) & P15 positive informant result & meta & \hyperref[def:p15-flt-uniform-family]{\ref*{def:p15-flt-uniform-family}} & \thepage & Individual recursive languages do not supply a uniformly total list; family indices and compiled program names are distinguished. Explicit displayed family; member types and common actual first-use exposed, not one atomic glyph. Family listing, not one atomic glyph. Author macro \texttt{\detokenize{D;H\_e}} not carried to the members. \\
\phantomsection\label{entry:symbol-registry-07452}\(H_e\) & \texttt{\detokenize{H_e}} & total joint membership decider and indexed family symbol family member \(H_e\) & P15 positive informant result & meta & \hyperref[def:p15-flt-uniform-family]{\ref*{def:p15-flt-uniform-family}} & \thepage & Individual recursive languages do not supply a uniformly total list; family indices and compiled program names are distinguished. Explicit displayed family; member types and common actual first-use exposed, not one atomic glyph. Family listing, not one atomic glyph. Author macro \texttt{\detokenize{D;H\_e}} not carried to the members. \\
\phantomsection\label{entry:symbol-registry-07453}\(d(F)\) & \texttt{\detokenize{d(F)}} & finite-set bit parameter and deterministic tester compiler symbol family member \(d(F)\) & P15 finite-text learner & meta & \hyperref[thm:p15-flt-finite-text]{\ref*{thm:p15-flt-finite-text}} & \thepage & Trace values are parameters not raw program indices. Explicit displayed family; member types and common actual first-use exposed, not one atomic glyph. Family listing, not one atomic glyph. Author macro \texttt{\detokenize{d;c}} not carried to the members. \\
\phantomsection\label{entry:symbol-registry-07454}\(c(d)\) & \texttt{\detokenize{c(d)}} & finite-set bit parameter and deterministic tester compiler symbol family member \(c(d)\) & P15 finite-text learner & meta & \hyperref[thm:p15-flt-finite-text]{\ref*{thm:p15-flt-finite-text}} & \thepage & Trace values are parameters not raw program indices. Explicit displayed family; member types and common actual first-use exposed, not one atomic glyph. Family listing, not one atomic glyph. Author macro \texttt{\detokenize{d;c}} not carried to the members. \\
\phantomsection\label{entry:symbol-registry-07455}\(p(e,n)\) & \texttt{\detokenize{p(e,}\allowbreak\detokenize{n)}} & primitive-recursive-description interpreter and source Boolean map symbol family member \(p(e,n)\) & P15 primitive-recursive positive results & meta & \hyperref[annex:p15-flt-enumeration]{\ref*{annex:p15-flt-enumeration}} & \thepage & Joint interpreter is total recursive not claimed primitive recursive; w selects exactly output one. Explicit displayed family; member types and common actual first-use exposed, not one atomic glyph. Family listing, not one atomic glyph. Author macro \texttt{\detokenize{p;w}} not carried to the members. \\
\phantomsection\label{entry:symbol-registry-07456}\(w(z)\) & \texttt{\detokenize{w(z)}} & primitive-recursive-description interpreter and source Boolean map symbol family member \(w(z)\) & P15 primitive-recursive positive results & meta & \hyperref[annex:p15-flt-enumeration]{\ref*{annex:p15-flt-enumeration}} & \thepage & Joint interpreter is total recursive not claimed primitive recursive; w selects exactly output one. Explicit displayed family; member types and common actual first-use exposed, not one atomic glyph. Family listing, not one atomic glyph. Author macro \texttt{\detokenize{p;w}} not carried to the members. \\
\phantomsection\label{entry:symbol-registry-07457}\(a_j\) & \texttt{\detokenize{a_j}} & repetition-free target enumeration prefix and diagonal block lengths symbol family member \(a_j\) & P15 superfinite diagonal & meta & \hyperref[thm:p15-flt-superfinite]{\ref*{thm:p15-flt-superfinite}} & \thepage & Block search compares actual hypothesis codes; no semantic equality test. Explicit displayed family; member types and common actual first-use exposed, not one atomic glyph. Family listing, not one atomic glyph. Author macro \texttt{\detokenize{a\_j;u;x;y}} not carried to the members. \\
\phantomsection\label{entry:symbol-registry-07458}\(u\) & \texttt{\detokenize{u}} & repetition-free target enumeration prefix and diagonal block lengths symbol family member \(u\) & P15 superfinite diagonal & meta & \hyperref[thm:p15-flt-superfinite]{\ref*{thm:p15-flt-superfinite}} & \thepage & Block search compares actual hypothesis codes; no semantic equality test. Explicit displayed family; member types and common actual first-use exposed, not one atomic glyph. Family listing, not one atomic glyph. Author macro \texttt{\detokenize{a\_j;u;x;y}} not carried to the members. \\
\phantomsection\label{entry:symbol-registry-07459}\(x\) & \texttt{\detokenize{x}} & repetition-free target enumeration prefix and diagonal block lengths symbol family member \(x\) & P15 superfinite diagonal & meta & \hyperref[thm:p15-flt-superfinite]{\ref*{thm:p15-flt-superfinite}} & \thepage & Block search compares actual hypothesis codes; no semantic equality test. Explicit displayed family; member types and common actual first-use exposed, not one atomic glyph. Family listing, not one atomic glyph. Author macro \texttt{\detokenize{a\_j;u;x;y}} not carried to the members. \\
\phantomsection\label{entry:symbol-registry-07460}\(y\) & \texttt{\detokenize{y}} & repetition-free target enumeration prefix and diagonal block lengths symbol family member \(y\) & P15 superfinite diagonal & meta & \hyperref[thm:p15-flt-superfinite]{\ref*{thm:p15-flt-superfinite}} & \thepage & Block search compares actual hypothesis codes; no semantic equality test. Explicit displayed family; member types and common actual first-use exposed, not one atomic glyph. Family listing, not one atomic glyph. Author macro \texttt{\detokenize{a\_j;u;x;y}} not carried to the members. \\
\phantomsection\label{entry:symbol-registry-07461}\(L^\dagger\) & \texttt{\detokenize{L^\dagger}} & diagonal recursive language escaping a purported uniform list & P15 indexing obstruction & object & \hyperref[annex:p15-flt-no-uniform-list]{\ref*{annex:p15-flt-no-uniform-list}} & \thepage & Not a claim of effective enumeration of every recursive language. \\
\phantomsection\label{entry:symbol-registry-07462}\(Q(i,s)\) & \texttt{\detokenize{Q}} & primitive-recursive accepting-computation witness predicate & P15 restricted PR-text existence proof & meta & \hyperref[annex:p15-flt-pr-text]{\ref*{annex:p15-flt-pr-text}} & \thepage & Uses existing normal form; fixed default word requires a nonempty target. \\
\phantomsection\label{entry:symbol-registry-07895}\(i_t\) & \texttt{\detokenize{i_t}} & positive-word observations finite prefix and positive stage symbol family member \(i_t\) & P15 primitive-recursive tester learning & object & \hyperref[def:p15-pr-learning-codes]{\ref*{def:p15-pr-learning-codes}} & \thepage & Gold1967 pp471–473 I.14–I.26; no negative labels or pause symbols; word coding agrees with frozen definition through an explicit positive-integer shift. Explicit displayed family; member types and common actual first-use exposed, not one atomic glyph. Family listing, not one atomic glyph. Author macro \texttt{\detokenize{i\_t;\sigma;t}} not carried to the members. \\
\phantomsection\label{entry:symbol-registry-07896}\(\sigma\) & \texttt{\detokenize{\sigma}} & positive-word observations finite prefix and positive stage symbol family member \(\sigma\) & P15 primitive-recursive tester learning & object & \hyperref[def:p15-pr-learning-codes]{\ref*{def:p15-pr-learning-codes}} & \thepage & Gold1967 pp471–473 I.14–I.26; no negative labels or pause symbols; word coding agrees with frozen definition through an explicit positive-integer shift. Explicit displayed family; member types and common actual first-use exposed, not one atomic glyph. Family listing, not one atomic glyph. Author macro \texttt{\detokenize{i\_t;\sigma;t}} not carried to the members. \\
\phantomsection\label{entry:symbol-registry-07897}\(t\) & \texttt{\detokenize{t}} & positive-word observations finite prefix and positive stage symbol family member \(t\) & P15 primitive-recursive tester learning & object & \hyperref[def:p15-pr-learning-codes]{\ref*{def:p15-pr-learning-codes}} & \thepage & Gold1967 pp471–473 I.14–I.26; no negative labels or pause symbols; word coding agrees with frozen definition through an explicit positive-integer shift. Explicit displayed family; member types and common actual first-use exposed, not one atomic glyph. Family listing, not one atomic glyph. Author macro \texttt{\detokenize{i\_t;\sigma;t}} not carried to the members. \\
\phantomsection\label{entry:symbol-registry-07898}\(Z_a(x)\) & \texttt{\detokenize{Z_a(x)}} & partial numerical tester program hypothesis code and word code symbol family member \(Z_a(x)\) & P15 native Gold tester interface & meta & \hyperref[def:p15-pr-learning-codes]{\ref*{def:p15-pr-learning-codes}} & \thepage & Native Z\_a from Gold1967 pp472–473; correct tester is total Boolean membership; intermediate programs may be partial or malformed. Explicit displayed family; member types and common actual first-use exposed, not one atomic glyph. Family listing, not one atomic glyph. Author macro \texttt{\detokenize{Z\_a;a;x}} not carried to the members. \\
\phantomsection\label{entry:symbol-registry-07899}\(a\) & \texttt{\detokenize{a}} & partial numerical tester program hypothesis code and word code symbol family member \(a\) & P15 native Gold tester interface & meta & \hyperref[def:p15-pr-learning-codes]{\ref*{def:p15-pr-learning-codes}} & \thepage & Native Z\_a from Gold1967 pp472–473; correct tester is total Boolean membership; intermediate programs may be partial or malformed. Explicit displayed family; member types and common actual first-use exposed, not one atomic glyph. Family listing, not one atomic glyph. Author macro \texttt{\detokenize{Z\_a;a;x}} not carried to the members. \\
\phantomsection\label{entry:symbol-registry-07900}\(x\) & \texttt{\detokenize{x}} & partial numerical tester program hypothesis code and word code symbol family member \(x\) & P15 native Gold tester interface & meta & \hyperref[def:p15-pr-learning-codes]{\ref*{def:p15-pr-learning-codes}} & \thepage & Native Z\_a from Gold1967 pp472–473; correct tester is total Boolean membership; intermediate programs may be partial or malformed. Explicit displayed family; member types and common actual first-use exposed, not one atomic glyph. Family listing, not one atomic glyph. Author macro \texttt{\detokenize{Z\_a;a;x}} not carried to the members. \\
\phantomsection\label{entry:symbol-registry-07901}\(G_0\) & \texttt{\detokenize{G_0}} & original total computable finite-prefix learner & P15 local normalization input & meta & \hyperref[def:p15-pr-learning-codes]{\ref*{def:p15-pr-learning-codes}} & \thepage & Local bridge separates the original rule from native G after normalization; no uniformly total enumeration of all computable learners. \\
\phantomsection\label{entry:symbol-registry-07902}\(T(a,x,y)\) & \texttt{\detokenize{T}} & primitive-recursive complete halting-history verification predicate & P15 native Gold computation interface & meta & \hyperref[def:p15-pr-learning-codes]{\ref*{def:p15-pr-learning-codes}} & \thepage & Source selector Gold1967 p471 I.16 and following definition; numerical three-place predicate distinct from text sequence T in frozen chapter and machine tape T in P4. \\
\phantomsection\label{entry:symbol-registry-07903}\(U(y)\) & \texttt{\detokenize{U(y)}} & primitive-recursive numerical output extraction and computation-history number symbol family member \(U(y)\) & P15 native Gold computation interface & meta & \hyperref[def:p15-pr-learning-codes]{\ref*{def:p15-pr-learning-codes}} & \thepage & Source selector Gold1967 p471 I.16 and following definition; y bound measures a history code not a step count; malformed inputs have a total default. Explicit displayed family; member types and common actual first-use exposed, not one atomic glyph. Family listing, not one atomic glyph. Author macro \texttt{\detokenize{U;y}} not carried to the members. \\
\phantomsection\label{entry:symbol-registry-07904}\(y\) & \texttt{\detokenize{y}} & primitive-recursive numerical output extraction and computation-history number symbol family member \(y\) & P15 native Gold computation interface & meta & \hyperref[def:p15-pr-learning-codes]{\ref*{def:p15-pr-learning-codes}} & \thepage & Source selector Gold1967 p471 I.16 and following definition; y bound measures a history code not a step count; malformed inputs have a total default. Explicit displayed family; member types and common actual first-use exposed, not one atomic glyph. Family listing, not one atomic glyph. Author macro \texttt{\detokenize{U;y}} not carried to the members. \\
\phantomsection\label{entry:symbol-registry-07905}\(\equiv\) & \texttt{\detokenize{\equiv}} & native predicate definition and Boolean conjunction symbol family member \(\equiv\) & P15 native Gold predicates & meta & \hyperref[def:p15-pr-learning-codes]{\ref*{def:p15-pr-learning-codes}} & \thepage & Actual source selectors Gold1967 p471 I.15–I.16; typed definitional equiv is not the house elementary-equivalence relation; source ampersand retained. Explicit displayed family; member types and common actual first-use exposed, not one atomic glyph. Family listing, not one atomic glyph. Author macro \texttt{\detokenize{\equiv;\&}} not carried to the members. \\
\phantomsection\label{entry:symbol-registry-07906}\(\&\) & \texttt{\detokenize{\&}} & native predicate definition and Boolean conjunction symbol family member \(\&\) & P15 native Gold predicates & meta & \hyperref[def:p15-pr-learning-codes]{\ref*{def:p15-pr-learning-codes}} & \thepage & Actual source selectors Gold1967 p471 I.15–I.16; typed definitional equiv is not the house elementary-equivalence relation; source ampersand retained. Explicit displayed family; member types and common actual first-use exposed, not one atomic glyph. Family listing, not one atomic glyph. Author macro \texttt{\detokenize{\equiv;\&}} not carried to the members. \\
\phantomsection\label{entry:symbol-registry-07907}\(\downarrow\) & \texttt{\detokenize{\downarrow}} & defined halting numerical output & P15 primitive-recursive tester learning & meta & \hyperref[def:p15-pr-learning-codes]{\ref*{def:p15-pr-learning-codes}} & \thepage & PR.1 expresses witness adequacy for the same Z\_a numbering; it is not a decision of termination. \\
\phantomsection\label{entry:symbol-registry-07908}\(\widehat G\) & \texttt{\detokenize{\widehat G}} & primitive-recursive serially normalized prefix rule & P15 local normalization output & meta & \hyperref[def:p15-pr-learning-relay]{\ref*{def:p15-pr-learning-relay}} & \thepage & Exact convergence and nonconvergence on every infinite sequence preserved with identical stabilized code; separately proved local construction not claimed native Gold1965 glyph. \\
\phantomsection\label{entry:symbol-registry-07909}\(m\) & \texttt{\detokenize{m}} & next prefix number and initial retained hypothesis symbol family member \(m\) & P15 serial normalization procedure & meta & \hyperref[def:p15-pr-learning-relay]{\ref*{def:p15-pr-learning-relay}} & \thepage & At most one completion per tick ensures m does not outrun available data. Explicit displayed family; member types and common actual first-use exposed, not one atomic glyph. Family listing, not one atomic glyph. Author macro \texttt{\detokenize{m;a\_{\mathrm{init}}}} not carried to the members. \\
\phantomsection\label{entry:symbol-registry-07910}\(a_{\mathrm{init}}\) & \texttt{\detokenize{a_{\mathrm{init}}}} & next prefix number and initial retained hypothesis symbol family member \(a_{\mathrm{init}}\) & P15 serial normalization procedure & meta & \hyperref[def:p15-pr-learning-relay]{\ref*{def:p15-pr-learning-relay}} & \thepage & At most one completion per tick ensures m does not outrun available data. Explicit displayed family; member types and common actual first-use exposed, not one atomic glyph. Family listing, not one atomic glyph. Author macro \texttt{\detokenize{m;a\_{\mathrm{init}}}} not carried to the members. \\
\phantomsection\label{entry:symbol-registry-07911}\(\longrightarrow a\) & \texttt{\detokenize{\longrightarrow}} & eventual constant equality to a hypothesis number & P15 normalization criterion & meta & \hyperref[thm:p15-pr-learning-normalization]{\ref*{thm:p15-pr-learning-normalization}} & \thepage & Natural-number syntactic stabilization not metric convergence; no known stabilization stage. \\
\phantomsection\label{entry:symbol-registry-07912}\(f\) & \texttt{\detokenize{f}} & total primitive-recursive positive-index range enumerator & P15 native Gold countertext & meta & \hyperref[thm:p15-pr-learning-range]{\ref*{thm:p15-pr-learning-range}} & \thepage & Source selector Gold1967 p471 before I.15; repetitions allowed; range language need not have primitive-recursive membership. \\
\phantomsection\label{entry:symbol-registry-07913}\(\mathcal C\) & \texttt{\detokenize{\mathcal C}} & target class and selected infinite language symbol family member \(\mathcal C\) & P15 tester obstruction & object & \hyperref[def:p15-pr-learning-native-diagonal]{\ref*{def:p15-pr-learning-native-diagonal}} & \thepage & Class includes every finite language and an infinite member; naming prerequisite is separate from presentation quantification. Explicit displayed family; member types and common actual first-use exposed, not one atomic glyph. Family listing, not one atomic glyph. Author macro \texttt{\detokenize{\mathcal C;L}} not carried to the members. \\
\phantomsection\label{entry:symbol-registry-07914}\(L\) & \texttt{\detokenize{L}} & target class and selected infinite language symbol family member \(L\) & P15 tester obstruction & object & \hyperref[def:p15-pr-learning-native-diagonal]{\ref*{def:p15-pr-learning-native-diagonal}} & \thepage & Class includes every finite language and an infinite member; naming prerequisite is separate from presentation quantification. Explicit displayed family; member types and common actual first-use exposed, not one atomic glyph. Family listing, not one atomic glyph. Author macro \texttt{\detokenize{\mathcal C;L}} not carried to the members. \\
\phantomsection\label{entry:symbol-registry-07915}\(G\) & \texttt{\detokenize{G}} & normalized primitive-recursive prefix learner and current hypothesis symbol family member \(G\) & P15 native Gold countertext & meta & \hyperref[def:p15-pr-learning-native-diagonal]{\ref*{def:p15-pr-learning-native-diagonal}} & \thepage & Source selector Gold1967 p471 I.14; native G retained after explicitly applying local normalization to G\_0. Explicit displayed family; member types and common actual first-use exposed, not one atomic glyph. Family listing, not one atomic glyph. Author macro \texttt{\detokenize{G;g\_t}} not carried to the members. \\
\phantomsection\label{entry:symbol-registry-07916}\(g_t\) & \texttt{\detokenize{g_t}} & normalized primitive-recursive prefix learner and current hypothesis symbol family member \(g_t\) & P15 native Gold countertext & meta & \hyperref[def:p15-pr-learning-native-diagonal]{\ref*{def:p15-pr-learning-native-diagonal}} & \thepage & Source selector Gold1967 p471 I.14; native G retained after explicitly applying local normalization to G\_0. Explicit displayed family; member types and common actual first-use exposed, not one atomic glyph. Family listing, not one atomic glyph. Author macro \texttt{\detokenize{G;g\_t}} not carried to the members. \\
\phantomsection\label{entry:symbol-registry-07917}\(X_t\) & \texttt{\detokenize{X_t}} & current position in the range enumeration & P15 native Gold countertext & meta & \hyperref[def:p15-pr-learning-native-diagonal]{\ref*{def:p15-pr-learning-native-diagonal}} & \thepage & Source selectors Gold1967 p471 I.17 and p472 I.19–I.20; advances by one only when current candidate was observed. \\
\phantomsection\label{entry:symbol-registry-07918}\(P_t\) & \texttt{\detokenize{P_t}} & already-observed candidate predicate & P15 native Gold countertext & meta & \hyperref[def:p15-pr-learning-native-diagonal]{\ref*{def:p15-pr-learning-native-diagonal}} & \thepage & Source selector Gold1967 p471 I.15; bounded disjunction over the actual finite information prefix. \\
\phantomsection\label{entry:symbol-registry-07919}\(Q_t\) & \texttt{\detokenize{Q_t}} & bounded rejecting-computation witness predicate & P15 native Gold countertext & meta & \hyperref[def:p15-pr-learning-native-diagonal]{\ref*{def:p15-pr-learning-native-diagonal}} & \thepage & Source selector Gold1967 p471 I.16; finite code bound \(y\leq t\); no totality oracle or unbounded rejection search. Mechanical TeX fix: the Unicode less-or-equal sign in the notes cell, which fails the symbol compile gate, is rewritten as inline math with the TeX command. \\
\phantomsection\label{entry:symbol-registry-07920}\(\varepsilon\) & \texttt{\detokenize{\varepsilon}} & empty word unary alphabet letter and repeated-letter word symbol family member \(\varepsilon\) & P15 delayed-candidate worked case & object & \hyperref[def:p15-pr-learning-candidate-case]{\ref*{def:p15-pr-learning-candidate-case}} & \thepage & Word names denote their declared positive-integer codes here; a is an alphabet symbol not native hypothesis number a. Explicit displayed family; member types and common actual first-use exposed, not one atomic glyph. Family listing, not one atomic glyph. Author macro \texttt{\detokenize{\varepsilon;a;a\^{}k}} not carried to the members. \\
\phantomsection\label{entry:symbol-registry-07921}\(a\) & \texttt{\detokenize{a}} & empty word unary alphabet letter and repeated-letter word symbol family member \(a\) & P15 delayed-candidate worked case & object & \hyperref[def:p15-pr-learning-candidate-case]{\ref*{def:p15-pr-learning-candidate-case}} & \thepage & Word names denote their declared positive-integer codes here; a is an alphabet symbol not native hypothesis number a. Explicit displayed family; member types and common actual first-use exposed, not one atomic glyph. Family listing, not one atomic glyph. Author macro \texttt{\detokenize{\varepsilon;a;a\^{}k}} not carried to the members. \\
\phantomsection\label{entry:symbol-registry-07922}\(a^k\) & \texttt{\detokenize{a^k}} & empty word unary alphabet letter and repeated-letter word symbol family member \(a^k\) & P15 delayed-candidate worked case & object & \hyperref[def:p15-pr-learning-candidate-case]{\ref*{def:p15-pr-learning-candidate-case}} & \thepage & Word names denote their declared positive-integer codes here; a is an alphabet symbol not native hypothesis number a. Explicit displayed family; member types and common actual first-use exposed, not one atomic glyph. Family listing, not one atomic glyph. Author macro \texttt{\detokenize{\varepsilon;a;a\^{}k}} not carried to the members. \\
\phantomsection\label{entry:symbol-registry-07923}\(\pi\) & \texttt{\detokenize{\pi}} & primitive-recursive pairing for finite-sequence coding & P15 local finite-simulation proof & meta & \hyperref[annex:p15-pr-learning-coding]{\ref*{annex:p15-pr-learning-coding}} & \thepage & Concrete local coding is declared separately; inverse projections obtained by bounded search. \\
\phantomsection\label{entry:symbol-registry-07924}\(B\) & \texttt{\detokenize{B}} & tape-symbol coding base and bounded simulated configuration symbol family member \(B\) & P15 local finite-simulation proof & meta & \hyperref[annex:p15-pr-learning-coding]{\ref*{annex:p15-pr-learning-coding}} & \thepage & Finite digit operations and primitive recursion establish bounded simulation only; not universal terminating evaluation. Explicit displayed family; member types and common actual first-use exposed, not one atomic glyph. Family listing, not one atomic glyph. Author macro \texttt{\detokenize{B;c\_s}} not carried to the members. \\
\phantomsection\label{entry:symbol-registry-07925}\(c_s\) & \texttt{\detokenize{c_s}} & tape-symbol coding base and bounded simulated configuration symbol family member \(c_s\) & P15 local finite-simulation proof & meta & \hyperref[annex:p15-pr-learning-coding]{\ref*{annex:p15-pr-learning-coding}} & \thepage & Finite digit operations and primitive recursion establish bounded simulation only; not universal terminating evaluation. Explicit displayed family; member types and common actual first-use exposed, not one atomic glyph. Family listing, not one atomic glyph. Author macro \texttt{\detokenize{B;c\_s}} not carried to the members. \\
\phantomsection\label{entry:symbol-registry-07926}\(l\) & \texttt{\detokenize{l}} & finite tape stacks right-tail code and written symbol symbol family member \(l\) & P15 local finite-simulation proof & meta & \hyperref[annex:p15-pr-learning-coding]{\ref*{annex:p15-pr-learning-coding}} & \thepage & Explicit left and right update formulas use the current right-stack top as scanned symbol; numerical input/output digit protocol is fixed separately from source history roles. Explicit displayed family; member types and common actual first-use exposed, not one atomic glyph. Family listing, not one atomic glyph. Author macro \texttt{\detokenize{l;r;r\_{\mathrm{tail}};w}} not carried to the members. \\
\phantomsection\label{entry:symbol-registry-07927}\(r\) & \texttt{\detokenize{r}} & finite tape stacks right-tail code and written symbol symbol family member \(r\) & P15 local finite-simulation proof & meta & \hyperref[annex:p15-pr-learning-coding]{\ref*{annex:p15-pr-learning-coding}} & \thepage & Explicit left and right update formulas use the current right-stack top as scanned symbol; numerical input/output digit protocol is fixed separately from source history roles. Explicit displayed family; member types and common actual first-use exposed, not one atomic glyph. Family listing, not one atomic glyph. Author macro \texttt{\detokenize{l;r;r\_{\mathrm{tail}};w}} not carried to the members. \\
\phantomsection\label{entry:symbol-registry-07928}\(r_{\mathrm{tail}}\) & \texttt{\detokenize{r_{\mathrm{tail}}}} & finite tape stacks right-tail code and written symbol symbol family member \(r_{\mathrm{tail}}\) & P15 local finite-simulation proof & meta & \hyperref[annex:p15-pr-learning-coding]{\ref*{annex:p15-pr-learning-coding}} & \thepage & Explicit left and right update formulas use the current right-stack top as scanned symbol; numerical input/output digit protocol is fixed separately from source history roles. Explicit displayed family; member types and common actual first-use exposed, not one atomic glyph. Family listing, not one atomic glyph. Author macro \texttt{\detokenize{l;r;r\_{\mathrm{tail}};w}} not carried to the members. \\
\phantomsection\label{entry:symbol-registry-07929}\(w\) & \texttt{\detokenize{w}} & finite tape stacks right-tail code and written symbol symbol family member \(w\) & P15 local finite-simulation proof & meta & \hyperref[annex:p15-pr-learning-coding]{\ref*{annex:p15-pr-learning-coding}} & \thepage & Explicit left and right update formulas use the current right-stack top as scanned symbol; numerical input/output digit protocol is fixed separately from source history roles. Explicit displayed family; member types and common actual first-use exposed, not one atomic glyph. Family listing, not one atomic glyph. Author macro \texttt{\detokenize{l;r;r\_{\mathrm{tail}};w}} not carried to the members. \\
\phantomsection\label{entry:symbol-registry-07930}\(s_m\) & \texttt{\detokenize{s_m}} & serial completion tick and corresponding original prefix answer symbol family member \(s_m\) & P15 normalization proof & meta & \hyperref[annex:p15-pr-learning-normalization]{\ref*{annex:p15-pr-learning-normalization}} & \thepage & Strictly increasing completion ticks display every answer; both implications in PR.2 use this invariant. Explicit displayed family; member types and common actual first-use exposed, not one atomic glyph. Family listing, not one atomic glyph. Author macro \texttt{\detokenize{s\_m;a\_m}} not carried to the members. \\
\phantomsection\label{entry:symbol-registry-07931}\(a_m\) & \texttt{\detokenize{a_m}} & serial completion tick and corresponding original prefix answer symbol family member \(a_m\) & P15 normalization proof & meta & \hyperref[annex:p15-pr-learning-normalization]{\ref*{annex:p15-pr-learning-normalization}} & \thepage & Strictly increasing completion ticks display every answer; both implications in PR.2 use this invariant. Explicit displayed family; member types and common actual first-use exposed, not one atomic glyph. Family listing, not one atomic glyph. Author macro \texttt{\detokenize{s\_m;a\_m}} not carried to the members. \\
\phantomsection\label{entry:symbol-registry-07932}\(a_L\) & \texttt{\detokenize{a_L}} & fixed positive default and accepting-history predicate symbol family member \(a_L\) & P15 primitive-recursive range proof & meta & \hyperref[annex:p15-pr-learning-range]{\ref*{annex:p15-pr-learning-range}} & \thepage & Nonempty target is required to choose a\_L; not a uniform emptiness decider. Explicit displayed family; member types and common actual first-use exposed, not one atomic glyph. Family listing, not one atomic glyph. Author macro \texttt{\detokenize{a\_L;R}} not carried to the members. \\
\phantomsection\label{entry:symbol-registry-07933}\(R(x,y)\) & \texttt{\detokenize{R(x,}\allowbreak\detokenize{y)}} & fixed positive default and accepting-history predicate symbol family member \(R(x,y)\) & P15 primitive-recursive range proof & meta & \hyperref[annex:p15-pr-learning-range]{\ref*{annex:p15-pr-learning-range}} & \thepage & Nonempty target is required to choose a\_L; not a uniform emptiness decider. Explicit displayed family; member types and common actual first-use exposed, not one atomic glyph. Family listing, not one atomic glyph. Author macro \texttt{\detokenize{a\_L;R}} not carried to the members. \\
\phantomsection\label{entry:symbol-registry-07934}\(F\) & \texttt{\detokenize{F}} & finite countertext content and eventual-case stages symbol family member \(F\) & P15 native Gold I.9 proof & meta & \hyperref[annex:p15-pr-learning-obstruction]{\ref*{annex:p15-pr-learning-obstruction}} & \thepage & Source selectors Gold1967 pp472–473 CaseII I.23–I.26 and CaseIII; stages assert existence and are not algorithmically known bounds. Explicit displayed family; member types and common actual first-use exposed, not one atomic glyph. Family listing, not one atomic glyph. Author macro \texttt{\detokenize{F;\alpha;\beta}} not carried to the members. \\
\phantomsection\label{entry:symbol-registry-07935}\(\alpha\) & \texttt{\detokenize{\alpha}} & finite countertext content and eventual-case stages symbol family member \(\alpha\) & P15 native Gold I.9 proof & meta & \hyperref[annex:p15-pr-learning-obstruction]{\ref*{annex:p15-pr-learning-obstruction}} & \thepage & Source selectors Gold1967 pp472–473 CaseII I.23–I.26 and CaseIII; stages assert existence and are not algorithmically known bounds. Explicit displayed family; member types and common actual first-use exposed, not one atomic glyph. Family listing, not one atomic glyph. Author macro \texttt{\detokenize{F;\alpha;\beta}} not carried to the members. \\
\phantomsection\label{entry:symbol-registry-07936}\(\beta\) & \texttt{\detokenize{\beta}} & finite countertext content and eventual-case stages symbol family member \(\beta\) & P15 native Gold I.9 proof & meta & \hyperref[annex:p15-pr-learning-obstruction]{\ref*{annex:p15-pr-learning-obstruction}} & \thepage & Source selectors Gold1967 pp472–473 CaseII I.23–I.26 and CaseIII; stages assert existence and are not algorithmically known bounds. Explicit displayed family; member types and common actual first-use exposed, not one atomic glyph. Family listing, not one atomic glyph. Author macro \texttt{\detokenize{F;\alpha;\beta}} not carried to the members. \\
\end{longtable}
\endgroup
\Needspace{.4\textheight}
\subsection*{Stable part P16 scopes}
\begingroup\scriptsize\setlength{\tabcolsep}{3pt}
\begin{longtable}{@{}p{0.9in}p{1.3in}p{1.0in}p{0.5in}p{0.3in}p{0.4in}p{0.35in}p{1.65in}@{}}
\toprule
\textbf{Glyph} & \textbf{Macro} & \textbf{Role} & \textbf{Scope} & \textbf{Prec./form} & \textbf{First} & \textbf{Reg. page} & \textbf{Principal uses / notes} \\
\midrule
\endfirsthead
\toprule
\textbf{Glyph} & \textbf{Macro} & \textbf{Role} & \textbf{Scope} & \textbf{Prec./form} & \textbf{First} & \textbf{Reg. page} & \textbf{Principal uses / notes} \\
\midrule
\endhead
\midrule\multicolumn{8}{r}{\emph{continued}}\\\endfoot
\bottomrule\endlastfoot
\phantomsection\label{entry:symbol-registry-00972}\(\mathcal P=(V,H,M,T)\) & \texttt{\detokenize{\mathcal P,}\allowbreak\detokenize{V,}\allowbreak\detokenize{H,}\allowbreak\detokenize{M,}\allowbreak\detokenize{T}} & finite logical-abduction input tuple & P16 finite-abduction & meta & \hyperref[sec:p16-finite-abduction]{\ref*{sec:p16-finite-abduction}} & \thepage & V is a finite atom alphabet, H its permitted hypothesis atoms, M its observation atoms, and T a finite background formula set. H and M may overlap. The extended input convention allows inconsistent T but its explanation family is then empty. \\
\phantomsection\label{entry:symbol-registry-00973}\(E,F,D,E_0\) & \texttt{\detokenize{E,}\allowbreak\detokenize{F,}\allowbreak\detokenize{D,}\allowbreak\detokenize{E_0}} & abduction candidates and local explanatory subsets & P16 finite-abduction & meta & \hyperref[sec:p16-finite-abduction]{\ref*{sec:p16-finite-abduction}} & \thepage & Candidates are subsets of H. Explanation additionally requires consistency and entailment of all observations. D and E0 are locally defined sets used in the minimality proofs. \\
\phantomsection\label{entry:symbol-registry-00974}\(\mathsf{Val}(V)\) & \texttt{\detokenize{\mathsf{Val}}} & valuations on a finite abduction alphabet & P16 finite-abduction & meta & \hyperref[sec:p16-finite-abduction]{\ref*{sec:p16-finite-abduction}} & \thepage & Functions from V to the existing Boolean truth-value set. The unique empty function is included when V is empty; restriction and extension connect them to global CPL valuations. \\
\phantomsection\label{entry:symbol-registry-00975}\(\operatorname{Mod}_V(\Gamma)\) & \texttt{\detokenize{\operatorname{Mod}_V}} & finite valuation models of an abduction theory & P16 finite-abduction & meta & \hyperref[sec:p16-finite-abduction]{\ref*{sec:p16-finite-abduction}} & \thepage & The formulas of Gamma use only atoms in V. Nonemptiness is the chapter's semantic consistency test; no model is required when Gamma is used merely as the antecedent of semantic consequence. \\
\phantomsection\label{entry:symbol-registry-00976}\(\Gamma\entails M\) & \texttt{\detokenize{\Gamma\entails M}} & simultaneous semantic entailment of observed atoms & P16 finite-abduction & meta & \hyperref[sec:p16-finite-abduction]{\ref*{sec:p16-finite-abduction}} & \thepage & Abbreviates Gamma entails m for every atom m in M. It is vacuous for empty M and for inconsistent Gamma; an explanation still requires a separate nonempty model set. \\
\phantomsection\label{entry:symbol-registry-00977}\(\operatorname{Expl}(\mathcal P)\) & \texttt{\detokenize{\operatorname{Expl}}} & family of all permitted logical explanations & P16 finite-abduction & meta & \hyperref[sec:p16-finite-abduction]{\ref*{sec:p16-finite-abduction}} & \thepage & Each member is a subset of H consistent with T and sufficient to entail every observation. This notation corresponds to Sol in the identified Eiter-Gottlob manuscript under the stated source translation. \\
\phantomsection\label{entry:symbol-registry-00978}\(\operatorname{MinExpl}(\mathcal P),\operatorname{CardExpl}(\mathcal P)\) & \texttt{\detokenize{\operatorname{MinExpl},}\allowbreak\detokenize{\operatorname{CardExpl}}} & subset-minimal and minimum-cardinality explanation families & P16 finite-abduction & meta & \hyperref[sec:p16-finite-abduction]{\ref*{sec:p16-finite-abduction}} & \thepage & Minimum cardinality compares all explanations, including incomparable sets. Distinct subset-minimal explanations can have different sizes; either family is empty when no explanation exists. \\
\phantomsection\label{entry:symbol-registry-00979}\(n=|V|,r=|H|\) & \texttt{\detokenize{n=|V|,}\allowbreak\detokenize{r=|H|}} & finite enumeration alphabet and hypothesis sizes & P16 finite-abduction & meta & \hyperref[sec:p16-finite-abduction]{\ref*{sec:p16-finite-abduction}} & \thepage & The exhaustive method inspects 2 to the power r+n candidate-valuation pairs. This count does not assign unit cost to formula evaluation or minimality filtering. \\
\phantomsection\label{entry:symbol-registry-00980}\(\mathcal L,\mathcal L_{\subseteq},\mathcal L_{\#}\) & \texttt{\detokenize{\mathcal L,}\allowbreak\detokenize{\mathcal L_{\subseteq},}\allowbreak\detokenize{\mathcal L_{\#}}} & exhaustive explanation and filtered output lists & P16 finite-abduction & meta & \hyperref[sec:p16-finite-abduction]{\ref*{sec:p16-finite-abduction}} & \thepage & The lists contain no duplicates and retain enumeration order. The output theorem views them as their member families; the empty list differs from a one-member list containing the empty explanation. \\
\phantomsection\label{entry:symbol-registry-00981}\(\mathcal P_{\rm inconsistent},\mathcal P_{\rm given},\mathcal P_{\rm empty},\mathcal P_{\rm sizes},\mathcal P_{\rm constraint},\mathcal P_{\rm permitted}\) & \texttt{\detokenize{\mathcal P_{\rm inconsistent},}\allowbreak\detokenize{\mathcal P_{\rm given},}\allowbreak\detokenize{\mathcal P_{\rm empty},}\allowbreak\detokenize{\mathcal P_{\rm sizes},}\allowbreak\detokenize{\mathcal P_{\rm constraint},}\allowbreak\detokenize{\mathcal P_{\rm permitted}}} & six complete finite abduction inputs & P16 finite-abduction & meta & \hyperref[sec:p16-finite-abduction]{\ref*{sec:p16-finite-abduction}} & \thepage & The tuples are defined in def:p16-abduction-examples. Their complete explanation families cover inconsistency, empty explanations and data, unequal minimal sizes, consistency constraints and hypotheses outside the permitted set. \\
\phantomsection\label{entry:symbol-registry-01166}\(N,n,G=(N,E),j\to i,\operatorname{pa}_G(i),\operatorname{de}_G(I)\) & \texttt{\detokenize{N,}\allowbreak\detokenize{n,}\allowbreak\detokenize{G,}\allowbreak\detokenize{E,}\allowbreak\detokenize{j\to i,}\allowbreak\detokenize{\operatorname{pa}_G,}\allowbreak\detokenize{\operatorname{de}_G}} & finite endogenous graph allowed edges parents and descendants & P16 peirce-abduction-and-inquiry & meta & \hyperref[sec:p16-peirce-abduction-and-inquiry]{\ref*{sec:p16-peirce-abduction-and-inquiry}} & \thepage & Parents specify allowed structural arguments; an edge may be inactive in a particular function. Descendants use strict directed paths and the graph is acyclic. \\
\phantomsection\label{entry:symbol-registry-01167}\(i_1,\ldots,i_n,I,J,O\) & \texttt{\detokenize{i_1,}\allowbreak\detokenize{\ldots,}\allowbreak\detokenize{i_n,}\allowbreak\detokenize{I,}\allowbreak\detokenize{J,}\allowbreak\detokenize{O}} & topological enumeration and intervention marginal and observation index sets & P16 peirce-abduction-and-inquiry & meta & \hyperref[sec:p16-peirce-abduction-and-inquiry]{\ref*{sec:p16-peirce-abduction-and-inquiry}} & \thepage & Topological order places parents before children. The coordinate sets may be empty; empty coordinate products have their unique empty tuple. \\
\phantomsection\label{entry:symbol-registry-01168}\(S_i,S_J,v,v_J,s,x,()\) & \texttt{\detokenize{S_i,}\allowbreak\detokenize{S_J,}\allowbreak\detokenize{v,}\allowbreak\detokenize{v_J,}\allowbreak\detokenize{s,}\allowbreak\detokenize{x,}\allowbreak\detokenize{()}} & finite value sets products and full projected intervention and empty assignments & P16 peirce-abduction-and-inquiry & meta & \hyperref[sec:p16-peirce-abduction-and-inquiry]{\ref*{sec:p16-peirce-abduction-and-inquiry}} & \thepage & Each coordinate value set is nonempty and finite. Full and projected assignments are indexed tuples; local binary example values are separately declared. \\
\phantomsection\label{entry:symbol-registry-01169}\(M,\Omega,q,\mathbb P,\mathcal U_i,U_i,f_i\) & \texttt{\detokenize{M,}\allowbreak\detokenize{\Omega,}\allowbreak\detokenize{q,}\allowbreak\detokenize{\mathbb P,}\allowbreak\detokenize{\mathcal U_i,}\allowbreak\detokenize{U_i,}\allowbreak\detokenize{f_i}} & finite structural model exogenous space law disturbances and equation functions & P16 peirce-abduction-and-inquiry & meta & \hyperref[sec:p16-peirce-abduction-and-inquiry]{\ref*{sec:p16-peirce-abduction-and-inquiry}} & \thepage & Structural functions are total. The exogenous law is a joint finite law, and its disturbances may be dependent unless mutual independence is expressly assumed. \\
\phantomsection\label{entry:symbol-registry-01170}\(V_i,V,V_J,p_M,p_{M,J}\) & \texttt{\detokenize{V_i,}\allowbreak\detokenize{V,}\allowbreak\detokenize{V_J,}\allowbreak\detokenize{p_M,}\allowbreak\detokenize{p_{M,J}}} & endogenous solutions and full and marginal observational mass functions & P16 peirce-abduction-and-inquiry & meta & \hyperref[sec:p16-peirce-abduction-and-inquiry]{\ref*{sec:p16-peirce-abduction-and-inquiry}} & \thepage & Acyclic topological evaluation supplies a unique solution at every exogenous outcome and its induced probability law. The law is independent of the chosen topological ordering. \\
\phantomsection\label{entry:symbol-registry-01171}\(M_{I=x},V^{I=x},p_M^{I=x},\operatorname{do}(V_I=x),\mathbb P_M(\cdot\mid\operatorname{do}(V_I=x))\) & \texttt{\detokenize{M_{I=x},}\allowbreak\detokenize{V^{I=x},}\allowbreak\detokenize{p_M^{I=x},}\allowbreak\detokenize{\operatorname{do},}\allowbreak\detokenize{\mathbb P_M(\cdot\mid\operatorname{do}(V_I=x))}} & equation-replaced model solution intervention law and do notation & P16 peirce-abduction-and-inquiry & meta & \hyperref[sec:p16-peirce-abduction-and-inquiry]{\ref*{sec:p16-peirce-abduction-and-inquiry}} & \thepage & Intervention replaces specified equations and retains the other functions and exogenous law. A bar accompanying do notation names this operation; it does not condition on a positive-probability event. \\
\phantomsection\label{entry:symbol-registry-01172}\(h,\mathbb E[h(V_J^{I=x})]\) & \texttt{\detokenize{h,}\allowbreak\detokenize{\mathbb E[h(V_J^{I=x})]}} & finite response function and its intervention mean & P16 peirce-abduction-and-inquiry & meta & \hyperref[sec:p16-peirce-abduction-and-inquiry]{\ref*{sec:p16-peirce-abduction-and-inquiry}} & \thepage & The response is a specified real function on the selected finite coordinate product. Its mean follows by summing the intervention law and makes no empirical estimation claim. \\
\phantomsection\label{entry:symbol-registry-01173}\(r_i,k_i(b\mid a)\) & \texttt{\detokenize{r_i,}\allowbreak\detokenize{k_i(b\mid a)}} & disturbance marginal and total structural kernel & P16 peirce-abduction-and-inquiry & meta & \hyperref[sec:p16-peirce-abduction-and-inquiry]{\ref*{sec:p16-peirce-abduction-and-inquiry}} & \thepage & Here a is a parent tuple and b an output value. Structural kernels exist at every parent input; their identification with observational conditional ratios requires a positive parent event. \\
\phantomsection\label{entry:symbol-registry-01174}\(\mathcal C,Q(M),M_1,M_2\) & \texttt{\detokenize{\mathcal C,}\allowbreak\detokenize{Q(M),}\allowbreak\detokenize{M_1,}\allowbreak\detokenize{M_2}} & model class target and two models in the identification criterion & P16 peirce-abduction-and-inquiry & meta & \hyperref[sec:p16-peirce-abduction-and-inquiry]{\ref*{sec:p16-peirce-abduction-and-inquiry}} & \thepage & Identification compares equal observational laws across the two models and requires equal targets throughout the stated class. It is different from estimating a target from a finite sample. \\
\phantomsection\label{entry:symbol-registry-01175}\(Z,X,Y,z,x,y,U_Z,U_X,U_Y\) & \texttt{\detokenize{Z,}\allowbreak\detokenize{X,}\allowbreak\detokenize{Y,}\allowbreak\detokenize{z,}\allowbreak\detokenize{x,}\allowbreak\detokenize{y,}\allowbreak\detokenize{U_Z,}\allowbreak\detokenize{U_X,}\allowbreak\detokenize{U_Y}} & named endogenous variables values and disturbances in causal examples & P16 peirce-abduction-and-inquiry & meta & \hyperref[sec:p16-peirce-abduction-and-inquiry]{\ref*{sec:p16-peirce-abduction-and-inquiry}} & \thepage & The source chain and the common-cause example each declare their own graph, equations and disturbance law before conditional or intervention calculations. \\
\phantomsection\label{entry:symbol-registry-01176}\(\oplus,\alpha,\beta,q_B,(z,a,b),(c,d)\) & \texttt{\detokenize{\oplus,}\allowbreak\detokenize{\alpha,}\allowbreak\detokenize{\beta,}\allowbreak\detokenize{q_B,}\allowbreak\detokenize{(z,}\allowbreak\detokenize{a,}\allowbreak\detokenize{b),}\allowbreak\detokenize{(c,}\allowbreak\detokenize{d)}} & exclusive-or binary masses and local exogenous coordinates & P16 peirce-abduction-and-inquiry & meta & \hyperref[sec:p16-peirce-abduction-and-inquiry]{\ref*{sec:p16-peirce-abduction-and-inquiry}} & \thepage & The mass alpha specifies the common-cause models, while beta specifies the separate direct-cause witness. Local a and b are bits, distinct from the general kernel's tuple and output placeholders. \\
\phantomsection\label{entry:symbol-registry-01177}\(M_A,M_A',M_B,M_C,M_D,\widetilde U_X,\widetilde U_Y\) & \texttt{\detokenize{M_A,}\allowbreak\detokenize{M_A',}\allowbreak\detokenize{M_B,}\allowbreak\detokenize{M_C,}\allowbreak\detokenize{M_D,}\allowbreak\detokenize{\widetilde U_X,}\allowbreak\detokenize{\widetilde U_Y}} & explicit causal and nonidentification witness models & P16 peirce-abduction-and-inquiry & meta & \hyperref[sec:p16-peirce-abduction-and-inquiry]{\ref*{sec:p16-peirce-abduction-and-inquiry}} & \thepage & The reduced common-cause disturbances share a bit and are dependent. A separate independent-disturbance pair locates failure at an observationally unsupported intervention value. \\
\phantomsection\label{entry:symbol-registry-01178}\(A_i(v),L,H,u_L\) & \texttt{\detokenize{A_i(v),}\allowbreak\detokenize{L,}\allowbreak\detokenize{H,}\allowbreak\detokenize{u_L}} & fixed-assignment fibers preceding vertices nondescendants and disturbance tuple & P16 peirce-abduction-and-inquiry & meta & \hyperref[sec:p16-peirce-abduction-and-inquiry]{\ref*{sec:p16-peirce-abduction-and-inquiry}} & \thepage & These local proof objects establish the rectangular independence factorization and invariance of nonintervened nondescendants. H is a vertex set, distinct from the prior chapter's abducible hypotheses. \\
\phantomsection\label{entry:symbol-registry-01247}\(\operatorname{an}_G(J),A,B,C,H,i_0,\ldots,i_m\) & \texttt{\detokenize{\operatorname{an}_G,}\allowbreak\detokenize{A,}\allowbreak\detokenize{B,}\allowbreak\detokenize{C,}\allowbreak\detokenize{H,}\allowbreak\detokenize{i_0,}\allowbreak\detokenize{\ldots,}\allowbreak\detokenize{i_m}} & reflexive ancestry endpoint conditioning and ancestral vertex sets and path positions & P16 peirce-abduction-and-inquiry & meta & \hyperref[sec:p16-peirce-abduction-and-inquiry]{\ref*{sec:p16-peirce-abduction-and-inquiry}} & \thepage & Ancestry includes the vertex itself. The previously defined strict descendants retain their separate meaning. \\
\phantomsection\label{entry:symbol-registry-01248}\(A\mathrel{\perp_G}B\mid C,G_H^{\mathrm m},G_H^{\mathrm m}\setminus C\) & \texttt{\detokenize{\perp_G,}\allowbreak\detokenize{G_H^{\mathrm m},}\allowbreak\detokenize{G_H^{\mathrm m}\setminus C}} & d-separation ancestral moral graph and deletion after moralization & P16 peirce-abduction-and-inquiry & meta & \hyperref[sec:p16-peirce-abduction-and-inquiry]{\ref*{sec:p16-peirce-abduction-and-inquiry}} & \thepage & The bridge proves both directions between active paths and separation in the moralized induced ancestral graph. \\
\phantomsection\label{entry:symbol-registry-01249}\(V_A\mathrel{\perp\!\!\!\perp}V_B\mid V_C,p_H,p_C\) & \texttt{\detokenize{\perp\!\!\!\perp,}\allowbreak\detokenize{p_H,}\allowbreak\detokenize{p_C}} & conditional independence and marginal mass functions & P16 peirce-abduction-and-inquiry & meta & \hyperref[sec:p16-peirce-abduction-and-inquiry]{\ref*{sec:p16-peirce-abduction-and-inquiry}} & \thepage & Conditional independence is required on every positive conditioning stratum; joint positivity is not assumed by the finite global Markov theorem. \\
\phantomsection\label{entry:symbol-registry-01250}\(D,F,\phi(d,c),\psi(f,c),g(a,c),h(b,c),g_C(c),h_C(c)\) & \texttt{\detokenize{D,}\allowbreak\detokenize{F,}\allowbreak\detokenize{\phi,}\allowbreak\detokenize{\psi,}\allowbreak\detokenize{g,}\allowbreak\detokenize{h,}\allowbreak\detokenize{g_C,}\allowbreak\detokenize{h_C}} & reachability partition factor products marginal sums and normalizers & P16 peirce-abduction-and-inquiry & meta & \hyperref[sec:p16-peirce-abduction-and-inquiry]{\ref*{sec:p16-peirce-abduction-and-inquiry}} & \thepage & These intermediate nonnegative factor products need not be normalized probability laws. Positive conditioning mass supplies the needed normalizers. \\
\phantomsection\label{entry:symbol-registry-01251}\(G_{\underline t},M_x^\dagger,W,R\) & \texttt{\detokenize{G_{\underline t},}\allowbreak\detokenize{M_x^\dagger,}\allowbreak\detokenize{W,}\allowbreak\detokenize{R}} & outgoing-edge deletion auxiliary model solution and treatment nondescendants & P16 peirce-abduction-and-inquiry & meta & \hyperref[sec:p16-peirce-abduction-and-inquiry]{\ref*{sec:p16-peirce-abduction-and-inquiry}} & \thepage & Outgoing-edge deletion is distinct from the incoming-edge deletion used for equation replacement. The back-door proof constructs the auxiliary model explicitly. \\
\phantomsection\label{entry:symbol-registry-01252}\(t,B,C,O,X,Y,Z,x,y,z,Y^{t=x},Z^{t=x}\) & \texttt{\detokenize{t,}\allowbreak\detokenize{B,}\allowbreak\detokenize{C,}\allowbreak\detokenize{O,}\allowbreak\detokenize{X,}\allowbreak\detokenize{Y,}\allowbreak\detokenize{Z,}\allowbreak\detokenize{x,}\allowbreak\detokenize{y,}\allowbreak\detokenize{z,}\allowbreak\detokenize{Y^{t=x},}\allowbreak\detokenize{Z^{t=x}}} & treatment vertex response adjustment observed sets variables values and intervention coordinates & P16 peirce-abduction-and-inquiry & meta & \hyperref[sec:p16-peirce-abduction-and-inquiry]{\ref*{sec:p16-peirce-abduction-and-inquiry}} & \thepage & The adjustment theorem has a single treatment vertex and arbitrary finite response and adjustment tuples. Chosen-treatment overlap is stated on the strata actually used. \\
\phantomsection\label{entry:symbol-registry-01253}\(\mathcal M,t,h,r,X,M,Y,x,m,y,x'\) & \texttt{\detokenize{\mathcal M,}\allowbreak\detokenize{t,}\allowbreak\detokenize{h,}\allowbreak\detokenize{r,}\allowbreak\detokenize{X,}\allowbreak\detokenize{M,}\allowbreak\detokenize{Y,}\allowbreak\detokenize{x,}\allowbreak\detokenize{m,}\allowbreak\detokenize{y,}\allowbreak\detokenize{x'}} & whole model treatment mediator response vertices their variables and values & P16 peirce-abduction-and-inquiry & meta & \hyperref[sec:p16-peirce-abduction-and-inquiry]{\ref*{sec:p16-peirce-abduction-and-inquiry}} & \thepage & The mediator M is a coordinate variable, distinct from the whole structural model. The front-door component has one treatment and one mediator. \\
\phantomsection\label{entry:symbol-registry-01254}\(G_{\overline J},M^x,Y^x,Y^m,Y^{x,m}\) & \texttt{\detokenize{G_{\overline J},}\allowbreak\detokenize{M^x,}\allowbreak\detokenize{Y^x,}\allowbreak\detokenize{Y^m,}\allowbreak\detokenize{Y^{x,m}}} & incoming-edge deletion and typed single and joint intervention outputs & P16 peirce-abduction-and-inquiry & meta & \hyperref[sec:p16-peirce-abduction-and-inquiry]{\ref*{sec:p16-peirce-abduction-and-inquiry}} & \thepage & Superscripts abbreviate declared vertex assignments. The joint intervention response is proved pointwise equal to the mediator-only response under directed-path interception. \\
\phantomsection\label{entry:symbol-registry-01255}\(\mathcal X_+,\mathcal H_x,g:S_r\to\mathbb R\) & \texttt{\detokenize{\mathcal X_+,}\allowbreak\detokenize{\mathcal H_x,}\allowbreak\detokenize{g}} & positive treatment strata mediator support at the chosen treatment and a response function & P16 peirce-abduction-and-inquiry & meta & \hyperref[sec:p16-peirce-abduction-and-inquiry]{\ref*{sec:p16-peirce-abduction-and-inquiry}} & \thepage & The front-door formula requires overlap only for the positive mediator and treatment strata named by these sets. \\
\phantomsection\label{entry:symbol-registry-01256}\(U,U_U,U_X,U_M,U_Y,q(u,a,b),\alpha(a),\beta(b),q_m\) & \texttt{\detokenize{U,}\allowbreak\detokenize{U_U,}\allowbreak\detokenize{U_X,}\allowbreak\detokenize{U_M,}\allowbreak\detokenize{U_Y,}\allowbreak\detokenize{q,}\allowbreak\detokenize{\alpha,}\allowbreak\detokenize{\beta,}\allowbreak\detokenize{q_m}} & latent root independent disturbances outcome masses and mediator-response probabilities & P16 peirce-abduction-and-inquiry & meta & \hyperref[sec:p16-peirce-abduction-and-inquiry]{\ref*{sec:p16-peirce-abduction-and-inquiry}} & \thepage & The numerical model keeps the latent root distinct from its exogenous disturbance. The response probability q\_m is distinct from the outcome mass q. \\
\phantomsection\label{entry:symbol-registry-01269}\(X,Y,Z,W,x,y,z,w,C=X\cup W\) & \texttt{\detokenize{X,}\allowbreak\detokenize{Y,}\allowbreak\detokenize{Z,}\allowbreak\detokenize{W,}\allowbreak\detokenize{x,}\allowbreak\detokenize{y,}\allowbreak\detokenize{z,}\allowbreak\detokenize{w,}\allowbreak\detokenize{C}} & four disjoint do-calculus vertex sets assignments and graphical conditioning set & P16 peirce-abduction-and-inquiry & meta & \hyperref[sec:p16-peirce-abduction-and-inquiry]{\ref*{sec:p16-peirce-abduction-and-inquiry}} & \thepage & These capital letters denote vertex sets in this section. Empty sets have the empty indexed assignment; positive conditional events are stated in each rule. \\
\phantomsection\label{entry:symbol-registry-01270}\(G_{\overline A,\underline B},K,L,H,Z(W)\) & \texttt{\detokenize{G_{\overline A,\underline B},}\allowbreak\detokenize{K,}\allowbreak\detokenize{L,}\allowbreak\detokenize{H,}\allowbreak\detokenize{Z(W)}} & combined graph deletions and the exact action-deletion ancestor exception & P16 peirce-abduction-and-inquiry & meta & \hyperref[sec:p16-peirce-abduction-and-inquiry]{\ref*{sec:p16-peirce-abduction-and-inquiry}} & \thepage & Z(W) is Z minus the ancestors of W in the graph after incoming X deletion. Replacing it by all Z is not valid. \\
\phantomsection\label{entry:symbol-registry-01271}\(p_x,p_{x,z},p_x(D),\operatorname{pr}(y\mid\check x,z,w)\) & \texttt{\detokenize{p_x,}\allowbreak\detokenize{p_{x,z},}\allowbreak\detokenize{p_x(D),}\allowbreak\detokenize{\operatorname{pr}}} & intervention laws event sums and the source intervention-marker notation & P16 peirce-abduction-and-inquiry & meta & \hyperref[sec:p16-peirce-abduction-and-inquiry]{\ref*{sec:p16-peirce-abduction-and-inquiry}} & \thepage & Conditionals are ratios only on positive events. The action-deletion rule explicitly requires positive W mass in both compared intervention regimes. \\
\phantomsection\label{entry:symbol-registry-01272}\(F_i,\mathcal F_Z,F,e,T_e,\widetilde K,\widetilde V,\widetilde\Omega,\widetilde{\mathbb P}\) & \texttt{\detokenize{F_i,}\allowbreak\detokenize{\mathcal F_Z,}\allowbreak\detokenize{F,}\allowbreak\detokenize{e,}\allowbreak\detokenize{T_e,}\allowbreak\detokenize{\widetilde K,}\allowbreak\detokenize{\widetilde V,}\allowbreak\detokenize{\widetilde\Omega,}\allowbreak\detokenize{\widetilde{\mathbb P}}} & regime roots their vertex set random and realized tuples switched vertices and augmented model & P16 peirce-abduction-and-inquiry & meta & \hyperref[sec:p16-peirce-abduction-and-inquiry]{\ref*{sec:p16-peirce-abduction-and-inquiry}} & \thepage & The added finite regime bits are mutually independent and independent of all old disturbances; their conditional laws implement exact simultaneous equation replacement. \\
\phantomsection\label{entry:symbol-registry-01308}\(V_J[I=x],V[I=x;J=y],\mathcal W,I_k,x^k,J_k,Q\) & \texttt{\detokenize{V_J[I=x],}\allowbreak\detokenize{V[I=x;J=y],}\allowbreak\detokenize{\mathcal W,}\allowbreak\detokenize{I_k,}\allowbreak\detokenize{x^k,}\allowbreak\detokenize{J_k,}\allowbreak\detokenize{Q}} & same-unit and sequential interventions world labels assignments and joint query tuple & P16 peirce-abduction-and-inquiry & meta & \hyperref[sec:p16-peirce-abduction-and-inquiry]{\ref*{sec:p16-peirce-abduction-and-inquiry}} & \thepage & All worlds use the same background outcome. Overlapping sequential assignments obey the stated override rule; composition substitutes values attained in the same intervention world. \\
\phantomsection\label{entry:symbol-registry-01309}\(E,D,q_E,U,\mathcal U,r_E,\widehat Q\) & \texttt{\detokenize{E,}\allowbreak\detokenize{D,}\allowbreak\detokenize{q_E,}\allowbreak\detokenize{U,}\allowbreak\detokenize{\mathcal U,}\allowbreak\detokenize{r_E,}\allowbreak\detokenize{\widehat Q}} & positive evidence query event posterior weights full disturbance tuple grouped posterior and background query & P16 peirce-abduction-and-inquiry & meta & \hyperref[sec:p16-peirce-abduction-and-inquiry]{\ref*{sec:p16-peirce-abduction-and-inquiry}} & \thepage & Conditioning the background can introduce dependence among its components. The posterior calculation does not redraw independent disturbances in each world. \\
\phantomsection\label{entry:symbol-registry-01310}\(R,R_i,T_i^k,r(u),v^k\) & \texttt{\detokenize{R,}\allowbreak\detokenize{R_i,}\allowbreak\detokenize{T_i^k,}\allowbreak\detokenize{r(u),}\allowbreak\detokenize{v^k}} & shared tuple-valued root components duplicated endogenous coordinates and copied-world assignments & P16 peirce-abduction-and-inquiry & meta & \hyperref[sec:p16-peirce-abduction-and-inquiry]{\ref*{sec:p16-peirce-abduction-and-inquiry}} & \thepage & The finite twin construction uses one root carrying the entire joint disturbance tuple, including any original dependence. \\
\phantomsection\label{entry:symbol-registry-01311}\(A,L,Y,Y(0),Y(1),Y_a,\Delta,\operatorname{ATE},\operatorname{ATT}\) & \texttt{\detokenize{A,}\allowbreak\detokenize{L,}\allowbreak\detokenize{Y,}\allowbreak\detokenize{Y(0),}\allowbreak\detokenize{Y(1),}\allowbreak\detokenize{Y_a,}\allowbreak\detokenize{\Delta,}\allowbreak\detokenize{\operatorname{ATE},}\allowbreak\detokenize{\operatorname{ATT}}} & binary treatment baseline covariates response potential outcomes contrast and two treatment-effect targets & P16 peirce-abduction-and-inquiry & meta & \hyperref[sec:p16-peirce-abduction-and-inquiry]{\ref*{sec:p16-peirce-abduction-and-inquiry}} & \thepage & ATE and ATT have different exchangeability and overlap requirements. The bridge proves consistency and a converse finite structural representation. \\
\phantomsection\label{entry:symbol-registry-01312}\(\ell,\mu_a(\ell),\Omega_0,q_0,a_0,\ell_0,y_0,y_1\) & \texttt{\detokenize{\ell,}\allowbreak\detokenize{\mu_a,}\allowbreak\detokenize{\Omega_0,}\allowbreak\detokenize{q_0,}\allowbreak\detokenize{a_0,}\allowbreak\detokenize{\ell_0,}\allowbreak\detokenize{y_0,}\allowbreak\detokenize{y_1}} & covariate values observed means and the abstract law used for the potential-outcome converse & P16 peirce-abduction-and-inquiry & meta & \hyperref[sec:p16-peirce-abduction-and-inquiry]{\ref*{sec:p16-peirce-abduction-and-inquiry}} & \thepage & Positive strata are specified before division. Separate potential-outcome exchangeability does not imply joint exchangeability. \\
\phantomsection\label{entry:symbol-registry-01313}\(\mathcal M_=,\mathcal M_{\ne},\kappa(i,j),p_0,p_1,b\) & \texttt{\detokenize{\mathcal M_=,}\allowbreak\detokenize{\mathcal M_{\ne},}\allowbreak\detokenize{\kappa,}\allowbreak\detokenize{p_0,}\allowbreak\detokenize{p_1,}\allowbreak\detokenize{b}} & coupling models latent pair law marginal success probabilities and probability of benefit & P16 peirce-abduction-and-inquiry & meta & \hyperref[sec:p16-peirce-abduction-and-inquiry]{\ref*{sec:p16-peirce-abduction-and-inquiry}} & \thepage & The two constructed models agree on every single-world intervention law on the treatment-response pair but differ on their joint potential outcomes. Sharp benefit bounds are proved. \\
\phantomsection\label{entry:symbol-registry-01314}\(N,\Lambda,U_\alpha,\mathcal U_\alpha,r_\alpha,C_\alpha,\Lambda_i\) & \texttt{\detokenize{N,}\allowbreak\detokenize{\Lambda,}\allowbreak\detokenize{U_\alpha,}\allowbreak\detokenize{\mathcal U_\alpha,}\allowbreak\detokenize{r_\alpha,}\allowbreak\detokenize{C_\alpha,}\allowbreak\detokenize{\Lambda_i}} & observed and latent indices independent latent inputs laws and actual child-incidence sets & P16 peirce-abduction-and-inquiry & meta & \hyperref[sec:p16-peirce-abduction-and-inquiry]{\ref*{sec:p16-peirce-abduction-and-inquiry}} & \thepage & A latent input may have more than two observed children. The model retains these incidence sets and does not assume an arbitrary projection or independent-pair replacement preserves its laws. \\
\phantomsection\label{entry:symbol-registry-01315}\(G,i\to j,i\leftrightarrow j,G[H],\mathcal D_H,\Lambda_D\) & \texttt{\detokenize{G,}\allowbreak\detokenize{\to,}\allowbreak\detokenize{\leftrightarrow,}\allowbreak\detokenize{G[H],}\allowbreak\detokenize{\mathcal D_H,}\allowbreak\detokenize{\Lambda_D}} & observed mixed graph edge types induced graph districts and district latent groups & P16 peirce-abduction-and-inquiry & meta & \hyperref[sec:p16-peirce-abduction-and-inquiry]{\ref*{sec:p16-peirce-abduction-and-inquiry}} & \thepage & Districts are maximal bidirected-connected components. Directed ancestry ignores the bidirected edges. \\
\phantomsection\label{entry:symbol-registry-01316}\(Q[H](v_H;v_{N\setminus H}),p,p_H,P_x(y)\) & \texttt{\detokenize{Q[H],}\allowbreak\detokenize{p,}\allowbreak\detokenize{p_H,}\allowbreak\detokenize{P_x(y)}} & intervention kernels with output and parameter roles observed laws and source effect notation & P16 peirce-abduction-and-inquiry & meta & \hyperref[sec:p16-peirce-abduction-and-inquiry]{\ref*{sec:p16-peirce-abduction-and-inquiry}} & \thepage & Equation-replacement kernels are defined at every fixed parameter. Identification by observed ratios has its separate positive-prefix domain. \\
\phantomsection\label{entry:symbol-registry-01317}\(h_j,H_j,\rho_j,\theta_j,m(D),\delta_i(u;v),g_{D,j}(v)\) & \texttt{\detokenize{h_j,}\allowbreak\detokenize{H_j,}\allowbreak\detokenize{\rho_j,}\allowbreak\detokenize{\theta_j,}\allowbreak\detokenize{m(D),}\allowbreak\detokenize{\delta_i,}\allowbreak\detokenize{g_{D,j}}} & topological prefixes prefix masses conditional ratios last district index equation indicators and partial factors & P16 peirce-abduction-and-inquiry & meta & \hyperref[sec:p16-peirce-abduction-and-inquiry]{\ref*{sec:p16-peirce-abduction-and-inquiry}} & \thepage & District recovery uses a positive prefix before the district's last position, allowing a zero final numerator while excluding zero denominators. \\
\phantomsection\label{entry:symbol-registry-01350}\(\mathsf{ID},N,R,F,X,Y,S,W,A,D,q_R,G[R],\Sigma_B,v[B\leftarrow b]\) & \texttt{\detokenize{\mathsf{ID},}\allowbreak\detokenize{N,}\allowbreak\detokenize{R,}\allowbreak\detokenize{F,}\allowbreak\detokenize{X,}\allowbreak\detokenize{Y,}\allowbreak\detokenize{S,}\allowbreak\detokenize{W,}\allowbreak\detokenize{A,}\allowbreak\detokenize{D,}\allowbreak\detokenize{q_R,}\allowbreak\detokenize{G[R],}\allowbreak\detokenize{\Sigma_B,}\allowbreak\detokenize{v[B\leftarrow b]}} & identification algorithm coordinate sets kernel graph coordinate summation and assignment replacement & P16 peirce-abduction-and-inquiry & meta & \hyperref[sec:p16-peirce-abduction-and-inquiry]{\ref*{sec:p16-peirce-abduction-and-inquiry}} & \thepage & Global random fixed action and response coordinates are explicit at each call. \\
\phantomsection\label{entry:symbol-registry-01351}\(\mathcal T(R;Y,X),\operatorname{Ext}^R_D,\Gamma,\widehat\Gamma,P^R_x,\mathcal Q_R[D]\) & \texttt{\detokenize{\mathcal T,}\allowbreak\detokenize{\operatorname{Ext}^R_D,}\allowbreak\detokenize{\Gamma,}\allowbreak\detokenize{\widehat\Gamma,}\allowbreak\detokenize{P^R_x,}\allowbreak\detokenize{\mathcal Q_R[D]}} & recursive target district extraction mixed graph explicit-root DAG and intervention kernels & P16 peirce-abduction-and-inquiry & meta & \hyperref[sec:p16-peirce-abduction-and-inquiry]{\ref*{sec:p16-peirce-abduction-and-inquiry}} & \thepage & Positive observed support supplies the proved denominator domain. Structural interventions remain separately defined. \\
\phantomsection\label{entry:symbol-registry-01352}\(L_i,A_D,B_i,D_i,E_i,a_D,b_i,\mathsf{FAIL},w^\circ,\nu(R,X)\) & \texttt{\detokenize{L_i,}\allowbreak\detokenize{A_D,}\allowbreak\detokenize{B_i,}\allowbreak\detokenize{D_i,}\allowbreak\detokenize{E_i,}\allowbreak\detokenize{a_D,}\allowbreak\detokenize{b_i,}\allowbreak\detokenize{\mathsf{FAIL},}\allowbreak\detokenize{w^\circ,}\allowbreak\detokenize{\nu(R,}\allowbreak\detokenize{X)}} & do-calculus partitions assignments raw failure reference action and recursion measure & P16 peirce-abduction-and-inquiry & meta & \hyperref[sec:p16-peirce-abduction-and-inquiry]{\ref*{sec:p16-peirce-abduction-and-inquiry}} & \thepage & Failures retain their parent traces; local query sets need not be subsets of the original query. \\
\phantomsection\label{entry:symbol-registry-01353}\(\mathbf S,\mathcal M^+(G,\mathbf S),C,N^*,X^*,Y^*,G^*,F,F',T,S,K,X_0,Y_0\) & \texttt{\detokenize{\mathbf S,}\allowbreak\detokenize{\mathcal M^+,}\allowbreak\detokenize{C,}\allowbreak\detokenize{N^*,}\allowbreak\detokenize{X^*,}\allowbreak\detokenize{Y^*,}\allowbreak\detokenize{G^*,}\allowbreak\detokenize{F,}\allowbreak\detokenize{F',}\allowbreak\detokenize{T,}\allowbreak\detokenize{S,}\allowbreak\detokenize{K,}\allowbreak\detokenize{X_0,}\allowbreak\detokenize{Y_0}} & fixed value sets positive model class constants reduced graph hedge forests vertices roots and original query & P16 peirce-abduction-and-inquiry & meta & \hyperref[sec:p16-peirce-abduction-and-inquiry]{\ref*{sec:p16-peirce-abduction-and-inquiry}} & \thepage & Singleton reduction preserves retained effects. Edge-subgraph forests have explicitly constructed common roots. \\
\phantomsection\label{entry:symbol-registry-01354}\(U_e,E,B_i,W_i,d_i,w_i,R_{i,b},H,J,C_i,L_i,I\) & \texttt{\detokenize{U_e,}\allowbreak\detokenize{E,}\allowbreak\detokenize{B_i,}\allowbreak\detokenize{W_i,}\allowbreak\detokenize{d_i,}\allowbreak\detokenize{w_i,}\allowbreak\detokenize{R_{i,b},}\allowbreak\detokenize{H,}\allowbreak\detokenize{J,}\allowbreak\detokenize{C_i,}\allowbreak\detokenize{L_i,}\allowbreak\detokenize{I}} & tree bits private parity prototype alphabets decoder fibers routing vertices sinks core and relay coordinates & P16 peirce-abduction-and-inquiry & meta & \hyperref[sec:p16-peirce-abduction-and-inquiry]{\ref*{sec:p16-peirce-abduction-and-inquiry}} & \thepage & Four values are required only at specified nonroot forest vertices traversed by response routes. \\
\phantomsection\label{entry:symbol-registry-01355}\(\Phi_{G,X,Y,\mathbf S},C,T,A,T',q_T,D,D_i,T_i,K\) & \texttt{\detokenize{\Phi_{G,X,Y,\mathbf S},}\allowbreak\detokenize{C,}\allowbreak\detokenize{T,}\allowbreak\detokenize{A,}\allowbreak\detokenize{T',}\allowbreak\detokenize{q_T,}\allowbreak\detokenize{D,}\allowbreak\detokenize{D_i,}\allowbreak\detokenize{T_i,}\allowbreak\detokenize{K}} & family-indexed identification function Tian component variables kernels outer districts and common roots & P16 peirce-abduction-and-inquiry & meta & \hyperref[sec:p16-peirce-abduction-and-inquiry]{\ref*{sec:p16-peirce-abduction-and-inquiry}} & \thepage & The graph-uniform quantifier ranges over all finite nonempty alphabet families. One fixed family remains a distinct question. \\
\phantomsection\label{entry:symbol-registry-01356}\(\mathsf{c\mbox{-}identify},\mathsf{identify},\operatorname{Ext}\) & \texttt{\detokenize{\mathsf{c\mbox{-}identify},}\allowbreak\detokenize{\mathsf{identify},}\allowbreak\detokenize{\operatorname{Ext}}} & modified Tian component and outer procedures and district extraction & P16 peirce-abduction-and-inquiry & meta & \hyperref[sec:p16-peirce-abduction-and-inquiry]{\ref*{sec:p16-peirce-abduction-and-inquiry}} & \thepage & Recursive inclusions permit T'=C and T'=A; only A must strictly shrink T. All outside coordinates remain parameters. \\
\phantomsection\label{entry:symbol-registry-01377}\(\mathcal A_{\mathrm{obs}},\mathcal A_{\mathrm{all}},B_X,\rho_j,\theta_j,h_D^\theta\) & \texttt{\detokenize{\mathcal A_{\mathrm{obs}},}\allowbreak\detokenize{\mathcal A_{\mathrm{all}},}\allowbreak\detokenize{B_X,}\allowbreak\detokenize{\rho_j,}\allowbreak\detokenize{\theta_j,}\allowbreak\detokenize{h_D^\theta}} & two action-support domains parent sets prefix masses conditional versions and district version products & P16 peirce-abduction-and-inquiry & meta & \hyperref[sec:p16-peirce-abduction-and-inquiry]{\ref*{sec:p16-peirce-abduction-and-inquiry}} & \thepage & Original-action support does not license every recursive extraction. The two explicit support quantifiers retain their distinct meanings at null parent tuples. \\
\phantomsection\label{entry:symbol-registry-01378}\(\mathcal O(q_R,D,j),J,B,\mathsf{Recover},\mathsf{Known},\mathsf{Unresolved}\) & \texttt{\detokenize{\mathcal O,}\allowbreak\detokenize{J,}\allowbreak\detokenize{B,}\allowbreak\detokenize{\mathsf{Recover},}\allowbreak\detokenize{\mathsf{Known},}\allowbreak\detokenize{\mathsf{Unresolved}}} & current complementary-action overlap complementary coordinates parent strata full-table recovery and scalar result types & P16 peirce-abduction-and-inquiry & meta & \hyperref[sec:p16-peirce-abduction-and-inquiry]{\ref*{sec:p16-peirce-abduction-and-inquiry}} & \thepage & An unresolved support record is not a probability or a pointwise nonidentification theorem. Only a proved causal zero permits omitting an unresolved factor. \\
\phantomsection\label{entry:symbol-registry-01399}\(P_x(y\mid z),X,Y,Z,x,y,z,\mathbf S,\widehat G,E(G;X,Y,Z),\mathsf{IDC},A,W\) & \texttt{\detokenize{P_x(y\mid z),}\allowbreak\detokenize{X,}\allowbreak\detokenize{Y,}\allowbreak\detokenize{Z,}\allowbreak\detokenize{x,}\allowbreak\detokenize{y,}\allowbreak\detokenize{z,}\allowbreak\detokenize{\mathbf S,}\allowbreak\detokenize{\widehat G,}\allowbreak\detokenize{E,}\allowbreak\detokenize{\mathsf{IDC},}\allowbreak\detokenize{A,}\allowbreak\detokenize{W}} & conditional effect disjoint query sets assignments value sets canonical root graph eligibility IDC and terminal exchange sets & P16 peirce-abduction-and-inquiry & meta & \hyperref[sec:p16-peirce-abduction-and-inquiry]{\ref*{sec:p16-peirce-abduction-and-inquiry}} & \thepage & Evidence belongs to the intervened world. Canonical roots are a separation device and do not replace the model's latent distributions. \\
\phantomsection\label{entry:symbol-registry-01400}\(U_i,L_i,\epsilon,t_1,t_2\) & \texttt{\detokenize{U_i,}\allowbreak\detokenize{L_i,}\allowbreak\detokenize{\epsilon,}\allowbreak\detokenize{t_1,}\allowbreak\detokenize{t_2}} & IDC active-path and passive relay coordinates noise and differing conditional evidence probabilities & P16 peirce-abduction-and-inquiry & meta & \hyperref[sec:p16-peirce-abduction-and-inquiry]{\ref*{sec:p16-peirce-abduction-and-inquiry}} & \thepage & The finite coordinate expansion is licensed by graph-uniform quantification. The masking proof preserves equal full observed laws and a nonzero two-by-two determinant. \\
\phantomsection\label{entry:symbol-registry-01419}\(T,S,C,t,s,c,\operatorname{Pa}^{+}_G(B),Q[B],\Sigma_B,\mathsf{Identify}\) & \texttt{\detokenize{T,}\allowbreak\detokenize{S,}\allowbreak\detokenize{C,}\allowbreak\detokenize{t,}\allowbreak\detokenize{s,}\allowbreak\detokenize{c,}\allowbreak\detokenize{\operatorname{Pa}^{+}_G,}\allowbreak\detokenize{Q,}\allowbreak\detokenize{\Sigma_B,}\allowbreak\detokenize{\mathsf{Identify}}} & Tian action response evidence sets values inclusive parents full intervention kernels coordinate summation and component subroutine & P16 peirce-abduction-and-inquiry & meta & \hyperref[sec:p16-peirce-abduction-and-inquiry]{\ref*{sec:p16-peirce-abduction-and-inquiry}} & \thepage & Inclusive parents contain the output vertices. Component calls retain their connected-set input domain and full outside intervention arguments. \\
\phantomsection\label{entry:symbol-registry-01420}\(D,F,D_i,H_i,\mathcal N,\mathcal I,\mathcal I_0,\mathcal I_1,H_0,F_0,F_1,H',U,K\) & \texttt{\detokenize{D,}\allowbreak\detokenize{F,}\allowbreak\detokenize{D_i,}\allowbreak\detokenize{H_i,}\allowbreak\detokenize{\mathcal N,}\allowbreak\detokenize{\mathcal I,}\allowbreak\detokenize{\mathcal I_0,}\allowbreak\detokenize{\mathcal I_1,}\allowbreak\detokenize{H_0,}\allowbreak\detokenize{F_0,}\allowbreak\detokenize{F_1,}\allowbreak\detokenize{H',}\allowbreak\detokenize{U,}\allowbreak\detokenize{K}} & ancestral target marginalized coordinates district families factor closure sets unknown factor and identified residual & P16 peirce-abduction-and-inquiry & meta & \hyperref[sec:p16-peirce-abduction-and-inquiry]{\ref*{sec:p16-peirce-abduction-and-inquiry}} & \thepage & The four-phase procedure closes interactions before cancelling an unknown positive response-independent factor. Graph and shared-expression complexity does not bound expanded numerical tables. \\
\phantomsection\label{entry:symbol-registry-01421}\(\mathcal R,D_{\mathcal R}\) & \texttt{\detokenize{\mathcal R,}\allowbreak\detokenize{D_{\mathcal R}}} & response-connected factor family and its retained vertex scope & P16 peirce-abduction-and-inquiry & meta & \hyperref[sec:p16-peirce-abduction-and-inquiry]{\ref*{sec:p16-peirce-abduction-and-inquiry}} & \thepage & The incidence and moral graph bridge retains the original conditional query. Completeness has the stated positive graph-uniform quantifiers. \\
\phantomsection\label{entry:symbol-registry-01422}\(K,\mathcal L_k,\mathcal A_k,\mathcal Y,O,\bar L_k,\bar A_k,H_k,h_k,\mathcal T,T_k(\bar a)\) & \texttt{\detokenize{K,}\allowbreak\detokenize{\mathcal L_k,}\allowbreak\detokenize{\mathcal A_k,}\allowbreak\detokenize{\mathcal Y,}\allowbreak\detokenize{O,}\allowbreak\detokenize{\bar L_k,}\allowbreak\detokenize{\bar A_k,}\allowbreak\detokenize{H_k,}\allowbreak\detokenize{h_k,}\allowbreak\detokenize{\mathcal T,}\allowbreak\detokenize{T_k}} & last decision time finite value sets observed trajectory prefixes history full potential table and one fixed world's future tuple & P16 peirce-abduction-and-inquiry & meta & \hyperref[sec:p16-peirce-abduction-and-inquiry]{\ref*{sec:p16-peirce-abduction-and-inquiry}} & \thepage & Times run from zero through K. Potential histories satisfy no anticipation and recursive consistency. A joint future tuple is distinct from the joint collection of all worlds. \\
\phantomsection\label{entry:symbol-registry-01423}\(g,g',\pi_k,R_{k,h},O^\pi,Y^\pi,\psi(\pi)\) & \texttt{\detokenize{g,}\allowbreak\detokenize{g',}\allowbreak\detokenize{\pi_k,}\allowbreak\detokenize{R_{k,h},}\allowbreak\detokenize{O^\pi,}\allowbreak\detokenize{Y^\pi,}\allowbreak\detokenize{\psi}} & deterministic policies reduced rules stochastic action rows independent selectors counterfactual trajectory response and policy value & P16 peirce-abduction-and-inquiry & meta & \hyperref[sec:p16-peirce-abduction-and-inquiry]{\ref*{sec:p16-peirce-abduction-and-inquiry}} & \thepage & Policy selectors are defined independently before random table evaluation. A policy uses only its available chronological history. \\
\phantomsection\label{entry:symbol-registry-01424}\(b,e_k,r_k,m_k,\mathrm{RP},q^\pi,\phi,B_k,V_k,W^\pi\) & \texttt{\detokenize{b,}\allowbreak\detokenize{e_k,}\allowbreak\detokenize{r_k,}\allowbreak\detokenize{m_k,}\allowbreak\detokenize{\mathrm{RP},}\allowbreak\detokenize{q^\pi,}\allowbreak\detokenize{\phi,}\allowbreak\detokenize{B_k,}\allowbreak\detokenize{V_k,}\allowbreak\detokenize{W^\pi}} & baseline law assignment rows transitions final response rows reachable positivity policy trajectory law test response backward expectations and policy weight & P16 peirce-abduction-and-inquiry & meta & \hyperref[sec:p16-peirce-abduction-and-inquiry]{\ref*{sec:p16-peirce-abduction-and-inquiry}} & \thepage & The g-formula identity has exact reachable-history support conditions. Zero-weight branches terminate before undefined downstream ratios. True assignment rows are distinct from fitted weights. \\
\phantomsection\label{entry:symbol-registry-01433}\(A,M,C,R,a_0,a_1,M_a,Y_{a,m},Y_{a,M_b},\theta_{ij}\) & \texttt{\detokenize{A,}\allowbreak\detokenize{M,}\allowbreak\detokenize{C,}\allowbreak\detokenize{R,}\allowbreak\detokenize{a_0,}\allowbreak\detokenize{a_1,}\allowbreak\detokenize{M_a,}\allowbreak\detokenize{Y_{a,m},}\allowbreak\detokenize{Y_{a,M_b},}\allowbreak\detokenize{\theta_{ij}}} & exposure mediator baseline response background exposure values potential mediator controlled and nested responses and nested means & P16 peirce-abduction-and-inquiry & meta & \hyperref[sec:p16-peirce-abduction-and-inquiry]{\ref*{sec:p16-peirce-abduction-and-inquiry}} & \thepage & The response family and potential mediators are jointly defined before random-index evaluation. Separate fixed-mediator null sets do not license an uncountable diagonal substitution. \\
\phantomsection\label{entry:symbol-registry-01434}\(\mathrm{CDE},\mathrm{TE},\mathrm{PNDE},\mathrm{TNDE},\mathrm{PNIE},\mathrm{TNIE},I\) & \texttt{\detokenize{\mathrm{CDE},}\allowbreak\detokenize{\mathrm{TE},}\allowbreak\detokenize{\mathrm{PNDE},}\allowbreak\detokenize{\mathrm{TNDE},}\allowbreak\detokenize{\mathrm{PNIE},}\allowbreak\detokenize{\mathrm{TNIE},}\allowbreak\detokenize{I}} & controlled direct total pure and total natural direct and indirect effects and interaction remainder & P16 peirce-abduction-and-inquiry & meta & \hyperref[sec:p16-peirce-abduction-and-inquiry]{\ref*{sec:p16-peirce-abduction-and-inquiry}} & \thepage & Each effect specifies its two exposure values and reference mediator distribution. The two additive decompositions retain their respective pure and total components. \\
\phantomsection\label{entry:symbol-registry-01435}\(\nu,\eta_b,\kappa,Q_a,q_a,L,L_a,e_a,\lambda_a,\rho_a,\mu_a,\mathrm{OV}\) & \texttt{\detokenize{\nu,}\allowbreak\detokenize{\eta_b,}\allowbreak\detokenize{\kappa,}\allowbreak\detokenize{Q_a,}\allowbreak\detokenize{q_a,}\allowbreak\detokenize{L,}\allowbreak\detokenize{L_a,}\allowbreak\detokenize{e_a,}\allowbreak\detokenize{\lambda_a,}\allowbreak\detokenize{\rho_a,}\allowbreak\detokenize{\mu_a,}\allowbreak\detokenize{\mathrm{OV}}} & baseline law mediator kernel response-background kernel controlled law and mean confounder potential confounder assignment and observation kernels and overlap condition & P16 peirce-abduction-and-inquiry & meta & \hyperref[sec:p16-peirce-abduction-and-inquiry]{\ref*{sec:p16-peirce-abduction-and-inquiry}} & \thepage & The general formulas use a common joint conditional measure and explicit absolute-continuity transfer to the target support. Exposure-atom conditions apply to the stated observational corollary. \\
\phantomsection\label{entry:symbol-registry-01436}\(k,G^k,J(a,k),J_{ij},W_0,W_1,W_2,W_3,g,\mathcal M_g^0\) & \texttt{\detokenize{k,}\allowbreak\detokenize{G^k,}\allowbreak\detokenize{J,}\allowbreak\detokenize{J_{ij},}\allowbreak\detokenize{W_0,}\allowbreak\detokenize{W_1,}\allowbreak\detokenize{W_2,}\allowbreak\detokenize{W_3,}\allowbreak\detokenize{g,}\allowbreak\detokenize{\mathcal M_g^0}} & policy kernel independent mediator draw policy response mean four adjustment sets edge subgraph and reference-input path model & P16 peirce-abduction-and-inquiry & meta & \hyperref[sec:p16-peirce-abduction-and-inquiry]{\ref*{sec:p16-peirce-abduction-and-inquiry}} & \thepage & Independent mediator draws define an interventional effect that need not equal a natural effect. The four-set criterion and full-parent specialization retain their separate graph scopes. \\
\phantomsection\label{entry:symbol-registry-01448}\(K_v,P_K,S_v,\operatorname{pr}(v),\operatorname{nd}(v),F,q,t\) & \texttt{\detokenize{K_v,}\allowbreak\detokenize{P_K,}\allowbreak\detokenize{S_v,}\allowbreak\detokenize{\operatorname{pr},}\allowbreak\detokenize{\operatorname{nd},}\allowbreak\detokenize{F,}\allowbreak\detokenize{q,}\allowbreak\detokenize{t}} & vertex kernels recursively constructed law coordinate spaces predecessors nondescendants measurable randomizer upper quantile and intervention value & P16 peirce-abduction-and-inquiry & meta & \hyperref[sec:p16-peirce-abduction-and-inquiry]{\ref*{sec:p16-peirce-abduction-and-inquiry}} & \thepage & The finite DAG construction uses probability kernels on standard Borel spaces without a product-density assumption. A vertex is excluded from its own predecessors and nondescendants. \\
\phantomsection\label{entry:symbol-registry-01449}\(D,R,R_0,R_1,R_2,F_0,F_1,F_2\) & \texttt{\detokenize{D,}\allowbreak\detokenize{R,}\allowbreak\detokenize{R_0,}\allowbreak\detokenize{R_1,}\allowbreak\detokenize{R_2,}\allowbreak\detokenize{F_0,}\allowbreak\detokenize{F_1,}\allowbreak\detokenize{F_2}} & response-side moral component private-root tuple three response-background fields and their jointly measurable response maps & P16 peirce-abduction-and-inquiry & meta & \hyperref[sec:p16-peirce-abduction-and-inquiry]{\ref*{sec:p16-peirce-abduction-and-inquiry}} & \thepage & The three fields supply distinct joint conditional identities for natural composition, exposure adjustment and mediator adjustment. They cover all parameter values through the same measurable response family. \\
\phantomsection\label{entry:symbol-registry-01450}\(C,S,T,\gamma_1,\gamma_3,\gamma_2^a,\rho_a,Q_a^{\rm obs},\eta_b,\mathrm{MS},G^b,\alpha\) & \texttt{\detokenize{C,}\allowbreak\detokenize{S,}\allowbreak\detokenize{T,}\allowbreak\detokenize{\gamma_1,}\allowbreak\detokenize{\gamma_3,}\allowbreak\detokenize{\gamma_2^a,}\allowbreak\detokenize{\rho_a,}\allowbreak\detokenize{Q_a^{\rm obs},}\allowbreak\detokenize{\eta_b,}\allowbreak\detokenize{\mathrm{MS},}\allowbreak\detokenize{G^b,}\allowbreak\detokenize{\alpha}} & baseline and adjustment tuples conditional covariate kernels mediator law observed response kernel reference mediator kernel domination condition independent reference draw and target measure & P16 peirce-abduction-and-inquiry & meta & \hyperref[sec:p16-peirce-abduction-and-inquiry]{\ref*{sec:p16-peirce-abduction-and-inquiry}} & \thepage & Repeated covariate coordinates are removed from tuples. One observed-kernel exceptional set transfers to the target measure by absolute continuity; null exposure points need their separately stated identifying premise. \\
\phantomsection\label{entry:symbol-registry-01506}\(C,c,a_c,S,O,W,\mathsf{Normal}(N),W^-_\Delta,W^=_\Delta,\mathsf D(W),A(\mathcal M)\) & \texttt{\detokenize{C,}\allowbreak\detokenize{c,}\allowbreak\detokenize{a_c,}\allowbreak\detokenize{S,}\allowbreak\detokenize{O,}\allowbreak\detokenize{W,}\allowbreak\detokenize{\mathsf{Normal}(N),}\allowbreak\detokenize{W^-_\Delta,}\allowbreak\detokenize{W^=_\Delta,}\allowbreak\detokenize{\mathsf D(W),}\allowbreak\detokenize{A(\mathcal M)}} & components abnormality sentences system observations normality assumptions withdrawal and exact theories diagnosis family and a model's fault set & P16 consistency-based-diagnosis & meta & \hyperref[sec:p16-consistency-based-diagnosis]{\ref*{sec:p16-consistency-based-diagnosis}} & \thepage & Minimal withdrawal and minimal exact assignments agree. Exact assignments need not be upward closed. \\
\phantomsection\label{entry:symbol-registry-01507}\(\mathcal K(W),\mathcal K_{\min}(W),\mathcal F,\mathcal G,H,\mathsf{Tr}(\mathcal F),K_c\) & \texttt{\detokenize{\mathcal K(W),}\allowbreak\detokenize{\mathcal K_{\min}(W),}\allowbreak\detokenize{\mathcal F,}\allowbreak\detokenize{\mathcal G,}\allowbreak\detokenize{H,}\allowbreak\detokenize{\mathsf{Tr}(\mathcal F),}\allowbreak\detokenize{K_c}} & conflict families finite set families hitting set minimal transversals and private witnesses & P16 consistency-based-diagnosis & meta & \hyperref[sec:p16-consistency-based-diagnosis]{\ref*{sec:p16-consistency-based-diagnosis}} & \thepage & The empty family has one empty minimal transversal; a family containing the empty conflict has none. \\
\phantomsection\label{entry:symbol-registry-01508}\(H,J,K,K',K^*,n,q,k,v\) & \texttt{\detokenize{H,}\allowbreak\detokenize{J,}\allowbreak\detokenize{K,}\allowbreak\detokenize{K',}\allowbreak\detokenize{K^*,}\allowbreak\detokenize{n,}\allowbreak\detokenize{q,}\allowbreak\detokenize{k,}\allowbreak\detokenize{v}} & shared graph identity successful subset witness old and new conflicts retained cache subset component count initial conflict count depth cutoff and common atom count & P16 consistency-based-diagnosis & meta & \hyperref[sec:p16-consistency-based-diagnosis]{\ref*{sec:p16-consistency-based-diagnosis}} & \thepage & Every incoming path to H has length equal to its cardinality. Storage bounds describe active nodes and edges; oracle termination is a separate premise. \\
\phantomsection\label{entry:symbol-registry-01509}\(X_1,X_2,A_1,A_2,O_1,x,u,v,y,z,M_1,M_2,M_3,x_1,x_2,x_3\) & \texttt{\detokenize{X_1,}\allowbreak\detokenize{X_2,}\allowbreak\detokenize{A_1,}\allowbreak\detokenize{A_2,}\allowbreak\detokenize{O_1,}\allowbreak\detokenize{x,}\allowbreak\detokenize{u,}\allowbreak\detokenize{v,}\allowbreak\detokenize{y,}\allowbreak\detokenize{z,}\allowbreak\detokenize{M_1,}\allowbreak\detokenize{M_2,}\allowbreak\detokenize{M_3,}\allowbreak\detokenize{x_1,}\allowbreak\detokenize{x_2,}\allowbreak\detokenize{x_3}} & Boolean adder and multiplier network components and local wire values & P16 consistency-based-diagnosis & meta & \hyperref[sec:p16-consistency-based-diagnosis]{\ref*{sec:p16-consistency-based-diagnosis}} & \thepage & The first network has Boolean wire values and ten conflicts. The second has integer wire values; its observation equations are explicit. \\
\phantomsection\label{entry:symbol-registry-01510}\(\pi,W',D',\mathcal T_W,\operatorname{Th},\Gamma,E,\alpha_i,A_E\) & \texttt{\detokenize{\pi,}\allowbreak\detokenize{W',}\allowbreak\detokenize{D',}\allowbreak\detokenize{\mathcal T_W,}\allowbreak\detokenize{\operatorname{Th},}\allowbreak\detokenize{\Gamma,}\allowbreak\detokenize{E,}\allowbreak\detokenize{\alpha_i,}\allowbreak\detokenize{A_E}} & measured sentence updated theory new diagnosis normality defaults closure Reiter operator extension consequents and licensed set & P16 consistency-based-diagnosis & meta & \hyperref[sec:p16-consistency-based-diagnosis]{\ref*{sec:p16-consistency-based-diagnosis}} & \thepage & The diagnosis-extension bijection requires consistent W. Inconsistent W has no diagnosis and the single inconsistent default extension. \\
\phantomsection\label{entry:symbol-registry-01511}\(D,M,R,M^+,R_m,p_d,o_m,\mathrm{Dis},\mathrm{Present},\mathrm{Observed},\mathrm{AB}\) & \texttt{\detokenize{D,}\allowbreak\detokenize{M,}\allowbreak\detokenize{R,}\allowbreak\detokenize{M^+,}\allowbreak\detokenize{R_m,}\allowbreak\detokenize{p_d,}\allowbreak\detokenize{o_m,}\allowbreak\detokenize{\mathrm{Dis},}\allowbreak\detokenize{\mathrm{Present},}\allowbreak\detokenize{\mathrm{Observed},}\allowbreak\detokenize{\mathrm{AB}}} & covering disorders manifestations incidence relation observations cause-candidate sets ground atoms and first-order predicates & P16 consistency-based-diagnosis & meta & \hyperref[sec:p16-consistency-based-diagnosis]{\ref*{sec:p16-consistency-based-diagnosis}} & \thepage & The incidence relation does not by itself supply intervention semantics. Empty observed candidate sets give inconsistency. \\
\phantomsection\label{entry:symbol-registry-01512}\(p_c,P(\Delta),E\) & \texttt{\detokenize{p_c,}\allowbreak\detokenize{P(\Delta),}\allowbreak\detokenize{E}} & fault probabilities exact-assignment event probability and a separately declared observed probability event & P16 consistency-based-diagnosis & meta & \hyperref[sec:p16-consistency-based-diagnosis]{\ref*{sec:p16-consistency-based-diagnosis}} & \thepage & Independent unequal priors do not rank incomparable diagnoses by cardinality. Posterior comparison also requires likelihoods. \\
\phantomsection\label{entry:symbol-registry-01535}\(L,\mathfrak A,T,C,E,\Gamma\) & \texttt{\detokenize{L,}\allowbreak\detokenize{\mathfrak A,}\allowbreak\detokenize{T,}\allowbreak\detokenize{C,}\allowbreak\detokenize{E,}\allowbreak\detokenize{\Gamma}} & relational language intended structure theory sentence initial conditions explanandum and basic-literal verification witness & P16 deductive-nomological-explanation & meta & \hyperref[sec:p16-deductive-nomological-explanation]{\ref*{sec:p16-deductive-nomological-explanation}} & \thepage & The final criterion requires satisfiable {T} union Gamma and separate failure of Gamma to entail E. \\
\phantomsection\label{entry:symbol-registry-01536}\(E_1,E_2,C_1,A,U\) & \texttt{\detokenize{E_1,}\allowbreak\detokenize{E_2,}\allowbreak\detokenize{C_1,}\allowbreak\detokenize{A,}\allowbreak\detokenize{U}} & common-content decomposition sentences local singular premise and generalized sentence & P16 deductive-nomological-explanation & meta & \hyperref[sec:p16-deductive-nomological-explanation]{\ref*{sec:p16-deductive-nomological-explanation}} & \thepage & The source's 7.7 decomposition does not impose the inconsistent joint countermodel condition. \\
\phantomsection\label{entry:symbol-registry-01537}\(\mathcal T,W,w,\mathsf{Pt},G,D,p,q\) & \texttt{\detokenize{\mathcal T,}\allowbreak\detokenize{W,}\allowbreak\detokenize{w,}\allowbreak\detokenize{\mathsf{Pt},}\allowbreak\detokenize{G,}\allowbreak\detokenize{D,}\allowbreak\detokenize{p,}\allowbreak\detokenize{q}} & available theory whole attribute whole object part relation structure description and local part objects & P16 deductive-nomological-explanation & meta & \hyperref[sec:p16-deductive-nomological-explanation]{\ref*{sec:p16-deductive-nomological-explanation}} & \thepage & Relative emergence depends on the fixed theory and available description; inconsistent theory makes deduction vacuous. \\
\phantomsection\label{entry:symbol-registry-01538}\(u,v,a,b,A,B,H,t_*,\varphi\) & \texttt{\detokenize{u,}\allowbreak\detokenize{v,}\allowbreak\detokenize{a,}\allowbreak\detokenize{b,}\allowbreak\detokenize{A,}\allowbreak\detokenize{B,}\allowbreak\detokenize{H,}\allowbreak\detokenize{t_*,}\allowbreak\detokenize{\varphi}} & two thermal components positive rates expansion coefficients net change turning time and golden ratio & P16 deductive-nomological-explanation & meta & \hyperref[sec:p16-deductive-nomological-explanation]{\ref*{sec:p16-deductive-nomological-explanation}} & \thepage & These are explicitly stipulated ODE parameters rather than measured coefficients. \\
\phantomsection\label{entry:symbol-registry-01539}\(B,V,D,K,O,A,G,U,X,Y\) & \texttt{\detokenize{B,}\allowbreak\detokenize{V,}\allowbreak\detokenize{D,}\allowbreak\detokenize{K,}\allowbreak\detokenize{O,}\allowbreak\detokenize{A,}\allowbreak\detokenize{G,}\allowbreak\detokenize{U,}\allowbreak\detokenize{X,}\allowbreak\detokenize{Y}} & local function and motive propositions and causal comparison variables & P16 deductive-nomological-explanation & meta & \hyperref[sec:p16-deductive-nomological-explanation]{\ref*{sec:p16-deductive-nomological-explanation}} & \thepage & Roles are local to their worked examples; observational equivalence does not imply intervention equivalence. \\
\phantomsection\label{entry:symbol-registry-03820}\(A,B,C,H\) & \texttt{\detokenize{A;B;C;H}} & italic propositional random variables & P16 Rationality4.2 & n/a & \hyperref[def:p16-rat42-notation]{\ref*{def:p16-rat42-notation}} & \thepage & Source note2 distinguishes italic variables from roman value events. \\
\phantomsection\label{entry:symbol-registry-03821}\(\mathrm A,\mathrm B,\mathrm C,\mathrm H\) & \texttt{\detokenize{\mathrm A;\mathrm B;\mathrm C;\mathrm H}} & positive value events of the corresponding variables & P16 Rationality4.2 & n/a & \hyperref[def:p16-rat42-notation]{\ref*{def:p16-rat42-notation}} & \thepage & Roman script is the source's value notation; negation denotes the other binary value. \\
\phantomsection\label{entry:symbol-registry-03822}\(P\) & \texttt{\detokenize{P}} & specified joint probability law & P16 Rationality4.2 & n/a & \hyperref[def:p16-rat42-notation]{\ref*{def:p16-rat42-notation}} & \thepage & Every computation fixes one law; reported empirical judgments are not identified with it without a model. \\
\phantomsection\label{entry:symbol-registry-03823}\(P'\) & \texttt{\detokenize{P'}} & posterior law after specified learning & P16 Rationality4.2 & n/a & \hyperref[def:p16-rat42-notation]{\ref*{def:p16-rat42-notation}} & \thepage & Distinct from the prior and from an intervention law. \\
\phantomsection\label{entry:symbol-registry-03824}\(\bar x\) & \texttt{\detokenize{\bar x}} & complement of a probability parameter equal to one minus x & P16 Rationality4.2 & n/a & \hyperref[def:p16-rat42-notation]{\ref*{def:p16-rat42-notation}} & \thepage & Published source shorthand; not event negation or algebraic closure. \\
\phantomsection\label{entry:symbol-registry-03825}\(pa(A_i)\) & \texttt{\detokenize{pa(A_i)}} & set of parent variables of a node & P16 Rationality4.2 & n/a & \hyperref[def:p16-rat42-bayes-net]{\ref*{def:p16-rat42-bayes-net}} & \thepage & Source parent-set notation; empty at a root. \\
\phantomsection\label{entry:symbol-registry-03826}\(k_i\) & \texttt{\detokenize{k_i}} & normalized local mass table at each parent assignment & P16 Rationality4.2 & n/a & \hyperref[def:p16-rat42-bayes-net]{\ref*{def:p16-rat42-bayes-net}} & \thepage & Tables on null parent strata are stipulated entries rather than conditionals on null events. \\
\phantomsection\label{entry:symbol-registry-03827}\(\perp\) & \texttt{\detokenize{\perp}} & joint conditional independence of variable tuples & P16 Rationality4.2 & relation & \hyperref[def:p16-rat42-bayes-net]{\ref*{def:p16-rat42-bayes-net}} & \thepage & Conditionals refer to positive conditioning strata; semigraphoid operations are proved. \\
\phantomsection\label{entry:symbol-registry-03828}\(a,b,\alpha,\beta,\gamma,\delta\) & \texttt{\detokenize{a;b;\alpha;\beta;\gamma;\delta}} & binary collider prior and conditional probability parameters & P16 Rationality4.2 & n/a & \hyperref[thm:p16-rat42-three-networks]{\ref*{thm:p16-rat42-three-networks}} & \thepage & Published source collider table; a later report randomization parameter has its explicit local definition. \\
\phantomsection\label{entry:symbol-registry-03829}\(D(Q\|P)\) & \texttt{\detokenize{D(Q\|P)}} & directed posterior-to-prior Kullback Leibler divergence & P16 Rationality4.2 & n/a & \hyperref[def:p16-rat42-kl]{\ref*{def:p16-rat42-kl}} & \thepage & Direction and zero conventions are explicit; not a metric. \\
\phantomsection\label{entry:symbol-registry-03830}\(s,u,v,t,s',d\) & \texttt{\detokenize{s;u;v;t;s';d}} & parameters of the finite conditional-constraint minimization & P16 Rationality4.2 & n/a & \hyperref[thm:p16-rat42-kl-update]{\ref*{thm:p16-rat42-kl-update}} & \thepage & Locally defined in the stated positive four-cell law; not latent empirical process parameters. \\
\phantomsection\label{entry:symbol-registry-03831}\(h,p,q,x\) & \texttt{\detokenize{h;p;q;x}} & prior and evidence likelihoods with inverse likelihood ratio q over p & P16 Rationality4.2 & n/a & \hyperref[def:p16-rat42-confirmation]{\ref*{def:p16-rat42-confirmation}} & \thepage & The source uses the inverse ratio; x is only defined when p is positive. \\
\phantomsection\label{entry:symbol-registry-03832}\(fp,fn\) & \texttt{\detokenize{fp;fn}} & false-positive and false-negative rates & P16 Rationality4.2 & n/a & \hyperref[sec:p16-rat42-reliability-confirmation]{\ref*{sec:p16-rat42-reliability-confirmation}} & \thepage & Their event definitions are retained separately from the posterior hypothesis probability. \\
\phantomsection\label{entry:symbol-registry-03833}\(Rel,Rep\) & \texttt{\detokenize{Rel;Rep}} & italic reliability and report variables & P16 Rationality4.2 & n/a & \hyperref[def:p16-rat42-source-model]{\ref*{def:p16-rat42-source-model}} & \thepage & Published source symbols; roman script specifies their positive value events. \\
\phantomsection\label{entry:symbol-registry-03834}\(\mathrm{Rel},\mathrm{Rep}\) & \texttt{\detokenize{\mathrm{Rel};\mathrm{Rep}}} & positive reliability and positive report events & P16 Rationality4.2 & n/a & \hyperref[def:p16-rat42-source-model]{\ref*{def:p16-rat42-source-model}} & \thepage & Not the variables themselves; event negation denotes their negative values. \\
\phantomsection\label{entry:symbol-registry-03835}\(r,a\) & \texttt{\detokenize{r;a}} & source reliability and unreliable-source positive-report probability & P16 Rationality4.2 & n/a & \hyperref[def:p16-rat42-source-model]{\ref*{def:p16-rat42-source-model}} & \thepage & The source's reliable-or-randomizing model is a stated idealization. \\
\phantomsection\label{entry:symbol-registry-03836}\(p_I,q_I,p_D,q_D\) & \texttt{\detokenize{p_I;q_I;p_D;q_D}} & two-positive-report likelihoods in separate and shared reliability models & P16 Rationality4.2 & n/a & \hyperref[thm:p16-rat42-witness]{\ref*{thm:p16-rat42-witness}} & \thepage & Report independence is conditional on the specified parents rather than unconditional. \\
\phantomsection\label{entry:symbol-registry-03837}\(Y\) & \texttt{\detokenize{Y}} & count variable mediating hypothesis and search outcome & P16 Rationality4.2 & n/a & \hyperref[def:p16-rat42-common-parent]{\ref*{def:p16-rat42-common-parent}} & \thepage & Finite suitable-alternative count in NAA; finite or countable good-argument count in NRF. \\
\phantomsection\label{entry:symbol-registry-03838}\(\mathrm Y_i\) & \texttt{\detokenize{\mathrm Y_i}} & value event that the mediator count is i & P16 Rationality4.2 & n/a & \hyperref[def:p16-rat42-common-parent]{\ref*{def:p16-rat42-common-parent}} & \thepage & Mutually exclusive exhaustive events; zero-weight values are permitted. \\
\phantomsection\label{entry:symbol-registry-03839}\(y_i,h_i,f_i\) & \texttt{\detokenize{y_i;h_i;f_i}} & mediator mass and conditional hypothesis and search-failure table entries & P16 Rationality4.2 & n/a & \hyperref[def:p16-rat42-common-parent]{\ref*{def:p16-rat42-common-parent}} & \thepage & On null mediator strata h\_i and f\_i are specified table entries not ordinary conditional probabilities. \\
\phantomsection\label{entry:symbol-registry-03840}\(F,\mathrm F\) & \texttt{\detokenize{F;\mathrm F}} & search-outcome variable and its positive failure-to-find event & P16 Rationality4.2 & n/a & \hyperref[def:p16-rat42-common-parent]{\ref*{def:p16-rat42-common-parent}} & \thepage & Unsuccessful search differs from nonexistence; P of the positive event must exceed zero for conditioning. \\
\phantomsection\label{entry:symbol-registry-03841}\(Exp,T,e,t\) & \texttt{\detokenize{Exp;T;e;t}} & expertise and honesty variables with their root probabilities & P16 Rationality4.2 & n/a & \hyperref[thm:p16-rat42-expertise-trust]{\ref*{thm:p16-rat42-expertise-trust}} & \thepage & Local extension distinguishes informed deceit from ignorance; not Hartmann's binary reliability variable. \\
\phantomsection\label{entry:symbol-registry-05175}\(\sigma\) & \texttt{\detokenize{\sigma}} & signature of finite-domain constants & \hyperref[sec:p16-nonmonotonic-causal-theories]{\ref*{sec:p16-nonmonotonic-causal-theories}} & meta & \hyperref[def:p16-nct-signature]{\ref*{def:p16-nct-signature}} & \thepage & May be set-sized; finiteness of each domain does not make the signature finite. \\
\phantomsection\label{entry:symbol-registry-05176}\(\operatorname{Dom}(c)\) & \texttt{\detokenize{\operatorname{Dom}(c)}} & specified nonempty finite value domain & \hyperref[sec:p16-nonmonotonic-causal-theories]{\ref*{sec:p16-nonmonotonic-causal-theories}} & meta & \hyperref[def:p16-nct-signature]{\ref*{def:p16-nct-signature}} & \thepage & Value symbols are not first-order individual terms. \\
\phantomsection\label{entry:symbol-registry-05177}\(c=v\) & \texttt{\detokenize{c=v}} & finite-domain value atom & \hyperref[sec:p16-nonmonotonic-causal-theories]{\ref*{sec:p16-nonmonotonic-causal-theories}} & object & \hyperref[def:p16-nct-signature]{\ref*{def:p16-nct-signature}} & \thepage & Exactly one domain value per constant; Boolean c abbreviates c=t. \\
\phantomsection\label{entry:symbol-registry-05178}\(I,J\) & \texttt{\detokenize{I,}\allowbreak\detokenize{J}} & finite-domain interpretations & \hyperref[sec:p16-nonmonotonic-causal-theories]{\ref*{sec:p16-nonmonotonic-causal-theories}} & meta & \hyperref[def:p16-nct-signature]{\ref*{def:p16-nct-signature}} & \thepage & I activates heads; J may test those same heads. \\
\phantomsection\label{entry:symbol-registry-05179}\(\operatorname{Mod}(\Gamma)\) & \texttt{\detokenize{\operatorname{Mod}(\Gamma)}} & set of interpretations satisfying a formula theory & \hyperref[sec:p16-nonmonotonic-causal-theories]{\ref*{sec:p16-nonmonotonic-causal-theories}} & meta & \hyperref[def:p16-nct-signature]{\ref*{def:p16-nct-signature}} & \thepage & Explanatory status requires this set to be exactly a singleton. \\
\phantomsection\label{entry:symbol-registry-05180}\(\equiv\) & \texttt{\detokenize{\equiv}} & object-formula biconditional in the selected source & \hyperref[sec:p16-nonmonotonic-causal-theories]{\ref*{sec:p16-nonmonotonic-causal-theories}} & infix & \hyperref[def:p16-nct-signature]{\ref*{def:p16-nct-signature}} & \thepage & Source-native local role; not elementary equivalence of structures. \\
\phantomsection\label{entry:symbol-registry-05181}\(\not\equiv\) & \texttt{\detokenize{\not\equiv}} & negated object biconditional & \hyperref[sec:p16-nonmonotonic-causal-theories]{\ref*{sec:p16-nonmonotonic-causal-theories}} & infix & \hyperref[def:p16-nct-signature]{\ref*{def:p16-nct-signature}} & \thepage & Boolean exclusive difference in the bowl's static causal body. \\
\phantomsection\label{entry:symbol-registry-05182}\(\supset\) & \texttt{\detokenize{\supset}} & material object implication & \hyperref[sec:p16-nonmonotonic-causal-theories]{\ref*{sec:p16-nonmonotonic-causal-theories}} & infix & \hyperref[def:p16-nct-signature]{\ref*{def:p16-nct-signature}} & \thepage & Distinct from the causal rule constructor. \\
\phantomsection\label{entry:symbol-registry-05183}\(\Leftarrow\) & \texttt{\detokenize{\Leftarrow}} & causal rule constructor & \hyperref[sec:p16-nonmonotonic-causal-theories]{\ref*{sec:p16-nonmonotonic-causal-theories}} & infix & \hyperref[def:p16-nct-theory]{\ref*{def:p16-nct-theory}} & \thepage & Head is caused when body holds; not a material connective or inference turnstile. \\
\phantomsection\label{entry:symbol-registry-05184}\(T^I\) & \texttt{\detokenize{T^I}} & activated-head formula theory & \hyperref[sec:p16-nonmonotonic-causal-theories]{\ref*{sec:p16-nonmonotonic-causal-theories}} & meta & \hyperref[def:p16-nct-theory]{\ref*{def:p16-nct-theory}} & \thepage & Selection uses I satisfying the body, correcting the isolated source p763 display. \\
\phantomsection\label{entry:symbol-registry-05185}\(\phi,\psi,F\) & \texttt{\detokenize{\phi,}\allowbreak\detokenize{\psi,}\allowbreak\detokenize{F}} & formula metavariables & \hyperref[sec:p16-nonmonotonic-causal-theories]{\ref*{sec:p16-nonmonotonic-causal-theories}} & meta & \hyperref[def:p16-nct-theory]{\ref*{def:p16-nct-theory}} & \thepage & Head/body formulas retain their different roles. \\
\phantomsection\label{entry:symbol-registry-05186}\(S,T,R\) & \texttt{\detokenize{S,}\allowbreak\detokenize{T,}\allowbreak\detokenize{R}} & causal theories and replacement context & \hyperref[sec:p16-nonmonotonic-causal-theories]{\ref*{sec:p16-nonmonotonic-causal-theories}} & meta & \hyperref[def:p16-nct-se]{\ref*{def:p16-nct-se}} & \thepage & Same fixed signature; contexts may be set-sized. \\
\phantomsection\label{entry:symbol-registry-05187}\((I,J)\) & \texttt{\detokenize{(I,}\allowbreak\detokenize{J)}} & SE-model pair & \hyperref[sec:p16-nonmonotonic-causal-theories]{\ref*{sec:p16-nonmonotonic-causal-theories}} & meta & \hyperref[def:p16-nct-se]{\ref*{def:p16-nct-se}} & \thepage & Both satisfy T\^{}I; not a logic-program here-and-there structure by definition. \\
\phantomsection\label{entry:symbol-registry-05188}\(c_t,a_t\) & \texttt{\detokenize{c_t,}\allowbreak\detokenize{a_t}} & time-indexed fluent and action constant families & \hyperref[sec:p16-nonmonotonic-causal-theories]{\ref*{sec:p16-nonmonotonic-causal-theories}} & object & \hyperref[sec:p16-nct-action-cases]{\ref*{sec:p16-nct-action-cases}} & \thepage & Subscripts distinguish constants; value metavariables range over whole declared domains. \\
\phantomsection\label{entry:symbol-registry-05189}\(upL_t,upR_t,sp_t\) & \texttt{\detokenize{upL_t,}\allowbreak\detokenize{upR_t,}\allowbreak\detokenize{sp_t}} & Boolean bowl fluent constants & \hyperref[sec:p16-nonmonotonic-causal-theories]{\ref*{sec:p16-nonmonotonic-causal-theories}} & object & \hyperref[sec:p16-nct-action-cases]{\ref*{sec:p16-nct-action-cases}} & \thepage & Renamed abbreviations retain source up-left/up-right/spilled roles. \\
\phantomsection\label{entry:symbol-registry-05190}\(liftL_0,liftR_0\) & \texttt{\detokenize{liftL_0,}\allowbreak\detokenize{liftR_0}} & Boolean handle-lifting action constants & \hyperref[sec:p16-nonmonotonic-causal-theories]{\ref*{sec:p16-nonmonotonic-causal-theories}} & object & \hyperref[sec:p16-nct-action-cases]{\ref*{sec:p16-nct-action-cases}} & \thepage & Their occurrence and nonoccurrence are both exogenously explained. \\
\phantomsection\label{entry:symbol-registry-05191}\(p_0,p_1\) & \texttt{\detokenize{p_0,}\allowbreak\detokenize{p_1}} & three-valued cyclic fluent constants & \hyperref[sec:p16-nonmonotonic-causal-theories]{\ref*{sec:p16-nonmonotonic-causal-theories}} & object & \hyperref[sec:p16-nct-action-cases]{\ref*{sec:p16-nct-action-cases}} & \thepage & Domain is {0,1,2}, correcting source p767's inconsistent declaration. \\
\phantomsection\label{entry:symbol-registry-05192}\(a_0\) & \texttt{\detokenize{a_0}} & Boolean cycling action constant & \hyperref[sec:p16-nonmonotonic-causal-theories]{\ref*{sec:p16-nonmonotonic-causal-theories}} & object & \hyperref[sec:p16-nct-action-cases]{\ref*{sec:p16-nct-action-cases}} & \thepage & Not an atom from a logic program. \\
\phantomsection\label{entry:symbol-registry-05193}\(v^{-}\) & \texttt{\detokenize{v^{-}}} & preceding value in the cycle & \hyperref[sec:p16-nonmonotonic-causal-theories]{\ref*{sec:p16-nonmonotonic-causal-theories}} & meta & \hyperref[sec:p16-nct-action-cases]{\ref*{sec:p16-nct-action-cases}} & \thepage & Order 0,1,2,0; defined only for the displayed three values. \\
\phantomsection\label{entry:symbol-registry-05194}\(coin_t\) & \texttt{\detokenize{coin_t}} & coin fluent with heads/tails values & \hyperref[sec:p16-nonmonotonic-causal-theories]{\ref*{sec:p16-nonmonotonic-causal-theories}} & object & \hyperref[sec:p16-nct-action-cases]{\ref*{sec:p16-nct-action-cases}} & \thepage & Nondeterministic successor support supplies no probabilities. \\
\phantomsection\label{entry:symbol-registry-05195}\(toss_0\) & \texttt{\detokenize{toss_0}} & Boolean tossing action constant & \hyperref[sec:p16-nonmonotonic-causal-theories]{\ref*{sec:p16-nonmonotonic-causal-theories}} & object & \hyperref[sec:p16-nct-action-cases]{\ref*{sec:p16-nct-action-cases}} & \thepage & Fixed source one-step example. \\
\phantomsection\label{entry:symbol-registry-05196}\(loc_t\) & \texttt{\detokenize{loc_t}} & location value constant & \hyperref[sec:p16-nonmonotonic-causal-theories]{\ref*{sec:p16-nonmonotonic-causal-theories}} & object & \hyperref[sec:p16-nct-action-cases]{\ref*{sec:p16-nct-action-cases}} & \thepage & Values 1 through k, k>1. \\
\phantomsection\label{entry:symbol-registry-05197}\(put(v)_0\) & \texttt{\detokenize{put(v)_0}} & Boolean put-action family & \hyperref[sec:p16-nonmonotonic-causal-theories]{\ref*{sec:p16-nonmonotonic-causal-theories}} & object & \hyperref[sec:p16-nct-action-cases]{\ref*{sec:p16-nct-action-cases}} & \thepage & Conflicting effects imply at most one action occurrence. \\
\phantomsection\label{entry:symbol-registry-05198}\(up(d)_t\) & \texttt{\detokenize{up(d)_t}} & Boolean domino fluent family & \hyperref[sec:p16-nonmonotonic-causal-theories]{\ref*{sec:p16-nonmonotonic-causal-theories}} & object & \hyperref[sec:p16-nct-action-cases]{\ref*{sec:p16-nct-action-cases}} & \thepage & Indices 1<=d<=k and 0<=t<=k. \\
\phantomsection\label{entry:symbol-registry-05199}\(tip_0\) & \texttt{\detokenize{tip_0}} & Boolean initial tip action & \hyperref[sec:p16-nonmonotonic-causal-theories]{\ref*{sec:p16-nonmonotonic-causal-theories}} & object & \hyperref[sec:p16-nct-action-cases]{\ref*{sec:p16-nct-action-cases}} & \thepage & No later tip actions are declared. \\
\phantomsection\label{entry:symbol-registry-05200}\(left_t\) & \texttt{\detokenize{left_t}} & Boolean pendulum fluent & \hyperref[sec:p16-nonmonotonic-causal-theories]{\ref*{sec:p16-nonmonotonic-causal-theories}} & object & \hyperref[sec:p16-nct-action-cases]{\ref*{sec:p16-nct-action-cases}} & \thepage & State times 0 through n. \\
\phantomsection\label{entry:symbol-registry-05201}\(hold_t\) & \texttt{\detokenize{hold_t}} & Boolean hold-action family & \hyperref[sec:p16-nonmonotonic-causal-theories]{\ref*{sec:p16-nonmonotonic-causal-theories}} & object & \hyperref[sec:p16-nct-action-cases]{\ref*{sec:p16-nct-action-cases}} & \thepage & Action times 0 through n-1; absent at horizon zero. \\
\phantomsection\label{entry:symbol-registry-05202}\(T(D,n)\) & \texttt{\detokenize{T(D,}\allowbreak\detokenize{n)}} & causal translation of displayed C+ description and horizon & \hyperref[sec:p16-nonmonotonic-causal-theories]{\ref*{sec:p16-nonmonotonic-causal-theories}} & meta & \hyperref[def:p16-nct-cplus-fragment]{\ref*{def:p16-nct-cplus-fragment}} & \thepage & Only the two selected local definite fragments are reconstructed as a full translation. \\
\phantomsection\label{entry:symbol-registry-05203}\(N,m\) & \texttt{\detokenize{N,}\allowbreak\detokenize{m}} & actual explainable-symbol list and typed alternative variables & \hyperref[sec:p16-nonmonotonic-causal-theories]{\ref*{sec:p16-nonmonotonic-causal-theories}} & meta & \hyperref[def:p16-nct-higher-order]{\ref*{def:p16-nct-higher-order}} & \thepage & Matching sorts/arities; nonexplainable symbols and object domains stay fixed. \\
\phantomsection\label{entry:symbol-registry-05204}\(T^*(m)\) & \texttt{\detokenize{T^*(m)}} & higher-order activated-head condition & \hyperref[sec:p16-nonmonotonic-causal-theories]{\ref*{sec:p16-nonmonotonic-causal-theories}} & meta & \hyperref[def:p16-nct-higher-order]{\ref*{def:p16-nct-higher-order}} & \thepage & Original bodies versus substituted heads; universally covers free-variable instances. \\
\phantomsection\label{entry:symbol-registry-05205}\(m=N\) & \texttt{\detokenize{m=N}} & componentwise extensional interpretation equality & \hyperref[sec:p16-nonmonotonic-causal-theories]{\ref*{sec:p16-nonmonotonic-causal-theories}} & relation & \hyperref[def:p16-nct-higher-order]{\ref*{def:p16-nct-higher-order}} & \thepage & Includes predicates/functions of matching sorts, not equality of their names. \\
\phantomsection\label{entry:symbol-registry-05206}\(\mathsf C\) & \texttt{\detokenize{\mathsf C}} & UCL caused modality & \hyperref[sec:p16-nonmonotonic-causal-theories]{\ref*{sec:p16-nonmonotonic-causal-theories}} & n/a & \hyperref[def:p16-nct-ucl]{\ref*{def:p16-nct-ucl}} & \thepage & S5 recursion plus unique-I-model selection; not an epistemic knowledge interpretation. \\
\phantomsection\label{entry:symbol-registry-05207}\(\mathcal S\) & \texttt{\detokenize{\mathcal S}} & set of worlds of an S5 structure & \hyperref[sec:p16-nonmonotonic-causal-theories]{\ref*{sec:p16-nonmonotonic-causal-theories}} & meta & \hyperref[def:p16-nct-ucl]{\ref*{def:p16-nct-ucl}} & \thepage & Contains I; distinct from a causal theory named S. \\
\phantomsection\label{entry:symbol-registry-05216}\(v\) & \texttt{\detokenize{v}} & relative volume in Peirce's empirical water-volume example & P16 Peirce empirical formula & object & \hyperref[def:p16-peirce-empirical-formula]{\ref*{def:p16-peirce-empirical-formula}} & \thepage & Published 1878 example notation; no observed dataset or exact physical law is supplied. \\
\phantomsection\label{entry:symbol-registry-05217}\(t\) & \texttt{\detokenize{t}} & temperature variable in Peirce's empirical water-volume example & P16 Peirce empirical formula & object & \hyperref[def:p16-peirce-empirical-formula]{\ref*{def:p16-peirce-empirical-formula}} & \thepage & Published 1878 notation; numerical values and coefficient values depend on measurement scales. \\
\phantomsection\label{entry:symbol-registry-05218}\(a\) & \texttt{\detokenize{a}} & fitted linear coefficient in Peirce's empirical cubic expression & P16 Peirce empirical formula & object & \hyperref[def:p16-peirce-empirical-formula]{\ref*{def:p16-peirce-empirical-formula}} & \thepage & Published notation; not an abducible atom or a measured value supplied here. \\
\phantomsection\label{entry:symbol-registry-05219}\(b\) & \texttt{\detokenize{b}} & fitted quadratic coefficient in Peirce's empirical cubic expression & P16 Peirce empirical formula & object & \hyperref[def:p16-peirce-empirical-formula]{\ref*{def:p16-peirce-empirical-formula}} & \thepage & Published notation; not an abducible atom or a measured value supplied here. \\
\phantomsection\label{entry:symbol-registry-05220}\(c\) & \texttt{\detokenize{c}} & fitted cubic coefficient in Peirce's empirical cubic expression & P16 Peirce empirical formula & object & \hyperref[def:p16-peirce-empirical-formula]{\ref*{def:p16-peirce-empirical-formula}} & \thepage & Published notation; not an abducible atom or a measured value supplied here. \\
\phantomsection\label{entry:symbol-registry-05356}\(\sigma^{fl},\sigma^{act}\) & \texttt{\detokenize{\sigma^{fl},}\allowbreak\detokenize{\sigma^{act}}} & fluent and action signature partitions & \hyperref[sec:p16-cplus-action-language]{\ref*{sec:p16-cplus-action-language}} & meta & \hyperref[def:p16-cplus-laws]{\ref*{def:p16-cplus-laws}} & \thepage & Action constants need not be Boolean. \\
\phantomsection\label{entry:symbol-registry-05357}\(\sigma^{sim},\sigma^{sd}\) & \texttt{\detokenize{\sigma^{sim},}\allowbreak\detokenize{\sigma^{sd}}} & simple and statically determined fluent partitions & \hyperref[sec:p16-cplus-action-language]{\ref*{sec:p16-cplus-action-language}} & meta & \hyperref[def:p16-cplus-laws]{\ref*{def:p16-cplus-laws}} & \thepage & Dynamic fluent heads cannot contain statically determined constants. \\
\phantomsection\label{entry:symbol-registry-05358}\(c=d\) & \texttt{\detokenize{c=d}} & finite-domain equality abbreviation & \hyperref[sec:p16-cplus-action-language]{\ref*{sec:p16-cplus-action-language}} & object & \hyperref[sec:p16-cplus-grammar]{\ref*{sec:p16-cplus-grammar}} & \thepage & Disjunction over shared domain values; empty disjunction is bottom. \\
\phantomsection\label{entry:symbol-registry-05359}\(D\) & \texttt{\detokenize{D}} & C+ action description & \hyperref[sec:p16-cplus-action-language]{\ref*{sec:p16-cplus-action-language}} & meta & \hyperref[def:p16-cplus-laws]{\ref*{def:p16-cplus-laws}} & \thepage & Set of three typed law classes; not a causal structural equation model. \\
\phantomsection\label{entry:symbol-registry-05360}\(F,G,H\) & \texttt{\detokenize{F,}\allowbreak\detokenize{G,}\allowbreak\detokenize{H}} & law formula metavariables & \hyperref[sec:p16-cplus-action-language]{\ref*{sec:p16-cplus-action-language}} & meta & \hyperref[def:p16-cplus-laws]{\ref*{def:p16-cplus-laws}} & \thepage & F head; G same-time condition; H origin state and event condition. \\
\phantomsection\label{entry:symbol-registry-05361}\(i:c\) & \texttt{\detokenize{i:c}} & time-stamped constant & \hyperref[sec:p16-cplus-action-language]{\ref*{sec:p16-cplus-action-language}} & object & \hyperref[def:p16-cplus-translation]{\ref*{def:p16-cplus-translation}} & \thepage & State fluent copies 0 through m; action copies 0 through m-1. \\
\phantomsection\label{entry:symbol-registry-05362}\(i:F\) & \texttt{\detokenize{i:F}} & simultaneous time substitution in a formula & \hyperref[sec:p16-cplus-action-language]{\ref*{sec:p16-cplus-action-language}} & meta & \hyperref[def:p16-cplus-translation]{\ref*{def:p16-cplus-translation}} & \thepage & All constant occurrences are stamped; not a temporal modality. \\
\phantomsection\label{entry:symbol-registry-05363}\(D_m\) & \texttt{\detokenize{D_m}} & finite-horizon causal translation & \hyperref[sec:p16-cplus-action-language]{\ref*{sec:p16-cplus-action-language}} & meta & \hyperref[def:p16-cplus-translation]{\ref*{def:p16-cplus-translation}} & \thepage & Only initial simple-fluent exogeneity is automatic. \\
\phantomsection\label{entry:symbol-registry-05364}\(\Leftarrow\) & \texttt{\detokenize{\Leftarrow}} & causal rule constructor & \hyperref[sec:p16-cplus-action-language]{\ref*{sec:p16-cplus-action-language}} & infix & \hyperref[sec:p16-cplus-grammar]{\ref*{sec:p16-cplus-grammar}} & \thepage & Body activation and unique head satisfaction; not material implication. \\
\phantomsection\label{entry:symbol-registry-05365}\(m,i,j\) & \texttt{\detokenize{m,}\allowbreak\detokenize{i,}\allowbreak\detokenize{j}} & finite horizon and time indices & \hyperref[sec:p16-cplus-action-language]{\ref*{sec:p16-cplus-action-language}} & object & \hyperref[def:p16-cplus-translation]{\ref*{def:p16-cplus-translation}} & \thepage & Horizon nonnegative; actions omitted at horizon zero. \\
\phantomsection\label{entry:symbol-registry-05366}\(\equiv\) & \texttt{\detokenize{\equiv}} & object-formula biconditional & \hyperref[sec:p16-cplus-action-language]{\ref*{sec:p16-cplus-action-language}} & infix & \hyperref[sec:p16-cplus-cases]{\ref*{sec:p16-cplus-cases}} & \thepage & Local source role; not elementary equivalence. \\
\phantomsection\label{entry:symbol-registry-05367}\(s,s',s_i\) & \texttt{\detokenize{s,}\allowbreak\detokenize{s',}\allowbreak\detokenize{s_i}} & complete fluent value assignments & \hyperref[sec:p16-cplus-action-language]{\ref*{sec:p16-cplus-action-language}} & meta & \hyperref[def:p16-cplus-system]{\ref*{def:p16-cplus-system}} & \thepage & Statehood additionally requires a D0 explanatory model. \\
\phantomsection\label{entry:symbol-registry-05368}\(e,e_i\) & \texttt{\detokenize{e,}\allowbreak\detokenize{e_i}} & complete event value assignments & \hyperref[sec:p16-cplus-action-language]{\ref*{sec:p16-cplus-action-language}} & meta & \hyperref[def:p16-cplus-system]{\ref*{def:p16-cplus-system}} & \thepage & Includes false Boolean actions and all non-Boolean attributes. \\
\phantomsection\label{entry:symbol-registry-05369}\(\langle s,e,s'\rangle\) & \texttt{\detokenize{\langle s,}\allowbreak\detokenize{e,}\allowbreak\detokenize{s'\rangle}} & labelled transition & \hyperref[sec:p16-cplus-action-language]{\ref*{sec:p16-cplus-action-language}} & relation & \hyperref[def:p16-cplus-system]{\ref*{def:p16-cplus-system}} & \thepage & Defined by the one-step explanatory theory; target statehood is proved. \\
\phantomsection\label{entry:symbol-registry-05370}\(R\) & \texttt{\detokenize{R}} & set of rigid fluent constants & \hyperref[sec:p16-cplus-action-language]{\ref*{sec:p16-cplus-action-language}} & meta & \hyperref[sec:p16-cplus-rigid]{\ref*{sec:p16-cplus-rigid}} & \thepage & Compression theorem also requires derived type and pure-R explaining laws. \\
\phantomsection\label{entry:symbol-registry-05371}\(D_m^R,I^R\) & \texttt{\detokenize{D_m^R,}\allowbreak\detokenize{I^R}} & compressed theory and assignment & \hyperref[sec:p16-cplus-action-language]{\ref*{sec:p16-cplus-action-language}} & meta & \hyperref[sec:p16-cplus-rigid]{\ref*{sec:p16-cplus-rigid}} & \thepage & Erase timestamps only on R; not unrestricted time identification. \\
\phantomsection\label{entry:symbol-registry-05372}\(\Delta_j,\Pi_0,A_i,B_j\) & \texttt{\detokenize{\Delta_j,}\allowbreak\detokenize{\Pi_0,}\allowbreak\detokenize{A_i,}\allowbreak\detokenize{B_j}} & activated-head coordinate blocks & \hyperref[sec:annex-p16-cplus-action-language]{\ref*{sec:annex-p16-cplus-action-language}} & meta & \hyperref[annex:p16.cplus-paths]{\ref*{annex:p16.cplus-paths}} & \thepage & Static, initial simple values, event heads, successor simple heads. \\
\phantomsection\label{entry:symbol-registry-05373}\(\Delta_R,\Gamma\) & \texttt{\detokenize{\Delta_R,}\allowbreak\detokenize{\Gamma}} & pure rigid and nonrigid reduct blocks & \hyperref[sec:annex-p16-cplus-action-language]{\ref*{sec:annex-p16-cplus-action-language}} & meta & \hyperref[annex:p16.cplus-rigid]{\ref*{annex:p16.cplus-rigid}} & \thepage & Blocks are formed after activation; bodies may mix only where permitted. \\
\phantomsection\label{entry:symbol-registry-05374}\(Ab_1(\alpha,\beta),Ab_2(\alpha)\) & \texttt{\detokenize{Ab_1(\alpha,}\allowbreak\detokenize{\beta),}\allowbreak\detokenize{Ab_2(\alpha)}} & source abnormality constants & \hyperref[sec:p16-cplus-action-language]{\ref*{sec:p16-cplus-action-language}} & object & \hyperref[sec:p16-cplus-cases]{\ref*{sec:p16-cplus-cases}} & \thepage & Ab1 derived fluent; Ab2 action; unless supplies negative default. \\
\phantomsection\label{entry:symbol-registry-05375}\(NextToBox,d\) & \texttt{\detokenize{NextToBox,}\allowbreak\detokenize{d}} & fresh derived Boolean definition constants & \hyperref[sec:p16-cplus-action-language]{\ref*{sec:p16-cplus-action-language}} & object & \hyperref[sec:p16-cplus-cases]{\ref*{sec:p16-cplus-cases}} & \thepage & No initial simple self-rule; unique extension depends on freshness. \\
\phantomsection\label{entry:symbol-registry-05376}\(None\) & \texttt{\detokenize{None}} & declared sentinel value in action or target domains & \hyperref[sec:p16-cplus-action-language]{\ref*{sec:p16-cplus-action-language}} & object & \hyperref[sec:p16-cplus-cases]{\ref*{sec:p16-cplus-cases}} & \thepage & Not Boolean false or an implicit language-wide no-event symbol. \\
\phantomsection\label{entry:symbol-registry-05377}\(Kind,Length,Publish\) & \texttt{\detokenize{Kind,}\allowbreak\detokenize{Length,}\allowbreak\detokenize{Publish}} & publication event attributes and action & \hyperref[sec:p16-cplus-action-language]{\ref*{sec:p16-cplus-action-language}} & object & \hyperref[sec:p16-cplus-cases]{\ref*{sec:p16-cplus-cases}} & \thepage & Kind and Length are non-Boolean; always constraints enforce sentinel iff no publication. \\
\phantomsection\label{entry:symbol-registry-05378}\(M,\alpha,\beta,l\) & \texttt{\detokenize{M,}\allowbreak\detokenize{\alpha,}\allowbreak\detokenize{\beta,}\allowbreak\detokenize{l}} & fixed monkey-name set and schema indices & \hyperref[sec:p16-cplus-action-language]{\ref*{sec:p16-cplus-action-language}} & meta & \hyperref[sec:p16-cplus-cases]{\ref*{sec:p16-cplus-cases}} & \thepage & Alpha and beta range over M; l ranges over three declared location values; these index constants rather than introduce first-order variables. \\
\phantomsection\label{entry:symbol-registry-05551}\(P,R,T,J,A,O,u,v\) & \texttt{\detokenize{P,}\allowbreak\detokenize{R,}\allowbreak\detokenize{T,}\allowbreak\detokenize{J,}\allowbreak\detokenize{A,}\allowbreak\detokenize{O,}\allowbreak\detokenize{u,}\allowbreak\detokenize{v}} & binary inputs outputs and contexts of stipulated printer contrast & P16 IBE printer example & n/a & \hyperref[def:p16-ibe-printer]{\ref*{def:p16-ibe-printer}} & \thepage & Power paper toner jam awake state successful output; u v specified contexts; one denotes presence; not probability measures. \\
\phantomsection\label{entry:symbol-registry-05552}\(R\) & \texttt{\detokenize{R}} & unary raven predicate in the logical comparison & P16 IBE raven logical predicates & n/a & \hyperref[thm:p16-ibe-raven-logic]{\ref*{thm:p16-ibe-raven-logic}} & \thepage & R and B predicates in logical statement; B separately denotes selected-object blackness event in probability calculation; H\_0 H\_1 exclusive population alternatives; no bird biology asserted. Root-authorized typed split; original grouped intake row remains unchanged as evidence. \\
\phantomsection\label{entry:symbol-registry-05553}\(B\) & \texttt{\detokenize{B}} & unary black predicate in the logical comparison & P16 IBE raven logical predicates & n/a & \hyperref[thm:p16-ibe-raven-logic]{\ref*{thm:p16-ibe-raven-logic}} & \thepage & R and B predicates in logical statement; B separately denotes selected-object blackness event in probability calculation; H\_0 H\_1 exclusive population alternatives; no bird biology asserted. Root-authorized typed split; original grouped intake row remains unchanged as evidence. \\
\phantomsection\label{entry:symbol-registry-05554}\(B\) & \texttt{\detokenize{B}} & selected-object blackness event in the sampling calculation & P16 IBE raven sampling events & n/a & \hyperref[def:p16-ibe-raven-sampling]{\ref*{def:p16-ibe-raven-sampling}} & \thepage & R and B predicates in logical statement; B separately denotes selected-object blackness event in probability calculation; H\_0 H\_1 exclusive population alternatives; no bird biology asserted. Root-authorized typed split; original grouped intake row remains unchanged as evidence. \\
\phantomsection\label{entry:symbol-registry-05555}\(H_0\) & \texttt{\detokenize{H_0}} & population alternative with both ravens black & P16 IBE raven sampling events & n/a & \hyperref[def:p16-ibe-raven-sampling]{\ref*{def:p16-ibe-raven-sampling}} & \thepage & R and B predicates in logical statement; B separately denotes selected-object blackness event in probability calculation; H\_0 H\_1 exclusive population alternatives; no bird biology asserted. Root-authorized typed split; original grouped intake row remains unchanged as evidence. \\
\phantomsection\label{entry:symbol-registry-05556}\(H_1\) & \texttt{\detokenize{H_1}} & population alternative with one black and one nonblack raven & P16 IBE raven sampling events & n/a & \hyperref[def:p16-ibe-raven-sampling]{\ref*{def:p16-ibe-raven-sampling}} & \thepage & R and B predicates in logical statement; B separately denotes selected-object blackness event in probability calculation; H\_0 H\_1 exclusive population alternatives; no bird biology asserted. Root-authorized typed split; original grouped intake row remains unchanged as evidence. \\
\phantomsection\label{entry:symbol-registry-05557}\(P(\cdot),B_0,E,H_i,\pi_i,\ell_i,Z\) & \texttt{\detokenize{P,}\allowbreak\detokenize{B_0,}\allowbreak\detokenize{E,}\allowbreak\detokenize{H_i,}\allowbreak\detokenize{pi_i,}\allowbreak\detokenize{ell_i,}\allowbreak\detokenize{Z}} & finite probability model and conditionalization inputs & P16 IBE Bayesian bridge & n/a & \hyperref[def:p16-ibe-bayes-input]{\ref*{def:p16-ibe-bayes-input}} & \thepage & P is probability law here; B\_0 background event; H\_i exhaustive event partition; positive priors pi\_i likelihoods ell\_i evidence normalizer Z. \\
\phantomsection\label{entry:symbol-registry-05558}\(I,C_I,a_i,q_i,r\) & \texttt{\detokenize{I,}\allowbreak\detokenize{C_I,}\allowbreak\detokenize{a_i,}\allowbreak\detokenize{q_i,}\allowbreak\detokenize{r}} & considered subset and absolute versus considered-set posterior & P16 IBE Bayesian bridge & n/a & \hyperref[def:p16-ibe-bayes-input]{\ref*{def:p16-ibe-bayes-input}} & \thepage & I finite index subset; C\_I its union event; a\_i absolute posterior; q\_i further conditioned posterior; r posterior mass of considered set. \\
\phantomsection\label{entry:symbol-registry-05559}\(b_i,a_i^*,D,c\) & \texttt{\detokenize{b_i,}\allowbreak\detokenize{a_i^*,}\allowbreak\detokenize{D,}\allowbreak\detokenize{c}} & positive posterior bonus altered mass and proof normalizer & P16 IBE bonus example & n/a & \hyperref[thm:p16-ibe-bonus]{\ref*{thm:p16-ibe-bonus}} & \thepage & b\_i not an explanatory-quality measure derived from Lipton; D weighted normalizer and c common positive weight; equality characterized only on positive posterior support. \\
\phantomsection\label{entry:symbol-registry-05560}\(H,K\) & \texttt{\detokenize{H,}\allowbreak\detokenize{K}} & events in conditional contradictory-ranking implication & P16 IBE contradictory-ranking lemma & n/a & \hyperref[thm:p16-ibe-contradictories]{\ref*{thm:p16-ibe-contradictories}} & \thepage & K subset of complement of H; probability comparisons are premises not aesthetic ranks. \\
\phantomsection\label{entry:symbol-registry-05561}\(X,k,s,S(X)\) & \texttt{\detokenize{X,}\allowbreak\detokenize{k,}\allowbreak\detokenize{s,}\allowbreak\detokenize{S(X)}} & binary string experiment and data-dependent selection & P16 IBE prediction example & n/a & \hyperref[def:p16-ibe-selection-input]{\ref*{def:p16-ibe-selection-input}} & \thepage & Uniform observed string X length k; fixed string s; selection map S; distinct from substantive theory truth. \\
\phantomsection\label{entry:symbol-registry-05739}\(\sigma,s,t,\oplus,\otimes\) & \texttt{\detokenize{sigma,}\allowbreak\detokenize{s,}\allowbreak\detokenize{t,}\allowbreak\detokenize{oplus,}\allowbreak\detokenize{otimes}} & sign abstraction and local operations & P16 qualitative signs & n/a & \hyperref[def:p16-qr-signs]{\ref*{def:p16-qr-signs}} & \thepage & Sum set-valued; product ordinary sign; not independent global magnitudes. \\
\phantomsection\label{entry:symbol-registry-05740}\(A(q),D(q),D_s(q),I^+,I^-,\propto_Q^+\) & \texttt{\detokenize{A,}\allowbreak\detokenize{D,}\allowbreak\detokenize{D_s,}\allowbreak\detokenize{I^+,}\allowbreak\detokenize{I^-,}\allowbreak\detokenize{propto_Q^+}} & native QPT amount derivative and influences & P16 QPT subsection & n/a & \hyperref[sec:p16-qr-processes]{\ref*{sec:p16-qr-processes}} & \thepage & Direct effects additive; indirect dependencies partial and context-qualified. \\
\phantomsection\label{entry:symbol-registry-05741}\(E,u,v\) & \texttt{\detokenize{E,}\allowbreak\detokenize{u,}\allowbreak\detokenize{v}} & heat stock and positive rates & P16 heat balance & n/a & \hyperref[def:p16-qr-balance]{\ref*{def:p16-qr-balance}} & \thepage & u heating v loss; complete influence set assumed. \\
\phantomsection\label{entry:symbol-registry-05742}\(f,g,H,M^+,M^-\) & \texttt{\detokenize{f,}\allowbreak\detokenize{g,}\allowbreak\detokenize{H,}\allowbreak\detokenize{M^+,}\allowbreak\detokenize{M^-}} & strong 1986 monotonic function constraints & P16 monotonic subsection & n/a & \hyperref[thm:p16-qr-monotone]{\ref*{thm:p16-qr-monotone}} & \thepage & f=H composed with g; first argument dependent output; nonzero interior derivative. \\
\phantomsection\label{entry:symbol-registry-05743}\(P,A,Q,dP,dA,dQ\) & \texttt{\detokenize{P,}\allowbreak\detokenize{A,}\allowbreak\detokenize{Q,}\allowbreak\detokenize{dP,}\allowbreak\detokenize{dA,}\allowbreak\detokenize{dQ}} & valve pressure aperture flow and changes & P16 valve example & n/a & \hyperref[thm:p16-qr-valve]{\ref*{thm:p16-qr-valve}} & \thepage & Changes at stated qualitative mode; not stock quantities. \\
\phantomsection\label{entry:symbol-registry-05744}\(F,G,H,V_i,h_i,\beta_i,q_F,q_H\) & \texttt{\detokenize{F,}\allowbreak\detokenize{G,}\allowbreak\detokenize{H,}\allowbreak\detokenize{V_i,}\allowbreak\detokenize{h_i,}\allowbreak\detokenize{beta_i,}\allowbreak\detokenize{q_F,}\allowbreak\detokenize{q_H}} & containers amounts heads gradients and net rates & P16 container example & n/a & \hyperref[def:p16-qr-containers]{\ref*{def:p16-qr-containers}} & \thepage & H container not function; signed net rates oriented out of G and zero at equal heads. \\
\phantomsection\label{entry:symbol-registry-05745}\(x,l,h,\epsilon,t_0\) & \texttt{\detokenize{x,}\allowbreak\detokenize{l,}\allowbreak\detokenize{h,}\allowbreak\detokenize{epsilon,}\allowbreak\detokenize{t_0}} & contact comparison gap and derivative bound & P16 contact lemma & n/a & \hyperref[thm:p16-qr-contact]{\ref*{thm:p16-qr-contact}} & \thepage & h=x-l in proof; epsilon positive; C1 needed for integration. \\
\phantomsection\label{entry:symbol-registry-05746}\(inc,std,dec,I2,I3,ACTIVE\) & \texttt{\detokenize{inc,}\allowbreak\detokenize{std,}\allowbreak\detokenize{dec,}\allowbreak\detokenize{I2,}\allowbreak\detokenize{I3,}\allowbreak\detokenize{ACTIVE}} & native directions transition labels and agenda & P16 QSIM subsection & n/a & \hyperref[sec:p16-qr-qsim]{\ref*{sec:p16-qr-qsim}} & \thepage & Labels not numerical time; point and interval states distinct. \\
\phantomsection\label{entry:symbol-registry-05747}\(x,v,H,{\cal E}\) & \texttt{\detokenize{x,}\allowbreak\detokenize{v,}\allowbreak\detokenize{H,}\allowbreak\detokenize{calE}} & restoring-model position velocity function and energy & P16 energy example & n/a & \hyperref[thm:p16-qr-energy]{\ref*{thm:p16-qr-energy}} & \thepage & v velocity not heat loss; no damping external work or global periodicity claim. \\
\phantomsection\label{entry:symbol-registry-05945}\(M,O,G\) & \texttt{\detokenize{M,}\allowbreak\detokenize{O,}\allowbreak\detokenize{G}} & behavior model observations and goals & P16 compositional task comparison & n/a & \hyperref[sec:p16-cm-tasks]{\ref*{sec:p16-cm-tasks}} & \thepage & Goals not automatically implied by normal component behavior. \\
\phantomsection\label{entry:symbol-registry-05946}\(C,c,B_c,ok_c,b_c,a,p\) & \texttt{\detokenize{C,}\allowbreak\detokenize{c,}\allowbreak\detokenize{B_c,}\allowbreak\detokenize{ok_c,}\allowbreak\detokenize{b_c,}\allowbreak\detokenize{a,}\allowbreak\detokenize{p}} & component modes and complete or partial assignments & P16 component model construction & n/a & \hyperref[def:p16-cm-library]{\ref*{def:p16-cm-library}} & \thepage & One mode per component; partial assignment p not probability. \\
\phantomsection\label{entry:symbol-registry-05947}\(L,S,F_{c,b}\) & \texttt{\detokenize{L,}\allowbreak\detokenize{S,}\allowbreak\detokenize{F_cb}} & library connections and guarded behavior fragments & P16 component composition & n/a & \hyperref[def:p16-cm-library]{\ref*{def:p16-cm-library}} & \thepage & S is structure here; state spaces S\_c are a later local role. \\
\phantomsection\label{entry:symbol-registry-05948}\(D,R_{c,b},R_a,V,P_a,\pi_V\) & \texttt{\detokenize{D,}\allowbreak\detokenize{R_cb,}\allowbreak\detokenize{R_a,}\allowbreak\detokenize{V,}\allowbreak\detokenize{P_a,}\allowbreak\detokenize{pi_V}} & global domain local and composite relations observable projection & P16 composition proof & n/a & \hyperref[thm:p16-cm-composition]{\ref*{thm:p16-cm-composition}} & \thepage & Join before exact projection; P\_a not probability. \\
\phantomsection\label{entry:symbol-registry-05949}\(\mathcal K,\neg p\) & \texttt{\detokenize{calK,}\allowbreak\detokenize{negp}} & minimal mode conflict clauses & P16 conflict characterization & n/a & \hyperref[thm:p16-cm-mode-conflicts]{\ref*{thm:p16-cm-mode-conflicts}} & \thepage & Relative to exhaustive exclusive mode-domain axioms. \\
\phantomsection\label{entry:symbol-registry-05950}\(N\) & \texttt{\detokenize{N}} & number of complete assignments & P16 finite candidate search & n/a & \hyperref[thm:p16-cm-search]{\ref*{thm:p16-cm-search}} & \thepage & Product of finite mode-set cardinalities; not polynomial bound. \\
\phantomsection\label{entry:symbol-registry-05951}\(B,R,T,H,w_i,K_i\) & \texttt{\detokenize{B,}\allowbreak\detokenize{R,}\allowbreak\detokenize{T,}\allowbreak\detokenize{H,}\allowbreak\detokenize{w_i,}\allowbreak\detokenize{K_i}} & battery rear light starter head light wires and conflict supports & P16 source circuit example & n/a & \hyperref[def:p16-cm-lights]{\ref*{def:p16-cm-lights}} & \thepage & R is rear light T starter H head light only in this example. \\
\phantomsection\label{entry:symbol-registry-05952}\(x_j,y_j\) & \texttt{\detokenize{x_j,}\allowbreak\detokenize{y_j}} & binary left and right branch port potentials & P16 explicit source-topology circuit & n/a & \hyperref[def:p16-cm-circuit-realization]{\ref*{def:p16-cm-circuit-realization}} & \thepage & Normality-only voltage-presence idealization; observations kept for faulty outputs. \\
\phantomsection\label{entry:symbol-registry-05953}\(A_0,C_0,D_0\) & \texttt{\detokenize{A_0,}\allowbreak\detokenize{C_0,}\allowbreak\detokenize{D_0}} & disjoint component groups for complete hitting-set calculation & P16 source circuit localization proof & n/a & \hyperref[thm:p16-cm-circuit]{\ref*{thm:p16-cm-circuit}} & \thepage & Nine pairs twelve triples and the rear-light/starter pair; twenty-one alternatives. \\
\phantomsection\label{entry:symbol-registry-05954}\(n,k,p,q\) & \texttt{\detokenize{n,}\allowbreak\detokenize{k,}\allowbreak\detokenize{p,}\allowbreak\detokenize{q}} & chain length faulty count positive and negative increments & P16 equal-component chain & n/a & \hyperref[thm:p16-cm-equal-chain]{\ref*{thm:p16-cm-equal-chain}} & \thepage & p q are integer counts here; distinct from assignment and probability roles. \\
\phantomsection\label{entry:symbol-registry-05955}\(W,\operatorname{Th},>,E\) & \texttt{\detokenize{W,}\allowbreak\detokenize{Th,}\allowbreak\detokenize{>,}\allowbreak\detokenize{E}} & background deductive closure local mode preference and extension & P16 preference-default proof & n/a & \hyperref[thm:p16-cm-default-sound]{\ref*{thm:p16-cm-default-sound}} & \thepage & E is an extension; event and test-set symbols have later typed scopes. \\
\phantomsection\label{entry:symbol-registry-05956}\(O_H,x,x'\) & \texttt{\detokenize{O_H,}\allowbreak\detokenize{x,}\allowbreak\detokenize{xprime}} & observation history state variable and time derivative & P16 temporal diagnosis & n/a & \hyperref[thm:p16-cm-temporal]{\ref*{thm:p16-cm-temporal}} & \thepage & Snapshot witnesses need not share a realizable trajectory. \\
\phantomsection\label{entry:symbol-registry-05957}\(S_c,s_{c,0},E_c,T_c\) & \texttt{\detokenize{S_c,}\allowbreak\detokenize{s_c0,}\allowbreak\detokenize{E_c,}\allowbreak\detokenize{T_c}} & component states initial state events and triggered-emitted transitions & P16 transition model & n/a & \hyperref[def:p16-cm-transition-model]{\ref*{def:p16-cm-transition-model}} & \thepage & Trigger and emitted event distinct; synchronization fixed before composition. \\
\phantomsection\label{entry:symbol-registry-05958}\(Q,q_0,F,J,\epsilon\) & \texttt{\detokenize{Q,}\allowbreak\detokenize{q_0,}\allowbreak\detokenize{F,}\allowbreak\detokenize{J,}\allowbreak\detokenize{epsilon}} & finite graph initial state fault labels accumulated fault set hidden event & P16 finite trace observer & n/a & \hyperref[def:p16-cm-transition-model]{\ref*{def:p16-cm-transition-model}} & \thepage & epsilon is silent edge; fault order and multiplicity not represented. \\
\phantomsection\label{entry:symbol-registry-05959}\(O_C,O_A\) & \texttt{\detokenize{O_C,}\allowbreak\detokenize{O_A}} & conditioning observations and symptoms to entail & P16 abductive diagnosis & n/a & \hyperref[sec:p16-cm-process-assessment]{\ref*{sec:p16-cm-process-assessment}} & \thepage & Commands need not be explained by faults. \\
\phantomsection\label{entry:symbol-registry-05960}\(u\) & \texttt{\detokenize{u}} & unknown binary supply state & P16 parallel-bulb example & n/a & \hyperref[def:p16-cm-parallel-bulbs]{\ref*{def:p16-cm-parallel-bulbs}} & \thepage & Different local role from test stimulus u. \\
\phantomsection\label{entry:symbol-registry-05961}\(I^+(x,y)\) & \texttt{\detokenize{I^+(x,}\allowbreak\detokenize{y)}} & Struss positive influence of x on derivative of y & P16 process assessment & n/a & \hyperref[def:p16-cm-process-fragment]{\ref*{def:p16-cm-process-fragment}} & \thepage & Source argument order differs from Forbus affected-quantity/rate notation. \\
\phantomsection\label{entry:symbol-registry-05962}\(U,Z,A,u,z,P_i,P_i(u),R_i\) & \texttt{\detokenize{U,}\allowbreak\detokenize{Z,}\allowbreak\detokenize{A,}\allowbreak\detokenize{u,}\allowbreak\detokenize{z,}\allowbreak\detokenize{P_i,}\allowbreak\detokenize{P_i(u),}\allowbreak\detokenize{R_i}} & stimuli responses admissibility and projected behavior relations & P16 relational testing & n/a & \hyperref[def:p16-cm-test-relations]{\ref*{def:p16-cm-test-relations}} & \thepage & Each admissible stimulus must have a response under each hypothesis. \\
\phantomsection\label{entry:symbol-registry-05963}\(D_{ij},B_{ij},E_{ij},T\) & \texttt{\detokenize{D_ij,}\allowbreak\detokenize{B_ij,}\allowbreak\detokenize{E_ij,}\allowbreak\detokenize{T}} & definite two-sided possible any informative stimuli and a test set & P16 test construction & n/a & \hyperref[thm:p16-cm-discrimination]{\ref*{thm:p16-cm-discrimination}} & \thepage & Two-sided quantifier distinguished from unilateral information. \\
\phantomsection\label{entry:symbol-registry-05964}\(p_i,q_i,r,H,p_{unk},m\) & \texttt{\detokenize{p_i,}\allowbreak\detokenize{q_i,}\allowbreak\detokenize{r,}\allowbreak\detokenize{H,}\allowbreak\detokenize{p_unk,}\allowbreak\detokenize{m}} & hypothesis weights response channels mixture entropy unknown mass and domain size & P16 information selection & n/a & \hyperref[thm:p16-cm-information]{\ref*{thm:p16-cm-information}} & \thepage & Finite response alphabet; probabilities are supplied not implied by relation. \\
\phantomsection\label{entry:symbol-registry-05965}\(O_P,s_A,s_G\) & \texttt{\detokenize{O_P,}\allowbreak\detokenize{s_A,}\allowbreak\detokenize{s_G}} & persistent observations actual and target component state assignments & P16 reconfiguration & n/a & \hyperref[def:p16-cm-reconfiguration]{\ref*{def:p16-cm-reconfiguration}} & \thepage & State domains exhaustive; persistence determined before goal checking. \\
\phantomsection\label{entry:symbol-registry-05966}\(d,R_i(d),G_i\) & \texttt{\detokenize{d,}\allowbreak\detokenize{R_i(d),}\allowbreak\detokenize{G_i}} & candidate intervention permitted outcomes and hypothesis-specific goal outcomes & P16 remedy proof & n/a & \hyperref[thm:p16-cm-remedies]{\ref*{thm:p16-cm-remedies}} & \thepage & Nonempty outcomes required; guarantees model-relative; plans need separate reachability. \\
\phantomsection\label{entry:symbol-registry-06223}\(Temperature(o),Place(o),On(a,b)\) & \texttt{\detokenize{Temperature,}\allowbreak\detokenize{Place,}\allowbreak\detokenize{On}} & native Davis numeric region and Boolean time fluents & P16 Davis architecture & n/a & \hyperref[sec:p16-phys-architectures]{\ref*{sec:p16-phys-architectures}} & \thepage & Parameter means numeric-valued fluent; not a propositional causal-rule head. \\
\phantomsection\label{entry:symbol-registry-06224}\(InCurrent,Voltage,R,C,u,v,v_0,T\) & \texttt{\detokenize{InCurrent,}\allowbreak\detokenize{Voltage,}\allowbreak\detokenize{R,}\allowbreak\detokenize{C,}\allowbreak\detokenize{u,}\allowbreak\detokenize{v,}\allowbreak\detokenize{v_0,}\allowbreak\detokenize{T}} & port currents voltages resistance capacitance input output and initial data & P16 Davis RC construction & n/a & \hyperref[def:p16-phys-rc]{\ref*{def:p16-phys-rc}} & \thepage & Incoming-current orientations explicit; R C positive; T positive horizon. \\
\phantomsection\label{entry:symbol-registry-06225}\(V,V_0,B,r\) & \texttt{\detokenize{V,}\allowbreak\detokenize{V_0,}\allowbreak\detokenize{B,}\allowbreak\detokenize{r}} & bucket volume initial volume capacity and stipulated net rate & P16 Davis process action event example & n/a & \hyperref[def:p16-phys-bucket]{\ref*{def:p16-phys-bucket}} & \thepage & Mandatory structural event at half capacity differs from optional tap action. \\
\phantomsection\label{entry:symbol-registry-06226}\(S_o,Q_o,b_o,I,d\) & \texttt{\detokenize{S_o,}\allowbreak\detokenize{Q_o,}\allowbreak\detokenize{b_o,}\allowbreak\detokenize{I,}\allowbreak\detokenize{d}} & shape proper rotation translation identity matrix and dimension & P16 rigid geometry & n/a & \hyperref[def:p16-phys-configuration]{\ref*{def:p16-phys-configuration}} & \thepage & I is matrix identity; rotations exclude reflection; d two or three. \\
\phantomsection\label{entry:symbol-registry-06227}\(Shape,Position,Place,\mathcal F,q_0,q_1\) & \texttt{\detokenize{Shape,}\allowbreak\detokenize{Position,}\allowbreak\detokenize{Place,}\allowbreak\detokenize{calF,}\allowbreak\detokenize{q_0,}\allowbreak\detokenize{q_1}} & source shape pose occupied-region functions and feasible configurations & P16 physical kinematics & n/a & \hyperref[def:p16-phys-configuration]{\ref*{def:p16-phys-configuration}} & \thepage & Position and Place are fluents; contact allowed; path component not arbitrary connected component. \\
\phantomsection\label{entry:symbol-registry-06228}\(M,\|\cdot\|\) & \texttt{\detokenize{M,}\allowbreak\detokenize{norm}} & shape radius bound and Euclidean/operator norms & P16 occupied-set continuity proof & n/a & \hyperref[annex:p16.phys-kinematics]{\ref*{annex:p16.phys-kinematics}} & \thepage & M local to Hausdorff estimate; compact shapes stipulated. \\
\phantomsection\label{entry:symbol-registry-06229}\(C_i,n\) & \texttt{\detokenize{C_i,}\allowbreak\detokenize{n}} & finite path-connected closed cell cover and its size & P16 exact configuration graph & n/a & \hyperref[thm:p16-phys-cells]{\ref*{thm:p16-phys-cells}} & \thepage & C\_i is a geometric subset not capacitance; actual intersections define edges. \\
\phantomsection\label{entry:symbol-registry-06230}\(m(t),q(t),x,y\) & \texttt{\detokenize{m,}\allowbreak\detokenize{q,}\allowbreak\detokenize{x,}\allowbreak\detokenize{y}} & commanded motion full configuration lift and toy constrained coordinates & P16 manipulator feasibility & n/a & \hyperref[thm:p16-phys-action-lift]{\ref*{thm:p16-phys-action-lift}} & \thepage & Pointwise frames do not ensure continuous path from initial state. \\
\phantomsection\label{entry:symbol-registry-06231}\(N,\mu,m,g\) & \texttt{\detokenize{N,}\allowbreak\detokenize{mu,}\allowbreak\detokenize{m,}\allowbreak\detokenize{g}} & specified normal preload friction coefficient mass and gravity & P16 nail static model & n/a & \hyperref[def:p16-phys-nail]{\ref*{def:p16-phys-nail}} & \thepage & m is mass here; preload not derived from geometry. \\
\phantomsection\label{entry:symbol-registry-06232}\(\Phi_t,K,E,D,\Omega,E_\infty\) & \texttt{\detokenize{Phi,}\allowbreak\detokenize{K,}\allowbreak\detokenize{E,}\allowbreak\detokenize{D,}\allowbreak\detokenize{Omega,}\allowbreak\detokenize{E_infty}} & continuous forward flow trap energy derivative limiting set and energy limit & P16 bounded exit proof & n/a & \hyperref[thm:p16-phys-exit]{\ref*{thm:p16-phys-exit}} & \thepage & Continuous-flow theorem; not arbitrary nonsmooth impact dynamics. \\
\phantomsection\label{entry:symbol-registry-06233}\(q,v,U,\gamma\) & \texttt{\detokenize{q,}\allowbreak\detokenize{v,}\allowbreak\detokenize{U,}\allowbreak\detokenize{gamma}} & position velocity potential and damping in sufficient progress model & P16 smooth damped model & n/a & \hyperref[thm:p16-phys-exit]{\ref*{thm:p16-phys-exit}} & \thepage & U potential energy here; not the later logical theory U. \\
\phantomsection\label{entry:symbol-registry-06234}\(x,x^2\) & \texttt{\detokenize{x,}\allowbreak\detokenize{x^2}} & abstract asymptotic exit coordinate and energy & P16 energy-only failure case & n/a & \hyperref[thm:p16-phys-exit]{\ref*{thm:p16-phys-exit}} & \thepage & x positive solves xprime equals minus x and never hits zero finitely. \\
\phantomsection\label{entry:symbol-registry-06235}\(Volume(l,r),Flow(l,b),Place(c),LiquidTrajectory(X,L)\) & \texttt{\detokenize{Volume,}\allowbreak\detokenize{Flow,}\allowbreak\detokenize{Place,}\allowbreak\detokenize{LiquidTrajectory}} & native region chunk and material-particle liquid representations & P16 liquid ontology & n/a & \hyperref[sec:p16-phys-liquid-ontology]{\ref*{sec:p16-phys-liquid-ontology}} & \thepage & Material chunk differs from moving observation region; chunk may be disconnected. \\
\phantomsection\label{entry:symbol-registry-06236}\(L(t),r(t),V,F_{rel},v,w,n_r,n_L\) & \texttt{\detokenize{L,}\allowbreak\detokenize{r,}\allowbreak\detokenize{V,}\allowbreak\detokenize{F_rel,}\allowbreak\detokenize{v,}\allowbreak\detokenize{w,}\allowbreak\detokenize{n_r,}\allowbreak\detokenize{n_L}} & material/control regions volume relative outflow velocities and normals & P16 moving liquid boundary & n/a & \hyperref[def:p16-phys-liquid-balance]{\ref*{def:p16-phys-liquid-balance}} & \thepage & v minus w gives relative flux; smooth transverse scope explicit. \\
\phantomsection\label{entry:symbol-registry-06237}\(\operatorname{vol},dA,\nabla\cdot,h\) & \texttt{\detokenize{vol,}\allowbreak\detokenize{dA,}\allowbreak\detokenize{div,}\allowbreak\detokenize{h}} & region measure surface area divergence and small transport increment & P16 liquid transport proof & n/a & \hyperref[annex:p16.phys-liquid-balance]{\ref*{annex:p16.phys-liquid-balance}} & \thepage & h local time increment not hydrostatic height; boundary-layer hypotheses declared. \\
\phantomsection\label{entry:symbol-registry-06238}\(p,\rho,g,e_z,h_i\) & \texttt{\detokenize{p,}\allowbreak\detokenize{rho,}\allowbreak\detokenize{g,}\allowbreak\detokenize{e_z,}\allowbreak\detokenize{h_i}} & fluid pressure constant density gravity upward vector and surface heights & P16 hydrostatics & n/a & \hyperref[thm:p16-phys-liquid-level]{\ref*{thm:p16-phys-liquid-level}} & \thepage & Pressure extends to relevant free surfaces; trapped gas retained. \\
\phantomsection\label{entry:symbol-registry-06239}\(A(h),a(z),V,h_0,h_r\) & \texttt{\detokenize{A,}\allowbreak\detokenize{a,}\allowbreak\detokenize{V,}\allowbreak\detokenize{h_0,}\allowbreak\detokenize{h_r}} & below-level capacity cross section volume bottom and lowest opening & P16 liquid capacity & n/a & \hyperref[thm:p16-phys-liquid-level]{\ref*{thm:p16-phys-liquid-level}} & \thepage & Strict integral increase and accessible open container; not inertial spilling proof. \\
\phantomsection\label{entry:symbol-registry-06240}\(U,V,A\) & \texttt{\detokenize{U,}\allowbreak\detokenize{V,}\allowbreak\detokenize{A}} & same-vocabulary stronger weaker and added-constraint theories & P16 physical abstraction & n/a & \hyperref[thm:p16-phys-abstraction]{\ref*{thm:p16-phys-abstraction}} & \thepage & Model inclusion not physical validity or efficiency; typed separately from volume/potential. \\
\phantomsection\label{entry:symbol-registry-06241}\(r,w,x\) & \texttt{\detokenize{r,}\allowbreak\detokenize{w,}\allowbreak\detokenize{x}} & disk radius corridor width and horizontal center & P16 point-object abstraction failure & n/a & \hyperref[thm:p16-phys-abstraction]{\ref*{thm:p16-phys-abstraction}} & \thepage & Contact allowed at w equals twice r; kinematic not force feasibility. \\
\phantomsection\label{entry:symbol-registry-06242}\(v_0,a,L,v_{min},v_{max},a_{min},a_{max},d_{min},d_{max}\) & \texttt{\detokenize{v_0,}\allowbreak\detokenize{a,}\allowbreak\detokenize{L,}\allowbreak\detokenize{vmin,}\allowbreak\detokenize{vmax,}\allowbreak\detokenize{amin,}\allowbreak\detokenize{amax,}\allowbreak\detokenize{dmin,}\allowbreak\detokenize{dmax}} & speed deceleration edge distance parameter bounds and stopping range & P16 quantitative cliff margin & n/a & \hyperref[def:p16-phys-braking]{\ref*{def:p16-phys-braking}} & \thepage & Independent rectangle and constant braking until stop; equality boundary retained. \\
\phantomsection\label{entry:symbol-registry-07937}\(SD\) & \texttt{\detokenize{SD}} & native library structure hypotheses persistent data and finite goals symbol family member \(SD\) & P16 common targets & n/a & \hyperref[def:p16-rem-target-input]{\ref*{def:p16-rem-target-input}} & \thepage & Fault modes fixed for target selection; contexts tagged. Explicit displayed family; member types and common actual first-use exposed, not one atomic glyph. Family listing, not one atomic glyph. Author macro \texttt{\detokenize{SD,}\allowbreak\detokenize{LIB,}\allowbreak\detokenize{STRUCTURE,}\allowbreak\detokenize{M\_i,}\allowbreak\detokenize{O\_Pi,}\allowbreak\detokenize{W\_i,}\allowbreak\detokenize{G\_i}} not carried to the members. \\
\phantomsection\label{entry:symbol-registry-07938}\(LIB\) & \texttt{\detokenize{LIB}} & native library structure hypotheses persistent data and finite goals symbol family member \(LIB\) & P16 common targets & n/a & \hyperref[def:p16-rem-target-input]{\ref*{def:p16-rem-target-input}} & \thepage & Fault modes fixed for target selection; contexts tagged. Explicit displayed family; member types and common actual first-use exposed, not one atomic glyph. Family listing, not one atomic glyph. Author macro \texttt{\detokenize{SD,}\allowbreak\detokenize{LIB,}\allowbreak\detokenize{STRUCTURE,}\allowbreak\detokenize{M\_i,}\allowbreak\detokenize{O\_Pi,}\allowbreak\detokenize{W\_i,}\allowbreak\detokenize{G\_i}} not carried to the members. \\
\phantomsection\label{entry:symbol-registry-07939}\(STRUCTURE\) & \texttt{\detokenize{STRUCTURE}} & native library structure hypotheses persistent data and finite goals symbol family member \(STRUCTURE\) & P16 common targets & n/a & \hyperref[def:p16-rem-target-input]{\ref*{def:p16-rem-target-input}} & \thepage & Fault modes fixed for target selection; contexts tagged. Explicit displayed family; member types and common actual first-use exposed, not one atomic glyph. Family listing, not one atomic glyph. Author macro \texttt{\detokenize{SD,}\allowbreak\detokenize{LIB,}\allowbreak\detokenize{STRUCTURE,}\allowbreak\detokenize{M\_i,}\allowbreak\detokenize{O\_Pi,}\allowbreak\detokenize{W\_i,}\allowbreak\detokenize{G\_i}} not carried to the members. \\
\phantomsection\label{entry:symbol-registry-07940}\(M_i\) & \texttt{\detokenize{M_i}} & native library structure hypotheses persistent data and finite goals symbol family member \(M_i\) & P16 common targets & n/a & \hyperref[def:p16-rem-target-input]{\ref*{def:p16-rem-target-input}} & \thepage & Fault modes fixed for target selection; contexts tagged. Explicit displayed family; member types and common actual first-use exposed, not one atomic glyph. Family listing, not one atomic glyph. Author macro \texttt{\detokenize{SD,}\allowbreak\detokenize{LIB,}\allowbreak\detokenize{STRUCTURE,}\allowbreak\detokenize{M\_i,}\allowbreak\detokenize{O\_Pi,}\allowbreak\detokenize{W\_i,}\allowbreak\detokenize{G\_i}} not carried to the members. \\
\phantomsection\label{entry:symbol-registry-07941}\(O_{P,i}\) & \texttt{\detokenize{O_{P,i}}} & native library structure hypotheses persistent data and finite goals symbol family member \(O_{P,i}\) & P16 common targets & n/a & \hyperref[def:p16-rem-target-input]{\ref*{def:p16-rem-target-input}} & \thepage & Fault modes fixed for target selection; contexts tagged. Explicit displayed family; member types and common actual first-use exposed, not one atomic glyph. Family listing, not one atomic glyph. Author macro \texttt{\detokenize{SD,}\allowbreak\detokenize{LIB,}\allowbreak\detokenize{STRUCTURE,}\allowbreak\detokenize{M\_i,}\allowbreak\detokenize{O\_Pi,}\allowbreak\detokenize{W\_i,}\allowbreak\detokenize{G\_i}} not carried to the members. \\
\phantomsection\label{entry:symbol-registry-07942}\(W_i\) & \texttt{\detokenize{W_i}} & native library structure hypotheses persistent data and finite goals symbol family member \(W_i\) & P16 common targets & n/a & \hyperref[def:p16-rem-target-input]{\ref*{def:p16-rem-target-input}} & \thepage & Fault modes fixed for target selection; contexts tagged. Explicit displayed family; member types and common actual first-use exposed, not one atomic glyph. Family listing, not one atomic glyph. Author macro \texttt{\detokenize{SD,}\allowbreak\detokenize{LIB,}\allowbreak\detokenize{STRUCTURE,}\allowbreak\detokenize{M\_i,}\allowbreak\detokenize{O\_Pi,}\allowbreak\detokenize{W\_i,}\allowbreak\detokenize{G\_i}} not carried to the members. \\
\phantomsection\label{entry:symbol-registry-07943}\(G_i\) & \texttt{\detokenize{G_i}} & native library structure hypotheses persistent data and finite goals symbol family member \(G_i\) & P16 common targets & n/a & \hyperref[def:p16-rem-target-input]{\ref*{def:p16-rem-target-input}} & \thepage & Fault modes fixed for target selection; contexts tagged. Explicit displayed family; member types and common actual first-use exposed, not one atomic glyph. Family listing, not one atomic glyph. Author macro \texttt{\detokenize{SD,}\allowbreak\detokenize{LIB,}\allowbreak\detokenize{STRUCTURE,}\allowbreak\detokenize{M\_i,}\allowbreak\detokenize{O\_Pi,}\allowbreak\detokenize{W\_i,}\allowbreak\detokenize{G\_i}} not carried to the members. \\
\phantomsection\label{entry:symbol-registry-07944}\(C_R\) & \texttt{\detokenize{C_R}} & reconfigurable domains targets and actual/common-state changes symbol family member \(C_R\) & P16 target minimality & n/a & \hyperref[def:p16-rem-change]{\ref*{def:p16-rem-change}} & \thepage & Common-entailed-state measure not actual change count. Explicit displayed family; member types and common actual first-use exposed, not one atomic glyph. Family listing, not one atomic glyph. Author macro \texttt{\detokenize{C\_R,}\allowbreak\detokenize{S\_c,}\allowbreak\detokenize{s,}\allowbreak\detokenize{s\_A,}\allowbreak\detokenize{s\_Mi,}\allowbreak\detokenize{s\_common,}\allowbreak\detokenize{Delta}} not carried to the members. \\
\phantomsection\label{entry:symbol-registry-07945}\(S_c\) & \texttt{\detokenize{S_c}} & reconfigurable domains targets and actual/common-state changes symbol family member \(S_c\) & P16 target minimality & n/a & \hyperref[def:p16-rem-change]{\ref*{def:p16-rem-change}} & \thepage & Common-entailed-state measure not actual change count. Explicit displayed family; member types and common actual first-use exposed, not one atomic glyph. Family listing, not one atomic glyph. Author macro \texttt{\detokenize{C\_R,}\allowbreak\detokenize{S\_c,}\allowbreak\detokenize{s,}\allowbreak\detokenize{s\_A,}\allowbreak\detokenize{s\_Mi,}\allowbreak\detokenize{s\_common,}\allowbreak\detokenize{Delta}} not carried to the members. \\
\phantomsection\label{entry:symbol-registry-07946}\(s\) & \texttt{\detokenize{s}} & reconfigurable domains targets and actual/common-state changes symbol family member \(s\) & P16 target minimality & n/a & \hyperref[def:p16-rem-change]{\ref*{def:p16-rem-change}} & \thepage & Common-entailed-state measure not actual change count. Explicit displayed family; member types and common actual first-use exposed, not one atomic glyph. Family listing, not one atomic glyph. Author macro \texttt{\detokenize{C\_R,}\allowbreak\detokenize{S\_c,}\allowbreak\detokenize{s,}\allowbreak\detokenize{s\_A,}\allowbreak\detokenize{s\_Mi,}\allowbreak\detokenize{s\_common,}\allowbreak\detokenize{Delta}} not carried to the members. \\
\phantomsection\label{entry:symbol-registry-07947}\(s_A\) & \texttt{\detokenize{s_A}} & reconfigurable domains targets and actual/common-state changes symbol family member \(s_A\) & P16 target minimality & n/a & \hyperref[def:p16-rem-change]{\ref*{def:p16-rem-change}} & \thepage & Common-entailed-state measure not actual change count. Explicit displayed family; member types and common actual first-use exposed, not one atomic glyph. Family listing, not one atomic glyph. Author macro \texttt{\detokenize{C\_R,}\allowbreak\detokenize{S\_c,}\allowbreak\detokenize{s,}\allowbreak\detokenize{s\_A,}\allowbreak\detokenize{s\_Mi,}\allowbreak\detokenize{s\_common,}\allowbreak\detokenize{Delta}} not carried to the members. \\
\phantomsection\label{entry:symbol-registry-07948}\(s_{M,i}\) & \texttt{\detokenize{s_{M,i}}} & reconfigurable domains targets and actual/common-state changes symbol family member \(s_{M,i}\) & P16 target minimality & n/a & \hyperref[def:p16-rem-change]{\ref*{def:p16-rem-change}} & \thepage & Common-entailed-state measure not actual change count. Explicit displayed family; member types and common actual first-use exposed, not one atomic glyph. Family listing, not one atomic glyph. Author macro \texttt{\detokenize{C\_R,}\allowbreak\detokenize{S\_c,}\allowbreak\detokenize{s,}\allowbreak\detokenize{s\_A,}\allowbreak\detokenize{s\_Mi,}\allowbreak\detokenize{s\_common,}\allowbreak\detokenize{Delta}} not carried to the members. \\
\phantomsection\label{entry:symbol-registry-07949}\(s_{common}\) & \texttt{\detokenize{s_{common}}} & reconfigurable domains targets and actual/common-state changes symbol family member \(s_{common}\) & P16 target minimality & n/a & \hyperref[def:p16-rem-change]{\ref*{def:p16-rem-change}} & \thepage & Common-entailed-state measure not actual change count. Explicit displayed family; member types and common actual first-use exposed, not one atomic glyph. Family listing, not one atomic glyph. Author macro \texttt{\detokenize{C\_R,}\allowbreak\detokenize{S\_c,}\allowbreak\detokenize{s,}\allowbreak\detokenize{s\_A,}\allowbreak\detokenize{s\_Mi,}\allowbreak\detokenize{s\_common,}\allowbreak\detokenize{Delta}} not carried to the members. \\
\phantomsection\label{entry:symbol-registry-07950}\(\Delta(s)\) & \texttt{\detokenize{\Delta(s)}} & reconfigurable domains targets and actual/common-state changes symbol family member \(\Delta(s)\) & P16 target minimality & n/a & \hyperref[def:p16-rem-change]{\ref*{def:p16-rem-change}} & \thepage & Common-entailed-state measure not actual change count. Explicit displayed family; member types and common actual first-use exposed, not one atomic glyph. Family listing, not one atomic glyph. Author macro \texttt{\detokenize{C\_R,}\allowbreak\detokenize{S\_c,}\allowbreak\detokenize{s,}\allowbreak\detokenize{s\_A,}\allowbreak\detokenize{s\_Mi,}\allowbreak\detokenize{s\_common,}\allowbreak\detokenize{Delta}} not carried to the members. \\
\phantomsection\label{entry:symbol-registry-07951}\(N\) & \texttt{\detokenize{N}} & target assignment count and remaining hypothesis count symbol family member \(N\) & P16 target search & n/a & \hyperref[thm:p16-rem-target-search]{\ref*{thm:p16-rem-target-search}} & \thepage & Finite product and total exact checks required. Explicit displayed family; member types and common actual first-use exposed, not one atomic glyph. Family listing, not one atomic glyph. Author macro \texttt{\detokenize{N,}\allowbreak\detokenize{k}} not carried to the members. \\
\phantomsection\label{entry:symbol-registry-07952}\(k\) & \texttt{\detokenize{k}} & target assignment count and remaining hypothesis count symbol family member \(k\) & P16 target search & n/a & \hyperref[thm:p16-rem-target-search]{\ref*{thm:p16-rem-target-search}} & \thepage & Finite product and total exact checks required. Explicit displayed family; member types and common actual first-use exposed, not one atomic glyph. Family listing, not one atomic glyph. Author macro \texttt{\detokenize{N,}\allowbreak\detokenize{k}} not carried to the members. \\
\phantomsection\label{entry:symbol-registry-07953}\(MH\) & \texttt{\detokenize{MH}} & native mode action observation and goal histories symbol family member \(MH\) & P16 source remedy criterion & n/a & \hyperref[sec:p16-rem-execution]{\ref*{sec:p16-rem-execution}} & \thepage & Intermediate goals and applicability retained. Explicit displayed family; member types and common actual first-use exposed, not one atomic glyph. Family listing, not one atomic glyph. Author macro \texttt{\detokenize{MH,}\allowbreak\detokenize{AH,}\allowbreak\detokenize{OBSH,}\allowbreak\detokenize{GOALSH}} not carried to the members. \\
\phantomsection\label{entry:symbol-registry-07954}\(AH\) & \texttt{\detokenize{AH}} & native mode action observation and goal histories symbol family member \(AH\) & P16 source remedy criterion & n/a & \hyperref[sec:p16-rem-execution]{\ref*{sec:p16-rem-execution}} & \thepage & Intermediate goals and applicability retained. Explicit displayed family; member types and common actual first-use exposed, not one atomic glyph. Family listing, not one atomic glyph. Author macro \texttt{\detokenize{MH,}\allowbreak\detokenize{AH,}\allowbreak\detokenize{OBSH,}\allowbreak\detokenize{GOALSH}} not carried to the members. \\
\phantomsection\label{entry:symbol-registry-07955}\(OBSH\) & \texttt{\detokenize{OBSH}} & native mode action observation and goal histories symbol family member \(OBSH\) & P16 source remedy criterion & n/a & \hyperref[sec:p16-rem-execution]{\ref*{sec:p16-rem-execution}} & \thepage & Intermediate goals and applicability retained. Explicit displayed family; member types and common actual first-use exposed, not one atomic glyph. Family listing, not one atomic glyph. Author macro \texttt{\detokenize{MH,}\allowbreak\detokenize{AH,}\allowbreak\detokenize{OBSH,}\allowbreak\detokenize{GOALSH}} not carried to the members. \\
\phantomsection\label{entry:symbol-registry-07956}\(GOALSH\) & \texttt{\detokenize{GOALSH}} & native mode action observation and goal histories symbol family member \(GOALSH\) & P16 source remedy criterion & n/a & \hyperref[sec:p16-rem-execution]{\ref*{sec:p16-rem-execution}} & \thepage & Intermediate goals and applicability retained. Explicit displayed family; member types and common actual first-use exposed, not one atomic glyph. Family listing, not one atomic glyph. Author macro \texttt{\detokenize{MH,}\allowbreak\detokenize{AH,}\allowbreak\detokenize{OBSH,}\allowbreak\detokenize{GOALSH}} not carried to the members. \\
\phantomsection\label{entry:symbol-registry-07957}\(X_i\) & \texttt{\detokenize{X_i}} & tagged execution states possibilities transitions safety and goals symbol family member \(X_i\) & P16 execution realization & n/a & \hyperref[def:p16-rem-execution-model]{\ref*{def:p16-rem-execution-model}} & \thepage & All relevant history memory included; repairs may change current modes. Explicit displayed family; member types and common actual first-use exposed, not one atomic glyph. Family listing, not one atomic glyph. Author macro \texttt{\detokenize{X\_i,}\allowbreak\detokenize{calX,}\allowbreak\detokenize{B\_0,}\allowbreak\detokenize{T\_a,}\allowbreak\detokenize{calS,}\allowbreak\detokenize{calG}} not carried to the members. \\
\phantomsection\label{entry:symbol-registry-07958}\(\mathcal X\) & \texttt{\detokenize{\mathcal X}} & tagged execution states possibilities transitions safety and goals symbol family member \(\mathcal X\) & P16 execution realization & n/a & \hyperref[def:p16-rem-execution-model]{\ref*{def:p16-rem-execution-model}} & \thepage & All relevant history memory included; repairs may change current modes. Explicit displayed family; member types and common actual first-use exposed, not one atomic glyph. Family listing, not one atomic glyph. Author macro \texttt{\detokenize{X\_i,}\allowbreak\detokenize{calX,}\allowbreak\detokenize{B\_0,}\allowbreak\detokenize{T\_a,}\allowbreak\detokenize{calS,}\allowbreak\detokenize{calG}} not carried to the members. \\
\phantomsection\label{entry:symbol-registry-07959}\(B_0\) & \texttt{\detokenize{B_0}} & tagged execution states possibilities transitions safety and goals symbol family member \(B_0\) & P16 execution realization & n/a & \hyperref[def:p16-rem-execution-model]{\ref*{def:p16-rem-execution-model}} & \thepage & All relevant history memory included; repairs may change current modes. Explicit displayed family; member types and common actual first-use exposed, not one atomic glyph. Family listing, not one atomic glyph. Author macro \texttt{\detokenize{X\_i,}\allowbreak\detokenize{calX,}\allowbreak\detokenize{B\_0,}\allowbreak\detokenize{T\_a,}\allowbreak\detokenize{calS,}\allowbreak\detokenize{calG}} not carried to the members. \\
\phantomsection\label{entry:symbol-registry-07960}\(T_a\) & \texttt{\detokenize{T_a}} & tagged execution states possibilities transitions safety and goals symbol family member \(T_a\) & P16 execution realization & n/a & \hyperref[def:p16-rem-execution-model]{\ref*{def:p16-rem-execution-model}} & \thepage & All relevant history memory included; repairs may change current modes. Explicit displayed family; member types and common actual first-use exposed, not one atomic glyph. Family listing, not one atomic glyph. Author macro \texttt{\detokenize{X\_i,}\allowbreak\detokenize{calX,}\allowbreak\detokenize{B\_0,}\allowbreak\detokenize{T\_a,}\allowbreak\detokenize{calS,}\allowbreak\detokenize{calG}} not carried to the members. \\
\phantomsection\label{entry:symbol-registry-07961}\(\mathcal S\) & \texttt{\detokenize{\mathcal S}} & tagged execution states possibilities transitions safety and goals symbol family member \(\mathcal S\) & P16 execution realization & n/a & \hyperref[def:p16-rem-execution-model]{\ref*{def:p16-rem-execution-model}} & \thepage & All relevant history memory included; repairs may change current modes. Explicit displayed family; member types and common actual first-use exposed, not one atomic glyph. Family listing, not one atomic glyph. Author macro \texttt{\detokenize{X\_i,}\allowbreak\detokenize{calX,}\allowbreak\detokenize{B\_0,}\allowbreak\detokenize{T\_a,}\allowbreak\detokenize{calS,}\allowbreak\detokenize{calG}} not carried to the members. \\
\phantomsection\label{entry:symbol-registry-07962}\(\mathcal G\) & \texttt{\detokenize{\mathcal G}} & tagged execution states possibilities transitions safety and goals symbol family member \(\mathcal G\) & P16 execution realization & n/a & \hyperref[def:p16-rem-execution-model]{\ref*{def:p16-rem-execution-model}} & \thepage & All relevant history memory included; repairs may change current modes. Explicit displayed family; member types and common actual first-use exposed, not one atomic glyph. Family listing, not one atomic glyph. Author macro \texttt{\detokenize{X\_i,}\allowbreak\detokenize{calX,}\allowbreak\detokenize{B\_0,}\allowbreak\detokenize{T\_a,}\allowbreak\detokenize{calS,}\allowbreak\detokenize{calG}} not carried to the members. \\
\phantomsection\label{entry:symbol-registry-07963}\(Post(B,a)\) & \texttt{\detokenize{Post(B,}\allowbreak\detokenize{a)}} & complete successors state count and action count symbol family member \(Post(B,a)\) & P16 strong plan search & n/a & \hyperref[def:p16-rem-strong-plan]{\ref*{def:p16-rem-strong-plan}} & \thepage & No filtering away inapplicable worlds. Explicit displayed family; member types and common actual first-use exposed, not one atomic glyph. Family listing, not one atomic glyph. Author macro \texttt{\detokenize{Post,}\allowbreak\detokenize{n,}\allowbreak\detokenize{m}} not carried to the members. \\
\phantomsection\label{entry:symbol-registry-07964}\(n\) & \texttt{\detokenize{n}} & complete successors state count and action count symbol family member \(n\) & P16 strong plan search & n/a & \hyperref[def:p16-rem-strong-plan]{\ref*{def:p16-rem-strong-plan}} & \thepage & No filtering away inapplicable worlds. Explicit displayed family; member types and common actual first-use exposed, not one atomic glyph. Family listing, not one atomic glyph. Author macro \texttt{\detokenize{Post,}\allowbreak\detokenize{n,}\allowbreak\detokenize{m}} not carried to the members. \\
\phantomsection\label{entry:symbol-registry-07965}\(m\) & \texttt{\detokenize{m}} & complete successors state count and action count symbol family member \(m\) & P16 strong plan search & n/a & \hyperref[def:p16-rem-strong-plan]{\ref*{def:p16-rem-strong-plan}} & \thepage & No filtering away inapplicable worlds. Explicit displayed family; member types and common actual first-use exposed, not one atomic glyph. Family listing, not one atomic glyph. Author macro \texttt{\detokenize{Post,}\allowbreak\detokenize{n,}\allowbreak\detokenize{m}} not carried to the members. \\
\phantomsection\label{entry:symbol-registry-07966}\(x\) & \texttt{\detokenize{x}} & main route and isolated backup switches symbol family member \(x\) & P16 power transfer & n/a & \hyperref[def:p16-rem-power-transfer]{\ref*{def:p16-rem-power-transfer}} & \thepage & Healthy active-source model; backup availability explicit. Explicit displayed family; member types and common actual first-use exposed, not one atomic glyph. Family listing, not one atomic glyph. Author macro \texttt{\detokenize{x,}\allowbreak\detokenize{y,}\allowbreak\detokenize{b}} not carried to the members. \\
\phantomsection\label{entry:symbol-registry-07967}\(y\) & \texttt{\detokenize{y}} & main route and isolated backup switches symbol family member \(y\) & P16 power transfer & n/a & \hyperref[def:p16-rem-power-transfer]{\ref*{def:p16-rem-power-transfer}} & \thepage & Healthy active-source model; backup availability explicit. Explicit displayed family; member types and common actual first-use exposed, not one atomic glyph. Family listing, not one atomic glyph. Author macro \texttt{\detokenize{x,}\allowbreak\detokenize{y,}\allowbreak\detokenize{b}} not carried to the members. \\
\phantomsection\label{entry:symbol-registry-07968}\(b\) & \texttt{\detokenize{b}} & main route and isolated backup switches symbol family member \(b\) & P16 power transfer & n/a & \hyperref[def:p16-rem-power-transfer]{\ref*{def:p16-rem-power-transfer}} & \thepage & Healthy active-source model; backup availability explicit. Explicit displayed family; member types and common actual first-use exposed, not one atomic glyph. Family listing, not one atomic glyph. Author macro \texttt{\detokenize{x,}\allowbreak\detokenize{y,}\allowbreak\detokenize{b}} not carried to the members. \\
\phantomsection\label{entry:symbol-registry-07969}\(h\) & \texttt{\detokenize{h}} & observation map alphabet branches and winning ranks symbol family member \(h\) & P16 observation policy & n/a & \hyperref[def:p16-rem-policy-search]{\ref*{def:p16-rem-policy-search}} & \thepage & Finite-depth strong criterion not probability or fairness. Explicit displayed family; member types and common actual first-use exposed, not one atomic glyph. Family listing, not one atomic glyph. Author macro \texttt{\detokenize{h,}\allowbreak\detokenize{Z,}\allowbreak\detokenize{B\_az,}\allowbreak\detokenize{Win\_j}} not carried to the members. \\
\phantomsection\label{entry:symbol-registry-07970}\(Z\) & \texttt{\detokenize{Z}} & observation map alphabet branches and winning ranks symbol family member \(Z\) & P16 observation policy & n/a & \hyperref[def:p16-rem-policy-search]{\ref*{def:p16-rem-policy-search}} & \thepage & Finite-depth strong criterion not probability or fairness. Explicit displayed family; member types and common actual first-use exposed, not one atomic glyph. Family listing, not one atomic glyph. Author macro \texttt{\detokenize{h,}\allowbreak\detokenize{Z,}\allowbreak\detokenize{B\_az,}\allowbreak\detokenize{Win\_j}} not carried to the members. \\
\phantomsection\label{entry:symbol-registry-07971}\(B_{a,z}\) & \texttt{\detokenize{B_{a,z}}} & observation map alphabet branches and winning ranks symbol family member \(B_{a,z}\) & P16 observation policy & n/a & \hyperref[def:p16-rem-policy-search]{\ref*{def:p16-rem-policy-search}} & \thepage & Finite-depth strong criterion not probability or fairness. Explicit displayed family; member types and common actual first-use exposed, not one atomic glyph. Family listing, not one atomic glyph. Author macro \texttt{\detokenize{h,}\allowbreak\detokenize{Z,}\allowbreak\detokenize{B\_az,}\allowbreak\detokenize{Win\_j}} not carried to the members. \\
\phantomsection\label{entry:symbol-registry-07972}\(Win_j\) & \texttt{\detokenize{Win_j}} & observation map alphabet branches and winning ranks symbol family member \(Win_j\) & P16 observation policy & n/a & \hyperref[def:p16-rem-policy-search]{\ref*{def:p16-rem-policy-search}} & \thepage & Finite-depth strong criterion not probability or fairness. Explicit displayed family; member types and common actual first-use exposed, not one atomic glyph. Family listing, not one atomic glyph. Author macro \texttt{\detokenize{h,}\allowbreak\detokenize{Z,}\allowbreak\detokenize{B\_az,}\allowbreak\detokenize{Win\_j}} not carried to the members. \\
\phantomsection\label{entry:symbol-registry-07973}\(q_0\) & \texttt{\detokenize{q_0}} & iron stocks derivatives threshold supply and rates symbol family member \(q_0\) & P16 therapy & n/a & \hyperref[def:p16-rem-therapy-model]{\ref*{def:p16-rem-therapy-model}} & \thepage & Before/after time indices; derivative recomputed. Explicit displayed family; member types and common actual first-use exposed, not one atomic glyph. Family listing, not one atomic glyph. Author macro \texttt{\detokenize{q0,}\allowbreak\detokenize{q1,}\allowbreak\detokenize{qprime0,}\allowbreak\detokenize{qprime1,}\allowbreak\detokenize{H,}\allowbreak\detokenize{Q0,}\allowbreak\detokenize{Q1,}\allowbreak\detokenize{p,}\allowbreak\detokenize{k}} not carried to the members. \\
\phantomsection\label{entry:symbol-registry-07974}\(q_1\) & \texttt{\detokenize{q_1}} & iron stocks derivatives threshold supply and rates symbol family member \(q_1\) & P16 therapy & n/a & \hyperref[def:p16-rem-therapy-model]{\ref*{def:p16-rem-therapy-model}} & \thepage & Before/after time indices; derivative recomputed. Explicit displayed family; member types and common actual first-use exposed, not one atomic glyph. Family listing, not one atomic glyph. Author macro \texttt{\detokenize{q0,}\allowbreak\detokenize{q1,}\allowbreak\detokenize{qprime0,}\allowbreak\detokenize{qprime1,}\allowbreak\detokenize{H,}\allowbreak\detokenize{Q0,}\allowbreak\detokenize{Q1,}\allowbreak\detokenize{p,}\allowbreak\detokenize{k}} not carried to the members. \\
\phantomsection\label{entry:symbol-registry-07975}\(q'_0\) & \texttt{\detokenize{q'_0}} & iron stocks derivatives threshold supply and rates symbol family member \(q'_0\) & P16 therapy & n/a & \hyperref[def:p16-rem-therapy-model]{\ref*{def:p16-rem-therapy-model}} & \thepage & Before/after time indices; derivative recomputed. Explicit displayed family; member types and common actual first-use exposed, not one atomic glyph. Family listing, not one atomic glyph. Author macro \texttt{\detokenize{q0,}\allowbreak\detokenize{q1,}\allowbreak\detokenize{qprime0,}\allowbreak\detokenize{qprime1,}\allowbreak\detokenize{H,}\allowbreak\detokenize{Q0,}\allowbreak\detokenize{Q1,}\allowbreak\detokenize{p,}\allowbreak\detokenize{k}} not carried to the members. \\
\phantomsection\label{entry:symbol-registry-07976}\(q'_1\) & \texttt{\detokenize{q'_1}} & iron stocks derivatives threshold supply and rates symbol family member \(q'_1\) & P16 therapy & n/a & \hyperref[def:p16-rem-therapy-model]{\ref*{def:p16-rem-therapy-model}} & \thepage & Before/after time indices; derivative recomputed. Explicit displayed family; member types and common actual first-use exposed, not one atomic glyph. Family listing, not one atomic glyph. Author macro \texttt{\detokenize{q0,}\allowbreak\detokenize{q1,}\allowbreak\detokenize{qprime0,}\allowbreak\detokenize{qprime1,}\allowbreak\detokenize{H,}\allowbreak\detokenize{Q0,}\allowbreak\detokenize{Q1,}\allowbreak\detokenize{p,}\allowbreak\detokenize{k}} not carried to the members. \\
\phantomsection\label{entry:symbol-registry-07977}\(H\) & \texttt{\detokenize{H}} & iron stocks derivatives threshold supply and rates symbol family member \(H\) & P16 therapy & n/a & \hyperref[def:p16-rem-therapy-model]{\ref*{def:p16-rem-therapy-model}} & \thepage & Before/after time indices; derivative recomputed. Explicit displayed family; member types and common actual first-use exposed, not one atomic glyph. Family listing, not one atomic glyph. Author macro \texttt{\detokenize{q0,}\allowbreak\detokenize{q1,}\allowbreak\detokenize{qprime0,}\allowbreak\detokenize{qprime1,}\allowbreak\detokenize{H,}\allowbreak\detokenize{Q0,}\allowbreak\detokenize{Q1,}\allowbreak\detokenize{p,}\allowbreak\detokenize{k}} not carried to the members. \\
\phantomsection\label{entry:symbol-registry-07978}\(Q_0\) & \texttt{\detokenize{Q_0}} & iron stocks derivatives threshold supply and rates symbol family member \(Q_0\) & P16 therapy & n/a & \hyperref[def:p16-rem-therapy-model]{\ref*{def:p16-rem-therapy-model}} & \thepage & Before/after time indices; derivative recomputed. Explicit displayed family; member types and common actual first-use exposed, not one atomic glyph. Family listing, not one atomic glyph. Author macro \texttt{\detokenize{q0,}\allowbreak\detokenize{q1,}\allowbreak\detokenize{qprime0,}\allowbreak\detokenize{qprime1,}\allowbreak\detokenize{H,}\allowbreak\detokenize{Q0,}\allowbreak\detokenize{Q1,}\allowbreak\detokenize{p,}\allowbreak\detokenize{k}} not carried to the members. \\
\phantomsection\label{entry:symbol-registry-07979}\(Q_1\) & \texttt{\detokenize{Q_1}} & iron stocks derivatives threshold supply and rates symbol family member \(Q_1\) & P16 therapy & n/a & \hyperref[def:p16-rem-therapy-model]{\ref*{def:p16-rem-therapy-model}} & \thepage & Before/after time indices; derivative recomputed. Explicit displayed family; member types and common actual first-use exposed, not one atomic glyph. Family listing, not one atomic glyph. Author macro \texttt{\detokenize{q0,}\allowbreak\detokenize{q1,}\allowbreak\detokenize{qprime0,}\allowbreak\detokenize{qprime1,}\allowbreak\detokenize{H,}\allowbreak\detokenize{Q0,}\allowbreak\detokenize{Q1,}\allowbreak\detokenize{p,}\allowbreak\detokenize{k}} not carried to the members. \\
\phantomsection\label{entry:symbol-registry-07980}\(p\) & \texttt{\detokenize{p}} & iron stocks derivatives threshold supply and rates symbol family member \(p\) & P16 therapy & n/a & \hyperref[def:p16-rem-therapy-model]{\ref*{def:p16-rem-therapy-model}} & \thepage & Before/after time indices; derivative recomputed. Explicit displayed family; member types and common actual first-use exposed, not one atomic glyph. Family listing, not one atomic glyph. Author macro \texttt{\detokenize{q0,}\allowbreak\detokenize{q1,}\allowbreak\detokenize{qprime0,}\allowbreak\detokenize{qprime1,}\allowbreak\detokenize{H,}\allowbreak\detokenize{Q0,}\allowbreak\detokenize{Q1,}\allowbreak\detokenize{p,}\allowbreak\detokenize{k}} not carried to the members. \\
\phantomsection\label{entry:symbol-registry-07981}\(k\) & \texttt{\detokenize{k}} & iron stocks derivatives threshold supply and rates symbol family member \(k\) & P16 therapy & n/a & \hyperref[def:p16-rem-therapy-model]{\ref*{def:p16-rem-therapy-model}} & \thepage & Before/after time indices; derivative recomputed. Explicit displayed family; member types and common actual first-use exposed, not one atomic glyph. Family listing, not one atomic glyph. Author macro \texttt{\detokenize{q0,}\allowbreak\detokenize{q1,}\allowbreak\detokenize{qprime0,}\allowbreak\detokenize{qprime1,}\allowbreak\detokenize{H,}\allowbreak\detokenize{Q0,}\allowbreak\detokenize{Q1,}\allowbreak\detokenize{p,}\allowbreak\detokenize{k}} not carried to the members. \\
\phantomsection\label{entry:symbol-registry-07982}\(p_{min}\) & \texttt{\detokenize{p_{min}}} & independent nonnegative introduction/removal bounds symbol family member \(p_{min}\) & P16 therapy calculation & n/a & \hyperref[thm:p16-rem-therapy]{\ref*{thm:p16-rem-therapy}} & \thepage & Corners admitted; complete influence set required. Explicit displayed family; member types and common actual first-use exposed, not one atomic glyph. Family listing, not one atomic glyph. Author macro \texttt{\detokenize{pmin,}\allowbreak\detokenize{pmax,}\allowbreak\detokenize{kmin,}\allowbreak\detokenize{kmax}} not carried to the members. \\
\phantomsection\label{entry:symbol-registry-07983}\(p_{max}\) & \texttt{\detokenize{p_{max}}} & independent nonnegative introduction/removal bounds symbol family member \(p_{max}\) & P16 therapy calculation & n/a & \hyperref[thm:p16-rem-therapy]{\ref*{thm:p16-rem-therapy}} & \thepage & Corners admitted; complete influence set required. Explicit displayed family; member types and common actual first-use exposed, not one atomic glyph. Family listing, not one atomic glyph. Author macro \texttt{\detokenize{pmin,}\allowbreak\detokenize{pmax,}\allowbreak\detokenize{kmin,}\allowbreak\detokenize{kmax}} not carried to the members. \\
\phantomsection\label{entry:symbol-registry-07984}\(k_{min}\) & \texttt{\detokenize{k_{min}}} & independent nonnegative introduction/removal bounds symbol family member \(k_{min}\) & P16 therapy calculation & n/a & \hyperref[thm:p16-rem-therapy]{\ref*{thm:p16-rem-therapy}} & \thepage & Corners admitted; complete influence set required. Explicit displayed family; member types and common actual first-use exposed, not one atomic glyph. Family listing, not one atomic glyph. Author macro \texttt{\detokenize{pmin,}\allowbreak\detokenize{pmax,}\allowbreak\detokenize{kmin,}\allowbreak\detokenize{kmax}} not carried to the members. \\
\phantomsection\label{entry:symbol-registry-07985}\(k_{max}\) & \texttt{\detokenize{k_{max}}} & independent nonnegative introduction/removal bounds symbol family member \(k_{max}\) & P16 therapy calculation & n/a & \hyperref[thm:p16-rem-therapy]{\ref*{thm:p16-rem-therapy}} & \thepage & Corners admitted; complete influence set required. Explicit displayed family; member types and common actual first-use exposed, not one atomic glyph. Family listing, not one atomic glyph. Author macro \texttt{\detokenize{pmin,}\allowbreak\detokenize{pmax,}\allowbreak\detokenize{kmin,}\allowbreak\detokenize{kmax}} not carried to the members. \\
\phantomsection\label{entry:symbol-registry-07986}\(DEC\) & \texttt{\detokenize{DEC}} & native action triggers persistent situation and therapy goals symbol family member \(DEC\) & P16 process therapy & n/a & \hyperref[sec:p16-rem-therapy]{\ref*{sec:p16-rem-therapy}} & \thepage & Location and transport not implicit. Explicit displayed family; member types and common actual first-use exposed, not one atomic glyph. Family listing, not one atomic glyph. Author macro \texttt{\detokenize{DEC,}\allowbreak\detokenize{SIT\_P,}\allowbreak\detokenize{GOALS\_T}} not carried to the members. \\
\phantomsection\label{entry:symbol-registry-07987}\(SIT_P\) & \texttt{\detokenize{SIT_P}} & native action triggers persistent situation and therapy goals symbol family member \(SIT_P\) & P16 process therapy & n/a & \hyperref[sec:p16-rem-therapy]{\ref*{sec:p16-rem-therapy}} & \thepage & Location and transport not implicit. Explicit displayed family; member types and common actual first-use exposed, not one atomic glyph. Family listing, not one atomic glyph. Author macro \texttt{\detokenize{DEC,}\allowbreak\detokenize{SIT\_P,}\allowbreak\detokenize{GOALS\_T}} not carried to the members. \\
\phantomsection\label{entry:symbol-registry-07988}\(GOALS_T\) & \texttt{\detokenize{GOALS_T}} & native action triggers persistent situation and therapy goals symbol family member \(GOALS_T\) & P16 process therapy & n/a & \hyperref[sec:p16-rem-therapy]{\ref*{sec:p16-rem-therapy}} & \thepage & Location and transport not implicit. Explicit displayed family; member types and common actual first-use exposed, not one atomic glyph. Family listing, not one atomic glyph. Author macro \texttt{\detokenize{DEC,}\allowbreak\detokenize{SIT\_P,}\allowbreak\detokenize{GOALS\_T}} not carried to the members. \\
\phantomsection\label{entry:symbol-registry-07989}\(CONF\) & \texttt{\detokenize{CONF}} & native configuration choices domain constraints and goals symbol family member \(CONF\) & P16 design configuration & n/a & \hyperref[sec:p16-rem-configuration]{\ref*{sec:p16-rem-configuration}} & \thepage & Structure is generated result not fixed diagnosis input. Explicit displayed family; member types and common actual first-use exposed, not one atomic glyph. Family listing, not one atomic glyph. Author macro \texttt{\detokenize{CONF,}\allowbreak\detokenize{PARS,}\allowbreak\detokenize{DD,}\allowbreak\detokenize{CONFGOALS}} not carried to the members. \\
\phantomsection\label{entry:symbol-registry-07990}\(PARS\) & \texttt{\detokenize{PARS}} & native configuration choices domain constraints and goals symbol family member \(PARS\) & P16 design configuration & n/a & \hyperref[sec:p16-rem-configuration]{\ref*{sec:p16-rem-configuration}} & \thepage & Structure is generated result not fixed diagnosis input. Explicit displayed family; member types and common actual first-use exposed, not one atomic glyph. Family listing, not one atomic glyph. Author macro \texttt{\detokenize{CONF,}\allowbreak\detokenize{PARS,}\allowbreak\detokenize{DD,}\allowbreak\detokenize{CONFGOALS}} not carried to the members. \\
\phantomsection\label{entry:symbol-registry-07991}\(DD\) & \texttt{\detokenize{DD}} & native configuration choices domain constraints and goals symbol family member \(DD\) & P16 design configuration & n/a & \hyperref[sec:p16-rem-configuration]{\ref*{sec:p16-rem-configuration}} & \thepage & Structure is generated result not fixed diagnosis input. Explicit displayed family; member types and common actual first-use exposed, not one atomic glyph. Family listing, not one atomic glyph. Author macro \texttt{\detokenize{CONF,}\allowbreak\detokenize{PARS,}\allowbreak\detokenize{DD,}\allowbreak\detokenize{CONFGOALS}} not carried to the members. \\
\phantomsection\label{entry:symbol-registry-07992}\(CONFGOALS\) & \texttt{\detokenize{CONFGOALS}} & native configuration choices domain constraints and goals symbol family member \(CONFGOALS\) & P16 design configuration & n/a & \hyperref[sec:p16-rem-configuration]{\ref*{sec:p16-rem-configuration}} & \thepage & Structure is generated result not fixed diagnosis input. Explicit displayed family; member types and common actual first-use exposed, not one atomic glyph. Family listing, not one atomic glyph. Author macro \texttt{\detokenize{CONF,}\allowbreak\detokenize{PARS,}\allowbreak\detokenize{DD,}\allowbreak\detokenize{CONFGOALS}} not carried to the members. \\
\phantomsection\label{entry:symbol-registry-07993}\(n_A\) & \texttt{\detokenize{n_A}} & counts of declared card types symbol family member \(n_A\) & P16 catalog case & n/a & \hyperref[def:p16-rem-card-configuration]{\ref*{def:p16-rem-card-configuration}} & \thepage & Stipulated units not industrial telephone specifications. Explicit displayed family; member types and common actual first-use exposed, not one atomic glyph. Family listing, not one atomic glyph. Author macro \texttt{\detokenize{n\_A,}\allowbreak\detokenize{n\_B}} not carried to the members. \\
\phantomsection\label{entry:symbol-registry-07994}\(n_B\) & \texttt{\detokenize{n_B}} & counts of declared card types symbol family member \(n_B\) & P16 catalog case & n/a & \hyperref[def:p16-rem-card-configuration]{\ref*{def:p16-rem-card-configuration}} & \thepage & Stipulated units not industrial telephone specifications. Explicit displayed family; member types and common actual first-use exposed, not one atomic glyph. Family listing, not one atomic glyph. Author macro \texttt{\detokenize{n\_A,}\allowbreak\detokenize{n\_B}} not carried to the members. \\
\phantomsection\label{entry:symbol-registry-07995}\(\Gamma_{F,SCEN}\) & \texttt{\detokenize{\Gamma_{F,SCEN}}} & fault/scenario model effect and scenario symbol family member \(\Gamma_{F,SCEN}\) & P16 FMEA & n/a & \hyperref[sec:p16-rem-fmea]{\ref*{sec:p16-rem-fmea}} & \thepage & Nonempty model required before effect classification. Explicit displayed family; member types and common actual first-use exposed, not one atomic glyph. Family listing, not one atomic glyph. Author macro \texttt{\detokenize{Gamma,}\allowbreak\detokenize{E,}\allowbreak\detokenize{M\_F,}\allowbreak\detokenize{SCEN}} not carried to the members. \\
\phantomsection\label{entry:symbol-registry-07996}\(E\) & \texttt{\detokenize{E}} & fault/scenario model effect and scenario symbol family member \(E\) & P16 FMEA & n/a & \hyperref[sec:p16-rem-fmea]{\ref*{sec:p16-rem-fmea}} & \thepage & Nonempty model required before effect classification. Explicit displayed family; member types and common actual first-use exposed, not one atomic glyph. Family listing, not one atomic glyph. Author macro \texttt{\detokenize{Gamma,}\allowbreak\detokenize{E,}\allowbreak\detokenize{M\_F,}\allowbreak\detokenize{SCEN}} not carried to the members. \\
\phantomsection\label{entry:symbol-registry-07997}\(M_F\) & \texttt{\detokenize{M_F}} & fault/scenario model effect and scenario symbol family member \(M_F\) & P16 FMEA & n/a & \hyperref[sec:p16-rem-fmea]{\ref*{sec:p16-rem-fmea}} & \thepage & Nonempty model required before effect classification. Explicit displayed family; member types and common actual first-use exposed, not one atomic glyph. Family listing, not one atomic glyph. Author macro \texttt{\detokenize{Gamma,}\allowbreak\detokenize{E,}\allowbreak\detokenize{M\_F,}\allowbreak\detokenize{SCEN}} not carried to the members. \\
\phantomsection\label{entry:symbol-registry-07998}\(SCEN\) & \texttt{\detokenize{SCEN}} & fault/scenario model effect and scenario symbol family member \(SCEN\) & P16 FMEA & n/a & \hyperref[sec:p16-rem-fmea]{\ref*{sec:p16-rem-fmea}} & \thepage & Nonempty model required before effect classification. Explicit displayed family; member types and common actual first-use exposed, not one atomic glyph. Family listing, not one atomic glyph. Author macro \texttt{\detokenize{Gamma,}\allowbreak\detokenize{E,}\allowbreak\detokenize{M\_F,}\allowbreak\detokenize{SCEN}} not carried to the members. \\
\phantomsection\label{entry:symbol-registry-07999}\(q_1\) & \texttt{\detokenize{q_1}} & route outputs and total service symbol family member \(q_1\) & P16 redundancy masking & n/a & \hyperref[def:p16-rem-redundancy]{\ref*{def:p16-rem-redundancy}} & \thepage & q here output not concentration; isolated open faults. Explicit displayed family; member types and common actual first-use exposed, not one atomic glyph. Family listing, not one atomic glyph. Author macro \texttt{\detokenize{q\_1,}\allowbreak\detokenize{q\_2,}\allowbreak\detokenize{s}} not carried to the members. \\
\phantomsection\label{entry:symbol-registry-08000}\(q_2\) & \texttt{\detokenize{q_2}} & route outputs and total service symbol family member \(q_2\) & P16 redundancy masking & n/a & \hyperref[def:p16-rem-redundancy]{\ref*{def:p16-rem-redundancy}} & \thepage & q here output not concentration; isolated open faults. Explicit displayed family; member types and common actual first-use exposed, not one atomic glyph. Family listing, not one atomic glyph. Author macro \texttt{\detokenize{q\_1,}\allowbreak\detokenize{q\_2,}\allowbreak\detokenize{s}} not carried to the members. \\
\phantomsection\label{entry:symbol-registry-08001}\(s\) & \texttt{\detokenize{s}} & route outputs and total service symbol family member \(s\) & P16 redundancy masking & n/a & \hyperref[def:p16-rem-redundancy]{\ref*{def:p16-rem-redundancy}} & \thepage & q here output not concentration; isolated open faults. Explicit displayed family; member types and common actual first-use exposed, not one atomic glyph. Family listing, not one atomic glyph. Author macro \texttt{\detokenize{q\_1,}\allowbreak\detokenize{q\_2,}\allowbreak\detokenize{s}} not carried to the members. \\
\phantomsection\label{entry:symbol-registry-08002}\(x\) & \texttt{\detokenize{x}} & integer input output repair increment and overfitting function symbol family member \(x\) & P16 debugging & n/a & \hyperref[def:p16-rem-debug-model]{\ref*{def:p16-rem-debug-model}} & \thepage & Old-version output not a post-revision frame. Explicit displayed family; member types and common actual first-use exposed, not one atomic glyph. Family listing, not one atomic glyph. Author macro \texttt{\detokenize{x,}\allowbreak\detokenize{y,}\allowbreak\detokenize{delta,}\allowbreak\detokenize{f}} not carried to the members. \\
\phantomsection\label{entry:symbol-registry-08003}\(y\) & \texttt{\detokenize{y}} & integer input output repair increment and overfitting function symbol family member \(y\) & P16 debugging & n/a & \hyperref[def:p16-rem-debug-model]{\ref*{def:p16-rem-debug-model}} & \thepage & Old-version output not a post-revision frame. Explicit displayed family; member types and common actual first-use exposed, not one atomic glyph. Family listing, not one atomic glyph. Author macro \texttt{\detokenize{x,}\allowbreak\detokenize{y,}\allowbreak\detokenize{delta,}\allowbreak\detokenize{f}} not carried to the members. \\
\phantomsection\label{entry:symbol-registry-08004}\(\delta\) & \texttt{\detokenize{\delta}} & integer input output repair increment and overfitting function symbol family member \(\delta\) & P16 debugging & n/a & \hyperref[def:p16-rem-debug-model]{\ref*{def:p16-rem-debug-model}} & \thepage & Old-version output not a post-revision frame. Explicit displayed family; member types and common actual first-use exposed, not one atomic glyph. Family listing, not one atomic glyph. Author macro \texttt{\detokenize{x,}\allowbreak\detokenize{y,}\allowbreak\detokenize{delta,}\allowbreak\detokenize{f}} not carried to the members. \\
\phantomsection\label{entry:symbol-registry-08005}\(f\) & \texttt{\detokenize{f}} & integer input output repair increment and overfitting function symbol family member \(f\) & P16 debugging & n/a & \hyperref[def:p16-rem-debug-model]{\ref*{def:p16-rem-debug-model}} & \thepage & Old-version output not a post-revision frame. Explicit displayed family; member types and common actual first-use exposed, not one atomic glyph. Family listing, not one atomic glyph. Author macro \texttt{\detokenize{x,}\allowbreak\detokenize{y,}\allowbreak\detokenize{delta,}\allowbreak\detokenize{f}} not carried to the members. \\
\end{longtable}
\endgroup
\Needspace{.4\textheight}
\subsection*{Stable part P17 scopes}
\begingroup\scriptsize\setlength{\tabcolsep}{3pt}
\begin{longtable}{@{}p{0.9in}p{1.3in}p{1.0in}p{0.5in}p{0.3in}p{0.4in}p{0.35in}p{1.65in}@{}}
\toprule
\textbf{Glyph} & \textbf{Macro} & \textbf{Role} & \textbf{Scope} & \textbf{Prec./form} & \textbf{First} & \textbf{Reg. page} & \textbf{Principal uses / notes} \\
\midrule
\endfirsthead
\toprule
\textbf{Glyph} & \textbf{Macro} & \textbf{Role} & \textbf{Scope} & \textbf{Prec./form} & \textbf{First} & \textbf{Reg. page} & \textbf{Principal uses / notes} \\
\midrule
\endhead
\midrule\multicolumn{8}{r}{\emph{continued}}\\\endfoot
\bottomrule\endlastfoot
\phantomsection\label{entry:symbol-registry-00983}\(X,x,y,z,R\) & \texttt{\detokenize{X,}\allowbreak\detokenize{x,}\allowbreak\detokenize{y,}\allowbreak\detokenize{z,}\allowbreak\detokenize{R}} & finite alternative set members and comparison relation & P17 finite-preferences & meta & \hyperref[sec:p17-finite-preferences]{\ref*{sec:p17-finite-preferences}} & \thepage & X may be empty. R is a binary relation on X; completeness and transitivity are explicit hypotheses when the finite preference theorems use them. \\
\phantomsection\label{entry:symbol-registry-00984}\(\succeq,\succ,\sim\) & \texttt{\detokenize{\succeq,}\allowbreak\detokenize{\succ,}\allowbreak\detokenize{\sim}} & weak preference strict preference and indifference & P17 finite-preferences & meta & \hyperref[sec:p17-finite-preferences]{\ref*{sec:p17-finite-preferences}} & \thepage & Completeness includes equal alternatives and implies reflexivity. Strict preference is one affirmative weak comparison and failure of its reverse; indifference requires both comparisons and differs from missing information. \\
\phantomsection\label{entry:symbol-registry-00985}\([x],\mathcal Q,X/{\sim}\) & \texttt{\detokenize{[x],}\allowbreak\detokenize{\mathcal Q,}\allowbreak\detokenize{X/}\allowbreak\detokenize{{\sim}}} & indifference class class family and preference quotient & P17 finite-preferences & meta & \hyperref[sec:p17-finite-preferences]{\ref*{sec:p17-finite-preferences}} & \thepage & Equivalence of indifference and the partition property are proved before the quotient is used. The class family is empty when X is empty. \\
\phantomsection\label{entry:symbol-registry-00986}\(A\unrhd B\) & \texttt{\detokenize{\unrhd}} & total comparison of indifference classes & P17 finite-preferences & meta & \hyperref[sec:p17-finite-preferences]{\ref*{sec:p17-finite-preferences}} & \thepage & The defined existential representative comparison is proved to agree for all representatives. A and B here denote nonempty classes, with a different local role from the later menu names. \\
\phantomsection\label{entry:symbol-registry-00987}\(u,v:X\to\mathbb{R}\) & \texttt{\detokenize{u,}\allowbreak\detokenize{v}} & real-valued ordinal representations & P17 finite-preferences & meta & \hyperref[sec:p17-finite-preferences]{\ref*{sec:p17-finite-preferences}} & \thepage & A representation preserves and reflects the complete weak comparison through greater-than-or-equal, including indifference as equality. It does not specify preference intensity or lottery values. \\
\phantomsection\label{entry:symbol-registry-00988}\(k,I_k,r,u_0\) & \texttt{\detokenize{k,}\allowbreak\detokenize{I_k,}\allowbreak\detokenize{r,}\allowbreak\detokenize{u_0}} & finite quotient size class ranks and canonical representation & P17 finite-preferences & meta & \hyperref[sec:p17-finite-preferences]{\ref*{sec:p17-finite-preferences}} & \thepage & I\_k contains the natural indices below k. Rank counts strictly lower classes; u0 assigns that rank to an alternative's class. The empty rank function represents the empty relation. \\
\phantomsection\label{entry:symbol-registry-00989}\(D,f,F,u(X)\) & \texttt{\detokenize{D,}\allowbreak\detokenize{f,}\allowbreak\detokenize{F,}\allowbreak\detokenize{u(X)}} & attained-value domain increasing scale map extension and image & P17 finite-preferences & meta & \hyperref[sec:p17-finite-preferences]{\ref*{sec:p17-finite-preferences}} & \thepage & The increasing bijection between two attained finite value sets is unique. A strictly increasing extension F to the whole real line exists but is not unique outside the attained values. \\
\phantomsection\label{entry:symbol-registry-00990}\(s_i,t_i\) & \texttt{\detokenize{s_i,}\allowbreak\detokenize{t_i}} & ordered attained values and their images in a scale extension & P17 finite-preferences & meta & \hyperref[sec:p17-finite-preferences]{\ref*{sec:p17-finite-preferences}} & \thepage & These finite increasing lists govern the displayed positive-slope interpolation and its exterior rays. Their local indices enumerate attained classes, not alternatives. \\
\phantomsection\label{entry:symbol-registry-00991}\(A,B,C_R(B)\) & \texttt{\detokenize{A,}\allowbreak\detokenize{B,}\allowbreak\detokenize{C_R}} & finite menus and their weak-maximum choice operator & P17 finite-preferences & meta & \hyperref[sec:p17-finite-preferences]{\ref*{sec:p17-finite-preferences}} & \thepage & Menus are subsets of X. The choice set contains the alternatives weakly above every member of the menu; it is nonempty for a nonempty finite menu under weak-preference assumptions and empty on the empty menu. \\
\phantomsection\label{entry:symbol-registry-00992}\(\succeq_{\rm tie},t,v\) & \texttt{\detokenize{\succeq_{\rm tie},}\allowbreak\detokenize{t,}\allowbreak\detokenize{v}} & complete tied preference example and its two score tables & P17 finite-preferences & meta & \hyperref[sec:p17-finite-preferences]{\ref*{sec:p17-finite-preferences}} & \thepage & The three indifference classes have scores 0,1,2 under t and 0,1,10 under v. Both represent the same ordering and all original menu choices; the chapter calculates their different gaps and averages. \\
\phantomsection\label{entry:symbol-registry-00993}\(X_{\rm cyc},R_{\rm cyc}\) & \texttt{\detokenize{X_{\rm cyc},}\allowbreak\detokenize{R_{\rm cyc}}} & complete three-alternative cyclic comparison & P17 finite-preferences & meta & \hyperref[sec:p17-finite-preferences]{\ref*{sec:p17-finite-preferences}} & \thepage & The six listed pairs include the diagonal and one directed three-cycle. Completeness and reflexivity hold, transitivity fails, and no real-valued ordinal representation exists. \\
\phantomsection\label{entry:symbol-registry-00994}\(\alpha t+\beta\) & \texttt{\detokenize{\alpha,}\allowbreak\detokenize{\beta}} & affine scale-change parameters in the ordinal boundary example & P17 finite-preferences & meta & \hyperref[sec:p17-finite-preferences]{\ref*{sec:p17-finite-preferences}} & \thepage & Positive affine changes are a narrower class than all increasing changes permitted by ordinal representation. The supplied two score tables are not related by any affine change on the original alternatives. \\
\phantomsection\label{entry:symbol-registry-01076}\(Z,\Delta(Z),\delta_z\) & \texttt{\detokenize{Z,}\allowbreak\detokenize{\Delta(Z),}\allowbreak\detokenize{\delta_z}} & finite prize set lottery simplex and certain lottery & P17 finite-expected-utility & meta & \hyperref[sec:p17-finite-expected-utility]{\ref*{sec:p17-finite-expected-utility}} & \thepage & The prize set is nonempty and finite. The simplex usually has infinitely many lotteries; the certain lottery places mass one at its indicated prize. \\
\phantomsection\label{entry:symbol-registry-01077}\(p,q,r,\alpha,\alpha_k,\beta_{k\ell}\) & \texttt{\detokenize{p,}\allowbreak\detokenize{q,}\allowbreak\detokenize{r,}\allowbreak\detokenize{\alpha,}\allowbreak\detokenize{\alpha_k,}\allowbreak\detokenize{\beta_{k\ell}}} & lottery vectors and finite mixture weights & P17 finite-expected-utility & meta & \hyperref[sec:p17-finite-expected-utility]{\ref*{sec:p17-finite-expected-utility}} & \thepage & Mixtures are coordinatewise probability sums. Each displayed weight family is nonnegative and normalized; zero-weight branches are allowed and do not supply a conditional ratio. \\
\phantomsection\label{entry:symbol-registry-01078}\(\succeq,\succ,\sim\) & \texttt{\detokenize{\succeq,}\allowbreak\detokenize{\succ,}\allowbreak\detokenize{\sim}} & weak strict and indifferent lottery comparisons & P17 finite-expected-utility & meta & \hyperref[sec:p17-finite-expected-utility]{\ref*{sec:p17-finite-expected-utility}} & \thepage & Completeness and transitivity apply on the whole lottery simplex. The finite-prize ordering alone does not establish a representation of its infinite lottery domain. \\
\phantomsection\label{entry:symbol-registry-01079}\(d(p,q)\) & \texttt{\detokenize{d}} & relative Euclidean distance between lottery vectors & P17 finite-expected-utility & meta & \hyperref[sec:p17-finite-expected-utility]{\ref*{sec:p17-finite-expected-utility}} & \thepage & The continuity axiom quantifies only over nearby points that remain in the simplex, including neighborhoods of certain lotteries at its boundary. \\
\phantomsection\label{entry:symbol-registry-01080}\(B_t,A,t_*\) & \texttt{\detokenize{B_t,}\allowbreak\detokenize{A,}\allowbreak\detokenize{t_*}} & binary lottery segment calibration set and its supremum & P17 finite-expected-utility & meta & \hyperref[sec:p17-finite-expected-utility]{\ref*{sec:p17-finite-expected-utility}} & \thepage & The segment joins a strictly better and a worse lottery. The proof obtains the unique indifference parameter from a nonempty closed subset of the unit interval, with endpoint cases explicit. \\
\phantomsection\label{entry:symbol-registry-01081}\(u,U_u,\mathbb E_p[u]\) & \texttt{\detokenize{u,}\allowbreak\detokenize{U_u,}\allowbreak\detokenize{\mathbb E_p[u]}} & consequence utility and its expected lottery score & P17 finite-expected-utility & meta & \hyperref[sec:p17-finite-expected-utility]{\ref*{sec:p17-finite-expected-utility}} & \thepage & Consequence values and completed lottery scores have different roles. Strictly increasing transformations of the latter preserve order; averaging transformed consequence values generally changes it. \\
\phantomsection\label{entry:symbol-registry-01082}\(M,m,z^+,z^-\) & \texttt{\detokenize{M,}\allowbreak\detokenize{m,}\allowbreak\detokenize{z^+,}\allowbreak\detokenize{z^-}} & best and worst certain lotteries and their prize labels & P17 finite-expected-utility & meta & \hyperref[sec:p17-finite-expected-utility]{\ref*{sec:p17-finite-expected-utility}} & \thepage & Finite weak preference supplies the extrema. If all certain prizes tie, all lotteries tie; otherwise calibration normalizes the extreme consequence utilities to one and zero. \\
\phantomsection\label{entry:symbol-registry-01083}\(a,b,c,\mathcal B\) & \texttt{\detokenize{a,}\allowbreak\detokenize{b,}\allowbreak\detokenize{c,}\allowbreak\detokenize{\mathcal B}} & locally typed lottery endpoints prize labels affine constants and menu & P17 finite-expected-utility & meta & \hyperref[sec:p17-finite-expected-utility]{\ref*{sec:p17-finite-expected-utility}} & \thepage & The mixture lemmas use a and b as lotteries; the counterpreferences use a, b and c as prize labels; affine uniqueness uses real constants with positive multiplier a. The separately declared \(\mathcal B\) is a finite lottery menu, with ties retained and no choice on the empty menu. \\
\phantomsection\label{entry:symbol-registry-01084}\(L_1,L_2,L_3,L_4,V,q_\varepsilon\) & \texttt{\detokenize{L_1,}\allowbreak\detokenize{L_2,}\allowbreak\detokenize{L_3,}\allowbreak\detokenize{L_4,}\allowbreak\detokenize{V,}\allowbreak\detokenize{q_\varepsilon}} & common-consequence lotteries and counterpreference witnesses & P17 finite-expected-utility & meta & \hyperref[sec:p17-finite-expected-utility]{\ref*{sec:p17-finite-expected-utility}} & \thepage & The four source lotteries have exactly related mixture differences. V is the largest-coordinate score in its separate two-prize example; the perturbed three-prize vector witnesses failed lexicographic continuity. \\
\phantomsection\label{entry:symbol-registry-01085}\(I,\Delta_f(I),\operatorname{supp}(p),\delta_x\) & \texttt{\detokenize{I,}\allowbreak\detokenize{\Delta_f,}\allowbreak\detokenize{\operatorname{supp},}\allowbreak\detokenize{\delta_x}} & monetary interval finite-support lottery domain support and certain consequence & P17 risk-aversion & meta & \hyperref[sec:p17-risk-aversion]{\ref*{sec:p17-risk-aversion}} & \thepage & An interval contains the finite means used as certain alternatives. The monetary utility is supplied on the whole interval; its existence is not inferred solely from finite-prize representation. \\
\phantomsection\label{entry:symbol-registry-01086}\(\mu(p),U_u(p),\succeq_u\) & \texttt{\detokenize{\mu,}\allowbreak\detokenize{U_u,}\allowbreak\detokenize{\succeq_u}} & monetary mean expected utility and represented preference & P17 risk-aversion & meta & \hyperref[sec:p17-risk-aversion]{\ref*{sec:p17-risk-aversion}} & \thepage & Both sums are finite. The mean is in monetary units, the expected utility uses the stipulated utility scale, and neither supplies a population estimate. \\
\phantomsection\label{entry:symbol-registry-01087}\(\succeq_{\rm FSD},F_p(t),G_p(t)\) & \texttt{\detokenize{\succeq_{\rm FSD},}\allowbreak\detokenize{F_p,}\allowbreak\detokenize{G_p}} & first-order stochastic dominance cumulative and weak upper tails & P17 risk-aversion & meta & \hyperref[sec:p17-risk-aversion]{\ref*{sec:p17-risk-aversion}} & \thepage & The lower tail includes its endpoint and the upper tail includes its endpoint. Their values at the same atom need not sum to one; the finite-support proof uses the correct adjacent-threshold identity. \\
\phantomsection\label{entry:symbol-registry-01088}\(\operatorname{CE}_u(p),\pi_u(p)\) & \texttt{\detokenize{\operatorname{CE}_u,}\allowbreak\detokenize{\pi_u}} & certainty equivalent and monetary risk premium & P17 risk-aversion & meta & \hyperref[sec:p17-risk-aversion]{\ref*{sec:p17-risk-aversion}} & \thepage & Continuity and strict increase give a unique certainty equivalent within the finite support's extreme values. The premium subtracts it from the same lottery mean. \\
\phantomsection\label{entry:symbol-registry-01089}\(u_1,u_2,\phi=u_1\circ u_2^{-1}\) & \texttt{\detokenize{u_1,}\allowbreak\detokenize{u_2,}\allowbreak\detokenize{\phi,}\allowbreak\detokenize{u_2^{-1}}} & two monetary utilities and their relative utility transformation & P17 risk-aversion & meta & \hyperref[sec:p17-risk-aversion]{\ref*{sec:p17-risk-aversion}} & \thepage & Both utilities share the same monetary interval and are continuous and strictly increasing. Concavity of phi on the attained utility interval is equivalent to the full certainty-equivalent ordering. \\
\phantomsection\label{entry:symbol-registry-01090}\(r_u(x)\) & \texttt{\detokenize{r_u}} & coefficient of absolute risk aversion & P17 risk-aversion & meta & \hyperref[sec:p17-risk-aversion]{\ref*{sec:p17-risk-aversion}} & \thepage & Defined as \(-u''(x)/u'(x)\) on the stated open interval with \(u'(x)>0\). A pointwise coefficient comparison must hold throughout the common domain to establish the global comparative relation. \\
\phantomsection\label{entry:symbol-registry-01091}\(\theta,\psi_t(x)\) & \texttt{\detokenize{\theta,}\allowbreak\detokenize{\psi_t}} & fixed binary probability and local indifference curve & P17 risk-aversion & meta & \hyperref[sec:p17-risk-aversion]{\ref*{sec:p17-risk-aversion}} & \thepage & The probability lies strictly between zero and one. The curve is defined near the certain pair by the inverse-utility formula; its two derivatives have the declared local meaning. \\
\phantomsection\label{entry:symbol-registry-01092}\(T_wp,w\) & \texttt{\detokenize{T_w,}\allowbreak\detokenize{w}} & translation of a monetary lottery and its wealth shift & P17 risk-aversion & meta & \hyperref[sec:p17-risk-aversion]{\ref*{sec:p17-risk-aversion}} & \thepage & The wealth-invariance theorem uses the complete real domain so every shift is admissible. Translating a lottery preserves probabilities and adds w to every prize. \\
\phantomsection\label{entry:symbol-registry-01093}\(A(w),h(w),k\) & \texttt{\detokenize{A,}\allowbreak\detokenize{h,}\allowbreak\detokenize{k}} & positive translation scale its logarithm and additive coefficient & P17 risk-aversion & meta & \hyperref[sec:p17-risk-aversion]{\ref*{sec:p17-risk-aversion}} & \thepage & The continuous wealth-invariance proof obtains a multiplicative positive A and an additive continuous logarithm. Rational approximation determines one coefficient on all real shifts, including negative shifts. \\
\phantomsection\label{entry:symbol-registry-01094}\(\Delta,D_n,q,r\) & \texttt{\detokenize{\Delta,}\allowbreak\detokenize{D_n,}\allowbreak\detokenize{q,}\allowbreak\detokenize{r}} & grid spacing increments indifference probability and geometric ratio & P17 risk-aversion & meta & \hyperref[sec:p17-risk-aversion]{\ref*{sec:p17-risk-aversion}} & \thepage & The positive spacing defines integer wealth grids. The ratio \((1-q)/q\) is positive and at most one, and the increment recurrence holds for negative as well as positive integer positions. \\
\phantomsection\label{entry:symbol-registry-01095}\(\rho,\cosh t,d\) & \texttt{\detokenize{\rho,}\allowbreak\detokenize{\cosh,}\allowbreak\detokenize{d}} & exponential risk coefficient hyperbolic cosine and symmetric risk size & P17 risk-aversion & meta & \hyperref[sec:p17-risk-aversion]{\ref*{sec:p17-risk-aversion}} & \thepage & The coefficient is positive and risk size nonnegative. The exact premium is \(\log\cosh(\rho d)/\rho\), with the zero-risk case and wealth independence stated explicitly. \\
\phantomsection\label{entry:symbol-registry-01117}\(S,K,x,y,z,\mathbf0,\mathbf1,e_i\) & \texttt{\detokenize{S,}\allowbreak\detokenize{K,}\allowbreak\detokenize{x,}\allowbreak\detokenize{y,}\allowbreak\detokenize{z,}\allowbreak\detokenize{\mathbf0,}\allowbreak\detokenize{\mathbf1,}\allowbreak\detokenize{e_i}} & finite state set payoff vectors certain vectors and state basis bets & P17 finite-subjective-probability & meta & \hyperref[sec:p17-finite-subjective-probability]{\ref*{sec:p17-finite-subjective-probability}} & \thepage & K is positive and every payoff vector is real. State-contingent additions may be arbitrary common vectors, not only sure cash shifts. \\
\phantomsection\label{entry:symbol-registry-01118}\(\geq_{\rm cw},d,\|\cdot\|_\infty\) & \texttt{\detokenize{\geq_{\rm cw},}\allowbreak\detokenize{d,}\allowbreak\detokenize{\|\cdot\|_\infty}} & coordinatewise payoff order Euclidean distance and maximum norm & P17 finite-subjective-probability & meta & \hyperref[sec:p17-finite-subjective-probability]{\ref*{sec:p17-finite-subjective-probability}} & \thepage & The source monotonicity requires strict improvement in all coordinates. Weak coordinatewise monotonicity is derived using continuity; the cash-equivalent maximum-norm bound is proved. \\
\phantomsection\label{entry:symbol-registry-01119}\(\succeq,\succ,\sim,C(x)\) & \texttt{\detokenize{\succeq,}\allowbreak\detokenize{\succ,}\allowbreak\detokenize{\sim,}\allowbreak\detokenize{C}} & bet comparisons and uniquely determined cash equivalent & P17 finite-subjective-probability & meta & \hyperref[sec:p17-finite-subjective-probability]{\ref*{sec:p17-finite-subjective-probability}} & \thepage & The cash-equivalent proof establishes additivity and real homogeneity, including negative scalars, before deriving the probability representation. \\
\phantomsection\label{entry:symbol-registry-01120}\(\pi,\pi\cdot x,\mathbb P_\pi,\mathbf1_A\) & \texttt{\detokenize{\pi,}\allowbreak\detokenize{\pi\cdot x,}\allowbreak\detokenize{\mathbb P_\pi,}\allowbreak\detokenize{\mathbf1_A}} & subjective probability vector payoff evaluation event measure and event bet & P17 finite-subjective-probability & meta & \hyperref[sec:p17-finite-subjective-probability]{\ref*{sec:p17-finite-subjective-probability}} & \thepage & The vector is nonnegative and normalized, with coordinates C(e\_i). Zero-weight states are allowed; weights are not asserted to be measured objective chances. \\
\phantomsection\label{entry:symbol-registry-01121}\(\mathcal U,\mathcal D\) & \texttt{\detokenize{\mathcal U,}\allowbreak\detokenize{\mathcal D}} & weakly acceptable and strictly negative bet halfspaces & P17 finite-subjective-probability & meta & \hyperref[sec:p17-finite-subjective-probability]{\ref*{sec:p17-finite-subjective-probability}} & \thepage & The separating probability vector is constructed explicitly. The first halfspace is closed and the second open, both convex, with zero threshold and orientation fixed by sure payments. \\
\phantomsection\label{entry:symbol-registry-01122}\(q(A),b_A,a_j\) & \texttt{\detokenize{q,}\allowbreak\detokenize{b_A,}\allowbreak\detokenize{a_j}} & quoted event price net event contract and signed portfolio quantity & P17 finite-subjective-probability & meta & \hyperref[sec:p17-finite-subjective-probability]{\ref*{sec:p17-finite-subjective-probability}} & \thepage & The finite contract model permits reverse transactions at the same quotes and settlement time. The theorem distinguishes a strict loss in every state from a loss confined to zero-weight states. \\
\phantomsection\label{entry:symbol-registry-01123}\(V(x)\) & \texttt{\detokenize{V}} & nonlinear two-state payoff evaluation in the additivity counterexample & P17 finite-subjective-probability & meta & \hyperref[sec:p17-finite-subjective-probability]{\ref*{sec:p17-finite-subjective-probability}} & \thepage & The exponential payoff score is continuous and monotone but violates invariance under arbitrary common state-dependent additions. It remains a transformed-payoff expected-utility example. \\
\phantomsection\label{entry:symbol-registry-01179}\(N,n\) & \texttt{\detokenize{N,}\allowbreak\detokenize{n}} & finite electorate and its positive size & P17 social-choice & meta & \hyperref[sec:p17-social-choice]{\ref*{sec:p17-social-choice}} & \thepage & Voters are indexed from one through n. Finiteness and positive size are explicit in the Arrow and binary-majority statements. \\
\phantomsection\label{entry:symbol-registry-01180}\(X,\ell,\#\) & \texttt{\detokenize{X,}\allowbreak\detokenize{\ell,}\allowbreak\detokenize{\#}} & alternative set finite size and counting notation & P17 social-choice & meta & \hyperref[sec:p17-social-choice]{\ref*{sec:p17-social-choice}} & \thepage & Arrow permits arbitrary alternative sets with at least three members. Finite rank and candidate-count procedures impose their own nonempty finite domain and use ell for its size. \\
\phantomsection\label{entry:symbol-registry-01181}\(\mathcal L(X),\mathcal W(X),P,P'\) & \texttt{\detokenize{\mathcal L(X),}\allowbreak\detokenize{\mathcal W(X),}\allowbreak\detokenize{P,}\allowbreak\detokenize{P'}} & strict ordering and weak preference classes and voter profiles & P17 social-choice & meta & \hyperref[sec:p17-social-choice]{\ref*{sec:p17-social-choice}} & \thepage & The unrestricted Arrow domain contains every strict-input profile. Social outputs are complete transitive weak relations and may initially allow ties. \\
\phantomsection\label{entry:symbol-registry-01182}\(F,\succeq_P,\succ_P,\sim_P\) & \texttt{\detokenize{F,}\allowbreak\detokenize{\succeq_P,}\allowbreak\detokenize{\succ_P,}\allowbreak\detokenize{\sim_P}} & social welfare map and its weak strict and indifferent comparisons & P17 social-choice & meta & \hyperref[sec:p17-social-choice]{\ref*{sec:p17-social-choice}} & \thepage & The output is a relation rather than a selected alternative. Dictatorship means one fixed voter's strict comparisons determine its reflexive weak closure across every profile and pair. \\
\phantomsection\label{entry:symbol-registry-01183}\(V_P(x,y),\succeq^M_P\) & \texttt{\detokenize{V_P,}\allowbreak\detokenize{\succeq^M_P}} & pairwise vote count and weak majority relation & P17 social-choice & meta & \hyperref[sec:p17-social-choice]{\ref*{sec:p17-social-choice}} & \thepage & The count concerns distinct alternatives; the weak relation includes the diagonal explicitly. Equal opposing counts create a tie, and completeness alone does not give transitivity. \\
\phantomsection\label{entry:symbol-registry-01184}\(C_P(x),B_P(x),w_k,r_i(x)\) & \texttt{\detokenize{C_P,}\allowbreak\detokenize{B_P,}\allowbreak\detokenize{w_k,}\allowbreak\detokenize{r_i}} & victory and rank scores weights and individual positions & P17 social-choice & meta & \hyperref[sec:p17-social-choice]{\ref*{sec:p17-social-choice}} & \thepage & Victory scores count strictly won pairs. Strictly decreasing rank weights produce the stated Pareto comparison; both score rules can violate profile-based IIA. \\
\phantomsection\label{entry:symbol-registry-01185}\(R,D(R,P)\) & \texttt{\detokenize{R,}\allowbreak\detokenize{D(R,}\allowbreak\detokenize{P)}} & candidate strict ordering and total disagreement count & P17 social-choice & meta & \hyperref[sec:p17-social-choice]{\ref*{sec:p17-social-choice}} & \thepage & The count is over ordered alternative pairs and voters. Minimization and most-common-ranking choices use an explicitly fixed tie priority on the finite ordering set. \\
\phantomsection\label{entry:symbol-registry-01186}\(P^m,m,m^*,b\) & \texttt{\detokenize{P^m,}\allowbreak\detokenize{m,}\allowbreak\detokenize{m^*,}\allowbreak\detokenize{b}} & extreme profiles step index pivotal voter and reference alternative & P17 social-choice & meta & \hyperref[sec:p17-social-choice]{\ref*{sec:p17-social-choice}} & \thepage & The pivotal index is the first socially top step. IIA fixes it independently of all other pair rankings; the proof does not assume monotonic persistence of top status. \\
\phantomsection\label{entry:symbol-registry-01187}\(A,C,a,c,d,e\) & \texttt{\detokenize{A,}\allowbreak\detokenize{C,}\allowbreak\detokenize{a,}\allowbreak\detokenize{c,}\allowbreak\detokenize{d,}\allowbreak\detokenize{e}} & extreme-comparison sets and distinct witness alternatives & P17 social-choice & meta & \hyperref[sec:p17-social-choice]{\ref*{sec:p17-social-choice}} & \thepage & The complete weak social relation supplies distinct witnesses around the reference alternative when it is not extreme. The proof includes possible social ties. \\
\phantomsection\label{entry:symbol-registry-01188}\(s,s',t,f\) & \texttt{\detokenize{s,}\allowbreak\detokenize{s',}\allowbreak\detokenize{t,}\allowbreak\detokenize{f}} & binary preference vectors balanced intermediate vector and aggregation rule & P17 social-choice & meta & \hyperref[sec:p17-social-choice]{\ref*{sec:p17-social-choice}} & \thepage & The binary domain permits individual indifference, encoded by zero. Both directions of positive responsiveness are part of the condition, including when testing its independence from neutrality. \\
\phantomsection\label{entry:symbol-registry-01189}\(\operatorname{sgn}\) & \texttt{\detokenize{\operatorname{sgn}}} & real sign function for binary aggregation & P17 social-choice & meta & \hyperref[sec:p17-social-choice]{\ref*{sec:p17-social-choice}} & \thepage & The function returns one, zero or minus one according to the argument's sign. The characterized majority rule applies it to the sum of the signed votes. \\
\phantomsection\label{entry:symbol-registry-01190}\(-s\) & \texttt{\detokenize{-s}} & reversal of a binary preference profile & P17 social-choice & meta & \hyperref[sec:p17-social-choice]{\ref*{sec:p17-social-choice}} & \thepage & Neutrality reverses the output under simultaneous sign reversal. A balanced profile is a permutation of its reversal; anonymity therefore fixes its output at zero. \\
\phantomsection\label{entry:symbol-registry-01204}\(X,\leq,<\) & \texttt{\detokenize{X,}\allowbreak\detokenize{\leq,}\allowbreak\detokenize{<}} & nonempty alternative set and its supplied common linear axis & P17 single-peaked-majority & meta & \hyperref[sec:p17-single-peaked-majority]{\ref*{sec:p17-single-peaked-majority}} & \thepage & The axis is distinct from each preference ordering. The source interval is preserved, and the stated generalization uses a supplied ordered set with actual peaks. \\
\phantomsection\label{entry:symbol-registry-01205}\(p_i,a_i^*\) & \texttt{\detokenize{p_i,}\allowbreak\detokenize{a_i^*}} & individual peak and its explicitly identified source notation & P17 single-peaked-majority & meta & \hyperref[sec:p17-single-peaked-majority]{\ref*{sec:p17-single-peaked-majority}} & \thepage & Preferences improve strictly toward the peak on either side of the common axis. The theorem adds no distance, continuity or numerical utility assumption to these comparisons. \\
\phantomsection\label{entry:symbol-registry-01206}\(\mathcal L_{\mathrm{SP}}(X,\leq),\mathcal D_{\mathrm{SP}}\) & \texttt{\detokenize{\mathcal L_{\mathrm{SP}},}\allowbreak\detokenize{\mathcal D_{\mathrm{SP}}}} & single-peaked strict ordering class and restricted profile domain & P17 single-peaked-majority & meta & \hyperref[sec:p17-single-peaked-majority]{\ref*{sec:p17-single-peaked-majority}} & \thepage & The individual rankings are strict. The weak-input convention of the binary May result is not inherited by this domain. \\
\phantomsection\label{entry:symbol-registry-01207}\(N,n,k\) & \texttt{\detokenize{N,}\allowbreak\detokenize{n,}\allowbreak\detokenize{k}} & finite electorate positive size and parity parameter & P17 single-peaked-majority & meta & \hyperref[sec:p17-single-peaked-majority]{\ref*{sec:p17-single-peaked-majority}} & \thepage & Odd size has the form \(2k+1\) with \(k\geq0\); even size has the form \(2k\) with \(k\geq1\). Repeated peaks count as separate voter occurrences. \\
\phantomsection\label{entry:symbol-registry-01208}\(V_P(x,y),\succ^M_P,\succeq^M_P,\sim^M_P\) & \texttt{\detokenize{V_P(x,}\allowbreak\detokenize{y),}\allowbreak\detokenize{\succ^M_P,}\allowbreak\detokenize{\succeq^M_P,}\allowbreak\detokenize{\sim^M_P}} & pairwise vote counts and majority comparisons on the restricted domain & P17 single-peaked-majority & meta & \hyperref[sec:p17-single-peaked-majority]{\ref*{sec:p17-single-peaked-majority}} & \thepage & Strict input comparisons and equal-count social ties retain the predecessor conventions. Odd-electorate transitivity is proved for the entire relation; even-electorate winner statements do not imply transitivity. \\
\phantomsection\label{entry:symbol-registry-01209}\(F_{\mathrm{SP}}\) & \texttt{\detokenize{F_{\mathrm{SP}}}} & majority social welfare function on the single-peaked domain & P17 single-peaked-majority & meta & \hyperref[sec:p17-single-peaked-majority]{\ref*{sec:p17-single-peaked-majority}} & \thepage & For odd electorates the proved majority relation gives a complete transitive social output. Its IIA condition compares only admissible profiles, so the unrestricted-domain premise of Arrow is absent. \\
\phantomsection\label{entry:symbol-registry-01210}\(p_{(j)},m,l,u,[l,u]_X\) & \texttt{\detokenize{p_{(j)},}\allowbreak\detokenize{m,}\allowbreak\detokenize{l,}\allowbreak\detokenize{u,}\allowbreak\detokenize{[l,}\allowbreak\detokenize{u]_X}} & sorted peak occurrences odd median even central peaks and their axis interval & P17 single-peaked-majority & meta & \hyperref[sec:p17-single-peaked-majority]{\ref*{sec:p17-single-peaked-majority}} & \thepage & Sorted positions are not voter identities. The even weak-winner interval is intersected with X, and a strict winner exists exactly when the two central peaks coincide. \\
\phantomsection\label{entry:symbol-registry-01211}\(q,a_j,r_i(j),h_i\) & \texttt{\detokenize{q,}\allowbreak\detokenize{a_j,}\allowbreak\detokenize{r_i(j),}\allowbreak\detokenize{h_i}} & finite axis size alternatives rank arrays and peak indices & P17 single-peaked-majority & meta & \hyperref[sec:p17-single-peaked-majority]{\ref*{sec:p17-single-peaked-majority}} & \thepage & The supplied-axis procedure validates a finite strict-ranking encoding. Its operation bounds use random-access arrays and unit-cost index/count operations; finite observations do not certify a whole interval preference. \\
\phantomsection\label{entry:symbol-registry-01212}\(w\) & \texttt{\detokenize{w}} & candidate strict or weak Condorcet winner & P17 single-peaked-majority & meta & \hyperref[sec:p17-single-peaked-majority]{\ref*{sec:p17-single-peaked-majority}} & \thepage & Winner comparisons are against every other alternative on the declared domain. Having a winner is distinguished from transitivity of the full majority relation. \\
\phantomsection\label{entry:symbol-registry-01213}\(d_p(t)\) & \texttt{\detokenize{d_p(t)}} & explicit real distance in the full-interval parity counterexample & P17 single-peaked-majority & meta & \hyperref[sec:p17-single-peaked-majority]{\ref*{sec:p17-single-peaked-majority}} & \thepage & Defined by \(d_p(t)=\max\{t-p,p-t\}\). Smaller distance is preferred, with the larger coordinate breaking a tie; this construction supplies one counterexample and is not a hypothesis of the main theorem. \\
\phantomsection\label{entry:symbol-registry-01257}\(S,Z,Y,n,r,t,P,Q,p_{sz},q_{sy}\) & \texttt{\detokenize{S,}\allowbreak\detokenize{Z,}\allowbreak\detokenize{Y,}\allowbreak\detokenize{n,}\allowbreak\detokenize{r,}\allowbreak\detokenize{t,}\allowbreak\detokenize{P,}\allowbreak\detokenize{Q,}\allowbreak\detokenize{p_{sz},}\allowbreak\detokenize{q_{sy}}} & finite state and signal sets sizes and row-stochastic experiment matrices & P17 value-of-information & meta & \hyperref[sec:p17-value-of-information]{\ref*{sec:p17-value-of-information}} & \thepage & The same state set is used in every comparison. Matrix entry q\_sy is distinguished from marginal signal mass q\_z. \\
\phantomsection\label{entry:symbol-registry-01258}\(\pi,S_\pi,q_z,\pi^z,\Delta(S)\) & \texttt{\detokenize{\pi,}\allowbreak\detokenize{S_\pi,}\allowbreak\detokenize{q_z,}\allowbreak\detokenize{\pi^z,}\allowbreak\detokenize{\Delta(S)}} & prior support signal probability posterior and simplex & P17 value-of-information & meta & \hyperref[sec:p17-value-of-information]{\ref*{sec:p17-value-of-information}} & \thepage & A posterior is defined only for a positive-probability signal. Zero-prior states remain in the experiment but contribute no expected payoff at that prior. \\
\phantomsection\label{entry:symbol-registry-01259}\(A,u,\delta,\widehat\delta,J,V,V_0,V_{\mathrm{PI}}\) & \texttt{\detokenize{A,}\allowbreak\detokenize{u,}\allowbreak\detokenize{\delta,}\allowbreak\detokenize{\widehat\delta,}\allowbreak\detokenize{J,}\allowbreak\detokenize{V,}\allowbreak\detokenize{V_0,}\allowbreak\detokenize{V_{\mathrm{PI}}}} & actions payoff policies expected payoff and information values & P17 value-of-information & meta & \hyperref[sec:p17-value-of-information]{\ref*{sec:p17-value-of-information}} & \thepage & Policies randomize using only the observed signal. All maxima are attained for nonempty finite action sets; signal acquisition costs are measured in the declared payoff units. \\
\phantomsection\label{entry:symbol-registry-01260}\(\operatorname{EVSI},\operatorname{EVPI},c\) & \texttt{\detokenize{\operatorname{EVSI},}\allowbreak\detokenize{\operatorname{EVPI},}\allowbreak\detokenize{c}} & expected sample and perfect information improvements and fixed cost & P17 value-of-information & meta & \hyperref[sec:p17-value-of-information]{\ref*{sec:p17-value-of-information}} & \thepage & The nonnegative value bound assumes the signal can be ignored and the original actions remain available. \\
\phantomsection\label{entry:symbol-registry-01261}\(M,P>Q,P\supset Q,\mathcal A,B(P,\mathcal A)\) & \texttt{\detokenize{M,}\allowbreak\detokenize{>,}\allowbreak\detokenize{\supset,}\allowbreak\detokenize{\mathcal A,}\allowbreak\detokenize{B}} & garbling matrix source sufficiency and informativeness comparisons and attainable loss set & P17 value-of-information & meta & \hyperref[sec:p17-value-of-information]{\ref*{sec:p17-value-of-information}} & \thepage & The source symbols > and superset name non-strict experiment comparisons. The loss-vector set A is nonempty closed bounded convex in the finite statewise loss space. \\
\phantomsection\label{entry:symbol-registry-01262}\(\langle H,J\rangle,d^2(H,J),\operatorname{Trace},\operatorname{diag},\operatorname{abs}\) & \texttt{\detokenize{\langle\cdot,}\allowbreak\detokenize{\cdot\rangle,}\allowbreak\detokenize{d^2,}\allowbreak\detokenize{\operatorname{Trace},}\allowbreak\detokenize{\operatorname{diag},}\allowbreak\detokenize{\operatorname{abs}}} & finite-array inner product squared distance trace diagonal and real absolute value & P17 value-of-information & meta & \hyperref[sec:p17-value-of-information]{\ref*{sec:p17-value-of-information}} & \thepage & The chapter defines these operations before its separation and matrix proofs. Ordinary set-cardinality bars are not repurposed for real absolute value. \\
\phantomsection\label{entry:symbol-registry-01263}\(\alpha,\beta,O_\varepsilon,\mathcal M,h(D,M),D\) & \texttt{\detokenize{\alpha,}\allowbreak\detokenize{\beta,}\allowbreak\detokenize{O_\varepsilon,}\allowbreak\detokenize{\mathcal M,}\allowbreak\detokenize{h,}\allowbreak\detokenize{D}} & game mixtures orthant garbling polytope and trace payoff & P17 value-of-information & meta & \hyperref[sec:p17-value-of-information]{\ref*{sec:p17-value-of-information}} & \thepage & The finite minimax proof derives the quantifier change for the cube of D arrays and the polytope of stochastic M arrays; it does not exchange quantifiers without that argument. \\
\phantomsection\label{entry:symbol-registry-01264}\(\mu_{P,\pi},\phi,\nu^z,\widetilde q_y,\widetilde\pi^y\) & \texttt{\detokenize{\mu_{P,\pi},}\allowbreak\detokenize{\phi,}\allowbreak\detokenize{\nu^z,}\allowbreak\detokenize{\widetilde q_y,}\allowbreak\detokenize{\widetilde\pi^y}} & posterior law convex test function specified posterior vectors and output posterior data & P17 value-of-information & meta & \hyperref[sec:p17-value-of-information]{\ref*{sec:p17-value-of-information}} & \thepage & Only positive-weight signals contribute to the posterior law. The converse garbling characterization fixes a full-support prior. \\
\phantomsection\label{entry:symbol-registry-01265}\(X,\mu_x,\nu_y,c,R,\varepsilon,\iota,\lambda_{xy},T_{yx}\) & \texttt{\detokenize{X,}\allowbreak\detokenize{\mu,}\allowbreak\detokenize{\nu,}\allowbreak\detokenize{c,}\allowbreak\detokenize{R,}\allowbreak\detokenize{\varepsilon,}\allowbreak\detokenize{\iota,}\allowbreak\detokenize{\lambda,}\allowbreak\detokenize{T}} & finite vector support masses common mean simplex embedding coupling and reverse transformation & P17 value-of-information & meta & \hyperref[sec:p17-value-of-information]{\ref*{sec:p17-value-of-information}} & \thepage & The affine embedding is explicit and injective. Every zero-weight input row of T is a point mass at that same vector, so the required row mean holds even on null rows. \\
\phantomsection\label{entry:symbol-registry-01279}\(k,\mathcal P_{P,k},\mathcal P^*_{P,k},M^f,P>_kQ\) & \texttt{\detokenize{k,}\allowbreak\detokenize{\mathcal P_{P,k},}\allowbreak\detokenize{\mathcal P^*_{P,k},}\allowbreak\detokenize{M^f,}\allowbreak\detokenize{>_k}} & decision bound obtainable output experiments deterministic transformations and restricted comparison & P17 restricted-decisions-binary-testing & meta & \hyperref[sec:p17-restricted-decisions-binary-testing]{\ref*{sec:p17-restricted-decisions-binary-testing}} & \thepage & The labelled partition permits empty cells; repeated actions implement lists of length at most k. The named source comparison is non-strict. \\
\phantomsection\label{entry:symbol-registry-01280}\(C,\operatorname{conv}C,w_f,\lambda_{zj},D\) & \texttt{\detokenize{C,}\allowbreak\detokenize{\operatorname{conv}C,}\allowbreak\detokenize{w_f,}\allowbreak\detokenize{\lambda_{zj},}\allowbreak\detokenize{D}} & loss generating set convex hull deterministic mixture weights signalwise coefficients and loss-row matrix & P17 restricted-decisions-binary-testing & meta & \hyperref[sec:p17-restricted-decisions-binary-testing]{\ref*{sec:p17-restricted-decisions-binary-testing}} & \thepage & The mixture weight multiplies over every signal. The finite-signal hull identity is proved for any nonempty generating set without a limiting purification assumption. \\
\phantomsection\label{entry:symbol-registry-01281}\(x,a,b,t_j,d_j,(v)_+\) & \texttt{\detokenize{x,}\allowbreak\detokenize{a,}\allowbreak\detokenize{b,}\allowbreak\detokenize{t_j,}\allowbreak\detokenize{d_j,}\allowbreak\detokenize{(v)_+}} & binary posterior coordinate affine part breakpoints nonnegative slope changes and positive part & P17 restricted-decisions-binary-testing & meta & \hyperref[sec:p17-restricted-decisions-binary-testing]{\ref*{sec:p17-restricted-decisions-binary-testing}} & \thepage & Every finite maximum of affine functions is explicitly decomposed into its affine part and hinges. Common posterior barycenters give equality of the affine expectations. \\
\phantomsection\label{entry:symbol-registry-01282}\(\delta,a_P,b_P,\mathcal R(P),f_P(t)\) & \texttt{\detokenize{\delta,}\allowbreak\detokenize{a_P,}\allowbreak\detokenize{b_P,}\allowbreak\detokenize{\mathcal R,}\allowbreak\detokenize{f_P}} & randomized test first and second error attainable region and its optimal lower boundary & P17 restricted-decisions-binary-testing & meta & \hyperref[sec:p17-restricted-decisions-binary-testing]{\ref*{sec:p17-restricted-decisions-binary-testing}} & \thepage & The defining first-error constraint is equality. The proof derives equivalence with a size-at-most optimum and reconstructs the full region by central symmetry. \\
\phantomsection\label{entry:symbol-registry-01283}\(p_z,q_z,\ell_z,q_*,A_j,B_j,\delta^*,c\) & \texttt{\detokenize{p_z,}\allowbreak\detokenize{q_z,}\allowbreak\detokenize{\ell_z,}\allowbreak\detokenize{q_*,}\allowbreak\detokenize{A_j,}\allowbreak\detokenize{B_j,}\allowbreak\detokenize{\delta^*,}\allowbreak\detokenize{c}} & null and alternative masses likelihood ratios singular alternative mass cumulative masses optimal test and threshold & P17 restricted-decisions-binary-testing & meta & \hyperref[sec:p17-restricted-decisions-binary-testing]{\ref*{sec:p17-restricted-decisions-binary-testing}} & \thepage & Zero null mass is treated before forming ratios. Tied ratios have the same optimal curve; threshold randomization gives every exact size. \\
\phantomsection\label{entry:symbol-registry-01284}\(P^{[m]},Q^{[m]},M^{[m]},z,y\) & \texttt{\detokenize{P^{[m]},}\allowbreak\detokenize{Q^{[m]},}\allowbreak\detokenize{M^{[m]},}\allowbreak\detokenize{z,}\allowbreak\detokenize{y}} & conditionally independent repeated experiments product garbling and signal tuples & P17 restricted-decisions-binary-testing & meta & \hyperref[sec:p17-restricted-decisions-binary-testing]{\ref*{sec:p17-restricted-decisions-binary-testing}} & \thepage & Every product uses one fixed state across its observations. Equal marginal experiments alone do not impose this conditional dependence structure. \\
\phantomsection\label{entry:symbol-registry-01369}\(\mu,\nu,X,d,\preceq_{\mathrm{cx}},K,\delta_x,\|\cdot\|_2,\operatorname{conv}\) & \texttt{\detokenize{\mu,}\allowbreak\detokenize{\nu,}\allowbreak\detokenize{X,}\allowbreak\detokenize{d,}\allowbreak\detokenize{\preceq_{\mathrm{cx}},}\allowbreak\detokenize{K,}\allowbreak\detokenize{\delta_x,}\allowbreak\detokenize{\|\cdot\|_2,}\allowbreak\detokenize{\operatorname{conv}}} & bounded vector laws Borel carrier dimension convex order kernel point mass Euclidean norm and convex hull & P17 general-blackwell-comparison & meta & \hyperref[sec:p17-general-blackwell-comparison]{\ref*{sec:p17-general-blackwell-comparison}} & \thepage & Mu is the less dispersed law in this chapter. Mean preservation is required at every input, with an explicit measurable repair on the null input set. \\
\phantomsection\label{entry:symbol-registry-01370}\(C,\phi,\lambda,E,H,(a,c),\gamma,\ell_x,\psi,r_k,\operatorname{dist},\omega_{f,K}\) & \texttt{\detokenize{C,}\allowbreak\detokenize{\phi,}\allowbreak\detokenize{\lambda,}\allowbreak\detokenize{E,}\allowbreak\detokenize{H,}\allowbreak\detokenize{a,}\allowbreak\detokenize{c,}\allowbreak\detokenize{\gamma,}\allowbreak\detokenize{\ell_x,}\allowbreak\detokenize{\psi,}\allowbreak\detokenize{r_k,}\allowbreak\detokenize{\operatorname{dist},}\allowbreak\detokenize{\omega_{f,K}}} & compact convex set test law truncated epigraph separation data affine minorants finite approximation distance and modulus & P17 general-blackwell-comparison & meta & \hyperref[sec:p17-general-blackwell-comparison]{\ref*{sec:p17-general-blackwell-comparison}} & \thepage & The affine approximation works even on a lower-dimensional compact convex set. The coefficient c of the epigraph height is proved strictly negative. \\
\phantomsection\label{entry:symbol-registry-01371}\(N,h_N,C_{N,a},c_{N,a},\mu_N,A_N,\lambda_{N,i},v_{N,a,\epsilon},w_{N,\epsilon},U_N,\nu_N,B_N\) & \texttt{\detokenize{N,}\allowbreak\detokenize{h_N,}\allowbreak\detokenize{C_{N,a},}\allowbreak\detokenize{c_{N,a},}\allowbreak\detokenize{\mu_N,}\allowbreak\detokenize{A_N,}\allowbreak\detokenize{\lambda_{N,i},}\allowbreak\detokenize{v_{N,a,\epsilon},}\allowbreak\detokenize{w_{N,\epsilon},}\allowbreak\detokenize{U_N,}\allowbreak\detokenize{\nu_N,}\allowbreak\detokenize{B_N}} & dyadic level mesh cells centroids lower laws supports normalized coordinates vertices weights splitting kernel and upper laws & P17 general-blackwell-comparison & meta & \hyperref[sec:p17-general-blackwell-comparison]{\ref*{sec:p17-general-blackwell-comparison}} & \thepage & Zero-mass cells have no centroid. Barycentric coordinates divide by the mesh, and shared faces satisfy the same kernel identity. \\
\phantomsection\label{entry:symbol-registry-01372}\(T_N,V_N,J_D,H_{N,D},W_N,D_j\) & \texttt{\detokenize{T_N,}\allowbreak\detokenize{V_N,}\allowbreak\detokenize{J_D,}\allowbreak\detokenize{H_{N,D},}\allowbreak\detokenize{W_N,}\allowbreak\detokenize{D_j}} & centroid refinement reverse grid refinement distant comparison transported comparison compatible kernels and diagonal indices & P17 general-blackwell-comparison & meta & \hyperref[sec:p17-general-blackwell-comparison]{\ref*{sec:p17-general-blackwell-comparison}} & \thepage & Kernel composition is in row order. Compatibility is W\_N=T\_N W\_(N+1) V\_(N+1), an equality of all finite rows. \\
\phantomsection\label{entry:symbol-registry-01373}\(L_N,R_N,L_*,R_*,\mathcal F_N,\mathcal G_N,\mathcal F_\infty,K_0,b(x),D\) & \texttt{\detokenize{L_N,}\allowbreak\detokenize{R_N,}\allowbreak\detokenize{L_*,}\allowbreak\detokenize{R_*,}\allowbreak\detokenize{\mathcal F_N,}\allowbreak\detokenize{\mathcal G_N,}\allowbreak\detokenize{\mathcal F_\infty,}\allowbreak\detokenize{K_0,}\allowbreak\detokenize{b,}\allowbreak\detokenize{D}} & two-sided process limits forward and reverse information conditional kernel barycenter and good-input set & P17 general-blackwell-comparison & meta & \hyperref[sec:p17-general-blackwell-comparison]{\ref*{sec:p17-general-blackwell-comparison}} & \thepage & The reverse information contains the entire remaining R tail. The terminal mean identity is proved through fixed forward information and then its generated limit. \\
\phantomsection\label{entry:symbol-registry-01374}\(\mathsf E,\mathsf F,P_s,Q_s,S,Z,Y,\mathcal Z,\mathcal Y,\pi,\rho,\eta,p,q,\mu_{\mathsf E,\pi},P_s^\mu\) & \texttt{\detokenize{\mathsf E,}\allowbreak\detokenize{\mathsf F,}\allowbreak\detokenize{P_s,}\allowbreak\detokenize{Q_s,}\allowbreak\detokenize{S,}\allowbreak\detokenize{Z,}\allowbreak\detokenize{Y,}\allowbreak\detokenize{\mathcal Z,}\allowbreak\detokenize{\mathcal Y,}\allowbreak\detokenize{\pi,}\allowbreak\detokenize{\rho,}\allowbreak\detokenize{\eta,}\allowbreak\detokenize{p,}\allowbreak\detokenize{q,}\allowbreak\detokenize{\mu_{\mathsf E,\pi},}\allowbreak\detokenize{P_s^\mu}} & measurable experiments state laws states observation spaces priors mixtures posterior maps laws and canonical experiments & P17 general-blackwell-comparison & meta & \hyperref[sec:p17-general-blackwell-comparison]{\ref*{sec:p17-general-blackwell-comparison}} & \thepage & Posterior reconstruction uses a full-support prior. Risk-set comparisons allow arbitrary measurable observations; full observation recovery has its stated kernel hypothesis. \\
\phantomsection\label{entry:symbol-registry-01375}\(D,u,\delta,J,V,B(\mathsf E,C),f,g,h,R,T,G,J(q,dy)\) & \texttt{\detokenize{D,}\allowbreak\detokenize{u,}\allowbreak\detokenize{\delta,}\allowbreak\detokenize{J,}\allowbreak\detokenize{V,}\allowbreak\detokenize{B,}\allowbreak\detokenize{f,}\allowbreak\detokenize{g,}\allowbreak\detokenize{h,}\allowbreak\detokenize{R,}\allowbreak\detokenize{T,}\allowbreak\detokenize{G,}\allowbreak\detokenize{J}} & action menu payoffs policies value risk set decisions posterior kernels observation simulation and recovery kernel & P17 general-blackwell-comparison & meta & \hyperref[sec:p17-general-blackwell-comparison]{\ref*{sec:p17-general-blackwell-comparison}} & \thepage & J denotes expected payoff in the definition and a separately typed recovery kernel in the simulation proof. Posterior mean-preserving transport runs opposite to observation garbling. \\
\phantomsection\label{entry:symbol-registry-01376}\(I_N,z_{N,a},\lambda_j,p_{j,N},p_N,j_l,\gamma_j,\gamma,D_t,\epsilon\) & \texttt{\detokenize{I_N,}\allowbreak\detokenize{z_{N,a},}\allowbreak\detokenize{\lambda_j,}\allowbreak\detokenize{p_{j,N},}\allowbreak\detokenize{p_N,}\allowbreak\detokenize{j_l,}\allowbreak\detokenize{\gamma_j,}\allowbreak\detokenize{\gamma,}\allowbreak\detokenize{D_t,}\allowbreak\detokenize{\epsilon}} & compact selection cells representatives laws cell masses diagonal subsequence joint limits and convex rectangle tests & P17 general-blackwell-comparison & meta & \hyperref[sec:p17-general-blackwell-comparison]{\ref*{sec:p17-general-blackwell-comparison}} & \thepage & The single selected subsequence works for every continuous test. Convex positive-part rectangle tests prove uniqueness of the standard measure. \\
\phantomsection\label{entry:symbol-registry-01519}\(N,A_i,\mathcal A,X,X_{-i},\mathcal A_{-i},x_i,P_x,u_i,U_i,x[j\leftarrow z_j]\) & \texttt{\detokenize{N,}\allowbreak\detokenize{A_i,}\allowbreak\detokenize{\mathcal A,}\allowbreak\detokenize{X,}\allowbreak\detokenize{X_{-i},}\allowbreak\detokenize{\mathcal A_{-i},}\allowbreak\detokenize{x_i,}\allowbreak\detokenize{P_x,}\allowbreak\detokenize{u_i,}\allowbreak\detokenize{U_i,}\allowbreak\detokenize{x[j\leftarrow z_j]}} & players actions pure and mixed profile spaces independent profile law cardinal payoffs expectations and sequential replacement & P17 finite-strategic-games & meta & \hyperref[sec:p17-finite-strategic-games]{\ref*{sec:p17-finite-strategic-games}} & \thepage & Empty opponent products equal one. Player-specific payoff scales need not be interpersonally comparable. \\
\phantomsection\label{entry:symbol-registry-01520}\(v_i,\operatorname{BR}_i,\mathcal B,\mathcal E\) & \texttt{\detokenize{v_i,}\allowbreak\detokenize{\operatorname{BR}_i,}\allowbreak\detokenize{\mathcal B,}\allowbreak\detokenize{\mathcal E}} & best pure value best-response face countering correspondence and equilibrium set & P17 finite-strategic-games & meta & \hyperref[sec:p17-finite-strategic-games]{\ref*{sec:p17-finite-strategic-games}} & \thepage & The three support and deviation criteria have both directions proved before existence is invoked. \\
\phantomsection\label{entry:symbol-registry-01521}\(\Delta_d,d,F_k,b_F,\lambda_j,\mu_k,\sigma\) & \texttt{\detokenize{\Delta_d,}\allowbreak\detokenize{d,}\allowbreak\detokenize{F_k,}\allowbreak\detokenize{b_F,}\allowbreak\detokenize{\lambda_j,}\allowbreak\detokenize{\mu_k,}\allowbreak\detokenize{\sigma}} & standard simplex dimension nested faces barycenters original and subdivision coordinates and their ordering & P17 finite-strategic-games & meta & \hyperref[sec:p17-finite-strategic-games]{\ref*{sec:p17-finite-strategic-games}} & \thepage & The dimension-zero case is explicit. Barycentric meshes shrink for positive dimension. \\
\phantomsection\label{entry:symbol-registry-01522}\(K,\pi_K,S,R,\operatorname{dist},\mathcal C,\eta,\delta,\epsilon,z_j,w_j,\theta_j,f_\epsilon\) & \texttt{\detokenize{K,}\allowbreak\detokenize{\pi_K,}\allowbreak\detokenize{S,}\allowbreak\detokenize{R,}\allowbreak\detokenize{\operatorname{dist},}\allowbreak\detokenize{\mathcal C,}\allowbreak\detokenize{\eta,}\allowbreak\detokenize{\delta,}\allowbreak\detokenize{\epsilon,}\allowbreak\detokenize{z_j,}\allowbreak\detokenize{w_j,}\allowbreak\detokenize{\theta_j,}\allowbreak\detokenize{f_\epsilon}} & compact convex domain projection enclosing simplex distance correspondence and approximate-selection data & P17 finite-strategic-games & meta & \hyperref[sec:p17-finite-strategic-games]{\ref*{sec:p17-finite-strategic-games}} & \thepage & Closed graph supplies upper approximation near the fixed limit point. No continuous exact selection is assumed. \\
\phantomsection\label{entry:symbol-registry-01523}\(g_{i,a},G_i,T_i,T,\sigma,h_i,C,C_i,v_i^+,v_i^-\) & \texttt{\detokenize{g_{i,a},}\allowbreak\detokenize{G_i,}\allowbreak\detokenize{T_i,}\allowbreak\detokenize{T,}\allowbreak\detokenize{\sigma,}\allowbreak\detokenize{h_i,}\allowbreak\detokenize{C,}\allowbreak\detokenize{C_i,}\allowbreak\detokenize{v_i^+,}\allowbreak\detokenize{v_i^-}} & positive gains normalization improvement map game automorphism and interchangeable sets factors and payoff extrema & P17 finite-strategic-games & meta & \hyperref[sec:p17-finite-strategic-games]{\ref*{sec:p17-finite-strategic-games}} & \thepage & Nash's solution-set terminology is distinct from an individual equilibrium and from an NM imputation solution. \\
\phantomsection\label{entry:symbol-registry-01524}\(M,f,v,X^*,Y^*,e_a,e_b,p,q,D_1,D_2,\rho\) & \texttt{\detokenize{M,}\allowbreak\detokenize{f,}\allowbreak\detokenize{v,}\allowbreak\detokenize{X^*,}\allowbreak\detokenize{Y^*,}\allowbreak\detokenize{e_a,}\allowbreak\detokenize{e_b,}\allowbreak\detokenize{p,}\allowbreak\detokenize{q,}\allowbreak\detokenize{D_1,}\allowbreak\detokenize{D_2,}\allowbreak\detokenize{\rho}} & zero-sum matrix payoff value optimal factors pure vectors example mixtures best-response differences and dominance perturbation & P17 finite-strategic-games & meta & \hyperref[sec:p17-finite-strategic-games]{\ref*{sec:p17-finite-strategic-games}} & \thepage & The source's dominance perturbation has its corrected direction. A pure saddle does not by itself imply strong solvability. \\
\phantomsection\label{entry:symbol-registry-01525}\(\mu,d_i(b_i\mid a_i),M,L,k,\varepsilon,r(x)\) & \texttt{\detokenize{\mu,}\allowbreak\detokenize{d_i(b_i\mid a_i),}\allowbreak\detokenize{M,}\allowbreak\detokenize{L,}\allowbreak\detokenize{k,}\allowbreak\detokenize{\varepsilon,}\allowbreak\detokenize{r(x)}} & joint recommendation law randomized deviation local payoff magnitude dimension bound grid accuracy and regret & P17 finite-strategic-games & meta & \hyperref[sec:p17-finite-strategic-games]{\ref*{sec:p17-finite-strategic-games}} & \thepage & Obedience uses weighted inequalities at zero-probability recommendations. The grid-count bound applies only when k is chosen. \\
\phantomsection\label{entry:symbol-registry-01655}\(S,A,A_t(s),A_t,X_t\) & \texttt{\detokenize{S,}\allowbreak\detokenize{A,}\allowbreak\detokenize{A_t(s),}\allowbreak\detokenize{A_t,}\allowbreak\detokenize{X_t}} & finite states common actions feasible actions and observed state and chosen action & P17 sequential-decisions & meta & \hyperref[def:p17-sequential-model]{\ref*{def:p17-sequential-model}} & \thepage & The chosen random action A\_t has no state argument; it differs from the feasible set A\_t(s). \\
\phantomsection\label{entry:symbol-registry-01656}\(P_t,r_t,\bar r_t,P,r,\bar r\) & \texttt{\detokenize{P_t,}\allowbreak\detokenize{r_t,}\allowbreak\detokenize{\bar r_t,}\allowbreak\detokenize{P,}\allowbreak\detokenize{r,}\allowbreak\detokenize{\bar r}} & transition tables realized rewards and expected one-step rewards & P17 sequential-decisions & meta & \hyperref[def:p17-sequential-model]{\ref*{def:p17-sequential-model}} & \thepage & Unsubscripted tables designate the homogeneous model; payoff timing follows the next-state draw. \\
\phantomsection\label{entry:symbol-registry-01657}\(h_t,\pi_t,\rho,f_t,f\) & \texttt{\detokenize{h_t,}\allowbreak\detokenize{\pi_t,}\allowbreak\detokenize{\rho,}\allowbreak\detokenize{f_t,}\allowbreak\detokenize{f}} & history policy kernel initial law Markov choice and stationary choice & P17 sequential-decisions & meta & \hyperref[def:p17-sequential-model]{\ref*{def:p17-sequential-model}} & \thepage & Policies specify actions at every admissible history including null histories; Markov marginal projection is for a fixed initial law. \\
\phantomsection\label{entry:symbol-registry-01658}\(N,g,\gamma,J_{t,N}^{\pi},v_t\) & \texttt{\detokenize{N,}\allowbreak\detokenize{g,}\allowbreak\detokenize{\gamma,}\allowbreak\detokenize{J_{t,N}^{\pi},}\allowbreak\detokenize{v_t}} & finite horizon terminal reward discount continuation return and backward value & P17 sequential-decisions & meta & \hyperref[def:p17-sequential-finite-return]{\ref*{def:p17-sequential-finite-return}} & \thepage & Finite-horizon gamma lies in [0,1]; gamma\^{}0=1 also at gamma=0. \\
\phantomsection\label{entry:symbol-registry-01659}\(J^\pi,V^*,V^f,R,T,T_f\) & \texttt{\detokenize{J^\pi,}\allowbreak\detokenize{V^*,}\allowbreak\detokenize{V^f,}\allowbreak\detokenize{R,}\allowbreak\detokenize{T,}\allowbreak\detokenize{T_f}} & discounted return optimal value fixed-policy value reward bound and optimality operators & P17 sequential-decisions & meta & \hyperref[def:p17-sequential-discounted]{\ref*{def:p17-sequential-discounted}} & \thepage & Infinite discounted stochastic method has finite state/action tables and gamma<1. \\
\phantomsection\label{entry:symbol-registry-01660}\(\|\cdot\|_\infty,E,B,u_k,w\) & \texttt{\detokenize{\|\cdot\|_\infty,}\allowbreak\detokenize{E,}\allowbreak\detokenize{B,}\allowbreak\detokenize{u_k,}\allowbreak\detokenize{w}} & supremum norm abstract domain contraction operator iterates and fixed point & P17 sequential-decisions & meta & \hyperref[thm:p17-sequential-contraction]{\ref*{thm:p17-sequential-contraction}} & \thepage & Here E can be an arbitrary nonempty set and B is an operator; B later also labels the bad maintenance state. \\
\phantomsection\label{entry:symbol-registry-01661}\(m,S_k,g_k,h_k,C,b,F\) & \texttt{\detokenize{m,}\allowbreak\detokenize{S_k,}\allowbreak\detokenize{g_k,}\allowbreak\detokenize{h_k,}\allowbreak\detokenize{C,}\allowbreak\detokenize{b,}\allowbreak\detokenize{F}} & Bellman transformation count state transformations immediate rewards continuation factors bounds and solution & P17 sequential-decisions & meta & \hyperref[thm:p17-bellman-bounded-transformations]{\ref*{thm:p17-bellman-bounded-transformations}} & \thepage & Original bounded-transformation theorem uses arbitrary nonempty domain; g\_k differs from the finite terminal reward g. \\
\phantomsection\label{entry:symbol-registry-01662}\(\delta(v),\eta,\epsilon\) & \texttt{\detokenize{\delta(v),}\allowbreak\detokenize{\eta,}\allowbreak\detokenize{\epsilon}} & Bellman residual approximate greedy tolerance and requested error tolerance & P17 sequential-decisions & meta & \hyperref[thm:p17-sequential-value-iteration]{\ref*{thm:p17-sequential-value-iteration}} & \thepage & The residual is a norm and is unrelated to a Dirac initial law. \\
\phantomsection\label{entry:symbol-registry-01663}\(P_f,r_f,B_f\) & \texttt{\detokenize{P_f,}\allowbreak\detokenize{r_f,}\allowbreak\detokenize{B_f}} & policy transition matrix expected reward vector and Neumann inverse & P17 sequential-decisions & meta & \hyperref[thm:p17-sequential-policy-evaluation]{\ref*{thm:p17-sequential-policy-evaluation}} & \thepage & B\_f is the sum of discounted transition powers and the inverse of I-gamma P\_f; finite rational elimination is supplied explicitly. \\
\phantomsection\label{entry:symbol-registry-01664}\(d,h\) & \texttt{\detokenize{d,}\allowbreak\detokenize{h}} & improvement quantity and Eaton-Zadeh switching policy & P17 sequential-decisions & meta & \hyperref[thm:p17-sequential-improvement]{\ref*{thm:p17-sequential-improvement}} & \thepage & These roles are local to the improvement construction; h is not a history h\_t. \\
\phantomsection\label{entry:symbol-registry-01665}\(G,B\) & \texttt{\detokenize{G,}\allowbreak\detokenize{B}} & good and bad maintenance states & P17 sequential-decisions & meta & \hyperref[def:p17-sequential-maintenance]{\ref*{def:p17-sequential-maintenance}} & \thepage & The worked example's state labels are separate from the abstract contraction operator. \\
\phantomsection\label{entry:symbol-registry-01666}\(\beta,U,\kappa_s,L_a\) & \texttt{\detokenize{\beta,}\allowbreak\detokenize{U,}\allowbreak\detokenize{\kappa_s,}\allowbreak\detokenize{L_a}} & Blackwell source discount generated operator action-erasure map and action functional & P17 sequential-decisions & meta & \hyperref[thm:p17-sequential-source-bridge]{\ref*{thm:p17-sequential-source-bridge}} & \thepage & The source-time and generated-action bridges are explicit; equal operators require expected rewards rather than equal realized reward functions. \\
\phantomsection\label{entry:symbol-registry-01680}\(S,Z,\Delta(Z),f,g,h,f_Eh\) & \texttt{\detokenize{S,}\allowbreak\detokenize{Z,}\allowbreak\detokenize{\Delta(Z),}\allowbreak\detokenize{f,}\allowbreak\detokenize{g,}\allowbreak\detokenize{h,}\allowbreak\detokenize{f_Eh}} & finite states consequences lotteries acts and event-pasted acts & P17 ambiguity-and-robust-choice & meta & \hyperref[def:p17-ambiguity-acts]{\ref*{def:p17-ambiguity-acts}} & \thepage & The full act domain contains every statewise consequence lottery; the pasting event is a subset of S. \\
\phantomsection\label{entry:symbol-registry-01681}\(u,x_f,\widehat u,p,E_p,H,L\) & \texttt{\detokenize{u,}\allowbreak\detokenize{x_f,}\allowbreak\detokenize{\widehat u,}\allowbreak\detokenize{p,}\allowbreak\detokenize{E_p,}\allowbreak\detokenize{H,}\allowbreak\detokenize{L}} & consequence utility state utility affine lottery utility prior expectation and ordered prizes & P17 ambiguity-and-robust-choice & meta & \hyperref[def:p17-ambiguity-acts]{\ref*{def:p17-ambiguity-acts}} & \thepage & One common nonconstant utility and one common prior govern the additive comparison; utility need not be linear in money. \\
\phantomsection\label{entry:symbol-registry-01682}\(R_{\rm I},B_{\rm I},R_{\rm II},B_{\rm II}\) & \texttt{\detokenize{R_{\rm I},}\allowbreak\detokenize{B_{\rm I},}\allowbreak\detokenize{R_{\rm II},}\allowbreak\detokenize{B_{\rm II}}} & original two-urn color bets and corresponding marked events & P17 ambiguity-and-robust-choice & meta & \hyperref[def:p17-ellsberg-two-urns]{\ref*{def:p17-ellsberg-two-urns}} & \thepage & The four paired menus differ from a menu allowing free joint selection of urn and color. \\
\phantomsection\label{entry:symbol-registry-01683}\(M_1,\ldots,M_6,A_1,\ldots,A_4,F_1,\ldots,F_4,R,B,Y\) & \texttt{\detokenize{M_1,}\allowbreak\detokenize{\ldots,}\allowbreak\detokenize{M_6,}\allowbreak\detokenize{A_1,}\allowbreak\detokenize{\ldots,}\allowbreak\detokenize{A_4,}\allowbreak\detokenize{F_1,}\allowbreak\detokenize{\ldots,}\allowbreak\detokenize{F_4,}\allowbreak\detokenize{R,}\allowbreak\detokenize{B,}\allowbreak\detokenize{Y}} & marked-urn acts Arrow acts three-color acts and consistent color coordinates & P17 ambiguity-and-robust-choice & meta & \hyperref[thm:p17-ellsberg-marked-arrow]{\ref*{thm:p17-ellsberg-marked-arrow}} & \thepage & Each act table has its stated experiment and state partition; the later source color permutation is made explicit. \\
\phantomsection\label{entry:symbol-registry-01684}\(\mathcal C,\mathcal M,V_{\mathcal C},\|\cdot\|_\infty\) & \texttt{\detokenize{\mathcal C,}\allowbreak\detokenize{\mathcal M,}\allowbreak\detokenize{V_{\mathcal C},}\allowbreak\detokenize{\|\cdot\|_\infty}} & supplied finite prior set menu maxmin value and maximum norm & P17 ambiguity-and-robust-choice & meta & \hyperref[thm:p17-finite-maxmin-properties]{\ref*{thm:p17-finite-maxmin-properties}} & \thepage & The finite method's forward properties do not establish the general representation converse. \\
\phantomsection\label{entry:symbol-registry-01685}\(\lambda,\alpha,\beta,\delta_i\) & \texttt{\detokenize{\lambda,}\allowbreak\detokenize{\alpha,}\allowbreak\detokenize{\beta,}\allowbreak\detokenize{\delta_i}} & lottery-mixture weight continuity bracketing weights and point priors & P17 ambiguity-and-robust-choice & meta & \hyperref[thm:p17-maxmin-mixtures]{\ref*{thm:p17-maxmin-mixtures}} & \thepage & Certainty independence differs from full mixture independence and from independence of random events. \\
\phantomsection\label{entry:symbol-registry-01686}\(\mathcal C_2,\mathcal C_3,\theta,d,a,b,\xi,\zeta\) & \texttt{\detokenize{\mathcal C_2,}\allowbreak\detokenize{\mathcal C_3,}\allowbreak\detokenize{\theta,}\allowbreak\detokenize{d,}\allowbreak\detokenize{a,}\allowbreak\detokenize{b,}\allowbreak\detokenize{\xi,}\allowbreak\detokenize{\zeta}} & urn prior sets composition parameter nonidentification parameter and auxiliary coupling coordinates & P17 ambiguity-and-robust-choice & meta & \hyperref[thm:p17-maxmin-urns]{\ref*{thm:p17-maxmin-urns}} & \thepage & Convex prior descriptions are distinguished from the finite physical ball compositions; endpoint reductions are proved. \\
\phantomsection\label{entry:symbol-registry-01687}\(w_j,\bar p,p_0,\rho,W_\rho,\mathcal C_\rho,D_u\) & \texttt{\detokenize{w_j,}\allowbreak\detokenize{\bar p,}\allowbreak\detokenize{p_0,}\allowbreak\detokenize{\rho,}\allowbreak\detokenize{W_\rho,}\allowbreak\detokenize{\mathcal C_\rho,}\allowbreak\detokenize{D_u}} & prior-averaging weights averaged prior best estimate confidence confidence-weighted value contracted prior set and utility span & P17 ambiguity-and-robust-choice & meta & \hyperref[thm:p17-ambiguity-randomization-averaging]{\ref*{thm:p17-ambiguity-randomization-averaging}} & \thepage & Source confidence is a specified preference input and is not an inferred statistical confidence level. \\
\phantomsection\label{entry:symbol-registry-01688}\(\operatorname{Reg}_{\mathcal M}\) & \texttt{\detokenize{\operatorname{Reg}_{\mathcal M}}} & maximum regret relative to a specified fixed menu & P17 ambiguity-and-robust-choice & meta & \hyperref[thm:p17-ambiguity-other-criteria]{\ref*{thm:p17-ambiguity-other-criteria}} & \thepage & Fixed-menu regret and pairwise-menu regret have different comparison rules. \\
\phantomsection\label{entry:symbol-registry-01712}\(x_0,m,D,D(x),\mathbf z,a,b,c,T,F_n,V,h,H,\eta\) & \texttt{\detokenize{x_0,}\allowbreak\detokenize{m,}\allowbreak\detokenize{D,}\allowbreak\detokenize{D(x),}\allowbreak\detokenize{\mathbf z,}\allowbreak\detokenize{a,}\allowbreak\detokenize{b,}\allowbreak\detokenize{c,}\allowbreak\detokenize{T,}\allowbreak\detokenize{F_n,}\allowbreak\detokenize{V,}\allowbreak\detokenize{h,}\allowbreak\detokenize{H,}\allowbreak\detokenize{\eta}} & shrinking-resource domain allocations reward successor contraction factor operator finite values canonical value and envelopes & P17 bellman-original-constructions & meta & \hyperref[def:p17-bellman-shrinking-model]{\ref*{def:p17-bellman-shrinking-model}} & \thepage & h is the maximum at exactly one resource level; H and eta control all smaller levels. The original grid condition is insufficient without that control. \\
\phantomsection\label{entry:symbol-registry-01713}\(r_i,d_i,K,(t)_+\) & \texttt{\detokenize{r_i,}\allowbreak\detokenize{d_i,}\allowbreak\detokenize{K,}\allowbreak\detokenize{(t)_+}} & linear allocation rewards successor fractions optimal coefficient and positive-part bump operation & P17 bellman-original-constructions & meta & \hyperref[thm:p17-bellman-shrinking]{\ref*{thm:p17-bellman-shrinking}} & \thepage & K is local to the linear allocation formula; the bump construction supplies the continuous counterexample. \\
\phantomsection\label{entry:symbol-registry-01714}\(\Delta_d,p_i,e_i,S_l,c_l,\tau,K\) & \texttt{\detokenize{\Delta_d,}\allowbreak\detokenize{p_i,}\allowbreak\detokenize{e_i,}\allowbreak\detokenize{S_l,}\allowbreak\detokenize{c_l,}\allowbreak\detokenize{\tau,}\allowbreak\detokenize{K}} & probability simplex coordinates basis vectors transformations mass-reduction factors absorption time and proper-policy bound & P17 bellman-original-constructions & meta & \hyperref[def:p17-bellman-probability-stopping]{\ref*{def:p17-bellman-probability-stopping}} & \thepage & Coordinates run from zero through d and e\_0 is absorbing. The source does not assume a global time discount or continuous transformations. \\
\phantomsection\label{entry:symbol-registry-01715}\(p_1,p_2,s_1,s_2,r_1,r_2,x,y,I_1,I_2,\bar p\) & \texttt{\detokenize{p_1,}\allowbreak\detokenize{p_2,}\allowbreak\detokenize{s_1,}\allowbreak\detokenize{s_2,}\allowbreak\detokenize{r_1,}\allowbreak\detokenize{r_2,}\allowbreak\detokenize{x,}\allowbreak\detokenize{y,}\allowbreak\detokenize{I_1,}\allowbreak\detokenize{I_2,}\allowbreak\detokenize{\bar p}} & two-resource continuation multipliers shrink factors reward coefficients resources indices and multiplier bound & P17 bellman-original-constructions & meta & \hyperref[thm:p17-bellman-two-resources]{\ref*{thm:p17-bellman-two-resources}} & \thepage & The bounded-function construction is applied on invariant rectangles and joined across the whole quadrant. \\
\phantomsection\label{entry:symbol-registry-01716}\(w,N_i(j),n_i(w),D_w,A_w,B_w\) & \texttt{\detokenize{w,}\allowbreak\detokenize{N_i(j),}\allowbreak\detokenize{n_i(w),}\allowbreak\detokenize{D_w,}\allowbreak\detokenize{A_w,}\allowbreak\detokenize{B_w}} & action word preceding action counts total counts prefix multiplier and linear coefficients & P17 bellman-original-constructions & meta & \hyperref[thm:p17-bellman-two-resources]{\ref*{thm:p17-bellman-two-resources}} & \thepage & Finite prefixes and infinite word limits describe the actual continuous resource dynamics. \\
\phantomsection\label{entry:symbol-registry-01717}\(q,u,v,k,\theta,z\) & \texttt{\detokenize{q,}\allowbreak\detokenize{u,}\allowbreak\detokenize{v,}\allowbreak\detokenize{k,}\allowbreak\detokenize{\theta,}\allowbreak\detokenize{z}} & common commensurate scale coprime exponents cone index fractional exponent and repeating action cycle & P17 bellman-original-constructions & meta & \hyperref[thm:p17-bellman-commensurate]{\ref*{thm:p17-bellman-commensurate}} & \thepage & The linear cones are countable and locally finite away from the axes; z is an action word rather than an allocation vector. \\
\phantomsection\label{entry:symbol-registry-01718}\(a,b,g,h,\lambda,A,B,K,c\) & \texttt{\detokenize{a,}\allowbreak\detokenize{b,}\allowbreak\detokenize{g,}\allowbreak\detokenize{h,}\allowbreak\detokenize{\lambda,}\allowbreak\detokenize{A,}\allowbreak\detokenize{B,}\allowbreak\detokenize{K,}\allowbreak\detokenize{c}} & split fractions reward functions power exponent coefficients and local bounds & P17 bellman-original-constructions & meta & \hyperref[thm:p17-bellman-convex-split]{\ref*{thm:p17-bellman-convex-split}} & \thepage & Convexity supplies endpoint reduction and a geometric envelope; concave rewards have a separate solved counterexample. \\
\phantomsection\label{entry:symbol-registry-01719}\(p,e_i,t_i,q_i,c_i,p^{(-i)},\sigma,F\) & \texttt{\detokenize{p,}\allowbreak\detokenize{e_i,}\allowbreak\detokenize{t_i,}\allowbreak\detokenize{q_i,}\allowbreak\detokenize{c_i,}\allowbreak\detokenize{p^{(-i)},}\allowbreak\detokenize{\sigma,}\allowbreak\detokenize{F}} & box prior unit prior attempt cost failure probability completed cost conditional prior inspection order and retrieval value & P17 bellman-original-constructions & meta & \hyperref[thm:p17-bellman-box-retrieval]{\ref*{thm:p17-bellman-box-retrieval}} & \thepage & Box coordinates start at one. Failure probabilities are conditional on past attempts and the ball location in the repeated-attempt interpretation. \\
\phantomsection\label{entry:symbol-registry-01720}\(L_l,L\) & \texttt{\detokenize{L_l,}\allowbreak\detokenize{L}} & cost with one specified uninspected box and optimal location cost & P17 bellman-original-constructions & meta & \hyperref[thm:p17-bellman-box-location]{\ref*{thm:p17-bellman-box-location}} & \thepage & Locating the ball and obtaining it have distinct boundary values and optimal policies. \\
\phantomsection\label{entry:symbol-registry-01721}\(a,f,C,k,\tau,n(x),g_n(x)\) & \texttt{\detokenize{a,}\allowbreak\detokenize{f,}\allowbreak\detokenize{C,}\allowbreak\detokenize{k,}\allowbreak\detokenize{\tau,}\allowbreak\detokenize{n(x),}\allowbreak\detokenize{g_n(x)}} & scalar contraction solution value at one minimizing integer threshold first crossing and affine candidates & P17 bellman-original-constructions & meta & \hyperref[thm:p17-bellman-threshold]{\ref*{thm:p17-bellman-threshold}} & \thepage & The absorbing value at zero is discontinuous from positive states; all minimizing-integer ties are retained. \\
\phantomsection\label{entry:symbol-registry-01741}\(S,X,\Sigma,Y=\Delta_f(X)\) & \texttt{\detokenize{S,}\allowbreak\detokenize{X,}\allowbreak\detokenize{\Sigma,}\allowbreak\detokenize{Y,}\allowbreak\detokenize{\Delta_f}} & arbitrary nonempty states outcomes state algebra and finite-support lotteries & P17 maxmin-representation & meta & \hyperref[def:p17-mmeu-domain]{\ref*{def:p17-mmeu-domain}} & \thepage & Sigma is an algebra; countable closure is not assumed. \\
\phantomsection\label{entry:symbol-registry-01742}\(L_0,L_c,y^*,f,g,h,\succeq,\succ,\sim\) & \texttt{\detokenize{L_0,}\allowbreak\detokenize{L_c,}\allowbreak\detokenize{y^*,}\allowbreak\detokenize{f,}\allowbreak\detokenize{g,}\allowbreak\detokenize{h,}\allowbreak\detokenize{\succeq,}\allowbreak\detokenize{\succ,}\allowbreak\detokenize{\sim}} & finite-step acts constant acts and preference relations & P17 maxmin-representation & meta & \hyperref[def:p17-mmeu-domain]{\ref*{def:p17-mmeu-domain}} & \thepage & A.1--A.6 retain the source strict independence and mixture continuity clauses. \\
\phantomsection\label{entry:symbol-registry-01743}\(v,u,K,J\) & \texttt{\detokenize{v,}\allowbreak\detokenize{u,}\allowbreak\detokenize{K,}\allowbreak\detokenize{J}} & consequence utility affine lottery utility its range and certainty-equivalent value & P17 maxmin-representation & meta & \hyperref[thm:p17-mmeu-utility]{\ref*{thm:p17-mmeu-utility}} & \thepage & J is introduced at thm:p17-mmeu-certainty. The outcome utility need not be bounded. \\
\phantomsection\label{entry:symbol-registry-01744}\(B_0,B,\|\cdot\|_\infty,I,c\mathbf1\) & \texttt{\detokenize{B_0,}\allowbreak\detokenize{B,}\allowbreak\detokenize{\|\cdot\|_\infty,}\allowbreak\detokenize{I,}\allowbreak\detokenize{c\mathbf1}} & simple functions uniform completion norm lower expectation and constant function & P17 maxmin-representation & meta & \hyperref[def:p17-mmeu-function-spaces]{\ref*{def:p17-mmeu-function-spaces}} & \thepage & I is constructed at thm:p17-mmeu-simple-functional without replacing Sigma by its generated sigma algebra. \\
\phantomsection\label{entry:symbol-registry-01745}\(E,p,\ell,L_b\) & \texttt{\detokenize{E,}\allowbreak\detokenize{p,}\allowbreak\detokenize{\ell,}\allowbreak\detokenize{L_b}} & real vector space sublinear functional and dominated or supporting linear functionals & P17 maxmin-representation & meta & \hyperref[thm:p17-mmeu-hahn-banach]{\ref*{thm:p17-mmeu-hahn-banach}} & \thepage & Local linear-extension roles; p is not a prior in this proof. \\
\phantomsection\label{entry:symbol-registry-01746}\(P,\mathcal P(\Sigma),\mathcal U\) & \texttt{\detokenize{P,}\allowbreak\detokenize{\mathcal P(\Sigma),}\allowbreak\detokenize{\mathcal U}} & finitely additive probability prior space and local ultrafilter & P17 maxmin-representation & meta & \hyperref[def:p17-mmeu-function-spaces]{\ref*{def:p17-mmeu-function-spaces}} & \thepage & Integration and compactness are constructed in annex:p17.mmeu.charges.compactness. \\
\phantomsection\label{entry:symbol-registry-01747}\(D_1,D_2,P_b,C\) & \texttt{\detokenize{D_1,}\allowbreak\detokenize{D_2,}\allowbreak\detokenize{P_b,}\allowbreak\detokenize{C}} & separating convex sets touching priors and their closed convex hull & P17 maxmin-representation & meta & \hyperref[thm:p17-mmeu-support]{\ref*{thm:p17-mmeu-support}} & \thepage & C also has the intrinsic dominance characterization proved there. \\
\phantomsection\label{entry:symbol-registry-01748}\(L_{\rm lev},L_u\) & \texttt{\detokenize{L_{\rm lev},}\allowbreak\detokenize{L_u}} & original bounded level-measurable acts and convex uniform-utility completion & P17 maxmin-representation & meta & \hyperref[def:p17-mmeu-extension-domain]{\ref*{def:p17-mmeu-extension-domain}} & \thepage & All original acts are retained. Equality of these domains uses a sigma algebra. \\
\phantomsection\label{entry:symbol-registry-01749}\(\operatorname{Ext},K,F\) & \texttt{\detokenize{\operatorname{Ext},}\allowbreak\detokenize{K,}\allowbreak\detokenize{F}} & extreme points compact convex set and its face & P17 maxmin-representation & meta & \hyperref[thm:p17-mmeu-extreme-geometry]{\ref*{thm:p17-mmeu-extreme-geometry}} & \thepage & Local convex-geometry K is distinct from the earlier utility range; exposed-point construction follows in thm:p17-mmeu-exposed-points. \\
\phantomsection\label{entry:symbol-registry-01750}\(S_i,\Sigma_i,C_i,P_1\otimes P_2,\Sigma_1\otimes\Sigma_2\) & \texttt{\detokenize{S_i,}\allowbreak\detokenize{\Sigma_i,}\allowbreak\detokenize{C_i,}\allowbreak\detokenize{P_1\otimes P_2,}\allowbreak\detokenize{\Sigma_1\otimes\Sigma_2}} & component states algebras prior sets product charge and rectangle algebra & P17 maxmin-representation & meta & \hyperref[def:p17-mmeu-product]{\ref*{def:p17-mmeu-product}} & \thepage & Both the common-minimizer and extreme-prior independence clauses are retained. \\
\phantomsection\label{entry:symbol-registry-01751}\(R_i,C',\Sigma_i^0\) & \texttt{\detokenize{R_i,}\allowbreak\detokenize{C',}\allowbreak\detokenize{\Sigma_i^0}} & marginal projection candidate joint prior set and finite subalgebra & P17 maxmin-representation & meta & \hyperref[thm:p17-mmeu-product-priors]{\ref*{thm:p17-mmeu-product-priors}} & \thepage & Marginal maps start in annex:p17.mmeu.product.priors; projection lifts and common minimizing faces are constructed in annex:p17.mmeu.product.representation. \\
\phantomsection\label{entry:symbol-registry-01752}\(X,Y,M(X),P(X),Q(Y\mid X),k,kv\) & \texttt{\detokenize{X,}\allowbreak\detokenize{Y,}\allowbreak\detokenize{M(X),}\allowbreak\detokenize{P(X),}\allowbreak\detokenize{Q(Y\mid X),}\allowbreak\detokenize{k,}\allowbreak\detokenize{kv}} & standard Borel spaces bounded Borel functions probabilities kernels and kernel integration & P17 blackwell-general-discounted & meta & \hyperref[def:p17-blackwell-borel-model]{\ref*{def:p17-blackwell-borel-model}} & \thepage & The source Baire terminology is translated to its actual Borel sigma algebra. \\
\phantomsection\label{entry:symbol-registry-01753}\(S,A,q,r,R,\beta,H_n,h_n,S_n,A_n\) & \texttt{\detokenize{S,}\allowbreak\detokenize{A,}\allowbreak\detokenize{q,}\allowbreak\detokenize{r,}\allowbreak\detokenize{R,}\allowbreak\detokenize{\beta,}\allowbreak\detokenize{H_n,}\allowbreak\detokenize{h_n,}\allowbreak\detokenize{S_n,}\allowbreak\detokenize{A_n}} & decision spaces transition reward reward bound discount histories and random coordinates & P17 blackwell-general-discounted & meta & \hyperref[def:p17-blackwell-borel-model]{\ref*{def:p17-blackwell-borel-model}} & \thepage & All kernels are specified at null histories as well as reachable ones. \\
\phantomsection\label{entry:symbol-registry-01754}\(\pi_n,f_n,f^\infty,K^\pi,I^\pi,\rho\) & \texttt{\detokenize{\pi_n,}\allowbreak\detokenize{f_n,}\allowbreak\detokenize{f^\infty,}\allowbreak\detokenize{K^\pi,}\allowbreak\detokenize{I^\pi,}\allowbreak\detokenize{\rho}} & history policies Markov maps stationary policy future law return and initial prior & P17 blackwell-general-discounted & meta & \hyperref[def:p17-blackwell-borel-model]{\ref*{def:p17-blackwell-borel-model}} & \thepage & Future law and return are constructed at thm:p17-blackwell-factorization; prior quantifiers are defined at def:p17-blackwell-optimality. \\
\phantomsection\label{entry:symbol-registry-01755}\(D,D_x,p,\tau_n,\Delta,d,d_O,\nu_x,F_x,O_n,c_n,z,f\) & \texttt{\detokenize{D,}\allowbreak\detokenize{D_x,}\allowbreak\detokenize{p,}\allowbreak\detokenize{\tau_n,}\allowbreak\detokenize{\Delta,}\allowbreak\detokenize{d,}\allowbreak\detokenize{d_O,}\allowbreak\detokenize{\nu_x,}\allowbreak\detokenize{F_x,}\allowbreak\detokenize{O_n,}\allowbreak\detokenize{c_n,}\allowbreak\detokenize{z,}\allowbreak\detokenize{f}} & Borel relation sections projection refined topologies diagonal metrics fiber measures and nested-ball selector data & P17 blackwell-general-discounted & meta & \hyperref[thm:p17-blackwell-positive-selection]{\ref*{thm:p17-blackwell-positive-selection}} & \thepage & Local scope annex:p17.blackwell.positive.selection; the projection p is a map. \\
\phantomsection\label{entry:symbol-registry-01756}\(a(x),m(x),\eta\) & \texttt{\detokenize{a(x),}\allowbreak\detokenize{m(x),}\allowbreak\detokenize{\eta}} & conditional mean essential supremum and strict selection tolerance & P17 blackwell-general-discounted & meta & \hyperref[thm:p17-blackwell-selection]{\ref*{thm:p17-blackwell-selection}} & \thepage & Local a is a real function and is distinct from a decision action. \\
\phantomsection\label{entry:symbol-registry-01757}\(\varepsilon,N,\kappa_i,V_{i+1},w_i\) & \texttt{\detokenize{\varepsilon,}\allowbreak\detokenize{N,}\allowbreak\detokenize{\kappa_i,}\allowbreak\detokenize{V_{i+1},}\allowbreak\detokenize{w_i}} & optimality tolerance truncation conditional action kernel continuation and discounted local return & P17 blackwell-general-discounted & meta & \hyperref[def:p17-blackwell-optimality]{\ref*{def:p17-blackwell-optimality}} & \thepage & Backward replacements and their prior-dependent conditional laws are proved in annex:p17.blackwell.prior.markov. \\
\phantomsection\label{entry:symbol-registry-01758}\(\mathcal N,e,\mathcal F,U,E,B_0,B,\phi,J\) & \texttt{\detokenize{\mathcal N,}\allowbreak\detokenize{e,}\allowbreak\detokenize{\mathcal F,}\allowbreak\detokenize{U,}\allowbreak\detokenize{E,}\allowbreak\detokenize{B_0,}\allowbreak\detokenize{B,}\allowbreak\detokenize{\phi,}\allowbreak\detokenize{J}} & Baire space prefix enumeration universal relation its projection diagonal and Borel coding data & P17 blackwell-general-discounted & meta & \hyperref[thm:p17-blackwell-projection-boundaries]{\ref*{thm:p17-blackwell-projection-boundaries}} & \thepage & Local scope annex:p17.blackwell.projection.boundaries. U and B here are set relations rather than value operators. \\
\phantomsection\label{entry:symbol-registry-01759}\(T_f,T_a,(f_n),E_n,U,u^*,B,u,d_i,\widehat q,\widehat r\) & \texttt{\detokenize{T_f,}\allowbreak\detokenize{T_a,}\allowbreak\detokenize{(f_n),}\allowbreak\detokenize{E_n,}\allowbreak\detokenize{U,}\allowbreak\detokenize{u^*,}\allowbreak\detokenize{B,}\allowbreak\detokenize{u,}\allowbreak\detokenize{d_i,}\allowbreak\detokenize{\widehat q,}\allowbreak\detokenize{\widehat r}} & stationary and generated operators choices partitions fixed point generic contraction terminal payoff losses and index-model data & P17 blackwell-general-discounted & meta & \hyperref[def:p17-blackwell-generated]{\ref*{def:p17-blackwell-generated}} & \thepage & Local generated-operator roles are distinct from the counterexample relations. \\
\phantomsection\label{entry:symbol-registry-01760}\((g,\pi),\pi^*_{s,a},V,v_n\) & \texttt{\detokenize{(g,}\allowbreak\detokenize{\pi),}\allowbreak\detokenize{\pi^*_{s,a},}\allowbreak\detokenize{V,}\allowbreak\detokenize{v_n}} & restarted policy parameterized continuation optimal return and approximating values & P17 blackwell-general-discounted & meta & \hyperref[thm:p17-blackwell-generated-operators]{\ref*{thm:p17-blackwell-generated-operators}} & \thepage & Restart semantics and jointly measurable continuation kernels are explicit in the generated and stationary-prior proofs. \\
\phantomsection\label{entry:symbol-registry-01761}\(\bar r_a,E_m,B_n,h\) & \texttt{\detokenize{\bar r_a,}\allowbreak\detokenize{E_m,}\allowbreak\detokenize{B_n,}\allowbreak\detokenize{h}} & expected immediate reward essential finite cover pieces maximizing pieces and switching map & P17 blackwell-general-discounted & meta & \hyperref[def:p17-blackwell-essential-actions]{\ref*{def:p17-blackwell-essential-actions}} & \thepage & Strong source equivalence and operator equivalence have separate hypotheses; h is introduced in the improvement result. \\
\phantomsection\label{entry:symbol-registry-01762}\(X,X_t,Y,Y_t,N,n,m,d,U,U_t,c_t,f_t,\pi_t\) & \texttt{\detokenize{X,}\allowbreak\detokenize{X_t,}\allowbreak\detokenize{Y,}\allowbreak\detokenize{Y_t,}\allowbreak\detokenize{N,}\allowbreak\detokenize{n,}\allowbreak\detokenize{m,}\allowbreak\detokenize{d,}\allowbreak\detokenize{U,}\allowbreak\detokenize{U_t,}\allowbreak\detokenize{c_t,}\allowbreak\detokenize{f_t,}\allowbreak\detokenize{\pi_t}} & finite hidden states observations horizon dimensions compact controls and policy objects & P17 partial-observation-decisions & meta & \hyperref[def:p17-astrom-model]{\ref*{def:p17-astrom-model}} & \thepage & The control set can be continuous; the belief state is the whole simplex. \\
\phantomsection\label{entry:symbol-registry-01763}\(P_t,p_{si,t},Q,q_{ij},g_t,J^\pi,\mu,p^0,u_0,P_0\) & \texttt{\detokenize{P_t,}\allowbreak\detokenize{p_{si,t},}\allowbreak\detokenize{Q,}\allowbreak\detokenize{q_{ij},}\allowbreak\detokenize{g_t,}\allowbreak\detokenize{J^\pi,}\allowbreak\detokenize{\mu,}\allowbreak\detokenize{p^0,}\allowbreak\detokenize{u_0,}\allowbreak\detokenize{P_0}} & transition and observation matrices costs expected cost and initialization data & P17 partial-observation-decisions & meta & \hyperref[def:p17-astrom-model]{\ref*{def:p17-astrom-model}} & \thepage & Departure-indexed transitions determine the initial observation law explicitly. \\
\phantomsection\label{entry:symbol-registry-01764}\(H_t,\mathcal I_t,\Delta_n,B_t,b,w,b^\dagger\) & \texttt{\detokenize{H_t,}\allowbreak\detokenize{\mathcal I_t,}\allowbreak\detokenize{\Delta_n,}\allowbreak\detokenize{B_t,}\allowbreak\detokenize{b,}\allowbreak\detokenize{w,}\allowbreak\detokenize{b^\dagger}} & observed history information sigma algebra simplex belief source alias and null-branch default & P17 partial-observation-decisions & meta & \hyperref[def:p17-astrom-model]{\ref*{def:p17-astrom-model}} & \thepage & Beliefs and the source w alias are defined at def:p17-astrom-belief. \\
\phantomsection\label{entry:symbol-registry-01765}\(a_i,z^j,\lambda_j,F_j,\lambda_j^0,F_j^0,\mathcal R_t\) & \texttt{\detokenize{a_i,}\allowbreak\detokenize{z^j,}\allowbreak\detokenize{\lambda_j,}\allowbreak\detokenize{F_j,}\allowbreak\detokenize{\lambda_j^0,}\allowbreak\detokenize{F_j^0,}\allowbreak\detokenize{\mathcal R_t}} & prediction unnormalized observation columns likelihood posterior initial filter and belief kernel & P17 partial-observation-decisions & meta & \hyperref[def:p17-astrom-belief]{\ref*{def:p17-astrom-belief}} & \thepage & Zero observation mass contributes zero regardless of the default posterior. \\
\phantomsection\label{entry:symbol-registry-01766}\(\ell_t,\mathcal H_t,V_t,\delta_t,W_t^\pi,\eta_t\) & \texttt{\detokenize{\ell_t,}\allowbreak\detokenize{\mathcal H_t,}\allowbreak\detokenize{V_t,}\allowbreak\detokenize{\delta_t,}\allowbreak\detokenize{W_t^\pi,}\allowbreak\detokenize{\eta_t}} & expected immediate cost Bellman objective value residual conditional future cost and local tolerance & P17 partial-observation-decisions & meta & \hyperref[thm:p17-astrom-information-state]{\ref*{thm:p17-astrom-information-state}} & \thepage & The recursion is defined at def:p17-astrom-recursion; calligraphic H is distinct from history H. \\
\phantomsection\label{entry:symbol-registry-01767}\(\kappa_t,\nu_t,C,C_0,h_C,H\) & \texttt{\detokenize{\kappa_t,}\allowbreak\detokenize{\nu_t,}\allowbreak\detokenize{C,}\allowbreak\detokenize{C_0,}\allowbreak\detokenize{h_C,}\allowbreak\detokenize{H}} & conditional action kernel belief marginal compact selector cubes restricted minimum and continuous objective & P17 partial-observation-decisions & meta & \hyperref[annex:p17.astrom.information.state]{\ref*{annex:p17.astrom.information.state}} & \thepage & The local compact-selector roles start in annex:p17.astrom.continuous.selection. \\
\phantomsection\label{entry:symbol-registry-01768}\(S_t,L_t,O_t,d_t,Z_t\) & \texttt{\detokenize{S_t,}\allowbreak\detokenize{L_t,}\allowbreak\detokenize{O_t,}\allowbreak\detokenize{d_t,}\allowbreak\detokenize{Z_t}} & perfect-state comparison values schedule value delivered prefix and augmented hidden path & P17 partial-observation-decisions & meta & \hyperref[def:p17-astrom-comparison-values]{\ref*{def:p17-astrom-comparison-values}} & \thepage & The delayed-observation prefix and path are defined at thm:p17-astrom-delays. \\
\phantomsection\label{entry:symbol-registry-01769}\(x,q,e,\alpha,\zeta,a,A,B,C,D,\Phi_2,\Phi_1,L,K,c\) & \texttt{\detokenize{x,}\allowbreak\detokenize{q,}\allowbreak\detokenize{e,}\allowbreak\detokenize{\alpha,}\allowbreak\detokenize{\zeta,}\allowbreak\detokenize{a,}\allowbreak\detokenize{A,}\allowbreak\detokenize{B,}\allowbreak\detokenize{C,}\allowbreak\detokenize{D,}\allowbreak\detokenize{\Phi_2,}\allowbreak\detokenize{\Phi_1,}\allowbreak\detokenize{L,}\allowbreak\detokenize{K,}\allowbreak\detokenize{c}} & binary-switch belief accuracy constants prediction observation branch entries and affine branch functions & P17 partial-observation-decisions & meta & \hyperref[def:p17-astrom-example-one]{\ref*{def:p17-astrom-example-one}} & \thepage & Local scope annex:p17.astrom.example.one; source q0 alpha0 beta0 aliases are explicit. \\
\phantomsection\label{entry:symbol-registry-01770}\(a_u,c_u,V_t,L_t,O_t,\mathcal H_t^{\rm open}\) & \texttt{\detokenize{a_u,}\allowbreak\detokenize{c_u,}\allowbreak\detokenize{V_t,}\allowbreak\detokenize{L_t,}\allowbreak\detokenize{O_t,}\allowbreak\detokenize{\mathcal H_t^{\rm open}}} & production prediction reward offset maximizing values and open-loop candidates & P17 partial-observation-decisions & meta & \hyperref[def:p17-astrom-example-two]{\ref*{def:p17-astrom-example-two}} & \thepage & These V L and O roles maximize rewards and reverse the earlier cost inequalities. \\
\phantomsection\label{entry:symbol-registry-01771}\(X,E,E_1,E_2,\Theta,\theta,p_{si,t}^{\theta}\) & \texttt{\detokenize{X,}\allowbreak\detokenize{E,}\allowbreak\detokenize{E_1,}\allowbreak\detokenize{E_2,}\allowbreak\detokenize{\Theta,}\allowbreak\detokenize{\theta,}\allowbreak\detokenize{p_{si,t}^{\theta}}} & boundary-example bits finite unknown parameter and parameter-indexed transition & P17 partial-observation-decisions & meta & \hyperref[thm:p17-astrom-information-failures]{\ref*{thm:p17-astrom-information-failures}} & \thepage & Local countermodels distinguish the one-time observation matrix from the full controlled joint law. \\
\phantomsection\label{entry:symbol-registry-01772}\(S=\mathbb R_+^n,\Delta_n,w,v,\|\cdot\|_1\) & \texttt{\detokenize{S,}\allowbreak\detokenize{\mathbb R_+^n,}\allowbreak\detokenize{\Delta_n,}\allowbreak\detokenize{w,}\allowbreak\detokenize{v,}\allowbreak\detokenize{\|\cdot\|_1}} & full nonnegative cone unit-mass simplex weight vectors and source norm & P17 astrom-incomplete-information-ii & meta & \hyperref[def:p17-astrom-ii-model]{\ref*{def:p17-astrom-ii-model}} & \thepage & Nonunit weights are not probability laws; the origin has value zero. \\
\phantomsection\label{entry:symbol-registry-01773}\(U,u,U_t,X_t,Y_t,H_t,P_t,Q_t,p_{si,t},q_{ij,t},g_t,G_t\) & \texttt{\detokenize{U,}\allowbreak\detokenize{u,}\allowbreak\detokenize{U_t,}\allowbreak\detokenize{X_t,}\allowbreak\detokenize{Y_t,}\allowbreak\detokenize{H_t,}\allowbreak\detokenize{P_t,}\allowbreak\detokenize{Q_t,}\allowbreak\detokenize{p_{si,t},}\allowbreak\detokenize{q_{ij,t},}\allowbreak\detokenize{g_t,}\allowbreak\detokenize{G_t}} & controls process history controlled transitions observations reward vector and bound & P17 astrom-incomplete-information-ii & meta & \hyperref[def:p17-astrom-ii-model]{\ref*{def:p17-astrom-ii-model}} & \thepage & The current control affects the next observation; Q continuity is an additional attainment condition. \\
\phantomsection\label{entry:symbol-registry-01774}\(A_{j,t}(u),V_t^+,V_t^-,\mathrm{ext},C_t,c\) & \texttt{\detokenize{A_{j,t}(u),}\allowbreak\detokenize{V_t^+,}\allowbreak\detokenize{V_t^-,}\allowbreak\detokenize{\mathrm{ext},}\allowbreak\detokenize{C_t,}\allowbreak\detokenize{c}} & positive observation maps maximum and minimum values extremum operator cumulative bound and mass scaling & P17 astrom-incomplete-information-ii & meta & \hyperref[def:p17-astrom-ii-model]{\ref*{def:p17-astrom-ii-model}} & \thepage & The source suppresses the control and time arguments; recursions are defined at def:p17-astrom-ii-recursions. \\
\phantomsection\label{entry:symbol-registry-01775}\(\mathcal P_A f,\lambda,\mu,\theta,a,b,c\) & \texttt{\detokenize{\mathcal P_A f,}\allowbreak\detokenize{\lambda,}\allowbreak\detokenize{\mu,}\allowbreak\detokenize{\theta,}\allowbreak\detokenize{a,}\allowbreak\detokenize{b,}\allowbreak\detokenize{c}} & perspective source mixture weights and local image masses & P17 astrom-incomplete-information-ii & meta & \hyperref[thm:p17-astrom-ii-perspective]{\ref*{thm:p17-astrom-ii-perspective}} & \thepage & Local proof c is total mixture mass; every zero-image branch is treated explicitly. \\
\phantomsection\label{entry:symbol-registry-01776}\(a_i^\pi,W_t^+,W_t^-,f_\alpha,F\) & \texttt{\detokenize{a_i^\pi,}\allowbreak\detokenize{W_t^+,}\allowbreak\detokenize{W_t^-,}\allowbreak\detokenize{f_\alpha,}\allowbreak\detokenize{F}} & fixed-policy coefficients policy extremal values and arbitrary convex-function family & P17 astrom-incomplete-information-ii & meta & \hyperref[annex:p17.astrom.ii.mass.bridge]{\ref*{annex:p17.astrom.ii.mass.bridge}} & \thepage & Policy extrema are established without presuming attainment; the convex-family roles occur in annex:p17.astrom.ii.maxima. \\
\phantomsection\label{entry:symbol-registry-01777}\(f_t,\widehat f_t,h(u)\) & \texttt{\detokenize{f_t,}\allowbreak\detokenize{\widehat f_t,}\allowbreak\detokenize{h(u)}} & Borel optimizer ray extension and boundary sensing-accuracy function & P17 astrom-incomplete-information-ii & meta & \hyperref[thm:p17-astrom-ii-continuity]{\ref*{thm:p17-astrom-ii-continuity}} & \thepage & The accuracy h is local to thm:p17-astrom-ii-boundaries and is distinct from a history. \\
\phantomsection\label{entry:symbol-registry-01852}\(I,C_i,C,C_{-i},u_i\) & \texttt{\detokenize{I,}\allowbreak\detokenize{C_i,}\allowbreak\detokenize{C,}\allowbreak\detokenize{C_{-i},}\allowbreak\detokenize{u_i}} & finite players choices joint profiles and common fixed payoffs & P17 epistemic-game-theory & meta & \hyperref[def:p17-egt-types]{\ref*{def:p17-egt-types}} & \thepage & One-player opponent products are singletons. \\
\phantomsection\label{entry:symbol-registry-01853}\(T_i,t_i,b_i,q_i,v_i,B_i\) & \texttt{\detokenize{T_i,}\allowbreak\detokenize{t_i,}\allowbreak\detokenize{b_i,}\allowbreak\detokenize{q_i,}\allowbreak\detokenize{v_i,}\allowbreak\detokenize{B_i}} & types subjective joint beliefs action marginals expected payoffs and optimal-choice sets & P17 epistemic-game-theory & meta & \hyperref[def:p17-egt-types]{\ref*{def:p17-egt-types}} & \thepage & Opponent conjectures may be correlated; types and choices remain distinct. \\
\phantomsection\label{entry:symbol-registry-01854}\(h_i^r,\mathcal H_i^r,\Delta_f,\pi_i^r\) & \texttt{\detokenize{h_i^r,}\allowbreak\detokenize{\mathcal H_i^r,}\allowbreak\detokenize{\Delta_f,}\allowbreak\detokenize{\pi_i^r}} & finite-level hierarchies their spaces finite-support probabilities and truncation maps & P17 epistemic-game-theory & meta & \hyperref[thm:p17-egt-hierarchies]{\ref*{thm:p17-egt-hierarchies}} & \thepage & Complete type-induced construction in annex:p17.egt.hierarchies; no terminal-type universality assertion. \\
\phantomsection\label{entry:symbol-registry-01855}\(H_i^r,Q_i^r,Q_i^\infty\) & \texttt{\detokenize{H_i^r,}\allowbreak\detokenize{Q_i^r,}\allowbreak\detokenize{Q_i^\infty}} & individual belief levels cumulative levels and common-belief types & P17 epistemic-game-theory & meta & \hyperref[def:p17-egt-common-belief]{\ref*{def:p17-egt-common-belief}} & \thepage & Individual levels need not be nested; cumulative intersections are explicit. \\
\phantomsection\label{entry:symbol-registry-01856}\((i,t_i),N\) & \texttt{\detokenize{(i,}\allowbreak\detokenize{t_i),}\allowbreak\detokenize{N}} & tagged support-graph vertex and total type count & P17 epistemic-game-theory & meta & \hyperref[thm:p17-egt-belief-algorithm]{\ref*{thm:p17-egt-belief-algorithm}} & \thepage & Bad vertices and their predecessors determine common belief. \\
\phantomsection\label{entry:symbol-registry-01857}\(R_i,R_i^k,R_i^\infty,x_i\) & \texttt{\detokenize{R_i,}\allowbreak\detokenize{R_i^k,}\allowbreak\detokenize{R_i^\infty,}\allowbreak\detokenize{x_i}} & restricted choices elimination stages final survivors and dominating mixture & P17 epistemic-game-theory & meta & \hyperref[def:p17-egt-reduction]{\ref*{def:p17-egt-reduction}} & \thepage & Simultaneous removal retains all original deviations in the epistemic interpretation. \\
\phantomsection\label{entry:symbol-registry-01858}\(D,p_i,v,z,M\) & \texttt{\detokenize{D,}\allowbreak\detokenize{p_i,}\allowbreak\detokenize{v,}\allowbreak\detokenize{z,}\allowbreak\detokenize{M}} & payoff-difference matrix belief witness minimax margin optimization variable and coefficient bound & P17 epistemic-game-theory & meta & \hyperref[thm:p17-egt-best-response-alternative]{\ref*{thm:p17-egt-best-response-alternative}} & \thepage & Local matrix computation; rational vertex enumeration is proved in the same annex item. \\
\phantomsection\label{entry:symbol-registry-01859}\(t_i^{m,a},p_i^{m,a},\tau_i,c_i^0\) & \texttt{\detokenize{t_i^{m,a},}\allowbreak\detokenize{p_i^{m,a},}\allowbreak\detokenize{\tau_i,}\allowbreak\detokenize{c_i^0}} & layered witness type supporting belief anchor type and anchor choice & P17 epistemic-game-theory & meta & \hyperref[annex:p17.egt.finite.depth]{\ref*{annex:p17.egt.finite.depth}} & \thepage & All witness beliefs are optimal against the full original choice set. \\
\phantomsection\label{entry:symbol-registry-01860}\(t_1^a,t_1^b,t_1^c,t_2^e,t_2^g,t_2^h\) & \texttt{\detokenize{t_1^a,}\allowbreak\detokenize{t_1^b,}\allowbreak\detokenize{t_1^c,}\allowbreak\detokenize{t_2^e,}\allowbreak\detokenize{t_2^g,}\allowbreak\detokenize{t_2^h}} & six types in the first source game & P17 epistemic-game-theory & meta & \hyperref[def:p17-egt-first-game]{\ref*{def:p17-egt-first-game}} & \thepage & Superscripts label the calculated unique optima in this example. \\
\phantomsection\label{entry:symbol-registry-01861}\(\sigma_i,x,y,z,p,q,r\) & \texttt{\detokenize{\sigma_i,}\allowbreak\detokenize{x,}\allowbreak\detokenize{y,}\allowbreak\detokenize{z,}\allowbreak\detokenize{p,}\allowbreak\detokenize{q,}\allowbreak\detokenize{r}} & equilibrium conjectures and local mixed-profile weights in the gap example & P17 epistemic-game-theory & meta & \hyperref[thm:p17-egt-nash-beliefs]{\ref*{thm:p17-egt-nash-beliefs}} & \thepage & The six scalar weights are introduced in annex:p17.egt.nash.gap; z there is a choice probability. \\
\phantomsection\label{entry:symbol-registry-01862}\(A,B,0,1\) & \texttt{\detokenize{A,}\allowbreak\detokenize{B,}\allowbreak\detokenize{0,}\allowbreak\detokenize{1}} & choices in the three-player correlation and disagreement countermodels & P17 epistemic-game-theory & meta & \hyperref[annex:p17.egt.many.player.boundaries]{\ref*{annex:p17.egt.many.player.boundaries}} & \thepage & Local model roles; the zero-payoff players' choices are all rational. \\
\phantomsection\label{entry:symbol-registry-01863}\(ae,af,be,bf,cg,ch,dg,dh,p_c,p_{dg},p_{dh}\) & \texttt{\detokenize{ae,}\allowbreak\detokenize{af,}\allowbreak\detokenize{be,}\allowbreak\detokenize{bf,}\allowbreak\detokenize{cg,}\allowbreak\detokenize{ch,}\allowbreak\detokenize{dg,}\allowbreak\detokenize{dh,}\allowbreak\detokenize{p_c,}\allowbreak\detokenize{p_{dg},}\allowbreak\detokenize{p_{dh}}} & contingent strategies and grouped conjecture probabilities in the four-move game & P17 epistemic-game-theory & meta & \hyperref[def:p17-egt-reny-game]{\ref*{def:p17-egt-reny-game}} & \thepage & Full strategies are distinguished from terminal outcomes; the conditional bound appears at thm:p17-egt-reny-reasoning. \\
\phantomsection\label{entry:symbol-registry-01875}\(N,n\) & \texttt{\detokenize{N,}\allowbreak\detokenize{n}} & finite player set and its size & P17 epistemic-nash-conditions & meta & \hyperref[def:p17-ab-system]{\ref*{def:p17-ab-system}} & \thepage & n is at least two for the conjecture-recovery theorems; the one-player boundary is explicit. \\
\phantomsection\label{entry:symbol-registry-01876}\(A_i,A,A^{-i}\) & \texttt{\detokenize{A_i,}\allowbreak\detokenize{A,}\allowbreak\detokenize{A^{-i}}} & finite action sets and products & P17 epistemic-nash-conditions & meta & \hyperref[def:p17-ab-system]{\ref*{def:p17-ab-system}} & \thepage & \(A_i\) is nonempty; superscript \(-i\) omits coordinate \(i\). \\
\phantomsection\label{entry:symbol-registry-01877}\(S_i,S,S^{-i}\) & \texttt{\detokenize{S_i,}\allowbreak\detokenize{S,}\allowbreak\detokenize{S^{-i}}} & measurable type spaces and state products & P17 epistemic-nash-conditions & meta & \hyperref[def:p17-ab-system]{\ref*{def:p17-ab-system}} & \thepage & S is the full Cartesian product including prior-null states. \\
\phantomsection\label{entry:symbol-registry-01878}\(\mathcal S_i,\mathcal S\) & \texttt{\detokenize{\mathcal S_i,}\allowbreak\detokenize{\mathcal S}} & type and product sigma-fields & P17 epistemic-nash-conditions & meta & \hyperref[def:p17-ab-system]{\ref*{def:p17-ab-system}} & \thepage & No standard-Borel or general disintegration-existence assumption is inserted. \\
\phantomsection\label{entry:symbol-registry-01879}\(s_i,t_i,s\) & \texttt{\detokenize{s_i,}\allowbreak\detokenize{t_i,}\allowbreak\detokenize{s}} & own types and full states & P17 epistemic-nash-conditions & meta & \hyperref[def:p17-ab-system]{\ref*{def:p17-ab-system}} & \thepage & Typed coordinates; actions and own payoffs are fixed by an own type. \\
\phantomsection\label{entry:symbol-registry-01880}\(a_i,a,a^{-i}\) & \texttt{\detokenize{a_i,}\allowbreak\detokenize{a,}\allowbreak\detokenize{a^{-i}}} & measurable actual-action maps & P17 epistemic-nash-conditions & meta & \hyperref[def:p17-ab-system]{\ref*{def:p17-ab-system}} & \thepage & Generic action arguments \(c_i\) and deviations \(d_i\) are separate from these maps. \\
\phantomsection\label{entry:symbol-registry-01881}\(g_i,g,g^*\) & \texttt{\detokenize{g_i,}\allowbreak\detokenize{g,}\allowbreak\detokenize{g^*}} & type-dependent payoff arrays and a specified game & P17 epistemic-nash-conditions & meta & \hyperref[def:p17-ab-system]{\ref*{def:p17-ab-system}} & \thepage & Each array has finitely many finite real coordinates; no prior-average payoff integrability is needed. \\
\phantomsection\label{entry:symbol-registry-01882}\(b_i\) & \texttt{\detokenize{b_i}} & opponent-type theory & P17 epistemic-nash-conditions & meta & \hyperref[def:p17-ab-system]{\ref*{def:p17-ab-system}} & \thepage & Probability on \(S^{-i}\); finite bridge separately uses predecessor \(T_i\) types. \\
\phantomsection\label{entry:symbol-registry-01883}\(p_i(E;s_i)\) & \texttt{\detokenize{p_i}} & extended own-coordinate probability kernel & P17 epistemic-nash-conditions & meta & \hyperref[def:p17-ab-system]{\ref*{def:p17-ab-system}} & \thepage & Defined at every type; event-wise measurability in own type is an explicit source hypothesis. \\
\phantomsection\label{entry:symbol-registry-01884}\([X=x]\) & \texttt{\detokenize{[X=x]}} & measurable equality event & P17 epistemic-nash-conditions & meta & \hyperref[def:p17-ab-system]{\ref*{def:p17-ab-system}} & \thepage & The variable and equality level must make a measurable event. \\
\phantomsection\label{entry:symbol-registry-01885}\(\phi^i,\phi,\varphi^i,\varphi\) & \texttt{\detokenize{\phi^i,}\allowbreak\detokenize{\phi,}\allowbreak\detokenize{\varphi^i,}\allowbreak\detokenize{\varphi}} & actual and specified full conjectures & P17 epistemic-nash-conditions & meta & \hyperref[def:p17-ab-conjecture]{\ref*{def:p17-ab-conjecture}} & \thepage & Each full conjecture concerns the joint opponents' action tuple; individual marginals are distinct. \\
\phantomsection\label{entry:symbol-registry-01886}\(R_i,R\) & \texttt{\detokenize{R_i,}\allowbreak\detokenize{R}} & individual and joint rationality events & P17 epistemic-nash-conditions & meta & \hyperref[def:p17-ab-conjecture]{\ref*{def:p17-ab-conjecture}} & \thepage & \(R_i\) is an own-type event; the literal \(R\) action in source Example 5.1 has its separately declared type. \\
\phantomsection\label{entry:symbol-registry-01887}\(K_i,K^1,CK,K^r\) & \texttt{\detokenize{K_i,}\allowbreak\detokenize{K^1,}\allowbreak\detokenize{CK,}\allowbreak\detokenize{K^r}} & probability-one individual mutual and common knowledge operators & P17 epistemic-nash-conditions & meta & \hyperref[def:p17-ab-knowledge]{\ref*{def:p17-ab-knowledge}} & \thepage & Original source operators retained. Common knowledge intersects positive-depth iterations; it is not generally factive. \\
\phantomsection\label{entry:symbol-registry-01888}\(P,P_i\) & \texttt{\detokenize{P,}\allowbreak\detokenize{P_i}} & common prior and type marginal & P17 epistemic-nash-conditions & meta & \hyperref[def:p17-ab-knowledge]{\ref*{def:p17-ab-knowledge}} & \thepage & Defined by the full conditional identity CP; null-type versions remain supplied. \\
\phantomsection\label{entry:symbol-registry-01889}\(\lambda\) & \texttt{\detokenize{\lambda}} & real probability proportionality constant & P17 epistemic-nash-conditions & meta & \hyperref[thm:p17-ab-measurability]{\ref*{thm:p17-ab-measurability}} & \thepage & Used for constant conditional probabilities and the source proportional-averaging identity. \\
\phantomsection\label{entry:symbol-registry-01890}\(\alpha\) & \texttt{\detokenize{\alpha}} & specified pure action profile & P17 epistemic Nash action profiles & meta & \hyperref[thm:p17-ab-preliminary]{\ref*{thm:p17-ab-preliminary}} & \thepage & Distinguished from actual state-dependent action a(s). \\
\phantomsection\label{entry:symbol-registry-01891}\(Q,Q_i,Q_{-i}\) & \texttt{\detokenize{Q,}\allowbreak\detokenize{Q_i,}\allowbreak\detokenize{Q_{-i}}} & conditional action law and its marginals & P17 epistemic-nash-conditions & meta & \hyperref[thm:p17-ab-product]{\ref*{thm:p17-ab-product}} & \thepage & In Theorem B Q is conditioned on the positive common-conjecture event F. \\
\phantomsection\label{entry:symbol-registry-01892}\(\sigma_i,\sigma\) & \texttt{\detokenize{\sigma_i,}\allowbreak\detokenize{\sigma}} & agreed marginal laws and resulting mixed profile & P17 epistemic-nash-conditions & meta & \hyperref[thm:p17-ab-many-player]{\ref*{thm:p17-ab-many-player}} & \thepage & Equilibrium support does not automatically contain every null-state actual action. \\
\phantomsection\label{entry:symbol-registry-01893}\(E,F,H,D\) & \texttt{\detokenize{E,}\allowbreak\detokenize{F,}\allowbreak\detokenize{H,}\allowbreak\detokenize{D}} & measurable event variables & P17 epistemic-nash-conditions & meta & \hyperref[def:p17-ab-knowledge]{\ref*{def:p17-ab-knowledge}} & \thepage & Local event bindings include \(F=\mathrm{CK}[\phi=\varphi]\) and \(H=[a_i=c_i]\cap F\); source action/type letters \(E,H,D\) are separately typed in their examples. \\
\phantomsection\label{entry:symbol-registry-01894}\(C,D,c,d,C_1,D_1,D_2,c_1,d_1,d_2\) & \texttt{\detokenize{C,}\allowbreak\detokenize{D,}\allowbreak\detokenize{c,}\allowbreak\detokenize{d,}\allowbreak\detokenize{C_1,}\allowbreak\detokenize{D_1,}\allowbreak\detokenize{D_2,}\allowbreak\detokenize{c_1,}\allowbreak\detokenize{d_1,}\allowbreak\detokenize{d_2}} & coordination-example action and type identifiers & P17 epistemic-nash-conditions & meta & \hyperref[thm:p17-ab-illustration]{\ref*{thm:p17-ab-illustration}} & \thepage & Source Figures 1-2 letters; capital C is an action in this model. \\
\phantomsection\label{entry:symbol-registry-01895}\(U,D,L,R,W,E\) & \texttt{\detokenize{U,}\allowbreak\detokenize{D,}\allowbreak\detokenize{L,}\allowbreak\detokenize{R,}\allowbreak\detokenize{W,}\allowbreak\detokenize{E}} & three-player source action and type identifiers & P17 epistemic-nash-conditions & meta & \hyperref[def:p17-ab-marginal-example]{\ref*{def:p17-ab-marginal-example}} & \thepage & Example 5.1 has one type per action; these letters denote finite alternatives locally. \\
\phantomsection\label{entry:symbol-registry-01896}\(H,T,h,t,W,E,H_r,T_r,h_r,t_r,W_r,E_1\) & \texttt{\detokenize{H,}\allowbreak\detokenize{T,}\allowbreak\detokenize{h,}\allowbreak\detokenize{t,}\allowbreak\detokenize{W,}\allowbreak\detokenize{E,}\allowbreak\detokenize{H_r,}\allowbreak\detokenize{T_r,}\allowbreak\detokenize{h_r,}\allowbreak\detokenize{t_r,}\allowbreak\detokenize{W_r,}\allowbreak\detokenize{E_1}} & chain-game actions and indexed types & P17 epistemic-nash-conditions & meta & \hyperref[def:p17-ab-chain-game]{\ref*{def:p17-ab-chain-game}} & \thepage & Source Figures 5-7; subscripted symbols identify types whose action is the corresponding letter. \\
\phantomsection\label{entry:symbol-registry-01897}\(m,L,r,w_r,Z,k\) & \texttt{\detokenize{m,}\allowbreak\detokenize{L,}\allowbreak\detokenize{r,}\allowbreak\detokenize{w_r,}\allowbreak\detokenize{Z,}\allowbreak\detokenize{k}} & finite-depth construction parameters and mass normalization & P17 epistemic-nash-conditions & meta & \hyperref[annex:p17.ab.chain.example]{\ref*{annex:p17.ab.chain.example}} & \thepage & Chain construction uses \(L=2m+5\) and \(w_r=3^{L-r}\); separate disagreement construction uses \(L=m+2\) and weights \(2^k\). Each \(Z\) is its explicitly stated finite positive normalizer. \\
\phantomsection\label{entry:symbol-registry-01898}\(x,y\) & \texttt{\detokenize{x,}\allowbreak\detokenize{y}} & matching-pennies mixture probabilities & P17 epistemic-nash-conditions & meta & \hyperref[annex:p17.ab.chain.example]{\ref*{annex:p17.ab.chain.example}} & \thepage & x is the probability of H and y of h in the complete equilibrium calculation. \\
\phantomsection\label{entry:symbol-registry-01899}\(L_\sigma\) & \texttt{\detokenize{L_\sigma}} & event that a fixed profile is Nash in the current game & P17 epistemic-nash-conditions & meta & \hyperref[annex:p17.ab.knowledge.of.equilibrium]{\ref*{annex:p17.ab.knowledge.of.equilibrium}} & \thepage & Not the event that varying conjectures at every possible state form equilibrium. \\
\phantomsection\label{entry:symbol-registry-01900}\(T_i,C_i\) & \texttt{\detokenize{T_i,}\allowbreak\detokenize{C_i}} & preceding finite-model type and action sets & P17 epistemic-nash-conditions & meta & \hyperref[thm:p17-ab-type-bridge]{\ref*{thm:p17-ab-type-bridge}} & \thepage & Bridge identifies predecessor \(C_i\) with source action set \(A_i\) and constructs source \(S_i=A_i\times T_i\). \\
\phantomsection\label{entry:symbol-registry-01993}\(N,n\) & \texttt{\detokenize{N,}\allowbreak\detokenize{n}} & finite electorate and its size & P17.judgment-aggregation & meta & \hyperref[def:p17-ja-agenda]{\ref*{def:p17-ja-agenda}} & \thepage & The main impossibility assumes n at least two; the one-person boundary is considered separately. \\
\phantomsection\label{entry:symbol-registry-01994}\(X,X_0\) & \texttt{\detokenize{X,}\allowbreak\detokenize{X_0}} & proposition agenda and the six-formula running agenda & P17.judgment-aggregation & meta & \hyperref[def:p17-ja-agenda]{\ref*{def:p17-ja-agenda}} & \thepage & \(X\) need not be finite or deductively closed; \(X_0\) consists of three opposing pairs. \\
\phantomsection\label{entry:symbol-registry-01995}\(p,q,r,c\) & \texttt{\detokenize{p,}\allowbreak\detokenize{q,}\allowbreak\detokenize{r,}\allowbreak\detokenize{c}} & propositional atoms and the abbreviation c for p and q & P17.judgment-aggregation & object & \hyperref[def:p17-ja-agenda]{\ref*{def:p17-ja-agenda}} & \thepage & p and q are distinct atoms; c is the conjunction p land q when introduced locally, not a new independent premise. \\
\phantomsection\label{entry:symbol-registry-01996}\(J,J_i\) & \texttt{\detokenize{J,}\allowbreak\detokenize{J_i}} & collective or individual sets of accepted agenda propositions & P17.judgment-aggregation & meta & \hyperref[def:p17-ja-rational]{\ref*{def:p17-ja-rational}} & \thepage & These replace source \(\Phi\) and \(\Phi_i\) by explicitly defined local notation; nonmembership means nonassent. \\
\phantomsection\label{entry:symbol-registry-01997}\(\mathcal J_X\) & \texttt{\detokenize{\mathcal J_X}} & complete weakly consistent agenda-relatively closed judgment sets & P17.judgment-aggregation & meta & \hyperref[def:p17-ja-rational]{\ref*{def:p17-ja-rational}} & \thepage & Weak consistency is literal pair exclusion, not semantic satisfiability by itself. \\
\phantomsection\label{entry:symbol-registry-01998}\(\operatorname{cl}_X\) & \texttt{\detokenize{\operatorname{cl}_X}} & semantic consequences restricted to the agenda & P17.judgment-aggregation & meta & \hyperref[def:p17-ja-rational]{\ref*{def:p17-ja-rational}} & \thepage & The argument may be any set of formulas; only conclusions belonging to X are retained. \\
\phantomsection\label{entry:symbol-registry-01999}\(v,v_i\) & \texttt{\detokenize{v,}\allowbreak\detokenize{v_i}} & classical propositional valuations & P17.judgment-aggregation & meta & \hyperref[thm:p17-ja-models]{\ref*{thm:p17-ja-models}} & \thepage & Input profiles in the impossibility proof are realized by full valuations. \\
\phantomsection\label{entry:symbol-registry-02000}\(\mathbf J\) & \texttt{\detokenize{\mathbf J}} & ordered profile of individual judgment sets & P17.judgment-aggregation & meta & \hyperref[def:p17-ja-profile]{\ref*{def:p17-ja-profile}} & \thepage & Individual identities are the coordinates before an anonymity permutation. \\
\phantomsection\label{entry:symbol-registry-02001}\(N_\varphi,s_\varphi\) & \texttt{\detokenize{N_\varphi,}\allowbreak\detokenize{s_\varphi}} & supporting coalition and number of supporters of a proposition & P17.judgment-aggregation & meta & \hyperref[def:p17-ja-profile]{\ref*{def:p17-ja-profile}} & \thepage & Both depend on the profile; their arguments may be suppressed locally in the appendix proof. \\
\phantomsection\label{entry:symbol-registry-02002}\(M\) & \texttt{\detokenize{M}} & strict proposition-wise majority function & P17.judgment-aggregation & meta & \hyperref[def:p17-ja-profile]{\ref*{def:p17-ja-profile}} & \thepage & Closure is a separate operation; odd electorates are specified when completeness is claimed. \\
\phantomsection\label{entry:symbol-registry-02003}\(F,f,g\) & \texttt{\detokenize{F,}\allowbreak\detokenize{f,}\allowbreak\detokenize{g}} & aggregation function and its systematic Boolean and anonymous count functions & P17.judgment-aggregation & meta & \hyperref[def:p17-ja-uas]{\ref*{def:p17-ja-uas}} & \thepage & One f applies to all propositions and all profiles; g is derived from anonymity and realizability of support vectors. \\
\phantomsection\label{entry:symbol-registry-02004}\(\mathcal D,\mathcal P(X)\) & \texttt{\detokenize{\mathcal D,}\allowbreak\detokenize{\mathcal P(X)}} & admissible profile domain and power set of the proposition agenda & P17.judgment-aggregation & meta & \hyperref[def:p17-ja-uas]{\ref*{def:p17-ja-uas}} & \thepage & The source universal domain is exactly the set of all rational profiles. \\
\phantomsection\label{entry:symbol-registry-02005}\(\delta_i(\varphi)\) & \texttt{\detokenize{\delta_i(\varphi)}} & indicator of individual assent to an agenda proposition & P17.judgment-aggregation & meta & \hyperref[def:p17-ja-uas]{\ref*{def:p17-ja-uas}} & \thepage & Values are zero and one; assent to a negation is distinct from absent assent without completeness. \\
\phantomsection\label{entry:symbol-registry-02006}\(\sigma\) & \texttt{\detokenize{\sigma}} & permutation of individuals & P17.judgment-aggregation & meta & \hyperref[def:p17-ja-uas]{\ref*{def:p17-ja-uas}} & \thepage & Used for anonymity, not for substitution into formulas. \\
\phantomsection\label{entry:symbol-registry-02007}\(Y,a,B_Y,F_Y\) & \texttt{\detokenize{Y,}\allowbreak\detokenize{a,}\allowbreak\detokenize{B_Y,}\allowbreak\detokenize{F_Y}} & independent priority questions selected truth assignment signed answers and their agenda closure & P17.judgment-aggregation & meta & \hyperref[def:p17-ja-priority-procedure]{\ref*{def:p17-ja-priority-procedure}} & \thepage & Independence means every assignment to Y has a classical extension; completeness additionally needs determination of every agenda formula. \\
\phantomsection\label{entry:symbol-registry-02008}\(U_\cap\) & \texttt{\detokenize{U_\cap}} & unanimous intersection of individual judgment sets & P17.judgment-aggregation & meta & \hyperref[thm:p17-ja-unanimity]{\ref*{thm:p17-ja-unanimity}} & \thepage & Not the symbol for the universal-domain condition U. \\
\phantomsection\label{entry:symbol-registry-02009}\(k,\theta_j\) & \texttt{\detokenize{k,}\allowbreak\detokenize{\theta_j}} & number of issue pairs and their positive representatives in the finite support theorem & P17.judgment-aggregation & meta & \hyperref[def:p17-ja-finite-issues]{\ref*{def:p17-ja-finite-issues}} & \thepage & The 2k formulas form disjoint opposing pairs; k counts issues rather than strings. \\
\phantomsection\label{entry:symbol-registry-02010}\(H_k,G_k\) & \texttt{\detokenize{H_k,}\allowbreak\detokenize{G_k}} & strict supermajority output and its agenda-relative closure & P17.judgment-aggregation & meta & \hyperref[def:p17-ja-finite-issues]{\ref*{def:p17-ja-finite-issues}} & \thepage & \(H_k\) requires support strictly greater than \((k-1)n/k\); \(G_k\) can fail systematicity. \\
\phantomsection\label{entry:symbol-registry-02011}\(\sim,f_C\) & \texttt{\detokenize{\sim,}\allowbreak\detokenize{f_C}} & equivalence of agenda propositions and a class-dependent Boolean aggregation function & P17.judgment-aggregation & meta & \hyperref[def:p17-ja-relative-systematicity]{\ref*{def:p17-ja-relative-systematicity}} & \thepage & A weaker systematicity convention from source note 4; the original impossibility is not asserted for every equivalence relation. \\
\phantomsection\label{entry:symbol-registry-02298}\(V,C,P,F,T,\mathcal N\) & \texttt{\detokenize{V,}\allowbreak\detokenize{C,}\allowbreak\detokenize{P,}\allowbreak\detokenize{F,}\allowbreak\detokenize{T,}\allowbreak\detokenize{\mathcal N}} & variable constant predicate primitive-task compound-task and label vocabularies & P17.HTN & meta & \hyperref[def:p17-htn-vocabulary]{\ref*{def:p17-htn-vocabulary}} & \thepage & C P F T are finite; variables and occurrence labels have fresh supplies. \\
\phantomsection\label{entry:symbol-registry-02299}\(s,I\) & \texttt{\detokenize{s,}\allowbreak\detokenize{I}} & world state and initial world state as sets of ground atoms & P17.HTN & meta & \hyperref[def:p17-htn-vocabulary]{\ref*{def:p17-htn-vocabulary}} & \thepage & Closed-world interpretation; unlisted ground atoms are false. \\
\phantomsection\label{entry:symbol-registry-02300}\(\operatorname{do},\operatorname{achieve},\operatorname{perform}\) & \texttt{\detokenize{\operatorname{do},}\allowbreak\detokenize{\operatorname{achieve},}\allowbreak\detokenize{\operatorname{perform}}} & three syntactic task constructors & P17.HTN & object & \hyperref[def:p17-htn-vocabulary]{\ref*{def:p17-htn-vocabulary}} & \thepage & Primitive action; literal goal; compound task. \\
\phantomsection\label{entry:symbol-registry-02301}\(\operatorname{add},\operatorname{del},\operatorname{apply}\) & \texttt{\detokenize{\operatorname{add},}\allowbreak\detokenize{\operatorname{del},}\allowbreak\detokenize{\operatorname{apply}}} & positive-atom effect sets and partial execution map & P17.HTN & meta & \hyperref[def:p17-htn-operator]{\ref*{def:p17-htn-operator}} & \thepage & Delete sets contain underlying atoms, not negative literal strings; add follows delete. \\
\phantomsection\label{entry:symbol-registry-02302}\(\sigma,a_i,s_i\) & \texttt{\detokenize{\sigma,}\allowbreak\detokenize{a_i,}\allowbreak\detokenize{s_i}} & ground action sequence its actions and execution states & P17.HTN & meta & \hyperref[def:p17-htn-operator]{\ref*{def:p17-htn-operator}} & \thepage & Action positions start at one; \(s_0\) is the state before the first action. \\
\phantomsection\label{entry:symbol-registry-02303}\(d,n_i,\alpha_i,\phi\) & \texttt{\detokenize{d,}\allowbreak\detokenize{n_i,}\allowbreak\detokenize{\alpha_i,}\allowbreak\detokenize{\phi}} & task network occurrence labels tasks and Boolean constraint formula & P17.HTN & meta & \hyperref[def:p17-htn-network]{\ref*{def:p17-htn-network}} & \thepage & Equal task expressions may have distinct labels. \\
\phantomsection\label{entry:symbol-registry-02304}\(\prec\) & \texttt{\detokenize{\prec}} & strict precedence between task occurrences or their resolved endpoints & P17.HTN & object & \hyperref[def:p17-htn-network]{\ref*{def:p17-htn-network}} & \thepage & After decomposition, a predecessor precedes every child and a successor follows every child. \\
\phantomsection\label{entry:symbol-registry-02305}\((\ell,n),(n,\ell),(n,\ell,n')\) & \texttt{\detokenize{(\ell,}\allowbreak\detokenize{n),}\allowbreak\detokenize{(n,}\allowbreak\detokenize{\ell),}\allowbreak\detokenize{(n,}\allowbreak\detokenize{\ell,}\allowbreak\detokenize{n')}} & before after and protected-interval constraint atoms & P17.HTN & object & \hyperref[def:p17-htn-network]{\ref*{def:p17-htn-network}} & \thepage & Protection checks state indices pi(n) less than or equal to e strictly below pi(n'). \\
\phantomsection\label{entry:symbol-registry-02306}\(\operatorname{first},\operatorname{last},\pi\) & \texttt{\detokenize{\operatorname{first},}\allowbreak\detokenize{\operatorname{last},}\allowbreak\detokenize{\pi}} & extreme occurrence positions and the execution-position map & P17.HTN & meta & \hyperref[def:p17-htn-completion]{\ref*{def:p17-htn-completion}} & \thepage & First and last have nonempty label arguments. \\
\phantomsection\label{entry:symbol-registry-02307}\(\theta\) & \texttt{\detokenize{\theta}} & consistent substitution used in grounding and method unification & P17.HTN & meta & \hyperref[def:p17-htn-method]{\ref*{def:p17-htn-method}} & \thepage & Methods are renamed apart before unification. \\
\phantomsection\label{entry:symbol-registry-02308}\(\mathcal D,\operatorname{Op},\operatorname{Me}\) & \texttt{\detokenize{\mathcal D,}\allowbreak\detokenize{\operatorname{Op},}\allowbreak\detokenize{\operatorname{Me}}} & planning domain and its finite operator and method sets & P17.HTN & meta & \hyperref[def:p17-htn-solutions]{\ref*{def:p17-htn-solutions}} & \thepage & Fixed domain fixes schemas; the instance may supply additional constants and initial facts. \\
\phantomsection\label{entry:symbol-registry-02309}\(\operatorname{comp},\operatorname{red},\operatorname{sol}_r,\operatorname{sol}\) & \texttt{\detokenize{\operatorname{comp},}\allowbreak\detokenize{\operatorname{red},}\allowbreak\detokenize{\operatorname{sol}_r,}\allowbreak\detokenize{\operatorname{sol}}} & primitive completion one-step reductions bounded derivability and all finite solutions & P17.HTN & meta & \hyperref[def:p17-htn-solutions]{\ref*{def:p17-htn-solutions}} & \thepage & The initial state is unchanged by syntactic decomposition; execution is checked at primitive completion. \\
\phantomsection\label{entry:symbol-registry-02310}\(k,\operatorname{level}\) & \texttt{\detokenize{k,}\allowbreak\detokenize{\operatorname{level}}} & finite task-level bound and decreasing ground-task rank & P17.HTN & meta & \hyperref[def:p17-htn-levels]{\ref*{def:p17-htn-levels}} & \thepage & Primitive tasks have rank zero; compound leaves without methods need not have solutions. \\
\phantomsection\label{entry:symbol-registry-02311}\(L,b,g,S\) & \texttt{\detokenize{L,}\allowbreak\detokenize{b,}\allowbreak\detokenize{g,}\allowbreak\detokenize{S}} & explicit input length relevant-ground-atom count ground-task count and summary-count bound & P17.HTN & meta & \hyperref[sec:p17-htn-restrictions]{\ref*{sec:p17-htn-restrictions}} & \thepage & The letter g for the guard atom is confined to the worked example; here g is a cardinal count in the tabulation proof. Irrelevant atom coordinates are projected away. \\
\phantomsection\label{entry:symbol-registry-02312}\(\rho,\eta,\tau,\Lambda\) & \texttt{\detokenize{\rho,}\allowbreak\detokenize{\eta,}\allowbreak\detokenize{\tau,}\allowbreak\detokenize{\Lambda}} & nonempty execution traces their finite summaries and true ground-literal sets & P17.HTN & meta & \hyperref[def:p17-htn-trace-summary]{\ref*{def:p17-htn-trace-summary}} & \thepage & A summary includes first/last boundary states and the length-one flag. \\
\phantomsection\label{entry:symbol-registry-02313}\(A(\rho),B(\rho),C(\rho),D(\rho)\) & \texttt{\detokenize{A(\rho),}\allowbreak\detokenize{B(\rho),}\allowbreak\detokenize{C(\rho),}\allowbreak\detokenize{D(\rho)}} & literal intersections over all states after-first before-last and internal states & P17.HTN & meta & \hyperref[def:p17-htn-trace-summary]{\ref*{def:p17-htn-trace-summary}} & \thepage & An empty state range has the full literal set as its intersection; these are local trace functions, distinct from vocabulary C and domain mathcal D. \\
\phantomsection\label{entry:symbol-registry-02314}\(h\) & \texttt{\detokenize{h}} & observation map erasing a distinguished bookkeeping action & P17.HTN & meta & \hyperref[def:p17-htn-strips]{\ref*{def:p17-htn-strips}} & \thepage & Erasure is not claimed to be symbol renaming. \\
\phantomsection\label{entry:symbol-registry-02560}\(A_i\) & \texttt{\detokenize{A_i}} & external action in the 1989 metareasoning presentation & \hyperref[sec:p17-metareasoning]{\ref*{sec:p17-metareasoning}} & object & \hyperref[def:p17-meta-source-notation]{\ref*{def:p17-meta-source-notation}} & \thepage & An action in the world; distinct from an executable computational block. \\
\phantomsection\label{entry:symbol-registry-02561}\(S_j\) & \texttt{\detokenize{S_j}} & computational action in the 1989 source notation & \hyperref[sec:p17-metareasoning]{\ref*{sec:p17-metareasoning}} & object & \hyperref[def:p17-meta-source-notation]{\ref*{def:p17-meta-source-notation}} & \thepage & A complete computation commits to action; a partial computation leaves another choice. \\
\phantomsection\label{entry:symbol-registry-02562}\(S,T\) & \texttt{\detokenize{S,}\allowbreak\detokenize{T}} & completed history and possible future complete computation trace & \hyperref[sec:p17-metareasoning]{\ref*{sec:p17-metareasoning}} & object & \hyperref[def:p17-meta-source-notation]{\ref*{def:p17-meta-source-notation}} & \thepage & Dot concatenation joins computational traces; a future trace is contingent on observed computation outcomes. \\
\phantomsection\label{entry:symbol-registry-02563}\([X,W],[X]\) & \texttt{\detokenize{[X,}\allowbreak\detokenize{W],}\allowbreak\detokenize{[X]}} & situation resulting from a specified action in the source notation & \hyperref[sec:p17-metareasoning]{\ref*{sec:p17-metareasoning}} & object & \hyperref[def:p17-meta-source-notation]{\ref*{def:p17-meta-source-notation}} & \thepage & The omitted initial situation is understood; both external and computational actions may occur. \\
\phantomsection\label{entry:symbol-registry-02564}\(U(W),\widehat U^S(W)\) & \texttt{\detokenize{U(W),}\allowbreak\detokenize{\widehat U^S(W)}} & actual external utility and a computation-stage estimate & \hyperref[sec:p17-metareasoning]{\ref*{sec:p17-metareasoning}} & object & \hyperref[def:p17-meta-source-notation]{\ref*{def:p17-meta-source-notation}} & \thepage & The evaluation snapshot is explicit; a heuristic estimate is not automatically calibrated. \\
\phantomsection\label{entry:symbol-registry-02565}\(\alpha(S)\) & \texttt{\detokenize{\alpha(S)}} & current external intention after computation history S & \hyperref[sec:p17-metareasoning]{\ref*{sec:p17-metareasoning}} & object & \hyperref[def:p17-meta-source-notation]{\ref*{def:p17-meta-source-notation}} & \thepage & A stated tie rule selects the named action whose current estimate is largest. \\
\phantomsection\label{entry:symbol-registry-02566}\(\alpha_c\) & \texttt{\detokenize{\alpha_c}} & external action recommended by complete computation c & \hyperref[sec:p17-metareasoning]{\ref*{sec:p17-metareasoning}} & object & \hyperref[def:p17-meta-complete-value]{\ref*{def:p17-meta-complete-value}} & \thepage & The action maximizes the new estimate vector; named actions can index their vector coordinates. \\
\phantomsection\label{entry:symbol-registry-02567}\(z\) & \texttt{\detokenize{z}} & metalevel information summary available to the computation controller & \hyperref[sec:p17-metareasoning]{\ref*{sec:p17-metareasoning}} & object & \hyperref[def:p17-meta-program-model]{\ref*{def:p17-meta-program-model}} & \thepage & Contains declared input or workspace features clock memory and retained computation results. \\
\phantomsection\label{entry:symbol-registry-02568}\(F_c\) & \texttt{\detokenize{F_c}} & actual executable-block transition on a modeled machine configuration & \hyperref[sec:p17-metareasoning]{\ref*{sec:p17-metareasoning}} & object & \hyperref[def:p17-meta-program-model]{\ref*{def:p17-meta-program-model}} & \thepage & Includes execution output running time and peak memory; deterministic execution may have an uncertain predicted result. \\
\phantomsection\label{entry:symbol-registry-02569}\(I\) & \texttt{\detokenize{I}} & map from modeled machine configurations to controller summaries & \hyperref[sec:p17-metareasoning]{\ref*{sec:p17-metareasoning}} & object & \hyperref[def:p17-meta-program-model]{\ref*{def:p17-meta-program-model}} & \thepage & Compression must retain a sufficient transition cost feasibility and stopping model. \\
\phantomsection\label{entry:symbol-registry-02570}\(\Omega,\omega\) & \texttt{\detokenize{\Omega,}\allowbreak\detokenize{\omega}} & finite class of modeled problem instances and one instance & \hyperref[sec:p17-metareasoning]{\ref*{sec:p17-metareasoning}} & object & \hyperref[def:p17-meta-program-model]{\ref*{def:p17-meta-program-model}} & \thepage & May include unresolved configuration or physical random-bit uncertainty; does not create randomness in a deterministic calculation. \\
\phantomsection\label{entry:symbol-registry-02571}\(\mu_z(\omega)\) & \texttt{\detokenize{\mu_z(\omega)}} & conditional predictive instance weights at a metalevel summary & \hyperref[sec:p17-metareasoning]{\ref*{sec:p17-metareasoning}} & object & \hyperref[def:p17-meta-program-model]{\ref*{def:p17-meta-program-model}} & \thepage & A population or subjective computation model; zero-probability successors require no posterior quotient. \\
\phantomsection\label{entry:symbol-registry-02572}\(K(z'\mid z,c)\) & \texttt{\detokenize{K(z'\mid z,}\allowbreak\detokenize{c)}} & successor-summary kernel induced by an executable computation model & \hyperref[sec:p17-metareasoning]{\ref*{sec:p17-metareasoning}} & object & \hyperref[def:p17-meta-program-model]{\ref*{def:p17-meta-program-model}} & \thepage & Its definition groups the actual program-transition outputs by their summaries. \\
\phantomsection\label{entry:symbol-registry-02573}\(X,p_c\) & \texttt{\detokenize{X,}\allowbreak\detokenize{p_c}} & joint vector of post-computation intrinsic estimates and its predictive distribution & \hyperref[sec:p17-metareasoning]{\ref*{sec:p17-metareasoning}} & object & \hyperref[def:p17-meta-complete-value]{\ref*{def:p17-meta-complete-value}} & \thepage & Several changing estimates need a joint law; subtree independence means one changed component. \\
\phantomsection\label{entry:symbol-registry-02574}\(\widehat{\operatorname{VOC}}_1(c)\) & \texttt{\detokenize{\widehat{\operatorname{VOC}}_1(c)}} & one-step estimated net value of a complete computation & \hyperref[sec:p17-metareasoning]{\ref*{sec:p17-metareasoning}} & object & \hyperref[def:p17-meta-complete-value]{\ref*{def:p17-meta-complete-value}} & \thepage & Compares new and old intentions using the same later estimate then charges the declared common delay cost once. \\
\phantomsection\label{entry:symbol-registry-02575}\(x_+\) & \texttt{\detokenize{x_+}} & positive part of a real scalar in computation-benefit formulas & \hyperref[thm:p17-meta-subtree]{\ref*{thm:p17-meta-subtree}} & object & \hyperref[thm:p17-meta-subtree]{\ref*{thm:p17-meta-subtree}} & \thepage & Defined as the maximum of x and zero. \\
\phantomsection\label{entry:symbol-registry-02576}\(C(z)\) & \texttt{\detokenize{C(z)}} & finite set of resource-feasible executable blocks at summary z & \hyperref[sec:p17-metareasoning]{\ref*{sec:p17-metareasoning}} & object & \hyperref[def:p17-meta-finite-control]{\ref*{def:p17-meta-finite-control}} & \thepage & May be empty; a deadline or memory violation excludes a block. \\
\phantomsection\label{entry:symbol-registry-02577}\(g(z)\) & \texttt{\detokenize{g(z)}} & model expected external payoff of committing to the implemented intention & \hyperref[sec:p17-metareasoning]{\ref*{sec:p17-metareasoning}} & object & \hyperref[def:p17-meta-finite-control]{\ref*{def:p17-meta-finite-control}} & \thepage & The stopping value does not optimize an uncomputed true-utility table. \\
\phantomsection\label{entry:symbol-registry-02578}\(k(z,c)\) & \texttt{\detokenize{k(z,}\allowbreak\detokenize{c)}} & expected direct resource cost of one executable block & \hyperref[sec:p17-metareasoning]{\ref*{sec:p17-metareasoning}} & object & \hyperref[def:p17-meta-finite-control]{\ref*{def:p17-meta-finite-control}} & \thepage & Delay belongs either to the time-dependent stopping payoff or to a separable cost; it is charged only once. \\
\phantomsection\label{entry:symbol-registry-02579}\(V_n(z),Q_n(z,c)\) & \texttt{\detokenize{V_n(z),}\allowbreak\detokenize{Q_n(z,}\allowbreak\detokenize{c)}} & optimal remaining-budget value and first-computation continuation value & \hyperref[sec:p17-metareasoning]{\ref*{sec:p17-metareasoning}} & object & \hyperref[thm:p17-meta-bellman]{\ref*{thm:p17-meta-bellman}} & \thepage & At budget zero the external stopping payoff grounds the backward recursion. \\
\phantomsection\label{entry:symbol-registry-02580}\(p_n(z),d_n(c\mid z),W_n(z)\) & \texttt{\detokenize{p_n(z),}\allowbreak\detokenize{d_n(c\mid z),}\allowbreak\detokenize{W_n(z)}} & stopping probability conditional computation policy and predicted policy return & \hyperref[sec:p17-metareasoning]{\ref*{sec:p17-metareasoning}} & object & \hyperref[def:p17-meta-policy-model]{\ref*{def:p17-meta-policy-model}} & \thepage & This models the actual controller without assuming that its choices are optimal. \\
\phantomsection\label{entry:symbol-registry-02581}\(\widetilde Q_n,\epsilon_n\) & \texttt{\detokenize{\widetilde Q_n,}\allowbreak\detokenize{\epsilon_n}} & approximate optimal-continuation scores and their uniform error bounds & \hyperref[sec:p17-metareasoning]{\ref*{sec:p17-metareasoning}} & object & \hyperref[thm:p17-meta-approximate]{\ref*{thm:p17-meta-approximate}} & \thepage & The policy-loss theorem requires an established bound at every state and remaining budget. \\
\phantomsection\label{entry:symbol-registry-02582}\(r\) & \texttt{\detokenize{r}} & additional action-independent utility cost of using the computation controller & \hyperref[thm:p17-meta-regress]{\ref*{thm:p17-meta-regress}} & object & \hyperref[thm:p17-meta-regress]{\ref*{thm:p17-meta-regress}} & \thepage & This is additional to the direct costs of the chosen computation blocks. \\
\phantomsection\label{entry:symbol-registry-02583}\((p,q),c_1,c_2\) & \texttt{\detokenize{(p,}\allowbreak\detokenize{q),}\allowbreak\detokenize{c_1,}\allowbreak\detokenize{c_2}} & input-bit pair and the two priced executable blocks in the worked program & \hyperref[def:p17-meta-two-block-program]{\ref*{def:p17-meta-two-block-program}} & object & \hyperref[def:p17-meta-two-block-program]{\ref*{def:p17-meta-two-block-program}} & \thepage & Bits are already in input memory; loading and evaluating them are computation actions not new external measurements. \\
\phantomsection\label{entry:symbol-registry-02972}\(Y,m,\epsilon\) & \texttt{\detokenize{Y;m;\epsilon}} & test response true conditional mean and centered test noise & P17 bounded-rationality squared-error comparison & n/a & \hyperref[thm:p17-br-bias-variance]{\ref*{thm:p17-br-bias-variance}} & \thepage & The fixed-input comparison stipulates a joint training and independent-test-noise law with finite second moments. \\
\phantomsection\label{entry:symbol-registry-02973}\(\widehat m\) & \texttt{\detokenize{\widehat m}} & trained prediction at the fixed input & P17 bounded-rationality squared-error comparison & n/a & \hyperref[thm:p17-br-bias-variance]{\ref*{thm:p17-br-bias-variance}} & \thepage & Its randomness comes from training; its second moment is finite. \\
\phantomsection\label{entry:symbol-registry-02974}\(\sigma^2\) & \texttt{\detokenize{\sigma^2}} & variance of the centered test noise & P17 bounded-rationality squared-error comparison & n/a & \hyperref[thm:p17-br-bias-variance]{\ref*{thm:p17-br-bias-variance}} & \thepage & The noise component is common to both compared procedures in the stipulated example. \\
\phantomsection\label{entry:symbol-registry-02975}\(a,v\) & \texttt{\detokenize{a;v}} & constant predictor and variance of an unbiased trained predictor & P17 bounded-rationality squared-error comparison & n/a & \hyperref[thm:p17-br-bias-variance]{\ref*{thm:p17-br-bias-variance}} & \thepage & These describe prediction procedures; a denotes an alternative only in the disjoint cue-comparison scope. \\
\phantomsection\label{entry:symbol-registry-02976}\(\mu\) & \texttt{\detokenize{\mu}} & mean of the trained predictor & P17 bounded-rationality squared-error proof & n/a & \hyperref[annex:p17.br.bias-variance]{\ref*{annex:p17.br.bias-variance}} & \thepage & Subtracting this mean separates the bias and variance terms in the full proof. \\
\phantomsection\label{entry:symbol-registry-02977}\(w_i\) & \texttt{\detokenize{w_i}} & oriented nonnegative linear-comparison weight & P17 bounded-rationality noncompensation comparison & n/a & \hyperref[thm:p17-br-noncompensation]{\ref*{thm:p17-br-noncompensation}} & \thepage & Negative original coefficients can be oriented by complementing binary attributes; zero-weight cues are omitted from the lexicographic rule. \\
\phantomsection\label{entry:symbol-registry-02978}\(d_i\) & \texttt{\detokenize{d_i}} & difference between two binary attribute values & P17 bounded-rationality noncompensation comparison & n/a & \hyperref[thm:p17-br-noncompensation]{\ref*{thm:p17-br-noncompensation}} & \thepage & Values are minus one zero or one; the first discriminating positive-weight cue governs the lemma under its strict inequalities. \\
\phantomsection\label{entry:symbol-registry-02979}\(r\) & \texttt{\detokenize{r}} & first discriminating positive-weight cue index & P17 bounded-rationality noncompensation proof & n/a & \hyperref[annex:p17.br.noncompensation]{\ref*{annex:p17.br.noncompensation}} & \thepage & Earlier weighted terms vanish; the later sum is bounded before its sign is concluded. \\
\phantomsection\label{entry:symbol-registry-02980}\(a,b\) & \texttt{\detokenize{a;b}} & specified alternatives with three binary coordinates & P17 bounded-rationality feasible-choice example & n/a & \hyperref[def:p17-br-cue-instance]{\ref*{def:p17-br-cue-instance}} & \thepage & Their coordinates are respectively one-zero-zero and zero-one-one; inspecting a cue reveals both alternatives' values. \\
\phantomsection\label{entry:symbol-registry-02981}\(N,C\) & \texttt{\detokenize{N;C}} & two stipulated environments with different evaluation weights & P17 bounded-rationality feasible-choice example & n/a & \hyperref[def:p17-br-cue-instance]{\ref*{def:p17-br-cue-instance}} & \thepage & N uses four-two-one and C uses one-two-four; they are analytical environments rather than estimated population distributions. \\
\phantomsection\label{entry:symbol-registry-03842}\(L\) & \texttt{\detokenize{L}} & finite described monetary lottery & P17 prospect theory & n/a & \hyperref[def:p17-pt-domain]{\ref*{def:p17-pt-domain}} & \thepage & Equal outcomes are combined by probability addition in distribution based CPT. \\
\phantomsection\label{entry:symbol-registry-03843}\(r\) & \texttt{\detokenize{r}} & monetary reference point & P17 prospect theory & n/a & \hyperref[def:p17-pt-domain]{\ref*{def:p17-pt-domain}} & \thepage & Shared across alternatives within a specified comparison. \\
\phantomsection\label{entry:symbol-registry-03844}\(x_i\) & \texttt{\detokenize{x_i}} & monetary deviation from reference & P17 prospect theory & n/a & \hyperref[def:p17-pt-domain]{\ref*{def:p17-pt-domain}} & \thepage & Not final wealth or an outcome probability. \\
\phantomsection\label{entry:symbol-registry-03845}\(v(x)\) & \texttt{\detokenize{v}} & value of a deviation & P17 prospect theory & n/a & \hyperref[def:p17-pt-power-value]{\ref*{def:p17-pt-power-value}} & \thepage & Distinct from the existing expected utility prize function and complete lottery score. \\
\phantomsection\label{entry:symbol-registry-03846}\(V(L;r)\) & \texttt{\detokenize{V}} & prospect evaluation score & P17 prospect theory & n/a & \hyperref[def:p17-pt-domain]{\ref*{def:p17-pt-domain}} & \thepage & An as if predictive index not a new normative obligation. \\
\phantomsection\label{entry:symbol-registry-03847}\(\alpha,\beta\) & \texttt{\detokenize{\alpha;\beta}} & gain and loss curvature exponents & P17 prospect theory & n/a & \hyperref[def:p17-pt-power-value]{\ref*{def:p17-pt-power-value}} & \thepage & Positive for monotonicity; strictly below one for the stated diminishing sensitivity case. \\
\phantomsection\label{entry:symbol-registry-03848}\(\lambda\) & \texttt{\detokenize{\lambda}} & loss value multiplier & P17 prospect theory & n/a & \hyperref[def:p17-pt-power-value]{\ref*{def:p17-pt-power-value}} & \thepage & A constant equal magnitude ratio only when alpha equals beta. \\
\phantomsection\label{entry:symbol-registry-03849}\(w^+,w^-\) & \texttt{\detokenize{w^+;w^-}} & increasing normalized gain and loss weighting functions & P17 prospect theory & n/a & \hyperref[def:p17-pt-binary-weighting]{\ref*{def:p17-pt-binary-weighting}} & \thepage & Not additive probability measures. \\
\phantomsection\label{entry:symbol-registry-03850}\(P_i^+,P_i^-\) & \texttt{\detokenize{P_i^+;P_i^-}} & upper gain and lower loss cumulative probabilities & P17 prospect theory & n/a & \hyperref[def:p17-pt-cumulative]{\ref*{def:p17-pt-cumulative}} & \thepage & Rank and sign specific tails with explicit strict comparison complements. \\
\phantomsection\label{entry:symbol-registry-03851}\(\pi_i^+,\pi_i^-\) & \texttt{\detokenize{\pi_i^+;\pi_i^-}} & rank dependent cumulative decision weights & P17 prospect theory & n/a & \hyperref[def:p17-pt-cumulative]{\ref*{def:p17-pt-cumulative}} & \thepage & Mixed nonzero outcome weights need not sum to one. \\
\phantomsection\label{entry:symbol-registry-03852}\(\gamma,\delta\) & \texttt{\detokenize{\gamma;\delta}} & gain and loss probability sensitivity parameters & P17 prospect theory & n/a & \hyperref[def:p17-pt-cumulative]{\ref*{def:p17-pt-cumulative}} & \thepage & The chapter identifies their roles; no unrestricted parameter family is assumed here. \\
\phantomsection\label{entry:symbol-registry-03853}\(K\) & \texttt{\detokenize{K}} & tails count before first head in the source payoff convention & P17 prospect theory St Petersburg & n/a & \hyperref[thm:p17-pt-source-calculations]{\ref*{thm:p17-pt-source-calculations}} & \thepage & Zero payoff for immediate head; differs from other indexing conventions. \\
\phantomsection\label{entry:symbol-registry-03854}\(N\) & \texttt{\detokenize{N}} & count of rare outcomes in a stipulated sampling calculation & P17 prospect theory sampling & n/a & \hyperref[thm:p17-pt-sampling]{\ref*{thm:p17-pt-sampling}} & \thepage & Separate from a participant sample size in the empirical framing report. \\
\phantomsection\label{entry:symbol-registry-03855}\(\pi(p)\) & \texttt{\detokenize{\pi}} & original 1979 probability decision weight & P17 original prospect theory & n/a & \hyperref[sec:p17-rationality-8-3]{\ref*{sec:p17-rationality-8-3}} & \thepage & Distinct from the cumulative rank specific pi weights; the two original evaluation branches are explicit. \\
\phantomsection\label{entry:symbol-registry-03987}\(S\) & \texttt{\detokenize{S}} & finite state space & P17 Rationality8.4 & n/a & \hyperref[def:p17-rat84-acts]{\ref*{def:p17-rat84-acts}} & \thepage & Each state resolves both urn draws in the source table. \\
\phantomsection\label{entry:symbol-registry-03988}\(X\) & \texttt{\detokenize{X}} & consequence space & P17 Rationality8.4 & n/a & \hyperref[def:p17-rat84-acts]{\ref*{def:p17-rat84-acts}} & \thepage & Consequences and their utility numbers remain distinct inputs. \\
\phantomsection\label{entry:symbol-registry-03989}\(f,g\) & \texttt{\detokenize{f;g}} & deterministic acts from states to consequences & P17 Rationality8.4 & n/a & \hyperref[def:p17-rat84-acts]{\ref*{def:p17-rat84-acts}} & \thepage & Hill's simplified domain differs from the general lottery-act domain of Gilboa Schmeidler. \\
\phantomsection\label{entry:symbol-registry-03990}\(u\) & \texttt{\detokenize{u}} & common state independent utility of consequences & P17 Rationality8.4 & n/a & \hyperref[def:p17-rat84-acts]{\ref*{def:p17-rat84-acts}} & \thepage & Fixed throughout the source comparison; not assumed linear in money. \\
\phantomsection\label{entry:symbol-registry-03991}\(v_f\) & \texttt{\detokenize{v_f}} & state utility vector of an act & P17 Rationality8.4 & n/a & \hyperref[def:p17-rat84-acts]{\ref*{def:p17-rat84-acts}} & \thepage & Explicit bridge for finite local calculations and the rich utility-vector domain. \\
\phantomsection\label{entry:symbol-registry-03992}\(x_Ey\) & \texttt{\detokenize{x_Ey}} & bet yielding utility x on event E and y otherwise & P17 Rationality8.4 & n/a & \hyperref[def:p17-rat84-acts]{\ref*{def:p17-rat84-acts}} & \thepage & Generalizes the source bet notation one subscript E zero; x and y label consequences of those utilities. \\
\phantomsection\label{entry:symbol-registry-03993}\(\Delta\) & \texttt{\detokenize{\Delta}} & simplex of additive first order state probabilities & P17 Rationality8.4 & n/a & \hyperref[def:p17-rat84-acts]{\ref*{def:p17-rat84-acts}} & \thepage & Distinct from a set of priors and from a second order probability over that simplex. \\
\phantomsection\label{entry:symbol-registry-03994}\(EU_p f\) & \texttt{\detokenize{EU_p f}} & expected utility under the specified first order prior & P17 Rationality8.4 & n/a & \hyperref[def:p17-rat84-acts]{\ref*{def:p17-rat84-acts}} & \thepage & Source notation with the common utility function retained. \\
\phantomsection\label{entry:symbol-registry-03995}\(\succeq\) & \texttt{\detokenize{\succeq}} & weak preference of acts & P17 Rationality8.4 & relation & \hyperref[def:p17-rat84-acts]{\ref*{def:p17-rat84-acts}} & \thepage & Scalar representation is not assumed for the unanimous relation. \\
\phantomsection\label{entry:symbol-registry-03996}\(B_k,W_k,B_u,W_u\) & \texttt{\detokenize{B_k;W_k;B_u;W_u}} & known and unknown urn color events & P17 Rationality8.4 & n/a & \hyperref[sec:p17-rat84-domain]{\ref*{sec:p17-rat84-domain}} & \thepage & Source subscripts distinguish uncertainty sources; complementary event pairs use common prizes. \\
\phantomsection\label{entry:symbol-registry-03997}\(\mathcal C\) & \texttt{\detokenize{\mathcal C}} & nonempty retained set of first order priors & P17 Rationality8.4 & n/a & \hyperref[def:p17-rat84-prior-rules]{\ref*{def:p17-rat84-prior-rules}} & \thepage & Compactness ensures extrema in this finite calculation; interpretation as evidence versus attitude is examined. \\
\phantomsection\label{entry:symbol-registry-03998}\(\alpha\) & \texttt{\detokenize{\alpha}} & fixed pessimistic weight in alpha maxmin & P17 Rationality8.4 & n/a & \hyperref[def:p17-rat84-prior-rules]{\ref*{def:p17-rat84-prior-rules}} & \thepage & One multiplies the minimum and zero the maximum; act dependent generalizations are separately flagged. \\
\phantomsection\label{entry:symbol-registry-03999}\(\mathcal C^*,\mathcal C_0,\mathcal C_1\) & \texttt{\detokenize{\mathcal C^*;\mathcal C_0;\mathcal C_1}} & source supplied or comparative prior sets & P17 Rationality8.4 & n/a & \hyperref[thm:p17-rat84-criteria]{\ref*{thm:p17-rat84-criteria}} & \thepage & Their interval restrictions and the comparative canonical-set roles are explicitly local. \\
\phantomsection\label{entry:symbol-registry-04000}\(\nu\) & \texttt{\detokenize{\nu}} & normalized monotone event capacity & P17 Rationality8.4 & n/a & \hyperref[def:p17-rat84-choquet]{\ref*{def:p17-rat84-choquet}} & \thepage & Not required additive; convex means the stated supermodular event inequality. \\
\phantomsection\label{entry:symbol-registry-04001}\(A_i\) & \texttt{\detokenize{A_i}} & upper utility level event in the Choquet formula & P17 Rationality8.4 & n/a & \hyperref[def:p17-rat84-choquet]{\ref*{def:p17-rat84-choquet}} & \thepage & Greedy core proof also uses initial ordered-state sets with the equivalent zero-increment convention for ties. \\
\phantomsection\label{entry:symbol-registry-04002}\(x_i\) & \texttt{\detokenize{x_i}} & ordered utility level & P17 Rationality8.4 & n/a & \hyperref[def:p17-rat84-choquet]{\ref*{def:p17-rat84-choquet}} & \thepage & Lowest baseline is explicit for signed utilities. \\
\phantomsection\label{entry:symbol-registry-04003}\(\operatorname{core}(\nu)\) & \texttt{\detokenize{\operatorname{core}(\nu)}} & additive probabilities dominating every capacity event value & P17 Rationality8.4 & n/a & \hyperref[def:p17-rat84-choquet]{\ref*{def:p17-rat84-choquet}} & \thepage & Nonemptiness and the minimum identity are proved for convex capacities. \\
\phantomsection\label{entry:symbol-registry-04004}\(\Xi\) & \texttt{\detokenize{\Xi}} & nested family of prior sets representing confidence levels & P17 Rationality8.4 & n/a & \hyperref[sec:p17-rat84-confidence]{\ref*{sec:p17-rat84-confidence}} & \thepage & Larger sets support fewer judgments and correspond to greater confidence in those surviving judgments; not a probability distribution. \\
\phantomsection\label{entry:symbol-registry-04005}\(D(f)\) & \texttt{\detokenize{D(f)}} & decision dependent selected confidence set & P17 Rationality8.4 & n/a & \hyperref[sec:p17-rat84-confidence]{\ref*{sec:p17-rat84-confidence}} & \thepage & A selector with no further restrictions need not preserve dominance. \\
\phantomsection\label{entry:symbol-registry-04006}\(c(p)\) & \texttt{\detokenize{c(p)}} & cardinal penalty in utility units & P17 Rationality8.4 & n/a & \hyperref[sec:p17-rat84-confidence]{\ref*{sec:p17-rat84-confidence}} & \thepage & Real-valued source display and extended indicator special case are distinguished. \\
\phantomsection\label{entry:symbol-registry-04007}\(\varphi(p)\) & \texttt{\detokenize{\varphi(p)}} & fuzzy confidence function on the prior simplex & P17 Rationality8.4 & n/a & \hyperref[sec:p17-rat84-confidence]{\ref*{sec:p17-rat84-confidence}} & \thepage & Published phi role with domain Delta and range zero to one; reciprocal zero and signed-utility boundaries are explicit. \\
\phantomsection\label{entry:symbol-registry-04008}\(w(p(E))\) & \texttt{\detokenize{w(p(E))}} & event probability weighting under a precise prior & P17 Rationality8.4 & n/a & \hyperref[sec:p17-rat84-smooth]{\ref*{sec:p17-rat84-smooth}} & \thepage & The same weighting function cannot distinguish equal-probability bets; source-specific extensions require additional data. \\
\phantomsection\label{entry:symbol-registry-04009}\(\mu\) & \texttt{\detokenize{\mu}} & second order probability over first order priors & P17 Rationality8.4 & n/a & \hyperref[sec:p17-rat84-smooth]{\ref*{sec:p17-rat84-smooth}} & \thepage & Local worked calculations use finite support and normalized weights. \\
\phantomsection\label{entry:symbol-registry-04010}\(\varphi(EU_p f)\) & \texttt{\detokenize{\varphi(EU_p f)}} & smooth ambiguity transform of expected utility & P17 Rationality8.4 & n/a & \hyperref[sec:p17-rat84-smooth]{\ref*{sec:p17-rat84-smooth}} & \thepage & The same published phi glyph now has utility-valued domain; not the earlier fuzzy confidence function on priors. \\
\phantomsection\label{entry:symbol-registry-04011}\(\bar p\) & \texttt{\detokenize{\bar p}} & reduced mean first order prior & P17 Rationality8.4 & n/a & \hyperref[thm:p17-rat84-reduction]{\ref*{thm:p17-rat84-reduction}} & \thepage & Finite mixture identity applies only to linear second order expected utility. \\
\phantomsection\label{entry:symbol-registry-04012}\(k\) & \texttt{\detokenize{k}} & Tversky Fox comparative source preference response-pattern index & P17 Rationality8.4 & n/a & \hyperref[sec:p17-rat84-evidence]{\ref*{sec:p17-rat84-evidence}} & \thepage & Empirical event-and-complement comparison index; not a probability or prior-set size. \\
\phantomsection\label{entry:symbol-registry-04151}\(T,R,P,S\) & \texttt{\detokenize{T,}\allowbreak\detokenize{R,}\allowbreak\detokenize{P,}\allowbreak\detokenize{S}} & temptation mutual-cooperation mutual-defection and sole-cooperator utility levels & P17.rationality-10-2.PD & meta & \hyperref[def:p17-rat102-pd]{\ref*{def:p17-rat102-pd}} & \thepage & Local cardinal payoff levels satisfying T>R>P>S; no interpersonal utility sum. \\
\phantomsection\label{entry:symbol-registry-04152}\(c,d\) & \texttt{\detokenize{c,}\allowbreak\detokenize{d}} & cooperate and defect pure actions & P17.rationality-10-2.PD & object & \hyperref[def:p17-rat102-pd]{\ref*{def:p17-rat102-pd}} & \thepage & The local d action differs from the later binary defense-spending indicator. \\
\phantomsection\label{entry:symbol-registry-04153}\(H,L\) & \texttt{\detokenize{H,}\allowbreak\detokenize{L}} & high and low coordination actions & P17.rationality-10-2.HiLo & object & \hyperref[def:p17-rat102-hilo]{\ref*{def:p17-rat102-hilo}} & \thepage & Not truth values or modal operators. \\
\phantomsection\label{entry:symbol-registry-04154}\(a,b\) & \texttt{\detokenize{a,}\allowbreak\detokenize{b}} & cardinal utilities of coordinated outcomes & P17.rationality-10-2.HiLo & meta & \hyperref[def:p17-rat102-hilo]{\ref*{def:p17-rat102-hilo}} & \thepage & a>b>0 with off-diagonal zero; ordinality alone does not determine mixing. \\
\phantomsection\label{entry:symbol-registry-04155}\(p,q\) & \texttt{\detokenize{p,}\allowbreak\detokenize{q}} & independent probabilities of high coordination action & P17.rationality-10-2.HiLo & meta & \hyperref[def:p17-rat102-hilo]{\ref*{def:p17-rat102-hilo}} & \thepage & The later p in the truth warning has a separate propositional role. \\
\phantomsection\label{entry:symbol-registry-04156}\(q,r\) & \texttt{\detokenize{q,}\allowbreak\detokenize{r}} & opponent cooperation probability and own cooperation mixture weight & P17.rationality-10-2.PD-annex & meta & \hyperref[annex:p17.rat102-pd]{\ref*{annex:p17.rat102-pd}} & \thepage & Distinct local role from HiLo probabilities and threshold. \\
\phantomsection\label{entry:symbol-registry-04157}\(r\) & \texttt{\detokenize{r}} & indifference threshold b/(a+b) & P17.rationality-10-2.HiLo-annex & meta & \hyperref[annex:p17.rat102-hilo]{\ref*{annex:p17.rat102-hilo}} & \thepage & Interior equilibrium probability of H. \\
\phantomsection\label{entry:symbol-registry-04158}\(x,y,z\) & \texttt{\detokenize{x,}\allowbreak\detokenize{y,}\allowbreak\detokenize{z}} & alternatives in majority cycle and paid cyclic exchanges & P17.rationality-10-2.cycle-pump & object & \hyperref[sec:p17-rat102-corporation]{\ref*{sec:p17-rat102-corporation}} & \thepage & Unpriced alternatives alone do not determine priced-pair comparisons. \\
\phantomsection\label{entry:symbol-registry-04159}\(w_0,w,w_j\) & \texttt{\detokenize{w_0,}\allowbreak\detokenize{w,}\allowbreak\detokenize{w_j}} & initial current and post-trade cash & P17.rationality-10-2.pump & meta & \hyperref[thm:p17-rat102-pump]{\ref*{thm:p17-rat102-pump}} & \thepage & Liquidity is an explicit premise. \\
\phantomsection\label{entry:symbol-registry-04160}\(\varepsilon\) & \texttt{\detokenize{\varepsilon}} & strictly positive price per exchange & P17.rationality-10-2.pump & meta & \hyperref[thm:p17-rat102-pump]{\ref*{thm:p17-rat102-pump}} & \thepage & Acceptance is stipulated not inferred from an ordinal cycle. \\
\phantomsection\label{entry:symbol-registry-04161}\(k,j\) & \texttt{\detokenize{k,}\allowbreak\detokenize{j}} & nonnegative round count and sequential trade count & P17.rationality-10-2.pump & meta & \hyperref[thm:p17-rat102-pump]{\ref*{thm:p17-rat102-pump}} & \thepage & Three trades per round; initial cash at least 3k epsilon. \\
\phantomsection\label{entry:symbol-registry-04162}\(t,d,o\) & \texttt{\detokenize{t,}\allowbreak\detokenize{d,}\allowbreak\detokenize{o}} & binary tax defense-spending and other-spending increase indicators & P17.rationality-10-2.policy & object & \hyperref[def:p17-rat102-policy]{\ref*{def:p17-rat102-policy}} & \thepage & Zero spending indicator means reduction; d is not the earlier PD action. \\
\phantomsection\label{entry:symbol-registry-04163}\(p\) & \texttt{\detokenize{p}} & possibly false proposition accepted unanimously & P17.rationality-10-2.truth-warning & object & \hyperref[def:p17-rat102-truth]{\ref*{def:p17-rat102-truth}} & \thepage & Not the local HiLo probability. \\
\phantomsection\label{entry:symbol-registry-04164}\(U(\cdot)\) & \texttt{\detokenize{U(\cdot)}} & expected utility of the selected player's pure action against a fixed opponent mixture & P17.rationality-10-2.PD-annex and HiLo-annex & meta & \hyperref[annex:p17.rat102-pd]{\ref*{annex:p17.rat102-pd}} & \thepage & Shorthand local to each independently specified game; no interpersonal utility aggregation. \\
\phantomsection\label{entry:symbol-registry-04165}\(H,A(H)\) & \texttt{\detokenize{H,}\allowbreak\detokenize{A(H)}} & information set and its available moves & p17/rationality-9-1 & n/a & \hyperref[def:p17-rat91-tree]{\ref*{def:p17-rat91-tree}} & \thepage & A(H) is the same move set at every node of H. \\
\phantomsection\label{entry:symbol-registry-04166}\(s_i,\sigma_i,b_i\) & \texttt{\detokenize{s_i,}\allowbreak\detokenize{\sigma_i,}\allowbreak\detokenize{b_i}} & pure contingent strategy; mixture of pure strategies; behavioral strategy & p17/rationality-9-1 & n/a & \hyperref[def:p17-rat91-strategies]{\ref*{def:p17-rat91-strategies}} & \thepage & Local strategic types; distinct from epistemic type beliefs b\_i elsewhere. \\
\phantomsection\label{entry:symbol-registry-04167}\(p\circ E\oplus q\circ F\) & \texttt{\detokenize{\circ,}\allowbreak\detokenize{\oplus}} & source lottery notation & p17/rationality-9-1 & n/a & \hyperref[def:p17-rat91-strategies]{\ref*{def:p17-rat91-strategies}} & \thepage & Lottery weights; no algebraic direct-sum meaning. \\
\phantomsection\label{entry:symbol-registry-04168}\(r_i(H),r_i(H,a)\) & \texttt{\detokenize{r_i(H),}\allowbreak\detokenize{r_i(H,}\allowbreak\detokenize{a)}} & own-history realization weights & p17/rationality-9-1 & n/a & \hyperref[thm:p17-rat91-realization]{\ref*{thm:p17-rat91-realization}} & \thepage & Perfect-recall mixtures; not actual unconditional node-reach probabilities. \\
\phantomsection\label{entry:symbol-registry-04169}\(M,N\) & \texttt{\detokenize{M,}\allowbreak\detokenize{N}} & source Battle-of-the-Sexes equilibrium lotteries & p17/rationality-9-1 & n/a & \hyperref[def:p17-rat91-outside]{\ref*{def:p17-rat91-outside}} & \thepage & M belongs to player 1; N belongs to player 2. \\
\phantomsection\label{entry:symbol-registry-04170}\(\mu\) & \texttt{\detokenize{\mu}} & conditional probability of left Horse-Game node & p17/rationality-9-1 & n/a & \hyperref[thm:p17-rat91-horse]{\ref*{thm:p17-rat91-horse}} & \thepage & Undefined by ordinary conditioning at a zero-probability information set. \\
\phantomsection\label{entry:symbol-registry-04171}\(\delta\) & \texttt{\detokenize{\delta}} & discount factor & p17/rationality-9-1 & n/a & \hyperref[thm:p17-rat91-repetition]{\ref*{thm:p17-rat91-repetition}} & \thepage & 0 <= delta < 1 in the explicitly infinite repetition comparison. \\
\phantomsection\label{entry:symbol-registry-04172}\(N\) & \texttt{\detokenize{N}} & finite number of strategic players & P17 Rationality9.3 & n/a & \hyperref[def:p17-rat93-domain]{\ref*{def:p17-rat93-domain}} & \thepage & Not the number of strategies in the one-population model. \\
\phantomsection\label{entry:symbol-registry-04173}\(S_i\) & \texttt{\detokenize{S_i}} & pure strategy set for strategic player i & P17 Rationality9.3 strategic domain & n/a & \hyperref[def:p17-rat93-domain]{\ref*{def:p17-rat93-domain}} & \thepage & Published source notation; the same indexed glyph denotes a pure strategy in the later one-population domain. \\
\phantomsection\label{entry:symbol-registry-04174}\(\Delta_i\) & \texttt{\detokenize{\Delta_i}} & individual mixed strategy simplex for player i & P17 Rationality9.3 strategic domain & n/a & \hyperref[def:p17-rat93-domain]{\ref*{def:p17-rat93-domain}} & \thepage & Distinct typed role from Delta\_n as a population-frequency simplex. \\
\phantomsection\label{entry:symbol-registry-04175}\(\sigma_i\) & \texttt{\detokenize{\sigma_i}} & individual mixed strategy of player i & P17 Rationality9.3 strategic domain & n/a & \hyperref[def:p17-rat93-domain]{\ref*{def:p17-rat93-domain}} & \thepage & Bridge to finite-strategic-games x\_i is explicit. \\
\phantomsection\label{entry:symbol-registry-04176}\(\vec\sigma\) & \texttt{\detokenize{\vec\sigma}} & profile of individual mixed strategies & P17 Rationality9.3 strategic domain & n/a & \hyperref[def:p17-rat93-domain]{\ref*{def:p17-rat93-domain}} & \thepage & Not a population-frequency vector. \\
\phantomsection\label{entry:symbol-registry-04177}\(\pi_i(\vec\sigma)\) & \texttt{\detokenize{\pi_i}} & player i expected payoff function on strategic profiles & P17 Rationality9.3 strategic domain & n/a & \hyperref[def:p17-rat93-domain]{\ref*{def:p17-rat93-domain}} & \thepage & Published source payoff notation bridge to U\_i; not the later numerical payoff vector pi. \\
\phantomsection\label{entry:symbol-registry-04178}\(S\) & \texttt{\detokenize{S}} & finite repertoire of one-population pure strategies & P17 Rationality9.3 population domain & n/a & \hyperref[def:p17-rat93-domain]{\ref*{def:p17-rat93-domain}} & \thepage & One-population model fixes role symmetry and random matching separately. \\
\phantomsection\label{entry:symbol-registry-04179}\(S_i\) & \texttt{\detokenize{S_i}} & pure strategy i in the single-population repertoire & P17 Rationality9.3 population domain & n/a & \hyperref[def:p17-rat93-domain]{\ref*{def:p17-rat93-domain}} & \thepage & Repeated published glyph with distinct role from the strategic player action set. \\
\phantomsection\label{entry:symbol-registry-04180}\(n\) & \texttt{\detokenize{n}} & number of pure strategies in a population repertoire & P17 Rationality9.3 population domain & n/a & \hyperref[def:p17-rat93-domain]{\ref*{def:p17-rat93-domain}} & \thepage & Delta\_n has n coordinates; no shift to n-dimensional simplex convention. \\
\phantomsection\label{entry:symbol-registry-04181}\(\Delta_n\) & \texttt{\detokenize{\Delta_n}} & simplex of mixed strategies or population frequencies under explicitly distinguished interpretations & P17 Rationality9.3 population domain & n/a & \hyperref[def:p17-rat93-domain]{\ref*{def:p17-rat93-domain}} & \thepage & Same vector algebra does not equate individual randomization and pure-user frequency. \\
\phantomsection\label{entry:symbol-registry-04182}\(A_{ij}\) & \texttt{\detokenize{A_{ij}}} & row user payoff of strategy i against strategy j & P17 Rationality9.3 matrix games & n/a & \hyperref[def:p17-rat93-domain]{\ref*{def:p17-rat93-domain}} & \thepage & Role-symmetric game has payoff pair A\_ij and A\_ji; A need not be symmetric as a matrix. \\
\phantomsection\label{entry:symbol-registry-04183}\(\pi(\sigma\mid\mu)\) & \texttt{\detokenize{\pi}} & bilinear payoff function of own and opponent distributions & P17 Rationality9.3 matrix games & n/a & \hyperref[def:p17-rat93-domain]{\ref*{def:p17-rat93-domain}} & \thepage & Published vertical-bar payoff syntax is not conditional probability. \\
\phantomsection\label{entry:symbol-registry-04184}\(\sigma,\mu,q\) & \texttt{\detokenize{\sigma;\mu;q}} & incumbent mutant or comparison probability distributions & P17 Rationality9.3 matrix games & n/a & \hyperref[def:p17-rat93-ess]{\ref*{def:p17-rat93-ess}} & \thepage & Each role is stated in its local comparison; source sigma may also be interpreted as population state. \\
\phantomsection\label{entry:symbol-registry-04185}\(\operatorname{supp}(\sigma)\) & \texttt{\detokenize{\operatorname{supp}}} & indices with positive mixture probabilities & P17 Rationality9.3 & n/a & \hyperref[def:p17-rat93-weak-dominance]{\ref*{def:p17-rat93-weak-dominance}} & \thepage & Tangent-vector support is defined separately by nonzero signed coordinates. \\
\phantomsection\label{entry:symbol-registry-04186}\(\varepsilon\) & \texttt{\detokenize{\varepsilon}} & mutant population fraction & P17 Rationality9.3 ESS definition & n/a & \hyperref[def:p17-rat93-ess]{\ref*{def:p17-rat93-ess}} & \thepage & Distinct local role from stability tolerance and final small-basin game parameter. \\
\phantomsection\label{entry:symbol-registry-04187}\(\varepsilon_\mu\) & \texttt{\detokenize{\varepsilon_\mu}} & mutant-dependent invasion threshold & P17 Rationality9.3 ESS definition & n/a & \hyperref[def:p17-rat93-ess]{\ref*{def:p17-rat93-ess}} & \thepage & Quantifier permits dependence on mutant; local-superiority equivalence is proved independently. \\
\phantomsection\label{entry:symbol-registry-04188}\(V,C\) & \texttt{\detokenize{V;C}} & resource benefit and injury cost in Hawk Dove & P17 Rationality9.3 Hawk Dove & n/a & \hyperref[def:p17-rat93-source-static-games]{\ref*{def:p17-rat93-source-static-games}} & \thepage & Standing assumptions zero less than V less than C. \\
\phantomsection\label{entry:symbol-registry-04189}\(\vec p(t)\) & \texttt{\detokenize{\vec p}} & time dependent vector of pure-user population shares & P17 Rationality9.3 population dynamics & n/a & \hyperref[def:p17-rat93-domain]{\ref*{def:p17-rat93-domain}} & \thepage & Not an individual lottery or a subjective conjecture; scalar p\_i is share. \\
\phantomsection\label{entry:symbol-registry-04190}\(W_i(\vec p)\) & \texttt{\detokenize{W_i}} & payoff or fitness of pure strategy i at population state & P17 Rationality9.3 population dynamics & n/a & \hyperref[def:p17-rat93-domain]{\ref*{def:p17-rat93-domain}} & \thepage & Linear W=Ap is random-matching specialization; Lipschitz nonlinear models require their own hypotheses. \\
\phantomsection\label{entry:symbol-registry-04191}\(\rho_{ij}(\pi,\vec p)\) & \texttt{\detokenize{\rho_{ij}}} & nonnegative conditional strategy-switch rate & P17 Rationality9.3 revision protocols & n/a & \hyperref[def:p17-rat93-revision]{\ref*{def:p17-rat93-revision}} & \thepage & Diagonal terms cancel; rates need not lie below one. \\
\phantomsection\label{entry:symbol-registry-04192}\(\pi\) & \texttt{\detokenize{\pi}} & numerical payoff vector supplied to a revision protocol & P17 Rationality9.3 revision protocols & n/a & \hyperref[def:p17-rat93-revision]{\ref*{def:p17-rat93-revision}} & \thepage & Repeated source glyph now in R\^{}n rather than a function; declared before use. \\
\phantomsection\label{entry:symbol-registry-04193}\(\phi(\vec p)\) & \texttt{\detokenize{\phi}} & mean population payoff & P17 Rationality9.3 population dynamics & n/a & \hyperref[def:p17-rat93-revision]{\ref*{def:p17-rat93-revision}} & \thepage & Weighted mean under current shares; not a revision-response function or welfare norm. \\
\phantomsection\label{entry:symbol-registry-04194}\([a]_+\) & \texttt{\detokenize{[a]_+}} & positive part of a real number & P17 Rationality9.3 revision protocols & n/a & \hyperref[def:p17-rat93-revision]{\ref*{def:p17-rat93-revision}} & \thepage & Maximum of a and zero. \\
\phantomsection\label{entry:symbol-registry-04195}\(\dot p_i\) & \texttt{\detokenize{\dot p_i}} & time derivative of population share & P17 Rationality9.3 population dynamics & n/a & \hyperref[def:p17-rat93-revision]{\ref*{def:p17-rat93-revision}} & \thepage & Continuous mean dynamic rather than finite stochastic sample path. \\
\phantomsection\label{entry:symbol-registry-04196}\(\phi_j\) & \texttt{\detokenize{\phi_j}} & sign-preserving strategy-specific switch response function & P17 Rationality9.3 comparison families & n/a & \hyperref[def:p17-rat93-comparison-families]{\ref*{def:p17-rat93-comparison-families}} & \thepage & Published primary glyph with scalar domain; independent of incumbent index but not incumbent payoff in the difference. \\
\phantomsection\label{entry:symbol-registry-04197}\(\phi\) & \texttt{\detokenize{\phi}} & common primary-paper scalar switch response function & P17 Rationality9.3 survival theorem & n/a & \hyperref[def:p17-rat93-comparison-families]{\ref*{def:p17-rat93-comparison-families}} & \thepage & Repeated source phi now has domain R rather than population simplex; common Lipschitz response with positive right derivative at zero. \\
\phantomsection\label{entry:symbol-registry-04198}\(\widetilde A_d\) & \texttt{\detokenize{\widetilde A_d}} & feeble twin payoff matrix & P17 Rationality9.3 survival construction & n/a & \hyperref[thm:p17-rat93-survival]{\ref*{thm:p17-rat93-survival}} & \thepage & Primary corrected symmetric-game construction; fourth row uniformly below third by d. \\
\phantomsection\label{entry:symbol-registry-04199}\(a,b,d\) & \texttt{\detokenize{a;b;d}} & win payoff loss magnitude and positive twin decrement & P17 Rationality9.3 survival construction & n/a & \hyperref[thm:p17-rat93-survival]{\ref*{thm:p17-rat93-survival}} & \thepage & Standing b greater than a greater than zero and sufficiently small positive d for survival theorem. \\
\phantomsection\label{entry:symbol-registry-04200}\(\varepsilon,\delta\) & \texttt{\detokenize{\varepsilon;\delta}} & distance tolerances in Lyapunov stability & P17 Rationality9.3 stability definition & n/a & \hyperref[def:p17-rat93-stability]{\ref*{def:p17-rat93-stability}} & \thepage & Local repeated source eps glyph is not mutant share or game parameter. \\
\phantomsection\label{entry:symbol-registry-04201}\(x\) & \texttt{\detokenize{x}} & signed zero-sum tangent vector on equilibrium support face & P17 Rationality9.3 regular ESS definition & n/a & \hyperref[def:p17-rat93-regular]{\ref*{def:p17-rat93-regular}} & \thepage & Not required nonnegative; support contained in rather than properly contained in equilibrium support. \\
\phantomsection\label{entry:symbol-registry-04202}\(p_i^R,p_i^C\) & \texttt{\detokenize{p_i^R;p_i^C}} & pure-user frequencies in Row and Column populations & P17 Rationality9.3 two-role model & n/a & \hyperref[def:p17-rat93-two-populations]{\ref*{def:p17-rat93-two-populations}} & \thepage & Independent role populations couple through expected payoffs. \\
\phantomsection\label{entry:symbol-registry-04203}\(W_i^R,W_i^C\) & \texttt{\detokenize{W_i^R;W_i^C}} & role-specific strategy payoffs against opposite population & P17 Rationality9.3 two-role model & n/a & \hyperref[def:p17-rat93-two-populations]{\ref*{def:p17-rat93-two-populations}} & \thepage & Row and Column fitness maps differ. \\
\phantomsection\label{entry:symbol-registry-04204}\(\phi^R,\phi^C\) & \texttt{\detokenize{\phi^R;\phi^C}} & role-specific mean payoffs depending on both populations & P17 Rationality9.3 two-role model & n/a & \hyperref[def:p17-rat93-two-populations]{\ref*{def:p17-rat93-two-populations}} & \thepage & Source abbreviated argument notation expanded to actual joint dependence. \\
\phantomsection\label{entry:symbol-registry-04205}\(x,y,z\) & \texttt{\detokenize{x;y;z}} & locally specified scalar population shares & P17 Rationality9.3 worked examples & n/a & \hyperref[thm:p17-rat93-static-calculations]{\ref*{thm:p17-rat93-static-calculations}} & \thepage & Local role declared in each example; not regular-ESS signed tangent vector x. \\
\phantomsection\label{entry:symbol-registry-04206}\(A_\varepsilon\) & \texttt{\detokenize{A_\varepsilon}} & three-strategy payoff matrix with small basin parameter & P17 Rationality9.3 basin example & n/a & \hyperref[def:p17-rat93-basins]{\ref*{def:p17-rat93-basins}} & \thepage & Parameter zero less than eps less than one half; distinct from mutant fraction and stability tolerance. \\
\phantomsection\label{entry:symbol-registry-04207}\(V(p)\) & \texttt{\detokenize{V}} & relative-entropy local Lyapunov function & P17 Rationality9.3 annex stability bridge & n/a & \hyperref[annex:p17.rat93-stability-bridges]{\ref*{annex:p17.rat93-stability-bridges}} & \thepage & Local function role distinct from Hawk Dove resource value V. \\
\phantomsection\label{entry:symbol-registry-04208}\(Z(t)\) & \texttt{\detokenize{Z}} & weighted share-ratio witness for strict mixed domination & P17 Rationality9.3 annex domination & n/a & \hyperref[annex:p17.rat93-replicator-dominance]{\ref*{annex:p17.rat93-replicator-dominance}} & \thepage & All used shares positive at finite times; exponential bound proves extinction. \\
\phantomsection\label{entry:symbol-registry-04209}\(C_I\) & \texttt{\detokenize{C_I}} & closed cone of feasible critical tangent directions & P17 Rationality9.3 annex local-superiority proof & n/a & \hyperref[annex:p17.rat93-ess-equivalence]{\ref*{annex:p17.rat93-ess-equivalence}} & \thepage & Compact normalized section gives uniform quadratic bound including boundary cases. \\
\phantomsection\label{entry:symbol-registry-04210}\(a,b\) & \texttt{\detokenize{a;b}} & incumbent-minus-mutant payoff differences against incumbent and mutant & P17 Rationality9.3 annex ESS equivalence & n/a & \hyperref[annex:p17.rat93-ess-equivalence]{\ref*{annex:p17.rat93-ess-equivalence}} & \thepage & Local scalar roles distinct from win and loss magnitudes in the survival matrix. \\
\phantomsection\label{entry:symbol-registry-04211}\(v,d_i,I\) & \texttt{\detokenize{v;d_i;I}} & incumbent self payoff pure-strategy disadvantage and tied-best-response index set & P17 Rationality9.3 annex local-superiority proof & n/a & \hyperref[annex:p17.rat93-ess-equivalence]{\ref*{annex:p17.rat93-ess-equivalence}} & \thepage & Support of incumbent is contained in I. \\
\phantomsection\label{entry:symbol-registry-04212}\(h,k,r,z\) & \texttt{\detokenize{h;k;r;z}} & signed nearby displacement critical-face component outside-face correction and tangent direction & P17 Rationality9.3 annex local-superiority proof & n/a & \hyperref[annex:p17.rat93-ess-equivalence]{\ref*{annex:p17.rat93-ess-equivalence}} & \thepage & Each local role is defined in the proof; z is not a population share here. \\
\phantomsection\label{entry:symbol-registry-04213}\(m,c,K,\alpha\) & \texttt{\detokenize{m;c;K;\alpha}} & outside-critical-face mass uniform quadratic constant norm-bound constant and positive feasible-ray scale & P17 Rationality9.3 annex local-superiority proof & n/a & \hyperref[annex:p17.rat93-ess-equivalence]{\ref*{annex:p17.rat93-ess-equivalence}} & \thepage & Finite dimensional compactness supplies constants; none is an empirical parameter. \\
\phantomsection\label{entry:symbol-registry-04214}\(\beta_i,B\) & \texttt{\detokenize{\beta_i;B}} & positive excess payoff and its total & P17 Rationality9.3 annex BNN rest proof & n/a & \hyperref[annex:p17.rat93-protocols]{\ref*{annex:p17.rat93-protocols}} & \thepage & Rest-point contradiction uses population averaging. \\
\phantomsection\label{entry:symbol-registry-04215}\(\eta\) & \texttt{\detokenize{\eta}} & positive uniform payoff gap or local lower growth bound & P17 Rationality9.3 annex dominance and stable rest proofs & n/a & \hyperref[annex:p17.rat93-replicator-dominance]{\ref*{annex:p17.rat93-replicator-dominance}} & \thepage & Compact domination gap and local continuity gap are distinct local instances. \\
\phantomsection\label{entry:symbol-registry-04216}\(L\) & \texttt{\detokenize{L}} & positive limiting ratio of weakly dominated share to dominating share & P17 Rationality9.3 annex weak dominance example & n/a & \hyperref[annex:p17.rat93-replicator-dominance]{\ref*{annex:p17.rat93-replicator-dominance}} & \thepage & Finite integral of disappearing third share keeps ratio positive. \\
\phantomsection\label{entry:symbol-registry-04217}\(r\) & \texttt{\detokenize{r}} & positive ratio of first to third strategy shares & P17 Rationality9.3 annex global-attraction example & n/a & \hyperref[annex:p17.rat93-global-unstable]{\ref*{annex:p17.rat93-global-unstable}} & \thepage & Local scalar role distinct from outside-face correction vector in ESS proof. \\
\phantomsection\label{entry:symbol-registry-04305}\(A,B,C,D\) & \texttt{\detokenize{A,}\allowbreak\detokenize{B,}\allowbreak\detokenize{C,}\allowbreak\detokenize{D}} & macro-conditions micro-conditions micro-outcomes and macro-outcomes as proposition classes & P17.rationality-10-4.Coleman & meta & \hyperref[def:p17-rat104-coleman]{\ref*{def:p17-rat104-coleman}} & \thepage & Not random events; actor/system levels are relative to the explanatory question. \\
\phantomsection\label{entry:symbol-registry-04306}\(\longrightarrow\) & \texttt{\detokenize{\longrightarrow}} & explanatory-model link between proposition classes & P17.rationality-10-4.Coleman & relation & \hyperref[def:p17-rat104-coleman]{\ref*{def:p17-rat104-coleman}} & \thepage & Bridge action transformation and macro relation are distinct; no causal identification by arrow alone. \\
\phantomsection\label{entry:symbol-registry-04307}\(N\) & \texttt{\detokenize{N}} & number of actors in symmetric volunteer game & P17.rationality-10-4.VOD & meta & \hyperref[def:p17-rat104-vod]{\ref*{def:p17-rat104-vod}} & \thepage & Integer at least two in the actual game; real extension declared separately for derivative analysis. \\
\phantomsection\label{entry:symbol-registry-04308}\(U,K\) & \texttt{\detokenize{U,}\allowbreak\detokenize{K}} & common cardinal consequence benefit and contribution cost & P17.rationality-10-4.VOD & meta & \hyperref[def:p17-rat104-vod]{\ref*{def:p17-rat104-vod}} & \thepage & U>K>0; not utility function and knowledge operator. \\
\phantomsection\label{entry:symbol-registry-04309}\(\mathrm{CONTR},\text{DON'T}\) & \texttt{\detokenize{\mathrm{CONTR},}\allowbreak\detokenize{\text{DON'T}}} & contribute and abstain pure actions & P17.rationality-10-4.VOD & object & \hyperref[def:p17-rat104-vod]{\ref*{def:p17-rat104-vod}} & \thepage & Mixed cooperation/defection labels later concern distributions rather than these pure actions. \\
\phantomsection\label{entry:symbol-registry-04310}\(p_i,p,p^*,p^{**}\) & \texttt{\detokenize{p_i,}\allowbreak\detokenize{p,}\allowbreak\detokenize{p^*,}\allowbreak\detokenize{p^{**}}} & individual contribution probabilities and two common-profile probabilities & P17.rationality-10-4.VOD & meta & \hyperref[def:p17-rat104-vod]{\ref*{def:p17-rat104-vod}} & \thepage & p* is the symmetric Nash probability; p** is the symmetric joint-improvement probability. \\
\phantomsection\label{entry:symbol-registry-04311}\(P,P^*,P^{**}\) & \texttt{\detokenize{P,}\allowbreak\detokenize{P^*,}\allowbreak\detokenize{P^{**}}} & system-level collective-good provision probabilities & P17.rationality-10-4.VOD & meta & \hyperref[def:p17-rat104-vod]{\ref*{def:p17-rat104-vod}} & \thepage & Different denominators and explanatory level from individual p. \\
\phantomsection\label{entry:symbol-registry-04312}\(E_i\) & \texttt{\detokenize{E_i}} & actor i expected consequence utility & P17.rationality-10-4.VOD & meta & \hyperref[def:p17-rat104-vod]{\ref*{def:p17-rat104-vod}} & \thepage & E\_i=UP-Kp\_i; no moral evaluation or observed-study estimate. \\
\phantomsection\label{entry:symbol-registry-04313}\(i,j\) & \texttt{\detokenize{i,}\allowbreak\detokenize{j}} & actor indices & P17.rationality-10-4.VOD & meta & \hyperref[def:p17-rat104-vod]{\ref*{def:p17-rat104-vod}} & \thepage & j attains maximum contribution probability in the Pareto proof. \\
\phantomsection\label{entry:symbol-registry-04314}\(a,q\) & \texttt{\detokenize{a,}\allowbreak\detokenize{q}} & abstention probability at p** and cost-factor shorthand & P17.rationality-10-4.cooperative-mixture & meta & \hyperref[thm:p17-rat104-cooperative-mixture]{\ref*{thm:p17-rat104-cooperative-mixture}} & \thepage & a=1-p** and q=1-(N-1)a/N; q is not a second actor's mixture. \\
\phantomsection\label{entry:symbol-registry-04315}\(r\) & \texttt{\detokenize{r}} & fixed cost-to-benefit ratio K/U & P17.rationality-10-4.VOD-proofs & meta & \hyperref[thm:p17-rat104-provision-correction]{\ref*{thm:p17-rat104-provision-correction}} & \thepage & 0<r<1; held fixed for all group-size comparisons. \\
\phantomsection\label{entry:symbol-registry-04316}\(F\) & \texttt{\detokenize{F}} & common expected-utility function of a symmetric independent contribution probability & P17.rationality-10-4.cooperative-annex & meta & \hyperref[annex:p17.rat104-cooperative-mixture]{\ref*{annex:p17.rat104-cooperative-mixture}} & \thepage & F(p)=U[1-(1-p)\^{}N]-Kp; not an aggregation rule. \\
\phantomsection\label{entry:symbol-registry-04317}\(s,m\) & \texttt{\detokenize{s,}\allowbreak\detokenize{m}} & symmetric utility maximizer and maximum contribution probability of a comparison profile & P17.rationality-10-4.cooperative-annex & meta & \hyperref[annex:p17.rat104-cooperative-mixture]{\ref*{annex:p17.rat104-cooperative-mixture}} & \thepage & s=p** and m=max\_i p\_i; roles are local to the Pareto argument. \\
\phantomsection\label{entry:symbol-registry-04318}\(N\) & \texttt{\detokenize{N}} & real group-size parameter in continuous extension only & P17.rationality-10-4.correction-annex & meta & \hyperref[annex:p17.rat104-provision-correction]{\ref*{annex:p17.rat104-provision-correction}} & \thepage & N>1; extension is an analytic device not a fractional actor count. \\
\phantomsection\label{entry:symbol-registry-04319}\(f,h\) & \texttt{\detokenize{f,}\allowbreak\detokenize{h}} & log failure probability and derivative-sign auxiliary functions & P17.rationality-10-4.correction-annex & meta & \hyperref[annex:p17.rat104-provision-correction]{\ref*{annex:p17.rat104-provision-correction}} & \thepage & f=N log(r/N)/(N-1); h=N-1-log(N/r). \\
\phantomsection\label{entry:symbol-registry-04320}\(\log,\exp,e\) & \texttt{\detokenize{\log,}\allowbreak\detokenize{\exp,}\allowbreak\detokenize{e}} & natural logarithm exponential and its base & P17.rationality-10-4.VOD-proofs & n/a & \hyperref[thm:p17-rat104-provision-correction]{\ref*{thm:p17-rat104-provision-correction}} & \thepage & Standard real analytic functions; logarithm arguments positive. \\
\phantomsection\label{entry:symbol-registry-07335}\(S\) & \texttt{\detokenize{S}} & state consequence and event-algebra domains symbol family member \(S\) & P17 Savage acts & object & \hyperref[def:p17-sa-acts]{\ref*{def:p17-sa-acts}} & \thepage & General domains; not finite monetary vectors. Explicit displayed family; member types and common actual first-use exposed, not one atomic glyph. Family listing, not one atomic glyph. Author macro \texttt{\detokenize{S;X;\mathcal E}} not carried to the members. \\
\phantomsection\label{entry:symbol-registry-07336}\(X\) & \texttt{\detokenize{X}} & state consequence and event-algebra domains symbol family member \(X\) & P17 Savage acts & object & \hyperref[def:p17-sa-acts]{\ref*{def:p17-sa-acts}} & \thepage & General domains; not finite monetary vectors. Explicit displayed family; member types and common actual first-use exposed, not one atomic glyph. Family listing, not one atomic glyph. Author macro \texttt{\detokenize{S;X;\mathcal E}} not carried to the members. \\
\phantomsection\label{entry:symbol-registry-07337}\(\mathcal E\) & \texttt{\detokenize{\mathcal E}} & state consequence and event-algebra domains symbol family member \(\mathcal E\) & P17 Savage acts & object & \hyperref[def:p17-sa-acts]{\ref*{def:p17-sa-acts}} & \thepage & General domains; not finite monetary vectors. Explicit displayed family; member types and common actual first-use exposed, not one atomic glyph. Family listing, not one atomic glyph. Author macro \texttt{\detokenize{S;X;\mathcal E}} not carried to the members. \\
\phantomsection\label{entry:symbol-registry-07338}\(f_Ag\) & \texttt{\detokenize{none}} & act pasting on A and its complement & P17 Savage acts & object & \hyperref[def:p17-sa-acts]{\ref*{def:p17-sa-acts}} & \thepage & Constants are acts; pasting differs from restriction. \\
\phantomsection\label{entry:symbol-registry-07339}\(\succeq_A\) & \texttt{\detokenize{\succeq_A}} & conditional preference and indifference symbol family member \(\succeq_A\) & P17 Savage acts & meta & \hyperref[def:p17-sa-acts]{\ref*{def:p17-sa-acts}} & \thepage & Background invariance requires P2. Explicit displayed family; member types and common actual first-use exposed, not one atomic glyph. Family listing, not one atomic glyph. Author macro \texttt{\detokenize{\succeq;\sim}} not carried to the members. \\
\phantomsection\label{entry:symbol-registry-07340}\(\sim_A\) & \texttt{\detokenize{\sim_A}} & conditional preference and indifference symbol family member \(\sim_A\) & P17 Savage acts & meta & \hyperref[def:p17-sa-acts]{\ref*{def:p17-sa-acts}} & \thepage & Background invariance requires P2. Explicit displayed family; member types and common actual first-use exposed, not one atomic glyph. Family listing, not one atomic glyph. Author macro \texttt{\detokenize{\succeq;\sim}} not carried to the members. \\
\phantomsection\label{entry:symbol-registry-07341}\(P1--P7\) & \texttt{\detokenize{none}} & source postulates & P17 Savage acts & meta & \hyperref[def:p17-sa-postulates]{\ref*{def:p17-sa-postulates}} & \thepage & P6 constant replacement and arbitrary-act replacement differ; P7 concerns extension. \\
\phantomsection\label{entry:symbol-registry-07342}\(\succeq_e\) & \texttt{\detokenize{\succeq_e}} & comparative event likelihood symbol family member \(\succeq_e\) & P17 CUP construction & meta & \hyperref[def:p17-sa-construction-axioms]{\ref*{def:p17-sa-construction-axioms}} & \thepage & Derived from stake-independent event bets. Explicit displayed family; member types and common actual first-use exposed, not one atomic glyph. Family listing, not one atomic glyph. Author macro \texttt{\detokenize{\succeq;\sim}} not carried to the members. \\
\phantomsection\label{entry:symbol-registry-07343}\(\sim_e\) & \texttt{\detokenize{\sim_e}} & comparative event likelihood symbol family member \(\sim_e\) & P17 CUP construction & meta & \hyperref[def:p17-sa-construction-axioms]{\ref*{def:p17-sa-construction-axioms}} & \thepage & Derived from stake-independent event bets. Explicit displayed family; member types and common actual first-use exposed, not one atomic glyph. Family listing, not one atomic glyph. Author macro \texttt{\detokenize{\succeq;\sim}} not carried to the members. \\
\phantomsection\label{entry:symbol-registry-07344}\(P6.1/P6.2\) & \texttt{\detokenize{none}} & 2020 Archimedean and event-solvability assumptions & P17 2020 atomless branch & meta & \hyperref[def:p17-sa-construction-axioms]{\ref*{def:p17-sa-construction-axioms}} & \thepage & Not renamed original P6; absence of atoms is explicit. \\
\phantomsection\label{entry:symbol-registry-07345}\([j]_A\) & \texttt{\detokenize{none}} & union of j CUP cells & P17 CUP construction & object & \hyperref[def:p17-sa-cup]{\ref*{def:p17-sa-cup}} & \thepage & Not an equivalence class. \\
\phantomsection\label{entry:symbol-registry-07346}\(I_A(E)\) & \texttt{\detokenize{I_A(E)}} & conservative intervals and coarse/fine copy counts symbol family member \(I_A(E)\) & P17 CUP consistency & meta & \hyperref[def:p17-sa-cup]{\ref*{def:p17-sa-cup}} & \thepage & Weak lower cuts and numerical bounds are explicit. Explicit displayed family; member types and common actual first-use exposed, not one atomic glyph. Family listing, not one atomic glyph. Author macro \texttt{\detokenize{I\_A;n;N;i;k;\ell}} not carried to the members. \\
\phantomsection\label{entry:symbol-registry-07347}\(n\) & \texttt{\detokenize{n}} & conservative intervals and coarse/fine copy counts symbol family member \(n\) & P17 CUP consistency & meta & \hyperref[def:p17-sa-cup]{\ref*{def:p17-sa-cup}} & \thepage & Weak lower cuts and numerical bounds are explicit. Explicit displayed family; member types and common actual first-use exposed, not one atomic glyph. Family listing, not one atomic glyph. Author macro \texttt{\detokenize{I\_A;n;N;i;k;\ell}} not carried to the members. \\
\phantomsection\label{entry:symbol-registry-07348}\(N\) & \texttt{\detokenize{N}} & conservative intervals and coarse/fine copy counts symbol family member \(N\) & P17 CUP consistency & meta & \hyperref[def:p17-sa-cup]{\ref*{def:p17-sa-cup}} & \thepage & Weak lower cuts and numerical bounds are explicit. Explicit displayed family; member types and common actual first-use exposed, not one atomic glyph. Family listing, not one atomic glyph. Author macro \texttt{\detokenize{I\_A;n;N;i;k;\ell}} not carried to the members. \\
\phantomsection\label{entry:symbol-registry-07349}\(i\) & \texttt{\detokenize{i}} & conservative intervals and coarse/fine copy counts symbol family member \(i\) & P17 CUP consistency & meta & \hyperref[def:p17-sa-cup]{\ref*{def:p17-sa-cup}} & \thepage & Weak lower cuts and numerical bounds are explicit. Explicit displayed family; member types and common actual first-use exposed, not one atomic glyph. Family listing, not one atomic glyph. Author macro \texttt{\detokenize{I\_A;n;N;i;k;\ell}} not carried to the members. \\
\phantomsection\label{entry:symbol-registry-07350}\(k\) & \texttt{\detokenize{k}} & conservative intervals and coarse/fine copy counts symbol family member \(k\) & P17 CUP consistency & meta & \hyperref[def:p17-sa-cup]{\ref*{def:p17-sa-cup}} & \thepage & Weak lower cuts and numerical bounds are explicit. Explicit displayed family; member types and common actual first-use exposed, not one atomic glyph. Family listing, not one atomic glyph. Author macro \texttt{\detokenize{I\_A;n;N;i;k;\ell}} not carried to the members. \\
\phantomsection\label{entry:symbol-registry-07351}\(\ell\) & \texttt{\detokenize{\ell}} & conservative intervals and coarse/fine copy counts symbol family member \(\ell\) & P17 CUP consistency & meta & \hyperref[def:p17-sa-cup]{\ref*{def:p17-sa-cup}} & \thepage & Weak lower cuts and numerical bounds are explicit. Explicit displayed family; member types and common actual first-use exposed, not one atomic glyph. Family listing, not one atomic glyph. Author macro \texttt{\detokenize{I\_A;n;N;i;k;\ell}} not carried to the members. \\
\phantomsection\label{entry:symbol-registry-07352}\(P\) & \texttt{\detokenize{P}} & constructed probability utility and conditional calibration measure symbol family member \(P\) & P17 simple-act representation & object & \hyperref[thm:p17-sa-utility]{\ref*{thm:p17-sa-utility}} & \thepage & T\_x(A)=U(x)P(A) is proved. Explicit displayed family; member types and common actual first-use exposed, not one atomic glyph. Family listing, not one atomic glyph. Author macro \texttt{\detokenize{P;U;T\_x}} not carried to the members. \\
\phantomsection\label{entry:symbol-registry-07353}\(U\) & \texttt{\detokenize{U}} & constructed probability utility and conditional calibration measure symbol family member \(U\) & P17 simple-act representation & object & \hyperref[thm:p17-sa-utility]{\ref*{thm:p17-sa-utility}} & \thepage & T\_x(A)=U(x)P(A) is proved. Explicit displayed family; member types and common actual first-use exposed, not one atomic glyph. Family listing, not one atomic glyph. Author macro \texttt{\detokenize{P;U;T\_x}} not carried to the members. \\
\phantomsection\label{entry:symbol-registry-07354}\(T_x\) & \texttt{\detokenize{T_x}} & constructed probability utility and conditional calibration measure symbol family member \(T_x\) & P17 simple-act representation & object & \hyperref[thm:p17-sa-utility]{\ref*{thm:p17-sa-utility}} & \thepage & T\_x(A)=U(x)P(A) is proved. Explicit displayed family; member types and common actual first-use exposed, not one atomic glyph. Family listing, not one atomic glyph. Author macro \texttt{\detokenize{P;U;T\_x}} not carried to the members. \\
\phantomsection\label{entry:symbol-registry-07355}\(I_P\) & \texttt{\detokenize{I_P}} & bounded finite-additive integral and indicator symbol family member \(I_P\) & P17 arbitrary-act converse & object & \hyperref[def:p17-sa-integral]{\ref*{def:p17-sa-integral}} & \thepage & No countable-additivity premise. Explicit displayed family; member types and common actual first-use exposed, not one atomic glyph. Family listing, not one atomic glyph. Author macro \texttt{\detokenize{I\_P;\mathbf1}} not carried to the members. \\
\phantomsection\label{entry:symbol-registry-07356}\(\mathbf1_A\) & \texttt{\detokenize{\mathbf1_A}} & bounded finite-additive integral and indicator symbol family member \(\mathbf1_A\) & P17 arbitrary-act converse & object & \hyperref[def:p17-sa-integral]{\ref*{def:p17-sa-integral}} & \thepage & No countable-additivity premise. Explicit displayed family; member types and common actual first-use exposed, not one atomic glyph. Family listing, not one atomic glyph. Author macro \texttt{\detokenize{I\_P;\mathbf1}} not carried to the members. \\
\phantomsection\label{entry:symbol-registry-07357}\(\lVert\cdot\rVert_\infty\) & \texttt{\detokenize{\lVert;\rVert}} & uniform bounded-function norm & P17 integral extension & object & \hyperref[thm:p17-sa-general-converse]{\ref*{thm:p17-sa-general-converse}} & \thepage & Not monetary vector additivity. \\
\phantomsection\label{entry:symbol-registry-07358}\(Q\) & \texttt{\detokenize{Q}} & alternative probability and state density symbol family member \(Q\) & P17 identification boundary & object & \hyperref[def:p17-sa-failure]{\ref*{def:p17-sa-failure}} & \thepage & State-dependent scales are outside uniqueness. Explicit displayed family; member types and common actual first-use exposed, not one atomic glyph. Family listing, not one atomic glyph. Author macro \texttt{\detokenize{Q;r}} not carried to the members. \\
\phantomsection\label{entry:symbol-registry-07359}\(r\) & \texttt{\detokenize{r}} & alternative probability and state density symbol family member \(r\) & P17 identification boundary & object & \hyperref[def:p17-sa-failure]{\ref*{def:p17-sa-failure}} & \thepage & State-dependent scales are outside uniqueness. Explicit displayed family; member types and common actual first-use exposed, not one atomic glyph. Family listing, not one atomic glyph. Author macro \texttt{\detokenize{Q;r}} not carried to the members. \\
\end{longtable}
\endgroup
\Needspace{.4\textheight}
\subsection*{Stable part P18 scopes}
\begingroup\scriptsize\setlength{\tabcolsep}{3pt}
\begin{longtable}{@{}p{0.9in}p{1.3in}p{1.0in}p{0.5in}p{0.3in}p{0.4in}p{0.35in}p{1.65in}@{}}
\toprule
\textbf{Glyph} & \textbf{Macro} & \textbf{Role} & \textbf{Scope} & \textbf{Prec./form} & \textbf{First} & \textbf{Reg. page} & \textbf{Principal uses / notes} \\
\midrule
\endfirsthead
\toprule
\textbf{Glyph} & \textbf{Macro} & \textbf{Role} & \textbf{Scope} & \textbf{Prec./form} & \textbf{First} & \textbf{Reg. page} & \textbf{Principal uses / notes} \\
\midrule
\endhead
\midrule\multicolumn{8}{r}{\emph{continued}}\\\endfoot
\bottomrule\endlastfoot
\phantomsection\label{entry:symbol-registry-00949}\(\mathcal F=(A,R)\) & \texttt{\detokenize{\mathcal F,}\allowbreak\detokenize{A,}\allowbreak\detokenize{R}} & finite argumentation framework and directed attack relation & P18 finite-argumentation & meta & \hyperref[sec:p18-finite-argumentation]{\ref*{sec:p18-finite-argumentation}} & \thepage & Arguments are identifiers and \(R\subseteq A\times A\) is the supplied attack relation. The argument set may be empty and self-attacks are permitted; no truth interpretation is implied. \\
\phantomsection\label{entry:symbol-registry-00950}\(\operatorname{Att}_{\mathcal F}(a),S^+_{\mathcal F}\) & \texttt{\detokenize{\operatorname{Att}_{\mathcal F},}\allowbreak\detokenize{S^+_{\mathcal F}}} & attackers of an argument and targets attacked by a set & P18 finite-argumentation & meta & \hyperref[sec:p18-finite-argumentation]{\ref*{sec:p18-finite-argumentation}} & \thepage & Attackers point into \(a\); the attacked set collects targets of arrows out of \(S\). Subscripts may be omitted only within a fixed framework. \\
\phantomsection\label{entry:symbol-registry-00951}\(\mathsf F_{\mathcal F}:2^A\to2^A\) & \texttt{\detokenize{\mathsf F_{\mathcal F}}} & argumentation characteristic operator & P18 finite-argumentation & meta & \hyperref[sec:p18-finite-argumentation]{\ref*{sec:p18-finite-argumentation}} & \thepage & Returns exactly the arguments defended by its input set. It is monotone; being a fixed point alone does not imply conflict-freeness. This is not a temporal future operator. \\
\phantomsection\label{entry:symbol-registry-00952}\(G_{\mathcal F},G_i\) & \texttt{\detokenize{G_{\mathcal F},}\allowbreak\detokenize{G_i}} & grounded extension and its finite iteration stages & P18 finite-argumentation & meta & \hyperref[sec:p18-finite-argumentation]{\ref*{sec:p18-finite-argumentation}} & \thepage & Start at the empty set and iterate the characteristic operator. Equality occurs by index \(|A|\), so at most \(|A|+1\) evaluations include its detection; this is not a constant-time evaluation claim. \\
\phantomsection\label{entry:symbol-registry-00953}\(\operatorname{Comp}(\mathcal F),\operatorname{Pref}(\mathcal F),\operatorname{Stab}(\mathcal F)\) & \texttt{\detokenize{\operatorname{Comp},}\allowbreak\detokenize{\operatorname{Pref},}\allowbreak\detokenize{\operatorname{Stab}}} & complete preferred and stable extension families & P18 finite-argumentation & meta & \hyperref[sec:p18-finite-argumentation]{\ref*{sec:p18-finite-argumentation}} & \thepage & These are families of argument sets. Preferred means inclusion-maximal admissible, while stability additionally attacks every outsider. The stable family can be empty; Comp here is not the arithmetic computation formula. \\
\phantomsection\label{entry:symbol-registry-00954}\(\mathcal F_{\rm chain},\mathcal F_{\rm pair},\mathcal F_{\rm self},\mathcal F_{\rm empty},\mathcal F_{\rm size},\mathcal F_{\rm shared}\) & \texttt{\detokenize{\mathcal F_{\rm chain},}\allowbreak\detokenize{\mathcal F_{\rm pair},}\allowbreak\detokenize{\mathcal F_{\rm self},}\allowbreak\detokenize{\mathcal F_{\rm empty},}\allowbreak\detokenize{\mathcal F_{\rm size},}\allowbreak\detokenize{\mathcal F_{\rm shared}}} & six explicit finite argumentation examples & P18 finite-argumentation & meta & \hyperref[sec:p18-finite-argumentation]{\ref*{sec:p18-finite-argumentation}} & \thepage & Each constant denotes the complete argument/attack pair listed in def:p18-example-frameworks. Their full extension lists and boundary properties are proved in the matching annex item. \\
\phantomsection\label{entry:symbol-registry-00955}\(\mathcal E\subseteq2^A\) & \texttt{\detokenize{\mathcal E}} & chosen extension family for acceptance queries & P18 finite-argumentation & meta & \hyperref[sec:p18-finite-argumentation]{\ref*{sec:p18-finite-argumentation}} & \thepage & Credulous acceptance uses membership in some extension. The stated skeptical query also requires a nonempty family; no extensions is distinguished from a family containing only the empty extension. \\
\phantomsection\label{entry:symbol-registry-01445}\(\operatorname{Adm},S_\alpha,G_\alpha,\operatorname{Cl},\kappa^+\) & \texttt{\detokenize{\operatorname{Adm},}\allowbreak\detokenize{S_\alpha,}\allowbreak\detokenize{G_\alpha,}\allowbreak\detokenize{\operatorname{Cl},}\allowbreak\detokenize{\kappa^+}} & admissible-set family transfinite seed and grounded stages complete closure and successor cardinal bound & P18 arbitrary-argumentation & meta & \hyperref[sec:p18-arbitrary-argumentation]{\ref*{sec:p18-arbitrary-argumentation}} & \thepage & Closure is defined on admissible seeds. The arbitrary-set proofs state their uses of choice and transfinite recursion; finite termination is a distinct specialization. \\
\phantomsection\label{entry:symbol-registry-01446}\(\mathsf N,L_n,I,B_E,\mathcal C,T_n,D_n,E_n\) & \texttt{\detokenize{\mathsf N,}\allowbreak\detokenize{L_n,}\allowbreak\detokenize{I,}\allowbreak\detokenize{B_E,}\allowbreak\detokenize{\mathcal C,}\allowbreak\detokenize{T_n,}\allowbreak\detokenize{D_n,}\allowbreak\detokenize{E_n}} & no-attack operator alternating levels indefeasible set residual argument set controversial relation threats defenders and seed stages & P18 arbitrary-argumentation & meta & \hyperref[sec:p18-arbitrary-argumentation]{\ref*{sec:p18-arbitrary-argumentation}} & \thepage & The correction uses all attackers of current threats. Nonempty-extension claims have a nonempty argument-set hypothesis, and alternating finite levels need not reach the transfinite grounded extension. \\
\phantomsection\label{entry:symbol-registry-01447}\((a,\epsilon),\mathcal F_\omega,\mathcal F_{\omega+1},\mathcal F_{\rm reverse},\mathcal F_{\rm shareddefender},\mathcal F_{\rm parity},\mathcal F_{\rm initial},\mathcal F_{\rm nomindefense}\) & \texttt{\detokenize{(a,}\allowbreak\detokenize{\epsilon),}\allowbreak\detokenize{\mathcal F_\omega,}\allowbreak\detokenize{\mathcal F_{\omega+1},}\allowbreak\detokenize{\mathcal F_{\rm reverse},}\allowbreak\detokenize{\mathcal F_{\rm shareddefender},}\allowbreak\detokenize{\mathcal F_{\rm parity},}\allowbreak\detokenize{\mathcal F_{\rm initial},}\allowbreak\detokenize{\mathcal F_{\rm nomindefense}}} & parity graph vertices and the seven specified example frameworks & P18 arbitrary-argumentation & meta & \hyperref[sec:p18-arbitrary-argumentation]{\ref*{sec:p18-arbitrary-argumentation}} & \thepage & Parity counts walks with repeated vertices and zero-length even walks. The minimal-defense counterexample retains every infinite-framework hypothesis used by the corrected construction. \\
\phantomsection\label{entry:symbol-registry-01468}\(H,Q=G_P,\operatorname{Lit}(H),l^*,I^+,I^-,\widehat M,T_D,\operatorname{LM},\Gamma\) & \texttt{\detokenize{H,}\allowbreak\detokenize{Q,}\allowbreak\detokenize{G_P,}\allowbreak\detokenize{\operatorname{Lit},}\allowbreak\detokenize{l^*,}\allowbreak\detokenize{I^+,}\allowbreak\detokenize{I^-,}\allowbreak\detokenize{\widehat M,}\allowbreak\detokenize{T_D,}\allowbreak\detokenize{\operatorname{LM},}\allowbreak\detokenize{\Gamma}} & Herbrand base ground program signed literals contrary truth and falsity components total interpretation positive consequence least model and reduct operator & P18 logic-program-argumentation & meta & \hyperref[sec:p18-logic-program-argumentation]{\ref*{sec:p18-logic-program-argumentation}} & \thepage & Ground clauses have finite bodies even when the ground program is infinite. Truth-component superscripts are distinct from attack targets. \\
\phantomsection\label{entry:symbol-registry-01469}\(\mathcal A_\infty,\mathcal A_{\rm c},\mathcal A_{\rm u},\nu_p,\operatorname{con},F_\infty,F_{\rm c},N_\infty,{\cal E}_\infty,{\cal E}_{\rm c}\) & \texttt{\detokenize{\mathcal A_\infty,}\allowbreak\detokenize{\mathcal A_{\rm c},}\allowbreak\detokenize{\mathcal A_{\rm u},}\allowbreak\detokenize{\nu_p,}\allowbreak\detokenize{\operatorname{con},}\allowbreak\detokenize{F_\infty,}\allowbreak\detokenize{F_{\rm c},}\allowbreak\detokenize{N_\infty,}\allowbreak\detokenize{{\cal E}_\infty,}\allowbreak\detokenize{{\cal E}_{\rm c}}} & literal support clause and used-support graphs negative units conclusions characteristic operators and extension constructors & P18 logic-program-argumentation & meta & \hyperref[sec:p18-logic-program-argumentation]{\ref*{sec:p18-logic-program-argumentation}} & \thepage & Unused assumptions remain in the literal source graph. Graph conditions transfer only after the proved consistency-guarded reduction. \\
\phantomsection\label{entry:symbol-registry-01470}\(U_Q,V_Q,W_Q,T_Q^*,Z_Q,\Phi_Q,\Psi_Q\) & \texttt{\detokenize{U_Q,}\allowbreak\detokenize{V_Q,}\allowbreak\detokenize{W_Q,}\allowbreak\detokenize{T_Q^*,}\allowbreak\detokenize{Z_Q,}\allowbreak\detokenize{\Phi_Q,}\allowbreak\detokenize{\Psi_Q}} & unfounded and well-founded operators positive saturation completion and original and accelerated Fitting operators & P18 logic-program-argumentation & meta & \hyperref[sec:p18-logic-program-argumentation]{\ref*{sec:p18-logic-program-argumentation}} & \thepage & The original and accelerated operators have the same fixed points; their successor stages need not coincide. Completion uses equality of all three truth values. \\
\phantomsection\label{entry:symbol-registry-01471}\(\mathcal L_0,\operatorname{Th},\bar\varphi,\mathcal T=(D,W),J,\operatorname{Jus},C_n(B),C(B),B(E),\Gamma\) & \texttt{\detokenize{\mathcal L_0,}\allowbreak\detokenize{\operatorname{Th},}\allowbreak\detokenize{\bar\varphi,}\allowbreak\detokenize{\mathcal T,}\allowbreak\detokenize{D,}\allowbreak\detokenize{W,}\allowbreak\detokenize{J,}\allowbreak\detokenize{\operatorname{Jus},}\allowbreak\detokenize{C_n,}\allowbreak\detokenize{C,}\allowbreak\detokenize{B(E),}\allowbreak\detokenize{\Gamma}} & closed objective sentences classical closure normalized contrary default theory justifications licensed stages and Reiter operator & P18 default-autoepistemic-argumentation & meta & \hyperref[sec:p18-default-autoepistemic-argumentation]{\ref*{sec:p18-default-autoepistemic-argumentation}} & \thepage & Justifications are tested individually. Empty justification lists use the unconditional closure for the consistency condition. \\
\phantomsection\label{entry:symbol-registry-01472}\(K,H(S),\operatorname{arg},\operatorname{flat},\operatorname{St}(E),\mathscr B(T),\tau,\Delta\) & \texttt{\detokenize{K,}\allowbreak\detokenize{H(S),}\allowbreak\detokenize{\operatorname{arg},}\allowbreak\detokenize{\operatorname{flat},}\allowbreak\detokenize{\operatorname{St},}\allowbreak\detokenize{\mathscr B,}\allowbreak\detokenize{\tau,}\allowbreak\detokenize{\Delta}} & finite supports support union argument and flattening maps modal-kernel completion introspective assignment direct transform and weak-extension operator & P18 default-autoepistemic-argumentation & meta & \hyperref[sec:p18-default-autoepistemic-argumentation]{\ref*{sec:p18-default-autoepistemic-argumentation}} & \thepage & The full default/argument correspondence allows infinite closed theories. The direct modal transform has its proved weak-extension meaning; its Reiter bijection requires the stated prerequisite-free scope. \\
\phantomsection\label{entry:symbol-registry-01493}\(N,v,x(K),I(v),\mathcal C,\mathcal A_v,S_\triangle,S_c\) & \texttt{\detokenize{N,}\allowbreak\detokenize{v,}\allowbreak\detokenize{x(K),}\allowbreak\detokenize{I(v),}\allowbreak\detokenize{\mathcal C,}\allowbreak\detokenize{\mathcal A_v,}\allowbreak\detokenize{S_\triangle,}\allowbreak\detokenize{S_c}} & players characteristic function coalition payoff imputations core domination graph and three-player solution families & P18 game-argumentation & meta & \hyperref[sec:p18-game-argumentation]{\ref*{sec:p18-game-argumentation}} & \thepage & Coalitional feasibility uses a weak payoff inequality. Coalitional game solutions here are distinct from Nash action profiles. \\
\phantomsection\label{entry:symbol-registry-01494}\(\succ_i,\mu,\mu^{-1},U\) & \texttt{\detokenize{\succ_i,}\allowbreak\detokenize{\mu,}\allowbreak\detokenize{\mu^{-1},}\allowbreak\detokenize{U}} & person-indexed strict preference matching partner inverse and comparison utility & P18 game-argumentation & meta & \hyperref[sec:p18-game-argumentation]{\ref*{sec:p18-game-argumentation}} & \thepage & Balanced two-sided matching uses separate roles. One-population physical matching uses unordered pairs of distinct people. \\
\phantomsection\label{entry:symbol-registry-01495}\(\widehat v,\ell,O,Q_j,A_\ell,B_\ell,c^{(2j)}\) & \texttt{\detokenize{\widehat v,}\allowbreak\detokenize{\ell,}\allowbreak\detokenize{O,}\allowbreak\detokenize{Q_j,}\allowbreak\detokenize{A_\ell,}\allowbreak\detokenize{B_\ell,}\allowbreak\detokenize{c^{(2j)}}} & superadditive partition cover Lucas characteristic function odd-player coalition source coalitions regions and core vertices & P18 game-argumentation & meta & \hyperref[sec:p18-game-argumentation]{\ref*{sec:p18-game-argumentation}} & \thepage & The cover preserves the imputation domination graph when the stated singleton and grand-coalition values are preserved. \\
\phantomsection\label{entry:symbol-registry-01496}\(A,B,C,H,\rho,(i,j,k),r,M_i,(p,q)\) & \texttt{\detokenize{A,}\allowbreak\detokenize{B,}\allowbreak\detokenize{C,}\allowbreak\detokenize{H,}\allowbreak\detokenize{\rho,}\allowbreak\detokenize{(i,}\allowbreak\detokenize{j,}\allowbreak\detokenize{k),}\allowbreak\detokenize{r,}\allowbreak\detokenize{M_i,}\allowbreak\detokenize{(p,}\allowbreak\detokenize{q)}} & Lucas source regions core odd-index set cyclic permutation cyclic triple successor index and source coalitions & P18 lucas-no-solution & meta & \hyperref[sec:p18-lucas-no-solution]{\ref*{sec:p18-lucas-no-solution}} & \thepage & The local set H is the odd-index triple. The exceptional pair coordinates are indexed 7 and 9. \\
\phantomsection\label{entry:symbol-registry-01497}\(E_i,E,F_{jk},F_p,F_{79},F_{135},F,R,e_i(t,u,v),G_h(z),\operatorname{Dom}_E,\operatorname{Pred}_E\) & \texttt{\detokenize{E_i,}\allowbreak\detokenize{E,}\allowbreak\detokenize{F_{jk},}\allowbreak\detokenize{F_p,}\allowbreak\detokenize{F_{79},}\allowbreak\detokenize{F_{135},}\allowbreak\detokenize{F,}\allowbreak\detokenize{R,}\allowbreak\detokenize{e_i(t,}\allowbreak\detokenize{u,}\allowbreak\detokenize{v),}\allowbreak\detokenize{G_h(z),}\allowbreak\detokenize{\operatorname{Dom}_E,}\allowbreak\detokenize{\operatorname{Pred}_E}} & exception regions excluded-region families remainder parametrization upper regions restricted domination and predecessor & P18 lucas-no-solution & meta & \hyperref[sec:p18-lucas-no-solution]{\ref*{sec:p18-lucas-no-solution}} & \thepage & The F subscripts are unordered where specified. The e\_i parameter cube still obeys the strict local conditions defining E\_i. \\
\phantomsection\label{entry:symbol-registry-01498}\(D,P_F^0,P_D,Q_F,H_F,H_{P_D},a^\circ,a^\bullet,B_D^+,B_D\) & \texttt{\detokenize{D,}\allowbreak\detokenize{P_F^0,}\allowbreak\detokenize{P_D,}\allowbreak\detokenize{Q_F,}\allowbreak\detokenize{H_F,}\allowbreak\detokenize{H_{P_D},}\allowbreak\detokenize{a^\circ,}\allowbreak\detokenize{a^\bullet,}\allowbreak\detokenize{B_D^+,}\allowbreak\detokenize{B_D}} & name domain literal source program guarded program sparse program ground bases acceptance defeat and positive and signed background & P18 argumentation-meta-interpreters & meta & \hyperref[sec:p18-argumentation-meta-interpreters]{\ref*{sec:p18-argumentation-meta-interpreters}} & \thepage & The literal source program requires its stated nonempty argument domain. The guarded program permits an empty argument set in a nonempty name domain. \\
\phantomsection\label{entry:symbol-registry-01499}\(m_Q,m_D,m_0,L_Q,L_D,L_0,\iota_D,d_X,\tau_\alpha,(T,Z,D,U)\) & \texttt{\detokenize{m_Q,}\allowbreak\detokenize{m_D,}\allowbreak\detokenize{m_0,}\allowbreak\detokenize{L_Q,}\allowbreak\detokenize{L_D,}\allowbreak\detokenize{L_0,}\allowbreak\detokenize{\iota_D,}\allowbreak\detokenize{d_X,}\allowbreak\detokenize{\tau_\alpha,}\allowbreak\detokenize{(T,}\allowbreak\detokenize{Z,}\allowbreak\detokenize{D,}\allowbreak\detokenize{U)}} & interpretation constructors full signed outputs embedding defeat map transfinite stages and local truth-falsity quadruple & P18 argumentation-meta-interpreters & meta & \hyperref[sec:p18-argumentation-meta-interpreters]{\ref*{sec:p18-argumentation-meta-interpreters}} & \thepage & The quadruple's D is the true defeat component. Negative background includes nonedges and extra names for the full guarded output. \\
\phantomsection\label{entry:symbol-registry-01500}\(n,m,\Phi,\Gamma\) & \texttt{\detokenize{n,}\allowbreak\detokenize{m,}\allowbreak\detokenize{\Phi,}\allowbreak\detokenize{\Gamma}} & finite argument and attack counts Fitting operator and reduct operator & P18 argumentation-meta-interpreters & meta & \hyperref[sec:p18-argumentation-meta-interpreters]{\ref*{sec:p18-argumentation-meta-interpreters}} & \thepage & Four original Fitting steps implement one defense step. Sparse complexity bounds do not count a full quadratic negative background as linear. \\
\phantomsection\label{entry:symbol-registry-01580}\(L,C,R_s,R_d,n,K_n,K_p,\mathcal A\) & \texttt{\detokenize{L,}\allowbreak\detokenize{C,}\allowbreak\detokenize{R_s,}\allowbreak\detokenize{R_d,}\allowbreak\detokenize{n,}\allowbreak\detokenize{K_n,}\allowbreak\detokenize{K_p,}\allowbreak\detokenize{\mathcal A}} & formula language contrary function strict and defeasible rules partial naming map axiom and ordinary premises and arguments & P18 structured-argumentation & meta & \hyperref[sec:p18-structured-argumentation]{\ref*{sec:p18-structured-argumentation}} & \thepage & Selected two-premise-sort tutorial model admits all finite argument trees and general contrary relations. \\
\phantomsection\label{entry:symbol-registry-01581}\(\operatorname{ant},\operatorname{head},\operatorname{Conc},\operatorname{Prem},\operatorname{Prem}_p,\operatorname{Sub}\) & \texttt{\detokenize{\operatorname{ant},}\allowbreak\detokenize{\operatorname{head},}\allowbreak\detokenize{\operatorname{Conc},}\allowbreak\detokenize{\operatorname{Prem},}\allowbreak\detokenize{\operatorname{Prem}_p,}\allowbreak\detokenize{\operatorname{Sub}}} & rule antecedents and consequent argument conclusion premises ordinary premises and subarguments & P18 structured-argumentation & meta & \hyperref[sec:p18-structured-argumentation]{\ref*{sec:p18-structured-argumentation}} & \thepage & Rule arrows are typed object data; premise preference applies only to ordinary premises. \\
\phantomsection\label{entry:symbol-registry-01582}\(\operatorname{Top},\operatorname{Def},\operatorname{Last}\) & \texttt{\detokenize{\operatorname{Top},}\allowbreak\detokenize{\operatorname{Def},}\allowbreak\detokenize{\operatorname{Last}}} & top rule defeasible rule set and last defeasible rule set & P18 structured-argumentation & meta & \hyperref[sec:p18-structured-argumentation]{\ref*{sec:p18-structured-argumentation}} & \thepage & Top is undefined for premise arguments; nullary rule arguments are distinct from premises. \\
\phantomsection\label{entry:symbol-registry-01583}\(\leq_p,\leq_d,<_{p},<_{d}\) & \texttt{\detokenize{\leq_p,}\allowbreak\detokenize{\leq_d,}\allowbreak\detokenize{<_{p},}\allowbreak\detokenize{<_{d}}} & ordinary-premise and rule priorities and their asymmetric strict parts & P18 structured-argumentation & meta & \hyperref[sec:p18-structured-argumentation]{\ref*{sec:p18-structured-argumentation}} & \thepage & Partial preorders need not permit replacing every weak comparison inside a set lift by its strict part. \\
\phantomsection\label{entry:symbol-registry-01584}\(\unlhd_{\rm E},\unlhd_{\rm D},\lhd,\preceq_{\rm L},\preceq_{\rm W},\prec_{\rm L},\prec_{\rm W}\) & \texttt{\detokenize{\unlhd_{\rm E},}\allowbreak\detokenize{\unlhd_{\rm D},}\allowbreak\detokenize{\lhd,}\allowbreak\detokenize{\preceq_{\rm L},}\allowbreak\detokenize{\preceq_{\rm W},}\allowbreak\detokenize{\prec_{\rm L},}\allowbreak\detokenize{\prec_{\rm W}}} & elitist and democratic set comparisons their strict part last-link and weakest-link argument comparisons and strict parts & P18 structured-argumentation & meta & \hyperref[sec:p18-structured-argumentation]{\ref*{sec:p18-structured-argumentation}} & \thepage & Strict comparison is weak-and-not-reverse; empty-set cases are applied before the quantified comparisons. \\
\phantomsection\label{entry:symbol-registry-01585}\(D,F_{\rm AS},\operatorname{Conc}(E),\operatorname{Cl}_s\) & \texttt{\detokenize{D,}\allowbreak\detokenize{F_{\rm AS},}\allowbreak\detokenize{\operatorname{Conc}(E),}\allowbreak\detokenize{\operatorname{Cl}_s}} & defeat relation induced abstract framework extension conclusions and strict closure & P18 structured-argumentation & meta & \hyperref[sec:p18-structured-argumentation]{\ref*{sec:p18-structured-argumentation}} & \thepage & Defeat preferences are evaluated at the attacked subargument; formula query quantifiers are explicit. \\
\phantomsection\label{entry:symbol-registry-01586}\(A_i,B_i,C_i,D_i,E_0\) & \texttt{\detokenize{A_i,}\allowbreak\detokenize{B_i,}\allowbreak\detokenize{C_i,}\allowbreak\detokenize{D_i,}\allowbreak\detokenize{E_0}} & twelve running-example arguments and the unranked extension & P18 structured-argumentation & meta & \hyperref[sec:p18-structured-argumentation]{\ref*{sec:p18-structured-argumentation}} & \thepage & The complete generated graph includes B4 and its undermining and axiom-rebuttal attackers. \\
\phantomsection\label{entry:symbol-registry-01607}\(T,D,\operatorname{Str}(A),\operatorname{Prem}_n(A),\vdash_s,M(A)\) & \texttt{\detokenize{T,}\allowbreak\detokenize{D,}\allowbreak\detokenize{\operatorname{Str}(A),}\allowbreak\detokenize{\operatorname{Prem}_n(A),}\allowbreak\detokenize{\vdash_s,}\allowbreak\detokenize{M(A)}} & structural attack defeat strict rules axiom premises strict consequence and maximal fallible frontier & P18 structured-rationality & meta & \hyperref[sec:p18-structured-rationality]{\ref*{sec:p18-structured-rationality}} & \thepage & Unrestricted and c-restricted argument universes are separate; profile continuation does not imply tree grafting. \\
\phantomsection\label{entry:symbol-registry-01608}\(\triangleleft_{\rm E},\triangleleft_{\rm D},\prec_{\rm L}^{\,c},\prec_{\rm W}^{\,c}\) & \texttt{\detokenize{\triangleleft_{\rm E},}\allowbreak\detokenize{\triangleleft_{\rm D},}\allowbreak\detokenize{\prec_{\rm L}^{\,c},}\allowbreak\detokenize{\prec_{\rm W}^{\,c}}} & corrected direct strict set comparisons and corrected last-link and weakest-link preferences & P18 structured-rationality & meta & \hyperref[sec:p18-structured-rationality]{\ref*{sec:p18-structured-rationality}} & \thepage & These use base strict comparisons directly and differ from the tutorial's asymmetric parts of weak comparisons. \\
\phantomsection\label{entry:symbol-registry-01609}\((R),C_i,C_{-i},W\) & \texttt{\detokenize{(R),}\allowbreak\detokenize{C_i,}\allowbreak\detokenize{C_{-i},}\allowbreak\detokenize{W}} & local returning-defeater condition profile family omitted-index continuation and returning argument & P18 structured-rationality & meta & \hyperref[sec:p18-structured-rationality]{\ref*{sec:p18-structured-rationality}} & \thepage & Finite sets and indexed families allowing repetitions impose different conditions; deriving R from priorities is a separate proof obligation. \\
\phantomsection\label{entry:symbol-registry-01639}\(<_p,<_d,\triangleleft_{\rm E},\triangleleft_{\rm D},s\) & \texttt{\detokenize{<_p,}\allowbreak\detokenize{<_d,}\allowbreak\detokenize{\triangleleft_{\rm E},}\allowbreak\detokenize{\triangleleft_{\rm D},}\allowbreak\detokenize{s}} & typed resource priorities direct strict set lifts and the selected elitist or democratic convention & P18 structured-priority-rationality & meta & \hyperref[sec:p18-structured-priority-rationality]{\ref*{sec:p18-structured-priority-rationality}} & \thepage & Finite resource sets and empty operands follow the declared typed comparisons. \\
\phantomsection\label{entry:symbol-registry-01640}\(\prec_{{\rm L},s},\prec_{{\rm W},s}\) & \texttt{\detokenize{\prec_{{\rm L},s},}\allowbreak\detokenize{\prec_{{\rm W},s}}} & four corrected last-link and weakest-link strict argument priorities & P18 structured-priority-rationality & meta & \hyperref[sec:p18-structured-priority-rationality]{\ref*{sec:p18-structured-priority-rationality}} & \thepage & Actual strict tree continuations are distinguished from the broader source profile relation. \\
\phantomsection\label{entry:symbol-registry-01641}\((R),\operatorname{Prem}(X)\cup\operatorname{Prem}(Y)\) & \texttt{\detokenize{(R),}\allowbreak\detokenize{\operatorname{Prem}(X)\cup\operatorname{Prem}(Y)}} & global return for unrestricted trees and the combined support governing restricted return & P18 structured-priority-rationality & meta & \hyperref[sec:p18-structured-priority-rationality]{\ref*{sec:p18-structured-priority-rationality}} & \thepage & Restricted return is proved when combined support is c-consistent; joint defense supplies that condition by a separate lemma. \\
\phantomsection\label{entry:symbol-registry-01667}\(E,S,A\) & \texttt{\detokenize{E,}\allowbreak\detokenize{S,}\allowbreak\detokenize{A}} & expert domain and asserted proposition & P18 argument-schemes-critical-questions & meta & \hyperref[sec:p18-scheme-expert]{\ref*{sec:p18-scheme-expert}} & \thepage & Roles are local to expert-opinion schemes; a source statement does not make its proposition true by definition. \\
\phantomsection\label{entry:symbol-registry-01668}\(W,A\) & \texttt{\detokenize{W,}\allowbreak\detokenize{A}} & witness and proposition testified to & P18 argument-schemes-critical-questions & meta & \hyperref[sec:p18-scheme-witness]{\ref*{sec:p18-scheme-witness}} & \thepage & Witness access truth-telling and statement are separate premise roles. \\
\phantomsection\label{entry:symbol-registry-01669}\(C_1,C_2,C_3\) & \texttt{\detokenize{C_1,}\allowbreak\detokenize{C_2,}\allowbreak\detokenize{C_3}} & analogy source case target case and countercase & P18 argument-schemes-critical-questions & meta & \hyperref[sec:p18-scheme-analogy]{\ref*{sec:p18-scheme-analogy}} & \thepage & The selected compact and revised schemes and their questions retain their separate source versions. \\
\phantomsection\label{entry:symbol-registry-01670}\(E_1,E_2\) & \texttt{\detokenize{E_1,}\allowbreak\detokenize{E_2}} & causal event variables & \hyperref[sec:p18-scheme-causal]{\ref*{sec:p18-scheme-causal}} & meta & \hyperref[sec:p18-scheme-causal]{\ref*{sec:p18-scheme-causal}} & \thepage & A semantic causal relation differs from the direction of an argument using it. \\
\phantomsection\label{entry:symbol-registry-01671}\(A,B_i\) & \texttt{\detokenize{A,}\allowbreak\detokenize{B_i}} & practical goal and candidate alternative actions & P18 argument-schemes-critical-questions & meta & \hyperref[sec:p18-scheme-practical]{\ref*{sec:p18-scheme-practical}} & \thepage & This local A is a goal rather than the expert-entry proposition role; candidate means and necessary means differ. \\
\phantomsection\label{entry:symbol-registry-01689}\((\tau,t),(\mathrm{prefer},a,b)\) & \texttt{\detokenize{(\tau,}\allowbreak\detokenize{t),}\allowbreak\detokenize{(\mathrm{prefer},}\allowbreak\detokenize{a,}\allowbreak\detokenize{b)}} & typed public commitment and speaker-owned action-preference entry & P18 dialogue-burden & meta & \hyperref[def:p18-dialogue-record]{\ref*{def:p18-dialogue-record}} & \thepage & The source's six content types are distinct; a factual type does not certify truth and a public commitment need not be a private belief. \\
\phantomsection\label{entry:symbol-registry-01690}\(M,D,g,c\) & \texttt{\detokenize{M,}\allowbreak\detokenize{D,}\allowbreak\detokenize{g,}\allowbreak\detokenize{c}} & mirror and desk service alternatives goal and cost constraint & P18 dialogue-burden & meta & \hyperref[def:p18-dialogue-library-exchange]{\ref*{def:p18-dialogue-library-exchange}} & \thepage & All costs documents and committee rules in the worked exchange are stipulated examples. \\
\phantomsection\label{entry:symbol-registry-01691}\(q,f_1,f_2,d,h,m\) & \texttt{\detokenize{q,}\allowbreak\detokenize{f_1,}\allowbreak\detokenize{f_2,}\allowbreak\detokenize{d,}\allowbreak\detokenize{h,}\allowbreak\detokenize{m}} & cost evaluation quoted cost added fee desk evaluation staff record and corrected mirror total & P18 dialogue-burden & meta & \hyperref[def:p18-dialogue-library-exchange]{\ref*{def:p18-dialogue-library-exchange}} & \thepage & Fact records and evaluations have separately declared types; the corrected total is 800+300=1100 credits. \\
\phantomsection\label{entry:symbol-registry-01692}\(+x,-x,\pi\) & \texttt{\detokenize{+x,}\allowbreak\detokenize{-x,}\allowbreak\detokenize{\pi}} & local store addition removal and preference for desk over mirror & P18 dialogue-burden & meta & \hyperref[sec:p18-dialogue-worked]{\ref*{sec:p18-dialogue-worked}} & \thepage & Only a speaker's own moves update that speaker's store; the chronological record also retains moves with zero store change. \\
\phantomsection\label{entry:symbol-registry-01707}\(C,E\) & \texttt{\detokenize{C,}\allowbreak\detokenize{E}} & relay-energization and indicator-light propositions in the causal-form example & P18 scheme-classification-diagnosis & meta & \hyperref[def:p18-scheme-causal-forms]{\ref*{def:p18-scheme-causal-forms}} & \thepage & Forward deduction contraposition explanatory hypothesis denial of antecedent and comparison retain their different premises. \\
\phantomsection\label{entry:symbol-registry-01708}\(A,G,V,s\) & \texttt{\detokenize{A,}\allowbreak\detokenize{G,}\allowbreak\detokenize{V,}\allowbreak\detokenize{s}} & action goal value and agent roles in practical schemes & P18 scheme-classification-diagnosis & meta & \hyperref[def:p18-scheme-consequence-value-forms]{\ref*{def:p18-scheme-consequence-value-forms}} & \thepage & The source's necessary-means sufficient-means and value-based forms remain separately identified. \\
\phantomsection\label{entry:symbol-registry-01709}\(A_0,\ldots,A_n\) & \texttt{\detokenize{A_0,}\allowbreak\detokenize{\ldots,}\allowbreak\detokenize{A_n}} & successive positions in a proposed slippery-slope argument & P18 scheme-classification-diagnosis & meta & \hyperref[def:p18-scheme-slope-diagnosis]{\ref*{def:p18-scheme-slope-diagnosis}} & \thepage & The stopping conditions and causal precedent or linguistic mechanism must be established rather than inferred from a slope label. \\
\phantomsection\label{entry:symbol-registry-01710}\(P_1,\ldots,P_{11}\) & \texttt{\detokenize{P_1,}\allowbreak\detokenize{\ldots,}\allowbreak\detokenize{P_{11}}} & propositions in the complete connected-argument example & P18 scheme-classification-diagnosis & meta & \hyperref[def:p18-scheme-net-diagnosis]{\ref*{def:p18-scheme-net-diagnosis}} & \thepage & The net contains five component inferences with separately recorded explicit contextual and needed premises. \\
\phantomsection\label{entry:symbol-registry-01711}\(X,Y,Z\) & \texttt{\detokenize{X,}\allowbreak\detokenize{Y,}\allowbreak\detokenize{Z}} & successive nodes along an inferential dependency path & P18 scheme-classification-diagnosis & meta & \hyperref[sec:p18-scheme-computational-boundary]{\ref*{sec:p18-scheme-computational-boundary}} & \thepage & A dependency path is not automatically transitive defeasible entailment. \\
\phantomsection\label{entry:symbol-registry-01801}\(\mathcal L,\overline l,r,\to,\Rightarrow\) & \texttt{\detokenize{\mathcal L,}\allowbreak\detokenize{\overline l,}\allowbreak\detokenize{r,}\allowbreak\detokenize{\to,}\allowbreak\detokenize{\Rightarrow}} & LIS language literal complement named strict rule and presumptive rule markers & P18 lis-burdens & meta & \hyperref[def:p18-lis-argument-interface]{\ref*{def:p18-lis-argument-interface}} & \thepage & Complementation cancels adjacent double negations; admitted rules and premises are inputs to this tree-argument interface. \\
\phantomsection\label{entry:symbol-registry-01802}\(\mathcal A,\operatorname{Conc}(A),\operatorname{Sub}(A)\) & \texttt{\detokenize{\mathcal A,}\allowbreak\detokenize{\operatorname{Conc}(A),}\allowbreak\detokenize{\operatorname{Sub}(A)}} & tree-argument universe root conclusion and subtree occurrences & P18 lis-burdens & meta & \hyperref[def:p18-lis-argument-interface]{\ref*{def:p18-lis-argument-interface}} & \thepage & Subarguments are addressed occurrences with unique parents; this is not the later IS rule-content convention. \\
\phantomsection\label{entry:symbol-registry-01803}\(\neg\operatorname{app}(r,\iota),\iota\) & \texttt{\detokenize{\neg\operatorname{app}(r,}\allowbreak\detokenize{\iota),}\allowbreak\detokenize{\iota}} & named application undercutter and application occurrence identifier & P18 lis-burdens & meta & \hyperref[def:p18-lis-argument-interface]{\ref*{def:p18-lis-argument-interface}} & \thepage & The formula denies applicability of the named presumptive application; it is not a classical negation of a strict implication. \\
\phantomsection\label{entry:symbol-registry-01804}\(D,D_s,t(C,A)\) & \texttt{\detokenize{D,}\allowbreak\detokenize{D_s,}\allowbreak\detokenize{t(C,}\allowbreak\detokenize{A)}} & adjudicated defeat strict defeat and unique smallest attacked occurrence & P18 lis-burdens & meta & \hyperref[def:p18-lis-argument-interface]{\ref*{def:p18-lis-argument-interface}} & \thepage & Strict defeat means no reverse defeat; the changing-role interface requires a unique minimal target and consistent target inheritance. \\
\phantomsection\label{entry:symbol-registry-01805}\(\mathord{\sim}l,\operatorname{Exc}(r)\) & \texttt{\detokenize{\mathord{\sim}l,}\allowbreak\detokenize{\operatorname{Exc}(r)}} & weak exception premise and named exception proposition in the compared representations & P18 lis-burdens & meta & \hyperref[sec:p18-lis-interface]{\ref*{sec:p18-lis-interface}} & \thepage & Weak assumptions ordinary negative premises and named applicability attacks have different response effects. \\
\phantomsection\label{entry:symbol-registry-01806}\(P,O\) & \texttt{\detokenize{P,}\allowbreak\detokenize{O}} & fixed-role proponent and opponent labels & P18 lis-burdens & meta & \hyperref[def:p18-lis-base-game]{\ref*{def:p18-lis-base-game}} & \thepage & P may not repeat argument content; O may repeat. These role labels differ from physical parties. \\
\phantomsection\label{entry:symbol-registry-01807}\(F,S,J,J_\alpha\) & \texttt{\detokenize{F,}\allowbreak\detokenize{S,}\allowbreak\detokenize{J,}\allowbreak\detokenize{J_\alpha}} & ordinary-defense operator support set grounded extension and ordinal stages & P18 lis-burdens & meta & \hyperref[sec:p18-lis-base-game]{\ref*{sec:p18-lis-base-game}} & \thepage & F uses ordinary D-defense on the fixed graph; the separate proof establishes correspondence to the strict-response game. \\
\phantomsection\label{entry:symbol-registry-01808}\(Q,\operatorname{Cn}_Q,\models_{\rm CL}\) & \texttt{\detokenize{Q,}\allowbreak\detokenize{\operatorname{Cn}_Q,}\allowbreak\detokenize{\models_{\rm CL}}} & issue vocabulary issue-relative classical closure and classical entailment & P18 lis-burdens & meta & \hyperref[def:p18-lis-burden-profile]{\ref*{def:p18-lis-burden-profile}} & \thepage & Q is complement-closed and excludes tautologies; this explicit corrected domain differs from unrestricted classical closure. \\
\phantomsection\label{entry:symbol-registry-01809}\(\pi,\delta,\bar p\) & \texttt{\detokenize{\pi,}\allowbreak\detokenize{\delta,}\allowbreak\detokenize{\bar p}} & physical parties and the other-party operation & P18 lis-burdens & meta & \hyperref[def:p18-lis-burden-profile]{\ref*{def:p18-lis-burden-profile}} & \thepage & Party identity remains fixed while dialectical roles may change by move and target. \\
\phantomsection\label{entry:symbol-registry-01810}\(\operatorname{Hist},h,\epsilon,B(h,p)\) & \texttt{\detokenize{\operatorname{Hist},}\allowbreak\detokenize{h,}\allowbreak\detokenize{\epsilon,}\allowbreak\detokenize{B(h,}\allowbreak\detokenize{p)}} & finite histories empty history and history-dependent burden profile & P18 lis-burdens & meta & \hyperref[def:p18-lis-burden-profile]{\ref*{def:p18-lis-burden-profile}} & \thepage & Profile consistency separation and retention constraints are explicit; computability and finite state are not assumed. \\
\phantomsection\label{entry:symbol-registry-01811}\(M_i,p_i,A_i,h_{i-1},\rho_i(U)\) & \texttt{\detokenize{M_i,}\allowbreak\detokenize{p_i,}\allowbreak\detokenize{A_i,}\allowbreak\detokenize{h_{i-1},}\allowbreak\detokenize{\rho_i(U)}} & candidate move party argument preceding history and role of a subtree occurrence & P18 lis-burdens & meta & \hyperref[def:p18-lis-roles]{\ref*{def:p18-lis-roles}} & \thepage & Historical roles are retained as assigned at the move; an explicit burden precedes inherited role. \\
\phantomsection\label{entry:symbol-registry-01812}\(H,\operatorname{Succ}(h),\operatorname{turn}(h)\) & \texttt{\detokenize{H,}\allowbreak\detokenize{\operatorname{Succ}(h),}\allowbreak\detokenize{\operatorname{turn}(h)}} & reachable history arena legal successors and next physical party & P18 lis-burdens & meta & \hyperref[def:p18-lis-attractor]{\ref*{def:p18-lis-attractor}} & \thepage & Histories with identical last arguments can have different roles and future moves. \\
\phantomsection\label{entry:symbol-registry-01813}\(T_p,\Phi_p,W_p,X_\alpha\) & \texttt{\detokenize{T_p,}\allowbreak\detokenize{\Phi_p,}\allowbreak\detokenize{W_p,}\allowbreak\detokenize{X_\alpha}} & winning terminals reachability operator winning set and ordinal iteration & P18 lis-burdens & meta & \hyperref[def:p18-lis-attractor]{\ref*{def:p18-lis-attractor}} & \thepage & The universal-successor clause requires nonempty successors; histories outside both winning sets are not assigned a terminal winner. \\
\phantomsection\label{entry:symbol-registry-01814}\(A,B,C,b\) & \texttt{\detokenize{A,}\allowbreak\detokenize{B,}\allowbreak\detokenize{C,}\allowbreak\detokenize{b}} & three tied-evidence arguments and their local issue & P18 lis-burdens & meta & \hyperref[sec:p18-lis-tied-evidence]{\ref*{sec:p18-lis-tied-evidence}} & \thepage & B is an argument in this example; B(h,p) remains the separately typed burden function. \\
\phantomsection\label{entry:symbol-registry-01815}\(A,B,C,D,a,w,e,d,r_d,r_c,r_D\) & \texttt{\detokenize{A,}\allowbreak\detokenize{B,}\allowbreak\detokenize{C,}\allowbreak\detokenize{D,}\allowbreak\detokenize{a,}\allowbreak\detokenize{w,}\allowbreak\detokenize{e,}\allowbreak\detokenize{d,}\allowbreak\detokenize{r_d,}\allowbreak\detokenize{r_c,}\allowbreak\detokenize{r_D}} & document example arguments application award exception document and evidence contents & P18 lis-burdens & meta & \hyperref[sec:p18-lis-document-issues]{\ref*{sec:p18-lis-document-issues}} & \thepage & The capital D here names a document argument; relational D is separately declared defeat data. \\
\phantomsection\label{entry:symbol-registry-01816}\(A_o,A_x,N_o,N_x,N_v,p,o,v,x,c\) & \texttt{\detokenize{A_o,}\allowbreak\detokenize{A_x,}\allowbreak\detokenize{N_o,}\allowbreak\detokenize{N_x,}\allowbreak\detokenize{N_v,}\allowbreak\detokenize{p,}\allowbreak\detokenize{o,}\allowbreak\detokenize{v,}\allowbreak\detokenize{x,}\allowbreak\detokenize{c}} & target-specific subarguments counterarguments and possession entitlement witness damage and compensation issues & P18 lis-burdens & meta & \hyperref[sec:p18-lis-target-specific]{\ref*{sec:p18-lis-target-specific}} & \thepage & The local witness premise v also denotes its premise argument where the complete defeat matrix says so. \\
\phantomsection\label{entry:symbol-registry-01817}\(U,E\) & \texttt{\detokenize{U,}\allowbreak\detokenize{E}} & protected production ruling and its evidential target & P18 lis-burdens & meta & \hyperref[thm:p18-lis-protected-production]{\ref*{thm:p18-lis-protected-production}} & \thepage & The ruling is assumed undefeated and strictly defeating; the theorem does not set a universal evidential threshold. \\
\phantomsection\label{entry:symbol-registry-01818}\(M,E,R,U\) & \texttt{\detokenize{M,}\allowbreak\detokenize{E,}\allowbreak\detokenize{R,}\allowbreak\detokenize{U}} & main claim exception response and undercutter across the four-stage example & P18 lis-burdens & meta & \hyperref[def:p18-lis-stage-example]{\ref*{def:p18-lis-stage-example}} & \thepage & The production event is retained while the tactical forecast changes; the chain is the full stipulated defeat relation. \\
\phantomsection\label{entry:symbol-registry-01819}\(\Gamma,S,\mathcal D,\prec\) & \texttt{\detokenize{\Gamma,}\allowbreak\detokenize{S,}\allowbreak\detokenize{\mathcal D,}\allowbreak\detokenize{\prec}} & fixed-priority rule theory strict rules defeasible rules and base priority order & P18 is1996 & meta & \hyperref[def:p18-is1996-theory]{\ref*{def:p18-is1996-theory}} & \thepage & The base order is a strict partial order on defeasible rules; ordinary strict contraposition is not silently added. \\
\phantomsection\label{entry:symbol-registry-01820}\(L,\overline L,\mathord{\sim}L,r,L_i,M_j,N,\to,\Rightarrow\) & \texttt{\detokenize{L,}\allowbreak\detokenize{\overline L,}\allowbreak\detokenize{\mathord{\sim}L,}\allowbreak\detokenize{r,}\allowbreak\detokenize{L_i,}\allowbreak\detokenize{M_j,}\allowbreak\detokenize{N,}\allowbreak\detokenize{\to,}\allowbreak\detokenize{\Rightarrow}} & strong literals complement weak premise rule name rule premises consequent and rule-status markers & P18 is1996 & meta & \hyperref[def:p18-is1996-theory]{\ref*{def:p18-is1996-theory}} & \thepage & Weak absence of support and strong contrary support are distinct; a schema denotes its ground instances. \\
\phantomsection\label{entry:symbol-registry-01821}\(\mathcal A_\Gamma,[],\operatorname{Con}(A),\operatorname{Ass}(A)\) & \texttt{\detokenize{\mathcal A_\Gamma,}\allowbreak\detokenize{[],}\allowbreak\detokenize{\operatorname{Con}(A),}\allowbreak\detokenize{\operatorname{Ass}(A)}} & legal rule-content arguments empty argument all conclusions and weak-assumption complements & P18 is1996 & meta & \hyperref[def:p18-is1996-argument]{\ref*{def:p18-is1996-argument}} & \thepage & Each content has a legal premise-before-conclusion ordering and one producer per head; Con includes every consequent rather than one root. \\
\phantomsection\label{entry:symbol-registry-01822}\(A+T,C_S(A)\) & \texttt{\detokenize{A+T,}\allowbreak\detokenize{C_S(A)}} & legal strict extension and strict saturation of an argument & P18 is1996 & meta & \hyperref[def:p18-is1996-strict-closure]{\ref*{def:p18-is1996-strict-closure}} & \thepage & Extensions retain unique producing heads; infinite strict rule sets use finite derivations without a termination claim. \\
\phantomsection\label{entry:symbol-registry-01823}\(U(A,B)\) & \texttt{\detokenize{U(A,}\allowbreak\detokenize{B)}} & source IS undercut of a weak antecedent & P18 is1996 & meta & \hyperref[def:p18-is1996-attack]{\ref*{def:p18-is1996-attack}} & \thepage & This local U is a relation and differs from the protected-ruling argument of P18.14. \\
\phantomsection\label{entry:symbol-registry-01824}\(R_L(E)\) & \texttt{\detokenize{R_L(E)}} & relevant defeasible frontier for one produced literal in one argument & P18 is1996 & meta & \hyperref[def:p18-is1996-frontier]{\ref*{def:p18-is1996-frontier}} & \thepage & The recursion uses the whole strict extension; strict facts have empty frontiers and different extensions are not merged. \\
\phantomsection\label{entry:symbol-registry-01825}\(R,R',\prec_*\) & \texttt{\detokenize{R,}\allowbreak\detokenize{R',}\allowbreak\detokenize{\prec_*}} & finite rule frontiers and existential-universal strict set lift & P18 is1996 & meta & \hyperref[def:p18-is1996-set-priority]{\ref*{def:p18-is1996-set-priority}} & \thepage & One left rule must be below every right rule; nonempty frontiers lie below the empty right frontier and the empty left frontier lies below none. \\
\phantomsection\label{entry:symbol-registry-01826}\(R(A,B),D_\Gamma,D_\Gamma^s\) & \texttt{\detokenize{R(A,}\allowbreak\detokenize{B),}\allowbreak\detokenize{D_\Gamma,}\allowbreak\detokenize{D_\Gamma^s}} & strict-extension rebuttal generated defeat and strict defeat & P18 is1996 & meta & \hyperref[def:p18-is1996-defeat]{\ref*{def:p18-is1996-defeat}} & \thepage & The empty incoherence clause can be a defeat without an attack; R(A,B) is a relation rather than a frontier set. \\
\phantomsection\label{entry:symbol-registry-01827}\(F_\Gamma,J_\Gamma\) & \texttt{\detokenize{F_\Gamma,}\allowbreak\detokenize{J_\Gamma}} & strict-defense operator and its least fixed point for fixed priorities & P18 is1996 & meta & \hyperref[sec:p18-is1996-semantics]{\ref*{sec:p18-is1996-semantics}} & \thepage & The ordinary-defense equality needed by the fixed-role LIS bridge is proved separately; source overruled labels use attack. \\
\phantomsection\label{entry:symbol-registry-01828}\(P,O\) & \texttt{\detokenize{P,}\allowbreak\detokenize{O}} & players in the corrected finite fixed-priority proof game & P18 is1996 & meta & \hyperref[def:p18-is1996-proof-game]{\ref*{def:p18-is1996-proof-game}} & \thepage & P uses strict defeat without repeating its content; O may repeat; global reply minimality is removed. \\
\phantomsection\label{entry:symbol-registry-01829}\(r_1,r_2,r_3,r_4,C\) & \texttt{\detokenize{r_1,}\allowbreak\detokenize{r_2,}\allowbreak\detokenize{r_3,}\allowbreak\detokenize{r_4,}\allowbreak\detokenize{C}} & four rules and protective combined content in the minimality countermodel & P18 is1996 & meta & \hyperref[def:p18-is1996-minimality-case]{\ref*{def:p18-is1996-minimality-case}} & \thepage & The countermodel covers every legal rule-content argument and disproves the globally minimal reply restriction. \\
\phantomsection\label{entry:symbol-registry-01830}\(r_1,\ldots,r_6\) & \texttt{\detokenize{r_1,}\allowbreak\detokenize{\ldots,}\allowbreak\detokenize{r_6}} & fresh rule names in the earlier-support priority example & P18 is1996 & meta & \hyperref[def:p18-is1996-earlier-support]{\ref*{def:p18-is1996-earlier-support}} & \thepage & Indices are local to this worked theory and do not reuse the preceding four-rule countermodel's rules. \\
\phantomsection\label{entry:symbol-registry-01831}\(\Gamma=(S,\mathcal D),x\triangleleft y\) & \texttt{\detokenize{\Gamma=(S,}\allowbreak\detokenize{\mathcal D),}\allowbreak\detokenize{x\triangleleft y}} & dynamic rule theory and object-language priority predicate & P18 is-dynamic & meta & \hyperref[def:p18-is-dynamic-language]{\ref*{def:p18-is-dynamic-language}} & \thepage & Only the four declared strict priority schemata mention this predicate in strict rules; other priority information is defeasible. \\
\phantomsection\label{entry:symbol-registry-01832}\(\mathcal A,X,E_X,<_X,\operatorname{Reb}_X,D_X,D_X^s\) & \texttt{\detokenize{\mathcal A,}\allowbreak\detokenize{X,}\allowbreak\detokenize{E_X,}\allowbreak\detokenize{<_X,}\allowbreak\detokenize{\operatorname{Reb}_X,}\allowbreak\detokenize{D_X,}\allowbreak\detokenize{D_X^s}} & argument universe support extracted comparisons frontier lift rebuttal defeat and strict defeat & P18 is-dynamic & meta & \hyperref[def:p18-is-dynamic-comparison]{\ref*{def:p18-is-dynamic-comparison}} & \thepage & Extraction reads actual conclusions; an arbitrary compatible support is not silently replaced by its transitive closure or by a merged argument. \\
\phantomsection\label{entry:symbol-registry-01833}\(G,J_\alpha,J\) & \texttt{\detokenize{G,}\allowbreak\detokenize{J_\alpha,}\allowbreak\detokenize{J}} & compatible-support operator ordinal justification stages and least compatible fixed point & P18 is-dynamic & meta & \hyperref[def:p18-is-dynamic-semantics]{\ref*{def:p18-is-dynamic-semantics}} & \thepage & Monotonicity is proved on conflict-free support; finitary means finitely many generated attackers per argument. \\
\phantomsection\label{entry:symbol-registry-01834}\(\mathcal D(A)\) & \texttt{\detokenize{\mathcal D(A)}} & finite defeasible-rule content of one argument & P18 is-dynamic & meta & \hyperref[annex:p18.is-dynamic-finitary]{\ref*{annex:p18.is-dynamic-finitary}} & \thepage & Strict extensions add no defeasible rules, so their frontier comparisons depend on finitely many priority bits. \\
\phantomsection\label{entry:symbol-registry-01835}\(P,O,A,B,C\) & \texttt{\detokenize{P,}\allowbreak\detokenize{O,}\allowbreak\detokenize{A,}\allowbreak\detokenize{B,}\allowbreak\detokenize{C}} & source neutralization-game players current claim challenge and single response & P18 is-dynamic & meta & \hyperref[def:p18-is-dynamic-game]{\ref*{def:p18-is-dynamic-game}} & \thepage & Opponent challenges use empty support; the original response uses its own priority conclusions; completeness is refuted in P18.17. \\
\phantomsection\label{entry:symbol-registry-01836}\(X_B,C_B\) & \texttt{\detokenize{X_B,}\allowbreak\detokenize{C_B}} & separate finite support and selected defender in the added certificate construction & P18 is-dynamic & meta & \hyperref[def:p18-is-dynamic-certificate]{\ref*{def:p18-is-dynamic-certificate}} & \thepage & This procedure is explicitly additional and does not call the support family one source argument. \\
\phantomsection\label{entry:symbol-registry-01837}\(D_\Gamma,D_\Gamma^s\) & \texttt{\detokenize{D_\Gamma,}\allowbreak\detokenize{D_\Gamma^s}} & fixed comparison used in the reversed-role game & P18 is-dynamic & meta & \hyperref[def:p18-is-dynamic-weak-game]{\ref*{def:p18-is-dynamic-weak-game}} & \thepage & The fixed-comparison result does not supply an unstated dynamic player-swap theorem. \\
\phantomsection\label{entry:symbol-registry-01838}\(r_i,X,Y\) & \texttt{\detokenize{r_i,}\allowbreak\detokenize{X,}\allowbreak\detokenize{Y}} & local indexed rules and nested supports in the incompatible-support counterexample & P18 is-dynamic & meta & \hyperref[def:p18-is-dynamic-outside-domain]{\ref*{def:p18-is-dynamic-outside-domain}} & \thepage & The larger support has conflicting priorities; subsequent priority examples redeclare their indexed rules locally. \\
\phantomsection\label{entry:symbol-registry-01839}\(R_P,R_M,K,T,K_1,T_1,T_2,v\) & \texttt{\detokenize{R_P,}\allowbreak\detokenize{R_M,}\allowbreak\detokenize{K,}\allowbreak\detokenize{T,}\allowbreak\detokenize{K_1,}\allowbreak\detokenize{T_1,}\allowbreak\detokenize{T_2,}\allowbreak\detokenize{v}} & building rules priority schemata distinct ground priority instances and building object & P18 is-dynamic & meta & \hyperref[def:p18-is-dynamic-building]{\ref*{def:p18-is-dynamic-building}} & \thepage & T2 is a different ground instance from T1; its factual premise is supplied and no rule asserts its own superiority. \\
\phantomsection\label{entry:symbol-registry-01840}\(E,L,E_0,L_0,P_0,B,I,q,h\) & \texttt{\detokenize{E,}\allowbreak\detokenize{L,}\allowbreak\detokenize{E_0,}\allowbreak\detokenize{L_0,}\allowbreak\detokenize{P_0,}\allowbreak\detokenize{B,}\allowbreak\detokenize{I,}\allowbreak\detokenize{q,}\allowbreak\detokenize{h}} & sale and local rule schemata selected ground instances product jurisdictions priority default and exception default & P18 is-dynamic & meta & \hyperref[def:p18-is-dynamic-exception]{\ref*{def:p18-is-dynamic-exception}} & \thepage & Only the declared ground instances and facts belong to this worked theory; the rule L is not a licensed list. \\
\phantomsection\label{entry:symbol-registry-01841}\(\operatorname{app}(r),\rho_r\) & \texttt{\detokenize{\operatorname{app}(r),}\allowbreak\detokenize{\rho_r}} & applicability atom and fresh positive-applicability default & P18 is-dynamic & meta & \hyperref[def:p18-is-dynamic-app-encoding]{\ref*{def:p18-is-dynamic-app-encoding}} & \thepage & Negative and positive premise translations differ; added helper defaults are not recursively gated and priorities must name the helper when it is the frontier. \\
\phantomsection\label{entry:symbol-registry-01842}\(\operatorname{Ex}_X,D_X^{\rm ex}\) & \texttt{\detokenize{\operatorname{Ex}_X,}\allowbreak\detokenize{D_X^{\rm ex}}} & direct exclusion and its modified defeat relation & P18 is-dynamic & meta & \hyperref[def:p18-is-dynamic-exclusion]{\ref*{def:p18-is-dynamic-exclusion}} & \thepage & The inapplicability argument attacks the rule use; priority-free and sensitive variants have different proved semantic scope. \\
\phantomsection\label{entry:symbol-registry-01843}\(r,s,k,\rho_r\) & \texttt{\detokenize{r,}\allowbreak\detokenize{s,}\allowbreak\detokenize{k,}\allowbreak\detokenize{\rho_r}} & local domain rule inapplicability rule priority rule and applicability helper in the boundary cases & P18 is-dynamic & meta & \hyperref[def:p18-is-dynamic-app-boundaries]{\ref*{def:p18-is-dynamic-app-boundaries}} & \thepage & Sensitive direct exclusion fails to preserve compatible accepted sets; the positive-premise example is a different theory. \\
\phantomsection\label{entry:symbol-registry-01844}\(a,b,c,d,u,v,t,z,m,n,p,q,U,V,P,Q,T,Z,A,B\) & \texttt{\detokenize{a,}\allowbreak\detokenize{b,}\allowbreak\detokenize{c,}\allowbreak\detokenize{d,}\allowbreak\detokenize{u,}\allowbreak\detokenize{v,}\allowbreak\detokenize{t,}\allowbreak\detokenize{z,}\allowbreak\detokenize{m,}\allowbreak\detokenize{n,}\allowbreak\detokenize{p,}\allowbreak\detokenize{q,}\allowbreak\detokenize{U,}\allowbreak\detokenize{V,}\allowbreak\detokenize{P,}\allowbreak\detokenize{Q,}\allowbreak\detokenize{T,}\allowbreak\detokenize{Z,}\allowbreak\detokenize{A,}\allowbreak\detokenize{B}} & countermodel atoms eight defeasible rules and their named argument contents & P18 is-dynamic-completeness & meta & \hyperref[def:p18-is-compression-countertheory]{\ref*{def:p18-is-compression-countertheory}} & \thepage & Capital P and Q are arguments; the two separately justified priority supports have incoherent combinations. All strict priority ground instances are included. \\
\phantomsection\label{entry:symbol-registry-01845}\(P_{\rm pl},O_{\rm pl},A_0\) & \texttt{\detokenize{P_{\rm pl},}\allowbreak\detokenize{O_{\rm pl},}\allowbreak\detokenize{A_0}} & sequential-game players and initial claim & P18 is-dynamic-completeness & meta & \hyperref[def:p18-is-licensing-state]{\ref*{def:p18-is-licensing-state}} & \thepage & Player subscripts distinguish the proponent from the counterexample argument P. \\
\phantomsection\label{entry:symbol-registry-01846}\(L,|L|,E_{|L|},E_L,D_L\) & \texttt{\detokenize{L,}\allowbreak\detokenize{|L|,}\allowbreak\detokenize{E_{|L|},}\allowbreak\detokenize{E_L,}\allowbreak\detokenize{D_L}} & licensed ordered list its entry set extracted priority context and context-relative defeat & P18 is-dynamic-completeness & meta & \hyperref[def:p18-is-licensing-state]{\ref*{def:p18-is-licensing-state}} & \thepage & E\_L and D\_L abbreviate extraction and defeat at the entry set; bars denote that set rather than cardinality and cooperative transcripts need not license only justified arguments. \\
\phantomsection\label{entry:symbol-registry-01847}\((C,B),H,\mathrm{Hear},\mathrm{Defend},\mathrm{Neutralize}\) & \texttt{\detokenize{(C,}\allowbreak\detokenize{B),}\allowbreak\detokenize{H,}\allowbreak\detokenize{Hear,}\allowbreak\detokenize{Defend,}\allowbreak\detokenize{Neutralize}} & active claim-challenge frame auxiliary or cited argument and the three move tags & P18 is-dynamic-completeness & meta & \hyperref[def:p18-is-licensing-moves]{\ref*{def:p18-is-licensing-moves}} & \thepage & Fresh hearings cannot repeat pending or licensed claims; closing a hearing returns to its parent and licensed support may be cited repeatedly. \\
\phantomsection\label{entry:symbol-registry-01848}\(\rho(C),\alpha,X,H_0,L'\) & \texttt{\detokenize{\rho(C),}\allowbreak\detokenize{\alpha,}\allowbreak\detokenize{X,}\allowbreak\detokenize{H_0,}\allowbreak\detokenize{L'}} & semantic entry rank predecessor stage separate finite witnesses defender and extended licensed list & P18 is-dynamic-completeness & meta & \hyperref[annex:p18.is-licensing-complete]{\ref*{annex:p18.is-licensing-complete}} & \thepage & Nested hearings have strictly smaller ordinal ranks; finite priority witnesses are heard separately rather than amalgamated. \\
\phantomsection\label{entry:symbol-registry-01849}\(N\) & \texttt{\detokenize{N}} & finite argument-universe size bounding fresh hearings & P18 is-dynamic-completeness & meta & \hyperref[annex:p18.is-licensing-finite]{\ref*{annex:p18.is-licensing-finite}} & \thepage & Every argument can be opened at most once in one play; finite grounding is distinguished from finitary infinite input. \\
\phantomsection\label{entry:symbol-registry-01850}\(r_a,r_{\neg a},r_b,r_s\) & \texttt{\detokenize{r_a,}\allowbreak\detokenize{r_{\neg a},}\allowbreak\detokenize{r_b,}\allowbreak\detokenize{r_s}} & four local rules in the cooperative-play boundary example & P18 is-dynamic-completeness & meta & \hyperref[def:p18-is-licensing-single-play]{\ref*{def:p18-is-licensing-single-play}} & \thepage & One completed transcript can have an unjustified root; soundness concerns a strategy covering every opponent choice. \\
\phantomsection\label{entry:symbol-registry-01851}\(A_0,B,J\) & \texttt{\detokenize{A_0,}\allowbreak\detokenize{B,}\allowbreak\detokenize{J}} & status target selected attack witness and source justified set in the overruling protocol & P18 is-dynamic-completeness & meta & \hyperref[def:p18-is-licensing-reversed]{\ref*{def:p18-is-licensing-reversed}} & \thepage & The challenger must select an actual attack witness; a designated empty incoherence defeat alone is not an attack. \\
\phantomsection\label{entry:symbol-registry-02038}\(\mathrm{Var},\mathrm{Form}\) & \texttt{\detokenize{\mathrm{Var},}\allowbreak\detokenize{\mathrm{Form}}} & countably infinite propositional variable set and generated formula set & P18 question evocation & meta & \hyperref[def:p18-qev-language]{\ref*{def:p18-qev-language}} & \thepage & Five classical connectives and parentheses; no primitive truth constants. \\
\phantomsection\label{entry:symbol-registry-02039}\(v\) & \texttt{\detokenize{v}} & classical valuation into zero and one & P18 question evocation & meta & \hyperref[def:p18-qev-language]{\ref*{def:p18-qev-language}} & \thepage & Source-local specialization to classical truth values. The source uses v for both the atomic assignment and its truth-functional extension; the global default writes the extension as hat-v. \\
\phantomsection\label{entry:symbol-registry-02040}\(?\{A_1,\ldots,A_n\}\) & \texttt{\detokenize{?\{A_1,\ldots,A_n\}}} & ordered object-language question expression & P18 question evocation & meta & \hyperref[def:p18-qev-language]{\ref*{def:p18-qev-language}} & \thepage & n greater than one; direct answers pairwise syntactically distinct, not necessarily inequivalent. \\
\phantomsection\label{entry:symbol-registry-02041}\(dQ\) & \texttt{\detokenize{dQ}} & set of direct answers to a question & P18 question evocation & meta & \hyperref[def:p18-qev-language]{\ref*{def:p18-qev-language}} & \thepage & Equal answer sets need not give the same question expression. \\
\phantomsection\label{entry:symbol-registry-02042}\(X,Y\) & \texttt{\detokenize{X,}\allowbreak\detokenize{Y}} & sets of declarative formulas & P18 question evocation & meta & \hyperref[def:p18-qev-entailment]{\ref*{def:p18-qev-entailment}} & \thepage & Semantics permits sets; erotetic-sequent antecedents are finite. \\
\phantomsection\label{entry:symbol-registry-02043}\(\models\) & \texttt{\detokenize{\models}} & ordinary classical single-conclusion entailment or satisfaction at a valuation & P18 question evocation & relation & \hyperref[def:p18-qev-entailment]{\ref*{def:p18-qev-entailment}} & \thepage & Operands distinguish the set-consequence and valuation-satisfaction roles. \\
\phantomsection\label{entry:symbol-registry-02044}\(\mathrel{\|=}\) & \texttt{\detokenize{\mathrel{\|=}}} & multiple-conclusion entailment & P18 question evocation & relation & \hyperref[def:p18-qev-entailment]{\ref*{def:p18-qev-entailment}} & \thepage & Every antecedent valuation satisfies at least one formula of the consequent set; no forcing role. \\
\phantomsection\label{entry:symbol-registry-02045}\(\mathrel{\|\!\triangleleft}\) & \texttt{\detokenize{\mathrel{\|\!\triangleleft}}} & proper multiple-conclusion entailment & P18 question evocation & relation & \hyperref[def:p18-qev-entailment]{\ref*{def:p18-qev-entailment}} & \thepage & Source double-bar plus left-pointing outlined triangle; nonempty consequent,coverage and non-entailment of every individual answer. \\
\phantomsection\label{entry:symbol-registry-02046}\(E(X,Q)\) & \texttt{\detokenize{E(X,}\allowbreak\detokenize{Q)}} & question evocation & P18 question evocation & meta & \hyperref[def:p18-qev-entailment]{\ref*{def:p18-qev-entailment}} & \thepage & Proper multiple-conclusion entailment of dQ. \\
\phantomsection\label{entry:symbol-registry-02047}\(X\vdash Q\) & \texttt{\detokenize{X\vdash Q}} & erotetic-sequent expression & P18 question evocation & relation & \hyperref[def:p18-qev-sequents]{\ref*{def:p18-qev-sequents}} & \thepage & Finite set antecedent and a question succedent; provability of this expression is stated separately. \\
\phantomsection\label{entry:symbol-registry-02048}\(X,B\) & \texttt{\detokenize{X,}\allowbreak\detokenize{B}} & set antecedent extended by one declarative & P18 question evocation & meta & \hyperref[def:p18-qev-sequents]{\ref*{def:p18-qev-sequents}} & \thepage & Abbreviates X union singleton B; repetitions do not create multiset multiplicities. \\
\phantomsection\label{entry:symbol-registry-02049}\(\vec A,\vec L,\vec R\) & \texttt{\detokenize{\vec A,}\allowbreak\detokenize{\vec L,}\allowbreak\detokenize{\vec R}} & ordered answer-list fragments in rule notation & P18 question evocation & meta & \hyperref[def:p18-qev-sequents]{\ref*{def:p18-qev-sequents}} & \thepage & Separate from sets; global question well-formedness conditions apply. \\
\phantomsection\label{entry:symbol-registry-02050}\(k,m_i\) & \texttt{\detokenize{k,}\allowbreak\detokenize{m_i}} & representative rank and clause counts & P18 question evocation & meta & \hyperref[def:p18-qev-normal-form]{\ref*{def:p18-qev-normal-form}} & \thepage & \(k=\sum_{i=1}^{n}(m_i-1)\); count chosen equivalent CNFs rather than extra syntactic markers. \\
\phantomsection\label{entry:symbol-registry-02051}\(I_X(A)\) & \texttt{\detokenize{I_X(A)}} & nested material-implication prefix for an enumeration of a finite antecedent & P18 question evocation & meta & \hyperref[thm:p18-qev-conditionalization]{\ref*{thm:p18-qev-conditionalization}} & \thepage & Empty enumeration returns A; fixed prefixes are syntactically injective. \\
\phantomsection\label{entry:symbol-registry-02052}\(?[A\mid n]\) & \texttt{\detokenize{?[A\mid n]}} & abbreviated ordered question with n answers & P18 question evocation & meta & \hyperref[def:p18-qev-rule-standing]{\ref*{def:p18-qev-rule-standing}} & \thepage & Source Section5 notation; right alpha/beta extensions explicitly include the one-common-answer cases. \\
\phantomsection\label{entry:symbol-registry-02053}\(\alpha,\alpha_1,\alpha_2,\beta,\beta_1,\beta_2\) & \texttt{\detokenize{\alpha,}\allowbreak\detokenize{\alpha_1,}\allowbreak\detokenize{\alpha_2,}\allowbreak\detokenize{\beta,}\allowbreak\detokenize{\beta_1,}\allowbreak\detokenize{\beta_2}} & selected conjunctive and disjunctive formula types and components & P18 question evocation & meta & \hyperref[def:p18-qev-alpha-beta]{\ref*{def:p18-qev-alpha-beta}} & \thepage & Exactly the three paired table rows; double negation is not an alpha form. \\
\phantomsection\label{entry:symbol-registry-02054}\(r\) & \texttt{\detokenize{r}} & number of distinct variables in a finite decision input & P18 question evocation & meta & \hyperref[thm:p18-qev-decision]{\ref*{thm:p18-qev-decision}} & \thepage & The \(2^r\) bound counts assignments, not unit-time evaluations. \\
\phantomsection\label{entry:symbol-registry-02387}\(C,D,\Gamma,\Delta\) & \texttt{\detokenize{C,}\allowbreak\detokenize{D,}\allowbreak\detokenize{\Gamma,}\allowbreak\detokenize{\Delta}} & target claim,critic conclusion and premise sets & P18 refutation diagnosis & meta & \hyperref[def:p18-rd-claim-record]{\ref*{def:p18-rd-claim-record}} & \thepage & The claim record fixes interpretation,domain,quantifiers,respect,relation,manner and time; letters can receive separately declared local example roles. \\
\phantomsection\label{entry:symbol-registry-02388}\(\mathcal K,\models_{\mathcal K}\) & \texttt{\detokenize{\mathcal K,}\allowbreak\detokenize{\models_{\mathcal K}}} & class of interpretations and consequence relative to it & P18 refutation diagnosis & relation & \hyperref[def:p18-rd-claim-record]{\ref*{def:p18-rd-claim-record}} & \thepage & A countermodel satisfies all premises and falsifies the proposed conclusion. \\
\phantomsection\label{entry:symbol-registry-02389}\(F,R,W\) & \texttt{\detokenize{F,}\allowbreak\detokenize{R,}\allowbreak\detokenize{W}} & financial-bank,river-bank and water-adjacency predicates & P18 refutation diagnosis bank example & n/a & \hyperref[def:p18-rd-bank]{\ref*{def:p18-rd-bank}} & \thepage & Local unary predicate roles; not the later binary parenthood predicate. \\
\phantomsection\label{entry:symbol-registry-02390}\(S,V,o,t\) & \texttt{\detokenize{S,}\allowbreak\detokenize{V,}\allowbreak\detokenize{o,}\allowbreak\detokenize{t}} & seeing relation,sight predicate,observer and target & P18 refutation diagnosis sight example & n/a & \hyperref[def:p18-rd-sight]{\ref*{def:p18-rd-sight}} & \thepage & Sight is licensed at the first argument of seeing,not its second. \\
\phantomsection\label{entry:symbol-registry-02391}\(S,W,\Diamond\) & \texttt{\detokenize{S,}\allowbreak\detokenize{W,}\allowbreak\detokenize{\Diamond}} & seated and walking propositions and possibility at an accessible world & P18 refutation diagnosis ability example & n/a & \hyperref[def:p18-rd-ability-scope]{\ref*{def:p18-rd-ability-scope}} & \thepage & These are local propositional roles in an explicitly defined two-world interpretation. \\
\phantomsection\label{entry:symbol-registry-02392}\(\mathrm{Odd}(n)\) & \texttt{\detokenize{\mathrm{Odd}(n)}} & oddness in the natural numbers & P18 refutation diagnosis arithmetic examples & n/a & \hyperref[def:p18-rd-formal-conventions]{\ref*{def:p18-rd-formal-conventions}} & \thepage & Defined by n=2k+1 for a natural k. \\
\phantomsection\label{entry:symbol-registry-02393}\(e_1,e_2,f,I\) & \texttt{\detokenize{e_1,}\allowbreak\detokenize{e_2,}\allowbreak\detokenize{f,}\allowbreak\detokenize{I}} & expression tokens,transcription map and interpretation map & P18 refutation diagnosis transcription example & meta & \hyperref[def:p18-rd-transcription]{\ref*{def:p18-rd-transcription}} & \thepage & Equal output of a noninjective transcription need not preserve interpretation. \\
\phantomsection\label{entry:symbol-registry-02394}\(H,A\) & \texttt{\detokenize{H,}\allowbreak\detokenize{A}} & shared-form relation and activity-denoting predicate on expressions & P18 refutation diagnosis category example & n/a & \hyperref[def:p18-rd-category-form]{\ref*{def:p18-rd-category-form}} & \thepage & Morphological form does not by itself imply semantic type. \\
\phantomsection\label{entry:symbol-registry-02395}\(H_0,H_1,|H|\) & \texttt{\detokenize{H_0,}\allowbreak\detokenize{H_1,}\allowbreak\detokenize{|H|}} & initial and retained possession sets and cardinality & P18 refutation diagnosis quantity example & meta & \hyperref[def:p18-rd-quantity-and-method]{\ref*{def:p18-rd-quantity-and-method}} & \thepage & Lost objects form the set difference; cardinality changes do not identify all initial items as lost. \\
\phantomsection\label{entry:symbol-registry-02396}\(L,D\) & \texttt{\detokenize{L,}\allowbreak\detokenize{D}} & learned and discovered predicates on items & P18 refutation diagnosis per-item-method example & n/a & \hyperref[def:p18-rd-quantity-and-method]{\ref*{def:p18-rd-quantity-and-method}} & \thepage & Local unary D differs from the critic's conclusion and the later arithmetic binary D. \\
\phantomsection\label{entry:symbol-registry-02397}\(O,F,a,d,b\) & \texttt{\detokenize{O,}\allowbreak\detokenize{F,}\allowbreak\detokenize{a,}\allowbreak\detokenize{d,}\allowbreak\detokenize{b}} & ownership and parenthood relations with named objects & P18 refutation diagnosis ownership example & n/a & \hyperref[def:p18-rd-ownership]{\ref*{def:p18-rd-ownership}} & \thepage & Both relation argument places are retained. \\
\phantomsection\label{entry:symbol-registry-02398}\(G(a,c),D(x,y)\) & \texttt{\detokenize{G(a,}\allowbreak\detokenize{c),}\allowbreak\detokenize{D(x,}\allowbreak\detokenize{y)}} & context-relative evaluation and arithmetic double-of relation & P18 refutation diagnosis qualification example & n/a & \hyperref[def:p18-rd-respects]{\ref*{def:p18-rd-respects}} & \thepage & G keeps its context argument; D(x,y) means x=2y. \\
\phantomsection\label{entry:symbol-registry-02399}\(p,q,h,u,w\) & \texttt{\detokenize{p,}\allowbreak\detokenize{q,}\allowbreak\detokenize{h,}\allowbreak\detokenize{u,}\allowbreak\detokenize{w}} & propositional letters in explicitly separate diagnostic examples & P18 refutation diagnosis & meta & \hyperref[def:p18-rd-formal-conventions]{\ref*{def:p18-rd-formal-conventions}} & \thepage & h is the disputed reductio hypothesis;u and w are source unreliability and report truth in the warning example. \\
\phantomsection\label{entry:symbol-registry-03856}\(H,H^c\) & \texttt{\detokenize{H,}\allowbreak\detokenize{H^c}} & toxicity hypothesis event and its complement & P18 rationality argument strength & n/a & \hyperref[def:p18-r55-detection-model]{\ref*{def:p18-r55-detection-model}} & \thepage & Events in a specified finite probability model; not argument identifiers or extension membership. \\
\phantomsection\label{entry:symbol-registry-03857}\(E,E^c\) & \texttt{\detokenize{E,}\allowbreak\detokenize{E^c}} & observed-effect event and its complement & P18 rationality argument strength & n/a & \hyperref[def:p18-r55-detection-model]{\ref*{def:p18-r55-detection-model}} & \thepage & Absence means the indicated experiment observed no effect; a test not performed is different information. \\
\phantomsection\label{entry:symbol-registry-03858}\(\pi\) & \texttt{\detokenize{\pi}} & prior probability of toxicity & P18 rationality argument strength & n/a & \hyperref[def:p18-r55-detection-model]{\ref*{def:p18-r55-detection-model}} & \thepage & Assumed strictly between zero and one for the displayed comparison theorem. \\
\phantomsection\label{entry:symbol-registry-03859}\(s,c\) & \texttt{\detokenize{s,}\allowbreak\detokenize{c}} & sensitivity and specificity of the stated detection experiment & P18 rationality argument strength & n/a & \hyperref[def:p18-r55-detection-model]{\ref*{def:p18-r55-detection-model}} & \thepage & Conditional probabilities given opposing hypothesis events; neither is a posterior. \\
\phantomsection\label{entry:symbol-registry-03860}\(p_+,q_-\) & \texttt{\detokenize{p_+,}\allowbreak\detokenize{q_-}} & toxicity-after-effect and nontoxicity-after-absence posteriors & P18 rationality argument strength & n/a & \hyperref[thm:p18-r55-detection]{\ref*{thm:p18-r55-detection}} & \thepage & Each is defined only for its positive-probability evidence event; their prior increases require sensitivity plus specificity greater than one. \\
\phantomsection\label{entry:symbol-registry-03861}\(h,r,p\) & \texttt{\detokenize{h,}\allowbreak\detokenize{r,}\allowbreak\detokenize{p}} & binary truth indicator and probability report under conditional quadratic loss & P18 rationality argument strength & n/a & \hyperref[thm:p18-r55-norm-boundaries]{\ref*{thm:p18-r55-norm-boundaries}} & \thepage & h is zero or one; p is its conditional chance of one and r the report being scored. \\
\phantomsection\label{entry:symbol-registry-03862}\(E_1,E_2\) & \texttt{\detokenize{E_1,}\allowbreak\detokenize{E_2}} & two identical report events in the duplication boundary & P18 rationality argument strength & n/a & \hyperref[thm:p18-r55-norm-boundaries]{\ref*{thm:p18-r55-norm-boundaries}} & \thepage & Their set equality is stipulated; differently worded or differently sourced reports are not automatically independent. \\
\phantomsection\label{entry:symbol-registry-03863}\(q\) & \texttt{\detokenize{q}} & unfavorable-outcome probability in a stipulated action comparison & P18 rationality argument strength & n/a & \hyperref[thm:p18-r55-norm-boundaries]{\ref*{thm:p18-r55-norm-boundaries}} & \thepage & Expected-utility threshold depends on the supplied positive and negative utilities, not only on posterior confidence. \\
\phantomsection\label{entry:symbol-registry-03864}\(\kappa\) & \texttt{\detokenize{\kappa}} & reported interrater agreement coefficient & P18 rationality argument strength & n/a & \hyperref[sec:p18-rationality-5-5]{\ref*{sec:p18-rationality-5-5}} & \thepage & Cohen's kappa from a bounded qualitative study record; not a probability of correct argumentation or a code-domain cardinal. \\
\phantomsection\label{entry:symbol-registry-03865}\(\mathcal F\) & \texttt{\detokenize{\mathcal F}} & attack-free one-argument abstract framework in a separation construction & P18 rationality argument strength & n/a & \hyperref[thm:p18-r55-graph-boundary]{\ref*{thm:p18-r55-graph-boundary}} & \thepage & The same extension verdict accompanies two coherent likelihood models with opposite confirmation directions. \\
\phantomsection\label{entry:symbol-registry-03866}\(D_+,D_-\) & \texttt{\detokenize{D_+,}\allowbreak\detokenize{D_-}} & positive evidence probabilities in the detection proof & P18 rationality argument strength & n/a & \hyperref[annex:p18.r55.detection]{\ref*{annex:p18.r55.detection}} & \thepage & Actual effect and absence probabilities; zero denominators are not assigned conditional values. \\
\phantomsection\label{entry:symbol-registry-03867}\(H\) & \texttt{\detokenize{H}} & proposition in the deductively valid circularity illustration & P18 rationality circularity illustration & n/a & \hyperref[sec:p18-r55-normative-argument]{\ref*{sec:p18-r55-normative-argument}} & \thepage & Locally declares the proposition repeated as premise and conclusion; the later detection model separately declares H as an event. \\
\phantomsection\label{entry:symbol-registry-03978}\(\mathsf{Spot}(a,\Gamma,\varphi)\) & \texttt{\detokenize{\mathsf{Spot}(a,}\allowbreak\detokenize{\Gamma,}\allowbreak\detokenize{\varphi)}} & agent-relative consequence-recognition question & P18 reasoning and argumentation distinction & n/a & \hyperref[def:p18-r56-three-consequence-roles]{\ref*{def:p18-r56-three-consequence-roles}} & \thepage & A notation for the selected distinction, not an empirical measure or a complete cognitive theory. \\
\phantomsection\label{entry:symbol-registry-03979}\(\mathsf{Draw}(a,\Gamma,\varphi)\) & \texttt{\detokenize{\mathsf{Draw}(a,}\allowbreak\detokenize{\Gamma,}\allowbreak\detokenize{\varphi)}} & agent-relative inferential-step question & P18 reasoning and argumentation distinction & n/a & \hyperref[def:p18-r56-three-consequence-roles]{\ref*{def:p18-r56-three-consequence-roles}} & \thepage & Actual drawing is distinct from static entailment and correct recognition; goals and response policies are additional. \\
\phantomsection\label{entry:symbol-registry-03980}\(K\) & \texttt{\detokenize{K}} & inconsistent finite belief set in a repair separation & P18 reasoning and argumentation distinction & n/a & \hyperref[thm:p18-r56-consequence-policy]{\ref*{thm:p18-r56-consequence-policy}} & \thepage & Three different maximal consistent subsets satisfy the same bare repair criterion; a preference is not supplied by consequence alone. \\
\phantomsection\label{entry:symbol-registry-03981}\(p,q_0,q_1,\ldots\) & \texttt{\detokenize{p,}\allowbreak\detokenize{q_0,}\allowbreak\detokenize{q_1,}\allowbreak\detokenize{\ldots}} & distinct propositional variables in the explicit-storage boundary & P18 reasoning and argumentation distinction & n/a & \hyperref[thm:p18-r56-explicit-closure]{\ref*{thm:p18-r56-explicit-closure}} & \thepage & The conclusion is about syntactically explicit formula records; it does not exclude a finite generator or selected queries. \\
\phantomsection\label{entry:symbol-registry-03982}\(m\) & \texttt{\detokenize{m}} & finite explicit formula-record capacity & P18 reasoning and argumentation distinction & n/a & \hyperref[thm:p18-r56-explicit-closure]{\ref*{thm:p18-r56-explicit-closure}} & \thepage & A stipulated finite store size, not a measured parameter of human memory. \\
\phantomsection\label{entry:symbol-registry-03983}\(C,G,W,B,R,Q\) & \texttt{\detokenize{C,}\allowbreak\detokenize{G,}\allowbreak\detokenize{W,}\allowbreak\detokenize{B,}\allowbreak\detokenize{R,}\allowbreak\detokenize{Q}} & claim grounds warrant backing rebuttal and qualifier roles in the selected Toulmin discussion & P18 Woods Toulmin presentation & n/a & \hyperref[def:p18-r56-toulmin-roles]{\ref*{def:p18-r56-toulmin-roles}} & \thepage & Published role letters retained; the printed Nigel example's R-labeled support sentence is distinguished from its later exception. \\
\phantomsection\label{entry:symbol-registry-03984}\(\mathsf{Red},\mathsf{Colored},\mathsf{RedHue}\) & \texttt{\detokenize{\mathsf{Red},}\allowbreak\detokenize{\mathsf{Colored},}\allowbreak\detokenize{\mathsf{RedHue}}} & predicates in an explicit lexical-representation construction & P18 reasoning and argumentation distinction & n/a & \hyperref[thm:p18-r56-lexical-representation]{\ref*{thm:p18-r56-lexical-representation}} & \thepage & Unconstrained predicates give a countermodel; a declared meaning postulate or decomposition restores the represented consequence. \\
\phantomsection\label{entry:symbol-registry-03985}\(t\) & \texttt{\detokenize{t}} & object-denoting constant in the shirt representation & P18 reasoning and argumentation distinction & n/a & \hyperref[thm:p18-r56-lexical-representation]{\ref*{thm:p18-r56-lexical-representation}} & \thepage & Its denotation is fixed within each first-order interpretation; no citizenship law is inferred from this mathematical example. \\
\phantomsection\label{entry:symbol-registry-03986}\(A_\Gamma,\chi_\Gamma\) & \texttt{\detokenize{A_\Gamma,}\allowbreak\detokenize{\chi_\Gamma}} & premise conjunction and conditional added in arbitrary deductivization & P18 reasoning and argumentation distinction & n/a & \hyperref[thm:p18-r56-conditional-repair]{\ref*{thm:p18-r56-conditional-repair}} & \thepage & The empty conjunction is true; each original countervaluation falsifies the added conditional. \\
\phantomsection\label{entry:symbol-registry-04124}\(p,q,r\) & \texttt{\detokenize{p,}\allowbreak\detokenize{q,}\allowbreak\detokenize{r}} & binary agenda propositions & P18 Social Epistemology & n/a & \hyperref[thm:p18-r101-aggregation-failures]{\ref*{thm:p18-r101-aggregation-failures}} & \thepage & Local court model declares r equivalent to the conjunction of p and q; not a universal legal condition. \\
\phantomsection\label{entry:symbol-registry-04125}\(a,b,c,a_1,a_2,a_3\) & \texttt{\detokenize{a,}\allowbreak\detokenize{b,}\allowbreak\detokenize{c,}\allowbreak\detokenize{a_1,}\allowbreak\detokenize{a_2,}\allowbreak\detokenize{a_3}} & ranked or plurality alternatives & P18 Social Epistemology aggregation failures & n/a & \hyperref[thm:p18-r101-aggregation-failures]{\ref*{thm:p18-r101-aggregation-failures}} & \thepage & Typed alternatives in two distinct counterexamples; not the worlds in the later probability pool. \\
\phantomsection\label{entry:symbol-registry-04126}\(\Omega,\omega,\nu\) & \texttt{\detokenize{\Omega,}\allowbreak\detokenize{\omega,}\allowbreak\detokenize{\nu}} & finite world space and member worlds & P18 Social Epistemology opinion pooling & n/a & \hyperref[def:p18-r101-pools]{\ref*{def:p18-r101-pools}} & \thepage & Nonempty finite space; strict mass positivity for the geometric pool. \\
\phantomsection\label{entry:symbol-registry-04127}\(P_i,w_i\) & \texttt{\detokenize{P_i,}\allowbreak\detokenize{w_i}} & individual probability mass and fixed pooling weights & P18 Social Epistemology opinion pooling & n/a & \hyperref[def:p18-r101-pools]{\ref*{def:p18-r101-pools}} & \thepage & Weights nonnegative with sum one; individual probability mass is not empirical competence. \\
\phantomsection\label{entry:symbol-registry-04128}\(L,G\) & \texttt{\detokenize{L,}\allowbreak\detokenize{G}} & linear and normalized geometric probability pools & P18 Social Epistemology opinion pooling & n/a & \hyperref[def:p18-r101-pools]{\ref*{def:p18-r101-pools}} & \thepage & Operations on a finite common event algebra; not general testimonial updating operators. \\
\phantomsection\label{entry:symbol-registry-04129}\(A,B,E,a,b,c,a_1,a_2\) & \texttt{\detokenize{A,}\allowbreak\detokenize{B,}\allowbreak\detokenize{E,}\allowbreak\detokenize{a,}\allowbreak\detokenize{b,}\allowbreak\detokenize{c,}\allowbreak\detokenize{a_1,}\allowbreak\detokenize{a_2}} & coarse events learned event and refined worlds & P18 Social Epistemology opinion pooling & n/a & \hyperref[thm:p18-r101-pooling]{\ref*{thm:p18-r101-pooling}} & \thepage & Disjoint local typing from the preceding alternatives; E has positive probability for all agents. \\
\phantomsection\label{entry:symbol-registry-04130}\(X_i,S_n,n\) & \texttt{\detokenize{X_i,}\allowbreak\detokenize{S_n,}\allowbreak\detokenize{n}} & correct-vote indicators sum and odd group size & P18 Social Epistemology jury model & n/a & \hyperref[def:p18-r101-jury-model]{\ref*{def:p18-r101-jury-model}} & \thepage & Binary correctness relative to one declared problem; independence is stochastic not absence of conversation. \\
\phantomsection\label{entry:symbol-registry-04131}\(t,M_n(t)\) & \texttt{\detokenize{t,}\allowbreak\detokenize{M_n(t)}} & common correct-vote probability and binomial majority reliability & P18 Social Epistemology jury model & n/a & \hyperref[def:p18-r101-jury-model]{\ref*{def:p18-r101-jury-model}} & \thepage & Conditional formal parameters; not fitted electorate accuracy or social legitimacy. \\
\phantomsection\label{entry:symbol-registry-04132}\(\Pi,\mu_j,t_j\) & \texttt{\detokenize{\Pi,}\allowbreak\detokenize{\mu_j,}\allowbreak\detokenize{t_j}} & finite common-problem variable masses and conditional competences & P18 Social Epistemology conditional jury model & n/a & \hyperref[thm:p18-r101-jury]{\ref*{thm:p18-r101-jury}} & \thepage & Reflected mass dominance not merely high mean competence; direct influence can defeat conditional independence. \\
\phantomsection\label{entry:symbol-registry-04133}\(m,d,\delta\) & \texttt{\detokenize{m,}\allowbreak\detokenize{d,}\allowbreak\detokenize{\delta}} & half-group index reflected offset and deviation threshold & P18 Social Epistemology jury proof & n/a & \hyperref[thm:p18-r101-jury]{\ref*{thm:p18-r101-jury}} & \thepage & Auxiliary proof quantities in the owned reciprocal annex; strict increase excludes degenerate endpoints. \\
\phantomsection\label{entry:symbol-registry-04134}\(D(X,f)\) & \texttt{\detokenize{D(X,}\allowbreak\detokenize{f)}} & performance of an action token by an individual at the fixed time & P18 Meggle communicative-action treatment & n/a & \hyperref[def:p18-r103-action-notation]{\ref*{def:p18-r103-action-notation}} & \thepage & Published qualitative action notation; not the P8 distributed-knowledge modality. \\
\phantomsection\label{entry:symbol-registry-04135}\(B(X,A)\) & \texttt{\detokenize{B(X,}\allowbreak\detokenize{A)}} & strong belief or conviction in the selected action explication & P18 Meggle communicative-action treatment & n/a & \hyperref[def:p18-r103-action-notation]{\ref*{def:p18-r103-action-notation}} & \thepage & Belief is not factive; closure is an explicit additional assumption in the reflexivity derivation. \\
\phantomsection\label{entry:symbol-registry-04136}\(W(X,A)\) & \texttt{\detokenize{W(X,}\allowbreak\detokenize{A)}} & wanting an outcome in the selected action explication & P18 Meggle communicative-action treatment & n/a & \hyperref[def:p18-r103-action-notation]{\ref*{def:p18-r103-action-notation}} & \thepage & Not a world set; the desire bridge is substantive and idealized. \\
\phantomsection\label{entry:symbol-registry-04137}\(I(X,f,A)\) & \texttt{\detokenize{I(X,}\allowbreak\detokenize{f,}\allowbreak\detokenize{A)}} & intention attributed to an actual performance & P18 Meggle communicative-action treatment & n/a & \hyperref[def:p18-r103-action-notation]{\ref*{def:p18-r103-action-notation}} & \thepage & Defined by action desire and believed material equivalence; not an intention existing without action. \\
\phantomsection\label{entry:symbol-registry-04138}\(X,Y,S,H\) & \texttt{\detokenize{X,}\allowbreak\detokenize{Y,}\allowbreak\detokenize{S,}\allowbreak\detokenize{H}} & individual performer interpreter speaker and addressed hearer & P18 Meggle communicative-action treatment & n/a & \hyperref[def:p18-r103-action-notation]{\ref*{def:p18-r103-action-notation}} & \thepage & Y may be a third person; X is an agent rather than the temporal next operator. \\
\phantomsection\label{entry:symbol-registry-04139}\(f,F\) & \texttt{\detokenize{f,}\allowbreak\detokenize{F}} & action token and action type & P18 Meggle communicative-action treatment & n/a & \hyperref[def:p18-r103-action-notation]{\ref*{def:p18-r103-action-notation}} & \thepage & Performance types are distinct from their product forms. \\
\phantomsection\label{entry:symbol-registry-04140}\(A,p,Q\) & \texttt{\detokenize{A,}\allowbreak\detokenize{p,}\allowbreak\detokenize{Q}} & proposition used as aim communicative content or schematic desired consequence & P18 Meggle communicative-action treatment & n/a & \hyperref[def:p18-r103-action-notation]{\ref*{def:p18-r103-action-notation}} & \thepage & A content-bearing attempt does not require its content to be true; Q is local to the desire-bridge schema. \\
\phantomsection\label{entry:symbol-registry-04141}\(t,t',B',D',K'\) & \texttt{\detokenize{t,}\allowbreak\detokenize{t',}\allowbreak\detokenize{B',}\allowbreak\detokenize{D',}\allowbreak\detokenize{K'}} & time and later-time attitude or performance reports & P18 Meggle communicative-action treatment & n/a & \hyperref[def:p18-r103-action-notation]{\ref*{def:p18-r103-action-notation}} & \thepage & Primes indicate a later time not a derivative or a new degree of belief. \\
\phantomsection\label{entry:symbol-registry-04142}\(K(X,A)\) & \texttt{\detokenize{K(X,}\allowbreak\detokenize{A)}} & source's provisional knowledge-as-true-conviction abbreviation & P18 Meggle communicative-action treatment & n/a & \hyperref[sec:p18-rationality-10-3]{\ref*{sec:p18-rationality-10-3}} & \thepage & The source acknowledges this is not an adequate general analysis of warranted knowledge. \\
\phantomsection\label{entry:symbol-registry-04143}\(U(Y,D(X,f)),U(Y,CA(S,H,f,p))\) & \texttt{\detokenize{U(Y,}\allowbreak\detokenize{D(X,}\allowbreak\detokenize{f)),}\allowbreak\detokenize{U(Y,}\allowbreak\detokenize{CA(S,}\allowbreak\detokenize{H,}\allowbreak\detokenize{f,}\allowbreak\detokenize{p))}} & understanding an action or communicative attempt in the selected definitions & P18 Meggle communicative-action treatment & n/a & \hyperref[sec:p18-rationality-10-3]{\ref*{sec:p18-rationality-10-3}} & \thepage & Recognition of content-bearing intent is not endorsement of content. \\
\phantomsection\label{entry:symbol-registry-04144}\(A^*\) & \texttt{\detokenize{A^*}} & source shorthand for all intended aims of a performance & P18 Meggle communicative-action treatment & n/a & \hyperref[sec:p18-rationality-10-3]{\ref*{sec:p18-rationality-10-3}} & \thepage & Combination of aims is a simplifying representation not a general existence theorem. \\
\phantomsection\label{entry:symbol-registry-04145}\(CA(S,H,f,p)\) & \texttt{\detokenize{CA(S,}\allowbreak\detokenize{H,}\allowbreak\detokenize{f,}\allowbreak\detokenize{p)}} & informative communicative attempt with content p & P18 Meggle communicative-action treatment & n/a & \hyperref[sec:p18-rationality-10-3]{\ref*{sec:p18-rationality-10-3}} & \thepage & AC1 is explicitly an adequacy condition rather than a circular definition. \\
\phantomsection\label{entry:symbol-registry-04146}\(I_0,I_n,I^*\) & \texttt{\detokenize{I_0,}\allowbreak\detokenize{I_n,}\allowbreak\detokenize{I^*}} & initial finite-level and total communicative-intention hierarchy & P18 Meggle communicative-action treatment & n/a & \hyperref[def:p18-r103-transparency-boundary]{\ref*{def:p18-r103-transparency-boundary}} & \thepage & The unbounded total is a source schematic idealization; no finitary calculus existence or closure is claimed. \\
\phantomsection\label{entry:symbol-registry-04147}\(P,\Sigma,M_0,M,K^*\) & \texttt{\detokenize{P,}\allowbreak\detokenize{\Sigma,}\allowbreak\detokenize{M_0,}\allowbreak\detokenize{M,}\allowbreak\detokenize{K^*}} & population situation type regular meaning intersubjective meaning and common knowledge & P18 Meggle group-meaning treatment & n/a & \hyperref[def:p18-r103-meaning-levels]{\ref*{def:p18-r103-meaning-levels}} & \thepage & K-star uses the finite-agent P8 common-knowledge semantics only in the precisely conditional AC2 proof. \\
\phantomsection\label{entry:symbol-registry-04148}\(d,b,r,s,q,c\) & \texttt{\detokenize{d,}\allowbreak\detokenize{b,}\allowbreak\detokenize{r,}\allowbreak\detokenize{s,}\allowbreak\detokenize{q,}\allowbreak\detokenize{c}} & local abbreviations for action uptake recognition speaker belief regularity and attempt propositions & P18 Meggle bounded annex derivations & n/a & \hyperref[thm:p18-r103-reflexivity]{\ref*{thm:p18-r103-reflexivity}} & \thepage & Typed only within the indicated derivations; not a new full communication model. \\
\phantomsection\label{entry:symbol-registry-04149}\(a\) & \texttt{\detokenize{a}} & hearer's later action in the source directive substitution & P18 Meggle communicative-action treatment & n/a & \hyperref[sec:p18-rationality-10-3]{\ref*{sec:p18-rationality-10-3}} & \thepage & Replacing a target belief by a target action does not identify all institutional felicity conditions. \\
\phantomsection\label{entry:symbol-registry-04150}\(C_P,R_P,w,u,v\) & \texttt{\detokenize{C_P,}\allowbreak\detokenize{R_P,}\allowbreak\detokenize{w,}\allowbreak\detokenize{u,}\allowbreak\detokenize{v}} & group common-knowledge modality its reachability relation and possible worlds in the bounded semantic proof & P18 Meggle common-meaning annex & n/a & \hyperref[thm:p18-r103-common-meaning]{\ref*{thm:p18-r103-common-meaning}} & \thepage & Reuses the actual P8 finite-agent epistemic semantics; R\_P is the reflexive-transitive closure of the union of agent relations. \\
\phantomsection\label{entry:symbol-registry-04668}\(n\) & \texttt{\detokenize{n}} & analysed participant count in the original METI studies & P18 public scientific trust & n/a & \hyperref[def:p18-r143-statistical-notation]{\ref*{def:p18-r143-statistical-notation}} & \thepage & Local scope anchored by sec:p18-rationality-14-3; completions and survey activations remain distinct. \\
\phantomsection\label{entry:symbol-registry-04669}\(r\) & \texttt{\detokenize{r}} & reported correlation between measured trustworthiness factors & P18 public scientific trust & n/a & \hyperref[def:p18-r143-statistical-notation]{\ref*{def:p18-r143-statistical-notation}} & \thepage & Not the reliability or truth probability of an actual scientist's claim. \\
\phantomsection\label{entry:symbol-registry-04670}\(\alpha\) & \texttt{\detokenize{\alpha}} & reported internal-consistency coefficient for METI items & P18 public scientific trust & n/a & \hyperref[def:p18-r143-statistical-notation]{\ref*{def:p18-r143-statistical-notation}} & \thepage & Item consistency does not establish external accuracy or construct uniqueness. \\
\phantomsection\label{entry:symbol-registry-04671}\(F\) & \texttt{\detokenize{F}} & reported repeated-measures comparison statistic with degrees of freedom & P18 public scientific trust & n/a & \hyperref[def:p18-r143-statistical-notation]{\ref*{def:p18-r143-statistical-notation}} & \thepage & The original study's corrected degrees of freedom and within-person dependence are preserved. \\
\phantomsection\label{entry:symbol-registry-04672}\(p\) & \texttt{\detokenize{p}} & reported significance level under the original analysis & P18 public scientific trust & n/a & \hyperref[def:p18-r143-statistical-notation]{\ref*{def:p18-r143-statistical-notation}} & \thepage & Not the probability that a source is truthful or an empirical public-decision efficacy rate. \\
\phantomsection\label{entry:symbol-registry-05994}qal va-homer & \texttt{\detokenize{\emph{qal va-homer}}} & Traditional name for the stronger-case argument discussed in the primary translation & P18 selected Talmudic arguments & n/a & \hyperref[sec:p18-talmudic-inference-principles]{\ref*{sec:p18-talmudic-inference-principles}} & \thepage & Transliterated terminology; no formal calculus or global severity scale Original precedence 'not an operator' retained as semantic descriptor; n/a is not an invented rank. \\
\phantomsection\label{entry:symbol-registry-05995}dayo & \texttt{\detokenize{\emph{dayo}}} & Traditional limiting principle rendered as the sufficient result & P18 selected Talmudic arguments & n/a & \hyperref[sec:p18-talmudic-inference-principles]{\ref*{sec:p18-talmudic-inference-principles}} & \thepage & Its disputed application and Tarphon's exception remain separate positions Original precedence 'not an operator' retained as semantic descriptor; n/a is not an invented rank. \\
\phantomsection\label{entry:symbol-registry-05996}horn & \texttt{\detokenize{\emph{horn}}} & Source-specific damage category in the specified comparison cases & P18 selected Talmudic arguments & n/a & \hyperref[sec:p18-talmudic-damage-premises]{\ref*{sec:p18-talmudic-damage-premises}} & \thepage & English category name, not a mathematical glyph; full case conditions remain necessary Original precedence 'not an operator' retained as semantic descriptor; n/a is not an invented rank. \\
\phantomsection\label{entry:symbol-registry-05997}tooth & \texttt{\detokenize{\emph{tooth}}} & Source-specific characteristic consumption category & P18 selected Talmudic arguments & n/a & \hyperref[sec:p18-talmudic-damage-premises]{\ref*{sec:p18-talmudic-damage-premises}} & \thepage & No damage liability is not absence of every payment for benefit Original precedence 'not an operator' retained as semantic descriptor; n/a is not an invented rank. \\
\phantomsection\label{entry:symbol-registry-05998}foot & \texttt{\detokenize{\emph{foot}}} & Source-specific characteristic walking-damage category & P18 selected Talmudic arguments & n/a & \hyperref[sec:p18-talmudic-damage-premises]{\ref*{sec:p18-talmudic-damage-premises}} & \thepage & Ordinary property damage and fatal-result consequence types are not interchangeable Original precedence 'not an operator' retained as semantic descriptor; n/a is not an invented rank. \\
\phantomsection\label{entry:symbol-registry-06198}R & \texttt{\detokenize{R}} & Reference class for the editorial example evaluation & p18 example-and-sign schemes & n/a & \hyperref[def:p18-example-extrapolation]{\ref*{def:p18-example-extrapolation}} & \thepage & Not a consequence relation; scope and sampling frame stated Original precedence descriptor 'set' preserved; n/a is not an invented rank. \\
\phantomsection\label{entry:symbol-registry-06199}F G a & \texttt{\detokenize{F G a}} & Condition outcome and observed individual in the example reconstruction & p18 example-and-sign example & n/a & \hyperref[def:p18-example-extrapolation]{\ref*{def:p18-example-extrapolation}} & \thepage & Separate from fact-set F and explanation E in the abductive source scheme Original precedence descriptor 'predicate and object' preserved; n/a is not an invented rank. \\
\phantomsection\label{entry:symbol-registry-06200}F E E' & \texttt{\detokenize{F E E'}} & Fact set proposed explanation and compared alternative in Walton's 2001 scheme & p18 example-and-sign diagnosis & n/a & \hyperref[sec:p18-sign-abductive-diagnosis]{\ref*{sec:p18-sign-abductive-diagnosis}} & \thepage & Source-specific roles; not arbitrary truth probabilities Original precedence descriptor 'metalinguistic content' preserved; n/a is not an invented rank. \\
\phantomsection\label{entry:symbol-registry-06201}A B & \texttt{\detokenize{A B}} & Observed finding and indicated proposition in the 2002 sign form & p18 example-and-sign sign entry & n/a & \hyperref[def:p18-sign-walton-reed-2002]{\ref*{def:p18-sign-walton-reed-2002}} & \thepage & Contextual presumptive warrant not universal material implication Original precedence descriptor 'propositional' preserved; n/a is not an invented rank. \\
\phantomsection\label{entry:symbol-registry-06202}H S & \texttt{\detokenize{H S}} & Hypothesis and sign events in the declared probability bridge & p18 example-and-sign probability & n/a & \hyperref[def:p18-sign-probability-model]{\ref*{def:p18-sign-probability-model}} & \thepage & Neither event's truth fixes the other's conditional probability Original precedence descriptor 'event' preserved; n/a is not an invented rank. \\
\phantomsection\label{entry:symbol-registry-06203}\ensuremath{\pi s f} & \texttt{\detokenize{\pi s f}} & Prior sensitivity and alternative-sign likelihood & p18 example-and-sign probability & n/a & \hyperref[def:p18-sign-probability-model]{\ref*{def:p18-sign-probability-model}} & \thepage & Same reference class; pi strictly between zero and one Original precedence descriptor 'scalar' preserved; n/a is not an invented rank. \\
\phantomsection\label{entry:symbol-registry-06204}\ensuremath{n p \varepsilon} & \texttt{\detokenize{n p \varepsilon}} & Trial count common success probability and population-test tolerance & p18 example-and-sign sampling & n/a & \hyperref[thm:p18-example-all-pass-bound]{\ref*{thm:p18-example-all-pass-bound}} & \thepage & Independence and equal Bernoulli probabilities explicit Original precedence descriptor 'scalar' preserved; n/a is not an invented rank. \\
\phantomsection\label{entry:symbol-registry-06205}\ensuremath{\tau} & \texttt{\detokenize{\tau}} & Posterior threshold in the declared probability bridge & p18 example-and-sign probability & n/a & \hyperref[thm:p18-sign-bayes-bridge]{\ref*{thm:p18-sign-bayes-bridge}} & \thepage & Strictly between zero and one; no universal rational acceptance threshold Original precedence descriptor 'scalar' preserved; n/a is not an invented rank. \\
\phantomsection\label{entry:symbol-registry-06206}\ensuremath{C U_X U_Y} & \texttt{\detokenize{C U_X U_Y}} & Load class and independent uniform inputs in the finite confounding countermodel & p18 example-and-sign confounding & n/a & \hyperref[def:p18-sign-confounding-process]{\ref*{def:p18-sign-confounding-process}} & \thepage & Failure equation depends on C and U-Y; installation equation separately depends on C and U-X Original precedence descriptor 'random variables' preserved; n/a is not an invented rank. \\
\phantomsection\label{entry:symbol-registry-06207}X Y & \texttt{\detokenize{X Y}} & Installation and failure events in the finite confounding countermodel & p18 example-and-sign confounding & n/a & \hyperref[def:p18-sign-confounding-process]{\ref*{def:p18-sign-confounding-process}} & \thepage & Setting X replaces only its defining equation; observational and intervention probabilities remain distinct Original precedence descriptor 'events' preserved; n/a is not an invented rank. \\
\phantomsection\label{entry:symbol-registry-06208}\ensuremath{X_i} & \texttt{\detokenize{X_i}} & Bernoulli success indicator for the ith trial & p18 example-and-sign sampling annex & n/a & \hyperref[annex:p18-example-all-pass-bound]{\ref*{annex:p18-example-all-pass-bound}} & \thepage & Distinct from installation event X in the separate confounding model Original precedence descriptor 'random variable' preserved; n/a is not an invented rank. \\
\phantomsection\label{entry:symbol-registry-06209}P & \texttt{\detokenize{P}} & Probability measure of the declared finite or Bernoulli model & p18 example-and-sign probability and sampling & n/a & \hyperref[thm:p18-example-all-pass-bound]{\ref*{thm:p18-example-all-pass-bound}} & \thepage & No probability is inferred merely from a critical-question answer or scheme label Original precedence descriptor 'function' preserved; n/a is not an invented rank. \\
\phantomsection\label{entry:symbol-registry-06382}\(D_L\) & \texttt{\detokenize{D_L}} & Declarative formula category of the fixed language & p18 erotetic implication & n/a & \hyperref[def:p18-ei-setting]{\ref*{def:p18-ei-setting}} & \thepage & Disjoint from questions; no first-order free variables in the stated wh direct answers Actual displayed TeX forms; original family descriptor retained as raw evidence. Original precedence descriptor 'set' retained; no numerical rank inferred. \\
\phantomsection\label{entry:symbol-registry-06383}\(dQ\) & \texttt{\detokenize{dQ}} & direct-answer set in the declared general fixed language L, possibly infinite (closed first-order answers in the wh instances) & P18 generalized direct-answer declaration & n/a & \hyperref[def:p18-ei-setting]{\ref*{def:p18-ei-setting}} & \thepage & At least two elements; possibly infinite; not restricted to true answers Domain/standing differs from the old finite-propositional PMCE record. The declaration-local catalog scope is chosen for existing resolver containment only; every later body/proof use is separately exposed as provenance. No extensional identity with an arbitrary old PMCE answer assignment is asserted. Actual displayed TeX forms; original family descriptor retained as raw evidence. Original precedence descriptor 'function application' retained; no numerical rank inferred. \\
\phantomsection\label{entry:symbol-registry-06384}\(P,T_P,U_P\) & \texttt{\detokenize{P,}\allowbreak\detokenize{T_P,}\allowbreak\detokenize{U_P}} & admissible-partition and component symbol family & p18 erotetic implication & n/a & \hyperref[def:p18-ei-setting]{\ref*{def:p18-ei-setting}} & \thepage & Admissible class fixed by declarative semantics; not arbitrary complex-formula assignments Family listing, not one atomic glyph: P admissible partition; T\_P designated-true declarative component; U\_P complementary untrue declarative component. Actual displayed TeX forms; original family descriptor retained as raw evidence. Original precedence descriptor 'semantic objects' retained; no numerical rank inferred. \\
\phantomsection\label{entry:symbol-registry-06385}\(\mathcal P\) & \texttt{\detokenize{\mathcal P}} & Fixed class of admissible partitions & p18 erotetic implication & n/a & \hyperref[def:p18-ei-setting]{\ref*{def:p18-ei-setting}} & \thepage & Must remain the same across composition legs Actual displayed TeX forms; original family descriptor retained as raw evidence. Original precedence descriptor 'class' retained; no numerical rank inferred. \\
\phantomsection\label{entry:symbol-registry-06386}\(\models_L\) & \texttt{\detokenize{\models_L}} & Single-conclusion semantic entailment & p18 erotetic implication & relation & \hyperref[def:p18-ei-setting]{\ref*{def:p18-ei-setting}} & \thepage & Not material implication and not question truth Actual displayed TeX forms; original family descriptor retained as raw evidence. \\
\phantomsection\label{entry:symbol-registry-06387}\(\mathrel{\|=}_L\) & \texttt{\detokenize{\mathrel{\|=}_L}} & Multiple-conclusion entailment & p18 erotetic implication & relation & \hyperref[def:p18-ei-setting]{\ref*{def:p18-ei-setting}} & \thepage & Truth of premises warrants at least one conclusion; empty conclusion entails unsatisfiable premises Actual displayed TeX forms; original family descriptor retained as raw evidence. \\
\phantomsection\label{entry:symbol-registry-06388}\(\vdash_\ell,\operatorname{Cn}_\ell\) & \texttt{\detokenize{\vdash_\ell,}\allowbreak\detokenize{\operatorname{Cn}_\ell}} & declarative-consequence relation/operator symbol family & p18 erotetic implication & n/a & \hyperref[def:p18-ei-consequence-setting]{\ref*{def:p18-ei-consequence-setting}} & \thepage & Overlap monotonicity cut for sets and Lb are explicit assumptions Family listing, not one atomic glyph: \(\vdash_\ell\) declarative consequence relation; \(\operatorname{Cn}_\ell\) declarative consequence-set operator. Actual displayed TeX forms; original family descriptor retained as raw evidence. Original precedence descriptor 'relation and function' retained; no numerical rank inferred. \\
\phantomsection\label{entry:symbol-registry-06389}\(\operatorname{Im}_L\) & \texttt{\detokenize{\operatorname{Im}_L}} & Two-clause question-to-question implication on a declarative basis & p18 erotetic implication & n/a & \hyperref[def:p18-ei-relation]{\ref*{def:p18-ei-relation}} & \thepage & Both soundness transmission and nonempty proper answer reduction Actual displayed TeX forms; original family descriptor retained as raw evidence. Original precedence descriptor 'ternary relation' retained; no numerical rank inferred. \\
\phantomsection\label{entry:symbol-registry-06390}\(E_L\) & \texttt{\detokenize{E_L}} & General semantic question evocation & p18 erotetic implication comparison & n/a & \hyperref[def:p18-ei-relation]{\ref*{def:p18-ei-relation}} & \thepage & Declarative-only premises; coverage and individual-answer non-entailment; not question-premise implication Actual displayed TeX forms; original family descriptor retained as raw evidence. Original precedence descriptor 'binary relation' retained; no numerical rank inferred. \\
\phantomsection\label{entry:symbol-registry-06391}\(\operatorname{RIm}_L\) & \texttt{\detokenize{\operatorname{RIm}_L}} & Editorial abbreviation for source regular erotetic implication & p18 erotetic implication & n/a & \hyperref[def:p18-ei-relation]{\ref*{def:p18-ei-relation}} & \thepage & Each target answer entails a single source answer on the basis; distinct from regularity of a question Actual displayed TeX forms; original family descriptor retained as raw evidence. Original precedence descriptor 'ternary relation' retained; no numerical rank inferred. \\
\phantomsection\label{entry:symbol-registry-06392}\(Y_B\) & \texttt{\detokenize{Y_B}} & first-leg nonempty proper answer-reduction witness family indexed by B & p18 erotetic implication & n/a & \hyperref[def:p18-ei-relation]{\ref*{def:p18-ei-relation}} & \thepage & Composition unions need a separate properness condition Split from the original Y\_B/V\_C family: Y\_B is actually declared here; V\_C is separately bound to the composition theorem. Actual displayed TeX forms; original family descriptor retained as raw evidence. Original precedence descriptor 'sets' retained; no numerical rank inferred. \\
\phantomsection\label{entry:symbol-registry-06393}\(\operatorname{Pres}(Q)\) & \texttt{\detokenize{\operatorname{Pres}(Q)}} & Presupposition set of a question & p18 erotetic implication & n/a & \hyperref[def:p18-ei-presupposition]{\ref*{def:p18-ei-presupposition}} & \thepage & Prospective presupposition normality and regularity separately defined Actual displayed TeX forms; original family descriptor retained as raw evidence. Original precedence descriptor 'set' retained; no numerical rank inferred. \\
\phantomsection\label{entry:symbol-registry-06394}\(F_C,Q^C\) & \texttt{\detokenize{F_C,}\allowbreak\detokenize{Q^C}} & conditional-prefix map and transformed-question symbol family & p18 erotetic implication & n/a & \hyperref[def:p18-ei-conditional-language]{\ref*{def:p18-ei-conditional-language}} & \thepage & Truth clause plus actual question closure; injectivity is syntactic Family listing, not one atomic glyph: F\_C conditional prefix map on declarative answer sets; Q\^{}C corresponding conditional-prefixed question. Actual displayed TeX forms; original family descriptor retained as raw evidence. Original precedence descriptor 'function and question' retained; no numerical rank inferred. \\
\phantomsection\label{entry:symbol-registry-06395}\(G_X,Q_X,R_X\) & \texttt{\detokenize{G_X,}\allowbreak\detokenize{Q_X,}\allowbreak\detokenize{R_X}} & conjunctive-prefix map and transformed-question symbol family & p18 erotetic implication & n/a & \hyperref[def:p18-ei-conjunct-language]{\ref*{def:p18-ei-conjunct-language}} & \thepage & Existence not inferred from semantic associativity or restricted one-step closure Family listing, not one atomic glyph: G\_X finite conjunctive prefix map on declaratives and their sets; Q\_X first exact conjunctively transformed question; R\_X second exact conjunctively transformed question. Actual displayed TeX forms; original family descriptor retained as raw evidence. Original precedence descriptor 'function and questions' retained; no numerical rank inferred. \\
\phantomsection\label{entry:symbol-registry-06396}\(?S(Ax),?U(Ax)\) & \texttt{\detokenize{?S(Ax),}\allowbreak\detokenize{?U(Ax)}} & source wh-question-form family & p18 erotetic implication & n/a & \hyperref[def:p18-ei-wh-language]{\ref*{def:p18-ei-wh-language}} & \thepage & Nonempty first-order domains; no unique names; empty infinite and unnamed-extension soundness limits retained Family listing, not one atomic glyph: ?S(Ax) named-witness wh-question form; ?U(Ax) nonempty finite complete-list wh-question form. Actual displayed TeX forms; original family descriptor retained as raw evidence. Original precedence descriptor 'questions' retained; no numerical rank inferred. \\
\phantomsection\label{entry:symbol-registry-06397}\(V_C\) & \texttt{\detokenize{V_C}} & second-leg nonempty proper answer-reduction witness family indexed by C & p18 erotetic implication & n/a & \hyperref[thm:p18-ei-composition]{\ref*{thm:p18-ei-composition}} & \thepage & V\_C is actually declared in the composition theorem as the second-leg witness subset of dR. Split from original raw Y\_B/V\_C record at line12; no fake first-use at the earlier relation definition. \\
\phantomsection\label{entry:symbol-registry-07390}\(\Phi\) & \texttt{\detokenize{\Phi}} & Finite derivation-path family & p18 erotetic search & n/a & \hyperref[def:p18-es-family]{\ref*{def:p18-es-family}} & \thepage & Fairness differs from truth and endpoint completeness Original precedence descriptor 'family' retained; no rank inferred. \\
\phantomsection\label{entry:symbol-registry-07391}\(e\) & \texttt{\detokenize{e}} & Sequence and indexed occurrence symbol family member \(e\) & p18 erotetic search & n/a & \hyperref[def:p18-es-derivation]{\ref*{def:p18-es-derivation}} & \thepage & Earlier need not mean immediately earlier Original precedence descriptor 'sequence and expression' retained; no rank inferred. Family listing, not one atomic glyph. Author macro \texttt{\detokenize{e \varphi\_k}} not carried to the members. \\
\phantomsection\label{entry:symbol-registry-07392}\(\varphi_k\) & \texttt{\detokenize{\varphi_k}} & Sequence and indexed occurrence symbol family member \(\varphi_k\) & p18 erotetic search & n/a & \hyperref[def:p18-es-derivation]{\ref*{def:p18-es-derivation}} & \thepage & Earlier need not mean immediately earlier Original precedence descriptor 'sequence and expression' retained; no rank inferred. Family listing, not one atomic glyph. Author macro \texttt{\detokenize{e \varphi\_k}} not carried to the members. \\
\phantomsection\label{entry:symbol-registry-07393}\(Q\) & \texttt{\detokenize{Q}} & Principal question and initial basis symbol family member \(Q\) & p18 erotetic search & n/a & \hyperref[def:p18-es-derivation]{\ref*{def:p18-es-derivation}} & \thepage & Direct answers disjoint from X; finite used initial set Original precedence descriptor 'question and set' retained; no rank inferred. Family listing, not one atomic glyph. Author macro \texttt{\detokenize{Q X}} not carried to the members. \\
\phantomsection\label{entry:symbol-registry-07394}\(X\) & \texttt{\detokenize{X}} & Principal question and initial basis symbol family member \(X\) & p18 erotetic search & n/a & \hyperref[def:p18-es-derivation]{\ref*{def:p18-es-derivation}} & \thepage & Direct answers disjoint from X; finite used initial set Original precedence descriptor 'question and set' retained; no rank inferred. Family listing, not one atomic glyph. Author macro \texttt{\detokenize{Q X}} not carried to the members. \\
\phantomsection\label{entry:symbol-registry-07395}\(?A\) & \texttt{\detokenize{?A}} & Polar question abbreviation & p18 erotetic search & n/a & \hyperref[sec:p18-erotetic-search-scenarios]{\ref*{sec:p18-erotetic-search-scenarios}} & \thepage & Direct answers A and not-A Original precedence descriptor 'question' retained; no rank inferred. \\
\phantomsection\label{entry:symbol-registry-07396}\(k\) & \texttt{\detokenize{k}} & Occurrence position & p18 erotetic search & n/a & \hyperref[def:p18-es-derivation]{\ref*{def:p18-es-derivation}} & \thepage & Query excludes k=1; literal fairness includes it Original precedence descriptor 'integer' retained; no rank inferred. \\
\phantomsection\label{entry:symbol-registry-07397}\(h_\Phi\) & \texttt{\detokenize{h_\Phi}} & Maximum path length & p18 erotetic search & n/a & \hyperref[annex:p18.es.golden-path]{\ref*{annex:p18.es.golden-path}} & \thepage & Successor bound not time bound Original precedence descriptor 'integer' retained; no rank inferred. \\
\phantomsection\label{entry:symbol-registry-07398}\(X_0\) & \texttt{\detokenize{X_0}} & Finite used initial-premise set & p18 erotetic search & n/a & \hyperref[annex:p18.es.canonical-concise]{\ref*{annex:p18.es.canonical-concise}} & \thepage & Move recipe requires query-answer/X disjointness Original precedence descriptor 'set' retained; no rank inferred. \\
\phantomsection\label{entry:symbol-registry-07399}\(Z\) & \texttt{\detokenize{Z}} & Retained grounds and replaced basis symbol family member \(Z\) & p18 erotetic search & n/a & \hyperref[thm:p18-es-basis-replacement]{\ref*{thm:p18-es-basis-replacement}} & \thepage & Z entails each Y-member Original precedence descriptor 'sets' retained; no rank inferred. Family listing, not one atomic glyph. Author macro \texttt{\detokenize{Z Y}} not carried to the members. \\
\phantomsection\label{entry:symbol-registry-07400}\(Y\) & \texttt{\detokenize{Y}} & Retained grounds and replaced basis symbol family member \(Y\) & p18 erotetic search & n/a & \hyperref[thm:p18-es-basis-replacement]{\ref*{thm:p18-es-basis-replacement}} & \thepage & Z entails each Y-member Original precedence descriptor 'sets' retained; no rank inferred. Family listing, not one atomic glyph. Author macro \texttt{\detokenize{Z Y}} not carried to the members. \\
\phantomsection\label{entry:symbol-registry-07401}\(H\) & \texttt{\detokenize{H}} & Universal predicate-link premise & p18 erotetic search & n/a & \hyperref[def:p18-es-predicate]{\ref*{def:p18-es-predicate}} & \thepage & Not tree height Original precedence descriptor 'declarative' retained; no rank inferred. \\
\phantomsection\label{entry:symbol-registry-07402}\(x_1\) & \texttt{\detokenize{x_1}} & Departure premise formula names symbol family member \(x_1\) & p18 erotetic search & n/a & \hyperref[def:p18-es-destination]{\ref*{def:p18-es-destination}} & \thepage & Names not extra equalities Original precedence descriptor 'declaratives' retained; no rank inferred. Family listing, not one atomic glyph. Author macro \texttt{\detokenize{x\_1 x\_2 x\_3}} not carried to the members. \\
\phantomsection\label{entry:symbol-registry-07403}\(x_2\) & \texttt{\detokenize{x_2}} & Departure premise formula names symbol family member \(x_2\) & p18 erotetic search & n/a & \hyperref[def:p18-es-destination]{\ref*{def:p18-es-destination}} & \thepage & Names not extra equalities Original precedence descriptor 'declaratives' retained; no rank inferred. Family listing, not one atomic glyph. Author macro \texttt{\detokenize{x\_1 x\_2 x\_3}} not carried to the members. \\
\phantomsection\label{entry:symbol-registry-07404}\(x_3\) & \texttt{\detokenize{x_3}} & Departure premise formula names symbol family member \(x_3\) & p18 erotetic search & n/a & \hyperref[def:p18-es-destination]{\ref*{def:p18-es-destination}} & \thepage & Names not extra equalities Original precedence descriptor 'declaratives' retained; no rank inferred. Family listing, not one atomic glyph. Author macro \texttt{\detokenize{x\_1 x\_2 x\_3}} not carried to the members. \\
\phantomsection\label{entry:symbol-registry-07405}\(\mapsto\) & \texttt{\detokenize{\mapsto}} & Response-action/endpoint edge & p18 erotetic search sheet & n/a & \hyperref[rule:p18-es-meeting]{\ref*{rule:p18-es-meeting}} & \thepage & Not material or erotetic implication Original precedence descriptor 'mapping' retained; no rank inferred. \\
\phantomsection\label{entry:symbol-registry-07406}\(\supset\) & \texttt{\detokenize{\supset}} & Source-native classical material conditional & P18 native search procedures & meta & \hyperref[def:p18-es-native-notation]{\ref*{def:p18-es-native-notation}} & \thepage & Wisniewski2001QuestionsInferences Figures6 and10; object-language connective, not a semantic or response-action arrow. \\
\phantomsection\label{entry:symbol-registry-07407}\(\equiv\) & \texttt{\detokenize{\equiv}} & Source-native classical biconditional & P18 native search procedures & meta & \hyperref[def:p18-es-native-notation]{\ref*{def:p18-es-native-notation}} & \thepage & Wisniewski2001QuestionsInferences Figures6,7 and12; not elementary equivalence of structures. \\
\phantomsection\label{entry:symbol-registry-07408}\(?\pm|A,B|\) & \texttt{\detokenize{?\pm|A,}\allowbreak\detokenize{B|}} & Source-native binary conjunctive question expression & P18 native search procedures & meta & \hyperref[def:p18-es-native-notation]{\ref*{def:p18-es-native-notation}} & \thepage & Four syntactically assigned conjunctive direct answers; preserves the source question marker and bars instead of a host B alias. Blank macro filled with the glyph's TeX spelling. \\
\phantomsection\label{entry:symbol-registry-07409}\(\otimes\) & \texttt{\detokenize{\otimes}} & Source-native Boolean connective-schema parameter & P18 native search procedures & meta & \hyperref[def:p18-es-native-notation]{\ref*{def:p18-es-native-notation}} & \thepage & Finite range is conjunction, disjunction, material conditional and biconditional; not linear tensor. Full source and calculus standing remain separate. \\
\phantomsection\label{entry:symbol-registry-08006}\(L_1\) & \texttt{\detokenize{L_1}} & Original classical propositional question language & p18 original erotetic patterns & n/a & \hyperref[def:p18-ep-notation]{\ref*{def:p18-ep-notation}} & \thepage & Wisniewski2001QuestionsInferences section6.1; all Boolean valuations Original precedence descriptor 'language' retained; no rank inferred. \\
\phantomsection\label{entry:symbol-registry-08007}\(L_2\) & \texttt{\detokenize{L_2}} & Original monadic equality question language & p18 original erotetic patterns & n/a & \hyperref[def:p18-ep-notation]{\ref*{def:p18-ep-notation}} & \thepage & Section6.2; constants and no function symbols; all nonempty-domain models Original precedence descriptor 'language' retained; no rank inferred. \\
\phantomsection\label{entry:symbol-registry-08008}\(\supset\) & \texttt{\detokenize{\supset}} & Source-native material conditional & p18 original erotetic patterns & n/a & \hyperref[def:p18-ep-notation]{\ref*{def:p18-ep-notation}} & \thepage & Not set inclusion; preserve original printed19-29 signs Original precedence descriptor 'Boolean binary connective' retained; no rank inferred. \\
\phantomsection\label{entry:symbol-registry-08009}\(\equiv\) & \texttt{\detokenize{\equiv}} & Source-native Boolean biconditional & p18 original erotetic patterns & n/a & \hyperref[def:p18-ep-notation]{\ref*{def:p18-ep-notation}} & \thepage & Local role distinct from global elementary equivalence; declared before source schemata Original precedence descriptor 'Boolean binary connective' retained; no rank inferred. \\
\phantomsection\label{entry:symbol-registry-08010}\(?\pm|A,B|\) & \texttt{\detokenize{?\pm|A,}\allowbreak\detokenize{B|}} & Source-native binary conjunctive question & p18 original erotetic patterns & n/a & \hyperref[def:p18-ep-notation]{\ref*{def:p18-ep-notation}} & \thepage & Four explicit conjunction answers; pm marker not arithmetic addition; source(6.4) Original precedence descriptor 'question' retained; no rank inferred. \\
\phantomsection\label{entry:symbol-registry-08011}\(\otimes\) & \texttt{\detokenize{\otimes}} & Metalinguistic connective parameter & p18 original erotetic patterns & n/a & \hyperref[def:p18-ep-notation]{\ref*{def:p18-ep-notation}} & \thepage & Ranges over land lor supset equiv; not linear-logic tensor; source(8.6)-(8.13) Original precedence descriptor 'parameter' retained; no rank inferred. \\
\phantomsection\label{entry:symbol-registry-08012}\(\alpha A\) & \texttt{\detokenize{\alpha A}} & Fixed positive-or-negative A-premise & p18 original erotetic patterns & n/a & \hyperref[def:p18-ep-notation]{\ref*{def:p18-ep-notation}} & \thepage & Source(8.12); A or neg A within each instance; not an alpha formula class Original precedence descriptor 'formula choice' retained; no rank inferred. \\
\phantomsection\label{entry:symbol-registry-08013}\(\alpha B\) & \texttt{\detokenize{\alpha B}} & Fixed positive-or-negative B-premise & p18 original erotetic patterns & n/a & \hyperref[def:p18-ep-notation]{\ref*{def:p18-ep-notation}} & \thepage & Source(8.13); B or neg B within each instance Original precedence descriptor 'formula choice' retained; no rank inferred. \\
\phantomsection\label{entry:symbol-registry-08014}\(A(x/c)\) & \texttt{\detokenize{A(x/}\allowbreak\detokenize{c)}} & Closed constant-substitution instance & p18 original erotetic patterns & n/a & \hyperref[def:p18-ep-notation]{\ref*{def:p18-ep-notation}} & \thepage & Slash is native substitution separator not division; x sole free variable Original precedence descriptor 'sentence' retained; no rank inferred. \\
\phantomsection\label{entry:symbol-registry-08015}\(S(Ax)\) & \texttt{\detokenize{S(Ax)}} & Source-native existential wh-answer set & p18 original erotetic patterns & n/a & \hyperref[def:p18-ep-notation]{\ref*{def:p18-ep-notation}} & \thepage & Metalanguage role; object-language question has ?S(Ax); extends existing actual wh definition Original precedence descriptor 'set of sentences' retained; no rank inferred. \\
\phantomsection\label{entry:symbol-registry-08016}\(U(Ax)\) & \texttt{\detokenize{U(Ax)}} & Source-native general wh-answer set & p18 original erotetic patterns & n/a & \hyperref[def:p18-ep-notation]{\ref*{def:p18-ep-notation}} & \thepage & Nonempty finite positive list plus exhaustion clause; no empty-list response Original precedence descriptor 'set of sentences' retained; no rank inferred. \\
\phantomsection\label{entry:symbol-registry-08017}\(\vdash_E\) & \texttt{\detokenize{\vdash_E}} & Abbreviation for semantic evocation & p18 original erotetic patterns & relation & \hyperref[def:p18-ep-notation]{\ref*{def:p18-ep-notation}} & \thepage & Not PMCE syntactic provability; source printed20 \\
\phantomsection\label{entry:symbol-registry-08018}\(\vdash_{\operatorname{Im}}\) & \texttt{\detokenize{\vdash_{\operatorname{Im}}}} & Abbreviation for semantic erotetic implication & p18 original erotetic patterns & relation & \hyperref[def:p18-ep-notation]{\ref*{def:p18-ep-notation}} & \thepage & Question and optional declarative basis; source printed27-29 \\
\phantomsection\label{entry:symbol-registry-08019}\(J_n\) & \texttt{\detokenize{J_n}} & Finite answer-index set & p18 original erotetic patterns & n/a & \hyperref[def:p18-ep-families]{\ref*{def:p18-ep-families}} & \thepage & Editorial proof abbreviation for 1 through n; not source extra hypothesis Original precedence descriptor 'set of indices' retained; no rank inferred. \\
\phantomsection\label{entry:symbol-registry-08020}\(D_n\) & \texttt{\detokenize{D_n}} & Full disjunction of finite answer list & p18 original erotetic patterns & n/a & \hyperref[def:p18-ep-families]{\ref*{def:p18-ep-families}} & \thepage & Editorial abbreviation expanded literally in source rule sheet Original precedence descriptor 'formula' retained; no rank inferred. \\
\phantomsection\label{entry:symbol-registry-08021}\(K_n\) & \texttt{\detokenize{K_n}} & Full conjunction of finite answer list & p18 original erotetic patterns & n/a & \hyperref[def:p18-ep-families]{\ref*{def:p18-ep-families}} & \thepage & Editorial abbreviation; fixed bracketing; source(8.16) Original precedence descriptor 'formula' retained; no rank inferred. \\
\phantomsection\label{entry:symbol-registry-08022}\(D^-_i\) & \texttt{\detokenize{D^-_i}} & Nonempty initial answer-block disjunction & p18 original erotetic patterns & n/a & \hyperref[def:p18-ep-families]{\ref*{def:p18-ep-families}} & \thepage & 1 <= i < n; source rule premises fully expanded Original precedence descriptor 'formula' retained; no rank inferred. \\
\phantomsection\label{entry:symbol-registry-08023}\(D^+_i\) & \texttt{\detokenize{D^+_i}} & Nonempty final answer-block disjunction & p18 original erotetic patterns & n/a & \hyperref[def:p18-ep-families]{\ref*{def:p18-ep-families}} & \thepage & Source(8.18)-(8.21); proper block because i < n Original precedence descriptor 'formula' retained; no rank inferred. \\
\phantomsection\label{entry:symbol-registry-08024}\(N_{\vec c}(x)\) & \texttt{\detokenize{N_{\vec c}(x)}} & Finite naming disjunction & p18 original erotetic patterns & n/a & \hyperref[def:p18-ep-families]{\ref*{def:p18-ep-families}} & \thepage & x=c1 or ... or x=ck; no unique-names assumption; sheet expands names Original precedence descriptor 'open formula' retained; no rank inferred. \\
\phantomsection\label{entry:symbol-registry-08025}\(H_{\vec c}\) & \texttt{\detokenize{H_{\vec c}}} & Exact named P-extension sentence & p18 original erotetic patterns & n/a & \hyperref[def:p18-ep-families]{\ref*{def:p18-ep-families}} & \thepage & Universal biconditional; each complete-list disjunct is expanded in the sheet Original precedence descriptor 'sentence' retained; no rank inferred. \\
\phantomsection\label{entry:symbol-registry-08026}\(W_{\vec c}\) & \texttt{\detokenize{W_{\vec c}}} & Positive finite complete-list answer & p18 original erotetic patterns & n/a & \hyperref[def:p18-ep-families]{\ref*{def:p18-ep-families}} & \thepage & Every positive instance and universal exhaustion retained Original precedence descriptor 'sentence' retained; no rank inferred. \\
\phantomsection\label{entry:symbol-registry-08027}\((\vec c,\vec t)\) & \texttt{\detokenize{(\vec c,}\allowbreak\detokenize{\vec t)}} & Displayed and candidate finite constant-name lists & p18 original erotetic patterns & n/a & \hyperref[def:p18-ep-families]{\ref*{def:p18-ep-families}} & \thepage & Grouped declared family, not atomic-glyph count; candidate lists may contain outside names Original precedence descriptor 'list family' retained; no rank inferred. \\
\phantomsection\label{entry:symbol-registry-08028}\(\mathcal H(Q)\) & \texttt{\detokenize{\mathcal H(Q)}} & Source-answer subformulas and their negations & p18 original erotetic patterns & n/a & \hyperref[def:p18-ep-analytic]{\ref*{def:p18-ep-analytic}} & \thepage & Editorial set abbreviation for exact source Definition11 analyticity; no semantic-equivalence closure Original precedence descriptor 'syntactic set' retained; no rank inferred. \\
\phantomsection\label{entry:symbol-registry-08029}\((C_j,T)\) & \texttt{\detokenize{(C_j,}\allowbreak\detokenize{T)}} & Alternative and candidate denotation sets in complete-list countermodel & p18 original erotetic patterns complete-list proof & n/a & \hyperref[annex:p18.ep.complete-list]{\ref*{annex:p18.ep.complete-list}} & \thepage & Grouped family; T is not an admissible-partition true component; both sets explicitly declared Original precedence descriptor 'set family' retained; no rank inferred. \\
\phantomsection\label{entry:symbol-registry-08030}\((11,10,01,00)\) & \texttt{\detokenize{(11,}\allowbreak\detokenize{10,}\allowbreak\detokenize{01,}\allowbreak\detokenize{00)}} & Four binary-answer truth-cell labels & p18 original erotetic patterns pure proof & n/a & \hyperref[annex:p18.ep.pure-patterns]{\ref*{annex:p18.ep.pure-patterns}} & \thepage & Grouped labels for actual conjunction formulas; not added answers or probability values Original precedence descriptor 'finite label family' retained; no rank inferred. \\
\phantomsection\label{entry:symbol-registry-08031}\(P\) & \texttt{\detokenize{P}} & Source object-language unary predicate symbol & p18 original erotetic patterns & n/a & \hyperref[def:p18-ep-notation]{\ref*{def:p18-ep-notation}} & \thepage & Source section6 evocation examples; not the question metavariable of an Im judgment Original precedence descriptor 'predicate' retained; no rank inferred. \\
\phantomsection\label{entry:symbol-registry-08032}\(R\) & \texttt{\detokenize{R}} & Source object-language unary predicate symbol & p18 original erotetic patterns & n/a & \hyperref[def:p18-ep-notation]{\ref*{def:p18-ep-notation}} & \thepage & Distinct from P; original section6; local type distinct from earlier further-question R Original precedence descriptor 'predicate' retained; no rank inferred. \\
\phantomsection\label{entry:symbol-registry-08033}\(P\) & \texttt{\detokenize{P}} & Extension of P in a concrete model description & p18 original erotetic patterns model descriptions & n/a & \hyperref[sec:p18-erotetic-patterns-proofs]{\ref*{sec:p18-erotetic-patterns-proofs}} & \thepage & Honest local predicate/extension role reuse; assignment P=V declares interpretation, not object-language equality Original precedence descriptor 'domain subset' retained; no rank inferred. \\
\phantomsection\label{entry:symbol-registry-08034}\(R\) & \texttt{\detokenize{R}} & Extension of R in a concrete model description & p18 original erotetic patterns model descriptions & n/a & \hyperref[sec:p18-erotetic-patterns-proofs]{\ref*{sec:p18-erotetic-patterns-proofs}} & \thepage & Same local assignment convention; not question R or object-language predicate occurrence Original precedence descriptor 'domain subset' retained; no rank inferred. \\
\end{longtable}
\endgroup
\Needspace{.4\textheight}
\subsection*{Stable part P19 scopes}
\begingroup\scriptsize\setlength{\tabcolsep}{3pt}
\begin{longtable}{@{}p{0.9in}p{1.3in}p{1.0in}p{0.5in}p{0.3in}p{0.4in}p{0.35in}p{1.65in}@{}}
\toprule
\textbf{Glyph} & \textbf{Macro} & \textbf{Role} & \textbf{Scope} & \textbf{Prec./form} & \textbf{First} & \textbf{Reg. page} & \textbf{Principal uses / notes} \\
\midrule
\endfirsthead
\toprule
\textbf{Glyph} & \textbf{Macro} & \textbf{Role} & \textbf{Scope} & \textbf{Prec./form} & \textbf{First} & \textbf{Reg. page} & \textbf{Principal uses / notes} \\
\midrule
\endhead
\midrule\multicolumn{8}{r}{\emph{continued}}\\\endfoot
\bottomrule\endlastfoot
\phantomsection\label{entry:symbol-registry-01009}\(A,B,A^c,\mathbb P\) & \texttt{\detokenize{A,}\allowbreak\detokenize{B,}\allowbreak\detokenize{A^c,}\allowbreak\detokenize{\mathbb P}} & probability events complement and measure in an empirical benchmark & P19 evidence-methods & meta & \hyperref[sec:p19-evidence-methods]{\ref*{sec:p19-evidence-methods}} & \thepage & Reuses the linked finite probability model. Complement and subset comparisons require the same outcome space, information and event interpretation; the notation supplies no empirical sampling design. \\
\phantomsection\label{entry:symbol-registry-01010}\(g,n_g,k_g\) & \texttt{\detokenize{g,}\allowbreak\detokenize{n_g,}\allowbreak\detokenize{k_g}} & response group denominator and coded-response count & P19 evidence-methods & meta & \hyperref[sec:p19-evidence-methods]{\ref*{sec:p19-evidence-methods}} & \thepage & The denominator belongs to the particular response estimate and is positive when its proportion is defined. Group labels do not imply random assignment. \\
\phantomsection\label{entry:symbol-registry-01011}\(\widehat p_g,\widehat\Delta\) & \texttt{\detokenize{\widehat p_g,}\allowbreak\detokenize{\widehat\Delta}} & observed response proportion and group difference & P19 evidence-methods & meta & \hyperref[sec:p19-evidence-methods]{\ref*{sec:p19-evidence-methods}} & \thepage & The proportion is \(k_g/n_g\) and the difference subtracts the specified group proportions. Rounded source percentages support approximate arithmetic, not invented exact success counts. \\
\phantomsection\label{entry:symbol-registry-01012}\(p\text{ in a }p\text{-value}\) & \texttt{\detokenize{p}} & reported statistical tail probability & P19 evidence-methods & meta & \hyperref[sec:p19-evidence-methods]{\ref*{sec:p19-evidence-methods}} & \thepage & Defined under a specified statistical model and extremeness rule. It is neither an observed response proportion nor a posterior probability that a hypothesis is true. \\
\phantomsection\label{entry:symbol-registry-01013}\(m,n_I,n_C,n_O\) & \texttt{\detokenize{m,}\allowbreak\detokenize{n_I,}\allowbreak\detokenize{n_C,}\allowbreak\detokenize{n_O}} & completed-item count and three response-category counts & P19 evidence-methods & meta & \hyperref[sec:p19-evidence-methods]{\ref*{sec:p19-evidence-methods}} & \thepage & The positive completed-item count is partitioned into designated intuitive errors, correct answers and other errors. The categories are response codes rather than directly observed cognitive stages. \\
\phantomsection\label{entry:symbol-registry-01014}\(\widehat a,\widehat c\) & \texttt{\detokenize{\widehat a,}\allowbreak\detokenize{\widehat c}} & non-designated-response and conditional-correctness proportions & P19 evidence-methods & meta & \hyperref[sec:p19-evidence-methods]{\ref*{sec:p19-evidence-methods}} & \thepage & The conditional score is defined only when \(n_C+n_O\) is positive. A source imputation for the zero-denominator case is a scoring convention, not the arithmetic value of zero divided by zero. \\
\phantomsection\label{entry:symbol-registry-01015}\(a,b,\overline Y_{ab}\) & \texttt{\detokenize{a,}\allowbreak\detokenize{b,}\allowbreak\detokenize{\overline Y_{ab}}} & hypothetical assignment-condition labels and cell means & P19 evidence-methods & meta & \hyperref[sec:p19-evidence-methods]{\ref*{sec:p19-evidence-methods}} & \thepage & The binary labels identify checking and aid conditions in the proposed design. Their role differs from the earlier hatted participant scores. A cell mean needs observed outcomes under its declared rule. \\
\phantomsection\label{entry:symbol-registry-01016}\(\widehat\Delta_{\rm check}\) & \texttt{\detokenize{\widehat\Delta_{\rm check}}} & secondary equal-weight descriptive contrast across aid conditions & P19 evidence-methods & meta & \hyperref[sec:p19-evidence-methods]{\ref*{sec:p19-evidence-methods}} & \thepage & The displayed average of two checking differences requires all four cell means. It answers the secondary question averaged over aid conditions and can have the opposite sign from the primary aid-present comparison. \\
\phantomsection\label{entry:symbol-registry-01017}\(\mathrm{CRT},\mathrm{BIC}\) & \texttt{\detokenize{CRT,}\allowbreak\detokenize{BIC}} & introduced empirical-method acronyms & P19 evidence-methods & meta & \hyperref[sec:p19-evidence-methods]{\ref*{sec:p19-evidence-methods}} & \thepage & Prose abbreviations for Cognitive Reflection Test and Bayesian information criterion. They are not logical constants; an instrument name or favorable model comparison does not establish a unique mental mechanism. \\
\phantomsection\label{entry:symbol-registry-01018}\(i,j,X_{ij}\) & \texttt{\detokenize{i,}\allowbreak\detokenize{j,}\allowbreak\detokenize{X_{ij}}} & participant item and classified response in the Study 1 model & P19 evidence-methods & meta & \hyperref[sec:p19-evidence-methods]{\ref*{sec:p19-evidence-methods}} & \thepage & The item index runs from one to five. \(X_{ij}\) records a response category rather than an observed psychological process. \\
\phantomsection\label{entry:symbol-registry-01019}\(\mathrm I,\mathrm C,\mathrm O\) & \texttt{\detokenize{\mathrm I,}\allowbreak\detokenize{\mathrm C,}\allowbreak\detokenize{\mathrm O}} & three classified item-response categories & P19 evidence-methods & meta & \hyperref[sec:p19-evidence-methods]{\ref*{sec:p19-evidence-methods}} & \thepage & Designated intuitive error, correct answer and other error. The first category leaves the model's second step missing rather than recording another observed failure. \\
\phantomsection\label{entry:symbol-registry-01020}\(r_i,c_i\) & \texttt{\detokenize{r_i,}\allowbreak\detokenize{c_i}} & person effects for the two response-model stages & P19 evidence-methods & meta & \hyperref[sec:p19-evidence-methods]{\ref*{sec:p19-evidence-methods}} & \thepage & Real log-odds deviations for reflection and calculation conditional on reaching its step. Model 2 allows separate variation and correlation; these effects are not raw response proportions. \\
\phantomsection\label{entry:symbol-registry-01021}\(\alpha_j,\beta_j,L(s)\) & \texttt{\detokenize{\alpha_j,}\allowbreak\detokenize{\beta_j,}\allowbreak\detokenize{L}} & item intercepts and the logistic function & P19 evidence-methods & meta & \hyperref[sec:p19-evidence-methods]{\ref*{sec:p19-evidence-methods}} & \thepage & L(s)=1/(1+exp(-s)) for real s. The five item intercepts at each of the two stages give ten item effects; a larger intercept increases the corresponding stage-success probability. \\
\phantomsection\label{entry:symbol-registry-01022}\(\pi^R_{ij},\pi^C_{ij},\Pr\) & \texttt{\detokenize{\pi^R_{ij},}\allowbreak\detokenize{\pi^C_{ij},}\allowbreak\detokenize{\Pr}} & conditional stage probabilities and response-model probability & P19 evidence-methods & meta & \hyperref[sec:p19-evidence-methods]{\ref*{sec:p19-evidence-methods}} & \thepage & The R and C superscripts name reflection and calculation stages. The three response probabilities are 1-piR, piR*piC and piR*(1-piC), conditional on person effects; calculation conditions on reaching its step. \\
\phantomsection\label{entry:symbol-registry-01023}\(\widehat r_i,\widehat c_i^{\rm mod}\) & \texttt{\detokenize{\widehat r_i,}\allowbreak\detokenize{\widehat c_i^{\rm mod}}} & conditional modal person scores used as Study 1 predictors & P19 evidence-methods & meta & \hyperref[sec:p19-evidence-methods]{\ref*{sec:p19-evidence-methods}} & \thepage & Source-described fitted log-odds deviations, distinct from proportion scores and from probabilities. The source does not specify joint versus marginal modal extraction; observed score SDs are distinct from model person-effect scales. \\
\phantomsection\label{entry:symbol-registry-01024}\(\widehat\Delta_{\rm aid}\) & \texttt{\detokenize{\widehat\Delta_{\rm aid}}} & primary descriptive checking contrast with numerical aid present & P19 evidence-methods & meta & \hyperref[sec:p19-evidence-methods]{\ref*{sec:p19-evidence-methods}} & \thepage & The difference \(\overline Y_{11}-\overline Y_{01}\) requires its two constituent means. It answers the primary aid-present question; a negative secondary average alone does not establish that this contrast is negative. \\
\phantomsection\label{entry:symbol-registry-01473}\(\mathrm{MP},\mathrm{MT},\mathrm{AC},\mathrm{DA},p,q,r,e,s,o,\Gamma,\Gamma',J\) & \texttt{\detokenize{\mathrm{MP},}\allowbreak\detokenize{\mathrm{MT},}\allowbreak\detokenize{\mathrm{AC},}\allowbreak\detokenize{\mathrm{DA},}\allowbreak\detokenize{p,}\allowbreak\detokenize{q,}\allowbreak\detokenize{r,}\allowbreak\detokenize{e,}\allowbreak\detokenize{s,}\allowbreak\detokenize{o,}\allowbreak\detokenize{\Gamma,}\allowbreak\detokenize{\Gamma',}\allowbreak\detokenize{J}} & conditional inference patterns task atoms literal and reconstructed premise sets and hypothetical valuation state & P19 conditional-reasoning-suppression & meta & \hyperref[sec:p19-conditional-reasoning-suppression]{\ref*{sec:p19-conditional-reasoning-suppression}} & \thepage & The benchmark fixes meanings and valuations. Reconstructed contexts are illustrations, not observed participant mental states. \\
\phantomsection\label{entry:symbol-registry-01474}\(N,F,\mathrm{MSe},p,\alpha\) & \texttt{\detokenize{N,}\allowbreak\detokenize{F,}\allowbreak\detokenize{\mathrm{MSe},}\allowbreak\detokenize{p,}\allowbreak\detokenize{\alpha}} & reported participant count analysis-of-variance statistic error mean square p-value and comparison threshold & P19 conditional-reasoning-suppression & meta & \hyperref[sec:p19-conditional-reasoning-suppression]{\ref*{sec:p19-conditional-reasoning-suppression}} & \thepage & These are reported empirical quantities with the source's repeated-response denominators. Rounded endorsement rates are not converted into invented participant counts. \\
\phantomsection\label{entry:symbol-registry-01475}\(p,q,r,\mathsf{Might}(r),W,[A],I,S,\operatorname{Poss}_I(A),C_I(A,B)\) & \texttt{\detokenize{p,}\allowbreak\detokenize{q,}\allowbreak\detokenize{r,}\allowbreak\detokenize{\mathsf{Might},}\allowbreak\detokenize{W,}\allowbreak\detokenize{[A],}\allowbreak\detokenize{I,}\allowbreak\detokenize{S,}\allowbreak\detokenize{\operatorname{Poss}_I,}\allowbreak\detokenize{C_I}} & task atoms English possibility abbreviation valuation universe truth set salient state subset possibility and conditional acceptance & P19 conditional-interpretation-experiments & meta & \hyperref[sec:p19-conditional-interpretation-experiments]{\ref*{sec:p19-conditional-interpretation-experiments}} & \thepage & The English might premise is not automatically identified with one psychological operation. Acceptance as a reasoning state requires nonemptiness. \\
\phantomsection\label{entry:symbol-registry-01476}\(E_A(S),E_{A,B}(S),I_0,J_{\rm wet},J_{\rm dry},K_{\rm wet},K_{\rm dry}\) & \texttt{\detokenize{E_A,}\allowbreak\detokenize{E_{A,B},}\allowbreak\detokenize{I_0,}\allowbreak\detokenize{J_{\rm wet},}\allowbreak\detokenize{J_{\rm dry},}\allowbreak\detokenize{K_{\rm wet},}\allowbreak\detokenize{K_{\rm dry}}} & restriction operations initial state and the two fully specified context expansions & P19 conditional-interpretation-experiments & meta & \hyperref[sec:p19-conditional-interpretation-experiments]{\ref*{sec:p19-conditional-interpretation-experiments}} & \thepage & Restriction operations remain defined on empty subsets. The two expansions satisfy the same possibility requirement but give different conclusion acceptance. \\
\phantomsection\label{entry:symbol-registry-01501}\(H,H^c,D,D^c,\pi,u,v,x,z,s,\theta,a,b,c,d,N\) & \texttt{\detokenize{H,}\allowbreak\detokenize{H^c,}\allowbreak\detokenize{D,}\allowbreak\detokenize{D^c,}\allowbreak\detokenize{\pi,}\allowbreak\detokenize{u,}\allowbreak\detokenize{v,}\allowbreak\detokenize{x,}\allowbreak\detokenize{z,}\allowbreak\detokenize{s,}\allowbreak\detokenize{\theta,}\allowbreak\detokenize{a,}\allowbreak\detokenize{b,}\allowbreak\detokenize{c,}\allowbreak\detokenize{d,}\allowbreak\detokenize{N}} & hypothesis evidence complements base rate conditional rates count-tree leaves total and posterior & P19 natural-frequencies-bayesian-reasoning & meta & \hyperref[sec:p19-natural-frequencies-bayesian-reasoning]{\ref*{sec:p19-natural-frequencies-bayesian-reasoning}} & \thepage & Natural frequencies share one population tree. A rounded displayed count tree and an exact rate model may give different benchmark fractions. \\
\phantomsection\label{entry:symbol-registry-01502}\(R,T,\varepsilon,r,h,m,\theta_{\rm rates},\mathrm{PS},\mathrm{PQ},\mathrm{FS},\mathrm{FQ},\mathrm{RF},\mathrm P\) & \texttt{\detokenize{R,}\allowbreak\detokenize{T,}\allowbreak\detokenize{\varepsilon,}\allowbreak\detokenize{r,}\allowbreak\detokenize{h,}\allowbreak\detokenize{m,}\allowbreak\detokenize{\theta_{\rm rates},}\allowbreak\detokenize{\mathrm{PS},}\allowbreak\detokenize{\mathrm{PQ},}\allowbreak\detokenize{\mathrm{FS},}\allowbreak\detokenize{\mathrm{FQ},}\allowbreak\detokenize{\mathrm{RF},}\allowbreak\detokenize{\mathrm P}} & reported estimates task totals tolerances locally declared shortcuts rate benchmark and study format codes & P19 natural-frequencies-bayesian-reasoning & meta & \hyperref[sec:p19-natural-frequencies-bayesian-reasoning]{\ref*{sec:p19-natural-frequencies-bayesian-reasoning}} & \thepage & PS/PQ/FS/FQ denote the first study formats. RF in the second study denotes relative frequency; it is not automatically a natural-frequency tree. \\
\phantomsection\label{entry:symbol-registry-01503}\(C,Y,w_c,m,\mu,\mathrm{Rel},\mathrm{Res},L,n_c,\widehat m,g,Q\) & \texttt{\detokenize{C,}\allowbreak\detokenize{Y,}\allowbreak\detokenize{w_c,}\allowbreak\detokenize{m,}\allowbreak\detokenize{\mu,}\allowbreak\detokenize{\mathrm{Rel},}\allowbreak\detokenize{\mathrm{Res},}\allowbreak\detokenize{L,}\allowbreak\detokenize{n_c,}\allowbreak\detokenize{\widehat m,}\allowbreak\detokenize{g,}\allowbreak\detokenize{Q}} & reported confidence correctness group weights conditional correctness mean reliability resolution squared loss group counts empirical correctness and recalibration maps & P19 calibration-extreme-confidence & meta & \hyperref[sec:p19-calibration-extreme-confidence]{\ref*{sec:p19-calibration-extreme-confidence}} & \thepage & Empirical calibration groups retain repeated responses within people. Calibration and discrimination are different properties. \\
\phantomsection\label{entry:symbol-registry-01504}\(B=h(C),a_B,b_B,v_B,\gamma_B,f_{ij},E_{ij},P_B\) & \texttt{\detokenize{B=h(C),}\allowbreak\detokenize{a_B,}\allowbreak\detokenize{b_B,}\allowbreak\detokenize{v_B,}\allowbreak\detokenize{\gamma_B,}\allowbreak\detokenize{f_{ij},}\allowbreak\detokenize{E_{ij},}\allowbreak\detokenize{P_B}} & coarse confidence bin mean forecast mean correctness within-bin forecast variance covariance category forecasts outcomes and original Brier score & P19 calibration-extreme-confidence & meta & \hyperref[sec:p19-calibration-extreme-confidence]{\ref*{sec:p19-calibration-extreme-confidence}} & \thepage & The original category-sum Brier score has binary range zero to two. Its one-coordinate counterpart differs by a factor of two. \\
\phantomsection\label{entry:symbol-registry-01505}\(t,A_t,k,w,b,Z_i,G,q_i\) & \texttt{\detokenize{t,}\allowbreak\detokenize{A_t,}\allowbreak\detokenize{k,}\allowbreak\detokenize{w,}\allowbreak\detokenize{b,}\allowbreak\detokenize{Z_i,}\allowbreak\detokenize{G,}\allowbreak\detokenize{q_i}} & confidence threshold selected answers odds ratio wrong-answer loss chip reward payoff game gain and subjective answer probability & P19 calibration-extreme-confidence & meta & \hyperref[sec:p19-calibration-extreme-confidence]{\ref*{sec:p19-calibration-extreme-confidence}} & \thepage & The certainty-game calculation uses the stated subjective chip probability as well as answer probabilities. Expected gain does not require independent trials. \\
\phantomsection\label{entry:symbol-registry-01526}\(N,G^2(d),F(d_1,d_2),t,p\) & \texttt{\detokenize{N,}\allowbreak\detokenize{G^2(d),}\allowbreak\detokenize{F(d_1,}\allowbreak\detokenize{d_2),}\allowbreak\detokenize{t,}\allowbreak\detokenize{p}} & participant counts likelihood-ratio statistic variance-ratio statistic mean-comparison statistic and reported p-value & P19 analogy-and-problem-transfer & meta & \hyperref[sec:p19-analogy-and-problem-transfer]{\ref*{sec:p19-analogy-and-problem-transfer}} & \thepage & Outcome percentages can use subgroup denominators. The two-condition binary dispersion comparison has one degree of freedom. \\
\phantomsection\label{entry:symbol-registry-01527}\(B,T,m,\tau_m\) & \texttt{\detokenize{B,}\allowbreak\detokenize{T,}\allowbreak\detokenize{m,}\allowbreak\detokenize{\tau_m}} & finite base and target object sets correspondence and induced expression substitution & P19 analogy-and-problem-transfer & meta & \hyperref[sec:p19-analogy-and-problem-transfer]{\ref*{sec:p19-analogy-and-problem-transfer}} & \thepage & The free-syntax model is a declared benchmark rather than a recovered participant representation. \\
\phantomsection\label{entry:symbol-registry-01528}\(A,R,S,P,\mathsf{Supports},e_j,F\) & \texttt{\detokenize{A,}\allowbreak\detokenize{R,}\allowbreak\detokenize{S,}\allowbreak\detokenize{P,}\allowbreak\detokenize{\mathsf{Supports},}\allowbreak\detokenize{e_j,}\allowbreak\detokenize{F}} & attribute object relations additional binary candidate relation higher-order constructor argument expressions and a generic constructor & P19 analogy-and-problem-transfer & meta & \hyperref[sec:p19-analogy-and-problem-transfer]{\ref*{sec:p19-analogy-and-problem-transfer}} & \thepage & Constructor names and argument positions remain fixed. Structural substitution supplies no unprovided causal or explanatory semantics. \\
\phantomsection\label{entry:symbol-registry-01529}\(h,q,k,x_j,\lambda\) & \texttt{\detokenize{h,}\allowbreak\detokenize{q,}\allowbreak\detokenize{k,}\allowbreak\detokenize{x_j,}\allowbreak\detokenize{\lambda}} & route capacity target requirement available routes allocation and positive common scaling & P19 analogy-and-problem-transfer & meta & \hyperref[sec:p19-analogy-and-problem-transfer]{\ref*{sec:p19-analogy-and-problem-transfer}} & \thepage & The threshold model assumes separate route loads and additive simultaneous arrival. These are explicit premises of its transfer proof. \\
\phantomsection\label{entry:symbol-registry-01530}\(\mathcal M_0,\mathcal M_1\) & \texttt{\detokenize{\mathcal M_0,}\allowbreak\detokenize{\mathcal M_1}} & two target structures agreeing on known relations and differing on the candidate binary relation & P19 analogy-and-problem-transfer & meta & \hyperref[sec:p19-analogy-and-problem-transfer]{\ref*{sec:p19-analogy-and-problem-transfer}} & \thepage & The countermodels demonstrate failure of entailment despite a common structural correspondence. \\
\phantomsection\label{entry:symbol-registry-01672}\(\Omega,\mathcal F,P,A,B,E,q\) & \texttt{\detokenize{\Omega,}\allowbreak\detokenize{\mathcal F,}\allowbreak\detokenize{P,}\allowbreak\detokenize{A,}\allowbreak\detokenize{B,}\allowbreak\detokenize{E,}\allowbreak\detokenize{q}} & probability space events fixed information and elicited judgment & Legacy pre-batch6 P19.7 reader meaning & meta & \hyperref[def:p19-conjunction-events]{\ref*{def:p19-conjunction-events}} & \thepage & P(E)>0 for displayed conditioning; q is a measured response and need not satisfy probability axioms.  Numeric chapter scope migrated to its frozen pre-batch6 source-labelled reader regions; glyph, role and precedence unchanged. \\
\phantomsection\label{entry:symbol-registry-01673}\(\operatorname{rank}(C)\) & \texttt{\detokenize{\operatorname{rank}(C)}} & elicited ordinal rank of a description & Legacy pre-batch6 P19.7 reader meaning & meta & \hyperref[def:p19-conjunction-events]{\ref*{def:p19-conjunction-events}} & \thepage & Smaller rank means greater judged probability; similarity and representativeness are separate measures.  Numeric chapter scope migrated to its frozen pre-batch6 source-labelled reader regions; glyph, role and precedence unchanged. \\
\phantomsection\label{entry:symbol-registry-01674}\(q_i,p_i,w_i\) & \texttt{\detokenize{q_i,}\allowbreak\detokenize{p_i,}\allowbreak\detokenize{w_i}} & individual conjunction and constituent judgments with common aggregation weights & Legacy pre-batch6 P19.7 reader meaning & meta & \hyperref[thm:p19-conjunction-aggregation]{\ref*{thm:p19-conjunction-aggregation}} & \thepage & Finitely many common respondents; weights are nonnegative and sum to one.  Numeric chapter scope migrated to its frozen pre-batch6 source-labelled reader regions; glyph, role and precedence unchanged. \\
\phantomsection\label{entry:symbol-registry-01675}\(T,F,A,J\) & \texttt{\detokenize{T,}\allowbreak\detokenize{F,}\allowbreak\detokenize{A,}\allowbreak\detokenize{J}} & teller feminist accountant and jazz-player task events & Legacy pre-batch6 P19.7 reader meaning & meta & \hyperref[sec:p19-conjunction-persons]{\ref*{sec:p19-conjunction-persons}} & \thepage & These event names are local to the Linda and Bill profiles and are not generic manuscript predicates.  Numeric chapter scope migrated to its frozen pre-batch6 source-labelled reader regions; glyph, role and precedence unchanged. \\
\phantomsection\label{entry:symbol-registry-01676}\(T^*,N\) & \texttt{\detokenize{T^*,}\allowbreak\detokenize{N}} & explicitly inclusive teller description and respondent sample size & Legacy pre-batch6 P19.7 reader meaning & meta & \hyperref[sec:p19-conjunction-transparent]{\ref*{sec:p19-conjunction-transparent}} & \thepage & N refers to its reported condition; a hypothetical count denominator is distinguished from the number of respondents.  Numeric chapter scope migrated to its frozen pre-batch6 source-labelled reader regions; glyph, role and precedence unchanged. \\
\phantomsection\label{entry:symbol-registry-01677}\(R,G\) & \texttt{\detokenize{R,}\allowbreak\detokenize{G}} & red and green dice outcomes & Legacy pre-batch6 P19.7 reader meaning & meta & \hyperref[sec:p19-conjunction-betting]{\ref*{sec:p19-conjunction-betting}} & \thepage & Fair six-face die has two red and four green faces; string occurrence is within 20 rolls.  Numeric chapter scope migrated to its frozen pre-batch6 source-labelled reader regions; glyph, role and precedence unchanged. \\
\phantomsection\label{entry:symbol-registry-01678}\(\delta(j,c)\) & \texttt{\detokenize{\delta(j,}\allowbreak\detokenize{c)}} & length of the longest retained word prefix after the next color & Legacy pre-batch6 P19.7 reader meaning & meta & \hyperref[annex:p19.conjunction-bets]{\ref*{annex:p19.conjunction-bets}} & \thepage & Finite prefix-state recurrence counts overlapping contiguous word occurrences.  Numeric chapter scope migrated to its frozen pre-batch6 source-labelled reader regions; glyph, role and precedence unchanged. \\
\phantomsection\label{entry:symbol-registry-01679}\(M,A,B\) & \texttt{\detokenize{M,}\allowbreak\detokenize{A,}\allowbreak\detokenize{B}} & profile added event and basic event in the source association diagrams & Legacy pre-batch6 P19.7 reader meaning & meta & \hyperref[sec:p19-conjunction-causes]{\ref*{sec:p19-conjunction-causes}} & \thepage & The source Figure 1 relations are associative descriptions, not identified structural causal models.  Numeric chapter scope migrated to its frozen pre-batch6 source-labelled reader regions; glyph, role and precedence unchanged. \\
\phantomsection\label{entry:symbol-registry-01693}\(T,F,E,P\) & \texttt{\detokenize{T,}\allowbreak\detokenize{F,}\allowbreak\detokenize{E,}\allowbreak\detokenize{P}} & teller and feminist events common information and probability measure & Legacy pre-batch6 P19.8 reader meaning & meta & \hyperref[def:p19-conjunction-language-record]{\ref*{def:p19-conjunction-language-record}} & \thepage & The intended mathematical comparison is distinguished from the elicited wording and response mode.  Numeric chapter scope migrated to its frozen pre-batch6 source-labelled reader regions; glyph, role and precedence unchanged. \\
\phantomsection\label{entry:symbol-registry-01694}\(E_1,E_2\) & \texttt{\detokenize{E_1,}\allowbreak\detokenize{E_2}} & mutually exclusive exhaustive information conditions in the eight-atom example & Legacy pre-batch6 P19.8 reader meaning & meta & \hyperref[def:p19-conjunction-profile-relevance]{\ref*{def:p19-conjunction-profile-relevance}} & \thepage & Each condition has probability one half; its four displayed conditional masses define the full model.  Numeric chapter scope migrated to its frozen pre-batch6 source-labelled reader regions; glyph, role and precedence unchanged. \\
\phantomsection\label{entry:symbol-registry-01695}\(f\) & \texttt{\detokenize{f}} & effect-size summary reported by the 1999 paper & Legacy pre-batch6 P19.8 reader meaning & meta & \hyperref[sec:p19-conjunction-language-study-two]{\ref*{sec:p19-conjunction-language-study-two}} & \thepage & The manuscript reports the source values and does not infer unreported confidence intervals or reanalyses.  Numeric chapter scope migrated to its frozen pre-batch6 source-labelled reader regions; glyph, role and precedence unchanged. \\
\phantomsection\label{entry:symbol-registry-01696}\(C(M),A,x,y,z\) & \texttt{\detokenize{C(M),}\allowbreak\detokenize{A,}\allowbreak\detokenize{x,}\allowbreak\detokenize{y,}\allowbreak\detokenize{z}} & choice correspondence menu subset and alternatives in the contraction-consistency example & Legacy pre-batch6 P19.8 reader meaning & meta & \hyperref[sec:p19-conjunction-language-limits]{\ref*{sec:p19-conjunction-language-limits}} & \thepage & x declines dessert y takes the eclair and z takes another pastry; full social consequences may depend on the menu.  Numeric chapter scope migrated to its frozen pre-batch6 source-labelled reader regions; glyph, role and precedence unchanged. \\
\phantomsection\label{entry:symbol-registry-01864}\(\theta,\delta\) & \texttt{\detokenize{\theta,}\allowbreak\detokenize{\delta}} & stipulated angle and shorter single-axis disparity & Legacy pre-batch6 P19.9 reader meaning & meta & \hyperref[def:p19-rotation-benchmark]{\ref*{def:p19-rotation-benchmark}} & \thepage & Local scope sec:p19-rotation-task. No disparity-to-congruence is assigned to incongruent mirror foils.  Numeric chapter scope migrated to its frozen pre-batch6 source-labelled reader regions; glyph, role and precedence unchanged. \\
\phantomsection\label{entry:symbol-registry-01865}\(m,\delta_i,t_i,\bar\delta,\bar t,a,b,S_{\delta\delta},\widehat t\) & \texttt{\detokenize{m,}\allowbreak\detokenize{\delta_i,}\allowbreak\detokenize{t_i,}\allowbreak\detokenize{\bar\delta,}\allowbreak\detokenize{\bar t,}\allowbreak\detokenize{a,}\allowbreak\detokenize{b,}\allowbreak\detokenize{S_{\delta\delta},}\allowbreak\detokenize{\widehat t}} & fictional summaries moments fitted coefficients and times & Legacy pre-batch6 P19.9 reader meaning & meta & \hyperref[thm:p19-rotation-fit-and-identification]{\ref*{thm:p19-rotation-fit-and-identification}} & \thepage & Local scope sec:p19-rotation-worked-example. b has units milliseconds per degree.  Numeric chapter scope migrated to its frozen pre-batch6 source-labelled reader regions; glyph, role and precedence unchanged. \\
\phantomsection\label{entry:symbol-registry-01866}\(j,\Delta,E\) & \texttt{\detokenize{j,}\allowbreak\detokenize{\Delta,}\allowbreak\detokenize{E}} & grid index angular step and common illustrative noise & Legacy pre-batch6 P19.9 reader meaning & meta & \hyperref[thm:p19-rotation-fit-and-identification]{\ref*{thm:p19-rotation-fit-and-identification}} & \thepage & Local scope annex:p19.rotation.fit.identification. The two processes are stipulated models, not observed mental operations.  Numeric chapter scope migrated to its frozen pre-batch6 source-labelled reader regions; glyph, role and precedence unchanged. \\
\phantomsection\label{entry:symbol-registry-01867}\(K,d',r\) & \texttt{\detokenize{K,}\allowbreak\detokenize{d',}\allowbreak\detokenize{r}} & reported binding capacity discrimination label and correlation & Legacy pre-batch6 P19.9 reader meaning & meta & \hyperref[sec:p19-xu-franconeri]{\ref*{sec:p19-xu-franconeri}} & \thepage & Local scope sec:p19-xu-franconeri. The manuscript cites but does not print the K estimator; Fig2D labels performance d-prime.  Numeric chapter scope migrated to its frozen pre-batch6 source-labelled reader regions; glyph, role and precedence unchanged. \\
\phantomsection\label{entry:symbol-registry-01868}\(N,k,d_k,h,f,g\) & \texttt{\detokenize{N,}\allowbreak\detokenize{k,}\allowbreak\detokenize{d_k,}\allowbreak\detokenize{h,}\allowbreak\detokenize{f,}\allowbreak\detokenize{g}} & ideal part and remembered counts detection hit false alarm and guessing probabilities & Legacy pre-batch6 P19.9 reader meaning & meta & \hyperref[thm:p19-feature-swap-benchmark]{\ref*{thm:p19-feature-swap-benchmark}} & \thepage & Local scope sec:p19-rotation-capacity-benchmark. This ideal counting model is not attributed as the published capacity estimator.  Numeric chapter scope migrated to its frozen pre-batch6 source-labelled reader regions; glyph, role and precedence unchanged. \\
\phantomsection\label{entry:symbol-registry-01869}\(n,q_i,y_i,\bar q,\bar y\) & \texttt{\detokenize{n,}\allowbreak\detokenize{q_i,}\allowbreak\detokenize{y_i,}\allowbreak\detokenize{\bar q,}\allowbreak\detokenize{\bar y}} & item count predicted recall binary criterion and mean prediction and performance & Legacy pre-batch4 P19.10 reader meaning & meta & \hyperref[def:p19-jol-record]{\ref*{def:p19-jol-record}} & \thepage & Local scope sec:p19-judgments-of-learning. The record concerns one participant in one declared condition; item counts are not independent participant counts.  Numeric chapter scope migrated to its frozen pre-batch4 source-labelled reader regions; glyph, role and precedence unchanged. \\
\phantomsection\label{entry:symbol-registry-01870}\(N_+,N_-,G\) & \texttt{\detokenize{N_+,}\allowbreak\detokenize{N_-,}\allowbreak\detokenize{G}} & concordant and discordant comparable-pair counts and Goodman-Kruskal gamma & Legacy pre-batch4 P19.10 reader meaning & meta & \hyperref[def:p19-jol-gamma]{\ref*{def:p19-jol-gamma}} & \thepage & Local scope sec:p19-judgments-of-learning. Pairs tied on either variable are excluded; zero denominator is undefined rather than zero accuracy.  Numeric chapter scope migrated to its frozen pre-batch4 source-labelled reader regions; glyph, role and precedence unchanged. \\
\phantomsection\label{entry:symbol-registry-01871}\(x,y\) & \texttt{\detokenize{x,}\allowbreak\detokenize{y}} & predicted and observed percentages in source calibration-line summaries & Legacy pre-batch4 P19.10 reader meaning & meta & \hyperref[sec:p19-jol-1991-main]{\ref*{sec:p19-jol-1991-main}} & \thepage & Local scope sec:p19-judgments-of-learning. Both axes use percentage units; intercepts must be divided by 100 if the equations are rewritten on a probability scale.  Numeric chapter scope migrated to its frozen pre-batch4 source-labelled reader regions; glyph, role and precedence unchanged. \\
\phantomsection\label{entry:symbol-registry-01872}\(C_{\rm norm}\) & \texttt{\detokenize{C_{\rm norm}}} & source-reported word-norm concreteness selection score & Legacy pre-batch4 P19.10 reader meaning & meta & \hyperref[sec:p19-jol-1992-design]{\ref*{sec:p19-jol-1992-design}} & \thepage & Local scope sec:p19-judgments-of-learning. The verified source criterion is at least 6.08; this is not the confidence variable C or a mean recall score.  Numeric chapter scope migrated to its frozen pre-batch4 source-labelled reader regions; glyph, role and precedence unchanged. \\
\phantomsection\label{entry:symbol-registry-01873}\(\operatorname{Res},m(C),\mu\) & \texttt{\detokenize{\operatorname{Res},}\allowbreak\detokenize{m(C),}\allowbreak\detokenize{\mu}} & inherited conditional calibration resolution and outcome-rate notation & Legacy pre-batch4 P19.10 reader meaning & meta & \hyperref[sec:p19-judgments-of-learning]{\ref*{sec:p19-judgments-of-learning}} & \thepage & Local scope sec:p19-judgments-of-learning. The calibration chapter's resolution is distinguished from within-person gamma; the two measures are not identified.  Numeric chapter scope migrated to its frozen pre-batch4 source-labelled reader regions; glyph, role and precedence unchanged. \\
\phantomsection\label{entry:symbol-registry-01901}\(D_1,D_2,D_3,S_1,S_2\) & \texttt{\detokenize{D_1,}\allowbreak\detokenize{D_2,}\allowbreak\detokenize{D_3,}\allowbreak\detokenize{S_1,}\allowbreak\detokenize{S_2}} & three stipulated disease propositions and two symptoms & Legacy pre-batch4 P19.11 reader meaning & meta & \hyperref[def:p19-explanation-task]{\ref*{def:p19-explanation-task}} & \thepage & The focal explanations invoke one cause versus two; not causing a symptom is distinguished from preventing it.  Numeric chapter scope migrated to its frozen pre-batch4 source-labelled reader regions; glyph, role and precedence unchanged. \\
\phantomsection\label{entry:symbol-registry-01902}\(T,H,M\) & \texttt{\detokenize{T,}\allowbreak\detokenize{H,}\allowbreak\detokenize{M}} & explicit closed background permitted hypothesis atoms and observations & Legacy pre-batch4 P19.11 reader meaning & meta & \hyperref[thm:p19-explanation-logical-menu]{\ref*{thm:p19-explanation-logical-menu}} & \thepage & This is a declared finite-abduction model; participant preferences do not establish its real-world truth.  Numeric chapter scope migrated to its frozen pre-batch4 source-labelled reader regions; glyph, role and precedence unchanged. \\
\phantomsection\label{entry:symbol-registry-01903}\(H_s,H_c,E,p_1,j,t\) & \texttt{\detokenize{H_s,}\allowbreak\detokenize{H_c,}\allowbreak\detokenize{E,}\allowbreak\detokenize{p_1,}\allowbreak\detokenize{j,}\allowbreak\detokenize{t}} & inclusive explanatory events observation event marginal joint and overlap probabilities & Legacy pre-batch4 P19.11 reader meaning & meta & \hyperref[def:p19-explanation-probability]{\ref*{def:p19-explanation-probability}} & \thepage & The explanatory events can overlap; t is their overlap mass rather than the later reported t-test statistic.  Numeric chapter scope migrated to its frozen pre-batch4 source-labelled reader regions; glyph, role and precedence unchanged. \\
\phantomsection\label{entry:symbol-registry-01904}\(R,q_0\) & \texttt{\detokenize{R,}\allowbreak\detokenize{q_0}} & explanatory-event probability ratio and normalized answer score & Legacy pre-batch4 P19.11 reader meaning & meta & \hyperref[thm:p19-explanation-probability]{\ref*{thm:p19-explanation-probability}} & \thepage & The normalized score differs from an actual event posterior when the inclusive explanations overlap.  Numeric chapter scope migrated to its frozen pre-batch4 source-labelled reader regions; glyph, role and precedence unchanged. \\
\phantomsection\label{entry:symbol-registry-01905}\(n_2,n_3\) & \texttt{\detokenize{n_2,}\allowbreak\detokenize{n_3}} & displayed Experiment 2 disease-frequency counts & Legacy pre-batch4 P19.11 reader meaning & meta & \hyperref[sec:p19-explanation-experiment-two]{\ref*{sec:p19-explanation-experiment-two}} & \thepage & Nominal ratio labels are distinguished from the exact fraction 37500/(n2*n3).  Numeric chapter scope migrated to its frozen pre-batch4 source-labelled reader regions; glyph, role and precedence unchanged. \\
\phantomsection\label{entry:symbol-registry-01906}\(a,b,q(R)\) & \texttt{\detokenize{a,}\allowbreak\detokenize{b,}\allowbreak\detokenize{q(R)}} & explicitly oriented intercept slope and simple-answer response probability & Legacy pre-batch4 P19.11 reader meaning & meta & \hyperref[def:p19-explanation-choice-curve]{\ref*{def:p19-explanation-choice-curve}} & \thepage & The positive simple-choice orientation is distinguished from negative constants printed without full coding information.  Numeric chapter scope migrated to its frozen pre-batch4 source-labelled reader regions; glyph, role and precedence unchanged. \\
\phantomsection\label{entry:symbol-registry-01907}\(X,Q_3,\operatorname{erfc}\) & \texttt{\detokenize{X,}\allowbreak\detokenize{Q_3,}\allowbreak\detokenize{\operatorname{erfc}}} & declared chi-square reference variable upper tail and complementary Gaussian integral & Legacy pre-batch4 P19.11 reader meaning & meta & \hyperref[def:p19-explanation-tail]{\ref*{def:p19-explanation-tail}} & \thepage & The computed .06406 uses the printed 7.26 statistic and three-degree reference; it does not reconstruct participant data.  Numeric chapter scope migrated to its frozen pre-batch4 source-labelled reader regions; glyph, role and precedence unchanged. \\
\phantomsection\label{entry:symbol-registry-01908}\(V,T(v),Z,A,Y\) & \texttt{\detokenize{V,}\allowbreak\detokenize{T(v),}\allowbreak\detokenize{Z,}\allowbreak\detokenize{A,}\allowbreak\detokenize{Y}} & fictional population coordinate threshold pre-choice estimate choice and later report & Legacy pre-batch4 P19.11 reader meaning & meta & \hyperref[annex:p19.explanation-choice-model]{\ref*{annex:p19.explanation-choice-model}} & \thepage & These constructed models establish non-unique process and causal interpretations; they are not additional empirical observations.  Numeric chapter scope migrated to its frozen pre-batch4 source-labelled reader regions; glyph, role and precedence unchanged. \\
\phantomsection\label{entry:symbol-registry-03005}\(\mathrm{FOR}\) & \texttt{\detokenize{\mathrm{FOR}}} & feeling of rightness accompanying an initial answer & P19 reasoning metacognition & n/a & \hyperref[def:p19-r86-judgment-types]{\ref*{def:p19-r86-judgment-types}} & \thepage & Judgment stage and rating scale remain distinct from final confidence. \\
\phantomsection\label{entry:symbol-registry-03006}\((t_0,\ldots,t_m)\) & \texttt{\detokenize{t_j}} & ordered monitoring times & P19 reasoning metacognition & n/a & \hyperref[def:p19-r86-stopping-model]{\ref*{def:p19-r86-stopping-model}} & \thepage & Finite observation grid with an explicit final time limit. \\
\phantomsection\label{entry:symbol-registry-03007}\(T\) & \texttt{\detokenize{T}} & maximum deliberation time & P19 reasoning metacognition & n/a & \hyperref[def:p19-r86-stopping-model]{\ref*{def:p19-r86-stopping-model}} & \thepage & The cap makes the finite stopping rule well defined. \\
\phantomsection\label{entry:symbol-registry-03008}\(q_j\) & \texttt{\detokenize{q_j}} & intermediate judgment trajectory & P19 reasoning metacognition & n/a & \hyperref[def:p19-r86-stopping-model]{\ref*{def:p19-r86-stopping-model}} & \thepage & A schematic model variable rather than a verified probability of correctness. \\
\phantomsection\label{entry:symbol-registry-03009}\(h_j\) & \texttt{\detokenize{h_j}} & acceptance criterion at a monitoring time & P19 reasoning metacognition & n/a & \hyperref[def:p19-r86-stopping-model]{\ref*{def:p19-r86-stopping-model}} & \thepage & Diminishing and constant criterion policies are distinct models. \\
\phantomsection\label{entry:symbol-registry-03010}\((q_E(t),q_H(t))\) & \texttt{\detokenize{q_E(t);q_H(t)}} & two constructed confidence trajectories & P19 reasoning metacognition & n/a & \hyperref[thm:p19-r86-stopping]{\ref*{thm:p19-r86-stopping}} & \thepage & No fitted psychological parameters or empirical optimality is claimed. \\
\phantomsection\label{entry:symbol-registry-03011}\(n\) & \texttt{\detokenize{n}} & number of scored reasoning problems in a trace & P19 reasoning metacognition & n/a & \hyperref[def:p19-r86-trace]{\ref*{def:p19-r86-trace}} & \thepage & The fictional calculations are descriptive over specified problems. \\
\phantomsection\label{entry:symbol-registry-03012}\((y_i^{(0)},y_i^{(1)})\) & \texttt{\detokenize{y_i^{(0)};y_i^{(1)}}} & initial and final criterion correctness values & P19 reasoning metacognition & n/a & \hyperref[def:p19-r86-trace]{\ref*{def:p19-r86-trace}} & \thepage & Criterion is fixed separately from confidence and duration. \\
\phantomsection\label{entry:symbol-registry-03013}\((c_i^{(0)},c_i^{(1)})\) & \texttt{\detokenize{c_i^{(0)};c_i^{(1)}}} & initial and final numerical correctness-confidence reports & P19 reasoning metacognition & n/a & \hyperref[def:p19-r86-trace]{\ref*{def:p19-r86-trace}} & \thepage & An ordinal FOR scale is not automatically this probability scale. \\
\phantomsection\label{entry:symbol-registry-03014}\(t_i\) & \texttt{\detokenize{t_i}} & extra deliberation time for a trace problem & P19 reasoning metacognition & n/a & \hyperref[def:p19-r86-trace]{\ref*{def:p19-r86-trace}} & \thepage & Seconds in the worked record; no causal manipulation is implied. \\
\phantomsection\label{entry:symbol-registry-03015}\((N_{01},N_{10})\) & \texttt{\detokenize{N_{01};N_{10}}} & corrections and harmful correctness changes & P19 reasoning metacognition & n/a & \hyperref[def:p19-r86-trace]{\ref*{def:p19-r86-trace}} & \thepage & Other answer changes can remain in the same correctness category. \\
\phantomsection\label{entry:symbol-registry-03016}\((C,Y)\) & \texttt{\detokenize{C;Y}} & confidence report and answer-correctness variables under a finite law & P19 reasoning metacognition & n/a & \hyperref[thm:p19-r86-calibrated-control]{\ref*{thm:p19-r86-calibrated-control}} & \thepage & Uses the existing calibration interface with a reasoning-specific event. \\
\phantomsection\label{entry:symbol-registry-03017}\(h\) & \texttt{\detokenize{h}} & selective-reporting confidence threshold & P19 reasoning metacognition & n/a & \hyperref[thm:p19-r86-calibrated-control]{\ref*{thm:p19-r86-calibrated-control}} & \thepage & Positive selected mass and same-law calibration are explicit assumptions. \\
\phantomsection\label{entry:symbol-registry-03018}\((p_0,p_1)\) & \texttt{\detokenize{p_0;p_1}} & belief-model correctness probabilities for current and additional procedures & P19 reasoning metacognition & n/a & \hyperref[def:p19-r86-step-value]{\ref*{def:p19-r86-step-value}} & \thepage & Prospective improvement is not inferred from current confidence alone. \\
\phantomsection\label{entry:symbol-registry-03019}\((R,L,\kappa)\) & \texttt{\detokenize{R;L;\kappa}} & correct-answer value incorrect-answer loss and extra-procedure cost & P19 reasoning metacognition & n/a & \hyperref[def:p19-r86-step-value]{\ref*{def:p19-r86-step-value}} & \thepage & Feasibility and deadlines are separate from the expected-value identity. \\
\phantomsection\label{entry:symbol-registry-03276}\(c,t\) & \texttt{\detokenize{c;t}} & circle and triangle Boolean propositions & P19 mental-models finite deduction benchmark & n/a & \hyperref[def:p19-r23-live-models]{\ref*{def:p19-r23-live-models}} & \thepage & Coordinate order is circle then triangle; negation represents an explicitly false proposition. \\
\phantomsection\label{entry:symbol-registry-03277}\(\Omega\) & \texttt{\detokenize{\Omega}} & complete valuations of the fixed finite Boolean language & P19 mental-models finite deduction benchmark & n/a & \hyperref[def:p19-r23-live-models]{\ref*{def:p19-r23-live-models}} & \thepage & A logical possibility domain rather than an empirical sample. \\
\phantomsection\label{entry:symbol-registry-03278}\(\Gamma\) & \texttt{\detokenize{\Gamma}} & interpreted premise family & P19 mental-models finite deduction benchmark & n/a & \hyperref[def:p19-r23-live-models]{\ref*{def:p19-r23-live-models}} & \thepage & Nonvacuous consequence separately requires its compatible valuation set to be nonempty. \\
\phantomsection\label{entry:symbol-registry-03279}\(\mathcal M(\Gamma)\) & \texttt{\detokenize{\mathcal M}} & premise-compatible complete valuation set & P19 mental-models finite deduction benchmark & n/a & \hyperref[def:p19-r23-live-models]{\ref*{def:p19-r23-live-models}} & \thepage & A complete formal benchmark need not coincide with an incomplete psychological representation. \\
\phantomsection\label{entry:symbol-registry-03280}\(\gamma\) & \texttt{\detokenize{\gamma}} & member of the interpreted premise family & P19 mental-models finite deduction benchmark & n/a & \hyperref[def:p19-r23-live-models]{\ref*{def:p19-r23-live-models}} & \thepage & Every premise must hold in a compatible valuation. \\
\phantomsection\label{entry:symbol-registry-03281}\(\models\) & \texttt{\detokenize{\models}} & Boolean valuation satisfaction & P19 mental-models finite deduction benchmark & n/a & \hyperref[def:p19-r23-live-models]{\ref*{def:p19-r23-live-models}} & \thepage & The source's extra nonvacuity criterion does not silently redefine classical consequence. \\
\phantomsection\label{entry:symbol-registry-03282}\(\Diamond\) & \texttt{\detokenize{\Diamond}} & ordinary object-language Kripke possibility & P19 mental-models modal counterexample & n/a & \hyperref[thm:p19-r23-possibility]{\ref*{thm:p19-r23-possibility}} & \thepage & Distinguished from existence of a compatible propositional valuation. \\
\phantomsection\label{entry:symbol-registry-03283}\(w\) & \texttt{\detokenize{w}} & sole world of the modal countermodel & P19 mental-models modal counterexample & n/a & \hyperref[annex:p19.r23.possibility]{\ref*{annex:p19.r23.possibility}} & \thepage & The sole accessibility pair is reflexive; circle is false and triangle true. \\
\phantomsection\label{entry:symbol-registry-03284}\(\oplus\) & \texttt{\detokenize{\oplus}} & exclusive disjunction of Boolean formulas & P19 mental-models exclusive-disjunction example & n/a & \hyperref[def:p19-r23-xor]{\ref*{def:p19-r23-xor}} & \thepage & Exactly one argument is true; not a numeric addition or vector operation. \\
\phantomsection\label{entry:symbol-registry-03285}\(u,w\) & \texttt{\detokenize{u;w}} & arguments of the exclusive-disjunction definition & P19 mental-models exclusive-disjunction example & n/a & \hyperref[def:p19-r23-xor]{\ref*{def:p19-r23-xor}} & \thepage & These formula metavariables are distinct from the modal-world scope. \\
\phantomsection\label{entry:symbol-registry-03286}\(k\) & \texttt{\detokenize{k}} & number of independent Boolean atoms in complete enumeration & P19 mental-models enumeration cost & n/a & \hyperref[sec:p19-r23-deduction]{\ref*{sec:p19-r23-deduction}} & \thepage & The full valuation list has two to this power entries; no complexity-class separation is claimed. \\
\phantomsection\label{entry:symbol-registry-03287}\(P\) & \texttt{\detokenize{P}} & normalized finite event probability & P19 mental-models binary joint assessments & n/a & \hyperref[sec:p19-r23-probability]{\ref*{sec:p19-r23-probability}} & \thepage & Probability is distinct from permission and from an empirical response frequency. \\
\phantomsection\label{entry:symbol-registry-03288}\(A,B\) & \texttt{\detokenize{A;B}} & binary events whose marginals and intersection are assessed & P19 mental-models binary joint assessments & n/a & \hyperref[thm:p19-r23-joint]{\ref*{thm:p19-r23-joint}} & \thepage & Negated event notation means complement within the declared sample space. \\
\phantomsection\label{entry:symbol-registry-03289}\(a,b,q\) & \texttt{\detokenize{a;b;q}} & two marginals and their intersection assessment & P19 mental-models binary joint assessments & n/a & \hyperref[thm:p19-r23-joint]{\ref*{thm:p19-r23-joint}} & \thepage & Cell algebra always normalizes; the bounds establish nonnegativity. \\
\phantomsection\label{entry:symbol-registry-03290}\(u\) & \texttt{\detokenize{u}} & union assessment & P19 mental-models binary joint assessments & n/a & \hyperref[thm:p19-r23-joint]{\ref*{thm:p19-r23-joint}} & \thepage & Intersection assessment is obtained by inclusion-exclusion. \\
\phantomsection\label{entry:symbol-registry-03291}\(r\) & \texttt{\detokenize{r}} & conditional assessment on the first event & P19 mental-models binary joint assessments & n/a & \hyperref[thm:p19-r23-joint]{\ref*{thm:p19-r23-joint}} & \thepage & The ordinary ratio requires positive conditioning-event probability. \\
\phantomsection\label{entry:symbol-registry-03292}\(r'\) & \texttt{\detokenize{r'}} & converse conditional assessment on the second event & P19 mental-models binary joint assessments & n/a & \hyperref[thm:p19-r23-joint]{\ref*{thm:p19-r23-joint}} & \thepage & Its ratio additionally requires a nonnull second event and the initial assessments must define an admissible joint distribution. \\
\phantomsection\label{entry:symbol-registry-03293}\(a,r,s\) & \texttt{\detokenize{a;r;s}} & first-event marginal and two conditional branch parameters & P19 mental-models admissible conditional kernel & n/a & \hyperref[thm:p19-r23-kernel]{\ref*{thm:p19-r23-kernel}} & \thepage & Parameters lie in the unit interval; a null branch is arbitrary and unidentifiable. \\
\phantomsection\label{entry:symbol-registry-03294}\(q_{11},q_{10},q_{01},q_{00}\) & \texttt{\detokenize{q_{11};q_{10};q_{01};q_{00}}} & four joint cells in the converse construction & P19 mental-models admissible conditional kernel & n/a & \hyperref[annex:p19.r23.kernel]{\ref*{annex:p19.r23.kernel}} & \thepage & Nonnegative cells sum to one; division is performed only on a nonnull row. \\
\phantomsection\label{entry:symbol-registry-03295}\(D,S\) & \texttt{\detokenize{D;S}} & disease and symptom events in a stipulated chance example & P19 mental-models chance-counting example & n/a & \hyperref[def:p19-r23-equiprobability]{\ref*{def:p19-r23-equiprobability}} & \thepage & Counts are invented task data from the source illustration rather than clinical prevalence. \\
\phantomsection\label{entry:symbol-registry-03296}\(R,G,B\) & \texttt{\detokenize{R;G;B}} & red green and blue marble propositions & P19 mental-models marble possibility example & n/a & \hyperref[thm:p19-r23-chances]{\ref*{thm:p19-r23-chances}} & \thepage & Boolean triples are ordered red green blue; letters are distinct from the binary joint-event scope. \\
\phantomsection\label{entry:symbol-registry-03297}\(L[S]R\) & \texttt{\detokenize{L;S;R}} & left siding and right words in the railway state & P19 mental-models railway reversal & n/a & \hyperref[def:p19-r23-railway]{\ref*{def:p19-r23-railway}} & \thepage & Movement rules preserve block order; labels may repeat. \\
\phantomsection\label{entry:symbol-registry-03298}\(T\) & \texttt{\detokenize{T}} & nonempty car block transferred by one allowed movement & P19 mental-models railway reversal & n/a & \hyperref[def:p19-r23-railway]{\ref*{def:p19-r23-railway}} & \thepage & Terminal and initial positions refer to the declared word orientation. \\
\phantomsection\label{entry:symbol-registry-03299}\(\epsilon\) & \texttt{\detokenize{\epsilon}} & empty car word & P19 mental-models railway reversal & n/a & \hyperref[def:p19-r23-railway]{\ref*{def:p19-r23-railway}} & \thepage & Not an error tolerance or infinitesimal. \\
\phantomsection\label{entry:symbol-registry-03300}\(a_i\) & \texttt{\detokenize{a_i}} & car label at original position i & P19 mental-models railway reversal & n/a & \hyperref[def:p19-r23-reversal]{\ref*{def:p19-r23-reversal}} & \thepage & Correctness works even when distinct physical cars bear repeated labels. \\
\phantomsection\label{entry:symbol-registry-03301}\(n\) & \texttt{\detokenize{n}} & finite input train length & P19 mental-models railway reversal & n/a & \hyperref[def:p19-r23-reversal]{\ref*{def:p19-r23-reversal}} & \thepage & The empty input is handled by a separate no-operation branch. \\
\phantomsection\label{entry:symbol-registry-03302}\(k\) & \texttt{\detokenize{k}} & number of cars already transferred to the right at a loop boundary & P19 mental-models railway reversal & n/a & \hyperref[thm:p19-r23-reversal]{\ref*{thm:p19-r23-reversal}} & \thepage & The siding length is the decreasing nonnegative termination measure. \\
\phantomsection\label{entry:symbol-registry-03685}\(m\) & \texttt{\detokenize{m}} & number of event categories & P19 heuristics exposure-selection model & n/a & \hyperref[def:p19-r24-exposure]{\ref*{def:p19-r24-exposure}} & \thepage & A positive finite category count; not a target mean. \\
\phantomsection\label{entry:symbol-registry-03686}\(i,j\) & \texttt{\detokenize{i;j}} & category indices & P19 heuristics exposure-selection model & n/a & \hyperref[def:p19-r24-exposure]{\ref*{def:p19-r24-exposure}} & \thepage & Indices range from one through m. \\
\phantomsection\label{entry:symbol-registry-03687}\(p_i,p\) & \texttt{\detokenize{p_i;p}} & strictly positive target proportions and their vector & P19 heuristics exposure-selection model & n/a & \hyperref[def:p19-r24-exposure]{\ref*{def:p19-r24-exposure}} & \thepage & A stipulated target law; not the empirical significance p in study summaries. \\
\phantomsection\label{entry:symbol-registry-03688}\(s_i\) & \texttt{\detokenize{s_i}} & category-conditional exposure admission probability & P19 heuristics exposure-selection model & n/a & \hyperref[def:p19-r24-exposure]{\ref*{def:p19-r24-exposure}} & \thepage & Weights lie in the unit interval and have positive total selected mass. \\
\phantomsection\label{entry:symbol-registry-03689}\(q_i,q\) & \texttt{\detokenize{q_i;q}} & conditional exposure-record proportions and their vector & P19 heuristics exposure-selection model & n/a & \hyperref[def:p19-r24-exposure]{\ref*{def:p19-r24-exposure}} & \thepage & An observed record law need not equal the target law. \\
\phantomsection\label{entry:symbol-registry-03690}\(I\) & \texttt{\detokenize{I}} & event of admission to the exposure record & P19 heuristics exposure-selection model & n/a & \hyperref[annex:p19.r24.exposure]{\ref*{annex:p19.r24.exposure}} & \thepage & An event under the stipulated selection law. \\
\phantomsection\label{entry:symbol-registry-03691}\(\lambda\) & \texttt{\detokenize{\lambda}} & positive normalizing scale constructing admissible selection weights & P19 heuristics exposure-selection model & n/a & \hyperref[annex:p19.r24.exposure]{\ref*{annex:p19.r24.exposure}} & \thepage & Chosen at most every positive p\_i/q\_i; no empirical selection parameter is estimated. \\
\phantomsection\label{entry:symbol-registry-03692}\(P\) & \texttt{\detokenize{P}} & probability under the declared target and selection law & P19 heuristics exposure-selection model & n/a & \hyperref[annex:p19.r24.exposure]{\ref*{annex:p19.r24.exposure}} & \thepage & Target and admitted-record probabilities remain distinct. \\
\phantomsection\label{entry:symbol-registry-03693}\(B,G\) & \texttt{\detokenize{B;G}} & boy and girl outcome labels for exact birth strings & P19 heuristics independent-string benchmark & n/a & \hyperref[thm:p19-r24-strings]{\ref*{thm:p19-r24-strings}} & \thepage & Labels describe model outcomes rather than demographic frequencies. \\
\phantomsection\label{entry:symbol-registry-03694}\(H,T\) & \texttt{\detokenize{H;T}} & head and tail outcome labels for exact coin strings & P19 heuristics independent-string benchmark & n/a & \hyperref[thm:p19-r24-strings]{\ref*{thm:p19-r24-strings}} & \thepage & Exact strings are not events specifying only a count of heads. \\
\phantomsection\label{entry:symbol-registry-03695}\(p\) & \texttt{\detokenize{p}} & common Bernoulli outcome probability & P19 heuristics independent-string benchmark & n/a & \hyperref[thm:p19-r24-strings]{\ref*{thm:p19-r24-strings}} & \thepage & Birth and coin illustrations each declare their own independent Bernoulli model; p is not a significance level. \\
\phantomsection\label{entry:symbol-registry-03696}\(P\) & \texttt{\detokenize{P}} & probability under the stipulated independent birth or coin law & P19 heuristics independent-string benchmark & n/a & \hyperref[thm:p19-r24-strings]{\ref*{thm:p19-r24-strings}} & \thepage & Fair marginal probabilities alone do not establish the product law. \\
\phantomsection\label{entry:symbol-registry-03697}\(N\) & \texttt{\detokenize{N}} & fixed number of independently sampled cases & P19 heuristics conditioned-count benchmark & n/a & \hyperref[thm:p19-r24-counts]{\ref*{thm:p19-r24-counts}} & \thepage & In the birth illustration N is the surveyed-family total; it is not given by the observed count 72. \\
\phantomsection\label{entry:symbol-registry-03698}\(X,Y\) & \texttt{\detokenize{X;Y}} & random counts of two disjoint exact-string categories & P19 heuristics conditioned-count benchmark & n/a & \hyperref[thm:p19-r24-counts]{\ref*{thm:p19-r24-counts}} & \thepage & Equal underlying category probabilities need not yield equal observed counts. \\
\phantomsection\label{entry:symbol-registry-03699}\(a,b\) & \texttt{\detokenize{a;b}} & probabilities of two disjoint categories & P19 heuristics conditioned-count benchmark & n/a & \hyperref[thm:p19-r24-counts]{\ref*{thm:p19-r24-counts}} & \thepage & Require zero at most a less than one and a+b at most one. \\
\phantomsection\label{entry:symbol-registry-03700}\(x,y\) & \texttt{\detokenize{x;y}} & attainable fixed category-count values & P19 heuristics conditioned-count benchmark & n/a & \hyperref[thm:p19-r24-counts]{\ref*{thm:p19-r24-counts}} & \thepage & Lowercase values distinguish conditioning from the random counts X and Y. \\
\phantomsection\label{entry:symbol-registry-03701}\(E\) & \texttt{\detokenize{E}} & expectation under the declared independent-case sampling law & P19 heuristics conditioned-count benchmark & n/a & \hyperref[thm:p19-r24-counts]{\ref*{thm:p19-r24-counts}} & \thepage & Conditional expectation depends on the fixed sample size N. \\
\phantomsection\label{entry:symbol-registry-03702}\(P\) & \texttt{\detokenize{P}} & probability under the declared independent-case sampling law & P19 heuristics conditioned-count benchmark & n/a & \hyperref[annex:p19.r24.counts]{\ref*{annex:p19.r24.counts}} & \thepage & The multinomial law supplies the conditional binomial calculation. \\
\phantomsection\label{entry:symbol-registry-03703}\(\operatorname{Bin}\) & \texttt{\detokenize{\operatorname{Bin}}} & binomial distribution & P19 heuristics conditioned-count benchmark & n/a & \hyperref[thm:p19-r24-counts]{\ref*{thm:p19-r24-counts}} & \thepage & Zero remaining trials yield a degenerate count of zero. \\
\phantomsection\label{entry:symbol-registry-03704}\(N,n\) & \texttt{\detokenize{N;n}} & total objects and recognized-object count & P19 heuristics recognition-mixture benchmark & n/a & \hyperref[def:p19-r24-recognition-law]{\ref*{def:p19-r24-recognition-law}} & \thepage & N is at least two and zero at most n at most N; not the independent-case sample size in the disjoint count benchmark. \\
\phantomsection\label{entry:symbol-registry-03705}\(\alpha,\beta\) & \texttt{\detokenize{\alpha;\beta}} & conditional accuracy of recognition and knowledge-based choice & P19 heuristics recognition-mixture benchmark & n/a & \hyperref[def:p19-r24-recognition-law]{\ref*{def:p19-r24-recognition-law}} & \thepage & Empty pair classes receive arbitrary versions of their conditional parameters with zero weight; comparisons explicitly fix both validities. \\
\phantomsection\label{entry:symbol-registry-03706}\(f(n)\) & \texttt{\detokenize{f}} & expected accuracy at recognition amount n & P19 heuristics recognition-mixture benchmark & n/a & \hyperref[thm:p19-r24-recognition-mixture]{\ref*{thm:p19-r24-recognition-mixture}} & \thepage & Uniform unordered pair selection weights the three pair classes. \\
\phantomsection\label{entry:symbol-registry-03707}\(p,F,t\) & \texttt{\detokenize{p;F;t}} & reported significance bound and F or t test statistic & P19 heuristics empirical source summaries & n/a & \hyperref[def:p19-r24-study-statistics]{\ref*{def:p19-r24-study-statistics}} & \thepage & Reported inferential statistics retain each source's task and design; p is not a Bernoulli or target-category probability. \\
\phantomsection\label{entry:symbol-registry-03786}\(T,K,W,M\) & \texttt{\detokenize{T;K;W;M}} & tap-tried heating-working hot-water and report events & P19 Bayesian report-selection illustration & n/a & \hyperref[def:p19-r45-report-selection]{\ref*{def:p19-r45-report-selection}} & \thepage & W is the intersection T and K; M is generated by declared reporting weights rather than bare observation of no water. \\
\phantomsection\label{entry:symbol-registry-03787}\(P\) & \texttt{\detokenize{P}} & probability law in the stipulated tap/report experiment & P19 Bayesian report-selection illustration & n/a & \hyperref[def:p19-r45-report-selection]{\ref*{def:p19-r45-report-selection}} & \thepage & Independent tap/heater priors and report-selection likelihoods are invented benchmark inputs. \\
\phantomsection\label{entry:symbol-registry-03788}\(A,B\) & \texttt{\detokenize{A;B}} & antecedent and consequent events & P19 Bayesian uncertain-MP coherence & n/a & \hyperref[thm:p19-r45-coherence]{\ref*{thm:p19-r45-coherence}} & \thepage & One event law and stable meanings govern the four cells; not the artist/beekeeper roles in the separate quantifier scope. \\
\phantomsection\label{entry:symbol-registry-03789}\(a,b\) & \texttt{\detokenize{a;b}} & conditional consequent probability and antecedent marginal & P19 Bayesian uncertain-MP coherence & n/a & \hyperref[thm:p19-r45-coherence]{\ref*{thm:p19-r45-coherence}} & \thepage & a is in the unit interval; b is strictly between zero and one for the sharp ratio calculation. \\
\phantomsection\label{entry:symbol-registry-03790}\(P\) & \texttt{\detokenize{P}} & joint event probability law for uncertain MP & P19 Bayesian uncertain-MP coherence & n/a & \hyperref[thm:p19-r45-coherence]{\ref*{thm:p19-r45-coherence}} & \thepage & The material conditional event and conditional-probability assessment remain different. \\
\phantomsection\label{entry:symbol-registry-03791}\(z\) & \texttt{\detokenize{z}} & free nonantecedent/consequent cell mass & P19 Bayesian uncertain-MP coherence & n/a & \hyperref[annex:p19.r45.coherence]{\ref*{annex:p19.r45.coherence}} & \thepage & Ranges from zero through one minus the antecedent marginal and establishes sharpness. \\
\phantomsection\label{entry:symbol-registry-03792}\(P_0,P_1\) & \texttt{\detokenize{P_0;P_1}} & old and Jeffrey-updated event laws & P19 Bayesian finite Jeffrey update & n/a & \hyperref[def:p19-r45-jeffrey]{\ref*{def:p19-r45-jeffrey}} & \thepage & Rigidity retains old conditional laws within a binary partition; dynamic revision is not determined by the new marginal alone. \\
\phantomsection\label{entry:symbol-registry-03793}\(A,B,D\) & \texttt{\detokenize{A;B;D}} & partition event consequent target and arbitrary event & P19 Bayesian finite Jeffrey update & n/a & \hyperref[def:p19-r45-jeffrey]{\ref*{def:p19-r45-jeffrey}} & \thepage & Both old partition cells are nonnull; new null-cell conditionals are not assigned ratio values. \\
\phantomsection\label{entry:symbol-registry-03794}\(t\) & \texttt{\detokenize{t}} & new partition-cell marginal probability & P19 Bayesian finite Jeffrey update & n/a & \hyperref[def:p19-r45-jeffrey]{\ref*{def:p19-r45-jeffrey}} & \thepage & t may be zero or one; this is not an empirical t statistic. \\
\phantomsection\label{entry:symbol-registry-03795}\(u\) & \texttt{\detokenize{u}} & arbitrary new consequent marginal in the nonidentification construction & P19 Bayesian finite Jeffrey update & n/a & \hyperref[annex:p19.r45.jeffrey]{\ref*{annex:p19.r45.jeffrey}} & \thepage & Every unit-interval value is attainable if only the new antecedent marginal is fixed. \\
\phantomsection\label{entry:symbol-registry-03796}\(A,B,C\) & \texttt{\detokenize{A;B;C}} & artist beekeeper and chemist events & P19 Bayesian quantifier-support bridge & n/a & \hyperref[thm:p19-r45-quantifiers]{\ref*{thm:p19-r45-quantifiers}} & \thepage & Local predicate interpretation differs from antecedent/consequent events in the coherence scope. \\
\phantomsection\label{entry:symbol-registry-03797}\(P\) & \texttt{\detokenize{P}} & probability measure on the stipulated object domain & P19 Bayesian quantifier-support bridge & n/a & \hyperref[thm:p19-r45-quantifiers]{\ref*{thm:p19-r45-quantifiers}} & \thepage & Positive mass and existential/nonexception support must not be conflated when zero-mass objects are permitted. \\
\phantomsection\label{entry:symbol-registry-03798}\(\Delta\) & \texttt{\detokenize{\Delta}} & declared few/most threshold & P19 Bayesian quantifier-support bridge & n/a & \hyperref[sec:p19-r45-quantifiers]{\ref*{sec:p19-r45-quantifiers}} & \thepage & Strictly between zero and one-half; most lies strictly between one-minus-Delta and one in this model version. \\
\phantomsection\label{entry:symbol-registry-03799}\(I(s)\) & \texttt{\detokenize{I}} & inverse-probability informativeness score & P19 Bayesian quantifier-information comparison & n/a & \hyperref[sec:p19-r45-quantifiers]{\ref*{sec:p19-r45-quantifiers}} & \thepage & The 2021 printed score 1/P(s) is distinguished from Shannon surprisal minus log-base-two P(s). \\
\phantomsection\label{entry:symbol-registry-03800}\(s,P(s)\) & \texttt{\detokenize{s;P}} & quantified statement and its assigned probability & P19 Bayesian quantifier-information comparison & n/a & \hyperref[sec:p19-r45-quantifiers]{\ref*{sec:p19-r45-quantifiers}} & \thepage & The score comparison requires positive probability and a supplied statement-probability model. \\
\phantomsection\label{entry:symbol-registry-03801}\(r,n\) & \texttt{\detokenize{r;n}} & prototype counterexample probability and number of independent draws & P19 Bayesian prototype-sampling boundary & n/a & \hyperref[thm:p19-r45-sample]{\ref*{thm:p19-r45-sample}} & \thepage & n is a positive integer; these are stipulated inputs rather than fitted or measured participant memory values. \\
\phantomsection\label{entry:symbol-registry-03802}\(h,g\) & \texttt{\detokenize{h;g}} & finite candidate-hypothesis indices & P19 Bayesian description-length bridge & n/a & \hyperref[thm:p19-r45-description]{\ref*{thm:p19-r45-description}} & \thepage & The candidate set and coding language are declared. \\
\phantomsection\label{entry:symbol-registry-03803}\(L(h)\) & \texttt{\detokenize{L}} & finite nonnegative code length & P19 Bayesian description-length bridge & n/a & \hyperref[thm:p19-r45-description]{\ref*{thm:p19-r45-description}} & \thepage & Measured in bits for the specified representation; not automatically the number of causes. \\
\phantomsection\label{entry:symbol-registry-03804}\(\pi\) & \texttt{\detokenize{\pi}} & hypothesis probability weights and their evidence-conditioned version & P19 Bayesian description-length bridge & n/a & \hyperref[thm:p19-r45-description]{\ref*{thm:p19-r45-description}} & \thepage & Prior normalization uses code lengths; posterior weights additionally use likelihoods. \\
\phantomsection\label{entry:symbol-registry-03805}\(Z\) & \texttt{\detokenize{Z}} & normalizer for description-length prior weights & P19 Bayesian description-length bridge & n/a & \hyperref[thm:p19-r45-description]{\ref*{thm:p19-r45-description}} & \thepage & Finite and strictly positive on a nonempty finite candidate set. \\
\phantomsection\label{entry:symbol-registry-03806}\(\ell(h)\) & \texttt{\detokenize{\ell}} & evidence-event likelihood under a candidate hypothesis & P19 Bayesian description-length bridge & n/a & \hyperref[thm:p19-r45-description]{\ref*{thm:p19-r45-description}} & \thepage & Zero likelihood gives posterior zero and infinite negative-log objective. \\
\phantomsection\label{entry:symbol-registry-03807}\(D,m\) & \texttt{\detokenize{D;m}} & evidence event and its prior-predictive probability & P19 Bayesian description-length bridge & n/a & \hyperref[thm:p19-r45-description]{\ref*{thm:p19-r45-description}} & \thepage & m is strictly positive in the posterior derivation. \\
\phantomsection\label{entry:symbol-registry-03808}\(H,h\) & \texttt{\detokenize{H;h}} & random hypothesis and its finite values & P19 Bayesian finite information experiment & n/a & \hyperref[def:p19-r45-information]{\ref*{def:p19-r45-information}} & \thepage & Hypothesis priors are positive; no unlisted alternative is supplied by the entropy formula. \\
\phantomsection\label{entry:symbol-registry-03809}\(a,y\) & \texttt{\detokenize{a;y}} & available test index and finite test outcome & P19 Bayesian finite information experiment & n/a & \hyperref[def:p19-r45-information]{\ref*{def:p19-r45-information}} & \thepage & a denotes a prospective test rather than the conditional assessment in the coherence scope. \\
\phantomsection\label{entry:symbol-registry-03810}\(\pi\) & \texttt{\detokenize{\pi}} & positive hypothesis-prior mass function & P19 Bayesian finite information experiment & n/a & \hyperref[def:p19-r45-information]{\ref*{def:p19-r45-information}} & \thepage & Fixed while comparing tests under a specified experiment. \\
\phantomsection\label{entry:symbol-registry-03811}\(k_a(y\mid h)\) & \texttt{\detokenize{k_a}} & test outcome likelihood kernel & P19 Bayesian finite information experiment & n/a & \hyperref[def:p19-r45-information]{\ref*{def:p19-r45-information}} & \thepage & Nonnegative and normalized over outcomes for each hypothesis. \\
\phantomsection\label{entry:symbol-registry-03812}\(m_a(y)\) & \texttt{\detokenize{m_a}} & prior-predictive test-outcome mass & P19 Bayesian finite information experiment & n/a & \hyperref[def:p19-r45-information]{\ref*{def:p19-r45-information}} & \thepage & Zero-marginal outcomes are excluded from posterior/conditional-entropy calculations. \\
\phantomsection\label{entry:symbol-registry-03813}\(\pi_a(h\mid y)\) & \texttt{\detokenize{\pi_a}} & Bayes posterior after a positive-marginal test outcome & P19 Bayesian finite information experiment & n/a & \hyperref[def:p19-r45-information]{\ref*{def:p19-r45-information}} & \thepage & A realized outcome may increase entropy even when expected gain is nonnegative. \\
\phantomsection\label{entry:symbol-registry-03814}\(S\) & \texttt{\detokenize{S}} & base-two finite entropy functional & P19 Bayesian finite information experiment & n/a & \hyperref[def:p19-r45-information]{\ref*{def:p19-r45-information}} & \thepage & Uses zero-log-zero convention; not the inverse-probability score I in the quantifier discussion. \\
\phantomsection\label{entry:symbol-registry-03815}\(G(a)\) & \texttt{\detokenize{G}} & expected information gain of a test & P19 Bayesian finite information experiment & n/a & \hyperref[def:p19-r45-information]{\ref*{def:p19-r45-information}} & \thepage & Measured in bits; test cost and downstream utility are not automatically included. \\
\phantomsection\label{entry:symbol-registry-03816}\(J,R\) & \texttt{\detokenize{J;R}} & joint and marginal-product probability laws & P19 Bayesian finite information experiment & n/a & \hyperref[annex:p19.r45.information]{\ref*{annex:p19.r45.information}} & \thepage & Their finite relative-entropy identity proves nonnegativity and its equality case. \\
\phantomsection\label{entry:symbol-registry-03817}\(b,c,q_1,q_0\) & \texttt{\detokenize{b;c;q_1;q_0}} & action benefit consequence disutility and action-specific consequence probabilities & P19 Bayesian warning-decision bridge & n/a & \hyperref[thm:p19-r45-slope]{\ref*{thm:p19-r45-slope}} & \thepage & The baseline q\_0 is not omitted; the planning laws are not identified by observational association alone. \\
\phantomsection\label{entry:symbol-registry-03818}\(P,Q,A\) & \texttt{\detokenize{P;Q;A}} & observational probability law consequence event and action-associated event & P19 Bayesian warning-decision bridge & n/a & \hyperref[thm:p19-r45-slope]{\ref*{thm:p19-r45-slope}} & \thepage & Observational P(Q given A) is distinguished from the stipulated action-specific planning probabilities. \\
\phantomsection\label{entry:symbol-registry-03819}\(F,p,t\) & \texttt{\detokenize{F;p;t}} & reported F statistic significance value and t statistic & P19 Bayesian empirical study paragraphs & n/a & \hyperref[def:p19-r45-study-statistics]{\ref*{def:p19-r45-study-statistics}} & \thepage & Restricted to explicitly cited empirical result paragraphs; excludes the formal probability parameters and Jeffrey weight. \\
\phantomsection\label{entry:symbol-registry-03929}\(S_1,S_2,X,\alpha\) & \texttt{\detokenize{S_1;S_2;X;alpha}} & cognitive predicates and content group and selected subgroup & P19 OBrien2021 natural-logic representation & n/a & \hyperref[sec:p19-r32-representation]{\ref*{sec:p19-r32-representation}} & \thepage & Not type-level S or a formula translation; pointwise formal benchmark is declared separately. \\
\phantomsection\label{entry:symbol-registry-03930}\(\mathrm{All},\mathrm{OR},\mathrm{PRO\!-!All}\) & \texttt{\detokenize{mathrm}} & cognitive quantifier and alternative and anaphoric scope annotation & P19 OBrien2021 natural-logic representation & n/a & \hyperref[sec:p19-r32-representation]{\ref*{sec:p19-r32-representation}} & \thepage & Source-native cognitive roles; not an unannounced classical connective or full formal calculus. \\
\phantomsection\label{entry:symbol-registry-03931}\(Z,G,A,X,P,Q\) & \texttt{\detokenize{Z;G;A;X;P;Q}} & unary predicates in stipulated quantified benchmarks & P19 natural-logic formal quantifier checks & n/a & \hyperref[thm:p19-r32-quantifier-benchmarks]{\ref*{thm:p19-r32-quantifier-benchmarks}} & \thepage & Distinct from source cognitive expressions and from sentential atoms in the separate chain. \\
\phantomsection\label{entry:symbol-registry-03932}\(N,P,H,Z,S,p,q\) & \texttt{\detokenize{N;P;H;Z;S;p;q}} & Boolean atoms in stipulated sentential checks & P19 natural-logic sentential benchmarks & n/a & \hyperref[thm:p19-r32-sentential-benchmarks]{\ref*{thm:p19-r32-sentential-benchmarks}} & \thepage & Source chain uses letter names; does not assert material truth-function semantics for ordinary-language if. \\
\phantomsection\label{entry:symbol-registry-03933}\(a,b,c\) & \texttt{\detokenize{a;b;c}} & unary task properties with middle property b & P19 r33 categorical task semantics & n/a & \hyperref[def:p19-r33-categorical-task]{\ref*{def:p19-r33-categorical-task}} & \thepage & End properties a and c; lower-case names do not denote probabilities in this scope. \\
\phantomsection\label{entry:symbol-registry-03934}\(U,U'\) & \texttt{\detokenize{U;U'}} & nonempty universes in categorical examples and countermodels & P19 r33 categorical task semantics & n/a & \hyperref[def:p19-r33-categorical-task]{\ref*{def:p19-r33-categorical-task}} & \thepage & Prime marks the witness-splitting interpretation; at most eight complete property types suffice. \\
\phantomsection\label{entry:symbol-registry-03935}\(a^U,b^U,c^U\) & \texttt{\detokenize{a^U;b^U;c^U}} & property extensions as subsets of the named universe & P19 r33 categorical task semantics & n/a & \hyperref[def:p19-r33-categorical-task]{\ref*{def:p19-r33-categorical-task}} & \thepage & Empty extensions allowed except where explicit all-three-class existence conditions are added. \\
\phantomsection\label{entry:symbol-registry-03936}\(A,I,E,O\) & \texttt{\detokenize{A;I;E;O}} & universal affirmative particular affirmative universal negative particular negative mood codes & P19 r33 categorical task semantics & n/a & \hyperref[def:p19-r33-categorical-task]{\ref*{def:p19-r33-categorical-task}} & \thepage & Aab and other two-term strings are sentence codes; capital letters are not property extensions. \\
\phantomsection\label{entry:symbol-registry-03937}\(\mathrm{NVC}\) & \texttt{\detokenize{\mathrm{NVC}}} & no licensed conclusion among eight requested end-term forms & P19 r33 response categories & n/a & \hyperref[def:p19-r33-existence]{\ref*{def:p19-r33-existence}} & \thepage & Does not deny all logical consequence; also distinguish a stochastic NVC output from a proved NVC benchmark. \\
\phantomsection\label{entry:symbol-registry-03938}\(N\) & \texttt{\detokenize{N}} & nominal dataset sample count & P19 r33 meta-analysis provenance table & n/a & \hyperref[tab:p19-r33-meta-samples]{\ref*{tab:p19-r33-meta-samples}} & \thepage & Twenty first-session and twenty second-session records reuse the same participants. \\
\phantomsection\label{entry:symbol-registry-03939}\(I^2,Q,p\) & \texttt{\detokenize{I^2;Q;p}} & reported heterogeneity measure Cochran statistic and tail probability & P19 r33 empirical meta-analysis record & n/a & \hyperref[sec:p19-r33-meta]{\ref*{sec:p19-r33-meta}} & \thepage & The I in I-squared is not the particular-affirmative task mood; empirical quantities retain reported model and metric. \\
\phantomsection\label{entry:symbol-registry-03940}\(S,M,P\) & \texttt{\detokenize{S;M;P}} & subject middle and predicate membership events & P19 r33 Hattori prototype scope & n/a & \hyperref[def:p19-r33-psm-procedure]{\ref*{def:p19-r33-psm-procedure}} & \thepage & Local event P is distinguished from the probability operator in P(M); properties a b c belong to a different declared scope. \\
\phantomsection\label{entry:symbol-registry-03941}\(P(\cdot)\) & \texttt{\detokenize{P(\cdot)}} & prototype probability operator & P19 r33 Hattori prototype scope & n/a & \hyperref[def:p19-r33-psm-procedure]{\ref*{def:p19-r33-psm-procedure}} & \thepage & Not the third membership event P; explicit probability-equipped representation assumptions rather than empirical response frequencies. \\
\phantomsection\label{entry:symbol-registry-03942}\(x,c\) & \texttt{\detokenize{x;c}} & middle-event occurrence and coverage parameters & P19 r33 Hattori prototype scope & n/a & \hyperref[def:p19-r33-psm-procedure]{\ref*{def:p19-r33-psm-procedure}} & \thepage & c is not the end property c of the task-semantic scope; displayed figure rounds cx-squared. \\
\phantomsection\label{entry:symbol-registry-03943}\(n\) & \texttt{\detokenize{n}} & finite independent draw count & P19 r33 Hattori prototype and sampling-boundary scope & n/a & \hyperref[def:p19-r33-psm-procedure]{\ref*{def:p19-r33-psm-procedure}} & \thepage & Not sample participant N; draws may realize identical membership regions. \\
\phantomsection\label{entry:symbol-registry-03944}\(u_A,u_I,u_E,u_O\) & \texttt{\detokenize{u_A;u_I;u_E;u_O}} & max-heuristic acceptance probabilities & P19 r33 Hattori prototype scope & n/a & \hyperref[def:p19-r33-psm-procedure]{\ref*{def:p19-r33-psm-procedure}} & \thepage & Gates in original Figure8 can produce NVC even though full ungated categorical tests cannot. \\
\phantomsection\label{entry:symbol-registry-03945}\(q\) & \texttt{\detokenize{q}} & prototype mass of sampled overlap & P19 r33 finite-sample boundary & n/a & \hyperref[thm:p19-r33-sampling]{\ref*{thm:p19-r33-sampling}} & \thepage & The independent miss probability is one minus q raised to n; not a human response proportion. \\
\phantomsection\label{entry:symbol-registry-03946}\(\lambda,\varepsilon,\sigma,\omega\) & \texttt{\detokenize{\lambda;\varepsilon;\sigma;\omega}} & size-distribution fuller-possibility search-engagement and conditional-reconclusion parameters & P19 r33 mReasoner2016 implementation & n/a & \hyperref[def:p19-r33-mreasoner-parameters]{\ref*{def:p19-r33-mreasoner-parameters}} & \thepage & Four local model roles; later journal version uses named counterparts and explicitly truncated Poisson size distribution. \\
\phantomsection\label{entry:symbol-registry-03947}\(z\) & \texttt{\detokenize{z}} & chosen existential witness & P19 r33 witness annex & n/a & \hyperref[annex:p19.r33.witness]{\ref*{annex:p19.r33.witness}} & \thepage & Same witness traverses inclusion; no fresh witness is created from an unconditioned universal premise. \\
\phantomsection\label{entry:symbol-registry-03948}\(d,m,w,e,x,y\) & \texttt{\detokenize{d;m;w;e;x;y}} & named individuals in complete counterinterpretations & P19 r33 formal countermodel annexes & n/a & \hyperref[annex:p19.r33.witness]{\ref*{annex:p19.r33.witness}} & \thepage & All memberships given by extensions; diagram blanks do not silently supply negative literals. \\
\phantomsection\label{entry:symbol-registry-03949}\(A,C\) & \texttt{\detokenize{A;C}} & antecedent and consequent events or source conditional clauses with the interpretation explicitly declared & \hyperref[sec:p19-rationality-4-6]{\ref*{sec:p19-rationality-4-6}} & meta & \hyperref[sec:p19-rat46-tasks]{\ref*{sec:p19-rat46-tasks}} & \thepage & The material comparison is written neg A or C; the source arrow is not silently treated as that event. \\
\phantomsection\label{entry:symbol-registry-03950}\(P(A\to C)\) & \texttt{\detokenize{P(A\to C)}} & source probability-conditional assessment in the Equation & \hyperref[sec:p19-rationality-4-6]{\ref*{sec:p19-rationality-4-6}} & meta & \hyperref[def:p19-rat46-equation]{\ref*{def:p19-rat46-equation}} & \thepage & Does not automatically denote probability of a fixed Boolean event. \\
\phantomsection\label{entry:symbol-registry-03951}\(a,b,e\) & \texttt{\detokenize{a;b;e}} & source joint-probability and exceptions parameters & \hyperref[sec:p19-rationality-4-6]{\ref*{sec:p19-rationality-4-6}} & meta & \hyperref[def:p19-rat46-joint]{\ref*{def:p19-rat46-joint}} & \thepage & a=P(A), b=P(C), e=1-P(C given A); positive denominators and feasibility constraints declared. \\
\phantomsection\label{entry:symbol-registry-03952}\(s,r,q,d_{11},d_{10},d_{01},d_{00}\) & \texttt{\detokenize{s;r;q;d_{11};d_{10};d_{01};d_{00}}} & temporary probability parameters and fixed-event atom membership coefficients & \hyperref[annex:p19.rat46.fixed-event]{\ref*{annex:p19.rat46.fixed-event}} & meta & \hyperref[annex:p19.rat46.fixed-event]{\ref*{annex:p19.rat46.fixed-event}} & \thepage & Three parameters lie in (0,1); each membership coefficient is zero or one. \\
\phantomsection\label{entry:symbol-registry-03953}Type 1; Type 2 & \texttt{\detokenize{Type~1;Type~2}} & source cognitive processing categories & P19 RAT2021 chapter 2.5 conceptual account & n/a & \hyperref[def:p19-r25-process-account]{\ref*{def:p19-r25-process-account}} & \thepage & Categories require an explicit account; neither speed nor normative correctness is adopted as a definition. \\
\phantomsection\label{entry:symbol-registry-03954}\(A,B,3,1\) & \texttt{\detokenize{A;B;3;1}} & visible card labels in the source selection task & P19 RAT2021 chapter 2.5 selection example & n/a & \hyperref[def:p19-r25-task-examples]{\ref*{def:p19-r25-task-examples}} & \thepage & Capital letters on one side and digits on the other; not probability events or numerical truth values. \\
\phantomsection\label{entry:symbol-registry-03955}\(p,q\) & \texttt{\detokenize{p;q}} & selection-rule propositional abbreviations & P19 RAT2021 chapter 2.5 selection proof & n/a & \hyperref[thm:p19-r25-benchmarks]{\ref*{thm:p19-r25-benchmarks}} & \thepage & p means card letter A; q means card digit3; p-to-q is the classical benchmark. Distinct from the probability function in the later odds paragraph. \\
\phantomsection\label{entry:symbol-registry-03956}\(c,h\) & \texttt{\detokenize{c;h}} & crying and happiness propositions in the stipulated-premise proof & P19 RAT2021 chapter 2.5 belief-task proof & n/a & \hyperref[thm:p19-r25-benchmarks]{\ref*{thm:p19-r25-benchmarks}} & \thepage & The cognitive process producing an answer is not established by the validity of this modus-ponens benchmark. \\
\phantomsection\label{entry:symbol-registry-03957}\(p(\cdot),D,\ell_P,\ell_N\) & \texttt{\detokenize{p;D;\ell_P;\ell_N}} & probability description event and two category likelihoods & P19 RAT2021 chapter 2.5 posterior comparison & n/a & \hyperref[thm:p19-r25-benchmarks]{\ref*{thm:p19-r25-benchmarks}} & \thepage & Finite Bayesian comparison uses stipulated positive likelihoods; the source figure gives category counts but not a likelihood ratio. \\
\phantomsection\label{entry:symbol-registry-03958}\(F,\eta_p^2,p_{\mathrm{rep}}\) & \texttt{\detokenize{F;\eta_p^2;p_{\mathrm{rep}}}} & reported ANOVA statistic partial eta squared and source inference convention & P19 RAT2021 chapter 2.5 DeNeys primary record & n/a & \hyperref[def:p19-r25-deneys-record]{\ref*{def:p19-r25-deneys-record}} & \thepage & p\_rep is not silently relabeled an ordinary p-value; no confidence intervals are manufactured. \\
\phantomsection\label{entry:symbol-registry-03959}FOR & \texttt{\detokenize{FOR}} & initial feeling-of-rightness report & P19 RAT2021 chapter 2.5 reviewed two-response studies & n/a & \hyperref[sec:p19-rationality-2-5]{\ref*{sec:p19-rationality-2-5}} & \thepage & A metacognitive rating associated with later change/time; not a directly observed processor or automatically a probability scale. \\
\phantomsection\label{entry:symbol-registry-03971}\((p,q,r)\) & \texttt{\detokenize{p,}\allowbreak\detokenize{q,}\allowbreak\detokenize{r}} & stipulated antecedent consequence and disabling-circumstance atoms & P19 rationality 5.4 & object & \hyperref[thm:p19-r54-disabler-boundary]{\ref*{thm:p19-r54-disabler-boundary}} & \thepage & These finite formal examples are not measured participant representations. \\
\phantomsection\label{entry:symbol-registry-03972}\(\Pr(q\mid p),\Pr(p\mid q)\) & \texttt{\detokenize{\Pr(q\mid p),}\allowbreak\detokenize{\Pr(p\mid q)}} & opposite conditional probabilities in the specified explanatory calculation & P19 rationality 5.4 & n/a & \hyperref[thm:p19-r54-probability-boundary]{\ref*{thm:p19-r54-probability-boundary}} & \thepage & Each ratio requires its positive conditioning denominator; response proportions are not automatically these probabilities. \\
\phantomsection\label{entry:symbol-registry-03973}\(\operatorname{Cn}(\Gamma)\) & \texttt{\detokenize{\operatorname{Cn}(\Gamma)}} & classical deductive closure of a finite-language formula set & P19 rationality 5.4 & object & \hyperref[thm:p19-r54-two-revisions]{\ref*{thm:p19-r54-two-revisions}} & \thepage & Distinguished from a participant's selected sentence or graded premise credibility. \\
\phantomsection\label{entry:symbol-registry-03974}\(\operatorname{Th}(X)\) & \texttt{\detokenize{\operatorname{Th}(X)}} & theory true at every member of a stipulated valuation set & P19 rationality 5.4 & object & \hyperref[thm:p19-r54-two-revisions]{\ref*{thm:p19-r54-two-revisions}} & \thepage & The theory of the empty valuation set is inconsistent; the finite proof gives all eight postulates for one fixed prior. \\
\phantomsection\label{entry:symbol-registry-03975}\(K,R_A,m_A,[A]\) & \texttt{\detokenize{K,}\allowbreak\detokenize{R_A,}\allowbreak\detokenize{m_A,}\allowbreak\detokenize{[A]}} & fixed prior theory revised theory selected input models and all input models & P19 rationality 5.4 & object & \hyperref[annex:p19.r54.two-revisions]{\ref*{annex:p19.r54.two-revisions}} & \thepage & The two declared total orders produce opposite repairs while satisfying the eight fixed-prior AGM conditions. \\
\phantomsection\label{entry:symbol-registry-03976}\(\prec\) & \texttt{\detokenize{\prec}} & one of two stipulated strict total orders on four valuations & P19 rationality 5.4 & relation & \hyperref[thm:p19-r54-two-revisions]{\ref*{thm:p19-r54-two-revisions}} & \thepage & The sole K-model is first in both; the differing remaining order is not estimated from the empirical samples. \\
\phantomsection\label{entry:symbol-registry-03977}\((LO,RO)\) & \texttt{\detokenize{LO,}\allowbreak\detokenize{RO}} & located and reference object roles in the spatial revision task & P19 rationality 5.4 & object & \hyperref[sec:p19-r54-revision-evidence]{\ref*{sec:p19-r54-revision-evidence}} & \thepage & The recorded Experiment 1 measures diagram choices among surviving trials; it is not unrestricted physical navigation. \\
\phantomsection\label{entry:symbol-registry-04059}\(N,a,c,v,u\) & \texttt{\detokenize{N,}\allowbreak\detokenize{a,}\allowbreak\detokenize{c,}\allowbreak\detokenize{v,}\allowbreak\detokenize{u}} & finite equally weighted task count with answered/correct counts and conditional correctness/applicability proportions & P19 rationality1.3 cognitive-success calculation & n/a & \hyperref[thm:p19-r13-success]{\ref*{thm:p19-r13-success}} & \thepage & Validity here is correctness among rendered inferences; a=0 has no defined conditional validity without another convention. \\
\phantomsection\label{entry:symbol-registry-04060}\(x,\epsilon,q\) & \texttt{\detokenize{x,}\allowbreak\detokenize{\epsilon,}\allowbreak\detokenize{q}} & stag-population proportion and invasion/mutant-mixture proportions in a stipulated symmetric game & P19 rationality1.3 evolutionary-stability calculation & n/a & \hyperref[thm:p19-r13-stability]{\ref*{thm:p19-r13-stability}} & \thepage & These are mathematical population parameters not measured human frequencies. \\
\phantomsection\label{entry:symbol-registry-04061}\(L_T,L_T^*\) & \texttt{\detokenize{L_T,}\allowbreak\detokenize{L_T^*}} & illustrative cumulative loss and comparator loss over T occasions & P19 rationality1.3 meta-induction boundary & n/a & \hyperref[thm:p19-r13-horizon]{\ref*{thm:p19-r13-horizon}} & \thepage & The convergence example is not a newly attributed algorithm or the full source optimality theorem. \\
\phantomsection\label{entry:symbol-registry-04062}\(\epsilon\) & \texttt{\detokenize{\epsilon}} & exceptional probability mass in a finite conditional countermodel & P19 rationality1.3 probability-logic comparison & n/a & \hyperref[thm:p19-r13-probabilistic-mt]{\ref*{thm:p19-r13-probabilistic-mt}} & \thepage & Distinct from the invasion proportion in the earlier stipulated game. \\
\phantomsection\label{entry:symbol-registry-04063}\(P(\text{if }A\text{ then }B)\) & \texttt{\detokenize{P(\text{if }A\text{ then }B)}} & assessment of an ordinary-language singular indicative conditional & \hyperref[sec:p19-rationality-6-2]{\ref*{sec:p19-rationality-6-2}} & meta & \hyperref[def:p19-r62-equation]{\ref*{def:p19-r62-equation}} & \thepage & Not assumed to be the probability of a fixed Boolean event. \\
\phantomsection\label{entry:symbol-registry-04064}\(A,B\) & \texttt{\detokenize{A;B}} & antecedent and consequent events in explicit finite comparisons & \hyperref[sec:p19-rationality-6-2]{\ref*{sec:p19-rationality-6-2}} & meta & \hyperref[def:p19-r62-equation]{\ref*{def:p19-r62-equation}} & \thepage & English clauses and material events are distinguished before formal calculations. \\
\phantomsection\label{entry:symbol-registry-04065}\(V\) & \texttt{\detokenize{V}} & void factual assessment in the source's de Finetti table & \hyperref[sec:p19-rationality-6-2]{\ref*{sec:p19-rationality-6-2}} & meta & \hyperref[def:p19-r62-values]{\ref*{def:p19-r62-values}} & \thepage & Distinct from falsity and from the numeric subjective value in the Jeffrey table. \\
\phantomsection\label{entry:symbol-registry-04066}\(a,b,r,x,z,v\) & \texttt{\detokenize{a;b;r;x;z;v}} & local finite probability parameters or atom masses & sec:p19-rationality-6-2;sec:annex-p19-rationality-6-2 & meta & \hyperref[thm:p19-r62-material]{\ref*{thm:p19-r62-material}} & \thepage & Each theorem and annex declares its marginal, joint or conditional role; no measured response counts. \\
\phantomsection\label{entry:symbol-registry-04067}\(X_r,E_P\) & \texttt{\detokenize{X_r;E_P}} & Jeffrey-table random quantity and finite expectation under P & \hyperref[sec:p19-rationality-6-2]{\ref*{sec:p19-rationality-6-2}} & meta & \hyperref[thm:p19-r62-jeffrey-value]{\ref*{thm:p19-r62-jeffrey-value}} & \thepage & Expectation fixed point is underdetermined at a null antecedent; not a general coherence theorem. \\
\phantomsection\label{entry:symbol-registry-04068}\(S,Q,L,E\) & \texttt{\detokenize{S;Q;L;E}} & short-route quickness and long-route enjoyment events & sec:p19-r62-supposition;sec:p19-r62-coherence & meta & \hyperref[sec:p19-r62-supposition]{\ref*{sec:p19-r62-supposition}} & \thepage & The event E is enjoyment; E\_P in a different explicitly typed expression denotes expectation. \\
\phantomsection\label{entry:symbol-registry-04069}\(C,N\) & \texttt{\detokenize{C;N}} & Canberra sunshine and London rain events & \hyperref[sec:p19-r62-centering]{\ref*{sec:p19-r62-centering}} & meta & \hyperref[def:p19-r62-missing-link]{\ref*{def:p19-r62-missing-link}} & \thepage & Keeps the source's defined London symbol N; source p.400 has one stray C-and-L occurrence. \\
\phantomsection\label{entry:symbol-registry-04070}\(P_1,P_2,F,H,t,q,s\) & \texttt{\detokenize{P_1;P_2;F;H;t;q;s}} & prior and posterior laws, finding and hypothesis events, update and prior cell parameters & sec:p19-r62-updating;annex:p19.r62.update & meta & \hyperref[def:p19-r62-invariance]{\ref*{def:p19-r62-invariance}} & \thepage & Time indices are not different assessments in one fixed law; rigidity is stipulated, not automatic. \\
\phantomsection\label{entry:symbol-registry-04071}\(F,p\) & \texttt{\detokenize{F;p}} & reported ANOVA and significance statistics in the original-study records & \hyperref[sec:p19-r62-explicit-inference]{\ref*{sec:p19-r62-explicit-inference}} & meta & \hyperref[sec:p19-r62-explicit-inference]{\ref*{sec:p19-r62-explicit-inference}} & \thepage & Statistical F is not the finding event of the later update model; significance p is not a model probability. \\
\phantomsection\label{entry:symbol-registry-04072}\(A,B,C\) & \texttt{\detokenize{A;B;C}} & temporary schematic category or ordered-item names in declared deductive task comparisons & \hyperref[sec:p19-rationality-1-4]{\ref*{sec:p19-rationality-1-4}} & meta & \hyperref[sec:p19-r14-logical-forms]{\ref*{sec:p19-r14-logical-forms}} & \thepage & Not brain states; inclusion and strict-ordering interpretations stated separately. \\
\phantomsection\label{entry:symbol-registry-04073}\(F_{2,33},F_{1,8},F_{2,16}\) & \texttt{\detokenize{F}} & reported source analysis-of-variance/covariance statistics with numerator and denominator degrees of freedom & \hyperref[sec:p19-rationality-1-4]{\ref*{sec:p19-rationality-1-4}} & meta & \hyperref[sec:p19-r14-lesion-record]{\ref*{sec:p19-r14-lesion-record}} & \thepage & Observed published statistics, not newly calculated estimates; distinct analysis samples and covariates retained. \\
\phantomsection\label{entry:symbol-registry-04074}\(p\) & \texttt{\detokenize{p}} & reported statistical tail-probability notation in the cited imaging and lesion analyses & \hyperref[sec:p19-rationality-1-4]{\ref*{sec:p19-rationality-1-4}} & meta & \hyperref[sec:p19-r14-meta]{\ref*{sec:p19-r14-meta}} & \thepage & FWE correction applies to the imaging meta-analysis; not posterior probability of a process or hypothesis. \\
\phantomsection\label{entry:symbol-registry-04075}\(A,B\) & \texttt{\detokenize{A,}\allowbreak\detokenize{B}} & antecedent and consequent propositions or corresponding events & P19 rationality-6-3 & n/a & \hyperref[def:p19-rat63-standing]{\ref*{def:p19-rat63-standing}} & \thepage & The active comparison declares whether a proposition or a probability event is intended. \\
\phantomsection\label{entry:symbol-registry-04076}\(AB,\neg A\neg B,\neg AB,A\neg B\) & \texttt{\detokenize{AB,}\allowbreak\detokenize{\neg A\neg B,}\allowbreak\detokenize{\neg AB,}\allowbreak\detokenize{A\neg B}} & schematic conjunctions of two propositional states & P19 rationality-6-3 & n/a & \hyperref[def:p19-rat63-standing]{\ref*{def:p19-rat63-standing}} & \thepage & Juxtaposition means conjunction; an unspecified initial-model component is not a negative assertion. \\
\phantomsection\label{entry:symbol-registry-04077}\(q\) & \texttt{\detokenize{q}} & conditional-probability assessment of B given not-A & P19 rationality-6-3 probability comparison & n/a & \hyperref[thm:p19-rat63-or-if]{\ref*{thm:p19-rat63-or-if}} & \thepage & Requires positive probability of not-A; not the probability of a material-implication event. \\
\phantomsection\label{entry:symbol-registry-04078}\(P_0,Q_0\) & \texttt{\detokenize{P_0,}\allowbreak\detokenize{Q_0}} & specified formulas in a compatibility counterexample & P19 rationality-6-3 compatibility & n/a & \hyperref[thm:p19-rat63-compatibility]{\ref*{thm:p19-rat63-compatibility}} & \thepage & P\_0 is the inclusive disjunction and Q\_0 the material conditional. \\
\phantomsection\label{entry:symbol-registry-04079}\(\mathcal T(\varphi)\) & \texttt{\detokenize{\mathcal T}} & set of two-proposition truth assignments satisfying a formula & P19 rationality-6-3 compatibility & n/a & \hyperref[thm:p19-rat63-compatibility]{\ref*{thm:p19-rat63-compatibility}} & \thepage & Nonempty intersection expresses compatibility; subset inclusion expresses entailment. \\
\phantomsection\label{entry:symbol-registry-04080}\(C[B=0]\) & \texttt{\detokenize{C[B=0]}} & counterfactual value after replacing B's equation & P19 rationality-6-3 world typing & n/a & \hyperref[def:p19-rat63-world-typing]{\ref*{def:p19-rat63-world-typing}} & \thepage & Reuses P16 same-background semantics; evidence B=1 remains factual. \\
\phantomsection\label{entry:symbol-registry-04081}\(U\) & \texttt{\detokenize{U}} & binary background unit in the stipulated causal illustration & P19 rationality-6-3 world typing & n/a & \hyperref[thm:p19-rat63-observe-intervene]{\ref*{thm:p19-rat63-observe-intervene}} & \thepage & Two values each have probability one-half; not an empirical participant model. \\
\phantomsection\label{entry:symbol-registry-04082}\(B,C\) & \texttt{\detokenize{B,}\allowbreak\detokenize{C}} & endogenous variables in the stipulated causal illustration & P19 rationality-6-3 world typing & n/a & \hyperref[thm:p19-rat63-observe-intervene]{\ref*{thm:p19-rat63-observe-intervene}} & \thepage & Locally typed variables rather than antecedent/consequent events or belief/common-knowledge operators. \\
\phantomsection\label{entry:symbol-registry-04083}\(R,P,L,Z\) & \texttt{\detokenize{R,}\allowbreak\detokenize{P,}\allowbreak\detokenize{L,}\allowbreak\detokenize{Z}} & rose poppy lily and unnamed category events & P19 rationality-6-3 category-complement annex & n/a & \hyperref[annex:p19.rat63-alternates]{\ref*{annex:p19.rat63-alternates}} & \thepage & Mutual exclusivity and exhaustiveness are explicit local assumptions; P here is not the probability function. \\
\phantomsection\label{entry:symbol-registry-04084}\(S\) & \texttt{\detokenize{S}} & single-category domain for the complement comparison & P19 rationality-6-3 category-complement annex & n/a & \hyperref[annex:p19.rat63-alternates]{\ref*{annex:p19.rat63-alternates}} & \thepage & Binary or three-category exhaustive domain; an unnamed category supplies the nonexhaustive countercase. \\
\phantomsection\label{entry:symbol-registry-04095}\(p,q\) & \texttt{\detokenize{p;q}} & antecedent action and consequent clause of a utility conditional & \hyperref[sec:p19-rationality-6-4]{\ref*{sec:p19-rationality-6-4}} & meta & \hyperref[def:p19-r64-utility-conditional]{\ref*{def:p19-r64-utility-conditional}} & \thepage & English clauses are not assumed automatically to be material propositions; significance p is separately typed in study reports. \\
\phantomsection\label{entry:symbol-registry-04096}\(a_1,a_0,r_1,r_0,v,c_1,c_0,U_1,U_0\) & \texttt{\detokenize{a_1;a_0;r_1;r_0;v;c_1;c_0;U_1;U_0}} & two available acts with consequence probabilities, additive outcome gain, costs and expected values & sec:p19-r64-nonstandard;annex:p19.r64.choice-gap & meta & \hyperref[def:p19-r64-binary-choice]{\ref*{def:p19-r64-binary-choice}} & \thepage & Stipulated bridge not participant-fit quantities or source-supplied numerical awards. \\
\phantomsection\label{entry:symbol-registry-04097}\(x,x',y,y',u,u'\) & \texttt{\detokenize{x;x';y;y';u;u'}} & utility-grid actors, affected targets and noncardinal utility signs & sec:p19-r64-nonstandard;sec:p19-r64-templates & meta & \hyperref[sec:p19-r64-nonstandard]{\ref*{sec:p19-r64-nonstandard}} & \thepage & Grid utility signs differ from additive v of the explicitly numeric bridge. \\
\phantomsection\label{entry:symbol-registry-04098}\(s,h,e,\omega,A\) & \texttt{\detokenize{s;h;e;\omega;A}} & speaker, hearer, other agent, neutral world agent and set of agents & sec:p19-r64-nonstandard;sec:p19-r64-templates & meta & \hyperref[sec:p19-r64-nonstandard]{\ref*{sec:p19-r64-nonstandard}} & \thepage & Zero known utility convention uses A as target; not a factual claim that every person is indifferent. \\
\phantomsection\label{entry:symbol-registry-04099}\(Z,D,Q\) & \texttt{\detokenize{Z;D;Q}} & background indicator, observed/chosen action and undesirable outcome in diagnostic counterexample & sec:p19-r64-causal-boundary;annex:p19.r64.diagnostic-failure & meta & \hyperref[thm:p19-r64-diagnostic-failure]{\ref*{thm:p19-r64-diagnostic-failure}} & \thepage & Finite structural model, not a measured representation or a fitted experimental task. \\
\phantomsection\label{entry:symbol-registry-04100}\(F,t,r,p,d,\eta_p^2\) & \texttt{\detokenize{F;t;r;p;d;\eta_p^2}} & reported ANOVA, t test, correlation, significance, effect size and partial eta squared & sec:p19-r64-revision;sec:p19-r64-templates;sec:p19-r64-causal-boundary & meta & \hyperref[sec:p19-r64-revision]{\ref*{sec:p19-r64-revision}} & \thepage & r statistic is distinct from outcome-probability r\_1,r\_0; d effect size distinct from intervention value d in annex. \\
\phantomsection\label{entry:symbol-registry-04102}\(C,E\) & \texttt{\detokenize{C,}\allowbreak\detokenize{E}} & binary target cause and effect random variables & p19/rationality-7-2 power comparison & n/a & \hyperref[sec:p19-rat72-dependency]{\ref*{sec:p19-rat72-dependency}} & \thepage & C is redeclared as a diamond variable only in its later local convention \\
\phantomsection\label{entry:symbol-registry-04103}\(q_0,q_1\) & \texttt{\detokenize{q_0,}\allowbreak\detokenize{q_1}} & conditional effect probabilities at absent and present target cause & p19/rationality-7-2 power comparison & n/a & \hyperref[sec:p19-rat72-dependency]{\ref*{sec:p19-rat72-dependency}} & \thepage & positive conditioning probabilities required \\
\phantomsection\label{entry:symbol-registry-04104}\(\Delta P\) & \texttt{\detokenize{\Delta P}} & observed conditional frequency contrast & p19/rationality-7-2 power comparison & n/a & \hyperref[sec:p19-rat72-dependency]{\ref*{sec:p19-rat72-dependency}} & \thepage & not automatically generative power or support for a link \\
\phantomsection\label{entry:symbol-registry-04105}\(K,T\) & \texttt{\detokenize{K,}\allowbreak\detokenize{T}} & independent binary background and target-success generators & p19/rationality-7-2 power comparison & n/a & \hyperref[thm:p19-rat72-power-benchmark]{\ref*{thm:p19-rat72-power-benchmark}} & \thepage & stipulated editorial model \\
\phantomsection\label{entry:symbol-registry-04106}\(r,s\) & \texttt{\detokenize{r,}\allowbreak\detokenize{s}} & background and target-success probabilities & p19/rationality-7-2 power comparison & n/a & \hyperref[thm:p19-rat72-power-benchmark]{\ref*{thm:p19-rat72-power-benchmark}} & \thepage & r separately redeclared as root probability in the later diamond convention \\
\phantomsection\label{entry:symbol-registry-04107}\(A,B,C,D\) & \texttt{\detokenize{A,}\allowbreak\detokenize{B,}\allowbreak\detokenize{C,}\allowbreak\detokenize{D}} & binary variables in the instructed diamond graph & p19/rationality-7-2 diamond comparison & n/a & \hyperref[def:p19-rat72-diamond-notation]{\ref*{def:p19-rat72-diamond-notation}} & \thepage & capital variable names distinguish graph roles from values \\
\phantomsection\label{entry:symbol-registry-04108}\(a,b,c,d\) & \texttt{\detokenize{a,}\allowbreak\detokenize{b,}\allowbreak\detokenize{c,}\allowbreak\detokenize{d}} & binary values indexing the diamond joint and kernels & p19/rationality-7-2 diamond comparison & n/a & \hyperref[def:p19-rat72-diamond-notation]{\ref*{def:p19-rat72-diamond-notation}} & \thepage & not propositions and not lowercase event-name shorthand \\
\phantomsection\label{entry:symbol-registry-04109}\(r,u_a,v_a,t_{bc}\) & \texttt{\detokenize{r,}\allowbreak\detokenize{u_a,}\allowbreak\detokenize{v_a,}\allowbreak\detokenize{t_{bc}}} & root and conditional success probabilities in the diamond & p19/rationality-7-2 diamond comparison & n/a & \hyperref[def:p19-rat72-diamond-notation]{\ref*{def:p19-rat72-diamond-notation}} & \thepage & local typed parameters \\
\phantomsection\label{entry:symbol-registry-04110}\(\ell_a(c)\) & \texttt{\detokenize{\ell_a(c)}} & likelihood of observed C value c given A value a & p19/rationality-7-2 diamond comparison & n/a & \hyperref[thm:p19-rat72-diamond-benchmark]{\ref*{thm:p19-rat72-diamond-benchmark}} & \thepage & conditional mass function \\
\phantomsection\label{entry:symbol-registry-04111}\(k_c,h_c,h\) & \texttt{\detokenize{k_c,}\allowbreak\detokenize{h_c,}\allowbreak\detokenize{h}} & observational root posterior observational B marginal and unchanged B marginal & p19/rationality-7-2 diamond comparison & n/a & \hyperref[thm:p19-rat72-diamond-benchmark]{\ref*{thm:p19-rat72-diamond-benchmark}} & \thepage & distinguish conditioning from replacement \\
\phantomsection\label{entry:symbol-registry-04112}\(\operatorname{do}(C=c)\) & \texttt{\detokenize{\operatorname{do}(C=c)}} & fixed replacement of the C mechanism & p19/rationality-7-2 diamond comparison & n/a & \hyperref[def:p19-rat72-diamond-notation]{\ref*{def:p19-rat72-diamond-notation}} & \thepage & inherits the P16 intervention interpretation rather than denoting a conditioning event \\
\phantomsection\label{entry:symbol-registry-04113}\(X,Y\) & \texttt{\detokenize{X,}\allowbreak\detokenize{Y}} & objects in the launching perceptual example & p19/rationality-7-2 disposition comparison & n/a & \hyperref[sec:p19-rat72-disposition]{\ref*{sec:p19-rat72-disposition}} & \thepage & object roles not the diamond random variables \\
\phantomsection\label{entry:symbol-registry-04114}\(C,E;c,e,\neg c,\neg e\) & \texttt{\detokenize{C,}\allowbreak\detokenize{E,}\allowbreak\detokenize{c,}\allowbreak\detokenize{e,}\allowbreak\detokenize{\neg c,}\allowbreak\detokenize{\neg e}} & binary candidate-cause and effect variables with their presence or absence events & P19 rationality-7-3 elementary diagnosis & n/a & \hyperref[def:p19-rat73-queries]{\ref*{def:p19-rat73-queries}} & \thepage & Uppercase denotes variables; lowercase denotes value events; not all probability relations are causal. \\
\phantomsection\label{entry:symbol-registry-04115}\(P,\widehat P,N\) & \texttt{\detokenize{P,}\allowbreak\detokenize{\widehat P,}\allowbreak\detokenize{N}} & specified probability law empirical table-relative frequency and cell count & P19 rationality-7-3 & n/a & \hyperref[sec:p19-rat73-elementary]{\ref*{sec:p19-rat73-elementary}} & \thepage & N counts fictitious learning cases rather than participants; a hat distinguishes observed proportion from population or model probability. \\
\phantomsection\label{entry:symbol-registry-04116}\(c\to e\) & \texttt{\detokenize{c\to e}} & noisy-OR productive-mechanism success event & P19 rationality-7-3 attribution & n/a & \hyperref[def:p19-rat73-queries]{\ref*{def:p19-rat73-queries}} & \thepage & Typed as C and X both present; neither material implication nor exclusive attribution or counterfactual necessity. \\
\phantomsection\label{entry:symbol-registry-04117}\(b_c,u,v,w_c,w_a\) & \texttt{\detokenize{b_c,}\allowbreak\detokenize{u,}\allowbreak\detokenize{v,}\allowbreak\detokenize{w_c,}\allowbreak\detokenize{w_a}} & cause base rate two effect likelihoods and target or background generative probabilities & P19 rationality-7-3 noisy-OR & n/a & \hyperref[thm:p19-rat73-elementary]{\ref*{thm:p19-rat73-elementary}} & \thepage & The inversion for w\_c requires v less than one and the generative range u at least v. \\
\phantomsection\label{entry:symbol-registry-04118}\(X,Z,F\) & \texttt{\detokenize{X,}\allowbreak\detokenize{Z,}\allowbreak\detokenize{F}} & candidate mechanism success background production and the production event & P19 rationality-7-3 noisy-OR body and annex & n/a & \hyperref[thm:p19-rat73-elementary]{\ref*{thm:p19-rat73-elementary}} & \thepage & Mutual independence of C X Z is stipulated; F is C and X in the reciprocal annex. \\
\phantomsection\label{entry:symbol-registry-04119}\(D,S_0,S_1,\theta_i\) & \texttt{\detokenize{D,}\allowbreak\detokenize{S_0,}\allowbreak\detokenize{S_1,}\allowbreak\detokenize{\theta_i}} & training contingency data two generating graph hypotheses and their own parameter tuples & P19 rationality-7-3 structure induction & n/a & \hyperref[def:p19-rat73-structure-method]{\ref*{def:p19-rat73-structure-method}} & \thepage & S0 omits a C-to-E link; S1 uses three parameters; no exhaustive universal graph space is asserted. \\
\phantomsection\label{entry:symbol-registry-04120}\(L_i,\pi_i,m_i,\rho_i,g_i,g\) & \texttt{\detokenize{L_i,}\allowbreak\detokenize{\pi_i,}\allowbreak\detokenize{m_i,}\allowbreak\detokenize{\rho_i,}\allowbreak\detokenize{g_i,}\allowbreak\detokenize{g}} & likelihood prior or posterior parameter measure marginal likelihood graph weight conditional-score integral and weighted diagnostic score & P19 rationality-7-3 structure induction & n/a & \hyperref[def:p19-rat73-structure-method]{\ref*{def:p19-rat73-structure-method}} & \thepage & All weights condition on training data only in the source rating rule; g is not silently the conditional of an integrated joint law. \\
\phantomsection\label{entry:symbol-registry-04121}\(\rho_i,q_i,t_i,P_i,H\) & \texttt{\detokenize{\rho_i,}\allowbreak\detokenize{q_i,}\allowbreak\detokenize{t_i,}\allowbreak\detokenize{P_i,}\allowbreak\detokenize{H}} & finite mixture model weight effect probability diagnostic conditional model law and arbitrary event & P19 rationality-7-3 mixture body and annex & n/a & \hyperref[thm:p19-rat73-averaging]{\ref*{thm:p19-rat73-averaging}} & \thepage & rho shares the graph-weight role in the finite comparison; all q\_i are positive. \\
\phantomsection\label{entry:symbol-registry-04122}\(A,B,C;a,b,c;p,r,s\) & \texttt{\detokenize{A,}\allowbreak\detokenize{B,}\allowbreak\detokenize{C,}\allowbreak\detokenize{a,}\allowbreak\detokenize{b,}\allowbreak\detokenize{c,}\allowbreak\detokenize{p,}\allowbreak\detokenize{r,}\allowbreak\detokenize{s}} & three network variables value events finite kernel mass and two independent cause probabilities & P19 rationality-7-3 networks and logical contrast & n/a & \hyperref[thm:p19-rat73-network]{\ref*{thm:p19-rat73-network}} & \thepage & Network C reuses the candidate-cause role locally; in the common-cause graph B is the shared cause and A C are effects. \\
\phantomsection\label{entry:symbol-registry-04123}\(r_{\mathrm{overall}},r_{\mathrm{trends}},\mathrm{RMSE},F,p,\eta^2\) & \texttt{\detokenize{r_{\mathrm{overall}},}\allowbreak\detokenize{r_{\mathrm{trends}},}\allowbreak\detokenize{\mathrm{RMSE},}\allowbreak\detokenize{F,}\allowbreak\detokenize{p,}\allowbreak\detokenize{\eta^2}} & reported model fit correlations root-mean-square error ANOVA statistic p-value and reported effect-size statistic & P19 rationality-7-3 empirical evidence & n/a & \hyperref[sec:p19-rat73-evidence]{\ref*{sec:p19-rat73-evidence}} & \thepage & Statistical F differs from the production event in the annex; p-value differs from network mass p; no partial-eta subtype is invented. \\
\phantomsection\label{entry:symbol-registry-04276}\(X\) & \texttt{\detokenize{X}} & nonempty alternative domain & P19 economics certainty & object & \hyperref[def:p19-r94-certainty]{\ref*{def:p19-r94-certainty}} & \thepage & Topology and outcome description are explicitly declared. \\
\phantomsection\label{entry:symbol-registry-04277}\(\succeq\) & \texttt{\detokenize{\succeq}} & weak preference & P19 economic comparison & relation & \hyperref[def:p19-r94-certainty]{\ref*{def:p19-r94-certainty}} & \thepage & Domains change explicitly between certainty lotteries and dated outcomes. \\
\phantomsection\label{entry:symbol-registry-04278}\(\succ\) & \texttt{\detokenize{\succ}} & strict preference & P19 economic comparison & relation & \hyperref[def:p19-r94-certainty]{\ref*{def:p19-r94-certainty}} & \thepage & Weak comparison with denial of the reverse comparison. \\
\phantomsection\label{entry:symbol-registry-04279}\(\sim\) & \texttt{\detokenize{\sim}} & preference indifference & P19 economic comparison & relation & \hyperref[def:p19-r94-certainty]{\ref*{def:p19-r94-certainty}} & \thepage & Two affirmative weak comparisons rather than incomparability. \\
\phantomsection\label{entry:symbol-registry-04280}\(u\) & \texttt{\detokenize{u}} & ordinal or consequence utility on a declared domain & P19 economic comparison & object & \hyperref[def:p19-r94-certainty]{\ref*{def:p19-r94-certainty}} & \thepage & Expected utility and dated-outcome uses have additional assumptions. \\
\phantomsection\label{entry:symbol-registry-04281}\((x_1,\ldots,x_n)\) & \texttt{\detokenize{x_i}} & finite prize list & P19 economics risk & object & \hyperref[sec:p19-r94-risk]{\ref*{sec:p19-r94-risk}} & \thepage & Numeric order alone does not prove monetary utility linearity. \\
\phantomsection\label{entry:symbol-registry-04282}\((p_1,\ldots,p_n)\) & \texttt{\detokenize{p_i}} & prize probabilities & P19 economics risk & object & \hyperref[sec:p19-r94-risk]{\ref*{sec:p19-r94-risk}} & \thepage & Nonnegative masses summing to one. \\
\phantomsection\label{entry:symbol-registry-04283}\((L,L_1,L_2,H)\) & \texttt{\detokenize{L;L_1;L_2;H}} & lottery alternatives and common mixture component & P19 economics risk & object & \hyperref[def:p19-r94-risk-axioms]{\ref*{def:p19-r94-risk-axioms}} & \thepage & Compound procedures are reduced under an explicit preference assumption. \\
\phantomsection\label{entry:symbol-registry-04284}\(q\) & \texttt{\detokenize{q}} & common mixture or endpoint-calibration weight & P19 economics risk & object & \hyperref[def:p19-r94-risk-axioms]{\ref*{def:p19-r94-risk-axioms}} & \thepage & Zero weight is excluded from the independence biconditional. \\
\phantomsection\label{entry:symbol-registry-04285}\((a_j,p_{ji})\) & \texttt{\detokenize{a_j;p_{ji}}} & branch weight and conditional prize probability & P19 economics risk & object & \hyperref[def:p19-r94-risk-axioms]{\ref*{def:p19-r94-risk-axioms}} & \thepage & Reduction uses products and finite sums; zero branches contribute nothing. \\
\phantomsection\label{entry:symbol-registry-04286}\(EU(L)\) & \texttt{\detokenize{EU}} & expected utility of a finite lottery & P19 economics risk & object & \hyperref[sec:p19-r94-risk]{\ref*{sec:p19-r94-risk}} & \thepage & Linear in probability weights not necessarily in monetary outcomes. \\
\phantomsection\label{entry:symbol-registry-04287}\((M,m)\) & \texttt{\detokenize{M;m}} & designated best and worst certain lotteries & P19 economics risk & object & \hyperref[thm:p19-r94-source-vnm]{\ref*{thm:p19-r94-source-vnm}} & \thepage & Nondegenerate endpoint ordering and normalization. \\
\phantomsection\label{entry:symbol-registry-04288}\((Z,\Gamma)\) & \texttt{\detokenize{Z;\Gamma}} & outcome interval and admissible date domain & P19 economics time & object & \hyperref[def:p19-r94-time-axioms]{\ref*{def:p19-r94-time-axioms}} & \thepage & Original domain A0 is preserved. \\
\phantomsection\label{entry:symbol-registry-04289}\((z,w,t,s,\tau)\) & \texttt{\detokenize{z;w;t;s;\tau}} & dated outcomes dates and their common difference & P19 economics time & object & \hyperref[def:p19-r94-time-axioms]{\ref*{def:p19-r94-time-axioms}} & \thepage & Every date in a stationarity comparison must belong to the domain. \\
\phantomsection\label{entry:symbol-registry-04290}\(\delta\) & \texttt{\detokenize{\delta}} & exponential factor & P19 economics time and planning & object & \hyperref[thm:p19-r94-time-representation]{\ref*{thm:p19-r94-time-representation}} & \thepage & Fixed in comparisons but not identified independently of dated-outcome utility. \\
\phantomsection\label{entry:symbol-registry-04291}\((S_i,y_i,c_i)\) & \texttt{\detokenize{S_i;y_i;c_i}} & savings income and consumption at date i & P19 economics planning & object & \hyperref[def:p19-r94-stream]{\ref*{def:p19-r94-stream}} & \thepage & Feasibility and terminal solvency are explicit. \\
\phantomsection\label{entry:symbol-registry-04292}\((T,r)\) & \texttt{\detokenize{T;r}} & finite horizon and nonnegative return parameter & P19 economics planning & object & \hyperref[def:p19-r94-stream]{\ref*{def:p19-r94-stream}} & \thepage & An illustrative mathematical planning model rather than investment advice. \\
\phantomsection\label{entry:symbol-registry-04293}\(V_i(s)\) & \texttt{\detokenize{V_i}} & finite continuation value at a savings state & P19 economics planning & object & \hyperref[thm:p19-r94-backward]{\ref*{thm:p19-r94-backward}} & \thepage & Maximum exists because feasible action sets are finite and nonempty. \\
\phantomsection\label{entry:symbol-registry-04294}\((\beta,U)\) & \texttt{\detokenize{\beta;U}} & present-bias factor and quasi-hyperbolic stream evaluation & P19 economics time & object & \hyperref[def:p19-r94-present-bias]{\ref*{def:p19-r94-present-bias}} & \thepage & Future re-evaluation changes which reward is current. \\
\phantomsection\label{entry:symbol-registry-04295}\((P,p,z,Z)\) & \texttt{\detokenize{P;p;z;Z}} & win probabilities and prizes in the preference-reversal pair & P19 economics reversal & object & \hyperref[sec:p19-r94-reversals]{\ref*{sec:p19-r94-reversals}} & \thepage & Distinct local roles from outcome-interval notation. \\
\phantomsection\label{entry:symbol-registry-04296}\((C_P,C_{\$})\) & \texttt{\detokenize{C_P;C_{\$}}} & certainty equivalents of the two bets & P19 economics reversal & object & \hyperref[thm:p19-r94-reversal]{\ref*{thm:p19-r94-reversal}} & \thepage & Genuine indifference and common consequence domain are required. \\
\phantomsection\label{entry:symbol-registry-04297}\((A,B,C,D)\) & \texttt{\detokenize{A;B;C;D}} & four declared Allais-type lotteries & P19 economics Allais contrast & object & \hyperref[def:p19-r94-allais]{\ref*{def:p19-r94-allais}} & \thepage & Worked mathematical illustration with no empirical frequency claim. \\
\phantomsection\label{entry:symbol-registry-04298}\((\chi^2,F,N,p)\) & \texttt{\detokenize{\chi^2;F;N;p}} & reported test statistics analysis denominator and significance value & P19 economics jam record & meta & \hyperref[sec:p19-r94-jam]{\ref*{sec:p19-r94-jam}} & \thepage & Do not identify a p-value with a lottery probability or effect interval. \\
\phantomsection\label{entry:symbol-registry-04299}\(\delta_{x_i}\) & \texttt{\detokenize{\delta_{x_i}}} & certain prize lottery & P19 economics risk & object & \hyperref[thm:p19-r94-source-vnm]{\ref*{thm:p19-r94-source-vnm}} & \thepage & Dirac lottery rather than the time-discount factor. \\
\phantomsection\label{entry:symbol-registry-04300}\((B_q,u_i)\) & \texttt{\detokenize{B_q;u_i}} & endpoint mixture and normalized prize utility & P19 economics risk & object & \hyperref[thm:p19-r94-source-vnm]{\ref*{thm:p19-r94-source-vnm}} & \thepage & Calibration weights identify normalized consequence values. \\
\phantomsection\label{entry:symbol-registry-04301}\(v\) & \texttt{\detokenize{v}} & alternative expected-utility prize representation & P19 economics risk & object & \hyperref[thm:p19-r94-source-vnm]{\ref*{thm:p19-r94-source-vnm}} & \thepage & Values on endpoint lotteries mean the utilities of their underlying certain prizes. \\
\phantomsection\label{entry:symbol-registry-04302}\((g,h)\) & \texttt{\detokenize{g;h}} & increasing utility transform and common date shift & P19 economics annex calculations & object & \hyperref[thm:p19-r94-continuous-converse]{\ref*{thm:p19-r94-continuous-converse}} & \thepage & Local objects in separate converse and time-comparison proofs. \\
\phantomsection\label{entry:symbol-registry-04303}\((x_i,x_{-i},u_m)\) & \texttt{\detokenize{x_i;x_{-i};u_m}} & own and other consumption and emotion-indexed utility & P19 economics auxiliary comparison & object & \hyperref[def:p19-r94-social-emotional]{\ref*{def:p19-r94-social-emotional}} & \thepage & Alternatives include relevant consequences; state changes need not create within-state cycles. \\
\phantomsection\label{entry:symbol-registry-04304}\(m\) & \texttt{\detokenize{m}} & emotional-state index & P19 economics auxiliary comparison & object & \hyperref[def:p19-r94-social-emotional]{\ref*{def:p19-r94-social-emotional}} & \thepage & Distinct local role from the designated worst lottery. \\
\phantomsection\label{entry:symbol-registry-04350}\(\Gamma\) & \texttt{\detokenize{\Gamma}} & finite interpreted relational premise set & P19 spatial-temporal relational reasoning & object & \hyperref[def:p19-rat132-task]{\ref*{def:p19-rat132-task}} & \thepage & Satisfiability is checked separately from necessity. \\
\phantomsection\label{entry:symbol-registry-04351}\(\mathcal M(\Gamma)\) & \texttt{\detokenize{\mathcal M(\Gamma)}} & satisfying configurations of the relational premises & P19 spatial-temporal relational reasoning & meta & \hyperref[def:p19-rat132-task]{\ref*{def:p19-rat132-task}} & \thepage & Not a claim about all models constructed by a person. \\
\phantomsection\label{entry:symbol-registry-04352}\(q\) & \texttt{\detokenize{q}} & queried relational conclusion & P19 spatial-temporal relational reasoning & meta & \hyperref[def:p19-rat132-task]{\ref*{def:p19-rat132-task}} & \thepage & Possible and necessary queries have different quantifiers. \\
\phantomsection\label{entry:symbol-registry-04353}\((c,m,g)\) & \texttt{\detokenize{c;m;g}} & computer moon-landing and genome event points & P19 spatial-temporal relational reasoning & object & \hyperref[thm:p19-rat132-order]{\ref*{thm:p19-rat132-order}} & \thepage & Local event names; m here is not the interval meets relation. \\
\phantomsection\label{entry:symbol-registry-04354}\((y_A,y_B,y_C)\) & \texttt{\detokenize{y_A;y_B;y_C}} & northward coordinates in a common spatial frame & P19 spatial-temporal relational reasoning & object & \hyperref[thm:p19-rat132-order]{\ref*{thm:p19-rat132-order}} & \thepage & Strictly increasing coordinate defines northward order in the formal example. \\
\phantomsection\label{entry:symbol-registry-04355}\((I,J,K)\) & \texttt{\detokenize{I;J;K}} & proper intervals in the endpoint benchmark & P19 spatial-temporal relational reasoning & object & \hyperref[def:p19-rat132-interval]{\ref*{def:p19-rat132-interval}} & \thepage & Can represent temporal or common-axis spatial intervals; not arbitrary planar regions. \\
\phantomsection\label{entry:symbol-registry-04356}\((i^-,i^+,j^-,j^+)\) & \texttt{\detokenize{i^-;i^+;j^-;j^+}} & real interval endpoints & P19 spatial-temporal relational reasoning & object & \hyperref[def:p19-rat132-interval]{\ref*{def:p19-rat132-interval}} & \thepage & Starts strictly precede ends; endpoints of different intervals may coincide. \\
\phantomsection\label{entry:symbol-registry-04357}\((b,m,o,d,s,f,e)\) & \texttt{\detokenize{b;m;o;d;s;f;e}} & Ragni interval before meets overlaps during starts finishes equals relations & P19 spatial-temporal relational reasoning & relation & \hyperref[def:p19-rat132-interval]{\ref*{def:p19-rat132-interval}} & \thepage & Local source syntax; before and equals translate to Allen less-than and equals symbols. \\
\phantomsection\label{entry:symbol-registry-04358}\((bi,mi,oi,di,si,fi)\) & \texttt{\detokenize{bi;mi;oi;di;si;fi}} & inverse interval base relations & P19 spatial-temporal relational reasoning & relation & \hyperref[def:p19-rat132-interval]{\ref*{def:p19-rat132-interval}} & \thepage & Inverses exchange intervals; equality is self-inverse. \\
\phantomsection\label{entry:symbol-registry-04359}\(I\{b,o\}J\) & \texttt{\detokenize{I\{b;o\}J}} & disjunction of permitted interval relations & P19 spatial-temporal relational reasoning & relation & \hyperref[def:p19-rat132-interval]{\ref*{def:p19-rat132-interval}} & \thepage & Relation sets have no probability interpretation. \\
\phantomsection\label{entry:symbol-registry-04360}\((R,S)\) & \texttt{\detokenize{R;S}} & sets of interval base relations & P19 spatial-temporal relational reasoning & object & \hyperref[sec:p19-rat132-intervals]{\ref*{sec:p19-rat132-intervals}} & \thepage & Typed separately from room interval R in the original switch example. \\
\phantomsection\label{entry:symbol-registry-04361}\(R\diamond S\) & \texttt{\detokenize{\diamond}} & qualitative composition of interval relation sets & P19 spatial-temporal relational reasoning & infix & \hyperref[sec:p19-rat132-intervals]{\ref*{sec:p19-rat132-intervals}} & \thepage & Union of base relations compatible with two consecutive interval relations. \\
\phantomsection\label{entry:symbol-registry-04362}\((N_{ij},N'_{ij})\) & \texttt{\detokenize{N_{ij};N'_{ij}}} & current and narrowed allowed interval-arc relation sets & P19 spatial-temporal relational reasoning & object & \hyperref[def:p19-rat132-propagation]{\ref*{def:p19-rat132-propagation}} & \thepage & Nonempty local fixed points do not certify global satisfiability. \\
\phantomsection\label{entry:symbol-registry-04363}\(R^{\smile}\) & \texttt{\detokenize{R^{\smile}}} & converse interval relation set & P19 spatial-temporal relational reasoning & meta & \hyperref[sheet:p19-rat132-interval-fragment]{\ref*{sheet:p19-rat132-interval-fragment}} & \thepage & Replace each base relation by its inverse. \\
\phantomsection\label{entry:symbol-registry-04364}\((S,L,R)\) & \texttt{\detokenize{S;L;R}} & switch light-on and room-presence intervals & P19 spatial-temporal relational reasoning & object & \hyperref[sec:p19-rat132-intervals]{\ref*{sec:p19-rat132-intervals}} & \thepage & Local original Allen example; not the relation-set variables R and S. \\
\phantomsection\label{entry:symbol-registry-04365}\((X,Y)\) & \texttt{\detokenize{X;Y}} & nonempty regular closed planar regions in the RCC illustration & P19 spatial-temporal relational reasoning & object & \hyperref[sec:p19-rat132-regions]{\ref*{sec:p19-rat132-regions}} & \thepage & Boundary and interior define the ordinary figure interpretation. \\
\phantomsection\label{entry:symbol-registry-04366}\((DC,EC,PO,EQ,TPP,NTPP,TPPi,NTPPi)\) & \texttt{\detokenize{DC;EC;PO;EQ;TPP;NTPP;TPPi;NTPPi}} & RCC-8 region base relations & P19 spatial-temporal relational reasoning & relation & \hyperref[sec:p19-rat132-regions]{\ref*{sec:p19-rat132-regions}} & \thepage & Not directional relations; inverse NTPP spelling repairs the source caption typo. \\
\phantomsection\label{entry:symbol-registry-04367}\((F,P,B,H,D)\) & \texttt{\detokenize{F;P;B;H;D}} & Ferrari Porsche Beetle Hummer Dodge tokens & P19 spatial-temporal relational reasoning & object & \hyperref[thm:p19-rat132-cars]{\ref*{thm:p19-rat132-cars}} & \thepage & Strict one-dimensional all-distinct car-order reconstruction. \\
\phantomsection\label{entry:symbol-registry-04368}\((X,r,Y)\) & \texttt{\detokenize{X;r;Y}} & PRISM located-object relation reference-object input triple & P19 spatial-temporal relational reasoning & object & \hyperref[sec:p19-rat132-prism]{\ref*{sec:p19-rat132-prism}} & \thepage & Separate type from RCC region names; interpretation is assumed before inference. \\
\phantomsection\label{entry:symbol-registry-04369}\((A,F,S)\) & \texttt{\detokenize{A;F;S}} & alarm fire and siren proper intervals & P19 spatial-temporal relational reasoning & object & \hyperref[thm:p19-rat132-durative]{\ref*{thm:p19-rat132-durative}} & \thepage & Separate type from Ferrari token F and Allen switch interval S. \\
\phantomsection\label{entry:symbol-registry-04370}\((a^-,a^+,f^-,f^+,s^-,s^+)\) & \texttt{\detokenize{a^-;a^+;f^-;f^+;s^-;s^+}} & alarm fire and siren interval endpoints & P19 spatial-temporal relational reasoning & object & \hyperref[thm:p19-rat132-durative]{\ref*{thm:p19-rat132-durative}} & \thepage & Strict containment is explicit rather than any-overlap during. \\
\phantomsection\label{entry:symbol-registry-04371}\(n\) & \texttt{\detokenize{n}} & number of intervals in finite endpoint enumeration & P19 spatial-temporal relational reasoning & object & \hyperref[thm:p19-rat132-endpoints]{\ref*{thm:p19-rat132-endpoints}} & \thepage & Zero intervals are separately covered in the complete annex proof. \\
\phantomsection\label{entry:symbol-registry-04372}\((p,z,\delta)\) & \texttt{\detokenize{p;z;\delta}} & reported test probability standardized statistic and Cliff effect size & P19 spatial-temporal relational reasoning & object & \hyperref[sec:p19-rat132-spatial-evidence]{\ref*{sec:p19-rat132-spatial-evidence}} & \thepage & Reported primary-study quantities; p first occurs in the spatial record and z and delta in the temporal record. \\
\phantomsection\label{entry:symbol-registry-04373}\(b\) & \texttt{\detokenize{b}} & fitted generalized logistic mixed-model coefficient & P19 spatial-temporal relational reasoning & object & \hyperref[sec:p19-rat132-temporal-evidence]{\ref*{sec:p19-rat132-temporal-evidence}} & \thepage & Typed separately from the interval before relation b. \\
\phantomsection\label{entry:symbol-registry-04374}\(L\) & \texttt{\detokenize{L}} & reported Page trend-test statistic & P19 spatial-temporal relational reasoning & object & \hyperref[sec:p19-rat132-spatial-evidence]{\ref*{sec:p19-rat132-spatial-evidence}} & \thepage & Typed separately from the light-on interval in Allen's example. \\
\phantomsection\label{entry:symbol-registry-04375}\(P\) & \texttt{\detokenize{P}} & raining proposition in the illustrative material conditional & P19 rationality improvement & object & \hyperref[sec:p19-r153-source-checks]{\ref*{sec:p19-r153-source-checks}} & \thepage & Fixed interpretation for the weather countermodel; not an empirical event-frequency estimate. \\
\phantomsection\label{entry:symbol-registry-04376}\(Q\) & \texttt{\detokenize{Q}} & wet-streets proposition in the illustrative material conditional & P19 rationality improvement & object & \hyperref[sec:p19-r153-source-checks]{\ref*{sec:p19-r153-source-checks}} & \thepage & Truth need not follow the occurrence of rain in reverse. \\
\phantomsection\label{entry:symbol-registry-04377}\(\to\) & \texttt{\detokenize{\to}} & material conditional between the weather propositions & P19 rationality improvement & infix & \hyperref[sec:p19-r153-source-checks]{\ref*{sec:p19-r153-source-checks}} & \thepage & Formal reading declared before the countermodel; not a causal arrow. \\
\phantomsection\label{entry:symbol-registry-04378}\(\neg\) & \texttt{\detokenize{\neg}} & Boolean negation in the weather countermodel & P19 rationality improvement & meta & \hyperref[sec:p19-r153-source-checks]{\ref*{sec:p19-r153-source-checks}} & \thepage & Propositional negation under the declared material-conditional semantics. \\
\phantomsection\label{entry:symbol-registry-04379}\(C\) & \texttt{\detokenize{C}} & cancer event in the stipulated screening population & P19 rationality improvement & object & \hyperref[sec:p19-r153-source-checks]{\ref*{sec:p19-r153-source-checks}} & \thepage & Hypothetical instructional event; not a current clinical population estimate. \\
\phantomsection\label{entry:symbol-registry-04380}\(C^c\) & \texttt{\detokenize{C^c}} & complementary no-cancer event in the screening population & P19 rationality improvement & object & \hyperref[sec:p19-r153-source-checks]{\ref*{sec:p19-r153-source-checks}} & \thepage & Cancer and no cancer form an exhaustive disjoint partition in the instructional law. \\
\phantomsection\label{entry:symbol-registry-04381}\(+\) & \texttt{\detokenize{+}} & positive-test event in the screening example & P19 rationality improvement & object & \hyperref[sec:p19-r153-source-checks]{\ref*{sec:p19-r153-source-checks}} & \thepage & An event name rather than an addition operator inside probability expressions. \\
\phantomsection\label{entry:symbol-registry-04382}\(\mathbb P\) & \texttt{\detokenize{\mathbb P}} & probability under the screening scenario law & P19 rationality improvement & n/a & \hyperref[sec:p19-r153-source-checks]{\ref*{sec:p19-r153-source-checks}} & \thepage & All inputs are stipulated and conditional denominator is positive. \\
\phantomsection\label{entry:symbol-registry-04383}\(u_M\) & \texttt{\detokenize{u_M}} & total future utility of the Michigan option & P19 rationality improvement & object & \hyperref[thm:p19-r153-example-checks]{\ref*{thm:p19-r153-example-checks}} & \thepage & Excludes common irrecoverable past cost; must include decision-related future experience. \\
\phantomsection\label{entry:symbol-registry-04384}\(u_W\) & \texttt{\detokenize{u_W}} & total future utility of the Wisconsin option & P19 rationality improvement & object & \hyperref[thm:p19-r153-example-checks]{\ref*{thm:p19-r153-example-checks}} & \thepage & Option ordering is a normative premise not a measured training outcome. \\
\phantomsection\label{entry:symbol-registry-04385}\(K\) & \texttt{\detokenize{K}} & common irrecoverable past cost under both choices & P19 rationality improvement & object & \hyperref[thm:p19-r153-example-checks]{\ref*{thm:p19-r153-example-checks}} & \thepage & Subtracting the same cost preserves the declared future-utility ordering. \\
\phantomsection\label{entry:symbol-registry-04386}\(F\) & \texttt{\detokenize{F}} & reported variance-ratio comparison statistic with degrees of freedom & P19 rationality improvement & n/a & \hyperref[sec:p19-r153-games]{\ref*{sec:p19-r153-games}} & \thepage & Primary articles' reported analyses retained without independent reanalysis. \\
\phantomsection\label{entry:symbol-registry-04387}\(t\) & \texttt{\detokenize{t}} & reported mean-comparison statistic with degrees of freedom & P19 rationality improvement & n/a & \hyperref[sec:p19-r153-games]{\ref*{sec:p19-r153-games}} & \thepage & Follow-up comparisons retain their actual analysed denominators. \\
\phantomsection\label{entry:symbol-registry-04388}\(p\) & \texttt{\detokenize{p}} & reported significance level under the original analysis & P19 rationality improvement & n/a & \hyperref[sec:p19-r153-games]{\ref*{sec:p19-r153-games}} & \thepage & Not a population effect size or probability that the hypothesis is true. \\
\phantomsection\label{entry:symbol-registry-04389}\(d\) & \texttt{\detokenize{d}} & reported standardized change effect size & P19 rationality improvement & n/a & \hyperref[sec:p19-r153-games]{\ref*{sec:p19-r153-games}} & \thepage & Within-group standardized changes do not alone identify a causal intervention effect. \\
\phantomsection\label{entry:symbol-registry-04390}\(P,Q,A\) & \texttt{\detokenize{P,}\allowbreak\detokenize{Q,}\allowbreak\detokenize{A}} & antecedent consequent and alternative-antecedent propositions in the developmental contrast & P19 rationality 15.1 & object & \hyperref[def:p19-r151-conditional-benchmark]{\ref*{def:p19-r151-conditional-benchmark}} & \thepage & These are benchmark roles not direct measurements of a child's internal representation. \\
\phantomsection\label{entry:symbol-registry-04391}\(N\) & \texttt{\detokenize{N}} & reported analysis denominator in the specified original study & p19-development-aging-statistics & object & \hyperref[sec:p19-r151-validity-record]{\ref*{sec:p19-r151-validity-record}} & \thepage & A denominator can count response patterns rather than independently recruited participants; each record specifies which. Reused for the qualified aging records with each participant/outcome denominator and source test kept explicit. \\
\phantomsection\label{entry:symbol-registry-04392}\(p\) & \texttt{\detokenize{p}} & reported statistical significance value & p19-development-aging-statistics & object & \hyperref[sec:p19-r151-validity-record]{\ref*{sec:p19-r151-validity-record}} & \thepage & Not the probability that a developmental theory is true; missing intervals are not invented. Reused for the qualified aging records with each participant/outcome denominator and source test kept explicit. \\
\phantomsection\label{entry:symbol-registry-04393}\(\chi^2\) & \texttt{\detokenize{\chi^2}} & reported chi-square statistic & P19 rationality 15.1 & object & \hyperref[sec:p19-r151-validity-record]{\ref*{sec:p19-r151-validity-record}} & \thepage & Response dependence is retained; no reanalysis is claimed. \\
\phantomsection\label{entry:symbol-registry-04394}\(F\) & \texttt{\detokenize{F}} & reported analysis-of-variance statistic & P19 rationality 15.1 & object & \hyperref[sec:p19-r151-2016-record]{\ref*{sec:p19-r151-2016-record}} & \thepage & Degrees of freedom are retained as printed rather than converted into a fabricated exclusion record. \\
\phantomsection\label{entry:symbol-registry-04395}\(r\) & \texttt{\detokenize{r}} & reported correlation in the specified preschool study & P19 rationality 15.1 & object & \hyperref[sec:p19-r151-2017-record]{\ref*{sec:p19-r151-2017-record}} & \thepage & Partial-correlation controls and sample boundary stated locally. \\
\phantomsection\label{entry:symbol-registry-04396}\(\eta^2\) & \texttt{\detokenize{\eta^2}} & partial eta-squared effect size in Study 2 & P19 rationality 15.1 & object & \hyperref[sec:p19-r151-2017-record]{\ref*{sec:p19-r151-2017-record}} & \thepage & The reported value is .132; no unreported interval is supplied. \\
\phantomsection\label{entry:symbol-registry-04611}\(p,q\) & \texttt{\detokenize{p,}\allowbreak\detokenize{q}} & interpreted antecedent and consequent conditions in the selection and perspective discussion & P19 sec:p19-r112-selection & object & \hyperref[sec:p19-r112-selection]{\ref*{sec:p19-r112-selection}} & \thepage & Roles change explicitly from abstract conditions to the tidying/permission example; not statistical p-values. \\
\phantomsection\label{entry:symbol-registry-04612}\(\neg\) & \texttt{\detokenize{\neg}} & denial of a stipulated selection condition & P19 sec:p19-r112-selection & 60 & \hyperref[sec:p19-r112-selection]{\ref*{sec:p19-r112-selection}} & \thepage & This benchmark connective does not assert that conditional permissions are material implications. \\
\phantomsection\label{entry:symbol-registry-04613}\(p\) & \texttt{\detokenize{p}} & reported statistical p-value in an original-study record & P19 sec:p19-r112-cummins and sec:p19-r112-moral-record & object & \hyperref[sec:p19-r112-cummins]{\ref*{sec:p19-r112-cummins}} & \thepage & Declared statistical role distinct from conditional antecedents; not the probability a theory is true. \\
\phantomsection\label{entry:symbol-registry-04614}\(G^2\) & \texttt{\detokenize{G^2}} & reported likelihood-ratio test statistic in Cummins1996 & P19 sec:p19-r112-cummins & object & \hyperref[sec:p19-r112-cummins]{\ref*{sec:p19-r112-cummins}} & \thepage & Counts and nominal age groups retain each experiment's actual denominators and exclusions. \\
\phantomsection\label{entry:symbol-registry-04615}\(F\) & \texttt{\detokenize{F}} & reported analysis-of-variance statistic in Elqayam2017 & P19 sec:p19-r112-moral-record & object & \hyperref[sec:p19-r112-moral-record]{\ref*{sec:p19-r112-moral-record}} & \thepage & Printed adjusted degrees of freedom are not silently converted into fabricated missing-data counts. \\
\phantomsection\label{entry:symbol-registry-04616}\(\eta_p^2\) & \texttt{\detokenize{\eta_p^2}} & reported partial eta-squared effect size & P19 sec:p19-r112-moral-record & object & \hyperref[sec:p19-r112-moral-record]{\ref*{sec:p19-r112-moral-record}} & \thepage & Subscript p marks partial effect size; it is neither the antecedent proposition nor a p-value. \\
\phantomsection\label{entry:symbol-registry-04657}\(\eta_p^2\) & \texttt{\detokenize{\eta_p^2}} & reported partial eta-squared effect size & P19 sec:p19-r151-2020-author-record & object & \hyperref[sec:p19-r151-2020-author-record]{\ref*{sec:p19-r151-2020-author-record}} & \thepage & The subscript identifies partial effect size rather than a p-value. \\
\phantomsection\label{entry:symbol-registry-04658}\(t\) & \texttt{\detokenize{t}} & reported paired t-test statistic for likelihood versus confidence & P19 sec:p19-r151-2008-record & object & \hyperref[sec:p19-r151-2008-record]{\ref*{sec:p19-r151-2008-record}} & \thepage & A comparison of rating questions is not a calibration coefficient. \\
\phantomsection\label{entry:symbol-registry-04673}\(g\) & \texttt{\detokenize{g}} & general-intelligence construct in the cited CHC account & P19 rationality and intelligence & object & \hyperref[def:p19-rat152-grounding]{\ref*{def:p19-rat152-grounding}} & \thepage & Not an observed universal rationality quotient or an individual clinical classification. \\
\phantomsection\label{entry:symbol-registry-04674}\((A,B,C,D,E,F,G)\) & \texttt{\detokenize{A;B;C;D;E;F;G}} & source identifiers of functional processing and initiation arrows & P19 rationality and intelligence & meta & \hyperref[sec:p19-rat152-functions]{\ref*{sec:p19-rat152-functions}} & \thepage & Diagram path labels rather than scalar variables; names retain the roles of source Figure 15.2.3. \\
\phantomsection\label{entry:symbol-registry-04675}\(\ell\) & \texttt{\detokenize{\ell}} & positive side length of the exact outer Euclidean square & P19.rationality-13-3.midpoint-square & meta & \hyperref[fig:p19-rat133-midpoint-square]{\ref*{fig:p19-rat133-midpoint-square}} & \thepage & Physical drawing scale is incidental; midpoint conditions are stipulated. \\
\phantomsection\label{entry:symbol-registry-04676}\(A_{\mathrm{out}},A_{\mathrm{in}}\) & \texttt{\detokenize{A_{\mathrm{out}},}\allowbreak\detokenize{A_{\mathrm{in}}}} & outer and inner square areas & P19.rationality-13-3.midpoint-square & meta & \hyperref[thm:p19-rat133-square-area]{\ref*{thm:p19-rat133-square-area}} & \thepage & Area is geometrical here unlike cardinality-neutral area in set diagrams. \\
\phantomsection\label{entry:symbol-registry-04677}\(\mathrm{RT}\) & \texttt{\detokenize{\mathrm{RT}}} & mean response latency for correct responses & P19.rationality-13-3.original-2002 & meta & \hyperref[tab:p19-rat133-original-outcomes]{\ref*{tab:p19-rat133-original-outcomes}} & \thepage & Seconds in this table; does not isolate a mental imagery operation or include all trials. \\
\phantomsection\label{entry:symbol-registry-04678}\(z,Z,p\) & \texttt{\detokenize{z,}\allowbreak\detokenize{Z,}\allowbreak\detokenize{p}} & reported standardized test statistics and significance probabilities & P19.rationality-13-3.empirical-records & meta & \hyperref[sec:p19-rat133-impedance]{\ref*{sec:p19-rat133-impedance}} & \thepage & Original statistic casing retained; p is not probability that a theory or internal representation is true. \\
\phantomsection\label{entry:symbol-registry-04679}\(V,X,Z\) & \texttt{\detokenize{V,}\allowbreak\detokenize{X,}\allowbreak\detokenize{Z}} & displayed letter tokens in the 2006 spatial inference and maintenance task & P19.rationality-13-3.original-2006 & object & \hyperref[sec:p19-rat133-neural]{\ref*{sec:p19-rat133-neural}} & \thepage & Z as a displayed token differs from Z as a reported test statistic. \\
\phantomsection\label{entry:symbol-registry-04680}\(t\) & \texttt{\detokenize{t}} & reported imaging-map statistic threshold & P19.rationality-13-3.original-2006-analysis & meta & \hyperref[sec:p19-rat133-neural]{\ref*{sec:p19-rat133-neural}} & \thepage & t=3.0 is a height threshold; cluster correction and contrast definitions are separate. \\
\phantomsection\label{entry:symbol-registry-04681}\(d(x,y)\) & \texttt{\detokenize{d(x,}\allowbreak\detokenize{y)}} & Minkowski distance between coordinate points & P19 cognitive-knowledge-representation metric & n/a & \hyperref[def:p19-ckr-distance]{\ref*{def:p19-ckr-distance}} & \thepage & Published equation1 role; active expression corrects missing absolute values and specifies finite r>=1. \\
\phantomsection\label{entry:symbol-registry-04682}\(x,y\) & \texttt{\detokenize{x,}\allowbreak\detokenize{y}} & coordinate points in finite real semantic space & P19 cognitive-knowledge-representation metric & object & \hyperref[def:p19-ckr-distance]{\ref*{def:p19-ckr-distance}} & \thepage & Published point roles; not the question-prefixed slots in source relational schemata. \\
\phantomsection\label{entry:symbol-registry-04683}\(x_i,y_i\) & \texttt{\detokenize{x_i,}\allowbreak\detokenize{y_i}} & coordinate values of the two points & P19 cognitive-knowledge-representation metric & n/a & \hyperref[def:p19-ckr-distance]{\ref*{def:p19-ckr-distance}} & \thepage & Published equation1 coordinate values. \\
\phantomsection\label{entry:symbol-registry-04684}\(N,i\) & \texttt{\detokenize{N,}\allowbreak\detokenize{i}} & positive dimension count and coordinate index & P19 cognitive-knowledge-representation metric & meta & \hyperref[def:p19-ckr-distance]{\ref*{def:p19-ckr-distance}} & \thepage & Published N dimension count and index i; not a participant population size. \\
\phantomsection\label{entry:symbol-registry-04685}\(r\) & \texttt{\detokenize{r}} & finite Minkowski exponent & P19 cognitive-knowledge-representation metric & n/a & \hyperref[def:p19-ckr-distance]{\ref*{def:p19-ckr-distance}} & \thepage & Published parameter r; mathematical metric guarantee requires r>=1. \\
\phantomsection\label{entry:symbol-registry-04686}\(X\) & \texttt{\detokenize{X}} & count of distinct items with all unordered pairs rated & P19 cognitive-knowledge-representation pair count & n/a & \hyperref[thm:p19-ckr-pairs]{\ref*{thm:p19-ckr-pairs}} & \thepage & Published equation2 role; nonnegative integer and one rating per unordered distinct pair. \\
\phantomsection\label{entry:symbol-registry-04687}\(?x,?y\) & \texttt{\detokenize{?x,}\allowbreak\detokenize{?y}} & argument slots in cognitive relational schemata & P19 cognitive-knowledge-representation schemata & object & \hyperref[sec:p19-ckr-structured]{\ref*{sec:p19-ckr-structured}} & \thepage & Published question-prefixed variables distinguish role fillers; no automatic first-order syntax declaration. \\
\phantomsection\label{entry:symbol-registry-04688}\(\operatorname{kind\!-\!of}\) & \texttt{\detokenize{\operatorname{kind\!-\!of}}} & binary kind-inclusion cognitive relation & P19 cognitive-knowledge-representation schemata & relation & \hyperref[sec:p19-ckr-structured]{\ref*{sec:p19-ckr-structured}} & \thepage & Published ordered-role name from printed43; kind inclusion is not symmetric. \\
\phantomsection\label{entry:symbol-registry-04689}\(\operatorname{red}\) & \texttt{\detokenize{\operatorname{red}}} & unary cognitive color attribute & P19 cognitive-knowledge-representation schemata & relation & \hyperref[sec:p19-ckr-structured]{\ref*{sec:p19-ckr-structured}} & \thepage & Published attribute at printed43; its argument fixes the object's property scope. \\
\phantomsection\label{entry:symbol-registry-04690}\(\operatorname{above}\) & \texttt{\detokenize{\operatorname{above}}} & binary ordered spatial relation in the cognitive scene graph & P19 cognitive-knowledge-representation schemata & relation & \hyperref[sec:p19-ckr-structured]{\ref*{sec:p19-ckr-structured}} & \thepage & Published Figure4.3 relation; a bound graph requires an actual query process. \\
\phantomsection\label{entry:symbol-registry-04691}\(\operatorname{revolve\!-\!around}\) & \texttt{\detokenize{\operatorname{revolve\!-\!around}}} & binary ordered relation in the source analogy schema & P19 cognitive-knowledge-representation schemata & relation & \hyperref[sec:p19-ckr-structured]{\ref*{sec:p19-ckr-structured}} & \thepage & Published orbiting versus orbited roles at46; not a modern physical guarantee. \\
\phantomsection\label{entry:symbol-registry-04692}\(\operatorname{greater}\) & \texttt{\detokenize{\operatorname{greater}}} & binary comparison of represented mass quantities & P19 cognitive-knowledge-representation schemata & relation & \hyperref[sec:p19-ckr-structured]{\ref*{sec:p19-ckr-structured}} & \thepage & Published source analogy relation between quantitative attributes. \\
\phantomsection\label{entry:symbol-registry-04693}\(\operatorname{mass}\) & \texttt{\detokenize{\operatorname{mass}}} & quantitative attribute function in the source analogy schema & P19 cognitive-knowledge-representation schemata & meta & \hyperref[sec:p19-ckr-structured]{\ref*{sec:p19-ckr-structured}} & \thepage & Published function maps represented objects to quantities; distinct from a relation on objects. \\
\phantomsection\label{entry:symbol-registry-04694}\(\operatorname{cause}\) & \texttt{\detokenize{\operatorname{cause}}} & higher-order causal schema taking relational contents as arguments & P19 cognitive-knowledge-representation schemata & relation & \hyperref[sec:p19-ckr-structured]{\ref*{sec:p19-ckr-structured}} & \thepage & Published source46 schema; not silently a well-formed first-order formula. \\
\phantomsection\label{entry:symbol-registry-04695}\(\mathrm{electron},\mathrm{nucleus},\mathrm{sun},\mathrm{planet}\) & \texttt{\detokenize{\mathrm{electron},}\allowbreak\detokenize{\mathrm{nucleus},}\allowbreak\detokenize{\mathrm{sun},}\allowbreak\detokenize{\mathrm{planet}}} & object role-fillers in the source analogy schemata & P19 cognitive-knowledge-representation schemata & object & \hyperref[sec:p19-ckr-structured]{\ref*{sec:p19-ckr-structured}} & \thepage & Published domain names preserve base/target and argument ordering. \\
\phantomsection\label{entry:symbol-registry-04696}\(\lvert a\rvert\) & \texttt{\detokenize{\lvert a\rvert}} & absolute magnitude of a real coordinate value & P19 cognitive-knowledge-representation metric & n/a & \hyperref[def:p19-ckr-distance]{\ref*{def:p19-ckr-distance}} & \thepage & Local typed absolute-value role introduced in the correction of source equation1; not cardinality. Also used in its proof. \\
\phantomsection\label{entry:symbol-registry-04697}\(\mathcal S,\mathcal R,\mathcal A\) & \texttt{\detokenize{\mathcal S,}\allowbreak\detokenize{\mathcal R,}\allowbreak\detokenize{\mathcal A}} & scene code and task-answer domains & P19 cognitive-knowledge-representation decoding & object & \hyperref[def:p19-ckr-decoding]{\ref*{def:p19-ckr-decoding}} & \thepage & Editorial finite information comparison, not measured brain states. \\
\phantomsection\label{entry:symbol-registry-04698}\(e,q,g\) & \texttt{\detokenize{e,}\allowbreak\detokenize{q,}\allowbreak\detokenize{g}} & encoding correct-answer map and decoder on used codes & P19 cognitive-knowledge-representation decoding & meta & \hyperref[def:p19-ckr-decoding]{\ref*{def:p19-ckr-decoding}} & \thepage & Editorial maps; decoder exists exactly when task answers are constant on encoding fibers. No efficient cognitive implementation follows. \\
\phantomsection\label{entry:symbol-registry-04699}\(s,t,r\) & \texttt{\detokenize{s,}\allowbreak\detokenize{t,}\allowbreak\detokenize{r}} & scene instances and an actually used encoding value in the decoder proof & P19 cognitive-knowledge-representation decoding & object & \hyperref[annex:p19.ckr-decoding]{\ref*{annex:p19.ckr-decoding}} & \thepage & Editorial proof variables; r here is a code rather than a Minkowski exponent. \\
\phantomsection\label{entry:symbol-registry-04700}\(a,b,c_i,u,v,s\) & \texttt{\detokenize{a,}\allowbreak\detokenize{b,}\allowbreak\detokenize{c_i,}\allowbreak\detokenize{u,}\allowbreak\detokenize{v,}\allowbreak\detokenize{s}} & finite vectors coordinate magnitudes Young variables and conjugate exponent in the metric proof & P19 cognitive-knowledge-representation metric proof & meta & \hyperref[annex:p19.ckr-metric]{\ref*{annex:p19.ckr-metric}} & \thepage & Editorial local proof variables; s=r/(r-1) only when r>1. \\
\phantomsection\label{entry:symbol-registry-04701}\(\lVert a\rVert_r\) & \texttt{\detokenize{\lVert a\rVert_r}} & finite coordinate r norm in the metric proof & P19 cognitive-knowledge-representation metric proof & n/a & \hyperref[annex:p19.ckr-metric]{\ref*{annex:p19.ckr-metric}} & \thepage & Editorial abbreviation for the corrected source coordinate sum. \\
\phantomsection\label{entry:symbol-registry-04702}\(r,r^{(i)}\) & \texttt{\detokenize{r,}\allowbreak\detokenize{r^{(i)}}} & real role vectors and indexed role vectors & P19 computational-cognitive-models & meta & \hyperref[def:p19-ccm-tensors]{\ref*{def:p19-ccm-tensors}} & \thepage & Source tensor role factors; distinct from relational predicate symbols. \\
\phantomsection\label{entry:symbol-registry-04703}\(f,f^{(i)},g\) & \texttt{\detokenize{f,}\allowbreak\detokenize{f^{(i)},}\allowbreak\detokenize{g}} & real filler vectors for finite product bindings & P19 computational-cognitive-models & meta & \hyperref[def:p19-ccm-tensors]{\ref*{def:p19-ccm-tensors}} & \thepage & Source filler factors; need not denote network activation functions. \\
\phantomsection\label{entry:symbol-registry-04704}\(m,n,s\) & \texttt{\detokenize{m,}\allowbreak\detokenize{n,}\allowbreak\detokenize{s}} & positive finite vector dimensions & P19 computational-cognitive-models & meta & \hyperref[def:p19-ccm-tensors]{\ref*{def:p19-ccm-tensors}} & \thepage & r has dimension m; f dimension n; g dimension s. \\
\phantomsection\label{entry:symbol-registry-04705}\(i,j,k\) & \texttt{\detokenize{i,}\allowbreak\detokenize{j,}\allowbreak\detokenize{k}} & finite coordinate and binding indices & P19 computational-cognitive-models & meta & \hyperref[def:p19-ccm-tensors]{\ref*{def:p19-ccm-tensors}} & \thepage & Coordinate indices in tensor equations; k also explicitly counts an orthonormal role family in the unbinding theorem. \\
\phantomsection\label{entry:symbol-registry-04706}\(\otimes\) & \texttt{\detokenize{\otimes}} & finite real outer product for role-filler binding & P19 computational-cognitive-models & infix & \hyperref[def:p19-ccm-tensors]{\ref*{def:p19-ccm-tensors}} & \thepage & Typed tensor product; not the linear-logic tensor connective. Source rf and rfg labels are made coordinate-explicit. \\
\phantomsection\label{entry:symbol-registry-04707}\(rf,rfg\) & \texttt{\detokenize{rf,}\allowbreak\detokenize{rfg}} & source shorthand labels for bound two-factor and three-factor tensors & P19 computational-cognitive-models & meta & \hyperref[def:p19-ccm-tensors]{\ref*{def:p19-ccm-tensors}} & \thepage & Printed equations1-2; not untyped multiplication of scalar variables. \\
\phantomsection\label{entry:symbol-registry-04708}\(T\) & \texttt{\detokenize{T}} & finite matrix sum of indexed role-filler outer products & P19 computational-cognitive-models & meta & \hyperref[def:p19-ccm-tensors]{\ref*{def:p19-ccm-tensors}} & \thepage & Editorial typed spelling of source summed bindings; shape m by n. \\
\phantomsection\label{entry:symbol-registry-04709}\(\cdot\) & \texttt{\detokenize{\cdot}} & Euclidean dot product of matching finite vectors or flattened tensors & P19 computational-cognitive-models & infix & \hyperref[thm:p19-ccm-tensor-similarity]{\ref*{thm:p19-ccm-tensor-similarity}} & \thepage & The source's scalar similarity product; distinguished from role binding. \\
\phantomsection\label{entry:symbol-registry-04710}\(\|\,\cdot\,\|_2\) & \texttt{\detokenize{\|\,}\allowbreak\detokenize{\cdot\,}\allowbreak\detokenize{\|_2}} & Euclidean norm of a finite real vector or flattened tensor & P19 computational-cognitive-models & meta & \hyperref[thm:p19-ccm-tensor-similarity]{\ref*{thm:p19-ccm-tensor-similarity}} & \thepage & Square root of self dot product; zero norm makes cosine undefined. \\
\phantomsection\label{entry:symbol-registry-04711}\(\operatorname{cos}\) & \texttt{\detokenize{\operatorname{cos}}} & cosine similarity of two nonzero matching vectors & P19 computational-cognitive-models & meta & \hyperref[thm:p19-ccm-tensor-similarity]{\ref*{thm:p19-ccm-tensor-similarity}} & \thepage & Dot product divided by product of Euclidean norms; not an angle measurement supplied by a simulation. \\
\phantomsection\label{entry:symbol-registry-04712}\(\mathsf T\) & \texttt{\detokenize{\mathsf T}} & matrix or vector transpose marker & P19 computational-cognitive-models & meta & \hyperref[thm:p19-ccm-unbinding]{\ref*{thm:p19-ccm-unbinding}} & \thepage & T with transpose multiplies the selected role to decode its filler under orthonormality. \\
\phantomsection\label{entry:symbol-registry-04828}\(\kappa\) & \texttt{\detokenize{\kappa}} & reported Cohen inter-rater agreement coefficient & P19 sec:p19-r151-kuhn-udell and sec:p19-r151-kuhn-crowell & object & \hyperref[sec:p19-r151-kuhn-udell]{\ref*{sec:p19-r151-kuhn-udell}} & \thepage & Agreement concerns a specified coding rubric rather than ethical truth or a universal skill measure. \\
\phantomsection\label{entry:symbol-registry-04829}\(t\) & \texttt{\detokenize{t}} & reported t-test statistic for the specified argumentation comparison & P19 sec:p19-r151-kuhn-udell and sec:p19-r151-kuhn-crowell & object & \hyperref[sec:p19-r151-kuhn-udell]{\ref*{sec:p19-r151-kuhn-udell}} & \thepage & Repeated and overlapping datasets are not summed as independent recruits. \\
\phantomsection\label{entry:symbol-registry-04830}\(\eta_p^2\) & \texttt{\detokenize{\eta_p^2}} & reported partial eta-squared effect size & P19 sec:p19-r151-kuhn-crowell & object & \hyperref[sec:p19-r151-kuhn-crowell]{\ref*{sec:p19-r151-kuhn-crowell}} & \thepage & Subscript p denotes partial effect size rather than the p-value. \\
\phantomsection\label{entry:symbol-registry-04831}\(H\) & \texttt{\detokenize{H}} & finite hypothesis space & P19 Bayesian cognitive inference & object & \hyperref[def:p19-bci-model]{\ref*{def:p19-bci-model}} & \thepage & Can contain durations categories graphs or grammars only when their representation is specified. \\
\phantomsection\label{entry:symbol-registry-04832}\(\mathcal D\) & \texttt{\detokenize{\mathcal D}} & finite data alphabet & P19 Bayesian cognitive inference & object & \hyperref[def:p19-bci-model]{\ref*{def:p19-bci-model}} & \thepage & Separate from source D observed structure-learning matrix. \\
\phantomsection\label{entry:symbol-registry-04833}\((h,d)\) & \texttt{\detokenize{h;d}} & hypothesis and observed datum & P19 Bayesian cognitive inference & object & \hyperref[def:p19-bci-model]{\ref*{def:p19-bci-model}} & \thepage & Data evidence must have positive prior predictive mass before ratio conditioning. \\
\phantomsection\label{entry:symbol-registry-04834}\(\pi(h)\) & \texttt{\detokenize{\pi(h)}} & prior hypothesis mass & P19 Bayesian cognitive inference & object & \hyperref[def:p19-bci-model]{\ref*{def:p19-bci-model}} & \thepage & Source P(h) in the finite benchmark; not a prior supplied by Bayes identity itself. \\
\phantomsection\label{entry:symbol-registry-04835}\(L(d\mid h)\) & \texttt{\detokenize{L(d\mid h)}} & hypothesis-conditioned data likelihood & P19 Bayesian cognitive inference & object & \hyperref[def:p19-bci-model]{\ref*{def:p19-bci-model}} & \thepage & Normalized over data for each hypothesis rather than over hypotheses after observation. \\
\phantomsection\label{entry:symbol-registry-04836}\(m(d)\) & \texttt{\detokenize{m(d)}} & prior predictive data mass & P19 Bayesian cognitive inference & object & \hyperref[def:p19-bci-model]{\ref*{def:p19-bci-model}} & \thepage & Zero mass leaves the displayed ratio undefined. \\
\phantomsection\label{entry:symbol-registry-04837}\(q_d(h)\) & \texttt{\detokenize{q_d(h)}} & posterior hypothesis mass at supported datum & P19 Bayesian cognitive inference & object & \hyperref[def:p19-bci-model]{\ref*{def:p19-bci-model}} & \thepage & Mathematical distribution distinct from a psychological selection rule. \\
\phantomsection\label{entry:symbol-registry-04838}\(K(y\mid h,d)\) & \texttt{\detokenize{K(y\mid h;d)}} & hypothesis-and-data-conditioned future probability kernel & P19 Bayesian cognitive inference & object & \hyperref[thm:p19-bci-bayes]{\ref*{thm:p19-bci-bayes}} & \thepage & Finite future alphabet and normalized conditional masses are explicit. \\
\phantomsection\label{entry:symbol-registry-04839}\(y\) & \texttt{\detokenize{y}} & future outcome in the finite predictive-kernel benchmark & P19 Bayesian cognitive inference & object & \hyperref[thm:p19-bci-bayes]{\ref*{thm:p19-bci-bayes}} & \thepage & Local role distinct from the response-bar magnitude in the original function-learning task. \\
\phantomsection\label{entry:symbol-registry-04840}\((F,G,D)\) & \texttt{\detokenize{F;G;D}} & source form grammar graph and observed data matrix & P19 Bayesian cognitive inference & object & \hyperref[sec:p19-bci-representations]{\ref*{sec:p19-bci-representations}} & \thepage & Source-specific structure-learning levels not one unstructured universal hypothesis space. \\
\phantomsection\label{entry:symbol-registry-04841}\(n\) & \texttt{\detokenize{n}} & number of distinct object leaves in tree restriction & P19 Bayesian cognitive inference & object & \hyperref[thm:p19-bci-tree-count]{\ref*{thm:p19-bci-tree-count}} & \thepage & All subsets and internal nonsingleton clades use different candidate conventions. \\
\phantomsection\label{entry:symbol-registry-04842}\(Q_{ij}\) & \texttt{\detokenize{Q_{ij}}} & iterated-learning next-i-given-current-j transition mass & P19 Bayesian cognitive inference & object & \hyperref[thm:p19-bci-iterated]{\ref*{thm:p19-bci-iterated}} & \thepage & Original destination-first column-stochastic orientation is retained. \\
\phantomsection\label{entry:symbol-registry-04843}\((i,j)\) & \texttt{\detokenize{i;j}} & next and current hypothesis indices & P19 Bayesian cognitive inference & object & \hyperref[thm:p19-bci-iterated]{\ref*{thm:p19-bci-iterated}} & \thepage & Separate from physical positions or individual learners' cognitive measurements. \\
\phantomsection\label{entry:symbol-registry-04844}\((x,y)\) & \texttt{\detokenize{x;y}} & function-learning input and produced response magnitudes & P19 Bayesian cognitive inference & object & \hyperref[sec:p19-bci-iterated]{\ref*{sec:p19-bci-iterated}} & \thepage & Source bar task values and rounding conditions are explicit. \\
\phantomsection\label{entry:symbol-registry-04845}\((C_1,C_2)\) & \texttt{\detokenize{C_1;C_2}} & hypothetical causal variable classes in source schema & P19 Bayesian cognitive inference & object & \hyperref[sec:p19-bci-hierarchies]{\ref*{sec:p19-bci-hierarchies}} & \thepage & Computational illustration not a validated clinical taxonomy. \\
\phantomsection\label{entry:symbol-registry-04846}\((\mu,\nu)\) & \texttt{\detokenize{\mu;\nu}} & probability columns over the finite hypothesis space & P19 Bayesian cognitive inference & object & \hyperref[annex:p19.bci.iterated]{\ref*{annex:p19.bci.iterated}} & \thepage & Used in the full geometric contraction proof rather than empirical sample distributions. \\
\phantomsection\label{entry:symbol-registry-04847}\(\varepsilon\) & \texttt{\detokenize{\varepsilon}} & positive common-column minorization mass & P19 Bayesian cognitive inference & object & \hyperref[annex:p19.bci.iterated]{\ref*{annex:p19.bci.iterated}} & \thepage & Equals the sum of rowwise minima of the strictly positive transition matrix. \\
\phantomsection\label{entry:symbol-registry-04848}\(\lVert v\rVert_1\) & \texttt{\detokenize{\lVert v\rVert_1}} & sum of absolute vector components & P19 Bayesian cognitive inference & meta & \hyperref[annex:p19.bci.iterated]{\ref*{annex:p19.bci.iterated}} & \thepage & Declared contraction distance for finite probability columns; not a psychological score. \\
\phantomsection\label{entry:symbol-registry-04849}\(J(h,d)\) & \texttt{\detokenize{J(h;d)}} & joint hypothesis-data mass in the finite model & P19 Bayesian cognitive inference & object & \hyperref[annex:p19.bci.bayes]{\ref*{annex:p19.bci.bayes}} & \thepage & Defined as prior mass times the likelihood before conditioning. \\
\phantomsection\label{entry:symbol-registry-04850}\(I\) & \texttt{\detokenize{I}} & number of internal nodes in a finite rooted full binary tree & P19 Bayesian cognitive inference & object & \hyperref[annex:p19.bci.tree-count]{\ref*{annex:p19.bci.tree-count}} & \thepage & Distinct from the named hypothesis indices i and j. \\
\phantomsection\label{entry:symbol-registry-04851}\((a_i,r_i)\) & \texttt{\detokenize{a_i;r_i}} & row minimum and its normalized common-column mass & P19 Bayesian cognitive inference & object & \hyperref[annex:p19.bci.iterated]{\ref*{annex:p19.bci.iterated}} & \thepage & Strict positivity gives positive minima and a probability column after normalization. \\
\phantomsection\label{entry:symbol-registry-04852}\(R_{ij}\) & \texttt{\detokenize{R_{ij}}} & residual column-stochastic transition mass & P19 Bayesian cognitive inference & object & \hyperref[annex:p19.bci.iterated]{\ref*{annex:p19.bci.iterated}} & \thepage & Defined only when the common-column mass is less than one; the equal-one case is handled separately. \\
\phantomsection\label{entry:symbol-registry-04853}\(k\) & \texttt{\detokenize{k}} & positive number of transmission transitions & P19 Bayesian cognitive inference & object & \hyperref[annex:p19.bci.iterated]{\ref*{annex:p19.bci.iterated}} & \thepage & Finite-chain convergence parameter rather than a human population measure. \\
\phantomsection\label{entry:symbol-registry-04854}\(s_\mu(d)\) & \texttt{\detokenize{s_\mu(d)}} & produced-data mass under hypothesis probability column mu & P19 Bayesian cognitive inference & object & \hyperref[annex:p19.bci.iterated]{\ref*{annex:p19.bci.iterated}} & \thepage & Likelihood mixture whose distance from prior predictive mass is bounded in the proof. \\
\phantomsection\label{entry:symbol-registry-04855}\(K\) & \texttt{\detokenize{K}} & visual-memory capacity estimate under the primary change-detection model & P19 neurocognitive-methods capacity estimator & object & \hyperref[def:p19-ncm-scores]{\ref*{def:p19-ncm-scores}} & \thepage & Published Vogel2005 quantity; not direct neural slot count. \\
\phantomsection\label{entry:symbol-registry-04856}\(S\) & \texttt{\detokenize{S}} & number of items in the memory array & P19 neurocognitive-methods capacity estimator & object & \hyperref[def:p19-ncm-scores]{\ref*{def:p19-ncm-scores}} & \thepage & Published positive array-size parameter,not independent participant count. \\
\phantomsection\label{entry:symbol-registry-04857}\(H,F\) & \texttt{\detokenize{H,}\allowbreak\detokenize{F}} & hit and false-alarm proportions & P19 neurocognitive-methods capacity estimator & object & \hyperref[def:p19-ncm-scores]{\ref*{def:p19-ncm-scores}} & \thepage & Published primary roles in K=S(H-F); F here is not four-item CDA amplitude. \\
\phantomsection\label{entry:symbol-registry-04858}\(\alpha\) & \texttt{\detokenize{\alpha}} & relative EEG filtering-efficiency score & P19 neurocognitive-methods filtering score & object & \hyperref[def:p19-ncm-scores]{\ref*{def:p19-ncm-scores}} & \thepage & Published Greek alpha visually verified on502; range not guaranteed when amplitude falls outside endpoints. \\
\phantomsection\label{entry:symbol-registry-04859}\(F,D,T\) & \texttt{\detokenize{F,}\allowbreak\detokenize{D,}\allowbreak\detokenize{T}} & four-item distractor-present and two-item mean CDA amplitudes & P19 neurocognitive-methods filtering score & object & \hyperref[def:p19-ncm-scores]{\ref*{def:p19-ncm-scores}} & \thepage & Published roles for the ratio (F-D)/(F-T); F is amplitude rather than a false-alarm proportion. \\
\phantomsection\label{entry:symbol-registry-04860}\(r\) & \texttt{\detokenize{r}} & reported correlation coefficient between behavioral capacity and filtering score & P19 neurocognitive-methods EEG reporting & object & \hyperref[sec:p19-ncm-measurements]{\ref*{sec:p19-ncm-measurements}} & \thepage & Primary reported0.69,0.62,0.57 correlations; not causal effects or cortical locations. \\
\phantomsection\label{entry:symbol-registry-04861}\(p\) & \texttt{\detokenize{p}} & reported statistical test tail probability & P19 neurocognitive-methods empirical reporting & object & \hyperref[sec:p19-ncm-measurements]{\ref*{sec:p19-ncm-measurements}} & \thepage & Source-reported significance values; no posterior probability or confidence interval interpretation. \\
\phantomsection\label{entry:symbol-registry-04862}\(F(1,19)\) & \texttt{\detokenize{F(1,}\allowbreak\detokenize{19)}} & reported analysis-of-variance contrast statistic and degrees of freedom & P19 neurocognitive-methods analogy reporting & object & \hyperref[sec:p19-ncm-analogy]{\ref*{sec:p19-ncm-analogy}} & \thepage & Published Morrison2004 statistic; does not by itself identify the independent sampling or analysis unit. \\
\phantomsection\label{entry:symbol-registry-04920}\(H_1\) & \texttt{\detokenize{H_1}} & first hypothesis & p19-hph-odds & object & \hyperref[def:p19-hph-odds-domain]{\ref*{def:p19-hph-odds-domain}} & \thepage & Positive-prior member of a two-hypothesis partition \\
\phantomsection\label{entry:symbol-registry-04921}\(H_2\) & \texttt{\detokenize{H_2}} & second hypothesis & p19-hph-odds & object & \hyperref[def:p19-hph-odds-domain]{\ref*{def:p19-hph-odds-domain}} & \thepage & Positive-prior member with positive likelihood for the evidence \\
\phantomsection\label{entry:symbol-registry-04922}\(D\) & \texttt{\detokenize{D}} & evidence event & p19-hph-odds & object & \hyperref[def:p19-hph-odds-domain]{\ref*{def:p19-hph-odds-domain}} & \thepage & Positive probability and common evidence in both likelihoods \\
\phantomsection\label{entry:symbol-registry-04923}\(P\) & \texttt{\detokenize{P}} & probability law & p19-hph-odds & object & \hyperref[def:p19-hph-odds-domain]{\ref*{def:p19-hph-odds-domain}} & \thepage & Normative law not an elicited response \\
\phantomsection\label{entry:symbol-registry-04924}\(o_0\) & \texttt{\detokenize{o_0}} & prior odds & p19-hph-odds & object & \hyperref[def:p19-hph-odds-domain]{\ref*{def:p19-hph-odds-domain}} & \thepage & Ratio P(H\_1)/P(H\_2) \\
\phantomsection\label{entry:symbol-registry-04925}\(\ell\) & \texttt{\detokenize{\ell}} & likelihood ratio & p19-hph-odds & object & \hyperref[def:p19-hph-odds-domain]{\ref*{def:p19-hph-odds-domain}} & \thepage & Ratio of the evidence likelihoods with positive denominator \\
\phantomsection\label{entry:symbol-registry-04926}\(x\) & \texttt{\detokenize{x}} & posterior probability used in the odds proof & p19-hph-odds & object & \hyperref[annex:p19.hph.odds]{\ref*{annex:p19.hph.odds}} & \thepage & Auxiliary x equals P(H\_1 given D) \\
\phantomsection\label{entry:symbol-registry-04927}\(p\) & \texttt{\detokenize{p}} & girl probability & p19-hph-patterns & object & \hyperref[thm:p19-hph-patterns]{\ref*{thm:p19-hph-patterns}} & \thepage & Fixed independent-birth parameter in (0;1) \\
\phantomsection\label{entry:symbol-registry-04928}\(P\) & \texttt{\detokenize{P}} & sampling law & p19-hph-patterns & object & \hyperref[annex:p19.hph.patterns]{\ref*{annex:p19.hph.patterns}} & \thepage & Each toy comparison supplies its own declared sampling assumptions \\
\phantomsection\label{entry:symbol-registry-04929}\(q\) & \texttt{\detokenize{q}} & red outcome probability & p19-hph-patterns & object & \hyperref[thm:p19-hph-patterns]{\ref*{thm:p19-hph-patterns}} & \thepage & Stipulated .7 and not a participant judgment \\
\phantomsection\label{entry:symbol-registry-04930}\(a\) & \texttt{\detokenize{a}} & red choice probability & p19-hph-patterns & object & \hyperref[thm:p19-hph-patterns]{\ref*{thm:p19-hph-patterns}} & \thepage & Choice independent of the trial outcome \\
\phantomsection\label{entry:symbol-registry-04931}\(R\) & \texttt{\detokenize{R}} & red outcome event & p19-hph-patterns & object & \hyperref[annex:p19.hph.patterns]{\ref*{annex:p19.hph.patterns}} & \thepage & Independent of the choice event C \\
\phantomsection\label{entry:symbol-registry-04932}\(C\) & \texttt{\detokenize{C}} & red choice event & p19-hph-patterns & object & \hyperref[annex:p19.hph.patterns]{\ref*{annex:p19.hph.patterns}} & \thepage & Not the calibration-confidence variable at the separate corpus interface \\
\phantomsection\label{entry:symbol-registry-04933}\(r\) & \texttt{\detokenize{r}} & positive red-draw reward & p19-hph-patterns & object & \hyperref[annex:p19.hph.patterns]{\ref*{annex:p19.hph.patterns}} & \thepage & Same reward in each urn comparison \\
\phantomsection\label{entry:symbol-registry-04934}\(A\) & \texttt{\detokenize{A}} & first description & p19-hph-support & object & \hyperref[def:p19-hph-support]{\ref*{def:p19-hph-support}} & \thepage & Support-theory description not an arbitrary extensional event \\
\phantomsection\label{entry:symbol-registry-04935}\(B\) & \texttt{\detokenize{B}} & second description & p19-hph-support & object & \hyperref[def:p19-hph-support]{\ref*{def:p19-hph-support}} & \thepage & Compared description with positive total support \\
\phantomsection\label{entry:symbol-registry-04936}\(s\) & \texttt{\detokenize{s}} & support score function & p19-hph-support & object & \hyperref[def:p19-hph-support]{\ref*{def:p19-hph-support}} & \thepage & Nonnegative scores under the descriptive model \\
\phantomsection\label{entry:symbol-registry-04937}\(P\) & \texttt{\detokenize{P}} & relative judgment function & p19-hph-support & object & \hyperref[def:p19-hph-support]{\ref*{def:p19-hph-support}} & \thepage & Source two-description syntax with locally distinct type from the coherent probability law \\
\phantomsection\label{entry:symbol-registry-04938}\(x,x'\) & \texttt{\detokenize{x;x'}} & first-description scores in the support proof & p19-hph-support & object & \hyperref[annex:p19.hph.support]{\ref*{annex:p19.hph.support}} & \thepage & Auxiliary score values not odds-proof posterior probabilities \\
\phantomsection\label{entry:symbol-registry-04939}\(y,y'\) & \texttt{\detokenize{y;y'}} & second-description scores in the support proof & p19-hph-support & object & \hyperref[annex:p19.hph.support]{\ref*{annex:p19.hph.support}} & \thepage & Auxiliary nonnegative score values \\
\phantomsection\label{entry:symbol-registry-05336}\(n\) & \texttt{\detokenize{n}} & number of students in the named experimental condition & P19 self-explanation and learning & object & \hyperref[tab:p19-self-explanation-diffusion]{\ref*{tab:p19-self-explanation-diffusion}} & \thepage & Participant count not number of coded utterances or test questions. \\
\phantomsection\label{entry:symbol-registry-05337}\(r\) & \texttt{\detokenize{r}} & reported within-sample correlation coefficient & P19 self-explanation and learning & object & \hyperref[sec:p19-self-explanation-physics]{\ref*{sec:p19-self-explanation-physics}} & \thepage & Association is distinguished from a randomized explanation-type effect. \\
\phantomsection\label{entry:symbol-registry-05338}\(t\) & \texttt{\detokenize{t}} & reported t-test statistic with source degrees of freedom & P19 self-explanation and learning & object & \hyperref[sec:p19-self-explanation-biology]{\ref*{sec:p19-self-explanation-biology}} & \thepage & Source-reported statistic; no raw-data reconstruction is claimed. \\
\phantomsection\label{entry:symbol-registry-05339}\(F\) & \texttt{\detokenize{F}} & reported analysis-of-variance or covariance statistic & P19 self-explanation and learning & object & \hyperref[sec:p19-self-explanation-biology]{\ref*{sec:p19-self-explanation-biology}} & \thepage & Source analysis and degrees of freedom remain attached to each comparison. \\
\phantomsection\label{entry:symbol-registry-05340}\(p\) & \texttt{\detokenize{p}} & source-reported p-value or threshold & P19 self-explanation and learning & object & \hyperref[sec:p19-self-explanation-physics]{\ref*{sec:p19-self-explanation-physics}} & \thepage & Not a probability that the mechanism or null hypothesis is true. \\
\phantomsection\label{entry:symbol-registry-05341}\(\eta^2\) & \texttt{\detokenize{\eta^2}} & reported eta-squared effect-size label & local:sec:p19-self-explanation-diffusion & object & \hyperref[tab:p19-self-explanation-diffusion]{\ref*{tab:p19-self-explanation-diffusion}} & \thepage & The source does not identify this as partial eta-squared; preserve its reported label and value. \\
\phantomsection\label{entry:symbol-registry-05342}\(\beta_1\) & \texttt{\detokenize{\beta_1}} & fitted age coefficient in the pilot-performance model & p19-aging-pilot-model & meta & \hyperref[sec:p19-aging-expertise]{\ref*{sec:p19-aging-expertise}} & \thepage & Age centered at the median; baseline-level and annual-slope models are distinct \\
\phantomsection\label{entry:symbol-registry-05343}\(\beta_2\) & \texttt{\detokenize{\beta_2}} & fitted ordinal-expertise coefficient in the pilot-performance model & p19-aging-pilot-model & meta & \hyperref[sec:p19-aging-expertise]{\ref*{sec:p19-aging-expertise}} & \thepage & Expertise coded minus one zero and one; not randomized training or an age-by-expertise interaction \\
\phantomsection\label{entry:symbol-registry-05562}\(p\) & \texttt{\detokenize{p}} & antecedent category property in the selection example & P19 motivated reasoning selection example & object & \hyperref[def:p19-mrg-selection-case]{\ref*{def:p19-mrg-selection-case}} & \thepage & Interpreted properties in an exceptionless descriptive rule. Root-authorized typed split; original grouped intake row remains unchanged as evidence. \\
\phantomsection\label{entry:symbol-registry-05563}\(q\) & \texttt{\detokenize{q}} & consequent category property in the selection example & P19 motivated reasoning selection example & object & \hyperref[def:p19-mrg-selection-case]{\ref*{def:p19-mrg-selection-case}} & \thepage & Interpreted properties in an exceptionless descriptive rule. Root-authorized typed split; original grouped intake row remains unchanged as evidence. \\
\phantomsection\label{entry:symbol-registry-05564}\(\neg p,\neg q\) & \texttt{\detokenize{\neg}} & denial of an interpreted category property & P19 motivated reasoning selection example & n/a & \hyperref[def:p19-mrg-selection-case]{\ref*{def:p19-mrg-selection-case}} & \thepage & Not preference rejection or an emotional response. Original precedence descriptor 'connective' preserved as semantic metadata; n/a is not an invented binding rank. \\
\phantomsection\label{entry:symbol-registry-05565}\(\chi^2\) & \texttt{\detokenize{\chi^2}} & reported Pearson comparison statistic & P19 motivated reasoning Wason evidence & meta & \hyperref[sec:p19-mrg-directional]{\ref*{sec:p19-mrg-directional}} & \thepage & Source statistic with displayed degrees of freedom and denominator; no recomputation. \\
\phantomsection\label{entry:symbol-registry-05566}\(N,n\) & \texttt{\detokenize{N;n}} & analysis total and group participant counts & P19 motivated reasoning empirical records & meta & \hyperref[sec:p19-mrg-directional]{\ref*{sec:p19-mrg-directional}} & \thepage & Trial counts are not participant counts. \\
\phantomsection\label{entry:symbol-registry-05567}\(p\) & \texttt{\detokenize{p}} & reported significance probability & P19 motivated reasoning empirical records & meta & \hyperref[sec:p19-mrg-directional]{\ref*{sec:p19-mrg-directional}} & \thepage & Distinct role from the selection-rule antecedent; not an effect size. \\
\phantomsection\label{entry:symbol-registry-05568}\(\beta\) & \texttt{\detokenize{\beta}} & source response-criterion statistic & P19 motivated reasoning recognition record & meta & \hyperref[def:p19-mrg-four-outcomes]{\ref*{def:p19-mrg-four-outcomes}} & \thepage & Below one favors yes and above one no in the source convention; not discrimination ability. \\
\phantomsection\label{entry:symbol-registry-05569}\(d'\) & \texttt{\detokenize{d'}} & reported recognition sensitivity score & P19 motivated reasoning recognition record & meta & \hyperref[def:p19-mrg-four-outcomes]{\ref*{def:p19-mrg-four-outcomes}} & \thepage & Reversed printed subtraction is not implemented or silently corrected. \\
\phantomsection\label{entry:symbol-registry-05570}\(F\) & \texttt{\detokenize{F}} & reported regression or analysis-of-variance statistic & P19 motivated reasoning recognition and closure records & meta & \hyperref[sec:p19-mrg-strategies]{\ref*{sec:p19-mrg-strategies}} & \thepage & Source model comparison with displayed degrees of freedom. \\
\phantomsection\label{entry:symbol-registry-05571}\(\mathrm{RT}\) & \texttt{\detokenize{\mathrm{RT}}} & reaction time used as an effort proxy & P19 motivated reasoning dot-motion record & meta & \hyperref[sec:p19-mrg-capacity]{\ref*{sec:p19-mrg-capacity}} & \thepage & Latency alone does not identify resource capacity or accuracy. \\
\phantomsection\label{entry:symbol-registry-05572}\(\eta^2\) & \texttt{\detokenize{\eta^2}} & partial variance-ratio effect-size notation & P19 motivated reasoning dot-motion reporting boundary & meta & \hyperref[sec:p19-mrg-capacity]{\ref*{sec:p19-mrg-capacity}} & \thepage & Printed 1.141 is outside its allowed range and is not adopted as valid. \\
\phantomsection\label{entry:symbol-registry-05748}\(90^\circ\) & \texttt{\detokenize{\circ}} & reported angular turn or disorientation-check threshold & P19 spatial updating & object & \hyperref[sec:p19-spu-updating]{\ref*{sec:p19-spu-updating}} & \thepage & Heading angle in degrees; not the angle of a rotated comparison object. \\
\phantomsection\label{entry:symbol-registry-05749}\(N,n\) & \texttt{\detokenize{N;n}} & recruited or analysed participant denominator & P19 spatial updating evidence & meta & \hyperref[sec:p19-spu-updating]{\ref*{sec:p19-spu-updating}} & \thepage & Group participant counts remain distinct from repeated trial counts. \\
\phantomsection\label{entry:symbol-registry-05750}\(p\) & \texttt{\detokenize{p}} & reported significance-test probability & P19 spatial updating evidence & meta & \hyperref[sec:p19-spu-updating]{\ref*{sec:p19-spu-updating}} & \thepage & Non-significance is not an equivalence test. \\
\phantomsection\label{entry:symbol-registry-05751}\(F\) & \texttt{\detokenize{F}} & reported F statistic with source degrees of freedom & P19 spatial updating evidence & meta & \hyperref[sec:p19-spu-updating]{\ref*{sec:p19-spu-updating}} & \thepage & No new fit or reanalysis is supplied. \\
\phantomsection\label{entry:symbol-registry-05752}\(\eta^2\) & \texttt{\detokenize{\eta}} & reported partial effect-size statistic & P19 Liu-Xiao evidence & meta & \hyperref[sec:p19-spu-updating]{\ref*{sec:p19-spu-updating}} & \thepage & The source reports partial eta squared for the selected latency contrast. \\
\phantomsection\label{entry:symbol-registry-05753}95\% CI & \texttt{\detokenize{none}} & reported confidence interval for signed error & P19 Liu-Xiao evidence & meta & \hyperref[tab:p19-spu-bias]{\ref*{tab:p19-spu-bias}} & \thepage & Intervals in degrees; distinct from the source bar charts described as plus/minus one standard error. \\
\phantomsection\label{entry:symbol-registry-05754}SEM & \texttt{\detokenize{none}} & reported standard error of a mean choice proportion & P19 virtual geometry evidence & meta & \hyperref[sec:p19-spu-geometry]{\ref*{sec:p19-spu-geometry}} & \thepage & Not a 95 percent confidence interval or variability among all individual navigators. \\
\phantomsection\label{entry:symbol-registry-05755}AIC & \texttt{\detokenize{none}} & reported model-selection criterion & P19 virtual geometry evidence & meta & \hyperref[sec:p19-spu-geometry]{\ref*{sec:p19-spu-geometry}} & \thepage & Printed numerical ordering conflicts with source prose; no repaired ranking claimed. \\
\phantomsection\label{entry:symbol-registry-05756}A/B/C/D/U & \texttt{\detokenize{none}} & corner labels in the L-shaped virtual room & P19 Ambosta transfer experiment & object & \hyperref[sec:p19-spu-geometry]{\ref*{sec:p19-spu-geometry}} & \thepage & Source-local labels; distinct from child-search outcome codes. \\
\phantomsection\label{entry:symbol-registry-05757}C/R/N/F & \texttt{\detokenize{none}} & correct rotationally equivalent near and far first-search codes & P19 Shusterman child-search evidence & object & \hyperref[def:p19-spu-search-codes]{\ref*{def:p19-spu-search-codes}} & \thepage & Defined relative to the hidden target; F here is not the statistical F symbol. \\
\phantomsection\label{entry:symbol-registry-05758}C+R & \texttt{\detokenize{none}} & geometry-consistent first-search proportion & P19 Shusterman child-search evidence & meta & \hyperref[def:p19-spu-search-codes]{\ref*{def:p19-spu-search-codes}} & \thepage & Can be high without identifying which opposite corner is correct. \\
\phantomsection\label{entry:symbol-registry-05759}C--R & \texttt{\detokenize{none}} & difference of correct and rotationally equivalent search proportions & P19 Shusterman child-search evidence & meta & \hyperref[def:p19-spu-search-codes]{\ref*{def:p19-spu-search-codes}} & \thepage & Not a conditional proportion given geometry-consistent responding. \\
\phantomsection\label{entry:symbol-registry-05760}\(t\) & \texttt{\detokenize{t}} & reported t statistic with source degrees of freedom & P19 Shusterman child-search evidence & meta & \hyperref[sec:p19-spu-landmarks]{\ref*{sec:p19-spu-landmarks}} & \thepage & Indirect-landmark subgroup contrast is explicitly one-tailed. \\
\phantomsection\label{entry:symbol-registry-06210}N & \texttt{\detokenize{N}} & Reported participant analysis/sample count & p19 insight empirical records & n/a & \hyperref[sec:p19-insight-matchstick-evidence]{\ref*{sec:p19-insight-matchstick-evidence}} & \thepage & Not number of trials or a natural-number domain Original precedence descriptor 'atomic' preserved; n/a is not an invented rank. \\
\phantomsection\label{entry:symbol-registry-06211}n & \texttt{\detokenize{n}} & Reported subgroup participant count & p19 insight empirical records & n/a & \hyperref[sec:p19-insight-matchstick-evidence]{\ref*{sec:p19-insight-matchstick-evidence}} & \thepage & Separate cohorts and imaging tails retained Original precedence descriptor 'atomic' preserved; n/a is not an invented rank. \\
\phantomsection\label{entry:symbol-registry-06212}F & \texttt{\detokenize{F}} & Reported ANOVA statistic with source degrees of freedom & p19 insight empirical records & n/a & \hyperref[sec:p19-insight-matchstick-evidence]{\ref*{sec:p19-insight-matchstick-evidence}} & \thepage & Not ranked alphabet or final-state set Original precedence descriptor 'statistic' preserved; n/a is not an invented rank. \\
\phantomsection\label{entry:symbol-registry-06213}H & \texttt{\detokenize{H}} & Reported Kruskal-Wallis statistic & p19 insight empirical records & n/a & \hyperref[sec:p19-insight-matchstick-evidence]{\ref*{sec:p19-insight-matchstick-evidence}} & \thepage & No participant-level recomputation Original precedence descriptor 'statistic' preserved; n/a is not an invented rank. \\
\phantomsection\label{entry:symbol-registry-06214}\ensuremath{\chi^2} & \texttt{\detokenize{\chi^2}} & Reported chi-square statistic & p19 insight empirical records & n/a & \hyperref[sec:p19-insight-matchstick-evidence]{\ref*{sec:p19-insight-matchstick-evidence}} & \thepage & Source comparison and sample bound retained Original precedence descriptor 'statistic' preserved; n/a is not an invented rank. \\
\phantomsection\label{entry:symbol-registry-06215}t & \texttt{\detokenize{t}} & Reported t statistic & p19 insight empirical records & n/a & \hyperref[sec:p19-insight-neural-evidence]{\ref*{sec:p19-insight-neural-evidence}} & \thepage & Not time variable Original precedence descriptor 'statistic' preserved; n/a is not an invented rank. \\
\phantomsection\label{entry:symbol-registry-06216}r & \texttt{\detokenize{r}} & Reported association coefficient & p19 insight empirical records & n/a & \hyperref[sec:p19-insight-affect-evidence]{\ref*{sec:p19-insight-affect-evidence}} & \thepage & No causal mediation interpretation Original precedence descriptor 'statistic' preserved; n/a is not an invented rank. \\
\phantomsection\label{entry:symbol-registry-06217}p & \texttt{\detokenize{p}} & Reported significance tail probability & p19 insight empirical records & n/a & \hyperref[sec:p19-insight-matchstick-evidence]{\ref*{sec:p19-insight-matchstick-evidence}} & \thepage & Not posterior hypothesis probability Original precedence descriptor 'statistic' preserved; n/a is not an invented rank. \\
\phantomsection\label{entry:symbol-registry-06218}PA & \texttt{\detokenize{PA}} & PANAS positive affect score & p19 insight affect record & n/a & \hyperref[sec:p19-insight-affect-evidence]{\ref*{sec:p19-insight-affect-evidence}} & \thepage & Assessed state, not assigned intervention Original precedence descriptor 'atomic' preserved; n/a is not an invented rank. \\
\phantomsection\label{entry:symbol-registry-06219}NA & \texttt{\detokenize{NA}} & PANAS negative affect score & p19 insight affect record & n/a & \hyperref[sec:p19-insight-affect-evidence]{\ref*{sec:p19-insight-affect-evidence}} & \thepage & Not assumed opposite pole of PA Original precedence descriptor 'atomic' preserved; n/a is not an invented rank. \\
\phantomsection\label{entry:symbol-registry-06220}x y & \texttt{\detokenize{x y}} & Cartesian coordinates in the geometric illustration & p19 insight geometric annex & n/a & \hyperref[annex:p19.insight-nine-dot]{\ref*{annex:p19.insight-nine-dot}} & \thepage & Only objective path check, not participant representation data Original precedence descriptor 'atomic' preserved; n/a is not an invented rank. \\
\phantomsection\label{entry:symbol-registry-06221}X Y Z & \texttt{\detokenize{X Y Z}} & Schematic arithmetic-statement placeholders & p19 insight matchstick mechanisms & n/a & \hyperref[sec:p19-insight-representation-change]{\ref*{sec:p19-insight-representation-change}} & \thepage & Not random variables; identity-chain validity requires actual matching values Original precedence descriptor 'atomic' preserved; n/a is not an invented rank. \\
\phantomsection\label{entry:symbol-registry-06222}\ensuremath{\mathrm{I} \mathrm{V} \mathrm{X}} & \texttt{\detokenize{\mathrm{I} \mathrm{V} \mathrm{X}}} & Roman numeral components in the matchstick task & p19 insight matchstick mechanisms & n/a & \hyperref[sec:p19-insight-representation-change]{\ref*{sec:p19-insight-representation-change}} & \thepage & Physical stroke/move constraints retain task-specific roles Original precedence descriptor 'atomic' preserved; n/a is not an invented rank. \\
\phantomsection\label{entry:symbol-registry-08035}\(N\) & \texttt{\detokenize{N}} & participant count in specified empirical population & P19 language categories empirical records & meta & \hyperref[def:p19-lc-evidence-record]{\ref*{def:p19-lc-evidence-record}} & \thepage & Not trial or item count; exclusions and reused comparison populations retained. \\
\phantomsection\label{entry:symbol-registry-08036}\(F\) & \texttt{\detokenize{F}} & reported analysis of variance statistic & P19 language categories empirical records & meta & \hyperref[def:p19-lc-evidence-record]{\ref*{def:p19-lc-evidence-record}} & \thepage & Source degrees of freedom recorded with each selected contrast. \\
\phantomsection\label{entry:symbol-registry-08037}\(t\) & \texttt{\detokenize{t}} & reported Student test statistic & P19 language categories empirical records & meta & \hyperref[def:p19-lc-evidence-record]{\ref*{def:p19-lc-evidence-record}} & \thepage & Not a time coordinate or latent cognitive variable. \\
\phantomsection\label{entry:symbol-registry-08038}\(p\) & \texttt{\detokenize{p}} & reported significance probability & P19 language categories empirical records & meta & \hyperref[def:p19-lc-evidence-record]{\ref*{def:p19-lc-evidence-record}} & \thepage & Non-significance is not a proof of equivalence. \\
\phantomsection\label{entry:symbol-registry-08039}\(r\) & \texttt{\detokenize{r}} & reported Pearson correlation between average item judgments & P19 language categories empirical records & meta & \hyperref[def:p19-lc-evidence-record]{\ref*{def:p19-lc-evidence-record}} & \thepage & Item correlations are not randomized language effects. \\
\phantomsection\label{entry:symbol-registry-08040}\(SEM\) & \texttt{\detokenize{none}} & reported standard error of condition mean & P19 Russian-blue latency table & meta & \hyperref[def:p19-lc-evidence-record]{\ref*{def:p19-lc-evidence-record}} & \thepage & Not standard deviation of participants or confidence interval. \\
\phantomsection\label{entry:symbol-registry-08041}\(x\) & \texttt{\detokenize{x}} & presented target in numerical units & P19 bounded number-line tasks & object & \hyperref[def:p19-nl-measures]{\ref*{def:p19-nl-measures}} & \thepage & Converted response is y; target values are not physical lengths. \\
\phantomsection\label{entry:symbol-registry-08042}\(y\) & \texttt{\detokenize{y}} & spatial mark converted to numerical estimate or response-model output in number units & P19 number-line response functions & object & \hyperref[def:p19-nl-measures]{\ref*{def:p19-nl-measures}} & \thepage & Type stated before use; proportion mappings themselves have unitless output. \\
\phantomsection\label{entry:symbol-registry-08043}\(U\) & \texttt{\detokenize{U}} & positive numerical upper endpoint on a zero-based line & P19 bounded number-line tasks & object & \hyperref[def:p19-nl-measures]{\ref*{def:p19-nl-measures}} & \thepage & Also the range denominator for the stated PAE; not an uncertainty bound. \\
\phantomsection\label{entry:symbol-registry-08044}\(\mathrm{PAE}\) & \texttt{\detokenize{\mathrm}} & reported or defined percentage absolute error & P19 number-line and quantitative tasks & meta & \hyperref[def:p19-nl-measures]{\ref*{def:p19-nl-measures}} & \thepage & Range-specific denominator; exact correct answers remain a separate measure. \\
\phantomsection\label{entry:symbol-registry-08045}\(a\) & \texttt{\detokenize{a}} & fitted affine response intercept & P19 number-line fit definition & object & \hyperref[def:p19-nl-measures]{\ref*{def:p19-nl-measures}} & \thepage & A fitted intercept is not assumed zero. \\
\phantomsection\label{entry:symbol-registry-08046}\(b\) & \texttt{\detokenize{b}} & fitted affine response slope & P19 number-line fit definition & object & \hyperref[def:p19-nl-measures]{\ref*{def:p19-nl-measures}} & \thepage & Distinct from the proportion exponent beta and not assumed one. \\
\phantomsection\label{entry:symbol-registry-08047}\(R^2\) & \texttt{\detokenize{R}} & reported explained-variance fit statistic & P19 empirical response fits & meta & \hyperref[def:p19-nl-measures]{\ref*{def:p19-nl-measures}} & \thepage & Group median fits differ from mean individual fits; not percent trials exactly correct. \\
\phantomsection\label{entry:symbol-registry-08048}\(R^2_{\rm lin}\) & \texttt{\detokenize{R}} & linear response-fit statistic in the Thompson-Siegler and related reports & P19 memory and arithmetic evidence & meta & \hyperref[sec:p19-nl-memory]{\ref*{sec:p19-nl-memory}} & \thepage & Retains source-specific linear-fit role rather than equating it with recall accuracy. \\
\phantomsection\label{entry:symbol-registry-08049}\(N\) & \texttt{\detokenize{N}} & total or condition participant denominator & P19 original study populations & meta & \hyperref[sec:p19-nl-siegler-opfer]{\ref*{sec:p19-nl-siegler-opfer}} & \thepage & Participant counts are distinct from targets and repeated trials. \\
\phantomsection\label{entry:symbol-registry-08050}\(NP\) & \texttt{\detokenize{none}} & number-to-position task code & P19 Siegler-Opfer original study & meta & \hyperref[sec:p19-nl-siegler-opfer]{\ref*{sec:p19-nl-siegler-opfer}} & \thepage & Source task label; written target maps to marked spatial location. \\
\phantomsection\label{entry:symbol-registry-08051}\(PN\) & \texttt{\detokenize{none}} & position-to-number task code & P19 Siegler-Opfer original study & meta & \hyperref[sec:p19-nl-siegler-opfer]{\ref*{sec:p19-nl-siegler-opfer}} & \thepage & Source reverse task label; marked spatial location maps to a stated numeral. \\
\phantomsection\label{entry:symbol-registry-08052}\(t\) & \texttt{\detokenize{t}} & reported Student test statistic with degrees of freedom & P19 empirical contrasts & meta & \hyperref[sec:p19-nl-siegler-opfer]{\ref*{sec:p19-nl-siegler-opfer}} & \thepage & Genuine glyph collision with the unitless proportion t; typed source/test notation retained. \\
\phantomsection\label{entry:symbol-registry-08053}\(p\) & \texttt{\detokenize{p}} & reported significance-test probability & P19 empirical contrasts & meta & \hyperref[sec:p19-nl-siegler-opfer]{\ref*{sec:p19-nl-siegler-opfer}} & \thepage & No inferred equivalence from non-significance; no invented multiplicity correction. \\
\phantomsection\label{entry:symbol-registry-08054}\(SD\) & \texttt{\detokenize{none}} & reported participant age or outcome standard deviation & P19 original study descriptions & meta & \hyperref[sec:p19-nl-siegler-opfer]{\ref*{sec:p19-nl-siegler-opfer}} & \thepage & Not a confidence interval or standard error. \\
\phantomsection\label{entry:symbol-registry-08055}\(\ln\) & \texttt{\detokenize{\ln}} & natural logarithm used in positive-domain candidate functions & P19 magnitude models & meta & \hyperref[thm:p19-nl-ratio-identification]{\ref*{thm:p19-nl-ratio-identification}} & \thepage & No logarithm of zero; offset models remain distinct. \\
\phantomsection\label{entry:symbol-registry-08056}\(t\) & \texttt{\detokenize{t}} & true unitless proportion in the interval from zero to one & P19 proportional response model & object & \hyperref[def:p19-nl-proportion-model]{\ref*{def:p19-nl-proportion-model}} & \thepage & Genuine glyph collision with reported test statistic t; type declared locally. \\
\phantomsection\label{entry:symbol-registry-08057}\(\beta\) & \texttt{\detokenize{\beta}} & positive power exponent controlling proportional response bias & P19 proportional response model & object & \hyperref[def:p19-nl-proportion-model]{\ref*{def:p19-nl-proportion-model}} & \thepage & Fitted parameter is not directly identified as a neural code; allows values above one. \\
\phantomsection\label{entry:symbol-registry-08058}\(P_\beta\) & \texttt{\detokenize{P_\beta}} & one-cycle unitless proportion mapping & P19 proportional response model & meta & \hyperref[def:p19-nl-proportion-model]{\ref*{def:p19-nl-proportion-model}} & \thepage & Explicit quotient formula including endpoints; source-inspired response construction. \\
\phantomsection\label{entry:symbol-registry-08059}\(Q_\beta\) & \texttt{\detokenize{Q_\beta}} & two-cycle unitless mapping using a midpoint & P19 proportional response model & meta & \hyperref[def:p19-nl-proportion-model]{\ref*{def:p19-nl-proportion-model}} & \thepage & Explicit piecewise construction; both pieces agree at one half. \\
\phantomsection\label{entry:symbol-registry-08060}\(\widehat y\) & \texttt{\detokenize{\widehat}} & predicted numerical estimate after scaling a proportion by U & P19 proportional response model & object & \hyperref[def:p19-nl-proportion-model]{\ref*{def:p19-nl-proportion-model}} & \thepage & Model prediction is distinct from an observed child response. \\
\phantomsection\label{entry:symbol-registry-08061}\(\alpha\) & \texttt{\detokenize{\alpha}} & power-function scaling factor & P19 unbounded response candidate & object & \hyperref[def:p19-nl-proportion-model]{\ref*{def:p19-nl-proportion-model}} & \thepage & Free versus fixed scaling assumptions must be specified. \\
\phantomsection\label{entry:symbol-registry-08062}\(\mathrm{AIC_c}\) & \texttt{\detokenize{\mathrm}} & reported small-sample-corrected Akaike model-selection criterion & P19 Slusser-Santiago-Barth comparison & meta & \hyperref[sec:p19-nl-proportion]{\ref*{sec:p19-nl-proportion}} & \thepage & Selected statistics are from Appendix C of the online-first presentation. \\
\phantomsection\label{entry:symbol-registry-08063}\(r\) & \texttt{\detokenize{r}} & reported Pearson correlation statistic & P19 memory evidence & meta & \hyperref[sec:p19-nl-memory]{\ref*{sec:p19-nl-memory}} & \thepage & Not a randomized treatment effect; genuine collision with odds ratio in the annex. \\
\phantomsection\label{entry:symbol-registry-08064}\(F\) & \texttt{\detokenize{F}} & reported F test statistic with degrees of freedom & P19 memory arithmetic and intervention evidence & meta & \hyperref[sec:p19-nl-memory]{\ref*{sec:p19-nl-memory}} & \thepage & Source degrees of freedom and outcome-specific denominator retained. \\
\phantomsection\label{entry:symbol-registry-08065}\(d\) & \texttt{\detokenize{d}} & reported standardized contrast effect-size statistic & P19 memory and intervention evidence & meta & \hyperref[sec:p19-nl-memory]{\ref*{sec:p19-nl-memory}} & \thepage & No independent re-estimation or confidence interval supplied. \\
\phantomsection\label{entry:symbol-registry-08066}\(\eta^2\) & \texttt{\detokenize{\eta}} & reported eta-squared effect-size statistic & P19 Booth-Siegler arithmetic analyses & meta & \hyperref[sec:p19-nl-arithmetic]{\ref*{sec:p19-nl-arithmetic}} & \thepage & Do not silently rename as partial eta squared when source prints eta squared. \\
\phantomsection\label{entry:symbol-registry-08067}\(17+29\) & \texttt{\detokenize{+}} & stipulated exact arithmetic addition used by the trained task & P19 arithmetic benchmark & infix & \hyperref[sec:p19-nl-arithmetic]{\ref*{sec:p19-nl-arithmetic}} & \thepage & Correct answer 46 is the task benchmark rather than inferred participant competence. \\
\phantomsection\label{entry:symbol-registry-08068}\(17\cdot29\) & \texttt{\detokenize{\cdot}} & multiplication induced by adding logarithms then exponentiating & P19 arithmetic code countercase & infix & \hyperref[sec:p19-nl-arithmetic]{\ref*{sec:p19-nl-arithmetic}} & \thepage & Demonstrates that a representation does not itself specify an addition procedure. \\
\phantomsection\label{entry:symbol-registry-08069}\(r\) & \texttt{\detokenize{r}} & positive odds ratio t divided by one minus t in the mathematical proof & P19 proportion-bias annex & meta & \hyperref[annex:p19-nl-proportion-bias]{\ref*{annex:p19-nl-proportion-bias}} & \thepage & Genuine collision with empirical Pearson correlation; local auxiliary defined before use. \\
\phantomsection\label{entry:symbol-registry-08070}\(s\) & \texttt{\detokenize{s}} & local upper-half proportion two t minus one & P19 proportion-bias annex & meta & \hyperref[annex:p19-nl-proportion-bias]{\ref*{annex:p19-nl-proportion-bias}} & \thepage & Auxiliary used only to prove the second cycle's bias. \\
\phantomsection\label{entry:symbol-registry-08071}\(z\) & \texttt{\detokenize{z}} & positive argument of the increasing odds-to-proportion transform & P19 proportion-bias annex & meta & \hyperref[annex:p19-nl-proportion-bias]{\ref*{annex:p19-nl-proportion-bias}} & \thepage & Auxiliary mathematical variable not a latent psychological score. \\
\phantomsection\label{entry:symbol-registry-08072}\(\kappa\) & \texttt{\detokenize{\kappa}} & positive logarithmic scale constant & P19 ratio-identification construction & object & \hyperref[thm:p19-nl-ratio-identification]{\ref*{thm:p19-nl-ratio-identification}} & \thepage & Stipulated model parameter rather than an empirical estimate. \\
\phantomsection\label{entry:symbol-registry-08073}\(\varepsilon\) & \texttt{\detokenize{\varepsilon}} & mean-zero real random comparison noise & P19 ratio-identification construction & object & \hyperref[thm:p19-nl-ratio-identification]{\ref*{thm:p19-nl-ratio-identification}} & \thepage & Indexed copies are independent with the same law; moments explicitly assumed. \\
\phantomsection\label{entry:symbol-registry-08074}\(\mathbb E\) & \texttt{\detokenize{\mathbb}} & expectation operator on stipulated random variables & P19 ratio-identification construction & meta & \hyperref[thm:p19-nl-ratio-identification]{\ref*{thm:p19-nl-ratio-identification}} & \thepage & Mathematical expectation rather than an experimental mean computed from participants. \\
\phantomsection\label{entry:symbol-registry-08075}\(L_x\) & \texttt{\detokenize{L}} & logarithmic latent comparison code & P19 ratio-identification construction & object & \hyperref[thm:p19-nl-ratio-identification]{\ref*{thm:p19-nl-ratio-identification}} & \thepage & Not an observed line placement. \\
\phantomsection\label{entry:symbol-registry-08076}\(m\) & \texttt{\detokenize{m}} & positive finite exponential-moment normalizer & P19 ratio-identification construction & meta & \hyperref[thm:p19-nl-ratio-identification]{\ref*{thm:p19-nl-ratio-identification}} & \thepage & Existence and positivity established in the reciprocal annex. \\
\phantomsection\label{entry:symbol-registry-08077}\(Z_x\) & \texttt{\detokenize{Z}} & positive mean-one multiplicative comparison factor & P19 ratio-identification construction & object & \hyperref[thm:p19-nl-ratio-identification]{\ref*{thm:p19-nl-ratio-identification}} & \thepage & Law derived from the same real noise used to couple comparison outcomes. \\
\phantomsection\label{entry:symbol-registry-08078}\(S_x\) & \texttt{\detokenize{S}} & linear-mean scalar-variability comparison code & P19 ratio-identification construction & object & \hyperref[thm:p19-nl-ratio-identification]{\ref*{thm:p19-nl-ratio-identification}} & \thepage & Distinct mean/noise decomposition with identical stipulated pair comparisons. \\
\phantomsection\label{entry:symbol-registry-08079}\(\operatorname{SD}\) & \texttt{\detokenize{\operatorname}} & standard-deviation operator for model random variables & P19 ratio-identification construction & meta & \hyperref[thm:p19-nl-ratio-identification]{\ref*{thm:p19-nl-ratio-identification}} & \thepage & Distinct from reported empirical SDs despite the shared abbreviation. \\
\phantomsection\label{entry:symbol-registry-08080}\(\operatorname{Var}\) & \texttt{\detokenize{\operatorname}} & variance operator for model random variables & P19 ratio-identification annex & meta & \hyperref[annex:p19-nl-ratio-identification]{\ref*{annex:p19-nl-ratio-identification}} & \thepage & Finite by the stated assumptions and derived exponential-moment bounds. \\
\phantomsection\label{entry:symbol-registry-08081}\(j\) & \texttt{\detokenize{j}} & positive integer indexing expectation lower-bound events & P19 ratio-identification annex & meta & \hyperref[annex:p19-nl-ratio-identification]{\ref*{annex:p19-nl-ratio-identification}} & \thepage & Auxiliary proof index; not a trial or participant index. \\
\phantomsection\label{entry:symbol-registry-08082}\(x\) & \texttt{\detokenize{x}} & first positive compared quantity in coupled latent-code models & P19 ratio-identification construction & object & \hyperref[thm:p19-nl-ratio-identification]{\ref*{thm:p19-nl-ratio-identification}} & \thepage & Source-model input role declared in the theorem; distinct from a target-estimate data pair. \\
\phantomsection\label{entry:symbol-registry-08083}\(y\) & \texttt{\detokenize{y}} & second positive compared quantity in coupled latent-code models & P19 ratio-identification construction & object & \hyperref[thm:p19-nl-ratio-identification]{\ref*{thm:p19-nl-ratio-identification}} & \thepage & Genuine collision with the later estimate-output y; both local types explicitly declared. \\
\phantomsection\label{entry:symbol-registry-08084}\(|y-x|\) & \texttt{\detokenize{|}} & absolute numerical difference for the source PAE measure & P19 number-line quantitative-error measure & object & \hyperref[def:p19-nl-measures]{\ref*{def:p19-nl-measures}} & \thepage & Absolute-value role retained honestly; not set cardinality or a count of exact responses. \\
\phantomsection\label{entry:symbol-registry-08085}\(x\) & \texttt{\detokenize{x}} & real input in algebraic identity and finite-data function constructions & P19 CTA identity and induction examples & object & \hyperref[def:p19-cta-like-terms]{\ref*{def:p19-cta-like-terms}} & \thepage & Defined locally; not an assertion of psychological representation. \\
\phantomsection\label{entry:symbol-registry-08086}\(y\) & \texttt{\detokenize{y}} & output in the finite function-induction data & P19 CTA function-induction example & object & \hyperref[def:p19-cta-function-induction]{\ref*{def:p19-cta-function-induction}} & \thepage & Four observations do not uniquely determine the intended function. \\
\phantomsection\label{entry:symbol-registry-08087}\(\lambda\) & \texttt{\detokenize{\lambda}} & real parameter in the interpolation counterexample family & P19 CTA function-induction example & object & \hyperref[def:p19-cta-function-induction]{\ref*{def:p19-cta-function-induction}} & \thepage & Constructed parameter not a fitted learner effect. \\
\phantomsection\label{entry:symbol-registry-08088}\(g_\lambda\) & \texttt{\detokenize{g}} & real polynomial family matching the four stipulated input-output pairs & P19 CTA function-induction example & object & \hyperref[def:p19-cta-function-induction]{\ref*{def:p19-cta-function-induction}} & \thepage & Construction gives distinct later values without claiming source-native experimental notation. \\
\phantomsection\label{entry:symbol-registry-08089}\(x\) & \texttt{\detokenize{x}} & hourly wage in the worked single-reference source equation & P19 CTA single-reference wage example & object & \hyperref[def:p19-cta-unwind]{\ref*{def:p19-cta-unwind}} & \thepage & Real monetary rate; distinct local role from identity input x. \\
\phantomsection\label{entry:symbol-registry-08090}\(X\) & \texttt{\detokenize{X}} & unknown monetary amount in the lottery and discount source cases & P19 CTA repeated-reference example & object & \hyperref[def:p19-cta-double-reference]{\ref*{def:p19-cta-double-reference}} & \thepage & Prize amount and original price are separately stipulated uses; rounding is not implicit. \\
\phantomsection\label{entry:symbol-registry-08091}\(b\) & \texttt{\detokenize{b}} & nonnegative whole boys count in the composition example & P19 CTA pupil substitution & object & \hyperref[def:p19-cta-students-substitution]{\ref*{def:p19-cta-students-substitution}} & \thepage & Defined independently of shortfall f. \\
\phantomsection\label{entry:symbol-registry-08092}\(f\) & \texttt{\detokenize{f}} & nonnegative whole shortfall of boys relative to 62 girls & P19 CTA pupil tasks & object & \hyperref[def:p19-cta-students-substitution]{\ref*{def:p19-cta-students-substitution}} & \thepage & Zero through 62; not the boys count itself. \\
\phantomsection\label{entry:symbol-registry-08093}\(b\) & \texttt{\detokenize{b}} & parental saved monetary amount in equal-sharing task & P19 CTA worked pooling example & object & \hyperref[sec:p19-cta-worked]{\ref*{sec:p19-cta-worked}} & \thepage & Real nonnegative amount; source-native glyph shared with the separately typed pupil count. \\
\phantomsection\label{entry:symbol-registry-08094}\(d\) & \texttt{\detokenize{d}} & amount or original price in grouped sharing and reduction examples & P19 CTA grouping and jacket examples & object & \hyperref[def:p19-cta-grouped-substitution]{\ref*{def:p19-cta-grouped-substitution}} & \thepage & Monetary scalar; not Cohen effect size d. \\
\phantomsection\label{entry:symbol-registry-08095}\(m\) & \texttt{\detokenize{m}} & amount spent before sharing the remainder & P19 CTA grouped spending example & object & \hyperref[def:p19-cta-grouped-substitution]{\ref*{def:p19-cta-grouped-substitution}} & \thepage & Source-native glyph; distinct from duration m in dock task. \\
\phantomsection\label{entry:symbol-registry-08096}\(m\) & \texttt{\detokenize{m}} & elapsed minutes during approach to a dock & P19 CTA worked dock example & object & \hyperref[sec:p19-cta-worked]{\ref*{sec:p19-cta-worked}} & \thepage & Real duration0through20 for the stated nonnegative remaining distance. \\
\phantomsection\label{entry:symbol-registry-08097}\(h\) & \texttt{\detokenize{h}} & duration measured in hours in earnings and travel source cases & P19 CTA grouped time and worked duration examples & object & \hyperref[def:p19-cta-grouped-substitution]{\ref*{def:p19-cta-grouped-substitution}} & \thepage & Local context states admissibility; travel h exceeds2. \\
\phantomsection\label{entry:symbol-registry-08098}\(t\) & \texttt{\detokenize{t}} & new travel duration in the grouped substitution illustration & P19 CTA grouped travel example & object & \hyperref[def:p19-cta-grouped-substitution]{\ref*{def:p19-cta-grouped-substitution}} & \thepage & Hours; not a reported t statistic. \\
\phantomsection\label{entry:symbol-registry-08099}\(c\) & \texttt{\detokenize{c}} & nonnegative whole number of CDs per box & P19 CTA worked deliveries example & object & \hyperref[sec:p19-cta-worked]{\ref*{sec:p19-cta-worked}} & \thepage & Not the monospace instructional-sequence code c. \\
\phantomsection\label{entry:symbol-registry-08100}\(N\) & \texttt{\detokenize{N}} & participant denominator for empirical study records & P19 CTA empirical measurements & meta & \hyperref[sec:p19-cta-representations]{\ref*{sec:p19-cta-representations}} & \thepage & Recruited and complete-case participants remain distinct. \\
\phantomsection\label{entry:symbol-registry-08101}\(F\) & \texttt{\detokenize{F}} & reported analysis-of-variance statistic & P19 CTA empirical measurements & meta & \hyperref[sec:p19-cta-representations]{\ref*{sec:p19-cta-representations}} & \thepage & Degrees of freedom retained; not algebraic functionality. \\
\phantomsection\label{entry:symbol-registry-08102}\(p\) & \texttt{\detokenize{p}} & reported significance probability & P19 CTA empirical measurements & meta & \hyperref[sec:p19-cta-representations]{\ref*{sec:p19-cta-representations}} & \thepage & Non-significance is not exact equivalence. \\
\phantomsection\label{entry:symbol-registry-08103}\(d\) & \texttt{\detokenize{d}} & reported standardized Cohen effect-size contrast & P19 CTA empirical measurements & meta & \hyperref[sec:p19-cta-representations]{\ref*{sec:p19-cta-representations}} & \thepage & Genuine typed role distinct from monetary amount d. \\
\phantomsection\label{entry:symbol-registry-08104}\(SD\) & \texttt{\detokenize{none}} & reported standard deviation of the outcome score & P19 CTA empirical measurements & meta & \hyperref[sec:p19-cta-representations]{\ref*{sec:p19-cta-representations}} & \thepage & Not a confidence interval or condition-mean standard error. \\
\phantomsection\label{entry:symbol-registry-08105}\(*\) & \texttt{\detokenize{none}} & source-spelled arithmetic multiplication operator & P19 CTA worked symbolization & infix & \hyperref[sec:p19-cta-worked]{\ref*{sec:p19-cta-worked}} & \thepage & Rational constants count as quantities in the source's operator-count analysis. \\
\phantomsection\label{entry:symbol-registry-08106}\(\texttt{Exp}\) & \texttt{\detokenize{\texttt}} & expression nonterminal label in the corrected single-category galley schematic & P19 CTA source grammar & meta & \hyperref[sec:p19-cta-production]{\ref*{sec:p19-cta-production}} & \thepage & Parenthesized expression uses the same category; source's isolated Expr spelling corrected with provenance. \\
\phantomsection\label{entry:symbol-registry-08107}\(\texttt{Op}\) & \texttt{\detokenize{\texttt}} & operator-token category in the galley expression schematic & P19 CTA source grammar & meta & \hyperref[sec:p19-cta-production]{\ref*{sec:p19-cta-production}} & \thepage & Four arithmetic tokens; schematic does not specify precedence or unique parsing. \\
\phantomsection\label{entry:symbol-registry-08108}\(\texttt{Num}\) & \texttt{\detokenize{\texttt}} & number-token category in the galley schematic & P19 CTA source grammar & meta & \hyperref[sec:p19-cta-production]{\ref*{sec:p19-cta-production}} & \thepage & Not the actual participant count N. \\
\phantomsection\label{entry:symbol-registry-08109}\(\texttt{Var}\) & \texttt{\detokenize{\texttt}} & variable-token category in the galley schematic & P19 CTA source grammar & meta & \hyperref[sec:p19-cta-production]{\ref*{sec:p19-cta-production}} & \thepage & Token category distinct from a variable's numerical value. \\
\phantomsection\label{entry:symbol-registry-08110}\(\texttt{RecExprProd}\) & \texttt{\detokenize{\texttt}} & proposed recursive-expression-production component label & P19 CTA component model & meta & \hyperref[sec:p19-cta-production]{\ref*{sec:p19-cta-production}} & \thepage & Source label not a proved indivisible mental module. \\
\phantomsection\label{entry:symbol-registry-08111}\(\texttt{EngToMathOps}\) & \texttt{\detokenize{\texttt}} & proposed English-to-mathematical-operation component label & P19 CTA component model & meta & \hyperref[sec:p19-cta-production]{\ref*{sec:p19-cta-production}} & \thepage & Source label not a formal calculus. \\
\phantomsection\label{entry:symbol-registry-08112}\(\texttt{K1}\) & \texttt{\detokenize{\texttt}} & shared component label in the overlap model & P19 CTA overlap comparison & meta & \hyperref[sec:p19-cta-overlap]{\ref*{sec:p19-cta-overlap}} & \thepage & TaskA needsK1; taskB needsK1andK2 within this proposed model. \\
\phantomsection\label{entry:symbol-registry-08113}\(\texttt{K2}\) & \texttt{\detokenize{\texttt}} & additional component label for taskB in overlap model & P19 CTA overlap comparison & meta & \hyperref[sec:p19-cta-overlap]{\ref*{sec:p19-cta-overlap}} & \thepage & Not inferred uniquely from a difficulty difference alone. \\
\phantomsection\label{entry:symbol-registry-08114}\(\texttt{Ka}\) & \texttt{\detokenize{\texttt}} & separate component label for taskA in nonoverlap model & P19 CTA overlap comparison & meta & \hyperref[sec:p19-cta-overlap]{\ref*{sec:p19-cta-overlap}} & \thepage & Alternative model's task-specific label. \\
\phantomsection\label{entry:symbol-registry-08115}\(\texttt{Kb}\) & \texttt{\detokenize{\texttt}} & separate component label for taskB in nonoverlap model & P19 CTA overlap comparison & meta & \hyperref[sec:p19-cta-overlap]{\ref*{sec:p19-cta-overlap}} & \thepage & Alternative model not collapsed with K2. \\
\phantomsection\label{entry:symbol-registry-08116}\(\texttt{a}\) & \texttt{\detokenize{\texttt}} & first one-step item code in the 2010 sequence & P19 CTA instructional sequence & meta & \hyperref[sec:p19-cta-intervention]{\ref*{sec:p19-cta-intervention}} & \thepage & Code not an algebraic coefficient. \\
\phantomsection\label{entry:symbol-registry-08117}\(\texttt{b}\) & \texttt{\detokenize{\texttt}} & second one-step item code in the 2010 sequence & P19 CTA instructional sequence & meta & \hyperref[sec:p19-cta-intervention]{\ref*{sec:p19-cta-intervention}} & \thepage & Code distinct from mathematical pupil or saved-amount variable b. \\
\phantomsection\label{entry:symbol-registry-08118}\(\texttt{c}\) & \texttt{\detokenize{\texttt}} & substitution-item code in the 2010 sequence & P19 CTA instructional sequence & meta & \hyperref[sec:p19-cta-intervention]{\ref*{sec:p19-cta-intervention}} & \thepage & Code not CD count c. \\
\phantomsection\label{entry:symbol-registry-08119}\(\texttt{0--3}\) & \texttt{\detokenize{\texttt}} & range of four two-step item codes in the 2010 sequence & P19 CTA instructional sequence & meta & \hyperref[sec:p19-cta-intervention]{\ref*{sec:p19-cta-intervention}} & \thepage & Range notation denotes 0, 1, 2, 3; one intake record not four atomically counted glyphs. \\
\end{longtable}
\endgroup
\end{landscape}


\subsection{Classical-logic core}
\label{sec:p01-classical-symbols}

The classical-logic core consists of global rows of the register: the
signature \(\mathcal L\), the object-language variables, the term, formula
and sentence metavariables, the formula contexts and theories, the
structures, valuations and assignments, the constants \(\bot,\top\), the
five connectives, the two quantifiers, and the signs for satisfaction,
truth in a structure, semantic consequence, derivability and sequents.
Table~\ref{tab:p01-classical-symbols}, the global-scope table of the
generated register above, prints these rows together with the other global
rows of \texttt{registry/symbols.csv}.  No second, hand-maintained extract
is kept, so a change to a core row changes this lookup table when the
tables are regenerated.  Propositional satisfaction by a valuation is not a
global row: it is recorded with the scope \emph{P2 prop-syntax-semantics}
and printed in Table~\ref{tab:p01-w02-p01-026-symbols-p2}.  The global
metavariable rows take their types from the active language, calculus and
semantics once these are fixed.
Table~\ref{tab:p01-w02b-p01-012-parts12} adds every registry row of global scope,
with a scope opening in Part~1 or Part~2, or with a scope label that names
Part~2 but is declared in Part~3, and so every connective, quantifier,
turnstile and metavariable row of those parts, with its recorded scope.

\section{Binding and precedence}
\label{sec:p01-binding-precedence}
\index{precedence}\index{binding}

Table~\ref{tab:p01-precedence}, generated from \texttt{registry/symbols.csv},
lists every registry row that carries a numerical rank, ordered by rank,
with its role and its recorded scope.  A row whose precedence column
holds a descriptor such as \texttt{meta} or \texttt{n/a} instead of a rank
is not listed there; the master register prints its descriptor.
\Cref{def:p01-w02-p01-026-precedence} gives the fixity, the association and
the reading rule of each classical rank and of the first nonclassical
extensions; unrelated subsystem symbols share a rank only where their
scopes are disjoint.  \Cref{def:p01-w02b-p01-014-quantifier-scope} states
how these ranks license omitted parentheses and how far a quantifier
reaches.  Fixity and association for every other registered rank, and the
rule for different operators that share a rank, are given in
\cref{def:p01-w02-p01-024-fixity} and
Table~\ref{tab:p01-w02-p01-024-ranks}.

\begin{hbconvention}{def:p01-w02-p01-026-precedence}{Fixity and association of the classical ranks}
The convention applies only after the active language has been fixed and
its unique readability established; where the grammar of that language
fixes a reading, the reading stands.  Smaller numerical ranks bind more
tightly.
\begin{enumerate}
\item \emph{Rank 1.}  Function and predicate application, and the term
  operations of rank~1 in Table~\ref{tab:p01-precedence}.  The arguments
  belong to the displayed function or predicate symbol, as its arity fixes
  them.
\item \emph{Rank 2.}  Identity \(=\) and the other declared atomic
  relations, such as \(\leq\) and \(>\).  They are nonassociative, and a
  relation combines terms into an atomic formula.
\item \emph{Rank 2.5.}  In the linear-logic scope the exponentials \(!A\)
  and \(?A\) are prefix and the linear negation \(A^\perp\) is postfix; in
  the da Costa scope the consistency sign \(A^\circ\) is postfix.  A prefix
  sign governs the immediately following formula and a postfix sign the
  immediately preceding one; a prefix and a postfix sign on one operand
  require parentheses.
\item \emph{Rank 3.}  \(\neg,\Box,\Diamond,\mathsf X,\mathsf F,\mathsf G\),
  \(K_i,E_G,C_G\) and the quantifier prefixes \(\forall x,\exists x\) are
  prefix.  The operator or quantifier governs its operand only
  (\cref{def:p01-w02b-p01-014-quantifier-scope}); parentheses extend its
  scope.
\item \emph{Rank 3.5.}  Temporal until \(A\mathbin{\mathsf U}B\) is infix
  and nonassociative.  It binds below the unary temporal operators and
  above the binary propositional connectives.
\item \emph{Rank 4.}  \(\land\), and \(\otimes,\bindnasrepma,\mathbin{\&}\)
  in the linear-logic scope, are infix and left associative for repetitions
  of one operator; different operators of this rank need parentheses.
  Conjunction-level operators bind above disjunction.
\item \emph{Rank 5.}  \(\lor\), and \(\oplus\) in the linear-logic scope,
  are infix and left associative for repetitions of one operator.
  Disjunction-level operators bind above implication.
\item \emph{Rank 6.}  \(\to\), and \(\multimap\) in the linear-logic
  scope, are infix and right associative.
\item \emph{Rank 7.}  \(\leftrightarrow\) is infix and nonassociative, so
  repeated biconditionals require parentheses.
\item \emph{Metalinguistic signs.}  \(\defeq,\Leftrightarrow,\entails,\derives\)
  and their declared subscripted forms carry the precedence descriptor
  \texttt{meta} instead of a rank.  In their editorial roles they are not
  object-language operators: they relate expressions or statements at the
  metalinguistic level.  The registered source-local roles of
  \(\Leftrightarrow\) as a connective of an object language, the defined
  biconditional of the internal higher-order language in
  \cref{sec:p10-n-d-023-source-03} and \cref{annex:p10.mlm.8} and the
  defined formula equivalence of \cref{sec:p01-bpi-compactness-bridge} and
  \cref{annex:p01-bpi-propositional-compactness}, are infix operators of
  those languages; \cref{tab:p01-w02b-p01-013-roles} lists them with their
  regions.
\end{enumerate}
\end{hbconvention}

\begin{hbconvention}{def:p01-w02b-p01-014-quantifier-scope}{Precedence and quantifier scope in abbreviated formulas}
The convention applies only after the active language has been fixed and
its unique readability established; it says which formula of that
language an abbreviated display denotes.  The grammar of the active
language remains primary: where that grammar fixes a reading, the reading
stands, and where it requires parentheses, they are written.
\begin{enumerate}
\item \emph{Ranks.}  A pair of parentheses may be omitted only where the
  ranks of Table~\ref{tab:p01-precedence} and the association entries of
  \cref{def:p01-w02-p01-026-precedence} and
  Table~\ref{tab:p01-w02-p01-024-ranks} restore it in exactly one way.
  Smaller numerical ranks bind more tightly: of two operators that compete
  for one operand, the operator of smaller rank takes it.  Repetitions of
  one binary operator are grouped by its association entry, and different
  operators of one rank are separated by written parentheses
  (\cref{def:p01-w02-p01-024-fixity}, clauses 1--3).
\item \emph{Quantifier scope.}  For every variable \(x\), the prefixes
  \(\forall x\) and \(\exists x\) are prefix operators of rank 3, like
  \(\neg\), \(\Box\), \(\Diamond\), \(K_i\), \(E_G\) and \(C_G\).  The
  \emph{operand} of an occurrence of \(\forall x\) or \(\exists x\) runs
  from the symbol immediately after the prefix up to, and not including,
  the first symbol that is either a binary connective of rank greater
  than 3 lying inside no pair of parentheses opened after the prefix, or a
  closing parenthesis matching a parenthesis opened before the prefix; if
  there is no such symbol, the operand runs to the end of the display.
  The quantifier occurrence governs its operand and nothing else, so
  every binary connective of rank greater than 3 written after the
  operand lies outside its scope.  A binary operator of
  rank 3 written after the prefix outside such parentheses is not
  licensed (\cref{def:p01-w02-p01-024-fixity}, clause 3), and parentheses
  are written.
\item \emph{Binding.}  The occurrence of \(\forall x\) or \(\exists x\)
  binds exactly those occurrences of \(x\) in its operand that are free in
  the operand; an occurrence of \(x\) written after the operand is not
  bound by it.  Scope and binding are defined occurrence by occurrence in
  \cref{def:p01-w02-p01-024-free-bound}.
\item \emph{Stacked prefixes.}  By clause 2 a prefix written directly
  after a quantifier lies inside its operand, so
  \(\forall x\,\exists y\,R(x,y)\) is \(\forall x\,(\exists y\,R(x,y))\)
  and \(\forall x\,\neg P(x)\) is \(\forall x\,(\neg P(x))\); by clause 4
  of \cref{def:p01-w02-p01-024-fixity}, \(\neg\forall x\,P(x)\) is
  \(\neg(\forall x\,P(x))\).
\item \emph{Extending a scope.}  To place a binary compound inside the
  scope of a quantifier, the compound is enclosed in parentheses:
  \(\forall x\,(P(x)\to Q(x))\) is written when the implication lies inside
  the scope.  Parentheses are also written whenever a reader would
  otherwise need to reconstruct a local grammar or a scope declaration.
\item \emph{Dotted binders.}  A binder written with a separating dot, such
  as \(\forall a.\alpha\), \(\exists r.C\) or \(\forall r.C\), is read by
  clause 5 of \cref{def:p01-w02-p01-024-fixity}, which leaves the body
  after the dot to the source grammar.  Clause 2 of this convention reads
  the undotted prefixes \(\forall x\) and \(\exists x\).
\end{enumerate}
Hence \(\forall x\,P(x)\to Q(x)\) means \((\forall x\,P(x))\to Q(x)\), and
\(\exists x\,P(x)\land Q(x)\) means \((\exists x\,P(x))\land Q(x)\).
\end{hbconvention}

% GENERATED by tools/symbols2tex.py from handbook/registry/symbols.csv -- do not hand-edit.
\begingroup\small\setlength{\tabcolsep}{3pt}
\begin{longtable}{@{}p{0.07\textwidth}p{0.30\textwidth}p{0.30\textwidth}p{0.27\textwidth}@{}}
\caption{Registered operator ranks, generated from \texttt{registry/symbols.csv}}\label{tab:p01-precedence}\\
\toprule
\textbf{Rank} & \textbf{Glyph} & \textbf{Role} & \textbf{Scope} \\
\midrule
\endfirsthead
\caption[]{Registered operator ranks, generated from \texttt{registry/symbols.csv} (continued)}\\
\toprule
\textbf{Rank} & \textbf{Glyph} & \textbf{Role} & \textbf{Scope} \\
\midrule
\endhead
\midrule\multicolumn{4}{r}{\emph{continued}}\\\endfoot
\bottomrule\endlastfoot
1 & \(N!\) & factorial & P4 decision-hierarchy \\
1 & \(b^{\complement_{\mathbb B}}\) & Boolean-algebra complement & P6 cardinals-choice \\
1 & \(\kappa^+\) & successor-cardinal operation & P6 cardinals-choice to large-cardinals-absoluteness \\
2 & \(\leq\) & canonical ordering & Parts 1–4 from set-prerequisites to universal-reductions \\
2 & \(\leq\) & object-language arithmetic order abbreviation & P4 decision-hierarchy \\
2 & \(\leq\) & canonical ordering & Legacy pre-batch5 P4.4-P4.6 reader meaning \\
2 & \(\leq\) & object-language arithmetic order abbreviation & P5 arithmetic-coding \\
2 & \(\leq\) & canonical ordering & Parts 5–19 from first-incompleteness \\
2 & \(a\leq_{\N}b\) & metalinguistic natural-number order & Parts 4–5 from decision-hierarchy to arithmetic-coding \\
2 & \(<\) & canonical strict ordering & Parts 1–4 from set-prerequisites to universal-reductions \\
2 & \(<\) & object-language arithmetic strict-order abbreviation & P4 decision-hierarchy \\
2 & \(<\) & canonical strict ordering & Parts 4–19 from frontiers-computability \\
2 & \(a<_{\N}b\) & metalinguistic natural-number strict order & P4 decision-hierarchy \\
2 & \(>\) & strict ordering & global \\
2 & \(<_{\mathcal O}\) & Kleene notation order & Legacy pre-batch5 P4.5 reader meaning \\
2 & \(q_1\subseteq q_2\) & conjunctive-query containment & P9 legacy unit 7 logic-programming \\
2 & \(x\doteq y\) & Kamp coincidence abbreviation & P9 legacy unit 4 temporal-automata \\
2 & \(\leq_{\mathcal O}\) & reflexive notation predecessor relation & P4 delta-one hyperarithmetical construction \\
2 & \(\leq\) & canonical ordering & P4 represented-space numerical ordering \\
2 & \(p\lhd q\) & canonical well-order on a square of ordinals & \hyperref[sec:p01-w02-p01-011-finite-coding]{\ref*{sec:p01-w02-p01-011-finite-coding}} \\
2 & \(a<_{\mathrm{lex}}b\) & lexicographic order on triples of ordinals & \hyperref[annex:p01-w02-p01-011-canonical-order]{\ref*{annex:p01-w02-p01-011-canonical-order}} \\
2 & \(m<n\) & order of natural numbers symbol family member \(m<n\) & \hyperref[sec:p01-w02-p01-008-naturals]{\ref*{sec:p01-w02-p01-008-naturals}} \\
2 & \(\ m\leq n\) & order of natural numbers symbol family member \(\ m\leq n\) & \hyperref[sec:p01-w02-p01-008-naturals]{\ref*{sec:p01-w02-p01-008-naturals}} \\
2 & \(D\mathbin{\triangle}E\) & symmetric difference of two measurable sets & \hyperref[sec:p15-w02b-p15-005-hahn-stieltjes]{\ref*{sec:p15-w02b-p15-005-hahn-stieltjes}} \\
2.5 & \(!A\) & bang exponential and its phase interpretation & P8 paraconsistent-and-nonmonotonic \\
2.5 & \(?A\) & why-not exponential & P8 paraconsistent-and-nonmonotonic \\
2.5 & \(A^\perp\) & formula-level linear negation & P8 paraconsistent-and-nonmonotonic \\
2.5 & \(A^\circ\) & da Costa C1 consistency formula & P8 further-reading \\
3 & \(\neg\varphi\) & object-language negation & global \\
3 & \(\forall x\,\varphi\) & universal quantifier & global \\
3 & \(\exists x\,\varphi\) & existential quantifier & global \\
3 & \(\Box\) & necessity operator & Parts 5–8 from arithmetic-coding to conditionals-deontic-free \\
3 & \(\Box A\) & Anderson--Belnap relevance abbreviation & P8 paraconsistent-and-nonmonotonic \\
3 & \(\Diamond\) & possibility operator & Parts 5–8 from arithmetic-coding to conditionals-deontic-free \\
3 & \(\mathsf X\) & temporal next & Parts 8–9 from conditionals-deontic-free \\
3 & \(\mathsf F\) & temporal eventually & Parts 8–9 from conditionals-deontic-free \\
3 & \(\mathsf G\) & temporal always & Parts 8–9 from conditionals-deontic-free \\
3 & \(K_i\) & knowledge operator (agent i) & P8 conditionals-deontic-free \\
3 & \(E_G\) & mutual knowledge & P8 conditionals-deontic-free \\
3 & \(C_G\) & common knowledge & P8 conditionals-deontic-free \\
3 & \(\lambda x{:}\alpha.B\) & typed lambda abstraction & Legacy pre-batch4 P2.3-P7.6 reader meaning \\
3 & \(\mathsf{Eq}_{\alpha},\forall_{\alpha}\) & typed equality and universal-quantifier constants & P2 prop-metatheory \\
3 & \(E!(t)\) & free-logic existence predicate & Parts 2–8 from fol-syntax-semantics to many-valued-paraconsistent \\
3 & \(\iota x\,\varphi\) & definite-description term former & Parts 2–8 from fol-syntax-semantics to many-valued-paraconsistent \\
3 & \(\lambda x.M\) & untyped lambda abstraction & P7 lambda-calculi-and-type-theories \\
3 & \(\forall a.\alpha\) & System-F universal type former & P7 lambda-calculi-and-type-theories \\
3 & \(\Lambda a.M\) & System-F type abstraction & P7 lambda-calculi-and-type-theories \\
3 & \(\Pi_{x:A}B,\Sigma_{x:A}B\) & MLTT dependent product and dependent sum type formers & P7 lambda-calculi-and-type-theories \\
3 & \(\alpha\mathbin{\times_{\mathrm{ty}}}\beta,\alpha\mathbin{+_{\mathrm{ty}}}\beta\) & product and sum type formers in the categorical STLC extension & P7 normalization-curry-howard \\
3 & \(\varepsilon x\,\varphi(x)\) & Hilbert epsilon term former & Legacy pre-batch4 P7.7 reader meaning \\
3 & \(A\alpha,\mathcal F\alpha\) & Davies--Humberstone actuality and fixedly operators & P8 modal-temporal-epistemic \\
3 & \((\forall^E x)\Phi,(\exists^E x)\Phi\) & Fitting actualist quantifier abbreviations & P8 modal-temporal-epistemic \\
3 & \(\mathsf E,\mathsf A\) & CTL and CTL-star path quantifiers & Parts 8–9 from conditionals-deontic-free \\
3 & \(\mathsf{EX},\mathsf{AX},\mathsf{EF},\mathsf{AF}\) & CTL next and eventuality operators & Parts 8–9 from conditionals-deontic-free \\
3 & \(Z,Y,\mu Z.A,\nu Z.A\) & modal fixed-point variables and binders & Parts 8–9 from conditionals-deontic-free \\
3 & \(a,b,\langle a\rangle A,[a]A\) & Kozen program constants and indexed modalities & P8 conditionals-deontic-free \\
3 & \(B_iA\) & agent-indexed belief operator & P8 conditionals-deontic-free \\
3 & \(D_GA\) & distributed-knowledge operator & P8 conditionals-deontic-free \\
3 & \(U\,A\,B,S\,A\,B\) & Kamp strict tense connectives & P8 conditionals-deontic-free \\
3 & \(\mathsf P,\mathsf H\) & unary past tense operators & P8 conditionals-deontic-free \\
3 & \([!\chi]A,\langle!\chi\rangle A\) & public-announcement modalities & P8 conditionals-deontic-free \\
3 & \(@_iA\) & hybrid satisfaction operator & P8 conditionals-deontic-free \\
3 & \(\mathsf{SDL},O A,\mathsf P_OA\) & standard deontic identifier obligation and permission & P8 many-valued-paraconsistent \\
3 & \(LA,\operatorname{Th}\) & autoepistemic belief and classical closure & P8 further-reading \\
3 & \(\exists r.C,\forall r.C\) & ALC role restrictions & P9 legacy unit 3 description-logics \\
3 & \(\mathsf Y A\) & temporal previous operator & P9 legacy unit 4 temporal-automata \\
3 & \(\mathsf P_{\bowtie q}[\psi],\mathsf P^{\min}_{\bowtie q}[\psi],\mathsf P^{\max}_{\bowtie q}[\psi]\) & probability-threshold modalities & P9 legacy unit 5 model-checking \\
3 & \(\mathsf S_{\bowtie q}[\Phi]\) & CSL long-run threshold modality & P9 legacy unit 5 model-checking \\
3 & \([\alpha]A\) & PDL universal program modality & P9 legacy unit 6 program-logics \\
3 & \(\mathsf{not}\ A\) & logic-programming default negation & P9 legacy unit 7 logic-programming \\
3 & \(?\varphi\) & question-mark operator & P8 inquisitive-propositional \\
3 & \(\langle m\rangle,[m]\) & indexed diamond and box & \hyperref[sec:p08-src-025-501]{\ref*{sec:p08-src-025-501}} \\
3 & \(E,A\) & universal diamond and box & \hyperref[sec:p08-bvb501-universal]{\ref*{sec:p08-bvb501-universal}} \\
3 & \(D\) & difference diamond & \hyperref[sec:p08-bvb501-universal]{\ref*{sec:p08-bvb501-universal}} \\
3 & \(@_i\) & hybrid satisfaction operator & \hyperref[sec:p08-bvb501-hybrid]{\ref*{sec:p08-bvb501-hybrid}} \\
3 & \(\mu p.,\nu p.\) & least and greatest propositional fixed-point binders & \hyperref[sec:p08-bvb501-mu]{\ref*{sec:p08-bvb501-mu}} \\
3 & \(\langle\ ]\varphi\) & source alternative modality with a neighborhood contained in the truth set & P8 Pacuit neighborhood \\
3 & \(\forall_{\cdot},\exists_{\cdot}\) & source dotted actual-domain quantifiers & P8 Scott Advice \\
3 & \(\Box,\boxed\rho,\boxed{\leftrightarrow}\) & source universal relational and real-line progressive modalities & P8 Scott Advice \\
3 & \((x),\exists x\) & source local-domain universal binder and its Boolean-dual existential & P8 Considerations supporting comparison \\
3 & \(K_i\) & knowledge operator & p08 rationality-2-6 \\
3 & \(C\) & common knowledge operator & p08 epistemic update \\
3 & \([!\phi]\) & truthful public announcement box & p08 rationality-2-6 \\
3 & \(B\) & belief operator & p08 soft information \\
3 & \(B^\chi\) & conditional belief & p08 soft information \\
3.5 & \(A\mathbin{\mathsf U}B\) & temporal until & Parts 8–9 from conditionals-deontic-free \\
3.5 & \(A\mathbin{\mathsf S}B\) & temporal since operator & P9 legacy unit 4 temporal-automata \\
3.5 & \(\mathsf U^{\leq k},\mathsf X^{I},\mathsf U^{I}\) & bounded-step and real-time temporal operators & P9 legacy unit 5 model-checking \\
3.5 & \(\mathsf U_s\) & strict future until operator & P8 linear temporal completeness \\
4 & \(\varphi\land\psi\) & object-language conjunction & global \\
4 & \(\otimes\) & linear tensor and its homomorphic semantic interpretation & P8 paraconsistent-and-nonmonotonic \\
4 & \(\bigwedge_{i\in I}\varphi_i\) & countable infinitary conjunction & P2 fol-syntax-semantics \\
4 & \(b\wedge_{\mathbb B}c\) & Boolean-algebra meet & P6 cardinals-choice \\
4 & \(\bigwedge\Gamma\) & finite context conjunction & P7 lk-lj \\
4 & \(x\mathbin{\ast_{\text{\L}}}y,x\mathbin{\ast_{\Pi}}y\) & Lukasiewicz and Product t-norms & P8 relevance-linear \\
4 & \(x\ast y\) & t-norm operation & P8 relevance-linear \\
4 & \(A\bindnasrepma B\) & multiplicative disjunction (par) & P8 paraconsistent-and-nonmonotonic \\
4 & \(A\mathbin{\&}B\) & additive conjunction (with) & P8 paraconsistent-and-nonmonotonic \\
4 & \(A\mathbin{\cdot}B\) & Lambek product & P8 paraconsistent-and-nonmonotonic \\
4 & \(A\mathbin{\&}B\) & strong t-norm conjunction & P8 relevance-linear \\
4 & \(P*Q\) & separating conjunction & P9 legacy unit 6 program-logics \\
5 & \(\varphi\lor\psi\) & object-language disjunction & global \\
5 & \(\bigvee_{i\in I}\varphi_i\) & countable infinitary disjunction & P2 fol-syntax-semantics \\
5 & \(b\vee_{\mathbb B}c\) & Boolean-algebra join & P6 cardinals-choice \\
5 & \(\bigvee\Delta\) & finite context disjunction & P7 lk-lj \\
5 & \(A\oplus B\) & additive disjunction (plus) & P8 paraconsistent-and-nonmonotonic \\
5 & \(\mathbin{\vee\!\vee}\) & inquisitive disjunction & P8 inquisitive-propositional \\
5 & \(\lor\) & defined classical disjunction & P8 inquisitive-propositional \\
5 & \(\otimes\) & tensor disjunction on information states & P8 inquisitive tensor extension \\
6 & \(\varphi\to\psi\) & object-language implication & global \\
6 & \(\multimap\) & linear implication and its residual interpretation & P8 paraconsistent-and-nonmonotonic \\
6 & \(\alpha\to_{\mathrm{ty}}\beta\) & object-language function-type former & Legacy pre-batch4 P2.3-P7.6 reader meaning \\
6 & \(E(x)\supset\Phi\) & Fitting object-language conditional & P8 modal-temporal-epistemic \\
6 & \(x\Rightarrow y\) & t-norm residuum and Lukasiewicz implication & P8 relevance-linear \\
6 & \(A\backslash B,B/A\) & Lambek residuals & P8 paraconsistent-and-nonmonotonic \\
6 & \(O(A\mid B)\) & dyadic obligation & P8 many-valued-paraconsistent \\
6 & \(A>B\) & conditional connective & P8 many-valued-paraconsistent \\
6 & \(A\Box\!\!\to B\) & Lewis variably strict conditional & P8 many-valued-paraconsistent \\
6 & \(\alpha\to_{\mathrm{ty}}\beta\) & object-language function-type former & P12 research-workflow \\
6 & \(C\sqsubseteq D\) & description-logic concept inclusion & P9 legacy unit 3 description-logics \\
6 & \(>\) & conditional connective & \hyperref[sec:p08-bvb501-conditional]{\ref*{sec:p08-bvb501-conditional}} \\
7 & \(\varphi\leftrightarrow\psi\) & object-language biconditional & global \\
60 & \(\neg\) & denial of a stipulated selection condition & P19 sec:p19-r112-selection \\
\end{longtable}
\endgroup


Under \cref{def:p01-w02-p01-026-precedence},
\cref{def:p01-w02b-p01-014-quantifier-scope} and
\cref{def:p01-w02-p01-024-fixity} the following displays have these
readings:
\[
\begin{aligned}
 \neg P\land Q\to R
   &\quad\text{means}\quad ((\neg P)\land Q)\to R,\\
 P\to Q\to R
   &\quad\text{means}\quad P\to(Q\to R),\\
 \forall x\,P(x)\to Q(x)
   &\quad\text{means}\quad (\forall x\,P(x))\to Q(x),\\
 !A\otimes B\multimap C
   &\quad\text{means}\quad ((!A)\otimes B)\multimap C,\\
 \mathsf X A\mathbin{\mathsf U}B\lor C
   &\quad\text{means}\quad (\mathsf X A\mathbin{\mathsf U}B)\lor C.
\end{aligned}
\]
The third line is clause 2 of
\cref{def:p01-w02b-p01-014-quantifier-scope}; the other reading is written
\(\forall x\,(P(x)\to Q(x))\) by clause 5, and
\cref{def:p01-w02b-p01-014-two-readings} shows that the two readings can
differ in truth value.  Free and bound
occurrences, the scope of a quantifier occurrence and the binding relation
are defined occurrence by occurrence in \cref{def:p01-w02-p01-024-free-bound},
and \cref{thm:p01-w02-p01-024-fv-occurrence} identifies \(\FV(\varphi)\)
with the variables that have a free occurrence.

\begin{hbexample}{def:p01-w02b-p01-014-two-readings}{Two readings of a quantified implication}
Let \(P\) and \(Q\) be unary relation symbols.  In the display
\(\forall x\,P(x)\to Q(x)\) the first binary connective after
\(\forall x\) that lies inside no pair of parentheses opened after the
prefix is \(\to\), of rank 6, and no parenthesis is opened before the
prefix, so by clause 2 of \cref{def:p01-w02b-p01-014-quantifier-scope} the
operand of \(\forall x\) is \(P(x)\) and the display denotes
\[
  \varphi_1=(\forall x\,P(x))\to Q(x),
  \qquad\text{not}\qquad
  \varphi_2=\forall x\,(P(x)\to Q(x)).
\]
The occurrence of \(x\) in \(Q(x)\) lies outside the operand, so it is not
in the scope of the quantifier occurrence
(\cref{def:p01-w02-p01-024-free-bound}); hence \(\FV(\varphi_1)=\{x\}\) and
\(\varphi_1\) is not a sentence, while \(\FV(\varphi_2)=\varnothing\).

The two formulas differ in truth value.  Let \(\mathcal M\) have
\(\dom\mathcal M=\{0,1\}\), \(P^{\mathcal M}=\{0\}\) and
\(Q^{\mathcal M}=\varnothing\), and let \(s\) be any assignment, with
satisfaction as in \cref{sec:fol-semantics-02}.  Since
\(1\notin P^{\mathcal M}\), \(\mathcal M,s[x\mapsto1]\not\entails P(x)\), so
\(\mathcal M,s\not\entails\forall x\,P(x)\) and therefore
\(\mathcal M,s\entails\varphi_1\).  Since \(0\in P^{\mathcal M}\) and
\(0\notin Q^{\mathcal M}\), \(\mathcal M,s[x\mapsto0]\entails P(x)\) and
\(\mathcal M,s[x\mapsto0]\not\entails Q(x)\), so
\(\mathcal M,s[x\mapsto0]\not\entails P(x)\to Q(x)\) and
\(\mathcal M,s\not\entails\varphi_2\).  Reading the display as
\(\varphi_2\) therefore gives the wrong truth value under every assignment
in \(\mathcal M\).
\end{hbexample}

\section{The sixteen binary Boolean truth functions}
\label{sec:p01-sixteen-functions}
\index{truth function!binary}\index{Sheffer function}\index{Peirce arrow}

For inputs \((p,q)=(0,0),(0,1),(1,0),(1,1)\), a binary Boolean truth
function is determined by its four output bits.  Table~\ref{tab:p01-sixteen}
lists all \(2^{2^2}=16\) functions.  The formulas in the fourth column are
definitions using the handbook's registered connectives; historical glyphs
mentioned in the last column are literature variants, not additional active
handbook operators.

\begingroup\scriptsize\setlength{\tabcolsep}{3pt}
\begin{longtable}{@{}c p{0.20\textwidth} cccc p{0.23\textwidth}p{0.22\textwidth}@{}}
\caption{All binary Boolean truth functions.}
\label{tab:p01-sixteen}\\
\toprule
Code & Standard name & 00 & 01 & 10 & 11 & Defining formula & Notes \\
\midrule
\endfirsthead
\caption[]{All binary Boolean truth functions (continued)}\\
\toprule
Code & Standard name & 00 & 01 & 10 & 11 & Defining formula & Notes \\
\midrule
\endhead
\midrule\multicolumn{8}{r}{\emph{continued}}\\\endfoot
\endlastfoot
0000 & contradiction / constant false & 0&0&0&0 & \(\bot\) & independent of both inputs \\
0001 & conjunction & 0&0&0&1 & \(p\land q\) & --- \\
0010 & material nonimplication & 0&0&1&0 & \(p\land\neg q\) & negation of \(p\to q\) \\
0011 & first projection & 0&0&1&1 & \(p\) & independent of \(q\) \\
0100 & converse nonimplication & 0&1&0&0 & \(\neg p\land q\) & negation of \(q\to p\) \\
0101 & second projection & 0&1&0&1 & \(q\) & independent of \(p\) \\
0110 & exclusive disjunction & 0&1&1&0 & \(\neg(p\leftrightarrow q)\) & often XOR or \(\oplus\) \\
0111 & disjunction & 0&1&1&1 & \(p\lor q\) & inclusive disjunction \\
1000 & joint denial / NOR & 1&0&0&0 & \(\neg(p\lor q)\) & Peirce arrow; often \(p\downarrow q\) \\
1001 & biconditional / equivalence & 1&0&0&1 & \(p\leftrightarrow q\) & Boolean equality \\
1010 & negation of second input & 1&0&1&0 & \(\neg q\) & independent of \(p\) \\
1011 & converse implication & 1&0&1&1 & \(q\to p\) & false only at \((0,1)\) \\
1100 & negation of first input & 1&1&0&0 & \(\neg p\) & independent of \(q\) \\
1101 & material implication & 1&1&0&1 & \(p\to q\) & false only at \((1,0)\) \\
1110 & alternative denial / NAND & 1&1&1&0 & \(\neg(p\land q)\) & Sheffer function; often \(p\uparrow q\) or \(p\mid q\) \\
1111 & tautology / constant true & 1&1&1&1 & \(\top\) & independent of both inputs \\
\bottomrule
\end{longtable}
\endgroup

The Sheffer stroke (NAND) and Peirce arrow (NOR) are each functionally
complete for positive arities.  Neither represents a nullary truth
function: a formula built from NAND or NOR alone always contains a
variable, so the two constants need separately declared constants such as
\(\bottom\) and \(\topv\)
(\cref{thm:p01-w02-p01-024-nand-nor,thm:p01-w02-p01-024-no-closed}).
\index{Sheffer stroke}\index{NAND}\index{NOR}
The vertical-bar glyph is unavailable as active handbook notation, so this
definability note uses eliminable metalinguistic names:
\(\operatorname{NAND}(p,q)\defeq\neg(p\land q)\) and
\(\operatorname{NOR}(p,q)\defeq\neg(p\lor q)\).  These names designate
exactly rows 1110 and 1000, are not object-language connectives, and occur
only here.  Thus
\[
\begin{aligned}
\neg p&\Leftrightarrow\operatorname{NAND}(p,p),\\
p\land q&\Leftrightarrow
 \operatorname{NAND}(\operatorname{NAND}(p,q),
                     \operatorname{NAND}(p,q));\\
\neg p&\Leftrightarrow\operatorname{NOR}(p,p),\\
p\lor q&\Leftrightarrow
 \operatorname{NOR}(\operatorname{NOR}(p,q),
                    \operatorname{NOR}(p,q)).
\end{aligned}
\]
The criterion for an arbitrary set of positive-arity connectives, with
Post's five maximal clones and the convention for constants, is
\cref{thm:p01-w02-p01-024-post}.

\section{Representative notation in the literature}
\label{sec:p01-notation-variants}
\index{notation!historical variants}

Table~\ref{tab:p01-notation-variants} records representative printed usage
from identified source traditions.  It is bibliographic rather than
normative: authors within a tradition and later editions vary, and the table
alone changes no live presentation.  When a chapter explicitly presents one
of those named sources, however, that source's own syntax governs locally.

\begingroup\scriptsize\setlength{\tabcolsep}{3pt}
\begin{longtable}{@{}p{0.16\textwidth}p{0.18\textwidth}p{0.18\textwidth}p{0.18\textwidth}p{0.18\textwidth}@{}}
\caption{Representative notation in four named primary sources.}
\label{tab:p01-notation-variants}\\
\toprule
Role (handbook key) & Whitehead--Russell, \emph{Principia Mathematica} I (1910) & {\L}ukasiewicz, \emph{Elementy} (1929) & Hilbert--Ackermann, \emph{Grundz\"uge}, 1st ed. (1928) & Gentzen, \emph{Untersuchungen I} (1935) \\
\midrule
\endfirsthead
\caption[]{Representative notation in four named primary sources (continued)}\\
\toprule
Role (handbook key) & Whitehead--Russell (1910) & {\L}ukasiewicz (1929) & Hilbert--Ackermann (1928) & Gentzen (1935) \\
\midrule
\endhead
\midrule\multicolumn{5}{r}{\emph{continued}}\\\endfoot
\endlastfoot
negation (\(\neg A\)) & \(\mathord{\sim}p\) & \(Np\) & \(\overline X\) & \(\neg A\) \\
conjunction (\(A\land B\)) & \(p\mathbin{\boldsymbol{\cdot}}q\) & \(Kpq\) & \(X\mathbin{\&}Y\) & \(A\mathbin{\&}B\) \\
inclusive disjunction (\(A\lor B\)) & \(p\lor q\) & \(Apq\) & \(X\lor Y\) & \(A\lor B\) \\
implication (\(A\to B\)) & \(p\supset q\) & \(Cpq\) & \(X\to Y\) & \(A\supset B\) \\
biconditional (\(A\leftrightarrow B\)) & \(p\equiv q\) & \(Epq\) & \(X\sim Y\) & \(A\supset\subset B\), defined abbreviation \\
universal quantification (\(\forall x\,A\)) & \((x).\,\phi x\) & \(\Pi p\,\varphi(p)\) & \((x)A(x)\) & \(\forall x\,A\) \\
existential quantification (\(\exists x\,A\)) & \((\exists x).\,\phi x\) & \(\Sigma p\,\varphi(p)\) & \((E x)A(x)\) & \(\exists x\,A\) \\
\bottomrule
\end{longtable}
\endgroup

The entries reproduce the displayed forms in the cited works and editions;
they compare glyphs and argument placement, not the types of bound variables,
primitive bases, or every notation used by the authors elsewhere.  In the
cited {\L}ukasiewicz passage, \(p\) is a propositional variable.  Gentzen's
\(A\supset\subset B\) is a defined abbreviation, not a primitive formula
former, and his sequent \(\to\) is an auxiliary sign rather than implication.
Primary-source locators are recorded in
\cite{WhiteheadRussell1910,Lukasiewicz1929,HilbertAckermann1928,Gentzen1934a}.

\section{Letter-family conventions}
\label{sec:p01-letter-conventions}
\index{Greek letters}\index{Fraktur letters}\index{blackboard bold}

Typeface alone never assigns standing.  Table~\ref{tab:p01-letter-families}
counts, for each letter family, the register rows whose glyph uses the
family's typeface command and the distinct scopes those rows record, and it
prints the first glyphs of the family in register order; it is generated
from \texttt{registry/symbols.csv}.  A glyph is counted under every family
whose command it contains, in a subscript or a superscript as well.  A
glyph written with a handbook shorthand such as \texttt{\textbackslash N}
in place of the typeface command is not counted, whereas the family
extracts of \cref{sec:p01-w02b-p01-012-selection} count those shorthands.
The generated master register above remains authoritative for each
individual symbol, role and scope.

% GENERATED by tools/symbols2tex.py from handbook/registry/symbols.csv -- do not hand-edit.
\begingroup\small\setlength{\tabcolsep}{3pt}
\begin{longtable}{@{}p{0.16\textwidth}p{0.08\textwidth}p{0.08\textwidth}p{0.58\textwidth}@{}}
\caption{Letter families in the registry, generated from \texttt{registry/symbols.csv}}\label{tab:p01-letter-families}\\
\toprule
\textbf{Family} & \textbf{Rows} & \textbf{Scopes} & \textbf{First registered glyphs} \\
\midrule
\endfirsthead
\caption[]{Letter families in the registry, generated from \texttt{registry/symbols.csv} (continued)}\\
\toprule
\textbf{Family} & \textbf{Rows} & \textbf{Scopes} & \textbf{First registered glyphs} \\
\midrule
\endhead
\midrule\multicolumn{4}{r}{\emph{continued}}\\\endfoot
\bottomrule\endlastfoot
Greek lowercase & 2694 & 1124 & \(\varphi,\psi,\chi\); \(\neg\varphi\); \(\varphi\land\psi\); \(\varphi\lor\psi\) \\
Greek uppercase & 682 & 396 & \(\Gamma,\Delta,\Sigma,\Pi\); \(\Sigma^0_n,\Pi^0_n\); \(\boldsymbol{\Sigma}^0_n,\boldsymbol{\Pi}^0_n,\boldsymbol{\Delta}^0_n\); \(\Gamma\models\varphi\) \\
Calligraphic & 1233 & 649 & \(\mathcal L\); \(\mathcal A_{\mathrm{in}},\mathcal A_{\mathrm{tape}}\); \(\mathcal M,\mathcal N\); \(\mathcal M,s\models\varphi\) \\
Blackboard bold & 167 & 107 & \(\mathbb N\); \(\mathbb{P}\); \(V^{\mathbb P}\); \(p\Vdash^{\mathrm{sem}}_{M,\mathbb P}\varphi\) \\
Fraktur & 463 & 263 & \(\mathfrak L\); \(\mathcal G,\mathfrak F\); \(\mathfrak c\); \(\mathfrak F\) \\
Sans serif & 566 & 192 & \(\mathsf{Proof}_T\); \(\mathsf{Prov}_T\); \(\mathsf{Prov}_{\PA}\); \(\mathsf{RProof}_T\) \\
Roman operator & 1132 & 560 & \(\mathcal A_{\mathrm{in}},\mathcal A_{\mathrm{tape}}\); \(\operatorname{dom}X\); \(\operatorname{FV}(\varphi),\operatorname{FV}(\Gamma),\operatorname{FV}(M),\operatorname{FV}(\alpha)\); \(\operatorname{Mod}(T)\) \\
\end{longtable}
\endgroup


The family-level defaults are these.  Once a language is fixed, the
lowercase Greek letters \(\varphi,\psi,\chi\) and \(\sigma,\tau\) are the
global formula and sentence metavariables; once a calculus is fixed, the
uppercase Greek letters \(\Gamma,\Delta,\Sigma,\Pi\) are the global formula
contexts, and the calculus declares whether a context is a set, a multiset
or a sequence.  The letters \(\alpha,\beta,\gamma\) are ordinal
metavariables or type metavariables only in the disjoint scopes their rows
record (\cref{def:p01-w02-p01-024-type-metavariables}), and
\(\theta,\rho,\eta\) are simultaneous term substitutions, not formula
metavariables, in the scope their row records; the other roles of these
letters are scoped as their own rows record.  Calligraphic \(\mathcal L\)
and \(\mathcal M,\mathcal N\) are the global signature and structure
metavariables; every other calligraphic letter, a machine alphabet for
instance, has the role of its own global or part-local row, and the
typeface by itself has no role.  Blackboard bold has global rows only for
the natural numbers (with \(0\)), the integers, the rationals and the
reals.  Roman and sans-serif names denote named operations and fixed
systems or operators, never generic variables.  Any other letter receives
a role only through its own registry row.  The family extracts of
\cref{sec:p01-w02b-p01-012-extracts} list the rows of the Greek, Fraktur
and blackboard families with their recorded roles and scopes.

No Fraktur or scoped blackboard letter has a global role: each registered
use is confined to the scope its own row records
(\cref{sec:p01-w02b-p01-012-selection}), and a further use requires a scoped
registry row before substantive use.  Typography alone cannot create a
mathematical category without a definition and a scoped registry entry.
